\subsubsection{时间变量的统一陈述（续）：跨尺度障碍、伴随闭包与可识别指数律}\label{subsubsec:xi-time-protocol-conclusions-part2}
本节续接 \S\ref{subsubsec:xi-time-protocol-conclusions}，在不引入新本体假设的前提下，给出一组可审计的形式化版本；其条目均可独立引用。

\begin{theorem}[尺度--时间双复形：全局统一时间的唯一障碍由可计算谱序列边界给出]\label{thm:xi-scale-time-bicomplex-obstruction}
设分辨率塔 \((X_m,p_{m+1,m})\) 诱导群胚 \(\mathcal G_m\)，并在每层给定严格可加余圈 \(\ell_m:\mathcal G_m\to\mathbb Z\)。
取双复形 \(C^{p,q}\)：横向微分为群胚上同调边界 \(\delta\)，纵向微分为尺度差分 \(\partial\)（由 \(p_{m+1,m}^\ast\) 拉回定义）。
记 \(\ell:=(\ell_m)_m\) 为由各层长度余圈组成的总上链，令总复形余圈条件为 \((\delta+\partial)\ell=0\)。
则存在自然谱序列
\[
E_r^{p,q}\Rightarrow H^{p+q}(\mathrm{Tot}(C^{\bullet,\bullet})),
\]
使“存在全局统一时间”当且仅当障碍类 \(\mathfrak o\in E_2^{2,0}\) 消失。
若 \(\mathfrak o\neq 0\)，最小时间中心扩张秩等于 \(\mathfrak o\) 在 \(H^2\) 中生成的阿贝尔子群秩。
因此统一时间问题归约为二阶边界消失判据；在有限呈示情形下，该判据可由有限数据计算并判定。
\end{theorem}

\begin{proof}
统一时间对应总余圈可被全局 \(1\)-上链校正为严格可加代表。
该可解性在双复形中由二阶边界消失刻画，即 \(E_2^{2,0}\) 障碍。
非零障碍必须外置中心坐标，最小秩由独立障碍生成元数给出。
\end{proof}

\begin{theorem}[时间--熵帕累托前沿：纤维内最优消歧策略必为 Gibbs 形式且温度被可组合性强制为全局常数]\label{thm:xi-time-entropy-pareto-gibbs-global-temperature}
设 \(\mathrm{Fold}:\Omega\to X\)，并给定纤维内消歧时间 \(\tau:\Omega\to\mathbb N\)。
对每个 \(x\in X\)，在条件分布 \(\pi(\cdot\mid x)\) 上同时最小化 \(\mathbb E[\tau\mid x]\) 与最大化 \(H(\pi(\cdot\mid x))\)。
则全部帕累托最优点唯一参数化为
\[
\pi_\lambda(\omega\mid x)
=
\frac{e^{-\lambda\tau(\omega)}}{\sum_{\omega'\in \mathrm{Fold}^{-1}(x)}e^{-\lambda\tau(\omega')}}\qquad(\lambda>0).
\]
若再要求该策略在态射复合下封闭（可组合最优性），则 \(\lambda\) 不能依赖于 \(x\)：存在全局常数 \(\lambda_\ast\) 使全部纤维同时取 \(\pi_{\lambda_\ast}\)。
因此时间统一在统计层面等价于温度参数的全局一致性。
\end{theorem}

\begin{proof}
固定 \(x\) 时双目标拉格朗日化给出唯一指数族解，即 Gibbs 形。
复合封闭要求跨纤维的对数势可加常数一致；若 \(\lambda\) 随 \(x\) 变化则破坏复合协变性。
故仅允许全局单一温度参数。
\end{proof}

\begin{theorem}[黄金底对数的刚性：Zeckendorf 扫描的时间标尺在正则性下唯一]\label{thm:xi-logphi-time-rigidity}
\leanverified{paper\_xi\_logphi\_time\_rigidity}
设规模函数 \(S(m)\)（例如 Zeckendorf/Fibonacci 规范语法下分辨率 \(m\) 的最大可表达规模）满足 Fibonacci 递推
\[
S(m+1)=S(m)+S(m-1),\qquad S(m)\sim c\varphi^m.
\]
设时间函数 \(t:\mathbb N\to\mathbb R_{\ge 0}\) 满足：
\begin{enumerate}
  \item 近似可加：\(|t(ab)-t(a)-t(b)|\le C\)；
  \item 局部正则：\(|t(n+1)-t(n)|\le L\)；
  \item 分辨率一致：\(t(S(m))=m+O(1)\)。
\end{enumerate}
则存在常数 \(B\) 使
\[
t(n)=\frac{\log n}{\log\varphi}+B+o(1)\qquad(n\to\infty).
\]
\end{theorem}

\begin{proof}
由 \(S(m)\sim c\varphi^m\) 与条件 (3) 得 \(t\) 在 Fibonacci 子序列上与 \(\log_\varphi\) 同步。
条件 (1)(2) 将该同步从稀疏子序列稳定外推到全体整数，误差仅保留有界漂移与次阶项。
\end{proof}

\begin{theorem}[时间--分辨率的 Galois 伴随：互反闭包与稳定类型不动点]\label{thm:xi-time-resolution-galois-adjunction}
\leanverified{paper\_xi\_time\_resolution\_galois\_adjunction}
在有序集 \((\mathcal R,\preceq)\) 上，令 \(T:\mathcal R\to\mathbb R_{\ge 0}\) 为最小实现时间，\(R:\mathbb R_{\ge 0}\to\mathcal R\) 为预算 \(t\) 内可达最细分辨率。
若 \(T\) 单调且次可加
\[
T(\rho\vee\rho')\le T(\rho)+T(\rho'),
\]
则成立伴随律
\[
\rho\preceq R(t)\ \Longleftrightarrow\ T(\rho)\le t.
\]
因此 \(R\circ T\) 与 \(T\circ R\) 分别是闭包算子；
\(\mathrm{Fix}(R\circ T)\) 构成完备格，每个不动点对应时间推进下的稳定类型。
\end{theorem}

\begin{proof}
由定义，\(R(t)\) 是约束 \(T(\rho)\le t\) 的上确界对象，故得到伴随充要关系。
伴随自然诱导双侧闭包与不动点格结构，稳定类型即闭包不再变化的分辨率类。
\end{proof}

\begin{theorem}[混合曲率零判据：加法与乘法在同一时间轴上的统一当且仅当交错方块缺陷为零]\label{thm:xi-mixed-curvature-zero-criterion}
\leanverified{paper\_xi\_mixed\_curvature\_zero\_criterion}
设两类可组合演化 \(\circ_+\)、\(\circ_\times\) 分别带长度余圈 \(\ell_+,\ell_\times\)。
对任意交错方块
\(
w\circ_+u\simeq u'\circ_\times w'
\)
定义混合曲率
\[
K(w,u):=\ell_+(w)+\ell_\times(u)-\ell_\times(u')-\ell_+(w').
\]
则：
\begin{enumerate}
  \item \(K\) 对方块粘合满足加性，构成 \(2\)-余圈；
  \item 存在单一时间函数 \(\widetilde\ell\) 同时对 \(\circ_+\)、\(\circ_\times\) 严格可加
  当且仅当 \(K\equiv 0\)。
\end{enumerate}
若 \(K\neq 0\)，最小统一修正是向中心寄存器加入由 \(\{K(w,u)\}\) 生成的阿贝尔群。
\end{theorem}

\begin{proof}
\(K\) 本质是交错方块两条路径长度差，粘合时差值可加。
若存在共同可加 \(\widetilde\ell\)，任一方块差必为零；反向零差给出路径无关长度，可定义统一 \(\widetilde\ell\)。
非零差只能外置中心坐标吸收。
\end{proof}

\begin{theorem}[临界切片的压力刻画：时间反演对称等价于零漂移与零 Lyapunov 的共同成立]\label{thm:xi-critical-slice-pressure-zero-drift-lyapunov}
\leanverified{paper\_xi\_critical\_slice\_pressure\_zero\_drift\_lyapunov}
设 \(\theta\mapsto\mathcal L_\theta\) 为转移算子族，时间压强
\(
P(\theta):=\log\rho(\mathcal L_\theta)
\)。
记漂移 \(v(\theta):=P'(\theta)\)（存在时）与主 Lyapunov 指数 \(\chi(\theta)\)。
若系统满足严格时间反演共轭 \(\theta\leftrightarrow-\theta\)，则 \(P(\theta)=P(-\theta)\) 且 \(v(0)=0\)。
在可审计微可逆条件下，
\[
\theta=0
\Longleftrightarrow
v(\theta)=0\ \text{且}\ \chi(\theta)=0,
\]
且 \(\theta=0\) 为唯一双零切片。
\end{theorem}

\begin{proof}
时间反演使压强偶对称，故中心点漂移为零。
微可逆下，非零 \(\theta\) 会引入净漂移或指数放大中的至少一项；双零仅在共轭固定切片同时成立。
\end{proof}

\begin{theorem}[时间 \(\zeta\)-数据的可识别性：加权行列式族决定协议的 Markov 同构类与时间余圈类]\label{thm:xi-time-zeta-identifiability-markov-cohomology}
设协议由有限标记图 \(\Gamma\) 与时间余圈 \(\widetilde\ell\) 给出。
对每个 \(s>0\) 定义加权邻接算子 \(A_s\)，并令
\[
\mathcal Z(s,z):=\det(I-zA_s).
\]
若两套协议 \((\Gamma,\widetilde\ell)\)、\((\Gamma',\widetilde\ell')\) 满足
\(
\mathcal Z\equiv \mathcal Z'
\)
（作为 \((s,z)\) 解析函数恒等），则存在图同构把 \(\Gamma\) 送到 \(\Gamma'\)，并使 \(\widetilde\ell,\widetilde\ell'\) 在该同构下同调等价（相差顶点函数余边界）。
\end{theorem}

\begin{proof}
全参数行列式族决定各阶加权闭轨迹计数，从而决定 Markov 结构同构类。
对时间余圈，仅边权在回路上的同调不变量可被行列式观测；恒等数据迫使两者仅差余边界。
\end{proof}

\begin{theorem}[有限时间可分辨阈值：不可区分微态比例按时间指数衰减，其指数恰为单位时间互信息率]\label{thm:xi-finite-time-distinguishability-mutual-information-rate}
设底层微态过程 \((\omega_t)\) 平稳遍历，观测为折叠输出 \(x_t=\mathrm{Fold}(\omega_t)\)，并令 \(\mathcal F_t:=\sigma(x_0,\dots,x_t)\)。
定义时间 \(t\) 下的不可区分关系：对两初值 \(\omega,\omega'\)，若它们诱导的 \(\mathcal F_t\)-条件分布一致，则记 \(\omega\sim_t\omega'\)。
记时间 \(t\) 下平均不可区分类大小的对数为 \(\Delta(t)\)。
在正熵与有限记忆可实现条件下，极限
\[
\lim_{t\to\infty}\frac1t\Delta(t)
=
H(\omega_0\mid x_0)-\lim_{t\to\infty}\frac1t I(\omega_0;x_0,\dots,x_t)
\]
存在；并且若 \(\omega,\omega'\) 独立取自稳态分布，则不可区分比例满足
\[
\mathbb P(\omega\sim_t\omega')\asymp e^{-t\mathcal I},
\]
其中 \(\mathcal I:=\lim_{t\to\infty}\frac1t I(\omega_0;x_0,\dots,x_t)\) 为单位时间互信息率。
故时间推进以指数速率 \(\mathcal I\) 消解纤维自由度。
\end{theorem}

\begin{proof}
由链式法则与平稳遍历性，信息增益率极限存在并等于单位时间互信息率。
不可区分类体积按条件熵速率收缩，故概率比呈指数律，指数即 \(\mathcal I\)。
\end{proof}

\begin{theorem}[时间边界的双重重建：时间滤过树的极限空间与协议同调类在双向上互定]\label{thm:xi-time-boundary-double-reconstruction}
\leanverified{paper\_xi\_time\_boundary\_double\_reconstruction}
沿定理 \ref{thm:xi-finite-time-distinguishability-mutual-information-rate} 的记号，令 \(\sim_t\) 为时间预算 \(t\) 下的不可区分关系，并定义分层商树
\[
\mathcal T_t:=\Omega/\!\sim_t,
\]
以及由粗到细的自然投影 \(\pi_{t_2,t_1}:\mathcal T_{t_2}\to \mathcal T_{t_1}\)（\(t_2>t_1\)）构成的逆系统。
定义\textbf{时间边界}为投影极限
\[
\partial_T:=\varprojlim_{t}\mathcal T_t.
\]
对任意 \(\lambda>0\)，对 \(\xi=(\xi_t)_t,\eta=(\eta_t)_t\in\partial_T\) 定义最大公共前缀深度
\[
\tau(\xi\wedge\eta):=\sup\{\,t:\ \xi_t=\eta_t\,\}\in[0,\infty]\qquad(\exp(-\lambda\cdot\infty):=0),
\]
并赋予超度量
\[
d_\lambda(\xi,\eta):=\exp\bigl(-\lambda\,\tau(\xi\wedge\eta)\bigr).
\]
则在“遍历 + 强连通 + 时间严格正”（例如存在全边严格正 gauge；命题 \ref{prop:xi-positive-time-cone-criterion}）条件下成立：
\begin{enumerate}
  \item \textbf{（边界几何的同调不变性）} \((\partial_T,d_\lambda)\) 的双 Lipschitz 共轭类不依赖具体实现细节；在 gauge 余边界扰动下，边界度量至多发生双 Lipschitz 等价变形。
  \item \textbf{（同调类的反向重建）} 给定任一实现产生的 \((\partial_T,d_\lambda)\) 及其自然柱集测度类（由遍历不变测度诱导的分层柱集质量），可在上同调意义下重建时间余圈类（差一个 gauge 余边界；等价于闭环周期和不变，见命题 \ref{prop:xi-cycle-sums-determine-time-cohomology}）。
\end{enumerate}
因此，时间既生成边界几何，又可由边界几何反演；“存在/可见性”在该意义下等价于 \(\partial_T\) 上的点。
\end{theorem}
\begin{proof}
由定义，\(\partial_T\) 是随 \(t\uparrow\infty\) 逐级细化的不可区分分割之相容极限；\(\tau(\xi\wedge\eta)\) 的定义给出树边界的标准前缀超度量，从而 \(d_\lambda\) 为超度量。
\par
对 (1)，gauge 余边界扰动只改变“在每一层把哪些对象视为等代价/等预算”的局部阈值，而不改变闭环周期和与其诱导的同调类；这等价于在树边界上作一类保持柱集层级的重标定，从而诱导双 Lipschitz 等价的边界几何。
\par
对 (2)，柱集测度在层级 \(t\) 上对 \(\mathcal T_t\) 的赋值给出沿边界射线的尺度信息；其对闭环的相容性（在遍历/强连通下可由有限循环基核验）迫使该尺度信息在同调意义下由唯一时间余圈类承载。
将该“周期和一致性”与有限循环基判别（命题 \ref{prop:xi-cycle-sums-determine-time-cohomology}）合并，即得到时间同调类的反演重建（差一个余边界）。证毕。
\end{proof}

\subsubsection{时间边界重建后的扩展结论簇}\label{subsubsec:xi-time-protocol-extended-cluster-a}

\paragraph{尺度联络、Holonomy 与时间特征类}

\begin{theorem}[尺度--时间联络的 Holonomy 判别：全局统一时间当且仅当离散曲率群消失]\label{thm:xi-scale-time-connection-holonomy-criterion}
\leanverified{paper\_xi\_scale\_time\_connection\_holonomy\_criterion}
设分辨率塔 \((X_m,p_{m+1,m})\) 在每一层诱导群胚 \(\mathcal G_m\)，并给定严格可加时间余圈 \(\ell_m:\mathcal G_m\to\mathbb Z\)。
定义尺度提升缺陷（离散联络的 \(1\)-形式）
\[
a_m(g):=\ell_{m+1}(g)-\ell_m\!\bigl(p_{m+1,m}(g)\bigr).
\]
对任意可交换的“尺度--时间小方块” \(g\Rightarrow g'\)（即在证书意义下可合流的平行态射对）定义其 Holonomy
\[
\mathrm{Hol}_m(g\Rightarrow g'):=a_m(g)-a_m(g')\in\mathbb Z.
\]
令 \(\mathfrak H\le \mathbb Z\) 为由全部 \(\mathrm{Hol}_m(g\Rightarrow g')\) 生成的子群。
则存在全局统一时间（等价地，存在一族 gauge 校正使各层时间余圈在投影下严格相容）当且仅当 \(\mathfrak H=\{0\}\)。
并且当 \(\mathfrak H\neq 0\) 时，使统一时间成立所需的最小中心扩张的中心秩等于 \(\mathrm{rank}(\mathfrak H)\)。
\end{theorem}

\begin{proof}
在定理 \ref{thm:xi-scale-time-bicomplex-obstruction} 的双复形中，\(a_m=\partial\ell\) 为纵向差分，而 \(\mathrm{Hol}_m\) 正是 \(\delta a_m=\delta\partial\ell\) 在基本交换方块上的取值。
因此 \(\mathfrak H\) 恰为障碍类 \(\mathfrak o\in E_2^{2,0}\) 的可观测生成子群；\(\mathfrak H=\{0\}\) 等价于 \(\mathfrak o=0\)，从而等价于全局统一时间的存在。
最小中心秩结论同理由 \(\mathfrak o\) 所生成阿贝尔子群的秩给出（定理 \ref{thm:xi-scale-time-bicomplex-obstruction}）。
\end{proof}

\begin{theorem}[尺度自洽时间的 Perron--Frobenius 刚性：跨尺度统一时间等价于正锥中的唯一本征余圈]\label{thm:xi-scale-coherent-time-pf-rigidity}
\leanverified{paper\_xi\_scale\_coherent\_time\_pf\_rigidity}
设分辨率提升函子 \(R:X_{m+1}\to X_m\)（例如由折叠/重整化诱导）在各层诱导群胚拉回
\(
R^\ast:=p_{m+1,m}^\ast:H^1(\mathcal G_m;\mathbb R)\to H^1(\mathcal G_{m+1};\mathbb R)
\)。
在每层上以“对所有有向闭环配对为正”定义时间正锥的内部 \( \mathcal C_{m,+}^\circ\subset H^1(\mathcal G_m;\mathbb R)\)（有限生成情形见命题 \ref{prop:xi-positive-time-cone-criterion}）。
若 \(R^\ast\) 在正锥上保持内点并在遍历分量上原始（存在某个幂次使其把任意非零正类推入内点），则存在唯一射影类 \([\ell_\infty]\in \mathbb P(\mathcal C_{+}^\circ)\) 与唯一常数 \(\Lambda>1\)（在固定时间单位归一化下），使得对所有足够大层级 \(m\) 有
\[
R^\ast[\ell_m]=\Lambda\,[\ell_{m+1}],
\]
并且任何其他跨尺度自洽的严格正时间余圈族（允许 gauge 余边界扰动）必与 \([\ell_\infty]\) 共线。
其含义是：跨尺度的“同一时间轴”并非可任意选择，而由重整化算子在正锥中的主本征方向刚性确定；\(\log\Lambda\) 同时给出分辨率增长常数与时间单位的唯一归一化标尺。
\end{theorem}
\begin{proof}
在正锥的 Hilbert 射影度量上，原始的锥保持算子为严格收缩，从而存在唯一吸引的不动射线（Perron--Frobenius/Krein--Rutman 型结论）。
将该射线记为 \([\ell_\infty]\)。
对任意严格正的自洽族 \([\ell_m]\)，反复拉回并在射影空间比较即可知其射影极限必须落在同一吸引射线上；对应的尺度因子在固定归一化下收敛到唯一常数 \(\Lambda\)。
gauge 余边界扰动不改变闭环配对与正锥射影类，故不影响结论。证毕。
\end{proof}

\begin{theorem}[时间特征类的积分性：复合异常与混合曲率统一为一个整数 \(H^2\) 类]\label{thm:xi-time-characteristic-class-integrality}
\leanverified{paper\_xi\_time\_characteristic\_class\_integrality}
设 \(\mathcal B\) 为承载两类可组合结构（例如加法层 \(\circ_+\) 与乘法层 \(\circ_\times\)）的 \(2\)-范畴，并假定每一层均赋予整数值长度余圈，使得：
复合异常由 \(2\)-余圈 \(\alpha\) 表示（如命题 \ref{prop:xi-time-length-cocycle}），交错方块缺陷由混合曲率 \(2\)-余圈 \(K\) 表示（如定理 \ref{thm:xi-mixed-curvature-zero-criterion}）。
则存在自然的整数 \(2\)-余圈 \(\Omega\)（由 \(\alpha\) 与 \(K\) 在相应双复形中的总微分合成得到），满足：
\[
[\Omega]\in H^2(\mathcal B;\mathbb Z)
\]
在 gauge 变换下不变；并且对任意闭 \(2\)-链 \(\Sigma\) 的配对
\[
\langle[\Omega],\Sigma\rangle\in\mathbb Z
\]
刻画 \(\Sigma\) 上的净“时间缺陷”。
因此，时间统一与加乘统一的全部障碍可压缩为同一个可积的整数特征类。
\end{theorem}

\begin{proof}
\(\alpha\) 与 \(K\) 均为整数值 \(2\)-余圈，其 gauge 变换只引入整数余边界项，故其在总复形中诱导的总余圈 \(\Omega\) 的同调类 \([\Omega]\) 为不变量。
对闭 \(2\)-链的配对由上同调基本性质给出并取整数值；当且仅当 \([\Omega]=0\) 时可同时消除复合异常与混合曲率，从而得到统一时间与统一复合结构。
\end{proof}

\begin{theorem}[时间索引定理：时间谱三元组的 \(K_1\) 配对等于时间 Holonomy]\label{thm:xi-time-index-theorem-k1-holonomy}
\leanverified{paper\_xi\_time\_index\_theorem\_k1\_holonomy}
在定理 \ref{thm:xi-time-spectral-triple-distance-volume} 的设定下，令 \((\mathcal A,\mathcal H,D)\) 为由时间余圈 \(\widetilde\ell:\mathcal G\to\mathbb Z\) 构造的谱三元组（\(D\delta_g=\widetilde\ell(g)\delta_g\)），并令 \(P:=\mathbf 1_{[0,\infty)}(D)\) 为 \(D\) 的非负谱投影。
对任意酉元 \(u\in\mathcal A\)（代表 \(K_1(\mathcal A)\) 中的闭环类），Fredholm 指数满足
\[
\mathrm{Index}(PuP)=\langle [u],[\tau_{\widetilde\ell}]\rangle,
\]
其中 \([\tau_{\widetilde\ell}]\) 为由时间余圈诱导的循环 \(1\)-余圈的同调类。
并且在群胚闭环表述下，该配对退化为沿闭环的净时间 Holonomy；等价地，它与定理 \ref{thm:xi-time-characteristic-class-integrality} 的时间特征类 \([\Omega]\) 在对应 \(2\)-链上的配对一致。
\end{theorem}

\begin{proof}
指数等式是奇型谱三元组的标准 Connes 配对公式：\(\mathrm{Index}(PuP)\) 由 \(K_1(\mathcal A)\) 与循环上同调中的 Chern 字符配对给出。
在本构造中，\(D\) 在群胚基上对角化，故对应循环余圈可取为时间余圈的群胚上同调代表；对闭环的求值即为时间余圈沿闭环的净增量，从而等价于 Holonomy。
与 \([\Omega]\) 的一致性来自“闭环边界为零”时，净增量可改写为二维特征类在填充 \(2\)-链上的积分配对。
\end{proof}

\paragraph{时间几何与可审计复杂度}

\begin{theorem}[审计等式问题的时间同伦不变量：证书 Dehn 函数由时间几何唯一决定]\label{thm:xi-audit-dehn-function-time-geometry}
设协议由有限生成收敛重写系统给出，\(\delta(n)\) 为 Dehn 型函数：对任意长度不超过 \(n\) 的两词 \(w\sim w'\)，令 \(\delta(n)\) 取其最短审计证书的最坏时间长度上确界。
若基本步时间严格可加且对 \(2\)-证书同伦不变，则 \(\delta(n)\) 与时间度量 Cayley 图的等周函数在拟等距意义下等价。
特别地，
\[
\delta(n)=O(n)
\quad\Longleftrightarrow\quad
\text{时间度量 Cayley 图满足线性等周不等式（从而为 Gromov 双曲）},
\]
并因此证书验证具有线性审计界。
\end{theorem}

\begin{proof}
\(\delta(n)\) 的定义正是把“词等价的最小证书面积/时间”外部化为等周函数；时间严格可加与同伦不变性保证该函数只依赖于时间度量的拟等距类型。
线性等周不等式与 Gromov 双曲性的等价性是几何群论中的标准结论；在本设定下，\(\delta(n)=O(n)\) 即给出线性等周界，从而推出双曲性与线性审计界。
\end{proof}

\paragraph{增长阈值、最大可见容量与消歧时间分解}

\begin{theorem}[时间相变阈值由增长指数决定，且中心缺陷秩控制平衡态族的维数]\label{thm:xi-time-phase-transition-growth-and-defect-rank}
在一遍历传递分量上，取由基本步生成的正向可组合半群（或正锥）
\[
\mathcal G_+:=\{g:\widetilde{\ell}(g)\in\mathbb N\},
\]
并定义时间球计数指数
\[
\kappa:=\lim_{t\to\infty}\frac{1}{t}\log \#\{g\in\mathcal G_+:\widetilde{\ell}(g)\le t\},
\]
并定义时间配分函数（按源端固定）
\[
Z(\lambda):=\sum_{g\in\mathcal G_+} e^{-\lambda \widetilde{\ell}(g)}.
\]
则存在临界值 \(\lambda_c=\kappa\) 使得 \(Z(\lambda)<\infty\) 当且仅当 \(\lambda>\lambda_c\)。
并且在 \(\lambda>\lambda_c\) 区间内，与该权重相容的平衡态为唯一；在临界点 \(\lambda=\lambda_c\) 处，平衡态的极端分解维数由尺度--时间 Holonomy 缺陷群 \(\mathfrak H\)（定理 \ref{thm:xi-scale-time-connection-holonomy-criterion}）的秩控制：其极端射线族的仿射维数等于 \(\mathrm{rank}(\mathfrak H)\)。
\end{theorem}

\begin{proof}
由时间球计数的指数增长与比较判别法，\(\sum_{g\in\mathcal G_+} e^{-\lambda\widetilde\ell(g)}\) 的收敛阈值由指数 \(\kappa\) 决定，故 \(\lambda_c=\kappa\)。
唯一性与临界退化的结构性来源于“异常/缺陷可否被中心吸收”的分岔：当 \(\lambda>\lambda_c\) 时，Gibbs 权重给出唯一平衡态；当 \(\lambda=\lambda_c\) 时，非平凡 Holonomy 产生等价类的中心参数自由度，使平衡态形成以 \(\mathfrak H\) 的 Pontryagin 对偶为参数的仿射族，其维数由 \(\mathrm{rank}(\mathfrak H)\) 给出。
\end{proof}

\begin{theorem}[Zeckendorf 禁 \(11\) 的极大性原理：局部可逆归一化类中的最大可见容量由黄金移位唯一实现]\label{thm:xi-zeckendorf-golden-shift-max-capacity}
\leanverified{paper\_xi\_zeckendorf\_golden\_shift\_max\_capacity}
考虑所有二元移位约束 \(Y\subset\{0,1\}^{\mathbb Z}\)，满足：存在有限半径、可逆且合流的局部归一化，使每个有限词可在局部规则下归一至唯一正规形，并使由正规形定义的时间长度对复合严格可加。
则在此类约束中，允许的拓扑熵 \(h_{\mathrm{top}}(Y)\) 存在最大值，且唯一（在共轭意义下）由黄金移位实现：其邻接矩阵为
\[
\begin{pmatrix}
1 & 1\\[2pt]
1 & 0
\end{pmatrix},
\]
等价于唯一禁词为 \(11\)。
因此，Zeckendorf 语法可从“时间可加 + 局部可逆 + 最大可见容量”三条原则推出，而非先验假设。
\end{theorem}

\begin{proof}
局部可逆与合流给出正规形语言的唯一性；时间严格可加把正规形长度提升为可组合余圈。
在二元局部归一化口径下，合法语法的计数增长必须与 Fibonacci 递推兼容（见命题 \ref{prop:xi-local-binary-normalization-fibonacci-rigidity} 的权重刚性口径），从而其熵率上界由 Perron--Frobenius 根 \(\varphi\) 锁定，达到上界当且仅当语法退化为黄金均值约束（禁 \(11\)）。
共轭唯一性由二元不可约 SFT 在给定 PF 根下的标准分类给出。
\end{proof}

\begin{proposition}[最小消歧时间的精确分解：几何项为 \(\mathbb E\lbrack \log\mathrm{Ind}\rbrack\)，统计项为纤维内相对熵]\label{prop:xi-min-disambiguation-time-entropy-decomposition}
设折叠 \(\mathrm{Fold}:\Omega\to X\) 诱导纤维指数 \(\mathrm{Ind}(x):=|\mathrm{Fold}^{-1}(x)|\)，并给定条件分布 \(\mu(\cdot\mid x)\)。
记纤维内均匀分布为 \(u_x\)。
则条件熵分解恒等式为
\[
H_\mu(\omega\mid x)
=\mathbb E_\mu[\log \mathrm{Ind}(x)]
-\mathbb E_\mu\!\bigl[D_{\mathrm{KL}}(\mu(\cdot\mid x)\,\|\,u_x)\bigr].
\]
进一步，对任意字母表 \(|\mathcal A|\) 的顺序读出协议，其最小期望消歧时间满足
\[
\inf \mathbb E_\mu[T]
=
\frac{H_\mu(\omega\mid x)}{\log|\mathcal A|}+O(1)
=
\frac{\mathbb E_\mu[\log \mathrm{Ind}(x)]}{\log|\mathcal A|}
-\frac{\mathbb E_\mu[D_{\mathrm{KL}}(\mu(\cdot\mid x)\,\|\,u_x)]}{\log|\mathcal A|}
+O(1).
\]
\end{proposition}

\begin{proof}
对每个 \(x\) 有
\(
D_{\mathrm{KL}}(\mu(\cdot\mid x)\,\|\,u_x)=\log\mathrm{Ind}(x)-H(\mu(\cdot\mid x))
\)；
对 \(x\) 取期望即得第一式。
第二式由定理 \ref{thm:xi-optimal-recovery-time-equals-conditional-entropy} 的最优停止时间界直接换算得到。
\end{proof}
\leanverified{paper\_xi\_min\_disambiguation\_time\_entropy\_decomposition}

\paragraph{时间序结构与实现刚性}

\begin{proposition}[时间正序化与无扭性：严格正时间余圈强制左序与无有限周期]\label{prop:xi-time-positive-cocycle-left-order-torsionfree}
设在某传递分量的各定点稳定子群（同伦群）上，时间余圈 \(\widetilde{\ell}\) 满足
\(\widetilde{\ell}(gh)=\widetilde{\ell}(g)+\widetilde{\ell}(h)\)，
且对任何非平凡元 \(g\neq e\) 有 \(\widetilde{\ell}(g)\neq 0\)。
定义正锥
\[
P:=\{g:\widetilde{\ell}(g)>0\}.
\]
则 \(P\) 为子半群且与其逆不交，从而稳定子群可左序；并且不存在扭元素（若 \(g^n=e\) 则 \(g=e\)）。
\end{proposition}

\begin{proof}
由可加性得 \(P\cdot P\subseteq P\)，且 \(g\in P\Rightarrow g^{-1}\notin P\)（因 \(\widetilde{\ell}(g^{-1})=-\widetilde{\ell}(g)\)）。
据此可定义左不变全序：\(g\prec h\) 当且仅当 \(h g^{-1}\in P\)。
若 \(g^n=e\)，则 \(0=\widetilde{\ell}(g^n)=n\widetilde{\ell}(g)\)，由非退化性得 \(\widetilde{\ell}(g)=0\)，从而 \(g=e\)。
\end{proof}
\leanverified{paper\_xi\_time\_positive\_cocycle\_left\_order\_torsionfree}

\begin{theorem}[时间共轭刚性：同一时间余圈决定有限状态实现的有界共轭类]\label{thm:xi-time-conjugacy-rigidity-bounded-class}
设两套有限状态实现 \((\Gamma,\widetilde{\ell})\)、\((\Gamma',\widetilde{\ell}')\) 表示同一协议（同一可见群胚的两种覆盖），并满足时间余圈同调等价：存在顶点函数 \(b\) 使
\[
\widetilde{\ell}'=\widetilde{\ell}+\partial b.
\]
则存在滑动块共轭 \(\Phi\)（有限记忆、有限前瞻）将 \(\Gamma\) 的路径空间送至 \(\Gamma'\) 的路径空间，并满足时间严格保持
\[
\widetilde{\ell}'(\Phi(\gamma))=\widetilde{\ell}(\gamma)
\quad\text{对几乎处处路径 \(\gamma\)}.
\]
此外，共轭窗口大小可由尺度--时间 Holonomy 生成元的最大绝对值（视作“时间曲率范数”）给出一致上界。
\end{theorem}

\begin{proof}
余圈同调等价意味着两实现的时间权重仅差一项顶点势的余边界；对路径上的累积时间，该余边界沿路径望远镜相消，从而在适当的状态重标定/状态分裂的 recoding 下可实现严格时间保持的共轭。
窗口大小的有界性来自 \(b\) 的有限振幅控制，而 \(b\) 的必要振幅由 holonomy（曲率）生成元的最大绝对值下界控制。
\end{proof}

\begin{proposition}[时间配分函数的可识别性：仅由 \(\{Z_m(\lambda)\}\) 可重建每层纤维指数多重集]\label{prop:xi-time-partition-identifiability-ind-multiset}
对每一尺度 \(m\)，令纤维指数多重集为 \(\{\mathrm{Ind}_m(x):x\in X_m\}\subset\mathbb N\)，并定义时间势 \(T_m(x)=\log \mathrm{Ind}_m(x)\)。
考虑 Mellin--Laplace 型配分函数
\[
Z_m(\lambda):=\sum_{x\in X_m} e^{-\lambda T_m(x)}=\sum_{x\in X_m}\mathrm{Ind}_m(x)^{-\lambda}.
\]
若在固定尺度 \(m\) 上 \(\mathrm{Ind}_m(x)\) 仅取有限多个整数值，则 \(Z_m(\lambda)\) 在任意含聚点的 \(\lambda\) 集合上的取值唯一决定该有限离散测度，从而唯一决定每个指数值的重数。
进一步，当 \(m\to\infty\) 存在压强极限
\(
P(\lambda)=\lim_{m\to\infty}\frac{1}{m}\log Z_m(\lambda)
\)
并在一区间可微时，\(P\) 的 Legendre 对偶唯一给出纤维时间的多重分形谱（见定理 \ref{thm:xi-fiber-time-multifractal-legendre}）。
\end{proposition}

\begin{proof}
在有限谱假设下，\(Z_m(\lambda)\) 是有限个函数 \(\lambda\mapsto n^{-\lambda}\) 的线性组合；这些函数在实解析意义下线性无关，故任意含聚点集合上的相等蕴含系数（即重数）逐一相等。
第二部分由压强极限与 G\"artner--Ellis 条件给出的大偏差--Legendre 对偶机制直接推出。
\end{proof}
\leanverified{paper\_xi\_time\_partition\_identifiability\_ind\_multiset}

\paragraph{有限状态实现下的时间像与离散化}

\begin{proposition}[有限状态实现下的时间像：离散量子化与稠密化二分]\label{prop:xi-finite-state-time-image-discrete-or-dense}
\leanverified{paper\_xi\_finite\_state\_time\_image\_discrete\_or\_dense}
设可见群胚 \(\mathcal G\) 具有有限生成的有向实现（有限状态、有限标记边集 \(E\)），并给定可加时间余圈 \(\widetilde{\ell}:\mathcal G\to\mathbb R\) 满足 \(\widetilde{\ell}(gh)=\widetilde{\ell}(g)+\widetilde{\ell}(h)\)。
若生成边上的取值集合 \(\{\widetilde{\ell}(e):e\in E\}\) 有限，则
\[
\widetilde{\ell}(\mathcal G)\subseteq
H:=\Bigl\langle \widetilde{\ell}(e):e\in E\Bigr\rangle_{\mathbb Z}\ \le\ (\mathbb R,+).
\]
并且 \(H\) 作为 \(\mathbb Z\)-模有限生成，故满足二分律：
\[
H=h\mathbb Z\ \text{（离散）}\qquad\text{或}\qquad \overline{H}=\mathbb R\ \text{（稠密）}.
\]
在离散情形中，\(h\) 可取为 \(H\cap(0,\infty)\) 的最小正元（等价于 \(\inf(H\cap(0,\infty))\)），并有 \(\widetilde{\ell}(\mathcal G)\subset h\mathbb Z\)。
此外，\(\widetilde{\ell}(\mathcal G)\) 至多可数，因而任何真正连续的时间像（不可数）只能通过状态数趋于无穷的极限获得。
\end{proposition}

\begin{proof}
因 \(\mathcal G\) 由有限边集生成，任意 \(g\in\mathcal G\) 的 \(\widetilde{\ell}(g)\) 为生成边取值的整数线性组合，故 \(\widetilde{\ell}(\mathcal G)\subseteq H\)。
而任意 \(\mathbb Z\)-有限生成子群 \(H\le\mathbb R\) 要么含有最小正元（从而必为该最小正元生成的循环群 \(h\mathbb Z\)），要么在 \(\mathbb R\) 中稠密；这是加法子群的标准分类。
最后，有限生成群 \(H\cong \mathbb Z^k\) 可数，故 \(\widetilde{\ell}(\mathcal G)\) 可数。
\end{proof}

\paragraph{欧拉乘积、Holonomy 与 RH 的同调化归约}

\begin{theorem}[时间欧拉分解定理：混合曲率零强制配分函数的素元乘积分解]\label{thm:xi-time-euler-product-under-zero-mixed-curvature}
设 \((\mathcal M_\times,\circ_\times)\subset \mathrm{Mor}(\mathbf{PW})\) 为乘法复合幺半群，并假设存在时间余圈 \(\widetilde{\ell}\) 对 \(\circ_\times\) 严格可加。
令 \(\mathcal P\subset\mathcal M_\times\) 为 \(\circ_\times\)-不可约元（素元）集合。
若满足：
\begin{enumerate}
  \item \(\mathcal M_\times\) 为交换、左消去且为素元的自由交换幺半群（等价地，每个元具有唯一的素元分解，按重数计）；
  \item 加乘交错方块的混合曲率恒为零（定理 \ref{thm:xi-mixed-curvature-zero-criterion}），从而时间在加/乘两层间无中心补偿缺陷；
\end{enumerate}
则对任意处于绝对收敛域的 \(\lambda>0\)，乘法配分函数
\[
Z_\times(\lambda):=\sum_{w\in\mathcal M_\times}e^{-\lambda \widetilde{\ell}(w)}
\]
具有欧拉乘积分解
\[
Z_\times(\lambda)=\prod_{p\in\mathcal P}\bigl(1-e^{-\lambda \widetilde{\ell}(p)}\bigr)^{-1}.
\]
反之，若对某开区间 \(\lambda\in(\lambda_0,\infty)\) 上 \(Z_\times\) 具有上述欧拉乘积且 \(\mathcal P\) 生成 \(\mathcal M_\times\)，则混合曲率必为零（等价于加法更新与乘法更新在时间意义下可交换且无需额外中心补偿）。
\end{theorem}

\begin{proof}
在唯一素元分解与可加性下，
\(
\widetilde{\ell}\bigl(\prod_p p^{n_p}\bigr)=\sum_p n_p\widetilde{\ell}(p)
\)。
对 \(\lambda\) 处于绝对收敛域时，按素元分解改写求和并交换求和与乘积，得到
\[
Z_\times(\lambda)
=\prod_{p\in\mathcal P}\sum_{n\ge 0}e^{-\lambda n\widetilde{\ell}(p)}
=\prod_{p\in\mathcal P}(1-e^{-\lambda\widetilde{\ell}(p)})^{-1}.
\]
反向方向中，欧拉乘积在一段开区间上成立迫使各阶系数满足素元独立的乘法型生成结构，从而与“交错方块无缺陷”的路径无关性相容；若混合曲率非零，则加/乘交错会引入不可由边界吸收的中心补偿项，破坏该乘法型分解，故必须有 \(K\equiv 0\)。
\end{proof}

\begin{theorem}[时间射线的素分解刚性：正锥的极端射线与协议的不可约乘积分解一一对应]\label{thm:xi-time-ray-prime-factorization-rigidity}
设 \(\mathcal G\) 为遍历强连通的可见群胚，并设时间余圈类 \([\widetilde\ell]\in H^1(\mathcal G;\mathbb R)\) 落在正锥内部（等价于存在全边严格正 gauge；见命题 \ref{prop:xi-positive-time-cone-criterion} 的有限图版本）。
若 \(\mathcal G\) 在 Morita 意义下存在非平凡乘积分解
\[
\mathcal G\ \simeq\ \mathcal G^{(1)}\times \mathcal G^{(2)},
\]
并且两因子在审计意义下独立（证书复合张成张量积、代价/迹/primitive 指纹按因子分解），
则时间同调类必分裂为两条非共线正余圈之和：存在 \([\widetilde\ell^{(i)}]\in H^1(\mathcal G^{(i)};\mathbb R)\)（各自为正锥内点）使
\[
[\widetilde\ell]\ =\ [\widetilde\ell^{(1)}]\oplus[\widetilde\ell^{(2)}],
\]
并且 \([\widetilde\ell^{(i)}]\) 分别由 \(\mathcal G^{(i)}\) 内禀决定（至 gauge 余边界）。
\par
反向地，若正锥内存在分裂
\[
[\widetilde\ell]=[\alpha]+[\beta],
\qquad
[\alpha],[\beta]\ \text{均为正锥内点且支撑互不纠缠（在证书层可交换到同伦）},
\]
则 \(\mathcal G\) 必在 Morita 意义下分解为两个非平凡因子，且对应的时间余圈类（至比例/余边界）即为 \([\alpha]\)、\([\beta]\)。
因此，协议的“不可约性”可由时间正锥的极端射线刻画：正锥的极端射线与 Morita-不可约因子之间存在自然对应。
\end{theorem}
\begin{proof}
Morita 乘积分解在群胚代数层面对应张量分解；在该分解下，任何可加余圈（时间）必为两因子余圈之和，且正性在因子上保持，从而得到第一段结论。
反向方向中，“支撑互不纠缠/可交换到同伦”确保证书复合与指纹输出在两族生成元上张量分裂；因此可按因子生成元分别闭合得到两个子群胚，并由普适性给出 Morita 分解。
极端射线刻画来自：若 \([\widetilde\ell]\) 可在正锥内写为两条非共线内点之和，则协议在上述意义下可分解；反之 Morita-不可约排除该类分裂。证毕。
\end{proof}

\begin{theorem}[时间统一推出 RH 的同调化归约：偏离临界线等价于非平凡时间 Holonomy]\label{thm:xi-rh-cohomological-reduction-holonomy}
假设存在完成化表示，使得经典完成化函数 \(\Xi(s)\) 可写为时间加权行列式族的限制：
\[
\Xi(s)=\det\!\bigl(I-e^{-s}\mathcal L_{\theta(s)}\bigr),
\qquad
\theta(s):=\Re(s)-\tfrac12,
\]
其中 \(\theta\mapsto\mathcal L_{\theta}\) 为由可审计局部演化与时间余圈诱导的转移算子族，并满足时间反演共轭
\(
\mathcal L_{\theta}\sim \mathcal L_{-\theta}
\)
对应于 \(\Xi(s)=\Xi(1-s)\)。
若进一步满足：
\begin{enumerate}
  \item 严格正时间：任意非平凡闭环的时间和非零；
  \item 尺度--时间联络的 Holonomy 群 \(\mathfrak H\) 平凡（定理 \ref{thm:xi-scale-time-connection-holonomy-criterion}）；
  \item 临界切片的压力判据成立（定理 \ref{thm:xi-critical-slice-pressure-zero-drift-lyapunov}），从而 \(\theta=0\) 为唯一的“双零”切片；
\end{enumerate}
则对任意 \(\theta\neq 0\) 必有 \(\rho(\mathcal L_{\theta})\neq 1\)，从而 \(\det(I-e^{-s}\mathcal L_{\theta})\neq 0\) 不可能发生在 \(\theta\neq 0\) 的切片上。
因此所有非平凡零点必满足 \(\theta(s)=0\)，即 \(\Re(s)=\tfrac12\)。
在上述表示假设下，RH 可被等价重述为“尺度--时间 Holonomy 消失”所强制的临界切片唯一性。
\end{theorem}

\begin{proof}
由定理 \ref{thm:xi-critical-slice-pressure-zero-drift-lyapunov}，在可审计微可逆与时间反演口径下，\(\theta=0\) 为唯一同时满足零漂移与零 Lyapunov 的切片。
Holonomy 平凡排除在跨尺度拼合中引入额外中心补偿的自由度，从而临界性不再可沿 \(\theta\neq 0\) 方向保持。
因此对 \(\theta\neq 0\) 必有 \(\rho(\mathcal L_{\theta})\neq 1\)，与行列式零点发生在非临界谱条件矛盾；结论遂得。
\end{proof}

\begin{proposition}[RH 的有限可判据：仅需检查有限生成的时间曲率矩阵的整秩消失]\label{prop:xi-rh-finite-curvature-matrix-criterion}
在定理 \ref{thm:xi-rh-cohomological-reduction-holonomy} 的表示假设下，若协议来自有限生成的收敛重写系统，则 Holonomy 群 \(\mathfrak H\) 由有限多个“尺度--时间临界方块”的缺陷生成。
将这些缺陷在某一中心坐标系（\(\mathbb Z^d\)）下记为整数向量，并排列成矩阵 \(B\in\mathbb Z^{r\times d}\)（行对应临界方块，列对应中心基向量），则
\[
\mathfrak H\cong \mathrm{im}(B)\subset\mathbb Z^d,
\qquad
\mathrm{rank}(\mathfrak H)=\mathrm{rank}_{\mathbb Z}(B).
\]
特别地在标量缺陷 \(d=1\) 情形，\(\mathfrak H=\{0\}\) 当且仅当 \(B=0\)。
因此在该框架内，“Holonomy 消失”归约为对有限整数矩阵的零性/秩判别，且与 gauge 选择无关。
\end{proposition}

\begin{proof}
收敛重写系统的所有高阶重叠由有限临界对生成；Holonomy 作为交换方块的缺陷累积，故由有限临界方块的缺陷生成。
把生成元在 \(\mathbb Z^d\) 基下坐标化即得矩阵描述；秩结论由 \(\mathbb Z\)-线性代数给出。
\end{proof}

\paragraph{正时间规范、Smith 归约与 PBW 可交换化}

\begin{proposition}[正时间的同调正锥判据：全边正 gauge 的存在性]\label{prop:xi-positive-time-cone-criterion}
\leanverified{paper\_xi\_positive\_time\_cone\_criterion}
设 \(\Gamma=(V,E)\) 为有限强连通有向图，\(t:E\to\mathbb R\) 为时间权，并作 gauge 变换 \(t\mapsto t+\partial b\)（\(b:V\to\mathbb R\)）。
则存在 \(b\) 使得 \(t+\partial b\) 在所有边上严格为正，当且仅当对每条有向闭环 \(\gamma\) 有
\[
\sum_{e\in\gamma} t(e)>0.
\]
等价地，\([t]\in H^1(\Gamma;\mathbb R)\) 落在由所有有向闭环定义的正锥
\[
\mathcal C_+:=\Bigl\{[c]\in H^1:\ \langle[c],[\gamma]\rangle>0\ \text{对一切非平凡闭环 }\gamma\Bigr\}
\]
的内部。
\end{proposition}

\begin{proof}
对任意闭环 \(\gamma\)，有 \(\sum_{e\in\gamma}\partial b(e)=0\)，故闭环和是同调不变量。
若存在 \(b\) 使各边正，则任意闭环和必正，得必要性。
反向为差分约束的可行性判据：不等式 \(t(e)+b(r(e))-b(s(e))>0\) 等价于严格差分约束系统，其可行性当且仅当不存在闭环总权 \(\le 0\)；在本设定中即为所有闭环和严格为正。
\end{proof}

\begin{theorem}[时间的热带化完备性：一切“可组合—可比较”的协议代价必因子化为非阿基米德赋值]\label{thm:xi-time-tropicalization-completeness}
设 \(\mathbf{PW}\) 为给定协议的可审计 \(2\)-范畴（定义 \ref{def:pom-projword-2cat}），并令 \(\mathfrak S\) 为其\emph{效果半环}：
以同源同靶 \(1\)-态射在 \(2\)-同伦下的等价类为元素，
取 \(\otimes\) 为可组合处的复合，
取 \(\oplus\) 为“并置/择一”生成的幂等和（在可比较对象上交换、结合且幂等），并令 \(0\) 为不可达元、\(1\) 为恒等元。
设函数
\[
T:\mathfrak S\to \overline{\mathbb R}_{\ge 0}:=\mathbb R_{\ge 0}\cup\{+\infty\}
\]
满足三条公理：
\begin{enumerate}
  \item \textbf{复合加性：}\(T(a\otimes b)=T(a)+T(b)\)；
  \item \textbf{择一取最小：}\(T(a\oplus b)=\min\{T(a),T(b)\}\)；
  \item \textbf{规范：}\(T(0)=+\infty,\ T(1)=0\)。
\end{enumerate}
则 \(T\) 必唯一因子化为到热带半环 \(\mathbb T=(\overline{\mathbb R}_{\ge 0},\min,+)\)（\S\ref{subsec:pom-cost-functor}）的半环同态；亦即 \(T\) 本身就是一个 \(\mathbb T\)-值的非阿基米德赋值（tropical valuation）。
\par
并且在有限生成实现中（以有向生成图 \(\Gamma\) 表示基本步），由 \(t(e):=T(e)\) 所诱导的 \(1\)-余圈类 \([t]\in H^1(\Gamma;\mathbb R)\) 落在正锥内部（命题 \ref{prop:xi-positive-time-cone-criterion}）。
反之，任意落在正锥内部的时间余圈类（允许 gauge 余边界扰动）都在上述构造下给出一族等价的 \(T\)（同调主类一致）。
因此，“时间统一”在代数层面等价于：协议的全部可比较资源结构在热带意义下均由同一非阿基米德赋值控制。
\end{theorem}
\begin{proof}
三条公理恰是“乘法送到 \(+\)、加法送到 \(\min\)、单位/零元保持”的半环同态条件；故 \(T:\mathfrak S\to\mathbb T\) 为半环同态。
唯一性来自：\(\mathbb T\) 的运算已被 \(+\) 与 \(\min\) 固定，任何满足三条公理的映射在生成元上取值一旦给定，其在 \(\mathfrak S\) 上的延拓即被复合加性与幂等和的最小律完全决定。
\par
在有限生成实现中，把 \(T\) 限制到基本边集得到边权 \(t:E\to\mathbb R\)；其在闭环上的和只依赖于同调类。
当 \([t]\) 落在正锥内部时，命题 \ref{prop:xi-positive-time-cone-criterion} 给出全边严格正 gauge，从而得到严格正时间的可比较赋值。
反向方向中，gauge 余边界扰动不改变闭环周期和，故不改变同调主类；而在同调主类固定下，\(T\) 的选择仅对应于在代表元层面的 gauge 规范与等价实现的选择。证毕。
\end{proof}

\begin{proposition}[时间缺陷群的 Smith 归约：中心扩张的自由度与扭结结构可被完全计算]\label{prop:xi-time-defect-smith-reduction}
沿命题 \ref{prop:xi-rh-finite-curvature-matrix-criterion} 的矩阵化描述，设缺陷矩阵为 \(B\in\mathbb Z^{r\times d}\) 并令 \(\mathfrak H:=\mathrm{im}(B)\subset\mathbb Z^d\)。
则存在可逆整数矩阵 \(U\in\mathrm{GL}_r(\mathbb Z)\)、\(V\in\mathrm{GL}_d(\mathbb Z)\) 使
\[
UBV=\mathrm{diag}(s_1,\dots,s_k,0,\dots,0),
\qquad
s_i\mid s_{i+1}.
\]
从而缺陷余商具有标准分解
\[
\mathbb Z^d/\mathfrak H
\ \cong\
\bigoplus_{i=1}^k\mathbb Z/s_i\mathbb Z\ \oplus\ \mathbb Z^{d-k}.
\]
其中自由秩 \(d-k\) 给出最小中心扩张的连续参数维数，扭因子 \(s_i\) 给出不可消去的离散“时间扭结”。
\end{proposition}

\begin{proof}
由 Smith 标准形定理，任意整数矩阵在左右可逆整数变换下可化为上述对角形；对角元给出像子群的结构不变量，从而导出商群分解。
\end{proof}

\begin{theorem}[缺陷的“谱指纹”：时间扭结的有限阶 torsion 必然外显为可见计数的准周期振荡]\label{thm:xi-time-torsion-spectral-fingerprint}
沿命题 \ref{prop:xi-finite-state-time-image-discrete-or-dense} 的离散化口径，设时间像落在 \(h\mathbb Z\)（某个 \(h>0\)）。
令 \(\mathcal T_t\) 为时间预算 \(t\) 下的不可区分商树（定理 \ref{thm:xi-time-boundary-double-reconstruction}），并记可见节点计数
\[
N(t):=\card(\mathcal T_t),\qquad t\in h\mathbb Z_{\ge 0}.
\]
在“有限生成 + 混合无漂移（临界切片）”条件下，存在常数 \(\kappa>0\)、\(r\ge 0\) 与有界函数 \(\Phi:\mathbb Z_{\ge 0}\to(0,\infty)\)，使得当 \(t\to\infty\)（沿 \(h\mathbb Z\)）有渐近分解
\[
\boxed{\;
N(t)=e^{\kappa t}\,t^{r}\,\bigl(\Phi(t/h)+o(1)\bigr)\; }.
\]
令
\[
q:=\exp\!\bigl(\mathrm{Tor}(\mathbb Z^d/\mathfrak H)\bigr)
\]
为扭结余商 \(\mathrm{Tor}(\mathbb Z^d/\mathfrak H)\) 的指数（由命题 \ref{prop:xi-time-defect-smith-reduction} 的 Smith 分解读出）。
则 \(\Phi\) 在充分大 \(n\) 上满足严格 \(q\)-准周期性：
\[
\Phi(n+q)=\Phi(n)\qquad(n\gg 1).
\]
并且若 \(q=1\)（torsion 平凡），则 \(\Phi\) 退化为常数。
因此 torsion 缺陷不需要直接访问同调计算即可被观测：它必然在指数主项之外留下可检出的有限周期振荡。
\end{theorem}
\begin{proof}
由命题 \ref{prop:xi-time-defect-smith-reduction}，余商 \(\mathbb Z^d/\mathfrak H\) 的 torsion 部分为有限阿贝尔群，其特征标（角色）集合有限；在离散时间 \(t\in h\mathbb Z\) 上，对应的计数/配分型生成函数可按有限角色分解为有限个“根为单位根的代数型分量”之和。
这些单位根分量在系数渐近中产生有限周期的振荡调制；其共同周期由 torsion 指数 \(q\) 控制。
若 torsion 平凡，则不存在非平凡单位根分量，调制函数只能为常数。证毕。
\end{proof}

\begin{proposition}[缺陷维数的统计可识别性：平衡态下时间协方差矩阵的核空间精确等于缺陷自由度]\label{prop:xi-defect-free-rank-identifiable-by-covariance-kernel}
在最小中心扩张后，令时间增量 cocycle 取值于 \(\mathbb Z^d\)，并在某个温度参数 \(\beta>\beta_c\) 的平衡态下令
\[
S_n:=\sum_{k=0}^{n-1}\Delta T(\omega_k)\in\mathbb R^{d}
\]
为 \(n\) 步时间增量和。
在“混合性 + 有界局部生成”条件下，渐近协方差极限
\[
\Sigma:=\lim_{n\to\infty}\frac{1}{n}\,\mathrm{Cov}(S_n)
\]
存在。
取命题 \ref{prop:xi-time-defect-smith-reduction} 的 Smith 坐标，使
\(
\mathbb Z^d/\mathfrak H\cong \mathrm{Tor}\oplus \mathbb Z^{d-k}
\)，并以该分解的自由直和因子在 \(\mathbb Z^d\) 中选定一个提升格，记为 \(\mathfrak H_{\mathrm{free}}\cong\mathbb Z^{d-k}\)。
则有
\[
\ker(\Sigma)=\mathrm{span}_{\mathbb R}\bigl(\mathfrak H_{\mathrm{free}}\bigr),
\qquad
\dim\ker(\Sigma)=\mathrm{rank}(\mathfrak H_{\mathrm{free}})=d-k.
\]
因此缺陷自由度可由纯统计量（平衡态下的时间涨落）直接读出：时间统一越彻底，协方差越非退化；任何非平凡缺陷自由度都会强制出现与之同维的零模态。
\end{proposition}
\begin{proof}
对平稳遍历的加性泛函，Gordin/鞅逼近给出：方向 \(v\in\mathbb R^d\) 上的渐近方差为零当且仅当标量泛函 \(v\cdot\Delta T\) 在平衡态下同调平凡（为余边界/共边界，从而可写为 \(f-f\circ\sigma\) 加常数）。
在最小中心扩张中，所有可被 gauge 吸收的“零涨落方向”恰对应于 Smith 分解中自由直和因子的提升格 \(\mathfrak H_{\mathrm{free}}\)；其余方向均携带不可消去的涨落。
故 \(\ker(\Sigma)\) 精确等于 \(\mathrm{span}_{\mathbb R}(\mathfrak H_{\mathrm{free}})\)，维数结论随即得到。证毕。
\end{proof}

\begin{theorem}[时间滤过的 PBW 型定理：混合曲率零当且仅当关联分级代数可交换]\label{thm:xi-time-filtration-pbw-graded-commutative}
令 \(\mathbb k[\mathbf{PW}]\) 为协议范畴代数（\(\mathbb k\) 为特征零域），并以时间长度 \(\widetilde{\ell}\) 定义滤过
\[
F_t:=\mathrm{span}\{w:\widetilde{\ell}(w)\le t\}.
\]
若 \(\widetilde{\ell}\) 对允许复合严格可加，则 \(\mathrm{gr}_{\widetilde{\ell}}\mathbb k[\mathbf{PW}]\) 为分级代数。
设系统存在两类生成（对应加法层与乘法层）并形成交错方块，则
\[
\mathrm{gr}_{\widetilde{\ell}}\mathbb k[\mathbf{PW}]\ \text{为交换代数}
\quad\Longleftrightarrow\quad
K\equiv 0,
\]
其中 \(K\) 为定理 \ref{thm:xi-mixed-curvature-zero-criterion} 的混合曲率。
\end{theorem}

\begin{proof}
若 \(K\equiv 0\)，交错方块两路径在时间主项上无缺陷，故交换关系在关联分级的首项中成立，从而分级代数可交换。
反之，若关联分级可交换，则所有交错生成元在首项中交换，故交换方块的时间缺陷必须为零，即 \(K\equiv 0\)。
\end{proof}

\begin{theorem}[稳定类型的谱表示：稳定分类等价于时间滤过关联分级半环的极小素谱]\label{thm:xi-stable-type-minspec-graded-semiring}
在定理 \ref{thm:xi-time-tropicalization-completeness} 的设定下，令 \((\mathfrak S,T)\) 为协议的效果半环与其热带时间赋值。
定义时间滤过
\[
F_t:=\{\,a\in\mathfrak S:\ T(a)\le t\,\},
\qquad
F_{<t}:=\bigcup_{u<t}F_u,
\]
并定义关联分级半环（Rees/分级商的幂等版本）
\[
\mathrm{gr}_T(\mathfrak S):=\bigoplus_{t\in\mathrm{im}(T)} F_t/F_{<t},
\]
其中 \(F_t/F_{<t}\) 表示将所有严格低阶元素 \(F_{<t}\) 在第 \(t\) 级上压缩为零元的自然商。
在“局部有限生成 + 审计合流”条件下，\(\mathrm{gr}_T(\mathfrak S)\) 为交换幂等半环。
则其极小素理想集合 \(\mathrm{MinSpec}(\mathrm{gr}_T(\mathfrak S))\) 与协议的稳定类型集合之间存在自然双射：每个微观态在时间极限下所收敛的稳定类型，恰对应于 \(\mathrm{gr}_T(\mathfrak S)\) 的一个极小素点。
\par
并且任何改变扫描细节但保持时间同调类（从而保持 \(T\) 的同调主类）的实现，都给出同构的 \(\mathrm{gr}_T(\mathfrak S)\)；
因此稳定类型在该口径下是时间的内禀不变量，而非外加定义。
\end{theorem}
\begin{proof}
时间滤过把“可组合—可比较”的协议演算压缩为按主导时间阶组织的层级结构；在幂等口径下，\(\oplus=\min\) 使“并置/择一”的比较语义与主阶提取相容，从而关联分级半环捕捉到在 \(t\to\infty\) 极限中稳定存活的首项组合结构。
在合流与局部有限生成条件下，首项组合结构不依赖具体规约路径，故给出实现无关的交换幂等半环。
极小素理想刻画不可再分的首项成分（不可约“谱块”），其点正对应于时间推进下的稳定极限类（稳定类型）。
同调主类不变性则来自：gauge 余边界扰动不改变闭环周期和与首项比较结构，因而不改变 \(\mathrm{gr}_T(\mathfrak S)\) 的同构类。证毕。
\end{proof}

\paragraph{时间边界与统计退化}

\begin{theorem}[时间边界的 Patterson--Sullivan 密度：KMS 平衡态与边界共形测度等价]\label{thm:xi-time-boundary-patterson-sullivan-kms}
若时间度量 Cayley 图满足线性等周界（定理 \ref{thm:xi-audit-dehn-function-time-geometry} 的双曲性条件），并设增长指数 \(\kappa\) 存在。
对 \(\lambda>\kappa\) 定义 Poincar\'e 级数
\[
\mathcal P_\lambda(x,y):=\sum_{g:x\to y}e^{-\lambda \widetilde{\ell}(g)}.
\]
则存在唯一（至尺度）\(\lambda\)-共形的边界密度族 \(\{\nu_x^{(\lambda)}\}\) 支持于 Gromov 边界 \(\partial\)，并满足由 Busemann--时间函数给出的 Radon--Nikodym 关系。
并且该边界密度诱导的极端态与由 Gibbs 权重 \(e^{-\lambda\widetilde{\ell}}\) 所定义的 KMS 平衡态之间存在自然一一对应。
\end{theorem}

\begin{proof}
在双曲度量下，Poincar\'e 级数收敛域与共形密度的存在性由 Patterson--Sullivan 构造给出，且增长指数为其临界参数。
KMS 态与共形密度的对应是由时间余圈诱导的动力系统（模流/时间流）在边界上的共形性所给出的标准对应；极端分解与边界测度的极端性一致。
\end{proof}

\begin{proposition}[缺陷秩控制大偏差退化：\(\mathrm{rank}(\mathfrak H)\) 等于速率函数的仿射平坦维数]\label{prop:xi-defect-rank-controls-ldp-flatness}
设压力函数 \(P(s)=\log\rho(A_s)\) 在 \(s>s_c\) 可微且严格凸，其中 \(A_s\) 为时间加权转移算子。
对 gauge 扰动 \(\widetilde{\ell}\mapsto \widetilde{\ell}+\partial b\) 有 \(P\) 不变；对更一般的中心缺陷方向（由 Holonomy 缺陷群 \(\mathfrak H\) 给出）压力沿某些线性方向保持不变。
令大偏差速率函数 \(I\) 为 \(P\) 的 Legendre 对偶，则
\[
\dim\bigl(\text{\(I\) 的最大仿射平坦子集}\bigr)=\mathrm{rank}(\mathfrak H).
\]
特别地，\(\mathfrak H=\{0\}\) 当且仅当 \(I\) 严格凸；一旦存在非平凡缺陷自由度，速率函数必出现相应维数的平坦“相共存”方向。
\end{proposition}

\begin{proof}
压力在缺陷方向上的不变性意味着 \(P\) 在相应线性子空间上具有平坦方向；Legendre 对偶将平坦方向对应为速率函数的仿射平坦面。
维数由该不变子空间（等价于缺陷自由度）的秩给出。
\end{proof}

\paragraph{周期和的有限判别性：同调类的完全重建}

\begin{proposition}[周期和决定时间同调类：有限基循环的时间和相等即余边界等价]\label{prop:xi-cycle-sums-determine-time-cohomology}
\leanverified{paper\_xi\_cycle\_sums\_determine\_time\_cohomology}
设 \(\Gamma=(V,E)\) 为强连通有限有向图，\(t,t':E\to h\mathbb Z\) 为两组时间权，视为 \(1\)-余圈。
则
\[
t-t'=\partial b\ \text{对某 }b:V\to h\mathbb Z
\quad\Longleftrightarrow\quad
\sum_{e\in\gamma}t(e)=\sum_{e\in\gamma}t'(e)\ \text{对所有有向闭环 }\gamma.
\]
并且“对所有闭环”的检验可缩减为对任意生成 \(H_1(\Gamma;\mathbb Z)\) 的有限基本循环族即可：取任一生成树 \(T\)，对每条弦边 \(e\in E\setminus T\) 所形成的基本循环 \(\gamma_e\) 构成整数循环空间一组基，检查有限多条 \(\{\gamma_e\}\) 的时间和相等即推出同调等价。
\end{proposition}

\begin{proof}
必要性由闭环上 \(\sum\partial b=0\) 直接得出。
反向方向是有限图上 \(1\)-上同调的标准刻画：若两余圈在一组循环基上的配对一致，则其差在 \(H^1\) 中为零，从而为余边界。
\end{proof}

\paragraph{码长 \texorpdfstring{$(1,2)$}{(1,2)} 的 KMS 锁定与黄金标度：模尺度群、重建阈值与两态纤维律的三域对齐}
在本文把时间固定为单生成扫描算子的迭代计数、并以二元窗读出与规范归一化层作为协议约束的口径下，码长 \((1,2)\) 的前缀分支（\(0\) 与 \(10\)）给出一个完全显式的算子论实现：其 KMS 逆温锁定为 \(\log\varphi\)，而折叠造成的信息缺口以 \(\log(2/\varphi)=\log 2-\log\varphi\) 量化；两者并非独立常数，而是在同一 KMS--纤维两态律下共同外部化的结构不变量。
\par
下述结论把该对齐写为可独立引用的定理链条。
\par

\begin{theorem}[码长 \texorpdfstring{$(1,2)$}{(1,2)} 的 KMS 因子与离散模尺度群：\texorpdfstring{$\varphi^{\mathbb Z}$}{phi^Z} 的刚性]\label{thm:xi-kms-code12-discrete-modular-scale-group}
记 Cuntz 代数 \(\mathcal{O}_2=C^\ast(\mathsf{S}_0,\mathsf{S}_1)\)，其中 \(\mathsf{S}_0,\mathsf{S}_1\) 为等距元并满足
\[
\mathsf{S}_i^\ast\mathsf{S}_j=\delta_{ij}1,
\qquad
\mathsf{S}_0\mathsf{S}_0^\ast+\mathsf{S}_1\mathsf{S}_1^\ast=1.
\]
在 \(\mathcal{O}_2\) 上定义准自由动力学 \((\alpha_t)_{t\in\mathbb R}\)：
\[
\alpha_t(\mathsf{S}_0)=e^{it}\mathsf{S}_0,
\qquad
\alpha_t(\mathsf{S}_1)=e^{2it}\mathsf{S}_1.
\]
则存在唯一 KMS 态 \(\omega\)，其逆温 \(\beta\) 由方程 \(e^{-\beta}+e^{-2\beta}=1\) 唯一确定，故
\[
\beta=\log\varphi.
\]
令 \((\pi_\omega,\mathcal H_\omega,\Omega_\omega)\) 为 GNS 表示并记 \(\mathcal M:=\pi_\omega(\mathcal{O}_2)''\)。
则 \(\omega\) 的模块算子 \(\Delta_\omega\) 的正谱离散且满足
\[
\mathrm{spec}(\Delta_\omega)\cap\mathbb R_{>0}=\varphi^{\mathbb Z},
\]
从而 Connes 的尺度不变量满足
\[
S_{\mathrm{Connes}}(\mathcal M)=\{0\}\cup\varphi^{\mathbb Z}.
\]
特别地，
\[
\mathbb Z[\varphi]_{>0}^\times=\varphi^{\mathbb Z},
\]
故该模尺度群与二次数域整数环 \(\mathbb Z[\varphi]\) 的正单位群一致。
\end{theorem}

\begin{corollary}[周期模流与 \texorpdfstring{$T_{\mathrm{Connes}}$}{T}-不变量：\texorpdfstring{$\log\varphi$}{log phi} 为基本模尺度]\label{cor:xi-kms-code12-connes-T-invariant}
在定理 \ref{thm:xi-kms-code12-discrete-modular-scale-group} 的设定下，
Connes 的 \(T\)-不变量为离散晶格
\[
T_{\mathrm{Connes}}(\mathcal M)=\frac{2\pi}{\log\varphi}\,\mathbb Z,
\]
从而模自同构群的内实现时刻形成基本间距由 \(\log\varphi\) 锁定的离散网格。
\end{corollary}

\begin{corollary}[类型刚性与流权：\texorpdfstring{$\mathcal M$}{M} 为超有限 \texorpdfstring{$\mathrm{III}_{1/\varphi}$}{III_1/phi} 因子]\label{cor:xi-kms-code12-type-rigidity}
在定理 \ref{thm:xi-kms-code12-discrete-modular-scale-group} 的设定下，
\(\mathcal M\) 为超有限 \(\mathrm{III}_{1/\varphi}\) 因子；
其 flow of weights 同构于圆周 \(\mathbb R/(\log\varphi)\mathbb Z\) 上的平移流。
因此该因子的尺度结构（至同构）不存在连续可调自由度，完全由 \(\log\varphi\) 一参决定。
\end{corollary}

\begin{proposition}[折叠纤维块的中心化逼近：有限块代数在 \texorpdfstring{$\omega$}{omega}-中心化子中 \texorpdfstring{$\sigma$}{sigma}-强稠密]\label{prop:xi-fold-microblock-centralizer-strong-density}
令 \(X_m\) 为禁止相邻 \(11\) 的长度 \(m\) 稳定词集，\(\Omega_m:=\{0,1\}^m\)，并取二进制区间折叠（bin-fold）
\(
\mathrm{Fold}^{\mathrm{bin}}_m:\Omega_m\twoheadrightarrow X_m
\)
（附录 \ref{app:fold-multiplicity}）。
记纤维多重度 \(d_m^{\mathrm{bin}}(x):=\card{(\mathrm{Fold}^{\mathrm{bin}}_m)^{-1}(x)}\)。
定义有限块代数
\[
H_m:=\bigoplus_{x\in X_m} M_{d_m^{\mathrm{bin}}(x)}(\mathbb C).
\]
在定理 \ref{thm:xi-kms-code12-discrete-modular-scale-group} 的 KMS 表示下，令
\[
\mathcal M_\omega:=\{a\in\mathcal M:\sigma_t^\omega(a)=a,\ \forall t\in\mathbb R\}
\]
为 \(\omega\) 的中心化子代数。
则存在一族（随 \(m\)）的 \(^\ast\)-单同态 \(\iota_m:H_m\hookrightarrow \mathcal M_\omega\)，使得
\[
\overline{\bigcup_{m\ge 1}\iota_m(H_m)}^{\ \sigma\text{-strong}}=\mathcal M_\omega.
\]
并且 \(\omega\) 在 \(\iota_m(H_m)\) 上的限制为各块归一迹的凸组合；当 \(m\to\infty\) 时，其权重向量收敛到由 \(\varphi\)-平衡测度诱导的唯一极限权重。
\end{proposition}

\begin{theorem}[可逆重建的极限最优成功率：常数偏置下的分段闭式]\label{thm:xi-optimal-reconstruction-success-probability-constant-bias}
\leanverified{paper\_xi\_optimal\_reconstruction\_success\_probability\_constant\_bias}
令 \(U_m\sim\mathrm{Unif}(\Omega_m)\)，并记宏观观测随机变量
\(
\mathsf{X}_m:=\mathrm{Fold}^{\mathrm{bin}}_m(U_m)\in X_m
\)。
允许附加寄存器取值集合大小为 \(M_m\in\mathbb N\)，并允许编码器在给定 \(\omega\in\Omega_m\) 与 \(x=\mathrm{Fold}^{\mathrm{bin}}_m(\omega)\) 时产生辅助标签 \(A\in\{1,\dots,M_m\}\)，解码器据 \((x,A)\) 输出估计 \(\widehat\omega\)。
在均匀输入下的平均最优成功率为
\[
P_m^\star(M_m):=\sup \mathbb P[\widehat\omega=U_m].
\]
则 \(P_m^\star(M_m)\) 具有闭式：
\[
P_m^\star(M_m)=\frac{1}{2^m}\sum_{x\in X_m}\min\!\bigl\{d_m^{\mathrm{bin}}(x),\,M_m\bigr\}.
\]
进一步取常数偏置尺度
\[
M_m=\exp\!\Bigl(m\log\!\Bigl(\frac{2}{\varphi}\Bigr)+c\Bigr),
\qquad c\in\mathbb R\ \text{固定}.
\]
则极限
\(
P_\star(c):=\lim_{m\to\infty}P_m^\star(M_m)
\)
存在且给出显式分段公式
\[
P_\star(c)=
\begin{cases}
1, & c\ge 0,\\[6pt]
\displaystyle \frac{\varphi}{\sqrt5}\,e^{c}+\frac{1}{\varphi\sqrt5}, & -\log\varphi\le c<0,\\[10pt]
\displaystyle \frac{\varphi^2}{\sqrt5}\,e^{c}, & c<-\log\varphi.
\end{cases}
\]
\end{theorem}

\begin{proof}
闭式 \(P_m^\star(M_m)=2^{-m}\sum_x\min\{d_m^{\mathrm{bin}}(x),M_m\}\) 来自：对每条纤维 \(x\) 最优策略是在 \(M_m\) 个微态上实现单射编码，其余必发生碰撞，故纤维内最优成功率为 \(\min\{1,M_m/d_m^{\mathrm{bin}}(x)\}\)，对均匀输入取平均即得。
\par
记输出分布 \(\mu_m(x):=d_m^{\mathrm{bin}}(x)/2^m\) 并令 \(W_m\sim\mu_m\)，则
\[
P_m^\star(M_m)=\mathbb E_{\mu_m}\!\left[\min\!\left\{1,\frac{M_m}{d_m^{\mathrm{bin}}(W_m)}\right\}\right].
\]
由推论 \ref{cor:fold-bin-two-point-limit-law}，
\(
(\varphi/2)^m d_m^{\mathrm{bin}}(W_m)\Rightarrow Z\in\{1,\varphi^{-1}\}
\)
且
\(
\mathbb P(Z=1)=\varphi/\sqrt5
\)、
\(
\mathbb P(Z=\varphi^{-1})=1/(\varphi\sqrt5)
\)。
又 \(M_m/d_m^{\mathrm{bin}}(W_m)=e^{c}/Z+o(1)\)（在分布意义下），把 \(\min\{1,e^{c}/Z\}\) 在两点 \(Z\) 上逐点计算并按 \(c\) 的阈值分段，即得所示极限公式。
\end{proof}

\begin{corollary}[基数预算的强反界与线性阈值斜率]\label{cor:xi-reconstruction-cardinality-strong-converse}
\leanverified{paper\_xi\_reconstruction\_cardinality\_strong\_converse}
在定理 \ref{thm:xi-optimal-reconstruction-success-probability-constant-bias} 的设定下，对任意 \(M_m\in\NN\) 有
\[
P_m^\star(M_m)\ \le\ \frac{M_m\,|X_m|}{2^m}.
\]
因此若要求存在方案使 \(\mathbb P[\widehat\omega=U_m]\ge \delta\)（\(\delta\in(0,1)\) 固定），则必有
\[
\log M_m\ \ge\ m\log 2-\log|X_m|+\log\delta
\ =\ m\log 2-\log F_{m+2}+\log\delta.
\]
等价地，在比特尺度 \(L:=\log_2 M_m\) 下，
\[
L\ \ge\ m-\log_2|X_m|+\log_2\delta
\ =\ m-\log_2F_{m+2}+\log_2\delta.
\]
由 Binet 渐近 \(F_{m+2}=\varphi^{m+2}/\sqrt5+O(\varphi^{-m})\) 得到线性阈值
\[
L\ \ge\ \Bigl(\log_2\frac{2}{\varphi}\Bigr)m+O(1).
\]
\end{corollary}
\begin{proof}
由定理 \ref{thm:xi-optimal-reconstruction-success-probability-constant-bias} 的闭式，
\(
P_m^\star(M_m)=2^{-m}\sum_{x\in X_m}\min\{d_m^{\mathrm{bin}}(x),M_m\}
\le 2^{-m}\sum_{x\in X_m}M_m
=M_m|X_m|/2^m
\)。
将其与成功率下界 \(\delta\) 联立并取对数即得前两条不等式。
最后一条由 \(\log_2F_{m+2}=(m+2)\log_2\varphi-\frac12\log_25+o(1)\) 直接整理得到。证毕。
\end{proof}

\begin{corollary}[常数项 \texorpdfstring{$b(\varepsilon)$}{b(eps)} 的闭式反演：零色散二阶项的精确量化]\label{cor:xi-constant-bias-inversion-b-eps}
令 \(L_m(\varepsilon)\) 为在平均错误率 \(\le\varepsilon\) 下所需的最小附加信息量（自然对数），即最小 \(L\) 使存在编码解码满足成功率 \(\ge 1-\varepsilon\) 且 \(|\mathrm{aux}|\le e^{L}\)。
则存在常数项 \(b(\varepsilon)\) 使
\[
L_m(\varepsilon)=m\log\!\Bigl(\frac{2}{\varphi}\Bigr)+b(\varepsilon)+o(1).
\]
令 \(s:=1-\varepsilon\) 并记阈值 \(s_c:=\varphi/\sqrt5\)。则
\[
b(\varepsilon)=
\begin{cases}
0, & s=1,\\[6pt]
\displaystyle \log\!\Bigl(\frac{\sqrt5\,s-\varphi^{-1}}{\varphi}\Bigr), & s_c\le s<1,\\[12pt]
\displaystyle \log\!\Bigl(\frac{\sqrt5\,s}{\varphi^2}\Bigr), & 0<s<s_c.
\end{cases}
\]
\end{corollary}

\begin{proof}
由定理 \ref{thm:xi-optimal-reconstruction-success-probability-constant-bias}，成功率极限仅依赖于偏置 \(c\)，且 \(P_\star(c)\) 连续严格递增。
因此 \(b(\varepsilon)\) 为方程 \(P_\star(c)=1-\varepsilon\) 的最小解；对两段线性形式逐段反解得所示闭式。
\end{proof}
\begin{conclusion}[平均错误率约束下最小侧信息的零色散二阶项：常数项 \texorpdfstring{$b(\varepsilon)$}{b(eps)} 的闭式]\label{con:xi-min-sideinfo-avg-error-constant-b-eps}
令 \(L_m(\varepsilon)\) 为在平均错误率 \(\le\varepsilon\) 下所需的最小附加信息量（自然对数），即最小 \(L\) 使存在编码解码满足成功率 \(\ge 1-\varepsilon\) 且 \(|\mathrm{aux}|\le e^{L}\)。
则存在常数项 \(b(\varepsilon)\) 使
\[
L_m(\varepsilon)=m\log\!\Bigl(\frac{2}{\varphi}\Bigr)+b(\varepsilon)+o(1).
\]
令 \(s:=1-\varepsilon\) 并记阈值 \(s_c:=\varphi/\sqrt5\)。则
\[
b(\varepsilon)=
\begin{cases}
0, & s=1,\\[6pt]
\displaystyle \log\!\Bigl(\frac{\sqrt5\,s-\varphi^{-1}}{\varphi}\Bigr), & s_c\le s<1,\\[12pt]
\displaystyle \log\!\Bigl(\frac{\sqrt5\,s}{\varphi^2}\Bigr), & 0<s<s_c.
\end{cases}
\]
\end{conclusion}
\begin{proof}[证明要点]
推论 \ref{cor:xi-constant-bias-inversion-b-eps}。证毕。
\end{proof}



\paragraph{折叠推前相对稳定层均匀基线的二原子极限账本}
令
\[
p^{\mathrm{bin}}_m(w):=\frac{d^{\mathrm{bin}}_m(w)}{2^m}\qquad(w\in X_m)
\]
为 bin-fold 在稳定类型层 \(X_m\) 上的推前分布（亦即推论 \ref{cor:fold-bin-two-point-limit-law} 的 \(\mu_m\)），并令
\[
u_m(w):=\frac{1}{|X_m|}
\qquad(w\in X_m)
\]
为稳定层均匀分布。
记相对密度
\[
R_m(w):=\frac{p^{\mathrm{bin}}_m(w)}{u_m(w)}=|X_m|\,p^{\mathrm{bin}}_m(w),
\]
并采用全变差距离的标准约定
\[
D_{\mathrm{TV}}(p\Vert q):=\frac12\sum_{w}\bigl|p(w)-q(w)\bigr|.
\]

\begin{theorem}[bin-fold 推前分布相对均匀基线的二原子极限：TV 常数、任意 \(f\)-散度与 \texorpdfstring{$L^p$}{Lp} 谱]\label{thm:xi-foldbin-pushforward-uniform-two-atom-ledger}
\leanverified{paper\_xi\_foldbin\_pushforward\_uniform\_two\_atom\_ledger}
在定理 \ref{thm:fold-bin-two-state-asymptotic} 的两态一致渐近口径下，令
\[
r_0:=\frac{\varphi^2}{\sqrt5},\qquad r_1:=\frac{\varphi}{\sqrt5}.
\]
则：
\begin{enumerate}
  \item \textbf{（全变差极限常数）}极限存在且为
  \[
  \lim_{m\to\infty}D_{\mathrm{TV}}\bigl(p^{\mathrm{bin}}_m\Vert u_m\bigr)
  =\frac12\Bigl(\frac{1}{\sqrt5}-\frac{1}{\varphi^3}\Bigr),
  \]
  并且误差为 \(O\bigl((\varphi/2)^m\bigr)\)。
  \item \textbf{（任意 \(f\)-散度的二原子闭式）}对任意凸函数 \(f:(0,\infty)\to\mathbb R\) 满足 \(f(1)=0\)，定义
  \[
  D_f\bigl(p^{\mathrm{bin}}_m\Vert u_m\bigr)
  :=\sum_{w\in X_m}u_m(w)\,f\!\left(\frac{p^{\mathrm{bin}}_m(w)}{u_m(w)}\right)
  =\sum_{w\in X_m}u_m(w)\,f\!\bigl(R_m(w)\bigr).
  \]
  则极限存在并满足二原子闭式
  \[
  \lim_{m\to\infty}D_f\bigl(p^{\mathrm{bin}}_m\Vert u_m\bigr)
  =\frac{1}{\varphi}\,f(r_0)+\frac{1}{\varphi^2}\,f(r_1),
  \]
  且误差为 \(O\bigl((\varphi/2)^m\bigr)\)。
  \item \textbf{（相对密度的 \(L^p\) 极限谱与动态范围）}对任意 \(p\in[1,\infty)\)，有
  \[
  \lim_{m\to\infty}\norm{R_m-1}_{L^p(u_m)}
  =
  \Bigl[\frac{1}{\varphi}\bigl|r_0-1\bigr|^p+\frac{1}{\varphi^2}\bigl|1-r_1\bigr|^p\Bigr]^{1/p}.
  \]
  并且
  \[
  \lim_{m\to\infty}\norm{R_m}_{L^\infty(u_m)}=r_0,\qquad
  \lim_{m\to\infty}\operatorname*{ess\,inf}_{u_m}R_m=r_1,
  \qquad
  \frac{r_0}{r_1}=\varphi.
  \]
  \item \textbf{（Pinsker 缺口的常数化）}令
  \(
  C_{\mathrm{TV}}:=\frac12(\frac{1}{\sqrt5}-\varphi^{-3})
  \)
  并令
  \(
  C_{\mathrm{KL}}:=\lim_{m\to\infty}D_{\mathrm{KL}}(p^{\mathrm{bin}}_m\Vert u_m)
  \)
  （其闭式见命题 \ref{prop:fold-bin-kl-constant}）。
  则
  \[
  \lim_{m\to\infty}\frac{2\,D_{\mathrm{TV}}\bigl(p^{\mathrm{bin}}_m\Vert u_m\bigr)^2}{D_{\mathrm{KL}}\bigl(p^{\mathrm{bin}}_m\Vert u_m\bigr)}
  =\frac{2\,C_{\mathrm{TV}}^2}{C_{\mathrm{KL}}}\in(0,1),
  \]
  从而 Pinsker 不等式在该 bin-fold 极限下严格不饱和。
\end{enumerate}
\end{theorem}

\begin{proof}
由定理 \ref{thm:fold-bin-two-state-asymptotic}，
\(
d^{\mathrm{bin}}_m(w)=\varphi^{-w_m}(2/\varphi)^m+O(1)
\)
且误差在 \(w\in X_m\) 上一致。
又 \(|X_m|=F_{m+2}=\varphi^{m+2}/\sqrt5+O(\varphi^{-m})\)（Binet 公式），故对所有 \(w\in X_m\) 有一致展开
\[
R_m(w)
=\frac{|X_m|}{2^m}\,d^{\mathrm{bin}}_m(w)
=\frac{\varphi^{m+2}}{\sqrt5}\cdot\frac{1}{\varphi^m}\cdot\varphi^{-w_m}+O\!\left(\Bigl(\frac{\varphi}{2}\Bigr)^m\right)
=\varphi^{-w_m}\,r_0+O\!\left(\Bigl(\frac{\varphi}{2}\Bigr)^m\right),
\]
从而 \(R_m(w)=r_{w_m}+O((\varphi/2)^m)\)（其中 \(r_{0}=r_0,r_{1}=r_1\)）且一致于 \(w\)。
另一方面，稳定层均匀分布下末位边际满足
\[
u_m(w_m=0)=\frac{F_{m+1}}{F_{m+2}}=\frac{1}{\varphi}+O(\varphi^{-m}),
\qquad
u_m(w_m=1)=\frac{F_{m}}{F_{m+2}}=\frac{1}{\varphi^2}+O(\varphi^{-m}).
\]
由此对任意连续函数 \(g\)（例如本定理各项出现的 \(g(t)=|t-1|\)、\(g(t)=f(t)\)、\(g(t)=|t-1|^p\)、\(g(t)=t\log t\)）可得
\[
\sum_{w\in X_m}u_m(w)\,g\!\bigl(R_m(w)\bigr)
=\frac{1}{\varphi}\,g(r_0)+\frac{1}{\varphi^2}\,g(r_1)+O\!\left(\Bigl(\frac{\varphi}{2}\Bigr)^m\right),
\]
从而 (2)(3) 成立。
(1) 取 \(g(t)=\frac12|t-1|\)，并用 \(r_0>1>r_1\) 得到所示常数；其化简使用恒等式
\(
\varphi-\varphi^{-1}=1
\)
与
\(
\varphi^{-2}-\varphi^{-1}=-\varphi^{-3}
\)。
(4) 由 (2) 取 \(f(t)=t\log t\) 得 \(C_{\mathrm{KL}}>0\)，再由 Pinsker 不等式与 \(p^{\mathrm{bin}}_m\neq u_m\)（充分大 \(m\)）得极限比值严格落在 \((0,1)\)。
证毕。
\end{proof}
\begin{conclusion}[推前分布相对稳定层均匀基线的二原子极限：TV 常数、任意 \(f\)-散度与 \texorpdfstring{$L^p$}{Lp} 谱闭式]\label{con:xi-foldbin-pushforward-uniform-two-atom-ledger}
令 \(p_m^{\mathrm{bin}}(w):=d_m^{\mathrm{bin}}(w)/2^m\) 为 bin-fold 在 \(X_m\) 上的推前分布，令 \(u_m(w):=1/|X_m|\) 为稳定层均匀基线，并记相对密度 \(R_m(w):=p_m^{\mathrm{bin}}(w)/u_m(w)\)。
在两态一致渐近口径下，取
\[
r_0:=\frac{\varphi^2}{\sqrt5},\qquad r_1:=\frac{\varphi}{\sqrt5}.
\]
则：
\begin{enumerate}
  \item \textbf{（全变差极限常数）}极限存在且
  \[
  \lim_{m\to\infty}D_{\mathrm{TV}}\bigl(p^{\mathrm{bin}}_m\Vert u_m\bigr)
  =\frac12\Bigl(\frac{1}{\sqrt5}-\frac{1}{\varphi^3}\Bigr).
  \]
  \item \textbf{（任意 \(f\)-散度的二原子闭式）}对任意凸函数 \(f:(0,\infty)\to\mathbb R\) 满足 \(f(1)=0\)，极限存在并满足
  \[
  \lim_{m\to\infty}D_f\bigl(p^{\mathrm{bin}}_m\Vert u_m\bigr)
  =\frac{1}{\varphi}\,f(r_0)+\frac{1}{\varphi^2}\,f(r_1).
  \]
  \item \textbf{（相对密度的 \(L^p\) 极限谱与动态范围）}对任意 \(p\in[1,\infty)\)，有
  \[
  \lim_{m\to\infty}\norm{R_m-1}_{L^p(u_m)}
  =
  \Bigl[\frac{1}{\varphi}\bigl|r_0-1\bigr|^p+\frac{1}{\varphi^2}\bigl|1-r_1\bigr|^p\Bigr]^{1/p},
  \]
  并且
  \[
  \lim_{m\to\infty}\norm{R_m}_{L^\infty(u_m)}=r_0,\qquad
  \lim_{m\to\infty}\operatorname*{ess\,inf}_{u_m}R_m=r_1,
  \qquad
  \frac{r_0}{r_1}=\varphi.
  \]
\end{enumerate}
因此，“非均匀性”在极限上被压缩为由二点账本 \((r_0,r_1)\) 完整刻画的常数层：所有可交换的散度与谱范数在极限中均由 \((r_0,r_1)\) 决定。
\end{conclusion}
\begin{proof}[证明要点]
定理 \ref{thm:xi-foldbin-pushforward-uniform-two-atom-ledger}。证毕。
\end{proof}



\begin{corollary}[bin-fold 推前分布的 Shannon 亏损常数闭式：由 \texorpdfstring{$\varphi$}{phi} 与 \texorpdfstring{$\sqrt5$}{sqrt5} 刚性锁定]\label{cor:xi-foldbin-shannon-deficit-constant-closed-form}
在定理 \ref{thm:xi-foldbin-pushforward-uniform-two-atom-ledger} 的设定下，KL 散度（相对熵）极限常数为
\[
C_{\mathrm{KL}}
:=\lim_{m\to\infty}D_{\mathrm{KL}}\bigl(p^{\mathrm{bin}}_m\Vert u_m\bigr)
=\frac{3\varphi-1}{\sqrt5}\log\varphi-\frac12\log 5\ >\ 0,
\]
并且收敛误差为 \(O((\varphi/2)^m)\)。
由于 \(u_m\) 为均匀分布，亦即
\[
D_{\mathrm{KL}}\bigl(p^{\mathrm{bin}}_m\Vert u_m\bigr)=\log|X_m|-H\bigl(p^{\mathrm{bin}}_m\bigr),
\]
故 Shannon 亏损 \(\log|X_m|-H(p^{\mathrm{bin}}_m)\) 收敛到同一闭式常数 \(C_{\mathrm{KL}}\)。
\end{corollary}
\begin{proof}
由定理 \ref{thm:xi-foldbin-pushforward-uniform-two-atom-ledger} (2) 取 \(f(t)=t\log t\)，得
\[
C_{\mathrm{KL}}
=\frac{1}{\varphi}\,r_0\log r_0+\frac{1}{\varphi^2}\,r_1\log r_1,
\qquad
r_0=\frac{\varphi^2}{\sqrt5},\ r_1=\frac{\varphi}{\sqrt5}.
\]
将 \(\log(\varphi^k/\sqrt5)=k\log\varphi-\tfrac12\log 5\) 代入，并用恒等式
\(\varphi^{-1}=\varphi-1\) 与 \(\varphi+\varphi^{-1}=\sqrt5\) 化简，即得所示闭式。
误差项同定理 \ref{thm:xi-foldbin-pushforward-uniform-two-atom-ledger} (2)。
最后，由 \(u_m(w)=1/|X_m|\) 得
\(
D_{\mathrm{KL}}(p\Vert u_m)=-H(p)-\sum_w p(w)\log u_m(w)=\log|X_m|-H(p)
\)，证毕。
\end{proof}

\begin{proposition}[全 Rényi 亏损谱的闭式：\texorpdfstring{$\alpha$}{alpha}-族散度在极限中完全可解]\label{prop:xi-foldbin-renyi-deficit-spectrum-closed-form}
在定理 \ref{thm:xi-foldbin-pushforward-uniform-two-atom-ledger} 的设定下，对任意 \(\alpha>0\)、\(\alpha\neq 1\)，定义 Rényi 散度
\[
D_{\alpha}\bigl(p^{\mathrm{bin}}_m\Vert u_m\bigr)
:=\frac{1}{\alpha-1}\log\sum_{w\in X_m}u_m(w)\,R_m(w)^{\alpha}
\qquad\bigl(R_m=p^{\mathrm{bin}}_m/u_m\bigr).
\]
则极限存在并满足完全闭式
\[
\lim_{m\to\infty}D_{\alpha}\bigl(p^{\mathrm{bin}}_m\Vert u_m\bigr)
=\frac{1}{\alpha-1}\log\Bigl(\frac{1}{\varphi}r_0^\alpha+\frac{1}{\varphi^2}r_1^\alpha\Bigr)
=\frac{1}{\alpha-1}\left(
\log(\varphi^{2\alpha-1}+\varphi^{\alpha-2})-\frac{\alpha}{2}\log 5
\right),
\]
其中 \(r_0=\varphi^2/\sqrt5\)、\(r_1=\varphi/\sqrt5\)。
并且端点阶满足
\[
\lim_{\alpha\to\infty}\ \lim_{m\to\infty} D_{\alpha}\bigl(p^{\mathrm{bin}}_m\Vert u_m\bigr)
=\log r_0
=2\log\varphi-\frac12\log 5,
\qquad
\lim_{\alpha\to 1}\ \lim_{m\to\infty} D_{\alpha}\bigl(p^{\mathrm{bin}}_m\Vert u_m\bigr)
=C_{\mathrm{KL}}.
\]
\end{proposition}
\begin{proof}
由定理 \ref{thm:xi-foldbin-pushforward-uniform-two-atom-ledger} 的证明，\(R_m(w)=r_{w_m}+O((\varphi/2)^m)\) 且一致于 \(w\)，并且 \(u_m(w_m=0)=\varphi^{-1}+o(1)\)、\(u_m(w_m=1)=\varphi^{-2}+o(1)\)。
因此
\[
\sum_{w}u_m(w)\,R_m(w)^\alpha
=\frac{1}{\varphi}r_0^\alpha+\frac{1}{\varphi^2}r_1^\alpha+o(1),
\]
代入定义即得第一条极限。
进一步代入 \(r_0=\varphi^2/\sqrt5\)、\(r_1=\varphi/\sqrt5\) 并整理为
\[
\frac{1}{\varphi}\Bigl(\frac{\varphi^2}{\sqrt5}\Bigr)^\alpha+\frac{1}{\varphi^2}\Bigl(\frac{\varphi}{\sqrt5}\Bigr)^\alpha
=\frac{\varphi^{2\alpha-1}+\varphi^{\alpha-2}}{5^{\alpha/2}},
\]
得到所示闭式。
\(\alpha\to\infty\) 时极限给出本质上确界散度 \(\log\max\{r_0,r_1\}=\log r_0\)；
\(\alpha\to 1\) 的连续性给出 KL 极限并与推论 \ref{cor:xi-foldbin-shannon-deficit-constant-closed-form} 一致。证毕。
\end{proof}

\begin{corollary}[bin-fold 熵率的刚性：每步熵趋于 \texorpdfstring{$\log\varphi$}{log phi}，亏损仅为常数级且指数收敛]\label{cor:xi-foldbin-entropy-rate-rigidity}
在定理 \ref{thm:xi-foldbin-pushforward-uniform-two-atom-ledger} 的设定下，
\[
\frac{1}{m}H\bigl(p^{\mathrm{bin}}_m\bigr)\ \longrightarrow\ \log\varphi,
\qquad
\log|X_m|-H\bigl(p^{\mathrm{bin}}_m\bigr)=C_{\mathrm{KL}}+O\!\left(\Bigl(\frac{\varphi}{2}\Bigr)^m\right).
\]
\end{corollary}
\begin{proof}
由推论 \ref{cor:xi-foldbin-shannon-deficit-constant-closed-form}，
\(
H(p^{\mathrm{bin}}_m)=\log|X_m|-D_{\mathrm{KL}}(p^{\mathrm{bin}}_m\Vert u_m)
\)。
再用 \(\log|X_m|=\log F_{m+2}=m\log\varphi+O(1)\) 与 \(D_{\mathrm{KL}}(p^{\mathrm{bin}}_m\Vert u_m)=C_{\mathrm{KL}}+O((\varphi/2)^m)\)，即可得到两条结论。证毕。
\end{proof}

\paragraph{稳定层均匀基线下的退化度极限与信息谱单点律}
除推前分布 \(p_m^{\mathrm{bin}}=\mu_m\) 自身的两点极限外，稳定层均匀基线与微观推前基线还分别诱导出两条更强的统计刚性：前者把尺度化退化度压缩为另一组二原子极限，后者则把全部主阶信息指数压缩到单一点。

\begin{theorem}[稳定层均匀基线下 bin-fold 退化度的二原子极限律]\label{thm:xi-foldbin-uniform-layer-degeneracy-two-atom-law}
\leanverified{paper\_xi\_foldbin\_uniform\_layer\_degeneracy\_two\_atom\_law}
令
\[
u_m(w):=\frac{1}{|X_m|}\qquad (w\in X_m),
\]
并取 \(\widetilde W_m\sim u_m\)。
定义尺度化退化度
\[
Y_m:=\left(\frac{\varphi}{2}\right)^m d_m^{\mathrm{bin}}(\widetilde W_m).
\]
则有弱收敛
\[
Y_m\Longrightarrow \nu
:=
\frac{1}{\varphi}\delta_{1}
+\frac{1}{\varphi^2}\delta_{\varphi^{-1}}.
\]
等价地，对任意有界连续函数 \(f\)，成立
\[
\lim_{m\to\infty}\mathbb E\bigl[f(Y_m)\bigr]
=
\frac{1}{\varphi}f(1)+\frac{1}{\varphi^2}f(\varphi^{-1}).
\]
特别地，对任意 \(a\ge 0\)，有矩极限
\[
\lim_{m\to\infty}\mathbb E[Y_m^a]
=
\frac{1}{\varphi}+\varphi^{-a-2}.
\]
\end{theorem}
\begin{proof}
由定理 \ref{thm:fold-bin-two-state-asymptotic} 的一致渐近，
\[
\left(\frac{\varphi}{2}\right)^m d_m^{\mathrm{bin}}(w)
=
\varphi^{-w_m}+O\!\left(\left(\frac{\varphi}{2}\right)^m\right)
\qquad (w\in X_m),
\]
且误差在 \(w\) 上一致。故存在 \(\varepsilon_m\to 0\) 使
\[
Y_m=\varphi^{-\widetilde W_m(m)}+\varepsilon_m,
\qquad
|\varepsilon_m|\le C\left(\frac{\varphi}{2}\right)^m.
\]
另一方面，
\[
u_m(\widetilde W_m(m)=0)=\frac{F_{m+1}}{F_{m+2}}\to\frac{1}{\varphi},
\qquad
u_m(\widetilde W_m(m)=1)=\frac{F_m}{F_{m+2}}\to\frac{1}{\varphi^2}.
\]
因此对任意有界连续函数 \(f\)，由一致连续性可得
\[
\mathbb E[f(Y_m)]
=
u_m(\widetilde W_m(m)=0)\,f(1)
+u_m(\widetilde W_m(m)=1)\,f(\varphi^{-1})
+o(1),
\]
从而得到所述极限。取 \(f(y)=y^a\) 即得矩公式。证毕。
\end{proof}

\begin{theorem}[bin-fold 统一层分布的精确两点指数族]\label{thm:xi-foldbin-uniform-layer-two-point-exponential-family}
沿用定理 \ref{thm:xi-foldbin-uniform-layer-degeneracy-two-atom-law} 的记号。对每个 \(q\ge 0\)，定义统一层矩函数
\[
Z_m(q):=\EE_{u_m}[Y_m^q]
=
\frac{1}{|X_m|}\sum_{w\in X_m}
\left(\frac{\varphi}{2}\right)^{mq}
\bigl(d_m^{\mathrm{bin}}(w)\bigr)^q.
\]
则极限
\[
\boxed{\;
Z(q):=\lim_{m\to\infty}Z_m(q)
=
\frac{1}{\varphi}+\varphi^{-q-2}\;}
\]
对一切 \(q\ge 0\) 都存在。若记
\[
\Lambda(q):=\log Z(q),
\qquad
\theta(q):=\frac{1}{1+\varphi^{q+1}},
\]
则成立：
\begin{enumerate}
  \item \(\Lambda\) 在 \([0,\infty)\) 上解析且严格凸；
  \item
  \[
  \boxed{\;
  -\Lambda'(q)=\theta(q)\log\varphi;\;}
  \]
  \item
  \[
  \boxed{\;
  \Lambda''(q)
  =
  (\log\varphi)^2\,\theta(q)\bigl(1-\theta(q)\bigr),\;}
  \]
  \[
  \boxed{\;
  \Lambda^{(3)}(q)
  =
  -(\log\varphi)^3\,\theta(q)\bigl(1-\theta(q)\bigr)\bigl(1-2\theta(q)\bigr).\;}
  \]
\end{enumerate}
\end{theorem}
\begin{proof}
由定理 \ref{thm:fold-bin-two-state-asymptotic} 的一致展开，
\[
Y_m(w)=\varphi^{-w_m}+O\!\left(\left(\frac{\varphi}{2}\right)^m\right)
\qquad (w\in X_m),
\]
故对每个固定 \(q\ge 0\)，有一致近似
\[
Y_m(w)^q=\varphi^{-q w_m}+O\!\left(\left(\frac{\varphi}{2}\right)^m\right).
\]
将稳定类型按末位拆分，得到
\[
Z_m(q)
=
\frac{F_{m+1}}{F_{m+2}}\cdot 1^q
+
\frac{F_m}{F_{m+2}}\cdot \varphi^{-q}
+
O\!\left(\left(\frac{\varphi}{2}\right)^m\right).
\]
再用
\[
\frac{F_{m+1}}{F_{m+2}}\to \frac{1}{\varphi},
\qquad
\frac{F_m}{F_{m+2}}\to \frac{1}{\varphi^2},
\]
便得到
\[
Z(q)=\frac{1}{\varphi}+\varphi^{-q-2}.
\]

于是
\[
\Lambda(q)=\log\!\left(\frac{1}{\varphi}+\varphi^{-q-2}\right).
\]
直接微分可得
\[
\Lambda'(q)
=
\frac{-\varphi^{-q-2}\log\varphi}{\varphi^{-1}+\varphi^{-q-2}}
=
-\frac{\log\varphi}{1+\varphi^{q+1}}
=
-\theta(q)\log\varphi,
\]
这就给出第二条。

再由
\[
\theta'(q)=-(\log\varphi)\,\theta(q)\bigl(1-\theta(q)\bigr),
\]
可得
\[
\Lambda''(q)=-(\log\varphi)\theta'(q)
=
(\log\varphi)^2\,\theta(q)\bigl(1-\theta(q)\bigr).
\]
继续微分即得
\[
\Lambda^{(3)}(q)
=
(\log\varphi)^2\,\theta'(q)\bigl(1-2\theta(q)\bigr)
=
-(\log\varphi)^3\,\theta(q)\bigl(1-\theta(q)\bigr)\bigl(1-2\theta(q)\bigr).
\]
由于 \(\theta(q)\in(0,1)\)，故 \(\Lambda''(q)>0\)，从而 \(\Lambda\) 严格凸。又 \(q\mapsto \varphi^{-q-2}\) 为整函数，且
\[
\frac{1}{\varphi}+\varphi^{-q-2}>0
\qquad(q\ge 0),
\]
故 \(\Lambda\) 在 \([0,\infty)\) 上解析。证毕。
\end{proof}

\begin{theorem}[bin-fold 信息谱的单点坍塌]\label{thm:xi-foldbin-information-spectrum-single-point-collapse}
\leanverified{paper\_xi\_foldbin\_information\_spectrum\_single\_point\_collapse}
沿用
\[
\mu_m(w):=\frac{d_m^{\mathrm{bin}}(w)}{2^m},
\qquad
w\in X_m,
\]
并令 \(W_m\sim\mu_m\)。
定义随机信息密度
\[
I_m(W_m):=-\frac{1}{m}\log \mu_m(W_m).
\]
则有
\[
I_m(W_m)=\log\varphi+O(m^{-1}),
\]
其中误差对 \(W_m\) 一致成立；特别地，
\[
I_m(W_m)\longrightarrow \log\varphi.
\]
此外，对一切 \(q>0\) 都有
\[
\frac{1}{m}H_q(\mu_m)\longrightarrow \log\varphi.
\]
因此，bin-fold 的主阶信息谱塌缩为单一点 \(\log\varphi\)，不存在非退化的多重信息指数展开。
\end{theorem}
\begin{proof}
由定义，
\[
I_m(W_m)
=
-\frac{1}{m}\log\frac{d_m^{\mathrm{bin}}(W_m)}{2^m}
=
\log 2-\frac{1}{m}\log d_m^{\mathrm{bin}}(W_m).
\]
而定理 \ref{thm:fold-bin-renyi-rate-collapse} 已给出
\[
\gamma_m:=\frac{1}{m}\log d_m^{\mathrm{bin}}(W_m)
=
\log\!\left(\frac{2}{\varphi}\right)+O(m^{-1}).
\]
故
\[
I_m(W_m)
=
\log 2-\log\!\left(\frac{2}{\varphi}\right)+O(m^{-1})
=
\log\varphi+O(m^{-1}),
\]
从而 \(I_m(W_m)\to\log\varphi\)。

同一定理还表明对全部 \(q>0\) 有
\(
H_q(\mu_m)/m\to\log\varphi
\)。
因此无论从点态信息密度还是从全体 Rényi 熵率观察，主阶信息指数都只能取同一常数 \(\log\varphi\)，故不存在非平凡的多重分形型信息谱。证毕。
\end{proof}


\begin{proposition}[可逆重建成功率曲线的精确折线律：离散导数等于纤维尾分布]\label{prop:xi-reconstruction-success-curve-discrete-derivative}
\leanverified{paper\_xi\_reconstruction\_success\_curve\_discrete\_derivative}
固定 \(m\)。
对任意寄存器大小 \(M\in\mathbb N\) 定义最优平均成功率
\[
P_m^\star(M):=\sup\mathbb P[\widehat\omega=U_m],
\]
其中上确界取遍允许使用 \(M\) 值附加寄存器的全部编码/解码方案（同定理 \ref{thm:xi-optimal-reconstruction-success-probability-constant-bias}，但此处把 \(M\) 视为自变量）。
则有闭式
\[
P_m^\star(M)=\frac{1}{2^m}\sum_{x\in X_m}\min\!\bigl\{d^{\mathrm{bin}}_m(x),\,M\bigr\}.
\]
因此 \(M\mapsto P_m^\star(M)\) 为离散凹折线，并且其离散导数满足
\[
P_m^\star(M+1)-P_m^\star(M)=\frac{1}{2^m}\#\{x\in X_m:\ d^{\mathrm{bin}}_m(x)\ge M+1\}.
\]
\end{proposition}

\begin{proof}
闭式即定理 \ref{thm:xi-optimal-reconstruction-success-probability-constant-bias} 的第一段（把 \(M_m\) 视为任意整数参数 \(M\)）。
差分恒等式由
\(
\min\{d,M+1\}-\min\{d,M\}=\ind_{\{d\ge M+1\}}
\)
逐项求和得到。
导数项随 \(M\) 单调不增，故 \(P_m^\star\) 离散凹。证毕。
\end{proof}

\begin{theorem}[强反定理与精确指数：低于临界速率时成功率指数塌缩]\label{thm:xi-strong-converse-exponent-below-critical-rate}
记附加寄存器速率 \(R_m:=\frac{1}{m}\log M_m\)。
若存在 \(\delta>0\) 使对充分大 \(m\) 有
\[
R_m\le \log\!\Bigl(\frac{2}{\varphi}\Bigr)-\delta,
\]
则任意方案的平均成功率满足指数上界
\[
P_m^\star(M_m)\ \le\ \frac{\varphi^2}{\sqrt5}\,e^{-\delta m}\,(1+o(1)).
\]
反之，若存在 \(\delta>0\) 使 \(R_m\ge \log(2/\varphi)+\delta\) 且 \(M_m\ge d^{\mathrm{bin}}_{\max}(m)\)，则对充分大 \(m\) 存在方案使错误率严格为 \(0\)。
因此可逆重建在指数尺度上表现为硬阈值，相变斜率恰为 \(\delta\)。
\end{theorem}

\begin{proof}
第一部分由定理 \ref{thm:xi-optimal-reconstruction-success-probability-constant-bias} 的第三段（\(c\to-\infty\)）取 \(c=-\delta m+o(m)\) 得到。
第二部分由“在每条纤维内选取单射编码当且仅当 \(M_m\ge d^{\mathrm{bin}}_{\max}(m)\)”的充要条件直接推出。
\end{proof}

\begin{proposition}[Parry 符号链的二点周边谱：相关衰减率 \texorpdfstring{$\varphi^{-2}$}{phi^{-2}} 的刚性]\label{prop:xi-parry-symbol-markov-peripheral-spectrum}
\leanverified{paper\_xi\_parry\_symbol\_markov\_peripheral\_spectrum}
在稳定边界 \(X_\infty=\{c\in\{0,1\}^{\mathbb N}:c_kc_{k+1}=0\}\) 上取 Parry（最大熵）平衡测度。
其首符号过程 \(C_0,C_1,\dots\in\{0,1\}\) 为二态马尔可夫链，转移矩阵为
\[
P_\varphi=
\begin{pmatrix}
\varphi^{-1} & \varphi^{-2}\\[2pt]
1 & 0
\end{pmatrix}.
\]
该矩阵具有谱 \(\{1,-\varphi^{-2}\}\)。
因此任意仅依赖有限前缀的可观测量，其相关函数的主衰减率被强迫为 \((\varphi^{-2})^{n}\)，并伴随符号交替（由负特征值导致）。
\end{proposition}

\begin{theorem}[KMS 传递算子在 \texorpdfstring{$L^2(\omega)$}{L2(omega)} 上的谱隙：非交换混合的显式常数]\label{thm:xi-kms-transfer-operator-spectral-gap-explicit}
在定理 \ref{thm:xi-kms-code12-discrete-modular-scale-group} 的设定下，定义 KMS-保态的完全正算子
\[
\Phi_{\varphi}(a):=\varphi^{-1}\mathsf{S}_0^\ast a\mathsf{S}_0+\varphi^{-2}\mathsf{S}_1^\ast a\mathsf{S}_1.
\]
则 \(\omega\circ\Phi_{\varphi}=\omega\)，且 \(\Phi_{\varphi}\) 在 \(L^2(\mathcal{O}_2,\omega)\) 上为自伴收缩。
其在对角子代数上的限制与命题 \ref{prop:xi-parry-symbol-markov-peripheral-spectrum} 的符号马尔可夫算子共轭，因而周边谱包含 \(\{1,-\varphi^{-2}\}\)。
并且在 \(\omega\)-中心化子正交补上存在统一谱隙：对任意 \(a\in\mathcal{O}_2\) 存在常数 \(C(a)>0\) 使
\[
\|\Phi_{\varphi}^n(a)-\omega(a)1\|_{2,\omega}\ \le\ C(a)\,\varphi^{-2n}\,\|a\|_{2,\omega},
\qquad n\ge 0.
\]
因此非交换动力学的混合速率被 \(\varphi^{-2}\) 唯一量化。
\end{theorem}

\begin{proposition}[有限分辨率谱证书的一致性：\texorpdfstring{$\Phi_\varphi$}{Phi}-有限截断谱的指数收敛]\label{prop:xi-finite-truncation-spectrum-converges}
令 \(\mathcal D_m\subset C(X_\infty)\) 为长度 \(m\) 的圆柱代数，并记 \(\Phi_{\varphi,m}\) 为 \(\Phi_\varphi\) 在 \(L^2(\mathcal D_m,\omega)\) 上的压缩。
则存在与 \(m\) 无关常数 \(C'>0\) 使得
\[
\mathrm{spec}(\Phi_{\varphi,m})\subset\bigl[-\varphi^{-2}-C'\varphi^{-2m},\,1\bigr],
\]
且第二特征值满足
\[
\lambda_2(\Phi_{\varphi,m})=-\varphi^{-2}+O(\varphi^{-2m}).
\]
因此 \(\varphi^{-2}\) 的混合常数可由有限 \(m\) 的谱计算以指数精度审计，给出完全有限维的可复现验证口径。
\end{proposition}

\begin{theorem}[尺度常数的刚性刻画：三条结构公理唯一确定 \texorpdfstring{$\varphi$}{phi}]\label{thm:xi-phi-rigidity-from-kms-entropy-density}
设存在一对前缀等距生成元 \((\mathsf{S}_0,\mathsf{S}_1)\) 与码长 \((\ell_0,\ell_1)=(1,2)\)，并满足：
\begin{enumerate}
  \item 存在 KMS 态且唯一；
  \item 相对熵密度
  \(\displaystyle \lim_{m\to\infty}\frac{1}{m}D_{\mathrm{KL}}\bigl(\mathrm{Unif}(\Omega_m)\,\Vert\,\nu_m\bigr)\)
  存在（其中 \(\nu_m\) 为 KMS 诱导的 \(m\)-圆柱边缘）；
  \item 模谱为离散乘法群 \(\lambda^{\mathbb Z}\)。
\end{enumerate}
则三者联合强迫
\[
\lambda=\frac{1}{\varphi},
\qquad
\beta=\log\varphi,
\]
并且相对熵密度必等于
\[
\log\!\Bigl(\frac{2}{\varphi}\Bigr)=\log 2-\beta.
\]
因此黄金比常数可由纯算子动力学（KMS 唯一性与离散模谱）与信息密度（相对熵率）三条结构条件唯一刻画，而不需先验引入 Fibonacci 组合学。
\end{theorem}

\begin{theorem}[可见--微观二层语义的同一化：重建阈值＝模尺度缺口＝最大纤维指数]\label{thm:xi-threshold-modular-gap-maxfiber-equivalence}
在 bin-fold 与定理 \ref{thm:xi-kms-code12-discrete-modular-scale-group} 的共同制度下，以下三项给出同一常数：
\[
\text{\upshape(A) 重建临界附加信息率}\quad
R_\star=\log\!\Bigl(\frac{2}{\varphi}\Bigr),
\qquad
\text{\upshape(B) 微观熵率与模尺度之差}\quad
\log 2-\beta,
\qquad
\text{\upshape(C) 最大纤维指数斜率}\quad
\lim_{m\to\infty}\frac{1}{m}\log d_{\max}^{\mathrm{bin}}(m).
\]
其中 \(\beta=\log\varphi\)，且 \(d_{\max}^{\mathrm{bin}}(m):=\max_{x\in X_m}d_m^{\mathrm{bin}}(x)\) 满足
\(\lim_{m\to\infty}m^{-1}\log d_{\max}^{\mathrm{bin}}(m)=\log(2/\varphi)\)（定理 \ref{thm:fold-bin-max-fiber-exponent}）。
因此，“折叠信息丢失的指数率”“可逆化所需的临界附加信息率”与“模尺度相对于微观熵率的缺口”被迫塌缩为同一数值不变量 \(\log(2/\varphi)\)，形成可对外引用的三域等价链：信息论阈值＝模论熵缺口＝组合最大退化指数。
\end{theorem}

\paragraph{Zeckendorf 前缀截断的边界细化：\texorpdfstring{$A_4$}{A4} 型 Bratteli 图、\texorpdfstring{$\ZZ[\varphi]$}{Z[phi]} 维数群与 Fibonacci 子因子}
令 \(X_m:=\{x\in\{0,1\}^m:\ \forall i,\ x_ix_{i+1}=0\}\) 为禁止相邻 \(11\) 的长度 \(m\) 稳定词集，并按末位划分
\[
X_m^{(0)}:=\{x\in X_m:\ x_m=0\},
\qquad
X_m^{(1)}:=\{x\in X_m:\ x_m=1\}.
\]
令 \(\pi_{m+1\to m}:X_{m+1}\to X_m\) 为前缀截断（丢弃最末一位）。

\begin{theorem}[边界细化的稳定关联矩阵与 \texorpdfstring{$A_4$}{A4} 型二部图]\label{thm:xi-zeckendorf-boundary-refinement-a4-bratteli}
\leanverified{paper\_xi\_zeckendorf\_boundary\_refinement\_a4\_bratteli}
在上述记号下，对任意 \(m\ge 1\) 与任意 \(x\in X_m\) 有
\[
\pi_{m+1\to m}^{-1}(x)=
\begin{cases}
\{x0,\ x1\}, & x\in X_m^{(0)},\\
\{x0\}, & x\in X_m^{(1)}.
\end{cases}
\]
因此由末位边界类型 \(\{X_m^{(0)},X_m^{(1)}\}\) 所诱导的边界细化在每一层都由同一关联矩阵
\[
K=\begin{pmatrix}1&1\\[2pt]1&0\end{pmatrix}
\]
描述；并且在把相邻两层的两类边界类型视为四个顶点的二部图时，其边连接恰形成四点链（与 Dynkin 图 \(A_4\) 同构）。
\end{theorem}

\begin{proof}
追加末位 \(0\) 总保持 \(11\)-禁词约束；追加末位 \(1\) 合法当且仅当当前末位为 \(0\)。
截断映射 \(\pi_{m+1\to m}\) 正是去除末位，故其纤维与上述合法延拓集合一致，得到纤维基数的两态律。
以 \(K_{bb'}\) 记从类型 \(b\) 延拓到类型 \(b'\) 的条数，直接得到 \(K_{00}=K_{01}=K_{10}=1\)、\(K_{11}=0\)，从而矩阵为所示 \(K\)。
二部图顶点为 \(X_m^{(0)},X_m^{(1)}\) 与 \(X_{m+1}^{(0)},X_{m+1}^{(1)}\)，边仅有三条且无重边，因而为四点链。证毕。
\end{proof}

\begin{theorem}[平稳 AF 代数的有序 \texorpdfstring{$K_0$}{K0} 与 \texorpdfstring{$\varphi$}{phi} 乘法细化]\label{thm:xi-zeckendorf-a4-af-k0-zphi}
\leanverified{paper\_xi\_zeckendorf\_a4\_af\_k0\_zphi}
令 \(\mathfrak{A}_K\) 为由关联矩阵 \(K\) 定义的平稳 AF 代数。
则
\[
K_1(\mathfrak{A}_K)=0,
\qquad
K_0(\mathfrak{A}_K)\cong \ZZ[\varphi]
\]
作为有序群同构，其中正锥由 \(\ZZ[\varphi]\subset\RR\) 的通常序诱导。
在该识别下，一级细化在 \(K_0\) 上的诱导自同态为
\[
x\longmapsto \varphi x.
\]
\end{theorem}

\begin{proof}
AF 代数满足 \(K_1=0\)，且 \(K_0\) 为维数群直极限
\(
K_0(\mathfrak{A}_K)\cong \varinjlim(\ZZ^2,K)
\)。
定义群同态
\[
\Psi:\ZZ^2\to \ZZ[\varphi],\qquad \Psi(a,b)=a\varphi+b.
\]
由 \(\varphi^2=\varphi+1\) 直接计算
\[
\Psi(K(a,b))=\Psi(a+b,a)=(a+b)\varphi+a=\varphi(a\varphi+b)=\varphi\,\Psi(a,b),
\]
故结构映射在 \(\Psi\) 下对应为乘以 \(\varphi\)，从而直极限与 \(\ZZ[\varphi]\) 同构，且细化自同态即乘 \(\varphi\)。
正锥由 Perron--Frobenius 正特征向量诱导并与 \(\ZZ[\varphi]\subset\RR\) 的通常非负锥一致。证毕。
\end{proof}

\begin{corollary}[唯一归一化迹的值域与 \texorpdfstring{$[0,1]$}{[0,1]} 内稠密性]\label{cor:xi-zeckendorf-a4-af-trace-range-dense}
令 \(\tau\) 为 \(\mathfrak{A}_K\) 的唯一归一化迹。
在定理 \ref{thm:xi-zeckendorf-a4-af-k0-zphi} 的识别下，有
\[
\tau\bigl(\mathrm{Proj}(\mathfrak{A}_K)\bigr)=\ZZ[\varphi]\cap[0,1],
\]
并且 \(\ZZ[\varphi]\cap[0,1]\) 在 \([0,1]\) 内稠密。
\end{corollary}

\begin{proof}
平稳原始 Bratteli 系统的迹态唯一性由 Perron--Frobenius 理论给出；迹在 \(K_0\) 上对应于嵌入 \(\ZZ[\varphi]\hookrightarrow\RR\) 的正同态并满足 \(\tau(1)=1\)。
因此投影迹值恰为 \(\ZZ[\varphi]\cap[0,1]\)。
稠密性来自 \(\ZZ[\varphi]=\{a+b\varphi:\ a,b\in\ZZ\}\) 为秩 \(2\) 的无理生成加法子群：集合 \(\{b\varphi\bmod 1:\ b\in\ZZ\}\) 在 \(\RR/\ZZ\) 中稠密，从而可在任意小区间内取到某个 \(a+b\varphi\)。证毕。
\end{proof}

\begin{theorem}[\texorpdfstring{$A_4$}{A4} 主图子因子与 Fibonacci 融合律]\label{thm:xi-zeckendorf-a4-subfactor-fib-fusion}
存在超有限 \(\mathrm{II}_1\) 子因子 \(N\subset M\)，其主图为 \(A_4\)，且 Jones 指数为
\[
[M:N]=\varphi^2.
\]
其偶部融合范畴的半单完备化仅含两类简单对象 \(\ind\) 与 \(\tau\)，并满足
\[
\tau\otimes\tau\ \cong\ \ind\ \oplus\ \tau,
\qquad
d(\tau)=\varphi.
\]
\end{theorem}

\begin{proof}
\(A_4\) 的邻接算子谱范数为 \(2\cos(\pi/5)=\varphi\)，故指数应为其平方 \(\varphi^2\)。
Temperley--Lieb--Jones 的 \(A_n\) 标准不变量给出以 \(A_4\) 为主图的超有限子因子实现，从而得到所示指数。
偶部简单对象与偶深度顶点对应，\(A_4\) 的偶深度顶点恰为两点；融合由路径拼接给出，直接两步计数得到 \(\tau\otimes\tau\simeq \ind\oplus\tau\)，其量子维数为融合矩阵 Perron 根，满足 \(d(\tau)^2=d(\tau)+1\)，取正根得 \(d(\tau)=\varphi\)。证毕。
\end{proof}


\paragraph{window-\texorpdfstring{$6$}{6} 的内禀自校准与边界证书：末位均值比率、四块谱数域、返时母函数与半群阈值}
本段在不引入外部标定的前提下，把若干与 window-\(6\) 锚点直接相连、且可由有限数据审计的结论压缩为一条可独立引用的链条：其一是 bin-fold 两态一致渐近所诱导的 \(\varphi\) 自校准统计量；其二是 boundary uplift 的四步平移同构所给出的家族数窗口量子化与最小窗唯一性；其三是四块粗粒化核的不可约三次谱数域与 \(S_3\) 伽罗瓦群；其四是边界块 \(U_R\) 的首返时间分布闭式与通量倒数律；其五是由 tail-offset 三元组 \(\{21,34,55\}\) 所生成数值半群的 Frobenius 阈值与由此导出的“维数筛法有限性”。
\par

\begin{theorem}[黄金比的内禀自校准：由末位条件均值比率恢复 \texorpdfstring{$\varphi$}{phi}]\label{thm:xi-foldbin-intrinsic-self-calibration-phi-lastbit-means}
\leanverified{paper\_xi\_foldbin\_intrinsic\_self\_calibration\_phi\_lastbit\_means}
在 bin-fold 两态一致渐近（定理 \ref{thm:fold-bin-two-state-asymptotic}）成立的口径下，
令 \(d_m^{\mathrm{bin}}(x):=\bigl|(\mathrm{Fold}^{\mathrm{bin}}_m)^{-1}(x)\bigr|\)（\(x\in X_m\)），并按末位划分
\[
X_m^{(0)}:=\{x\in X_m:\ x_m=0\},
\qquad
X_m^{(1)}:=\{x\in X_m:\ x_m=1\}.
\]
定义可观测统计量
\[
\widehat\varphi_m
\ :=\
\frac{\frac{1}{|X_m^{(0)}|}\sum_{x\in X_m^{(0)}} d_m^{\mathrm{bin}}(x)}{\frac{1}{|X_m^{(1)}|}\sum_{x\in X_m^{(1)}} d_m^{\mathrm{bin}}(x)}.
\]
则存在与 \(m\) 无关常数 \(C>0\) 使得
\[
\bigl|\widehat\varphi_m-\varphi\bigr|\ \le\ C\left(\frac{\varphi}{2}\right)^m,
\]
从而 \(\widehat\varphi_m\) 以指数精度收敛到 \(\varphi\)。
该结论给出一条无需外部标定的自校准机制：仅凭折叠纤维多重度与末位可见位即可恢复增长常数。
\end{theorem}

\begin{proof}
由定理 \ref{thm:fold-bin-two-state-asymptotic}，存在常数 \(C_0>0\) 使对充分大 \(m\) 与任意 \(x\in X_m\) 有一致估计
\[
\Bigl|d_m^{\mathrm{bin}}(x)-\varphi^{-x_m}\Bigl(\frac{2}{\varphi}\Bigr)^m\Bigr|\le C_0.
\]
对 \(X_m^{(0)}\)、\(X_m^{(1)}\) 分别取均值，得
\[
\frac{1}{|X_m^{(0)}|}\sum_{x\in X_m^{(0)}} d_m^{\mathrm{bin}}(x)
=\Bigl(\frac{2}{\varphi}\Bigr)^m+O(1),
\qquad
\frac{1}{|X_m^{(1)}|}\sum_{x\in X_m^{(1)}} d_m^{\mathrm{bin}}(x)
=\varphi^{-1}\Bigl(\frac{2}{\varphi}\Bigr)^m+O(1),
\]
其中 \(O(1)\) 的常数与 \(m\) 无关。
两式相除并用 \((2/\varphi)^{-m}=(\varphi/2)^m\) 得
\[
\widehat\varphi_m=\varphi\cdot\frac{1+O((\varphi/2)^m)}{1+O((\varphi/2)^m)}
=\varphi+O\!\left(\left(\frac{\varphi}{2}\right)^m\right),
\]
从而存在常数 \(C\) 满足所示误差界。证毕。
\end{proof}

\begin{theorem}[家族数的窗口量子化与最小窗唯一性：\texorpdfstring{$N_f=|X_m^{\mathrm{bdry}}|$}{Nf=|Xbdry|}]\label{thm:xi-family-number-window-quantization-minimal-window}
\leanverified{paper\_xi\_family\_number\_window\_quantization\_minimal\_window}
固定每代 \(SU(5)\) 手征计数 \(15\) 为外部输入，并要求 “\(SO(10)\)-型最小补全” 在同一窗口 \(m\) 上由边界 singlet 通道完成，即满足离散恒等式
\[
15N_f+|X_m^{\mathrm{bdry}}|=16N_f.
\]
则必有
\[
N_f=|X_m^{\mathrm{bdry}}|=F_{m-2},
\]
从而家族数在该制度下只能取 Fibonacci 值。
特别地，若要求 \(N_f=3\)，则唯一解为 \(m=6\)（最小窗唯一性），并与 window-\(6\) uplift 三分支中 \(F_9=34\) 分支的离散对齐一致（表 \ref{tab:window6_family_uplift_lock} 与 \eqref{eq:window6_family_singlet_completion}）。
\end{theorem}

\begin{proof}
由所给离散恒等式立得 \(N_f=|X_m^{\mathrm{bdry}}|\)。
另一方面，由定理 \ref{thm:boundary-shift4-uplift-isomorphism} 的 uplift 同构 \(X_{m-4}\simeq X_m^{\mathrm{bdry}}\)（\(m\ge 6\)）与 \(|X_n|=F_{n+2}\) 得
\(|X_m^{\mathrm{bdry}}|=|X_{m-4}|=F_{m-2}\)，从而得到 \(N_f=F_{m-2}\)。
若 \(N_f=3\)，则 \(F_{m-2}=3\) 强制 \(m-2=4\)，即 \(m=6\)，且该解唯一。证毕。
\end{proof}

\begin{theorem}[window-\(6\) 四块粗粒化核的谱数域刚性：不可约三次域与 \texorpdfstring{$S_3$}{S3} 伽罗瓦群]\label{thm:xi-window6-four-block-kernel-cubic-field-s3}
\leanverified{paper\_xi\_window6\_four\_block\_kernel\_cubic\_field\_s3}
令四块骨架的粗粒化 Markov 核 \(T\) 如 \eqref{eq:window6_edge_flux_skeleton}，
并记其特征多项式分解为
\[
\chi_T(t)=(t-1)\cdot p(t),
\qquad
p(t)=\frac{48114t^3+7263t^2-506t-55}{48114}.
\]
则：
\begin{enumerate}
  \item 三次因子 \(p(t)\) 在 \(\mathbb Q[t]\) 上不可约；
  \item 其判别式 \(\Delta=\mathrm{Disc}(48114t^3+7263t^2-506t-55)\) 非平方（见 \eqref{eq:window6_edge_flux_skeleton_arithmetic}）；
  \item 因而分裂域的伽罗瓦群为 \(S_3\)，并且 \(T\) 的全部非平凡本征值均为代数次数 \(3\) 的实数。
\end{enumerate}
因此非平凡谱生成一个不可交换的三次数域 \(K_T=\mathbb Q(\lambda)\)，并且“粗粒化混合”所需的最小代数复杂度在数域意义下为三次且不可降阶。
\end{theorem}

\begin{proof}
对三次多项式，\(\mathbb Q[t]\) 上不可约当且仅当无有理根。
把整系数多项式
\(
f(t):=48114t^3+7263t^2-506t-55
\)
模 \(13\) 化简，得
\(
f(t)\equiv t^3+9t^2+t+10\pmod{13}
\)。
直接检验 \(t\in\mathbb F_{13}\) 的全部取值可知该三次多项式在 \(\mathbb F_{13}\) 上无根，故在 \(\mathbb F_{13}[t]\) 上不可约，从而 \(f\) 在 \(\mathbb Q[t]\) 上不可约，\(p\) 亦不可约。
\par
由 \eqref{eq:window6_edge_flux_skeleton_arithmetic}，判别式 \(\Delta\) 非平方。
对不可约三次多项式，其伽罗瓦群为 \(A_3\) 当且仅当判别式为平方；否则为 \(S_3\)。
因此本例伽罗瓦群为 \(S_3\)。
又因 \(\Delta>0\)，三根互异且均为实数，从而 \(T\) 的三个非平凡本征值均为代数次数 \(3\) 的实数。证毕。
\end{proof}

\begin{corollary}[判别式分解的审计素数：\texorpdfstring{$23$}{23} 的不可消去出现]\label{cor:xi-window6-four-block-kernel-disc-audit-prime-23}
沿定理 \ref{thm:xi-window6-four-block-kernel-cubic-field-s3} 的整系数多项式
\(
f(t)=48114t^3+7263t^2-506t-55
\)。
其判别式满足
\[
\boxed{\;
\Disc(f)
=2^{2}\cdot 3^{2}\cdot 11^{2}\cdot 13^{2}\cdot 23\;}
\]
并且 \(\det(tI-T)\) 的常数项分解含因子
\(
23\cdot 11\cdot 5
\)
（见证书 \eqref{eq:window6_edge_flux_skeleton_arithmetic}）。
因此素数 \(23\) 作为四块粗粒化核的代数数论指纹在该口径下不可被规避或坐标变换消去：其出现由判别式与常数项同时锁定。
\end{corollary}

\begin{corollary}[粗粒化谱数域与黄金尺度域的互素分离：合成域次数为 \(6\)]\label{cor:xi-window6-field-disjointness-degree-six}
令 \(K_T=\mathbb Q(\lambda)\) 为定理 \ref{thm:xi-window6-four-block-kernel-cubic-field-s3} 的非平凡谱所生成三次数域，
并令 \(K_\varphi=\mathbb Q(\varphi)\) 为黄金二次数域。
则必有
\[
K_T\cap K_\varphi=\mathbb Q,
\]
从而合成域次数满足
\[
[K_TK_\varphi:\mathbb Q]=[K_T:\mathbb Q]\,[K_\varphi:\mathbb Q]=6.
\]
\end{corollary}

\begin{proof}
交域 \(K_T\cap K_\varphi\) 是 \(K_T\) 的子域，其次数必须整除 \([K_T:\mathbb Q]=3\)。
另一方面，\(K_T\cap K_\varphi\subseteq K_\varphi\) 且 \([K_\varphi:\mathbb Q]=2\)，故其次数亦整除 \(2\)。
于是 \([K_T\cap K_\varphi:\mathbb Q]\) 同时整除 \(2\) 与 \(3\)，只能为 \(1\)，即交域为 \(\mathbb Q\)。
线性无关性推出合成域次数为积。证毕。
\end{proof}

\begin{corollary}[混合指数与模尺度的非共振：\texorpdfstring{$-\log\Lambda/\log\varphi\notin\mathbb Q$}{-logLambda/logphi notin Q}]\label{cor:xi-window6-logLambda-logphi-irrational}
令
\[
\Lambda:=\max\{|\lambda|:\ \lambda\in\mathrm{spec}(T)\setminus\{1\}\}\in(0,1)
\]
为四块粗粒化核 \(T\) 的非平凡谱半径（取绝对值的谱半径）。
则
\[
\frac{-\log\Lambda}{\log\varphi}\notin\mathbb Q.
\]
\end{corollary}

\begin{proof}
由定理 \ref{thm:xi-window6-four-block-kernel-cubic-field-s3}，\(T\) 的非平凡本征值为代数次数 \(3\) 的实数；
因此 \(\Lambda\) 等于某个非平凡实本征值 \(\lambda_\star\) 的绝对值，故 \(\Lambda\in K_T\setminus\mathbb Q\)。
若 \(-\log\Lambda/\log\varphi=p/q\in\mathbb Q\)（既约，\(q\ge 1\)），则
\(
\Lambda^q=\varphi^{-p}\in K_\varphi
\)。
但 \(\Lambda^q\in K_T\)，故 \(\Lambda^q\in K_T\cap K_\varphi=\mathbb Q\)（推论 \ref{cor:xi-window6-field-disjointness-degree-six}）。
由于 \(\Lambda\in(0,1)\) 得 \(p\neq 0\)，而 \(\varphi^{-p}\notin\mathbb Q\)，矛盾。证毕。
\end{proof}

\begin{proposition}[Baker 型下界在本体系中的专门化：尺度拍频的有效 Diophantine 约束]\label{prop:xi-window6-baker-effective-beat-bound}
在推论 \ref{cor:xi-window6-logLambda-logphi-irrational} 的设定下，令 \(\alpha_1=\Lambda\)、\(\alpha_2=\varphi\)。
则由对数线性形式下界理论（Baker--W\"ustholz/Matveev 型结果）推出：存在可有效取值的常数 \(C>0\)，使对任意整数对 \((n,k)\neq(0,0)\) 有
\[
\bigl|n\log\Lambda+k\log\varphi\bigr|
\ \ge\
\exp\!\bigl(-C\log(1+\max\{|n|,|k|\})\bigr).
\]
因此任意“混合尺度 \(\Lambda^n\)”与“模尺度 \(\varphi^{-k}\)”之间的近共振误差至多以多对数速度出现，排除了指数级同步。
\end{proposition}

\begin{proof}
由推论 \ref{cor:xi-window6-logLambda-logphi-irrational}，\(\Lambda\) 与 \(\varphi\) 乘法独立。
对数线性形式下界给出对任意非零整数线性组合 \(n\log\Lambda+k\log\varphi\) 的有效下界，其量级为 \((1+\max\{|n|,|k|\})^{-C}\)（等价写作所示指数形式）；详见 \cite{BakerWustholz2007LogarithmicForms,Matveev2000LinearFormsLogs}。证毕。
\end{proof}

\paragraph{四块商图的 Kirchhoff 复杂度与粗粒化 Green 核的分母障碍}
在四块骨架证书 \eqref{eq:window6_edge_flux_skeleton} 的记号下，令 \(\mathcal{G}_{\mathrm{blk}}\) 为加权无向多重图：顶点为 \(\{U_2,U_1,U_L,U_R\}\)，且当 \(\alpha\neq\beta\) 时边重数为 \(E_{\alpha\beta}\)（忽略块内自环）。
记其加权 Laplace 矩阵为 \(L_{\mathrm{blk}}\)（对角为度数、非对角为负边重数），并令 \(L_{\mathrm{blk}}^{(\alpha)}\) 为删去同一行列所得主余子式。

\begin{theorem}[回路格的电阻型 Gram 行列式与生成树权]\label{thm:xi-flow-lattice-gram-determinant-tree-weight}
\leanverified{paper\_xi\_flow\_lattice\_gram\_determinant\_tree\_weight}
令 \(G=(V,E)\) 为有限连通（允许多重边）无向图，并固定一组定向以定义去掉一行的有向关联矩阵
\(
B\in\ZZ^{(|V|-1)\times|E|}
\)
（于是 \(\mathrm{rank}\,B=|V|-1\)）。
给定边权 \(w_e>0\)，令 \(W:=\mathrm{diag}(w_e)_{e\in E}\)。
定义整数流格
\[
\Lambda_{\mathrm{fl}}:=\ker(B)\cap\ZZ^{E}\subset \RR^E,
\]
并在 \(\RR^E\) 上取电阻型二次型
\[
\langle x,y\rangle_{W^{-1}}:=x^{\mathsf T}W^{-1}y.
\]
记 \(\tau_w(G)\) 为加权生成树总权（对每棵生成树取边权乘积并求和）。
则 \(\Lambda_{\mathrm{fl}}\) 在该二次型下的 Gram 行列式满足精确恒等式
\[
\det\!\bigl(\mathrm{Gram}_{W^{-1}}(\Lambda_{\mathrm{fl}})\bigr)=\frac{\tau_w(G)}{\det(W)}.
\]
等价地，其协体积为
\[
\mathrm{covol}_{W^{-1}}(\Lambda_{\mathrm{fl}})=\sqrt{\tau_w(G)/\det(W)}.
\]
\end{theorem}
\begin{proof}
取任一生成树 \(T\subseteq E\)，令弦边集 \(C:=E\setminus T\)，并按 \(E=T\sqcup C\) 重排列边。
相应写 \(B=[B_T\ B_C]\) 与 \(W=\mathrm{diag}(W_T,W_C)\)。
由于 \(T\) 为生成树，\(B_T\in\ZZ^{(|V|-1)\times(|V|-1)}\) 为可逆整矩阵且 \(\det(B_T)=\pm 1\)。
令
\[
C_*:=\begin{pmatrix}-B_T^{-1}B_C\\ I\end{pmatrix}\in\ZZ^{E\times|C|},
\]
则 \(BC_*=0\)，且其列向量给出 \(\Lambda_{\mathrm{fl}}\) 的一组整基。
故 Gram 矩阵为
\[
G_*:=C_*^{\mathsf T}W^{-1}C_*
=W_C^{-1}+B_C^{\mathsf T}B_T^{-\mathsf T}W_T^{-1}B_T^{-1}B_C.
\]
由矩阵行列式引理与 Sylvester 恒等式 \(\det(I+AB)=\det(I+BA)\)，得
\[
\det(G_*)
=\det(W_C^{-1})\cdot
\det\!\Bigl(I+(B_TW_TB_T^{\mathsf T})^{-1}B_CW_CB_C^{\mathsf T}\Bigr).
\]
另一方面，加权约化拉普拉斯为
\[
L_w:=BW B^{\mathsf T}=B_TW_TB_T^{\mathsf T}+B_CW_CB_C^{\mathsf T},
\]
并有分解
\[
\det(L_w)=\det(B_TW_TB_T^{\mathsf T})\cdot
\det\!\Bigl(I+(B_TW_TB_T^{\mathsf T})^{-1}B_CW_CB_C^{\mathsf T}\Bigr).
\]
由于 \(\det(B_TW_TB_T^{\mathsf T})=\det(W_T)\)（因 \(\det(B_T)=\pm1\)），比较两式得到
\[
\det(G_*)
=\det(W_C^{-1})\cdot\frac{\det(L_w)}{\det(W_T)}
=\frac{\det(L_w)}{\det(W_T)\det(W_C)}
=\frac{\det(L_w)}{\det(W)}.
\]
加权矩阵树定理给出 \(\det(L_w)=\tau_w(G)\)，从而结论成立。证毕。
\end{proof}

\begin{corollary}[固定边界偏差后的整数流扭量协体积]\label{cor:xi-affine-flow-torsor-covolume}
\leanverified{paper\_xi\_affine\_flow\_torsor\_covolume}
在定理 \ref{thm:xi-flow-lattice-gram-determinant-tree-weight} 的设定下，给定任意 \(b\in\ZZ^{|V|-1}\) 使得同余方程 \(BF=b\) 在 \(\ZZ^E\) 中可解。
取一解 \(F^{(0)}\in\ZZ^E\)，并令
\[
S(b):=\{F\in\ZZ^E:\ BF=b\}=F^{(0)}+\Lambda_{\mathrm{fl}}.
\]
则 \(S(b)\) 为以 \(\Lambda_{\mathrm{fl}}\) 为平移群的仿射格。
因此在二次型 \(\langle\cdot,\cdot\rangle_{W^{-1}}\) 下，解集的自由度密度由
\[
\mathrm{covol}_{W^{-1}}(\Lambda_{\mathrm{fl}})=\sqrt{\tau_w(G)/\det(W)}
\]
精确控制。
\end{corollary}
\begin{proof}
若 \(F\in S(b)\)，则 \(B(F-F^{(0)})=0\)，即 \(F-F^{(0)}\in\Lambda_{\mathrm{fl}}\)，反之亦然，故 \(S(b)=F^{(0)}+\Lambda_{\mathrm{fl}}\)。
协体积等式由定理 \ref{thm:xi-flow-lattice-gram-determinant-tree-weight} 直接给出。证毕。
\end{proof}

\begin{corollary}[window-\(6\) 四块商图的加权树数分解与审计素数 \texorpdfstring{$p_\star=571$}{p*=571}]\label{cor:xi-window6-block-tree-number-audit-prime}
\leanverified{paper\_xi\_window6\_block\_tree\_number\_audit\_prime}
对任意块 \(\alpha\in\{2,1,L,R\}\)，主余子式的行列式一致且等于同一加权树数
\[
\det\!\bigl(L_{\mathrm{blk}}^{(\alpha)}\bigr)
=\tau(\mathcal{G}_{\mathrm{blk}})
=123336
=2^{3}\cdot 3^{3}\cdot 571.
\]
因此在该四块几何的 Kirchhoff 复杂度中出现唯一素因子 \(p_\star:=571\) 满足 \(\gcd(p_\star,6)=1\)；该素数可作为块级树复杂度意义下的算术指纹。
\end{corollary}

\begin{proof}
矩阵树定理给出 \(\det(L_{\mathrm{blk}}^{(\alpha)})=\tau(\mathcal{G}_{\mathrm{blk}})\) 且与 \(\alpha\) 无关。
数值与分解由证书 \eqref{eq:window6_edge_flux_skeleton_arithmetic} 给出。证毕。
\end{proof}

\begin{proposition}[粗粒化 Green 核的算术闭包与 \texorpdfstring{$p_\star$}{p*} 的最小性]\label{prop:xi-window6-green-kernel-arithmetic-audit-prime}
在四块粗粒化核 \(T\)（\eqref{eq:window6_edge_flux_skeleton}）与平稳分布
\(
\pi_\alpha:=|U_\alpha|/64
\)
下，定义基本 Green 核
\[
Z:=(I-T+\mathbf 1\pi^\top)^{-1}.
\]
则 \(Z\) 的全部矩阵元属于 \(\mathbb{Z}[1/2,1/3,1/p_\star]\)，并且在既约分母中除 \(2,3\) 外不出现其他素因子。
此外，
\[
\det(I-T+\mathbf 1\pi^\top)=\frac{2^{4}\,p_\star}{3^{6}\cdot 11}
\]
（证书 \eqref{eq:window6_green_kernel_audit_prime_571}），从而 \(p_\star\) 在 \(Z\) 中不可被整体消去：存在矩阵元的既约分母含素因子 \(p_\star\)（见表 \ref{tab:window6_green_kernel_Z}）。
\end{proposition}

\begin{proof}
由脚本审计输出 \eqref{eq:window6_green_kernel_audit_prime_571} 与表 \ref{tab:window6_green_kernel_Z}，\(Z\) 的分母谱仅含 \(2,3,p_\star\)。
又由 Cramer 法则，\(Z\) 的任一矩阵元为 \(\mathrm{adj}(I-T+\mathbf 1\pi^\top)\) 的余子式除以 \(\det(I-T+\mathbf 1\pi^\top)\)。
由于该行列式既约分子含 \(p_\star\)，且表 \ref{tab:window6_green_kernel_Z} 中确有既约分母含 \(p_\star\) 的矩阵元，故 \(p_\star\) 不能在 Green 核层面被消去。证毕。
\end{proof}

\begin{theorem}[Markov--导数行列式恒等式：坏素数当且仅当特征多项式在 \(1\) 处模 \(p\) 重根]\label{thm:xi-markov-derivative-determinant-bad-prime}
\leanverified{paper\_xi\_markov\_derivative\_determinant\_bad\_prime}
设 \(T\) 为 \(N\times N\) 行随机矩阵，\(\pi^\top\) 为其平稳分布（\(\pi^\top T=\pi^\top\)，\(\pi^\top\mathbf 1=1\)），并假设特征值 \(1\) 在 \(\mathbb Q\) 上为单根。
定义
\[
A:=I-T+\mathbf 1\pi^\top,\qquad Z:=A^{-1},
\qquad
\chi_T(t):=\det(tI-T).
\]
则：
\begin{enumerate}
  \item 存在严格恒等式
  \[
  \det(A)=\chi_T'(1).
  \]
  \item 若 \(T,\pi\) 的分量均为有理数，取一正整数 \(D\) 使 \(DA\in M_N(\mathbb Z)\)，则对任意素数 \(p\nmid D\) 以下条件等价：
  \[
  DA\ \text{在 }\mathbb F_p\text{ 上奇异}
  \quad\Longleftrightarrow\quad
  \chi_T(t)\equiv (t-1)^2q(t)\pmod p.
  \]
  并且一旦成立，则 \(p\) 必出现在 \(Z\) 的既约分母中，即 \(Z\notin M_N(\mathbb Z[1/D])\)。
\end{enumerate}
\end{theorem}
\begin{proof}
令 \(V=\mathbb Q^N\)。
由 \(\pi^\top T=\pi^\top\) 可知 \(\ker(\pi^\top)\) 为 \(T\)-不变子空间；且 \(V=\mathbb Q\mathbf 1\oplus \ker(\pi^\top)\) 为直和分解。
对 \(v\in \ker(\pi^\top)\)，有 \(\mathbf 1\pi^\top v=0\)，故 \(Av=(I-T)v\)；
对 \(\mathbf 1\)，有 \(A\mathbf 1=\mathbf 1-T\mathbf 1+\mathbf 1\pi^\top\mathbf 1=\mathbf 1\)。
因此
\[
\det(A)=\det(I-T\mid_{\ker(\pi^\top)}).
\]
另一方面，由上述直和分解，特征多项式分解为
\[
\chi_T(t)=(t-1)\cdot \det(tI-T\mid_{\ker(\pi^\top)}),
\]
从而
\(
\chi_T'(1)=\det(I-T\mid_{\ker(\pi^\top)})=\det(A)
\)，得到 (1)。
\par
对 (2)，当 \(p\nmid D\) 时，\(DA\) 在 \(\mathbb F_p\) 上奇异当且仅当 \(\det(A)\equiv 0\pmod p\)。
由 (1) 这等价于 \(\chi_T'(1)\equiv 0\pmod p\)，又因 \(\chi_T(1)=0\)，故等价于 \(t=1\) 在 \(\chi_T(t)\) 中为模 \(p\) 的重根。
最后 \(Z=A^{-1}=\mathrm{adj}(A)/\det(A)\)，若 \(\det(A)\) 的既约分子被 \(p\) 整除，则 \(Z\) 的既约分母必含素因子 \(p\)。证毕。
\end{proof}

\begin{table}[H]
\centering
\scriptsize
\setlength{\tabcolsep}{6pt}
\caption{Window-6 four-block fundamental matrix $Z=(I-T+\mathbf 1\pi^\top)^{-1}$ (exact rationals).}
\label{tab:window6_green_kernel_Z}
\begin{tabular}{c c c c c}
\toprule
 & $U_2$ & $U_1$ & $U_L$ & $U_R$ \\
\midrule
$U_2$ & $(\frac{546339}{584704})$ & $(\frac{32153}{877056})$ & $(\frac{1441}{584704})$ & $(\frac{23233}{877056})$ \\
$U_1$ & $(\frac{26307}{584704})$ & $(\frac{288211}{292352})$ & $(\frac{9345}{584704})$ & $(-\frac{13685}{292352})$ \\
$U_L$ & $(\frac{4323}{584704})$ & $(\frac{34265}{877056})$ & $(\frac{548833}{584704})$ & $(\frac{13057}{877056})$ \\
$U_R$ & $(\frac{69699}{584704})$ & $(-\frac{150535}{877056})$ & $(\frac{13057}{584704})$ & $(\frac{903457}{877056})$ \\
\bottomrule
\end{tabular}
\end{table}


\begin{proposition}[块内停留指数四元组与 \texorpdfstring{$22/7$}{22/7} 的刚性出现]\label{prop:xi-window6-stay-indices-22-7}
\leanverified{paper\_xi\_window6\_stay\_indices\_22\_7}
令块内停留概率 \(s_\alpha:=T_{\alpha\alpha}\)，并定义停留指数 \(I_\alpha:=s_\alpha^{-1}\)。
则
\[
(I_2,I_1,I_L,I_R)=\left(\frac{81}{28},\ \frac{22}{7},\ \frac{27}{2},\ 9\right).
\]
特别地，\(I_1=22/7\) 以严格等式出现，且仅由离散计数 \((|U_1|,E_{11})=(22,21)\) 与定义 \(s_1=2E_{11}/(6|U_1|)\) 强制，不依赖任何连续调参。
\end{proposition}

\begin{proof}
由 \eqref{eq:window6_edge_flux_skeleton}，
\(
s_2=\frac{28}{81},\ s_1=\frac{7}{22},\ s_L=\frac{2}{27},\ s_R=\frac{1}{9}
\)。
逐项取倒数即得。证毕。
\end{proof}

\leanverified{paper\_xi\_window6\_kemeny\_constant\_rational}
\begin{proposition}[Kemeny 常数的有理闭式与 \texorpdfstring{$p'(1)/p(1)$}{p'(1)/p(1)} 恒等式]\label{prop:xi-window6-kemeny-constant-rational}
记特征多项式分解 \(\chi_T(t)=(t-1)\,p(t)\)，其中
\[
p(t)=\frac{48114t^{3}+7263t^{2}-506t-55}{48114}.
\]
则 Kemeny 常数
\(
K:=\sum_{\lambda\in\mathrm{Spec}(T)\setminus\{1\}}\frac{1}{1-\lambda}
\)
满足严格恒等式
\[
K=\frac{p'(1)}{p(1)}=\frac{79181}{27408}.
\]
其既约分母含素因子 \(p_\star\)（见 \eqref{eq:window6_green_kernel_audit_prime_571}），与推论 \ref{cor:xi-window6-block-tree-number-audit-prime} 的树数指纹同源。
\end{proposition}

\begin{proof}
设 \(p(t)=\prod_{j=1}^{3}(t-\lambda_j)\)（\(\lambda_j\) 为非平凡特征值）。
则 \(\frac{p'(t)}{p(t)}=\sum_{j=1}^{3}\frac{1}{t-\lambda_j}\)，取 \(t=1\) 得 \(K=\frac{p'(1)}{p(1)}\)。
将 \(p\) 的显式表达代入并化简得到 \(\frac{79181}{27408}\)（亦见证书 \eqref{eq:window6_green_kernel_audit_prime_571}）。证毕。
\end{proof}

\begin{corollary}[均值首达时间的分母谱量子化：首达时间落入 \texorpdfstring{$\mathbb{Z}[1/3,1/p_\star]$}{Z[1/3,1/p*]}]\label{cor:xi-window6-mean-first-passage-denominator-quantization}
令 \(m_{\alpha\to\beta}\) 为粗粒化链 \(T\) 上从块 \(\alpha\) 首达块 \(\beta\) 的期望步数（\(\alpha\neq\beta\)）。
则全部 \(m_{\alpha\to\beta}\) 为有理数，且其既约分母仅含素因子 \(3\) 与 \(p_\star\)；更具体地，在本 window-\(6\) 骨架中有
\[
m_{\alpha\to\beta}\in \frac{1}{3^{2}p_\star}\mathbb Z,\qquad \alpha\neq\beta,
\]
其全量闭式由表 \ref{tab:window6_mean_first_passage_times} 给出。
\end{corollary}

\begin{proof}
均值首达时间可由基本矩阵 \(Z\) 闭式表示：对 \(\alpha\neq\beta\) 有
\(
m_{\alpha\to\beta}=(Z_{\beta\beta}-Z_{\alpha\beta})/\pi_\beta
\)。
由命题 \ref{prop:xi-window6-green-kernel-arithmetic-audit-prime}，\(Z\) 的分母谱仅含 \(2,3,p_\star\)；
结合平稳分布 \(\pi_\beta=|U_\beta|/64\)（\(|U_\beta|\in\{27,22,9,6\}\)）并化简，得到 \(\alpha\neq\beta\) 时分母仅含 \(3,p_\star\)。
具体分母量子化与全量闭式由表 \ref{tab:window6_mean_first_passage_times} 可审计核验。证毕。
\end{proof}

\begin{table}[H]
\centering
\scriptsize
\setlength{\tabcolsep}{6pt}
\caption{Mean first passage times $m_{\alpha\to\beta}$ for the window-6 four-block coarse chain $T$ (exact rationals; diagonal omitted).}
\label{tab:window6_mean_first_passage_times}
\begin{tabular}{c c c c c}
\toprule
 & $U_2$ & $U_1$ & $U_L$ & $U_R$ \\
\midrule
$U_2$ & -- & $(\frac{4730}{1713})$ & $(\frac{11404}{1713})$ & $(\frac{18338}{1713})$ \\
$U_1$ & $(\frac{10834}{5139})$ & -- & $(\frac{33718}{5139})$ & $(\frac{59032}{5139})$ \\
$U_L$ & $(\frac{3764}{1713})$ & $(\frac{4718}{1713})$ & -- & $(\frac{18550}{1713})$ \\
$U_R$ & $(\frac{3310}{1713})$ & $(\frac{5768}{1713})$ & $(\frac{11162}{1713})$ & -- \\
\bottomrule
\end{tabular}
\end{table}


% AUTO-GENERATED by scripts/exp_window6_green_kernel_audit_prime_571.py
\begin{equation}\label{eq:window6_green_kernel_audit_prime_571}
\begin{aligned}
p_\star&:=571,\\
\det\bigl(I-T+\mathbf 1\pi^\top\bigr)&=\frac{9136}{8019}=\frac{2^{4} \cdot 571}{3^{6} \cdot 11},\\
K:=\sum_{\lambda\in\mathrm{Spec}(T)\setminus\{1\}}\frac{1}{1-\lambda}=\frac{79181}{27408}=\frac{79181}{2^{4} \cdot 3 \cdot 571}.
\end{aligned}
\end{equation}


\begin{proposition}[边界块 \texorpdfstring{$U_R$}{UR} 的首返时间分布完全可解：有理型母函数与显式二阶矩]\label{prop:xi-window6-UR-first-return-time-closed-form}
\leanverified{paper\_xi\_window6\_ur\_first\_return\_time\_closed\_form}
以上述四状态核 \(T\) 为动力学（\eqref{eq:window6_edge_flux_skeleton}），令 \(\tau_R\) 为从状态 \(U_R\) 出发的首返时间（首次回到 \(U_R\) 的步数，取值于 \(\mathbb N\)）。
则其概率母函数 \(F_R(z):=\mathbb E[z^{\tau_R}]\) 为有理函数，并满足闭式证书
% AUTO-GENERATED by scripts/exp_window6_ur_first_return_time.py
\begin{equation}\label{eq:window6_ur_first_return_time_pgf}
F_R(z)=\frac{z\,(110z^3+1525z^2-1221z-10692)}{623z^3+14317z^2+71010z-96228}.
\end{equation}

从而严格矩量为
\[
\mathbb E[\tau_R]=\frac{32}{3},
\qquad
\mathrm{Var}(\tau_R)=\frac{1695268}{15417}.
\]
此外，\(F_R\) 的收敛半径由分母多项式的唯一正根 \(\rho_R>1\) 决定；因此 \(\tau_R\) 具有严格指数型尾部衰减。
\end{proposition}

\begin{corollary}[首返尾率与模尺度的二重非共振：\texorpdfstring{$\log\rho_R/\log\varphi\notin\mathbb Q$}{log rhoR/logphi notin Q}]\label{cor:xi-window6-logrhoR-logphi-irrational}
在命题 \ref{prop:xi-window6-UR-first-return-time-closed-form} 的设定下，令
\[
q_R(z):=623z^3+14317z^2+71010z-96228\in\mathbb Z[z],
\]
并令 \(\rho_R>1\) 为 \(q_R\) 的唯一正根（即 \(F_R\) 的收敛半径）。
则 \(\rho_R\) 为代数次数 \(3\) 的实数，且满足
\[
\frac{\log\rho_R}{\log\varphi}\notin\mathbb Q.
\]
并且同样存在 Baker 型有效下界：存在常数 \(C_R>0\) 使对任意整数对 \((n,k)\neq(0,0)\) 有
\[
\bigl|n\log\rho_R-k\log\varphi\bigr|
\ \ge\
\exp\!\bigl(-C_R\log(1+\max\{|n|,|k|\})\bigr).
\]
\end{corollary}

\begin{proof}
由模 \(5\) 检验可知 \(q_R(z)\equiv 3z^3+2z^2+2\ (\mathrm{mod}\ 5)\) 在 \(\mathbb F_5\) 上无根，从而 \(q_R\) 在 \(\mathbb Q[z]\) 上不可约；因此 \(\rho_R\) 为三次代数数。
令 \(K_R:=\mathbb Q(\rho_R)\)，则 \([K_R:\mathbb Q]=3\)。
与推论 \ref{cor:xi-window6-field-disjointness-degree-six} 同理可得 \(K_R\cap K_\varphi=\mathbb Q\)。
若 \(\log\rho_R/\log\varphi=p/q\in\mathbb Q\)（既约），则 \(\rho_R^q=\varphi^{p}\in K_\varphi\) 且 \(\rho_R^q\in K_R\)，故 \(\rho_R^q\in K_R\cap K_\varphi=\mathbb Q\)，矛盾（因 \(\rho_R>1\Rightarrow p\neq 0\) 且 \(\varphi^{p}\notin\mathbb Q\)）。
最后一式由 Baker--W\"ustholz/Matveev 的线性对数形式下界直接推出（\cite{BakerWustholz2007LogarithmicForms,Matveev2000LinearFormsLogs}）。证毕。
\end{proof}

\begin{corollary}[三底对数的有理线性独立性与二维环面稠密性]\label{cor:xi-window6-logrhoR-logphi-logL-independent}
在推论 \ref{cor:xi-window6-logrhoR-logphi-irrational} 的设定下，对任意整数 \(L\ge 2\)，三数
\(
\log\rho_R,\ \log\varphi,\ \log L
\)
在 \(\QQ\) 上线性无关：若整数 \(m,n,k\in\ZZ\) 满足
\[
m\log\rho_R+n\log\varphi+k\log L=0,
\]
则必有 \(m=n=k=0\)。
\par
因此 \(\log\rho_R/\log L\notin\QQ\)，从而一参数轨道
\[
\bigl\{\,t(\log\rho_R,\log L)\bmod 2\pi:\ t\in\RR\,\bigr\}
\]
在二维环面 \((\RR/2\pi\ZZ)^2\) 上稠密；等价地，任何同时依赖两底相位 \(t\log\rho_R\) 与 \(t\log L\) 的有限型谱构造都不可能出现非零的共同纯虚周期。
\end{corollary}
\begin{proof}
若 \(m\log\rho_R+n\log\varphi+k\log L=0\)，指数化得
\(
\rho_R^{m}\varphi^{n}L^{k}=1
\)，即
\(
\rho_R^{m}=\varphi^{-n}L^{-k}\in\QQ(\varphi)
\)。
左端属于 \(\QQ(\rho_R)\)，故 \(\rho_R^{m}\in \QQ(\rho_R)\cap\QQ(\varphi)=\QQ\)（推论 \ref{cor:xi-window6-logrhoR-logphi-irrational} 的证明已给出交域为 \(\QQ\)）。
由于 \(\rho_R\) 为三次代数数且 \(\rho_R\notin\QQ\)，只能有 \(m=0\)。
此时 \(\varphi^{n}L^{k}=1\)。又 \(L^{k}\in\QQ\)，故 \(\varphi^{n}\in\QQ\)，从而 \(n=0\)，继而 \(k=0\)。
稠密性为 Kronecker--Weyl 定理在二维情形的直接推论：\(\log\rho_R/\log L\notin\QQ\) 等价于二维旋转稠密。证毕。
\end{proof}

\begin{corollary}[有效有理逼近下界与幂律锁定排除]\label{cor:xi-window6-alphaR-rational-approx-bound}
在推论 \ref{cor:xi-window6-logrhoR-logphi-irrational} 的设定下，令
\[
\alpha_R:=\frac{\log\rho_R}{\log\varphi}.
\]
则对任意整数 \(n\ge 1\) 与任意 \(k\in\ZZ\)，成立显式的有理逼近下界
\[
\Bigl\lvert \alpha_R-\frac{k}{n}\Bigr\rvert
=\frac{\bigl|n\log\rho_R-k\log\varphi\bigr|}{n\log\varphi}
\;\ge\;
\frac{1}{n\log\varphi}\exp\!\Bigl(-C_R\log\bigl(1+\max\{n,|k|\}\bigr)\Bigr).
\]
从而不存在非零整数对 \((n,k)\) 使 \(\rho_R^n=\varphi^k\)；
并且任何试图在黄金尺度 \(\varphi^k\) 与首返尾率尺度 \(\rho_R^n\) 之间建立精确幂律锁定（或误差可忽略的近幂律锁定）的方案，均受该有效非共振下界约束。
\end{corollary}
\begin{proof}
由推论 \ref{cor:xi-window6-logrhoR-logphi-irrational} 的 Baker 型下界，直接除以 \(n\log\varphi\) 得到所示不等式。
若存在 \((n,k)\neq(0,0)\) 使 \(\rho_R^n=\varphi^k\)，则 \(n\log\rho_R-k\log\varphi=0\)，与下界矛盾，故不存在幂律锁定。证毕。
\end{proof}

\begin{corollary}[\texorpdfstring{$(1+3+8)$}{(1+3+8)} 的倒数通量：边界稳态通量常数恰为 \texorpdfstring{$1/\dim\mathfrak g_{\mathrm{SM}}$}{1/dim gSM}]\label{cor:xi-window6-reciprocal-flux-boson12}
\leanverified{paper\_xi\_window6\_reciprocal\_flux\_boson12}
在 window-\(6\) 的四块骨架证书中（\eqref{eq:window6_edge_flux_skeleton}），边界块 \(U_R\) 的割大小与稳态通量常数满足
\[
\Phi_R=\frac{|\partial U_R|}{64\cdot 6}=\frac{1}{12}.
\]
若将标准模型规范玻色子计数解释为 \(\dim\mathfrak g_{\mathrm{SM}}=1+3+8=12\)，则得到严格倒数律
\[
\Phi_R=\frac{1}{\dim\mathfrak g_{\mathrm{SM}}}.
\]
\end{corollary}

\begin{corollary}[复位事件时钟的分支频率刚性：\texorpdfstring{$\{34,55\}$}{\{34,55\}} 的极限频率由 \texorpdfstring{$\varphi^{-1}$}{phi^{-1}} 唯一决定]\label{cor:xi-window6-reset-clock-gap-frequencies}
\leanverified{paper\_xi\_window6\_reset\_clock\_gap\_frequencies}
对 window-\(6\) 的复位事件集合 \(\mathcal R_6\)（定义 \ref{def:terminal-reset-events}），其返回时间二值谱为
\(\{F_9,F_{10}\}=\{34,55\}\)（定理 \ref{thm:terminal-reset-events-sturmian}；并见表 \ref{tab:fold_zeck_reset_event_gaps_m6}）。
记字母 \(A:=55\)、\(B:=34\)，则差分字母序列为 Fibonacci 替换 \(\sigma(A)=AB,\sigma(B)=A\) 的不动点；因此两间隔的极限频率存在且为
\[
\lim_{N\to\infty}\frac{\#\{n\le N:\ r_{n+1}-r_n=55\}}{N}=\frac{1}{\varphi},
\qquad
\lim_{N\to\infty}\frac{\#\{n\le N:\ r_{n+1}-r_n=34\}}{N}=\frac{1}{\varphi^2}.
\]
从而 “\(SO(10)\) 尾部偏移 \(34\)” 在无限时钟尺度上获得不可调的概率权重，为统一分支的概率化选择给出可审计的确定性频率定律。
\end{corollary}

\begin{theorem}[尾部偏移半群的 Frobenius 阈值：\texorpdfstring{$\langle 21,34,55\rangle$}{<21,34,55>} 的导通界]\label{thm:xi-window6-tail-semigroup-frobenius-threshold}
\leanverified{paper\_xi\_window6\_tail\_semigroup\_frobenius\_threshold}
在 window-\(6\) uplift 的离散尾部偏移口径下，非零偏移仅可能取自
\(\{21,34,55\}=\{F_8,F_9,F_{10}\}\)（见 \S\ref{subsubsec:window6-terminal-propositions} 的相应结论链）。
令数值半群
\[
S_6:=\langle 21,34,55\rangle\subset\mathbb N.
\]
则 \(S_6\) 为余有限半群，且其 Frobenius 数与导通数满足
% AUTO-GENERATED by scripts/exp_window6_tail_semigroup_frobenius.py
\begin{equation}\label{eq:window6_tail_semigroup_frobenius}
g(S_6)=659,\qquad c(S_6)=660.
\end{equation}

换言之，对任意整数 \(n\ge 660\) 都存在 \(a,b,c\in\mathbb N\) 使
\(
n=21a+34b+55c
\)，而 \(659\) 为最大不可表示数。
\end{theorem}

\begin{corollary}[维数筛法的有限性：在尾部可加假设下 \texorpdfstring{$D>671$}{D>671} 不能再由维数否证]\label{cor:xi-window6-dimension-sieve-finite}
\leanverified{paper\_xi\_window6\_dimension\_sieve\_finite}
若在“额外规范模态仅通过尾部偏移的可加叠加进入”这一可审计假设下，将可达玻色子总数集合写为
\[
\mathcal D_{\mathrm{reach}}:=12+S_6,
\]
则由定理 \ref{thm:xi-window6-tail-semigroup-frobenius-threshold} 的导通数立得：
所有 \(D\ge 12+660=672\) 均属于 \(\mathcal D_{\mathrm{reach}}\)，而不可达维数集合 \(\mathbb N\setminus\mathcal D_{\mathrm{reach}}\) 为有限集，且其最大元为 \(12+659=671\)。
因此，超过 \(671\) 的维数不能再仅由尾部可达性被否证；必须引入更细的签名 \(\Sigma\)、退化谱、或谱穿越证书等更高分辨率不变量。
\end{corollary}

\begin{proposition}[从半群阈值到可证伪预测：硬否证与软筛选的分层]\label{prop:xi-window6-falsifiability-layering-from-semigroup}
\leanverified{paper\_xi\_window6\_falsifiability\_layering\_from\_semigroup}
在区间 \(D\le 671\) 内，\(\mathcal D_{\mathrm{reach}}\) 之外的维数点构成绝对否证点：任何声称由 window-\(6\) 尾部可加机制产生的统一候选若具有这些维数，则与协议的可达性约束直接矛盾。
相反，对 \(\mathcal D_{\mathrm{reach}}\cap[0,671]\) 内的维数点，维数本身不再构成否证，必须进一步满足 Zeckendorf--Fibonacci 分辨率签名的一致审计（例如禁相邻约束、边界层激活的无冗余性以及谱/PSD 证书的兼容性）。
因此可证伪性在该框架中被精确分割为：低维有限区间内由半群阈值提供的硬否证，与同区间内由签名/谱证书提供的软筛选两类。
\end{proposition}

\paragraph{尾部可加闭包的 Ap\'ery 规约：二生成正规形、同余截断账本与复位时钟矩不变量}
本段在定理 \ref{thm:xi-window6-tail-semigroup-frobenius-threshold} 的尾部可加机制之上，把可达性从“生成元集合”进一步规约为一个同余阈值族，并给出与之等价的正规形、不可达集合的有限截断剖分、以及复位间隔时钟的矩不变量。
\par

\begin{theorem}[尾部可加机制的二生成规约与规范解]\label{thm:xi-window6-tail-semigroup-two-generator-normal-form}
在 window-\(6\) uplift 的尾部偏移可加闭包口径下，设可加生成元集合为 \(\{21,34,55\}\)。
由于 \(55=21+34\)，相应可达半群满足
\[
S_6=\langle 21,34,55\rangle=\langle 21,34\rangle\subset\mathbb N.
\]
对任意 \(n\in\mathbb N\)，令 \(r\equiv n\ (\mathrm{mod}\ 21)\)。
定义
\[
b_r:=13r\ (\mathrm{mod}\ 21)\in\{0,1,\dots,20\},
\qquad
w_r:=34\,b_r,
\]
则 \(n\in\langle 21,34\rangle\) 当且仅当 \(n\ge w_r\)。
并且在该可达情形下，存在唯一一组
\(
(a,b)\in\mathbb N\times\{0,1,\dots,20\}
\)
满足
\[
n=21a+34b,\qquad b=b_r,\qquad a=\frac{n-w_r}{21}.
\]
因此可达性在同余类 \(r\) 上具有常数时间判据，并诱导唯一的“模 \(21\) 规范形”（把 \(b\) 规约到 \(\{0,\dots,20\}\)）。
\end{theorem}
\leanverified{paper\_xi\_window6\_tail\_semigroup\_two\_generator\_normal\_form}

\begin{proof}
恒等式 \(55=21+34\) 说明 \(\langle 21,34,55\rangle=\langle 21,34\rangle\)。
又因 \(34\equiv 13\ (\mathrm{mod}\ 21)\) 且 \(13^2\equiv 1\ (\mathrm{mod}\ 21)\)，故 \(13\) 为 \(34\) 在模 \(21\) 下的逆元，从而
\(
34\,b_r\equiv r\ (\mathrm{mod}\ 21)
\)
成立，且 \(b_r\) 在 \(\{0,\dots,20\}\) 内唯一。
\par
若 \(n=21a+34b\) 可达，则必有 \(b\equiv b_r\ (\mathrm{mod}\ 21)\)，从而 \(b=b_r+21k\)（\(k\in\mathbb N\)）并推出
\(
n=w_r+21(a+34k)\ge w_r
\)。
反之，若 \(n\equiv r\ (\mathrm{mod}\ 21)\) 且 \(n\ge w_r\)，则 \(n-w_r\) 为 \(21\) 的非负倍数，取 \(a=(n-w_r)/21\ge 0\) 与 \(b=b_r\) 即得分解。
唯一性由 \(b\in\{0,\dots,20\}\) 的规约条件强制：若 \(21a+34b=21a'+34b'\) 且 \(b,b'\in\{0,\dots,20\}\)，则 \(34(b-b')\) 为 \(21\) 的倍数，因 \(\gcd(34,21)=1\) 得 \(b=b'\)，进而 \(a=a'\)。证毕。
\end{proof}

\begin{proposition}[Ap\'ery 阈值族与不可达计数的同余分解]\label{prop:xi-window6-tail-semigroup-apery-thresholds}
\leanverified{paper\_xi\_window6\_tail\_semigroup\_apery\_thresholds}
令 \(w_r\) 如定理 \ref{thm:xi-window6-tail-semigroup-two-generator-normal-form}。
则 \(\{w_r\}_{r=0}^{20}\) 构成 \(\langle 21,34\rangle\) 关于模 \(21\) 的 Ap\'ery 阈值族：
\[
w_r=\min\{s\in\langle 21,34\rangle:\ s\equiv r\ (\mathrm{mod}\ 21)\}.
\]
并且同余类 \(r\) 内的不可达点精确个数为
\[
g_r:=\frac{w_r-r}{21}\in\mathbb N,
\]
从而不可达集合可写为有限截断并
\(
\bigcup_{r=0}^{20}\{r,\,r+21,\dots,r+21(g_r-1)\}
\)。
其全量数据由表 \ref{tab:window6_tail_semigroup_apery_21_34} 给出，且总不可达计数满足
\[
\sum_{r=0}^{20}g_r=330.
\]
\end{proposition}

% AUTO-GENERATED by scripts/exp_window6_tail_semigroup_apery_21_34.py
\begin{table}[H]
\centering
\small
\setlength{\tabcolsep}{7pt}
\renewcommand{\arraystretch}{1.12}
\caption{Ap\'ery thresholds for the numerical semigroup $\langle 21,34\rangle$ modulo $21$: for each residue $r\in\{0,\dots,20\}$ we list the canonical $b_r\in\{0,\dots,20\}$ with $34b_r\equiv r\pmod{21}$, the Ap\'ery element $w_r=34b_r$, and the gap count $g_r=(w_r-r)/21$.}
\label{tab:window6_tail_semigroup_apery_21_34}
\begin{tabular}{r r r r}
\toprule
$r$ & $b_r$ & $w_r$ & $g_r$ \\
\midrule
0 & 0 & 0 & 0 \\
1 & 13 & 442 & 21 \\
2 & 5 & 170 & 8 \\
3 & 18 & 612 & 29 \\
4 & 10 & 340 & 16 \\
5 & 2 & 68 & 3 \\
6 & 15 & 510 & 24 \\
7 & 7 & 238 & 11 \\
8 & 20 & 680 & 32 \\
9 & 12 & 408 & 19 \\
10 & 4 & 136 & 6 \\
11 & 17 & 578 & 27 \\
12 & 9 & 306 & 14 \\
13 & 1 & 34 & 1 \\
14 & 14 & 476 & 22 \\
15 & 6 & 204 & 9 \\
16 & 19 & 646 & 30 \\
17 & 11 & 374 & 17 \\
18 & 3 & 102 & 4 \\
19 & 16 & 544 & 25 \\
20 & 8 & 272 & 12 \\
\bottomrule
\end{tabular}
\end{table}


\begin{proposition}[维数可达集合的有限截断同余并完全刻画]\label{prop:xi-window6-dimension-reachability-congruence-union}
令总维数写为 \(D=12+n\)（其中 \(n\) 为尾部可加增量）。
则在定理 \ref{thm:xi-window6-tail-semigroup-two-generator-normal-form} 的口径下：
\begin{enumerate}
  \item \(D\) 可达当且仅当存在 \(r\in\{0,\dots,20\}\) 使 \(n\equiv r\ (\mathrm{mod}\ 21)\) 且 \(n\in w_r+21\mathbb N\)；
  等价地，
  \[
  D\in\bigcup_{r=0}^{20}\bigl(12+w_r+21\mathbb N\bigr).
  \]
  \item \(D\) 不可达当且仅当存在 \(r\) 使 \(n=r+21k\) 且 \(0\le k<g_r\)；
  因此不可达集合是 \(21\) 条算术级数的有限截断并，其截断长度由 \(\{g_r\}\) 完全给出。
\end{enumerate}
\end{proposition}
\leanverified{paper\_xi\_window6\_dimension\_reachability\_congruence\_union}

\begin{proposition}[最小可达代表与 Zeckendorf 签名的单射编码]\label{prop:xi-window6-minrep-zeckendorf-signature-injection}
对每个 \(r\in\{0,\dots,20\}\)，定义同余类的最小可达维数
\[
D_r^{\min}:=12+w_r=12+34\,b_r.
\]
令 \(S(D)\subset\mathbb N\) 表示整数 \(D\) 的 Zeckendorf 指标集合（即 \(D=\sum_{k\in S(D)}F_k\)，且 \(S(D)\) 不含相邻指标）。
则 \(\{S(D_r^{\min})\}_{r=0}^{20}\) 两两不同；
并且对全部 \(r\)，有 \(2\in S(D_r^{\min})\) 且 \(3\notin S(D_r^{\min})\)。
其全量数据由表 \ref{tab:window6_tail_semigroup_minrep_zeckendorf_signatures} 给出。
因此，“同余类 \(r\)”在最小代表层面可被 \(S(D_r^{\min})\) 无歧义编码。
\end{proposition}
\leanverified{paper\_xi\_window6\_minrep\_zeckendorf\_signature\_injection}

% AUTO-GENERATED by scripts/exp_window6_tail_semigroup_apery_21_34.py
\begin{table}[H]
\centering
\small
\setlength{\tabcolsep}{7pt}
\renewcommand{\arraystretch}{1.12}
\caption{Minimal reachable representatives $D_r^{\min}=12+w_r$ for each residue class $r\pmod{21}$, together with their Zeckendorf index sets $S(D_r^{\min})$ in Fibonacci indices (so $D_r^{\min}=\sum_{k\in S}F_k$).}
\label{tab:window6_tail_semigroup_minrep_zeckendorf_signatures}
\begin{tabular}{r r l}
\toprule
$r$ & $D_r^{\min}$ & $S(D_r^{\min})$ \\
\midrule
0 & 12 & $\{2,4,6\}$ \\
1 & 454 & $\{2,8,10,14\}$ \\
2 & 182 & $\{2,4,9,12\}$ \\
3 & 624 & $\{2,7,15\}$ \\
4 & 352 & $\{2,6,8,11,13\}$ \\
5 & 80 & $\{2,4,8,10\}$ \\
6 & 522 & $\{2,12,14\}$ \\
7 & 250 & $\{2,4,7,13\}$ \\
8 & 692 & $\{2,5,8,10,15\}$ \\
9 & 420 & $\{2,6,9,14\}$ \\
10 & 148 & $\{2,4,12\}$ \\
11 & 590 & $\{2,7,10,12,14\}$ \\
12 & 318 & $\{2,6,8,10,13\}$ \\
13 & 46 & $\{2,4,6,9\}$ \\
14 & 488 & $\{2,8,11,14\}$ \\
15 & 216 & $\{2,4,7,10,12\}$ \\
16 & 658 & $\{2,7,9,15\}$ \\
17 & 386 & $\{2,6,14\}$ \\
18 & 114 & $\{2,4,8,11\}$ \\
19 & 556 & $\{2,9,12,14\}$ \\
20 & 284 & $\{2,4,7,9,13\}$ \\
\bottomrule
\end{tabular}
\end{table}


\begin{proposition}[同余置换的自对偶性：\texorpdfstring{$r\mapsto b_r$}{r->br} 为对合]\label{prop:xi-window6-congruence-involution}
映射 \(r\mapsto b_r=13r\ (\mathrm{mod}\ 21)\) 满足 \(13^2\equiv 1\ (\mathrm{mod}\ 21)\)，从而为自反置换：
\[
b_{b_r}=r.
\]
因此 Ap\'ery 阈值族满足互换律
\[
w_{b_r}=34\,b_{b_r}=34\,r.
\]
\end{proposition}
\leanverified{paper\_xi\_window6\_congruence\_involution}

\begin{proposition}[阈值族的一阶守恒：平均最小代表的闭式]\label{prop:xi-window6-apery-first-moment}
由于 \(r\mapsto b_r\) 为 \(\{0,\dots,20\}\) 上的置换，得
\[
\sum_{r=0}^{20}b_r=\sum_{k=0}^{20}k=210,
\qquad
\sum_{r=0}^{20}w_r=34\sum_{r=0}^{20}b_r=7140.
\]
从而
\[
\frac{1}{21}\sum_{r=0}^{20}w_r=340,
\qquad
\frac{1}{21}\sum_{r=0}^{20}D_r^{\min}=12+340=352.
\]
\end{proposition}
\leanverified{paper\_xi\_window6\_apery\_first\_moment}

\begin{proposition}[两间隔复位时钟的显式矩不变量]\label{prop:xi-window6-reset-gap-mgf-moments}
在结论 \ref{cor:xi-window6-reset-clock-gap-frequencies} 的设定下，
令复位间隔随机变量 \(G\) 取值于 \(\{34,55\}\)，且
\[
\mathbb P(G=55)=\frac{1}{\varphi},\qquad \mathbb P(G=34)=\frac{1}{\varphi^2}.
\]
则其矩母函数为
\[
M_G(t):=\mathbb E[e^{tG}]
=\frac{1}{\varphi}e^{55t}+\frac{1}{\varphi^2}e^{34t},
\]
并且均值与方差满足闭式
\[
\mathbb E[G]=\frac{55}{\varphi}+\frac{34}{\varphi^2}=13+21\varphi,
\qquad
\mathrm{Var}(G)=441(2\varphi-3).
\]
\end{proposition}

\begin{corollary}[Fibonacci 表达的漂移耦合]\label{cor:xi-window6-reset-gap-fib-coupling}
记 \(F_7=13\)、\(F_8=21\)。
则命题 \ref{prop:xi-window6-reset-gap-mgf-moments} 的漂移与涨落可重写为
\[
\mathbb E[G]=F_7+F_8\varphi,
\qquad
\mathrm{Var}(G)=F_8^2(2\varphi-3),
\]
从而复位漂移与涨落在整数层面由 \((F_7,F_8)\) 与同一无理常数 \(\varphi\) 完全确定。
\end{corollary}

\begin{proposition}[不可达集合的质量不变量与二阶一致性账本]\label{prop:xi-window6-gap-ledger-invariants}
令 \(\mathcal G\subset\mathbb N\) 为 \(\langle 21,34\rangle\) 的不可达增量集合（即命题 \ref{prop:xi-window6-tail-semigroup-apery-thresholds} 的有限截断并对应点集）。
则其基数、总权重以及由截断长度 \(\{g_r\}\) 导出的二阶组合量满足严格恒等式
% AUTO-GENERATED by scripts/exp_window6_tail_semigroup_gap_ledger_21_34.py
\begin{equation}\label{eq:window6_tail_semigroup_gap_ledger_21_34}
\begin{aligned}
|\mathcal{G}|&=330,
&\sum_{n\in\mathcal{G}} n&=75460,
&\frac{1}{|\mathcal{G}|}\sum_{n\in\mathcal{G}} n&=\frac{686}{3},\\
\sum_{r=0}^{20} g_r&=330,
&\sum_{r=0}^{20}\binom{g_r}{2}&=3432,
&\sum_{r=0}^{20} g_r^2&=7194,
&\sum_{r=0}^{20} r\,g_r&=3388.
\end{aligned}
\end{equation}

因此，维数层面的不可达点集 \(\{12+n:\ n\in\mathcal G\}\) 的均值为 \(12+\frac{686}{3}=\frac{722}{3}\)。
上述恒等式与表 \ref{tab:window6_tail_semigroup_apery_21_34} 的阈值族共同构成一组可核验的一致性账本；其任一破坏都等价于同余阈值族或“尾部可加”机制假设的失配。
\end{proposition}
\leanverified{paper\_xi\_window6\_gap\_ledger\_invariants}

\paragraph{Sturmian 复位时钟的有界漂移与常数时间可计算性}

\begin{proposition}[两间隔 Sturmian 复位时钟的有界漂移：抖动上界等于间隔差]\label{prop:xi-window6-reset-clock-bounded-drift}
\leanverified{paper\_xi\_window6\_reset\_clock\_bounded\_drift}
令 window-\(6\) 的复位事件序列如定理 \ref{thm:terminal-reset-events-sturmian}，并记相邻复位间隔为
\[
G_n:=r_{n+1}-r_n\in\{34,55\}\qquad(n\ge 1).
\]
定义部分和（复位时刻增量）
\[
T_n:=\sum_{j=1}^{n}G_j\qquad(n\ge 1),
\]
并令均值漂移常数
\[
\mu:=\mathbb E[G]=\frac{55}{\varphi}+\frac{34}{\varphi^2}=13+21\varphi
\]
如命题 \ref{prop:xi-window6-reset-gap-mgf-moments}。
则存在截距参数 \(\theta\in[0,1)\)，使得长间隔计数
\[
S_n:=\#\{1\le j\le n:\ G_j=55\}
\]
满足 \(S_n=\lfloor n/\varphi+\theta\rfloor\)，并且漂移偏差具有统一有界性
\[
T_n-n\mu=21\Bigl(S_n-\frac{n}{\varphi}\Bigr),
\qquad
\bigl|T_n-n\mu\bigr|\le 21
\quad(\forall n\ge 1).
\]
因此该复位时钟在所有尺度上具有确定性的“零扩散”抖动：偏离均值线性漂移的误差不随 \(n\) 增长，且其最坏幅度恰等于两间隔之差 \(55-34=21\)。
\end{proposition}

\begin{proof}
由定理 \ref{thm:terminal-reset-events-sturmian}，指示序列 \(\ind_{\{G_j=55\}}\) 为 slope \(\varphi^{-1}\) 的 Sturmian 序列（等价于无理旋转的区间窗读出）。
Sturmian 序列的平衡性给出：存在 \(\theta\in[0,1)\) 使得
\(
S_n=\lfloor n/\varphi+\theta\rfloor
\)
且 \(|S_n-n/\varphi|\le 1\)。
另一方面，恒等式 \(T_n=55S_n+34(n-S_n)=34n+21S_n\) 与 \(\mu=34+21/\varphi\) 推出
\(
T_n-n\mu=21(S_n-n/\varphi)
\)；
由 \(|S_n-n/\varphi|\le 1\) 得 \(|T_n-n\mu|\le 21\)。证毕。
\end{proof}

\begin{corollary}[复位时钟的常数时间可计算性：Beatty 取整与无扫描更新]\label{cor:xi-window6-reset-clock-O1-computable}
在命题 \ref{prop:xi-window6-reset-clock-bounded-drift} 的记号下，给定 \(n\) 可在常数时间计算 \(T_n\)：
令 \(S_n=\lfloor n/\varphi+\theta\rfloor\)，则
\[
T_n=34n+21S_n.
\]
从而复位时刻集合 \(\{T_n\}\) 可由一次 Beatty 取整实现，不需枚举 \((G_1,\dots,G_n)\)。
\end{corollary}

\begin{remark}[漂移余差的旋转表示与同余不变量]\label{rem:xi-window6-reset-clock-residual-rotation-congruence}
由 \(S_n=\lfloor n/\varphi+\theta\rfloor\) 得
\[
S_n-\frac{n}{\varphi}=\theta-\Bigl\{\frac{n}{\varphi}+\theta\Bigr\}\in(\theta-1,\theta],
\qquad
T_n-n\mu=21\Bigl(\theta-\Bigl\{\frac{n}{\varphi}+\theta\Bigr\}\Bigr),
\]
其中 \(\{\cdot\}\) 为小数部分。
由于 \(1/\varphi\) 无理，\(\{\frac{n}{\varphi}+\theta\}\) 在 \([0,1)\) 上稠密，从而漂移余差 \(\{T_n-n\mu\}_{n\ge 1}\) 在某个长度为 \(21\) 的区间上稠密。
另一方面，\(\theta-\{\frac{n}{\varphi}+\theta\}\) 的符号仅由
\(\{\frac{n}{\varphi}+\theta\}\in[0,\theta]\) 或 \((\theta,1)\) 两端区间决定；
因此漂移余差的\emph{符号语义}在二元预算下即可压缩（对应 Sturmian 二元窗读出下的局部模式二分）。
另一方面，由 \(34\equiv 55\equiv 13\ (\mathrm{mod}\ 21)\) 得
\[
T_n\equiv 13n\ (\mathrm{mod}\ 21),
\]
故复位时刻在模 \(21\) 的同余类上呈现完全确定的有限剩余类结构。
\end{remark}

\begin{corollary}[复位时钟与尾部半群的可达性耦合：有限障碍仅发生于浅层]\label{cor:xi-window6-reset-clock-semigroup-coupling}
令 \(S=\langle 21,34\rangle\)。
则对任意整数 \(t\ge 715\)，存在 \(n\in\mathbb N\) 与 \(s\in S\) 使
\[
t=T_n+s
\qquad\text{且}\qquad
\bigl|t-(n\mu+s)\bigr|\le 21,
\]
其中 \(\mu\) 与 \(\{T_n\}\) 如命题 \ref{prop:xi-window6-reset-clock-bounded-drift}。
因此在深尺度上，“复位时钟（Sturmian）”与“尾部可达半群（Ap\'ery）”形成确定性的近共终结构；真正的不可达障碍仅可能出现在有限多个浅层尺度上。
\end{corollary}

\begin{proof}
由定理 \ref{thm:xi-window6-tail-semigroup-two-generator-normal-form}（或等价地由 \(\gcd(21,34)=1\) 的 Frobenius 余有限性），有 \(n\ge 660\Rightarrow n\in S\)。
取任意 \(t\ge 715\)。
注意 \(T_1\in\{34,55\}\)，故 \(s:=t-T_1\ge 660\)，从而 \(s\in S\) 且 \(t=T_1+s\)。
另一方面由命题 \ref{prop:xi-window6-reset-clock-bounded-drift}（取 \(n=1\)）得
\(
|T_1-\mu|\le 21
\)，于是
\(
|t-(\mu+s)|=|T_1-\mu|\le 21
\)。
证毕。
\end{proof}

\paragraph{尾部三元组与审计素数的算术闭环}

\begin{corollary}[GUT 三元尾部偏移与最大退化三元组的双源同一：\texorpdfstring{$\{21,34,55\}$}{\{21,34,55\}} 的自洽闭环]\label{cor:xi-tail-offsets-maxfiber-triple-closure}
在 window-\(6\) uplift 的尾部偏移口径下，非零偏移仅可能取自
\[
\Delta_{\mathrm{tail}}(6)\subseteq \{F_8,F_9,F_{10}\}=\{21,34,55\}
\]
（定理 \ref{thm:xi-window6-tail-semigroup-frobenius-threshold}）。
另一方面，由最大纤维闭式（定理 \ref{thm:pom-max-fiber}）在偶数分辨率上有 \(D_{2k}=F_{k+2}\)，从而
\[
D_{12}=F_8=21,\qquad D_{14}=F_9=34,\qquad D_{16}=F_{10}=55.
\]
因此同一 Fibonacci 三元组 \(\{21,34,55\}\) 以两条机制同时出现：一条来自 window-\(6\) 的尾部跃迁字母表，另一条来自更高分辨率的最大退化规模；该双源同一在纯离散层面封闭了“统一三分支”的自洽核对链条。
\end{corollary}

\begin{proof}
第一部分由定理 \ref{thm:xi-window6-tail-semigroup-frobenius-threshold}。
第二部分由定理 \ref{thm:pom-max-fiber} 取 \(k=6,7,8\) 并代入 \(F_8=21,F_9=34,F_{10}=55\)。
合并即得。证毕。
\end{proof}

\begin{corollary}[最大纤维的奇偶闭式与最大达到者 hidden-bit 分裂的有限态谱]\label{cor:xi-max-fiber-parity-hiddenbit-finite-spectrum}
令 \(D_m:=\max_{x\in X_m}d_m(x)\)。
则最大纤维满足奇偶闭式
\[
D_{2k}=F_{k+2},\qquad D_{2k+1}=2F_{k+1}
\qquad(k\ge 1)
\]
（定理 \ref{thm:pom-max-fiber}）。
对任意最大达到者 \(x\in X_m\)（即 \(d_m(x)=D_m\)），按 hidden-bit 分裂定义 \((d_{m,0}(x),d_{m,1}(x))\)（注 \ref{rem:pom-max-fiber-achievers-bsplit}），则在稳定阈值区间内分裂对满足闭式有限态结构：
\begin{itemize}
  \item \textbf{偶数 \(m=2k\ge 10\)：}
  \[
  (d_{2k,0}(x),d_{2k,1}(x))\in\{(F_{k+1},F_k),(F_k,F_{k+1})\}.
  \]
  \item \textbf{奇数 \(m=2k+1\ge 13\)：}
  \[
  (d_{2k+1,0}(x),d_{2k+1,1}(x))\in\{(F_{k+1},F_{k+1}),\ (F_{k+1}\pm F_{k-2},\,F_{k+1}\mp F_{k-2})\}.
  \]
\end{itemize}
特别地，对奇数窗 \(m=2k+1\ge 13\) 定义
\(
p_1(x):=d_{m,1}(x)/D_m
\)。
则三态谱满足
\[
p_1(x)\in\left\{\frac12,\ \frac12\pm\frac{F_{k-2}}{2F_{k+1}}\right\},
\qquad
\Delta_k:=\frac{F_{k-2}}{2F_{k+1}}=\frac12\varphi^{-3}+O(\varphi^{-2k}),
\]
从而 \(p_1\) 的极限谱为 \(\{\varphi^{-2},\tfrac12,\varphi^{-1}\}\)，并以指数速度收敛。
\end{corollary}

\begin{proof}
最大纤维的奇偶闭式由定理 \ref{thm:pom-max-fiber} 给出。
达到者分裂对的闭式有限态结构见注 \ref{rem:pom-max-fiber-achievers-bsplit}，而三态谱与极限谱见推论 \ref{cor:pom-max-fiber-achievers-hidden-bit-prob-spectrum}。
最后的渐近由 Binet 公式给出：\(F_{k-2}/F_{k+1}=\varphi^{-3}+O(\varphi^{-2k})\)。证毕。
\end{proof}

\begin{corollary}[审计素数 \texorpdfstring{$p_\star=571$}{p*=571} 的尾部半群 Ap\'ery 正规形与唯一性]\label{cor:xi-audit-prime-571-apery-normal-form}
令 \(p_\star:=571\) 为四块商图的树数指纹素数（推论 \ref{cor:xi-window6-block-tree-number-audit-prime}），并令尾部半群 \(S=\langle 21,34\rangle\)。
则 \(p_\star\in S\)，且在定理 \ref{thm:xi-window6-tail-semigroup-two-generator-normal-form} 的模 \(21\) 规范形约束 \(b\in\{0,\dots,20\}\) 下具有唯一表示
\[
p_\star=21\cdot 11+34\cdot 10.
\]
其中系数 \(b=10\) 正是同余类 \(p_\star\ (\mathrm{mod}\ 21)\) 的 Ap\'ery 正规形系数（即 \(34\cdot 10\) 为该同余类在 \(S\) 中的最小代表）。
\end{corollary}

\begin{proof}
计算 \(p_\star\equiv 571\equiv 4\ (\mathrm{mod}\ 21)\)。
由定理 \ref{thm:xi-window6-tail-semigroup-two-generator-normal-form}，规范系数为
\(
b_{4}=13\cdot 4\ (\mathrm{mod}\ 21)=10
\)，从而同余类阈值 \(w_4=34\,b_4=340\)。
于是
\(
p_\star-w_4=571-340=231=21\cdot 11
\)，得表示 \(571=21\cdot 11+34\cdot 10\)。
唯一性由同一定理的“模 \(21\) 规范形”唯一性直接推出。证毕。
\end{proof}

\paragraph{两概念相干判别的 Lee--Yang 量子化、椭圆化与高维显式分解}
在既有时间--前缀语义中，二概念竞争可写为带 Gödel 权重的复分区函数。以下命题链把零点谱列、后验阈值、有限多对一投影误差、最小传输子椭圆化以及超立方体核的 Walsh--Joukowsky 解析化置于同一可审计口径。

\begin{theorem}[两概念相干判别的 Lee--Yang 量子化谱列]\label{thm:xi-two-concept-leyang-quantization}
设字母表 \(\Sigma=\{0,1\}\)。
给定两个概念 \(C_0,C_1\)，各自产生可计算无限序列 \(s^{(0)},s^{(1)}\in\Sigma^{\mathbb N}\)，并赋予前缀自由 Gödel 码长 \(\ell_0,\ell_1\in\mathbb N\)。
对任意观测前缀 \(x_{1:n}\in\Sigma^n\)，定义失配能量
\[
E_i(n):=\#\{1\le t\le n:\ s_t^{(i)}\neq x_t\}\qquad(i=0,1),
\]
并记
\[
\Delta_n:=E_1(n)-E_0(n),
\qquad
Z_n(\beta):=2^{-\ell_0}e^{-\beta E_0(n)}+2^{-\ell_1}e^{-\beta E_1(n)}.
\]
若 \(\Delta_n\neq 0\)，则有：
\begin{enumerate}
  \item Lee--Yang 零点全集显式为
  \[
  \mathcal Z_n:=\{\beta\in\mathbb C:\ Z_n(\beta)=0\}
  =\left\{
  -\frac{(\ell_1-\ell_0)\ln 2+(2k+1)\pi i}{\Delta_n}
  :\ k\in\mathbb Z
  \right\}.
  \]
  \item 对任意实 \(\beta\in\mathbb R\)，后验权重
  \[
  \pi_{n,\beta}(0):=\frac{2^{-\ell_0}e^{-\beta E_0(n)}}{Z_n(\beta)},
  \qquad
  \pi_{n,\beta}(1):=1-\pi_{n,\beta}(0)
  \]
  满足严格 logistic 形式
  \[
  \pi_{n,\beta}(0)
  =\frac{1}{1+\exp\!\left(-\bigl[\beta\Delta_n+(\ell_1-\ell_0)\ln 2\bigr]\right)}.
  \]
  \item 对任意 \(\varepsilon\in(0,\tfrac12)\)，条件 \(\pi_{n,\beta}(0)\ge 1-\varepsilon\) 等价于
  \[
  \beta\Delta_n+(\ell_1-\ell_0)\ln 2
  \ge
  \ln\frac{1-\varepsilon}{\varepsilon}.
  \]
  并且全部零点 \(\beta_{n,k}\in\mathcal Z_n\) 具有共同实部
  \[
  \operatorname{Re}\beta_{n,k}
  =-\frac{(\ell_1-\ell_0)\ln 2}{\Delta_n},
  \]
  即后验平衡面 \(\pi_{n,\beta}(0)=\tfrac12\) 与 Lee--Yang 零点垂直谱线一致。
\end{enumerate}
\end{theorem}

\begin{proof}
因式分解得
\[
Z_n(\beta)
=
2^{-\ell_0}e^{-\beta E_0(n)}
\Bigl(1+2^{-(\ell_1-\ell_0)}e^{-\beta\Delta_n}\Bigr).
\]
故 \(Z_n(\beta)=0\) 当且仅当
\[
e^{-\beta\Delta_n}=-2^{\ell_1-\ell_0}.
\]
取复对数得
\[
-\beta\Delta_n=(\ell_1-\ell_0)\ln 2+(2k+1)\pi i,\quad k\in\mathbb Z,
\]
即得零点谱列与共同实部。
再由
\[
\pi_{n,\beta}(0)=\Bigl(1+2^{-(\ell_1-\ell_0)}e^{-\beta\Delta_n}\Bigr)^{-1}
\]
直接化为所示 logistic 形式。
阈值不等式来自
\[
\frac{\pi_{n,\beta}(1)}{\pi_{n,\beta}(0)}
=
2^{-(\ell_1-\ell_0)}e^{-\beta\Delta_n}
\le
\frac{\varepsilon}{1-\varepsilon}
\]
并取对数。证毕。
\end{proof}

\begin{corollary}[相干时间的描述差--失配密度定量律]\label{cor:xi-two-concept-coherence-time-density-law}
在定理 \ref{thm:xi-two-concept-leyang-quantization} 的设定下，进一步假设观测流由 \(C_0\) 生成，即 \(E_0(n)=0\) 且 \(\Delta_n=E_1(n)\)。
记
\[
\underline\delta:=\liminf_{n\to\infty}\frac{E_1(n)}{n}.
\]
若 \(\underline\delta>0\)，则对任意固定 \(\beta>0\) 与 \(\varepsilon\in(0,\tfrac12)\)，存在 \(N\) 使得一切 \(n\ge N\) 皆有
\[
\pi_{n,\beta}(0)\ge 1-\varepsilon.
\]
定义最小相干时间
\[
\tau_\varepsilon(\beta):=\inf\{n\ge 1:\ \pi_{n,\beta}(0)\ge 1-\varepsilon\},
\]
则有渐近上界
\[
\tau_\varepsilon(\beta)
\le
\frac{(\ell_1-\ell_0)\ln 2+\ln\frac{1-\varepsilon}{\varepsilon}}{\beta\,\underline\delta}
+o(1),
\]
其中 \(o(1)\) 由 \(\liminf\) 外推误差诱导。
\end{corollary}

\begin{proof}
由定理 \ref{thm:xi-two-concept-leyang-quantization}，\(\pi_{n,\beta}(0)\ge 1-\varepsilon\) 等价于
\[
\beta E_1(n)+(\ell_1-\ell_0)\ln 2
\ge
\ln\frac{1-\varepsilon}{\varepsilon}.
\]
任取 \(\eta\in(0,\underline\delta)\)，当 \(n\) 充分大时有 \(E_1(n)\ge (\underline\delta-\eta)n\)，从而
\[
n
\ge
\frac{(\ell_1-\ell_0)\ln 2+\ln\frac{1-\varepsilon}{\varepsilon}}
{\beta(\underline\delta-\eta)}
\Longrightarrow
\pi_{n,\beta}(0)\ge 1-\varepsilon.
\]
令 \(\eta\downarrow 0\) 即得所示上界。证毕。
\end{proof}

\begin{conclusion}[前缀自由族上的统一相干证书：可计算停时与显式 I 类错误上界]\label{con:xi-prefixfree-universal-coherence-certificate}
设 \(b\ge 2\)，\(\Sigma_b:=\{0,1,\dots,b-1\}\)，并固定一台输出字母表为 \(\Sigma_b\) 的前缀自由通用机 \(U\)。
对任意显著性水平 \(\alpha\in(0,1)\)，存在一个可计算的顺序判别过程 \(\mathcal T_\alpha\)，其输入为无限序列 \(x\in\Sigma_b^{\NN}\)，并满足：
\begin{enumerate}
  \item 若零假设为 \(x\sim \mathrm{Unif}(\Sigma_b)^{\otimes \NN}\)（独立同分布均匀），则
  \[
  \PP\bigl(\mathcal T_\alpha\ \text{在有限时刻停止}\bigr)\le \alpha.
  \]
  \item 若存在程序 \(p\in\mathrm{dom}(U)\) 使得 \(U(p)=x\)（逐符号输出 \(x\)），则 \(\mathcal T_\alpha\) 在有限时刻停止，并输出某个与 \(x\) 前缀一致的程序（可取停时刻的最短者）。
  更具体地，令
  \[
  n_\alpha(p):=|p|+\left\lceil \log_b\frac1\alpha\right\rceil,
  \]
  则过程在读入 \(x\upharpoonright n_\alpha(p)\) 后必停止。
\end{enumerate}
\end{conclusion}
\begin{proof}
对每个程序 \(p\in\mathrm{dom}(U)\) 定义事件
\[
E_p:=\bigl\{x\in\Sigma_b^{\NN}:\ x\upharpoonright n_\alpha(p)=U(p)\upharpoonright n_\alpha(p)\bigr\},
\]
其中若 \(U(p)\) 在输出 \(n_\alpha(p)\) 个符号前停机则视 \(E_p=\varnothing\)。
在零假设下，\(U(p)\upharpoonright n_\alpha(p)\) 为确定的长度 \(n_\alpha(p)\) 字串，故
\[
\PP(E_p)=b^{-n_\alpha(p)}
=b^{-|p|}\cdot b^{-\lceil\log_b(1/\alpha)\rceil}
\le \alpha\,b^{-|p|}.
\]
由并合界与前缀自由的 Kraft 不等式 \(\sum_{p\in\mathrm{dom}(U)} b^{-|p|}\le 1\) 得
\[
\PP\Bigl(\bigcup_{p\in\mathrm{dom}(U)}E_p\Bigr)
\le
\sum_{p\in\mathrm{dom}(U)}\PP(E_p)
\le
\alpha\sum_{p\in\mathrm{dom}(U)}b^{-|p|}
\le \alpha,
\]
从而可取 \(\mathcal T_\alpha\) 为“首次命中事件族 \(\{E_p\}\)”的顺序检验：当读入 \(x\upharpoonright n\) 时，并行枚举并模拟全部 \(p\) 并尝试生成 \(U(p)\upharpoonright n\)，一旦存在 \(p\) 使得 \(n\ge n_\alpha(p)\) 且 \(U(p)\upharpoonright n=x\upharpoonright n\)，即停止并输出其中 \(|p|\) 最小者。上述构造可计算且其停止事件包含于 \(\bigcup_p E_p\)，故第 (1) 条成立。
若 \(U(p)=x\)，则恒有 \(x\upharpoonright n_\alpha(p)=U(p)\upharpoonright n_\alpha(p)\)，故当过程读到长度 \(n_\alpha(p)\) 时，枚举到 \(p\) 并完成到该长度的模拟比对即可验证并停止，从而第 (2) 条成立。证毕。
\end{proof}

\begin{proposition}[有限多对一投影下的全息聚合界]\label{prop:xi-finite-to-one-holographic-aggregation-bound}
设 \(\mathcal C\) 为微观概念族，带前缀自由码长 \(\ell:\mathcal C\to\mathbb N\)。
设投影 \(\Pi:\mathcal C\to\bar{\mathcal C}\) 满足统一纤维上界
\[
\#\Pi^{-1}(\bar C)\le M<\infty\qquad(\forall \bar C\in\bar{\mathcal C}).
\]
定义聚合权重与等效码长
\[
\bar w(\bar C):=\sum_{C:\Pi(C)=\bar C}2^{-\ell(C)},
\qquad
\bar\ell(\bar C):=-\log_2\bar w(\bar C).
\]
则对任意 \(\bar C\in\bar{\mathcal C}\) 有
\[
\min_{C:\Pi(C)=\bar C}\ell(C)
\le
\bar\ell(\bar C)
\le
\min_{C:\Pi(C)=\bar C}\ell(C)+\log_2 M.
\]
因此，任意以 \(\bar\ell\) 进入分区函数的相干判别中，投影引入的描述代价仅为不随时间增长的常数级扰动，幅度至多 \(\log_2 M\)。
\leanverified{paper\_xi\_finite\_to\_one\_holographic\_aggregation\_bound}
\end{proposition}

\begin{proof}
记
\[
m:=\min_{C:\Pi(C)=\bar C}\ell(C).
\]
对纤维内任意 \(C\) 有 \(\ell(C)\ge m\)，故
\[
\bar w(\bar C)
=
\sum_{C:\Pi(C)=\bar C}2^{-\ell(C)}
\le
\sum_{C:\Pi(C)=\bar C}2^{-m}
\le
M\cdot 2^{-m}.
\]
另一方面存在 \(C_\ast\) 使 \(\ell(C_\ast)=m\)，故 \(\bar w(\bar C)\ge 2^{-m}\)。
对上述双边不等式取 \(-\log_2\) 即得结论。证毕。
\end{proof}

\begin{theorem}[最小相干传输子的谱曲线与 \(j=0\) 的 CM 椭圆曲线]\label{thm:xi-minimal-coherence-propagator-cm-j0}
考虑两顶点有向多重图 \(G\)（顶点集 \(\{0,1\}\)），其边按失配代价赋权：
\[
0\to 0:\ 0,\quad
0\to 1:\ 0,\quad
1\to 0:\ 1,\quad
1\to 1:\ 0\ \text{与}\ 2.
\]
令 \(u\in\mathbb C\) 为失配赝配分参数，定义加权邻接矩阵
\[
M(u):=
\begin{pmatrix}
1 & 1\\
u & 1+u^2
\end{pmatrix}.
\]
再令 \(z\) 为时间生成变量，并定义谱判别式
\[
\Delta(z,u):=\det(I-zM(u))
=1-(u^2+2)z+(u^2-u+1)z^2.
\]
则曲线
\[
E:\ \Delta(z,u)=0\subset \mathbb P^1_z\times\mathbb P^1_u
\]
的正规化是属 \(1\) 的光滑射影曲线；进一步，它与 Weierstrass 曲线
\[
E_0:\quad Y^2=4X^3+1
\]
双有理同构，故 \(j(E_0)=0\)，并具有 \(\mathrm{CM}(\mathbb Z[\omega])\)（\(\omega=e^{2\pi i/3}\)）。
\end{theorem}

\begin{proof}
把 \(\Delta(z,u)=0\) 视作关于 \(z\) 的二次方程，其判别式为
\[
D(u):=(u^2+2)^2-4(u^2-u+1)=u^4+4u=u(u^3+4).
\]
于是 \(E\) 的正规化与超椭圆曲线
\[
C:\ y^2=D(u)=u^4+4u
\]
双有理等价。
因 \(D(u)\) 在 \(\mathbb C\) 上有四个互异根（\(u=0\) 与 \(u^3=-4\) 的三根），故二重覆盖 \(C\to\mathbb P^1_u\) 有四个简单分歧点。
由 Riemann--Hurwitz 公式得 \(g(C)=1\)。
再作双有理变换
\[
X:=\frac{1}{u},
\qquad
Y:=\frac{y}{u^2},
\]
可得
\[
Y^2
=
\frac{u^4+4u}{u^4}
=
1+4X^3,
\]
即 \(C\dashrightarrow E_0\)。
经典结果给出 \(Y^2=4X^3+1\) 的 \(j\)-不变量为 \(0\)，故具 Eisenstein 型复乘法。证毕。
\end{proof}

\begin{proposition}[相干消相时刻的根单位判据]\label{prop:xi-minimal-propagator-root-of-unity-criterion}
\leanverified{paper\_xi\_minimal\_propagator\_root\_of\_unity\_criterion}
沿用定理 \ref{thm:xi-minimal-coherence-propagator-cm-j0} 的 \(M(u)\)。
其特征值为
\[
\lambda_\pm(u)
=
\frac{\operatorname{tr}M(u)\pm\sqrt{D(u)}}{2}
=
\frac{(u^2+2)\pm\sqrt{u^4+4u}}{2}.
\]
对任意 \(n\ge 1\)，有
\[
\operatorname{tr}(M(u)^n)=\lambda_+(u)^n+\lambda_-(u)^n.
\]
在 \(\lambda_-(u)\neq 0\) 的点上，
\[
\operatorname{tr}(M(u)^n)=0
\iff
\left(\frac{\lambda_+(u)}{\lambda_-(u)}\right)^n=-1.
\]
记
\[
r:=\frac{\lambda_+}{\lambda_-},
\]
其在 \(C:y^2=u^4+4u\) 上可写为
\[
r(u,y)=\frac{(u^2+2)+y}{(u^2+2)-y},
\]
则“第 \(n\) 步完全消相”的参数集合满足
\[
\{u:\ \operatorname{tr}(M(u)^n)=0\}
=
\pi_u\!\left(
r^{-1}\!\left(\left\{e^{(2k+1)\pi i/n}:k=0,\dots,n-1\right\}\right)
\right),
\]
其中 \(\pi_u:C\to\mathbb P^1_u\) 为 \(u\)-坐标投影。
\end{proposition}

\begin{proof}
由谱分解（或 Cayley--Hamilton）立得
\(
\operatorname{tr}(M^n)=\lambda_+^n+\lambda_-^n
\)。
当 \(\lambda_-\neq 0\) 时，
\[
\operatorname{tr}(M^n)=0
\iff
\lambda_+^n=-\lambda_-^n
\iff
\left(\frac{\lambda_+}{\lambda_-}\right)^n=-1.
\]
把 \(\sqrt{D(u)}\) 提升为曲线 \(C\) 上的坐标 \(y\)，即得 \(r(u,y)\) 的显式式与根单位原像描述。证毕。
\end{proof}

\begin{theorem}[\(k\)-维超立方体相干核的 Walsh--Joukowsky 分解]\label{thm:xi-hypercube-walsh-joukowsky-leyang-explicit}
令 \(k\in\mathbb N\)，并记 \(\mathcal H_k:=\{0,1\}^k\)。
赋予任意权函数 \(w:\mathcal H_k\to\mathbb C\)（例如 \(w(a)=2^{-\ell(a)}\)）。
对任意观测地址 \(x\in\mathcal H_k\) 与参数 \(u\in\mathbb C\)，定义
\[
Z_x(u):=\sum_{a\in\mathcal H_k}w(a)\,u^{d_H(a,x)}.
\]
记 \([k]:=\{1,\dots,k\}\)，Walsh 字符
\[
\chi_S(y):=(-1)^{\sum_{j\in S}y_j},
\qquad
S\subseteq[k],
\]
及 Walsh 系数
\[
\widehat w(S):=\sum_{a\in\mathcal H_k}w(a)\chi_S(a).
\]
则恒有
\[
Z_x(u)
=
\left(\frac{1+u}{2}\right)^k
\sum_{S\subseteq[k]}
\widehat w(S)
\left(\frac{1-u}{1+u}\right)^{|S|}
\chi_S(x).
\]
等价地，令
\[
t:=\frac{1-u}{1+u},
\qquad
u=\frac{1-t}{1+t},
\]
并定义
\[
P_x(t):=\sum_{S\subseteq[k]}\widehat w(S)\,t^{|S|}\chi_S(x),
\]
则
\[
Z_x(u)=\left(\frac{1+u}{2}\right)^k P_x\!\left(\frac{1-u}{1+u}\right).
\]
因此，除显式因子 \((1+u)^k\) 外，Lee--Yang 零点完全由 \(P_x(t)=0\) 的根经 Möbius 变换 \(u=(1-t)/(1+t)\) 给出；特别地，\(t\)-平面的虚轴被送到 \(u\)-平面的单位圆。
\end{theorem}
\leanverified{paper\_xi\_hypercube\_walsh\_joukowsky\_leyang\_explicit}

\begin{proof}
对每个坐标 \(j\)，置
\[
s_j:=(-1)^{a_j\oplus x_j}=\chi_{\{j\}}(a)\chi_{\{j\}}(x).
\]
则
\[
\mathbf 1_{\{a_j=x_j\}}=\frac{1+s_j}{2},
\qquad
\mathbf 1_{\{a_j\neq x_j\}}=\frac{1-s_j}{2},
\]
从而
\[
u^{d_H(a,x)}
=
\prod_{j=1}^k
\left(
\frac{1+u}{2}+\frac{1-u}{2}s_j
\right).
\]
展开得
\[
u^{d_H(a,x)}
=
\left(\frac{1+u}{2}\right)^k
\sum_{S\subseteq[k]}
\left(\frac{1-u}{1+u}\right)^{|S|}
\prod_{j\in S}s_j.
\]
又
\[
\prod_{j\in S}s_j=\chi_S(a)\chi_S(x).
\]
代回并对 \(a\) 求和即得首式；以 \(t=(1-u)/(1+u)\) 复写即得第二式。证毕。
\end{proof}

% Fourier--Walsh boundary holography: coefficients and Parseval energy on the hypercube.

\begin{theorem}[Fourier--Walsh 全息 Parseval 恒等式：方差作为边界流能量]\label{thm:xi-hypercube-fourier-walsh-boundary-parseval}
令 \(k\in\NN\)，并记 \(\mathcal H_k:=\{0,1\}^k\)。
取任意子集 \(B\subseteq\mathcal H_k\) 与权向量 \(\mathbf w=(w_1,\dots,w_k)\in\RR^k\)，并对 \(\omega\in\mathcal H_k\) 记
\[
w(\omega):=\sum_{i=1}^k w_i\omega_i.
\]
对任意非空 \(I\subseteq[k]\) 令 Walsh 字符 \(\chi_I(\omega):=(-1)^{\sum_{i\in I}\omega_i}\)，并采用归一化 Walsh 系数
\[
\widehat{\ind_B}(I):=2^{-k}\sum_{\omega\in\mathcal H_k}\ind_B(\omega)\chi_I(\omega).
\]
令 \(\partial\overrightarrow{B}:=\{(\omega,\omega^{(i)}):\ \omega\in B,\ \omega^{(i)}\notin B,\ 1\le i\le k\}\) 为有向边界，其中 \(\omega^{(i)}\) 表示翻转第 \(i\) 位后的顶点。
当 \(\sum_{i\in I}w_i\neq 0\) 时，在 \(\partial\overrightarrow{B}\) 上定义 \(1\)-形式
\[
\alpha_{I,\mathbf w}(\omega,\omega^{(i)}):=\frac{w(\omega)-w(\omega^{(i)})}{\sum_{j\in I}w_j}.
\]
则对每个满足 \(\sum_{i\in I}w_i\neq 0\) 的非空 \(I\subseteq[k]\) 成立边界读出恒等式
\[
2^k\,\widehat{\ind_B}(I)=\sum_{(\omega,\nu)\in\partial\overrightarrow{B}}\alpha_{I,\mathbf w}(\omega,\nu).
\]
因此指示函数方差具有严格的“边界能量”表示：
\[
\Var(\ind_B)
=\sum_{\varnothing\neq I\subseteq[k]}\widehat{\ind_B}(I)^2
=2^{-2k}\sum_{\varnothing\neq I\subseteq[k]}
\left(\sum_{(\omega,\nu)\in\partial\overrightarrow{B}}\alpha_{I,\mathbf w}(\omega,\nu)\right)^2,
\]
其中首个等号为 Walsh 基下的 Parseval 恒等式。
\end{theorem}
\begin{proof}
第一式为超立方体上加权边界读出的恒等式，其证明可由对 \(\chi_I\) 的离散散度分解与对 \(B\) 内外消去得到；由定义可知 \(\alpha_{I,\mathbf w}\) 仅依赖于沿边的 \(w\)-差分与 \(\sum_{i\in I}w_i\) 的归一化。
\par
对第二式，由 Walsh--Parseval 恒等式有
\[
\sum_{I\subseteq[k]}\widehat{\ind_B}(I)^2=\EE[\ind_B^2]=\EE[\ind_B],
\]
从而
\(
\sum_{\varnothing\neq I}\widehat{\ind_B}(I)^2=\EE[\ind_B]-\EE[\ind_B]^2=\Var(\ind_B)
\)。
再以第一式将每个 \(\widehat{\ind_B}(I)\) 写为边界求和并整理系数即得。证毕。
\end{proof}

\endinput


\begin{corollary}[Lee--Yang 核的 Stokes 全息：零点由边界决定]\label{cor:xi-hypercube-leyang-stokes-holography}
\leanverified{paper\_xi\_hypercube\_leyang\_stokes\_holography}
在定理 \ref{thm:xi-hypercube-walsh-joukowsky-leyang-explicit} 的设定下，若权函数取为集合 \(B\subseteq\mathcal H_k\) 的指示函数 \(w=\ind_B\)，则对每个 \(x\in\mathcal H_k\)，函数 \(u\mapsto Z_x(u)\) 的 Taylor 系数（等价地 \(t\mapsto P_x(t)\) 的全部系数）可由 \(B\) 的有向边界 \(\partial\overrightarrow{B}\) 的通量数据唯一确定（定理 \ref{thm:xi-hypercube-fourier-walsh-boundary-parseval}）。
因此，对每个固定的 \(x\)，Lee--Yang 零点集合 \(\{u:\ Z_x(u)=0\}\) 亦由 \(\partial\overrightarrow{B}\) 的边界通量唯一确定。
\end{corollary}
\begin{proof}
当 \(w=\ind_B\) 时，
\[
\widehat w(S)=\sum_{a\in\mathcal H_k}\ind_B(a)\chi_S(a)=\sum_{a\in B}\chi_S(a)
=2^k\,\widehat{\ind_B}(S),
\]
其中 \(\widehat{\ind_B}(S)\) 为定理 \ref{thm:xi-hypercube-fourier-walsh-boundary-parseval} 的归一化 Walsh 系数。
取权向量 \(\mathbf 1=(1,\dots,1)\)，则对任意非空 \(S\subseteq[k]\) 有 \(\sum_{i\in S}1\neq 0\)，从而由定理 \ref{thm:xi-hypercube-fourier-walsh-boundary-parseval} 得到
\(
\widehat w(S)=\sum_{(\omega,\nu)\in\partial\overrightarrow{B}}\alpha_{S,\mathbf 1}(\omega,\nu)
\)，即每个非平凡 Walsh 系数为边界通量。
将其代入定理 \ref{thm:xi-hypercube-walsh-joukowsky-leyang-explicit} 的展开式，得到对每个 \(x\) 的 \(Z_x(u)\)（从而 \(P_x(t)\)）完全由边界通量决定。
解析函数由其系数唯一确定，故其零点集合亦由边界通量唯一确定。证毕。
\end{proof}

\begin{theorem}[仿射对合不变权下的 Lee--Yang 多项式二层分解]\label{thm:xi-hypercube-leyang-affine-orbifold-two-layer}
\leanverified{paper\_xi\_hypercube\_leyang\_affine\_orbifold\_two\_layer}
沿用定理 \ref{thm:xi-hypercube-walsh-joukowsky-leyang-explicit} 的设定，取 \(1\le i<j\le k\)，并在 \(\mathcal H_k\) 上定义仿射对合
\[
\sigma_{i,j}(y)_i:=1-y_j,\qquad
\sigma_{i,j}(y)_j:=1-y_i,\qquad
\sigma_{i,j}(y)_r:=y_r\ (r\neq i,j).
\]
若权函数满足 \(\sigma_{i,j}\)-不变性 \(w\circ\sigma_{i,j}=w\)，则对任意 \(x\in\mathcal H_k\) 与 \(t\in\CC\)（其中 \(t=(1-u)/(1+u)\)），多项式
\[
P_x(t)=\sum_{S\subseteq[k]}\widehat w(S)\,t^{|S|}\chi_S(x)
\]
具有严格分解
\[
P_x(t)=P_x^{\mathrm{even}}(t)+\bigl(1-\chi_{\{i,j\}}(x)\bigr)\,Q_x(t),
\]
其中
\[
P_x^{\mathrm{even}}(t):=\sum_{\substack{S\subseteq[k]\\ |S\cap\{i,j\}|\ \mathrm{even}}}\widehat w(S)\,t^{|S|}\chi_S(x),
\]
以及
\[
Q_x(t):=-\sum_{J\subseteq[k]\setminus\{i,j\}}
\widehat w(J\cup\{i\})\,t^{|J|+1}\,\chi_J(x)\,\chi_{\{j\}}(x).
\]
特别地，若 \(x_i=x_j\)（等价于 \(\chi_{\{i,j\}}(x)=1\)），则 \(P_x(t)=P_x^{\mathrm{even}}(t)\)；
若 \(x_i\neq x_j\)（等价于 \(\chi_{\{i,j\}}(x)=-1\)），则 \(P_x(t)=P_x^{\mathrm{even}}(t)+2Q_x(t)\)。
\end{theorem}
\begin{proof}
记 \(R=[k]\setminus\{i,j\}\)。
由 \(w\circ\sigma_{i,j}=w\) 及变元替换 \(a=\sigma_{i,j}(b)\) 得
\[
\widehat w(J\cup\{i\})
=\sum_{a\in\mathcal H_k}w(a)\chi_{J\cup\{i\}}(a)
=\sum_{b\in\mathcal H_k}w(b)\chi_{J\cup\{i\}}(\sigma_{i,j}b)
=-\sum_{b\in\mathcal H_k}w(b)\chi_{J\cup\{j\}}(b)
=-\widehat w(J\cup\{j\}),
\]
其中第三步使用了恒等式 \(\chi_{J\cup\{i\}}\circ\sigma_{i,j}=-\chi_{J\cup\{j\}}\)。
\par
将 \(P_x(t)\) 按 \(|S\cap\{i,j\}|\) 的奇偶分组。
偶层即给出 \(P_x^{\mathrm{even}}(t)\)。
对奇层，配对 \(S=J\cup\{i\}\) 与 \(S'=J\cup\{j\}\)（\(J\subseteq R\)）并用 \(\widehat w(S')=-\widehat w(S)\)，得到
\[
\widehat w(S)t^{|S|}\chi_S(x)+\widehat w(S')t^{|S'|}\chi_{S'}(x)
=\widehat w(J\cup\{i\})\,t^{|J|+1}\,\chi_J(x)\bigl(\chi_{\{i\}}(x)-\chi_{\{j\}}(x)\bigr).
\]
注意
\(
\chi_{\{i\}}(x)-\chi_{\{j\}}(x)=-(1-\chi_{\{i,j\}}(x))\,\chi_{\{j\}}(x)
\)，
将其代回并对 \(J\) 求和即得奇层可因式分解为 \((1-\chi_{\{i,j\}}(x))Q_x(t)\)，从而得到所示分解。
两种取值情形由 \(\chi_{\{i,j\}}(x)\in\{\pm 1\}\) 立得。证毕。
\end{proof}

\endinput


\paragraph{坐标权超立方体的热迹与树复杂度：欧拉型因子化与可逆读出}
定理 \ref{thm:xi-hypercube-walsh-joukowsky-leyang-explicit} 以 Walsh 展开刻画了 Hamming 代价核的 Möbius 归一化零点结构。
同一组 Walsh 字符亦对角化坐标权拉普拉斯，从而将热迹与树复杂度压缩为严格乘积，并在权重取对数素数时呈现欧拉型分解。

\begin{theorem}[坐标权拉普拉斯的 Walsh 谱分解与热迹因子化]\label{thm:xi-hypercube-weighted-laplacian-heat-trace-factorization}
\leanverified{paper\_xi\_hypercube\_weighted\_laplacian\_heat\_trace\_factorization}
沿用定理 \ref{thm:xi-hypercube-walsh-joukowsky-leyang-explicit} 的记号，取权向量 \(\mathbf w=(w_1,\dots,w_k)\in(0,\infty)^k\)。
在 \(\mathcal H_k\) 上定义坐标权拉普拉斯算子
\[
(\Delta_{\mathbf w}f)(\omega):=\sum_{i=1}^{k}w_i\bigl(f(\omega)-f(\omega^{(i)})\bigr),
\]
其中 \(\omega^{(i)}\) 表示把 \(\omega\) 的第 \(i\) 位翻转后的顶点。
则 \(\{\chi_S\}_{S\subseteq[k]}\) 为 \(\Delta_{\mathbf w}\) 的共同特征基，并满足
\[
\Delta_{\mathbf w}\chi_S=2\Bigl(\sum_{i\in S}w_i\Bigr)\chi_S.
\]
从而对任意 \(t>0\) 有热迹完全因子化
\[
\Tr\bigl(e^{-t\Delta_{\mathbf w}}\bigr)
=\sum_{S\subseteq[k]}\exp\!\left(-2t\sum_{i\in S}w_i\right)
=\prod_{i=1}^{k}\bigl(1+e^{-2tw_i}\bigr).
\]
\end{theorem}
\begin{proof}
对任意 \(S\subseteq[k]\)，若 \(i\notin S\) 则 \(\chi_S(\omega^{(i)})=\chi_S(\omega)\)；
若 \(i\in S\) 则 \(\chi_S(\omega^{(i)})=-\chi_S(\omega)\)。
代入定义得到
\(
(\Delta_{\mathbf w}\chi_S)(\omega)=\sum_{i\in S}2w_i\,\chi_S(\omega)
=2(\sum_{i\in S}w_i)\chi_S(\omega)
\)，从而给出谱分解。
热迹为全部特征值的指数和，故
\[
\Tr\bigl(e^{-t\Delta_{\mathbf w}}\bigr)
=\sum_{S\subseteq[k]}\exp\!\left(-2t\sum_{i\in S}w_i\right).
\]
取 \(a_i:=e^{-2tw_i}\) 并用子集展开恒等式
\(
\sum_{S\subseteq[k]}\prod_{i\in S}a_i=\prod_{i=1}^k(1+a_i)
\)
即得乘积式。证毕。
\end{proof}

\begin{corollary}[素数对数权的欧拉型热迹]\label{cor:xi-prime-log-weighted-hypercube-heat-trace-euler}
\leanverified{paper\_xi\_prime\_log\_weighted\_hypercube\_heat\_trace\_euler}
在定理 \ref{thm:xi-hypercube-weighted-laplacian-heat-trace-factorization} 的设定下，若取 \(w_i=\log p_i\)（\(p_i\) 为第 \(i\) 个素数），则
\[
\Tr\bigl(e^{-t\Delta_{\log p}}\bigr)=\prod_{i=1}^{k}\bigl(1+p_i^{-2t}\bigr).
\]
\end{corollary}
\begin{proof}
由 \(e^{-2t\log p_i}=p_i^{-2t}\) 直接代入乘积式即得。证毕。
\end{proof}

\begin{theorem}[热迹的权重刚性：边界配分函数唯一解码多重集]\label{thm:xi-hypercube-heat-trace-weight-rigidity}
在定理 \ref{thm:xi-hypercube-weighted-laplacian-heat-trace-factorization} 的设定下，令
\[
Z_{\mathbf w}(t):=\Tr\bigl(e^{-t\Delta_{\mathbf w}}\bigr)=\prod_{i=1}^{k}\bigl(1+e^{-2tw_i}\bigr),\qquad t>0.
\]
则函数 \(Z_{\mathbf w}:(0,\infty)\to(1,\infty)\) 唯一决定多重集 \(\{w_1,\dots,w_k\}\)（含重数）。
具体地，若存在 \(\mathbf w'=(w'_1,\dots,w'_m)\in(0,\infty)^m\) 使 \(Z_{\mathbf w}(t)\equiv Z_{\mathbf w'}(t)\) 对一切 \(t>0\) 成立，则必有 \(m=k\)，且两组权重一致到重排。
\end{theorem}
\begin{proof}
令 \(a:=\min_i w_i\)，并记其重数为 \(r:=\#\{i:\ w_i=a\}\)。
展开
\(
Z_{\mathbf w}(t)=\prod_i(1+e^{-2tw_i})
\)
得
\[
Z_{\mathbf w}(t)-1=\sum_{i=1}^{k}e^{-2tw_i}+O\!\left(e^{-2t(w_{(1)}+w_{(2)})}\right)\qquad(t\to\infty),
\]
其中 \(w_{(1)}\le w_{(2)}\le\cdots\) 为从小到大排序。
故
\[
\lim_{t\to\infty}-\frac{1}{2t}\log\bigl(Z_{\mathbf w}(t)-1\bigr)=a,
\qquad
\lim_{t\to\infty}e^{2at}\bigl(Z_{\mathbf w}(t)-1\bigr)=r.
\]
因此 \(a\) 与 \(r\) 可由 \(Z_{\mathbf w}\) 唯一恢复。
令
\[
\widetilde Z(t):=\frac{Z_{\mathbf w}(t)}{(1+e^{-2at})^{r}},
\]
则 \(\widetilde Z(t)=\prod_{w_i>a}(1+e^{-2tw_i})\)。
对 \(\widetilde Z\) 重复同一提取过程即可递归恢复全部权重及其重数，从而得到结论。证毕。
\end{proof}

\begin{theorem}[加权超立方体生成树的谱乘积闭式]\label{thm:xi-hypercube-weighted-spanning-tree-spectrum-product}
\leanverified{paper\_xi\_hypercube\_weighted\_spanning\_tree\_spectrum\_product}
令 \(Q_k\) 为顶点集 \(\mathcal H_k\) 的 \(k\)-维超立方体图，并令每条第 \(i\) 坐标方向的边赋权 \(w_i>0\)。
记 \(\tau_{\mathbf w}(Q_k)\) 为按 Kirchhoff--Matrix Tree 口径定义的加权生成树总权（对每棵生成树取边权乘积并求和）。
则有完全闭式
\[
\tau_{\mathbf w}(Q_k)
=\frac{1}{2^k}\prod_{\varnothing\neq S\subseteq[k]}2\sum_{i\in S}w_i
=2^{2^k-k-1}\prod_{\varnothing\neq S\subseteq[k]}\Bigl(\sum_{i\in S}w_i\Bigr).
\]
若进一步取 \(w_i=\log p_i\)，则
\[
\tau_{\log p}(Q_k)
=2^{2^k-k-1}\prod_{\varnothing\neq S\subseteq[k]}\log\!\Bigl(\prod_{i\in S}p_i\Bigr).
\]
\end{theorem}
\begin{proof}
由定理 \ref{thm:xi-hypercube-weighted-laplacian-heat-trace-factorization}，
\(\Delta_{\mathbf w}\) 的非零特征值为
\(
\lambda_S=2\sum_{i\in S}w_i
\)
（\(\varnothing\neq S\subseteq[k]\)）。
加权矩阵树定理给出
\[
\tau_{\mathbf w}(Q_k)
=\frac{1}{|\mathcal H_k|}\prod_{\lambda_S\neq 0}\lambda_S
=\frac{1}{2^k}\prod_{\varnothing\neq S\subseteq[k]}2\sum_{i\in S}w_i,
\]
化简常数因子即得第一式。
素数特例由
\(
\sum_{i\in S}\log p_i=\log\prod_{i\in S}p_i
\)
直接代入。证毕。
\end{proof}

\paragraph{Cartesian 幂模型的双原子凝聚、切向极限与椭圆输运}
以下把迭代层级 \(d\) 视为唯一外参，并以 \(d\to\infty\) 的尺度极限刻画单位圆 Lee--Yang 零点云的两尺度几何：宏观上原子化，切向上复苏为一维连续测度；在 Joukowsky--Gödel 椭圆输运下，该切向高斯极限进一步诱导各向异性的二次像。

\begin{theorem}[Cartesian 幂的 Lee--Yang 零点双原子凝聚与切向中心极限定理]\label{thm:xi-cartesian-power-leyang-double-atom-tangent-clt}
设 \(G\) 为有限无向图，邻接矩阵为 \(A\in\RR^{N\times N}\)，其谱（按重数）为 \(\{\lambda_j\}_{j=1}^{N}\subset\RR\)。
记谱半径与归一化谱测度
\[
\rho:=\max_{1\le j\le N}|\lambda_j|,
\qquad
\nu_G:=\frac1N\sum_{j=1}^{N}\delta_{\lambda_j},
\qquad
\sigma^2:=\int \lambda^2\,d\nu_G(\lambda)=\frac1N\sum_{j=1}^{N}\lambda_j^2.
\]
对每个 \(d\in\NN\)，令 \(G^{\square d}\) 为 \(d\) 次 Cartesian 幂，记其邻接矩阵为 \(A_d\)。
定义缩放后的特征多项式
\[
P_{G,d}(t):=\det\!\Bigl(t I_{N^d}-\frac{2}{d\rho}A_d\Bigr),
\]
并作 Lee--Yang 提升
\[
L_{G,d}(z):=z^{N^d}P_{G,d}(z+z^{-1}).
\]
令 \(\mathcal Z_{G,d}\subset\TT\) 为 \(L_{G,d}\) 的零点多重集，并定义其经验测度
\[
\mu_{G,d}:=\frac1{|\mathcal Z_{G,d}|}\sum_{\zeta\in\mathcal Z_{G,d}}\delta_{\zeta}.
\]
则当 \(d\to\infty\) 时成立：
\begin{enumerate}
  \item \textup{（双原子凝聚）}\(\mu_{G,d}\Longrightarrow \tfrac12\delta_{i}+\tfrac12\delta_{-i}\)。
  \item \textup{（切向 CLT）}在上半圆以 \(\zeta=e^{i\theta}\)（\(\theta\in(0,\pi)\)）作角坐标，并定义切向放大映射
  \[
  B_d(e^{i\theta}):=\sqrt d\Bigl(\theta-\frac{\pi}{2}\Bigr)\in\RR.
  \]
  记上半圆条件化测度
  \[
  \mu_{G,d}^{+}
  :=\frac{\mu_{G,d}\restriction_{\theta\in(0,\pi)}}{\mu_{G,d}(\theta\in(0,\pi))},
  \]
  并置 \(\nu_{G,d}:=(B_d)_*\mu_{G,d}^{+}\)，则
  \[
  \nu_{G,d}\Longrightarrow \mathcal N\!\Bigl(0,\frac{\sigma^2}{\rho^2}\Bigr).
  \]
  同理，在凝聚点 \(-i=e^{i3\pi/2}\) 处以相同尺度展开亦成立。
\end{enumerate}
\leanverified{paper\_xi\_cartesian\_power\_leyang\_double\_atom\_tangent\_clt}
\end{theorem}

\begin{proof}
由 Cartesian 积的谱性质，\(A_d\) 的全部特征值多重集恰为
\[
\left\{\lambda_{j_1}+\cdots+\lambda_{j_d}:\ 1\le j_k\le N\right\},
\]
其按重数归一化的谱测度为 \(\nu_G^{*d}\)。
因此 \(P_{G,d}\) 的根可写为
\[
t=\frac{2}{d\rho}(\lambda_{j_1}+\cdots+\lambda_{j_d})\in[-2,2],
\]
从而方程 \(z+z^{-1}=t\) 在单位圆上给出一对根 \(z=e^{\pm i\theta}\)（其中 \(t=2\cos\theta\)）。
等价地，令 \(X_1,X_2,\dots\overset{\mathrm{i.i.d.}}{\sim}\nu_G\) 并置
\[
U_d:=\frac{1}{d\rho}\sum_{k=1}^{d}X_k,
\]
则 \(\mu_{G,d}\) 是把 \(U_d\) 经 \(u\mapsto e^{\pm i\arccos(u)}\) 推送后的对称化像测度。
注意 \(\EE[X_k]=\tfrac1N\mathrm{tr}(A)=0\) 且 \(\mathrm{Var}(X_k)=\sigma^2\)，故经典中心极限定理给出
\[
\sqrt d\,U_d=\frac{1}{\rho\sqrt d}\sum_{k=1}^{d}X_k
\Longrightarrow
\mathcal N\!\Bigl(0,\frac{\sigma^2}{\rho^2}\Bigr),
\]
从而 \(U_d\to 0\) 依概率收敛。
由 \(\arccos(u)\to \pi/2\)（\(u\to 0\)）及正负号对称性，立得零点宏观凝聚到 \(\pm i\) 且权重各半，得到第一条。
\par
进一步，\(\arccos\) 在 \(0\) 处可微且 \(\arccos'(0)=-1\)，由 Delta 方法
\[
\sqrt d\left(\arccos(U_d)-\frac{\pi}{2}\right)
=-\sqrt d\,U_d+o_{\PP}(1)
\Longrightarrow
\mathcal N\!\Bigl(0,\frac{\sigma^2}{\rho^2}\Bigr).
\]
将 \(\theta=\arccos(U_d)\) 视为上半圆角坐标，并把上述收敛转写为 \((B_d)_*\) 的测度收敛，即得第二条。证毕。
\end{proof}

\begin{theorem}[Joukowsky--Gödel 椭圆输运下的各向异性二次极限律]\label{thm:xi-cartesian-power-joukowsky-godel-anisotropic-parabola}
\leanverified{paper\_xi\_cartesian\_power\_joukowsky\_godel\_anisotropic\_parabola}
取 \(r>1\)，定义 Joukowsky--Gödel 变换
\[
J_r(z):=rz+r^{-1}z^{-1}.
\]
在定理 \ref{thm:xi-cartesian-power-leyang-double-atom-tangent-clt} 的设定下，令 \(\zeta_d\) 依上半圆条件化测度 \(\mu_{G,d}^{+}\) 抽样，并置 \(W_d:=J_r(\zeta_d)\)。
记椭圆上端点 \(w_\star:=i(r-r^{-1})\)。
则当 \(d\to\infty\) 时存在 \(Y\sim\mathcal N(0,\sigma^2/\rho^2)\) 使得
\[
\Bigl(\ \sqrt d\,\Re(W_d-w_\star),\ d\,\Im(W_d-w_\star)\ \Bigr)
\Longrightarrow
\Bigl(\ -(r+r^{-1})Y,\ -\frac{r-r^{-1}}{2}Y^2\Bigr).
\]
特别地，极限支持于抛物线
\[
\Im=\ -\frac{r-r^{-1}}{2(r+r^{-1})^2}\,(\Re)^2.
\]
\end{theorem}

\begin{proof}
写 \(\zeta_d=e^{i\theta_d}\)，并令 \(\Delta_d:=\theta_d-\pi/2\) 为 \(i\) 处的角偏差。
由定理 \ref{thm:xi-cartesian-power-leyang-double-atom-tangent-clt} 的切向 CLT，有 \(\sqrt d\,\Delta_d\Rightarrow Y\)。
对 \(\theta=\pi/2+\Delta\) 作泰勒展开：
\[
\cos(\pi/2+\Delta)=-\Delta+O(\Delta^3),
\qquad
\sin(\pi/2+\Delta)=1-\Delta^2/2+O(\Delta^4).
\]
又
\[
J_r(e^{i\theta})=(r+r^{-1})\cos\theta+i(r-r^{-1})\sin\theta,
\]
故
\[
W_d-w_\star
=(r+r^{-1})\cos\theta_d+i(r-r^{-1})(\sin\theta_d-1),
\]
从而
\[
\Re(W_d-w_\star)=-(r+r^{-1})\Delta_d+O(\Delta_d^3),
\qquad
\Im(W_d-w_\star)=-(r-r^{-1})\frac{\Delta_d^2}{2}+O(\Delta_d^4).
\]
两式分别乘以 \(\sqrt d\) 与 \(d\) 并用 \(\Delta_d=O_{\PP}(d^{-1/2})\) 即得
\[
\sqrt d\,\Re(W_d-w_\star)=-(r+r^{-1})\sqrt d\,\Delta_d+o_{\PP}(1),
\qquad
d\,\Im(W_d-w_\star)=-\frac{r-r^{-1}}{2}(\sqrt d\,\Delta_d)^2+o_{\PP}(1).
\]
由连续映射定理与 Slutsky 定理给出联合收敛。证毕。
\end{proof}

\begin{theorem}[超立方体维数的边界二阶统计可读出性]\label{thm:xi-hypercube-dimension-readable-from-leyang-angular-variance}
\leanverified{paper\_xi\_hypercube\_dimension\_readable\_from\_leyang\_angular\_variance}
令 \(H_d:=K_2^{\square d}\) 为 \(d\) 维超立方体图，其邻接谱为 \(\{d-2k\}_{k=0}^{d}\)（重数 \(\binom{d}{k}\)）。
沿用定理 \ref{thm:xi-cartesian-power-leyang-double-atom-tangent-clt} 的缩放口径（此处基图谱半径为 \(1\)，故 \(\rho=1\)），相应的单位圆 Lee--Yang 角可取
\[
\theta_{d,k}:=\arccos\!\Bigl(1-\frac{2k}{d}\Bigr)\in[0,\pi],
\]
并按权重 \(\binom{d}{k}/2^d\) 出现（单位圆上再与 \(\pm\theta_{d,k}\) 配对）。
定义边界二阶统计量
\[
M_d:=\sum_{k=0}^{d}\frac{\binom{d}{k}}{2^d}\left(\theta_{d,k}-\frac{\pi}{2}\right)^2.
\]
则当 \(d\to\infty\) 时有渐近展开
\[
d\,M_d=1+\frac1d+O\!\Bigl(\frac1{d^2}\Bigr),
\qquad
\frac1{M_d}=d-1+O\!\Bigl(\frac1d\Bigr).
\]
因此，超立方体维数 \(d\) 可由 Lee--Yang 零点角云的单一二阶边界统计量在渐近意义下反演。
\end{theorem}

\begin{proof}
令 \(K\sim\mathrm{Bin}(d,\tfrac12)\)，并置
\[
U_d:=1-\frac{2K}{d}\in[-1,1],
\qquad
\Theta_d:=\arccos(U_d),
\]
则 \(M_d=\EE\bigl[(\Theta_d-\pi/2)^2\bigr]\)。
由二项分布的中心矩计算得
\[
\EE[U_d^2]=\frac1d,
\qquad
\EE[U_d^4]=\frac{3}{d^2}-\frac{2}{d^3},
\qquad
\EE[U_d^6]=\frac{15}{d^3}+O\!\Bigl(\frac1{d^4}\Bigr),
\]
并且由 \(|U_d|\le 1\) 有 \(\EE[U_d^8]=O(d^{-4})\)。
另一方面，\(\arccos\) 在 \(0\) 处的泰勒展开为
\[
\arccos(u)-\frac{\pi}{2}
=-u-\frac{u^3}{6}-\frac{3u^5}{40}+O(u^7),
\]
故平方后
\[
\left(\arccos(u)-\frac{\pi}{2}\right)^2
=u^2+\frac13u^4+\frac{8}{45}u^6+O(u^8).
\]
取期望并代入上述矩，得到
\[
M_d=\frac1d+\frac1{d^2}+O\!\Bigl(\frac1{d^3}\Bigr),
\]
从而 \(dM_d=1+\frac1d+O(d^{-2})\)。
再由 \(\frac1{1+x}=1-x+O(x^2)\)（\(x\to 0\)）得到 \(\frac1{M_d}=d-1+O(d^{-1})\)。证毕。
\end{proof}

\leanverified{paper\_xi\_leyang\_tangent\_dimension\_resurrection}
\begin{corollary}[圆维的两尺度跃迁：切向维数复苏]\label{cor:xi-leyang-tangent-dimension-resurrection}
在定理 \ref{thm:xi-cartesian-power-leyang-double-atom-tangent-clt} 的设定下，\(\mu_{G,d}\) 的宏观弱极限为纯原子测度 \(\tfrac12(\delta_i+\delta_{-i})\)，其 Hausdorff 维数为 \(0\)。
而在凝聚点 \(i\) 处的切向推前测度
\[
\nu_{G,d}:=(B_d)_*\mu_{G,d}^{+}
\]
收敛到高斯测度 \(\mathcal N(0,\sigma^2/\rho^2)\)，从而任一非平凡切向极限测度皆与该高斯测度绝对连续等价并为 exact-dimensional，且其 Hausdorff 维数为 \(1\)。
\end{corollary}

\begin{proof}
第一条由定理 \ref{thm:xi-cartesian-power-leyang-double-atom-tangent-clt} 的双原子凝聚立得。
第二条由同一定理的切向 CLT 立得；高斯测度在 \(\RR\) 上绝对连续且具有正的有限密度，故为 exact-dimensional 且 Hausdorff 维数为 \(1\)。证毕。
\end{proof}


\begin{conjecture}[Gödel 权重下相干零点的普适圆周律与时间读出谱测]\label{conj:xi-godel-weighted-leyang-universal-circle-law}
设 \(\mathfrak C\) 为可枚举概念类，每个概念给出可计算生成子，并以前缀自由 Gödel 码长 \(\ell\) 赋权 \(2^{-\ell}\)。
给定观测流 \(x_{1:\infty}\)，对每个 \(n\) 定义
\[
Z_n(\beta):=
\sum_{C\in\mathfrak C}
2^{-\ell(C)}\exp\!\bigl(-\beta E_C(n)\bigr),
\]
其中 \(E_C(n)\) 为概念 \(C\) 对前缀 \(x_{1:n}\) 的可加失配能量。
猜想：当观测由某一正熵可计算遍历概念生成时，随 \(n\to\infty\)，\(Z_n(\beta)\) 在复 \(\beta\)-平面的 Lee--Yang 零点经适当 Möbius 归一化后收敛到单位圆上的普适测度；该测度角密度由“描述长度--失配率”的联合大偏差谱唯一确定，从而把时间推进 \(n\) 转写为可观测的零点输运流。
\end{conjecture}

\paragraph{Replay 可逆化的 prime-register 终端结构：Gödel 编码、CRT 预算与坏素门槛}
以下结果在折叠--外置范式内，把 replay 证书、Gödel--prime 整数化、分位 CRT 预算、分层投影合成与 Lee--Yang 坏素门槛组织为同一条可审计链条。

\begin{theorem}[Replay-可逆化证书的 Gödel--Prime 终端性]\label{thm:xi-replay-godel-prime-terminality}
\leanverified{paper\_xi\_replay\_godel\_prime\_terminality}
设 \(\mathcal E\) 为可数事件字母表，\(\mathrm{code}:\mathcal E\to\mathbb N_{\ge1}\) 为单射，\((p_t)_{t\ge1}\) 为素数序列。
定义
\[
\mathrm{Enc}_{\mathrm{PR}}:\mathcal E^{<\omega}\to \mathbb N^{(\mathbb P)},
\qquad
\mathrm{Enc}_{\mathrm{PR}}(e_1,\dots,e_T)(p_t):=
\begin{cases}
\mathrm{code}(e_t),&1\le t\le T,\\
0,&t>T,
\end{cases}
\]
并记
\[
\mathrm{PR}(\mathcal E):=\operatorname{Im}\!\bigl(\mathrm{Enc}_{\mathrm{PR}}\bigr),
\qquad
\mathrm{Dec}_{\mathrm{PR}}:=\left.\mathrm{Enc}_{\mathrm{PR}}^{-1}\right|_{\mathrm{PR}(\mathcal E)}.
\]

设 \(\pi:\Omega\to X\) 为多对一投影，给定规范截面 \(s:X\to\Omega\) 与原子逆操作族 \(\{\psi_e:\Omega\to\Omega\}_{e\in\mathcal E}\)。
称 \((R,\rho,\mathrm{Dec})\) 为 replay-可逆化证书系统，若 \(\rho:\Omega\to R\)、\(\mathrm{Dec}:R\to\mathcal E^{<\omega}\) 且
\[
\mathrm{UNLIFT}(x,c):=\psi_{e_T}\circ\cdots\circ\psi_{e_1}\bigl(s(x)\bigr),
\qquad
(e_1,\dots,e_T)=\mathrm{Dec}(c),
\]
满足
\[
\omega=\mathrm{UNLIFT}\bigl(\pi(\omega),\rho(\omega)\bigr)
\quad(\forall\,\omega\in\Omega),
\]
并且 \(\omega\mapsto(\pi(\omega),\rho(\omega))\) 为单射。

则映射
\[
\Phi:R\to \mathrm{PR}(\mathcal E),
\qquad
\Phi(c):=\mathrm{Enc}_{\mathrm{PR}}\bigl(\mathrm{Dec}(c)\bigr)
\]
是满足
\[
\mathrm{Dec}_{\mathrm{PR}}\circ\Phi=\mathrm{Dec}
\]
的唯一映射。
令 \(c\sim c'\iff \mathrm{Dec}(c)=\mathrm{Dec}(c')\)，则 \(\Phi\) 在商集上诱导双射
\[
\overline\Phi:\ R/{\sim}\ \xrightarrow{\ \cong\ }\ \operatorname{Im}\!\bigl(\mathrm{Enc}_{\mathrm{PR}}\circ\mathrm{Dec}\bigr).
\]
若进一步满足 replay 完备性 \(\mathrm{Dec}(R)=\mathcal E^{<\omega}\)，则
\[
\overline\Phi:\ R/{\sim}\ \xrightarrow{\ \cong\ }\ \mathrm{PR}(\mathcal E).
\]
\end{theorem}

\begin{proof}
由定义直接有
\[
\mathrm{Dec}_{\mathrm{PR}}(\Phi(c))
=
\mathrm{Dec}_{\mathrm{PR}}\!\left(\mathrm{Enc}_{\mathrm{PR}}(\mathrm{Dec}(c))\right)
=
\mathrm{Dec}(c),
\]
故 \(\mathrm{Dec}_{\mathrm{PR}}\circ\Phi=\mathrm{Dec}\)。

若 \(c\sim c'\)，则 \(\mathrm{Dec}(c)=\mathrm{Dec}(c')\)，故 \(\Phi(c)=\Phi(c')\)，从而 \(\Phi\) 可因子化到商集。
由像集定义，\(\overline\Phi\) 对 \(\operatorname{Im}(\mathrm{Enc}_{\mathrm{PR}}\circ\mathrm{Dec})\) 显然满射。
若 \(\overline\Phi([c])=\overline\Phi([c'])\)，则
\[
\mathrm{Enc}_{\mathrm{PR}}(\mathrm{Dec}(c))
=
\mathrm{Enc}_{\mathrm{PR}}(\mathrm{Dec}(c')),
\]
由 \(\mathrm{Enc}_{\mathrm{PR}}\) 单射得 \(\mathrm{Dec}(c)=\mathrm{Dec}(c')\)，即 \([c]=[c']\)，故单射。

唯一性：若 \(\Phi':R\to\mathrm{PR}(\mathcal E)\) 也满足 \(\mathrm{Dec}_{\mathrm{PR}}\circ\Phi'=\mathrm{Dec}\)，则对任意 \(c\),
\[
\Phi'(c)
=
\mathrm{Enc}_{\mathrm{PR}}\!\left(\mathrm{Dec}_{\mathrm{PR}}(\Phi'(c))\right)
=
\mathrm{Enc}_{\mathrm{PR}}(\mathrm{Dec}(c))
=
\Phi(c),
\]
故 \(\Phi'=\Phi\)。
若 \(\mathrm{Dec}(R)=\mathcal E^{<\omega}\)，则
\(\operatorname{Im}(\mathrm{Enc}_{\mathrm{PR}}\circ\mathrm{Dec})=\mathrm{PR}(\mathcal E)\)。
证毕。
\end{proof}

\begin{corollary}[算法--素数确定性]\label{cor:xi-replay-algorithm-prime-determinism}
\leanverified{paper\_xi\_replay\_algorithm\_prime\_determinism}
在定理 \ref{thm:xi-replay-godel-prime-terminality} 的设定下，任何证书 \(c\in R\) 的可逆贡献仅由 \(\mathrm{Dec}(c)\) 决定，并与其 prime-register 规范形式 \(\Phi(c)\) 完全等价。
因此在 replay-可逆化语义中，投影算法轨迹与素数寄存器位置化 Gödel 代码构成同构表述。
\end{corollary}

\begin{remark}[与折叠系统实现的对齐]\label{rem:xi-replay-godel-prime-alignment-fold}
上述抽象结构在定义 \ref{def:pom-fold-prime-lift} 与命题 \ref{prop:pom-fold-prime-lift-injective} 中得到具体实现：\(\omega\mapsto(\Fold_m(\omega),r_{\mathsf S}(\omega))\) 的单射化正是 replay-可逆化证书系统的实例化。
\end{remark}

\begin{theorem}[可寻址 Gödel 整数化的精确乘积下界与极小化器分类]\label{thm:xi-addressable-godel-product-lower-bound-classification}
\leanverified{paper\_xi\_addressable\_godel\_product\_lower\_bound\_classification}
设 \(\Sigma\) 为有限字母表，\(|\Sigma|=q\ge 2\)。固定整数 \(T\ge 1\) 与互异素数
\[
p_1,\dots,p_T.
\]
设
\[
E_T:\Sigma^T\hookrightarrow \mathbb N_{\ge 1}
\]
满足如下可寻址性：对每个 \(1\le t\le T\)，存在函数
\[
\delta_t:\mathbb N\to \Sigma
\]
使得对任意 \(s=(s_1,\dots,s_T)\in\Sigma^T\) 都有
\[
s_t=\delta_t\!\bigl(v_{p_t}(E_T(s))\bigr).
\]
则全部码字的乘积满足严格下界
\[
\prod_{s\in\Sigma^T}E_T(s)
\ge
\left(\prod_{t=1}^{T}p_t\right)^{q^T(q-1)/2},
\]
等价地，
\[
\frac{1}{q^T}\sum_{s\in\Sigma^T}\log E_T(s)
\ge
\frac{q-1}{2}\sum_{t=1}^{T}\log p_t.
\]
并且取等当且仅当：对每个 \(t\) 都存在双射
\[
\iota_t:\Sigma\xrightarrow{\sim}\{0,1,\dots,q-1\}
\]
使得
\[
E_T(s)=\prod_{t=1}^{T}p_t^{\,\iota_t(s_t)}
\qquad
(\forall\,s\in\Sigma^T).
\]
换言之，全部极小化器恰为逐坐标、逐素数、无冗余因子的纯估值编码。
\end{theorem}
\begin{proof}
对每个 \(t\) 与 \(a\in\Sigma\)，记
\[
\mathcal V_{t,a}:=
\{v_{p_t}(E_T(s)):\ s_t=a\}\subset \ZZ_{\ge 0},
\qquad
m_{t,a}:=\min \mathcal V_{t,a}.
\]
由于 \(s_t\) 由 \(v_{p_t}(E_T(s))\) 经 \(\delta_t\) 单独解码，不同字母 \(a\neq b\) 的估值集合 \(\mathcal V_{t,a},\mathcal V_{t,b}\) 必两两不交。故
\[
\{m_{t,a}:a\in\Sigma\}
\]
是 \(q\) 个互异的非负整数，从而
\[
\sum_{a\in\Sigma}m_{t,a}\ge 0+1+\cdots+(q-1)=\frac{q(q-1)}{2}.
\]

对固定的 \(t\)，把所有长度 \(T\) 的词按第 \(t\) 位分组，得到
\[
\sum_{s\in\Sigma^T}v_{p_t}(E_T(s))
=
\sum_{a\in\Sigma}\sum_{s_t=a}v_{p_t}(E_T(s))
\ge
q^{T-1}\sum_{a\in\Sigma}m_{t,a}
\ge
q^T\frac{q-1}{2}.
\]
于是
\[
\begin{aligned}
\sum_{s\in\Sigma^T}\log E_T(s)
&=
\sum_{s\in\Sigma^T}\sum_{p}v_p(E_T(s))\log p\\
&\ge
\sum_{t=1}^{T}\log p_t\sum_{s\in\Sigma^T}v_{p_t}(E_T(s))\\
&\ge
q^T\frac{q-1}{2}\sum_{t=1}^{T}\log p_t.
\end{aligned}
\]
指数化即得乘积下界，除以 \(q^T\) 即得平均对数下界。

若取等，则上述每一步均须取等。首先，求和中被忽略的其它素数贡献必须全部为零，因此任一码字都不含 \(p_1,\dots,p_T\) 之外的额外素因子。其次，对每个固定的 \(t,a\)，必有
\[
v_{p_t}(E_T(s))=m_{t,a}
\qquad
(\forall\,s\in\Sigma^T\text{ with }s_t=a),
\]
即第 \(t\) 个素数的估值仅依赖于第 \(t\) 个字母。最后，由
\[
\sum_{a\in\Sigma}m_{t,a}=\frac{q(q-1)}{2}
\]
且这些 \(m_{t,a}\) 是 \(q\) 个互异非负整数，必知
\[
\{m_{t,a}:a\in\Sigma\}=\{0,1,\dots,q-1\}.
\]
故存在双射 \(\iota_t:\Sigma\to\{0,1,\dots,q-1\}\) 使
\[
v_{p_t}(E_T(s))=\iota_t(s_t).
\]
再结合“无额外素因子”，便得到
\[
E_T(s)=\prod_{t=1}^{T}p_t^{\,\iota_t(s_t)}.
\]
反之，此类编码显然满足逐坐标可寻址性，并在每个坐标上以最小可能估值集 \(\{0,\dots,q-1\}\) 取到全部下界等号。证毕。
\end{proof}

\begin{corollary}[可寻址 Gödel 整数化的平均复杂度最优律与二元最坏情形闭式]\label{cor:xi-addressable-godel-average-binary-optimality}
\leanverified{paper\_xi\_addressable\_godel\_average\_binary\_optimality}
沿用定理 \ref{thm:xi-addressable-godel-product-lower-bound-classification} 的记号。则：
\begin{enumerate}
  \item 全部可寻址编码的最小平均对数复杂度精确为
  \[
  \inf_{E_T}
  \frac{1}{q^T}\sum_{s\in\Sigma^T}\log E_T(s)
  =
  \frac{q-1}{2}\sum_{t=1}^{T}\log p_t.
  \]
  \item 存在绝对常数 \(c>0\)，使得
  \[
  \inf_{E_T}
  \frac{1}{q^T}\sum_{s\in\Sigma^T}\log E_T(s)
  \ge
  \frac{q-1}{2}\,c\,T\log T.
  \]
  \item 当 \(q=2\) 时，最小最坏情形复杂度也被完全锁定：
  \[
  \inf_{E_T}\ \max_{s\in\Sigma^T}\log E_T(s)
  =
  \sum_{t=1}^{T}\log p_t.
  \]
\end{enumerate}
\end{corollary}
\begin{proof}
第 \(1\) 条由定理 \ref{thm:xi-addressable-godel-product-lower-bound-classification} 的下界与极小化器分类直接得到。任选双射
\[
\iota_t:\Sigma\xrightarrow{\sim}\{0,1,\dots,q-1\},
\]
则纯估值编码
\[
E_T^{\min}(s):=\prod_{t=1}^{T}p_t^{\,\iota_t(s_t)}
\]
达到该平均值。

对第 \(2\) 条，由 \(p_t\ge t+1\) 得
\[
\sum_{t=1}^{T}\log p_t
\ge
\sum_{t=1}^{T}\log(t+1)
=
\log((T+1)!).
\]
再用 Stirling 下界可知存在绝对常数 \(c>0\)，使对充分大的 \(T\) 有
\[
\log((T+1)!)\ge c\,T\log T.
\]
吸收有限个小 \(T\) 的情形后即得所示统一下界。

最后处理 \(q=2\)。由定理 \ref{thm:xi-addressable-godel-product-lower-bound-classification} 的极小化器分类，取
\[
\iota_t:\Sigma\xrightarrow{\sim}\{0,1\}
\]
即可得到
\[
\max_{s\in\Sigma^T}\log E_T^{\min}(s)=\sum_{t=1}^{T}\log p_t.
\]
反之，任取任一可寻址编码 \(E_T\)。对每个 \(t\)，两字母对应的估值集合是两个互不相交的非空非负整数集，因此总有某个字母 \(a_t^\star\in\Sigma\) 满足
\[
m_{t,a_t^\star}\ge 1.
\]
取
\[
s^\star:=(a_1^\star,\dots,a_T^\star)\in\Sigma^T,
\]
则对每个 \(t\) 都有
\[
v_{p_t}(E_T(s^\star))\ge m_{t,a_t^\star}\ge 1,
\]
从而
\[
\log E_T(s^\star)\ge \sum_{t=1}^{T}\log p_t.
\]
于是
\[
\max_{s\in\Sigma^T}\log E_T(s)\ge \sum_{t=1}^{T}\log p_t.
\]
上下界合并即得第 \(3\) 条。证毕。
\end{proof}

\begin{theorem}[可寻址 Gödel 外置的最坏相对冗余比无界发散]\label{thm:xi-addressable-godel-worstcase-redundancy-diverges}
\leanverified{paper\_xi\_addressable\_godel\_worstcase\_redundancy\_diverges}
沿用定理 \ref{thm:xi-addressable-godel-product-lower-bound-classification} 的记号。对任意满足逐坐标可寻址性的编码
\[
E_T:\Sigma^T\hookrightarrow \NN_{\ge 1},
\qquad
|\Sigma|=q\ge 2,
\]
定义最坏对数复杂度与相对冗余比
\[
L_T(E_T):=\max_{s\in\Sigma^T}\log E_T(s),
\qquad
\mathscr R_T(E_T):=\frac{L_T(E_T)}{T\log q}.
\]
则存在仅依赖于 \(q\) 的常数 \(c_q>0\)，使得对一切 \(T\ge 1\) 都有
\[
\boxed{\;
L_T(E_T)\ge c_q\,T\log T.\;}
\]
因而不存在任何可寻址编码族使
\[
L_T(E_T)=O(T).
\]
更进一步，
\[
\boxed{\;
\mathscr R_T(E_T)\ge \frac{c_q}{\log q}\log T\longrightarrow +\infty.\;}
\]
也就是说，一旦要求每个坐标都能被单独 prime-address 解码，则最坏情形开销相对于原始信息量 \(T\log q\) 的膨胀因子必然无界增长。

特别地，当 \(q=2\) 时，推论 \ref{cor:xi-addressable-godel-average-binary-optimality} 还给出精确最优值
\[
\inf_{E_T}L_T(E_T)=\sum_{t=1}^{T}\log p_t=\Theta(T\log T).
\]
因此，把链上规范投影通过
\[
\mathcal I_n\cong \{0,1\}^{\,n-1}
\]
识别为二元标签空间后，任何逐位可寻址的 Gödel 算术化都必须支付
\[
\Omega(n\log n)
\]
的最坏情形对数代价；再结合推论 \ref{cor:xi-chain-interior-canonical-arithmeticization-optimal-order}，这一阶数同时又是 sharp 的。
\end{theorem}
\begin{proof}
由定理 \ref{thm:xi-addressable-godel-product-lower-bound-classification}，
\[
\frac{1}{q^T}\sum_{s\in\Sigma^T}\log E_T(s)
\ge
\frac{q-1}{2}\sum_{t=1}^{T}\log p_t.
\]
而平均值不超过最大值，故
\[
L_T(E_T)\ge \frac{q-1}{2}\sum_{t=1}^{T}\log p_t.
\]
由素数定理的标准推论，存在绝对常数 \(c>0\) 使
\[
\sum_{t=1}^{T}\log p_t\ge c\,T\log T
\]
对充分大的 \(T\) 成立；吸收有限个小 \(T\) 后可得对全部 \(T\ge1\) 都成立的常数
\[
c_q:=\frac{q-1}{2}c>0.
\]
于是
\[
L_T(E_T)\ge c_q\,T\log T,
\]
这立刻排除了 \(L_T(E_T)=O(T)\) 的可能性。

再除以 \(T\log q\)，得到
\[
\mathscr R_T(E_T)\ge \frac{c_q}{\log q}\log T\to+\infty,
\]
即最坏相对冗余比无界发散。

当 \(q=2\) 时，推论 \ref{cor:xi-addressable-godel-average-binary-optimality} 已给出
\[
\inf_{E_T}\max_{s\in\Sigma^T}\log E_T(s)=\sum_{t=1}^{T}\log p_t,
\]
故最优最坏复杂度恰为 \(\Theta(T\log T)\)。

最后，对链上规范投影，由定理 \ref{thm:xi-chain-interior-boolean-flag-closed-form} 有
\[
|\mathcal I_n|=2^{n-1},
\]
并由固定点集识别得到
\[
\mathcal I_n\cong \{0,1\}^{\,n-1}.
\]
于是取 \(T=n-1\)、\(q=2\) 即知任何逐位可寻址的 Gödel 外置都满足最坏情形下界 \(\Omega(n\log n)\)；而推论 \ref{cor:xi-chain-interior-canonical-arithmeticization-optimal-order} 又给出 matching 的 \(O(n\log n)\) 上界，从而该阶数是 sharp 的。证毕。
\end{proof}

\paragraph{可寻址 Gödel 外置的 Chebyshev 码长主项与见证位膨胀障碍}
在逐坐标可寻址的 Gödel 整数化已经获得精确乘积下界、平均复杂度最优律与最坏冗余发散之后，还可把最坏情形码长进一步精化为由 Chebyshev 函数控制的显式主项。

\begin{theorem}[可寻址 Gödel 外置的最坏码长下界由 Chebyshev 主项精确控制]\label{thm:xi-addressable-godel-worstcase-chebyshev-sharp-lower-bound}
\leanverified{paper\_xi\_addressable\_godel\_worstcase\_chebyshev\_sharp\_lower\_bound}
沿用定理 \ref{thm:xi-addressable-godel-product-lower-bound-classification} 与定理 \ref{thm:xi-addressable-godel-worstcase-redundancy-diverges} 的记号。对任意满足逐坐标可寻址性的单射
\[
E_T:\Sigma^T\hookrightarrow \mathbb N_{\ge 1},
\qquad
|\Sigma|=q\ge 2,
\]
记
\[
L_T(E_T):=\max_{s\in\Sigma^T}\log E_T(s).
\]
若把定理 \ref{thm:xi-addressable-godel-product-lower-bound-classification} 中出现的互异素数按升序重排为
\[
\pi_1<\cdots<\pi_T,
\]
并记 \(\mathfrak p_t\) 为第 \(t\) 个素数，则
\[
L_T(E_T)
\ge
\sum_{t=1}^{T}\log \pi_t
\ge
\sum_{t=1}^{T}\log \mathfrak p_t
=
\vartheta(\mathfrak p_T).
\]
从而
\[
L_T(E_T)
\ge
T\log T+T\log\log T+O(T).
\]
\end{theorem}
\begin{proof}
沿用定理 \ref{thm:xi-addressable-godel-product-lower-bound-classification} 中的估值族记号
\[
\mathcal V_{t,a}:=
\{v_{p_t}(E_T(s)):\ s_t=a\},
\qquad
m_{t,a}:=\min \mathcal V_{t,a}.
\]
由于不同字母 \(a\neq b\) 的集合 \(\mathcal V_{t,a},\mathcal V_{t,b}\) 两两不交，故最小值 \(m_{t,a}\) 也彼此不同。特别地，在这 \(q\ge 2\) 个互异非负整数中，至多有一个等于 \(0\)。于是对每个 \(t\)，都可选取某个字母 \(a_t^\star\in\Sigma\) 使
\[
m_{t,a_t^\star}\ge 1.
\]
令
\[
s^\star:=(a_1^\star,\dots,a_T^\star)\in\Sigma^T.
\]
则对每个 \(t\) 都有
\[
v_{p_t}(E_T(s^\star))\ge m_{t,a_t^\star}\ge 1.
\]
因此
\[
L_T(E_T)
\ge
\log E_T(s^\star)
\ge
\sum_{t=1}^{T}\log p_t
=
\sum_{t=1}^{T}\log \pi_t.
\]

再由 \(\pi_1<\cdots<\pi_T\) 为 \(T\) 个互异素数，知其逐项满足
\[
\pi_t\ge \mathfrak p_t
\qquad (1\le t\le T),
\]
故
\[
\sum_{t=1}^{T}\log \pi_t
\ge
\sum_{t=1}^{T}\log \mathfrak p_t
=
\vartheta(\mathfrak p_T).
\]
这就得到第一条下界。

最后，由经典素数定理的精化形式，
\[
\mathfrak p_T
=
T\log T+T\log\log T+O(T),
\]
并且
\[
\vartheta(x)=x+o(x)\qquad (x\to\infty).
\]
代入 \(x=\mathfrak p_T\) 即得
\[
\vartheta(\mathfrak p_T)
=
T\log T+T\log\log T+O(T).
\]
于是所示精化下界成立。证毕。
\end{proof}

\leanverified{paper\_xi\_addressable\_witness\_no\_asymptotically\_linear\_bitlength}
\begin{corollary}[逐坐标 prime-addressable 见证外置不存在渐近线性位长度]\label{cor:xi-addressable-witness-no-asymptotically-linear-bitlength}
设 \(T=T(n)\to\infty\)，并对每个 \(n\) 给定满足逐坐标可寻址性的单射
\[
E_{T(n)}:\Sigma^{T(n)}\hookrightarrow \mathbb N_{\ge 1},
\qquad
|\Sigma|\ge 2.
\]
则
\[
\max_{s\in\Sigma^{T(n)}}\log_2 E_{T(n)}(s)
\ge
T(n)\log_2 T(n)+T(n)\log_2\log T(n)+O(T(n)).
\]
特别地，
\[
\max_{s\in\Sigma^{T(n)}}\log_2 E_{T(n)}(s)\neq O(T(n)).
\]
若进一步
\[
T(n)=n^c
\qquad (c>0),
\]
则更具体地有
\[
\max_{s\in\Sigma^{T(n)}}\log_2 E_{T(n)}(s)
\ge
c\,n^c\log_2 n+n^c\log_2\log n+O(n^c).
\]
\end{corollary}
\begin{proof}
把定理 \ref{thm:xi-addressable-godel-worstcase-chebyshev-sharp-lower-bound} 应用于 \(T=T(n)\)，并除以 \(\log 2\)，得到
\[
\max_{s\in\Sigma^{T(n)}}\log_2 E_{T(n)}(s)
\ge
\frac{1}{\log 2}
\Bigl(
T(n)\log T(n)+T(n)\log\log T(n)+O(T(n))
\Bigr),
\]
即首个不等式。由于右端严格超线性于 \(T(n)\)，故不可能属于 \(O(T(n))\)。

若 \(T(n)=n^c\)，则
\[
T(n)\log_2 T(n)=c\,n^c\log_2 n,
\]
且
\[
T(n)\log_2\log T(n)
=
n^c\log_2(c\log n)
=
n^c\log_2\log n+O(n^c).
\]
代回即可得到所示具体下界。证毕。
\end{proof}

\begin{theorem}[Gödel--Prime 外置的不可约比特复杂度下界]\label{thm:xi-godel-prime-bit-complexity-lower-bound}
在定理 \ref{thm:xi-replay-godel-prime-terminality} 的设定下，对任意 \(r\in\mathrm{PR}(\mathcal E)\)，记
\[
T:=\bigl|\mathrm{Dec}_{\mathrm{PR}}(r)\bigr|,
\qquad
G(r):=\prod_{t=1}^{T}p_t^{\,r(p_t)+1}.
\]
则成立点态下界
\[
\log G(r)\ \ge\ \sum_{t=1}^{T}\log p_t
=
\vartheta(p_T)
=
(1+o(1))\,T\log T.
\]

进一步，在折叠系统中令 \(r=r_{\mathsf S}(\omega)\)（定义 \ref{def:pom-fold-prime-lift}），\(T=T_{m+1}^{\mathsf S}(\omega)\)，并取 \(\omega\sim\mathrm{Unif}(\Omega_{m+1})\)。
则有
\[
\mathbb E\!\left[\log G\!\bigl(r_{\mathsf S}(\omega)\bigr)\right]
\ge
\left(\frac{4}{27}+o(1)\right)m\log m.
\]
\end{theorem}

\begin{proof}
由 \(r(p_t)+1\ge1\) 得
\[
\log G(r)=\sum_{t=1}^{T}\bigl(r(p_t)+1\bigr)\log p_t
\ge
\sum_{t=1}^{T}\log p_t
=
\vartheta(p_T).
\]
由素数定理的标准推论 \(\vartheta(p_T)\sim p_T\sim T\log T\)，得首个渐近式。

对期望下界，固定 \(\eta\in(0,1)\)，存在 \(T_0\) 使 \(T\ge T_0\Rightarrow \vartheta(p_T)\ge (1-\eta)T\log T\)。
因此
\[
\mathbb E[\log G(r_{\mathsf S})]
\ge
(1-\eta)\,\mathbb E[T\log T]-O(1).
\]
又 \(x\mapsto x\log x\) 在 \((0,\infty)\) 上凸，故
\[
\mathbb E[T\log T]\ge \mathbb E[T]\log\mathbb E[T].
\]
由推论 \ref{cor:pom-prime-residual-linear-lb}，
\(
\mathbb E[T]\ge(\frac{4}{27}+o(1))m
\)。
代回即得
\[
\mathbb E[\log G(r_{\mathsf S})]
\ge
\left(\frac{4}{27}+o(1)\right)m\log m.
\]
证毕。
\end{proof}

\begin{proposition}[分位-CRT 外置下的最小素数轴数量标度律]\label{prop:xi-quantile-crt-prime-axis-scaling}
沿用定义 \ref{def:pom-gamma-quantile} 与推论 \ref{cor:pom-prime-register-quantile-budget} 的阈值
\[
P_m>2B_m(\varepsilon),
\qquad
B_m(\varepsilon):=\exp\!\bigl(m\gamma_m(\varepsilon)\bigr).
\]
现约束模轴均为素数：取 \(q_1<\cdots<q_K\) 并记
\[
P_{m,K}:=\prod_{i=1}^{K}q_i.
\]
在“最小化最大模数 \(q_K\)”的设计准则下，最优方案可取 \(q_i=p_i\)（前 \(K\) 个素数）。

若 \(m\gamma_m(\varepsilon)\to\infty\)，则满足
\(
\prod_{i=1}^{K}p_i>2B_m(\varepsilon)
\)
的最小 \(K=K(m,\varepsilon)\) 具有渐近
\[
K(m,\varepsilon)
=
\left(1+o(1)\right)
\frac{m\gamma_m(\varepsilon)}{\log\!\bigl(m\gamma_m(\varepsilon)\bigr)},
\]
\[
p_{K(m,\varepsilon)}
=
\left(1+o(1)\right)\,m\gamma_m(\varepsilon).
\]
\end{proposition}

\begin{proof}
固定 \(K\) 时，任意递增素数组 \(q_1<\cdots<q_K\) 都满足 \(q_i\ge p_i\)，故
\[
\prod_{i=1}^{K}q_i\ge \prod_{i=1}^{K}p_i=:p_K\#.
\]
因此要以最小 \(q_K\) 达到给定乘积阈值，取 \(q_i=p_i\) 为最优。

阈值条件写为
\[
\log(p_K\#)>m\gamma_m(\varepsilon)+\log 2.
\]
由 \(\log(p_K\#)=\vartheta(p_K)\sim p_K\)，且 \(m\gamma_m(\varepsilon)\to\infty\)，得
\[
p_K\sim m\gamma_m(\varepsilon).
\]
再由 \(p_K\sim K\log K\) 得
\[
K\log K\sim m\gamma_m(\varepsilon),
\]
从而
\[
K=
\left(1+o(1)\right)\frac{m\gamma_m(\varepsilon)}{\log(m\gamma_m(\varepsilon))}.
\]
证毕。
\end{proof}

\begin{theorem}[多级 CRT 阈值的加性控制律：最大素数由总对数预算决定]\label{thm:xi-multistage-crt-max-prime-additive-control}
\leanverified{paper\_xi\_multistage\_crt\_max\_prime\_additive\_control}
给定整数 \(k\ge 1\) 与阈值 \(A_1,\dots,A_k>1\)。
考虑选择两两不交的素数集合 \(S_1,\dots,S_k\) 使得
\[
\prod_{p\in S_i}p\ \ge\ A_i\qquad(1\le i\le k),
\]
并最小化全局最大素数
\[
p_{\max}:=\max\!\Bigl(\bigcup_{i=1}^{k}S_i\Bigr).
\]
记 Chebyshev 函数
\(
\vartheta(x):=\sum_{p\le x}\log p
\)。
则最优值满足精确刻画
\[
\boxed{\;
p_{\max}^{\ast}
\ =\
\inf\!\left\{x\ge 2:\ \vartheta(x)\ge \sum_{i=1}^{k}\log A_i\right\}.
\;}
\]
并且当 \(\sum_{i=1}^{k}\log A_i\to\infty\) 时，由 \(\vartheta(x)\sim x\) 得
\[
p_{\max}^{\ast}\sim \sum_{i=1}^{k}\log A_i.
\]
\end{theorem}
\begin{proof}
必要性：任意可行解满足
\[
\sum_{i=1}^{k}\log A_i
\le
\sum_{i=1}^{k}\sum_{p\in S_i}\log p
=
\sum_{p\in\cup_i S_i}\log p
\le
\sum_{p\le p_{\max}}\log p
=\vartheta(p_{\max}),
\]
故若 \(\vartheta(x)<\sum_i\log A_i\)，则不存在可行解，得到下界。
\par
充分性：取 \(x\) 使 \(\vartheta(x)\ge\sum_i\log A_i\)，并将所有 \(\le x\) 的素数按从小到大排列，依次分配到 \(S_1,S_2,\dots,S_k\)，使得对每个 \(i\) 都在首次达到
\(\sum_{p\in S_i}\log p\ge \log A_i\)
处停止并转入下一组。
由于总和不少于目标总和，该过程在用完素数前必完成全部 \(k\) 个阈值；且构造保证集合两两不交并满足 \(p_{\max}\le x\)。
于是最优值等于所示下确界。
\par
渐近式由 \(\vartheta(x)\sim x\) 立得。证毕。
\end{proof}

\begin{corollary}[多级阈值下的素数轴数量标度]\label{cor:xi-multistage-crt-axis-count-scaling}
\leanverified{paper\_xi\_multistage\_crt\_axis\_count\_scaling}
在定理 \ref{thm:xi-multistage-crt-max-prime-additive-control} 的设定下，令
\[
L:=\sum_{i=1}^{k}\log A_i,
\]
并假设 \(L\to\infty\)。
令 \(K^\ast\) 为达到最优最大素数 \(p_{\max}^\ast\) 所需的总素数轴数量，则
\[
K^\ast\sim \pi(p_{\max}^\ast)\sim \frac{L}{\log L}.
\]
\end{corollary}
\begin{proof}
由定理 \ref{thm:xi-multistage-crt-max-prime-additive-control}，\(p_{\max}^\ast\sim L\)。
而素数轴数量为不超过 \(p_{\max}^\ast\) 的素数个数，即 \(K^\ast=\pi(p_{\max}^\ast)\)。
由素数定理 \(\pi(x)\sim x/\log x\) 得 \(K^\ast\sim L/\log L\)。证毕。
\end{proof}

\begin{theorem}[素数轴分解下的可逆化合成与素数张量积结构]\label{thm:xi-prime-axis-tensorized-replay-composition}
设有两级投影
\[
\Omega\xrightarrow{\ \pi_1\ }X\xrightarrow{\ \pi_2\ }Y,
\]
并分别给定 replay 提升
\[
\Omega\ni\omega\longmapsto\bigl(\pi_1(\omega),r_1(\omega)\bigr),
\qquad
X\ni x\longmapsto\bigl(\pi_2(x),r_2(x)\bigr),
\]
二者均为单射。

把素数集合固定划分为
\[
\mathbb P_{\mathrm{odd}}:=\{p_{2t-1}\}_{t\ge1},
\qquad
\mathbb P_{\mathrm{even}}:=\{p_{2t}\}_{t\ge1}.
\]
定义重编号嵌入
\[
\iota_{\mathrm{odd}}(r)(p_{2t-1})=r(p_t),\ \ \iota_{\mathrm{odd}}(r)(p_{2t})=0,
\]
\[
\iota_{\mathrm{even}}(r)(p_{2t})=r(p_t),\ \ \iota_{\mathrm{even}}(r)(p_{2t-1})=0.
\]
令
\[
R(\omega):=
\iota_{\mathrm{odd}}\!\bigl(r_1(\omega)\bigr)
\oplus
\iota_{\mathrm{even}}\!\bigl(r_2(\pi_1(\omega))\bigr),
\]
其中 \(\oplus\) 为按坐标并合（两项支持互不相交）。
则组合提升
\[
\Omega\ni\omega\longmapsto\bigl(\pi_2(\pi_1(\omega)),R(\omega)\bigr)
\]
为单射；并存在显式重建：先由偶素数子支持恢复 \(\pi_1(\omega)\)，再由奇素数子支持恢复 \(\omega\)。
\end{theorem}

\begin{proof}
设 \(\omega,\omega'\in\Omega\) 满足
\[
\pi_2(\pi_1(\omega))=\pi_2(\pi_1(\omega')),
\qquad
R(\omega)=R(\omega').
\]
取偶素数子支持，得
\[
\iota_{\mathrm{even}}\!\bigl(r_2(\pi_1(\omega))\bigr)
=
\iota_{\mathrm{even}}\!\bigl(r_2(\pi_1(\omega'))\bigr),
\]
故
\[
r_2(\pi_1(\omega))=r_2(\pi_1(\omega')).
\]
结合 \(\pi_2(\pi_1(\omega))=\pi_2(\pi_1(\omega'))\) 与第二级提升单射性，得
\[
\pi_1(\omega)=\pi_1(\omega').
\]
再取奇素数子支持，得
\[
r_1(\omega)=r_1(\omega').
\]
结合第一级提升单射性，推出 \(\omega=\omega'\)。
故组合提升单射，且重建顺序即由上述逆推过程给出。证毕。
\end{proof}
\leanverified{paper\_xi\_prime\_axis\_tensorized\_replay\_composition}

\begin{theorem}[Prime-register 诱导的算法超度量与柱集拓扑]\label{thm:xi-prime-register-ultrametric-cylinder-topology}
\leanverified{paper\_xi\_prime\_register\_ultrametric\_cylinder\_topology}
固定确定性策略 \(S\)，并取一状态类 \(\Omega^{(S)}\subseteq\Omega\)，使映射
\[
r^S:\Omega^{(S)}\to \mathrm{PR}(\mathcal E)
\]
为单射。
定义
\[
d_S(\omega,\omega') :=
\begin{cases}
0, & \omega=\omega',\\[2pt]
2^{-k}, & k=\min\{t\ge1:\ r^S(\omega)(p_t)\neq r^S(\omega')(p_t)\}.
\end{cases}
\]
则：
\begin{enumerate}
  \item \(d_S\) 是超度量；
  \item \((\Omega^{(S)},d_S)\) 的开球恰为 prime-register 有限前缀一致类（柱集），并形成 clopen 基；
  \item 其完备化同胚于
  \[
  \overline{\mathrm{Dec}_{\mathrm{PR}}\!\left(r^S(\Omega^{(S)})\right)}
  \subseteq
  \mathcal E^{\mathbb N}
  \]
  的乘积拓扑闭包。
  若进一步满足前缀满性（任意有限前缀均可达），则完备化同胚于 \(\mathcal E^{\mathbb N}\) 的 Baire 拓扑。
\end{enumerate}
\end{theorem}

\begin{proof}
设 \(k_{ij}\) 为 \(\omega_i,\omega_j\) 的首次差异位。
若 \(\omega_1,\omega_2\) 与 \(\omega_2,\omega_3\) 在前 \(K\) 位一致，则 \(\omega_1,\omega_3\) 亦在前 \(K\) 位一致，故
\[
\min\{k_{12},k_{23}\}\le k_{13}.
\]
于是
\[
d_S(\omega_1,\omega_3)=2^{-k_{13}}
\le
\max\{2^{-k_{12}},2^{-k_{23}}\}
=
\max\{d_S(\omega_1,\omega_2),d_S(\omega_2,\omega_3)\},
\]
得超三角不等式。

对任意 \(K\) 与前缀 \((a_1,\dots,a_K)\)，集合
\[
\mathcal C(a_1,\dots,a_K):=
\{\omega\in\Omega^{(S)}:\ r^S(\omega)(p_t)=a_t,\ 1\le t\le K\}
\]
正是半径 \(2^{-K}\) 级别的球；反向任一球都由某前缀一致条件给出，故柱集即 clopen 基。

把 \(\omega\) 映到序列 \(\mathrm{Dec}_{\mathrm{PR}}(r^S(\omega))\) 得到等距嵌入到 \(\mathcal E^{\mathbb N}\)（把有限序列补零/终止符后的标准嵌入）。
Cauchy 列等价于前缀稳定列，极限对应逐前缀极限序列，故完备化同胚于该像集闭包。
若前缀满性成立，则闭包即全空间 \(\mathcal E^{\mathbb N}\)。证毕。
\end{proof}

\begin{conjecture}[Lee--Yang 椭圆化中不可消去坏素数的算法门槛]\label{conj:xi-leyang-badprime37-algorithmic-threshold}
沿用定理 \ref{thm:xi-terminal-zm-leyang-elliptic-structure} 的四重覆盖 \(\Pi(\lambda,y)=0\) 与既有坏素一致性猜想 \ref{conj:xi-terminal-zm-spectral-compatible-externalization-badprime-minimal}。
在“外置证书轴 + 可审计重建”的语义下，考虑所有对 \(y\)-线去分歧域给出全局一致分支选择的算法族；若其证书轴被限制为 CRT/prime-register 类型的离散残余轴，则该算法族的必要模轴集合在结构上必须包含素数 \(37\)。

等价表述为：\(37\) 是该类全局分支选择算法不可消去的证书素数门槛。
\end{conjecture}

\paragraph{Lee--Yang 算术动力学与时间素数轴的一致化刻画}
在前述椭圆化、分歧结构与 prime-register 语义的共同框架下，可把 Lee--Yang 覆盖、Latt\`es 动力学、三次分歧体与时间素数轴的零点统计组织为一条统一的可核验命题链。

\begin{theorem}[Lee--Yang 倍点 Latt\`es 迭代的固定点分解与动力学 \texorpdfstring{$\zeta$}{zeta} 闭式]\label{thm:xi-time-part9c-leyang-lattes-fix-zeta}
令
\[
E:\quad Y^2=X^3-X+\frac14
\]
并记 \(x:E\to\PP^1\) 为 \(X\)-坐标投影，\([2]:E\to E\) 为倍点映射。设 \(\Phi\) 为 \([2]\) 在 \(x\)-坐标上的诱导有理自映射
\[
\Phi(X)=\frac{X^4+2X^2-2X+1}{4X^3-4X+1}\in\QQ(X).
\]
则对任意 \(n\ge 1\)：
\begin{enumerate}
  \item 固定点集合满足
  \[
  \mathrm{Fix}(\Phi^{\circ n})
  =x\!\bigl(E[2^n-1]\cup E[2^n+1]\bigr).
  \]
  \item 固定点数（不计重数）满足
  \[
  \card{\mathrm{Fix}(\Phi^{\circ n})}=4^n+1,
  \]
  因而 \(a_n:=\card{\mathrm{Fix}(\Phi^{\circ n})}\) 服从递推
  \[
  a_{n+1}=4a_n-3,\qquad a_1=5.
  \]
  \item Artin--Mazur 动力学 \(\zeta\) 函数
  \[
  \zeta_\Phi(t):=\exp\!\left(\sum_{n\ge 1}\frac{\card{\mathrm{Fix}(\Phi^{\circ n})}}{n}t^n\right)
  \]
  具有闭式
  \[
  \zeta_\Phi(t)=\frac{1}{(1-t)(1-4t)}.
  \]
\end{enumerate}
\end{theorem}
\begin{proof}
由定义 \(\Phi^{\circ n}(x(P))=x([2^n]P)\)。故 \(x_0\in\PP^1(\CC)\) 为 \(\Phi^{\circ n}\) 的固定点，当且仅当存在 \(P\in E(\CC)\) 使 \(x(P)=x_0\) 且
\[
x([2^n]P)=x(P).
\]
在椭圆曲线上，对奇阶挠点纤维有 \(x(P)=x(Q)\iff Q=\pm P\)，从而
\[
x([2^n]P)=x(P)
\iff [2^n]P=\pm P
\iff (2^n-1)P=0\ \text{或}\ (2^n+1)P=0.
\]
即得
\(
\mathrm{Fix}(\Phi^{\circ n})=x(E[2^n-1]\cup E[2^n+1])
\)。

又 \(\gcd(2^n-1,2^n+1)=1\)，故
\[
E[2^n-1]\cap E[2^n+1]=\{O\}.
\]
在 \(\CC\) 上 \(\card{E[m]}=m^2\)，故
\[
\card{E[2^n-1]\cup E[2^n+1]}
=(2^n-1)^2+(2^n+1)^2-1
=2\cdot 4^n+1.
\]
除 \(O\) 外，奇阶挠点满足 \(P\neq -P\)，且 \(x\) 把 \(\{P,-P\}\) 识别为一点，故
\[
\card{\mathrm{Fix}(\Phi^{\circ n})}
=1+\frac{(2\cdot 4^n+1)-1}{2}
=4^n+1.
\]
于是
\[
4^{n+1}+1=4(4^n+1)-3
\]
给出递推式。

最后
\[
\log\zeta_\Phi(t)
=\sum_{n\ge 1}\frac{4^n+1}{n}t^n
=\sum_{n\ge 1}\frac{(4t)^n}{n}+\sum_{n\ge 1}\frac{t^n}{n}
=-\log(1-4t)-\log(1-t),
\]
指数化即得
\(
\zeta_\Phi(t)=((1-t)(1-4t))^{-1}
\)。
\end{proof}

\begin{theorem}[Lee--Yang 四层覆盖的几何单值群与满 \texorpdfstring{$S_4$}{S4} 有理特化]\label{thm:xi-time-part9c-leyang-cover-s4-specialization}
沿用
\[
y:=X^2+Y-\frac12\in\QQ(E),\qquad \pi_y:E\to\PP^1_y
\]
所定义的次数 \(4\) 覆盖。设其有限分歧值恰为
\[
y=0,\qquad y=1,\qquad P_{\mathrm{LY}}(y)=0
\]
的三根，且每个有限分歧值之上恰有一个简单分歧点；并设 \(y=\infty\) 为唯一四重极点对应的全分歧点。则：
\begin{enumerate}
  \item \(\pi_y\) 的几何单值群同构于 \(S_4\)；
  \item 存在无穷多个 \(y_0\in\QQ\)（避开分歧值）使四次特化 \(\Pi(\lambda,y_0)\in\QQ[\lambda]\) 的伽罗瓦群仍同构于 \(S_4\)。
\end{enumerate}
\end{theorem}
\begin{proof}
第一条即定理 \ref{thm:xi-terminal-zm-leyang-monodromy-s4} 的结论。
第二条由 Hilbert 不可约定理对该正则 \(S_4\)-扩张的直接应用得到。
\end{proof}

\begin{theorem}[\texorpdfstring{$S_4$}{S4} 正则专门化诱导的二维 Artin 表示与权一模性]\label{thm:xi-time-part9c-leyang-s4-specialization-artin-weight1}
沿用定理 \ref{thm:xi-time-part9c-leyang-cover-s4-specialization} 的记号，并取满足其第二条的某个 \(y_0\in\QQ\)。
令 \(K_{y_0}/\QQ\) 为四次多项式 \(\Pi(\lambda,y_0)\) 的分裂域，并设 \(\Gal(K_{y_0}/\QQ)\cong S_4\)。
令 \(\rho_2\) 为 \(S_4\) 的二维不可约复表示，取其为复合表示
\[
S_4\twoheadrightarrow S_4/V_4\cong S_3\xrightarrow{\ \rho_{\mathrm{std}}\ }\GL_2(\CC),
\]
其中 \(V_4\) 为 Klein 四元正规子群、\(\rho_{\mathrm{std}}\) 为 \(S_3\) 的标准二维表示。
则由自然满射 \(\varphi_{y_0}:\Gal(\overline{\QQ}/\QQ)\twoheadrightarrow \Gal(K_{y_0}/\QQ)\cong S_4\) 得到二维 Artin 表示
\[
\varrho_{y_0}:=\rho_2\circ \varphi_{y_0}:\ \Gal(\overline{\QQ}/\QQ)\longrightarrow \GL_2(\CC),
\]
其像为可解群。
若进一步满足 \(\Pi(\lambda,y_0)\) 具有非实根（等价于 \(K_{y_0}\) 至少有一处复嵌入），则 \(\varrho_{y_0}\) 为奇表示；从而存在权 \(1\) 本原尖点新形式 \(f_{y_0}\) 使
\[
L(s,f_{y_0})=L(s,\varrho_{y_0}).
\]
并且对任意不分歧素数 \(p\)，其 Hecke 本征值满足
\[
a_p(f_{y_0})=\tr\bigl(\varrho_{y_0}(\mathrm{Frob}_p)\bigr),
\]
从而由 \(\Pi(\lambda,y_0)\bmod p\) 的分解型（即 \(\mathrm{Frob}_p\) 在 \(S_4\) 中的共轭类）确定。
\end{theorem}
\begin{proof}
由构造，\(\varrho_{y_0}\) 的像包含于 \(\rho_2(S_4)=\rho_{\mathrm{std}}(S_3)\)，故为可解群。
当 \(\Pi(\lambda,y_0)\) 具有非实根时，复共轭在 \(\Gal(K_{y_0}/\QQ)\) 中为某个阶 \(2\) 元素；其在 \(S_3\) 商上的像为换位，从而 \(\det(\varrho_{y_0}(c))=-1\)，即 \(\varrho_{y_0}\) 为奇表示。
由 Deligne--Serre 与 Langlands--Tunnell，二维奇且像为可解群的 Artin 表示对应于权 \(1\) 本原尖点模形式，并给出 \(L\)-函数恒等式。
对不分歧素数，Hecke 本征值与 Frobenius 迹的等同性为 Artin 模性对应的标准结论；而 \(\mathrm{Frob}_p\) 的共轭类由 \(\Pi(\lambda,y_0)\bmod p\) 的分解型确定。证毕。
\end{proof}

\begin{theorem}[Lee--Yang 三次分歧体的判别式分解、整基指数与权一模形式级别]\label{thm:xi-time-part9c-leyang-cubic-artin-weight1}
令
\[
P_{\mathrm{LY}}(y)=256y^3+411y^2+165y+32\in\ZZ[y].
\]
设 \(y_{\mathrm{LY}}\) 为其任一根，\(K:=\QQ(y_{\mathrm{LY}})\)。并记
\[
\Disc(P_{\mathrm{LY}})=-3^9\cdot 31^2\cdot 37,\qquad
\Disc(K)=-999=-3^3\cdot 37.
\]
则成立：
\begin{enumerate}
  \item 整基指数满足
  \[
  [\mathcal O_K:\ZZ[y_{\mathrm{LY}}]]=3^3\cdot 31=837.
  \]
  因而 \(31\) 在 \(K\) 中不分歧，其出现仅来自非整基指数；\(K\) 的分歧素数恰为 \(\{3,37\}\)。
  \item 设 \(L\) 为 \(K\) 的伽罗瓦闭包，则 \(\Gal(L/\QQ)\cong S_3\)。令 \(\rho_{\mathrm{std}}\) 为 \(S_3\) 的标准二维不可约表示，则
  \[
  L(s,\rho_{\mathrm{std}})=\frac{\zeta_K(s)}{\zeta(s)}.
  \]
  \item \(\rho_{\mathrm{std}}\) 的 Artin 导子为 \(999\)，且
  \[
  \det(\rho_{\mathrm{std}})=\chi_{-111},
  \]
  其中 \(\chi_{-111}\) 对应二次分辨域 \(\QQ(\sqrt{-111})\)。
  因而存在权 \(1\) 本原尖点新形式
  \[
  f_{\mathrm{LY}}\in S_1(\Gamma_0(999),\chi_{-111}),
  \qquad
  L(s,f_{\mathrm{LY}})=L(s,\rho_{\mathrm{std}}).
  \]
  对任意 \(p\nmid 999\)，有
  \[
  a_p(f_{\mathrm{LY}})=\operatorname{tr}\bigl(\rho_{\mathrm{std}}(\mathrm{Frob}_p)\bigr)\in\{2,0,-1\},
  \]
  并由 \(P_{\mathrm{LY}}\bmod p\) 的分解型决定：
  \[
  (1,1,1)\Rightarrow a_p=2,\qquad
  (1,2)\Rightarrow a_p=0,\qquad
  (3)\Rightarrow a_p=-1.
  \]
\end{enumerate}
在 \(p\le 191\) 且 \(p\nmid 999\) 的样本上，可取
\begin{align*}
&a_{73}=a_{89}=a_{131}=a_{139}=a_{191}=2,\\
&a_{5}=a_{7}=a_{17}=a_{23}=a_{29}=a_{59}=a_{67}=a_{113}=a_{127}=a_{151}=a_{157}=a_{167}=a_{179}=a_{181}=-1,\\
&a_{11}=a_{13}=a_{19}=a_{41}=a_{43}=a_{47}=a_{53}=a_{61}=a_{71}=a_{79}=a_{83}=a_{97}=a_{101}=a_{103}=a_{107}=a_{109}=a_{137}=a_{149}=a_{163}=a_{173}=0.
\end{align*}
\end{theorem}
\begin{proof}
对任意三次数域生成元 \(\theta\) 有标准恒等式
\[
\Disc(\ZZ[\theta])=\Disc(K)\,[\mathcal O_K:\ZZ[\theta]]^2.
\]
取 \(\theta=y_{\mathrm{LY}}\) 并代入
\(
\Disc(\ZZ[y_{\mathrm{LY}}])=\Disc(P_{\mathrm{LY}})
\)
与 \(\Disc(K)=-999\)，得
\[
[\mathcal O_K:\ZZ[y_{\mathrm{LY}}]]^2
=\frac{3^9\cdot 31^2\cdot 37}{3^3\cdot 37}
=3^6\cdot 31^2
=(3^3\cdot 31)^2.
\]
故指数为 \(3^3\cdot 31=837\)。分歧素数由 \(\Disc(K)\) 的素因子给出，故恰为 \(\{3,37\}\)。

对第二条，设 \(H\le S_3\) 为稳定一根的子群，\([S_3:H]=3\)，且 \(K=L^H\)。
置换表示 \(V:=\Ind_H^{S_3}\mathbf 1\) 分解为
\[
V\cong \mathbf 1\oplus \rho_{\mathrm{std}}.
\]
Artin \(L\)-函数对直和可乘，故
\[
\zeta_K(s)=L(s,V)=L(s,\mathbf 1)\,L(s,\rho_{\mathrm{std}})
=\zeta(s)\,L(s,\rho_{\mathrm{std}}),
\]
即
\(
L(s,\rho_{\mathrm{std}})=\zeta_K(s)/\zeta(s)
\)。

第三条中，\(\det(\rho_{\mathrm{std}})\) 为 \(S_3\) 的符号特征，故对应唯一二次子域 \(\QQ(\sqrt{-111})\)，即 \(\chi_{-111}\)。
并且三次非伽罗瓦扩张的标准二维 Artin 表示导子等于 \(|\Disc(K)|=999\)。
由 Deligne--Serre 与 Langlands--Tunnell，对应得到权 \(1\) 新形式
\(
f_{\mathrm{LY}}\in S_1(\Gamma_0(999),\chi_{-111})
\) 并满足 \(L\)-函数恒等。
当 \(p\nmid 999\) 时，\(\mathrm{Frob}_p\) 的共轭类由 \(P_{\mathrm{LY}}\bmod p\) 分解型决定；标准表示在单位元、换位、三循环上的迹分别为 \(2,0,-1\)，故得到所述取值判据与样本。
\end{proof}

\begin{theorem}[\texorpdfstring{$\dd y$}{dy} 的显式分解与分子函数的精确除子]\label{thm:xi-leyang-ed-dy-explicit-divisor}
\leanverified{paper\_xi\_leyang\_ed\_dy\_explicit\_divisor}
在仿射坐标 \((m,u)\) 下考虑椭圆曲线
\[
E_D/\QQ:\quad u^2=-4m^3-16m^2+16m+65,
\]
并定义有理函数
\[
y
:=\frac{u-(4m^2-17)}{2(m-2)(m+2)^2}\in\QQ(E_D).
\]
令 \(\omega:=\dd m/u\) 为 \(E_D\) 的不变微分。
则成立严格恒等式
\[
\dd y=\frac{F}{2(m-2)^2(m+2)^3}\,\omega,
\]
其中
\[
F=\bigl(4m^3-8m^2-19m+34\bigr)u+\bigl(6m^4+24m^3-48m^2-99m+98\bigr)\in\QQ[m,u].
\]
并且将 \(F\) 视为 \(E_D\) 上的有理函数，其除子满足
\[
\operatorname{div}(F)
=2(2,-1)+3(-2,-1)+(-4,-1)+\sum_{i=1}^{3}R_{g,i}-9O,
\]
其中 \(O\) 为无穷远点，\(R_{g,i}\in E_D(\overline\QQ)\) 为映射 \(y:E_D\to\PP^1\) 在三条 Lee--Yang 分支值之上的三处有限分歧点。
\end{theorem}
\begin{proof}
对 \(y=\frac{u-(4m^2-17)}{2(m-2)(m+2)^2}\) 作全导数，并以曲线方程隐式微分
\[
2u\,\dd u=\bigl(-12m^2-32m+16\bigr)\dd m
\]
消去 \(\dd u\)，即可化简得到所示分解式 \(\dd y=\frac{F}{2(m-2)^2(m+2)^3}\omega\)。
\par
为刻画 \(F=0\) 的零点与重数，注意对任意点满足 \(u^2=f(m)\)（其中 \(f(m)=-4m^3-16m^2+16m+65\)），有消元恒等式
\begin{equation}\label{eq:xi-leyang-ed-F-elimination}
\bigl(6m^4+24m^3-48m^2-99m+98\bigr)^2
-f(m)\bigl(4m^3-8m^2-19m+34\bigr)^2
=4(m-2)^2(m+2)^3(m+4)\,h(m),
\end{equation}
其中
\(
h(m)=16m^3-87m^2+186m-128
\)。
因此 \(F=0\) 在 \(E_D\) 上的解由 \(m=2\)（重数 \(2\)）、\(m=-2\)（重数 \(3\)）、\(m=-4\)（重数 \(1\)）及 \(h(m)=0\) 的三根（各重数 \(1\)）完全给出；相应的点坐标由 \(u^2=f(m)\) 唯一确定到 \(\pm\) 号分支。
\par
另一方面，由 \(F\) 的次数估计（或以局部参数在 \(O\) 处展开）可得 \(F\) 在 \(O\) 处具有九阶极点，从而其零点总次数（计重数）亦为 \(9\)。
结合 \eqref{eq:xi-leyang-ed-F-elimination} 所给出的全部有限零点计数与重数，得到所述除子分解，其中 \(h(m)=0\) 所对应的三点即为三条 Lee--Yang 分支值之上的分歧点 \(R_{g,i}\)。
\end{proof}

\begin{theorem}[Lee--Yang 三个分歧点由一条抛物线截出]\label{thm:xi-leyang-ed-parabola-cuts-ramification}
在定理 \ref{thm:xi-leyang-ed-dy-explicit-divisor} 的设定下，定义
\[
\Phi:=u-8m^2+29m-31\in\QQ(E_D).
\]
则
\[
\operatorname{div}(\Phi)=Q_0'+\sum_{i=1}^{3}R_{g,i}-4O,
\qquad
Q_0'=\Bigl(\frac74,\frac{19}4\Bigr)\in E_D(\QQ).
\]
等价地，仿射二次曲线 \(u=8m^2-29m+31\) 与三次曲线 \(E_D\) 的交点（按重数计）恰为 \(Q_0'\) 与三个 \(R_{g,i}\)，其余交数集中在无穷远点 \(O\)（交重数为 \(4\)）。
\end{theorem}
\begin{proof}
由 \(\Phi=0\) 得 \(u=8m^2-29m+31\)。
代回 \(E_D\) 的方程并化简，得到
\[
(8m^2-29m+31)^2-(-4m^3-16m^2+16m+65)=(4m-7)\,h(m),
\]
其中 \(h(m)=16m^3-87m^2+186m-128\) 与 \eqref{eq:xi-leyang-ed-F-elimination} 相同。
故 \(\Phi=0\) 在 \(E_D\) 上的仿射解由 \(m=\frac74\)（对应 \(Q_0'\)）与 \(h(m)=0\) 的三根（对应三点 \(R_{g,i}\)）构成。
\par
又 \(\Phi\) 的主部由 \(-8m^2\) 主导，故在 \(O\) 处具有四阶极点，从而除子为
\(\operatorname{div}(\Phi)=Q_0'+\sum_i R_{g,i}-4O\)。
\end{proof}

\begin{corollary}[Lee--Yang 三个分歧点的迹为有理点]\label{cor:xi-leyang-ed-trace-of-ramification-is-rational}
记
\[
Q_0=\Bigl(\frac74,-\frac{19}4\Bigr)\in E_D(\QQ),
\qquad
Q_0'=-Q_0.
\]
则在 \(E_D\) 的群律下有
\[
R_{g,1}+R_{g,2}+R_{g,3}=Q_0.
\]
等价地，度为 \(3\) 的有效除子 \(D_g:=\sum_{i=1}^3R_{g,i}\) 满足线性等价
\[
D_g\sim Q_0+2O.
\]
\end{corollary}
\begin{proof}
由定理 \ref{thm:xi-leyang-ed-parabola-cuts-ramification} 得
\(
D_g+Q_0'-4O\sim 0
\)。
另一方面，函数 \(m-\frac74\) 的除子为
\[
\operatorname{div}\Bigl(m-\frac74\Bigr)=Q_0+Q_0'-2O,
\]
故 \(Q_0'\sim 2O-Q_0\)。
代回得到 \(D_g\sim Q_0+2O\)。
在椭圆曲线同构 \(\mathrm{Pic}^0(E_D)\simeq E_D\) 下，度零除子类对应于群和，遂得 \(R_{g,1}+R_{g,2}+R_{g,3}=Q_0\)。
\end{proof}

\begin{theorem}[Lee--Yang 的非分歧同伴点与另一条抛物截面]\label{thm:xi-leyang-ed-parabola-cuts-unramified-companions}
设 \(U_{g,i}\in E_D(\overline\QQ)\)（\(i=1,2,3\)）为满足 \(P_{\mathrm{LY}}(y)=0\) 且在 \(y\) 上不分歧的三点（即每条 Lee--Yang 分支值之上除去分歧点后的剩余点）。
则它们满足如下两条显式描述。
\begin{enumerate}
  \item \(U_{g,i}\) 的 \(m\)-坐标是三次多项式
  \[
  \tilde h(m):=4m^3-15m^2+132m+148
  \]
  的三根。
  \item 定义二次截面函数
  \[
  \Psi:=u+m^2-\frac{19}{4}m+\frac12\in\QQ(E_D),
  \]
  则
  \[
  \operatorname{div}(\Psi)=Q_0'+\sum_{i=1}^3U_{g,i}-4O.
  \]
  等价地，仿射二次曲线 \(u=-m^2+\frac{19}{4}m-\frac12\) 与 \(E_D\) 的交点（按重数计）恰由 \(Q_0'\) 与三点 \(U_{g,i}\) 给出。
\end{enumerate}
并且有判别式恒等式
\[
\operatorname{disc}(\tilde h)=\operatorname{disc}(P_{\mathrm{LY}})=-2^6\cdot 3^3\cdot 31^2\cdot 37,
\]
从而 \(\tilde h\) 的分裂域与 Lee--Yang 三次域一致。
\end{theorem}
\begin{proof}
对 \(\Psi=0\) 解得 \(u=-m^2+\frac{19}{4}m-\frac12\)。
代回 \(E_D\) 并化简可得
\[
\Bigl(-m^2+\frac{19}{4}m-\frac12\Bigr)^2-(-4m^3-16m^2+16m+65)=\frac1{16}(4m-7)\,\tilde h(m),
\]
从而 \(\Psi=0\) 的仿射解由 \(m=\frac74\)（对应 \(Q_0'\)）与 \(\tilde h(m)=0\) 的三根（对应三点）组成。
由于 \(y:E_D\to\PP^1\) 为次数 \(3\) 的映射，而每条 Lee--Yang 分支值纤维型为 \((2,1)\)，故在每条分支值之上除分歧点外恰有一个剩余点；三条分支值总计给出三点，必与上述三点一致，即为 \(\{U_{g,1},U_{g,2},U_{g,3}\}\)。
\par
判别式恒等式与分裂域一致性由直接计算给出。
\end{proof}

\begin{proposition}[同一分支纤维中的群律约束]\label{prop:xi-leyang-ed-fiber-group-law-constraint}
设 \(y_0\in\overline\QQ\) 为 \(y\) 的任一分支值，且纤维型为 \((2,1)\)。
记分歧点为 \(R_{y_0}\)，同纤维剩余点为 \(U_{y_0}\)。
则在 \(E_D\) 的群律下恒有
\[
2R_{y_0}+U_{y_0}=Q_0,
\qquad
Q_0=\Bigl(\frac74,-\frac{19}4\Bigr).
\]
特别地，对三条 Lee--Yang 分支值满足
\[
U_{g,i}=Q_0-2R_{g,i}\quad(i=1,2,3),
\qquad
\sum_{i=1}^3U_{g,i}=Q_0.
\]
\end{proposition}
\begin{proof}
函数 \(y-y_0\) 与 \(y\) 具有相同极点除子。
记 \(P_-:=(-2,1)\)、\(P_+:=(2,1)\)，则由 \(y\) 的极点结构可取
\(
(y)_\infty=2P_-+P_+
\)。
另一方面，\(y-y_0\) 在纤维上零点除子为 \(2R_{y_0}+U_{y_0}\)，故
\[
\operatorname{div}(y-y_0)=2R_{y_0}+U_{y_0}-(2P_-+P_+).
\]
由主除子度零与 \(\mathrm{Pic}^0(E_D)\simeq E_D\) 的识别，得到群律恒等式
\(
2R_{y_0}+U_{y_0}=2P_-+P_+
\)。
又由既定有理点数据可核验 \(2P_-+P_+=Q_0\)，从而得到第一式。
对三条 Lee--Yang 分支值逐点代入并求和，结合推论 \ref{cor:xi-leyang-ed-trace-of-ramification-is-rational}，即得最后两式。
\end{proof}

\begin{theorem}[翻倍--平移对应的 \texorpdfstring{$m$}{m}-坐标有理提升与 \texorpdfstring{$h\mapsto \tilde h$}{h->htilde} 的显式恒等式]\label{thm:xi-leyang-ed-double-translate-varphi-h-to-htilde}
沿用定理 \ref{thm:xi-leyang-ed-parabola-cuts-ramification} 的分歧点 \(R_{g,i}\) 与命题 \ref{prop:xi-leyang-ed-fiber-group-law-constraint} 的群律约束。
对任意满足
\[
u=8m^2-29m+31,\qquad h(m)=0
\]
的分歧点 \(R=(m,u)=R_{g,i}\)，定义同纤维未分歧同伴点
\[
U:=Q_0-2R\in E_{\mathrm D}(\overline\QQ),
\qquad
Q_0=\left(\frac74,-\frac{19}{4}\right).
\]
则 \(U\) 的 \(m\)-坐标由 \(m\) 的一个显式有理函数 \(\varphi(m)\) 给出，并且 \(\varphi\) 与 \(\tilde h\) 之间满足严格代数恒等式：
% Auto-generated by scripts/exp_xi_ed_h_to_htilde_varphi_audit.py
\[
h(m)=16m^3-87m^2+186m-128,\qquad \tilde h(m)=4m^3-15m^2+132m+148.
\]
\[
\varphi(m)=\frac{9824488844415016105 m^{2} - 44675850650870527750 m + 41624375182502178688}{1420238224376141276 m^{2} - 3136065797377302248 m + 637207999589057024}.
\]
\[
\Theta(m)=515399676328725406597427896646032 m^{3} - 3106151691496474514988931964936439 m^{2} + 5911701997465751789317317241415802 m - 3698028848970378941250154426183808.
\]
\[
\tilde h\bigl(\varphi(m)\bigr)=\frac{579237207545457760272384\,\Theta(m)\,h(m)}{2864729308668308675877656138721721055151754409239472576 m^{6} - 18977054939368374677243458195835754861366342033979138944 m^{5} + 45759631807736146060786116929096274837252472691986377984 m^{4} - 47871516386496307287039786938396501599226136228675029504 m^{3} + 20530642638454232143126621527438620223524794967372169216 m^{2} - 3820048346760370730704956504283420216193616154453868544 m + 258728135041932407896792521915285771003530228818509824}.
\]

从而 \(h(m)=0\Rightarrow \tilde h(\varphi(m))=0\)，并且对每个 \(i\) 有
\[
m(U_{g,i})=\varphi\bigl(m(R_{g,i})\bigr).
\]
\end{theorem}
\begin{proof}
由命题 \ref{prop:xi-leyang-ed-fiber-group-law-constraint}，对每条 Lee--Yang 分歧值之上的分歧点 \(R_{g,i}\) 与未分歧点 \(U_{g,i}\) 恒有
\(
U_{g,i}=Q_0-2R_{g,i}
\)。
将椭圆曲线的加法/倍点公式写为 \((m,u)\)-坐标的有理表达式，并在约束 \(u=8m^2-29m+31,\ h(m)=0\) 下化简，
即可得到 \(m(U)\) 作为 \(m\) 的有理函数 \(\varphi(m)\)；
进一步把 \(\varphi\) 代入 \(\tilde h\) 并清分母，即得到所示恒等式。
上述化简与恒等式核验可由相应脚本逐项复现。证毕。
\end{proof}

\begin{corollary}[\texorpdfstring{$2$}{2}-进判别式指数的跃迁]\label{cor:xi-leyang-ed-disc-tildeh-vs-h}
\leanverified{paper\_xi\_leyang\_ed\_disc\_tildeh\_vs\_h}
沿用 \(\Disc(h)=-2^{2}\cdot 3^{3}\cdot 31^{2}\cdot 37\)（定理 \ref{thm:xi-leyang-branchfield-equals-hm-splittingfield}）与
\(\Disc(\tilde h)=-2^{6}\cdot 3^{3}\cdot 31^{2}\cdot 37\)（定理 \ref{thm:xi-leyang-ed-parabola-cuts-unramified-companions}），则
\[
\Disc(\tilde h)=2^{4}\,\Disc(h).
\]
\end{corollary}
\begin{proof}
两条判别式分解式相除即得。证毕。
\end{proof}

\begin{theorem}[两条 Tschirnhaus 消元恒等式锁定 Lee--Yang 三次]\label{thm:xi-leyang-ed-tschirnhaus-resultants-lock-PLY}
令
\[
P_{\mathrm{LY}}(y)=256y^3+411y^2+165y+32\in\ZZ[y],
\]
并令
\[
h(m):=16m^3-87m^2+186m-128,\qquad
\tilde h(m):=4m^3-15m^2+132m+148.
\]
则在 \(\QQ[y]\) 中成立两条结式恒等式（\(\operatorname{Res}_m\) 表示关于变量 \(m\) 的结式）：
\begin{enumerate}
  \item
  \[
  \operatorname{Res}_m\!\left(
  h(m),\ 2(m-2)(m+2)^2\,y-(4m^2-29m+48)
  \right)
  =-714432\cdot P_{\mathrm{LY}}(y).
  \]
  \item
  \[
  \operatorname{Res}_m\!\left(
  \tilde h(m),\ 8(m-2)(m+2)^2\,y+(20m^2-19m-66)
  \right)
  =-33226752\cdot P_{\mathrm{LY}}(y).
  \]
\end{enumerate}
因此 Lee--Yang 三次域可由两套完全显式的有理参数化同时生成：
\[
y=\frac{4m^2-29m+48}{2(m-2)(m+2)^2}\quad\text{其中 }h(m)=0;
\qquad
y=\frac{-20m^2+19m+66}{8(m-2)(m+2)^2}\quad\text{其中 }\tilde h(m)=0.
\]
\leanverified{paper\_xi\_leyang\_ed\_tschirnhaus\_resultants\_lock\_ply}
\end{theorem}
\begin{proof}
对给定二元方程组分别对变量 \(m\) 取结式消元即可得到一元多项式。
所得结式非零且次数为 \(3\)，并且其零点集合与 \(P_{\mathrm{LY}}(y)\) 的三根一致，故二者在 \(\QQ[y]\) 中必相差一个非零常数因子。
该常数可由首项系数（或在任取一点处代入比较）确定，从而得到所示恒等式。
\end{proof}


\begin{conjecture}[Newman--Lee--Yang 双谱 Artin 参数的无交叉性]\label{conj:xi-newman-leyang-artin-no-cross}
设 \(L_{\mathrm{New}}/\QQ\) 为定理 \ref{thm:sync-kernel-weighted-newman-Q-galois-s37} 的次数 \(37\) Newman 多项式 \(Q\) 的分裂域（故 \(\Gal(L_{\mathrm{New}}/\QQ)\cong S_{37}\)），并设 \(L_{\mathrm{LY}}/\QQ\) 为定理 \ref{thm:xi-time-part9c-leyang-cubic-artin-weight1} 的 Lee--Yang 三次分歧体 \(K\) 的伽罗瓦闭包（故 \(\Gal(L_{\mathrm{LY}}/\QQ)\cong S_3\)）。
令 \(\rho\) 为 \(\Gal(L_{\mathrm{New}}/\QQ)\) 的任意非平凡不可约复表示，并把其视为 \(\Gal(\overline\QQ/\QQ)\) 的有限像 Artin 表示；
令 \(\rho_{\mathrm{std}}\) 为 \(\Gal(L_{\mathrm{LY}}/\QQ)\) 的标准二维 Artin 表示（定理 \ref{thm:xi-time-part9c-leyang-cubic-artin-weight1}）。
猜想在 Artin/Langlands 自守提升存在的前提下，Rankin--Selberg 卷积
\[
L(s,\rho\otimes\rho_{\mathrm{std}})
\]
不出现由共同二次特征所诱导的系统性因子化或异常极点；其解析与算术复杂度由两表示的有限分歧素集合的并所控制，而不通过同一二次子域产生结构性耦合。
该猜想可通过分歧素集合计算、局部因子比较与数值零点统计把判据操作化。
\end{conjecture}

\begin{theorem}[时间素数轴上的 Gödel--Lee--Yang 圆律]\label{thm:xi-time-part9c-godel-leyang-circle}
固定整数 \(m\ge 2\)，令 \(p_t\) 为第 \(t\) 个素数。对任意长度 \(T\) 的时间词
\(
a=(a_1,\dots,a_T)\in\{0,1,\dots,m-1\}^T
\)，定义
\[
G_T(a):=\prod_{t=1}^{T}p_t^{a_t},
\qquad
Z_T(\beta):=\sum_{a\in\{0,\dots,m-1\}^T}G_T(a)^{-\beta}.
\]
则：
\begin{enumerate}
  \item 有乘法分解
  \[
  Z_T(\beta)=\prod_{t=1}^{T}\sum_{k=0}^{m-1}p_t^{-\beta k}
  =\prod_{t=1}^{T}\frac{1-p_t^{-m\beta}}{1-p_t^{-\beta}}.
  \]
  \item 全部零点均位于虚轴，且
  \[
  Z_T(\beta)=0
  \iff
  \exists\,t\in\{1,\dots,T\},\ \exists\,j\in\{1,\dots,m-1\},\ \exists\,\ell\in\ZZ:
  \quad
  \beta=\frac{2\pi i\,(j+m\ell)}{m\log p_t}.
  \]
  特别地，任一零点都满足 \(\operatorname{Re}(\beta)=0\)。
\end{enumerate}
\end{theorem}
\begin{proof}
由定义直接分离变量：
\[
Z_T(\beta)
=\sum_{a_1=0}^{m-1}\cdots\sum_{a_T=0}^{m-1}\prod_{t=1}^Tp_t^{-\beta a_t}
=\prod_{t=1}^{T}\sum_{k=0}^{m-1}p_t^{-\beta k},
\]
即得第一式。
对固定 \(t\)，令 \(z:=p_t^{-\beta}\)，则
\[
\sum_{k=0}^{m-1}z^k=\frac{1-z^m}{1-z}.
\]
其零点恰为 \(z^m=1\) 且 \(z\neq 1\)，即
\[
z=e^{2\pi i(j+m\ell)/m},
\qquad
j=1,\dots,m-1,\ \ell\in\ZZ.
\]
回代 \(z=p_t^{-\beta}\) 即得零点显式表达。
由 \(|z|=1\) 立得 \(\operatorname{Re}(\beta)=0\)。
\end{proof}

\begin{theorem}[Gödel--Lee--Yang 零点计数的 \texorpdfstring{$T\log T$}{TlogT} 增长律]\label{thm:xi-time-part9c-godel-leyang-zero-count}
沿用定理 \ref{thm:xi-time-part9c-godel-leyang-circle} 的记号。对固定 \(H>0\)，定义（按重数计）
\[
N_T(H):=\#\{\beta\in\CC:\ Z_T(\beta)=0,\ |\operatorname{Im}\beta|\le H\}.
\]
则当 \(T\to\infty\) 时有
\[
N_T(H)\sim \frac{(m-1)H}{\pi}\sum_{t=1}^{T}\log p_t.
\]
从而
\[
N_T(H)\sim \frac{(m-1)H}{\pi}\,T\log T.
\]
\end{theorem}
\begin{proof}
由定理 \ref{thm:xi-time-part9c-godel-leyang-circle}，每个因子对应的零点为
\[
\beta_{t,j,\ell}
=\frac{2\pi i\,(j+m\ell)}{m\log p_t},
\qquad
j=1,\dots,m-1,\ \ell\in\ZZ.
\]
按重数计数可逐因子求和。令
\[
B_t:=\frac{Hm\log p_t}{2\pi}.
\]
条件 \(|\operatorname{Im}\beta_{t,j,\ell}|\le H\) 等价于
\[
|j+m\ell|\le B_t.
\]
固定 \(j\) 时满足该不等式的整数 \(\ell\) 个数为
\(
\frac{2B_t}{m}+O(1)
\)。
再对 \(j=1,\dots,m-1\) 求和，得到第 \(t\) 个因子的贡献
\[
\frac{2(m-1)}{m}B_t+O(1)
=\frac{(m-1)H}{\pi}\log p_t+O(1).
\]
对 \(t=1,\dots,T\) 求和：
\[
N_T(H)
=\frac{(m-1)H}{\pi}\sum_{t=1}^{T}\log p_t+O(T).
\]
故得首个渐近式。又
\[
\sum_{t=1}^{T}\log p_t=\vartheta(p_T)\sim p_T\sim T\log T,
\]
遂得
\(
N_T(H)\sim \frac{(m-1)H}{\pi}T\log T
\)。
\end{proof}

\begin{theorem}[圆维--素数账本编码的一致误差信息下界]\label{thm:xi-time-part9c-cdim-prime-ledger-info-lb}
设 \(K\) 为紧致交换 Lie 群，记
\[
r:=\cdim(K)
\]
为其环面维数（等于连通部分维数），\(d\) 为与拓扑相容的度量。
给定 \(\varepsilon\in(0,\frac12)\)，考虑任意编码--解码对
\[
\mathrm{enc}:K\to\{0,1\}^b\times(\ZZ/P\ZZ),\qquad
\mathrm{dec}:\{0,1\}^b\times(\ZZ/P\ZZ)\to K,
\]
其中 \(b\in\NN\)、\(P\) 为奇素数。若满足一致误差
\[
d\bigl(x,\mathrm{dec}(\mathrm{enc}(x))\bigr)\le\varepsilon
\qquad(\forall x\in K),
\]
则存在仅依赖于 \((K,d)\) 的常数 \(C_{K,d}\) 使
\[
b+\log_2P\ \ge\ r\log_2\frac1\varepsilon\ -\ C_{K,d}.
\]
特别地，在有限生成交换群的离散词度量球口径中，取 \(\varepsilon\asymp N^{-1}\) 得
\[
b+\log_2P\ \ge\ r\log_2N-C,
\]
即前缀--素数账本守恒型下界。
\end{theorem}
\begin{proof}
取 \(S\subset K\) 为一组极大 \(2\varepsilon\)-分离集，即任意不同 \(x,y\in S\) 满足 \(d(x,y)>2\varepsilon\)。
由紧致 \(r\)-维 Lie 群的标准打包估计，存在常数 \(c_{K,d}>0\) 使
\[
\card{S}\ge c_{K,d}\,\varepsilon^{-r}.
\]

若存在 \(x\neq y\in S\) 使 \(\mathrm{enc}(x)=\mathrm{enc}(y)\)，则两者解码点相同，故
\[
d(x,y)
\le d\bigl(x,\mathrm{dec}(\mathrm{enc}(x))\bigr)
+d\bigl(\mathrm{dec}(\mathrm{enc}(y)),y\bigr)
\le 2\varepsilon,
\]
与 \(S\) 的 \(2\varepsilon\)-分离性矛盾。
故 \(\mathrm{enc}\) 在 \(S\) 上必单射，进而编码空间大小必须满足
\[
2^bP\ge \card{S}\ge c_{K,d}\,\varepsilon^{-r}.
\]
取二进制对数得
\[
b+\log_2P
\ge r\log_2\frac1\varepsilon-\log_2\frac1{c_{K,d}}.
\]
令 \(C_{K,d}:=\log_2(c_{K,d}^{-1})\) 即证。
\end{proof}

\paragraph{Clark 密度下 Jensen 缺陷的反熵桥与离线证书化}
在既有 Jensen 缺陷账本与 Poisson--KL 耗散链之间，以下给出一条 Schur--Carath\'eodory 专化的精确桥接：
先把 Jensen 缺陷分解为反向相对熵与严格负项之和，再把反向相对熵压缩为 Fourier 能量与有限 Toeplitz 数据的可计算上界。

\begin{theorem}[Clark--Jensen 反熵分解恒等式]\label{thm:xi-time-part9d-clark-jensen-reversekl-identity}
设 \(f\) 为单位圆盘 \(\mathbb D\) 上的 Schur 函数，满足 \(f(0)=0\)。定义
\[
F(z):=1-f(z),\qquad
C(z):=\frac{1+f(z)}{1-f(z)}.
\]
则 \(C\) 为 Carath\'eodory 函数，\(\Re C(z)>0\)；并存在唯一概率测度 \(\mu\) 于 \(\TT\) 使 \(C\) 满足 Herglotz 表示。
对任意 \(0<r<1\)，记
\[
p_r(\theta):=\Re C(re^{i\theta}),\qquad
m(\dd\theta):=\frac{\dd\theta}{2\pi}.
\]
则 \(p_r>0\)、\(\int_{\TT}p_r\,\dd m=1\)，且 Jensen 缺陷
\[
D_F(r):=\int_{\TT}\log|F(re^{i\theta})|\,\dd m(\theta)-\log|F(0)|
\]
满足精确恒等式
\begin{equation}\label{eq:xi-time-part9d-clark-jensen-identity}
\boxed{\;
D_F(r)
=\frac12\,D_{\mathrm{KL}}\!\bigl(m\mid p_rm\bigr)
+\frac12\int_{\TT}\log\!\bigl(1-|f(re^{i\theta})|^2\bigr)\,\dd m(\theta)\;}
\end{equation}
其中
\[
D_{\mathrm{KL}}\!\bigl(m\mid p_rm\bigr)
:=\int_{\TT}\log\!\frac{1}{p_r(\theta)}\,\dd m(\theta).
\]
特别地，
\begin{equation}\label{eq:xi-time-part9d-clark-jensen-upper-bridge}
0\le D_F(r)\le \frac12\,D_{\mathrm{KL}}\!\bigl(m\mid p_rm\bigr),\qquad 0<r<1.
\end{equation}
\end{theorem}
\begin{proof}
由 \(\Re C(z)>0\) 立得 \(p_r(\theta)>0\)。
又 \(\Re C\) 为调和函数，圆周平均等于中心值，故
\[
\int_{\TT}p_r(\theta)\,\dd m(\theta)=\Re C(0)=1.
\]
由代数恒等式
\[
\Re\frac{1+f}{1-f}=\frac{1-|f|^2}{|1-f|^2}
\]
可得
\[
p_r(\theta)=\frac{1-|f(re^{i\theta})|^2}{|1-f(re^{i\theta})|^2}.
\]
取对数并整理：
\[
\log|1-f(re^{i\theta})|
=\frac12\log\!\bigl(1-|f(re^{i\theta})|^2\bigr)-\frac12\log p_r(\theta).
\]
对 \(\theta\) 以 \(m\) 平均，并利用 \(f(0)=0\Rightarrow F(0)=1\Rightarrow \log|F(0)|=0\)，即得
\eqref{eq:xi-time-part9d-clark-jensen-identity}。

又 \(0<r<1\) 时 \(f(re^{i\theta})\) 在紧圆周上满足 \(|f|<1\)，故
\(
\log(1-|f|^2)\le 0
\)。
结合 Jensen 不等式给出的 \(D_F(r)\ge0\)，得到
\eqref{eq:xi-time-part9d-clark-jensen-upper-bridge}。
\end{proof}

\begin{theorem}[反向 KL 的显式 \texorpdfstring{$L^2$}{L2} 控制与 Fourier 能量恒等式]\label{thm:xi-time-part9d-reversekl-l2-fourier}
设 \(C\) 为 Carath\'eodory 函数且 \(C(0)=1\)，写作
\[
C(z)=1+2\sum_{n\ge1}c_n z^n.
\]
记
\[
p_r(\theta):=\Re C(re^{i\theta}),\qquad 0<r<1.
\]
则：
\begin{enumerate}
  \item 对一切 \(\theta\) 有一致下界
  \[
  p_r(\theta)\ge a_r:=\frac{1-r}{1+r}.
  \]
  \item 反向相对熵满足显式上界
  \[
  D_{\mathrm{KL}}\!\bigl(m\mid p_rm\bigr)
  \le \frac{1}{2a_r^2}\,\|p_r-1\|_{L^2(m)}^2.
  \]
  \item Fourier 能量恒等式
  \[
  \|p_r-1\|_{L^2(m)}^2
  =2\sum_{n\ge1}|c_n|^2r^{2n}.
  \]
\end{enumerate}
从而
\begin{equation}\label{eq:xi-time-part9d-reversekl-fourier-explicit}
D_{\mathrm{KL}}\!\bigl(m\mid p_rm\bigr)
\le \left(\frac{1+r}{1-r}\right)^2\sum_{n\ge1}|c_n|^2r^{2n}.
\end{equation}
\end{theorem}
\begin{proof}
由 Herglotz 表示可写 \(p_r=P_r*\mu\)，其中 \(\mu\) 为概率测度。Poisson 核最小值为 \((1-r)/(1+r)\)，故 \(p_r\ge a_r\)。

令 \(\phi(x):=-\log x\)。在区间 \([a_r,\infty)\) 上有 \(\phi''(x)=1/x^2\le 1/a_r^2\)。
围绕 \(x=1\) 作二阶 Taylor 估计，得对任意 \(x\ge a_r\),
\[
-\log x\le (1-x)+\frac{(x-1)^2}{2a_r^2}.
\]
取 \(x=p_r(\theta)\) 并积分，利用 \(\int_{\TT}p_r\,\dd m=1\)，得到
\[
D_{\mathrm{KL}}(m\mid p_rm)=\int_{\TT}-\log p_r\,\dd m
\le \frac{1}{2a_r^2}\int_{\TT}(p_r-1)^2\,\dd m.
\]
即第二条。

再由
\[
p_r(\theta)
=1+\sum_{n\ge1}\bigl(c_nr^n e^{in\theta}+\overline{c_n}r^n e^{-in\theta}\bigr),
\]
其 Fourier 系数为 \(a_n=c_nr^n\)（\(n\ge1\)）与 \(a_{-n}=\overline{c_n}r^n\)。
Parseval 恒等式给出
\[
\|p_r-1\|_{L^2(m)}^2
=\sum_{n\ne0}|a_n|^2
=2\sum_{n\ge1}|c_n|^2r^{2n}.
\]
合并即得 \eqref{eq:xi-time-part9d-reversekl-fourier-explicit}。
\end{proof}

\begin{corollary}[Jensen 缺陷的 Carath\'eodory 系数上界证书]\label{cor:xi-time-part9d-jensen-defect-coefficient-certificate}
在定理 \ref{thm:xi-time-part9d-clark-jensen-reversekl-identity} 的设定下，
令
\[
\frac{1+f(z)}{1-f(z)}=1+2\sum_{n\ge1}c_n z^n.
\]
则对任意 \(0<r<1\) 有
\begin{equation}\label{eq:xi-time-part9d-jensen-defect-coefficient-certificate}
D_{1-f}(r)
\le \frac12\left(\frac{1+r}{1-r}\right)^2\sum_{n\ge1}|c_n|^2r^{2n}.
\end{equation}
\end{corollary}
\begin{proof}
由 \eqref{eq:xi-time-part9d-clark-jensen-upper-bridge} 与
\eqref{eq:xi-time-part9d-reversekl-fourier-explicit} 直接代入即得。
\end{proof}

\leanverified{paper\_xi\_time\_part9d\_finite\_toeplitz\_certificate}
\begin{theorem}[有限 Toeplitz 数据的反熵/缺陷上界与误差预算]\label{thm:xi-time-part9d-finite-toeplitz-certificate}
沿用定理 \ref{thm:xi-time-part9d-reversekl-l2-fourier} 的记号，并令
\[
E_N(r):=\sum_{n=1}^{N}|c_n|^2r^{2n}\qquad(N\ge1).
\]
则对任意 \(0<r<1\) 有有限数据上界
\begin{equation}\label{eq:xi-time-part9d-finite-toeplitz-certificate-main}
D_{1-f}(r)
\le \frac12\left(\frac{1+r}{1-r}\right)^2
\left(E_N(r)+\frac{r^{2(N+1)}}{1-r^2}\right).
\end{equation}
并且给定 \(\varepsilon\in(0,1)\)，若
\begin{equation}\label{eq:xi-time-part9d-finite-toeplitz-N-choice}
N\ \ge\ \frac{1}{1-r}\log\frac{1}{\varepsilon},
\end{equation}
则
\[
\frac{r^{2(N+1)}}{1-r^2}\le \frac{\varepsilon^2}{1-r^2},
\]
从而
\begin{equation}\label{eq:xi-time-part9d-finite-toeplitz-certificate-eps}
D_{1-f}(r)
\le \frac12\left(\frac{1+r}{1-r}\right)^2E_N(r)
+\frac12\left(\frac{1+r}{1-r}\right)^2\frac{\varepsilon^2}{1-r^2}.
\end{equation}
\end{theorem}
\begin{proof}
由推论 \ref{cor:xi-time-part9d-jensen-defect-coefficient-certificate},
\[
D_{1-f}(r)\le \frac12\left(\frac{1+r}{1-r}\right)^2
\sum_{n\ge1}|c_n|^2r^{2n}.
\]
拆分前 \(N\) 项与尾项，并利用 Herglotz 系数满足 \(|c_n|\le1\)，得
\[
\sum_{n\ge N+1}|c_n|^2r^{2n}
\le \sum_{n\ge N+1}r^{2n}
=\frac{r^{2(N+1)}}{1-r^2},
\]
即得 \eqref{eq:xi-time-part9d-finite-toeplitz-certificate-main}。

再由 \(\log r\le -(1-r)\)（\(0<r<1\)），得到
\[
r^{2(N+1)}
=\exp\!\bigl(2(N+1)\log r\bigr)
\le \exp\!\bigl(-2(N+1)(1-r)\bigr).
\]
若 \(N\) 满足 \eqref{eq:xi-time-part9d-finite-toeplitz-N-choice}，则上式 \(\le \varepsilon^2\)，故得
\eqref{eq:xi-time-part9d-finite-toeplitz-certificate-eps}。
\end{proof}

\begin{proposition}[模 \(4\) Fourier 正交通道与 \texorpdfstring{\(\chi_4\)}{chi4} 能量审计]\label{prop:xi-time-part9d-mod4-fourier-channel-energy}
在定理 \ref{thm:xi-time-part9d-reversekl-l2-fourier} 的设定下，定义
\[
\mathcal H_j:=\overline{\mathrm{span}}\{e^{in\theta}:n\in\ZZ,\ n\equiv j\!\!\!\pmod 4\},
\qquad j=0,1,2,3,
\]
并令 \(\mathsf P_j:L^2(\TT,m)\to\mathcal H_j\) 为正交投影。则：
\begin{enumerate}
  \item 正交分解
  \[
  \|p_r-1\|_{L^2(m)}^2=\sum_{j=0}^3\|\mathsf P_j(p_r-1)\|_{L^2(m)}^2.
  \]
  \item 通道能量显式公式
  \[
  \|\mathsf P_0(p_r-1)\|_{L^2(m)}^2
  =2\sum_{k\ge1}|c_{4k}|^2r^{8k},
  \]
  \[
  \|\mathsf P_2(p_r-1)\|_{L^2(m)}^2
  =2\sum_{k\ge0}|c_{4k+2}|^2r^{2(4k+2)},
  \]
  \[
  \|\mathsf P_1(p_r-1)\|_{L^2(m)}^2
  =\|\mathsf P_3(p_r-1)\|_{L^2(m)}^2
  =\sum_{k\ge0}|c_{4k+1}|^2r^{2(4k+1)}
  +\sum_{k\ge0}|c_{4k+3}|^2r^{2(4k+3)}.
  \]
  \item 定义 \(\chi_4\)-通道能量
  \[
  \mathcal E_0(r):=\sum_{k\ge1}|c_{4k}|^2r^{8k},\qquad
  \mathcal E_j(r):=\sum_{k\ge0}|c_{4k+j}|^2r^{2(4k+j)}\ (j=1,2,3),
  \]
  则 Jensen 缺陷满足通道化证书
  \[
  D_{1-f}(r)
  \le \frac12\left(\frac{1+r}{1-r}\right)^2\sum_{j=0}^3\mathcal E_j(r)
  =\frac14\left(\frac{1+r}{1-r}\right)^2\sum_{j=0}^3\|\mathsf P_j(p_r-1)\|_{L^2(m)}^2.
  \]
\end{enumerate}
\end{proposition}
\begin{proof}
第一条由 \(L^2(\TT,m)=\bigoplus_{j=0}^3\mathcal H_j\) 的正交直和分解得到。
第二条由 Fourier 系数
\(
\widehat{(p_r-1)}(n)=c_nr^n\ (n\ge1),\ 
\widehat{(p_r-1)}(-n)=\overline{c_n}r^n
\)
和 Parseval 恒等式逐类求和即得。
第三条由推论 \ref{cor:xi-time-part9d-jensen-defect-coefficient-certificate}
与 \(\sum_{n\ge1}|c_n|^2r^{2n}=\sum_{j=0}^3\mathcal E_j(r)\) 直接得到。
\end{proof}

\begin{remark}[与 Poisson--KL 耗散链的接口]
命题 \ref{prop:xi-time-part9d-mod4-fourier-channel-energy} 给出纯离线可计算的 Jensen 缺陷上界；
若结合猜想 \ref{conj:xi-poisson-kl-exact-alignment} 或
\eqref{eq:xi-golden-chain-alignment-inequality} 的对齐接口，
即可把 \(\chi_4\)-通道能量预算直接接入
结论 \ref{con:xi-poisson-second-order-two-closed-constants} 与
结论 \ref{con:xi-poisson-moment-vanishing-kl-step-closed-form}
的 Poisson--KL 耗散常数链。
\end{remark}

\paragraph{递推 Hankel 反演与窗口 Lie 闭合的代数审计链}
下述陈述均在本文已给出的定义、记号与编号体系内推导；其与外部文献的同一性或新颖性比对不作为本段的审计目标。

\begin{theorem}[递推参数化的 Jacobian--Hankel 恒等式]\label{thm:xi-time-part9e-recurrence-jacobian-hankel}
设 \(d\ge 2\)，\(k\) 为特征 \(0\) 的交换域，初值
\[
(a_0,\dots,a_{d-1})\in k^d
\]
固定。对参数
\(
c=(c_0,\dots,c_{d-1})\in k^d
\)
定义递推序列
\[
a_n(c):=-\sum_{j=0}^{d-1}c_j\,a_{n-d+j}(c),\qquad n\ge d.
\]
记
\[
b_n(c):=a_{d+n}(c)\quad(0\le n\le d-1),\qquad
\Phi_{d,2d-1}(c):=(b_0(c),\dots,b_{d-1}(c)),
\]
以及主 Hankel 块
\[
H(c):=\bigl(a_{i+j}(c)\bigr)_{0\le i,j\le d-1}.
\]
则有严格恒等式
\[
\det\!\bigl(D\Phi_{d,2d-1}(c)\bigr)=(-1)^d\det H(c).
\]
\leanverified{paper\_xi\_time\_part9e\_recurrence\_jacobian\_hankel}
\end{theorem}
\begin{proof}
把 \(b=(b_0,\dots,b_{d-1})\) 视为独立变量，并令 \(a_n(c,b)\) 由同一递推延拓，且
\(
a_{d+n}(c,b)=b_n
\)
（\(0\le n\le d-1\)）。
定义
\[
F_n(c,b):=a_{n+d}(c,b)+\sum_{j=0}^{d-1}c_j\,a_{n+j}(c,b),\qquad 0\le n\le d-1.
\]
由构造有 \(F(c,\Phi_{d,2d-1}(c))=0\)。
隐函数求导给出
\[
D\Phi_{d,2d-1}(c)
=-\bigl(D_bF(c,b)\bigr)^{-1}D_cF(c,b)\big|_{b=\Phi_{d,2d-1}(c)}.
\]
其中 \(D_cF=H(c)\)；
\(D_bF\) 为单位下三角矩阵，故 \(\det(D_bF)=1\)。
两边取行列式即得
\[
\det\!\bigl(D\Phi_{d,2d-1}(c)\bigr)
=(-1)^d\det H(c).
\]
\end{proof}

\begin{corollary}[递推系数的显式有理反演与双有理性]\label{cor:xi-time-part9e-recurrence-rational-inversion}
沿用定理 \ref{thm:xi-time-part9e-recurrence-jacobian-hankel} 的记号。
设
\[
b=\Phi_{d,2d-1}(c)=(b_0,\dots,b_{d-1}).
\]
定义矩阵 \(\widehat H(b)=(\widehat h_{ij}(b))_{0\le i,j\le d-1}\) 为
\[
\widehat h_{ij}(b):=
\begin{cases}
a_{i+j}, & i+j\le d-1,\\[2mm]
b_{i+j-d}, & d\le i+j\le 2d-2,
\end{cases}
\]
并记列向量 \(\mathbf b:=(b_0,\dots,b_{d-1})^{\mathsf T}\)。
若 \(\det H(c)\neq 0\)（等价地 \(\det\widehat H(b)\neq 0\)），则
\[
c=-\,\widehat H(b)^{-1}\mathbf b,
\]
从而 \(\Phi_{d,2d-1}\) 在开集 \(\{\det H\neq 0\}\) 上为双有理同构，其分歧由 \(\det H=0\) 切出。
\end{corollary}
\begin{proof}
递推关系在 \(n=0,\dots,d-1\) 上可写为
\[
\sum_{j=0}^{d-1}a_{n+j}(c)\,c_j=-a_{n+d}(c)=-b_n.
\]
系数矩阵即 \(\widehat H(b)\)，右端即 \(-\mathbf b\)。
当 \(\det\widehat H(b)\neq 0\) 时解唯一，得到所示反演公式。
由 \(\widehat H(b)\) 的条目只依赖固定初值与 \(b_0,\dots,b_{d-2}\)，该反演为有理映射；与 \(c\mapsto b\) 互逆即得双有理性。
分歧结论由定理 \ref{thm:xi-time-part9e-recurrence-jacobian-hankel} 直接推出。
\end{proof}

\begin{theorem}[全主 Hankel 行列式的消去闭式与正权符号律]\label{thm:xi-time-part9e-hankel-resultant-closed-form}
\leanverified{paper\_xi\_time\_part9e\_hankel\_resultant\_closed\_form}
设 \(y_+(z)\) 为次数 \(\kappa+1\) 的首一多项式，根为互异实数 \(b_0,\dots,b_\kappa\)。
定义
\[
c_m:=\frac{y(b_m)}{y_+'(b_m)},\qquad
\mu_n:=\sum_{m=0}^{\kappa}c_m b_m^n\quad(n\ge 0),
\]
并令
\[
H_\kappa:=(\mu_{i+j})_{0\le i,j\le \kappa}.
\]
则成立
\[
\det(H_\kappa)=(-1)^{\kappa(\kappa+1)/2}\Res(y,y_+).
\]
若进一步 \(c_m>0\)（正权情形），则 \(H_\kappa\) 正定，因而
\[
(-1)^{\kappa(\kappa+1)/2}\Res(y,y_+)>0.
\]
\end{theorem}
\begin{proof}
令 Vandermonde 矩阵
\[
V=(b_m^{\,i})_{0\le i,m\le \kappa},\qquad
C=\mathrm{diag}(c_0,\dots,c_\kappa).
\]
由矩定义有 \(H_\kappa=VCV^{\mathsf T}\)，从而
\[
\det(H_\kappa)=(\det V)^2\prod_{m=0}^{\kappa}c_m.
\]
又
\[
\prod_{m=0}^{\kappa}c_m
=\frac{\prod_m y(b_m)}{\prod_m y_+'(b_m)}
=\frac{\Res(y,y_+)}{\prod_m y_+'(b_m)}.
\]
利用恒等式
\[
(\det V)^2=\Disc(y_+),\qquad
\prod_m y_+'(b_m)=(-1)^{\kappa(\kappa+1)/2}\Disc(y_+),
\]
消去 \(\Disc(y_+)\) 即得主式。
当 \(c_m>0\) 时，\(H_\kappa\) 为正权 Gram 矩阵，故正定，行列式为正，符号律随之成立。
\end{proof}

\begin{theorem}[移位 Hankel 线性笔的节点对角化与 \texorpdfstring{\(H_\kappa\)}{Hk}-自伴随]\label{thm:xi-time-part9e-shifted-hankel-pencil-diagonalization}
在定理 \ref{thm:xi-time-part9e-hankel-resultant-closed-form} 的设定下，定义
\[
H_\kappa^+:=(\mu_{i+j+1})_{0\le i,j\le \kappa},\qquad
A_\kappa:=H_\kappa^{-1}H_\kappa^+.
\]
则有相似对角化
\[
A_\kappa=(V^{\mathsf T})^{-1}\mathrm{diag}(b_0,\dots,b_\kappa)V^{\mathsf T},
\]
故
\[
\chi_{A_\kappa}(z)=\prod_{m=0}^{\kappa}(z-b_m)=y_+(z).
\]
此外，
\[
A_\kappa^{\mathsf T}H_\kappa=H_\kappa A_\kappa.
\]
\end{theorem}
\begin{proof}
由 \(\mu_{i+j+1}=\sum_m c_m b_m^{i+j+1}\) 得
\[
H_\kappa^+=V\,\mathrm{diag}(c_0b_0,\dots,c_\kappa b_\kappa)\,V^{\mathsf T}.
\]
代入 \(A_\kappa=H_\kappa^{-1}H_\kappa^+\) 并用 \(H_\kappa=VCV^{\mathsf T}\) 直接化简，得到所示对角化公式与特征多项式结论。
又因 \(H_\kappa,H_\kappa^+\) 对称，
\[
A_\kappa^{\mathsf T}H_\kappa
=(H_\kappa^{-1}H_\kappa^+)^{\mathsf T}H_\kappa
=H_\kappa^+
=H_\kappa(H_\kappa^{-1}H_\kappa^+)
=H_\kappa A_\kappa.
\]
\end{proof}

\begin{theorem}[Fibonacci 幂 Hankel 行列式的完全乘积公式]\label{thm:xi-time-part9e-fibonacci-power-hankel-product}
\leanverified{paper\_xi\_time\_part9e\_fibonacci\_power\_hankel\_product}
对整数 \(q\ge 1\)，定义
\[
H_q:=\bigl(F_{i+j+3}^{\,q}\bigr)_{0\le i,j\le q},
\]
其中 \((F_n)\) 为 Fibonacci 数列。
则
\[
\det(H_q)=
\left(\prod_{i=0}^{q}\binom{q}{i}\right)
\left(\prod_{k=1}^{q}F_k^{\,2(q+1-k)}\right).
\]
右端为整数乘积闭式。
\end{theorem}
\begin{proof}
由 Binet 表达 \(F_{n+3}=\alpha\varphi^n+\beta\psi^n\) 得
\[
F_{n+3}^{\,q}
=\sum_{i=0}^{q}\binom{q}{i}\alpha^{q-i}\beta^i
\bigl(\varphi^{q-i}\psi^i\bigr)^n
=\sum_{i=0}^{q}c_i r_i^{\,n}.
\]
故 \(H_q\) 可分解为 Vandermonde 型乘积 \(H_q=V^{\mathsf T}CV\)，于是
\[
\det(H_q)=(\det V)^2\prod_{i=0}^{q}c_i.
\]
把 \((\det V)^2\) 按节点差显式展开，并使用
\[
1-(-\varphi^{-2})^k=\sqrt5\,F_k/\varphi^k
\]
逐项化简，可将所有代数扩张因子与 \(\prod_i c_i\) 中对应项精确抵消，最终得到所示整数乘积式。
\end{proof}

\begin{proposition}[\texorpdfstring{\(\det(H_q)\)}{det(Hq)} 的 \texorpdfstring{\(p\)}{p}-进赋值闭式与坏素数有限支撑]\label{prop:xi-time-part9e-fibonacci-hankel-padic-valuation}
沿用定理 \ref{thm:xi-time-part9e-fibonacci-power-hankel-product}，对任意素数 \(p\) 有
\[
v_p\!\bigl(\det(H_q)\bigr)
=\sum_{i=0}^{q}v_p\!\binom{q}{i}
+\sum_{k=1}^{q}2(q+1-k)\,v_p(F_k).
\]
特别地，素因子支撑包含于有限集合
\[
\{p:\exists\,i,\ p\mid \tbinom{q}{i}\}
\cup
\{p:\exists\,k,\ p\mid F_k\}.
\]
\end{proposition}
\begin{proof}
对定理 \ref{thm:xi-time-part9e-fibonacci-power-hankel-product} 的乘积公式逐项取 \(v_p\)，并用 \(v_p(xy)=v_p(x)+v_p(y)\) 即得。
支撑有限性由固定 \(q\) 时参与因子仅有限个直接推出。
\end{proof}

\begin{theorem}[Smith 计数律：核与余核在 \texorpdfstring{\(p^k\)}{p^k} 尺度上的精确公式]\label{thm:xi-time-part9e-smith-kernel-cokernel-count}
\leanverified{paper\_xi\_time\_part9e\_smith\_kernel\_cokernel\_count}
设 \(H_0\in M_d(\ZZ)\) 满秩，Smith 正规形为
\[
U H_0 V=\mathrm{diag}(s_1,\dots,s_d),
\qquad
U,V\in\GL_d(\ZZ),\quad s_i\mid s_{i+1}.
\]
固定素数 \(p\)，定义
\[
\mathcal H_p(k):=\sum_{i=1}^{d}\min\{v_p(s_i),k\},\qquad k\ge 1.
\]
令 \(\overline H_0:=H_0\bmod p^k\) 诱导映射
\[
\overline H_0:(\ZZ/p^k\ZZ)^d\to(\ZZ/p^k\ZZ)^d.
\]
则有
\[
|\ker(\overline H_0)|=p^{\mathcal H_p(k)},
\qquad
|\mathrm{coker}(\overline H_0)|=p^{\mathcal H_p(k)},
\]
并且
\[
\lim_{k\to\infty}\mathcal H_p(k)=v_p(\det H_0).
\]
\end{theorem}
\begin{proof}
因 \(U,V\) 在 \(\ZZ/p^k\ZZ\) 上仍可逆，核与余核基数可化归到对角映射
\[
(x_1,\dots,x_d)\mapsto(s_1x_1,\dots,s_dx_d).
\]
第 \(i\) 坐标方程 \(s_ix_i\equiv 0\pmod{p^k}\) 的解数为 \(p^{\min(v_p(s_i),k)}\)，故核基数为其乘积，即 \(p^{\mathcal H_p(k)}\)。
对有限环上满秩方阵，核与余核基数相同，故余核同式。
\(\mathcal H_p(k)\) 单调并在大 \(k\) 稳定到 \(\sum_i v_p(s_i)=v_p(\det H_0)\)。
\end{proof}

\begin{theorem}[模素数审计下 Hankel 正规形的可定义性与秩保持]\label{thm:xi-time-part9e-hankel-modp-rank-preservation}
设 \((a_n)_{n\ge 0}\subset\ZZ\) 在 \(\QQ\) 上的 Hankel 秩为 \(d\)。
取偏移 \(\ell\ge 0\)，令
\[
H_\ell:=\bigl(a_{\ell+i+j}\bigr)_{0\le i,j\le d-1}\in M_d(\ZZ),
\qquad
\Delta_\ell:=\det(H_\ell)\neq 0.
\]
对素数 \(p\) 定义 \(\bar a_n:=a_n\bmod p\in\FF_p\) 与 \(\bar H_\ell:=H_\ell\bmod p\)。
则：
\begin{enumerate}
  \item 若 \(p\nmid \Delta_\ell\)，则 \(\bar H_\ell\in\GL_d(\FF_p)\)，故正规形恢复（例如 \(\bar A:=\bar H_\ell^{-1}\bar H_{\ell+1}\)）在 \(\FF_p\) 上可定义，且 \((\bar a_n)\) 的 Hankel 秩仍为 \(d\)；
  \item 坏素数集合 \(\{p:\ p\mid \Delta_\ell\}\) 有限，因此除有限多个 \(p\) 外，模 \(p\) 的 Hankel 秩保持且正规形恢复始终可定义。
\end{enumerate}
\end{theorem}
\begin{proof}
当 \(p\nmid\Delta_\ell\) 时，\(\bar H_\ell\) 可逆，直接给出 \(\mathrm{rank}_{\FF_p}H(\bar a)\ge d\)。
另一方面，矩阵 \(\bar A=\bar H_\ell^{-1}\bar H_{\ell+1}\) 提供维数 \(d\) 的线性实现，故 \(\mathrm{rank}_{\FF_p}H(\bar a)\le d\)。
于是秩恰为 \(d\)，且恢复映射可定义。
坏素数有限性由非零整数 \(\Delta_\ell\) 仅有有限素因子得到。
\end{proof}

\begin{theorem}[窗口 \texorpdfstring{\(m=6\)}{m=6} 的 Lie 闭合达到最大正交与特殊线性]\label{thm:xi-time-part9e-window6-lie-maximal-closure}
令 \(m=6\)，则 \(|X_6|=F_8=21\)。记 \(V:=\RR^{X_6}\)，并令 \(L_1,\dots,L_6\) 为窗口 \(6\) 折叠/投影结构诱导的整数线性算子。
定义
\[
\mathfrak g_6:=\Lie\langle L_1,\dots,L_6\rangle\subset\mathfrak{gl}(V),
\qquad
\mathfrak k_6:=[\mathfrak g_6,\mathfrak g_6].
\]
则
\[
\mathfrak k_6=\mathfrak{so}(V)\cong\mathfrak{so}(21),
\qquad
\mathfrak g_6=\mathfrak{sl}(V)\cong\mathfrak{sl}(21).
\]
\end{theorem}
\begin{proof}
由构造，\(L_i\) 迹为零，故 \(\mathfrak g_6\subset\mathfrak{sl}(V)\)，于是
\[
\dim\mathfrak g_6\le \dim\mathfrak{sl}(21)=440.
\]
同时由对称性，\([L_i,L_j]\) 落在反对称端，故 \(\mathfrak k_6\subset\mathfrak{so}(V)\)，于是
\[
\dim\mathfrak k_6\le \dim\mathfrak{so}(21)=210.
\]
窗口 \(6\) 的模素数审计脚本给出可复算维数证书：
% AUTO-GENERATED by scripts/exp_window6_lie_envelope.py
\begin{equation}\label{eq:window6_lie_envelope_invariants}
\begin{aligned}
\dim\mathfrak{g}_6&=440,\\
\dim[\mathfrak{g}_6,\mathfrak{g}_6]&=440,\\
\dim Z(\mathfrak{g}_6)&=0,\\
\mathrm{rank}\,\kappa_6&=440,\\
\dim\,\mathrm{Comm}(\mathrm{ad}(\mathfrak{g}_6))&=1.
\end{aligned}
\end{equation}

% AUTO-GENERATED by scripts/exp_window6_lie_envelope_orthogonal_part.py
\begin{equation}\label{eq:window6_lie_envelope_orthogonal_part_dim}
\dim\mathfrak{k}_6=210.
\end{equation}

因此 \(\dim\mathfrak g_6=440\)、\(\dim\mathfrak k_6=210\)。
结合上界即得
\(
\mathfrak g_6=\mathfrak{sl}(21)
\)
与
\(
\mathfrak k_6=\mathfrak{so}(21)
\)。
\end{proof}

\begin{conjecture}[窗口族的普遍最大正交闭合]\label{conj:xi-time-part9e-window-general-maximal-closure}
对任意 \(m\ge 6\)，令 \(V_m:=\RR^{X_m}\)（\(|X_m|=F_{m+2}\)），并定义
\[
\mathfrak g_m:=\Lie\langle L_1,\dots,L_m\rangle\subset\mathfrak{gl}(V_m),
\qquad
\mathfrak k_m:=[\mathfrak g_m,\mathfrak g_m].
\]
猜想存在阈值 \(m_0\)（经验上 \(m_0=6\)），使对一切 \(m\ge m_0\)：
\[
\mathfrak k_m=\mathfrak{so}(V_m)\cong\mathfrak{so}(F_{m+2}),
\qquad
\mathfrak g_m=\mathfrak{sl}(V_m)\cong\mathfrak{sl}(F_{m+2}).
\]
若该猜想成立，则局部 Zeckendorf 约束诱导生成元在 Lie 意义下将给出普遍最大简单闭包，为后续谱分解、算术审计与几何分层提供统一刚性骨架。
\end{conjecture}

\paragraph{折叠纤维不可压缩性的编码、状态与步数统一律}
以下陈述在本文既定符号与接口公理下成立。
默认 \(m\ge 1\)，\(\Omega_m=\{0,1\}^m\)，\(\Fold_m:\Omega_m\twoheadrightarrow X_m\)，
\[
d_m(x):=\bigl|\Fold_m^{-1}(x)\bigr|,\qquad x\in X_m.
\]

\begin{theorem}[熵账本诱导的侧信息期望码长下界]\label{thm:xi-time-part9f-sideinfo-expected-length-lower-bound}
\leanverified{paper\_xi\_time\_part9f\_sideinfo\_expected\_length\_lower\_bound}
取有限字母表 \(\Sigma\) 且 \(|\Sigma|=M\ge 2\)。
设
\[
\Psi_m:\Omega_m\to X_m\times\Sigma^{\ast},\qquad
\Psi_m(\omega)=\bigl(\Fold_m(\omega),R(\omega)\bigr),
\]
满足：
\begin{enumerate}
  \item \(\Psi_m\) 为单射；
  \item 对每个 \(x\in X_m\)，码字族
  \(
  \{R(\omega):\Fold_m(\omega)=x\}
  \)
  在给定 \(x\) 的条件解码意义下为前缀码。
\end{enumerate}
则对任意 \(\mu\in\Delta(\Omega_m)\)，记 \(\pi:=(\Fold_m)_{\ast}\mu\)，并定义纤维均匀提升
\[
\mu^{\mathrm e}(\omega):=
\frac{\pi(x)}{d_m(x)},
\qquad x=\Fold_m(\omega),
\]
有
\[
\mathbb E_{\omega\sim\mu}\!\bigl[|R(\omega)|\bigr]
\ge
\frac{H(\mu)-H(\pi)}{\log M}
=
\frac{\mathbb E_{X\sim\pi}\!\bigl[\log d_m(X)\bigr]-D_{\mathrm{KL}}(\mu\Vert\mu^{\mathrm e})}{\log M}.
\]
\end{theorem}
\begin{proof}
固定 \(x\in X_m\)，令 \(l(\omega):=|R(\omega)|\)。
由条件前缀码的 Kraft 不等式与标准熵--码长下界，
\[
\mathbb E\!\bigl[l(\omega)\mid X=x\bigr]
\ge
\frac{H(\mu(\cdot\mid X=x))}{\log M}.
\]
对 \(x\sim\pi\) 取期望得到
\[
\mathbb E[l(\omega)]
\ge
\frac{H(\mu\mid X)}{\log M}
=
\frac{H(\mu)-H(\pi)}{\log M}.
\]
另一方面，
\[
\log\mu^{\mathrm e}(\omega)=\log\pi(X)-\log d_m(X),
\]
故
\[
D_{\mathrm{KL}}(\mu\Vert\mu^{\mathrm e})
=
\mathbb E_{\mu}\!\left[\log\frac{\mu(\omega)}{\mu^{\mathrm e}(\omega)}\right]
=
-H(\mu)+H(\pi)+\mathbb E_{\pi}\!\bigl[\log d_m(X)\bigr].
\]
整理即得第二个等式。
\end{proof}

\begin{corollary}[可逆外置平均码长的线性下界]\label{cor:xi-time-part9f-no-constant-memory}
\leanverified{paper\_xi\_time\_part9f\_no\_constant\_memory}
沿用定理 \ref{thm:xi-time-part9f-sideinfo-expected-length-lower-bound}。
取均匀微态基线 \(\mu=u_{\Omega_m}\)，则 \(H(\mu)=m\log 2\)。
若 \(|X_m|=F_{m+2}\)，则
\[
\mathbb E_{\omega\sim u_{\Omega_m}}|R(\omega)|
\ge
\frac{m\log 2-\log F_{m+2}}{\log M}
=
\frac{m\log(2/\varphi)+O(1)}{\log M},
\]
其中 \(\varphi\) 为黄金比例。
\end{corollary}
\begin{proof}
由定理 \ref{thm:xi-time-part9f-sideinfo-expected-length-lower-bound} 第一式，
\[
\mathbb E|R|
\ge
\frac{H(\mu)-H(\pi)}{\log M}.
\]
又 \(H(\pi)\le \log|X_m|=\log F_{m+2}\)，代入即得第一不等式。
由 \(F_{m+2}=\Theta(\varphi^m)\) 得到渐近式。
\end{proof}

\begin{theorem}[算法信息坍缩的纤维对数体积上界与指数尾]\label{thm:xi-time-part9f-kolmogorov-fiber-saturation}
设存在可计算外置
\[
\widetilde F_{m,S}(\omega):=
\bigl(\Fold_m(\omega),\,r_{m,S}(\omega)\bigr),
\]
以及可计算恢复映射 \(R_{m,S}\) 使
\(
R_{m,S}\circ \widetilde F_{m,S}=\mathrm{id}_{\Omega_m}
\)。
则存在常数 \(c_S\)（仅依赖模板 \(S\)）使对所有 \(\omega\in\Omega_m\)：
\[
K\!\bigl(\omega\mid X\bigr)
\le
\left\lceil\log_2 d_m(X)\right\rceil+c_S,
\qquad
X:=\Fold_m(\omega).
\]
并且在 \(\omega\sim\mathrm{Unif}(\Omega_m)\) 下，对任意 \(t\ge 0\) 有
\[
\mathbb P\!\left(
K(\omega\mid X)\le \log_2 d_m(X)-t
\right)
\le
2^{-t+c_S}.
\]
\end{theorem}
\begin{proof}
由可计算恢复，给定 \((X,r_{m,S}(\omega))\) 可重构 \(\omega\)，故
\[
K(\omega\mid X)\le K(r_{m,S}(\omega)\mid X)+O_S(1).
\]
对固定 \(X=x\)，可区分候选数最多 \(d_m(x)\)，于是
\[
K(r_{m,S}(\omega)\mid X=x)\le \lceil \log_2 d_m(x)\rceil+O_S(1),
\]
得到上界。
对尾界，固定 \(x\) 时满足
\(
K(\omega\mid x)\le \log_2 d_m(x)-t
\)
的对象个数至多 \(2^{\log_2 d_m(x)-t+O_S(1)}\)，在纤维均匀分布下概率至多 \(2^{-t+O_S(1)}\)。
再对 \(x\) 按 \(\mathrm{Unif}(\Omega_m)\) 诱导分布平均，得到全局界。
\end{proof}

\begin{proposition}[线性预算下的字母表指数阈值]\label{prop:xi-time-part9f-alphabet-threshold}
\leanverified{paper\_xi\_time\_part9f\_alphabet\_threshold}
设存在可逆提升
\[
\Psi_m:\Omega_m\to X_m\times\Sigma^{\le L_m},
\qquad
\mathrm{pr}_{X_m}\circ\Psi_m=\Fold_m,
\]
且 \(\Psi_m\) 单射。
若要求最坏情形长度满足
\[
L_m\le (\alpha+o(1))m,\qquad m\to\infty,
\]
则必有
\[
M\ge \exp\!\left(\frac{\log(2/\varphi)}{\alpha}\right).
\]
\end{proposition}
\begin{proof}
由 \(|R(\omega)|\le L_m\) 得 \(\mathbb E|R|\le L_m\)。
结合推论 \ref{cor:xi-time-part9f-no-constant-memory}：
\[
\frac{L_m}{m}\ge \frac{\log(2/\varphi)}{\log M}+o(1).
\]
取 \(m\to\infty\) 并移项即得结论。
\end{proof}

\begin{theorem}[微正则模型空间中的有限描述搜索失效率]\label{thm:xi-time-part9f-microcanonical-search-failure}
固定纤维谱 \(d\)，令
\[
\mathcal M_m(d):=
\left\{
F:\Omega_m\twoheadrightarrow X_m:\ |F^{-1}(x)|=d(x)\ \forall x
\right\}.
\]
取先验 \(F\sim\mathrm{Unif}(\mathcal M_m(d))\)。
记
\(
N:=|\Omega_m|,\ n:=|X_m|,\ w(x):=d(x)/N
\)。
若满足
\[
\log|\mathcal M_m(d)|=N\,H(w)+O(n\log N),\qquad
n\log N=o(N),
\]
且 \(H(w)\) 不退化（例如 \(\inf_m H(w)>0\)），则对任意候选族 \(\mathcal C_m\subseteq\mathcal M_m(d)\) 满足
\(
\log|\mathcal C_m|=\mathrm{poly}(m)
\)
有
\[
\mathbb P(F\in\mathcal C_m)\le \exp\!\bigl(-\Theta(N)\bigr).
\]
在主线 \(N=2^m\) 下，失败概率为 \(\exp(-\Theta(2^m))\)。
\end{theorem}
\begin{proof}
均匀先验下
\[
\mathbb P(F\in\mathcal C_m)=\frac{|\mathcal C_m|}{|\mathcal M_m(d)|}.
\]
取对数并代入规模展开：
\[
\log\mathbb P(F\in\mathcal C_m)
\le
\log|\mathcal C_m|-N\,H(w)+O(n\log N).
\]
由 \(n\log N=o(N)\)、\(\log|\mathcal C_m|=\mathrm{poly}(m)\)、\(\inf_m H(w)>0\)，右端为 \(-\Theta(N)\)。
指数化即得结论。
\end{proof}

\begin{theorem}[有限状态流式实现的指数稀有性与典型最小宽度]\label{thm:xi-time-part9f-streaming-width-rarity}
\leanverified{paper\_xi\_time\_part9f\_streaming\_width\_rarity}
设 \(N^{\mathrm{bp}}_{m,S}(n)\) 为长度 \(m\)、宽度至多 \(S\) 的只读一次流式分支程序在 \(n\) 元输出上的可实现函数数，且
\[
N^{\mathrm{bp}}_{m,S}(n)\le n^{S}S^{2Sm}.
\]
在定理 \ref{thm:xi-time-part9f-microcanonical-search-failure} 的微正则先验下，
\[
\mathbb P\!\left(
F\ \text{可由某个宽度}\le S\ \text{程序实现}
\right)
\le
\exp\!\Bigl(
2Sm\log S+S\log n-\log|\mathcal M_m(d)|
\Bigr).
\]
特别地，在 \(N=2^m\) 主线下，存在常数 \(c>0\) 使
\[
\mathbb P\!\left(
\min\{S:\ F\ \text{可由宽度}\le S\ \text{程序实现}\}
\le
c\,\frac{2^m}{m^2}
\right)
\le
\exp(-\Theta(2^m)),
\]
即典型最小宽度为 \(\Omega(2^m/m^2)\)。
\end{theorem}
\begin{proof}
第一式由可实现函数计数与并联合并直接得到。
在 \(N=2^m\) 下取 \(S=c2^m/m^2\)。
分子指数项 \(2Sm\log S+S\log n\) 为 \(O(c\,2^m)\)；
分母主项 \(\log|\mathcal M_m(d)|\) 为 \(\Theta(2^m)\)。
取 \(c\) 充分小即可使总指数为 \(-\Theta(2^m)\)，得到所述概率界与宽度下界。
\end{proof}

\begin{proposition}[高维极限下的边界全息饱和]\label{prop:xi-time-part9f-boundary-holographic-saturation}
\leanverified{paper\_xi\_time\_part9f\_boundary\_holographic\_saturation}
设 \(A\subset[0,1]^n\)，体 Gödel 维与边界 Gödel 维分别记为 \(d_G(A)\)、\(d_{\partial G}(A)\)。
若满足
\[
\frac{n-1}{n}d_G(A)\le d_{\partial G}(A)\le d_G(A),
\]
则
\[
0\le d_G(A)-d_{\partial G}(A)\le \frac{1}{n}d_G(A).
\]
故在 \(d_G(A)\) 有界族上，\(n\to\infty\) 时边界尺度对体维的双对数信息损失消失。
\end{proposition}
\begin{proof}
由夹逼式两端同减 \(d_{\partial G}(A)\) 立得不等式。
\end{proof}

\paragraph{dyadic 体--界全息的精确屏幕、线性阴影与有限矩恢复}
在命题 \ref{prop:xi-time-part9f-boundary-holographic-saturation} 已把边界尺度对体维的双对数损失压缩到 \(1/n\) 量级之后，固定分辨率 dyadic 体--界全息链还可继续抽取出下列五条有限对象结论。它们分别把精确屏幕的最小化、审计缺口的拓扑来源、模 \(2\) 阴影的编码结构、有限矩盒的完备恢复以及边界 Gödel 码的中值几何统一到同一组可直接调用的终端表述。

\begin{theorem}[精确屏幕的最小规模恰等于顶维胞腔数]\label{thm:xi-time-part9fa-exact-screen-minimal-size-matroid-basis}
固定 \(m,n\ge 1\)，令 \(\mathcal Q_m\) 为 \([0,1]^n\) 的 \(m\)-级 dyadic 立方剖分，并记
\[
\partial_n:C_n(\mathcal Q_m;\ZZ)\longrightarrow C_{n-1}(\mathcal Q_m;\ZZ)
\]
为顶维边界算子。对任意屏幕
\[
S\subset \overrightarrow{\mathcal F_m^{n-1}},
\qquad
f_S:=\Pi_S\circ\partial_n,
\]
再记
\[
R_{m,n}:=\rank_{\QQ}C_n(\mathcal Q_m;\QQ)=2^{mn},
\qquad
r(S):=\rank_{\QQ}\operatorname{im}(f_S\otimes\QQ).
\]
则
\[
\ker f_S=0
\iff
r(S)=R_{m,n}.
\]
并且在全部满足 \(\ker f_S=0\) 的屏幕中，最小可能基数恰为
\[
2^{mn},
\]
达到该最小值的全部屏幕恰构成边界行拟阵的基族。
\leanverified{paper\_xi\_time\_part9fa\_exact\_screen\_minimal\_size\_matroid\_basis}
\end{theorem}
\begin{proof}
由定理 \ref{thm:spg-dyadic-cubical-boundary-injective}，\(\partial_n\otimes\QQ\) 为单射，故
\[
\rank_{\QQ}\operatorname{im}(\partial_n\otimes\QQ)=\dim_{\QQ}C_n(\mathcal Q_m;\QQ)=2^{mn}=R_{m,n}.
\]
另一方面，定理 \ref{thm:spg-screen-rank-matroid-supermodularity} 给出
\[
\delta(S)=R_{m,n}-r(S)=\rank_{\QQ}\ker(f_S\otimes\QQ).
\]
由于 \(C_n(\mathcal Q_m;\ZZ)\) 自由，\(\ker f_S\) 亦自由，从而
\[
\ker f_S=0
\iff
\rank_{\QQ}\ker(f_S\otimes\QQ)=0
\iff
r(S)=R_{m,n}.
\]
这证明了第一部分。

再由定理 \ref{thm:spg-exact-screen-minima-are-matroid-bases}，极小精确屏幕恰为上述可表示行拟阵的全部基；而每一组基的基数都等于整条边界行空间的秩 \(R_{m,n}=2^{mn}\)。故最小规模恰为 \(2^{mn}\)，并且达到该极小值者恰是全部拟阵基。证毕。
\end{proof}

\begin{theorem}[屏幕核由未观测边界子复形的顶维约化同调完全刻画]\label{thm:xi-time-part9fa-screen-kernel-top-homology}
\leanverified{paper\_xi\_time\_part9fa\_screen\_kernel\_top\_homology}
沿用定理 \ref{thm:xi-time-part9fa-exact-screen-minimal-size-matroid-basis} 的记号。令 \(A_{\neg S}\) 为删去 \(S\) 所对应底层 \((n-1)\)-胞腔后得到的未观测边界子复形。则存在自然同构
\[
\ker f_S
\cong
H_n(\mathcal Q_m,A_{\neg S};\ZZ)
\cong
\widetilde H_{n-1}(A_{\neg S};\ZZ).
\]
因此若记最小审计补全成本为
\[
\mathrm{Audit}(S):=\rank_{\ZZ}\ker f_S,
\]
则
\[
\mathrm{Audit}(S)=\rank_{\ZZ}\widetilde H_{n-1}(A_{\neg S};\ZZ).
\]
特别地，
\[
\ker f_S=0
\iff
\widetilde H_{n-1}(A_{\neg S};\ZZ)=0.
\]
\end{theorem}
\begin{proof}
定理 \ref{thm:spg-partial-boundary-screen-kernel-relative-homology} 已给出
\[
\ker f_S\cong H_n(\mathcal Q_m,A_{\neg S};\ZZ).
\]
又因为 \(|\mathcal Q_m|=[0,1]^n\) 可缩，故其约化同调全部消失。对配对 \((\mathcal Q_m,A_{\neg S})\) 取约化长正合列，得到
\[
H_n(\mathcal Q_m,A_{\neg S};\ZZ)\cong \widetilde H_{n-1}(A_{\neg S};\ZZ),
\]
从而得到第二个自然同构。再取自由阿贝尔群的秩便得
\[
\mathrm{Audit}(S)=\rank_{\ZZ}\widetilde H_{n-1}(A_{\neg S};\ZZ),
\]
其零秩情形即给出最后一个等价关系。证毕。
\end{proof}

\begin{theorem}[dyadic 边界像天然形成一族 LDPC 码]\label{thm:xi-time-part9fa-dyadic-boundary-image-ldpc}
\leanverified{paper\_xi\_time\_part9fa\_dyadic\_boundary\_image\_ldpc}
令
\[
\mathcal C_{m,n}
:=
\operatorname{im}\!\bigl(
\partial_n:
C_n(\mathcal Q_m;\FF_2)\to C_{n-1}(\mathcal Q_m;\FF_2)
\bigr).
\]
则 \(\mathcal C_{m,n}\) 是长度
\[
N_{m,n}
:=
\#\mathcal F_{n-1}(\mathcal Q_m)
=
n(2^m+1)\,2^{m(n-1)}
\]
的线性码，并且
\[
\dim_{\FF_2}\mathcal C_{m,n}=2^{mn},
\qquad
R_{m,n}:=\frac{\dim \mathcal C_{m,n}}{N_{m,n}}
=
\frac{2^m}{n(2^m+1)}
\xrightarrow[m\to\infty]{}
\frac1n.
\]
此外，其最小距离满足精确闭式
\[
d_{\min}(\mathcal C_{m,n})=2n.
\]
并且以 \(\partial_{n-1}\) 为校验算子时，每条校验方程恰涉及 \(4\) 个 \((n-1)\)-胞腔坐标、每个变量节点度为 \(2(n-1)\)；因而 \(\mathcal C_{m,n}\) 构成一族规则 LDPC 码。
\end{theorem}
\begin{proof}
长度计数与维数公式正是定理 \ref{thm:spg-dyadic-boundary-image-ldpc} 的第一部分：每个坐标方向有 \(2^m+1\) 个切片，每个切片上含 \(2^{m(n-1)}\) 个 \((n-1)\)-胞腔，从而得到总长度 \(N_{m,n}\)；而顶维边界算子在 \(\FF_2\) 上仍保持单射，故
\[
\dim_{\FF_2}\mathcal C_{m,n}=2^{mn}.
\]
码率公式随即化简得到，并趋于 \(1/n\)。

最小距离的精确值由推论 \ref{cor:spg-dyadic-boundary-image-min-distance} 给出，即
\[
d_{\min}(\mathcal C_{m,n})=2n.
\]
最后，定理 \ref{thm:spg-dyadic-boundary-image-ldpc} 还表明：\(\partial_{n-1}\) 的每一条校验行只含四个非零项，而每个 \((n-1)\)-胞腔位于 \(2(n-1)\) 个 \((n-2)\)-胞腔边界之中，故该码族确为规则 LDPC 码。证毕。
\end{proof}

\begin{theorem}[固定分辨率 dyadic polycube 的有限矩盒唯一恢复]\label{thm:xi-time-part9fa-dyadic-finite-moment-completeness}
\leanverified{paper\_xi\_time\_part9fa\_dyadic\_finite\_moment\_completeness}
固定 \(m,n\ge 1\)，并令 \(N:=2^m\)。对任意分辨率 \(m\) 的 dyadic polycube
\[
U\subset [0,1]^n
\]
及任意多重指标 \(\alpha=(\alpha_1,\dots,\alpha_n)\in\{0,\dots,N-1\}^n\)，定义单项式矩
\[
\mu_\alpha(U):=\int_U x^\alpha\,\dd x
=
\int_U x_1^{\alpha_1}\cdots x_n^{\alpha_n}\,\dd x.
\]
则映射
\[
U\longmapsto \bigl(\mu_\alpha(U)\bigr)_{\alpha\in\{0,\dots,N-1\}^{n}}
\]
在全部 \(m\)-级 dyadic polycube 类上为单射。换言之，固定分辨率下仅用有限矩盒
\[
\{\mu_\alpha(U):0\le \alpha_i\le 2^m-1\}
\]
便足以唯一确定 \(U\)。
\end{theorem}
\begin{proof}
这是定理 \ref{thm:spg-dyadic-finite-moment-completeness} 的直接改写。其证明把分辨率 \(m\) 的 dyadic 小立方体示性函数作为基底，从而得到维数为 \(N^n\) 的分段常数空间 \(\mathcal V_m\)。对应的矩映射
\[
\mathcal M_m(f):=
\left(\int_{I^n} f(x)x^\alpha\,\dd x\right)_{\alpha\in\{0,\dots,N-1\}^{n}}
\]
在该基下分解为一维矩阵 \(M\) 的张量积 \(M^{\otimes n}\)，其中
\[
M_{j,k}:=\int_{k/N}^{(k+1)/N}x^j\,\dd x.
\]
由落阶幂与 Abel 型整理，一维矩阵 \(M\) 可逆，故 \(M^{\otimes n}\) 亦可逆；于是 \(\mathcal M_m\) 为线性同构。代入 \(f=\ind_U\) 即得所述单射性。证毕。
\end{proof}

\leanverified{paper\_xi\_time\_part9fa\_boundary\_godel\_median\_lipschitz}
\begin{theorem}[边界 Gödel 码的中值几何与体积读出的 Lipschitz 控制]\label{thm:xi-time-part9fa-boundary-godel-median-lipschitz}
为每一条取向 \((n-1)\)-面 \(\vec f\in\overrightarrow{\mathcal F_m^{n-1}}\) 赋予一枚互异素数 \(q_{\vec f}\)，并对任意 dyadic polycube \(U\) 定义边界 Gödel 整数
\[
H_m(U):=
\prod_{\vec f}q_{\vec f}^{\varepsilon_{\vec f}},
\qquad
\partial_n[U]=\sum_{\vec f}\varepsilon_{\vec f}\,\vec f,
\quad
\varepsilon_{\vec f}\in\{0,1\}.
\]
再定义
\[
\mathfrak Q_m^{\partial}
:=
\left\{
\prod_{\vec f\in\overrightarrow{\mathcal F_m^{n-1}}}
q_{\vec f}^{\varepsilon_{\vec f}}
:\ \varepsilon_{\vec f}\in\{0,1\}
\right\},
\]
以及算术差分距离
\[
d_m^{\mathrm{arith}}(A,B):=
\omega\!\left(\frac{AB}{\gcd(A,B)^2}\right),
\qquad
A,B\in\mathfrak Q_m^{\partial},
\]
其中 \(\omega\) 表示互异素因子个数。则：
\begin{enumerate}
  \item \((\mathfrak Q_m^{\partial},d_m^{\mathrm{arith}})\) 与有限 Hamming 超立方体等距；
  \item 因而它是一个中值度量空间，且其中值满足显式闭式
        \[
        \operatorname{med}(A,B,C)=\operatorname{rad}\!\bigl(\gcd(AB,AC,BC)\bigr);
        \]
        其中 \(\operatorname{rad}\) 表示平方自由核；
  \item 对任意 dyadic polycube \(U,V\)，有
        \[
        d_m^{\mathrm{arith}}\bigl(H_m(U),H_m(V)\bigr)
        =
        \bigl\|\partial_n[U]-\partial_n[V]\bigr\|_{\ell^1},
        \]
        并且体积读出满足严格 Lipschitz 控制
        \[
        \bigl|\operatorname{Vol}(U)-\operatorname{Vol}(V)\bigr|
        \le
        2^{-m(n-1)}\,
        d_m^{\mathrm{arith}}\bigl(H_m(U),H_m(V)\bigr).
        \]
\end{enumerate}
\end{theorem}
\begin{proof}
定理 \ref{thm:spg-boundary-godel-arithmetic-hypercube-median} 已证明：把平方自由整数 \(A\in\mathfrak Q_m^{\partial}\) 送到其素因子支撑的 \(0/1\) 向量，得到 \(\mathfrak Q_m^{\partial}\) 与 Hamming 立方体之间的等距同构；而逐坐标多数表决翻译回平方自由整数后，恰得到
\[
\operatorname{med}(A,B,C)=\operatorname{rad}\!\bigl(\gcd(AB,AC,BC)\bigr).
\]
这给出前两条。

对第 \(3\) 条，定理 \ref{thm:spg-godel-boundary-gcd-lipschitz-stability} 给出
\[
\omega\!\left(\frac{H_m(U)\,H_m(V)}{\gcd(H_m(U),H_m(V))^2}\right)
=
\bigl\|\partial_n[U]-\partial_n[V]\bigr\|_{\ell^1},
\]
也即
\[
d_m^{\mathrm{arith}}\bigl(H_m(U),H_m(V)\bigr)
=
\bigl\|\partial_n[U]-\partial_n[V]\bigr\|_{\ell^1}.
\]
同一定理还给出
\[
\bigl|\operatorname{Vol}(U)-\operatorname{Vol}(V)\bigr|
\le
2^{-m(n-1)}
\omega\!\left(\frac{H_m(U)\,H_m(V)}{\gcd(H_m(U),H_m(V))^2}\right),
\]
代入定义即得所示 Lipschitz 控制。证毕。
\end{proof}

\noindent
因此，dyadic 体--界全息在结论层形成了一条闭合链：精确屏幕的极小性由边界行拟阵基完全控制，审计缺口被未观测边界子复形的顶维洞精确吸收，其模 \(2\) 阴影又自动生成一族码率趋于 \(1/n\)、最小距离恒为 \(2n\) 的 LDPC 码；与此同时，固定分辨率 dyadic polycube 仍可由有限矩盒唯一恢复，而边界 Gödel 码则不仅保留全部边界信息，还天然携带 Hamming 型中值几何与体积读出的严格 Lipschitz 稳定。

\endinput


\begin{theorem}[局部归一化步数的线性下界与指数左尾]\label{thm:xi-time-part9f-local-normalization-step-lower-bound}
\leanverified{paper\_xi\_time\_part9f\_local\_normalization\_step\_lower\_bound}
设 \(G_m(\omega)\) 为长度 \(m\) 的规范差总量，并假定存在密度极限
\[
g^{\ast}:=\lim_{m\to\infty}\frac{\mathbb E[G_m]}{m}=\frac49.
\]
设局部重写算法 \(S\) 在每一步至多消去 \(R\) 个单位规范差，且总步数 \(T_m(S)\) 满足
\[
G_m(\omega)\le R\,T_m(S)(\omega)\qquad(\forall\,\omega,m).
\]
若对任意 \(b<g^{\ast}\) 存在 \(c_b>0\) 使
\[
\mathbb P\!\left(\frac{G_m}{m}\le b\right)\le e^{-c_b m},
\]
则
\[
\liminf_{m\to\infty}\frac{T_m(S)}{m}\ge \frac{g^{\ast}}{R}
=\frac{4}{9R}
\quad\text{a.s.},
\]
并且对任意 \(a<4/(9R)\) 有
\[
\limsup_{m\to\infty}\frac1m\log
\mathbb P\!\left(\frac{T_m(S)}{m}\le a\right)<0.
\]
\end{theorem}
\begin{proof}
由 \(G_m\le RT_m\) 得事件包含
\[
\left\{\frac{T_m}{m}\le a\right\}\subseteq
\left\{\frac{G_m}{m}\le aR\right\}.
\]
若 \(a<g^{\ast}/R\)，则 \(aR<g^{\ast}\)，由左尾指数界得
\[
\mathbb P\!\left(\frac{T_m}{m}\le a\right)\le e^{-c_{aR}m},
\]
从而得到第二式。
又由 \(\sum_m e^{-c_{aR}m}<\infty\) 与 Borel--Cantelli，几乎处处仅有限多个 \(m\) 使 \(T_m/m\le a\)。
令 \(a\uparrow g^{\ast}/R\) 即得 \(\liminf\) 下界。
\end{proof}

\paragraph{统一解释}
在折叠纤维语义下，受限于有限字母表、有限状态与局部重写半径的实现机制，未被消除的微观自由度将以侧信息密度、状态宽度与归一化步数三类量纲被强制外置，且其主尺度均由纤维体积与其熵账本决定。

\paragraph{纤维 Rényi 全谱、二元吸引子测度与 Euler 乘积分离率}
延续前述折叠语义，以下条目给出一组互相兼容的全谱定量结论。

\begin{theorem}[纤维均匀提升的全阶 Rényi 极大性与熵缺口指数表示]\label{thm:xi-time-part9g-renyi-maximizer-fiber-uniform}
设 \(\Fold_m:\Omega_m\to X_m\) 为有限纤维映射，
\[
d_m(x):=\bigl|\Fold_m^{-1}(x)\bigr|,\qquad
\mathcal L(\pi):=\{\mu\in\mathcal P(\Omega_m):(\Fold_m)_{\ast}\mu=\pi\},
\]
其中 \(\pi\in\mathcal P(X_m)\) 给定。定义纤维条件均匀提升
\[
\pi^{\mathrm e}(\omega):=\frac{\pi(x)}{d_m(x)},
\qquad \omega\in \Fold_m^{-1}(x).
\]
对 \(\alpha>0\)，\(\alpha\neq 1\)，令
\[
H_\alpha(\mu):=\frac{1}{1-\alpha}\log\sum_{\omega\in\Omega_m}\mu(\omega)^\alpha,
\]
并取 \(H_1:=H\) 为 Shannon 熵。则：
\begin{enumerate}
  \item 对每个 \(\alpha>0\)（含 \(\alpha=1\)），\(\pi^{\mathrm e}\) 是 \(\mathcal L(\pi)\) 上 \(H_\alpha\) 的唯一极大元；
  \item 对 \(\alpha\neq 1\) 与任意 \(\mu\in\mathcal L(\pi)\)，若 \(\mu_x\) 为 \(\mu\) 在纤维 \(\Fold_m^{-1}(x)\) 上的条件分布、\(U_x\) 为该纤维均匀分布，则
  \[
  H_\alpha(\mu)
  =\frac{1}{1-\alpha}
  \log\!\sum_{x\in X_m}
  \pi(x)^\alpha d_m(x)^{1-\alpha}
  \exp\!\bigl((\alpha-1)D_\alpha(\mu_x\Vert U_x)\bigr),
  \]
  从而
  \[
  H_\alpha(\pi^{\mathrm e})-H_\alpha(\mu)
  =\frac{1}{\alpha-1}
  \log\!\sum_{x\in X_m}
  w_\alpha(x)\exp\!\bigl((\alpha-1)D_\alpha(\mu_x\Vert U_x)\bigr),
  \]
  其中
  \[
  w_\alpha(x):=
  \frac{\pi(x)^\alpha d_m(x)^{1-\alpha}}
  {\sum_{y\in X_m}\pi(y)^\alpha d_m(y)^{1-\alpha}}.
  \]
\end{enumerate}
\end{theorem}
\begin{proof}
对 \(\mu\in\mathcal L(\pi)\) 作纤维分解：
\(
\mu(\omega)=\pi(x)\mu_x(\omega)
\)
（\(\omega\in\Fold_m^{-1}(x)\)）。
于是
\[
\sum_{\omega}\mu(\omega)^\alpha
=\sum_{x}\pi(x)^\alpha\sum_{\omega\in\Fold_m^{-1}(x)}\mu_x(\omega)^\alpha.
\]
对固定 \(x\)，函数 \(\sum \mu_x^\alpha\) 在单纯形上于 \(\alpha>1\) 时严格凸、于 \(0<\alpha<1\) 时严格凹；结合 \(H_\alpha\) 前系数 \(1/(1-\alpha)\) 的符号，可知两种情形均在且仅在 \(\mu_x=U_x\) 时达到最大熵。故 \(\pi^{\mathrm e}\) 为唯一极大元。 \(\alpha=1\) 时由 Shannon 熵在纤维条件分布上的严格凹性同样得到唯一性。

再由
\[
D_\alpha(\mu_x\Vert U_x)
=\frac{1}{\alpha-1}
\log\!\left(
\sum_{\omega\in\Fold_m^{-1}(x)}
\mu_x(\omega)^\alpha U_x(\omega)^{1-\alpha}
\right)
\]
及 \(U_x(\omega)=1/d_m(x)\) 得
\[
\sum_{\omega\in\Fold_m^{-1}(x)}\mu_x(\omega)^\alpha
=d_m(x)^{1-\alpha}\exp\!\bigl((\alpha-1)D_\alpha(\mu_x\Vert U_x)\bigr),
\]
代回即得表示式。再与 \(\mu=\pi^{\mathrm e}\) 情形相减即得熵缺口公式。
\end{proof}

\begin{theorem}[熵窗提升的 Rényi 全谱塌缩]\label{thm:xi-time-part9g-renyi-spectrum-holography}
设
\[
L_m(\gamma,\varepsilon):=
\left\{
x\in X_m:
e^{m(\gamma-\varepsilon)}\le d_m(x)\le e^{m(\gamma+\varepsilon)}
\right\},
\]
\(\pi_{m,\gamma,\varepsilon}\) 为 \(L_m(\gamma,\varepsilon)\) 上均匀分布，
\[
\mu^{\mathrm e}_{m,\gamma,\varepsilon}:=
Q_m\pi_{m,\gamma,\varepsilon},
\]
并记
\[
f(\gamma):=
\lim_{\varepsilon\downarrow 0}\lim_{m\to\infty}
\frac1m\log|L_m(\gamma,\varepsilon)|
\]
（极限存在时）。则对任意 \(\alpha>0\)：
\[
\lim_{\varepsilon\downarrow 0}\lim_{m\to\infty}
\frac1m H_\alpha\!\left(\mu^{\mathrm e}_{m,\gamma,\varepsilon}\right)
=f(\gamma)+\gamma,
\]
且极限与 \(\alpha\) 无关。
\end{theorem}
\begin{proof}
记 \(L_m:=L_m(\gamma,\varepsilon)\)、\(\mu^{\mathrm e}:=\mu^{\mathrm e}_{m,\gamma,\varepsilon}\)。
对 \(\alpha\neq1\)：
\[
\sum_{\omega}(\mu^{\mathrm e}(\omega))^\alpha
=|L_m|^{-\alpha}\sum_{x\in L_m}d_m(x)^{1-\alpha}.
\]
由 \(x\in L_m\) 的定义得到
\[
\left|
\frac{1}{1-\alpha}
\log\!\left(\frac1{|L_m|}\sum_{x\in L_m}d_m(x)^{1-\alpha}\right)
-m\gamma
\right|
\le m\varepsilon.
\]
故
\[
\left|
H_\alpha(\mu^{\mathrm e})-\bigl(\log|L_m|+m\gamma\bigr)
\right|
\le m\varepsilon.
\]
对 \(\alpha=1\)，有
\[
H(\mu^{\mathrm e})
=H(\pi_{m,\gamma,\varepsilon})+\mathbb E_{\pi_{m,\gamma,\varepsilon}}[\log d_m(X)]
=\log|L_m|+\mathbb E[\log d_m(X)],
\]
且 \(\log d_m(X)\in[m(\gamma-\varepsilon),m(\gamma+\varepsilon)]\)，同样得
\(
|H(\mu^{\mathrm e})-(\log|L_m|+m\gamma)|\le m\varepsilon
\)。
两式除以 \(m\)，先令 \(m\to\infty\) 再令 \(\varepsilon\downarrow0\) 即得结论。
\end{proof}

\begin{theorem}[二元数字吸引子的临界 Hausdorff 测度与停机深度]\label{thm:xi-time-part9g-godel-leyang-hausdorff}
令 \(B:=2^L\)，并在 \(T\simeq\ZZ_2^2\) 上取标准方向 \(e=(1,0)\)。
定义数字集
\[
D_{M,x}:=\{0,\alpha_{M,x}e\},
\]
其中 \(\alpha_{M,x}\in\ZZ_2\) 满足：
非停机时 \(\alpha_{M,x}=0\)；若停机深度为 \(N\)，则 \(\alpha_{M,x}=-B^N\)。
令 \(K_{M,x}\subset T\) 为对应二元数字吸引子。设 \(s:=1/L\)。则
\[
\mathcal H^s(K_{M,x})=
\begin{cases}
0, & \alpha_{M,x}=0,\\[1mm]
2^{-N}, & \alpha_{M,x}=-B^N.
\end{cases}
\]
因此停机深度可由
\[
N=-\log_2\mathcal H^{1/L}(K_{M,x})
\]
精确恢复（停机情形）。
\end{theorem}
\begin{proof}
若 \(\alpha_{M,x}=0\)，则 \(D_{M,x}\) 退化为单点，\(K_{M,x}\) 为单点集，故 \(\mathcal H^s(K_{M,x})=0\)（\(s>0\)）。

若 \(\alpha_{M,x}=-B^N\)，记
\[
C:=\left\{\sum_{k\ge0}u_k B^k:\ u_k\in\{0,1\}\right\}\subset\ZZ_2.
\]
则
\[
K_{M,x}=-B^N(Ce).
\]
在 \(2\)-进超距量下，伸缩 \(z\mapsto B^N z\) 的相似比为 \(|B^N|_2=2^{-LN}\)，故
\[
\mathcal H^s(K_{M,x})=(2^{-LN})^s\,\mathcal H^s(Ce)=2^{-N}\mathcal H^s(C).
\]
又 \(e\) 给出 \(\ZZ_2\to\ZZ_2^2\) 的等距嵌入，故 \(\mathcal H^s(Ce)=\mathcal H^s(C)\)。

对任意 \(m\)，集合 \(C\) 被前缀长度 \(m\) 的 \(2^m\) 个互不相交圆柱块覆盖，每块半径 \(2^{-Lm}\)。
故上内容量给出
\[
\mathcal H^s(C)\le \lim_{m\to\infty}2^m(2^{-Lm})^s=1.
\]
反向，由这些圆柱块两两不交，任意半径不大于 \(2^{-Lm}\) 的覆盖至少需 \(2^m\) 个球，故
\[
\mathcal H^s(C)\ge \liminf_{m\to\infty}2^m(2^{-Lm})^s=1.
\]
于是 \(\mathcal H^s(C)=1\)，从而 \(\mathcal H^s(K_{M,x})=2^{-N}\)。
\end{proof}

\begin{proposition}[二进扭子链中的首次偏离深度判据]\label{prop:xi-time-part9g-first-deviation-depth}
设 \(E_{\min}\) 为固定椭圆曲线，\(P_i\in E_{\min}\)。
取兼容链 \((P_{i,n})_{n\ge0}\) 满足 \([2]P_{i,n}=P_{i,n-1}\)，并令
\[
\Pi_{i,\infty}:=\varprojlim_n\Pi_{i,n},
\qquad
\Pi_{i,n}:=\{Q\in E_{\min}:[2^n]Q=P_i\}.
\]
选定同构 \(\Pi_{i,\infty}\simeq P_i+T_2(E_{\min})\)，取 \(e\in T_2(E_{\min})\)。
定义
\[
Q_{M,x}:=P_i+\alpha_{M,x}e,\qquad
Q_{M,x}^{(n)}:=\mathrm{pr}_n(Q_{M,x}).
\]
若 \(\alpha_{M,x}=0\)（非停机），则 \(Q_{M,x}^{(n)}=P_{i,n}\) 对全部 \(n\) 成立。
若 \(\alpha_{M,x}=-2^{LN}\)，则
\[
Q_{M,x}^{(n)}=P_{i,n}\ \ (n\le LN),\qquad
Q_{M,x}^{(LN+1)}\neq P_{i,LN+1},
\]
并且
\[
N=\frac{1}{L}\left(\min\{n\ge1:Q_{M,x}^{(n)}\neq P_{i,n}\}-1\right).
\]
\end{proposition}
\begin{proof}
在所选扭子同构下，\(\mathrm{pr}_n\) 等价于模 \(2^n\) 投影。
因此
\[
Q_{M,x}^{(n)}=P_{i,n}
\Longleftrightarrow
\alpha_{M,x}e\in 2^nT_2(E_{\min})
\Longleftrightarrow
v_2(\alpha_{M,x})\ge n.
\]
非停机时 \(\alpha_{M,x}=0\) 对所有 \(n\) 成立。
停机深度 \(N\) 时 \(\alpha_{M,x}=-2^{LN}\) 恰有 \(v_2(\alpha_{M,x})=LN\)，故结论立即得到。
\end{proof}

\begin{proposition}[有限层扭子投影的可计算性与首次偏离深度的停机不可计算性]\label{prop:xi-time-part9g-deviation-depth-computability-separation}
沿命题 \ref{prop:xi-time-part9g-first-deviation-depth} 的记号，并定义首次偏离深度
\[
\delta(M,x):=\min\{n\ge 1:\ Q_{M,x}^{(n)}\neq P_{i,n}\}\in\NN\cup\{\infty\},
\]
其中若集合为空则约定 \(\delta(M,x)=\infty\)。
则：
\begin{enumerate}
\normalfont
  \item \textbf{有限层投影可计算：}映射 \((M,x,n)\mapsto Q_{M,x}^{(n)}\in\Pi_{i,n}\) 为全可计算函数；
  \item \textbf{首次偏离深度不可计算：}映射 \((M,x)\mapsto \delta(M,x)\) 不是全可计算函数；
  \item \textbf{统一验证深度的 Busy--Beaver 下界：}
  若存在函数 \(\Phi:\NN\to\NN\) 使对所有满足 \(|(M,x)|\le \ell\) 的停机实例都有 \(\delta(M,x)\le \Phi(\ell)\)，
  则存在常数 \(c\)（仅依赖于编码）使对所有 \(\ell\) 有
  \[
  \Phi(\ell)\ \ge\ L\cdot \mathrm{BB}(\ell-c)+1,
  \]
  其中 \(\mathrm{BB}\) 为 Busy--Beaver 函数（定理 \ref{thm:conclusion-finite-test-depth-busybeaver}），从而 \(\Phi\) 不可计算。
\end{enumerate}
\end{proposition}
\begin{proof}
对 (1)，给定 \((M,x,n)\)。
令 \(t:=\lfloor (n-1)/L\rfloor\)。
若 \(M(x)\) 在不超过 \(t\) 步内停机，则首次停机时间 \(N\le t\) 可由有限步模拟得到，进而由命题 \ref{prop:xi-time-part9g-first-deviation-depth} 确定 \(\alpha_{M,x}=-2^{LN}\) 及 \(\Pi_{i,n}\) 中的点 \(Q_{M,x}^{(n)}=\mathrm{pr}_n(P_i+\alpha_{M,x}e)\)。
若 \(M(x)\) 在 \(t\) 步内不停机，则无论其是否在更晚时刻停机，都有 \(n\le LN\)（停机时）或 \(\alpha_{M,x}=0\)（不停机时），从而由命题 \ref{prop:xi-time-part9g-first-deviation-depth} 得 \(Q_{M,x}^{(n)}=P_{i,n}\)。
因此 \(Q_{M,x}^{(n)}\) 可由有限步模拟与有限群运算构造性求出，故为全可计算函数。
\par
对 (2)，由命题 \ref{prop:xi-time-part9g-first-deviation-depth}，
\[
\delta(M,x)<\infty\quad\Longleftrightarrow\quad M(x)\ \text{停机}.
\]
若 \(\delta\) 全可计算，则停机问题可判定，矛盾。
\par
对 (3)，若 \(M(x)\) 首次停机时间为 \(N\)，则由命题 \ref{prop:xi-time-part9g-first-deviation-depth} 有 \(\delta(M,x)=LN+1\)，从而
\[
N\le \frac{\Phi(|(M,x)|)-1}{L}.
\]
于是 \(\Phi\) 给出了按描述长度统一控制的停机时间上界，违背定理 \ref{thm:conclusion-finite-test-depth-busybeaver} 所给出的 Busy--Beaver 下界；整理常数即得所示不等式。证毕。
\end{proof}

\paragraph{\texorpdfstring{$2^L$}{2^L}-进仿射半群的固定点刚性}
沿用推论 \ref{cor:xi-terminal-zm-leyang-branchpoint-tate-torsor-decomposition} 给出的识别
\[
T_2(E_{\min})\cong \varprojlim_n E_{\min}[2^n]\cong \ZZ_2^2.
\]
固定正整数 \(L\)，并记 \(B:=2^L\)。
设 \(R\) 为交换幺半群，\(\Sigma:R\to R\) 为幺半群端同态，
\[
H:=R\rtimes_{\Sigma}\NN
\]
为乘法
\[
(r,k)(r',k'):=\bigl(r+\Sigma^k r',\,k+k'\bigr)
\]
定义的半直积幺半群。
再设 \(\operatorname{ev}_B:R\to T_2(E_{\min})\) 为加法映射，满足
\[
\operatorname{ev}_B(\Sigma^k r)=B^k\operatorname{ev}_B(r)
\qquad (r\in R,\ k\in\NN),
\]
并令
\[
\Psi(r,k)(x):=B^k x+\operatorname{ev}_B(r)
\qquad \bigl(x\in T_2(E_{\min})\bigr)
\]
为 \(H\) 在 \(T_2(E_{\min})\) 上的忠实 \(2^L\)-进仿射表示。

\begin{theorem}[正层 \texorpdfstring{$2^L$}{2^L}-进仿射元的唯一不动点]\label{thm:xi-time-part9g-affine-semigroup-unique-fixed-point}
\leanverified{paper\_xi\_time\_part9g\_affine\_semigroup\_unique\_fixed\_point}
沿用上段记号。设 \(g=(r,k)\in H\)。
\begin{enumerate}
  \item 若 \(k\ge 1\)，则 \(\Psi(g)\) 存在唯一不动点
  \[
  x_g=(1-B^k)^{-1}\operatorname{ev}_B(r)\in T_2(E_{\min}).
  \]
  \item 若 \(k=0\)，则 \(\Psi(r,0)\) 有不动点当且仅当 \((r,0)\) 为 \(H\) 的单位元。
\end{enumerate}
\end{theorem}
\begin{proof}
在 \(T_2(E_{\min})\cong \ZZ_2^2\) 的识别下，若 \(k\ge 1\) 则
\[
1-B^k=1-2^{Lk}
\]
为奇数，因而是 \(\ZZ_2\) 中的单位元；于是乘法
\(
x\mapsto (1-B^k)x
\)
是 \(T_2(E_{\min})\) 上的自同构。
不动点方程
\[
x=B^k x+\operatorname{ev}_B(r)
\]
等价于
\[
(1-B^k)x=\operatorname{ev}_B(r),
\]
故唯一解即为
\(
x_g=(1-B^k)^{-1}\operatorname{ev}_B(r)
\)。

若 \(k=0\)，则
\[
\Psi(r,0)(x)=x+\operatorname{ev}_B(r).
\]
其存在不动点当且仅当 \(\operatorname{ev}_B(r)=0\)，亦即 \(\Psi(r,0)=\Id\)。
由于表示 \(\Psi\) 忠实，且 \(\Psi(0,0)=\Id\)，故此时且仅此时 \((r,0)=(0,0)\)，即 \((r,0)\) 为 \(H\) 的单位元。
\end{proof}

\begin{theorem}[不动点族的仿射 \texorpdfstring{$1$}{1}-余循环律]\label{thm:xi-time-part9g-affine-fixed-point-cocycle}
\leanverified{paper\_xi\_time\_part9g\_affine\_fixed\_point\_cocycle}
沿用上段记号。设
\[
g=(r,k),\qquad h=(r',k')\in H,
\qquad k,k'\ge 1,
\]
并记各自唯一不动点为 \(x_g\)、\(x_h\)。
则乘积 \(gh\) 的唯一不动点满足严格闭式
\[
x_{gh}
=
\frac{1-B^k}{1-B^{k+k'}}\,x_g
+\frac{B^k(1-B^{k'})}{1-B^{k+k'}}\,x_h.
\]
从而映射
\[
g\longmapsto x_g
\]
在正层子半群
\(
\{(r,k)\in H:\ k\ge 1\}
\)
上满足严格的 \(2\)-进仿射 \(1\)-余循环律。
\end{theorem}
\begin{proof}
由半直积乘法，
\[
gh=(r+\Sigma^k r',\,k+k').
\]
因此
\[
x_{gh}
=(1-B^{k+k'})^{-1}\operatorname{ev}_B(r+\Sigma^k r').
\]
再由 \(\operatorname{ev}_B\) 的加法性与移位相容性，
\[
\operatorname{ev}_B(r+\Sigma^k r')
=\operatorname{ev}_B(r)+B^k\operatorname{ev}_B(r').
\]
于是
\[
x_{gh}
=(1-B^{k+k'})^{-1}
\bigl(\operatorname{ev}_B(r)+B^k\operatorname{ev}_B(r')\bigr).
\]
将
\[
\operatorname{ev}_B(r)=(1-B^k)x_g,
\qquad
\operatorname{ev}_B(r')=(1-B^{k'})x_h
\]
代入即可得到所示公式。
\end{proof}

\begin{theorem}[固定点的指数吸引与同层可分辨性]\label{thm:xi-time-part9g-affine-fixed-point-contraction-rigidity}
沿用上段记号。设 \(g=(r,k)\in H\) 且 \(k\ge 1\)，并记其唯一不动点为 \(x_g\)。
则对任意 \(x\in T_2(E_{\min})\) 与任意 \(n\ge 1\)，有
\[
\Psi(g)^n(x)-x_g=B^{kn}(x-x_g).
\]
因此 \(2\)-进范数满足
\[
\bigl|\Psi(g)^n(x)-x_g\bigr|_2
=
2^{-Lkn}\,\bigl|x-x_g\bigr|_2.
\]
特别地，每个正层元 \(\Psi(g)\) 都是严格 \(2\)-进压缩映射，其全部轨道以指数速率收敛到唯一不动点。

进一步，若
\[
g=(r,k),\qquad g'=(r',k)\in H
\]
具有相同层级 \(k\ge 1\)，且 \(x_g=x_{g'}\)，则 \(g=g'\)。
换言之，不动点在每个固定层级上完全可分辨。
\end{theorem}
\begin{proof}
由 \(x_g\) 的定义，
\[
(1-B^k)x_g=\operatorname{ev}_B(r).
\]
故
\[
\Psi(g)(x)-x_g
=B^k x+\operatorname{ev}_B(r)-x_g
=B^k x+(1-B^k)x_g-x_g
=B^k(x-x_g).
\]
迭代 \(n\) 次即得
\[
\Psi(g)^n(x)-x_g=B^{kn}(x-x_g).
\]
又因 \(|B|_2=|2^L|_2=2^{-L}\)，故
\[
\bigl|B^{kn}(x-x_g)\bigr|_2
=2^{-Lkn}\,\bigl|x-x_g\bigr|_2,
\]
从而严格压缩与指数吸引成立。

若 \(x_g=x_{g'}\)，则
\[
\operatorname{ev}_B(r)
=(1-B^k)x_g
=(1-B^k)x_{g'}
=\operatorname{ev}_B(r').
\]
于是
\[
\Psi(g)(x)=B^k x+\operatorname{ev}_B(r)
=B^k x+\operatorname{ev}_B(r')
=\Psi(g')(x)
\]
对所有 \(x\) 成立，即 \(\Psi(g)=\Psi(g')\)。
由于 \(\Psi\) 忠实，遂得 \(g=g'\)。
\end{proof}

\paragraph{规范 Gödel 半直积的 \texorpdfstring{$B$}{B}-进固定点、前缀地址与 Fibonacci 截面}
在上述唯一不动点、仿射余循环与压缩刚性的基础上，进一步回到定理 \ref{thm:killo-godel-semidir-affine-2adic} 的规范编码情形：取
\[
\mathcal R=\NN^{(\mathbb P)},
\qquad
B=2^L>M,
\qquad
H_{\mathrm{can}}
:=
\{(r,k)\in \mathcal R\rtimes_{\Sigma}\NN:\ \mathrm{supp}(r)\subseteq\{p_1,\dots,p_k\},\ 0\le r(p_t)\le M\}.
\]
此时每个规范元的固定点都可写成严格纯周期的 \(B\)-进地址，而不同规范块之间的重合恰由 primitive 周期核控制；同一地址系统还同时锁定规范忠实性的最小进位阈值、无限事件流的前缀拼接律以及 Zeckendorf--Gödel 有限截面的 Fibonacci 仿射递推。

\begin{theorem}[规范编码子单子与有限字母自由单子的严格同构]\label{thm:xi-time-part9g-canonical-submonoid-free-monoid}
\leanverified{paper\_xi\_time\_part9g\_canonical\_submonoid\_free\_monoid}
沿用上段记号，并记
\[
\{0,\dots,M\}^{\ast}:=\bigsqcup_{k\ge 0}\{0,\dots,M\}^{k}
\]
为字母表 \(\{0,\dots,M\}\) 上的自由单子，其乘法为词连接。定义映射
\[
\Theta:H_{\mathrm{can}}\longrightarrow \{0,\dots,M\}^{\ast},
\qquad
\Theta(r,k):=\bigl(r(p_1),\dots,r(p_k)\bigr),
\]
并约定 \(\Theta(0,0)=\varnothing\) 为单位空词。则 \(H_{\mathrm{can}}\) 是
\(
\mathcal R\rtimes_{\Sigma}\NN
\)
的子单子，并且 \(\Theta\) 是单子同构。换言之，规范编码子集并非半直积中的任意稀薄子集，而恰是字母表大小 \(M+1\) 的自由单子。
\end{theorem}
\begin{proof}
先证封闭性。设
\[
(r,k),(s,\ell)\in H_{\mathrm{can}}.
\]
由定义，
\[
\mathrm{supp}(r)\subseteq\{p_1,\dots,p_k\},
\qquad
\mathrm{supp}(s)\subseteq\{p_1,\dots,p_{\ell}\}.
\]
故移位后有
\[
\mathrm{supp}(\Sigma^k s)\subseteq\{p_{k+1},\dots,p_{k+\ell}\},
\]
从而
\[
\mathrm{supp}(r+\Sigma^k s)\subseteq\{p_1,\dots,p_{k+\ell}\}.
\]
并且两段支撑区间不交，因此对 \(1\le t\le k+\ell\) 有
\[
(r+\Sigma^k s)(p_t)=
\begin{cases}
r(p_t),& 1\le t\le k,\\[2pt]
s(p_{t-k}),& k<t\le k+\ell.
\end{cases}
\]
故仍满足
\(
0\le (r+\Sigma^k s)(p_t)\le M
\)。
于是
\[
(r,k)\cdot (s,\ell)=\bigl(r+\Sigma^k s,\ k+\ell\bigr)\in H_{\mathrm{can}},
\]
且单位元 \((0,0)\) 也显然属于 \(H_{\mathrm{can}}\)。因此 \(H_{\mathrm{can}}\) 为子单子。

再证 \(\Theta\) 为双射。若
\[
\Theta(r,k)=\Theta(s,\ell),
\]
则两侧词长相同，故 \(k=\ell\)；再逐坐标比较可得
\[
r(p_t)=s(p_t)\qquad (1\le t\le k).
\]
而 \(r,s\) 在其余素数坐标上都为零，故 \(r=s\)。因此 \(\Theta\) 单射。

任取一词
\[
(a_1,\dots,a_k)\in\{0,\dots,M\}^{k},
\]
定义 \(r\in\mathcal R\) 为
\[
r(p_t)=a_t\quad(1\le t\le k),
\qquad
r(p)=0\quad(p\notin\{p_1,\dots,p_k\}).
\]
则 \((r,k)\in H_{\mathrm{can}}\) 且
\(
\Theta(r,k)=(a_1,\dots,a_k)
\)。
故 \(\Theta\) 满射。

最后比较乘法。对任意 \((r,k),(s,\ell)\in H_{\mathrm{can}}\)，由上面的分段公式，
\[
\Theta\bigl((r,k)\cdot(s,\ell)\bigr)
=
\Theta\bigl(r+\Sigma^k s,\ k+\ell\bigr)
=
\bigl(r(p_1),\dots,r(p_k),s(p_1),\dots,s(p_{\ell})\bigr),
\]
这正是词连接
\(
\Theta(r,k)\Theta(s,\ell)
\)。
因此 \(\Theta\) 为单子同构。证毕。
\end{proof}

\begin{theorem}[规范仿射元的固定点展开与纯周期 \texorpdfstring{$B$}{B}-进实现]\label{thm:xi-time-part9g-canonical-fixed-point-badic-periodic}
沿用前述 \(\Psi\)、\(\operatorname{ev}_B\)、\(T_2(E_{\min})\) 与 \(v\) 的记号。设
\[
h=(r,k)\in H_{\mathrm{can}},
\qquad
k\ge 1.
\]
则定理 \ref{thm:xi-time-part9g-affine-semigroup-unique-fixed-point} 给出的唯一不动点
\[
x_h=(1-B^k)^{-1}\operatorname{ev}_B(r)
\]
满足收敛展开
\[
x_h
=
\sum_{j\ge 0}B^{jk}\operatorname{ev}_B(r)
=
\sum_{j\ge 0}\sum_{t=1}^{k} r(p_t)\,B^{t-1+jk}v.
\]
因此 \(x_h\) 恰由长度 \(k\) 的数字块
\[
\bigl(r(p_1),\dots,r(p_k)\bigr)\in\{0,\dots,M\}^{k}
\]
作无限纯周期重复得到。

反之，对任意长度 \(k\) 的数字块
\[
(a_1,\dots,a_k)\in\{0,\dots,M\}^{k},
\]
若定义 \(r\in\mathcal R\) 满足
\[
r(p_t)=a_t\quad (1\le t\le k),
\qquad
r(p)=0\quad (p\notin\{p_1,\dots,p_k\}),
\]
则 \(h=(r,k)\in H_{\mathrm{can}}\)，且
\[
x=\sum_{j\ge 0}\sum_{t=1}^{k} a_t\,B^{t-1+jk}v
\]
正是 \(\Psi(h)\) 的唯一不动点。
\end{theorem}
\begin{proof}
由定理 \ref{thm:xi-time-part9g-affine-semigroup-unique-fixed-point}，
\[
x_h=(1-B^k)^{-1}\operatorname{ev}_B(r).
\]
由于 \(B=2^L\) 且 \(k\ge 1\)，有 \(|B^k|_2<1\)，故几何级数在 \(\ZZ_2\) 中收敛并给出
\[
(1-B^k)^{-1}=\sum_{j\ge 0}B^{jk}.
\]
另一方面，由 \(h\in H_{\mathrm{can}}\) 的支撑约束，
\[
\operatorname{ev}_B(r)=\sum_{t=1}^{k}r(p_t)\,B^{t-1}v.
\]
将两式相乘即得
\[
x_h
=
\sum_{j\ge 0}B^{jk}\operatorname{ev}_B(r)
=
\sum_{j\ge 0}\sum_{t=1}^{k} r(p_t)\,B^{t-1+jk}v,
\]
这正是对应数字块的纯周期 \(B\)-进展开。

反向地，给定任意 \((a_1,\dots,a_k)\in\{0,\dots,M\}^{k}\)，按题述定义 \(r\) 即有
\((r,k)\in H_{\mathrm{can}}\)，并满足
\[
\operatorname{ev}_B(r)=\sum_{t=1}^{k}a_t\,B^{t-1}v.
\]
于是同一几何级数计算给出
\[
(1-B^k)^{-1}\operatorname{ev}_B(r)
=
\sum_{j\ge 0}\sum_{t=1}^{k} a_t\,B^{t-1+jk}v=x.
\]
再由定理 \ref{thm:xi-time-part9g-affine-semigroup-unique-fixed-point}，\(x\) 确为 \(\Psi(h)\) 的唯一不动点。证毕。
\end{proof}

\begin{corollary}[固定点重合的 primitive word 分类与最小周期]\label{cor:xi-time-part9g-fixed-point-primitive-word-classification}
沿用定理 \ref{thm:xi-time-part9g-canonical-fixed-point-badic-periodic} 的记号。设
\[
h=(r,k),\qquad h'=(s,\ell)\in H_{\mathrm{can}},
\]
并记对应数字块为
\[
\mathbf r:=\bigl(r(p_1),\dots,r(p_k)\bigr),
\qquad
\mathbf s:=\bigl(s(p_1),\dots,s(p_\ell)\bigr).
\]
则下列条件等价：
\begin{enumerate}
  \item \(x_h=x_{h'}\)；
  \item 两条无限数字序列逐位相同，即
  \[
  \mathbf r^\infty=\mathbf s^\infty;
  \]
  \item 存在某个 primitive word \(u\)，以及正整数 \(a,b\)，使得
  \[
  \mathbf r=u^a,\qquad \mathbf s=u^b.
  \]
\end{enumerate}
特别地，若 \(\mathbf r\) 本身为 primitive word，则其对应固定点的最小周期恰为 \(k\)；因此不存在任何 \(\ell<k\) 的规范元素 \(h'\in H_{\mathrm{can}}\) 满足 \(x_{h'}=x_h\)。
\end{corollary}
\begin{proof}
由定理 \ref{thm:xi-time-part9g-canonical-fixed-point-badic-periodic}，
\[
x_h=\sum_{n\ge 0}a_n B^n v,
\qquad
x_{h'}=\sum_{n\ge 0}a_n' B^n v,
\]
其中 \((a_n)_{n\ge 0}\)、\((a_n')_{n\ge 0}\) 分别由 \(\mathbf r,\mathbf s\) 作无限纯周期重复得到，且所有 digits 都落在 \(\{0,\dots,M\}\subset[0,B-1]\) 内。

若 \(x_h=x_{h'}\)，则
\[
\sum_{n\ge 0}(a_n-a_n')B^n v=0.
\]
由于 \(v\neq 0\) 且 \(T_2(E_{\min})\cong \ZZ_2^2\) 对 \(\ZZ_2\) 无挠，上式等价于
\[
\sum_{n\ge 0}(a_n-a_n')B^n=0
\]
在 \(\ZZ_2\) 中成立。由 digits 取值于 \([0,B-1]\) 的 \(B\)-进展开唯一性，必有 \(a_n=a_n'\) 对一切 \(n\) 成立，即 \(\mathbf r^\infty=\mathbf s^\infty\)。反向蕴含显然，故 (1) 与 (2) 等价。

而对有限字的标准组合论事实表明：两条纯周期无限字相同，当且仅当它们都可写成同一个 primitive word 的整数次幂。故 (2) 与 (3) 等价。

最后，若 \(\mathbf r\) 为 primitive word，则 \(\mathbf r^\infty\) 的最小周期正是 \(k\)。因而它不可能同时等于任何长度 \(\ell<k\) 的纯周期块重复；等价地，不存在 \(\ell<k\) 的规范元素与 \(h\) 共享固定点。证毕。
\end{proof}

\begin{theorem}[规范账本表示的最小进位阈值]\label{thm:xi-time-part9g-canonical-faithful-minimal-radix-threshold}
沿用定理 \ref{thm:killo-godel-semidir-affine-2adic} 的记号，但仅假设 \(B=2^L\ge 2\)。
则
\[
\Psi|_{H_{\mathrm{can}}}\ \text{为忠实表示}
\quad\Longleftrightarrow\quad
B>M.
\]
\end{theorem}
\begin{proof}
充分性正是定理 \ref{thm:killo-godel-semidir-affine-2adic} 的结论。
若 \(B\le M\)，则取 \(k=2\)，并定义 \(r,s\in\mathcal R\) 为
\[
r(p_1)=B,\qquad r(p_2)=0,
\]
\[
s(p_1)=0,\qquad s(p_2)=1,
\]
其余素数坐标全为 \(0\)。
于是 \((r,2),(s,2)\in H_{\mathrm{can}}\)，且
\[
\operatorname{ev}_B(r)=Bv=\operatorname{ev}_B(s).
\]
故对任意 \(x\in T_2(E_{\min})\)，都有
\[
\Psi(r,2)(x)=B^2x+Bv=\Psi(s,2)(x).
\]
但 \((r,2)\neq (s,2)\)，故 \(\Psi|_{H_{\mathrm{can}}}\) 不忠实。
因此规范忠实性当且仅当 \(B>M\)。
\end{proof}

\begin{theorem}[无限事件流地址的前缀拼接与精确超度量]\label{thm:xi-time-part9g-holographic-prefix-concatenation-ultrametric}
沿用推论 \ref{cor:killo-infinite-stream-2adic-holographic-point} 的记号，并进一步假设
\[
v\notin 2T_2(E_{\min}).
\]
对任意有限事件词
\[
\mathbf a=(a_1,\dots,a_k)\in\mathcal E^k,
\]
定义 \(r_{\mathbf a}\in\mathcal R\) 为
\[
r_{\mathbf a}(p_t)=\mathrm{code}(a_t)\quad (1\le t\le k),
\qquad
r_{\mathbf a}(p)=0\quad (p\notin\{p_1,\dots,p_k\}).
\]
则对任意无限事件流 \(\mathbf e\in\mathcal E^{\NN}\) 都有
\[
\Psi(r_{\mathbf a},k)\bigl(X(\mathbf e)\bigr)=X(\mathbf a\mathbf e),
\]
其中 \(\mathbf a\mathbf e\) 表示把 \(\mathbf a\) 作为前缀接到 \(\mathbf e\) 之前。

进一步，若 \(\mathbf e,\mathbf e'\in\mathcal E^{\NN}\) 的最长公共前缀长度恰为 \(n\)，则
\[
X(\mathbf e)-X(\mathbf e')\in B^nT_2(E_{\min})\setminus B^{n+1}T_2(E_{\min}).
\]
等价地，对每个 \(n\ge 0\) 都有
\[
X(\mathbf e)\equiv X(\mathbf e')\pmod{B^nT_2(E_{\min})}
\quad\Longleftrightarrow\quad
\mathrm{pref}_n(\mathbf e)=\mathrm{pref}_n(\mathbf e').
\]
因此，全息地址的 \(B\)-进分辨率层级与事件流的前缀超度量完全一致。
\end{theorem}
\begin{proof}
由定义，
\[
\Psi(r_{\mathbf a},k)\bigl(X(\mathbf e)\bigr)
=
B^k\sum_{u\ge 1}\mathrm{code}(e_u)\,B^{u-1}v
+\sum_{t=1}^{k}\mathrm{code}(a_t)\,B^{t-1}v.
\]
整理得
\[
\Psi(r_{\mathbf a},k)\bigl(X(\mathbf e)\bigr)
=
\sum_{t=1}^{k}\mathrm{code}(a_t)\,B^{t-1}v
+\sum_{u\ge 1}\mathrm{code}(e_u)\,B^{u+k-1}v
=
X(\mathbf a\mathbf e).
\]

对第二部分，由编码的单射性，\(\mathbf e,\mathbf e'\) 的最长公共前缀长度与其编码序列的最长公共前缀长度相同。
又因 \(v\notin 2T_2(E_{\min})\) 蕴含 \(v\notin BT_2(E_{\min})\)，故可将定理 \ref{thm:killo-2adic-holographic-exact-cylinder-separation} 应用于编码序列，立得
\[
X(\mathbf e)-X(\mathbf e')\in B^nT_2(E_{\min})\setminus B^{n+1}T_2(E_{\min}).
\]
由同一定理中的柱分离等价式，也立即得到模 \(B^nT_2(E_{\min})\) 同余与长度 \(n\) 前缀相等之间的等价关系。
\end{proof}

\begin{theorem}[无限事件流地址在自然地址线上的严格前缀等距性]\label{thm:xi-time-part9g-holographic-prefix-isometry-on-line}
沿用推论 \ref{cor:killo-infinite-stream-2adic-holographic-point} 的设定，并记
\[
T_v:=\ZZ_2 v\subset T_2(E_{\min}).
\]
对任意 \(\mathbf e,\mathbf e'\in\mathcal E^{\NN}\)，定义首个分歧位置
\[
N(\mathbf e,\mathbf e')
:=
\min\{t\ge 1:\ e_t\neq e_t'\}\in\NN\cup\{\infty\}.
\]
若 \(\mathbf e\neq \mathbf e'\)，则
\[
\boxed{\;
X(\mathbf e)-X(\mathbf e')
\in
B^{N(\mathbf e,\mathbf e')-1}T_v\setminus B^{N(\mathbf e,\mathbf e')}T_v.\;}
\]
因此，若定义前缀超度量
\[
d_{\mathrm{pref}}(\mathbf e,\mathbf e')
:=
\begin{cases}
0,& \mathbf e=\mathbf e',\\[2pt]
|B|_2^{\,N(\mathbf e,\mathbf e')-1},& \mathbf e\neq \mathbf e',
\end{cases}
\]
以及 \(T_v\) 上的 \(B\)-进超度量
\[
d_B(x,y)
:=
\begin{cases}
0,& x=y,\\[2pt]
|B|_2^{\,\nu_B(x-y)},& x\neq y,
\end{cases}
\qquad
\nu_B(z):=\max\{n\ge 0:\ z\in B^nT_v\},
\]
则
\[
\boxed{\;
d_B\bigl(X(\mathbf e),X(\mathbf e')\bigr)=d_{\mathrm{pref}}(\mathbf e,\mathbf e').\;}
\]
进一步，对任意有限前缀
\[
a=(e_1,\dots,e_n)\in\mathcal E^n
\]
记
\[
[a]:=\{\mathbf e\in\mathcal E^{\NN}:\ e_1,\dots,e_n\ \text{固定}\},
\qquad
X_n(a):=\sum_{t=1}^{n}\mathrm{code}(e_t)\,B^{t-1}v.
\]
则
\[
\boxed{\;
X([a])=\bigl(X_n(a)+B^nT_v\bigr)\cap X(\mathcal E^{\NN}).\;}
\]
因此，长度 \(n\) 的 cylinder 在商模 \(T_v/B^nT_v\) 中给出两两不交的剩余类。
\end{theorem}
\begin{proof}
把 \(\mathbf e,\mathbf e'\) 通过编码识别为 digit 序列
\[
a_t:=\mathrm{code}(e_t),
\qquad
a_t':=\mathrm{code}(e_t').
\]
若 \(\mathbf e\neq \mathbf e'\)，记 \(N:=N(\mathbf e,\mathbf e')\)。则
\[
X(\mathbf e)-X(\mathbf e')
=
\sum_{t\ge N}(a_t-a_t')B^{t-1}v
=
B^{N-1}uv,
\]
其中
\[
u:=(a_N-a_N')+\sum_{t>N}(a_t-a_t')B^{t-N}\in\ZZ_2.
\]
由于 \(0\le a_t,a_t'\le M<B\) 且 \(a_N\neq a_N'\)，有
\[
u\equiv a_N-a_N'\not\equiv 0\pmod B.
\]
故 \(u\notin B\ZZ_2\)，于是
\[
uv\in T_v\setminus B T_v.
\]
从而
\[
X(\mathbf e)-X(\mathbf e')
\in
B^{N-1}T_v\setminus B^N T_v,
\]
即得第一条。\(d_B\) 与 \(d_{\mathrm{pref}}\) 的等式只是该 \(B\)-进阶数结论的直接重述。

若 \(\mathbf e\in[a]\)，则由推论 \ref{cor:killo-infinite-stream-2adic-holographic-point} 的前缀截断公式，
\[
X(\mathbf e)-X_n(a)=\sum_{t>n}\mathrm{code}(e_t)\,B^{t-1}v\in B^nT_v,
\]
故
\[
X([a])\subseteq X_n(a)+B^nT_v.
\]
反之，设
\[
X(\mathbf e)\in \bigl(X_n(a)+B^nT_v\bigr)\cap X(\mathcal E^{\NN}).
\]
任取 \(\mathbf f\in[a]\)，则同样有
\[
X(\mathbf f)-X_n(a)\in B^nT_v,
\]
从而
\[
X(\mathbf e)-X(\mathbf f)\in B^nT_v.
\]
若 \(\mathbf e\notin[a]\)，则 \(\mathbf e\) 与 \(\mathbf f\) 在某个 \(j\le n\) 处首次分歧；由第一条可知
\[
X(\mathbf e)-X(\mathbf f)\in B^{j-1}T_v\setminus B^jT_v,
\]
这与其属于 \(B^nT_v\) 矛盾。故 \(\mathbf e\in[a]\)，于是
\[
X([a])=\bigl(X_n(a)+B^nT_v\bigr)\cap X(\mathcal E^{\NN}).
\]
若两个长度 \(n\) 的前缀 \(a,b\) 不同，则任取 \(\mathbf e\in[a]\)、\(\mathbf e'\in[b]\)，它们在某个 \(j\le n\) 处首次分歧，故
\[
X(\mathbf e)-X(\mathbf e')\notin B^nT_v.
\]
因此对应的 \(B^nT_v\)-剩余类必然互不相交。证毕。
\end{proof}

\begin{theorem}[Zeckendorf--Gödel 有限截面的 Fibonacci 仿射递推]\label{thm:xi-time-part9g-zg-primorial-fibonacci-affine-recursion}
取二值事件集
\[
\mathcal E=\{0,1\},
\qquad
\mathrm{code}(0)=0,\qquad \mathrm{code}(1)=1,
\]
并固定 \(B=2^L>1\) 与原始向量
\[
v\in T_2(E_{\min})\setminus 2T_2(E_{\min}).
\]
对 \(n\ge 1\) 定义 admissible 前缀集
\[
\mathcal A_n
:=
\bigl\{
(z_1,\dots,z_n)\in\{0,1\}^n:\ z_tz_{t+1}=0\ \ (1\le t<n)
\bigr\},
\]
\leanverified{paper\_xi\_time\_part9g\_zg\_primorial\_fibonacci\_affine\_recursion}
并约定 \(\mathcal A_0:=\{\varnothing\}\)。
再定义
\[
\mathcal D_n:=
\left\{
\prod_{t=1}^{n}p_t^{z_t}:\ (z_1,\dots,z_n)\in\mathcal A_n
\right\}\subset \mathrm{ZG},
\]
\[
\mathcal X_n:=
\left\{
\sum_{t=1}^{n}z_t\,B^{t-1}v:\ (z_1,\dots,z_n)\in\mathcal A_n
\right\}\subset T_2(E_{\min}),
\]
并约定 \(\mathcal D_0:=\{1\}\)、\(\mathcal X_0:=\{0\}\)。
则成立：
\begin{enumerate}
  \item 对每个 \(n\ge 0\)，对应
  \[
  \prod_{t=1}^{n}p_t^{z_t}
  \longmapsto
  \sum_{t=1}^{n}z_t\,B^{t-1}v
  \]
  给出 \(\mathcal D_n\) 与 \(\mathcal X_n\) 的双射；
  \item 对每个 \(n\ge 2\)，有严格不交并分解
  \[
  \mathcal X_n
  =
  \mathcal X_{n-1}\ \sqcup\ \bigl(B^{n-1}v+\mathcal X_{n-2}\bigr);
  \]
  \item 因而
  \[
  |\mathcal D_n|=|\mathcal X_n|=F_{n+2}
  \qquad (n\ge 0).
  \]
\end{enumerate}
\end{theorem}
\begin{proof}
由唯一素因子分解，映射
\[
(z_1,\dots,z_n)\longmapsto \prod_{t=1}^{n}p_t^{z_t}
\]
给出 \(\mathcal A_n\) 与 \(\mathcal D_n\) 的双射。

再看地址像。
对任意 \((z_1,\dots,z_n)\in\mathcal A_n\)，附加零尾得到 \(\widetilde{\mathbf z}\in\{0,1\}^{\NN}\)。
由于 \(M=1<B\) 且 \(v\notin 2T_2(E_{\min})\)，定理 \ref{thm:xi-time-part9g-holographic-prefix-concatenation-ultrametric} 表明不同前缀对应不同的全息地址。
因此
\[
(z_1,\dots,z_n)\longmapsto \sum_{t=1}^{n}z_t\,B^{t-1}v
\]
在 \(\mathcal A_n\) 上为单射，从而给出 \(\mathcal A_n\) 与 \(\mathcal X_n\) 的双射。
与上一段合并即得第 \(1\) 条。

现证递推。
对任意 admissible 词 \((z_1,\dots,z_n)\in\mathcal A_n\)，若 \(z_n=0\)，则 \((z_1,\dots,z_{n-1})\in\mathcal A_{n-1}\)，其地址落在 \(\mathcal X_{n-1}\)；
若 \(z_n=1\)，则 admissibility 强制 \(z_{n-1}=0\)，故 \((z_1,\dots,z_{n-2})\in\mathcal A_{n-2}\)，并且相应地址恰为
\[
B^{n-1}v+\mathcal X_{n-2}.
\]
于是
\[
\mathcal X_n=\mathcal X_{n-1}\cup \bigl(B^{n-1}v+\mathcal X_{n-2}\bigr).
\]
该并实际上不交：若二者有公共点，则会给出两个不同的 admissible 前缀产生同一地址，这与第 \(1\) 条的单射性矛盾。

取基数得
\[
|\mathcal X_n|=|\mathcal X_{n-1}|+|\mathcal X_{n-2}|
\qquad (n\ge 2),
\]
并且
\[
|\mathcal X_0|=1,\qquad |\mathcal X_1|=2.
\]
故 \(|\mathcal X_n|=F_{n+2}\)。
再由第 \(1\) 条，\(|\mathcal D_n|=|\mathcal X_n|\)，于是第 \(3\) 条成立。
\end{proof}

\endinput

\paragraph{规范纯周期固定点的层计数、primitive 周期核与地址分形维数}
在定理 \ref{thm:xi-time-part9g-canonical-fixed-point-badic-periodic} 与推论 \ref{cor:xi-time-part9g-fixed-point-primitive-word-classification} 已把规范固定点的纯周期展开及其碰撞机制锁定之后，还可把同一地址系统的层计数、最小周期、全字母闭包与 prime-register 预算进一步压缩为如下五条闭合结论。

\begin{theorem}[规范纯周期固定点层的精确基数与前缀球分层]\label{thm:xi-time-part9gb-canonical-fixedpoint-count-prefix-stratification}
\leanverified{paper\_xi\_time\_part9gb\_canonical\_fixedpoint\_count\_prefix\_stratification}
沿用定理 \ref{thm:killo-godel-semidir-affine-2adic} 与定理
\ref{thm:xi-time-part9g-canonical-fixed-point-badic-periodic} 的记号。记
\[
T_v:=\ZZ_2v\subset T_2(E_{\min}),
\qquad
H_{\mathrm{can}}(k):=\{(r,k)\in H_{\mathrm{can}}\}
\qquad (k\ge 1).
\]
对每个
\[
h=(r,k)\in H_{\mathrm{can}}(k),
\]
记对应长度 \(k\) 的数字块为
\[
a(h):=\bigl(r(p_1),\dots,r(p_k)\bigr)\in\{0,\dots,M\}^k.
\]
则成立：
\begin{enumerate}
  \item \(\Psi(h)\) 在 \(T_v\) 上有且仅有一个固定点
  \[
  x_h=(1-B^k)^{-1}\operatorname{ev}_B(r)
  =
  \sum_{j\ge 0}\sum_{t=1}^{k}r(p_t)\,B^{jk+t-1}v.
  \]
  换言之，\(x_h\) 的 \(B\)-进展开恰为块 \(a(h)\) 的无限纯周期重复。
  \item 映射
  \[
  H_{\mathrm{can}}(k)\longrightarrow T_v,
  \qquad
  h\longmapsto x_h
  \]
  为单射，且
  \[
  \#\{x_h:\ h\in H_{\mathrm{can}}(k)\}=(M+1)^k.
  \]
  \item 对任意 \(0\le n\le k\) 及任意 \(h,h'\in H_{\mathrm{can}}(k)\)，有
  \[
  x_h-x_{h'}\in B^nT_v
  \quad\Longleftrightarrow\quad
  a(h)\ \text{与}\ a(h')\ \text{的前 }n\text{ 位相同}.
  \]
\end{enumerate}
\end{theorem}
\begin{proof}
第 \(1\) 条正是定理 \ref{thm:xi-time-part9g-canonical-fixed-point-badic-periodic} 的内容。

对第 \(2\) 条，由定理 \ref{thm:xi-time-part9g-affine-semigroup-unique-fixed-point}，在固定层级 \(k\) 上若 \(x_h=x_{h'}\)，则必有 \(h=h'\)，故 \(h\mapsto x_h\) 为单射。另一方面，\(H_{\mathrm{can}}(k)\) 的每个元素恰由 \(k\) 个彼此独立的 digits
\[
r(p_t)\in\{0,\dots,M\}
\qquad (1\le t\le k)
\]
唯一决定，因此
\[
|H_{\mathrm{can}}(k)|=(M+1)^k.
\]
于是固定点层的基数正为 \((M+1)^k\)。

再证第 \(3\) 条。对 \(h=(r,k)\)，记
\[
a(h)^\infty=(a_1,a_2,\dots)\in\{0,\dots,M\}^{\NN}
\]
为块 \(a(h)\) 的无限纯周期重复，则由第 \(1\) 条，
\[
x_h=\sum_{m\ge 1}a_mB^{m-1}v.
\]
对 \(h'\) 同理。应用定理 \ref{thm:xi-time-part9g-holographic-prefix-isometry-on-line} 于两条无限 digit 序列 \(a(h)^\infty\) 与 \(a(h')^\infty\)，得到
\[
x_h-x_{h'}\in B^nT_v
\quad\Longleftrightarrow\quad
a(h)^\infty,\ a(h')^\infty\ \text{的前 }n\text{ 位相同}.
\]
由于 \(0\le n\le k\)，上述前 \(n\) 位恰好就是块 \(a(h)\) 与 \(a(h')\) 的前 \(n\) 位，故结论成立。证毕。
\end{proof}

\begin{theorem}[规范纯周期固定点碰撞的 primitive 核与 Möbius 计数]\label{thm:xi-time-part9gb-fixedpoint-primitive-kernel-mobius}
\leanverified{paper\_xi\_time\_part9gb\_fixedpoint\_primitive\_kernel\_mobius}
沿用定理 \ref{thm:xi-time-part9gb-canonical-fixedpoint-count-prefix-stratification} 的记号。对任意长度 \(k\ge 1\) 的数字块
\[
a=(a_1,\dots,a_k)\in\{0,\dots,M\}^k,
\]
定义其对应纯周期点
\[
x(a):=\sum_{j\ge 0}\sum_{t=1}^{k}a_t\,B^{jk+t-1}v\in T_v.
\]
再记
\[
\mathcal F_k:=\{x_h:\ h\in H_{\mathrm{can}}(k)\}\subset T_v,
\]
以及最小周期层
\[
\Pi_k:=
\left\{
x\in \mathcal F_k:\ 
x\notin \mathcal F_d\ \text{对一切真因子 }d\mid k
\right\}.
\]
则：
\begin{enumerate}
  \item 对任意
  \[
  a\in\{0,\dots,M\}^k,
  \qquad
  b\in\{0,\dots,M\}^{\ell},
  \]
  有
  \[
  x(a)=x(b)
  \iff
  a^\infty=b^\infty
  \iff
  \exists\,u\ \text{为 primitive word},\ \exists\,r,s\ge 1:
  \ a=u^r,\ b=u^s.
  \]
  \item 对每个 \(k\ge 1\)，有不交并分解
  \[
  \mathcal F_k=\bigsqcup_{d\mid k}\Pi_d,
  \qquad
  |\mathcal F_k|=\sum_{d\mid k}|\Pi_d|.
  \]
  \item 若记
  \[
  P_M(k):=\#\{u\in\{0,\dots,M\}^k:\ u\ \text{为 primitive word}\},
  \]
  则
  \[
  |\Pi_k|=P_M(k)=\sum_{d\mid k}\mu(d)(M+1)^{k/d}.
  \]
  特别地，最小周期恰为 \(k\) 的规范纯周期固定点个数正是长度 \(k\) 的 primitive 字个数。
\end{enumerate}
\end{theorem}
\begin{proof}
第 \(1\) 条正是推论 \ref{cor:xi-time-part9g-fixed-point-primitive-word-classification} 的内容。

现证第 \(2\) 条。任取 \(x\in \mathcal F_k\)。由定理
\ref{thm:xi-time-part9gb-canonical-fixedpoint-count-prefix-stratification}，\(x\) 对应某个长度 \(k\) 的纯周期块 \(a\)。取 \(d\) 为无限字 \(a^\infty\) 的最小周期，则 \(d\mid k\)，并存在长度 \(d\) 的 primitive word \(u\) 使
\[
a=u^{k/d}.
\]
于是 \(x=x(u)\in \Pi_d\)。这表明
\[
\mathcal F_k\subseteq \bigsqcup_{d\mid k}\Pi_d.
\]
反向地，若 \(x\in \Pi_d\) 且 \(d\mid k\)，则 \(x=x(u)\) 对某个长度 \(d\) 的 primitive word \(u\) 成立；将其写成长度 \(k\) 的块 \(u^{k/d}\)，便得 \(x=x(u^{k/d})\in \mathcal F_k\)。故上述包含反向亦成立，且不同 \(d\) 的 \(\Pi_d\) 由最小周期定义两两不交，从而得到不交并分解及计数恒等式。

对第 \(3\) 条，定理 \ref{thm:xi-time-part9gb-canonical-fixedpoint-count-prefix-stratification} 已给出
\[
|\mathcal F_k|=(M+1)^k.
\]
与第 \(2\) 条合并得
\[
(M+1)^k=\sum_{d\mid k}|\Pi_d|.
\]
对该整除卷积施行 Möbius 反演，即得
\[
|\Pi_k|=\sum_{d\mid k}\mu(d)(M+1)^{k/d}.
\]
另一方面，第 \(1\) 条表明最小周期恰为 \(k\) 的点与长度 \(k\) 的 primitive words 一一对应，故 \(|\Pi_k|=P_M(k)\)。证毕。
\end{proof}

\begin{theorem}[全字母地址像的 \texorpdfstring{$2$}{2}-进自相似 Cantor 几何]\label{thm:xi-time-part9gb-full-alphabet-address-cantor-dimension}
沿用前述记号，定义全字母地址像
\[
\mathcal A_M
:=
\left\{
\sum_{t\ge 1}a_tB^{t-1}v:\ a_t\in\{0,\dots,M\}
\right\}
\subset T_v\subset T_2(E_{\min}),
\]
以及规范固定点全集
\[
\operatorname{Fix}_{\mathrm{can}}
:=
\bigcup_{k\ge 1}\mathcal F_k.
\]
则成立：
\begin{enumerate}
  \item \(\operatorname{Fix}_{\mathrm{can}}\) 在 \(\mathcal A_M\) 中稠密，亦即
  \[
  \overline{\operatorname{Fix}_{\mathrm{can}}}=\mathcal A_M.
  \]
  \item \(\mathcal A_M\) 是由 \((M+1)\) 个仿射收缩
  \[
  f_d(x):=dv+Bx
  \qquad (d\in\{0,\dots,M\})
  \]
  生成的严格自相似紧集，并满足不交分解
  \[
  \mathcal A_M=\bigsqcup_{d=0}^{M}f_d(\mathcal A_M).
  \]
  \item 若 \(M\ge 1\)，则 \(\mathcal A_M\) 是 \(T_v\) 中的紧、全不连通且无孤立点的 Cantor 型子集；并且在 \(T_2(E_{\min})\) 中的 Hausdorff 维数与 Minkowski 维数一致，满足
  \[
  \dim_{\mathrm H}(\mathcal A_M)
  =
  \dim_{\mathrm M}(\mathcal A_M)
  =
  \frac{\log(M+1)}{L\log 2}.
  \]
  此外，\(\mathcal A_M\) 在 \(T_2(E_{\min})\) 的 Haar 测度下为零测集。
\end{enumerate}
\end{theorem}
\begin{proof}
先证第 \(1\) 条。任取
\[
x=\sum_{t\ge 1}a_tB^{t-1}v\in \mathcal A_M.
\]
对每个 \(k\ge 1\)，把前 \(k\) 个 digits 周期化，定义
\[
x^{(k)}
:=
\sum_{j\ge 0}\sum_{t=1}^{k}a_t\,B^{jk+t-1}v.
\]
由定理 \ref{thm:xi-time-part9gb-canonical-fixedpoint-count-prefix-stratification}，\(x^{(k)}\in\mathcal F_k\subset \operatorname{Fix}_{\mathrm{can}}\)。同时 \(x^{(k)}\) 与 \(x\) 的前 \(k\) 个 \(B\)-进 digits 完全一致，故
\[
x^{(k)}-x\in B^kT_v.
\]
由于 \(k\to\infty\) 时 \(B^kT_v\) 在 \(2\)-进拓扑下收缩到 \(\{0\}\)，遂得 \(x^{(k)}\to x\)。这表明 \(\mathcal A_M\subseteq \overline{\operatorname{Fix}_{\mathrm{can}}}\)。反向包含显然，因为每个规范固定点本身都属于 \(\mathcal A_M\)。故闭包等式成立。

对第 \(2\) 与第 \(3\) 条，只需把定理 \ref{thm:xi-time-part62d-hologram-image-cantor-homeomorphism} 与定理 \ref{thm:xi-time-part62d-hologram-image-zero-haar-dimension} 应用于有限事件集
\[
E=\{0,\dots,M\},
\qquad
\mathrm{code}=\mathrm{id}_E.
\]
此时对应像集 \(\Lambda_{E,v}\) 恰为 \(\mathcal A_M\)。于是严格自相似分解、紧性、全不连通性、\(M\ge 1\) 时的无孤立点、维数公式以及 Haar 零测性同时成立。证毕。
\end{proof}

\begin{theorem}[Zeckendorf--Gödel 地址像的 Fibonacci 分形维数]\label{thm:xi-time-part9gb-zg-address-fractal-dimension}
取二值字母表 \(\{0,1\}\)，并定义 Zeckendorf--Gödel 地址像
\[
\mathcal A_{\mathrm{ZG}}
:=
\left\{
\sum_{t\ge 1}z_tB^{t-1}v:\ 
z_t\in\{0,1\},\ z_tz_{t+1}=0
\right\}
\subset T_v.
\]
则 \(\mathcal A_{\mathrm{ZG}}\) 正是禁止相邻 \(1\) 子移位在地址线上的 \(2\)-进自仿射吸引子，并满足：
\begin{enumerate}
  \item \(\mathcal A_{\mathrm{ZG}}\) 满足严格自仿射方程
  \[
  \mathcal A_{\mathrm{ZG}}
  =
  B\mathcal A_{\mathrm{ZG}}
  \ \sqcup\
  \bigl(v+B^2\mathcal A_{\mathrm{ZG}}\bigr).
  \]
  \item 在尺度 \(B^{-n}\) 上，\(\mathcal A_{\mathrm{ZG}}\) 恰分解为 \(F_{n+2}\) 个两两不交的 cylinder，其中 \(F_n\) 为 Fibonacci 数列。
  \item \(\mathcal A_{\mathrm{ZG}}\) 在 \(T_2(E_{\min})\) 中的 Hausdorff 维数与 Minkowski 维数一致，且
  \[
  \dim_{\mathrm H}(\mathcal A_{\mathrm{ZG}})
  =
  \dim_{\mathrm M}(\mathcal A_{\mathrm{ZG}})
  =
  \frac{\log\varphi}{L\log 2}.
  \]
\end{enumerate}
\end{theorem}
\begin{proof}
第 \(1\) 条正是定理 \ref{thm:xi-zg-address-binary-self-affine} 的自仿射分解。

对第 \(2\) 与第 \(3\) 条，定理 \ref{thm:xi-zg-address-minkowski-dimension-content} 已证明：在地址线 \(T_v\) 上，\(\mathcal A_{\mathrm{ZG}}\) 于尺度 \(B^{-n}\) 的最少球覆盖数恰为
\[
N_{B,v}(n)=F_{n+2},
\]
并且
\[
\dim_{\mathrm H}(\mathcal A_{\mathrm{ZG}})
=
\dim_{\mathrm M}(\mathcal A_{\mathrm{ZG}})
=
\frac{\log\varphi}{\log B}
=
\frac{\log\varphi}{L\log 2}.
\]
最后，由定理 \ref{thm:xi-time-part62d-hologram-image-zero-haar-dimension} 证明中关于 rank-one 地址线嵌入的双 Lipschitz 论证可知，\(\mathcal A_{\mathrm{ZG}}\) 在环境空间 \(T_2(E_{\min})\) 中的 Hausdorff 与 Minkowski 维数，等于其在地址线 \(T_v\) 中的相应维数。因此上述数值在 \(T_2(E_{\min})\) 中保持不变，结论成立。证毕。
\end{proof}

\begin{corollary}[prime-register 最坏层成本的地址维数补码解释]\label{cor:xi-time-part9gb-prime-register-dimension-complement}
若满足定理 \ref{thm:xi-layered-prime-register-sharp-natural-extension} 的最坏情形假设，即每一层都存在真实分支，则分层 prime-register 外置的最小层成本恰为每层一枚有效 prime slice。

把每一层 \(B=2^L\) 的地址槽位视为总宽度 \(L\) 比特，则：
\begin{enumerate}
  \item 对全字母地址像 \(\mathcal A_M\)，长度 \(k\) 的 cylinder 数为 \((M+1)^k\)，故层信息率为
  \[
  \lim_{k\to\infty}\frac{1}{k}\log_2 (M+1)^k=\log_2(M+1).
  \]
  由定理 \ref{thm:xi-time-part9gb-full-alphabet-address-cantor-dimension}，
  \[
  L\,\dim_{\mathrm H}(\mathcal A_M)=\log_2(M+1),
  \]
  从而其地址维数补码恰为
  \[
  L\bigl(1-\dim_{\mathrm H}(\mathcal A_M)\bigr)
  =
  L-\log_2(M+1).
  \]
  \item 对 Zeckendorf--Gödel 地址像 \(\mathcal A_{\mathrm{ZG}}\)，长度 \(k\) 的允许词数为 \(F_{k+2}\)，故层信息率为
  \[
  \lim_{k\to\infty}\frac{1}{k}\log_2 F_{k+2}=\log_2\varphi.
  \]
  由定理 \ref{thm:xi-time-part9gb-zg-address-fractal-dimension}，
  \[
  L\,\dim_{\mathrm H}(\mathcal A_{\mathrm{ZG}})=\log_2\varphi,
  \]
  因而其地址维数补码恰为
  \[
  L\bigl(1-\dim_{\mathrm H}(\mathcal A_{\mathrm{ZG}})\bigr)
  =
  L-\log_2\varphi.
  \]
\end{enumerate}
因此，在最坏情形下必须逐层引入的 prime-register 外置，可在地址几何上解释为：它所补偿的正是每个 \(B\)-进层片中未被内部符号分裂吸收的维数缺口。
\end{corollary}
\begin{proof}
定理 \ref{thm:xi-layered-prime-register-sharp-natural-extension} 已给出最坏情形下一层一枚有效 prime slice 的必要且可达下界。

对 \(\mathcal A_M\)，长度 \(k\) 的可用地址前缀恰为全部 \(k\)-位 blocks，故其每层信息率为 \(\log_2(M+1)\) 比特；再由定理 \ref{thm:xi-time-part9gb-full-alphabet-address-cantor-dimension} 的维数公式，
\[
\dim_{\mathrm H}(\mathcal A_M)=\frac{\log(M+1)}{L\log 2},
\]
两式等价于
\[
L\,\dim_{\mathrm H}(\mathcal A_M)=\log_2(M+1).
\]
于是补码公式立即成立。

对 \(\mathcal A_{\mathrm{ZG}}\)，定理 \ref{thm:xi-time-part9gb-zg-address-fractal-dimension} 的第 \(2\) 条给出长度 \(k\) 的允许词数为 \(F_{k+2}\)，由 Fibonacci 增长率
\[
F_{k+2}\sim \frac{\varphi^{k+2}}{\sqrt5}
\]
可得其每层信息率为 \(\log_2\varphi\)。再把同一定理中的维数公式
\[
\dim_{\mathrm H}(\mathcal A_{\mathrm{ZG}})=\frac{\log\varphi}{L\log 2}
\]
改写为
\[
L\,\dim_{\mathrm H}(\mathcal A_{\mathrm{ZG}})=\log_2\varphi,
\]
便得到第二条补码公式。证毕。
\end{proof}

\endinput

在规范 \(B\)-进实现的忠实性与纯周期地址化之外，同一作用还满足一个逆极限意义下的唯一延拓性：有限层上不存在任何额外的非线性自由度可在 Tate 模块上重新出现。

\begin{theorem}[跨分辨率仿射刚性]\label{thm:xi-time-part9g-cross-resolution-affine-rigidity}
\leanverified{paper\_xi\_time\_part9g\_cross\_resolution\_affine\_rigidity}
沿用定理 \ref{thm:killo-godel-semidir-affine-2adic} 与定理
\ref{thm:xi-time-part9g-canonical-submonoid-free-monoid} 的记号。记
\[
T:=T_2(E_{\min})=\varprojlim_{n}E_{\min}[2^n],
\qquad
\Psi(r,k)(x)=B^k x+\operatorname{ev}_B(r)
\]
为 \(H_{\mathrm{can}}\) 在 \(T\) 上的规范仿射作用。设
\[
\widetilde\Psi:H_{\mathrm{can}}\longrightarrow \operatorname{Homeo}(T)
\]
是另一连续单子作用，并满足对每个 \(n\ge 1\)：
\begin{enumerate}
  \item \(\widetilde\Psi\) 在有限层 \(E_{\min}[2^n]\) 上诱导出作用
  \(\widetilde\Psi_n\)；
  \item 该诱导作用与截断映射 \(T\to E_{\min}[2^n]\) 交换；
  \item \(\widetilde\Psi_n=\Psi_n\)，其中 \(\Psi_n\) 表示 \(\Psi\) 在
  \(E_{\min}[2^n]\) 上的约化。
\end{enumerate}
则
\[
\widetilde\Psi=\Psi.
\]
换言之，与整条分辨率塔相容且在每个有限层都与规范仿射作用一致的任何连续动力学，在逆极限上都只能等于该规范 \(2\)-进仿射系统。
\end{theorem}
\begin{proof}
记自然投影为
\[
\pi_n:T\longrightarrow E_{\min}[2^n].
\]
任取 \(h\in H_{\mathrm{can}}\) 与 \(x\in T\)。由假设，
\[
\pi_n\bigl(\widetilde\Psi(h)(x)\bigr)
=
\widetilde\Psi_n(h)\bigl(\pi_n(x)\bigr)
=
\Psi_n(h)\bigl(\pi_n(x)\bigr)
=
\pi_n\bigl(\Psi(h)(x)\bigr)
\qquad (n\ge 1).
\]
因此 \(\widetilde\Psi(h)(x)\) 与 \(\Psi(h)(x)\) 在全部有限层坐标上完全一致。

而 \(T=\varprojlim_n E_{\min}[2^n]\) 的元素由所有有限层坐标唯一确定，故
\[
\widetilde\Psi(h)(x)=\Psi(h)(x).
\]
由于 \(h\) 与 \(x\) 任意，遂得 \(\widetilde\Psi=\Psi\)。证毕。
\end{proof}

\endinput

\paragraph{分辨率塔的 Haar 极限与 Tate 模块中的秩分裂}
在跨分辨率仿射刚性已经锁定之后，H.159 的 dyadic 分辨率塔与 H.160 的规范 Gödel 半直积仿射表示在同一 Tate 模块中还强制出一条更细的测度--秩--几何分裂链：有限扭点层的均匀分布在逆极限上闭合为 full-rank Haar 几何，而任何规范语义地址像仍只能坍缩到由单一地址向量生成的一条 rank-one \(B\)-进 Cantor 子系统。下列诸条把这一分裂压缩为可直接调用的闭式结论。

\begin{theorem}[分辨率塔的 Haar 极限律]\label{thm:xi-time-part9gc-resolution-tower-haar-limit}
设
\[
T:=T_2(E_{\min})
\cong
\varprojlim_n E_{\min}[2^n]
\cong
\ZZ_2^2
\]
为推论 \ref{cor:killo-2adic-tate-hypercube-limit} 给出的逆极限识别。对每个 \(n\ge 1\)，记 \(\mu_n\) 为有限群 \(E_{\min}[2^n]\) 上的均匀概率测度。则：
\begin{enumerate}
  \item 在乘二投影
  \[
  [2]:E_{\min}[2^{n+1}]\longrightarrow E_{\min}[2^n]
  \]
  下，测度族 \((\mu_n)_{n\ge 1}\) 严格相容；
  \item 存在唯一 Borel 概率测度 \(\mu_\infty\) 于 \(T\) 上，使得对每个 \(n\) 与每个模 \(2^n\) 余类
  \[
  C(a,b;n):=
  \left\{
  x\in T:\ x\equiv (a,b)\pmod{2^nT}
  \right\},
  \qquad
  (a,b)\in (\ZZ/2^n\ZZ)^2,
  \]
  都有
  \[
  \mu_\infty\bigl(C(a,b;n)\bigr)=4^{-n};
  \]
  \item \(\mu_\infty\) 正是紧群 \(T\cong \ZZ_2^2\) 的归一化 Haar 概率测度。
\end{enumerate}
\leanverified{paper\_xi\_time\_part9gc\_resolution\_tower\_haar\_limit}
\end{theorem}
\begin{proof}
由推论 \ref{cor:killo-2adic-tate-hypercube-limit}，每一层都可坐标化为
\[
E_{\min}[2^n]\cong (\ZZ/2^n\ZZ)^2,
\]
且乘二映射恰对应模约化
\[
(a\bmod 2^{n+1},\,b\bmod 2^{n+1})
\longmapsto
(a\bmod 2^n,\,b\bmod 2^n).
\]
因此，对每个深度 \(n\) 的余类 \(C(a,b;n)\)，其在下一层恰有 \(4\) 个 lift；于是
\[
\mu_{n+1}\bigl([2]^{-1}\{(a,b)\}\bigr)
=
4\cdot 4^{-(n+1)}
=
4^{-n}
=
\mu_n\bigl(\{(a,b)\}\bigr),
\]
这就证明了相容性。

现以 cylinder 质量
\[
\mu_\infty\bigl(C(a,b;n)\bigr):=4^{-n}
\]
定义逆极限上的柱测度。该定义是相容的，因为每个深度 \(n\) 的 cylinder 正好分裂为 \(4\) 个深度 \(n+1\) 的 cylinder。故由 cylinder 扩张唯一得到 \(T\) 上的 Borel 概率测度 \(\mu_\infty\)。

最后，由 \(T\cong \ZZ_2^2\) 可知每个 \(C(a,b;n)\) 都是开子群 \(2^nT\) 的一个陪集，而这类陪集在任意平移下彼此置换。由于 \(\mu_\infty\) 在同一深度的全部 \(4^n\) 个陪集上取相同质量 \(4^{-n}\)，故 \(\mu_\infty\) 在全部基本 clopen 陪集上都满足平移不变性。此类集合生成 \(T\) 的 Borel \(\sigma\)-代数，故 \(\mu_\infty\) 即为归一化 Haar 概率测度。证毕。
\end{proof}

\begin{corollary}[分辨率塔的 \texorpdfstring{$2$}{2}-进秩二刚性]\label{cor:xi-time-part9gc-resolution-tower-rank-two-rigidity}
\leanverified{paper\_xi\_time\_part9gc\_resolution\_tower\_rank\_two\_rigidity}
设 \(M\) 为自由 \(\ZZ_2\)-模，且存在对一切 \(n\ge 1\) 都与模 \(2^n\) 商相容的层同构
\[
M/2^nM\cong E_{\min}[2^n].
\]
若 \(M\cong \ZZ_2^r\)，则必有
\[
r=2.
\]
特别地，任何试图以单一 \(2\)-进地址轴 \(\ZZ_2\) 忠实实现整条分辨率塔的方案都不可能成立。
\end{corollary}
\begin{proof}
若 \(M\cong \ZZ_2^r\)，则
\[
\#(M/2^nM)=2^{rn}.
\]
另一方面，由推论 \ref{cor:killo-2adic-tate-hypercube-limit}，
\[
E_{\min}[2^n]\cong (\ZZ/2^n\ZZ)^2,
\]
故
\[
\#E_{\min}[2^n]=2^{2n}.
\]
逐层比较得 \(2^{rn}=2^{2n}\) 对全部 \(n\) 成立，只能有 \(r=2\)。若 \(r=1\)，则 \(\#(M/2^nM)=2^n\) 与 \(\#E_{\min}[2^n]=4^n\) 矛盾，故单轴实现不可能。证毕。
\end{proof}

\begin{theorem}[规范 Gödel 编码的 \texorpdfstring{$B$}{B}-进 Cantor 吸引子与前缀轨道闭包]\label{thm:xi-time-part9gc-canonical-godel-cantor-attractor-orbit-closure}
沿用定理 \ref{thm:killo-godel-semidir-affine-2adic} 与定理
\ref{thm:xi-time-part9g-canonical-submonoid-free-monoid} 的记号。固定 primitive 地址向量
\[
v\in T\setminus 2T,
\qquad
B:=2^L>M,
\qquad
T_v:=\ZZ_2v\subset T.
\]
定义
\[
\Xi_{B,M}:\{0,\dots,M\}^{\NN}\longrightarrow T_v,
\qquad
\Xi_{B,M}(a_1,a_2,\dots):=\sum_{t\ge 1}a_tB^{t-1}v,
\]
以及其像
\[
K_{B,M}(v):=\Xi_{B,M}\bigl(\{0,\dots,M\}^{\NN}\bigr).
\]
则成立：
\begin{enumerate}
  \item \(\Xi_{B,M}\) 是到 \(K_{B,M}(v)\) 的拓扑同胚；等价地，每个 \(x\in K_{B,M}(v)\) 都具有唯一的 \(B\)-进 digit 展开；
  \item \(K_{B,M}(v)\) 满足严格不交并
  \[
  K_{B,M}(v)=\bigsqcup_{a=0}^{M}\bigl(av+B K_{B,M}(v)\bigr);
  \]
  \item 若 \(H_{\mathrm{can}}\) 表示规范编码子单子，则
  \[
  \overline{\{\Psi(r,k)(0):\ (r,k)\in H_{\mathrm{can}}\}}
  =
  K_{B,M}(v).
  \]
\end{enumerate}
\end{theorem}
\begin{proof}
取有限事件集
\[
E:=\{0,\dots,M\},
\qquad
\mathrm{code}:E\hookrightarrow \NN,\quad \mathrm{code}(a):=a.
\]
则由定理 \ref{thm:xi-time-part62d-hologram-image-cantor-homeomorphism}，地址映射 \(X:E^{\NN}\to \Lambda_{E,v}\) 在此正好就是 \(\Xi_{B,M}\)，且其像满足严格不交并
\[
\Lambda_{E,v}
=
\bigsqcup_{a=0}^{M}\bigl(av+B\Lambda_{E,v}\bigr).
\]
于是第 \(1\) 与第 \(2\) 条成立，并且 \(K_{B,M}(v)=\Lambda_{E,v}\)。

现证第 \(3\) 条。任取
\[
x=\Xi_{B,M}(a_1,a_2,\dots)\in K_{B,M}(v).
\]
对每个 \(k\ge 1\)，定义 \(r_k\in \mathcal R\) 为
\[
r_k(p_t):=
\begin{cases}
a_t,& 1\le t\le k,\\
0,& t>k.
\end{cases}
\]
则 \((r_k,k)\in H_{\mathrm{can}}\)，并且
\[
\Psi(r_k,k)(0)=\sum_{t=1}^{k}a_tB^{t-1}v.
\]
由于 \(|B|_2<1\)，上述有限前缀和在 \(k\to\infty\) 时收敛到 \(x\)。故
\[
K_{B,M}(v)\subseteq
\overline{\{\Psi(r,k)(0):\ (r,k)\in H_{\mathrm{can}}\}}.
\]
反向包含是显然的，因为每个 \(\Psi(r,k)(0)\) 都是某条无限 digit 序列在尾部补零后得到的地址点，故本身属于 \(K_{B,M}(v)\)。证毕。
\end{proof}

\begin{theorem}[规范 Gödel 吸引子的精确维数与 Haar 测度二歧律]\label{thm:xi-time-part9gc-canonical-godel-attractor-dimension-haar-dichotomy}
沿用定理 \ref{thm:xi-time-part9gc-canonical-godel-cantor-attractor-orbit-closure} 的记号。把 \(T_v\cong \ZZ_2\) 视作一维 \(2\)-进紧群，则
\[
\dim_{\mathrm H}K_{B,M}(v)
=
\dim_{\mathrm M}K_{B,M}(v)
=
\frac{\log(M+1)}{\log B}
=
\frac{\log_2(M+1)}{L}.
\]
此外：
\begin{enumerate}
  \item 若 \(M=B-1\)，则
  \[
  K_{B,M}(v)=T_v;
  \]
  \item 若 \(M\le B-2\)，则
  \[
  \mu_{T_v}\bigl(K_{B,M}(v)\bigr)=0,
  \qquad
  \operatorname{Int}_{T_v}\bigl(K_{B,M}(v)\bigr)=\varnothing,
  \]
  其中 \(\mu_{T_v}\) 为 \(T_v\) 的归一化 Haar 概率测度。
\end{enumerate}
\end{theorem}
\begin{proof}
由定理 \ref{thm:xi-time-part9gc-canonical-godel-cantor-attractor-orbit-closure}，集合 \(K_{B,M}(v)\) 与
\[
\Lambda_{E,v}
\qquad
\bigl(E=\{0,\dots,M\},\ \mathrm{code}=\mathrm{id}\bigr)
\]
完全一致。因而，结论 \ref{con:xi-time-part62b-rankone-hologram-selfsimilar-measure-dimension} 的第 \(3\) 条给出维数公式，而结论 \ref{con:xi-time-part62c-godel-tate-minimal-bitwidth-threshold} 的第 \(2\)、\(3\) 条给出满地址线与真 Cantor 子集的二歧律。于是
\[
\dim_{\mathrm H}K_{B,M}(v)
=
\dim_{\mathrm M}K_{B,M}(v)
=
\frac{\log|E|}{\log B}
=
\frac{\log(M+1)}{\log B}.
\]
若 \(M=B-1\)，则结论 \ref{con:xi-time-part62c-godel-tate-minimal-bitwidth-threshold} 的第 \(2\) 条给出 \(K_{B,M}(v)=T_v\)。若 \(M\le B-2\)，则其第 \(3\) 条给出 Haar 零测；再由结论 \ref{con:xi-time-part62b-rankone-hologram-selfsimilar-measure-dimension} 的第 \(2\) 条，内部必为空。证毕。
\end{proof}

\begin{theorem}[规范仿射字值的忠实阈值恰为 \texorpdfstring{$C>M$}{C>M}]\label{thm:xi-time-part9gc-canonical-affine-word-faithful-threshold}
\leanverified{paper\_xi\_time\_part9gc\_canonical\_affine\_word\_faithful\_threshold}
设 \(C=2^s\ge 2\)，并定义有限字集
\[
\mathcal W_M:=\bigsqcup_{k\ge 0}\{0,\dots,M\}^k.
\]
对每个字
\[
\mathbf a=(a_1,\dots,a_k)\in \{0,\dots,M\}^k
\]
定义仿射字值
\[
\Xi_C(\mathbf a)(x):=
C^k x+\sum_{t=1}^{k}a_tC^{t-1}v
\qquad (x\in T_v).
\]
则映射
\[
\Xi_C:\mathcal W_M\longrightarrow \operatorname{Aff}(T_v)
\]
忠实当且仅当
\[
C>M.
\]
\end{theorem}
\begin{proof}
先设 \(C>M\)。若
\[
\Xi_C(a_1,\dots,a_k)=\Xi_C(b_1,\dots,b_\ell)
\]
为同一仿射映射，则分别在 \(0\) 与 \(v\) 处取值并相减，得到
\[
C^k v=C^\ell v.
\]
由于 \(T_v\) 对 \(\ZZ_2\) 无挠且 \(v\neq 0\)，必有 \(C^k=C^\ell\)，从而 \(k=\ell\)。再由
\[
\sum_{t=1}^{k}(a_t-b_t)C^{t-1}v=0
\]
可知
\[
\sum_{t=1}^{k}(a_t-b_t)C^{t-1}=0
\qquad\text{于 }\ZZ_2.
\]
若存在首个分歧位置 \(j\) 使 \(a_j\neq b_j\)，则上式可写为
\[
C^{j-1}\bigl((a_j-b_j)+Cu\bigr)=0
\]
对某个 \(u\in \ZZ_2\) 成立。由于
\[
0<|a_j-b_j|\le M<C,
\]
有 \((a_j-b_j)\notin C\ZZ_2\)，故 \((a_j-b_j)+Cu\neq 0\)，矛盾。于是所有 digits 都相同，故 \(\Xi_C\) 忠实。

反之，若 \(C\le M\)，则长度 \(2\) 的两字
\[
(C,0),
\qquad
(0,1)
\]
都属于 \(\{0,\dots,M\}^2\)，并且对一切 \(x\in T_v\) 都有
\[
\Xi_C(C,0)(x)=C^2x+C v=\Xi_C(0,1)(x).
\]
故 \(\Xi_C\) 不忠实。证毕。
\end{proof}

\begin{theorem}[规范 Gödel 仿射层的精确字增长指数]\label{thm:xi-time-part9gc-canonical-affine-semigroup-word-growth}
\leanverified{paper\_xi\_time\_part9gc\_canonical\_affine\_semigroup\_word\_growth}
沿用前述记号。对每个 \(k\ge 0\)，记
\[
H_{\mathrm{can}}(k):=\{(r,k)\in H_{\mathrm{can}}\},
\qquad
\mathcal N_k:=\bigl|\{\Psi(r,k):\ (r,k)\in H_{\mathrm{can}}(k)\}\bigr|.
\]
则有精确公式
\[
\mathcal N_k=(M+1)^k.
\]
因而规范仿射层的字增长熵
\[
h_{\mathrm{grow}}
:=
\lim_{k\to\infty}\frac{1}{k}\log \mathcal N_k
\]
满足
\[
h_{\mathrm{grow}}=\log(M+1).
\]
特别地，结合定理 \ref{thm:xi-time-part9gc-canonical-godel-attractor-dimension-haar-dichotomy}，
\[
\dim_{\mathrm H}K_{B,M}(v)=\frac{h_{\mathrm{grow}}}{\log B}.
\]
\end{theorem}
\begin{proof}
由定理 \ref{thm:xi-time-part9g-canonical-submonoid-free-monoid}，每个 \(H_{\mathrm{can}}(k)\) 的元素都唯一对应一个长度 \(k\) 的 digit 字
\[
(a_1,\dots,a_k)\in \{0,\dots,M\}^k.
\]
而在 \(C=B\) 的情形下，定理 \ref{thm:xi-time-part9gc-canonical-affine-word-faithful-threshold} 表明该长度 \(k\) 的字值映射是忠实的。故不同 digit 字给出不同仿射映射，遂有
\[
\mathcal N_k=\bigl|\{0,\dots,M\}^k\bigr|=(M+1)^k.
\]
两边取对数并除以 \(k\)，即得
\[
h_{\mathrm{grow}}=\log(M+1).
\]
最后一式只是把定理 \ref{thm:xi-time-part9gc-canonical-godel-attractor-dimension-haar-dichotomy} 的维数公式重写为
\[
\frac{\log(M+1)}{\log B}
=
\frac{h_{\mathrm{grow}}}{\log B}.
\qedhere
\]
\end{proof}

\begin{theorem}[同一 Tate 环境中的 Haar--Cantor 分裂]\label{thm:xi-time-part9gc-haar-cantor-splitting}
记
\[
T:=T_2(E_{\min})\cong \ZZ_2^2,
\qquad
\mu_{\mathrm{res}}:=\mu_\infty,
\]
其中 \(\mu_\infty\) 为定理 \ref{thm:xi-time-part9gc-resolution-tower-haar-limit} 给出的 Haar 极限测度。对任意 \(M<B\)，设 \(\nu\) 为任意支撑在 \(K_{B,M}(v)\) 上的 Borel 概率测度。则
\[
\nu\perp \mu_{\mathrm{res}}.
\]
\end{theorem}
\begin{proof}
由定理 \ref{thm:xi-time-part9gc-canonical-godel-cantor-attractor-orbit-closure}，
\[
K_{B,M}(v)\subseteq T_v=\ZZ_2v.
\]
因此只需证明
\[
\mu_{\mathrm{res}}(T_v)=0.
\]
对每个 \(n\ge 1\)，记自然约化映射为
\[
\rho_n:T\longrightarrow T/2^nT\cong (\ZZ/2^n\ZZ)^2.
\]
由于 \(T_v\) 是由单一向量 \(v\) 生成的循环 \(\ZZ_2\)-子模，其像 \(\rho_n(T_v)\) 在有限群 \(T/2^nT\) 中仍是循环子群，因而
\[
\#\rho_n(T_v)\le 2^n.
\]
另一方面，由定理 \ref{thm:xi-time-part9gc-resolution-tower-haar-limit}，每个模 \(2^nT\) 陪集的 \(\mu_{\mathrm{res}}\)-质量都等于 \(4^{-n}\)。故
\[
\mu_{\mathrm{res}}(T_v)
\le
\#\rho_n(T_v)\cdot 4^{-n}
\le
2^n\cdot 4^{-n}
=
2^{-n}.
\]
令 \(n\to\infty\) 得 \(\mu_{\mathrm{res}}(T_v)=0\)。

现在取可测集 \(A:=T_v\)。则
\[
\nu(A)=1,
\qquad
\mu_{\mathrm{res}}(A)=0.
\]
按互相奇异的定义，遂得 \(\nu\perp \mu_{\mathrm{res}}\)。证毕。
\end{proof}

\begin{corollary}[语义 Cantor 子系统对 full-rank 采样空间的仿射不可恢复性]\label{cor:xi-time-part9gc-affine-nonrecoverability-fullrank-space}
沿用定理 \ref{thm:xi-time-part9gc-haar-cantor-splitting} 的记号。对任意连续仿射自映射
\[
A:T\longrightarrow T,
\qquad
A(x)=Lx+t,
\]
都有
\[
A\bigl(K_{B,M}(v)\bigr)\subseteq t+L(T_v).
\]
其中 \(L(T_v)\) 是秩至多为 \(1\) 的闭 \(\ZZ_2\)-子模。因而
\[
A\bigl(K_{B,M}(v)\bigr)
\]
在 \(T\) 中始终 Haar 零测且内部为空，特别不可能等于 \(T\)。
\end{corollary}
\begin{proof}
由定理 \ref{thm:xi-time-part9gc-canonical-godel-cantor-attractor-orbit-closure}，
\[
K_{B,M}(v)\subseteq T_v.
\]
故
\[
A\bigl(K_{B,M}(v)\bigr)\subseteq A(T_v)=t+L(T_v).
\]
而 \(T_v\) 是循环 \(\ZZ_2\)-模，其在线性映射 \(L\) 下的像 \(L(T_v)\) 仍是循环 \(\ZZ_2\)-模；又因 \(L\) 连续且 \(T_v\) 紧，\(L(T_v)\) 闭。故 \(L(T_v)\) 的秩至多为 \(1\)。

若 \(L(T_v)=\{0\}\)，则 \(A(K_{B,M}(v))\) 为单点，结论显然。若 \(L(T_v)\neq \{0\}\)，则与定理 \ref{thm:xi-time-part9gc-haar-cantor-splitting} 的证明完全相同，可得任意 rank-one 闭子模在 \(T\cong \ZZ_2^2\) 中都是 Haar 零测。于是其仿射平移 \(t+L(T_v)\) 亦为 Haar 零测。

在紧群 \(T\) 中，任一非空开集都具有正 Haar 测度；因此 Haar 零测集不可能含有内点。故
\[
A\bigl(K_{B,M}(v)\bigr)
\subseteq
t+L(T_v)
\]
也必内部为空，特别不可能等于整个 \(T\)。证毕。
\end{proof}

\begin{theorem}[binary Gödel 吸引子上的 \texorpdfstring{$\varphi$}{phi}-logistic 维数曲线]\label{thm:xi-time-part9gc-binary-godel-logistic-dimension-curve}
沿用定理 \ref{thm:xi-time-part9gc-canonical-godel-cantor-attractor-orbit-closure} 的记号，并取二元字母情形 \(M=1\)。记
\[
\Xi_B:=\Xi_{B,1}:\{0,1\}^{\NN}\longrightarrow K_{B,1}(v).
\]
对每个 \(q\ge 0\)，定义
\[
\theta(q):=\frac{1}{1+\varphi^{q+1}},
\]
并令 \(\beta_q\) 为 \(\{0,1\}^{\NN}\) 上的 Bernoulli 乘积测度，
\[
\beta_q(\omega_1=1)=\theta(q),
\qquad
\beta_q(\omega_1=0)=1-\theta(q).
\]
再记
\[
\nu_q:=(\Xi_B)_*\beta_q,
\qquad
h(t):=-t\log t-(1-t)\log(1-t).
\]
则 \(\nu_q\) 为 exact-dimensional 测度，并且
\[
\dim_{\mathrm H}\nu_q=\frac{h(\theta(q))}{\log B}.
\]
此外，对任意 \(p,q\ge 0\)，escort 族的极限 KL 散度满足
\[
\lim_{m\to\infty}D_{\mathrm{KL}}(\pi_{m,q}\Vert \pi_{m,p})
=
D_{\mathrm{KL}}\!\bigl(\mathrm{Bernoulli}(\theta(q))\Vert \mathrm{Bernoulli}(\theta(p))\bigr).
\]
\end{theorem}
\begin{proof}
在 \(E=\{0,1\}\)、\(\mathrm{code}=\mathrm{id}\) 的情形，定理
\ref{thm:xi-time-part9gc-canonical-godel-cantor-attractor-orbit-closure} 表明
\[
K_{B,1}(v)=\Xi_B(\{0,1\}^{\NN}),
\]
且 \(\Xi_B\) 为拓扑同胚。

对每个 \(q\)，把
\[
\mathbf p_q:=(1-\theta(q),\,\theta(q))
\]
视为 \(\{0,1\}\) 上的概率向量。则 \(\beta_q\) 正是定理 \ref{thm:xi-time-part62dc-hologram-bernoulli-exact-dimensional} 中的 Bernoulli 乘积测度 \(\nu_{\mathbf p_q}\)，而 \(\nu_q\) 正是其地址推前 \(\mu_{\mathbf p_q}\)。故同一定理立即给出 \(\nu_q\) 的 exact-dimensional 性以及
\[
\dim_{\mathrm H}\nu_q
=
\frac{H(\mathbf p_q)}{\log B}.
\]
由于
\[
H(\mathbf p_q)
=
-\theta(q)\log\theta(q)
-(1-\theta(q))\log(1-\theta(q))
=
h(\theta(q)),
\]
遂得
\[
\dim_{\mathrm H}\nu_q=\frac{h(\theta(q))}{\log B}.
\]

另一方面，定理 \ref{thm:xi-time-part9s-escort-logistic-expfamily-bregman} 已把
\[
q\longmapsto \theta(q)=\frac{1}{1+\varphi^{q+1}}
\]
识别为 escort 极限族的 logistic 指数坐标，而定理 \ref{thm:xi-time-part9od-escort-escort-fdiv-binary-closure} 在 \(f(u)=u\log u\) 时恰给出
\[
\lim_{m\to\infty}D_{\mathrm{KL}}(\pi_{m,q}\Vert \pi_{m,p})
=
\bigl(1-\theta(q)\bigr)\log\frac{1-\theta(q)}{1-\theta(p)}
+\theta(q)\log\frac{\theta(q)}{\theta(p)}.
\]
右端正是
\[
D_{\mathrm{KL}}\!\bigl(\mathrm{Bernoulli}(\theta(q))\Vert \mathrm{Bernoulli}(\theta(p))\bigr).
\]
故第二条成立。证毕。
\end{proof}

\noindent
定理 \ref{thm:xi-time-part9gc-resolution-tower-haar-limit} 至定理 \ref{thm:xi-time-part9gc-binary-godel-logistic-dimension-curve} 共同表明，H.159 的分辨率塔极限并不与 H.160 的规范语义地址像混同：前者在 \(T_2(E_{\min})\cong \ZZ_2^2\) 上携带 full-rank 的 Haar 几何，后者则始终困在由单一地址向量生成的 rank-one \(B\)-进 Cantor 子系统内。于是，同一 Tate 环境中的分辨率采样协议与 Gödel 语义协议不是简单的坐标改写关系，而是两类在秩、测度与仿射可恢复性上同时分裂的读出机制。

\endinput


\begin{theorem}[Lee--Yang 三次核的判别式分解与 \texorpdfstring{\(S_3\)}{S3} 分裂域]\label{thm:xi-time-part9g-leyang-cubic-discriminant}
定义
\[
P_{\mathrm{LY}}(y):=256y^3+411y^2+165y+32\in\ZZ[y].
\]
则
\[
\Disc(P_{\mathrm{LY}})
=-3^9\cdot 31^2\cdot 37
=-(3\cdot 37)\,(3^4\cdot 31)^2.
\]
从而 \(P_{\mathrm{LY}}\) 在 \(\QQ\) 上可分离且不可约，分裂域的 Galois 群为 \(S_3\)，其唯一二次子域为
\[
\QQ\bigl(\sqrt{-111}\bigr).
\]
\end{theorem}
\leanverified{paper\_xi\_time\_part9g\_leyang\_cubic\_discriminant}
\begin{proof}
对三次多项式 \(ay^3+by^2+cy+d\) 用判别式公式
\[
\Delta=b^2c^2-4ac^3-4b^3d-27a^2d^2+18abcd
\]
代入 \(a=256,b=411,c=165,d=32\) 得
\[
\Delta=-699868431=-3^9\cdot 31^2\cdot 37.
\]
故可分离。
再取模 \(5\)：
\[
P_{\mathrm{LY}}(y)\equiv y^3+y^2+2\pmod 5,
\]
在 \(\FF_5\) 上无根，故在 \(\QQ\) 上不可约。
不可约三次多项式的 Galois 群仅可能为 \(A_3\) 或 \(S_3\)；判别式非平方，故为 \(S_3\)。
其二次子域即 \(\QQ(\sqrt{\Delta})=\QQ(\sqrt{-111})\)。
\end{proof}

\begin{theorem}[Langlands--Euler 乘积的停机分离率与素数截断复杂度]\label{thm:xi-time-part9g-langlands-euler-separation}
取 \(\sigma>1\)。
设非停机时 \(g_{M,x}=I_2\)；停机深度为 \(N\) 时
\[
g_{M,x}=I_N:=\mathrm{diag}(e^{i\theta_N},e^{-i\theta_N}),
\qquad
\theta_N:=\frac{\pi}{2^N}.
\]
定义局部与全局因子
\[
L_p(\sigma,g):=\det(I-p^{-\sigma}g)^{-1},
\qquad
L(\sigma;M,x):=\prod_p L_p(\sigma,g_{M,x}),
\]
并记
\[
\zeta(\sigma)^2=\prod_p(1-p^{-\sigma})^{-2}.
\]
则存在仅依赖 \(\sigma\) 的常数 \(0<c_\sigma\le C_\sigma<\infty\)，使停机深度为 \(N\) 时
\[
c_\sigma\,4^{-N}
\le
\zeta(\sigma)^2-L(\sigma;M,x)
\le
C_\sigma\,4^{-N}.
\]
此外，若用素数截断
\[
L_{\le P}(\sigma;M,x):=\prod_{p\le P}L_p(\sigma,g_{M,x}),\qquad
\zeta_{\le P}(\sigma)^2:=\prod_{p\le P}(1-p^{-\sigma})^{-2},
\]
且某判别方案总误差为 \(o(4^{-N})\)，则必有
\[
P\ge K_\sigma\,4^{\,N/(\sigma-1)}
\]
对某 \(K_\sigma>0\) 成立。
\end{theorem}
\begin{proof}
停机情形记 \(q:=p^{-\sigma}\in(0,1)\)，则
\[
L_p(\sigma,I_N)=\frac{1}{1-2q\cos\theta_N+q^2},
\qquad
(1-q)^{-2}=\frac{1}{1-2q+q^2}.
\]
定义
\[
R_p(N):=\frac{L_p(\sigma,I_N)}{(1-q)^{-2}}
=\frac{(1-q)^2}{1-2q\cos\theta_N+q^2}\in(0,1),
\]
并令 \(R(N):=\prod_p R_p(N)\)。
则
\[
\zeta(\sigma)^2-L(\sigma;M,x)=\zeta(\sigma)^2\,(1-R(N)).
\]
又
\[
1-R_p(N)
=\frac{2q(1-\cos\theta_N)}{1-2q\cos\theta_N+q^2}.
\]
利用
\[
(1-q)^2\le 1-2q\cos\theta_N+q^2\le (1+q)^2
\]
得到
\[
\frac{2q(1-\cos\theta_N)}{(1+q)^2}
\le
1-R_p(N)
\le
\frac{2q(1-\cos\theta_N)}{(1-q)^2}.
\]
下界侧，由 \(R(N)\le R_2(N)\) 得
\[
1-R(N)\ge 1-R_2(N)\ge c'_\sigma(1-\cos\theta_N).
\]
上界侧，用 \(1-\prod_p r_p\le \sum_p(1-r_p)\)（\(0<r_p\le1\)）得
\[
1-R(N)\le \sum_p(1-R_p(N))
\le C'_\sigma(1-\cos\theta_N)\sum_p\frac{p^{-\sigma}}{(1-p^{-\sigma})^2}
\le C''_\sigma(1-\cos\theta_N).
\]
并且 \(\theta_N=\pi/2^N\) 满足 \(1-\cos\theta_N\asymp \theta_N^2\asymp 4^{-N}\)，故第一组双边界成立。

再看截断复杂度。对 \(\sigma>1\)，有
\[
\log\frac{\zeta(\sigma)^2}{\zeta_{\le P}(\sigma)^2}
=\sum_{p>P}O(p^{-\sigma})
=O(P^{1-\sigma}),
\]
同理
\[
\log\frac{L(\sigma;M,x)}{L_{\le P}(\sigma;M,x)}
=O(P^{1-\sigma}),
\]
常数仅依赖 \(\sigma\)。
故若总误差需为 \(o(4^{-N})\)，必要条件是
\(
P^{1-\sigma}=o(4^{-N})
\)，等价于
\(
P\ge K_\sigma 4^{N/(\sigma-1)}
\)（常数重标定后）。
\end{proof}

\paragraph{\texorpdfstring{\(\chi\)}{chi} 扇区极值冻结与谱税漂移耦合律}
本段把极大纤维的 \(\mathbb Z_2^2\)-扇区结构与 near-\(1\) 极点预算律并置到同一审计链条中。

\begin{theorem}[极大纤维的稀有扇区承载律]\label{thm:xi-time-part9h-maxfiber-rare-sector-carrying}
\leanverified{paper\_xi\_time\_part9h\_maxfiber\_rare\_sector\_carrying}
令
\[
D_m:=\max_{x\in X_m}d_m(x),\qquad
\mathcal M_m:=\{x\in X_m:\ d_m(x)=D_m\},
\]
并定义
\[
\chi(x):=\bigl(\operatorname{sgn}(c_x),\operatorname{sgn}(\iota_x)\bigr)\in\{\pm1\}^2.
\]
记 \(\chi_m\) 为 \(\mathcal M_m\) 上的常值（见定理 \ref{thm:pom-maxfiber-orientation-charge-periodic}），并令
\(
\mathsf U_m:=\mathrm{Unif}(X_m)
\)。
则：
\begin{enumerate}
  \item \(\chi_m\) 良定义，且 \(\mathcal M_m\subseteq\{x:\chi(x)=\chi_m\}\)；
  \item 若 \(\chi_m\neq(+,+)\)，则存在常数 \(\kappa>0\) 使
  \[
  \mathsf U_m(\mathcal M_m)
  \le
  \mathsf U_m\!\bigl(\chi=\chi_m\bigr)
  \le
  2e^{-\kappa m}
  \]
  对充分大 \(m\) 成立；
  \item 以下同余类上必有 \(\chi_m\neq(+,+)\)，从而 \(\mathsf U_m(\mathcal M_m)\) 指数衰减：
  \[
  m\ \text{奇且}\ m\equiv 1,3\ (\mathrm{mod}\ 6),
  \]
  \[
  m=2k\ \text{且}\ k\not\equiv 0,5,10,11\ (\mathrm{mod}\ 12).
  \]
\end{enumerate}
\end{theorem}
\begin{proof}
第一条与第三条直接由定理 \ref{thm:pom-maxfiber-orientation-charge-periodic} 的常值性与周期分类得到。
若 \(\chi_m\neq(+,+)\)，则至少有一个坐标为 \(-1\)，故
\[
\{\chi=\chi_m\}
\subseteq
\{\operatorname{sgn}(c_x)=-1\}\cup\{\operatorname{sgn}(\iota_x)=-1\}.
\]
由定理 \ref{thm:pom-toggle-odd-sector-exponential-thinness}，
\[
\mathsf U_m\!\bigl(\operatorname{sgn}(c_x)=-1\bigr)\le e^{-\kappa m},\qquad
\mathsf U_m\!\bigl(\operatorname{sgn}(\iota_x)=-1\bigr)\le e^{-\kappa m}
\]
（调整 \(\kappa\) 取两者较小值）。并联合并即得
\(
\mathsf U_m(\chi=\chi_m)\le 2e^{-\kappa m}
\)。
再用 \(\mathcal M_m\subseteq\{\chi=\chi_m\}\) 得到对 \(\mathsf U_m(\mathcal M_m)\) 的同界。
\end{proof}

\begin{theorem}[\texorpdfstring{\(\chi\)}{chi} 扇区冻结与 escort 极限的两比特坍缩]\label{thm:xi-time-part9h-escort-sector-freezing}
\leanverified{paper\_xi\_time\_part9h\_escort\_sector\_freezing}
对 \(q>0\) 定义 escort 分布
\[
w_{m,q}(x):=\frac{d_m(x)^q}{S_q(m)},\qquad
S_q(m):=\sum_{y\in X_m}d_m(y)^q.
\]
令
\[
D_m:=\max_{x\in X_m}d_m(x),\qquad
\mathcal M_m:=\{x:\ d_m(x)=D_m\},
\]
并设
\[
D_m^{(2)}:=
\max_{x\in X_m\setminus\mathcal M_m}d_m(x)
\]
（若补集为空则取 \(D_m^{(2)}=0\)）。
则对任意 \(m,q\) 有显式界
\[
w_{m,q}\!\bigl(\chi\neq\chi_m\bigr)
\le
w_{m,q}(X_m\setminus\mathcal M_m)
\le
|X_m|\left(\frac{D_m^{(2)}}{D_m}\right)^q.
\]
进一步，若满足
\[
\liminf_{m\to\infty}\frac1m\log\frac{D_m}{D_m^{(2)}}>0,\qquad
|X_m|=\exp(O(m)),
\]
则对任意序列 \(q_m\to\infty\) 有
\[
w_{m,q_m}\!\bigl(\chi=\chi_m\bigr)\longrightarrow 1\qquad(m\to\infty).
\]
\end{theorem}
\begin{proof}
由定理 \ref{thm:pom-maxfiber-orientation-charge-periodic}，\(\mathcal M_m\subseteq\{\chi=\chi_m\}\)，故
\[
w_{m,q}(\chi\neq\chi_m)\le w_{m,q}(X_m\setminus\mathcal M_m).
\]
并且
\[
w_{m,q}(X_m\setminus\mathcal M_m)
=
\frac{\sum_{x\notin\mathcal M_m}d_m(x)^q}{\sum_{y}d_m(y)^q}
\le
\frac{|X_m|(D_m^{(2)})^q}{D_m^q},
\]
得显式界。

在额外假设下，存在 \(\delta>0\) 与 \(m_0\) 使 \(D_m^{(2)}/D_m\le e^{-\delta m}\)（\(m\ge m_0\)）。
故
\[
|X_m|\left(\frac{D_m^{(2)}}{D_m}\right)^{q_m}
\le
\exp\!\bigl(Cm-\delta m q_m\bigr)\xrightarrow[m\to\infty]{}0
\]
（因 \(q_m\to\infty\)）。
因此 \(w_{m,q_m}(X_m\setminus\mathcal M_m)\to0\)，等价于
\(
w_{m,q_m}(\chi=\chi_m)\to1
\)。
\end{proof}

\begin{theorem}[\texorpdfstring{\(\mathbb Z_2^2\)}{Z2^2} 扇区矩的 Hadamard 可逆层析]\label{thm:xi-time-part9h-z2x2-hadamard-tomography}
\leanverified{paper\_xi\_time\_part9h\_z2x2\_hadamard\_tomography}
固定 \(m\) 与 \(q\ge0\)。
对 \(s,t\in\{0,1\}\) 定义
\[
\mathbf M_{s,t}^{(q)}(m)
:=
2^m\,\tau_{m,1}\!\left(
\varepsilon_{\mathrm{scan}}^{\,s}\varepsilon_{\mathrm{rev}}^{\,t}\,
\mathrm{IndW}(E_{m,1})^q
\right),
\]
并对 \((\alpha,\beta)\in\{\pm1\}^2\) 定义扇区矩
\[
\mathbf S_{\alpha,\beta}^{(q+1)}(m)
:=
\sum_{\substack{x\in X_m\\ \chi(x)=(\alpha,\beta)}} d_m(x)^{q+1}.
\]
则有双向显式变换
\[
\mathbf M_{s,t}^{(q)}(m)
=
\sum_{\alpha,\beta\in\{\pm1\}}
\alpha^s\beta^t\,
\mathbf S_{\alpha,\beta}^{(q+1)}(m),
\]
\[
\mathbf S_{\alpha,\beta}^{(q+1)}(m)
=
\frac14\sum_{s,t\in\{0,1\}}
\alpha^s\beta^t\,
\mathbf M_{s,t}^{(q)}(m).
\]
\end{theorem}
\begin{proof}
由定理 \ref{thm:fold-groupoid-z2x2-joint-spectral-measure}，
\[
\mathbf M_{s,t}^{(q)}(m)
=
\sum_{x\in X_m}
\operatorname{sgn}(c_x)^s\operatorname{sgn}(\iota_x)^t
d_m(x)^{q+1}.
\]
按 \(\chi(x)=(\alpha,\beta)\) 分组即得第一式。
第二式由 \(\{\pm1\}^2\) 角色正交性给出的 Fourier--Hadamard 反演直接得到。
\end{proof}

\begin{theorem}[near-\texorpdfstring{$1$}{1} 极点的漂移--混合不确定性关系]\label{thm:xi-time-part9h-drift-mix-uncertainty}
\leanverified{paper\_xi\_time\_part9h\_drift\_mix\_uncertainty}
在谱税框架（注 \ref{rem:pom-spectral-tax-workload}）下，记
\[
r_m:=\frac{\Lambda_m}{\lambda_m},\qquad
\Delta_m:=-\log r_m,\qquad
\tau_{\mathrm{mix}}(m):=\frac{1}{\Delta_m},
\]
\[
W(m):=\frac{\log|G_m|}{\Delta_m}
=\log|G_m|\,\tau_{\mathrm{mix}}(m).
\]
取最坏扭曲方向 \(\vartheta_\star\in\Lambda(m)\setminus\{0\}\)，并记
\[
\Delta(\vartheta):=-\log\!\left(\frac{\rho(M_{i\vartheta})}{\lambda}\right),\qquad
s(\vartheta):=\frac{\log\rho(M_{i\vartheta})}{\log\lambda}.
\]
在小角域内（推论 \ref{cor:pom-multi-axis-near1-pole-barrier}）有
\[
1-\Re s(\vartheta)=\frac{\Delta(\vartheta)}{\log\lambda}.
\]
再令
\[
H:=\nabla^2\log\mathfrak M(0),\qquad
\mathcal W(v):=\frac{v^{\mathsf T}Hv}{v^{\mathsf T}\Sigma v},\qquad
v(\vartheta):=\frac{\vartheta}{\|\vartheta\|},
\]
\[
\Delta_{\mathfrak M}(\vartheta):=
\log\mathfrak M(0)-\Re\log\mathfrak M(i\vartheta).
\]
则（推论 \ref{cor:pom-multi-axis-pole-drift-binding}）
\[
\Delta_{\mathfrak M}(\vartheta)
=(\log\lambda)\,\mathcal W(v(\vartheta))\,(1-\Re s(\vartheta))
+O(\|\vartheta\|^4).
\]
特别地，在 \(\vartheta=\vartheta_\star\) 且 \(\Delta_m=\Delta(\vartheta_\star)\) 时，
\[
\tau_{\mathrm{mix}}(m)\,\Delta_{\mathfrak M}(\vartheta_\star)
=
\mathcal W\!\bigl(v(\vartheta_\star)\bigr)
+O\!\bigl(q(m)\bigr),
\]
其中
\[
q(m):=\min_{\vartheta\in\Lambda(m)\setminus\{0\}}\vartheta^{\mathsf T}\Sigma\vartheta.
\]
并且由 Minkowski 障碍（结论 \ref{con:xi-terminal-minkowski-budget-barrier} 与推论 \ref{cor:pom-multi-axis-near1-pole-barrier}）有 \(q(m)=O(B^{-2/d})\)。
\end{theorem}
\begin{proof}
由
\(
1-\Re s(\vartheta)=\Delta(\vartheta)/\log\lambda
\)
代入漂移换算式，得
\[
\Delta_{\mathfrak M}(\vartheta)
=
\mathcal W(v(\vartheta))\,\Delta(\vartheta)+O(\|\vartheta\|^4).
\]
取 \(\vartheta=\vartheta_\star\)：
\[
\tau_{\mathrm{mix}}(m)\Delta_{\mathfrak M}(\vartheta_\star)
=
\frac{\Delta_{\mathfrak M}(\vartheta_\star)}{\Delta_m}
=
\mathcal W(v(\vartheta_\star))
+O\!\left(\frac{\|\vartheta_\star\|^4}{\Delta_m}\right).
\]
又在小角二次近似下
\(
\Delta_m=\tfrac12 q(m)+O(q(m)^2)
\)
且 \(\|\vartheta_\star\|^2\asymp q(m)\)，故余项为 \(O(q(m))\)。
最后由 \(q(m)=O(B^{-2/d})\) 得尺度结论。
\end{proof}

\begin{corollary}[谱税--漂移乘积律]\label{cor:xi-time-part9h-spectral-tax-drift-product}
\leanverified{paper\_xi\_time\_part9h\_spectral\_tax\_drift\_product}
在定理 \ref{thm:xi-time-part9h-drift-mix-uncertainty} 设定下，
\[
W(m)\,\Delta_{\mathfrak M}(\vartheta_\star)
=
\log|G_m|\cdot \mathcal W\!\bigl(v(\vartheta_\star)\bigr)
+O\!\bigl(\log|G_m|\,q(m)\bigr).
\]
因而在 \(q(m)\to0\) 的高预算极限中，乘积主导阶由 \(\log|G_m|\) 线性控制；若要求 \(|\Delta_{\mathfrak M}(\vartheta_\star)|=o(1)\)，则必须伴随 \(\tau_{\mathrm{mix}}(m)\to\infty\)，从而 \(W(m)\to\infty\)。
\end{corollary}
\begin{proof}
将定理 \ref{thm:xi-time-part9h-drift-mix-uncertainty} 的关系两边乘以 \(\log|G_m|\)，并代入
\(
W(m)=\log|G_m|\tau_{\mathrm{mix}}(m)
\)
即得第一式。
第二句由该式主项与误差阶直接推出。
\end{proof}

\paragraph{取向 torsor 的覆盖同调、传递律与逆极限不可有限化障碍}
沿用“二值跳跃由取向函子 \(\mathrm{Or}\) 统一承载”的语义轴（定理 \ref{thm:xi-binary-jump-orientation-unification}），并使用符号单值化在 \(\ZZ_2\) 上的同调表示。

\begin{theorem}[二值跳跃类的上同调刻画与可消去性判别]\label{thm:xi-time-part9i-orientation-local-system-cohomology}
\leanverified{paper\_xi\_time\_part9i\_orientation\_local\_system\_cohomology}
设 \(Y\) 连通、局部道路连通且半局部单连通，\(\pi:\widetilde Y\to Y\) 为次数 \(d\ge 2\) 的连通有限覆盖。
对每个 \(y\in Y\) 记纤维 \(S_y:=\pi^{-1}(y)\)。
则存在与 \(\pi\) 典范关联的 \(\ZZ_2\)-局部系统 \(\mathcal L_\pi\to Y\)，满足：
\begin{enumerate}
  \item 对每个 \(y\in Y\)，\((\mathcal L_\pi)_y\) 与 \(\mathrm{Or}(S_y)\) 在 \(\ZZ_2\)-torsor 意义下自然同构；
  \item \(\mathcal L_\pi\) 可平凡化当且仅当 \(\pi\) 的 monodromy 结构群可约化到 \(\mathfrak A_d\)；
  \item 若 \(\rho:\pi_1(Y,y_0)\to\mathfrak S_d\) 为 \(\pi\) 的 monodromy 表示，则
  \[
  [\mathcal L_\pi]=[\operatorname{sgn}\circ\rho]
  =: \omega(\pi)\in H^1(Y;\ZZ_2),
  \]
  且该类完全刻画 \(\mathcal L_\pi\) 的同构类。
\end{enumerate}
\end{theorem}
\begin{proof}
由命题 \ref{prop:bdry-orientation-torsor-determinant-realization-jumpclass}，
\(
y\mapsto \mathrm{Or}(S_y)
\)
组装成主 \(\ZZ_2\)-丛，且其同调类为
\(
j(\pi)=w_1(\det(\pi_\ast\RR))
\)。
该类对应于 \(\operatorname{sgn}\circ\rho\) 的表示，故可取
\(
\mathcal L_\pi:=\mathrm{Or}(S_\bullet)
\)
并记 \(\omega(\pi):=j(\pi)\)。
纤维辨识与自然性即得第 (1) 条。

第 (2) 与第 (3) 的“平凡化 \(\Leftrightarrow\) \(\mathfrak A_d\)-约化”及“由 \(\operatorname{sgn}\circ\rho\) 完全分类”正是推论
\ref{cor:bdry-orientation-local-system-w1-alternating-reduction}
的内容。
\end{proof}

\begin{corollary}[符号局部系统的完备性（无外加寄存器的二值跳跃）]\label{cor:xi-time-part9i-sign-local-system-completeness}
沿用定理 \ref{thm:xi-time-part9i-orientation-local-system-cohomology} 的设定，并进一步设 \(d\ge 3\)。
令 \(\mathcal J\to Y\) 为任意 \(\ZZ_2\)-局部系统，且其单值化同态 \(\chi_{\mathcal J}:\pi_1(Y,y_0)\to\ZZ_2\) 仅通过覆叠单值化置换表示 \(\rho:\pi_1(Y,y_0)\to\mathfrak S_d\) 因子化：
存在群同态 \(\chi:\mathfrak S_d\to\ZZ_2\) 使 \(\chi_{\mathcal J}=\chi\circ\rho\)。
则在同构意义下仅有两种可能：
\[
\chi\equiv 0\ \Longrightarrow\ \mathcal J\ \text{平凡},\qquad
\chi=\operatorname{sgn}\ \Longrightarrow\ \mathcal J\simeq \mathcal L_\pi.
\]
等价地，同调类 \([\mathcal J]\in H^1(Y;\ZZ_2)\) 必为 \(0\) 或 \(\omega(\pi)\)。
\end{corollary}
\begin{proof}
当 \(d\ge 3\) 时 \(\mathfrak S_d^{\mathrm{ab}}\cong \ZZ_2\)，故 \(\mathrm{Hom}(\mathfrak S_d,\ZZ_2)\) 仅含平凡同态与 \(\operatorname{sgn}\) 两者。
若 \(\chi\equiv 0\) 则 \(\chi_{\mathcal J}\equiv 0\)，从而 \(\mathcal J\) 平凡。
若 \(\chi=\operatorname{sgn}\) 则 \(\chi_{\mathcal J}=\operatorname{sgn}\circ\rho\)，而定理 \ref{thm:xi-time-part9i-orientation-local-system-cohomology} 已给出 \(\mathcal L_\pi\) 的同构类正由该同态唯一决定，故 \(\mathcal J\simeq \mathcal L_\pi\)。证毕。
\end{proof}

\begin{theorem}[最小定向化二重基变与普适因子化]\label{thm:xi-time-part9i-minimal-orientation-double-cover}
沿用定理 \ref{thm:xi-time-part9i-orientation-local-system-cohomology}，记
\[
\chi_\pi:=\operatorname{sgn}\circ\rho:\pi_1(Y,y_0)\to\ZZ_2.
\]
令 \(p:Y^{\mathrm{or}}\to Y\) 为对应于子群 \(\ker(\chi_\pi)\le \pi_1(Y,y_0)\) 的（可能平凡的）覆盖。
则：
\begin{enumerate}
  \item 拉回覆盖 \(p^\ast(\pi):\widetilde Y\times_Y Y^{\mathrm{or}}\to Y^{\mathrm{or}}\) 的 monodromy 全落在 \(\mathfrak A_d\)，从而 \(p^\ast(\mathcal L_\pi)\) 平凡；
  \item 对任意连续映射 \(f:Z\to Y\)，若 \(f^\ast(\pi)\) 的 monodromy 落在 \(\mathfrak A_d\)，则 \(f\) 唯一经由 \(p\) 因子化：存在唯一 \(\widetilde f:Z\to Y^{\mathrm{or}}\) 使 \(p\circ\widetilde f=f\)。
\end{enumerate}
\end{theorem}
\begin{proof}
由构造，\(\chi_\pi\) 在 \(\pi_1(Y^{\mathrm{or}})\) 上为零，故 \(p^\ast(\pi)\) 的符号单值化平凡，等价于 monodromy 落在 \(\ker(\operatorname{sgn})=\mathfrak A_d\)；于是 \(p^\ast(\mathcal L_\pi)\) 平凡（推论 \ref{cor:bdry-orientation-local-system-w1-alternating-reduction}）。

对第 (2) 条，条件“\(f^\ast(\pi)\) 的 monodromy 落在 \(\mathfrak A_d\)”等价于
\(
\chi_\pi\circ f_\ast=0
\)，即
\(
f_\ast(\pi_1(Z))\subseteq \ker(\chi_\pi)
\)。
覆盖分类定理遂给出唯一提升 \(\widetilde f\) 到 \(Y^{\mathrm{or}}\)。
\end{proof}

\begin{theorem}[覆盖塔二值跳跃传递律：楔积符号与 transfer 加法公式]\label{thm:xi-time-part9i-cover-tower-transfer-law}
\leanverified{paper\_xi\_time\_part9i\_cover\_tower\_transfer\_law}
设
\[
X\xleftarrow{\ \pi\ }Y\xleftarrow{\ \rho\ }Z
\]
为连通有限覆盖，\(\deg(\pi)=n\)、\(\deg(\rho)=m\)。
记
\[
\omega(\pi)\in H^1(X;\ZZ_2),\qquad
\omega(\rho)\in H^1(Y;\ZZ_2)
\]
为定理 \ref{thm:xi-time-part9i-orientation-local-system-cohomology} 的二值跳跃类。
定义 transfer 映射
\[
\pi_!:H^1(Y;\ZZ_2)\to H^1(X;\ZZ_2)
\]
如下：对闭路类 \([\gamma]\in\pi_1(X,x_0)\)，令 \(\{\widetilde\gamma_y\}_{y\in\pi^{-1}(x_0)}\) 为从各纤维点起始的全部提升闭路，则
\[
(\pi_!\alpha)([\gamma])
:=
\sum_{y\in\pi^{-1}(x_0)}\alpha([\widetilde\gamma_y])\pmod 2.
\]
则有恒等式
\[
\omega(\rho\circ\pi)
=
(m\bmod 2)\,\omega(\pi)+\pi_!(\omega(\rho))
\qquad\text{于 }H^1(X;\ZZ_2).
\]
\end{theorem}
\begin{proof}
固定 \(\gamma\in\pi_1(X,x_0)\)。
设 \(\tau_\gamma\in\mathfrak S_n\) 为 \(\pi\) 沿 \(\gamma\) 的 monodromy；
对每个 \(y_i\in\pi^{-1}(x_0)\)，设 \(\sigma_{\gamma,i}\in\mathfrak S_m\) 为 \(\rho\) 沿提升闭路 \(\widetilde\gamma_{y_i}\) 的 monodromy。
则复合覆盖的 monodromy 位于 \(\mathfrak S_m\wr\mathfrak S_n\)，其符号满足
\[
\operatorname{sgn}_{mn}\bigl((\sigma_{\gamma,i})_{i=1}^n,\tau_\gamma\bigr)
=
\operatorname{sgn}_n(\tau_\gamma)^m
\prod_{i=1}^n\operatorname{sgn}_m(\sigma_{\gamma,i}).
\]
转为 \(\ZZ_2\) 加法即
\[
\omega(\rho\circ\pi)([\gamma])
=
m\,\omega(\pi)([\gamma])
\;+\;
\sum_{i=1}^n\omega(\rho)([\widetilde\gamma_{y_i}])
\pmod 2.
\]
右端第二项正是 \((\pi_!\omega(\rho))([\gamma])\)，故结论成立。
\end{proof}

\begin{corollary}[偶数复制的二值跳跃湮灭律]\label{cor:xi-time-part9i-even-stacking-annihilation}
\leanverified{paper\_xi\_time\_part9i\_even\_stacking\_annihilation}
在定理 \ref{thm:xi-time-part9i-orientation-local-system-cohomology} 设定下，
对任意 \(r\ge 1\) 定义 \(r\)-份平凡复制覆盖
\[
\pi^{(r)}:\widetilde Y\times [r]\to Y.
\]
则
\[
\omega(\pi^{(r)})=(r\bmod 2)\,\omega(\pi)\in H^1(Y;\ZZ_2).
\]
特别地：
\begin{enumerate}
  \item 若 \(r\) 为偶数，则 \(\omega(\pi^{(r)})=0\)；
  \item 若 \(r\) 为奇数，则 \(\omega(\pi^{(r)})=\omega(\pi)\)。
\end{enumerate}
\end{corollary}
\begin{proof}
令 \(\epsilon_r:Y\times[r]\to Y\) 为平凡 \(r\)-重覆盖，则
\(
\widetilde Y\times_Y(Y\times[r])\cong \widetilde Y\times[r]
\)
且 \(j(\epsilon_r)=0\)。
将命题 \ref{prop:bdry-orientation-jumpclass-fiberproduct-multiplicativity} 应用于 \((\pi,\epsilon_r)\) 即得
\[
\omega(\pi^{(r)})
=(r\bmod 2)\omega(\pi)+(\deg\pi\bmod2)\omega(\epsilon_r)
=(r\bmod 2)\omega(\pi).
\]
\end{proof}

\begin{theorem}[幺半兼容二值跳跃理论的刚性分类]\label{thm:xi-time-part9i-monoidal-rigidity-by-two-points}
令 \(\mathbf{FinSet}^{\cong}\) 为有限集合群胚（张量为 \(\sqcup\)），
\(\mathbf{Tors}(\ZZ_2)\) 为 \(\ZZ_2\)-torsor 的 Picard 群胚（张量为 \(\otimes\)）。
考虑对称幺半函子
\[
J:(\mathbf{FinSet}^{\cong},\sqcup)\to(\mathbf{Tors}(\ZZ_2),\otimes),
\qquad J(\varnothing)\ \text{平凡}.
\]
则在对称幺半自然同构意义下，仅有两类：
\begin{enumerate}
  \item 处处平凡类；
  \item 与取向函子 \(\mathrm{Or}\) 等价的非平凡类。
\end{enumerate}
并且该分类完全由 \(J\) 在二点集 \([2]\) 上诱导同态
\(
\chi_2:S_2\to\ZZ_2
\)
决定：
\[
\chi_2\equiv 0\ \Longrightarrow\ J\ \text{平凡},\qquad
\chi_2\not\equiv 0\ \Longrightarrow\ J\simeq \mathrm{Or}.
\]
\end{theorem}
\begin{proof}
记 \(\chi_n:S_n\to\Aut(J([n]))\cong\ZZ_2\) 为 \(J\) 的自同构表示。
幺半自然性给出
\[
\chi_{m+n}(\sigma\oplus\tau)=\chi_m(\sigma)+\chi_n(\tau)\pmod2.
\]
取 \((12)\in S_2\) 与恒等元得
\[
\chi_n\bigl((12)\oplus\mathrm{id}_{n-2}\bigr)=\chi_2((12)).
\]
又 \(S_n\) 中任意换位与 \((12)\oplus\mathrm{id}_{n-2}\) 共轭，且目标 \(\ZZ_2\) 交换，故所有换位像相同。
于是：

若 \(\chi_2((12))=0\)，则换位全映到 \(0\)，因换位生成 \(S_n\)，故 \(\chi_n\equiv0\)（\(n\ge2\)），从而 \(J\) 为平凡类。

若 \(\chi_2((12))=1\)，则 \(S_n\) 全体换位均映到 \(1\)，故 \(\chi_n=\operatorname{sgn}\)（\(n\ge2\)）。
由定理 \ref{thm:bdry-binary-jump-orientation-functor-uniqueness} 的唯一性，
每个固定阶层上的非平凡类唯一且由 \(\mathrm{Or}\) 实现；再结合命题
\ref{prop:bdry-orientation-functor-symmetric-monoidal-koszul}
给出的对称幺半结构，得到全局 \(J\simeq\mathrm{Or}\)。
\end{proof}

\begin{conjecture}[Fold 逆极限上的二值跳跃类不可有限化]\label{conj:xi-time-part9i-fold-inverse-limit-nonfinitizability}
沿用 \(\Fold_m\) 的尺度塔与逆极限对象 \(X_\infty=\varprojlim X_m\)（定理 \ref{thm:xi-suspension-flow-universality}）。
对每层 \(m\) 由取向函子 \(\mathrm{Or}\) 所诱导的二值跳跃类记为
\(
\omega_m\in H^1(X_m;\ZZ_2)
\)
（语义锚点见定理 \ref{thm:xi-binary-jump-orientation-unification}）。
则猜想：
\begin{enumerate}
  \item 兼容族 \((\omega_m)_m\) 在逆系统上给出非平凡的 pro-\(\ZZ_2\) 不变量 \(\omega_\infty\)；
  \item 该不变量不可由任何有限层截断决定：对任意 \(m\)，存在 \(n>m\) 与 \(\gamma_n\in\pi_1(X_n)\)，使
  \[
  (p_{n,m})_\ast(\gamma_n)=1,\qquad
  \langle \omega_n,\gamma_n\rangle=1\in\ZZ_2,
  \]
  其中 \(p_{n,m}:X_n\to X_m\) 为层间投影。
\end{enumerate}
等价地，任何试图以有限窗口实现全局单值化的外置方案，在逆极限上都将出现新的奇回路并触发取向翻转。
\end{conjecture}

\paragraph{固定正温软序门的语义障碍与零温相变图景}
沿用 FRACTRAN 一步语义（定义 \ref{def:pom-fractran}、命题 \ref{prop:pom-fractran-prime-translation}）与软 first-fit 族（注 \ref{rem:pom-fractran-zero-temp}、命题 \ref{prop:pom-fractran-soft-firstfit-analytic-dominance}）。

\begin{theorem}[固定正温软序门对硬 first-fit 的零误差实现不可能性]\label{thm:xi-time-part9j-fixed-temp-impossibility}
固定 \(u_0\in(0,1)\)。
考虑二指令 FRACTRAN 程序
\[
f_1=\frac1p,\qquad f_2=p,
\]
其中 \(p\) 为任意素数。
记 \(r=v(N)\) 为素数指数向量，\(I(r):=\{i:\ r\ge \beta_i\}\)，\(i^\ast(r):=\min I(r)\)（若 \(I(r)=\varnothing\) 则停机）。
则有
\[
I\bigl(v(1)\bigr)=\{2\},\qquad I\bigl(v(p)\bigr)=\{1,2\},
\]
故硬 first-fit 在 \(N=1\) 上输出 \(2\)，在 \(N=p\) 上输出 \(1\)。

设 \(\mathcal A\) 为如下控制层：给定输入 \(r\)，其对 \(r\) 的随机访问仅通过黑盒调用
\[
\mathrm{PROJ}_{u_0}(\cdot\mid r),
\]
其余更新为确定性，并在 \(I(r)\neq\varnothing\) 时几乎处处终止并输出一个索引。
若 \(\mathcal A\) 在 \(N=1\) 上以概率 \(1\) 输出 \(2\)，则 \(\mathcal A\) 在 \(N=p\) 上必以正概率输出错误索引 \(2\)；因此不存在在全部输入上零误差复现硬 first-fit 的此类控制层。
\end{theorem}
\begin{proof}
在 \(N=1\) 上，因 \(I(v(1))=\{2\}\)，每次调用 \(\mathrm{PROJ}_{u_0}(\cdot\mid v(1))\) 都确定返回 \(2\)。
于是 \(\mathcal A\) 所见黑盒输出串为确定串，且由“几乎处处终止”可得存在确定调用次数 \(n_0\in\NN\)，使其在前缀 \(2^{n_0}\) 后终止并输出 \(2\)。

在 \(N=p\) 上，有
\[
\Pr_{u_0}(2\mid v(p))
=
\frac{u_0^2}{u_0+u_0^2}
=
\frac{u_0}{1+u_0}.
\]
故事件
\[
E_{n_0}:=\{\text{前 }n_0\text{ 次黑盒调用均返回 }2\}
\]
满足
\[
\Pr(E_{n_0})=\left(\frac{u_0}{1+u_0}\right)^{n_0}>0.
\]
在黑盒查询模型下，\(\mathcal A\) 的控制分支仅依赖已观测输出前缀；因此在事件 \(E_{n_0}\) 上，其执行与 \(N=1\) 情形共享同一长度-\(n_0\) 前缀并输出 \(2\)。
但硬 first-fit 在 \(N=p\) 上应输出 \(1\)。
故 \(\mathcal A\) 在 \(N=p\) 上有正错误概率，矛盾于“全输入零误差”。
\end{proof}

\begin{corollary}[固定温度猜想在 soft first-fit 黑盒实例中的成立]\label{cor:xi-time-part9j-soft-firstfit-conjecture-instance}
\leanverified{paper\_xi\_time\_part9j\_soft\_firstfit\_conjecture\_instance}
在定理 \ref{thm:xi-time-part9j-fixed-temp-impossibility} 的黑盒语义下，猜想
\ref{conj:pom-soft-order-fixed-temp-not-deterministic}
的核心断言在 soft first-fit 实例中成立：固定 \(u_0\in(0,1)\) 时，无法在全部输入上实现硬 first-fit 的零误差语义。
此外，若要求输入无关且几乎处处终止，则不存在把总失败概率压到 \(0\) 的固定正温控制层。
\end{corollary}
\begin{proof}
第一部分即定理 \ref{thm:xi-time-part9j-fixed-temp-impossibility}。
第二部分来自同一定理给出的单输入正失败下界：一旦某输入的单步失败概率为正，则“全输入零失败”不可能。
\end{proof}

\begin{proposition}[固定温度下的可靠性--资源对数下界]\label{prop:xi-time-part9j-log-budget-lower-bound}
沿用定理 \ref{thm:xi-time-part9j-fixed-temp-impossibility} 的记号，设 \(\mathcal A\) 在 \(N=1\) 上对应确定调用次数 \(n_0\)。
则在输入 \(N=p\) 上有显式下界
\[
\Pr\bigl(\mathcal A\text{ 输出 }2\bigr)
\ge
\left(\frac{u_0}{1+u_0}\right)^{n_0}.
\]
因此，对任意目标 \(\varepsilon\in(0,1)\)，若要求该错误概率不超过 \(\varepsilon\)，必须满足
\[
n_0\ \ge\ \frac{\log(1/\varepsilon)}{\log\!\bigl((1+u_0)/u_0\bigr)}.
\]
特别地，在固定 \(u_0>0\) 且 \(n_0<\infty\) 时错误概率严格为正；零误差仅可能在 \(u_0=0\) 的硬化极限中实现。
\end{proposition}
\begin{proof}
第一式为定理 \ref{thm:xi-time-part9j-fixed-temp-impossibility} 中事件 \(E_{n_0}\) 的概率下界。
第二式由
\(
(u_0/(1+u_0))^{n_0}\le \varepsilon
\)
取对数并整理得到。
最后一条由 \((u_0/(1+u_0))^{n_0}>0\)（\(u_0>0,n_0<\infty\)）直接推出。
\end{proof}
\leanverified{paper\_xi\_time\_part9j\_log\_budget\_lower\_bound}

\begin{theorem}[二指令分支循环的长期正确率指数塌缩]\label{thm:xi-time-part9j-two-cycle-exponential-collapse}
沿用程序 \((1/p,p)\) 与温度 \(u_0\in(0,1)\)。
硬 first-fit 轨道从 \(N_0=1\) 出发满足
\[
1\longleftrightarrow p
\]
的周期 \(2\) 循环。
设每一步都独立调用同一黑盒控制层 \(\mathcal A\)，并令
\[
\varepsilon:=\left(\frac{u_0}{1+u_0}\right)^{n_0},
\]
其中 \(n_0\) 为定理 \ref{thm:xi-time-part9j-fixed-temp-impossibility} 中的确定调用次数。
若 \(\widehat N_t\) 为 \(\mathcal A\) 产生的随机轨道，则对任意 \(T\ge 1\),
\[
\Pr\!\bigl(\forall\,0\le t<T,\ \widehat N_t=N_t\bigr)
\le
(1-\varepsilon)^{\lfloor T/2\rfloor}.
\]
特别地，
\[
\Pr\!\bigl(\forall\,t\ge 0,\ \widehat N_t=N_t\bigr)=0.
\]
\end{theorem}
\begin{proof}
每次访问状态 \(N_t=p\) 时，定理 \ref{thm:xi-time-part9j-fixed-temp-impossibility} 给出“该步输出错误索引 \(2\)”的条件概率下界 \(\varepsilon\)。
因此，在一次 \(p\)-访问上保持与硬轨道一致的概率至多 \(1-\varepsilon\)。
前 \(T\) 步中访问 \(p\) 的次数为 \(\lfloor T/2\rfloor\)。
在“每步独立调用”假设下，连乘即得
\[
\Pr(\text{前 }T\text{ 步全程一致})
\le
(1-\varepsilon)^{\lfloor T/2\rfloor}.
\]
令 \(T\to\infty\)，右端收敛到 \(0\)，得到全程一致事件概率为 \(0\)。
\end{proof}
\leanverified{paper\_xi\_time\_part9j\_two\_cycle\_exponential\_collapse}

\begin{conjecture}[正温软序门的零误差语义相变]\label{conj:xi-time-part9j-soft-order-phase-transition}
在最小生成片段 \(\mathbf{PW}^{\ast}\)（注 \ref{rem:pom-pw-star-fragment}）中，考虑“输入无关、几乎处处终止、全输入零误差”的控制流语义。
猜想存在位于 \(u=0\) 的结构相变：
\begin{enumerate}
  \item 当 \(u>0\) 固定时，\(\mathrm{PROJ}_u\) 不可承载可复用的无界可靠优先级控制；硬 first-fit 语义不可在该约束下被零误差实现；
  \item 当 \(u=0\)（硬化）可用时，序门 \(P_{\le}\) 恢复确定优先级选择，并与 FRACTRAN-step 的控制流外置接口相容。
\end{enumerate}
定理 \ref{thm:xi-time-part9j-fixed-temp-impossibility}--定理 \ref{thm:xi-time-part9j-two-cycle-exponential-collapse} 给出该相变在 soft first-fit 最小二指令回路上的证据链。
\end{conjecture}

\begin{remark}[门深预算与温度障碍的同阶化接口]\label{rem:xi-time-part9j-log-budget-gate-depth}
命题 \ref{prop:xi-time-part9j-log-budget-lower-bound} 给出的对数级预算下界可与序门瓶颈刻画（\S\ref{subsec:pom-order-bottleneck}）并置为同一资源不等式：固定正温下，可靠性提升所需调用深度与失败预算的对数阶同向增长，且不可由确定性后处理消去。
\end{remark}

\paragraph{有限状态前向实现的不可逆预算、纤维根几何与 prime-register 视界}
在定理 \ref{thm:online-delay-fold} 与命题 \ref{prop:pom-delay-min} 的口径下，\(\Fold_\infty\) 的前向归一化可由最小延迟 \(3\) 的有限状态在线核实现。
下述结果说明：这种前向有限状态性并不导致后向可逆性的常数内存化；相反，纤维体积、层计数根几何与互素账本位长共同给出不可压缩下界。

\begin{theorem}[Fold 的前向有限状态化不推出常数内存可逆性]\label{thm:xi-time-part9k-fold-forward-online-not-constant-memory-reversible}
\leanverified{paper\_xi\_time\_part9k\_fold\_forward\_online\_not\_constant\_memory\_reversible}
固定 \(m\ge 1\)。
令
\[
\mathcal D_m:X_m\times\{0,1\}^{B}\to\Omega_m
\]
为任意译码器，并允许旁信息映射
\(
a_m:\Omega_m\to\{0,1\}^{B}
\)
依赖于输入微态。
在均匀先验 \(\omega\sim\mathrm{Unif}(\Omega_m)\) 下，其恢复成功概率满足
\[
\mathbb P\!\Bigl(\mathcal D_m\bigl(\Fold_m(\omega),a_m(\omega)\bigr)=\omega\Bigr)
\le
\frac{|X_m|\,2^{B}}{2^m}.
\]
若进一步要求对所有 \(\omega\in\Omega_m\) 都精确恢复，则必要条件为
\[
2^{B}\ge D_m:=\max_{x\in X_m}|\Fold_m^{-1}(x)|.
\]
特别地，由 \(|X_m|=F_{m+2}\) 与定理 \ref{thm:pom-max-fiber}，
\begin{align*}
\mathbb P(\text{success})\ge p>0\ \text{与}\ m\ \text{无关}
&\Longrightarrow
B\ge m\log_2\!\frac{2}{\varphi}+O(1),\\
\text{最坏情形精确可逆}
&\Longrightarrow
B\ge \frac{m}{2}\log_2\varphi+O(1).
\end{align*}
\end{theorem}
\begin{proof}
对每一对 \((x,b)\in X_m\times\{0,1\}^{B}\)，译码器至多输出一个候选微态 \(\mathcal D_m(x,b)\)。
因此全部可能被正确恢复的微态总数至多为 \(|X_m|2^B\)。
在均匀先验下，成功概率即不超过该覆盖比例，得到第一式。

若要求对所有 \(\omega\) 都正确，则对任意固定 \(x\in X_m\)，纤维 \(\Fold_m^{-1}(x)\) 中的不同微态必须对应不同的旁信息值；否则在同一 \((x,b)\) 上译码输出唯一，无法同时恢复两个不同微态。
故 \(d_m(x)\le 2^B\) 对所有 \(x\) 必要，特别地 \(D_m\le 2^B\)。

最后，\(|X_m|=F_{m+2}=\Theta(\varphi^m)\) 给出平均情形下界，而定理 \ref{thm:pom-max-fiber} 的闭式
\(
D_m=\Theta(\varphi^{m/2})
\)
给出最坏情形下界。
将两式取 \(\log_2\) 即得所述渐近结论。证毕。
\end{proof}

\begin{theorem}[Fold 纤维层计数多项式的全实负零点与强对数凹性]\label{thm:xi-time-part9k-fold-fiber-layer-count-realroot-logconcavity}
对任意 \(x\in X_m\)，记
\[
Z_x(t)=\sum_{j\ge 0}N_j(x)t^j
\]
为纤维 \(\Fold_m^{-1}(x)\) 的逆代换层计数多项式。
若 \(\Gamma_x\cong\bigsqcup_{i=1}^{r}P_{\ell_i}\) 是定理 \ref{thm:pom-fiber-indset-factorization} 的路径分量分解，则
\[
Z_x(t)=\prod_{i=1}^{r}I_{\ell_i}(t),
\]
且每个路径因子 \(I_{\ell}(t)\) 的零点显式为
\[
t_{\ell,k}=-\frac{1}{4\cos^2\!\left(\frac{\pi k}{\ell+2}\right)},
\qquad
1\le k\le \left\lfloor\frac{\ell+1}{2}\right\rfloor.
\]
从而 \(Z_x(t)\) 的全部零点均为严格负实数。
特别地，系数列 \((N_j(x))_{j\ge 0}\) 满足 Newton 型不等式；在支持内部的可行指标上，
\[
N_j(x)^2> N_{j-1}(x)\,N_{j+1}(x),
\]
因而 \((N_j(x))_{j\ge 0}\) 严格对数凹并单峰。
\end{theorem}
\begin{proof}
第一条分解由定理 \ref{thm:pom-fiber-indset-factorization} 给出。
路径因子的零点闭式由定理 \ref{thm:pom-path-indset-leyang-cyclotomic-param} 给出，故每个 \(I_{\ell}\) 的全部零点都位于负实轴。
乘积 \(Z_x(t)=\prod_i I_{\ell_i}(t)\) 因而同样仅有负实零点。

随后由定理 \ref{thm:pom-fiber-rewrite-realroot-ulc}，\(Z_x(t)\) 的系数列为 ultra log-concave，特别满足 Newton 型不等式，并推出内部支持上的严格对数凹与单峰性。
证毕。
\end{proof}

\begin{theorem}[互素寻址的常数内存视界]\label{thm:xi-time-part9k-coprime-ledger-memory-horizon}
\leanverified{paper\_xi\_time\_part9k\_coprime\_ledger\_memory\_horizon}
沿用定理 \ref{thm:pom-coprime-ledger-primorial-optimality} 的设定。
若某一互素时间索引账本在任意时刻均以至多 \(B\) 比特精确存储其全局整数 \(G_T\)，则必有
\[
B\ge M\log_2(p_T\#),
\qquad
p_T\#:=\prod_{t=1}^{T}p_t.
\]
进一步，由 \(p_t\ge t+1\) 得
\[
B\ge M\log_2\bigl((T+1)!\bigr)
\ge
\frac{M}{\log 2}\Bigl((T+1)\log(T+1)-(T+1)+1\Bigr).
\]
因此 \(T\log T=O(B/M)\)；
特别地，当 \(B=O(1)\) 时，允许的时间长度 \(T\) 被绝对上界束缚，无法在常数内存下无限延展地维持互素寻址。
\end{theorem}
\begin{proof}
若账本由 \(B\) 比特精确存储，则最坏情形满足 \(G_T\le 2^B\)，即
\[
\log_2 G_T\le B.
\]
另一方面，定理 \ref{thm:pom-coprime-ledger-primorial-optimality} 给出
\[
\log G_T\ge M\log(p_T\#),
\]
等价于
\[
\log_2 G_T\ge M\log_2(p_T\#).
\]
两式合并即得第一条下界。

再由 \(p_t\ge t+1\) 可得
\(
p_T\#\ge (T+1)!
\)。
最后用积分比较
\[
\log((T+1)!)=\sum_{k=1}^{T+1}\log k
\ge
\int_{1}^{T+1}\log x\,dx
\]
即可得到显式下界。证毕。
\end{proof}

\begin{conjecture}[prime Euler lift 的 pro-Čech 障碍增长律]\label{conj:xi-time-part9k-prime-euler-procech-growth}
设 \(\{(\mathcal G_m,M_m,[\alpha_m])\}_{m\ge 1}\) 是一条细化塔，其中每个 \(M_m\) 产生有限层 Euler 对象
\[
\zeta_m(s):=\det(I-e^{-s}M_m)^{-1},
\]
而 \([\alpha_m]\in H^2(\mathcal G_m;\mathrm{Anom}_m)\) 表示结论 \ref{con:xi-implementation-moduli-h2-anom} 与定理 \ref{thm:xi-null-z2-cech-transfer-bockstein} 所刻画的有限层 strictification/Čech 障碍。
记 \(G_m\) 为该层的可见 lift 群（其最大 strictification 商见结论 \ref{con:xi-terminal-strictification-maximal-cp-visible-quotient}）。
定义最小有限 strictification 预算
\[
\mathfrak h_m
:=
\min\Bigl\{|Q|:\ \text{存在有限可见商 }G_m\twoheadrightarrow Q\ \text{使 }[\alpha_m]\text{ 在 }Q\text{-层上可 strictify}\Bigr\}.
\]
若 \(\zeta_m\) 在某个右半平面上经适当归一化收敛到 prime Euler 积
\[
\prod_{p}(1-p^{-s})^{-1},
\]
则必有 \(\mathfrak h_m\to\infty\)；
更强地，预期存在常数 \(c>0\) 使
\[
\log \mathfrak h_m\ge c\,m
\]
沿任意此类细化塔成立。
反之，若 \(\sup_m\mathfrak h_m<\infty\)，则该极限在严格化后仍停留在统一有限状态的 \(\zeta\)-语义层内，从而依定理 \ref{thm:operator-finite-state-zeta-2pii-periodic-separation} 必保留 \(2\pi i\)-周期零极晶格；这与经典 \(\zeta\) 在 \(s=1\) 的孤立极点结构不相容。
\end{conjecture}

\endinput

\paragraph{折叠频谱与输出位势的闭式协调：算术能量、幂散度坐标、整除传播与单比特实验}
在折叠同余的模半环表示、多重度谱的 Parseval 恒等式、bin-fold 的两态一致渐近以及真实输入 \(40\) 状态核的输出位势扭曲谱已经建立之后，还可抽出下列六条彼此兼容的结构结论。

\begin{conclusion}[多重度 Fourier 能量的算术下界与完全平坦性的充要条件]\label{con:xi-time-part9l-fold-fourier-energy-arithmetic-lower-bound}
令
\[
M:=F_{m+2},\qquad N:=2^m,\qquad \overline d_m:=\frac{N}{M},
\]
并将 Euclidean 除法写作
\[
N=qM+s,\qquad 0\le s<M.
\]
则折叠多重度谱满足
\[
\frac{1}{M}\sum_{r\in\ZZ/(M\ZZ)}\bigl(d_m(r)-\overline d_m\bigr)^2
\ge
\frac{s(M-s)}{M^2}.
\]
等价地，
\[
\sum_{k\ne 0}\bigl|\widehat d_m(k)\bigr|^2\ge s(M-s).
\]
并且取等当且仅当 \(d_m(r)\) 仅取两值 \(q\) 与 \(q+1\)，其中恰有 \(s\) 个剩余类取 \(q+1\)，其余 \(M-s\) 个剩余类取 \(q\)。

特别地，
\[
d_m(r)\equiv \overline d_m
\]
成立当且仅当 \(M\mid 2^m\)；否则非零 Fourier 模的总能量必为正，且至少为 \(s(M-s)\)。
\end{conclusion}
\begin{proof}
由定理 \ref{thm:fold-multiplicity-energy}，
\[
\frac{1}{M}\sum_r\bigl(d_m(r)-\overline d_m\bigr)^2
=
\frac{1}{M^2}\sum_{k\ne 0}\bigl|\widehat d_m(k)\bigr|^2,
\]
故只需最小化左端。

设 \((v_r)_{r\in\ZZ/(M\ZZ)}\in\ZZ_{\ge 0}^{M}\) 满足 \(\sum_r v_r=N\)。若存在 \(r_1,r_2\) 使
\[
v_{r_1}\ge v_{r_2}+2,
\]
则把一单位质量从 \(r_1\) 移到 \(r_2\) 后，平方偏差和减少
\[
\bigl(v_{r_1}-\overline d_m\bigr)^2+\bigl(v_{r_2}-\overline d_m\bigr)^2
-\bigl(v_{r_1}-1-\overline d_m\bigr)^2
-\bigl(v_{r_2}+1-\overline d_m\bigr)^2
=2\bigl(v_{r_1}-v_{r_2}-1\bigr)>0.
\]
故极小配置只能出现在任意两坐标相差至多 \(1\) 的情形。由 \(N=qM+s\)，此时必恰有 \(s\) 个坐标等于 \(q+1\)，其余 \(M-s\) 个坐标等于 \(q\)。于是
\[
\sum_r\bigl(d_m(r)-\overline d_m\bigr)^2
=
s\Bigl(1-\frac{s}{M}\Bigr)^2+(M-s)\Bigl(\frac{s}{M}\Bigr)^2
=
\frac{s(M-s)}{M}.
\]
再除以 \(M\) 即得第一式，代回 Parseval 形式得到第二式。取等条件由上述调平过程的严格性给出。最后，\(d_m(r)\equiv \overline d_m\) 等价于 \(s=0\)，亦即 \(M\mid 2^m\)。证毕。
\end{proof}

\begin{corollary}[碰撞概率与均匀性缺口的算术量子化下界]\label{cor:xi-time-part9l-collision-gap-arithmetic-quantization}
沿用结论 \ref{con:xi-time-part9l-fold-fourier-energy-arithmetic-lower-bound} 与结论
\ref{con:xi-fold-collision-spectral-trace-parseval} 的记号。令
\[
M:=F_{m+2},\qquad N:=2^m=qM+s,\qquad 0\le s<M,
\]
并定义碰撞概率
\[
\mathsf{Col}_m:=\sum_{r\in\ZZ/(M\ZZ)}\mu_m(r)^2.
\]
则有精确下界
\[
\boxed{\;
\mathsf{Col}_m
\ge
\frac{1}{F_{m+2}}
+
\frac{s(F_{m+2}-s)}{2^{2m}F_{m+2}}.\;}
\]
等价地，
\[
\boxed{\;
\mathsf{Col}_m-\frac{1}{F_{m+2}}
\ge
\frac{s(F_{m+2}-s)}{2^{2m}F_{m+2}}.\;}
\]
并且取等当且仅当多重度剖面仅取两种相邻整数值
\[
d_m(r)\in\{q,q+1\},
\]
其中恰有 \(s\) 个剩余类取 \(q+1\)，其余 \(F_{m+2}-s\) 个剩余类取 \(q\)。
\end{corollary}
\leanverified{paper\_xi\_time\_part9l\_collision\_gap\_arithmetic\_quantization}

\begin{proof}
由 \(\mu_m=d_m/2^m\) 与结论 \ref{con:xi-fold-collision-spectral-trace-parseval}，
\[
\mathsf{Col}_m-\frac1M
=
\frac{1}{M}\sum_{k\ne 0}\bigl|\widehat\mu_m(k)\bigr|^2
=
\frac{1}{2^{2m}M}\sum_{k\ne 0}\bigl|\widehat d_m(k)\bigr|^2.
\]
再应用结论 \ref{con:xi-time-part9l-fold-fourier-energy-arithmetic-lower-bound} 的 Fourier 能量下界
\[
\sum_{k\ne 0}\bigl|\widehat d_m(k)\bigr|^2\ge s(M-s),
\]
即得
\[
\mathsf{Col}_m-\frac1M\ge \frac{s(M-s)}{2^{2m}M}.
\]
把 \(M=F_{m+2}\) 代回便得到所示不等式。取等条件亦与前述 Fourier 能量下界的取等条件完全等价。证毕。
\end{proof}

\begin{conclusion}[单一 \texorpdfstring{$\chi^2$}{chi2} 极限常数对幂散度谱与输出累积量域的统一坐标化]\label{con:xi-time-part9l-chi2-generates-power-divergence-and-cumulants}
令
\[
f_\alpha(t):=t^\alpha-1-\alpha(t-1)\qquad(\alpha>1),
\]
并记 bin-fold 的极限 \(f\)-散度为 \(D_{f_\alpha}^\infty\)。其中
\[
f_2(t)=(t-1)^2,
\]
故 \(D_{f_2}^\infty\) 正是 bin-fold 相对稳定层均匀基线的 \(\chi^2\)-极限常数。则有
\[
D_{f_2}^\infty=\frac{2\varphi-3}{5},
\qquad
\varphi=\frac{5D_{f_2}^\infty+3}{2},
\qquad
\sqrt5=5D_{f_2}^\infty+2.
\]
从而对任意 \(\alpha>1\)，全部幂散度极限都可写成 \(D_{f_2}^\infty\) 的显式代数函数：
\[
D_{f_\alpha}^\infty
=
5^{-\alpha/2}
\left[
\left(\frac{5D_{f_2}^\infty+3}{2}\right)^{2\alpha-1}
+
\left(\frac{5D_{f_2}^\infty+3}{2}\right)^{\alpha-2}
\right]-1.
\]

进一步，若
\[
P(\theta):=\log\lambda(e^\theta)
\]
为真实输入 \(40\) 状态核的输出位势，则对每个 \(k\ge 1\) 都存在 \(a_k,b_k\in\QQ\) 使
\[
P^{(k)}(0)=a_k+b_k D_{f_2}^\infty.
\]
特别地，
\[
P'(0)=\frac{9}{20}-\frac{7}{4}D_{f_2}^\infty,
\qquad
P''(0)=\frac{989}{1000}-\frac{3983}{200}D_{f_2}^\infty.
\]
\end{conclusion}
\begin{proof}
命题 \ref{prop:fold-bin-phi-identifiable-power-divergence} 已给出
\[
D_{f_\alpha}^\infty
=
\frac{1}{5^{\alpha/2}}\bigl(\varphi^{2\alpha-1}+\varphi^{\alpha-2}\bigr)-1,
\]
且该标量对 \(\varphi\) 严格单调可辨识。取 \(\alpha=2\) 得
\[
D_{f_2}^\infty
=
\frac{\varphi^3+1}{5}-1
=
\frac{\varphi^3-4}{5}.
\]
再用 \(\varphi^2=\varphi+1\) 推得 \(\varphi^3=2\varphi+1\)，于是
\[
D_{f_2}^\infty=\frac{2\varphi-3}{5},
\]
从而得到 \(\varphi\) 与 \(\sqrt5\) 的线性反演公式。将 \(\varphi=(5D_{f_2}^\infty+3)/2\) 代回一般 \(\alpha\) 的闭式即得所示表达式。

另一方面，命题 \ref{prop:real-input-40-output-cumulants-closed} 给出 \(P^{(k)}(0)\in\QQ(\sqrt5)\) 对一切 \(k\ge 1\) 成立。由于
\[
\sqrt5=5D_{f_2}^\infty+2,
\]
任意 \(a+b\sqrt5\in\QQ(\sqrt5)\) 都可唯一写成
\[
(a+2b)+5b\,D_{f_2}^\infty,
\]
故每个 \(P^{(k)}(0)\) 都具有所述仿射表示。表 \ref{tab:real_input_40_output_potential_cumulants_closed} 中前两阶闭式在代入 \(\sqrt5=5D_{f_2}^\infty+2\) 后，恰化为最后两式。证毕。
\end{proof}

\begin{theorem}[bin-fold 幂散度极限谱的二指数刚性与严格对数凸]\label{thm:xi-time-part9l-power-divergence-spectrum-biexponential-rigidity}
沿用命题 \ref{prop:fold-bin-phi-identifiable-power-divergence} 的记号。对 \(\alpha>1\) 定义重整化谱函数
\[
\mathcal E(\alpha):=5^{\alpha/2}\bigl(1+D_{f_\alpha}^\infty\bigr).
\]
则有显式闭式
\[
\boxed{\;
\mathcal E(\alpha)=\varphi^{2\alpha-1}+\varphi^{\alpha-2}
=
\varphi^{-1}e^{2\alpha\log\varphi}+\varphi^{-2}e^{\alpha\log\varphi}.\;}
\]
从而：
\begin{enumerate}
  \item 若记 \(L:=\log\varphi\)，则 \(\mathcal E\) 满足精确二阶常系数微分方程
  \[
  \boxed{\;
  \mathcal E''(\alpha)-3L\,\mathcal E'(\alpha)+2L^2\,\mathcal E(\alpha)=0.\;}
  \]
  \item \(\mathcal E\) 严格对数凸，且
  \[
  \boxed{\;
  \frac{d^2}{d\alpha^2}\log \mathcal E(\alpha)
  =
  (\log\varphi)^2\,
  \frac{\varphi^{\alpha+1}}{(1+\varphi^{\alpha+1})^2}>0.\;}
  \]
  \item 对任意固定 \(\alpha_0>1\) 与任意 \(n\ge 0\)，有
  \[
  \boxed{\;
  \mathcal E^{(n)}(\alpha_0)
  =
  (2\log\varphi)^n\varphi^{2\alpha_0-1}
  +(\log\varphi)^n\varphi^{\alpha_0-2}.\;}
  \]
  因而 \(\mathcal E\) 在任意点的全部 Taylor 喷射都由单一参数 \(\varphi\) 完全决定；结合命题 \ref{prop:fold-bin-phi-identifiable-power-divergence} 的严格可辨识性，任意一个固定 \(\alpha>1\) 上的标量 \(D_{f_\alpha}^\infty\) 已足以唯一锁定整条重整化谱。
\end{enumerate}
\end{theorem}
\begin{proof}
由命题 \ref{prop:fold-bin-phi-identifiable-power-divergence}，
\[
D_{f_\alpha}^\infty
=
\frac{1}{5^{\alpha/2}}
\bigl(\varphi^{2\alpha-1}+\varphi^{\alpha-2}\bigr)-1.
\]
故按定义
\[
\mathcal E(\alpha)=\varphi^{2\alpha-1}+\varphi^{\alpha-2}
=
\varphi^{-1}e^{2\alpha\log\varphi}+\varphi^{-2}e^{\alpha\log\varphi}.
\]
这说明 \(\mathcal E\) 是两个指数函数的线性组合，其特征指数为 \(L\) 与 \(2L\)，于是
\[
(\partial_\alpha-L)(\partial_\alpha-2L)\mathcal E=0,
\]
即
\[
\mathcal E''-3L\,\mathcal E'+2L^2\,\mathcal E=0.
\]

再将 \(\mathcal E\) 写成
\[
\mathcal E(\alpha)=\varphi^{-2}e^{\alpha\log\varphi}\bigl(1+\varphi^{\alpha+1}\bigr),
\]
便有
\[
\log \mathcal E(\alpha)
=
-2\log\varphi+\alpha\log\varphi+\log\bigl(1+\varphi^{\alpha+1}\bigr).
\]
对 \(\alpha\) 连续求导两次得到
\[
\frac{d^2}{d\alpha^2}\log \mathcal E(\alpha)
=
(\log\varphi)^2\,
\frac{\varphi^{\alpha+1}}{(1+\varphi^{\alpha+1})^2}>0,
\]
故 \(\mathcal E\) 严格对数凸。

第三条只需对
\[
\mathcal E(\alpha)=\varphi^{-1}e^{2\alpha\log\varphi}+\varphi^{-2}e^{\alpha\log\varphi}
\]
逐次求导并在 \(\alpha=\alpha_0\) 处代入即可。于是任意点处的全部导数都由 \(\varphi\) 单独决定。最后，命题 \ref{prop:fold-bin-phi-identifiable-power-divergence} 已证明：对每个固定 \(\alpha>1\)，标量 \(D_{f_\alpha}^\infty\) 对 \(\varphi\) 严格单调可辨识，故一旦该标量确定，整条 \(\mathcal E\)-谱与等价的 \(D_{f_\alpha}^\infty\)-谱便都被唯一固定。证毕。
\end{proof}

\begin{conclusion}[Dirichlet 扭曲误差指数沿整除链单调传播]\label{con:xi-time-part9l-dirichlet-twist-divisibility-monotonicity}
对真实输入 \(40\) 状态核，沿用命题 \ref{prop:real-input-40-output-dirichlet-nonrh} 的记号。若 \(m,n\ge 2\) 且 \(m\mid n\)，则对每个 \(1\le j\le m-1\)，令
\[
J:=j\,\frac{n}{m}.
\]
则 \(1\le J\le n-1\)，并且
\[
\omega_n^{J}=\omega_m^j,
\qquad
\rho\!\bigl(M(\omega_n^J)\bigr)=\rho\!\bigl(M(\omega_m^j)\bigr).
\]
特别地，
\[
\rho_n\ge \rho_m,
\qquad
\theta_n\ge \theta_m.
\]
因此，一旦存在某个模数 \(m_0\ge 2\) 与某个 \(1\le j_0\le m_0-1\) 满足
\[
\rho\!\bigl(M(\omega_{m_0}^{j_0})\bigr)>\sqrt{\lambda},
\]
则对每个倍数 \(n\in m_0\NN\)，取
\[
J:=j_0\,\frac{n}{m_0},
\]
便有
\[
\rho\!\bigl(M(\omega_n^J)\bigr)>\sqrt{\lambda}.
\]
从而对应模数 \(n\) 仍为坏模数。等价地，只要平方根门槛在某一模数上失效一次，就会沿倍数塔无穷遗传。作为聚合版本，若某个模数 \(m_0\) 满足
\[
\theta_{m_0}>\frac12,
\]
则其全部倍数 \(n\in m_0\NN\) 也都满足
\[
\theta_n>\frac12.
\]
换言之，平方根门槛的失效不会局限于单一模数，而会沿整除链向上传播为一条无限倍数序列。
\end{conclusion}
\begin{proof}
记 \(\omega_\ell:=e^{2\pi i/\ell}\)。若 \(m\mid n\)，则对任意 \(1\le j\le m-1\) 有
\[
\omega_m^j=\omega_n^{jn/m},
\]
其中 \(jn/m\) 是 \(1,\dots,n-1\) 之间的整数。故每个非平凡 \(m\)-次单位根扭曲都出现在 \(n\)-次单位根扭曲集合之内，并且逐点保持同一矩阵取值：
\[
\rho\!\bigl(M(\omega_n^{jn/m})\bigr)=\rho\!\bigl(M(\omega_m^j)\bigr).
\]
于是
\[
\rho_n
=
\max_{1\le \ell\le n-1}\rho\!\bigl(M(\omega_n^\ell)\bigr)
\ge
\max_{1\le j\le m-1}\rho\!\bigl(M(\omega_m^j)\bigr)
=
\rho_m.
\]
又因 \(\lambda=\rho(M(1))=\varphi^2>1\)，函数 \(x\mapsto \log x/\log\lambda\) 单调递增，故
\[
\theta_n=\frac{\log\rho_n}{\log\lambda}\ge \frac{\log\rho_m}{\log\lambda}=\theta_m.
\]
若某个特定 \(j_0\) 已满足 \(\rho(M(\omega_{m_0}^{j_0}))>\sqrt{\lambda}\)，则对任意倍数 \(n\in m_0\NN\) 取 \(J=j_0n/m_0\) 即保持同一单位根值，故坏谱半径逐层原封不动地嵌入所有倍数模数。最后一条聚合版表述只是将该结论应用于 \(\theta_{m_0}=\log\rho_{m_0}/\log\lambda\) 的定义。证毕。
\end{proof}

\begin{conclusion}[互素 CRT 分解下的多重度相互作用能与可加分离判据]\label{con:xi-time-part9l-crt-interaction-energy-additive-separation}
设
\[
F_{m+2}=uv,\qquad \gcd(u,v)=1.
\]
借助定理 \ref{thm:fold-congruence-mod-semirings} 与中国剩余定理，可将折叠多重度谱视为定义在
\[
\ZZ/u\ZZ\times \ZZ/v\ZZ
\]
上的函数 \(d_m(a,b)\)。定义二维 Fourier 变换
\[
\widehat d_m(\xi,\eta)
:=
\sum_{a\bmod u}\sum_{b\bmod v}
d_m(a,b)\,e^{-2\pi i(\xi a/u+\eta b/v)}.
\]
再令 \(\overline d_m:=2^m/F_{m+2}\)，并定义三种能量
\[
\mathcal E_u:=\frac{1}{F_{m+2}^2}\sum_{\xi\ne 0}\bigl|\widehat d_m(\xi,0)\bigr|^2,
\qquad
\mathcal E_v:=\frac{1}{F_{m+2}^2}\sum_{\eta\ne 0}\bigl|\widehat d_m(0,\eta)\bigr|^2,
\]
\[
\mathcal E_{\mathrm{int}}
:=
\frac{1}{F_{m+2}^2}\sum_{\xi\ne 0,\ \eta\ne 0}\bigl|\widehat d_m(\xi,\eta)\bigr|^2.
\]
则有精确分解
\[
\frac{1}{F_{m+2}}\sum_{a,b}\bigl(d_m(a,b)-\overline d_m\bigr)^2
=
\mathcal E_u+\mathcal E_v+\mathcal E_{\mathrm{int}}.
\]
并且以下两条严格等价：
\begin{enumerate}
  \item \(\mathcal E_{\mathrm{int}}=0\)；
  \item 存在函数 \(f:\ZZ/u\ZZ\to\CC\) 与 \(g:\ZZ/v\ZZ\to\CC\)，使得
  \[
  d_m(a,b)=f(a)+g(b)-\overline d_m
  \qquad
  (\forall\,a,b).
  \]
\end{enumerate}
因此 \(\mathcal E_{\mathrm{int}}\) 精确刻画了多重度谱在互素分量之间是否存在不可分离的混合耦合。
\end{conclusion}
\begin{proof}
由中国剩余定理，\(\ZZ/(F_{m+2}\ZZ)\cong \ZZ/u\ZZ\times \ZZ/v\ZZ\)。在该乘积群上对函数 \(d_m-\overline d_m\) 应用离散 Parseval 恒等式，得到
\[
\frac{1}{F_{m+2}}\sum_{a,b}\bigl(d_m(a,b)-\overline d_m\bigr)^2
=
\frac{1}{F_{m+2}^2}
\sum_{(\xi,\eta)\ne (0,0)}
\bigl|\widehat d_m(\xi,\eta)\bigr|^2.
\]
将非零频率按 \((\xi,0)\)、\((0,\eta)\) 与 \((\xi,\eta)\)（两者皆非零）三类分组，即得所示分解。

若 \(\mathcal E_{\mathrm{int}}=0\)，则所有混合频率 \(\widehat d_m(\xi,\eta)\) 在 \(\xi\ne 0,\eta\ne 0\) 时全为零。由逆 Fourier 展开，
\[
d_m(a,b)
=
\overline d_m
+
\frac{1}{F_{m+2}}\sum_{\xi\ne 0}\widehat d_m(\xi,0)e^{2\pi i\xi a/u}
+
\frac{1}{F_{m+2}}\sum_{\eta\ne 0}\widehat d_m(0,\eta)e^{2\pi i\eta b/v},
\]
故 \(d_m(a,b)\) 必可写成 \(\overline d_m\) 加上仅依赖于 \(a\) 与仅依赖于 \(b\) 的两项之和。将常数分别吸收到两项中即可得到
\[
d_m(a,b)=f(a)+g(b)-\overline d_m.
\]

反之，若 \(d_m(a,b)=f(a)+g(b)-\overline d_m\)，则对任意 \(\xi\ne 0,\eta\ne 0\) 有
\[
\widehat d_m(\xi,\eta)
=
\sum_{a,b}f(a)e^{-2\pi i(\xi a/u+\eta b/v)}
+
\sum_{a,b}g(b)e^{-2\pi i(\xi a/u+\eta b/v)}
-\overline d_m\sum_{a,b}e^{-2\pi i(\xi a/u+\eta b/v)}.
\]
其中第一项含有因子 \(\sum_b e^{-2\pi i\eta b/v}=0\)，第二项含有因子 \(\sum_a e^{-2\pi i\xi a/u}=0\)，第三项同样为零，故 \(\widehat d_m(\xi,\eta)=0\)。于是 \(\mathcal E_{\mathrm{int}}=0\)。证毕。
\end{proof}

\begin{conclusion}[escort 统计实验的单比特渐近充分性]\label{con:xi-time-part9l-escort-single-bit-sufficiency}
对每个固定 \(q\ge 0\)，令 \(\pi_{m,q}\) 为 bin-fold 的 escort 分布，并记
\[
\theta(q):=\frac{1}{1+\varphi^{q+1}}.
\]
将稳定词空间按末位分成
\[
X_m^{(b)}:=\{w\in X_m:\ w_m=b\},
\qquad
b\in\{0,1\},
\]
并记 \(U_{m,b}\) 为 \(X_m^{(b)}\) 上的均匀分布。再定义混合分布
\[
\Lambda_{m,q}:=(1-\theta(q))\,U_{m,0}+\theta(q)\,U_{m,1}.
\]
则有
\[
D_{\mathrm{TV}}\bigl(\pi_{m,q},\Lambda_{m,q}\bigr)
=
O\!\left(\Bigl(\frac{\varphi}{2}\Bigr)^m\right).
\]
等价地，若 \(L_m\) 为由 \(L_m(b,\cdot)=U_{m,b}\) 给出的两点到 \(X_m\) 的 Markov 核，则 \(\Lambda_{m,q}\) 正是 \(L_m\) 作用于 \(\mathrm{Bernoulli}(\theta(q))\) 的输出分布；因此参数族 \(\{\pi_{m,q}\}_{q\ge 0}\) 在指数精度上由单一末位比特 \(w_m\) 所充分化。
\end{conclusion}
\begin{proof}
记 \(\varepsilon_m:=(\varphi/2)^m\)。由定理 \ref{thm:fold-bin-two-state-asymptotic} 与命题 \ref{prop:fold-bin-power-sum-asymptotic}，
\[
\pi_{m,q}(w)
=
\frac{\bigl(d_m^{\mathrm{bin}}(w)\bigr)^q}{S_q^{\mathrm{bin}}(m)}
=
\frac{\varphi^{-q w_m}}{F_{m+1}+\varphi^{-q}F_m}\bigl(1+O(\varepsilon_m)\bigr),
\qquad
w\in X_m,
\]
其中误差在 \(w\) 上一致。故在固定末位 \(w_m=b\) 的切片内，\(\pi_{m,q}(w)\) 仅通过公共因子 \(\varphi^{-qb}\) 依赖于 \(w\)，从而
\[
D_{\mathrm{TV}}\bigl(\pi_{m,q}(\,\cdot\,\mid w_m=b),U_{m,b}\bigr)=O(\varepsilon_m).
\]
若记
\[
\theta_m(q):=\pi_{m,q}(X_m^{(1)}),
\]
则由条件分布分解
\[
\pi_{m,q}
=(1-\theta_m(q))\,\pi_{m,q}(\,\cdot\,\mid w_m=0)+\theta_m(q)\,\pi_{m,q}(\,\cdot\,\mid w_m=1)
\]
可得
\[
D_{\mathrm{TV}}\Bigl(\pi_{m,q},(1-\theta_m(q))U_{m,0}+\theta_m(q)U_{m,1}\Bigr)
=
O(\varepsilon_m).
\]
另一方面，命题 \ref{prop:fold-bin-escort-lastbit} 给出
\[
\theta_m(q)=\theta(q)+O(\varepsilon_m).
\]
由于 \(U_{m,0}\) 与 \(U_{m,1}\) 支撑互不相交，
\[
D_{\mathrm{TV}}\Bigl((1-\theta_m(q))U_{m,0}+\theta_m(q)U_{m,1},\Lambda_{m,q}\Bigr)
=
\bigl|\theta_m(q)-\theta(q)\bigr|
=
O(\varepsilon_m).
\]
由三角不等式即得
\[
D_{\mathrm{TV}}\bigl(\pi_{m,q},\Lambda_{m,q}\bigr)=O(\varepsilon_m).
\]
证毕。
\end{proof}

\begin{corollary}[bin-fold escort 温度族的双块均匀正则形]\label{cor:xi-time-part9l-escort-two-block-normal-form}
沿用结论 \ref{con:xi-time-part9l-escort-single-bit-sufficiency} 的记号，并记
\[
A_0:=X_m^{(0)},
\qquad
A_1:=X_m^{(1)},
\qquad
u_{A_b}:=U_{m,b}\quad (b=0,1).
\]
则对任意固定 \(q\ge 0\)，存在总质量为零的有符号测度 \(\varepsilon_{m,q}\) 使
\[
\pi_{m,q}
=
(1-\theta(q))u_{A_0}
+
\theta(q)u_{A_1}
+
\varepsilon_{m,q},
\]
并且
\[
\|\varepsilon_{m,q}\|_{\mathrm{TV}}
=
O\!\left(\Bigl(\frac{\varphi}{2}\Bigr)^m\right).
\]
换言之，escort 温度族在全分布层面并不只把末位边际压缩为一维参数，而是以指数小误差退化为两个末位块上的均匀分布凸组合。
\end{corollary}
\begin{proof}
记
\[
\Lambda_{m,q}:=(1-\theta(q))u_{A_0}+\theta(q)u_{A_1}.
\]
由结论 \ref{con:xi-time-part9l-escort-single-bit-sufficiency}，
\[
D_{\mathrm{TV}}(\pi_{m,q},\Lambda_{m,q})
=
O\!\left(\Bigl(\frac{\varphi}{2}\Bigr)^m\right).
\]
取
\[
\varepsilon_{m,q}:=\pi_{m,q}-\Lambda_{m,q},
\]
则 \(\varepsilon_{m,q}\) 为总质量为零的有符号测度，且在有符号测度的标准总变差范数下
\[
\|\varepsilon_{m,q}\|_{\mathrm{TV}}
=
2D_{\mathrm{TV}}(\pi_{m,q},\Lambda_{m,q})
=
O\!\left(\Bigl(\frac{\varphi}{2}\Bigr)^m\right).
\]
证毕。
\end{proof}

\paragraph{bin-fold 压力谱与 fold 共振缺口的统一后果}
在 bin-fold 的两态一致渐近、escort 一比特 Gibbs 化以及一般 fold 多重度谱的 Parseval--Fourier 表示已经建立之后，还可抽取出下列四条彼此衔接的后果。

\begin{conclusion}[bin-fold 正矩阶压力谱的全局线性化与阈值退化]\label{con:xi-time-part9m-binfold-linear-pressure-no-positive-threshold}
对每个整数 \(a\ge 1\)，定义
\[
S_a^{\mathrm{bin}}(m):=\sum_{w\in X_m}\bigl(d_m^{\mathrm{bin}}(w)\bigr)^a,
\qquad
P_a^{\mathrm{bin}}:=\lim_{m\to\infty}\frac{1}{m}\log S_a^{\mathrm{bin}}(m).
\]
则极限存在并满足闭式
\[
P_a^{\mathrm{bin}}
=
a\log\!\Bigl(\frac{2}{\varphi}\Bigr)+\log\varphi.
\]
更精确地，
\[
S_a^{\mathrm{bin}}(m)
=
\Bigl(\frac{2}{\varphi}\Bigr)^{am}
\bigl(F_{m+1}+\varphi^{-a}F_m\bigr)
\bigl(1+O\!\bigl((\varphi/2)^m\bigr)\bigr).
\]
因此 \(a\mapsto P_a^{\mathrm{bin}}\) 在正整数半轴上严格线性，其 Legendre--Fenchel 对偶退化为对单一斜率 \(\log(2/\varphi)\) 的单点支撑。换言之，bin-fold 的正温压力轴上不存在非平凡冻结阈值；若仍采用冻结语言，则相应临界点已退化到原点。
\end{conclusion}
\begin{proof}
命题 \ref{prop:fold-bin-power-sum-asymptotic} 已给出
\[
S_a^{\mathrm{bin}}(m)
=
\Bigl(\frac{2}{\varphi}\Bigr)^{am}
\bigl(F_{m+1}+\varphi^{-a}F_m\bigr)
\bigl(1+O\!\bigl(a(\varphi/2)^m\bigr)\bigr)
\]
对每个固定整数 \(a\ge 1\) 成立。由于 \(a\) 固定，误差即为 \(O((\varphi/2)^m)\)。再由 Binet 渐近
\[
F_{m+1}+\varphi^{-a}F_m=\Theta(\varphi^m),
\]
取对数并除以 \(m\) 即得
\[
P_a^{\mathrm{bin}}
=
a\log\!\Bigl(\frac{2}{\varphi}\Bigr)+\log\varphi.
\]
线性性由闭式立即推出。对偶只保留单一斜率也随之成立，故正整数半轴上不再出现新的冻结转折。证毕。
\end{proof}

\begin{conclusion}[bin-fold escort 族的末位偏置与 \texorpdfstring{$\log\varphi$}{log phi} 型 min-entropy 常数]\label{con:xi-time-part9m-binfold-escort-soft-freezing-minentropy}
对每个整数 \(a\ge 1\)，定义 escort 分布
\[
\pi_{m,a}^{\mathrm{bin}}(w):=
\frac{\bigl(d_m^{\mathrm{bin}}(w)\bigr)^a}{S_a^{\mathrm{bin}}(m)},
\qquad
w\in X_m.
\]
则末位边际满足
\[
\pi_{m,a}^{\mathrm{bin}}(w_m=0)\longrightarrow \frac{\varphi^{a+1}}{\varphi^{a+1}+1},
\qquad
\pi_{m,a}^{\mathrm{bin}}(w_m=1)\longrightarrow \frac{1}{\varphi^{a+1}+1}.
\]
因此当 \(a\to\infty\) 时，escort 质量渐近集中到末位 \(0\) 扇区；然而由结论 \ref{con:xi-time-part9m-binfold-linear-pressure-no-positive-threshold}，此种扇区偏置并不伴随正温压力谱中的非线性转折。

进一步，定义 min-entropy 率
\[
h_{\infty}^{\mathrm{bin}}(a):=
\lim_{m\to\infty}
-\frac{1}{m}\log\max_{w\in X_m}\pi_{m,a}^{\mathrm{bin}}(w).
\]
则该极限存在，并且对全部整数 \(a\ge 1\) 满足
\[
h_{\infty}^{\mathrm{bin}}(a)=\log\varphi.
\]
\end{conclusion}
\begin{proof}
由命题 \ref{prop:fold-bin-escort-lastbit}，
\[
\pi_{m,a}^{\mathrm{bin}}(w_m=1)\to \theta(a):=\frac{1}{1+\varphi^{a+1}},
\qquad
\pi_{m,a}^{\mathrm{bin}}(w_m=0)\to 1-\theta(a)=\frac{\varphi^{a+1}}{1+\varphi^{a+1}}.
\]
当 \(a\to\infty\) 时，\(\theta(a)\to 0\)，故质量渐近集中于末位 \(0\) 扇区。

对最大点质量，由定理 \ref{thm:fold-bin-two-state-asymptotic} 与命题 \ref{prop:fold-bin-power-sum-asymptotic}，
\[
\pi_{m,a}^{\mathrm{bin}}(w)
=
\frac{\varphi^{-aw_m}}{F_{m+1}+\varphi^{-a}F_m}
\Bigl(1+O\!\bigl((\varphi/2)^m\bigr)\Bigr)
\]
在 \(w\in X_m\) 上一致成立。故最大值总在 \(w_m=0\) 扇区取得，且
\[
\max_{w\in X_m}\pi_{m,a}^{\mathrm{bin}}(w)
=
\frac{1+O((\varphi/2)^m)}{F_{m+1}+\varphi^{-a}F_m}.
\]
由于 \(F_{m+1}+\varphi^{-a}F_m=\Theta(\varphi^m)\)，于是
\[
-\frac{1}{m}\log\max_{w\in X_m}\pi_{m,a}^{\mathrm{bin}}(w)\to \log\varphi.
\]
这就给出 \(h_\infty^{\mathrm{bin}}(a)=\log\varphi\)。最后，由结论 \ref{con:xi-time-part9m-binfold-linear-pressure-no-positive-threshold} 的线性压力谱可知，上述 escort 偏置只改变有限温度下的扇区权重，而不引发新的压力拐点。证毕。
\end{proof}

\begin{conclusion}[碰撞盈余、峰值过剩与非平凡 Fourier 能量的双向夹逼]\label{con:xi-time-part9m-collision-peak-fourier-sandwich}
令
\[
F:=F_{m+2},
\qquad
\overline d_m:=\frac{2^m}{F},
\qquad
M_m:=\max_{r}d_m(r),
\]
\[
S_m:=\max_{k\ne 0}\bigl|\widehat d_m(k)\bigr|,
\qquad
E_m:=\sum_{k\ne 0}\bigl|\widehat d_m(k)\bigr|^2,
\qquad
\mathsf{Col}_m:=\frac{1}{2^{2m}}\sum_r d_m(r)^2.
\]
则有
\[
2^m\sqrt{\frac{\mathsf{Col}_m-\frac1F}{F}}
\ \le\
M_m-\overline d_m
\ \le\
2^m\sqrt{\mathsf{Col}_m-\frac1F},
\]
以及
\[
\frac{S_m}{F}
\ \le\
M_m-\overline d_m
\ \le\
\sqrt{\frac{E_m}{F}}.
\]
特别地，碰撞盈余 \(\mathsf{Col}_m-1/F\)、峰值过剩 \(M_m-\overline d_m\) 与非平凡 Fourier 能量 \(E_m\) 被同一组显式不等式联结；任一量的定量偏离都会强制其余两量发生相应偏离。
\end{conclusion}
\begin{proof}
由定理 \ref{thm:fold-collision-spectrum}，
\[
\mathsf{Col}_m-\frac1F=\frac{E_m}{2^{2m}F}.
\]
另一方面，定理 \ref{thm:fold-max-fiber-spectrum} 给出
\[
M_m-\overline d_m\ge \frac{\sqrt{E_m}}{F}.
\]
将 \(E_m=2^{2m}F(\mathsf{Col}_m-1/F)\) 代回，即得第一组不等式的左端。

再由
\[
\bigl(M_m-\overline d_m\bigr)^2
\le
\sum_r\bigl(d_m(r)-\overline d_m\bigr)^2
\]
以及定理 \ref{thm:fold-multiplicity-energy} 的方差恒等式
\[
\sum_r\bigl(d_m(r)-\overline d_m\bigr)^2=\frac{E_m}{F},
\]
得到
\[
M_m-\overline d_m\le \sqrt{\frac{E_m}{F}}
=
2^m\sqrt{\mathsf{Col}_m-\frac1F},
\]
这给出第一组不等式的右端及第二组不等式的右端。

最后，由定理 \ref{thm:fold-max-fiber-fourier}，
\[
M_m-\overline d_m\ge \frac{S_m}{F},
\]
即得第二组不等式的左端。证毕。
\end{proof}

\begin{conclusion}[黄金共振常数对碰撞缺口与峰值偏置的持久强制]\label{con:xi-time-part9m-cphi-forces-collision-gap-and-peak-bias}
记
\[
a_m(k):=\frac{|\widehat d_m(k)|}{2^m},
\]
并令 \(C_\varphi>0\) 为定理 \ref{thm:fold-critical-resonance-constant} 所定义的黄金共振常数，即
\[
a_m(F_m)\to C_\varphi,
\qquad
a_m(F_{m+1})\to C_\varphi.
\]
则有
\[
\liminf_{m\to\infty}
F_{m+2}\Bigl(\mathsf{Col}_m-\frac{1}{F_{m+2}}\Bigr)
\ge
2C_\varphi^2,
\]
以及
\[
\liminf_{m\to\infty}\frac{M_m}{\overline d_m}
\ge
1+C_\varphi.
\]
因此，\(C_\varphi\) 并非孤立的 Fourier 指纹，而会在碰撞盈余与最大纤维偏置这两类宏观统计量中留下持久且不可消去的下界。
\end{conclusion}
\begin{proof}
由定理 \ref{thm:fold-critical-resonance-constant}，
\[
|\widehat d_m(F_m)|=2^m\bigl(C_\varphi+o(1)\bigr),
\qquad
|\widehat d_m(F_{m+1})|=2^m\bigl(C_\varphi+o(1)\bigr).
\]
于是
\[
E_m=\sum_{k\ne 0}|\widehat d_m(k)|^2
\ge
|\widehat d_m(F_m)|^2+|\widehat d_m(F_{m+1})|^2
=
2\cdot 2^{2m}\bigl(C_\varphi^2+o(1)\bigr).
\]
代入结论 \ref{con:xi-time-part9m-collision-peak-fourier-sandwich} 中
\[
\mathsf{Col}_m-\frac1F=\frac{E_m}{2^{2m}F},
\]
即得
\[
F_{m+2}\Bigl(\mathsf{Col}_m-\frac{1}{F_{m+2}}\Bigr)\ge 2C_\varphi^2+o(1).
\]

另一方面，由同一结论的 Fourier 峰值下界
\[
M_m-\overline d_m\ge \frac{S_m}{F_{m+2}},
\]
且
\[
S_m\ge |\widehat d_m(F_m)|=2^m\bigl(C_\varphi+o(1)\bigr).
\]
故
\[
M_m-\overline d_m
\ge
\frac{2^m}{F_{m+2}}\bigl(C_\varphi+o(1)\bigr)
=
\overline d_m\bigl(C_\varphi+o(1)\bigr).
\]
两边除以 \(\overline d_m\) 即得
\[
\frac{M_m}{\overline d_m}\ge 1+C_\varphi+o(1).
\]
证毕。
\end{proof}

\paragraph{最大纤维偏置、信息缺口与 Fourier 零模稀疏性}
黄金共振常数已经被证明会同时强制碰撞盈余与最大纤维偏置；沿这条链继续推进，还可把峰值偏差依次传递到 \(\ell^\infty\)、碰撞缺口、总变差、相对熵以及非平凡 Fourier 模的非熄灭结论，并在反向上由 KL 缺口给出最大纤维超额的确定性上界。

\begin{theorem}[最大纤维偏置给出的统一非均匀性下界链]\label{thm:xi-time-part9m-maxfiber-uniformity-gap-chain}
\leanpartial{paper\_xi\_time\_part9m\_maxfiber\_uniformity\_gap\_chain}{the stated Lean signature is false without additional nonnegativity hypotheses; from \((C-\varepsilon)/M \le \mathrm{maxBias}\) alone one cannot square the bound or derive the KL estimate honestly}
记
\[
M:=F_{m+2},
\qquad
p_m(r):=\frac{d_m(r)}{2^m},
\qquad
u(r):=\frac{1}{M}
\qquad
\bigl(r\in \ZZ/(M\ZZ)\bigr),
\]
并令
\[
\overline d_m:=\frac{2^m}{M},
\qquad
M_m:=\max_r d_m(r),
\qquad
\mathsf{Col}_m:=\sum_{r}p_m(r)^2.
\]
设 \(C_\varphi>0\) 为定理 \ref{thm:fold-critical-resonance-constant} 的黄金共振常数。则对任意 \(\varepsilon\in(0,C_\varphi)\)，存在 \(m_0(\varepsilon)\)，使得对所有 \(m\ge m_0(\varepsilon)\) 同时成立
\[
\boxed{\;
\|p_m-u\|_{\ell^\infty}\ge \frac{C_\varphi-\varepsilon}{M},\;}
\]
\[
\boxed{\;
\mathsf{Col}_m-\frac{1}{M}
=
\sum_{r\in \ZZ/(M\ZZ)}\bigl(p_m(r)-u(r)\bigr)^2
\ge
\frac{(C_\varphi-\varepsilon)^2}{M^2},\;}
\]
\[
\boxed{\;
\|p_m-u\|_{\mathrm{TV}}
:=
\frac12\sum_{r\in \ZZ/(M\ZZ)}|p_m(r)-u(r)|
\ge
\frac{C_\varphi-\varepsilon}{M},\;}
\]
\[
\boxed{\;
D_{\mathrm{KL}}(p_m\|u)\ge \frac{2(C_\varphi-\varepsilon)^2}{M^2}.\;}
\]
因此，bin-fold 推前分布相对均匀分布的偏离不可能快于 \(M^{-1}\) 消失，而碰撞缺口与相对熵缺口亦不可能快于 \(M^{-2}\) 消失。
\end{theorem}
\begin{proof}
由结论 \ref{con:xi-time-part9m-cphi-forces-collision-gap-and-peak-bias}，
\[
\liminf_{m\to\infty}\frac{M_m}{\overline d_m}\ge 1+C_\varphi.
\]
故对任意 \(\varepsilon\in(0,C_\varphi)\)，当 \(m\) 充分大时有
\[
M_m-\overline d_m\ge (C_\varphi-\varepsilon)\overline d_m.
\]
取达到最大值的 \(r_\ast\)，便得
\[
p_m(r_\ast)-u(r_\ast)
=
\frac{M_m}{2^m}-\frac{1}{M}
=
\frac{M_m-\overline d_m}{2^m}
\ge
\frac{(C_\varphi-\varepsilon)\overline d_m}{2^m}
=
\frac{C_\varphi-\varepsilon}{M},
\]
从而得到 \(\ell^\infty\) 下界。

再由 \(\sum_r p_m(r)=\sum_r u(r)=1\)，直接展开可得
\[
\sum_r\bigl(p_m(r)-u(r)\bigr)^2
=
\sum_r p_m(r)^2-\frac{1}{M}
=
\mathsf{Col}_m-\frac{1}{M}.
\]
该平方和至少不小于任一坐标偏差的平方，于是
\[
\mathsf{Col}_m-\frac{1}{M}
\ge
\|p_m-u\|_{\ell^\infty}^2
\ge
\frac{(C_\varphi-\varepsilon)^2}{M^2}.
\]

另一方面，若某个坐标上有正偏差 \(p_m(r)-u(r)=\eta>0\)，则负偏差总量至少为 \(\eta\)，故
\[
\|p_m-u\|_{\mathrm{TV}}
=
\frac12\sum_r|p_m(r)-u(r)|
\ge \eta.
\]
对 \(r=r_\ast\) 取 \(\eta=p_m(r_\ast)-u(r_\ast)\)，便得总变差下界。

最后由 Pinsker 不等式
\[
D_{\mathrm{KL}}(p_m\|u)\ge 2\|p_m-u\|_{\mathrm{TV}}^2
\]
与前述总变差估计相结合，即得最后一条。证毕。
\end{proof}

\begin{theorem}[KL 缺口对最大纤维超额的确定性控制]\label{thm:xi-time-part9m-kl-controls-maxfiber-excess}
沿用定理 \ref{thm:xi-time-part9m-maxfiber-uniformity-gap-chain} 的记号。则对每个 \(m\ge 1\)，都有
\[
\boxed{\;
M_m-\overline d_m
\le
2^m\sqrt{\frac12\,D_{\mathrm{KL}}(p_m\|u)}.\;}
\]
等价地，
\[
\boxed{\;
M_m-\overline d_m
\le
2^m\sqrt{\frac12\,\bigl(\log M-H(p_m)\bigr)}.\;}
\]
因此，一旦相对熵缺口被控制，最大纤维相对于平均纤维的超额便被同一尺度上的确定性不等式锁定。
\end{theorem}
\begin{proof}
取达到最大值的 \(r_\ast\)。则
\[
\frac{M_m-\overline d_m}{2^m}
=
p_m(r_\ast)-u(r_\ast)
\le
\max_r |p_m(r)-u(r)|.
\]
而对任意概率分布 \(p,u\)，都有
\[
\max_r |p(r)-u(r)|
\le
\|p-u\|_{\mathrm{TV}}.
\]
再由 Pinsker 不等式，
\[
\|p_m-u\|_{\mathrm{TV}}
\le
\sqrt{\frac12\,D_{\mathrm{KL}}(p_m\|u)}.
\]
三式合并即得
\[
\frac{M_m-\overline d_m}{2^m}
\le
\sqrt{\frac12\,D_{\mathrm{KL}}(p_m\|u)}.
\]
两边乘以 \(2^m\) 即得第一条不等式。第二条只需使用恒等式
\[
D_{\mathrm{KL}}(p_m\|u)=\log M-H(p_m).
\]
证毕。
\end{proof}

\begin{corollary}[Fourier 零模不可能在大尺度上完全饱和]\label{cor:xi-time-part9m-fourier-zero-set-nonsaturation}
沿用定理 \ref{thm:xi-time-part9m-maxfiber-uniformity-gap-chain} 的记号，并定义 Fourier 零点集合
\[
\mathcal Z_m:=\{k\in \ZZ/(M\ZZ):\ \widehat d_m(k)=0\}.
\]
则对任意 \(\varepsilon\in(0,C_\varphi)\)，存在 \(m_1(\varepsilon)\)，使得对所有 \(m\ge m_1(\varepsilon)\) 都有
\[
\boxed{\;
M-1-|\mathcal Z_m|
\ge
\frac{(C_\varphi-\varepsilon)^2}{M}.\;}
\]
特别地，当 \(m\) 充分大时，
\[
\boxed{\;
|\mathcal Z_m|\le M-2.\;}
\]
亦即：除平凡频率外，非平凡 Fourier 模不可能在充分大尺度上全部同时熄灭。
\end{corollary}
\begin{proof}
由定义 \ref{def:fold-normalized-amplitude}，记
\[
a_m(k):=\frac{|\widehat d_m(k)|}{2^m}.
\]
则 \(0\le a_m(k)\le 1\)，且 \(a_m(k)=0\) 当且仅当 \(k\in\mathcal Z_m\)。再由定理 \ref{thm:fold-collision-spectrum} 的 Parseval 恒等式，
\[
\mathsf{Col}_m-\frac{1}{M}
=
\frac{1}{M}\sum_{k\ne 0}a_m(k)^2
\le
\frac{1}{M}\#\{k\ne 0:\ a_m(k)\ne 0\}
=
\frac{M-1-|\mathcal Z_m|}{M}.
\]
另一方面，定理 \ref{thm:xi-time-part9m-maxfiber-uniformity-gap-chain} 已给出
\[
\mathsf{Col}_m-\frac{1}{M}
\ge
\frac{(C_\varphi-\varepsilon)^2}{M^2}
\]
对充分大的 \(m\) 成立。合并两式即得
\[
M-1-|\mathcal Z_m|
\ge
\frac{(C_\varphi-\varepsilon)^2}{M}.
\]
由于左侧是非负整数，而右侧对充分大的 \(m\) 为正，遂必有
\[
M-1-|\mathcal Z_m|\ge 1,
\]
即 \(|\mathcal Z_m|\le M-2\)。证毕。
\end{proof}

\begin{theorem}[固定有限频率证书对碰撞尺度的渐近失效]\label{thm:xi-time-part9m-fixed-frequency-collision-certificate-fails}
沿用定理 \ref{thm:xi-time-part9m-maxfiber-uniformity-gap-chain} 的记号，并令
\[
\pi_m:\ZZ\twoheadrightarrow \ZZ/(M\ZZ)
\]
为自然投影。对任意有限集合 \(S\subset\ZZ\)，定义其对应的有限频率截断碰撞量
\[
\mathsf{Col}_m^{(S)}
:=
\frac{1}{M}\sum_{\bar k\in \pi_m(S)}a_m(\bar k)^2.
\]
则有
\[
M\bigl(\mathsf{Col}_m-\mathsf{Col}_m^{(S)}\bigr)\longrightarrow +\infty
\qquad (m\to\infty),
\]
并且
\[
\frac{\mathsf{Col}_m^{(S)}}{\mathsf{Col}_m}\longrightarrow 0.
\]
换言之，任意固定有限频率列表在渐近上都不可能饱和碰撞尺度；主导贡献必然来自频率数目随 \(m\) 增长的无限共振族。
\end{theorem}
\begin{proof}
由定理 \ref{thm:fold-collision-spectrum} 的 Parseval 恒等式，
\[
M\,\mathsf{Col}_m=\sum_{\bar k\in \ZZ/(M\ZZ)}a_m(\bar k)^2.
\]
因此
\[
\begin{aligned}
M\bigl(\mathsf{Col}_m-\mathsf{Col}_m^{(S)}\bigr)
&=
\sum_{\bar k\notin \pi_m(S)}a_m(\bar k)^2\\
&=
M\,\mathsf{Col}_m-\sum_{\bar k\in \pi_m(S)}a_m(\bar k)^2.
\end{aligned}
\]
又由于 \(0\le a_m(\bar k)\le 1\)，有
\[
\sum_{\bar k\in \pi_m(S)}a_m(\bar k)^2\le |\pi_m(S)|\le |S|.
\]
故
\[
M\bigl(\mathsf{Col}_m-\mathsf{Col}_m^{(S)}\bigr)\ge M\,\mathsf{Col}_m-|S|.
\]
而附录中的推论 \ref{cor:fold-collision-scale-diverges} 已证明
\[
M\,\mathsf{Col}_m=F_{m+2}\,\mathsf{Col}_m\longrightarrow +\infty,
\]
从而第一条成立。

另一方面，同样由 \(0\le a_m(\bar k)\le 1\) 得
\[
\mathsf{Col}_m^{(S)}
\le
\frac{|\pi_m(S)|}{M}
\le
\frac{|S|}{M}.
\]
因此
\[
0\le
\frac{\mathsf{Col}_m^{(S)}}{\mathsf{Col}_m}
\le
\frac{|S|}{M\,\mathsf{Col}_m},
\]
而右端再次由推论 \ref{cor:fold-collision-scale-diverges} 趋于 \(0\)。故第二条成立。证毕。
\end{proof}

\paragraph{bin-fold 位率阈值、尖锐阈值窗、缩放 Mellin 读出与临界 defect 原子谱}
在 bin-fold 的两态一致渐近、临界容量双折线与末位两块化已经稳定之后，还可把精确恢复位率、临界比特窗的显式跃迁曲线、缩放多重度的 Mellin 结构以及极限容量曲线的原子二阶导数继续压缩为下列四条终端结论。

\begin{theorem}[bin-fold 离散预言机恢复的精确位率阈值]\label{thm:xi-binfold-oracle-success-bitrate-threshold-exact}
\leanverified{paper\_xi\_binfold\_oracle\_success\_bitrate\_threshold\_exact}
对整数 \(B\ge 0\)，定义 bin-fold 的离散预言机容量与最优恢复成功率
\[
C_m^{\mathrm{bin}}(B)
:=
\sum_{w\in X_m}\min\bigl(d_m^{\mathrm{bin}}(w),2^B\bigr),
\qquad
\mathrm{Succ}_m^{\mathrm{bin}}(B):=\frac{C_m^{\mathrm{bin}}(B)}{2^m}.
\]
再记
\[
\beta_c:=\log_2\!\Bigl(\frac{2}{\varphi}\Bigr).
\]
对任意 \(\beta\ge 0\) 令 \(B_m:=\lfloor \beta m\rfloor\)。则有：
\begin{enumerate}
\item 若 \(\beta<\beta_c\)，则
\[
\lim_{m\to\infty}\frac{1}{m}\log \mathrm{Succ}_m^{\mathrm{bin}}(B_m)
=
-(\beta_c-\beta)\log 2<0.
\]
特别地，\(\mathrm{Succ}_m^{\mathrm{bin}}(B_m)\) 指数衰减到 \(0\)。
\item 若 \(\beta>\beta_c\)，则存在 \(m_0(\beta)\)，使得对一切 \(m\ge m_0(\beta)\) 都有
\[
\mathrm{Succ}_m^{\mathrm{bin}}(B_m)=1.
\]
\end{enumerate}
因此 \(\beta_c\) 是 bin-fold 精确微态恢复的严格位率阈值。
\end{theorem}
\begin{proof}
记
\[
c_m:=\Bigl(\frac{2}{\varphi}\Bigr)^m,
\qquad
\varepsilon_m:=\Bigl(\frac{\varphi}{2}\Bigr)^m.
\]
由定理 \ref{thm:fold-bin-two-state-asymptotic}，存在常数 \(C>0\) 使得对充分大 \(m\) 与一切 \(w\in X_m\) 都有一致展开
\[
d_m^{\mathrm{bin}}(w)=c_m\varphi^{-w_m}\bigl(1+\delta_{m,w}\bigr),
\qquad
|\delta_{m,w}|\le C\varepsilon_m.
\]
于是可取常数 \(0<c_1<c_2<\infty\)，使对充分大 \(m\) 与全部 \(w\in X_m\) 都有
\[
c_1c_m\le d_m^{\mathrm{bin}}(w)\le c_2c_m.
\]

若 \(\beta<\beta_c\)，则
\[
\frac{2^{B_m}}{c_m}\longrightarrow 0.
\]
故当 \(m\) 充分大时，\(2^{B_m}<d_m^{\mathrm{bin}}(w)\) 对全部 \(w\in X_m\) 成立，从而
\[
C_m^{\mathrm{bin}}(B_m)=2^{B_m}|X_m|.
\]
又由 \(|X_m|=F_{m+2}\)，便有
\[
\begin{aligned}
\frac{1}{m}\log \mathrm{Succ}_m^{\mathrm{bin}}(B_m)
&=
\frac{B_m}{m}\log 2-\log 2+\frac{1}{m}\log F_{m+2}\\
&\longrightarrow
\beta\log 2-\log 2+\log\varphi\\
&=
\beta\log 2-\log\!\Bigl(\frac{2}{\varphi}\Bigr)
=
-(\beta_c-\beta)\log 2.
\end{aligned}
\]

若 \(\beta>\beta_c\)，则
\[
\frac{2^{B_m}}{c_m}\longrightarrow \infty.
\]
故对充分大 \(m\) 有 \(2^{B_m}>c_2c_m\ge d_m^{\mathrm{bin}}(w)\) 对全部 \(w\in X_m\) 成立。因此
\[
C_m^{\mathrm{bin}}(B_m)=\sum_{w\in X_m}d_m^{\mathrm{bin}}(w)=2^m,
\]
亦即
\[
\mathrm{Succ}_m^{\mathrm{bin}}(B_m)=1.
\]
证毕。
\end{proof}

\begin{theorem}[bin-fold 预言机可逆的尖锐阈值窗与显式极限曲线]\label{thm:xi-time-part9m2-binfold-sharp-threshold-window-explicit-curve}
\leanverified{paper\_xi\_time\_part9m2\_binfold\_sharp\_threshold\_window\_explicit\_curve}
把定理 \ref{thm:xi-binfold-oracle-success-bitrate-threshold-exact} 的预算沿 dyadic 尺度连续化为实参数 \(B\in\mathbb R\)，定义
\[
C_m(B):=
\sum_{w\in X_m}\min\bigl(d_m^{\mathrm{bin}}(w),2^B\bigr),
\qquad
\mathrm{Succ}_m(B):=\frac{C_m(B)}{2^m}.
\]
记临界比特率
\[
\beta_\star:=\beta_c=\log_2\!\Bigl(\frac{2}{\varphi}\Bigr).
\]
则对任意固定 \(b\in\mathbb R\)，在临界缩放
\[
B=\beta_\star m+b
\]
下有极限
\[
\lim_{m\to\infty}\mathrm{Succ}_m(\beta_\star m+b)=\mathcal S(b),
\]
其中
\[
\mathcal S(b)=
\begin{cases}
\dfrac{\varphi^2\,2^b}{\sqrt5}, & b\le -\log_2\varphi,\\[1.2ex]
\dfrac{\varphi\,2^b+\varphi^{-1}}{\sqrt5}, & -\log_2\varphi\le b\le 0,\\[1.2ex]
1, & b\ge 0.
\end{cases}
\]
特别地，跃迁窗宽度恰为
\[
\log_2\varphi,
\]
因为 \(\mathcal S(b)\) 只在
\[
b\in[-\log_2\varphi,0]
\]
上经历非平凡双段跃迁，并从
\[
\mathcal S(-\log_2\varphi)=\frac{\varphi}{\sqrt5}
\]
单调上升到
\[
\mathcal S(0)=1.
\]
\end{theorem}
\begin{proof}
记
\[
s:=2^b.
\]
由 \(\beta_\star=\log_2(2/\varphi)\) 可知
\[
2^{\beta_\star m+b}=s\left(\frac{2}{\varphi}\right)^m.
\]
因此
\[
\mathrm{Succ}_m(\beta_\star m+b)
=
2^{-m}\,\mathcal C_m^{\mathrm{bin,cont}}\!\left(
s\left(\frac{2}{\varphi}\right)^m
\right)
=
\mathcal S_m(s),
\]
其中 \(\mathcal S_m\) 正是推论
\ref{cor:xi-foldbin-critical-scale-success-curve-inversion}
中的连续化临界成功曲线。于是
\[
\mathrm{Succ}_m(\beta_\star m+b)\longrightarrow \mathcal S_\infty(s).
\]
再把推论
\ref{cor:xi-foldbin-critical-scale-success-curve-inversion}
给出的分段极限
\[
\mathcal S_\infty(s)=
\begin{cases}
\dfrac{\varphi^2}{\sqrt5}\,s, & 0\le s\le \varphi^{-1},\\[6pt]
\dfrac{\varphi}{\sqrt5}\,s+\dfrac{1}{\varphi\sqrt5}, & \varphi^{-1}\le s\le 1,\\[8pt]
1, & s\ge 1
\end{cases}
\]
代入 \(s=2^b\)，便得到所示 \(\mathcal S(b)\)。

最后，
\[
2^b\le \varphi^{-1}
\iff
b\le -\log_2\varphi,
\]
\[
\varphi^{-1}\le 2^b\le 1
\iff
-\log_2\varphi\le b\le 0,
\]
\[
2^b\ge 1
\iff
b\ge 0,
\]
故非平凡跃迁恰发生在区间 \([-\log_2\varphi,0]\)，其宽度即为 \(\log_2\varphi\)。证毕。
\end{proof}

\begin{theorem}[bin-fold 缩放多重度的 Mellin 极限与对数累积量]\label{thm:xi-binfold-scaled-multiplicity-mellin-cumulants}
令
\[
U_m\sim \mathrm{Unif}\{0,1,\dots,2^m-1\},
\qquad
W_m:=\mathrm{Fold}^{\mathrm{bin}}_m(U_m),
\]
并定义缩放多重度随机变量
\[
Z_m:=\Bigl(\frac{\varphi}{2}\Bigr)^m d_m^{\mathrm{bin}}(W_m).
\]
则存在固定紧区间 \(I\Subset(0,\infty)\)，使得当 \(m\) 充分大时 \(Z_m\in I\) 几乎处处成立，并且对任意连续函数 \(f\in C(I)\) 都有
\[
\lim_{m\to\infty}\mathbf E[f(Z_m)]
=
\frac{\varphi}{\sqrt5}f(1)+\frac{1}{\varphi\sqrt5}f(\varphi^{-1}).
\]
特别地，对每个固定复数 \(s\in\CC\)，都有
\[
\lim_{m\to\infty}\mathbf E[Z_m^{\,s}]
=
M(s):=\frac{\varphi+\varphi^{-(s+1)}}{\sqrt5}.
\]
其中 \(M\) 在整个复平面上解析。又因 \(M(0)=1\)，故可在 \(\theta=0\) 的某邻域内取对数支并定义
\[
K(\theta):=\log M(\theta)
=
\log\!\Bigl(\frac{\varphi+\varphi^{-(\theta+1)}}{\sqrt5}\Bigr).
\]
则 \(K\) 在 \(\theta=0\) 邻域内全纯，且 \(\log Z_m\) 的极限对数矩母函数正是 \(K\)；因此全部极限对数累积量由 \(K^{(r)}(0)\) 显式给出。
\end{theorem}
\begin{proof}
仍记
\[
c_m:=\Bigl(\frac{2}{\varphi}\Bigr)^m,
\qquad
\varepsilon_m:=\Bigl(\frac{\varphi}{2}\Bigr)^m.
\]
由定理 \ref{thm:fold-bin-two-state-asymptotic}，对充分大 \(m\) 与全部 \(w\in X_m\) 有
\[
Z_m
=
\Bigl(\frac{\varphi}{2}\Bigr)^m d_m^{\mathrm{bin}}(w)
=
\varphi^{-w_m}\bigl(1+O(\varepsilon_m)\bigr),
\]
其中误差在 \(w\) 上一致。故当 \(m\) 充分大时，
\[
Z_m\in \Bigl[\frac{\varphi^{-1}}{2},2\Bigr]=:I
\]
几乎处处成立。

把稳定词按末位分成两块
\[
X_m^{(b)}:=\{w\in X_m:\ w_m=b\},
\qquad b\in\{0,1\}.
\]
则对 \(w\in X_m^{(b)}\)，有一致极限
\[
Z_m(w)=\varphi^{-b}+O(\varepsilon_m).
\]
另一方面，\(W_m\) 的分布为 bin-fold 推前分布，因此
\[
\mathbf P(W_m\in X_m^{(b)})
=
\frac{1}{2^m}\sum_{w\in X_m^{(b)}}d_m^{\mathrm{bin}}(w).
\]
再由两态一致展开与
\[
\#X_m^{(0)}=F_{m+1},
\qquad
\#X_m^{(1)}=F_m,
\]
得到
\[
\mathbf P(W_m\in X_m^{(0)})
=
\frac{F_{m+1}}{\varphi^m}\bigl(1+O(\varepsilon_m)\bigr)
\longrightarrow
\frac{\varphi}{\sqrt5},
\]
\[
\mathbf P(W_m\in X_m^{(1)})
=
\frac{F_m}{\varphi^{m+1}}\bigl(1+O(\varepsilon_m)\bigr)
\longrightarrow
\frac{1}{\varphi\sqrt5}.
\]
由于 \(f\) 在紧区间 \(I\) 上一致连续，于是
\[
\sup_{w\in X_m^{(b)}}|f(Z_m(w))-f(\varphi^{-b})|\longrightarrow 0
\qquad(b=0,1).
\]
将期望按两块拆分并令 \(m\to\infty\)，便得到
\[
\lim_{m\to\infty}\mathbf E[f(Z_m)]
=
\frac{\varphi}{\sqrt5}f(1)+\frac{1}{\varphi\sqrt5}f(\varphi^{-1}).
\]

对每个固定 \(s\in\CC\)，在 \(I\subset(0,\infty)\) 上取实对数主支，函数 \(z\mapsto z^s=e^{s\log z}\) 连续有界，故可取 \(f(z)=z^s\)。于是
\[
\begin{aligned}
\lim_{m\to\infty}\mathbf E[Z_m^{\,s}]
&=
\frac{\varphi}{\sqrt5}\cdot 1^s
\;+\;
\frac{1}{\varphi\sqrt5}\cdot (\varphi^{-1})^s\\
&=
\frac{\varphi+\varphi^{-(s+1)}}{\sqrt5}.
\end{aligned}
\]
这就给出整个函数 \(M\)。又因 \(M(0)=1\neq 0\)，故 \(M\) 在 \(0\) 的某邻域内无零点，从而 \(K=\log M\) 在该邻域内全纯。最后，
\[
\mathbf E[e^{\theta\log Z_m}]=\mathbf E[Z_m^\theta]\longrightarrow M(\theta)
\]
在 \(\theta=0\) 邻域内成立，于是极限对数矩母函数即为 \(K\)，其各阶导数 \(K^{(r)}(0)\) 正是极限对数累积量。证毕。
\end{proof}

\begin{theorem}[临界容量曲线的分布二阶导数等于极限 defect law 的加权原子谱]\label{thm:xi-binfold-critical-capacity-defect-atomic-law}
定义归一化临界容量曲线
\[
\mathcal C_\infty(z)
:=
\frac{\varphi}{\sqrt5}\min\{1,z\}
+
\frac{1}{\sqrt5}\min\{\varphi^{-1},z\}
\qquad(z>0).
\]
它等价于定理 \ref{thm:xi-foldbin-critical-capacity-broken-line} 的极限轮廓在变量 \(z=e^s\) 下的写法。再令
\[
\nu
:=
\frac{\varphi}{\sqrt5}\delta_1
+
\frac{1}{\varphi\sqrt5}\delta_{\varphi^{-1}}
\]
为定理 \ref{thm:xi-binfold-scaled-multiplicity-mellin-cumulants} 给出的极限缩放多重度分布。则有精确表示
\[
\mathcal C_\infty(z)
=
\int_{(0,\infty)}\min\!\Bigl(1,\frac{z}{u}\Bigr)\,\nu(\dd u).
\]
并且在分布意义下，
\[
-\frac{\dd^2}{\dd z^2}\mathcal C_\infty
=
\frac{1}{\sqrt5}\,\delta_{\varphi^{-1}}
+
\frac{\varphi}{\sqrt5}\,\delta_1
=
z^{-1}\nu(\dd z).
\]
因此 \(\mathcal C_\infty\) 不仅给出临界容量轮廓，而且完整编码了极限 defect spectrum；其分布二阶导数可无损恢复 \(\nu\)。
\end{theorem}
\begin{proof}
由 \(\nu\) 的两原子表达，
\[
\begin{aligned}
\int_{(0,\infty)}\min\!\Bigl(1,\frac{z}{u}\Bigr)\,\nu(\dd u)
&=
\frac{\varphi}{\sqrt5}\min(1,z)
\;+\;
\frac{1}{\varphi\sqrt5}\min(1,\varphi z)\\
&=
\frac{\varphi}{\sqrt5}\min(1,z)
\;+\;
\frac{1}{\sqrt5}\min(\varphi^{-1},z),
\end{aligned}
\]
即得第一式。

再固定 \(u>0\)，考虑函数
\[
h_u(z):=\min\!\Bigl(1,\frac{z}{u}\Bigr).
\]
则 \(h_u\) 的分布一阶导数为
\[
\frac{\dd}{\dd z}h_u(z)=u^{-1}\mathbf 1_{(0,u)}(z),
\]
从而二阶分布导数满足
\[
\frac{\dd^2}{\dd z^2}h_u(z)=-u^{-1}\delta_u.
\]
对 \(\nu\) 积分便得到
\[
-\frac{\dd^2}{\dd z^2}\mathcal C_\infty
=
\int_{(0,\infty)}u^{-1}\delta_u\,\nu(\dd u)
=
z^{-1}\nu(\dd z).
\]
最后把 \(\nu\) 的显式形式代回，便有
\[
z^{-1}\nu(\dd z)
=
\frac{\varphi}{\sqrt5}\delta_1
\;+\;
\frac{1}{\varphi\sqrt5}\cdot\varphi\,\delta_{\varphi^{-1}}
=
\frac{\varphi}{\sqrt5}\delta_1+\frac{1}{\sqrt5}\delta_{\varphi^{-1}}.
\]
故 \(-\mathcal C_\infty''\) 的原子位置与权重恰与 \(\nu\) 一一对应，从而 \(\nu\) 可由 \(\mathcal C_\infty\) 的分布二阶导数唯一恢复。证毕。
\end{proof}

\endinput

\paragraph{归一化可见度、残余相干、经典化成本与解析腔室刚性}
在多重度剖面的 Fourier 能量下界、连续容量曲线的 layer-cake 对偶、实阶矩的 Mellin--Stieltjes 反演、全微态精确反演位率以及实参数无退化窗口已经建立之后，还可进一步把非零模可见度的算术障碍、残余相干的微分结构、经典化成本的可见度下界与解析读出的腔室刚性压缩为下列若干条结果。下文固定分辨率 \(m\)，并记
\[
M:=F_{m+2}=|X_m|.
\]
\[
G_m:=\ZZ/(M\ZZ),
\qquad
p_m(r):=\frac{d_m(r)}{2^m}
\qquad
(r\in G_m).
\]

\begin{theorem}[全模可见度同时湮灭的整除障碍]\label{thm:xi-time-part9m3-all-nonzero-visibility-extinction-divisibility-obstruction}
\leanverified{paper\_xi\_time\_part9m3\_all\_nonzero\_visibility\_extinction\_divisibility\_obstruction}
在该循环群上定义归一化字符振幅
\[
\mathcal A_m(k):=\sum_{r\in G_m}p_m(r)e^{-2\pi i kr/M}
=
\frac{\widehat d_m(k)}{2^m}
\qquad
(k\in G_m).
\]
则下列两条命题等价：
\[
\mathcal A_m(k)=0
\qquad
\text{对一切 }k\in G_m\setminus\{0\},
\]
\[
d_m(r)\equiv \frac{2^m}{M}
\qquad
\text{对一切 }r\in G_m.
\]
特别地，若 \(M\nmid 2^m\)，则必存在某个非零频率 \(k\in G_m\setminus\{0\}\) 使
\[
\mathcal A_m(k)\neq 0.
\]
进一步，当
\[
m\not\equiv 1\pmod 3
\]
时，\(F_{m+2}\) 为奇数，故必不整除 \(2^m\)；于是全部非零模同时湮灭在该分辨率族上被纯算术所禁止。
\end{theorem}
\begin{proof}
有限循环群上的 Fourier 反演给出
\[
p_m(r)=\frac1M\sum_{k\in G_m}\mathcal A_m(k)e^{2\pi i kr/M}.
\]
若对一切 \(k\neq 0\) 都有 \(\mathcal A_m(k)=0\)，则因
\[
\mathcal A_m(0)=\sum_{r\in G_m}p_m(r)=1,
\]
可得
\[
p_m(r)\equiv \frac1M,
\]
从而
\[
d_m(r)\equiv \frac{2^m}{M}.
\]
反向方向是平凡的：若 \(d_m\) 常值，则 \(p_m\) 为 \(G_m\) 上的均匀分布，其全部非平凡 Fourier 系数均为零。

若 \(M\nmid 2^m\)，则 \(\frac{2^m}{M}\notin\ZZ\)，从而上式常值条件不可能成立，因此至少有一个非零频率满足 \(\mathcal A_m(k)\neq 0\)。最后，Fibonacci 数列的奇偶性以周期 \(3\) 循环，故 \(F_n\) 为偶数当且仅当 \(3\mid n\)。于是当 \(m\not\equiv 1\pmod 3\) 时，\(M=F_{m+2}\) 为奇数；任何大于 \(1\) 的奇数都不可能整除 \(2^m\)，故结论成立。证毕。
\end{proof}

\begin{theorem}[归一化均方可见度的算术下界]\label{thm:xi-time-part9m3-mean-square-visibility-arithmetic-floor}
\leanverified{paper\_xi\_time\_part9m3\_mean\_square\_visibility\_arithmetic\_floor}
沿用定理 \ref{thm:xi-time-part9m3-all-nonzero-visibility-extinction-divisibility-obstruction} 的记号，并定义
\[
\Sigma_m^2
:=
\frac{1}{M-1}\sum_{k\in G_m\setminus\{0\}}|\mathcal A_m(k)|^2.
\]
将 \(2^m\) 除以 \(M\) 写成
\[
2^m=qM+s,
\qquad
0\le s<M.
\]
则有精确下界
\[
\Sigma_m^2\ge \frac{s(M-s)}{(M-1)2^{2m}}.
\]
并且等号成立，当且仅当多重度剖面达到唯一的最均匀整数配置：
\[
d_m(r)\in\{q,q+1\}
\qquad
(r\in G_m),
\]
且恰有 \(s\) 个剩余类满足 \(d_m(r)=q+1\)。
\end{theorem}
\begin{proof}
由定义
\[
|\mathcal A_m(k)|=\frac{|\widehat d_m(k)|}{2^m},
\]
故
\[
\Sigma_m^2
=
\frac{1}{(M-1)2^{2m}}\sum_{k\ne 0}|\widehat d_m(k)|^2.
\]
而结论 \ref{con:xi-time-part9l-fold-fourier-energy-arithmetic-lower-bound} 已给出精确能量下界
\[
\sum_{k\ne 0}|\widehat d_m(k)|^2\ge s(M-s),
\]
并且取等当且仅当 \(d_m\) 呈 \(q/q+1\) 两值配置，其中恰有 \(s\) 个剩余类取 \(q+1\)。代回上式即得所述不等式与取等刻画。证毕。
\end{proof}

\leanverified{paper\_xi\_time\_part9m3\_residual\_coherence\_differential\_calculus}
\begin{theorem}[残余相干曲线的微分结构]\label{thm:xi-time-part9m3-residual-coherence-differential-calculus}
\leanverified{paper\_xi\_time\_part9m3\_residual\_coherence\_differential\_calculus}
定义残余相干曲线
\[
\mathscr R_m(T):=
\frac{1}{2^m}\sum_{x\in X_m}(d_m(x)-T)_+^2
\qquad
(T\ge 0).
\]
则在几乎处处意义下有
\[
\mathscr R_m'(T)
=
-2^{1-m}\bigl(2^m-\mathcal C_m^{\mathrm{cont}}(T)\bigr),
\]
\[
\mathscr R_m''(T)=2^{1-m}N_m(T).
\]
从而
\[
\mathscr R_m(T)
=
2^{1-m}\int_T^\infty
\bigl(2^m-\mathcal C_m^{\mathrm{cont}}(s)\bigr)\,\dd s
\]
\[
=
2^{1-m}\int_T^\infty\int_s^\infty N_m(t)\,\dd t\,\dd s.
\]
此外，若记
\[
D_m^{\max}:=\max_{x\in X_m}d_m(x),
\qquad
\mathsf{Col}_m:=\sum_{r\in G_m}p_m(r)^2,
\]
则
\[
\mathscr R_m(0)=2^m\mathsf{Col}_m,
\qquad
\mathscr R_m(T)=0
\quad
\text{对一切 }T\ge D_m^{\max}.
\]
\end{theorem}
\begin{proof}
对任意固定整数 \(d\ge 0\)，函数
\[
f_d(T):=(d-T)_+^2
\]
满足
\[
f_d'(T)=-2(d-T)_+,
\qquad
f_d''(T)=2\,\mathbf 1_{\{T<d\}}
\]
在几乎处处意义下成立。逐项求和并除以 \(2^m\) 得
\[
\mathscr R_m'(T)
=
-2^{1-m}\sum_{x\in X_m}(d_m(x)-T)_+.
\]
另一方面，
\[
\sum_{x\in X_m}(d_m(x)-T)_+
=
\sum_{x\in X_m}\bigl(d_m(x)-\min(d_m(x),T)\bigr)
\]
\[
=
\sum_{x\in X_m}d_m(x)-\mathcal C_m^{\mathrm{cont}}(T)
=
2^m-\mathcal C_m^{\mathrm{cont}}(T),
\]
故第一条微分公式成立。

再由定理 \ref{thm:xi-capacity-tail-histogram-moment-complete-statistic}，
\[
\partial_T^+\mathcal C_m^{\mathrm{cont}}(T)=N_m(T)
\qquad
\text{a.e.},
\]
对上式再求导便得到
\[
\mathscr R_m''(T)=2^{1-m}N_m(T)
\qquad
\text{a.e.}
\]

由于 \(\mathscr R_m(T)=0\) 当 \(T\ge D_m^{\max}\) 时成立，故 \(\mathscr R_m(\infty)=0\)。把第一条公式从 \(T\) 积分到 \(+\infty\) 即得
\[
\mathscr R_m(T)
=
2^{1-m}\int_T^\infty
\bigl(2^m-\mathcal C_m^{\mathrm{cont}}(s)\bigr)\,\dd s.
\]
再由定理 \ref{thm:xi-time-part57b-capacity-moment-exact-mellin-stieltjes} 证明中已经出现的恒等式
\[
2^m-\mathcal C_m^{\mathrm{cont}}(s)
=
\int_s^\infty N_m(t)\,\dd t,
\]
代回即可得到双重积分表达。

最后，
\[
\mathscr R_m(0)
=
\frac{1}{2^m}\sum_{x\in X_m}d_m(x)^2
=
2^m\mathsf{Col}_m,
\]
而 \(T\ge D_m^{\max}\) 时每一项 \((d_m(x)-T)_+\) 都为零，因此 \(\mathscr R_m(T)=0\)。证毕。
\end{proof}

\begin{theorem}[容量、尾谱与残余相干的完全层析]\label{thm:xi-time-part9m3-capacity-residual-coherence-complete-tomography}
固定 \(m\)。记直方图
\[
h_m(n):=\#\{x\in X_m:\ d_m(x)=n\}
\qquad
(n\in\ZZ_{\ge 1}).
\]
则下列四类对象彼此一一决定：
\[
\{d_m(x):x\in X_m\},
\qquad
N_m,
\qquad
\mathcal C_m^{\mathrm{cont}},
\qquad
\mathscr R_m.
\]
更具体地，
\[
h_m(n)=N_m(n)-N_m(n+1),
\]
\[
N_m(T)=\partial_T^+\mathcal C_m^{\mathrm{cont}}(T)
\qquad
\text{a.e.},
\]
\[
N_m(T)=2^{m-1}\mathscr R_m''(T)
\qquad
\text{a.e.}
\]
因此，仅由 \(\mathcal C_m^{\mathrm{cont}}\) 或 \(\mathscr R_m\) 中任一条曲线，便可在固定分辨率 \(m\) 上唯一恢复全部纤维多重度多重集。进一步，对任意 \(q>0\)，都有
\[
S_q(m)=q\int_0^\infty t^{q-1}N_m(t)\,\dd t,
\]
故全部矩谱亦随之唯一恢复。
\leanverified{paper\_xi\_time\_part9m3\_capacity\_residual\_coherence\_complete\_tomography}
\end{theorem}
\begin{proof}
定理 \ref{thm:xi-capacity-tail-histogram-moment-complete-statistic} 已经给出
\[
\mathcal C_m^{\mathrm{cont}},
\qquad
N_m,
\qquad
h_m
\]
三者的完全互相恢复，并且同时给出
\[
S_q(m)=q\int_0^\infty t^{q-1}N_m(t)\,\dd t
\qquad(q>0).
\]
另一方面，定理 \ref{thm:xi-time-part9m3-residual-coherence-differential-calculus} 给出
\[
N_m(T)=2^{m-1}\mathscr R_m''(T)
\qquad
\text{a.e.}
\]
由于 \(N_m\) 是右连续、整数值且紧支撑的阶梯函数，其几乎处处取值已足以唯一确定整个函数；因此 \(\mathscr R_m\) 也唯一决定 \(N_m\)。

由
\[
h_m(n)=N_m(n)-N_m(n+1)
\]
可恢复直方图，而直方图与纤维多重度多重集显然互相等价。综上，四类对象彼此一一决定；再结合上式的 Mellin 表达，全部矩谱亦被唯一恢复。证毕。
\end{proof}

\begin{corollary}[最大纤维指数的精确经典化速率解释]\label{cor:xi-time-part9m3-exact-classicalization-maxfiber-rate}
\leanverified{paper\_xi\_time\_part9m3\_exact\_classicalization\_maxfiber\_rate}
将定理 \ref{thm:xi-full-microstate-exact-inversion-bitrate-threshold} 中的全微态精确反演预算解释为精确经典化成本。定义
\[
B_m^{\mathrm{cl}}
:=
\min\Bigl\{
B\in\ZZ_{\ge 0}:\ 2^B\ge d_m(x)\ \text{对一切 }x\in X_m
\Bigr\},
\]
并记
\[
D_m^{\max}:=\max_{x\in X_m}d_m(x).
\]
则
\[
B_m^{\mathrm{cl}}=\left\lceil \log_2 D_m^{\max}\right\rceil.
\]
若最大纤维指数
\[
\alpha_*:=\lim_{m\to\infty}\frac1m\log D_m^{\max}
\]
存在，则
\[
\lim_{m\to\infty}\frac{B_m^{\mathrm{cl}}}{m}
=
\frac{\alpha_*}{\log 2}.
\]
\end{corollary}
\begin{proof}
由定义，
\[
B_m^{\mathrm{cl}}
=
\min\{B\in\ZZ_{\ge 0}:2^B\ge D_m^{\max}\}
=
\left\lceil \log_2 D_m^{\max}\right\rceil.
\]
若 \(\alpha_*\) 存在，则
\[
\frac1m\log_2 D_m^{\max}
=
\frac{1}{m\log 2}\log D_m^{\max}
\longrightarrow
\frac{\alpha_*}{\log 2}.
\]
而上取整仅引入 \(O(1)\) 误差，故结论成立。证毕。
\end{proof}

\begin{theorem}[可见度与经典化成本的互补下界]\label{thm:xi-time-part9m3-visibility-classicalization-complementarity}
\leanverified{paper\_xi\_time\_part9m3\_visibility\_classicalization\_complementarity}
记 Rényi--\(2\) 有效宽度
\[
W_m^{(2)}:=\frac{1}{\mathsf{Col}_m}.
\]
则精确经典化成本满足
\[
B_m^{\mathrm{cl}}\ge m-\log_2 W_m^{(2)}.
\]
等价地，若用定理 \ref{thm:xi-time-part9m3-mean-square-visibility-arithmetic-floor} 中的均方可见度 \(\Sigma_m^2\) 表示，则
\[
B_m^{\mathrm{cl}}
\ge
m-\log_2 M+\log_2\bigl(1+(M-1)\Sigma_m^2\bigr).
\]
\end{theorem}
\begin{proof}
令
\[
p_{\max}:=\max_{r\in G_m}p_m(r).
\]
则
\[
\mathsf{Col}_m=\sum_{r\in G_m}p_m(r)^2
\le
p_{\max}\sum_{r\in G_m}p_m(r)
=
p_{\max}.
\]
因此
\[
D_m^{\max}=2^m p_{\max}\ge 2^m\mathsf{Col}_m=\frac{2^m}{W_m^{(2)}}.
\]
应用推论 \ref{cor:xi-time-part9m3-exact-classicalization-maxfiber-rate} 即得
\[
B_m^{\mathrm{cl}}
=
\left\lceil \log_2 D_m^{\max}\right\rceil
\ge
\log_2 D_m^{\max}
\ge
m-\log_2 W_m^{(2)}.
\]

另一方面，由 \(\mathcal A_m(0)=1\) 与 Parseval 恒等式，
\[
\sum_{k\in G_m}|\mathcal A_m(k)|^2
=
M\mathsf{Col}_m,
\]
故
\[
\Sigma_m^2
=
\frac{M\mathsf{Col}_m-1}{M-1}.
\]
解得
\[
\mathsf{Col}_m=\frac{1+(M-1)\Sigma_m^2}{M},
\]
代回第一式即得第二种写法。证毕。
\end{proof}

\begin{theorem}[容量亏损与残余相干的 Mellin 表示]\label{thm:xi-time-part9m3-capacity-defect-residual-coherence-mellin}
对任意实数 \(q>1\)，有
\[
S_q(m)
=
q(q-1)\int_0^\infty
T^{q-2}\bigl(2^m-\mathcal C_m^{\mathrm{cont}}(T)\bigr)\,\dd T.
\]
若进一步 \(q>2\)，则还有
\[
S_q(m)
=
q(q-1)(q-2)\,2^{m-1}
\int_0^\infty T^{q-3}\mathscr R_m(T)\,\dd T.
\]
\end{theorem}
\begin{proof}
第一条公式正是定理 \ref{thm:xi-time-part57b-capacity-moment-exact-mellin-stieltjes} 的第二条积分表达。

再由定理 \ref{thm:xi-time-part9m3-residual-coherence-differential-calculus}，
\[
2^m-\mathcal C_m^{\mathrm{cont}}(T)
=
-2^{m-1}\mathscr R_m'(T)
\qquad
\text{a.e.}
\]
故对 \(q>2\)，
\[
\begin{aligned}
S_q(m)
&=
q(q-1)\int_0^\infty
T^{q-2}\bigl(2^m-\mathcal C_m^{\mathrm{cont}}(T)\bigr)\,\dd T\\
&=
-q(q-1)\,2^{m-1}\int_0^\infty T^{q-2}\mathscr R_m'(T)\,\dd T.
\end{aligned}
\]
对最后一个积分分部积分。由于 \(\mathscr R_m\) 在 \([D_m^{\max},\infty)\) 上恒为零，且 \(q>2\) 时有
\[
T^{q-2}\mathscr R_m(T)\longrightarrow 0
\qquad
(T\downarrow 0),
\]
边界项全部消失，于是得到
\[
S_q(m)
=
q(q-1)(q-2)\,2^{m-1}
\int_0^\infty T^{q-3}\mathscr R_m(T)\,\dd T.
\]
证毕。
\end{proof}

\begin{proposition}[解析读出腔室内的折叠谱刚性]\label{prop:xi-time-part9m3-analytic-chamber-rigidity-sampling-readout}
\leanverified{paper\_xi\_time\_part9m3\_analytic\_chamber\_rigidity\_sampling\_readout}
固定 \(m\)，设 \(I\subset\RR\) 为区间。对每个微态 \(\omega\in\Omega_m\) 与每个标签 \(x\in X_m\)，设
\[
s_{\omega,x}:I\to\RR
\]
为实解析函数，并假定对每个 \(\theta\in I\) 都存在唯一最大者
\[
\Fold_{m,\theta}(\omega):=
\arg\max_{x\in X_m}s_{\omega,x}(\theta).
\]
定义对应纤维计数
\[
d_{m,\theta}(x):=
\#\{\omega\in\Omega_m:\ \Fold_{m,\theta}(\omega)=x\}.
\]
则存在离散集 \(E_m\subset I\)，使得对 \(I\setminus E_m\) 的每个连通分支 \(J\)，函数
\[
\theta\longmapsto d_{m,\theta}(x)
\]
对每个 \(x\in X_m\) 都在 \(J\) 上常值。因而
\[
D_{m,\theta}^{\max}:=\max_{x\in X_m}d_{m,\theta}(x),
\]
\[
\mathcal C_{m,\theta}^{\mathrm{cont}}(T):=
\sum_{x\in X_m}\min\bigl(d_{m,\theta}(x),T\bigr),
\]
\[
\mathscr R_{m,\theta}(T):=
\frac1{2^m}\sum_{x\in X_m}(d_{m,\theta}(x)-T)_+^2,
\]
\[
S_q(m;\theta):=\sum_{x\in X_m}d_{m,\theta}(x)^q
\qquad
(q>0)
\]
都在 \(J\) 上保持常值。

进一步，若参数 \(\theta\) 取自真实输入 \(40\) 状态核输出位势的实对数坐标，则注 \ref{rem:real-input-40-output-potential-degeneracy-phase} 给出统一窗口
\[
|\theta|<\theta_0,
\qquad
\theta_0=0.0533477827\ldots,
\]
其中不会跨越正实退化点。因而任何可归约为上述解析 \(\arg\max\) 读出的实参数变形，只有在击中离散交叉集 \(E_m\) 时才可能改变折叠谱与由其诱导的容量、残余相干及矩谱。
\end{proposition}
\begin{proof}
对每个三元组 \((\omega,x,y)\)（其中 \(x\neq y\)），考虑解析差值函数
\[
f_{\omega,x,y}(\theta):=s_{\omega,x}(\theta)-s_{\omega,y}(\theta).
\]
若 \(f_{\omega,x,y}\not\equiv 0\)，则其零点在 \(I\) 内离散。定义
\[
E_m:=
\bigcup_{\omega\in\Omega_m}
\bigcup_{\substack{x,y\in X_m\\ x\neq y\\ f_{\omega,x,y}\not\equiv 0}}
\{\theta\in I:\ f_{\omega,x,y}(\theta)=0\}.
\]
由于 \(\Omega_m\) 与 \(X_m\) 都有限，\(E_m\) 是有限个离散集之并，故仍为离散集。

取 \(I\setminus E_m\) 的任一连通分支 \(J\)。若对某个固定 \(\omega\) 而言，\(\Fold_{m,\theta}(\omega)\) 在 \(J\) 上发生变化，则存在 \(\theta_1<\theta_2\) 属于 \(J\) 与两个不同标签 \(x\neq y\)，使得
\[
\Fold_{m,\theta_1}(\omega)=x,
\qquad
\Fold_{m,\theta_2}(\omega)=y.
\]
由连续性，介值定理迫使某个内点 \(\theta_\ast\in J\) 上有
\[
s_{\omega,x}(\theta_\ast)=s_{\omega,y}(\theta_\ast),
\]
即 \(f_{\omega,x,y}(\theta_\ast)=0\)。若该差值函数不恒为零，则 \(\theta_\ast\in E_m\)，与 \(J\cap E_m=\varnothing\) 矛盾；若该差值函数恒为零，则 \(x\) 与 \(y\) 在整个 \(I\) 上始终并列，从而违背“唯一最大者”的假设。故对每个 \(\omega\)，标签 \(\Fold_{m,\theta}(\omega)\) 在 \(J\) 上恒定。

于是对每个 \(x\in X_m\)，计数 \(d_{m,\theta}(x)\) 在 \(J\) 上恒定。其余各量都只是 \(d_{m,\theta}(x)\) 的确定性函数，故亦在 \(J\) 上恒定。最后一段只是在真实输入 \(40\) 输出位势的实参数坐标中调用注 \ref{rem:real-input-40-output-potential-degeneracy-phase} 所给出的无退化窗口。证毕。
\end{proof}

\endinput


\paragraph{属 \(49\) 的 \texorpdfstring{$S_4$}{S4}--Hurwitz 商塔上的商曲线亏格、总分歧预算、Prym Weil 型、局部 \texorpdfstring{$\zeta$}{zeta} 分解与 Hodge 类传递}
沿用定理 \ref{thm:killo-s4-even-subgroup-free-action-unramified-a4-lift}、\ref{thm:killo-s4-burnside-kani-rosen-prym-square}、\ref{thm:killo-s4-genus49-jacobian-computable-isogeny} 与命题 \ref{prop:killo-two-prym-polarization-types} 的记号，记
\[
X:=\mathcal X,\qquad
H:=\mathcal H=X/A_4,\qquad
C_6:=\mathcal C_6=X/V_4,
\]
\[
Z:=X/D_8,\qquad
C_F:=\mathcal C_F=X/S_3,\qquad
Y:=X/C_4,
\]
并置
\[
A:=\operatorname{Prym}(C_6/H),\qquad
P:=\operatorname{Prym}(Y/Z).
\]

在这一记号下，商塔各层亏格、总分歧预算、分支纤维的固定点几何、Chevalley--Weil 重数的可逆恢复与两类 Prym 极化缺陷的素数分裂都可进一步完全显式化。

\begin{conclusion}[属 \(49\) 的 \texorpdfstring{$S_4$}{S4}--Hurwitz 曲线上四类非平凡循环商的亏格光谱]\label{con:xi-time-part9n-s4-cyclic-quotient-genus-spectrum}
设
\[
V:=H^0(X,\Omega_X^1)
\cong
5\cdot \mathrm{sgn}\ \oplus\ 4\cdot V_2\ \oplus\ 3\cdot V_3\ \oplus\ 9\cdot V_3',
\]
其中平凡表示不出现。若 \(\tau,\delta,\sigma_3,\sigma_4\in S_4\) 分别表示换位、双换位、\(3\)-循环与 \(4\)-循环，则
\[
g(X/\langle \tau\rangle)=19,\qquad
g(X/\langle \delta\rangle)=25,\qquad
g(X/\langle \sigma_3\rangle)=17,\qquad
g(X/\langle \sigma_4\rangle)=13.
\]
\end{conclusion}
\begin{proof}
对任意子群 \(U\le S_4\)，恒有
\[
g(X/U)=\dim V^U.
\]
因此只需计算 \(V\) 在四类非平凡循环子群下的不变维数。取 \(S_4\) 的角色值
\[
\begin{array}{c|cccc}
 & \tau & \delta & \sigma_3 & \sigma_4\\
\hline
\mathrm{sgn} & -1 & 1 & 1 & -1\\
V_2 & 0 & 2 & -1 & 0\\
V_3 & 1 & -1 & 0 & -1\\
V_3' & -1 & -1 & 0 & 1
\end{array}
\]
并用角色平均公式
\[
\dim W^{\langle \tau\rangle}=\frac{\chi_W(1)+\chi_W(\tau)}{2},\qquad
\dim W^{\langle \delta\rangle}=\frac{\chi_W(1)+\chi_W(\delta)}{2},
\]
\[
\dim W^{\langle \sigma_3\rangle}=\frac{\chi_W(1)+2\chi_W(\sigma_3)}{3},\qquad
\dim W^{\langle \sigma_4\rangle}=\frac{\chi_W(1)+\chi_W(\sigma_4)+\chi_W(\delta)+\chi_W(\sigma_4^{-1})}{4},
\]
可得
\[
\dim(\mathrm{sgn}^{\langle\tau\rangle},V_2^{\langle\tau\rangle},V_3^{\langle\tau\rangle},(V_3')^{\langle\tau\rangle})
=(0,1,2,1),
\]
\[
\dim(\mathrm{sgn}^{\langle\delta\rangle},V_2^{\langle\delta\rangle},V_3^{\langle\delta\rangle},(V_3')^{\langle\delta\rangle})
=(1,2,1,1),
\]
\[
\dim(\mathrm{sgn}^{\langle\sigma_3\rangle},V_2^{\langle\sigma_3\rangle},V_3^{\langle\sigma_3\rangle},(V_3')^{\langle\sigma_3\rangle})
=(1,0,1,1),
\]
\[
\dim(\mathrm{sgn}^{\langle\sigma_4\rangle},V_2^{\langle\sigma_4\rangle},V_3^{\langle\sigma_4\rangle},(V_3')^{\langle\sigma_4\rangle})
=(0,1,0,1).
\]
于是分别得到
\[
g(X/\langle \tau\rangle)=5\cdot0+4\cdot1+3\cdot2+9\cdot1=19,
\]
\[
g(X/\langle \delta\rangle)=5\cdot1+4\cdot2+3\cdot1+9\cdot1=25,
\]
\[
g(X/\langle \sigma_3\rangle)=5\cdot1+4\cdot0+3\cdot1+9\cdot1=17,
\]
\[
g(X/\langle \sigma_4\rangle)=5\cdot0+4\cdot1+3\cdot0+9\cdot1=13.
\]
证毕。
\end{proof}

\begin{conclusion}[分支纤维的六对换位固定点分解与其余三类循环元的自由作用]\label{con:xi-time-part9n-s4-branch-fiber-transposition-fixedpairs}
设
\[
f:X\longrightarrow \PP^1_u
\]
为属 \(49\) 的 \(S_4\)-正则闭包覆盖。由结论 \ref{con:fold-gauge-anomaly-s4-hurwitz-tower-genus}，其分支集合为
\[
\{u=0,1\}\cup\{P_{10}(u)=0\},
\]
共 \(12\) 个分支值，且每个分支点的惯性群都由换位生成。于是：
\begin{enumerate}
  \item 对任意分支值 \(u_0\)，纤维 \(f^{-1}(u_0)\) 恰有 \(12\) 个点；
  \item 固定任一换位 \(\tau\in S_4\)。则在每个分支纤维上，\(\tau\) 恰固定 \(2\) 个点；因此 \(\tau\) 在整条曲线上恰固定 \(24\) 个点；
  \item 每个分支纤维因而可分裂为六对两两不交的换位固定点，每一对对应于 \(S_4\) 的一条换位；
  \item 任一双换位、\(3\)-循环或 \(4\)-循环在 \(X\) 上都无固定点。
\end{enumerate}
\end{conclusion}
\begin{proof}
取任一分支值 \(u_0\)，其惯性群记为
\[
I_{u_0}=\langle \tau_0\rangle\simeq C_2.
\]
由于 \(f\) 为正则 \(S_4\)-覆盖，有标准陪集识别
\[
f^{-1}(u_0)\cong S_4/I_{u_0},
\]
故
\[
\#f^{-1}(u_0)=|S_4:I_{u_0}|=24/2=12.
\]

现固定一条换位 \(\tau\)。对任意分支纤维，由正则覆盖的陪集公式，
\[
\#\bigl(\operatorname{Fix}_X(\tau)\cap f^{-1}(u_0)\bigr)
=
\frac{\#\{x\in S_4:\ x^{-1}\tau x\in I_{u_0}\}}{|I_{u_0}|}.
\]
而 \(I_{u_0}\) 的唯一非平凡元即 \(\tau_0\)，故上式条件等价于
\[
x^{-1}\tau x=\tau_0.
\]
满足此式的 \(x\) 构成一个 \(C_{S_4}(\tau)\)-左陪集；又 \(|C_{S_4}(\tau)|=4\)，从而
\[
\#\bigl(\operatorname{Fix}_X(\tau)\cap f^{-1}(u_0)\bigr)=4/2=2.
\]
对全部 \(12\) 个分支值求和即得
\[
\#\operatorname{Fix}_X(\tau)=12\cdot 2=24.
\]
另一方面，每个分支纤维中任一点的稳定子都等于某个换位子群 \(xI_{u_0}x^{-1}\)，故该点只可能被唯一的非平凡换位固定。由 \(6\) 条换位各贡献 \(2\) 个固定点，整个分支纤维遂分解成六对彼此不交的换位固定点。

最后，对双换位 \(\delta\)、\(3\)-循环 \(\sigma_3\) 与 \(4\)-循环 \(\sigma_4\)，由结论 \ref{con:xi-time-part9n-s4-cyclic-quotient-genus-spectrum} 及 Riemann--Hurwitz 公式分别有
\[
96=2\bigl(2g(X/\langle\delta\rangle)-2\bigr)+R_\delta
=2(2\cdot 25-2)+R_\delta,
\]
\[
96=3\bigl(2g(X/\langle\sigma_3\rangle)-2\bigr)+R_{\sigma_3}
=3(2\cdot 17-2)+R_{\sigma_3},
\]
\[
96=4\bigl(2g(X/\langle\sigma_4\rangle)-2\bigr)+R_{\sigma_4}
=4(2\cdot 13-2)+R_{\sigma_4},
\]
故 \(R_\delta=R_{\sigma_3}=R_{\sigma_4}=0\)。因此这三类循环元都自由作用于 \(X\)。证毕。
\end{proof}

\begin{conclusion}[五个循环商亏格对 Chevalley--Weil 重数向量的线性刚性]\label{con:xi-time-part9n-s4-quotient-genera-recover-multiplicities}
设
\[
H^0(X,\Omega_X^1)\cong
m_{\mathbf 1}\cdot \mathbf 1\oplus
m_{\mathrm{sgn}}\cdot \mathrm{sgn}\oplus
m_2\cdot V_2\oplus
m_3\cdot V_3\oplus
m_3'\cdot V_3'.
\]
则仅由
\[
g(X/S_4)=0,\qquad
g(X/\langle \tau\rangle)=19,\qquad
g(X/\langle \delta\rangle)=25,
\]
\[
g(X/\langle \sigma_3\rangle)=17,\qquad
g(X/\langle \sigma_4\rangle)=13
\]
这五个 genus 数据，便可唯一恢复
\[
(m_{\mathbf 1},m_{\mathrm{sgn}},m_2,m_3,m_3')=(0,5,4,3,9).
\]
\end{conclusion}
\begin{proof}
由 \(g(X/S_4)=\dim H^0(X,\Omega_X^1)^{S_4}\) 立即得到
\[
m_{\mathbf 1}=0.
\]
再由结论 \ref{con:xi-time-part9n-s4-cyclic-quotient-genus-spectrum} 的四类循环商 genus 与相应不变维数公式，得到线性系统
\[
m_2+2m_3+m_3'=19,
\]
\[
m_{\mathrm{sgn}}+2m_2+m_3+m_3'=25,
\]
\[
m_{\mathrm{sgn}}+m_3+m_3'=17,
\]
\[
m_2+m_3'=13.
\]
由第一式与第四式相减，得
\[
2m_3=6,\qquad m_3=3.
\]
代回第三式得
\[
m_{\mathrm{sgn}}+m_3'=14.
\]
再将第四式代入第二式，可得
\[
m_{\mathrm{sgn}}+m_2=12.
\]
与
\[
m_{\mathrm{sgn}}+m_3'=14,\qquad m_2+m_3'=13
\]
联立，推出
\[
m_2=4,\qquad m_3'=9,\qquad m_{\mathrm{sgn}}=5.
\]
于是得到唯一解
\[
(m_{\mathbf 1},m_{\mathrm{sgn}},m_2,m_3,m_3')=(0,5,4,3,9).
\]
证毕。
\end{proof}

\begin{conclusion}[八个商曲线的属谱完全表]\label{con:xi-time-part9n-s4-all-subgroup-genera}
设
\[
V:=H^0(X,\Omega_X^1)
\cong 5\cdot \mathrm{sgn}\oplus 4\cdot V_2\oplus 3\cdot V_3\oplus 9\cdot V_3'
\]
为属 \(49\) 的 \texorpdfstring{$S_4$}{S4}--Hurwitz 曲线 \(X\) 的微分表示。记
\[
H:=X/A_4,\qquad
C_6:=X/V_4,\qquad
Z:=X/D_8,\qquad
C_F:=X/S_3,\qquad
Y:=X/C_4,
\]
并记
\[
W_3:=X/C_3,\qquad W_2:=X/C_2.
\]
其中 \(C_2=\langle\tau\rangle\)、\(C_3=\langle\sigma_3\rangle\)、\(C_4=\langle\sigma_4\rangle\) 分别表示由换位、\(3\)-循环与 \(4\)-循环生成的循环子群。则所有这些商曲线的亏格都由 \(S_4\)-微分表示唯一决定，并满足
\[
g(X/S_4)=0,\qquad
g(H)=g(X/A_4)=5,\qquad
g(C_6)=g(X/V_4)=13,\qquad
g(Z)=g(X/D_8)=4,
\]
\[
g(C_F)=g(X/S_3)=3,\qquad
g(Y)=g(X/C_4)=13,\qquad
g(W_3)=g(X/C_3)=17,\qquad
g(W_2)=g(X/C_2)=19.
\]
特别地，\(W_3\) 与 \(W_2\) 的亏格无需额外 Hurwitz 计算，已由同一组 Chevalley--Weil 重数完全锁定。
\end{conclusion}
\begin{proof}
对任意子群 \(K\le S_4\)，商曲线亏格满足
\[
g(X/K)=\dim V^K.
\]
记
\[
\rho_K:=\operatorname{Ind}_{K}^{S_4}\mathbf 1.
\]
由 Frobenius 互反，
\[
\dim(\pi^K)=\langle \pi,\rho_K\rangle_{S_4}
\]
对每个不可约表示 \(\pi\) 都成立。另一方面，定理 \ref{thm:killo-s4-burnside-kani-rosen-prym-square} 给出
\[
\rho_{A_4}=\mathbf 1\oplus \mathrm{sgn},\qquad
\rho_{V_4}=\mathbf 1\oplus \mathrm{sgn}\oplus 2V_2,\qquad
\rho_{D_8}=\mathbf 1\oplus V_2,
\]
\[
\rho_{S_3}=\mathbf 1\oplus V_3,\qquad
\rho_{C_4}=\mathbf 1\oplus V_2\oplus V_3',\qquad
\rho_{C_2}=\mathbf 1\oplus V_2\oplus 2V_3\oplus V_3',
\]
\[
\rho_{C_3}=\mathbf 1\oplus \mathrm{sgn}\oplus V_3\oplus V_3',\qquad
\rho_{S_4}=\mathbf 1.
\]
将这些分解与 \(V\) 的 Chevalley--Weil 分解逐项配对，便得到
\[
g(X/A_4)=5,\quad
g(X/V_4)=5+2\cdot4=13,\quad
g(X/D_8)=4,\quad
g(X/S_3)=3,
\]
\[
g(X/C_4)=4+9=13,\quad
g(X/C_3)=5+3+9=17,\quad
g(X/C_2)=4+2\cdot3+9=19.
\]
又因 \(V\) 中不含平凡表示，故
\[
g(X/S_4)=\dim V^{S_4}=0.
\]
证毕。
\end{proof}

\begin{conclusion}[属 \texorpdfstring{$49$}{49} 的 \texorpdfstring{$S_4$}{S4}--Hurwitz 链中间商映射总分歧度的 \texorpdfstring{$24$}{24} 倍数谱]\label{con:xi-time-part9n-s4-total-ramification-spectrum}
对上述各子群 \(K\le S_4\)，记
\[
\pi_K:X\to X/K
\]
为自然商映射，并令
\[
R_K:=\sum_{P\in X}\bigl(e_P(\pi_K)-1\bigr)
\]
为其总分歧度。则有显式公式
\[
R_{S_4}=144,\qquad
R_{A_4}=0,\qquad
R_{V_4}=0,\qquad
R_{D_8}=48,
\]
\[
R_{S_3}=72,\qquad
R_{C_4}=0,\qquad
R_{C_3}=0,\qquad
R_{C_2}=24.
\]
因此全部非零总分歧度恰落在
\[
24\cdot\{1,2,3,6\},
\]
并形成严格链
\[
24<48<72<144.
\]
\end{conclusion}
\begin{proof}
由上一定理与 Chevalley--Weil 分解可知 \(g(X)=49\)，故
\[
2g(X)-2=96.
\]
对任意 \(K\le S_4\)，Riemann--Hurwitz 公式给出
\[
2g(X)-2=|K|\bigl(2g(X/K)-2\bigr)+R_K,
\]
即
\[
R_K=96-|K|\bigl(2g(X/K)-2\bigr).
\]
将结论 \ref{con:xi-time-part9n-s4-all-subgroup-genera} 的亏格值与
\[
|S_4|=24,\quad |A_4|=12,\quad |V_4|=4,\quad |D_8|=8,\quad
|S_3|=6,\quad |C_4|=4,\quad |C_3|=3,\quad |C_2|=2
\]
逐一代入，立即得到
\[
R_{S_4}=144,\quad R_{A_4}=0,\quad R_{V_4}=0,\quad R_{D_8}=48,
\]
\[
R_{S_3}=72,\quad R_{C_4}=0,\quad R_{C_3}=0,\quad R_{C_2}=24.
\]
其中 \(R_{A_4}=R_{V_4}=R_{C_3}=0\) 与换位惯性对子群 \(A_4,V_4,C_3\) 的平凡交相一致，而 \(R_{C_4}=0\) 则对应于 \(C_4\) 中不含任何换位惯性元。证毕。
\end{proof}

\begin{conclusion}[属 \texorpdfstring{$49$}{49} 的 \texorpdfstring{$S_4$}{S4}--Hurwitz Jacobian 的 \texorpdfstring{$13+9+9+9+9$}{13+9+9+9+9} 型块压缩]\label{con:xi-time-part9n-jacobian-c6-compression}
设
\[
J(X)\sim_{\QQ}J(H)\times J(Z)^2\times J(C_F)^3\times P^3,
\qquad
P=\operatorname{Prym}(Y/Z),
\]
并沿上文记号令 \(C_6=X/V_4\)。则有更压缩的同源分解
\[
\boxed{\;
J(X)\sim_{\QQ} J(C_6)\times J(C_F)^3\times P^3.
\;}
\]
等价地，原来的
\[
5+2\cdot 4+3\cdot 3+3\cdot 9=49
\]
维分解可压缩为
\[
13+9+9+9+9=49
\]
的块结构，其中 \(13\) 维块 \(J(C_6)\) 已吸收 \(J(H)\) 与 \(J(Z)^2\) 的全部几何信息。
\end{conclusion}
\begin{proof}
定理 \ref{thm:killo-s4-burnside-kani-rosen-prym-square} 给出
\[
J(C_6)\sim_{\QQ}J(H)\times J(Z)^2,
\]
其几何内容正是无分歧循环三重覆盖 \(C_6\to H\) 的 Prym 分裂
\[
\operatorname{Prym}(C_6/H)\sim_{\QQ}J(Z)^2.
\]
另一方面，定理 \ref{thm:killo-s4-genus49-jacobian-computable-isogeny} 给出
\[
J(X)\sim_{\QQ}J(H)\times J(Z)^2\times J(C_F)^3\times P^3.
\]
将前一式代入后一式，立即得到
\[
J(X)\sim_{\QQ}J(C_6)\times J(C_F)^3\times P^3.
\]
又由
\[
g(C_6)=13,\qquad 3g(C_F)=9,\qquad \dim P^3=27=9+9+9
\]
可知该分解正对应于 \(13+9+9+9+9\) 的块压缩。证毕。
\end{proof}

\begin{theorem}[属 \texorpdfstring{$49$}{49} 的 \texorpdfstring{$S_4$}{S4}--Hurwitz Jacobian 具有显式半单非交换自同态子代数]\label{thm:xi-time-part9n-jacobian-endomorphism-lower-bound}
\leanverified{paper\_xi\_time\_part9n\_jacobian\_endomorphism\_lower\_bound}
沿用结论 \ref{con:xi-time-part9n-jacobian-c6-compression} 的记号，并令
\[
A:=J(H),\qquad
B:=J(Z),\qquad
C:=J(C_F),\qquad
D:=P.
\]
则有理自同态代数满足显式包含
\[
\boxed{\;
\operatorname{End}^0_{\QQ}\!\bigl(J(X)\bigr)
\supseteq
\QQ\times M_2(\QQ)\times M_3(\QQ)\times M_3(\QQ).\;}
\]
特别地，
\[
\boxed{\;
\dim_{\QQ}\operatorname{End}^0_{\QQ}\!\bigl(J(X)\bigr)\ge 1+2^2+3^2+3^2=23.\;}
\]
并且在 \(\QQ\)-同源意义下，\(J(X)\) 既不简单，也不等型。
\end{theorem}
\begin{proof}
由定理 \ref{thm:killo-s4-genus49-jacobian-computable-isogeny}，
\[
J(X)\sim_{\QQ}A\times B^2\times C^3\times D^3.
\]
同源不改变有理自同态代数，故
\[
\operatorname{End}^0_{\QQ}\!\bigl(J(X)\bigr)
\cong
\operatorname{End}^0_{\QQ}\!\bigl(A\times B^2\times C^3\times D^3\bigr).
\]
对任意阿贝尔簇 \(E\) 与整数 \(n\ge 1\)，标量矩阵给出自然嵌入
\[
M_n(\QQ)\hookrightarrow M_n\!\bigl(\operatorname{End}^0_{\QQ}(E)\bigr)
\cong
\operatorname{End}^0_{\QQ}(E^n).
\]
分别取 \(n=1,2,3,3\) 并作用于四个直积块，即得块对角嵌入
\[
\QQ\times M_2(\QQ)\times M_3(\QQ)\times M_3(\QQ)
\hookrightarrow
\operatorname{End}^0_{\QQ}\!\bigl(J(X)\bigr).
\]
于是
\[
\dim_{\QQ}\operatorname{End}^0_{\QQ}\!\bigl(J(X)\bigr)\ge 1+4+9+9=23.
\]

另一方面，四个可见因子的维数分别为
\[
\dim A=5,\qquad \dim B=4,\qquad \dim C=3,\qquad \dim D=9.
\]
不同维数的阿贝尔簇不可能互相同源，故 \(A,B,C,D\) 代表四个互不同源的等型块。因而 \(J(X)\) 在 \(\QQ\)-同源意义下至少分裂为四个互异块，既不简单，也不等型。证毕。
\end{proof}

\begin{theorem}[\texorpdfstring{$S_4$}{S4}--Hurwitz 链的圆维预算完全闭合]\label{thm:xi-time-part9n-s4-hurwitz-circle-dimension-budget}
\leanverified{paper\_xi\_time\_part9n\_s4\_hurwitz\_circle\_dimension\_budget}
沿用结论 \ref{con:xi-time-part9n-s4-all-subgroup-genera} 与
\ref{con:xi-time-part9n-jacobian-c6-compression} 的记号。对任意复光滑射影曲线 \(C\) 与复阿贝尔簇 \(A\)，记
\[
\Lambda(C):=H_1\!\bigl(C(\CC),\ZZ\bigr),
\qquad
\Lambda(A):=H_1\!\bigl(A(\CC),\ZZ\bigr).
\]
则
\[
\cdim\!\bigl(\Lambda(H)\bigr)=10,\qquad
\cdim\!\bigl(\Lambda(Z)\bigr)=8,\qquad
\cdim\!\bigl(\Lambda(C_F)\bigr)=6,\qquad
\cdim\!\bigl(\Lambda(P)\bigr)=18,
\]
\[
\cdim\!\bigl(\Lambda(C_6)\bigr)=26,
\qquad
\cdim\!\bigl(\Lambda(Y)\bigr)=26.
\]
并且整同调圆维预算满足严格闭合公式
\[
\boxed{\;
\cdim\!\bigl(\Lambda(X)\bigr)
=
\cdim\!\bigl(\Lambda(H)\bigr)
+2\,\cdim\!\bigl(\Lambda(Z)\bigr)
+3\,\cdim\!\bigl(\Lambda(C_F)\bigr)
+3\,\cdim\!\bigl(\Lambda(P)\bigr)
=98.\;}
\]
等价地，
\[
\boxed{\;
\cdim\!\bigl(\Lambda(X)\bigr)=98=10+2\cdot 8+3\cdot 6+3\cdot 18.\;}
\]
因此，属 \(49\) 的 Jacobian 在圆维层面不存在任何隐藏自由秩；文中给出的四块分解
\[
J(H),\qquad J(Z)^2,\qquad J(C_F)^3,\qquad P^3
\]
恰好耗尽全部 \(98\) 维整同调预算。
\end{theorem}
\begin{proof}
由定理 \ref{thm:killo-cdim-short-exact-additivity}，对任意有限生成离散交换群 \(G\) 有
\[
\cdim(G)=\operatorname{rk}_{\ZZ}(G)=\dim_{\QQ}(G\otimes_{\ZZ}\QQ),
\]
并且短正合列下圆维可加。

对任意亏格为 \(g(C)\) 的复光滑射影曲线 \(C\)，其一维整同调群是秩 \(2g(C)\) 的自由阿贝尔群，故
\[
\cdim\!\bigl(\Lambda(C)\bigr)=2g(C).
\]
对任意维数为 \(d\) 的复阿贝尔簇 \(A\)，同理有
\[
\cdim\!\bigl(\Lambda(A)\bigr)=2d.
\]
由结论 \ref{con:xi-time-part9n-s4-all-subgroup-genera}，
\[
g(H)=5,\qquad g(Z)=4,\qquad g(C_F)=3,\qquad g(C_6)=13,\qquad g(Y)=13.
\]
再由结论 \ref{con:xi-time-part9n-jacobian-c6-compression} 的记号，\(\dim P=9\)。于是
\[
\cdim\!\bigl(\Lambda(H)\bigr)=10,\qquad
\cdim\!\bigl(\Lambda(Z)\bigr)=8,\qquad
\cdim\!\bigl(\Lambda(C_F)\bigr)=6,
\]
\[
\cdim\!\bigl(\Lambda(P)\bigr)=18,\qquad
\cdim\!\bigl(\Lambda(C_6)\bigr)=26,\qquad
\cdim\!\bigl(\Lambda(Y)\bigr)=26.
\]

另一方面，由定理 \ref{thm:killo-s4-genus49-jacobian-computable-isogeny}，
\[
J(X)\sim_{\QQ}J(H)\times J(Z)^2\times J(C_F)^3\times P^3.
\]
同源在有理一维同调上诱导直和分解
\[
H_1\!\bigl(J(X)(\CC),\QQ\bigr)
\cong
H_1\!\bigl(J(H)(\CC),\QQ\bigr)
\oplus
H_1\!\bigl(J(Z)(\CC),\QQ\bigr)^{\oplus 2}
\oplus
H_1\!\bigl(J(C_F)(\CC),\QQ\bigr)^{\oplus 3}
\oplus
H_1\!\bigl(P(\CC),\QQ\bigr)^{\oplus 3}.
\]
又因 \(H_1(J(C)(\CC),\ZZ)\) 与 \(H_1(C(\CC),\ZZ)\) 具有相同秩，故
\[
\cdim\!\bigl(\Lambda(X)\bigr)
=
\cdim\!\bigl(\Lambda(H)\bigr)
+2\,\cdim\!\bigl(\Lambda(Z)\bigr)
+3\,\cdim\!\bigl(\Lambda(C_F)\bigr)
+3\,\cdim\!\bigl(\Lambda(P)\bigr).
\]
代入上式即得
\[
\cdim\!\bigl(\Lambda(X)\bigr)=10+2\cdot 8+3\cdot 6+3\cdot 18=98.
\]
另一方面，\(g(X)=49\)，故
\[
\cdim\!\bigl(\Lambda(X)\bigr)=2g(X)=98.
\]
两边一致，命题得证。
\end{proof}

\begin{conclusion}[缺失商 Jacobian 的显式因子化]\label{con:xi-time-part9n-c2-c3-quotient-jacobian-complete-splitting}
沿用上述 \(S_4\)-Hurwitz 商塔记号，并记
\[
W_3:=X/C_3,\qquad
W_2:=X/C_2,\qquad
P:=\operatorname{Prym}(Y/Z).
\]
则有显式同源分解
\[
\boxed{\;
J(W_2)\sim_{\QQ}J(Z)\times J(C_F)^2\times P,
\;}
\]
\[
\boxed{\;
J(W_3)\sim_{\QQ}J(H)\times J(C_F)\times P.
\;}
\]
特别地，
\[
g(W_2)=4+2\cdot 3+9=19,
\qquad
g(W_3)=5+3+9=17,
\]
与结论 \ref{con:xi-time-part9n-s4-all-subgroup-genera} 的亏格恢复完全一致。
\end{conclusion}
\begin{proof}
先证 \(C_2\)-商。由定理 \ref{thm:killo-s4-burnside-kani-rosen-prym-square}，
\[
J(X/C_2)\sim_{\QQ}J(C_F)^2\times J(Y).
\]
另一方面，\(Y\to Z\) 是二重覆盖，而 \(P=\operatorname{Prym}(Y/Z)\)。对二重覆盖的标准 Jacobian--Prym 分解给出
\[
J(Y)\sim_{\QQ}J(Z)\times P.
\]
代入即得
\[
J(X/C_2)\sim_{\QQ}J(Z)\times J(C_F)^2\times P.
\]

再证 \(C_3\)-商。由定理 \ref{thm:killo-s4-genus49-jacobian-computable-isogeny}，
\[
J(X)\sim_{\QQ}J(H)\times J(Z)^2\times J(C_F)^3\times P^3,
\]
其中四个可见因子分别对应于 \(\mathrm{sgn},V_2,V_3,V_3'\) 四个非同构有理不可约表示块。对任意子群 \(K\le S_4\)，商曲线的 holomorphic differential 空间满足
\[
H^0(X/K,\Omega_{X/K}^1)\cong H^0(X,\Omega_X^1)^K,
\]
故每个等型块在 \(J(X/K)\) 中出现的次数正是相应不可约表示的 \(K\)-不变维数。

对 \(K=C_3\)，由
\[
\rho_{C_3}=\mathbf 1\oplus \mathrm{sgn}\oplus V_3\oplus V_3'
\]
可知
\[
\bigl(\dim \mathrm{sgn}^{C_3},\dim V_2^{C_3},\dim V_3^{C_3},\dim (V_3')^{C_3}\bigr)=(1,0,1,1).
\]
因此 \(J(X/C_3)\) 恰由一个 \(\mathrm{sgn}\)-块、一个 \(V_3\)-块与一个 \(V_3'\)-块组成，从而
\[
J(X/C_3)\sim_{\QQ}J(H)\times J(C_F)\times P.
\]
最后，两式的维数分别为 \(4+2\cdot 3+9=19\) 与 \(5+3+9=17\)，恰与结论 \ref{con:xi-time-part9n-s4-all-subgroup-genera} 一致。证毕。
\end{proof}

\begin{conclusion}[良素处换位商与三循环商局部 \texorpdfstring{$L$}{L} 因子的完全分裂]\label{con:xi-time-part9n-c2-c3-local-lfactor-complete-splitting}
设 \(v\) 为一处使
\[
H,\ Z,\ C_F,\ Y,\ X/C_2,\ X/C_3
\]
及 Prym 因子 \(P=\operatorname{Prym}(Y/Z)\) 同时良约化的有限素位。记各自的局部 \(L\)-多项式分子为
\[
P_H(v,T),\quad P_Z(v,T),\quad P_{C_F}(v,T),\quad
P_P(v,T),\quad P_{X/C_2}(v,T),\quad P_{X/C_3}(v,T).
\]
则有完全乘积分解
\[
\boxed{\;
P_{X/C_2}(v,T)=P_Z(v,T)\,P_{C_F}(v,T)^2\,P_P(v,T),
\;}
\]
\[
\boxed{\;
P_{X/C_3}(v,T)=P_H(v,T)\,P_{C_F}(v,T)\,P_P(v,T).
\;}
\]
\end{conclusion}
\begin{proof}
由结论 \ref{con:xi-time-part9n-c2-c3-quotient-jacobian-complete-splitting}，
\[
J(X/C_2)\sim_{\QQ}J(Z)\times J(C_F)^2\times P,
\]
\[
J(X/C_3)\sim_{\QQ}J(H)\times J(C_F)\times P.
\]
对共同良约化素位 \(v\)，同源阿贝尔簇的 \(\ell\)-进 Tate 模半单化具有相同的 Frobenius 特征多项式，而直积对应特征多项式相乘。故两式立即推出
\[
P_{X/C_2}(v,T)=P_Z(v,T)\,P_{C_F}(v,T)^2\,P_P(v,T),
\]
\[
P_{X/C_3}(v,T)=P_H(v,T)\,P_{C_F}(v,T)\,P_P(v,T).
\]
证毕。
\end{proof}

\begin{conclusion}[两类 Prym 极化缺陷诱导的最小同源度刚性]\label{con:xi-time-part9n-prym-polarization-minimal-isogeny-degree}
记
\[
A:=\operatorname{Prym}(C_6/H),
\qquad
P:=\operatorname{Prym}(Y/Z),
\]
并令 \(\lambda_A,\lambda_P\) 为命题 \ref{prop:killo-two-prym-polarization-types} 给出的诱导极化。于是
\[
\mathrm{type}(\lambda_A)=(1,1,1,1,3,3,3,3),
\qquad
\mathrm{type}(\lambda_P)=(1,1,1,1,2,2,2,2,2).
\]
则成立：
\begin{enumerate}
  \item 若 \((B,\lambda_B)\) 为八维主极化阿贝尔簇，且存在同源
  \[
  f:B\to A
  \]
  与整数 \(n\ge 1\) 使
  \[
  f^*\lambda_A=n\lambda_B,
  \]
  则
  \[
  \deg f\ge 3^4=81,
  \]
  且等号当且仅当 \(n=3\)。
  \item 若 \((C,\lambda_C)\) 为九维主极化阿贝尔簇，且存在同源
  \[
  g:C\to P
  \]
  与整数 \(m\ge 1\) 使
  \[
  g^*\lambda_P=m\lambda_C,
  \]
  则
  \[
  \deg g\ge 2^4=16,
  \]
  且等号当且仅当 \(m=2\)。
\end{enumerate}
\end{conclusion}
\begin{proof}
由极化型公式
\[
\deg(\lambda)=\prod_{i=1}^{g}d_i^2,
\]
得
\[
\deg(\lambda_A)=3^8,
\qquad
\deg(\lambda_P)=2^{10}.
\]

先看 \(A\) 的情形。由于 \(\dim B=8\) 且 \(\lambda_B\) 为主极化，
\[
\deg(n\lambda_B)=n^{16}.
\]
另一方面，极化对同源的拉回满足
\[
\deg(f^*\lambda_A)=(\deg f)^2\deg(\lambda_A).
\]
由 \(f^*\lambda_A=n\lambda_B\) 得
\[
(\deg f)^2\cdot 3^8=n^{16},
\qquad\text{即}\qquad
(\deg f)^2=\frac{n^{16}}{3^8}.
\]
右端为整数平方，故其 \(3\)-进赋值必须为非负偶数。设 \(v_3(n)=a\)，则
\[
v_3\!\bigl((\deg f)^2\bigr)=16a-8.
\]
因此必有 \(a\ge 1\)，即 \(3\mid n\)。写 \(n=3r\)，则
\[
(\deg f)^2=3^8r^{16},
\qquad
\deg f=3^4r^8\ge 3^4.
\]
等号成立当且仅当 \(r=1\)，亦即 \(n=3\)。

再看 \(P\) 的情形。由于 \(\dim C=9\) 且 \(\lambda_C\) 为主极化，
\[
\deg(m\lambda_C)=m^{18}.
\]
同理，
\[
(\deg g)^2\cdot 2^{10}=m^{18},
\qquad\text{即}\qquad
(\deg g)^2=\frac{m^{18}}{2^{10}}.
\]
设 \(v_2(m)=b\)，则
\[
v_2\!\bigl((\deg g)^2\bigr)=18b-10.
\]
其为非负偶整数，故必有 \(b\ge 1\)，即 \(2\mid m\)。写 \(m=2s\)，则
\[
(\deg g)^2=2^8s^{18},
\qquad
\deg g=2^4s^9\ge 2^4.
\]
等号成立当且仅当 \(s=1\)，亦即 \(m=2\)。证毕。
\end{proof}

在 \(P\) 本身的极化缺陷之外，\(V_3'\)-等型块 \(A_{3'}\sim_{\QQ} P^3\) 的主极化化代价、全部额外有理端同态以及正维子块类型也可由同一分解立即锁定。

\begin{theorem}[Prym 三重等型块的主极化化、条件性矩阵端代数与子块分类]\label{thm:xi-time-part9n-prym-cubed-principalization-endomorphism}
\leanverified{paper\_xi\_time\_part9n\_prym\_cubed\_principalization\_endomorphism}
沿用结论 \ref{con:xi-time-part9n-prym-polarization-minimal-isogeny-degree} 的记号，记
\[
P:=\operatorname{Prym}(Y/Z),
\qquad
A_{3'}\sim_{\QQ} P^3,
\qquad
\mathrm{type}(\lambda_P)=(1^4,2^5).
\]
令 \(\lambda_{P^3}\) 表示 \(P^3\) 上的乘积极化。则成立：
\begin{enumerate}
  \item \(\lambda_{P^3}\) 的极化型为
  \[
  \mathrm{type}(\lambda_{P^3})=(1^{12},2^{15}).
  \]
  \item 若 \((B,\Theta)\) 为主极化阿贝尔簇，且
  \[
  f:P^3\to B
  \]
  为同源并满足
  \[
  f^*\Theta=\lambda_{P^3},
  \]
  则
  \[
  \deg f=2^{15}.
  \]
  特别地，\(P^3\) 的主极化化最小次数恰为 \(2^{15}\)。
  \item 若猜想 \ref{conj:killo-minimal-extra-endomorphism-large-prym} 的第一条成立，即 \(P\) 在 \(\overline{\QQ}\) 上绝对单，且
  \[
  \operatorname{End}_{\overline{\QQ}}(P)=\ZZ,
  \]
  则
  \[
  \operatorname{End}^0_{\QQ}(A_{3'})\cong M_3(\QQ),
  \qquad
  Z\!\bigl(\operatorname{End}^0_{\QQ}(A_{3'})\bigr)=\QQ.
  \]
  并且 \(A_{3'}\) 的任意正维 \(\QQ\)-阿贝尔子簇都同源于
  \[
  P^r,\qquad 1\le r\le 3.
  \]
  特别地，\(A_{3'}\) 的任意 \(9\) 维简单 \(\QQ\)-因子都唯一同源于 \(P\)。
\end{enumerate}
\end{theorem}
\begin{proof}
乘积极化的 elementary divisors 由各因子的 elementary divisors 直接并列给出，故
\[
\mathrm{type}(\lambda_{P^3})
=
(1^4,2^5)^{\times 3}
=(1^{12},2^{15}).
\]
这证明了第 \(1\) 条。

由极化型公式，
\[
\deg(\lambda_{P^3})=(2^{15})^2.
\]
另一方面，若 \(\Theta\) 为主极化且 \(f^*\Theta=\lambda_{P^3}\)，则
\[
\deg(\lambda_{P^3})
=
\deg(f^*\Theta)
=
(\deg f)^2\deg(\Theta)
=
(\deg f)^2.
\]
故
\[
\deg f=2^{15}.
\]
再者，\(\ker(\lambda_{P^3})\) 作为有限辛群含有阶为 \(2^{15}\) 的极大各向同性子群 \(K\)。对商映射
\[
q:P^3\to P^3/K
\]
应用极化下降的标准构造，可得某个主极化 \(\Theta_K\) 满足
\[
q^*\Theta_K=\lambda_{P^3}.
\]
因此 \(P^3\) 的主极化化确可由次数 \(2^{15}\) 的同源实现，从而该次数既是必要下界，也是可达下界，第 \(2\) 条成立。

若猜想 \ref{conj:killo-minimal-extra-endomorphism-large-prym} 的第一条成立，则
\[
\operatorname{End}^0_{\overline{\QQ}}(P)
=
\operatorname{End}_{\overline{\QQ}}(P)\otimes_{\ZZ}\QQ
=
\QQ.
\]
因而
\[
\operatorname{End}^0_{\QQ}(P)=\QQ.
\]
另一方面，由结论 \ref{con:xi-terminal-zm-s4-jacobian-a3p-isog-prym3-cubed}，
\[
A_{3'}\sim_{\QQ} P^3.
\]
同源阿贝尔簇的有理端同态代数在 \(\QQ\)-同源范畴中同构，而任意阿贝尔簇 \(A\) 都满足
\[
\operatorname{End}^0_{\QQ}(A^3)\cong
M_3\!\bigl(\operatorname{End}^0_{\QQ}(A)\bigr).
\]
取 \(A=P\) 即得
\[
\operatorname{End}^0_{\QQ}(A_{3'})
\cong
\operatorname{End}^0_{\QQ}(P^3)
\cong
M_3(\QQ).
\]
其中心显然等于 \(\QQ\)。

再设 \(C\subseteq A_{3'}\) 为任意正维 \(\QQ\)-阿贝尔子簇。经由 \(A_{3'}\sim_{\QQ} P^3\) 的同源识别，\(C\) 在 \(\QQ\)-同源范畴中对应于某个幂等元
\[
e\in \operatorname{End}^0_{\QQ}(P^3)\cong M_3(\QQ)
\]
的像。线性代数表明，\(M_3(\QQ)\) 中任一幂等元都按共轭同价于
\[
\operatorname{diag}(I_r,0)\qquad (0\le r\le 3),
\]
其中 \(r=\operatorname{rank}(e)\) 唯一确定。于是
\[
\operatorname{Im}(e)\sim_{\QQ} P^r.
\]
由于 \(C\) 正维，故 \(1\le r\le 3\)。这就给出了全部正维 \(\QQ\)-阿贝尔子簇的同源类型分类。

最后，若 \(C\) 还是简单且 \(\dim C=9=\dim P\)，则只能有 \(r=1\)，从而
\[
C\sim_{\QQ} P.
\]
故第 \(3\) 条全部成立。证毕。
\end{proof}

\endinput

在 \(V_3'\)-等型三重块 \(P^3\) 的主极化化最小次数与矩阵端代数已经锁定之后，还可把这一个 \(27\)-维可见块在 \(\QQ\)-同源范畴中的子等型与 Rosati 对称端代数完全显式化。

\begin{theorem}[\texorpdfstring{$27$}{27} 维 Prym 等型块的 \texorpdfstring{$\QQ$}{Q}-子等型完全分类]\label{thm:xi-time-part9n1e-prym-cubed-q-subisogeny-rosati-classification}
沿用
\[
J(X)\sim_{\QQ}J(H)\times J(Z)^2\times J(C_F)^3\times P^3,
\qquad
P:=\operatorname{Prym}(Y/Z),
\qquad
\dim P=9
\]
的记号，并进一步假设猜想 \ref{conj:killo-minimal-extra-endomorphism-large-prym} 的第一条成立，即
\[
P\ \text{在 }\overline{\QQ}\text{ 上绝对单，且}\ \operatorname{End}_{\overline{\QQ}}(P)=\ZZ.
\]
则成立：
\begin{enumerate}
  \item \textbf{（有理端代数）}
        \[
        \operatorname{End}_{\QQ}^0(P^3)\cong M_3(\QQ).
        \]
  \item \textbf{（\(\QQ\)-子等型完全分类）}
        \(P^3\) 的一切 \(\QQ\)-阿贝尔子簇，在 \(\QQ\)-同源意义下恰为
        \[
        P^r\qquad (r=0,1,2,3)
        \]
        这四种类型；不存在任何额外的、非幂次型的 \(\QQ\)-子等型。
  \item \textbf{（Rosati 对称端代数）}
        若在 \(P^3\) 上取乘积极化 \(\lambda_{P^3}\)，并记相应 Rosati 反自同态为 \(\dagger\)，则
        \[
        \operatorname{End}_{\QQ}^0(P^3)^{\dagger}
        \cong
        \operatorname{Sym}_3(\QQ).
        \]
        因而其有理极化锥恰由正定对称 \(3\times 3\) 有理矩阵参数化。
\end{enumerate}
\end{theorem}
\begin{proof}
由假设
\[
\operatorname{End}_{\overline{\QQ}}(P)=\ZZ
\]
可知
\[
\operatorname{End}_{\overline{\QQ}}^0(P)=\QQ.
\]
而 \(\operatorname{End}_{\QQ}(P)\) 自然嵌入 \(\operatorname{End}_{\overline{\QQ}}(P)\)，故
\[
\QQ\subseteq \operatorname{End}_{\QQ}^0(P)\subseteq \operatorname{End}_{\overline{\QQ}}^0(P)=\QQ,
\]
从而
\[
\operatorname{End}_{\QQ}^0(P)=\QQ.
\]
对任意阿贝尔簇 \(A\) 与整数 \(m\ge 1\)，标准等型公式给出
\[
\operatorname{End}_{\QQ}^0(A^m)\cong M_m\!\bigl(\operatorname{End}_{\QQ}^0(A)\bigr).
\]
取 \(A=P\)、\(m=3\) 即得
\[
\operatorname{End}_{\QQ}^0(P^3)\cong M_3(\QQ),
\]
证明了第 \(1\) 条。

现证第 \(2\) 条。阿贝尔簇的 \(\QQ\)-同源范畴是半单伪阿贝尔范畴，因此 \(P^3\) 的 \(\QQ\)-阿贝尔子簇的同源类与 \(\operatorname{End}_{\QQ}^0(P^3)\) 中的幂等元一一对应。借助第 \(1\) 条，可把此问题转化为 \(M_3(\QQ)\) 中幂等矩阵的分类。

对任意 \(e\in M_3(\QQ)\) 满足 \(e^2=e\)，其最小多项式整除 \(t(t-1)\)，故无重根，从而 \(e\) 在 \(\QQ\) 上可对角化。于是 \(e\) 必与某个
\[
\operatorname{diag}(I_r,0)\qquad (r=0,1,2,3)
\]
在 \(\operatorname{GL}_3(\QQ)\) 下共轭。对应的像对象在 \(\QQ\)-同源意义下正是 \(P^r\)。反过来，这四个秩值显然都可实现。故 \(P^3\) 的 \(\QQ\)-子等型只有
\[
0,\ P,\ P^2,\ P^3
\]
四类，而不存在任何额外的非幂次型子等型。

最后处理第 \(3\) 条。由乘积极化 \(\lambda_{P^3}=\lambda_P^{\times 3}\) 所诱导的 Rosati 反自同态，在
\[
\operatorname{End}_{\QQ}^0(P^3)\cong M_3\!\bigl(\operatorname{End}_{\QQ}^0(P)\bigr)=M_3(\QQ)
\]
上的作用与“系数域上的 Rosati 对合加矩阵转置”一致。由于 \(\operatorname{End}_{\QQ}^0(P)=\QQ\) 且 \(\QQ\) 上唯一正反演为恒等，故 \(\dagger\) 在 \(M_3(\QQ)\) 上即与转置等价。因此
\[
\operatorname{End}_{\QQ}^0(P^3)^{\dagger}
\cong
\{M\in M_3(\QQ):M^{\mathsf T}=M\}
=
\operatorname{Sym}_3(\QQ).
\]
在此识别下，Rosati 正元正对应于正定对称矩阵，于是有理极化锥恰由正定对称 \(3\times 3\) 有理矩阵参数化。证毕。
\end{proof}

\noindent
因此，在大维 Prym 因子的绝对单纯性与端同态极小化假设下，\(P^3\) 不仅具有矩阵型有理端代数，而且其 \(\QQ\)-子等型、Rosati 对称端代数与极化锥都被 \(M_3(\QQ)\) 的线性代数完全刚性化。

\endinput

\paragraph{规范化纤维积、Hodge 张量分层与局部 Euler 因子的平方/立方刚性}
在 \(S_4\)-Hurwitz 商塔的亏格、总分歧、Jacobian 同源分解与 Prym 极化缺陷已经完全确定之后，还可继续把 \(A_4\)、\(D_8\)、\(V_4\) 三层商曲线之间的几何拼接、全纯微分的表示论来源以及良素处局部 Euler 因子的幂次结构进一步压缩为下列四条后果。

\begin{conclusion}[\texorpdfstring{$V_4$}{V4} 商曲线的规范化纤维积刻画]\label{con:xi-time-part9n-s4-c6-normalized-fiberproduct}
沿用 \(S_4\)-Hurwitz 商塔记号，并令
\[
H\longrightarrow \PP^1_u,\qquad Z\longrightarrow \PP^1_u
\]
为由 \(X\to \PP^1_u\) 诱导的自然映射。则有规范化纤维积识别
\[
\boxed{\;
C_6\ \cong\ \widetilde{H\times_{\PP^1_u} Z},
\;}
\]
其中右侧表示纤维积的规范化。

并且两条自然投影满足
\[
\deg\!\bigl(C_6\to H\bigr)=3,\qquad \deg\!\bigl(C_6\to Z\bigr)=2.
\]
其中 \(C_6\to H\) 为无分歧循环三重覆盖，而 \(C_6\to Z\) 的总分歧度恰为
\[
\boxed{\;12.\;}
\]
\end{conclusion}
\begin{proof}
在 \(S_4\) 中有群论恒等式
\[
A_4\cap D_8=V_4,\qquad A_4D_8=S_4.
\]
因此对正则 \(S_4\)-覆盖 \(X\to \PP^1_u\)，其函数域满足
\[
\QQ(H)=\QQ(X)^{A_4},\qquad
\QQ(Z)=\QQ(X)^{D_8},
\]
而在 \(\QQ(X)\) 中的复合域为
\[
\QQ(H)\,\QQ(Z)
=
\QQ(X)^{A_4\cap D_8}
=
\QQ(X)^{V_4}
=
\QQ(C_6).
\]
另一方面，\(\widetilde{H\times_{\PP^1_u} Z}\) 的函数域正是 \(\QQ(H)\) 与 \(\QQ(Z)\) 在 \(\QQ(X)\) 中的复合域，故
\[
\widetilde{H\times_{\PP^1_u} Z}\cong C_6.
\]

两条投影的次数由子群指数直接给出：
\[
[A_4:V_4]=3,\qquad [D_8:V_4]=2.
\]
又由定理 \ref{thm:killo-s4-even-subgroup-free-action-unramified-a4-lift}，
\[
C_6\to H
\]
本身就是无分歧循环三重覆盖。

对另一条二重映射，由结论 \ref{con:xi-time-part9n-s4-all-subgroup-genera}，
\[
g(C_6)=13,\qquad g(Z)=4.
\]
代入 Riemann--Hurwitz 公式
\[
2g(C_6)-2=2\bigl(2g(Z)-2\bigr)+R
\]
即得
\[
24=12+R,
\]
故 \(R=12\)。证毕。
\end{proof}

\begin{conclusion}[全纯微分空间的 \texorpdfstring{$S_4$}{S4} 张量型 Hodge 分层]\label{con:xi-time-part9n-s4-hodge-tensor-factorization}
记
\[
V:=H^0(X,\Omega_X^1).
\]
则 \(V\) 作为 \(S_4\)-表示与 Hodge 余切空间同时满足严格张量分解
\[
\boxed{\;
H^0(X,\Omega_X^1)
\cong
\mathrm{sgn}\otimes H^0(H,\Omega_H^1)
\oplus
V_2\otimes H^0(Z,\Omega_Z^1)
\oplus
V_3\otimes H^0(C_F,\Omega_{C_F}^1)
\oplus
V_3'\otimes H^0(P,\Omega_P^1).
\;}
\]
特别地，\(V_3'\)-等型部分全部由九维 Prym 因子 \(P\) 承载，不再混入任何来自曲线商 \(H\)、\(Z\)、\(C_F\) 的贡献。
\end{conclusion}
\begin{proof}
由定理 \ref{thm:killo-s4-genus49-jacobian-computable-isogeny}，
\[
J(X)\sim_{\QQ}J(H)\times J(Z)^2\times J(C_F)^3\times P^3.
\]
对任意阿贝尔簇同源分解取单位元处余切空间，并用 Jacobian 的标准识别
\[
H^0\!\bigl(J(C),\Omega_{J(C)}^1\bigr)\cong H^0(C,\Omega_C^1),
\]
得到向量空间分解
\[
H^0(X,\Omega_X^1)
\cong
H^0(H,\Omega_H^1)
\oplus
H^0(Z,\Omega_Z^1)^{\oplus 2}
\oplus
H^0(C_F,\Omega_{C_F}^1)^{\oplus 3}
\oplus
H^0(P,\Omega_P^1)^{\oplus 3}.
\]
其中
\[
\dim H^0(H,\Omega_H^1)=5,\qquad
\dim H^0(Z,\Omega_Z^1)=4,
\]
\[
\dim H^0(C_F,\Omega_{C_F}^1)=3,\qquad
\dim H^0(P,\Omega_P^1)=9.
\]

另一方面，结论 \ref{con:xi-time-part9n-s4-all-subgroup-genera} 给出的 Chevalley--Weil 分解为
\[
H^0(X,\Omega_X^1)
\cong
5\cdot \mathrm{sgn}\oplus 4\cdot V_2\oplus 3\cdot V_3\oplus 9\cdot V_3',
\]
故四个等型块的总维数分别为
\[
1\cdot 5,\qquad 2\cdot 4,\qquad 3\cdot 3,\qquad 3\cdot 9.
\]
它们恰与
\[
J(H),\qquad J(Z)^2,\qquad J(C_F)^3,\qquad P^3
\]
在余切空间中的维数分块完全一致。由于 \(\mathrm{sgn}\)、\(V_2\)、\(V_3\)、\(V_3'\) 两两不同构，等型分解的重数空间唯一，于是只能分别识别为
\[
H^0(H,\Omega_H^1),\qquad
H^0(Z,\Omega_Z^1),\qquad
H^0(C_F,\Omega_{C_F}^1),\qquad
H^0(P,\Omega_P^1).
\]
这就得到所示张量型分解。最后一条断言只是把 \(V_3'\)-块与 \(P^3\) 的唯一对应关系重新表述。证毕。
\end{proof}

\begin{conclusion}[良素处局部 Hasse--Weil 分子的平方/立方块与点数整除律]\label{con:xi-time-part9n-s4-local-euler-square-cube}
设 \(p\) 为一个使
\[
X,\ H,\ C_6,\ Z,\ C_F
\]
及阿贝尔簇 \(P\) 同时良约化的素数。记各自的局部 Hasse--Weil 分子为
\[
P_{H,p}(T),\quad P_{C_6,p}(T),\quad P_{Z,p}(T),\quad
P_{C_F,p}(T),\quad P_{X,p}(T),\quad P_{P,p}(T).
\]
则有严格因子分解
\[
\boxed{\;
P_{C_6,p}(T)=P_{H,p}(T)\,P_{Z,p}(T)^2,
\;}
\]
\[
\boxed{\;
P_{X,p}(T)=P_{H,p}(T)\,P_{Z,p}(T)^2\,P_{C_F,p}(T)^3\,P_{P,p}(T)^3.
\;}
\]
等价地，
\[
P_{X,p}(T)=P_{H,p}(T)\,P_{Z,p}(T)^2
\bigl(P_{C_F,p}(T)P_{P,p}(T)\bigr)^3,
\]
因此 \(C_6\) 的局部分子必带严格平方块，而 \(X\) 的局部分子必带严格立方块。

若记
\[
t_p(P):=
\operatorname{Tr}\!\Bigl(
\mathrm{Frob}_p\mid H^1_{\text{ét}}(P_{\overline{\FF}_p},\QQ_\ell)
\Bigr),
\]
则进一步有点数公式
\[
\boxed{\;
\#C_6(\FF_p)=\#H(\FF_p)+2\,\#Z(\FF_p)-2(p+1),
\;}
\]
\[
\boxed{\;
\#X(\FF_p)=\#H(\FF_p)+2\,\#Z(\FF_p)+3\,\#C_F(\FF_p)-5(p+1)-3\,t_p(P).
\;}
\]
特别地，
\[
\boxed{\;
\#X(\FF_p)-\#H(\FF_p)-2\,\#Z(\FF_p)-3\,\#C_F(\FF_p)+5(p+1)\equiv 0\pmod 3.
\;}
\]
\end{conclusion}
\begin{proof}
由定理 \ref{thm:killo-s4-burnside-kani-rosen-prym-square}，
\[
J(C_6)\sim_{\QQ}J(H)\times J(Z)^2,
\]
而由定理 \ref{thm:killo-s4-genus49-jacobian-computable-isogeny}，
\[
J(X)\sim_{\QQ}J(H)\times J(Z)^2\times J(C_F)^3\times P^3.
\]
在共同良约化素数 \(p\) 处，同源阿贝尔簇的 \(\ell\)-进 \(H^1\) 具有相同 Frobenius 特征多项式，而直积对应特征多项式相乘，故立即得到两条局部分子分解。

对任意曲线 \(C/\FF_p\)，写
\[
P_{C,p}(T)=1-a_p(C)T+\cdots,\qquad a_p(C)=p+1-\#C(\FF_p).
\]
对阿贝尔簇 \(P\)，同理有
\[
P_{P,p}(T)=1-t_p(P)T+\cdots.
\]
比较
\[
P_{C_6,p}(T)=P_{H,p}(T)\,P_{Z,p}(T)^2
\]
的一次项，得到
\[
a_p(C_6)=a_p(H)+2a_p(Z),
\]
即
\[
\#C_6(\FF_p)=\#H(\FF_p)+2\,\#Z(\FF_p)-2(p+1).
\]

再比较
\[
P_{X,p}(T)=P_{H,p}(T)\,P_{Z,p}(T)^2\,P_{C_F,p}(T)^3\,P_{P,p}(T)^3
\]
的一次项，得到
\[
a_p(X)=a_p(H)+2a_p(Z)+3a_p(C_F)+3t_p(P).
\]
把 \(a_p(C)=p+1-\#C(\FF_p)\) 代回并整理，即得
\[
\#X(\FF_p)=\#H(\FF_p)+2\,\#Z(\FF_p)+3\,\#C_F(\FF_p)-5(p+1)-3\,t_p(P).
\]
最后一条同余式只是把上式移项后按模 \(3\) 约化。证毕。
\end{proof}

\begin{conclusion}[双平方块上的 \texorpdfstring{$\QQ(\zeta_3)$}{Q(zeta3)}-矩阵乘法与 cyclotomic norm-square]\label{con:xi-time-part9n-s4-jz2-cyclotomic-norm-square}
记
\[
K:=\QQ(\zeta_3),\qquad A:=\operatorname{Prym}(C_6/H).
\]
则存在 \(\QQ\)-代数嵌入
\[
\boxed{\;
K\hookrightarrow \operatorname{End}^0\!\bigl(J(Z)^2\bigr)
\cong
M_2\!\bigl(\operatorname{End}^0(J(Z))\bigr).
\;}
\]
因此亏格 \(4\) 因子 \(J(Z)\) 的平方块带有必然的 cyclotomic matrix multiplication。

进一步地，对任意 \(Z\) 的良素数 \(p\)，存在次数 \(8\) 的多项式
\[
Q_{Z,p}(T)\in K[T]
\]
使得
\[
\boxed{\;
P_{Z,p}(T)^2=
\operatorname{Norm}_{K/\QQ}\!\bigl(Q_{Z,p}(T)\bigr).
\;}
\]
换言之，\(P_{Z,p}(T)^2\) 不是任意平方，而是一个 cyclotomic norm-square。
\end{conclusion}
\begin{proof}
由命题 \ref{prop:killo-two-prym-polarization-types}，
\[
K=\QQ(\zeta_3)\subset \operatorname{End}^0(A).
\]
再由定理 \ref{thm:killo-s4-burnside-kani-rosen-prym-square}，
\[
A\sim_{\QQ}J(Z)^2.
\]
沿任一 \(\QQ\)-同源把端同态代数共轭传送，便得到
\[
K\hookrightarrow \operatorname{End}^0\!\bigl(J(Z)^2\bigr).
\]
而
\[
\operatorname{End}^0\!\bigl(J(Z)^2\bigr)\cong M_2\!\bigl(\operatorname{End}^0(J(Z))\bigr)
\]
是标准矩阵代数识别，故第一条结论成立。

现取 \(Z\) 的一处良素数 \(p\) 与 \(\ell\neq p\)。由于 \(K\) 的作用来自代数端同态，它与 Frobenius 作用对易，因此
\[
V_\ell\!\bigl(J(Z)^2\bigr):=
H^1_{\text{ét}}\!\bigl(J(Z)^2_{\overline{\FF}_p},\QQ_\ell\bigr)
\]
可视为一个 \(K\otimes_{\QQ}\QQ_\ell\)-线性表示。其 \(\QQ_\ell\)-维数为 \(16\)，故作为 \(K\)-线性空间的维数为 \(8\)。记
\[
Q_{Z,p}(T)\in K[T]
\]
为 Frobenius 在该 \(K\)-线性表示上的特征多项式，则 \(\deg Q_{Z,p}=8\)。

按定义，\(\operatorname{Norm}_{K/\QQ}(Q_{Z,p}(T))\) 正是 Frobenius 在同一空间上作为 \(\QQ_\ell\)-线性变换的特征多项式。另一方面，
\[
H^1_{\text{ét}}\!\bigl(J(Z)^2_{\overline{\FF}_p},\QQ_\ell\bigr)
\cong
H^1_{\text{ét}}\!\bigl(Z_{\overline{\FF}_p},\QQ_\ell\bigr)^{\oplus 2},
\]
故该 \(\QQ_\ell\)-线性特征多项式恰为 \(P_{Z,p}(T)^2\)。于是
\[
P_{Z,p}(T)^2=
\operatorname{Norm}_{K/\QQ}\!\bigl(Q_{Z,p}(T)\bigr).
\]
证毕。
\end{proof}

\begin{conclusion}[Hurwitz 塔中的素数分离型 Prym 刚性]\label{con:xi-time-part9n-prym-prime-separated-rigidity}
记
\[
A_{(3)}:=\operatorname{Prym}(C_6/H),
\qquad
A_{(2)}:=P=\operatorname{Prym}(Y/Z).
\]
并记 \(\lambda_{A_{(3)}}\)、\(\lambda_{A_{(2)}}\) 为其自然诱导极化。则三次层与二次层的 Prym 极化缺陷严格按素数分离：
\[
\boxed{\;
A_{(3)}\sim_{\QQ}J(Z)^2,\qquad
\mathrm{type}(\lambda_{A_{(3)}})=(1,1,1,1,3,3,3,3),\qquad
\QQ(\zeta_3)\subset \operatorname{End}^0(A_{(3)}).\;}
\]
\[
\boxed{\;
\mathrm{type}(\lambda_{A_{(2)}})=(1,1,1,1,2,2,2,2,2).\;}
\]
特别地，三次分支所产生的全部非主极化缺陷都局限于 \(3\)-初等方向，而二次分支所产生的全部非主极化缺陷都局限于 \(2\)-初等方向；两者在极化型与端同态约束上互不混杂。
\end{conclusion}
\begin{proof}
由定理 \ref{thm:killo-s4-burnside-kani-rosen-prym-square}，
\[
\operatorname{Prym}(C_6/H)\sim_{\QQ}J(Z)^2.
\]
再由命题 \ref{prop:killo-two-prym-polarization-types}，
\[
\mathrm{type}(\lambda_{A_{(3)}})=(1,1,1,1,3,3,3,3),
\qquad
\mathrm{type}(\lambda_{A_{(2)}})=(1,1,1,1,2,2,2,2,2),
\]
并且
\[
\QQ(\zeta_3)\subset \operatorname{End}^0\!\bigl(\operatorname{Prym}(C_6/H)\bigr).
\]
这与结论 \ref{con:xi-time-part9n-s4-jz2-cyclotomic-norm-square} 中由 \(J(Z)^2\) 读出的 cyclotomic matrix multiplication 完全一致。故 \(A_{(3)}\) 的全部非单位初等因子都是 \(3\)，而 \(A_{(2)}\) 的全部非单位初等因子都是 \(2\)，从而两类 Prym 的非主极化缺陷严格分裂为互素的两块。证毕。
\end{proof}

\begin{theorem}[\texorpdfstring{$\operatorname{Prym}(C_6/H)\to J(Z)^2$}{Prym(C6/H) to J(Z)^2} 的极化兼容同源度刚性]\label{thm:xi-time-part9n1a-prym-c6h-jz2-polarized-isogeny-degree}
\leanverified{paper\_xi\_time\_part9n1a\_prym\_c6h\_jz2\_polarized\_isogeny\_degree}
记
\[
A:=\operatorname{Prym}(C_6/H),
\qquad
B:=J(Z)^2.
\]
令 \(\lambda_A\) 为 \(A\) 的自然诱导极化，\(\lambda_B\) 为 \(B\) 上的乘积主极化。若
\[
f:A\to B
\]
是满足
\[
f^*\lambda_B=\lambda_A
\]
的同源，则
\[
\deg f=3^4.
\]
\end{theorem}
\begin{proof}
由定理 \ref{thm:killo-s4-burnside-kani-rosen-prym-square}，
\[
A\sim_{\QQ}B.
\]
再由命题 \ref{prop:killo-two-prym-polarization-types}，
\[
\mathrm{type}(\lambda_A)=(1,1,1,1,3,3,3,3).
\]
因此
\[
\deg(\lambda_A)
=
(1\cdot1\cdot1\cdot1\cdot3\cdot3\cdot3\cdot3)^2
=
3^8,
\qquad
\deg(\lambda_B)=1.
\]
若 \(f^*\lambda_B=\lambda_A\)，则极化拉回的度数公式给出
\[
\deg(\lambda_A)
=
\deg(f^*\lambda_B)
=
(\deg f)^2\deg(\lambda_B)
=
(\deg f)^2.
\]
故
\[
(\deg f)^2=3^8,
\]
从而
\[
\deg f=3^4.
\]
证毕。
\end{proof}

\begin{corollary}[\texorpdfstring{$V_2$}{V2}-等型 Jacobian 块的 Eisenstein 乘法]\label{cor:xi-time-part9n1a-v2-block-eisenstein-multiplication}
记
\[
A_{V_2}:=J(Z)^2.
\]
则有
\[
\QQ(\zeta_3)\subset \operatorname{End}^0\!\bigl(A_{V_2}\bigr)
\subset
\operatorname{End}^0_{\QQ}\!\bigl(J(X)\bigr).
\]
特别地，属 \(49\) 的 \(S_4\)-Hurwitz Jacobian 分解中，对应 \(V_2\)-等型的八维块天然携带一条 Eisenstein 乘法结构，而不再是纯实型块。
\end{corollary}
\begin{proof}
第一重包含由结论 \ref{con:xi-time-part9n-s4-jz2-cyclotomic-norm-square} 直接给出。另一方面，定理 \ref{thm:xi-time-part9n1b-jacobian-projector-idempotent-algebra} 已把 \(J(Z)^2\) 识别为 \(J(X)\) 的一个规范块对角等型分量，故其有理端同态代数自然嵌入
\[
\operatorname{End}^0_{\QQ}\!\bigl(J(X)\bigr).
\]
由于 \(\QQ(\zeta_3)\) 是虚二次域，所述 \(V_2\)-块遂带有 Eisenstein 乘法。证毕。
\end{proof}

\endinput

\paragraph{Prym--Chevalley--Weil 互反、商曲线 genus 反演与投影代数刚性}
在 \(S_4\)-Hurwitz 商塔的亏格恢复、Hodge 张量分层、Burnside--Kani--Rosen 同源关系与 Jacobian 分块已经明确之后，还可继续把其中若干隐含的互反与刚性压缩为下列五条后果。

\begin{theorem}[\texorpdfstring{$S_4$}{S4}-几何表示的无余因子实现]\label{thm:xi-time-part9n1b-prym-chevalley-weil-reciprocity}
记
\[
V:=H^0(X,\Omega_X^1),
\qquad
A_{\mathrm{sgn}}:=J(H),\quad
A_{V_2}:=J(Z),\quad
A_{V_3}:=J(C_F),\quad
A_{V_3'}:=P.
\]
则四个非平凡有理不可约表示
\[
\mathrm{sgn},\qquad V_2,\qquad V_3,\qquad V_3'
\]
分别且唯一地由几何因子
\[
J(H),\qquad J(Z),\qquad J(C_F),\qquad P
\]
承载。更具体地，表示重数与几何维数满足精确互反
\[
\boxed{\;
V\cong
g(H)\cdot \mathrm{sgn}\ \oplus\
g(Z)\cdot V_2\ \oplus\
g(C_F)\cdot V_3\ \oplus\
\dim(P)\cdot V_3'.\;}
\]
并且 \(J(X)\) 的有理同源分解可重写为
\[
\boxed{\;
J(X)\sim_{\QQ}
A_{\mathrm{sgn}}^{\dim(\mathrm{sgn})}\times
A_{V_2}^{\dim(V_2)}\times
A_{V_3}^{\dim(V_3)}\times
A_{V_3'}^{\dim(V_3')}. \;}
\]
因此，\(S_4\)-微分表示的全部非平凡部分已被“商 Jacobian 因子 \(+\) 单个 Prym 因子”完全穷尽，不存在任何额外的隐藏等型几何块。
\end{theorem}
\begin{proof}
由结论 \ref{con:xi-time-part9n-s4-hodge-tensor-factorization}，
\[
V\cong
\mathrm{sgn}\otimes H^0(H,\Omega_H^1)
\oplus
V_2\otimes H^0(Z,\Omega_Z^1)
\oplus
V_3\otimes H^0(C_F,\Omega_{C_F}^1)
\oplus
V_3'\otimes H^0(P,\Omega_P^1).
\]
对四个重数空间分别取维数，并用
\[
\dim H^0(H,\Omega_H^1)=g(H),\quad
\dim H^0(Z,\Omega_Z^1)=g(Z),\quad
\dim H^0(C_F,\Omega_{C_F}^1)=g(C_F),\quad
\dim H^0(P,\Omega_P^1)=\dim(P),
\]
便得到第一式。

另一方面，由定理 \ref{thm:killo-s4-genus49-jacobian-computable-isogeny}，
\[
J(X)\sim_{\QQ}J(H)\times J(Z)^2\times J(C_F)^3\times P^3.
\]
再代入
\[
\dim(\mathrm{sgn})=1,\qquad \dim(V_2)=2,\qquad \dim(V_3)=\dim(V_3')=3,
\]
即得第二式。由于 \(S_4\) 的非平凡有理不可约表示恰止于这四类，而相应几何因子的维数分别为 \(5,4,3,9\)，总维数
\[
1\cdot 5+2\cdot 4+3\cdot 3+3\cdot 9=49=g(X)
\]
正好充满 \(J(X)\) 的全部 \(H^1\) 贡献，故不可能再有第五个隐藏等型块。证毕。
\end{proof}

\begin{corollary}[四个 genus 数据对非平凡 \(S_4\)-Hodge 向量的唯一恢复]\label{cor:xi-time-part9n1b-four-genera-recover-hodge-vector}
记
\[
V\cong
m_{\mathrm{sgn}}\cdot \mathrm{sgn}\oplus
m_{V_2}\cdot V_2\oplus
m_{V_3}\cdot V_3\oplus
m_{V_3'}\cdot V_3'.
\]
则仅由四个几何量
\[
g(X),\qquad g(H),\qquad g(C_6),\qquad g(C_F)
\]
便可唯一恢复重数向量 \((m_{\mathrm{sgn}},m_{V_2},m_{V_3},m_{V_3'})\)，其显式公式为
\[
\boxed{\;
m_{\mathrm{sgn}}=g(H),\qquad
m_{V_2}=\frac{g(C_6)-g(H)}{2},\qquad
m_{V_3}=g(C_F),\qquad
m_{V_3'}=\frac{g(X)-g(C_6)-3g(C_F)}{3}. \;}
\]
特别地，在本文数据
\[
g(X)=49,\qquad g(H)=5,\qquad g(C_6)=13,\qquad g(C_F)=3
\]
下，有
\[
(m_{\mathrm{sgn}},m_{V_2},m_{V_3},m_{V_3'})=(5,4,3,9).
\]
\end{corollary}
\begin{proof}
由定理 \ref{thm:killo-s4-burnside-kani-rosen-prym-square}，
\[
J(C_6)\sim_{\QQ}J(H)\times J(Z)^2,
\]
故取维数得到
\[
g(C_6)=g(H)+2g(Z),
\qquad
g(Z)=\frac{g(C_6)-g(H)}{2}.
\]
再由定理 \ref{thm:xi-time-part9n1b-prym-chevalley-weil-reciprocity}，
\[
m_{\mathrm{sgn}}=g(H),\qquad
m_{V_2}=g(Z),\qquad
m_{V_3}=g(C_F).
\]
最后仍由该定理，
\[
g(X)=\dim V=m_{\mathrm{sgn}}+2m_{V_2}+3m_{V_3}+3m_{V_3'},
\]
代入前三项即可解得
\[
m_{V_3'}=\frac{g(X)-g(C_6)-3g(C_F)}{3}.
\]
将 \((49,5,13,3)\) 代入即得 \((5,4,3,9)\)。证毕。
\end{proof}

\begin{theorem}[Burnside--Kani--Rosen 关系所迫定的两个隐含商曲线 genus]\label{thm:xi-time-part9n1b-kani-rosen-hidden-quotient-genera}
\leanverified{paper\_xi\_time\_part9n1b\_kani\_rosen\_hidden\_quotient\_genera}
沿用记号
\[
H:=X/A_4,\qquad C_6:=X/V_4,\qquad C_F:=X/S_3,\qquad Y:=X/C_4.
\]
若
\[
g(H)=5,\qquad g(C_6)=13,\qquad g(C_F)=3,\qquad g(Y)=13,
\]
则 Burnside--Kani--Rosen 关系在维数层面强制
\[
\boxed{\;
g(X/C_2)=19,\qquad g(X/C_3)=17.\;}
\]
\end{theorem}
\begin{proof}
由定理 \ref{thm:killo-s4-burnside-kani-rosen-prym-square}，
\[
J(X/C_2)\sim_{\QQ}J(C_F)^2\times J(Y).
\]
取维数即得
\[
g(X/C_2)=2g(C_F)+g(Y)=2\cdot 3+13=19.
\]

同一 Kani--Rosen 关系还给出
\[
J(C_6)\times J(X/C_3)^2\sim_{\QQ}J(H)^3\times J(C_F)^2\times J(Y)^2.
\]
再次取维数，得到
\[
g(C_6)+2g(X/C_3)=3g(H)+2g(C_F)+2g(Y).
\]
代入
\[
g(C_6)=13,\qquad g(H)=5,\qquad g(C_F)=3,\qquad g(Y)=13
\]
便有
\[
13+2g(X/C_3)=3\cdot 5+2\cdot 3+2\cdot 13=47,
\]
从而
\[
g(X/C_3)=17.
\]
证毕。
\end{proof}

\begin{theorem}[Jacobian 投影代数的显式幂等分层]\label{thm:xi-time-part9n1b-jacobian-projector-idempotent-algebra}
\leanverified{paper\_xi\_time\_part9n1b\_jacobian\_projector\_idempotent\_algebra}
有理自同态代数 \(\operatorname{End}^0_{\QQ}(J(X))\) 含有规范半单 \(\QQ\)-子代数
\[
\boxed{\;
\mathcal A_X:=\QQ\times M_2(\QQ)\times M_3(\QQ)\times M_3(\QQ)
\subseteq \operatorname{End}^0_{\QQ}(J(X)).\;}
\]
因而
\[
\boxed{\;\dim_{\QQ}\mathcal A_X=1+4+9+9=23.\;}
\]
此外，在 \(\mathcal A_X\) 中可取 \(9\) 个两两正交、非零且和为 \(1\) 的幂等元；并且存在 \(4\) 个中心幂等元，分别投影到
\[
J(H),\qquad J(Z)^2,\qquad J(C_F)^3,\qquad P^3
\]
这四个可见等型块。
\end{theorem}
\begin{proof}
定理 \ref{thm:xi-time-part9n-jacobian-endomorphism-lower-bound} 已给出块对角嵌入
\[
\QQ\times M_2(\QQ)\times M_3(\QQ)\times M_3(\QQ)
\hookrightarrow
\operatorname{End}^0_{\QQ}(J(X)),
\]
故 \(\mathcal A_X\) 的存在与 \(23\) 维结论立即成立。

现取标准矩阵单位
\[
E_{11}^{(2)},E_{22}^{(2)}\in M_2(\QQ),
\qquad
E_{11}^{(3)},E_{22}^{(3)},E_{33}^{(3)}\in M_3(\QQ),
\]
并在第二个 \(M_3(\QQ)\) 块上再取一份同样的对角矩阵单位。于是
\[
(1,0,0,0),
\quad
(0,E_{11}^{(2)},0,0),\ (0,E_{22}^{(2)},0,0),
\]
\[
(0,0,E_{11}^{(3)},0),\ (0,0,E_{22}^{(3)},0),\ (0,0,E_{33}^{(3)},0),
\]
\[
(0,0,0,E_{11}^{(3)}),\ (0,0,0,E_{22}^{(3)}),\ (0,0,0,E_{33}^{(3)})
\]
是 \(9\) 个两两正交的非零幂等元，并且其和恰为 \(\mathcal A_X\) 的单位元。

另一方面，四个块单位
\[
e_{\mathrm{sgn}}:=(1,0,0,0),\qquad
e_{V_2}:=(0,I_2,0,0),\qquad
e_{V_3}:=(0,0,I_3,0),\qquad
e_{V_3'}:=(0,0,0,I_3)
\]
位于直积代数的中心，并满足
\[
e_{\mathrm{sgn}}+e_{V_2}+e_{V_3}+e_{V_3'}=1.
\]
在定理 \ref{thm:killo-s4-genus49-jacobian-computable-isogeny} 给出的分块
\[
J(X)\sim_{\QQ}J(H)\times J(Z)^2\times J(C_F)^3\times P^3
\]
下，它们分别投影到所述四个可见等型块。证毕。
\end{proof}

\begin{corollary}[\texorpdfstring{$\operatorname{Prym}(C_6/H)$}{Prym(C6/H)} 在可见因子表中的唯一平方根]\label{cor:xi-time-part9n1b-prym-c6h-unique-square-root}
在几何因子表
\[
\{\,J(H),\ J(Z),\ J(C_F),\ P\,\}
\]
中，若某个因子 \(A\) 满足
\[
A^2\sim_{\QQ}\operatorname{Prym}(C_6/H),
\]
则必有
\[
\boxed{\;A\sim_{\QQ}J(Z).\;}
\]
\end{corollary}
\begin{proof}
由定理 \ref{thm:killo-s4-burnside-kani-rosen-prym-square}，
\[
\operatorname{Prym}(C_6/H)\sim_{\QQ}J(Z)^2.
\]
故
\[
\dim\operatorname{Prym}(C_6/H)=2g(Z)=8.
\]
若 \(A^2\sim_{\QQ}\operatorname{Prym}(C_6/H)\)，则必有
\[
2\dim A=8,
\qquad
\dim A=4.
\]
另一方面，由结论 \ref{con:xi-time-part9n-s4-all-subgroup-genera} 与结论 \ref{con:xi-time-part9n-s4-hodge-tensor-factorization}，
\[
\dim J(H)=g(H)=5,\qquad
\dim J(Z)=g(Z)=4,
\]
\[
\dim J(C_F)=g(C_F)=3,\qquad
\dim P=9.
\]
在这四个可见因子中，只有 \(J(Z)\) 的维数等于 \(4\)。故 \(A\sim_{\QQ}J(Z)\) 为唯一可能。证毕。
\end{proof}

\endinput

\paragraph{prime-register solenoid 对 Hurwitz 链 Jacobian/Prym 的输入阻断与输出代价}
在局部化数系对偶 \(\Sigma_S=\widehat{G_S}\) 的 Pontryagin 刚性与 \(S_4\)-Hurwitz 商塔的 Jacobian/Prym 同源分解都已确定之后，还可把这两条主线压缩为一组统一的态射结论：正维紧复环面指向有限 prime-register solenoid 直积的输入完全消失，而反向输出则仅按因子个数贡献实相位方向；因此 Hurwitz 链中各 Jacobian/Prym 因子的 prime-register 最小输出代价可被精确求值，并对 Kani--Rosen 同源分解严格可加。

\begin{theorem}[有限 prime-register solenoid 积对正维紧复环面不存在非零输入态射]\label{thm:xi-time-part9n1d-prime-register-input-obstruction}
沿用定理 \ref{thm:killo-localization-circle-dimension-prime-ledger-splitting} 的记号。设
\[
S_1,\dots,S_r\subset \mathbb P
\]
为有限非空素数集，并记
\[
G_{S_i}:=\ZZ[S_i^{-1}],
\qquad
\Sigma_{S_i}:=\widehat{G_{S_i}},
\qquad
K:=\prod_{i=1}^r \Sigma_{S_i}.
\]
则对任意复维 \(g\ge 1\) 的紧复环面 \(A\)，都有
\[
\boxed{\;
\operatorname{Hom}_{\mathrm{cts}}(A,K)=0.\;}
\]
特别地，对任意 Jacobian 或 Prym 都成立。
\end{theorem}
\begin{proof}
由于 \(A\) 是复维 \(g\) 的紧复环面，其 Pontryagin 对偶满足
\[
\widehat A\cong \ZZ^{2g}.
\]
又因为 Pontryagin 对偶把有限直积送为直和，故
\[
\widehat K\cong \bigoplus_{i=1}^r G_{S_i}.
\]
任取连续群同态
\[
f:A\longrightarrow K.
\]
对偶化得到离散群同态
\[
\widehat f:\bigoplus_{i=1}^r G_{S_i}\longrightarrow \ZZ^{2g}.
\]
只需证明任意群同态
\[
\phi:G_{S_i}\longrightarrow \ZZ^{2g}
\]
都恒为零。固定 \(i\) 并取任意 \(p\in S_i\)。对每个 \(n\ge 1\)，由于 \(p^{-n}\in G_{S_i}\)，有
\[
p^n\phi(p^{-n})=\phi(1).
\]
因此
\[
\phi(1)\in \bigcap_{n\ge 1}p^n\ZZ^{2g}=\{0\},
\]
故 \(\phi(1)=0\)。再取任意
\[
x=\frac{a}{b}\in G_{S_i},
\]
其中 \(a\in \ZZ\)，且 \(b\) 的全部素因子都属于 \(S_i\)。则
\[
b\,\phi(x)=\phi(a)=a\,\phi(1)=0.
\]
由于 \(\ZZ^{2g}\) 无挠，只能有 \(\phi(x)=0\)。故 \(\phi=0\)，从而 \(\widehat f=0\)。再对偶化便得 \(f=0\)。证毕。
\end{proof}

\begin{theorem}[有限 prime-register solenoid 积对紧交换 Lie 环面的输出秩与满射阈值]\label{thm:xi-time-part9n1d-prime-register-output-rank-threshold}
\leanverified{paper\_xi\_time\_part9n1d\_prime\_register\_output\_rank\_threshold}
沿用定理 \ref{thm:xi-time-part9n1d-prime-register-input-obstruction} 的记号。设 \(A\) 为实维 \(d\ge 1\) 的紧连通交换 Lie 群。则 \(A\cong \TT^d\)，并且对任意连续群同态
\[
f:K\longrightarrow A
\]
成立：
\begin{enumerate}
  \item \(f(K)\) 是 \(A\) 的一个实维数至多为 \(r\) 的子环面；
  \item 存在连续满射
  \[
  K\twoheadrightarrow A
  \]
  当且仅当 \(r\ge d\)。
\end{enumerate}
等价地，
\[
\boxed{\;
\min\left\{
r\ge 1:\ \exists\ S_1,\dots,S_r\subset \mathbb P\ \text{有限非空},
\prod_{i=1}^r \Sigma_{S_i}\twoheadrightarrow A
\right\}
=d.\;}
\]
对任意复维 \(g\) 的 Jacobian 或 Prym，上式右端即 \(2g\)。
\end{theorem}
\begin{proof}
记
\[
B:=f(K).
\]
每个 \(G_{S_i}\subset \QQ\) 都是无挠群，因此其 Pontryagin 对偶 \(\Sigma_{S_i}\) 连通；故 \(K\) 连通，从而 \(B\) 也是 \(A\) 的连通紧子群。由于 \(A\cong \TT^d\)，可知 \(B\) 必为一个子环面，因而
\[
\widehat B\cong \ZZ^{\dim_{\mathbb R}B}.
\]

由满射
\[
K\twoheadrightarrow B
\]
对偶化得到单射
\[
\widehat B\hookrightarrow \widehat K\cong \bigoplus_{i=1}^r G_{S_i}.
\]
另一方面，定理 \ref{thm:killo-localization-circle-dimension-prime-ledger-splitting} 给出
\[
G_{S_i}\otimes_{\ZZ}\QQ\cong \QQ,
\]
故
\[
\left(\bigoplus_{i=1}^r G_{S_i}\right)\otimes_{\ZZ}\QQ\cong \QQ^r.
\]
因此 \(\widehat B\) 的秩至多为 \(r\)，即
\[
\dim_{\mathbb R}B=\operatorname{rk}_{\ZZ}(\widehat B)\le r.
\]
第 \(1\) 条得证。

若 \(f\) 为满射，则 \(B=A\)，于是
\[
d=\dim_{\mathbb R}A\le r.
\]
反之，若 \(r\ge d\)，则由定理 \ref{thm:killo-localization-circle-dimension-prime-ledger-splitting}，每个 \(\Sigma_{S_i}\) 都带有规范满射
\[
\pi_{S_i}:\Sigma_{S_i}\twoheadrightarrow \TT.
\]
取直积得到
\[
\prod_{i=1}^r\pi_{S_i}:K\twoheadrightarrow \TT^r.
\]
再投影到前 \(d\) 个坐标，即得
\[
K\twoheadrightarrow \TT^d\cong A.
\]
故第 \(2\) 条成立。最后一个等价表述只是把第 \(2\) 条重写为最小因子数公式。证毕。
\end{proof}

\begin{corollary}[Hurwitz 链各 Jacobian/Prym 因子的 prime-register 最小输出代价]\label{cor:xi-time-part9n1d-prime-register-hurwitz-minimal-output-cost}
沿用定理 \ref{thm:killo-s4-genus49-jacobian-computable-isogeny} 与结论 \ref{con:xi-time-part9n-s4-all-subgroup-genera} 的记号。对任意正维紧复环面 \(A\)，定义
\[
r_{\min}(A):=
\min\left\{
r\ge 1:\ \exists\ S_1,\dots,S_r\subset \mathbb P\ \text{有限非空},
\prod_{i=1}^r \Sigma_{S_i}\twoheadrightarrow A
\right\}.
\]
则
\[
\boxed{\;
\begin{aligned}
r_{\min}(J(H))&=10,\qquad r_{\min}(J(C_6))=26,\qquad r_{\min}(J(Z))=8,\\
r_{\min}(J(C_F))&=6,\qquad r_{\min}(P)=18,\qquad r_{\min}(J(X))=98.
\end{aligned}
\;}
\]
\end{corollary}
\begin{proof}
由定理 \ref{thm:xi-time-part9n1d-prime-register-output-rank-threshold}，任意正维紧复环面 \(A\) 的最小 prime-register 输出代价满足
\[
r_{\min}(A)=2\dim_{\mathbb C}A.
\]
而由结论 \ref{con:xi-time-part9n-s4-all-subgroup-genera} 与定理 \ref{thm:killo-s4-genus49-jacobian-computable-isogeny}，
\[
g(H)=5,\qquad g(C_6)=13,\qquad g(Z)=4,\qquad g(C_F)=3,
\]
\[
\dim_{\mathbb C}P=9,\qquad g(X)=49.
\]
对 Jacobian 有 \(\dim_{\mathbb C}J(C)=g(C)\)。将上述维数代入即可得到所示六个精确数值。证毕。
\end{proof}

\begin{theorem}[Kani--Rosen 同源分解上的 prime-register 输出代价严格可加]\label{thm:xi-time-part9n1d-prime-register-output-cost-additivity}
\leanverified{paper\_xi\_time\_part9n1d\_prime\_register\_output\_cost\_additivity}
沿用定理 \ref{thm:killo-s4-genus49-jacobian-computable-isogeny} 的记号。对任意正维紧复环面 \(A\)，定义
\[
\rho(A):=
\min\left\{
r\ge 1:\ \exists\ S_1,\dots,S_r\subset \mathbb P\ \text{有限非空},
\prod_{i=1}^r \Sigma_{S_i}\twoheadrightarrow A
\right\}.
\]
则对同源分解
\[
J(X)\sim_{\QQ}J(H)\times J(Z)^2\times J(C_F)^3\times P^3
\]
有严格可加公式
\[
\boxed{\;
\rho(J(X))
=
\rho(J(H))
+2\,\rho(J(Z))
+3\,\rho(J(C_F))
+3\,\rho(P).\;}
\]
特别地，
\[
\boxed{\;
98=10+2\cdot 8+3\cdot 6+3\cdot 18.\;}
\]
\end{theorem}
\begin{proof}
由定理 \ref{thm:xi-time-part9n1d-prime-register-output-rank-threshold}，任意正维紧复环面 \(A\) 都满足
\[
\rho(A)=2\dim_{\mathbb C}A.
\]
另一方面，由同源分解的维数可加性，
\[
\dim_{\mathbb C}J(X)
=
\dim_{\mathbb C}J(H)
+2\,\dim_{\mathbb C}J(Z)
+3\,\dim_{\mathbb C}J(C_F)
+3\,\dim_{\mathbb C}P.
\]
将上式两边同时乘以 \(2\)，即得 \(\rho\) 的严格可加性。再代入
\[
49=5+2\cdot 4+3\cdot 3+3\cdot 9
\]
便得到数值公式
\[
98=10+2\cdot 8+3\cdot 6+3\cdot 18.
\]
证毕。
\end{proof}

\begin{proposition}[九维 Prym 因子的单寄存器输入消失与单相位输出上界]\label{prop:xi-time-part9n1d-prym-single-register-asymmetry}
沿用定理 \ref{thm:killo-s4-genus49-jacobian-computable-isogeny} 的记号，记
\[
P:=\operatorname{Prym}(Y/Z).
\]
任取有限非空素数集 \(S\subset \mathbb P\)，并记
\[
\Sigma_S:=\widehat{\ZZ[S^{-1}]}.
\]
则有
\[
\boxed{\;
\operatorname{Hom}_{\mathrm{cts}}(P,\Sigma_S)=0,
\qquad
\dim_{\mathbb R}\operatorname{im}(g)\le 1
\quad
(\forall\, g:\Sigma_S\to P).\;}
\]
等价地，单一 prime-register 对 \(P\) 只能输出至多一个实相位方向，而不能从 \(P\) 读取任何非零连续输入。
\end{proposition}
\begin{proof}
由定理 \ref{thm:killo-s4-genus49-jacobian-computable-isogeny}，\(\dim_{\mathbb C}P=9\)。因此把定理 \ref{thm:xi-time-part9n1d-prime-register-input-obstruction} 应用于 \(r=1\) 与 \(A=P\)，立即得到
\[
\operatorname{Hom}_{\mathrm{cts}}(P,\Sigma_S)=0.
\]
再把定理 \ref{thm:xi-time-part9n1d-prime-register-output-rank-threshold} 应用于 \(r=1\) 与 \(A=P\)，便得任意连续群同态
\[
g:\Sigma_S\to P
\]
的像都是 \(P\) 中实维数至多为 \(1\) 的子环面。证毕。
\end{proof}

\noindent
上述五条结果表明，局部化数系 prime-register solenoid 与 \(S_4\)-Hurwitz 链 Jacobian/Prym 的接口在输入与输出两端呈现严格非对称性：输入方向不存在任何非零连续态射，而输出方向中每个 prime-register 恰只贡献一个实相位。因而 prime-register 最小输出代价不仅对各个 Jacobian/Prym 因子可精确求值，而且与 Kani--Rosen 同源分解逐块严格兼容。

\endinput


\noindent
结合结论 \ref{con:xi-time-part9n-s4-all-subgroup-genera} 与以上诸条结论，\(S_4\)-Hurwitz 商塔中这条以 Chevalley--Weil 重数为起点的链条已在几何、表示、算术与 prime-register 态射预算四侧同时闭合：一方面，\(C_6\) 作为 \(H\) 与 \(Z\) 的规范化纤维积被唯一识别，全纯微分空间按四个非平凡不可约表示张量分层，八个商曲线的 genus 以及 \(X/C_2\)、\(X/C_3\) 等此前缺失商的 Jacobian 同源型都由 Burnside--Kani--Rosen 关系与同一组表示论数据唯一补全，从而四个非平凡不可约表示的几何承载因子不再遗留任何隐藏等型块；另一方面，\(J(X)\) 的有理投影代数下界、\(\operatorname{Prym}(C_6/H)\) 的平方根刚性、良素处局部 Euler 因子的平方/立方块、两类 Prym 的最小同源代价及其 \(3\)-初等与 \(2\)-初等极化缺陷，以及 \(V_3'\)-等型三重块 \(P^3\) 的主极化化次数、\(\QQ\)-子等型分类与 Rosati 对称端代数，继续由群论分块与极化缺陷严格锁定；此外，有限 prime-register solenoid 对正维 Jacobian/Prym 的输入态射完全消失，而其输出代价恰等于目标实维数，并对 Kani--Rosen 同源分解逐块保持严格可加。

\endinput


\begin{conclusion}[\texorpdfstring{$\operatorname{Prym}(C_6/H)$}{Prym(C6/H)} 的 \texorpdfstring{$\QQ(\zeta_3)$}{Q(zeta3)}-Weil 型结构与签名 \texorpdfstring{$(4,4)$}{(4,4)}]\label{con:xi-time-part9n-prym-c6h-weil-type-signature}
阿贝尔簇 \(A\) 是以
\[
K:=\QQ(\zeta_3)
\]
为虚二次域的 Weil 型阿贝尔八维簇，其 \(K\)-signature 为
\[
(4,4).
\]
并且 \(A\) 的诱导极化型为
\[
(1,1,1,1,3,3,3,3),
\]
且有
\[
\QQ(\zeta_3)\subset \operatorname{End}^0(A),\qquad
A\sim_{\QQ}J(Z)^2.
\]
特别地，存在自然二维 \(\QQ\)-向量空间
\[
W_K:=\bigwedge\nolimits_{K}^{8}H^1(A,\QQ)\subset H^8(A,\QQ)\cap H^{4,4}(A),
\]
其由 Weil classes 给出；经同源 \(A\sim_{\QQ}J(Z)^2\)，该二维 Hodge 类空间亦传递到
\[
H^8\!\bigl(J(Z)^2,\QQ\bigr)\cap H^{4,4}\!\bigl(J(Z)^2\bigr).
\]
\end{conclusion}
\begin{proof}
由定理 \ref{thm:killo-s4-even-subgroup-free-action-unramified-a4-lift}，覆盖
\[
\pi:C_6\to H
\]
为无分歧循环三重覆盖，且
\[
g(H)=5,\qquad g(C_6)=13.
\]
故
\[
\dim A=g(C_6)-g(H)=8.
\]
再由定理 \ref{thm:killo-s4-burnside-kani-rosen-prym-square} 与命题 \ref{prop:killo-two-prym-polarization-types}，
\[
A\sim_{\QQ}J(Z)^2,
\]
并且 \(A\) 的诱导极化型为 \((1^4,3^4)\)，同时
\[
\QQ(\zeta_3)\subset \operatorname{End}^0(A).
\]

为确定 signature，记 \(\sigma\) 为 \(\pi\) 的甲板变换生成元。由于 \(\pi\) 为循环三重无分歧覆盖，存在非常值 \(3\)-扭线丛
\[
L\in \operatorname{Pic}^0(H),\qquad L^{\otimes 3}\simeq \mathcal O_H,
\]
使得
\[
\pi_*\mathcal O_{C_6}\simeq \mathcal O_H\oplus L^{-1}\oplus L^{-2}.
\]
由投影公式，
\[
\pi_*\omega_{C_6}
\simeq
\omega_H\oplus (\omega_H\otimes L)\oplus (\omega_H\otimes L^2).
\]
于是 \(H^0(C_6,\Omega_{C_6}^1)\) 按 \(C_3=\langle \sigma\rangle\) 的特征分解为
\[
H^0(C_6,\Omega_{C_6}^1)=U_0\oplus U_{\chi}\oplus U_{\bar\chi},
\qquad
\chi(\sigma)=\zeta_3,
\]
其中
\[
U_0\simeq H^0(H,\omega_H),\qquad
U_{\chi}\simeq H^0(H,\omega_H\otimes L),\qquad
U_{\bar\chi}\simeq H^0(H,\omega_H\otimes L^2).
\]
由于 \(L\) 与 \(L^2\) 均为非平凡零次数线丛，Riemann--Roch 与 Serre 对偶给出
\[
h^0(H,\omega_H\otimes L)=h^0(H,\omega_H\otimes L^2)=g(H)-1=4,
\]
而
\[
h^0(H,\omega_H)=g(H)=5.
\]
故
\[
\dim U_0=5,\qquad \dim U_{\chi}=\dim U_{\bar\chi}=4.
\]
Prym 余切空间正对应于非平凡特征部分
\[
H^0(A,\Omega_A^1)\simeq U_{\chi}\oplus U_{\bar\chi},
\]
因此 \(K\)-作用在 \(H^{1,0}(A)\) 上的两条复嵌入重数均为 \(4\)，即 signature 为 \((4,4)\)。

另一方面，甲板变换群 \(C_3\) 通过保持 Prym 诱导极化的自同构作用于 \(A\)，故 Rosati 对合把 \(\sigma\) 送到 \(\sigma^{-1}\)，也即在 \(K=\QQ(\zeta_3)\) 上对应复共轭。于是 \(A\) 在该极化下确为 \(K\)-Weil 型阿贝尔八维簇。对签名 \((4,4)\) 的 Weil 型阿贝尔簇，一般理论表明
\[
W_K=\bigwedge\nolimits_K^{8}H^1(A,\QQ)
\]
是二维 \(\QQ\)-向量空间，并落在 \(H^8(A,\QQ)\cap H^{4,4}(A)\) 中。最后，同源 \(A\sim_{\QQ}J(Z)^2\) 在有理上同调上诱导 Hodge 结构同构，故 \(W_K\) 传递到 \(H^8(J(Z)^2,\QQ)\) 中相应的 \((4,4)\)-部分。证毕。
\end{proof}

\begin{conclusion}[两类 Prym 缺陷核在素数支撑上的完全分离]\label{con:xi-time-part9n-prym-polarization-primary-splitting}
记 \(\lambda_A\) 与 \(\lambda_P\) 分别为命题 \ref{prop:killo-two-prym-polarization-types} 在
\[
A=\operatorname{Prym}(C_6/H),\qquad P=\operatorname{Prym}(Y/Z)
\]
上给出的诱导极化。则有
\[
\ker(\lambda_A)\cong (\ZZ/3\ZZ)^8,
\qquad
\ker(\lambda_P)\cong (\ZZ/2\ZZ)^{10}.
\]
因此，在可计算同源分解
\[
J(X)\sim_{\QQ}J(H)\times J(Z)^2\times J(C_F)^3\times P^3
\]
中，显式出现的非主极化缺陷按素数严格解耦：\(3\)-初缺陷全部锚定在
\[
A\sim_{\QQ}J(Z)^2
\]
这一 Prym 平方方向，而 \(2\)-初缺陷全部锚定在 \(P\) 方向；相对地，\(J(H)\) 与 \(J(C_F)\) 作为 Jacobian 因子各自携带其自然的主极化。
\end{conclusion}
\begin{proof}
对任意极化型
\[
(d_1,\dots,d_g),
\]
其极化同态的核满足标准结构公式
\[
\ker(\lambda)\cong \bigoplus_{i=1}^{g}(\ZZ/d_i\ZZ)^2.
\]
由命题 \ref{prop:killo-two-prym-polarization-types}，
\[
\mathrm{type}(\lambda_A)=(1,1,1,1,3,3,3,3),
\]
故
\[
\ker(\lambda_A)\cong (\ZZ/3\ZZ)^8.
\]
同理，
\[
\mathrm{type}(\lambda_P)=(1,1,1,1,2,2,2,2,2),
\]
从而
\[
\ker(\lambda_P)\cong (\ZZ/2\ZZ)^{10}.
\]
再由定理 \ref{thm:killo-s4-burnside-kani-rosen-prym-square} 与
\ref{thm:killo-s4-genus49-jacobian-computable-isogeny}，
\[
A\sim_{\QQ}J(Z)^2,\qquad
J(X)\sim_{\QQ}J(H)\times J(Z)^2\times J(C_F)^3\times P^3.
\]
故与 \(A\) 对应的缺陷群是纯 \(3\)-群，而与 \(P\) 对应的缺陷群是纯 \(2\)-群，二者在素数层面彼此互素。最后，\(J(H)\) 与 \(J(C_F)\) 作为 Jacobian，本身都带有其规范主极化，故不引入新的极化缺陷素因子。证毕。
\end{proof}

\begin{conclusion}[Burnside--Kani--Rosen 同源关系诱导的局部 \texorpdfstring{$\zeta$}{zeta} 因子分解与点数恒等式]\label{con:xi-time-part9n-burnside-kani-rosen-local-zeta-pointcount}
设 \(v\) 为一处使
\[
X,\ H,\ C_6,\ Z,\ C_F,\ Y,\ X/C_2,\ X/C_3
\]
及阿贝尔簇 \(P\) 同时良约化的有限素位，记残域大小为 \(q_v\)。对任意光滑射影曲线 \(C/\FF_{q_v}\)，定义
\[
Z(C/\FF_{q_v},T)=\frac{P_C(v,T)}{(1-T)(1-q_vT)},
\]
其中 \(P_C(v,T)\) 为 Frobenius 在 \(H^1_{\text{ét}}(C_{\overline{\FF}_{q_v}},\QQ_\ell)\) 上的特征多项式；并记
\[
P_P(v,T):=
\det\!\Bigl(1-T\,\mathrm{Frob}_v\mid
H^1_{\text{ét}}(P_{\overline{\FF}_{q_v}},\QQ_\ell)\Bigr).
\]
则有
\[
P_{C_6}(v,T)=P_H(v,T)\,P_Z(v,T)^2,
\]
\[
P_{X/C_2}(v,T)=P_{C_F}(v,T)^2\,P_Y(v,T),
\]
\[
P_{C_6}(v,T)\,P_{X/C_3}(v,T)^2
=
P_H(v,T)^3\,P_{C_F}(v,T)^2\,P_Y(v,T)^2,
\]
以及
\[
P_X(v,T)=P_H(v,T)\,P_Z(v,T)^2\,P_{C_F}(v,T)^3\,P_P(v,T)^3.
\]

等价地，对任意 \(n\ge 1\)，若记
\[
N_C(v,n):=\#C(\FF_{q_v^n}),
\qquad
a_C(v,n):=q_v^n+1-N_C(v,n),
\]
则有精确恒等式
\[
a_{C_6}(v,n)=a_H(v,n)+2a_Z(v,n),
\]
\[
a_{X/C_2}(v,n)=2a_{C_F}(v,n)+a_Y(v,n),
\]
\[
a_{C_6}(v,n)+2a_{X/C_3}(v,n)
=
3a_H(v,n)+2a_{C_F}(v,n)+2a_Y(v,n).
\]
\end{conclusion}
\begin{proof}
由定理 \ref{thm:killo-s4-burnside-kani-rosen-prym-square}，
\[
J(C_6)\sim_{\QQ}J(H)\times J(Z)^2,
\]
\[
J(X/C_2)\sim_{\QQ}J(C_F)^2\times J(Y),
\]
\[
J(C_6)\times J(X/C_3)^2
\sim_{\QQ}
J(H)^3\times J(C_F)^2\times J(Y)^2.
\]
再由定理 \ref{thm:killo-s4-genus49-jacobian-computable-isogeny}，
\[
J(X)\sim_{\QQ}J(H)\times J(Z)^2\times J(C_F)^3\times P^3.
\]
在共同良约化处，同源阿贝尔簇的 \(\ell\)-进 \(H^1\) 具有相同的 Frobenius 特征多项式，而直积对应特征多项式相乘，故立即得到四条 \(P\)-因子恒等式。

对任意曲线 \(C/\FF_{q_v}\)，标准展开给出
\[
\log P_C(v,T)=-\sum_{n\ge 1}\frac{a_C(v,n)}{n}T^n.
\]
将前三条 \(P\)-因子恒等式取对数并比较 \(T^n\) 系数，即得三条点数恒等式。证毕。
\end{proof}

\begin{conclusion}[九维 Prym 因子的 Frobenius 迹整除律]\label{con:xi-time-part9n-prym9-frobenius-trace-divisibility}
对任意共同良约化素位 \(v\) 及任意 \(n\ge 1\)，有严格公式
\[
\operatorname{Tr}\!\Bigl(
\mathrm{Frob}_v^n\mid
H^1_{\text{ét}}(P_{\overline{\FF}_{q_v}},\QQ_\ell)
\Bigr)
=
\frac{
a_X(v,n)-a_H(v,n)-2a_Z(v,n)-3a_{C_F}(v,n)
}{3}.
\]
特别地，
\[
a_X(v,n)-a_H(v,n)-2a_Z(v,n)-3a_{C_F}(v,n)\equiv 0\pmod 3.
\]
\end{conclusion}
\begin{proof}
由定理 \ref{thm:killo-s4-genus49-jacobian-computable-isogeny}，
\[
J(X)\sim_{\QQ}J(H)\times J(Z)^2\times J(C_F)^3\times P^3.
\]
因此在共同良约化处，
\[
\operatorname{Tr}\!\Bigl(
\mathrm{Frob}_v^n\mid H^1_{\text{ét}}(J(X)_{\overline{\FF}_{q_v}},\QQ_\ell)
\Bigr)
=
a_H(v,n)+2a_Z(v,n)+3a_{C_F}(v,n)
\]
\[
\qquad\qquad
+3\,\operatorname{Tr}\!\Bigl(
\mathrm{Frob}_v^n\mid H^1_{\text{ét}}(P_{\overline{\FF}_{q_v}},\QQ_\ell)
\Bigr).
\]
又对曲线 Jacobian 有
\[
H^1_{\text{ét}}(J(X)_{\overline{\FF}_{q_v}},\QQ_\ell)
\cong
H^1_{\text{ét}}(X_{\overline{\FF}_{q_v}},\QQ_\ell),
\]
故左端正是 \(a_X(v,n)\)。移项并除以 \(3\) 即得所示公式。右端既等于实际的 Frobenius 迹，因而必为整数，从而得到模 \(3\) 整除律。证毕。
\end{proof}

\begin{conclusion}[属 \(49\) Jacobian 上由 \texorpdfstring{$\QQ(\zeta_3)$}{Q(zeta3)}-Weil 型强制的余维 \texorpdfstring{$4$}{4} Hodge 类]\label{con:xi-time-part9n-jx-codim4-weil-hodge-classes}
在 \(H^8(J(X),\QQ)\) 中存在自然二维 \(\QQ\)-子空间
\[
W_{K,X}\subset H^8(J(X),\QQ)\cap H^{4,4}(J(X)),
\qquad K=\QQ(\zeta_3),
\]
它由结论 \ref{con:xi-time-part9n-prym-c6h-weil-type-signature} 中 \(A=\operatorname{Prym}(C_6/H)\) 的 Weil classes 经同源分解
\[
J(X)\sim_{\QQ}J(H)\times J(Z)^2\times J(C_F)^3\times P^3
\]
传递而来。换言之，属 \(49\) 的 Jacobian \(J(X)\) 含有一组由无分歧循环三重覆盖 \(C_6\to H\) 强制产生的余维 \(4\) Hodge 类。
\end{conclusion}
\begin{proof}
由结论 \ref{con:xi-time-part9n-prym-c6h-weil-type-signature}，存在二维 \(\QQ\)-向量空间
\[
W_K\subset H^8(A,\QQ)\cap H^{4,4}(A),
\qquad A=\operatorname{Prym}(C_6/H).
\]
同一结论与定理 \ref{thm:killo-s4-burnside-kani-rosen-prym-square} 给出
\[
A\sim_{\QQ}J(Z)^2.
\]
故 \(W_K\) 可视为
\[
H^8\!\bigl(J(Z)^2,\QQ\bigr)\cap H^{4,4}\!\bigl(J(Z)^2\bigr)
\]
中的二维 Hodge 子结构。

另一方面，定理 \ref{thm:killo-s4-genus49-jacobian-computable-isogeny} 给出
\[
J(X)\sim_{\QQ}J(H)\times J(Z)^2\times J(C_F)^3\times P^3.
\]
该同源分解在有理上同调上诱导 Hodge 结构同构，因此
\[
H^8\!\bigl(J(Z)^2,\QQ\bigr)
\]
自然识别为 \(H^8(J(X),\QQ)\) 的一个有理 Hodge 直和因子。将上述 \(W_K\) 沿此同构拉回，即得
\[
W_{K,X}\subset H^8(J(X),\QQ)\cap H^{4,4}(J(X)),
\]
且 \(\dim_{\QQ}W_{K,X}=2\)。证毕。
\end{proof}

同一组 \(S_4\)-Hurwitz 商塔数据还进一步固定了权一表示的等型重数、若干未显式写出的商曲线亏格、八维 Prym 的端同态约束以及 \(J(C_6)\) 在 Siegel 模空间中的位置；相应地可压缩为下列四条结论。

\begin{conclusion}[\texorpdfstring{$S_4$}{S4} 权一表示的完全分解与全部商曲线亏格的统一恢复]\label{con:xi-time-part9n-s4-hodge-quotient-genera-recovery}
令
\[
V:=H^0(X,\Omega_X^1).
\]
则有有理 \(S_4\)-表示的完全分解
\[
\boxed{\;
V\cong
\mathrm{sgn}^{\oplus 5}\oplus
V_2^{\oplus 4}\oplus
V_3^{\oplus 3}\oplus
\bigl(V_3'\bigr)^{\oplus 9}.\;}
\]
并且对任意子群 \(K\le S_4\)，若记
\[
\rho_K:=\Ind_K^{S_4}\mathbf 1,
\]
则对应商曲线亏格满足
\[
g(X/K)=\dim V^K=\langle V,\rho_K\rangle_{S_4}.
\]
特别地，
\[
\boxed{\;
g(X/C_2)=19,\qquad g(X/C_3)=17.\;}
\]
与此同时，文中已出现的四个商曲线亏格
\[
g(H)=5,\qquad g(Z)=4,\qquad g(C_F)=3,\qquad g(Y)=13
\]
亦由同一角色公式统一恢复。
\end{conclusion}
\begin{proof}
把定理 \ref{thm:killo-s4-genus49-jacobian-computable-isogeny} 所用的等型分块记为
\[
J(X)\sim A_{\mathrm{sgn}}\times A_2\times A_3\times A_{3'}.
\]
由同一定理及其证明中的商曲线识别，
\[
A_{\mathrm{sgn}}\sim J(H),\qquad
A_2\sim J(Z)^2,\qquad
A_3\sim J(C_F)^3,\qquad
A_{3'}\sim P^3.
\]
因此
\[
\dim A_{\mathrm{sgn}}=g(H)=5,\qquad
\dim A_2=2g(Z)=8,
\]
\[
\dim A_3=3g(C_F)=9,\qquad
\dim A_{3'}=3\dim P=27.
\]
若 \(m_\rho\) 表示不可约有理表示 \(\rho\) 在 \(V=H^0(X,\Omega_X^1)\) 中的重数，则对应等型阿贝尔因子的维数满足
\[
\dim A_\rho=(\dim \rho)\,m_\rho.
\]
于是
\[
m_{\mathrm{sgn}}=5,\qquad
m_{V_2}=4,\qquad
m_{V_3}=3,\qquad
m_{V_3'}=9.
\]
又因 \(X/S_4\simeq \PP^1\)，故平凡表示重数为 \(0\)。遂得所示完全分解。

对任意 \(K\le S_4\)，Frobenius 互反给出
\[
\dim V^K=\langle V,\rho_K\rangle_{S_4},
\qquad
\rho_K=\Ind_K^{S_4}\mathbf 1.
\]
将
\[
\rho_{A_4}=\mathbf 1\oplus \mathrm{sgn},\qquad
\rho_{D_8}=\mathbf 1\oplus V_2,\qquad
\rho_{S_3}=\mathbf 1\oplus V_3,
\]
\[
\rho_{C_4}=\mathbf 1\oplus V_2\oplus V_3',\qquad
\rho_{C_2}=\mathbf 1\oplus V_2\oplus 2V_3\oplus V_3',
\]
\[
\rho_{C_3}=\mathbf 1\oplus \mathrm{sgn}\oplus V_3\oplus V_3'
\]
代入即得
\[
g(H)=5,\qquad g(Z)=4,\qquad g(C_F)=3,\qquad g(Y)=4+9=13,
\]
并且
\[
g(X/C_2)=4+2\cdot 3+9=19,
\qquad
g(X/C_3)=5+3+9=17.
\]
证毕。
\end{proof}

\begin{conclusion}[\texorpdfstring{$\operatorname{Prym}(C_6/H)$}{Prym(C6/H)} 的最小端同态表述与平方同源不相容]\label{con:xi-time-part9n-prym-c6h-endomorphism-obstruction}
若将猜想 \ref{conj:killo-minimal-extra-endomorphism-large-prym} 的第二条按字面解释为
\[
\operatorname{End}^0_{\overline{\QQ}}\!\bigl(\operatorname{Prym}(C_6/H)\bigr)=\QQ(\zeta_3),
\]
则该表述与定理 \ref{thm:killo-s4-burnside-kani-rosen-prym-square} 不相容。更精确地，
\[
\boxed{\;
\dim_{\QQ}
\operatorname{End}^0_{\overline{\QQ}}\!\bigl(\operatorname{Prym}(C_6/H)\bigr)
\ge 4.\;}
\]
因此，八维 Prym 的“最小额外自同态”不可能直接在整个端同态代数层面收缩为二次数域 \(\QQ(\zeta_3)\)。
\end{conclusion}
\begin{proof}
由定理 \ref{thm:killo-s4-burnside-kani-rosen-prym-square}，
\[
\operatorname{Prym}(C_6/H)\sim_{\QQ} J(Z)^2.
\]
同源阿贝尔簇的有理端同态代数经由准同源传递同构，因此
\[
\operatorname{End}^0_{\overline{\QQ}}\!\bigl(\operatorname{Prym}(C_6/H)\bigr)
\cong
\operatorname{End}^0_{\overline{\QQ}}\!\bigl(J(Z)^2\bigr).
\]
而对任意正维阿贝尔簇 \(B\)，都有标准恒等式
\[
\operatorname{End}^0_{\overline{\QQ}}(B^2)\cong
M_2\!\bigl(\operatorname{End}^0_{\overline{\QQ}}(B)\bigr).
\]
取 \(B=J(Z)\)，即得
\[
\operatorname{End}^0_{\overline{\QQ}}\!\bigl(\operatorname{Prym}(C_6/H)\bigr)
\cong
M_2\!\bigl(\operatorname{End}^0_{\overline{\QQ}}(J(Z))\bigr).
\]
由于恒等元给出
\[
\dim_{\QQ}\operatorname{End}^0_{\overline{\QQ}}(J(Z))\ge 1,
\]
故
\[
\dim_{\QQ}
\operatorname{End}^0_{\overline{\QQ}}\!\bigl(\operatorname{Prym}(C_6/H)\bigr)
=
4\,
\dim_{\QQ}\operatorname{End}^0_{\overline{\QQ}}(J(Z))
\ge 4.
\]
另一方面，
\[
\dim_{\QQ}\QQ(\zeta_3)=2.
\]
因此不可能有
\[
\operatorname{End}^0_{\overline{\QQ}}\!\bigl(\operatorname{Prym}(C_6/H)\bigr)=\QQ(\zeta_3).
\]
证毕。
\end{proof}

\begin{conclusion}[\texorpdfstring{$V_2$}{V2} 等型块的平衡 \texorpdfstring{$\QQ(\zeta_3)$}{Q(zeta3)}-Weil 结构]\label{con:xi-time-part9n-v2-block-balanced-weil}
令 \(A_2\) 表示 \(J(X)\) 的 \(V_2\)-等型块。则
\[
A_2\sim_{\QQ}J(Z)^2\sim_{\QQ}\operatorname{Prym}(C_6/H).
\]
并且经由同源识别
\[
A_2\sim_{\QQ}\operatorname{Prym}(C_6/H)
\]
所传递的 \(C_3\)-作用满足
\[
H^0(A_2,\Omega_{A_2}^1)\cong \chi^{\oplus 4}\oplus \bar\chi^{\oplus 4},
\]
其中 \(\chi,\bar\chi\) 为 \(C_3\) 的两个非平凡复特征。因而 \(A_2\) 是一个以
\[
K=\QQ(\zeta_3)
\]
为乘法域的八维 Weil 型阿贝尔簇，其 \(K\)-signature 恰为
\[
\boxed{\;(4,4).\;}
\]
此外，\(A_2\) 的诱导极化型为
\[
(1,1,1,1,3,3,3,3).
\]
\end{conclusion}
\begin{proof}
结论 \ref{con:xi-time-part9n-prym-c6h-weil-type-signature} 已证明
\[
\operatorname{Prym}(C_6/H)
\]
是 \(K=\QQ(\zeta_3)\)-Weil 型阿贝尔八维簇，且其签名为 \((4,4)\)，并具有极化型 \((1^4,3^4)\)。另一方面，定理 \ref{thm:killo-s4-burnside-kani-rosen-prym-square} 给出
\[
\operatorname{Prym}(C_6/H)\sim_{\QQ}J(Z)^2,
\]
而定理 \ref{thm:killo-s4-genus49-jacobian-computable-isogeny} 正是将 \(J(Z)^2\) 识别为 \(J(X)\) 的 \(V_2\)-等型块 \(A_2\)。故所述 Weil 结构、特征分解与极化型全部传递到 \(A_2\)。证毕。
\end{proof}

\begin{conclusion}[\texorpdfstring{$J(C_6)$}{J(C6)} 落在余维至少 \texorpdfstring{$11$}{11} 的特异分解模子簇中]\label{con:xi-time-part9n-c6-special-jacobian-locus}
设
\[
\mathcal D_{5,(4)^2}
:=
\left\{
A\in \mathcal A_{13}:\ 
A\sim B_5\times B_4^2
\text{ for some }B_5\in\mathcal A_5,\ B_4\in\mathcal A_4
\right\}.
\]
则
\[
J(C_6)\in \mathcal D_{5,(4)^2},
\]
并且
\[
\dim \mathcal D_{5,(4)^2}\le 25.
\]
因此，\(J(C_6)\) 在属 \(13\) 的 Jacobian 模子簇
\[
\mathcal J_{13}\subset \mathcal A_{13}
\]
中落在余维至少
\[
\boxed{\;36-25=11\;}
\]
的特异子簇内。换言之，\(C_6\) 并非属 \(13\) 的泛性曲线，而被强制压入一个高度特异的同源分解子簇。
\end{conclusion}
\begin{proof}
由定理 \ref{thm:killo-s4-burnside-kani-rosen-prym-square}，
\[
J(C_6)\sim J(H)\times J(Z)^2,
\]
其中 \(g(H)=5\)、\(g(Z)=4\)。故 \(J(C_6)\) 的确属于 \(\mathcal D_{5,(4)^2}\)。

再看维数。Siegel 模空间满足
\[
\dim \mathcal A_g=\frac{g(g+1)}{2},
\]
故
\[
\dim \mathcal A_5=15,\qquad \dim \mathcal A_4=10.
\]
映射
\[
\mathcal A_5\times \mathcal A_4\longrightarrow \mathcal A_{13},
\qquad
(B_5,B_4)\longmapsto B_5\times B_4^2
\]
的像维数至多为 \(15+10=25\)。允许“\(\sim\)”意义下的同源，只会产生该乘积像的可数个 Hecke 对应平移，因此相应同源子簇的维数仍至多为 \(25\)。于是
\[
\dim \mathcal D_{5,(4)^2}\le 25.
\]

另一方面，属 \(13\) 的 Jacobian 子簇 \(\mathcal J_{13}\) 与曲线模空间 \(\mathcal M_{13}\) 具有相同维数
\[
\dim \mathcal J_{13}=\dim \mathcal M_{13}=3\cdot 13-3=36.
\]
故 \(J(C_6)\) 所在的分解子簇在 \(\mathcal J_{13}\) 中的余维至少为
\[
36-25=11.
\]
证毕。
\end{proof}

\paragraph{bin-fold escort 温度族的信息几何闭合与纯碰撞单位根扭曲的二阶弛豫律}
在 bin-fold escort 分布的末位充分化、温度间 KL 极限以及纯碰撞单位根扭曲的小角主曲线已经建立之后，还可进一步抽出下列若干条彼此衔接的结构结论。

\begin{theorem}[escort 温度族的 Le Cam 二元塌缩]\label{con:xi-time-part9o-escort-blackwell-collapse}
对任意紧区间
\[
I\subset [0,\infty),
\]
记
\[
A_i^{(m)}:=\{w\in X_m:w_m=i\}
\qquad (i\in\{0,1\}),
\]
定义统计实验
\[
\mathcal E_m(I):=\{\pi_{m,q}:q\in I\}\quad\text{定义在 }X_m\text{ 上},
\]
以及两点实验
\[
\mathcal B(I):=\{\mathrm{Bernoulli}(\theta(q)):q\in I\}\quad\text{定义在 }\{0,1\}\text{ 上},
\qquad
\theta(q):=\frac{1}{1+\varphi^{q+1}}.
\]
则
\[
\Delta_{\mathrm{LeCam}}\bigl(\mathcal E_m(I),\mathcal B(I)\bigr)
=
O_I\!\left(\Bigl(\frac{\varphi}{2}\Bigr)^m\right).
\]
换言之，bin-fold 的整个 escort 温度族在统计实验意义下并不渐近逼近一个高维模型，而是以指数精度塌缩为由末位比特承载的一维 Bernoulli 族。
\end{theorem}
\begin{proof}
取 \(Q\ge 0\) 使 \(I\subset[0,Q]\)。定义前向核
\[
K_m(w):=w_m\in\{0,1\},
\]
以及反向核
\[
L_m(i,\cdot):=\mathrm{Unif}\bigl(A_i^{(m)}\bigr)
\qquad (i\in\{0,1\}).
\]
由定理 \ref{thm:killo-fold-bin-escort-onebit-fisher}，对 \(q\in I\) 一致有
\[
\|\pi_{m,q}-L_{m\#}\mathrm{Bernoulli}(\theta_m(q))\|_{\mathrm{TV}}
=
O_I\!\left(\Bigl(\frac{\varphi}{2}\Bigr)^m\right),
\]
其中
\[
\theta_m(q):=
\frac{\varphi^{-q}F_m}{F_{m+1}+\varphi^{-q}F_m}.
\]
另一方面，
\[
K_{m\#}\pi_{m,q}
=
\mathrm{Bernoulli}\!\bigl(\theta_m(q)\bigr)
\]
是精确恒等式。又由
\[
\theta_m(q)
=
\frac{\varphi^{-q}}{F_{m+1}/F_m+\varphi^{-q}},
\qquad
\theta(q)
=
\frac{\varphi^{-q}}{\varphi+\varphi^{-q}},
\]
以及 \(F_{m+1}/F_m=\varphi+O(\varphi^{-2m})\)，可得
\[
\sup_{q\in I}|\theta_m(q)-\theta(q)|
=
O_I\!\left(\Bigl(\frac{\varphi}{2}\Bigr)^m\right).
\]
故
\[
\sup_{q\in I}
\|K_{m\#}\pi_{m,q}-\mathrm{Bernoulli}(\theta(q))\|_{\mathrm{TV}}
=
O_I\!\left(\Bigl(\frac{\varphi}{2}\Bigr)^m\right).
\]
按 deficiency 的定义，前两式分别给出
\[
\delta\!\bigl(\mathcal E_m(I),\mathcal B(I)\bigr)
=
O_I\!\left(\Bigl(\frac{\varphi}{2}\Bigr)^m\right),
\qquad
\delta\!\bigl(\mathcal B(I),\mathcal E_m(I)\bigr)
=
O_I\!\left(\Bigl(\frac{\varphi}{2}\Bigr)^m\right).
\]
取两者最大值即得所示 Le Cam 距离估计。证毕。
\end{proof}

\begin{theorem}[全部 \texorpdfstring{$f$}{f}-散度的二元闭包]\label{con:xi-time-part9o-escort-fdiv-collapse}
设 \(f:(0,\infty)\to\RR\) 为凸函数并满足 \(f(1)=0\)。对任意 \(p,q\ge 0\)，定义
\[
D_f(P\Vert Q):=\sum_x Q(x)\,f\!\Bigl(\frac{P(x)}{Q(x)}\Bigr).
\]
并记
\[
\beta_s:=\mathrm{Bernoulli}(\theta(s)),
\qquad
\theta(s)=\frac{1}{1+\varphi^{s+1}}
\qquad (s\ge 0).
\]
则极限
\[
\lim_{m\to\infty}D_f(\pi_{m,q}\Vert \pi_{m,p})
\]
存在，并满足
\[
\lim_{m\to\infty}D_f(\pi_{m,q}\Vert \pi_{m,p})
=
D_f(\beta_q\Vert \beta_p).
\]
等价地，
\[
\lim_{m\to\infty}D_f(\pi_{m,q}\Vert \pi_{m,p})
=
\bigl(1-\theta(p)\bigr)\,
f\!\Bigl(\frac{1-\theta(q)}{1-\theta(p)}\Bigr)
\;+\;
\theta(p)\,
f\!\Bigl(\frac{\theta(q)}{\theta(p)}\Bigr),
\]
特别地，
\[
\lim_{m\to\infty}D_{\mathrm{TV}}(\pi_{m,q},\pi_{m,p})=|\theta(q)-\theta(p)|,
\]
\[
\lim_{m\to\infty}\sum_{w\in X_m}\sqrt{\pi_{m,q}(w)\pi_{m,p}(w)}
=
\sqrt{\theta(q)\theta(p)}+\sqrt{(1-\theta(q))(1-\theta(p))}.
\]
\end{theorem}
\begin{proof}
记
\[
A_b^{(m)}:=\{w\in X_m:w_m=b\},
\qquad b\in\{0,1\}.
\]
由定理 \ref{thm:fold-bin-two-state-asymptotic} 与命题 \ref{prop:fold-bin-power-sum-asymptotic}，对每个固定参数 \(s\ge 0\) 有一致展开
\[
\pi_{m,s}(w)
=
\frac{\varphi^{-s w_m}}{F_{m+1}+\varphi^{-s}F_m}
\Bigl(1+O\!\left(\Bigl(\frac{\varphi}{2}\Bigr)^m\right)\Bigr),
\qquad w\in X_m.
\]
因此在每个扇区 \(A_b^{(m)}\) 内，比值
\[
\frac{\pi_{m,q}(w)}{\pi_{m,p}(w)}
\]
一致趋于常数；更精确地，
\[
\sup_{w\in A_0^{(m)}}\left|
\frac{\pi_{m,q}(w)}{\pi_{m,p}(w)}
-\frac{1-\theta(q)}{1-\theta(p)}
\right|
\longrightarrow 0,
\]
\[
\sup_{w\in A_1^{(m)}}\left|
\frac{\pi_{m,q}(w)}{\pi_{m,p}(w)}
-\frac{\theta(q)}{\theta(p)}
\right|
\longrightarrow 0.
\]
另一方面，由命题 \ref{prop:fold-bin-escort-lastbit}，
\[
\pi_{m,p}(A_0^{(m)})\to 1-\theta(p),
\qquad
\pi_{m,p}(A_1^{(m)})\to \theta(p).
\]

将 \(D_f(\pi_{m,q}\Vert \pi_{m,p})\) 按两扇区拆开：
\[
D_f(\pi_{m,q}\Vert \pi_{m,p})
=
\sum_{b=0}^1\ \sum_{w\in A_b^{(m)}}
\pi_{m,p}(w)\,
f\!\Bigl(\frac{\pi_{m,q}(w)}{\pi_{m,p}(w)}\Bigr).
\]
由于两扇区上的比值均落入某个与 \(m\) 无关的紧区间，凸函数 \(f\) 在该紧区间上一致连续，故可把每个扇区内的 \(f\)-值替换为其极限常数，误差 \(o(1)\)；再结合扇区总质量的极限，即得所示两点闭式。

特别地，取
\[
f(u)=\frac12|u-1|
\]
便得到全变差极限。再取
\[
f(u)=1-\sqrt u,
\]
则
\[
D_f(P\Vert Q)=1-\sum_x\sqrt{P(x)Q(x)},
\]
故上一极限公式立即化为所示 Bhattacharyya 系数闭式。证毕。
\end{proof}

上述两点模型还可进一步定量化：不仅末位边际收敛到 \(\theta(q)\)，而且整个 escort 温度族在总变差意义下以指数精度塌缩到双块均匀混合；据此，escort--escort 散度的全族闭式也可统一压缩为 Bernoulli 两点模型。

\begin{theorem}[bin-fold escort 族在总变差意义下指数塌缩为二块均匀混合]\label{thm:xi-time-part9od-escort-tv-collapse-block-uniform}
\leanverified{paper\_xi\_time\_part9od\_escort\_tv\_collapse\_block\_uniform}
沿用结论 \ref{con:xi-time-part9l-escort-single-bit-sufficiency} 与推论 \ref{cor:xi-time-part9l-escort-two-block-normal-form} 的记号。记
\[
A_0:=\{w\in X_m:\ w_m=0\},
\qquad
A_1:=\{w\in X_m:\ w_m=1\},
\]
以及
\[
\overline\pi_{m,q}:=(1-\theta(q))u_{A_0}+\theta(q)u_{A_1}.
\]
则对任意固定 \(q\ge 0\)，存在常数 \(C_q>0\) 使得
\[
D_{\mathrm{TV}}(\pi_{m,q},\overline\pi_{m,q})
\le
C_q\Bigl(\frac{\varphi}{2}\Bigr)^m.
\]
特别地，
\[
\pi_{m,q}(A_1)=\theta(q)+O\!\Bigl(\Bigl(\frac{\varphi}{2}\Bigr)^m\Bigr),
\qquad
\pi_{m,q}(A_0)=1-\theta(q)+O\!\Bigl(\Bigl(\frac{\varphi}{2}\Bigr)^m\Bigr).
\]
因此，对每个固定 \(q\)，整个 escort 温度族在大尺度上以指数精度塌缩为由末位坐标 \(w_m\) 所决定的二元统计实验。
\end{theorem}
\begin{proof}
记
\[
\theta_m(q):=\pi_{m,q}(A_1).
\]
由结论 \ref{con:xi-time-part9l-escort-single-bit-sufficiency}，存在常数 \(C_{1,q}>0\) 使
\[
D_{\mathrm{TV}}\!\Bigl(
\pi_{m,q},
(1-\theta_m(q))u_{A_0}+\theta_m(q)u_{A_1}
\Bigr)
\le
C_{1,q}\Bigl(\frac{\varphi}{2}\Bigr)^m.
\]
另一方面，由命题 \ref{prop:fold-bin-escort-lastbit}，存在常数 \(C_{2,q}>0\) 使
\[
|\theta_m(q)-\theta(q)|
\le
C_{2,q}\Bigl(\frac{\varphi}{2}\Bigr)^m.
\]
由于 \(u_{A_0}\) 与 \(u_{A_1}\) 的支撑互不相交，
\[
D_{\mathrm{TV}}\!\Bigl(
(1-\theta_m(q))u_{A_0}+\theta_m(q)u_{A_1},
\overline\pi_{m,q}
\Bigr)
=
|\theta_m(q)-\theta(q)|.
\]
由三角不等式，
\[
D_{\mathrm{TV}}(\pi_{m,q},\overline\pi_{m,q})
\le
(C_{1,q}+C_{2,q})\Bigl(\frac{\varphi}{2}\Bigr)^m,
\]
取 \(C_q:=C_{1,q}+C_{2,q}\) 即得所示估计。块质量公式则是
\(\pi_{m,q}(A_1)=\theta_m(q)\) 与上式的直接重述。证毕。
\end{proof}

\begin{proposition}[bin-fold escort 族按真实块权重的块内均匀化]\label{prop:xi-time-part9od-escort-tv-collapse-exact-block-weights}
沿用定理 \ref{thm:xi-time-part9od-escort-tv-collapse-block-uniform} 的记号。对任意固定 \(q\ge 0\)，定义精确块权重
\[
p_{m,q}(i):=\pi_{m,q}(A_i),
\qquad i\in\{0,1\},
\]
并记
\[
\nu_{m,q}:=
p_{m,q}(0)\,u_{A_0}+p_{m,q}(1)\,u_{A_1}.
\]
则存在常数 \(C_q'>0\)，使得
\[
D_{\mathrm{TV}}(\pi_{m,q},\nu_{m,q})
\le
C_q'\Bigl(\frac{\varphi}{2}\Bigr)^m.
\]
换言之，escort 分布在全空间层面不仅把末位边际压缩为两点模型，而且以指数精度退化为“按末位分块后块内完全均匀”的结构。
\leanverified{paper\_xi\_time\_part9od\_escort\_tv\_collapse\_exact\_block\_weights}
\end{proposition}
\begin{proof}
在定理 \ref{thm:xi-time-part9od-escort-tv-collapse-block-uniform} 的证明中，已得到
\[
D_{\mathrm{TV}}\!\Bigl(
\pi_{m,q},
(1-\theta_m(q))u_{A_0}+\theta_m(q)u_{A_1}
\Bigr)
\le
C_{1,q}\Bigl(\frac{\varphi}{2}\Bigr)^m,
\]
其中
\[
\theta_m(q)=\pi_{m,q}(A_1)=p_{m,q}(1),
\qquad
1-\theta_m(q)=\pi_{m,q}(A_0)=p_{m,q}(0).
\]
因此右端的双块均匀混合恰为 \(\nu_{m,q}\) 本身，所需估计立刻成立。证毕。
\end{proof}

\begin{theorem}[escort--escort 散度族的统一二元化与 Rényi--KL 闭式极限]\label{thm:xi-time-part9od-escort-escort-fdiv-binary-closure}
对任意 \(p,q\ge 0\)，记
\[
\beta_q:=\bigl(1-\theta(q),\theta(q)\bigr)
\qquad\text{定义在 }\{0,1\}\text{ 上}.
\]
若 \(f:(0,\infty)\to\RR\) 在有限集合
\[
\left\{
\frac{1-\theta(q)}{1-\theta(p)},
\frac{\theta(q)}{\theta(p)}
\right\}
\]
的某个开邻域上连续且凸，并满足 \(f(1)=0\)，则
\[
\lim_{m\to\infty}D_f(\pi_{m,q}\Vert \pi_{m,p})
=
D_f(\beta_q\Vert \beta_p)
\]
\[
=
\bigl(1-\theta(p)\bigr)
f\!\Bigl(\frac{1-\theta(q)}{1-\theta(p)}\Bigr)
\;+\;
\theta(p)\,
f\!\Bigl(\frac{\theta(q)}{\theta(p)}\Bigr).
\]
特别地，对任意 \(\alpha>0\)、\(\alpha\neq 1\)，Rényi 散度满足
\[
\lim_{m\to\infty}D_\alpha(\pi_{m,q}\Vert \pi_{m,p})
=
\frac{1}{\alpha-1}
\log
\frac{\varphi^{\alpha(q+1)+(1-\alpha)(p+1)}+1}
{(1+\varphi^{q+1})^\alpha(1+\varphi^{p+1})^{1-\alpha}}.
\]
而 Kullback--Leibler 散度满足
\[
\lim_{m\to\infty}D_{\mathrm{KL}}(\pi_{m,q}\Vert \pi_{m,p})
=
\bigl(1-\theta(q)\bigr)\log\frac{1-\theta(q)}{1-\theta(p)}
+\theta(q)\log\frac{\theta(q)}{\theta(p)}.
\]
\end{theorem}
\begin{proof}
第一条正是结论 \ref{con:xi-time-part9o-escort-fdiv-collapse} 在两点模型 \(\beta_q\) 的记号下的改写。

对 Rényi 散度，二点分布的标准公式给出
\[
D_\alpha(\beta_q\Vert \beta_p)
=
\frac{1}{\alpha-1}
\log\!\Bigl(
(1-\theta(p))^{1-\alpha}(1-\theta(q))^\alpha
+\theta(p)^{1-\alpha}\theta(q)^\alpha
\Bigr).
\]
再代入
\[
\theta(s)=\frac{1}{1+\varphi^{s+1}},
\qquad
1-\theta(s)=\frac{\varphi^{s+1}}{1+\varphi^{s+1}},
\]
直接化简便得到所示闭式。

对 KL 散度，只需在同一两点模型上取
\[
f(u)=u\log u,
\]
即可得到
\[
D_{\mathrm{KL}}(\beta_q\Vert \beta_p)
=
\bigl(1-\theta(q)\bigr)\log\frac{1-\theta(q)}{1-\theta(p)}
+\theta(q)\log\frac{\theta(q)}{\theta(p)}.
\]
再与第一条合并便得结论。证毕。
\end{proof}
\leanverified{paper\_xi\_time\_part9od\_escort\_escort\_fdiv\_binary\_closure}

\endinput


\begin{corollary}[escort 温度间判别几何的显式二元闭式]\label{cor:xi-time-part9o-escort-testing-geometry-closed}
对任意 \(p,q\ge 0\)，记
\[
\theta_p:=\theta(p),
\qquad
\theta_q:=\theta(q).
\]
则有
\[
\lim_{m\to\infty}D_{\mathrm{TV}}(\pi_{m,q},\pi_{m,p})
=
|\theta_q-\theta_p|,
\]
\[
\lim_{m\to\infty}\sum_{w\in X_m}\sqrt{\pi_{m,q}(w)\pi_{m,p}(w)}
=
\sqrt{\theta_q\theta_p}+\sqrt{(1-\theta_q)(1-\theta_p)},
\]
\[
\lim_{m\to\infty}H^2(\pi_{m,q},\pi_{m,p})
=
1-\sqrt{\theta_q\theta_p}-\sqrt{(1-\theta_q)(1-\theta_p)},
\]
其中平方 Hellinger 距离定义为
\[
H^2(P,Q):=1-\sum_x\sqrt{P(x)Q(x)}.
\]
若两假设取等先验，则相应 Bayes 最小判别误差满足
\[
\lim_{m\to\infty}P_e^\ast(\pi_{m,q},\pi_{m,p})
=
\frac12\bigl(1-|\theta_q-\theta_p|\bigr).
\]
因此 escort 温度流在极限中完全退化为单参数 Bernoulli 判别几何，其最优测试、Bhattacharyya 系数、Hellinger 结构与总变差距离都由 logistic 坐标 \(\theta(q)\) 显式支配。
\end{corollary}
\begin{proof}
全变差极限已包含在结论 \ref{con:xi-time-part9o-escort-fdiv-collapse} 中。
对平方 Hellinger 距离，在同一结论中取
\[
f(u):=\frac12(\sqrt u-1)^2,
\]
便得到两点闭式
\[
H^2(\pi_{m,q},\pi_{m,p})
\longrightarrow
\frac12\left[
(1-\theta_p)\left(\sqrt{\frac{1-\theta_q}{1-\theta_p}}-1\right)^2
+
\theta_p\left(\sqrt{\frac{\theta_q}{\theta_p}}-1\right)^2
\right],
\]
化简即得
\[
1-\sqrt{\theta_q\theta_p}-\sqrt{(1-\theta_q)(1-\theta_p)}.
\]
等价地，Bhattacharyya 系数满足
\[
\sum_{w\in X_m}\sqrt{\pi_{m,q}(w)\pi_{m,p}(w)}
\longrightarrow
\sqrt{\theta_q\theta_p}+\sqrt{(1-\theta_q)(1-\theta_p)}.
\]
最后，对等先验简单假设检验，Helstrom 公式给出
\[
P_e^\ast(P,Q)=\frac12\bigl(1-D_{\mathrm{TV}}(P,Q)\bigr),
\]
将 \(P=\pi_{m,q}\)、\(Q=\pi_{m,p}\) 并令 \(m\to\infty\)，即得最后一个极限公式。证毕。
\end{proof}

\begin{theorem}[escort 缩放残差两点律的矩闭式与单矩反演]\label{thm:xi-time-part9o-escort-residual-moment-identifiability}
\leanverified{paper\_xi\_time\_part9o\_escort\_residual\_moment\_identifiability}
对任意 \(q\ge 0\) 令 \(W_m^{(q)}\sim \pi_{m,q}\)，并定义缩放残差变量
\[
Z_m^{(q)}:=\left(\frac{\varphi}{2}\right)^m d_m^{\mathrm{bin}}(W_m^{(q)}).
\]
则对每个整数 \(r\ge 1\)，极限矩
\[
M_r(q):=\lim_{m\to\infty}\mathbf E_{\pi_{m,q}}\!\left[(Z_m^{(q)})^r\right]
\]
存在且满足闭式
\[
\boxed{\;
M_r(q)=1-\bigl(1-\varphi^{-r}\bigr)\theta(q),
\qquad
\theta(q)=\frac{1}{1+\varphi^{q+1}}.\;}
\]
并且对固定的 \(r\ge 1\)，该矩可显式反演温度参数：
\[
\boxed{\;
q=\log_\varphi\!\left(\frac{M_r(q)-\varphi^{-r}}{1-M_r(q)}\right)-1.\;}
\]
特别地，任意一个非平凡极限矩都唯一确定 \(q\)，因而唯一确定 escort 温度族的全部极限 Bernoulli 判别几何。
\end{theorem}
\begin{proof}
令
\[
\varepsilon_m:=\left(\frac{\varphi}{2}\right)^m.
\]
由定理 \ref{thm:fold-bin-two-state-asymptotic}，在 \(X_m\) 上一致有
\[
Z_m^{(q)}=\varphi^{-W_m^{(q)}(m)}\bigl(1+O(\varepsilon_m)\bigr).
\]
故对固定 \(r\ge 1\)，仍一致有
\[
\bigl(Z_m^{(q)}\bigr)^r
=
\varphi^{-rW_m^{(q)}(m)}\bigl(1+O(\varepsilon_m)\bigr).
\]
另一方面，由命题 \ref{prop:fold-bin-escort-lastbit}，
\[
\mathbf P_{\pi_{m,q}}\bigl(W_m^{(q)}(m)=1\bigr)=\theta(q)+O(\varepsilon_m),
\]
\[
\mathbf P_{\pi_{m,q}}\bigl(W_m^{(q)}(m)=0\bigr)=1-\theta(q)+O(\varepsilon_m).
\]
于是
\[
\mathbf E_{\pi_{m,q}}\!\left[(Z_m^{(q)})^r\right]
=
\bigl(1-\theta(q)\bigr)+\theta(q)\varphi^{-r}+O(\varepsilon_m),
\]
令 \(m\to\infty\) 即得
\[
M_r(q)=1-\bigl(1-\varphi^{-r}\bigr)\theta(q).
\]

再由
\[
\theta(q)=\frac{1}{1+\varphi^{q+1}}
\]
可得
\[
\varphi^{q+1}=\frac{1-\theta(q)}{\theta(q)}.
\]
而前一公式又给出
\[
\theta(q)=\frac{1-M_r(q)}{1-\varphi^{-r}},
\]
故
\[
\frac{1-\theta(q)}{\theta(q)}
=
\frac{M_r(q)-\varphi^{-r}}{1-M_r(q)}.
\]
两边取 \(\log_\varphi\) 即得所示反演公式。最后，由 \(q\mapsto \theta(q)\) 的严格递减性与 \(1-\varphi^{-r}>0\)，可知 \(q\mapsto M_r(q)\) 严格单调，从而单个极限矩确能唯一确定 \(q\)。证毕。
\end{proof}

\begin{corollary}[escort 温度可由纤维信息的 \texorpdfstring{$O(1)$}{O(1)} 常数项唯一反演]\label{cor:xi-time-part9o-escort-fiberinfo-constantterm-inversion}
\leanverified{paper\_xi\_time\_part9o\_escort\_fiberinfo\_constantterm\_inversion}
沿用命题 \ref{prop:fold-bin-escort-lastbit} 的记号。定义
\[
c(q)
:=
\lim_{m\to\infty}
\left(
m\log\!\frac{2}{\varphi}
-\kappa_m^{\mathrm{bin}}(q)
\right).
\]
则极限存在，且满足显式闭式
\[
\boxed{\;
c(q)
=
\theta(q)\log\varphi
=
\frac{\log\varphi}{1+\varphi^{q+1}}.\;}
\]
特别地，映射 \(q\mapsto c(q)\) 在 \([0,\infty)\) 上严格递减，因而可逆；其反演公式为
\[
\boxed{\;
q
=
\log_\varphi\!\left(\frac{\log\varphi}{c(q)}-1\right)-1.\;}
\]
因此，在 bin-fold 的 escort 温度族中，单个 \(\,O(1)\,\) 级纤维信息常数项已足以唯一恢复温度参数 \(q\)。
\end{corollary}
\begin{proof}
由命题 \ref{prop:fold-bin-escort-lastbit}，
\[
\kappa_m^{\mathrm{bin}}(q)
=
m\log\!\Bigl(\frac{2}{\varphi}\Bigr)-\theta(q)\log\varphi
+O\!\left(\Bigl(\frac{\varphi}{2}\Bigr)^m\right),
\qquad
\theta(q)=\frac{1}{1+\varphi^{q+1}}.
\]
将主阶项移到左端并令 \(m\to\infty\)，便得到
\[
c(q)=\theta(q)\log\varphi=\frac{\log\varphi}{1+\varphi^{q+1}}.
\]
又因 \(q\mapsto \varphi^{q+1}\) 严格递增，故 \(q\mapsto c(q)\) 严格递减，从而可逆。

最后解代数方程
\[
c=\frac{\log\varphi}{1+\varphi^{q+1}}
\]
即得
\[
1+\varphi^{q+1}=\frac{\log\varphi}{c},
\qquad
\varphi^{q+1}=\frac{\log\varphi}{c}-1,
\]
再取 \(\log_\varphi\) 便得到所示反演公式。证毕。
\end{proof}

\endinput


\begin{theorem}[bin-fold escort 温度族在指数几何下的渐近闭合与仿射自然参数律]\label{thm:xi-time-part9o-escort-egeodesic-affine-closure}
\leanverified{paper\_xi\_time\_part9o\_escort\_egeodesic\_affine\_closure}
沿用结论 \ref{con:xi-time-part9o-escort-fdiv-collapse} 的记号。对任意 \(p,q\ge 0\) 与 \(\lambda\in[0,1]\)，定义归一化 pointwise geometric mean
\[
\Gamma_{m,\lambda}^{p,q}(w)
:=
\frac{\pi_{m,p}(w)^\lambda\pi_{m,q}(w)^{1-\lambda}}
{\sum_{u\in X_m}\pi_{m,p}(u)^\lambda\pi_{m,q}(u)^{1-\lambda}},
\qquad
w\in X_m,
\]
并记
\[
r:=\lambda p+(1-\lambda)q.
\]
则有精确的指数几何闭合律
\[
\boxed{\;
D_{\mathrm{TV}}\!\bigl(\Gamma_{m,\lambda}^{p,q},\pi_{m,r}\bigr)
=
O_{p,q,\lambda}\!\Bigl(\Bigl(\frac{\varphi}{2}\Bigr)^m\Bigr).\;}
\]
换言之，bin-fold escort 温度族虽然在有限 \(m\) 时仍带有指数小残差，但在主阶上已经形成一条真正的 \((e)\)-测地闭合曲线：归一化几何平均所对应的温度参数恰是仿射组合
\[
r=\lambda p+(1-\lambda)q.
\]
等价地，若定义极限自然参数
\[
\eta(q):=-q\log\varphi,
\]
则有精确仿射加法律
\[
\boxed{\;
\eta(r)=\lambda\eta(p)+(1-\lambda)\eta(q).\;}
\]
\end{theorem}
\begin{proof}
由定理 \ref{thm:fold-bin-two-state-asymptotic} 与命题 \ref{prop:fold-bin-power-sum-asymptotic}，对每个固定 \(s\ge 0\) 都有一致展开
\[
\pi_{m,s}(w)
=
\frac{\varphi^{-sw_m}}{F_{m+1}+\varphi^{-s}F_m}
\Bigl(1+\varepsilon_{m,s}(w)\Bigr),
\qquad
\sup_{w\in X_m}|\varepsilon_{m,s}(w)|
=
O_s\!\Bigl(\Bigl(\frac{\varphi}{2}\Bigr)^m\Bigr).
\]
将 \(s=p,q\) 代入并相乘，得到
\[
\pi_{m,p}(w)^\lambda\pi_{m,q}(w)^{1-\lambda}
=
\frac{\varphi^{-rw_m}}
{\bigl(F_{m+1}+\varphi^{-p}F_m\bigr)^\lambda
\bigl(F_{m+1}+\varphi^{-q}F_m\bigr)^{1-\lambda}}
\Bigl(1+O_{p,q,\lambda}\!\Bigl(\Bigl(\frac{\varphi}{2}\Bigr)^m\Bigr)\Bigr),
\]
其中误差在 \(w\in X_m\) 上一致。对 \(w\) 求和后，分母化为
\[
\frac{F_{m+1}+\varphi^{-r}F_m}
{\bigl(F_{m+1}+\varphi^{-p}F_m\bigr)^\lambda
\bigl(F_{m+1}+\varphi^{-q}F_m\bigr)^{1-\lambda}}
\Bigl(1+O_{p,q,\lambda}\!\Bigl(\Bigl(\frac{\varphi}{2}\Bigr)^m\Bigr)\Bigr).
\]
因此
\[
\Gamma_{m,\lambda}^{p,q}(w)
=
\frac{\varphi^{-rw_m}}{F_{m+1}+\varphi^{-r}F_m}
\Bigl(1+O_{p,q,\lambda}\!\Bigl(\Bigl(\frac{\varphi}{2}\Bigr)^m\Bigr)\Bigr),
\]
并且该误差仍在 \(w\in X_m\) 上一致。

另一方面，对参数 \(r\) 应用同一展开，得到
\[
\pi_{m,r}(w)
=
\frac{\varphi^{-rw_m}}{F_{m+1}+\varphi^{-r}F_m}
\Bigl(1+O_r\!\Bigl(\Bigl(\frac{\varphi}{2}\Bigr)^m\Bigr)\Bigr).
\]
两式相减，并对 \(w\in X_m\) 求和，即得
\[
D_{\mathrm{TV}}\!\bigl(\Gamma_{m,\lambda}^{p,q},\pi_{m,r}\bigr)
=
O_{p,q,\lambda}\!\Bigl(\Bigl(\frac{\varphi}{2}\Bigr)^m\Bigr).
\]
最后，
\[
\eta(r)=-(\lambda p+(1-\lambda)q)\log\varphi
=
\lambda\eta(p)+(1-\lambda)\eta(q),
\]
故自然参数满足所述仿射加法律。证毕。
\end{proof}

\endinput


\begin{conclusion}[escort 温度线的极限 Fisher--Rao 几何完全闭式]\label{con:xi-time-part9o-escort-fisher-rao-closed}
定义 escort 温度参数直线
\[
q\longmapsto \pi_{m,q},
\qquad q\ge 0.
\]
则其极限 Fisher 信息存在，并满足
\[
\mathcal I(q)
:=
\lim_{\substack{h\to 0\\ q+h\ge 0}}\lim_{m\to\infty}\frac{2}{h^2}
D_{\mathrm{KL}}(\pi_{m,q+h}\Vert \pi_{m,q})
=
(\log\varphi)^2\,\theta(q)\bigl(1-\theta(q)\bigr),
\]
其中
\[
\theta(q)=\frac{1}{1+\varphi^{q+1}}.
\]
因此极限 Fisher--Rao 度量为
\[
\mathrm ds^2
=
(\log\varphi)^2\,\theta(q)\bigl(1-\theta(q)\bigr)\,\mathrm dq^2.
\]
相应测地距离具有完全闭式
\[
d_{\mathrm{FR}}(q_1,q_2)
=
2\left|
\arcsin\sqrt{\theta(q_2)}-\arcsin\sqrt{\theta(q_1)}
\right|
\]
\[
=
2\left|
\arctan\!\bigl(\varphi^{-(q_2+1)/2}\bigr)
\;-\;
\arctan\!\bigl(\varphi^{-(q_1+1)/2}\bigr)
\right|.
\]
\end{conclusion}
\begin{proof}
由结论 \ref{con:xi-time-part9o-escort-fdiv-collapse} 在
\[
f(u)=u\log u
\]
处的专化，
\[
\lim_{m\to\infty}D_{\mathrm{KL}}(\pi_{m,q+h}\Vert \pi_{m,q})
=
D_{\mathrm{KL}}\!\bigl(\mathrm{Bernoulli}(\theta(q+h))\Vert \mathrm{Bernoulli}(\theta(q))\bigr).
\]
因此 \(\mathcal I(q)\) 只需求取右端 Bernoulli KL 在 \(h=0\) 的二阶项。对任意光滑两点族 \(q\mapsto \mathrm{Bernoulli}(\theta(q))\)，Fisher 信息为
\[
\mathcal I(q)=\frac{(\theta'(q))^2}{\theta(q)(1-\theta(q))}.
\]
而
\[
\theta'(q)=-(\log\varphi)\,\theta(q)(1-\theta(q)),
\]
故
\[
\mathcal I(q)=
(\log\varphi)^2\,\theta(q)(1-\theta(q)).
\]

于是 Fisher--Rao 测地距离为
\[
d_{\mathrm{FR}}(q_1,q_2)
=
\left|
\int_{q_1}^{q_2}\sqrt{\mathcal I(q)}\,\mathrm dq
\right|
=
|\log\varphi|
\left|
\int_{q_1}^{q_2}\sqrt{\theta(q)(1-\theta(q))}\,\mathrm dq
\right|.
\]
再由
\[
\frac{\mathrm d\theta}{\mathrm dq}
=
-(\log\varphi)\,\theta(1-\theta)
\]
可得
\[
\sqrt{\mathcal I(q)}\,\mathrm dq
=
-\,\frac{\mathrm d\theta}{\sqrt{\theta(1-\theta)}}.
\]
积分即得
\[
d_{\mathrm{FR}}(q_1,q_2)
=
2\left|
\arcsin\sqrt{\theta(q_2)}-\arcsin\sqrt{\theta(q_1)}
\right|.
\]
最后设
\[
t(q):=\varphi^{-(q+1)/2},
\]
则
\[
\theta(q)=\frac{1}{1+\varphi^{q+1}}=\frac{t(q)^2}{1+t(q)^2},
\]
故
\[
\arcsin\sqrt{\theta(q)}=\arctan t(q),
\]
从而得到第二个闭式表达。证毕。
\end{proof}

\paragraph{bin-fold escort 温度流形的一维 Bernoulli 塌缩、有限弧长相图与全阶累积量闭包}
在末位充分化、两点 \texorpdfstring{$f$}{f}-散度闭式与 Fisher--Rao 距离都已确定之后，bin-fold 的 escort 温度族还可进一步被压缩为一条真正的一维极限统计流形：其尾分位阈值与高阶涨落都由同一 Bernoulli 末位律完全锁定。

\begin{theorem}[末位比特把 escort 温度族识别为一维 Bernoulli \texorpdfstring{$f$}{f}-散度流形]\label{thm:xi-time-part9ob-escort-lastbit-bernoulli-fdiv}
沿用结论 \ref{con:xi-time-part9o-escort-fdiv-collapse} 与推论 \ref{cor:xi-time-part9l-escort-two-block-normal-form} 的记号，并记
\[
\ell(w):=w_m,
\qquad
\theta(q):=\frac{1}{1+\varphi^{q+1}}.
\]
则对任意 \(p,q\ge 0\) 与任意凸函数 \(f:(0,\infty)\to \mathbb R\)（满足 \(f(1)=0\)），都有
\[
\lim_{m\to\infty}D_f(\pi_{m,q}\Vert \pi_{m,p})
=
\bigl(1-\theta(p)\bigr)\,
f\!\Bigl(\frac{1-\theta(q)}{1-\theta(p)}\Bigr)
\;+\;
\theta(p)\,
f\!\Bigl(\frac{\theta(q)}{\theta(p)}\Bigr).
\]
换言之，整个 escort 温度族在全部 Csisz\'ar--Morimoto \(f\)-散度意义下都塌缩到一维 Bernoulli 统计流形
\[
q\longmapsto \mathrm{Bernoulli}(\theta(q)),
\]
而末位比特 \(\ell\) 正是该族的极限充分统计量。
\end{theorem}
\begin{proof}
上述 \(f\)-散度闭式正是结论 \ref{con:xi-time-part9o-escort-fdiv-collapse} 的内容。

另一方面，推论 \ref{cor:xi-time-part9l-escort-two-block-normal-form} 给出
\[
\pi_{m,q}
=
\bigl(1-\theta(q)\bigr)u_{A_0}
\;+\;
\theta(q)u_{A_1}
\;+\;
\varepsilon_{m,q},
\]
其中
\[
A_b:=\{w\in X_m:\ w_m=b\},
\qquad
u_{A_b}:=U_{m,b},
\qquad
\|\varepsilon_{m,q}\|_{\mathrm{TV}}
=
O\!\left(\Bigl(\frac{\varphi}{2}\Bigr)^m\right).
\]
因此，对每个固定 \(q\ge 0\)，\(\pi_{m,q}\) 以指数小误差退化为由末位读出 \(\ell\) 所控制的两块均匀混合，而该混合正是 \(\mathrm{Bernoulli}(\theta(q))\) 经提升核 \(L_m(b,\cdot)=u_{A_b}\) 的输出。故末位比特 \(\ell\) 在极限中成为充分统计量；再与前述 \(f\)-散度闭式合并，即得所述一维 Bernoulli 统计流形解释。证毕。
\end{proof}

\begin{theorem}[escort 温度流形的 Fisher--Rao 完整几何]\label{thm:xi-time-part9ob-escort-fisher-total-length}
沿用结论 \ref{con:xi-time-part9o-escort-fisher-rao-closed} 的记号。定义极限 Fisher 信息密度
\[
\mathcal I_\infty(q)
:=
\lim_{p\to q}
\frac{2\,D_{\mathrm{KL}}(\mathrm{Bernoulli}(\theta(q))\Vert \mathrm{Bernoulli}(\theta(p)))}{(q-p)^2}.
\]
则
\[
\boxed{\;
\mathcal I_\infty(q)
=
(\log\varphi)^2\,\theta(q)\bigl(1-\theta(q)\bigr)
=
(\log\varphi)^2\frac{\varphi^{q+1}}{(1+\varphi^{q+1})^2}.\;}
\]
相应的 Fisher--Rao 度量可写成
\[
\mathrm ds^2=\mathcal I_\infty(q)\,\mathrm dq^2,
\]
\[
d_{\mathrm{FR}}(q_1,q_2)
=
2\left|
\arcsin\sqrt{\theta(q_2)}-\arcsin\sqrt{\theta(q_1)}
\right|.
\]
进一步，整条温度射线 \(q\in[0,\infty)\) 的总信息长度有限，并且
\[
\mathcal L_{[0,\infty)}
:=
\int_0^\infty \sqrt{\mathcal I_\infty(q)}\,\mathrm dq
=
2\arctan(\varphi^{-1/2})
=
2\arcsin(\varphi^{-1}).
\]
\leanpartial{paper\_xi\_time\_part9ob\_escort\_fisher\_total\_length}{The requested Lean signature quantifies over arbitrary functions and constants, so \texttt{paper\_xi\_time\_part9ob\_escort\_fisher\_total\_length} is not derivable without additional hypotheses.}
\end{theorem}
\begin{proof}
由结论 \ref{con:xi-time-part9o-escort-fisher-rao-closed}，Bernoulli 极限模型的 Fisher--Rao 度量正由
\[
\frac{(\theta'(q))^2}{\theta(q)(1-\theta(q))}
\]
给出，因此
\[
\mathcal I_\infty(q)
=
\frac{(\theta'(q))^2}{\theta(q)(1-\theta(q))}.
\]
再由
\[
\theta(q)=\frac{1}{1+\varphi^{q+1}}
\]
直接求导可得
\[
\theta'(q)
=
-(\log\varphi)\,\theta(q)\bigl(1-\theta(q)\bigr),
\]
故
\[
\mathcal I_\infty(q)
=(\log\varphi)^2\,\theta(q)\bigl(1-\theta(q)\bigr)
=(\log\varphi)^2\frac{\varphi^{q+1}}{(1+\varphi^{q+1})^2}.
\]
这就给出 Fisher 信息密度闭式。测地距离公式则与结论 \ref{con:xi-time-part9o-escort-fisher-rao-closed} 完全一致。

对总长度，令
\[
a:=\varphi^{q+1},
\qquad
\mathrm dq=\frac{\mathrm da}{a\log\varphi}.
\]
于是
\[
\sqrt{\mathcal I_\infty(q)}
=(\log\varphi)\frac{\sqrt a}{1+a},
\]
从而
\[
\mathcal L_{[0,\infty)}
=
\int_0^\infty \sqrt{\mathcal I_\infty(q)}\,\mathrm dq
=
\int_{\varphi}^{\infty}\frac{a^{-1/2}}{1+a}\,\mathrm da.
\]
再令 \(a=t^2\)，则 \(\mathrm da=2t\,\mathrm dt\)，故
\[
\mathcal L_{[0,\infty)}
=
2\int_{\sqrt\varphi}^{\infty}\frac{1}{1+t^2}\,\mathrm dt
=
\pi-2\arctan(\sqrt\varphi).
\]
由恒等式
\[
\arctan x+\arctan(x^{-1})=\frac{\pi}{2}
\qquad (x>0)
\]
可得
\[
\pi-2\arctan(\sqrt\varphi)=2\arctan(\varphi^{-1/2}).
\]
另一方面，由
\[
\sin\!\bigl(2\arctan(\varphi^{-1/2})\bigr)
=
\frac{2\varphi^{-1/2}}{1+\varphi^{-1}}
=
\varphi^{-1},
\]
且 \(2\arctan(\varphi^{-1/2})\in(0,\pi/2)\)，故
\[
2\arctan(\varphi^{-1/2})=2\arcsin(\varphi^{-1}).
\]
证毕。
\end{proof}

\begin{theorem}[escort 尾分位阈值的完整相图由显式相变曲线控制]\label{thm:xi-time-part9ob-escort-quantile-phase-diagram}
对任意 \(q\ge 0\)，令
\[
W_m^{(q)}\sim \pi_{m,q},
\]
并定义尾部分位指数
\[
\gamma_{m,q}^{\mathrm{bin}}(\varepsilon)
:=
\inf\left\{
\gamma:\ 
\mathbf P_{\pi_{m,q}}\!\left(
\frac{1}{m}\log d_m^{\mathrm{bin}}(W_m^{(q)})>\gamma
\right)\le \varepsilon
\right\},
\qquad
0<\varepsilon<1,
\]
以及相应阈值
\[
B_{m,q}^{\mathrm{bin}}(\varepsilon)
:=
\exp\!\bigl(m\gamma_{m,q}^{\mathrm{bin}}(\varepsilon)\bigr).
\]
再定义显式相变曲线
\[
\varepsilon_{\mathrm c}(q)
:=
1-\theta(q)
=
\frac{\varphi^{q+1}}{1+\varphi^{q+1}}.
\]
则对任意固定 \(q\ge 0\) 与任意
\[
\varepsilon\in(0,1)\setminus\{\varepsilon_{\mathrm c}(q)\},
\]
都有
\[
\boxed{\;
\left(\frac{\varphi}{2}\right)^m B_{m,q}^{\mathrm{bin}}(\varepsilon)
\longrightarrow
\begin{cases}
1, & 0<\varepsilon<\varepsilon_{\mathrm c}(q),\\[4pt]
\varphi^{-1}, & \varepsilon_{\mathrm c}(q)<\varepsilon<1.
\end{cases}
\;}
\]
此外，相变曲线完全显式，且严格递增：
\[
\varepsilon_{\mathrm c}'(q)
=
\frac{\log\varphi\ \varphi^{q+1}}{(1+\varphi^{q+1})^2}
>
0.
\]
当 \(q=1\) 时，上式退化为推论 \ref{cor:fold-bin-quantile-threshold-constant} 的原始 bin-fold 阈值常数律。
\leanpartial{paper\_xi\_time\_part9ob\_escort\_quantile\_phase\_diagram}{The requested Lean signature quantifies over arbitrary functions and constants, so \texttt{paper\_xi\_time\_part9ob\_escort\_quantile\_phase\_diagram} is not derivable without additional hypotheses.}
\end{theorem}
\begin{proof}
定义缩放随机变量
\[
Z_m^{(q)}
:=
\left(\frac{\varphi}{2}\right)^m d_m^{\mathrm{bin}}(W_m^{(q)}).
\]
由定理 \ref{thm:fold-bin-two-state-asymptotic} 的一致展开，
\[
Z_m^{(q)}
=
\varphi^{-W_m^{(q)}(m)}
\left(1+O\!\left(\Bigl(\frac{\varphi}{2}\Bigr)^m\right)\right),
\]
其中误差在 \(X_m\) 上一致。另一方面，由命题 \ref{prop:fold-bin-escort-lastbit}，
\[
\mathbf P_{\pi_{m,q}}\bigl(W_m^{(q)}(m)=1\bigr)
=
\theta(q)+O\!\left(\Bigl(\frac{\varphi}{2}\Bigr)^m\right),
\]
\[
\mathbf P_{\pi_{m,q}}\bigl(W_m^{(q)}(m)=0\bigr)
=
1-\theta(q)+O\!\left(\Bigl(\frac{\varphi}{2}\Bigr)^m\right).
\]
故
\[
Z_m^{(q)}\Longrightarrow Z^{(q)},
\]
其中
\[
\mathbf P\bigl(Z^{(q)}=1\bigr)=1-\theta(q)=\varepsilon_{\mathrm c}(q),
\qquad
\mathbf P\bigl(Z^{(q)}=\varphi^{-1}\bigr)=\theta(q).
\]

而
\[
\left(\frac{\varphi}{2}\right)^m B_{m,q}^{\mathrm{bin}}(\varepsilon)
\]
正是随机变量 \(Z_m^{(q)}\) 在同一尾部分位定义下的阈值。由于
\[
\varphi^{-1}<1,
\]
极限两点分布的上尾量在 \(\varepsilon=\varepsilon_{\mathrm c}(q)\) 处发生唯一跳变：若 \(\varepsilon<\varepsilon_{\mathrm c}(q)\)，则必须把原子 \(1\) 一并压掉，故极限阈值为 \(1\)；若 \(\varepsilon>\varepsilon_{\mathrm c}(q)\)，则容许丢弃原子 \(1\) 的全部质量，故极限阈值降为 \(\varphi^{-1}\)。这便得到所示分段极限。

最后，
\[
\varepsilon_{\mathrm c}(q)=\frac{\varphi^{q+1}}{1+\varphi^{q+1}}
\]
直接求导即得
\[
\varepsilon_{\mathrm c}'(q)
=
\frac{\log\varphi\ \varphi^{q+1}}{(1+\varphi^{q+1})^2}>0.
\]
当 \(q=1\) 时，
\[
\varepsilon_{\mathrm c}(1)
=
\frac{\varphi^2}{1+\varphi^2}
=
\frac{\varphi}{\sqrt5},
\]
故恰好恢复推论 \ref{cor:fold-bin-quantile-threshold-constant}。证毕。
\end{proof}

\begin{theorem}[escort 纤维信息极限的全阶累积量闭式与永久负偏斜]\label{thm:xi-time-part9ob-escort-loginfo-cumulants-negative-skew}
对任意 \(q\ge 0\)，令 \(W_m^{(q)}\sim \pi_{m,q}\)，并定义中心化纤维信息
\[
Y_m^{(q)}
:=
\log d_m^{\mathrm{bin}}(W_m^{(q)})
-m\log\!\Bigl(\frac{2}{\varphi}\Bigr).
\]
则
\[
Y_m^{(q)}\Longrightarrow Y^{(q)},
\]
其中极限分布是两点律
\[
Y^{(q)}=
\begin{cases}
0,& \text{概率 }1-\theta(q),\\
-\log\varphi,& \text{概率 }\theta(q).
\end{cases}
\]
因此其累积母函数为
\[
K_q(t)
:=
\log \mathbf E\!\left[e^{tY^{(q)}}\right]
=
\log\!\bigl(1-\theta(q)+\theta(q)\varphi^{-t}\bigr),
\qquad t\in\mathbb R.
\]
若记
\[
\kappa_r^{(q)}:=K_q^{(r)}(0),
\]
则前四阶累积量满足
\[
\kappa_1^{(q)}=-\theta(q)\log\varphi,
\]
\[
\kappa_2^{(q)}=(\log\varphi)^2\,\theta(q)\bigl(1-\theta(q)\bigr),
\]
\[
\kappa_3^{(q)}=-(\log\varphi)^3\,\theta(q)\bigl(1-\theta(q)\bigr)\bigl(1-2\theta(q)\bigr),
\]
\[
\kappa_4^{(q)}=(\log\varphi)^4\,\theta(q)\bigl(1-\theta(q)\bigr)\bigl(1-6\theta(q)+6\theta(q)^2\bigr).
\]
特别地，对所有 \(q\ge 0\) 都有
\[
\kappa_3^{(q)}<0.
\]
当 \(q=1\) 时，上述闭式恰化为定理 \ref{thm:xi-fold-hidden-logmultiplicity-laplace-cumulant-closed} 的均匀基线公式。
\end{theorem}
\begin{proof}
由定理 \ref{thm:fold-bin-two-state-asymptotic} 的一致对数展开，
\[
\log d_m^{\mathrm{bin}}(w)
=
m\log\!\Bigl(\frac{2}{\varphi}\Bigr)-w_m\log\varphi
+O\!\left(\Bigl(\frac{\varphi}{2}\Bigr)^m\right)
\]
在 \(w\in X_m\) 上一致成立，因此
\[
Y_m^{(q)}
=
-W_m^{(q)}(m)\log\varphi
+O_{\mathbf P}\!\left(\Bigl(\frac{\varphi}{2}\Bigr)^m\right).
\]
再与命题 \ref{prop:fold-bin-escort-lastbit} 的末位极限
\[
\mathbf P_{\pi_{m,q}}\bigl(W_m^{(q)}(m)=1\bigr)\to \theta(q)
\]
合并，即得所示两点极限分布。

于是对任意实数 \(t\)，有
\[
\mathbf E\!\left[e^{tY^{(q)}}\right]
=
\bigl(1-\theta(q)\bigr)+\theta(q)\varphi^{-t},
\]
取对数便得到 \(K_q(t)\)。

其后各阶累积量只需在 \(t=0\) 处逐次求导；等价地，也可写成
\[
Y^{(q)}=-(\log\varphi)\,B_q,
\qquad
B_q\sim \mathrm{Bernoulli}(\theta(q)),
\]
并直接调用 Bernoulli 两点律的标准累积量公式，即得所列四式。

最后，由
\[
0<\theta(q)\le \theta(0)=\frac{1}{1+\varphi}=\varphi^{-2}<\frac12
\]
可知
\[
1-2\theta(q)>0,
\]
故
\[
\kappa_3^{(q)}
=
-(\log\varphi)^3\,\theta(q)\bigl(1-\theta(q)\bigr)\bigl(1-2\theta(q)\bigr)<0.
\]
这便给出永久负偏斜性。证毕。
\end{proof}

\begin{theorem}[信息谱的二层量子化]\label{thm:xi-time-part9ob-escort-surprisal-two-level-quantization}
对任意 \(q\ge 0\)，令
\[
W_m^{(q)}\sim \pi_{m,q},
\qquad
\ell_m(w):=w_m,
\qquad
I_m^{(q)}:=-\log \pi_{m,q}(W_m^{(q)}).
\]
并记
\[
c(q):=\log\!\left(\frac{\varphi+\varphi^{-q}}{\sqrt5}\right).
\]
则有概率展开
\[
I_m^{(q)}
=
m\log\varphi+c(q)+q\,\ell_m(W_m^{(q)})\log\varphi+o_{\mathbf P}(1).
\]
因此
\[
I_m^{(q)}-m\log\varphi-c(q)
\Longrightarrow
q\log\varphi\cdot B_q,
\qquad
B_q\sim \mathrm{Bernoulli}(\theta(q)).
\]
特别地，在原始分布 \(\mu_m=\pi_{m,1}\) 下，
\[
I_m^{(1)}-m\log\varphi
\Longrightarrow
\log\varphi\cdot B,
\]
\[
\mathbf P(B=1)=\frac{1}{\varphi\sqrt5},
\qquad
\mathbf P(B=0)=\frac{\varphi}{\sqrt5}.
\]
\end{theorem}
\begin{proof}
由定理 \ref{thm:fold-bin-two-state-asymptotic} 与命题 \ref{prop:fold-bin-power-sum-asymptotic}，对 \(w\in X_m\) 一致有
\[
\pi_{m,q}(w)
=
\frac{\varphi^{-q\ell_m(w)}}{F_{m+1}+\varphi^{-q}F_m}
\Bigl(1+O\!\left(\Bigl(\frac{\varphi}{2}\Bigr)^m\right)\Bigr).
\]
取负对数得
\[
-\log \pi_{m,q}(w)
=
\log(F_{m+1}+\varphi^{-q}F_m)
+q\,\ell_m(w)\log\varphi
+O\!\left(\Bigl(\frac{\varphi}{2}\Bigr)^m\right),
\]
且误差在 \(w\) 上一致。另一方面，由 Binet 公式，
\[
F_{m+1}+\varphi^{-q}F_m
=
\frac{\varphi^m}{\sqrt5}\bigl(\varphi+\varphi^{-q}\bigr)+O(\varphi^{-m}),
\]
从而
\[
\log(F_{m+1}+\varphi^{-q}F_m)
=
m\log\varphi+c(q)+o(1).
\]
将两式合并并代入 \(w=W_m^{(q)}\)，即得
\[
I_m^{(q)}
=
m\log\varphi+c(q)+q\,\ell_m(W_m^{(q)})\log\varphi+o_{\mathbf P}(1).
\]
再由命题 \ref{prop:fold-bin-escort-lastbit}，
\[
\ell_m(W_m^{(q)})\Longrightarrow \mathrm{Bernoulli}(\theta(q)),
\]
故第二个极限立即成立。

当 \(q=1\) 时，
\[
c(1)=\log\!\left(\frac{\varphi+\varphi^{-1}}{\sqrt5}\right)=0,
\]
且
\[
\theta(1)=\frac{1}{1+\varphi^2}=\frac{1}{\varphi\sqrt5},
\qquad
1-\theta(1)=\frac{\varphi}{\sqrt5}.
\]
于是得到最后两式。证毕。
\end{proof}

\endinput

\paragraph{escort 温度族的一维指数封闭、局部正态性与阈值重建预算}
在末位两块正则形、温度间二点散度闭式、\((e)\)-测地仿射闭合与 Fisher--Rao 几何已经明确之后，还可继续把统计模型、局部几何、假设检验与阈值预算压缩为同一条单比特指数链。

\begin{theorem}[escort 温度族的一致对数指数分解与末位渐近充分性]\label{thm:xi-time-part9ob-escort-loglik-onebit-expfamily}
\leanverified{paper\_xi\_time\_part9ob\_escort\_loglik\_onebit\_expfamily}
对任意 \(q\ge 0\)，记
\[
\varepsilon_m:=\Bigl(\frac{\varphi}{2}\Bigr)^m,
\qquad
A_0:=\{w\in X_m:\ w_m=0\},
\qquad
A_1:=\{w\in X_m:\ w_m=1\}.
\]
定义显式两块权重
\[
\vartheta_m(q):=
\frac{\varphi^{-q}F_m}{F_{m+1}+\varphi^{-q}F_m},
\]
以及块内均匀模型
\[
\bar\pi_{m,q}(w):=
\begin{cases}
\dfrac{1-\vartheta_m(q)}{F_{m+1}},& w\in A_0,\\[8pt]
\dfrac{\vartheta_m(q)}{F_m},& w\in A_1.
\end{cases}
\]
则对每个固定 \(q\ge 0\)，有一致对数逼近
\[
\sup_{w\in X_m}\left|
\log\frac{\pi_{m,q}(w)}{\bar\pi_{m,q}(w)}
\right|
=
O(\varepsilon_m).
\]
进一步，对任意固定 \(p,q\ge 0\)，有一致似然比展开
\[
\log\frac{\pi_{m,q}(w)}{\pi_{m,p}(w)}
=
A_m(q,p)
-(q-p)w_m\log\varphi
+O(\varepsilon_m),
\]
其中
\[
A_m(q,p):=
\log\frac{F_{m+1}+\varphi^{-p}F_m}{F_{m+1}+\varphi^{-q}F_m}.
\]
特别地，若记
\[
T_m(w):=w_m\in\{0,1\},
\qquad
\eta(q):=-q\log\varphi,
\]
则 \(\bar\pi_{m,q}\) 具有精确指数族表达
\[
\bar\pi_{m,q}(w)
=
\exp\!\bigl(\eta(q)T_m(w)-\Psi_m(\eta(q))\bigr),
\qquad
\Psi_m(\eta):=\log\!\bigl(F_{m+1}+e^{\eta}F_m\bigr),
\]
故 escort 温度族在渐近上被同一末位统计量 \(T_m\) 充分化为一维 Bernoulli 指数族。
\end{theorem}
\begin{proof}
由定理 \ref{thm:fold-bin-two-state-asymptotic} 与命题 \ref{prop:fold-bin-power-sum-asymptotic}，对每个固定 \(q\ge 0\) 都有一致展开
\[
\pi_{m,q}(w)
=
\frac{\varphi^{-q w_m}}{F_{m+1}+\varphi^{-q}F_m}
\bigl(1+O(\varepsilon_m)\bigr),
\qquad
w\in X_m.
\]
而由 \(|A_0|=F_{m+1}\)、\(|A_1|=F_m\) 可直接计算
\[
\bar\pi_{m,q}(w)=
\begin{cases}
\dfrac{1}{F_{m+1}+\varphi^{-q}F_m},& w\in A_0,\\[8pt]
\dfrac{\varphi^{-q}}{F_{m+1}+\varphi^{-q}F_m},& w\in A_1.
\end{cases}
\]
因此
\[
\frac{\pi_{m,q}(w)}{\bar\pi_{m,q}(w)}=1+O(\varepsilon_m)
\]
在 \(w\in X_m\) 上一致成立。由于右端对充分大 \(m\) 落在某个固定紧区间 \([1/2,3/2]\) 内，取对数便得到第一条式子。

对第二条，只需将上式分别用于参数 \(q\) 与 \(p\)：
\[
\log\frac{\pi_{m,q}(w)}{\pi_{m,p}(w)}
=
\log\frac{\bar\pi_{m,q}(w)}{\bar\pi_{m,p}(w)}
+O(\varepsilon_m).
\]
而
\[
\log\frac{\bar\pi_{m,q}(w)}{\bar\pi_{m,p}(w)}
=
\log\frac{F_{m+1}+\varphi^{-p}F_m}{F_{m+1}+\varphi^{-q}F_m}
+\,
\log\!\bigl(\varphi^{-(q-p)w_m}\bigr),
\]
这就给出所示似然比展开。

最后，由 \(\eta(q)=-q\log\varphi\) 可把 \(\bar\pi_{m,q}\) 精确重写为
\[
\bar\pi_{m,q}(w)
=
\frac{e^{\eta(q)w_m}}{F_{m+1}+e^{\eta(q)}F_m}
=
\exp\!\bigl(\eta(q)T_m(w)-\Psi_m(\eta(q))\bigr).
\]
因此，全部参数依赖只通过 \(T_m(w)=w_m\) 出现，故 escort 温度族在上述一致对数误差意义下确已塌缩为一维指数族。证毕。
\end{proof}

\begin{theorem}[Neyman--Pearson 判别的末位坍缩]\label{cor:xi-time-part9ob-escort-np-lastbit-collapse}
\leanverified{paper\_xi\_time\_part9ob\_escort\_np\_lastbit\_collapse}
沿用定理 \ref{thm:xi-time-part9ob-escort-loglik-onebit-expfamily} 的记号。对任意 \(0\le p<q\)，记
\[
\ell_m(w):=w_m,
\qquad
C_m(p,q):=
\log\frac{F_{m+1}+\varphi^{-p}F_m}{F_{m+1}+\varphi^{-q}F_m}.
\]
则有统一似然比展开
\[
\log\frac{\pi_{m,q}(w)}{\pi_{m,p}(w)}
=
C_m(p,q)-(q-p)\log\varphi\cdot \ell_m(w)
+O(\varepsilon_m),
\]
其中误差对 \(w\in X_m\) 一致。因而在 \(m\to\infty\) 极限中，任何 Neyman--Pearson 最强检验都退化为对单个比特 \(\ell_m(w)\) 的阈值检验。

再令
\[
\theta(s):=\frac{1}{1+\varphi^{s+1}}
\qquad (s\ge 0).
\]
则有
\[
\|\pi_{m,q}-\pi_{m,p}\|_{\mathrm{TV}}
=
|\theta(q)-\theta(p)|+O(\varepsilon_m),
\]
从而
\[
\lim_{m\to\infty}\|\pi_{m,q}-\pi_{m,p}\|_{\mathrm{TV}}
=
|\theta(q)-\theta(p)|.
\]
等先验下，相应 Bayes 最小判别误差满足
\[
\lim_{m\to\infty}P_e^\ast(\pi_{m,q},\pi_{m,p})
=
\frac12\bigl(1-|\theta(q)-\theta(p)|\bigr).
\]
更具体地，等先验 Bayes 最小错误率在极限中完全由二元参数差 \(|\theta(q)-\theta(p)|\) 支配。
\end{theorem}
\begin{proof}
由定理 \ref{thm:xi-time-part9ob-escort-loglik-onebit-expfamily}，
\[
\|\pi_{m,q}-\bar\pi_{m,q}\|_{\mathrm{TV}}=O(\varepsilon_m),
\qquad
\|\pi_{m,p}-\bar\pi_{m,p}\|_{\mathrm{TV}}=O(\varepsilon_m).
\]
因此由三角不等式，
\[
\Bigl|
\|\pi_{m,q}-\pi_{m,p}\|_{\mathrm{TV}}
-\|\bar\pi_{m,q}-\bar\pi_{m,p}\|_{\mathrm{TV}}
\Bigr|
=
O(\varepsilon_m).
\]
而 \(\bar\pi_{m,q}\)、\(\bar\pi_{m,p}\) 在相同分块 \(A_0\sqcup A_1\) 上块内均匀，故
\[
\|\bar\pi_{m,q}-\bar\pi_{m,p}\|_{\mathrm{TV}}
=
|\vartheta_m(q)-\vartheta_m(p)|.
\]
又由 \(F_{m+1}/F_m\to\varphi\)，有
\[
\vartheta_m(s)\to\theta(s)
\qquad (s\ge 0),
\]
于是得到总变差展开与极限。Bayes 最小错误概率由 Helstrom 公式
\[
P_e^\ast(P,Q)=\frac12\bigl(1-\|P-Q\|_{\mathrm{TV}}\bigr)
\]
立得。

最后，由定理 \ref{thm:xi-time-part9ob-escort-loglik-onebit-expfamily} 的似然比公式，
\[
\log\frac{\pi_{m,q}(w)}{\pi_{m,p}(w)}
=
C_m(p,q)-(q-p)\log\varphi\cdot \ell_m(w)+O(\varepsilon_m).
\]
因此两块上的似然比差恰为
\[
(p-q)\log\varphi+O(\varepsilon_m).
\]
由于此处 \(q>p\)，上式对充分大 \(m\) 为正，故似然比在 \(A_0\) 上严格更大。于是 Neyman--Pearson 最优检验不再需要读取整个 \(w\)，只需读取末位并在对应块上作决定。证毕。
\end{proof}

\begin{theorem}[escort 温度流的双向局部 KL 曲率闭式]\label{thm:xi-time-part9ob-escort-two-sided-local-kl}
对任意固定 \(q\ge 0\)，当 \(h\to 0\) 且 \(q+h\ge 0\) 时，有
\[
\lim_{m\to\infty}D_{\mathrm{KL}}(\pi_{m,q+h}\Vert \pi_{m,q})
=
\frac12 I(q)h^2+O(h^3),
\]
\[
\lim_{m\to\infty}D_{\mathrm{KL}}(\pi_{m,q}\Vert \pi_{m,q+h})
=
\frac12 I(q)h^2+O(h^3),
\]
其中
\[
I(q):=
(\log\varphi)^2\,\theta(q)\bigl(1-\theta(q)\bigr)
=
(\log\varphi)^2\frac{\varphi^{q+1}}{(1+\varphi^{q+1})^2}.
\]
因此 escort 温度流的局部 KL 曲率在正反两个方向上都由同一 Fisher 密度 \(I(q)\) 完全控制。
\end{theorem}
\begin{proof}
由结论 \ref{con:xi-time-part9o-escort-fdiv-collapse} 在
\[
f(u)=u\log u
\]
处的专化，
\[
\lim_{m\to\infty}D_{\mathrm{KL}}(\pi_{m,q+h}\Vert \pi_{m,q})
=
D_{\mathrm{KL}}\!\bigl(\mathrm{Bernoulli}(\theta(q+h))\Vert \mathrm{Bernoulli}(\theta(q))\bigr).
\]
对 Bernoulli 散度的标准 Taylor 展开给出
\[
D_{\mathrm{KL}}\!\bigl(\mathrm{Bernoulli}(\theta+\delta)\Vert \mathrm{Bernoulli}(\theta)\bigr)
=
\frac{\delta^2}{2\theta(1-\theta)}+O(\delta^3),
\]
\[
D_{\mathrm{KL}}\!\bigl(\mathrm{Bernoulli}(\theta)\Vert \mathrm{Bernoulli}(\theta+\delta)\bigr)
=
\frac{\delta^2}{2\theta(1-\theta)}+O(\delta^3).
\]
另一方面，
\[
\theta'(q)
=
-(\log\varphi)\,\theta(q)\bigl(1-\theta(q)\bigr),
\]
故
\[
\theta(q+h)-\theta(q)=\theta'(q)h+O(h^2).
\]
将
\[
\delta=\theta(q+h)-\theta(q)
\]
代回前两式，即得
\[
\frac{(\theta'(q))^2}{2\theta(q)(1-\theta(q))}h^2+O(h^3)
=
\frac12(\log\varphi)^2\theta(q)(1-\theta(q))h^2+O(h^3),
\]
这正是所示公式。证毕。
\end{proof}

\begin{theorem}[亚临界样本规模下 escort 乘积试验与显式 Bernoulli 乘积试验的 Le Cam 等价]\label{thm:xi-time-part9ob-escort-product-lecam}
\leanverified{paper\_xi\_time\_part9ob\_escort\_product\_lecam}
设 \(Q\subset[0,\infty)\) 为任意有限集合，\((n_m)_{m\ge 1}\) 为正整数序列。定义 escort 乘积试验
\[
\mathcal E_m(Q):=\{\pi_{m,q}^{\otimes n_m}:q\in Q\}
\qquad\text{定义在 }X_m^{n_m}\text{ 上},
\]
以及显式 Bernoulli 乘积试验
\[
\mathcal B_m(Q):=
\{\mathrm{Bernoulli}(\vartheta_m(q))^{\otimes n_m}:q\in Q\}
\qquad\text{定义在 }\{0,1\}^{n_m}\text{ 上}.
\]
则存在与 \(q\) 无关的 Markov 核
\[
L_m:X_m^{n_m}\rightsquigarrow \{0,1\}^{n_m},
\qquad
K_m:\{0,1\}^{n_m}\rightsquigarrow X_m^{n_m},
\]
使得对某个仅依赖于 \(Q\) 的常数 \(C_Q>0\)，有
\[
\sup_{q\in Q}
D_{\mathrm{TV}}\!\bigl(L_{m\#}\pi_{m,q}^{\otimes n_m},
\mathrm{Bernoulli}(\vartheta_m(q))^{\otimes n_m}\bigr)
\le
C_Q n_m\varepsilon_m,
\]
\[
\sup_{q\in Q}
D_{\mathrm{TV}}\!\bigl(K_{m\#}\mathrm{Bernoulli}(\vartheta_m(q))^{\otimes n_m},
\pi_{m,q}^{\otimes n_m}\bigr)
\le
C_Q n_m\varepsilon_m.
\]
因而 Le Cam 距离满足
\[
\Delta\!\bigl(\mathcal E_m(Q),\mathcal B_m(Q)\bigr)
\le
C_Q n_m\varepsilon_m.
\]
特别地，若
\[
n_m=o\!\left(\Bigl(\frac{2}{\varphi}\Bigr)^m\right),
\]
则
\[
\Delta\!\bigl(\mathcal E_m(Q),\mathcal B_m(Q)\bigr)\longrightarrow 0.
\]
\end{theorem}
\begin{proof}
对 \((w^{(1)},\dots,w^{(n_m)})\in X_m^{n_m}\)，令
\[
L_m(w^{(1)},\dots,w^{(n_m)})
:=
\bigl(w^{(1)}_m,\dots,w^{(n_m)}_m\bigr).
\]
再令 \(U_{m,b}\) 为 \(A_b\) 上的均匀分布，并定义重构核
\[
K_m(b_1,\dots,b_{n_m},\cdot)
:=
U_{m,b_1}\otimes\cdots\otimes U_{m,b_{n_m}}.
\]

先证前向估计。记
\[
p_{m,q}:=\pi_{m,q}(A_1).
\]
则
\[
L_{m\#}\pi_{m,q}^{\otimes n_m}
=
\mathrm{Bernoulli}(p_{m,q})^{\otimes n_m}.
\]
另一方面，由定理 \ref{thm:xi-time-part9ob-escort-loglik-onebit-expfamily}，
\[
|p_{m,q}-\vartheta_m(q)|
\le
\|\pi_{m,q}-\bar\pi_{m,q}\|_{\mathrm{TV}}
=
O_Q(\varepsilon_m),
\]
其中由于 \(Q\) 有限，常数可对 \(q\in Q\) 统一。对乘积 Bernoulli 分布用望远镜估计，得到
\[
D_{\mathrm{TV}}\!\bigl(
\mathrm{Bernoulli}(p_{m,q})^{\otimes n_m},
\mathrm{Bernoulli}(\vartheta_m(q))^{\otimes n_m}
\bigr)
\le
n_m|p_{m,q}-\vartheta_m(q)|
\le
C_Q n_m\varepsilon_m.
\]

再看反向估计。由构造，
\[
K_{m\#}\mathrm{Bernoulli}(\vartheta_m(q))^{\otimes n_m}
=
\bar\pi_{m,q}^{\otimes n_m}.
\]
又由乘积测度的望远镜估计与定理
\ref{thm:xi-time-part9ob-escort-loglik-onebit-expfamily}，
\[
D_{\mathrm{TV}}\!\bigl(\bar\pi_{m,q}^{\otimes n_m},\pi_{m,q}^{\otimes n_m}\bigr)
\le
n_m D_{\mathrm{TV}}(\bar\pi_{m,q},\pi_{m,q})
\le
C_Q n_m\varepsilon_m.
\]
故后向估计亦成立。两侧 deficiency 都由上述界控制，遂得 Le Cam 距离估计；若 \(n_m=o((2/\varphi)^m)\)，则右端趋于 \(0\)。证毕。
\end{proof}

\begin{theorem}[escort 乘积试验的局部渐近正态性与末位计数中心序列]\label{thm:xi-time-part9ob-escort-lan-lastbit}
设 \((n_m)_{m\ge 1}\) 满足
\[
n_m\to\infty,
\qquad
n_m=o\!\left(\Bigl(\frac{2}{\varphi}\Bigr)^m\right).
\]
固定 \(q>0\)，并记
\[
P_{m,q}:=\pi_{m,q}^{\otimes n_m}.
\]
对 \(W^{(1)},\dots,W^{(n_m)}\sim P_{m,q}\)，定义末位计数
\[
N_m:=\sum_{j=1}^{n_m}W^{(j)}_m,
\]
以及中心序列
\[
\Delta_{m,q}:=
-\frac{\log\varphi}{\sqrt{n_m}}
\Bigl(N_m-n_m\theta(q)\Bigr).
\]
再记
\[
I(q):=(\log\varphi)^2\theta(q)\bigl(1-\theta(q)\bigr).
\]
则对任意固定 \(h\in\RR\)，都有
\[
\log\frac{dP_{m,q+h/\sqrt{n_m}}}{dP_{m,q}}
=
h\Delta_{m,q}
-\frac12 h^2 I(q)
+o_{P_{m,q}}(1),
\]
并且
\[
\Delta_{m,q}\Rightarrow \mathcal N\!\bigl(0,I(q)\bigr).
\]
因此 escort 乘积试验在每个内点 \(q>0\) 上都满足局部渐近正态性，而其全部局部信息只由末位计数 \(N_m\) 承载。
\end{theorem}
\begin{proof}
记
\[
\delta_m:=\frac{h}{\sqrt{n_m}}.
\]
由定理 \ref{thm:fold-bin-two-state-asymptotic} 与命题 \ref{prop:fold-bin-power-sum-asymptotic} 的同一推导，在 \(q\) 的一个固定紧邻域内可取统一常数，从而
\[
\log\frac{\pi_{m,q+\delta_m}(w)}{\pi_{m,q}(w)}
=
A_m(q+\delta_m,q)-\delta_m w_m\log\varphi+O(\varepsilon_m)
\]
在 \(w\in X_m\) 上一致成立。对 \(n_m\) 个样本求和得
\[
\log\frac{dP_{m,q+\delta_m}}{dP_{m,q}}
=
n_mA_m(q+\delta_m,q)-\delta_m(\log\varphi)N_m
+O(n_m\varepsilon_m).
\]

令
\[
D_m(s):=F_{m+1}+\varphi^{-s}F_m,
\qquad
\vartheta_m(s):=\frac{\varphi^{-s}F_m}{D_m(s)}.
\]
则
\[
A_m(q+\delta,q)=\log D_m(q)-\log D_m(q+\delta).
\]
直接求导得到
\[
\partial_\delta A_m(q+\delta,q)\big|_{\delta=0}
=
(\log\varphi)\vartheta_m(q),
\]
\[
\partial_\delta^2 A_m(q+\delta,q)\big|_{\delta=0}
=
-(\log\varphi)^2\vartheta_m(q)\bigl(1-\vartheta_m(q)\bigr).
\]
由于第三导数在该紧邻域内一致有界，Taylor 展开给出
\[
n_mA_m(q+\delta_m,q)
=
h(\log\varphi)\sqrt{n_m}\,\vartheta_m(q)
-\frac12 h^2 I_m(q)
+O(n_m^{-1/2}),
\]
其中
\[
I_m(q):=
(\log\varphi)^2\vartheta_m(q)\bigl(1-\vartheta_m(q)\bigr).
\]
代回得
\[
\log\frac{dP_{m,q+\delta_m}}{dP_{m,q}}
=
h\widetilde\Delta_{m,q}
-\frac12 h^2 I_m(q)
+O(n_m\varepsilon_m)+O(n_m^{-1/2}),
\]
其中
\[
\widetilde\Delta_{m,q}
:=
-\frac{\log\varphi}{\sqrt{n_m}}
\Bigl(N_m-n_m\vartheta_m(q)\Bigr).
\]

在 \(P_{m,q}\) 下，随机变量 \(W^{(j)}_m\) 相互独立，且其末位 \(W^{(j)}_m\in\{0,1\}\) 服从 Bernoulli\((p_{m,q})\)，其中
\[
p_{m,q}:=\pi_{m,q}(A_1)=\vartheta_m(q)+O(\varepsilon_m)
\]
由定理 \ref{thm:xi-time-part9ob-escort-loglik-onebit-expfamily} 得到。因此
\[
N_m\sim \mathrm{Binomial}(n_m,p_{m,q}).
\]
由经典中心极限定理，
\[
-\frac{\log\varphi}{\sqrt{n_m}}
\Bigl(N_m-n_mp_{m,q}\Bigr)
\Rightarrow
\mathcal N\!\bigl(0,(\log\varphi)^2p_{m,q}(1-p_{m,q})\bigr).
\]
再用
\[
\sqrt{n_m}\,|p_{m,q}-\vartheta_m(q)|
=
O(\sqrt{n_m}\varepsilon_m)\to 0
\]
与
\[
\sqrt{n_m}\,|\vartheta_m(q)-\theta(q)|
=
O(\sqrt{n_m}\varepsilon_m)\to 0,
\]
便得到
\[
\widetilde\Delta_{m,q}-\Delta_{m,q}\to 0
\qquad\text{于 }P_{m,q}\text{-概率},
\]
同时 \(I_m(q)\to I(q)\)。又因
\[
n_m\varepsilon_m\to 0,
\qquad
n_m^{-1/2}\to 0,
\]
所以上述对数似然比展开化为
\[
\log\frac{dP_{m,q+h/\sqrt{n_m}}}{dP_{m,q}}
=
h\Delta_{m,q}
-\frac12 h^2 I(q)
+o_{P_{m,q}}(1).
\]
这就证明了 LAN 公式与中心序列的 Gaussian 极限。证毕。
\end{proof}

\begin{theorem}[bin-fold 配额阈值存在严格二相常数律，且相变幅度恰为黄金单位]\label{thm:xi-time-part9ob-binfold-sharp-threshold-budget}
记 \(\mu_m:=\pi_{m,1}\)，即
\[
\mu_m(w)=\frac{d_m^{\mathrm{bin}}(w)}{2^m}.
\]
对 \(\varepsilon\in(0,1)\)，定义临界重建预算
\[
P_m^{\mathrm{crit}}(\varepsilon):=
\min\Bigl\{
P\in\NN:\ 
\mathbf P_{\mu_m}\!\Bigl(d_m^{\mathrm{bin}}(W_m)>\frac{P}{2}\Bigr)\le \varepsilon
\Bigr\}.
\]
再令
\[
\varepsilon_{\mathrm c}:=\frac{\varphi}{\sqrt5}.
\]
则有精确夹逼
\[
2B_m^{\mathrm{bin}}(\varepsilon)
\le
P_m^{\mathrm{crit}}(\varepsilon)
\le
2B_m^{\mathrm{bin}}(\varepsilon)+2,
\]
从而
\[
\frac{P_m^{\mathrm{crit}}(\varepsilon)}{2B_m^{\mathrm{bin}}(\varepsilon)}\to 1.
\]
特别地，对任意固定 \(\varepsilon\in(0,\varepsilon_{\mathrm c})\)，有
\[
\boxed{\;
P_m^{\mathrm{crit}}(\varepsilon)
\sim
2\left(\frac{2}{\varphi}\right)^m.\;}
\]
而对任意固定 \(\varepsilon\in(\varepsilon_{\mathrm c},1)\)，有
\[
\boxed{\;
P_m^{\mathrm{crit}}(\varepsilon)
\sim
2\varphi^{-1}\left(\frac{2}{\varphi}\right)^m.\;}
\]
等价地，
\[
\log P_m^{\mathrm{crit}}(\varepsilon)
=
m\log\!\Bigl(\frac{2}{\varphi}\Bigr)
+\begin{cases}
\log 2+o(1),& 0<\varepsilon<\varepsilon_{\mathrm c},\\[4pt]
\log 2-\log\varphi+o(1),& \varepsilon_{\mathrm c}<\varepsilon<1.
\end{cases}
\]
若以 \(P_m^{\mathrm{crit}}(\varepsilon_{\mathrm c}^-)\)、\(P_m^{\mathrm{crit}}(\varepsilon_{\mathrm c}^+)\) 记左右两相的渐近延拓，则
\[
\boxed{\;
\frac{P_m^{\mathrm{crit}}(\varepsilon_{\mathrm c}^+)}
{P_m^{\mathrm{crit}}(\varepsilon_{\mathrm c}^-)}
=\varphi^{-1},\qquad
\log P_m^{\mathrm{crit}}(\varepsilon_{\mathrm c}^-)
-\log P_m^{\mathrm{crit}}(\varepsilon_{\mathrm c}^+)
=\log\varphi.\;}
\]
\end{theorem}
\begin{proof}
由推论 \ref{cor:fold-bin-quantile-threshold-constant}，
\[
B_m^{\mathrm{bin}}(\varepsilon)
=
\inf\Bigl\{B>0:\ 
\mathbf P_{\mu_m}\bigl(d_m^{\mathrm{bin}}(W_m)>B\bigr)\le \varepsilon
\Bigr\}.
\]
因此，任意满足
\[
\mathbf P_{\mu_m}\!\Bigl(d_m^{\mathrm{bin}}(W_m)>\frac{P}{2}\Bigr)\le \varepsilon
\]
的整数 \(P\) 都必须满足 \(P/2\ge B_m^{\mathrm{bin}}(\varepsilon)\)，从而
\[
P_m^{\mathrm{crit}}(\varepsilon)\ge 2B_m^{\mathrm{bin}}(\varepsilon).
\]
反过来，取
\[
P=\lfloor 2B_m^{\mathrm{bin}}(\varepsilon)\rfloor+2,
\]
则 \(P/2>B_m^{\mathrm{bin}}(\varepsilon)\)。又因
\[
B\longmapsto \mathbf P_{\mu_m}\bigl(d_m^{\mathrm{bin}}(W_m)>B\bigr)
\]
随 \(B\) 单调递减，满足尾概率约束的阈值集合必为一条上射线，故该 \(P\) 已满足所要求的尾概率约束，于是
\[
P_m^{\mathrm{crit}}(\varepsilon)\le 2B_m^{\mathrm{bin}}(\varepsilon)+2.
\]
这就给出所示夹逼。

又由推论 \ref{cor:fold-bin-quantile-threshold-constant}，
\[
\Bigl(\frac{\varphi}{2}\Bigr)^m B_m^{\mathrm{bin}}(\varepsilon)\longrightarrow
\begin{cases}
1,& 0<\varepsilon<\varepsilon_{\mathrm c},\\[4pt]
\varphi^{-1},& \varepsilon_{\mathrm c}<\varepsilon<1.
\end{cases}
\]
故 \(B_m^{\mathrm{bin}}(\varepsilon)\) 以指数速度趋于 \(+\infty\)，并且
\[
\frac{2B_m^{\mathrm{bin}}(\varepsilon)+2}{2B_m^{\mathrm{bin}}(\varepsilon)}\to 1.
\]
与前述夹逼合并便得
\[
\frac{P_m^{\mathrm{crit}}(\varepsilon)}{2B_m^{\mathrm{bin}}(\varepsilon)}\to 1.
\]
再把分段极限代入，即得两条乘性渐近。取对数后立得对应的分段对数渐近。

最后，左右两相的主常数分别为 \(2\) 与 \(2\varphi^{-1}\)，故
\[
\frac{2\varphi^{-1}}{2}=\varphi^{-1},
\qquad
\log 2-\log(2\varphi^{-1})=\log\varphi,
\]
这就给出相变幅度的精确黄金单位。证毕。
\end{proof}

\begin{corollary}[预算相变阈值恰等于主导末位块的渐近质量]\label{cor:xi-time-part9ob-binfold-threshold-equals-main-block-mass}
沿用定理 \ref{thm:xi-time-part9ob-binfold-sharp-threshold-budget} 的记号，并记
\[
A_0:=\{w\in X_m:\ w_m=0\},
\qquad
A_1:=\{w\in X_m:\ w_m=1\}.
\]
则
\[
\mu_m(A_0)\to \frac{\varphi}{\sqrt5},
\qquad
\mu_m(A_1)\to \frac{1}{\varphi\sqrt5},
\]
从而临界阈值具有完全内禀的块质量解释
\[
\boxed{\;
\varepsilon_{\mathrm c}
=
\frac{\varphi}{\sqrt5}
=
\lim_{m\to\infty}\mu_m(A_0).\;}
\]
换言之，二相预算律的相变点并非外加参数，而恰由主导末位块 \(A_0\) 的渐近质量强制给出；一旦容许失败率跨越该质量，临界预算的主常数便精确降低为原先的 \(\varphi^{-1}\) 倍。
\end{corollary}
\begin{proof}
由推论 \ref{cor:fold-bin-two-point-limit-law}，在 \(\mu_m\) 下有
\[
\mu_m(A_0)\to \frac{\varphi}{\sqrt5},
\qquad
\mu_m(A_1)\to \frac{1}{\varphi\sqrt5}.
\]
而定理 \ref{thm:xi-time-part9ob-binfold-sharp-threshold-budget} 中的相变阈值正是
\(
\varepsilon_{\mathrm c}=\varphi/\sqrt5
\)。
二者相等即得所示块质量解释。最后一句只是同一定理中两相主常数之比
\(
(2\varphi^{-1})/2=\varphi^{-1}
\)
的直接改写。证毕。
\end{proof}

\noindent 这组结果表明，escort 温度族的塌缩并不只发生在熵率层面。统一末位坐标 \(w_m\) 同时给出：单样本对数似然的指数族因子分解、两温度间最优假设检验的全部非平凡拒绝域、局部 KL 曲率与 Fisher 信息、亚临界样本规模下乘积试验的 Le Cam 等价，以及 \(q=1\) 切片上的 sharp 阈值重建预算。于是，最新主线中关于统计模型、局部几何、判别理论与资源预算的全部主阶信息，便在同一个一维 Bernoulli 指数族内闭合；在读出口径上，这正是一种经末位中心投影实现的相位盲两态读出。

\endinput


\endinput


\begin{conclusion}[纯碰撞单位根扭曲的素模弛豫时间二阶展开]\label{con:xi-time-part9o-prime-modulus-relaxation-second-order}
沿用命题 \ref{prop:arity-collision-quadratic-closed} 与注记 \ref{cor:arity-335-diffusive-mixing} 的记号，记
\[
r_p:=\frac{\rho_{3,3,p}}{\lambda},
\qquad
\tau_p:=\tau_{\mathrm{mix}}(p):=\frac{1}{-\log r_p}.
\]
则当奇素数 \(p\to\infty\) 时有二阶渐近展开
\[
\tau_p
=
\frac{p^2}{4\pi^2\kappa_\infty}
\;+\;
\frac{\beta}{\kappa_\infty^2}
\;+\;
O(p^{-2}),
\]
其中
\[
\kappa_\infty=\frac{6\sqrt5-5}{250},
\qquad
\beta=\frac{P_\xi^{(4)}(0)}{24}=\frac{7+24\sqrt5}{75000}.
\]
等价地，
\[
\tau_p
=
\frac{125}{2\pi^2(6\sqrt5-5)}\,p^2
\;+\;
\frac{\frac{7+24\sqrt5}{75000}}
{\left(\frac{6\sqrt5-5}{250}\right)^2}
\;+\;
O(p^{-2}).
\]
特别地，prime-modulus 弛豫时间呈扩散型二次标度
\[
\tau_p\sim c_\varphi\,p^2,
\qquad
c_\varphi:=\frac{125}{2\pi^2(6\sqrt5-5)}.
\]
\end{conclusion}
\begin{proof}
由命题 \ref{prop:arity-collision-quadratic-closed}，
\[
\log r_p
=
-\kappa_\infty\left(\frac{2\pi}{p}\right)^2
\;+\;
\beta\left(\frac{2\pi}{p}\right)^4
\;+\;
O(p^{-6}),
\]
其中 \(\kappa_\infty\) 与 \(\beta\) 如上。
于是
\[
\log\frac{\lambda}{\rho_{3,3,p}}
=
\kappa_\infty\left(\frac{2\pi}{p}\right)^2
\;-\;
\beta\left(\frac{2\pi}{p}\right)^4
\;+\;
O(p^{-6}).
\]
记
\[
X_p:=\kappa_\infty\left(\frac{2\pi}{p}\right)^2,
\qquad
Y_p:=\beta\left(\frac{2\pi}{p}\right)^4+O(p^{-6}),
\]
则
\[
\tau_p=\frac{1}{X_p-Y_p}
=
\frac{1}{X_p}\cdot \frac{1}{1-Y_p/X_p}.
\]
由于 \(X_p\asymp p^{-2}\) 且 \(Y_p/X_p=O(p^{-2})\)，对几何级数展开可得
\[
\tau_p
=
\frac{1}{X_p}\left(1+\frac{Y_p}{X_p}+O(p^{-4})\right)
=
\frac{1}{X_p}+\frac{Y_p}{X_p^2}+O(p^{-2}).
\]
逐项化简得到
\[
\frac{1}{X_p}
=
\frac{p^2}{4\pi^2\kappa_\infty},
\qquad
\frac{Y_p}{X_p^2}
=
\frac{\beta}{\kappa_\infty^2}+O(p^{-2}),
\]
从而
\[
\tau_p
=
\frac{p^2}{4\pi^2\kappa_\infty}
\;+\;
\frac{\beta}{\kappa_\infty^2}
\;+\;
O(p^{-2}).
\]
把 \(\kappa_\infty=(6\sqrt5-5)/250\) 与 \(\beta=(7+24\sqrt5)/75000\) 代回，即得显式系数式。证毕。
\end{proof}
在二阶弛豫展开之外，同一 slow mode 还可进一步压缩为一个只保留主导标度的扩散寿命律。

\begin{theorem}[周期波纹的扩散寿命律]\label{thm:xi-time-part9o-periodic-ripple-diffusive-lifetime}
\leanverified{paper\_xi\_time\_part9o\_periodic\_ripple\_diffusive\_lifetime}
沿用推论 \ref{cor:real-input-40-collision-twist-fourth} 与结论
\ref{con:xi-time-part9o-prime-modulus-relaxation-second-order} 的记号。记
\[
\lambda:=\rho(1)=\varphi^2,
\qquad
\rho(t):=\rho(e^{it}),
\qquad
t_p:=\frac{2\pi}{p}
\]
以及
\[
L(p):=\frac{1}{-\log(\rho(t_p)/\lambda)}
\qquad
(p\ge 3\ \text{为奇素数}).
\]
则当 \(p\to\infty\) 时，
\[
L(p)
=
\frac{250}{(6\sqrt5-5)(2\pi)^2}\,p^2+O(1)
=
0.7524081626\ldots\,p^2+O(1).
\]
因此，最慢周期波纹从离散振荡过渡到平滑包络的主导寿命尺度具有严格扩散型 \(p^2\) 标度。
\end{theorem}
\begin{proof}
结论 \ref{con:xi-time-part9o-prime-modulus-relaxation-second-order} 已给出
\[
L(p)
=
\frac{125}{2\pi^2(6\sqrt5-5)}\,p^2
\;+\;
\frac{\frac{7+24\sqrt5}{75000}}
{\left(\frac{6\sqrt5-5}{250}\right)^2}
\;+\;
O(p^{-2}).
\]
把常数项并入 \(O(1)\) 后，即得
\[
L(p)=\frac{250}{(6\sqrt5-5)(2\pi)^2}\,p^2+O(1).
\]
数值系数由直接代入计算得到。证毕。
\end{proof}

\paragraph{小角扭曲的连续时间尺度展开}
在单位根扭曲的四阶谱隙展开、素模弛豫时间的二阶渐近以及扩散寿命律都已建立之后，同一 slow mode 还可进一步写成一个连续角度口径下的显式时间尺度公式。

\begin{theorem}[小角扭曲模态的谱混合时间具有显式二阶展开]\label{thm:xi-time-part9oe-small-angle-twist-mixingtime-second-order}
沿用推论 \ref{cor:real-input-40-collision-twist-fourth} 的记号，记
\[
\lambda:=\rho(1)=\varphi^2,
\qquad
\rho(t):=\rho(e^{it}),
\]
并定义小角扭曲模态的谱混合时间
\[
\tau(t):=\left(-\log\frac{\rho(t)}{\lambda}\right)^{-1}.
\]
则当 \(t\to 0\) 时，有严格展开
\[
\boxed{\;
\tau(t)
=
\frac{250}{6\sqrt5-5}\,\frac{1}{t^2}
+
\frac{5(7+24\sqrt5)}{6(6\sqrt5-5)^2}
+
O(t^2).\;}
\]
进一步，对奇素数模的最小非平凡角
\[
t_p:=\frac{2\pi}{p},
\]
有
\[
\boxed{\;
\tau(t_p)
=
\frac{125}{2\pi^2(6\sqrt5-5)}\,p^2
+
\frac{5(7+24\sqrt5)}{6(6\sqrt5-5)^2}
+
O(p^{-2}).\;}
\]
\leanverified{paper\_xi\_time\_part9oe\_small\_angle\_twist\_mixingtime\_second\_order}
\end{theorem}
\begin{proof}
由推论 \ref{cor:real-input-40-collision-twist-fourth}，
\[
\log\frac{\rho(t)}{\lambda}
=
-\frac{6\sqrt5-5}{250}\,t^2
+
\frac{7+24\sqrt5}{75000}\,t^4
+
O(t^6).
\]
记
\[
A:=\frac{6\sqrt5-5}{250},
\qquad
B:=\frac{7+24\sqrt5}{75000},
\]
则
\[
-\log\frac{\rho(t)}{\lambda}=A t^2-B t^4+O(t^6).
\]
于是
\[
\tau(t)
=
\frac{1}{A t^2}\cdot\frac{1}{1-(B/A)t^2+O(t^4)}.
\]
对第二因子作几何级数展开，得到
\[
\frac{1}{1-(B/A)t^2+O(t^4)}
=
1+\frac{B}{A}t^2+O(t^4),
\]
从而
\[
\tau(t)=\frac{1}{A t^2}+\frac{B}{A^2}+O(t^2).
\]
把 \(A,B\) 代回，即得第一条公式。

再取 \(t=t_p=2\pi/p\)，便有
\[
\frac{1}{A t_p^2}
=
\frac{1}{A}\cdot\frac{p^2}{4\pi^2}
=
\frac{125}{2\pi^2(6\sqrt5-5)}\,p^2,
\]
而常数项 \(\frac{B}{A^2}\) 保持不变，余项 \(O(t_p^2)\) 则化为 \(O(p^{-2})\)。证毕。
\end{proof}

\begin{theorem}[单位根扭曲的局部谱隙严格高于高斯二阶模型]\label{thm:xi-time-part9oe-twist-gap-strictly-above-gaussian}
沿用推论 \ref{cor:real-input-40-collision-twist-fourth} 的记号，记
\[
a:=\frac{6\sqrt5-5}{250}>0,
\qquad
b:=\frac{7+24\sqrt5}{75000}>0.
\]
则存在 \(\varepsilon_0>0\)，使得对一切
\[
0<|t|<\varepsilon_0,
\]
都有
\[
\frac{\rho(t)}{\lambda}>\exp(-a t^2).
\]
更精确地，
\[
\log\frac{\rho(t)}{\lambda}+a t^2
=
b t^4+O(t^6)
\sim b t^4>0.
\]
因此，任何只保留二阶项的高斯型局部谱隙模型都会系统性低估真实的扭曲 Perron 根；其首个不可约修正严格发生在四次阶，且符号恒为正。
\end{theorem}
\begin{proof}
由推论 \ref{cor:real-input-40-collision-twist-fourth}，
\[
\log\frac{\rho(t)}{\lambda}
=
-a t^2+b t^4+O(t^6).
\]
移项即得
\[
\log\frac{\rho(t)}{\lambda}+a t^2=b t^4+O(t^6).
\]
由于 \(b>0\)，可取 \(\varepsilon_0>0\) 充分小，使得当 \(0<|t|<\varepsilon_0\) 时右端严格为正。于是
\[
\log\frac{\rho(t)}{\lambda}>-a t^2.
\]
两边指数化便得到
\[
\frac{\rho(t)}{\lambda}>\exp(-a t^2).
\]
而
\[
\log\frac{\rho(t)}{\lambda}+a t^2\sim b t^4>0
\]
则直接给出四次正修正的主项。证毕。
\end{proof}

\endinput


\endinput


\paragraph{碰撞频率代数曲线、奇素数 shadow 与 Rényi--2 刚性}
在碰撞势的三次压力方程、核版 RH/Ramanujan 窗口、分歧晶格以及 Fibonacci 共振常数已经建立之后，还可把频率几何、素数扭曲与均匀性障碍进一步压缩为下列五条结果。

\begin{theorem}[碰撞频率的代数曲线与精确大偏差参数化]\label{thm:xi-time-part9p-collision-frequency-algebraic-ldp}
\leanverified{paper\_xi\_time\_part9p\_collision\_frequency\_algebraic\_ldp}
沿用命题 \ref{prop:real-input-40-collision-pressure} 的记号，令
\[
u=e^\theta,\qquad
P_2(\theta):=\log \rho(e^\theta),\qquad
\alpha(\theta):=P_2'(\theta).
\]
则在正实扭曲支 \(u>0\) 上成立：
\begin{enumerate}
  \item 碰撞频率 \(\alpha\) 与扭曲参数 \(u\) 满足显式三次代数曲线
  \begin{equation}\label{eq:xi-time-part9p-collision-frequency-algebraic-curve}
  \bigl(4u^3+25u^2+88u+8\bigr)\bigl(\alpha^3-\alpha^2\bigr)
  +\bigl(u^3+10u^2+24u\bigr)\alpha-2u^2=0.
  \end{equation}
  \item 以 Perron 根 \(\rho\) 为参数，可写成
  \begin{equation}\label{eq:xi-time-part9p-collision-frequency-rho-param}
  u(\rho)=\frac{\rho(\rho^2-2\rho-1)}{\rho-1},
  \qquad
  \alpha(\rho)=\frac{(\rho^2-2\rho-1)(\rho-1)}{2\rho^3-5\rho^2+4\rho+1},
  \qquad \rho\in(1+\sqrt2,\infty).
  \end{equation}
  \item 记
  \[
  I_2(a):=\sup_{\theta\in\RR}\Bigl\{a\theta-\bigl(P_2(\theta)-P_2(0)\bigr)\Bigr\},
  \]
  则在参数点 \(a=\alpha(\rho)\) 处有闭式
  \begin{equation}\label{eq:xi-time-part9p-collision-frequency-rate-param}
  I_2\bigl(\alpha(\rho)\bigr)
  =
  \alpha(\rho)\log\!\frac{\rho(\rho^2-2\rho-1)}{\rho-1}
  -\log\!\frac{\rho}{\varphi^2}.
  \end{equation}
\end{enumerate}
\end{theorem}
\begin{proof}
记
\[
f(\rho,u):=\rho^3-2\rho^2-(u+1)\rho+u.
\]
由 \(f(\rho(u),u)=0\) 对 \(u\) 作隐函数求导，得到
\[
\frac{\mathrm d\rho}{\mathrm du}
=-\frac{\partial_u f}{\partial_\rho f}
=\frac{\rho-1}{3\rho^2-4\rho-(u+1)}.
\]
又因 \(u=e^\theta\)，故
\[
\alpha(\theta)=\frac{\mathrm d}{\mathrm d\theta}\log\rho(e^\theta)
=\frac{u}{\rho}\frac{\mathrm d\rho}{\mathrm du}.
\]
由 \(f(\rho,u)=0\) 解出
\[
u(\rho)=\frac{\rho(\rho^2-2\rho-1)}{\rho-1}.
\]
对正实扭曲支 \(u>0\)，必有 \(\rho>1+\sqrt2\)，故上式分母不为零且右端为正。将 \(u(\rho)\) 代回 \(\alpha(\theta)\) 的表达式并化简，即得
\eqref{eq:xi-time-part9p-collision-frequency-rho-param} 中的 \(\alpha(\rho)\)。

为证明 \eqref{eq:xi-time-part9p-collision-frequency-algebraic-curve}，只需消去参数 \(\rho\)。把
\[
u(\rho-1)-\rho(\rho^2-2\rho-1)=0
\]
与
\[
\alpha(2\rho^3-5\rho^2+4\rho+1)-(\rho^2-2\rho-1)(\rho-1)=0
\]
视为关于 \(\rho\) 的两条代数方程，取 Sylvester 结果式即可得到
\eqref{eq:xi-time-part9p-collision-frequency-algebraic-curve}（结果式只差一个非零常数因子，不影响零点集）。

最后，由命题 \ref{prop:real-input-40-collision-all-real} 的全实谱性与有限状态 Perron 压力的一般谱理论，\(P_2\) 在实轴上实解析且严格凸；因此 Legendre--Fenchel 对偶在驻点
\(
a=P_2'(\theta)
\)
处取到。代入
\[
\theta=\log u(\rho),
\qquad
P_2(\theta)-P_2(0)=\log\rho-\log\varphi^2
\]
便得到 \eqref{eq:xi-time-part9p-collision-frequency-rate-param}。证毕。
\end{proof}

\begin{theorem}[碰撞势 Perron 支的正实有理参数化与单调双射]\label{thm:xi-time-part9p-collision-pressure-rational-monotone-branch}
\leanverified{paper\_xi\_time\_part9p\_collision\_pressure\_rational\_monotone\_branch}
沿用定理 \ref{thm:xi-time-part9p-collision-frequency-algebraic-ldp} 的记号。正实 Perron 支满足
\[
\boxed{\;
u(\rho)=\frac{\rho(\rho^2-2\rho-1)}{\rho-1}.\;}
\]
并且
\[
\boxed{\;
u>0
\iff
\rho>1+\sqrt2.\;}
\]
进一步，
\[
\boxed{\;
u'(\rho)=\frac{2\rho^3-5\rho^2+4\rho+1}{(\rho-1)^2}>0
\qquad(\rho>1+\sqrt2).\;}
\]
因此
\[
u:(1+\sqrt2,\infty)\longrightarrow (0,\infty)
\]
为严格递增双射，其逆映射正是正实扭曲参数 \(u\mapsto \rho(u)\) 的 Perron 支。
\end{theorem}
\begin{proof}
定理 \ref{thm:xi-time-part9p-collision-frequency-algebraic-ldp} 已给出
\[
u(\rho)=\frac{\rho(\rho^2-2\rho-1)}{\rho-1}.
\]
在 \(\rho>1\) 上，分母为正，故 \(u(\rho)>0\) 当且仅当
\[
\rho^2-2\rho-1>0,
\]
也即
\[
\rho>1+\sqrt2.
\]
这就得到了正实支的精确参数区间。

再对 \(u(\rho)\) 直接求导，得到
\[
u'(\rho)=\frac{2\rho^3-5\rho^2+4\rho+1}{(\rho-1)^2}.
\]
而分子可重写为
\[
2\rho^3-5\rho^2+4\rho+1
=
(\rho^2-2\rho-1)(2\rho-1)+4\rho.
\]
当 \(\rho>1+\sqrt2\) 时，上式右端各项都严格为正，故 \(u'(\rho)>0\)。

最后，
\[
\lim_{\rho\downarrow 1+\sqrt2}u(\rho)=0,
\qquad
\lim_{\rho\to\infty}u(\rho)=\infty,
\]
再结合严格单调性，即得 \(u:(1+\sqrt2,\infty)\to(0,\infty)\) 为双射。证毕。
\end{proof}

\begin{theorem}[碰撞势压力的显式导数公式与全轴严格凸性]\label{thm:xi-time-part9p-collision-pressure-derivatives-global-convexity}
沿用定理 \ref{thm:xi-time-part9p-collision-frequency-algebraic-ldp} 与
\ref{thm:xi-time-part9p-collision-pressure-rational-monotone-branch} 的记号。若
\[
\rho=\rho(e^\theta)\in(1+\sqrt2,\infty),
\]
则
\[
\boxed{\;
P_2'(\theta)=
\frac{(\rho-1)(\rho^2-2\rho-1)}
{2\rho^3-5\rho^2+4\rho+1}.\;}
\]
并且
\[
\boxed{\;
P_2''(\theta)=
\frac{\rho(\rho-1)(\rho^2-2\rho-1)\bigl(\rho^4+4\rho^3-10\rho^2+4\rho-3\bigr)}
{\bigl(2\rho^3-5\rho^2+4\rho+1\bigr)^3}.\;}
\]
特别地，对全部 \(\theta\in\RR\) 都有
\[
\boxed{\;
P_2''(\theta)>0.\;}
\]
因此 \(P_2\) 在整个实轴上严格凸。又在 \(\theta=0\) 时 \(u=1\) 且 \(\rho=\varphi^2\)，故
\[
P_2'(0)=\frac{3-\sqrt5}{10},
\qquad
P_2''(0)=\frac{6\sqrt5-5}{125},
\]
与零点处的既有 \(\sqrt5\)-闭式完全一致。
\end{theorem}
\begin{proof}
第一式正是定理 \ref{thm:xi-time-part9p-collision-frequency-algebraic-ldp} 中的
\(\alpha(\rho)\) 参数化，因为
\[
P_2'(\theta)=\alpha(\theta)=\alpha(\rho).
\]

再对该有理函数按 \(\rho\) 求导，可得
\[
\frac{\dd}{\dd\rho}P_2'(\theta(\rho))
=
\frac{\rho^4+4\rho^3-10\rho^2+4\rho-3}
{\bigl(2\rho^3-5\rho^2+4\rho+1\bigr)^2}.
\]
另一方面，由
\[
\theta=\log u(\rho)
\]
与定理 \ref{thm:xi-time-part9p-collision-pressure-rational-monotone-branch}，
\[
\frac{\dd}{\dd\theta}
=
\frac{u(\rho)}{u'(\rho)}\frac{\dd}{\dd\rho}
=
\frac{\rho(\rho-1)(\rho^2-2\rho-1)}
{2\rho^3-5\rho^2+4\rho+1}\frac{\dd}{\dd\rho}.
\]
将两式相乘，便得到 \(P_2''(\theta)\) 的闭式表达。

为证正性，只需检查各因子在 \(\rho>1+\sqrt2\) 上都严格为正。由上一条定理，
\[
\rho>0,
\qquad
\rho-1>0,
\qquad
\rho^2-2\rho-1>0,
\qquad
2\rho^3-5\rho^2+4\rho+1>0.
\]
对新出现的四次因子，可写成
\[
\rho^4+4\rho^3-10\rho^2+4\rho-3
=
(\rho^2-2\rho-1)(\rho^2+6\rho+3)+16\rho,
\]
故在 \(\rho>1+\sqrt2\) 上亦严格为正。因此 \(P_2''(\theta)>0\) 对全部 \(\theta\in\RR\) 成立，严格凸性随之得到。

最后，在 \(\theta=0\) 时 \(u=e^\theta=1\)。由
\[
u(\rho)=\frac{\rho(\rho^2-2\rho-1)}{\rho-1}
\]
与正实 Perron 支的唯一性，必有 \(\rho=\varphi^2\)。将 \(\rho=\varphi^2\) 代入上两式并化简，即得
\[
P_2'(0)=\frac{3-\sqrt5}{10},
\qquad
P_2''(0)=\frac{6\sqrt5-5}{125}.
\]
证毕。
\end{proof}

\begin{theorem}[碰撞势上尾物理支的大偏差率函数单分支统一化]\label{thm:xi-time-part9p-collision-frequency-physical-branch-monotonicity}
沿用定理 \ref{thm:xi-time-part9p-collision-frequency-algebraic-ldp} 的记号，并记
\[
\alpha_\ast:=P_2'(0)=\frac{3-\sqrt5}{10}.
\]
在经过 \(u=1\) 的物理上尾支
\[
\rho\ge \varphi^2
\]
上，
\[
\alpha(\rho)
=
\frac{(\rho^2-2\rho-1)(\rho-1)}{2\rho^3-5\rho^2+4\rho+1}
=
\frac{\rho^3-3\rho^2+\rho+1}{2\rho^3-5\rho^2+4\rho+1}
\]
严格单调递增。因而
\[
\alpha\bigl([\varphi^2,\infty)\bigr)
=
\left[\alpha_\ast,\frac12\right),
\]
并且对每个
\(
a\in[\alpha_\ast,\frac12)
\)
恰有唯一的 \(\rho\in[\varphi^2,\infty)\) 使 \(a=\alpha(\rho)\)。对应速率函数由单一实代数参数 \(\rho\) 统一给出：
\[
I_2\bigl(\alpha(\rho)\bigr)
=
\alpha(\rho)\log\!\frac{\rho(\rho^2-2\rho-1)}{\rho-1}
-\log\!\frac{\rho}{\varphi^2},
\qquad \rho\in[\varphi^2,\infty).
\]
特别地，物理上尾支不存在内部回折或分支切换。
\end{theorem}
\begin{proof}
由定理 \ref{thm:xi-time-part9p-collision-frequency-algebraic-ldp}，
\[
\alpha(\rho)=\frac{\rho^3-3\rho^2+\rho+1}{2\rho^3-5\rho^2+4\rho+1}.
\]
直接求导得
\[
\alpha'(\rho)
=
\frac{\rho^4+4\rho^3-10\rho^2+4\rho-3}{(2\rho^3-5\rho^2+4\rho+1)^2}.
\]
记
\[
q(\rho):=\rho^4+4\rho^3-10\rho^2+4\rho-3.
\]
则
\[
q(2)=13>0,
\qquad
q'(\rho)=4\rho^3+12\rho^2-20\rho+4>0
\qquad(\rho\ge 2).
\]
故 \(q(\rho)>0\) 对所有 \(\rho\ge 2\) 成立。另一方面，记
\[
d(\rho):=2\rho^3-5\rho^2+4\rho+1.
\]
则
\[
d'(\rho)=2(3\rho-2)(\rho-1)>0
\qquad(\rho\ge \varphi^2),
\]
并且
\[
d(\varphi^2)=5\varphi^2>0.
\]
因此 \(d(\rho)>0\) 在 \([\varphi^2,\infty)\) 上恒成立，分母平方严格为正，从而
\[
\alpha'(\rho)>0
\qquad(\rho\ge \varphi^2).
\]

再由 \(u(\varphi^2)=1\) 与
\(
\alpha_\ast=P_2'(0)
\)
可知
\[
\alpha(\varphi^2)=\alpha_\ast=\frac{3-\sqrt5}{10}.
\]
同时，
\[
\alpha(\rho)
=
\frac{\rho^3-3\rho^2+\rho+1}{2\rho^3-5\rho^2+4\rho+1}
\longrightarrow \frac12
\qquad(\rho\to\infty).
\]
故其像区间恰为 \([\alpha_\ast,\frac12)\)，且每个可达 \(a\) 对应唯一参数 \(\rho\)。最后，速率函数公式正是
\eqref{eq:xi-time-part9p-collision-frequency-rate-param} 在该区间上的限制。证毕。
\end{proof}

\begin{theorem}[实正扭曲分支的全局单调坐标化]\label{thm:xi-time-part9p-positive-branch-global-coordinate}
\leanverified{paper\_xi\_time\_part9p\_positive\_branch\_global\_coordinate}
记
\[
\rho_0:=1+\sqrt2.
\]
则 \eqref{eq:xi-time-part9p-collision-frequency-rho-param} 中的参数函数在区间 \((\rho_0,\infty)\) 上满足：
\begin{enumerate}
  \item \(u(\rho)\) 严格递增，并给出双射
  \[
  u:(\rho_0,\infty)\xrightarrow{\sim}(0,\infty);
  \]
  \item \(\alpha(\rho)\) 严格递增，并给出双射
  \[
  \alpha:(\rho_0,\infty)\xrightarrow{\sim}\left(0,\frac12\right).
  \]
\end{enumerate}
从而对每个实正扭曲参数 \(u>0\)，存在唯一 \(\rho>\rho_0\) 使
\[
u=\frac{\rho(\rho^2-2\rho-1)}{\rho-1},
\]
并且存在唯一碰撞频率 \(\alpha(u)\in(0,\tfrac12)\) 与之对应。亦即，\(\rho\)、\(u\) 与物理实轴上的碰撞频率在实正支上给出同一条全局一维坐标。
\end{theorem}
\begin{proof}
由 \eqref{eq:xi-time-part9p-collision-frequency-rho-param}，
\[
u'(\rho)=\frac{2\rho^3-5\rho^2+4\rho+1}{(\rho-1)^2}.
\]
记
\[
q(\rho):=2\rho^3-5\rho^2+4\rho+1.
\]
则
\[
q'(\rho)=6\rho^2-10\rho+4=2(3\rho-2)(\rho-1)>0
\qquad(\rho>1),
\]
故 \(q\) 在 \((1,\infty)\) 上严格递增。又
\[
q(\rho_0)=4+4\sqrt2>0,
\qquad
u(\rho_0)=0,
\qquad
u(\rho)\sim \rho^2\to+\infty
\quad(\rho\to+\infty),
\]
从而对一切 \(\rho>\rho_0\) 都有 \(u'(\rho)>0\)，并且 \(u\) 把 \((\rho_0,\infty)\) 严格单调地双射到 \((0,\infty)\)。

另一方面，
\[
\alpha'(\rho)=
\frac{\rho^4+4\rho^3-10\rho^2+4\rho-3}
{(2\rho^3-5\rho^2+4\rho+1)^2}.
\]
记
\[
p(\rho):=\rho^4+4\rho^3-10\rho^2+4\rho-3.
\]
则
\[
p''(\rho)=12\rho^2+24\rho-20>0
\qquad(\rho\ge \rho_0),
\]
并且
\[
p'(\rho_0)=48+24\sqrt2>0,
\qquad
p(\rho_0)=16+16\sqrt2>0.
\]
因此 \(p'(\rho)>0\) 且 \(p(\rho)>0\) 对所有 \(\rho>\rho_0\) 成立，故 \(\alpha'(\rho)>0\)。
再由
\[
\alpha(\rho_0)=0,
\qquad
\lim_{\rho\to+\infty}\alpha(\rho)=\frac12,
\]
可知 \(\alpha\) 把 \((\rho_0,\infty)\) 严格单调地双射到 \((0,\tfrac12)\)。
两条单调双射合并后，唯一性与全局坐标化结论立即成立。上述导数因子分解与端点值可由脚本 \texttt{scripts/exp\_real\_input\_40\_collision\_frequency\_prime\_shadow\_audit.py} 逐项复算。证毕。
\end{proof}

\begin{corollary}[RH/Ramanujan 窗口的碰撞频率端点为五次代数数]\label{cor:xi-time-part9p-rh-window-frequency-endpoint}
令 \(u_R>1\) 为命题 \ref{prop:real-input-40-collision-rhk-window} 给出的唯一临界参数，并记
\[
\alpha_R:=\alpha(\log u_R).
\]
则 \(\alpha_R\) 是区间 \((0,\tfrac12)\) 内唯一实根，满足五次方程
\begin{equation}\label{eq:xi-time-part9p-alphaR-quintic}
191\alpha^5-701\alpha^4+1043\alpha^3-748\alpha^2+252\alpha-28=0.
\end{equation}
数值上，
\[
\alpha_R\approx 0.204230263243173.
\]
并且 RH/Ramanujan 窗口中的全部可实现碰撞频率恰为
\[
[0,\alpha_R].
\]
\end{corollary}
\begin{proof}
命题 \ref{prop:real-input-40-collision-rhk-window} 已给出 \(u_R\) 是核版 RH/Ramanujan 条件在 \(u>0\) 上的唯一正端点。另一方面，由定理
\ref{thm:xi-time-part9p-positive-branch-global-coordinate}，\(\alpha\) 作为 \(u>0\) 的函数严格递增，并把 \((0,\infty)\) 双射到 \((0,\tfrac12)\)；因此 \(u\in(0,u_R]\) 所对应的频率闭包恰为 \([0,\alpha_R]\)。

为求 \(\alpha_R\) 的代数方程，把
\eqref{eq:xi-time-part9p-collision-frequency-algebraic-curve} 与五次方程
\[
u^5+4u^4+3u^3-96u^2-42u-14=0
\]
对 \(u\) 消元。所得 \(15\) 次多项式分解为三个五次因子；其中与正实 Perron 分支兼容、并且在 \((0,\tfrac12)\) 内有实根者，恰为
\eqref{eq:xi-time-part9p-alphaR-quintic}。由上一段的单调性，该根只能是 \(\alpha_R\)，并且在 \((0,\tfrac12)\) 内唯一。数值近似由同一代数方程直接求根即得。证毕。
\end{proof}

\begin{theorem}[奇素数扭曲的 prime-shadow Euler 乘积绝对收敛]\label{thm:xi-time-part9p-odd-prime-shadow-product}
沿用命题 \ref{prop:real-input-40-collision-pressure} 的记号，并记
\[
\lambda:=\rho(1)=\varphi^2,
\qquad
t_p:=\frac{2\pi}{p}\qquad (p\ge 3\ \text{为奇素数}).
\]
由于 \(t_p<2\pi/3<R_\theta\)（推论 \ref{cor:real-input-40-collision-branch-radius}），Perron 分支 \(\rho(e^{it_p})\) 沿实轴良定义。定义
\[
\mathfrak P_{\mathrm{odd}}
:=
\prod_{p\ge 3}
\left|
\frac{\rho(e^{it_p})}{\lambda}
\right|.
\]
则该乘积绝对收敛到一个正实常数，并且
\begin{equation}\label{eq:xi-time-part9p-odd-prime-shadow-expansion}
\log \mathfrak P_{\mathrm{odd}}
=
-\frac{6\sqrt5-5}{250}(2\pi)^2P_{\mathrm{odd}}(2)
+\frac{7+24\sqrt5}{75000}(2\pi)^4P_{\mathrm{odd}}(4)
+O\!\bigl(P_{\mathrm{odd}}(6)\bigr),
\end{equation}
其中
\[
P_{\mathrm{odd}}(s):=\sum_{p\ge 3}p^{-s}.
\]
特别地，
\[
0<\mathfrak P_{\mathrm{odd}}<1.
\]
\end{theorem}
\begin{proof}
由命题 \ref{prop:arity-collision-quadratic-closed}，当 \(t\to 0\) 时有
\[
\log\left|\frac{\rho(e^{it})}{\lambda}\right|
=
-\frac{6\sqrt5-5}{250}\,t^2
+\frac{7+24\sqrt5}{75000}\,t^4
+O(t^6).
\]
取 \(t=t_p=2\pi/p\) 即得
\[
\log\left|\frac{\rho(e^{it_p})}{\lambda}\right|
=
-\frac{6\sqrt5-5}{250}\left(\frac{2\pi}{p}\right)^2
+\frac{7+24\sqrt5}{75000}\left(\frac{2\pi}{p}\right)^4
+O(p^{-6}).
\]
由于 \(\sum_{p\ge 3}p^{-2}\)、\(\sum_{p\ge 3}p^{-4}\) 与 \(\sum_{p\ge 3}p^{-6}\) 全部收敛，级数
\[
\sum_{p\ge 3}\log\left|\frac{\rho(e^{it_p})}{\lambda}\right|
\]
绝对收敛，并且逐项求和即得到
\eqref{eq:xi-time-part9p-odd-prime-shadow-expansion}。因此 Euler 乘积收敛到某个正实常数。

为证明严格上界 \(\mathfrak P_{\mathrm{odd}}<1\)，注意 \(|W(e^{it})|=A\) 且 \(A\) 为原始非负矩阵。由谱半径比较，矩阵 \(W(e^{it})\) 的谱半径不超过 \(\rho(A)=\lambda\)，故作为其本征值的 \(\rho(e^{it})\) 满足
\[
\left|\rho(e^{it})\right|\le \lambda.
\]
若某个 \(t\notin 2\pi\ZZ\) 取等，则 Wielandt 型等号情形迫使 \(W(e^{it})\) 与 \(A\) 只差一个全局相位及对角酉共轭。比较长度 \(1\) 的闭路 \(00\to 00\) 与长度 \(2\) 的闭路 \(00\to 11\to 00\) 可知该全局相位必为 \(1\)，同时又必须满足 \(e^{it}=1\)，矛盾。故对全部奇素数 \(p\) 都有
\[
\left|\rho(e^{it_p})\right|<\lambda.
\]
于是每个乘子都位于 \((0,1)\)，从而极限常数亦满足 \(0<\mathfrak P_{\mathrm{odd}}<1\)。证毕。
\end{proof}

\begin{theorem}[Rényi--2 几何下的均匀化障碍]\label{thm:xi-time-part9p-renyi2-obstruction}
\leanverified{paper\_xi\_time\_part9p\_renyi2\_obstruction}
令 \(p_m^{\mathrm{bin}}\) 为 bin-fold 推前分布，\(u_m\) 为 \(X_m\) 上的均匀分布。则 Rényi--\(2\) 散度满足精确恒等式
\[
D_2\bigl(p_m^{\mathrm{bin}}\Vert u_m\bigr)
=
\log\!\bigl(F_{m+2}\,\mathsf{Col}_m\bigr),
\]
并且
\[
D_2\bigl(p_m^{\mathrm{bin}}\Vert u_m\bigr)
\ge
\log\!\bigl(1+2C_\varphi^2\bigr)+o(1).
\]
因此
\[
\liminf_{m\to\infty}D_2\bigl(p_m^{\mathrm{bin}}\Vert u_m\bigr)
\ge
\log\!\bigl(1+2C_\varphi^2\bigr)
>
0,
\]
从而 bin-fold 推前分布在 \(\chi^2\)/Rényi--\(2\) 意义下不可能趋于均匀分布。
\end{theorem}
\begin{proof}
由 \(u_m(w)\equiv 1/F_{m+2}\) 与 Rényi--\(2\) 散度定义，
\[
D_2\bigl(p_m^{\mathrm{bin}}\Vert u_m\bigr)
=
\log\sum_{w\in X_m}\frac{\bigl(p_m^{\mathrm{bin}}(w)\bigr)^2}{u_m(w)}
=
\log\!\left(F_{m+2}\sum_{w\in X_m}\bigl(p_m^{\mathrm{bin}}(w)\bigr)^2\right).
\]
而
\[
\sum_{w\in X_m}\bigl(p_m^{\mathrm{bin}}(w)\bigr)^2
=
\frac{1}{2^{2m}}\sum_{w\in X_m}\bigl(d_m^{\mathrm{bin}}(w)\bigr)^2
=
\mathsf{Col}_m,
\]
故第一条恒等式成立。

再由推论 \ref{cor:fold-critical-resonance-gap}，
\[
\mathsf{Col}_m
\ge
\frac{1+2C_\varphi^2+o(1)}{F_{m+2}}.
\]
代回即可得到
\[
D_2\bigl(p_m^{\mathrm{bin}}\Vert u_m\bigr)
\ge
\log\!\bigl(1+2C_\varphi^2+o(1)\bigr)
=
\log\!\bigl(1+2C_\varphi^2\bigr)+o(1).
\]
取下极限便得结论。证毕。
\end{proof}

\begin{proposition}[Fibonacci 共振双频给出二阶碰撞的最小刚性证书]\label{prop:xi-time-part9p-fibonacci-resonance-cut-collision}
沿用定义 \ref{def:fold-normalized-amplitude} 的记号，
\[
a_m(k):=\frac{|\widehat d_m(k)|}{2^m},
\qquad
\mathsf{Col}_m=\frac{1}{F_{m+2}}\sum_{k\in\ZZ/(F_{m+2}\ZZ)}a_m(k)^2.
\]
定义去掉两条 Fibonacci 共振频率后的截断碰撞量
\[
\mathsf{Col}_m^{\mathrm{cut}}
:=
\frac{1}{F_{m+2}}
\sum_{\substack{k\in\ZZ/(F_{m+2}\ZZ)\\ k\notin\{F_m,F_{m+1}\}}}
a_m(k)^2.
\]
则有精确差值律
\[
\mathsf{Col}_m-\mathsf{Col}_m^{\mathrm{cut}}
=
\frac{a_m(F_m)^2+a_m(F_{m+1})^2}{F_{m+2}}
=
\frac{2C_\varphi^2+o(1)}{F_{m+2}}.
\]
因此，任何忽略 \(\{F_m,F_{m+1}\}\) 的二阶碰撞估计都会在主尺度上损失一个不可恢复的常数份额 \(2C_\varphi^2\)。
\end{proposition}
\begin{proof}
由 \(\mathsf{Col}_m\) 与 \(\mathsf{Col}_m^{\mathrm{cut}}\) 的定义，二者之差恰为被删去的两项：
\[
\mathsf{Col}_m-\mathsf{Col}_m^{\mathrm{cut}}
=
\frac{a_m(F_m)^2+a_m(F_{m+1})^2}{F_{m+2}}.
\]
再由定理 \ref{thm:fold-critical-resonance-constant}，
\[
a_m(F_m)\to C_\varphi,
\qquad
a_m(F_{m+1})\to C_\varphi.
\]
将其代回即得
\[
\mathsf{Col}_m-\mathsf{Col}_m^{\mathrm{cut}}
=
\frac{2C_\varphi^2+o(1)}{F_{m+2}}.
\]
证毕。
\end{proof}

\begin{remark}[碰撞频率端点与奇素数 shadow 的数值锚点]\label{rem:xi-time-part9p-frequency-prime-shadow-audit}
对 \(P\ge 3\)，记截断 odd-prime zeta 和
\[
P_{\mathrm{odd}}^{(\le P)}(s):=\sum_{\substack{3\le p\le P\\ p\ \mathrm{prime}}}p^{-s}.
\]
脚本 \texttt{scripts/exp\_real\_input\_40\_collision\_frequency\_prime\_shadow\_audit.py} 对 \(\alpha_R\) 与 odd-prime partial product 给出如下核验片段：
% Auto-generated; do not edit by hand.
\begin{equation}\label{eq:real_input_40_collision_frequency_prime_shadow_audit}
\begin{aligned}
\alpha_R &\approx 0.204230263243172496,\\
191\alpha_R^5-701\alpha_R^4+1043\alpha_R^3-748\alpha_R^2+252\alpha_R-28 &\approx 1.3494013e-79,\\
\mathfrak P_{\mathrm{odd}}^{(\le 2000)}:=\prod_{\substack{3\le p\le 2000\\ p\ \mathrm{prime}}}\left|\frac{\rho(e^{2\pi i/p})}{\varphi^2}\right| &\approx 0.787090286807732362,\\
\log \mathfrak P_{\mathrm{odd}}^{(\le 2000)} &\approx -0.239412314390460902,\\
-\frac{6\sqrt5-5}{250}(2\pi)^2P_{\mathrm{odd}}^{(\le 2000)}(2)+\frac{7+24\sqrt5}{75000}(2\pi)^4P_{\mathrm{odd}}^{(\le 2000)}(4) &\approx -0.250451591524040363.
\end{aligned}
\end{equation}

\end{remark}

\paragraph{Fourier conductor 壳层、真商模下降与 full-modulus 原始性}
在折叠多重度剖面的 Parseval 恒等式、群环乘积分解以及 Fibonacci 共振双频的常数级下界已经建立之后，还可把“是否能够下降到真商模”完全转写为 Fourier 支撑的 conductor 截断问题。

\begin{theorem}[折叠多重度剖面下降到真商模的 Fourier--conductor 充要条件]\label{thm:xi-fold-multiplicity-quotient-descent-fourier-conductor}
记
\[
M:=F_{m+2},
\qquad
d_m:\ZZ/(M\ZZ)\to\NN,
\qquad
\widehat d_m(k):=\sum_{r\in\ZZ/(M\ZZ)}d_m(r)e^{-2\pi i kr/M}.
\]
固定约数 \(D\mid M\)。则下列三条命题等价：
\begin{enumerate}
  \item 存在函数
  \[
  f_D:\ZZ/(D\ZZ)\to\NN
  \]
  使得
  \[
  d_m(r)=f_D(r\bmod D)
  \qquad
  (\forall r\in\ZZ/(M\ZZ));
  \]
  \item 对每个频率 \(k\in\ZZ/(M\ZZ)\)，只要
  \[
  k\not\equiv 0\pmod{M/D},
  \]
  就有
  \[
  \widehat d_m(k)=0;
  \]
  \item 若定义导出模
  \[
  \operatorname{cond}(k):=\frac{M}{\gcd(k,M)}
  \]
  及各壳层能量
  \[
  E_m(q):=\sum_{\operatorname{cond}(k)=q}|\widehat d_m(k)|^2
  \qquad (q\mid M),
  \]
  则
  \[
  E_m(q)=0
  \qquad
  \text{对一切 }q\nmid D.
  \]
\end{enumerate}
此外，不加任何额外假设时总有精确分解
\[
\mathsf{Col}_m
=
\frac1M+\frac{1}{2^{2m}M}\sum_{\substack{q\mid M\\ q>1}}E_m(q),
\]
而在上述等价条件成立时，上式退化为
\[
\mathsf{Col}_m
=
\frac1M+\frac{1}{2^{2m}M}\sum_{\substack{q\mid D\\ q>1}}E_m(q).
\]
\end{theorem}
\begin{proof}
先证 \((1)\Rightarrow(2)\)。若
\[
d_m(r)=f_D(r\bmod D),
\]
则把 \(r\) 唯一写成
\[
r=a+Db,
\qquad
0\le a<D,
\qquad
0\le b<\frac{M}{D}.
\]
于是
\[
\begin{aligned}
\widehat d_m(k)
&=
\sum_{a=0}^{D-1}\sum_{b=0}^{M/D-1}
f_D(a)\,e^{-2\pi i k(a+Db)/M}\\
&=
\left(\sum_{a=0}^{D-1}f_D(a)e^{-2\pi i ka/M}\right)
\left(\sum_{b=0}^{M/D-1}e^{-2\pi i kb/(M/D)}\right).
\end{aligned}
\]
第二个括号是长度 \(M/D\) 的几何和；当且仅当
\[
\frac{M}{D}\mid k
\]
时非零，否则为零。故 \((2)\) 成立。

再证 \((2)\Rightarrow(1)\)。Fourier 反演给出
\[
d_m(r)=\frac1M\sum_{k\in\ZZ/(M\ZZ)}\widehat d_m(k)e^{2\pi i kr/M}.
\]
在 \((2)\) 的假设下，只保留 \(k=(M/D)\ell\) 的项，可得
\[
\begin{aligned}
d_m(r)
&=
\frac1M\sum_{\ell=0}^{D-1}\widehat d_m\!\left(\frac{M}{D}\ell\right)
e^{2\pi i \ell r/D}.
\end{aligned}
\]
右端仅依赖于 \(r\bmod D\)，故 \(d_m\) 下降到 \(\ZZ/(D\ZZ)\)。

现证 \((2)\Leftrightarrow(3)\)。若
\[
k=\frac{M}{D}\ell,
\]
则
\[
\gcd(k,M)=\frac{M}{D}\gcd(\ell,D),
\qquad
\operatorname{cond}(k)=\frac{D}{\gcd(\ell,D)}\mid D.
\]
反之，若 \(\operatorname{cond}(k)\mid D\)，记
\[
g:=\gcd(k,M),
\qquad
\operatorname{cond}(k)=\frac{M}{g}.
\]
由 \(\frac{M}{g}\mid D\) 且 \(D\mid M\)，可写
\[
D=\frac{M}{g}\,t
\]
对某个整数 \(t\)。于是
\[
\frac{M}{D}=\frac{g}{t},
\]
从而 \(\frac{M}{D}\mid g\)。由于 \(g\mid k\)，遂有
\[
\frac{M}{D}\mid k.
\]
故
\[
\frac{M}{D}\mid k
\iff
\operatorname{cond}(k)\mid D,
\]
于是 \((2)\) 与 \((3)\) 等价。

最后证明碰撞分解。由定理 \ref{thm:fold-collision-spectrum}，
\[
\mathsf{Col}_m
=
\frac{1}{2^{2m}M}\sum_{k\in\ZZ/(M\ZZ)}|\widehat d_m(k)|^2.
\]
其中 \(k=0\) 一项恰给出 \(1/M\)，而非零频率按 \(\operatorname{cond}(k)\) 分组，便得到
\[
\mathsf{Col}_m
=
\frac1M+\frac{1}{2^{2m}M}\sum_{\substack{q\mid M\\ q>1}}E_m(q).
\]
若 \((1)\)--\((3)\) 成立，则所有 \(q\nmid D\) 的壳层均为零，故立即退化为
\[
\mathsf{Col}_m
=
\frac1M+\frac{1}{2^{2m}M}\sum_{\substack{q\mid D\\ q>1}}E_m(q).
\]
证毕。
\end{proof}
\leanverified{paper\_xi\_fold\_multiplicity\_quotient\_descent\_fourier\_conductor}

\begin{corollary}[大尺度下折叠多重度剖面的 full-modulus 原始性]\label{cor:xi-fold-multiplicity-fullmodulus-primitive}
存在整数 \(m_0\)，使得对所有 \(m\ge m_0\)，折叠多重度剖面
\[
d_m:\ZZ/(F_{m+2}\ZZ)\to\NN
\]
都不可能下降到任何真商模
\[
\ZZ/(D\ZZ),
\qquad
D\mid F_{m+2},
\qquad
D<F_{m+2}.
\]
等价地，对充分大的 \(m\)，其最小可见模数正是 \(F_{m+2}\) 本身。
\end{corollary}
\begin{proof}
由定理 \ref{thm:fold-critical-resonance-constant}，
\[
a_m(F_m):=\frac{|\widehat d_m(F_m)|}{2^m}\longrightarrow C_\varphi>0.
\]
故存在 \(m_0\)，使得当 \(m\ge m_0\) 时，
\[
\widehat d_m(F_m)\neq 0.
\]
另一方面，由 Fibonacci 最大公因子恒等式，
\[
\gcd(F_m,F_{m+2})=F_{\gcd(m,m+2)}=1,
\]
从而
\[
\operatorname{cond}(F_m)=\frac{F_{m+2}}{\gcd(F_m,F_{m+2})}=F_{m+2}.
\]
若 \(d_m\) 能下降到某个真商模 \(D<F_{m+2}\)，则由定理
\ref{thm:xi-fold-multiplicity-quotient-descent-fourier-conductor} 的条件 \((3)\)，所有
\[
q\nmid D
\]
的 conductor 壳层都必须消失。特别地，full-conductor 壳层
\[
q=F_{m+2}
\]
必须为零，这与 \(\widehat d_m(F_m)\neq 0\) 矛盾。故不存在这样的真商模 \(D\)。证毕。
\end{proof}

\begin{theorem}[full-conductor 壳层的碰撞缺口下界由 Fibonacci 共振双频承担]\label{thm:xi-full-conductor-shell-collision-gap-fibonacci-pair}
记
\[
M:=F_{m+2},
\qquad
E_m^{\mathrm{full}}:=E_m(M)=\sum_{\gcd(k,M)=1}|\widehat d_m(k)|^2,
\qquad
a_m(k):=\frac{|\widehat d_m(k)|}{2^m}.
\]
则有
\[
E_m^{\mathrm{full}}
\ge
|\widehat d_m(F_m)|^2+|\widehat d_m(F_{m+1})|^2
=
2^{2m}\bigl(a_m(F_m)^2+a_m(F_{m+1})^2\bigr).
\]
从而
\[
\mathsf{Col}_m-\frac1M
\ge
\frac{E_m^{\mathrm{full}}}{2^{2m}M}
\ge
\frac{a_m(F_m)^2+a_m(F_{m+1})^2}{M}
=
\frac{2C_\varphi^2+o(1)}{M}.
\]
特别地，碰撞缺口的主尺度下界已经完全由 full-conductor 壳层内的共轭 Fibonacci 双频承担。
\end{theorem}
\begin{proof}
由定理 \ref{thm:xi-fold-multiplicity-quotient-descent-fourier-conductor} 的壳层分解，
\[
\mathsf{Col}_m-\frac1M
=
\frac{1}{2^{2m}M}\sum_{\substack{q\mid M\\ q>1}}E_m(q)
\ge
\frac{E_m^{\mathrm{full}}}{2^{2m}M}.
\]

又由 Fibonacci 最大公因子恒等式，
\[
\gcd(F_m,F_{m+2})=\gcd(F_{m+1},F_{m+2})=1,
\]
故频率 \(F_m\) 与 \(F_{m+1}\) 都落在 full-conductor 壳层内。于是
\[
E_m^{\mathrm{full}}
\ge
|\widehat d_m(F_m)|^2+|\widehat d_m(F_{m+1})|^2.
\]
再用
\[
|\widehat d_m(F_m)|=2^m a_m(F_m),
\qquad
|\widehat d_m(F_{m+1})|=2^m a_m(F_{m+1}),
\]
便得到第二个不等式。

最后，由定理 \ref{thm:fold-critical-resonance-constant}，
\[
a_m(F_m)\to C_\varphi,
\qquad
a_m(F_{m+1})\to C_\varphi,
\]
故
\[
\frac{a_m(F_m)^2+a_m(F_{m+1})^2}{M}
=
\frac{2C_\varphi^2+o(1)}{M}.
\]
证毕。
\end{proof}

\begin{theorem}[Fourier 幅度的中点门级联实现与临界共振半径]\label{thm:xi-time-part9p1-midpoint-gate-cascade-critical-radius}
沿用定理 \ref{thm:fold-spectrum-factorization} 的记号。对每个 \(m\ge 1\) 与 \(k\in\ZZ/(F_{m+2}\ZZ)\)，定义角序列
\[
\vartheta_j^{(m)}(k):=
2\pi k\frac{F_{j+1}}{F_{m+2}}
\qquad (1\le j\le m),
\]
以及复平面上的中点门
\[
\mathcal M_{\vartheta}(z):=\frac{z+e^{-\mathrm{i}\vartheta}z}{2}.
\]
再置
\[
z_0^{(m,k)}:=1,
\qquad
z_j^{(m,k)}:=\mathcal M_{\vartheta_j^{(m)}(k)}\bigl(z_{j-1}^{(m,k)}\bigr)
\qquad (1\le j\le m).
\]
则成立：
\begin{enumerate}
  \item 级联终值满足
  \[
  \boxed{\;
  z_m^{(m,k)}
  =
  2^{-m}\widehat d_m(k),\;}
  \]
  从而
  \[
  \boxed{\;
  |z_m^{(m,k)}|=a_m(k).\;}
  \]
  因此 \(a_m(k)\) 精确等同于角序列 \(\{\vartheta_j^{(m)}(k)\}_{j=1}^{m}\) 所驱动的中点门级联收缩半径；
  \item 有
  \[
  \widehat d_m(k)=0
  \iff
  z_m^{(m,k)}=0
  \iff
  \exists\, j\in\{1,\dots,m\}\ \text{使}\ 
  \vartheta_j^{(m)}(k)\equiv \pi\pmod{2\pi}.
  \]
  换言之，Fourier 零点恰对应于某一级中点门把一点与其对径旋转副本取平均而坍缩到原点；
  \item 对 Fibonacci 补频对 \(k=F_m\) 与 \(k=F_{m+1}\)，有
  \[
  |z_m^{(m,F_m)}|\longrightarrow C_\varphi,
  \qquad
  |z_m^{(m,F_{m+1})}|\longrightarrow C_\varphi.
  \]
  因而临界共振常数 \(C_\varphi\) 可解释为一条中点门级联在 Fibonacci 角驱动下保持的严格正极限半径。
\end{enumerate}
\end{theorem}
\begin{proof}
由定义，
\[
\mathcal M_{\vartheta}(z)=\frac{1+e^{-\mathrm{i}\vartheta}}{2}\,z.
\]
故递推展开给出
\[
z_m^{(m,k)}
=
\prod_{j=1}^{m}\frac{1+e^{-\mathrm{i}\vartheta_j^{(m)}(k)}}{2}
=
2^{-m}\prod_{j=1}^{m}\left(1+e^{-2\pi \mathrm{i}kF_{j+1}/F_{m+2}}\right).
\]
再由定理 \ref{thm:fold-spectrum-factorization}，
\[
\widehat d_m(k)=\prod_{j=1}^{m}\left(1+e^{-2\pi \mathrm{i}kF_{j+1}/F_{m+2}}\right),
\]
便得到
\[
z_m^{(m,k)}=2^{-m}\widehat d_m(k).
\]
取模即得
\[
|z_m^{(m,k)}|=\frac{|\widehat d_m(k)|}{2^m}=a_m(k).
\]

第二条由乘积分解立即推出：\(\widehat d_m(k)=0\) 当且仅当某一因子
\[
1+e^{-\mathrm{i}\vartheta_j^{(m)}(k)}
\]
为零，而这又等价于
\[
\vartheta_j^{(m)}(k)\equiv \pi\pmod{2\pi}.
\]
此时
\[
\mathcal M_{\vartheta_j^{(m)}(k)}(z)=0
\qquad (\forall z\in\CC),
\]
故级联在该步后永久坍缩到原点。

第三条正是定理 \ref{thm:fold-critical-resonance-constant} 的结论，因为
\[
|z_m^{(m,F_m)}|=a_m(F_m),
\qquad
|z_m^{(m,F_{m+1})}|=a_m(F_{m+1}).
\]
证毕。
\end{proof}

\endinput


在 Rényi--\(2\) 的均匀化障碍之外，还可把同一下界直接改写为 \(\chi^2\) 几何中的二维共振地板。

\begin{theorem}[均匀基线下 \texorpdfstring{$\chi^2$}{chi2} 偏离的 rank-two 共振地板]\label{thm:xi-time-part9p-chi2-ranktwo-resonant-floor}
\leanverified{paper\_xi\_time\_part9p\_chi2\_ranktwo\_resonant\_floor}
沿用定理 \ref{thm:xi-time-part9p-renyi2-obstruction} 与命题 \ref{prop:xi-time-part9p-fibonacci-resonance-cut-collision} 的记号。定义
\[
\chi^2\bigl(p_m^{\mathrm{bin}}\Vert u_m\bigr)
:=
\sum_{w\in X_m}
\frac{\bigl(p_m^{\mathrm{bin}}(w)-u_m(w)\bigr)^2}{u_m(w)}.
\]
则有恒等式
\[
\chi^2\bigl(p_m^{\mathrm{bin}}\Vert u_m\bigr)
=
F_{m+2}\,\mathsf{Col}_m-1
=
\frac{1}{2^{2m}}
\sum_{k\ne 0}\bigl|\widehat d_m(k)\bigr|^2.
\]
若进一步定义去掉两条 Fibonacci 共振频率后的截断量
\[
\chi^2_{\mathrm{cut}}(m)
:=
\frac{1}{2^{2m}}
\sum_{k\notin\{0,F_m,F_{m+1}\}}
\bigl|\widehat d_m(k)\bigr|^2,
\]
则
\[
\chi^2\bigl(p_m^{\mathrm{bin}}\Vert u_m\bigr)-\chi^2_{\mathrm{cut}}(m)
=
a_m(F_m)^2+a_m(F_{m+1})^2
=
2C_\varphi^2+o(1).
\]
特别地，
\[
\chi^2\bigl(p_m^{\mathrm{bin}}\Vert u_m\bigr)\ge 2C_\varphi^2+o(1).
\]
因此，bin-fold 推前分布对均匀分布的非平凡 \(\chi^2\) 偏离，已经由频率 \(F_m\) 与 \(F_{m+1}\) 所张成的二维共振子空间单独强制。
\end{theorem}
\begin{proof}
由定理 \ref{thm:xi-time-part9p-renyi2-obstruction}，
\[
D_2\bigl(p_m^{\mathrm{bin}}\Vert u_m\bigr)
=
\log\!\bigl(F_{m+2}\,\mathsf{Col}_m\bigr).
\]
另一方面，对均匀基线 \(u_m(w)\equiv 1/F_{m+2}\) 直接计算可得
\[
1+\chi^2\bigl(p_m^{\mathrm{bin}}\Vert u_m\bigr)
=
\sum_{w\in X_m}\frac{\bigl(p_m^{\mathrm{bin}}(w)\bigr)^2}{u_m(w)}
=
F_{m+2}\,\mathsf{Col}_m.
\]
因此
\[
\chi^2\bigl(p_m^{\mathrm{bin}}\Vert u_m\bigr)
=
F_{m+2}\,\mathsf{Col}_m-1.
\]

再由定理 \ref{thm:fold-collision-spectrum} 的 Parseval 恒等式，
\[
\mathsf{Col}_m
=
\frac{1}{2^{2m}F_{m+2}}
\sum_{k\in\ZZ/(F_{m+2}\ZZ)}
\bigl|\widehat d_m(k)\bigr|^2.
\]
其中 \(k=0\) 项满足 \(\widehat d_m(0)=2^m\)，故
\[
F_{m+2}\,\mathsf{Col}_m-1
=
\frac{1}{2^{2m}}
\sum_{k\ne 0}\bigl|\widehat d_m(k)\bigr|^2,
\]
得到第一组恒等式。

最后，由命题 \ref{prop:xi-time-part9p-fibonacci-resonance-cut-collision}，
\[
\mathsf{Col}_m-\mathsf{Col}_m^{\mathrm{cut}}
=
\frac{a_m(F_m)^2+a_m(F_{m+1})^2}{F_{m+2}}
=
\frac{2C_\varphi^2+o(1)}{F_{m+2}}.
\]
两边同乘 \(F_{m+2}\) 即得
\[
\chi^2\bigl(p_m^{\mathrm{bin}}\Vert u_m\bigr)-\chi^2_{\mathrm{cut}}(m)
=
a_m(F_m)^2+a_m(F_{m+1})^2
=
2C_\varphi^2+o(1),
\]
从而特别推出
\[
\chi^2\bigl(p_m^{\mathrm{bin}}\Vert u_m\bigr)\ge 2C_\varphi^2+o(1).
\]
证毕。
\end{proof}

\endinput

\paragraph{核版 RH 临界五次域、分歧三次域与均匀性缺口的闭合接口}
在碰撞位势的核版 RH/Ramanujan 窗口、分歧晶格与 Fibonacci 共振常数已经建立之后，还可继续抽取出下列五条彼此衔接的算术--解析后果。它们分别刻画临界参数的数域次数、与分歧三次域的线性无交、RH 边界的纯谱性质、RH 窗口内部仍不可消除的均匀性缺口，以及最近分歧格点对高阶累积量增长与振荡的解析支配。

\begin{theorem}[核版 RH 临界参数生成五次域]\label{thm:xi-time-part9pb-rh-critical-quintic-field}
\leanverified{paper\_xi\_time\_part9pb\_rh\_critical\_quintic\_field}
沿用命题 \ref{prop:real-input-40-collision-rhk-window} 的记号，记
\[
Q(u):=u^5+4u^4+3u^3-96u^2-42u-14.
\]
则 \(Q\) 在 \(\QQ[u]\) 上不可约。因而若 \(u_R>1\) 为命题
\ref{prop:real-input-40-collision-rhk-window} 给出的唯一正实临界参数，则
\[
\boxed{\ [\QQ(u_R):\QQ]=5.\ }
\]
\end{theorem}
\begin{proof}
先把 \(Q\) 模 \(5\) 约化，得到
\[
\overline Q(u)=u^5-u^4-2u^3-u^2-2u+1\in \FF_5[u].
\]
对 \(a\in\FF_5=\{0,1,2,3,4\}\) 逐一代入，可知 \(\overline Q(a)\neq 0\)，故
\(\overline Q\) 没有一次因子。再直接检验 \(\FF_5[u]\) 中全部不可约二次因子，均不整除 \(\overline Q\)。由于五次多项式若在域上可约，则必出现 \(1+4\) 或 \(2+3\) 型分解，故 \(\overline Q\) 在 \(\FF_5[u]\) 上不可约。

又因 \(Q\) 为首一整系数多项式，Gauss 引理推出 \(Q\) 在 \(\QQ[u]\) 上亦不可约。命题 \ref{prop:real-input-40-collision-rhk-window} 已给出 \(u_R\) 满足 \(Q(u_R)=0\)，故 \(u_R\) 的极小多项式只能是 \(Q\) 本身，于是
\[
[\QQ(u_R):\QQ]=\deg Q=5.
\]
证毕。
\end{proof}

\begin{theorem}[临界五次域与分歧三次域线性无交]\label{thm:xi-time-part9pb-rh-branch-cubic-linear-disjoint}
沿用推论 \ref{cor:real-input-40-collision-branch-radius} 与定理
\ref{thm:xi-time-part9pb-rh-critical-quintic-field} 的记号。记
\[
\Delta(u):=4u^3+25u^2+88u+8.
\]
任取 \(\Delta(u)=0\) 的一个根 \(u_{\mathrm{br}}\)，并设
\[
K_{\mathrm{br}}:=\QQ(u_{\mathrm{br}}),
\qquad
K_{\mathrm{RH}}:=\QQ(u_R).
\]
则
\[
\boxed{\ K_{\mathrm{br}}\cap K_{\mathrm{RH}}=\QQ,\qquad
[K_{\mathrm{br}}K_{\mathrm{RH}}:\QQ]=15.\ }
\]
特别地，\(K_{\mathrm{br}}\) 与 \(K_{\mathrm{RH}}\) 在 \(\QQ\) 上线性无交。
\end{theorem}
\begin{proof}
先证 \(\Delta\) 在 \(\QQ[u]\) 上不可约。由有理根判别，其有理根只可能来自
\[
\pm 1,\ \pm 2,\ \pm 4,\ \pm 8,\ \pm \frac12,\ \pm \frac14,\ \pm \frac18.
\]
逐一代入均不为零，故 \(\Delta\) 没有有理根。由于 \(\deg \Delta=3\)，这已推出 \(\Delta\) 在 \(\QQ[u]\) 上不可约，从而
\[
[K_{\mathrm{br}}:\QQ]=3.
\]

另一方面，定理 \ref{thm:xi-time-part9pb-rh-critical-quintic-field} 已给出
\[
[K_{\mathrm{RH}}:\QQ]=5.
\]
若记
\[
E:=K_{\mathrm{br}}\cap K_{\mathrm{RH}},
\]
则塔式定理给出
\[
[E:\QQ]\mid [K_{\mathrm{br}}:\QQ]=3,
\qquad
[E:\QQ]\mid [K_{\mathrm{RH}}:\QQ]=5.
\]
因 \(3\) 与 \(5\) 互素，只能有
\[
[E:\QQ]=1,
\]
即
\[
K_{\mathrm{br}}\cap K_{\mathrm{RH}}=\QQ.
\]
于是
\[
[K_{\mathrm{br}}K_{\mathrm{RH}}:\QQ]
=
\frac{[K_{\mathrm{br}}:\QQ]\,[K_{\mathrm{RH}}:\QQ]}
{[K_{\mathrm{br}}\cap K_{\mathrm{RH}}:\QQ]}
=
\frac{3\cdot 5}{1}
=
15.
\]
证毕。
\end{proof}

\begin{theorem}[核版 RH 的临界边界是解析谱相变而非分歧碰撞]\label{thm:xi-time-part9pb-rh-transition-analytic-not-branch}
沿用推论 \ref{cor:real-input-40-collision-branch-radius} 与命题
\ref{prop:real-input-40-collision-rhk-window} 的记号。记
\[
\Delta(u)=4u^3+25u^2+88u+8,
\]
\[
\theta_R:=\log u_R,
\qquad
R_\theta:=\min_j\min_{k\in\ZZ}\bigl|\log u_j+2\pi i k\bigr|,
\]
其中 \(u_j\) 遍历 \(\Delta(u)=0\) 的三个根。则：
\begin{enumerate}
  \item 正实轴上不存在任何代数分歧点；
  \item 核版 RH 的临界点 \(u_R\) 位于 Perron 分支的解析域内部，而不在分歧边界上；
  \item 其到最近复分歧格点的解析裕量满足
  \[
  \boxed{\ R_\theta-\theta_R
  =
  R_\theta-\log u_R
  \approx 1.48342064735593>0.\ }
  \]
\end{enumerate}
因此，核版 RH 的临界边界来自解析域内部的谱支穿越，而不是由代数分歧碰撞触发。
\end{theorem}
\begin{proof}
对任意 \(u>0\)，多项式 \(\Delta(u)=4u^3+25u^2+88u+8\) 的各项都严格为正，故
\[
\Delta(u)>0\qquad (u>0).
\]
因此 \(\Delta(u)=0\) 的分歧点不可能落在正实轴上，第一条成立。

再由推论 \ref{cor:real-input-40-collision-branch-radius}，\(P_2(\theta)=\log\rho(e^\theta)\) 在 \(\theta=0\) 的 Taylor 收敛半径正是 \(R_\theta\)，而全部代数分歧点恰为
\[
\theta=\log u_j+2\pi i k,\qquad k\in\ZZ.
\]
另一方面，命题 \ref{prop:real-input-40-collision-rhk-window} 给出
\[
\textsf{RH}(u)\Longleftrightarrow 0<u\le u_R,
\qquad
u_R\approx 3.592621813443688,
\]
故
\[
\theta_R=\log u_R\approx 1.27888224610494.
\]
将其与
\[
R_\theta\approx 2.76230289346087
\]
比较，得到
\[
R_\theta-\theta_R\approx 1.48342064735593>0.
\]
于是 \(\theta_R\) 严格位于解析圆盘 \(\{|\theta|<R_\theta\}\) 的内部，从而第二条与第三条同时成立。

最后，在 \(u=u_R\) 处命题 \ref{prop:real-input-40-collision-rhk-window} 给出临界等号
\[
\Lambda(u_R)=\lambda_1(u_R)^{1/2}.
\]
故 RH 边界由正则谱支之间的不等式穿越所决定，而非由分歧方程 \(\Delta(u)=0\) 的塌缩所决定。证毕。
\end{proof}

\begin{theorem}[RH 窗口与常数级均匀性缺口严格共存]\label{thm:xi-time-part9pb-rh-window-uniform-gap-coexistence}
沿用命题 \ref{prop:real-input-40-collision-rhk-window}、推论
\ref{cor:fold-critical-resonance-gap} 与定理 \ref{thm:fold-critical-resonance-constant} 的记号。则在未加权点 \(u=1\) 有
\[
\textsf{RH}(1)\ \text{成立},
\]
但同时
\[
\boxed{\ 
\liminf_{m\to\infty}
F_{m+2}\left(\mathsf{Col}_m-\frac{1}{F_{m+2}}\right)
\ge
2C_\varphi^2
>0.\ }
\]
等价地，
\[
\boxed{\ 
\mathsf{Col}_m
\ge
\frac{1+2C_\varphi^2+o(1)}{F_{m+2}},
\qquad
1+2C_\varphi^2\approx 1.1007267225141177.\ }
\]
其中
\[
C_\varphi\approx 0.2244178274047294,
\qquad
2C_\varphi^2\approx 0.1007267225141179.
\]
因此，尽管 \(u=1\) 位于 kernel RH/Ramanujan 窗口内部，碰撞概率相对于均匀基线 \(1/F_{m+2}\) 仍保留一个不可消除的常数级缺口。
\end{theorem}
\begin{proof}
由命题 \ref{prop:real-input-40-collision-rhk-window}，
\[
\textsf{RH}(u)\Longleftrightarrow 0<u\le u_R,
\qquad
u_R>1.
\]
故 \(u=1\) 的确位于 RH 窗口内部，从而 \(\textsf{RH}(1)\) 成立。

另一方面，推论 \ref{cor:fold-critical-resonance-gap} 已给出
\[
\liminf_{m\to\infty}
F_{m+2}\left(\mathsf{Col}_m-\frac{1}{F_{m+2}}\right)
\ge
2C_\varphi^2,
\]
而定理 \ref{thm:fold-critical-resonance-constant} 给出 \(C_\varphi>0\)，故右端严格为正。将不等式改写，即得
\[
\mathsf{Col}_m
\ge
\frac{1+2C_\varphi^2+o(1)}{F_{m+2}}.
\]
代入现有数值近似
\[
C_\varphi\approx 0.2244178274047294
\]
便得到
\[
2C_\varphi^2\approx 0.1007267225141179,
\qquad
1+2C_\varphi^2\approx 1.1007267225141177.
\]
证毕。
\end{proof}

\begin{theorem}[最近分歧格点控制碰撞累积量的 Gevrey--1 增长与振荡渐近]\label{thm:xi-time-part9pb-branch-lattice-cumulant-gevrey-oscillation}
沿用推论 \ref{cor:real-input-40-collision-branch-radius} 与定理
\ref{thm:xi-time-part9pb-rh-transition-analytic-not-branch} 的记号。记 \(P_2(\theta)\) 为
\(\theta=0\) 处由 Perron 根定义的解析芽；在实轴上它与
\[
\log\rho(e^\theta)
\]
一致。再记
\[
\kappa_n:=P_2^{(n)}(0),
\qquad
R_\theta:=\min_j\min_{k\in\ZZ}\bigl|\log u_j+2\pi i k\bigr|.
\]
设最近分歧格点由一对共轭点
\[
\theta_\star,\ \overline{\theta_\star}
\]
实现，其中
\[
|\theta_\star|=R_\theta,
\qquad
\Im \theta_\star\neq 0,
\]
并且 Perron 解析支在这两个点处都呈简单平方根分歧。则：
\begin{enumerate}
  \item 存在常数 \(C>0\)，使得对一切 \(n\ge 1\) 都有
  \[
  \boxed{\;
  |\kappa_{2n}|
  \le
  C\,(2n)!\,R_\theta^{-2n}.\;}
  \]
  \item 存在复常数 \(c_\star\neq 0\)，使得当 \(n\to\infty\) 时，
  \[
  \boxed{\;
  \kappa_{2n}
  =
  2\Re\!\Bigl(
  c_\star\,(2n)!\,\theta_\star^{-2n}\,n^{-3/2}
  \Bigr)
  +O\!\bigl((2n)!\,R_\theta^{-2n}\,n^{-5/2}\bigr).\;}
  \]
\end{enumerate}
从而偶阶累积量的符号与幅度都呈现由 \(\arg\theta_\star\) 决定的确定性振荡；最近分歧晶格在实轴喷射上的干涉条纹由此直接外显。
\end{theorem}
\begin{proof}
由推论 \ref{cor:real-input-40-collision-branch-radius}，\(P_2\) 在圆盘
\[
|\theta|<R_\theta
\]
上全纯。于是对任意 \(0<r<R_\theta\)，由 Cauchy 估计可得
\[
|\kappa_{2n}|
=
\bigl|P_2^{(2n)}(0)\bigr|
\le
\frac{(2n)!}{r^{2n}}\sup_{|\theta|=r}|P_2(\theta)|.
\]
令 \(r\uparrow R_\theta\)，并把圆周上的上确界吸收到常数中，即得第一条 Gevrey--1 上界。

再看第二条。由简单平方根分歧假设，存在 \(P_2\) 在两处分歧点附近的局部展开
\[
P_2(\theta)
=
H_+(\theta)+b_+\Bigl(1-\frac{\theta}{\theta_\star}\Bigr)^{1/2},
\]
\[
P_2(\theta)
=
H_-(\theta)+b_-\Bigl(1-\frac{\theta}{\overline{\theta_\star}}\Bigr)^{1/2},
\]
其中 \(H_\pm\) 全纯，且 \(b_\pm\neq 0\)。由于这对共轭点是全部最近奇点，标准分支点 Darboux 传递给出 Taylor 系数渐近
\[
[\theta^N]P_2(\theta)
=
\frac{b_+}{\Gamma(-1/2)}\,N^{-3/2}\theta_\star^{-N}
+\frac{b_-}{\Gamma(-1/2)}\,N^{-3/2}\overline{\theta_\star}^{-N}
+O(R_\theta^{-N}N^{-5/2}).
\]
又因 \(P_2\) 在实轴上取实值，其 Taylor 系数皆为实数，故 \(b_-=\overline{b_+}\)。将 \(N=2n\) 代入并把因子 \(2^{-3/2}\) 吸收到常数中，令
\[
c_\star:=\frac{2^{-3/2}b_+}{\Gamma(-1/2)},
\]
再用
\[
\kappa_{2n}=(2n)!\,[\theta^{2n}]P_2(\theta),
\]
便得到
\[
\kappa_{2n}
=
2\Re\!\Bigl(
c_\star\,(2n)!\,\theta_\star^{-2n}\,n^{-3/2}
\Bigr)
+O\!\bigl((2n)!\,R_\theta^{-2n}\,n^{-5/2}\bigr).
\]
这就证明了第二条；振荡性则由 \(\theta_\star^{-2n}=R_\theta^{-2n}e^{-2in\arg\theta_\star}\) 直接给出。证毕。
\end{proof}

\noindent
由此，纯碰撞位势的 RH 临界参数不但生成一个独立的五次域，而且与分歧三次域在线性代数层面完全分离；同时，这一 RH 边界在解析几何上属于谱支穿越而非分歧塌缩，并且即便在最自然的未加权点 \(u=1\) 处，RH 窗口也并不迫使碰撞统计回到均匀基线。进一步，最近复分歧格点还把高阶累积量的 Gevrey--1 增长与振荡喷射直接锁定在同一解析半径上。于是算术次数、复分歧、局部喷射与宏观非均匀性四者在同一接口上获得了严格闭合。

\endinput

\paragraph{碰撞频率饱和端点的显式速率尺度与有限边界代价}
沿用定理 \ref{thm:xi-time-part9p-collision-frequency-algebraic-ldp}、
\ref{thm:xi-time-part9p-collision-frequency-physical-branch-monotonicity} 与
\ref{thm:xi-time-part9p-positive-branch-global-coordinate} 的记号。下述两条把物理上尾支 \(\alpha\uparrow \frac12\) 的趋近机制压缩为显式的指数尺度与边界代价公式。

\begin{theorem}[碰撞频率饱和端点的精确 \texorpdfstring{$e^{-\theta/2}$}{e^(-theta/2)} 收敛律]\label{thm:xi-time-part9pa-collision-saturation-expansion}
在物理上尾支上，记
\[
u=e^\theta,
\qquad
\alpha(\theta):=P_2'(\theta).
\]
则当 \(\theta\to+\infty\) 时，有显式展开
\[
\boxed{\;
\alpha(\theta)
=
\frac12-\frac14e^{-\theta/2}-e^{-\theta}+O(e^{-3\theta/2}).\;}
\]
等价地，
\[
\boxed{\;
\frac12-\alpha(\theta)
=
\frac14e^{-\theta/2}\Bigl(1+4e^{-\theta/2}+O(e^{-\theta})\Bigr).\;}
\]
\end{theorem}
\begin{proof}
由定理 \ref{thm:xi-time-part9p-positive-branch-global-coordinate}，\(\theta\to+\infty\) 等价于
\[
u=e^\theta\to+\infty
\qquad\Longleftrightarrow\qquad
\rho\to+\infty.
\]
由定理 \ref{thm:xi-time-part9p-collision-frequency-algebraic-ldp} 的参数公式，
\[
u(\rho)=\frac{\rho(\rho^2-2\rho-1)}{\rho-1}
=
\rho^2-\rho-2-\frac{2}{\rho-1},
\]
从而
\[
u(\rho)=\rho^2-\rho-2-\frac{2}{\rho}+O(\rho^{-2}).
\]
同理，
\[
\alpha(\rho)
=
\frac{\rho^3-3\rho^2+\rho+1}{2\rho^3-5\rho^2+4\rho+1}
=
\frac12-\frac{1}{4\rho}-\frac{9}{8\rho^2}+O(\rho^{-3}).
\]

把第一式写成
\[
u(\rho)=\rho^2\Bigl(1-\rho^{-1}-2\rho^{-2}+O(\rho^{-3})\Bigr),
\]
则
\[
u(\rho)^{-1/2}
=
\rho^{-1}\Bigl(1-\rho^{-1}-2\rho^{-2}+O(\rho^{-3})\Bigr)^{-1/2}
=
\rho^{-1}+\frac12\rho^{-2}+O(\rho^{-3}).
\]
反解得
\[
\rho^{-1}=u^{-1/2}-\frac12u^{-1}+O(u^{-3/2}).
\]
代回 \(\alpha(\rho)\) 的展开，可得
\[
\alpha
=
\frac12-\frac14\left(u^{-1/2}-\frac12u^{-1}\right)-\frac98u^{-1}+O(u^{-3/2})
=
\frac12-\frac14u^{-1/2}-u^{-1}+O(u^{-3/2}).
\]
最后以 \(u=e^\theta\) 代回即得
\[
\alpha(\theta)=\frac12-\frac14e^{-\theta/2}-e^{-\theta}+O(e^{-3\theta/2}).
\]
第二个盒装式只是把主项提取为公因子后的重写。证毕。
\end{proof}

\begin{theorem}[碰撞率函数饱和端点的有限边界代价与对数首修正]\label{thm:xi-time-part9pa-collision-boundary-rate-finite}
沿用定理 \ref{thm:xi-time-part9p-collision-frequency-algebraic-ldp} 的记号，并记
\[
\lambda:=\rho(1)=\varphi^2.
\]
当 \(\alpha\uparrow \frac12\) 时，有边界展开
\[
\boxed{\;
I_2(\alpha)
=
\log\lambda
+2\Bigl(\frac12-\alpha\Bigr)\log\!\Bigl(4\Bigl(\frac12-\alpha\Bigr)\Bigr)
+O\!\Bigl(\frac12-\alpha\Bigr).\;}
\]
特别地，
\[
\boxed{\;
\lim_{\alpha\uparrow 1/2}I_2(\alpha)=\log\lambda=2\log\varphi.\;}
\]
因此，碰撞频率逼近极端值 \(\frac12\) 的代价并不发散，而是趋向一个有限边界值；其首个非平凡修正为带对数的线性项。
\end{theorem}
\begin{proof}
令
\[
\varepsilon:=\frac12-\alpha.
\]
由定理 \ref{thm:xi-time-part9pa-collision-saturation-expansion}，
\[
\varepsilon=\frac14u^{-1/2}+O(u^{-1}),
\]
故
\[
u^{-1/2}=4\varepsilon+O(\varepsilon^2),
\qquad
\log u=-2\log(4\varepsilon)+O(\varepsilon).
\]
另一方面，由
\[
u(\rho)=\rho^2-\rho-2+O(\rho^{-1})
\]
可得
\[
\log\rho=\frac12\log u+O(\rho^{-1})
=
-\log(4\varepsilon)+O(\varepsilon),
\]
因为 \(\rho^{-1}=O(u^{-1/2})=O(\varepsilon)\)。

再由定理 \ref{thm:xi-time-part9p-collision-frequency-algebraic-ldp} 的参数公式，
\[
I_2\bigl(\alpha(\rho)\bigr)
=
\alpha(\rho)\log u(\rho)-\log\rho+\log\lambda.
\]
把 \(\alpha=\frac12-\varepsilon\) 代入，得到
\[
I_2(\alpha)
=
\Bigl(\frac12-\varepsilon\Bigr)\log u-\log\rho+\log\lambda.
\]
再代入上面的两条展开，
\[
I_2(\alpha)
=
\Bigl(\frac12-\varepsilon\Bigr)\bigl(-2\log(4\varepsilon)+O(\varepsilon)\bigr)
-\bigl(-\log(4\varepsilon)+O(\varepsilon)\bigr)
+\log\lambda.
\]
其中主项 \(-\log(4\varepsilon)\) 完全抵消，故
\[
I_2(\alpha)
=
\log\lambda+2\varepsilon\log(4\varepsilon)+O(\varepsilon).
\]
将 \(\varepsilon=\frac12-\alpha\) 代回即得所示展开。令 \(\varepsilon\downarrow 0\) 即得
\[
\lim_{\alpha\uparrow 1/2}I_2(\alpha)=\log\lambda=2\log\varphi.
\]
证毕。
\end{proof}

\endinput

\paragraph{时间边界维数、临界退化同一、零知识时间不变性与 dyadic 零点禁逸}
在时间边界的双重重建、Patterson--Sullivan 共形密度、Holonomy 缺陷控制的临界退化、视界零知识的条件独立刻画以及 dyadic 零点稳定与边界零点原理已经建立之后，还可进一步压缩为下列四条跨模块结论。

\begin{theorem}[时间边界上的 Patterson--Sullivan 共形密度维数公式]\label{thm:xi-time-part9q-patterson-sullivan-boundary-dimension}\leanverified{paper\_xi\_time\_part9q\_patterson\_sullivan\_boundary\_dimension}
沿用定理 \ref{thm:xi-time-boundary-double-reconstruction} 的时间边界 \(\partial_T\) 与前缀超度量构造，并沿用定理 \ref{thm:xi-time-boundary-patterson-sullivan-kms} 的 \(\lambda\)-共形边界密度族 \(\{\nu_x^{(\lambda)}\}\)。
对任意 \(a>0\) 定义
\[
d_a(\xi,\eta):=\exp\!\bigl(-a\,\tau(\xi\wedge\eta)\bigr),
\qquad
C_t(\xi):=\{\eta\in\partial_T:\ \tau(\xi\wedge\eta)\ge t\}.
\]
若存在常数 \(0<c\le C<\infty\) 使得对任意 \(x\)、\(\xi\in\partial_T\) 与 \(t\ge 0\) 都有
\[
c\,e^{-\lambda t}\le \nu_x^{(\lambda)}(C_t(\xi))\le C\,e^{-\lambda t},
\]
则 \(\nu_x^{(\lambda)}\) 为 exact-dimensional 测度，并且对 \(\nu_x^{(\lambda)}\)-几乎处处的 \(\xi\in\partial_T\) 有
\[
\dim_{\mathrm{loc}}\bigl(\nu_x^{(\lambda)},\xi\bigr)=\frac{\lambda}{a}.
\]
若进一步 \(\nu_x^{(\lambda)}\) 在 \(\partial_T\) 上满支撑，则
\[
\dim_{\mathrm H}(\partial_T,d_a)=\frac{\lambda}{a}.
\]
\end{theorem}
\begin{proof}
由定理 \ref{thm:xi-time-boundary-double-reconstruction}，\((\partial_T,d_a)\) 是由最大公共前缀深度给出的超度量空间，且半径 \(e^{-at}\) 的闭球恰为 cylinder \(C_t(\xi)\)。
因此若记 \(r_t:=e^{-at}\)，则
\[
B_{d_a}(\xi,r_t)=C_t(\xi).
\]
由假设得到
\[
\log c-\lambda t
\le
\log \nu_x^{(\lambda)}\bigl(B_{d_a}(\xi,r_t)\bigr)
\le
\log C-\lambda t,
\qquad
\log r_t=-at.
\]
于是
\[
\frac{\lambda t-\log C}{at}
\le
\frac{\log \nu_x^{(\lambda)}(B_{d_a}(\xi,r_t))}{\log r_t}
\le
\frac{\lambda t-\log c}{at}.
\]
令 \(t\to\infty\) 即得局部维数极限为 \(\lambda/a\)，从而 \(\nu_x^{(\lambda)}\) 为 exact-dimensional。

若再有满支撑，则同一双边估计可重写为
\[
c\,r^{\lambda/a}\le \nu_x^{(\lambda)}\bigl(B_{d_a}(\xi,r)\bigr)\le C\,r^{\lambda/a}
\]
在 dyadic 半径序列上统一成立；故 \((\partial_T,d_a,\nu_x^{(\lambda)})\) 为 \(\lambda/a\)-Ahlfors 正则，进而
\(
\dim_{\mathrm H}(\partial_T,d_a)=\lambda/a
\)。
证毕。
\end{proof}

\begin{theorem}[临界热力学退化与大偏差平坦维数的 Holonomy 同一]\label{thm:xi-time-part9q-critical-flatness-holonomy-identity}
\leanverified{paper\_xi\_time\_part9q\_critical\_flatness\_holonomy\_identity}
沿用定理 \ref{thm:xi-time-phase-transition-growth-and-defect-rank} 的记号，记 \(\mathrm{ExtEq}_{\kappa}\) 为临界点 \(\lambda=\kappa\) 时平衡态锥的极端射线族；又记 \(\mathrm{Flat}_{\max}(I)\) 为命题 \ref{prop:xi-defect-rank-controls-ldp-flatness} 中大偏差速率函数 \(I\) 的最大仿射平坦子集。
则有严格恒等式
\[
\dim_{\mathrm{aff}}\bigl(\mathrm{ExtEq}_{\kappa}\bigr)
=
\dim\bigl(\mathrm{Flat}_{\max}(I)\bigr)
=
\mathrm{rank}(\mathfrak H).
\]
特别地，
\[
\mathfrak H=\{0\}
\quad\Longleftrightarrow\quad
\text{临界平衡态唯一}
\quad\Longleftrightarrow\quad
I\ \text{严格凸}.
\]
\end{theorem}
\begin{proof}
由定理 \ref{thm:xi-time-phase-transition-growth-and-defect-rank}，临界平衡态极端射线族的仿射维数恰为 \(\mathrm{rank}(\mathfrak H)\)。
另一方面，命题 \ref{prop:xi-defect-rank-controls-ldp-flatness} 给出
\[
\dim\bigl(\mathrm{Flat}_{\max}(I)\bigr)=\mathrm{rank}(\mathfrak H).
\]
两式合并即得所述同一。

若 \(\mathfrak H=\{0\}\)，则上述两种维数同时为 \(0\)，故临界平衡态退化为唯一极端态，且 \(I\) 严格凸。
反向地，若临界平衡态唯一或 \(I\) 严格凸，则对应维数均为 \(0\)，从而 \(\mathrm{rank}(\mathfrak H)=0\)。
再由命题 \ref{prop:xi-time-defect-smith-reduction} 中 \(\mathfrak H\subseteq \ZZ^{d}\) 的 Smith 归约可知 \(\mathfrak H\) 为自由阿贝尔群，故秩为零即 \(\mathfrak H=\{0\}\)。
证毕。
\end{proof}

\begin{theorem}[完美零知识下最小消歧时间的严格不变性]\label{thm:xi-time-part9q-perfect-zk-no-disambiguation-gain}
\leanverified{paper\_xi\_time\_part9q\_perfect\_zk\_no\_disambiguation\_gain}
沿用定义 \ref{def:xi-horizon-perfect-zero-knowledge} 与定理 \ref{thm:xi-horizon-perfect-zk-equivalences} 的记号，令 \(y=\pi(\omega)\)，并假设每条可见纤维 \(\pi^{-1}(y)\) 有限。
记
\[
\mathrm{Ind}(y):=\bigl|\pi^{-1}(y)\bigr|,
\]
并记 \(u_y\) 为纤维 \(\pi^{-1}(y)\) 上的均匀分布。
若协议相对于 \(\pi\) 完美零知识，则把转录 \(\tau\) 作为附加可见输入不会改善任何最优顺序读出协议的最小期望消歧时间；更确切地，对任意字母表 \(\mathcal A\) 有
\[
\inf_{\Pi_{y,\tau}}\mathbb E_\mu[T_{\Pi_{y,\tau}}]
=
\frac{H_\mu(\omega\mid y,\tau)}{\log|\mathcal A|}+O(1)
=
\frac{H_\mu(\omega\mid y)}{\log|\mathcal A|}+O(1).
\]
并且
\[
\inf_{\Pi_{y,\tau}}\mathbb E_\mu[T_{\Pi_{y,\tau}}]
=
\frac{\mathbb E_\mu[\log \mathrm{Ind}(y)]
-\mathbb E_\mu\!\bigl[D_{\mathrm{KL}}(\mu(\cdot\mid y)\,\|\,u_y)\bigr]}{\log|\mathcal A|}
+O(1).
\]
特别地，若 \(\mu(\cdot\mid y)=u_y\) 几乎处处成立，则
\[
\inf_{\Pi_{y,\tau}}\mathbb E_\mu[T_{\Pi_{y,\tau}}]
=
\frac{\mathbb E_\mu[\log \mathrm{Ind}(y)]}{\log|\mathcal A|}+O(1).
\]
\end{theorem}
\begin{proof}
由定理 \ref{thm:xi-horizon-perfect-zk-equivalences}，完美零知识等价于条件互信息消失
\[
I(\omega;\tau\mid y)=0,
\]
也等价于条件独立 \(\omega\perp\tau\mid y\)。
因此对几乎处处的 \((y,\tau)\) 有
\[
\mu(\cdot\mid y,\tau)=\mu(\cdot\mid y),
\qquad
H_\mu(\omega\mid y,\tau)=H_\mu(\omega\mid y).
\]

把定理 \ref{thm:xi-optimal-recovery-time-equals-conditional-entropy} 的条件编码论证逐条件块应用到可见输入 \((y,\tau)\)，即可得
\[
\inf_{\Pi_{y,\tau}}\mathbb E_\mu[T_{\Pi_{y,\tau}}]
=
\frac{H_\mu(\omega\mid y,\tau)}{\log|\mathcal A|}+O(1)
=
\frac{H_\mu(\omega\mid y)}{\log|\mathcal A|}+O(1).
\]
这表明转录 \(\tau\) 在主项上不能带来额外消歧收益。

再由命题 \ref{prop:xi-min-disambiguation-time-entropy-decomposition} 对可见因子 \(y\) 的熵分解，
\[
H_\mu(\omega\mid y)
=
\mathbb E_\mu[\log \mathrm{Ind}(y)]
-\mathbb E_\mu\!\bigl[D_{\mathrm{KL}}(\mu(\cdot\mid y)\,\|\,u_y)\bigr],
\]
代入即得第二个显示公式。
当 \(\mu(\cdot\mid y)=u_y\) 时，相对熵项为零，遂得最后一式。
证毕。
\end{proof}

\begin{theorem}[边界零点原理与 dyadic 逼近给出的零点禁逸与跟踪]\label{thm:xi-time-part9q-dyadic-boundary-zero-no-escape}\leanverified{paper\_xi\_time\_part9q\_dyadic\_boundary\_zero\_no\_escape}
沿用 \S\ref{subsubsec:xi-theta-kernel-dyadic-scale-decomposition} 的 dyadic 分解
\[
\Xi(t)=\sum_{k\ge 0}\Xi^{[k]}(t),
\qquad
F_K(t):=\sum_{k=0}^{K-1}\Xi^{[k]}(t),
\qquad
R_K(t):=\Xi(t)-F_K(t),
\]
并假设在边界切片上存在表示
\[
\Xi(t)=\mathcal D_T(u(t)),
\qquad
\mathcal D_T(u):=\det(I-T(u)),
\]
其中 \(t\mapsto u(t)\) 参数化 \(|u|=1\)，且 \(T(u)\) 满足定理 \ref{thm:xi-contractive-boundary-zero} 的假设。
若 \(\gamma\in\RR\) 为 \(\Xi\) 的简单零点，则存在 \(r>0\)、\(K_0\) 与常数 \(C_\gamma>0\) 使得对一切 \(K\ge K_0\)：
\begin{enumerate}
  \item \(F_K\) 在区间 \((\gamma-r,\gamma+r)\) 内恰有一个实零点 \(\gamma_K\)，并满足 \(\gamma_K\to\gamma\)；
  \item 该零点具有双指数跟踪界
  \[
  |\gamma_K-\gamma|
  \le
  C_\gamma\,(K+1)\,2^{\frac54 K}e^{-\pi 2^K};
  \]
  \item 精确 Fredholm 行列式 \(\mathcal D_T(u)\) 在 \(|u|\neq 1\) 上无零点；因此 dyadic 零点链所追踪到的真实零点只能来自边界切片，而不存在向非临界模长的真实外逸。
\end{enumerate}
\end{theorem}
\begin{proof}
由推论 \ref{cor:xi-theta-kernel-dyadic-zero-chain}，对充分大的 \(K\) 存在唯一零点 \(\gamma_K\) 落在某个固定邻域 \(U\ni\gamma\) 内，并且
\[
|\gamma_K-\gamma|
\le
C_\gamma\,(K+1)\,2^{\frac54 K}e^{-\pi 2^K}.
\]
取 \(r>0\) 使 \((\gamma-r,\gamma+r)\subset U\)，便得到前两条。

另一方面，定理 \ref{thm:xi-contractive-boundary-zero} 给出
\[
\mathcal D_T(u)\neq 0
\qquad(|u|\neq 1).
\]
由于 \(\Xi(t)=\mathcal D_T(u(t))\) 只是该 Fredholm 行列式在边界参数化 \(|u|=1\) 上的限制，故精确对象的任何真零点都只能位于边界切片，不能在 \(|u|\neq 1\) 的区域中出现。
因此，dyadic 逼近所追踪的真实零点不存在向非临界模长的外逸现象。
证毕。
\end{proof}

\endinput

\paragraph{window-\texorpdfstring{$6$}{6} 径向球面设计、moment-kernel 谱根尺度与三能级 escort 闭式}
在定理 \ref{thm:xi-window6-cyclic-kernel-bc3-root-datum} 的 \(BC_3\) 组合根数据、附录 \ref{app:pom-projective-pressure-operator} 的整数 moment-kernel 退化，以及命题 \ref{prop:xi-window6-complete-finite-scale-spectrum-closed-form} 的 window-\(6\) 完整有限尺度谱已经固定之后，还可进一步抽取出下列五条彼此闭合的后果。

\begin{theorem}[window-\texorpdfstring{$6$}{6} 的径向 \texorpdfstring{$B_3$}{B3} 壳层恰为球面 \texorpdfstring{$3$}{3}-设计而非 \texorpdfstring{$4$}{4}-设计]\label{thm:xi-time-part9r-window6-radial-spherical-3design}
\leanverified{paper\_xi\_time\_part9r\_window6\_radial\_spherical\_3design}
沿用定理 \ref{thm:xi-window6-cyclic-kernel-bc3-root-datum} 的 \(B_3\) 根字典，记
\[
R_B:=
\{\pm e_1,\pm e_2,\pm e_3\}
\cup
\{\pm e_i\pm e_j:\ 1\le i<j\le 3\},
\]
\[
\Sigma_6:=
\left\{
\frac{r}{\norm{r}}:\ r\in R_B
\right\}
=
\{\pm e_1,\pm e_2,\pm e_3\}
\cup
\left\{
\frac{\pm e_i\pm e_j}{\sqrt2}:\ 1\le i<j\le 3
\right\}
\subset S^2.
\]
则 \(\Sigma_6\) 是 \(S^2\) 上的精确球面 \(3\)-设计，但不是球面 \(4\)-设计；特别地，它不可能是任何 \(t\ge 4\) 的球面设计。
\end{theorem}
\begin{proof}
记 \(\sigma\) 为 \(S^2\) 上的归一化旋转不变概率测度。先计算二阶矩。六个轴向点给出
\[
\sum_{u\in\{\pm e_1,\pm e_2,\pm e_3\}}uu^{\mathsf T}=2I_3.
\]
对每一对 \(1\le i<j\le 3\)，四个对角方向满足
\[
\sum_{\varepsilon_i,\varepsilon_j\in\{\pm1\}}
\frac{(\varepsilon_i e_i+\varepsilon_j e_j)(\varepsilon_i e_i+\varepsilon_j e_j)^{\mathsf T}}{2}
=
2(e_ie_i^{\mathsf T}+e_je_j^{\mathsf T}),
\]
因为混合项按符号对称相消。对三对 \((i,j)\) 求和得到
\[
\sum_{u\in\Sigma_6}uu^{\mathsf T}=6I_3,
\qquad
\frac1{18}\sum_{u\in\Sigma_6}uu^{\mathsf T}=\frac13 I_3.
\]
这正是 \(S^2\) 上归一化球面测度的二阶矩。

另一方面，\(\Sigma_6\) 为严格对心对称集合：若 \(u\in\Sigma_6\)，则 \(-u\in\Sigma_6\)。因此一切奇次单项式在 \(\Sigma_6\) 上的平均都为零，与 \(\sigma\) 一致。常数项显然一致，二次项由上式给出，故 \(\Sigma_6\) 的 \(0,1,2,3\) 阶矩全部与 \(\sigma\) 相同，从而 \(\Sigma_6\) 为球面 \(3\)-设计。

再检验四阶矩。取坐标函数 \(x_1\)。在 \(\Sigma_6\) 上，\(x_1=\pm 1\) 的点恰有 \(2\) 个，\(x_1=\pm 1/\sqrt2\) 的点恰有 \(8\) 个，其余 \(8\) 个点满足 \(x_1=0\)。因此
\[
\frac1{18}\sum_{u\in\Sigma_6}u_1^4
=
\frac{2\cdot 1^4+8\cdot (1/\sqrt2)^4}{18}
=
\frac{2+8/4}{18}
=
\frac29.
\]
而归一化球面测度满足
\[
\int_{S^2}x_1^4\,\dd\sigma=\frac15.
\]
两者不相等，故 \(\Sigma_6\) 不是球面 \(4\)-设计。于是更高阶 \(t\ge 4\) 的设计也不可能成立。证毕。
\end{proof}

\begin{proposition}[window-\texorpdfstring{$6$}{6} 径向两轨道的最小正次数 Weyl 分离不变量恰为四次幂和]\label{prop:xi-time-part9r-window6-radial-minimal-weyl-separator}
沿用定理 \ref{thm:xi-time-part9r-window6-radial-spherical-3design} 的记号，并设
\[
\Sigma_{\mathrm{ax}}:=\{\pm e_1,\pm e_2,\pm e_3\},
\qquad
\Sigma_{\mathrm{diag}}:=
\left\{
\frac{\pm e_i\pm e_j}{\sqrt2}:\ 1\le i<j\le 3
\right\}.
\]
则：
\begin{enumerate}
  \item \(W(B_3)\) 在 \(\Sigma_6\) 上恰有两条轨道 \(\Sigma_{\mathrm{ax}}\) 与 \(\Sigma_{\mathrm{diag}}\)，其大小分别为 \(6\) 与 \(12\)；
  \item 任意 \(W(B_3)\)-不变齐次二次多项式在 \(\Sigma_6\) 上都为常数，因而不能分离这两条轨道；
  \item 四次不变量
  \[
  Q_4(x):=x_1^4+x_2^4+x_3^4
  \]
  满足
  \[
  Q_4|_{\Sigma_{\mathrm{ax}}}=1,
  \qquad
  Q_4|_{\Sigma_{\mathrm{diag}}}=\frac12.
  \]
\end{enumerate}
因此 \(Q_4\) 是分离 \(\Sigma_6\) 两条径向 Weyl 轨道的最小正次数不变量。
\end{proposition}
\begin{proof}
第一条由定理 \ref{thm:xi-window6-cyclic-kernel-bc3-root-datum} 的 \(6+12\) 轨道分裂立即得到。由于 \(W(B_3)\) 包含全部坐标置换与独立符号翻转，任意 \(W(B_3)\)-不变齐次二次多项式都只能是
\[
c(x_1^2+x_2^2+x_3^2)
\]
的倍数。可是在 \(\Sigma_6\subset S^2\) 上有 \(x_1^2+x_2^2+x_3^2=1\)，故这类二次不变量在 \(\Sigma_6\) 上恒等于常数 \(c\)，无法区分两条轨道。

另一方面，对任一轴向点（例如 \(e_1\)）有
\[
Q_4(e_1)=1,
\]
而对任一对角点（例如 \((e_1+e_2)/\sqrt2\)）有
\[
Q_4\!\left(\frac{e_1+e_2}{\sqrt2}\right)
=
\left(\frac1{\sqrt2}\right)^4+\left(\frac1{\sqrt2}\right)^4
=
\frac12.
\]
因此 \(Q_4\) 的确分离两条轨道。既然正次数 \(2\) 的不变量无法分离，而 \(4\) 次已经可以分离，故 \(Q_4\) 即为最小正次数分离不变量。证毕。
\end{proof}

\begin{theorem}[window-\texorpdfstring{$6$}{6} 的 Hamming 权已完全决定两张根字典中的 Weyl 轨道]\label{thm:xi-time-part9r-window6-hamming-complete-weyl-orbit-invariant}
沿用定理 \ref{thm:xi-window6-cyclic-kernel-bc3-root-datum} 的 \(B_3/C_3\) 根字典，并记
\[
\iota_B:X_6^{\mathrm{cyc}}\to \Phi(B_3),
\qquad
\iota_C:X_6^{\mathrm{cyc}}\to \Phi(C_3),
\]
其中
\[
\Phi(B_3)=\{\pm e_i,\ \pm e_i\pm e_j:\ 1\le i<j\le 3\},
\qquad
\Phi(C_3)=\{\pm 2e_i,\ \pm e_i\pm e_j:\ 1\le i<j\le 3\}.
\]
对 \(w\in X_6^{\mathrm{cyc}}\) 记
\[
s(w):=\mathrm{wt}(w)\in\{0,1,2,3\}.
\]
则 \(X_6^{\mathrm{cyc}}\) 的 Weyl 轨道恰由条件 \(s(w)=1\) 与 \(s(w)\neq 1\) 两层完全决定。更具体地：
\begin{enumerate}
  \item 在 \(B_3\) 字典中，
  \[
  \{w\in X_6^{\mathrm{cyc}}:\ s(w)=1\}
  \]
  恰对应六条短根，而其补集恰对应十二条长根；
  \item 在 \(C_3\) 字典中，同一集合恰对应六条长根，而其补集恰对应十二条短根；
  \item Weyl 轨道投影已由唯一的三次插值多项式显式给出：
  \[
  \mathbf 1_{\{s=1\}}(w)
  =
  \frac12\,s(w)\bigl(s(w)-2\bigr)\bigl(s(w)-3\bigr).
  \]
\end{enumerate}
\end{theorem}
\begin{proof}
定理 \ref{thm:xi-window6-cyclic-kernel-bc3-root-datum} 已给出 \(B_3\) 字典中的两轨道分裂：六个 weight-\(1\) 元恰对应
\[
\{\pm e_1,\pm e_2,\pm e_3\},
\]
其余十二个元素恰对应
\[
\{\pm e_i\pm e_j:\ 1\le i<j\le 3\}.
\]
这正是 \(B_3\) 的短根轨道与长根轨道。另一方面，定理 \ref{thm:xi-window6-c3-hamming-root-length-split} 已给出同一 Hamming 壳在 \(C_3\) 字典下恰对应六条长根
\[
\{\pm 2e_1,\pm 2e_2,\pm 2e_3\},
\]
其补集对应十二条短根
\[
\{\pm e_i\pm e_j:\ 1\le i<j\le 3\}.
\]
于是前两条成立。

最后，\(s(w)\) 只可能取 \(0,1,2,3\) 四值，而我们要求一个次数不超过 \(3\) 的多项式在 \(s=1\) 处取 \(1\)，在 \(s=0,2,3\) 处取 \(0\)。由 Lagrange 插值唯一性，
\[
\mathbf 1_{\{s=1\}}
=
\frac{(s-0)(s-2)(s-3)}{(1-0)(1-2)(1-3)}
=
\frac12\,s(s-2)(s-3).
\]
证毕。
\end{proof}

\begin{theorem}[window-\texorpdfstring{$6$}{6} 的 \texorpdfstring{$B_3/C_3$}{B3/C3} 字典诱导两类不同的根多胞体]\label{thm:xi-time-part9r-window6-b3c3-root-polytopes}
沿用定理 \ref{thm:xi-time-part9r-window6-hamming-complete-weyl-orbit-invariant} 的字典 \(\iota_B,\iota_C\)，定义
\[
P_B:=\operatorname{conv}\bigl(\iota_B(X_6^{\mathrm{cyc}})\bigr),
\qquad
P_C:=\operatorname{conv}\bigl(\iota_C(X_6^{\mathrm{cyc}})\bigr).
\]
则成立：
\begin{enumerate}
  \item
  \[
  P_B
  =
  \bigl\{x\in\RR^3:\ |x_i|\le 1,\ |x_1|+|x_2|+|x_3|\le 2\bigr\},
  \]
  \[
  \operatorname{Vol}(P_B)=\frac{20}{3}.
  \]
  因而 \(P_B\) 即标准 cuboctahedron；其 \(12\) 个顶点恰为 \(B_3\) 的长根，而六个短根 \(\{\pm e_1,\pm e_2,\pm e_3\}\) 恰位于六个方形面的中心。
  \item
  \[
  P_C
  =
  \bigl\{x\in\RR^3:\ |x_1|+|x_2|+|x_3|\le 2\bigr\},
  \]
  \[
  \operatorname{Vol}(P_C)=\frac{32}{3}.
  \]
  因而 \(P_C\) 即尺度 \(2\) 的正交叉多胞体；其 \(6\) 个顶点恰为 \(C_3\) 的长根 \(\{\pm 2e_1,\pm 2e_2,\pm 2e_3\}\)，而 \(12\) 个短根 \(\{\pm e_i\pm e_j\}\) 恰为这些顶点之间边的中点。
\end{enumerate}
\end{theorem}
\begin{proof}
由定理 \ref{thm:xi-time-part9r-window6-hamming-complete-weyl-orbit-invariant}，\(B_3\) 字典下的点集恰为
\[
\{\pm e_1,\pm e_2,\pm e_3\}\cup\{\pm e_i\pm e_j:\ 1\le i<j\le 3\}.
\]
记
\[
K_B:=\bigl\{x\in\RR^3:\ |x_i|\le 1,\ |x_1|+|x_2|+|x_3|\le 2\bigr\}.
\]
在第一卦限中，\(K_B\) 化为
\[
\bigl\{x_i\ge 0,\ x_i\le 1,\ x_1+x_2+x_3\le 2\bigr\},
\]
它正是单位立方体 \([0,1]^3\) 去掉顶角四面体
\[
\bigl\{x_i\ge 0,\ x_1+x_2+x_3\ge 2\bigr\}
\]
之后所得区域。因此 \(K_B\) 的极点恰为
\[
(1,1,0),\ (1,0,1),\ (0,1,1)
\]
及其所有符号与坐标置换，也就是 \(B_3\) 的 \(12\) 个长根。六个短根 \(\pm e_i\) 满足 \(|x_i|=1\) 且其余坐标为 \(0\)，故恰位于相应方形面的中心而不是新的顶点。于是 \(P_B=K_B\)。又第一卦限中的体积为
\[
1-\frac16=\frac56,
\]
乘以八个卦限即得
\[
\operatorname{Vol}(P_B)=8\cdot \frac56=\frac{20}{3}.
\]

对 \(C_3\) 字典，由同一定理，点集恰为
\[
\{\pm 2e_1,\pm 2e_2,\pm 2e_3\}\cup\{\pm e_i\pm e_j:\ 1\le i<j\le 3\}.
\]
记
\[
K_C:=\bigl\{x\in\RR^3:\ |x_1|+|x_2|+|x_3|\le 2\bigr\}.
\]
该集合的顶点正是六个向量 \(\pm 2e_i\)，故 \(K_C\) 为尺度 \(2\) 的正交叉多胞体。另一方面，每个点 \(\pm e_i\pm e_j\) 都恰是对应符号下两顶点 \(\pm 2e_i\) 与 \(\pm 2e_j\) 的中点，因此全部 \(12\) 个短根都落在 \(K_C\) 的边上而不产生新顶点。故 \(P_C=K_C\)。

在第一卦限中，\(K_C\) 退化为四面体
\[
\bigl\{x_i\ge 0,\ x_1+x_2+x_3\le 2\bigr\},
\]
其体积为
\[
\frac{2^3}{3!}=\frac43.
\]
乘以八个卦限即得
\[
\operatorname{Vol}(P_C)=8\cdot \frac43=\frac{32}{3}.
\]
证毕。
\end{proof}

\begin{proposition}[window-\texorpdfstring{$6$}{6} 的 \texorpdfstring{$C_3$}{C3} 字典由 \texorpdfstring{$B_3$}{B3} 字典经 Hamming 分层径向伸缩唯一确定]\label{prop:xi-time-part9r-window6-b3-to-c3-weight-radial-dilation}
沿用定理 \ref{thm:xi-time-part9r-window6-hamming-complete-weyl-orbit-invariant} 的字典 \(\iota_B,\iota_C\)。则存在唯一函数
\[
\lambda:\{0,1,2,3\}\to \RR_{>0}
\]
使得对每个 \(w\in X_6^{\mathrm{cyc}}\) 都有
\[
\iota_C(w)=\lambda\bigl(\mathrm{wt}(w)\bigr)\,\iota_B(w).
\]
并且该函数恰为
\[
\lambda(s)=1+\frac12\,s(s-2)(s-3)
\qquad
\bigl(s\in\{0,1,2,3\}\bigr).
\]
特别地，
\[
\lambda(1)=2,
\qquad
\lambda(0)=\lambda(2)=\lambda(3)=1.
\]
\end{proposition}
\begin{proof}
由定理 \ref{thm:xi-time-part9r-window6-hamming-complete-weyl-orbit-invariant}，若 \(\mathrm{wt}(w)=1\)，则
\[
\iota_B(w)\in\{\pm e_1,\pm e_2,\pm e_3\},
\qquad
\iota_C(w)\in\{\pm 2e_1,\pm 2e_2,\pm 2e_3\},
\]
故 \(\iota_C(w)=2\,\iota_B(w)\)。若 \(\mathrm{wt}(w)\neq 1\)，则两张字典都把 \(w\) 送到同一向量 \(\pm e_i\pm e_j\)，从而 \(\iota_C(w)=\iota_B(w)\)。因此任何满足题设的函数都必须取值
\[
\lambda(1)=2,
\qquad
\lambda(0)=\lambda(2)=\lambda(3)=1,
\]
故唯一性成立。

再由前述四个插值条件，定义在 \(\{0,1,2,3\}\) 上的唯一次数不超过 \(3\) 的插值多项式为
\[
\lambda(s)=1+\frac12\,s(s-2)(s-3).
\]
它恰在 \(s=1\) 处取值 \(2\)，在 \(s=0,2,3\) 处取值 \(1\)，故所示公式成立。证毕。
\end{proof}

\begin{theorem}[附录整数退化 moment-kernel 的谱半径满足次乘法律]\label{thm:xi-time-part9r-moment-kernel-radius-submultiplicative}
\leanverified{paper\_xi\_time\_part9r\_moment\_kernel\_radius\_submultiplicative}
沿用定义 \ref{def:pom-projective-transfer-operator}、定义 \ref{def:pom-homogeneous-polynomials-simplex} 与定理 \ref{thm:pom-projective-operator-degenerates-to-moment-kernel} 的记号。对每个整数 \(q\ge 2\)，记
\[
r_q:=\rho(\mathcal L_q|_{\mathcal H_q})=\rho(\widetilde A_q).
\]
则对任意整数 \(a,b\ge 2\) 都有
\[
r_{a+b}\le r_a\,r_b.
\]
等价地，序列 \(\bigl(\log r_q\bigr)_{q\ge 2}\) 为次可加序列。
\end{theorem}
\begin{proof}
定理 \ref{thm:xi-projective-moment-radius-submultiplicative} 已经证明
\[
\rho(\mathcal L_{a+b}|_{\mathcal H_{a+b}})
\le
\rho(\mathcal L_a|_{\mathcal H_a})\,
\rho(\mathcal L_b|_{\mathcal H_b})
\qquad(a,b\ge 2).
\]
再由定理 \ref{thm:pom-projective-operator-degenerates-to-moment-kernel} 的整数退化识别
\[
\rho(\mathcal L_q|_{\mathcal H_q})=\rho(\widetilde A_q)=r_q,
\]
代入即可得到所述不等式。对两边取对数即得次可加性。证毕。
\end{proof}

\begin{corollary}[moment-kernel 的谱根尺度极限存在]\label{cor:xi-time-part9r-moment-kernel-root-scale}
沿用定理 \ref{thm:xi-time-part9r-moment-kernel-radius-submultiplicative} 的记号。定义
\[
\rho_{\infty}^{\mathrm{mk}}
:=
\lim_{q\to\infty} r_q^{1/q}.
\]
则该极限存在，并满足
\[
\rho_{\infty}^{\mathrm{mk}}
=
\inf_{q\ge 2} r_q^{1/q}.
\]
\end{corollary}
\begin{proof}
由推论 \ref{cor:xi-projective-moment-fekete-slope}，
\[
\tau_\ast:=\lim_{q\to\infty}\frac{1}{q}\log r_q
=
\inf_{q\ge 2}\frac{1}{q}\log r_q.
\]
对上式两边取指数即可得到结论。证毕。
\end{proof}

\begin{theorem}[window-\texorpdfstring{$6$}{6} 的 escort 族构成精确可解的三能级模型]\label{thm:xi-time-part9r-window6-escort-threelevel-exact}
沿用 \(\pi_{m,q}(w)\propto (d_m^{\mathrm{bin}}(w))^q\) 的 escort 记号。对 \(m=6\) 记
\[
Z_6(q):=S_q(6),
\qquad
\pi_{6,q}(d=j):=
\sum_{\substack{w\in X_6\\ d_6^{\mathrm{bin}}(w)=j}}\pi_{6,q}(w)
\qquad (j\in\{2,3,4\}).
\]
则对一切 \(q\ge 0\) 成立：
\begin{enumerate}
  \item 配分函数具有闭式
  \[
  Z_6(q)=8\cdot 2^q+4\cdot 3^q+9\cdot 4^q;
  \]
  \item 三个纤维层级的 escort 质量分别为
  \[
  \pi_{6,q}(d=2)=\frac{8\cdot 2^q}{Z_6(q)},
  \qquad
  \pi_{6,q}(d=3)=\frac{4\cdot 3^q}{Z_6(q)},
  \qquad
  \pi_{6,q}(d=4)=\frac{9\cdot 4^q}{Z_6(q)};
  \]
  \item 存在唯一交叉参数 \(q_6^\dagger\in(0,1)\) 使
  \[
  \pi_{6,q_6^\dagger}(d=4)=\frac12,
  \]
  等价地，
  \[
  9\cdot 4^{q_6^\dagger}=8\cdot 2^{q_6^\dagger}+4\cdot 3^{q_6^\dagger};
  \]
  \item 当 \(q\to+\infty\) 时，
  \[
  1-\pi_{6,q}(d=4)
  =
  \frac49\left(\frac34\right)^q
  +O\!\left(\left(\frac{9}{16}\right)^q\right).
  \]
\end{enumerate}
因此 window-\(6\) 的 escort 族在有限尺度上恰形成一个精确可解的三能级模型。
\end{theorem}
\begin{proof}
由命题 \ref{prop:xi-window6-complete-finite-scale-spectrum-closed-form}，window-\(6\) 的纤维直方图为
\[
2:8,\qquad 3:4,\qquad 4:9,
\]
且
\[
S_q(6)=8\cdot 2^q+4\cdot 3^q+9\cdot 4^q.
\]
这立即给出 \(Z_6(q)\) 的闭式以及三条 escort 质量公式。

对交叉参数，令
\[
G(q):=9\cdot 2^q-4\cdot \left(\frac32\right)^q-8.
\]
则 \(G(q)=0\) 与
\(
9\cdot 4^q=8\cdot 2^q+4\cdot 3^q
\)
等价。又有
\[
G(0)=-3,
\qquad
G(1)=4.
\]
并且
\[
G'(q)
=
9\log 2\cdot 2^q
-4\log\!\left(\frac32\right)\left(\frac32\right)^q
=
4\log\!\left(\frac32\right)\left(\frac32\right)^q
\left(
\frac{9\log 2}{4\log(3/2)}\left(\frac43\right)^q-1
\right).
\]
由于
\[
\frac{9\log 2}{4\log(3/2)}>1,
\]
且 \((4/3)^q\ge 1\) 对 \(q\ge 0\) 恒成立，故 \(G'(q)>0\) 在 \([0,\infty)\) 上处处成立。于是 \(G\) 严格递增，从而在 \((0,1)\) 内恰有唯一零点 \(q_6^\dagger\)。

其数值审计为
% AUTO-GENERATED by scripts/exp_window6_escort_threelevel_crossover_audit.py
\begin{equation}\label{eq:window6_escort_threelevel_crossover_audit}
\begin{aligned}
q_6^\dagger &\approx 0.524783184425175131,\\
9\cdot 4^{q_6^\dagger}-8\cdot 2^{q_6^\dagger}-4\cdot 3^{q_6^\dagger} &\approx 1.7763568394002505e-15,\\
Z_6(q_6^\dagger) &\approx 37.258336829036636573,\\
\pi_{6,q_6^\dagger}(d=4) &\approx 0.500000000000000000.
\end{aligned}
\end{equation}


最后，对冻结尾项，写
\[
a_q:=\left(\frac34\right)^q,
\qquad
b_q:=\left(\frac12\right)^q.
\]
则
\[
1-\pi_{6,q}(d=4)
=
\frac{8\cdot 2^q+4\cdot 3^q}{8\cdot 2^q+4\cdot 3^q+9\cdot 4^q}
=
\frac{4a_q+8b_q}{9+4a_q+8b_q}.
\]
注意 \(b_q=O(a_q^2)\)，因为
\[
\frac{b_q}{a_q^2}
=
\left(\frac{8}{9}\right)^q
\longrightarrow 0.
\]
因此
\[
\frac{4a_q+8b_q}{9+4a_q+8b_q}
=
\frac{4a_q+O(a_q^2)}{9+O(a_q)}
=
\frac49 a_q+O(a_q^2)
=
\frac49\left(\frac34\right)^q
+O\!\left(\left(\frac{9}{16}\right)^q\right).
\]
证毕。
\end{proof}

\noindent
由此，window-\(6\) 的径向根几何、附录 \(E\) 的整数 moment-kernel 谱族与有限窗 escort 冻结在同一条终端链上闭合：径向方向集在三阶球面矩处精确封闭而在四阶首次分叉，Weyl 轨道投影已由 Hamming 权的三次多项式显式给出，\(B_3/C_3\) 两张字典的凸包分别固定为 cuboctahedron 与尺度 \(2\) 的 octahedron，且二者之间的转换不是全局线性统一而是由重量分层径向伸缩唯一实现；另一方面，整数退化谱根满足 Fekete 型次乘规律并强制出单一谱根尺度，有限窗 escort 则坍缩为一个精确可解的三能级模型。于是，离散根壳、根多胞体、谱根尺度与冻结支撑函数可被置于同一条可显式审计的径向变分主线。

\paragraph{二阶有效支撑、共振整函数临界饱和与 escort 温度流的差分闭式}
在二阶碰撞底噪、共振整函数的 Cartwright 精确型、escort 温度族的末位 Bernoulli 塌缩以及 fold-gauge 常数的 Bernoulli 模态消去均已建立之后，还可进一步抽取出下列六条彼此衔接的后果。它们分别刻画 bin-fold 二阶有效支撑的绝对比例坍缩、黄金共振整函数在零点密度与虚轴增长上的双重临界饱和、整数采样越过平方可和阈值、escort 温度流的极限 logistic 指数几何、单样本最优判别的末位完备性，以及 Stirling--Bernoulli 层级的移位算子加速表述。

\begin{theorem}[bin-fold 二阶有效支撑的绝对比例坍缩]\label{thm:xi-time-part9s-renyi2-effective-support-collapse}
\leanverified{paper\_xi\_time\_part9s\_renyi2\_effective\_support\_collapse}
沿用
\[
\mu_m(w):=\frac{d_m^{\mathrm{bin}}(w)}{2^m}
\qquad (w\in X_m)
\]
的记号，并定义
\[
H_2(\mu_m):=-\log\sum_{w\in X_m}\mu_m(w)^2
=-\log \mathsf{Col}_m,
\qquad
\mathcal N_2(\mu_m):=e^{H_2(\mu_m)}=\frac{1}{\mathsf{Col}_m}.
\]
则有
\[
\log |X_m|-H_2(\mu_m)
=
\log\!\bigl(F_{m+2}\mathsf{Col}_m\bigr)
\longrightarrow +\infty.
\]
等价地，
\[
\frac{\mathcal N_2(\mu_m)}{|X_m|}
=
\frac{1}{F_{m+2}\mathsf{Col}_m}
\longrightarrow 0.
\]
另一方面，
\[
\frac{1}{m}H_2(\mu_m)\longrightarrow \log\varphi.
\]
因此，\(\mu_m\) 的二阶有效支撑虽然在指数率上仍与 \(|X_m|\) 同阶，却在绝对比例上坍缩到零。
\end{theorem}
\begin{proof}
由定理 \ref{thm:xi-time-part9p-renyi2-obstruction}，
\[
\log |X_m|-H_2(\mu_m)
=
\log\!\bigl(F_{m+2}\mathsf{Col}_m\bigr).
\]
而推论 \ref{cor:fold-collision-scale-diverges} 给出
\[
F_{m+2}\mathsf{Col}_m\longrightarrow +\infty,
\]
故第一条极限成立。将上式指数化并重排，即得
\[
\frac{\mathcal N_2(\mu_m)}{|X_m|}
=
\frac{1}{F_{m+2}\mathsf{Col}_m}
\longrightarrow 0.
\]

另一方面，定理 \ref{thm:xi-foldbin-information-spectrum-single-point-collapse} 说明对一切 \(q>0\) 都有
\[
\frac{1}{m}H_q(\mu_m)\longrightarrow \log\varphi.
\]
取 \(q=2\) 即得
\[
\frac{1}{m}H_2(\mu_m)\longrightarrow \log\varphi.
\]
故
\(
\mathcal N_2(\mu_m)=e^{H_2(\mu_m)}
\)
仍具有 \(\varphi^m\) 级的指数尺度，而其相对比例却趋于 \(0\)。证毕。
\end{proof}

\begin{theorem}[黄金共振整函数的零点密度与虚轴增长同时达到临界型]\label{thm:xi-time-part9s-resonance-entire-imaginary-axis-saturation}
沿用定理 \ref{thm:fold-resonance-entire-lp} 的记号，定义双侧零点计数函数
\[
N_{\mathcal F_\varphi}(R):=
\#\bigl(\mathcal Z(\mathcal F_\varphi)\cap[-R,R]\bigr).
\]
则有
\[
N_{\mathcal F_\varphi}(R)=2R+O(\log R)
\qquad (R\to\infty),
\]
并且
\[
\log|\mathcal F_\varphi(\mathrm{i}y)|
=
\pi|y|+O(\log(2+|y|))
\qquad (|y|\to\infty).
\]
因而 \(\mathcal F_\varphi\) 的指数型恰为 \(\pi\)，且定理
\ref{thm:fold-resonance-entire-lp} 中的增长上界
\(
|\mathcal F_\varphi(z)|\le e^{\pi|\Im z|}
\)
在虚轴上被渐近饱和。
\end{theorem}
\begin{proof}
第一条正是定理 \ref{thm:xi-fold-resonance-entire-zero-counting-linear} 的双侧计数公式。

再证第二条。由于 \(\mathcal F_\varphi\) 为偶函数，只需考虑 \(y\ge 2\)。由乘积表达式，
\[
\mathcal F_\varphi(\mathrm{i}y)
=
\prod_{n=2}^{\infty}\cosh(\pi y\varphi^{-n}).
\]
记
\[
N:=\lfloor \log_\varphi y\rfloor.
\]
则当 \(2\le n\le N\) 时，\(\pi y\varphi^{-n}\ge \pi/\varphi\)，故
\[
\log\cosh(\pi y\varphi^{-n})
=
\pi y\varphi^{-n}-\log 2+O(e^{-c y\varphi^{-n}})
\]
对某个绝对常数 \(c>0\) 成立。于是
\[
\sum_{n=2}^{N}\log\cosh(\pi y\varphi^{-n})
=
\pi y\sum_{n=2}^{N}\varphi^{-n}+O(N).
\]
由于
\[
\sum_{n=2}^{\infty}\varphi^{-n}=1
\]
且 \(y\varphi^{-N}=O(1)\)，便得到
\[
\sum_{n=2}^{N}\log\cosh(\pi y\varphi^{-n})
=
\pi y+O(\log y).
\]

对尾和 \(n>N\)，有 \(\pi y\varphi^{-n}\le \pi\varphi^{-1}\)，从而
\[
\log\cosh(\pi y\varphi^{-n})=O(y^2\varphi^{-2n}).
\]
于是
\[
\sum_{n>N}\log\cosh(\pi y\varphi^{-n})
=
O\!\left(y^2\sum_{n>N}\varphi^{-2n}\right)
=
O(1).
\]
两段合并即得
\[
\log|\mathcal F_\varphi(\mathrm{i}y)|
=
\pi y+O(\log y)
\qquad (y\to+\infty).
\]
由偶性可推广为
\[
\log|\mathcal F_\varphi(\mathrm{i}y)|
=
\pi|y|+O(\log(2+|y|)).
\]
最后结合定理 \ref{thm:fold-resonance-entire-lp} 的上界
\(
|\mathcal F_\varphi(z)|\le e^{\pi|\Im z|}
\)
即可看出其指数型恰为 \(\pi\)，且该上界在虚轴上达到主阶等号。证毕。
\end{proof}

\begin{corollary}[黄金共振整函数的整数采样越过平方可和阈值]\label{cor:xi-time-part9s-resonance-entire-integer-sampling-not-l2}
沿用定理 \ref{thm:xi-time-part9s-resonance-entire-imaginary-axis-saturation} 的记号，则整数采样列
\[
\bigl(\mathcal F_\varphi(n)\bigr)_{n\in\ZZ}
\]
不属于 \(\ell^2(\ZZ)\)。特别地，尽管 \(\mathcal F_\varphi\) 属于 Laguerre--P\'olya 类且指数型恰为 \(\pi\)，其整数轨道采样仍不落入平方可和采样制度。
\end{corollary}
\begin{proof}
由定理 \ref{thm:fold-pisot-bernoulli-convolution-representation}，
\[
\mathcal F_\varphi(u)
=
\prod_{n=2}^{\infty}\cos(\pi u\varphi^{-n})
=
\widehat{\mu_q^{(2)}}(\pi u)
\qquad (u\in\ZZ),
\]
故
\[
|\mathcal F_\varphi(u)|=C_\varphi(u).
\]
另一方面，由定理 \ref{thm:fold-sigma-phi-diverges}，
\[
1+2\sum_{u=1}^{\infty}C_\varphi(u)^2=+\infty.
\]
于是
\[
\sum_{n\in\ZZ}|\mathcal F_\varphi(n)|^2
=
|\mathcal F_\varphi(0)|^2+2\sum_{u=1}^{\infty}|\mathcal F_\varphi(u)|^2
=
1+2\sum_{u=1}^{\infty}C_\varphi(u)^2
=
\infty.
\]
故 \((\mathcal F_\varphi(n))_{n\in\ZZ}\notin\ell^2(\ZZ)\)。证毕。
\end{proof}

\begin{theorem}[escort 极限族的 logistic 指数几何]\label{thm:xi-time-part9s-escort-logistic-expfamily-bregman}
\leanverified{paper\_xi\_time\_part9s\_escort\_logistic\_expfamily\_bregman}
沿用 \(\pi_{m,q}\) 的记号，并定义
\[
\theta(q):=\frac{1}{1+\varphi^{q+1}},
\qquad
\eta(q):=-(q+1)\log\varphi,
\qquad
A(\eta):=\log(1+e^\eta).
\]
则
\[
\theta(q)=A'(\eta(q)).
\]
对任意 \(p,q\ge 0\)，有
\[
\lim_{m\to\infty}D_{\mathrm{KL}}(\pi_{m,q}\Vert \pi_{m,p})
=
A(\eta(p))-A(\eta(q))-A'(\eta(q))\bigl(\eta(p)-\eta(q)\bigr).
\]
因此，bin-fold escort 温度族在极限上恰形成一条一维 logistic 指数族。其 Fisher 信息满足
\[
\lim_{h\to 0}\lim_{m\to\infty}
\frac{2}{h^2}D_{\mathrm{KL}}(\pi_{m,q+h}\Vert \pi_{m,q})
=
(\log\varphi)^2\,\theta(q)\bigl(1-\theta(q)\bigr).
\]
\end{theorem}
\begin{proof}
由定义，
\[
\theta(q)
=
\frac{1}{1+\varphi^{q+1}}
=
\frac{e^{\eta(q)}}{1+e^{\eta(q)}}
=
A'(\eta(q)).
\]

另一方面，由结论 \ref{con:xi-time-part9o-escort-fdiv-collapse} 在
\(
f(u)=u\log u
\)
处的特化，
\[
\lim_{m\to\infty}D_{\mathrm{KL}}(\pi_{m,q}\Vert \pi_{m,p})
=
\theta(q)\log\frac{\theta(q)}{\theta(p)}
+
\bigl(1-\theta(q)\bigr)
\log\frac{1-\theta(q)}{1-\theta(p)}.
\]
而这是 Bernoulli 指数族的标准 Bregman 形式，恰可改写为
\[
A(\eta(p))-A(\eta(q))-A'(\eta(q))\bigl(\eta(p)-\eta(q)\bigr).
\]
故极限族确为 logistic 指数族。

最后，由定理 \ref{thm:xi-time-part9ob-escort-two-sided-local-kl}，
\[
\lim_{m\to\infty}D_{\mathrm{KL}}(\pi_{m,q+h}\Vert \pi_{m,q})
=
\frac12(\log\varphi)^2\,\theta(q)\bigl(1-\theta(q)\bigr)h^2+O(h^3).
\]
两边乘以 \(2/h^2\) 并令 \(h\to 0\)，即得 Fisher 信息闭式。证毕。
\end{proof}

\begin{theorem}[escort 单样本最优判别的末位完备性]\label{thm:xi-time-part9s-escort-bayes-lastbit-completeness}
对任意 \(p,q\ge 0\) 且 \(p\neq q\)，记等先验下判别 \(\pi_{m,p}\) 与 \(\pi_{m,q}\) 的最优单样本 Bayes 错误概率为
\[
P_{\mathrm e}^\ast(m;p,q).
\]
则有极限
\[
P_{\mathrm e}^\ast(m;p,q)
\longrightarrow
\frac12\Bigl(1-\bigl|\theta(q)-\theta(p)\bigr|\Bigr),
\qquad
\theta(r)=\frac{1}{1+\varphi^{r+1}}.
\]
并且任意渐近最优检验都只依赖于末位统计量
\[
\ell_m(w):=w_m.
\]
等价地，在极限统计实验
\[
\{\mathrm{Bernoulli}(\theta(r)):r\ge 0\}
\]
中，末位读出已成为完备且充分的一维坐标。
\end{theorem}
\begin{proof}
由定理 \ref{cor:xi-time-part9ob-escort-np-lastbit-collapse}，
\[
\|\pi_{m,q}-\pi_{m,p}\|_{\mathrm{TV}}
\longrightarrow
\bigl|\theta(q)-\theta(p)\bigr|,
\]
并且任何 Neyman--Pearson 最强检验在极限中都退化为仅依赖 \(\ell_m(w)=w_m\) 的二块判别。由等先验下的 Helstrom 公式，
\[
P_{\mathrm e}^\ast(m;p,q)
=
\frac12\bigl(1-\|\pi_{m,q}-\pi_{m,p}\|_{\mathrm{TV}}\bigr),
\]
立即得到所示 Bayes 错误极限。

再证完备性。设 \(g:\{0,1\}\to\RR\) 满足
\[
\mathbb E_{\theta(r)}[g(B)]=0
\qquad
\text{对所有 }r\ge 0,
\]
其中 \(B\sim\mathrm{Bernoulli}(\theta(r))\)。则
\[
g(0)+\theta(r)\bigl(g(1)-g(0)\bigr)=0
\]
对全部 \(r\ge 0\) 成立。由于 \(\theta(r)\) 取遍区间 \((0,\varphi^{-2}]\) 中的无穷多点，上式作为 \(\theta\) 的仿射函数只能恒为零，故
\[
g(0)=g(1)=0.
\]
因此极限二元实验中的身份统计量是完备的，而其参数依赖显然也完全由该二元坐标承载，故同时充分。证毕。
\end{proof}

\begin{theorem}[fold-gauge 任意阶加速的移位算子表述]\label{thm:xi-time-part9s-gauge-shift-operator-accelerator}
沿用推论 \ref{cor:fold-bin-recover-2pi} 的规范常数列 \((C_m)_{m\ge 1}\)，定义
\[
f_m:=\log C_m,
\qquad
q_r:=\Bigl(\frac{\varphi}{2}\Bigr)^{2r-1}
\qquad (r\ge 1),
\]
并令移位算子 \(E\) 作用于序列 \(f=(f_m)_{m\ge 1}\) 为
\[
(Ef)_m:=f_{m+1}.
\]
对任意 \(R\ge 1\)，定义
\[
\mathcal A_R:=
\prod_{r=1}^{R}\frac{E-q_r}{1-q_r}.
\]
则有
\[
(\mathcal A_R f)_m
=
\log(2\pi)
+O\!\left(\Bigl(\frac{\varphi}{2}\Bigr)^{(2R+1)m}\right)
\qquad (m\to\infty).
\]
相应地，若定义
\[
C_m^{[R]}:=\exp\bigl((\mathcal A_R f)_m\bigr),
\]
则
\[
C_m^{[R]}
=
2\pi\left(
1+O\!\left(\Bigl(\frac{\varphi}{2}\Bigr)^{(2R+1)m}\right)
\right).
\]
因此，fold-gauge 常数列不仅内生恢复 \(2\pi\)，而且可由有限差分算子族实现任意阶指数加速。
\end{theorem}
\begin{proof}
由定理 \ref{thm:fold-bin-gauge-constant-stirling-bernoulli-hierarchy}，对任意固定 \(R\ge 1\) 有展开
\[
f_m
=
\log(2\pi)+\sum_{s=1}^{R+1}a_s q_s^m+O(q_{R+2}^m),
\]
其中
\[
a_s:=
\frac{B_{2s}}{s(2s-1)}
\bigl(\varphi^{-1}+\varphi^{2s-3}\bigr).
\]
对任意几何模态 \(m\mapsto q_s^m\)，有
\[
\bigl((E-q_r)q_s^{\,\bullet}\bigr)_m
=
q_s^{m+1}-q_r q_s^m
=
(q_s-q_r)q_s^m.
\]
故当 \(s=r\) 时，
\[
\bigl((E-q_r)q_r^{\,\bullet}\bigr)_m=0,
\]
即第 \(r\) 个模态被精确消去；而对常数序列 \(1\) 则有
\[
\frac{E-q_r}{1-q_r}\,1=1.
\]
因此应用全部 \(R\) 个因子后，前 \(R\) 个 Bernoulli 模态全部消失，得到
\[
(\mathcal A_R f)_m
=
\log(2\pi)
+a_{R+1}
\prod_{r=1}^{R}\frac{q_{R+1}-q_r}{1-q_r}\,q_{R+1}^m
+O(q_{R+2}^m).
\]
由于
\[
q_{R+1}=\Bigl(\frac{\varphi}{2}\Bigr)^{2R+1},
\qquad
q_{R+2}=q_{R+1}\Bigl(\frac{\varphi}{2}\Bigr)^2,
\]
便得到
\[
(\mathcal A_R f)_m
=
\log(2\pi)
+O\!\left(\Bigl(\frac{\varphi}{2}\Bigr)^{(2R+1)m}\right).
\]
最后对该式指数化即可得到 \(C_m^{[R]}\) 的逼近公式。证毕。
\end{proof}

\endinput

\paragraph{escort 温度半轴的平坦化、圆弧嵌入与固定样本二项塌缩}
在 escort 极限族的一维 logistic 指数几何、Fisher--Rao 有限弧长、Hellinger--Chernoff 母核以及乘积试验的单比特塌缩已经确立之后，该温度半轴还可继续被压缩为一条带唯一端点完备化的平坦圆弧模型。下列结果分别给出显式平坦坐标、零温端点的度量完备化、平方根嵌入的单位圆弧表述、Chernoff 极小点的闭式反演、固定样本下的二项实验塌缩，以及顺序对数似然步长的严格二点格结构。

\begin{theorem}[bin-fold escort 温度半轴的 Fisher--Rao 度量存在显式平坦坐标]\label{thm:xi-time-part9sb-escort-fisher-flat-coordinate}
定义
\[
\theta(q):=\frac{1}{1+\varphi^{q+1}},
\qquad
\mathcal I(q):=(\log\varphi)^2\,\theta(q)\bigl(1-\theta(q)\bigr),
\qquad q\ge 0,
\]
以及坐标函数
\[
\psi(q):=2\arctan\!\bigl(\varphi^{(q+1)/2}\bigr)
=
\pi-2\arctan\!\bigl(\varphi^{-(q+1)/2}\bigr).
\]
则
\[
\psi'(q)=\sqrt{\mathcal I(q)}.
\]
从而对任意 \(p,q\ge 0\)，escort 温度半轴上的 Fisher--Rao 距离满足闭式
\[
d_{\mathrm{FR}}(p,q)
=
\left|\psi(q)-\psi(p)\right|
=
\left|
2\arctan\!\bigl(\varphi^{(q+1)/2}\bigr)
-
2\arctan\!\bigl(\varphi^{(p+1)/2}\bigr)
\right|.
\]
特别地，\(\psi\) 把整个温度半轴 \([0,\infty)\) 等距拉直到 Euclidean 线段
\[
\psi([0,\infty))
=
\bigl[2\arctan(\sqrt\varphi),\,\pi\bigr).
\]
\end{theorem}
\begin{proof}
由结论 \ref{con:xi-time-part9o-escort-fisher-rao-closed}，极限 Fisher 密度正是
\[
\mathcal I(q)=(\log\varphi)^2\,\theta(q)\bigl(1-\theta(q)\bigr).
\]
另一方面，
\[
\psi'(q)
=
\frac{\log\varphi\,\varphi^{(q+1)/2}}{1+\varphi^{q+1}}
=
(\log\varphi)\sqrt{\frac{\varphi^{q+1}}{(1+\varphi^{q+1})^2}}
=
(\log\varphi)\sqrt{\theta(q)\bigl(1-\theta(q)\bigr)}
=
\sqrt{\mathcal I(q)}.
\]
故
\[
d_{\mathrm{FR}}(p,q)
=
\left|
\int_p^q \sqrt{\mathcal I(r)}\,\dd r
\right|
=
\left|\psi(q)-\psi(p)\right|.
\]
又因 \(\psi\) 严格递增，且
\[
\psi(0)=2\arctan(\sqrt\varphi),
\qquad
\lim_{q\to\infty}\psi(q)=\pi,
\]
遂得所示像区间。证毕。
\end{proof}

\begin{corollary}[escort 温度半轴的 Fisher--Rao 完备化只需添加一个零温端点]\label{cor:xi-time-part9sb-escort-fisher-completion-zero-endpoint}
沿用定理 \ref{thm:xi-time-part9sb-escort-fisher-flat-coordinate} 的记号。则对任意 \(q\ge 0\)，有
\[
d_{\mathrm{FR}}(q,\infty)
:=
\lim_{r\to\infty}d_{\mathrm{FR}}(q,r)
=
\pi-\psi(q)
=
2\arctan\!\bigl(\varphi^{-(q+1)/2}\bigr).
\]
特别地，
\[
\operatorname{diam}_{\mathrm{FR}}[0,\infty)
=
2\arctan(\varphi^{-1/2}).
\]
因此，Fisher--Rao 完备化只需添加唯一端点 \(q=\infty\)；在极限一比特实验中，该端点对应于退化分布 \(\delta_0\)。
\end{corollary}
\begin{proof}
由定理 \ref{thm:xi-time-part9sb-escort-fisher-flat-coordinate}，
\[
d_{\mathrm{FR}}(q,r)=|\psi(r)-\psi(q)|.
\]
令 \(r\to\infty\) 即得
\[
d_{\mathrm{FR}}(q,\infty)=\pi-\psi(q).
\]
再用恒等式
\[
\pi-2\arctan x=2\arctan(x^{-1})
\qquad (x>0)
\]
并取 \(x=\varphi^{(q+1)/2}\)，便得到
\[
d_{\mathrm{FR}}(q,\infty)=2\arctan\!\bigl(\varphi^{-(q+1)/2}\bigr).
\]
取 \(q=0\) 即得总直径。

另一方面，
\[
\theta(q)=\frac{1}{1+\varphi^{q+1}}\longrightarrow 0
\qquad (q\to\infty),
\]
故极限 Bernoulli 模型退化为 \(\mathrm{Bernoulli}(0)=\delta_0\)。在 \(\psi\) 坐标下，\([0,\infty)\) 的像是半开区间 \([2\arctan(\sqrt\varphi),\pi)\)，其度量完备化只需补入端点 \(\pi\)，从而原半轴只需添加唯一零温端点。证毕。
\end{proof}

\begin{theorem}[escort 温度族的平方根嵌入是一段单位圆弧，Hellinger 亲和量满足精确余弦律]\label{thm:xi-time-part9sb-escort-sqrt-circle-hellinger}
定义平方根嵌入
\[
\Psi(q):=\bigl(\sqrt{1-\theta(q)},\,\sqrt{\theta(q)}\bigr)\in\RR^2.
\]
则有参数化
\[
\Psi(q)
=
\left(
\sin\frac{\psi(q)}{2},
\,
\cos\frac{\psi(q)}{2}
\right),
\qquad
\|\Psi(q)\|_2=1.
\]
因此 \(\Psi([0,\infty))\) 是第一象限中的一段单位圆弧。进一步，对任意 \(p,q\ge 0\)，若记
\[
\mathcal H(p,q)
:=
\sqrt{(1-\theta(p))(1-\theta(q))}
+
\sqrt{\theta(p)\theta(q)},
\]
则
\[
\mathcal H(p,q)
=
\left\langle \Psi(p),\Psi(q)\right\rangle
=
\cos\frac{d_{\mathrm{FR}}(p,q)}{2}.
\]
相应地，
\[
\|\Psi(p)-\Psi(q)\|_2^2
=
2-2\mathcal H(p,q)
=
4\sin^2\frac{d_{\mathrm{FR}}(p,q)}{4}.
\]
按结论 \ref{cor:xi-time-part9o-escort-testing-geometry-closed} 所采用的归一化，平方 Hellinger 距离
\[
H_{\mathrm{hel}}^2(p,q):=1-\mathcal H(p,q)
\]
因此满足
\[
H_{\mathrm{hel}}^2(p,q)=2\sin^2\frac{d_{\mathrm{FR}}(p,q)}{4}.
\]
\end{theorem}
\begin{proof}
令
\[
x_q:=\varphi^{(q+1)/2}.
\]
则
\[
\theta(q)=\frac{1}{1+x_q^2},
\qquad
1-\theta(q)=\frac{x_q^2}{1+x_q^2},
\qquad
\psi(q)=2\arctan x_q.
\]
于是半角公式给出
\[
\cos\frac{\psi(q)}{2}
=
\frac{1}{\sqrt{1+x_q^2}}
=
\sqrt{\theta(q)},
\]
\[
\sin\frac{\psi(q)}{2}
=
\frac{x_q}{\sqrt{1+x_q^2}}
=
\sqrt{1-\theta(q)}.
\]
故 \(\Psi(q)\) 确为单位圆上的上述参数化。

再计算内积：
\[
\left\langle \Psi(p),\Psi(q)\right\rangle
=
\sin\frac{\psi(p)}{2}\sin\frac{\psi(q)}{2}
+
\cos\frac{\psi(p)}{2}\cos\frac{\psi(q)}{2}
=
\cos\frac{\psi(q)-\psi(p)}{2}.
\]
由定理 \ref{thm:xi-time-part9sb-escort-fisher-flat-coordinate}，
\[
d_{\mathrm{FR}}(p,q)=|\psi(q)-\psi(p)|,
\]
而余弦为偶函数，故
\[
\mathcal H(p,q)=\cos\frac{d_{\mathrm{FR}}(p,q)}{2}.
\]
最后，由单位向量的弦长恒等式
\[
\|\Psi(p)-\Psi(q)\|_2^2
=
\|\Psi(p)\|_2^2+\|\Psi(q)\|_2^2-2\langle \Psi(p),\Psi(q)\rangle
=
2-2\mathcal H(p,q)
\]
即可得到
\[
\|\Psi(p)-\Psi(q)\|_2^2
=
4\sin^2\frac{d_{\mathrm{FR}}(p,q)}{4}.
\]
将 \(\mathcal H(p,q)=\cos(d_{\mathrm{FR}}(p,q)/2)\) 代回
\(
H_{\mathrm{hel}}^2(p,q)=1-\mathcal H(p,q)
\)
即得最后一式。证毕。
\end{proof}

\begin{theorem}[不同温度之间的 Chernoff 极小点具有显式闭式]\label{thm:xi-time-part9sb-escort-chernoff-saddle}
\leanverified{paper\_xi\_time\_part9sb\_escort\_chernoff\_saddle}
对 \(r\ge 0\) 定义
\[
A(r):=\varphi+\varphi^{-r}.
\]
对任意 \(p\neq q\) 与 \(s\in[0,1]\)，记
\[
\mathcal M_s(p,q)
:=
\frac{A(sp+(1-s)q)}{A(p)^sA(q)^{1-s}}.
\]
再定义 Chernoff 信息
\[
C_{\mathrm{Ch}}(p,q):=
-\log\inf_{0\le s\le 1}\mathcal M_s(p,q).
\]
若记
\[
c_{p,q}
:=
\frac{\log A(q)-\log A(p)}{(p-q)\log\varphi}
\in
\bigl(\theta(\max\{p,q\}),\,\theta(\min\{p,q\})\bigr)\subset(0,1),
\]
则
\[
r_*(p,q)
:=
\log_\varphi\!\Bigl(\frac{1-c_{p,q}}{c_{p,q}}\Bigr)-1
\in
\bigl(\min\{p,q\},\,\max\{p,q\}\bigr),
\]
并且 \(s\mapsto \mathcal M_s(p,q)\) 在 \([0,1]\) 上存在唯一极小点
\[
s_*(p,q)=\frac{r_*(p,q)-q}{p-q}\in(0,1).
\]
进一步，
\[
C_{\mathrm{Ch}}(p,q)
=
s_*(p,q)\log A(p)
+
\bigl(1-s_*(p,q)\bigr)\log A(q)
-
\log A(r_*(p,q)).
\]
\end{theorem}
\begin{proof}
由定理 \ref{thm:xi-time-part65h-escort-hellinger-chernoff-kernel}，
\[
\log \mathcal M_s(p,q)
=
\log A(r_s)-s\log A(p)-(1-s)\log A(q),
\qquad
r_s:=sp+(1-s)q.
\]
因此
\[
\frac{\dd}{\dd s}\log \mathcal M_s(p,q)
=
(p-q)\frac{A'(r_s)}{A(r_s)}-\log A(p)+\log A(q).
\]
又因为
\[
A'(r)=-(\log\varphi)\varphi^{-r},
\qquad
-\frac{A'(r)}{(\log\varphi)A(r)}
=
\frac{\varphi^{-r}}{\varphi+\varphi^{-r}}
=
\frac{1}{1+\varphi^{r+1}}
=
\theta(r),
\]
临界点条件等价于
\[
\theta(r_s)=c_{p,q}.
\]

现在令 \(g(r):=\log A(r)\)。由均值定理，存在 \(\xi\) 介于 \(p\) 与 \(q\) 之间，使得
\[
g(q)-g(p)=g'(\xi)(q-p).
\]
于是
\[
c_{p,q}
=
\frac{g(q)-g(p)}{(p-q)\log\varphi}
=
-\frac{g'(\xi)}{\log\varphi}
=
\theta(\xi),
\]
从而
\[
c_{p,q}\in
\bigl(\theta(\max\{p,q\}),\,\theta(\min\{p,q\})\bigr).
\]
由于 \(\theta\) 在 \([0,\infty)\) 上严格递减，方程 \(\theta(r)=c_{p,q}\) 存在唯一解，且该解恰落在 \((\min\{p,q\},\max\{p,q\})\) 中。由
\[
\theta(r)=\frac{1}{1+\varphi^{r+1}}
\]
直接反解，便得到
\[
r_*(p,q)=\log_\varphi\!\Bigl(\frac{1-c_{p,q}}{c_{p,q}}\Bigr)-1.
\]
再由 \(r_s=sp+(1-s)q\) 得
\[
s_*(p,q)=\frac{r_*(p,q)-q}{p-q}\in(0,1).
\]

最后，对二阶导数求值：
\[
\frac{\dd^2}{\dd s^2}\log \mathcal M_s(p,q)
=
(p-q)^2\frac{\dd}{\dd r}\!\left(\frac{A'(r)}{A(r)}\right)_{r=r_s}
=
-(p-q)^2(\log\varphi)\theta'(r_s).
\]
而
\[
\theta'(r)=-(\log\varphi)\theta(r)\bigl(1-\theta(r)\bigr),
\]
故
\[
\frac{\dd^2}{\dd s^2}\log \mathcal M_s(p,q)
=
(p-q)^2(\log\varphi)^2\theta(r_s)\bigl(1-\theta(r_s)\bigr)>0.
\]
因此 \(\log\mathcal M_s(p,q)\) 严格凸，极小点唯一，且
\[
C_{\mathrm{Ch}}(p,q)
=
-\log \mathcal M_{s_*}(p,q)
=
s_*(p,q)\log A(p)
+
\bigl(1-s_*(p,q)\bigr)\log A(q)
-
\log A(r_*(p,q)).
\]
证毕。
\end{proof}

\begin{theorem}[固定样本数下 escort 乘积试验在 Blackwell 意义下指数塌缩为二项试验]\label{thm:xi-time-part9sb-escort-fixed-n-binomial-blackwell}
设 \(Q<\infty\) 与 \(n\in\NN\)。定义 \(n\)-样本 escort 试验族
\[
\mathcal E_m^{(n,Q)}:=\{\pi_{m,q}^{\otimes n}:0\le q\le Q\},
\]
以及二项极限试验族
\[
\mathcal G^{(n,Q)}:=\{\mathrm{Binomial}(n,\theta(q)):0\le q\le Q\},
\qquad
\theta(q)=\frac{1}{1+\varphi^{q+1}}.
\]
则存在常数 \(C_{n,Q}>0\) 使得
\[
\Delta\bigl(\mathcal E_m^{(n,Q)},\mathcal G^{(n,Q)}\bigr)
\le
C_{n,Q}\Bigl(\frac{\varphi}{2}\Bigr)^m.
\]
特别地，前向塌缩核正是末位计数
\[
N_n(w^{(1)},\dots,w^{(n)})
:=
\sum_{j=1}^n w_m^{(j)}\in\{0,\dots,n\}.
\]
对任意固定 \(p,q\in[0,Q]\)，极限对数似然比仅依赖于 \(N_n\)：
\[
\Lambda_n(p,q)
=
N_n\log\frac{\theta(q)}{\theta(p)}
+
\bigl(n-N_n\bigr)\log\frac{1-\theta(q)}{1-\theta(p)}.
\]
等价地，在与 \(\mathcal G^{(n,Q)}\) 对应的 Bernoulli 乘积极限试验中，\(N_n\) 承载了全部参数依赖。
\end{theorem}
\begin{proof}
记
\[
A_i:=\{w\in X_m:w_m=i\},
\qquad
U_{m,i}:=\mathrm{Unif}(A_i)
\qquad (i=0,1).
\]
并记单样本末位读出
\[
K_m(w):=w_m\in\{0,1\}.
\]
定义前向核
\[
R_m(w^{(1)},\dots,w^{(n)})
:=
\sum_{j=1}^n w_m^{(j)}.
\]
再定义反向核
\[
S_m(k,\cdot)
:=
\frac{1}{\binom{n}{k}}
\sum_{\substack{J\subset\{1,\dots,n\}\\ |J|=k}}
\bigotimes_{j\in J}U_{m,1}
\otimes
\bigotimes_{j\notin J}U_{m,0},
\qquad 0\le k\le n.
\]

由定理 \ref{thm:xi-time-part65h-escort-blackwell-onebit-equivalence} 的证明，单样本末位分布精确满足
\[
K_{m\#}\pi_{m,q}=\mathrm{Bernoulli}(\theta_m(q)),
\qquad
\theta_m(q):=
\frac{\varphi^{-q}F_m}{F_{m+1}+\varphi^{-q}F_m},
\]
并且
\[
\sup_{0\le q\le Q}|\theta_m(q)-\theta(q)|
=
O_Q\!\left(\Bigl(\frac{\varphi}{2}\Bigr)^m\right).
\]
因此在 \(n\) 个独立样本下，
\[
R_{m\#}\pi_{m,q}^{\otimes n}
=
\mathrm{Binomial}(n,\theta_m(q)).
\]
又因为计数映射不会增加总变差，而 Bernoulli 乘积分布满足望远镜估计，
\[
D_{\mathrm{TV}}\!\bigl(
\mathrm{Binomial}(n,\theta_m(q)),
\mathrm{Binomial}(n,\theta(q))
\bigr)
\le
D_{\mathrm{TV}}\!\bigl(
\mathrm{Bernoulli}(\theta_m(q))^{\otimes n},
\mathrm{Bernoulli}(\theta(q))^{\otimes n}
\bigr)
\le
n\,|\theta_m(q)-\theta(q)|.
\]
故前向 deficiency 由
\[
\sup_{0\le q\le Q}
D_{\mathrm{TV}}\!\bigl(R_{m\#}\pi_{m,q}^{\otimes n},\mathrm{Binomial}(n,\theta(q))\bigr)
\le
C_{n,Q}\Bigl(\frac{\varphi}{2}\Bigr)^m
\]
控制。

再看反向方向。对任意 \(\theta\in[0,1]\)，由二项分布与均匀放置 \(k\) 个 \(1\) 的构造，
\[
S_{m\#}\mathrm{Binomial}(n,\theta)
=
\bigl((1-\theta)U_{m,0}+\theta U_{m,1}\bigr)^{\otimes n}.
\]
特别地，
\[
S_{m\#}\mathrm{Binomial}(n,\theta(q))
=
\bigl((1-\theta(q))U_{m,0}+\theta(q)U_{m,1}\bigr)^{\otimes n}.
\]
而定理 \ref{thm:xi-time-part65h-escort-blackwell-onebit-equivalence} 已给出
\[
\sup_{0\le q\le Q}
D_{\mathrm{TV}}\!\bigl(
(1-\theta(q))U_{m,0}+\theta(q)U_{m,1},
\pi_{m,q}
\bigr)
=
O_Q\!\left(\Bigl(\frac{\varphi}{2}\Bigr)^m\right).
\]
对乘积测度再次应用望远镜估计，得到
\[
\sup_{0\le q\le Q}
D_{\mathrm{TV}}\!\bigl(
S_{m\#}\mathrm{Binomial}(n,\theta(q)),
\pi_{m,q}^{\otimes n}
\bigr)
\le
C_{n,Q}\Bigl(\frac{\varphi}{2}\Bigr)^m.
\]
两侧 deficiency 同阶受控，故所示 Le Cam 距离估计成立。

最后，二项族的似然函数可写为
\[
\mathbf P_q(N_n=k)
=
\binom{n}{k}\theta(q)^k\bigl(1-\theta(q)\bigr)^{n-k},
\]
于是对参数 \(p,q\) 的对数似然比恰为
\[
\Lambda_n(p,q)
=
k\log\frac{\theta(q)}{\theta(p)}
+
\bigl(n-k\bigr)\log\frac{1-\theta(q)}{1-\theta(p)}.
\]
这表明与 \(\mathcal G^{(n,Q)}\) 等价的 Bernoulli 乘积极限试验完全由计数 \(N_n\) 承载。证毕。
\end{proof}

\begin{corollary}[等先验温度判别的最优误差指数与顺序对数似然步长都具有闭式]\label{cor:xi-time-part9sb-escort-error-exponent-lattice-walk}
\leanverified{paper\_xi\_time\_part9sb\_escort\_error\_exponent\_lattice\_walk}
固定 \(p\neq q\)。令
\[
P_{e,m}^{\ast}(n;p,q)
\]
表示在等先验 \((1/2,1/2)\) 下、基于 \(n\) 个独立样本区分 \(\pi_{m,p}\) 与 \(\pi_{m,q}\) 的最优 Bayes 错误概率。则
\[
\lim_{n\to\infty}
-\frac1n
\log\!\Bigl(
\lim_{m\to\infty}P_{e,m}^{\ast}(n;p,q)
\Bigr)
=
C_{\mathrm{Ch}}(p,q),
\]
其中 \(C_{\mathrm{Ch}}(p,q)\) 为定理 \ref{thm:xi-time-part9sb-escort-chernoff-saddle} 中的 Chernoff 信息。

再记
\[
\beta_s:=\mathrm{Bernoulli}(\theta(s)),
\qquad s\ge 0,
\]
并在真参数 \(q\) 下定义一步对数似然增量
\[
L_{p,q}:=\log\frac{\dd\beta_q}{\dd\beta_p}(B),
\qquad
B\sim\beta_q.
\]
则
\[
L_{p,q}
=
\begin{cases}
\log\dfrac{\theta(q)}{\theta(p)}, & \text{概率 }\theta(q),\\[1.2ex]
\log\dfrac{1-\theta(q)}{1-\theta(p)}, & \text{概率 }1-\theta(q).
\end{cases}
\]
因此
\[
\mathbb E_q[L_{p,q}]
=
D_{\mathrm{KL}}(\beta_q\Vert \beta_p),
\]
\[
\operatorname{Var}_q(L_{p,q})
=
\theta(q)\bigl(1-\theta(q)\bigr)
\left[
\log\frac{\theta(q)(1-\theta(p))}{\theta(p)(1-\theta(q))}
\right]^2.
\]
更具体地，若 \((L_{p,q}^{(j)})_{j\ge 1}\) 为独立同分布拷贝，则
\[
\sum_{j=1}^n L_{p,q}^{(j)}
=
n\log\frac{1-\theta(q)}{1-\theta(p)}
+
N_n\log\frac{\theta(q)(1-\theta(p))}{\theta(p)(1-\theta(q))},
\]
故其取值落在一条严格仿射格上。
\end{corollary}
\begin{proof}
取 \(Q>\max\{p,q\}\)。由定理
\ref{thm:xi-time-part9sb-escort-fixed-n-binomial-blackwell}，对每个固定 \(n\)，\(n\)-样本 escort 试验与极限二项试验在 \(m\to\infty\) 时的 Le Cam 距离趋于 \(0\)。因此，等先验 Bayes 最优风险收敛到二项极限试验中的对应风险：
\[
\lim_{m\to\infty}P_{e,m}^{\ast}(n;p,q)
=
P_{e,\infty}^{\ast}(n;p,q),
\]
其中右端为区分 \(\mathrm{Binomial}(n,\theta(p))\) 与 \(\mathrm{Binomial}(n,\theta(q))\) 的最优 Bayes 错误概率。对该 i.i.d. Bernoulli 模型应用经典 Chernoff 定理，即得
\[
\lim_{n\to\infty}
-\frac1n\log P_{e,\infty}^{\ast}(n;p,q)
=
C_{\mathrm{Ch}}(p,q).
\]
这就证明了第一式。

再看一步对数似然增量。由于 \(B\in\{0,1\}\)，直接计算可得
\[
\log\frac{\dd\beta_q}{\dd\beta_p}(1)=\log\frac{\theta(q)}{\theta(p)},
\qquad
\log\frac{\dd\beta_q}{\dd\beta_p}(0)=\log\frac{1-\theta(q)}{1-\theta(p)}.
\]
故 \(L_{p,q}\) 只取上述两值，其分布即由 \(B\sim\beta_q\) 给出。于是
\[
\mathbb E_q[L_{p,q}]
=
\theta(q)\log\frac{\theta(q)}{\theta(p)}
+
\bigl(1-\theta(q)\bigr)\log\frac{1-\theta(q)}{1-\theta(p)}
=
D_{\mathrm{KL}}(\beta_q\Vert \beta_p).
\]
若记
\[
a:=\log\frac{\theta(q)}{\theta(p)},
\qquad
b:=\log\frac{1-\theta(q)}{1-\theta(p)},
\]
则二点分布方差公式给出
\[
\operatorname{Var}_q(L_{p,q})
=
\theta(q)\bigl(1-\theta(q)\bigr)(a-b)^2,
\]
即
\[
\operatorname{Var}_q(L_{p,q})
=
\theta(q)\bigl(1-\theta(q)\bigr)
\left[
\log\frac{\theta(q)(1-\theta(p))}{\theta(p)(1-\theta(q))}
\right]^2.
\]
最后，由定理 \ref{thm:xi-time-part9sb-escort-fixed-n-binomial-blackwell} 中的似然比公式，
\[
\sum_{j=1}^n L_{p,q}^{(j)}
=
N_n a+(n-N_n)b
=
n b+N_n(a-b),
\]
故其状态空间是仿射格
\[
n\log\frac{1-\theta(q)}{1-\theta(p)}
+
\log\frac{\theta(q)(1-\theta(p))}{\theta(p)(1-\theta(q))}\,\ZZ.
\]
因此极限顺序对数似然过程是一条严格二点格随机游走。证毕。
\end{proof}

由此，escort 温度半轴的极限模型不再只是一个抽象的一维 logistic 族；在 Fisher--Rao 度量下，它是一条带唯一完备端点的平坦线段，在平方根嵌入下，它是一段单位圆弧，而在判别论上又同时呈现 Chernoff 极小点闭式、固定样本二项充分化与严格二点格对数似然过程。于是几何、Blackwell 缺损与顺序检验三者被压缩到同一组显式公式之中。

\endinput


\endinput

\paragraph{原始 bin-fold 切片的显式二点实验、\texorpdfstring{$\chi^2$}{chi2} 反演与二点评价闭包}
在 \(u_m=\pi_{m,0}\) 与 \(\mu_m=\pi_{m,1}\) 的末位充分化、两点似然比展开以及幂散度极限闭式都已建立之后，还可把原始切片与稳定层均匀基线之间的统计几何进一步压缩为一个固定二点实验。其后果是：Le Cam 等价、最优判别、\(\chi^2\) 常数反演、全阶 Rényi 谱以及全部 \(f\)-散度极限，都可由同一末位二相结构统一闭合。

\begin{theorem}[原始 bin-fold 切片与显式二点实验的指数等价]\label{thm:xi-time-part9sa-uniform-baseline-fixed-binary-experiment}
\leanverified{paper\_xi\_time\_part9sa\_uniform\_baseline\_fixed\_binary\_experiment}
设 \(u_m=\pi_{m,0}\) 为 \(X_m\) 上的均匀分布，\(\mu_m=\pi_{m,1}\) 为原始 bin-fold 分布。记
\[
A_{m,b}:=\{w\in X_m:\ w_m=b\},
\qquad
\ell_m(w):=w_m\in\{0,1\},
\qquad
b\in\{0,1\}.
\]
在 \(\{0,1\}\) 上定义固定二点实验
\[
Q(0)=\frac1\varphi,
\qquad
Q(1)=\frac1{\varphi^2},
\]
\[
R(0)=\frac{\varphi}{\sqrt5},
\qquad
R(1)=\frac{1}{\varphi\sqrt5}.
\]
再记
\[
Q_m^{\mathrm{bit}}:=(\ell_m)_\#u_m,
\qquad
R_m^{\mathrm{bit}}:=(\ell_m)_\#\mu_m,
\]
以及从 \(\{0,1\}\) 回到 \(X_m\) 的 Markov 核
\[
L_m(b,\cdot):=u_m(\cdot\mid A_{m,b})
\qquad (b=0,1).
\]
则有
\[
D_{\mathrm{TV}}\bigl(Q_m^{\mathrm{bit}},Q\bigr)=O(\varphi^{-2m}),
\qquad
D_{\mathrm{TV}}\bigl(R_m^{\mathrm{bit}},R\bigr)
=
O\!\left(\left(\frac{\varphi}{2}\right)^m\right),
\]
并且
\[
u_m=L_{m\#}Q_m^{\mathrm{bit}},
\qquad
D_{\mathrm{TV}}\bigl(\mu_m,L_{m\#}R_m^{\mathrm{bit}}\bigr)
=
O\!\left(\left(\frac{\varphi}{2}\right)^m\right).
\]
从而把 \((u_m,\mu_m)\) 与 \((Q,R)\) 视为分别定义在 \(X_m\) 与 \(\{0,1\}\) 上的二元统计实验时，有
\[
\Delta_{\mathrm{LeCam}}\bigl((u_m,\mu_m),(Q,R)\bigr)
=
O\!\left(\left(\frac{\varphi}{2}\right)^m\right).
\]
\end{theorem}
\begin{proof}
由 \(|A_{m,0}|=F_{m+1}\)、\(|A_{m,1}|=F_m\) 与 \(|X_m|=F_{m+2}\) 得
\[
Q_m^{\mathrm{bit}}(0)=\frac{F_{m+1}}{F_{m+2}},
\qquad
Q_m^{\mathrm{bit}}(1)=\frac{F_m}{F_{m+2}}.
\]
Binet 公式给出
\[
\frac{F_{m+1}}{F_{m+2}}=\frac1\varphi+O(\varphi^{-2m}),
\qquad
\frac{F_m}{F_{m+2}}=\frac1{\varphi^2}+O(\varphi^{-2m}),
\]
故第一条全变差估计成立。

另一方面，定理 \ref{thm:xi-time-part64b-fold-information-geometry-two-point-collapse} 的证明给出一致展开
\[
\frac{\mu_m(w)}{u_m(w)}
=
\frac{\varphi^{2-\ell_m(w)}}{\sqrt5}
+
O\!\left(\left(\frac{\varphi}{2}\right)^m\right)
\qquad (w\in X_m).
\]
对每个 \(b\in\{0,1\}\) 在块 \(A_{m,b}\) 上按 \(u_m\) 求和，得到
\[
R_m^{\mathrm{bit}}(b)
=
\frac{\varphi^{2-b}}{\sqrt5}\,Q_m^{\mathrm{bit}}(b)
+
O\!\left(\left(\frac{\varphi}{2}\right)^m\right).
\]
再代入 \(Q_m^{\mathrm{bit}}(0)=\varphi^{-1}+O(\varphi^{-2m})\) 与
\(Q_m^{\mathrm{bit}}(1)=\varphi^{-2}+O(\varphi^{-2m})\)，便得
\[
R_m^{\mathrm{bit}}(0)=\frac{\varphi}{\sqrt5}
+
O\!\left(\left(\frac{\varphi}{2}\right)^m\right),
\qquad
R_m^{\mathrm{bit}}(1)=\frac{1}{\varphi\sqrt5}
+
O\!\left(\left(\frac{\varphi}{2}\right)^m\right),
\]
从而
\[
D_{\mathrm{TV}}\bigl(R_m^{\mathrm{bit}},R\bigr)
=
O\!\left(\left(\frac{\varphi}{2}\right)^m\right).
\]

由于 \(u_m\) 在每个 \(A_{m,b}\) 上本来就是均匀分布，故
\[
u_m=L_{m\#}Q_m^{\mathrm{bit}}
\]
为精确恒等式。再对 \(w\in A_{m,b}\) 计算
\[
L_{m\#}R_m^{\mathrm{bit}}(w)
=
\frac{R_m^{\mathrm{bit}}(b)}{u_m(A_{m,b})}\,u_m(w)
=
u_m(w)\left(
\frac{\varphi^{2-b}}{\sqrt5}
+
O\!\left(\left(\frac{\varphi}{2}\right)^m\right)
\right).
\]
与上面的似然比一致展开比较，即得
\[
L_{m\#}R_m^{\mathrm{bit}}(w)
=
\mu_m(w)
+
u_m(w)\,O\!\left(\left(\frac{\varphi}{2}\right)^m\right)
\]
在 \(w\in X_m\) 上一致成立。对 \(w\) 求和便得
\[
D_{\mathrm{TV}}\bigl(\mu_m,L_{m\#}R_m^{\mathrm{bit}}\bigr)
=
O\!\left(\left(\frac{\varphi}{2}\right)^m\right).
\]

最后，以前向核 \(\ell_m\) 和后向核 \(L_m\) 分别估计两侧 deficiency。由 Markov 核对总变差的压缩性，
\[
\delta\bigl((u_m,\mu_m),(Q,R)\bigr)
\le
\max\Bigl\{
D_{\mathrm{TV}}\bigl(Q_m^{\mathrm{bit}},Q\bigr),
D_{\mathrm{TV}}\bigl(R_m^{\mathrm{bit}},R\bigr)
\Bigr\},
\]
\[
\delta\bigl((Q,R),(u_m,\mu_m)\bigr)
\le
\max\Bigl\{
D_{\mathrm{TV}}\bigl(L_{m\#}Q,u_m\bigr),
D_{\mathrm{TV}}\bigl(L_{m\#}R,\mu_m\bigr)
\Bigr\}.
\]
而
\[
D_{\mathrm{TV}}\bigl(L_{m\#}Q,u_m\bigr)
=
D_{\mathrm{TV}}\bigl(L_{m\#}Q,L_{m\#}Q_m^{\mathrm{bit}}\bigr)
\le
D_{\mathrm{TV}}\bigl(Q,Q_m^{\mathrm{bit}}\bigr),
\]
\[
D_{\mathrm{TV}}\bigl(L_{m\#}R,\mu_m\bigr)
\le
D_{\mathrm{TV}}\bigl(L_{m\#}R,L_{m\#}R_m^{\mathrm{bit}}\bigr)
+
D_{\mathrm{TV}}\bigl(L_{m\#}R_m^{\mathrm{bit}},\mu_m\bigr)
\le
D_{\mathrm{TV}}\bigl(R,R_m^{\mathrm{bit}}\bigr)
+
O\!\left(\left(\frac{\varphi}{2}\right)^m\right).
\]
故两侧 deficiency 均为 \(O((\varphi/2)^m)\)，结论成立。
\end{proof}

\begin{corollary}[原始 bin-fold 与稳定层均匀基线的最优判别由末位比特完全决定]\label{cor:xi-time-part9sa-uniform-baseline-lastbit-neyman-pearson}
沿用定理 \ref{thm:xi-time-part9sa-uniform-baseline-fixed-binary-experiment} 的记号。对等先验的一样本判别，充分大 \(m\) 时 Neyman--Pearson 最强检验恰为
\[
\delta_m(w)=
\begin{cases}
\text{判为 }\mu_m,& w_m=0,\\
\text{判为 }u_m,& w_m=1.
\end{cases}
\]
并且
\[
D_{\mathrm{TV}}(\mu_m,u_m)
=
1-\frac{2}{\sqrt5}
+
O\!\left(\left(\frac{\varphi}{2}\right)^m\right),
\]
\[
P_{\mathrm e}^\ast(m)
=
\frac1{\sqrt5}
+
O\!\left(\left(\frac{\varphi}{2}\right)^m\right).
\]
特别地，
\[
D_{\mathrm{TV}}(\mu_m,u_m)\longrightarrow 1-\frac{2}{\sqrt5},
\qquad
P_{\mathrm e}^\ast(m)\longrightarrow \frac1{\sqrt5}.
\]
\end{corollary}
\begin{proof}
由定理 \ref{thm:xi-time-part64b-fold-information-geometry-two-point-collapse} 的一致展开，
\[
\frac{\mu_m(w)}{u_m(w)}
=
\frac{\varphi^{2-w_m}}{\sqrt5}
+
O\!\left(\left(\frac{\varphi}{2}\right)^m\right),
\]
而
\[
\frac{\varphi^2}{\sqrt5}>1>\frac{\varphi}{\sqrt5}.
\]
因此对充分大的 \(m\)，似然比是否超过 \(1\) 完全由 \(w_m\) 决定，从而最强检验恰如所述。

再记
\[
\bar\mu_m:=L_{m\#}R_m^{\mathrm{bit}}.
\]
由定理 \ref{thm:xi-time-part9sa-uniform-baseline-fixed-binary-experiment}，
\[
D_{\mathrm{TV}}(\mu_m,\bar\mu_m)
=
O\!\left(\left(\frac{\varphi}{2}\right)^m\right).
\]
另一方面，\(\bar\mu_m\) 与 \(u_m=L_{m\#}Q_m^{\mathrm{bit}}\) 在相同分块
\(A_{m,0}\sqcup A_{m,1}\) 上都具有块内均匀的条件分布，故
\[
D_{\mathrm{TV}}(\bar\mu_m,u_m)
=
D_{\mathrm{TV}}\bigl(R_m^{\mathrm{bit}},Q_m^{\mathrm{bit}}\bigr).
\]
于是
\[
D_{\mathrm{TV}}(\mu_m,u_m)
=
D_{\mathrm{TV}}\bigl(R_m^{\mathrm{bit}},Q_m^{\mathrm{bit}}\bigr)
+
O\!\left(\left(\frac{\varphi}{2}\right)^m\right)
=
D_{\mathrm{TV}}(R,Q)
+
O\!\left(\left(\frac{\varphi}{2}\right)^m\right).
\]
对二点实验有
\[
D_{\mathrm{TV}}(R,Q)=|R(1)-Q(1)|
=
\frac{1}{\varphi^2}-\frac{1}{\varphi\sqrt5}
=
1-\frac{2}{\sqrt5}.
\]
故得到总变差展开。最后由 Helstrom 公式
\[
P_{\mathrm e}^\ast(m)
=
\frac12\bigl(1-D_{\mathrm{TV}}(\mu_m,u_m)\bigr)
\]
立即得到 Bayes 最小错误率的展开与极限。
\end{proof}

\begin{theorem}[bin-fold 的 \texorpdfstring{$\chi^2$}{chi2} 极限常数对黄金比例的仿射反演]\label{thm:xi-time-part9sa-chi2-affine-phi-inversion}
\leanverified{paper\_xi\_time\_part9sa\_chi2\_affine\_phi\_inversion}
记
\[
D_\infty^{\chi^2}
:=
\lim_{m\to\infty}\chi^2(\mu_m\Vert u_m).
\]
则有
\[
D_\infty^{\chi^2}
=
\chi^2(R\Vert Q)
=
\frac{2\varphi-3}{5}
=
\frac{\sqrt5-2}{5},
\]
从而黄金比例可由单个 \(\chi^2\) 极限常数线性反演为
\[
\varphi=\frac{5D_\infty^{\chi^2}+3}{2}.
\]
\end{theorem}
\begin{proof}
由定理 \ref{thm:xi-time-part67-lastbit-terminal-csiszar-collapse} 在
\[
f(t)=(t-1)^2
\]
处的专化，可得
\[
D_\infty^{\chi^2}
=
D_f(R\Vert Q)
=
\chi^2(R\Vert Q).
\]
直接计算
\[
\chi^2(R\Vert Q)
=
\frac{R(0)^2}{Q(0)}+\frac{R(1)^2}{Q(1)}-1
=
\frac{1}{5}\bigl(\varphi^3+1\bigr)-1.
\]
再用恒等式
\[
\varphi^2=\varphi+1,
\qquad
\varphi^3=2\varphi+1,
\]
便得到
\[
D_\infty^{\chi^2}
=
\frac{2\varphi-3}{5}.
\]
最后直接解出 \(\varphi\) 即得
\[
\varphi=\frac{5D_\infty^{\chi^2}+3}{2}.
\]
\end{proof}

\begin{theorem}[bin-fold 极限 Rényi 谱与显式 Bernoulli--Bernoulli 实验的逐阶同一]\label{thm:xi-time-part9sa-renyi-spectrum-explicit-bernoulli-experiment}
\leanverified{paper\_xi\_time\_part9sa\_renyi\_spectrum\_explicit\_bernoulli\_experiment}
沿用定理 \ref{thm:xi-time-part9sa-uniform-baseline-fixed-binary-experiment} 的记号。定义二点分布
\[
\nu_*:=
\left(
\frac{\varphi}{\sqrt5},
\frac{1}{\varphi\sqrt5}
\right),
\qquad
\eta_*:=
\left(
\frac{1}{\varphi},
\frac{1}{\varphi^2}
\right).
\]
对任意 \(\alpha>0\)，其中 \(\alpha=1\) 按连续延拓理解为 \(D_1=D_{\mathrm{KL}}\)，记
\[
D_\alpha^\infty:=\lim_{m\to\infty}D_\alpha(\mu_m\Vert u_m).
\]
则有精确恒等式
\[
D_\alpha^\infty=D_\alpha(\nu_*\Vert \eta_*).
\]
换言之，bin-fold 输出分布相对于稳定层均匀基线的整个极限 Rényi 散度族，与一个显式 Bernoulli--Bernoulli 实验严格同一。
\end{theorem}
\begin{proof}
由定理 \ref{thm:xi-time-part9sa-uniform-baseline-fixed-binary-experiment}，
\[
R(0)=\frac{\varphi}{\sqrt5},
\qquad
R(1)=\frac{1}{\varphi\sqrt5},
\qquad
Q(0)=\frac{1}{\varphi},
\qquad
Q(1)=\frac{1}{\varphi^2},
\]
故 \(\nu_*=R\)、\(\eta_*=Q\)。

当 \(\alpha\neq 1\) 时，由命题 \ref{prop:xi-foldbin-renyi-deficit-spectrum-closed-form}，
\[
D_\alpha^\infty
=
\frac{1}{\alpha-1}\log\left(
\frac{1}{\varphi}\Bigl(\frac{\varphi^2}{\sqrt5}\Bigr)^\alpha
+
\frac{1}{\varphi^2}\Bigl(\frac{\varphi}{\sqrt5}\Bigr)^\alpha
\right).
\]
另一方面，二点 Rényi 散度直接计算为
\[
D_\alpha(\nu_*\Vert \eta_*)
=
\frac{1}{\alpha-1}\log\left(
\nu_*(0)^\alpha\eta_*(0)^{1-\alpha}
+
\nu_*(1)^\alpha\eta_*(1)^{1-\alpha}
\right),
\]
代入 \(\nu_*,\eta_*\) 的显式表达后，恰得到同一闭式。因此
\[
D_\alpha^\infty=D_\alpha(\nu_*\Vert \eta_*)
\qquad
(\alpha>0,\ \alpha\neq 1).
\]

当 \(\alpha=1\) 时，二点 KL 散度为
\[
D_{\mathrm{KL}}(\nu_*\Vert \eta_*)
=
\frac{\varphi}{\sqrt5}\log\frac{\varphi^2}{\sqrt5}
+
\frac{1}{\varphi\sqrt5}\log\frac{\varphi}{\sqrt5},
\]
这与推论 \ref{cor:xi-foldbin-shannon-deficit-constant-closed-form} 给出的 \(D_1^\infty\) 完全一致。证毕。
\end{proof}

\begin{corollary}[全部 \texorpdfstring{$f$}{f}-散度极限的二点评价闭包]\label{cor:xi-time-part9sa-fdiv-two-point-evaluation-closure}\leanverified{paper\_xi\_time\_part9sa\_fdiv\_two\_point\_evaluation\_closure}
设
\[
a:=\frac{\varphi^2}{\sqrt5},
\qquad
b:=\frac{\varphi}{\sqrt5},
\]
并对任意满足本文 \(f\)-散度可定义条件的生成元 \(f\) 记
\[
D_f^\infty:=\lim_{m\to\infty}D_f(\mu_m\Vert u_m).
\]
则有
\[
D_f^\infty
=
D_f(R\Vert Q)
=
\frac1\varphi\,f(a)+\frac1{\varphi^2}\,f(b).
\]
因此若 \(f(a)=g(a)\) 且 \(f(b)=g(b)\)，则
\[
D_f^\infty=D_g^\infty.
\]
进一步，若 \(\mathcal G\) 是任意包含本文所用生成元的实线性空间，并记
\[
\mathcal K_{\mathcal G}:=\{h\in\mathcal G:\ h(a)=h(b)=0\},
\]
则极限映射
\[
f\longmapsto D_f^\infty
\]
在商空间 \(\mathcal G/\mathcal K_{\mathcal G}\) 上由线性泛函
\[
[f]\longmapsto \frac1\varphi\,f(a)+\frac1{\varphi^2}\,f(b)
\]
完全决定。
\end{corollary}
\begin{proof}
定理 \ref{thm:xi-time-part67-lastbit-terminal-csiszar-collapse} 已给出
\[
D_f^\infty=\frac1\varphi\,f(a)+\frac1{\varphi^2}\,f(b).
\]
另一方面，由
\[
\frac{R(0)}{Q(0)}=a,
\qquad
\frac{R(1)}{Q(1)}=b,
\]
可知
\[
D_f(R\Vert Q)
=
Q(0)f(a)+Q(1)f(b)
=
\frac1\varphi\,f(a)+\frac1{\varphi^2}\,f(b).
\]
故第一条等式成立。其余结论只是该显式公式的直接推论。
\end{proof}

\paragraph{Fourier 非均匀性底噪、容量 Laplace 变分与 bin-fold 正温双折线}
在 Fibonacci 共振模的非熄灭、均匀层对数多重度的大偏差接口以及 bin-fold 临界容量双折线都已建立之后，还可进一步抽取出下列六条彼此衔接的后果。它们分别把单一非衰减 Fourier 模提升为相对均匀基线的常数全变差缺口与正相对熵地板，把均匀层上尾大偏差转写为连续容量指数的 Laplace 变分公式，把最大纤维能级间隙转写为容量曲线的精确仿射段，并把 bin-fold 的正实压力轴与临界资源窗一并闭式化。

\begin{theorem}[非衰减 Fibonacci Fourier 模强制相对均匀基线的常数全变差缺口]\label{thm:xi-time-part9t-fourier-tv-gap-positive-floor}
记
\[
M:=F_{m+2},
\qquad
p_m(r):=\frac{d_m(r)}{2^m},
\qquad
u(r):=\frac{1}{M}
\qquad
\bigl(r\in \ZZ/(M\ZZ)\bigr),
\]
并定义 Fourier 系数
\[
\widehat p_m(k):=
\sum_{r\in \ZZ/(M\ZZ)}p_m(r)e^{-2\pi i kr/M}.
\]
设 \(C_\varphi>0\) 为定理 \ref{thm:fold-critical-resonance-constant} 的无限乘积常数。则有
\[
\liminf_{m\to\infty}\|p_m-u\|_{\mathrm{TV}}
\ge
\frac{C_\varphi}{2}.
\]
\end{theorem}
\begin{proof}
对任意 \(k\in\ZZ/(M\ZZ)\)，令
\[
f_k(r):=e^{-2\pi i kr/M},
\qquad |f_k(r)|=1.
\]
于是
\[
\widehat p_m(k)-\widehat u(k)
=
\mathbb E_{p_m}[f_k]-\mathbb E_u[f_k].
\]
由总变差与有界测试函数的标准对偶估计，
\[
\bigl|\mathbb E_{p_m}[f_k]-\mathbb E_u[f_k]\bigr|
\le
2\|p_m-u\|_{\mathrm{TV}}.
\]
若 \(k\not\equiv 0\pmod M\)，则均匀分布的 Fourier 系数满足 \(\widehat u(k)=0\)。因此
\[
\|p_m-u\|_{\mathrm{TV}}
\ge
\frac12\,|\widehat p_m(k)|.
\]
现取 \(k=F_m\)。由于 \(0<F_m<F_{m+2}=M\)，故 \(F_m\not\equiv 0\pmod M\)。再由定义 \ref{def:fold-normalized-amplitude} 与定理 \ref{thm:fold-critical-resonance-constant}，
\[
|\widehat p_m(F_m)|
=
\frac{|\widehat d_m(F_m)|}{2^m}
=
a_m(F_m)
\longrightarrow
C_\varphi.
\]
将此代入前述估计，即得
\[
\liminf_{m\to\infty}\|p_m-u\|_{\mathrm{TV}}
\ge
\frac{C_\varphi}{2}.
\]
证毕。
\end{proof}

\begin{corollary}[相对熵与 Shannon 熵缺口的非消失常数下界]\label{cor:xi-time-part9t-fourier-kl-shannon-gap-positive-floor}
在定理 \ref{thm:xi-time-part9t-fourier-tv-gap-positive-floor} 的记号下，按自然对数约定有
\[
D_{\mathrm{KL}}(p_m\|u)=\log M-H(p_m).
\]
并且
\[
\liminf_{m\to\infty}D_{\mathrm{KL}}(p_m\|u)
\ge
\frac{C_\varphi^2}{2},
\]
\[
\liminf_{m\to\infty}\bigl(\log M-H(p_m)\bigr)
\ge
\frac{C_\varphi^2}{2}.
\]
特别地，
\[
\limsup_{m\to\infty}\bigl(\log M-H(p_m)\bigr)
\ge
\frac{C_\varphi^2}{2}.
\]
\end{corollary}
\begin{proof}
由 Pinsker 不等式，
\[
D_{\mathrm{KL}}(p_m\|u)\ge 2\|p_m-u\|_{\mathrm{TV}}^2.
\]
再由定理 \ref{thm:xi-time-part9t-fourier-tv-gap-positive-floor}，
\[
\liminf_{m\to\infty}D_{\mathrm{KL}}(p_m\|u)
\ge
2\left(\frac{C_\varphi}{2}\right)^2
=
\frac{C_\varphi^2}{2}.
\]
另一方面，均匀分布 \(u\) 满足
\[
\log M-H(p_m)=D_{\mathrm{KL}}(p_m\|u),
\]
故同一下界立即转移到 Shannon 熵缺口。最后，\(\limsup\ge\liminf\)，遂得最后一式。证毕。
\end{proof}

\begin{theorem}[均匀层上尾大偏差给出的连续容量指数的 Laplace--LDP 公式]\label{thm:xi-time-part9t-capacity-laplace-ldp-tail-variational}
沿用定理 \ref{thm:xi-time-part57eb-uniform-logmultiplicity-ldp} 的记号。令
\[
u_m(x):=\frac{1}{|X_m|},
\qquad
Y_m:=\frac1m\log d_m(X),
\qquad
X\sim u_m.
\]
并记
\[
|X_m|=\exp(m\log\varphi+o(m)).
\]
设 \((Y_m)\) 在 \(u_m\) 下满足速度 \(m\) 的大偏差原理，其好速率函数为 \(I_u\)。再设对任意紧区间 \(I\subset[0,\log 2]\) 都有一致上尾指数渐近
\[
\mathbb P_{u_m}(Y_m\ge y)
=
\exp\!\bigl(-mI_u^\uparrow(y)+o(m)\bigr)
\qquad (y\in I),
\]
其中
\[
I_u^\uparrow(y):=\inf_{z\ge y}I_u(z)
\qquad (0\le y\le \log 2),
\]
并约定 \(I_u^\uparrow(y):=+\infty\) 当 \(y>\log 2\)。
则对每个 \(\beta\ge 0\)，连续容量
\[
\mathcal C_m^{\mathrm{cont}}(T):=\sum_{x\in X_m}\min\bigl(d_m(x),T\bigr)
\]
满足
\[
\lim_{m\to\infty}\frac1m\log \mathcal C_m^{\mathrm{cont}}(e^{\beta m})
=
\sup_{0\le y\le \beta}
\Bigl\{
y+\log\varphi-I_u^\uparrow(y)
\Bigr\}.
\]
\end{theorem}
\begin{proof}
定义尾计数函数
\[
N_m(T):=\#\{x\in X_m:\ d_m(x)\ge T\}.
\]
则对任意 \(y\ge 0\)，由 \(X\sim u_m\) 可得
\[
N_m(e^{ym})
=
|X_m|\,\mathbb P_{u_m}(Y_m\ge y).
\]
又因 \(|X_m|=\exp(m\log\varphi+o(m))\)，故在每个紧区间 \(I\subset[0,\log 2]\) 上都成立
\[
N_m(e^{ym})
=
\exp\!\Bigl(
m\bigl(\log\varphi-I_u^\uparrow(y)\bigr)+o(m)
\Bigr)
\qquad (y\in I),
\]
其中 \(o(m)\) 对 \(y\in I\) 一致。

另一方面，由定理 \ref{thm:xi-capacity-tail-histogram-moment-complete-statistic}，
\[
\mathcal C_m^{\mathrm{cont}}(e^{\beta m})
=
\int_0^{e^{\beta m}}N_m(t)\,\dd t.
\]
作变量代换 \(t=e^{ym}\)，得到
\[
\mathcal C_m^{\mathrm{cont}}(e^{\beta m})
=
m\int_{-\infty}^{\beta}e^{ym}N_m(e^{ym})\,\dd y.
\]
当 \(y<0\) 时有 \(e^{ym}<1\)，而所有纤维多重度都至少为 \(1\)，故
\[
N_m(e^{ym})=|X_m|=\exp(m\log\varphi+o(m)).
\]
于是该区间上的指数至多为 \(y+\log\varphi\le \log\varphi\)，不会优于 \(y=0\) 处的候选值。

再由 \(d_m(x)\le 2^m\) 可知 \(N_m(e^{ym})=0\) 对一切 \(y>\log 2\) 最终成立，故积分主项只可能来自紧区间 \([0,\min\{\beta,\log 2\}]\)。在该区间上，积分核具有统一指数形
\[
e^{ym}N_m(e^{ym})
=
\exp\!\Bigl(
m\bigl(y+\log\varphi-I_u^\uparrow(y)\bigr)+o(m)
\Bigr).
\]
对该有限区间应用 Laplace 原理，即得
\[
\lim_{m\to\infty}\frac1m\log \mathcal C_m^{\mathrm{cont}}(e^{\beta m})
=
\sup_{0\le y\le \beta}
\Bigl\{
y+\log\varphi-I_u^\uparrow(y)
\Bigr\},
\]
其中 \(y>\log 2\) 的部分按上文自动不给出更大贡献。证毕。
\end{proof}

\begin{proposition}[最大纤维能级间隙强制容量曲线的精确仿射段]\label{prop:xi-time-part9t-top-gap-affine-capacity-segment}
\leanverified{paper\_xi\_time\_part9t\_top\_gap\_affine\_capacity\_segment}
定义
\[
M_m:=\max_{x\in X_m}d_m(x),
\qquad
K_m:=\#\{x\in X_m:\ d_m(x)=M_m\},
\]
并令
\[
M_m^{(2)}:=
\max\{d_m(x):\ d_m(x)<M_m\},
\]
若所有纤维都达到同一最大值，则约定 \(M_m^{(2)}:=0\)。再定义连续容量
\[
\mathcal C_m^{\mathrm{cont}}(T):=
\sum_{x\in X_m}\min\bigl(d_m(x),T\bigr)
\qquad (T\ge 0).
\]
若某个 \(T\) 满足
\[
M_m^{(2)}<T<M_m,
\]
则有精确恒等式
\[
\mathcal C_m^{\mathrm{cont}}(T)
=
2^m-K_m(M_m-T).
\]
因此 \(\mathcal C_m^{\mathrm{cont}}\) 在区间 \((M_m^{(2)},M_m)\) 上为仿射函数，且
\[
\frac{\dd}{\dd T}\mathcal C_m^{\mathrm{cont}}(T)=K_m,
\qquad
M_m=\frac{2^m-\mathcal C_m^{\mathrm{cont}}(T)}{K_m}+T.
\]
特别地，若整数预算 \(B\ge 0\) 满足
\[
M_m^{(2)}<2^B<M_m,
\]
则离散预言机容量
\[
C_m(B):=\sum_{x\in X_m}\min\bigl(d_m(x),2^B\bigr)
\]
满足
\[
C_m(B)=2^m-K_m(M_m-2^B).
\]
\end{proposition}
\begin{proof}
当 \(M_m^{(2)}<T<M_m\) 时，对所有满足 \(d_m(x)<M_m\) 的纤维都有
\[
d_m(x)\le M_m^{(2)}<T,
\]
故这些纤维在 \(\mathcal C_m^{\mathrm{cont}}(T)\) 中不被截断；而所有最大纤维都贡献 \(T\)。于是
\[
\mathcal C_m^{\mathrm{cont}}(T)
=
\sum_{d_m(x)<M_m}d_m(x)+K_mT.
\]
再由总质量恒等式
\[
\sum_{x\in X_m}d_m(x)=2^m,
\]
得到
\[
\sum_{d_m(x)<M_m}d_m(x)=2^m-K_mM_m.
\]
代回即得
\[
\mathcal C_m^{\mathrm{cont}}(T)=2^m-K_m(M_m-T).
\]
因此 \(\mathcal C_m^{\mathrm{cont}}\) 在该区间上为斜率 \(K_m\) 的仿射函数，导数公式与 \(M_m\) 的重建公式随即成立。令 \(T=2^B\) 即得到离散预算版本。证毕。
\end{proof}

\begin{theorem}[bin-fold 正实压力轴的全阶线性化]\label{thm:xi-time-part9t-binfold-positive-axis-linear-pressure}
\leanverified{paper\_xi\_time\_part9t\_binfold\_positive\_axis\_linear\_pressure}
对任意实数 \(a>0\)，定义
\[
S_a^{\mathrm{bin}}(m):=\sum_{w\in X_m}\bigl(d_m^{\mathrm{bin}}(w)\bigr)^a,
\qquad
P_a^{\mathrm{bin}}:=\lim_{m\to\infty}\frac1m\log S_a^{\mathrm{bin}}(m).
\]
则当 \(m\to\infty\) 时有
\[
S_a^{\mathrm{bin}}(m)
=
\Bigl(\frac{2}{\varphi}\Bigr)^{am}
\Bigl(F_{m+1}+\varphi^{-a}F_m\Bigr)
\Bigl(1+O\!\bigl(a(\varphi/2)^m\bigr)\Bigr),
\]
从而极限存在并满足闭式
\[
P_a^{\mathrm{bin}}
=
\log\varphi+a\log\!\Bigl(\frac{2}{\varphi}\Bigr).
\]
因此，\(a\mapsto P_a^{\mathrm{bin}}\) 在整个正实半轴 \((0,\infty)\) 上都是仿射函数；任何正阶“冻结阈值”都只能退化到原点。
\end{theorem}
\begin{proof}
命题 \ref{prop:fold-bin-power-sum-asymptotic} 已经对每个固定 \(a\ge 0\) 给出
\[
S_a^{\mathrm{bin}}(m)
=
\Bigl(\frac{2}{\varphi}\Bigr)^{am}
\Bigl(F_{m+1}+\varphi^{-a}F_m\Bigr)
\Bigl(1+O\!\bigl(a(\varphi/2)^m\bigr)\Bigr).
\]
对固定 \(a>0\)，误差即为 \(O((\varphi/2)^m)\)。再由 Binet 渐近，
\[
F_{m+1}+\varphi^{-a}F_m=\Theta(\varphi^m),
\]
故
\[
\frac1m\log S_a^{\mathrm{bin}}(m)
=
a\log\!\Bigl(\frac{2}{\varphi}\Bigr)+\log\varphi+o(1).
\]
令 \(m\to\infty\) 即得
\[
P_a^{\mathrm{bin}}
=
\log\varphi+a\log\!\Bigl(\frac{2}{\varphi}\Bigr).
\]
于是结论 \ref{con:xi-time-part9m-binfold-linear-pressure-no-positive-threshold} 中的正整数线性化可直接延拓到整个正实半轴。证毕。
\end{proof}

\begin{corollary}[bin-fold 临界 dyadic 窗的普适双折点标度函数]\label{cor:xi-time-part9t-binfold-critical-window-two-break-scaling}
沿用定理 \ref{thm:xi-foldbin-dyadic-capacity-critical-window-limit} 的 dyadic 预言机容量记号，并记
\[
\lambda:=\frac{2}{\varphi}.
\]
若整数预算序列 \((B_m)_{m\ge 1}\) 满足
\[
\frac{2^{B_m}}{\lambda^m}\longrightarrow 2^b
\qquad (b\in\RR),
\]
则有极限
\[
\lim_{m\to\infty}\frac{C_m^{\mathrm{bin}}(B_m)}{2^m}
=
f(b),
\]
其中
\[
f(b)
:=
\frac{\varphi\min\{1,2^b\}+\min\{\varphi^{-1},2^b\}}{\sqrt5}.
\]
等价地，
\[
f(b)=
\begin{cases}
\dfrac{\varphi^2}{\sqrt5}\,2^b,
& b\le -\log_2\varphi,\\[8pt]
\dfrac{\varphi\,2^b+\varphi^{-1}}{\sqrt5},
& -\log_2\varphi\le b\le 0,\\[10pt]
1,
& b\ge 0.
\end{cases}
\]
因此，bin-fold 的临界资源窗恰有两处折点，对应于
\[
b=-\log_2\varphi,
\qquad
b=0,
\]
两阈值相距 \(\log_2\varphi\)；临界窗宽度保持为 \(O(1)\)，不随 \(m\) 增长。
\end{corollary}
\begin{proof}
由定理 \ref{thm:xi-foldbin-dyadic-capacity-critical-window-limit}，在
\[
\zeta:=\lim_{m\to\infty}\frac{2^{B_m}}{\lambda^m}=2^b
\]
存在时有
\[
\frac{C_m^{\mathrm{bin}}(B_m)}{2^m}
=
\mathrm{Succ}_m^{\mathrm{bin}}(B_m)
\longrightarrow
\begin{cases}
\dfrac{\varphi^2}{\sqrt5}\,\zeta,
& 0<\zeta\le \varphi^{-1},\\[8pt]
\dfrac{\varphi}{\sqrt5}\,\zeta+\dfrac{1}{\varphi\sqrt5},
& \varphi^{-1}\le \zeta\le 1,\\[10pt]
1,
& \zeta\ge 1.
\end{cases}
\]
将 \(\zeta=2^b\) 代入，即得所示分段公式。再把三段写成最小值形式，便得到
\[
f(b)=\frac{\varphi\min\{1,2^b\}+\min\{\varphi^{-1},2^b\}}{\sqrt5}.
\]
折点位置分别对应于
\[
2^b=\varphi^{-1},
\qquad
2^b=1,
\]
亦即 \(b=-\log_2\varphi\) 与 \(b=0\)。由于折点位置与 \(m\) 无关，临界窗宽度保持为 \(O(1)\)。证毕。
\end{proof}

\endinput

\paragraph{冻结相 escort 刚性、偶模暗模与纯碰撞临界窗口}
在冻结相的基态质量指数集中、折叠多重度谱的乘积分解以及纯碰撞误差指数的平方根律已经确立之后，还可进一步抽取下列四条彼此衔接的后果。它们分别把单温 escort 集中提升为整族成对塌缩，把对数多重度的自平均提升为有限尺度矩母函数刚性，把 \((-1)\)-频点的暗模现象压缩为偶模算术判据，并把纯碰撞次谱与 Perron 主谱之间的竞争时间尺度识别为 \(\sqrt u\) 级临界窗口。

\begin{theorem}[冻结相 escort 族的成对指数塌缩]\label{thm:xi-time-part9u-frozen-escort-pairwise-tv-collapse}
沿用定理 \ref{thm:xi-fixed-freezing-escort-renyi-spectrum-collapse} 的记号，并记
\[
U_m^{\max}:=\mathrm{Unif}(A_m^{\max}).
\]
若 \(a,b>a_{\mathrm c}\)，则有
\[
\left\|\pi_{m,a}-\pi_{m,b}\right\|_{\mathrm{TV}}
\le
\exp\!\bigl(-m\Delta(a)+o(m)\bigr)
+
\exp\!\bigl(-m\Delta(b)+o(m)\bigr).
\]
特别地，
\[
\left\|\pi_{m,a}-\pi_{m,b}\right\|_{\mathrm{TV}}
=
O\!\Bigl(
\exp\!\bigl(-m\min\{\Delta(a),\Delta(b)\}+o(m)\bigr)
\Bigr).
\]
从而冻结相内任意两条 escort 温度支都以指数速度共同塌缩到同一基态均匀分布 \(U_m^{\max}\)。
\end{theorem}
\begin{proof}
由定理 \ref{thm:xi-fixed-freezing-escort-renyi-spectrum-collapse} 第四条，
\[
\left\|\pi_{m,a}-U_m^{\max}\right\|_{\mathrm{TV}}
\le
\exp\!\bigl(-m\Delta(a)+o(m)\bigr),
\]
\[
\left\|\pi_{m,b}-U_m^{\max}\right\|_{\mathrm{TV}}
\le
\exp\!\bigl(-m\Delta(b)+o(m)\bigr).
\]
由三角不等式，
\[
\left\|\pi_{m,a}-\pi_{m,b}\right\|_{\mathrm{TV}}
\le
\left\|\pi_{m,a}-U_m^{\max}\right\|_{\mathrm{TV}}
+
\left\|U_m^{\max}-\pi_{m,b}\right\|_{\mathrm{TV}},
\]
代入上述两式即得第一条估计；第二条只是把两项指数界压缩到较小指数率的写法。证毕。
\end{proof}

\begin{theorem}[冻结相中对数多重度涨落的矩母函数刚性]\label{thm:xi-time-part9u-frozen-escort-logmultiplicity-mgf}
\leanverified{paper\_xi\_time\_part9u\_frozen\_escort\_logmultiplicity\_mgf}
在严格基态间隙
\[
\alpha_*^{(2)}<\alpha_*
\]
下，沿用定理 \ref{thm:xi-fixed-freezing-escort-renyi-spectrum-collapse} 的记号，并定义
\[
Y_m(x):=\frac1m\log d_m(x)-\frac1m\log M_m
\qquad (x\in X_m).
\]
则 \(Y_m\le 0\)。若 \(a>a_{\mathrm c}\) 且 \(t\in\RR\) 满足 \(a+t>a_{\mathrm c}\)，则
\[
\mathbb E_{\pi_{m,a}}\!\left[e^{tmY_m}\right]
=
\frac{S_{a+t}(m)}{M_m^tS_a(m)}
=
1+O\!\Bigl(
\exp\!\bigl(-m\delta(a,t)+o(m)\bigr)
\Bigr),
\]
其中
\[
\delta(a,t):=\min\{\Delta(a),\Delta(a+t)\}>0.
\]
进一步，对每个整数 \(r\ge 1\) 都存在常数 \(C_r(a)>0\)，使得
\[
\mathbb E_{\pi_{m,a}}\!\left[|Y_m|^r\right]
\le
C_r(a)\,
\exp\!\bigl(-m\Delta(a)+o(m)\bigr).
\]
因此，若 \(X\sim \pi_{m,a}\)，则 \(\frac1m\log d_m(X)\) 在每个固定 \(L^r\) 中都以指数速度收敛到 \(\alpha_*\)。
\end{theorem}
\begin{proof}
若次大层不存在，则约定 \(M_m^{(2)}:=1\)。由定义直接计算得
\[
\mathbb E_{\pi_{m,a}}\!\left[e^{tmY_m}\right]
=
\sum_{x\in X_m}\frac{d_m(x)^a}{S_a(m)}
\frac{d_m(x)^t}{M_m^t}
=
\frac{1}{M_m^tS_a(m)}\sum_{x\in X_m}d_m(x)^{a+t}
=
\frac{S_{a+t}(m)}{M_m^tS_a(m)}.
\]
这给出等式部分。

对任意 \(u>a_{\mathrm c}\)，写
\[
S_u(m)=K_mM_m^u+R_u(m),
\qquad
0\le R_u(m)\le |X_m|\,(M_m^{(2)})^u.
\]
于是
\[
\frac{R_u(m)}{K_mM_m^u}
\le
\frac{|X_m|\,(M_m^{(2)})^u}{K_mM_m^u}.
\]
利用
\[
\log|X_m|=m\log\varphi+o(m),
\qquad
\log K_m=g_*m+o(m),
\]
\[
\log M_m=\alpha_*m+o(m),
\qquad
\log M_m^{(2)}\le \alpha_*^{(2)}m+o(m),
\]
便得到
\[
\frac{R_u(m)}{K_mM_m^u}
\le
\exp\!\bigl(-m\Delta(u)+o(m)\bigr),
\]
其中
\[
\Delta(u):=u(\alpha_*-\alpha_*^{(2)})-(\log\varphi-g_*)>0.
\]
故
\[
S_u(m)=K_mM_m^u
\Bigl(
1+O\!\bigl(e^{-m\Delta(u)+o(m)}\bigr)
\Bigr).
\]
分别取 \(u=a+t\) 与 \(u=a\)，再作比值，即得
\[
\frac{S_{a+t}(m)}{M_m^tS_a(m)}
=
\frac{
1+O\!\bigl(e^{-m\Delta(a+t)+o(m)}\bigr)
}{
1+O\!\bigl(e^{-m\Delta(a)+o(m)}\bigr)
}
=
1+O\!\Bigl(
e^{-m\delta(a,t)+o(m)}
\Bigr).
\]

再证矩估计。由于 \(Y_m=0\) 在 \(A_m^{\max}\) 上恒成立，故
\[
\mathbb E_{\pi_{m,a}}\!\left[|Y_m|^r\right]
=
\sum_{x\notin A_m^{\max}}\pi_{m,a}(x)\,|Y_m(x)|^r.
\]
另一方面，对全部 \(x\in X_m\) 都有 \(1\le d_m(x)\le M_m\)，于是
\[
0\le |Y_m(x)|
\le
\frac1m\log M_m
=
\alpha_*+o(1).
\]
对充分大的 \(m\) 可取
\[
|Y_m(x)|\le \alpha_*+1.
\]
因此
\[
\mathbb E_{\pi_{m,a}}\!\left[|Y_m|^r\right]
\le
(\alpha_*+1)^r\,\pi_{m,a}(X_m\setminus A_m^{\max}).
\]
再由定理 \ref{thm:xi-fixed-freezing-escort-renyi-spectrum-collapse} 第四条，
\[
\pi_{m,a}(X_m\setminus A_m^{\max})
\le
\exp\!\bigl(-m\Delta(a)+o(m)\bigr),
\]
遂得所示界。最后，由
\[
\frac1m\log d_m(X)=\frac1m\log M_m+Y_m(X),
\qquad
\frac1m\log M_m\to \alpha_*,
\]
即可推出 \(\frac1m\log d_m(X)\) 在每个固定 \(L^r\) 中都以指数速度收敛到 \(\alpha_*\)。证毕。
\end{proof}

\begin{theorem}[偶模 \texorpdfstring{$(-1)$}{(-1)} 频点的精确暗模判据]\label{thm:xi-time-part9u-minusone-darkmode-even-modulus}
\leanverified{paper\_xi\_time\_part9u\_minusone\_darkmode\_even\_modulus}
记
\[
M:=F_{m+2},
\qquad
D_m^{\sharp}(x):=\sum_{r=0}^{M-1}d_m(r)x^r\in \CC[x].
\]
则下列三条命题等价：
\begin{enumerate}
  \item \(3\mid (m+2)\)；
  \item \(M\) 为偶数；
  \item 频率 \(M/2\in \ZZ/(M\ZZ)\) 有意义且
  \[
  \widehat d_m(M/2)=0.
  \]
\end{enumerate}
在这些等价条件下，还成立
\[
D_m^{\sharp}(-1)=0,
\]
从而 \(x+1\) 整除标准代表 \(D_m^{\sharp}(x)\)。
\end{theorem}
\begin{proof}
定理 \ref{thm:xi-time-part60ab2-exact-dark-mode-arithmetic-criterion} 已经给出：(2) 与 (3) 等价，并且二者都等价于 \(3\mid (m+2)\)。因此前三条彼此等价。

在这些等价条件下，\(M\) 为偶数，故取值 \(x=-1\) 在商环 \(\CC[x]/(x^M-1)\) 上良定。再由定理 \ref{thm:xi-fold-multiplicity-strict-parity-balance}，
\[
D_m^{\sharp}(-1)=\widehat d_m(M/2)=0.
\]
因此普通多项式 \(D_m^{\sharp}(x)\) 在 \(x=-1\) 处有根，等价于 \(x+1\) 整除 \(D_m^{\sharp}(x)\)。证毕。
\end{proof}

\noindent
需要强调的是：当 \(M\) 为奇数时，取值 \(x=-1\) 并不在商环 \(\CC[x]/(x^M-1)\) 上定义为群特征，因此普通代表在 \(x=-1\) 处的消失不再由 Fourier 暗模自动控制。

\begin{theorem}[纯碰撞立方的 \texorpdfstring{$\sqrt u$}{sqrt(u)} 临界竞争窗口]\label{thm:xi-time-part9u-pure-collision-sqrtu-critical-window}
沿用推论 \ref{cor:real-input-40-collision-error-phase} 的记号。对 \(u>1\) 定义
\[
R(u):=\frac{\Lambda(u)}{\lambda_+(u)}\in(0,1).
\]
则当 \(u\to+\infty\) 时有
\[
\log R(u)
=
-\frac{1}{\sqrt u}
+O\!\left(\frac1u\right).
\]
进一步，对任意固定 \(\tau>0\)，若取整数序列
\[
n(u):=\lfloor \tau\sqrt u\rfloor,
\]
则有极限
\[
R(u)^{\,n(u)}\longrightarrow e^{-\tau}.
\]
换言之，纯碰撞立方中次谱与 Perron 主谱之间的有效竞争并不发生在对数时间尺度，而是恰好压缩为 \(n\asymp \sqrt u\) 的临界窗口。
\end{theorem}
\begin{proof}
令 \(t:=\sqrt u\)。由推论 \ref{cor:real-input-40-collision-error-phase}，
\[
\lambda_+(u)=t+\frac12+\frac{9}{8t}+O(t^{-2}),
\qquad
\Lambda(u)=t-\frac12+\frac{9}{8t}+O(t^{-2}).
\]
因此
\[
\lambda_+(u)
=
t\left(
1+\frac{1}{2t}+\frac{9}{8t^2}+O(t^{-3})
\right),
\]
\[
\Lambda(u)
=
t\left(
1-\frac{1}{2t}+\frac{9}{8t^2}+O(t^{-3})
\right).
\]
两边取对数并相减，得到
\[
\log R(u)
=
\log\!\left(
1-\frac{1}{2t}+\frac{9}{8t^2}+O(t^{-3})
\right)
-
\log\!\left(
1+\frac{1}{2t}+\frac{9}{8t^2}+O(t^{-3})
\right)
=
-\frac{1}{t}+O(t^{-2}).
\]
亦即
\[
\log R(u)=-\frac{1}{\sqrt u}+O\!\left(\frac1u\right).
\]

再由
\[
n(u)=\tau t+O(1),
\]
可得
\[
n(u)\log R(u)
=
\bigl(\tau t+O(1)\bigr)
\left(-\frac{1}{t}+O(t^{-2})\right)
=
-\tau+O(t^{-1})
\longrightarrow
-\tau.
\]
指数化即得
\[
R(u)^{\,n(u)}\to e^{-\tau}.
\]
证毕。
\end{proof}

\noindent
由此，冻结相 escort、fold 多重度频谱与纯碰撞慢模在同一条终端链上闭合：前者在统计实验层面塌缩为共同基态，频谱侧把 \((-1)\)-暗模锁定为偶模共振，而纯碰撞侧则把次谱竞争压缩为一条可连续审计的 \(\sqrt u\)-时间窗口。

\endinput

\paragraph{临界切片刚性、共尾 prime-twist 慢模与双不可压缩核}\label{subsubsec:xi-time-part9v-critical-slice-cofinal-prime-twist-rigidity}
在 Mellin--Plancherel 幺正切片的唯一等距性、离切片信息的正交记录轴外置纪律、有限 prime-stage solenoid 核的函子恢复，以及真实输入 \(40\) 状态核在单位根邻域的四阶慢模展开都已确立之后，还可进一步抽取出下列八条彼此衔接的结论。它们分别把临界线的 Hilbert 刚性、离切片补偿的算术代价，以及共尾 prime-twist 分支的不可有限化承载统一为同一条证明复杂度障碍链。

\begin{theorem}[离开临界切片的有限余维不可修复性]\label{thm:xi-time-part9v-offcritical-finite-codim-nonrepairability}
\leanverified{paper\_xi\_time\_part9v\_offcritical\_finite\_codim\_nonrepairability}
沿用定理 \ref{thm:mellin-unitary-slice-unique} 的记号。对任意
\[
\sigma\in\RR\setminus\left\{\frac12\right\},
\]
不存在有限余维闭子空间
\[
H\subset L^2_\times
\]
使得限制读出
\[
\mathcal M_\sigma|_H:H\longrightarrow L^2\!\left(\RR,\frac{\dd t}{2\pi}\right)
\]
为等距算子。
\end{theorem}
\begin{proof}
定义乘法权算子
\[
(W_\sigma f)(x):=x^{\sigma-\frac12}f(x).
\]
由定理 \ref{thm:mellin-unitary-slice-unique} 的恒等式 \eqref{eq:mellin-sigma-weight-twist},
\[
\mathcal M_\sigma=\mathcal M_{1/2}\circ W_\sigma,
\]
且 \(\mathcal M_{1/2}\) 幺正。故 \(\mathcal M_\sigma|_H\) 等距当且仅当 \(W_\sigma|_H\) 等距。

先设 \(\sigma>\tfrac12\)。取
\[
E:=[2,4],
\qquad
V:=L^2(E,\dd^\times x).
\]
则 \(V\) 是 \(L^2_\times\) 的无限维闭子空间。由于 \(H\) 余维有限，商空间 \(L^2_\times/H\) 有限维；因此商映射 \(V\to L^2_\times/H\) 的像亦有限维，从而
\[
V\cap H\neq\{0\}.
\]
任取 \(0\neq f\in V\cap H\)。在 \(E\) 上有
\[
x^{2\sigma-1}\ge 2^{2\sigma-1}>1,
\]
于是
\[
\|W_\sigma f\|_{L^2_\times}^2
=
\int_E x^{2\sigma-1}|f(x)|^2\,\dd^\times x
\ge
2^{2\sigma-1}\|f\|_{L^2_\times}^2
>
\|f\|_{L^2_\times}^2,
\]
与等距性矛盾。

若 \(\sigma<\tfrac12\)，则取
\[
E:=\left[\frac14,\frac12\right].
\]
此时 \(2\sigma-1<0\)，故对全部 \(x\in E\) 有
\[
x^{2\sigma-1}\ge \left(\frac12\right)^{2\sigma-1}>1.
\]
完全同理可得矛盾。故所述 \(H\) 不存在。证毕。
\end{proof}

\begin{corollary}[离临界切片不存在有限秩 strictification]\label{cor:xi-time-part9v-offcritical-no-finite-rank-strictification}
沿用定理 \ref{thm:xi-time-part9v-offcritical-finite-codim-nonrepairability} 的记号。若
\[
\sigma\neq\frac12,
\]
则不存在有限秩算子
\[
K:L^2_\times\longrightarrow L^2\!\left(\RR,\frac{\dd t}{2\pi}\right)
\]
使得
\[
\mathcal M_\sigma+K
\]
成为等距算子。
\end{corollary}
\begin{proof}
若 \(K\) 有限秩，则 \(\ker K\subset L^2_\times\) 为有限余维闭子空间。若 \(\mathcal M_\sigma+K\) 等距，则对任意 \(f\in\ker K\) 有
\[
\|f\|_{L^2_\times}
=
\|(\mathcal M_\sigma+K)f\|
=
\|\mathcal M_\sigma f\|,
\]
即 \(\mathcal M_\sigma\) 在 \(\ker K\) 上等距。这与定理 \ref{thm:xi-time-part9v-offcritical-finite-codim-nonrepairability} 矛盾。证毕。
\end{proof}

\begin{theorem}[Hurwitz--Carath\'eodory 证书链的临界切片刚性]\label{thm:xi-time-part9v-hurwitz-carath-critical-slice-rigidity}
考虑一条将 slice 读出直接送入圆盘完成化、Hurwitz 零点禁入以及 Toeplitz--Carath\'eodory 正性终端的证明管线。若该管线始终固定在某一切片
\[
\Re(s)=\sigma,
\]
并且协议拒绝引入额外正交记录轴，只允许在该切片内部施加有限秩修正或转入有限余维载体，则必有
\[
\sigma=\frac12.
\]
\end{theorem}
\begin{proof}
定理 \ref{thm:xi-slice-horizon-carath-toeplitz} 把 classical \(\Xi\)-切片的终端证书写成单切片上的 Carath\'eodory 正性与 Toeplitz--PSD 层级。若 \(\sigma\neq\tfrac12\)，则由定理 \ref{thm:xi-time-part9v-offcritical-finite-codim-nonrepairability} 与推论 \ref{cor:xi-time-part9v-offcritical-no-finite-rank-strictification}，任何有限余维 restriction 或有限秩 strictification 都不能把 \(\mathcal M_\sigma\) 修补为等距载体。

另一方面，定理 \ref{thm:xi-extension-preserving-initial-object} 与命题 \ref{prop:xi-offcritical-falsifiable-restatement} 表明：若仍欲在闭环内保留离切片信息，则它只能经由额外正交记录轴外置；若拒绝该外置，则由定理 \ref{thm:xi-time-part60eb-pure-phase-offslice-bifurcation}，相关负载只能退化为 \(\mathsf{NULL}\) 或失败见证，而不能作为切片内部的正性证书继续传递到 Toeplitz--Carath\'eodory 终端。

因此，在“不增设正交记录轴且只允许有限缺陷内部修补”的前提下，此类单切片直送管线只能驻留于唯一等距切片 \(\Re(s)=\tfrac12\)。证毕。
\end{proof}

\begin{theorem}[有限 prime-stage 对近单位根慢模分支的结构失明]\label{thm:xi-time-part9v-finite-primestage-nearunitroot-blindness}
沿用定理 \ref{thm:xi-time-part56-solenoid-terminal-prime-register-recovery} 与推论 \ref{cor:real-input-40-collision-twist-fourth} 的记号。设 \(S\subset\mathbb P\) 为有限素数集，并记
\[
G_S:=\ZZ[S^{-1}],
\qquad
\Sigma_S:=\widehat{G_S},
\qquad
K_S^{\mathrm{reg}}:=\ker(\pi_S)\cong\prod_{q\in S}\ZZ_q,
\]
其中 \(\pi_S:\Sigma_S\twoheadrightarrow\TT\) 为标准投影。再记
\[
\lambda:=\rho(1)=\varphi^2,
\qquad
\rho(t):=\rho(e^{it}).
\]
则对任意 \(\varepsilon>0\)，都存在奇素数
\[
p\notin S
\]
使得
\[
\frac{\log\rho(2\pi/p)}{\log\lambda}>1-\varepsilon,
\]
但同时
\[
\bigl(K_S^{\mathrm{reg}}\bigr)_p=0.
\]
因而该 \(p\)-twist slow mode 在固定有限 prime-stage 的账本核上不具结构性地址；特别地，任何固定有限 prime-stage 都不能承载整条共尾的近单位根慢模分支。
\end{theorem}
\begin{proof}
由定理 \ref{thm:xi-time-part56-solenoid-terminal-prime-register-recovery}，
\[
S=\{\,q\in\mathbb P:\ (K_S^{\mathrm{reg}})_q\neq 0\,\}.
\]
因此一旦 \(p\notin S\)，便有 \((K_S^{\mathrm{reg}})_p=0\)，即该 prime-stage 的核不含任何 \(p\)-主账本方向。

另一方面，由推论 \ref{cor:real-input-40-collision-twist-fourth}，当 \(t\to 0\) 时
\[
\log\frac{\rho(t)}{\lambda}
=
-\frac{6\sqrt5-5}{250}\,t^2
+\frac{7+24\sqrt5}{75000}\,t^4
+O(t^6).
\]
代入
\[
t_p:=\frac{2\pi}{p},
\]
得
\[
\log\frac{\rho(t_p)}{\lambda}
=
-\frac{6\sqrt5-5}{250}\left(\frac{2\pi}{p}\right)^2
+O(p^{-4}).
\]
又因 \(\log\lambda>0\)，故
\[
\frac{\log\rho(t_p)}{\log\lambda}
=
1-\frac{6\sqrt5-5}{250\,\log\lambda}\left(\frac{2\pi}{p}\right)^2
+O(p^{-4})
\longrightarrow 1
\qquad(p\to\infty,\ p\ \text{奇素数}).
\]
于是可取充分大的奇素数 \(p\notin S\)，使
\[
\frac{\log\rho(2\pi/p)}{\log\lambda}>1-\varepsilon.
\]
然而此时 \((K_S^{\mathrm{reg}})_p=0\)，故对应的 \(p\)-twist slow mode 不在该有限 prime-stage 的核上被结构性寻址。命题成立。证毕。
\end{proof}

\begin{theorem}[可回放离切片补偿协议的 Gödel 位长下界]\label{thm:xi-time-part9v-replayable-offslice-godel-bitlength-lower-bound}
设 \(T\in\NN\)、\(|\Sigma|\ge 2\)，并把
\[
s=(s_1,\dots,s_T)\in\Sigma^T
\]
视为长度 \(T\) 的离切片 defect word。若存在单射
\[
E_T:\Sigma^T\hookrightarrow \mathbb N_{\ge 1}
\]
以及两两互异的素数 \(p_1,\dots,p_T\)，使得对每个 \(t\) 而言，坐标 \(s_t\) 都能由估值
\[
v_{p_t}\!\bigl(E_T(s)\bigr)
\]
独立恢复，则存在某个输入 \(s^\star\in\Sigma^T\) 与常数 \(c>0\) 使
\[
\log E_T(s^\star)\ge c\,T\log T.
\]
因此，一切要求“可回放且逐坐标 prime-addressable”的离切片补偿协议，在最坏情形下都至少具有 \(T\log T\) 级别的算术位长。
\end{theorem}
\begin{proof}
上述假设恰把 defect word 的 replay 语义编译为逐坐标 prime-addressable 的 Gödel 外置。故可直接应用定理 \ref{thm:xi-addressable-godel-worstcase-chebyshev-sharp-lower-bound}：存在某个 \(s^\star\in\Sigma^T\) 满足
\[
\log E_T(s^\star)\ge \vartheta(\mathfrak p_T),
\]
其中 \(\mathfrak p_T\) 为第 \(T\) 个素数，\(\vartheta\) 为 Chebyshev 函数。再由该定理中的精化渐近
\[
\vartheta(\mathfrak p_T)=T\log T+T\log\log T+O(T),
\]
立即得到某个常数 \(c>0\) 使
\[
\log E_T(s^\star)\ge c\,T\log T.
\]
证毕。
\end{proof}

\begin{theorem}[普适离切片补偿账本的有限生成不可能性]\label{thm:xi-time-part9v-universal-offslice-ledger-not-finitely-generated}
设 \(L\) 为交换群。若存在单射幺半群同态
\[
\Phi:\mathbb N_{\ge 1}\longrightarrow \ZZ\oplus L,
\qquad
\Phi(ab)=\Phi(a)+\Phi(b),
\]
并把 \(\Phi(n)\) 解释为 prime-addressed defect label \(n\) 的可组合外置账本，则 \(L\) 不可能有限生成。
\end{theorem}
\begin{proof}
这正是定理 \ref{thm:xi-time-part64b-local-additivity-nonlocal-multiplicativity-polarization} 的有限秩障碍在 defect-ledger 语义下的直接应用。该定理已经证明：不存在任何有限生成交换群 \(L\) 使得
\[
\Phi:\mathbb N_{\ge 1}\to \ZZ\oplus L
\]
成为单射幺半群同态。故普适且可组合的离切片补偿账本必然超出有限生成范围。证毕。
\end{proof}

\begin{corollary}[离切片 finite-defect strictification 与 finite prime-stage ledger 的双截断 no-go]\label{cor:xi-time-part9v-dual-truncation-no-go}
考虑一条试图闭合定理 \ref{thm:xi-slice-horizon-carath-toeplitz} 终端的 classical RH 审计协议。若其补偿制度同时满足以下两项约束：
\[
\text{可见读出锚定在某个 }\Re(s)=\sigma\neq\frac12\text{，且切片内部只允许有限缺陷修补；}
\]
\[
\text{一切外置 prime-ledger 若出现，都只经由某个固定有限 prime-stage }\Sigma_S\text{ 因子化，}
\]
则该协议不可能完成终端证书闭合。
\end{corollary}
\begin{proof}
由定理 \ref{thm:xi-time-part9v-hurwitz-carath-critical-slice-rigidity}，离临界切片上的任何 finite-defect strictification 都不能在不引入额外正交记录轴的前提下把读出送达 Carath\'eodory 终端。因而若协议仍欲补偿离切片负载，唯一剩余路径只能转向外置账本。

然而一旦外置 prime-ledger 被限制为固定有限阶段 \(\Sigma_S\)，定理 \ref{thm:xi-time-part9v-finite-primestage-nearunitroot-blindness} 便给出：存在任意接近单位根的 slow modes，其对应 prime-twist 方向落在 \(S\)-外且不在 \(K_S^{\mathrm{reg}}\) 上可寻址。故有限阶段账本不能承载整条共尾慢模分支。

两类截断因此分别在切片内部与账本外部同时失效，终端证书闭合不可能成立。证毕。
\end{proof}

\begin{conjecture}[classical RH 证书的双不可压缩核]\label{conj:xi-time-part9v-rh-double-incompressible-core}
在本文的 solenoidal--Mellin 语义中，任何真正决定 classical RH 终端真值的证明系统都必须同时保留两类不可压缩对象：
\[
\mathfrak U:=\text{临界幺正切片载体 }(\Re(s)=\tfrac12),
\]
\[
\mathfrak P:=\text{共尾 prime-twist 慢模账本}.
\]
更精确地，若某证明系统对 \(\mathfrak U\) 施行有限缺陷 strictification，或把 \(\mathfrak P\) 截断到固定有限 prime-stage，则其全部证书最终都只能因子化到有限视界 shadow，而不能决定 RH 的终端真假。
\end{conjecture}

\noindent
以上八条结果表明，classical RH 路线中真正不可压缩的并不是抽象的无限维背景本身，而是两类彼此独立的核：一类是承载 Carath\'eodory--Toeplitz 终端的临界幺正切片，另一类是记录近单位根共尾慢模的 prime-twist 账本。前者拒绝任何有限缺陷 strictification，后者拒绝任何固定有限 prime-stage 截断；两者合并后，离切片修补与有限账本压缩便在同一结构层同时失效。与前述 faithful compression、window-\(6\) 尾谱与 tail--twist 型障碍相比，此处的 no-go 不再停留于单一表示层，而是把临界载体与共尾分支的双重不可压缩性并置为同一个终端约束。

\endinput

\paragraph{有限局部化加法账本、输出端点非高斯曲率与严格基态间隙的精确平顶}
在有限局部化 rank-one 轴的全局障碍、输出位势右端点的对数尖点、固定单位根模态的四阶展开，以及冻结相中的严格基态间隙与容量仿射段都已建立之后，还可继续抽取出下列几条派生结论。它们分别把有限局部化加法账本的结构失效、饱和前沿的非高斯曲率发散、基本单位根慢模对 Perron 指数的刚性贴近，以及顶端指数带中的尾计数精确平顶与容量精确肩部统一到同一条终端链上。

\leanverified{paper\_xi\_time\_part9w\_finite\_localized\_additive\_ledger\_obstruction}
\begin{theorem}[有限局部化圆因子与有限加法寄存器不足以忠实线性化完整乘法]\label{thm:xi-time-part9w-finite-localized-additive-ledger-obstruction}
设 \(S_1,\dots,S_r\subset\mathbb P\) 为有限素数集，并记
\[
G_{S_i}:=\ZZ[S_i^{-1}].
\]
则对任意 \(k\in\ZZ_{\ge 0}\)，都不存在交换幺半群单射
\[
\Phi:\NN_{\times}\hookrightarrow
\Bigl(\bigoplus_{i=1}^{r}G_{S_i}\Bigr)\oplus(\NN^k,+).
\]
\end{theorem}
\begin{proof}
设
\[
\iota:(\NN^k,+)\hookrightarrow(\ZZ^k,+)
\]
为自然嵌入。若存在所述 \(\Phi\)，则
\[
(\mathrm{id}\oplus\iota)\circ\Phi:
\NN_{\times}\hookrightarrow
\Bigl(\bigoplus_{i=1}^{r}G_{S_i}\Bigr)\oplus\ZZ^k
\]
仍为幺半群单射。把右端记为
\[
\mathbf S:=(S_1,\dots,S_r),
\qquad
\mathcal L_{k,\mathbf S}:=\ZZ^k\oplus G_{S_1}\oplus\cdots\oplus G_{S_r},
\]
便与定理 \ref{thm:xi-time-part62debb-finite-localization-axes-no-full-multiplicative-host} 直接矛盾。故所述 \(\Phi\) 不存在。证毕。
\end{proof}

\begin{corollary}[输出饱和前沿不存在有限曲率的 Gaussian 切触模型]\label{cor:xi-time-part9w-output-frontier-no-finite-gaussian-tangent}
沿用定理 \ref{thm:xi-output-potential-rate-right-endpoint-log-cusp} 与推论 \ref{cor:xi-output-microcanonical-endpoint-slope-curvature-diverge} 的记号。记
\[
I_{\mathrm{out}}(\alpha):=I(\alpha),
\qquad
h_{\mathrm{out}}(\alpha):=h(\alpha)=\log\lambda(1)-I(\alpha).
\]
则当 \(\alpha\uparrow\tfrac12\) 时，
\[
I_{\mathrm{out}}''(\alpha)=\frac{2}{1/2-\alpha}+O(1),
\qquad
h_{\mathrm{out}}''(\alpha)=-\frac{2}{1/2-\alpha}+O(1).
\]
特别地，\(I_{\mathrm{out}}\) 与 \(h_{\mathrm{out}}\) 在饱和端点都不具有有限二阶曲率；因此任何把输出饱和前沿视为局部 Gaussian 或有限阶 Edgeworth 二次切触的近似，都不可能在端点处成立。
\end{corollary}
\begin{proof}
推论 \ref{cor:xi-output-microcanonical-endpoint-slope-curvature-diverge} 已给出
\[
h_{\mathrm{out}}''(\alpha)=-\frac{2}{1/2-\alpha}+O(1).
\]
又由定义
\[
h_{\mathrm{out}}(\alpha)=\log\lambda(1)-I_{\mathrm{out}}(\alpha),
\]
两边二次求导便得
\[
I_{\mathrm{out}}''(\alpha)=-h_{\mathrm{out}}''(\alpha)
=\frac{2}{1/2-\alpha}+O(1).
\]
于是两式同时成立。证毕。
\end{proof}

\leanverified{paper\_xi\_time\_part9w\_basic\_root\_unity\_error\_exponent\_to\_one}
\begin{theorem}[基本单位根通道的误差指数以 \texorpdfstring{$m^{-2}$}{m^-2} 速度贴近 Perron 指数]\label{thm:xi-time-part9w-basic-root-unity-error-exponent-to-one}
沿用定理 \ref{thm:xi-output-potential-dirichlet-fixed-j-fourth-order} 的记号。记
\[
\vartheta_{m,1}:=\theta_{m,1}
=\frac{\log\rho_{m,1}}{\log\lambda},
\qquad
\lambda=\varphi^2.
\]
则
\[
\vartheta_{m,1}
=
1-A_{\mathrm{out}}m^{-2}
+B_{\mathrm{out}}m^{-4}
+O(m^{-6}),
\]
其中
\[
A_{\mathrm{out}}
:=
\frac{\pi^2(8955-3983\sqrt5)}{500\log(\varphi^2)}
>0,
\]
\[
B_{\mathrm{out}}
:=
\frac{\pi^4(1354428303-605718497\sqrt5)}{150000\log(\varphi^2)}.
\]
因而
\[
\vartheta_{m,1}\longrightarrow 1
\qquad(m\to\infty).
\]
换言之，基本单位根慢模的误差指数并不趋向平方根门槛 \(\tfrac12\)，而是以显式 \(m^{-2}\) 速率刚性贴近 Perron 主指数 \(1\)。
\end{theorem}
\begin{proof}
定理 \ref{thm:xi-output-potential-dirichlet-fixed-j-fourth-order} 在 \(j=1\) 时给出
\[
\vartheta_{m,1}
=
1-\frac{2\pi^2\kappa_2}{\log\lambda}\,m^{-2}
+\frac{2\pi^4\kappa_4}{3\log\lambda}\,m^{-4}
+O(m^{-6}),
\]
其中 \(\kappa_2=P''(0)\)、\(\kappa_4=P^{(4)}(0)\)。再由定理 \ref{thm:xi-real-input-40-output-edgeworth-nongaussian},
\[
\kappa_2=\frac{1791}{200}-\frac{3983\sqrt5}{1000}
=\frac{8955-3983\sqrt5}{1000},
\]
\[
\kappa_4=\frac{1354428303}{100000}-\frac{605718497\sqrt5}{100000}.
\]
直接代入即可得到所示常数 \(A_{\mathrm{out}},B_{\mathrm{out}}\) 的闭式。由于 \(A_{\mathrm{out}}>0\)，首阶修正为负且量级为 \(m^{-2}\)，故 \(\vartheta_{m,1}\to1\)。证毕。
\end{proof}

\begin{theorem}[严格基态间隙强制顶端指数带上的尾计数精确平顶]\label{thm:xi-time-part9w-strict-gap-tailcount-exact-plateau}
记
\[
M_m:=\max_{x\in X_m}d_m(x),
\qquad
K_m:=\#\{x\in X_m:\ d_m(x)=M_m\},
\]
并令
\[
M_m^{(2)}:=
\max\{d_m(x):\ d_m(x)<M_m\},
\]
若次大层不存在，则约定 \(M_m^{(2)}:=1\)。进一步假设存在常数
\[
\alpha_*^{(2)}<\alpha_*,
\qquad
g_*,
\]
使
\[
\log M_m=\alpha_*m+o(m),
\qquad
\log K_m=g_*m+o(m),
\qquad
\log M_m^{(2)}\le \alpha_*^{(2)}m+o(m).
\]
这正是定理 \ref{thm:xi-time-part9u-frozen-escort-logmultiplicity-mgf} 证明中所使用的 strict-gap 数据。若阈值序列 \(T_m>0\) 满足
\[
\alpha_*^{(2)}
<
\liminf_{m\to\infty}\frac{1}{m}\log T_m
\le
\limsup_{m\to\infty}\frac{1}{m}\log T_m
<
\alpha_*,
\]
则对一切充分大的 \(m\)，都有
\[
N_m(T_m):=\#\{x\in X_m:\ d_m(x)\ge T_m\}=K_m.
\]
从而
\[
\lim_{m\to\infty}\frac1m\log N_m(T_m)=g_*.
\]
\end{theorem}
\begin{proof}
取 \(\varepsilon>0\) 使
\[
\alpha_*^{(2)}+\varepsilon
<
\liminf_{m\to\infty}\frac1m\log T_m
\le
\limsup_{m\to\infty}\frac1m\log T_m
<
\alpha_*-\varepsilon.
\]
于是对充分大的 \(m\)，一方面由上述指数假设可得
\[
M_m^{(2)}\le \exp\bigl((\alpha_*^{(2)}+\varepsilon)m\bigr),
\qquad
M_m\ge \exp\bigl((\alpha_*-\varepsilon)m\bigr),
\]
另一方面又有
\[
\exp\bigl((\alpha_*^{(2)}+\varepsilon)m\bigr)
<T_m<
\exp\bigl((\alpha_*-\varepsilon)m\bigr).
\]
故最终
\[
M_m^{(2)}<T_m<M_m.
\]
于是所有非最大纤维都严格落在阈值之下，而全部最大纤维都严格超过阈值，因此
\[
N_m(T_m)=\#\{x\in X_m:\ d_m(x)=M_m\}=K_m
\]
对充分大的 \(m\) 成立。再用 \(\log K_m=g_*m+o(m)\) 即得极限式。证毕。
\end{proof}

\begin{corollary}[严格基态间隙下连续容量曲线的精确仿射肩部]\label{cor:xi-time-part9w-strict-gap-capacity-exact-affine-shoulder}
沿用定理 \ref{thm:xi-time-part9w-strict-gap-tailcount-exact-plateau} 的记号。定义连续容量
\[
\mathcal C_m^{\mathrm{cont}}(T):=
\sum_{x\in X_m}\min\bigl(d_m(x),T\bigr),
\]
并记
\[
A_m^{\max}:=\{x\in X_m:\ d_m(x)=M_m\},
\qquad
\mu_m(A):=\frac1{2^m}\sum_{x\in A}d_m(x).
\]
则对定理 \ref{thm:xi-time-part9w-strict-gap-tailcount-exact-plateau} 中任意阈值序列 \(T_m\)，一切充分大的 \(m\) 都满足
\[
\mathcal C_m^{\mathrm{cont}}(T_m)
=
2^m-K_m(M_m-T_m).
\]
等价地，
\[
\frac{\mathcal C_m^{\mathrm{cont}}(T_m)}{2^m}
=
1-\mu_m(A_m^{\max})\Bigl(1-\frac{T_m}{M_m}\Bigr).
\]
因此在顶端指数带内，连续容量曲线的全部非平凡弯曲最终完全消失，只留下由最大层 \(A_m^{\max}\) 单独控制的精确仿射肩部。
\end{corollary}
\begin{proof}
由定理 \ref{thm:xi-time-part9w-strict-gap-tailcount-exact-plateau}，对充分大的 \(m\) 有
\[
M_m^{(2)}<T_m<M_m.
\]
于是可直接应用命题 \ref{prop:xi-time-part9t-top-gap-affine-capacity-segment}，得到
\[
\mathcal C_m^{\mathrm{cont}}(T_m)=2^m-K_m(M_m-T_m).
\]
两边同除以 \(2^m\)，并注意
\[
\mu_m(A_m^{\max})=\frac{K_mM_m}{2^m},
\]
即得第二式。证毕。
\end{proof}

\noindent
这几条结果表明，有限局部化加法账本、输出饱和前沿与冻结相顶端层之间共享同一种有限化失效机制：在结构层，有限个局部化轴与有限个加法寄存器仍不足以承载完整乘法；在解析层，饱和前沿的二阶曲率在端点发散，排斥一切有限曲率的 Gaussian 切触；而在组合层，一旦严格基态间隙出现，顶端指数带便从渐近平坦提升为最终精确平顶，其对应容量曲线也塌缩为由最大层单独控制的仿射肩部。与前述临界切片刚性与共尾 prime-twist 账本障碍合并来看，有限化失效不再只发生在 prime-ledger 或 Hilbert 载体层，而是同时渗入加法线性化、端点几何与顶端纤维结构三种不同口径。

\endinput
\paragraph{有限局部化加法账本、输出端点非高斯曲率与严格基态间隙的精确平顶}
在有限局部化 rank-one 轴的全局障碍、输出位势右端点的对数尖点、固定单位根模态的四阶展开，以及冻结相中的严格基态间隙与容量仿射段都已建立之后，还可继续抽取出下列几条派生结论。它们分别把有限局部化加法账本的结构失效、饱和前沿的非高斯曲率发散、基本单位根慢模对 Perron 指数的刚性贴近，以及顶端指数带中的尾计数精确平顶与容量精确肩部统一到同一条终端链上。

\begin{theorem}[有限局部化圆因子与有限加法寄存器不足以忠实线性化完整乘法]\label{thm:xi-time-part9w-finite-localized-additive-ledger-obstruction}
设 \(S_1,\dots,S_r\subset\mathbb P\) 为有限素数集，并记
\[
G_{S_i}:=\ZZ[S_i^{-1}].
\]
则对任意 \(k\in\ZZ_{\ge 0}\)，都不存在交换幺半群单射
\[
\Phi:\NN_{\times}\hookrightarrow
\Bigl(\bigoplus_{i=1}^{r}G_{S_i}\Bigr)\oplus(\NN^k,+).
\]
\end{theorem}
\begin{proof}
设
\[
\iota:(\NN^k,+)\hookrightarrow(\ZZ^k,+)
\]
为自然嵌入。若存在所述 \(\Phi\)，则
\[
(\mathrm{id}\oplus\iota)\circ\Phi:
\NN_{\times}\hookrightarrow
\Bigl(\bigoplus_{i=1}^{r}G_{S_i}\Bigr)\oplus\ZZ^k
\]
仍为幺半群单射。把右端改写为
\[
\mathbf S:=(S_1,\dots,S_r),
\qquad
\mathcal L_{k,\mathbf S}:=\ZZ^k\oplus G_{S_1}\oplus\cdots\oplus G_{S_r},
\]
便与定理 \ref{thm:xi-time-part62debb-finite-localization-axes-no-full-multiplicative-host} 直接矛盾。故所述 \(\Phi\) 不存在。证毕。
\end{proof}

\begin{corollary}[输出饱和前沿不存在有限曲率的 Gaussian 切触模型]\label{cor:xi-time-part9w-output-frontier-no-finite-gaussian-tangent}
沿用定理 \ref{thm:xi-output-potential-rate-right-endpoint-log-cusp} 与推论 \ref{cor:xi-output-microcanonical-endpoint-slope-curvature-diverge} 的记号。记
\[
I_{\mathrm{out}}(\alpha):=I(\alpha),
\qquad
h_{\mathrm{out}}(\alpha):=h(\alpha)=\log\lambda(1)-I(\alpha).
\]
则当 \(\alpha\uparrow\tfrac12\) 时，
\[
I_{\mathrm{out}}''(\alpha)=\frac{2}{1/2-\alpha}+O(1),
\qquad
h_{\mathrm{out}}''(\alpha)=-\frac{2}{1/2-\alpha}+O(1).
\]
特别地，\(I_{\mathrm{out}}\) 与 \(h_{\mathrm{out}}\) 在饱和端点都不具有有限二阶曲率；因此任何把输出饱和前沿视为局部 Gaussian 或有限阶 Edgeworth 二次切触的近似，都不可能在端点处成立。
\end{corollary}
\begin{proof}
推论 \ref{cor:xi-output-microcanonical-endpoint-slope-curvature-diverge} 已给出
\[
h_{\mathrm{out}}''(\alpha)=-\frac{2}{1/2-\alpha}+O(1).
\]
又由定义
\[
h_{\mathrm{out}}(\alpha)=\log\lambda(1)-I_{\mathrm{out}}(\alpha),
\]
两边二次求导便得
\[
I_{\mathrm{out}}''(\alpha)=-h_{\mathrm{out}}''(\alpha)
=\frac{2}{1/2-\alpha}+O(1).
\]
于是两式同时成立。证毕。
\end{proof}

\begin{theorem}[基本单位根通道的误差指数以 \texorpdfstring{$m^{-2}$}{m^-2} 速度贴近 Perron 指数]\label{thm:xi-time-part9w-basic-root-unity-error-exponent-to-one}
沿用定理 \ref{thm:xi-output-potential-dirichlet-fixed-j-fourth-order} 的记号。记
\[
\vartheta_{m,1}:=\theta_{m,1}
=\frac{\log\rho_{m,1}}{\log\lambda},
\qquad
\lambda=\varphi^2.
\]
则
\[
\vartheta_{m,1}
=
1-A_{\mathrm{out}}m^{-2}
+B_{\mathrm{out}}m^{-4}
+O(m^{-6}),
\]
其中
\[
A_{\mathrm{out}}
:=
\frac{\pi^2(8955-3983\sqrt5)}{500\log(\varphi^2)}
>0,
\]
\[
B_{\mathrm{out}}
:=
\frac{\pi^4(1354428303-605718497\sqrt5)}{150000\log(\varphi^2)}.
\]
因而
\[
\vartheta_{m,1}\longrightarrow 1
\qquad(m\to\infty).
\]
换言之，基本单位根慢模的误差指数并不趋向平方根门槛 \(\tfrac12\)，而是以显式 \(m^{-2}\) 速率刚性贴近 Perron 主指数 \(1\)。
\end{theorem}
\begin{proof}
定理 \ref{thm:xi-output-potential-dirichlet-fixed-j-fourth-order} 在 \(j=1\) 时给出
\[
\vartheta_{m,1}
=
1-\frac{2\pi^2\kappa_2}{\log\lambda}\,m^{-2}
+\frac{2\pi^4\kappa_4}{3\log\lambda}\,m^{-4}
+O(m^{-6}),
\]
其中 \(\kappa_2=P''(0)\)、\(\kappa_4=P^{(4)}(0)\)。再由定理 \ref{thm:xi-real-input-40-output-edgeworth-nongaussian},
\[
\kappa_2=\frac{1791}{200}-\frac{3983\sqrt5}{1000}
=\frac{8955-3983\sqrt5}{1000},
\]
\[
\kappa_4=\frac{1354428303}{100000}-\frac{605718497\sqrt5}{100000}.
\]
直接代入即可得到所示常数 \(A_{\mathrm{out}},B_{\mathrm{out}}\) 的闭式。由于 \(A_{\mathrm{out}}>0\)，首阶修正为负且量级为 \(m^{-2}\)，故 \(\vartheta_{m,1}\to1\)。证毕。
\end{proof}

\begin{theorem}[严格基态间隙强制顶端指数带上的尾计数精确平顶]\label{thm:xi-time-part9w-strict-gap-tailcount-exact-plateau}
记
\[
M_m:=\max_{x\in X_m}d_m(x),
\qquad
K_m:=\#\{x\in X_m:\ d_m(x)=M_m\},
\]
并令
\[
M_m^{(2)}:=
\max\{d_m(x):\ d_m(x)<M_m\},
\]
若次大层不存在，则约定 \(M_m^{(2)}:=1\)。进一步假设存在常数
\[
\alpha_*^{(2)}<\alpha_*,
\qquad
g_*,
\]
使
\[
\log M_m=\alpha_*m+o(m),
\qquad
\log K_m=g_*m+o(m),
\qquad
\log M_m^{(2)}\le \alpha_*^{(2)}m+o(m).
\]
这正是定理 \ref{thm:xi-time-part9u-frozen-escort-logmultiplicity-mgf} 中所使用的 strict-gap 数据。若阈值序列 \(T_m>0\) 满足
\[
\alpha_*^{(2)}
<
\liminf_{m\to\infty}\frac{1}{m}\log T_m
\le
\limsup_{m\to\infty}\frac{1}{m}\log T_m
<
\alpha_*,
\]
则对一切充分大的 \(m\)，都有
\[
N_m(T_m):=\#\{x\in X_m:\ d_m(x)\ge T_m\}=K_m.
\]
从而
\[
\lim_{m\to\infty}\frac1m\log N_m(T_m)=g_*.
\]
\end{theorem}
\begin{proof}
取 \(\varepsilon>0\) 使
\[
\alpha_*^{(2)}+\varepsilon
<
\liminf_{m\to\infty}\frac1m\log T_m
\le
\limsup_{m\to\infty}\frac1m\log T_m
<
\alpha_*-\varepsilon.
\]
于是对充分大的 \(m\)，一方面由上述指数假设可得
\[
M_m^{(2)}\le \exp\bigl((\alpha_*^{(2)}+\varepsilon)m\bigr),
\qquad
M_m\ge \exp\bigl((\alpha_*-\varepsilon)m\bigr),
\]
另一方面又有
\[
\exp\bigl((\alpha_*^{(2)}+\varepsilon)m\bigr)
<T_m<
\exp\bigl((\alpha_*-\varepsilon)m\bigr).
\]
故最终
\[
M_m^{(2)}<T_m<M_m.
\]
于是所有非最大纤维都严格落在阈值之下，而全部最大纤维都严格超过阈值，因此
\[
N_m(T_m)=\#\{x\in X_m:\ d_m(x)=M_m\}=K_m
\]
对充分大的 \(m\) 成立。再用 \(\log K_m=g_*m+o(m)\) 即得极限式。证毕。
\end{proof}

\begin{corollary}[严格基态间隙下连续容量曲线的精确仿射肩部]\label{cor:xi-time-part9w-strict-gap-capacity-exact-affine-shoulder}
沿用定理 \ref{thm:xi-time-part9w-strict-gap-tailcount-exact-plateau} 的记号。定义连续容量
\[
\mathcal C_m^{\mathrm{cont}}(T):=
\sum_{x\in X_m}\min\bigl(d_m(x),T\bigr),
\]
并记
\[
A_m^{\max}:=\{x\in X_m:\ d_m(x)=M_m\},
\qquad
\mu_m(A):=\frac1{2^m}\sum_{x\in A}d_m(x).
\]
则对定理 \ref{thm:xi-time-part9w-strict-gap-tailcount-exact-plateau} 中任意阈值序列 \(T_m\)，一切充分大的 \(m\) 都满足
\[
\mathcal C_m^{\mathrm{cont}}(T_m)
=
2^m-K_m(M_m-T_m).
\]
等价地，
\[
\frac{\mathcal C_m^{\mathrm{cont}}(T_m)}{2^m}
=
1-\mu_m(A_m^{\max})\Bigl(1-\frac{T_m}{M_m}\Bigr).
\]
因此在顶端指数带内，连续容量曲线的全部非平凡弯曲最终完全消失，只留下由最大层 \(A_m^{\max}\) 单独控制的精确仿射肩部。
\end{corollary}
\begin{proof}
由定理 \ref{thm:xi-time-part9w-strict-gap-tailcount-exact-plateau}，对充分大的 \(m\) 有
\[
M_m^{(2)}<T_m<M_m.
\]
于是可直接应用命题 \ref{prop:xi-time-part9t-top-gap-affine-capacity-segment}，得到
\[
\mathcal C_m^{\mathrm{cont}}(T_m)=2^m-K_m(M_m-T_m).
\]
两边同除以 \(2^m\)，并注意
\[
\mu_m(A_m^{\max})=\frac{K_mM_m}{2^m},
\]
即得第二式。证毕。
\end{proof}

\noindent
这几条结果表明，有限局部化加法账本、输出饱和前沿与冻结相顶端层之间共享同一种有限化失效机制：在结构层，有限个局部化轴与有限个加法寄存器仍不足以承载完整乘法；在解析层，饱和前沿的二阶曲率在端点发散，排斥一切有限曲率的 Gaussian 切触；而在组合层，一旦严格基态间隙出现，顶端指数带便从渐近平坦提升为最终精确平顶，其对应容量曲线也塌缩为由最大层单独控制的仿射肩部。与前述临界切片刚性与共尾 prime-twist 账本障碍合并来看，有限化失效不再只发生在 prime-ledger 或 Hilbert 载体层，而是同时渗入加法线性化、端点几何与顶端纤维结构三种不同口径。

\endinput

\paragraph{Perron 贴近、最大层肩部与 core/full 的非对称可见性}
在输出端点的对数尖点、严格基态间隙所强制的顶端仿射肩部，以及单原子 Euler 因子的 full/core 剥离都已建立之后，还可把这几条链进一步压缩为同一组终端后果。其共同结论是：平方根门槛的失效与顶端容量肩部都由 core 侧的真实指数结构决定，而原子因子本身只在 residue 与 finite-part 层留下显式而刚性的可见壳层。

\begin{corollary}[基本单位根通道最终一致越过平方根门槛]\label{cor:xi-time-part9wa-basic-root-beyond-square-root}
\leanverified{paper\_xi\_time\_part9wa\_basic\_root\_beyond\_square\_root}
沿用定理 \ref{thm:xi-time-part9w-basic-root-unity-error-exponent-to-one} 的记号。则存在 \(m_0\)，使得当 \(m\ge m_0\) 时，
\[
\vartheta_{m,1}>\frac12.
\]
更精确地，
\[
\vartheta_{m,1}
=
1-A_{\mathrm{out}}m^{-2}
+B_{\mathrm{out}}m^{-4}
+O(m^{-6}),
\qquad
m^2(1-\vartheta_{m,1})=A_{\mathrm{out}}+O(m^{-2}).
\]
因此，基本单位根通道并非只在零散窗口中越过平方根门槛，而是最终整体贴近 Perron 指数 \(1\)。
\end{corollary}
\begin{proof}
定理 \ref{thm:xi-time-part9w-basic-root-unity-error-exponent-to-one} 已给出
\[
\vartheta_{m,1}
=
1-A_{\mathrm{out}}m^{-2}
+B_{\mathrm{out}}m^{-4}
+O(m^{-6}),
\qquad
A_{\mathrm{out}}>0.
\]
于是 \(\vartheta_{m,1}\to 1\)，从而存在 \(m_0\) 使得当 \(m\ge m_0\) 时 \(\vartheta_{m,1}>\tfrac12\)。把上式改写为
\[
m^2(1-\vartheta_{m,1})=A_{\mathrm{out}}+O(m^{-2})
\]
便得第二条量化式。证毕。
\end{proof}

\begin{corollary}[顶端指数带完全由最大层控制]\label{cor:xi-time-part9wa-top-band-max-layer-sufficiency}
沿用定理 \ref{thm:xi-time-part9w-strict-gap-tailcount-exact-plateau} 与推论 \ref{cor:xi-time-part9w-strict-gap-capacity-exact-affine-shoulder} 的记号。若阈值序列 \(T_m\) 落在严格基态间隙所允许的顶端指数带内，则对一切充分大的 \(m\) 都有
\[
N_m(T_m)=K_m,
\qquad
\mathcal C_m^{\mathrm{cont}}(T_m)=2^m-K_m(M_m-T_m).
\]
因而在整个顶端指数带中，所有次极大层都对尾计数与连续容量失去可见性，容量曲线最终塌缩为由最大层单独决定的精确仿射肩部。
\end{corollary}
\begin{proof}
第一式正是定理 \ref{thm:xi-time-part9w-strict-gap-tailcount-exact-plateau} 的结论，第二式则由推论 \ref{cor:xi-time-part9w-strict-gap-capacity-exact-affine-shoulder} 直接给出。证毕。
\end{proof}

\begin{theorem}[RH-sharp 门槛完全由 core 生成，而 residue 壳保持刚性]\label{thm:xi-time-part9wa-core-generated-threshold-residue-shell}
对任意 \(m\ge 2\) 与 \(1\le j\le m-1\)，沿用定理 \ref{thm:xi-time-part57b-atomic-factor-zero-twisted-exponent} 的记号，则
\[
\theta_{m,j}^{\mathrm{full}}=\max\{0,\theta_{m,j}^{\mathrm{core}}\},
\qquad
\theta_{m,j}^{\mathrm{full}}>\frac12
\iff
\theta_{m,j}^{\mathrm{core}}>\frac12.
\]
另一方面，对任意 residue 类 \(r\in \ZZ/m\ZZ\)，沿用定理 \ref{thm:conclusion-realinput40-residue-fourier-recovery-delta-trajectory} 与定理 \ref{thm:conclusion-realinput40-primitive-residue-monomial-rigidity} 的记号，有
\[
A_{m,r}^{\mathrm{full}}(n)
=
A_{m,r}^{\mathrm{core}}(n)
+2\,\mathbf 1_{\{2\mid n\}}\mathbf 1_{\{r\equiv n/2\;(\mathrm{mod}\,m)\}},
\]
\[
\mathcal P_{m,r}^{\mathrm{full}}(z)
=
\mathcal P_{m,r}^{\mathrm{core}}(z)
+z^2\,\mathbf 1_{\{r\equiv 1\;(\mathrm{mod}\,m)\}},
\qquad
\mathcal P_{m,r}^{\star}(z):=\sum_{n\ge 1}P_{m,r}^{\star}(n)\,z^n.
\]
因此，单原子 Euler 因子只在 periodic 层产生一条沿 \(r\equiv n/2\pmod m\) 移动的 Dirac 壳，并在 primitive 层产生一个唯一的二次单项式；它既不能制造新的正 twisted 指数，也不能改变任何 residue-resolved primitive 生成函数的非多项式奇点、收敛半径与指数增长率。凡真正的 RH-sharp 门槛失效，均严格来自 core。
\end{theorem}
\begin{proof}
前两条关于 \(\theta_{m,j}^{\mathrm{full}}\) 与 \(\theta_{m,j}^{\mathrm{core}}\) 的恒等式与平方根门槛等价性，分别由定理 \ref{thm:xi-time-part57b-atomic-factor-zero-twisted-exponent} 与推论 \ref{cor:xi-time-part57b-core-full-square-root-barrier-equivalence} 给出。

residue 差分公式
\[
A_{m,r}^{\mathrm{full}}(n)
=
A_{m,r}^{\mathrm{core}}(n)
+2\,\mathbf 1_{\{2\mid n\}}\mathbf 1_{\{r\equiv n/2\;(\mathrm{mod}\,m)\}}
\]
正是定理 \ref{thm:conclusion-realinput40-residue-fourier-recovery-delta-trajectory} 的 periodic 部分；primitive 生成函数的 monomial 刚性
\[
\mathcal P_{m,r}^{\mathrm{full}}(z)
=
\mathcal P_{m,r}^{\mathrm{core}}(z)
+z^2\,\mathbf 1_{\{r\equiv 1\;(\mathrm{mod}\,m)\}}
\]
则由定理 \ref{thm:conclusion-realinput40-primitive-residue-monomial-rigidity} 直接成立。

由于 full/core primitive residue 生成函数之差只是次数 \(2\) 的多项式，二者的全部非多项式奇点与收敛半径完全一致；再由 Cauchy--Hadamard 公式，它们的系数指数增长率亦一致。于是，任何由主导奇点、收敛半径或指数增长率定义的 residue-resolved primitive 阈值都不会被该原子因子改变。结合 twisted 指数的平方根门槛等价性，便得到最后的结论。证毕。
\end{proof}

\begin{corollary}[压力透明而有限部记忆不可压平]\label{cor:xi-time-part9wa-pressure-blind-finitepart-memory}
沿用定理 \ref{thm:xi-atomic-witt-surgery-pressure-ldp-abel-invariance} 与推论 \ref{cor:conclusion-realinput40-zero-temp-finitepart-displacement} 的记号。则
\[
P(\theta)=P_{\mathrm{core}}(\theta),
\qquad
\Delta_{\mathfrak M}(\theta)
:=
\log\mathfrak M(e^\theta)-\log\mathfrak M_{\mathrm{core}}(e^\theta)
=
e^\theta\lambda(e^\theta)^{-2}
>0,
\]
并且当 \(\theta\to+\infty\) 时，
\[
\Delta_{\mathfrak M}(\theta)
=
c^{-2}\left(1-\frac{2d}{c}e^{-\theta/2}+O(e^{-\theta})\right),
\qquad
\lim_{\theta\to+\infty}\Delta_{\mathfrak M}(\theta)=\rho(B_\infty)^{-1}.
\]
特别地，不存在任何仅依赖于压力 germ 数据
\[
\bigl(P^{(k)}(\theta)\bigr)_{k\ge 0}
\]
的标量反项，能够在 full 与 core 两侧同时消去 Abel 有限部的差额。换言之，原子因子在压力层完全不可见，但其 finite-part 记忆不仅严格为正，而且在零温端点仍保留非零极限。
\end{corollary}
\begin{proof}
由定理 \ref{thm:xi-atomic-witt-surgery-pressure-ldp-abel-invariance}，
\[
P(\theta)=P_{\mathrm{core}}(\theta),
\qquad
\Delta_{\mathfrak M}(\theta)=e^\theta\lambda(e^\theta)^{-2}>0.
\]
再由推论 \ref{cor:conclusion-realinput40-zero-temp-finitepart-displacement}，
\[
\Delta_{\mathfrak M}(\theta)
=
c^{-2}\left(1-\frac{2d}{c}e^{-\theta/2}+O(e^{-\theta})\right),
\qquad
\lim_{\theta\to+\infty}\Delta_{\mathfrak M}(\theta)=\rho(B_\infty)^{-1}.
\]

现设 \(\mathcal R\) 是任意一个仅依赖压力 germ 数据 \(\bigl(P^{(k)}(\theta)\bigr)_{k\ge 0}\) 的标量反项。由于
\[
P(\theta)=P_{\mathrm{core}}(\theta)
\]
恒成立，故 full 与 core 对应的反项完全相同。若 \(\mathcal R\) 能同时消去两侧 Abel 有限部差额，则应有
\[
\log\mathfrak M(e^\theta)-\mathcal R
=
\log\mathfrak M_{\mathrm{core}}(e^\theta)-\mathcal R,
\]
从而推出 \(\Delta_{\mathfrak M}(\theta)=0\)，这与 \(\Delta_{\mathfrak M}(\theta)>0\) 矛盾。故此类仅由压力 germ 数据决定的反项不存在。证毕。
\end{proof}

\noindent
因此，输出端点与原子剥离这两条看似分离的链在终端处实际上闭合为同一图景：一方面，基本单位根通道最终整体越过平方根门槛，而严格基态间隙又把顶端指数带压成最大层单独控制的精确肩部；另一方面，所有真正的 RH-sharp 门槛失效都由 core 侧生成，被剥离的原子因子既不能制造新的正 twisted 指数，也不能改变 residue-resolved primitive 的非多项式奇性；但它同时在 finite-part 层留下严格正且零温不灭的记忆。于是，pressure 层与 finite-part 层在 full/core 分裂下表现出根本不同的可见性：前者完全透明，后者则刚性记忆原子缺陷。

\endinput

\paragraph{共振整函数的 Shannon 非可寻址性、局部 germ 传输与 Mellin 桥接猜想}
在共振整函数的 Cartwright 精确型、零点自相似递推与对数 Taylor 级数已经确定，而 fold-\(\pi\) 常数尾的奇指数 Bernoulli 层级亦已固定之后，还可把这两条主线继续压缩为下列七条接口结果。它们分别刻画整数栅格下的 Shannon 非可寻址性、零点计数的精确重整化、线性密度律、局部共振 germ 与 fold-\(\pi\) 奇层尾之间的 Bernoulli-free 传输、逐阶 Bernoulli 消去、有限共振射线族对 bulk 碰撞的不可压缩性，以及可能的 germ--Mellin 唯一桥接原理。

\begin{theorem}[共振整函数的 Shannon--Paley--Wiener 非可寻址性]\label{thm:xi-time-part9v-resonance-shannon-paley-wiener-obstruction}
记 \(\mathrm{PW}_\pi\) 为经典 Shannon--Paley--Wiener 空间。则
\[
\mathcal F_\varphi\notin \mathrm{PW}_\pi.
\]
更一般地，\(\mathcal F_\varphi\) 不属于任何以整数栅格 \(\ZZ\) 为稳定完全采样集的带限 \(L^2\) 整函数类。因而，整数共振样值
\[
u\longmapsto C_\varphi(u)=|\mathcal F_\varphi(u)|
\]
不是 Shannon 意义下可寻址的频谱对象。
\end{theorem}
\begin{proof}
由定理 \ref{thm:fold-pisot-bernoulli-convolution-representation} 的整数采样公式，
\[
|\mathcal F_\varphi(u)|=C_\varphi(u)
\qquad (u\in\ZZ).
\]
又由于 \(\mathcal F_\varphi\) 为偶函数且 \(\mathcal F_\varphi(0)=1\)，故
\[
\sum_{u\in\ZZ}|\mathcal F_\varphi(u)|^2
=
1+2\sum_{u=1}^{\infty}C_\varphi(u)^2.
\]
而定理 \ref{thm:fold-sigma-phi-diverges} 给出
\[
1+2\sum_{u=1}^{\infty}C_\varphi(u)^2=+\infty.
\]
因此
\[
\bigl(\mathcal F_\varphi(u)\bigr)_{u\in\ZZ}\notin \ell^2(\ZZ).
\]
然而，对任意 \(f\in \mathrm{PW}_\pi\)，经典 Shannon--Parseval 理论都给出
\(
(f(n))_{n\in\ZZ}\in \ell^2(\ZZ)
\)；
故 \(\mathcal F_\varphi\notin \mathrm{PW}_\pi\)。

更一般地，若 \(\mathcal B\) 是任意以 \(\ZZ\) 为稳定完全采样集的带限 \(L^2\) 整函数类，则样值算子
\[
f\longmapsto (f(n))_{n\in\ZZ}
\]
必连续取值于 \(\ell^2(\ZZ)\)。由于 \(\mathcal F_\varphi\) 的整数样值列不在 \(\ell^2(\ZZ)\)，故 \(\mathcal F_\varphi\notin \mathcal B\)。最后由 \(C_\varphi(u)=|\mathcal F_\varphi(u)|\) 即得结论。证毕。
\end{proof}

\begin{theorem}[共振零点计数的精确重整化方程]\label{thm:xi-time-part9v-resonance-zero-count-renormalization}
\leanverified{paper\_xi\_time\_part9v\_resonance\_zero\_count\_renormalization}
记
\[
N_\varphi(R):=\#\bigl(\mathcal Z(\mathcal F_\varphi)\cap[-R,R]\bigr)
\qquad (R\ge 0).
\]
则对每个 \(R\ge 0\) 都有严格恒等式
\[
N_\varphi(R)
=
N_\varphi\!\left(\frac{R}{\varphi}\right)
+
2\left\lfloor \frac{R}{\varphi^2}+\frac12\right\rfloor.
\]
等价地，
\[
N_\varphi(\varphi R)
=
N_\varphi(R)
+
2\left\lfloor \frac{R}{\varphi}+\frac12\right\rfloor.
\]
\end{theorem}
\begin{proof}
当 \(R=0\) 时结论显然成立。以下设 \(R>0\)。由定理 \ref{thm:fold-resonance-entire-lp} 的函数方程，
\[
\mathcal F_\varphi(z)
=
\cos\!\Bigl(\frac{\pi z}{\varphi^2}\Bigr)\,
\mathcal F_\varphi\!\Bigl(\frac{z}{\varphi}\Bigr).
\]
第一因子的零点恰为
\[
z=\left(k+\tfrac12\right)\varphi^2,
\qquad
k\in\ZZ,
\]
其落在区间 \([-R,R]\) 内的个数为
\[
2\left\lfloor \frac{R}{\varphi^2}+\frac12\right\rfloor.
\]

另一方面，定理 \ref{thm:fold-resonance-entire-lp} 还给出
\[
\mathcal Z(\mathcal F_\varphi)
=
\Bigl\{\left(k+\tfrac12\right)\varphi^n:\ k\in\ZZ,\ n\ge 2\Bigr\},
\]
故 \(\cos(\pi z/\varphi^2)\) 的零点对应 \(n=2\) 层，而 \(\mathcal F_\varphi(z/\varphi)\) 的零点对应 \(n\ge 3\) 层；不同层零点互不相交。于是 \([-R,R]\) 内的全部零点精确分解为两部分：一部分来自第一因子，另一部分来自第二因子。

对第二因子，\(\mathcal F_\varphi(z/\varphi)=0\) 当且仅当 \(z/\varphi\in \mathcal Z(\mathcal F_\varphi)\)。因此
\[
\#\Bigl(\mathcal Z\bigl(\mathcal F_\varphi(z/\varphi)\bigr)\cap[-R,R]\Bigr)
=
N_\varphi\!\left(\frac{R}{\varphi}\right).
\]
两部分相加即得
\[
N_\varphi(R)
=
N_\varphi\!\left(\frac{R}{\varphi}\right)
+
2\left\lfloor \frac{R}{\varphi^2}+\frac12\right\rfloor.
\]
再把 \(R\) 替换为 \(\varphi R\)，便得到等价的前向重整化式
\[
N_\varphi(\varphi R)
=
N_\varphi(R)
+
2\left\lfloor \frac{R}{\varphi}+\frac12\right\rfloor.
\]
证毕。
\end{proof}

\begin{corollary}[共振零点的线性密度律]\label{cor:xi-time-part9v-resonance-zero-density-linear}
\leanverified{paper\_xi\_time\_part9v\_resonance\_zero\_density\_linear}
沿用定理 \ref{thm:xi-time-part9v-resonance-zero-count-renormalization} 的记号，则
\[
N_\varphi(R)=2R+O(\log R)
\qquad (R\to\infty),
\]
并且
\[
\lim_{R\to\infty}\frac{N_\varphi(R)}{2R}=1.
\]
\end{corollary}
\begin{proof}
记
\[
E(R):=N_\varphi(R)-2R.
\]
由定理 \ref{thm:xi-time-part9v-resonance-zero-count-renormalization}，
\[
E(\varphi R)
=
E(R)+
2\left\lfloor \frac{R}{\varphi}+\frac12\right\rfloor
-
\frac{2R}{\varphi}.
\]
因此存在绝对常数 \(C>0\)，使得
\[
|E(\varphi R)-E(R)|\le C
\qquad (R>0).
\]
取 \(J:=\lfloor \log_\varphi R\rfloor\)，反复迭代直到 \(R\varphi^{-J}=O(1)\)，可得
\[
|E(R)|\le |E(R\varphi^{-J})|+CJ=O(\log R).
\]
于是
\[
N_\varphi(R)=2R+O(\log R).
\]
两边同除以 \(2R\) 并令 \(R\to\infty\)，便得到极限密度公式。证毕。
\end{proof}

\begin{theorem}[共振 germ 与 fold-\(\pi\) 奇层尾的 Bernoulli-free 传输律]\label{thm:xi-time-part9v-resonance-germ-foldpi-bernoulli-free-transfer}
在圆盘
\[
|z|<\frac{\varphi^2}{2}
\]
内写
\[
\log \mathcal F_\varphi(z)=-\sum_{r=1}^{\infty}b_r z^{2r}.
\]
再记 \(a_r\) 为定理 \ref{thm:fold-bin-gauge-constant-stirling-bernoulli-hierarchy} 中
\(
\log C_m-\log(2\pi)
\)
的第 \(r\) 个奇指数 Bernoulli 系数，即
\[
a_r=\frac{B_{2r}}{r(2r-1)}
\bigl(\varphi^{-1}+\varphi^{2r-3}\bigr),
\]
并定义
\[
\widetilde a_r:=(-1)^{r-1}a_r>0.
\]
则对每个 \(r\ge 1\)，都有
\[
b_r=
\frac{2^{2r}(2^{2r}-1)|B_{2r}|\pi^{2r}}{2r(2r)!}\,
\frac{\varphi^{-2r}}{\varphi^{2r}-1},
\]
并且满足严格比值恒等式
\[
\frac{\widetilde a_r}{b_r}
=
\frac{2(2r)!}{2^{2r}(2^{2r}-1)(2r-1)\pi^{2r}}\,
\bigl(\varphi^{-1}+\varphi^{2r-3}\bigr)\,
\varphi^{2r}(\varphi^{2r}-1).
\]
特别地，共振整函数在原点的局部 germ 与 fold-\(\pi\) 奇层尾共享同一 Bernoulli 骨架；二者之差完全由 \(\varphi\)-几何因子与 \(\pi^{-2r}\) 尺度转移给出。
\end{theorem}
\begin{proof}
由定理 \ref{thm:xi-time-part71-resonance-entire-log-taylor-zero-moments}，
\[
\log \mathcal F_\varphi(z)
=
-\sum_{r=1}^{\infty}\frac{\sigma_r}{r}\,z^{2r},
\qquad
\sigma_r=
(2^{2r}-1)\zeta(2r)\frac{\varphi^{-2r}}{\varphi^{2r}-1}.
\]
故 \(b_r=\sigma_r/r\)。再用 Euler 公式
\[
\zeta(2r)=\frac{2^{2r-1}|B_{2r}|\pi^{2r}}{(2r)!},
\]
便得到
\[
b_r=
\frac{2^{2r}(2^{2r}-1)|B_{2r}|\pi^{2r}}{2r(2r)!}\,
\frac{\varphi^{-2r}}{\varphi^{2r}-1}.
\]
另一方面，定理 \ref{thm:fold-bin-gauge-constant-stirling-bernoulli-hierarchy} 给出
\[
a_r=\frac{B_{2r}}{r(2r-1)}
\bigl(\varphi^{-1}+\varphi^{2r-3}\bigr).
\]
由于
\[
(-1)^{r-1}B_{2r}=|B_{2r}|,
\]
故
\[
\widetilde a_r=
\frac{|B_{2r}|}{r(2r-1)}
\bigl(\varphi^{-1}+\varphi^{2r-3}\bigr).
\]
将 \(\widetilde a_r\) 与 \(b_r\) 直接相除，Bernoulli 因子完全消去，即得所示比值公式。证毕。
\end{proof}

\begin{corollary}[Bernoulli 消去原理]\label{cor:xi-time-part9v-bernoulli-elimination-principle}
沿用定理 \ref{thm:xi-time-part9v-resonance-germ-foldpi-bernoulli-free-transfer} 的记号，对每个 \(r\ge 1\) 都有
\[
\frac{\widetilde a_r}{-[z^{2r}]\log \mathcal F_\varphi(z)}
\in
\QQ(\varphi,\pi).
\]
换言之，当 fold-\(\pi\) 奇层系数与共振整函数的局部 Maclaurin 系数逐阶相除时，全部 Bernoulli 数与全部偶值 \(\zeta(2r)\) 信息同时消失；残余完全退化为代数--几何尺度因子。
\end{corollary}
\begin{proof}
由定义
\[
-[z^{2r}]\log \mathcal F_\varphi(z)=b_r.
\]
再由定理 \ref{thm:xi-time-part9v-resonance-germ-foldpi-bernoulli-free-transfer}，
\(
\widetilde a_r/b_r
\)
已显式写成仅含 \(r,\varphi,\pi\) 的表达式，故结论成立。证毕。
\end{proof}

\begin{theorem}[有限共振射线族对 bulk 碰撞的不可压缩性]\label{thm:xi-time-part9v-finite-ray-family-bulk-collision-incompressibility}
设 \(U_m\subset\NN\) 为一列整数集合，并满足
\[
\sup_{m\ge 1}|U_m|<\infty.
\]
定义由该射线族承载的截断共振能量
\[
E_m(U_m):=1+2\sum_{u\in U_m}C_\varphi(u)^2.
\]
则有
\[
F_{m+2}\,\mathsf{Col}_m-E_m(U_m)\longrightarrow +\infty
\qquad (m\to\infty).
\]
特别地，任何只保留一致有界条整数共振射线的方案，都不可能在主尺度上饱和 bulk 碰撞。
\end{theorem}
\begin{proof}
记
\[
K:=\sup_{m\ge 1}|U_m|<\infty.
\]
由于 \(0\le C_\varphi(u)\le 1\)，对一切 \(m\) 都有
\[
E_m(U_m)\le 1+2K.
\]
另一方面，推论 \ref{cor:fold-collision-scale-diverges} 给出
\[
F_{m+2}\,\mathsf{Col}_m\longrightarrow +\infty.
\]
于是
\[
F_{m+2}\,\mathsf{Col}_m-E_m(U_m)
\ge
F_{m+2}\,\mathsf{Col}_m-(1+2K)
\longrightarrow +\infty.
\]
证毕。
\end{proof}

\begin{conjecture}[共振 germ 与 fold-\(\pi\) 奇层尾的 Mellin 唯一桥接原理]\label{conj:xi-time-part9v-resonance-germ-mellin-unique-bridge}
在本文框架内，若存在一条严格、无外参且可递阶延拓的桥接，把共振侧与 fold-\(\pi\) 侧直接联系起来，则该桥接不可能来自
\begin{enumerate}
  \item 整数栅格上的 Shannon 型采样；
  \item 任意基数一致有界的整数共振射线族；
  \item 任意把共振样值闭包压入 \(\ell^2\)-可控离散频谱的方案。
\end{enumerate}
与现有结果相容的候选机制应当是：先取
\[
-\log \mathcal F_\varphi(z)=\sum_{r=1}^{\infty}b_r z^{2r},
\]
再施加某个仅依赖于 \(\varphi\) 的 Mellin 型重整化算子 \(\mathfrak M_\varphi\)，使
\[
\mathfrak M_\varphi\bigl((b_r)_{r\ge 1}\bigr)
=
(a_r)_{r\ge 1}.
\]
\end{conjecture}
\begin{proof}[证据]
定理 \ref{thm:xi-time-part9v-resonance-shannon-paley-wiener-obstruction} 排除了 Shannon--Paley--Wiener 型桥接；定理 \ref{thm:xi-time-part9v-finite-ray-family-bulk-collision-incompressibility} 排除了基数一致有界的共振射线压缩；而定理 \ref{thm:xi-time-part9v-resonance-germ-foldpi-bernoulli-free-transfer} 与推论 \ref{cor:xi-time-part9v-bernoulli-elimination-principle} 又表明，在局部系数层面已经存在逐阶、显式且 Bernoulli 因子完全消失的传输比。因而，若严格桥接存在，则它不应发生在整数样值层，而应发生在局部 germ 经 Mellin 型重整化后的系数层。证毕。
\end{proof}

\noindent
上述结果表明，共振整函数 \(\mathcal F_\varphi\) 一方面处在 Laguerre--P\'olya/Cartwright 的临界边界上，另一方面却在整数栅格上失去 Shannon 可寻址性，并在 bulk 碰撞尺度上拒绝任何基数一致有界的共振射线压缩。与此同时，fold-\(\pi\) 奇层尾的系数与 \(-\log \mathcal F_\varphi\) 的局部系数又共享同一 Bernoulli 骨架，其差异完全退化为 \(\varphi\)-几何因子与 \(\pi^{-2r}\) 型尺度转移。因而，两侧之间真正可能的解析接口不再是整数样值的离散封装，而是局部 germ 经 Mellin 型重整化后的系数传输。

\endinput

\paragraph{共振 Stieltjes 影子、对数导数极点塔与 fold-\(\pi\) 奇层的有限深度障碍}
在定理 \ref{thm:xi-time-part9v-resonance-germ-foldpi-bernoulli-free-transfer} 已把共振 germ 与 fold-\(\pi\) 奇层尾之间的逐阶 Bernoulli-free 传输写成显式比值，而猜想 \ref{conj:xi-time-part9v-resonance-germ-mellin-unique-bridge} 又已指出任何严格桥接都必须离开整数采样层之后，还可把两侧之间的解析断裂进一步压缩为下列五条结论。它们分别刻画共振 Stieltjes 影子的精确重整化方程与唯一性、有限步迭代下的极点生成律、共振对数导数极点塔的有限深度不可消除性、从共振解析 germ 到 Bernoulli 形式尾的有限局部解析变换障碍，以及 fold 标量影子与共振标量影子之间的终止 continued-fraction 深度分叉。

\begin{theorem}[共振 Stieltjes 影子的精确重整化方程与解析唯一性]\label{thm:xi-time-part9vk-resonance-stieltjes-shadow-renormalization-uniqueness}
\leanverified{paper\_xi\_time\_part9vk\_resonance\_stieltjes\_shadow\_renormalization\_uniqueness}
定义
\[
R_\varphi(z):=
\sum_{n=2}^{\infty}\frac{\varphi^{-2n}}{1-\varphi^{-2n}z},
\qquad
z\in\CC\setminus\{\varphi^{2n}:n\ge 2\}.
\]
则该级数在避开极点集的紧集上一致收敛，从而定义一个亚纯函数，并满足严格恒等式
\[
R_\varphi(z)=\frac{\varphi^{-4}}{1-\varphi^{-4}z}
+\varphi^{-2}R_\varphi(\varphi^{-2}z).
\]
更强地，在所有于 \(z=0\) 解析的 germ 中，上述仿射缩放方程只有唯一解；该唯一解即为 \(R_\varphi\) 在 \(0\) 邻域的解析 germ。
\end{theorem}
\begin{proof}
对任意避开极点集
\[
\{\varphi^{2n}:n\ge 2\}
\]
的紧集 \(K\subset\CC\)，存在常数 \(\delta_K>0\) 使
\[
|1-\varphi^{-2n}z|\ge \delta_K
\qquad (z\in K,\ n\ge 2).
\]
因此
\[
\sum_{n=2}^{\infty}\sup_{z\in K}
\left|
\frac{\varphi^{-2n}}{1-\varphi^{-2n}z}
\right|
\le
\delta_K^{-1}\sum_{n=2}^{\infty}\varphi^{-2n}<\infty,
\]
故级数在 \(K\) 上绝对且一致收敛，于是 \(R_\varphi\) 在极点集外亚纯。

把定义式拆出 \(n=2\) 项并重编号，得到
\[
R_\varphi(z)
=
\frac{\varphi^{-4}}{1-\varphi^{-4}z}
+\sum_{n=3}^{\infty}\frac{\varphi^{-2n}}{1-\varphi^{-2n}z}
\]
\[
=
\frac{\varphi^{-4}}{1-\varphi^{-4}z}
+\varphi^{-2}\sum_{m=2}^{\infty}
\frac{\varphi^{-2m}}{1-\varphi^{-2m}(\varphi^{-2}z)},
\]
这正是
\[
R_\varphi(z)=\frac{\varphi^{-4}}{1-\varphi^{-4}z}
+\varphi^{-2}R_\varphi(\varphi^{-2}z).
\]

现证 germ 唯一性。设 \(H\) 为 \(0\) 邻域解析函数，满足齐次方程
\[
H(z)=\varphi^{-2}H(\varphi^{-2}z).
\]
写作幂级数
\[
H(z)=\sum_{m\ge 0}h_m z^m.
\]
代回得
\[
\sum_{m\ge 0}h_m z^m
=
\sum_{m\ge 0}\varphi^{-2m-2}h_m z^m.
\]
比较系数便有
\[
\bigl(1-\varphi^{-2m-2}\bigr)h_m=0
\qquad (m\ge 0).
\]
由于 \(\varphi^{-2m-2}\neq 1\)，故对一切 \(m\) 都有 \(h_m=0\)。因此齐次方程只有零解，从而非齐次方程在解析 germ 类中至多有一个解。又 \(R_\varphi\) 已给出一个解析 germ 解，故唯一性成立。证毕。
\end{proof}

\begin{corollary}[共振影子的有限步极点生成律]\label{cor:xi-time-part9vk-resonance-shadow-finite-step-pole-generation}
\leanverified{paper\_xi\_time\_part9vk\_resonance\_shadow\_finite\_step\_pole\_generation}
定义仿射重整化算子
\[
(\mathscr T_\varphi f)(z):=
\frac{\varphi^{-4}}{1-\varphi^{-4}z}
+\varphi^{-2}f(\varphi^{-2}z).
\]
则对每个整数 \(N\ge 1\)，都有显式闭式
\[
\mathscr T_\varphi^{\circ N}(0)(z)
=
\sum_{n=2}^{N+1}\frac{\varphi^{-2n}}{1-\varphi^{-2n}z}.
\]
因此成立：
\begin{enumerate}
  \item \(\mathscr T_\varphi^{\circ N}(0)\) 恰有 \(N\) 个有限简单极点，位置为
  \[
  z=\varphi^4,\varphi^6,\dots,\varphi^{2N+2};
  \]
  \item \(R_\varphi\) 是 \(\mathscr T_\varphi\) 在解析 germ 类中的唯一固定点；
  \item \(R_\varphi\) 不能在任何有限步重整化后闭合为终止对象，而只能作为“每一步新增一个极点”的无限重整化极限出现。
\end{enumerate}
\end{corollary}
\begin{proof}
对 \(N=1\)，有
\[
\mathscr T_\varphi(0)(z)=\frac{\varphi^{-4}}{1-\varphi^{-4}z},
\]
结论成立。若对某个 \(N\ge 1\) 已知
\[
\mathscr T_\varphi^{\circ N}(0)(z)
=
\sum_{n=2}^{N+1}\frac{\varphi^{-2n}}{1-\varphi^{-2n}z},
\]
则
\[
\mathscr T_\varphi^{\circ(N+1)}(0)(z)
=
\frac{\varphi^{-4}}{1-\varphi^{-4}z}
+\varphi^{-2}\sum_{n=2}^{N+1}
\frac{\varphi^{-2n}}{1-\varphi^{-2n}(\varphi^{-2}z)}
\]
\[
=
\frac{\varphi^{-4}}{1-\varphi^{-4}z}
+\sum_{n=3}^{N+2}\frac{\varphi^{-2n}}{1-\varphi^{-2n}z}
=
\sum_{n=2}^{N+2}\frac{\varphi^{-2n}}{1-\varphi^{-2n}z}.
\]
故归纳成立。

这立刻给出第 \(1\) 条：每一步迭代都会新增分母
\[
1-\varphi^{-2(N+1)}z,
\]
因而恰新增一个有限简单极点。

若某解析 germ \(f\) 满足 \(\mathscr T_\varphi f=f\)，则它正满足定理
\ref{thm:xi-time-part9vk-resonance-stieltjes-shadow-renormalization-uniqueness}
中的同一方程，因此固定点在解析 germ 类中唯一，且必等于 \(R_\varphi\)。这证明了第 \(2\) 条。

最后，由上式可见 \(\mathscr T_\varphi^{\circ N}(0)\) 始终只具有有限个极点，而 \(R_\varphi\) 的极点集却是
\[
\{\varphi^{2n}:n\ge 2\},
\]
为无穷集，故 \(R_\varphi\) 不可能在任何有限步达到。证毕。
\end{proof}

\begin{theorem}[共振对数导数的有限深度极点消除不可能性]\label{thm:xi-time-part9vk-logderivative-finite-depth-pole-cancellation-impossible}
沿用定理 \ref{thm:xi-fold-resonance-logderivative-selfsimilar} 的记号，并记
\[
J_\varphi(z):=\Psi_\varphi(z)=\frac{\mathcal F_\varphi'(z)}{\mathcal F_\varphi(z)}.
\]
则有严格自相似差分方程
\[
J_\varphi(z)=
-\frac{\pi}{\varphi^2}\tan\!\Bigl(\frac{\pi z}{\varphi^2}\Bigr)
+\varphi^{-1}J_\varphi\!\Bigl(\frac{z}{\varphi}\Bigr).
\]
现固定整数 \(N\ge 0\)，取任意复系数 \(c_0,\dots,c_N\)，并定义
\[
\mathscr L_N(z):=
\sum_{j=0}^{N}c_j\,J_\varphi\!\Bigl(\frac{z}{\varphi^j}\Bigr).
\]
则
\[
\mathscr L_N \text{ 为整函数 }
\Longleftrightarrow
c_0=\cdots=c_N=0.
\]
换言之，任意有限深度的共振对数导数重整化都不可能把极点塔完全消除。
\end{theorem}
\begin{proof}
第一式正是定理 \ref{thm:xi-fold-resonance-logderivative-selfsimilar} 的递归恒等式。

再证消除不可能性。由定理 \ref{thm:fold-resonance-entire-lp}，
\[
\mathcal Z(\mathcal F_\varphi)=
\left\{
\left(k+\tfrac12\right)\varphi^m:\ k\in\ZZ,\ m\ge 2
\right\},
\]
且全部为实单零点。故 \(J_\varphi=\mathcal F_\varphi'/\mathcal F_\varphi\) 在每个
\[
\xi_{m,k}:=\left(k+\frac12\right)\varphi^m
\qquad (m\ge 2,\ k\in\ZZ)
\]
处都有留数 \(1\) 的简单极点。

考虑缩放后的项 \(J_\varphi(z/\varphi^j)\)。它在点 \(\xi_{m,k}\) 处出现极点，当且仅当
\[
\frac{\xi_{m,k}}{\varphi^j}
=
\left(k+\frac12\right)\varphi^{m-j}
\]
仍落在 \(\mathcal Z(\mathcal F_\varphi)\) 中，即当且仅当 \(m-j\ge 2\)，也就是 \(j\le m-2\)。此时由简单变量替换的留数公式，
\[
\operatorname{Res}_{z=\xi_{m,k}}
J_\varphi\!\Bigl(\frac{z}{\varphi^j}\Bigr)
=
\varphi^j.
\]
故 \(\mathscr L_N\) 在 \(\xi_{m,k}\) 处的留数为
\[
\operatorname{Res}_{z=\xi_{m,k}}\mathscr L_N
=
\sum_{j=0}^{\min(N,m-2)}c_j\varphi^j.
\]

若 \(\mathscr L_N\) 为整函数，则上述留数对一切 \(m\ge 2\) 与 \(k\in\ZZ\) 都必须为零。取 \(m=2\) 得
\[
c_0=0.
\]
再取 \(m=3\) 得
\[
c_0+c_1\varphi=0,
\]
故 \(c_1=0\)。归纳地，若已知
\[
c_0=\cdots=c_{r-1}=0,
\]
则取 \(m=r+2\)，便得到
\[
\sum_{j=0}^{r}c_j\varphi^j=c_r\varphi^r=0,
\]
故 \(c_r=0\)。于是 \(c_0,\dots,c_N\) 全部为零。

反向蕴含显然。故结论成立。证毕。
\end{proof}

\begin{theorem}[从共振 germ 到 fold-\(\pi\) 奇层 formal tail 的有限局部解析变换不可能性]\label{thm:xi-time-part9vk-no-finite-local-analytic-transform-from-resonance-germ-to-foldpi-tail}
记
\[
G_\varphi(s):=-\log \mathcal F_\varphi(s),
\]
其中对数取 \(s=0\) 邻域的解析支。由 \(\mathcal F_\varphi(0)=1\) 及定理 \ref{thm:fold-resonance-entire-lp} 的零点公式可知，\(G_\varphi\) 在圆盘
\[
|s|<\frac{\varphi^2}{2}
\]
内解析。再记 fold-\(\pi\) 的奇层 formal tail 为
\[
\mathcal E_\varphi(s):=
\sum_{r\ge 1}
\frac{B_{2r}}{r(2r-1)}
\bigl(\varphi^{-1}+\varphi^{2r-3}\bigr)s^{2r-1}.
\]
则不存在任何由下列有限次局部操作生成的表达式，能够把 \(G_\varphi\) 精确变为 \(\mathcal E_\varphi\)：
\begin{enumerate}
  \item 有限次微分 \(f\mapsto f^{(n)}\)；
  \item 有限次缩放 \(f(s)\mapsto f(s/\varphi^j)\)；
  \item 乘以 \(s=0\) 邻域解析函数；
  \item 有限线性组合。
\end{enumerate}
亦即，共振侧的解析 germ 不可能经任何有限局部解析演算，被转化为 fold-\(\pi\) 的零半径 Bernoulli formal tail。
\end{theorem}
\begin{proof}
先证所有允许操作都保持“正收敛半径”这一性质。若 \(f\) 在某圆盘 \(|s|<R\) 中解析，则：
\begin{enumerate}
  \item \(f^{(n)}\) 仍在 \(|s|<R\) 中解析；
  \item \(f(s/\varphi^j)\) 在 \(|s|<\varphi^jR\) 中解析；
  \item 若 \(h\) 在 \(|s|<R_h\) 中解析，则 \(h(s)f(s)\) 在 \(|s|<\min(R,R_h)\) 中解析；
  \item 有限个正半径解析 germ 的线性组合仍具有正收敛半径。
\end{enumerate}
因此，从 \(G_\varphi\) 出发，经任意有限次这四类操作得到的结果，仍必是 \(s=0\) 邻域中的解析 germ，因而其幂级数收敛半径严格为正。

另一方面，\(\mathcal E_\varphi\) 的系数
\[
e_r:=
\frac{B_{2r}}{r(2r-1)}
\bigl(\varphi^{-1}+\varphi^{2r-3}\bigr)
\]
对应于幂次 \(s^{2r-1}\)。由 Euler 的 Bernoulli--\(\zeta\) 恒等式
\[
|B_{2r}|=
\frac{2(2r)!}{(2\pi)^{2r}}\zeta(2r),
\]
且 \(\zeta(2r)\to 1\)，可知存在常数 \(c>0\) 与 \(r_0\) 使得当 \(r\ge r_0\) 时，
\[
|e_r|
\ge
c\,
\frac{(2r)!}{r(2r-1)}
\left(\frac{\varphi^2}{4\pi^2}\right)^r.
\]
由 Stirling 公式便得
\[
|e_r|^{1/(2r-1)}
\longrightarrow +\infty,
\]
故级数
\[
\mathcal E_\varphi(s)=\sum_{r\ge 1}e_r\,s^{2r-1}
\]
的收敛半径等于 \(0\)。

若存在所述有限局部解析变换 \(T\) 使得
\[
T[G_\varphi]=\mathcal E_\varphi
\]
作为 formal series 成立，则二者具有同一系数列，从而 \(T[G_\varphi]\) 也必须具有零收敛半径；这与前半部分“任意此类 \(T[G_\varphi]\) 都具有正收敛半径”矛盾。故所述有限局部解析变换不存在。证毕。
\end{proof}

\begin{proposition}[标量 continued-fraction 深度上的精确分叉]\label{prop:xi-time-part9vk-scalar-continued-fraction-depth-splitting}
\leanverified{paper\_xi\_time\_part9vk\_scalar\_continued\_fraction\_depth\_splitting}
定义 fold 侧的标量影子
\[
S_\varphi(z):=
\sum_{k\ge 1}\bigl(\varphi^{-1}+\varphi^{2k-3}\bigr)z^{k-1}
=
\frac{\varphi^{-1}}{1-z}+\frac{\varphi^{-1}}{1-\varphi^2 z},
\]
以及定理 \ref{thm:xi-time-part9vk-resonance-stieltjes-shadow-renormalization-uniqueness} 中的共振侧标量影子 \(R_\varphi(z)\)。则成立：
\begin{enumerate}
  \item \(S_\varphi\) 属于恰好深度 \(2\) 的终止标量 continued fraction 类；
  \item \(R_\varphi\) 不属于任何终止标量 continued fraction 类；
  \item 因而，fold 标量影子与共振标量影子之间存在严格的有限深度标量 continued fraction 屏障。
\end{enumerate}
\end{proposition}
\begin{proof}
首先，\(S_\varphi\) 是有理函数，并且其分母为二次多项式。对任意在 \(0\) 邻域解析的有理函数，Euclid 递除算法都给出一条终止标量 continued fraction 展开；因此 \(S_\varphi\) 的终止深度至多为 \(2\)。另一方面，若其终止深度为 \(1\)，则 \(S_\varphi\) 必为 Möbius 函数，但由上式可见它具有两个不同的有限极点 \(z=1\) 与 \(z=\varphi^{-2}\)，故不可能是 Möbius 函数。于是其最小终止深度恰为 \(2\)。

其次，任意终止标量 continued fraction 都是有限次代数嵌套所得，因而必为有理函数。然由推论 \ref{cor:xi-time-part9vk-resonance-shadow-finite-step-pole-generation}，\(R_\varphi\) 的极点集是无穷集
\[
\{\varphi^{2n}:n\ge 2\},
\]
故 \(R_\varphi\) 不是有理函数，从而不可能属于任何终止标量 continued fraction 类。

综上，fold 影子仍位于有限深度标量 continued fraction 范畴之内，而共振影子则已严格越出该范畴。证毕。
\end{proof}

\noindent
这五条结果表明，共振侧与 fold-\(\pi\) 侧之间的分离并不只体现在先前已经出现的 Bernoulli-free 比值、Shannon 非可寻址性或 Mellin 候选桥接上；它还在“有限深度”这一层发生更尖锐的断裂。共振侧的标量影子 \(R_\varphi\) 是由仿射缩放方程唯一钉住的解析固定点，并且只能通过“每一步新增一个极点”的无限重整化极限出现；共振对数导数的极点塔又拒绝任何有限深度的留数消除。相对地，fold-\(\pi\) 的奇层尾 \(\mathcal E_\varphi\) 则是一个零收敛半径的 Bernoulli formal object，其标量影子 \(S_\varphi\) 仍停留在深度 \(2\) 的终止 continued-fraction 范畴。因而，两侧之间不存在任何有限深度的极点消除、有限局部解析变换或终止标量 continued fraction 桥接；若严格桥接确有存在，它就不能保持通常的解析半径制度，而只能表现为穿越解析半径边界的奇异重整化机制。

\endinput

\paragraph{解析分歧半径、二维全息承载、可见 Jacobian 赤字与稀疏 prime-language 的共尾障碍}
在核版 RH 的临界边界、\(2\)-进全息地址、\(S_4\)-Hurwitz 商塔的 Jacobian/Prym 分解以及 Zeckendorf--G\"odel 素数语言的密度理论都已确定之后，还可把这些看似分散的接口继续压缩为下列七条终端结论。它们分别刻画 RH-sharp 阈值与解析分歧半径的严格分离、全息地址中首差深度的精确赋值律、前缀柱集与 \(2^L\)-进仿射球的严格对应、固定二维 \(2\)-进 ambient space 与有限秩加法账本不可能性之间的结构分裂、前缀恢复的正线性码率下界、商曲线可见因子对九维 Prym 三重块的严格维数赤字，以及稀疏 prime-register 语言在强局部约束下仍保持的正密度共尾性。

\begin{theorem}[解析分歧半径与 RH-sharp 阈值的严格分离]\label{thm:xi-time-part9w-analytic-branch-radius-rhsharp-separation}
记
\[
P_2(\theta):=\log \rho(e^\theta),
\qquad
R_\theta:=\text{\(P_2\) 在 \(\theta=0\) 处的 Taylor 收敛半径},
\]
并设 \(u_R>1\) 为 kernel RH/Ramanujan 条件的唯一正实临界参数，
\[
\theta_R:=\log u_R.
\]
则有严格不等式
\[
0<\theta_R<R_\theta.
\]
更具体地，记
\[
f(\lambda;u):=\lambda^3-2\lambda^2-(u+1)\lambda+u
\]
为碰撞核的控制性三次特征因子，则对每个 \(u>0\)，其判别式
\[
\Delta(u)=4u^3+25u^2+88u+8
\]
恒满足 \(\Delta(u)>0\)，并且三条实谱支始终保持固定根型
\[
\lambda_+(u)>1,\qquad 0<\lambda_0(u)<1,\qquad \lambda_-(u)<0.
\]
再记
\[
\Lambda(u):=\max\{|\lambda_0(u)|,|\lambda_-(u)|,1\}.
\]
则
\[
\textsf{RH}(u)\iff
\begin{cases}
\text{恒真},& 0<u\le 1,\\[1mm]
(-\lambda_-(u))^2\le \lambda_+(u),& u>1,
\end{cases}
\]
并且在 \(u=u_R\) 处满足
\[
\Lambda(u_R)=\lambda_+(u_R)^{1/2}.
\]
因此，RH-sharp 失效并非由分歧点逼近、复谱生成或根碰撞触发，而是发生在解析域内部的纯实谱阈值穿越。
\end{theorem}
\begin{proof}
由定理 \ref{thm:xi-time-part9pb-rh-transition-analytic-not-branch}，
\[
\theta_R\approx 1.27888224610494,
\qquad
R_\theta\approx 2.76230289346087,
\]
故
\[
0<\theta_R<R_\theta.
\]

另一方面，定理 \ref{thm:xi-time-part9pb-rh-transition-analytic-not-branch} 的证明已给出
\[
\Delta(u)=4u^3+25u^2+88u+8>0
\qquad (u>0),
\]
从而正实轴上不会发生代数分歧塌缩。再由命题
\ref{prop:real-input-40-collision-all-real}，三次因子 \(f(\lambda;u)\) 对每个 \(u>0\) 都具有三条互异实根，并按
\[
\lambda_+(u)>1,\qquad 0<\lambda_0(u)<1,\qquad \lambda_-(u)<0
\]
排序。

定理 \ref{thm:xi-time-part62d-collision-rh-negative-branch-control} 进一步说明：\(\lambda_0(u)\) 从不决定 RH 窗口边界；当 \(u>1\) 时，
\[
\Lambda(u)=|\lambda_-(u)|=-\lambda_-(u),
\]
于是
\[
\textsf{RH}(u)\iff \Lambda(u)\le \lambda_+(u)^{1/2}
\iff
(-\lambda_-(u))^2\le \lambda_+(u).
\]
特别地，在 \(u=u_R\) 处有
\[
\Lambda(u_R)=\lambda_+(u_R)^{1/2}.
\]
因此，RH-sharp 边界既不由解析边界触发，也不由复共轭谱支生成触发，而只由解析域内部的纯实谱不等式翻转触发。证毕。
\end{proof}

\begin{theorem}[\texorpdfstring{$2^L$}{2^L}-进全息地址的首差--赋值恒等式]\label{thm:xi-time-part9w-holographic-first-difference-valuation}
沿用定理 \ref{thm:killo-godel-semidir-affine-2adic} 的记号。设 \(E\) 为有限事件集，
\[
\mathrm{code}:E\hookrightarrow \NN,
\qquad
M:=\max_{e\in E}\mathrm{code}(e),
\qquad
B:=2^L>M,
\]
并取 primitive 地址向量
\[
v\in T:=T_2(E_{\min})\setminus 2T.
\]
对任意无限事件流
\[
\mathbf e=(e_1,e_2,\dots)\in E^{\NN},
\]
定义其全息地址
\[
X(\mathbf e):=\sum_{t\ge 1}\mathrm{code}(e_t)\,B^{t-1}v\in T.
\]
若 \(\mathbf e\neq \mathbf e'\)，并记其首个分歧位置为
\[
\tau(\mathbf e,\mathbf e')
:=
\min\{t\ge 1:\ e_t\neq e_t'\},
\]
则有精确赋值恒等式
\[
X(\mathbf e)-X(\mathbf e')
\in
B^{\tau(\mathbf e,\mathbf e')-1}T\setminus B^{\tau(\mathbf e,\mathbf e')}T.
\]
\leanverified{paper\_xi\_time\_part9w\_holographic\_first\_difference\_valuation}
\end{theorem}
\begin{proof}
记 \(\tau:=\tau(\mathbf e,\mathbf e')\)，并置
\[
\delta:=\mathrm{code}(e_\tau)-\mathrm{code}(e_\tau')\neq 0.
\]
由于 \(\mathrm{code}\) 取值于 \(\{0,\dots,M\}\) 且 \(M<B\)，有
\[
|\delta|<B.
\]
把差值从首个分歧项起展开，得到
\[
\begin{aligned}
X(\mathbf e)-X(\mathbf e')
&=
\sum_{t\ge \tau}
\bigl(\mathrm{code}(e_t)-\mathrm{code}(e_t')\bigr)B^{t-1}v\\
&=
B^{\tau-1}
\left(
\delta v+
\sum_{t>\tau}
\bigl(\mathrm{code}(e_t)-\mathrm{code}(e_t')\bigr)B^{t-\tau}v
\right).
\end{aligned}
\]
括号内第二项属于 \(BT\)，故
\[
X(\mathbf e)-X(\mathbf e')\in B^{\tau-1}T.
\]

现证它不属于 \(B^\tau T\)。若反设
\[
X(\mathbf e)-X(\mathbf e')\in B^\tau T,
\]
则必有
\[
\delta v+By\in BT
\]
对某个 \(y\in T\) 成立，从而 \(\delta v\in BT\)。设
\[
\delta=2^s\delta_0,
\qquad
0\le s<L,
\qquad
\delta_0\in \ZZ_2^\times.
\]
则
\[
2^s\delta_0 v\in 2^L T.
\]
由于 \(\delta_0\) 在 \(\ZZ_2\) 中为单位，可消去得
\[
v\in 2^{L-s}T\subset 2T,
\]
与 \(v\notin 2T\) 矛盾。故
\[
X(\mathbf e)-X(\mathbf e')\notin B^\tau T.
\]
于是所示精确赋值恒等式成立。证毕。
\end{proof}

\begin{corollary}[前缀柱集在全息像中的仿射球等式]\label{cor:xi-time-part9w-cylinder-affine-ball-identity}
沿用定理 \ref{thm:xi-time-part9w-holographic-first-difference-valuation} 的记号。固定长度 \(n\) 的前缀
\[
a=(a_1,\dots,a_n)\in E^n,
\]
并记
\[
x_a:=\sum_{t=1}^{n}\mathrm{code}(a_t)\,B^{t-1}v.
\]
定义 cylinder
\[
[a]
:=
\{\mathbf e\in E^{\NN}:\ e_1=a_1,\dots,e_n=a_n\}.
\]
则有精确仿射球公式
\[
X([a])
=
\bigl(x_a+B^nT\bigr)\cap X(E^{\NN}).
\]
因此，长度 \(n\) 的事件前缀并非仅仅近似决定地址，而是恰好切出一个 \(B^n\)-级 \(2\)-进仿射球。
\end{corollary}
\begin{proof}
若 \(\mathbf e\in[a]\)，则
\[
X(\mathbf e)-x_a
=
\sum_{t>n}\mathrm{code}(e_t)\,B^{t-1}v
\in B^nT,
\]
故
\[
X([a])\subseteq \bigl(x_a+B^nT\bigr)\cap X(E^{\NN}).
\]

反之，设 \(X(\mathbf e)\in (x_a+B^nT)\cap X(E^{\NN})\)。任取一条固定尾流
\[
\boldsymbol\eta=(\eta_1,\eta_2,\dots)\in E^{\NN},
\]
并令
\[
\mathbf f:=(a_1,\dots,a_n,\eta_1,\eta_2,\dots)\in[a].
\]
则 \(X(\mathbf f)-x_a\in B^nT\)，因此
\[
X(\mathbf e)-X(\mathbf f)\in B^nT.
\]
若 \(\mathbf e\) 的前 \(n\) 位中有某一位偏离 \(a\)，则 \(\mathbf e\) 与 \(\mathbf f\) 的首个分歧位置必不超过 \(n\)。由定理
\ref{thm:xi-time-part9w-holographic-first-difference-valuation}，
\[
X(\mathbf e)-X(\mathbf f)\in B^{t-1}T\setminus B^tT
\]
对某个 \(t\le n\) 成立，这与 \(X(\mathbf e)-X(\mathbf f)\in B^nT\) 矛盾。故 \(\mathbf e\) 的前 \(n\) 位必恰为 \(a\)，亦即 \(\mathbf e\in[a]\)。证毕。
\end{proof}

\begin{theorem}[有限二维全息实现与有限秩账本不可能性的严格分离]\label{thm:xi-time-part9w-twoadic-ambient-vs-additive-ledger-separation}
沿用定理 \ref{thm:killo-godel-semidir-affine-2adic} 的记号。任意 canonical 有限时深 prime-register 子单子
\[
H_{\mathrm{can}}\subset \mathcal R\rtimes_\Sigma \NN
\]
都可忠实实现为固定二维 \(2\)-进 Tate 模块
\[
T_2(E_{\min})\cong \ZZ_2^2
\]
上的仿射变换单子。另一方面，正整数乘法幺半群 \(\NN_{\times}\) 不能单射同态嵌入任何有限生成交换可消去加法幺半群，因而也不能嵌入任何有限生成加法账本或有限生成交换群。再者，对任意有限素数集
\[
S\subset \mathbb P,
\qquad
G_S:=\ZZ[S^{-1}],
\qquad
\Sigma_S:=\widehat{G_S},
\]
局部化相位采样仅能唯一恢复该有限素数集 \(S\)，等价地仅能恢复对偶短正合列中的 profinite 核
\[
\prod_{p\in S}\ZZ_p.
\]
因此，在本文框架中，真正的障碍并不是缺乏有限维 ambient space，而是 cofinal prime-multiplicative 结构无法被 finite-rank additive ledger 保持。
\end{theorem}
\begin{proof}
定理 \ref{thm:xi-time-part75-faithful-finite-dimensional-realization-affine-necessity} 已证明：当 \(B=2^L>M\) 时，存在忠实单子嵌入
\[
\Psi:H_{\mathrm{can}}\hookrightarrow \mathrm{Aff}\!\bigl(T_2(E_{\min})\bigr),
\qquad
T_2(E_{\min})\cong \ZZ_2^2.
\]
故任意有限时深的 canonical prime-register 机制都可以放入固定二维 \(2\)-进 ambient space 中。

同一定理还证明：\(\NN_{\times}\) 不能嵌入任何有限生成交换可消去幺半群，因此尤其不能嵌入 \((\NN^k,+)\) 或任何有限生成交换群的加法账本。故一切有限秩 additive bookkeeping 都不足以忠实保存 prime-multiplicative 结构。

最后，由定理 \ref{thm:xi-localized-phase-sampling-closed-form}，局部化相位采样函数 \(\Sigma^{\mathrm{ph}}_{G_S}\) 唯一恢复素数集 \(S\)，并进一步唯一恢复 Pontryagin 对偶短正合列
\[
0\to \prod_{p\in S}\ZZ_p\to \Sigma_S\to \TT\to 0
\]
中的 profinite 核。由于 \(S\) 有限，任何固定 prime-stage 只能携带有限多个 prime 方向，不可能承载全体素数的共尾结构。综上得到所述严格分离。证毕。
\end{proof}

\begin{corollary}[\texorpdfstring{$2^L$}{2^L}-进前缀恢复的最小码率下界]\label{cor:xi-time-part9w-prefix-recovery-minimal-rate}
沿用定理 \ref{thm:xi-time-part9w-holographic-first-difference-valuation} 的全息地址记号，并沿用定理 \ref{thm:xi-optimal-reconstruction-success-probability-constant-bias} 的微态恢复口径。若在分辨率 \(m\) 上，一种恢复协议只允许读取全息地址 \(X(\mathbf e)\) 的前 \(n\) 个 \(B\)-进位，并要求在均匀微态输入下达到成功率
\[
\mathrm{Succ}_m\ge 1-\varepsilon,
\qquad
0<\varepsilon<1,
\]
则必有
\[
n\log_2 B
\ge
m-\log_2|X_m|+\log_2(1-\varepsilon).
\]
特别地，由 \(|X_m|=F_{m+2}\sim \varphi^m\) 可得
\[
\liminf_{m\to\infty}\frac{n}{m}
\ge
\frac{\log_2(2/\varphi)}{\log_2 B}.
\]
\end{corollary}
\begin{proof}
读取前 \(n\) 个 \(B\)-进位，至多只能区分
\[
B^n
\]
种不同前缀标签；换言之，该协议所使用的附加字母表大小 \(M_m\) 至多满足
\[
M_m\le B^n.
\]
因此，其可用辅助信息量
\[
L_m:=\log_2 M_m
\]
满足
\[
L_m\le n\log_2 B.
\]

另一方面，推论 \ref{cor:xi-reconstruction-cardinality-strong-converse} 给出：任何在均匀输入下达到成功率 \(\delta\in(0,1)\) 的恢复协议都必须满足
\[
L_m\ge m-\log_2|X_m|+\log_2\delta.
\]
取 \(\delta=1-\varepsilon\)，便得到
\[
n\log_2 B
\ge
L_m
\ge
m-\log_2|X_m|+\log_2(1-\varepsilon),
\]
即第一式。

再由 \(|X_m|=F_{m+2}\) 与 Binet 渐近
\[
\log_2 F_{m+2}=m\log_2\varphi+O(1),
\]
两边同除以 \(m\log_2 B\) 并取下极限，即得
\[
\liminf_{m\to\infty}\frac{n}{m}
\ge
\frac{\log_2(2/\varphi)}{\log_2 B}.
\]
证毕。
\end{proof}

\begin{theorem}[商曲线可见因子的维数赤字与 Prym 严格多数]\label{thm:xi-time-part9w-visible-jacobian-deficit-prym-majority}
沿用定理 \ref{thm:killo-s4-genus49-jacobian-computable-isogeny} 与结论 \ref{con:xi-time-part9n-s4-all-subgroup-genera} 的记号。记
\[
J_{\mathrm{vis}}:=J(H)\times J(Z)^2\times J(C_F)^3,
\qquad
J_{\mathrm{hid}}:=P^3,
\]
其中
\[
g(H)=5,\qquad g(Z)=4,\qquad g(C_F)=3,\qquad \dim P=9.
\]
则
\[
\dim J_{\mathrm{vis}}=5+2\cdot 4+3\cdot 3=22,
\qquad
\dim J_{\mathrm{hid}}=3\cdot 9=27.
\]
特别地，
\[
27>22.
\]
因此，九维 Prym 因子的三重块不仅非平凡，而且在维数上严格超过全部商曲线可见因子的总和。于是，任何只通过商曲线
\[
X/A_4,\qquad X/D_8,\qquad X/S_3
\]
所对应的可见 Jacobian 乘积来读取证书的阿贝尔读出函子，都必然遗漏 \(J(X)\) 的严格多数部分。
\end{theorem}
\begin{proof}
由定理 \ref{thm:killo-s4-genus49-jacobian-computable-isogeny}，
\[
J(X)\sim_{\QQ}J(H)\times J(Z)^2\times J(C_F)^3\times P^3.
\]
再由结论 \ref{con:xi-time-part9n-s4-all-subgroup-genera} 与相关后续结论，
\[
g(H)=5,\qquad g(Z)=4,\qquad g(C_F)=3,\qquad \dim P=9.
\]
对 Jacobian 有 \(\dim J(C)=g(C)\)，故
\[
\dim J_{\mathrm{vis}}=5+2\cdot 4+3\cdot 3=22,
\qquad
\dim J_{\mathrm{hid}}=3\cdot 9=27.
\]
于是
\[
22+27=49=\dim J(X),
\]
并且 \(27>22\)。

此外，结论 \ref{con:xi-time-part9n-s4-cyclic-quotient-genus-spectrum} 已给出
\[
H^0(X,\Omega_X^1)\cong
5\cdot \mathrm{sgn}\oplus 4\cdot V_2\oplus 3\cdot V_3\oplus 9\cdot V_3',
\]
其中平凡表示不出现。故缺失的 \(27\) 维部分并不是某个可由不变平凡方向补回的低维尾项，而是几何上真正主导的 Prym 扇区。凡是仅经由 \(J_{\mathrm{vis}}\) 因子化的阿贝尔读出，都无法探测 \(P^3\) 这一严格多数块。证毕。
\end{proof}

\begin{theorem}[稀疏 prime-register 语言的正密度共尾性]\label{thm:xi-time-part9w-zg-positive-density-cofinal-tail}
定义
\[
\mathrm{ZG}
:=
\left\{
\prod_{k\ge 1}p_k^{z_k}:\ 
z_k\in\{0,1\},\ 
z_kz_{k+1}=0,\ 
z_k=0\ \text{对充分大 }k
\right\},
\]
即平方自由且不含相邻素因子的 Zeckendorf--G\"odel 素数语言。则：
\begin{enumerate}
  \item \(\mathrm{ZG}\) 具有正的自然密度
  \[
  \delta(\mathrm{ZG})\in(0,1);
  \]
  \item 对任意有限素数集 \(S\subset\mathbb P\)，都存在 \(n\in \mathrm{ZG}\) 使其素因子支撑满足
  \[
  \operatorname{supp}(n)\subset \mathbb P\setminus S;
  \]
  \item 因而，“禁止相邻素因子”这一强局部稀疏约束并不会把 prime-register 语言压成可忽略尾集；它仍保持正密度与共尾素数支撑。
\end{enumerate}
\end{theorem}
\begin{proof}
第一条由定理 \ref{thm:xi-zg-abel-residue-log-density} 直接得到：该定理已证明
\[
\delta_{\mathrm{ZG}}=\delta(\mathrm{ZG})\in(0,1).
\]

对第二条，任取有限素数集 \(S\subset\mathbb P\)。由于 \(\mathbb P\setminus S\) 非空，可取某个素数 \(p_j\notin S\)。令
\[
z_j=1,
\qquad
z_k=0\quad (k\neq j).
\]
则对应整数正是
\[
n=p_j\in \mathrm{ZG},
\]
因为它显然平方自由，且仅激活一个素数指标，不可能违反相邻约束；同时
\[
\operatorname{supp}(n)=\{p_j\}\subset \mathbb P\setminus S.
\]
故第二条成立。

第三条只是前两条的直接合并。再结合定理 \ref{thm:xi-localized-phase-sampling-closed-form}，任意固定有限 prime-stage 只恢复其自身有限素数支撑 \(S\)，因而不可能穷尽 \(\mathrm{ZG}\) 的共尾支撑。证毕。
\end{proof}

\paragraph{有限维 Lie 相位读出、dyadic solenoid 圆因子与零温四点骨架}
在整数共振梯的 Shannon 非可寻址性、有限二维 \(2\)-进 ambient space 的忠实实现分裂以及真实输入 \(40\) 状态核的镜像慢模已经确定之后，还可进一步抽取出下列五条后果。它们分别把有限维紧 Lie 相位上的 \(L^2\) Fourier 读出障碍、dyadic solenoid 圆因子的绝对连续排斥、显式原子 Witt 因子剥离下的慢混合主阶不变性，以及零温归一化谱中的四点骨架与双层分裂统一到同一条相位--谱层级链上。

\begin{theorem}[有限维 Lie 相位的 \texorpdfstring{$L^2$}{L2} 共振读出障碍]\label{thm:xi-time-part9w-lie-phase-l2-resonance-readout-obstruction}
设 \(G\) 为有限维紧阿贝尔 Lie 群，\(\widehat G\) 为其角色格，\(\gamma\in\widehat G\) 为无限阶角色。则不存在 \(f\in L^2(G)\) 使
\[
\widehat f(u\gamma)=C_\varphi(u)\qquad(u\in\ZZ).
\]
因而黄金 bulk 共振梯不可能是任何有限维紧 Lie 相位空间上有限能量可观测量的单频 Fourier 读出。
\end{theorem}
\begin{proof}
由 Plancherel 定理，
\[
\sum_{\chi\in\widehat G}|\widehat f(\chi)|^2=\|f\|_{L^2(G)}^2<\infty.
\]
由于 \(\gamma\) 具有无限阶，映射 \(u\mapsto u\gamma\) 在 \(\ZZ\) 上单射，从而
\[
\sum_{u\in\ZZ}|\widehat f(u\gamma)|^2\le \sum_{\chi\in\widehat G}|\widehat f(\chi)|^2<\infty.
\]
另一方面，定理 \ref{thm:fold-sigma-phi-diverges} 已给出
\[
\sum_{u\in\ZZ}C_\varphi(u)^2=+\infty.
\]
故若 \(\widehat f(u\gamma)=C_\varphi(u)\)，便与上式矛盾。证毕。
\end{proof}

\begin{corollary}[dyadic solenoid 圆因子的绝对连续障碍]\label{cor:xi-time-part9w-dyadic-solenoid-circle-factor-nonac}
记 \(\Sigma_{\{2\}}:=\widehat{\ZZ[1/2]}\)，并令
\[
0\longrightarrow \ZZ_2\longrightarrow \Sigma_{\{2\}}
\xrightarrow{\pi_{\{2\}}}\TT\longrightarrow 0
\]
为其标准圆商短正合列。则不存在 \(\Sigma_{\{2\}}\) 上的 Borel 概率测度 \(\nu\)，使得
\[
\eta:=\pi_{\{2\},*}\nu
\]
同时满足
\[
\widehat\eta(n)=C_\varphi(n)\qquad(n\in\ZZ)
\]
与 \(\eta\ll m_{\TT}\)。等价地，这样的圆因子测度不能写成任何 \(L^1(\TT)\) 密度；并且它也不可能来自任何 \(L^2(\TT)\) 密度。
\end{corollary}
\begin{proof}
若 \(\eta\ll m_{\TT}\)，则在有限测度空间 \(\TT\) 上可写成 \(d\eta=h\,dm_{\TT}\) 且 \(h\in L^1(\TT)\)。由 Riemann--Lebesgue 引理，
\[
\widehat\eta(n)\longrightarrow 0 \qquad (|n|\to\infty).
\]
但推论 \ref{cor:fold-resonance-ladder-fib-luc-limits} 与推论 \ref{cor:fold-resonance-ladder-no-decay-no-limit} 表明，\(C_\varphi(n)\) 沿 Fibonacci 与 Lucas 子列分别趋于两个不同的正常数，尤其 \(\limsup_{|n|\to\infty}C_\varphi(n)>0\)。故 \(\widehat\eta(n)=C_\varphi(n)\) 不可能成立。
若进一步有 \(h\in L^2(\TT)\)，则 Parseval 定理给出 \((\widehat\eta(n))_{n\in\ZZ}\in\ell^2(\ZZ)\)；而定理 \ref{thm:fold-sigma-phi-diverges} 表明 \((C_\varphi(n))_{n\in\ZZ}\notin\ell^2(\ZZ)\)。故 \(L^2\) 密度同样不可能。证毕。
\end{proof}

\begin{theorem}[原子 Witt 因子剥离下的慢混合主阶不变性]\label{thm:xi-time-part9w-atomic-witt-surgery-slowmixing-invariance}
沿用推论 \ref{cor:real-input-40-factor-sqrtu-eigs}、命题 \ref{prop:real-input-40-output-mirror-branch} 与定理 \ref{thm:xi-output-potential-endpoint-near-degeneracy-sqrtu-relaxation} 的记号。设
\[
\Delta(z,u)=\det(I-zM(u))=(1-uz^2)F(z,u),
\qquad
\zeta_{\mathrm{core}}(z,u):=F(z,u)^{-1}.
\]
记 \(\lambda_1(u)\) 为 Perron 根，\(\Lambda_2(u)\) 为全谱最大非主模特征值，\(\lambda_{\mathrm{mir}}(u)<0\) 为核心块的镜像慢模，并定义
\[
\tau_{\mathrm{full}}(u):=\left[-\log\!\left(\frac{\Lambda_2(u)}{\lambda_1(u)}\right)\right]^{-1},
\qquad
\tau_{\mathrm{core}}(u):=\left[-\log\!\left(\frac{|\lambda_{\mathrm{mir}}(u)|}{\lambda_1(u)}\right)\right]^{-1}.
\]
则
\[
\frac{\tau_{\mathrm{full}}(u)}{\tau_{\mathrm{core}}(u)}\longrightarrow 1
\qquad (u\to+\infty).
\]
因此剥离显式原子 Witt 因子 \((1-uz^2)^{-1}\) 以后，主导慢混合的首阶时间尺度保持不变。
\end{theorem}
\begin{proof}
由定理 \ref{thm:xi-output-potential-endpoint-near-degeneracy-sqrtu-relaxation}，
\[
1-\frac{\Lambda_2(u)}{\lambda_1(u)}=\frac{2d}{c}u^{-1/2}+O(u^{-1}).
\]
另一方面，命题 \ref{prop:real-input-40-output-mirror-branch} 给出
\[
1-\frac{|\lambda_{\mathrm{mir}}(u)|}{\lambda_1(u)}=\frac{2d}{c}u^{-1/2}+O(u^{-1}),
\qquad
\frac{\sqrt u}{\lambda_1(u)}\longrightarrow \frac1c<1.
\]
故显式原子分支 \(\pm\sqrt u\) 在归一化后始终严格位于镜像慢模之内，而
\[
-\log\!\left(\frac{\Lambda_2(u)}{\lambda_1(u)}\right)
\quad\text{与}\quad
-\log\!\left(\frac{|\lambda_{\mathrm{mir}}(u)|}{\lambda_1(u)}\right)
\]
具有相同的主阶展开 \(\frac{2d}{c}u^{-1/2}+O(u^{-1})\)。两边取倒数即得所示极限。证毕。
\end{proof}

\begin{theorem}[零温归一化谱的四点骨架]\label{thm:xi-time-part9w-zero-temp-normalized-spectrum-fourpoint-skeleton}
沿用命题 \ref{prop:real-input-40-output-mirror-branch} 的常数 \(c\)，并置
\[
b:=\frac1{c^2}\in(0,1).
\]
则 \(b\) 满足
\[
1-6b+9b^2-b^3-b^4=0.
\]
对真实输入 \(40\) 状态核，定义
\[
\mathrm{Spec}_{\mathrm{norm}}(u):=\left\{\frac{\lambda}{\lambda_1(u)}:\lambda\in\mathrm{Spec}(M(u))\right\}.
\]
若 \(\mathcal L\) 是 \(\mathrm{Spec}_{\mathrm{norm}}(u)\) 沿某子列 \(u\to+\infty\) 的 Hausdorff 极限点集，则
\[
\{1,-1,\sqrt b,-\sqrt b\}\subseteq \mathcal L.
\]
\end{theorem}
\begin{proof}
归一化 Perron 根恒等于 \(1\)，故 \(1\in\mathcal L\)。由命题 \ref{prop:real-input-40-output-mirror-branch}，
\[
\frac{\lambda_{\mathrm{mir}}(u)}{\lambda_1(u)}
=
-\left(1-\frac{2d}{c}u^{-1/2}+O(u^{-1})\right)
\longrightarrow -1,
\]
故 \(-1\in\mathcal L\)。再由推论 \ref{cor:real-input-40-factor-sqrtu-eigs}，谱中始终含有 \(\lambda_\pm(u)=\pm\sqrt u\)，而定理 \ref{thm:xi-time-part73-perron-vs-atomic-sqrtu-separation} 给出
\[
\frac{\sqrt u}{\lambda_1(u)}\longrightarrow \frac1c=\sqrt b.
\]
因此
\[
\frac{\lambda_\pm(u)}{\lambda_1(u)}\longrightarrow \pm\sqrt b,
\]
从而 \(\pm\sqrt b\in\mathcal L\)。证毕。
\end{proof}

\begin{corollary}[输运层与保护层的双层谱分裂]\label{cor:xi-time-part9w-transport-protection-two-layer-splitting}
沿用定理 \ref{thm:xi-time-part9w-zero-temp-normalized-spectrum-fourpoint-skeleton} 的记号。零温极限 \(u\to+\infty\) 中，归一化谱至少保留四点骨架
\[
\underbrace{\{1,-1\}}_{\text{输运层}}
\sqcup
\underbrace{\{\sqrt b,-\sqrt b\}}_{\text{保护层}},
\qquad 0<\sqrt b<1.
\]
其中外层对 \(\{\pm1\}\) 控制慢松弛，而内层对 \(\{\pm\sqrt b\}\) 由显式原子 Witt 因子给出。去除因子 \((1-uz^2)^{-1}\) 会抹去内层 \(\{\pm\sqrt b\}\)，但不改变首阶慢混合尺度；反之，仅保留内层而忽略镜像慢模，则无法控制 \(\tau_{\mathrm{full}}(u)\) 的主阶发散。
\end{corollary}
\begin{proof}
由 \(b=1/c^2\in(0,1)\) 可知 \(0<\sqrt b<1\)。四点同时存在由定理 \ref{thm:xi-time-part9w-zero-temp-normalized-spectrum-fourpoint-skeleton} 给出，而内层 \(\{\pm\sqrt b\}\) 正是显式特征值 \(\pm\sqrt u\) 归一化后的极限。定理 \ref{thm:xi-time-part9w-atomic-witt-surgery-slowmixing-invariance} 表明去除原子 Witt 因子后，慢混合主阶保持不变。另一方面，若只跟踪 \(\pm\sqrt b\)，则对应的归一化模始终固定为 \(\sqrt b<1\)，其关联时间尺度 \((-\log\sqrt b)^{-1}\) 保持有界；而定理 \ref{thm:xi-output-potential-endpoint-near-degeneracy-sqrtu-relaxation} 给出 \(\tau_{\mathrm{full}}(u)\asymp \sqrt u\to+\infty\)。故内层本身不能控制全谱慢松弛。证毕。
\end{proof}

\noindent
上述结果表明，本文对象中的连续 Lie 相位、\(2\)-进全息地址、Hurwitz 链 Jacobian/Prym 分裂、prime-language 共尾支撑与真实输入 \(40\) 状态核的零温谱学并不处于同一承载层级。有限二维 \(2\)-进 ambient space 足以忠实实现有限时深地址动力学，但黄金 bulk 共振既不能在任何有限维紧 Lie 相位上作为 \(L^2\) Fourier 读出出现，也不能经 dyadic solenoid 的圆因子下降为绝对连续相位密度；与此同时，零温慢混合由镜像慢模而非显式原子分支主导，而归一化谱则稳定留下 \(\{\pm1,\pm\sqrt b\}\) 的四点骨架。因而，在本文语义中，连通相位、原子保护、镜像输运与算术共振必须被视为彼此不可压平的不同层级，而任何试图把它们同时压缩到单一有限维 Lie 图景的方案都会丢失关键的证明复杂度与动力学信息。

\endinput

\paragraph{visible 条件期望的有限指数谱、对称溢价与 hidden sector 的零模约束}
在纤维群胚代数的中心值指数元、window-\(6\) 的 oracle 几何完全可解、boundary 奇偶子格在交换化中的规范直和实现，以及微态 Hilbert 空间的 visible/hidden 分解已经固定之后，还可把交换函数代数上的纤维平均条件期望组织为另一条有限尺度闭链。以下只调用 window-\(6\) 的纤维直方图 \(2:8,\ 3:4,\ 4:9\)、规范群分解 \(G_6^{\mathrm{bin}}\cong S_2^{\,8}\times S_3^{\,4}\times S_4^{\,9}\)、boundary/full center 的精确秩，以及离散容量闭式；由此可同时得到可见条件期望的有限指数谱、输出 KL 缺口的指数元公式、精确反演与异常维数的裂分，以及对称保持 completion 的刚性比特下界。该链条显示出一条清晰的分离：hidden sector 的可逆代价由最大纤维控制，而对称代价则严格更高。

\begin{theorem}[window-\texorpdfstring{$6$}{6} 可见条件期望的有限指数谱]\label{thm:xi-time-part9x-visible-condexp-finite-index-spectrum}
记
\[
\Omega_6:=\{0,1,\dots,63\},
\qquad
\pi:=\Fold_6^{\mathrm{bin}}:\Omega_6\to X_6,
\qquad
F_w:=\pi^{-1}(w),
\qquad
d(w):=|F_w|.
\]
再令
\[
A_6:=\ell^\infty(\Omega_6),
\qquad
B_6:=\pi^\ast \ell^\infty(X_6)\subset A_6.
\]
定义纤维平均映射
\[
(E_6f)(n):=\frac{1}{|F_{\pi(n)}|}\sum_{m\in F_{\pi(n)}}f(m)
\qquad (f\in A_6,\ n\in \Omega_6).
\]
则 \(E_6:A_6\to B_6\) 是忠实条件期望，并且满足：
\begin{enumerate}
  \item 其 Watatani 指数元为
  \[
  \mathrm{Ind}_{\mathrm W}(E_6)
  =
  \sum_{w\in X_6}d(w)\,\mathbf 1_{F_w},
  \qquad
  \sigma\!\bigl(\mathrm{Ind}_{\mathrm W}(E_6)\bigr)=\{2,3,4\};
  \]
  \item 其 Pimsner--Popa 指数为
  \[
  \mathrm{Ind}_{\mathrm{PP}}(E_6)=4;
  \]
  \item 对均匀微态迹
  \[
  \tau_\Omega(f):=\frac{1}{64}\sum_{n\in\Omega_6}f(n),
  \]
  有
  \[
  \tau_\Omega\!\bigl(\mathrm{Ind}_{\mathrm W}(E_6)\bigr)
  =
  \frac{1}{64}\sum_{w\in X_6}d(w)^2
  =
  \frac{53}{16},
  \]
  \[
  \tau_\Omega\!\bigl(\log \mathrm{Ind}_{\mathrm W}(E_6)\bigr)
  =
  \frac{1}{64}\sum_{w\in X_6}d(w)\log d(w)
  =
  \kappa_6
  =
  \frac{11}{8}\log 2+\frac{3}{16}\log 3.
  \]
\end{enumerate}
\end{theorem}
\begin{proof}
首先，\(E_6\) 显然保持单位元与正性，并且对 \(g\in B_6\) 有
\[
E_6(g)=g,
\]
因为 \(g\) 在每条纤维上常值。又由于 \(E_6(f)\) 只依赖于 \(\pi(n)\)，故 \(E_6(A_6)\subset B_6\)。若 \(f\ge 0\) 且 \(E_6(f)=0\)，则对每个 \(w\in X_6\) 都有
\[
\sum_{m\in F_w}f(m)=0.
\]
由于每一项都非负，遂得 \(f|_{F_w}=0\)；对全部 \(w\) 合并即 \(f=0\)。故 \(E_6\) 为忠实条件期望。

对每个 \(n\in \Omega_6\)，记
\[
d_n:=d(\pi(n)),
\qquad
u_n:=\sqrt{d_n}\,\delta_n\in A_6,
\]
其中 \(\delta_n\) 为点质量函数。对任意 \(f\in A_6\)，若 \(m\in \Omega_6\)，则
\[
E_6(u_nf)(m)
=
\frac{1}{d(\pi(m))}
\sum_{\ell\in F_{\pi(m)}}u_n(\ell)f(\ell).
\]
当 \(\pi(m)\neq \pi(n)\) 时，上式为 \(0\)；当 \(\pi(m)=\pi(n)\) 时，只有 \(\ell=n\) 一项非零，因此
\[
E_6(u_nf)(m)=d_n^{-1/2}f(n).
\]
于是
\[
u_n(m)\,E_6(u_nf)(m)=f(n)\delta_n(m),
\]
从而
\[
\sum_{n\in\Omega_6}u_nE_6(u_nf)=f.
\]
由于 \(A_6\) 交换，同理也有
\[
\sum_{n\in\Omega_6}E_6(fu_n)u_n=f.
\]
故 \((u_n)_{n\in\Omega_6}\) 构成 \(E_6\) 的一个 quasi-basis。因而
\[
\mathrm{Ind}_{\mathrm W}(E_6)
=
\sum_{n\in\Omega_6}u_nu_n^\ast
=
\sum_{n\in\Omega_6}d_n\delta_n
=
\sum_{w\in X_6}d(w)\,\mathbf 1_{F_w}.
\]
由 window-\(6\) 直方图 \(2:8,\ 3:4,\ 4:9\)，其谱恰为 \(\{2,3,4\}\)。

对任意 \(a\in A_6^+\) 与任意 \(n\in F_w\)，有
\[
E_6(a)(n)
=
\frac{1}{d(w)}\sum_{m\in F_w}a(m)
\ge
\frac{1}{d(w)}a(n)
\ge
\frac14 a(n).
\]
故 \(E_6(a)\ge \frac14 a\) 对一切 \(a\ge 0\) 成立，从而
\[
\mathrm{Ind}_{\mathrm{PP}}(E_6)\le 4.
\]
若 \(n\) 位于某条四点纤维上并取 \(a=\delta_n\)，则
\[
E_6(a)(n)=\frac14.
\]
因此任何小于 \(4\) 的常数都不能满足 Pimsner--Popa 不等式，故
\[
\mathrm{Ind}_{\mathrm{PP}}(E_6)=4.
\]

最后，把直方图代入即可得到
\[
\tau_\Omega\!\bigl(\mathrm{Ind}_{\mathrm W}(E_6)\bigr)
=
\frac{1}{64}\bigl(8\cdot 2^2+4\cdot 3^2+9\cdot 4^2\bigr)
=
\frac{212}{64}
=
\frac{53}{16},
\]
\[
\tau_\Omega\!\bigl(\log \mathrm{Ind}_{\mathrm W}(E_6)\bigr)
=
\frac{1}{64}\bigl(8\cdot 2\log 2+4\cdot 3\log 3+9\cdot 4\log 4\bigr)
=
\frac{11}{8}\log 2+\frac{3}{16}\log 3.
\]
这与定理 \ref{thm:xi-foldbin-fkdet-logvolume-kappa} 在 \(m=6\) 时给出的 \(\kappa_6\) 完全一致。证毕。
\end{proof}

\begin{theorem}[指数元的相对熵恒等式]\label{thm:xi-time-part9x-index-relative-entropy-identity}
\leanverified{paper\_xi\_time\_part9x\_index\_relative\_entropy\_identity}
记
\[
\mu_6:=\pi_\ast \tau_\Omega,
\qquad
u_6(w):=\frac{1}{21}
\quad (w\in X_6),
\]
其中 \(u_6\) 为 \(X_6\) 上的均匀分布。再定义
\[
\tau_X(g):=\frac{1}{21}\sum_{w\in X_6}g(w)
\qquad (g\in \ell^\infty(X_6)).
\]
则 \(\mu_6(w)=d(w)/64\)，并且有精确恒等式
\[
D_{\mathrm{KL}}(\mu_6\Vert u_6)
=
\tau_\Omega\!\bigl(\log \mathrm{Ind}_{\mathrm W}(E_6)\bigr)
-\log \tau_X\!\bigl(\mathrm{Ind}_{\mathrm W}(E_6)\bigr).
\]
更显式地，
\[
D_{\mathrm{KL}}(\mu_6\Vert u_6)
=
\kappa_6-\log\frac{64}{21}
=
\log\!\Bigl(\frac{21}{64}\,2^{11/8}3^{3/16}\Bigr)
=
0.044706531758946\ldots
\]
\end{theorem}
\begin{proof}
由定义，
\[
\mu_6(w)=\tau_\Omega(\mathbf 1_{F_w})=\frac{d(w)}{64}.
\]
因此
\[
\begin{aligned}
D_{\mathrm{KL}}(\mu_6\Vert u_6)
&=
\sum_{w\in X_6}\frac{d(w)}{64}
\log\frac{d(w)/64}{1/21}\\
&=
\frac{1}{64}\sum_{w\in X_6}d(w)\log d(w)-\log\frac{64}{21}.
\end{aligned}
\]
第一项正是定理 \ref{thm:xi-time-part9x-visible-condexp-finite-index-spectrum} 中的
\[
\tau_\Omega\!\bigl(\log \mathrm{Ind}_{\mathrm W}(E_6)\bigr)=\kappa_6.
\]
另一方面，
\[
\tau_X\!\bigl(\mathrm{Ind}_{\mathrm W}(E_6)\bigr)
=
\frac{1}{21}\sum_{w\in X_6}d(w)
=
\frac{64}{21},
\]
故
\[
\log \tau_X\!\bigl(\mathrm{Ind}_{\mathrm W}(E_6)\bigr)=\log\frac{64}{21}.
\]
合并即得所示恒等式。由于 \(\mu_6=p_6^{\mathrm{bin}}\)，这也与命题
\ref{prop:xi-window6-renyi-divergence-parity-charge-redundancy} 的 KL 闭式完全一致。证毕。
\end{proof}

\begin{corollary}[hidden multiplicity 的三重平均严格分离]\label{cor:xi-time-part9x-hidden-multiplicity-three-means-separation}
定义
\[
m_{\mathrm{geo}}
:=
\exp \tau_\Omega\!\bigl(\log \mathrm{Ind}_{\mathrm W}(E_6)\bigr)
=
2^{11/8}3^{3/16},
\]
\[
m_{\mathrm{quad}}
:=
\tau_\Omega\!\bigl(\mathrm{Ind}_{\mathrm W}(E_6)\bigr)
=
\frac{53}{16},
\qquad
m_{\max}:=\bigl\|\mathrm{Ind}_{\mathrm W}(E_6)\bigr\|_\infty=4.
\]
则有严格不等式
\[
2^{11/8}3^{3/16}
=
3.186959021689957\ldots
<
\frac{53}{16}
=
3.3125
<
4.
\]
\end{corollary}
\begin{proof}
第二个不等式是显然的。对第一个不等式，由定理
\ref{thm:xi-time-part9x-visible-condexp-finite-index-spectrum}，
\[
\sigma\!\bigl(\mathrm{Ind}_{\mathrm W}(E_6)\bigr)=\{2,3,4\},
\]
故 \(\mathrm{Ind}_{\mathrm W}(E_6)\) 非常数。于是由 \(\log\) 的严格凹性与 Jensen 不等式，
\[
\tau_\Omega\!\bigl(\log \mathrm{Ind}_{\mathrm W}(E_6)\bigr)
<
\log \tau_\Omega\!\bigl(\mathrm{Ind}_{\mathrm W}(E_6)\bigr).
\]
两边取指数便得
\[
\exp \tau_\Omega\!\bigl(\log \mathrm{Ind}_{\mathrm W}(E_6)\bigr)
<
\tau_\Omega\!\bigl(\mathrm{Ind}_{\mathrm W}(E_6)\bigr).
\]
代入前述显式值即得结论。证毕。
\end{proof}

\begin{theorem}[window-\texorpdfstring{$6$}{6} 的可逆容量阈值与二比特精确反演]\label{thm:xi-time-part9x-reversible-capacity-two-bit-inversion}
\leanverified{paper\_xi\_time\_part9x\_reversible\_capacity\_two\_bit\_inversion}
定义
\[
b_{\mathrm{inv}}
:=
\min\{B\in \ZZ_{\ge 0}:\ C_6(B)=64\},
\]
其中 \(C_6(B)\) 为离散 dyadic 预言机容量。则
\[
b_{\mathrm{inv}}
=
\left\lceil \log_2 \mathrm{Ind}_{\mathrm{PP}}(E_6)\right\rceil
=
2.
\]
更具体地，
\[
C_6(B)
=
8\min(2,2^B)+4\min(3,2^B)+9\min(4,2^B),
\]
从而
\[
C_6(0)=21,
\qquad
C_6(1)=42,
\qquad
C_6(B)=64
\quad (B\ge 2).
\]
对应的最优成功率
\[
\mathrm{Succ}_6^{\mathrm{bin}}(B):=\frac{C_6(B)}{64}
\]
满足
\[
\mathrm{Succ}_6^{\mathrm{bin}}(0)=\frac{21}{64},
\qquad
\mathrm{Succ}_6^{\mathrm{bin}}(1)=\frac{21}{32},
\qquad
\mathrm{Succ}_6^{\mathrm{bin}}(B)=1
\quad (B\ge 2).
\]
\end{theorem}
\begin{proof}
预言机容量闭式给出
\[
C_6(B)=\sum_{w\in X_6}\min(d(w),2^B).
\]
把直方图 \(2:8,\ 3:4,\ 4:9\) 代入，便得到
\[
C_6(B)
=
8\min(2,2^B)+4\min(3,2^B)+9\min(4,2^B).
\]
于是
\[
C_6(0)=8+4+9=21,
\]
\[
C_6(1)=8\cdot 2+4\cdot 2+9\cdot 2=42,
\]
\[
C_6(2)=8\cdot 2+4\cdot 3+9\cdot 4=64.
\]
又由于最大纤维大小为 \(4\)，当 \(B\ge 2\) 时有 \(2^B\ge 4\)，从而
\[
C_6(B)=\sum_{w\in X_6}d(w)=64.
\]
故 \(b_{\mathrm{inv}}=2\)。另一方面，定理
\ref{thm:xi-time-part9x-visible-condexp-finite-index-spectrum} 已证明
\[
\mathrm{Ind}_{\mathrm{PP}}(E_6)=4,
\]
因此
\[
b_{\mathrm{inv}}
=
\left\lceil \log_2 4\right\rceil
=
\left\lceil \log_2 \mathrm{Ind}_{\mathrm{PP}}(E_6)\right\rceil
=
2.
\]
最后，两边除以 \(64\) 即得成功率公式。证毕。
\end{proof}

\begin{theorem}[解码维数与异常维数的十九比特裂分]\label{thm:xi-time-part9x-decoding-anomaly-nineteen-bit-gap}
\leanverified{paper\_xi\_time\_part9x\_decoding\_anomaly\_nineteen\_bit\_gap}
定义
\[
b_{\mathrm{anom}}
:=
\dim_{\FF_2}\bigl(G_6^{\mathrm{bin}}\bigr)^{\mathrm{ab}}.
\]
则
\[
b_{\mathrm{anom}}=21,
\qquad
b_{\mathrm{inv}}=2,
\]
从而
\[
b_{\mathrm{anom}}-b_{\mathrm{inv}}=19,
\qquad
\frac{2^{b_{\mathrm{anom}}}}{2^{b_{\mathrm{inv}}}}=2^{19}.
\]
\end{theorem}
\begin{proof}
由定理 \ref{thm:xi-foldbin-gauge-representation-distribution-law-m6-spectrum}，
\[
\bigl(G_6^{\mathrm{bin}}\bigr)^{\mathrm{ab}}
\cong
(\ZZ/2\ZZ)^{21},
\]
故
\[
b_{\mathrm{anom}}=21.
\]
另一方面，定理 \ref{thm:xi-time-part9x-reversible-capacity-two-bit-inversion} 已给出
\[
b_{\mathrm{inv}}=2.
\]
代入即可得到
\[
b_{\mathrm{anom}}-b_{\mathrm{inv}}=19,
\qquad
\frac{2^{b_{\mathrm{anom}}}}{2^{b_{\mathrm{inv}}}}=2^{19}.
\]
证毕。
\end{proof}

\noindent
该裂分表明，window-\(6\) 的局部精确反演仅受最大纤维顶谱控制，而交换化异常则由全部 \(21\) 条纤维的符号坐标共同编码；因此二者属于不同层级的有限尺度不变量，不能彼此压缩为同一类局部标签。

\endinput

\paragraph{固定分辨率全息完备性、屏幕相对同调与 Artin 通道商擦除}
在固定分辨率 dyadic 体--界全息、部分屏幕观测的相对同调核、Zeckendorf 折叠的确定化歧义壳以及非阿贝尔 Peter--Weyl/Artin 分块都已建立之后，还可进一步抽取出下列七条后果。它们共同表明：困难并不首先来自 bulk 无法在有限层面被完整编码；真正拒绝统一有限压缩的，是部分屏幕留下的相对同调赤字，以及群商压缩所擦除的非阿贝尔 Artin 通道。

\begin{theorem}[固定分辨率 Stokes--Gödel 全息完备性的强形式]\label{thm:xi-time-part9x-fixed-resolution-stokes-godel-strong-completeness}
固定 \(m,n\ge 1\)，令 \(\mathcal U_{m,n}\) 为 \(I^n=[0,1]^n\) 上全部分辨率 \(m\) 的 dyadic polycube，\(\mathcal Q_m\) 为对应的 \(m\)-级 dyadic 立方复形。对每个 \(U\in\mathcal U_{m,n}\)，记 \(u_U\in C_n(\mathcal Q_m;\FF_2)\) 为其占据链，并令
\[
\Phi(U):=(G\circ\partial_n)(u_U),
\]
其中 \(G\) 为任意在像上可计算可逆的平方自由 boundary Gödel 整数化。则存在常数 \(C>0\)（仅依赖于所选通用前缀机与编码方案）使得对一切 \(U\in\mathcal U_{m,n}\)，都有
\[
\bigl|K(U\mid m,n)-K(\Phi(U)\mid m,n)\bigr|\le C,
\qquad
K(U\mid \Phi(U),m,n)\le C.
\]
此外，\(\Phi(U)\) 唯一确定 \(U\)，并且唯一确定有限矩盒
\[
\{\mu_\alpha(U):0\le \alpha_i\le 2^m-1\}.
\]
因而在固定分辨率上，boundary Gödel code 是复杂度守恒且信息完备的全息坐标。
\end{theorem}
\begin{proof}
在给定 \((m,n)\) 时，占据链 \(u_U\) 与 polycube \(U\) 之间存在可计算双向转换，故
\[
K(U\mid m,n)=K(u_U\mid m,n)+O(1).
\]
由定理 \ref{thm:spg-stokes-godel-algorithmic-holographic-completeness}，
\[
\bigl|K(u_U\mid m,n)-K((G\circ\partial_n)(u_U)\mid m,n)\bigr|\le O(1),
\qquad
K(u_U\mid \Phi(U),m,n)\le O(1).
\]
把 \(U\leftrightarrow u_U\) 的编码常数吸收进同一 \(O(1)\) 项，即得复杂度结论。

另一方面，推论 \ref{cor:spg-boundary-godel-finite-moment-completeness} 已表明 \(\Phi(U)\) 唯一决定 \(U\)。更具体地，由定理 \ref{thm:spg-boundary-godel-moment-readout}，\(\Phi(U)\) 决定全部单项式矩；再由定理 \ref{thm:spg-dyadic-finite-moment-completeness}，其中有限矩盒已足以唯一恢复 \(U\)。证毕。
\end{proof}

\begin{theorem}[部分屏幕观测的相对同调赤字闭式]\label{thm:xi-time-part9x-partial-screen-relative-homology-deficit-closed}
沿用定理 \ref{thm:spg-partial-boundary-screen-kernel-relative-homology} 的记号。对任意屏幕
\[
S\subset \overrightarrow{\mathcal F_{m}^{n-1}},
\]
其部分边界观测
\[
f_S:=\Pi_S\circ\partial_n:\ C_n(\mathcal Q_m;\ZZ)\longrightarrow \ZZ^S
\]
满足自然同构
\[
\ker f_S\cong H_n(\mathcal Q_m,A_{\neg S};\ZZ)\cong \ZZ^{\,c(\Gamma_S)-1},
\]
其中 \(\Gamma_S\) 为屏幕图，\(c(\Gamma_S)\) 为其连通分量数。特别地，
\[
\rank_{\ZZ}\ker f_S=c(\Gamma_S)-1.
\]
\end{theorem}
\begin{proof}
前一同构由定理 \ref{thm:spg-partial-boundary-screen-kernel-relative-homology} 给出，后一同构由定理 \ref{thm:spg-screen-kernel-connected-components} 给出。复合即得结论。证毕。
\end{proof}
\leanverified{paper\_xi\_time\_part9x\_partial\_screen\_relative\_homology\_deficit\_closed}

\begin{corollary}[屏幕补全成本的面积律 sharp 化]\label{cor:xi-time-part9x-screen-completion-cost-sharp-area-law}
沿用定理 \ref{thm:xi-time-part9x-partial-screen-relative-homology-deficit-closed} 的记号，定义最小审计补全成本
\[
\mathsf{Cost}(S):=\min_{(R,r_S)}\cdim(R)=\rank_{\ZZ}\ker f_S.
\]
则
\[
\mathsf{Cost}(S)=c(\Gamma_S)-1.
\]
特别地，对轴向屏幕 \(S^{(1)}\)（命题 \ref{prop:spg-axial-screen-area-law-audit-cost} 的记号），有精确公式
\[
\mathsf{Cost}(S^{(1)})=2^{m(n-1)}.
\]
\end{corollary}
\begin{proof}
由推论 \ref{cor:spg-screen-injectivity-and-audit-cost-components}，
\[
\mathsf{Cost}(S)=\rank_{\ZZ}\ker f_S=c(\Gamma_S)-1.
\]
对轴向屏幕 \(S^{(1)}\)，命题 \ref{prop:spg-axial-screen-area-law-audit-cost} 给出
\[
c(\Gamma_{S^{(1)}})=2^{m(n-1)}+1.
\]
代入第一式即可。证毕。
\end{proof}

\begin{theorem}[瞬态壳与周期谱的严格正交]\label{thm:xi-time-part9x-transient-shell-periodic-spectrum-orthogonality}
对 Zeckendorf 折叠的右确定化表示 \(H_m^{\det}\)（\(m\ge 3\)），按“歧义态在前、单点态在后”的顺序排列状态后，其邻接矩阵可写为
\[
A_m^{\det}
=
\begin{pmatrix}
N_m & R_m\\
0 & A_m^{\mathrm{ess}}
\end{pmatrix},
\]
其中 \(N_m\) 为 ambiguity shell 块，\(A_m^{\mathrm{ess}}\) 为 singleton skeleton 块。则成立：
\[
N_m^{\,m}=0,
\qquad
\det(I-zA_m^{\det})=\det(I-zA_m^{\mathrm{ess}}),
\qquad
D_{\mathrm{sync}}(m)=m-1.
\]
因此歧义壳既不承载任何周期轨道，也不改变表示图层面的 \(\zeta\)-因子；它只以有限同步延迟进入时间复杂度。

更一般地，在一般 sofic 口径下，若歧义壳子系统 \(Y_m^{\mathrm{amb}}\neq\varnothing\) 且其最大本质分量的周期为 \(p\ge 1\)，则存在常数族
\[
\{c_r(m)\}_{r=0}^{p-1}\subset(0,\infty)
\]
使得
\[
\nu_m(T_m>n)=c_{n\bmod p}(m)e^{-\varepsilon_m n}(1+o(1)),
\qquad
\varepsilon_m:=\log 2-\log\rho_{\mathrm{amb}}(m).
\]
\end{theorem}
\begin{proof}
前两式由定理 \ref{thm:Ym-ambiguity-shell-zeta-factor-trivial} 直接给出，其中 \(N_m\) 幂零且 \(\det(I-zN_m)=1\)。同步延迟闭式由推论 \ref{cor:Ym-sync-delay} 给出。最后一式正是定理 \ref{thm:Ym-sync-tail-pf-sharp} 的结论。故周期/Fredholm 数据与同步预算在结构上正交：前者只感知单点骨架，后者则由歧义壳的瞬态或低熵尾主导。证毕。
\end{proof}

\begin{theorem}[非阿贝尔 Artin 通道在商压缩下的 determinant--energy 双擦除]\label{thm:xi-time-part9x-nonabelian-artin-channel-quotient-erasure}
沿用命题 \ref{prop:kernel-peter-weyl-block-diagonalization} 与命题 \ref{prop:kernel-artin-factorization} 的记号。设 \(G\) 为有限群，\(K\in\CC[G]^{V\times V}\) 为群标记核，\(\mathcal K\in\End(\CC^V\otimes\CC[G])\) 为其右卷积实现，并记
\[
\mathcal K\cong \bigoplus_{\rho\in\widehat G}\bigl(K_\rho\otimes I_{d_\rho}\bigr),
\qquad
\det(I-z\mathcal K)=\prod_{\rho\in\widehat G}\det(I-zK_\rho)^{d_\rho},
\]
其中 \(d_\rho:=\dim\rho\)。对任意满同态
\[
\pi:G\twoheadrightarrow H,
\]
令 \(\pi_*K\in\CC[H]^{V\times V}\) 为逐矩阵项推前的商核，\(\mathcal K_H\) 为其右卷积实现，并定义
\[
\mathrm{Inf}_\pi(\widehat H):=\{\sigma\circ\pi:\sigma\in\widehat H\}\subseteq \widehat G.
\]
则对每个 \(\sigma\in\widehat H\) 都有
\[
(\pi_*K)_\sigma=K_{\sigma\circ\pi},
\]
从而
\[
\det(I-z\mathcal K_H)=\prod_{\sigma\in\widehat H}\det(I-zK_{\sigma\circ\pi})^{d_\sigma},
\]
其中 \(d_\sigma:=\dim\sigma\)。特别地，一切不属于 \(\mathrm{Inf}_\pi(\widehat H)\) 的 irreducible Artin channel 都从商核的 determinant 因子中完全消失。

此外，对任意零均值类函数 \(\varepsilon\in\mathrm{Cl}(G)\)，其纤维平均推前满足精确能量损失公式
\[
\|\varepsilon\|_{L^2(G)}^2-\|\pi_*\varepsilon\|_{L^2(H)}^2
=
\frac1{|G|^2}
\sum_{\rho\in\widehat G\setminus\mathrm{Inf}_\pi(\widehat H)}
\bigl|B_G(\rho)\bigr|^2.
\]
因此群商压缩会把所有非 inflation 通道同时从 determinant 数据与非阿贝尔能量账本中不可逆擦除。
\end{theorem}
\begin{proof}
对任意群代数元素
\[
a=\sum_{g\in G}a_g\,g\in\CC[G]
\]
及任意 \(\sigma\in\widehat H\)，逐项推前定义给出
\[
\sum_{h\in H}(\pi_*a)_h\,\sigma(h)
=
\sum_{h\in H}\left(\sum_{\pi(g)=h}a_g\right)\sigma(h)
=
\sum_{g\in G}a_g\,\sigma(\pi(g))
=
\sum_{g\in G}a_g\,(\sigma\circ\pi)(g).
\]
逐矩阵项应用此恒等式，即得
\[
(\pi_*K)_\sigma=K_{\sigma\circ\pi}.
\]
再对 \(\pi_*K\) 应用命题 \ref{prop:kernel-artin-factorization}，便得到所示 determinant 分解；其中只出现 inflated 通道，因此所有 \(\rho\notin\mathrm{Inf}_\pi(\widehat H)\) 的因子都在商核侧消失。

能量损失公式则正是定理 \ref{thm:xi-time-part64-nonabelian-quotient-energy-loss-exact}。最后，由 Peter--Weyl 直和分解可知：若固定全部 inflated 通道而改变某个 \(\rho\notin\mathrm{Inf}_\pi(\widehat H)\) 的分块 \(K_\rho\)，则 \(\mathcal K_H\) 保持不变，而全 determinant 因子化与上式右端一般随之改变。故这些通道无法由商压缩后的数据重建。证毕。
\end{proof}

\begin{corollary}[阿贝尔化证书的严格不足性]\label{cor:xi-time-part9x-abelianized-certificate-strict-insufficiency}
令
\[
q:G\twoheadrightarrow G^{\mathrm{ab}}:=G/[G,G]
\]
为阿贝尔化商，并记 \(\mathcal K\) 为 \(K\) 的右卷积实现。若存在某个 \(\rho\in\widehat G\) 满足
\[
\dim\rho>1,
\qquad
\det(I-zK_\rho)\not\equiv 1,
\]
则任何仅依赖于 abelianized labels 的 Dirichlet/Artin 证书都不可能恢复完整 determinant
\[
\det(I-z\mathcal K),
\]
也不可能恢复由全部 Artin 因子共同编码的 twisted primitive-orbit law。
\end{corollary}
\begin{proof}
对阿贝尔群 \(G^{\mathrm{ab}}\)，全部不可约表示均为一维表示。因此 \(\mathrm{Inf}_q(\widehat{G^{\mathrm{ab}}})\) 恰好是 \(G\) 的全部一维不可约表示。由定理 \ref{thm:xi-time-part9x-nonabelian-artin-channel-quotient-erasure}，所有满足 \(\dim\rho>1\) 的通道都会从 \(q_*K\) 的 determinant 因子化中消失。若其中某个通道的 Artin 因子 \(\det(I-zK_\rho)\) 非平凡，则它不可能由 abelianized system 重建；故完整 determinant 与相应的 twisted primitive-orbit 统计也都无法恢复。证毕。
\end{proof}

\begin{conjecture}[classical RH 证书的 screen--Artin 双障碍]\label{conj:xi-time-part9x-rh-screen-artin-double-obstruction}
在本文的 solenoidal--Mellin--Fredholm 语义中，任何可审计的 classical RH 终端证书若要在共尾分辨率上保持完备，必须同时保留两类对象：
\[
\mathfrak H_{\mathrm{scr}}
:=
\{\,H_n(\mathcal Q_m,A_{\neg S_m};\ZZ)\,\}_{m\ge 1},
\qquad
\mathfrak H_{\mathrm{Art}}
:=
\{\,K_{\rho,m}:\rho\in\widehat G_m\setminus\mathrm{Inf}_{\pi_m}(\widehat H_m)\,\}_{m\ge 1}.
\]
其中 \(K_{\rho,m}\) 表示第 \(m\) 级群标记核的 \(\rho\)-通道块，而 \(\pi_m:G_m\twoheadrightarrow H_m\) 表示相应的商压缩。
更具体地说：
\begin{enumerate}
  \item 任何试图把 bulk 证书压缩到一族部分屏幕 \(S_m\) 且在共尾分辨率上保持统一有界补全成本的方案，只要 \(\rank H_n(\mathcal Q_m,A_{\neg S_m};\ZZ)\) 不最终消失，就必在某一级失去相对同调核信息；
  \item 任何试图把群标签统一压缩到固定有限商或纯阿贝尔化的方案，只要仍存在非 inflation 的活跃 Artin 通道，就必丢失相应的 determinant 因子与通道能量；
  \item 因而，同时施行这两类 finite-horizon 压缩的正规化，不应能够闭合 classical RH 所需的完整终端证书链。
\end{enumerate}
\end{conjecture}

\noindent
上述结果表明，本文对象中的困难并不首先来自“无限维对象不可构造”。在固定分辨率上，bulk 甚至可以由单个 boundary Gödel 整数复杂度守恒地编码并唯一恢复。真正不可统一有限化的，是部分 screen 观测留下的相对同调赤字，以及群商压缩所抹去的非阿贝尔 Artin determinant 因子。因而，更合适的终端障碍语言应当同时记录 screen completion cost 与 nonabelian channel loss；在此意义下，screen--Artin obstruction theory 比单纯的“有限见证不足”命题更接近本文对象的真实不可压缩核心。

\endinput

\paragraph{各向异性 Serrin 射线的 Fibonacci 同余骨架、局部化账本与六倍窗边界阈值}
在 Figalli--Zhang 对 ADR 与全局 \(\beta\)-square-function 假设下的各向异性 Serrin 刚性已经把可实现粗糙域压缩为 Wulff 形的平移与伸缩，并且折叠同余的模半环表示、window-\(6\) 的 \(18\oplus 3\) 刚性分裂、boundary sheet-flip 的规范中心、容量曲线与纤维直方图的完备统计、局部化数系的 Pontryagin 对偶分解以及六倍窗零点块下界都已建立之后，还可进一步得到如下七码接口结果。它们分别刻画未归一化 Serrin 射线的有限分辨率模半环、该射线上的精确容量律、局部化 Minkowski 射线的圆因子--素数账本分解、window-\(6\) boundary 三点集的算术钉扎、boundary 扇区的 Weyl 完成阈值、六倍窗算术零点块在 Serrin 射线骨架上的传送，以及单参数各向异性 Serrin 观测的 Fibonacci 同余普适性。以下命题均是对前述结果的推导性转写。 \cite{FigalliZhang2026AnisotropicSerrinGeneralDomains}

\begin{theorem}[未归一化各向异性 Serrin 射线的有限分辨率模半环与 Fibonacci 同余完备化]\label{thm:xi-time-part9x-serrin-ray-fibonacci-semirings}
固定满足 Figalli--Zhang 假设的均匀凸各向异性体 \(K\)，记对应 Wulff 形为 \(W_K\)。模去平移并固定重心于原点后，添入形式零元
\[
0\cdot W_K:=\{0\},
\]
并定义整数尺度射线
\[
\mathfrak W_K^{(0)}:=\{\rho W_K:\ \rho\in \NN_0\}.
\]
在其上定义两种运算
\[
(\rho W_K)\boxplus(\sigma W_K):=(\rho+\sigma)W_K,
\qquad
(\rho W_K)\boxtimes(\sigma W_K):=(\rho\sigma)W_K.
\]
对每个 \(m\ge 1\)，再定义同余
\[
\rho W_K\equiv_m^{(K)}\sigma W_K
\iff
\rho\equiv \sigma \pmod{F_{m+2}}.
\]
则映射
\[
[\rho W_K]_{\equiv_m^{(K)}}\longmapsto [\rho]_{F_{m+2}}
\]
给出规范半环同构
\[
\bigl(\mathfrak W_K^{(0)}/\!\equiv_m^{(K)},\boxplus,\boxtimes\bigr)
\cong
\bigl(\ZZ/(F_{m+2}\ZZ),+,\times\bigr).
\]

再记 \(\mathfrak W_K^{\mathrm{gp}}\) 为 \((\mathfrak W_K^{(0)},\boxplus)\) 的 Grothendieck 群完成。则
\[
\mathfrak W_K^{\mathrm{gp}}\cong \ZZ.
\]
若赋予 \(\mathfrak W_K^{\mathrm{gp}}\) 由子群
\[
F_{m+2}\,\mathfrak W_K^{\mathrm{gp}}
\qquad(m\ge 1)
\]
生成的 Fibonacci 同余拓扑，则其完备化满足规范拓扑群同构
\[
\widehat{\mathfrak W}_{K,F}^{\mathrm{gp}}
\cong
\widehat{\ZZ}_F
\cong
\widehat{\ZZ}.
\]
\end{theorem}
\begin{proof}
由 Figalli--Zhang 的主定理，模去平移并固定重心后，一切可实现各向异性 Serrin 域恰为单参数族
\[
\{\rho W_K:\ \rho>0\}.
\]
对同心 Wulff 形，Minkowski 和满足
\[
\rho W_K+\sigma W_K=(\rho+\sigma)W_K,
\]
而尺度复合显然给出
\[
\rho(\sigma W_K)=(\rho\sigma)W_K.
\]
故映射
\[
\iota:\mathfrak W_K^{(0)}\longrightarrow \NN_0,
\qquad
\iota(\rho W_K)=\rho
\]
把 \((\mathfrak W_K^{(0)},\boxplus,\boxtimes)\) 识别为 \((\NN_0,+,\times)\)。

于是对每个 \(m\ge 1\)，同余 \(\equiv_m^{(K)}\) 正是模 \(F_{m+2}\) 同余在 \(\NN_0\) 上的拉回，故
\[
[\rho W_K]_{\equiv_m^{(K)}}\longmapsto [\rho]_{F_{m+2}}
\]
立即给出所述半环同构。

另一方面，\((\NN_0,+)\) 的 Grothendieck 群完成是 \(\ZZ\)，故
\[
\mathfrak W_K^{\mathrm{gp}}\cong \ZZ.
\]
再把由 \(F_{m+2}\mathfrak W_K^{\mathrm{gp}}\) 生成的拓扑经上述同构输运到 \(\ZZ\)，便恰得到定理 \ref{thm:xi-visible-arithmetic-fibonacci-adic-profinite-coincidence} 中的 Fibonacci 同余拓扑 \(\tau_F\)。因此其完备化满足
\[
\widehat{\mathfrak W}_{K,F}^{\mathrm{gp}}
\cong
\widehat{\ZZ}_F
\cong
\widehat{\ZZ}.
\]
证毕。
\end{proof}

\begin{corollary}[Wulff 射线的精确预言机容量律与恢复位率常数]\label{cor:xi-time-part9x-serrin-ray-exact-capacity}
在定理 \ref{thm:xi-time-part9x-serrin-ray-fibonacci-semirings} 的设定下，固定分辨率 \(m\ge 1\)，并只观察尺度窗口
\[
0\le \rho<2^m.
\]
定义残余类读出
\[
r_m(\rho W_K):=[\rho]_{F_{m+2}}
\in
\ZZ/(F_{m+2}\ZZ),
\]
以及对应纤维基数
\[
d_m^{\mathrm{ray}}(r)
:=
\#\{\,0\le \rho<2^m:\ [\rho]_{F_{m+2}}=r\,\}.
\]
再记
\[
q_m:=\left\lfloor\frac{2^m}{F_{m+2}}\right\rfloor,
\qquad
a_m:=2^m-q_mF_{m+2}.
\]
则对每个 \(r\in \ZZ/(F_{m+2}\ZZ)\) 都有
\[
d_m^{\mathrm{ray}}(r)\in\{q_m,q_m+1\},
\]
且恰有 \(a_m\) 条纤维取值 \(q_m+1\)，其余 \(F_{m+2}-a_m\) 条纤维取值 \(q_m\)。

定义该射线上的 dyadic 预言机容量
\[
C_m^{\mathrm{ray}}(B)
:=
\sum_{r\in \ZZ/(F_{m+2}\ZZ)}
\min\bigl(d_m^{\mathrm{ray}}(r),2^B\bigr),
\qquad
B\ge 0.
\]
则有严格闭式
\[
C_m^{\mathrm{ray}}(B)
=
a_m\min(q_m+1,2^B)
+
\bigl(F_{m+2}-a_m\bigr)\min(q_m,2^B).
\]
相应最优成功率为
\[
\mathrm{Succ}_m^{\mathrm{ray}}(B)
=
\frac{C_m^{\mathrm{ray}}(B)}{2^m}.
\]
若要求在窗口 \(0\le \rho<2^m\) 内精确恢复全部尺度，则所需最小附加比特数满足
\[
B_m^{\mathrm{exact}}
=
\left\lceil \log_2(q_m+1)\right\rceil.
\]
并且存在极限
\[
\lim_{m\to\infty}\frac{B_m^{\mathrm{exact}}}{m}
=
\log_2\!\Bigl(\frac{2}{\varphi}\Bigr).
\]
\end{corollary}
\begin{proof}
将 \(2^m\) 作 Euclid 分解
\[
2^m=q_mF_{m+2}+a_m,
\qquad
0\le a_m<F_{m+2}.
\]
于是模 \(F_{m+2}\) 的每个剩余类在区间 \([0,2^m)\) 中出现 \(q_m\) 次或 \(q_m+1\) 次；并且恰有 \(a_m\) 个剩余类多出一次。这就得到纤维直方图的二值结构。

将该直方图代入与定理 \ref{thm:xi-capacity-tail-histogram-moment-complete-statistic} 相同的容量泛函
\[
C(B)=\sum_x\min(d(x),2^B),
\]
便立得所示闭式。

若要求在给定残余类之外再加 \(B\) 比特以恢复窗口内全部尺度，则每条纤维内部至多可用 \(2^B\) 个不同附加标签。故若
\[
2^B<q_m+1,
\]
则某条最大纤维中必有两个尺度无法区分，精确恢复不可能。反之，只要
\[
2^B\ge q_m+1,
\]
就可在每条纤维内部赋予互异附加标签，从而实现精确恢复。于是
\[
B_m^{\mathrm{exact}}
=
\left\lceil \log_2(q_m+1)\right\rceil.
\]

最后，由 Binet 公式
\[
F_{m+2}\sim \frac{\varphi^{m+2}}{\sqrt5}
\]
得到
\[
q_m
\sim
\frac{2^m}{F_{m+2}}
\sim
\frac{\sqrt5}{\varphi^2}
\Bigl(\frac{2}{\varphi}\Bigr)^m.
\]
故
\[
\log_2(q_m+1)
=
m\log_2\!\Bigl(\frac{2}{\varphi}\Bigr)+O(1),
\]
从而
\[
\frac{B_m^{\mathrm{exact}}}{m}\longrightarrow \log_2\!\Bigl(\frac{2}{\varphi}\Bigr).
\]
证毕。
\end{proof}

\begin{theorem}[局部化 Wulff Minkowski 射线的圆因子--素数账本分解]\label{thm:xi-time-part9x-serrin-ray-localized-ledger}
取有限素数集 \(S\subset\mathbb P\)，并定义同心局部化尺度半群
\[
\mathfrak W_{K,S}^{+}
:=
\{\rho W_K:\ \rho\in \ZZ[S^{-1}]_{\ge 0}\},
\]
其运算取为 Minkowski 和 \(\boxplus\)。记
\[
\mathfrak W_{K,S}^{\mathrm{gp}}
\]
为其 Grothendieck 群完成。则存在规范群同构
\[
\mathfrak W_{K,S}^{\mathrm{gp}}\cong \ZZ[S^{-1}].
\]
因此其 Pontryagin 对偶满足短正合列
\[
0\longrightarrow \prod_{p\in S}\ZZ_p
\longrightarrow
\widehat{\mathfrak W_{K,S}^{\mathrm{gp}}}
\longrightarrow
\TT
\longrightarrow 0.
\]
并且
\[
\cdim_{\mathrm{fr}}\bigl(\mathfrak W_{K,S}^{\mathrm{gp}}\bigr)=1,
\]
\[
p\in S
\iff
\text{核 }\prod_{q\in S}\ZZ_q\text{ 的 \(p\)-主因子非平凡}.
\]
\end{theorem}
\begin{proof}
对同心 Wulff 形有
\[
\rho W_K\boxplus \sigma W_K=(\rho+\sigma)W_K.
\]
因此映射
\[
\iota_S:\mathfrak W_{K,S}^{+}\longrightarrow \ZZ[S^{-1}]_{\ge 0},
\qquad
\iota_S(\rho W_K)=\rho
\]
是加法半群同构，进而其 Grothendieck 群完成满足
\[
\mathfrak W_{K,S}^{\mathrm{gp}}\cong \ZZ[S^{-1}].
\]

再由定理 \ref{thm:killo-localization-circle-dimension-prime-ledger-splitting} 作用于 \(G_S=\ZZ[S^{-1}]\)，立得 Pontryagin 对偶短正合列
\[
0\to \prod_{p\in S}\ZZ_p\to \widehat{G_S}\to \TT\to 0,
\]
以及
\[
\cdim_{\mathrm{fr}}(G_S)=1,
\qquad
p\in S
\iff
\text{核的 \(p\)-主因子非平凡}.
\]
经由 \(\mathfrak W_{K,S}^{\mathrm{gp}}\cong G_S\) 传回即得结论。证毕。
\end{proof}

\begin{proposition}[window-\texorpdfstring{$6$}{6} boundary 三点集的算术钉扎与非陪集性]\label{prop:xi-time-part9x-window6-boundary-arithmetic-pinning}
在定理 \ref{thm:xi-window6-cyclic-kernel-bc3-root-datum} 的分裂
\[
X_6=X_6^{\mathrm{cyc}}\sqcup X_6^{\mathrm{bdry}},
\qquad
|X_6|=21,\quad |X_6^{\mathrm{cyc}}|=18,\quad |X_6^{\mathrm{bdry}}|=3
\]
下，boundary 扇区满足
\[
X_6^{\mathrm{bdry}}
=
\{\texttt{100001},\texttt{100101},\texttt{101001}\}.
\]
通过定理 \ref{thm:fold-congruence-mod-semirings} 的同构
\[
\Omega_6/\!\equiv_6\cong \ZZ/(F_8\ZZ)=\ZZ/(21\ZZ),
\]
该三点集精确对应于剩余类集合
\[
\mathcal B_6=\{14,17,19\}\subset \ZZ/(21\ZZ).
\]
此外：
\[
\mathcal B_6\ \text{不是}\ \ZZ/(21\ZZ)\ \text{中的加法陪集};
\]
\[
\mathcal B_6\ \text{也不是单位群}\ (\ZZ/(21\ZZ))^\times\ \text{中任何乘法子群的陪集}.
\]
并且，对每个 \(w\in X_6^{\mathrm{bdry}}\)，其 dyadic 二层预像都是一个两点纤维，且两点差值恒为
\[
34=F_9.
\]
\end{proposition}
\begin{proof}
对稳定词 \(w=(w_1,\dots,w_6)\in X_6\)，其值映射为
\[
V_6(w):=\sum_{k=1}^{6}w_kF_{k+1}.
\]
故
\[
V_6(\texttt{100001})=F_2+F_7=1+13=14,
\]
\[
V_6(\texttt{101001})=F_2+F_4+F_7=1+3+13=17,
\]
\[
V_6(\texttt{100101})=F_2+F_5+F_7=1+5+13=19.
\]
于是
\[
\mathcal B_6=\{14,17,19\}.
\]

若 \(\mathcal B_6\) 是 \(\ZZ/(21\ZZ)\) 中的加法陪集，则因其基数为 \(3\)，必为唯一三阶加法子群
\[
\{0,7,14\}
\]
的某个陪集，因而任意两点差都应为 \(7\) 的倍数。但
\[
17-14=3,\qquad 19-17=2,\qquad 19-14=5,
\]
矛盾。故 \(\mathcal B_6\) 不是加法陪集。

另一方面，任何 \((\ZZ/(21\ZZ))^\times\) 中乘法子群的陪集都只由单位组成；但
\[
\gcd(14,21)=7\neq 1,
\]
故 \(14\notin(\ZZ/(21\ZZ))^\times\)。因此 \(\mathcal B_6\) 不可能是单位群中任何乘法子群的陪集。

最后，推论 \ref{cor:fold-bin-boundary-center-z2} 已给出：对每个
\[
w\in X_6^{\mathrm{bdry}},
\]
纤维 \((\Fold_6^{\mathrm{bin}})^{-1}(w)\) 都是二点集，并且两点差值固定为
\[
34=F_9.
\]
证毕。
\end{proof}

\begin{theorem}[boundary 扇区的 Weyl 完成阈值]\label{thm:xi-time-part9x-window6-boundary-weyl-threshold}
\leanverified{paper\_xi\_time\_part9x\_window6\_boundary\_weyl\_threshold}
称集合 \(Y\subset X_6\) 为 boundary 扇区的一个 Weyl 完成，若
\[
X_6^{\mathrm{bdry}}\subset Y
\]
且 \(Y\cap X_6^{\mathrm{cyc}}\) 在定理 \ref{thm:xi-time-part9r-window6-hamming-complete-weyl-orbit-invariant} 的 \(B_3/C_3\) 根字典下对 Weyl 群作用不变。则必有
\[
|Y|\in\{3,9,15,21\}.
\]
特别地，任何非平凡 Weyl 完成都满足
\[
|Y|\ge 9.
\]
\end{theorem}
\begin{proof}
由定理 \ref{thm:xi-time-part9r-window6-hamming-complete-weyl-orbit-invariant}，\(X_6^{\mathrm{cyc}}\) 在 \(B_3/C_3\) 两张字典下都恰分裂为两条 Weyl 轨道，其基数分别为
\[
6,\qquad 12.
\]
因此 \(X_6^{\mathrm{cyc}}\) 的 Weyl 不变子集只能是
\[
\varnothing,
\qquad
\text{六元轨道},
\qquad
\text{十二元轨道},
\qquad
X_6^{\mathrm{cyc}},
\]
其基数依次为
\[
0,\ 6,\ 12,\ 18.
\]

若 \(Y\) 含全部 boundary 三点，而 \(Y\cap X_6^{\mathrm{cyc}}\) Weyl 不变，则
\[
|Y|
=
|X_6^{\mathrm{bdry}}|+|Y\cap X_6^{\mathrm{cyc}}|
=
3+|Y\cap X_6^{\mathrm{cyc}}|
\in
\{3,9,15,21\}.
\]
若要求 completion 非平凡，即
\[
Y\cap X_6^{\mathrm{cyc}}\neq\varnothing,
\]
则最小可能值对应于加入六元轨道，故
\[
|Y|\ge 3+6=9.
\]
证毕。
\end{proof}

\begin{theorem}[六倍窗上 Wulff 射线算术骨架的半阶零点块]\label{thm:xi-time-part9x-serrin-ray-arithmetic-zero-block}
记
\[
\mathcal Z_m\subset \ZZ/(F_{m+2}\ZZ)
\]
为附录 \ref{app:fold-multiplicity} 的谱零点集合。通过定理
\ref{thm:xi-time-part9x-serrin-ray-fibonacci-semirings} 的有限分辨率同构，把它传送为 Wulff 射线算术骨架上的零点类集合
\[
\mathcal Z_{m,K}^{\mathrm{arith}}
:=
\bigl\{[\rho W_K]_{\equiv_m^{(K)}}:\ [\rho]_{F_{m+2}}\in \mathcal Z_m\bigr\}
\subset
\mathfrak W_K^{(0)}/\!\equiv_m^{(K)}.
\]
若
\[
m+2\equiv 0\pmod 6,
\qquad
d=\frac{m+2}{2},
\]
则
\[
\bigl|\mathcal Z_{m,K}^{\mathrm{arith}}\bigr|
\ge
F_d
=
F_{(m+2)/2}.
\]
特别地，在最小三轴窗
\[
m=58
\]
处，有精确值
\[
\bigl|\mathcal Z_{58,K}^{\mathrm{arith}}\bigr|
=
839415
\ge
F_{30}=832040.
\]
\end{theorem}
\begin{proof}
由定理 \ref{thm:xi-time-part9x-serrin-ray-fibonacci-semirings}，集合
\[
\mathfrak W_K^{(0)}/\!\equiv_m^{(K)}
\]
与
\[
\ZZ/(F_{m+2}\ZZ)
\]
之间存在规范双射，故
\[
\bigl|\mathcal Z_{m,K}^{\mathrm{arith}}\bigr|
=
|\mathcal Z_m|.
\]

若 \(m+2\equiv 0\pmod 6\)，则由推论 \ref{cor:fold-zero-half-index-multiple6}，
\[
|\mathcal Z_m|
\ge
F_{(m+2)/2}.
\]
于是
\[
\bigl|\mathcal Z_{m,K}^{\mathrm{arith}}\bigr|
\ge
F_{(m+2)/2}.
\]

当 \(m=58\) 时，注 \ref{rem:fold-zero-lb-m58} 给出
\[
|\mathcal Z_{58}|=839415.
\]
故同样有
\[
\bigl|\mathcal Z_{58,K}^{\mathrm{arith}}\bigr|=839415.
\]
证毕。
\end{proof}

\begin{conjecture}[单参数各向异性 Serrin 观测的 Fibonacci 同余普适性]\label{conj:xi-time-part9x-serrin-ray-fibonacci-universality}
设
\[
q_m:\mathfrak W_K^{(0)}\twoheadrightarrow Q_m
\qquad(m\ge 1)
\]
是一族附着于整数尺度 Wulff 射线的有限分辨率观测，并满足：
\begin{enumerate}
  \item \(Q_m\) 为有限交换半环，且
  \[
  |Q_m|=F_{m+2};
  \]
  \item \(q_m\) 精确实现 Minkowski 和与尺度复合，即
  \[
  q_m(x\boxplus y)=q_m(x)+q_m(y),
  \qquad
  q_m(x\boxtimes y)=q_m(x)\,q_m(y);
  \]
  \item 这些有限商在加法 Grothendieck 包络上相容，并诱导与
  \[
  \{F_{m+2}\ZZ:\ m\ge 1\}
  \]
  相同的 Fibonacci 同余拓扑。
\end{enumerate}
则应存在唯一一族半环同构
\[
Q_m\cong \ZZ/(F_{m+2}\ZZ)
\qquad(m\ge 1),
\]
使之与全部相容结构交换。等价地，任何满足上述三条公理的单参数各向异性 Serrin 有限分辨率观测，其完备化都应退化为同一个 Fibonacci 同余完备化
\[
\widehat{\ZZ}_F\cong \widehat{\ZZ}.
\]
若该猜想成立，则
\[
\log_2\!\Bigl(\frac{2}{\varphi}\Bigr)
\]
将成为单参数可实现几何的普适精确恢复指数，而定理
\ref{thm:xi-time-part9x-serrin-ray-arithmetic-zero-block} 的半阶零点块则给出其在六倍窗上的普适算术盲区指纹。
\end{conjecture}
\begin{proof}[证据]
定理 \ref{thm:xi-time-part9x-serrin-ray-fibonacci-semirings} 已经给出一条显式模型：整数尺度 Wulff 射线在每个分辨率 \(m\) 上都产生模 \(F_{m+2}\) 半环。推论 \ref{cor:xi-time-part9x-serrin-ray-exact-capacity} 说明，一旦这一算术骨架被固定，精确恢复阈值便被强制为
\[
\log_2(2/\varphi).
\]
定理 \ref{thm:xi-time-part9x-serrin-ray-localized-ledger} 进一步说明，该骨架的局部化群完成在对偶端必裂解为一个圆因子与完全可恢复的 \(p\)-adic 账本核；而定理 \ref{thm:xi-time-part9x-serrin-ray-arithmetic-zero-block} 则表明，在六倍窗子序列上同一骨架还携带一个确定性的半阶零点类块。

因此，上述猜想所断言的是：对单参数各向异性 Serrin 观测而言，一旦有限分辨率对象数、两种运算以及诱导拓扑同时被固定，就不再存在第二种与之不同的算术骨架。换言之，Fibonacci 模数族不只是一条可实现模型，而应是这一单参数几何在有限分辨率端的唯一刚性终型。证毕。
\end{proof}

\noindent
上述链条表明，未归一化 Wulff 射线在有限分辨率上并不退化为平凡的单参数残余；相反，它携带一条精确可算的 Fibonacci 同余骨架。该骨架在局部化群完成后裂解为一个圆因子与一个可完全反演的 \(p\)-adic 账本核，在 window-\(6\) 上把 boundary 三点集精确钉扎为 \(\ZZ/(21\ZZ)\) 中一个既非加法陪集亦非乘法陪集的三点残余，并在六倍窗子序列上承载一个半阶量级的算术零点块。因而，对单参数可实现几何而言，有限分辨率容量律、局部化对偶核、boundary 阈值与六倍窗盲区不再是彼此分离的现象，而应被视为同一条 Fibonacci 同余主线的不同投影。

\endinput

\paragraph{Wulff 射线的 Schur 极小谱、window-\texorpdfstring{$6$}{6} 的 convex-tax 与审计偶窗缺口}
在定理 \ref{thm:xi-time-part9x-serrin-ray-fibonacci-semirings} 与推论 \ref{cor:xi-time-part9x-serrin-ray-exact-capacity} 已把 Figalli--Zhang 的各向异性 Serrin 单参数几何压缩为一条有限分辨率上的 Fibonacci 同余 Wulff 射线之后，还可把同一射线继续提升为一个严格的有限维变分基准。其核心观察是：一旦分辨率 \(m\) 固定，可实现纯几何对象只剩 \(M_m=F_{m+2}\) 个剩余类箱与总质量 \(N_m=2^m\) 的分配问题，而定理 \ref{thm:xi-time-part64ba-fold-multiplicity-majorization-balancing} 中的 Robin--Hood 调平机制便把这一问题推向唯一的 Schur 极小端点。由此得到的平衡谱
\[
d_m^{W}=\bigl((q_m+1)^{a_m},\,q_m^{M_m-a_m}\bigr)
\]
可视为各向异性 Wulff 射线在分辨率 \(m\) 上的纯几何基准谱；相对于这一基准的任何偏离，都自动转化为可显式计量的 convex-tax。

以下固定
\[
M_m:=F_{m+2},\qquad
N_m:=2^m,\qquad
N_m=q_mM_m+a_m,\qquad
0\le a_m<M_m.
\]

\begin{theorem}[Wulff--Schur 极小谱定理]\label{thm:xi-time-part9xaa-wulff-schur-minimal-spectrum}
\leanverified{paper\_xi\_time\_part9xaa\_wulff\_schur\_minimal\_spectrum}
固定 \(m\ge 1\)。设
\[
T:\{0,1,\dots,N_m-1\}\twoheadrightarrow Y,
\qquad
|Y|=M_m,
\]
并记其纤维向量为
\[
d(T)=(d_1,\dots,d_{M_m})\in\NN^{M_m},
\qquad
\sum_{i=1}^{M_m}d_i=N_m.
\]
定义 Wulff 平衡谱
\[
d_m^{W}:=\bigl((q_m+1)^{a_m},\,q_m^{M_m-a_m}\bigr).
\]
则 \(d(T)\) 可经有限次单位平衡变换
\[
(\ldots,x,\ldots,y,\ldots)\longmapsto (\ldots,x-1,\ldots,y+1,\ldots)
\qquad (x-y\ge 2)
\]
下降到 \(d_m^{W}\) 的一个排列。等价地，在按非增序重排后有主序关系
\[
d(T)^{\downarrow}\succ (d_m^{W})^{\downarrow}.
\]
因此 \(d_m^{W}\) 是所有 \(M_m\)-分箱、总质量 \(N_m\) 的正整数纤维谱中唯一的 Schur 极小元。
\end{theorem}
\begin{proof}
这正是定理 \ref{thm:xi-time-part64ba-fold-multiplicity-majorization-balancing} 在纯几何 Wulff 射线口径下的专化；这里给出直接证明以固定记号。取任意
\[
d=(d_1,\dots,d_{M_m})\in\NN^{M_m},
\qquad
\sum_i d_i=N_m.
\]
若存在 \(i,j\) 使
\[
d_i-d_j\ge 2,
\]
则执行一次单位平衡变换
\[
(d_i,d_j)\mapsto(d_i-1,d_j+1).
\]
该操作保持正整数性与总和不变；同时平方和严格下降，因为
\[
(d_i-1)^2+(d_j+1)^2-d_i^2-d_j^2
=
-2(d_i-d_j)+2
\le -2.
\]
故这一过程只能在有限步后终止。

终止时任意两坐标之差都至多为 \(1\)。由于
\[
N_m=q_mM_m+a_m,
\qquad
0\le a_m<M_m,
\]
满足这一性质的正整数向量只能是恰有 \(a_m\) 个分量等于 \(q_m+1\)、其余 \(M_m-a_m\) 个分量等于 \(q_m\) 的排列，也就是 \(d_m^{W}\) 的一个排列。

最后，由 Hardy--Littlewood--Pólya 判据，一个整数向量可经有限次 Robin--Hood 平衡化被送到另一个向量，当且仅当前者主序支配后者。故
\[
d(T)^{\downarrow}\succ (d_m^{W})^{\downarrow}.
\]
唯一的 Schur 极小元随之成立。证毕。
\end{proof}

\begin{corollary}[纯几何 Wulff 射线给出一切 separable 凸量的统一下包络]\label{cor:xi-time-part9xaa-wulff-universal-convex-envelope}
\leanverified{paper\_xi\_time\_part9xaa\_wulff\_universal\_convex\_envelope}
在定理 \ref{thm:xi-time-part9xaa-wulff-schur-minimal-spectrum} 的设定下，对任意严格凸函数
\[
\Phi:\NN\to\RR
\]
都有
\[
\sum_{i=1}^{M_m}\Phi(d_i)\ge (M_m-a_m)\Phi(q_m)+a_m\Phi(q_m+1),
\]
且等号当且仅当 \(d(T)\) 与 \(d_m^{W}\) 仅差一个排列。

特别地，若定义
\[
\mathsf{Col}(T):=\frac{1}{N_m^2}\sum_{i=1}^{M_m}d_i^2,
\]
\[
\kappa(T):=\frac1{N_m}\sum_{i=1}^{M_m}d_i\log d_i,
\]
\[
\mathcal G(T):=\prod_{i=1}^{M_m}\mathfrak S_{d_i},
\qquad
|\mathcal G(T)|=\prod_{i=1}^{M_m}d_i!,
\]
\[
\mathcal G_m^{W}:=\mathfrak S_{q_m+1}^{\,a_m}\times \mathfrak S_{q_m}^{\,M_m-a_m},
\qquad
|\mathcal G_m^{W}|=((q_m+1)!)^{a_m}(q_m!)^{M_m-a_m},
\]
则同时成立
\[
\sum_i d_i^s\ge (M_m-a_m)q_m^s+a_m(q_m+1)^s
\qquad (s>1),
\]
\[
\kappa(T)\ge \kappa_m^{W}
:=
\frac1{N_m}\Bigl[(M_m-a_m)q_m\log q_m+a_m(q_m+1)\log(q_m+1)\Bigr],
\]
\[
\log|\mathcal G(T)|\ge \log|\mathcal G_m^{W}|.
\]

另一方面，对任意严格凹函数
\[
\Psi:\NN\to\RR
\]
有
\[
\sum_{i=1}^{M_m}\Psi(d_i)\le (M_m-a_m)\Psi(q_m)+a_m\Psi(q_m+1).
\]
特别地，若定义
\[
p_i:=\frac{d_i}{N_m},
\qquad
u_i:=\frac1{M_m},
\qquad
p^{W}:=\frac{d_m^{W}}{N_m},
\]
则 Shannon 熵在 \(d_m^{W}\) 处取得唯一极大，而
\[
D_{\mathrm{KL}}(p\|u)
\]
在 \(d_m^{W}\) 处取得唯一极小。若再定义 dyadic 容量
\[
C_T(B):=\sum_{i=1}^{M_m}\min(d_i,2^B),
\qquad B\ge 0,
\]
则有统一上包络
\[
C_T(B)\le (M_m-a_m)\min(q_m,2^B)+a_m\min(q_m+1,2^B).
\]
\end{corollary}
\begin{proof}
由定理 \ref{thm:xi-time-part9xaa-wulff-schur-minimal-spectrum}，
\[
d(T)\succ d_m^{W}.
\]
对严格凸 \(\Phi\)，Karamata 不等式给出所示下界与等号刚性；对严格凹 \(\Psi\)，不等式方向反转。

取
\[
\Phi(t)=t^s\quad (s>1),
\qquad
\Phi(t)=t\log t,
\qquad
\Phi(t)=\log\Gamma(t+1),
\]
便分别得到幂和、\(\kappa\) 与 \(\log|\mathcal G(T)|\) 的下界。

对熵与 KL，记
\[
H(p):=-\sum_i p_i\log p_i.
\]
由于
\[
H(p)=\log N_m-\frac1{N_m}\sum_i d_i\log d_i
=
\log N_m-\kappa(T),
\]
而 \(D_{\mathrm{KL}}(p\|u)=\log M_m-H(p)\)，故熵极大与 KL 极小都等价于 \(\kappa(T)\) 取唯一极小，因而等价于 \(d(T)\sim d_m^{W}\)。

最后，函数
\[
t\longmapsto \min(t,2^B)
\]
在 \(\NN\) 上是凹函数，故再次由 Karamata 得到 \(C_T(B)\) 的统一上包络。证毕。
\end{proof}

\begin{theorem}[单泛函刚性]\label{thm:xi-time-part9xaa-single-functional-rigidity}
在定理 \ref{thm:xi-time-part9xaa-wulff-schur-minimal-spectrum} 的设定下，以下任一等式单独成立，均足以推出
\[
d(T)\sim d_m^{W}
\]
（仅差排列）：
\[
\sum_i d_i^2=(M_m-a_m)q_m^2+a_m(q_m+1)^2,
\]
\[
\log|\mathcal G(T)|=\log|\mathcal G_m^{W}|,
\]
\[
D_{\mathrm{KL}}(p\|u)=D_{\mathrm{KL}}(p^{W}\|u).
\]
换言之，碰撞率极小、规范体积极小与熵缺口极小三者中的任一者，已足以单独刻画纯尺度 Wulff 射线。
\end{theorem}
\begin{proof}
第一式对应严格凸函数
\[
\Phi(t)=t^2
\]
的唯一极小；第二式对应
\[
\Phi(t)=\log\Gamma(t+1)
\]
的唯一极小；第三式则等价于 Shannon 熵取唯一极大，亦即 \(\kappa(T)\) 取唯一极小。于是由推论 \ref{cor:xi-time-part9xaa-wulff-universal-convex-envelope} 的等号刚性，任一条件都迫使
\[
d(T)\sim d_m^{W}.
\]
证毕。
\end{proof}

\begin{theorem}[低纤维异常税的一般公式]\label{thm:xi-time-part9xaa-lowfiber-anomaly-tax}
假设
\[
q_m\ge 2,
\]
并固定整数
\[
0\le r\le \left\lfloor \frac{M_m-a_m}{2}\right\rfloor.
\]
考虑所有满足
\[
\sum_{i=1}^{M_m}d_i=N_m
\]
且至少有 \(r\) 个坐标不高于 \(q_m-1\) 的正整数向量 \(d\in\NN^{M_m}\)。则对任意严格凸函数
\[
\Phi:\NN\to\RR,
\]
该约束下的唯一极小元为
\[
d_m^{(r)}=\bigl((q_m-1)^r,\ q_m^{\,M_m-a_m-2r},\ (q_m+1)^{\,a_m+r}\bigr),
\]
并且极小值相对于 Wulff 基准的超额精确等于
\[
\sum_i \Phi(d_i)-\sum_i \Phi\bigl((d_m^{W})_i\bigr)
=
r\Bigl[\Phi(q_m-1)-2\Phi(q_m)+\Phi(q_m+1)\Bigr].
\]

特别地，取
\[
\Phi(t)=t^2
\]
得到
\[
\sum_i d_i^2\ge \sum_i (d_i^{W})^2+2r;
\]
取
\[
\Phi(t)=\log\Gamma(t+1)
\]
得到
\[
\log|\mathcal G(T)|\ge \log|\mathcal G_m^{W}|+r\log\frac{q_m+1}{q_m};
\]
取
\[
\Phi(t)=t\log t
\]
得到
\[
\kappa(T)\ge \kappa_m^{W}
+
\frac r{N_m}\Bigl[(q_m-1)\log(q_m-1)-2q_m\log q_m+(q_m+1)\log(q_m+1)\Bigr].
\]
\end{theorem}
\begin{proof}
设 \(d\) 是约束类中的极小元。若存在两个坐标分别满足
\[
d_i\le q_m-2,
\qquad
d_j\ge q_m+2,
\]
则执行一次单位平衡变换
\[
(d_j,d_i)\mapsto (d_j-1,d_i+1).
\]
由于 \(d_i+1\le q_m-1\)，低纤维个数不会减少；由于 \(d_j-1\ge q_m+1\)，也不会新产生低纤维。故该操作不破坏“至少 \(r\) 个低纤维”约束。同时由严格凸性，\(\sum_i \Phi(d_i)\) 严格下降，矛盾。

因此极小元的所有坐标只能取值于
\[
\{q_m-1,\ q_m,\ q_m+1\}.
\]
设这三类坐标个数分别为
\[
x,\ y,\ z.
\]
则
\[
x+y+z=M_m,
\]
\[
(q_m-1)x+q_my+(q_m+1)z=N_m=q_mM_m+a_m.
\]
联立解得
\[
z=a_m+x,
\qquad
y=M_m-a_m-2x.
\]
由可行性
\[
y\ge 0
\]
得
\[
x\le \left\lfloor\frac{M_m-a_m}{2}\right\rfloor,
\]
而约束“至少 \(r\) 个低纤维”则给出
\[
x\ge r.
\]

把这些关系代回目标函数，可得
\[
\sum_i \Phi(d_i)
=
x\Phi(q_m-1)+(M_m-a_m-2x)\Phi(q_m)+(a_m+x)\Phi(q_m+1).
\]
它关于 \(x\) 的离散增量为
\[
\Phi(q_m-1)-2\Phi(q_m)+\Phi(q_m+1)>0,
\]
故最小值在最小可行 \(x=r\) 处取得，且唯一。于是唯一极小元正是
\[
d_m^{(r)}=\bigl((q_m-1)^r,\ q_m^{\,M_m-a_m-2r},\ (q_m+1)^{\,a_m+r}\bigr),
\]
而相对 Wulff 基准 \(d_m^W=((q_m+1)^{a_m},q_m^{M_m-a_m})\) 的超额恰为
\[
r\Bigl[\Phi(q_m-1)-2\Phi(q_m)+\Phi(q_m+1)\Bigr].
\]

对三种特例，只需直接代入：
\[
(q_m-1)^2-2q_m^2+(q_m+1)^2=2,
\]
\[
\log\Gamma(q_m)-2\log\Gamma(q_m+1)+\log\Gamma(q_m+2)
=
\log\frac{q_m+1}{q_m},
\]
以及
\[
(q_m-1)\log(q_m-1)-2q_m\log q_m+(q_m+1)\log(q_m+1).
\]
证毕。
\end{proof}

\begin{theorem}[window-\texorpdfstring{$6$}{6} 审计直方图的变分刻画]\label{thm:xi-time-part9xaa-window6-audit-histogram-variational-characterization}
\leanverified{paper\_xi\_time\_part9xaa\_window6\_audit\_histogram\_variational\_characterization}
由表 \ref{tab:fold_bin_degeneracy_spectrum_m6_m12}，window-\(6\) 的 bin-fold 审计数据满足
\[
|X_6|=21,
\qquad
d_6^{\mathrm{bin}}:\#\{d:\#w\}=2:8,\ 3:4,\ 4:9.
\]
另一方面，
\[
N_6=64=3\cdot 21+1,
\]
故纯尺度 Wulff 基准谱为
\[
d_6^{W}=(4^1,3^{20}).
\]
于是，在所有 \(21\) 维正整数向量、总和为 \(64\)、且恰有 \(8\) 个二重纤维的模型中，审计直方图
\[
(2^8,3^4,4^9)
\]
是一切严格凸 separable 泛函
\[
\sum_i \Phi(d_i)
\]
的唯一极小元。
\end{theorem}
\begin{proof}
由表 \ref{tab:fold_bin_degeneracy_spectrum_m6_m12}，实际审计直方图确为
\[
(2^8,3^4,4^9).
\]
在 \(m=6\) 时有
\[
q_6=3,\qquad a_6=1.
\]
对定理 \ref{thm:xi-time-part9xaa-lowfiber-anomaly-tax} 取
\[
r=8,
\]
便得到在“至少 \(8\) 个坐标不高于 \(2\)”这一更宽约束下的唯一极小元为
\[
(2^8,3^{\,21-1-16},4^{\,1+8})=(2^8,3^4,4^9).
\]
因此在更窄的约束类“恰有 \(8\) 个二重纤维”中，它仍是唯一极小元。证毕。
\end{proof}

\begin{corollary}[boundary center 的不可消除代价]\label{cor:xi-time-part9xaa-boundary-center-irreducible-tax}
由推论 \ref{cor:fold-bin-boundary-center-z2}，window-\(6\) 的三条 boundary 二点纤维在规范群中给出一个规范中心嵌入
\[
(\ZZ_2)^3\hookrightarrow Z(\mathcal G_6^{\mathrm{bin}}).
\]
对任意此类模型的纤维向量 \(d=(d_1,\dots,d_{21})\)，再记
\[
\mathsf{Col}(d):=\frac1{64^2}\sum_{i=1}^{21}d_i^2,
\qquad
\mathcal G(d):=\prod_{i=1}^{21}\mathfrak S_{d_i},
\qquad
\kappa(d):=\frac1{64}\sum_{i=1}^{21}d_i\log d_i.
\]
若该模型仍以三个独立 \(S_2\)-因子保留上述 boundary sheet-parity 中心，则至少拥有 \(r=3\) 个二重纤维；因而由定理 \ref{thm:xi-time-part9xaa-lowfiber-anomaly-tax} 必有
\[
\sum_i d_i^2\ge 196+6=202,
\qquad
\mathsf{Col}(d)\ge \frac{202}{64^2}=\frac{101}{2048},
\]
\[
\log|\mathcal G(d)|\ge \log|\mathcal G_6^{W}|+3\log\frac43,
\]
\[
\kappa(d)\ge \kappa_6^{W}
+
\frac{3}{64}\Bigl(2\log2-6\log3+4\log4\Bigr)
=
\kappa_6^{W}+\frac{3}{32}\log\frac{32}{27}.
\]
换言之，window-\(6\) 的 boundary 中心不是“免费异常”，而携带严格的 convex-tax。
\end{corollary}
\begin{proof}
由推论 \ref{cor:fold-bin-boundary-center-z2}，保留 canonical boundary sheet-parity 中心意味着必须保留三个彼此独立的 \(S_2\) 因子；故至少存在三条二重纤维。对 \(m=6\) 有
\[
q_6=3,\qquad a_6=1,
\]
故由定理 \ref{thm:xi-time-part9xaa-lowfiber-anomaly-tax} 取 \(r=3\)，平方和超额至少为 \(2r=6\)，群体积超额至少为 \(3\log\frac43\)，\(\kappa\)-超额至少为
\[
\frac{3}{64}\Bigl[(2)\log2-2\cdot 3\log3+4\log4\Bigr].
\]
而 Wulff 基准
\[
d_6^{W}=(4^1,3^{20})
\]
的平方和为
\[
4^2+20\cdot 3^2=16+180=196,
\]
从而得到
\[
\mathsf{Col}(d)\ge \frac{202}{64^2}=\frac{101}{2048}.
\]
其余各式同理。证毕。
\end{proof}
\endinput

在上述 Wulff 基准、低纤维凸税与 boundary 中心代价都已固定之后，window-\(6\) 的审计谱与审计偶窗的有限窗缺口还可进一步写成完全显式的闭式。

\begin{theorem}[window-\texorpdfstring{$6$}{6} 审计谱的精确异常税恒等式]\label{thm:xi-time-part9xaa-window6-audited-spectrum-exact-tax}
\leanverified{paper\_xi\_time\_part9xaa\_window6\_audited\_spectrum\_exact\_tax}
记
\[
d_6^{W}=(4^1,3^{20}),
\qquad
\mathcal G_6^{W}:=\mathfrak S_4\times \mathfrak S_3^{20},
\qquad
g_6^{W}:=\frac1{64}\log|\mathcal G_6^{W}|,
\]
\[
\kappa_6^{W}:=\frac1{64}\bigl(4\log4+20\cdot 3\log3\bigr),
\qquad
\mathsf{Col}_6^{W}:=\frac{4^2+20\cdot 3^2}{64^2}=\frac{49}{1024}.
\]
另一方面，由表 \ref{tab:fold_bin_degeneracy_spectrum_m6_m12}、表 \ref{tab:fold_bin_auditable_invariants_m6_m12} 与表 \ref{tab:fold_bin_gauge_m6_invariants}，审计值满足
\[
d_6^{\mathrm{bin}}=(2^8,3^4,4^9),
\qquad
\mathsf{Col}_6^{\mathrm{bin}}=\frac{53}{1024},
\]
\[
g_6^{\mathrm{bin}}=0.645542\ldots,
\qquad
\kappa_6=1.159067\ldots.
\]
则成立精确恒等式
\[
\mathsf{Col}_6^{\mathrm{bin}}-\mathsf{Col}_6^{W}=\frac1{256},
\]
\[
\log|\mathcal G_6^{\mathrm{bin}}|-\log|\mathcal G_6^{W}|=8\log\frac43,
\]
\[
g_6^{\mathrm{bin}}-g_6^{W}=\frac18\log\frac43,
\]
\[
\kappa_6-\kappa_6^{W}=\frac14\log\frac{32}{27}.
\]
因此，window-\(6\) 的实际审计谱并非任意偏离 Wulff 基准；它恰好实现了“具有八个二重纤维”这一约束下的最小可付异常税。
\end{theorem}
\begin{proof}
由定理 \ref{thm:xi-time-part9xaa-window6-audit-histogram-variational-characterization}，审计直方图正是 \(m=6\)、\(r=8\) 下的唯一凸极小元。于是把定理 \ref{thm:xi-time-part9xaa-lowfiber-anomaly-tax} 取 \(r=8\) 并代入 \(q_6=3,a_6=1\) 即得全部精确超额。

对平方和，
\[
\sum_i d_i^2-\sum_i \bigl((d_6^{W})_i\bigr)^2=2r=16,
\]
故
\[
\mathsf{Col}_6^{\mathrm{bin}}-\mathsf{Col}_6^{W}
=
\frac{16}{64^2}
=
\frac1{256}.
\]

对群体积，
\[
\log|\mathcal G_6^{\mathrm{bin}}|-\log|\mathcal G_6^{W}|
=
8\log\frac{4}{3},
\]
再除以 \(64\)，得到
\[
g_6^{\mathrm{bin}}-g_6^{W}
=
\frac{8}{64}\log\frac43
=
\frac18\log\frac43.
\]

对 \(\kappa\)，由定理 \ref{thm:xi-time-part9xaa-lowfiber-anomaly-tax}，
\[
\kappa_6-\kappa_6^{W}
=
\frac{8}{64}\Bigl(2\log2-6\log3+4\log4\Bigr).
\]
化简得
\[
\kappa_6-\kappa_6^{W}
=
\frac18(10\log2-6\log3)
=
\frac14\log\frac{32}{27}.
\]
证毕。
\end{proof}

\begin{proposition}[审计偶窗始终严格高于 Wulff--Schur 地板]\label{prop:xi-time-part9xaa-audited-even-windows-above-wulff-floor}
对审计偶窗 \(m\in\{8,10,12\}\)，记
\[
d_m^{W}=\bigl((q_m+1)^{a_m},\,q_m^{M_m-a_m}\bigr),
\qquad
\mathsf{Col}_m^{W}:=\frac{(M_m-a_m)q_m^2+a_m(q_m+1)^2}{2^{2m}},
\]
\[
g_m^{W}:=\frac1{2^m}\log\!\Bigl(((q_m+1)!)^{a_m}(q_m!)^{M_m-a_m}\Bigr).
\]
则由表 \ref{tab:fold_bin_auditable_invariants_m6_m12} 的审计值与上述 Wulff 地板比较，可得严格正的有限窗缺口：
\[
\mathsf{Col}_8^{\mathrm{bin}}-\mathsf{Col}_8^{W}
=
\frac{323}{16384}-\frac{301}{16384}
=
\frac{11}{8192},
\]
\[
\mathsf{Col}_{10}^{\mathrm{bin}}-\mathsf{Col}_{10}^{W}
=
\frac{1925}{262144}-\frac{57}{8192}
=
\frac{101}{262144},
\]
\[
\mathsf{Col}_{12}^{\mathrm{bin}}-\mathsf{Col}_{12}^{W}
=
\frac{11617}{4194304}-\frac{22273}{8388608}
=
\frac{961}{8388608},
\]
并且
\[
g_8^{\mathrm{bin}}-g_8^{W}\approx 0.0346852,
\qquad
g_{10}^{\mathrm{bin}}-g_{10}^{W}\approx 0.0268796,
\qquad
g_{12}^{\mathrm{bin}}-g_{12}^{W}\approx 0.0213913.
\]
因此，已审计偶窗上的 bin-fold 从未落到纯尺度 Wulff 射线的 Schur 极小地板上；每一窗都携带额外的非尺度纤维结构。
\end{proposition}
\begin{proof}
由表 \ref{tab:fold_bin_auditable_invariants_m6_m12}，
\[
\mathsf{Col}_8^{\mathrm{bin}}=\frac{323}{16384},
\qquad
\mathsf{Col}_{10}^{\mathrm{bin}}=\frac{1925}{262144},
\qquad
\mathsf{Col}_{12}^{\mathrm{bin}}=\frac{11617}{4194304}.
\]
另一方面，
\[
(q_8,a_8)=(4,36),\qquad
(q_{10},a_{10})=(7,16),\qquad
(q_{12},a_{12})=(10,326),
\]
故
\[
\mathsf{Col}_8^{W}=\frac{19\cdot 4^2+36\cdot 5^2}{256^2}=\frac{301}{16384},
\]
\[
\mathsf{Col}_{10}^{W}=\frac{128\cdot 7^2+16\cdot 8^2}{1024^2}=\frac{57}{8192},
\]
\[
\mathsf{Col}_{12}^{W}=\frac{51\cdot 10^2+326\cdot 11^2}{4096^2}=\frac{22273}{8388608}.
\]
逐窗相减即得三条碰撞率缺口。

对 \(g\)-缺口，只需把表 \ref{tab:fold_bin_degeneracy_spectrum_m6_m12} 的直方图直接代入
\[
g_m^{\mathrm{bin}}
=
\frac1{2^m}\sum_{d}\#\{w\in X_m:\ d_m^{\mathrm{bin}}(w)=d\}\,\log(d!).
\]
于是
\[
g_8^{\mathrm{bin}}
=
\frac1{256}\bigl(21\log 3!+11\log 5!+23\log 6!\bigr),
\]
\[
g_{10}^{\mathrm{bin}}
=
\frac1{1024}\bigl(55\log 5!+52\log 8!+37\log 9!\bigr),
\]
\[
g_{12}^{\mathrm{bin}}
=
\frac1{4096}\bigl(144\log 8!+85\log 12!+148\log 13!\bigr).
\]
再与相应的 \(g_m^{W}\) 比较，即得
\[
g_8^{\mathrm{bin}}-g_8^{W}\approx 0.0346852,\qquad
g_{10}^{\mathrm{bin}}-g_{10}^{W}\approx 0.0268796,\qquad
g_{12}^{\mathrm{bin}}-g_{12}^{W}\approx 0.0213913.
\]
证毕。
\end{proof}

\begin{conjecture}[偶窗 Fibonacci 底层的 Schur 刚性]\label{conj:xi-time-part9xaa-even-window-fibonacci-schur-rigidity}
对每个偶数 \(m\ge 6\)，bin-fold 的完整多重度谱应是如下变分问题的唯一解：在所有 \(M_m\)-维正整数向量 \(d\) 中，满足
\[
\sum_i d_i=2^m,
\qquad
\#\{i:\ d_i=F_{m/2}\}=F_m,
\]
使得对某个严格凸 separable 泛函
\[
\sum_i \Phi(d_i)
\]
取极小。等价地，上述“某个”可替换为“任意一个”严格凸 separable 泛函。

若该猜想成立，则定理 \ref{thm:fold-bin-min-degeneracy-fib-audited} 不再只是审计偶窗上的有限样本现象，而会升级为一个全偶窗的 Schur 刚性原理：一旦最底层 Fibonacci 扇区
\[
\#\{d=d_{\min}(m)\}=F_m,
\qquad
d_{\min}(m)=F_{m/2}
\]
被固定，bin-fold 的全部高层直方图便被唯一强制决定。
\end{conjecture}
\begin{proof}[证据]
定理 \ref{thm:fold-bin-min-degeneracy-fib-audited} 已在审计偶窗 \(m\in\{6,8,10,12\}\) 上给出
\[
d_{\min}(m)=F_{m/2},
\qquad
\bigl|\mathcal S_{m,d_{\min}(m)}\bigr|=F_m=|X_{m-2}|.
\]
另一方面，定理 \ref{thm:xi-time-part9xaa-lowfiber-anomaly-tax} 表明：一旦“低纤维负载”被固定，严格凸 separable 泛函的极小元便被唯一地推进到一个三层平衡谱。对 \(m=6\) 而言，定理 \ref{thm:xi-time-part9xaa-window6-audit-histogram-variational-characterization} 更进一步表明，审计直方图 \((2^8,3^4,4^9)\) 恰是“八个二重纤维”约束下的唯一凸极小元。

因此，现有审计数据已经显示：最底层 Fibonacci 扇区并非被动统计量，而极可能是整个多重度谱的变分锚点。该猜想所断言的正是这一现象在全偶窗上的延续。证毕。
\end{proof}

\noindent
由此，Figalli--Zhang 所给出的 Wulff 单参数几何，在有限分辨率上获得了一条此前尚未显式抽出的变分主线：纯尺度 Wulff 射线不只给出一个可实现的 Fibonacci 同余骨架，而且给出同一窗口内唯一的 Schur 极小基准；一切低纤维、boundary 中心或审计异常都因此获得了精确可付的 convex-tax。特别地，window-\(6\) 的实际审计谱并非任意偏离纯几何地板，而恰好是在“八个二重纤维”这一约束下达到凸极小的唯一谱型。于是，Wulff 射线、Robin--Hood 调平、boundary sheet-parity 与审计直方图第一次被统一压缩进同一条有限分辨率 Schur 变分链。

\paragraph{同一 Fibonacci 同余骨架上的频谱实现分离}
前述链条已经把 Figalli--Zhang 的各向异性 Serrin 可实现类压缩为 Fibonacci 同余骨架上的纯尺度 Wulff 射线，并把该射线识别为同窗口内唯一的 Schur 极小地板。然而，共享同一有限分辨率半环并不意味着共享同一 Fourier 实现：bin-fold 的非平凡频谱由定理 \ref{thm:fold-spectrum-factorization} 所给出的 Fibonacci 频率乘积分解支配，而纯尺度 Wulff 射线对应的只是模 \(F_{m+2}\) 的截断区间计数。下列结果把这一区别压缩为零集结构、最大纤维归一化偏置与谱核类型三条彼此独立的可审计分离律。

\begin{theorem}[纯尺度 Wulff 射线的非平凡 Fourier 零集公式]\label{thm:xi-time-part9xaab-wulff-ray-fourier-zero-formula}
\leanverified{paper\_xi\_time\_part9xaab\_wulff\_ray\_fourier\_zero\_formula}
沿用定理 \ref{thm:xi-time-part9x-serrin-ray-fibonacci-semirings} 与推论 \ref{cor:xi-time-part9x-serrin-ray-exact-capacity} 的记号，记
\[
M_m:=F_{m+2},
\qquad
2^m=q_mM_m+a_m,
\qquad
0\le a_m<M_m.
\]
对纯尺度 Wulff 射线定义纤维计数函数
\[
d_m^{W,\mathrm{ray}}(r)
:=
\#\bigl\{0\le \rho<2^m:\ [\rho]_{M_m}=r\bigr\}
\qquad
\bigl(r\in \ZZ/(M_m\ZZ)\bigr).
\]
则有显式剖面
\[
d_m^{W,\mathrm{ray}}(r)=q_m+\mathbf 1_{I_m}(r),
\qquad
I_m:=\{0,1,\dots,a_m-1\}\subset \ZZ/(M_m\ZZ).
\]
定义其离散 Fourier 变换
\[
\widehat d_m^{W,\mathrm{ray}}(k)
:=
\sum_{r\in \ZZ/(M_m\ZZ)}
d_m^{W,\mathrm{ray}}(r)\,
e^{-2\pi i kr/M_m},
\]
以及非平凡零集
\[
\mathcal Z_m^{W,\mathrm{ray}}
:=
\bigl\{k\in \ZZ/(M_m\ZZ)\setminus\{0\}:\ \widehat d_m^{W,\mathrm{ray}}(k)=0\bigr\}.
\]
则对每个 \(k\neq 0\) 都有
\[
\widehat d_m^{W,\mathrm{ray}}(k)
=
\sum_{r=0}^{a_m-1}e^{-2\pi i kr/M_m}
=
\frac{1-e^{-2\pi i ka_m/M_m}}{1-e^{-2\pi i k/M_m}}.
\]
若再记
\[
g_m:=\gcd(M_m,a_m),
\]
则
\[
\mathcal Z_m^{W,\mathrm{ray}}
=
\left\{
j\frac{M_m}{g_m}\ \ (\mathrm{mod}\ M_m):\ 1\le j\le g_m-1
\right\},
\qquad
\bigl|\mathcal Z_m^{W,\mathrm{ray}}\bigr|=g_m-1.
\]
特别地，\(\{0\}\cup \mathcal Z_m^{W,\mathrm{ray}}\) 恰为 \(\ZZ/(M_m\ZZ)\) 的一个循环子群，故纯尺度 Wulff 射线的非平凡零集完全由 \(\gcd(M_m,a_m)\) 控制，与各向异性体 \(K\) 的具体形状无关。
\end{theorem}
\begin{proof}
由 Euclid 除法
\[
2^m=q_mM_m+a_m
\]
可知区间 \(\{0,1,\dots,2^m-1\}\) 中每个模 \(M_m\) 的余类恰出现 \(q_m\) 次，而前 \(a_m\) 个余类各多出一次，故
\[
d_m^{W,\mathrm{ray}}(r)=q_m+\mathbf 1_{I_m}(r).
\]
对 \(k\neq 0\)，常值部分的 Fourier 系数消失，于是
\[
\widehat d_m^{W,\mathrm{ray}}(k)
=
\sum_{r=0}^{a_m-1}e^{-2\pi i kr/M_m},
\]
并由几何级数公式得到所示闭式。

该和式为零当且仅当
\[
e^{-2\pi i ka_m/M_m}=1
\qquad\text{且}\qquad
e^{-2\pi i k/M_m}\neq 1,
\]
也即
\[
M_m\mid ka_m,
\qquad
M_m\nmid k.
\]
写
\[
M_m=g_mM_m',
\qquad
a_m=g_ma_m',
\qquad
\gcd(M_m',a_m')=1,
\]
则 \(M_m\mid ka_m\) 等价于 \(M_m'\mid k\)。因此所有非零解恰为模 \(M_m\) 下的非零倍数
\[
\frac{M_m}{g_m},
\ 2\frac{M_m}{g_m},
\ \dots,
\ (g_m-1)\frac{M_m}{g_m},
\]
共有 \(g_m-1\) 个。于是 \(\{0\}\cup \mathcal Z_m^{W,\mathrm{ray}}\) 正是由 \(M_m/g_m\) 生成的循环子群。证毕。
\end{proof}

\begin{corollary}[纯尺度 Wulff 零集的二进赋值相图]\label{cor:xi-time-part9xaab-wulff-ray-zero-2adic-phase}
在定理 \ref{thm:xi-time-part9xaab-wulff-ray-fourier-zero-formula} 的记号下，
\[
g_m=\gcd(F_{m+2},2^m)=2^{\nu_2(F_{m+2})}.
\]
因此
\[
\bigl|\mathcal Z_m^{W,\mathrm{ray}}\bigr|
=
\begin{cases}
0,
& m+2\equiv 1,2,4,5 \pmod 6,
\\[1mm]
1,
& m+2\equiv 3 \pmod 6,
\\[1mm]
2^{\nu_2(m+2)+2}-1,
& m+2\equiv 0 \pmod 6.
\end{cases}
\]
特别地，沿六倍窗子序列 \(m+2\equiv 0\pmod 6\) 有线性上界
\[
\bigl|\mathcal Z_m^{W,\mathrm{ray}}\bigr|
\le
4(m+2)-1.
\]
\end{corollary}
\begin{proof}
由 \(a_m\equiv 2^m\pmod{M_m}\) 与定理 \ref{thm:xi-time-part9xaab-wulff-ray-fourier-zero-formula}，
\[
g_m=\gcd(M_m,a_m)=\gcd(F_{m+2},2^m).
\]
因为 \(2^m\) 是纯 \(2\)-幂，故
\[
g_m=2^{\nu_2(F_{m+2})}.
\]
再用 Fibonacci 数的标准 \(2\)-adic 赋值公式
\[
\nu_2(F_n)=
\begin{cases}
0,
& n\equiv 1,2,4,5 \pmod 6,
\\
1,
& n\equiv 3 \pmod 6,
\\
\nu_2(n)+2,
& n\equiv 0 \pmod 6,
\end{cases}
\]
便得到三分情形。最后由
\[
\bigl|\mathcal Z_m^{W,\mathrm{ray}}\bigr|=g_m-1
\]
与 \(2^{\nu_2(m+2)}\le m+2\) 推得六倍窗上的线性上界。证毕。
\end{proof}

\begin{theorem}[六倍窗上半阶零点块与纯尺度零集的严格分离]\label{thm:xi-time-part9xaab-wulff-ray-zero-block-separation}
记
\[
\mathcal Z_m^{\mathrm{bin}}:=\mathcal Z_m
\]
为 bin-fold 的非平凡谱零点集合，\(\mathcal Z_m^{W,\mathrm{ray}}\) 为定理 \ref{thm:xi-time-part9xaab-wulff-ray-fourier-zero-formula} 的纯尺度 Wulff 零点集合。若
\[
m+2\equiv 0\pmod 6,
\]
则有严格分离估计
\[
\frac{|\mathcal Z_m^{\mathrm{bin}}|}{|\mathcal Z_m^{W,\mathrm{ray}}|}
\ge
\frac{F_{(m+2)/2}}{4(m+2)-1}.
\]
因此，bin-fold 在六倍窗上的半阶大陪集零点块不可能由任何单参数 Wulff 尺度自由度单独产生。

特别地，在最小三轴窗 \(m=58\) 上，
\[
|\mathcal Z_{58}^{\mathrm{bin}}|=839415,
\qquad
|\mathcal Z_{58}^{W,\mathrm{ray}}|=15,
\qquad
\frac{|\mathcal Z_{58}^{\mathrm{bin}}|}{|\mathcal Z_{58}^{W,\mathrm{ray}}|}=55961.
\]
\end{theorem}
\begin{proof}
若 \(m+2\equiv 0\pmod 6\)，则由推论 \ref{cor:fold-zero-half-index-multiple6}，
\[
|\mathcal Z_m^{\mathrm{bin}}|
=
|\mathcal Z_m|
\ge
F_{(m+2)/2}.
\]
另一方面，推论 \ref{cor:xi-time-part9xaab-wulff-ray-zero-2adic-phase} 给出
\[
|\mathcal Z_m^{W,\mathrm{ray}}|
\le
4(m+2)-1.
\]
两式相除即得一般比值下界。

当 \(m=58\) 时，推论 \ref{cor:xi-fold-zero-m58-three-coset-exhaustion} 给出
\[
|\mathcal Z_{58}^{\mathrm{bin}}|=839415,
\]
而 \(m+2=60\) 且 \(\nu_2(60)=2\)，故由推论 \ref{cor:xi-time-part9xaab-wulff-ray-zero-2adic-phase}
\[
|\mathcal Z_{58}^{W,\mathrm{ray}}|
=
2^{2+2}-1
=15.
\]
因此
\[
\frac{|\mathcal Z_{58}^{\mathrm{bin}}|}{|\mathcal Z_{58}^{W,\mathrm{ray}}|}
=
\frac{839415}{15}
=55961.
\]
证毕。
\end{proof}

\begin{theorem}[一阶指数普适而二阶归一化偏置分离]\label{thm:xi-time-part9xaab-wulff-ray-first-second-order-separation}
记
\[
M_m^{W,\mathrm{ray}}
:=
\max_{r}d_m^{W,\mathrm{ray}}(r)=q_m+1,
\qquad
\overline d_m:=\frac{2^m}{F_{m+2}}.
\]
则纯尺度 Wulff 射线满足
\[
\lim_{m\to\infty}\frac{1}{m}\log M_m^{W,\mathrm{ray}}
=
\log\frac{2}{\varphi},
\qquad
\lim_{m\to\infty}\frac{M_m^{W,\mathrm{ray}}}{\overline d_m}=1.
\]
另一方面，bin-fold 的最大纤维 \(d_{\max}^{\mathrm{bin}}(m)\) 满足
\[
\lim_{m\to\infty}\frac{1}{m}\log d_{\max}^{\mathrm{bin}}(m)
=
\log\frac{2}{\varphi},
\]
但同时
\[
\liminf_{m\to\infty}\frac{d_{\max}^{\mathrm{bin}}(m)}{\overline d_m}
\ge
1+C_\varphi
\approx
1.2244178274047294
>1.
\]
因此，\(\log(2/\varphi)\) 只是一阶基数失配常数；真正区分“纯尺度几何”与“bin-fold 异常纤维”的，是二阶归一化 \(L^\infty\)-偏置系数。
\end{theorem}
\begin{proof}
由
\[
2^m=q_mF_{m+2}+a_m,
\qquad
0\le a_m<F_{m+2},
\]
以及
\[
F_{m+2}\sim \frac{\varphi^{m+2}}{\sqrt5},
\]
可得
\[
q_m
\sim
\frac{2^m}{F_{m+2}}
\sim
\frac{\sqrt5}{\varphi^2}\left(\frac{2}{\varphi}\right)^m.
\]
故
\[
M_m^{W,\mathrm{ray}}=q_m+1
\]
满足
\[
\frac{1}{m}\log M_m^{W,\mathrm{ray}}\to \log\frac{2}{\varphi}.
\]
另一方面，
\[
\frac{M_m^{W,\mathrm{ray}}}{\overline d_m}
=
\frac{q_m+1}{q_m+a_m/F_{m+2}}
=
1+\frac{1-a_m/F_{m+2}}{q_m+a_m/F_{m+2}}
\longrightarrow 1,
\]
因为 \(q_m\to\infty\)。

对 bin-fold，指数极限由定理 \ref{thm:fold-bin-max-fiber-exponent} 给出，而归一化偏置下界正是推论 \ref{cor:fold-max-fiber-constant-bias} 的内容。证毕。
\end{proof}

\leanverified{paper\_xi\_time\_part9xaab\_wulff\_ray\_fejer\_kernel}
\begin{theorem}[纯尺度 Wulff 射线中心化缺陷的离散 Fej\'er 核表示]\label{thm:xi-time-part9xaab-wulff-ray-fejer-kernel}
沿用定理 \ref{thm:xi-time-part9xaab-wulff-ray-fourier-zero-formula} 的记号，并定义中心化缺陷
\[
\delta_m^{W,\mathrm{ray}}(r)
:=
d_m^{W,\mathrm{ray}}(r)-\frac{2^m}{F_{m+2}}
=
\mathbf 1_{I_m}(r)-\frac{a_m}{M_m}.
\]
则
\[
\widehat{\delta_m^{W,\mathrm{ray}}}(0)=0,
\]
且对每个 \(k\neq 0\) 都有
\[
\widehat{\delta_m^{W,\mathrm{ray}}}(k)
=
\frac{1-e^{-2\pi i ka_m/M_m}}{1-e^{-2\pi i k/M_m}}.
\]
从而
\[
\bigl|\widehat{\delta_m^{W,\mathrm{ray}}}(k)\bigr|^2
=
\frac{\sin^2(\pi ka_m/M_m)}{\sin^2(\pi k/M_m)}
\qquad(k\neq 0).
\]
亦即，纯尺度 Wulff 射线的全部非平凡谱能恰是一个离散 Fej\'er 核。并且 Parseval 质量满足
\[
\frac{1}{M_m}
\sum_{k\neq 0}
\bigl|\widehat{\delta_m^{W,\mathrm{ray}}}(k)\bigr|^2
=
a_m\left(1-\frac{a_m}{M_m}\right).
\]
\end{theorem}
\begin{proof}
由定理 \ref{thm:xi-time-part9xaab-wulff-ray-fourier-zero-formula}，
\[
d_m^{W,\mathrm{ray}}(r)=q_m+\mathbf 1_{I_m}(r),
\qquad
\frac{2^m}{M_m}=q_m+\frac{a_m}{M_m},
\]
故
\[
\delta_m^{W,\mathrm{ray}}(r)=\mathbf 1_{I_m}(r)-\frac{a_m}{M_m}.
\]
于是 \(\widehat{\delta_m^{W,\mathrm{ray}}}(0)=0\) 显然成立。对 \(k\neq 0\)，常值项 Fourier 系数消失，因此
\[
\widehat{\delta_m^{W,\mathrm{ray}}}(k)
=
\sum_{r=0}^{a_m-1}e^{-2\pi i kr/M_m},
\]
并由几何级数公式得到第一条闭式。再用
\[
|1-e^{-i\theta}|=2\bigl|\sin(\theta/2)\bigr|
\]
即可得到 Fej\'er 核表达式。

最后由 Parseval 恒等式，
\[
\frac{1}{M_m}\sum_{k}\bigl|\widehat{\delta_m^{W,\mathrm{ray}}}(k)\bigr|^2
=
\sum_{r}\bigl|\delta_m^{W,\mathrm{ray}}(r)\bigr|^2.
\]
右端直接化简为
\[
a_m\left(1-\frac{a_m}{M_m}\right)^2
+
(M_m-a_m)\left(\frac{a_m}{M_m}\right)^2
=
a_m\left(1-\frac{a_m}{M_m}\right).
\]
又因 \(\widehat{\delta_m^{W,\mathrm{ray}}}(0)=0\)，故得到所示非平凡频率 Parseval 质量。证毕。
\end{proof}

\begin{theorem}[Fibonacci--Riesz 与 Fej\'er 类在最小三轴窗上的不可同构性]\label{thm:xi-time-part9xaab-binfold-fejer-riesz-nonisomorphism}
记
\[
\delta_m^{\mathrm{bin}}(r):=d_m(r)-\frac{2^m}{F_{m+2}},
\qquad
\delta_m^{W,\mathrm{ray}}(r):=d_m^{W,\mathrm{ray}}(r)-\frac{2^m}{F_{m+2}}.
\]
则在非零频率上，
\[
\widehat{\delta_m^{\mathrm{bin}}}(k)=\widehat d_m(k)
\qquad (k\neq 0),
\]
故由定理 \ref{thm:fold-spectrum-factorization} 与注 \ref{rem:fold-zero-sparse-necessity}，bin-fold 的归一化非平凡谱能是 Fibonacci 频率驱动的有限 Riesz 乘积；另一方面，由定理 \ref{thm:xi-time-part9xaab-wulff-ray-fejer-kernel}，纯尺度 Wulff 射线的归一化非平凡谱能是单区间诱导的离散 Fej\'er 核。

更强地，在最小三轴窗 \(m=58\) 上，不存在任何双射
\[
\sigma:\ZZ/(F_{60}\ZZ)\setminus\{0\}\longrightarrow \ZZ/(F_{60}\ZZ)\setminus\{0\}
\]
使得
\[
\bigl|\widehat{\delta_{58}^{\mathrm{bin}}}(k)\bigr|
=
\bigl|\widehat{\delta_{58}^{W,\mathrm{ray}}}(\sigma(k))\bigr|
\qquad(\forall\,k\neq 0).
\]
因此，Fibonacci--Riesz 类与 Fej\'er 类在最小三轴窗上已严格分离，不能通过任何频率重标号互相同构。
\end{theorem}
\begin{proof}
若存在这样的双射 \(\sigma\)，则两侧非零频率零集的基数必须相同。由推论 \ref{cor:xi-fold-zero-m58-three-coset-exhaustion}，
\[
\bigl|\mathcal Z_{58}^{\mathrm{bin}}\bigr|=839415.
\]
而由推论 \ref{cor:xi-time-part9xaab-wulff-ray-zero-2adic-phase}，
\[
\bigl|\mathcal Z_{58}^{W,\mathrm{ray}}\bigr|=15.
\]
两者矛盾，故所述双射不存在。于是两类频谱模型在最小三轴窗上已不可同构。证毕。
\end{proof}

\begin{conjecture}[Fej\'er 刚性]\label{conj:xi-time-part9xaab-fejer-rigidity}
对每个 \(m\ge 1\)，令
\[
G_m:=\ZZ/(F_{m+2}\ZZ),
\qquad
\overline d_m:=\frac{2^m}{F_{m+2}}.
\]
设 \(d_m^\sharp:G_m\to \NN\) 是一列满足
\[
\sum_{r\in G_m}d_m^\sharp(r)=2^m
\]
的纤维函数，并记
\[
M_m^\sharp:=\max_{r\in G_m}d_m^\sharp(r),
\qquad
\delta_m^\sharp:=d_m^\sharp-\overline d_m.
\]
我猜想下列三条合起来足以刻画纯尺度 Wulff 射线：
\begin{enumerate}
  \item
  \[
  \lim_{m\to\infty}\frac{1}{m}\log M_m^\sharp
  =
  \log\frac{2}{\varphi};
  \]
  \item
  \[
  \lim_{m\to\infty}\frac{M_m^\sharp}{\overline d_m}=1;
  \]
  \item 对每个 \(m\)，集合
  \[
  \{0\}\cup
  \bigl\{k\in G_m:\ \widehat{\delta_m^\sharp}(k)=0\bigr\}
  \]
  是 \(G_m\) 的一个循环子群。
\end{enumerate}
若此猜想成立，则对每个 \(m\) 都存在群自同构 \(\sigma_m\in \mathrm{Aut}(G_m)\)，使
\[
d_m^\sharp\circ \sigma_m^{-1}
=
q_m+\mathbf 1_{I_m},
\qquad
I_m=\{0,1,\dots,a_m-1\}.
\]
等价地，
\[
\delta_m^\sharp\circ \sigma_m^{-1}
=
\mathbf 1_{I_m}-\frac{a_m}{F_{m+2}}.
\]
换言之，在同一 Fibonacci 同余骨架上，满足上述三条条件的有限分辨率对象必然就是纯尺度 Wulff 射线。
\end{conjecture}
\begin{proof}[证据]
定理 \ref{thm:xi-time-part9xaab-wulff-ray-fourier-zero-formula}、
定理 \ref{thm:xi-time-part9xaab-wulff-ray-first-second-order-separation}
与定理 \ref{thm:xi-time-part9xaab-wulff-ray-fejer-kernel}
表明，纯尺度 Wulff 射线确实满足这三条条件。另一方面，bin-fold 虽与之共享一阶指数常数 \(\log(2/\varphi)\)，但定理 \ref{thm:xi-time-part9xaab-wulff-ray-first-second-order-separation} 表明其归一化最大纤维偏置始终高于 \(1\)，而定理 \ref{thm:xi-time-part9xaab-binfold-fejer-riesz-nonisomorphism} 又表明其零集结构与 Fej\'er 类严格不符。故 bin-fold 至少在第二、第三条上系统性失效。证毕。
\end{proof}

\noindent
综上，Figalli--Zhang 的单参数各向异性 Serrin 可实现类与本文的 bin-fold 模型虽然共享同一 Fibonacci 同余半环骨架，并在一阶恢复指数上同现 \(\log(2/\varphi)\)，但二者的频谱实现已在三条互不依赖的尺度上严格分离：纯尺度 Wulff 射线的非平凡零集只是一个由 \(\gcd(F_{m+2},2^m)\) 控制的单循环子群，其中心化谱能是离散 Fej\'er 核，且最大纤维相对均值比趋于 \(1\)；bin-fold 则保留半阶大陪集零点块、Fibonacci 频率驱动的有限 Riesz 乘积与常数级 \(L^\infty\) 偏置。由此，Wulff 射线在本文内部的角色被精确定位为同骨架下的纯尺度比较对象，而非 fold 异常的频谱来源。

\endinput


\endinput

\paragraph{fixed-resolution 全息、screen--Artin 双缺陷与 Wulff 凸税的终端综合}
在定理 \ref{thm:xi-time-part9x-fixed-resolution-stokes-godel-strong-completeness}--推论
\ref{cor:xi-time-part9x-screen-completion-cost-sharp-area-law}、
定理 \ref{thm:xi-time-part9x-nonabelian-artin-channel-quotient-erasure}--推论
\ref{cor:xi-time-part9x-abelianized-certificate-strict-insufficiency}、
定理 \ref{thm:xi-time-part9x-serrin-ray-fibonacci-semirings}--猜想
\ref{conj:xi-time-part9x-serrin-ray-fibonacci-universality}、
以及定理 \ref{thm:xi-time-part9xaa-wulff-schur-minimal-spectrum}--命题
\ref{prop:xi-time-part9xaa-audited-even-windows-above-wulff-floor} 已经分别建立之后，还可把这三条链再压缩为一组终端综合命题。它们共同表明：fixed-resolution 全息的实质性障碍并不在于 bulk 无法被精确编码，而在于 screen 赤字、non-inflation Artin 通道损失与低纤维强制 convex tax 这三类彼此正交的不可压缩界面。

\begin{theorem}[全边界完备与部分屏幕赤字的严格二分]\label{thm:xi-time-part9xaad-fullboundary-partialscreen-sharp-dichotomy}
固定 \(m,n\ge 1\)。记完整 boundary Gödel 全息码为
\[
\Phi(U):=(G\circ\partial_n)(u_U),
\]
并对任意屏幕 \(S\subset \overrightarrow{\mathcal F_m^{n-1}}\) 记部分读出为
\[
f_S:=\Pi_S\circ\partial_n.
\]
则成立严格二分：
\begin{enumerate}
  \item 全边界层上，\(\Phi(U)\) 在 \(O(1)\) 描述复杂度损失下唯一恢复 \(U\)，并唯一恢复有限矩盒
  \[
  \mathcal M_m(U)=\bigl(\mu_\alpha(U)\bigr)_{0\le \alpha_i\le 2^m-1};
  \]
  \item 部分屏幕层上，由 \(f_S\) 精确补全 bulk 所需的最小外置秩恰为
  \[
  \mathsf{Cost}(S)=\rank_\ZZ H_n(\mathcal Q_m,A_{\neg S};\ZZ)=c(\Gamma_S)-1.
  \]
\end{enumerate}
特别地，部分屏幕精确恢复当且仅当 \(c(\Gamma_S)=1\)；对一切 \(c(\Gamma_S)>1\) 的断裂屏幕，都存在严格正整数型相对同调赤字。
\end{theorem}
\begin{proof}
第 \(1\) 条正是定理 \ref{thm:xi-time-part9x-fixed-resolution-stokes-godel-strong-completeness} 的内容。第 \(2\) 条由定理 \ref{thm:xi-time-part9x-partial-screen-relative-homology-deficit-closed} 与推论 \ref{cor:xi-time-part9x-screen-completion-cost-sharp-area-law} 直接给出：
\[
\ker f_S\cong H_n(\mathcal Q_m,A_{\neg S};\ZZ)\cong \ZZ^{\,c(\Gamma_S)-1},
\qquad
\mathsf{Cost}(S)=\rank_\ZZ\ker f_S.
\]
因此 \(\mathsf{Cost}(S)=0\iff c(\Gamma_S)=1\)，而 \(c(\Gamma_S)>1\) 时赤字严格为正。证毕。
\end{proof}

\begin{corollary}[轴向屏幕的指数不可修补性]\label{cor:xi-time-part9xaad-axial-screen-exponential-irreparability}
在定理 \ref{thm:xi-time-part9xaad-fullboundary-partialscreen-sharp-dichotomy} 的设定下，对轴向屏幕 \(S^{(1)}\) 有
\[
\mathsf{Cost}(S^{(1)})=2^{m(n-1)}.
\]
因此任意随 \(m\to\infty\) 的轴向屏幕家族都不可能在统一有界外置秩下保持 exact bulk recovery。
\end{corollary}
\begin{proof}
推论 \ref{cor:xi-time-part9x-screen-completion-cost-sharp-area-law} 已给出
\[
\mathsf{Cost}(S^{(1)})=2^{m(n-1)}.
\]
故若要求一族轴向屏幕在全部 \(m\) 上以统一有界外置秩精确补全，则必与上式的指数增长矛盾。证毕。
\end{proof}

\begin{theorem}[screen--Artin 双缺陷的完备性判据]\label{thm:xi-time-part9xaad-screen-artin-double-defect-completeness}
设 \(S\subset \overrightarrow{\mathcal F_m^{n-1}}\)，并设
\[
\pi:G\twoheadrightarrow H
\]
为群标签核的一个满商。定义 Artin 缺陷簇
\[
\mathfrak A(\pi)
:=
\bigoplus_{\substack{
\rho\in\widehat G\setminus \mathrm{Inf}_\pi(\widehat H)\\
\det(I-zK_\rho)\not\equiv 1}}
\rho,
\]
以及双缺陷对象
\[
\mathfrak D(S,\pi)
:=
H_n(\mathcal Q_m,A_{\neg S};\ZZ)\oplus \mathfrak A(\pi).
\]
则由屏幕 \(S\) 与商标签 \(\pi_*K\) 构成的联合证书系统，对 fixed-resolution bulk 与完整 determinant 规律同时完备，当且仅当
\[
\mathfrak D(S,\pi)=0.
\]
\end{theorem}
\begin{proof}
若
\[
H_n(\mathcal Q_m,A_{\neg S};\ZZ)\neq 0,
\]
则由定理 \ref{thm:xi-time-part9xaad-fullboundary-partialscreen-sharp-dichotomy}，\(\ker f_S\neq 0\)，因而部分屏幕不可能唯一恢复 fixed-resolution bulk。

若存在
\[
\rho\in\widehat G\setminus \mathrm{Inf}_\pi(\widehat H),
\qquad
\det(I-zK_\rho)\not\equiv 1,
\]
则由定理 \ref{thm:xi-time-part9x-nonabelian-artin-channel-quotient-erasure}，该活跃通道在商核 determinant 分解中被完全擦除，故 \(\pi_*K\) 不能恢复完整 determinant。

反之，若 \(\mathfrak D(S,\pi)=0\)，则第一分量为零意味着 \(\ker f_S=0\)，从而屏幕侧对 bulk 完备；第二分量为零意味着一切 determinant-活跃通道都来自 \(\mathrm{Inf}_\pi(\widehat H)\)，故商核保留全部活跃 Artin 因子。两者合并即得联合证书完备。证毕。
\end{proof}

\begin{corollary}[阿贝尔化证书的完备性充要条件]\label{cor:xi-time-part9xaad-abelianized-certificate-iff-onedim}
设
\[
q:G\twoheadrightarrow G^{\mathrm{ab}}:=G/[G,G]
\]
为阿贝尔化商。则仅依赖 abelianized labels 的 Dirichlet/Artin 证书能够恢复完整 determinant 与全部 twisted primitive-orbit law，当且仅当所有 determinant-活跃不可约 Artin 通道都为一维表示。等价地，只要存在某个
\[
\rho\in\widehat G,
\qquad
\dim\rho>1,
\qquad
\det(I-zK_\rho)\not\equiv 1,
\]
则阿贝尔化证书必然严格不完备。
\end{corollary}
\begin{proof}
对阿贝尔群 \(G^{\mathrm{ab}}\)，全部不可约表示都是一维表示；并且任一 \(G\) 的一维表示都经由阿贝尔化因子化。故
\[
\mathrm{Inf}_q(\widehat{G^{\mathrm{ab}}})
\]
恰为 \(G\) 的全部一维不可约表示。于是由定理 \ref{thm:xi-time-part9xaad-screen-artin-double-defect-completeness}，阿贝尔化证书完备当且仅当不存在非 inflation 的活跃通道；这又当且仅当全部活跃通道都是一维的。证毕。
\end{proof}

\begin{theorem}[boundary 算术钉扎与 Weyl 完成的不可兼容性]\label{thm:xi-time-part9xaad-boundary-pinning-weyl-incompatibility}
令
\[
\mathcal B_6=\{14,17,19\}\subset \ZZ/(21\ZZ)
\]
为 window-\(6\) boundary 三点在算术骨架上的像。则：
\begin{enumerate}
  \item \(\mathcal B_6\) 既不是 \(\ZZ/(21\ZZ)\) 的加法陪集，也不是 \((\ZZ/(21\ZZ))^\times\) 中任何乘法子群的陪集；
  \item 任一非平凡 boundary Weyl 完成 \(Y\) 都满足
  \[
  |Y|\in\{9,15,21\},
  \qquad
  |Y|\ge 9.
  \]
\end{enumerate}
因此，boundary 的 arithmetic pinning 与 Weyl closure 属于两种互不相容的封闭机制：它既不能由单一同余陪集解释，也不能通过少于一条完整六点 Weyl 轨道的方式被连续补全。
\end{theorem}
\begin{proof}
第 \(1\) 条正是命题 \ref{prop:xi-time-part9x-window6-boundary-arithmetic-pinning}。第 \(2\) 条由定理 \ref{thm:xi-time-part9x-window6-boundary-weyl-threshold} 直接给出：一般 boundary 完成的规模只可能属于 \(\{3,9,15,21\}\)，而非平凡完成必须含非空 cyclic 轨道，故排除 \(3\) 即得结论。证毕。
\end{proof}

\begin{theorem}[局部化 Wulff 射线的连通观测一维刚性]\label{thm:xi-time-part9xaad-localized-wulff-connected-observer-rigidity}
固定有限素数集 \(S\subset\mathbb P\)，并记
\[
X_{K,S}:=\widehat{\mathfrak W_{K,S}^{\mathrm{gp}}}.
\]
则存在短正合列
\[
0\longrightarrow \prod_{p\in S}\ZZ_p
\longrightarrow X_{K,S}
\longrightarrow \TT
\longrightarrow 0,
\]
并且 \(X_{K,S}\) 的每一个非平凡连通紧 Lie 商都同构于 \(\TT\)。因此，一切真正的局部算术记忆都被压入 totally disconnected 的 \(p\)-adic 核，而任何 connected finite-dimensional observer 都只能看到单一圆相位。
\end{theorem}
\begin{proof}
短正合列正是定理 \ref{thm:xi-time-part9x-serrin-ray-localized-ledger} 的内容。记
\[
P_S:=\prod_{p\in S}\ZZ_p.
\]
由于 \(P_S\) 全不连通，其任一连通子群都只能是平凡群。设 \(q:X_{K,S}\twoheadrightarrow Q\) 为一个非平凡连通紧 Lie 商。则 \(q(X_{K,S}^\circ)=Q\)，且
\[
X_{K,S}^\circ\cap P_S=\{0\}.
\]
另一方面，短正合列中的商 \(\TT\) 连通，故 \(X_{K,S}^\circ\to \TT\) 亦为满同态。于是 \(X_{K,S}^\circ\cong \TT\)。因此 \(Q\) 是 \(\TT\) 的一个连通紧 Lie 商，只能是 \(\TT\) 本身。证毕。
\end{proof}

\begin{theorem}[Wulff 平衡谱的单泛函刚性]\label{thm:xi-time-part9xaad-wulff-single-functional-rigidity}
沿用定理 \ref{thm:xi-time-part9xaa-wulff-schur-minimal-spectrum} 的记号。则下列任一单一标量等式都足以推出
\[
d(T)\sim d_m^{W}
\]
（仅差排列）：
\[
\sum_i d_i^2=\sum_i \bigl(d_i^{W}\bigr)^2,
\]
\[
\log|\mathcal G(T)|=\log|\mathcal G_m^{W}|,
\]
\[
D_{\mathrm{KL}}(p\|u)=D_{\mathrm{KL}}(p^{W}\|u).
\]
因此，collision、group-volume 与 KL 缺口三种看似异质的审计器，在 Wulff 射线基准上给出同一个唯一极小谱。
\end{theorem}
\begin{proof}
这正是定理 \ref{thm:xi-time-part9xaa-single-functional-rigidity} 的内容。证毕。
\end{proof}

\begin{theorem}[低纤维异常税的线性量子化]\label{thm:xi-time-part9xaad-lowfiber-linear-quantization}
假设 \(q_m\ge 2\)，并固定
\[
0\le r\le \left\lfloor\frac{M_m-a_m}{2}\right\rfloor.
\]
则在所有满足“至少有 \(r\) 个坐标不高于 \(q_m-1\)”的 \(M_m\)-维正整数谱中，唯一严格凸极小元为
\[
d_m^{(r)}
=
\bigl((q_m-1)^r,\ q_m^{\,M_m-a_m-2r},\ (q_m+1)^{\,a_m+r}\bigr),
\]
并且相对 Wulff 基准的任意 separable 严格凸税额都精确量子化为
\[
\Delta_\Phi(r)
=
r\Bigl[\Phi(q_m-1)-2\Phi(q_m)+\Phi(q_m+1)\Bigr].
\]
在 \(m=6\) 时，\(q_6=3\)，故每增加一条二重纤维，单位增量恰为
\[
\Delta\!\Bigl(\sum_i d_i^2\Bigr)=2,
\qquad
\Delta\log|\mathcal G|=\log\frac43,
\]
\[
\Delta\kappa=\frac{1}{64}\log\frac{1024}{729},
\qquad
\Delta\mathsf{Col}=\frac{1}{2048}.
\]
因此，low-fiber defect 不是连续漂移量，而是按 \(r\) 严格线性量子化的 convex tax。
\end{theorem}
\begin{proof}
前半部分正是定理 \ref{thm:xi-time-part9xaa-lowfiber-anomaly-tax}。对 \(m=6\) 代入 \(q_6=3\) 即得
\[
\Phi(t)=t^2:\quad (2)^2-2\cdot 3^2+4^2=2,
\]
\[
\Phi(t)=\log\Gamma(t+1):\quad
\log\Gamma(3)-2\log\Gamma(4)+\log\Gamma(5)=\log\frac43,
\]
\[
\Phi(t)=t\log t:\quad
\Delta\kappa
=\frac1{64}\bigl(2\log2-6\log3+4\log4\bigr)
=\frac1{64}\log\frac{1024}{729}.
\]
再除以 \(N_6^2=64^2\) 即得
\[
\Delta\mathsf{Col}=\frac{2}{64^2}=\frac1{2048}.
\]
证毕。
\end{proof}

\begin{theorem}[window-\texorpdfstring{$6$}{6} 审计谱的受约束 Wulff 饱和]\label{thm:xi-time-part9xaad-window6-constrained-wulff-saturation}
\leanverified{paper\_xi\_time\_part9xaad\_window6\_constrained\_wulff\_saturation}
window-\(6\) 的实际审计直方图
\[
d_6^{\mathrm{bin}}=(2^8,3^4,4^9)
\]
恰是“总质量 \(64\)、长度 \(21\)、且恰有 \(8\) 条二重纤维”这一约束下的唯一严格凸极小元。故其全部异常税都不是经验偏移，而是由 \(r=8\) 强制出来的最小代价：
\[
\mathsf{Col}_6^{\mathrm{bin}}-\mathsf{Col}_6^{W}=\frac1{256},
\]
\[
\log|\mathcal G_6^{\mathrm{bin}}|-\log|\mathcal G_6^{W}|=8\log\frac43,
\]
\[
g_6^{\mathrm{bin}}-g_6^{W}=\frac18\log\frac43,
\]
\[
\kappa_6-\kappa_6^{W}=\frac14\log\frac{32}{27}.
\]
与此同时，对审计偶窗 \(m\in\{8,10,12\}\)，bin-fold 谱仍全部严格高于各自的 unconstrained Wulff--Schur 地板。因而 window-\(6\) 不是平衡射线的任意样本，而是一个受 \(r=8\) 约束精确饱和的异常端点。
\end{theorem}
\begin{proof}
唯一极小元这一断言正是定理 \ref{thm:xi-time-part9xaa-window6-audit-histogram-variational-characterization}。四条精确差额由定理 \ref{thm:xi-time-part9xaa-window6-audited-spectrum-exact-tax} 直接给出。对 \(m=8,10,12\) 的严格正缺口，则由命题 \ref{prop:xi-time-part9xaa-audited-even-windows-above-wulff-floor} 给出。证毕。
\end{proof}

\begin{conjecture}[单参数各向异性 Serrin 观测的 Fibonacci 终型刚性]\label{conj:xi-time-part9xaad-serrin-fibonacci-terminal-rigidity}
作为猜想 \ref{conj:xi-time-part9x-serrin-ray-fibonacci-universality} 的终型表述，设
\[
q_m:\mathfrak W_K^{(0)}\twoheadrightarrow Q_m
\qquad (m\ge 1)
\]
是一族单参数各向异性 Serrin 有限分辨率观测，并满足：
\begin{enumerate}
  \item
  \[
  |Q_m|=F_{m+2};
  \]
  \item \(q_m\) 精确实现 Minkowski 和与尺度复合的半环结构；
  \item 其 Grothendieck 完成所诱导的同余拓扑与 Fibonacci 同余拓扑一致。
\end{enumerate}
则必存在唯一半环同构
\[
Q_m\cong \ZZ/(F_{m+2}\ZZ)
\qquad (m\ge 1).
\]
若该猜想成立，则下列两条性质将成为普适结论而非模型特例：
\[
\lim_{m\to\infty}\frac{B_m^{\mathrm{exact}}}{m}
=
\log_2\!\Bigl(\frac{2}{\varphi}\Bigr),
\]
并且当 \(m+2\equiv 0\pmod 6\) 时，
\[
\bigl|\mathcal Z_{m,K}^{\mathrm{arith}}\bigr|
\ge
F_{(m+2)/2}.
\]
\end{conjecture}

\noindent
综上，前述链条可被压缩为三条终端结构：
\[
\begin{aligned}
&\text{full boundary Gödel holography}
\Longrightarrow
\text{exact bulk recovery},\\
&\text{partial screens}
\Longrightarrow
\text{exact deficit }c(\Gamma_S)-1,\\
&\text{screen}+\text{quotient}
\Longrightarrow
\text{screen--Artin dual obstruction};
\end{aligned}
\]
\[
\begin{aligned}
&\mathfrak W_K^{(0)}/\!\equiv_m^{(K)}
\cong
\ZZ/(F_{m+2}\ZZ)
\Longrightarrow
d_m^{W}\text{ 为唯一 Schur 地板},\\
&\text{low-fiber defect}
\Longrightarrow
\text{all convex taxes are quantized},\\
&d_6^{\mathrm{bin}}
\Longrightarrow
\text{the exact constrained minimizer with }r=8;
\end{aligned}
\]
\[
\begin{aligned}
&\mathcal B_6=\{14,17,19\}
\text{ 既非加法陪集亦非乘法陪集},\\
&\text{nontrivial Weyl completion}
\Longrightarrow
|Y|\ge 9,\\
&\widehat{\mathfrak W_{K,S}^{\mathrm{gp}}}
\Longrightarrow
\text{its only nontrivial connected Lie quotient is }\TT.
\end{aligned}
\]
因而，fixed-resolution 全息在本文对象中的真正不可压缩界面，并不位于 bulk 是否可编码这一层，而位于 screen 赤字、non-inflation Artin 通道损失与低纤维 convex tax 这三条彼此正交的终端障碍。

\endinput

\paragraph{等变 completion 的刚性下界与 visible sector 的零模图像}
在可见条件期望的有限指数谱、二比特精确反演与十九比特异常裂分已经确定之后，还可进一步把局部标签的对称代价及其对连续极限零模结构的约束组织为下列三条结果。其核心现象是：window-\(6\) 的 hidden sector 虽然可被有限 bit 精确解码，但一旦要求保留 boundary/full-center 对称，局部 completion 的开销便出现严格抬升。

\begin{theorem}[对称保持 completion 的刚性比特下界]\label{thm:xi-time-part9x-equivariant-completion-bit-lower-bound}
\leanverified{paper\_xi\_time\_part9x\_equivariant\_completion\_bit\_lower\_bound}
沿用定理 \ref{thm:xi-time-part9x-visible-condexp-finite-index-spectrum} 的记号。
对整数 \(B\ge 0\)，记
\[
Y_B:=\{0,1\}^B.
\]
称映射
\[
c:\Omega_6\to Y_B
\]
为一个 \(B\)-bit completion，若
\[
(\pi,c):\Omega_6\longrightarrow X_6\times Y_B
\]
为单射。再称其对 \(H\le G_6^{\mathrm{bin}}\) 等变，若存在群同态
\[
\rho:H\to \mathfrak S_{2^B}
\]
使
\[
c(h\omega)=\rho(h)c(\omega)
\qquad
(\forall\, h\in H,\ \omega\in \Omega_6).
\]
若 \(H\cong (\ZZ/2\ZZ)^r\) 为初等阿贝尔 \(2\)-群，且 \(c\) 是一个 \(H\)-等变 completion，则必有
\[
r\le 2^{B-1}.
\]
特别地：
\begin{enumerate}
  \item 对 boundary 中心
  \[
  Z_{\partial}\cong (\ZZ/2\ZZ)^3,
  \]
  任意 \(Z_{\partial}\)-等变 completion 都满足
  \[
  B\ge 3;
  \]
  \item 对 full center
  \[
  Z\!\bigl(G_6^{\mathrm{bin}}\bigr)\cong (\ZZ/2\ZZ)^8,
  \]
  任意 \(Z(G_6^{\mathrm{bin}})\)-等变 completion 都满足
  \[
  B\ge 4.
  \]
\end{enumerate}
\end{theorem}
\begin{proof}
先证 \(\rho\) 必忠实。若 \(1\neq h\in H\)，则存在某个 \(\omega\in \Omega_6\) 使 \(h\omega\neq \omega\)。又因 \(H\le G_6^{\mathrm{bin}}=\Aut(\pi)\)，有
\[
\pi(h\omega)=\pi(\omega).
\]
由于 \((\pi,c)\) 为单射，遂必有
\[
c(h\omega)\neq c(\omega).
\]
而
\[
c(h\omega)=\rho(h)c(\omega),
\]
因此 \(\rho(h)\neq 1\)。故 \(\rho\) 为忠实表示。

现将 \(Y_B\) 分解为 \(H\)-轨道
\[
Y_B=\bigsqcup_j O_j,
\qquad
|O_j|=2^{k_j}.
\]
这里 \(k_j\ge 0\)，因为 \(H\) 是 \(2\)-群，故每个轨道大小必为 \(2\) 的幂。对每个 \(j\)，记
\[
K_j:=\ker\bigl(H\curvearrowright O_j\bigr),
\qquad
H_j:=H/K_j.
\]
则 \(H_j\) 在 \(O_j\) 上忠实且传递。由于 \(H_j\) 阿贝尔，任一点稳定子都是正规子群；而传递性使所有稳定子相同，忠实性又迫使其交为平凡，因此稳定子只能是平凡群。故 \(H_j\) 在 \(O_j\) 上正则作用，从而
\[
|H_j|=|O_j|=2^{k_j},
\qquad
\rank_{\FF_2}H_j=k_j.
\]

另一方面，由 \(\rho\) 的忠实性，
\[
\ker \rho=\bigcap_j K_j=\{1\}.
\]
因此自然映射
\[
H\longrightarrow \prod_j H_j
\]
为单射，于是
\[
r=\rank_{\FF_2}H\le \sum_j \rank_{\FF_2}H_j=\sum_j k_j.
\]
对每个 \(k_j\ge 1\)，有初等不等式
\[
k_j\le 2^{k_j-1},
\]
而当 \(k_j=0\) 时两边都为 \(0\)。故
\[
r\le \sum_j k_j\le \sum_j 2^{k_j-1}
=
\frac12\sum_j |O_j|
=
\frac12\,|Y_B|
=
2^{B-1}.
\]
一般下界得证。

最后，把
\[
r=3
\qquad\text{与}\qquad
r=8
\]
分别代入即可得到
\[
3\le 2^{B-1}\Longrightarrow B\ge 3,
\qquad
8\le 2^{B-1}\Longrightarrow B\ge 4.
\]
其中 boundary 中心与 full center 的秩分别由定理
\ref{thm:xi-window6-boundary-parity-canonical-direct-summand} 与定理
\ref{thm:xi-window6-center-exact-splitting-bdry-cyclic} 给出。证毕。
\end{proof}
\leanverified{paper\_xi\_time\_part9x\_equivariant\_completion\_bit\_lower\_bound}

\begin{corollary}[不存在 boundary-equivariant 的四值伴随场]\label{cor:xi-time-part9x-no-boundary-equivariant-four-valued-companion-field}
不存在任何映射
\[
c:\Omega_6\to \{1,2,3,4\}
\]
使
\[
(\pi,c):\Omega_6\to X_6\times \{1,2,3,4\}
\]
为单射，且 \(c\) 同时对 boundary sheet parity 中心
\[
Z_{\partial}\cong (\ZZ/2\ZZ)^3
\]
等变。换言之，在保留 canonical boundary parity 的前提下，window-\(6\) 的 hidden sector 不能被压缩成一个四值局域伴随场。
\end{corollary}
\begin{proof}
四值标签等价于 \(B=2\) 的 completion。若它还要求 \(Z_{\partial}\)-等变，则定理
\ref{thm:xi-time-part9x-equivariant-completion-bit-lower-bound} 强制
\[
B\ge 3.
\]
这与 \(B=2\) 矛盾，故所述四值伴随场不存在。证毕。
\end{proof}
\leanverified{paper\_xi\_time\_part9x\_no\_boundary\_equivariant\_four\_valued\_companion\_field}

\begin{conjecture}[有限指数--对称溢价原则：GFF 质量零模仅驻留于 visible sector]\label{conj:xi-time-part9x-finite-index-symmetry-premium-visible-gff}
沿用定理 \ref{thm:xi-time-part9x-visible-condexp-finite-index-spectrum} 的条件期望 \(E_6\)。
设某一临界六顶点系综中，高度函数在适当归一化下收敛到标量高斯自由场。记 \(\mathcal O_\delta\) 为任意由 window-\(6\) 提升得到的局部可观测量，并定义
\[
\mathcal O_\delta^{\mathrm{vis}}:=E_6(\mathcal O_\delta),
\qquad
\mathcal O_\delta^{\mathrm{hid}}:=\mathcal O_\delta-\mathcal O_\delta^{\mathrm{vis}}.
\]
若 \(\mathcal O_\delta\) 的构造保持 canonical boundary parity 对称，则在任何使 visible 高度涨落保持非平凡的 GFF 归一化下，应有：
\begin{enumerate}
  \item \(\mathcal O_\delta^{\mathrm{vis}}\) 与六顶点高度场产生同一质量零模；
  \item \(\mathcal O_\delta^{\mathrm{hid}}\) 只能贡献 tight 的有限内部标签涨落，而不能再生新的独立高斯自由场；
  \item 任何试图把 \(\mathcal O_\delta^{\mathrm{hid}}\) 解释为局域四值伴随场的方案，都会被推论 \ref{cor:xi-time-part9x-no-boundary-equivariant-four-valued-companion-field} 在离散层面排除。
\end{enumerate}
\end{conjecture}
\begin{proof}[证据]
定理 \ref{thm:xi-time-part9x-visible-condexp-finite-index-spectrum} 给出
\[
\mathrm{Ind}_{\mathrm{PP}}(E_6)=4,
\]
而定理 \ref{thm:xi-time-part9x-reversible-capacity-two-bit-inversion} 说明精确反演只需
\[
b_{\mathrm{inv}}=2
\]
个 bit。换言之，裸的可逆性仅由最大纤维顶谱控制。

然而一旦要求保持对称，定理
\ref{thm:xi-time-part9x-equivariant-completion-bit-lower-bound} 立即把局部 completion 的最小开销抬升为
\[
B\ge 3
\quad\text{(boundary parity)},
\qquad
B\ge 4
\quad\text{(full center)}.
\]
再结合定理 \ref{thm:xi-time-part9x-decoding-anomaly-nineteen-bit-gap} 的
\[
\dim_{\FF_2}\bigl(G_6^{\mathrm{bin}}\bigr)^{\mathrm{ab}}=21,
\]
可见 hidden sector 的交换化异常与局部可逆标签之间存在显著尺度裂分。

因此，window-\(6\) 的 hidden sector 在有限尺度上表现为一个有限指数、有限标签而且具有严格对称溢价的部门：它可以被精确解码，但任何保持 boundary/full-center 结构的实现都必须支付额外局部比特成本。这一图像与“hidden 方向只留下 tight 的内部标签涨落，而不再产生新的独立质量零模”相一致；相反，真正的质量零模应当集中在由 \(E_6\) 保留下来的 visible sector 中。证毕。
\end{proof}

\noindent
综上，window-\(6\) 的可见条件期望在交换函数代数侧给出了一条与群胚代数中心指数元平行而互补的结构链：
\[
E_6
\Longrightarrow
\mathrm{Ind}_{\mathrm W}(E_6)\ \text{的谱 } \{2,3,4\}
\Longrightarrow
\kappa_6\ \text{与}\ D_{\mathrm{KL}}(\mu_6\Vert u_6)
\Longrightarrow
b_{\mathrm{inv}}=2
\Longrightarrow
b_{\mathrm{anom}}-b_{\mathrm{inv}}=19
\Longrightarrow
\text{对称保持 completion 的刚性比特下界}.
\]
由此，hidden sector 的有限性并不意味着它在所有语义下都可由低比特局域标签实现；恰相反，window-\(6\) 在可逆性与对称性之间展示出严格可审计的溢价裂分。

\endinput

\paragraph{window-\(6\) hidden fiber 的 cubical 规范形、Bernoulli 因子化与 UV alias 超平面}
在定理 \ref{thm:xi-time-part9x-visible-condexp-finite-index-spectrum} 与定理 \ref{thm:xi-time-part9x-equivariant-completion-bit-lower-bound} 已经把 window-\(6\) 的 visible/hidden 裂分、局部可逆容量与对称 completion 溢价固定之后，还可仅凭 \(m=6\) 的显式审计数据，把 hidden fiber 的局部组合形、Fourier 零模与连续极限下的 alias 机制压缩为另一条有限尺度闭链。以下沿用
\[
\Omega_6:=\{0,\dots,63\},
\qquad
\pi:=\Fold_6^{\mathrm{bin}}:\Omega_6\to X_6,
\qquad
F_w:=\pi^{-1}(w),
\qquad
d(w):=|F_w|,
\]
并采用全文的稳定词记号
\[
X_m:=\{u\in\{0,1\}^m:\ u_i u_{i+1}=0\ \ (1\le i<m)\}.
\]

\begin{theorem}[window-\texorpdfstring{$6$}{6} 最小退化扇区的 rigid \texorpdfstring{$X_4$}{X4} 子覆盖]\label{thm:xi-time-part9xab-window6-minsector-rigid-x4-subcover}
\leanverified{paper\_xi\_time\_part9xab\_window6\_minsector\_rigid\_x4\_subcover}
定义
\[
S_{6,2}:=\{w\in X_6:\ d(w)=2\},
\qquad
\iota_4:X_4\hookrightarrow X_6,
\qquad
\iota_4(v):=(v,0,1).
\]
则有精确恒等式
\[
S_{6,2}=\iota_4(X_4).
\]
换言之，window-\(6\) 的全部最小退化态恰由后缀固定为 \(01\) 的稳定词组成，并与 \(X_4\) 逐词同构。
\end{theorem}
\begin{proof}
由定理 \ref{thm:xi-foldbin-zeckendorf-monotone-staircase-law} 在 \(m=6\) 时的三档公式，
\[
d(w)=2
\quad\Longleftrightarrow\quad
w_6=1.
\]
若 \(w\in X_6\) 满足 \(w_6=1\)，则稳定性约束 \(w_5w_6=0\) 强制 \(w_5=0\)，故 \(w\) 唯一写成 \(w=(v,0,1)\)，其中前四位 \(v\in X_4\)。反之，任取 \(v\in X_4\)，由于 \(v_4\cdot 0=0\) 且 \(0\cdot 1=0\)，故 \((v,0,1)\) 仍属 \(X_6\)。于是
\[
\{w\in X_6:\ w_6=1\}=\iota_4(X_4).
\]
与上式合并即得
\[
S_{6,2}=\iota_4(X_4).
\]
证毕。
\end{proof}

\begin{theorem}[window-\texorpdfstring{$6$}{6} hidden fiber 的仿射 Boolean 方块标准形]\label{thm:xi-time-part9xab-window6-hidden-fiber-affine-boolean-normal-form}
\leanverified{paper\_xi\_time\_part9xab\_window6\_hidden\_fiber\_affine\_boolean\_normal\_form}
定义
\[
\beta:\{0,1\}^2\to \ZZ,
\qquad
\beta(a,b):=21a+34b,
\]
并令
\[
I_2:=\{(0,0),(0,1)\},
\qquad
I_3:=\{(0,0),(1,0),(0,1)\},
\qquad
I_4:=\{0,1\}^2.
\]
若对任意 \(J\subseteq \{0,1\}^2\) 记
\[
\beta(J):=\{\beta(u):\ u\in J\},
\]
则对每个 \(w\in X_6\) 都有
\[
F_w=V_6(w)+\beta(I_{d(w)}).
\]
特别地，
\[
\beta(I_2)=\{0,34\},
\qquad
\beta(I_3)=\{0,21,34\},
\qquad
\beta(I_4)=\{0,21,34,55\},
\]
并且 \(I_2,I_3,I_4\) 都是 Boolean 方块 \(\{0,1\}^2\) 在乘积序下的序理想。
\end{theorem}
\begin{proof}
由定理 \ref{thm:xi-foldbin-zeckendorf-monotone-staircase-law} 的三档分类与表 \ref{tab:fold6_bin_uplift_choice_collapse}，
\[
F_w=
\begin{cases}
\{V_6(w),\,V_6(w)+34\}, & d(w)=2,\\[4pt]
\{V_6(w),\,V_6(w)+21,\,V_6(w)+34\}, & d(w)=3,\\[4pt]
\{V_6(w),\,V_6(w)+21,\,V_6(w)+34,\,V_6(w)+55\}, & d(w)=4.
\end{cases}
\]
又因为
\[
55=21+34=\beta(1,1),
\]
故三类偏移集合恰分别等于 \(\beta(I_2),\beta(I_3),\beta(I_4)\)，从而
\[
F_w=V_6(w)+\beta(I_{d(w)}).
\]
最后，在乘积序 \((a,b)\le (a',b')\iff a\le a',\ b\le b'\) 下，
\[
I_2=\{u\in\{0,1\}^2:\ u\le (0,1)\},
\]
\[
I_3=\{u\in\{0,1\}^2:\ u\le (1,0)\ \text{或}\ u\le (0,1)\},
\qquad
I_4=\{u\in\{0,1\}^2:\ u\le (1,1)\},
\]
均为下闭集，故皆为序理想。证毕。
\end{proof}

\begin{corollary}[任意有限体积下 hidden fiber 的全局可缩性]\label{cor:xi-time-part9xab-window6-hidden-fiber-global-contractibility}
对任意有限指标集 \(\Lambda\)，定义
\[
\pi_\Lambda:=\pi^\Lambda:\Omega_6^\Lambda\to X_6^\Lambda.
\]
固定 \(w=(w_x)_{x\in\Lambda}\in X_6^\Lambda\)，并定义 cubical realization
\[
\mathfrak F(w):=\prod_{x\in\Lambda} I_{d(w_x)}
\subset ([0,1]^2)^\Lambda.
\]
则 \(\pi_\Lambda^{-1}(w)\) 与 \(\mathfrak F(w)\) 的顶点集存在规范双射，而且 \(\mathfrak F(w)\) 可缩。因而
\[
\widetilde H_k(\mathfrak F(w);\ZZ)=0
\qquad (k\ge 0).
\]
\end{corollary}
\begin{proof}
由定理 \ref{thm:xi-time-part9xab-window6-hidden-fiber-affine-boolean-normal-form}，
\[
\pi_\Lambda^{-1}(w)
=
\prod_{x\in\Lambda}\bigl(V_6(w_x)+\beta(I_{d(w_x)})\bigr).
\]
由于 \(\beta\) 在 \(I_2,I_3,I_4\) 上均为单射，故每个 \(n_x\in V_6(w_x)+\beta(I_{d(w_x)})\) 唯一对应于某个顶点 \(u_x\in I_{d(w_x)}\)。逐坐标合并，得到 \(\pi_\Lambda^{-1}(w)\) 与 \(\mathfrak F(w)\) 顶点集之间的规范双射。

再者，把 \(I_2,I_3,I_4\) 视为 \([0,1]^2\) 中的有限 cubical complex，则 \(I_2\) 是闭区间，\(I_3\) 是一个 \(L\) 形树，\(I_4\) 是单位正方形，三者都可缩。有限乘积仍可缩，因此 \(\mathfrak F(w)\) 可缩。可缩空间的约化同调全消失，故
\[
\widetilde H_k(\mathfrak F(w);\ZZ)=0
\qquad (k\ge 0).
\]
证毕。
\end{proof}

\begin{theorem}[四点纤维的 rank-two Bernoulli 因子化]\label{thm:xi-time-part9xab-window6-fourpoint-fiber-ranktwo-bernoulli-factorization}
\leanverified{paper\_xi\_time\_part9xab\_window6\_fourpoint\_fiber\_ranktwo\_bernoulli\_factorization}
对任意 \(w\in X_6\)，定义局部偏移 Laplace 变换
\[
L_w(t):=\frac1{d(w)}
\sum_{n\in F_w} e^{\,t(n-V_6(w))}.
\]
则有
\[
L_w(t)=\frac{1+e^{34t}}{2}
\qquad (d(w)=2),
\]
\[
L_w(t)=\frac{1+e^{21t}+e^{34t}}{3}
\qquad (d(w)=3),
\]
\[
L_w(t)=\frac{1+e^{21t}+e^{34t}+e^{55t}}{4}
=
\frac{1+e^{21t}}2\cdot \frac{1+e^{34t}}2
\qquad (d(w)=4).
\]
因此当 \(d(w)=4\) 时，均匀 hidden 偏移分布与
\[
21B_1+34B_2,
\qquad
B_1,B_2\stackrel{\mathrm{i.i.d.}}{\sim}\mathrm{Bernoulli}\!\left(\tfrac12\right),
\]
完全同分布。
\end{theorem}
\begin{proof}
由定理 \ref{thm:xi-time-part9xab-window6-hidden-fiber-affine-boolean-normal-form}，
\[
F_w-V_6(w)=
\begin{cases}
\{0,34\}, & d(w)=2,\\
\{0,21,34\}, & d(w)=3,\\
\{0,21,34,55\}, & d(w)=4.
\end{cases}
\]
把这三种偏移集合代入 \(L_w(t)\) 的定义，即得前三个闭式。

当 \(d(w)=4\) 时，
\[
\{0,21,34,55\}=\{0,21\}+\{0,34\},
\]
并且四个和都恰好出现一次。因此均匀分布于 \(\{0,21,34,55\}\) 的随机变量，与两个独立公平 Bernoulli 变量经线性映射
\[
(a,b)\longmapsto 21a+34b
\]
的像分布完全一致。Laplace 变换遂分解为
\[
\frac{1+e^{21t}}2\cdot \frac{1+e^{34t}}2.
\]
证毕。
\end{proof}

\begin{theorem}[window-\texorpdfstring{$6$}{6} 局部 Fourier 乘子的精确零模集]\label{thm:xi-time-part9xab-window6-local-fourier-multipliers-exact-zero-set}
\leanverified{paper\_xi\_time\_part9xab\_window6\_local\_fourier\_multipliers\_exact\_zero\_set}
令
\[
\omega:=e^{2\pi i/64}.
\]
对三类局部模板定义归一化 Fourier 乘子
\[
m_2(t):=\frac{1+\omega^{34t}}2,
\qquad
m_3(t):=\frac{1+\omega^{21t}+\omega^{34t}}3,
\]
\[
m_4(t):=\frac{1+\omega^{21t}+\omega^{34t}+\omega^{55t}}4
=
\frac{(1+\omega^{21t})(1+\omega^{34t})}{4},
\qquad
t\in \ZZ/64\ZZ.
\]
则其零点集精确为
\[
Z(m_2)=\{16,48\},
\qquad
Z(m_3)=\varnothing,
\qquad
Z(m_4)=\{16,32,48\}.
\]
\end{theorem}
\begin{proof}
先看 \(m_2\)。有
\[
m_2(t)=0
\quad\Longleftrightarrow\quad
\omega^{34t}=-1
\quad\Longleftrightarrow\quad
34t\equiv 32 \pmod{64}.
\]
两边同除以 \(2\)，得到
\[
17t\equiv 16 \pmod{32}.
\]
由于 \(17^{-1}\equiv 17\pmod{32}\)，故
\[
t\equiv 16 \pmod{32},
\]
于是
\[
Z(m_2)=\{16,48\}.
\]

再看 \(m_4\)。由因子化公式，
\[
m_4(t)=0
\quad\Longleftrightarrow\quad
\omega^{21t}=-1\ \text{或}\ \omega^{34t}=-1.
\]
第二种情形已给出 \(t\in\{16,48\}\)。第一种情形等价于
\[
21t\equiv 32 \pmod{64}.
\]
由于 \(21^{-1}\equiv 61\pmod{64}\)，遂得
\[
t\equiv 61\cdot 32\equiv 32 \pmod{64}.
\]
故
\[
Z(m_4)=\{16,32,48\}.
\]

最后考察 \(m_3\)。若 \(m_3(t)=0\)，则存在单位复数
\[
\zeta:=\omega^{21t},
\qquad
\eta:=\omega^{34t}
\]
满足
\[
1+\zeta+\eta=0.
\]
于是 \(\eta=-(1+\zeta)\)，从而
\[
1=|\eta|^2=|1+\zeta|^2=2+2\Re(\zeta).
\]
故
\[
\Re(\zeta)=-\frac12.
\]
这迫使 \(\zeta\) 为本原三次单位根。然而 \(\zeta\) 同时又是 \(64\) 次单位根，其阶数必须整除 \(64\)，与 \(3\nmid 64\) 矛盾。因此 \(m_3\) 无零点，即
\[
Z(m_3)=\varnothing.
\]
证毕。
\end{proof}

\begin{theorem}[全局 hidden averaging 的 alias 超平面]\label{thm:xi-time-part9xab-window6-hidden-averaging-alias-hyperplanes}
固定有限指标集 \(\Lambda\) 与可见构型 \(w=(w_x)_{x\in\Lambda}\in X_6^\Lambda\)。把每个全局原像 \((n_x)_{x\in\Lambda}\in \pi_\Lambda^{-1}(w)\) 唯一写成
\[
n_x=V_6(w_x)+\delta_x,
\qquad
\delta_x\in \beta(I_{d(w_x)}).
\]
定义 offsets 上的均匀平均算子
\[
E_w f
:=
\frac{1}{\prod_{x\in\Lambda} d(w_x)}
\sum_{(\delta_x)_x\in \prod_{x\in\Lambda}\beta(I_{d(w_x)})}
f\bigl((\delta_x)_{x\in\Lambda}\bigr).
\]
对任意 Fourier 角色
\[
\chi_t(\delta):=
\exp\!\Bigl(\frac{2\pi i}{64}\sum_{x\in\Lambda} t_x\delta_x\Bigr),
\qquad
t=(t_x)_{x\in\Lambda}\in (\ZZ/64\ZZ)^\Lambda,
\]
有
\[
E_w\chi_t=\prod_{x\in\Lambda} m_{d(w_x)}(t_x),
\]
其中 \(m_2,m_3,m_4\) 由定理 \ref{thm:xi-time-part9xab-window6-local-fourier-multipliers-exact-zero-set} 定义。因而精确零模集合
\[
\mathcal Z_w:=\{t\in (\ZZ/64\ZZ)^\Lambda:\ E_w\chi_t=0\}
\]
满足
\[
\mathcal Z_w
=
\bigcup_{x:\,d(w_x)=2}\{t:\ t_x\in\{16,48\}\}
\ \cup\
\bigcup_{x:\,d(w_x)=4}\{t:\ t_x\in\{16,32,48\}\}.
\]
特别地，三点纤维不贡献任何精确零超平面。
\end{theorem}
\begin{proof}
由定理 \ref{thm:xi-time-part9xab-window6-hidden-fiber-affine-boolean-normal-form}，全局 hidden fiber 按坐标分解为直积
\[
\pi_\Lambda^{-1}(w)-V_6(w)
=
\prod_{x\in\Lambda}\beta(I_{d(w_x)}),
\]
其中 \(V_6(w):=(V_6(w_x))_{x\in\Lambda}\)。因此对角色 \(\chi_t\) 的均匀平均可完全分离：
\[
\begin{aligned}
E_w\chi_t
&=
\frac{1}{\prod_{x\in\Lambda} d(w_x)}
\sum_{(\delta_x)_x\in \prod_{x\in\Lambda}\beta(I_{d(w_x)})}
\prod_{x\in\Lambda}
\exp\!\Bigl(\frac{2\pi i}{64}t_x\delta_x\Bigr)\\
&=
\prod_{x\in\Lambda}
\left(
\frac1{d(w_x)}
\sum_{\delta_x\in \beta(I_{d(w_x)})}
e^{2\pi i t_x\delta_x/64}
\right)\\
&=
\prod_{x\in\Lambda} m_{d(w_x)}(t_x).
\end{aligned}
\]
于是 \(E_w\chi_t=0\) 当且仅当某个坐标因子为零。再将定理 \ref{thm:xi-time-part9xab-window6-local-fourier-multipliers-exact-zero-set} 的精确零点集逐坐标代入，即得
\[
\mathcal Z_w
=
\bigcup_{x:\,d(w_x)=2}\{t:\ t_x\in\{16,48\}\}
\ \cup\
\bigcup_{x:\,d(w_x)=4}\{t:\ t_x\in\{16,32,48\}\}.
\]
而 \(Z(m_3)=\varnothing\)，故三点纤维不贡献任何精确零超平面。证毕。
\end{proof}

\begin{conjecture}[sitewise window-\texorpdfstring{$6$}{6} lift 的 UV alias 原理]\label{conj:xi-time-part9xab-window6-lift-uv-alias-principle}
沿用猜想 \ref{conj:xi-time-part9x-finite-index-symmetry-premium-visible-gff} 的设定，设某一临界六顶点系综中 visible 高度场在适当归一化下收敛到标量高斯自由场。则任一 canonical sitewise window-\(6\) lift 的 hidden averaging 在连续极限中只可能留下两类效应：
\[
\text{(i) 有限局域熵势或接触项重整化},\qquad
\text{(ii) 仅支撑于 } \{16,32,48\}\text{ 的 UV alias 缺口}.
\]
更精确地，hidden averaging 不会创造新的红外零模，也不会改变 visible sector 所决定的极限 Green 核；它在离散频域中的全部精确零模，均由定理 \ref{thm:xi-time-part9xab-window6-hidden-averaging-alias-hyperplanes} 的坐标超平面生成，而三点纤维对应的局部模板在频域上保持严格零缺口。
\end{conjecture}
\begin{proof}[证据]
推论 \ref{cor:xi-time-part9xab-window6-hidden-fiber-global-contractibility} 表明，任意有限体积上的 hidden fiber 都可由若干可缩 cubical 因子的乘积实现，从而纤维本体不携带新的局域高阶拓扑。定理 \ref{thm:xi-time-part9xab-window6-fourpoint-fiber-ranktwo-bernoulli-factorization} 又把全部四点纤维精确裂解为两个独立 Bernoulli 模，而定理 \ref{thm:xi-time-part9xab-window6-local-fourier-multipliers-exact-zero-set} 说明：局部精确零模只可能出现在 \(d=2\) 与 \(d=4\) 模板上，频率严格限制在
\[
\{16,48\}\quad\text{与}\quad \{16,32,48\},
\]
其中 \(d=3\) 模板完全零自由。定理 \ref{thm:xi-time-part9xab-window6-hidden-averaging-alias-hyperplanes} 进一步把这一现象推广到任意有限体积：全部精确零模只是坐标方向上的 alias 超平面，与体积增长无关，也不产生新的低频连续族。

因此，在猜想 \ref{conj:xi-time-part9x-finite-index-symmetry-premium-visible-gff} 所要求的 GFF 归一化下，hidden averaging 的频域效应应当只能停留在离散 UV 层，并在连续极限中塌缩为局域接触型修正；真正控制红外质量零模的仍是 visible sector。证毕。
\end{proof}

\noindent
由此，window-\(6\) 的 hidden sector 不仅在纤维计数上是有限的，而且在局部组合与频域上都被完全 rigidify：最小二重点层恰是 \(X_4\) 的 rigid 后缀子覆盖；任一局部纤维都由二维 Boolean 方块的仿射序理想给出；四点纤维裂解为两条独立 Bernoulli 模，而全部精确频域零模又只落在 \(\{16,32,48\}\) 的 alias 超平面上。于是，若 visible sector 的高斯自由场极限成立，则 hidden averaging 的剩余作用便应当只能停留在局域 UV 修正层面。

\endinput

\paragraph{window-\texorpdfstring{$6$}{6} hidden sector 的两寄存器局域气体、唯一 \texorpdfstring{$D_2$}{D2} 切片相互作用与 Gaussian 相容判据}
在定理 \ref{thm:xi-time-part9xab-window6-hidden-fiber-affine-boolean-normal-form}、定理 \ref{thm:xi-time-part9xab-window6-fourpoint-fiber-ranktwo-bernoulli-factorization} 与猜想 \ref{conj:xi-time-part9xab-window6-lift-uv-alias-principle} 已经把 window-\(6\) 的 hidden fiber 压缩为二维 Boolean/cubical 规范形、rank-\(2\) Bernoulli 裂解与 UV alias 超平面之后，还可把同一 hidden sector 再进一步组织为一条更细的局域气体链。其要点是：每个站点的 hidden offset 实际只由两枚二值寄存器控制，而三档纤维 \(d=2,3,4\) 仅对应三种 on-site 势垒；其中唯一非乘积、唯一带奇阶累积量的局域 defect，恰全部压缩在 \(B_3\) 字典的单个 \(D_2=A_1\times A_1\) 长根切片上。由此，window-\(6\) 与六顶点/GFF 接口可被进一步精炼为一组完全显式的必要判据。

\begin{theorem}[window-\texorpdfstring{$6$}{6} hidden sector 的两寄存器正规形]\label{thm:xi-time-part9xaba-window6-hidden-two-register-normal-form}
\leanverified{paper\_xi\_time\_part9xaba\_window6\_hidden\_two\_register\_normal\_form}
对固定 \(w\in X_6\)，在纤维 \(F_w\) 上取均匀随机元 \(n\)，并定义局部 hidden offset
\[
\Delta_w:=n-V_6(w).
\]
再定义
\[
\mathcal A_2:=\{(0,0),(0,1)\},\qquad
\mathcal A_3:=\{(0,0),(1,0),(0,1)\},\qquad
\mathcal A_4:=\{0,1\}^2.
\]
则存在一对二值寄存器
\[
(\eta_w,\xi_w)\in\{0,1\}^2
\]
使得
\[
\Delta_w=21\eta_w+34\xi_w,
\]
并且在纤维均匀分布下，
\[
(\eta_w,\xi_w)\in \mathcal A_{d(w)},
\qquad
\mathbf P\bigl((\eta_w,\xi_w)=a\mid w\bigr)=\frac{1}{d(w)}
\quad (a\in \mathcal A_{d(w)}).
\]
特别地：
\begin{enumerate}
  \item \(34\)-寄存器 \(\xi\) 在三种纤维中全部存在，因而构成 universal hidden channel；
  \item \(21\)-寄存器 \(\eta\) 在 \(d=2\) 层被钉死为 \(0\)；
  \item 在 \(d=3\) 层，唯一局域约束是 hard-core 排斥
  \[
  \eta_w\xi_w=0;
  \]
  \item 在 \(d=4\) 层，\((\eta_w,\xi_w)\) 退化为两枚独立公平比特。
\end{enumerate}
\end{theorem}
\begin{proof}
由定理 \ref{thm:xi-time-part9xab-window6-hidden-fiber-affine-boolean-normal-form}，存在
\[
\beta:\{0,1\}^2\to\ZZ,
\qquad
\beta(a,b)=21a+34b,
\]
以及序理想
\[
I_2=\{(0,0),(0,1)\},\qquad
I_3=\{(0,0),(1,0),(0,1)\},\qquad
I_4=\{0,1\}^2
\]
使得
\[
F_w=V_6(w)+\beta(I_{d(w)}).
\]
把 \(\mathcal A_d\) 定义为上述 \(I_d\) 即可。由于 \(\beta\) 在 \(\mathcal A_2,\mathcal A_3,\mathcal A_4\) 上均为单射，故对纤维中的均匀随机元 \(n\) 可唯一写成
\[
n-V_6(w)=\beta(\eta_w,\xi_w)=21\eta_w+34\xi_w
\]
并定义出二值寄存器 \((\eta_w,\xi_w)\)。由纤维均匀分布的推前，\((\eta_w,\xi_w)\) 在 \(\mathcal A_{d(w)}\) 上严格均匀。

四条结论遂立即得到：第 \(1\) 条来自三种 \(\mathcal A_d\) 都允许 \(\xi=1\)；第 \(2\) 条来自 \(\mathcal A_2=\{(0,0),(0,1)\}\)；第 \(3\) 条来自 \(\mathcal A_3\) 中唯一缺失的态正是 \((1,1)\)；第 \(4\) 条则因 \(\mathcal A_4=\{0,1\}^2\) 上的均匀律恰是两枚独立公平比特的乘积分布。证毕。
\end{proof}

\begin{theorem}[条件 hidden 生成多项式的实稳定性]\label{thm:xi-time-part9xaba-window6-hidden-generating-polynomial-real-stable}
设 \(\Lambda\) 为有限站点集，\(w=(w_x)_{x\in\Lambda}\in X_6^\Lambda\) 为固定可见构型，并记
\[
N_d(w):=\#\{x\in\Lambda:\ d(w_x)=d\}
\qquad (d=2,3,4).
\]
定义条件 hidden 配分多项式
\[
\mathcal Z_w(u,v):=
\prod_{x\in\Lambda}
\sum_{(\eta,\xi)\in\mathcal A_{d(w_x)}}u^\eta v^\xi.
\]
则有精确分解
\[
\mathcal Z_w(u,v)
=
(1+v)^{N_2(w)}(1+u+v)^{N_3(w)}\bigl((1+u)(1+v)\bigr)^{N_4(w)}.
\]
更强地，按站点细化的多仿射生成多项式
\[
\mathscr Z_w\bigl((u_x),(v_x)\bigr)
:=
\prod_{x:\,d(w_x)=2}(1+v_x)
\prod_{x:\,d(w_x)=3}(1+u_x+v_x)
\prod_{x:\,d(w_x)=4}(1+u_x)(1+v_x)
\]
是实稳定的。因而其对角专门化 \(\mathcal Z_w(u,v)\) 同样是实稳定多项式；亦即，当
\[
\Im u>0,\qquad \Im v>0
\]
时恒有
\[
\mathcal Z_w(u,v)\neq 0.
\]
\end{theorem}
\begin{proof}
由定理 \ref{thm:xi-time-part9xaba-window6-hidden-two-register-normal-form}，
\[
\sum_{(\eta,\xi)\in\mathcal A_2}u^\eta v^\xi=1+v,
\]
\[
\sum_{(\eta,\xi)\in\mathcal A_3}u^\eta v^\xi=1+u+v,
\]
\[
\sum_{(\eta,\xi)\in\mathcal A_4}u^\eta v^\xi=(1+u)(1+v).
\]
逐站点相乘便得到
\[
\mathcal Z_w(u,v)
=
(1+v)^{N_2(w)}(1+u+v)^{N_3(w)}\bigl((1+u)(1+v)\bigr)^{N_4(w)}.
\]

再证实稳定性。对细化多项式 \(\mathscr Z_w\)，只需检查其基本因子：若 \(\Im v_x>0\)，则 \(1+v_x\neq 0\)；若 \(\Im u_x>0\) 且 \(\Im v_x>0\)，则
\[
\Im(1+u_x+v_x)=\Im u_x+\Im v_x>0,
\]
故 \(1+u_x+v_x\neq 0\)；而 \((1+u_x)(1+v_x)\neq 0\) 显然。于是 \(\mathscr Z_w\) 在所有上半平面变量同时取值时均不为零，即为实稳定多项式。

最后，把全部 \(u_x\) 同时取为 \(u\)，全部 \(v_x\) 同时取为 \(v\)，便回到 \(\mathcal Z_w(u,v)\)。实稳定性在这种对角专门化下保持，因此 \(\mathcal Z_w(u,v)\) 同样实稳定。证毕。
\end{proof}

\noindent
因此，相应归一化后的多仿射条件生成多项式在通常术语下落入 Strongly-Rayleigh 范畴。window-\(6\) 的 hidden sector 因而从一开始就被排除在任何依赖长程正相关凝聚的机制之外。

\begin{proposition}[唯一的局域相互作用源来自 \texorpdfstring{$D_2$}{D2} 切片]\label{prop:xi-time-part9xaba-window6-d2-slice-unique-local-interaction}
\leanverified{paper\_xi\_time\_part9xaba\_window6\_d2\_slice\_unique\_local\_interaction}
设 \(\Lambda\) 为有限站点集，\(w=(w_x)_{x\in\Lambda}\in X_6^\Lambda\) 为固定可见构型，并在定理 \ref{thm:xi-time-part9z-window6-canonical-lift-conditional-uniformization} 的条件分布
\[
\widehat\mu_\Lambda(\,\cdot\mid \pi^\Lambda=w)
\]
下，按定理 \ref{thm:xi-time-part9xaba-window6-hidden-two-register-normal-form} 在每个站点定义寄存器 \((\eta_x,\xi_x)\)。则有：
\begin{enumerate}
  \item 若 \(x\neq y\)，则
  \[
  (\eta_x,\xi_x)\ \text{与}\ (\eta_y,\xi_y)
  \quad\text{在给定 }w\text{ 后条件独立};
  \]
  \item 同一站点上的交叉协方差满足
  \[
  \operatorname{Cov}(\eta_x,\xi_x\mid w)=
  \begin{cases}
  0,& d(w_x)=2,\\[4pt]
  -\dfrac19,& d(w_x)=3,\\[6pt]
  0,& d(w_x)=4;
  \end{cases}
  \]
  \item 因而 \(d=3\) 扇区是 window-\(6\) hidden sector 中唯一产生真正局域相互作用的层；其余两层均为乘积型。
\end{enumerate}
此外，\(d=3\) 扇区在 \(m=6\) 上恰对应四个 cyclic 词
\[
\texttt{100010},\ \texttt{010010},\ \texttt{001010},\ \texttt{101010},
\]
并在 \(B_3\) 字典下恰对应
\[
L_y=\{(\pm 1,0,\pm 1)\}=\{\pm e_1\pm e_3\},
\]
即平面 \(y=0\) 内的 \(D_2=A_1\times A_1\) 长根切片。
\end{proposition}
\begin{proof}
由定理 \ref{thm:xi-time-part9z-window6-canonical-lift-conditional-uniformization}，给定可见构型 \(w\) 以后，hidden 变量在各站点上条件独立，并且在各自纤维上严格均匀。因此由定理 \ref{thm:xi-time-part9xaba-window6-hidden-two-register-normal-form} 的寄存器重编码，条件联合律可写为
\[
\mathbf P\!\left((\eta_x,\xi_x)=a_x\ \forall x\in\Lambda\mid w\right)
=
\prod_{x\in\Lambda}
\frac{\mathbf 1_{\mathcal A_{d(w_x)}}(a_x)}{d(w_x)},
\]
其中 \(a_x\in\{0,1\}^2\)。这立即给出第 \(1\) 条。

再逐层计算同点协方差。若 \(d(w_x)=2\)，则 \(\eta_x\equiv 0\)，故
\[
\operatorname{Cov}(\eta_x,\xi_x\mid w)=0.
\]
若 \(d(w_x)=4\)，则 \((\eta_x,\xi_x)\) 是两枚独立公平比特，故协方差仍为零。若 \(d(w_x)=3\)，则 \((\eta_x,\xi_x)\) 在
\[
\{(0,0),(1,0),(0,1)\}
\]
上均匀，因此
\[
\mathbf E(\eta_x\mid w)=\mathbf E(\xi_x\mid w)=\frac13,
\qquad
\mathbf E(\eta_x\xi_x\mid w)=0,
\]
从而
\[
\operatorname{Cov}(\eta_x,\xi_x\mid w)
=
0-\frac13\cdot\frac13
=
-\frac19.
\]
第 \(2\) 条与第 \(3\) 条遂同时成立。

最后，由定理 \ref{thm:xi-window6-c3-cyclic-fiber-trisection}，
\[
\mathcal R_3=\{\texttt{100010},\texttt{010010},\texttt{001010},\texttt{101010}\},
\]
而定理 \ref{thm:xi-window6-dbin-coordinate-quadruple-stratification} 给出在 \(B_3\) 字典下
\[
\iota_B(\mathcal R_3)=L_y=\{(\pm 1,0,\pm 1)\}.
\]
这正是平面 \(y=0\) 中的四点 \(D_2=A_1\times A_1\) 长根切片。证毕。
\end{proof}

\begin{theorem}[odd 累积量完全支撑于 \texorpdfstring{$D_2$}{D2} 切片]\label{thm:xi-time-part9xaba-window6-odd-cumulants-supported-on-d2-slice}
对单站点 law，记
\[
\widetilde\Delta:=21\eta+34\xi-\mathbf E(21\eta+34\xi\mid d),
\]
其中 \((\eta,\xi)\) 在 \(\mathcal A_d\) 上均匀。相应局部矩母函数记为
\[
M_d(t):=\mathbf E(e^{t\widetilde\Delta}\mid d).
\]
则三层公式精确为
\[
M_2(t)=\cosh(17t),
\]
\[
M_3(t)=\frac{e^{-55t/3}+e^{8t/3}+e^{47t/3}}{3},
\]
\[
M_4(t)=\cosh\!\Bigl(\frac{21t}{2}\Bigr)\cosh(17t).
\]
因此：
\begin{enumerate}
  \item \(d=2\) 与 \(d=4\) 两层的全部奇阶累积量恒为零；
  \item 一切非平凡奇阶累积量全部由 \(d=3\) 层承担；
  \item 特别地，第三阶累积量满足
  \[
  \operatorname{cum}_3^{(2)}=0,\qquad
  \operatorname{cum}_3^{(3)}=-\frac{20680}{27},\qquad
  \operatorname{cum}_3^{(4)}=0.
  \]
\end{enumerate}

更一般地，对任意有限 \(\Lambda\)、可见背景 \(w\in X_6^\Lambda\) 与实权
\[
a=(a_x)_{x\in\Lambda}\in\RR^\Lambda,
\]
定义
\[
H_a:=\sum_{x\in\Lambda}a_x\,\widetilde\Delta_x,
\]
其中 \(\widetilde\Delta_x\) 是第 \(x\) 个站点的中心化 hidden offset。则在给定 \(w\) 的条件下有
\[
\operatorname{cum}_2(H_a\mid w)
=
289\sum_{x:\,d(w_x)=2}a_x^2
+
\frac{1766}{9}\sum_{x:\,d(w_x)=3}a_x^2
+
\frac{1597}{4}\sum_{x:\,d(w_x)=4}a_x^2,
\]
\[
\operatorname{cum}_3(H_a\mid w)
=
-\frac{20680}{27}
\sum_{x:\,d(w_x)=3}a_x^3.
\]
并且更一般地，对每个 \(m\ge 1\) 都有
\[
\operatorname{cum}_{2m+1}(H_a\mid w)
=
c_{2m+1}^{(3)}
\sum_{x:\,d(w_x)=3}a_x^{2m+1},
\]
其中
\[
c_{2m+1}^{(3)}:=\operatorname{cum}_{2m+1}(\widetilde\Delta\mid d=3).
\]
\end{theorem}
\begin{proof}
由定理 \ref{thm:xi-time-part9xaba-window6-hidden-two-register-normal-form} 可直接写出三种局部中心化 law。

当 \(d=2\) 时，\((\eta,\xi)\) 在 \(\{(0,0),(0,1)\}\) 上均匀，因此
\[
21\eta+34\xi\in\{0,34\},
\qquad
\mathbf E(21\eta+34\xi\mid d=2)=17.
\]
故
\[
\widetilde\Delta\in\{-17,17\},
\qquad
M_2(t)=\frac{e^{-17t}+e^{17t}}2=\cosh(17t).
\]
这是偶函数，所以全部奇阶累积量为零。

当 \(d=4\) 时，\((\eta,\xi)\) 是两枚独立公平比特，于是
\[
\widetilde\Delta
=
21\Bigl(\eta-\frac12\Bigr)+34\Bigl(\xi-\frac12\Bigr),
\]
从而
\[
M_4(t)
=
\mathbf E\!\left(e^{21t(\eta-1/2)}\right)
\mathbf E\!\left(e^{34t(\xi-1/2)}\right)
=
\cosh\!\Bigl(\frac{21t}{2}\Bigr)\cosh(17t).
\]
它同样是偶函数，故全部奇阶累积量仍为零。

当 \(d=3\) 时，\((\eta,\xi)\) 在
\[
\{(0,0),(1,0),(0,1)\}
\]
上均匀，因此
\[
21\eta+34\xi\in\{0,21,34\},
\qquad
\mathbf E(21\eta+34\xi\mid d=3)=\frac{55}{3}.
\]
故
\[
\widetilde\Delta\in\left\{-\frac{55}{3},\frac{8}{3},\frac{47}{3}\right\},
\]
并得到
\[
M_3(t)=\frac{e^{-55t/3}+e^{8t/3}+e^{47t/3}}{3}.
\]
由于 \(M_3\) 不是偶函数，故奇阶累积量可能非零；特别地，第三阶累积量等于中心三阶矩：
\[
\operatorname{cum}_3^{(3)}
=
\frac{1}{3}\left[
\left(-\frac{55}{3}\right)^3
+
\left(\frac{8}{3}\right)^3
+
\left(\frac{47}{3}\right)^3
\right]
=
-\frac{20680}{27}.
\]
于是前三条均证。

再看加权和 \(H_a\)。由命题 \ref{prop:xi-time-part9xaba-window6-d2-slice-unique-local-interaction} 第 \(1\) 条，给定可见背景 \(w\) 后，各站点独立，故条件累积母函数相加：
\[
\log \mathbf E(e^{tH_a}\mid w)
=
\sum_{x\in\Lambda}\log M_{d(w_x)}(a_xt).
\]
于是各阶条件累积量可由 \(t=0\) 处逐次求导得到。第二阶只需代入三种局部方差
\[
\operatorname{cum}_2^{(2)}=17^2=289,
\]
\[
\operatorname{cum}_2^{(3)}
=
\frac{1}{3}\left[
\left(-\frac{55}{3}\right)^2
+
\left(\frac{8}{3}\right)^2
+
\left(\frac{47}{3}\right)^2
\right]
=
\frac{1766}{9},
\]
\[
\operatorname{cum}_2^{(4)}
=
\frac{21^2}{4}+17^2
=
\frac{1597}{4},
\]
便得到所示
\(\operatorname{cum}_2(H_a\mid w)\) 公式。

第三阶由局部奇阶支撑性直接给出：
\[
\operatorname{cum}_3(H_a\mid w)
=
\sum_{x\in\Lambda}\operatorname{cum}_3^{(d(w_x))}a_x^3
=
-\frac{20680}{27}\sum_{x:\,d(w_x)=3}a_x^3.
\]
同理，对任意 \(m\ge 1\)，因为 \(d=2,4\) 两层的全部奇阶累积量均为零，故
\[
\operatorname{cum}_{2m+1}(H_a\mid w)
=
\sum_{x:\,d(w_x)=3}\operatorname{cum}_{2m+1}^{(3)}a_x^{2m+1}
=
c_{2m+1}^{(3)}
\sum_{x:\,d(w_x)=3}a_x^{2m+1}.
\]
证毕。
\end{proof}

\begin{proposition}[Gaussian 相容的显式必要条件]\label{prop:xi-time-part9xaba-window6-gaussian-compatibility-necessary-condition}
\leanverified{paper\_xi\_time\_part9xaba\_window6\_gaussian\_compatibility\_necessary\_condition}
设 \((\Lambda_\delta,w^{(\delta)})\) 为一列有限体积可见背景，\(a^{(\delta)}=(a_x^{(\delta)})_{x\in\Lambda_\delta}\) 为实权，并定义
\[
Y_\delta
:=
\frac{1}{b_\delta}
\sum_{x\in\Lambda_\delta}
a_x^{(\delta)}\,\widetilde\Delta_x.
\]
若在条件可见背景 \(w^{(\delta)}\) 下，
\[
Y_\delta\Longrightarrow \mathcal N(0,\sigma^2)
\]
收敛到某个中心高斯律，则必有
\[
\frac{1}{b_\delta^3}
\sum_{x:\,d(w_x^{(\delta)})=3}
\bigl(a_x^{(\delta)}\bigr)^3
\longrightarrow 0.
\]
更一般地，对每个 \(m\ge 1\) 都有精确公式
\[
\operatorname{cum}_{2m+1}(Y_\delta\mid w^{(\delta)})
=
\frac{c_{2m+1}^{(3)}}{b_\delta^{2m+1}}
\sum_{x:\,d(w_x^{(\delta)})=3}
\bigl(a_x^{(\delta)}\bigr)^{2m+1},
\]
其中 \(c_{2m+1}^{(3)}\) 为定理 \ref{thm:xi-time-part9xaba-window6-odd-cumulants-supported-on-d2-slice} 中 \(d=3\) 单点 law 的对应 odd 累积量。
\end{proposition}
\begin{proof}
对中心高斯极限而言，第三阶及以上的奇阶累积量必须在极限中消失。由定理 \ref{thm:xi-time-part9xaba-window6-odd-cumulants-supported-on-d2-slice}，
\[
\operatorname{cum}_3(Y_\delta\mid w^{(\delta)})
=
\frac{1}{b_\delta^3}
\operatorname{cum}_3\!\left(
\sum_{x\in\Lambda_\delta}
a_x^{(\delta)}\widetilde\Delta_x
\Bigm|\,
w^{(\delta)}
\right)
=
-\frac{20680}{27}\,
\frac{1}{b_\delta^3}
\sum_{x:\,d(w_x^{(\delta)})=3}
\bigl(a_x^{(\delta)}\bigr)^3.
\]
因此要与中心高斯律相容，右端必须趋于 \(0\)，这便得到所示必要条件。

更一般地，同一定理已经给出全部奇阶条件累积量的支撑公式：
\[
\operatorname{cum}_{2m+1}\!\left(
\sum_{x\in\Lambda_\delta}
a_x^{(\delta)}\widetilde\Delta_x
\Bigm|\,
w^{(\delta)}
\right)
=
c_{2m+1}^{(3)}
\sum_{x:\,d(w_x^{(\delta)})=3}
\bigl(a_x^{(\delta)}\bigr)^{2m+1}.
\]
再除以 \(b_\delta^{2m+1}\) 即得
\[
\operatorname{cum}_{2m+1}(Y_\delta\mid w^{(\delta)})
=
\frac{c_{2m+1}^{(3)}}{b_\delta^{2m+1}}
\sum_{x:\,d(w_x^{(\delta)})=3}
\bigl(a_x^{(\delta)}\bigr)^{2m+1}.
\]
证毕。
\end{proof}

\begin{conjecture}[六顶点/GFF 的 \texorpdfstring{$D_2$}{D2}-slice washout 原理]\label{conj:xi-time-part9xaba-window6-sixvertex-gff-d2slice-washout}
设某一临界六顶点系综中，可见高度场在适当归一化下收敛到高斯自由场。则对任意 canonical window-\(6\) lift，其 hidden sector 的连续极限应满足：
\begin{enumerate}
  \item universal \(34\)-channel 只能贡献局域的、解析可重整的一体修正；
  \item 自由 \(21\)-channel 在解冻的乘积层上只能贡献对称的接触型噪声；
  \item 唯一非乘积且带 odd 累积量的 hard-core \(D_2\)-slice 必全部 wash out 为纯接触项；
  \item 因而 window-\(6\) 不会生成第二个独立的 massless boson，其全部非平凡连续信息仍由 visible height/GFF 承担。
\end{enumerate}
\end{conjecture}
\begin{proof}[证据]
定理 \ref{thm:xi-time-part9xaba-window6-hidden-generating-polynomial-real-stable} 表明，按站点细化后的条件 hidden law 已落入实稳定/Strongly-Rayleigh 范畴，因此不存在可支持长程正相关凝聚的机制。命题 \ref{prop:xi-time-part9xaba-window6-d2-slice-unique-local-interaction} 又把唯一的局域相互作用源严格定位在
\[
L_y=\{\pm e_1\pm e_3\},
\]
即单个 \(D_2\) 切片上；其余 \(d=2,4\) 两层均为乘积型。再由定理 \ref{thm:xi-time-part9xaba-window6-odd-cumulants-supported-on-d2-slice}，全部 odd 累积量都只由该切片承担；而命题 \ref{prop:xi-time-part9xaba-window6-gaussian-compatibility-necessary-condition} 则进一步说明，任何与 Gaussian 极限相容的标度都必须把这一 \(D_2\)-slice 的偏斜源压到极限中消失。

于是，一旦该 skew source 在目标标度上被洗出，剩余 hidden 部门就全部退回到实稳定的乘积/半乘积范畴，只能留下局域一体项与接触型噪声，而不可能孕育第二个独立的无质量连续场。这正是题述 washout 原理的证据。证毕。
\end{proof}

\noindent
由此，window-\(6\) 的 hidden sector 被进一步压缩为一个两寄存器、三势垒的局域气体：\(34\)-channel 是 universal hidden 通道，\(21\)-channel 按 \(d=2,3,4\) 三层逐步解冻，而唯一非乘积、唯一反相关、唯一携带 odd 累积量的 defect 则全部集中在 \(B_3\) 字典的单个 \(D_2\)-slice 上。于是，window-\(6\) 与六顶点/GFF 接口中最锋利的局域障碍，不再只是抽象的“hidden 只有短程”，而被精确定位为这四个 \(d=3\) 词所承载的 hard-core skew source。

\endinput

\paragraph{\texorpdfstring{$2$}{2} 进全息地址、二维刚性谱支与 Bernoulli--\texorpdfstring{$\zeta$}{zeta} 反演的 genus-zero 约束}
在事件词的自由仿射单子、全息像集的严格 cylinder 分解、fold-gauge 常数的 Stirling--Bernoulli 层级、真实输入 \(40\) 状态核的 \((2+8)\) 谱分裂，以及碰撞势的代数参数化已经建立之后，还可把这三条此前分散的主线进一步压缩为下列七条后果。它们分别把有限前缀的地址化作用提升为忠实仿射表示，把刚性谱支提升为奇偶矩分裂与压力平坦化机制，并把规范常数、局部 Cram\'er germ 与碰撞大偏差曲线统一到一条可递归反演且受 genus-zero 几何约束的代数骨架之中。

\begin{theorem}[有限前缀单词在 \texorpdfstring{$2$}{2} 进全息地址上的忠实仿射实现]\label{thm:xi-time-part9xb-prefix-words-faithful-affine-realization}
\leanverified{paper\_xi\_time\_part9xb\_prefix\_words\_faithful\_affine\_realization}
沿用定理 \ref{thm:killo-godel-semidir-affine-2adic} 的记号。设 \(E\) 为有限原子事件集，
\[
B=2^L>\max_{e\in E}\mathrm{code}(e),
\qquad
T:=T_2(E_{\min}),
\qquad
v\in T\setminus 2T.
\]
对任意有限单词
\[
w=(e_1,\dots,e_n)\in E^\ast
\]
定义
\[
r_w(p_t):=
\begin{cases}
\mathrm{code}(e_t),&1\le t\le n,\\
0,&t>n,
\end{cases}
\qquad
X_n(w):=\sum_{t=1}^{n}\mathrm{code}(e_t)B^{t-1}v.
\]
则成立：
\begin{enumerate}
  \item 对任意尾流 \(f\in E^{\NN}\)，有精确前缀公式
  \[
  X(wf)=X_n(w)+B^nX(f)=\Psi(r_w,n)\bigl(X(f)\bigr);
  \]
  \item 赋值
  \[
  \Theta:E^\ast\longrightarrow \mathrm{Aff}(T),
  \qquad
  \Theta(w):=\Psi(r_w,|w|)
  \]
  是自由单子 \(E^\ast\) 到 \(\mathrm{Aff}(T)\) 的单射同态。
\end{enumerate}
\end{theorem}
\begin{proof}
由 \(\operatorname{ev}_B(r_w)=X_n(w)\)，对任意 \(x\in T\) 都有
\[
\Psi(r_w,n)(x)=B^n x+X_n(w).
\]
因此第 \(1\) 条正是定理
\ref{thm:xi-time-part62cc-eventword-free-affine-submonoid} 中
\[
X(wf)=F_w\bigl(X(f)\bigr)
\]
的记号重写。

再令
\[
F_w(x):=X_n(w)+B^{|w|}x.
\]
同一定理已经证明
\[
F_u\circ F_v=F_{uv}
\qquad\text{且}\qquad
w\longmapsto F_w
\]
是从 \(E^\ast\) 到 \(\mathrm{Aff}(T)\) 的单射。由于
\[
F_w=\Psi(r_w,|w|)=\Theta(w),
\]
故 \(\Theta\) 亦为单射同态。证毕。
\end{proof}

\begin{corollary}[全息像集的严格自仿射分解与精确分支数]\label{cor:xi-time-part9xb-holographic-image-strict-self-affine-decomposition}
沿用定理 \ref{thm:xi-time-part9xb-prefix-words-faithful-affine-realization} 的记号，并记
\[
\mathcal X:=X(E^{\NN})\subset T.
\]
则对任意 \(n\ge 1\) 都有严格分解
\[
\mathcal X=\bigsqcup_{w\in E^n}\bigl(X_n(w)+B^n\mathcal X\bigr).
\]
特别地，\(n\) 级全息 cylinder 恰有
\[
|E|^n
\]
个，且彼此不交。
\end{corollary}
\begin{proof}
任取 \(x=X(\mathbf e)\in \mathcal X\)，把 \(\mathbf e\) 写成长度 \(n\) 前缀 \(w\) 与尾流 \(f\) 的连接 \(\mathbf e=wf\)。由定理
\ref{thm:xi-time-part9xb-prefix-words-faithful-affine-realization}，
\[
x=X(wf)=X_n(w)+B^nX(f)\in X_n(w)+B^n\mathcal X,
\]
故得覆盖。

反向包含同样由
\[
X_n(w)+B^nX(f)=X(wf)\in \mathcal X
\]
立即成立。若
\[
\bigl(X_n(w)+B^n\mathcal X\bigr)\cap
\bigl(X_n(w')+B^n\mathcal X\bigr)\neq\varnothing,
\]
则存在 \(f,f'\in E^{\NN}\) 使 \(X(wf)=X(w'f')\)。由定理
\ref{thm:xi-time-part62d-hologram-image-cantor-homeomorphism} 的单射性，遂有
\[
wf=w'f',
\]
从而 \(w=w'\)。故上述并为严格不交并。块数显然等于长度 \(n\) 单词数 \(|E|^n\)。证毕。
\end{proof}

\begin{theorem}[由折叠规范常数序列递归反演全部 \texorpdfstring{$\zeta(2r)$}{zeta(2r)} 的显式公式]\label{thm:xi-time-part9xb-gauge-constants-recover-even-zeta}
沿用定理 \ref{thm:fold-bin-gauge-constant-stirling-bernoulli-hierarchy} 的规范常数列 \((C_m)\)。对每个 \(r\ge 1\)，记
\[
a_r:=
\frac{B_{2r}}{r(2r-1)}
\bigl(\varphi^{-1}+\varphi^{2r-3}\bigr),
\]
并定义
\[
L_r:=
\lim_{m\to\infty}
\Bigl(\frac{2}{\varphi}\Bigr)^{(2r-1)m}
\left[
\log C_m-\log(2\pi)-\sum_{s=1}^{r-1}a_s
\Bigl(\frac{\varphi}{2}\Bigr)^{(2s-1)m}
\right].
\]
则极限存在且满足
\[
L_r=a_r.
\]
从而对每个 \(r\ge 1\)，有显式反演公式
\[
\zeta(2r)
=
(-1)^{r-1}\frac{(2\pi)^{2r}}{2(2r)!}\,
\frac{r(2r-1)}{\varphi^{-1}+\varphi^{2r-3}}\,
L_r.
\]
\end{theorem}
\begin{proof}
这正是定理 \ref{thm:xi-time-part60aa-gauge-constant-explicit-even-zeta-recovery} 的重写，其中该定理的 \(A_r\) 与此处的 \(a_r\) 完全一致。证毕。
\end{proof}

\begin{theorem}[二维刚性谱支的奇偶矩完全分裂]\label{thm:xi-time-part9xb-rigid-branches-odd-even-moment-splitting}
沿用定理 \ref{thm:real-input-40-det-uv-2plus8} 与推论 \ref{cor:real-input-40-zeta-uv-sqrtv-eigs} 的记号。记
\[
\Delta(z;u,v)=\det(I-zM_{u,v})=(1-vz^2)\mathcal Q_8(z;u,v),
\]
并把
\[
\mathcal Q_8(z;u,v):=-P_8(z;u,v)=\prod_{j=1}^{8}\bigl(1-\Lambda_j(u,v)z\bigr)
\]
视为按代数重数计的八条 core 谱支分解；若次数下降，则允许若干 \(\Lambda_j(u,v)=0\)。则对每个 \(n\ge 1\) 都有精确谱矩公式
\[
\operatorname{Tr}\bigl(M_{u,v}^n\bigr)
=
(\sqrt v)^n+(-\sqrt v)^n+\sum_{j=1}^{8}\Lambda_j(u,v)^n.
\]
特别地，对任意 \(m\ge 0\)，有
\[
\operatorname{Tr}\bigl(M_{u,v}^{2m+1}\bigr)
=
\sum_{j=1}^{8}\Lambda_j(u,v)^{2m+1},
\]
而
\[
\operatorname{Tr}\bigl(M_{u,v}^{2m}\bigr)
=
2v^m+\sum_{j=1}^{8}\Lambda_j(u,v)^{2m}.
\]
\end{theorem}
\begin{proof}
由定理 \ref{thm:real-input-40-det-uv-2plus8}，矩阵 \(M_{u,v}\) 的特征多项式分解为
\[
\det(I-zM_{u,v})=(1-vz^2)\mathcal Q_8(z;u,v).
\]
再由推论 \ref{cor:real-input-40-zeta-uv-sqrtv-eigs}，因子 \(1-vz^2\) 对应两条刚性谱支
\[
\lambda_\pm=\pm\sqrt v.
\]
其余特征值恰由 \(\mathcal Q_8\) 的八个根给出，即 \(\Lambda_1,\dots,\Lambda_8\)（按代数重数计，并允许零根）。因此 \(M_{u,v}\) 的全部特征值正是
\[
\sqrt v,\ -\sqrt v,\ \Lambda_1(u,v),\dots,\Lambda_8(u,v).
\]
对有限矩阵的幂迹展开便得到
\[
\operatorname{Tr}\bigl(M_{u,v}^n\bigr)
=
(\sqrt v)^n+(-\sqrt v)^n+\sum_{j=1}^{8}\Lambda_j(u,v)^n.
\]
当 \(n=2m+1\) 为奇数时，前两项相消；当 \(n=2m\) 为偶数时，前两项相加为 \(2v^m\)。证毕。
\end{proof}

\begin{corollary}[刚性谱支主导区的联合累积量全湮灭]\label{cor:xi-time-part9xb-rigid-branch-dominant-cumulant-vanishing}
沿用定理 \ref{thm:xi-time-part9xb-rigid-branches-odd-even-moment-splitting} 的记号。设 \(U\subset \RR_{>0}^2\) 为开集，并假定在 \(U\) 上处处成立
\[
\max_{1\le j\le 8}|\Lambda_j(u,v)|<\sqrt v.
\]
令
\[
\widetilde U:=\{(\theta,\eta)\in\RR^2:\ (e^\theta,e^\eta)\in U\},
\qquad
P(\theta,\eta):=\log \rho(M_{e^\theta,e^\eta}).
\]
则在 \(\widetilde U\) 上恒有
\[
P(\theta,\eta)=\frac{\eta}{2}.
\]
从而
\[
\partial_\theta P\equiv 0,
\qquad
\partial_\eta P\equiv \frac12,
\]
并且对任意满足 \(a+b\ge 2\) 的多重指标 \((a,b)\)，都有
\[
\partial_\theta^a\partial_\eta^b P\equiv 0.
\]
\end{corollary}
\begin{proof}
对任意 \((u,v)\in U\)，由定理
\ref{thm:xi-time-part9xb-rigid-branches-odd-even-moment-splitting} 可知 \(M_{u,v}\) 的全部特征值由两条刚性谱支 \(\pm\sqrt v\) 与八条 core 谱支 \(\Lambda_j(u,v)\) 组成，因此
\[
\rho(M_{u,v})
=
\max\left\{\sqrt v,\ \max_{1\le j\le 8}|\Lambda_j(u,v)|\right\}.
\]
所给假设保证主导谱支恒为 \(\sqrt v\)，故在 \(\widetilde U\) 上
\[
P(\theta,\eta)=\log\sqrt{e^\eta}=\frac{\eta}{2}.
\]
所述偏导公式随即成立。证毕。
\end{proof}

\begin{proposition}[真实输入输出位势的局部 Cram\'er 级数闭包于 \texorpdfstring{$\QQ(\sqrt5)$}{Q(sqrt5)}]\label{prop:xi-time-part9xb-output-cramer-series-quadratic-field}
沿用命题 \ref{prop:real-input-40-output-cumulants-closed} 的记号，记
\[
\kappa_j:=P^{(j)}(0)\qquad (j\ge 1),
\qquad
\alpha_0:=\kappa_1.
\]
设 \(I(\alpha)\) 为 \(P\) 的 Legendre 对偶在 \(\alpha_0\) 邻域内的解析支。则存在唯一级数
\[
I(\alpha_0+\delta)=\sum_{n\ge 2}c_n\delta^n
\]
并且对全部 \(n\ge 2\) 都有
\[
c_n\in \QQ(\sqrt5).
\]
特别地，
\[
c_2=\frac{1}{2\kappa_2},
\qquad
c_3=-\frac{\kappa_3}{6\kappa_2^3},
\qquad
c_4=\frac{3\kappa_3^2-\kappa_2\kappa_4}{24\kappa_2^5}.
\]
\end{proposition}
\begin{proof}
定理 \ref{thm:xi-output-potential-local-legendre-quadratic-field} 已证明 \(I\) 在 \(\alpha_0\) 邻域内解析，且其全部 Taylor 系数都属于 \(\QQ(\sqrt5)\)。再把定理 \ref{thm:xi-real-input-cramer-germ-quadratic-field-rigidity} 中的展开
\[
I(\alpha_0+\delta)
=
\frac{\delta^2}{2\kappa_2}
-\frac{\kappa_3}{6\kappa_2^3}\delta^3
+\frac{3\kappa_3^2-\kappa_2\kappa_4}{24\kappa_2^5}\delta^4
+O(\delta^5)
\]
与
\[
I(\alpha_0+\delta)=\sum_{n\ge 2}c_n\delta^n
\]
逐项比较，即得所示前三个系数公式；它们显然仍落在 \(\QQ(\sqrt5)\) 中。证毕。
\end{proof}

\begin{theorem}[碰撞大偏差约束曲线的有理性]\label{thm:xi-time-part9xb-collision-ldp-constraint-curve-rationality}
定义有理函数
\[
u(\rho)=\frac{\rho(\rho^2-2\rho-1)}{\rho-1},
\qquad
\alpha(\rho)=\frac{(\rho^2-2\rho-1)(\rho-1)}{2\rho^3-5\rho^2+4\rho+1}.
\]
设 \(\mathcal C\subset \mathbb A^2_{u,\alpha}\) 为有理映射
\[
\PP^1\dashrightarrow \mathbb A^2_{u,\alpha},
\qquad
\rho\longmapsto \bigl(u(\rho),\alpha(\rho)\bigr)
\]
之像的 Zariski 闭包。则 \(\mathcal C\) 的规范化是一条有理曲线；等价地，\(\mathcal C\) 的几何 genus 为 \(0\)。此外，在物理支
\[
\rho>1+\sqrt2
\]
上有速率函数表示
\[
I_2\bigl(\alpha(\rho)\bigr)
=
\alpha(\rho)\log u(\rho)-\log \rho+\log(\varphi^2).
\]
因此，碰撞大偏差的约束部分由 genus-zero 的有理曲线完全承载，而速率函数只是在该曲线上附加一个初等对数势。
\end{theorem}
\begin{proof}
定理 \ref{thm:xi-time-part9p-collision-frequency-algebraic-ldp} 已给出 \((u,\alpha)\) 的显式有理参数化
\[
\rho\longmapsto \bigl(u(\rho),\alpha(\rho)\bigr),
\]
故其像的 Zariski 闭包正是某条不可约代数曲线 \(\mathcal C\)。记 \(\overline{\mathcal C}\) 为 \(\mathcal C\) 的射影闭包，\(\widetilde{\mathcal C}\) 为其规范化。上述参数化诱导出一个从 \(\PP^1\) 到 \(\overline{\mathcal C}\) 的非常值有理映射，因此也给出一个非常值有理映射
\[
\nu:\PP^1\dashrightarrow \widetilde{\mathcal C}.
\]
由于 \(\PP^1\) 完备，\(\nu\) 实际延拓为非常值态射。对该态射应用 Riemann--Hurwitz 公式，得
\[
-2=\deg(\nu)\bigl(2g(\widetilde{\mathcal C})-2\bigr)+R,
\qquad
R\ge 0.
\]
因此只能有
\[
g(\widetilde{\mathcal C})=0,
\]
故 \(\widetilde{\mathcal C}\) 为有理曲线。

最后，同一物理支上的速率函数表示
\[
I_2\bigl(\alpha(\rho)\bigr)
=
\alpha(\rho)\log u(\rho)-\log \rho+\log(\varphi^2)
\]
正是定理 \ref{thm:xi-time-part9p-collision-frequency-algebraic-ldp} 的结论。于是约束部分完全由上述有理参数曲线控制，而作用部分只额外携带 \(\log u(\rho)\) 与 \(\log\rho\) 两个初等对数项。证毕。
\end{proof}

\noindent
由此，有限前缀的自由单子、二维核的刚性谱支与碰撞/规范常数的代数反演链不再只是分散的局部现象：前者在 \(2\)-进地址空间中形成严格不交的自仿射 cylinder 塔，中间谱层把奇偶矩与压力曲率强制分裂，而后者则把 \((C_m)\)、局部 Cram\'er germ 与碰撞约束曲线同时压缩到可递归恢复并受 genus-zero 几何约束的统一骨架。

\endinput

\paragraph{Serrin 可实现支上的全息因子化、尺度残余与有限通量坐标}
在固定分辨率 Stokes--Gödel 全息强完备性（定理 \ref{thm:xi-time-part9x-fixed-resolution-stokes-godel-strong-completeness}）、精确屏幕的 sharp 极小规模（定理 \ref{thm:xi-time-part9fa-exact-screen-minimal-size-matroid-basis}）、各向异性 Serrin 射线的 Wulff 压缩（定理 \ref{thm:xi-time-part9x-serrin-ray-fibonacci-semirings}）、以及椭球与对角 \(\ell_p\)-超椭球的有限边界通量反演（定理 \ref{thm:spg-ellipsoid-boundary-flux-reconstruction} 与定理 \ref{thm:spg-lp-superellipsoid-boundary-volume-decoding}）都已建立之后，还可把这些接口进一步压缩为一条纯粹的因子化链。其要点在于：一般 dyadic 类中的高维 exactification 自由度，在 Serrin 可实现支上不再表现为拟阵基交换，而只剩平移与尺度两类几何参数；模去平移后，一切 fixed-resolution 边界全息坐标都只能读取单一尺度残余；而在已有有限通量反演的显式 Wulff 家族上，这一残余又可被有限阶边界通量直接坐标化。

\begin{theorem}[Serrin 可实现支上的商全息坐标必经尺度因子]\label{thm:xi-time-part9xbb-serrin-quotient-holography-factors-through-scale}
固定 \(n\ge 2\)、分辨率 \(m\ge 1\) 与某个 Wulff 形 \(W\subset\RR^n\)。记
\[
\mathfrak S_W:=\{aW+t:\ a>0,\ t\in\RR^n\},
\qquad
\mathfrak S_{W,m,n}^{\mathrm{tr}}
:=
(\mathfrak S_W\cap \mathcal U_{m,n})/\RR^n.
\]
定义尺度映射
\[
\sigma_W:\mathfrak S_{W,m,n}^{\mathrm{tr}}\longrightarrow \RR_{>0},
\qquad
\sigma_W([aW+t]):=a.
\]
则 \(\sigma_W\) 良定义且为单射。因而对任意集合 \(Z_m\) 与任意映射
\[
\overline{\mathfrak I}_{W,m}:\mathfrak S_{W,m,n}^{\mathrm{tr}}\to Z_m,
\]
都存在唯一映射
\[
\Theta_{W,m}:\sigma_W(\mathfrak S_{W,m,n}^{\mathrm{tr}})\to Z_m
\]
使得
\[
\overline{\mathfrak I}_{W,m}
=
\Theta_{W,m}\circ \sigma_W.
\]
特别地，一切在平移商上定义的 fixed-resolution 边界全息坐标都至多保留一个 genuinely visible 的尺度模参数。
\end{theorem}
\begin{proof}
由 Figalli--Zhang 的各向异性 Serrin 刚性 \cite{FigalliZhang2026AnisotropicSerrinGeneralDomains}，\(\mathfrak S_W\) 的每个元素都可写成 \(aW+t\)。若
\[
[aW+t]=[a'W+t']
\]
于平移商中相等，则 \(aW\) 与 \(a'W\) 互为平移。对两边取体积即得
\[
a^n|W|=a'^n|W|,
\]
故 \(a=a'\)。这说明 \(\sigma_W\) 良定义且单射。

既然 \(\sigma_W\) 把 \(\mathfrak S_{W,m,n}^{\mathrm{tr}}\) 识别为其像中的一个子集，则对任意 \(\overline{\mathfrak I}_{W,m}\) 只需定义
\[
\Theta_{W,m}(a):=\overline{\mathfrak I}_{W,m}([aW])
\qquad
\bigl(a\in \sigma_W(\mathfrak S_{W,m,n}^{\mathrm{tr}})\bigr),
\]
便有
\[
\overline{\mathfrak I}_{W,m}([aW+t])
=
\Theta_{W,m}(a)
=
(\Theta_{W,m}\circ\sigma_W)([aW+t]).
\]
唯一性由 \(\sigma_W\) 的单射性立即得到。证毕。
\end{proof}

\begin{corollary}[从 \texorpdfstring{$2^{mn}$}{2^(mn)} 级 exactification 到单尺度残余的 sharp 塌缩]\label{cor:xi-time-part9xbb-sharp-exactification-collapse-to-one-scale}
在一般分辨率 \(m\) 的 dyadic 体链 \(\mathcal U_{m,n}\) 上，极小精确屏幕的最小规模恰为
\[
2^{mn}
\]
（定理 \ref{thm:xi-time-part9fa-exact-screen-minimal-size-matroid-basis}）。另一方面，在 Serrin 可实现支的平移商 \(\mathfrak S_{W,m,n}^{\mathrm{tr}}\) 上，任意商边界全息坐标都必经由单一尺度参数 \(\sigma_W\)。因此，一般 dyadic 类中的 \(2^{mn}\) 级 exactification 在 Serrin 支上 sharp 地塌缩为至多一维的尺度残余：
\[
2^{mn}\ \leadsto\ 1.
\]
若再施加体积归一化
\[
|U|=|W|,
\]
则 \(a=1\)，从而进一步塌缩为
\[
2^{mn}\ \leadsto\ 0.
\]
\end{corollary}
\begin{proof}
第一条结论正是定理 \ref{thm:xi-time-part9fa-exact-screen-minimal-size-matroid-basis}。第二条结论由定理 \ref{thm:xi-time-part9xbb-serrin-quotient-holography-factors-through-scale} 立即得到。若再加体积归一化，则对任意 \(U=aW+t\in\mathfrak S_W\) 有
\[
|U|=a^n|W|=|W|,
\]
故 \(a=1\)。于是平移商上的唯一剩余尺度亦被固定，因而可见模空间退化为单点。证毕。
\end{proof}

\begin{theorem}[显式 Wulff 家族上的有限通量坐标化]\label{thm:xi-time-part9xbb-explicit-wulff-finite-flux-coordinate-factorization}
\leanverified{paper\_xi\_time\_part9xbb\_explicit\_wulff\_finite\_flux\_coordinate\_factorization}
设 \(\mathfrak W_{\mathrm{exp}}\) 为文中任一显式 Wulff 家族，并设已有一个有限边界通量坐标
\[
\mathcal F=(\mathcal F_1,\dots,\mathcal F_k):\mathfrak W_{\mathrm{exp}}\to\RR^k
\]
在该家族上为单射。则对任意集合 \(Z_m\) 与任意映射
\[
\mathfrak I_m:\mathfrak W_{\mathrm{exp}}\to Z_m,
\]
都存在唯一映射
\[
\Phi_m:\mathcal F(\mathfrak W_{\mathrm{exp}})\to Z_m
\]
使得
\[
\mathfrak I_m=\Phi_m\circ \mathcal F.
\]
特别地，只要
\[
\mathcal F(U)=\mathcal F(V)
\qquad
(U,V\in\mathfrak W_{\mathrm{exp}}),
\]
便必有
\[
\mathfrak I_m(U)=\mathfrak I_m(V).
\]

椭球家族与对角 \(\ell_p\)-超椭球家族分别由定理 \ref{thm:spg-ellipsoid-boundary-flux-reconstruction} 与定理 \ref{thm:spg-lp-superellipsoid-boundary-volume-decoding} 提供这样的有限通量坐标。
\end{theorem}
\begin{proof}
由假设，\(\mathcal F\) 在 \(\mathfrak W_{\mathrm{exp}}\) 上为单射，故其像 \(\mathcal F(\mathfrak W_{\mathrm{exp}})\) 上可定义逆映射
\[
\mathcal F^{-1}:\mathcal F(\mathfrak W_{\mathrm{exp}})\to\mathfrak W_{\mathrm{exp}}.
\]
于是只需令
\[
\Phi_m(y):=\mathfrak I_m\bigl(\mathcal F^{-1}(y)\bigr)
\qquad
\bigl(y\in \mathcal F(\mathfrak W_{\mathrm{exp}})\bigr),
\]
便有
\[
(\Phi_m\circ\mathcal F)(U)
=
\Phi_m(\mathcal F(U))
=
\mathfrak I_m(U)
\qquad
(U\in\mathfrak W_{\mathrm{exp}}).
\]
唯一性同样由 \(\mathcal F\) 的单射性得到。证毕。
\end{proof}

\begin{corollary}[归一化 Serrin 类上的 fixed-resolution 全息平凡化]\label{cor:xi-time-part9xbb-normalized-serrin-holography-trivialization}
固定 Wulff 形 \(W\subset\RR^n\)，并设
\[
\mathfrak S_W^{\mathrm{nor}}
:=
\{\,U\in\mathfrak S_W:\ x_c(U)=0,\ |U|=|W|\,\},
\]
其中 \(x_c(U)\) 表示重心。则 \(\mathfrak S_W^{\mathrm{nor}}\) 在平移与尺度意义下只含一个对象。因而，对任意集合 \(Z_m\) 与任意映射
\[
\mathfrak J_m:\mathfrak S_W^{\mathrm{nor}}\to Z_m,
\]
都有 \(\mathfrak J_m\) 为常值。

更一般地，任何足以同时固定平移与尺度的归一化条件都会强制同样的常值化；若 \(W\) 属于定理 \ref{thm:xi-time-part9xbb-explicit-wulff-finite-flux-coordinate-factorization} 的显式家族，则也可用一个正一阶齐次、能够固定尺度的边界通量坐标替代体积归一化。于是 exact screen、有限矩盒、边界 Gödel 坐标以及整数 moment-kernel / 射影转移读出在归一化 Serrin 类上全部退化为单值。
\end{corollary}
\begin{proof}
任取 \(U\in\mathfrak S_W^{\mathrm{nor}}\)。由 \(U\in \mathfrak S_W\)，可写
\[
U=aW+t.
\]
由重心归一化 \(x_c(U)=0\) 可解出平移 \(t\)；再由体积归一化
\[
|U|=a^n|W|=|W|
\]
知 \(a=1\)。故 \(U\) 被唯一强制为 \(W\) 的归一化代表，从而 \(\mathfrak S_W^{\mathrm{nor}}\) 仅含一个对象。定义在单点集上的任意映射必为常值。

若改用任意其他足以同时固定平移与尺度的归一化条件，论证完全相同；而在显式通量家族中，尺度若可由一个正一阶齐次边界通量坐标固定，则定理 \ref{thm:xi-time-part9xbb-explicit-wulff-finite-flux-coordinate-factorization} 表明全部读出都只剩该已被冻结的通量参数，故仍为常值。证毕。
\end{proof}

\begin{conjecture}[Serrin 支上的 screen--RT--moment--Gödel 因子化原则]\label{conj:xi-time-part9xbb-serrin-screen-rt-moment-godel-factorization}
设 \(\mathfrak G_{W,m}\) 为任意 fixed-resolution 边界函数量，并假定它与下列接口同时兼容：精确屏幕极小基、平方自由 RT 影子、有限矩完备恢复、边界 Gödel 全息，以及整数 moment-kernel / 射影转移算子。则在 Serrin 可实现支上，\(\mathfrak G_{W,m}\) 应当满足如下因子化链：
\[
\mathfrak G_{W,m}
\quad\text{经由}\quad
(a,t)\in \RR_{>0}\times\RR^n
\quad\text{因子化};
\]
模去平移后，仅经由单一尺度参数 \(a\) 因子化；而在任一已知承认有限通量反演的显式 Wulff 家族上，又进一步经由有限维边界通量坐标 \(\mathcal F\) 因子化。

等价地，exact screen、平方自由 RT cut、finite moments、boundary Gödel holography 与 moment-kernel 动力学在 Serrin 支上不应再被视为并列的五种独立机制，而应被视为同一低维几何轨道的五种读出。
\end{conjecture}
\begin{proof}[证据]
固定分辨率上的强完备边界全息由定理 \ref{thm:xi-time-part9x-fixed-resolution-stokes-godel-strong-completeness}、定理 \ref{thm:spg-dyadic-finite-moment-completeness} 与推论 \ref{cor:spg-boundary-godel-finite-moment-completeness} 已经建立；一般 dyadic 类上的 exact-screen 复杂度由定理 \ref{thm:xi-time-part9fa-exact-screen-minimal-size-matroid-basis} 锁定为 \(2^{mn}\)。另一方面，定理 \ref{thm:xi-time-part9xbb-serrin-quotient-holography-factors-through-scale} 与推论 \ref{cor:xi-time-part9xbb-sharp-exactification-collapse-to-one-scale} 已表明：一旦限制到 Serrin 可实现支并模去平移，高维自由度便塌缩为单尺度残余；而定理 \ref{thm:xi-time-part9xbb-explicit-wulff-finite-flux-coordinate-factorization} 又说明，在椭球与对角 \(\ell_p\)-超椭球等显式 Wulff 家族中，这一残余进一步压缩到有限阶边界通量坐标。现有全部链内证据都与该因子化原则同向。证毕。
\end{proof}

\[
2^{mn}\text{ 级 exactification}
\Longrightarrow
\text{Serrin 支上的单尺度残余}
\Longrightarrow
\text{显式 Wulff 家族上的有限通量坐标化}
\Longrightarrow
\text{归一化后的全息平凡化}.
\]

\begin{theorem}[Serrin 尺度纤维上的随机包冻结与共振寄存器下界]\label{distill:bourgain-random_to_deterministic_certificate_extraction-serrin-scale-fibre-resonance-register}
固定 \(n\ge2\)、分辨率 \(m\ge1\) 与 Wulff 形 \(W\subset\RR^n\)。设 \(\mathcal P\) 为一族有限 screen-packet 候选，给定几何落点映射
\[
U:\mathcal P\longrightarrow \mathfrak S_{W,m,n}^{\mathrm{tr}},
\qquad
s:=\sigma_W\circ U:\mathcal P\longrightarrow \RR_{>0},
\]
以及有限算术标签映射
\[
\tau:\mathcal P\longrightarrow \mathcal A,
\]
其中 \(\tau\) 记录不能由 Serrin 尺度坐标本身读取的同余、denominator spectrum 或 arithmetic resonance 类型。令 \(\mathcal G\subseteq\mathcal P\) 为大偏差抽样可能命中的 good packet 集合。称有限寄存器
\[
r:\mathcal G\longrightarrow R
\]
在 Serrin 尺度纤维上对 \(\tau\) 忠实，若对任意 \(p,q\in\mathcal G\)，
\[
s(p)=s(q),\qquad r(p)=r(q)
\quad\Longrightarrow\quad
\tau(p)=\tau(q).
\]
则任何这样的寄存器都满足 sharp 下界
\[
\card R\ge
\max_{a\in s(\mathcal G)}
\card\tau\bigl(\mathcal G\cap s^{-1}(a)\bigr).
\]
反过来，存在达到该下界的尺度纤维内寄存器。特别地，若某个尺度纤维中同时出现共振与非共振标签，则只保存 Serrin 尺度 \(s(p)\)，或只保存任意经由 \(s\) 因子化的边界全息读出，都不能作为 arithmetic resonance exclusion 的确定性证书。
\end{theorem}
\begin{proof}
先证明下界。固定 \(a\in s(\mathcal G)\)，并令
\[
T_a:=\tau\bigl(\mathcal G\cap s^{-1}(a)\bigr).
\]
对每个 \(\alpha\in T_a\)，选取一个代表 \(p_\alpha\in\mathcal G\cap s^{-1}(a)\) 满足 \(\tau(p_\alpha)=\alpha\)。若 \(\alpha\ne\beta\)，则 \(p_\alpha,p_\beta\) 位于同一尺度纤维；由寄存器忠实性，必有
\[
r(p_\alpha)\ne r(p_\beta).
\]
因此 \(\alpha\mapsto r(p_\alpha)\) 是从 \(T_a\) 到 \(R\) 的单射，故
\[
\card R\ge \card T_a.
\]
对所有 \(a\in s(\mathcal G)\) 取最大值即得所需下界。

再证 sharp 性。令
\[
M:=\max_{a\in s(\mathcal G)}\card T_a.
\]
取一个固定有限集 \(R_0\) 满足 \(\card R_0=M\)。对每个尺度值 \(a\)，选择一个单射
\[
\iota_a:T_a\hookrightarrow R_0.
\]
定义
\[
r_0(p):=\iota_{s(p)}(\tau(p))\qquad(p\in\mathcal G).
\]
若 \(s(p)=s(q)=a\) 且 \(r_0(p)=r_0(q)\)，则
\[
\iota_a(\tau(p))=\iota_a(\tau(q)).
\]
由于 \(\iota_a\) 为单射，得到 \(\tau(p)=\tau(q)\)。故 \(r_0\) 忠实，且 \(\card R_0=M\)，下界 sharp。

最后说明为何该下界是确定性证书约束，而不是概率估计约束。若概率测度 \(\nu\in\Delta(\mathcal P)\) 满足 \(\nu(\mathcal G)>0\)，则抽样只证明存在某个 \(p\in\mathcal G\)。一旦要把该命中用于 Omega 的对象层审计，必须冻结一个确定包 \(p\) 或冻结某个足以在相关尺度纤维中恢复 \(\tau(p)\) 的寄存器值。另一方面，由定理 \ref{thm:xi-time-part9xbb-serrin-quotient-holography-factors-through-scale}，任意定义在 Serrin 平移商上的 fixed-resolution 边界全息坐标都经由 \(\sigma_W\) 因子化；因此其拉回到 \(\mathcal G\) 后只能经由 \(s\) 读出。若同一 \(s\)-纤维中存在两个 good packets 具有不同 \(\tau\)-标签，则所有只经由 \(s\) 的几何读出在二者上相同，却无法判定其 arithmetic resonance 类型。于是随机抵消命中与算术共振排除必须作为两条正交账本同时登记，尺度残余或有限通量坐标本身不能替代共振寄存器。证毕。
\end{proof}


\endinput

\paragraph{归一化 Serrin 几何因子的累积量真空、中心化 \texorpdfstring{$\zeta$}{zeta} 平坦性与 Hankel 秩一极限}
接受如下唯一外部刚性输入：对有界、不可分解、有限周长且满足 ADR 与全局 \(\beta\)-square-function 条件的粗糙域，Figalli--Zhang 的各向异性 Serrin 定理强制弱解域恰为某个 Wulff 形的平移与伸缩；特别地，该结论覆盖 Lipschitz 域。将此刚性压入本文的 weighted transfer--moment-kernel--weighted zeta--Hankel 接口后，若再施加体积归一化、重心归一化并模去平移，则相应归一化几何因子
\[
\mathfrak G_{K}^{\mathrm{nor}}
\]
退化为单点
\[
\mathfrak G_{K}^{\mathrm{nor}}=\{W_K\}.
\]
以下七条结果表明：在该单状态几何支上，本文框架中的高阶累积量、混合压力曲率、大偏差相图、primitive 长度分层、温度分歧与动力学慢模全部同时塌缩为真空态；因此凡在真实输入核、碰撞核或纤维外置层中观测到的非平凡曲率、冻结、谱隙与 Hankel 高秩现象，都不可能归属于纯归一化 Serrin 几何本体。

\begin{theorem}[归一化 Serrin 几何因子的全阶累积量湮灭]\label{thm:xi-time-part9xc-serrin-geometry-all-cumulants-vanish}
沿用定义 \ref{def:pom-projective-transfer-operator} 的 weighted transfer 口径。设
\[
\psi:\mathfrak G_{K}^{\mathrm{nor}}\to \RR
\]
为任意实值可观测量，并定义加权转移算子
\[
\mathcal L_{\theta,\psi}:C(\mathfrak G_{K}^{\mathrm{nor}})\to C(\mathfrak G_{K}^{\mathrm{nor}}),
\qquad
(\mathcal L_{\theta,\psi}f)(x):=e^{\theta\psi(x)}f(x).
\]
记其 Perron 根与压力为
\[
\lambda_\psi(\theta):=\rho(\mathcal L_{\theta,\psi}),
\qquad
P_\psi(\theta):=\log\lambda_\psi(\theta).
\]
则有严格闭式
\[
\lambda_\psi(\theta)=e^{\theta\psi(W_K)},
\qquad
P_\psi(\theta)=\theta\psi(W_K).
\]
特别地，对一切 \(r\ge 2\)，都有
\[
P_\psi^{(r)}(0)=0.
\]
\end{theorem}
\begin{proof}
由于 \(\mathfrak G_{K}^{\mathrm{nor}}=\{W_K\}\) 为单状态空间，\(C(\mathfrak G_{K}^{\mathrm{nor}})\) 是一维空间。对任意 \(f\in C(\mathfrak G_{K}^{\mathrm{nor}})\)，有
\[
(\mathcal L_{\theta,\psi}f)(W_K)=e^{\theta\psi(W_K)}f(W_K).
\]
因此 \(\mathcal L_{\theta,\psi}\) 在该一维空间上的矩阵表示就是标量
\[
[\,e^{\theta\psi(W_K)}\,].
\]
其唯一特征值亦即 Perron 根正是 \(e^{\theta\psi(W_K)}\)，故
\[
P_\psi(\theta)=\log e^{\theta\psi(W_K)}=\theta\psi(W_K).
\]
于是所有高于一阶的导数全部为零。证毕。
\end{proof}
\leanverified{paper\_xi\_time\_part9xc\_serrin\_geometry\_all\_cumulants\_vanish}

\begin{theorem}[联合压力面的仿射退化]\label{thm:xi-time-part9xc-serrin-joint-pressure-affine-degeneration}
\leanverified{paper\_xi\_time\_part9xc\_serrin\_joint\_pressure\_affine\_degeneration}
设
\[
\psi_1,\psi_2:\mathfrak G_{K}^{\mathrm{nor}}\to\RR
\]
为任意两条可观测轴，并定义联合加权转移算子
\[
(\mathcal L_{\theta_1,\theta_2}f)(x)
:=
e^{\theta_1\psi_1(x)+\theta_2\psi_2(x)}f(x).
\]
记其 Perron 根与二维压力面为
\[
\lambda(\theta_1,\theta_2):=\rho(\mathcal L_{\theta_1,\theta_2}),
\qquad
P(\theta_1,\theta_2):=\log\lambda(\theta_1,\theta_2).
\]
则
\[
P(\theta_1,\theta_2)
=
\theta_1\psi_1(W_K)+\theta_2\psi_2(W_K),
\]
并且任意总阶至少为 \(2\) 的混合偏导都恒为零：
\[
\partial_{\theta_1}^a\partial_{\theta_2}^bP(\theta_1,\theta_2)\equiv 0
\qquad (a+b\ge 2).
\]
\end{theorem}
\begin{proof}
仍由单状态性，
\[
(\mathcal L_{\theta_1,\theta_2}f)(W_K)
=
e^{\theta_1\psi_1(W_K)+\theta_2\psi_2(W_K)}f(W_K).
\]
故其 Perron 根为
\[
\lambda(\theta_1,\theta_2)=e^{\theta_1\psi_1(W_K)+\theta_2\psi_2(W_K)}.
\]
取对数即可得到
\[
P(\theta_1,\theta_2)=\theta_1\psi_1(W_K)+\theta_2\psi_2(W_K),
\]
于是全部二阶及以上偏导恒等消失。证毕。
\end{proof}

\begin{theorem}[大偏差退化与冻结线真空]\label{thm:xi-time-part9xc-serrin-ldp-vacuum}
设 \(T:\mathfrak G_K^{\mathrm{nor}}\to\mathfrak G_K^{\mathrm{nor}}\) 为归一化几何因子上的规范演化。对任意实值可观测量
\[
\psi:\mathfrak G_K^{\mathrm{nor}}\to\RR
\]
定义 Birkhoff 和
\[
S_n\psi:=\sum_{t=0}^{n-1}\psi\circ T^t.
\]
则对一切 \(n\ge 1\)，逐点恒有
\[
S_n\psi=n\,\psi(W_K).
\]
相应经验均值的速率函数满足
\[
I_\psi(a)=
\begin{cases}
0,& a=\psi(W_K),\\[2mm]
+\infty,& a\neq \psi(W_K).
\end{cases}
\]
\leanverified{paper\_xi\_time\_part9xc\_serrin\_ldp\_vacuum}
\end{theorem}
\begin{proof}
由于 \(\mathfrak G_K^{\mathrm{nor}}\) 只有唯一状态，规范演化只能满足
\[
T(W_K)=W_K.
\]
因此
\[
S_n\psi(W_K)=\sum_{t=0}^{n-1}\psi(T^tW_K)=\sum_{t=0}^{n-1}\psi(W_K)=n\,\psi(W_K).
\]
故经验均值 \(S_n\psi/n\) 的分布对每个 \(n\) 都是点质量 \(\delta_{\psi(W_K)}\)。按大偏差定义，其速率函数仅在该点为零，其余处为 \(+\infty\)。证毕。
\end{proof}

\begin{theorem}[整数观测的确定同余律与长度一原子性]\label{thm:xi-time-part9xc-serrin-congruence-dirac-length-one}
设
\[
E:\mathfrak G_K^{\mathrm{nor}}\to\ZZ
\]
为整数值可观测量，并记
\[
e_0:=E(W_K).
\]
则对每个 \(n\ge 1\)，都有
\[
S_nE=n e_0.
\]
因此，对任意模数 \(m\ge 2\)，长度 \(n\) 的模 \(m\) 统计分布是严格的 Dirac 质量
\[
\mu_{n,m}=\delta_{\overline{n e_0}},
\qquad
\overline{n e_0}\in\ZZ/m\ZZ.
\]
其 Fourier 系数满足
\[
\widehat\mu_{n,m}(j)=e^{2\pi i j n e_0/m}
\qquad (0\le j\le m-1).
\]
若进一步按文稿既有的 weighted periodic/primitive 规范记周期数据为 \(a_n(u)\)、primitive 数据为 \(p_n(u)\)，则有
\[
a_n(u)=u^{n e_0},
\qquad
p_1(u)=u^{e_0},
\qquad
p_n(u)=0\quad (n\ge 2).
\]
\leanverified{paper\_xi\_time\_part9xc\_serrin\_congruence\_dirac\_length\_one}
\end{theorem}
\begin{proof}
由定理 \ref{thm:xi-time-part9xc-serrin-ldp-vacuum}，
\[
S_nE=nE(W_K)=n e_0.
\]
故模 \(m\) 残差只可能取到唯一值 \(\overline{n e_0}\in \ZZ/m\ZZ\)，从而
\[
\mu_{n,m}=\delta_{\overline{n e_0}}.
\]
对 Dirac 质量的 Fourier 变换立即给出
\[
\widehat\mu_{n,m}(j)=e^{2\pi i j n e_0/m}.
\]

另一方面，在唯一状态上长度 \(n\) 的 weighted periodic 贡献只有一项，且其总权为
\[
u^{S_nE}=u^{n e_0},
\]
故
\[
a_n(u)=u^{n e_0}.
\]
再由标准 Witt--M\"obius 反演，序列 \((a_n(u))_{n\ge 1}\) 只对应一个长度 \(1\) 的 primitive 轨道，其权为 \(u^{e_0}\)；于是
\[
p_1(u)=u^{e_0},
\qquad
p_n(u)=0\quad(n\ge 2).
\]
证毕。
\end{proof}

\begin{corollary}[中心化 weighted \texorpdfstring{$\zeta$}{zeta} 的温度平坦性与分歧真空]\label{cor:xi-time-part9xc-serrin-centered-zeta-temperature-flatness}
沿用定理 \ref{thm:xi-time-part9xc-serrin-congruence-dirac-length-one}，定义中心化观测
\[
\widetilde E:=E-e_0.
\]
则
\[
S_n\widetilde E=0
\qquad(\forall n\ge 1).
\]
因而相应的 weighted zeta 与完成化行列式满足
\[
\zeta_{\widetilde E}(z,u)=\frac{1}{1-z},
\qquad
\Delta_{\widetilde E}(z,u)=1-z.
\]
特别地，作为温度参数 \(u\) 的函数，\(\zeta_{\widetilde E}\) 与 \(\Delta_{\widetilde E}\) 都是常值整函数；因此不存在任何有限温度分歧点、判别式零簇、Newman 型阈值或有限温度临界参数。
\end{corollary}
\begin{proof}
由 \(\widetilde E(W_K)=0\)，对每个 \(n\ge 1\) 都有
\[
S_n\widetilde E=n\widetilde E(W_K)=0.
\]
因此按 weighted zeta 的定义，
\[
\log \zeta_{\widetilde E}(z,u)
=
\sum_{n\ge 1}\frac{u^{S_n\widetilde E}}{n}z^n
=
\sum_{n\ge 1}\frac{z^n}{n}
=
-\log(1-z).
\]
指数化便得
\[
\zeta_{\widetilde E}(z,u)=\frac{1}{1-z}.
\]
其倒数完成化行列式于是满足
\[
\Delta_{\widetilde E}(z,u)=\zeta_{\widetilde E}(z,u)^{-1}=1-z.
\]
两式都与 \(u\) 无关，因此关于温度变量 \(u\) 的一切分歧、判别式和有限半径现象全部消失。证毕。
\end{proof}

\begin{theorem}[Hankel--Prony 秩一刚性]\label{thm:xi-time-part9xc-serrin-hankel-prony-rank-one}
设
\[
\psi:\mathfrak G_K^{\mathrm{nor}}\to\RR
\]
为任意实值可观测量，并定义矩序列
\[
M_n:=e^{S_n\psi}=e^{n\psi(W_K)}
\qquad (n\ge 0).
\]
对任意 \(N\ge 0\)，令 Hankel 矩阵
\[
H^{(N)}:=(M_{i+j})_{0\le i,j\le N}.
\]
则恒有
\[
\operatorname{rank} H^{(N)}=1.
\]
等价地，一切 \(2\times 2\) 子式都严格消失：
\[
M_{i+j}M_{k+\ell}-M_{i+\ell}M_{k+j}=0
\qquad
(\forall i,j,k,\ell\ge 0).
\]
\leanverified{paper\_xi\_time\_part9xc\_serrin\_hankel\_prony\_rank\_one}
\end{theorem}
\begin{proof}
记
\[
a_i:=e^{i\psi(W_K)}.
\]
则
\[
M_{i+j}=e^{(i+j)\psi(W_K)}=a_i a_j.
\]
因此
\[
H^{(N)}=(a_i a_j)_{0\le i,j\le N}=a\,a^{\mathsf T},
\]
其中
\[
a=(a_0,\dots,a_N)^{\mathsf T}.
\]
故 \(H^{(N)}\) 是非零秩一外积矩阵，遂有 \(\operatorname{rank}H^{(N)}=1\)。秩一矩阵的一切 \(2\times 2\) 子式都为零，所示行列式恒等式随之成立。证毕。
\end{proof}

\begin{proposition}[归一化几何信道的动力学真空]\label{prop:xi-time-part9xc-serrin-normalized-transfer-vacuum}
沿用定理 \ref{thm:xi-time-part9xc-serrin-geometry-all-cumulants-vanish} 的记号。以 Perron 根归一化转移算子：
\[
\widehat{\mathcal L}_{\theta,\psi}
:=
e^{-\theta\psi(W_K)}\mathcal L_{\theta,\psi}.
\]
则
\[
\widehat{\mathcal L}_{\theta,\psi}=I
\]
在 \(C(\mathfrak G_K^{\mathrm{nor}})\) 上恒成立。因而对任意可观测量 \(f,g\) 与任意 \(n\ge 0\)，都有
\[
f(W_K)\,(\widehat{\mathcal L}_{\theta,\psi}^{\,n}g)(W_K)
=
f(W_K)g(W_K),
\]
即相关函数恒为常数。进一步，\(\widehat{\mathcal L}_{\theta,\psi}\) 的谱仅含单点 \(\{1\}\)；因此任何由主谱/次谱比、谱隙或次谱半径构造的混合时间、弛豫时间与慢模尺度，都只能退化为平凡量。
\end{proposition}
\begin{proof}
对任意 \(f\in C(\mathfrak G_K^{\mathrm{nor}})\)，由定理
\ref{thm:xi-time-part9xc-serrin-geometry-all-cumulants-vanish}，
\[
(\mathcal L_{\theta,\psi}f)(W_K)=e^{\theta\psi(W_K)}f(W_K).
\]
两边同乘 \(e^{-\theta\psi(W_K)}\) 即得
\[
(\widehat{\mathcal L}_{\theta,\psi}f)(W_K)=f(W_K),
\]
故 \(\widehat{\mathcal L}_{\theta,\psi}=I\)。

于是对任意 \(n\ge 0\) 都有
\[
\widehat{\mathcal L}_{\theta,\psi}^{\,n}=I,
\]
从而
\[
f(W_K)\,(\widehat{\mathcal L}_{\theta,\psi}^{\,n}g)(W_K)
=
f(W_K)g(W_K).
\]
最后，恒等算子的谱正是 \(\{1\}\)，故不存在任何非平凡次谱结构；相应的谱隙、弛豫与混合尺度皆失去非平凡内容。证毕。
\end{proof}

\noindent
综上，归一化 Serrin 几何因子在本文的 weighted transfer--pressure--large-deviation--weighted zeta--Hankel 全链条上表现为一个严格的热力学真空：高阶累积量全部消失，联合压力面完全仿射，大偏差退化为单点支撑，primitive 层只剩长度 \(1\) 原子，中心化 \(\zeta\) 对温度完全失明，Hankel 矩阵永远停留在秩一射线，而归一化动力学则塌缩为恒等算子。因而，真实输入 \(40\)-状态核中已经明确出现的非零局部累积量与 Edgeworth 常数（定理 \ref{thm:xi-real-input-cramer-germ-quadratic-field-rigidity}）、冻结边界与基态简并指数（定理 \ref{thm:xi-time-part65c-escort-renyi-finite-size-mainterm}、定理 \ref{thm:xi-time-part57f-legendre-hard-wall-boundary-cost}）、以及近简并慢模与 \(\sqrt u\)-级弛豫时间（定理 \ref{thm:xi-output-potential-endpoint-near-degeneracy-sqrtu-relaxation}、结论 \ref{con:xi-time-part9o-prime-modulus-relaxation-second-order}），都只能来自几何之外的算术、同步或纤维外置层，而不能归属于纯归一化 Serrin 几何本体。

\endinput

\paragraph{共动 Jensen 缺陷、有限缺陷层析、位复杂度差额与 \texorpdfstring{\v{C}ech}{Cech} 谱阶的统一接口}
在共动 Jensen 缺陷的无高度偏移界、有限缺陷扫描剖面的 Fourier--Laplace 注入性、边界扫描的有限秩核化、端点二次压缩所诱导的显式位复杂度门槛，以及 unitary-slice 闭包下 \(\mathsf{NULL}\) 的 \v{C}ech--cyclotomic 量子化都已建立之后，还可继续抽取出下列五条交叉后果。它们分别把一致缺陷界重写为非平凡零点的显式模数化集中，把一维扫描、二维 Fourier--Laplace 指纹与有限秩 RKHS 核压缩为同一恢复对象，把单频等距采样的最小层析窗写成严格的 Prony 型 \(2\kappa\) 阈值，把固定图表与精确共动局部化之间的位复杂度差额闭合为显式对数公式，并把 \v{C}ech 型空值分支压缩为一个由单位根阶读出的最小谱不变量。

\begin{theorem}[临界线集中的显式模数化判别]\label{thm:xi-time-part9xd-explicit-critical-line-concentration-modulus}
设
\[
(\varrho_n)_{n\ge 1}\subset(0,1),
\qquad
(\varepsilon_n)_{n\ge 1}\subset[0,\infty)
\]
满足
\[
\varrho_n\uparrow 1,
\qquad
\varepsilon_n\downarrow 0,
\]
并且对每个 \(n\) 都有一致缺陷界
\[
\sup_{x_0\in\RR}D_{\zeta,x_0}(\varrho_n)\le \varepsilon_n.
\]
定义
\[
\Delta_n:=\frac{1-\varrho_n e^{-\varepsilon_n}}{1+\varrho_n e^{-\varepsilon_n}}.
\]
则任意非平凡零点 \(\rho=\beta+\mathrm{i}\gamma\) 都满足
\[
\beta\ge \frac12-\Delta_n
\qquad (\forall n).
\]
特别地，\(\Delta_n\to 0\)，从而全部非平凡零点都被压到临界线。反之，RH 成立时必存在这样的 \((\varrho_n,\varepsilon_n)\)。因此，RH 等价于存在一列显式模数
\[
\Delta_n\downarrow 0
\]
使一切非平凡零点都满足统一集中不等式
\[
\beta\ge \frac12-\Delta_n
\qquad (\forall n).
\]
\end{theorem}

\begin{proof}
设 \(\rho=\beta+\mathrm{i}\gamma\) 为任意非平凡零点。若 \(\beta\neq \frac12\)，则由函数方程对称性可不失一般性设
\[
\delta:=\frac12-\beta>0.
\]
对任意 \(n\)，取共动中心 \(x_0=\gamma\)。由假设
\[
D_{\zeta,\gamma}(\varrho_n)\le \sup_{x_0\in\RR}D_{\zeta,x_0}(\varrho_n)\le \varepsilon_n,
\]
再应用定理 \ref{thm:xi-comoving-defect-implies-delta-bound}，得到
\[
\delta\le \frac{1-\varrho_n e^{-\varepsilon_n}}{1+\varrho_n e^{-\varepsilon_n}}=\Delta_n.
\]
于是
\[
\beta=\frac12-\delta\ge \frac12-\Delta_n
\qquad (\forall n).
\]
由 \(\varrho_n\uparrow 1\) 与 \(\varepsilon_n\downarrow 0\) 可知 \(\varrho_n e^{-\varepsilon_n}\to 1\)，从而 \(\Delta_n\to 0\)。故 \(\beta<\frac12\) 不可能发生；再由 \(\rho\mapsto 1-\rho\) 的对称性，\(\beta>\frac12\) 亦被排除，遂得 RH。

反向命题正是定理 \ref{thm:xi-rh-iff-uniform-comoving-defect} 的 \((i)\Rightarrow(ii)\) 方向。证毕。
\end{proof}

\begin{theorem}[有限缺陷扫描的对角刚性与二维 Fourier--Laplace 指纹的唯一提升]\label{thm:xi-time-part9xd-finite-defect-diagonal-rigidity-fingerprint-kernel}
设有限缺陷测度
\[
\nu=\sum_{j=1}^{\kappa}m_j\,\delta_{(\gamma_j,\delta_j)},
\qquad
m_j\in\NN,\ \delta_j\in(0,1),
\]
并定义一维扫描剖面
\[
H_\nu(x):=\sum_{j=1}^{\kappa}\frac{4m_j\delta_j}{(x-\gamma_j)^2+(1+\delta_j)^2}.
\]
再定义二维 Fourier--Laplace 指纹
\[
F_\nu(\xi,s):=
4\pi\sum_{j=1}^{\kappa}
m_j\frac{\delta_j}{1+\delta_j}\,
e^{-s\delta_j}e^{-\mathrm{i}\gamma_j\xi}
\qquad
\bigl((\xi,s)\in\RR\times[0,\infty)\bigr),
\]
以及有限秩核
\[
K_\nu(x,y):=
\sum_{j=1}^{\kappa}
\frac{4m_j\delta_j}
{(x-\gamma_j+\mathrm{i}(1+\delta_j))(y-\gamma_j-\mathrm{i}(1+\delta_j))}.
\]
则对任意非空开区间 \(I\subset\RR\)，下列四类数据彼此唯一决定：
\[
\nu,
\qquad
\left.H_\nu\right|_I,
\qquad
F_\nu,
\qquad
K_\nu.
\]
特别地，有限缺陷口径下仅知道核的对角限制
\[
K_\nu(x,x)=H_\nu(x)
\]
在任意非空开区间上的取值，已经足以唯一恢复整个有限秩核 \(K_\nu\) 与全部二维缺陷几何 \(\nu\)。
\end{theorem}

\begin{proof}
首先，把重复支点合并后可设 \((\gamma_j,\delta_j)\) 两两不同。

由定理 \ref{thm:xi-comoving-horizon-scan-fourier-inversion} 的 Fourier 闭式，
\[
\widehat H_\nu(\xi)
=
4\pi\sum_{j=1}^{\kappa}
m_j\frac{\delta_j}{1+\delta_j}\,
e^{-(1+\delta_j)|\xi|}e^{-\mathrm{i}\gamma_j\xi}
=
e^{-|\xi|}F_\nu(\xi,|\xi|).
\]
因此若两组有限缺陷数据在某个非空开区间 \(I\) 上给出同一 \(H_\nu\)，则由实解析性知两者在 \(\RR\) 上处处一致，从而其 Fourier 变换处处一致。限制到 \(\xi>0\) 后得到一条有限指数和恒等式；由互异指数函数的线性无关性（等价地，Vandermonde 可逆性），支持点与权重完全一致，故 \(\left.H_\nu\right|_I\) 唯一决定 \(\nu\)。反向由定义显然。

其次，令
\[
\widetilde\nu:=
4\pi\sum_{j=1}^{\kappa}
m_j\frac{\delta_j}{1+\delta_j}\,
\delta_{(\gamma_j,\delta_j)}.
\]
则
\[
F_\nu(\xi,s)=\int_{\RR\times(0,1)}e^{-\mathrm{i}\gamma\xi}e^{-s\delta}\,\dd\widetilde\nu(\gamma,\delta),
\]
即 \(F_\nu\) 正是有限正测度 \(\widetilde\nu\) 的 Fourier--Laplace 变换。由 Fourier--Laplace 变换在有限正测度类上的唯一性（同命题 \ref{prop:xi-defect-density-cone-choquet-extreme} 证明中所用机制），\(F_\nu\) 唯一决定 \(\widetilde\nu\)，因而唯一决定 \(\nu\)；反向则由定义成立。

最后，对核 \(K_\nu\) 应用定理 \ref{thm:xi-basepoint-scan-finite-rank-rkhs}，可得
\[
K_\nu(x,x)=H_\nu(x)\qquad (x\in\RR).
\]
因此 \(K_\nu\) 的对角函数等于扫描剖面。另一方面，命题 \ref{prop:xi-basepoint-scan-kernel-uniqueness} 表明：在同一 Cauchy 极点骨架类型的有限秩核类中，对角函数已经唯一决定整个核。故
\[
K_\nu\Longleftrightarrow H_\nu\Longleftrightarrow \nu,
\]
再与上一步合并即得
\[
\nu
\Longleftrightarrow
\left.H_\nu\right|_I
\Longleftrightarrow
F_\nu
\Longleftrightarrow
K_\nu.
\]
证毕。
\end{proof}

\begin{theorem}[单频有限采样的精确层析重构]\label{thm:xi-time-part9xd-single-frequency-sampling-exact-tomography}
\leanverified{paper\_xi\_time\_part9xd\_single\_frequency\_sampling\_exact\_tomography}
设
\[
\nu=\sum_{j=1}^{\kappa}a_j\,\delta_{(\gamma_j,\delta_j)},
\qquad
a_j>0,\ \delta_j\in(0,1),
\]
并假设 \(\delta_1,\dots,\delta_\kappa\) 两两不同，且存在先验窗口
\[
|\gamma_j|\le G\qquad(1\le j\le \kappa).
\]
固定
\[
0<|\xi|<\frac{\pi}{G},
\qquad
\Delta s>0,
\qquad
s_0>0,
\]
并定义
\[
F_\nu(\xi,s):=
4\pi\sum_{j=1}^{\kappa}
a_j\frac{\delta_j}{1+\delta_j}\,
e^{-s\delta_j}e^{-\mathrm{i}\gamma_j\xi}.
\]
记
\[
y_n:=F_\nu(\xi,s_0+n\Delta s)
\qquad (n=0,1,\dots,2\kappa-1).
\]
则有限样本窗
\[
(y_0,\dots,y_{2\kappa-1})
\]
已经唯一决定全部原子参数
\[
\{(\gamma_j,\delta_j,a_j)\}_{j=1}^{\kappa}.
\]
换言之，在 \(\delta_j\) 互异且 \(\gamma_j\) 有界的泛型口径下，\(\kappa\) 个缺陷原子无需二维连续扫描；单一频率上的 \(2\kappa\) 个等距样本即足以完成精确重构。
\end{theorem}

\begin{proof}
把样本写成
\[
y_n=
\sum_{j=1}^{\kappa}c_j\lambda_j^{\,n},
\qquad
\lambda_j:=e^{-\Delta s\,\delta_j},
\qquad
c_j:=
4\pi a_j\frac{\delta_j}{1+\delta_j}\,
e^{-s_0\delta_j}e^{-\mathrm{i}\gamma_j\xi}.
\]
由于 \(\delta_j\) 两两不同，\(\lambda_j\) 两两不同；又因 \(a_j>0\) 与 \(\delta_j>0\)，故 \(c_j\neq 0\)。于是 \((y_n)\) 构成一般位置上的 Prony 序列。由推论 \ref{cor:xi-prony-decision-vs-reconstruction-one-sample-gap}，长度 \(2\kappa\) 的样本窗 \((y_0,\dots,y_{2\kappa-1})\) 唯一决定节点--权重多重集
\[
\{(\lambda_j,c_j)\}_{j=1}^{\kappa}
\]
（至多差一个排列）。

从而
\[
\delta_j=-\frac{1}{\Delta s}\log \lambda_j
\]
被唯一恢复。再由
\[
|c_j|=
4\pi a_j\frac{\delta_j}{1+\delta_j}e^{-s_0\delta_j}
\]
可唯一恢复 \(a_j\)。

最后，由 \(c_j/|c_j|=e^{-\mathrm{i}\gamma_j\xi}\) 得到相位。若
\[
e^{-\mathrm{i}\gamma\xi}=e^{-\mathrm{i}\gamma'\xi}
\qquad
(\gamma,\gamma'\in[-G,G]),
\]
则 \((\gamma-\gamma')\xi\in 2\pi\ZZ\)。但由 \(0<|\xi|<\pi/G\) 可知
\[
|(\gamma-\gamma')\xi|\le 2G|\xi|<2\pi,
\]
故只能有 \((\gamma-\gamma')\xi=0\)，即 \(\gamma=\gamma'\)。因此
\[
\gamma\longmapsto e^{-\mathrm{i}\gamma\xi}
\]
在 \([-G,G]\) 上单射，遂每个 \(\gamma_j\) 亦被唯一恢复。证毕。
\end{proof}

\begin{theorem}[二次红移代价与共动前缀收益的精确差额公式]\label{thm:xi-time-part9xd-redshift-tax-comoving-prefix-gap}
\leanverified{paper\_xi\_time\_part9xd\_redshift\_tax\_comoving\_prefix\_gap}
固定高度窗 \(|\gamma|\le T\) 与偏移下界 \(\delta\ge \delta_0\in(0,\tfrac12)\)。定义固定图表制度下的最小必要端点精度为
\[
p_{\mathrm{fix}}^{\mathrm{nec}}(T,\delta_0)
:=
\inf\Bigl\{
p\in\RR:\ 2^{-p}\ \text{足以在固定 Cayley 图表中统一分离全部候选}
\Bigr\}.
\]
若允许外置精确共动前缀并取 \(x_0=\gamma\)，定义可实现证书精度
\[
p_{\mathrm{mov}}^{\mathrm{cert}}(\delta_0)
:=
\log_2\!\Bigl(\frac{(1+\delta_0)^2}{4\delta_0}\Bigr).
\]
则有显式差额下界
\[
p_{\mathrm{fix}}^{\mathrm{nec}}(T,\delta_0)
-p_{\mathrm{mov}}^{\mathrm{cert}}(\delta_0)
\ge
\log_2\!\Bigl(1+\frac{T^2}{(1+\delta_0)^2}\Bigr).
\]
特别地，
\[
p_{\mathrm{fix}}^{\mathrm{nec}}(T,\delta_0)
-p_{\mathrm{mov}}^{\mathrm{cert}}(\delta_0)
=
2\log_2 T+O_{\delta_0}(1)
\qquad (T\to\infty).
\]
因此，共动前缀并不消去制度成本，而是恰好剥离固定图表中的二次红移主项；其可回收位数的主导阶正是 \(2\log_2 T\)。
\end{theorem}

\begin{proof}
由定理 \ref{thm:xi-offcritical-endpoint-resolution-lower-bound}，固定图表下任何统一分离策略都必须满足
\[
p_{\mathrm{fix}}^{\mathrm{nec}}(T,\delta_0)
\ge
\log_2\!\Bigl(\frac{T^2+(1+\delta_0)^2}{4\delta_0}\Bigr).
\]

另一方面，若外置精确共动前缀并取 \(x_0=\gamma\)，则由附录定理 \ref{thm:app-comoving-horizon-localization} 及推论 \ref{cor:xi-comoving-prefix-bit-budget-gain}，
\[
1-|w_{\rho,\gamma}|^2
=
\frac{4\delta}{(1+\delta)^2}
\ge
\frac{4\delta_0}{(1+\delta_0)^2}.
\]
因而取任意阈值
\[
2^{-p}\le \frac{4\delta_0}{(1+\delta_0)^2}
\]
即可在精确共动中心上形成统一证书，这正对应于
\[
p_{\mathrm{mov}}^{\mathrm{cert}}(\delta_0)
=
\log_2\!\Bigl(\frac{(1+\delta_0)^2}{4\delta_0}\Bigr).
\]

两式相减得到
\[
p_{\mathrm{fix}}^{\mathrm{nec}}(T,\delta_0)
-p_{\mathrm{mov}}^{\mathrm{cert}}(\delta_0)
\ge
\log_2\!\Bigl(
\frac{T^2+(1+\delta_0)^2}{(1+\delta_0)^2}
\Bigr)
=
\log_2\!\Bigl(1+\frac{T^2}{(1+\delta_0)^2}\Bigr).
\]
令 \(T\to\infty\) 即得
\[
\log_2\!\Bigl(1+\frac{T^2}{(1+\delta_0)^2}\Bigr)
=
2\log_2 T+O_{\delta_0}(1).
\]
证毕。
\end{proof}

\begin{theorem}[\v{C}ech 空值分支的最小谱阶不变量]\label{thm:xi-time-part9xd-cech-null-minimal-spectral-order}
设 \(A\) 为有限阿贝尔群，\([\alpha]\in \check H^2(\{b_i\},A)\) 为前缀站点上的二阶粘合障碍类。对每个角色 \(\chi\in\widehat A\)，记
\[
\chi_*([\alpha])\in \check H^2(\{b_i\},\TT)
\]
为系数推前，并记其阶为
\[
\ord\!\bigl(\chi_*([\alpha])\bigr)\in \NN.
\]
定义最小谱阶
\[
\mathfrak o([\alpha]):=
\begin{cases}
\min\bigl\{\ord(\chi_*([\alpha])):\ \chi\in\widehat A,\ \chi_*([\alpha])\neq 0\bigr\},& [\alpha]\neq 0,\\[3pt]
\infty,& [\alpha]=0.
\end{cases}
\]
再记 \(B_\chi\) 为相应的 \(\chi\)-扭曲转移矩阵，\(\lambda:=\rho(B_{\ind})\) 为非扭曲 Perron 根。则：
\begin{enumerate}
  \item 若 \([\alpha]\neq 0\)，则 \(\mathfrak o([\alpha])\) 恰等于所有非平凡归一化行列式
  \[
  \det\!\Bigl(I-\frac{t}{\lambda}B_\chi\Bigr)
  \qquad (\chi\neq \ind)
  \]
  在单位圆 \(|t|=1\) 上所出现的本原单位根零点的最小阶；
  \item
  \[
  [\alpha]=0
  \iff
  \text{对每个非平凡 }\chi\neq \ind,\ 
  \det\!\Bigl(I-\frac{t}{\lambda}B_\chi\Bigr)
  \text{ 在 }|t|=1\text{ 上无非平凡本原单位根零点};
  \]
  \item 在 unitary-slice 协议下，该单位根谱签名只能实现 \(\mathrm{Cech}\) 型失败见证，而不能被 \(\mathrm{Out}\) 型或 \(\mathrm{Hard}\) 型失败见证模拟。
\end{enumerate}
\end{theorem}

\begin{proof}
由定理 \ref{thm:cdim-cech-cyclotomic-quantization}，对每个非平凡角色 \(\chi\neq \ind\) 与每个整数 \(r>1\)，都有严格对应
\[
\ord(\chi_*([\alpha]))=r
\iff
\det\!\Bigl(I-\frac{t}{\lambda}B_\chi\Bigr)
\text{ 在 }|t|=1\text{ 上出现一个 }r\text{ 阶本原单位根零点}.
\]
因此当 \([\alpha]\neq 0\) 时，对全部 \(\chi\neq \ind\) 取最小值，立得
\[
\mathfrak o([\alpha])
=
\min\Bigl\{
\text{单位圆上出现的本原单位根零点阶}
\Bigr\},
\]
这就是第 \(1\) 条。

第 \(2\) 条的“\(\Rightarrow\)”方向显然：若 \([\alpha]=0\)，则一切 \(\chi_*([\alpha])\) 都为零，故由同一定理不存在非平凡本原单位根零点。反之，若对所有非平凡 \(\chi\) 都不存在该类零点，则定理 \ref{thm:cdim-cech-cyclotomic-quantization} 给出 \(\chi_*([\alpha])=0\)。有限阿贝尔群的角色族分离元素；对任一 \(A\)-值 \(2\)-余圈代表逐点施加全部角色，可知所有推前类皆零只能意味着原类本身为零，因此 \([\alpha]=0\)。

最后，由定义 \ref{def:cdim-null-witness-tag}，\([\alpha]\neq 0\) 的谱签名对应的正是 \(\mathrm{Cech}\) 型见证；而命题 \ref{prop:cdim-null-witness-nonsubstitutable} 与定理 \ref{thm:xi-null-orthogonal-trichotomy-resource-nonsubstitutable} 说明 \(\mathrm{Out}\)、\(\mathrm{Cech}\)、\(\mathrm{Hard}\) 三类见证在体系内正交且不可替代。故该单位根谱签名不能由 \(\mathrm{Out}\) 型或 \(\mathrm{Hard}\) 型失败机制模拟。证毕。
\end{proof}

上述五条结果表明，共动 Jensen 缺陷的一致小化、有限缺陷扫描的一维对角、二维 Fourier--Laplace 指纹、有限秩 Cauchy--RKHS 核与 \v{C}ech 型单位根谱签名在本文并非彼此分离的接口。对角数据经解析延拓与核唯一性恢复整个有限秩核，核与有限原子测度互定，测度进一步决定单频采样的 Prony 窗口，而由单位根零点阶读出的最小谱阶则把纯二阶粘合障碍压缩为一个可审计整数不变量；与之并行，固定图表与共动图表之间的位复杂度差额则给出这条恢复链在制度层面的精确红移成本。

\endinput

\paragraph{末位二相、对数似然曲率、escort 涨落与 \texorpdfstring{$\zeta$}{zeta} 提取的统一刚性}
在原始 bin-fold 切片的固定二点实验、escort 半轴的 Fisher 平坦化、fold-gauge 常数的 Bernoulli 层级以及单一幂散度的严格可辨识性已经确立之后，还可把这些分散接口进一步压缩为下列六条综合后果。它们分别刻画块均匀化模型上的末位精确充分化、对数似然比的极限二点律、escort Fisher 密度与似然矩势之间的曲率同一、残差序参量的涨落--响应恒等式、规范常数对偶整数 \(\zeta\) 值的逐阶提取，以及单一幂散度标量对全部极限信息几何的统一重构。

\begin{theorem}[块均匀化 bin-fold 模型的末位精确充分化]\label{thm:xi-time-part9xe-block-uniformized-lastbit-exact-fdiv}
\leanverified{paper\_xi\_time\_part9xe\_block\_uniformized\_lastbit\_exact\_fdiv}
沿用定理 \ref{thm:xi-fold-lastbit-asymptotic-sufficiency} 与定理 \ref{thm:xi-time-part9sa-uniform-baseline-fixed-binary-experiment} 的记号。定义末位推前
\[
\mu_m^{\mathrm{bit}}:=(\ell_m)_\#\mu_m,
\qquad
u_m^{\mathrm{bit}}:=(\ell_m)_\#u_m,
\]
以及块均匀化分布
\[
\overline\mu_m:=K_{m\#}\mu_m^{\mathrm{bit}},
\qquad
K_m(i,\cdot):=u_m(\cdot\mid A_i)
\qquad (i=0,1).
\]
则对任意凸函数 \(f:(0,\infty)\to\RR\) 满足 \(f(1)=0\)，都有精确恒等式
\[
D_f(\overline\mu_m\Vert u_m)
=
D_f(\mu_m^{\mathrm{bit}}\Vert u_m^{\mathrm{bit}}).
\]
特别地，对任意 \(\alpha>0\)，有精确幂和恒等式
\[
\sum_{w\in X_m}\overline\mu_m(w)^\alpha u_m(w)^{1-\alpha}
=
\sum_{i=0}^{1}\bigl(\mu_m^{\mathrm{bit}}(i)\bigr)^\alpha
\bigl(u_m^{\mathrm{bit}}(i)\bigr)^{1-\alpha}.
\]
此外，
\[
D_f(\mu_m\Vert u_m)
=
D_f(\mu_m^{\mathrm{bit}}\Vert u_m^{\mathrm{bit}})
+O_f\!\left(\left(\frac{\varphi}{2}\right)^m\right).
\]
等价地，块均匀化模型相对稳定层均匀基线的全部有限尺度可辨识信息已经严格压缩到末位一比特；而原始模型只在指数小误差内偏离这一精确二相结构。
\end{theorem}
\begin{proof}
对 \(w\in A_i\)，由定义有
\[
\overline\mu_m(w)
=
\frac{\mu_m^{\mathrm{bit}}(i)}{u_m^{\mathrm{bit}}(i)}\,u_m(w),
\]
因此似然比
\[
\frac{\overline\mu_m(w)}{u_m(w)}
\]
在每个块 \(A_i\) 上为常数，恰等于
\[
\frac{\mu_m^{\mathrm{bit}}(i)}{u_m^{\mathrm{bit}}(i)}.
\]
于是
\[
\begin{aligned}
D_f(\overline\mu_m\Vert u_m)
&=
\sum_{i=0}^{1}\sum_{w\in A_i}
u_m(w)\,
f\!\left(\frac{\mu_m^{\mathrm{bit}}(i)}{u_m^{\mathrm{bit}}(i)}\right)\\
&=
\sum_{i=0}^{1}
u_m^{\mathrm{bit}}(i)\,
f\!\left(\frac{\mu_m^{\mathrm{bit}}(i)}{u_m^{\mathrm{bit}}(i)}\right)
=
D_f(\mu_m^{\mathrm{bit}}\Vert u_m^{\mathrm{bit}}),
\end{aligned}
\]
这就证明了第一式。对任意 \(\alpha>0\)，同样按块分组便得到
\[
\begin{aligned}
\sum_{w\in X_m}\overline\mu_m(w)^\alpha u_m(w)^{1-\alpha}
&=
\sum_{i=0}^{1}\sum_{w\in A_i}
\left(
\frac{\mu_m^{\mathrm{bit}}(i)}{u_m^{\mathrm{bit}}(i)}
\right)^\alpha
u_m(w)\\
&=
\sum_{i=0}^{1}
\bigl(\mu_m^{\mathrm{bit}}(i)\bigr)^\alpha
\bigl(u_m^{\mathrm{bit}}(i)\bigr)^{1-\alpha}.
\end{aligned}
\]
另一方面，结论 \ref{con:xi-terminal-foldbin-lastbit-sufficient-statistic} 已指出 \(\overline\mu_m\) 亦可精确写成
\[
\overline\mu_m(w)=Z_m^{-1}u_m(w)e^{-\beta_m\ell_m(w)},
\]
故上述精确充分化也可视为该单参数指数倾斜的直接后果。最后，由定理 \ref{thm:xi-time-part67-lastbit-terminal-csiszar-collapse}，
\[
D_f(\mu_m\Vert u_m)
=
D_f\!\bigl((\ell_m)_\#\mu_m\Vert(\ell_m)_\#u_m\bigr)
+O_f\!\left(\left(\frac{\varphi}{2}\right)^m\right),
\]
而
\[
(\ell_m)_\#\mu_m=\mu_m^{\mathrm{bit}},
\qquad
(\ell_m)_\#u_m=u_m^{\mathrm{bit}},
\]
故得到最后一式。证毕。
\end{proof}

\begin{theorem}[原始切片对数似然比的极限二点律与全阶累积母函数]\label{thm:xi-time-part9xe-loglikelihood-two-point-cgf}
沿用定理 \ref{thm:xi-time-part9sa-uniform-baseline-fixed-binary-experiment} 的记号，定义对数似然比
\[
L_m(w):=\log\frac{\mu_m(w)}{u_m(w)}.
\]
则在 \(u_m\)-测度下，
\[
L_m\Longrightarrow L_u,
\]
其中 \(L_u\) 为二点分布
\[
\mathbf P(L_u=\ell_0)=\frac1\varphi,
\qquad
\mathbf P(L_u=\ell_1)=\frac1{\varphi^2},
\]
\[
\ell_0:=2\log\varphi-\frac12\log 5,
\qquad
\ell_1:=\log\varphi-\frac12\log 5.
\]
因此其矩母函数显式为
\[
M_L(t):=\mathbf E\!\left[e^{tL_u}\right]
=
\frac1\varphi\left(\frac{\varphi^2}{\sqrt5}\right)^t
+\frac1{\varphi^2}\left(\frac{\varphi}{\sqrt5}\right)^t
=
5^{-t/2}\bigl(\varphi^{2t-1}+\varphi^{t-2}\bigr).
\]
从而对任意 \(\alpha>0\)、\(\alpha\neq 1\)，有
\[
D_\alpha^\infty
:=
\lim_{m\to\infty}D_\alpha(\mu_m\Vert u_m)
=
\frac{\log M_L(\alpha)}{\alpha-1}.
\]
进一步，在 \(\mu_m\)-测度下亦有
\[
L_m\Longrightarrow L_\mu,
\]
其中 \(L_\mu\) 仍取值于 \(\{\ell_0,\ell_1\}\)，但概率变为
\[
\mathbf P(L_\mu=\ell_0)=\frac{\varphi}{\sqrt5},
\qquad
\mathbf P(L_\mu=\ell_1)=\frac{1}{\varphi\sqrt5}.
\]
特别地，
\[
\lim_{m\to\infty}\operatorname{Var}_{u_m}(L_m)
=
\frac{(\log\varphi)^2}{\varphi^3},
\qquad
\lim_{m\to\infty}\operatorname{Var}_{\mu_m}(L_m)
=
\frac{(\log\varphi)^2}{5}.
\]
\end{theorem}
\begin{proof}
由定理 \ref{thm:xi-time-part64b-fold-information-geometry-two-point-collapse} 的证明，
\[
\frac{\mu_m(w)}{u_m(w)}
=
\frac{\varphi^{2-w_m}}{\sqrt5}
+O\!\left(\left(\frac{\varphi}{2}\right)^m\right)
\qquad (w\in X_m)
\]
一致成立。取对数后得到
\[
L_m(w)=\ell_{w_m}+O\!\left(\left(\frac{\varphi}{2}\right)^m\right)
\qquad (w\in X_m),
\]
其中 \(\ell_0,\ell_1\) 如上定义。

另一方面，
\[
u_m(w_m=0)=\frac{F_{m+1}}{F_{m+2}}\longrightarrow \frac1\varphi,
\qquad
u_m(w_m=1)=\frac{F_m}{F_{m+2}}\longrightarrow \frac1{\varphi^2},
\]
故 \(u_m\)-下的极限分布即为所述 \(L_u\)。矩母函数由定义直接给出
\[
M_L(t)
=
\frac1\varphi e^{t\ell_0}+\frac1{\varphi^2}e^{t\ell_1}
=
\frac1\varphi\left(\frac{\varphi^2}{\sqrt5}\right)^t
+\frac1{\varphi^2}\left(\frac{\varphi}{\sqrt5}\right)^t,
\]
整理即得闭式
\[
M_L(t)=5^{-t/2}\bigl(\varphi^{2t-1}+\varphi^{t-2}\bigr).
\]
又由命题 \ref{prop:xi-foldbin-renyi-deficit-spectrum-closed-form}，
\[
\exp\!\bigl((\alpha-1)D_\alpha^\infty\bigr)
=
\frac1\varphi\left(\frac{\varphi^2}{\sqrt5}\right)^\alpha
+\frac1{\varphi^2}\left(\frac{\varphi}{\sqrt5}\right)^\alpha
=
M_L(\alpha),
\]
故得到 Rényi 公式。

对 \(\mu_m\)-测度，只需使用命题 \ref{prop:fold-bin-escort-lastbit} 在 \(q=1\) 处的末位极限：
\[
\mu_m(w_m=0)\longrightarrow \frac{\varphi}{\sqrt5},
\qquad
\mu_m(w_m=1)\longrightarrow \frac{1}{\varphi\sqrt5}.
\]
再与同一一致展开合并，便得 \(L_m\Rightarrow L_\mu\)。

最后，由
\[
\ell_0-\ell_1=\log\varphi,
\]
二点分布方差公式给出
\[
\operatorname{Var}(L_u)
=
\frac1\varphi\cdot\frac1{\varphi^2}(\log\varphi)^2
=
\frac{(\log\varphi)^2}{\varphi^3},
\]
\[
\operatorname{Var}(L_\mu)
=
\frac{\varphi}{\sqrt5}\cdot\frac{1}{\varphi\sqrt5}(\log\varphi)^2
=
\frac{(\log\varphi)^2}{5}.
\]
由分布收敛与一致有界性，这两者分别就是所述方差极限。证毕。
\end{proof}

\begin{theorem}[escort Fisher 密度是对数似然矩势的曲率且具有显式测地坐标]\label{thm:xi-time-part9xe-escort-fisher-likelihood-curvature}
沿用定理 \ref{thm:xi-time-part9xe-loglikelihood-two-point-cgf} 的记号。对 \(q\ge 0\) 定义
\[
\mathcal I(q):=\frac{\dd^2}{\dd q^2}\log M_L(q).
\]
则有显式公式
\[
\mathcal I(q)
=
(\log\varphi)^2\frac{\varphi^{q+1}}{(1+\varphi^{q+1})^2}.
\]
该量恰与 escort 极限族的 Fisher 密度一致；换言之，escort 温度线的极限 Fisher--Rao 度量可写为
\[
\mathrm ds^2=\mathcal I(q)\,\mathrm dq^2.
\]
更进一步，单位速度坐标可显式取为
\[
s(q):=2\arctan\!\bigl(\varphi^{(q+1)/2}\bigr),
\]
并满足
\[
s'(q)=\sqrt{\mathcal I(q)}.
\]
因此对任意 \(p,q\ge 0\)，escort 温度线的内在距离为
\[
d_{\mathrm{FR}}(p,q)
=
\left|
2\arctan\!\bigl(\varphi^{(q+1)/2}\bigr)
-
2\arctan\!\bigl(\varphi^{(p+1)/2}\bigr)
\right|.
\]
\end{theorem}
\begin{proof}
由定理 \ref{thm:xi-time-part9xe-loglikelihood-two-point-cgf}，
\[
\log M_L(q)
=
q\left(2\log\varphi-\frac12\log 5\right)
+\log\!\left(\frac1\varphi+\varphi^{-q-2}\right).
\]
仿射项的二阶导数为零，因此
\[
\mathcal I(q)
=
\frac{\dd^2}{\dd q^2}
\log\!\left(\frac1\varphi+\varphi^{-q-2}\right).
\]
直接微分得到
\[
\frac{\dd}{\dd q}
\log\!\left(\frac1\varphi+\varphi^{-q-2}\right)
=
-\frac{\log\varphi}{1+\varphi^{q+1}},
\]
再次微分便得
\[
\mathcal I(q)
=
(\log\varphi)^2\frac{\varphi^{q+1}}{(1+\varphi^{q+1})^2}.
\]
而定理 \ref{thm:xi-time-part9sb-escort-fisher-flat-coordinate} 已给出 escort 极限族的 Fisher 密度正是同一闭式，故
\[
\mathrm ds^2=\mathcal I(q)\,\mathrm dq^2.
\]
再由同一定理，
\[
s(q)=2\arctan\!\bigl(\varphi^{(q+1)/2}\bigr)
\]
满足
\[
s'(q)=\sqrt{\mathcal I(q)},
\]
距离公式随之成立。证毕。
\end{proof}

\begin{corollary}[escort 残差序参量的精确涨落--响应恒等式]\label{cor:xi-time-part9xe-escort-fluctuation-response}
沿用定理 \ref{thm:xi-time-part9ob-escort-quantile-phase-diagram} 的极限残差变量
\[
Z^{(q)}\in\{1,\varphi^{-1}\},
\qquad
\mathbf P\!\left(Z^{(q)}=\varphi^{-1}\right)=\theta(q),
\qquad
\theta(q)=\frac{1}{1+\varphi^{q+1}},
\]
并定义对数序参量
\[
X^{(q)}:=-\log Z^{(q)}\in\{0,\log\varphi\}.
\]
则
\[
\mathbf E\!\left[X^{(q)}\right]
=
\theta(q)\log\varphi
=
-\frac{\dd}{\dd q}
\left[
\log M_L(q)-q\left(2\log\varphi-\frac12\log 5\right)
\right],
\]
并且
\[
\operatorname{Var}\!\left(X^{(q)}\right)
=
\theta(q)\bigl(1-\theta(q)\bigr)(\log\varphi)^2
=
\mathcal I(q)
=
-\frac{\dd}{\dd q}\mathbf E\!\left[X^{(q)}\right].
\]
因此 escort 温度线上的 Fisher 密度恰等于缩放残差序参量的方差，同时也是其平均序参量对温度参数的负响应率。
\end{corollary}
\begin{proof}
由定理 \ref{thm:xi-time-part9ob-escort-quantile-phase-diagram}，
\[
X^{(q)}=
\begin{cases}
0,& \text{概率 }1-\theta(q),\\
\log\varphi,& \text{概率 }\theta(q).
\end{cases}
\]
故
\[
\mathbf E\!\left[X^{(q)}\right]=\theta(q)\log\varphi,
\qquad
\operatorname{Var}\!\left(X^{(q)}\right)
=
\theta(q)\bigl(1-\theta(q)\bigr)(\log\varphi)^2.
\]
另一方面，定理 \ref{thm:xi-time-part9xe-loglikelihood-two-point-cgf} 的公式给出
\[
\log M_L(q)-q\left(2\log\varphi-\frac12\log 5\right)
=
\log\!\left(\frac1\varphi+\varphi^{-q-2}\right),
\]
对其求导便得到
\[
-\frac{\dd}{\dd q}
\left[
\log M_L(q)-q\left(2\log\varphi-\frac12\log 5\right)
\right]
=
\frac{\log\varphi}{1+\varphi^{q+1}}
=
\theta(q)\log\varphi.
\]
再由定理 \ref{thm:xi-time-part9xe-escort-fisher-likelihood-curvature}，
\[
\mathcal I(q)
=
(\log\varphi)^2\theta(q)\bigl(1-\theta(q)\bigr),
\]
从而
\[
\operatorname{Var}\!\left(X^{(q)}\right)=\mathcal I(q).
\]
最后，由
\[
\theta'(q)=-(\log\varphi)\,\theta(q)\bigl(1-\theta(q)\bigr)
\]
可得
\[
\frac{\dd}{\dd q}\mathbf E\!\left[X^{(q)}\right]
=
\theta'(q)\log\varphi
=
-\mathcal I(q).
\]
证毕。
\end{proof}

\begin{theorem}[校正规范常数残差的 Bernoulli--\texorpdfstring{$\zeta$}{zeta} 提取范式]\label{thm:xi-time-part9xe-gauge-residual-even-zeta-extraction}
沿用定理 \ref{thm:fold-bin-gauge-constant-stirling-bernoulli-hierarchy} 的记号，设
\[
a_r:=
\frac{B_{2r}}{r(2r-1)}
\bigl(\varphi^{-1}+\varphi^{2r-3}\bigr)
\qquad (r\ge 1),
\]
并对每个 \(r\ge 1\) 定义残差重整化序列
\[
\mathcal E_r(m):=
\left(\frac{2}{\varphi}\right)^{(2r-1)m}
\left[
\log C_m-\log(2\pi)-\sum_{s=1}^{r-1}a_s
\left(\frac{\varphi}{2}\right)^{(2s-1)m}
\right].
\]
则极限存在并满足
\[
\lim_{m\to\infty}\mathcal E_r(m)=a_r.
\]
因而对每个 \(r\ge 1\) 都有
\[
\zeta(2r)
=
(-1)^{r-1}\frac{(2\pi)^{2r}}{2(2r)!}\,
\frac{r(2r-1)}{\varphi^{-1}+\varphi^{2r-3}}\,
\lim_{m\to\infty}\mathcal E_r(m).
\]
特别地，前两级提取公式为
\[
\lim_{m\to\infty}
\left(\frac{2}{\varphi}\right)^m
\bigl(\log C_m-\log(2\pi)\bigr)
=
\frac{1}{3\varphi},
\]
\[
\lim_{m\to\infty}
\left(\frac{2}{\varphi}\right)^{3m}
\left[
\log C_m-\log(2\pi)-\frac{1}{3\varphi}\left(\frac{\varphi}{2}\right)^m
\right]
=
-\frac{\sqrt5}{180}.
\]
\end{theorem}
\begin{proof}
定理 \ref{thm:xi-time-part9xb-gauge-constants-recover-even-zeta} 已直接给出：对每个 \(r\ge 1\)，上式中的极限存在且恰为 \(a_r\)。因此第一条立即成立。

再由 Euler 的 Bernoulli--\(\zeta\) 恒等式
\[
B_{2r}
=
(-1)^{r-1}\frac{2(2r)!}{(2\pi)^{2r}}\zeta(2r),
\]
将其代回 \(a_r\) 的定义并整理，便得到所示 \(\zeta(2r)\) 提取公式。

当 \(r=1\) 时，
\[
a_1=\frac{B_2}{1\cdot 1}\bigl(\varphi^{-1}+\varphi^{-1}\bigr)
=
\frac{1}{6}\cdot\frac{2}{\varphi}
=
\frac{1}{3\varphi}.
\]
当 \(r=2\) 时，
\[
a_2
=
\frac{B_4}{2\cdot 3}\bigl(\varphi^{-1}+\varphi\bigr)
=
-\frac{1}{180}\bigl(\varphi+\varphi^{-1}\bigr)
=
-\frac{\sqrt5}{180},
\]
其中使用了
\[
B_2=\frac16,
\qquad
B_4=-\frac1{30},
\qquad
\varphi+\varphi^{-1}=\sqrt5.
\]
证毕。
\end{proof}

\begin{theorem}[单一幂散度标量对极限信息几何的统一重构]\label{thm:xi-time-part9xe-single-power-divergence-unifies-limit-geometry}
对任意固定 \(\alpha>1\)，定义
\[
f_\alpha(t):=t^\alpha-1-\alpha(t-1),
\qquad
\mathcal D_\alpha:=\lim_{m\to\infty}D_{f_\alpha}(\mu_m\Vert u_m).
\]
则标量 \(\mathcal D_\alpha\) 唯一决定黄金参数 \(\varphi\)。因此，它进一步唯一决定：
\begin{enumerate}
  \item 全部 Rényi 极限
  \[
  D_\beta^\infty:=\lim_{m\to\infty}D_\beta(\mu_m\Vert u_m)
  \qquad (\beta>0,\ \beta\neq 1);
  \]
  \item 全部 Csisz\'ar 极限
  \[
  D_f^\infty:=\lim_{m\to\infty}D_f(\mu_m\Vert u_m)
  \qquad
  \bigl(f\ \text{凸且 }f(1)=0\bigr);
  \]
  \item 定理 \ref{thm:xi-time-part9xe-loglikelihood-two-point-cgf} 的两点对数似然律 \(L_u,L_\mu\) 及其完整矩母函数 \(M_L\)；
  \item 定理 \ref{thm:xi-time-part9xe-escort-fisher-likelihood-curvature} 与推论 \ref{cor:xi-time-part9xe-escort-fluctuation-response} 的 escort Fisher 几何与涨落--响应律。
\end{enumerate}
换言之，极限散度锥、对数似然二点律与 escort 信息几何并非彼此独立的对象，而是一条由 \(\varphi\) 单独参数化的一维刚性轨道。
\end{theorem}
\begin{proof}
由定理 \ref{thm:xi-time-part75a-single-power-divergence-determines-limit-tower}，对每个固定 \(\alpha>1\)，标量 \(\mathcal D_\alpha\) 对 \(\varphi\) 严格单调可辨识，故 \(\mathcal D_\alpha\) 唯一确定 \(\varphi\)。

一旦 \(\varphi\) 已知，命题 \ref{prop:xi-foldbin-renyi-deficit-spectrum-closed-form} 即给出全部 \(D_\beta^\infty\) 的闭式，而命题 \ref{prop:xi-foldbin-f-divergence-two-point-experiment} 则给出全部 \(D_f^\infty\) 的二点实验闭式。这证明了前两条。

第三条来自定理 \ref{thm:xi-time-part9xe-loglikelihood-two-point-cgf}：其中 \(L_u,L_\mu\) 的原子位置、权重以及 \(M_L(t)\) 都是 \(\varphi\) 的显式函数。第四条则来自定理 \ref{thm:xi-time-part9xe-escort-fisher-likelihood-curvature} 与推论 \ref{cor:xi-time-part9xe-escort-fluctuation-response}，因为
\[
\mathcal I(q)
=
(\log\varphi)^2\frac{\varphi^{q+1}}{(1+\varphi^{q+1})^2},
\qquad
s(q)=2\arctan\!\bigl(\varphi^{(q+1)/2}\bigr),
\qquad
\theta(q)=\frac{1}{1+\varphi^{q+1}},
\]
同样都是 \(\varphi\) 的显式函数。故所列全部极限信息几何对象都由同一 \(\mathcal D_\alpha\) 唯一锁定。证毕。
\end{proof}

\noindent
这组六条结果共同表明，bin-fold 主线中的各类极限对象虽分别以散度、似然、温度几何与规范常数残差的形式出现，但其真正的结构核心始终只有同一个末位二相机制。块均匀化模型在有限尺度上已经把全部 \(f\)-几何严格压缩到末位一比特；原始切片的对数似然比则把整条 Rényi 曲线收束为同一二点随机变量的矩势；escort 温度线的 Fisher 密度又与该矩势的曲率以及残差序参量的方差完全一致；最后，fold-gauge 常数的 Bernoulli 层级则把同一刚性进一步传输到偶整数 \(\zeta\) 值的逐阶提取。于是，极限散度锥、escort 温度几何与规范常数残差不再是平行现象，而是同一黄金参数轨道在不同读出口径下的等价显现。

\endinput

\paragraph{对称幂矩谱、尾谱完备统计、冻结双坐标与 \texorpdfstring{$S_4$}{S4} Jacobian 分层}
在附录射影压力算子的整数退化、固定分辨率上的容量--尾谱反演、冻结相的双指数控制、金属平均族的严格 gap，以及属 \(49\) 的 \(S_4\)-Hurwitz 商塔的 Chevalley--Weil/Jacobian/Prym 分解都已建立之后，还可把这些分散接口进一步压缩为下列六条综合色结论。它们分别把整数矩谱塔识别为一次层有限核在对称代数上的分次影像，把固定窗口上的全部纯标量统计压缩为唯一尾谱，把冻结相的宏观热力学收束为两枚指数坐标，把黄金分支锁定为常数型金属平均族中唯一的对数成本极小点，并把属 \(49\) 的 \(S_4\)-全对称中性层消失与 Jacobian 的平方--立方分层一并固定。

\begin{theorem}[整数矩谱塔的对称幂封闭与一次层共轭不变性]\label{thm:xi-time-part9xf-moment-tower-symmetric-power-closure}
沿用定义 \ref{def:pom-projective-transfer-operator}、定义 \ref{def:pom-homogeneous-polynomials-simplex}、注 \ref{rem:pom-moment-symmetric-power-substitution-operator} 与定理 \ref{thm:pom-projective-operator-degenerates-to-moment-kernel} 的记号，并记
\[
\mathcal H:=\bigoplus_{q\ge 0}\mathcal H_q,
\qquad
r_q:=\rho(\mathcal L_q|_{\mathcal H_q})=\rho(\widetilde A_q)
\qquad (q\ge 2).
\]
则对每个整数 \(q\ge 2\)，\(\mathcal L_q|_{\mathcal H_q}\) 正是由一次层有限核 \(\{B_b\}_{b\in A_{\mathrm{out}}}\) 在对称代数上的 \(q\) 次齐次分量所诱导出的作用，而 \(\widetilde A_q\) 恰为该作用在单项式基下的矩阵表示。因而整数矩谱塔
\[
\{r_q\}_{q\ge 2}
\]
并非彼此无关的标量族，而是同一个一次层线性核在对称代数分次表示上的谱半径序列。

进一步，若两套确定性在线核 \(\mathcal T,\mathcal T'\) 满足：存在某个线性同构 \(C\) 使得
\[
B_b(\mathcal T')=C^{-1}B_b(\mathcal T)C
\qquad (b\in A_{\mathrm{out}})
\]
同时成立，则对每个整数 \(q\ge 2\) 都有
\[
\mathcal L_q(\mathcal T')|_{\mathcal H_q}
\cong
\mathcal L_q(\mathcal T)|_{\mathcal H_q},
\qquad
\widetilde A_q(\mathcal T')
\cong
\widetilde A_q(\mathcal T),
\]
从而
\[
r_q(\mathcal T')=r_q(\mathcal T).
\]
换言之，一次层上的公共共轭会通过对称幂函子自动传递到全部整数矩层。
\end{theorem}
\begin{proof}
由定理 \ref{thm:pom-projective-operator-degenerates-to-moment-kernel}，对每个整数 \(q\ge 2\)，\(\mathcal L_q\) 在 \(\mathcal H_q\) 上的限制在任一单项式基下的矩阵表示恰为 \(\widetilde A_q\)，并且
\[
\rho(\mathcal L_q|_{\mathcal H_q})=\rho(\widetilde A_q)=r_q.
\]
另一方面，注 \ref{rem:pom-moment-symmetric-power-substitution-operator} 已指出：该有限维作用来自一次层线性替换 \(x\mapsto B_b^{\mathsf T}x\) 在对称代数上的 \(q\) 次齐次分量。于是 \(\widetilde A_q\) 与 \(\mathcal L_q|_{\mathcal H_q}\) 不再是独立构造，而是同一个一次层有限核在 \(\mathcal H\) 上分次表示的 \(q\) 级坐标实现。这就证明了第一部分。

若再有
\[
B_b(\mathcal T')=C^{-1}B_b(\mathcal T)C
\qquad (b\in A_{\mathrm{out}}),
\]
则一次层线性替换由公共共轭 \(C\) 联系起来。对称幂表示是函子性的，故在每个 \(q\) 次齐次层上都会自动诱导一个由 \(C\) 决定的线性同构，把 \(\mathcal T\) 与 \(\mathcal T'\) 的 \(q\) 级作用互相共轭。于是
\[
\mathcal L_q(\mathcal T')|_{\mathcal H_q}
\cong
\mathcal L_q(\mathcal T)|_{\mathcal H_q},
\]
而在单项式基下这正意味着
\[
\widetilde A_q(\mathcal T')
\cong
\widetilde A_q(\mathcal T).
\]
谱半径作为共轭不变量因而相等，即
\[
r_q(\mathcal T')=r_q(\mathcal T).
\]
证毕。
\end{proof}

\begin{theorem}[固定分辨率纤维尾谱的完备统计原则]\label{thm:xi-time-part9xf-tail-spectrum-complete-statistic}
固定分辨率 \(m\)。定义纤维尾函数、连续容量曲线与全矩谱
\[
N_m(t):=\#\{x\in X_m:\ d_m(x)\ge t\},
\qquad
\mathcal C_m^{\mathrm{cont}}(T):=\sum_{x\in X_m}\min\bigl(d_m(x),T\bigr),
\]
\[
S_q(m):=\sum_{x\in X_m}d_m(x)^q
\qquad (q>0).
\]
则
\[
N_m
\longleftrightarrow
\mathcal C_m^{\mathrm{cont}}
\longleftrightarrow
\{S_q(m)\}_{q>0}
\]
两两唯一决定。更准确地说，这三类数据都等价于同一个纤维直方图 \(h_m(d)=\#\{x:d_m(x)=d\}\)。

因而，在固定窗口层级上，任意仅依赖于多重度多重集
\[
\{d_m(x):x\in X_m\}
\]
的纯标量统计都不再携带超出 \(N_m\) 之外的独立不变量；特别地，离散 oracle 容量
\[
C_m(B)=\mathcal C_m^{\mathrm{cont}}(2^B),
\]
全部碰撞矩 \(S_q(m)\)，以及由同一多重度多重集定义的任意对称标量泛函，都已被 \(N_m\) 完备编码。
\end{theorem}
\begin{proof}
定理 \ref{thm:xi-capacity-tail-histogram-moment-complete-statistic} 已经给出：固定 \(m\) 时，连续容量曲线 \(\mathcal C_m^{\mathrm{cont}}\)、尾计数函数 \(N_m\)、纤维直方图 \(h_m\) 与全矩谱 \(\{S_q(m)\}_{q>0}\) 彼此唯一决定。定理 \ref{thm:xi-time-part62df-capacity-curve-complete-inversion} 又把这一点写成完全可逆的显式反演公式；特别地，
\[
N_m(t)=\partial_T^+\mathcal C_m^{\mathrm{cont}}(t)\quad\text{a.e.},
\qquad
h_m(d)=N_m(d)-N_m(d+1),
\]
并且
\[
S_q(m)=q\int_0^\infty t^{q-1}N_m(t)\,\dd t
\qquad (q>0).
\]
因此
\[
N_m
\longleftrightarrow
\mathcal C_m^{\mathrm{cont}}
\longleftrightarrow
h_m
\longleftrightarrow
\{S_q(m)\}_{q>0}
\]
完全等价。

另一方面，定理 \ref{thm:xi-time-part62df-capacity-curve-complete-inversion} 的第三条已说明：对任意函数 \(F:\ZZ_{\ge 1}\to\RR\)，都有
\[
\sum_{x\in X_m}F\bigl(d_m(x)\bigr)
=
\sum_{d\ge 1}h_m(d)F(d),
\]
故任意仅依赖于多重度多重集的对称标量泛函都由 \(h_m\) 唯一恢复，亦即由 \(N_m\) 唯一恢复。再注意到离散 oracle 容量只是连续容量在 \(T=2^B\) 处的取值，
\[
C_m(B)=\mathcal C_m^{\mathrm{cont}}(2^B),
\]
而碰撞矩就是 \(S_q(m)\) 本身，于是所述完备编码原则随之成立。证毕。
\leanverified{paper\_xi\_time\_part9xf\_tail\_spectrum\_complete\_statistic}
\end{proof}

\begin{theorem}[冻结相宏观热力学的双指数坐标化]\label{thm:xi-time-part9xf-frozen-phase-two-coordinate-thermodynamics}
沿用定理 \ref{thm:fold-pressure-freezing-threshold}、定理 \ref{thm:xi-fixed-freezing-escort-renyi-spectrum-collapse} 与定理 \ref{thm:xi-fixed-freezing-pressure-address-cost-residual-decomposition} 的记号。设
\[
M_m:=\max_{x\in X_m}d_m(x),
\qquad
A_m^{\max}:=\{x\in X_m:\ d_m(x)=M_m\},
\qquad
K_m:=\#A_m^{\max},
\]
并记次大纤维规模为 \(M_m^{(2)}\)。定义
\[
\alpha_*:=\lim_{m\to\infty}\frac1m\log M_m,
\qquad
g_*:=\lim_{m\to\infty}\frac1m\log K_m,
\qquad
\alpha_*^{(2)}:=\limsup_{m\to\infty}\frac1m\log M_m^{(2)}.
\]
若存在严格基态间隙
\[
\alpha_*^{(2)}<\alpha_*,
\]
并定义冻结阈值
\[
a_{\mathrm c}:=\frac{\log\varphi-g_*}{\alpha_*-\alpha_*^{(2)}},
\]
则对每个实数 \(a>a_{\mathrm c}\)，压力极限
\[
P_a:=\lim_{m\to\infty}\frac1m\log S_a(m),
\qquad
\pi_{m,a}(x):=\frac{d_m(x)^a}{S_a(m)},
\qquad
h_\infty(a):=\lim_{m\to\infty}\frac1m\Bigl(-\log\max_x\pi_{m,a}(x)\Bigr)
\]
满足三条同步刚性律：
\[
P_a=a\alpha_*+g_*,
\]
\[
\pi_{m,a}(A_m^{\max})
\ge
1-\exp\!\bigl(-m\Delta(a)+o(m)\bigr),
\qquad
\Delta(a):=a(\alpha_*-\alpha_*^{(2)})-(\log\varphi-g_*),
\]
\[
h_\infty(a)=g_*.
\]
因此，冻结相一旦越过阈值 \(a_{\mathrm c}\)，其完整宏观热力学便精确收束为两枚指数坐标 \((\alpha_*,g_*)\)：前者记录最大纤维的能量增长，后者记录基态简并的指数增长。
\end{theorem}
\begin{proof}
由定理 \ref{thm:fold-pressure-freezing-threshold}，对每个实数 \(a>a_{\mathrm c}\) 都有
\[
P_a=a\alpha_*+g_*.
\]
这给出第一条。

再由定理 \ref{thm:xi-fixed-freezing-escort-renyi-spectrum-collapse} 的第四条，冻结区 escort 分布对基态均匀分布
\[
U_m^{\max}:=\mathrm{Unif}(A_m^{\max})
\]
满足指数型总变差塌缩
\[
\|\pi_{m,a}-U_m^{\max}\|_{\mathrm{TV}}
\le
\exp\!\bigl(-m\Delta(a)+o(m)\bigr),
\]
其中
\[
\Delta(a)=a(\alpha_*-\alpha_*^{(2)})-(\log\varphi-g_*).
\]
由于 \(U_m^{\max}(A_m^{\max})=1\)，故
\[
1-\pi_{m,a}(A_m^{\max})
\le
\|\pi_{m,a}-U_m^{\max}\|_{\mathrm{TV}},
\]
于是得到第二条。

最后，定理 \ref{thm:xi-fixed-freezing-pressure-address-cost-residual-decomposition} 已给出冻结相的 min-entropy 率满足
\[
h_\infty(a)=g_*.
\]
这就证明了第三条。

三式联立可见：在整个冻结区 \(a>a_{\mathrm c}\) 上，压力只剩斜率 \(\alpha_*\) 与截距 \(g_*\)，escort 质量几乎全部集中到基态集合 \(A_m^{\max}\)，而最大点概率的指数率也恰由 \(g_*\) 控制。故冻结相的宏观描述已被 \((\alpha_*,g_*)\) 完全锁定。证毕。
\end{proof}

\begin{theorem}[黄金分支在金属平均族中的对数成本极小性]\label{thm:xi-time-part9xf-golden-metallic-logcost-minimizer}
沿用定义 \ref{def:const-type} 与定理 \ref{thm:metallic-gap} 的记号。对每个实数 \(A\ge 1\)，记
\[
\alpha_A=[0;A,A,A,\dots],
\qquad
\lambda_A:=\frac{A+\sqrt{A^2+4}}{2},
\]
并定义对数成本系数
\[
\kappa(A):=\frac{A}{\log\lambda_A}.
\]
则函数 \(A\mapsto \kappa(A)\) 在区间 \([1,\infty)\) 上严格递增。特别地，对一切整数 \(A\ge 2\)，都有
\[
\kappa(A)>\kappa(1)=\frac{1}{\log\varphi}.
\]
因而在常数型金属平均族中，黄金端点 \(A=1\) 是唯一的对数成本极小点；等价地，它也是单位符号预算下指数效率的唯一极大点。
\end{theorem}
\begin{proof}
定理 \ref{thm:metallic-gap} 已经证明
\[
\kappa(A)=\frac{A}{\log\lambda_A}
\]
在 \([1,\infty)\) 上严格递增，故其最小值只能在左端点 \(A=1\) 取得。又由
\[
\lambda_1=\varphi
\]
可得
\[
\kappa(1)=\frac{1}{\log\varphi}.
\]
因此对每个整数 \(A\ge 2\) 都有
\[
\kappa(A)>\kappa(1)=\frac{1}{\log\varphi}.
\]

另一方面，定理 \ref{thm:xi-golden-metallic-symbol-budget-entropy-optimality} 已将其倒数解释为单位符号预算下的分母熵率
\[
\eta(A)=\kappa(A)^{-1}=\frac{\log\lambda_A}{A}.
\]
故 \(\kappa(A)\) 的唯一极小性与 \(\eta(A)\) 的唯一极大性是完全等价的两种表述。证毕。
\end{proof}

\begin{theorem}[属 \texorpdfstring{$49$}{49} 的 \texorpdfstring{$S_4$}{S4}--Hurwitz 微分表示没有平凡分量]\label{thm:xi-time-part9xf-s4-neutral-layer-vanishing}
设 \(X\) 为正文中的属 \(49\) 的 \(S_4\)-Hurwitz 曲线，并记
\[
V:=H^0(X,\Omega_X^1).
\]
则
\[
V^{S_4}=0.
\]
等价地，\(V\) 的 Chevalley--Weil 分解中平凡表示不出现。因而
\[
\bigl(T_0J(X)\bigr)^{S_4}=0,
\]
从而 \(J(X)\) 不存在任何正维的 \(S_4\)-中性阿贝尔商；特别地，全对称层上不存在非平凡的一维环面相位方向。
\end{theorem}
\begin{proof}
由结论 \ref{con:xi-time-part9n-s4-all-subgroup-genera}，属 \(49\) 的 \(S_4\)-Hurwitz 曲线满足
\[
V
\cong
5\cdot \mathrm{sgn}\oplus 4\cdot V_2\oplus 3\cdot V_3\oplus 9\cdot V_3',
\]
其中平凡表示 \(\mathbf 1\) 不出现。故立即有
\[
V^{S_4}=0.
\]
又因为 \(T_0J(X)\cong V^\vee\)，于是
\[
\bigl(T_0J(X)\bigr)^{S_4}=0.
\]

若反设存在一个正维的 \(S_4\)-中性阿贝尔商
\[
f:J(X)\to A,
\]
即 \(S_4\) 在 \(A\) 上作用为平凡，则拉回映射给出一个单射
\[
f^*:H^0(A,\Omega_A^1)\hookrightarrow H^0(X,\Omega_X^1)=V.
\]
由于 \(A\) 上的 \(S_4\)-作用平凡，其像必包含在 \(V^{S_4}\) 中；但 \(V^{S_4}=0\)，这与 \(A\) 正维矛盾。故不存在任何正维的 \(S_4\)-中性阿贝尔商。证毕。
\end{proof}

\begin{theorem}[属 \texorpdfstring{$49$}{49} Jacobian 的平方--立方分层与三次 Prym 平方块]\label{thm:xi-time-part9xf-genus49-jacobian-square-cube-stratification}
沿用定理 \ref{thm:killo-s4-genus49-jacobian-computable-isogeny}、定理 \ref{thm:xi-time-part9n1b-prym-chevalley-weil-reciprocity} 与结论 \ref{con:xi-time-part9n-prym-c6h-weil-type-signature} 的记号，记
\[
H:=X/A_4,\qquad
Z:=X/D_8,\qquad
C_F:=X/S_3,\qquad
Y:=X/C_4,
\]
\[
P:=\operatorname{Prym}(Y/Z),
\qquad
A:=\operatorname{Prym}(C_6/H).
\]
则有精确同源分解
\[
J(X)\sim_{\QQ}J(H)\times J(Z)^2\times J(C_F)^3\times P^3,
\]
并且
\[
g(H)=5,
\qquad
g(Z)=4,
\qquad
g(C_F)=3,
\qquad
\dim P=9.
\]
此外，
\[
A\sim_{\QQ}J(Z)^2,
\qquad
\mathrm{type}(\lambda_A)=(1^4,3^4),
\qquad
\QQ(\zeta_3)\subset \operatorname{End}^0(A).
\]
因此，属 \(49\) Jacobian 的连续阿贝尔几何严格裂解为一个五维基块、一个八维平方块与两个立方块；其中无分歧三重 Prym 平方块 \(A\) 具有非主极化并携带三次 cyclotomic 乘法。
\end{theorem}
\begin{proof}
定理 \ref{thm:killo-s4-genus49-jacobian-computable-isogeny} 已给出
\[
J(X)\sim_{\QQ}J(H)\times J(Z)^2\times J(C_F)^3\times P^3.
\]
这就是所述平方--立方分层的同源形式。

再由定理 \ref{thm:xi-time-part9n1b-prym-chevalley-weil-reciprocity}，
\[
V\cong
g(H)\cdot \mathrm{sgn}\oplus
g(Z)\cdot V_2\oplus
g(C_F)\cdot V_3\oplus
\dim(P)\cdot V_3'.
\]
而结论 \ref{con:xi-time-part9n-s4-all-subgroup-genera} 已给出
\[
g(H)=5,
\qquad
g(Z)=4,
\qquad
g(C_F)=3.
\]
将这些数值代回前式，便得到
\[
\dim P=9.
\]

最后，结论 \ref{con:xi-time-part9n-prym-c6h-weil-type-signature} 直接给出
\[
A=\operatorname{Prym}(C_6/H)\sim_{\QQ}J(Z)^2,
\]
并且
\[
\mathrm{type}(\lambda_A)=(1^4,3^4),
\qquad
\QQ(\zeta_3)\subset \operatorname{End}^0(A).
\]
于是 \(A\) 的确是一个八维平方块，同时带有非主极化与三次 cyclotomic 乘法。综合这些事实，\(J(X)\) 的连续阿贝尔几何分层不再只是维数分拆，而是被同源分解、极化型与端同态代数同时强制出来的刚性结构。证毕。
\end{proof}

\noindent
这六条结果共同表明，附录射影压力的整数 moment-kernel、固定分辨率上的纤维尾谱、冻结端面的双指数坐标、金属平均族的端点极值，以及属 \(49\) 的 \(S_4\)-Hurwitz Jacobian/Prym 分层，并非彼此分离的结构片段。整数矩谱塔只是一次层有限核在对称代数上的分次影像；固定窗口上的全部纯标量统计已被唯一尾谱完备编码；冻结区宏观热力学仅保留 \((\alpha_*,g_*)\) 两枚指数坐标；常数型金属平均族的对数成本在黄金端点取得唯一极小；而 \(S_4\)-全对称层的平凡方向完全消失，使得连续 Abelian 自由度只能分配到商曲线与 Prym 因子的平方块、立方块之中。

\endinput

\paragraph{金属平均族标量化价值函数的端点二阶相变}
前述链条已经把金属平均族上的两条竞争目标
\[
h_Z(A)=\log\!\frac{A+1}{\lambda_A},
\qquad
\kappa(A)=\frac{A}{\log\lambda_A},
\qquad
\lambda_A=\frac{A+\sqrt{A^2+4}}{2}
\]
分别压缩为压缩率极值、帕累托前沿与线性标量化阈值。
再把它们进一步收束到价值函数层面时，可得到一条严格的端点相变律：最优尺度在 \(\beta_c\) 处从解析内点相切换到黄金端点线性相，而价值函数本身则保持 \(C^1\) 连续但失去二阶光滑性。

\begin{theorem}[金属平均族标量化价值函数的 \texorpdfstring{$C^1$}{C1} 非 \texorpdfstring{$C^2$}{C2} 相变]\label{thm:xi-time-part9xfa-metallic-value-c1-nonc2-transition}
\leanverified{paper\_xi\_time\_part9xfa\_metallic\_value\_c1\_nonc2\_transition}
沿用命题 \ref{prop:metallic-linear-scalarization-threshold} 与推论
\ref{cor:metallic-pareto-scale-law} 的记号。
对每个 \(\beta\ge 0\) 定义
\[
J_\beta(A):=h_Z(A)-\beta\,\kappa(A),
\qquad
V(\beta):=\sup_{A\ge 1}J_\beta(A),
\]
并记 \(A_\beta\) 为 \(J_\beta\) 在 \(A\ge 1\) 上的唯一最大化点。
再记
\[
\beta_c=
\frac{\left(\frac12-\frac1{\sqrt5}\right)(\log\varphi)^2}
{\log\varphi-\frac1{\sqrt5}}.
\]
则有：
\begin{enumerate}
  \item 对每个 \(\beta\ge 0\)，最优尺度 \(A_\beta\) 唯一存在，且
  \[
  A_\beta\in\left[1,\frac32\right].
  \]
  \item 当 \(0\le \beta<\beta_c\) 时，
  \[
  A_\beta\in\left(1,\frac32\right]
  \]
  为唯一内点解，并且映射 \(\beta\mapsto A_\beta\) 在 \([0,\beta_c)\) 上实解析且严格递减。
  \item 当 \(\beta\ge \beta_c\) 时，
  \[
  A_\beta=1,
  \qquad
  V(\beta)=h_Z(1)-\beta\,\kappa(1)
  =\log\!\frac{2}{\varphi}-\frac{\beta}{\log\varphi},
  \]
  故价值函数进入严格线性相。
  \item \(V\) 在 \([0,\infty)\) 上属于 \(C^1\)，但在 \(\beta=\beta_c\) 处不属于 \(C^2\)。
  更准确地，
  \[
  \lim_{\beta\uparrow \beta_c}V''(\beta)
  =
  -\frac{\kappa'(1)^2}{h_Z''(1)-\beta_c\kappa''(1)}
  >0,
  \qquad
  \lim_{\beta\downarrow \beta_c}V''(\beta)=0.
  \]
  \item 特别地，
  \[
  V(0)=h_Z\!\left(\frac32\right)=\log\!\frac54.
  \]
\end{enumerate}
\end{theorem}
\begin{proof}
命题 \ref{prop:metallic-linear-scalarization-threshold} 与推论
\ref{cor:metallic-pareto-scale-law} 已经给出：对每个 \(\beta\ge 0\)，最大化点 \(A_\beta\) 唯一存在并落在 \([1,3/2]\)；当 \(0\le \beta<\beta_c\) 时，\(A_\beta\in(1,3/2]\) 为唯一内点解；当 \(\beta\ge \beta_c\) 时，端点 \(A_\beta=1\) 被强制选中。
这就得到 (1) 与 (3) 的第一部分。

对 \(0\le \beta<\beta_c\)，定义
\[
F(A,\beta):=h_Z'(A)-\beta\,\kappa'(A).
\]
由推论 \ref{cor:metallic-pareto-scale-law} 的证明，在解点处有
\[
\partial_A F(A_\beta,\beta)=h_Z''(A_\beta)-\beta\,\kappa''(A_\beta)<0.
\]
由于 \(h_Z\) 与 \(\kappa\) 在 \(A>1\) 上实解析，隐函数定理推出 \(\beta\mapsto A_\beta\) 在 \([0,\beta_c)\) 上实解析。
同时
\[
A'_\beta
=
\frac{\kappa'(A_\beta)}{h_Z''(A_\beta)-\beta\,\kappa''(A_\beta)}
<0,
\]
故 \(A_\beta\) 严格递减。
这就证明了 (2)。

当 \(\beta\ge \beta_c\) 时，由 \(A_\beta=1\) 立得
\[
V(\beta)=J_\beta(1)=h_Z(1)-\beta\,\kappa(1).
\]
又因为
\[
h_Z(1)=\log\!\frac{2}{\varphi},
\qquad
\kappa(1)=\frac{1}{\log\varphi},
\]
遂得到 (3) 的线性公式。

再证光滑性。
在内点相 \(0\le \beta<\beta_c\) 上，由链式法则和一阶条件 \(J_\beta'(A_\beta)=0\)，
\[
V'(\beta)
=
\frac{\partial}{\partial\beta}J_\beta(A_\beta)
+J_\beta'(A_\beta)A'_\beta
=
-\kappa(A_\beta).
\]
因此
\[
V''(\beta)=-\kappa'(A_\beta)A'_\beta
=
-\frac{\kappa'(A_\beta)^2}{h_Z''(A_\beta)-\beta\,\kappa''(A_\beta)}
>0.
\]
当 \(\beta\ge \beta_c\) 时，\(V(\beta)\) 为线性函数，故
\[
V'(\beta)=-\kappa(1),
\qquad
V''(\beta)=0.
\]

又由推论 \ref{cor:metallic-pareto-scale-law}，当 \(\beta\uparrow\beta_c\) 时 \(A_\beta\to 1\)，于是
\[
\lim_{\beta\uparrow\beta_c}V'(\beta)
=
-\kappa(1)
=
\lim_{\beta\downarrow\beta_c}V'(\beta).
\]
故 \(V\in C^1([0,\infty))\)。

另一方面，由上式与 \(A_\beta\to 1\) 的连续性，
\[
\lim_{\beta\uparrow\beta_c}V''(\beta)
=
-\frac{\kappa'(1)^2}{h_Z''(1)-\beta_c\kappa''(1)}.
\]
而命题 \ref{prop:metallic-linear-scalarization-threshold} 的证明已给出
\[
h_Z''(A)-\beta\,\kappa''(A)<0
\qquad
\left(1\le A\le \frac32,\ \beta\ge 0\right),
\]
故上述左极限严格为正；而右侧二阶导恒等于 \(0\)。
因此 \(V\) 在 \(\beta_c\) 处不属于 \(C^2\)。
这就证明了 (4)。

最后，当 \(\beta=0\) 时，
\[
V(0)=\sup_{A\ge 1}h_Z(A).
\]
由命题 \ref{prop:metallic-compression-extremum}，唯一极大点为 \(A=3/2\)，并且
\[
h_Z\!\left(\frac32\right)=\log\!\frac54.
\]
故得到 (5)。证毕。
\end{proof}

\endinput

\paragraph{平方自由证书、二次符号层与连续压力补桥的有限盲核}
在链内算子格的平方自由算术化、Smith 模核大小的逐素数可恢复性、双域 Chebotarev 指纹的强乘法律，以及附录 \(E\) 的实参数射影压力插值都已建立之后，还可把这三条彼此分离的接口进一步压缩为一组有限盲核命题。它们分别表明：平方自由 gcd/lcm 证书无法恢复任何 multiplicity-sensitive 的 Smith 指数高度；\(S_{10}\) 的纯二次符号层对 \(S_3\) 的全部局部分裂证书保持严格中性；而有限多个整数 \(q\)-层 moment-kernel 样本也不足以闭合连续实参数压力曲线本身。

\begin{theorem}[平方自由 gcd/lcm 证书对 Smith 指数谱的严格盲性]\label{thm:xi-time-part9xg-squarefree-certificate-blindness-to-smith-heights}
把定理 \ref{thm:xi-time-part65d-chain-interior-boolean-arithmeticization} 所给出的 \(\kappa\)-像视为平方自由 gcd/lcm 证书模型。设 \(A\in M_{m\times n}(\ZZ)\) 满行秩，其 Smith 不变量为
\[
d_1\mid \cdots \mid d_m,
\qquad
e_i(p):=v_p(d_i).
\]
定义其平方自由支撑为
\[
\operatorname{sqf}(A):=\prod_{p:\,\max_i e_i(p)>0}p.
\]
若某个证书赋值 \(\mathfrak C\) 仅通过平方自由支撑因子化，即存在映射 \(F\) 使
\[
\mathfrak C(A)=F\!\bigl(\operatorname{sqf}(A)\bigr),
\]
并且其可用闭包运算只由 gcd/lcm 驱动，则 \(\mathfrak C\) 不可能恢复完整的 Smith 指数谱 \(\{e_i(p)\}_{p,i}\)。更强地，存在两个满秩整数矩阵 \(A^{(1)},A^{(2)}\) 使得
\[
\mathfrak C(A^{(1)})=\mathfrak C(A^{(2)}),
\]
但它们的局部模核大小函数
\[
K_{A,p}(k):=\#\ker(A\bmod p^k)
\]
并不相同。
\leanverified{paper\_xi\_time\_part9xg\_squarefree\_certificate\_blindness\_to\_smith\_heights}
\end{theorem}
\begin{proof}
由定理 \ref{thm:xi-time-part65d-chain-interior-boolean-arithmeticization}，\(\kappa\)-像严格落在平方自由整除格内，并满足
\[
\kappa(I\wedge J)=\gcd(\kappa(I),\kappa(J)),
\qquad
\kappa(I\vee J)=\mathrm{lcm}(\kappa(I),\kappa(J)).
\]
因此，任何经由同类平方自由 gcd/lcm 证书模型因子化的赋值都只能记录“哪些素数出现”，而无法记录“每个素数以何种高度出现”。

现取最简单的两个 \(1\times 1\) 满秩矩阵
\[
A^{(1)}=[p],
\qquad
A^{(2)}=[p^2].
\]
它们满足
\[
\operatorname{sqf}(A^{(1)})=\operatorname{sqf}(A^{(2)})=p,
\]
故一切仅依赖平方自由支撑的证书都无法区分二者，即
\[
\mathfrak C(A^{(1)})=\mathfrak C(A^{(2)}).
\]

然而由推论 \ref{cor:killo-kernel-size-general-modulus}，
\[
\#\ker(A_b)=b^{n-m}\prod_{i=1}^{m}\gcd(d_i,b).
\]
对上述两个矩阵，\(m=n=1\)。取 \(b=p^2\)，便得到
\[
K_{A^{(1)},p}(2)=\#\ker(A^{(1)}\bmod p^2)=\gcd(p,p^2)=p,
\]
\[
K_{A^{(2)},p}(2)=\#\ker(A^{(2)}\bmod p^2)=\gcd(p^2,p^2)=p^2.
\]
故它们的局部模核大小函数不同。于是完整的 Smith 指数谱不可能由此类平方自由证书恢复。证毕。
\end{proof}

\begin{corollary}[平方自由 gcd/lcm 正规化不能闭合 multiplicity-sensitive 亏损律]\label{cor:xi-time-part9xg-squarefree-normalization-cannot-close-loss-law}
\leanverified{paper\_xi\_time\_part9xg\_squarefree\_normalization\_cannot\_close\_loss\_law}
对任意满行秩整数矩阵 \(A\in M_{m\times n}(\ZZ)\)，若把
\[
b=\prod_p p^{k_p}
\]
分解为素数幂，则其模核大小满足
\[
\log \#\ker(A_b)
=
(n-m)\log b
+\sum_p\left(\sum_{i=1}^{m}\min\bigl(k_p,e_i(p)\bigr)\right)\log p.
\]
因此，任何把 Smith 数据正规化到平方自由 gcd/lcm 证书代数的方案，都不可能恢复完整的 multiplicity-sensitive 信息亏损律。
\end{corollary}
\begin{proof}
推论 \ref{cor:killo-kernel-size-general-modulus} 给出
\[
\#\ker(A_b)=b^{n-m}\prod_{i=1}^{m}\gcd(d_i,b).
\]
把 \(b=\prod_p p^{k_p}\) 与 \(d_i=\prod_p p^{e_i(p)}\) 的素因子分解代入，可得
\[
\gcd(d_i,b)=\prod_p p^{\min(k_p,e_i(p))},
\]
从而
\[
\#\ker(A_b)
=
b^{n-m}\prod_p p^{\sum_{i=1}^{m}\min(k_p,e_i(p))}.
\]
两边取对数便得到所示公式。

另一方面，定理 \ref{thm:xi-time-part9xg-squarefree-certificate-blindness-to-smith-heights} 已表明：平方自由 gcd/lcm 正规化会完全丢失各素数方向的高度数据 \(e_i(p)\)。而上式恰恰通过
\[
\min\bigl(k_p,e_i(p)\bigr)
\]
逐素数叠加这些高度。因此该正规化不可能恢复完整亏损律。证毕。
\end{proof}

\begin{theorem}[\texorpdfstring{$S_{10}$}{S10} 二次符号层对 \texorpdfstring{$S_3$}{S3} 分裂证书的完全中性]\label{thm:xi-time-part9xg-s10-parity-neutrality-for-s3-splitting}
\leanverified{paper\_xi\_time\_part9xg\_s10\_parity\_neutrality\_for\_s3\_splitting}
仍在共同不分歧素数的密度概率空间上，设 \(D\subset S_3\) 为任意共轭类并，并记
\[
E_D:=\{\,p:\ \sigma_3(p)\in D\,\},
\qquad
\chi_{10}(p):=\operatorname{sgn}\!\bigl(\sigma_{10}(p)\bigr)\in\{\pm1\}.
\]
则对任意 \(\varepsilon\in\{\pm1\}\)，都有
\[
\delta\!\bigl(E_D\mid \chi_{10}=\varepsilon\bigr)=\delta(E_D)=\frac{|D|}{6}.
\]
换言之，\(S_{10}\) 侧的纯二次符号条件对 \(S_3\) 侧的任意局部分裂证书都不产生信息增益。
\end{theorem}
\begin{proof}
在定理 \ref{thm:xi-time-part62ea-dualfingerprint-boolean-strong-product-law} 中，取
\[
E:=\{\chi_{10}=\varepsilon\}\in \mathcal A_{10},
\qquad
F:=E_D\in \mathcal A_3.
\]
该强乘法律直接给出
\[
\delta(E\cap F)=\delta(E)\,\delta(F).
\]
又因为 \(\delta(E)=\delta(\chi_{10}=\varepsilon)=1/2>0\)，故
\[
\delta\!\bigl(E_D\mid \chi_{10}=\varepsilon\bigr)
=
\frac{\delta(E\cap F)}{\delta(E)}
=
\delta(F)
=
\delta(E_D).
\]
最后，由 \(S_3\)-侧的 Chebotarev 密度公式，任意共轭类并 \(D\subset S_3\) 的密度为
\[
\delta(E_D)=\frac{|D|}{|S_3|}=\frac{|D|}{6}.
\]
证毕。
\end{proof}

\begin{corollary}[二次符号对不能恢复完整的 \texorpdfstring{$S_3$}{S3} 分裂型]\label{cor:xi-time-part9xg-quadratic-sign-pair-cannot-recover-full-s3-type}
\leanverified{paper\_xi\_time\_part9xg\_quadratic\_sign\_pair\_cannot\_recover\_full\_s3\_type}
任何仅依赖于二次符号对
\[
\bigl(\chi_{10}(p),\chi_{\mathrm{LY}}(p)\bigr)\in\{\pm1\}^2
\]
的证书，都不可能恢复 \(g\bmod p\) 的完整 \(S_3\)-分裂型。更精确地，对任意固定 \(\varepsilon\in\{\pm1\}\)，都有
\[
\delta\!\bigl((1)(1)(1)\mid \chi_{10}=\varepsilon,\ \chi_{\mathrm{LY}}=+1\bigr)=\frac13,
\]
\[
\delta\!\bigl((3)\mid \chi_{10}=\varepsilon,\ \chi_{\mathrm{LY}}=+1\bigr)=\frac23.
\]
因此，在同一二次符号层
\[
\{\chi_{10}=\varepsilon,\ \chi_{\mathrm{LY}}=+1\}
\]
上，“完全分裂”与“完全不可约”两种局部行为仍以正条件密度同时出现。
\end{corollary}
\begin{proof}
定理 \ref{thm:xi-time-part62e-chi-ly-sole-controller} 已经给出：在条件 \(\chi_{\mathrm{LY}}=+1\) 下，\(g\bmod p\) 只可能呈现两种分裂型
\[
(1)(1)(1),\qquad (3),
\]
并且
\[
\delta\!\bigl((1)(1)(1)\mid \chi_{\mathrm{LY}}=+1\bigr)=\frac13,
\qquad
\delta\!\bigl((3)\mid \chi_{\mathrm{LY}}=+1\bigr)=\frac23.
\]
另一方面，定理 \ref{thm:xi-time-part9xg-s10-parity-neutrality-for-s3-splitting} 说明：附加条件 \(\chi_{10}=\varepsilon\) 不会改变任何 \(S_3\)-分裂事件的条件密度。于是
\[
\delta\!\bigl((1)(1)(1)\mid \chi_{10}=\varepsilon,\ \chi_{\mathrm{LY}}=+1\bigr)=\frac13,
\]
\[
\delta\!\bigl((3)\mid \chi_{10}=\varepsilon,\ \chi_{\mathrm{LY}}=+1\bigr)=\frac23.
\]
两种不同分裂型在同一二次符号层上保持正条件密度，故不能由二次符号对唯一恢复。证毕。
\end{proof}

\begin{theorem}[奇偶扭曲谱壁与共尾奇素数慢模的严格分离]\label{thm:xi-time-part9xg-parity-wall-versus-cofinal-prime-twists}
记
\[
\lambda:=\rho(1)=\varphi^2,
\qquad
\rho_-:=\rho\!\bigl(W(-1)\bigr),
\qquad
\rho_p:=\left|\rho\!\left(e^{2\pi i/p}\right)\right|
\quad (p\ge 3\ \text{为奇素数}).
\]
则成立：
\begin{enumerate}
  \item \(\rho_-\) 是三次方程
  \[
  x^3-2x^2-1=0
  \]
  的最大实根，并且
  \[
  \frac{\rho_-}{\lambda}\approx 0.8424525579<1.
  \]
  \item 当 \(p\to\infty\) 且 \(p\) 为奇素数时，
  \[
  \log\frac{\rho_p}{\lambda}
  =
  -\frac{6\sqrt5-5}{250}\left(\frac{2\pi}{p}\right)^2
  +\frac{7+24\sqrt5}{75000}\left(\frac{2\pi}{p}\right)^4
  +O(p^{-6}),
  \]
  从而
  \[
  \frac{\rho_p}{\lambda}\to 1.
  \]
  \item 特别地，存在奇素数阈值 \(p_0\)，使得对一切奇素数 \(p\ge p_0\) 都有
  \[
  \rho_p>\rho_-.
  \]
\end{enumerate}
因此，模 \(2\) 的奇偶扭曲只提供一个固定的谱壁，而不能代表共尾最慢的非平凡扭曲分支。
\end{theorem}
\begin{proof}
推论 \ref{cor:xi-time-part62dg-parity-vs-large-prime-twist-separation} 已经同时给出三件事实：其一，\(\rho_-\) 正是三次方程
\[
x^3-2x^2-1=0
\]
的最大实根，并满足
\[
\frac{\rho_-}{\lambda}\approx 0.8424525579<1,
\]
其二，对奇素数 \(p\) 的扭曲分支有四阶展开
\[
\log\frac{r_p}{\lambda}
=
-\frac{6\sqrt5-5}{250}\left(\frac{2\pi}{p}\right)^2
+\frac{7+24\sqrt5}{75000}\left(\frac{2\pi}{p}\right)^4
+O(p^{-6}),
\]
其三，存在奇素数阈值 \(p_0\) 使得所有 \(p\ge p_0\) 都满足 \(r_p>\rho_-\)。把其中记号 \(r_p\) 仅改写为 \(\rho_p\)，便得到本定理。证毕。
\end{proof}

\begin{corollary}[有限模数模板不能捕捉最慢扭曲分支]\label{cor:xi-time-part9xg-finite-modulus-template-misses-slowest-twist}
\leanverified{paper\_xi\_time\_part9xg\_finite\_modulus\_template\_misses\_slowest\_twist}
设
\[
\mathcal M\subset \NN_{\ge 2}
\]
为任意有限模数集，并定义
\[
r(\mathcal M):=\max_{m\in\mathcal M}\frac{\rho(e^{2\pi i/m})}{\lambda}.
\]
则
\[
r(\mathcal M)<1,
\]
并且存在无穷多个奇素数 \(p\notin\mathcal M\) 满足
\[
\frac{\rho(e^{2\pi i/p})}{\lambda}>r(\mathcal M).
\]
因此，任何固定有限模数模板都不可能捕捉碰撞通道中的共尾最慢扭曲分支。
\end{corollary}
\begin{proof}
首先，对每个 \(m\in\mathcal M\)，都有 \(e^{2\pi i/m}\neq 1\)。而定理 \ref{thm:xi-output-potential-dirichlet-global-worst-mode-j1} 的证明已经给出对应的 Wielandt 型严格不等式：对全部非平凡单位根扭曲 \(t\neq 0\)，皆有
\[
\rho(e^{it})<\lambda.
\]
因此有限集合 \(\mathcal M\) 上的最大值满足
\[
r(\mathcal M)<1.
\]

另一方面，由定理 \ref{thm:xi-time-part9xg-parity-wall-versus-cofinal-prime-twists}，
\[
\frac{\rho(e^{2\pi i/p})}{\lambda}\to 1
\qquad
(p\to\infty,\ p\ \text{为奇素数}).
\]
故可取充分大的奇素数 \(p\notin\mathcal M\)，使
\[
\frac{\rho(e^{2\pi i/p})}{\lambda}>r(\mathcal M).
\]
又因为这样的充分大奇素数有无穷多枚，结论随之成立。证毕。
\end{proof}

\begin{theorem}[有限整数 moment-kernel 样本不能唯一闭合连续压力桥]\label{thm:xi-time-part9xg-finite-integer-moment-samples-do-not-determine-analytic-bridge}
\leanverified{paper\_xi\_time\_part9xg\_finite\_integer\_moment\_samples\_do\_not\_determine\_analytic\_bridge}
沿用定理 \ref{thm:pom-real-q-pressure-analytic-interpolation} 的设定。设 \(J\subset [1,\infty)\) 为该实参数 Perron 分支的解析定义区间，并记
\[
\Lambda_0(q):=\log \lambda(q)
\qquad (q\in J),
\]
其中 \(\lambda(q)\) 为射影压力算子 \(\mathcal L_q\) 的 Perron 主特征值。于是对每个整数 \(q\in J\cap \ZZ_{\ge 2}\)，都有
\[
\Lambda_0(q)=\log r_q.
\]

固定任意有限非空集合
\[
Q_0\subset J\cap \ZZ_{\ge 2}.
\]
则仅凭样本值
\[
\{\log r_q\}_{q\in Q_0}
\]
并不能在“实解析函数”这一类中唯一确定一条连续候选压力曲线。更精确地，对任意包含 \(Q_0\) 的开区间 \(I\subset J\)，都存在两条不同的实解析函数
\[
\Lambda_+,\Lambda_-:I\to\RR
\]
使得
\[
\Lambda_+(q)=\Lambda_-(q)=\log r_q
\qquad (\forall q\in Q_0),
\]
但
\[
\Lambda_+\not\equiv \Lambda_-
\qquad\text{于 }I.
\]
\end{theorem}
\begin{proof}
由定理 \ref{thm:pom-real-q-pressure-analytic-interpolation}，\(\Lambda_0(q)=\log\lambda(q)\) 在区间 \(J\) 上实解析，并满足
\[
\Lambda_0(q)=\log r_q
\qquad
(q\in J\cap \ZZ_{\ge 2}).
\]
特别地，它在 \(Q_0\) 上实现了所需的离散插值。

现在定义非零多项式
\[
P(t):=\prod_{q\in Q_0}(t-q)^2.
\]
则 \(P\) 在每个 \(q\in Q_0\) 处都为零，但 \(P\not\equiv 0\) 于 \(I\)。任取 \(\varepsilon>0\)，定义
\[
\Lambda_\pm(t):=\Lambda_0(t)\pm \varepsilon P(t)
\qquad (t\in I).
\]
由于 \(\Lambda_0\) 与 \(P\) 都实解析，\(\Lambda_\pm\) 亦为 \(I\) 上的实解析函数；并且对每个 \(q\in Q_0\)，都有
\[
\Lambda_\pm(q)=\Lambda_0(q)\pm \varepsilon P(q)=\Lambda_0(q)=\log r_q.
\]
但因 \(P\not\equiv 0\)，故
\[
\Lambda_+-\Lambda_-=2\varepsilon P\not\equiv 0,
\]
于是 \(\Lambda_+\not\equiv \Lambda_-\) 于 \(I\)。

因此，有限多个整数层样本不足以在实解析类中唯一闭合连续压力桥。换言之，定理 \ref{thm:xi-projective-moment-tower-anchors-normalized-pressure} 所给出的规范锚定依赖于完整算子族及其解析制度，而不能由有限节点审计单独替代。证毕。
\end{proof}

\begin{conjecture}[平方自由--二次符号--离散样本的三重盲核]\label{conj:xi-time-part9xg-squarefree-quadratic-analytic-blind-core}
设某个 finitary normalization 方案试图在本文框架中同时只保留下列三类有限证书：
\begin{enumerate}
  \item 由定理 \ref{thm:xi-time-part65d-chain-interior-boolean-arithmeticization} 所模型化的平方自由 gcd/lcm 证书代数；
  \item 由有限个二次 Galois 商所诱导出的符号层
  \[
  (\chi_1,\dots,\chi_r)\in\{\pm1\}^r;
  \]
  \item 有限个整数节点上的 moment-kernel 样本
  \[
  \{\log r_q:\ q\in Q_0\},
  \qquad
  Q_0\subset \ZZ_{\ge 2}\ \text{有限}.
  \]
\end{enumerate}
则该方案必然留下一个不可消除的共同盲核，其缺失恰对应于三条彼此独立的数据：
\begin{enumerate}
  \item Smith 指数高度 \(e_i(p)\) 的 multiplicity-sensitive 信息；
  \item 高维 Artin 分裂型，尤其是 \(S_3\) 层中的完整局部分裂类别；
  \item 连续压力曲线在整数节点之间的解析形状。
\end{enumerate}
等价地，任何真正试图闭合本文 RH 主线的证明，都不可能同时完全因子化到这三类有限证书之中。
\end{conjecture}

\noindent
这八条结果共同指向同一事实：本文框架中的终端障碍并不首先来自 ambient object 的抽象增大，而来自三种有限压缩的同时失真。平方自由 gcd/lcm 证书只能记录素数支撑而不能记录指数高度；二次符号层只能记录 abelianized 的奇偶信息而不能恢复高维 Artin 分裂型；有限多个整数 moment-kernel 样本只能提供离散锚点而不能闭合节点之间的解析压力桥。于是，真正需要排除的并不是“对象过大”这一粗糙直觉，而是这三种 finitary normalization 同时施加时所形成的共同盲核。

\endinput

\paragraph{window-\texorpdfstring{$6$}{6} 的 quartic gauge 量子、黄金密度缺陷与可见奇张量选轴}
在 window-\(6\) 的 bin-fold gauge 分解、审计偶数窗的最小退化 Fibonacci 律与补层--极大纤维锁定，以及 \(X_6^{\mathrm{cyc}}\) 在 \(B_3\) 字典下的坐标四重态分层已经固定之后，还可进一步把 binary gauge 体积、最小退化扇区密度与可见奇张量压缩为下列六条闭式后果。它们表明：同一组有限审计数据同时控制 gauge 体积的纯 quartic 幂、黄金极限的双边有限窗缺陷，以及奇阶可见张量在单一 Cartan 轴上的 rank-\(1\) 选轴。

\begin{theorem}[window-\texorpdfstring{$6$}{6} 的 binary gauge 体积出现纯 quartic 幂塌缩]\label{thm:xi-time-part9y-window6-bin-gauge-quartic-collapse}
\leanverified{paper\_xi\_time\_part9y\_window6\_bin\_gauge\_quartic\_collapse}
沿用推论 \ref{cor:fold-bin-gauge-m6} 的记号，记
\[
\mathcal G^{\mathrm{bin}}_6
\cong
S_2^{\,8}\times S_3^{\,4}\times S_4^{\,9}.
\]
再记
\[
\mathcal S_{6,2}:=\mathcal S_{6,d_{\min}(6)},
\qquad
D_n:=\max_{x\in X_n}d_n(x).
\]
则有精确恒等式
\[
|\mathcal G^{\mathrm{bin}}_6|
=
(2!)^8(3!)^4(4!)^9
=
24^{13}
=
(4!)^{\,|X_6\setminus \mathcal S_{6,2}|}
=
(4!)^{\,D_{10}}.
\]
换言之，window-\(6\) 的异质纤维谱 \(2{:}8,3{:}4,4{:}9\) 在总体 gauge 体积层面并不留下三种独立尺度，而是完全塌缩为 quartic 基元 \(4!\) 的第 \(13\) 次幂。
\end{theorem}
\begin{proof}
由推论 \ref{cor:fold-bin-gauge-m6}，
\[
|\mathcal G^{\mathrm{bin}}_6|
=
(2!)^8(3!)^4(4!)^9.
\]
直接化简得
\[
(2!)^8(3!)^4(4!)^9
=
2^8\cdot 6^4\cdot 24^9
=
2^8\cdot \left(\frac{24}{4}\right)^4\cdot 24^9
=
24^{13}.
\]
另一方面，定理 \ref{thm:fold-bin-min-degeneracy-fib-audited} 给出 \(d_{\min}(6)=2\) 与 \(|\mathcal S_{6,2}|=F_6=8\)，而稳定类型计数律给出 \(|X_6|=F_8=21\)。故
\[
|X_6\setminus \mathcal S_{6,2}|=21-8=13.
\]
再由推论 \ref{cor:fold-bin-minsector-complement-maxfiber-audited}，
\[
|X_6\setminus \mathcal S_{6,2}|=D_{10}.
\]
于是 \(D_{10}=13\)，代回即得
\[
|\mathcal G^{\mathrm{bin}}_6|
=
24^{13}
=
(4!)^{\,|X_6\setminus \mathcal S_{6,2}|}
=
(4!)^{\,D_{10}}.
\]
证毕。
\end{proof}

\begin{theorem}[审计偶数窗最小退化扇区的双边黄金误差律]\label{thm:xi-time-part9y-audited-even-minsector-golden-defect-law}
对审计窗口
\[
m\in\{6,8,10,12\},
\]
沿用定义 \ref{def:fold-bin-degeneracy-sectors} 的记号，则有精确恒等式
\[
\frac{|\mathcal S_{m,d_{\min}(m)}|}{|X_m|}
=
\frac{F_m}{F_{m+2}}
=
\varphi^{-2}-\frac{1}{\varphi^{m+2}F_{m+2}},
\]
\[
\frac{|X_m\setminus \mathcal S_{m,d_{\min}(m)}|}{|X_m|}
=
\frac{F_{m+1}}{F_{m+2}}
=
\varphi^{-1}+\frac{1}{\varphi^{m+2}F_{m+2}}
=
\frac{D_{2m-2}}{|X_m|}.
\]
因此，最小退化扇区与其补层相对于 \((\varphi^{-2},\varphi^{-1})\) 的偏差大小完全相同而符号相反。
\end{theorem}
\begin{proof}
由定理 \ref{thm:fold-bin-min-degeneracy-fib-audited} 与推论 \ref{cor:fold-bin-minsector-complement-maxfiber-audited}，
\[
|\mathcal S_{m,d_{\min}(m)}|=F_m,
\qquad
|X_m\setminus \mathcal S_{m,d_{\min}(m)}|=D_{2m-2}=F_{m+1},
\qquad
|X_m|=F_{m+2}.
\]
记
\[
\psi:=-\varphi^{-1}.
\]
由 Binet 公式 \(F_n=(\varphi^n-\psi^n)/\sqrt5\)，可得
\[
F_{m+2}=\varphi^2F_m+\psi^m,
\qquad
F_{m+2}=\varphi F_{m+1}+\psi^{m+1}.
\]
由于 \(m\) 为偶数，故
\[
\psi^m=\varphi^{-m},
\qquad
-\psi^{m+1}=\varphi^{-m-1}.
\]
于是
\[
\varphi^{-2}-\frac{F_m}{F_{m+2}}
=
\frac{F_{m+2}-\varphi^2F_m}{\varphi^2F_{m+2}}
=
\frac{\psi^m}{\varphi^2F_{m+2}}
=
\frac{1}{\varphi^{m+2}F_{m+2}},
\]
从而
\[
\frac{F_m}{F_{m+2}}
=
\varphi^{-2}-\frac{1}{\varphi^{m+2}F_{m+2}}.
\]
同理，
\[
\frac{F_{m+1}}{F_{m+2}}-\varphi^{-1}
=
\frac{\varphi F_{m+1}-F_{m+2}}{\varphi F_{m+2}}
=
-\frac{\psi^{m+1}}{\varphi F_{m+2}}
=
\frac{1}{\varphi^{m+2}F_{m+2}},
\]
故
\[
\frac{F_{m+1}}{F_{m+2}}
=
\varphi^{-1}+\frac{1}{\varphi^{m+2}F_{m+2}}.
\]
最后代入 \(F_{m+1}=D_{2m-2}\) 即得全部公式。证毕。
\end{proof}

\begin{corollary}[window-\texorpdfstring{$6$}{6} 的 base-\texorpdfstring{$24$}{24} gauge 体积密度与补层密度同一]\label{cor:xi-time-part9y-window6-base24-volume-density-identity}
\leanverified{paper\_xi\_time\_part9y\_window6\_base24\_volume\_density\_identity}
沿用定理 \ref{thm:xi-time-part9y-window6-bin-gauge-quartic-collapse} 的记号，则
\[
\frac{\log_{24}|\mathcal G^{\mathrm{bin}}_6|}{|X_6|}
=
\frac{|X_6\setminus \mathcal S_{6,2}|}{|X_6|}
=
\frac{13}{21}
=
\varphi^{-1}+\frac{1}{21\varphi^8},
\]
并且
\[
1-\frac{\log_{24}|\mathcal G^{\mathrm{bin}}_6|}{|X_6|}
=
\frac{|\mathcal S_{6,2}|}{|X_6|}
=
\frac{8}{21}
=
\varphi^{-2}-\frac{1}{21\varphi^8}.
\]
因此，window-\(6\) 的 base-\(24\) gauge 体积密度并非额外统计量，而恰等于非最小退化层的几何密度。
\end{corollary}
\begin{proof}
由定理 \ref{thm:xi-time-part9y-window6-bin-gauge-quartic-collapse}，
\[
|\mathcal G^{\mathrm{bin}}_6|=24^{13},
\qquad
|X_6\setminus \mathcal S_{6,2}|=13.
\]
故
\[
\log_{24}|\mathcal G^{\mathrm{bin}}_6|=13.
\]
再由 \(|X_6|=21\)，立即得到
\[
\frac{\log_{24}|\mathcal G^{\mathrm{bin}}_6|}{|X_6|}
=
\frac{|X_6\setminus \mathcal S_{6,2}|}{|X_6|}
=
\frac{13}{21}.
\]
将定理 \ref{thm:xi-time-part9y-audited-even-minsector-golden-defect-law} 取 \(m=6\) 即得
\[
\frac{13}{21}
=
\varphi^{-1}+\frac{1}{21\varphi^8},
\qquad
\frac{8}{21}
=
\varphi^{-2}-\frac{1}{21\varphi^8}.
\]
证毕。
\end{proof}

\begin{theorem}[window-\texorpdfstring{$6$}{6} 可见奇张量的极端壳选择律]\label{thm:xi-time-part9y-window6-visible-odd-tensor-extremal-selector}
沿用定理 \ref{thm:xi-window6-cyclic-kernel-bc3-root-datum} 与定理 \ref{thm:xi-window6-dbin-coordinate-quadruple-stratification} 的记号，记
\[
\iota_B:X_6^{\mathrm{cyc}}\xrightarrow{\sim}\Phi(B_3)\subset \RR^3,
\qquad
\widetilde d_6:=d_6^{\mathrm{bin}}\circ \iota_B^{-1}:\Phi(B_3)\to\{2,3,4\}.
\]
则对任意函数 \(f:\{2,3,4\}\to\CC\)，均有精确恒等式
\[
\sum_{\beta\in\Phi(B_3)} f\bigl(\widetilde d_6(\beta)\bigr)\,\beta
=
\bigl(f(4)-f(2)\bigr)e_1,
\]
\[
\sum_{\beta\in\Phi(B_3)} f\bigl(\widetilde d_6(\beta)\bigr)\,\beta^{\otimes 3}
=
\bigl(f(4)-f(2)\bigr)e_1^{\otimes 3}.
\]
因此，在一阶与三阶奇张量层面，可见壳谱并不读取整个三层退化分布，而只读取极大壳与极小壳之间的差值 \(f(4)-f(2)\)；中间壳 \(d=3\) 对这两条奇通道均严格透明。
\end{theorem}
\begin{proof}
由定理 \ref{thm:xi-window6-dbin-coordinate-quadruple-stratification}，
\[
\Phi(B_3)
=
\Phi_2\sqcup \Phi_3\sqcup \Phi_4,
\]
其中
\[
\Phi_2
=
L_x\cup\{-e_1\}
=
\{-e_1\}\cup\{(0,\pm1,\pm1)\},
\]
\[
\Phi_3
=
L_y
=
\{(\pm1,0,\pm1)\},
\]
\[
\Phi_4
=
L_z\cup\{e_1,\pm e_2,\pm e_3\}
=
\{e_1,\pm e_2,\pm e_3\}\cup\{(\pm1,\pm1,0)\}.
\]
对一阶和式，四点集 \(\{(0,\pm1,\pm1)\}\)、\(\{(\pm1,0,\pm1)\}\)、\(\{(\pm1,\pm1,0)\}\) 都按中心对称成对相消，而 \(\pm e_2,\pm e_3\) 亦相互抵消，故
\[
\sum_{\beta\in\Phi_2}\beta=-e_1,
\qquad
\sum_{\beta\in\Phi_3}\beta=0,
\qquad
\sum_{\beta\in\Phi_4}\beta=e_1.
\]
对三阶张量和，同样由
\[
(-\beta)^{\otimes 3}=-\,\beta^{\otimes 3}
\]
知每一对中心对称向量的贡献相消，因此
\[
\sum_{\beta\in\Phi_2}\beta^{\otimes 3}=-e_1^{\otimes 3},
\qquad
\sum_{\beta\in\Phi_3}\beta^{\otimes 3}=0,
\qquad
\sum_{\beta\in\Phi_4}\beta^{\otimes 3}=e_1^{\otimes 3}.
\]
于是
\[
\sum_{\beta\in\Phi(B_3)} f\bigl(\widetilde d_6(\beta)\bigr)\,\beta
=
f(2)\sum_{\beta\in\Phi_2}\beta
+f(3)\sum_{\beta\in\Phi_3}\beta
+f(4)\sum_{\beta\in\Phi_4}\beta
=
\bigl(f(4)-f(2)\bigr)e_1,
\]
\[
\sum_{\beta\in\Phi(B_3)} f\bigl(\widetilde d_6(\beta)\bigr)\,\beta^{\otimes 3}
=
f(2)\sum_{\beta\in\Phi_2}\beta^{\otimes 3}
+f(3)\sum_{\beta\in\Phi_3}\beta^{\otimes 3}
+f(4)\sum_{\beta\in\Phi_4}\beta^{\otimes 3}
=
\bigl(f(4)-f(2)\bigr)e_1^{\otimes 3}.
\]
证毕。
\end{proof}

\begin{corollary}[window-\texorpdfstring{$6$}{6} 的可见奇张量读出塌缩为 rank-\texorpdfstring{$1$}{1}]\label{cor:xi-time-part9y-window6-visible-odd-tensor-rankone}
沿用定理 \ref{thm:xi-time-part9y-window6-visible-odd-tensor-extremal-selector} 的记号。对 \(r\in\{1,3\}\)，定义线性映射
\[
\mathcal O_r:\CC^{\{2,3,4\}}\to (\CC^3)^{\otimes r},
\qquad
\mathcal O_r(f):=\sum_{\beta\in\Phi(B_3)} f\bigl(\widetilde d_6(\beta)\bigr)\,\beta^{\otimes r}.
\]
则
\[
\operatorname{rank}\mathcal O_1=\operatorname{rank}\mathcal O_3=1,
\qquad
\ker\mathcal O_1=\ker\mathcal O_3=\{f:\ f(4)=f(2)\}.
\]
特别地，
\[
\sum_{\beta\in\Phi(B_3)}\beta=0,
\qquad
\sum_{\beta\in\Phi(B_3)}\beta^{\otimes 3}=0,
\]
但
\[
\sum_{\beta\in\Phi(B_3)}\widetilde d_6(\beta)\,\beta=2e_1,
\qquad
\sum_{\beta\in\Phi(B_3)}\widetilde d_6(\beta)\,\beta^{\otimes 3}=2e_1^{\otimes 3}.
\]
因此，window-\(6\) 的可见 pinning 并非仅出现于一阶偶极通道，而是在三阶奇张量通道上以同一系数 \(2\) 锁定到同一 Cartan 轴。
\end{corollary}
\begin{proof}
由定理 \ref{thm:xi-time-part9y-window6-visible-odd-tensor-extremal-selector}，
\[
\mathcal O_1(f)=\bigl(f(4)-f(2)\bigr)e_1,
\qquad
\mathcal O_3(f)=\bigl(f(4)-f(2)\bigr)e_1^{\otimes 3}.
\]
故两者像空间都一维，且核都由条件 \(f(4)=f(2)\) 刻画。
取 \(f\equiv 1\)，便得
\[
\sum_{\beta\in\Phi(B_3)}\beta=0,
\qquad
\sum_{\beta\in\Phi(B_3)}\beta^{\otimes 3}=0.
\]
再取 \(f(d)=d\)，便有
\[
\mathcal O_1(f)=(4-2)e_1=2e_1,
\qquad
\mathcal O_3(f)=(4-2)e_1^{\otimes 3}=2e_1^{\otimes 3}.
\]
证毕。
\end{proof}

\begin{conjecture}[偶数窗口可见奇张量的 extremal-shell rank-\texorpdfstring{$1$}{1} 原理]\label{conj:xi-time-part9y-even-window-visible-odd-tensor-extremal-shell-rankone}
设 \(m\) 为偶数，并假定存在与 window-\(6\) 兼容的 pinned 可见根数据
\[
\alpha_m:X_m^{\mathrm{cyc}}\to \mathfrak h_m^\ast
\]
以及由二进制 uplift 诱导的壳分解
\[
X_m=\bigsqcup_d \mathcal S_{m,d}.
\]
则对每个 \(r\ge 0\)，应存在一条可见主轴 \(\nu_m\in \mathfrak h_m^\ast\) 与线性泛函
\[
\Lambda_{m,r}:\CC^{\{d:\ \mathcal S_{m,d}\neq\emptyset\}}\to \CC
\]
使得对任意壳函数 \(f\) 都有
\[
\sum_{w\in X_m^{\mathrm{cyc}}}
f\bigl(d_m^{\mathrm{bin}}(w)\bigr)\,
\alpha_m(w)^{\otimes(2r+1)}
=
\Lambda_{m,r}(f)\,\nu_m^{\otimes(2r+1)}.
\]
并且 \(\Lambda_{m,r}\) 只依赖于 \(f\) 在极小壳 \(d_{\min}(m)\) 与极大壳 \(d_{\max}(m)\) 上的取值；所有中间壳在奇阶可见张量通道上均严格透明。
在 \(m=6\) 时，定理 \ref{thm:xi-time-part9y-window6-visible-odd-tensor-extremal-selector} 与推论 \ref{cor:xi-time-part9y-window6-visible-odd-tensor-rankone} 给出
\[
\nu_6=e_1,
\qquad
\Lambda_{6,0}(f)=\Lambda_{6,1}(f)=f(4)-f(2),
\]
从而 window-\(6\) 正是该原理的最小审计样本。
\end{conjecture}

\noindent
上述结果表明，window-\(6\) 的异质纤维谱 \(2{:}8,3{:}4,4{:}9\) 在三个彼此不同的层级上同时发生严格压缩：在 gauge 体积层面，它塌缩为纯 quartic 幂 \(24^{13}\)，且其指数同时等于最小退化补层基数与 ordinary Fold 在窗口 \(10\) 的极大纤维数；在审计偶数窗的密度层面，最小扇区与补层相对于 \((\varphi^{-2},\varphi^{-1})\) 的有限窗偏差呈现完全对称的双边修正；在可见奇张量层面，整个 cyclic 壳谱又只留下极端壳差值所驱动的单轴 rank-\(1\) 读出。因而，binary gauge 体积、最小退化几何、黄金密度缺陷与奇张量选轴并不是彼此独立的派生量，而是沿同一条四次 gauge 量子--黄金密度缺陷--秩一奇张量选择律共同闭合。

\endinput

\paragraph{\texorpdfstring{$S_4$}{S4}--Hurwitz 链的几何--表示刚化、\texorpdfstring{$2$}{2}-进全息 Cantor 共轭与 ZG 密度的标量化递推}
在 \(S_4\)-Hurwitz 商塔的 Jacobian/Prym 同源分解、Chevalley--Weil 重数刚性、良素处局部 Euler 因子分裂、Gödel--Tate 地址像的严格自相似与移位共轭，以及 ZG 自然密度的 prime \(J\)-分式和无穷乘积表示都已建立之后，还可继续抽取出下列八条闭合后果。前四条把属 \(49\) 的 \(S_4\)-Hurwitz 链压缩为一套几何--表示完全对齐的骨架分解；中间两条把 \(2\)-进全息地址重新表述为 Cantor 刚性与一侧满移位共轭；最后两条则把 ZG 密度从二维矩阵递推压缩为一条一维素数递推与一个绝对收敛的无穷乘积。

\begin{theorem}[属 \texorpdfstring{$49$}{49} Jacobian 的骨架压缩同源分解]\label{thm:xi-time-part9y-s4-jacobian-skeleton-compression}
设
\[
H:=X/A_4,\qquad
C_6:=X/V_4,\qquad
Z:=X/D_8,\qquad
C_F:=X/S_3,
\]
\[
Y:=X/C_4,
\qquad
P:=\operatorname{Prym}(Y/Z).
\]
则有有理同源分解
\[
J(X)\sim_{\QQ}J(C_6)\times J(C_F)^3\times P^3.
\]
\end{theorem}
\begin{proof}
这正是结论 \ref{con:xi-time-part9n-jacobian-c6-compression} 的直接转写。具体地，由定理 \ref{thm:killo-s4-burnside-kani-rosen-prym-square}，
\[
J(C_6)\sim_{\QQ}J(H)\times J(Z)^2,
\qquad
\operatorname{Prym}(C_6/H)\sim_{\QQ}J(Z)^2.
\]
另一方面，由定理 \ref{thm:killo-s4-genus49-jacobian-computable-isogeny}，
\[
J(X)\sim_{\QQ}J(H)\times J(Z)^2\times J(C_F)^3\times P^3.
\]
将第一式代入第二式，立即得到
\[
J(X)\sim_{\QQ}J(C_6)\times J(C_F)^3\times P^3.
\]
证毕。
\end{proof}

\begin{theorem}[微分空间的几何--表示刚化]\label{thm:xi-time-part9y-s4-hodge-geometric-representation-rigidification}
记 \(S_4\) 的四个非平凡不可约复表示为
\[
\mathrm{sgn},\qquad V_2,\qquad V_3,\qquad V_3'.
\]
则 \(H^0(X,\Omega_X^1)\) 在复 \(S_4\)-模范畴中满足规范张量分解
\[
H^0(X,\Omega_X^1)
\cong
\mathrm{sgn}\otimes H^0(H,\Omega_H^1)
\oplus
V_2\otimes H^0(Z,\Omega_Z^1)
\]
\[
\oplus
V_3\otimes H^0(C_F,\Omega_{C_F}^1)
\oplus
V_3'\otimes H^0(P,\Omega_P^1).
\]
\end{theorem}
\begin{proof}
这是结论 \ref{con:xi-time-part9n-s4-hodge-tensor-factorization} 的直接重述。其证明只需把定理 \ref{thm:killo-s4-genus49-jacobian-computable-isogeny} 给出的同源分解
\[
J(X)\sim_{\QQ}J(H)\times J(Z)^2\times J(C_F)^3\times P^3
\]
在单位元处取余切空间，并与结论 \ref{con:xi-time-part9n-s4-all-subgroup-genera} 的 Chevalley--Weil 重数分解
\[
H^0(X,\Omega_X^1)\cong
5\cdot \mathrm{sgn}\oplus
4\cdot V_2\oplus
3\cdot V_3\oplus
9\cdot V_3'
\]
逐块比较即可。由于
\[
g(H)=5,\qquad g(Z)=4,\qquad g(C_F)=3,\qquad \dim P=9,
\]
四个重数空间恰分别与
\[
H^0(H,\Omega_H^1),\quad H^0(Z,\Omega_Z^1),\quad H^0(C_F,\Omega_{C_F}^1),\quad H^0(P,\Omega_P^1)
\]
同维，并由不可约型的两两不同构性强制唯一配对，故得所示张量分解。证毕。
\end{proof}

\begin{corollary}[\texorpdfstring{$S_4$}{S4}-等变块的严格正交性与分解唯一性]\label{cor:xi-time-part9y-s4-isotypic-orthogonality-uniqueness}
在定理 \ref{thm:xi-time-part9y-s4-hodge-geometric-representation-rigidification} 的记号下，记 \(J(X)\) 中与四个不可约型对应的 \(S_4\)-等型阿贝尔子簇分别为
\[
A_{\mathrm{sgn}},\qquad
A_{V_2},\qquad
A_{V_3},\qquad
A_{V_3'}.
\]
则它们分别满足
\[
A_{\mathrm{sgn}}\sim_{\QQ}J(H),\qquad
A_{V_2}\sim_{\QQ}J(Z)^2,\qquad
A_{V_3}\sim_{\QQ}J(C_F)^3,\qquad
A_{V_3'}\sim_{\QQ}P^3.
\]
若 \(\rho\neq \rho'\) 属于
\[
\{\mathrm{sgn},V_2,V_3,V_3'\},
\]
则有
\[
\operatorname{Hom}_{S_4}^0(A_\rho,A_{\rho'})=0.
\]
因此
\[
J(X)\sim_{\QQ}A_{\mathrm{sgn}}\times A_{V_2}\times A_{V_3}\times A_{V_3'}
\]
是按有理 \(S_4\)-型的唯一等变同源分解，并且任意 \(S_4\)-等变有理自同态在此分解下必为块对角形式。
\end{corollary}
\begin{proof}
任取一个 \(S_4\)-等变同态
\[
f:A_{\rho}\longrightarrow A_{\rho'}.
\]
在单位元处，它诱导 \(S_4\)-等变余切映射
\[
f^\ast:H^0(A_{\rho'},\Omega_{A_{\rho'}}^1)\longrightarrow H^0(A_{\rho},\Omega_{A_{\rho}}^1).
\]
由定理 \ref{thm:xi-time-part9y-s4-hodge-geometric-representation-rigidification}，
\[
H^0(A_{\rho},\Omega_{A_{\rho}}^1)\cong \rho\otimes M_\rho,
\qquad
H^0(A_{\rho'},\Omega_{A_{\rho'}}^1)\cong \rho'\otimes M_{\rho'}
\]
对相应重数空间 \(M_\rho,M_{\rho'}\) 成立。若 \(\rho\not\cong \rho'\)，则由 Schur 引理，
\[
\operatorname{Hom}_{S_4}\!\bigl(\rho'\otimes M_{\rho'},\rho\otimes M_\rho\bigr)=0,
\]
从而 \(f^\ast=0\)。而阿贝尔簇同态若余切微分为零，则同态本身为零，故
\[
\operatorname{Hom}_{S_4}^0(A_\rho,A_{\rho'})=0.
\]

于是不同等型块之间不存在 \(S_4\)-等变耦合；因此任意 \(S_4\)-等变自同态只能分别作用于四个等型块，自然呈块对角形式。等型分解的唯一性亦随之成立。证毕。
\end{proof}

\begin{theorem}[骨架压缩后的局部 \texorpdfstring{$L$}{L}-因子分裂]\label{thm:xi-time-part9y-s4-local-lfactor-skeleton-splitting}
设 \(X,C_6,C_F\) 与阿贝尔簇 \(P\) 在某共同定义域上给出，并取一处同时良约化的有限素点 \(v\)。记其 \(H^1\) 的局部 Weil 多项式分别为
\[
P_v(X,T),\qquad
P_v(C_6,T),\qquad
P_v(C_F,T),\qquad
P_v(P,T).
\]
则有精确因式分解
\[
P_v(X,T)=P_v(C_6,T)\,P_v(C_F,T)^3\,P_v(P,T)^3.
\]
等价地，在去掉有限多个坏因子后，全局 Hasse--Weil \(L\)-函数满足
\[
L(X,s)=L(C_6,s)\,L(C_F,s)^3\,L(P,s)^3.
\]
\end{theorem}
\begin{proof}
由定理 \ref{thm:xi-time-part9y-s4-jacobian-skeleton-compression}，
\[
J(X)\sim_{\QQ}J(C_6)\times J(C_F)^3\times P^3.
\]
在共同定义域上，同源阿贝尔簇的 \(\ell\)-进 Tate 模在任一良约化素点 \(v\) 处给出半单同构，而直积对应 Frobenius 作用的直和。因此 Frobenius 特征多项式满足
\[
P_v\!\bigl(J(X),T\bigr)
=
P_v\!\bigl(J(C_6),T\bigr)\,
P_v\!\bigl(J(C_F),T\bigr)^3\,
P_v(P,T)^3.
\]
对曲线 \(X\) 与 \(C_6,C_F\) 而言，\(H^1\) 的局部 Weil 多项式与其 Jacobian 的相应多项式一致，故
\[
P_v(X,T)=P_v(C_6,T)\,P_v(C_F,T)^3\,P_v(P,T)^3.
\]
将这些局部因子在所有良素点上相乘，便得到去掉有限多个坏因子后的全局 \(L\)-函数分裂。证毕。
\end{proof}

\begin{theorem}[\texorpdfstring{$2$}{2}-进全息像集的 Cantor 刚性]\label{thm:xi-time-part9y-hologram-cantor-rigidity}
设 \(E\) 为有限原子事件集，\(\mathrm{code}:E\hookrightarrow \NN\) 为单射，并记
\[
M:=\max_{e\in E}\mathrm{code}(e),
\qquad
B:=2^L>M,
\qquad
T:=T_2(E_{\min}).
\]
取任意 primitive 地址向量
\[
v\in T\setminus 2T.
\]
对每个无限事件流
\[
\mathbf e=(e_t)_{t\ge 1}\in E^{\NN},
\]
定义其全息地址
\[
\pi(\mathbf e):=X(\mathbf e):=\sum_{t\ge 1}\mathrm{code}(e_t)\,B^{t-1}v,
\qquad
\mathcal X:=\pi(E^{\NN})\subseteq T.
\]
则成立：
\begin{enumerate}
  \item \(\pi\) 是到其像 \(\mathcal X\) 的拓扑同胚；
  \item 若 \(|E|\ge 2\)，则 \(\mathcal X\) 是 \(T\) 中的 Cantor 子集；
  \item 对任意长度 \(n\) 的前缀
  \[
  w=(e_1,\dots,e_n)\in E^n,
  \]
  记
  \[
  X_n(w):=\sum_{t=1}^{n}\mathrm{code}(e_t)\,B^{t-1}v.
  \]
  若 \([w]\subset E^{\NN}\) 表示对应 cylinder，则有精确公式
  \[
  \pi([w])=X_n(w)+B^n\mathcal X.
  \]
\end{enumerate}
\end{theorem}
\begin{proof}
这正是定理 \ref{thm:xi-time-part65d-godeltate-selfsimilar-shift-cylinder} 的直接专门化。该定理已证明：在 \(B>M\) 且 \(\mathrm{code}\) 单射时，地址映射
\[
\Xi:E^{\NN}\longrightarrow \mathcal X_E
\]
是拓扑同胚，并满足对任意事件前缀 \(\mathbf a=(e_1,\dots,e_n)\)，
\[
\Xi([\mathbf a])
=
\left(\sum_{t=1}^{n}\mathrm{code}(e_t)\,B^{t-1}v\right)+B^n\mathcal X_E.
\]
将 \(\Xi,\mathcal X_E\) 分别改记为 \(\pi,\mathcal X\)，便得到第 \(1\) 条与第 \(3\) 条。

若 \(|E|\ge 2\)，则积空间 \(E^{\NN}\) 紧、全不连通且无孤立点，故为 Cantor 空间；其同胚像 \(\mathcal X\) 亦具有同样性质，从而是 \(T\) 中的 Cantor 子集。证毕。
\end{proof}

\begin{theorem}[全息地址动力学与一侧满移位的共轭]\label{thm:xi-time-part9y-hologram-full-shift-conjugacy}
沿用定理 \ref{thm:xi-time-part9y-hologram-cantor-rigidity} 的记号。对每个 \(a\in E\)，定义仿射压缩
\[
F_a:T\longrightarrow T,
\qquad
F_a(x):=\mathrm{code}(a)\,v+Bx.
\]
则：
\begin{enumerate}
  \item \(\mathcal X\) 是唯一满足
  \[
  \mathcal X=\bigsqcup_{a\in E}F_a(\mathcal X)
  \]
  的非空紧集；
  \item 记左移
  \[
  \sigma:E^{\NN}\to E^{\NN},
  \qquad
  \sigma(e_1,e_2,\dots)=(e_2,e_3,\dots),
  \]
  并在 \(\mathcal X\) 上定义分段映射
  \[
  R:\mathcal X\to \mathcal X,
  \qquad
  R(F_a(x)):=x.
  \]
  则 \(R\) 良定义且连续，并满足
  \[
  R\circ \pi=\pi\circ \sigma.
  \]
\end{enumerate}
因此动力系统 \((\mathcal X,R)\) 与一侧满移位 \((E^{\NN},\sigma)\) 拓扑共轭。
\end{theorem}
\begin{proof}
在 \(2\)-进度量下，
\[
|B|_2=2^{-L}<1,
\]
故每个 \(F_a\) 都是严格压缩。又由定理 \ref{thm:xi-time-part65d-godeltate-selfsimilar-shift-cylinder}，
\[
\mathcal X=\bigsqcup_{a\in E}F_a(\mathcal X)
\]
已成立。于是由 Hutchinson 不动点原理，\(\mathcal X\) 是该压缩迭代系统的唯一非空紧不动点，得第 \(1\) 条。

对任意 \(\mathbf e=(e_1,e_2,\dots)\in E^{\NN}\)，按首位拆解地址和可得
\[
\pi(\mathbf e)=\mathrm{code}(e_1)\,v+B\,\pi(\sigma\mathbf e)=F_{e_1}(\pi(\sigma\mathbf e)).
\]
因此 \(\mathcal X\) 在各片
\[
F_a(\mathcal X)\qquad (a\in E)
\]
上的分解严格不交，映射
\[
R|_{F_a(\mathcal X)}:=F_a^{-1}
\]
良定义。由各片在 \(\mathcal X\) 中均为 clopen，\(R\) 连续。再由上式立即得到
\[
R(\pi(\mathbf e))
=
R\!\bigl(F_{e_1}(\pi(\sigma\mathbf e))\bigr)
=
\pi(\sigma\mathbf e),
\]
即
\[
R\circ \pi=\pi\circ \sigma.
\]
结合定理 \ref{thm:xi-time-part9y-hologram-cantor-rigidity} 中 \(\pi\) 为同胚，便得到所述拓扑共轭。证毕。
\end{proof}

\begin{theorem}[Zeckendorf--Gödel 密度的标量化素数递推]\label{thm:xi-time-part9y-zg-scalar-prime-recursion}
\leanverified{paper\_xi\_time\_part9y\_zg\_scalar\_prime\_recursion}
设
\[
A_p:=
\begin{pmatrix}
\dfrac{p}{p+1} & \dfrac{p}{p+1}\\[1ex]
\dfrac{1}{p+1} & 0
\end{pmatrix},
\qquad
\binom{x_n}{y_n}:=
A_{p_n}\cdots A_{p_1}\binom10,
\]
其中 \(p_n\) 为第 \(n\) 个素数。定义
\[
r_n:=\frac{y_n}{x_n}.
\]
则成立：
\begin{enumerate}
  \item
  \[
  r_1=\frac1{p_1},
  \qquad
  r_{n+1}=\frac{1}{p_{n+1}(1+r_n)}
  \qquad (n\ge 1);
  \]
  \item 因而
  \[
  r_n=
  \cfrac{1}{p_n\!\left(1+\cfrac{1}{p_{n-1}\!\left(1+\cfrac{1}{\ddots\left(1+\cfrac1{p_1}\right)}\right)}\right)};
  \]
  \item 对一切 \(n\ge 1\)，有
  \[
  0<r_n\le \frac1{p_n},
  \]
  且对 \(n\ge 2\) 时严格不等式成立；从而 \(r_n\to 0\)；
  \item 记 \(\delta_{\mathrm{ZG}}:=\delta(\mathrm{ZG})\)。则
  \[
  \delta_{\mathrm{ZG}}=
  \frac1{\zeta(2)}
  \lim_{n\to\infty}(x_n+y_n).
  \]
\end{enumerate}
\end{theorem}
\begin{proof}
这是定理 \ref{thm:xi-time-part62dgd-zg-density-prime-jfraction-positive} 的记号改写。该定理中记
\[
\binom{A_n}{B_n}:=
\left(\prod_{i=1}^{n}
\begin{pmatrix}
\dfrac{p_i}{p_i+1} & \dfrac{p_i}{p_i+1}\\[1ex]
\dfrac{1}{p_i+1} & 0
\end{pmatrix}
\right)e_1,
\]
并设
\[
t_n:=A_n+B_n,\qquad r_n:=\frac{B_n}{A_n}.
\]
将 \(A_n,B_n\) 改记为 \(x_n,y_n\) 后，定理 \ref{thm:xi-time-part62dgd-zg-density-prime-jfraction-positive} 的第一条即给出
\[
r_n=\frac{1}{p_n(1+r_{n-1})},
\qquad
r_0:=0.
\]
特别地，
\[
r_1=\frac{1}{p_1}.
\]
递归展开便得到所示连分式表示。

再由同一定理证明中的估计，
\[
0<r_n\le \frac1{p_n},
\]
且因 \(r_{n-1}>0\)，实际上对 \(n\ge 2\) 有严格不等式
\[
r_n<\frac1{p_n}.
\]
故 \(r_n\to 0\)。

最后，由同一定理的第二条，
\[
\delta_{\mathrm{ZG}}
=
\frac1{\zeta(2)}\lim_{n\to\infty}t_n.
\]
而 \(t_n=x_n+y_n\)，故得到
\[
\delta_{\mathrm{ZG}}=
\frac1{\zeta(2)}
\lim_{n\to\infty}(x_n+y_n).
\]
证毕。
\end{proof}

\begin{theorem}[Zeckendorf--Gödel 密度的无穷乘积化与正性机制]\label{thm:xi-time-part9y-zg-infinite-product-positivity}
在定理 \ref{thm:xi-time-part9y-zg-scalar-prime-recursion} 的记号下，约定
\[
r_0:=0.
\]
则
\[
\delta_{\mathrm{ZG}}
=
\frac1{\zeta(2)}
\prod_{n\ge 1}
\left(
1-\frac{r_{n-1}}{(p_n+1)(1+r_{n-1})}
\right).
\]
该无穷乘积绝对收敛，且每个因子都位于 \((0,1]\) 内；因此
\[
\delta_{\mathrm{ZG}}>0.
\]
\end{theorem}
\begin{proof}
记
\[
s_n:=x_n+y_n.
\]
由矩阵递推，
\[
x_n=\frac{p_n}{p_n+1}(x_{n-1}+y_{n-1})
=
\frac{p_n}{p_n+1}s_{n-1},
\qquad
y_n=\frac{1}{p_n+1}x_{n-1}.
\]
又因
\[
x_{n-1}=\frac{s_{n-1}}{1+r_{n-1}},
\]
故
\[
s_n
=
\frac{p_n}{p_n+1}s_{n-1}
+
\frac{1}{p_n+1}\frac{s_{n-1}}{1+r_{n-1}}
=
s_{n-1}\left(1-\frac{r_{n-1}}{(p_n+1)(1+r_{n-1})}\right).
\]
迭代得到
\[
s_n
=
s_0\prod_{j=1}^{n}
\left(1-\frac{r_{j-1}}{(p_j+1)(1+r_{j-1})}\right).
\]
由于 \(s_0=1\) 且
\[
\delta_{\mathrm{ZG}}=\frac1{\zeta(2)}\lim_{n\to\infty}s_n
\]
由定理 \ref{thm:xi-time-part9y-zg-scalar-prime-recursion} 成立，遂得所示无穷乘积公式。

再由同一定理，
\[
0<r_{j-1}<\frac1{p_{j-1}}
\qquad (j\ge 2),
\]
故对 \(j\ge 2\) 有
\[
0<
\frac{r_{j-1}}{(p_j+1)(1+r_{j-1})}
<
\frac{1}{p_{j-1}(p_j+1)}.
\]
右端级数收敛，故
\[
\sum_{j\ge 2}\frac{r_{j-1}}{(p_j+1)(1+r_{j-1})}<\infty.
\]
由正项无穷乘积的标准判据，所示乘积收敛；又每个因子都严格位于 \((0,1]\) 内，故其极限严格为正。乘上常数 \(\zeta(2)^{-1}>0\)，得到
\[
\delta_{\mathrm{ZG}}>0.
\]
证毕。
\end{proof}

\noindent
上述八条结果把两条此前分离的主线同时压缩为可直接横向调用的终端包：一方面，属 \(49\) 的 \(S_4\)-Hurwitz 链在 Jacobian、Hodge 余切空间与局部 \(L\)-因子三侧完成了同一骨架压缩，并且不同不可约型之间的等变耦合被 Schur 刚性完全切断；另一方面，\(2\)-进全息地址像不仅是一个 Cantor 型刚性对象，而且其自然动力学与一侧满移位严格共轭；再一方面，ZG 密度又不再需要以二维矩阵极限保留，而可被压缩为一条一维素数 Riccati 递推与一个绝对收敛的正无穷乘积。因此，\(S_4\)-几何骨架、非阿基米德地址动力学与 prime-index hard-core 密度机制在这一层面上都获得了可审计的标量化终端表达。

\endinput

\paragraph{window-\texorpdfstring{$6$}{6} 的三原子 moment 测度、Hankel 刚性与 involutive thermodynamics}
在 window-\(6\) 的完整有限尺度谱、bin-fold 纤维群胚代数的 Wedderburn 分解、输出熵账本以及三能级 escort 模型已经固定之后，还可把稳定类型、微态、碰撞、矩阵代数与有限热力学进一步压缩为单一的三原子 Stieltjes 对象。记
\[
d_6(w):=d_6^{\mathrm{bin}}(w),
\qquad
\mu_6:=\sum_{w\in X_6}\delta_{d_6(w)}=8\delta_2+4\delta_3+9\delta_4,
\]
并沿用本文关于 involutive groupoid algebra 的记号
\[
\mathcal A_6:=\CC[\mathcal R_6^{\mathrm{bin}}],
\qquad
\mathcal A_6^{\alpha,\pm}:=\{a\in \mathcal A_6:\ \alpha(a)=\pm a\}.
\]
下列结果表明，window-\(6\) 的多层对象并非平行排布：\(\mu_6\) 的矩序列具有三阶湮没多项式，其 Hankel 秩恰等于支撑点数，其前三个 moment 在 \((\mathcal A_6,\alpha)\) 中分别实现为中心、反不动扇区与全代数，而输出 Rényi 谱、dyadic 可逆容量与 gauge 型热力学量又全部由同一测度通过显式积分恢复。

\begin{theorem}[window-\texorpdfstring{$6$}{6} 纤维多重度测度的湮没三次式与有理母函数]\label{thm:xi-time-part9z-window6-threeatom-moment-ogf}
\leanverified{paper\_xi\_time\_part9z\_window6\_threeatom\_moment\_ogf}
定义矩序列
\[
M_n:=\int t^n\,\dd\mu_6(t)
=
\sum_{w\in X_6}d_6(w)^n
=
8\cdot 2^n+4\cdot 3^n+9\cdot 4^n
\qquad (n\ge 0).
\]
则其 ordinary generating function 满足
\[
\mathcal M_6(z):=\sum_{n\ge 0}M_n z^n
=
\frac{21-125z+182z^2}{(1-2z)(1-3z)(1-4z)}.
\]
特别地，\((M_n)_{n\ge 0}\) 满足精确线性递推
\[
M_{n+3}=9M_{n+2}-26M_{n+1}+24M_n
\qquad (n\ge 0),
\]
其对应特征多项式为
\[
p_6(\lambda):=(\lambda-2)(\lambda-3)(\lambda-4)
=
\lambda^3-9\lambda^2+26\lambda-24.
\]
\end{theorem}
\begin{proof}
由 \(\mu_6=8\delta_2+4\delta_3+9\delta_4\)，对每个 \(n\ge 0\) 直接有
\[
M_n=8\cdot 2^n+4\cdot 3^n+9\cdot 4^n.
\]
因此
\[
\mathcal M_6(z)
=
\frac{8}{1-2z}+\frac{4}{1-3z}+\frac{9}{1-4z},
\]
通分后得到
\[
\mathcal M_6(z)
=
\frac{21-125z+182z^2}{(1-2z)(1-3z)(1-4z)}.
\]
又
\[
(1-2z)(1-3z)(1-4z)=1-9z+26z^2-24z^3,
\]
故生成函数分母对应的递推正是
\[
M_{n+3}-9M_{n+2}+26M_{n+1}-24M_n=0.
\]
把 \(\lambda\) 视为移位算子的特征值变量，即得所示特征多项式。证毕。
\end{proof}

\begin{theorem}[window-\texorpdfstring{$6$}{6} 三原子 moment 测度的 Hankel 平坦化]\label{thm:xi-time-part9z-window6-threeatom-hankel-flatness}
对 \(r\ge 0\)，记
\[
H_r^{(6)}:=(M_{i+j})_{0\le i,j\le r}.
\]
则有
\[
\det H_2^{(6)}=1152>0,
\qquad
\det H_3^{(6)}=0.
\]
因此 \((M_n)\) 的 Hankel 秩恰为 \(3\)。等价地，\(\mu_6\) 恰为三原子测度，而其支撑由定理 \ref{thm:xi-time-part9z-window6-threeatom-moment-ogf} 的湮没多项式唯一恢复为
\[
\operatorname{supp}\mu_6=\{2,3,4\}.
\]
\leanverified{paper\_xi\_time\_part9z\_window6\_threeatom\_hankel\_flatness}
\end{theorem}
\begin{proof}
对每个 \(r\ge 0\)，记
\[
V_r:=
\begin{pmatrix}
1&1&1\\
2&3&4\\
\vdots&\vdots&\vdots\\
2^r&3^r&4^r
\end{pmatrix},
\qquad
D:=\operatorname{diag}(8,4,9).
\]
由
\[
M_{i+j}=8\cdot 2^{i+j}+4\cdot 3^{i+j}+9\cdot 4^{i+j},
\]
可得 Hankel 分解
\[
H_r^{(6)}=V_r D V_r^{\mathsf T}.
\]
于是对所有 \(r\) 都有 \(\operatorname{rank}H_r^{(6)}\le 3\)，从而
\[
\det H_3^{(6)}=0.
\]

当 \(r=2\) 时，\(V_2\) 为 \(3\times 3\) Vandermonde 矩阵，故
\[
\det H_2^{(6)}
=
(\det V_2)^2\det D
=
\bigl((3-2)(4-2)(4-3)\bigr)^2\cdot 8\cdot 4\cdot 9
=
1152.
\]
因此 \(\operatorname{rank}H_2^{(6)}=3\)，遂 \((M_n)\) 的 Hankel 秩恰为 \(3\)。

另一方面，三原子测度的矩序列最小递推阶恰等于其 Hankel 秩；故定理 \ref{thm:xi-time-part9z-window6-threeatom-moment-ogf} 中的三次湮没多项式已达到最小可能阶数。其根 \(2,3,4\) 两两不同，因此正是 \(\mu_6\) 的支撑点。证毕。
\end{proof}

\begin{theorem}[window-\texorpdfstring{$6$}{6} 的 \texorpdfstring{$0$}{0}--\texorpdfstring{$1$}{1}--\texorpdfstring{$2$}{2} moment 在 involutive algebra 中的实现]\label{thm:xi-time-part9z-window6-moment-involutive-algebra-realization}
window-\(6\) 的前三个纤维 moment 满足
\[
M_0=\dim Z(\mathcal A_6)=21,
\]
\[
M_1=\dim \mathcal A_6^{\alpha,-}=64,
\]
\[
M_2=\dim \mathcal A_6=212.
\]
因此，稳定类型数、微态总数与有序碰撞对总数并非三套彼此平行的量，而分别是同一测度 \(\mu_6\) 的 \(0\)、\(1\)、\(2\) 阶 moment，并在 involutive groupoid algebra 中依次实现为中心、反不动扇区与全代数。
\end{theorem}
\begin{proof}
由命题 \ref{prop:fold-bin-groupoid-algebra-wedderburn} 与推论 \ref{cor:fold-bin-m6-fiber-hist}，
\[
\mathcal A_6
\cong
M_2(\CC)^{\oplus 8}\oplus M_3(\CC)^{\oplus 4}\oplus M_4(\CC)^{\oplus 9}.
\]
于是
\[
\dim Z(\mathcal A_6)=8+4+9=21=M_0,
\]
且
\[
\dim \mathcal A_6
=
8\cdot 2^2+4\cdot 3^2+9\cdot 4^2
=
212
=
M_2.
\]

对 \(M_1\)，一方面，纤维分割覆盖全部 \(2^6\) 个微态，故
\[
M_1=\sum_{w\in X_6}d_6(w)=2^6=64.
\]
另一方面，本文已审计数据给出
\[
\dim \mathcal A_6^{\alpha,-}=64.
\]
遂得第二式。证毕。
\end{proof}

\begin{corollary}[window-\texorpdfstring{$6$}{6} 的 \texorpdfstring{$\alpha$}{alpha}-固定扇区恰计数非对角碰撞]\label{cor:xi-time-part9z-window6-fixed-sector-distinct-collision}
有精确恒等式
\[
\dim \mathcal A_6^{\alpha,+}
=
M_2-M_1
=
\sum_{w\in X_6}d_6(w)\bigl(d_6(w)-1\bigr)
=
148.
\]
换言之，\(\alpha\)-固定扇区并不计数全部有序碰撞对，而恰计数有序互异碰撞对；全部对角配对 \((\omega,\omega)\) 的总数 \(64\) 则完全落在 \(\alpha\)-反不动扇区。
\end{corollary}
\begin{proof}
由于 \(\alpha\) 为对合，\(\mathcal A_6=\mathcal A_6^{\alpha,+}\oplus \mathcal A_6^{\alpha,-}\)。故由定理 \ref{thm:xi-time-part9z-window6-moment-involutive-algebra-realization}，
\[
\dim \mathcal A_6^{\alpha,+}
=
\dim \mathcal A_6-\dim \mathcal A_6^{\alpha,-}
=
212-64
=
148.
\]
另一方面，
\[
\sum_{w\in X_6}d_6(w)\bigl(d_6(w)-1\bigr)
=
\sum_{w\in X_6}d_6(w)^2-\sum_{w\in X_6}d_6(w)
=
M_2-M_1.
\]
合并两式即得结论。证毕。
\end{proof}

\begin{theorem}[window-\texorpdfstring{$6$}{6} 的前三个 moment 唯一恢复完整纤维直方图]\label{thm:xi-time-part9z-window6-histogram-from-three-moments}
记
\[
N_2:=\#\{w\in X_6:\ d_6(w)=2\},
\qquad
N_3:=\#\{w\in X_6:\ d_6(w)=3\},
\qquad
N_4:=\#\{w\in X_6:\ d_6(w)=4\}.
\]
则 \((N_2,N_3,N_4)\) 可由 \((M_0,M_1,M_2)\) 唯一恢复，且显式公式为
\[
N_2=6M_0-\frac72 M_1+\frac12 M_2,
\]
\[
N_3=-8M_0+6M_1-M_2,
\]
\[
N_4=3M_0-\frac52 M_1+\frac12 M_2.
\]
代入
\[
(M_0,M_1,M_2)=(21,64,212)
\]
即得
\[
(N_2,N_3,N_4)=(8,4,9).
\]
\leanverified{paper\_xi\_time\_part9z\_window6\_histogram\_from\_three\_moments}
\end{theorem}
\begin{proof}
按定义，
\[
N_2+N_3+N_4=M_0,
\]
\[
2N_2+3N_3+4N_4=M_1,
\]
\[
4N_2+9N_3+16N_4=M_2.
\]
由前两式得
\[
N_3+2N_4=M_1-2M_0.
\]
再用第三式减去第二式的两倍，得到
\[
N_3+4N_4=M_2-2M_1.
\]
两式相减即得
\[
2N_4=M_2-5M_1+6M_0,
\]
故
\[
N_4=3M_0-\frac52M_1+\frac12M_2.
\]
代回
\(
N_3+2N_4=M_1-2M_0
\)
得
\[
N_3=-8M_0+6M_1-M_2.
\]
最后由 \(N_2=M_0-N_3-N_4\)，得到
\[
N_2=6M_0-\frac72M_1+\frac12M_2.
\]
代入 \((M_0,M_1,M_2)=(21,64,212)\) 便得 \((8,4,9)\)。证毕。
\end{proof}

\begin{theorem}[window-\texorpdfstring{$6$}{6} 输出分布的精确 Rényi 谱由三原子测度闭合]\label{thm:xi-time-part9z-window6-threeatom-renyi-spectrum}
令
\[
p_6(w):=\frac{d_6(w)}{64}
\qquad (w\in X_6)
\]
为 bin-fold 在稳定类型层上的输出分布。则对任意 \(q\in\RR\setminus\{1\}\)，其 Rényi 熵满足
\[
H_q(p_6)
=
\frac{1}{1-q}
\Bigl[
\log\bigl(8\cdot 2^q+4\cdot 3^q+9\cdot 4^q\bigr)-6q\log 2
\Bigr].
\]
相应极限值为
\[
H_0(p_6)=\log 21,
\]
\[
H_1(p_6)=\frac{37}{8}\log 2-\frac{3}{16}\log 3,
\]
\[
H_2(p_6)=10\log 2-\log 53,
\]
\[
H_\infty(p_6)=4\log 2.
\]
\end{theorem}
\begin{proof}
由 \(\mu_6=8\delta_2+4\delta_3+9\delta_4\)，有
\[
\sum_{w\in X_6}p_6(w)^q
=
\frac{8\cdot 2^q+4\cdot 3^q+9\cdot 4^q}{64^q}.
\]
代入 Rényi 熵定义即得第一式。

\(q\to 0\) 时，有效支撑大小为 \(|X_6|=21\)，故
\[
H_0(p_6)=\log 21.
\]
\(q\to 1\) 时，Shannon 熵为
\[
H_1(p_6)
=
-\sum_{w\in X_6}p_6(w)\log p_6(w)
=
6\log 2-\frac{1}{64}\sum_{w\in X_6}d_6(w)\log d_6(w),
\]
从而
\[
H_1(p_6)
=
6\log 2-\frac{1}{64}\Bigl(8\cdot 2\log 2+4\cdot 3\log 3+9\cdot 4\log 4\Bigr)
=
\frac{37}{8}\log 2-\frac{3}{16}\log 3.
\]

对 \(q=2\)，有
\[
\sum_{w\in X_6}p_6(w)^2
=
\frac{M_2}{64^2}
=
\frac{212}{4096}
=
\frac{53}{1024},
\]
故
\[
H_2(p_6)=-\log\frac{53}{1024}=10\log 2-\log 53.
\]
最后，
\[
\max_{w\in X_6}p_6(w)=\frac{4}{64}=\frac{1}{16},
\]
故
\[
H_\infty(p_6)=-\log\frac{1}{16}=4\log 2.
\]
证毕。
\end{proof}

\begin{theorem}[window-\texorpdfstring{$6$}{6} escort 壳层族的射影可识别性与四室相图]\label{thm:xi-time-part9z-window6-escort-projective-identifiability}
\leanverified{paper\_xi\_time\_part9z\_window6\_escort\_projective\_identifiability}
对任意 \(q\in\RR\)，记
\[
Z_6(q):=8\cdot 2^q+4\cdot 3^q+9\cdot 4^q,
\]
并定义三条壳层 escort 概率
\[
\pi_2(q):=\frac{8\cdot 2^q}{Z_6(q)},
\qquad
\pi_3(q):=\frac{4\cdot 3^q}{Z_6(q)},
\qquad
\pi_4(q):=\frac{9\cdot 4^q}{Z_6(q)}.
\]
则参数 \(q\) 可由单一交叉比精确恢复：
\[
q=
\frac{
\log\!\Bigl(\dfrac{9\,\pi_3(q)^2}{2\,\pi_2(q)\pi_4(q)}\Bigr)
}{\log(9/8)}.
\]
因此，window-\(6\) 的 escort 壳层族不是仅凭数值拟合辨识的一维曲线，而是严格射影可识别的参数族。

此外，三壳层权重的次序随 \(q\) 恰发生三次交换；相应临界温度为
\[
q_{3=4}:=\frac{\log(4/9)}{\log(4/3)}\approx -2.8188416793,
\]
\[
q_{2=4}:=\frac{\log(8/9)}{\log 2}\approx -0.1699250014,
\]
\[
q_{2=3}:=\frac{\log 2}{\log(3/2)}\approx 1.7095112914.
\]
从而三壳层相图严格分裂为四个 chamber：
\[
q<q_{3=4}:\qquad \pi_2(q)>\pi_3(q)>\pi_4(q),
\]
\[
q_{3=4}<q<q_{2=4}:\qquad \pi_2(q)>\pi_4(q)>\pi_3(q),
\]
\[
q_{2=4}<q<q_{2=3}:\qquad \pi_4(q)>\pi_2(q)>\pi_3(q),
\]
\[
q>q_{2=3}:\qquad \pi_4(q)>\pi_3(q)>\pi_2(q).
\]
\end{theorem}
\begin{proof}
直接计算得
\[
\frac{\pi_3(q)^2}{\pi_2(q)\pi_4(q)}
=
\frac{(4\cdot 3^q)^2}{(8\cdot 2^q)(9\cdot 4^q)}
=
\frac{2}{9}\left(\frac{9}{8}\right)^q,
\]
整理即得
\[
q=
\frac{
\log\!\bigl(9\pi_3(q)^2/(2\pi_2(q)\pi_4(q))\bigr)
}{\log(9/8)}.
\]
故 \(q\) 已由单一交叉比唯一恢复。

再看壳层次序。三条成对比值满足
\[
\frac{\pi_2(q)}{\pi_3(q)}=2\left(\frac{2}{3}\right)^q,
\qquad
\frac{\pi_2(q)}{\pi_4(q)}=\frac{8}{9}\,2^{-q},
\qquad
\frac{\pi_3(q)}{\pi_4(q)}=\frac{4}{9}\left(\frac{3}{4}\right)^q.
\]
它们都是 \(q\) 的严格递减函数，因此每一对壳层只会发生一次次序交换。令各比值等于 \(1\) 即得三条临界曲线
\[
8\cdot 2^q=4\cdot 3^q,
\qquad
8\cdot 2^q=9\cdot 4^q,
\qquad
4\cdot 3^q=9\cdot 4^q,
\]
解得 \(q_{2=3}\)、\(q_{2=4}\)、\(q_{3=4}\) 三式。

由其数值顺序
\[
q_{3=4}<q_{2=4}<q_{2=3}
\]
可知：当 \(q<q_{3=4}\) 时，三条比值均大于 \(1\)，故 \(\pi_2>\pi_3>\pi_4\)；越过 \(q_{3=4}\) 后，只有 \(\pi_3/\pi_4\) 变为小于 \(1\)，故 \(\pi_2>\pi_4>\pi_3\)；越过 \(q_{2=4}\) 后，再有 \(\pi_2/\pi_4<1\)，故 \(\pi_4>\pi_2>\pi_3\)；最后越过 \(q_{2=3}\) 后，\(\pi_2/\pi_3<1\)，故 \(\pi_4>\pi_3>\pi_2\)。证毕。
\end{proof}

\begin{theorem}[window-\texorpdfstring{$6$}{6} 可逆容量的 Stieltjes 势表示与 Mellin--Stieltjes 对偶]\label{thm:xi-time-part9z-window6-stieltjes-capacity-duality}
对整数 \(B\ge 0\)，定义 dyadic 预言机下的最大可逆恢复容量
\[
C_6(B):=\sum_{w\in X_6}\min(2^B,d_6(w)).
\]
则有精确表示
\[
C_6(B)=\int \min(2^B,t)\,\dd\mu_6(t).
\]
若再定义尾函数
\[
T_6(s):=\mu_6([s,\infty)),
\]
则
\[
C_6(B)=\int_0^{2^B}T_6(s)\,\dd s.
\]
进一步，把整数矩沿实参数延拓为
\[
M_s:=\int t^s\,\dd\mu_6(t)
=
8\cdot 2^s+4\cdot 3^s+9\cdot 4^s
\qquad (s>0),
\]
则有精确 Mellin--Stieltjes 恒等式
\[
M_s=s\int_0^\infty u^{s-1}T_6(u)\,\dd u
\qquad (s>0).
\]
对 window-\(6\) 而言，
\[
T_6(s)=
\begin{cases}
21,& 0<s\le 2,\\
13,& 2<s\le 3,\\
9,& 3<s\le 4,\\
0,& s>4,
\end{cases}
\]
从而
\[
C_6(0)=21,\qquad
C_6(1)=42,\qquad
C_6(B)=64\ \ (B\ge 2).
\]
\end{theorem}
\begin{proof}
由 \(\mu_6=\sum_{w\in X_6}\delta_{d_6(w)}\)，直接有
\[
\int \min(2^B,t)\,\dd\mu_6(t)
=
\sum_{w\in X_6}\min(2^B,d_6(w))
=
C_6(B).
\]

对任意 \(a>0\)，恒等式
\[
\min(a,t)=\int_0^a \mathbf 1_{\{t\ge s\}}\,\dd s
\]
成立。将其对 \(\mu_6\) 积分，得到
\[
\int \min(a,t)\,\dd\mu_6(t)
=
\int_0^a \mu_6([s,\infty))\,\dd s
=
\int_0^a T_6(s)\,\dd s.
\]
取 \(a=2^B\) 即得容量公式。

对 \(s>0\)，对有限正测度施行 Stieltjes 分部积分，得到
\[
\int t^s\,\dd\mu_6(t)
=
s\int_0^\infty u^{s-1}\mu_6([u,\infty))\,\dd u
=
s\int_0^\infty u^{s-1}T_6(u)\,\dd u.
\]
最后，由 \(\mu_6=8\delta_2+4\delta_3+9\delta_4\) 可立即读出 \(T_6\) 的分段表达；逐段积分便得
\[
C_6(0)=\int_0^1 21\,\dd s=21,
\qquad
C_6(1)=\int_0^2 21\,\dd s=42,
\]
而当 \(B\ge 2\) 时，\(2^B\ge 4\)，故
\[
C_6(B)
=
\int_0^4 T_6(s)\,\dd s
=
21\cdot 2+13\cdot 1+9\cdot 1
=
64.
\]
证毕。
\end{proof}
\endinput

\paragraph{window-\texorpdfstring{$6$}{6} 规范站点抬升的条件均匀化、局域熵势与 GFF 标度不变性}
在 window-\(6\) 的可见/隐藏 Hilbert 分解、bin-fold 规范群的直积分解、纤维多重度的三档阈值分类、\(B_3/C_3\) 根字典以及隐藏部门的 \(2,3,4\) 三质量量子化都已建立之后，还可把同一窗口上的 canonical sitewise lift 再压缩为一组完全显式的有限体积恒等式与连续极限后果。以下固定
\[
\pi:=\Fold_6^{\mathrm{bin}}:\Omega_6\longrightarrow X_6,
\qquad
F_w:=\pi^{-1}(w),
\qquad
d(w):=|F_w|.
\]
由定理 \ref{thm:xi-window6-cyclic-kernel-bc3-root-datum}、推论 \ref{cor:fold-bin-m6-fiber-hist}、定理 \ref{thm:xi-foldbin-lastbit-tail-spectrum-single-jump-staircase} 与定理 \ref{thm:xi-foldbin-gauge-representation-distribution-law-m6-spectrum} 可知：
\[
|X_6|=21,
\qquad
X_6=X_6^{\mathrm{cyc}}\sqcup X_6^{\mathrm{bdry}},
\qquad
|X_6^{\mathrm{cyc}}|=18,
\qquad
|X_6^{\mathrm{bdry}}|=3,
\]
\[
d(w)\in\{2,3,4\},
\qquad
\#\{d=2\}=8,\quad
\#\{d=3\}=4,\quad
\#\{d=4\}=9,
\]
\[
d(w)=2\iff w_6=1,
\qquad
d(w)=4\iff w_6=0,\ V_6(w)\le 8,
\qquad
d(w)=3\iff w_6=0,\ V_6(w)>8,
\]
并且
\[
G_6^{\mathrm{bin}}
\cong
\prod_{w\in X_6}\mathfrak S_{d(w)}.
\]
从现在起，对任意有限位点集 \(\Lambda\)，任取可见可容集
\[
\mathcal W_\Lambda\subset X_6^\Lambda,
\]
并定义其 canonical sitewise lift
\[
\widehat{\mathcal W}_\Lambda
:=
(\pi^\Lambda)^{-1}(\mathcal W_\Lambda)
\subset \Omega_6^\Lambda,
\qquad
\mathcal G_\Lambda:=
\bigl(G_6^{\mathrm{bin}}\bigr)^\Lambda.
\]

\begin{theorem}[规范抬升的条件均匀化]\label{thm:xi-time-part9z-window6-canonical-lift-conditional-uniformization}
设 \(\widehat\mu_\Lambda\) 是 \(\widehat{\mathcal W}_\Lambda\) 上的概率测度，并满足对一切
\[
g\in \mathcal G_\Lambda,\qquad \omega\in \widehat{\mathcal W}_\Lambda,
\]
都有
\[
\widehat\mu_\Lambda(g\omega)=\widehat\mu_\Lambda(\omega).
\]
记
\[
\mu_\Lambda:=(\pi^\Lambda)_*\widehat\mu_\Lambda.
\]
则对每个 \(w=(w_x)_{x\in\Lambda}\in \mathcal W_\Lambda\)，都有精确条件分解
\[
\widehat\mu_\Lambda(\omega\mid \pi^\Lambda=w)
=
\prod_{x\in\Lambda}
\frac{\mathbf 1_{F_{w_x}}(\omega_x)}{d(w_x)}.
\]
特别地，给定可见构型以后，全部 hidden 变量在各站点上条件独立，并且在各自纤维内严格均匀。
\end{theorem}
\begin{proof}
固定 \(w\in \mathcal W_\Lambda\)。其整条抬升纤维为
\[
E_w:=(\pi^\Lambda)^{-1}(w)=\prod_{x\in\Lambda}F_{w_x}.
\]
考虑 \(\mathcal G_\Lambda\) 中保持 \(E_w\) 的稳定子
\[
(\mathcal G_\Lambda)_w.
\]
由于 \(\mathcal G_\Lambda\) 按位点作 sitewise 作用，而
\[
G_6^{\mathrm{bin}}\cong \prod_{u\in X_6}\mathfrak S_{d(u)},
\]
故限制到第 \(x\) 个位点后，稳定子在该位点上恰给出
\[
\mathfrak S(F_{w_x})\cong \mathfrak S_{d(w_x)}.
\]
于是
\[
(\mathcal G_\Lambda)_w\cong \prod_{x\in\Lambda}\mathfrak S(F_{w_x}),
\]
它在笛卡尔积 \(E_w=\prod_x F_{w_x}\) 上显然传递：任取 \(\omega,\omega'\in E_w\)，对每个 \(x\) 可选择
\[
\sigma_x\in \mathfrak S(F_{w_x})
\]
使得 \(\sigma_x(\omega_x)=\omega'_x\)，再把这些 \(\sigma_x\) 组装即可。

由 \(\widehat\mu_\Lambda\) 对整个 \(\mathcal G_\Lambda\) 的不变性可知，其条件分布
\[
\widehat\mu_\Lambda(\,\cdot\mid \pi^\Lambda=w)
\]
对稳定子 \((\mathcal G_\Lambda)_w\) 同样不变。由于该稳定子在有限集合 \(E_w\) 上传递，故条件分布只能是 \(E_w\) 上的常数测度，即
\[
\widehat\mu_\Lambda(\omega\mid \pi^\Lambda=w)
=
\frac{\mathbf 1_{E_w}(\omega)}{|E_w|}.
\]
再由
\[
|E_w|=\prod_{x\in\Lambda}d(w_x),
\qquad
\mathbf 1_{E_w}(\omega)=\prod_{x\in\Lambda}\mathbf 1_{F_{w_x}}(\omega_x),
\]
便得到所示公式。乘积形式立即给出条件独立与各纤维内严格均匀性。证毕。
\leanverified{paper\_xi\_time\_part9z\_window6\_canonical\_lift\_conditional\_uniformization}
\end{proof}

\begin{theorem}[三层局域熵势与源项整化]\label{thm:xi-time-part9z-window6-local-entropy-source-renormalization}
设
\[
H_\Lambda:\mathcal W_\Lambda\to \RR
\]
为任意可见 Hamiltonian，并在 \(\widehat{\mathcal W}_\Lambda\) 上定义 canonical lifted Gibbs 权
\[
\widehat\mu_\Lambda(\omega)
\propto
\exp\!\bigl(-H_\Lambda(\pi^\Lambda\omega)\bigr).
\]
则 lifted 配分函数满足精确消隐公式
\[
\widehat Z_\Lambda
=
\sum_{w\in\mathcal W_\Lambda}
\exp\!\Bigl(
-H_\Lambda(w)+\sum_{x\in\Lambda}s(w_x)
\Bigr),
\qquad
s(w):=\log d(w).
\]
因此，hidden sector 对可见理论的全部新增作用恰为纯一体项
\[
H_\Lambda^{\mathrm{eff}}(w)
=
H_\Lambda(w)-\sum_{x\in\Lambda}s(w_x).
\]
在 window-\(6\) 上，
\[
s(w)=
\begin{cases}
\log 2,& w_6=1,\\[4pt]
\log 4,& w_6=0,\ V_6(w)\le 8,\\[4pt]
\log 3,& w_6=0,\ V_6(w)>8.
\end{cases}
\]

更一般地，对任意局部 observable \(h:\Omega_6\to\RR\) 与任意实源
\[
t=(t_x)_{x\in\Lambda}\in\RR^\Lambda,
\]
定义
\[
\widehat Z_\Lambda(t)
:=
\sum_{\omega\in\widehat{\mathcal W}_\Lambda}
\exp\!\Bigl(
-H_\Lambda(\pi^\Lambda\omega)+\sum_{x\in\Lambda}t_x h(\omega_x)
\Bigr).
\]
再定义
\[
M_h(t;w):=
\frac1{d(w)}\sum_{n\in F_w}e^{t h(n)}.
\]
则 \(M_h(\,\cdot\,;w)\) 对 \(t\) 为整函数，并且对实 \(t\) 可写成
\[
K_h(t;w):=\log M_h(t;w),
\]
其中 \(K_h(\,\cdot\,;w)\) 在实轴上实解析，并在 \(t=0\) 的某复邻域内可取解析对数支。此时有
\[
\widehat Z_\Lambda(t)
=
\sum_{w\in\mathcal W_\Lambda}
\exp\!\Bigl(
-H_\Lambda(w)+\sum_{x\in\Lambda}\bigl[s(w_x)+K_h(t_x;w_x)\bigr]
\Bigr).
\]
因此，消去 hidden 变量以后，不会新增任何二体、长程或非局域项；hidden sector 只留下显式的一体熵势与一体解析源项。
\end{theorem}
\begin{proof}
先看无源情形。按可见构型分组求和可得
\[
\widehat Z_\Lambda
=
\sum_{\omega\in\widehat{\mathcal W}_\Lambda}
e^{-H_\Lambda(\pi^\Lambda\omega)}
=
\sum_{w\in\mathcal W_\Lambda}
e^{-H_\Lambda(w)}\,
\bigl|(\pi^\Lambda)^{-1}(w)\bigr|.
\]
而
\[
\bigl|(\pi^\Lambda)^{-1}(w)\bigr|
=
\prod_{x\in\Lambda}|F_{w_x}|
=
\prod_{x\in\Lambda}d(w_x)
=
\exp\!\Bigl(\sum_{x\in\Lambda}s(w_x)\Bigr),
\]
故
\[
\widehat Z_\Lambda
=
\sum_{w\in\mathcal W_\Lambda}
\exp\!\Bigl(
-H_\Lambda(w)+\sum_{x\in\Lambda}s(w_x)
\Bigr).
\]
于是 hidden sector 对有效 Hamiltonian 的唯一新增贡献便是
\[
-\sum_{x\in\Lambda}s(w_x).
\]

在 window-\(6\) 上，\(s(w)\) 的三层公式直接来自
\[
d(w)=2\iff w_6=1,
\qquad
d(w)=4\iff w_6=0,\ V_6(w)\le 8,
\qquad
d(w)=3\iff w_6=0,\ V_6(w)>8.
\]

再看有源情形。同样按 \(w=\pi^\Lambda(\omega)\) 分组，
\[
\widehat Z_\Lambda(t)
=
\sum_{w\in\mathcal W_\Lambda}
e^{-H_\Lambda(w)}
\sum_{\omega:\,\pi^\Lambda\omega=w}
\exp\!\Bigl(\sum_{x\in\Lambda}t_x h(\omega_x)\Bigr).
\]
由于纤维分解为直积，
\[
(\pi^\Lambda)^{-1}(w)=\prod_{x\in\Lambda}F_{w_x},
\]
内层和相应分裂为
\[
\sum_{\omega:\,\pi^\Lambda\omega=w}
e^{\sum_x t_x h(\omega_x)}
=
\prod_{x\in\Lambda}\sum_{n\in F_{w_x}}e^{t_x h(n)}.
\]
写成
\[
\sum_{n\in F_{w_x}}e^{t_x h(n)}
=
d(w_x)\,M_h(t_x;w_x),
\]
便得到
\[
\widehat Z_\Lambda(t)
=
\sum_{w\in\mathcal W_\Lambda}
\exp\!\Bigl(
-H_\Lambda(w)+\sum_{x\in\Lambda}\bigl[s(w_x)+K_h(t_x;w_x)\bigr]
\Bigr).
\]
最后，\(M_h(\,\cdot\,;w)\) 是有限个指数函数之和，故为整函数；对实 \(t\)，其值严格为正，因此 \(K_h(\,\cdot\,;w)\) 在实轴上实解析，并在 \(t=0\) 的某复邻域内可取解析对数支。证毕。
\end{proof}

\leanverified{paper\_xi\_time\_part9z\_window6\_local\_entropy\_source\_renormalization}
\leanverified{paper\_xi\_time\_part9z\_window6\_isolated\_site\_annihilation}
\begin{theorem}[孤立站点湮没律]\label{thm:xi-time-part9z-window6-isolated-site-annihilation}
对任意 \(f:\Omega_6\to\RR\)，定义其纤维平均
\[
\bar f(w):=
\frac1{d(w)}\sum_{n\in F_w}f(n)
\]
以及纤维中心化部分
\[
f^\circ(n):=f(n)-\bar f(\pi(n)).
\]
则在任意 \(\mathcal G_\Lambda\)-不变 lift \(\widehat\mu_\Lambda\) 下，若
\[
A:\mathcal W_\Lambda\to \RR
\]
为任意有界可见 observable，并且在位点列
\[
x_1,\dots,x_m\in \Lambda
\]
中至少存在一个位点只出现一次，则
\[
\mathbf E_{\widehat\mu_\Lambda}\!\left[
A(\pi^\Lambda)\prod_{j=1}^m f^\circ(\omega_{x_j})
\right]=0.
\]
\end{theorem}
\begin{proof}
取一处在 \((x_1,\dots,x_m)\) 中恰出现一次的位点 \(x_\ast\)。记
\[
\mathcal F:=
\sigma\!\bigl(\pi^\Lambda,(\omega_y)_{y\neq x_\ast}\bigr).
\]
由定理 \ref{thm:xi-time-part9z-window6-canonical-lift-conditional-uniformization}，给定 \(\pi^\Lambda\) 以后，各站点 hidden 变量条件独立，且 \(\omega_{x_\ast}\) 在
\[
F_{(\pi^\Lambda\omega)_{x_\ast}}
\]
上严格均匀。因此
\[
\mathbf E_{\widehat\mu_\Lambda}\!\bigl[f^\circ(\omega_{x_\ast})\mid \mathcal F\bigr]
=
\frac1{d(w_{x_\ast})}\sum_{n\in F_{w_{x_\ast}}}
\bigl(f(n)-\bar f(w_{x_\ast})\bigr)
=0,
\]
其中
\[
w:=\pi^\Lambda(\omega).
\]
另一方面，因 \(x_\ast\) 仅出现一次，
\[
A(\pi^\Lambda)\prod_{j\ne j_\ast}f^\circ(\omega_{x_j})
\]
是 \(\mathcal F\)-可测。故由条件期望塔性质，
\[
\mathbf E_{\widehat\mu_\Lambda}\!\left[
A(\pi^\Lambda)\prod_{j=1}^m f^\circ(\omega_{x_j})
\right]
=0.
\]
证毕。
\end{proof}

\begin{theorem}[异点 \texorpdfstring{$n$}{n}-点函数的完全可见化]\label{thm:xi-time-part9z-window6-distinct-point-complete-visibility}
\leanverified{paper\_xi\_time\_part9z\_window6\_distinct\_point\_complete\_visibility}
设
\[
x_1,\dots,x_n\in \Lambda
\]
两两不同，且
\[
F_1,\dots,F_n:\Omega_6\to\RR
\]
为任意局部 observable。记
\[
\bar F_j(w):=\frac1{d(w)}\sum_{m\in F_w}F_j(m).
\]
则在任意 \(\mathcal G_\Lambda\)-不变 lift 下，都有
\[
\mathbf E_{\widehat\mu_\Lambda}\!\left[\prod_{j=1}^n F_j(\omega_{x_j})\right]
=
\mathbf E_{\mu_\Lambda}\!\left[\prod_{j=1}^n \bar F_j(w_{x_j})\right].
\]
特别地，一切异点截断相关函数都已完全包含在可见 sector 之中；hidden sector 对任何真正的长程 \(n\)-点结构都没有独立贡献。
\end{theorem}
\begin{proof}
把每个 \(F_j\) 分解为
\[
F_j=\bar F_j\circ \pi + F_j^\circ,
\qquad
F_j^\circ:=F_j-\bar F_j\circ\pi.
\]
将乘积
\[
\prod_{j=1}^n F_j(\omega_{x_j})
\]
展开为 \(2^n\) 项。若某一项含有至少一个中心化因子 \(F_j^\circ(\omega_{x_j})\)，由于 \(x_1,\dots,x_n\) 两两不同，该项中必存在某个位点只出现一次，于是由定理 \ref{thm:xi-time-part9z-window6-isolated-site-annihilation}，其期望为零。

因此唯一幸存项只能是纯可见项
\[
\prod_{j=1}^n \bar F_j(\pi(\omega_{x_j}))
=
\prod_{j=1}^n \bar F_j(w_{x_j}),
\]
对 \(\widehat\mu_\Lambda\) 取期望便得到
\[
\mathbf E_{\widehat\mu_\Lambda}\!\left[\prod_{j=1}^n F_j(\omega_{x_j})\right]
=
\mathbf E_{\mu_\Lambda}\!\left[\prod_{j=1}^n \bar F_j(w_{x_j})\right].
\]
证毕。
\end{proof}

\noindent
尤其，对由表 \ref{tab:fold6_b3c3_root_dictionary_b3} 与表 \ref{tab:fold6_b3c3_root_dictionary_c3} 所诱导并经 \(\pi\) 拉回到 \(\Omega_6\) 的任意 root-valued local observable，任意两两不同位点上的 \(n\)-点函数在抬升前后完全一致；因此，\(B_3/C_3\) 的长程相关结构已经在 visible quotient 中封闭。

\begin{theorem}[隐藏部门的严格局域性]\label{thm:xi-time-part9z-window6-hidden-sector-strict-locality}
\leanverified{paper\_xi\_time\_part9z\_window6\_hidden\_sector\_strict\_locality}
设
\[
h^\circ:\Omega_6\to\RR
\]
有界，并且在每条纤维上均值为零，即
\[
\frac1{d(w)}\sum_{n\in F_w}h^\circ(n)=0
\qquad
(w\in X_6).
\]
则在任意 \(\mathcal G_\Lambda\)-不变 lift 下成立：
\begin{enumerate}
  \item 对任意 \(x\neq y\) 都有
  \[
  \operatorname{Cov}_{\widehat\mu_\Lambda}\!\bigl(h^\circ(\omega_x),h^\circ(\omega_y)\bigr)=0.
  \]
  \item 令 \(\Lambda_\delta\subset D\cap \delta\ZZ^2\) 逼近有界区域 \(D\)，并定义面积型测试函数配对
  \[
  \Phi_\delta^\circ(\varphi)
  :=
  \delta^2\sum_{x\in\Lambda_\delta}\varphi(x)\,h^\circ(\omega_x),
  \qquad
  \varphi\in C_c^\infty(D).
  \]
  则
  \[
  \Var\!\bigl(\Phi_\delta^\circ(\varphi)\bigr)
  \le
  \|\varphi\|_\infty^2\|h^\circ\|_\infty^2
  \bigl(|D|\,\delta^2+o(\delta^2)\bigr),
  \]
  从而
  \[
  \Phi_\delta^\circ(\varphi)\longrightarrow 0
  \qquad\text{于 }L^2.
  \]
  \item 若改用白噪声型归一化
  \[
  \Psi_\delta^\circ(\varphi)
  :=
  \delta\sum_{x\in\Lambda_\delta}\varphi(x)\,h^\circ(\omega_x),
  \]
  则对支撑互不相交的 \(\varphi,\psi\in C_c^\infty(D)\)，有
  \[
  \operatorname{Cov}\!\bigl(\Psi_\delta^\circ(\varphi),\Psi_\delta^\circ(\psi)\bigr)=0
  \]
  对充分小的 \(\delta\) 恒成立。因而任何存在的子序列极限，其协方差泛函都只能支撑在对角线上；隐藏部门至多产生接触项或局域白噪声，而绝不可能产生 Green 核型的长程场。
\end{enumerate}
\end{theorem}
\begin{proof}
先证第 \(1\) 条。由中心化条件与定理
\ref{thm:xi-time-part9z-window6-canonical-lift-conditional-uniformization}，
\[
\mathbf E_{\widehat\mu_\Lambda}\!\bigl[h^\circ(\omega_x)\mid \pi^\Lambda\bigr]=0
\]
对每个 \(x\) 都成立，因此
\[
\mathbf E_{\widehat\mu_\Lambda}\!\bigl[h^\circ(\omega_x)\bigr]=0.
\]
当 \(x\neq y\) 时，定理
\ref{thm:xi-time-part9z-window6-isolated-site-annihilation} 在 \(m=2\)、\(A\equiv 1\)、\(f=h^\circ\) 的情形立即给出
\[
\mathbf E_{\widehat\mu_\Lambda}\!\bigl[h^\circ(\omega_x)h^\circ(\omega_y)\bigr]=0.
\]
故
\[
\operatorname{Cov}_{\widehat\mu_\Lambda}\!\bigl(h^\circ(\omega_x),h^\circ(\omega_y)\bigr)=0.
\]

再证第 \(2\) 条。由第 \(1\) 条，面积型配对的方差化为纯对角和：
\[
\Var\!\bigl(\Phi_\delta^\circ(\varphi)\bigr)
=
\delta^4\sum_{x\in\Lambda_\delta}\varphi(x)^2
\Var\!\bigl(h^\circ(\omega_x)\bigr).
\]
由于
\[
\Var\!\bigl(h^\circ(\omega_x)\bigr)\le \|h^\circ\|_\infty^2,
\qquad
\varphi(x)^2\le \|\varphi\|_\infty^2,
\]
并且
\[
|\Lambda_\delta|=|D|\,\delta^{-2}+o(\delta^{-2}),
\]
便得到
\[
\Var\!\bigl(\Phi_\delta^\circ(\varphi)\bigr)
\le
\delta^4|\Lambda_\delta|\|\varphi\|_\infty^2\|h^\circ\|_\infty^2
\le
\|\varphi\|_\infty^2\|h^\circ\|_\infty^2
\bigl(|D|\,\delta^2+o(\delta^2)\bigr).
\]
于是 \(\Phi_\delta^\circ(\varphi)\to 0\) 于 \(L^2\)。

最后看第 \(3\) 条。若 \(\operatorname{supp}\varphi\cap \operatorname{supp}\psi=\varnothing\)，则对充分小的 \(\delta\)，任何 \(x\in\Lambda_\delta\) 都不会同时落在两者支撑中。因此
\[
\operatorname{Cov}\!\bigl(\Psi_\delta^\circ(\varphi),\Psi_\delta^\circ(\psi)\bigr)
=
\delta^2\sum_{x,y\in\Lambda_\delta}
\varphi(x)\psi(y)\,
\operatorname{Cov}\!\bigl(h^\circ(\omega_x),h^\circ(\omega_y)\bigr)
\]
中，所有 \(x\neq y\) 项由第 \(1\) 条消失，而 \(x=y\) 的对角项又因支撑不交而为零，故总协方差恒为零。由此，任何可能的子序列极限都只能带有对角支撑协方差，也就是接触项或局域白噪声型行为。证毕。
\end{proof}

\endinput

在上述三原子 Stieltjes 描述已经固定之后，involutive algebra、有限热力学与一般窗口刚性之间的闭合关系还可进一步收束为如下推论与猜想。

\begin{corollary}[window-\texorpdfstring{$6$}{6} 的有限信息热力学由 involutive algebra 单独闭合]\label{cor:xi-time-part9z-window6-involutive-thermodynamics-closure}
\leanverified{paper\_xi\_time\_part9z\_window6\_involutive\_thermodynamics\_closure}
若已知输入仅为 involutive algebra \((\mathcal A_6,\alpha)\) 的对象层数据
\[
\dim Z(\mathcal A_6),\qquad
\dim \mathcal A_6^{\alpha,-},\qquad
\dim \mathcal A_6,
\]
则下列对象全部由它唯一决定：
\[
\mu_6=8\delta_2+4\delta_3+9\delta_4,
\]
\[
M_n=\int t^n\,\dd\mu_6(t)
\qquad (n\ge 0),
\]
\[
H_q(p_6)
=
\frac{1}{1-q}
\Bigl[
\log\bigl(8\cdot 2^q+4\cdot 3^q+9\cdot 4^q\bigr)-6q\log 2
\Bigr]
\qquad (q\in\RR\setminus\{1\}),
\]
\[
C_6(B)=\int \min(2^B,t)\,\dd\mu_6(t)
\qquad (B\in\ZZ_{\ge 0}),
\]
\[
g_6
=
\frac{1}{64}\int \log\Gamma(t+1)\,\dd\mu_6(t)
=
\frac{1}{64}\Bigl(8\log(2!)+4\log(3!)+9\log(4!)\Bigr),
\]
\[
\kappa_6
=
\frac{1}{64}\int t\log t\,\dd\mu_6(t)
=
\frac{1}{64}\Bigl(8\cdot 2\log 2+4\cdot 3\log 3+9\cdot 4\log 4\Bigr).
\]
换言之，window-\(6\) 的有限信息热力学不是附着在几何之上的附属层，而是 involutive groupoid algebra 的一个确定性函子像。
\end{corollary}
\begin{proof}
由定理 \ref{thm:xi-time-part9z-window6-moment-involutive-algebra-realization}，\((\mathcal A_6,\alpha)\) 先给出
\[
(M_0,M_1,M_2)=\bigl(\dim Z(\mathcal A_6),\ \dim \mathcal A_6^{\alpha,-},\ \dim \mathcal A_6\bigr).
\]
再由定理 \ref{thm:xi-time-part9z-window6-histogram-from-three-moments}，上述三元组唯一恢复 \((N_2,N_3,N_4)=(8,4,9)\)，从而唯一恢复三原子测度 \(\mu_6\)。

一旦 \(\mu_6\) 已知，定理 \ref{thm:xi-time-part9z-window6-threeatom-moment-ogf} 给出全部整数矩，定理 \ref{thm:xi-time-part9z-window6-threeatom-renyi-spectrum} 给出全部 Rényi 谱，定理 \ref{thm:xi-time-part9z-window6-stieltjes-capacity-duality} 给出全部 dyadic 容量。最后，\(g_6\) 与 \(\kappa_6\) 只是对 \(\mu_6\) 施加固定对数核所得的积分，因此同样由 \(\mu_6\) 唯一决定。证毕。
\end{proof}

\begin{conjecture}[三原子窗口的第一 moment--反不动刚性]\label{conj:xi-time-part9z-threeatom-window-firstmoment-antifixed-rigidity}
设一般窗口 \(m\) 的 canonical groupoid algebra 为 \(\mathcal A_m\)，几何对合为 \(\alpha_m\)，并记
\[
\mathcal A_m^{\alpha_m,\pm}:=\{a\in \mathcal A_m:\ \alpha_m(a)=\pm a\},
\qquad
\mu_m:=\sum_{x\in X_m}\delta_{d_m(x)}.
\]
则有理由猜想
\[
\dim \mathcal A_m^{\alpha_m,-}
=
\int t\,\dd\mu_m(t)
=
\sum_{x\in X_m}d_m(x)
=
2^m.
\]
若该猜想成立，则当 \(\mu_m\) 的支撑恰有三个点时，三元组
\[
\bigl(\dim Z(\mathcal A_m),\ \dim \mathcal A_m^{\alpha_m,-},\ \dim \mathcal A_m\bigr)
\]
必唯一决定完整纤维直方图。window-\(6\) 正是这一刚性机制的首个非平凡已审计样本。
\end{conjecture}

\noindent
上述结果表明，window-\(6\) 的深层统一并不止于根系字典或 Lie 完成。真正闭合的对象是三原子 Stieltjes 测度
\[
\mu_6=8\delta_2+4\delta_3+9\delta_4.
\]
它的 \(0/1/2\) 阶矩分别实现为中心维数 \(21\)、反不动维数 \(64\) 与全代数维数 \(212\)；相应差值 \(M_2-M_1=148\) 又精确等于有序互异碰撞对总数。与此同时，Hankel 秩 \(3\) 说明前三层代数阴影已足以反演完整纤维直方图；在概率侧，同一测度同时生成全部 output Rényi 熵与 escort 腔室相图；在资源侧，它又通过 layer-cake 与 Mellin--Stieltjes 对偶给出全部 dyadic 可逆容量；在热侧，\(g_6\) 与 \(\kappa_6\) 只不过是对该测度施加两类固定对数核所得的积分。因而，stable type、microstate、collision pair、matrix algebra 与 finite thermodynamics 在 window-\(6\) 上并不是五套并列描述，而是同一三原子 moment 载体的五个严格等价投影。
\endinput

\begin{corollary}[面积型配对下抬升观测量的塌缩]\label{cor:xi-time-part9za-window6-area-pairing-lifted-collapse}
对任意有界局部 observable \(F:\Omega_6\to\RR\)，定义
\[
\mathcal F_\delta(\varphi)
:=
\delta^2\sum_{x\in\Lambda_\delta}\varphi(x)\,F(\omega_x),
\]
\[
\overline{\mathcal F}_\delta(\varphi)
:=
\delta^2\sum_{x\in\Lambda_\delta}\varphi(x)\,\bar F(w_x),
\qquad
\bar F(w):=\frac1{d(w)}\sum_{n\in F_w}F(n).
\]
则在任意 \(\mathcal G_{\Lambda_\delta}\)-不变 lift 下，
\[
\mathcal F_\delta(\varphi)-\overline{\mathcal F}_\delta(\varphi)\longrightarrow 0
\qquad\text{于 }L^2.
\]
\end{corollary}
\begin{proof}
把 \(F\) 写成
\[
F=\bar F\circ \pi+F^\circ,
\qquad
F^\circ:=F-\bar F\circ \pi.
\]
则 \(F^\circ\) 在每条纤维上均值为零，且
\[
\mathcal F_\delta(\varphi)-\overline{\mathcal F}_\delta(\varphi)
=
\delta^2\sum_{x\in\Lambda_\delta}\varphi(x)\,F^\circ(\omega_x).
\]
右端正是定理 \ref{thm:xi-time-part9z-window6-hidden-sector-strict-locality} 第 \(2\) 条中的 \(\Phi_\delta^\circ(\varphi)\)，故其 \(L^2\)-极限为 \(0\)。证毕。
\end{proof}
\leanverified{paper\_xi\_time\_part9za\_window6\_area\_pairing\_lifted\_collapse}

\begin{theorem}[GFF 标度下的抬升不变性]\label{thm:xi-time-part9za-window6-gff-scaling-lift-invariance}
设 \(F:\Omega_6\to\RR\) 为任意有界局部 observable，并沿用推论
\ref{cor:xi-time-part9za-window6-area-pairing-lifted-collapse} 的记号。考虑一列
\[
\mathcal G_{\Lambda_\delta}\text{-不变 lifts }\widehat\mu_{\Lambda_\delta},
\qquad
\mu_{\Lambda_\delta}:=(\pi^{\Lambda_\delta})_*\widehat\mu_{\Lambda_\delta}.
\]
再设 \(a_\delta>0\) 满足
\[
\frac{a_\delta}{\delta}\longrightarrow \infty.
\]
若对任意测试函数 \(\varphi\in C_c^\infty(D)\)，可见配对满足
\[
a_\delta^{-1}
\Bigl(
\overline{\mathcal F}_\delta(\varphi)
-\mathbf E_{\mu_{\Lambda_\delta}}\,\overline{\mathcal F}_\delta(\varphi)
\Bigr)
\Longrightarrow
\Gamma(\varphi),
\]
其中 \(\Gamma\) 是某个极限随机分布的测试函数读出，则同一 observable 的任意 canonical gauge-invariant lift 也满足
\[
a_\delta^{-1}
\Bigl(
\mathcal F_\delta(\varphi)
-\mathbf E_{\widehat\mu_{\Lambda_\delta}}\,\mathcal F_\delta(\varphi)
\Bigr)
\Longrightarrow
\Gamma(\varphi).
\]
特别地，只要 visible sector 的某个局部 observable 进入 GFF 标度，任何 canonical window-\(6\) 抬升都不会改变其连续极限。
\end{theorem}
\begin{proof}
由推论 \ref{cor:xi-time-part9za-window6-area-pairing-lifted-collapse}，
\[
\left\|
\mathcal F_\delta(\varphi)-\overline{\mathcal F}_\delta(\varphi)
\right\|_{L^2}
=
O(\delta).
\]
于是
\[
\left\|
a_\delta^{-1}
\bigl(
\mathcal F_\delta(\varphi)-\overline{\mathcal F}_\delta(\varphi)
\bigr)
\right\|_{L^2}
=
O\!\left(\frac{\delta}{a_\delta}\right)
\longrightarrow 0.
\]
因此
\[
a_\delta^{-1}
\Bigl(
\mathcal F_\delta(\varphi)-\mathbf E_{\widehat\mu_{\Lambda_\delta}}\,\mathcal F_\delta(\varphi)
\Bigr)
-
a_\delta^{-1}
\Bigl(
\overline{\mathcal F}_\delta(\varphi)-\mathbf E_{\mu_{\Lambda_\delta}}\,\overline{\mathcal F}_\delta(\varphi)
\Bigr)
\longrightarrow 0
\]
于概率。由 Slutsky 定理，两者具有同一极限分布，即
\[
a_\delta^{-1}
\Bigl(
\mathcal F_\delta(\varphi)-\mathbf E_{\widehat\mu_{\Lambda_\delta}}\,\mathcal F_\delta(\varphi)
\Bigr)
\Longrightarrow
\Gamma(\varphi).
\]
证毕。
\end{proof}
\leanverified{paper\_xi\_time\_part9za\_window6\_gff\_scaling\_lift\_invariance}

\noindent
公开文献中，六顶点高度函数的去局域化与对数方差已见
\href{https://research.abo.fi/en/publications/delocalization-of-the-height-function-of-the-six-vertex-model/?utm_source=chatgpt.com}{Åbo Akademi University}；而严格的 Gaussian free field 收敛目前仍只在
\[
-1\le \Delta\le -\frac12
\]
的区间内得到，对应可参见
\href{https://arxiv.org/abs/2603.06268}{arXiv:2603.06268}。因此，对 square ice \((a=b=c=1)\) 的 window-\(6\) 连续极限这里只保持猜想表述。

\begin{conjecture}[square ice 的 window-\texorpdfstring{$6$}{6} 局域熵重整化原理]\label{conj:xi-time-part9za-square-ice-window6-local-entropy-renormalization}
设 \(\nu_\delta^{\mathrm{ice}}\) 为 square ice \((a=b=c=1)\) 在离散域 \(D_\delta\) 上诱导的 window-\(6\) 可见 Gibbs 律，\(\widehat\nu_\delta\) 为其 canonical gauge lift。则在适当边界坡度调谐下，应有：
\begin{enumerate}
  \item 由定理 \ref{thm:xi-time-part9z-window6-local-entropy-source-renormalization} 给出的显式三层熵势
  \[
  s(w)=\log d(w)\in\{\log 2,\log 3,\log 4\}
  \]
  仍应是 hidden sector 对可见有效理论的唯一新增形变；除这一体项外，不应再有任何新的非局域有效作用在临界缩放下存活；
  \item 中心化 hidden 部门在白噪声型归一化下至多留下局域接触噪声，而在面积型 GFF 配对归一化下完全消失；
  \item 因而 square-ice 的 massless fixed point 与加入 \(s(w)\) 后的 visible marginal 应属于同一 Gaussian universality class；window-\(6\) 不会改变其 GFF 协方差核，只可能改写局域 surface tension 的解析部分。
\end{enumerate}
\end{conjecture}

\noindent
与猜想 \ref{conj:xi-time-part55d-window6-sixvertex-gff-visible-hidden-separation} 相比，这里的新增内容不在于再次排除第二个无质量 Gaussian 模，而在于把 hidden sector 在连续极限中的全部可见余量唯一钉扎为定理 \ref{thm:xi-time-part9z-window6-local-entropy-source-renormalization} 的显式一体熵势 \(s=\log d\) 与定理 \ref{thm:xi-time-part9z-window6-hidden-sector-strict-locality} 所允许的对角接触噪声。定理 \ref{thm:xi-time-part9z-window6-canonical-lift-conditional-uniformization} 至定理 \ref{thm:xi-time-part9za-window6-gff-scaling-lift-invariance} 已经把除此两类机制之外的全部通道逐项排除。

\endinput

\paragraph{fold-\texorpdfstring{$\pi$}{pi} 奇层尾与共振 germ 的矩谱分岔、Hankel 精确分界与块 Jacobi 桥接猜想}
在 \(\arctan z\) 的经典 Gauss continued fraction 以及固定阶 hypergeometric/minimal-solution 原型所代表的标量 continued-fraction 谱系之外，项目内部还同时携带两类彼此独立的 \(\pi\)-相关对象：其一是 fold-\(\pi\) 通道的奇层 Bernoulli 尾，其二是 Fibonacci 共振整函数的局部 germ。前者在结论小节中已经被识别为零收敛半径的形式芽，后者在 \ref{thm:xi-time-part9v-resonance-germ-foldpi-bernoulli-free-transfer} 中已被写成与同一 Bernoulli 骨架耦合的局部 Taylor 系数。这里不再沿 \(m\)-方向考察有限窗深度，也不再沿 \(s\)-方向考察零半径误差 germ，而是把两侧都重新归一化为“系数指标 \(k\)”上的矩列，并在矩问题、Hankel 秩与 holonomicity 三侧比较其范畴类型。因而，下述 Hankel 结论并不与定理 \ref{thm:xi-logcm-hankel-vandermonde-mainterm} 的深度方向 Hankel 非退化矛盾，也不与定理 \ref{thm:conclusion-foldgauge-pi-zero-radius-odd-germ} 的零半径 germ 障碍矛盾；三者作用于彼此不同的变量方向。

\begin{theorem}[fold-\texorpdfstring{$\pi$}{pi} Bernoulli 尾的二原子矩刻画]\label{thm:xi-time-part9zb-foldpi-two-atomic-moment-characterization}
\leanverified{paper\_xi\_time\_part9zb\_foldpi\_two\_atomic\_moment\_characterization}
沿用定理 \ref{thm:xi-time-part9v-resonance-germ-foldpi-bernoulli-free-transfer} 的记号，记
\[
a_k=\frac{B_{2k}}{k(2k-1)}
\bigl(\varphi^{-1}+\varphi^{2k-3}\bigr),
\qquad
\widetilde a_k:=(-1)^{k-1}a_k>0.
\]
定义规范化矩列
\[
f_k:=
\frac{k(2k-1)}{|B_{2k}|}\,\widetilde a_k
\qquad (k\ge 1).
\]
则
\[
f_k=\varphi^{-1}+\varphi^{2k-3}
=
\int x^k\,\dd\mu_{\mathrm{fold}}(x),
\]
其中
\[
\mu_{\mathrm{fold}}
:=
\varphi^{-1}\delta_1+\varphi^{-3}\delta_{\varphi^2}.
\]
特别地，
\[
\sum_{k\ge 1}f_k z^{k-1}
=
\frac{\varphi^{-1}}{1-z}
+
\frac{\varphi^{-1}}{1-\varphi^2 z},
\]
故 \((f_k)\) 满足二阶常系数递推
\[
f_{k+2}=(1+\varphi^2)f_{k+1}-\varphi^2 f_k
\qquad (k\ge 1),
\]
并且其 Hankel 秩恰为 \(2\)。
\end{theorem}
\begin{proof}
由定理 \ref{thm:xi-time-part9v-resonance-germ-foldpi-bernoulli-free-transfer}，
\[
\widetilde a_k
=
\frac{|B_{2k}|}{k(2k-1)}
\bigl(\varphi^{-1}+\varphi^{2k-3}\bigr),
\]
故
\[
f_k=\varphi^{-1}+\varphi^{2k-3}
=
\varphi^{-1}\cdot 1^k+\varphi^{-3}\cdot (\varphi^2)^k.
\]
这正是二原子测度
\[
\mu_{\mathrm{fold}}=\varphi^{-1}\delta_1+\varphi^{-3}\delta_{\varphi^2}
\]
的 \(k\)-阶矩。

于是其生成函数立即为
\[
\sum_{k\ge 1}f_k z^{k-1}
=
\varphi^{-1}\sum_{k\ge 1}z^{k-1}
+
\varphi^{-3}\sum_{k\ge 1}\varphi^{2k}z^{k-1}
=
\frac{\varphi^{-1}}{1-z}
+
\frac{\varphi^{-1}}{1-\varphi^2 z}.
\]
分母化简为
\[
(1-z)(1-\varphi^2 z)=1-(1+\varphi^2)z+\varphi^2 z^2,
\]
故系数列满足所示二阶递推。

最后，二原子矩列的 Hankel 秩恰等于其支撑点个数，而
\[
\{1,\varphi^2\}
\]
恰含两个互异点，故 Hankel 秩恰为 \(2\)。证毕。
\end{proof}

\begin{theorem}[共振 germ 的纯点 Stieltjes 矩刻画]\label{thm:xi-time-part9zb-resonance-purepoint-stieltjes-moment-characterization}
\leanverified{paper\_xi\_time\_part9zb\_resonance\_purepoint\_stieltjes\_moment\_characterization}
沿用定理 \ref{thm:xi-time-part9v-resonance-germ-foldpi-bernoulli-free-transfer} 的记号，在 \(z=0\) 邻域写
\[
-\log \mathcal F_\varphi(z)=\sum_{k\ge 1}b_k z^{2k}.
\]
定义规范化矩列
\[
r_k:=
\frac{2k(2k)!}{2^{2k}(2^{2k}-1)|B_{2k}|\pi^{2k}}\,b_k
\qquad (k\ge 1).
\]
则
\[
r_k
=
\sum_{n=2}^{\infty}\varphi^{-2kn}
=
\frac{\varphi^{-2k}}{\varphi^{2k}-1}
=
\int x^k\,\dd\mu_{\mathrm{res}}(x),
\]
其中
\[
\mu_{\mathrm{res}}
:=
\sum_{n=2}^{\infty}\delta_{\varphi^{-2n}}
\]
是一个在 \(0\) 处聚点的纯点正测度，并且其所有正阶矩都有限。因而其 Stieltjes 生成函数为
\[
R(z):=\sum_{k\ge 1}r_k z^{k-1}
=
\sum_{n=2}^{\infty}\frac{\varphi^{-2n}}{1-\varphi^{-2n}z},
\]
并在每个
\[
z=\varphi^{2n}
\qquad (n\ge 2)
\]
处具有有限简单极点。
\end{theorem}
\begin{proof}
由定理 \ref{thm:xi-time-part9v-resonance-germ-foldpi-bernoulli-free-transfer}，
\[
b_k=
\frac{2^{2k}(2^{2k}-1)|B_{2k}|\pi^{2k}}{2k(2k)!}\,
\frac{\varphi^{-2k}}{\varphi^{2k}-1}.
\]
故归一化后立即得到
\[
r_k=\frac{\varphi^{-2k}}{\varphi^{2k}-1}.
\]
再用几何级数求和，
\[
\frac{\varphi^{-2k}}{\varphi^{2k}-1}
=
\frac{\varphi^{-4k}}{1-\varphi^{-2k}}
=
\sum_{n=2}^{\infty}\varphi^{-2kn}.
\]
因此
\[
r_k=
\sum_{n=2}^{\infty}\bigl(\varphi^{-2n}\bigr)^k
=
\int x^k\,\dd\mu_{\mathrm{res}}(x),
\qquad
\mu_{\mathrm{res}}=\sum_{n=2}^{\infty}\delta_{\varphi^{-2n}}.
\]
对每个固定 \(k\ge 1\)，该级数绝对收敛，故正阶矩都有限。

再对 \(k\) 求和，可得
\[
R(z)
=
\sum_{k\ge 1}\sum_{n=2}^{\infty}\varphi^{-2kn}z^{k-1}
=
\sum_{n=2}^{\infty}\sum_{k\ge 1}\varphi^{-2kn}z^{k-1}
=
\sum_{n=2}^{\infty}\frac{\varphi^{-2n}}{1-\varphi^{-2n}z},
\]
其中在足够小的 \(|z|\) 下交换求和次序由绝对收敛保证。于是每个分式项都在
\[
z=\varphi^{2n}
\]
给出一个简单极点，且这些极点彼此不同。证毕。
\end{proof}

\begin{theorem}[三阶 Hankel 障碍是 fold 与共振的最小精确分界]\label{thm:xi-time-part9zb-third-order-hankel-obstruction}
\leanverified{paper\_xi\_time\_part9zb\_third\_order\_hankel\_obstruction}
定义 Hankel 主子式
\[
\Delta_N^{\mathrm{fold}}
:=
\det\bigl(f_{i+j-1}\bigr)_{1\le i,j\le N},
\qquad
\Delta_N^{\mathrm{res}}
:=
\det\bigl(r_{i+j-1}\bigr)_{1\le i,j\le N}.
\]
则
\[
\Delta_1^{\mathrm{fold}}>0,
\qquad
\Delta_2^{\mathrm{fold}}=1,
\qquad
\Delta_N^{\mathrm{fold}}=0
\quad (N\ge 3),
\]
而
\[
\Delta_N^{\mathrm{res}}>0
\qquad (\forall\,N\ge 1).
\]
因此，fold-\(\pi\) 奇层尾与共振 germ 之间的首个精确不可约失配，恰发生在三阶 Hankel。等价地，任何试图以“二原子 / 二状态 / depth-two scalar”机制同时精确承载这两侧归一化数据的方案，在 \(N=3\) 时即必失败。
\end{theorem}
\begin{proof}
对 fold 侧，定理 \ref{thm:xi-time-part9zb-foldpi-two-atomic-moment-characterization} 已表明 \((f_k)\) 是二原子矩列，故其 Hankel 秩为 \(2\)。因此
\[
\Delta_N^{\mathrm{fold}}=0
\qquad (N\ge 3).
\]
又因
\[
\Delta_1^{\mathrm{fold}}=f_1=\frac{2}{\varphi}>0,
\]
而对二原子测度
\[
\mu_{\mathrm{fold}}=w_1\delta_{x_1}+w_2\delta_{x_2}
\]
有标准公式
\[
\Delta_2^{\mathrm{fold}}=w_1w_2x_1x_2(x_2-x_1)^2.
\]
在这里
\[
w_1=\varphi^{-1},\qquad
w_2=\varphi^{-3},\qquad
x_1=1,\qquad
x_2=\varphi^2,
\]
故
\[
\Delta_2^{\mathrm{fold}}
=
\varphi^{-1}\varphi^{-3}\varphi^2(\varphi^2-1)^2
=
\varphi^{-2}\varphi^2
=1,
\]
其中用到了
\[
\varphi^2-1=\varphi.
\]

对共振侧，任取 \(N\ge 1\) 与非零向量
\[
c=(c_1,\dots,c_N)^{\mathsf T}\in\RR^N.
\]
定义多项式
\[
P(x):=\sum_{j=1}^{N}c_j x^{j-1}.
\]
则由定理 \ref{thm:xi-time-part9zb-resonance-purepoint-stieltjes-moment-characterization}，
\[
c^{\mathsf T}\bigl(r_{i+j-1}\bigr)_{1\le i,j\le N}c
=
\sum_{n=2}^{\infty}\varphi^{-2n}P(\varphi^{-2n})^2.
\]
由于
\[
\{\varphi^{-2n}:n\ge 2\}
\]
是一个无穷互异点集，而非零多项式 \(P\) 只能有有限个零点，故上式严格大于 \(0\)。于是每个有限阶 Hankel 主块都正定，从而
\[
\Delta_N^{\mathrm{res}}>0
\qquad (\forall\,N\ge 1).
\]
证毕。
\end{proof}

\begin{theorem}[holonomic 分岔：fold 有限阶而 resonance 非 \texorpdfstring{$D$}{D}-有限]\label{thm:xi-time-part9zb-holonomic-bifurcation-fold-vs-resonance}
在定理 \ref{thm:xi-time-part9zb-foldpi-two-atomic-moment-characterization} 与定理 \ref{thm:xi-time-part9zb-resonance-purepoint-stieltjes-moment-characterization} 的记号下，定义
\[
F_{\mathrm{fold}}(z):=\sum_{k\ge 1}f_k z^{k-1},
\qquad
R(z):=\sum_{k\ge 1}r_k z^{k-1}.
\]
则 \(F_{\mathrm{fold}}(z)\) 是有理函数，因而 \(D\)-有限；其最小线性递推阶恰为 \(2\)。另一方面，
\[
R(z)=\sum_{n=2}^{\infty}\frac{\varphi^{-2n}}{1-\varphi^{-2n}z}
\]
不是 \(D\)-有限；特别地，系数列 \((r_k)_{k\ge 1}\) 不是 \(P\)-递归序列。
\end{theorem}
\begin{proof}
前半部分已由定理 \ref{thm:xi-time-part9zb-foldpi-two-atomic-moment-characterization} 给出：\(F_{\mathrm{fold}}\) 的分母为
\[
(1-z)(1-\varphi^2 z),
\]
故其最小递推阶恰为 \(2\)。

对后半部分，若 \(R(z)\) 为 \(D\)-有限函数，则与定理 \ref{thm:fold-gauge-anomaly-singularity-lattice-nonholonomic} 的证明中所用有限奇点判据相同，\(R\) 在复平面上的有限奇点至多有限多个；关于这一判据的一般背景，亦可参见 \href{https://algo.inria.fr/flajolet/Publications/FlGeSa05.pdf?utm_source=chatgpt.com}{Flajolet--Gerhold--Salvy}。另一方面，定理 \ref{thm:xi-time-part9zb-resonance-purepoint-stieltjes-moment-characterization} 已表明 \(R(z)\) 在每个
\[
z=\varphi^{2n}
\qquad (n\ge 2)
\]
处都具有简单极点，故其有限奇点集为无穷集合。这一矛盾说明 \(R(z)\) 非 \(D\)-有限。

最后，对 ordinary generating function，\(D\)-有限性与系数列 \(P\)-递归性是标准等价命题；故 \((r_k)\) 亦非 \(P\)-递归。证毕。
\end{proof}

\begin{corollary}[固定阶 scalar continued-fraction / minimal-solution 范式的单侧相容性]\label{cor:xi-time-part9zb-fixed-order-scalar-cf-one-sided-compatibility}
在 classical Gauss continued fraction for \(\arctan z\) 的标量核、以及固定阶 hypergeometric/minimal-solution 原型所代表的 scalar continued-fraction 范式内，项目内部只有定理 \ref{thm:xi-time-part9zb-foldpi-two-atomic-moment-characterization} 的 fold 规范化侧与该范式相容；定理 \ref{thm:xi-time-part9zb-resonance-purepoint-stieltjes-moment-characterization} 的共振规范化侧不可能属于同一 scalar 范式。参见 \href{https://dlmf.nist.gov/3.10?utm_source=chatgpt.com}{DLMF 3.10} 与 \cite{Wang2026RamanujanMachinePi}。
\end{corollary}
\begin{proof}
经典 \(\arctan z\) 的 Gauss continued fraction 对应一个 \(D\)-有限的标量函数核；而 \cite{Wang2026RamanujanMachinePi} 所覆盖的 Ramanujan Machine 型固定阶原型同样归约到 fixed-order 的 hypergeometric/minimal-solution 标量框架。在此处所采用的 coefficient-direction ordinary generating function 口径下，此类固定阶 scalar 核都落在 holonomic 范畴。

由定理 \ref{thm:xi-time-part9zb-holonomic-bifurcation-fold-vs-resonance}，fold 规范化侧的
\[
F_{\mathrm{fold}}(z)
\]
是二极点有理函数，故不仅 \(D\)-有限，而且处在最小二阶情形；与之相对，共振规范化侧的
\[
R(z)
\]
不是 \(D\)-有限，因而不可能属于同一 scalar fixed-order 范式。

需要指出的是，这一“单侧相容性”仅讨论定理 \ref{thm:xi-time-part9zb-foldpi-two-atomic-moment-characterization} 的系数方向规范化生成函数 \(F_{\mathrm{fold}}\)，并不与定理 \ref{thm:conclusion-foldgauge-pi-zero-radius-odd-germ} 讨论的零收敛半径误差 germ \(\widehat E(s)\) 相冲突。前者是对 odd Bernoulli jet 作矩归一化后的 \(k\)-方向对象，后者则是 \(s\)-方向的局部解析 germ；二者属于不同变量口径。证毕。
\end{proof}

\begin{conjecture}[完整 \texorpdfstring{$\pi$}{pi} 桥接的唯一候选必须是块 Jacobi 或算子值 continued fraction]\label{conj:xi-time-part9zb-pi-bridge-unique-block-jacobi-candidate}
任何试图把项目内部完整的 \(\pi\)-机制压缩为单一 scalar continued fraction 的方案都应失败。若严格桥接存在，则其最低可能结构应是一个块 Jacobi 或算子值 continued fraction：存在某个块三对角算子
\[
\mathcal J_\varphi
\]
及其 Weyl--Titchmarsh 函数
\[
M_\varphi(z),
\]
使得在适当分块下：
\[
\text{有限秩块}
\quad\Longleftrightarrow\quad
\text{定理 \ref{thm:xi-time-part9zb-foldpi-two-atomic-moment-characterization} 的二原子矩结构},
\]
\[
\text{无限纯点块}
\quad\Longleftrightarrow\quad
\text{定理 \ref{thm:xi-time-part9zb-resonance-purepoint-stieltjes-moment-characterization} 的纯点 Stieltjes 结构}.
\]
更具体地，\(M_\varphi(z)\) 的某一分量应退化为定理 \ref{thm:xi-time-part9zb-foldpi-two-atomic-moment-characterization} 的二极点有理函数，另一分量应退化为定理 \ref{thm:xi-time-part9zb-resonance-purepoint-stieltjes-moment-characterization} 的纯点 Stieltjes 变换，而常数 \(2\pi\) 则作为两块之间唯一允许的归一化耦合常数出现。
\end{conjecture}
\begin{proof}[证据]
定理 \ref{thm:xi-time-part9zb-foldpi-two-atomic-moment-characterization} 至推论 \ref{cor:xi-time-part9zb-fixed-order-scalar-cf-one-sided-compatibility} 已经给出四重排除链。其一，fold 规范化侧的 Hankel 秩仅为 \(2\)；其二，共振规范化侧的全部有限 Hankel 主子式均严格为正，因而保持全阶正定；其三，fold 侧生成函数是有理的，而共振侧生成函数非 \(D\)-有限；其四，公开文献中当前可直接调用的 continued-fraction/minimal-solution 原型仍停留在 fixed-order scalar 核的范围内。

因此，若要在单一谱对象中同时容纳这两侧数据而不在三阶 Hankel 处立即崩塌，就不能继续停留在 scalar architecture。二原子有限秩与无限纯点正定这两种矩谱类型的并置，强制最低结构成本从 scalar 核跃迁到块化或算子值核。换言之，这里的“块化”不是表述上的加强，而是 rank-\(2\)/rank-\(\infty\) 分岔在矩谱层面的最低代价实现。证毕。
\end{proof}

\noindent
上述六条结果把项目内部两类 \(\pi\)-相关数据的差异提升为一条范畴级分界：fold-\(\pi\) 奇层尾在系数方向上归约为二原子、有限秩、holonomic 的矩谱；共振 germ 在同一方向上则归约为无限纯点、全阶正定、非 holonomic 的矩谱。于是，真正决定“何种 continued-fraction/Pad\'e/minimal-solution 架构可以承载该数据”的，不再只是常数 \(\pi\) 本身，而是其背后矩列的 Hankel 秩、极点几何与 holonomicity 类型。若未来存在统一桥接，它便不能继续停留在 scalar Pad\'e/continued-fraction 的边界之内，而必须进入块 Jacobi 或算子值 continued fraction 的层级。

\endinput

\paragraph{fold-\texorpdfstring{$\pi$}{pi} 奇层 tail 的黄金双尺度 Stirling--Borel 几何与共尺度刚性}
在推论 \ref{cor:xi-time-part9zb-fixed-order-scalar-cf-one-sided-compatibility} 已经把经典 \(\arctan z\) 的 Gauss continued fraction 与 Ramanujan Machine 所覆盖的 contiguous \({}_2F_1\) 比值统一归入 fixed-order scalar 机制，而定理 \ref{thm:conclusion-foldgauge-pi-zero-radius-odd-germ} 又把 fold-\(\pi\) 奇层尾钉为零收敛半径 Bernoulli formal germ 之后，还可从 Stirling 余项与 Borel 平面几何重新观察同一对象。由此得到的并不是另一条单核标量近似，而是一条由两条无理相关尺度同时承载的 Stirling--Gamma/Borel 对象；参见 \href{https://dlmf.nist.gov/3.10}{DLMF 3.10}、\href{https://dlmf.nist.gov/5.11}{DLMF 5.11} 与 \cite{Wang2026RamanujanMachinePi}。以下记经典 Stirling 余项为
\[
\mu_{\mathrm{St}}(z):=
\log\Gamma(z)-\Bigl(z-\frac12\Bigr)\log z+z-\frac12\log(2\pi),
\]
则
\[
\mu_{\mathrm{St}}(z)\sim \sum_{r\ge 1}\frac{B_{2r}}{2r(2r-1)}\,z^{-(2r-1)}
\qquad (z\to\infty),
\]
并沿用定理 \ref{thm:xi-time-part9vk-no-finite-local-analytic-transform-from-resonance-germ-to-foldpi-tail} 的 odd formal tail
\[
\mathcal E_\varphi(s):=
\sum_{r\ge 1}
\frac{B_{2r}}{r(2r-1)}
\bigl(\varphi^{-1}+\varphi^{2r-3}\bigr)s^{2r-1}.
\]

\begin{theorem}[有限奇矩到有限 Stirling 和的 formal 转译]\label{thm:xi-time-part9zbb-finite-stirling-sum-translation}
\leanverified{paper\_xi\_time\_part9zbb\_finite\_stirling\_sum\_translation}
设 \(m\ge 1\)、\(\lambda_1,\dots,\lambda_m>0\)、\(w_1,\dots,w_m\in\CC\)，并令
\[
c_r=\sum_{j=1}^{m}w_j\lambda_j^{2r-1}
\qquad (r\ge 1).
\]
定义相应的 Bernoulli 奇层 formal series
\[
\mathcal E_c(s):=
\sum_{r\ge 1}\frac{B_{2r}}{r(2r-1)}\,c_r\,s^{2r-1}.
\]
则在 \(s=0\) 的 formal 渐近意义下有严格恒等式
\[
\mathcal E_c(s)\sim 2\sum_{j=1}^{m}w_j\,\mu_{\mathrm{St}}\!\bigl((\lambda_j s)^{-1}\bigr).
\]
\end{theorem}
\begin{proof}
由 Stirling 级数，
\[
\mu_{\mathrm{St}}\!\bigl((\lambda_j s)^{-1}\bigr)
\sim
\sum_{r\ge 1}\frac{B_{2r}}{2r(2r-1)}(\lambda_j s)^{2r-1}.
\]
两边乘以 \(2w_j\) 并对 \(j\) 求和，得到
\[
2\sum_{j=1}^{m}w_j\,\mu_{\mathrm{St}}\!\bigl((\lambda_j s)^{-1}\bigr)
\sim
\sum_{r\ge 1}
\frac{B_{2r}}{r(2r-1)}
\left(\sum_{j=1}^{m}w_j\lambda_j^{2r-1}\right)s^{2r-1}.
\]
右端按 \(c_r\) 的定义正是 \(\mathcal E_c(s)\)。证毕。
\end{proof}

\begin{corollary}[fold-\texorpdfstring{$\pi$}{pi} 奇层的黄金双尺度 Stirling 分解]\label{cor:xi-time-part9zbb-foldpi-golden-twoscale-stirling}
\leanverified{paper\_xi\_time\_part9zbb\_foldpi\_golden\_twoscale\_stirling}
fold-\(\pi\) 奇层 tail 满足
\[
\mathcal E_\varphi(s)
\sim
2\varphi^{-1}\mu_{\mathrm{St}}(s^{-1})
+
2\varphi^{-2}\mu_{\mathrm{St}}\!\bigl((\varphi s)^{-1}\bigr).
\]
特别地，记
\[
q:=\frac{\varphi}{2},
\]
又因定理 \ref{thm:conclusion-foldgauge-pi-zero-radius-odd-germ} 中的
\(\widehat E(s)\) 与此处 \(\mathcal E_\varphi(s)\) 逐项相同，故
\[
\log C_m-\log(2\pi)\sim \mathcal E_\varphi(q^m),
\]
可得
\[
\log C_m-\log(2\pi)
\sim
2\varphi^{-1}\mu_{\mathrm{St}}\!\left(\left(\frac{2}{\varphi}\right)^m\right)
+
2\varphi^{-2}\mu_{\mathrm{St}}\!\left(\frac{2^m}{\varphi^{m+1}}\right).
\]
\end{corollary}
\begin{proof}
注意到
\[
\varphi^{-1}+\varphi^{2r-3}
=
\varphi^{-1}\cdot 1^{2r-1}
+
\varphi^{-2}\cdot \varphi^{2r-1}.
\]
因此可将定理 \ref{thm:xi-time-part9zbb-finite-stirling-sum-translation} 应用于
\[
(w_1,\lambda_1)=(\varphi^{-1},1),
\qquad
(w_2,\lambda_2)=(\varphi^{-2},\varphi),
\]
便得第一式。再把 \(s=q^m=(\varphi/2)^m\) 代入，即得第二式。证毕。
\end{proof}

\begin{theorem}[fold-\texorpdfstring{$\pi$}{pi} 奇层的 Borel 闭式与双晶格奇点几何]\label{thm:xi-time-part9zbb-foldpi-borel-double-lattice}
对 odd formal series 定义 Borel 变换
\[
\mathcal B\!\left[\sum_{r\ge 1}a_r s^{2r-1}\right](\xi)
:=
\sum_{r\ge 1}a_r\frac{\xi^{2r-2}}{(2r-2)!}.
\]
再定义
\[
H(\xi):=\frac{\xi\coth(\xi/2)-2}{\xi^2}.
\]
则
\[
\widehat{\mathcal E}_\varphi(\xi):=\mathcal B[\mathcal E_\varphi](\xi)
=
\varphi^{-1}\bigl(H(\xi)+H(\varphi\xi)\bigr).
\]
因此 \(\widehat{\mathcal E}_\varphi\) 是亚纯函数，其奇点集合恰为
\[
\Omega_\varphi=
2\pi i\,\ZZ_{\neq 0}
\ \cup\
\frac{2\pi i}{\varphi}\,\ZZ_{\neq 0},
\]
并且由于 \(\varphi\notin\QQ\)，这两条晶格在 \(0\) 之外互不相交。相应留数为
\[
\operatorname*{Res}_{\xi=2\pi i n}\widehat{\mathcal E}_\varphi(\xi)
=
\frac{\varphi^{-1}}{\pi i\,n},
\qquad
\operatorname*{Res}_{\xi=2\pi i n/\varphi}\widehat{\mathcal E}_\varphi(\xi)
=
\frac{\varphi^{-2}}{\pi i\,n}
\qquad (n\in\ZZ_{\neq 0}).
\]
特别地，\(\widehat{\mathcal E}_\varphi\) 在正实轴上无奇点，且
\[
\widehat{\mathcal E}_\varphi(\xi)=O(\xi^{-1})
\qquad (\xi\to+\infty,\ \xi\in\RR_+),
\]
故 \(\mathcal E_\varphi\) 沿 \(\RR_+\) 可 Borel 求和。
\end{theorem}
\begin{proof}
由 \href{https://dlmf.nist.gov/24.2}{DLMF 24.2} 中的 Bernoulli 生成函数，
\[
\frac{\xi}{2}\coth\!\Bigl(\frac{\xi}{2}\Bigr)
=
1+\sum_{r\ge 1}B_{2r}\frac{\xi^{2r}}{(2r)!},
\]
故
\[
H(\xi)
=
\sum_{r\ge 1}2B_{2r}\frac{\xi^{2r-2}}{(2r)!}.
\]
另一方面，
\[
\widehat{\mathcal E}_\varphi(\xi)
=
\sum_{r\ge 1}
\frac{B_{2r}}{r(2r-1)}
\bigl(\varphi^{-1}+\varphi^{2r-3}\bigr)
\frac{\xi^{2r-2}}{(2r-2)!}.
\]
利用
\[
\frac{1}{r(2r-1)(2r-2)!}=\frac{2}{(2r)!},
\]
化简得到
\[
\widehat{\mathcal E}_\varphi(\xi)
=
\varphi^{-1}\sum_{r\ge 1}2B_{2r}\frac{\xi^{2r-2}}{(2r)!}
+
\varphi^{-1}\sum_{r\ge 1}2B_{2r}\frac{(\varphi\xi)^{2r-2}}{(2r)!},
\]
也即
\[
\widehat{\mathcal E}_\varphi(\xi)=\varphi^{-1}\bigl(H(\xi)+H(\varphi\xi)\bigr).
\]

再看奇点。由于 \(\coth z\) 在 \(z=\pi i n\) 处具有留数 \(1\)，函数 \(\coth(\xi/2)\) 在 \(\xi=2\pi i n\) 处具有留数 \(2\)。故 \(H(\xi)\) 在 \(\xi=2\pi i n\) 处有简单极点，且
\[
\operatorname*{Res}_{\xi=2\pi i n}H(\xi)
=
\frac{2\cdot 2\pi i n}{(2\pi i n)^2}
=
\frac{1}{\pi i\,n}.
\]
于是第一族留数为
\[
\operatorname*{Res}_{\xi=2\pi i n}\widehat{\mathcal E}_\varphi(\xi)
=
\frac{\varphi^{-1}}{\pi i\,n}.
\]
第二族来自 \(H(\varphi\xi)\) 的缩放：若 \(\eta_n=2\pi i n/\varphi\)，则
\[
\operatorname*{Res}_{\xi=\eta_n}H(\varphi\xi)
=
\frac{1}{\varphi}\operatorname*{Res}_{u=2\pi i n}H(u)
=
\frac{1}{\varphi\,\pi i\,n},
\]
故
\[
\operatorname*{Res}_{\xi=\eta_n}\widehat{\mathcal E}_\varphi(\xi)
=
\frac{\varphi^{-2}}{\pi i\,n}.
\]

由于 \(\varphi\notin\QQ\)，若
\[
2\pi i n=\frac{2\pi i m}{\varphi},
\]
则 \(\varphi=m/n\in\QQ\)，矛盾；故两族极点在 \(0\) 外互不相交。

最后，正实轴不与纯虚奇点相交，而当 \(\xi\to+\infty\) 时 \(\coth(\xi/2)=1+O(e^{-\xi})\)，因此
\[
H(\xi)=\frac{1}{\xi}+O(\xi^{-2}),
\qquad
H(\varphi\xi)=\frac{1}{\varphi\xi}+O(\xi^{-2}),
\]
从而 \(\widehat{\mathcal E}_\varphi(\xi)=O(\xi^{-1})\) 于 \(\RR_+\)。故其正向 Laplace 变换收敛，\(\mathcal E_\varphi\) 沿 \(\RR_+\) 可 Borel 求和。证毕。
\end{proof}

\begin{proposition}[正虚射线上的双频 Stokes 跳跃]\label{prop:xi-time-part9zbb-foldpi-stokes-two-frequency}
记 \(\mathfrak S_{\pi/2,\pm}\mathcal E_\varphi\) 为沿 Borel 平面方向 \(\arg\xi=\pi/2\) 的上下侧向 Borel--Laplace 和。则在半平面
\[
\Re(i/s)>0
\]
内，有
\[
\bigl(\mathfrak S_{\pi/2,+}-\mathfrak S_{\pi/2,-}\bigr)\mathcal E_\varphi(s)
=
-2\varphi^{-1}\log\!\bigl(1-e^{-2\pi i/s}\bigr)
-2\varphi^{-2}\log\!\bigl(1-e^{-2\pi i/(\varphi s)}\bigr),
\]
其中对数取主支。因而 fold-\(\pi\) 的 Stokes 自动同构并非单频 Fourier 因子，而是两条无理相关频率 \(1\) 与 \(\varphi^{-1}\) 的对数叠加。
\end{proposition}
\begin{proof}
由定理 \ref{thm:xi-time-part9zbb-foldpi-borel-double-lattice}，正虚轴上的全部 Borel 奇点为
\[
\omega_n=2\pi i n,
\qquad
\eta_n=\frac{2\pi i n}{\varphi}
\qquad (n\ge 1),
\]
且留数分别为
\[
\frac{\varphi^{-1}}{\pi i\,n},
\qquad
\frac{\varphi^{-2}}{\pi i\,n}.
\]
故侧向 Laplace 积分之差由留数定理给出
\[
\bigl(\mathfrak S_{\pi/2,+}-\mathfrak S_{\pi/2,-}\bigr)\mathcal E_\varphi(s)
=
2\pi i\sum_{n\ge 1}
\left(
\frac{\varphi^{-1}}{\pi i\,n}e^{-\omega_n/s}
+
\frac{\varphi^{-2}}{\pi i\,n}e^{-\eta_n/s}
\right),
\]
即
\[
=
2\varphi^{-1}\sum_{n\ge 1}\frac{e^{-2\pi i n/s}}{n}
+
2\varphi^{-2}\sum_{n\ge 1}\frac{e^{-2\pi i n/(\varphi s)}}{n}.
\]
在 \(\Re(i/s)>0\) 时，
\[
\left|e^{-2\pi i/s}\right|<1,
\qquad
\left|e^{-2\pi i/(\varphi s)}\right|<1,
\]
故两级数绝对收敛。再用恒等式
\[
\sum_{n\ge 1}\frac{x^n}{n}=-\log(1-x)
\qquad (|x|<1),
\]
便得到所示闭式。证毕。
\end{proof}

\leanverified{paper\_xi\_time\_part9zbb\_no\_rational\_coscale\_stirling\_collapse}
\begin{theorem}[有理共尺度 Stirling--Gamma 塌缩不可能性]\label{thm:xi-time-part9zbb-no-rational-coscale-stirling-collapse}
不存在有限个常数 \(c_1,\dots,c_N\in\CC\)、尺度 \(\lambda_1,\dots,\lambda_N>0\) 与收敛 germ \(A(s)\in\CC\{s\}\)，使得在 \(s=0\) 的 formal 渐近口径下
\[
\mathcal E_\varphi(s)
\sim
\sum_{j=1}^{N}c_j\,\mu_{\mathrm{St}}\!\bigl((\lambda_j s)^{-1}\bigr)
+A(s),
\]
并且所有尺度都两两有理共尺度：
\[
\frac{\lambda_i}{\lambda_j}\in\QQ
\qquad (1\le i,j\le N).
\]
特别地，\(\mathcal E_\varphi\) 不可能由任何单尺度 Stirling--Gamma germ 重建；更不可能由任何仅依赖有限次 duplication 或 Gauss multiplication 所生成的有限超叠加重建。
\end{theorem}
\begin{proof}
先把 \(A\) 替换为其奇部，不妨假设
\[
A(s)\in s\,\CC\{s^2\}.
\]
则其 Borel 变换 \(\widehat A(\xi)\) 为整函数。

另一方面，对单一尺度 \(\lambda>0\)，定理
\ref{thm:xi-time-part9zbb-finite-stirling-sum-translation} 的 \(m=1\) 情形表明
\[
2\,\mu_{\mathrm{St}}\!\bigl((\lambda s)^{-1}\bigr)
\sim
\sum_{r\ge 1}\frac{B_{2r}}{r(2r-1)}\lambda^{2r-1}s^{2r-1},
\]
故其 Borel 变换为
\[
\lambda H(\lambda\xi).
\]
于是所设分解将推出
\[
\widehat{\mathcal E}_\varphi(\xi)
=
\frac12\sum_{j=1}^{N}c_j\lambda_j H(\lambda_j\xi)+\widehat A(\xi).
\]
因此 \(\widehat{\mathcal E}_\varphi\) 的全部奇点都必须落在
\[
\bigcup_{j=1}^{N}\frac{2\pi i}{\lambda_j}\,\ZZ_{\neq 0}.
\]

现利用有理共尺度假设。任选 \(\lambda_1\) 为基准，并写
\[
\lambda_j=\lambda_1\frac{a_j}{b_j}
\qquad (a_j,b_j\in\NN,\ \gcd(a_j,b_j)=1).
\]
令
\[
L:=\operatorname{lcm}(a_1,\dots,a_N),
\qquad
\Lambda:=\lambda_1 L.
\]
则
\[
\frac{\Lambda}{\lambda_j}=L\frac{b_j}{a_j}\in\NN.
\]
从而对每个 \(j\) 都有
\[
\frac{2\pi i}{\lambda_j}\,\ZZ_{\neq 0}
\subset
\frac{2\pi i}{\Lambda}\,\ZZ_{\neq 0}.
\]
也就是说，全部 Borel 奇点都必须落在单一有理晶格
\[
\frac{2\pi i}{\Lambda}\,\ZZ_{\neq 0}
\]
上。

然而定理 \ref{thm:xi-time-part9zbb-foldpi-borel-double-lattice} 已经给出
\[
\Omega_\varphi=
2\pi i\,\ZZ_{\neq 0}
\ \cup\
\frac{2\pi i}{\varphi}\,\ZZ_{\neq 0},
\]
这两条晶格之比为 \(\varphi\notin\QQ\)，不可能同时嵌入任何单一有理晶格。矛盾。故所设分解不存在。

最后，\href{https://dlmf.nist.gov/5.5}{DLMF 5.5} 中的 duplication 与 Gauss multiplication 公式只引入整数倍尺度，因此经任意有限次此类操作得到的全部尺度仍然两两有理共尺度，故也自动落入上述排除情形。证毕。
\end{proof}

\begin{conjecture}[黄金尺度半群的桥接刚性]\label{conj:xi-time-part9zbb-golden-scale-semigroup-rigidity}
设
\[
\mathcal F_\varphi(z)=\prod_{n=2}^{\infty}\cos(\pi z\varphi^{-n})
\]
为 Fibonacci 共振整函数，而 \(\mathcal E_\varphi(s)\) 为 fold-\(\pi\) 的 Bernoulli 奇层 tail。若存在一条把 \(\mathcal F_\varphi\) 与 \(\mathcal E_\varphi\) 严格桥接起来、并与尺度变换相容的解析机制，则其最小可行乘法尺度闭包应为
\[
\varphi^{\ZZ}=\{\varphi^n:\ n\in\ZZ\}.
\]
换言之，不存在任何仅由有限个有理伸缩闭合的桥接范畴；真正可行的桥接若存在，必须尊重黄金尺度半群而非任何有限生成的有理伸缩群。
\end{conjecture}
\begin{proof}[证据]
一方面，直接把无穷乘积拆出首项便得到
\[
\mathcal F_\varphi(z)
=
\cos\!\Bigl(\frac{\pi z}{\varphi^2}\Bigr)\,
\mathcal F_\varphi\!\Bigl(\frac{z}{\varphi}\Bigr),
\]
故共振侧的内禀伸缩已由 \(\varphi^{-1}\) 锁定。另一方面，推论
\ref{cor:xi-time-part9zbb-foldpi-golden-twoscale-stirling}
把 fold-\(\pi\) 奇层尾严格分解为两条 Stirling 尺度 \(1\) 与 \(\varphi\)；定理
\ref{thm:xi-time-part9zbb-foldpi-borel-double-lattice}
又把它们升级为两条无理相关的 Borel 晶格；而定理
\ref{thm:xi-time-part9zbb-no-rational-coscale-stirling-collapse}
进一步排除了全部有限有理共尺度塌缩。

因此，若严格桥接存在，则它不可能只依赖某个有限生成的 rational scaling closure。包含 \(1\) 与 \(\varphi\) 并在放大/缩小下闭合的最小乘法结构，只能是
\[
\varphi^{\ZZ}.
\]
这正是所述刚性猜想的证据。证毕。
\end{proof}

\noindent
由此，fold-\(\pi\) 奇层 tail 的本征复杂度不再表现为单纯的收敛速度或截断误差，而表现为一条不可约的 Borel 奇点几何：它可以沿正实方向求和，却同时携带两条无理相关的纯虚晶格、对应两频 Stokes 跳跃，并拒绝一切有限有理共尺度的 Stirling--Gamma 塌缩。于是，和 classical continued fraction / hypergeometric 路线的单核标量结构不同，这里的 \(\pi\)-对象类型首先由其双尺度 Stirling 结构与双晶格 Borel 几何决定。

\endinput

\paragraph{fold-\texorpdfstring{$\pi$}{pi} 奇层 tail 的 Hausdorff 矩刚性、正 \texorpdfstring{$S$}{S}-fraction 与 Jacobi 尾猜想}
与定理 \ref{thm:xi-time-part9zb-foldpi-two-atomic-moment-characterization} 中把 Bernoulli 因子完全剥离后得到的二原子矩列不同，这里保留 odd formal tail 在 Borel 方向上的阶乘归一化与交错符号。
这样得到的对象不再是秩 \(2\) 的几何阴影，而是一个支撑于有界区间上的严格 Hausdorff moment sequence，并因而内生地产生一条 canonical scalar continued fraction。
这并不与推论 \ref{cor:xi-time-part9zb-fixed-order-scalar-cf-one-sided-compatibility} 冲突：后者排除的是 fixed-order scalar kernel；这里出现的 \(S\)-fraction 虽然唯一，但其 Jacobi 尾部并不预期落入任何 eventually periodic 或固定阶超几何轨道。

沿用推论 \ref{cor:xi-time-part9zbb-foldpi-golden-twoscale-stirling} 与定理 \ref{thm:xi-time-part9zbb-foldpi-borel-double-lattice} 的记号。
把 odd formal tail 写成
\[
\mathcal E_\varphi(s)
=
\sum_{k\ge 0}\mathfrak a_{k+1}^{(\pi)}\,s^{2k+1},
\qquad
\mathfrak a_{k+1}^{(\pi)}
:=
\frac{B_{2k+2}}{(k+1)(2k+1)}
\bigl(\varphi^{-1}+\varphi^{2k-1}\bigr).
\]
并继续记其 Borel 变换为
\[
\widehat{\mathcal E}_\varphi(\xi)
:=
\sum_{k\ge 0}
\mathfrak a_{k+1}^{(\pi)}
\frac{\xi^{2k}}{(2k)!}.
\]

\begin{theorem}[fold-\texorpdfstring{$\pi$}{pi} Borel 影子的有限正测度表象]\label{thm:xi-time-part9zbc-foldpi-borel-finite-positive-measure}
定义
\[
\rho_\ast:=\frac{\varphi^2}{4\pi^2}\in(0,1),
\]
以及有限正离散测度
\[
\nu_{\varphi}^{\mathrm{odd}}
:=
\sum_{n\ge 1}
\frac{\varphi^{-1}}{\pi^2n^2}
\left(
\delta_{(4\pi^2n^2)^{-1}}
+\delta_{\varphi^2(4\pi^2n^2)^{-1}}
\right).
\]
则 \(\nu_{\varphi}^{\mathrm{odd}}\) 支撑于
\[
\operatorname{supp}\nu_{\varphi}^{\mathrm{odd}}
\subset (0,\rho_\ast],
\]
其总质量为
\[
\nu_{\varphi}^{\mathrm{odd}}\bigl((0,\rho_\ast]\bigr)=\frac{1}{3\varphi}.
\]
并且对一切 \(\xi\in\CC\setminus \Omega_\varphi\) 成立严格恒等式
\[
\widehat{\mathcal E}_\varphi(\xi)
=
\int_{0}^{\rho_\ast}\frac{\dd\nu_{\varphi}^{\mathrm{odd}}(y)}{1+\xi^2 y}
=
4\sum_{n\ge 1}
\left(
\frac{\varphi^{-1}}{\xi^2+4\pi^2n^2}
+\frac{\varphi^{-3}}{\xi^2+4\pi^2n^2/\varphi^2}
\right).
\]
因此，H.111 所对应的 Borel 影子不只是抽象亚纯函数，而是一个有限正测度的 Hausdorff--Stieltjes 变换。
\end{theorem}
\begin{proof}
由定理 \ref{thm:xi-time-part9zbb-foldpi-borel-double-lattice}，
\[
\widehat{\mathcal E}_\varphi(\xi)=\varphi^{-1}\bigl(H(\xi)+H(\varphi\xi)\bigr),
\qquad
H(\xi)=\frac{\xi\coth(\xi/2)-2}{\xi^2}.
\]
再由 \href{https://dlmf.nist.gov/4.36}{DLMF 4.36} 中 \(\coth\) 的部分分式展开，
\[
H(\xi)=4\sum_{n\ge 1}\frac{1}{\xi^2+4\pi^2n^2},
\]
故
\[
H(\varphi\xi)
=
4\sum_{n\ge 1}\frac{1}{\varphi^2\xi^2+4\pi^2n^2}
=
4\varphi^{-2}\sum_{n\ge 1}\frac{1}{\xi^2+4\pi^2n^2/\varphi^2}.
\]
于是
\[
\widehat{\mathcal E}_\varphi(\xi)
=
4\sum_{n\ge 1}
\left(
\frac{\varphi^{-1}}{\xi^2+4\pi^2n^2}
+\frac{\varphi^{-3}}{\xi^2+4\pi^2n^2/\varphi^2}
\right).
\]

把每一项重写为
\[
\frac{1}{\xi^2+t}=\frac{t^{-1}}{1+\xi^2t^{-1}},
\]
便得到原子位置
\[
y=\frac{1}{4\pi^2n^2},
\qquad
y=\frac{\varphi^2}{4\pi^2n^2},
\]
及相应权重
\[
4\varphi^{-1}\frac{1}{4\pi^2n^2}
=
4\varphi^{-3}\frac{\varphi^2}{4\pi^2n^2}
=
\frac{\varphi^{-1}}{\pi^2n^2}.
\]
因此 \(\widehat{\mathcal E}_\varphi\) 正是所示测度的 Stieltjes 变换。

又因为
\[
\frac{\varphi^2}{4\pi^2n^2}\le \frac{\varphi^2}{4\pi^2}=\rho_\ast<1,
\]
故支撑包含于 \((0,\rho_\ast]\)。
最后，
\[
\nu_{\varphi}^{\mathrm{odd}}\bigl((0,\rho_\ast]\bigr)
=
\sum_{n\ge 1}\frac{2\varphi^{-1}}{\pi^2n^2}
=
\frac{2\varphi^{-1}}{\pi^2}\zeta(2)
=
\frac{2\varphi^{-1}}{\pi^2}\cdot\frac{\pi^2}{6}
=
\frac{1}{3\varphi}.
\]
证毕。
\end{proof}

\begin{theorem}[H.111 奇层系数的正矩公式]\label{thm:xi-time-part9zbc-foldpi-oddlayer-positive-moment-formula}
对每个 \(k\ge 0\) 定义
\[
M_k:=\int_{0}^{\rho_\ast}y^k\,\dd\nu_{\varphi}^{\mathrm{odd}}(y).
\]
则有显式闭式
\[
M_k
=
4(2\pi)^{-2k-2}\zeta(2k+2)\bigl(\varphi^{-1}+\varphi^{2k-1}\bigr),
\]
并且 odd formal tail 的阶乘--交错规范化系数满足
\[
u_k
:=
(-1)^k\frac{\mathfrak a_{k+1}^{(\pi)}}{(2k)!}
=
M_k
>0.
\]
从而对任意固定 \(R\ge 1\)，若记 \(q:=\varphi/2\)，则有矩展开
\[
\log C_m-\log(2\pi)
=
\sum_{k=0}^{R-1}
(-1)^k(2k)!\,M_k\,q^{(2k+1)m}
+O\!\bigl(q^{(2R+1)m}\bigr).
\]
换言之，fold-\(\pi\) 的 odd formal tail 在该规范化下恰是 \(\nu_{\varphi}^{\mathrm{odd}}\) 的普通矩序列。
\end{theorem}
\begin{proof}
由定理 \ref{thm:xi-time-part9zbc-foldpi-borel-finite-positive-measure} 中的测度定义，
\[
M_k
=
\sum_{n\ge 1}\frac{\varphi^{-1}}{\pi^2n^2}
\left(
\frac{1}{(4\pi^2n^2)^k}
+\frac{\varphi^{2k}}{(4\pi^2n^2)^k}
\right).
\]
提取公共因子可得
\[
M_k
=
\frac{\varphi^{-1}+\varphi^{2k-1}}{4^k\pi^{2k+2}}
\sum_{n\ge 1}n^{-2k-2}
=
\frac{\varphi^{-1}+\varphi^{2k-1}}{4^k\pi^{2k+2}}
\zeta(2k+2),
\]
亦即
\[
M_k
=
4(2\pi)^{-2k-2}\zeta(2k+2)\bigl(\varphi^{-1}+\varphi^{2k-1}\bigr).
\]

再用标准恒等式
\[
\zeta(2k+2)
=
(-1)^k\frac{B_{2k+2}(2\pi)^{2k+2}}{2(2k+2)!},
\]
得到
\[
M_k
=
2(-1)^k\frac{B_{2k+2}}{(2k+2)!}
\bigl(\varphi^{-1}+\varphi^{2k-1}\bigr).
\]
因此
\[
(2k)!M_k
=
(-1)^k
\frac{B_{2k+2}}{(k+1)(2k+1)}
\bigl(\varphi^{-1}+\varphi^{2k-1}\bigr)
=
(-1)^k\mathfrak a_{k+1}^{(\pi)},
\]
即
\[
u_k=(-1)^k\frac{\mathfrak a_{k+1}^{(\pi)}}{(2k)!}=M_k>0.
\]

最后，由定理 \ref{thm:conclusion-foldgauge-pi-zero-radius-odd-germ} 的有限截断表述，对任意固定 \(R\ge 1\) 有
\[
\log C_m-\log(2\pi)
=
\sum_{k=0}^{R-1}\mathfrak a_{k+1}^{(\pi)}q^{(2k+1)m}
+O\!\bigl(q^{(2R+1)m}\bigr).
\]
把
\(
\mathfrak a_{k+1}^{(\pi)}=(-1)^k(2k)!M_k
\)
逐项代回，并截断到 \(k=R-1\) 即得所示展开。证毕。
\end{proof}

\begin{theorem}[矩 Hankel 矩阵的严格正定性]\label{thm:xi-time-part9zbc-foldpi-hankel-strict-positivity}
\leanverified{paper\_xi\_time\_part9zbc\_foldpi\_hankel\_strict\_positivity}
对每个 \(N\ge 0\)，定义 Hankel 矩阵
\[
H_N:=(M_{i+j})_{0\le i,j\le N}.
\]
则 \(H_N\) 严格正定；特别地，
\[
\det H_N>0
\qquad (N\ge 0).
\]
因此，矩列 \((M_k)_{k\ge 0}\) 具有无限 Hankel 秩。
\end{theorem}
\begin{proof}
任取非零向量
\[
c=(c_0,\dots,c_N)^{\mathsf T}\in\RR^{N+1}\setminus\{0\},
\]
并定义多项式
\[
P_c(y):=\sum_{j=0}^{N}c_jy^j.
\]
则
\[
c^{\mathsf T}H_Nc
=
\sum_{i,j=0}^{N}c_ic_jM_{i+j}
=
\int_{0}^{\rho_\ast}
\left(\sum_{i=0}^{N}c_iy^i\right)^2
\dd\nu_{\varphi}^{\mathrm{odd}}(y)
=
\int_{0}^{\rho_\ast}P_c(y)^2\,\dd\nu_{\varphi}^{\mathrm{odd}}(y).
\]
由于 \(\nu_{\varphi}^{\mathrm{odd}}\) 是正测度，且其支撑由无穷多个互异点组成，非零多项式 \(P_c\) 不可能在整个支撑上恒零，故上式严格大于 \(0\)。
于是 \(H_N\) 严格正定，从而 \(\det H_N>0\)。

若 Hankel 秩有限，则高阶 Hankel 主子式终将为零，与 \(\det H_N>0\) 对所有 \(N\) 成立矛盾。故 Hankel 秩必为无限。证毕。
\end{proof}

\begin{corollary}[阶乘归一化奇层系数的严格 Hausdorff 性]\label{cor:xi-time-part9zbc-foldpi-strict-hausdorff-logconvex}
对 \(k\ge 0\) 记
\[
u_k:=(-1)^k\frac{\mathfrak a_{k+1}^{(\pi)}}{(2k)!}=M_k.
\]
则 \((u_k)_{k\ge 0}\) 是 \([0,\rho_\ast]\subset[0,1)\) 上的严格 Hausdorff moment sequence。
更具体地，对任意 \(k,\ell\ge 0\) 定义前向差分 \(\Delta u_k:=u_{k+1}-u_k\)，则
\[
(-1)^\ell\Delta^\ell u_k>0.
\]
此外，
\[
u_{k+1}^2<u_ku_{k+2}
\qquad (k\ge 0),
\]
故 \((u_k)\) 严格对数凸。
\end{corollary}
\begin{proof}
由定理 \ref{thm:xi-time-part9zbc-foldpi-oddlayer-positive-moment-formula}，
\[
u_k=\int_0^{\rho_\ast}y^k\,\dd\nu_{\varphi}^{\mathrm{odd}}(y).
\]
对任意 \(\ell\ge 0\)，迭代差分得
\[
\Delta^\ell u_k
=
\int_0^{\rho_\ast}y^k(y-1)^\ell\,\dd\nu_{\varphi}^{\mathrm{odd}}(y).
\]
由于
\[
\operatorname{supp}\nu_{\varphi}^{\mathrm{odd}}\subset(0,\rho_\ast]\subset(0,1),
\]
故
\[
(-1)^\ell\Delta^\ell u_k
=
\int_0^{\rho_\ast}y^k(1-y)^\ell\,\dd\nu_{\varphi}^{\mathrm{odd}}(y)>0.
\]
这正是严格 Hausdorff 性。

再由 Cauchy--Schwarz 不等式，
\[
u_{k+1}^2
=
\left(\int y^{k+1}\,\dd\nu_{\varphi}^{\mathrm{odd}}(y)\right)^2
\le
\left(\int y^k\,\dd\nu_{\varphi}^{\mathrm{odd}}(y)\right)
\left(\int y^{k+2}\,\dd\nu_{\varphi}^{\mathrm{odd}}(y)\right)
=
u_ku_{k+2}.
\]
而支撑上 \(y\) 不是常数，故等号不可能成立，于是得到严格不等式。证毕。
\end{proof}

\begin{proposition}[canonical scalar \texorpdfstring{$S$}{S}-fraction 的存在唯一性]\label{prop:xi-time-part9zbc-foldpi-canonical-positive-sfraction}
定义
\[
\mathcal M_\varphi(z):=\sum_{k\ge 0}u_k z^k
=
\int_0^{\rho_\ast}\frac{\dd\nu_{\varphi}^{\mathrm{odd}}(y)}{1-yz},
\qquad |z|<\rho_\ast^{-1}=\frac{4\pi^2}{\varphi^2}.
\]
则 \(\mathcal M_\varphi\) 属于 Hausdorff--Stieltjes 类，并因而存在唯一的正 \(S\)-fraction
\[
\mathcal M_\varphi(z)
=
\cfrac{\alpha_0}{1-\cfrac{\alpha_1 z}{1-\cfrac{\alpha_2 z}{1-\cfrac{\alpha_3 z}{\ddots}}}},
\qquad
\alpha_j>0.
\]
其收敛子分母与分子分别由 \(\nu_{\varphi}^{\mathrm{odd}}\) 的 monic orthogonal polynomials 及其 first associated family 给出。
特别地，H.111 的 odd formal tail 在该规范化下内生地产生了一条唯一的标量 continued fraction。
\end{proposition}
\begin{proof}
定理 \ref{thm:xi-time-part9zbc-foldpi-oddlayer-positive-moment-formula} 与推论 \ref{cor:xi-time-part9zbc-foldpi-strict-hausdorff-logconvex} 表明 \((u_k)\) 是支撑于有界区间 \([0,\rho_\ast]\) 上的严格 Hausdorff moment sequence。
于是 \(\mathcal M_\varphi\) 的积分表示成立，并且对应正测度唯一。
关于 bounded interval 上正交测度的唯一性，可参见 \href{https://dlmf.nist.gov/18.2.viii}{DLMF 18.2(viii)}；关于 \(S\)-fraction 的 formal 定义与 qd 算法，可参见 \href{https://dlmf.nist.gov/3.10}{DLMF 3.10}；而 monic denominator / numerator polynomials 与 Markov 定理的 continued-fraction 接口，则见 \href{https://dlmf.nist.gov/18.2.x}{DLMF 18.2(x)}。

将这些经典结果施用于 \(\nu_{\varphi}^{\mathrm{odd}}\) 即可得到唯一正 \(S\)-fraction 展开。
又因为支撑不是有限集合，相关 Jacobi 参数全部为正，故 \(\alpha_j>0\) 对所有 \(j\) 成立。证毕。
\end{proof}

\begin{theorem}[Bernoulli--\texorpdfstring{$\zeta$}{zeta} 算术扭曲的秩跃迁]\label{thm:xi-time-part9zbc-foldpi-zeta-rank-jump}
\leanverified{paper\_xi\_time\_part9zbc\_foldpi\_zeta\_rank\_jump}
定义纯几何影子
\[
g_k:=\varphi^{-1}+\varphi^{2k-1}
\qquad (k\ge 0).
\]
则
\[
u_k
=
4(2\pi)^{-2k-2}\zeta(2k+2)\,g_k.
\]
并且：
\begin{enumerate}
  \item \((g_k)_{k\ge 0}\) 具有精确 Hankel 秩 \(2\)；
  \item \((u_k)_{k\ge 0}\) 具有无限 Hankel 秩，且其全部有限 Hankel 主子式都严格为正。
\end{enumerate}
因此，\(\zeta(2k+2)\) 并非可被抽出的无害权重；它正是把项目中黄金两态几何从有限秩代数影子抬升为无限深 \(S\)-fraction / Jacobi 对象的秩跃迁机制。
\end{theorem}
\begin{proof}
第一式只是定理 \ref{thm:xi-time-part9zbc-foldpi-oddlayer-positive-moment-formula} 的显式式改写。

又因为
\[
g_k=f_{k+1},
\]
其中 \(f_k\) 正是定理 \ref{thm:xi-time-part9zb-foldpi-two-atomic-moment-characterization} 中的 Bernoulli-free 规范化矩列，故 \((g_k)\) 的 Hankel 秩至多为 \(2\)。
另一方面，
\[
g_0=\frac{2}{\varphi},
\qquad
g_1=\varphi^{-1}+\varphi,
\qquad
g_2=\varphi^{-1}+\varphi^3,
\]
并且
\[
\det
\begin{pmatrix}
g_0 & g_1\\
g_1 & g_2
\end{pmatrix}
=
g_0g_2-g_1^2
=1>0,
\]
故其 Hankel 秩恰为 \(2\)。

第二条则由定理 \ref{thm:xi-time-part9zbc-foldpi-hankel-strict-positivity} 直接给出：\((u_k)\) 的每个有限 Hankel 主块都严格正定，故其 Hankel 秩无限。证毕。
\end{proof}

\begin{conjecture}[Jacobi 参数的黄金平方重整化]\label{conj:xi-time-part9zbc-foldpi-jacobi-golden-square-renormalization}
令 \((\alpha_n)_{n\ge 0}\) 为命题 \ref{prop:xi-time-part9zbc-foldpi-canonical-positive-sfraction} 中 canonical \(S\)-fraction 的正系数；等价地，令
\[
(\mathfrak a_n,\mathfrak b_n)_{n\ge 1}
\]
为测度 \(\nu_{\varphi}^{\mathrm{odd}}\) 的 Jacobi 参数。
则该参数序列的尾部既不可能 eventually periodic，也不可能退化为固定阶超几何 continued fraction 的系数轨道；其唯一自然的尾部自相似应当由
\[
y\longmapsto \varphi^2 y
\]
所生成。
换言之，H.111 所导出的 scalar continued-fraction 对象应当是一条 \(\varphi^2\)-重整化固定点型 Jacobi 轨道，而不是周期尾型轨道。
\end{conjecture}
\begin{proof}[证据]
定理 \ref{thm:xi-time-part9zbc-foldpi-borel-finite-positive-measure} 表明，底层正测度的支撑恰为两条 reciprocal-square ladders
\[
\left\{\frac{1}{4\pi^2n^2}\right\}_{n\ge 1}
\quad\text{与}\quad
\left\{\frac{\varphi^2}{4\pi^2n^2}\right\}_{n\ge 1},
\]
两者的比例由无理数 \(\varphi^2\) 锁定。
另一方面，定理 \ref{thm:xi-time-part9zbc-foldpi-zeta-rank-jump} 表明，正是 \(\zeta(2k+2)\) 的算术扭曲把定理 \ref{thm:xi-time-part9zb-foldpi-two-atomic-moment-characterization} 中秩 \(2\) 的几何影子抬升为无限 Hankel 深度。

因此，若 Jacobi 尾部 eventually periodic，则对应 Weyl 函数将落入有限-gap 代数型类；若其退化为固定阶超几何 continued fraction，则又会重新返回 fixed-order scalar 范畴。
这两类尾部机制都与 \(\nu_{\varphi}^{\mathrm{odd}}\) 的双梯支撑及其无限正定 Hankel 结构在机制上不相容。
在现有审计链下，唯一仍与底层支撑自相似相匹配的乘法重整化，只剩
\[
y\mapsto \varphi^2 y.
\]
证毕。
\end{proof}

\noindent
于是，H.111 的 odd formal tail 不再只意味着“存在一个 Borel 变换”。
更强地，它在阶乘与交错符号的规范化之后，严格落入有界正测度的 Hausdorff moment problem；相应的 scalar \(S\)-fraction 与 Jacobi 轨道由此被唯一决定。
这也说明，真正持续承载该项目 \(\pi\)-tail 的对象，并不是另一条快收敛标量级数，而是一个由正矩问题、正交多项式与 continued fraction 共同锁定的 Jacobi 范畴对象。

\endinput

\paragraph{scaled multiplicity law、双折点容量骨架与 fold-gauge 的正 Laplace 传输}
在 H.81 的最优可逆容量、H.95 的两尺度一致渐近、H.110 的精确 Stirling 余项表示以及 H.111 的 Bernoulli 奇层 tail 都已建立之后，还可把这四条此前分散出现的接口压缩为同一对象：一个支撑在两点上的 scaled multiplicity law。
其一侧经 \(x\mapsto \min(x,s)\) 变换给出临界预言机预算下的 recoverability 相图，另一侧经 \(x\mapsto x^{-(2r-1)}\) 给出 fold-\(\pi\) tail 的 Bernoulli 权重，而 H.110 的余项恒等式则把同一纤维谱进一步送入一个严格正的 Laplace 传输核。

\begin{theorem}[临界预言机预算的双折点相变律]\label{thm:xi-time-part9zbd-scaled-multiplicity-double-kink}
设
\[
\mathrm{Succ}_m(B):=
2^{-m}\sum_{w\in X_m}\min\!\bigl(d_m^{\mathrm{bin}}(w),2^B\bigr)
\qquad (B\in\ZZ_{\ge 0}),
\]
并令整数序列 \((B_m)_{m\ge 1}\) 满足
\[
2^{B_m}\left(\frac{\varphi}{2}\right)^m\longrightarrow s\ge 0.
\]
则
\[
\mathrm{Succ}_m(B_m)\longrightarrow \mathsf S_\varphi(s),
\]
其中
\[
\mathsf S_\varphi(s):=
\frac{\varphi^2}{\sqrt5}
\left(
\frac1\varphi \min\{1,s\}
+
\frac1{\varphi^2}\min\{\varphi^{-1},s\}
\right).
\]
等价地，
\[
\mathsf S_\varphi(s)=
\begin{cases}
\dfrac{\varphi^2}{\sqrt5}\,s, & 0\le s\le \varphi^{-1},\\[6pt]
\dfrac{\varphi s+\varphi^{-1}}{\sqrt5}, & \varphi^{-1}\le s\le 1,\\[6pt]
1, & s\ge 1.
\end{cases}
\]
因此，在临界 bit-budget 尺度
\[
B_m\asymp m\log_2\!\left(\frac{2}{\varphi}\right)
\]
上，recoverability 极限恰有两个且仅两个不可消去折点，位置精确为
\[
s=\varphi^{-1},
\qquad
s=1.
\]
\end{theorem}
\begin{proof}
记
\[
A_m:=\left(\frac{2}{\varphi}\right)^m,
\qquad
\Psi_m(\tau):=
2^{-m}\mathcal C_m^{\mathrm{bin,cont}}(\tau A_m).
\]
由定义，
\[
\mathrm{Succ}_m(B_m)=\Psi_m(s_m),
\qquad
s_m:=2^{B_m}\left(\frac{\varphi}{2}\right)^m.
\]
另一方面，定理 \ref{thm:xi-time-part70b-center-capacity-two-kink-skeleton} 已给出：对每个固定 \(\tau\ge 0\)，都有
\[
\Psi_m(\tau)\longrightarrow
\frac{\varphi}{\sqrt5}\min\{1,\tau\}
+
\frac{1}{\sqrt5}\min\{\varphi^{-1},\tau\}
=
\mathsf S_\varphi(\tau).
\]
又因 \(\mathcal C_m^{\mathrm{bin,cont}}\) 对阈值变量的斜率至多为 \(|X_m|=F_{m+2}\)，故
\[
\abs{\Psi_m(\tau_1)-\Psi_m(\tau_2)}
\le
\frac{A_mF_{m+2}}{2^m}\abs{\tau_1-\tau_2}
=
\frac{F_{m+2}}{\varphi^m}\abs{\tau_1-\tau_2}.
\]
由 Binet 公式，\(\sup_m F_{m+2}\varphi^{-m}<2\)，故 \((\Psi_m)\) 在每个有界区间上一致 Lipschitz。
于是
\[
\abs{\Psi_m(s_m)-\mathsf S_\varphi(s)}
\le
\abs{\Psi_m(s_m)-\Psi_m(s)}
+
\abs{\Psi_m(s)-\mathsf S_\varphi(s)}
\longrightarrow 0,
\]
因为 \(s_m\to s\)。最后把 \(\mathsf S_\varphi\) 逐段化简，便得到所示分段闭式。证毕。
\end{proof}

\begin{theorem}[同一两原子极限律同时控制容量相图与 Bernoulli 奇层]\label{thm:xi-time-part9zbd-scaled-multiplicity-min-oddmoment}
定义两原子概率测度
\[
\nu_\varphi^{\mathrm{sc}}
:=
\frac1\varphi\,\delta_1+\frac1{\varphi^2}\,\delta_{\varphi^{-1}}.
\]
则它同时满足以下两条表示：
\begin{enumerate}
  \item 临界成功率曲线可写为
  \[
  \mathsf S_\varphi(s)
  =
  \frac{\varphi^2}{\sqrt5}
  \int \min\{x,s\}\,\dd\nu_\varphi^{\mathrm{sc}}(x).
  \]
  \item 对任意整数 \(r\ge 1\)，其负奇矩满足
  \[
  \int x^{-(2r-1)}\,\dd\nu_\varphi^{\mathrm{sc}}(x)
  =
  \varphi^{-1}+\varphi^{2r-3}.
  \]
\end{enumerate}
因此，对任意固定 \(R\ge 1\)，fold-gauge 常数的奇层 tail 可重写为
\[
\log C_m
=
\log(2\pi)
+
\sum_{r=1}^{R}
\frac{B_{2r}}{r(2r-1)}
\left(
\int x^{-(2r-1)}\,\dd\nu_\varphi^{\mathrm{sc}}(x)
\right)
\left(\frac{\varphi}{2}\right)^{(2r-1)m}
+
O\!\left(\left(\frac{\varphi}{2}\right)^{(2R+1)m}\right).
\]
换言之，项目中的临界可逆性相图与 H.111 的 Bernoulli 奇层 tail 并非两条并列结构，而是同一两原子尺度律的 \(\min\)-变换与 inverse-odd-moment 变换。
\end{theorem}
\begin{proof}
第一式只是把定理 \ref{thm:xi-time-part9zbd-scaled-multiplicity-double-kink} 中的闭式写成测度积分：
\[
\frac1\varphi \min\{1,s\}+\frac1{\varphi^2}\min\{\varphi^{-1},s\}
=
\int \min\{x,s\}\,\dd\nu_\varphi^{\mathrm{sc}}(x).
\]
第二式直接计算得
\[
\int x^{-(2r-1)}\,\dd\nu_\varphi^{\mathrm{sc}}(x)
=
\frac1\varphi\cdot 1^{-(2r-1)}
+
\frac1{\varphi^2}\cdot(\varphi^{-1})^{-(2r-1)}
=
\varphi^{-1}+\varphi^{2r-3}.
\]
将该恒等式代回定理 \ref{thm:fold-bin-gauge-constant-stirling-bernoulli-hierarchy} 的系数表达，便得到最后一式。证毕。
\end{proof}

\begin{theorem}[有限原子尺度律的折点数与指数 ladder 刚性]\label{thm:xi-time-part9zbd-finite-atomic-kink-ladder-rigidity}
设
\[
\nu=\sum_{j=1}^{J}p_j\delta_{\sigma_j}
\]
是任意有限原子概率测度，其中
\[
0<\sigma_1<\cdots<\sigma_J,
\qquad
p_j>0,
\qquad
\sum_{j=1}^{J}p_j=1.
\]
对某个正常数 \(\kappa\) 定义
\[
\mathsf S_\nu(s):=
\kappa\int \min\{x,s\}\,\dd\nu(x)
=
\kappa\sum_{j=1}^{J}p_j\min\{\sigma_j,s\},
\]
以及负奇矩序列
\[
M_r(\nu):=
\int x^{-(2r-1)}\,\dd\nu(x)
=
\sum_{j=1}^{J}p_j\sigma_j^{-(2r-1)}
\qquad (r\ge 1).
\]
则成立：
\begin{enumerate}
  \item \(\mathsf S_\nu\) 在每个 \(\sigma_j\) 处恰有一个 kink；计入末端饱和区后，\(\mathsf S_\nu\) 恰有 \(J+1\) 个 affine 线性区间。
  \item 序列 \((M_r(\nu))_{r\ge 1}\) 恰是 \(J\) 条互异指数 ladder 的和。
  \item 因而，在有限原子范畴内，若某一模型的临界容量曲线恰有两个折点，而其 Bernoulli 奇层权重又恰是两项指数型和，则其 scaled multiplicity law 必为两原子测度。
\end{enumerate}
特别地，定理 \ref{thm:xi-time-part9zbd-scaled-multiplicity-double-kink} 与定理 \ref{thm:xi-time-part9zbd-scaled-multiplicity-min-oddmoment} 联合表明，项目中的两尺度律并非近似简化，而是结构刚性：任意带正质量的第三尺度都会同时强制产生一个新的容量折点与一条新的 Bernoulli 指数 ladder。
\end{theorem}
\begin{proof}
对 \(s\notin\{\sigma_1,\dots,\sigma_J\}\)，逐段求导得到
\[
\frac{\dd}{\dd s}\mathsf S_\nu(s)
=
\kappa\,\nu((s,\infty))
=
\kappa\sum_{\sigma_j>s}p_j.
\]
故 \(\mathsf S_\nu\) 在每个 \(\sigma_j\) 处的斜率跳跃量正为 \(\kappa p_j\)，于是每个原子对应一个且仅一个 kink，而线性区间则为
\[
(0,\sigma_1),\ (\sigma_1,\sigma_2),\ \dots,\ (\sigma_J,\infty),
\]
共 \(J+1\) 段。

另一方面，
\[
M_r(\nu)=\sum_{j=1}^{J}p_j(\sigma_j^{-1})^{2r-1},
\]
这正是 \(J\) 条互异指数序列之和。若再增加一条新尺度 \(\sigma_{J+1}\) 且其质量为正，则 \(\mathsf S_\nu\) 中必新增一个 kink，而 \(M_r(\nu)\) 中必新增一条指数 ladder。
因此，在有限原子范畴内，“折点数”与“指数 ladder 数”都精确记录原子个数。将此一般论断应用于前两条定理，即得最后结论。证毕。
\end{proof}

\begin{theorem}[fold-gauge 缺陷的精确 Binet--Laplace 传输]\label{thm:xi-time-part9zbd-foldgauge-binet-laplace-transport}
定义 Stirling 余项
\[
\varepsilon(x):=
\log\Gamma(x+1)-x\log x+x-\frac12\log(2\pi x)
\qquad (x>0),
\]
以及纤维 multiplicity spectrum 的 Laplace 轮廓
\[
\mathcal L_m(t):=
\frac1{|X_m|}\sum_{w\in X_m}e^{-d_m^{\mathrm{bin}}(w)t}
\qquad (t>0).
\]
则有精确恒等式
\[
\log\frac{C_m}{2\pi}
=
\int_0^\infty K(t)\,\mathcal L_m(t)\,\dd t,
\]
其中
\[
K(t):=
2\left(\frac12-\frac1t+\frac1{e^t-1}\right)\frac1t
=
\frac{\coth(t/2)-2/t}{t}
>0
\qquad (t>0).
\]
特别地，\(\log(C_m/2\pi)\) 不是一条后验拟合的渐近尾，而是纤维 multiplicity spectrum 的精确正 Laplace 传输。进一步，\(\varepsilon\) 是完全单调函数，即对任意整数 \(n\ge 0\) 与任意 \(x>0\)，都有
\[
(-1)^n\varepsilon^{(n)}(x)>0.
\]
\end{theorem}
\begin{proof}
主文 H.110 的精确 Stirling 余项表示给出
\[
\log C_m
=
\log(2\pi)
+
\frac{2}{|X_m|}\sum_{w\in X_m}\varepsilon\bigl(d_m^{\mathrm{bin}}(w)\bigr).
\]
另一方面，由 Binet 公式可得（见 \href{https://dlmf.nist.gov/5.9}{DLMF 5.9}）
\[
\varepsilon(x)
=
\int_0^\infty
\left(\frac12-\frac1t+\frac1{e^t-1}\right)e^{-xt}\frac{\dd t}{t}
\qquad (x>0).
\]
把该表示代回并交换有限求和与积分，得到
\[
\log\frac{C_m}{2\pi}
=
\frac{2}{|X_m|}
\sum_{w\in X_m}
\int_0^\infty
\left(\frac12-\frac1t+\frac1{e^t-1}\right)e^{-d_m^{\mathrm{bin}}(w)t}\frac{\dd t}{t}.
\]
于是
\[
\log\frac{C_m}{2\pi}
=
\int_0^\infty
2\left(\frac12-\frac1t+\frac1{e^t-1}\right)\frac1t
\cdot
\frac1{|X_m|}\sum_{w\in X_m}e^{-d_m^{\mathrm{bin}}(w)t}\,\dd t,
\]
即所求第一式。

再用恒等式
\[
\frac1{e^t-1}=\frac12\coth(t/2)-\frac12
\]
便得
\[
K(t)=\frac{\coth(t/2)-2/t}{t}.
\]
对 \(u>0\) 恒有 \(\coth u>1/u\)，故 \(K(t)>0\)。

最后，把
\[
B(t):=
\left(\frac12-\frac1t+\frac1{e^t-1}\right)\frac1t
\]
记作正核，则
\[
\varepsilon(x)=\int_0^\infty B(t)e^{-xt}\,\dd t.
\]
逐次对 \(x\) 求导即得
\[
(-1)^n\varepsilon^{(n)}(x)
=
\int_0^\infty t^n B(t)e^{-xt}\,\dd t
>0.
\]
因此 \(\varepsilon\) 完全单调。证毕。
\end{proof}

\begin{corollary}[fold-gauge 常数对 \texorpdfstring{$2\pi$}{2pi} 的严格上逼近与非渐近双边夹逼]\label{cor:xi-time-part9zbd-foldgauge-jensen-sandwich}
记平均纤维大小与最小纤维大小分别为
\[
\bar d_m:=\frac{2^m}{|X_m|},
\qquad
d_{\min}(m):=\min_{w\in X_m}d_m^{\mathrm{bin}}(w).
\]
则对每个 \(m\ge 1\) 都有
\[
C_m>2\pi,
\]
并且满足完全显式的双边夹逼
\[
2\,\varepsilon(\bar d_m)
\le
\log\frac{C_m}{2\pi}
\le
\frac{1}{6|X_m|}\sum_{w\in X_m}d_m^{\mathrm{bin}}(w)^{-1}
\le
\frac{1}{6\,d_{\min}(m)}.
\]
因此，项目中的 \(\pi\)-恢复并不是在 \(2\pi\) 的两侧摇摆逼近，而是对 \(2\pi\) 的严格上侧校正；其偏离量同时受平均纤维尺度的 Jensen 下界与最小纤维尺度的倒数上界控制。
\end{corollary}
\begin{proof}
由定理 \ref{thm:xi-time-part9zbd-foldgauge-binet-laplace-transport}，\(\varepsilon(x)>0\) 对一切 \(x>0\) 成立，且
\[
\log\frac{C_m}{2\pi}
=
\frac{2}{|X_m|}\sum_{w\in X_m}\varepsilon\bigl(d_m^{\mathrm{bin}}(w)\bigr).
\]
故立得 \(C_m>2\pi\)。

再由完全单调性，\(\varepsilon\) 尤其凸。于是 Jensen 不等式给出
\[
\frac{2}{|X_m|}\sum_{w\in X_m}\varepsilon\bigl(d_m^{\mathrm{bin}}(w)\bigr)
\ge
2\,\varepsilon\!\left(
\frac1{|X_m|}\sum_{w\in X_m}d_m^{\mathrm{bin}}(w)
\right)
=
2\,\varepsilon(\bar d_m),
\]
其中用到
\[
\sum_{w\in X_m}d_m^{\mathrm{bin}}(w)=2^m.
\]

上界则由经典 Stirling 余项界
\[
0<\varepsilon(x)\le \frac{1}{12x}
\qquad (x>0)
\]
得到
\[
\log\frac{C_m}{2\pi}
\le
\frac{2}{|X_m|}\sum_{w\in X_m}\frac{1}{12\,d_m^{\mathrm{bin}}(w)}
=
\frac{1}{6|X_m|}\sum_{w\in X_m}d_m^{\mathrm{bin}}(w)^{-1}.
\]
最后，由 \(d_m^{\mathrm{bin}}(w)\ge d_{\min}(m)\) 对全部 \(w\) 成立，
\[
\frac{1}{|X_m|}\sum_{w\in X_m}d_m^{\mathrm{bin}}(w)^{-1}
\le
d_{\min}(m)^{-1}.
\]
三式合并即得结论。证毕。
\end{proof}

\noindent
由此，项目中的 \((2\pi)\) 不再只是一个被动出现的极限常数。
同一条 scaled multiplicity law \(\nu_\varphi^{\mathrm{sc}}\) 一方面给出临界 bit-budget 下的双折点 recoverability 相图，另一方面给出 H.111 的 Bernoulli 奇层权重；而 H.110 的精确 Stirling 余项表示又把有限窗规范缺陷送入一个严格正的 Laplace 传输核。
于是，可逆资源临界相变、fold-\(\pi\) 奇层 tail 与规范常数缺陷三者便不再是并列现象，而是同一纤维 multiplicity spectrum 在 \(\min\)、inverse-odd-moment 与 Laplace 三种变换下的不同投影。

\endinput

\paragraph{fold-\texorpdfstring{$\pi$}{pi} 奇层 tail 的 compact Jacobi 谱宿主、operator \texorpdfstring{$\zeta$}{zeta} 与 Fredholm 双正弦}
沿用定理 \ref{thm:xi-time-part9zbc-foldpi-borel-finite-positive-measure}、
\ref{thm:xi-time-part9zbc-foldpi-oddlayer-positive-moment-formula}
与命题 \ref{prop:xi-time-part9zbc-foldpi-canonical-positive-sfraction} 的记号。
记
\[
\mu_\varphi:=\nu_{\varphi}^{\mathrm{odd}}
=
\sum_{n\ge 1}
\frac{\varphi^{-1}}{\pi^2n^2}
\left(
\delta_{(4\pi^2n^2)^{-1}}
+\delta_{\varphi^2(4\pi^2n^2)^{-1}}
\right),
\]
并继续写
\[
u_k=\int_0^{\rho_\ast}y^k\,\dd\mu_\varphi(y)
=
4(2\pi)^{-2k-2}\zeta(2k+2)\bigl(\varphi^{-1}+\varphi^{2k-1}\bigr)
\qquad (k\ge 0).
\]
这样，H.111 的 odd formal tail 已不再只被视为一个 Hausdorff moment sequence 与唯一正 \(S\)-fraction，而是进一步拥有一个正、紧、自伴的 Jacobi 谱宿主；相应的 operator \(\zeta\) 与 Fredholm 行列式都可完全显式写出。

\begin{theorem}[fold-\texorpdfstring{$\pi$}{pi} 奇层的 canonical compact Jacobi realization]\label{thm:xi-time-part9zbe-foldpi-canonical-compact-jacobi}
存在唯一的 Jacobi 算子
\[
J_\varphi=
\begin{pmatrix}
\beta_0 & \alpha_1 \\
\alpha_1 & \beta_1 & \alpha_2 \\
& \alpha_2 & \beta_2 & \alpha_3 \\
&& \ddots & \ddots
\end{pmatrix},
\qquad
\alpha_n>0,\ \beta_n\in\RR,
\]
以及一个循环向量 \(v_0\)，使得
\[
u_k=\langle v_0,J_\varphi^k v_0\rangle
\qquad (k\ge 0).
\]
并且 \(J_\varphi\) 为正、紧、自伴算子，其谱恰为
\[
\sigma(J_\varphi)
=
\{0\}
\cup
\left\{\frac{1}{4\pi^2n^2}:n\ge 1\right\}
\cup
\left\{\frac{\varphi^2}{4\pi^2n^2}:n\ge 1\right\}.
\]
\end{theorem}
\begin{proof}
记
\[
K:=\operatorname{supp}\mu_\varphi\cup\{0\}\subset [0,\rho_\ast].
\]
由于 \(K\) 为紧集且多项式代数在 \(K\) 上分离点并含常数，Stone--Weierstrass 定理表明多项式在 \(C(K)\) 中稠密，因而也在 \(L^2(\mu_\varphi)\) 中稠密。
对
\[
1,\ y,\ y^2,\ \dots
\]
作 Gram--Schmidt 正交化，得到一组正交归一多项式 \((p_n)_{n\ge 0}\)。

在 \(L^2(\mu_\varphi)\) 上定义乘法算子
\[
(M_yf)(y):=yf(y).
\]
因为 \(0\le y\le \rho_\ast\)，故 \(M_y\) 是有界正自伴算子。
在正交基 \((p_n)\) 下，Favard 三项递推给出
\[
yp_n=\alpha_{n+1}p_{n+1}+\beta_n p_n+\alpha_n p_{n-1},
\qquad
\alpha_n>0,\ \beta_n\in\RR,
\]
从而得到唯一 Jacobi 矩阵 \(J_\varphi\)。
若令
\[
U:L^2(\mu_\varphi)\longrightarrow \ell^2(\NN_0),
\qquad
U(p_n)=e_n,
\]
则
\[
J_\varphi=UM_yU^{-1}.
\]
取
\[
v_0:=U(1),
\]
便有
\[
\langle v_0,J_\varphi^k v_0\rangle
=
\langle 1,M_y^k1\rangle_{L^2(\mu_\varphi)}
=
\int_0^{\rho_\ast}y^k\,\dd\mu_\varphi(y)
=
u_k.
\]
又因为多项式在 \(L^2(\mu_\varphi)\) 中稠密，而
\[
\operatorname{span}\{M_y^k1:\ k\ge 0\}
\]
正是多项式所张成的子空间，故 \(v_0\) 为循环向量。

下面证明紧性。
对任意 \(\varepsilon>0\)，定义
\[
M_y^{(\varepsilon)}:=M_y\,\mathbf 1_{\{y\ge \varepsilon\}}.
\]
由于 \(\mu_\varphi\) 的支撑由两条 \(n^{-2}\) 梯构成，在区间 \([\varepsilon,\rho_\ast]\) 上只有有限多个原子，因此 \(M_y^{(\varepsilon)}\) 为有限秩算子。
另一方面，
\[
\|M_y-M_y^{(\varepsilon)}\|
\le
\sup_{0\le y<\varepsilon}y
\le
\varepsilon.
\]
故 \(M_y\) 是有限秩算子的范数极限，从而紧；于是 \(J_\varphi=UM_yU^{-1}\) 亦紧。

最后，由 \(\mu_\varphi\) 的纯点性，
\[
\sigma(M_y)=\overline{\operatorname{supp}\mu_\varphi}
=
\{0\}
\cup
\left\{\frac{1}{4\pi^2n^2}:n\ge 1\right\}
\cup
\left\{\frac{\varphi^2}{4\pi^2n^2}:n\ge 1\right\}.
\]
酉等价保持谱，故 \(\sigma(J_\varphi)=\sigma(M_y)\)。证毕。
\end{proof}

\begin{corollary}[Jacobi 参数的消失律与最终周期尾不可能性]\label{cor:xi-time-part9zbe-foldpi-jacobi-vanishing-no-periodic-tail}
若
\[
J_\varphi e_n=\alpha_n e_{n-1}+\beta_n e_n+\alpha_{n+1}e_{n+1},
\qquad
\alpha_0:=0,
\]
则有
\[
\alpha_n\longrightarrow 0,
\qquad
\beta_n\longrightarrow 0.
\]
特别地，\((\alpha_n,\beta_n)\) 不可能 eventually periodic。
因此，猜想 \ref{conj:xi-time-part9zbc-foldpi-jacobi-golden-square-renormalization} 中的 eventual-periodic 分支被严格排除；现有审计链下仅剩 \(\varphi^2\)-重整化型尾部仍与该 compact Jacobi 模型相容。
\end{corollary}
\begin{proof}
由于 \(J_\varphi\) 紧且标准基向量 \(e_n\rightharpoonup 0\) 弱收敛，故
\[
\|J_\varphi e_n\|\to 0.
\]
而
\[
J_\varphi e_n=\alpha_n e_{n-1}+\beta_n e_n+\alpha_{n+1}e_{n+1},
\]
故
\[
\|J_\varphi e_n\|^2=\alpha_n^2+\beta_n^2+\alpha_{n+1}^2\longrightarrow 0.
\]
于是
\[
\alpha_n\to 0,
\qquad
\beta_n\to 0.
\]

若 \((\alpha_n,\beta_n)\) 自某处起 eventually periodic 且尾部不恒零，则存在常数 \(c>0\) 与无穷多个 \(n\) 使
\[
\|J_\varphi e_n\|\ge c.
\]
取其中一个足够稀疏的子序列 \((n_k)\)，使得各向量 \(J_\varphi e_{n_k}\) 的支撑
\[
\{e_{n_k-1},e_{n_k},e_{n_k+1}\}
\]
两两不交。
则
\[
\|J_\varphi e_{n_k}-J_\varphi e_{n_\ell}\|^2
=
\|J_\varphi e_{n_k}\|^2+\|J_\varphi e_{n_\ell}\|^2
\ge
2c^2
\qquad (k\neq \ell),
\]
从而 \((J_\varphi e_{n_k})\) 不可能有范数收敛子列，这与紧性矛盾。
故 eventually periodic 的尾部只能是 eventually zero。

但若 \((\alpha_n,\beta_n)\) eventually zero，则 \(J_\varphi\) 为有限秩算子。
这又与定理 \ref{thm:xi-time-part9zbe-foldpi-canonical-compact-jacobi} 中 \(J_\varphi\) 拥有无穷多个非零谱点矛盾。
故 eventually periodic 不可能发生。证毕。
\end{proof}

\begin{theorem}[compact Jacobi 模型的精确谱计数与 Weyl 型平方反比律]\label{thm:xi-time-part9zbe-foldpi-jacobi-counting-weyl}
对 \(\lambda>0\)，定义特征值计数函数
\[
N_{J_\varphi}(\lambda)
:=
\#\bigl\{\lambda'\in \sigma(J_\varphi)\setminus\{0\}:\ \lambda'\ge \lambda\bigr\}.
\]
则有严格闭式
\[
N_{J_\varphi}(\lambda)
=
\left\lfloor \frac{1}{2\pi\sqrt{\lambda}}\right\rfloor
+
\left\lfloor \frac{\varphi}{2\pi\sqrt{\lambda}}\right\rfloor.
\]
从而当 \(\lambda\downarrow 0\) 时，
\[
N_{J_\varphi}(\lambda)
=
\frac{1+\varphi}{2\pi}\lambda^{-1/2}+O(1)
=
\frac{\varphi^2}{2\pi}\lambda^{-1/2}+O(1).
\]
若把非零特征值按降序写为
\[
\lambda_1(J_\varphi)\ge \lambda_2(J_\varphi)\ge \cdots >0,
\]
则
\[
\lambda_n(J_\varphi)
=
\frac{\varphi^4}{4\pi^2}\,n^{-2}+O(n^{-3}).
\]
\end{theorem}
\begin{proof}
由定理 \ref{thm:xi-time-part9zbe-foldpi-canonical-compact-jacobi}，
\[
\sigma(J_\varphi)\setminus\{0\}
=
\left\{\frac{1}{4\pi^2n^2}:n\ge 1\right\}
\cup
\left\{\frac{\varphi^2}{4\pi^2n^2}:n\ge 1\right\}.
\]
两列互不相交；若有
\[
\frac{1}{4\pi^2n^2}=\frac{\varphi^2}{4\pi^2m^2},
\]
则 \(m^2=\varphi^2n^2\)，与 \(\varphi^2\notin\QQ\) 矛盾。
因此
\[
N_{J_\varphi}(\lambda)
=
\#\left\{n\ge 1:\ \frac{1}{4\pi^2n^2}\ge \lambda\right\}
+
\#\left\{n\ge 1:\ \frac{\varphi^2}{4\pi^2n^2}\ge \lambda\right\},
\]
直接化为所示 floor 公式。

再由
\[
\lfloor x\rfloor=x+O(1),
\qquad
1+\varphi=\varphi^2,
\]
得到
\[
N_{J_\varphi}(\lambda)
=
\frac{\varphi^2}{2\pi}\lambda^{-1/2}+O(1).
\]

令
\[
C:=\frac{\varphi^2}{2\pi}.
\]
把 \(\lambda=\lambda_n(J_\varphi)\) 代入计数式，则
\[
n\le C\,\lambda_n^{-1/2}
\qquad\text{且}\qquad
n>C\,\lambda_n^{-1/2}-2,
\]
因为两个 floor 项之和与其去掉 floor 后的和相差严格小于 \(2\)。
于是
\[
\frac{n}{C}\le \lambda_n^{-1/2}<\frac{n+2}{C},
\]
等价于
\[
\frac{C^2}{(n+2)^2}<\lambda_n\le \frac{C^2}{n^2}.
\]
两端与 \(C^2n^{-2}\) 的差都为 \(O(n^{-3})\)，故
\[
\lambda_n=\frac{C^2}{n^2}+O(n^{-3})
=
\frac{\varphi^4}{4\pi^2}\,n^{-2}+O(n^{-3}).
\]
证毕。
\end{proof}

\begin{theorem}[Schatten 阈值与 operator \texorpdfstring{$\zeta$}{zeta} 的闭式]\label{thm:xi-time-part9zbe-foldpi-jacobi-schatten-zeta}
对 \(\Re s>1/2\)，定义
\[
\zeta_{J_\varphi}(s):=\Tr(J_\varphi^s).
\]
则有严格闭式
\[
\zeta_{J_\varphi}(s)
=
(4\pi^2)^{-s}\bigl(1+\varphi^{2s}\bigr)\zeta(2s).
\]
因此：
\begin{enumerate}
  \item \(J_\varphi\in \mathcal S_p\) 当且仅当 \(p>1/2\)；
  \item \(\zeta_{J_\varphi}\) 亚纯延拓到整个复平面，并且只在 \(s=\tfrac12\) 处具有一个简单极点，其留数为
  \[
  \operatorname*{Res}_{s=1/2}\zeta_{J_\varphi}(s)
  =
  \frac{1+\varphi}{4\pi}
  =
  \frac{\varphi^2}{4\pi};
  \]
  \item 特别地，
  \[
  \Tr(J_\varphi)=\frac{1+\varphi^2}{24},
  \qquad
  \|J_\varphi\|_{\mathcal S_2}^2=\frac{1+\varphi^4}{1440}.
  \]
\end{enumerate}
\end{theorem}
\begin{proof}
由定理 \ref{thm:xi-time-part9zbe-foldpi-canonical-compact-jacobi} 的显式谱，
\[
\zeta_{J_\varphi}(s)
=
\sum_{n\ge 1}\left(\frac{1}{4\pi^2n^2}\right)^s
+
\sum_{n\ge 1}\left(\frac{\varphi^2}{4\pi^2n^2}\right)^s
=
(4\pi^2)^{-s}\bigl(1+\varphi^{2s}\bigr)\zeta(2s),
\]
在 \(\Re s>1/2\) 时绝对收敛。

于是
\[
J_\varphi\in \mathcal S_p
\iff
\sum_{n\ge 1}\lambda_n(J_\varphi)^p<\infty
\iff
\zeta(2p)<\infty
\iff
p>\frac12.
\]
这证明第一条。

第二条由 \(\zeta(2s)\) 的标准亚纯延拓立即得到。
在 \(s=1/2\) 邻域，
\[
\zeta(2s)\sim \frac{1}{2(s-1/2)},
\]
故
\[
\operatorname*{Res}_{s=1/2}\zeta_{J_\varphi}(s)
=
(4\pi^2)^{-1/2}(1+\varphi)\cdot \frac12
=
\frac{1+\varphi}{4\pi}
=
\frac{\varphi^2}{4\pi}.
\]

最后把 \(s=1\) 与 \(s=2\) 代入闭式即得
\[
\Tr(J_\varphi)=\zeta_{J_\varphi}(1)
=(4\pi^2)^{-1}(1+\varphi^2)\zeta(2)
=
\frac{1+\varphi^2}{24},
\]
以及
\[
\|J_\varphi\|_{\mathcal S_2}^2
=
\Tr(J_\varphi^2)
=
\zeta_{J_\varphi}(2)
=(4\pi^2)^{-2}(1+\varphi^4)\zeta(4)
=
\frac{1+\varphi^4}{1440}.
\]
证毕。
\end{proof}

\begin{theorem}[golden double-sine Fredholm determinant]\label{thm:xi-time-part9zbe-foldpi-fredholm-double-sine}
由于 \(J_\varphi\) 为 trace-class，Fredholm 行列式
\[
D_\varphi(z):=\det(I-zJ_\varphi)
\]
定义良好，并满足显式闭式
\[
D_\varphi(z)
=
\prod_{n\ge 1}
\left(1-\frac{z}{4\pi^2n^2}\right)
\left(1-\frac{\varphi^2 z}{4\pi^2n^2}\right)
=
\frac{\sin(\sqrt z/2)}{\sqrt z/2}\,
\frac{\sin(\varphi\sqrt z/2)}{\varphi\sqrt z/2}.
\]
因此 \(D_\varphi\) 是一条阶为 \(1/2\) 的 entire function，其零点集恰为
\[
\mathcal Z(D_\varphi)
=
\{4\pi^2n^2:\ n\ge 1\}
\cup
\left\{\frac{4\pi^2n^2}{\varphi^2}:\ n\ge 1\right\}.
\]
\end{theorem}
\begin{proof}
由定理 \ref{thm:xi-time-part9zbe-foldpi-jacobi-schatten-zeta}，\(J_\varphi\in\mathcal S_1\)。
故 Fredholm 行列式可按谱乘积写成
\[
D_\varphi(z)=\prod_{\lambda\in \sigma(J_\varphi)\setminus\{0\}}(1-z\lambda).
\]
再代入定理 \ref{thm:xi-time-part9zbe-foldpi-canonical-compact-jacobi} 的谱式，得到
\[
D_\varphi(z)
=
\prod_{n\ge 1}
\left(1-\frac{z}{4\pi^2n^2}\right)
\left(1-\frac{\varphi^2 z}{4\pi^2n^2}\right).
\]

对 Euler sine product
\[
\frac{\sin x}{x}
=
\prod_{n\ge 1}\left(1-\frac{x^2}{\pi^2n^2}\right)
\]
分别取
\[
x=\frac{\sqrt z}{2},
\qquad
x=\frac{\varphi\sqrt z}{2},
\]
便得到所示闭式。
由显式表达可知 \(D_\varphi\) 是两个阶 \(1/2\) entire functions 的乘积，因此自身也为阶 \(1/2\) 的 entire function。
零点集则直接由两条正弦因子给出。证毕。
\end{proof}

\begin{theorem}[H.111 奇层系数的 compact determinant 采样公式]\label{thm:xi-time-part9zbe-foldpi-compact-determinant-sampling}
对 \(r\ge 1\) 记 odd formal tail 系数
\[
\mathfrak a_r^{(\pi)}
:=
\frac{B_{2r}}{r(2r-1)}
\bigl(\varphi^{-1}+\varphi^{2r-3}\bigr),
\qquad
\mathcal E_\varphi(s)=\sum_{r\ge 1}\mathfrak a_r^{(\pi)}s^{2r-1}.
\]
则在
\[
|z|<\|J_\varphi\|^{-1}=\frac{4\pi^2}{\varphi^2}
\]
内有
\[
-\log D_\varphi(z)
=
\sum_{r\ge 1}\frac{\zeta_{J_\varphi}(r)}{r}\,z^r.
\]
并且对每个 \(r\ge 1\)，成立精确公式
\[
\mathfrak a_r^{(\pi)}
=
(-1)^{r-1}\,
4(2r-2)!\,
\frac{\varphi^{-1}+\varphi^{2r-3}}{1+\varphi^{2r}}\,
\zeta_{J_\varphi}(r),
\]
等价地，
\[
\mathfrak a_r^{(\pi)}
=
(-1)^{r-1}\,
4r(2r-2)!\,
\frac{\varphi^{-1}+\varphi^{2r-3}}{1+\varphi^{2r}}\,
[z^r]\bigl(-\log D_\varphi(z)\bigr).
\]
因此，H.111 的全部 Bernoulli odd layer 不仅形成一个 formal tail，而且事实上是一个 trace-class compact determinant 经 factorial 采样后的精确像。
\end{theorem}
\begin{proof}
由 Fredholm 行列式的对数展开，
\[
-\log D_\varphi(z)
=
-\sum_{\lambda\in \sigma(J_\varphi)\setminus\{0\}}\log(1-z\lambda)
=
\sum_{r\ge 1}\frac{z^r}{r}\sum_{\lambda}\lambda^r
=
\sum_{r\ge 1}\frac{\zeta_{J_\varphi}(r)}{r}\,z^r,
\]
在 \(|z|<\|J_\varphi\|^{-1}\) 内绝对收敛。

另一方面，由 Euler 公式
\[
B_{2r}
=
(-1)^{r-1}\frac{2(2r)!}{(2\pi)^{2r}}\zeta(2r),
\]
可得
\[
\mathfrak a_r^{(\pi)}
=
(-1)^{r-1}\,
4(2r-2)!\,
(4\pi^2)^{-r}
\bigl(\varphi^{-1}+\varphi^{2r-3}\bigr)\zeta(2r).
\]
再由定理 \ref{thm:xi-time-part9zbe-foldpi-jacobi-schatten-zeta}，
\[
\zeta_{J_\varphi}(r)
=(4\pi^2)^{-r}(1+\varphi^{2r})\zeta(2r),
\]
消去 \(\zeta(2r)\) 后便得到
\[
\mathfrak a_r^{(\pi)}
=
(-1)^{r-1}\,
4(2r-2)!\,
\frac{\varphi^{-1}+\varphi^{2r-3}}{1+\varphi^{2r}}\,
\zeta_{J_\varphi}(r).
\]
而
\[
[z^r]\bigl(-\log D_\varphi(z)\bigr)=\frac{\zeta_{J_\varphi}(r)}{r},
\]
代回即得第二式。证毕。
\end{proof}

\begin{conjecture}[exact \texorpdfstring{$C_m$}{Cm} 的 Écalle--Fredholm 提升]\label{conj:xi-time-part9zbe-foldpi-ecalle-fredholm-lift}
定义线性变换
\[
\mathcal T_\varphi\!\left(\sum_{r\ge 1}c_r z^r\right)
:=
\sum_{r\ge 1}
(-1)^{r-1}\,
4r(2r-2)!\,
\frac{\varphi^{-1}+\varphi^{2r-3}}{1+\varphi^{2r}}\,
c_r\,s^{2r-1}.
\]
则由定理 \ref{thm:xi-time-part9zbe-foldpi-compact-determinant-sampling}，在 formal 层面有严格恒等式
\[
\mathcal E_\varphi(s)
=
\mathcal T_\varphi\!\bigl(-\log D_\varphi\bigr)(s).
\]
猜想 H.111 所对应的精确常数序列 \((C_m)_{m\ge 1}\) 不仅拥有该 formal tail，而且满足
\[
\log C_m-\log(2\pi)
=
\left.
\mathscr S_+\!\left[\mathcal T_\varphi\!\bigl(-\log D_\varphi\bigr)\right](s)
\right|_{s=(\varphi/2)^m},
\]
其中 \(\mathscr S_+\) 表示沿正实方向的 canonical Écalle--Borel summation。
若此猜想成立，则项目中的 \(\pi\)-机制便不是 hypergeometric、modular 或 BBP 型 kernel，而是一个 explicit trace-class compact determinant 经过单一 factorial transmutation 后的 resurgent uplift。
\end{conjecture}
\begin{proof}[证据]
定理 \ref{thm:xi-time-part9zbe-foldpi-canonical-compact-jacobi}、
\ref{thm:xi-time-part9zbe-foldpi-jacobi-schatten-zeta}
与 \ref{thm:xi-time-part9zbe-foldpi-fredholm-double-sine}
已经把 H.111 odd tail 的底层对象严格压缩为一个具有 \(n^{-2}\) 特征值衰减、Schatten 临界指数 \(1/2\) 以及黄金双正弦 Fredholm 行列式的紧谱算子。
而定理 \ref{thm:xi-time-part9zbe-foldpi-compact-determinant-sampling} 又说明，全部 odd-layer Bernoulli 系数正是 \(-\log D_\varphi\) 的谱系数经过显式 factorial 变换后的像。

因此，若 exact \((C_m)\) 的余项结构确由 H.111 的 formal tail 唯一控制，则最自然的非形式提升应当不是重新寻找另一条快收敛标量级数，而是把 \(-\log D_\varphi\) 先经 \(\mathcal T_\varphi\) 转写到 odd formal 层，再由 Écalle--Borel 求和恢复 exact 常数序列。
这正是上式所表达的 Fredholm 型 resurgent 提升。证毕。
\end{proof}

\noindent
于是，H.111 的 odd formal tail 在现有审计链下已获得比 “Hausdorff moment sequence” 更强的谱论定位：它来自一个正、紧、自伴 Jacobi 算子 \(J_\varphi\)，其非零谱由两条 reciprocal-square ladders 组成；其 operator \(\zeta\) 完全显式，Fredholm 行列式退化为黄金双正弦乘积，而 odd-layer Bernoulli 系数则恰是 \(-\log D_\varphi\) 的 factorially weighted 采样值。
这说明该项目 \(\pi\)-tail 的真正范畴宿主，并不是另一条标量加速公式，而是一个可压缩到 trace-class compact operator 的谱宿主。

\endinput

\paragraph{fold-\texorpdfstring{$\pi$}{pi} compact Jacobi 宿主的 Wythoff 序、cyclic Nevanlinna 数据与 periodic-tail 排除}
沿用定理 \ref{thm:xi-time-part9zbe-foldpi-canonical-compact-jacobi}、
\ref{thm:xi-time-part9zbe-foldpi-jacobi-counting-weyl}、
\ref{thm:xi-time-part9zbe-foldpi-jacobi-schatten-zeta}
与 \ref{thm:xi-time-part9zbe-foldpi-fredholm-double-sine} 的记号。
记
\[
A_n:=\frac{\varphi^2}{4\pi^2n^2},
\qquad
B_n:=\frac{1}{4\pi^2n^2}
\qquad (n\ge 1),
\]
则
\[
\sigma(J_\varphi)\setminus\{0\}=\{A_n:n\ge 1\}\cup\{B_n:n\ge 1\}.
\]
上一节已经确定了这两条 reciprocal-square ladders、对应的 Fredholm 双正弦行列式以及奇层 Bernoulli 系数的 determinant 采样公式。进一步地，谱的排序、极点的 Nevanlinna 数据以及 periodic-tail 背景的排除也都可以在同一 compact 宿主上被完全显式化。

\begin{theorem}[compact Jacobi 谱的 Wythoff 精确排序]\label{thm:xi-time-part9zbf-foldpi-wythoff-spectral-order}
把 \(J_\varphi\) 的非零特征值按严格降序排列为
\[
\lambda_1>\lambda_2>\cdots>0.
\]
则对每个 \(n\ge 1\)，成立精确索引公式
\[
\lambda_{\lfloor n\varphi\rfloor}=A_n=\frac{\varphi^2}{4\pi^2n^2},
\qquad
\lambda_{\lfloor n\varphi^2\rfloor}=B_n=\frac{1}{4\pi^2n^2}.
\]
因此，\(J_\varphi\) 的有序正谱并非只是两条 reciprocal-square ladders 的并，而是恰由 Wythoff 分拆
\[
\NN
=
\{\lfloor n\varphi\rfloor:\ n\ge 1\}
\sqcup
\{\lfloor n\varphi^2\rfloor:\ n\ge 1\}
\]
严格控制。
\end{theorem}
\begin{proof}
先求 \(A_n\) 的序位。
高于 \(A_n\) 的 \(A\)-谱点恰有 \(n-1\) 个。
另一方面，
\[
B_m>A_n
\iff
\frac{1}{m^2}>\frac{\varphi^2}{n^2}
\iff
m<\frac{n}{\varphi}.
\]
由于 \(n/\varphi\notin\ZZ\)，故高于 \(A_n\) 的 \(B\)-谱点恰有
\[
\left\lfloor \frac{n}{\varphi}\right\rfloor
\]
个。
于是
\[
\operatorname{pos}(A_n)
=
n+\left\lfloor \frac{n}{\varphi}\right\rfloor.
\]
再由
\[
\varphi=1+\varphi^{-1}
\]
以及 \(n/\varphi\notin\ZZ\)，得到
\[
n+\left\lfloor \frac{n}{\varphi}\right\rfloor
=
\lfloor n\varphi\rfloor.
\]
故
\[
\lambda_{\lfloor n\varphi\rfloor}=A_n.
\]

同理，高于 \(B_n\) 的 \(B\)-谱点有 \(n-1\) 个，而
\[
A_m>B_n
\iff
\frac{\varphi^2}{m^2}>\frac{1}{n^2}
\iff
m<\varphi n.
\]
由于 \(\varphi n\notin\ZZ\)，故高于 \(B_n\) 的 \(A\)-谱点恰有
\[
\lfloor \varphi n\rfloor
\]
个。
于是
\[
\operatorname{pos}(B_n)
=
n+\lfloor \varphi n\rfloor
=
\lfloor n\varphi^2\rfloor,
\]
其中使用了
\[
\varphi^2=\varphi+1.
\]
故
\[
\lambda_{\lfloor n\varphi^2\rfloor}=B_n.
\]
两式合起来即得所述 Wythoff 排序。证毕。
\end{proof}

\begin{corollary}[golden double-sine Fredholm 行列式零点的 reciprocal Wythoff 排序]\label{cor:xi-time-part9zbf-foldpi-fredholm-zero-wythoff-order}
把
\[
D_\varphi(z)=\det(I-zJ_\varphi)
\]
的正实零点按严格升序排列为
\[
0<\zeta_1<\zeta_2<\cdots.
\]
则对每个 \(n\ge 1\)，有
\[
\zeta_{\lfloor n\varphi\rfloor}
=
\frac{4\pi^2n^2}{\varphi^2},
\qquad
\zeta_{\lfloor n\varphi^2\rfloor}
=
4\pi^2n^2.
\]
故 \(D_\varphi\) 的零点序并非任意交错的二列平方点集，而是与定理 \ref{thm:xi-time-part9zbf-foldpi-wythoff-spectral-order} 的特征值序完全 reciprocal-locked 的 Wythoff 编码。
\end{corollary}
\begin{proof}
由定理 \ref{thm:xi-time-part9zbe-foldpi-fredholm-double-sine}，
\[
\mathcal Z(D_\varphi)\cap (0,\infty)
=
\left\{\frac{1}{\lambda_k}:\ k\ge 1\right\}.
\]
故
\[
\zeta_k=\lambda_k^{-1}
\qquad (k\ge 1).
\]
将定理 \ref{thm:xi-time-part9zbf-foldpi-wythoff-spectral-order} 的两条索引公式逐项取倒数，便得到
\[
\zeta_{\lfloor n\varphi\rfloor}
=
\frac{4\pi^2n^2}{\varphi^2},
\qquad
\zeta_{\lfloor n\varphi^2\rfloor}
=
4\pi^2n^2.
\]
证毕。
\end{proof}

\begin{theorem}[canonical cyclic Nevanlinna 函数的双 cotangent 闭式与 branchwise 常留数律]\label{thm:xi-time-part9zbf-foldpi-cyclic-nevanlinna-double-cotangent}
定义与循环向量 \(v_0\) 对应的 resolvent 标量函数
\[
m_\varphi(z)
:=
\langle v_0,(I-zJ_\varphi)^{-1}v_0\rangle
=
\sum_{k\ge 0}u_k z^k
=
\int_0^{\rho_\ast}\frac{\dd\mu_\varphi(y)}{1-yz}.
\]
则在
\[
z\in
\CC\setminus
\left(
\{4\pi^2n^2:\ n\ge 1\}
\cup
\left\{\frac{4\pi^2n^2}{\varphi^2}:\ n\ge 1\right\}
\right)
\]
上成立显式闭式
\[
m_\varphi(z)
=
\frac{\varphi^{-1}}{z}
\left(
2-\sqrt z\,\cot\frac{\sqrt z}{2}
\right)
+
\frac{\varphi^{-3}}{z}
\left(
2-\varphi\sqrt z\,\cot\frac{\varphi\sqrt z}{2}
\right),
\]
其中 \(\sqrt z\) 取主支。
等价地，
\[
m_\varphi(z)
=
-2\varphi^{-1}\,\partial_z
\log\!\Bigl(\frac{\sin(\sqrt z/2)}{\sqrt z/2}\Bigr)
-2\varphi^{-3}\,\partial_z
\log\!\Bigl(\frac{\sin(\varphi\sqrt z/2)}{\varphi\sqrt z/2}\Bigr).
\]
特别地，\(m_\varphi\) 的极点与推论 \ref{cor:xi-time-part9zbf-foldpi-fredholm-zero-wythoff-order} 中的零点 obey 相同的 Wythoff 排序，并且两条分支上的留数分别为常数
\[
\operatorname*{Res}_{z=4\pi^2n^2}m_\varphi(z)=-4\varphi^{-1},
\qquad
\operatorname*{Res}_{z=4\pi^2n^2/\varphi^2}m_\varphi(z)=-4\varphi^{-3}.
\]
\end{theorem}
\begin{proof}
由 \(\mu_\varphi\) 的显式原子分解，
\[
m_\varphi(z)
=
\sum_{n\ge 1}\frac{\varphi^{-1}}{\pi^2n^2}
\left(
\frac{1}{1-z/(4\pi^2n^2)}
+
\frac{1}{1-\varphi^2z/(4\pi^2n^2)}
\right).
\]
对第一项，令
\[
a:=\frac{\sqrt z}{2\pi},
\]
则
\[
\frac{1}{1-z/(4\pi^2n^2)}
=
\frac{n^2}{n^2-a^2},
\]
从而
\[
\sum_{n\ge 1}\frac{\varphi^{-1}}{\pi^2n^2}
\frac{1}{1-z/(4\pi^2n^2)}
=
\frac{\varphi^{-1}}{\pi^2}
\sum_{n\ge 1}\frac{1}{n^2-a^2}.
\]
由 \(\cot\) 的标准 Mittag--Leffler 展开，
\[
\sum_{n\ge 1}\frac{1}{n^2-a^2}
=
\frac{1}{2a^2}-\frac{\pi}{2a}\cot(\pi a),
\]
参见 \href{https://dlmf.nist.gov/4.36}{DLMF 4.36} 与 \href{https://dlmf.nist.gov/1.10}{DLMF 1.10}。
代回并用 \(a=\sqrt z/(2\pi)\) 化简，即得
\[
\sum_{n\ge 1}\frac{\varphi^{-1}}{\pi^2n^2}
\frac{1}{1-z/(4\pi^2n^2)}
=
\frac{\varphi^{-1}}{z}
\left(
2-\sqrt z\,\cot\frac{\sqrt z}{2}
\right).
\]

对第二项，令
\[
b:=\frac{\varphi\sqrt z}{2\pi},
\]
同理得到
\[
\sum_{n\ge 1}\frac{\varphi^{-1}}{\pi^2n^2}
\frac{1}{1-\varphi^2z/(4\pi^2n^2)}
=
\frac{\varphi^{-1}}{\pi^2}
\sum_{n\ge 1}\frac{1}{n^2-b^2}
=
\frac{\varphi^{-3}}{z}
\left(
2-\varphi\sqrt z\,\cot\frac{\varphi\sqrt z}{2}
\right).
\]
两式相加即得第一闭式。

对第二闭式，只需分别对
\[
\log\!\Bigl(\frac{\sin(\sqrt z/2)}{\sqrt z/2}\Bigr),
\qquad
\log\!\Bigl(\frac{\sin(\varphi\sqrt z/2)}{\varphi\sqrt z/2}\Bigr)
\]
求导并代回前式即可。

最后，\(m_\varphi\) 的极点正是原子位置的倒数，因而与
\[
\mathcal Z(D_\varphi)\cap (0,\infty)
\]
完全一致；由推论 \ref{cor:xi-time-part9zbf-foldpi-fredholm-zero-wythoff-order} 便得到同一 Wythoff 排序。
而若某原子位置为 \(y\)、质量为 \(w\)，则
\[
\frac{w}{1-yz}
=
-\frac{w/y}{z-y^{-1}}+O(1)
\qquad (z\to y^{-1}),
\]
故相应极点留数为 \(-w/y\)。
代入
\[
y=\frac{1}{4\pi^2n^2},
\qquad
w=\frac{\varphi^{-1}}{\pi^2n^2},
\]
得
\[
\operatorname*{Res}_{z=4\pi^2n^2}m_\varphi(z)
=
-4\varphi^{-1},
\]
再代入
\[
y=\frac{\varphi^2}{4\pi^2n^2},
\qquad
w=\frac{\varphi^{-1}}{\pi^2n^2},
\]
得
\[
\operatorname*{Res}_{z=4\pi^2n^2/\varphi^2}m_\varphi(z)
=
-4\varphi^{-3}.
\]
证毕。
\end{proof}

\begin{theorem}[finite-gap 标量背景的强排除]\label{thm:xi-time-part9zbf-foldpi-no-finite-rank-periodic-tail-bridge}
不存在周期 Jacobi 算子 \(P\) 与有限秩自伴算子 \(K\)，使得
\[
J_\varphi=P+K.
\]
更强地，\(J_\varphi\) 不能是任何 eventually periodic scalar \(J\)-fraction 背景的有限秩修正。
因此，H.111 所导出的 canonical continued-fraction 宿主不属于 finite-gap 或 periodic-tail 范畴。
\end{theorem}
\begin{proof}
由定理 \ref{thm:xi-time-part9zbe-foldpi-canonical-compact-jacobi}，\(J_\varphi\) 为紧自伴算子，故
\[
\sigma_{\mathrm{ess}}(J_\varphi)=\{0\}.
\]
反设
\[
J_\varphi=P+K,
\]
其中 \(P\) 为周期 Jacobi 算子而 \(K\) 为有限秩自伴算子。
有限秩算子是紧算子，而紧扰动保持本质谱不变，因此
\[
\sigma_{\mathrm{ess}}(J_\varphi)=\sigma_{\mathrm{ess}}(P).
\]
另一方面，周期 Jacobi 算子的本质谱是有限条闭区间之并；在真正的周期尾情形下，这一集合不可能退化为单点 \(\{0\}\)。
于是得到矛盾。

若把“周期”放宽为 eventually periodic，则对应算子只是某个周期 Jacobi 算子的有限秩扰动；同样由本质谱不变性，仍与
\[
\sigma_{\mathrm{ess}}(J_\varphi)=\{0\}
\]
矛盾。
故所述 periodic-tail 背景不可能存在。证毕。
\end{proof}

\begin{theorem}[Wythoff--Nevanlinna 刚性]\label{thm:xi-time-part9zbf-foldpi-wythoff-nevanlinna-rigidity}
设 \(J\) 为任意正、紧、自伴 Jacobi 算子，\(v\) 为其循环向量，并定义
\[
m_{J,v}(z):=\langle v,(I-zJ)^{-1}v\rangle.
\]
若 \(m_{J,v}\) 的正实极点按升序排列为
\[
0<p_1<p_2<\cdots,
\]
并满足
\[
p_{\lfloor n\varphi\rfloor}=\frac{4\pi^2n^2}{\varphi^2},
\qquad
p_{\lfloor n\varphi^2\rfloor}=4\pi^2n^2
\qquad (n\ge 1),
\]
且对应留数按两条 Wythoff 分支分别恒为
\[
\operatorname*{Res}_{z=p_{\lfloor n\varphi\rfloor}}m_{J,v}(z)=-4\varphi^{-3},
\qquad
\operatorname*{Res}_{z=p_{\lfloor n\varphi^2\rfloor}}m_{J,v}(z)=-4\varphi^{-1},
\]
则循环对 \((J,v)\) 与 \((J_\varphi,v_0)\) 酉等价；特别地，\(J\) 与 \(J_\varphi\) 酉等价。

换言之，对该 fold-\(\pi\) compact 宿主而言，“两条 reciprocal-square ladders 的 Wythoff 排序与 branchwise 常留数”已经构成完备的有序 Nevanlinna 不变量。
\end{theorem}
\begin{proof}
由 \(J\) 的正紧自伴性，\(m_{J,v}\) 具有纯点 Stieltjes 表示
\[
m_{J,v}(z)=\int\frac{\dd\mu_{J,v}(y)}{1-yz},
\]
其中 \(\mu_{J,v}\) 是支撑于 \([0,\|J\|]\) 的有限正离散测度。
若极点 \(p>0\) 的留数为 \(r<0\)，则对应原子位置为 \(y=p^{-1}\)，且质量满足
\[
r=-\frac{w}{y},
\qquad
w=-\frac{r}{p}.
\]

对第一条 Wythoff 分支，
\[
p_{\lfloor n\varphi\rfloor}=\frac{4\pi^2n^2}{\varphi^2},
\qquad
r_{\lfloor n\varphi\rfloor}=-4\varphi^{-3},
\]
故
\[
w_{\lfloor n\varphi\rfloor}
=
-\frac{r_{\lfloor n\varphi\rfloor}}{p_{\lfloor n\varphi\rfloor}}
=
\frac{\varphi^{-1}}{\pi^2n^2}.
\]
对应原子位置为
\[
y_{\lfloor n\varphi\rfloor}
=
\frac{\varphi^2}{4\pi^2n^2}.
\]

对第二条 Wythoff 分支，
\[
p_{\lfloor n\varphi^2\rfloor}=4\pi^2n^2,
\qquad
r_{\lfloor n\varphi^2\rfloor}=-4\varphi^{-1},
\]
故
\[
w_{\lfloor n\varphi^2\rfloor}
=
-\frac{r_{\lfloor n\varphi^2\rfloor}}{p_{\lfloor n\varphi^2\rfloor}}
=
\frac{\varphi^{-1}}{\pi^2n^2},
\]
而对应原子位置为
\[
y_{\lfloor n\varphi^2\rfloor}
=
\frac{1}{4\pi^2n^2}.
\]

于是
\[
\mu_{J,v}
=
\sum_{n\ge 1}\frac{\varphi^{-1}}{\pi^2n^2}
\left(
\delta_{(4\pi^2n^2)^{-1}}
+
\delta_{\varphi^2(4\pi^2n^2)^{-1}}
\right)
=
\mu_\varphi.
\]
而由定理 \ref{thm:xi-time-part9zbe-foldpi-canonical-compact-jacobi}，\(\mu_\varphi\) 所生成的循环 Jacobi 对唯一。
因此 \((J,v)\) 与 \((J_\varphi,v_0)\) 酉等价。
特别地，\(J\) 与 \(J_\varphi\) 酉等价。证毕。
\end{proof}

\begin{conjecture}[fold-\texorpdfstring{$\pi$}{pi} 机制的有序 Nevanlinna 刚性]\label{conj:xi-time-part9zbf-foldpi-ordered-nevanlinna-rigidity}
若某一外部标量 \(\pi\)-核希望与 H.111 的 odd Bernoulli tail 精确桥接，并同时承载一个 compact spectral host，则其谱论刚性不应仅表现为谱集重现
\[
\left\{\frac{\varphi^2}{4\pi^2n^2}:n\ge 1\right\}
\cup
\left\{\frac{1}{4\pi^2n^2}:n\ge 1\right\},
\]
而应当在更细的有序 Nevanlinna 数据层面同时重现：
\begin{enumerate}
  \item 两条 reciprocal-square ladders 的 Wythoff 排序；
  \item 与该排序兼容的双常留数律；
  \item 由此唯一确定的 cyclic Stieltjes 测度与 Jacobi 宿主。
\end{enumerate}
因此，对本项目的 compact \(\pi\)-宿主而言，真正的谱不变量不在 support equality，而在 ordered Nevanlinna data。
\end{conjecture}
\begin{proof}[证据]
定理 \ref{thm:xi-time-part9zbf-foldpi-wythoff-spectral-order} 把双平方倒数谱梯提升为一条精确的 Wythoff 有序谱序；
推论 \ref{cor:xi-time-part9zbf-foldpi-fredholm-zero-wythoff-order} 与定理 \ref{thm:xi-time-part9zbf-foldpi-cyclic-nevanlinna-double-cotangent} 则把同一排序传递到 Fredholm 零点序与 cyclic Nevanlinna 极点序，并把两条分支上的留数钉死为常数；
最后，定理 \ref{thm:xi-time-part9zbf-foldpi-wythoff-nevanlinna-rigidity} 说明，正是“Wythoff 序位 + 双常留数”这一对 ordered Nevanlinna 数据，把整个 compact Jacobi 宿主唯一恢复出来。

因此，若某外部 \(\pi\)-核只复现谱集，却不能同时复现这一有序极点数据，则它至多与该项目对象谱同胚，而不可能与其 Jacobi 宿主同构。证毕。
\end{proof}

\noindent
由此，H.111 odd formal tail 的 compact 谱宿主在现有审计链下已经不再只由“两个 reciprocal-square 谱梯”刻画；其真正的刚性对象乃是一个更细的三元组：双平方倒数梯、Wythoff 序位与 branchwise 常留数。换言之，决定该 fold-\(\pi\) Jacobi 宿主的，并非单纯的谱集合论数据，而是完整的有序 Nevanlinna 轮廓。

\endinput

\paragraph{fold-\texorpdfstring{$\pi$}{pi} compact 宿主的 finite-rank 耗尽塔与 finite-depth exactness obstruction}
沿用定理 \ref{thm:xi-time-part9zbc-foldpi-oddlayer-positive-moment-formula}、
\ref{thm:xi-time-part9zbc-foldpi-hankel-strict-positivity}、
\ref{thm:xi-time-part9zbe-foldpi-canonical-compact-jacobi}
与 \ref{thm:xi-time-part9zbe-foldpi-compact-determinant-sampling} 的记号。
记
\[
\alpha_r:=\mathfrak a_r^{(\pi)}
=
\frac{B_{2r}}{r(2r-1)}
\bigl(\varphi^{-1}+\varphi^{2r-3}\bigr),
\qquad
u_k:=(-1)^k\frac{\alpha_{k+1}}{(2k)!}
\qquad (r\ge 1,\ k\ge 0).
\]
由前述结果，\((u_k)_{k\ge 0}\) 正是 compact Jacobi 宿主 \(J_\varphi\) 的矩列，因而 H.111 的 odd Bernoulli tail 不仅对应一个无限秩谱宿主，而且还可被一列 finite-rank scalar hosts 从下方单调耗尽。关键点在于：这一耗尽塔对每一层都严格逼近，但在任何有限深度都不会稳定。

\begin{theorem}[H.111 奇层矩的 finite-rank 耗尽塔]\label{thm:xi-time-part9zbg-foldpi-finite-rank-moment-exhaustion}
对每个 \(N\ge 1\)，定义截断测度
\[
\mu_{\varphi,N}
:=
\sum_{n=1}^{N}\frac{\varphi^{-1}}{\pi^2n^2}
\left(
\delta_{(4\pi^2n^2)^{-1}}
+
\delta_{\varphi^2(4\pi^2n^2)^{-1}}
\right),
\]
以及其矩
\[
u_k^{(N)}:=\int y^k\,\dd\mu_{\varphi,N}(y)
\qquad (k\ge 0).
\]
则对任意 \(k\ge 0\)，有显式闭式
\[
u_k^{(N)}
=
4(2\pi)^{-2k-2}
\bigl(\varphi^{-1}+\varphi^{2k-1}\bigr)
\sum_{n=1}^{N}n^{-2k-2},
\]
并且
\[
0<u_k^{(1)}<u_k^{(2)}<\cdots<u_k^{(N)}\uparrow u_k.
\]
更精确地，
\[
u_k-u_k^{(N)}
=
4(2\pi)^{-2k-2}
\bigl(\varphi^{-1}+\varphi^{2k-1}\bigr)
\sum_{n>N}n^{-2k-2}.
\]
\leanverified{paper\_xi\_time\_part9zbg\_foldpi\_finite\_rank\_moment\_exhaustion}
\end{theorem}
\begin{proof}
由 \(\mu_{\varphi,N}\) 的定义，
\[
u_k^{(N)}
=
\sum_{n=1}^{N}\frac{\varphi^{-1}}{\pi^2n^2}
\left(
\frac{1}{(4\pi^2n^2)^k}
+
\frac{\varphi^{2k}}{(4\pi^2n^2)^k}
\right).
\]
整理即得
\[
u_k^{(N)}
=
4(2\pi)^{-2k-2}
\bigl(\varphi^{-1}+\varphi^{2k-1}\bigr)
\sum_{n=1}^{N}n^{-2k-2}.
\]
又因为
\[
\mu_{\varphi,N+1}-\mu_{\varphi,N}
=
\frac{\varphi^{-1}}{\pi^2(N+1)^2}
\left(
\delta_{(4\pi^2(N+1)^2)^{-1}}
+
\delta_{\varphi^2(4\pi^2(N+1)^2)^{-1}}
\right)
\]
是非零正测度，故对每个 \(k\ge 0\) 都有
\[
u_k^{(N+1)}-u_k^{(N)}
=
\int y^k\,\dd(\mu_{\varphi,N+1}-\mu_{\varphi,N})(y)>0.
\]
于是 \((u_k^{(N)})_{N\ge 1}\) 严格递增。
其上界 \(u_k\) 则由定理 \ref{thm:xi-time-part9zbc-foldpi-oddlayer-positive-moment-formula} 给出，而极限与尾项公式都由
\[
\mu_\varphi-\mu_{\varphi,N}
=
\sum_{n>N}\frac{\varphi^{-1}}{\pi^2n^2}
\left(
\delta_{(4\pi^2n^2)^{-1}}
+
\delta_{\varphi^2(4\pi^2n^2)^{-1}}
\right)
\]
直接得到。证毕。
\end{proof}

\begin{theorem}[finite-depth scalar host 的 Jacobi--Fredholm 实现]\label{thm:xi-time-part9zbg-foldpi-finite-rank-jacobi-fredholm-realization}
对每个 \(N\ge 1\)，存在唯一的 \(2N\times 2N\) Jacobi 矩阵
\[
J_{\varphi,N}
=
\begin{pmatrix}
\beta_0^{(N)} & \alpha_1^{(N)} \\
\alpha_1^{(N)} & \beta_1^{(N)} & \alpha_2^{(N)} \\
& \ddots & \ddots & \ddots \\
&& \alpha_{2N-1}^{(N)} & \beta_{2N-1}^{(N)}
\end{pmatrix},
\qquad
\alpha_j^{(N)}>0,\ \beta_j^{(N)}\in\RR,
\]
以及循环向量 \(v_0^{(N)}\)，使得
\[
u_k^{(N)}=\langle v_0^{(N)},(J_{\varphi,N})^k v_0^{(N)}\rangle
\qquad (k\ge 0).
\]
其谱精确为
\[
\sigma(J_{\varphi,N})
=
\left\{\frac{1}{4\pi^2n^2}:1\le n\le N\right\}
\cup
\left\{\frac{\varphi^2}{4\pi^2n^2}:1\le n\le N\right\}.
\]
相应的 cyclic resolvent 为
\[
m_{\varphi,N}(z)
:=
\langle v_0^{(N)},(I-zJ_{\varphi,N})^{-1}v_0^{(N)}\rangle
=
\int\frac{\dd\mu_{\varphi,N}(y)}{1-yz},
\]
即
\[
m_{\varphi,N}(z)
=
\sum_{n=1}^{N}\frac{\varphi^{-1}}{\pi^2n^2}
\left(
\frac{1}{1-z/(4\pi^2n^2)}
+
\frac{1}{1-\varphi^2z/(4\pi^2n^2)}
\right).
\]
因此 \(m_{\varphi,N}\) 是恰有 \(2N\) 个极点的有理 Stieltjes 函数，其 canonical scalar continued fraction 属于 terminating Stieltjes class；关于有限点支撑正交测度与对应 Jacobi/continued-fraction 实现的标准背景，可参见 \href{https://dlmf.nist.gov/18.2}{DLMF 18.2}。
\end{theorem}
\begin{proof}
在有限维 Hilbert 空间 \(L^2(\mu_{\varphi,N})\) 上考虑乘法算子
\[
(M_yf)(y):=y f(y).
\]
由于 \(\mu_{\varphi,N}\) 支撑于 \(2N\) 个互异点，\(L^2(\mu_{\varphi,N})\) 的维数即为 \(2N\)，且
\[
1,\ y,\ \dots,\ y^{2N-1}
\]
张成整个空间。
对这组函数作 Gram--Schmidt 正交化，得到一组正交归一多项式基。
在该基下，\(M_y\) 表现为一唯一的 \(2N\times 2N\) Jacobi 矩阵 \(J_{\varphi,N}\)。

若 \(U_N:L^2(\mu_{\varphi,N})\to \RR^{2N}\) 把这组正交基送到标准基，则
\[
J_{\varphi,N}=U_NM_yU_N^{-1},
\qquad
v_0^{(N)}:=U_N(1).
\]
于是
\[
u_k^{(N)}
=
\langle 1,M_y^k1\rangle_{L^2(\mu_{\varphi,N})}
=
\langle v_0^{(N)},(J_{\varphi,N})^k v_0^{(N)}\rangle.
\]
由于纯点乘法算子的谱恰为支撑点集，故
\[
\sigma(J_{\varphi,N})=\operatorname{supp}\mu_{\varphi,N},
\]
即为所述两条有限 reciprocal-square ladders。
最后，
\[
m_{\varphi,N}(z)
=
\langle 1,(I-zM_y)^{-1}1\rangle_{L^2(\mu_{\varphi,N})}
=
\int\frac{\dd\mu_{\varphi,N}(y)}{1-yz},
\]
再把 \(\mu_{\varphi,N}\) 展开即可得到有理 Stieltjes 闭式。
证毕。
\end{proof}

\begin{theorem}[黄金双正弦行列式的 finite-rank 局部一致截断]\label{thm:xi-time-part9zbg-foldpi-finite-rank-fredholm-local-uniform}
定义
\[
D_{\varphi,N}(z)
:=
\det(I-zJ_{\varphi,N})
=
\prod_{n=1}^{N}
\left(1-\frac{z}{4\pi^2n^2}\right)
\left(1-\frac{\varphi^2 z}{4\pi^2n^2}\right).
\]
则 \(D_{\varphi,N}\to D_\varphi\) 在 \(\CC\) 的每个紧集上一致收敛。
更精确地，若 \(K\Subset \CC\setminus \mathcal Z(D_\varphi)\) 且
\[
R:=\sup_{z\in K}|z|,
\]
则存在 \(N_0(K)\) 使得对一切 \(N\ge N_0(K)\)，可在 \(K\) 上取一致的对数分支并有
\[
\sup_{z\in K}
\left|
\log D_\varphi(z)-\log D_{\varphi,N}(z)
\right|
\le
\frac{R(1+\varphi^2)}{2\pi^2}\,\frac{1}{N}.
\]
同时，对每个整数 \(r\ge 1\)，有限秩谱 \(\zeta\) 函数满足
\[
\zeta_{J_{\varphi,N}}(r):=\Tr(J_{\varphi,N}^r)
=
(4\pi^2)^{-r}(1+\varphi^{2r})\sum_{n=1}^{N}n^{-2r},
\]
并且
\[
0<
\zeta_{J_\varphi}(r)-\zeta_{J_{\varphi,N}}(r)
\le
(4\pi^2)^{-r}(1+\varphi^{2r})\frac{N^{1-2r}}{2r-1}.
\]
\end{theorem}
\begin{proof}
第一条乘积公式由定理 \ref{thm:xi-time-part9zbg-foldpi-finite-rank-jacobi-fredholm-realization} 的显式谱立即得到。
对紧集 \(K\subset\CC\)，记 \(R:=\sup_{z\in K}|z|\)。
则
\[
\sum_{n>N}\sup_{z\in K}
\left(
\left|\frac{z}{4\pi^2n^2}\right|
+
\left|\frac{\varphi^2 z}{4\pi^2n^2}\right|
\right)
\le
\frac{R(1+\varphi^2)}{4\pi^2}
\sum_{n>N}n^{-2},
\]
而右端随 \(N\to\infty\) 收敛到 \(0\)。
于是尾积
\[
\prod_{n>N}
\left(1-\frac{z}{4\pi^2n^2}\right)
\left(1-\frac{\varphi^2 z}{4\pi^2n^2}\right)
\]
在 \(K\) 上一致收敛到 \(1\)，故 \(D_{\varphi,N}\to D_\varphi\) 局部一致。

若进一步 \(K\Subset\CC\setminus \mathcal Z(D_\varphi)\)，则可取 \(N_0(K)\) 使对 \(N\ge N_0(K)\) 与一切 \(n>N\) 都有
\[
\sup_{z\in K}
\left(
\left|\frac{z}{4\pi^2n^2}\right|
\vee
\left|\frac{\varphi^2 z}{4\pi^2n^2}\right|
\right)
\le
\frac12.
\]
于是利用
\[
|\log(1-w)|\le 2|w|
\qquad (|w|\le 1/2),
\]
便得到
\[
\sup_{z\in K}
\left|
\log\frac{D_\varphi(z)}{D_{\varphi,N}(z)}
\right|
\le
2R\sum_{n>N}
\left(
\frac{1}{4\pi^2n^2}
+
\frac{\varphi^2}{4\pi^2n^2}
\right)
\le
\frac{R(1+\varphi^2)}{2\pi^2}\,\frac1N.
\]

对谱 \(\zeta\) 函数，只需把定理 \ref{thm:xi-time-part9zbg-foldpi-finite-rank-jacobi-fredholm-realization} 的显式谱代入，得到
\[
\zeta_{J_{\varphi,N}}(r)
=
\sum_{n=1}^{N}\left(\frac{1}{4\pi^2n^2}\right)^r
+
\sum_{n=1}^{N}\left(\frac{\varphi^2}{4\pi^2n^2}\right)^r
=
(4\pi^2)^{-r}(1+\varphi^{2r})\sum_{n=1}^{N}n^{-2r}.
\]
而
\[
\zeta_{J_\varphi}(r)-\zeta_{J_{\varphi,N}}(r)
=
(4\pi^2)^{-r}(1+\varphi^{2r})\sum_{n>N}n^{-2r}.
\]
再用积分判别
\[
\sum_{n>N}n^{-2r}
\le
\int_N^\infty x^{-2r}\,\dd x
=
\frac{N^{1-2r}}{2r-1}
\]
即得尾项估计。证毕。
\end{proof}

\begin{theorem}[H.111 Bernoulli 奇层的 finite-rank 单调逼近]\label{thm:xi-time-part9zbg-foldpi-finite-rank-oddlayer-monotone-approx}
\leanverified{paper\_xi\_time\_part9zbg\_foldpi\_finite\_rank\_oddlayer\_monotone\_approx}
对每个 \(r\ge 1\) 与 \(N\ge 1\)，定义 finite-rank odd-layer 系数
\[
\alpha_r^{(N)}
:=
(-1)^{r-1}\,
4(2r-2)!\,
\frac{\varphi^{-1}+\varphi^{2r-3}}{1+\varphi^{2r}}\,
\zeta_{J_{\varphi,N}}(r).
\]
则有显式闭式
\[
\alpha_r^{(N)}
=
(-1)^{r-1}\,
4(2r-2)!\,(4\pi^2)^{-r}
\bigl(\varphi^{-1}+\varphi^{2r-3}\bigr)
\sum_{n=1}^{N}n^{-2r},
\]
并且对每个固定 \(r\) 都有
\[
0<
(-1)^{r-1}\alpha_r^{(1)}
<
(-1)^{r-1}\alpha_r^{(2)}
<
\cdots
<
(-1)^{r-1}\alpha_r,
\]
其中 \(\alpha_r\) 正是 H.111 的 exact odd-layer coefficient。
更精确地，
\[
0<
(-1)^{r-1}\bigl(\alpha_r-\alpha_r^{(N)}\bigr)
\le
\frac{4(2r-2)!}{(4\pi^2)^r}
\frac{\varphi^{-1}+\varphi^{2r-3}}{2r-1}\,
N^{1-2r}.
\]
\end{theorem}
\begin{proof}
将定理 \ref{thm:xi-time-part9zbg-foldpi-finite-rank-fredholm-local-uniform} 的显式公式
\[
\zeta_{J_{\varphi,N}}(r)
=
(4\pi^2)^{-r}(1+\varphi^{2r})\sum_{n=1}^{N}n^{-2r}
\]
代入定义即可得到 \(\alpha_r^{(N)}\) 的闭式。

另一方面，由定理 \ref{thm:xi-time-part9zbe-foldpi-compact-determinant-sampling}，
\[
\alpha_r
=
(-1)^{r-1}\,
4(2r-2)!\,
\frac{\varphi^{-1}+\varphi^{2r-3}}{1+\varphi^{2r}}\,
\zeta_{J_\varphi}(r).
\]
再由定理 \ref{thm:xi-time-part9zbg-foldpi-finite-rank-fredholm-local-uniform} 中的严格尾差
\[
\zeta_{J_\varphi}(r)-\zeta_{J_{\varphi,N}}(r)>0,
\]
立得
\[
0<
(-1)^{r-1}\bigl(\alpha_r-\alpha_r^{(N)}\bigr).
\]
同时利用同一定理中的尾界便得到
\[
0<
(-1)^{r-1}\bigl(\alpha_r-\alpha_r^{(N)}\bigr)
\le
\frac{4(2r-2)!}{(4\pi^2)^r}
\frac{\varphi^{-1}+\varphi^{2r-3}}{2r-1}\,
N^{1-2r}.
\]
故 \(((-1)^{r-1}\alpha_r^{(N)})_{N\ge 1}\) 对固定 \(r\) 严格递增并收敛到 \((-1)^{r-1}\alpha_r\)。
证毕。
\end{proof}

\begin{theorem}[finite-rank 塔中的 Hankel 正性阈值与单调收敛]\label{thm:xi-time-part9zbg-foldpi-finite-rank-hankel-threshold}
对固定 \(n\ge 0\)，定义截断 Hankel 矩阵
\[
H_n^{(N)}:=(u_{i+j}^{(N)})_{0\le i,j\le n},
\qquad
H_n:=(u_{i+j})_{0\le i,j\le n}.
\]
则有
\[
H_n^{(N)}\preceq H_n^{(N+1)}\preceq H_n
\qquad (\forall N\ge 1),
\]
并且只要
\[
2N\ge n+1,
\]
就有 \(H_n^{(N)}\) 严格正定，从而
\[
\det H_n^{(N)}>0.
\]
此外，
\[
\det H_n^{(N)}\uparrow \det H_n>0
\qquad (N\to\infty).
\]
\end{theorem}
\begin{proof}
令
\[
v_n(y):=(1,y,\dots,y^n)^{\mathsf T}.
\]
则
\[
H_n^{(N+1)}-H_n^{(N)}
=
\int v_n(y)v_n(y)^{\mathsf T}\,\dd(\mu_{\varphi,N+1}-\mu_{\varphi,N})(y)\succeq 0,
\]
同理
\[
H_n-H_n^{(N)}
=
\int v_n(y)v_n(y)^{\mathsf T}\,\dd(\mu_\varphi-\mu_{\varphi,N})(y)\succeq 0.
\]
故得到 Loewner 单调性
\[
H_n^{(N)}\preceq H_n^{(N+1)}\preceq H_n.
\]

若 \(2N\ge n+1\)，则 \(\mu_{\varphi,N}\) 至少含有 \(n+1\) 个互异支撑点。
任取 \(c=(c_0,\dots,c_n)\neq 0\)，令
\[
P_c(y):=\sum_{j=0}^{n}c_j y^j.
\]
由于 \(\deg P_c\le n\)，该非零多项式不可能在 \(n+1\) 个互异点上同时为零，故
\[
c^{\mathsf T}H_n^{(N)}c
=
\int P_c(y)^2\,\dd\mu_{\varphi,N}(y)>0.
\]
从而 \(H_n^{(N)}\) 严格正定。

最后，由定理 \ref{thm:xi-time-part9zbg-foldpi-finite-rank-moment-exhaustion}，每个矩条目
\[
u_{i+j}^{(N)}\uparrow u_{i+j},
\]
故
\[
H_n^{(N)}\longrightarrow H_n
\]
逐项收敛。
由行列式的连续性与上面的 Loewner 单调性，得到
\[
\det H_n^{(N)}\uparrow \det H_n.
\]
而 \(\det H_n>0\) 则由定理 \ref{thm:xi-time-part9zbc-foldpi-hankel-strict-positivity} 给出。
证毕。
\end{proof}

\begin{theorem}[finite-depth exactness obstruction]\label{thm:xi-time-part9zbg-foldpi-finite-depth-exactness-obstruction}
\leanverified{paper\_xi\_time\_part9zbg\_foldpi\_finite\_depth\_exactness\_obstruction}
不存在任何 finite-rank scalar spectral host，能够精确承载 H.111 的全体 odd Bernoulli layers。
更具体地，对任意有限 \(N\) 与任意 \(r\ge 1\)，恒有严格不等式
\[
(-1)^{r-1}\alpha_r^{(N)}<(-1)^{r-1}\alpha_r.
\]
因此，exact projective \(\pi\)-tail 不是某个 finite continued fraction 或 finite Jacobi kernel 的改写；其谱宿主必然是无限秩对象。
\end{theorem}
\begin{proof}
前述严格不等式已经由定理 \ref{thm:xi-time-part9zbg-foldpi-finite-rank-oddlayer-monotone-approx} 给出。

更强地，反设存在某个 finite-rank scalar spectral host \(J\) 与循环向量 \(v\)，能够精确承载 H.111 的全部 odd Bernoulli layers。
则通过定理 \ref{thm:xi-time-part9zbe-foldpi-compact-determinant-sampling} 所建立的同一 factorial 采样关系，可得到与 exact tail 对应的矩列仍然必须是
\[
u_k=(-1)^k\frac{\alpha_{k+1}}{(2k)!}
\qquad (k\ge 0).
\]
然而 finite-rank 正自伴 Jacobi 宿主只对应有限支撑测度，因此其矩列的 Hankel 秩有限。
这与定理 \ref{thm:xi-time-part9zbc-foldpi-hankel-strict-positivity} 以及定理 \ref{thm:xi-time-part9zbc-foldpi-zeta-rank-jump} 所给出的“\((u_k)\) 具有无限 Hankel 秩且全部有限 Hankel 主子式严格为正”直接矛盾。
故 exact host 不可能是 finite-rank 对象。

于是，虽然 \((J_{\varphi,N})_{N\ge 1}\) 构成一座从下方单调逼近 exact odd tail 的 finite-depth 塔，但这一塔在任何有限深度都不会稳定为 H.111 的 exact object。
证毕。
\end{proof}

\noindent
因此，H.111 odd Bernoulli tail 在现有审计链下不只是“属于 compact Jacobi/Fredholm 范畴”；更准确地说，它是一个由 finite-depth scalar hosts 严格单调耗尽、却永不在任何有限深度闭合的无限秩对象。每个 \(\mu_{\varphi,N}\) 只给出一个 terminating Stieltjes kernel，每个 \(J_{\varphi,N}\) 只携带有限条 reciprocal-square ladders，每个 \(D_{\varphi,N}\) 只拥有有限零点集，而每个 \(\alpha_r^{(N)}\) 只提供 exact odd tail 的严格下侧逼近。唯有在 \(N\to\infty\) 的严格归纳极限中，exact \(2\pi\)-恢复与其全部 Bernoulli--\(\zeta(2r)\) 谱信息才同时出现。

\endinput

\paragraph{fold-\texorpdfstring{$\pi$}{pi} compact 宿主的 quarter-line completion、RH-blind 局部影子与 positive determinant cone}
沿用定理 \ref{thm:xi-time-part9zbe-foldpi-jacobi-schatten-zeta}、
\ref{thm:xi-time-part9zbe-foldpi-fredholm-double-sine}、
\ref{thm:xi-time-part9zbe-foldpi-compact-determinant-sampling}
与定理 \ref{thm:xi-time-part9zbg-foldpi-finite-depth-exactness-obstruction} 的记号。
上一节已经把 H.111 的 odd Bernoulli tail 压缩为一个 canonical compact Jacobi 宿主，并证明了任何 finite-depth scalar host 都只能给出严格下侧逼近。
进一步地，这个 compact 宿主一方面具有完全显式的 \(\zeta\)-正则化谱不变量，另一方面又允许把 RH 精确改写为一个 quarter-line symmetric completion 的正 determinant 判据。
因此，H.111 的 odd tail 在该链条中同时呈现出两条互补特征：finite local data 完全 RH-blind，而全局 completion 则落入一个严格的 positive trace-class determinant cone。

\begin{theorem}[canonical compact Jacobi 宿主的 \texorpdfstring{$\zeta$}{zeta}-正则化行列式]\label{thm:xi-time-part9zbh-foldpi-zeta-regularized-determinant}
\leanverified{paper\_xi\_time\_part9zbh\_foldpi\_zeta\_regularized\_determinant}
设
\[
\det_\zeta(J_\varphi):=\exp\!\bigl(-\zeta_{J_\varphi}'(0)\bigr),
\]
其中
\[
\zeta_{J_\varphi}(s)=\operatorname{Tr}(J_\varphi^s)
=(4\pi^2)^{-s}(1+\varphi^{2s})\zeta(2s).
\]
则有严格恒等式
\[
\det_\zeta(J_\varphi)=\varphi.
\]
\end{theorem}
\begin{proof}
令
\[
A(s):=(4\pi^2)^{-s}(1+\varphi^{2s}),
\]
则
\[
\zeta_{J_\varphi}(s)=A(s)\zeta(2s).
\]
于是
\[
\zeta_{J_\varphi}'(0)
=
A'(0)\zeta(0)+A(0)\cdot 2\zeta'(0).
\]
直接计算得到
\[
A(0)=2,
\qquad
A'(0)=-2\log(4\pi^2)+2\log\varphi.
\]
再代入经典值
\[
\zeta(0)=-\frac12,
\qquad
\zeta'(0)=-\frac12\log(2\pi),
\]
便有
\[
\zeta_{J_\varphi}'(0)
=
\bigl[-2\log(4\pi^2)+2\log\varphi\bigr]\Bigl(-\frac12\Bigr)
+2\Bigl(-\frac12\log(2\pi)\Bigr)\cdot 2
=
-\log\varphi.
\]
因此
\[
\det_\zeta(J_\varphi)
=
\exp\!\bigl(-\zeta_{J_\varphi}'(0)\bigr)
=
\varphi.
\]
证毕。有关 \(\zeta(0)\)、\(\zeta'(0)\) 与 completed \(\xi\)-函数的标准记号，可参见 \href{https://dlmf.nist.gov/25.6}{DLMF 25.6}。
\end{proof}

\begin{theorem}[quarter-line completed spectral host 的对称化]\label{thm:xi-time-part9zbh-foldpi-quarter-line-completion}
\leanverified{paper\_xi\_time\_part9zbh\_foldpi\_quarter\_line\_completion}
定义
\[
\xi(s):=\frac12\,s(s-1)\pi^{-s/2}\Gamma\!\Bigl(\frac s2\Bigr)\zeta(s),
\]
以及
\[
\Xi_\varphi(s):=(1+\varphi^{2s})(1+\varphi^{1-2s})\,\xi(2s).
\]
则 \(\Xi_\varphi\) 为整函数，并满足精确对称律
\[
\Xi_\varphi(s)=\Xi_\varphi\!\Bigl(\frac12-s\Bigr).
\]
此外，显式因子恰产生两列算术零点
\[
\mathcal A_\varphi^{(0)}
:=
\left\{
\frac{(2k+1)\pi i}{2\log\varphi}:\ k\in\ZZ
\right\},
\]
\[
\mathcal A_\varphi^{(1/2)}
:=
\left\{
\frac12+\frac{(2k+1)\pi i}{2\log\varphi}:\ k\in\ZZ
\right\},
\]
它们分别位于 \(\Re s=0\) 与 \(\Re s=\frac12\) 上，并互为关于 \(\Re s=\frac14\) 的镜像。
\end{theorem}
\begin{proof}
由 \(\xi(s)=\xi(1-s)\) 可得
\[
\Xi_\varphi\!\Bigl(\frac12-s\Bigr)
=
(1+\varphi^{1-2s})(1+\varphi^{2s})\,\xi(1-2s)
=
\Xi_\varphi(s).
\]
整性由 \(\xi\) 的整性与两个显式指数因子的整性立即得到。

若 \(1+\varphi^{2s}=0\)，则
\[
\varphi^{2s}=-1=e^{(2k+1)\pi i},
\]
因此
\[
s=\frac{(2k+1)\pi i}{2\log\varphi}\in \mathcal A_\varphi^{(0)}.
\]
同理，若 \(1+\varphi^{1-2s}=0\)，则
\[
1-2s=\frac{(2k+1)\pi i}{\log\varphi},
\]
从而
\[
s=\frac12+\frac{(2k+1)\pi i}{2\log\varphi}\in \mathcal A_\varphi^{(1/2)}.
\]
这两列零点关于 \(s\mapsto \frac12-s\) 互相对换，故亦关于 \(\Re s=\frac14\) 镜像。
证毕。有关 \(\xi\)-函数的函数方程与零点分布背景，可参见 \href{https://dlmf.nist.gov/25.4}{DLMF 25.4} 与 \href{https://dlmf.nist.gov/25.10}{DLMF 25.10}。
\end{proof}

\begin{theorem}[quarter-line centered completion 的精确分解]\label{thm:xi-time-part9zbh-foldpi-centered-completion-factorization}
记
\[
L:=\log\varphi,
\qquad
\Theta_\varphi(z):=\Xi_\varphi\!\Bigl(\frac14+z\Bigr).
\]
则有严格恒等式
\[
\Theta_\varphi(z)
=
4\varphi^{1/2}
\cosh\!\Bigl(L\bigl(z+\tfrac14\bigr)\Bigr)
\cosh\!\Bigl(L\bigl(z-\tfrac14\bigr)\Bigr)
\,\xi\!\Bigl(\frac12+2z\Bigr).
\]
相应地，去算术化后的 centered completion
\[
\Theta_\varphi^\sharp(z)
:=
\frac{\Theta_\varphi(z)}
{4\varphi^{1/2}\cosh(L(z+\tfrac14))\cosh(L(z-\tfrac14))}
=
\xi\!\Bigl(\frac12+2z\Bigr)
\]
是偶整函数：
\[
\Theta_\varphi^\sharp(z)=\Theta_\varphi^\sharp(-z).
\]
\end{theorem}
\begin{proof}
由定义，
\[
\Theta_\varphi(z)
=
\bigl(1+\varphi^{1/2+2z}\bigr)
\bigl(1+\varphi^{1/2-2z}\bigr)
\,\xi\!\Bigl(\frac12+2z\Bigr).
\]
先化简显式因子。
注意到
\[
1+\varphi^{1/2}e^{2Lz}
=
2\varphi^{1/4}e^{Lz}\cosh\!\Bigl(L\bigl(z+\tfrac14\bigr)\Bigr),
\]
而
\[
1+\varphi^{1/2}e^{-2Lz}
=
2\varphi^{1/4}e^{-Lz}\cosh\!\Bigl(L\bigl(z-\tfrac14\bigr)\Bigr).
\]
两式相乘便得到
\[
\bigl(1+\varphi^{1/2+2z}\bigr)\bigl(1+\varphi^{1/2-2z}\bigr)
=
4\varphi^{1/2}
\cosh\!\Bigl(L\bigl(z+\tfrac14\bigr)\Bigr)
\cosh\!\Bigl(L\bigl(z-\tfrac14\bigr)\Bigr).
\]
因此
\[
\Theta_\varphi^\sharp(z)=\xi\!\Bigl(\frac12+2z\Bigr).
\]
再由 \(\xi(s)=\xi(1-s)\)，有
\[
\Theta_\varphi^\sharp(-z)
=
\xi\!\Bigl(\frac12-2z\Bigr)
=
\xi\!\Bigl(\frac12+2z\Bigr)
=
\Theta_\varphi^\sharp(z).
\]
故 \(\Theta_\varphi^\sharp\) 为偶整函数。证毕。
\end{proof}

\begin{theorem}[RH 的 quarter-line 与 centered-imaginary 精确改写]\label{thm:xi-time-part9zbh-foldpi-rh-quarter-line-equivalence}
\leanverified{paper\_xi\_time\_part9zbh\_foldpi\_rh\_quarter\_line\_equivalence}
下列命题严格等价：
\begin{enumerate}
  \item Riemann 假设成立；
  \item \(\Xi_\varphi\) 的全部零点中，除去显式算术零点
  \[
  \mathcal A_\varphi^{(0)}\cup \mathcal A_\varphi^{(1/2)},
  \]
  其余零点全部落在 quarter-line \(\Re s=\frac14\) 上；
  \item \(\Theta_\varphi^\sharp(z)=\xi(\frac12+2z)\) 的全部零点都落在虚轴上。
\end{enumerate}
\end{theorem}
\begin{proof}
\((1)\Leftrightarrow(2)\) 由定理 \ref{thm:xi-time-part9zbh-foldpi-quarter-line-completion} 直接得到。
事实上，\(\xi(2s)\) 的零点恰为
\[
s=\frac{\rho}{2},
\]
其中 \(\rho\) 取遍 \(\zeta(s)\) 的全部非平凡零点。
因此
\[
\Re \rho=\frac12
\iff
\Re \frac{\rho}{2}=\frac14.
\]
而 \(\Xi_\varphi\) 除 \(\xi(2s)\) 的零点之外，其余零点恰由显式因子给出，即
\[
\mathcal A_\varphi^{(0)}\cup \mathcal A_\varphi^{(1/2)}.
\]
于是 RH 当且仅当剔除这两列显式算术零点后，\(\Xi_\varphi\) 的其余零点全部落在 \(\Re s=\frac14\) 上。

另一方面，由定理 \ref{thm:xi-time-part9zbh-foldpi-centered-completion-factorization}，
\[
\Theta_\varphi^\sharp(z)=\xi\!\Bigl(\frac12+2z\Bigr).
\]
故 \(\Theta_\varphi^\sharp(z)\) 的零点恰为
\[
z=\frac{\rho-\frac12}{2}.
\]
因此
\[
\Re \rho=\frac12
\iff
\Re z=0.
\]
这便给出 \((1)\Leftrightarrow(3)\)。证毕。
\end{proof}

\begin{theorem}[H.111 odd tail 的 finite local data 的 RH-blindness]\label{thm:xi-time-part9zbh-foldpi-finite-local-rh-blindness}
\leanverified{paper\_xi\_time\_part9zbh\_foldpi\_finite\_local\_rh\_blindness}
任取固定整数 \(R,N\ge 1\)。
则下列有限数据都仅依赖于有限多个显式偶值
\[
\zeta(2),\zeta(4),\dots,\zeta(2M)
\]
中的一段，因而不包含 RH-sensitive 信息：
\begin{enumerate}
  \item odd tail 的前 \(R\) 个 Bernoulli 层系数 \(\mathfrak a_1^{(\pi)},\dots,\mathfrak a_R^{(\pi)}\)；
  \item 前 \(2N\) 个矩 \(u_0,\dots,u_{2N-1}\)；
  \item 前 \(N\) 个 Jacobi 参数及其对应的任意有限 principal Jacobi section；
  \item \(m_\varphi(z)\) 的任意有限阶 Pad\'e 或 terminating \(S\)-fraction 截断；
  \item \(-\log D_\varphi(z)\) 的任意有限 Taylor jet。
\end{enumerate}
特别地，任何仅由这些 finite local data 所构造的判据都不可能与 RH 等价。
\end{theorem}
\begin{proof}
由定理 \ref{thm:xi-time-part9zbe-foldpi-compact-determinant-sampling}，
\[
\mathfrak a_r^{(\pi)}
=
\frac{B_{2r}}{r(2r-1)}
\bigl(\varphi^{-1}+\varphi^{2r-3}\bigr),
\]
故前 \(R\) 个奇层系数只依赖于
\[
B_2,B_4,\dots,B_{2R}.
\]
再由 Euler 的偶值公式
\[
\zeta(2n)=(-1)^{n+1}\frac{B_{2n}(2\pi)^{2n}}{2(2n)!},
\]
可知这些系数只依赖于显式偶值
\[
\zeta(2),\dots,\zeta(2R).
\]

由
\[
u_k=(-1)^k\frac{\mathfrak a_{k+1}^{(\pi)}}{(2k)!}
\]
立得前 \(2N\) 个矩也只依赖于
\[
\zeta(2),\dots,\zeta(4N).
\]

由 Favard 递推与 Gram--Schmidt 过程，前 \(N\) 个 Jacobi 参数由有限个矩唯一确定，因此也只依赖于有限多个显式偶值；对应的任意有限 principal Jacobi section 亦然。
同样，有限阶 Pad\'e 逼近与 terminating \(S\)-fraction 截断都由有限 Hankel 数据决定，所以仍只依赖于有限矩，因而只依赖于有限多个偶值。

最后，由定理 \ref{thm:xi-time-part9zbe-foldpi-fredholm-double-sine} 与定理 \ref{thm:xi-time-part9zbe-foldpi-jacobi-schatten-zeta}，
\[
-\log D_\varphi(z)
=
\sum_{r\ge 1}\frac{\zeta_{J_\varphi}(r)}{r}z^r
=
\sum_{r\ge 1}\frac{(4\pi^2)^{-r}(1+\varphi^{2r})\zeta(2r)}{r}\,z^r.
\]
因此任何有限 Taylor jet 也只依赖于有限多个显式偶值。

而 RH 是关于 \(\zeta(s)\) 全部非平凡零点几何的全局命题，不可能由上述只含有限显式偶值的数据得到等价刻画。
故任意仅依赖这些 finite local data 的判据都不可能与 RH 等价。证毕。有关 \(\zeta(2n)\) 的显式闭式可参见 \href{https://dlmf.nist.gov/25.6}{DLMF 25.6}。
\end{proof}

\begin{theorem}[de-arithmetized centered completion 的 Stieltjes 对数导数判据]\label{thm:xi-time-part9zbh-foldpi-stieltjes-log-derivative-criterion}
RH 成立，当且仅当存在正纯点测度
\[
\nu_{\mathrm{RH}}
=
\sum_{j\ge 1}m_j\,\delta_{\tau_j}
\qquad
(m_j\in\NN,\ \tau_j>0)
\]
使得
\[
\frac1z\frac{\dd}{\dd z}\log \Theta_\varphi^\sharp(z)
=
2\int_0^\infty \frac{\dd\nu_{\mathrm{RH}}(t)}{z^2+t}.
\]
在 RH 成立时，可以取
\[
\nu_{\mathrm{RH}}
=
\sum_{\gamma>0}m_\gamma\,\delta_{\gamma^2/4},
\]
其中 \(m_\gamma\) 为零点 \(\zeta(\frac12+i\gamma)=0\) 的重数。
\end{theorem}
\begin{proof}
先设 RH 成立。
由定理 \ref{thm:xi-time-part9zbh-foldpi-rh-quarter-line-equivalence}，
\[
\Theta_\varphi^\sharp(z)=\xi\!\Bigl(\frac12+2z\Bigr)
\]
的全部零点恰为
\[
z=\pm i\frac{\gamma}{2},
\qquad
\gamma>0,
\]
并按相同重数 \(m_\gamma\) 计数。
由于 \(\Theta_\varphi^\sharp\) 是偶整函数，故其 canonical product 可写为
\[
\Theta_\varphi^\sharp(z)
=
\Theta_\varphi^\sharp(0)
\prod_{\gamma>0}\left(1+\frac{4z^2}{\gamma^2}\right)^{m_\gamma}.
\]
对数求导后除以 \(z\)，得到
\[
\frac1z\frac{\dd}{\dd z}\log \Theta_\varphi^\sharp(z)
=
\sum_{\gamma>0}\frac{8m_\gamma}{\gamma^2+4z^2}
=
2\sum_{\gamma>0}\frac{m_\gamma}{z^2+\gamma^2/4}.
\]
这正是所述 Stieltjes 形，其中
\[
\nu_{\mathrm{RH}}
=
\sum_{\gamma>0}m_\gamma\,\delta_{\gamma^2/4}.
\]

反过来，若存在正纯点测度 \(\nu_{\mathrm{RH}}\) 使上式成立，则右端全部极点只能位于
\[
z=\pm i\sqrt t,
\qquad
t\in \operatorname{supp}\nu_{\mathrm{RH}}\subset (0,\infty),
\]
且极点重数由原子质量 \(m_j\in\NN\) 给出。
而对数导数的极点恰对应于 \(\Theta_\varphi^\sharp\) 的零点。
故 \(\Theta_\varphi^\sharp\) 的全部零点都落在虚轴上。
再由定理 \ref{thm:xi-time-part9zbh-foldpi-rh-quarter-line-equivalence}，这与 RH 等价。
证毕。
\end{proof}

\begin{corollary}[RH 下 centered completion 的完全单调 log-slope]\label{cor:xi-time-part9zbh-foldpi-completely-monotone-log-slope}
若 RH 成立，定义
\[
\Lambda_\varphi(u)
:=
\frac{1}{2\sqrt u}
\frac{\dd}{\dd x}\log \Theta_\varphi^\sharp(x)\Big|_{x=\sqrt u}
\qquad (u>0).
\]
则
\[
\Lambda_\varphi(u)
=
2\sum_{\gamma>0}\frac{m_\gamma}{u+\gamma^2/4},
\]
因而 \(\Lambda_\varphi\) 是 Stieltjes 函数，并满足
\[
(-1)^n\Lambda_\varphi^{(n)}(u)>0
\qquad (u>0,\ n\ge 0).
\]
特别地，
\begin{enumerate}
  \item \(\Theta_\varphi^\sharp(x)\) 在 \(x>0\) 上严格递增；
  \item 函数
  \[
  x\longmapsto \frac1x\frac{\dd}{\dd x}\log \Theta_\varphi^\sharp(x)
  \]
  在 \(x>0\) 上严格递减。
\end{enumerate}
\end{corollary}
\begin{proof}
在 RH 前提下，定理 \ref{thm:xi-time-part9zbh-foldpi-stieltjes-log-derivative-criterion} 给出
\[
\frac1x\frac{\dd}{\dd x}\log \Theta_\varphi^\sharp(x)
=
2\sum_{\gamma>0}\frac{m_\gamma}{x^2+\gamma^2/4}.
\]
令 \(u=x^2\)，即得
\[
\Lambda_\varphi(u)
=
2\sum_{\gamma>0}\frac{m_\gamma}{u+\gamma^2/4}.
\]
逐次求导可得
\[
\Lambda_\varphi^{(n)}(u)
=
2(-1)^n n!\sum_{\gamma>0}\frac{m_\gamma}{(u+\gamma^2/4)^{n+1}},
\]
从而
\[
(-1)^n\Lambda_\varphi^{(n)}(u)>0.
\]
又因为 \(\Theta_\varphi^\sharp(0)=\xi(\frac12)>0\)，且
\[
\frac{\dd}{\dd x}\log \Theta_\varphi^\sharp(x)
=
2x\,\Lambda_\varphi(x^2)>0
\qquad (x>0),
\]
故 \(\Theta_\varphi^\sharp(x)\) 在 \(x>0\) 上严格递增。
再由 \(\Lambda_\varphi\) 严格递减，立得
\[
x\longmapsto \frac1x\frac{\dd}{\dd x}\log \Theta_\varphi^\sharp(x)=2\Lambda_\varphi(x^2)
\]
严格递减。证毕。
\end{proof}

\begin{theorem}[RH 的 positive trace-class determinant 判据]\label{thm:xi-time-part9zbh-foldpi-positive-trace-class-rh}
RH 成立，当且仅当存在正 trace-class 算子 \(K_{\mathrm{RH}}\) 使得
\[
\Theta_\varphi^\sharp(z)
=
\Theta_\varphi^\sharp(0)\,
\det(I+4z^2K_{\mathrm{RH}}).
\]
在 RH 成立时，\(K_{\mathrm{RH}}\) 的非零谱恰为
\[
\sigma(K_{\mathrm{RH}})\setminus\{0\}
=
\left\{
\gamma^{-2}:\ \zeta\!\left(\frac12+i\gamma\right)=0,\ \gamma>0
\right\},
\]
并按零点重数计数。
该算子在酉等价意义下唯一。
\end{theorem}
\begin{proof}
先设 RH 成立。
由定理 \ref{thm:xi-time-part9zbh-foldpi-rh-quarter-line-equivalence}，
\[
\Theta_\varphi^\sharp(z)=\xi\!\Bigl(\frac12+2z\Bigr)
\]
的零点恰为 \(\pm i\gamma/2\)。
因此
\[
\Theta_\varphi^\sharp(z)
=
\Theta_\varphi^\sharp(0)
\prod_{\gamma>0}\left(1+\frac{4z^2}{\gamma^2}\right)^{m_\gamma}.
\]
再由 Riemann 零点计数的标准估计 \(N(T)=O(T\log T)\)，可知
\[
\sum_{\gamma>0}\frac{m_\gamma}{\gamma^2}<\infty.
\]
于是按重数排列非零特征值
\[
\gamma_1^{-2},\gamma_2^{-2},\dots
\]
可定义正对角算子
\[
K_{\mathrm{RH}}
:=
\operatorname{diag}(\gamma_1^{-2},\gamma_2^{-2},\dots),
\]
其为 trace class，并满足
\[
\det(I+4z^2K_{\mathrm{RH}})
=
\prod_{\gamma>0}\left(1+\frac{4z^2}{\gamma^2}\right)^{m_\gamma}.
\]
故 determinant 表示成立。

反过来，若存在正 trace-class 算子 \(K_{\mathrm{RH}}\) 使所示公式成立，则 Fredholm 行列式的零点只能为
\[
z=\pm \frac{i}{2\sqrt\lambda},
\qquad
\lambda\in \sigma(K_{\mathrm{RH}})\setminus\{0\},
\]
因而全部位于虚轴上。
于是 \(\Theta_\varphi^\sharp(z)\) 的零点全在虚轴上，再由定理 \ref{thm:xi-time-part9zbh-foldpi-rh-quarter-line-equivalence}，得到 RH。

唯一性方面，determinant 表示唯一确定 \(K_{\mathrm{RH}}\) 的非零特征值多重集，而正紧自伴算子由该谱多重集在酉等价意义下唯一确定。
证毕。有关 \(\xi\)-函数的零点分布与 Hadamard 型表示，可参见 \href{https://dlmf.nist.gov/25.10}{DLMF 25.10}。
\end{proof}

\begin{theorem}[project centered completion 的 arithmetic--spectral exact decoupling]\label{thm:xi-time-part9zbh-foldpi-arithmetic-spectral-decoupling}
在 RH 成立的前提下，\(\Theta_\varphi\) 具有唯一分解
\[
\Theta_\varphi(z)
=
A_\varphi(z)\,
\Theta_\varphi^\sharp(0)\,
\det(I+4z^2K_{\mathrm{RH}}),
\]
其中
\[
A_\varphi(z)
:=
4\varphi^{1/2}
\cosh\!\Bigl(L\bigl(z+\tfrac14\bigr)\Bigr)
\cosh\!\Bigl(L\bigl(z-\tfrac14\bigr)\Bigr),
\qquad
L=\log\varphi,
\]
而 \(K_{\mathrm{RH}}\) 为定理 \ref{thm:xi-time-part9zbh-foldpi-positive-trace-class-rh} 中的正 trace-class 算子。
因此：
\begin{enumerate}
  \item 全部算术零点由 \(A_\varphi\) 精确承担；
  \item 全部 RH-sensitive 零点由 \(\det(I+4z^2K_{\mathrm{RH}})\) 精确承担；
  \item 两者之间不是近似解耦，而是严格乘法解耦。
\end{enumerate}
\end{theorem}
\begin{proof}
由定理 \ref{thm:xi-time-part9zbh-foldpi-centered-completion-factorization}，
\[
\Theta_\varphi(z)=A_\varphi(z)\Theta_\varphi^\sharp(z).
\]
在 RH 前提下，再由定理 \ref{thm:xi-time-part9zbh-foldpi-positive-trace-class-rh}，
\[
\Theta_\varphi^\sharp(z)
=
\Theta_\varphi^\sharp(0)\det(I+4z^2K_{\mathrm{RH}}).
\]
两式合并即得所示分解。
唯一性来自 \(A_\varphi\) 的显式固定性以及 \(K_{\mathrm{RH}}\) 的酉等价唯一性。
证毕。
\end{proof}

\begin{conjecture}[从 RH-blind compact 宿主到 positive determinant cone 的非局域桥梁]\label{conj:xi-time-part9zbh-foldpi-nonlocal-bridge-rh}
若项目框架内部确能生成一条通向 RH 的严格证明链，则该证明链必定经过定理 \ref{thm:xi-time-part9zbh-foldpi-quarter-line-completion} 与定理 \ref{thm:xi-time-part9zbh-foldpi-positive-trace-class-rh} 所揭示的 global completion 层，而不可能来自任何 finite local mechanism。
更具体地说，这条桥梁不可能仅由下列对象给出：
\begin{enumerate}
  \item H.111 odd tail 的任意有限截断；
  \item canonical compact Jacobi 宿主 \(J_\varphi\) 的任意有限压缩或 finite-rank section；
  \item 对应的任意 finite \(S\)-fraction、Pad\'e 截断或 continued-fraction tail；
  \item 任意有限阶 Fredholm--Taylor jet
  \[
  [z^k]\bigl(-\log D_\varphi(z)\bigr).
  \]
\end{enumerate}
更进一步，任何真正把 RH-blind 的 compact 宿主 \(J_\varphi\) 提升为 RH-sensitive 算子 \(K_{\mathrm{RH}}\) 的函子，都不可能是 \(J_\varphi\) 上的 polynomial、rational 或 finite-rank functional calculus；它必须是 genuinely nonlocal 的，并在谱层面把局部 continued-fraction 数据提升为一个全局 determinant 对象。
\end{conjecture}

\noindent
因此，H.111 odd tail 在该项目链中的 RH 接口具有严格的双重性：在 local/finite 层面，它只编码显式偶值 \(\zeta(2n)\)，因而对 RH 完全失明；在 global/completed 层面，它却又把 RH 精确转写为 centered completion 的正 determinant 判据。由此，项目与 RH 的可能交汇点不在 tail coefficient algebra、finite continued fraction 或 finite Jacobi section，而只能位于一个 genuinely nonlocal 的 global completion functor 及其对应的 positive trace-class determinant cone。

\endinput

\paragraph{fold-\texorpdfstring{$\pi$}{pi} 两原子 Mellin 尺度律的 completion 统一、平方 Bernstein 刚性与 finite local algebra 的零点几何盲性}
沿用定理~\ref{thm:xi-time-part9zbd-scaled-multiplicity-min-oddmoment}、
\ref{thm:xi-time-part9zbe-foldpi-compact-determinant-sampling}、
\ref{thm:xi-time-part9zbh-foldpi-quarter-line-completion}、
\ref{thm:xi-time-part9zbh-foldpi-centered-completion-factorization}
与定理~\ref{thm:xi-time-part9zbh-foldpi-stieltjes-log-derivative-criterion} 的记号。
前文已经分别得到：临界容量曲线来自两原子 scaled multiplicity law，H.111 的 odd Bernoulli tail 是同一尺度律的奇 Mellin 采样，而 quarter-line completion 则把 centered \(\xi\)-函数嵌入一条显式自对偶完成化链。
下面可把这三条结构进一步压缩为同一 Mellin 宿主上的三种投影：\(\min\)-变换、奇 Mellin 取样与自对偶 Mellin 自相关。
由此，capacity、odd tail 与 quarter-line completion 之间不再只是并行对应，而形成一条精确的统一链；同时，RH 敏感性与 finite local algebra 之间的分界也可写成一条平方变量上的 Bernstein 刚性判据。

\begin{theorem}[两原子尺度律的 min--Mellin--completion 三重统一]\label{thm:xi-time-part9zbha-twoatomic-mellin-triple-unification}
定义两原子尺度律
\[
\nu_\varphi^{\mathrm{sc}}
:=
\varphi^{-1}\delta_1+\varphi^{-2}\delta_{\varphi^{-1}},
\qquad
M_\varphi(z)
:=
\int x^{-z}\,\dd\nu_\varphi^{\mathrm{sc}}(x)
=
\varphi^{-1}+\varphi^{z-2}.
\]
则它同时满足下列三条精确表示：
\[
\mathsf S_\varphi(s)
=
\frac{\varphi^2}{\sqrt5}
\int \min\{x,s\}\,\dd\nu_\varphi^{\mathrm{sc}}(x),
\]
\[
\mathfrak a_r^{(\pi)}
=
\frac{B_{2r}}{r(2r-1)}\,M_\varphi(2r-1)
\qquad (r\ge 1),
\]
\[
\Xi_\varphi(s)
=
\varphi^3\,M_\varphi(2s)\,M_\varphi(1-2s)\,\xi(2s).
\]
因此，临界容量相图、odd Bernoulli tail 与 quarter-line completion 并非三条彼此独立的机制，而是同一两原子尺度律在 \(\min\)-变换、奇 Mellin 取样与自对偶 Mellin 自相关下的三种投影。
\end{theorem}
\begin{proof}
第一式正是定理~\ref{thm:xi-time-part9zbd-scaled-multiplicity-min-oddmoment} 的测度表示。
又由
\[
M_\varphi(2r-1)
=
\varphi^{-1}+\varphi^{2r-3},
\]
结合
\[
\mathfrak a_r^{(\pi)}
=
\frac{B_{2r}}{r(2r-1)}
\bigl(\varphi^{-1}+\varphi^{2r-3}\bigr),
\]
便得第二式。

对于第三式，由
\[
M_\varphi(2s)=\varphi^{-1}+\varphi^{2s-2},
\qquad
M_\varphi(1-2s)=\varphi^{-1}+\varphi^{-1-2s},
\]
可得
\[
\varphi^3M_\varphi(2s)M_\varphi(1-2s)
=
\varphi^3\bigl(\varphi^{-1}+\varphi^{2s-2}\bigr)\bigl(\varphi^{-1}+\varphi^{-1-2s}\bigr)
=
\bigl(1+\varphi^{2s}\bigr)\bigl(1+\varphi^{1-2s}\bigr).
\]
代回定理~\ref{thm:xi-time-part9zbh-foldpi-quarter-line-completion} 中
\[
\Xi_\varphi(s)
=
\bigl(1+\varphi^{2s}\bigr)\bigl(1+\varphi^{1-2s}\bigr)\xi(2s),
\]
便得所示分解。证毕。
\end{proof}

\begin{theorem}[completion 因子的二值分支与容量曲线的唯一选枝]\label{thm:xi-time-part9zbha-completion-branch-capacity-selection}
\leanverified{paper\_xi\_time\_part9zbha\_completion\_branch\_capacity\_selection}
对任意 \(p\in(0,1)\)，定义归一化两原子概率测度
\[
\nu_{\varphi,p}
:=
p\,\delta_1+(1-p)\,\delta_{\varphi^{-1}},
\qquad
M_{\varphi,p}(z)
:=
\int x^{-z}\,\dd\nu_{\varphi,p}(x)
=
p+(1-p)\varphi^z.
\]
则存在常数 \(\kappa>0\) 使得
\[
\kappa\,M_{\varphi,p}(2s)M_{\varphi,p}(1-2s)
=
\bigl(1+\varphi^{2s}\bigr)\bigl(1+\varphi^{1-2s}\bigr)
\qquad (\forall s\in\CC)
\]
当且仅当
\[
p=\frac12
\qquad\text{或}\qquad
p=\varphi^{-1}.
\]
相应地，
\[
(p,\kappa)=\left(\frac12,4\right)
\qquad\text{或}\qquad
(p,\kappa)=\left(\varphi^{-1},\varphi^3\right).
\]

进一步，若要求其 \(\min\)-变换给出的双折点容量曲线在两条线性段上的斜率比恰与定理~\ref{thm:xi-time-part9zbd-scaled-multiplicity-double-kink} 中的真实容量相图一致，则唯一允许的分支是
\[
p=\varphi^{-1},
\qquad
\nu_{\varphi,p}=\nu_\varphi^{\mathrm{sc}}.
\]
因此，quarter-line completion factor 并不单独唯一决定 projective 两原子律；真正消解二值歧义的，是容量曲线的折点斜率比。
\end{theorem}
\begin{proof}
展开得
\[
\kappa\,M_{\varphi,p}(2s)M_{\varphi,p}(1-2s)
=
\kappa\Bigl(
p^2+p(1-p)\varphi^{2s}+p(1-p)\varphi^{1-2s}+\varphi(1-p)^2
\Bigr).
\]
与
\[
\bigl(1+\varphi^{2s}\bigr)\bigl(1+\varphi^{1-2s}\bigr)
=
\varphi^2+\varphi^{2s}+\varphi^{1-2s}
\]
逐项比较，得到
\[
\kappa p(1-p)=1,
\qquad
\kappa\bigl(p^2+\varphi(1-p)^2\bigr)=\varphi^2.
\]
消去 \(\kappa\) 后得到
\[
p^2+\varphi(1-p)^2=\varphi^2p(1-p),
\]
等价于
\[
2\varphi^2\left(p-\frac12\right)\left(p-\varphi^{-1}\right)=0.
\]
故
\[
p=\frac12
\qquad\text{或}\qquad
p=\varphi^{-1}.
\]
再由 \(\kappa p(1-p)=1\) 即得相应 \(\kappa\) 的取值。

现记
\[
\mathsf S_{\varphi,p}(s)
:=
c\int \min\{x,s\}\,\dd\nu_{\varphi,p}(x),
\]
其中 \(c>0\) 为任意归一化常数。
则对 \(0<s<\varphi^{-1}\) 有
\[
\mathsf S_{\varphi,p}'(s)=c,
\]
对 \(\varphi^{-1}<s<1\) 有
\[
\mathsf S_{\varphi,p}'(s)=cp.
\]
因此折点 \(s=\varphi^{-1}\) 处的左右斜率比为
\[
\frac{\mathsf S_{\varphi,p}'(\varphi^{-1}+0)}{\mathsf S_{\varphi,p}'(\varphi^{-1}-0)}=p.
\]
另一方面，由定理~\ref{thm:xi-time-part9zbd-scaled-multiplicity-double-kink} 的闭式，
\[
\frac{\varphi/\sqrt5}{\varphi^2/\sqrt5}=\varphi^{-1}
\]
正是目标容量相图的真实斜率比。
故只能取 \(p=\varphi^{-1}\)。证毕。
\end{proof}

\begin{theorem}[de-arithmetized centered completion 的平方 complete Bernstein 判据]\label{thm:xi-time-part9zbha-centered-completion-complete-bernstein-rh}
定义
\[
H_\varphi(u)
:=
\log\frac{\Theta_\varphi^\sharp(\sqrt u)}{\Theta_\varphi^\sharp(0)}
\qquad (u>0),
\]
其中 \(\sqrt u\) 取正支。
则下列命题严格等价：
\begin{enumerate}
  \item Riemann 假设成立；
  \item \(H_\varphi\) 是 complete Bernstein 函数；
  \item \(H_\varphi'(u)\) 是 Stieltjes 函数。
\end{enumerate}
因此，quarter-line centered completion 的 RH 接口可被重新压缩为一条平方变量上的 Bernstein 正性判据。
\end{theorem}
\begin{proof}
先证 \((1)\Rightarrow(3)\)。
由定理~\ref{thm:xi-time-part9zbh-foldpi-stieltjes-log-derivative-criterion}，RH 成立当且仅当
\[
\frac1z\frac{\dd}{\dd z}\log \Theta_\varphi^\sharp(z)
=
2\sum_{\gamma>0}\frac{m_\gamma}{z^2+\gamma^2/4},
\]
其中 \(m_\gamma\) 为零点 \(\zeta(\frac12+i\gamma)=0\) 的重数。
令 \(u=z^2\)，则
\[
H_\varphi'(u)
=
\frac{1}{2\sqrt u}\frac{\dd}{\dd z}\log \Theta_\varphi^\sharp(z)\Big|_{z=\sqrt u}
=
\sum_{\gamma>0}\frac{m_\gamma}{u+\gamma^2/4},
\]
这正是一条 Stieltjes 表示，故 \(H_\varphi'\) 为 Stieltjes 函数。

再证 \((3)\Rightarrow(1)\)。
若 \(H_\varphi'\) 是 Stieltjes 函数，则存在正测度 \(\mu\) 使得
\[
H_\varphi'(u)=\int_0^\infty \frac{\dd\mu(t)}{u+t}.
\]
于是对 \(z\neq 0\)，
\[
\frac1z\frac{\dd}{\dd z}\log \Theta_\varphi^\sharp(z)
=
2H_\varphi'(z^2)
=
2\int_0^\infty \frac{\dd\mu(t)}{z^2+t},
\]
从而定理~\ref{thm:xi-time-part9zbh-foldpi-stieltjes-log-derivative-criterion} 立即推出 RH。

最后，\((2)\Leftrightarrow(3)\) 是 complete Bernstein 函数的标准微分刻画：函数 \(f\) 在 \((0,\infty)\) 上为 complete Bernstein，当且仅当 \(f'\) 为 Stieltjes 函数。
将 \(f=H_\varphi\) 代入即可。证毕。
\end{proof}

\begin{corollary}[RH 下 centered completion 的移位 jet 为 Hausdorff 矩序列]\label{cor:xi-time-part9zbha-centered-jet-hausdorff}
在 RH 成立的前提下，对任意固定 \(u_0>0\)，定义
\[
m_n(u_0)
:=
\frac{(-1)^n}{n!}\,H_\varphi^{(n+1)}(u_0)
\qquad (n\ge 0).
\]
则 \(\bigl(m_n(u_0)\bigr)_{n\ge 0}\) 是支撑于区间 \([0,u_0^{-1}]\) 上的 Hausdorff 矩序列。
更精确地，
\[
m_n(u_0)=\int_0^{u_0^{-1}}x^n\,\dd\mu_{u_0}(x),
\]
其中
\[
\mu_{u_0}
:=
\sum_{\gamma>0}
m_\gamma\,(u_0+\gamma^2/4)^{-1}\,
\delta_{(u_0+\gamma^2/4)^{-1}}.
\]
特别地，对每个 \(k\ge 0\)，Hankel 矩阵
\[
\bigl(m_{i+j}(u_0)\bigr)_{0\le i,j\le k}
\]
与 localizing 矩阵
\[
\bigl(u_0^{-1}m_{i+j}(u_0)-m_{i+j+1}(u_0)\bigr)_{0\le i,j\le k}
\]
都正半定。
\end{corollary}
\begin{proof}
在 RH 前提下，定理~\ref{thm:xi-time-part9zbha-centered-completion-complete-bernstein-rh} 的证明已经给出
\[
H_\varphi'(u)
=
\sum_{\gamma>0}\frac{m_\gamma}{u+\gamma^2/4}.
\]
逐次求导得
\[
m_n(u_0)
=
\sum_{\gamma>0}\frac{m_\gamma}{(u_0+\gamma^2/4)^{n+1}}.
\]
记
\[
x_\gamma:=(u_0+\gamma^2/4)^{-1}\in (0,u_0^{-1}],
\]
并定义
\[
\mu_{u_0}
:=
\sum_{\gamma>0}m_\gamma x_\gamma\,\delta_{x_\gamma}.
\]
则
\[
\int_0^{u_0^{-1}} x^n\,\dd\mu_{u_0}(x)
=
\sum_{\gamma>0}m_\gamma x_\gamma^{n+1}
=
m_n(u_0).
\]
故 \(\bigl(m_n(u_0)\bigr)\) 为 Hausdorff 矩序列。
Hankel 与 localizing 的正半定性则是 Hausdorff 矩问题的标准推论。证毕。
\end{proof}

\begin{theorem}[finite local scalar algebra 的显式偶值因子化与零点几何盲性]\label{thm:xi-time-part9zbha-finite-local-algebra-even-zeta-factorization}
\leanverified{paper\_xi\_time\_part9zbha\_finite\_local\_algebra\_even\_zeta\_factorization}
对每个 \(M\ge 1\)，定义域
\[
\mathcal K_M
:=
\mathbb Q\bigl(\varphi,\pi,\zeta(2),\zeta(4),\dots,\zeta(2M)\bigr).
\]
再令 \(\mathcal A_M\) 表示由下列 finite local scalar data 经过有理运算所生成的代数：
\begin{enumerate}
  \item \(-\log D_\varphi(z)\) 在 \(z=0\) 的前 \(M\) 个 Taylor 系数；
  \item H.111 odd tail 的前 \(M\) 个系数 \(\mathfrak a_1^{(\pi)},\dots,\mathfrak a_M^{(\pi)}\)；
  \item 前 \(M\) 个矩 \(u_0,\dots,u_{M-1}\)；
  \item 以及任何由这 \(M\) 个矩唯一决定的有限 principal Jacobi section 的标量不变量。
\end{enumerate}
则有
\[
\mathcal A_M\subseteq \mathcal K_M.
\]
因此，fold-\(\pi\) compact 宿主的全部 finite local scalar algebra 都完全因子化到有限多个显式偶值 \(\zeta(2),\dots,\zeta(2M)\) 上；非平凡零点的全局几何在任何固定深度 \(M\) 上都不可见。
\end{theorem}
\begin{proof}
由定理~\ref{thm:xi-time-part9zbe-foldpi-fredholm-double-sine} 与定理~\ref{thm:xi-time-part9zbe-foldpi-jacobi-schatten-zeta}，
\[
-\log D_\varphi(z)
=
\sum_{r\ge 1}\frac{(4\pi^2)^{-r}(1+\varphi^{2r})\zeta(2r)}{r}\,z^r.
\]
故其前 \(M\) 个 Taylor 系数都属于 \(\mathcal K_M\)。

再由定理~\ref{thm:xi-time-part9zbe-foldpi-compact-determinant-sampling}，
\[
\mathfrak a_r^{(\pi)}
=
\frac{B_{2r}}{r(2r-1)}\bigl(\varphi^{-1}+\varphi^{2r-3}\bigr),
\]
结合 Euler 的偶值公式
\[
B_{2r}
=
(-1)^{r-1}\frac{2(2r)!}{(2\pi)^{2r}}\zeta(2r),
\]
可知
\[
\mathfrak a_1^{(\pi)},\dots,\mathfrak a_M^{(\pi)}\in \mathcal K_M.
\]

又由
\[
u_k=(-1)^k\frac{\mathfrak a_{k+1}^{(\pi)}}{(2k)!},
\]
得到
\[
u_0,\dots,u_{M-1}\in \mathcal K_M.
\]

最后，按照定义，有限 principal Jacobi section 的这类标量不变量仅由上述矩数据经有理运算产生，因而仍属于 \(\mathcal K_M\)。
故 \(\mathcal A_M\subseteq \mathcal K_M\)。

另一方面，\(\mathcal K_M\) 只由有限多个显式偶值生成，而 RH 敏感对象由定理~\ref{thm:xi-time-part9zbh-foldpi-quarter-line-completion}、
\ref{thm:xi-time-part9zbh-foldpi-rh-quarter-line-equivalence}
与定理~\ref{thm:xi-time-part9zbha-centered-completion-complete-bernstein-rh} 所揭示的 centered completion 全局零点几何承担。
因此 finite local scalar algebra 不承载该零点几何。证毕。
\end{proof}

\begin{corollary}[finite-depth \texorpdfstring{$\pi$}{pi}-tail formalism 中不存在 RH 敏感的局部证书]\label{cor:xi-time-part9zbha-no-local-rh-certificate}
\leanverified{paper\_xi\_time\_part9zbha\_no\_local\_rh\_certificate}
设 \(F\) 是任意一个只允许依赖下列有限 local quantities 的标量表达式：
\begin{enumerate}
  \item 有限个 odd-layer 系数；
  \item 有限个 moments；
  \item \(-\log D_\varphi\) 的有限阶 Taylor jet；
  \item 有限 principal Jacobi section 的标量不变量；
\end{enumerate}
并且 \(F\) 仅经有理运算组合而成。
则存在 \(M\ge 1\) 使得
\[
F\in \mathcal K_M
=
\mathbb Q\bigl(\varphi,\pi,\zeta(2),\dots,\zeta(2M)\bigr).
\]
因此，任何真正对 nontrivial-zero geometry 敏感的判据，都不可能在 finite-depth \(\pi\)-tail formalism 内闭合；它必须越出 finite local algebra，进入 centered completion 的全局层。
\end{corollary}
\begin{proof}
这只是定理~\ref{thm:xi-time-part9zbha-finite-local-algebra-even-zeta-factorization} 的直接重述。证毕。
\end{proof}

\begin{conjecture}[最小全局完成化原理]\label{conj:xi-time-part9zbha-minimal-global-completion-principle}
设 \(\nu\) 为 \((0,1]\) 上紧支撑正概率测度，并记其 Mellin 变换为
\[
M_\nu(s):=\int x^{-s}\,\dd\nu(x).
\]
设又存在常数 \(\kappa_\nu>0\)，使得
\[
\Xi_\nu(s):=\kappa_\nu\,M_\nu(2s)M_\nu(1-2s)\,\xi(2s)
\]
给出一条自对偶 completion。
猜想：若 \(\nu\) 同时满足下列三项要求，
\begin{enumerate}
  \item 其 \(\min\)-变换精确复现临界双折点容量曲线，即
  \[
  \frac{\varphi^2}{\sqrt5}\int \min\{x,s\}\,\dd\nu(x)=\mathsf S_\varphi(s);
  \]
  \item 其奇 Mellin 取样精确复现 H.111 odd Bernoulli 权重，即
  \[
  \mathfrak a_r^{(\pi)}
  =
  \frac{B_{2r}}{r(2r-1)}\,M_\nu(2r-1)
  \qquad (r\ge 1);
  \]
  \item \(\Xi_\nu\) 的去算术化 centered completion 在平方变量化后满足定理~\ref{thm:xi-time-part9zbha-centered-completion-complete-bernstein-rh} 的 complete Bernstein 判据，
\end{enumerate}
则必有
\[
\nu=\nu_\varphi^{\mathrm{sc}}
=
\varphi^{-1}\delta_1+\varphi^{-2}\delta_{\varphi^{-1}},
\qquad
\Xi_\nu(s)=\Xi_\varphi(s).
\]
亦即，项目中的两原子尺度律并不是方便选取的一组参数，而是同时兼容容量相图、odd tail 与 RH-completion 的唯一最小正测度宿主。
\end{conjecture}

\noindent
这条猜想已有三重证据支撑。
定理~\ref{thm:xi-time-part9zbha-twoatomic-mellin-triple-unification} 已证明 \(\nu_\varphi^{\mathrm{sc}}\) 的确同时生成容量、odd tail 与 quarter-line completion 三条结构。
定理~\ref{thm:xi-time-part9zbha-completion-branch-capacity-selection} 进一步表明，即便局限在最小的两原子归一化族内，completion 因子仍存在一条真正的二值歧义，而容量曲线的折点斜率比恰好唯一消去这条歧义。
最后，定理~\ref{thm:xi-time-part9zbha-finite-local-algebra-even-zeta-factorization} 与推论~\ref{cor:xi-time-part9zbha-no-local-rh-certificate} 表明，任何候选替代若想承载 RH 敏感信息，便必然是 genuinely global 的 completion 对象，而不可能闭合在 finite local algebra 内。
因此，fold-\(\pi\) 机制中的 RH-sensitive 全球层并不是由 compact host 单独决定，而是由 compact host 与两原子尺度律之间的自对偶 Mellin 拼接共同决定。

\endinput

\paragraph{fold-\texorpdfstring{$\pi$}{pi} exact Jacobi 宿主的主子压缩、最小 exact compressor 与 Ritz--Wythoff 极点逼近}
与上一节 finite-support truncation 塔不同，这里不再截断谱支撑，而是直接对 exact compact Jacobi 宿主 \(J_\varphi\) 取 principal compression。
这样得到的 finite-dimensional scalar host 不再只是从下方单调逼近 odd Bernoulli tail，而是在一个精确长度阈值之前完全无损。
更进一步，这条 principal compression 层级具有最优维数、首个失配项的唯一回路公式，以及沿 Wythoff 双梯的外侧单调 pole convergence。

沿用定理~\ref{thm:xi-time-part9zbe-foldpi-canonical-compact-jacobi}、
\ref{thm:xi-time-part9zbe-foldpi-compact-determinant-sampling}、
\ref{thm:xi-time-part9zbf-foldpi-wythoff-spectral-order}
与定理~\ref{thm:xi-time-part9zbg-foldpi-finite-depth-exactness-obstruction} 的记号。
把 \(J_\varphi\) 写成其标准 Jacobi 形式
\[
J_\varphi e_j=\alpha_{j+1}e_{j+1}+\beta_j e_j+\alpha_j e_{j-1},
\qquad
\alpha_j>0,\ \alpha_0:=0,\ \beta_j\in\RR,
\]
其中 \((e_j)_{j\ge 0}\) 为 \(\ell^2(\NN_0)\) 的标准正交基。
由定理~\ref{thm:xi-time-part9zbe-foldpi-canonical-compact-jacobi} 的构造，循环向量满足
\[
v_0=U(1)=\sqrt{u_0}\,e_0,
\qquad
u_k=\langle v_0,J_\varphi^k v_0\rangle
\qquad (k\ge 0).
\]
因此 H.111 的 odd Bernoulli tail 仍可写为
\[
\mathcal E_\varphi(s)
=
\sum_{k\ge 0}(-1)^k(2k)!\,u_k\,s^{2k+1}.
\]

\begin{theorem}[principal compression 的 \texorpdfstring{$(2n-1)$}{(2n-1)} 阶精确矩保持]\label{thm:xi-time-part9zbi-principal-compression-exact-moments}
对每个 \(n\ge 1\)，定义
\[
E_n:=\operatorname{span}\{e_0,\dots,e_{n-1}\},
\qquad
P_n:\ell^2(\NN_0)\to E_n
\]
为正交投影，并令
\[
J_\varphi^{[n]}:=P_nJ_\varphi P_n|_{E_n}.
\]
则对每个 \(n\ge 1\) 与每个整数 \(0\le k\le 2n-1\)，都有严格恒等式
\[
\langle v_0,(J_\varphi^{[n]})^k v_0\rangle
=
\langle v_0,J_\varphi^k v_0\rangle
=
u_k.
\]
换言之，\(n\times n\) 的 principal compression 已能精确保留 exact host 的前 \(2n\) 个 odd Bernoulli layers。
\end{theorem}
\begin{proof}
因为 \(v_0=\sqrt{u_0}\,e_0\)，只需证明
\[
\langle e_0,(J_\varphi^{[n]})^k e_0\rangle
=
\langle e_0,J_\varphi^k e_0\rangle
\qquad (0\le k\le 2n-1).
\]

由 Jacobi 结构，每作用一次 \(J_\varphi\) 最多只会把基指标改变 \(1\)，亦可保持不变。
因此
\[
\langle e_0,J_\varphi^k e_0\rangle
\]
可解释为所有长度为 \(k\) 的带权闭路之和，其中每条路径都从层 \(0\) 出发并回到层 \(0\)。

若某条闭路曾到达层 \(n\)，则至少需要 \(n\) 次上行与 \(n\) 次下行，因此总长度至少为 \(2n\)。
于是当 \(k\le 2n-1\) 时，一切有贡献的闭路都完全落在
\[
\{0,1,\dots,n-1\}
\]
内。
而在这一区间上，\(J_\varphi^{[n]}\) 与 \(J_\varphi\) 的边权和留点权完全一致，因此这些闭路的权重与总和都不变。
故
\[
\langle e_0,(J_\varphi^{[n]})^k e_0\rangle
=
\langle e_0,J_\varphi^k e_0\rangle
\]
对 \(0\le k\le 2n-1\) 成立。
两边再同乘 \(u_0\) 即得
\[
\langle v_0,(J_\varphi^{[n]})^k v_0\rangle=u_k.
\]
证毕。
\end{proof}

\begin{corollary}[H.111 奇层的 principal compression 精确重建]\label{cor:xi-time-part9zbi-principal-compression-odd-tail-reconstruction}
定义 principal compression 的有理 Borel 影子
\[
\widehat{\mathcal E}_\varphi^{[n]}(\xi)
:=
\langle v_0,(I+\xi^2J_\varphi^{[n]})^{-1}v_0\rangle,
\]
以及相应的 factorial lift
\[
\mathcal E_\varphi^{[n]}(s)
:=
\sum_{k\ge 0}
(-1)^k(2k)!\,
\langle v_0,(J_\varphi^{[n]})^k v_0\rangle\,
s^{2k+1}.
\]
则当 \(\xi\to 0\) 与 \(s\to 0\) 时分别有
\[
\widehat{\mathcal E}_\varphi^{[n]}(\xi)-\widehat{\mathcal E}_\varphi(\xi)
=
O(\xi^{4n}),
\]
\[
\mathcal E_\varphi^{[n]}(s)-\mathcal E_\varphi(s)
=
O(s^{4n+1}).
\]
因此，\(\mathcal E_\varphi^{[n]}\) 与 exact odd tail 在
\[
s,\ s^3,\ \dots,\ s^{4n-1}
\]
诸项上完全一致；亦即前 \(2n\) 个 H.111 奇层系数被 \(n\)-维主子压缩无损重建。
\end{corollary}
\begin{proof}
由 Neumann 展开，
\[
\widehat{\mathcal E}_\varphi(\xi)
=
\sum_{k\ge 0}(-1)^k u_k\,\xi^{2k},
\]
而
\[
\widehat{\mathcal E}_\varphi^{[n]}(\xi)
=
\sum_{k\ge 0}
(-1)^k
\langle v_0,(J_\varphi^{[n]})^k v_0\rangle
\xi^{2k}.
\]
由定理~\ref{thm:xi-time-part9zbi-principal-compression-exact-moments}，两者在 \(k=0,\dots,2n-1\) 上系数完全一致，故差从 \(\xi^{4n}\) 起始。

同理，
\[
\mathcal E_\varphi(s)
=
\sum_{k\ge 0}(-1)^k(2k)!\,u_k\,s^{2k+1},
\]
\[
\mathcal E_\varphi^{[n]}(s)
=
\sum_{k\ge 0}
(-1)^k(2k)!\,
\langle v_0,(J_\varphi^{[n]})^k v_0\rangle
s^{2k+1}.
\]
由同一矩保持结论，二者在 \(k=0,\dots,2n-1\) 上仍逐项相同，因此差从
\[
s^{2(2n)+1}=s^{4n+1}
\]
开始。
证毕。
\end{proof}

\begin{theorem}[\texorpdfstring{$n$}{n}-维 exact compressor 的最优性]\label{thm:xi-time-part9zbi-minimal-exact-compressor}
\leanverified{paper\_xi\_time\_part9zbi\_minimal\_exact\_compressor}
设 \(K\) 为任意自伴算子，\(v\) 为其向量。
若对某个 \(n\ge 1\) 有
\[
\langle v,K^k v\rangle=u_k
\qquad (0\le k\le 2n-1),
\]
则必有
\[
\dim \operatorname{span}\{v,Kv,\dots,K^{n-1}v\}\ge n.
\]
特别地，不存在任何维数小于 \(n\) 的标量谱宿主，能够精确承载 H.111 的前 \(2n\) 个 odd Bernoulli layers；因此 principal compression \(J_\varphi^{[n]}\) 在维数意义下是最小 exact compressor。
\end{theorem}
\begin{proof}
由定理~\ref{thm:xi-time-part9zbc-foldpi-hankel-strict-positivity}，矩 Hankel 矩阵
\[
H_{n-1}:=(u_{i+j})_{0\le i,j\le n-1}
\]
严格正定。
事实上，任取 \(c=(c_0,\dots,c_{n-1})\neq 0\)，令
\[
P_c(y):=\sum_{j=0}^{n-1}c_j y^j.
\]
由于 \(\mu_\varphi\) 的支撑为无穷互异点集，非零多项式 \(P_c\) 不可能在其支撑上恒零，因此
\[
c^\ast H_{n-1}c
=
\int |P_c(y)|^2\,\dd\mu_\varphi(y)>0.
\]

另一方面，若
\[
\dim \operatorname{span}\{v,Kv,\dots,K^{n-1}v\}<n,
\]
则这些向量线性相关。
其 Gram 矩阵为
\[
\bigl(\langle K^i v,K^j v\rangle\bigr)_{0\le i,j\le n-1}
=
\bigl(\langle v,K^{i+j}v\rangle\bigr)_{0\le i,j\le n-1}
=
H_{n-1},
\]
其中用到了 \(K\) 的自伴性及假设中的矩匹配。
但线性相关将迫使该 Gram 矩阵奇异，这与 \(H_{n-1}\) 的严格正定矛盾。
故必有
\[
\dim \operatorname{span}\{v,Kv,\dots,K^{n-1}v\}\ge n.
\]

对 \(J_\varphi^{[n]}\) 而言，\(E_n\) 本身即由
\[
e_0,\ J_\varphi^{[n]}e_0,\ \dots,\ (J_\varphi^{[n]})^{n-1}e_0
\]
生成，因此其维数正好为 \(n\)。
故它是最小 exact compressor。证毕。
\end{proof}

\begin{proposition}[principal rational Borel shadow 的极点几何与 natural strip 保持]\label{prop:xi-time-part9zbi-principal-rational-shadow-natural-strip}
设 \(J_\varphi^{[n]}\) 的正特征值按降序排列为
\[
\lambda_1^{[n]}>\cdots>\lambda_n^{[n]}>0,
\]
相应归一化特征向量为 \(v_1^{[n]},\dots,v_n^{[n]}\)，并记
\[
w_j^{[n]}:=|\langle v_0,v_j^{[n]}\rangle|^2.
\]
则 \(w_j^{[n]}>0\)，并且
\[
\widehat{\mathcal E}_\varphi^{[n]}(\xi)
=
\sum_{j=1}^{n}\frac{w_j^{[n]}}{1+\lambda_j^{[n]}\xi^2}.
\]
因此 \(\widehat{\mathcal E}_\varphi^{[n]}\) 的全部奇点恰为
\[
\xi=\pm \frac{i}{\sqrt{\lambda_j^{[n]}}}
\qquad (1\le j\le n),
\]
并满足
\[
|\Im \xi|\ge \frac{2\pi}{\varphi}.
\]
于是，每一个 finite principal rational shadow 都不会在 natural Borel strip
\[
|\Im \xi|<\frac{2\pi}{\varphi}
\]
内引入伪奇点。
\end{proposition}
\begin{proof}
有限维谱定理给出
\[
(I+\xi^2J_\varphi^{[n]})^{-1}
=
\sum_{j=1}^{n}\frac{1}{1+\lambda_j^{[n]}\xi^2}
\,
|v_j^{[n]}\rangle\langle v_j^{[n]}|.
\]
两侧取 \(v_0\) 的矩，即得
\[
\widehat{\mathcal E}_\varphi^{[n]}(\xi)
=
\sum_{j=1}^{n}\frac{|\langle v_0,v_j^{[n]}\rangle|^2}{1+\lambda_j^{[n]}\xi^2}.
\]

下面证明 \(w_j^{[n]}>0\)。
由 \(v_0=\sqrt{u_0}\,e_0\)，只需证明每个特征向量的第一坐标非零。
若某个归一化特征向量 \(v=(x_0,\dots,x_{n-1})\) 满足 \(x_0=0\)，则由特征方程首行
\[
\alpha_1 x_1+\beta_0 x_0=\lambda x_0
\]
得到 \(x_1=0\)。
再逐行递推，由 \(\alpha_j>0\) 可依次推出 \(x_2=\cdots=x_{n-1}=0\)，与 \(v\neq 0\) 矛盾。
故每个特征向量第一坐标皆非零，因而 \(w_j^{[n]}>0\)。

又因为 \(J_\varphi^{[n]}\) 是正压缩，
\[
0<\lambda_j^{[n]}\le \|J_\varphi^{[n]}\|\le \|J_\varphi\|=\rho_\ast=\frac{\varphi^2}{4\pi^2}.
\]
于是
\[
\frac{1}{\sqrt{\lambda_j^{[n]}}}
\ge
\frac{1}{\sqrt{\rho_\ast}}
=
\frac{2\pi}{\varphi}.
\]
故全部极点均位于边界之外。
证毕。
\end{proof}

\begin{theorem}[唯一越界回路给出的首个失配项精确公式]\label{thm:xi-time-part9zbi-first-mismatch-unique-excursion}
对每个 \(n\ge 1\)，成立精确恒等式
\[
u_{2n}-\langle v_0,(J_\varphi^{[n]})^{2n}v_0\rangle
=
u_0(\alpha_1\alpha_2\cdots \alpha_n)^2.
\]
等价地，
\[
\mathcal E_\varphi(s)-\mathcal E_\varphi^{[n]}(s)
=
(4n)!\,u_0(\alpha_1\alpha_2\cdots \alpha_n)^2\,s^{4n+1}
+O(s^{4n+3}).
\]
因此，principal compression 对 exact odd tail 的第一次失配，不仅阶数精确可知，而且其系数也被一个单一、唯一、显式的越界回路权重完全控制。
\end{theorem}
\begin{proof}
由定理~\ref{thm:xi-time-part9zbi-principal-compression-exact-moments}，所有 \(k\le 2n-1\) 的矩完全一致。
现在考察 \(k=2n\)。

在
\[
\langle e_0,J_\varphi^{2n}e_0\rangle
\]
的闭路展开中，若一条长度 \(2n\) 的闭路曾到达层 \(n\)，则它至少需要 \(n\) 次上行和 \(n\) 次下行。
由于总长度恰等于 \(2n\)，这迫使该闭路不能包含任何停留，也不能在到达层 \(n\) 之前发生额外回撤。
故这样的闭路唯一，即
\[
0\to 1\to 2\to \cdots\to n\to n-1\to \cdots\to 1\to 0.
\]
其权重恰为
\[
\alpha_1\alpha_2\cdots\alpha_n\cdot \alpha_n\cdots\alpha_2\alpha_1
=
(\alpha_1\alpha_2\cdots\alpha_n)^2.
\]

其余所有长度 \(2n\) 的闭路若不触及层 \(n\)，则完全落在
\[
\{0,1,\dots,n-1\}
\]
内，因此在压缩后仍完整保留。
于是
\[
\langle e_0,J_\varphi^{2n}e_0\rangle
-
\langle e_0,(J_\varphi^{[n]})^{2n}e_0\rangle
=
(\alpha_1\cdots\alpha_n)^2.
\]
两边同乘 \(u_0\) 即得
\[
u_{2n}-\langle v_0,(J_\varphi^{[n]})^{2n}v_0\rangle
=
u_0(\alpha_1\cdots\alpha_n)^2.
\]

最后，把这一差值代入
\[
\mathcal E_\varphi(s)
=
\sum_{k\ge 0}(-1)^k(2k)!\,u_k\,s^{2k+1},
\]
并注意在 \(k=2n\) 时符号为正，便得到
\[
\mathcal E_\varphi(s)-\mathcal E_\varphi^{[n]}(s)
=
(4n)!\,u_0(\alpha_1\cdots\alpha_n)^2\,s^{4n+1}
+O(s^{4n+3}).
\]
证毕。
\end{proof}

\begin{theorem}[Ritz--Wythoff 收敛定理]\label{thm:xi-time-part9zbi-ritz-wythoff-convergence}
设
\[
\lambda_1>\lambda_2>\cdots>0
\]
为 exact host \(J_\varphi\) 的非零特征值降序排列；
对每个 \(n\ge 1\)，设
\[
\lambda_1^{[n]}>\cdots>\lambda_n^{[n]}>0
\]
为 principal compression \(J_\varphi^{[n]}\) 的正特征值降序排列。
则对每个固定 \(j\ge 1\)，当 \(n\ge j\) 时有
\[
\lambda_j^{[n]}\uparrow \lambda_j
\qquad (n\to\infty).
\]
从而，若定义 principal rational shadow 的正虚轴极点
\[
p_j^{[n]}:=\frac{1}{\sqrt{\lambda_j^{[n]}}},
\qquad
p_j:=\frac{1}{\sqrt{\lambda_j}},
\]
则
\[
p_j^{[n]}\downarrow p_j.
\]
再结合定理~\ref{thm:xi-time-part9zbf-foldpi-wythoff-spectral-order}，可得
\[
p_{\lfloor m\varphi\rfloor}=\frac{2\pi m}{\varphi},
\qquad
p_{\lfloor m\varphi^2\rfloor}=2\pi m
\qquad (m\ge 1).
\]
故主子压缩所产生的 finite rational Borel shadows 并非任意趋近于 exact pole lattice，而是沿 Wythoff 序从外侧单调逼近该双梯 pole geometry。
\end{theorem}
\begin{proof}
由 Courant--Fischer 极小极大原理，
\[
\lambda_j^{[n]}
=
\max_{\substack{L\subset E_n\\ \dim L=j}}
\ \min_{\substack{x\in L\\ \|x\|=1}}
\langle J_\varphi x,x\rangle
\qquad (n\ge j),
\]
而
\[
\lambda_j
=
\max_{\substack{L\subset \ell^2(\NN_0)\\ \dim L=j}}
\ \min_{\substack{x\in L\\ \|x\|=1}}
\langle J_\varphi x,x\rangle.
\]
由于 \(E_n\subset E_{n+1}\subset \ell^2(\NN_0)\)，立得
\[
\lambda_j^{[n]}\le \lambda_j^{[n+1]}\le \lambda_j.
\]

下面证明极限等于 \(\lambda_j\)。
任取 \(\varepsilon>0\)。
由 \(J_\varphi\) 的极小极大表征，可取某个 \(j\)-维子空间 \(L_\varepsilon\subset \ell^2(\NN_0)\)，使
\[
\min_{\substack{x\in L_\varepsilon\\ \|x\|=1}}
\langle J_\varphi x,x\rangle
>
\lambda_j-\varepsilon.
\]
而 \(\bigcup_n E_n\) 在 \(\ell^2(\NN_0)\) 中稠密，因此对充分大的 \(n\)，可在 \(E_n\) 中找到某个 \(j\)-维子空间 \(L_{\varepsilon,n}\) 近似 \(L_\varepsilon\)，使其最小 Rayleigh 商仍满足
\[
\min_{\substack{x\in L_{\varepsilon,n}\\ \|x\|=1}}
\langle J_\varphi x,x\rangle
>
\lambda_j-2\varepsilon.
\]
于是
\[
\lambda_j^{[n]}\ge \lambda_j-2\varepsilon
\]
对充分大 \(n\) 成立。
结合上面的单调夹逼
\[
\lambda_j^{[n]}\le \lambda_j,
\]
便得到
\[
\lambda_j^{[n]}\uparrow \lambda_j.
\]

对极点 \(p_j^{[n]}\) 的结论只需应用函数 \(x\mapsto x^{-1/2}\) 在 \((0,\infty)\) 上的严格单调递减性。
最后，由定理~\ref{thm:xi-time-part9zbf-foldpi-wythoff-spectral-order}，
\[
\lambda_{\lfloor m\varphi\rfloor}=\frac{\varphi^2}{4\pi^2m^2},
\qquad
\lambda_{\lfloor m\varphi^2\rfloor}=\frac{1}{4\pi^2m^2},
\]
逐项取倒数平方根即得
\[
p_{\lfloor m\varphi\rfloor}=\frac{2\pi m}{\varphi},
\qquad
p_{\lfloor m\varphi^2\rfloor}=2\pi m.
\]
证毕。
\end{proof}

\begin{conjecture}[Ritz 误差的唯一越界回路主导律]\label{conj:xi-time-part9zbi-ritz-error-unique-excursion}
对每个固定 \(j\ge 1\)，存在常数 \(c_j>0\)，使得
\[
\lambda_j-\lambda_j^{[n]}
\sim
c_j\,u_0(\alpha_1\alpha_2\cdots\alpha_n)^2
\qquad (n\to\infty),
\]
等价地，
\[
p_j^{[n]}-p_j
\sim
\frac{c_j\,u_0}{2\lambda_j^{3/2}}
(\alpha_1\alpha_2\cdots\alpha_n)^2.
\]
这里的唯一自然小参数，正是定理~\ref{thm:xi-time-part9zbi-first-mismatch-unique-excursion} 中首个失配项所暴露的最短越界回路权重。
\end{conjecture}

\noindent
因此，H.111 odd Bernoulli tail 除了拥有一条由 finite-support truncation 构成的单调耗尽塔之外，还内生出另一条性质更强的压缩层级：principal compression 不仅在长度 \(2n-1\) 之前完全无损，而且在维数上是最小的 exact compressor；其第一次失配由唯一越界回路显式控制，而其 pole geometry 又沿 Wythoff 双梯从外侧单调逼近 exact 宿主。由此，项目中的 compact Jacobi 谱宿主不再只是一个抽象的 infinite-rank realization，而是同时携带一条最优维数、局部 exact、误差可定量、并与 Wythoff pole ordering 严格兼容的 Ritz--Jacobi 压缩层级。

\endinput

\paragraph{fold-\texorpdfstring{$\pi$}{pi} Jacobi 宿主的 Hausdorff 正性、Laguerre--P\'olya 刚性与纯点波迹}
沿用定理~\ref{thm:xi-time-part9zbe-foldpi-canonical-compact-jacobi}、
\ref{thm:xi-time-part9zbe-foldpi-jacobi-schatten-zeta}、
\ref{thm:xi-time-part9zbe-foldpi-fredholm-double-sine}、
\ref{thm:xi-time-part9zbe-foldpi-compact-determinant-sampling}
与定理~\ref{thm:xi-time-part9zbh-foldpi-zeta-regularized-determinant} 的记号。
前几节已经确定：H.111 的 odd tail 由一个 compact Jacobi 宿主 \(J_\varphi\) 精确承载，其 operator \(\zeta\)、Fredholm 行列式以及 \(\zeta\)-正则化行列式都完全显式。
在此基础上，还可继续抽出一条此前未被终端化的谱正性链：shifted spectral zeta samples 本身形成 Hausdorff 矩序列；负轴 Fredholm determinant 落入 Laguerre--P\'olya / PF\(_\infty\) 与 complete Bernstein 锥；而 Weyl 余项与波迹则分别出现精确三角极限律与纯点 Dirac-comb 公式。

\leanverified{paper\_xi\_time\_part9zbj\_shifted\_spectral\_zeta\_hausdorff}
\begin{theorem}[shifted spectral \texorpdfstring{$\zeta$}{zeta} samples 的 Hausdorff 矩实现]\label{thm:xi-time-part9zbj-shifted-spectral-zeta-hausdorff}
对每个 \(r\ge 0\)，定义
\[
M_r:=\zeta_{J_\varphi}(r+1).
\]
则
\[
M_r
=
(-1)^r\,
\frac{1+\varphi^{2r+2}}
{4(2r)!\,(\varphi^{-1}+\varphi^{2r-1})}
\,
\mathfrak a_{r+1}^{(\pi)}.
\]
并且若定义有限正离散测度
\[
\nu_\varphi^{\mathrm{spec}}
:=
\sum_{n\ge 1}\frac{1}{4\pi^2n^2}\,
\delta_{(4\pi^2n^2)^{-1}}
+
\sum_{n\ge 1}\frac{\varphi^2}{4\pi^2n^2}\,
\delta_{\varphi^2(4\pi^2n^2)^{-1}},
\]
则
\[
M_r=\int_0^{\rho_\ast} y^r\,\dd\nu_\varphi^{\mathrm{spec}}(y)
\qquad (r\ge 0),
\qquad
\rho_\ast=\frac{\varphi^2}{4\pi^2}.
\]
换言之，H.111 的 odd Bernoulli layer 经这一显式重标定后，得到的不是任意正数列，而是一条 Hausdorff 矩序列。
\end{theorem}
\begin{proof}
由定理~\ref{thm:xi-time-part9zbe-foldpi-canonical-compact-jacobi}，
\[
\sigma(J_\varphi)\setminus\{0\}
=
\left\{\frac{1}{4\pi^2n^2}:n\ge 1\right\}
\cup
\left\{\frac{\varphi^2}{4\pi^2n^2}:n\ge 1\right\}.
\]
再由定理~\ref{thm:xi-time-part9zbe-foldpi-jacobi-schatten-zeta}，
\[
M_r=\zeta_{J_\varphi}(r+1)
=
\sum_{\lambda\in \sigma(J_\varphi)\setminus\{0\}}\lambda^{r+1}.
\]
因此若取
\[
\nu_\varphi^{\mathrm{spec}}
:=
\sum_{\lambda\in \sigma(J_\varphi)\setminus\{0\}}\lambda\,\delta_\lambda,
\]
便有
\[
\int_0^{\rho_\ast} y^r\,\dd\nu_\varphi^{\mathrm{spec}}(y)
=
\sum_{\lambda}\lambda^{r+1}
=
M_r.
\]
这正给出所述 Hausdorff 矩表示。
又因为
\[
\nu_\varphi^{\mathrm{spec}}\bigl([0,\rho_\ast]\bigr)
=
\sum_{\lambda}\lambda
=
\operatorname{Tr}(J_\varphi)<\infty,
\]
故 \(\nu_\varphi^{\mathrm{spec}}\) 为有限正测度。

最后，由定理~\ref{thm:xi-time-part9zbe-foldpi-compact-determinant-sampling}，
\[
\mathfrak a_{r+1}^{(\pi)}
=
(-1)^r\,
4(2r)!\,
\frac{\varphi^{-1}+\varphi^{2r-1}}{1+\varphi^{2r+2}}\,
\zeta_{J_\varphi}(r+1),
\]
反解即得
\[
M_r
=
(-1)^r\,
\frac{1+\varphi^{2r+2}}
{4(2r)!\,(\varphi^{-1}+\varphi^{2r-1})}
\,
\mathfrak a_{r+1}^{(\pi)}.
\]
证毕。
\end{proof}

\begin{corollary}[shifted spectral moments 的 Hankel--localizing 双正性]\label{cor:xi-time-part9zbj-shifted-spectral-zeta-hankel-localizing}
\leanverified{paper\_xi\_time\_part9zbj\_shifted\_spectral\_zeta\_hankel\_localizing}
对任意 \(m,\ell\ge 0\)，定义
\[
H_{m,\ell}:=(M_{\ell+i+j})_{0\le i,j\le m},
\qquad
L_{m,\ell}:=(\rho_\ast M_{\ell+i+j}-M_{\ell+i+j+1})_{0\le i,j\le m}.
\]
则
\[
H_{m,\ell}\succeq 0,
\qquad
L_{m,\ell}\succeq 0.
\]
特别地，
\[
M_r^2\le M_{r-1}M_{r+1}
\qquad (r\ge 1),
\]
故这条重标定后的 odd-layer 序列自动落入 Hausdorff 矩锥。
\end{corollary}
\begin{proof}
对任意 \(u=(u_0,\dots,u_m)\in\RR^{m+1}\)，记
\[
P_u(y):=\sum_{k=0}^{m}u_k y^k.
\]
则
\[
u^\top H_{m,\ell}u
=
\sum_{i,j=0}^{m}u_i u_j M_{\ell+i+j}
=
\int_0^{\rho_\ast} y^\ell |P_u(y)|^2\,\dd\nu_\varphi^{\mathrm{spec}}(y)\ge 0.
\]
同理，
\[
u^\top L_{m,\ell}u
=
\int_0^{\rho_\ast} y^\ell(\rho_\ast-y)|P_u(y)|^2\,\dd\nu_\varphi^{\mathrm{spec}}(y)\ge 0.
\]
故 \(H_{m,\ell}\) 与 \(L_{m,\ell}\) 皆正半定。

取 \(m=1\) 与 \(\ell=r-1\)，便得
\[
\det
\begin{pmatrix}
M_{r-1} & M_r\\
M_r & M_{r+1}
\end{pmatrix}
\ge 0,
\]
即
\[
M_r^2\le M_{r-1}M_{r+1}.
\]
证毕。
\end{proof}

\begin{theorem}[negative-axis Fredholm determinant 的 Laguerre--P\'olya / PF\(_\infty\) 刚性]\label{thm:xi-time-part9zbj-lp-pf-rigidity}
\leanverified{paper\_xi\_time\_part9zbj\_lp\_pf\_rigidity}
定义
\[
G_\varphi(z):=D_\varphi(-z)
=
\frac{\sinh(\sqrt z/2)}{\sqrt z/2}
\,
\frac{\sinh(\varphi\sqrt z/2)}{\varphi\sqrt z/2}
=
\sum_{m\ge 0}c_m z^m.
\]
则 \(G_\varphi\) 属于 Laguerre--P\'olya I 类。
因此：
\begin{enumerate}
  \item \(G_\varphi\) 的全部零点都位于 \((-\infty,0)\)；
  \item 系数列 \((c_m)_{m\ge 0}\) 是 P\'olya frequency 序列 PF\(_\infty\)；
  \item \(G_\varphi\) 的一切 Jensen 多项式均为超双曲；
  \item 特别地，所有 Toeplitz 小行列式皆非负，且
  \[
  c_m^2\ge c_{m-1}c_{m+1}
  \qquad (m\ge 1).
  \]
\end{enumerate}
此外，系数满足显式公式
\[
c_m
=
\frac{1}{4^m}
\sum_{j=0}^{m}
\frac{\varphi^{2j}}
{(2j+1)!\,(2m-2j+1)!}.
\]
\end{theorem}
\begin{proof}
由定理~\ref{thm:xi-time-part9zbe-foldpi-fredholm-double-sine}，
\[
G_\varphi(z)
=
\prod_{n\ge 1}
\left(1+\frac{z}{4\pi^2n^2}\right)
\left(1+\frac{\varphi^2 z}{4\pi^2n^2}\right).
\]
而
\[
\sum_{n\ge 1}
\left(
\frac{1}{4\pi^2n^2}
+
\frac{\varphi^2}{4\pi^2n^2}
\right)
<\infty,
\]
故该 genus \(0\) 乘积在紧集上一致收敛，并且每个有限截断多项式都只具有实负零点。
因此 \(G_\varphi\) 是一列实根仅位于负轴的多项式的局部一致极限，从而属于 Laguerre--P\'olya I 类。

由 Laguerre--P\'olya I 类与 PF\(_\infty\) 生成函数之间的经典 Aissen--Schoenberg--Whitney--Edrei 对应，以及 Jensen 多项式的超双曲性判据，立得 2--4。

最后，由
\[
\frac{\sinh(\sqrt z/2)}{\sqrt z/2}
=
\sum_{a\ge 0}\frac{z^a}{4^a(2a+1)!},
\qquad
\frac{\sinh(\varphi\sqrt z/2)}{\varphi\sqrt z/2}
=
\sum_{b\ge 0}\frac{\varphi^{2b}z^b}{4^b(2b+1)!},
\]
作 Cauchy 乘积便得到
\[
c_m
=
\frac{1}{4^m}
\sum_{j=0}^{m}
\frac{\varphi^{2j}}
{(2j+1)!\,(2m-2j+1)!}.
\]
证毕。
\end{proof}

\begin{theorem}[negative-axis spectral action 的 complete Bernstein 实现]\label{thm:xi-time-part9zbj-complete-bernstein-spectral-action}
定义
\[
F_\varphi(x):=-\log D_\varphi(-x)
\qquad (x>0).
\]
则 \(F_\varphi\) 是 complete Bernstein 函数。
更精确地，
\[
F_\varphi'(x)
=
\sum_{n\ge 1}\frac{1}{x+4\pi^2n^2}
+
\sum_{n\ge 1}\frac{1}{x+4\pi^2n^2/\varphi^2},
\]
故 \(F_\varphi'\) 是 Stieltjes 函数。
并且 \(F_\varphi\) 具有显式 L\'evy 表示
\[
F_\varphi(x)
=
\int_0^\infty (1-e^{-xt})\,\eta_\varphi(t)\,\dd t,
\]
其中
\[
\eta_\varphi(t)
=
\frac{1}{t}
\sum_{n\ge 1}
\left(
e^{-4\pi^2n^2 t}
+
e^{-4\pi^2n^2 t/\varphi^2}
\right)
=
\frac{\vartheta_3(0,e^{-4\pi^2 t})+\vartheta_3(0,e^{-4\pi^2 t/\varphi^2})-2}{2t}.
\]
\end{theorem}
\begin{proof}
由定理~\ref{thm:xi-time-part9zbe-foldpi-fredholm-double-sine} 的谱乘积，
\[
F_\varphi(x)
=
\sum_{\lambda\in \sigma(J_\varphi)\setminus\{0\}}\log(1+x\lambda).
\]
逐项求导即得
\[
F_\varphi'(x)
=
\sum_{\lambda}\frac{\lambda}{1+x\lambda}
=
\sum_{\lambda}\frac{1}{x+\lambda^{-1}}.
\]
由于 \(\lambda^{-1}\) 恰为两列 reciprocal-square ladders
\[
\{4\pi^2n^2:\ n\ge 1\}
\cup
\left\{\frac{4\pi^2n^2}{\varphi^2}:\ n\ge 1\right\},
\]
于是
\[
F_\varphi'(x)
=
\sum_{n\ge 1}\frac{1}{x+4\pi^2n^2}
+
\sum_{n\ge 1}\frac{1}{x+4\pi^2n^2/\varphi^2}.
\]
这给出 Stieltjes 表示，因此 \(F_\varphi\) 为 complete Bernstein 函数。

另一方面，对每个 \(\lambda>0\) 都有
\[
\log(1+x\lambda)
=
\int_0^\infty (1-e^{-xt})\,e^{-t/\lambda}\,\frac{\dd t}{t}.
\]
逐项求和便得
\[
F_\varphi(x)
=
\int_0^\infty (1-e^{-xt})
\left[
\frac{1}{t}
\sum_{\lambda\in \sigma(J_\varphi)\setminus\{0\}}e^{-t/\lambda}
\right]\dd t.
\]
再把两列谱显式展开，即得
\[
\eta_\varphi(t)
=
\frac{1}{t}
\sum_{n\ge 1}
\left(
e^{-4\pi^2n^2 t}
+
e^{-4\pi^2n^2 t/\varphi^2}
\right).
\]
最后用
\[
\vartheta_3(0,q)=1+2\sum_{n\ge 1}q^{n^2}
\]
改写即可。证毕。
\end{proof}

\begin{theorem}[Weyl 线性坐标中的精确三角涨落律]\label{thm:xi-time-part9zbj-weyl-triangular-fluctuation}
定义
\[
R_\varphi(x)
:=
N_{J_\varphi}\!\bigl((2\pi x)^{-2}\bigr)-\varphi^2 x
\qquad (x>0).
\]
则
\[
R_\varphi(x)=-\{x\}-\{\varphi x\}\in [-2,0].
\]
并且对任意有界连续函数 \(f\)，有
\[
\lim_{T\to\infty}\frac{1}{T}\int_0^T f(R_\varphi(x))\,\dd x
=
\int_{-2}^{0}f(r)\rho_\triangle(r)\,\dd r,
\]
其中
\[
\rho_\triangle(r)=
\begin{cases}
r+2,& -2\le r\le -1,\\[1mm]
-r,& -1\le r\le 0,\\[1mm]
0,& \text{otherwise}.
\end{cases}
\]
故该 host 的 Weyl 余项不只是 \(O(1)\) 有界，而且具有精确极限分布；其均值与方差分别为
\[
\mathbb E(R_\varphi)=-1,
\qquad
\operatorname{Var}(R_\varphi)=\frac16.
\]
\end{theorem}
\begin{proof}
由定理~\ref{thm:xi-time-part9zbe-foldpi-jacobi-schatten-zeta} 中的精确计数公式，
\[
N_{J_\varphi}(\lambda)
=
\left\lfloor\frac{1}{2\pi\sqrt\lambda}\right\rfloor
+
\left\lfloor\frac{\varphi}{2\pi\sqrt\lambda}\right\rfloor.
\]
令 \(\lambda=(2\pi x)^{-2}\)，便得到
\[
N_{J_\varphi}\!\bigl((2\pi x)^{-2}\bigr)
=
\lfloor x\rfloor+\lfloor \varphi x\rfloor.
\]
再由 \(\varphi^2=1+\varphi\)，有
\[
R_\varphi(x)
=
\lfloor x\rfloor+\lfloor \varphi x\rfloor-(1+\varphi)x
=
-\{x\}-\{\varphi x\}.
\]
故 \(R_\varphi(x)\in[-2,0]\)。

由于 \(1\) 与 \(\varphi\) 在 \(\QQ\) 上线性无关，Kronecker--Weyl 定理表明
\[
x\longmapsto (x,\varphi x)\pmod 1
\]
在 \(\TT^2\) 上等分布。
于是 \(R_\varphi\) 的极限分布就是单位方形上的 Haar 测度经映射
\[
(u,v)\longmapsto -(u+v)
\]
的推前。
该推前密度正是两份 \([0,1]\) 上均匀分布卷积后的三角律，故得到
\[
\rho_\triangle(r)=
\begin{cases}
r+2,& -2\le r\le -1,\\[1mm]
-r,& -1\le r\le 0,\\[1mm]
0,& \text{otherwise}.
\end{cases}
\]
其均值与方差是该三角分布的一阶、二阶矩，分别为 \(-1\) 与 \(1/6\)。
证毕。
\end{proof}

\begin{proposition}[pure-point Poisson--Jacobi 波迹公式]\label{prop:xi-time-part9zbj-pure-point-poisson-jacobi-wave-trace}
\leanverified{paper\_xi\_time\_part9zbj\_pure\_point\_poisson\_jacobi\_wave\_trace}
在 tempered distribution 意义下，定义
\[
\mathcal W_\varphi(x)
:=
2\sum_{\lambda\in \sigma(J_\varphi)\setminus\{0\}}
\cos\!\left(\frac{x}{\sqrt\lambda}\right).
\]
则成立严格恒等式
\[
\mathcal W_\varphi(x)
=
\sum_{m\in\ZZ}\delta(x-m)
+
\varphi\sum_{m\in\ZZ}\delta(x-m\varphi)
-2.
\]
因此 fold-\(\pi\) Jacobi host 的 wave trace 是纯点型的；其奇支撑恰由两条基本长度格 \(\ZZ\) 与 \(\varphi\ZZ\) 构成，而不存在连续背景项。
\end{proposition}
\begin{proof}
由定理~\ref{thm:xi-time-part9zbe-foldpi-canonical-compact-jacobi}，
\[
\sigma(J_\varphi)\setminus\{0\}
=
\left\{\frac{1}{4\pi^2n^2}:n\ge 1\right\}
\cup
\left\{\frac{\varphi^2}{4\pi^2n^2}:n\ge 1\right\}.
\]
故
\[
\mathcal W_\varphi(x)
=
2\sum_{n\ge 1}\cos(2\pi n x)
+
2\sum_{n\ge 1}\cos\!\left(\frac{2\pi n x}{\varphi}\right).
\]
对第一项，应用标准 Dirac comb 恒等式
\[
\sum_{n\in\ZZ}e^{2\pi i n x}
=
\sum_{m\in\ZZ}\delta(x-m),
\]
即可得到
\[
2\sum_{n\ge 1}\cos(2\pi n x)
=
\sum_{m\in\ZZ}\delta(x-m)-1.
\]
对第二项，同样有
\[
\sum_{n\in\ZZ}e^{2\pi i n x/\varphi}
=
\varphi\sum_{m\in\ZZ}\delta(x-m\varphi),
\]
因此
\[
2\sum_{n\ge 1}\cos\!\left(\frac{2\pi n x}{\varphi}\right)
=
\varphi\sum_{m\in\ZZ}\delta(x-m\varphi)-1.
\]
两式相加即得所述分布恒等式。
证毕。
\end{proof}

\begin{conjecture}[complete-Bernstein 型 Écalle--Fredholm 提升]\label{conj:xi-time-part9zbj-complete-bernstein-ecalle-fredholm}
沿用猜想~\ref{conj:xi-time-part9zbe-foldpi-ecalle-fredholm-lift} 中的 formal 变换 \(\mathcal T_\varphi\)。
存在一个 complete Bernstein 函数
\[
B_\varphi:(0,\infty)\to \RR
\]
使得
\[
B_\varphi(s)\sim \mathcal T_\varphi\!\bigl[-\log D_\varphi\bigr](s)
\qquad (s\downarrow 0),
\]
并且对一切 \(m\ge 1\)，都有精确采样公式
\[
B_\varphi\!\left(\left(\frac{\varphi}{2}\right)^m\right)
=
\log C_m-\log(2\pi).
\]
若此猜想成立，则 exact \(\pi\)-tail 的非形式提升将不再是任意 resurgent germ，而必定落在 Pick--Bernstein 锥内。
特别地，exact 常数序列将自动继承 complete Bernstein 函数的解析正性：
\[
(-1)^{k-1}B_\varphi^{(k)}(s)\ge 0,
\qquad
\text{Stieltjes--L\'evy 表示},
\qquad
\text{Hankel 正性}.
\]
\end{conjecture}

\noindent
因此，H.111 odd tail 在现有 compact Jacobi 主链下不仅具有一个黄金 \(\zeta\)-正则化行列式 \(\det_\zeta(J_\varphi)=\varphi\)，而且还同时落入三种经典正性几何：shifted spectral zeta samples 落入 Hausdorff 矩锥，negative-axis Fredholm determinant 落入 Laguerre--P\'olya / PF\(_\infty\) 锥，negative-axis spectral action 则落入 complete Bernstein 锥；与此同时，其 Weyl 余项与波迹分别显示出精确三角涨落律与纯点 Dirac-comb 结构。由此，fold-\(\pi\) Jacobi host 的谱对象不再只是一个 compact determinant 模型，而是进一步呈现出一条相互兼容的 positivity--rigidity--wave-trace 链。

\endinput

\paragraph{fold-\texorpdfstring{$\pi$}{pi} odd tail 的 Gaussian compressor、Pad\'e shadow 与 exactness ceiling}
沿用 \(\mu_\varphi\)、\(u_k\)、\(J_\varphi\)、\(J_\varphi^{[n]}\) 与 \(\rho_\ast=\varphi^2/(4\pi^2)\) 的记号。
进一步把与 \(\mu_\varphi\) 对应的 Stieltjes 变换写为
\[
m_\varphi(z)
:=
\int_{0}^{\rho_\ast}\frac{\dd\mu_\varphi(y)}{1-yz}
=
\langle v_0,(I-zJ_\varphi)^{-1}v_0\rangle,
\]
并定义 principal compression 的有理 shadow
\[
m_{\varphi,n}(z)
:=
\langle v_0,(I-zJ_\varphi^{[n]})^{-1}v_0\rangle.
\]
再记 \(\pi_{\varphi,n}\) 为 \(\mu_\varphi\) 的首一 \(n\) 次正交多项式。
上一节已经证明：\(J_\varphi^{[n]}\) 在前 \(2n\) 个 odd Bernoulli layers 上精确保留 exact host 的矩数据，并且在维数上是最小的 exact compressor。
这里进一步把这一事实提升为一条 Gaussian quadrature / Pad\'e 刚性链：给定极点预算 \(n\)，project odd tail 所允许的正 Stieltjes shadow 不是许多，而是唯一的 Gaussian--Pad\'e shadow；其可达精确阶数也不是模糊增长，而恰由 \(2n\)-moment ceiling 锁定。

\begin{theorem}[principal compression 的 Gaussian quadrature 唯一性]\label{thm:xi-time-part9zbk-gaussian-compressor-uniqueness}
\leanverified{paper\_xi\_time\_part9zbk\_gaussian\_compressor\_uniqueness}
对每个 \(n\ge 1\)，存在唯一的正 \(n\)-原子测度
\[
\nu_{\varphi,n}
=
\sum_{j=1}^{n}w_j^{[n]}\,\delta_{\lambda_j^{[n]}},
\qquad
w_j^{[n]}>0,
\qquad
0<\lambda_n^{[n]}<\cdots<\lambda_1^{[n]}\le \rho_\ast,
\]
使得
\[
\int y^k\,\dd\nu_{\varphi,n}(y)=u_k
\qquad (0\le k\le 2n-1).
\]
并且该唯一测度正是主子压缩 \(J_\varphi^{[n]}\) 在循环向量 \(v_0\) 处的谱测度：
\[
\nu_{\varphi,n}(B)=\langle v_0,\mathbf 1_B(J_\varphi^{[n]})v_0\rangle.
\]
换言之，Ritz--Jacobi principal compression 正是 H.111 odd tail 矩问题在节点预算 \(n\) 下的唯一 Gaussian compressor。
\end{theorem}
\begin{proof}
由定理~\ref{thm:xi-time-part9zbi-principal-compression-exact-moments}，
\[
\langle v_0,(J_\varphi^{[n]})^k v_0\rangle=u_k
\qquad (0\le k\le 2n-1).
\]
再由有限维谱定理，
\[
\nu_{\varphi,n}(B):=\langle v_0,\mathbf 1_B(J_\varphi^{[n]})v_0\rangle
\]
是一条正 \(n\)-原子测度，并满足
\[
\int y^k\,\dd\nu_{\varphi,n}(y)
=
\langle v_0,(J_\varphi^{[n]})^k v_0\rangle
=
u_k
\qquad (0\le k\le 2n-1).
\]
由命题~\ref{prop:xi-time-part9zbi-principal-rational-shadow-natural-strip}，各特征权重 \(w_j^{[n]}\) 全部严格为正，故存在性成立。

下面证明唯一性。
设
\[
\nu=\sum_{j=1}^{n}c_j\delta_{x_j}
\]
是任意正 \(n\)-原子测度，满足
\[
\int y^k\,\dd\nu(y)=u_k
\qquad (0\le k\le 2n-1).
\]
由于 \(\pi_{\varphi,n}\) 是 \(\mu_\varphi\) 的首一 \(n\) 次正交多项式，故对每个 \(0\le k\le n-1\) 都有
\[
\int y^k\pi_{\varphi,n}(y)\,\dd\mu_\varphi(y)=0.
\]
由矩匹配立得
\[
\int y^k\pi_{\varphi,n}(y)\,\dd\nu(y)=0
\qquad (0\le k\le n-1).
\]

现在取次数 \(<n\) 的插值多项式 \(q\)，使
\[
q(x_j)=\pi_{\varphi,n}(x_j)
\qquad (1\le j\le n).
\]
则 \(\deg(q\pi_{\varphi,n})\le 2n-1\)，因而
\[
0=\int q(y)\pi_{\varphi,n}(y)\,\dd\nu(y)
=
\sum_{j=1}^{n}c_j\,\pi_{\varphi,n}(x_j)^2.
\]
由于 \(c_j>0\)，只能有
\[
\pi_{\varphi,n}(x_j)=0
\qquad (1\le j\le n).
\]
因此 \(\{x_1,\dots,x_n\}\) 必为 \(\pi_{\varphi,n}\) 的零点集。
节点既定后，权重由前 \(n\) 个矩条件
\[
\int y^k\,\dd\nu(y)=u_k
\qquad (0\le k\le n-1)
\]
给出的 Vandermonde 线性系统唯一决定。
故 \(\nu=\nu_{\varphi,n}\)。
证毕。
\end{proof}

\begin{theorem}[positive \texorpdfstring{$n$}{n}-atomic exactness ceiling]\label{thm:xi-time-part9zbk-positive-natomic-exactness-ceiling}
\leanverified{paper\_xi\_time\_part9zbk\_positive\_natomic\_exactness\_ceiling}
不存在任何正 \(n\)-原子测度 \(\nu\) 满足
\[
\int y^k\,\dd\nu(y)=u_k
\qquad (0\le k\le 2n).
\]
等价地，不存在任何具有 \(n\) 个正极点的 Stieltjes 型有理 shadow，能够精确复制 H.111 odd tail 的前 \(2n+1\) 个矩层。
\end{theorem}
\begin{proof}
反设存在这样的正 \(n\)-原子测度 \(\nu\)。
由定理~\ref{thm:xi-time-part9zbk-gaussian-compressor-uniqueness}，匹配到 \(2n-1\) 阶矩已经迫使其支撑必须等于 \(\pi_{\varphi,n}\) 的零点集。
于是
\[
\int \pi_{\varphi,n}(y)^2\,\dd\nu(y)=0.
\]
但 \(\deg(\pi_{\varphi,n}^2)=2n\)，而反设又给出直到 \(2n\) 阶的矩匹配，因此
\[
\int \pi_{\varphi,n}(y)^2\,\dd\mu_\varphi(y)
=
\int \pi_{\varphi,n}(y)^2\,\dd\nu(y)
=
0.
\]
这与 \(\mu_\varphi\) 为正测度且支撑无穷而 \(\pi_{\varphi,n}\not\equiv 0\) 矛盾。
故这样的 \(\nu\) 不存在。
证毕。
\end{proof}

\begin{theorem}[Christoffel--Gauss 平方余项公式]\label{thm:xi-time-part9zbk-christoffel-gauss-square-remainder}
\leanverified{paper\_xi\_time\_part9zbk\_christoffel\_gauss\_square\_remainder}
对每个 \(n\ge 1\) 与每个实数
\[
x\in(-\rho_\ast^{-1},\infty)\setminus\{0\}
\]
成立精确恒等式
\[
m_\varphi(-x)-m_{\varphi,n}(-x)
=
\frac{1}{\pi_{\varphi,n}(-1/x)^2}
\int_0^{\rho_\ast}
\frac{\pi_{\varphi,n}(y)^2}{1+xy}\,\dd\mu_\varphi(y).
\]
因此，principal compression 的误差不是抽象 \(O(\cdot)\) 余项，而是一个正核平方积分。
\end{theorem}
\begin{proof}
记
\[
\sigma_n:=\mu_\varphi-\nu_{\varphi,n}.
\]
由定理~\ref{thm:xi-time-part9zbk-gaussian-compressor-uniqueness}，
\[
\int y^k\,\dd\sigma_n(y)=0
\qquad (0\le k\le 2n-1).
\]

对固定的 \(x\in(-\rho_\ast^{-1},\infty)\setminus\{0\}\)，定义
\[
R_{n,x}(y)
:=
\frac{\pi_{\varphi,n}(y)^2-\pi_{\varphi,n}(-1/x)^2}{y+1/x}.
\]
由于分子在 \(y=-1/x\) 处为零，\(R_{n,x}\) 是次数至多 \(2n-1\) 的多项式。
于是
\[
\int R_{n,x}(y)\,\dd\sigma_n(y)=0.
\]
再利用
\[
y+\frac1x=\frac{1+xy}{x},
\]
便得到
\[
0=
\int
\frac{\pi_{\varphi,n}(y)^2-\pi_{\varphi,n}(-1/x)^2}{1+xy}
\,
\dd\sigma_n(y).
\]
整理得
\[
\pi_{\varphi,n}(-1/x)^2
\int\frac{\dd\sigma_n(y)}{1+xy}
=
\int\frac{\pi_{\varphi,n}(y)^2}{1+xy}\,\dd\sigma_n(y).
\]

另一方面，由定理~\ref{thm:xi-time-part9zbk-gaussian-compressor-uniqueness}，\(\nu_{\varphi,n}\) 的支撑正是 \(\pi_{\varphi,n}\) 的零点集，因此
\[
\int\frac{\pi_{\varphi,n}(y)^2}{1+xy}\,\dd\nu_{\varphi,n}(y)=0.
\]
于是右端化为
\[
\int\frac{\pi_{\varphi,n}(y)^2}{1+xy}\,\dd\mu_\varphi(y).
\]
左端则正是
\[
\pi_{\varphi,n}(-1/x)^2\bigl(m_\varphi(-x)-m_{\varphi,n}(-x)\bigr).
\]
两式合并便得所述公式。
证毕。
\end{proof}

\begin{corollary}[pre-singular 轴上的严格单侧逼近]\label{cor:xi-time-part9zbk-presingular-one-sided-approximation}
记
\[
\widehat{\mathcal E}_\varphi(\xi):=m_\varphi(-\xi^2),
\qquad
\widehat{\mathcal E}_{\varphi,n}(\xi):=m_{\varphi,n}(-\xi^2).
\]
则有：
\begin{enumerate}
  \item 对任意实 \(\xi\neq 0\)，
  \[
  0<\widehat{\mathcal E}_{\varphi,n}(\xi)<\widehat{\mathcal E}_\varphi(\xi);
  \]
  \item 对任意实 \(\tau\) 满足
  \[
  |\tau|<\frac{2\pi}{\varphi},
  \]
  有
  \[
  0<\widehat{\mathcal E}_{\varphi,n}(i\tau)<\widehat{\mathcal E}_\varphi(i\tau).
  \]
\end{enumerate}
因此，principal rational shadows 在 natural Borel strip 中对 exact shadow 给出严格单侧下逼近，而不只是收敛。
\end{corollary}
\begin{proof}
对实 \(\xi\neq 0\)，取 \(x=\xi^2>0\)。
由定理~\ref{thm:xi-time-part9zbk-christoffel-gauss-square-remainder}，
\[
\widehat{\mathcal E}_\varphi(\xi)-\widehat{\mathcal E}_{\varphi,n}(\xi)
=
\frac{1}{\pi_{\varphi,n}(-1/\xi^2)^2}
\int_0^{\rho_\ast}
\frac{\pi_{\varphi,n}(y)^2}{1+\xi^2 y}\,\dd\mu_\varphi(y).
\]
因为 \(1+\xi^2y>0\) 对 \(y\in[0,\rho_\ast]\) 成立，且被积函数非负而不恒零，右端严格为正，故得到第一条。

对纯虚轴，取 \(\xi=i\tau\)，则
\[
\widehat{\mathcal E}_\varphi(i\tau)=m_\varphi(\tau^2),
\qquad
\widehat{\mathcal E}_{\varphi,n}(i\tau)=m_{\varphi,n}(\tau^2).
\]
若 \(|\tau|<2\pi/\varphi\)，则
\[
\tau^2<\rho_\ast^{-1},
\]
因此可把定理~\ref{thm:xi-time-part9zbk-christoffel-gauss-square-remainder} 应用于 \(x=-\tau^2\in(-\rho_\ast^{-1},0)\)，得到
\[
\widehat{\mathcal E}_\varphi(i\tau)-\widehat{\mathcal E}_{\varphi,n}(i\tau)
=
\frac{1}{\pi_{\varphi,n}(1/\tau^2)^2}
\int_0^{\rho_\ast}
\frac{\pi_{\varphi,n}(y)^2}{1-\tau^2 y}\,\dd\mu_\varphi(y).
\]
由于 \(1-\tau^2 y>0\) 在 \([0,\rho_\ast]\) 上仍然成立，右端严格为正。
于是第二条亦成立。
证毕。
\end{proof}

\begin{theorem}[Pad\'e--Jacobi 恒等定理]\label{thm:xi-time-part9zbk-pade-jacobi-identity}
\leanverified{paper\_xi\_time\_part9zbk\_pade\_jacobi\_identity}
对每个 \(n\ge 1\)，\(m_{\varphi,n}\) 正是 \(m_\varphi\) 在 \(z=0\) 处唯一的 \([n-1/n]\)-Pad\'e 逼近。
也就是说，存在唯一一对多项式
\[
P_{n-1},Q_n,
\qquad
\deg P_{n-1}\le n-1,\ \deg Q_n\le n,\ Q_n(0)=1,
\]
使得
\[
m_\varphi(z)-\frac{P_{n-1}(z)}{Q_n(z)}=O(z^{2n}),
\]
并且该有理函数恰为
\[
\frac{P_{n-1}(z)}{Q_n(z)}=m_{\varphi,n}(z).
\]
\end{theorem}
\begin{proof}
由定理~\ref{thm:xi-time-part9zbk-gaussian-compressor-uniqueness}，
\[
m_{\varphi,n}(z)
=
\int\frac{\dd\nu_{\varphi,n}(y)}{1-yz}
=
\sum_{k\ge 0}\left(\int y^k\,\dd\nu_{\varphi,n}(y)\right)z^k
=
\sum_{k=0}^{2n-1}u_k z^k+O(z^{2n}).
\]
而
\[
m_\varphi(z)=\sum_{k\ge 0}u_k z^k.
\]
因此
\[
m_\varphi(z)-m_{\varphi,n}(z)=O(z^{2n}).
\]
又因为 \(\nu_{\varphi,n}\) 是 \(n\)-原子测度，\(m_{\varphi,n}\) 的分母次数至多为 \(n\)，分子次数至多为 \(n-1\)，故存在性成立。

再证唯一性。
若
\[
m_\varphi(z)-\frac{P_{n-1}(z)}{Q_n(z)}=O(z^{2n}),
\qquad
m_\varphi(z)-\frac{\widetilde P_{n-1}(z)}{\widetilde Q_n(z)}=O(z^{2n}),
\]
其中 \(Q_n(0)=\widetilde Q_n(0)=1\)，则相减得
\[
\frac{P_{n-1}\widetilde Q_n-\widetilde P_{n-1}Q_n}{Q_n\widetilde Q_n}
=
O(z^{2n}).
\]
但分子是次数至多 \(2n-1\) 的多项式，却在 \(z=0\) 处具有至少 \(2n\) 阶零点，因此只能恒等为零。
故
\[
P_{n-1}\widetilde Q_n-\widetilde P_{n-1}Q_n\equiv 0,
\]
从而两有理函数相同。
证毕。
\end{proof}

\begin{corollary}[finite-depth scalar continued fraction 的刚性坍缩]\label{cor:xi-time-part9zbk-continued-fraction-rigidity-collapse}
设 \(R_n(z)\) 是任意深度 \(n\) 的 terminating scalar continued fraction，其值在 \(z=0\) 邻域解析，并化简后满足
\[
R_n(z)=\frac{P_{n-1}(z)}{Q_n(z)},
\qquad
\deg P_{n-1}\le n-1,\ \deg Q_n\le n,\ Q_n(0)=1.
\]
若
\[
m_\varphi(z)-R_n(z)=O(z^{2n}),
\]
则必有
\[
R_n(z)=m_{\varphi,n}(z).
\]
因此，在 fixed depth \(n\) 的标量 continued-fraction 范畴内，凡达到项目 odd tail 最优局部阶数的算法，全都坍缩到同一 rational shadow；不存在第二条不同的 depth-\(n\) 路径。
\end{corollary}
\begin{proof}
由定理~\ref{thm:xi-time-part9zbk-pade-jacobi-identity}，满足上述 Pad\'e 条件的有理函数唯一。
故 \(R_n\) 只能等于 \(m_{\varphi,n}\)。
证毕。
\end{proof}
\leanverified{paper\_xi\_time\_part9zbk\_continued\_fraction\_rigidity\_collapse}

\begin{theorem}[odd Bernoulli layer 的 pole-count 最优律]\label{thm:xi-time-part9zbk-pole-count-optimality}
\leanpartial{paper\_xi\_time\_part9zbk\_pole\_count\_optimality}{the requested Lean signature proves uniqueness from \(2n\) matched moments, but it omits the separate hypothesis that \texttt{shadow n} itself matches those moments, so the reverse implication is not derivable as stated}
任取正 \(n\)-极点的 Stieltjes 型有理 shadow
\[
R_n(z)=\int\frac{\dd\nu(y)}{1-yz},
\qquad
\nu=\sum_{j=1}^{n}c_j\delta_{x_j},
\qquad
c_j>0,\ x_j>0.
\]
则以下两条严格等价：
\begin{enumerate}
  \item \(R_n\) 与 \(m_\varphi\) 在 \(z=0\) 处的 Taylor 系数匹配到 \(z^{2n-1}\)，即
  \[
  m_\varphi(z)-R_n(z)=O(z^{2n});
  \]
  \item \(R_n=m_{\varphi,n}\)。
\end{enumerate}
并且，不存在任何这样的 \(R_n\) 使匹配进一步提升到 \(z^{2n}\)。

等价地，在 H.111 的 factorial lift
\[
\mathcal E_\varphi(s)=\sum_{k\ge 0}(-1)^k(2k)!\,u_k\,s^{2k+1}
\]
下，正 \(n\)-极点标量 shadow 所能达到的最大 exact odd-layer 数恰为 \(2n\)，即最高精确项为
\[
s,\ s^3,\ \dots,\ s^{4n-1},
\]
而 \(s^{4n+1}\) 必然失配。
该上界由 principal compression 唯一达到。
\end{theorem}
\begin{proof}
“\(1\Rightarrow 2\)” 正是定理~\ref{thm:xi-time-part9zbk-pade-jacobi-identity}。
“\(2\Rightarrow 1\)” 亦显然，因为 \(m_{\varphi,n}\) 正是通过前 \(2n\) 个矩构造出来的。

若还能匹配到 \(z^{2n}\)，则
\[
\int y^k\,\dd\nu(y)=u_k
\qquad (0\le k\le 2n),
\]
这与定理~\ref{thm:xi-time-part9zbk-positive-natomic-exactness-ceiling} 矛盾。
故不可能进一步提升。

把这一结论转写到 factorial lift，只需注意 \(z^k\) 与 odd tail 中的 \(s^{2k+1}\) 一一对应。
因此最高可精确到 \(k=2n-1\)，亦即
\[
s^{4n-1},
\]
而对应 \(k=2n\) 的
\[
s^{4n+1}
\]
必然失配。
唯一达到该上界者即 \(m_{\varphi,n}\)，亦即 principal compression 的 Gaussian--Pad\'e shadow。
证毕。
\leanverified{paper\_xi\_time\_part9zbk\_pole\_count\_optimality}
\end{proof}

\noindent
因此，H.111 odd tail 与 classical rational-approximation machinery 的真正接缝并不在于再造另一条公式，而在于一条更细的最优性原则：给定极点预算 \(n\)，项目 odd tail 所允许的正标量有理 shadow 不是很多，而是唯一的 Gaussian--Pad\'e shadow；其可达阶数也不是模糊增长，而恰好由 \(2n\)-moment exactness ceiling 锁定。由此，fold-\(\pi\) compact Jacobi 宿主除了 infinite-rank realization、principal compression 与 wave-trace 正性之外，又额外携带了一条 Gaussian exactness、Pad\'e rigidity 与 strict one-sided pre-singular approximation 同步闭合的有限极点理论。

\endinput

\paragraph{fold-\texorpdfstring{$\pi$}{pi} Jacobi host 的算子级黄金分裂、双 cotangent Weyl--Stieltjes 核与双尺度模变换热迹}
沿用定理~\ref{thm:xi-time-part9zbe-foldpi-canonical-compact-jacobi}、
\ref{thm:xi-time-part9zbe-foldpi-jacobi-schatten-zeta}、
\ref{thm:xi-time-part9zbe-foldpi-fredholm-double-sine}
与定理~\ref{thm:xi-time-part9zbb-foldpi-borel-double-lattice} 的记号。
前几节已经确定 \(J_\varphi\) 的显式谱、operator \(\zeta\) 与 Fredholm 行列式。
在此基础上可进一步确定其算子级直和结构、谱测度的 Stieltjes 变换闭式以及 inverse host 的 Jacobi--theta 型热迹。
由此产生的并非原有结果的改写，而是一条将 compact Jacobi host 的谱函数论推到 Weyl--Stieltjes 核与 modular 对偶层面的闭合链。

\begin{theorem}[fold-\texorpdfstring{$\pi$}{pi} Jacobi host 的算子级黄金分裂]\label{thm:xi-time-part9zbl-foldpi-golden-operator-splitting}
\leanverified{paper\_xi\_time\_part9zbl\_foldpi\_golden\_operator\_splitting}
令
\[
S:=\operatorname{diag}\!\left(\frac{1}{4\pi^2n^2}\right)_{n\ge 1}
\]
作用于 \(\ell^2(\NN)\)。则 \(J_\varphi\) 与 \(S\oplus \varphi^2 S\) 酉等价：
\[
J_\varphi \simeq S\oplus \varphi^2 S.
\]
因而其 operator \(\zeta\)、Fredholm 行列式分别分裂为
\[
\zeta_{J_\varphi}(s)=(1+\varphi^{2s})\,\zeta_S(s),
\qquad
\det(I-zJ_\varphi)=\det(I-zS)\,\det(I-\varphi^2 zS),
\]
其中
\[
\det(I-zS)
=\prod_{n\ge 1}\left(1-\frac{z}{4\pi^2n^2}\right)
=\frac{\sin(\sqrt z/2)}{\sqrt z/2}.
\]
\end{theorem}
\begin{proof}
由定理~\ref{thm:xi-time-part9zbe-foldpi-canonical-compact-jacobi}，
\[
\sigma(J_\varphi)\setminus\{0\}
=
\left\{\frac{1}{4\pi^2n^2}:n\ge 1\right\}
\cup
\left\{\frac{\varphi^2}{4\pi^2n^2}:n\ge 1\right\},
\]
且每个非零谱点为简单特征值。
另一方面，\(S\oplus \varphi^2 S\) 的非零谱连同重数恰为同一集合，因为 \(\varphi^2\notin\QQ\) 排除了两列之间的重合。
紧自伴算子由其非零特征值及重数唯一决定（至多差一个作用于零空间上的酉变换），
故
\[
J_\varphi\simeq S\oplus\varphi^2 S.
\]
其后两式分别由酉等价保持 Schatten 迹与 Fredholm 行列式的直和乘法性立得；最后一式即 Euler sine product。证毕。
\end{proof}

\begin{theorem}[fold-\texorpdfstring{$\pi$}{pi} host 的 Weyl--Stieltjes 函数闭式]\label{thm:xi-time-part9zbl-foldpi-weyl-stieltjes-closed-form}
\leanverified{paper\_xi\_time\_part9zbl\_foldpi\_weyl\_stieltjes\_closed\_form}
对 \(z\in\CC\setminus[0,\infty)\)，取 \(\sqrt z\) 的主值支，定义谱测度 \(\mu_\varphi\) 的 Stieltjes 变换
\[
M_\varphi(z):=\int\frac{\dd\mu_\varphi(y)}{y-z}.
\]
则成立显式恒等式
\[
M_\varphi(z)
=
\frac{\varphi^{-1}}{\sqrt z}\cot\!\left(\frac{1}{2\sqrt z}\right)
+
\frac{\varphi^{-2}}{\sqrt z}\cot\!\left(\frac{\varphi}{2\sqrt z}\right)
-2(\varphi^{-1}+\varphi^{-3}).
\]
因此 \(J_\varphi\) 的 Jacobi 参数所对应的整条 continued-fraction 宿主由一个双 cotangent Herglotz 核完全控制。
\end{theorem}
\begin{proof}
记
\[
a_n=\frac{1}{4\pi^2n^2},
\qquad
b_n=\frac{\varphi^2}{4\pi^2n^2}.
\]
则
\[
M_\varphi(z)
=
\sum_{n\ge 1}\frac{\varphi^{-1}}{\pi^2n^2}
\left(\frac{1}{a_n-z}+\frac{1}{b_n-z}\right)
=
4\varphi^{-1}\sum_{n\ge 1}\frac{1}{1-4\pi^2n^2z}
+
4\varphi^{-3}\sum_{n\ge 1}\frac{1}{1-4\pi^2n^2z/\varphi^2}.
\]
由 cotangent 部分分式展开的经典推论，
\[
\sum_{n\ge 1}\frac{1}{1-a^2n^2}
=
-\frac{1}{2}+\frac{\pi}{2a}\cot\!\left(\frac{\pi}{a}\right).
\]
分别取 \(a=2\pi\sqrt z\) 与 \(a=2\pi\sqrt z/\varphi\) 代入，得到
\[
\sum_{n\ge 1}\frac{1}{1-4\pi^2n^2z}
=
-\frac{1}{2}+\frac{1}{4\sqrt z}\cot\!\left(\frac{1}{2\sqrt z}\right),
\]
\[
\sum_{n\ge 1}\frac{1}{1-4\pi^2n^2z/\varphi^2}
=
-\frac{1}{2}+\frac{\varphi}{4\sqrt z}\cot\!\left(\frac{\varphi}{2\sqrt z}\right).
\]
代回并整理常数项即得所示闭式。证毕。
\end{proof}

\begin{theorem}[inverse host 的双尺度 Jacobi--theta 模变换公式]\label{thm:xi-time-part9zbl-foldpi-inverse-host-theta-heat-trace}
\leanverified{paper\_xi\_time\_part9zbl\_foldpi\_inverse\_host\_theta\_heat\_trace}
在 \(\ker(J_\varphi)^\perp\) 上定义无界正算子 \(L_\varphi:=J_\varphi^{-1}\)。其热迹
\[
\Theta_\varphi(t):=\Tr(e^{-tL_\varphi}),
\qquad t>0,
\]
满足
\[
\Theta_\varphi(t)
=
\sum_{n\ge 1}e^{-4\pi^2n^2 t}
+\sum_{n\ge 1}e^{-4\pi^2n^2 t/\varphi^2}
=
\frac{1}{2}\Bigl(\vartheta_3(0,e^{-4\pi^2t})+\vartheta_3(0,e^{-4\pi^2t/\varphi^2})\Bigr)-1.
\]
经 Jacobi 模变换后得到对偶表示
\[
\Theta_\varphi(t)
=
\frac{1}{4\sqrt{\pi t}}\,\vartheta_3(0,e^{-1/(4t)})
+
\frac{\varphi}{4\sqrt{\pi t}}\,\vartheta_3(0,e^{-\varphi^2/(4t)})
-1,
\]
从而当 \(t\downarrow 0\) 时有小时间展开
\[
\Theta_\varphi(t)
=
\frac{\varphi^2}{4\sqrt{\pi t}}-1+O(e^{-1/(4t)}).
\]
因此 fold-\(\pi\) host 的高能端给出的不仅是一个 Weyl 主项，而是一条二尺度模变换精确式，其中 \(\varphi\) 以热时间对偶的方式同时进入主项与瞬子尾。
\end{theorem}
\begin{proof}
由定理~\ref{thm:xi-time-part9zbe-foldpi-canonical-compact-jacobi}，
\[
\sigma(L_\varphi)
=
\{4\pi^2n^2:n\ge 1\}\cup\{4\pi^2n^2/\varphi^2:n\ge 1\}.
\]
故第一式立即成立。
再用 Jacobi theta 函数的定义
\[
\vartheta_3(0,q)=1+2\sum_{n\ge 1}q^{n^2}
\]
得第二式。
对两个 theta 因子分别应用 Jacobi 模变换（参见 \href{https://dlmf.nist.gov/20.7}{DLMF 20.7}）
\[
\vartheta_3(0,e^{-\pi u})=u^{-1/2}\,\vartheta_3(0,e^{-\pi/u}),
\]
取 \(u=4\pi t\) 得到
\[
\vartheta_3(0,e^{-4\pi^2 t})
=
\frac{1}{2\sqrt{\pi t}}\,\vartheta_3(0,e^{-1/(4t)}),
\]
取 \(u=4\pi t/\varphi^2\) 得到
\[
\vartheta_3(0,e^{-4\pi^2 t/\varphi^2})
=
\frac{\varphi}{2\sqrt{\pi t}}\,\vartheta_3(0,e^{-\varphi^2/(4t)}).
\]
代入第二式即得对偶表示。
当 \(t\downarrow 0\) 时，\(e^{-1/(4t)}\to 0\) 与 \(e^{-\varphi^2/(4t)}\to 0\) 使得两项 \(\vartheta_3\) 在对偶形式中均趋于 \(1\)，
故
\[
\Theta_\varphi(t)
=
\frac{1+\varphi}{4\sqrt{\pi t}}-1+O(e^{-1/(4t)})
=
\frac{\varphi^2}{4\sqrt{\pi t}}-1+O(e^{-1/(4t)}).
\]
证毕。
\end{proof}

\noindent
因此，fold-\(\pi\) Jacobi host 除已有的 Fredholm 双正弦与 Schatten \(\zeta\) 闭式之外，还同时具备一个以 \(S\oplus \varphi^2 S\) 为典范形的算子级黄金分裂、一个由双 cotangent 核驱动的 Weyl--Stieltjes 函数完全显式化，以及一条连接 Jacobi theta 与 modular 对偶的双尺度热迹公式。这三层结构为后续 Hadamard 增长几何与零点斜率刚性分析提供了算子基础。

\endinput

\paragraph{Fredholm 行列式的 Hadamard 增长几何、零点斜率 \texorpdfstring{$\varphi^4$}{phi4}-刚性与 resurgent--wave Fourier 对偶}
沿用定理~\ref{thm:xi-time-part9zbe-foldpi-fredholm-double-sine}
与定理~\ref{thm:xi-time-part9zbl-foldpi-golden-operator-splitting} 的记号。
Fredholm 行列式 \(D_\varphi(z)\) 已被确定为阶 \(1/2\) 的 entire function，其零点集由两条 reciprocal-square 梯构成。
在此基础上可进一步刻画 \(D_\varphi\) 的 Phragm\'en--Lindel\"of 指示函数、两族零点上的导数逐点倍比律以及归一化斜率的概率极限律，
并将 Borel 奇点格与波迹谱之间的联系推至精确 Fourier 对偶。

\begin{theorem}[黄金双正弦 Fredholm 行列式的 Phragm\'en--Lindel\"of 指示函数]\label{thm:xi-time-part9zbm-foldpi-hadamard-indicator}
\leanverified{paper\_xi\_time\_part9zbm\_foldpi\_hadamard\_indicator}
Fredholm 行列式
\[
D_\varphi(z)
=
\frac{\sin(\sqrt z/2)}{\sqrt z/2}\,
\frac{\sin(\varphi\sqrt z/2)}{\varphi\sqrt z/2}
\]
是阶 \(\rho=1/2\) 的 entire function，其指示函数为
\[
h_{D_\varphi}(\theta)
=
\limsup_{r\to\infty}r^{-1/2}\log|D_\varphi(re^{i\theta})|
=
\frac{\varphi^2}{2}\,|\sin(\theta/2)|.
\]
其 type 为
\[
\tau(D_\varphi)=\max_\theta h_{D_\varphi}(\theta)=\frac{\varphi^2}{2},
\]
并且沿负实轴有精确渐近
\[
D_\varphi(-x)\sim \frac{e^{\varphi^2\sqrt x/2}}{\varphi\, x},
\qquad x\to+\infty.
\]
黄金参数不仅控制零点位置，也控制整个 Fredholm 行列式的 Hadamard 增长几何。
\end{theorem}
\begin{proof}
令 \(z=re^{i\theta}\)，\(w=\sqrt z/2=(r^{1/2}/2)e^{i\theta/2}\)。
由 \(\sin u=(e^{iu}-e^{-iu})/(2i)\) 可得
\[
\log|\sin u|=|\Im u|+O(1)\qquad (|u|\to\infty),
\]
从而
\[
\log\left|\frac{\sin w}{w}\right|
=
\frac{r^{1/2}}{2}|\sin(\theta/2)|+O(\log r),
\]
\[
\log\left|\frac{\sin(\varphi w)}{\varphi w}\right|
=
\frac{\varphi\, r^{1/2}}{2}|\sin(\theta/2)|+O(\log r).
\]
相加即得
\[
r^{-1/2}\log|D_\varphi(re^{i\theta})|
\to
\frac{1+\varphi}{2}|\sin(\theta/2)|
=
\frac{\varphi^2}{2}|\sin(\theta/2)|.
\]
沿负实轴 \(\sqrt{-x}=i\sqrt x\)，故
\[
D_\varphi(-x)
=
\frac{\sinh(\sqrt x/2)}{\sqrt x/2}\,
\frac{\sinh(\varphi\sqrt x/2)}{\varphi\sqrt x/2}
\sim
\frac{e^{\sqrt x/2}}{\sqrt x}\,
\frac{e^{\varphi\sqrt x/2}}{\varphi\sqrt x}
=
\frac{e^{\varphi^2\sqrt x/2}}{\varphi\, x}.
\]
证毕。
\end{proof}

\begin{theorem}[双零点梯的逐点斜率刚性]\label{thm:xi-time-part9zbm-foldpi-pointwise-slope-rigidity}
\leanverified{paper\_xi\_time\_part9zbm\_foldpi\_pointwise\_slope\_rigidity}
记 \(D_\varphi\) 的两族零点为
\[
z_n^-:=\frac{4\pi^2n^2}{\varphi^2},
\qquad
z_n^+:=4\pi^2n^2
\qquad (n\ge 1).
\]
则有精确斜率公式
\[
D_\varphi'(z_n^+)
=
\frac{\sin(\pi n/\varphi)}{8\pi^3\varphi\,n^3},
\]
\[
D_\varphi'(z_n^-)
=
\frac{\varphi^3\sin(\pi n\varphi)}{8\pi^3 n^3}
=
(-1)^n\varphi^4 D_\varphi'(z_n^+).
\]
特别地，
\[
|D_\varphi'(z_n^-)|=\varphi^4 |D_\varphi'(z_n^+)|
\qquad (n\ge 1).
\]
两条零点梯之间的局部横截率满足逐点 \(\varphi^4\)-刚性倍比律。
\end{theorem}
\begin{proof}
记
\[
s(z):=\frac{\sin(\sqrt z/2)}{\sqrt z/2},
\qquad
D_\varphi(z)=s(z)\,s(\varphi^2 z).
\]
令 \(u=\sqrt z/2\)，则
\[
s'(z)=\frac{u\cos u-\sin u}{8u^3}.
\]
故
\[
s'(4\pi^2n^2)=\frac{(-1)^n}{8\pi^2n^2}.
\]

在整数支 \(z_n^+\) 处 \(s(z_n^+)=0\)，因此
\[
D_\varphi'(z_n^+)
=
s'(z_n^+)\,s(\varphi^2 z_n^+)
=
\frac{(-1)^n}{8\pi^2n^2}\,\frac{\sin(\pi n\varphi)}{\pi n\varphi}.
\]
利用 \(\varphi=1+\varphi^{-1}\) 可得
\[
\sin(\pi n\varphi)=\sin(\pi n+\pi n/\varphi)=(-1)^n\sin(\pi n/\varphi),
\]
代入即得第一式。

在黄金支 \(z_n^-\) 处 \(s(\varphi^2 z_n^-)=0\)，因此
\[
D_\varphi'(z_n^-)
=
s(z_n^-)\,\varphi^2 s'(\varphi^2 z_n^-)
=
\frac{\sin(\pi n/\varphi)}{\pi n/\varphi}\,\varphi^2\,\frac{(-1)^n}{8\pi^2n^2}
=
\frac{(-1)^n\varphi^3\sin(\pi n/\varphi)}{8\pi^3n^3}.
\]
再用 \(\sin(\pi n/\varphi)=(-1)^n\sin(\pi n\varphi)\) 便得第二式。
两式之比为 \((-1)^n\varphi^4\)，故绝对值满足逐点 \(\varphi^4\)-倍比。证毕。
\end{proof}

\begin{corollary}[归一化节点斜率的 arcsine 极限律]\label{cor:xi-time-part9zbm-foldpi-arcsine-slope-law}
\leanverified{paper\_xi\_time\_part9zbm\_foldpi\_arcsine\_slope\_law}
定义归一化斜率
\[
s_n^+ := 8\pi^3\varphi\,n^3 |D_\varphi'(z_n^+)|,
\qquad
s_n^-:= \frac{8\pi^3}{\varphi^3}\,n^3 |D_\varphi'(z_n^-)|.
\]
则
\[
s_n^+=s_n^-= |\sin(\pi n/\varphi)|.
\]
当 \(N\to\infty\) 时，两列的经验分布同时收敛到 \([0,1]\) 上的 arcsine 律
\[
\dd\mu_{\mathrm{arc}}(u)
=
\frac{2}{\pi\sqrt{1-u^2}}\,\mathbf{1}_{(0,1)}(u)\,\dd u.
\]
等价地，对每个连续函数 \(f\in C([0,1])\) 有
\[
\frac{1}{N}\sum_{n=1}^N f(s_n^\pm)
\longrightarrow
\int_0^1 f(u)\,\frac{2\,\dd u}{\pi\sqrt{1-u^2}}.
\]
特别地，对任意 \(m>-1\) 有矩公式
\[
\lim_{N\to\infty}\frac{1}{N}\sum_{n=1}^N (s_n^\pm)^m
=
\frac{\Gamma\!\bigl(\frac{m+1}{2}\bigr)}
{\sqrt\pi\,\Gamma\!\bigl(\frac{m+2}{2}\bigr)},
\]
且几何均值满足
\[
\lim_{N\to\infty}\exp\!\left(\frac{1}{N}\sum_{n=1}^N \log s_n^\pm\right)
=
\frac{1}{2}.
\]
双正弦 Fredholm 行列式的零点斜率遵循一条由 irrational rotation 强制出来的 arcsine 概率律。
\end{corollary}
\begin{proof}
由定理~\ref{thm:xi-time-part9zbm-foldpi-pointwise-slope-rigidity} 立得
\[
s_n^\pm = |\sin(\pi n/\varphi)|.
\]
因 \(1/\varphi\notin\QQ\)，Weyl 等分布定理保证序列 \(\{n/\varphi\}\) 在 \([0,1)\) 上等分布。
故对任意连续 \(g\in C([0,1])\) 有
\[
\frac{1}{N}\sum_{n=1}^N g(\{n/\varphi\}) \to \int_0^1 g(x)\,\dd x.
\]
取 \(g(x)=f(|\sin\pi x|)\)，便得
\[
\frac{1}{N}\sum_{n=1}^N f(s_n^\pm) \to \int_0^1 f(|\sin\pi x|)\,\dd x.
\]
在 \([0,1/2]\) 上令 \(u=\sin(\pi x)\) 换元并利用 \([0,1]\) 上的对称性，即得推前密度
\[
\frac{2}{\pi\sqrt{1-u^2}}\,\dd u.
\]
矩公式由 Beta 积分
\[
\int_0^1 u^m\,\frac{2\,\dd u}{\pi\sqrt{1-u^2}}
=
\frac{\Gamma\bigl(\frac{m+1}{2}\bigr)}
{\sqrt\pi\,\Gamma\bigl(\frac{m+2}{2}\bigr)}
\]
直接给出。对数平均则由经典积分
\[
\int_0^1 \log(\sin\pi x)\,\dd x=-\log 2
\]
得出。证毕。
\end{proof}

\begin{theorem}[Borel 奇点格与波迹谱的精确 Fourier 对偶]\label{thm:xi-time-part9zbm-foldpi-borel-wave-fourier-duality}
记
\[
\Sigma_\varphi
:=
2\pi\ZZ_{\neq 0}\cup \frac{2\pi}{\varphi}\ZZ_{\neq 0}.
\]
则定理~\ref{thm:xi-time-part9zbb-foldpi-borel-double-lattice} 中的 Borel 奇点集合满足
\[
\Omega_\varphi=i\,\Sigma_\varphi.
\]
对命题~\ref{prop:xi-time-part9zbj-pure-point-poisson-jacobi-wave-trace} 中的波迹
\[
\mathcal W_\varphi(x)
=
2\sum_{\lambda\in\sigma(J_\varphi)\setminus\{0\}}
\cos\!\left(\frac{x}{\sqrt\lambda}\right)
\]
取 Fourier 变换（约定 \(\widehat f(\xi)=\int_{\RR}e^{-ix\xi}f(x)\,\dd x\)），成立分布恒等式
\[
\widehat{\mathcal W_\varphi}(\xi)
=
2\pi\sum_{\omega\in\Sigma_\varphi}\delta(\xi-\omega).
\]
特别地，
\[
\operatorname{supp}\widehat{\mathcal W_\varphi}=\Sigma_\varphi,
\qquad
\Omega_\varphi=i\,\operatorname{supp}\widehat{\mathcal W_\varphi}.
\]
fold-\(\pi\) 奇层的 resurgent singular support 与 compact Jacobi host 的 spectral comb 之间存在一条精确的 Fourier 对偶律：前者是后者的纯虚频率化，后者是前者的实频投影。
\end{theorem}
\begin{proof}
由命题~\ref{prop:xi-time-part9zbj-pure-point-poisson-jacobi-wave-trace} 已经建立的分布恒等式，
\[
\mathcal W_\varphi(x)
=
\sum_{m\in\ZZ}\delta(x-m)
+
\varphi\sum_{m\in\ZZ}\delta(x-m\varphi)
-2.
\]
对第一个 Dirac comb 取 Fourier 变换，Poisson 求和公式给出
\[
\mathcal F\!\left[\sum_{m\in\ZZ}\delta(x-m)\right](\xi)
=
2\pi\sum_{k\in\ZZ}\delta(\xi-2\pi k).
\]
对第二个 Dirac comb 同理，
\[
\mathcal F\!\left[\varphi\sum_{m\in\ZZ}\delta(x-m\varphi)\right](\xi)
=
2\pi\sum_{k\in\ZZ}\delta(\xi-2\pi k/\varphi).
\]
常数项 \(-2\) 的 Fourier 变换为 \(-4\pi\delta(\xi)\)。
将三项相加，\(k=0\) 处的两项 \(\delta(\xi)\) 之和恰为 \(4\pi\delta(\xi)\)，与常数项的贡献精确抵消，于是
\[
\widehat{\mathcal W_\varphi}(\xi)
=
2\pi\sum_{k\neq 0}\delta(\xi-2\pi k)
+
2\pi\sum_{k\neq 0}\delta(\xi-2\pi k/\varphi)
=
2\pi\sum_{\omega\in\Sigma_\varphi}\delta(\xi-\omega).
\]
再由定理~\ref{thm:xi-time-part9zbb-foldpi-borel-double-lattice}，\(\widehat{\mathcal E}_\varphi\) 的奇点集合正是
\(2\pi i\,\ZZ_{\neq 0}\cup (2\pi i/\varphi)\,\ZZ_{\neq 0}=i\,\Sigma_\varphi\)。
因此
\[
\Omega_\varphi=i\,\operatorname{supp}\widehat{\mathcal W_\varphi}.
\]
证毕。
\end{proof}

\noindent
综上，fold-\(\pi\) Jacobi host 在现有闭式框架下显现出一条完整的结构链：
\[
\text{黄金双尺度 Borel 奇点}
\;\Longleftrightarrow\;
\text{cotangent 型 Weyl--Stieltjes 核}
\;\Longleftrightarrow\;
\text{双尺度模变换热迹}
\;\Longleftrightarrow\;
\text{Hadamard 增长指标}
\;\Longleftrightarrow\;
\text{arcsine 斜率统计}
\;\Longleftrightarrow\;
\text{Fourier 纯点波迹}.
\]
黄金参数 \(\varphi\) 从 Borel 奇点几何（双晶格）经由谱函数论（Weyl--Stieltjes 核、Jacobi--theta 热迹）穿透到整函数增长（Hadamard 指标 \(\varphi^2/2\)）与零点横截（\(\varphi^4\)-刚性），最终由等分布强制出现于概率极限律（arcsine），并通过 Fourier 对偶封闭回到 Borel 层面。
这给出的不是同一对象的多种表述，而是一个闭合的 resurgence--spectral--probabilistic package。

\endinput

\paragraph{window-\texorpdfstring{$6$}{6} 微态偏置纤维律的 Cauchy--Stieltjes 闭式、Cartan--Lanczos 共轭与 Toda 轨道}
在 window-\(6\) 的三档纤维直方图、\(B_3/C_3\) 根字典、projective long sector 的纯 Cartan 晶场分解，以及极小紧 Lie 完成被唯一锁定到 \(\mathrm{Spin}(7)\) 之后，还可把同一离散对象进一步压缩为一条完全闭合的谱测度链：先由微态均匀分布推出纤维大小的偏置概率律，再由该三原子测度生成唯一的 Jacobi--Lanczos Hamiltonian，继而把它识别为同一 \(3\)-能级 Cartan 可观测量在特定量子态下的 Born 分布，最后再沿有限非周期 Toda 流把这一 Hamiltonian 放入同一 regular semisimple 轨道的等谱异宿链。由此，window-\(6\) 的直方图、Cartan 能级、Lanczos 链、正交多项式与 Toda 轨道不再是彼此并列的读法，而是同一谱测度对象的五种等价坐标。

\begin{theorem}[微态偏置纤维律的 Cauchy--Stieltjes/Herglotz 闭式]\label{thm:xi-time-part9zc-window6-microstate-biased-fiberlaw-herglotz-closedform}
沿用推论 \ref{cor:fold-bin-m6-fiber-hist} 的记号，定义 window-\(6\) 的微态偏置纤维律
\[
\nu_6
:=
\frac1{64}\sum_{w\in X_6}d_6(w)\,\delta_{d_6(w)}.
\]
则
\[
\nu_6=\frac14\,\delta_2+\frac3{16}\,\delta_3+\frac9{16}\,\delta_4.
\]
定义其 Cauchy--Stieltjes 变换
\[
G_6(z):=\int_{\RR}\frac{\dd\nu_6(t)}{z-t},
\qquad
M_6(z):=-G_6(z)=\int_{\RR}\frac{\dd\nu_6(t)}{t-z}.
\]
则 \(M_6\) 在上半平面上是 Herglotz 函数，而 \(G_6\) 具有精确闭式
\[
G_6(z)=\frac{16z^2-91z+126}{16(z-2)(z-3)(z-4)}.
\]
\leanverified{paper\_xi\_time\_part9zc\_window6\_microstate\_biased\_fiberlaw\_herglotz\_closedform}
等价地，若
\[
m_n:=\int_{\RR} t^n\,\dd\nu_6(t)
\qquad (n\ge 0),
\]
则
\[
m_n=\frac{8\cdot 2^{n+1}+4\cdot 3^{n+1}+9\cdot 4^{n+1}}{64},
\]
并且在 \(z=\infty\) 邻域内有 Laurent 展开
\[
G_6(z)=\sum_{n\ge 0}\frac{m_n}{z^{n+1}}.
\]
\end{theorem}
\begin{proof}
由推论 \ref{cor:fold-bin-m6-fiber-hist}，window-\(6\) 的纤维大小只取 \(2,3,4\) 三个值，且相应纤维个数分别为 \(8,4,9\)。总微态数为
\[
8\cdot 2+4\cdot 3+9\cdot 4=64,
\]
故微态均匀分布向纤维大小层的推前正是
\[
\nu_6(\{2\})=\frac{8\cdot 2}{64}=\frac14,\qquad
\nu_6(\{3\})=\frac{4\cdot 3}{64}=\frac3{16},\qquad
\nu_6(\{4\})=\frac{9\cdot 4}{64}=\frac9{16}.
\]
于是
\[
G_6(z)=\frac{1/4}{z-2}+\frac{3/16}{z-3}+\frac{9/16}{z-4}.
\]
将三项通分，得到
\[
G_6(z)=\frac{4(z-3)(z-4)+3(z-2)(z-4)+9(z-2)(z-3)}{16(z-2)(z-3)(z-4)}
=\frac{16z^2-91z+126}{16(z-2)(z-3)(z-4)}.
\]

另一方面，对任意 \(n\ge 0\)，
\[
m_n=\frac14\,2^n+\frac3{16}\,3^n+\frac9{16}\,4^n
=\frac{8\cdot 2^{n+1}+4\cdot 3^{n+1}+9\cdot 4^{n+1}}{64}.
\]
再由几何级数展开
\[
\frac1{z-t}=\sum_{n\ge 0}\frac{t^n}{z^{n+1}}
\qquad (|z|>|t|),
\]
逐项积分即得 \(G_6\) 在 \(z=\infty\) 邻域内的 Laurent 展开。最后，\(\Im z>0\) 时
\[
\Im M_6(z)=\int_{\RR}\frac{\Im z}{|t-z|^2}\,\dd\nu_6(t)>0,
\]
故 \(M_6\) 为 Herglotz 函数。证毕。
\end{proof}

\begin{theorem}[window-\texorpdfstring{$6$}{6} 微态偏置律的正交多项式闭合与三次终止]\label{thm:xi-time-part9zc-window6-orthogonal-polynomials-cubic-termination}
\leanverified{paper\_xi\_time\_part9zc\_window6\_orthogonal\_polynomials\_cubic\_termination}
以 \(\nu_6\) 为内积
\[
\langle f,g\rangle_{\nu_6}:=\int_{\RR}f(t)g(t)\,\dd\nu_6(t).
\]
令 \(P_0,P_1,P_2\) 为 \(\nu_6\) 的 monic 正交多项式，并令 \(P_3\) 为唯一与所有次数至多 \(2\) 的多项式正交的 monic 三次多项式。则
\[
P_0(x)=1,
\]
\[
P_1(x)=x-\frac{53}{16},
\]
\[
P_2(x)=x^2-\frac{371}{61}x+\frac{516}{61},
\]
\[
P_3(x)=(x-2)(x-3)(x-4).
\]
其非退化平方范数为
\[
h_0:=\langle P_0,P_0\rangle_{\nu_6}=1,\qquad
h_1:=\langle P_1,P_1\rangle_{\nu_6}=\frac{183}{256},
\qquad
h_2:=\langle P_2,P_2\rangle_{\nu_6}=\frac9{61},
\]
而
\[
\langle P_3,P_3\rangle_{\nu_6}=0.
\]
此外，\(P_2\) 的两个零点为
\[
x_{\pm}=\frac{371\pm 11\sqrt{97}}{122},
\]
并满足严格交错
\[
2<x_-<3<x_+<4.
\]
\end{theorem}
\begin{proof}
由定理 \ref{thm:xi-time-part9zc-window6-microstate-biased-fiberlaw-herglotz-closedform}，
\[
m_0=1,\qquad
m_1=\frac{53}{16},\qquad
m_2=\frac{187}{16},\qquad
m_3=\frac{689}{16}.
\]
故
\[
P_1(x)=x-m_1=x-\frac{53}{16},
\]
并且
\[
h_1=m_2-m_1^2=\frac{187}{16}-\left(\frac{53}{16}\right)^2=\frac{183}{256}.
\]

设
\[
P_2(x)=x^2+ux+v.
\]
由 \(\langle P_2,1\rangle_{\nu_6}=0\) 与 \(\langle P_2,x\rangle_{\nu_6}=0\)，得到线性方程组
\[
m_2+u m_1+v=0,
\qquad
m_3+u m_2+v m_1=0.
\]
代入上面的矩并求解，得到
\[
u=-\frac{371}{61},
\qquad
v=\frac{516}{61}.
\]
故
\[
P_2(x)=x^2-\frac{371}{61}x+\frac{516}{61}.
\]
再直接在支撑点 \(2,3,4\) 上代入，得到
\[
P_2(2)=\frac{18}{61},\qquad
P_2(3)=-\frac{48}{61},\qquad
P_2(4)=\frac8{61}.
\]
于是
\[
h_2
=
\frac14\left(\frac{18}{61}\right)^2
+
\frac3{16}\left(\frac{48}{61}\right)^2
+
\frac9{16}\left(\frac8{61}\right)^2
=\frac9{61}.
\]

由于 \(\operatorname{supp}\nu_6=\{2,3,4\}\) 恰含三点，唯一的 monic 三次多项式
\[
P_3(x):=(x-2)(x-3)(x-4)
\]
在支撑上处处为零，故自动与所有次数至多 \(2\) 的多项式正交，且其平方范数为零。

最后，对方程
\[
61x^2-371x+516=0
\]
求根即可得到
\[
x_{\pm}=\frac{371\pm 11\sqrt{97}}{122}.
\]
又因
\[
P_2(2)>0,\qquad P_2(3)<0,\qquad P_2(4)>0,
\]
故两根必分别落在区间 \((2,3)\) 与 \((3,4)\) 内。证毕。
\end{proof}

\begin{theorem}[window-\texorpdfstring{$6$}{6} 的规范信息 Hamiltonian]\label{thm:xi-time-part9zc-window6-canonical-information-hamiltonian}
定义三对角矩阵
\[
J_6:=
\begin{pmatrix}
\frac{53}{16} & \frac{\sqrt{183}}{16} & 0\\[1mm]
\frac{\sqrt{183}}{16} & \frac{2703}{976} & \frac{16\sqrt3}{61}\\[1mm]
0 & \frac{16\sqrt3}{61} & \frac{178}{61}
\end{pmatrix}.
\]
则：
\begin{enumerate}
  \item \(\chi_{J_6}(\lambda)=\det(\lambda I_3-J_6)=(\lambda-2)(\lambda-3)(\lambda-4)\)；
  \item \(\nu_6\) 正是 \(J_6\) 在循环向量 \(e_1=(1,0,0)^{\mathsf T}\) 下的谱测度；亦即对任意 Borel 集 \(B\subset\RR\),
  \[
  \nu_6(B)=\left\langle e_1,\mathbf 1_B(J_6)e_1\right\rangle;
  \]
  \item 等价地，对一切 \(n\ge 0\),
  \[
  \int_{\RR} t^n\,\dd\nu_6(t)=\left\langle e_1,J_6^n e_1\right\rangle.
  \]
\end{enumerate}
因此，window-\(6\) 的全部微态偏置信息可被唯一压缩为一个 \(3\)-站点 Jacobi Hamiltonian，而不损失任何矩数据。
\end{theorem}
\begin{proof}
由定理 \ref{thm:xi-time-part9zc-window6-orthogonal-polynomials-cubic-termination}，
\[
P_1(x)=\left(x-\frac{53}{16}\right)P_0(x),
\]
\[
P_2(x)=\left(x-\frac{2703}{976}\right)P_1(x)-\frac{183}{256}P_0(x),
\]
\[
P_3(x)=\left(x-\frac{178}{61}\right)P_2(x)-\frac{768}{3721}P_1(x).
\]
这三式分别把三项递推系数读出为
\[
a_0=\frac{53}{16},\qquad
\beta_1=\frac{183}{256},\qquad
a_1=\frac{2703}{976},\qquad
\beta_2=\frac{768}{3721},\qquad
a_2=\frac{178}{61}.
\]
于是对应的 Jacobi 矩阵恰为
\[
J_6=
\begin{pmatrix}
a_0 & \sqrt{\beta_1} & 0\\
\sqrt{\beta_1} & a_1 & \sqrt{\beta_2}\\
0 & \sqrt{\beta_2} & a_2
\end{pmatrix},
\]
即题设矩阵。

对有限支撑正测度的 Favard 理论表明，在正交归一基下，“乘以 \(x\)” 的算子由该 Jacobi 矩阵表示，因此其特征多项式正是递推终止多项式 \(P_3\)。由定理 \ref{thm:xi-time-part9zc-window6-orthogonal-polynomials-cubic-termination}，
\[
P_3(\lambda)=(\lambda-2)(\lambda-3)(\lambda-4),
\]
从而第一条成立。

同一谱定理同时给出
\[
\left\langle e_1,J_6^n e_1\right\rangle
=
\int_{\RR}t^n\,\dd\nu_6(t)
\qquad (n\ge 0),
\]
于是 \(\nu_6\) 正是 \(J_6\) 在 \(e_1\) 下的谱测度。证毕。
\end{proof}

\begin{theorem}[spin-\texorpdfstring{$1$}{1} Cartan 可观测量的 Born 实现]\label{thm:xi-time-part9zc-window6-spin1-cartan-born-realization}
沿用定理 \ref{thm:conclusion-window6-projective-long-sector-spin1-crystal-field} 在 \(q=1\) 时的记号，定义
\[
C_6:=3I_3+J_z=\operatorname{diag}(2,3,4),
\qquad
J_z=\operatorname{diag}(-1,0,1).
\]
再定义单位向量
\[
\psi_6:=
\begin{pmatrix}
\frac12\\[1mm]
\frac{\sqrt3}{4}\\[1mm]
\frac34
\end{pmatrix},
\qquad
\|\psi_6\|=1.
\]
则 \(\nu_6\) 恰为 \(C_6\) 在态 \(\psi_6\) 下的谱律：
\[
\nu_6(B)=\left\langle \psi_6,\mathbf 1_B(C_6)\psi_6\right\rangle
\qquad (B\subset\RR\ \text{Borel}),
\]
亦即
\[
\int_{\RR} t^n\,\dd\nu_6(t)=\left\langle \psi_6,C_6^n\psi_6\right\rangle
\qquad (n\ge 0).
\]
因此，window-\(6\) 的微态偏置纤维律严格等于一个 spin-\(1\) 系统对 Cartan 可观测量 \(3I_3+J_z\) 的 Born 测量分布。
\end{theorem}
\begin{proof}
由于 \(C_6\) 已在标准基下对角化，其三个谱投影分别投到三条坐标轴。因此
\[
\left\langle \psi_6,\mathbf 1_{\{2\}}(C_6)\psi_6\right\rangle
=
\left|\frac12\right|^2
=
\frac14,
\]
\[
\left\langle \psi_6,\mathbf 1_{\{3\}}(C_6)\psi_6\right\rangle
=
\left|\frac{\sqrt3}{4}\right|^2
=
\frac3{16},
\]
\[
\left\langle \psi_6,\mathbf 1_{\{4\}}(C_6)\psi_6\right\rangle
=
\left|\frac34\right|^2
=
\frac9{16}.
\]
这恰与定理 \ref{thm:xi-time-part9zc-window6-microstate-biased-fiberlaw-herglotz-closedform} 中的 \(\nu_6\) 一致。矩恒等式再由谱定理立即推出。证毕。
\end{proof}

\begin{theorem}[Cartan--Lanczos 共轭定理]\label{thm:xi-time-part9zc-window6-cartan-lanczos-conjugacy}
\leanverified{paper\_xi\_time\_part9zc\_window6\_cartan\_lanczos\_conjugacy}
存在唯一正交矩阵 \(U\in SO(3)\)，其第一行等于 \(\psi_6^{\mathsf T}\)，并满足
\[
J_6=U\,C_6\,U^{\mathsf T}.
\]
更具体地，可取
\[
U=
\begin{pmatrix}
\frac12 & \frac{\sqrt3}{4} & \frac34\\[1mm]
-\frac{7\sqrt{183}}{122} & -\frac{5\sqrt{61}}{244} & \frac{11\sqrt{183}}{244}\\[1mm]
\frac{3\sqrt{61}}{61} & -\frac{4\sqrt{183}}{61} & \frac{2\sqrt{61}}{61}
\end{pmatrix}\in SO(3),
\]
从而
\[
J_6=U\,\operatorname{diag}(2,3,4)\,U^{\mathsf T}.
\]
因此，定理 \ref{thm:conclusion-window6-projective-long-sector-spin1-crystal-field} 中的 Cartan 势与定理 \ref{thm:xi-time-part9zc-window6-canonical-information-hamiltonian} 中的 Jacobi Hamiltonian 并非两件不同对象，而是同一 regular semisimple 轨道在 Cartan 基与 Krylov--Lanczos 基中的两种坐标表达。
\end{theorem}
\begin{proof}
令
\[
\phi_0(x)=1,
\qquad
\phi_1(x)=\frac{P_1(x)}{\sqrt{h_1}},
\qquad
\phi_2(x)=\frac{P_2(x)}{\sqrt{h_2}},
\]
其中 \(P_k,h_k\) 见定理 \ref{thm:xi-time-part9zc-window6-orthogonal-polynomials-cubic-termination}。再记
\[
(\lambda_1,\lambda_2,\lambda_3)=(2,3,4),
\qquad
(w_1,w_2,w_3)=\left(\frac14,\frac3{16},\frac9{16}\right).
\]
定义
\[
U_{ij}:=\sqrt{w_j}\,\phi_{i-1}(\lambda_j)
\qquad (1\le i,j\le 3).
\]
由于 \((\phi_0,\phi_1,\phi_2)\) 在 \(\nu_6\) 下正交归一，故 \(U\) 为正交矩阵。其第一行为
\[
(\sqrt{w_1},\sqrt{w_2},\sqrt{w_3})
=
\left(\frac12,\frac{\sqrt3}{4},\frac34\right)
=
\psi_6^{\mathsf T}.
\]

另一方面，在正交多项式基下，“乘以 \(x\)” 的矩阵表示正是 \(J_6\)，而在谱基下则是 \(\operatorname{diag}(2,3,4)\)。因此
\[
J_6=U\,\operatorname{diag}(2,3,4)\,U^{\mathsf T}.
\]
将 \(\phi_1,\phi_2\) 在三点 \(2,3,4\) 上的取值显式代入，便得到题设矩阵。再直接计算可得 \(\det U=1\)，故 \(U\in SO(3)\)。

最后，由于 \(J_6\) 的谱简单，且各列是按特征值 \(2,3,4\) 排列并要求第一行全部取正，因此每一列的符号都已被唯一固定，从而满足题设条件的正交矩阵只能有这一组。证毕。
\end{proof}

\begin{corollary}[三层自能的精确 \texorpdfstring{$J$}{J}-fraction]\label{cor:xi-time-part9zc-window6-three-layer-selfenergy-jfraction}
定理 \ref{thm:xi-time-part9zc-window6-microstate-biased-fiberlaw-herglotz-closedform} 中的 Cauchy--Stieltjes 变换 \(G_6\) 具有精确 Jacobi continued fraction
\[
G_6(z)
=
\frac{1}{
z-\frac{53}{16}
-\cfrac{\frac{183}{256}}{
z-\frac{2703}{976}
-\cfrac{\frac{768}{3721}}{
z-\frac{178}{61}
}}}.
\]
等价地，若定义
\[
\Sigma_2(z):=\frac{\frac{768}{3721}}{z-\frac{178}{61}},
\qquad
\Sigma_1(z):=\frac{\frac{183}{256}}{z-\frac{2703}{976}-\Sigma_2(z)},
\]
则
\[
G_6(z)=\frac{1}{z-\frac{53}{16}-\Sigma_1(z)}.
\]
因此，window-\(6\) 的微态偏置纤维律可被解释为一个严格三层终止的自能递归。
\end{corollary}
\begin{proof}
由定理 \ref{thm:xi-time-part9zc-window6-canonical-information-hamiltonian},
\[
G_6(z)=\left\langle e_1,(zI_3-J_6)^{-1}e_1\right\rangle.
\]
对三对角矩阵 \(zI_3-J_6\) 作两次 Schur 补，得到
\[
G_6(z)
=
\frac{1}{
z-a_0-\cfrac{\beta_1}{z-a_1-\cfrac{\beta_2}{z-a_2}}},
\]
其中
\[
a_0=\frac{53}{16},\qquad
a_1=\frac{2703}{976},\qquad
a_2=\frac{178}{61},\qquad
\beta_1=\frac{183}{256},\qquad
\beta_2=\frac{768}{3721}.
\]
代入即可得到所示 continued fraction。再把它化简，便回到定理 \ref{thm:xi-time-part9zc-window6-microstate-biased-fiberlaw-herglotz-closedform} 的有理闭式。证毕。
\end{proof}
\endinput

\noindent
在微态偏置纤维律已被压缩为三原子 Cauchy--Stieltjes 核、唯一 Jacobi Hamiltonian 与 Cartan--Lanczos 共轭之后，同一谱测度链还可继续沿动力学与刚性两侧闭合：其一，\(J_6\) 落在一条完全可积的三站点 Toda 轨道上；其二，该轨道又可被重新解释为 regular semisimple 轨道上的唯一正 Jacobi 截面；其三，Lie 完成的秩、谱支撑层数与 Krylov 闭包维数在这一窗口上发生严格一致。

\begin{theorem}[window-\texorpdfstring{$6$}{6} 的三站点 Toda 分子]\label{thm:xi-time-part9zd-window6-three-site-toda-molecule}
\leanverified{paper\_xi\_time\_part9zd\_window6\_three\_site\_toda\_molecule}
记
\[
(\lambda_1,\lambda_2,\lambda_3)=(2,3,4),
\qquad
(w_1,w_2,w_3)=\left(\frac14,\frac3{16},\frac9{16}\right).
\]
定义
\[
\tau_0(t)\equiv 1,
\]
\[
\tau_1(t)
=
\frac14e^{4t}+\frac3{16}e^{6t}+\frac9{16}e^{8t}
=
\frac{e^{4t}}{16}\bigl(4+3e^{2t}+9e^{4t}\bigr),
\]
\[
\tau_2(t)
=
\frac3{64}e^{10t}+\frac9{16}e^{12t}+\frac{27}{256}e^{14t}
=
\frac{3e^{10t}}{256}\bigl(4+48e^{2t}+9e^{4t}\bigr),
\]
\[
\tau_3(t)=\frac{27}{256}e^{18t}.
\]
再定义
\[
a_k(t):=\frac12\frac{\dd}{\dd t}\log\frac{\tau_k(t)}{\tau_{k-1}(t)}
\qquad (k=1,2,3),
\]
\[
b_k(t):=\left(\frac{\tau_{k+1}(t)\tau_{k-1}(t)}{\tau_k(t)^2}\right)^{1/2}>0
\qquad (k=1,2),
\]
并置
\[
J_6(t):=
\begin{pmatrix}
a_1(t)&b_1(t)&0\\
b_1(t)&a_2(t)&b_2(t)\\
0&b_2(t)&a_3(t)
\end{pmatrix}.
\]
则 \(J_6(t)\) 满足三粒子非周期 Toda--Flaschka 方程
\[
\dot a_1=2b_1^2,\qquad
\dot a_2=2(b_2^2-b_1^2),\qquad
\dot a_3=-2b_2^2,
\]
\[
\dot b_1=b_1(a_2-a_1),\qquad
\dot b_2=b_2(a_3-a_2),
\]
并且
\[
J_6(0)=J_6,
\qquad
\operatorname{spec}J_6(t)=\{2,3,4\}
\quad (\forall\,t\in\RR).
\]
因此，window-\(6\) 的规范信息 Hamiltonian 并非静态对象，而是一个完全可积的三站点 Toda 轨道上的精确初值点。
\end{theorem}
\begin{proof}
对有限 Jacobi 矩阵的逆谱表示，给定离散谱数据 \((\lambda_j,w_j)_{j=1}^3\) 后，开放 Toda 链的标准 \(\tau\)-函数正是
\[
\tau_1(t)=\sum_{j=1}^3 w_j e^{2\lambda_j t},
\]
\[
\tau_2(t)=\sum_{1\le i<j\le 3}w_iw_j(\lambda_j-\lambda_i)^2e^{2(\lambda_i+\lambda_j)t},
\]
\[
\tau_3(t)=w_1w_2w_3\prod_{1\le i<j\le 3}(\lambda_j-\lambda_i)^2e^{2(\lambda_1+\lambda_2+\lambda_3)t}.
\]
将
\[
(\lambda_1,\lambda_2,\lambda_3)=(2,3,4),
\qquad
\left(w_1,w_2,w_3\right)=\left(\frac14,\frac3{16},\frac9{16}\right)
\]
代入，立即得到题设三式。由同一逆谱构造，
\[
a_k(t)=\frac12\frac{\dd}{\dd t}\log\frac{\tau_k}{\tau_{k-1}},
\qquad
b_k(t)^2=\frac{\tau_{k+1}\tau_{k-1}}{\tau_k^2}
\]
自动满足 Toda--Flaschka 方程，并保持谱数据不变，因此
\[
\operatorname{spec}J_6(t)=\{2,3,4\}
\qquad (\forall\,t\in\RR).
\]

最后直接代入 \(t=0\)：
\[
\tau_1(0)=1,\qquad
\tau_2(0)=\frac{183}{256},\qquad
\tau_3(0)=\frac{27}{256}.
\]
故
\[
a_1(0)=\frac12\tau_1'(0)=\frac{53}{16},
\]
\[
b_1(0)^2=\tau_2(0)=\frac{183}{256},
\]
\[
a_2(0)=\frac12\frac{\dd}{\dd t}\log\frac{\tau_2}{\tau_1}\Big|_{t=0}
=\frac{2703}{976},
\]
\[
b_2(0)^2=\frac{\tau_3(0)\tau_1(0)}{\tau_2(0)^2}
=\frac{768}{3721},
\]
\[
a_3(0)=\frac12\frac{\dd}{\dd t}\log\frac{\tau_3}{\tau_2}\Big|_{t=0}
=\frac{178}{61}.
\]
因此 \(J_6(0)\) 恰等于定理 \ref{thm:xi-time-part9zc-window6-canonical-information-hamiltonian} 中的 \(J_6\)。证毕。
\end{proof}

\begin{corollary}[Cartan 排序之间的异宿轨道]\label{cor:xi-time-part9zd-window6-cartan-ordering-heteroclinic}
在定理 \ref{thm:xi-time-part9zd-window6-three-site-toda-molecule} 的 Toda 轨道上，有极限
\[
\lim_{t\to-\infty}J_6(t)=\operatorname{diag}(2,3,4)=C_6,
\]
\[
\lim_{t\to+\infty}J_6(t)=\operatorname{diag}(4,3,2),
\]
并且
\[
\lim_{t\to\pm\infty}b_1(t)=\lim_{t\to\pm\infty}b_2(t)=0.
\]
因此，\(J_6\) 精确位于一条连接两种 Cartan 排序的异宿轨道上。
\end{corollary}
\begin{proof}
当 \(t\to-\infty\) 时，
\[
\tau_1(t)\sim \frac14e^{4t},\qquad
\tau_2(t)\sim \frac3{64}e^{10t},\qquad
\tau_3(t)\sim \frac{27}{256}e^{18t}.
\]
故
\[
a_1(t)\to 2,\qquad
a_2(t)\to 3,\qquad
a_3(t)\to 4.
\]
同时
\[
b_1(t)^2=\frac{\tau_2\tau_0}{\tau_1^2}\sim \frac34 e^{2t}\to 0,
\qquad
b_2(t)^2=\frac{\tau_3\tau_1}{\tau_2^2}\sim 3e^{2t}\to 0.
\]

当 \(t\to+\infty\) 时，
\[
\tau_1(t)\sim \frac9{16}e^{8t},\qquad
\tau_2(t)\sim \frac{27}{256}e^{14t},\qquad
\tau_3(t)\sim \frac{27}{256}e^{18t}.
\]
于是
\[
a_1(t)\to 4,\qquad
a_2(t)\to 3,\qquad
a_3(t)\to 2,
\]
并且
\[
b_1(t)^2\sim \frac13 e^{-2t}\to 0,
\qquad
b_2(t)^2\sim \frac{16}{3} e^{-2t}\to 0.
\]
证毕。
\end{proof}

\begin{theorem}[regular semisimple 轨道上的唯一正 Jacobi 截面]\label{thm:xi-time-part9zd-window6-unique-positive-jacobi-section}
记
\[
\mathcal O_{2,3,4}
:=
\left\{
Q\,\operatorname{diag}(2,3,4)\,Q^{\mathsf T}:Q\in SO(3)
\right\}
\subset \operatorname{Sym}_3(\RR).
\]
则在该 regular semisimple 轨道上，满足下列两条的矩阵唯一存在：
\begin{enumerate}
  \item 它是 Jacobi 型，即对称三对角且两条次对角元都严格为正；
  \item 它在循环向量 \(e_1\) 下的谱测度恰为 \(\nu_6\)。
\end{enumerate}
该唯一矩阵正是定理 \ref{thm:xi-time-part9zc-window6-canonical-information-hamiltonian} 中的 \(J_6\)。
\end{theorem}
\begin{proof}
由定理 \ref{thm:xi-time-part9zc-window6-canonical-information-hamiltonian} 与定理 \ref{thm:xi-time-part9zc-window6-cartan-lanczos-conjugacy}，\(J_6\in \mathcal O_{2,3,4}\)，且其两条次对角元
\[
\frac{\sqrt{183}}{16},
\qquad
\frac{16\sqrt3}{61}
\]
都严格为正，并具有谱测度 \(\nu_6\)。

反过来，设 \(J\) 是任一满足两条条件的 Jacobi 矩阵。则其 \(e_1\)-谱测度为 \(\nu_6\)。对有限正测度，谱测度通过 Favard 三项递推唯一决定 Jacobi 系数；而定理 \ref{thm:xi-time-part9zc-window6-orthogonal-polynomials-cubic-termination} 已经把 \(\nu_6\) 的 monic 正交多项式与递推系数全部显式写出。因此 \(J\) 的对角元与正次对角元都被唯一确定，只能等于定理 \ref{thm:xi-time-part9zc-window6-canonical-information-hamiltonian} 中的 \(J_6\)。证毕。
\end{proof}

\begin{theorem}[Lie rank、支撑层数与 Krylov rank 的严格一致]\label{thm:xi-time-part9zd-window6-lie-support-krylov-rank-coincidence}
\leanverified{paper\_xi\_time\_part9zd\_window6\_lie\_support\_krylov\_rank\_coincidence}
沿用定理 \ref{thm:conclusion-window6-minuscule-eightfold-b3-spin7}、定理 \ref{thm:conclusion-window6-projective-long-sector-spin1-crystal-field} 与定理 \ref{thm:xi-time-part9zc-window6-spin1-cartan-born-realization} 的记号，则有
\[
\operatorname{rank}\mathrm{Spin}(7)
=
\left|\operatorname{supp}\nu_6\right|
=
\dim \operatorname{span}\{\psi_6,C_6\psi_6,C_6^2\psi_6\}.
\]
等价地，
\[
\text{Lie rank}
=
\text{微态偏置纤维律的支撑层数}
=
\text{Lanczos 链长度}.
\]
因此，window-\(6\) 中同一个整数 \(3\) 同时刻画了 Lie 完成的环面秩、规范信息 Hamiltonian 的谱层数以及从 Cartan 可观测量出发的最小 Krylov 闭包维数。
\end{theorem}
\begin{proof}
由定理 \ref{thm:conclusion-window6-minimal-compact-lie-completion-circle-rigidity}，极小紧 Lie 完成具有秩 \(3\)；再由定理 \ref{thm:conclusion-window6-minuscule-eightfold-b3-spin7}，其单连通紧完成唯一为 \(\mathrm{Spin}(7)\)。故
\[
\operatorname{rank}\mathrm{Spin}(7)=3.
\]
另一方面，由定理 \ref{thm:xi-time-part9zc-window6-microstate-biased-fiberlaw-herglotz-closedform},
\[
\operatorname{supp}\nu_6=\{2,3,4\},
\]
故
\[
\left|\operatorname{supp}\nu_6\right|=3.
\]

最后，由定理 \ref{thm:xi-time-part9zc-window6-spin1-cartan-born-realization},
\[
C_6=\operatorname{diag}(2,3,4),
\qquad
\psi_6=\left(\frac12,\frac{\sqrt3}{4},\frac34\right)^{\mathsf T},
\]
且 \(\psi_6\) 在三个特征值对应的特征方向上投影都非零。因此 \(C_6\) 在循环子空间
\[
\mathcal K:=\operatorname{span}\{\psi_6,C_6\psi_6,C_6^2\psi_6,\dots\}
\]
上的最小多项式正是
\[
(\lambda-2)(\lambda-3)(\lambda-4),
\]
其次数为 \(3\)。故 \(\dim\mathcal K=3\)，而在三维空间中
\[
\mathcal K=\operatorname{span}\{\psi_6,C_6\psi_6,C_6^2\psi_6\}.
\]
证毕。
\end{proof}

\begin{conjecture}[Cartan--Lanczos 刚性原理]\label{conj:xi-time-part9zd-window6-cartan-lanczos-rigidity-principle}
对任意偶数窗口 \(m\)，设其 projective long sector 已在某一不可约紧 Lie 完成 \(K_m\) 中线性化为 Cartan 型可观测量
\[
C_m,
\]
并令
\[
\nu_m:=2^{-m}\sum_{x\in X_m}d_m(x)\,\delta_{d_m(x)}
\]
为微态偏置纤维律。若满足
\[
\left|\operatorname{supp}\nu_m\right|=\operatorname{rank}K_m,
\]
则应有：
\begin{enumerate}
  \item \(\nu_m\) 等于 \(C_m\) 在某一规范态 \(\psi_m\) 下的谱律；
  \item \(\nu_m\) 的 Jacobi--Lanczos 矩阵 \(J_m\) 与 \(C_m\) 正交共轭；
  \item \(J_m\) 给出该窗口的规范信息 Hamiltonian，而其非周期 Toda 轨道在不同 Cartan 排序之间建立一条异宿插值。
\end{enumerate}
window-\(6\) 提供了该原理的首个严格非平凡实例：由定理 \ref{thm:conclusion-window6-minuscule-eightfold-b3-spin7} 与定理 \ref{thm:conclusion-window6-minimal-compact-lie-completion-circle-rigidity}，
\[
K_{6,\mathrm{sc}}\cong \mathrm{Spin}(7),
\qquad
\operatorname{rank}K_{6,\mathrm{sc}}=\operatorname{rank}K_6=3,
\]
而由定理 \ref{thm:xi-time-part9zc-window6-microstate-biased-fiberlaw-herglotz-closedform},
\[
\operatorname{supp}\nu_6=\{2,3,4\}.
\]
定理 \ref{thm:xi-time-part9zc-window6-spin1-cartan-born-realization} 至定理 \ref{thm:xi-time-part9zd-window6-lie-support-krylov-rank-coincidence} 已经完成了 \(m=6\) 的全部验证。
\end{conjecture}

\noindent
上述结果表明，window-\(6\) 的微态偏置纤维律并非对直方图 \(8:4:9\) 的平凡重写，而是同一个 regular semisimple 轨道在五种不同语言下的完全等价外显：它首先表现为三原子 Cauchy--Stieltjes/Herglotz 核，其后在正交多项式语言中闭合为唯一的三站点 Jacobi Hamiltonian，再在量子测量语言中化为 spin-\(1\) Cartan 可观测量 \(3I_3+J_z\) 于态 \(\psi_6\) 下的 Born 分布，又在 Krylov--Lanczos 语言中转写为与 Cartan 基正交共轭的信息 Hamiltonian，最后在可积系统语言中延展为连接两种 Cartan 排序的 Toda 异宿轨道。因而，Lie 理论、量子测量、有限支撑矩问题、Jacobi 逆谱与有限信息热力学在这一窗口上并不是并列结构，而是同一对象的五种严格同构坐标。

\endinput

\paragraph{window-\texorpdfstring{$6$}{6} 的有限指数链、零温冻结与中心熵损失}
在 window-\(6\) 的三原子矩谱、Cartan--Lanczos 共轭、Toda 等谱轨道以及中心谱指数元已经完全闭合之后，同一份 rigid split 数据还可继续在条件期望、零温冻结与有限维熵损失三侧显影。以下沿用命题 \ref{prop:fold-bin-groupoid-algebra-wedderburn} 的 Wedderburn 分解、推论 \ref{cor:fold-bin-m6-fiber-hist} 的直方图 \(2:8,\ 3:4,\ 4:9\)、定理 \ref{thm:xi-foldbin-fkdet-logvolume-kappa} 的 \(\kappa_m\) 记号、定理 \ref{thm:xi-time-part9zc-window6-spin1-cartan-born-realization} 与定理 \ref{thm:xi-time-part9zc-window6-cartan-lanczos-conjugacy} 中的 \(C_6,J_6\)，并记
\[
d_m(x):=d_m^{\mathrm{bin}}(x),
\qquad
\mathcal A_m:=\bigoplus_{x\in X_m}M_{d_m(x)}(\CC),
\qquad
d_{\max}(m):=\max_{x\in X_m}d_m(x),
\qquad
d_{\min}(m):=\min_{x\in X_m}d_m(x),
\]
\[
\mathcal S_{6,2}:=\{w\in X_6:\ d_6(w)=2\},
\qquad
\mathcal S_{6,4}:=\{w\in X_6:\ d_6(w)=4\}.
\]
对任意保单位正映射 \(E\)，定义其 Pimsner--Popa 常数与指数
\[
\lambda(E):=\sup\{\lambda>0:\ E(a)\ge \lambda a\ \forall\,a\ge 0\},
\qquad
\operatorname{Ind}_{\mathrm{PP}}(E):=\lambda(E)^{-1}.
\]

\begin{theorem}[fold 代数中心条件期望的 Pimsner--Popa 指数公式]\label{thm:xi-time-part9ze-center-condexp-pimsner-popa-index}
\leanverified{paper\_xi\_time\_part9ze\_center\_condexp\_pimsner\_popa\_index}
对任意有限分辨率 \(m\)，定义
\[
E_m^Z:\mathcal A_m\to Z(\mathcal A_m),
\qquad
E_m^Z\bigl((A_x)_{x\in X_m}\bigr)
:=
\left(
\frac{\operatorname{Tr}(A_x)}{d_m(x)}I_{d_m(x)}
\right)_{x\in X_m}.
\]
则
\[
\operatorname{Ind}_{\mathrm{PP}}(E_m^Z)=d_{\max}(m).
\]
特别地，在 window-\(6\) 上有
\[
\operatorname{Ind}_{\mathrm{PP}}(E_6^Z)=4.
\]
\end{theorem}
\begin{proof}
该映射正是定理 \ref{thm:xi-foldbin-minimal-center-valued-index-spectrum-top} 中的逐块归一化迹条件期望。对任意
\[
a=(A_x)_{x\in X_m}\in \mathcal A_m^+,
\]
逐块都有
\[
E_m^Z(a)_x
=
\frac{\operatorname{Tr}(A_x)}{d_m(x)}I_{d_m(x)}
\ge
\frac{1}{d_m(x)}A_x
\ge
\frac{1}{d_{\max}(m)}A_x,
\]
因为正矩阵总满足 \(A_x\le \operatorname{Tr}(A_x)I_{d_m(x)}\)。故
\[
E_m^Z(a)\ge \frac{1}{d_{\max}(m)}a,
\]
从而
\[
\lambda(E_m^Z)\ge \frac{1}{d_{\max}(m)}.
\]

反过来，取某个满足 \(d_m(x_0)=d_{\max}(m)\) 的块，并在该块中取秩一投影 \(p\)。则
\[
E_m^Z(p)=\frac{1}{d_{\max}(m)}I_{d_{\max}(m)}.
\]
若 \(\lambda>d_{\max}(m)^{-1}\)，则
\[
E_m^Z(p)-\lambda p
\]
在 \(\operatorname{Ran}(p)\) 上的特征值为 \(d_{\max}(m)^{-1}-\lambda<0\)，故不再为正。于是
\[
\lambda(E_m^Z)=\frac{1}{d_{\max}(m)},
\qquad
\operatorname{Ind}_{\mathrm{PP}}(E_m^Z)=d_{\max}(m).
\]
对 \(m=6\)，由推论 \ref{cor:fold-bin-m6-fiber-hist} 的直方图得 \(d_{\max}(6)=4\)，从而结论成立。
\end{proof}

\begin{theorem}[正零温冻结的有限窗精确式]\label{thm:xi-time-part9ze-positive-zero-temperature-freezing}
定义有限窗配分函数
\[
Z_m(q):=\sum_{x\in X_m}d_m(x)^q,
\qquad q\in \RR,
\]
并记
\[
K_m^{\max}:=\#\{x\in X_m:\ d_m(x)=d_{\max}(m)\}.
\]
则对任意固定 \(m\) 都有
\[
\lim_{q\to+\infty}\frac{1}{q}\log Z_m(q)
=
\log \operatorname{Ind}_{\mathrm{PP}}(E_m^Z),
\]
以及
\[
\lim_{q\to+\infty}
e^{-q\log \operatorname{Ind}_{\mathrm{PP}}(E_m^Z)}\,Z_m(q)
=
K_m^{\max}.
\]

在 window-\(6\) 上，
\[
Z_6(q)=8\cdot 2^q+4\cdot 3^q+9\cdot 4^q
=
4^q\left(
9+4\left(\frac{3}{4}\right)^q+8\left(\frac{1}{2}\right)^q
\right),
\]
故
\[
\log Z_6(q)
=
q\log 4+\log 9+
\log\!\left(
1+\frac{4}{9}\left(\frac{3}{4}\right)^q
+\frac{8}{9}\left(\frac{1}{2}\right)^q
\right).
\]
\leanverified{paper\_xi\_time\_part9ze\_positive\_zero\_temperature\_freezing}
\end{theorem}
\begin{proof}
把 \(d_{\max}(m)^q\) 提出即可得到
\[
Z_m(q)
=
d_{\max}(m)^q
\left(
K_m^{\max}
+
\sum_{d_m(x)<d_{\max}(m)}
\left(
\frac{d_m(x)}{d_{\max}(m)}
\right)^q
\right).
\]
当 \(q\to+\infty\) 时，括号内第二项趋于 \(0\)。由定理
\ref{thm:xi-time-part9ze-center-condexp-pimsner-popa-index},
\[
d_{\max}(m)=\operatorname{Ind}_{\mathrm{PP}}(E_m^Z),
\]
于是两条极限式同时成立。window-\(6\) 的展开只需把直方图
\[
(2:8,\ 3:4,\ 4:9)
\]
代入即可。
\end{proof}

\begin{theorem}[负零温冻结与冷基态计数]\label{thm:xi-time-part9ze-negative-zero-temperature-freezing}
\leanverified{paper\_xi\_time\_part9ze\_negative\_zero\_temperature\_freezing}
记
\[
K_m^{\min}:=\#\{x\in X_m:\ d_m(x)=d_{\min}(m)\}.
\]
则
\[
\lim_{q\to-\infty}\frac{1}{q}\log Z_m(q)=\log d_{\min}(m),
\]
并且
\[
\lim_{q\to-\infty}e^{-q\log d_{\min}(m)}\,Z_m(q)=K_m^{\min}.
\]

在 window-\(6\) 上，
\[
d_{\min}(6)=2,
\qquad
K_6^{\min}=8,
\]
且有精确分解
\[
Z_6(q)
=
2^q\left(
8+4\left(\frac{3}{2}\right)^q+9\cdot 2^q
\right).
\]
因此
\[
\log Z_6(q)
=
q\log 2+\log 8+
\log\!\left(
1+\frac{1}{2}\left(\frac{3}{2}\right)^q+\frac{9}{8}2^q
\right)
\qquad (q\to-\infty).
\]
\end{theorem}
\begin{proof}
同样把 \(d_{\min}(m)^q\) 提出，得到
\[
Z_m(q)
=
d_{\min}(m)^q
\left(
K_m^{\min}
+
\sum_{d_m(x)>d_{\min}(m)}
\left(
\frac{d_m(x)}{d_{\min}(m)}
\right)^q
\right).
\]
当 \(q\to-\infty\) 时，对一切 \(d_m(x)>d_{\min}(m)\) 都有
\[
\left(
\frac{d_m(x)}{d_{\min}(m)}
\right)^q\to 0,
\]
故结论成立。window-\(6\) 的公式再次由直方图直接代入得出。
\end{proof}

\begin{theorem}[标量条件期望的全地址指数与加权中心指数]\label{thm:xi-time-part9ze-scalar-weighted-center-condexp-indices}
定义 canonical microstate trace
\[
\tau_m\bigl((A_x)_{x\in X_m}\bigr)
:=
2^{-m}\sum_{x\in X_m}\operatorname{Tr}(A_x),
\]
以及标量条件期望
\[
E_m^{\mathrm{sc}}(a):=\tau_m(a)\,1_{\mathcal A_m}.
\]
则
\[
\operatorname{Ind}_{\mathrm{PP}}(E_m^{\mathrm{sc}})=2^m.
\]

另一方面，把中心识别为 \(Z(\mathcal A_m)\cong \CC^{X_m}\)，并定义加权中心标量期望
\[
\varepsilon_m\bigl((\lambda_x)_{x\in X_m}\bigr)
:=
\left(
\sum_{x\in X_m}\frac{d_m(x)}{2^m}\lambda_x
\right)1.
\]
则
\[
\operatorname{Ind}_{\mathrm{PP}}(\varepsilon_m)=\frac{2^m}{d_{\min}(m)}.
\]

特别地，在 window-\(6\) 上，
\[
\operatorname{Ind}_{\mathrm{PP}}(E_6^{\mathrm{sc}})=64,
\qquad
\operatorname{Ind}_{\mathrm{PP}}(\varepsilon_6)=32.
\]
\end{theorem}
\begin{proof}
先证 \(E_m^{\mathrm{sc}}\)。任取
\[
a=(A_x)_{x\in X_m}\in \mathcal A_m^+.
\]
对任意固定的块 \(x\)，有
\[
\tau_m(a)\,I_{d_m(x)}
=
2^{-m}\sum_{y\in X_m}\operatorname{Tr}(A_y)\,I_{d_m(x)}
\ge
2^{-m}\operatorname{Tr}(A_x)\,I_{d_m(x)}
\ge
2^{-m}A_x.
\]
故逐块并合得到
\[
E_m^{\mathrm{sc}}(a)\ge 2^{-m}a,
\]
从而
\[
\lambda(E_m^{\mathrm{sc}})\ge 2^{-m}.
\]
若 \(\lambda>2^{-m}\)，取任意块中的秩一投影 \(p\)，则
\[
E_m^{\mathrm{sc}}(p)=2^{-m}1_{\mathcal A_m},
\]
而
\[
E_m^{\mathrm{sc}}(p)-\lambda p
\]
在 \(p\) 的像上取值 \(2^{-m}-\lambda<0\)。故
\[
\lambda(E_m^{\mathrm{sc}})=2^{-m},
\qquad
\operatorname{Ind}_{\mathrm{PP}}(E_m^{\mathrm{sc}})=2^m.
\]

再证 \(\varepsilon_m\)。设
\[
p_x:=\frac{d_m(x)}{2^m},
\qquad
p_{\min}:=\min_{x\in X_m}p_x=\frac{d_{\min}(m)}{2^m}.
\]
对任意正元
\[
f=(f_x)_{x\in X_m}\in Z(\mathcal A_m)^+\cong \CC_+^{X_m},
\]
任取 \(y\in X_m\)，都有
\[
\varepsilon_m(f)=\sum_{x\in X_m}p_x f_x\ge p_y f_y\ge p_{\min} f_y.
\]
因此
\[
\varepsilon_m(f)\,1\ge p_{\min} f,
\]
从而
\[
\lambda(\varepsilon_m)\ge p_{\min}.
\]
若 \(\lambda>p_{\min}\)，取某个满足 \(d_m(y)=d_{\min}(m)\) 的坐标，并记 \(z_y\in Z(\mathcal A_m)\) 为对应的中心原始幂等元。则
\[
\varepsilon_m(z_y)=\frac{d_{\min}(m)}{2^m}1=p_{\min}1.
\]
于是
\[
\varepsilon_m(z_y)-\lambda z_y
\]
在 \(y\)-坐标上等于 \(p_{\min}-\lambda<0\)，不再为正。故
\[
\lambda(\varepsilon_m)=\frac{d_{\min}(m)}{2^m},
\qquad
\operatorname{Ind}_{\mathrm{PP}}(\varepsilon_m)=\frac{2^m}{d_{\min}(m)}.
\]
对 \(m=6\)，由 \(d_{\min}(6)=2\) 立即得到
\[
\operatorname{Ind}_{\mathrm{PP}}(E_6^{\mathrm{sc}})=64,
\qquad
\operatorname{Ind}_{\mathrm{PP}}(\varepsilon_6)=32.
\]
\end{proof}

\begin{corollary}[window-\texorpdfstring{$6$}{6} 的有限指数 bit 梯级]\label{cor:xi-time-part9ze-window6-bit-ladder}
\leanverified{paper\_xi\_time\_part9ze\_window6\_bit\_ladder}
window-\(6\) 上存在严格的三层 dyadic 指数链
\[
\log_2\operatorname{Ind}_{\mathrm{PP}}(E_6^Z)=2,
\qquad
\log_2\operatorname{Ind}_{\mathrm{PP}}(\varepsilon_6)=5,
\qquad
\log_2\operatorname{Ind}_{\mathrm{PP}}(E_6^{\mathrm{sc}})=6.
\]
其中第一层恢复精确反演所需的最小 dyadic 预算，第二层恢复 microstate law
\[
p_6(w):=\frac{d_6(w)}{64}
\]
下稳定类型的最坏情形信息代价
\[
\max_{w\in X_6}\bigl(-\log_2 p_6(w)\bigr)=5,
\]
且该最大值只在 \(\mathcal S_{6,2}\) 上达到；第三层则恢复全部微态地址长度 \(\log_2 64=6\)。
\end{corollary}
\begin{proof}
由定理 \ref{thm:xi-time-part9ze-center-condexp-pimsner-popa-index} 与定理
\ref{thm:xi-time-part9ze-scalar-weighted-center-condexp-indices},
\[
\operatorname{Ind}_{\mathrm{PP}}(E_6^Z)=4,
\qquad
\operatorname{Ind}_{\mathrm{PP}}(\varepsilon_6)=32,
\qquad
\operatorname{Ind}_{\mathrm{PP}}(E_6^{\mathrm{sc}})=64,
\]
故三条对数恒等式立即成立。

再由
\[
p_6(w)=\frac{d_6(w)}{64},
\qquad
d_6(w)\in\{2,3,4\},
\]
可得
\[
\max_{w\in X_6}\bigl(-\log_2 p_6(w)\bigr)
=
-\log_2\frac{d_{\min}(6)}{64}
=
-\log_2\frac{2}{64}
=
5.
\]
这一最大值当且仅当 \(d_6(w)=2\) 时达到，也即只在 \(\mathcal S_{6,2}\) 上达到。另一方面，定理
\ref{thm:xi-time-part9x-reversible-capacity-two-bit-inversion} 已把第一层解释为精确反演所需的二比特预算；第三层则只是
\[
2^6=64
\]
的对数写法。证毕。
\end{proof}

\begin{theorem}[中心条件期望指数与 Cartan--Lanczos 谱半径一致]\label{thm:xi-time-part9ze-index-spectral-radius-consistency}
\leanverified{paper\_xi\_time\_part9ze\_index\_spectral\_radius\_consistency}
沿用定理 \ref{thm:xi-time-part9zc-window6-spin1-cartan-born-realization} 的
\[
C_6=3I_3+J_z=\operatorname{diag}(2,3,4)
\]
与定理 \ref{thm:xi-time-part9zc-window6-canonical-information-hamiltonian} 的 \(J_6\)。则
\[
\operatorname{Ind}_{\mathrm{PP}}(E_6^Z)=\rho(C_6)=\rho(J_6)=4.
\]
因此，window-\(6\) 的局域不可逆度顶端与 Cartan 场的最大能级、以及对应 Jacobi Hamiltonian 的谱半径严格同一。
\end{theorem}
\begin{proof}
由定理 \ref{thm:xi-time-part9zc-window6-cartan-lanczos-conjugacy},
\[
J_6=U\,C_6\,U^{\mathsf T}
\]
对某个 \(U\in SO(3)\) 成立，故
\[
\rho(J_6)=\rho(C_6).
\]
而
\[
C_6=\operatorname{diag}(2,3,4)
\]
直接给出
\[
\rho(C_6)=4.
\]
再由定理 \ref{thm:xi-time-part9ze-center-condexp-pimsner-popa-index},
\[
\operatorname{Ind}_{\mathrm{PP}}(E_6^Z)=4.
\]
合并即得结论。
\end{proof}

\begin{proposition}[Pimsner--Popa 临界射线的最大壳层局域化]\label{prop:xi-time-part9ze-window6-jones-ray-maxshell-localization}
令 \(\mathcal P_4\) 为 \(\mathcal A_6\) 中全部位于 \(4\times 4\) 块内的秩一投影之集合。则有
\[
\forall\,\lambda>\frac14,\ \exists\,p\in \mathcal P_4
\ \text{使得}\ 
E_6^Z(p)\not\ge \lambda p.
\]
反之，若 \(x\in \mathcal A_6^+\) 只支撑于 \(d_6=2\) 与 \(d_6=3\) 的块上，则
\[
E_6^Z(x)\ge \frac13 x.
\]
因而，从 \(1/3\) 降到 \(1/4\) 的全部 Pimsner--Popa 障碍都严格局域于最大壳层 \(\mathcal S_{6,4}\)。
\end{proposition}
\begin{proof}
若 \(p\in \mathcal P_4\)，则它落在某个 \(4\times 4\) 块内，故
\[
E_6^Z(p)=\frac14 I_4
\]
于该块上成立。若 \(\lambda>1/4\)，则
\[
E_6^Z(p)-\lambda p
=
\frac14 I_4-\lambda p
\]
在 \(\operatorname{Ran}(p)\) 上的特征值为 \(1/4-\lambda<0\)，故不为正。

另一方面，若 \(x\ge 0\) 只支撑于 \(d_6=2\) 与 \(d_6=3\) 的块上，则对每个非零块 \(w\) 都有
\[
E_6^Z(x)_w
=
\frac{\operatorname{Tr}(x_w)}{d_6(w)}I_{d_6(w)}
\ge
\frac{1}{d_6(w)}x_w
\ge
\frac13 x_w.
\]
逐块并合便得
\[
E_6^Z(x)\ge \frac13 x.
\]
这表明全局最坏常数 \(1/4\) 只能由 \(d_6=4\) 的块，也即 \(\mathcal S_{6,4}\)，给出。
\end{proof}

\begin{theorem}[\texorpdfstring{$\kappa_6$}{kappa6} 的算子代数熵解释]\label{thm:xi-time-part9ze-kappa6-vn-entropy-loss}
令 \(\tau_6\) 为定理 \ref{thm:xi-time-part9ze-scalar-weighted-center-condexp-indices} 中的 canonical microstate trace，并以
\[
S(\omega):=-\operatorname{Tr}(\rho_\omega\log \rho_\omega)
\]
记有限维状态 \(\omega\) 的 von Neumann 熵。则
\[
S(\tau_6)=6\log 2.
\]
把 \(\tau_6\) 限制到中心 \(Z(\mathcal A_6)\cong \CC^{X_6}\) 后，得到经典分布
\[
p_6(w)=\frac{d_6(w)}{64},
\]
其熵满足
\[
S(\tau_6|_{Z(\mathcal A_6)})
=
-\sum_{w\in X_6}\frac{d_6(w)}{64}\log\frac{d_6(w)}{64}
=
6\log 2-\kappa_6.
\]
特别地，
\[
\kappa_6=\frac{11}{8}\log 2+\frac{3}{16}\log 3=1.159067\ldots,
\]
从而
\[
S(\tau_6|_{Z(\mathcal A_6)})
=
\frac{37}{8}\log 2-\frac{3}{16}\log 3.
\]
因此
\[
S(\tau_6)-S(\tau_6|_{Z(\mathcal A_6)})=\kappa_6.
\]
\end{theorem}
\begin{proof}
\(\tau_6\) 在
\[
\mathcal A_6=\bigoplus_{w\in X_6}M_{d_6(w)}(\CC)
\]
上的密度矩阵为
\[
\rho_6=\bigoplus_{w\in X_6}\frac{1}{64}I_{d_6(w)}.
\]
由于
\[
\sum_{w\in X_6}d_6(w)=64,
\]
\(\rho_6\) 的全部非零特征值恰有 \(64\) 个，且都等于 \(1/64\)。于是
\[
S(\tau_6)
=
-64\cdot \frac{1}{64}\log\frac{1}{64}
=
6\log 2.
\]

再看中心限制。对每个中心原始幂等元 \(z_w\)，有
\[
\tau_6(z_w)=\frac{d_6(w)}{64}=p_6(w),
\]
故 \(\tau_6|_{Z(\mathcal A_6)}\) 正对应于概率分布 \((p_6(w))_{w\in X_6}\)。因此
\[
S(\tau_6|_{Z(\mathcal A_6)})
=
-\sum_{w\in X_6}p_6(w)\log p_6(w)
=
-\sum_{w\in X_6}\frac{d_6(w)}{64}
\log\frac{d_6(w)}{64}.
\]
展开后得到
\[
S(\tau_6|_{Z(\mathcal A_6)})
=
6\log 2-\frac{1}{64}\sum_{w\in X_6}d_6(w)\log d_6(w)
=
6\log 2-\kappa_6,
\]
其中最后一步正是定理 \ref{thm:xi-foldbin-fkdet-logvolume-kappa} 的 \(m=6\) 情形。把直方图
\[
(2:8,\ 3:4,\ 4:9)
\]
代入即可得到 \(\kappa_6\) 与中心熵的显式值。证毕。
\end{proof}

\endinput

\paragraph{Wulff 射线的 Schur 极小谱、低纤维异常税与 window-\texorpdfstring{$6$}{6} 审计刚性}
在定理 \ref{thm:xi-time-part9x-serrin-ray-fibonacci-semirings} 已把未归一化各向异性 Serrin 射线的分辨率 \(m\) 观测识别为模 \(F_{m+2}\) 的有限半环、推论 \ref{cor:xi-time-part9x-serrin-ray-exact-capacity} 已把窗口 \(0\le \rho<2^m\) 上的残余纤维写成 \(\{q_m,q_m+1\}\) 二值直方图、而 window-\(6\) 的 bin-fold 纤维谱、boundary sheet-parity 中心与规范热力学亦已固定之后，还可进一步把这条 Wulff 射线提升为有限分辨率上的 Schur 极小基准。其关键点在于：Figalli--Zhang 的各向异性 Serrin 刚性在模去平移后只留下单参数 Wulff 射线，因此纯几何侧在每个分辨率上只允许一条 residue-count 基准谱；而一旦把该基准谱置入主序与可分凸泛函框架，它便自动成为全部 \(F_{m+2}\)-纤维模型中的唯一 Schur 极小端点。由此，任何额外的低纤维、boundary 中心或审计异常都不再只是“偏离均衡”的质性描述，而获得一条精确的 convex-tax 定律。 \cite{FigalliZhang2026AnisotropicSerrinGeneralDomains}

\begin{theorem}[Wulff 射线的 Schur 极小谱]\label{thm:xi-time-part9ze-wulff-schur-minimal-spectrum}
\leanverified{paper\_xi\_time\_part9ze\_wulff\_schur\_minimal\_spectrum}
在本段记
\[
M_m:=F_{m+2},\qquad
N_m:=2^m,\qquad
N_m=q_mM_m+a_m,\qquad
0\le a_m<M_m.
\]
把推论 \ref{cor:xi-time-part9x-serrin-ray-exact-capacity} 中 Wulff 射线 residue map 的纤维多重集按非增序重排，记为
\[
d_m^W:=\bigl((q_m+1)^{a_m},q_m^{M_m-a_m}\bigr)\in \mathbb N^{M_m}.
\]
则对任意满射
\[
T:\{0,1,\dots,N_m-1\}\twoheadrightarrow Y,
\qquad |Y|=M_m,
\]
若记其纤维向量为
\[
d(T):=(d_1,\dots,d_{M_m})\in \mathbb N^{M_m},
\qquad
\sum_{i=1}^{M_m}d_i=N_m,
\]
则 \(d(T)\) 经有限次单位平衡步
\[
(\ldots,x,\ldots,y,\ldots)\longmapsto
(\ldots,x-1,\ldots,y+1,\ldots)
\qquad (x-y\ge 2)
\]
必到达且只可能到达 \(d_m^W\) 的某个排列。等价地，
\[
d(T)^\downarrow \succ d_m^W,
\]
故 \(d_m^W\) 是全部 \(M_m\)-分箱、总质量 \(N_m\) 的正整数纤维谱中唯一的 Schur 极小元。
\end{theorem}
\begin{proof}
由推论 \ref{cor:xi-time-part9x-serrin-ray-exact-capacity}，Wulff 射线在窗口 \(0\le \rho<2^m\) 上的纯几何余类计数恰有 \(a_m\) 条纤维取值 \(q_m+1\)、其余 \(M_m-a_m\) 条纤维取值 \(q_m\)，故其非增序纤维谱正是
\[
d_m^W=\bigl((q_m+1)^{a_m},q_m^{M_m-a_m}\bigr).
\]

现取任意 \(d=d(T)\)。若存在 \(i,j\) 使 \(d_i-d_j\ge 2\)，则作一次单位平衡步
\[
(d_i,d_j)\mapsto (d_i-1,d_j+1).
\]
该变换保持总和 \(\sum_i d_i=N_m\) 与正整数性不变，并且平方和满足
\[
(d_i-1)^2+(d_j+1)^2-d_i^2-d_j^2
=-2(d_i-d_j-1)\le -2.
\]
故每做一步，\(\sum_i d_i^2\) 至少下降 \(2\)。由于平方和非负，该过程必在有限步后终止。

终止时所有坐标差都不超过 \(1\)。而满足
\[
\sum_{i=1}^{M_m}d_i=N_m=q_mM_m+a_m
\]
且任意两坐标差至多为 \(1\) 的正整数向量，只能由 \(a_m\) 个 \(q_m+1\) 与 \(M_m-a_m\) 个 \(q_m\) 组成，故终点必为 \(d_m^W\) 的一个排列。于是 \(d\) 经有限次 Robin--Hood 平衡步到达 \(d_m^W\)，Hardy--Littlewood--Pólya 判别遂给出
\[
d(T)^\downarrow\succ d_m^W.
\]
唯一性亦随之成立。证毕。
\end{proof}

\begin{corollary}[Wulff 基准给出可分凸量的统一下包络]\label{cor:xi-time-part9ze-wulff-schur-convex-envelope}
在定理 \ref{thm:xi-time-part9ze-wulff-schur-minimal-spectrum} 的设定下，对任意严格凸函数
\[
\Phi:[1,\infty)\to \RR
\]
都有
\[
\sum_{i=1}^{M_m}\Phi(d_i)\ge
(M_m-a_m)\Phi(q_m)+a_m\Phi(q_m+1),
\]
并且等号当且仅当 \(d(T)\) 与 \(d_m^W\) 仅差一个排列。
\leanverified{paper\_xi\_time\_part9ze\_wulff\_schur\_convex\_envelope}

令
\[
p_i:=\frac{d_i}{N_m},
\qquad
p^W:=\frac{d_m^W}{N_m},
\qquad
u_i:=\frac{1}{M_m},
\]
\[
\operatorname{Col}(T):=\frac{1}{N_m^2}\sum_{i=1}^{M_m}d_i^2,
\qquad
\kappa(T):=\frac{1}{N_m}\sum_{i=1}^{M_m}d_i\log d_i,
\]
\[
G(T):=\prod_{i=1}^{M_m}\mathfrak S_{d_i},
\qquad
g(T):=\frac{1}{N_m}\log|G(T)|.
\]
再记
\[
\operatorname{Col}_m^W:=
\frac{(M_m-a_m)q_m^2+a_m(q_m+1)^2}{N_m^2},
\]
\[
\kappa_m^W:=
\frac{(M_m-a_m)q_m\log q_m+a_m(q_m+1)\log(q_m+1)}{N_m},
\]
\[
G_m^W:=\mathfrak S_{q_m}^{\,M_m-a_m}\times \mathfrak S_{q_m+1}^{\,a_m},
\qquad
g_m^W:=\frac{1}{N_m}\log|G_m^W|.
\]
则有
\[
\operatorname{Col}(T)\ge \operatorname{Col}_m^W,
\qquad
\kappa(T)\ge \kappa_m^W,
\qquad
g(T)\ge g_m^W.
\]
对任意严格凹函数
\[
\Psi:[1,\infty)\to \RR
\]
则不等式反向。特别地，
\[
H(p):=-\sum_{i=1}^{M_m}p_i\log p_i
=\log N_m-\kappa(T)
\]
在 \(d_m^W\) 处取得唯一极大，而
\[
D_{\mathrm{KL}}(p\Vert u)
=
\kappa(T)-\log\frac{N_m}{M_m}
\]
在 \(d_m^W\) 处取得唯一极小。对 dyadic 预言机容量
\[
C_T(B):=\sum_{i=1}^{M_m}\min(d_i,2^B)
\qquad (B\ge 0)
\]
亦有
\[
C_T(B)\le
(M_m-a_m)\min(q_m,2^B)+a_m\min(q_m+1,2^B).
\]
\end{corollary}
\begin{proof}
由定理 \ref{thm:xi-time-part9ze-wulff-schur-minimal-spectrum}，
\[
d(T)^\downarrow\succ d_m^W.
\]
对严格凸 \(\Phi\)，Karamata 不等式给出
\[
\sum_i\Phi(d_i)\ge \sum_i\Phi(d_i^W),
\]
且等号当且仅当 \(d(T)\) 与 \(d_m^W\) 仅差排列；对严格凹 \(\Psi\) 则不等式反向。

取 \(\Phi(t)=t^2\) 即得碰撞率下界；取 \(\Phi(t)=t\log t\) 即得 \(\kappa\) 下界；取 \(\Phi(t)=\log\Gamma(t+1)\) 并用
\[
\log|G(T)|=\sum_i\log(d_i!)
\]
即得 \(g(T)\ge g_m^W\)。

再由
\[
H(p)= -\sum_i\frac{d_i}{N_m}\log\frac{d_i}{N_m}
=\log N_m-\kappa(T)
\]
可知，Shannon 熵的唯一极大与 \(\kappa(T)\) 的唯一极小等价。进一步，
\[
D_{\mathrm{KL}}(p\Vert u)
=\sum_i\frac{d_i}{N_m}\log\frac{d_i/N_m}{1/M_m}
=\kappa(T)-\log\frac{N_m}{M_m},
\]
故 KL 缺口的唯一极小与 \(\kappa(T)\) 的唯一极小同义。

最后，函数
\[
t\longmapsto \min(t,2^B)
\]
在 \([1,\infty)\) 上是凹函数，故同样由 Karamata 给出容量上包络。证毕。
\end{proof}

\begin{theorem}[单泛函刚性]\label{thm:xi-time-part9ze-single-functional-rigidity}
\leanverified{paper\_xi\_time\_part9ze\_single\_functional\_rigidity}
在定理 \ref{thm:xi-time-part9ze-wulff-schur-minimal-spectrum} 的设定下，下列任一条件单独成立，均已足以推出 \(d(T)\) 与 \(d_m^W\) 仅差一个排列：
\[
\operatorname{Col}(T)=\operatorname{Col}_m^W,
\]
\[
g(T)=g_m^W,
\]
\[
D_{\mathrm{KL}}(p\Vert u)=
\kappa_m^W-\log\frac{N_m}{M_m}.
\]
换言之，碰撞率极小、规范体积密度极小与相对均匀基线的 KL 缺口极小中的任一者，已单独刻画纯尺度 Wulff 射线。
\end{theorem}
\begin{proof}
第一式对应推论 \ref{cor:xi-time-part9ze-wulff-schur-convex-envelope} 中的严格凸函数 \(\Phi(t)=t^2\)；第二式对应 \(\Phi(t)=\log\Gamma(t+1)\)；第三式则因
\[
D_{\mathrm{KL}}(p\Vert u)=\kappa(T)-\log\frac{N_m}{M_m}
\]
而等价于 \(\kappa(T)=\kappa_m^W\)，也就是 \(\Phi(t)=t\log t\) 的严格凸极小情形。三种情形都由同一等号刚性推出
\[
d(T)\sim d_m^W.
\]
证毕。
\end{proof}

\begin{theorem}[低纤维异常税的一般公式]\label{thm:xi-time-part9ze-lowfiber-anomaly-tax}
\leanverified{paper\_xi\_time\_part9ze\_lowfiber\_anomaly\_tax}
假设 \(q_m\ge 2\)（等价地 \(m\ge 4\)），并取整数
\[
0\le r\le \left\lfloor \frac{M_m-a_m}{2}\right\rfloor.
\]
考虑全部满足
\[
d=(d_1,\dots,d_{M_m})\in \mathbb N^{M_m},
\qquad
\sum_{i=1}^{M_m}d_i=N_m,
\qquad
\#\{i:d_i\le q_m-1\}\ge r
\]
的向量。则对任意严格凸函数
\[
\Phi:[1,\infty)\to \RR
\]
而言，该约束下的唯一极小元为
\[
d_m^{(r)}=
\bigl((q_m-1)^r,\ q_m^{M_m-a_m-2r},\ (q_m+1)^{a_m+r}\bigr).
\]
并且其相对于 Wulff 基准的超额精确等于
\[
\sum_{i=1}^{M_m}\Phi\bigl(d_{m,i}^{(r)}\bigr)
-
\sum_{i=1}^{M_m}\Phi\bigl(d_{m,i}^{W}\bigr)
=
r\Bigl[\Phi(q_m-1)-2\Phi(q_m)+\Phi(q_m+1)\Bigr].
\]
特别地，若对任意可行 \(d\) 仍记
\[
\operatorname{Col}(d):=\frac{1}{N_m^2}\sum_i d_i^2,
\qquad
\kappa(d):=\frac{1}{N_m}\sum_i d_i\log d_i,
\]
\[
G(d):=\prod_i\mathfrak S_{d_i},
\qquad
g(d):=\frac{1}{N_m}\log|G(d)|,
\]
则有
\[
\sum_i d_i^2\ge \sum_i(d_{m,i}^{W})^2+2r,
\]
\[
\log|G(d)|\ge \log|G_m^W|+r\log\frac{q_m+1}{q_m},
\]
\[
\kappa(d)\ge
\kappa_m^W+\frac{r}{N_m}
\Bigl[(q_m-1)\log(q_m-1)-2q_m\log q_m+(q_m+1)\log(q_m+1)\Bigr].
\]
\end{theorem}
\begin{proof}
记
\[
x:=\#\{i:d_i\le q_m-1\}.
\]
由于约束要求 \(x\ge r\)，只需先固定 \(x\) 考察在“恰有 \(x\) 个低纤维”的层内极小元。

设 \(d\) 为该层上的极小元。若某个低纤维满足 \(d_i\le q_m-2\)，则由于总和
\[
N_m=q_mM_m+a_m
\]
严格大于
\[
x(q_m-1)+(M_m-x)q_m
=N_m-(a_m+x),
\]
必存在另一个坐标 \(d_j\ge q_m+1\)。于是一步平衡
\[
(d_i,d_j)\mapsto (d_i+1,d_j-1)
\]
保持“至少 \(x\) 个坐标不高于 \(q_m-1\)”这一约束，并因 \(\Phi\) 严格凸而严格降低 \(\sum_i\Phi(d_i)\)，矛盾。故层内极小元的全部低纤维都必须恰等于 \(q_m-1\)。

现在把这 \(x\) 个坐标固定为 \(q_m-1\)。剩余 \(M_m-x\) 个坐标之和为
\[
N_m-x(q_m-1)=q_m(M_m-x)+(a_m+x).
\]
因此，在剩余块上我们处于一个新的“总质量 \(q_m(M_m-x)+(a_m+x)\)、箱数 \(M_m-x\)”的平衡问题。由定理 \ref{thm:xi-time-part9ze-wulff-schur-minimal-spectrum}，其唯一 Schur 极小元正是
\[
\bigl(q_m^{M_m-a_m-2x},(q_m+1)^{a_m+x}\bigr),
\]
故在固定 \(x\) 的层内，唯一极小元为
\[
d_m^{(x)}=
\bigl((q_m-1)^x,\ q_m^{M_m-a_m-2x},\ (q_m+1)^{a_m+x}\bigr).
\]
其中可行性要求
\[
M_m-a_m-2x\ge 0
\iff
x\le \left\lfloor \frac{M_m-a_m}{2}\right\rfloor.
\]

于是只需比较单变量函数
\[
\mathcal F(x):=
x\Phi(q_m-1)+(M_m-a_m-2x)\Phi(q_m)+(a_m+x)\Phi(q_m+1).
\]
其离散增量为
\[
\mathcal F(x+1)-\mathcal F(x)
=
\Phi(q_m-1)-2\Phi(q_m)+\Phi(q_m+1)>0,
\]
故 \(\mathcal F(x)\) 严格递增，最小值在最小可行 \(x=r\) 处取得，且唯一极小元正是 \(d_m^{(r)}\)。

把
\[
\mathcal F(r)-\mathcal F(0)
\]
写开，即得相对 Wulff 基准的精确超额公式。再分别代入
\[
\Phi(t)=t^2,\qquad
\Phi(t)=\log\Gamma(t+1),\qquad
\Phi(t)=t\log t
\]
便得到平方和、\(\log|G|\) 与 \(\kappa\) 的三条下界。证毕。
\end{proof}

\begin{theorem}[window-\texorpdfstring{$6$}{6} 审计谱的变分极小性]\label{thm:xi-time-part9ze-window6-audit-variational-characterization}
\leanverified{paper\_xi\_time\_part9ze\_window6\_audit\_variational\_characterization}
记
\[
d_6^{\mathrm{audit}}:=(2^8,3^4,4^9).
\]
由推论 \ref{cor:fold-bin-m6-fiber-hist}，这正是 window-\(6\) 的 bin-fold 审计直方图。另一方面，
\[
N_6=64=3\cdot 21+1,
\qquad
M_6=F_8=21,
\]
故相应 Wulff 基准谱为
\[
d_6^W=(4^1,3^{20}).
\]
于是 \(d_6^{\mathrm{audit}}\) 是所有满足
\[
d\in \mathbb N^{21},
\qquad
\sum_{i=1}^{21}d_i=64,
\qquad
\#\{i:d_i\le 2\}\ge 8
\]
的向量中，对任意严格凸函数
\[
\Phi:[1,\infty)\to \RR
\]
的唯一极小元。特别地，由于 \(d_6^{\mathrm{audit}}\) 恰有八个二重纤维，故在更强约束
\[
\#\{i:d_i=2\}=8
\]
下，它仍是唯一极小元。
\end{theorem}
\begin{proof}
由定理 \ref{thm:xi-time-part9ze-lowfiber-anomaly-tax} 取
\[
q_6=3,\qquad a_6=1,\qquad r=8
\]
即得唯一极小元
\[
d_6^{(8)}=
\bigl(2^8,3^{21-1-16},4^{1+8}\bigr)
=(2^8,3^4,4^9).
\]
而推论 \ref{cor:fold-bin-m6-fiber-hist} 已给出 window-\(6\) 的审计直方图正是该向量，故结论成立。证毕。
\end{proof}

\begin{corollary}[window-\texorpdfstring{$6$}{6} boundary 中心的不可消除凸税]\label{cor:xi-time-part9ze-window6-boundary-center-convex-tax}
设 \(d\in \mathbb N^{21}\) 满足
\[
\sum_{i=1}^{21}d_i=64,
\]
并记
\[
G(d):=\prod_{i=1}^{21}\mathfrak S_{d_i}.
\]
若该 \(21\)-纤维模型保留 boundary sheet-parity 规范中心
\[
(\ZZ/2\ZZ)^3\hookrightarrow Z\!\bigl(G(d)\bigr),
\]
则必至少有三条二重纤维，从而
\[
\sum_i d_i^2\ge 202,
\qquad
\operatorname{Col}(d)\ge \frac{202}{64^2}=\frac{101}{2048},
\]
\[
\log|G(d)|\ge \log|G_6^W|+3\log\frac43,
\qquad
g(d)\ge g_6^W+\frac{3}{64}\log\frac43,
\]
\[
\kappa(d)\ge
\kappa_6^W+\frac{3}{64}\Bigl(2\log 2-6\log 3+4\log 4\Bigr)
=
\kappa_6^W+\frac{3}{32}\log\frac{32}{27}.
\]
因此，window-\(6\) 的 boundary 中心不是无代价异常，而携带一个严格的凸税。
\end{corollary}
\begin{proof}
由推论 \ref{cor:fold-bin-boundary-center-z2}，window-\(6\) 的 boundary 三点
\[
\texttt{100001},\ \texttt{100101},\ \texttt{101001}
\]
各自产生一个独立的中心 involution，因而保留该中心结构必要求至少三条纤维满足
\[
d_i=2.
\]
特别地，
\[
\#\{i:d_i\le 2\}\ge 3.
\]
将定理 \ref{thm:xi-time-part9ze-lowfiber-anomaly-tax} 应用于
\[
q_6=3,\qquad a_6=1,\qquad r=3
\]
即可得到三条不等式。证毕。
\end{proof}

\begin{theorem}[window-\texorpdfstring{$6$}{6} 审计谱相对于 Wulff 地板的精确异常税]\label{thm:xi-time-part9ze-window6-exact-tax-identity}
记
\[
d_6^{\mathrm{bin}}:=(2^8,3^4,4^9),
\qquad
d_6^W:=(4^1,3^{20}),
\]
\[
G_6^W:=\mathfrak S_4\times \mathfrak S_3^{20},
\qquad
g_6^W:=\frac{1}{64}\log|G_6^W|,
\]
\[
\kappa_6^W:=\frac{1}{64}\bigl(60\log 3+4\log 4\bigr),
\qquad
\operatorname{Col}_6^W:=\frac{49}{1024}.
\]
则相对于纯尺度 Wulff 基准，window-\(6\) 的实际审计谱满足精确恒等式
\[
\operatorname{Col}(d_6^{\mathrm{bin}})-\operatorname{Col}_6^W=\frac{1}{256},
\]
\[
\log|G_6^{\mathrm{bin}}|-\log|G_6^W|=8\log\frac43,
\qquad
g_6^{\mathrm{bin}}-g_6^W=\frac18\log\frac43,
\]
\[
\kappa_6-\kappa_6^W=\frac14\log\frac{32}{27}.
\]
换言之，window-\(6\) 的审计谱并非任意偏离 Wulff 地板；它恰好实现了“具有八条低纤维”这一约束下的最小可付异常税。
\end{theorem}
\begin{proof}
由定理 \ref{thm:xi-time-part9ze-window6-audit-variational-characterization}，
\[
d_6^{\mathrm{bin}}=d_6^{(8)}.
\]
故把定理 \ref{thm:xi-time-part9ze-lowfiber-anomaly-tax} 中
\[
q_6=3,\qquad a_6=1,\qquad r=8
\]
代入，即得所有精确超额。

更具体地，Wulff 基准的平方和为
\[
20\cdot 3^2+1\cdot 4^2=196,
\]
而审计谱的平方和为
\[
8\cdot 2^2+4\cdot 3^2+9\cdot 4^2=212,
\]
故
\[
\operatorname{Col}(d_6^{\mathrm{bin}})-\operatorname{Col}_6^W
=\frac{212-196}{64^2}
=\frac{16}{4096}
=\frac{1}{256}.
\]
同理，
\[
\log|G_6^{\mathrm{bin}}|-\log|G_6^W|
=8\log\frac43,
\]
两边同除以 \(64\) 即得
\[
g_6^{\mathrm{bin}}-g_6^W=\frac18\log\frac43.
\]
对 \(\kappa\) 则有
\[
\kappa_6-\kappa_6^W
=
\frac{8}{64}\Bigl(2\log 2-6\log 3+4\log 4\Bigr)
=\frac14\log\frac{32}{27}.
\]
证毕。
\end{proof}

\begin{proposition}[审计偶窗始终严格高于 Wulff--Schur 地板]\label{prop:xi-time-part9ze-audited-even-window-above-wulff-floor}
\leanverified{paper\_xi\_time\_part9ze\_audited\_even\_window\_above\_wulff\_floor}
在审计偶窗
\[
m\in\{8,10,12\}
\]
上，bin-fold 纤维直方图由表 \ref{tab:fold_bin_degeneracy_spectrum_m6_m12} 给出；记相应多重度向量为 \(d_m^{\mathrm{bin}}\)。对应的 Wulff 基准谱为
\[
d_8^W=(5^{36},4^{19}),
\qquad
d_{10}^W=(8^{16},7^{128}),
\qquad
d_{12}^W=(11^{326},10^{51}).
\]
这里约定
\[
\operatorname{Col}_m^{\mathrm{bin}}:=\operatorname{Col}(d_m^{\mathrm{bin}}),
\qquad
\operatorname{Col}_m^W:=\operatorname{Col}(d_m^W).
\]
于是碰撞率严格高于 Wulff 地板：
\[
\operatorname{Col}_8^{\mathrm{bin}}-\operatorname{Col}_8^W
=\frac{11}{8192},
\]
\[
\operatorname{Col}_{10}^{\mathrm{bin}}-\operatorname{Col}_{10}^W
=\frac{101}{262144},
\]
\[
\operatorname{Col}_{12}^{\mathrm{bin}}-\operatorname{Col}_{12}^W
=\frac{961}{8388608}.
\]
并且规范体积密度同样保留严格正缺口：
\[
g_8^{\mathrm{bin}}-g_8^W=0.0346850\ldots,
\]
\[
g_{10}^{\mathrm{bin}}-g_{10}^W=0.0268797\ldots,
\]
\[
g_{12}^{\mathrm{bin}}-g_{12}^W=0.0213913\ldots.
\]
因此，已审计到的偶窗 bin-fold 从未落到纯尺度 Wulff 射线的 Schur 极小地板；每一窗都含有额外的非尺度纤维结构。
\end{proposition}
\begin{proof}
由表 \ref{tab:fold_bin_degeneracy_spectrum_m6_m12}，
\[
d_8^{\mathrm{bin}}=(3^{21},5^{11},6^{23}),
\qquad
d_{10}^{\mathrm{bin}}=(5^{55},8^{52},9^{37}),
\qquad
d_{12}^{\mathrm{bin}}=(8^{144},12^{85},13^{148}).
\]
另一方面，
\[
256=4\cdot 55+36,
\qquad
1024=7\cdot 144+16,
\qquad
4096=10\cdot 377+326,
\]
故 Wulff 基准分别为
\[
d_8^W=(5^{36},4^{19}),
\qquad
d_{10}^W=(8^{16},7^{128}),
\qquad
d_{12}^W=(11^{326},10^{51}).
\]

直接计算平方和得到
\[
21\cdot 3^2+11\cdot 5^2+23\cdot 6^2=1292,
\qquad
19\cdot 4^2+36\cdot 5^2=1204,
\]
\[
55\cdot 5^2+52\cdot 8^2+37\cdot 9^2=7700,
\qquad
128\cdot 7^2+16\cdot 8^2=7296,
\]
\[
144\cdot 8^2+85\cdot 12^2+148\cdot 13^2=46468,
\qquad
51\cdot 10^2+326\cdot 11^2=44546.
\]
除以 \(256^2,1024^2,4096^2\) 即得三条碰撞率缺口。

再由
\[
g_m^{\mathrm{bin}}=\frac{1}{2^m}\sum_{w\in X_m}\log\bigl(d_m^{\mathrm{bin}}(w)!\bigr)
\]
与 Wulff 基准的同型定义，直接化简得到
\[
g_8^{\mathrm{bin}}-g_8^W
=
\frac{23\log 3-19\log 2-2\log 5}{256}
=0.0346850\ldots,
\]
\[
g_{10}^{\mathrm{bin}}-g_{10}^W
=
\frac{164\log 2+19\log 3-55\log 7}{1024}
=0.0268797\ldots,
\]
\[
g_{12}^{\mathrm{bin}}-g_{12}^W
=
\frac{322\log 2-55\log 3-144\log 5-93\log 11+148\log 13}{4096}
=0.0213913\ldots.
\]
因此三窗都严格高于 Wulff 地板。证毕。
\end{proof}

\begin{conjecture}[偶窗 Fibonacci 底层的 Schur 刚性]\label{conj:xi-time-part9ze-even-window-fibonacci-schur-rigidity}
由定理 \ref{thm:fold-bin-min-degeneracy-fib-audited}，在审计偶窗
\[
m\in\{6,8,10,12\}
\]
上已有
\[
d_{\min}(m)=F_{m/2},
\qquad
\bigl|\mathcal S_{m,d_{\min}(m)}\bigr|=F_m=|X_{m-2}|.
\]
这提示如下推广：对每个偶数 \(m\ge 6\)，bin-fold 的完整多重度谱应当是下述变分问题的唯一解：

在所有向量
\[
d\in \mathbb N^{M_m},
\qquad
\sum_{i=1}^{M_m}d_i=N_m,
\qquad
\#\{i:d_i=F_{m/2}\}=F_m
\]
中，使得对某个（等价地，对任意一个）严格凸可分泛函
\[
\sum_{i=1}^{M_m}\Phi(d_i)
\]
取得极小。

若该猜想成立，则审计偶窗上的 Fibonacci 底层不再只是有限样本中的经验稳定扇区，而将升级为一条全偶窗的 Schur 刚性律：最低 Fibonacci 扇区一旦被固定，完整的高层直方图便被唯一强制决定。
\end{conjecture}
\begin{proof}[证据]
window-\(6\) 的情形已由定理 \ref{thm:xi-time-part9ze-window6-audit-variational-characterization} 完成：其八条二重纤维
\[
8=F_6,\qquad 2=F_3
\]
已经足以强制整个直方图为
\[
(2^8,3^4,4^9).
\]
另一方面，定理 \ref{thm:fold-bin-min-degeneracy-fib-audited} 给出的审计偶窗数据
\[
(d_{\min},|\mathcal S_{m,d_{\min}}|)
=(2,8),\ (3,21),\ (5,55),\ (8,144)
\]
全部落在
\[
\bigl(F_{m/2},F_m\bigr)
\]
这条 Fibonacci 律上，显示最低扇区的规模与底值都具有非偶然的刚性形态。该猜想断言，这一 Fibonacci 底层不仅控制最小扇区本身，而且应继续以 Schur 极小性的方式锁定整个多重度谱。证毕。
\end{proof}

\noindent
综上，Figalli--Zhang 所给出的单参数 Wulff 几何在有限分辨率端并不只产生一条同余骨架与容量律；它还进一步诱导出一个严格的 Schur 极小地板。该地板统一控制碰撞率、KL 缺口、规范体积密度与一切严格凸可分泛函的最优下包络，而任何低纤维异常都只能以精确可算的凸税形式偏离它。window-\(6\) 的审计谱则表明，这一变分框架并非抽象的平衡原理，而可在具体审计窗口上实现为一条完全刚性的异常税恒等式：边界中心的三条二重纤维强制一个不可消除的最小代价，而实际八条低纤维又恰好达到相应约束下的唯一凸极小端点。由此，Wulff 射线、低纤维异常与有限窗审计谱被压缩进同一条可审计的主序--变分主线。

\endinput

\paragraph{window-\texorpdfstring{$6$}{6} 零温冻结的几何分裂与 Jones--熵缺陷刚性原理}
在有限指数链与中心熵损失已经固定之后，window-\(6\) 的 escort 零温极限还可继续在根几何与 boundary--interior 分裂两侧精确闭合，并进一步导向一个统一的 Jones--冻结--熵缺陷刚性图景。

\begin{theorem}[正零温极限的几何分裂]\label{thm:xi-time-part9ze-positive-zero-temp-geometric-splitting}
\leanverified{paper\_xi\_time\_part9ze\_positive\_zero\_temp\_geometric\_splitting}
定义 escort 分布
\[
\pi_{6,q}(w):=\frac{d_6(w)^q}{Z_6(q)},
\qquad w\in X_6.
\]
则
\[
\pi_{6,q}\xrightarrow[q\to+\infty]{\mathrm{TV}}
\frac{1}{9}\mathbf 1_{\mathcal S_{6,4}},
\]
即它按总变差收敛到最大壳层上的均匀分布。

进一步，沿用定理 \ref{thm:xi-time-part9r-window6-hamming-complete-weyl-orbit-invariant} 的 \(B_3\) 字典 \(\iota_B\)，记
\[
\mathcal S_{6,4}^{\mathrm{long}}
:=
\iota_B^{-1}\{\pm e_1\pm e_2\},
\]
\[
\mathcal S_{6,4}^{\mathrm{short}}
:=
\iota_B^{-1}\{e_1,\pm e_2,\pm e_3\}.
\]
则
\[
|\mathcal S_{6,4}^{\mathrm{long}}|=4,
\qquad
|\mathcal S_{6,4}^{\mathrm{short}}|=5,
\]
并且在正零温极限中有精确质量分裂
\[
\pi_{6,+\infty}\bigl(\mathcal S_{6,4}^{\mathrm{long}}\bigr)=\frac49,
\qquad
\pi_{6,+\infty}\bigl(\mathcal S_{6,4}^{\mathrm{short}}\bigr)=\frac59.
\]
\end{theorem}
\begin{proof}
由定理 \ref{thm:xi-time-part9ze-positive-zero-temperature-freezing},
\[
Z_6(q)=4^q\left(9+4\left(\frac34\right)^q+8\left(\frac12\right)^q\right).
\]
若 \(w\in \mathcal S_{6,4}\)，则 \(d_6(w)=4\)，故
\[
\pi_{6,q}(w)=\frac{4^q}{Z_6(q)}
\longrightarrow \frac19.
\]
若 \(w\notin \mathcal S_{6,4}\)，则 \(d_6(w)\in\{2,3\}\)，从而
\[
\pi_{6,q}(w)=\frac{d_6(w)^q}{Z_6(q)}\longrightarrow 0.
\]
由于 \(X_6\) 有限，这就等价于总变差收敛。

再由定理 \ref{thm:xi-window6-c3-cyclic-fiber-trisection},
\[
\iota_C(\mathcal S_{6,4})
=
\{2e_1,\pm 2e_2,\pm 2e_3,\pm e_1\pm e_2\}.
\]
命题 \ref{prop:xi-time-part9r-window6-b3-to-c3-weight-radial-dilation} 表明，\(C_3\) 与 \(B_3\) 字典只在 weight-\(1\) 轴根上相差一个因子 \(2\)，而其余非轴向根完全不变。故
\[
\iota_B(\mathcal S_{6,4})
=
\{e_1,\pm e_2,\pm e_3,\pm e_1\pm e_2\}.
\]
其中前五点是 \(B_3\) 的短根，后四点是 \(B_3\) 的长根，于是
\[
|\mathcal S_{6,4}^{\mathrm{short}}|=5,
\qquad
|\mathcal S_{6,4}^{\mathrm{long}}|=4.
\]
极限分布在 \(\mathcal S_{6,4}\) 上均匀，故两类总质量分别为 \(5/9\) 与 \(4/9\)。证毕。
\end{proof}

\begin{theorem}[负零温极限的 boundary--interior 分裂]\label{thm:xi-time-part9ze-negative-zero-temp-boundary-interior-splitting}
\leanverified{paper\_xi\_time\_part9ze\_negative\_zero\_temp\_boundary\_interior\_splitting}
仍记
\[
\pi_{6,q}(w)=\frac{d_6(w)^q}{Z_6(q)}.
\]
则
\[
\pi_{6,q}\xrightarrow[q\to-\infty]{\mathrm{TV}}
\frac{1}{8}\mathbf 1_{\mathcal S_{6,2}},
\]
即它按总变差收敛到最小壳层上的均匀分布。

再记
\[
X_{6,\mathrm{int}}^{(2)}:=X_{6,\mathrm{cyc},2}.
\]
由定理 \ref{thm:xi-window6-center-exact-splitting-bdry-cyclic},
\[
\mathcal S_{6,2}
=
X_6^{\mathrm{bdry}}\sqcup X_{6,\mathrm{int}}^{(2)},
\qquad
|X_6^{\mathrm{bdry}}|=3,
\qquad
|X_{6,\mathrm{int}}^{(2)}|=5.
\]
因此负零温极限中的质量精确分裂为
\[
\pi_{6,-\infty}\bigl(X_6^{\mathrm{bdry}}\bigr)=\frac38,
\qquad
\pi_{6,-\infty}\bigl(X_{6,\mathrm{int}}^{(2)}\bigr)=\frac58.
\]
换言之，boundary 中心因子对应的三条二点纤维只在冷基态极限中被热冻结显式选中。
\end{theorem}
\begin{proof}
由定理 \ref{thm:xi-time-part9ze-negative-zero-temperature-freezing},
\[
Z_6(q)=2^q\left(8+4\left(\frac32\right)^q+9\cdot 2^q\right).
\]
若 \(w\in \mathcal S_{6,2}\)，则 \(d_6(w)=2\)，故
\[
\pi_{6,q}(w)=\frac{2^q}{Z_6(q)}\longrightarrow \frac18.
\]
若 \(w\notin \mathcal S_{6,2}\)，则 \(d_6(w)\in\{3,4\}\)，从而
\[
\pi_{6,q}(w)\longrightarrow 0.
\]
由于定义域有限，这就等价于总变差收敛。

再由定理 \ref{thm:xi-window6-center-exact-splitting-bdry-cyclic},
\[
\mathcal S_{6,2}=X_6^{\mathrm{bdry}}\sqcup X_{6,\mathrm{cyc},2},
\qquad
|X_6^{\mathrm{bdry}}|=3,
\qquad
|X_{6,\mathrm{cyc},2}|=5.
\]
把 \(X_{6,\mathrm{int}}^{(2)}:=X_{6,\mathrm{cyc},2}\) 代入即可得到
\[
\pi_{6,-\infty}(X_6^{\mathrm{bdry}})=\frac38,
\qquad
\pi_{6,-\infty}(X_{6,\mathrm{int}}^{(2)})=\frac58.
\]
证毕。
\end{proof}

\begin{conjecture}[Jones--冻结--熵缺陷刚性原理]\label{conj:xi-time-part9ze-jones-freeze-entropy-rigidity}
对任意偶数窗口 \(m\)，设
\[
\mathcal A_m=\bigoplus_{x\in X_m}M_{d_m(x)}(\CC)
\]
为其 fold algebra，并假设 cyclic sector 满足猜想
\ref{conj:xi-time-part9zd-window6-cartan-lanczos-rigidity-principle}
中的 rigid Lie completion 与 Cartan--Lanczos 假设。则应同时成立：
\[
\operatorname{Ind}_{\mathrm{PP}}(E_m^Z)=d_{\max}(m)=\rho(J_m),
\]
\[
\lim_{q\to+\infty}\frac1q\log\sum_{x\in X_m}d_m(x)^q
=
\log\operatorname{Ind}_{\mathrm{PP}}(E_m^Z),
\]
\[
S(\tau_m)-S(\tau_m|_{Z(\mathcal A_m)})=\kappa_m,
\]
并且正、负零温极限的基态壳层分别恰等于中心条件期望 \(E_m^Z\) 与加权中心期望 \(\varepsilon_m\) 的 Pimsner--Popa 障碍层
\[
\{x\in X_m:\ d_m(x)=d_{\max}(m)\},
\qquad
\{x\in X_m:\ d_m(x)=d_{\min}(m)\}.
\]

window-\(6\) 给出了该原理的首个严格实例：其中心指数、加权中心指数与全地址标量指数分别为
\[
4,\qquad 32,\qquad 64,
\]
其正零温基态壳层为九点最大壳层，负零温基态壳层为八点最小壳层，而其中心熵损失恰等于 \(\kappa_6\)。
\end{conjecture}

\noindent
以上结果把 window-\(6\) 的群胚代数进一步压缩为一条严格的有限指数--冻结--熵缺陷链：中心条件期望给出指数 \(4\)，其 dyadic 对数恢复精确反演的两比特预算；加权中心标量期望给出指数 \(32\)，其 dyadic 对数刻画稳定类型的最坏情形信息代价；全代数标量条件期望给出指数 \(64\)，其 dyadic 对数正好恢复全部微态地址长度。与此同时，正零温与负零温极限分别冻结到九点最大壳层与八点最小壳层，并在 \(B_3/C_3\) 根几何与 boundary/cyclic 二点分裂下继续细化。最终，\(\kappa_6\) 被识别为 canonical microstate trace 向中心坍缩时的精确 von Neumann 熵损失，从而有限指数、热冻结、根几何、boundary parity 与熵缺陷在同一组条件期望公式下被锁定为同一对象的五个截面。

\endinput

\paragraph{window-\texorpdfstring{$6$}{6} 的 gauge/groupoid 半单闭包与信息谱单点坍缩}
在前述 Wulff 射线的 Schur 极小地板已经把纤维多重度的凸变分端点固定之后，window-\(6\) 的 bin-fold 纤维侧还可继续沿群、群胚与信息谱三层闭合。其要点在于：纤维直方图 \(2{:}8,\ 3{:}4,\ 4{:}9\) 不仅决定规范群的直积分解，也同步决定纤维群胚代数的 Wedderburn 块结构；另一方面，主阶信息密度与全族 Rényi 熵率又同时塌缩到同一指数 \(\log\varphi\)。因此，window-\(6\) 的离散纤维侧并不只是若干局部不变量的并列堆叠，而是一个在有限群论、半单代数与信息论三端完全闭合的单一对象层。

\begin{theorem}[window-\texorpdfstring{$6$}{6} bin-gauge 群的完全群论剖分]\label{thm:xi-time-part9zf-window6-bingauge-complete-group-dissection}
\leanverified{paper\_xi\_time\_part9zf\_window6\_bingauge\_complete\_group\_dissection}
记
\[
G_6^{\mathrm{bin}}:=\Aut(\Fold_6^{\mathrm{bin}}).
\]
则有群同构
\[
G_6^{\mathrm{bin}}\cong S_2^{\,8}\times S_3^{\,4}\times S_4^{\,9}.
\]
从而
\[
Z\!\bigl(G_6^{\mathrm{bin}}\bigr)\cong (\ZZ/2\ZZ)^8,
\qquad
\bigl[G_6^{\mathrm{bin}},G_6^{\mathrm{bin}}\bigr]\cong A_3^{\,4}\times A_4^{\,9},
\qquad
\bigl(G_6^{\mathrm{bin}}\bigr)^{\mathrm{ab}}\cong (\ZZ/2\ZZ)^{21}.
\]
再记
\[
X_6^{\mathrm{bdry}}=
\{\texttt{100001},\texttt{100101},\texttt{101001}\},
\]
\[
\mathcal I_6:=
\{\texttt{000001},\texttt{000101},\texttt{001001},\texttt{010001},\texttt{010101}\}.
\]
则 \(\mathcal I_6\) 恰为 window-\(6\) 的 interior 二点纤维集合，并且中心具有规范直和分解
\[
Z\!\bigl(G_6^{\mathrm{bin}}\bigr)
\cong
(\ZZ/2\ZZ)^{X_6^{\mathrm{bdry}}}
\oplus
(\ZZ/2\ZZ)^{\mathcal I_6}
\cong
(\ZZ/2\ZZ)^3\oplus(\ZZ/2\ZZ)^5.
\]
\end{theorem}
\begin{proof}
由推论 \ref{cor:fold-bin-gauge-m6}，
\[
G_6^{\mathrm{bin}}\cong S_2^{\,8}\times S_3^{\,4}\times S_4^{\,9},
\qquad
\bigl(G_6^{\mathrm{bin}}\bigr)^{\mathrm{ab}}\cong (\ZZ/2\ZZ)^{21}.
\]
这已给出第一式与阿贝尔化结论。

再由
\[
Z(S_2)=S_2\cong \ZZ/2\ZZ,
\qquad
Z(S_d)=1\quad(d\ge 3),
\]
得到
\[
Z\!\bigl(G_6^{\mathrm{bin}}\bigr)\cong (\ZZ/2\ZZ)^8.
\]
另一方面，
\[
[S_2,S_2]=1,
\qquad
[S_3,S_3]=A_3,
\qquad
[S_4,S_4]=A_4,
\]
而直积群的交换子群按因子逐项分裂，故
\[
\bigl[G_6^{\mathrm{bin}},G_6^{\mathrm{bin}}\bigr]
\cong
A_3^{\,4}\times A_4^{\,9}.
\]

最后，由推论 \ref{cor:killo-foldbin6-center-boundary-interior}，
\[
X_6^{(2),\mathrm{int}}
=
\{\texttt{000001},\texttt{000101},\texttt{001001},\texttt{010001},\texttt{010101}\},
\]
且
\[
Z\!\bigl(G_6^{\mathrm{bin}}\bigr)
\cong
(\ZZ/2\ZZ)^{X_6^{\mathrm{bdry}}}\oplus(\ZZ/2\ZZ)^{X_6^{(2),\mathrm{int}}}.
\]
取 \(\mathcal I_6:=X_6^{(2),\mathrm{int}}\) 即得所述规范直和分解。证毕。
\end{proof}

\begin{theorem}[window-\texorpdfstring{$6$}{6} 纤维群胚代数的 Wedderburn 光谱闭式]\label{thm:xi-time-part9zf-window6-groupoid-wedderburn-closedform}
设 \(\mathcal R_6^{\mathrm{bin}}\) 为定义 \ref{def:fold-bin-fiber-groupoid} 的 bin-fold 纤维等价关系群胚。则其复群胚代数满足
\[
\CC[\mathcal R_6^{\mathrm{bin}}]
\cong
M_2(\CC)^{\oplus 8}\oplus M_3(\CC)^{\oplus 4}\oplus M_4(\CC)^{\oplus 9}.
\]
因此
\[
\dim_{\CC}\CC[\mathcal R_6^{\mathrm{bin}}]
=
8\cdot 2^2+4\cdot 3^2+9\cdot 4^2
=212,
\]
\[
\dim_{\CC}Z\!\bigl(\CC[\mathcal R_6^{\mathrm{bin}}]\bigr)=8+4+9=21,
\qquad
K_0\!\bigl(\CC[\mathcal R_6^{\mathrm{bin}}]\bigr)\cong \ZZ^{21}.
\]
特别地，其简单模维数分布可由单一生成多项式完全封装：
\[
\mathcal P_{\mathcal R_6}(t):=
\sum_{V\in\operatorname{Irr}(\CC[\mathcal R_6^{\mathrm{bin}}])}t^{\dim V}
=8t^2+4t^3+9t^4.
\]
\end{theorem}
\begin{proof}
由命题 \ref{prop:fold-bin-groupoid-algebra-wedderburn}，
\[
\CC[\mathcal R_6^{\mathrm{bin}}]
\cong
\bigoplus_{w\in X_6}M_{d_6^{\mathrm{bin}}(w)}(\CC).
\]
再由推论 \ref{cor:fold-bin-m6-fiber-hist}，
\[
\#\{d_6^{\mathrm{bin}}=2\}=8,\qquad
\#\{d_6^{\mathrm{bin}}=3\}=4,\qquad
\#\{d_6^{\mathrm{bin}}=4\}=9,
\]
故立得
\[
\CC[\mathcal R_6^{\mathrm{bin}}]
\cong
M_2(\CC)^{\oplus 8}\oplus M_3(\CC)^{\oplus 4}\oplus M_4(\CC)^{\oplus 9}.
\]

有限维半单复代数的总维数等于各 Wedderburn 块维数平方和，故
\[
\dim_{\CC}\CC[\mathcal R_6^{\mathrm{bin}}]
=8\cdot 4+4\cdot 9+9\cdot 16
=212.
\]
其中心恰由各矩阵块的标量中心直和给出，故中心维数等于块数
\[
8+4+9=21.
\]
又因
\[
K_0\!\left(\bigoplus_{j=1}^{21}M_{n_j}(\CC)\right)\cong \ZZ^{21},
\]
可得
\[
K_0\!\bigl(\CC[\mathcal R_6^{\mathrm{bin}}]\bigr)\cong \ZZ^{21}.
\]
最后，\(\operatorname{Irr}(\CC[\mathcal R_6^{\mathrm{bin}}])\) 的简单模维数正是这些矩阵块的尺寸多重集，故
\[
\mathcal P_{\mathcal R_6}(t)=8t^2+4t^3+9t^4.
\]
证毕。
\end{proof}

\begin{theorem}[bin-fold 信息维谱的单点坍缩与高阶矩收缩]\label{thm:xi-time-part9zf-binfold-information-spectrum-moment-collapse}
令
\[
\mu_m(w):=\frac{d_m^{\mathrm{bin}}(w)}{2^m},
\qquad
w\in X_m,
\]
并记 \(W_m\sim \mu_m\)。则成立：
\begin{enumerate}
  \item 对一切 \(q>0\)，Rényi 熵率满足统一极限
  \[
  \lim_{m\to\infty}\frac{1}{m}H_q(\mu_m)=\log\varphi;
  \]
  \item 点态信息密度具有一致收缩
  \[
  -\frac{1}{m}\log \mu_m(W_m)=\log\varphi+O(m^{-1});
  \]
  \item 因而对任意 \(\varepsilon>0\)，当 \(m\) 充分大时，
  \[
  \mu_m\!\left(
  \left\{
  w\in X_m:
  \left|
  -\frac{1}{m}\log \mu_m(w)-\log\varphi
  \right|>\varepsilon
  \right\}
  \right)=0;
  \]
  \item 更一般地，对每个 \(r\ge 1\)，有
  \[
  \EE_{\mu_m}\left|
  -\frac{1}{m}\log \mu_m(W_m)-\log\varphi
  \right|^r
  =O(m^{-r}).
  \]
\end{enumerate}
换言之，bin-fold 的离散信息谱并不承载非平凡的多重信息指数层次；其全部质量在主阶上塌缩到唯一指数 \(\log\varphi\)。
\end{theorem}
\begin{proof}
定理 \ref{thm:xi-foldbin-information-spectrum-single-point-collapse} 已经给出两条关键事实：
\[
-\frac{1}{m}\log \mu_m(W_m)=\log\varphi+O(m^{-1}),
\]
其中误差对 \(W_m\) 一致成立；并且对全部 \(q>0\) 都有
\[
\frac{1}{m}H_q(\mu_m)\longrightarrow \log\varphi.
\]
这就得到前两条。

现取常数 \(C>0\) 使得对全部 \(m\) 与全部 \(w\in X_m\) 都有
\[
\left|
-\frac{1}{m}\log \mu_m(w)-\log\varphi
\right|
\le \frac{C}{m}.
\]
若给定 \(\varepsilon>0\)，只要 \(m>C/\varepsilon\)，则上式右端小于 \(\varepsilon\)，故异常集合为空集，从而第三条成立。

第四条同样由这一逐点界直接推出：
\[
\EE_{\mu_m}\left|
-\frac{1}{m}\log \mu_m(W_m)-\log\varphi
\right|^r
\le
\left(\frac{C}{m}\right)^r
=O(m^{-r}).
\]
证毕。
\end{proof}

\noindent
由此，window-\(6\) 的纤维侧在三重口径下达到严格闭合：在群论上，全部规范自由度被 \(S_2^{8}\times S_3^{4}\times S_4^{9}\) 的直积完全穷尽，并进一步裂解为 boundary 三维与 interior 五维的中心直和；在代数上，同一纤维谱又把群胚代数压缩为 \(21\) 个矩阵块的有限维半单对象；而在信息论上，所有点态与 Rényi 级别的主阶指数最终都塌缩到唯一常数 \(\log\varphi\)。因此，window-\(6\) 的 gauge/groupoid 结构与其信息谱并不是彼此平行的三套语言，而是由同一纤维多重度直方图唯一统摄的三个同构投影。

\endinput

\paragraph{有限局部化的 Pontryagin 刚性与 Chebotarev 指纹的零信息耦合}
与前述纤维侧的 gauge/groupoid 闭链并行，算术侧也出现一条同样刚性的双层结构。一方面，有限素数局部化 \(G_S=\ZZ[S^{-1}]\) 在 Pontryagin 对偶后并不遗失素数账本；相反，全部局部化信息被严格编码为一个唯一的圆连通分量与一个可完全恢复的 \(p\)-adic component 账本。另一方面，当 \(L_{10}\) 与 Lee--Yang 共同三次数域的算术外置性、以及 \(L_{10}\)、\(L_{\mathrm{LY}}\)、\(L_4\) 的直积 Galois 结构已经固定后，各域 Frobenius 指纹之间的耦合也彻底退化为乘积分布，从而在信息论意义下给出严格的零互信息律。由此，prime-ledger 的对偶刚性与多域指纹的统计独立性在同一算术层上闭合。

\begin{theorem}[有限素数局部化的 Pontryagin 刚性分类]\label{thm:xi-time-part9zg-finite-localization-pontryagin-rigidity}
\leanverified{paper\_xi\_time\_part9zg\_finite\_localization\_pontryagin\_rigidity}
对任意有限素数集 \(S\subset \mathbb P\)，记
\[
G_S:=\ZZ[S^{-1}],
\qquad
\Sigma_S:=\widehat{G_S}.
\]
则：
\begin{enumerate}
  \item \(\Sigma_S\) 的单位连通分量与分量群满足
  \[
  \Sigma_S^\circ\cong \TT,
  \qquad
  \pi_0(\Sigma_S):=\Sigma_S/\Sigma_S^\circ\cong \prod_{p\in S}\ZZ_p;
  \]
  \item 对任意有限 \(S,T\subset\mathbb P\)，下列三条件等价：
  \[
  \Sigma_S\cong \Sigma_T
  \quad\text{（作为紧拓扑群）},
  \]
  \[
  G_S\cong G_T
  \quad\text{（作为离散交换群）},
  \]
  \[
  S=T.
  \]
\end{enumerate}
换言之，有限素数局部化的全部数据并未在对偶化中丢失，而被精确编码为“圆分量 \(+\) \(p\)-adic 分量群”的二层结构。
\end{theorem}
\begin{proof}
由定理 \ref{thm:killo-localization-circle-dimension-prime-ledger-splitting}，
\[
0\longrightarrow K_S:=\prod_{p\in S}\ZZ_p
\longrightarrow
\Sigma_S
\xrightarrow{q}
\TT
\longrightarrow 0
\]
为一条紧交换群短正合列。记 \(C:=\Sigma_S^\circ\)。由于 \(K_S\) 全不连通，故
\[
C\cap K_S=1.
\]
另一方面，\(q(C)\) 是 \(\TT\) 的连通闭子群，故只可能是 \(\{1\}\) 或 \(\TT\)。若 \(q(C)=\{1\}\)，则 \(q\) 经过商群
\[
\Sigma_S/C=\pi_0(\Sigma_S)
\]
因子化；但 \(\pi_0(\Sigma_S)\) 全不连通，而全不连通紧群的连续像仍全不连通，不可能连续满射到连通且非平凡的 \(\TT\)。故必有 \(q(C)=\TT\)。于是
\[
q|_C:C\longrightarrow \TT
\]
是连续双射；因 \(C\) 紧而 \(\TT\) Hausdorff，它实际上是拓扑群同构。令
\[
s:=(q|_C)^{-1}:\TT\to C.
\]
对任意 \(x\in \Sigma_S\)，有
\[
x-\!s(q(x))\in K_S,
\]
故 \(x\) 唯一写成 \(K_S\) 与 \(C\) 中元素之和。于是
\[
\Sigma_S\cong K_S\times C\cong \left(\prod_{p\in S}\ZZ_p\right)\times \TT,
\]
从而
\[
\pi_0(\Sigma_S)=\Sigma_S/C\cong K_S\cong \prod_{p\in S}\ZZ_p.
\]
这证明了第 \(1\) 条。

若 \(\Sigma_S\cong \Sigma_T\)，则 Pontryagin 对偶给出
\[
G_S\cong G_T.
\]
反向同样成立。再由推论 \ref{cor:cdim-localized-gs-dual-complete-classification}，
\[
\Sigma_S\cong \Sigma_T
\qquad\Longleftrightarrow\qquad
S=T.
\]
于是第 \(2\) 条三者等价。证毕。
\end{proof}

\begin{theorem}[有限局部化格与紧对偶商格的反同构]\label{thm:xi-time-part9zg-finite-localization-dual-quotient-lattice}
对有限素数集 \(S,T\subset\mathbb P\)，下列条件等价：
\[
T\subseteq S;
\]
\[
\exists\ \text{群单射 }G_T\hookrightarrow G_S;
\]
\[
\exists\ \text{连续满同态 }\Sigma_S\twoheadrightarrow \Sigma_T.
\]
因此，有限局部化的包含关系并非仅是集合论偏序；它恰与紧对偶侧的连续商结构构成反同构。
\end{theorem}
\begin{proof}
若 \(T\subseteq S\)，则
\[
G_T=\ZZ[T^{-1}]\subseteq \ZZ[S^{-1}]=G_S
\]
为自然包含，故存在单射 \(G_T\hookrightarrow G_S\)。

反之，设 \(\phi:G_T\hookrightarrow G_S\) 为群单射。任取 \(p\in T\)。由于 \(1/p\in G_T\)，乘法映射
\[
[p]:G_T\to G_T,\qquad x\mapsto px
\]
是满射，故 \(G_T\) 是 \(p\)-可除群，从而 \(\phi(G_T)\) 也是 \(G_S\) 的 \(p\)-可除子群。

现断言：若 \(p\notin S\)，则 \(G_S\) 中不存在非零 \(p\)-可除元。事实上，任取
\[
x=\frac{a}{d}\in G_S\setminus\{0\},
\]
其中 \(d\) 的全部素因子都属于 \(S\)。若 \(x\) 是 \(p\)-可除的，则对每个 \(n\ge 1\) 都存在 \(y_n\in G_S\) 使
\[
p^ny_n=x.
\]
写 \(y_n=b_n/d_n\) 且 \(d_n\) 只含 \(S\)-素因子，则
\[
p^n b_n d=a d_n.
\]
右端的 \(p\)-进赋值恒等于 \(v_p(a)\)，而左端至少为 \(n\)。当 \(n>v_p(a)\) 时矛盾。故 \(p\notin S\) 时 \(G_S\) 无非零 \(p\)-可除子群。

于是 \(\phi(G_T)\neq 0\) 迫使 \(p\in S\)。因 \(p\in T\) 任意，得
\[
T\subseteq S.
\]

最后，Pontryagin 对偶把离散交换群单射送为紧交换群连续满射，反之亦然，故第二条与第三条等价。证毕。
\end{proof}

\begin{theorem}[十次谱分歧塔与 Lee--Yang 共同三次数域的线性互素性]\label{thm:xi-time-part9zg-p10-leyang-cubic-linear-disjointness}
沿用命题 \ref{prop:fold-gauge-anomaly-P10-galois-discriminant}、定理 \ref{thm:xi-terminal-zm-stokes-leyang-shared-artin-representation} 与推论 \ref{cor:xi-terminal-zm-stokes-leyang-common-quadratic-resolvent} 的记号。记
\[
K_*:=\QQ(u)=\QQ(b)=\QQ(\lambda_*),
\]
并令 \(L_*\) 为三次数域 \(K_*/\QQ\) 的伽罗瓦闭包，\(L_{10}\) 为十次分歧多项式 \(P_{10}\) 的分裂域。则
\[
\Gal(L_*/\QQ)\cong S_3,
\qquad
L_{10}\cap L_*=\QQ,
\qquad
\Gal(L_{10}L_*/\QQ)\cong S_{10}\times S_3.
\]
等价地，Lee--Yang 分支系数 \(u\)、振幅平方 \(b=B^2\) 与分歧谱值 \(\lambda_*\) 所生成的共同三次数域在算术上严格外置于十次谱分歧塔之外。
\end{theorem}
\begin{proof}
由定理 \ref{thm:xi-terminal-zm-stokes-leyang-shared-artin-representation}，
\[
K_*=\QQ(u)=\QQ(b)=\QQ(\lambda_*),
\]
并且 \(K_*/\QQ\) 为非伽罗瓦三次扩张，其伽罗瓦闭包满足
\[
\Gal(L_*/\QQ)\cong S_3.
\]
再由推论 \ref{cor:xi-terminal-zm-stokes-leyang-common-quadratic-resolvent}，\(L_*/\QQ\) 的唯一非平凡二次子域为
\[
\QQ(\sqrt{-111}).
\]

另一方面，由命题 \ref{prop:fold-gauge-anomaly-P10-galois-discriminant}，
\[
\Gal(L_{10}/\QQ)\cong S_{10},
\]
且 \(L_{10}/\QQ\) 的唯一非平凡二次子域为
\[
L_{10}^{A_{10}}=\QQ(\sqrt{66269989}).
\]
由于 \(L_{10}/\QQ\) 与 \(L_*/\QQ\) 都是伽罗瓦扩张，交域
\[
F:=L_{10}\cap L_*
\]
亦为 \(\QQ\) 上伽罗瓦扩张。作为 \(L_{10}\) 的伽罗瓦子扩张，\(F\) 只能是 \(\QQ\) 或 \(\QQ(\sqrt{66269989})\)；作为 \(L_*\) 的伽罗瓦子扩张，\(F\) 只能是 \(\QQ\) 或 \(\QQ(\sqrt{-111})\)。但
\[
\QQ(\sqrt{66269989})\neq \QQ(\sqrt{-111}),
\]
故
\[
L_{10}\cap L_*=\QQ.
\]
于是两扩张线性互素，并得到
\[
\Gal(L_{10}L_*/\QQ)\cong S_{10}\times S_3.
\]
证毕。
\end{proof}

\begin{corollary}[十次谱分歧塔上 Lee--Yang 三次分支不变量的不可约保持]\label{cor:xi-time-part9zg-p10-branch-cubic-irreducibility}
在定理 \ref{thm:xi-time-part9zg-p10-leyang-cubic-linear-disjointness} 的设定下，
\[
L_{10}\cap \QQ(u)=
L_{10}\cap \QQ(b)=
L_{10}\cap \QQ(\lambda_*)=
\QQ.
\]
因此
\[
[L_{10}\QQ(u):L_{10}]
=
[L_{10}\QQ(b):L_{10}]
=
[L_{10}\QQ(\lambda_*):L_{10}]
=3.
\]
从而 \(u\)、\(b\) 与 \(\lambda_*\) 在 \(L_{10}\) 上都保持三次不可约。
\end{corollary}
\begin{proof}
由定理 \ref{thm:xi-time-part9zg-p10-leyang-cubic-linear-disjointness}，
\[
L_{10}\cap L_*=\QQ.
\]
而
\[
\QQ(u),\ \QQ(b),\ \QQ(\lambda_*)\subseteq L_*,
\]
故它们与 \(L_{10}\) 的交也只能是 \(\QQ\)。再由定理 \ref{thm:xi-terminal-zm-kappa-square-cubic-field-s3} 与定理 \ref{thm:xi-terminal-zm-stokes-leyang-shared-artin-representation}，
\[
[\QQ(u):\QQ]=[\QQ(b):\QQ]=[\QQ(\lambda_*):\QQ]=3.
\]
于是塔式公式给出
\[
[L_{10}\QQ(u):L_{10}]=3,
\qquad
[L_{10}\QQ(b):L_{10}]=3,
\qquad
[L_{10}\QQ(\lambda_*):L_{10}]=3.
\]
这恰表明三者在 \(L_{10}\) 上仍为三次不可约。证毕。
\end{proof}

\begin{theorem}[Frobenius 类函数的张量期望分裂]\label{thm:xi-time-part9zg-frobenius-classfunction-tensor-factorization}
沿用命题 \ref{prop:fold-gauge-anomaly-P10-leyang-linear-disjointness} 与定理 \ref{thm:killo-leyang-lattes-triple-product-galois} 的记号。对共同不分歧素数 \(p\)，记
\[
\sigma_{10}(p)\in S_{10},
\qquad
\sigma_3(p)\in S_3,
\qquad
\tau_4(p):=\bigl(\sigma_4(p),\varepsilon_4(p)\bigr)\in S_4\times C_2.
\]
把这些 Frobenius 指纹按自然密度视为随机变量，并以 \(\langle\cdot\rangle\) 记其相应期望。

则对任意类函数
\[
f:S_{10}\to\CC,
\qquad
g:S_3\to\CC,
\]
记群平均
\[
\langle f\rangle_{S_{10}}:=\frac{1}{10!}\sum_{\tau\in S_{10}}f(\tau),
\qquad
\langle g\rangle_{S_3}:=\frac{1}{6}\sum_{\eta\in S_3}g(\eta),
\]
都有精确分裂
\[
\left\langle f(\sigma_{10})\,g(\sigma_3)\right\rangle
=
\langle f\rangle_{S_{10}}\,
\langle g\rangle_{S_3}.
\]
等价地，
\[
\left\langle
\bigl(f-\langle f\rangle_{S_{10}}\bigr)(\sigma_{10})\,
\bigl(g-\langle g\rangle_{S_3}\bigr)(\sigma_3)
\right\rangle
=0.
\]

若再取任意类函数
\[
h:S_4\times C_2\to\CC,
\]
则三重张量分裂成立：
\[
\left\langle
f(\sigma_{10})\,g(\sigma_3)\,h(\tau_4)
\right\rangle
=
\langle f\rangle_{S_{10}}\,
\langle g\rangle_{S_3}\,
\langle h\rangle_{S_4\times C_2},
\]
其中
\[
\langle h\rangle_{S_4\times C_2}:=
\frac{1}{48}\sum_{\xi\in S_4\times C_2}h(\xi).
\]
\end{theorem}
\begin{proof}
两域情形正是定理 \ref{thm:fold-gauge-anomaly-P10-leyang-class-function-tensor-independence} 的内容，其结论表明共同不分歧素数上的极限分布为
\[
\mu_{10}\otimes \mu_3,
\]
其中 \(\mu_{10},\mu_3\) 分别为 \(S_{10}\) 与 \(S_3\) 上的均匀共轭类分布，故第一条公式成立。

三域情形由定理 \ref{thm:xi-time-part62d-triple-frobenius-classfunction-tensor-independence} 直接给出。该定理断言三元 Frobenius 观测的极限分布为
\[
\mu_{10}\otimes \mu_3\otimes \mu_4,
\]
其中 \(\mu_4\) 是 \(S_4\times C_2\) 上的均匀共轭类分布，从而得到三重张量分裂。

中心化后的等式只需把 \(f,g\) 或 \(f,g,h\) 分别替换为去均值类函数即可，因为它们的群平均为零。证毕。
\end{proof}

\begin{theorem}[十次谱分歧域与 Lee--Yang 三次域的联合分拆型密度完全显式]\label{thm:xi-time-part9zg-p10-leyang-joint-splitting-partition-density-table}
沿用定理 \ref{thm:xi-time-part9zg-p10-leyang-cubic-linear-disjointness} 与定理 \ref{thm:xi-time-part9zg-frobenius-classfunction-tensor-factorization} 的记号。设
\[
P_{\mathrm{LY}}(y):=256y^3+411y^2+165y+32.
\]
对任意分拆
\[
\lambda=1^{m_1}2^{m_2}\cdots\vdash 10,
\qquad
\mu=1^{n_1}2^{n_2}3^{n_3}\vdash 3,
\]
定义
\[
z_\lambda:=\prod_{j\ge 1}j^{m_j}m_j!,
\qquad
z_\mu:=\prod_{j\ge 1}j^{n_j}n_j!.
\]
则对共同不分歧素数 \(p\)，同时满足
\[
P_{10}\bmod p\ \text{的分解型为 }\lambda,
\qquad
P_{\mathrm{LY}}\bmod p\ \text{的分解型为 }\mu
\]
的素数集合具有自然密度
\[
\boxed{\;
\delta(\lambda,\mu)=\frac{1}{z_\lambda z_\mu}.\;}
\]
\end{theorem}
\begin{proof}
记
\[
C_\lambda\subset S_{10},
\qquad
D_\mu\subset S_3
\]
分别为循环型为 \(\lambda\) 与 \(\mu\) 的共轭类。由对称群的标准共轭类大小公式，
\[
|C_\lambda|=\frac{10!}{z_\lambda},
\qquad
|D_\mu|=\frac{3!}{z_\mu}.
\]
另一方面，由定理 \ref{thm:xi-time-part9zg-frobenius-classfunction-tensor-factorization}，取类函数
\[
f=\mathbf 1_{C_\lambda},
\qquad
g=\mathbf 1_{D_\mu},
\]
便知联合密度等于
\[
\frac{|C_\lambda|}{10!}\cdot \frac{|D_\mu|}{3!}.
\]
代入上式即得
\[
\delta(\lambda,\mu)=\frac{1}{z_\lambda z_\mu}.
\]
又对不分歧素数，Frobenius 循环型与模素分解型一致，故结论成立。证毕。
\end{proof}

\begin{corollary}[条件于 \texorpdfstring{$P_{10}\bmod p$}{P10 mod p} 不可约后，Lee--Yang 三态分布仍保持严格独立]\label{cor:xi-time-part9zg-p10-irreducible-leyang-tristate-independence}
沿用定理 \ref{thm:xi-time-part9zg-p10-leyang-joint-splitting-partition-density-table} 的记号。记事件
\[
\mathcal I_{10}:=\{\text{\(P_{10}\bmod p\) 不可约}\},
\]
\[
\mathcal S:=\{\text{\(P_{\mathrm{LY}}\bmod p\) 完全分裂}\},
\qquad
\mathcal M:=\{\text{\(P_{\mathrm{LY}}\bmod p\) 分解为 \(1+2\)}\},
\qquad
\mathcal J:=\{\text{\(P_{\mathrm{LY}}\bmod p\) 保持三次不可约}\}.
\]
则
\[
\delta(\mathcal I_{10}\cap \mathcal S)=\frac{1}{60},
\qquad
\delta(\mathcal I_{10}\cap \mathcal M)=\frac{1}{20},
\qquad
\delta(\mathcal I_{10}\cap \mathcal J)=\frac{1}{30}.
\]
特别地，条件于 \(P_{10}\bmod p\) 不可约之后，\(P_{\mathrm{LY}}\bmod p\) 的三态分布仍为
\[
\boxed{\;
\mathbb P(\mathcal S\mid \mathcal I_{10})=\frac16,
\qquad
\mathbb P(\mathcal M\mid \mathcal I_{10})=\frac12,
\qquad
\mathbb P(\mathcal J\mid \mathcal I_{10})=\frac13.\;}
\]
\end{corollary}
\begin{proof}
由定理 \ref{thm:xi-time-part9zg-p10-leyang-joint-splitting-partition-density-table}，事件 \(\mathcal I_{10}\) 对应分拆 \((10)\)，其
\[
z_{(10)}=10.
\]
而三次分解型 \((1+1+1)\)、\((1+2)\)、\((3)\) 分别满足
\[
z_{(1+1+1)}=3!,
\qquad
z_{(1+2)}=2,
\qquad
z_{(3)}=3.
\]
故联合密度分别为
\[
\frac{1}{10\cdot 6}=\frac{1}{60},
\qquad
\frac{1}{10\cdot 2}=\frac{1}{20},
\qquad
\frac{1}{10\cdot 3}=\frac{1}{30}.
\]
又因为
\[
\delta(\mathcal I_{10})=\frac{1}{10},
\]
把上述三式同时除以 \(1/10\) 即得条件分布
\[
\left(\frac16,\frac12,\frac13\right).
\]
这说明即便在“十次不可约”这一最粗条件下，Lee--Yang 三态分布仍与其边际分布完全一致。证毕。
\end{proof}

\begin{corollary}[算术指纹的零互信息律]\label{cor:xi-time-part9zg-arithmetic-fingerprint-zero-mutual-information}
在定理 \ref{thm:xi-time-part9zg-frobenius-classfunction-tensor-factorization} 的自然密度极限测度下，记
\[
X_{10}:=\sigma_{10}(p),\qquad
X_3:=\sigma_3(p),\qquad
X_4:=\tau_4(p).
\]
又记
\[
\chi_{10}(p):=\operatorname{sgn}\!\bigl(\sigma_{10}(p)\bigr),
\qquad
\chi_{\mathrm{LY}}(p):=\operatorname{sgn}\!\bigl(\sigma_3(p)\bigr),
\]
并令 \(\mathcal T_{\mathrm{LY}}(p)\) 表示 \(g\bmod p\) 的分解型随机变量。则：
\begin{enumerate}
  \item 两域情形下，
  \[
  X_{10}\perp X_3,
  \qquad
  H(X_{10},X_3)=H(X_{10})+H(X_3),
  \qquad
  I(X_{10};X_3)=0;
  \]
  \item 三域情形下，\((X_{10},X_3,X_4)\) 互相独立，故
  \[
  H(X_{10},X_3,X_4)=H(X_{10})+H(X_3)+H(X_4);
  \]
  \item 二元组 \((\chi_{10},\chi_{\mathrm{LY}})\) 在 \(\{\pm1\}^2\) 上严格等分布，因而
  \[
  I(\chi_{10};\chi_{\mathrm{LY}})=0;
  \]
  \item 任意条件 \(\chi_{10}(p)=\pm1\) 都不改变 \(\mathcal T_{\mathrm{LY}}\) 的分布：
  \[
  \delta\!\bigl(\mathcal T_{\mathrm{LY}}=(1)(2)\mid \chi_{10}=\pm1\bigr)=\frac12,
  \]
  \[
  \delta\!\bigl(\mathcal T_{\mathrm{LY}}=(1)(1)(1)\mid \chi_{10}=\pm1\bigr)=\frac16,
  \]
  \[
  \delta\!\bigl(\mathcal T_{\mathrm{LY}}=(3)\mid \chi_{10}=\pm1\bigr)=\frac13.
  \]
\end{enumerate}
因此，双指纹与三指纹体系在 Chebotarev 极限测度下并非仅为弱相关，而是具有严格的零信息耦合。
\end{corollary}
\begin{proof}
第 \(1\) 条与第 \(2\) 条只是定理 \ref{thm:xi-time-part9zg-frobenius-classfunction-tensor-factorization} 的概率论改写：期望按张量分裂即意味着极限分布为乘积分布，而乘积分布的 Shannon 熵可加、互信息为零。

第 \(3\) 条与第 \(4\) 条由定理 \ref{thm:killo-leyang-quadratic-character-decoupling} 直接给出。该定理已经证明
\[
(\chi_{10},\chi_{\mathrm{LY}})
\]
在 \(\{\pm1\}^2\) 上等分布，每个符号盒密度均为 \(1/4\)；并且在任一条件 \(\chi_{10}=\pm1\) 下，\(g\bmod p\) 的三种分解型密度仍分别保持
\[
\frac12,\qquad \frac16,\qquad \frac13.
\]
于是 \(I(\chi_{10};\chi_{\mathrm{LY}})=0\)，并得到条件分布刚性。证毕。
\end{proof}

\begin{theorem}[双二次扇区上的条件张量平均公式]\label{thm:xi-time-part9zg-quadratic-sector-conditional-tensor-average}
沿用定理 \ref{thm:xi-time-part9zg-frobenius-classfunction-tensor-factorization}、
定理 \ref{thm:killo-leyang-joint-chebotarev-product-law}
与定理 \ref{thm:killo-leyang-quadratic-character-decoupling} 的记号。
对任意类函数
\[
f:S_{10}\to\CC,
\qquad
g:S_3\to\CC,
\]
定义群平均与符号配对
\[
\langle f\rangle_{S_{10}}:=\frac{1}{10!}\sum_{\tau\in S_{10}}f(\tau),
\qquad
\langle g\rangle_{S_3}:=\frac{1}{6}\sum_{\upsilon\in S_3}g(\upsilon),
\]
\[
\langle f,\operatorname{sgn}_{10}\rangle_{S_{10}}
:=
\frac{1}{10!}\sum_{\tau\in S_{10}}f(\tau)\operatorname{sgn}(\tau),
\qquad
\langle g,\operatorname{sgn}_{3}\rangle_{S_3}
:=
\frac{1}{6}\sum_{\upsilon\in S_3}g(\upsilon)\operatorname{sgn}(\upsilon).
\]
对每个符号对 \((\varepsilon,\delta)\in\{\pm1\}^2\)，记
\[
\pi_{\varepsilon,\delta}(x)
:=
\#\Bigl\{
p\le x:\ p\nmid \Disc(P_{10})\Disc(P_{\mathrm{LY}}),\
\chi_{10}(p)=\varepsilon,\
\chi_{\mathrm{LY}}(p)=\delta
\Bigr\}.
\]
则有极限恒等式
\[
\boxed{\;
\lim_{x\to\infty}
\frac{1}{\pi_{\varepsilon,\delta}(x)}
\sum_{\substack{p\le x\\ p\nmid \Disc(P_{10})\Disc(P_{\mathrm{LY}})\\
\chi_{10}(p)=\varepsilon,\ \chi_{\mathrm{LY}}(p)=\delta}}
f(\sigma_{10}(p))\,g(\sigma_3(p))
=
\bigl(\langle f\rangle_{S_{10}}+\varepsilon\langle f,\operatorname{sgn}_{10}\rangle_{S_{10}}\bigr)
\bigl(\langle g\rangle_{S_3}+\delta\langle g,\operatorname{sgn}_{3}\rangle_{S_3}\bigr).\;}
\]
\leanverified{paper\_xi\_time\_part9zg\_quadratic\_sector\_conditional\_tensor\_average}
换言之，条件于任意双二次扇区 \((\chi_{10},\chi_{\mathrm{LY}})=(\varepsilon,\delta)\) 后，全部类函数统计都被严格压缩到
\[
\operatorname{span}\{1,\operatorname{sgn}_{10}\}\otimes
\operatorname{span}\{1,\operatorname{sgn}_{3}\}
\]
这一四维阿贝尔谐波子空间上。
\end{theorem}
\begin{proof}
定义符号投影类函数
\[
e_{\varepsilon}^{(10)}(\tau):=\frac{1+\varepsilon\,\operatorname{sgn}(\tau)}{2},
\qquad
e_{\delta}^{(3)}(\upsilon):=\frac{1+\delta\,\operatorname{sgn}(\upsilon)}{2}.
\]
则对任意共同不分歧素数 \(p\)，有
\[
e_{\varepsilon}^{(10)}(\sigma_{10}(p))
=
\mathbf 1_{\{\chi_{10}(p)=\varepsilon\}},
\qquad
e_{\delta}^{(3)}(\sigma_3(p))
=
\mathbf 1_{\{\chi_{\mathrm{LY}}(p)=\delta\}}.
\]
因此条件平均的分子可写为
\[
\sum_{p\le x}^{\ast}
f(\sigma_{10}(p))\,e_{\varepsilon}^{(10)}(\sigma_{10}(p))\,
g(\sigma_3(p))\,e_{\delta}^{(3)}(\sigma_3(p)),
\]
其中 \(\sum^{\ast}\) 表示对共同不分歧素数求和。
再记
\[
\pi_{10,\mathrm{LY}}^\ast(x)
:=
\#\{p\le x:\ p\nmid \Disc(P_{10})\Disc(P_{\mathrm{LY}})\}.
\]
由定理 \ref{thm:xi-time-part9zg-frobenius-classfunction-tensor-factorization}，
\[
\lim_{x\to\infty}
\frac{1}{\pi_{10,\mathrm{LY}}^\ast(x)}
\sum_{p\le x}^{\ast}
f(\sigma_{10}(p))\,e_{\varepsilon}^{(10)}(\sigma_{10}(p))\,
g(\sigma_3(p))\,e_{\delta}^{(3)}(\sigma_3(p))
=
\langle f e_{\varepsilon}^{(10)}\rangle_{S_{10}}\,
\langle g e_{\delta}^{(3)}\rangle_{S_3}.
\]
另一方面，定理 \ref{thm:killo-leyang-quadratic-character-decoupling} 给出
\[
\lim_{x\to\infty}\frac{\pi_{\varepsilon,\delta}(x)}{\pi_{10,\mathrm{LY}}^\ast(x)}=\frac14.
\]
故相除后得到
\[
\lim_{x\to\infty}
\frac{1}{\pi_{\varepsilon,\delta}(x)}
\sum_{\substack{p\le x\\ p\nmid \Disc(P_{10})\Disc(P_{\mathrm{LY}})\\
\chi_{10}(p)=\varepsilon,\ \chi_{\mathrm{LY}}(p)=\delta}}
f(\sigma_{10}(p))\,g(\sigma_3(p))
=
4\langle f e_{\varepsilon}^{(10)}\rangle_{S_{10}}\,
\langle g e_{\delta}^{(3)}\rangle_{S_3}.
\]
最后由
\[
4\langle f e_{\varepsilon}^{(10)}\rangle_{S_{10}}
=
4\left\langle
f\cdot \frac{1+\varepsilon\,\operatorname{sgn}_{10}}{2}
\right\rangle_{S_{10}}
=
\langle f\rangle_{S_{10}}+\varepsilon\langle f,\operatorname{sgn}_{10}\rangle_{S_{10}},
\]
以及 \(S_3\) 上的同样公式，即得结论。证毕。
\end{proof}

\leanverified{paper\_xi\_time\_part9zg\_quadratic\_sector\_conditional\_splitting\_density}
\begin{theorem}[双二次扇区中的联合分解型条件密度全公式]\label{thm:xi-time-part9zg-quadratic-sector-conditional-splitting-density}
沿用定理 \ref{thm:xi-time-part9zg-p10-leyang-joint-splitting-partition-density-table} 的记号。
对分拆
\[
\lambda\vdash 10,
\qquad
\mu\vdash 3,
\]
记 \(C_\lambda\subset S_{10}\)、\(C_\mu\subset S_3\) 为相应共轭类，并定义其置换符号
\[
\operatorname{sgn}(\lambda):=(-1)^{10-\ell(\lambda)},
\qquad
\operatorname{sgn}(\mu):=(-1)^{3-\ell(\mu)}.
\]
对每个 \((\varepsilon,\delta)\in\{\pm1\}^2\)，记
\[
\delta_{\varepsilon,\delta}(\lambda,\mu)
\]
为在扇区条件
\[
\chi_{10}(p)=\varepsilon,
\qquad
\chi_{\mathrm{LY}}(p)=\delta
\]
下，\(P_{10}\bmod p\) 的分解型为 \(\lambda\) 且 \(P_{\mathrm{LY}}\bmod p\) 的分解型为 \(\mu\) 的条件自然密度。
则
\[
\boxed{\;
\delta_{\varepsilon,\delta}(\lambda,\mu)=
\begin{cases}
\displaystyle
4\,\frac{|C_\lambda|}{10!}\frac{|C_\mu|}{3!},
& \operatorname{sgn}(\lambda)=\varepsilon,\ \operatorname{sgn}(\mu)=\delta,\\[8pt]
0,
& \text{否则}.
\end{cases}\;}
\]
换言之，进入任一双二次扇区之后，允许出现的联合分解型恰为符号匹配的那些；在允许子集上，条件分布仅是无条件乘法密度的统一放大因子 \(4\)。
\end{theorem}
\begin{proof}
在定理 \ref{thm:xi-time-part9zg-quadratic-sector-conditional-tensor-average} 中取
\[
f=\mathbf 1_{C_\lambda},
\qquad
g=\mathbf 1_{C_\mu}.
\]
则
\[
\langle f\rangle_{S_{10}}=\frac{|C_\lambda|}{10!},
\qquad
\langle g\rangle_{S_3}=\frac{|C_\mu|}{3!}.
\]
又由于符号在共轭类上为常数，
\[
\langle f,\operatorname{sgn}_{10}\rangle_{S_{10}}
=
\operatorname{sgn}(\lambda)\frac{|C_\lambda|}{10!},
\qquad
\langle g,\operatorname{sgn}_{3}\rangle_{S_3}
=
\operatorname{sgn}(\mu)\frac{|C_\mu|}{3!}.
\]
代回定理 \ref{thm:xi-time-part9zg-quadratic-sector-conditional-tensor-average}，即得
\[
\delta_{\varepsilon,\delta}(\lambda,\mu)
=
\Bigl(1+\varepsilon\operatorname{sgn}(\lambda)\Bigr)\frac{|C_\lambda|}{10!}
\cdot
\Bigl(1+\delta\operatorname{sgn}(\mu)\Bigr)\frac{|C_\mu|}{3!}.
\]
若符号不匹配，则某一括号为 \(0\)，密度即为 \(0\)；若两边都匹配，则两个括号都等于 \(2\)，从而得到所示公式。证毕。
\end{proof}

\begin{corollary}[Lee--Yang 二次角色是“一根情形”的完美探测器]\label{cor:xi-time-part9zg-leyang-quadratic-character-perfect-single-root-detector}
沿用定理 \ref{thm:xi-time-part9zg-quadratic-sector-conditional-splitting-density} 的记号。
记
\[
\mathcal R:=\{\text{\(P_{\mathrm{LY}}\bmod p\) 恰有一个 \(\FF_p\)-根}\},
\]
\[
\mathcal S:=\{\text{\(P_{\mathrm{LY}}\bmod p\) 完全分裂}\},
\qquad
\mathcal I:=\{\text{\(P_{\mathrm{LY}}\bmod p\) 在 \(\FF_p[y]\) 上不可约}\}.
\]
则：
\begin{enumerate}
  \item 条件于 \(\chi_{\mathrm{LY}}(p)=-1\) 时，
  \[
  \delta(\mathcal R\mid \chi_{\mathrm{LY}}=-1)=1;
  \]
  \item 条件于 \(\chi_{\mathrm{LY}}(p)=+1\) 时，
  \[
  \delta(\mathcal R\mid \chi_{\mathrm{LY}}=+1)=0,
  \qquad
  \delta(\mathcal I\mid \chi_{\mathrm{LY}}=+1)=\frac23,
  \qquad
  \delta(\mathcal S\mid \chi_{\mathrm{LY}}=+1)=\frac13.
  \]
\end{enumerate}
因此，二次角色
\[
\chi_{\mathrm{LY}}(p)=\left(\frac{-111}{p}\right)
\]
对 Lee--Yang 三次模 \(p\) 的“单根 / 非单根”机制给出完全判别：负号扇区上单根必然发生，正号扇区上单根绝不发生。
\end{corollary}
\begin{proof}
在 \(S_3\) 中，分解型 \((1)(2)\) 对应换位类，大小为 \(3\)，其符号为 \(-1\)；分解型 \((1)(1)(1)\) 与 \((3)\) 分别对应恒等类与 \(3\)-循环类，大小为 \(1\) 与 \(2\)，其符号均为 \(+1\)。
于是由定理 \ref{thm:xi-time-part9zg-quadratic-sector-conditional-splitting-density}，条件于 \(\chi_{\mathrm{LY}}=-1\) 时，只允许 \((1)(2)\) 型出现，故
\[
\delta(\mathcal R\mid \chi_{\mathrm{LY}}=-1)=1.
\]
同理，条件于 \(\chi_{\mathrm{LY}}=+1\) 时，只允许 \((1)(1)(1)\) 与 \((3)\) 两类出现，因此
\[
\delta(\mathcal R\mid \chi_{\mathrm{LY}}=+1)=0.
\]
再对两种允许类使用同一定理，
\[
\delta(\mathcal I\mid \chi_{\mathrm{LY}}=+1)
=
2\cdot \frac{|(3)|}{3!}
=
2\cdot\frac{2}{6}
=\frac23,
\]
\[
\delta(\mathcal S\mid \chi_{\mathrm{LY}}=+1)
=
2\cdot \frac{|(1)(1)(1)|}{3!}
=
2\cdot\frac{1}{6}
=\frac13.
\]
证毕。
\end{proof}

\leanverified{paper\_xi\_time\_part9zg\_quadratic\_sector\_character\_harmonic\_collapse}
\begin{theorem}[四扇区条件下的角色谱完全坍缩]\label{thm:xi-time-part9zg-quadratic-sector-character-harmonic-collapse}
设 \(\chi\) 为 \(S_{10}\) 的不可约特征标，\(\psi\) 为 \(S_3\) 的不可约特征标。
对每个双二次扇区 \((\varepsilon,\delta)\in\{\pm1\}^2\)，定义条件角色平均
\[
\mathcal A_{\varepsilon,\delta}(\chi,\psi)
:=
\lim_{x\to\infty}
\frac{1}{\pi_{\varepsilon,\delta}(x)}
\sum_{\substack{p\le x\\ p\nmid \Disc(P_{10})\Disc(P_{\mathrm{LY}})\\
\chi_{10}(p)=\varepsilon,\ \chi_{\mathrm{LY}}(p)=\delta}}
\chi(\sigma_{10}(p))\,\psi(\sigma_3(p)).
\]
则有精确公式
\[
\boxed{\;
\mathcal A_{\varepsilon,\delta}(\chi,\psi)=
\begin{cases}
1,& \chi=\mathbf 1,\ \psi=\mathbf 1,\\[2pt]
\varepsilon,& \chi=\operatorname{sgn}_{10},\ \psi=\mathbf 1,\\[2pt]
\delta,& \chi=\mathbf 1,\ \psi=\operatorname{sgn}_{3},\\[2pt]
\varepsilon\delta,& \chi=\operatorname{sgn}_{10},\ \psi=\operatorname{sgn}_{3},\\[2pt]
0,& \text{其余一切情形}.
\end{cases}\;}
\]
换言之，双二次扇区条件把联合 Frobenius 分布的角色谱完全压缩到
\[
\{\mathbf 1,\operatorname{sgn}_{10}\}\otimes
\{\mathbf 1,\operatorname{sgn}_{3}\}
\]
这四维阿贝尔子空间中，而所有真正的非阿贝尔谐波在每个扇区内都严格平均为零。
\end{theorem}
\begin{proof}
由定理 \ref{thm:xi-time-part9zg-quadratic-sector-conditional-tensor-average}，
\[
\mathcal A_{\varepsilon,\delta}(\chi,\psi)
=
\bigl(\langle \chi\rangle_{S_{10}}+\varepsilon\langle \chi,\operatorname{sgn}_{10}\rangle_{S_{10}}\bigr)
\bigl(\langle \psi\rangle_{S_3}+\delta\langle \psi,\operatorname{sgn}_{3}\rangle_{S_3}\bigr).
\]
对有限群不可约特征标的标准正交关系给出
\[
\langle \chi\rangle_{S_{10}}=
\begin{cases}
1,& \chi=\mathbf 1,\\
0,& \chi\neq \mathbf 1,
\end{cases}
\qquad
\langle \chi,\operatorname{sgn}_{10}\rangle_{S_{10}}=
\begin{cases}
1,& \chi=\operatorname{sgn}_{10},\\
0,& \chi\neq \operatorname{sgn}_{10},
\end{cases}
\]
以及
\[
\langle \psi\rangle_{S_3}=
\begin{cases}
1,& \psi=\mathbf 1,\\
0,& \psi\neq \mathbf 1,
\end{cases}
\qquad
\langle \psi,\operatorname{sgn}_{3}\rangle_{S_3}=
\begin{cases}
1,& \psi=\operatorname{sgn}_{3},\\
0,& \psi\neq \operatorname{sgn}_{3}.
\end{cases}
\]
代回即可得到分段公式。证毕。
\end{proof}


\noindent
因此，有限局部化与多域 Frobenius 指纹在算术侧形成了两条彼此衔接的终端闭环：其一，\(\ZZ[S^{-1}]\) 的 Pontryagin 对偶不仅保留对象的同构类型，而且把素数包含偏序精确反译为紧对偶侧的连续商偏序；其二，\(S_{10}\)、\(S_3\) 与 \(S_4\times C_2\) 三个 Frobenius 因子的全部混合统计在密度极限中完全张量分解，从而把判别式符号层、分解型层与第三分歧域层之间的一切互信息压缩为零。于是，prime-ledger 的拓扑刚性与 arithmetic fingerprint 的信息解耦便在同一对偶框架内获得统一刻画。

\endinput

\paragraph{dyadic bulk 地址的线性精度、prime-register 边界账本的超线性代价与奇偶谱壁分离}
在全息 \(2\)-进地址、prime-register 外置下界与单位根扭曲 slow mode 已经分别获得闭式之后，还可把 bulk 与 boundary 的制度分工进一步压缩为同一条刚性链。其要点在于：bulk 侧的 \(B\)-进地址在固定二维 Tate 模块上只需线性精度深度即可恢复有限前缀，而一旦要求 boundary 侧仍保持逐坐标可寻址，prime-register 账本便不可避免地出现 \(T\log T\) 级的最坏情形位长；同时，奇素数扭曲会在 \(p\to\infty\) 时软化到 Perron 主支，而模 \(2\) 奇偶缺陷则始终停留在一条固定的硬谱壁之上。由此，线性 bulk 地址、超线性 boundary 外置、soft prime modes 与 hard parity defect 便在同一层面完成分离。

\begin{theorem}[dyadic bulk 的线性精度深度与 prime-register 边界账本的超线性代价分离]\label{thm:xi-time-part9zh-dyadic-bulk-linear-depth-vs-prime-ledger}
\leanverified{paper\_xi\_time\_part9zh\_dyadic\_bulk\_linear\_depth\_vs\_prime\_ledger}
设 \(E\) 为满足 \(|E|\ge 2\) 的有限事件集，并固定注入编码
\[
\mathrm{code}:E\hookrightarrow \NN,
\qquad
M:=\max_{e\in E}\mathrm{code}(e),
\qquad
B:=2^L>M.
\]
再取非零原始向量
\[
\mathcal T:=T_2(E_{\min}),
\qquad
v\in \mathcal T\setminus 2\mathcal T.
\]
对任意长度 \(T\) 的事件前缀
\[
e^{(T)}=(e_1,\dots,e_T)\in E^T
\]
定义其 dyadic bulk 地址
\[
X_T\!\bigl(e^{(T)}\bigr)
:=
\sum_{t=1}^{T}\mathrm{code}(e_t)B^{t-1}v\in \mathcal T.
\]
并定义最小 \(B\)-进精度深度
\[
\ell_B\!\bigl(X_T(e^{(T)})\bigr)
:=
\min\Bigl\{n\ge 1:\ 
X_T(e^{(T)})\bmod B^n\mathcal T
\text{ 在 }E^T\text{ 中唯一确定 }e_1,\dots,e_T
\Bigr\}.
\]
则对每个 \(e^{(T)}\in E^T\) 都有
\[
\ell_B\!\bigl(X_T(e^{(T)})\bigr)=T.
\]

反之，存在绝对常数 \(c>0\) 与 \(T_0\)，使得对每个 \(T\ge T_0\)，若 \(|\Sigma|\ge 2\) 且
\[
E_T:\Sigma^T\hookrightarrow \NN_{\ge 1}
\]
是任意逐坐标可寻址的 prime-register 外置编码，则有
\[
\max_{s\in\Sigma^T}\log E_T(s)\ge c\,T\log T.
\]
因此
\[
\frac{\max_{s\in\Sigma^T}\log E_T(s)}
{\sup_{e^{(T)}\in E^T}\ell_B(X_T(e^{(T)}))}
\ge c\log T.
\]
换言之，在保持逐位置可寻址性的前提下，boundary prime-register 账本相对于 dyadic bulk 地址必然多出一个不可消去的对数量级。
\end{theorem}
\begin{proof}
固定任意参考尾流 \(\boldsymbol\eta=(\eta_1,\eta_2,\dots)\in E^{\NN}\)。对每个前缀 \(e^{(T)}\in E^T\)，记其补尾无限流为
\[
\mathbf e:=\bigl(e_1,\dots,e_T,\eta_1,\eta_2,\dots\bigr)\in E^{\NN}.
\]
由定理 \ref{thm:xi-time-part9xb-prefix-words-faithful-affine-realization}，
\[
X(\mathbf e)=X_T(e^{(T)})+B^T X(\boldsymbol\eta).
\]
因此，对任意两条长度 \(T\) 的前缀 \(e^{(T)},f^{(T)}\in E^T\)，都有
\[
X(\mathbf e)-X(\mathbf f)
=
X_T(e^{(T)})-X_T(f^{(T)}).
\]

若 \(e^{(T)}\neq f^{(T)}\)，记它们的首个分歧位置为
\[
\tau:=\min\{1\le t\le T:\ e_t\neq f_t\}.
\]
由定理 \ref{thm:xi-time-part9w-holographic-first-difference-valuation}，
\[
X_T(e^{(T)})-X_T(f^{(T)})
\in
B^{\tau-1}\mathcal T\setminus B^\tau\mathcal T.
\]
由于 \(\tau\le T\)，便有
\[
X_T(e^{(T)})-X_T(f^{(T)})\notin B^T\mathcal T.
\]
故不同长度 \(T\) 前缀在商 \(\mathcal T/B^T\mathcal T\) 中具有不同像，从而模 \(B^T\) 的信息足以唯一恢复整条长度 \(T\) 前缀。这说明
\[
\ell_B\!\bigl(X_T(e^{(T)})\bigr)\le T.
\]

另一方面，由 \(|E|\ge 2\)，可取两个不同事件 \(a\neq b\in E\)。固定任意共同前缀
\[
(e_1,\dots,e_{T-1})\in E^{T-1},
\]
并定义
\[
u^{(T)}:=(e_1,\dots,e_{T-1},a),
\qquad
v^{(T)}:=(e_1,\dots,e_{T-1},b).
\]
它们的首个分歧位置恰为 \(T\)，故再由定理
\ref{thm:xi-time-part9w-holographic-first-difference-valuation}，
\[
X_T(u^{(T)})-X_T(v^{(T)})
\in
B^{T-1}\mathcal T\setminus B^T\mathcal T.
\]
于是
\[
X_T(u^{(T)})\equiv X_T(v^{(T)})\pmod{B^{T-1}\mathcal T},
\]
故仅知道模 \(B^{T-1}\) 的类不足以区分长度 \(T\) 前缀。于是
\[
\ell_B\!\bigl(X_T(e^{(T)})\bigr)\ge T.
\]
合并得
\[
\ell_B\!\bigl(X_T(e^{(T)})\bigr)=T.
\]

再看 prime-register 外置编码。由定理
\ref{thm:xi-addressable-godel-worstcase-chebyshev-sharp-lower-bound}，对任意逐坐标可寻址单射
\[
E_T:\Sigma^T\hookrightarrow \NN_{\ge 1},
\qquad |\Sigma|\ge 2,
\]
都满足
\[
\max_{s\in\Sigma^T}\log E_T(s)\ge T\log T+T\log\log T+O(T),
\]
从而存在绝对常数 \(c>0\) 与 \(T_0\)，使得当 \(T\ge T_0\) 时，
\[
\max_{s\in\Sigma^T}\log E_T(s)\ge c\,T\log T.
\]
又因前半部分已证明
\[
\sup_{e^{(T)}\in E^T}\ell_B(X_T(e^{(T)}))=T,
\]
故最后的不等式立即成立。证毕。
\end{proof}

\begin{corollary}[边界 admissibility 的正密度并不削弱最坏情形地址代价]\label{cor:xi-time-part9zh-zg-positive-density-does-not-weaken-address-cost}
定义
\[
\mathrm{ZG}
:=
\left\{
\prod_{k\ge 1}p_k^{z_k}:\ 
z_k\in\{0,1\},\ 
z_kz_{k+1}=0,\ 
z_k=0\ \text{对充分大 }k
\right\},
\]
即平方自由且不含相邻素因子的整数集。若 \(|\Sigma|\ge 2\) 且一族逐坐标可寻址编码
\[
E_T:\Sigma^T\hookrightarrow \mathrm{ZG}
\]
的像完全落在 \(\mathrm{ZG}\) 中，则存在绝对常数 \(c>0\) 与 \(T_0\)，使得对每个 \(T\ge T_0\) 都有
\[
\max_{s\in\Sigma^T}\log E_T(s)\ge c\,T\log T.
\]
与此同时，
\[
\delta(\mathrm{ZG})\in(0,1).
\]
故 prime-register 外置的代价爆炸并非来自 admissible 码字在 \(\NN\) 中过于稀薄；真正的障碍是逐坐标可寻址性本身，而非 admissibility 的体积稀缺。
\end{corollary}
\begin{proof}
由定理 \ref{thm:xi-time-part9w-zg-positive-density-cofinal-tail}，
\[
\delta(\mathrm{ZG})\in(0,1).
\]
另一方面，定理
\ref{thm:xi-addressable-godel-worstcase-chebyshev-sharp-lower-bound}
适用于任意逐坐标可寻址单射
\[
E_T:\Sigma^T\hookrightarrow \NN_{\ge 1},
\]
其证明仅依赖于素数估值可逐坐标恢复，而并不依赖像集在 \(\NN\) 中的稠密度。故把像进一步限制在 \(\mathrm{ZG}\subset\NN_{\ge 1}\) 内，并不会改变同一下界。于是仍存在绝对常数 \(c>0\) 与 \(T_0\)，使得当 \(T\ge T_0\) 时都有
\[
\max_{s\in\Sigma^T}\log E_T(s)\ge c\,T\log T.
\]
证毕。
\end{proof}

\begin{theorem}[奇素数扭曲的软化极限与模 \texorpdfstring{$2$}{2} 缺陷的硬隔离]\label{thm:xi-time-part9zh-oddprime-softening-vs-parity-hard-wall}
记
\[
\lambda:=\rho(1)=\varphi^2,
\qquad
\rho_{(2)}:=\rho_-:=\rho\!\bigl(W(-1)\bigr),
\qquad
\rho_p:=\left|\rho\!\left(e^{2\pi i/p}\right)\right|
\quad (p\ge 3\ \text{为奇素数}).
\]
则有四阶渐近
\[
\log\frac{\rho_p}{\lambda}
=
-\frac{6\sqrt5-5}{250}\left(\frac{2\pi}{p}\right)^2
+\frac{7+24\sqrt5}{75000}\left(\frac{2\pi}{p}\right)^4
+O(p^{-6}),
\]
从而
\[
\frac{\rho_p}{\lambda}\longrightarrow 1
\qquad (p\to\infty,\ p\ \text{奇素数}).
\]
另一方面，
\[
\frac{\rho_{(2)}}{\lambda}\approx 0.8424525579<1.
\]
于是存在 \(p_0\)，使得对所有奇素数 \(p\ge p_0\) 都有
\[
\rho_p>\rho_{(2)}.
\]
因此，充分大的奇素数 congruence 扭曲都比模 \(2\) 奇偶扭曲更接近 Perron 主支；模 \(2\) 缺陷在谱意义上构成一条硬隔离分支，而奇素数扭曲则形成一族软化分支。
\end{theorem}
\begin{proof}
定理 \ref{thm:xi-time-part9xg-parity-wall-versus-cofinal-prime-twists} 已经给出全部所需结论：其一，
\[
\rho_-=\rho\!\bigl(W(-1)\bigr)
\]
是三次方程
\[
x^3-2x^2-1=0
\]
的最大实根，且
\[
\frac{\rho_-}{\lambda}\approx 0.8424525579<1;
\]
其二，对奇素数 \(p\) 的扭曲分支有四阶展开
\[
\log\frac{\rho_p}{\lambda}
=
-\frac{6\sqrt5-5}{250}\left(\frac{2\pi}{p}\right)^2
+\frac{7+24\sqrt5}{75000}\left(\frac{2\pi}{p}\right)^4
+O(p^{-6}),
\]
故 \(\rho_p/\lambda\to 1\)；其三，存在奇素数阈值 \(p_0\) 使得所有 \(p\ge p_0\) 都满足
\[
\rho_p>\rho_-.
\]
把 \(\rho_-\) 记为 \(\rho_{(2)}\) 即得结论。证毕。
\end{proof}

\begin{corollary}[奇素数扭曲弛豫时间的 \texorpdfstring{$p^2$}{p2} 律与 parity-defect 的有限寿命]\label{cor:xi-time-part9zh-oddprime-p2-relaxation-vs-parity-lifetime}
沿用定理 \ref{thm:xi-time-part9zh-oddprime-softening-vs-parity-hard-wall} 的记号，定义奇素数扭曲的弛豫时间
\[
\tau_p:=\frac{1}{-\log(\rho_p/\lambda)},
\]
以及 parity-defect 的弛豫时间
\[
\tau_{(2)}:=\frac{1}{-\log(\rho_{(2)}/\lambda)}.
\]
则
\[
\tau_p=
\frac{250}{(6\sqrt5-5)(2\pi)^2}\,p^2\bigl(1+O(p^{-2})\bigr),
\qquad
\tau_{(2)}<\infty.
\]
特别地，
\[
\frac{\tau_p}{\tau_{(2)}}\longrightarrow\infty
\qquad (p\to\infty,\ p\ \text{奇素数}).
\]
故奇素数扭曲形成一族可任意变慢的 soft modes，而 parity-defect 始终对应一条寿命有限的 hard mode；二者不属于同一低能动力学普适类。
\end{corollary}
\begin{proof}
由定理 \ref{thm:xi-time-part9zh-oddprime-softening-vs-parity-hard-wall}，
\[
-\log\frac{\rho_p}{\lambda}
=
\frac{6\sqrt5-5}{250}\left(\frac{2\pi}{p}\right)^2
+O(p^{-4}).
\]
对其取倒数，得到
\[
\tau_p
=
\frac{1}{\left(\frac{6\sqrt5-5}{250}\right)(2\pi)^2}\,p^2
\bigl(1+O(p^{-2})\bigr),
\]
即所示常数。

另一方面，定理
\ref{thm:xi-time-part9zh-oddprime-softening-vs-parity-hard-wall}
已给出
\[
0<\frac{\rho_{(2)}}{\lambda}<1,
\]
故
\[
\tau_{(2)}=\frac{1}{-\log(\rho_{(2)}/\lambda)}
\]
是有限正常数。于是
\[
\frac{\tau_p}{\tau_{(2)}}\to\infty.
\]
证毕。
\end{proof}

\noindent
由此，bulk 与 boundary、soft 与 hard 两组分工同时被锁定：在地址侧，dyadic bulk 只需线性精度深度 \(T\) 便足以恢复长度 \(T\) 的有限前缀，而 boundary 侧一旦要求逐坐标 prime-addressability，就不可避免地付出 \(T\log T\) 级的最坏位长；在 slow-mode 侧，奇素数扭曲可以连续软化到 Perron 主支，而模 \(2\) 奇偶缺陷却始终停留在固定的 \(O(1)\) 谱壁之上。于是，线性 bulk 地址、超线性 prime-ledger 外置、以及 soft prime modes 与 hard parity defect 的制度分层便在同一条主线中同时显现。

\endinput

\paragraph{最佳几何均匀化长度上的共振保真与传输预算失配}
在黄金分支的 \(W_1\) 传输奇偶双相律与 bulk 共振梯的 Fibonacci 极限常数已经分别固定之后，还可把几何均匀化与算术共振之间的制度断裂进一步压缩为两条直接后果。其核心是：即便沿最佳逼近长度 \(N=F_n\) 观察，一维几何误差也只按 \(F_n^{-1}\) 衰减，而同一子列上的共振梯却保留严格正的宏观幅度；因此，几何意义下的准均匀化并不会迫使算术共振同步消失。进一步，对全部频率的累积共振预算 \(\sum_{u\le N}C_\varphi(u)^2\) 而言，其增长至少为对数量级，因而不可能由仅依赖一维传输差异的任何上界机制统一控制。

\begin{theorem}[最佳逼近长度上的几何均匀化与算术共振不相容]\label{thm:xi-time-part9zi-fibonacci-transport-resonance-incompatibility}
\leanverified{paper\_xi\_time\_part9zi\_fibonacci\_transport\_resonance\_incompatibility}
记
\[
T_N(\varphi^{-1}):=\mathsf T_N(\varphi^{-1})
=
W_1(\mu_N,\lambda),
\]
其中 \(\mu_N\) 为黄金 Kronecker 轨道的经验分布，\(\lambda\) 为 Lebesgue 测度。又记 \(C_\varphi(u)\) 为 bulk 共振梯。则沿 Fibonacci 长度 \(N=F_n\) 有
\[
F_n\,T_{F_n}(\varphi^{-1})
\longrightarrow
\begin{cases}
\displaystyle \frac12+\frac{1}{2\sqrt5}, & n\ \text{为偶},\\[2mm]
\displaystyle \frac12-\frac{1}{2\sqrt5}+\frac1{15}, & n\ \text{为奇},
\end{cases}
\]
并且存在严格正常数
\[
c_{\mathrm{Fib}}:=\lim_{n\to\infty}C_\varphi(F_n)>0.
\]
因此
\[
T_{F_n}(\varphi^{-1})=\Theta(F_n^{-1}),
\qquad
C_\varphi(F_n)=\Theta(1),
\]
并且
\[
\liminf_{n\to\infty}
F_n\,T_{F_n}(\varphi^{-1})\,C_\varphi(F_n)>0.
\]
故在最优一维准均匀化子序列上，算术共振并不随几何误差同步消失；最佳几何分布与共振湮灭在黄金分支上不能同时发生。
\end{theorem}
\begin{proof}
由定理 \ref{thm:xi-golden-fibonacci-w1-star-strict-parity-law}，
\[
\lim_{n\to\infty,\ n\ \mathrm{even}}
F_n\,T_{F_n}(\varphi^{-1})
=
\frac12+\frac{1}{2\sqrt5},
\]
\[
\lim_{n\to\infty,\ n\ \mathrm{odd}}
F_n\,T_{F_n}(\varphi^{-1})
=
\frac12-\frac{1}{2\sqrt5}+\frac1{15}.
\]
这两个极限常数都严格为正，故
\[
F_n\,T_{F_n}(\varphi^{-1})=\Theta(1),
\qquad
T_{F_n}(\varphi^{-1})=\Theta(F_n^{-1}).
\]

另一方面，由推论 \ref{cor:fold-resonance-ladder-fib-luc-limits}，
\[
C_\varphi(F_n)\longrightarrow c_{\mathrm{Fib}}>0,
\]
从而
\[
C_\varphi(F_n)=\Theta(1).
\]
于是
\[
\liminf_{n\to\infty}
F_n\,T_{F_n}(\varphi^{-1})\,C_\varphi(F_n)
\ge
\left(\frac12-\frac{1}{2\sqrt5}+\frac1{15}\right)c_{\mathrm{Fib}}
>0.
\]
这便得到所述不相容性。证毕。
\end{proof}

\begin{corollary}[传输正则性不能约束共振预算]\label{cor:xi-time-part9zi-transport-regularity-cannot-bound-resonance-budget}
\leanverified{paper\_xi\_time\_part9zi\_transport\_regularity\_cannot\_bound\_resonance\_budget}
定义累积共振预算
\[
\mathcal E(N):=\sum_{u=1}^{N}C_\varphi(u)^2.
\]
则存在常数 \(c>0\) 与 \(N_0\)，使得对一切 \(N\ge N_0\) 都有
\[
\mathcal E(N)\ge c\log N.
\]
与此同时，沿 Fibonacci 长度 \(N=F_n\) 有
\[
F_n\,T_{F_n}(\varphi^{-1})=O(1),
\qquad
\mathcal E(F_n)\longrightarrow\infty.
\]
特别地，不存在任何仅依赖一维传输差异 \(T_N(\varphi^{-1})\) 的上界机制，可以在最佳逼近长度上控制累积共振预算。等价地，几何意义下的准均匀采样并不驯化 bulk 共振的累积增长。
\end{corollary}
\begin{proof}
由定理 \ref{thm:xi-resonance-closure-log-lower-bound-fib-luc}，
\[
\liminf_{N\to\infty}
\frac{1}{\log N}\sum_{u=1}^{N}C_\varphi(u)^2
\ge
\frac{c_{\mathrm{Fib}}^2+c_{\mathrm{Luc}}^2}{\log\varphi}>0.
\]
因此可取某个常数 \(c>0\) 与 \(N_0\)，使得对一切 \(N\ge N_0\) 都有
\[
\mathcal E(N)\ge c\log N.
\]

再由定理 \ref{thm:xi-time-part9zi-fibonacci-transport-resonance-incompatibility}，
\[
F_n\,T_{F_n}(\varphi^{-1})=O(1).
\]
而由于 \(F_n\to\infty\)，前式中的对数下界给出
\[
\mathcal E(F_n)\ge c\log F_n\longrightarrow\infty.
\]

若存在某个仅依赖传输差异的控制机制，即存在函数
\[
\Phi:[0,\infty)\to[0,\infty)
\]
在有界区间上有界，并满足对全部 \(N\) 都有
\[
\mathcal E(N)\le \Phi\!\bigl(N\,T_N(\varphi^{-1})\bigr),
\]
则沿 \(N=F_n\) 应得到右端有界而左端发散的矛盾。故不存在此类控制机制。证毕。
\end{proof}

\noindent
因此，黄金分支在最佳几何均匀化尺度上的表现并不是“几何误差越小，算术记忆便越弱”的单调图景。相反，Fibonacci 长度同时暴露出两种彼此独立的极限：一方面，\(W_1\) 传输误差确实达到 \(F_n^{-1}\) 级的最优衰减；另一方面，bulk 共振梯在同一子列上却保留严格正的宏观幅度，并使累积预算至少按 \(\log N\) 增长。于是，transport regularity 与 resonance budget 之间不存在可相互代偿的统一读出准则。

\endinput

\paragraph{RH-sharp 五次阈值域、五元语言门槛、整数射影桥与奇偶阴影塌缩}
在 RH-sharp 临界五次域、Lee--Yang 分支不变量的共同三次数域、有限字母表的外置预算下界、window-\(6\) 群胚代数的 Wedderburn/自同构分解，以及附录射影压力的整数退化已经固定之后，还可把这些接口进一步压缩为下列八条后果。它们分别刻画 RH-sharp 端点与 Lee--Yang 局部奇性域的算术正交、线性密度可逆语言门槛与 window-\(6\) Lie 因子禁阻的错位、整数 \(q\)-桥上的分母湮灭刚性、输出分桶矩阵对整个射影--moment 塔的充分性、六倍窗零点相对密度对碰撞尺度的有界传输律失效，以及审计共振窗中 parity 阴影的联合 \(16\)-周期塌缩与矩阶别名。

\begin{theorem}[RH-sharp 临界五次域与 Lee--Yang 局部奇性三次数域的算术正交]\label{thm:xi-time-part9zj-rhsharp-leyang-localfield-arithmetic-orthogonality}
沿用定理 \ref{thm:xi-time-part9pb-rh-critical-quintic-field} 与定理 \ref{thm:xi-time-part65b-leyang-branch-b-birational-uniformization} 的记号。令
\[
K_{\mathrm{LY}}:=\QQ(y_*,\lambda_*,\kappa_*^2,B^2).
\]
则有
\[
[\QQ(u_R):\QQ]=5,
\qquad
[K_{\mathrm{LY}}:\QQ]=3,
\qquad
\QQ(u_R)\cap K_{\mathrm{LY}}=\QQ.
\]
特别地，
\[
u_R\notin K_{\mathrm{LY}},
\qquad
[\QQ(u_R)K_{\mathrm{LY}}:\QQ]=15.
\]
\end{theorem}
\begin{proof}
由定理 \ref{thm:xi-time-part9pb-rh-critical-quintic-field}，
\[
[\QQ(u_R):\QQ]=5.
\]
另一方面，定理 \ref{thm:xi-time-part65b-leyang-branch-b-birational-uniformization} 已给出
\[
\QQ(y_*)=\QQ(\lambda_*)=\QQ(\kappa_*^2)=\QQ(B^2),
\]
且该共同域为三次数域，故
\[
[K_{\mathrm{LY}}:\QQ]=3.
\]
令
\[
E:=\QQ(u_R)\cap K_{\mathrm{LY}}.
\]
则塔式定理给出
\[
[E:\QQ]\mid 5,
\qquad
[E:\QQ]\mid 3.
\]
由于 \(3\) 与 \(5\) 互素，只能有
\[
[E:\QQ]=1,
\]
即
\[
\QQ(u_R)\cap K_{\mathrm{LY}}=\QQ.
\]
于是 \(u_R\notin K_{\mathrm{LY}}\)，并且
\[
[\QQ(u_R)K_{\mathrm{LY}}:\QQ]
=
\frac{[\QQ(u_R):\QQ]\,[K_{\mathrm{LY}}:\QQ]}
{[\QQ(u_R)\cap K_{\mathrm{LY}}:\QQ]}
=15.
\]
证毕。
\end{proof}

\begin{theorem}[五元事件语言门槛与 window-\texorpdfstring{$6$}{6} Lie 因子禁阻的错位]\label{thm:xi-time-part9zj-quinary-language-window6-lie-mismatch}
设存在可逆提升
\[
\Psi_m:\Omega_m\to X_m\times \Sigma^{\le L_m},
\qquad
\mathrm{pr}_{X_m}\circ\Psi_m=\Fold_m,
\]
其中 \(|\Sigma|=M\ge 2\)，并且最坏情形长度满足
\[
L_m\le \left(\frac{4}{27}+o(1)\right)m.
\]
则必有
\[
M\ge 5.
\]

另一方面，记
\[
A_6:=\CC[\mathcal R_6^{\mathrm{bin}}],
\]
并令 \(\mathcal K_6\subset \Aut_{\CC\text{-alg}}(A_6)^\circ\) 为任一极大紧连通子群。则
\[
\Lie(\mathcal K_6)
\cong
\mathfrak{su}(2)^{\oplus 8}\oplus
\mathfrak{su}(3)^{\oplus 4}\oplus
\mathfrak{su}(4)^{\oplus 9}.
\]
特别地，window-\(6\) 的连通 Lie 来源只可能出现
\[
\mathfrak{su}(2),\qquad \mathfrak{su}(3),\qquad \mathfrak{su}(4),
\]
而不可能出现 \(\mathfrak{su}(5)\)。
\end{theorem}
\begin{proof}
由命题 \ref{prop:xi-time-part9f-alphabet-threshold}，当
\(
L_m\le (\alpha+o(1))m
\)
时必有
\[
M\ge \exp\!\left(\frac{\log(2/\varphi)}{\alpha}\right).
\]
取 \(\alpha=\frac{4}{27}\)，得到
\[
M\ge \exp\!\left(\frac{27}{4}\log(2/\varphi)\right)\approx 4.181055178\ldots,
\]
故因 \(M\) 为整数，只能有
\[
M\ge 5.
\]

对 Lie 侧，由定理 \ref{thm:xi-time-part60acb-binfold-algebra-automorphism-semidirect}，
\[
\Aut_{\CC\text{-alg}}(A_6)^\circ
\cong
\operatorname{PGL}_2(\CC)^8\times
\operatorname{PGL}_3(\CC)^4\times
\operatorname{PGL}_4(\CC)^9.
\]
任取其极大紧连通子群 \(\mathcal K_6\)，则标准 Lie 群理论给出
\[
\mathcal K_6
\cong
\operatorname{PU}(2)^8\times
\operatorname{PU}(3)^4\times
\operatorname{PU}(4)^9,
\]
从而
\[
\Lie(\mathcal K_6)
\cong
\mathfrak{su}(2)^{\oplus 8}\oplus
\mathfrak{su}(3)^{\oplus 4}\oplus
\mathfrak{su}(4)^{\oplus 9}.
\]
因此其紧简单 Lie 因子只可能来自 \(\mathfrak{su}(2),\mathfrak{su}(3),\mathfrak{su}(4)\)，绝无 \(\mathfrak{su}(5)\) 的出现。证毕。
\end{proof}

\begin{theorem}[整数 \texorpdfstring{$q$}{q} 层射影桥上的分母湮灭刚性]\label{thm:xi-time-part9zj-integer-q-projective-bridge-denominator-annihilation}
\leanverified{paper\_xi\_time\_part9zj\_integer\_q\_projective\_bridge\_denominator\_annihilation}
沿用定义 \ref{def:pom-projective-transfer-operator} 与定义 \ref{def:pom-homogeneous-polynomials-simplex} 的记号。对任意整数 \(q\ge 1\)、任意 \(p\in\mathcal H_q\) 与任意 \(w\in\Delta\)，都有严格恒等式
\[
(\mathcal L_q p)(w)=\sum_{b\in A_{\mathrm{out}}}p(B_bw).
\]
特别地，整数 \(q\)-层上的 projective bridge 不再保留任何归一化分母
\[
g_b(w)^{-q},
\]
而只依赖原始输出分桶矩阵 \(B_b\)。
\end{theorem}
\begin{proof}
当 \(g_b(w)>0\) 时，由齐次性与
\[
\Phi_b(w)=\frac{B_bw}{g_b(w)}
\]
可得
\[
p(\Phi_b(w))
=
p\!\left(\frac{B_bw}{g_b(w)}\right)
=
\frac{p(B_bw)}{g_b(w)^q}.
\]
故
\[
g_b(w)^q\,p(\Phi_b(w))=p(B_bw).
\]
若 \(g_b(w)=0\)，则 \(B_b\ge 0\)、\(w\ge 0\) 且 \(\langle\mathbf 1,B_bw\rangle=0\) 蕴含 \(B_bw=0\)，从而
\[
p(B_bw)=0.
\]
把两种情形代回定义 \ref{def:pom-projective-transfer-operator} 即得
\[
(\mathcal L_q p)(w)=\sum_{b\in A_{\mathrm{out}}}p(B_bw).
\]
证毕。
\end{proof}

\begin{corollary}[输出分桶矩阵对整个射影--moment 塔的充分性]\label{cor:xi-time-part9zj-output-bucket-full-tower-sufficiency}
沿用定义 \ref{def:pom-projective-transfer-operator} 的记号。设 \(\mathcal T,\mathcal T'\) 为两套确定性在线核，具有同一状态集 \(Q\) 与同一输出字母表 \(A_{\mathrm{out}}\)，并满足
\[
B_b(\mathcal T)=B_b(\mathcal T')
\qquad
(b\in A_{\mathrm{out}}).
\]
则对每个实参数 \(q\ge 1\) 都有
\[
\mathcal L_q(\mathcal T)=\mathcal L_q(\mathcal T')
\qquad\text{作为 }\Delta\text{ 上的射影转移算子}.
\]
进一步，对每个整数 \(q\ge 2\) 都有
\[
\widetilde A_q(\mathcal T)=\widetilde A_q(\mathcal T'),
\qquad
r_q(\mathcal T)=r_q(\mathcal T').
\]
因此，在 projective--moment 口径下，原始输出分桶矩阵族 \(\{B_b\}_{b\in A_{\mathrm{out}}}\) 已经充分决定整条连续射影桥与全部整数 moment-kernel 塔。
\end{corollary}
\begin{proof}
由定义 \ref{def:pom-projective-transfer-operator}，
\[
g_b(w)=\langle\mathbf 1,B_bw\rangle,
\qquad
\Phi_b(w)=\frac{B_bw}{\langle\mathbf 1,B_bw\rangle}
\quad(g_b(w)>0),
\]
故对每个实 \(q\ge 1\)，\(\mathcal L_q\) 完全由矩阵 \(B_b\) 决定。于是
\[
\mathcal L_q(\mathcal T)=\mathcal L_q(\mathcal T').
\]
对整数 \(q\ge 2\)，结论进一步由定理
\ref{thm:xi-projective-output-bucket-determines-integer-moment-tower}
立即推出。证毕。
\end{proof}

\begin{theorem}[六倍窗零点相对密度对碰撞尺度不存在有界闭包]\label{thm:xi-time-part9zj-zero-density-no-bounded-closure-collision-scale}
记
\[
\eta_m:=\frac{|\mathcal Z_m|}{F_{m+2}}.
\]
则不存在任何在 \(0\) 的某邻域内有界的函数
\[
\Psi:[0,1]\to [0,\infty)
\]
使得沿全部充分大的六倍窗
\[
m+2\equiv 0\pmod 6
\]
恒有
\[
F_{m+2}\,\mathsf{Col}_m\le \Psi(\eta_m).
\]
\end{theorem}
\begin{proof}
由定理 \ref{thm:xi-time-part70c-zero-density-exact-ratio-exponent-multiple6}，当
\[
m+2\equiv 0\pmod 6
\]
时有
\[
\eta_m=\frac{|\mathcal Z_m|}{F_{m+2}}
=
\varphi^{-m/2+o(m)}
\longrightarrow 0.
\]
另一方面，推论 \ref{cor:fold-collision-scale-diverges} 给出
\[
F_{m+2}\,\mathsf{Col}_m\longrightarrow +\infty.
\]
若存在所述 \(\Psi\)，且其在 \(0\) 的某邻域内有界，则当 \(m\) 充分大时，由 \(\eta_m\to 0\) 可知 \(\Psi(\eta_m)\) 必有统一上界；这与
\[
F_{m+2}\,\mathsf{Col}_m\to +\infty
\]
矛盾。证毕。
\end{proof}

\begin{theorem}[审计共振窗 mod-\texorpdfstring{$2$}{2} 阴影的联合 \texorpdfstring{$16$}{16} 周期塌缩]\label{thm:xi-time-part9zj-resonance-window-mod2-joint-period16-collapse}
\leanverified{paper\_xi\_time\_part9zj\_resonance\_window\_mod2\_joint\_period16\_collapse}
对每个 \(q\in\{9,\dots,17\}\)，取一个整数平移 \(a_q\ge 0\)，使得平移后序列
\[
\mathcal S_q(m):=S_q(m+a_q)\pmod 2
\qquad (m\ge 0)
\]
已经进入纯周期阶段。定义联合向量过程
\[
\mathcal V(m):=
\bigl(
\mathcal S_9(m),\mathcal S_{10}(m),\dots,\mathcal S_{17}(m)
\bigr)\in \FF_2^9.
\]
则对一切 \(m\ge 0\) 都有
\[
\mathcal V(m+16)=\mathcal V(m).
\]
换言之，整个审计共振窗的 parity 阴影可同时压缩到一个周期长度至多为 \(16\) 的有限状态系统。
\end{theorem}
\begin{proof}
结论 \ref{con:xi-terminal-high-order-collision-mod2-automaticity} 已证明：对 \(q=9,\dots,17\)，序列 \(S_q(m)\bmod 2\) 都最终周期化，并且其最终周期分别整除
\[
4\quad (q=9,10),\qquad
8\quad (q=11,\dots,16),\qquad
16\quad (q=17).
\]
选取各个 \(a_q\) 使相应预周期被完全吸收后，每个 \(\mathcal S_q\) 都成为纯周期序列，且其周期仍分别整除上述数值。于是联合向量过程 \(\mathcal V(m)\) 的周期整除
\[
\mathrm{lcm}(4,8,16)=16.
\]
故
\[
\mathcal V(m+16)=\mathcal V(m)
\qquad (m\ge 0).
\]
证毕。
\end{proof}

\begin{proposition}[审计共振窗的 mod-\texorpdfstring{$2$}{2} 矩阶别名]\label{prop:xi-time-part9zj-resonance-window-mod2-order-aliasing}
\leanverified{paper\_xi\_time\_part9zj\_resonance\_window\_mod2\_order\_aliasing}
在审计共振窗 \(q=9,\dots,17\) 中，不同矩阶可以具有完全相同的 mod-\(2\) 特征阴影。具体地，
\[
P_{13}(\lambda)\equiv P_{15}(\lambda)\equiv \lambda^3(\lambda+1)^8\pmod 2.
\]
因此，mod-\(2\) 特征阴影并不忠实记录 moment order。
\end{proposition}
\begin{proof}
这正是表 \ref{tab:fold_collision_charpoly_mod2_shadow_q9_17} 的直接读数：对 \(q=13\) 与 \(q=15\)，表中均给出
\[
P_q(\lambda)\ (\mathrm{mod}\ 2)=\lambda^3(\lambda+1)^8.
\]
另一方面，定理 \ref{thm:pom-projective-operator-degenerates-to-moment-kernel} 说明每个整数 \(q\ge 2\) 都对应一条独立的整数 moment-kernel 层。因此，mod-\(2\) 阴影在 \(q\)-阶上发生了真正的矩阶别名，而不是记号层面的巧合。证毕。
\end{proof}

\begin{conjecture}[五次阈值域、五元语言门槛与 parity 别名的三重不可压缩性]\label{conj:xi-time-part9zj-quintic-quinary-parity-triple-hardcore}
任一试图在本文语义中闭合 classical RH 的 finitary normalization，若同时施行下列三种压缩：
\begin{enumerate}
  \item 把 RH-sharp 阈值数据压缩到 Lee--Yang 局部奇性三次数域 \(K_{\mathrm{LY}}\)；
  \item 把可逆重写压缩到线性密度级的有限事件语言，并把门槛理解为某个自动对应的 \(SU(5)\)-几何来源；
  \item 把高阶 moment 信息压缩到审计共振窗的 mod-\(2\) 阴影。
\end{enumerate}
则该正规化不可能闭合。其不可消除的缺失分别对应于：
\begin{enumerate}
  \item RH-sharp 端点 \(u_R\) 所在的五次域与 \(K_{\mathrm{LY}}\) 算术正交；
  \item 五元语言门槛与 window-\(6\) 的 Lie 因子禁阻彼此错位，后者的紧简单因子至多停留在 \(\mathfrak{su}(4)\)；
  \item parity 阴影在审计共振窗内同时发生联合 \(16\)-周期塌缩与矩阶别名。
\end{enumerate}
等价地，任何同时试图有限化阈值域、语言资源与高阶矩谱的证明链，都会分别丢失 quintic threshold、quinary language 与 moment-order faithfulness 这三类彼此不等价的刚性。
\end{conjecture}

\noindent
由此，整数 \(q\)-桥上的分母湮灭说明：在 projective--moment 口径下，原始输出分桶矩阵族已经吸收了整条离散矩谱塔的全部可见结构；而六倍窗零点相对密度对碰撞尺度不存在任何有界传输律闭包，又表明频域稀疏率并不足以单独控制真实碰撞强度。在此基础上，真正拒绝共同有限压缩的三条硬核链已经显现：RH-sharp 端点属于与 Lee--Yang 局部奇性三次数域算术正交的五次域；线性密度可逆语言达到五元门槛时，window-\(6\) 连通 Lie 因子却仍止于 \(\mathfrak{su}(4)\)；而审计共振窗中的 parity 阴影既塌缩为联合 \(16\)-周期系统，又在不同矩阶之间发生特征阴影别名。于是，端点域、有限事件语言与奇偶谱影子三者之间不存在共同的忠实有限压缩通道。

\endinput

\paragraph{临界成功相图、六倍窗半尺度缺口、局部 Cram\'er 展开与 escort 截断律}
在 bin-fold 临界资源窗的双折点极限、六倍窗零点块的半阶密度律、真实输入 \(40\)-状态核输出位势的局部 Legendre 正规形，以及 escort 温度族的两点极限已经固定之后，还可进一步抽取出下列七条彼此兼容的后果。它们分别刻画临界 bit 尺度下最优反演成功率的精确三相相图、给定目标成功率时连续预算与整数预算的反演公式、六倍窗零点块的 Lucas 半尺度律、碰撞缺口与熵缺口共享的算术半尺度缺口、典型输出密度附近的四阶 Cram\'er--Legendre 展开及其左右不对称性，以及 escort 温度族在固定双尺度 \(\{\varphi^{-1},1\}\) 上的截断极限。

\begin{theorem}[临界 bit 尺度下最优反演成功率的精确三相相图]\label{thm:xi-time-part9zk-critical-bit-success-phase-diagram}
定义 bin-fold 的最优成功率
\[
\mathrm{Succ}_m(B):=\mathrm{Succ}_m^{\mathrm{bin}}(B)
=
2^{-m}C_m^{\mathrm{bin}}(B)
\qquad (B\in \ZZ_{\ge 0}).
\]
若整数序列 \((B_m)_{m\ge 1}\) 满足
\[
s_m:=2^{B_m}\left(\frac{\varphi}{2}\right)^m\longrightarrow s\in(0,\infty),
\]
则存在极限
\[
\mathrm{Succ}_m(B_m)\longrightarrow \mathcal S_\infty(s),
\]
其中
\[
\mathcal S_\infty(s)
:=
\frac{\varphi}{\sqrt5}\min\{1,s\}
\;+\;
\frac{1}{\sqrt5}\min\{\varphi^{-1},s\}.
\]
等价地，
\[
\mathcal S_\infty(s)=
\begin{cases}
\dfrac{\varphi^2}{\sqrt5}\,s,
& 0<s\le \varphi^{-1},\\[8pt]
\dfrac{\varphi}{\sqrt5}\,s+\dfrac{1}{\varphi\sqrt5},
& \varphi^{-1}\le s\le 1,\\[10pt]
1,
& s\ge 1.
\end{cases}
\]
因此，在临界位率
\[
m\log_2\!\left(\frac{2}{\varphi}\right)
\]
的 \(O(1)\)-窗口内，最优恢复并不表现为模糊过渡，而是由
\[
s=\varphi^{-1},
\qquad
s=1
\]
两处刚性折点所控制的三相极限曲线。
\end{theorem}
\begin{proof}
由推论 \ref{cor:xi-foldbin-critical-scale-success-curve-inversion}，定义
\[
\mathcal S_m(s):=
2^{-m}\,\mathcal C_m^{\mathrm{bin,cont}}\!\left(
s\left(\frac{2}{\varphi}\right)^m
\right),
\]
则 \(\mathcal S_m\) 在每个紧集 \(K\Subset(0,\infty)\) 上一致收敛到
\[
\mathcal S_\infty(s)
=
\frac{\varphi}{\sqrt5}\min\{1,s\}
\;+\;
\frac{1}{\sqrt5}\min\{\varphi^{-1},s\}.
\]
另一方面，由定义
\[
\mathrm{Succ}_m(B_m)
=
2^{-m}C_m^{\mathrm{bin}}(B_m)
=
2^{-m}\,\mathcal C_m^{\mathrm{bin,cont}}(2^{B_m})
=
\mathcal S_m(s_m).
\]
取紧区间 \(K\Subset(0,\infty)\) 使得 \(s\in K\) 且对充分大的 \(m\) 都有 \(s_m\in K\)。于是
\[
\bigl|\mathrm{Succ}_m(B_m)-\mathcal S_\infty(s)\bigr|
\le
\bigl|\mathcal S_m(s_m)-\mathcal S_\infty(s_m)\bigr|
+
\bigl|\mathcal S_\infty(s_m)-\mathcal S_\infty(s)\bigr|.
\]
第一项由 \(\mathcal S_m\to \mathcal S_\infty\) 在 \(K\) 上的一致收敛趋于 \(0\)，第二项由 \(\mathcal S_\infty\) 的连续性和 \(s_m\to s\) 也趋于 \(0\)。故
\[
\mathrm{Succ}_m(B_m)\to \mathcal S_\infty(s).
\]
将 \(\mathcal S_\infty\) 逐段展开，即得所示分段公式。证毕。
\end{proof}

\begin{theorem}[给定目标成功率的连续临界预算反演公式]\label{thm:xi-time-part9zk-continuous-critical-budget-inversion}
对任意 \(p\in(0,1)\)，定义连续临界预算
\[
b_m^\star(p)
:=
\inf\left\{
b\in [0,\infty):\
2^{-m}\,\mathcal C_m^{\mathrm{bin,cont}}(2^b)\ge p
\right\}.
\]
则有渐近展开
\[
b_m^\star(p)
=
m\log_2\!\left(\frac{2}{\varphi}\right)
+
\log_2 s^\star(p)
+
o(1),
\]
其中反演函数 \(s^\star:(0,1)\to(0,1)\) 显式为
\[
s^\star(p)
=
\begin{cases}
\dfrac{\sqrt5}{\varphi^2}\,p,
& 0<p\le \dfrac{\varphi}{\sqrt5},\\[8pt]
\dfrac{\sqrt5\,p-\varphi^{-1}}{\varphi},
& \dfrac{\varphi}{\sqrt5}\le p<1.
\end{cases}
\]
相应的整数最小预算
\[
B_m^\star(p)
:=
\min\bigl\{B\in\ZZ_{\ge 0}:\ \mathrm{Succ}_m(B)\ge p\bigr\}
\]
满足严格等式
\[
B_m^\star(p)=\left\lceil b_m^\star(p)\right\rceil.
\]
因此，目标成功率
\[
p=\frac{\varphi}{\sqrt5}
\]
构成严格的预算转折点；其下侧与上侧分别对应两条不同斜率的线性反演律。
\end{theorem}
\begin{proof}
定义由连续临界预算所诱导的尺度变量
\[
s_m^\star(p)
:=
2^{b_m^\star(p)}\left(\frac{\varphi}{2}\right)^m.
\]
其中 \(\mathcal S_m\) 仍取自定理
\ref{thm:xi-time-part9zk-critical-bit-success-phase-diagram}
的证明。由
\[
2^{-m}\,\mathcal C_m^{\mathrm{bin,cont}}(2^b)
=
\mathcal S_m\!\left(2^b\left(\frac{\varphi}{2}\right)^m\right)
\]
可知
\[
b_m^\star(p)
=
m\log_2\!\left(\frac{2}{\varphi}\right)
+
\log_2 s_m^\star(p).
\]
又由于
\[
b\longmapsto 2^{-m}\,\mathcal C_m^{\mathrm{bin,cont}}(2^b)
\]
在 \([0,\infty)\) 上连续且单调不减，定义中的下确界必被取到，从而
\[
\mathcal S_m\bigl(s_m^\star(p)\bigr)=p.
\]

再令 \(s^\star(p)\) 为极限曲线 \(\mathcal S_\infty\) 的最小资源广义逆：
\[
s^\star(p):=
\inf\{s\ge 0:\ \mathcal S_\infty(s)\ge p\}.
\]
由定理 \ref{thm:xi-time-part9zk-critical-bit-success-phase-diagram}，\(\mathcal S_\infty\) 在 \([0,1]\) 上连续且严格递增，而 \(p\in(0,1)\) 时有 \(s^\star(p)\in(0,1)\)。

任取
\[
0<\varepsilon<\min\{s^\star(p),1-s^\star(p)\}.
\]
由严格单调性，
\[
\mathcal S_\infty(s^\star(p)-\varepsilon)<p<\mathcal S_\infty(s^\star(p)+\varepsilon).
\]
再由推论 \ref{cor:xi-foldbin-critical-scale-success-curve-inversion} 的紧集上一致收敛，可知当 \(m\) 充分大时，
\[
\mathcal S_m(s^\star(p)-\varepsilon)<p<\mathcal S_m(s^\star(p)+\varepsilon).
\]
结合 \(\mathcal S_m(s_m^\star(p))=p\) 与 \(\mathcal S_m\) 的单调性，得到
\[
s^\star(p)-\varepsilon
<
s_m^\star(p)
\le
s^\star(p)+\varepsilon.
\]
令 \(\varepsilon\downarrow 0\)，便得
\[
s_m^\star(p)\longrightarrow s^\star(p).
\]
将其代回 \(b_m^\star(p)\) 的表达式，即得
\[
b_m^\star(p)
=
m\log_2\!\left(\frac{2}{\varphi}\right)
+
\log_2 s^\star(p)
+
o(1).
\]

显式公式则由定理 \ref{thm:xi-time-part9zk-critical-bit-success-phase-diagram} 中 \(\mathcal S_\infty\) 的分段线性表达逐段求逆而得。

最后，函数
\[
b\longmapsto 2^{-m}\,\mathcal C_m^{\mathrm{bin,cont}}(2^b)
\]
在 \([0,\infty)\) 上连续且单调不减，而 \(\mathrm{Succ}_m(B)\) 正是其在整数点 \(B\in\ZZ_{\ge 0}\) 的取值。因此达到水平 \(p\) 的最小整数预算恰为连续阈值的上整化：
\[
B_m^\star(p)=\left\lceil b_m^\star(p)\right\rceil.
\]
证毕。
\end{proof}

\begin{theorem}[六倍窗零点块的精确半尺度律]\label{thm:xi-time-part9zk-sixfold-zero-block-exact-halfscale}
\leanverified{paper\_xi\_time\_part9zk\_sixfold\_zero\_block\_exact\_halfscale}
记
\[
M_m:=F_{m+2}.
\]
在六倍窗子列
\[
m+2=2d\equiv 0\pmod 6
\]
上，零点块 \(\mathcal Z_m\) 的相对规模满足精确双边界
\[
\frac{1}{L_d}
\le
\frac{|\mathcal Z_m|}{M_m}
\le
\tau_{\mathrm{div}}(m+2)\,\frac{1}{L_d},
\]
其中 \(L_d\) 为第 \(d\) 个 Lucas 数，\(\tau_{\mathrm{div}}\) 为约数函数。并且
\[
\frac{1}{L_d}
=
\varphi^{-d}+O(\varphi^{-3d})
=
\varphi^{-(m+2)/2}
+
O\!\left(\varphi^{-3(m+2)/2}\right).
\]
特别地，
\[
\frac{|\mathcal Z_m|}{M_m}
=
\exp\!\left(
-\frac{m}{2}\log\varphi+o(m)
\right),
\]
等价地，
\[
\lim_{\substack{m\to\infty\\ m+2\equiv 0\ (6)}}
\frac1m\log\frac{M_m}{|\mathcal Z_m|}
=
\frac12\log\varphi.
\]
因此，在六倍窗上，零点块并非仅仅指数稀薄，而是以由 Lucas 因子精确锁定的半尺度斜率衰减。
\end{theorem}
\begin{proof}
由推论 \ref{cor:fold-zero-half-index-multiple6}，
\[
|\mathcal Z_m|\ge F_d.
\]
而 \(M_m=F_{m+2}=F_{2d}\)，并且 Fibonacci--Lucas 恒等式给出
\[
F_{2d}=F_dL_d.
\]
因此
\[
\frac{|\mathcal Z_m|}{M_m}
\ge
\frac{F_d}{F_{2d}}
=
\frac{1}{L_d}.
\]

另一方面，由定理 \ref{thm:fold-zero-density-sparse}，
\[
|\mathcal Z_m|
\le
\tau_{\mathrm{div}}(m+2)\,F_{\lfloor (m+2)/2\rfloor}
=
\tau_{\mathrm{div}}(m+2)\,F_d.
\]
再除以 \(M_m=F_{2d}=F_dL_d\)，便得
\[
\frac{|\mathcal Z_m|}{M_m}
\le
\tau_{\mathrm{div}}(m+2)\,\frac{1}{L_d}.
\]

Lucas 数的 Binet 公式为
\[
L_d=\varphi^d+(-\varphi)^{-d},
\]
故
\[
\frac{1}{L_d}
=
\varphi^{-d}\Bigl(1+O(\varphi^{-2d})\Bigr)^{-1}
=
\varphi^{-d}+O(\varphi^{-3d}).
\]
再利用
\[
\log \tau_{\mathrm{div}}(n)=o(n),
\]
便得到
\[
\frac{|\mathcal Z_m|}{M_m}
=
\exp\!\left(-\frac{m}{2}\log\varphi+o(m)\right),
\]
以及相应的对数极限。证毕。
\end{proof}

\begin{theorem}[六倍窗上的碰撞缺口与熵缺口共享 Lucas 半尺度缺口]\label{thm:xi-time-part9zk-sixfold-notch-collision-entropy}
沿用定理 \ref{thm:xi-time-part9zk-sixfold-zero-block-exact-halfscale} 的记号，并令 \(u_m\) 为 \(\ZZ/(M_m\ZZ)\) 上的均匀分布，
\[
p_m(r):=\frac{d_m(r)}{2^m}.
\]
则在
\[
m+2=2d\equiv 0\pmod 6
\]
上，有显式下界
\[
1-\mathsf{Col}_m
\ge
\frac{1}{L_d},
\]
以及
\[
\log M_m-D_{\mathrm{KL}}(p_m\|u_m)
\ge
-\log\!\left(1-\frac{1}{L_d}\right)
\ge
\frac{1}{L_d}.
\]
特别地，两种缺口都被同一个 Lucas 半尺度
\[
\frac{1}{L_d}
=
\varphi^{-d}+O(\varphi^{-3d})
\]
从下方共同钉扎。换言之，六倍窗上的算术零点块在碰撞与熵两侧同时强制注入一个确定性的半尺度缺口。
\end{theorem}
\begin{proof}
由定理 \ref{thm:xi-time-part60ab-window6-zero-block-halfscale-information-improvement}，
\[
1-\mathsf{Col}_m
\ge
\frac{F_d}{F_{m+2}},
\qquad
\log F_{m+2}-D_{\mathrm{KL}}(p_m\|u_m)
\ge
\log\frac{F_{m+2}}{F_{m+2}-F_d}.
\]
由于 \(F_{m+2}=F_{2d}=F_dL_d\)，故
\[
\frac{F_d}{F_{m+2}}=\frac{1}{L_d},
\]
以及
\[
\log\frac{F_{m+2}}{F_{m+2}-F_d}
=
-\log\!\left(1-\frac{1}{L_d}\right).
\]
又因对 \(0\le x<1\) 恒有
\[
-\log(1-x)\ge x,
\]
遂得
\[
\log M_m-D_{\mathrm{KL}}(p_m\|u_m)
\ge
\frac{1}{L_d}.
\]
最后代入定理 \ref{thm:xi-time-part9zk-sixfold-zero-block-exact-halfscale} 中
\[
\frac{1}{L_d}=\varphi^{-d}+O(\varphi^{-3d}),
\]
即得结论。证毕。
\end{proof}

\begin{theorem}[典型输出密度附近的四阶 Cram\'er--Legendre 展开]\label{thm:xi-time-part9zk-output-cramer-legendre-fourth-order}
沿用定理 \ref{thm:xi-real-input-cramer-germ-quadratic-field-rigidity} 的记号。记
\[
P(\theta)=\log\lambda(e^\theta),
\qquad
\alpha_0:=P'(0),
\qquad
\kappa_j:=P^{(j)}(0)\ \ (j\ge 2),
\]
并令 \(I(\alpha)\) 为其 Legendre 对偶速率函数。则当 \(\delta\to 0\) 时，
\[
I(\alpha_0+\delta)
=
\frac{\delta^2}{2\kappa_2}
-\frac{\kappa_3}{6\kappa_2^3}\delta^3
+
\frac{3\kappa_3^2-\kappa_2\kappa_4}{24\kappa_2^5}\delta^4
+
O(\delta^5).
\]
将定理 \ref{thm:xi-real-input-cramer-germ-quadratic-field-rigidity} 中的闭式常数
\[
\kappa_2=\frac{1791}{200}-\frac{3983\sqrt5}{1000},
\qquad
\kappa_3=\frac{1416369}{5000}-\frac{126691\sqrt5}{1000},
\]
\[
\kappa_4=\frac{1354428303}{100000}-\frac{605718497\sqrt5}{100000}
\]
代入，可得数值展开
\[
I(\alpha_0+\delta)
=
10.2582524032149\,\delta^2
+
22.8681939097743\,\delta^3
+
72.7346495328052\,\delta^4
+
O(\delta^5).
\]
因此，Gaussian 二次主项之后的首个非平凡修正已经在三阶出现，并且四阶系数仍保持显著正值。
\end{theorem}
\begin{proof}
第一式正是定理 \ref{thm:xi-real-input-cramer-germ-quadratic-field-rigidity} 的局部展开。将其中给出的 \(\kappa_2,\kappa_3,\kappa_4\) 闭式代入
\[
\frac{1}{2\kappa_2},
\qquad
-\frac{\kappa_3}{6\kappa_2^3},
\qquad
\frac{3\kappa_3^2-\kappa_2\kappa_4}{24\kappa_2^5},
\]
直接化简即可得到所示数值系数。证毕。
\end{proof}

\begin{corollary}[典型点附近的大偏差左右不对称性]\label{cor:xi-time-part9zk-local-rate-left-right-asymmetry}
沿用定理 \ref{thm:xi-time-part9zk-output-cramer-legendre-fourth-order} 的记号。由于
\[
\kappa_3<0,
\]
故当 \(\delta\downarrow 0\) 时有
\[
I(\alpha_0+\delta)-I(\alpha_0-\delta)
=
-\frac{\kappa_3}{3\kappa_2^3}\,\delta^3
+
O(\delta^5).
\]
特别地，对充分小的 \(\delta>0\)，成立严格不等式
\[
I(\alpha_0+\delta)>I(\alpha_0-\delta).
\]
因此，在典型输出密度附近，向上偏离比同量级的向下偏离具有更大的指数代价；局部速率函数并不以 \(\alpha_0\) 为偶对称中心。
\end{corollary}
\begin{proof}
把定理 \ref{thm:xi-time-part9zk-output-cramer-legendre-fourth-order} 的展开分别代入 \(\delta\) 与 \(-\delta\)。偶次项相消，奇次项翻转，从而
\[
I(\alpha_0+\delta)-I(\alpha_0-\delta)
=
2\left(-\frac{\kappa_3}{6\kappa_2^3}\right)\delta^3
+
O(\delta^5)
=
-\frac{\kappa_3}{3\kappa_2^3}\,\delta^3
+
O(\delta^5).
\]
又因 \(\kappa_2>0\) 且 \(\kappa_3<0\)，故主导系数为正，结论成立。证毕。
\end{proof}

\begin{theorem}[escort 温度族的两尺度截断律]\label{thm:xi-time-part9zk-escort-two-scale-clip-law}
对任意 \(q\ge 0\)，令
\[
W_m^{(q)}\sim \pi_{m,q},
\]
并定义尺度化截断泛函
\[
\mathcal C_{m,q}^{\mathrm{clip}}(s)
:=
\mathbf E_{\pi_{m,q}}\!\left[
\min\!\left\{
\left(\frac{\varphi}{2}\right)^m d_m^{\mathrm{bin}}(W_m^{(q)}),\,s
\right\}
\right],
\qquad
s\ge 0.
\]
则当 \(m\to\infty\) 时，对每个固定 \(s\ge 0\) 都有
\[
\mathcal C_{m,q}^{\mathrm{clip}}(s)
\longrightarrow
\mathcal C_\infty^{(q)}(s),
\]
其中
\[
\mathcal C_\infty^{(q)}(s)
=
\bigl(1-\theta(q)\bigr)\min\{1,s\}
+
\theta(q)\min\{\varphi^{-1},s\},
\qquad
\theta(q)=\frac{1}{1+\varphi^{q+1}}.
\]
特别地，对所有温度 \(q\) 而言，折点始终固定在
\[
s=\varphi^{-1},
\qquad
s=1,
\]
而 \(q\) 仅通过 \(\theta(q)\) 改变这两条刚性尺度支之间的权重分配，而不改变其位置。
\end{theorem}
\begin{proof}
由定理 \ref{thm:xi-time-part9ob-escort-quantile-phase-diagram} 的证明，缩放随机变量
\[
Z_m^{(q)}
:=
\left(\frac{\varphi}{2}\right)^m d_m^{\mathrm{bin}}(W_m^{(q)})
\]
满足弱收敛
\[
Z_m^{(q)}\Longrightarrow Z^{(q)},
\]
其中
\[
\mathbf P\bigl(Z^{(q)}=1\bigr)=1-\theta(q),
\qquad
\mathbf P\bigl(Z^{(q)}=\varphi^{-1}\bigr)=\theta(q).
\]
对固定 \(s\ge 0\)，函数
\[
x\longmapsto \min\{x,s\}
\]
在 \([0,\infty)\) 上有界且连续。因此由弱收敛与连续映射定理，
\[
\mathcal C_{m,q}^{\mathrm{clip}}(s)
=
\mathbf E\!\left[\min\{Z_m^{(q)},s\}\right]
\longrightarrow
\mathbf E\!\left[\min\{Z^{(q)},s\}\right].
\]
而右端对该两点分布直接计算即得
\[
\mathbf E\!\left[\min\{Z^{(q)},s\}\right]
=
\bigl(1-\theta(q)\bigr)\min\{1,s\}
+
\theta(q)\min\{\varphi^{-1},s\}.
\]
证毕。
\end{proof}

\noindent
由此，critical-scale 预言机恢复、六倍窗零点块、典型密度附近的局部大偏差以及 escort 温度重加权之间又出现了一条新的共同骨架：在资源侧，成功率相图与预算反演都由 \(\{\varphi^{-1},1\}\) 两个折点完全决定；在频域侧，六倍窗零点块把碰撞缺口与熵缺口共同钉扎在 Lucas 半尺度 \(L_d^{-1}\) 上；在输出位势侧，局部速率函数的三阶项已经强制右尾代价严格大于左尾；而在 escort 侧，任意温度 \(q\) 只会重新分配两条固定尺度支的权重，而不会移动其位置。于是，阈值反演、暗带稀疏、非高斯偏斜与温度重加权四条接口便在同一组黄金双尺度与半尺度缺口律之下闭合。

\endinput

\paragraph{链上算术高度、有限局部化核正交、双域中心去相关与熵账本线段可达域}
在链上内算子的平方自由算术化、有限局部化 solenoid 的 prime-register 核、Lee--Yang / RH 双域的 Chebotarev 乘法律以及固定宏观边缘下的最大熵提升账本已经建立之后，还可进一步抽取出下列六条彼此兼容的后果。它们依次把链上规范投影的格高度完全译为 Gödel 整数的素因子高度，把任意 \(\gcd/\mathrm{lcm}\) 编译器的最坏寄存器规模压到 primorial 量级，把有限局部化核之间的连续通道收束为严格的同素数对角块，并把 Frobenius 中心可观测量与 lift 族的熵--相对熵可达域分别压缩为完全去相关律与一维仿射线段。

\begin{theorem}[链上规范投影的完全算术化与精确高度公式]\label{thm:xi-time-part9zl-chain-arithmetic-height-formula}
\leanverified{paper\_xi\_time\_part9zl\_chain\_arithmetic\_height\_formula}
沿用定理 \ref{thm:killo-chain-interior-godel-gcd-lcm} 与定理 \ref{thm:xi-chain-interior-boolean-flag-closed-form} 的记号。记
\[
\mathcal I_n:=\mathrm{Int}(C_n),
\qquad
J_n:=\prod_{i=0}^{n-1}q_i,
\qquad
\mathcal D_n:=\{d\mid J_n:\ q_0\mid d\},
\]
并以
\[
\omega_{\mathrm{arith}}(N):=\#\{p\ \text{prime}: p\mid N\}
\]
表示整数 \(N\) 的不同素因子个数。则映射
\[
\kappa:\bigl(\mathcal I_n,\wedge,\vee,\preceq\bigr)
\xrightarrow{\sim}
\bigl(\mathcal D_n,\gcd,\mathrm{lcm},\mid\bigr)
\]
为格同构。并且对任意 \(I\in\mathcal I_n\)，其在格 \((\mathcal I_n,\preceq)\) 中的高度满足精确公式
\[
\operatorname{ht}(I)
=
|\mathrm{Fix}(I)|-1
=
\omega_{\mathrm{arith}}(\kappa(I))-1.
\]
\end{theorem}
\begin{proof}
由定理 \ref{thm:killo-chain-interior-godel-gcd-lcm}，映射
\[
F\longmapsto I_F
\]
给出
\[
\mathcal F_n:=\{F\subseteq[n]:0\in F\}
\]
与 \(\mathcal I_n\) 的格同构，且
\[
\mathrm{Fix}(I_F)=F,
\qquad
\kappa(I_F)=\prod_{i\in F}q_i.
\]
由于 \(J_n\) 本身平方自由，\(\kappa\) 的像恰为所有含素因子 \(q_0\) 的因子，即 \(\mathcal D_n\)。并且 meet/join 在 \(\mathcal F_n\) 上对应集合交并，而在 \(\mathcal D_n\) 上对应 \(\gcd/\mathrm{lcm}\)，故 \(\kappa\) 为格同构。

另一方面，定理 \ref{thm:xi-chain-interior-boolean-flag-closed-form} 已表明 \((\mathcal I_n,\preceq)\) 是秩 \(n-1\) 的 Boolean 格，且其秩函数即
\[
\operatorname{rk}(I)=|\mathrm{Fix}(I)|-1.
\]
在有限分次格中，元素自最小元起的高度等于其秩，故
\[
\operatorname{ht}(I)=|\mathrm{Fix}(I)|-1.
\]
又因 \(\kappa(I)\) 为平方自由整数，且其素因子恰由 \(\mathrm{Fix}(I)\) 给出，遂有
\[
\omega_{\mathrm{arith}}(\kappa(I))=|\mathrm{Fix}(I)|.
\]
三式合并即得所述精确高度公式。证毕。
\end{proof}

\begin{theorem}[任意 \texorpdfstring{$\gcd/\mathrm{lcm}$}{gcd/lcm} 编译器的 primorial 下界]\label{thm:xi-time-part9zl-gcd-lcm-compiler-primorial-lower-bound}
沿用定理 \ref{thm:killo-chain-interior-godel-gcd-lcm} 的记号，记
\[
\mathcal I_n:=\mathrm{Int}(C_n).
\]
设单射
\[
\eta:\mathcal I_n\hookrightarrow \NN_{\ge 1}
\]
满足对任意 \(I,J\in\mathcal I_n\) 都有
\[
\eta(I\wedge J)=\gcd(\eta(I),\eta(J)),
\qquad
\eta(I\vee J)=\mathrm{lcm}(\eta(I),\eta(J)).
\]
若 \(p_1<p_2<\cdots\) 为递增素数序列，则必有
\[
\max_{I\in\mathcal I_n}\eta(I)\ \ge\ \prod_{j=1}^{n-1}p_j.
\]
因而
\[
\max_{I\in\mathcal I_n}\log \eta(I)
\ \ge\
\sum_{j=1}^{n-1}\log p_j
=
\vartheta(p_{n-1})
=
(1+o(1))\,n\log n.
\]
\leanverified{paper\_xi\_time\_part9zl\_gcd\_lcm\_compiler\_primorial\_lower\_bound}
\end{theorem}
\begin{proof}
当 \(n=1\) 时，\(\mathcal I_1\) 仅含一个元素，结论平凡。以下设 \(n\ge 2\)。

记最小元与原子为
\[
\widehat 0:=I_{\{0\}},
\qquad
A_i:=I_{\{0,i\}}
\qquad(1\le i\le n-1).
\]
由 \(\widehat 0\wedge I=\widehat 0\) 可得
\[
\eta(\widehat 0)=\gcd(\eta(\widehat 0),\eta(I)),
\]
故 \(\eta(\widehat 0)\mid \eta(I)\) 对所有 \(I\in\mathcal I_n\) 成立。于是
\[
\widetilde\eta(I):=\frac{\eta(I)}{\eta(\widehat 0)}\in\NN
\]
定义良好；而且 \(\widetilde\eta\) 仍为单射，并继续保持 \(\gcd/\mathrm{lcm}\)。

由定理 \ref{thm:killo-chain-interior-godel-gcd-lcm}，\(\mathcal I_n\) 是包含 \(0\) 的子集 Boolean 格，故 \(A_1,\dots,A_{n-1}\) 正是全部原子，并满足
\[
A_i\wedge A_j=\widehat 0
\qquad(i\neq j).
\]
于是
\[
\gcd\bigl(\widetilde\eta(A_i),\widetilde\eta(A_j)\bigr)=1
\qquad(i\neq j).
\]
又因 \(\widetilde\eta\) 单射且 \(A_i\neq \widehat 0\)，故 \(\widetilde\eta(A_i)>1\)。因此可为每个 \(i\) 选取素数
\[
r_i\mid \widetilde\eta(A_i).
\]
上式的两两互素性迫使 \(r_1,\dots,r_{n-1}\) 两两不同。

再记最大元为
\[
\widehat 1:=I_{[n]}.
\]
由于
\[
\widehat 1=A_1\vee\cdots\vee A_{n-1},
\]
保持 join 的性质给出
\[
\widetilde\eta(\widehat 1)
=
\mathrm{lcm}\bigl(\widetilde\eta(A_1),\dots,\widetilde\eta(A_{n-1})\bigr).
\]
故 \(\widetilde\eta(\widehat 1)\) 必被 \(\prod_{i=1}^{n-1}r_i\) 整除，从而
\[
\max_{I\in\mathcal I_n}\eta(I)
\ge
\eta(\widehat 1)
\ge
\widetilde\eta(\widehat 1)
\ge
\prod_{i=1}^{n-1}r_i.
\]
而 \(n-1\) 个两两不同素数的最小乘积正是 \(\prod_{j=1}^{n-1}p_j\)，故
\[
\max_{I\in\mathcal I_n}\eta(I)\ge \prod_{j=1}^{n-1}p_j.
\]
取对数后，再用 Chebyshev 函数的素数定理形式
\[
\sum_{j=1}^{n-1}\log p_j=\vartheta(p_{n-1})=(1+o(1))\,n\log n
\]
即得所述渐近下界。证毕。
\end{proof}

\begin{theorem}[有限局部化 solenoid 核之间的素数正交律]\label{thm:xi-time-part9zl-localized-solenoid-prime-orthogonality}
对任意有限素数集 \(S,T\subset\mathbb P\)，记
\[
\mathfrak K_S:=\prod_{p\in S}\ZZ_p,
\qquad
\mathfrak K_T:=\prod_{q\in T}\ZZ_q.
\]
则存在自然同构
\[
\operatorname{Hom}_{\mathrm{cont}}(\mathfrak K_S,\mathfrak K_T)
\cong
\prod_{p\in S\cap T}\ZZ_p.
\]
特别地，若 \(S\cap T=\varnothing\)，则
\[
\operatorname{Hom}_{\mathrm{cont}}(\mathfrak K_S,\mathfrak K_T)=0.
\]
\end{theorem}
\begin{proof}
由于 \(S,T\) 有限，任意连续同态
\[
f:\mathfrak K_S\to\mathfrak K_T
\]
由其坐标分量
\[
f_{q,p}:=\operatorname{pr}_q\circ f\circ \operatorname{in}_p:\ZZ_p\to\ZZ_q
\qquad(p\in S,\ q\in T)
\]
唯一决定，从而
\[
\operatorname{Hom}_{\mathrm{cont}}(\mathfrak K_S,\mathfrak K_T)
\cong
\prod_{q\in T}\prod_{p\in S}
\operatorname{Hom}_{\mathrm{cont}}(\ZZ_p,\ZZ_q).
\]

对 \(p\neq q\)，定理 \ref{thm:cdim-star-prime-coordinate-orthogonality} 已给出
\[
\operatorname{Hom}_{\mathrm{cont}}(\ZZ_p,\ZZ_q)=0.
\]
因此仅剩对角块 \(p=q\)。

现取任意连续加法自同态
\[
u:\ZZ_p\to\ZZ_p.
\]
记 \(a:=u(1)\in\ZZ_p\)。则对任意整数 \(n\in\ZZ\) 都有
\[
u(n)=na.
\]
另一方面，\(\ZZ\) 在 \(\ZZ_p\) 中稠密，而乘以 \(a\) 的映射
\[
m_a:\ZZ_p\to\ZZ_p,\qquad x\mapsto ax
\]
连续，且与 \(u\) 在稠密子群 \(\ZZ\) 上一致，故由连续性得
\[
u(x)=ax
\qquad(\forall x\in\ZZ_p).
\]
这表明
\[
\operatorname{End}_{\mathrm{cont}}(\ZZ_p)\cong \ZZ_p.
\]

将以上各素数方向合并，即得
\[
\operatorname{Hom}_{\mathrm{cont}}(\mathfrak K_S,\mathfrak K_T)
\cong
\prod_{p\in S\cap T}\operatorname{End}_{\mathrm{cont}}(\ZZ_p)
\cong
\prod_{p\in S\cap T}\ZZ_p.
\]
若 \(S\cap T=\varnothing\)，右端为空积，故同态群为零。证毕。
\end{proof}

\begin{corollary}[有限局部化的圆维塌缩与支撑刚性]\label{cor:xi-time-part9zl-localization-circle-collapse-support-rigidity}
对任意有限非空素数集 \(S,T\subset\mathbb P\)，记
\[
G_S:=\ZZ[S^{-1}],
\qquad
\Sigma_S:=\widehat{G_S}.
\]
则有
\[
\cdim_{\mathrm{fr}}(G_S)=\cdim_{\mathrm{fr}}(G_T)=1.
\]
与此同时，
\[
\Sigma_S\cong \Sigma_T
\quad\Longleftrightarrow\quad
S=T.
\]
\end{corollary}
\begin{proof}
定理 \ref{thm:killo-localization-circle-dimension-prime-ledger-splitting} 给出
\[
\cdim_{\mathrm{fr}}(G_S)=1
\]
对每个有限非空 \(S\subset\mathbb P\) 恒成立，并给出短正合列
\[
0\longrightarrow \prod_{p\in S}\ZZ_p
\longrightarrow \Sigma_S
\longrightarrow \TT
\longrightarrow 0.
\]
同一定理还说明核的 \(p\)-主因子非平凡当且仅当 \(p\in S\)，故该核完全恢复 \(S\)。于是
\[
\Sigma_S\cong \Sigma_T
\]
必推出二者核同构，继而推出 \(S=T\)；反向蕴含显然。证毕。
\end{proof}

\begin{theorem}[Lee--Yang / RH 双分歧域上的中心可观测量完全去相关]\label{thm:xi-time-part9zl-leyang-rh-central-observable-decorrelation}
沿用命题 \ref{prop:fold-gauge-anomaly-P10-leyang-linear-disjointness} 的记号。对每个共同不分歧素数 \(p\)，记
\[
\sigma_{10}(p):=\mathrm{Frob}_p(L_{10}/\QQ)\in S_{10},
\qquad
\sigma_3(p):=\mathrm{Frob}_p(L_{\mathrm{LY}}/\QQ)\in S_3,
\]
并记
\[
\pi_{\mathrm{ur}}(x)
:=
\#\{p\le x:\ p\nmid \Disc(P_{10})\Disc(P_{\mathrm{LY}})\}.
\]
对任意类函数
\[
f:S_{10}\to\CC,
\qquad
g:S_3\to\CC,
\]
定义
\[
\mathbb E_x[f,g]
:=
\frac{1}{\pi_{\mathrm{ur}}(x)}
\sum_{\substack{p\le x\\ p\nmid \Disc(P_{10})\Disc(P_{\mathrm{LY}})}}
f\bigl(\sigma_{10}(p)\bigr)\,
g\bigl(\sigma_3(p)\bigr).
\]
则极限存在且满足
\[
\lim_{x\to\infty}\mathbb E_x[f,g]
=
\left(\frac{1}{10!}\sum_{\tau\in S_{10}}f(\tau)\right)
\left(\frac{1}{6}\sum_{\eta\in S_3}g(\eta)\right).
\]
特别地，若
\[
\sum_{\tau\in S_{10}}f(\tau)=0,
\qquad
\sum_{\eta\in S_3}g(\eta)=0,
\]
则有
\[
\lim_{x\to\infty}\mathbb E_x[f,g]=0.
\]
\end{theorem}
\begin{proof}
把 \(f\) 与 \(g\) 分别按共轭类指标函数展开：
\[
f=\sum_C a_C\,\mathbf 1_C,
\qquad
g=\sum_D b_D\,\mathbf 1_D,
\]
其中 \(C\subset S_{10}\)、\(D\subset S_3\) 取遍共轭类。于是
\[
f\bigl(\sigma_{10}(p)\bigr)\,g\bigl(\sigma_3(p)\bigr)
=
\sum_{C,D}a_Cb_D\,
\mathbf 1_{\{\sigma_{10}(p)\in C,\ \sigma_3(p)\in D\}}.
\]
对共同不分歧素数取平均并交换有限求和后，问题归结为各联合指标事件的自然密度。由推论 \ref{cor:fold-gauge-anomaly-P10-leyang-chebotarev-product-law}，
\[
\delta\bigl(\sigma_{10}(p)\in C,\ \sigma_3(p)\in D\bigr)
=
\frac{|C|}{10!}\cdot \frac{|D|}{6}.
\]
故
\[
\lim_{x\to\infty}\mathbb E_x[f,g]
=
\sum_{C,D}a_Cb_D\frac{|C|}{10!}\frac{|D|}{6}
=
\left(\sum_C a_C\frac{|C|}{10!}\right)
\left(\sum_D b_D\frac{|D|}{6}\right),
\]
这恰是所述群平均的乘积。若二者群平均均为零，则右端立即退化为 \(0\)。证毕。
\end{proof}

\begin{theorem}[固定宏观边缘下的熵--相对熵可达域恰为一条闭线段]\label{thm:xi-time-part9zl-lift-entropy-kl-line-segment}
固定 \(m\ge 1\) 与宏观边缘分布 \(\pi\in\Delta(X_m)\)。记
\[
\mathcal L_m(\pi)
:=
\{\mu\in\Delta(\Omega_m):(\Fold_m)_*\mu=\pi\},
\]
并以 \(\pi^{\mathrm e}\) 表示 \eqref{eq:xi-null-maxent-lift} 中的纤维均匀 lift。则有
\[
\bigl\{
\bigl(H(\mu),D_{\mathrm{KL}}(\mu\|\pi^{\mathrm e})\bigr):
\mu\in\mathcal L_m(\pi)
\bigr\}
\]
\[
=
\Bigl\{
\bigl(h,\ H(\pi)+\kappa_m(\pi)-h\bigr):
\ H(\pi)\le h\le H(\pi)+\kappa_m(\pi)
\Bigr\}.
\]
等价地，
\[
\{D_{\mathrm{KL}}(\mu\|\pi^{\mathrm e}):\mu\in\mathcal L_m(\pi)\}
=
[0,\kappa_m(\pi)],
\]
并且
\[
\{H(\mu):\mu\in\mathcal L_m(\pi)\}
=
[H(\pi),\,H(\pi)+\kappa_m(\pi)].
\]
\end{theorem}
\begin{proof}
由定理 \ref{thm:xi-null-fiber-entropy-window}，对每个 \(\mu\in\mathcal L_m(\pi)\) 都有严格账本恒等式
\[
H(\mu)=H(\pi)+\kappa_m(\pi)-D_{\mathrm{KL}}(\mu\|\pi^{\mathrm e}).
\]
因此任一可达点都落在仿射直线
\[
h+d=H(\pi)+\kappa_m(\pi)
\]
上。

同一定理还给出两个端点的可达性：纤维均匀 lift \(\pi^{\mathrm e}\) 唯一实现最大熵，并满足
\[
H(\pi^{\mathrm e})=H(\pi)+\kappa_m(\pi),
\qquad
D_{\mathrm{KL}}(\pi^{\mathrm e}\|\pi^{\mathrm e})=0;
\]
同时存在纤维点质量 lift \(\mu_\star\in\mathcal L_m(\pi)\) 使
\[
H(\mu_\star)=H(\pi),
\qquad
D_{\mathrm{KL}}(\mu_\star\|\pi^{\mathrm e})=\kappa_m(\pi).
\]
故上述仿射线段的两个端点均被实现。

再定义
\[
\mu_t:=(1-t)\mu_\star+t\pi^{\mathrm e}
\qquad(0\le t\le 1).
\]
由于 \(\mathcal L_m(\pi)\) 是概率单纯形中的仿射切片，故其为凸集，从而 \(\mu_t\in\mathcal L_m(\pi)\)。另一方面，Shannon 熵在有限单纯形上连续，故
\[
t\longmapsto H(\mu_t)
\]
连续，并取遍连接
\[
H(\mu_\star)=H(\pi)
\qquad\text{与}\qquad
H(\pi^{\mathrm e})=H(\pi)+\kappa_m(\pi)
\]
的整个闭区间。再代回账本恒等式，即得
\[
D_{\mathrm{KL}}(\mu_t\|\pi^{\mathrm e})
=
H(\pi)+\kappa_m(\pi)-H(\mu_t),
\]
从而可达点恰充满整条闭线段。证毕。
\end{proof}

上述六条后果表明，链上 Boolean 算术化、有限局部化的 prime-ledger 核、双域 Frobenius 指纹以及固定边缘 lift 的熵账本虽然分属组合、对偶、Chebotarev 与信息论四个层面，其末端结构却都压缩为一类一维或对角对象：高度由素因子数精确给出，连续通道仅沿同素数方向保留，中心相关在极限平均下完全消失，而全部可达热力学域则退化为一条闭线段。

\paragraph{歧义壳精确记忆指数、零谱一阶激活与解析域内部的 RH 相变}
在确定化歧义壳的 DAG 瞬态、真实输入 \(40\) 状态核的零谱 Jordan 结构、碰撞位势的核版 RH/Ramanujan 窗口以及模 \(2\) 扭曲的最小证人性质已经建立之后，还可进一步抽取出下列八条彼此兼容的后果。前三条把 sofic 歧义壳从“有限同步延迟”压缩为精确幂零指数、精确有限记忆深度与最小统一滞后的三联锁；其后五条则把零谱扇区的加权激活、RH 阈值与解析半径的错位，以及最小奇偶角色触发的平方根律失效统一为一条发生在解析域内部的纯谱相变链。

\begin{theorem}[歧义壳邻接块的精确幂零指数]\label{thm:xi-time-part9zm-ambiguity-shell-exact-nilpotent-index}
\leanverified{paper\_xi\_time\_part9zm\_ambiguity\_shell\_exact\_nilpotent\_index}
沿用定理 \ref{thm:Ym-ambiguity-shell-zeta-factor-trivial} 的记号。对每个 \(m\ge 3\)，把确定化右解析表示 \(H_m^{\mathrm{det}}\) 的状态按“非单点歧义态在前、单点本质态在后”的顺序排列，则其邻接矩阵可写为
\[
A_m^{\mathrm{det}}
=
\begin{pmatrix}
N_m & R_m\\
0 & A_m^{\mathrm{ess}}
\end{pmatrix}.
\]
则有
\[
N_m^{\,m}=0,
\qquad
N_m^{\,m-1}\neq 0.
\]
因而歧义壳块 \(N_m\) 的幂零指数恰为 \(m\)。
\end{theorem}
\begin{proof}
定理 \ref{thm:Ym-ambiguity-shell-zeta-factor-trivial} 已直接给出
\[
N_m^{\,m}=0.
\]
因此只需证明 \(N_m^{\,m-1}\neq 0\)。

由推论 \ref{cor:Ym-sync-delay}，
\[
D_{\mathrm{sync}}(m)=m-1.
\]
故存在某个 \(y\in Y_m\)，使得从初始候选集 \(S_0:=V_m\) 出发递推
\[
S_{t+1}:=\delta_m(S_t,y_t)
\]
时满足
\[
|S_{m-1}|>1.
\]
而命题 \ref{prop:det-candidate-set-monotone} 说明候选集基数随时间单调不增，因此对每个 \(0\le t\le m-1\) 都有
\[
|S_t|>1.
\]
换言之，
\[
S_0\to S_1\to\cdots\to S_{m-1}
\]
是一条完全落在歧义态诱导子图 \(H_m^{\mathrm{amb}}\) 内的长度 \(m-1\) 路径。于是邻接块 \(N_m\) 的某个 \((m-1)\)-步通达项非零，即
\[
N_m^{\,m-1}\neq 0.
\]
两式合并即得幂零指数恰为 \(m\)。证毕。
\end{proof}

\begin{corollary}[瞬态修正的精确记忆深度]\label{cor:xi-time-part9zm-ambiguity-shell-exact-memory-depth}
在定理 \ref{thm:xi-time-part9zm-ambiguity-shell-exact-nilpotent-index} 的设定下，有
\[
(I-zN_m)^{-1}
=
\sum_{j=0}^{m-1} z^jN_m^j,
\]
并且该矩阵多项式的次数恰为 \(m-1\)。此外，
\[
(I-zA_m^{\mathrm{det}})^{-1}
=
\begin{pmatrix}
(I-zN_m)^{-1} &
(I-zN_m)^{-1}\,zR_m\,(I-zA_m^{\mathrm{ess}})^{-1}\\[4pt]
0 &
(I-zA_m^{\mathrm{ess}})^{-1}
\end{pmatrix}.
\]
因此，歧义壳对 resolvent 的全部修正都压缩为一段且仅一段长度为 \(m-1\) 的有限记忆前缀；它不产生任何新极点。
\end{corollary}
\begin{proof}
由定理 \ref{thm:xi-time-part9zm-ambiguity-shell-exact-nilpotent-index}，
\[
N_m^{\,m}=0,
\qquad
N_m^{\,m-1}\neq 0.
\]
因此 Neumann 级数在第 \(m\) 项处严格终止，给出
\[
(I-zN_m)^{-1}
=
\sum_{j=0}^{m-1} z^jN_m^j,
\]
且由于最高项 \(N_m^{\,m-1}\) 非零，其次数不能再降。

另一方面，
\[
I-zA_m^{\mathrm{det}}
=
\begin{pmatrix}
I-zN_m & -zR_m\\
0 & I-zA_m^{\mathrm{ess}}
\end{pmatrix}
\]
是块上三角矩阵。对其应用标准块逆公式，即得所示表达式。

最后，定理 \ref{thm:Ym-ambiguity-shell-zeta-factor-trivial} 已给出
\[
\det(I-zA_m^{\mathrm{det}})=\det(I-zA_m^{\mathrm{ess}}).
\]
故歧义壳只通过有限多项式因子修正 resolvent，而不会改变其极点集合。证毕。
\end{proof}

\begin{theorem}[有限状态右逆的最小统一滞后]\label{thm:xi-time-part9zm-finite-state-right-inverse-minimal-lag}
设 \(m\ge 3\)。称有限状态转导器 \(\mathcal T\) 为 \(\Phi_m\) 的一个统一滞后为 \(\ell\) 的右逆，若对任意 \(y\in Y_m\) 与任意 \(t\ge 0\)，它在读入输出前缀
\[
y_0,\dots,y_{t+\ell}
\]
后即可唯一恢复与 \(y\) 对应的输入序列 \(b\) 的第 \(t\) 位 \(b_t\)。则最小可能统一滞后满足
\[
\ell_{\min}(m)=m-1.
\]
\end{theorem}
\begin{proof}
\textbf{上界。}
定理 \ref{thm:Phi_m-conjugacy-threshold} 已给出一个逆码
\[
\Psi_m:Y_m\to\{0,1\}^{\mathbb Z},
\]
其记忆为 \(m-1\)、预期性为 \(0\)。这意味着：在双边移位口径下，每一位输入都由一个长度 \(m\) 的输出窗口决定；而在单边在线解码口径下，初始时缺失的左侧 \(m-1\) 个坐标恰可等价地转写为统一启动滞后 \(m-1\)。再结合推论 \ref{cor:Ym-singleton-essential-finite-delay}，一旦轨道进入单点本质分量，后续各位都可零额外预期性在线恢复。故存在一个统一滞后为 \(m-1\) 的有限状态右逆，从而
\[
\ell_{\min}(m)\le m-1.
\]

\textbf{下界。}
同一定理的证明以及表 \ref{tab:phi_m_conjugacy_threshold} 明确给出 sharpness 证书：当 \(m\ge 3\) 时，
\[
\text{non-singleton ends }(m-1)>0,
\]
并因此“通常需要完整的长度 \(m\) 输出块”才能实现唯一 lift。换言之，长度 \(m-1\) 的输出窗口一般不足以完成统一反演。若存在某个 \(\ell<m-1\) 的统一滞后右逆，则所有坐标的恢复都只需长度 \(\ell+1\le m-1\) 的输出窗口，这与上述 sharpness 证书矛盾。故
\[
\ell_{\min}(m)\ge m-1.
\]
综上
\[
\ell_{\min}(m)=m-1.
\]
证毕。
\end{proof}

\begin{theorem}[零谱扇区的一阶激活律]\label{thm:xi-time-part9zm-zero-sector-first-order-activation}
设 \(M(u)\) 为真实输入 \(40\) 状态核在输出位势下的一参数加权矩阵。则在无权点 \(u=1\) 处，零谱扇区并非高阶惰性残余，而是可被任意非平凡权重立即激活的一条真实谱通道：存在特征值分支 \(\lambda_{\mathrm{act}}(u)\) 满足
\[
\lambda_{\mathrm{act}}(1)=0,
\qquad
\lambda_{\mathrm{act}}(u)=(u-1)+O\!\bigl((u-1)^2\bigr),
\]
从而
\[
\lambda_{\mathrm{act}}(u)\sim u-1
\qquad (u\to 1).
\]
特别地，该分支以一阶而非更高阶方式从零点脱离。
\end{theorem}
\begin{proof}
命题 \ref{prop:real-input-40-nilpotent-index} 表明：在 \(u=1\) 的无权核上，\(0\)-谱的不可对角化性全部集中于零特征值，其最大 Jordan 块大小在 essential 核上为 \(3\)，在全 \(40\) 状态核上为 \(5\)。这说明零谱扇区在无权点确实承载全部瞬态幂零结构。

另一方面，备注 \ref{rem:real-input-40-zero-activation} 指出：在
\[
\Delta(z,u)=(1-uz^2)F(z,u)
\]
中，\(F\) 的最高次项系数为
\[
u^3(1-u)z^8.
\]
因此 \(u=1\) 时该项消失，次数由 \(8\) 降为 \(7\)；而 \(u\neq 1\) 时次数恢复为 \(8\)。这已表明无权核中的一条零特征方向在加权后被激活为非零分支。

更精确地，命题 \ref{prop:real-input-40-activated-branch-linear} 已给出唯一小根分支
\[
\lambda_{\mathrm{act}}(u)
\]
满足
\[
\lambda_{\mathrm{act}}(u)=(u-1)+O\!\bigl((u-1)^2\bigr).
\]
故其临界指数恰为 \(1\)，即
\[
\lambda_{\mathrm{act}}(u)\sim u-1.
\]
证毕。
\end{proof}

\begin{theorem}[RH 阈值先于解析障碍]\label{thm:xi-time-part9zm-rh-threshold-before-analytic-obstruction}
对碰撞位势压力
\[
P_2(\theta):=\log\rho(e^\theta),
\]
记其核版 RH 临界点为
\[
\theta_R:=\log u_R,
\]
其中 \(u_R\) 为命题 \ref{prop:real-input-40-collision-rhk-window} 的唯一正临界根；并记
\[
R_\theta:=\operatorname{rad}_0 P_2
\]
为 \(P_2\) 在 \(\theta=0\) 处 Taylor 级数的收敛半径。则有严格不等式
\[
\theta_R<R_\theta.
\]
更具体地，
\[
\theta_R\approx 1.27888,
\qquad
R_\theta\approx 2.76230289346087.
\]
\end{theorem}
\begin{proof}
推论 \ref{cor:real-input-40-collision-branch-radius} 给出
\[
R_\theta\approx 2.76230289346087.
\]
命题 \ref{prop:real-input-40-collision-rhk-window} 则给出
\[
u_R\approx 3.59262181344,
\qquad
\theta_R=\log u_R\approx 1.27888.
\]
直接比较即得
\[
\theta_R<R_\theta.
\]
证毕。
\end{proof}

\begin{corollary}[解析超临界窗的存在]\label{cor:xi-time-part9zm-analytic-supercritical-window}
在定理 \ref{thm:xi-time-part9zm-rh-threshold-before-analytic-obstruction} 的设定下，存在非空开区间
\[
(\theta_R,R_\theta)
\]
使得对每个 \(\theta\) 落于其中，压力函数 \(P_2(\theta)\) 仍解析，而 prime-orbit 误差指数却满足
\[
\beta(e^\theta)>\frac12.
\]
等价地，在 \(u\)-变量上存在非空开区间
\[
(u_R,e^{R_\theta})
\approx
(3.59262181344,\ 15.8362702312)
\]
使体系同时具有“局部解析可 perturb”与“平方根指数失效”这两种性质。
\end{corollary}
\begin{proof}
由定理 \ref{thm:xi-time-part9zm-rh-threshold-before-analytic-obstruction}，
\[
0<\theta_R<R_\theta,
\]
故 \((\theta_R,R_\theta)\) 为非空开区间。若 \(\theta\in(\theta_R,R_\theta)\)，则一方面
\[
|\theta|<R_\theta,
\]
故 \(P_2(\theta)\) 仍位于其 Taylor 解析圆盘内；另一方面
\[
u=e^\theta>u_R.
\]
由推论 \ref{cor:real-input-40-collision-error-phase}，
\[
u>u_R\Longrightarrow \beta(u)>\frac12.
\]
代入 \(u=e^\theta\) 即得
\[
\beta(e^\theta)>\frac12.
\]
再把端点指数化，得到对应的 \(u\)-区间
\[
(u_R,e^{R_\theta}).
\]
利用
\[
u_R\approx 3.59262181344,
\qquad
R_\theta\approx 2.76230289346087
\]
可得
\[
e^{R_\theta}\approx 15.8362702312.
\]
故所述 \(u\)-区间确为非空。证毕。
\end{proof}

\begin{theorem}[最小非平凡角色已足以打破平方根律]\label{thm:xi-time-part9zm-parity-character-breaks-square-root-barrier}
仍在碰撞位势三次方程的设定下，取最小非平凡单位根扭曲 \(u=-1\)。则对应谱半径 \(\rho(-1)\) 是三次方程
\[
x^3-2x^2-1=0
\]
的最大实根，数值上
\[
\rho(-1)\approx 2.2055694304.
\]
若记
\[
\lambda:=\rho(1)=\varphi^2,
\qquad
\beta_{\mathrm{par}}
:=
\frac{\log \rho(-1)}{\log\lambda},
\]
则
\[
\rho(-1)>\lambda^{1/2}=\varphi,
\qquad
\beta_{\mathrm{par}}\approx 0.82186854047>\frac12.
\]
因此，仅仅奇偶角色便已足以严格突破平方根误差门槛。
\end{theorem}
\begin{proof}
备注 \ref{rem:real-input-40-collision-parity} 已给出：当 \(u=-1\) 时，碰撞三次方程退化为
\[
\rho^3-2\rho^2-1=0,
\]
其最大实根即为 \(\rho(-1)\)，且
\[
\rho(-1)\approx 2.2055694304.
\]

另一方面，定理 \ref{thm:xi-real-input-40-collision-minimal-modulus-two-witness} 已证明
\[
\sqrt{\lambda}<\rho_-<1+\sqrt2<\lambda,
\qquad
\theta_2:=\frac{\log \rho_-}{\log\lambda}>\frac12,
\]
其中 \(\rho_-=\rho(-1)\)。因此
\[
\rho(-1)>\lambda^{1/2}=\varphi.
\]
再定义
\[
\beta_{\mathrm{par}}:=\frac{\log \rho(-1)}{\log\lambda},
\]
便有
\[
\beta_{\mathrm{par}}=\theta_2>\frac12.
\]
结合现有数值
\[
\rho(-1)\approx 2.2055694304,
\qquad
\lambda=\varphi^2\approx 2.6180339887,
\]
即可得到
\[
\beta_{\mathrm{par}}\approx 0.82186854047.
\]
证毕。
\end{proof}

\begin{corollary}[局部累积量体系无法以收敛半径失效来预告 RH 相变]\label{cor:xi-time-part9zm-local-cumulant-radius-cannot-detect-rh-threshold}
设
\[
P_2(\theta)=\sum_{n\ge 0} a_n\theta^n
\]
为碰撞位势压力在 \(\theta=0\) 处的 Taylor 展开。则不存在任何仅以“收敛半径耗尽”为判据的局部解析机制，能够检测 RH 阈值 \(\theta_R\) 的到达；因为存在整个区间
\[
\theta\in(\theta_R,R_\theta)
\]
使得该级数仍收敛，而 prime-orbit 误差指数已满足
\[
\beta(e^\theta)>\frac12.
\]
\end{corollary}
\begin{proof}
由定理 \ref{thm:xi-time-part9zm-rh-threshold-before-analytic-obstruction} 与推论 \ref{cor:xi-time-part9zm-analytic-supercritical-window}，
\[
(\theta_R,R_\theta)\neq\varnothing,
\]
并且对每个 \(\theta\) 落在该区间中，都同时满足
\[
|\theta|<R_\theta
\qquad\text{与}\qquad
\beta(e^\theta)>\frac12.
\]
前者说明 \(P_2(\theta)\) 在 \(\theta=0\) 的 Taylor 级数仍处于收敛圆盘内，后者说明 RH 型平方根律已经失效。故任何把 RH 相变等同于“局部幂级数开始失效”的判据都与该区间矛盾。检测相变所必需的不是收敛性本身，而是非主谱与 Perron 谱之间的序关系。证毕。
\end{proof}

上述八条后果表明，Zeckendorf 折叠的 sofic 歧义壳与碰撞位势的 RH 型阈值虽然分属时间层同步与谱层解析两个方向，其末端结构却都可压缩为严格的精确量：前者由幂零指数 \(m\)、有限记忆深度 \(m-1\) 与最小统一滞后 \(m-1\) 完全锁定，后者则表现为一场发生在解析圆盘内部、并已由最小奇偶角色触发的纯谱序重排。

\paragraph{window-\texorpdfstring{$6$}{6} 推送代数的全矩阵化、可见模单纯性与宏观标量化原则}
在 window-\(6\) 的可见 \(21\) 态空间、刚性切分
\[
X_6=X_6^{\mathrm{cyc}}\sqcup X_6^{\mathrm{bdry}},
\qquad
|X_6^{\mathrm{cyc}}|=18,
\quad
|X_6^{\mathrm{bdry}}|=3,
\]
以及 bin-fold 纤维直方图
\[
d_6^{\mathrm{bin}}(w)\in\{2,3,4\},
\qquad
\#\{d=2\}=8,\ \#\{d=3\}=4,\ \#\{d=4\}=9
\]
已经固定之后，若再把另一条审计链给出的交换子包络
\[
\mathfrak{k}_6=\mathfrak{so}(V_{\RR}),
\qquad
\mathfrak{g}_6=\mathfrak{sl}(V_{\RR}),
\qquad
\dim V_{\RR}=|X_6|=21
\]
真正并入同一线性动力学口径，则还可继续抽取出下列六条后果。它们依次表明：六个局部推送方向在可见 \(21\) 维空间上已经生成全部线性算子；对应可见模因此完全不可压缩；\(18\oplus 3\) 的 cyclic/boundary 切分并非动力学扇区分解；任一非零推送方向在 \(\mathfrak{so}_{21}\) 的伴随作用下即可饱和全部 \(230\) 维对称无迹扇区；任何与六个推送方向完全等变的二次型都只能是标量；而一旦这一 microscopic full matrixization 可向连续极限传递，则 macroscopic covariance 也应被迫标量化。

\begin{theorem}[window-\texorpdfstring{$6$}{6} 推送代数的全矩阵化]\label{thm:xi-time-part9zn-window6-pushforward-algebra-full-matrixization}
\leanverified{paper\_xi\_time\_part9zn\_window6\_pushforward\_algebra\_full\_matrixization}
设
\[
V_{\RR}:=\RR\langle X_6\rangle,
\qquad
V_{\CC}:=V_{\RR}\otimes_{\RR}\CC\cong \CC^{X_6},
\]
并令 \(L_1,\dots,L_6\in \End_{\CC}(V_{\CC})\) 为定义 \ref{def:window6-coordinate-pushforward-lie-envelope} 中推送算子的复线性延拓。记
\[
\mathcal A_6:=\CC\langle L_1,\dots,L_6\rangle\subset \End_{\CC}(V_{\CC})
\]
为其生成的复结合代数。则
\[
\mathcal A_6=\End_{\CC}(V_{\CC})\cong M_{21}(\CC).
\]
等价地，window-\(6\) 的六个局部推送方向在可见 \(21\) 维空间上已经生成了全部线性可观测量。
\end{theorem}
\begin{proof}
由定理 \ref{thm:xi-time-part53da-root-weight-common-refinement-five-atoms}，
\[
|X_6|=21,
\]
故 \(\dim_{\CC}V_{\CC}=21\)。

另一方面，定理 \ref{thm:window6-lie-envelope-sl21} 已给出
\[
\mathfrak g_6=\Lie_{\RR}\langle L_1,\dots,L_6\rangle=\mathfrak{sl}(V_{\RR}).
\]
由于任意结合代数对交换子封闭，\(\mathcal A_6\) 作为复结合代数必包含上述实 Lie 包络。再对其作复线性张成，得到
\[
\mathfrak{sl}(V_{\CC})\subset \mathcal A_6.
\]

取 \(V_{\CC}\) 的标准基 \((e_w)_{w\in X_6}\)，并记矩阵单位为
\[
E_{uv}\qquad (u,v\in X_6).
\]
对任意 \(u\neq v\)，矩阵单位 \(E_{uv}\) 迹为零，故
\[
E_{uv}\in \mathfrak{sl}(V_{\CC})\subset \mathcal A_6.
\]
固定 \(u\in X_6\)，任选 \(v\neq u\)。则
\[
E_{uu}=E_{uv}E_{vu}\in \mathcal A_6.
\]
因此全部对角矩阵单位也都属于 \(\mathcal A_6\)。于是
\[
I=\sum_{u\in X_6}E_{uu}\in \mathcal A_6,
\]
并且全体矩阵单位 \(\{E_{uv}\}_{u,v\in X_6}\) 都落在 \(\mathcal A_6\) 中。故
\[
\mathcal A_6=\End_{\CC}(V_{\CC})\cong M_{21}(\CC).
\]
证毕。
\end{proof}

\begin{corollary}[可见 \texorpdfstring{$21$}{21} 维模的单纯性与线性不可压缩性]\label{cor:xi-time-part9zn-window6-visible-module-simplicity-incompressibility}
把 \(V_{\CC}\) 视为 \(\mathcal A_6\)-模，则 \(V_{\CC}\) 为单纯模。因而：
\begin{enumerate}
  \item 任意 \(\mathcal A_6\)-不变子空间只能是 \(0\) 或 \(V_{\CC}\)；
  \item 若 \(W\) 为任意 \(\mathcal A_6\)-模，且
  \[
  T:V_{\CC}\to W
  \]
  为非零 intertwiner，则 \(T\) 必为单射；特别地，当 \(W\) 有限维时，
  \[
  \dim W\ge 21;
  \]
  \item 若 \(W\) 为任意 \(\mathcal A_6\)-模，且
  \[
  S:W\to V_{\CC}
  \]
  为非零 intertwiner，则 \(S\) 必为满射；
  \item 不存在任何到维数 \(<21\) 的非平凡 pushforward-equivariant 线性压缩。
\end{enumerate}
\end{corollary}
\begin{proof}
由定理 \ref{thm:xi-time-part9zn-window6-pushforward-algebra-full-matrixization}，
\[
\mathcal A_6=\End_{\CC}(V_{\CC}).
\]
设 \(0\neq U\subseteq V_{\CC}\) 为 \(\mathcal A_6\)-不变子空间。取任意非零向量
\[
\xi=\sum_{w\in X_6}\xi_w e_w\in U.
\]
选取某个 \(u\in X_6\) 使 \(\xi_u\neq 0\)。对任意 \(v\in X_6\)，矩阵单位 \(E_{vu}\in \mathcal A_6\) 给出
\[
E_{vu}\xi=\xi_u e_v\in U.
\]
故所有基向量 \(e_v\) 都落在 \(U\) 中，从而 \(U=V_{\CC}\)。这证明了 \(V_{\CC}\) 的单纯性与第 \(1\) 条。

现在令
\[
T:V_{\CC}\to W
\]
为非零 intertwiner，则 \(\ker T\) 为 \(\mathcal A_6\)-不变子空间。由第 \(1\) 条，
\[
\ker T=0
\quad\text{或}\quad
\ker T=V_{\CC}.
\]
若 \(T\neq 0\)，后者不可能，故 \(\ker T=0\)，从而 \(T\) 单射。这证明了第 \(2\) 条。

同理，对任意非零 intertwiner
\[
S:W\to V_{\CC},
\]
其像 \(\operatorname{im}S\) 是 \(V_{\CC}\) 的非零 \(\mathcal A_6\)-不变子空间，由第 \(1\) 条只能等于 \(V_{\CC}\)。故 \(S\) 满射，得到第 \(3\) 条。

第 \(4\) 条是第 \(2\) 条的直接推论：若存在到某个有限维 \(W\) 且 \(\dim W<21\) 的非零等变线性压缩，则该映射既非零又不可能单射，矛盾。证毕。
\end{proof}

\begin{theorem}[cyclic/boundary 切分的动力学非不变性与精确可达性]\label{thm:xi-time-part9zn-window6-cyclic-boundary-noninvariance-reachability}
\leanverified{paper\_xi\_time\_part9zn\_window6\_cyclic\_boundary\_noninvariance\_reachability}
沿用定理 \ref{thm:xi-time-part53da-root-weight-common-refinement-five-atoms} 的记号，记
\[
V_{\CC}^{\mathrm{cyc}}:=\CC^{X_6^{\mathrm{cyc}}},
\qquad
V_{\CC}^{\mathrm{bdry}}:=\CC^{X_6^{\mathrm{bdry}}},
\qquad
V_{\CC}=V_{\CC}^{\mathrm{cyc}}\oplus V_{\CC}^{\mathrm{bdry}},
\]
其中
\[
\dim V_{\CC}^{\mathrm{cyc}}=18,
\qquad
\dim V_{\CC}^{\mathrm{bdry}}=3.
\]
则：
\begin{enumerate}
  \item \(V_{\CC}^{\mathrm{cyc}}\) 与 \(V_{\CC}^{\mathrm{bdry}}\) 都不是 \(\mathcal A_6\)-不变子空间；
  \item 对任意 \(u,v\in X_6\)，存在一个六变量非交换多项式 \(P_{vu}\) 使
  \[
  P_{vu}(L_1,\dots,L_6)e_u=e_v;
  \]
  \item 特别地，从任一 boundary 基矢可以精确送达任一 cyclic 基矢，反之亦然。
\end{enumerate}
\end{theorem}
\begin{proof}
由推论 \ref{cor:xi-time-part9zn-window6-visible-module-simplicity-incompressibility}，\(V_{\CC}\) 作为 \(\mathcal A_6\)-模是单纯的。
若 \(V_{\CC}^{\mathrm{cyc}}\) 为 \(\mathcal A_6\)-不变子空间，则必有
\[
V_{\CC}^{\mathrm{cyc}}=0
\quad\text{或}\quad
V_{\CC}^{\mathrm{cyc}}=V_{\CC},
\]
这与
\[
\dim V_{\CC}^{\mathrm{cyc}}=18
\]
矛盾。故 \(V_{\CC}^{\mathrm{cyc}}\) 不变性不成立。对 \(V_{\CC}^{\mathrm{bdry}}\) 亦同理，得到第 \(1\) 条。

再取任意 \(u,v\in X_6\)。由定理 \ref{thm:xi-time-part9zn-window6-pushforward-algebra-full-matrixization}，
\[
E_{vu}\in \mathcal A_6=\CC\langle L_1,\dots,L_6\rangle.
\]
因此存在一个有限非交换多项式 \(P_{vu}\) 使
\[
E_{vu}=P_{vu}(L_1,\dots,L_6).
\]
作用在基向量 \(e_u\) 上即得
\[
P_{vu}(L_1,\dots,L_6)e_u=E_{vu}e_u=e_v.
\]
这证明了第 \(2\) 条，而第 \(3\) 条只是把 \(u\) 与 \(v\) 分别取在 \(X_6^{\mathrm{bdry}}\) 与 \(X_6^{\mathrm{cyc}}\) 中的特例。证毕。
\end{proof}

\begin{theorem}[\texorpdfstring{$\mathfrak{so}_{21}$}{so21} 轨道饱和：任一非零推送方向生成全部对称无迹扇区]\label{thm:xi-time-part9zn-window6-so21-orbit-saturation-sym0}
回到实可见空间
\[
V_{\RR}:=\RR\langle X_6\rangle.
\]
记
\[
\mathcal S_0(V_{\RR})
:=
\{T\in \End_{\RR}(V_{\RR}):T^\top=T,\ \operatorname{tr}T=0\}.
\]
则有正交分解
\[
\mathfrak{sl}(V_{\RR})
=
\mathfrak{so}(V_{\RR})\oplus \mathcal S_0(V_{\RR}),
\qquad
\dim \mathcal S_0(V_{\RR})=\frac{21\cdot 22}{2}-1=230.
\]
并且对每个 \(i\in\{1,\dots,6\}\)，只要 \(L_i\neq 0\)，便有
\[
U(\mathfrak{so}(V_{\RR}))\cdot L_i=\mathcal S_0(V_{\RR}),
\]
其中作用由交换子给出：
\[
A\cdot T:=[A,T],
\qquad
A\in \mathfrak{so}(V_{\RR}),\ T\in \mathcal S_0(V_{\RR}).
\]
换言之，任取一个非零推送方向，其在 \(\mathfrak{so}_{21}\) 的伴随作用下便可生成全部 \(230\) 维对称无迹可观测量。
\end{theorem}
\begin{proof}
由引理 \ref{lem:window6-li-symmetric} 与引理 \ref{lem:window6-li-traceless}，每个 \(L_i\) 都满足
\[
L_i^\top=L_i,
\qquad
\operatorname{tr}L_i=0,
\]
故
\[
L_i\in \mathcal S_0(V_{\RR}).
\]
而备注 \ref{rem:window6-sl-decomposition-so-sym0} 已给出
\[
\mathfrak{sl}(V_{\RR})
=
\mathfrak{so}(V_{\RR})\oplus \mathcal S_0(V_{\RR}),
\qquad
\dim \mathcal S_0(V_{\RR})=230.
\]

现取任意非零
\[
T\in \mathcal S_0(V_{\RR}).
\]
由于 \(T\) 是实对称矩阵，可取一组正交基使
\[
T=\operatorname{diag}(\lambda_1,\dots,\lambda_{21}),
\qquad
\sum_{a=1}^{21}\lambda_a=0,
\]
且不全为零。因此存在 \(a\neq b\) 使
\[
\lambda_a\neq \lambda_b.
\]
记
\[
K_{ab}:=E_{ab}-E_{ba}\in\mathfrak{so}(V_{\RR}),
\qquad
S_{ab}:=E_{ab}+E_{ba},
\qquad
H_{ab}:=E_{aa}-E_{bb}.
\]
则直接计算给出
\[
[K_{ab},T]=(\lambda_b-\lambda_a)S_{ab}.
\]
故由 \(T\) 生成的 \(\mathfrak{so}(V_{\RR})\)-模包含某个非零 \(S_{ab}\)。再有
\[
[K_{ab},S_{ab}]=2H_{ab},
\]
故该模也包含某个非零对角差。

现在对任意 \(c\notin\{a,b\}\)，直接计算得
\[
[K_{bc},S_{ab}]=-S_{ac},
\qquad
[K_{ac},S_{ab}]=S_{bc}.
\]
因此该模同时包含全部与指标 \(a\) 或 \(b\) 相连的对称非对角元。重复同样的计算，可逐步得到任意
\[
S_{ij}\qquad (i<j).
\]
再由
\[
[K_{ij},S_{ij}]=2(E_{ii}-E_{jj})=2H_{ij}
\]
得到全部对角差 \(H_{ij}\)。

由于 \(\mathcal S_0(V_{\RR})\) 正由
\[
\{S_{ij}:1\le i<j\le 21\}
\quad\text{与}\quad
\{H_{ij}:1\le i<j\le 21\}
\]
张成，故由任意非零 \(T\in \mathcal S_0(V_{\RR})\) 生成的 \(\mathfrak{so}(V_{\RR})\)-模就是整个 \(\mathcal S_0(V_{\RR})\)。最后取 \(T=L_i\neq 0\) 即得
\[
U(\mathfrak{so}(V_{\RR}))\cdot L_i=\mathcal S_0(V_{\RR}).
\]
证毕。
\end{proof}

\begin{theorem}[pushforward-equivariant 二次型的标量刚性]\label{thm:xi-time-part9zn-window6-pushforward-equivariant-quadratic-scalar-rigidity}
\leanverified{paper\_xi\_time\_part9zn\_window6\_pushforward\_equivariant\_quadratic\_scalar\_rigidity}
设
\[
\Sigma\in \End_{\CC}(V_{\CC})
\]
满足
\[
[\Sigma,L_i]=0
\qquad (i=1,\dots,6).
\]
则必存在某个标量 \(c\in\CC\) 使
\[
\Sigma=c\,I_{V_{\CC}}.
\]
特别地，若 \(\Sigma\) 还是正半定协方差型算子，则
\[
\Sigma=c\,I_{V_{\CC}},
\qquad
c\ge 0.
\]
\end{theorem}
\begin{proof}
由假设，\(\Sigma\) 与所有生成元 \(L_1,\dots,L_6\) 对易，因此也与由它们生成的整个结合代数 \(\mathcal A_6\) 对易。故
\[
\Sigma\in \mathcal A_6'.
\]
而由定理 \ref{thm:xi-time-part9zn-window6-pushforward-algebra-full-matrixization}，
\[
\mathcal A_6=\End_{\CC}(V_{\CC})\cong M_{21}(\CC).
\]
因此
\[
\mathcal A_6'=\CC I_{V_{\CC}},
\]
从而
\[
\Sigma=c\,I_{V_{\CC}}
\]
对某个 \(c\in\CC\) 成立。若 \(\Sigma\) 进一步正半定，则其唯一特征值 \(c\) 非负，故 \(c\ge 0\)。证毕。
\end{proof}

\begin{conjecture}[full matrixization 导致的宏观标量化原则]\label{conj:xi-time-part9zn-window6-full-matrixization-macroscopic-scalarization}
设 \((\mathcal Y_\delta)\) 是一族 \(V_{\CC}\)-值的 window-\(6\) 可见场，来源于局部六顶点 observable 经 pushforward-polynomial 函子化后的提升。假定存在归一化 \(a_\delta\to\infty\)，使得
\[
\frac{\mathcal Y_\delta}{a_\delta}
\]
在分布意义下具有非平凡高斯极限。再记其双点协方差核为
\[
\Sigma_\delta(x,y)\in \End_{\CC}(V_{\CC}).
\]
若对任意测试函数对 \((\phi,\psi)\)，压缩后的算子
\[
\Sigma_\delta[\phi,\psi]
:=
\iint \phi(x)\,\Sigma_\delta(x,y)\,\psi(y)\,dx\,dy
\]
满足渐近等变性
\[
\max_{1\le i\le 6}
\bigl\|
[L_i,\Sigma_\delta[\phi,\psi]]
\bigr\|
\longrightarrow 0,
\]
则应存在一个标量核 \(K\) 使
\[
\Sigma_\delta(x,y)\longrightarrow K(x,y)\,I_{V_{\CC}}
\]
（等价地，对任意测试函数对都有
\(\Sigma_\delta[\phi,\psi]\to K[\phi,\psi]\,I_{V_{\CC}}\)）。

因此，在任何与六顶点 height/GFF 连续极限相容、且不显式破坏 pushforward 包络的 window-\(6\) 提升中：
\begin{enumerate}
  \item \(B_3/C_3\) 标签不能作为额外独立的 massless 协方差方向存活；
  \item boundary 三态不能作为独立连续张量因子存活；
  \item 全部可见长程涨落在 \(21\) 维标签空间上都必须发生宏观标量化。
\end{enumerate}
\end{conjecture}

上述六条后果表明，window-\(6\) 的 \(B_3/C_3\) 几何、\(18\oplus 3\) 字典切分与 boundary 三态虽然在离散坐标组织上都具有鲜明结构，但一旦转入由六个推送方向生成的真正动力学包络，这些结构便不再对应独立的线性扇区：在 microscopic 层，\(\mathcal A_6\) 已达到 \(M_{21}(\CC)\) 的全矩阵复杂度；若这一全矩阵化机制在连续极限中仍以协方差等变性的形式保留，则 macroscopic 层也应被迫退化为单标量核。这一链条的新增内容不再诉诸局域性、紧致性或经验近似，而把“多分量自由场为何不能出现”压缩为一个纯代数判据：microscopic full matrixization forces macroscopic scalarization。

\paragraph{prime grammar、window-\texorpdfstring{$6$}{6} quotient 与 Lee--Yang 五重分支的三重不可压缩核}
在 Zeckendorf--G\"odel 稀疏素数语言的矩阵 Euler 刚性、window-\(6\) 的非自由群商障碍以及 Lee--Yang 分支塔的五个 \(T_2(E_{\min})\)-torsor 都已确定之后，还可把三者进一步压缩为一条统一的有限层阻碍链。其核心并不在于“无穷维极限尚未充分展开”，而在于下列三类有限层刚性已先行闭合：prime grammar 不能被标量 Euler 化，window-\(6\) 的 fold quotient 不能被自由轨道化，而 Lee--Yang 局部分支又不能被单链化为单一 dyadic sheet。

\begin{theorem}[window-\texorpdfstring{$6$}{6} 奇数残差的根系闭包障碍]\label{thm:xi-time-part9zp-window6-odd-residual-root-closure-obstruction}
\leanverified{paper\_xi\_time\_part9zp\_window6\_odd\_residual\_root\_closure\_obstruction}
沿用定理~\ref{thm:xi-window6-parity-charge-trigrading-b3} 的记号。记
\[
X_6=X_6^{\mathrm{cyc}}\sqcup X_6^{\mathrm{bdry}},
\qquad
|X_6|=21,\quad |X_6^{\mathrm{cyc}}|=18,\quad |X_6^{\mathrm{bdry}}|=3,
\]
并取表~\ref{tab:fold6_b3c3_root_dictionary_b3} 与表~\ref{tab:fold6_b3c3_root_dictionary_c3} 给出的 \(B_3/C_3\) 根字典，把 \(X_6^{\mathrm{cyc}}\) 与 \(18\) 个根一一对应。则：
\begin{enumerate}
  \item 全体 \(X_6\) 不可能识别为任何有限约化根系；
  \item 更强地，不存在任何有限约化根系 \(R\) 使
  \[
  X_6^{\mathrm{cyc}}\subseteq R,
  \qquad
  |R\setminus X_6^{\mathrm{cyc}}|=3.
  \]
\end{enumerate}
\end{theorem}
\begin{proof}
有限约化根系满足中心对称性
\[
\alpha\in R\Longrightarrow -\alpha\in R,
\]
且 \(0\notin R\)，故其基数必为偶数。

若 \(X_6\) 本身是某个有限约化根系，则其基数应为偶数；但
\[
|X_6|=21
\]
为奇数，矛盾。

再设存在有限约化根系 \(R\) 使
\[
X_6^{\mathrm{cyc}}\subseteq R,
\qquad
|R\setminus X_6^{\mathrm{cyc}}|=3.
\]
则
\[
|R|=|X_6^{\mathrm{cyc}}|+3=18+3=21,
\]
仍为奇数，同样矛盾。证毕。
\end{proof}

\begin{corollary}[window-\texorpdfstring{$6$}{6} 折叠商的非轨道化]\label{cor:xi-time-part9zp-window6-quotient-intrinsically-groupoidal}
设
\[
\omega\sim\eta
\iff
\Fold_6^{\mathrm{bin}}(\omega)=\Fold_6^{\mathrm{bin}}(\eta).
\]
则该等价关系既不能由任何有限群在 \(\Omega_6\) 上的自由作用轨道划分实现，也不能来自任何阿贝尔群的线性陪集划分。因而，window-\(6\) 的最小折叠商在有限层上已本质地属于群胚型压缩，而非自由群商。
\end{corollary}
\begin{proof}
定理~\ref{thm:xi-time-part67-window6-nonfree-group-quotient-obstruction} 已经证明：\(F_6=\Fold_6^{\mathrm{bin}}\) 的纤维分解既不能由自由群作用解释，也不能压平为阿贝尔线性陪集划分。故 \(\sim\) 的最小自然宿主只能允许稳定子随点变化，其群论语言必为群胚而非主群商。证毕。
\end{proof}

\begin{corollary}[Zeckendorf--G\"odel 素数语法的非标量 Euler 刚性]\label{cor:xi-time-part9zp-zg-prime-grammar-matrixization}
\leanverified{paper\_xi\_time\_part9zp\_zg\_prime\_grammar\_matrixization}
记
\[
\mathrm{ZG}
:=
\left\{
\prod_{i\in S}p_i:\ S\subset \NN\ \text{有限且不含相邻指标}
\right\},
\qquad
D_{\mathrm{ZG}}(s):=\sum_{n\in \mathrm{ZG}}n^{-s}.
\]
则 \(D_{\mathrm{ZG}}(s)\) 不能表示为任何纯局部标量 Euler 积
\[
\prod_{i\ge 1}F_i(p_i^{-s}),
\]
其中每个 \(F_i(z)\) 在 \(z=0\) 邻域解析且满足 \(F_i(0)=1\)。因此，prime-adjacency 约束迫使 \(\mathrm{ZG}\) 的 Euler 结构本质上是有序矩阵型，而不可能压缩为独立素点的标量局部乘积。
\end{corollary}
\begin{proof}
这正是定理~\ref{thm:xi-time-part62dgc-zg-matrix-euler-no-scalar-product} 的第三条。其证明表明：由于
\[
p_i\in \mathrm{ZG},
\qquad
p_{i+1}\in \mathrm{ZG},
\qquad
{p_i p_{i+1}}\notin \mathrm{ZG},
\]
对应指示函数 \(1_{\mathrm{ZG}}\) 不是乘法函数，故不存在保持同一 Dirichlet 系数的标量 Euler 积分解。证毕。
\end{proof}

\begin{theorem}[全支撑种子的面积--时间立方下界]\label{thm:xi-time-part9zp-allones-area-time-cubic-lower-bound}
\leanverified{paper\_xi\_time\_part9zp\_allones\_area\_time\_cubic\_lower\_bound}
设
\[
\Psi_m:\Omega_m\to X_m\times \Sigma^{\le L_m}
\]
为任意确定性可逆提升，字母表大小记为
\[
|\Sigma|=M\ge 2.
\]
再设
\[
\omega^{(m)}:=11\cdots 1\in\Omega_m
\]
是全 \(1\) 输入，并令 \(T(m)\) 为定理~\ref{thm:all-ones-reduction-length} 中 \(\omega^{(m)}\) 在规范化程序 \(\Pi_{m+1}\) 下的停机步数。则
\[
T(m)=\Bigl\lfloor \frac{m^2}{4}\Bigr\rfloor,
\]
并且
\[
\liminf_{m\to\infty}\frac{L_m\,T(m)}{m^3}
\ \ge\
\frac{1}{4}\cdot \frac{\log(2/\varphi)}{\log M}.
\]
\end{theorem}
\begin{proof}
由定理~\ref{thm:all-ones-reduction-length}，
\[
T(m)=\Bigl\lfloor \frac{m^2}{4}\Bigr\rfloor,
\qquad
\lim_{m\to\infty}\frac{T(m)}{m^2}=\frac14.
\]

另一方面，\(\Psi_m\) 为单射且满足
\[
\operatorname{pr}_{X_m}\circ \Psi_m=\Fold_m,
\]
故
\[
2^m=|\Omega_m|
\le
|X_m|\cdot |\Sigma^{\le L_m}|
=
F_{m+2}\sum_{\ell=0}^{L_m}M^\ell
<
F_{m+2}\,M^{L_m+1}.
\]
取 \(\log_M\) 得
\[
L_m\ge \left\lceil \log_M\!\Bigl(\frac{2^m}{F_{m+2}}\Bigr)\right\rceil-1,
\]
从而
\[
\liminf_{m\to\infty}\frac{L_m}{m}
\ge
\frac{\log(2/\varphi)}{\log M},
\]
因为 \(F_{m+2}=\Theta(\varphi^m)\)。

于是
\[
\liminf_{m\to\infty}\frac{L_m\,T(m)}{m^3}
\ge
\left(\liminf_{m\to\infty}\frac{L_m}{m}\right)
\left(\lim_{m\to\infty}\frac{T(m)}{m^2}\right)
\ge
\frac{\log(2/\varphi)}{\log M}\cdot \frac14.
\]
证毕。
\end{proof}

\begin{corollary}[Lee--Yang 分支塔的五重二维 dyadic 全移位分解]\label{cor:xi-time-part9zp-leyang-five-two-dimensional-fullshifts}
沿用定理~\ref{thm:xi-time-part62d-leyang-branch-inverse-limit-five-torsors} 的记号。记
\[
\Pi_n^{\mathrm{fin}}
=
\bigsqcup_{i=0}^{4}\Pi_{i,n},
\qquad
\Pi_{i,n}=P_i+E_{\min}[2^n].
\]
则：
\begin{enumerate}
  \item 对每个 \(i\) 与 \(n\)，\(\Pi_{i,n}\) 都与 \((\ZZ/(2^n\ZZ))^{2}\) 同构，因而 \(|\Pi_{i,n}|=4^n=2^{2n}\)；
  \item 过渡映射 \([2]:\Pi_{i,n+1}\to \Pi_{i,n}\) 在该识别下正是二坐标 dyadic 地址的“去最高位截断”；
  \item 分支逆极限 \(\mathcal R_\infty=\bigsqcup_{i=0}^{4}\mathcal R_\infty^{(i)}\) 非规范地分裂为五份 \(\ZZ_2^2\)；
  \item 若定义层熵
  \[
  h_{\mathrm{LY}}
  :=
  \lim_{n\to\infty}\frac1n\log \deg R(y_n),
  \]
  则
  \[
  h_{\mathrm{LY}}=\log 4=2\log 2.
  \]
\end{enumerate}
因此，Lee--Yang 局部分支动力学并非单一 dyadic 链，而是五份彼此平行的二维 dyadic 全移位。
\end{corollary}
\begin{proof}
由定理~\ref{thm:xi-time-part62d-leyang-branch-inverse-limit-five-torsors}，每个 \(\Pi_{i,n}\) 都是 \(E_{\min}[2^n]\) 的平移余类。又
\[
E_{\min}[2^n]\cong (\ZZ/(2^n\ZZ))^{2},
\]
故第 \(1\) 条成立。

推论~\ref{cor:xi-terminal-zm-leyang-finite-branch-regular-4ary-address} 已给出 \([2]\) 在有限分支层上的 regular \(4\)-ary 地址解释；在二坐标表示下，这正是两条独立 dyadic 坐标同时去最高位的截断映射，故得第 \(2\) 条。

第 \(3\) 条则是定理~\ref{thm:xi-time-part62d-leyang-branch-inverse-limit-five-torsors} 的非规范 torsor 识别
\[
\mathcal R_\infty^{(i)}\cong T_2(E_{\min})\cong \ZZ_2^2
\]
逐分支求不交并的结果。

最后，由
\[
\deg R(y_n)=5\cdot 4^n
\]
（见定理~\ref{thm:xi-time-part62d-leyang-branch-inverse-limit-five-torsors} 及定理~\ref{thm:xi-time-part62dgc-leyang-fivecore-dyadic-core-collapse}），直接计算得到
\[
h_{\mathrm{LY}}
=
\lim_{n\to\infty}\frac1n\log(5\cdot 4^n)
=
\log 4.
\]
证毕。
\end{proof}

\begin{conjecture}[grammar--quotient--sheet 三重不可压缩核]\label{conj:xi-time-part9zp-grammar-quotient-sheet-incompressible-kernel}
设某条沿 PRIMEGAME / Fold / Lee--Yang / moment-kernel 路线逼近 classical \(\zeta/\Xi\) 的有限化正规化，同时施行以下三类压缩：
\begin{enumerate}
  \item 将 prime grammar 标量化为独立 Euler 因子；
  \item 将 fold quotient 轨道化为自由群商或阿贝尔线性陪集商；
  \item 将 Lee--Yang 局部分支的五重 sheet 标号压缩为单分量 dyadic 链。
\end{enumerate}
则该正规化必然失败。其不可消去的根源分别在于：
\begin{enumerate}
  \item prime adjacency 迫使 Euler 结构本质上矩阵化；
  \item window-\(6\) 的折叠等价关系本质上不是自由轨道商，也不是线性陪集商；
  \item Lee--Yang 局部分支动力学本质上是五份平行的二维 dyadic 全移位，而非单链。
\end{enumerate}
若此猜想成立，则 classical \(\zeta/\Xi\) 路线中的真正困难，不首先表现为“尚未找到足够大的 ambient space”，而首先表现为：任何同时压缩 prime grammar、fold quotient 与局部分支标签的有限化范式，都会精确丢失三类彼此独立的刚性数据。
\end{conjecture}

\noindent
由此，PRIMEGAME/FRACTRAN 侧的非标量语法、window-\(6\) 侧的非自由 quotient，以及 Lee--Yang 侧的五重分支标签便不再只是三条互不相干的离散现象，而是同一条 obstruction chain 的三个有限层面：prime grammar 不能被单点 Euler 化，fold quotient 不能被主群商化，而局部分支又不能被单链化。若要继续沿 RH 路线推进，则更自然的方向已不再是寻找一个更大的单一算子去包揽全部数据，而是系统建立一条 grammar--quotient--sheet obstruction theory。

\endinput

\paragraph{Jensen 递归阈值、primitive twist 资源律与 prime-slice 支撑障碍}
在 Jensen 缺陷的帽核二阶差分、圆盘完成化的 RH 证书、单位根扭曲的四阶慢模展开以及 prime-slice 外置与局部化 solenoid 对偶都已建立之后，还可把这些接口进一步压缩为一条共同的有限层刚性链。其核心现象分别是：中心单零点在 dyadic 递归中具有 sharp 的噪声阈值；主 primitive twist 的抑制代价必服从二次资源律；而一切处处分支的前史链一旦要求忠实 prime-slice 外置，便不可能压入固定有限支持的局部化 solenoid。

\begin{theorem}[中心单零点的 dyadic Jensen 递归阈值与自证化]\label{thm:xi-time-part9zq-jensen-central-root-threshold}
\leanverified{paper\_xi\_time\_part9zq\_jensen\_central\_root\_threshold}
在定理~\ref{thm:xi-jensen-defect-hat-kernel-second-difference} 与定理~\ref{thm:xi-jensen-recursive-resolution-noise-coupling} 的设定下，设 \(x_\ast\) 是 \(J_F\) 对应的一个中心单零点位置，并对某个尺度 \(h>0\) 给出带噪三点观测
\[
J^\varepsilon(x_\ast-h),\qquad
J^\varepsilon(x_\ast),\qquad
J^\varepsilon(x_\ast+h),
\]
满足
\[
\left|J^\varepsilon(\xi)-J_F(\xi)\right|\le \varepsilon
\qquad
\bigl(\xi\in\{x_\ast-h,x_\ast,x_\ast+h\}\bigr).
\]
记
\[
W_h^\varepsilon(x_\ast)
:=
J^\varepsilon(x_\ast+h)-2J^\varepsilon(x_\ast)+J^\varepsilon(x_\ast-h).
\]
则成立：
\begin{enumerate}
  \item 若
  \[
  \varepsilon<\frac{h}{4},
  \]
  则
  \[
  W_h^\varepsilon(x_\ast)>0;
  \]
  因而以 \(W_h^\varepsilon>0\) 为触发规则的中心化 dyadic 递归在该层必然继续；
  \item 若更强地
  \[
  \varepsilon<\frac{h}{8},
  \]
  则
  \[
  W_h^\varepsilon(x_\ast)>4\varepsilon,
  \]
  因而该层已构成一个有限三点零点证书，足以推出窗口 \((x_\ast-h,x_\ast+h)\) 内确有盘内零点；
  \item 常数 \(1/4\) 对“继续触发”这一目的而言是尖锐的。
\end{enumerate}
\end{theorem}
\begin{proof}
由定理~\ref{thm:xi-jensen-recursive-resolution-noise-coupling}，在中心单零点情形下，
\[
\varepsilon<\frac{h}{4}
\quad\Longrightarrow\quad
W_h^\varepsilon(x_\ast)>0,
\]
且该门槛对触发规则 \(W_h^\varepsilon>0\) 是尖锐的。这给出第 \(1\) 条与第 \(3\) 条。

若 \(\varepsilon<h/8\)，则
\[
W_h^\varepsilon(x_\ast)\ge h-4\varepsilon>4\varepsilon.
\]
于是可应用推论~\ref{cor:xi-jensen-defect-noisy-threepoint-certificate}，得到
\[
\mu_F((x_\ast-h,x_\ast+h))\ge 1,
\]
即对应半径带内必有盘内零点。证毕。
\end{proof}

\begin{corollary}[无穷深度 dyadic 细化的临界噪声律]\label{cor:xi-time-part9zq-jensen-dyadic-critical-noise}
沿用定理~\ref{thm:xi-time-part9zq-jensen-central-root-threshold} 的设定。令
\[
h_n:=2^{-n}h_0,
\qquad
\varepsilon_n\ge 0
\qquad(n\ge 0).
\]
则：
\begin{enumerate}
  \item 若
  \[
  \sup_{n\ge 0}\frac{4\varepsilon_n}{h_n}<1,
  \]
  则中心化递归在全部尺度上永不熄灭；
  \item 若更强地
  \[
  \sup_{n\ge 0}\frac{8\varepsilon_n}{h_n}<1,
  \]
  则全部尺度同时自证化：每一层都给出一个有限三点零点证书；
  \item 反之，只要存在某层 \(n_0\) 使
  \[
  4\varepsilon_{n_0}\ge h_{n_0},
  \]
  则不能对一切合法噪声实现统一保证该层继续触发。
\end{enumerate}
\end{corollary}
\begin{proof}
前两条逐层应用定理~\ref{thm:xi-time-part9zq-jensen-central-root-threshold} 即得。第三条则是同一定理中 sharpness 结论在第 \(n_0\) 层的直接重述。证毕。
\end{proof}

\begin{theorem}[单个三点 Jensen 见证给出的有限 RH 否证器]\label{thm:xi-time-part9zq-threepoint-jensen-rh-disproof}
\leanverified{paper\_xi\_time\_part9zq\_threepoint\_jensen\_rh\_disproof}
令
\[
F_\zeta(w):=\Xi_\zeta(s(w))
\]
为圆盘完成化读出函数，并记其 Jensen 缺陷为 \(J_\zeta\)。若对某个中心 \(x\)、半宽 \(h\) 与噪声水平 \(\varepsilon\) 已知
\[
\left|J_\zeta^\varepsilon(\xi)-J_\zeta(\xi)\right|\le \varepsilon
\qquad
\bigl(\xi\in\{x-h,x,x+h\}\bigr),
\]
且
\[
W_h^\varepsilon(x)
:=
J_\zeta^\varepsilon(x+h)-2J_\zeta^\varepsilon(x)+J_\zeta^\varepsilon(x-h)
>
4\varepsilon,
\]
则对应半径窗口内必存在 \(F_\zeta\) 的盘内零点。等价地，视界对数导数
\[
\mathscr C_\zeta(w):=-\frac{\Xi_\zeta'(s(w))}{\Xi_\zeta(s(w))}
\]
在圆盘内部存在极点。因而 RH 失效。
\end{theorem}
\begin{proof}
由推论~\ref{cor:xi-jensen-defect-noisy-threepoint-certificate}，条件
\[
W_h^\varepsilon(x)>4\varepsilon
\]
直接推出该半径窗口对应的对数半径带中至少存在一个盘内零点，即 \(F_\zeta\) 在 \(\mathbb D\) 内有零点。

再由定理~\ref{thm:xi-disk-zero-pole-index-reduction}，
\[
F_\zeta(w_0)=0
\iff
\mathscr C_\zeta \text{ 在 }w_0\text{ 处有极点},
\]
而附录定理~\ref{thm:app-rh-iff-disk-zero-free} 又给出 RH 等价于 \(F_\zeta\) 在单位圆盘内无零点。故一旦存在上述三点见证，便得到 RH 失效。证毕。
\end{proof}

\begin{corollary}[单个三点 Jensen 见证强制负平方指数为正]\label{cor:xi-time-part9zq-threepoint-jensen-negative-square-index}
在定理~\ref{thm:xi-time-part9zq-threepoint-jensen-rh-disproof} 的设定下，一旦存在某个三点噪声见证满足
\[
W_h^\varepsilon(x)>4\varepsilon,
\]
则
\[
\mathrm{sq}_-(\mathscr C_\zeta)\ge 1.
\]
\end{corollary}
\begin{proof}
由定理~\ref{thm:xi-time-part9zq-threepoint-jensen-rh-disproof}，该见证推出 RH 失效。再由定理~\ref{thm:xi-disk-zero-pole-index-reduction}，
\[
\text{RH}\iff \mathrm{sq}_-(\mathscr C_\zeta)=0.
\]
故 RH 失效即推出
\[
\mathrm{sq}_-(\mathscr C_\zeta)\neq 0.
\]
又负平方指数取值于非负整数，遂有
\[
\mathrm{sq}_-(\mathscr C_\zeta)\ge 1.
\]
证毕。
\end{proof}

\begin{theorem}[主 primitive 扭曲的二次资源律]\label{thm:xi-time-part9zq-primitive-twist-quadratic-resource}
记
\[
\lambda:=\rho(1)=\varphi^2,
\qquad
r_p:=\frac{\rho(e^{2\pi i/p})}{\lambda}
\]
为奇素数模 \(p\) 上主 primitive 扭曲通道的归一化谱半径。定义
\[
A_\varphi:=\frac{6\sqrt5-5}{250}{(2\pi)}^2,
\qquad
B_\varphi:=\frac{7+24\sqrt5}{75000}{(2\pi)}^4.
\]
则当 \(p\to\infty\) 时，
\[
\log r_p=-A_\varphi p^{-2}+B_\varphi p^{-4}+O(p^{-6}),
\]
从而
\[
1-r_p
=
A_\varphi p^{-2}
\Bigl(-B_\varphi-\frac{A_\varphi^2}{2}\Bigr)p^{-4}
O(p^{-6}).
\]
特别地，若要求
\[
1-r_p\le \varepsilon,
\]
则必有
\[
p\ge \sqrt{\frac{A_\varphi}{\varepsilon}}\,(1+o(1))
\qquad
(\varepsilon\downarrow 0).
\]
因此，主 primitive 扭曲的抑制成本在素数模度量上必然服从二次资源律。
\end{theorem}
\begin{proof}
由推论~\ref{cor:real-input-40-collision-twist-fourth}，对
\[
t_p:=\frac{2\pi}{p}
\]
有
\[
\log\frac{\rho(t_p)}{\lambda}
=
-\frac{6\sqrt5-5}{250}\,t_p^2
+\frac{7+24\sqrt5}{75000}\,t_p^4
+O(t_p^6),
\]
即
\[
\log r_p=-A_\varphi p^{-2}+B_\varphi p^{-4}+O(p^{-6}).
\]
对其指数化，得到
\[
r_p
=
\exp\!\bigl(-A_\varphi p^{-2}+B_\varphi p^{-4}+O(p^{-6})\bigr)
=
1-A_\varphi p^{-2}
+\Bigl(B_\varphi+\frac{A_\varphi^2}{2}\Bigr)p^{-4}
+O(p^{-6}),
\]
从而
\[
1-r_p
=
A_\varphi p^{-2}
+\Bigl(-B_\varphi-\frac{A_\varphi^2}{2}\Bigr)p^{-4}
+O(p^{-6}).
\]
最后由
\[
1-r_p=A_\varphi p^{-2}+O(p^{-4})
\]
反解即可得到
\[
p\ge \sqrt{\frac{A_\varphi}{\varepsilon}}\,(1+o(1)).
\]
证毕。
\end{proof}

\begin{corollary}[四阶补偿把主 primitive 审计复杂度指数从 \texorpdfstring{$1/4$}{1/4} 降到 \texorpdfstring{$1/6$}{1/6}]\label{cor:xi-time-part9zq-primitive-twist-fourth-order-compensation}
沿用定理~\ref{thm:xi-time-part9zq-primitive-twist-quadratic-resource} 的记号。定义四阶补偿预测器
\[
\widehat r_p^{(4)}:=\exp\!\bigl(-A_\varphi p^{-2}+B_\varphi p^{-4}\bigr),
\]
以及二阶高斯预测器
\[
\widehat r_p^{(2)}:=\exp(-A_\varphi p^{-2}).
\]
则有
\[
\left|r_p-\widehat r_p^{(4)}\right|=O(p^{-6}),
\qquad
\left|r_p-\widehat r_p^{(2)}\right|=O(p^{-4}).
\]
因此，为达到误差阈值 \(\eta\)，四阶补偿仅需
\[
p\gg \eta^{-1/6},
\]
而二阶高斯近似则需
\[
p\gg \eta^{-1/4}.
\]
\end{corollary}
\begin{proof}
由定理~\ref{thm:xi-time-part9zq-primitive-twist-quadratic-resource}，
\[
r_p
=
\exp\!\bigl(-A_\varphi p^{-2}+B_\varphi p^{-4}+O(p^{-6})\bigr).
\]
故
\[
r_p-\widehat r_p^{(4)}
=
\widehat r_p^{(4)}\bigl(e^{O(p^{-6})}-1\bigr)
=
O(p^{-6}),
\]
而
\[
r_p-\widehat r_p^{(2)}
=
\widehat r_p^{(2)}\bigl(e^{B_\varphi p^{-4}+O(p^{-6})}-1\bigr)
=
O(p^{-4}).
\]
对这两个余项分别反解即得
\[
p\gg \eta^{-1/6},
\qquad
p\gg \eta^{-1/4}.
\]
证毕。
\end{proof}

\begin{theorem}[有限层 prime-slice 账本的活跃片数线性下界]\label{thm:xi-time-part9zq-active-primeslice-linear-growth}
设
\[
X_0\xrightarrow{\pi_0}X_1\xrightarrow{\pi_1}\cdots\xrightarrow{\pi_{k-1}}X_k
\]
为一条有限纤维满射塔，并假定每层都存在某个 \(y_j\in X_{j+1}\) 使
\[
\#\pi_j^{-1}(y_j)\ge 2
\qquad
(0\le j\le k-1).
\]
设
\[
\Psi(x_0)=\bigl(x_k,H(x_0)\bigr),
\qquad
H(x_0)=\prod_{j=0}^{k-1}h_j(x_0),
\]
其中每个 \(h_j(x_0)\) 的全部素因子均落在预先固定的 prime-slice \(\mathcal P_j\) 中，并设 \(\Psi\) 为单射。定义活跃片集合
\[
\mathcal A_k
:=
\left\{
0\le j\le k-1:\ \exists\,x_0\in X_0,\ h_j(x_0)\neq 1
\right\}.
\]
则
\[
|\mathcal A_k|\ge k.
\]
更强地，分层 prime-slice 最优构造在最坏情形下达到等号。
\end{theorem}
\begin{proof}
由假设，每一层都是非平凡分支层，故
\[
J=
\left\{
0,\dots,k-1
\right\}
\]
在定理~\ref{thm:xi-prime-slice-nontrivial-layer-exact-minimality} 的记号下满足 \(|J|=k\)。该定理给出任何这类按层分解的单射 prime-slice 外置都必须满足
\[
N_{\mathrm{slice}}(\Psi)\ge |J|=k.
\]
而这里
\[
N_{\mathrm{slice}}(\Psi)=|\mathcal A_k|,
\]
故
\[
|\mathcal A_k|\ge k.
\]

另一方面，定理~\ref{thm:xi-time-part53ac-primeslice-truncated-history-inverse-limit} 已给出恰在每层写入一枚有效素数的 sharp 构造，故在最坏情形下等号可达。证毕。
\end{proof}

\begin{corollary}[自然扩张一旦处处分支，活跃 prime-slice 必为无穷]\label{cor:xi-time-part9zq-natural-extension-infinite-primeslice-support}
在定理~\ref{thm:xi-time-part9zq-active-primeslice-linear-growth} 的设定下，设
\[
X_0\xleftarrow{\pi_0}X_1\xleftarrow{\pi_1}X_2\xleftarrow{\pi_2}\cdots
\]
是一条在每一深度都存在非平凡分支的自然扩张链。则任一与有限截断相容的 injective prime-slice 账本，在无限深度上必然具有无限活跃片支撑。特别地，若各 prime-slice \(\mathcal P_j\) 两两不交，则所用 prime-slice 的总数必为无穷。
\end{corollary}
\begin{proof}
对任意 \(k\ge 1\)，把账本限制到前 \(k\) 层截断，便得到定理~\ref{thm:xi-time-part9zq-active-primeslice-linear-growth} 的有限深度情形，因此前 \(k\) 层的活跃片数至少为 \(k\)。令 \(k\to\infty\)，便知无限深度上活跃片支撑不可能有限。

这也可视为定理~\ref{thm:xi-time-part63b-infinite-depth-finite-prime-register-impossibility} 的直接重述：一旦每层都真实分支，就不存在只用有限多个 prime-slice 的全历史精确外置。故若各 \(\mathcal P_j\) 两两不交，则所使用的 prime-slice 总数必为无穷。证毕。
\end{proof}

\begin{theorem}[固定有限支持局部化 solenoid 不可能忠实承载无穷分支历史]\label{thm:xi-time-part9zq-no-fixed-finite-support-solenoid-host}
设 \(S\subset\mathbb P\) 为有限素数集，并记
\[
G_S:=\ZZ[S^{-1}],
\qquad
\Sigma_S:=\widehat{G_S}.
\]
考虑一条在每一深度都存在非平凡分支的自然扩张链。则不存在任何忠实的 prime-support 外置模型，能够把其全部无穷分支历史压入某个固定有限支持的局部化系统 \(G_S\) 或其 Pontryagin 对偶 \(\Sigma_S\) 中。
\end{theorem}
\begin{proof}
由定理~\ref{thm:killo-localization-circle-dimension-prime-ledger-splitting}，固定有限支持 \(S\) 的 Pontryagin 对偶满足短正合列
\[
0\to \prod_{p\in S}\ZZ_p\to \widehat{G_S}\to \TT\to 0,
\]
其全不连通核的 prime-support 恰为有限集 \(S\)，并且 \(S\) 可由该核完全恢复。故任何通过 \(G_S\) 或 \(\Sigma_S\) 实现的 prime-support 外置模型，至多只能动用有限多个 prime-slice 方向。

另一方面，推论~\ref{cor:xi-time-part9zq-natural-extension-infinite-primeslice-support} 表明：只要自然扩张在每一深度都有真实分支，则任何忠实账本都必须具有无限活跃片支撑。二者矛盾，故固定有限支持的局部化 solenoid 不可能忠实承载全部无穷分支历史。证毕。
\end{proof}

\noindent
由此，Jensen 递归阈值、primitive twist 慢模与 prime-slice/solenoid 外置便在同一条有限层刚性链上闭合：中心单零点在 dyadic 细化中存在 sharp 的触发与自证阈值；主 primitive twist 的抑制代价被 \(p^{-2}\) 资源律刚性锁定；而一切处处分支的前史链又不能被任何固定有限支持的 prime-ledger 或 solenoid 宿主忠实容纳。若继续沿 RH 路线推进，则更自然的工作方向不再是寻找单一有限宿主去同时吸收这些现象，而是分别建立 Jensen 阈值、twist 资源律与 prime-support 外置障碍之间的统一结构理论。

\endinput

\paragraph{非可缩顶层、Latt\`es 透明切片、预算断崖与 Lee--Yang 五重超立方体谱}
在最大非可缩纤维的模 \(6\) 相图、三域 Frobenius 直积的张量分解、bin-fold 配额阈值的两相常数律以及 Lee--Yang 五分量 dyadic branch 图的显式谱学都已建立之后，还可把这几条此前分散出现的接口继续压缩为下列五条闭合后果。它们分别对应：拓扑非可缩筛落只需下降到前三层、Latt\`es 全分裂切片对双判别特征与 Lee--Yang 局部分解型的统计透明、临界预算在容错阈值处发生精确黄金断崖、有限层 Lee--Yang branch 图的五重超立方体谱包，以及单位根扭曲的素数二次律与模 \(2\) 奇偶扭曲的谱孤立。

\begin{theorem}[非可缩顶层的二级筛落与模 \(6\) 六相极限律]\label{thm:xi-time-part9zr-noncontractible-sixphase-law}
\leanverified{paper\_xi\_time\_part9zr\_noncontractible\_sixphase\_law}
沿用定理 \ref{thm:xi-time-part62d-max-noncontractible-fiber-gap-four-constants} 与推论 \ref{cor:xi-time-part62d-max-noncontractible-fiber-tax-mod6} 的记号。记
\[
D_m:=\max_{x\in X_m} d_m(x),
\qquad
\widetilde D_m:=\max\{d_m(x): |K_x|\not\simeq *\}.
\]
则对一切充分大的 \(m\)，非可缩顶层永远不需要下降到第三层以下：
\[
\widetilde D_m\in \{D_m,D_m^{(2)},D_m^{(3)}\}.
\]
更精确地，存在严格的模 \(6\) 相图
\[
\widetilde D_m=
\begin{cases}
D_m, & m\equiv 0,4 \pmod 6,\\[2pt]
D_m^{(2)}, & m\equiv 1,5 \pmod 6,\\[2pt]
D_m^{(3)}, & m\equiv 2,3 \pmod 6,
\end{cases}
\]
并且六个同余类上的极限比值显式存在：
\[
\lim_{\ell\to\infty}\frac{\widetilde D_{6\ell}}{D_{6\ell}}=1,
\qquad
\lim_{\ell\to\infty}\frac{\widetilde D_{6\ell+1}}{D_{6\ell+1}}
=
\lim_{\ell\to\infty}\frac{\widetilde D_{6\ell+5}}{D_{6\ell+5}}
=
1-\frac{1}{2}\varphi^{-5},
\]
\[
\lim_{\ell\to\infty}\frac{\widetilde D_{6\ell+2}}{D_{6\ell+2}}
=
1-\varphi^{-5},
\qquad
\lim_{\ell\to\infty}\frac{\widetilde D_{6\ell+3}}{D_{6\ell+3}}
=
1-\frac{1}{2}\bigl(\varphi^{-5}+\varphi^{-9}\bigr),
\]
\[
\lim_{\ell\to\infty}\frac{\widetilde D_{6\ell+4}}{D_{6\ell+4}}=1.
\]
特别地，
\[
\lim_{m\to\infty}\frac{1}{m}\log \widetilde D_m
=
\lim_{m\to\infty}\frac{1}{m}\log D_m
=
\frac{1}{2}\log\varphi.
\]
\end{theorem}
\begin{proof}
前半部分正是定理 \ref{thm:xi-time-part62d-max-noncontractible-fiber-gap-four-constants} 与推论 \ref{cor:xi-time-part62d-max-noncontractible-fiber-tax-mod6} 的直接合并：\(\widetilde D_m\) 的最大层只能出现在 \(D_m\)、\(D_m^{(2)}\)、\(D_m^{(3)}\) 三层之一，且由模 \(6\) 剩余类完全决定。

四类极限比值亦已在推论 \ref{cor:xi-time-part62d-max-noncontractible-fiber-tax-mod6} 中逐类给出，只需按 \(m\equiv 0,1,2,3,4,5\pmod 6\) 重新写成六相口径即可。

最后，由上述六个比值极限都趋于正有限常数，可知
\[
\log \widetilde D_m=\log D_m+O(1)
\qquad (m\to\infty),
\]
故二者具有相同的指数增长率。证毕。
\end{proof}

\begin{theorem}[Latt\`es 全分裂切片上的双判别特征透明性]\label{thm:xi-time-part9zr-lattes-fullsplit-transparency}
沿用定理 \ref{thm:xi-time-part64a-l4-event-complete-conditional-independence} 与推论 \ref{cor:xi-time-part9zg-arithmetic-fingerprint-zero-mutual-information} 的记号。设
\[
E_{\mathrm{split}}
:=
\{\text{\(C(X)\bmod p\) 全分裂}\}.
\]
则对任意 \(\varepsilon_1,\varepsilon_2\in\{\pm1\}\)，有精确密度
\[
\delta\!\bigl(
E_{\mathrm{split}}\cap\{\chi_{10}=\varepsilon_1\}\cap\{\chi_{\mathrm{LY}}=\varepsilon_2\}
\bigr)
=
\frac{1}{48}\cdot \frac14
=
\frac{1}{192}.
\]
更强地，在任意固定条件
\[
E_{\mathrm{split}}\cap\{\chi_{10}=\varepsilon\}
\]
之下，Lee--Yang 三次多项式 \(g(y)=256y^3+411y^2+165y+32\) 的三种分解型密度保持不变：
\[
\delta\!\bigl(\mathcal T_{\mathrm{LY}}=(1)(1)(1)\mid E_{\mathrm{split}},\chi_{10}=\varepsilon\bigr)=\frac16,
\]
\[
\delta\!\bigl(\mathcal T_{\mathrm{LY}}=(1)(2)\mid E_{\mathrm{split}},\chi_{10}=\varepsilon\bigr)=\frac12,
\]
\[
\delta\!\bigl(\mathcal T_{\mathrm{LY}}=(3)\mid E_{\mathrm{split}},\chi_{10}=\varepsilon\bigr)=\frac13.
\]
换言之，Latt\`es 全分裂切片对十次判别特征与 Lee--Yang 局部分解型都是统计透明的。
\end{theorem}
\begin{proof}
定理 \ref{thm:xi-time-part64a-l4-event-complete-conditional-independence} 已证明：\(E_{\mathrm{split}}\) 只依赖第三个 Frobenius 因子，且对任意 \((\varepsilon_1,\varepsilon_2)\in\{\pm1\}^2\) 有
\[
\delta\!\bigl(
E_{\mathrm{split}}\cap\{\chi_{10}=\varepsilon_1,\ \chi_{\mathrm{LY}}=\varepsilon_2\}
\bigr)
=
\frac14\cdot\frac{1}{48}
=
\frac{1}{192}.
\]
这就是第一式。

再看第二部分。由三域 Frobenius 张量分解，任意只依赖 \(S_{10}\times S_3\) 两个因子的事件与 \(E_{\mathrm{split}}\) 严格独立；因此对 Lee--Yang 三次的任一分解型事件 \(\mathcal T\) 都有
\[
\delta\!\bigl(\mathcal T\mid E_{\mathrm{split}},\chi_{10}=\varepsilon\bigr)
=
\delta\!\bigl(\mathcal T\mid \chi_{10}=\varepsilon\bigr).
\]
而推论 \ref{cor:xi-time-part9zg-arithmetic-fingerprint-zero-mutual-information} 已给出在固定 \(\chi_{10}=\varepsilon\) 条件下，三种 Lee--Yang 分解型的条件密度恰为
\[
\frac16,\qquad \frac12,\qquad \frac13.
\]
代入即得结论。证毕。
\end{proof}

\begin{theorem}[重建预算的 \texorpdfstring{$\varphi$}{phi} 因子断崖]\label{thm:xi-time-part9zr-reconstruction-budget-phi-cliff}
沿用定理 \ref{thm:xi-time-part9ob-binfold-sharp-threshold-budget} 的记号。记
\[
P_m^\star(\varepsilon):=P_m^{\mathrm{crit}}(\varepsilon),
\qquad
\varepsilon_{\mathrm c}:=\frac{\varphi}{\sqrt5}.
\]
则 \(P_m^\star(\varepsilon)\) 具有严格的两相渐近：
\[
\log P_m^\star(\varepsilon)
=
m\log\!\Bigl(\frac{2}{\varphi}\Bigr)+\log 2+o(1),
\qquad
0<\varepsilon<\varepsilon_{\mathrm c},
\]
\[
\log P_m^\star(\varepsilon)
=
m\log\!\Bigl(\frac{2}{\varphi}\Bigr)+\log 2-\log\varphi+o(1),
\qquad
\varepsilon_{\mathrm c}<\varepsilon<1.
\]
因此，对任意
\[
0<\varepsilon_-<\varepsilon_{\mathrm c}<\varepsilon_+<1,
\]
都成立
\[
\frac{P_m^\star(\varepsilon_-)}{P_m^\star(\varepsilon_+)}
\longrightarrow
\varphi.
\]
亦即：一旦容错跨越 \(\varepsilon_{\mathrm c}\)，最优重建预算在指数主项不变的前提下发生精确的 \(\varphi\)-倍断崖式跌落。
\end{theorem}
\begin{proof}
这只是定理 \ref{thm:xi-time-part9ob-binfold-sharp-threshold-budget} 的直接改写。该定理已给出
\[
P_m^\star(\varepsilon)
\sim
2\left(\frac{2}{\varphi}\right)^m
\qquad
(0<\varepsilon<\varepsilon_{\mathrm c}),
\]
以及
\[
P_m^\star(\varepsilon)
\sim
2\varphi^{-1}\left(\frac{2}{\varphi}\right)^m
\qquad
(\varepsilon_{\mathrm c}<\varepsilon<1).
\]
取对数即得两条对数渐近；再相除便得到
\[
\frac{P_m^\star(\varepsilon_-)}{P_m^\star(\varepsilon_+)}
\longrightarrow
\frac{2}{2\varphi^{-1}}
=
\varphi.
\]
证毕。
\end{proof}

\begin{theorem}[Lee--Yang 分支图的五重超立方体谱定理]\label{thm:xi-time-part9zr-leyang-fivecube-spectrum}
沿用定理 \ref{thm:xi-time-part62d-leyang-branch-inverse-limit-five-torsors} 的记号。对每个 \(n\ge 1\)，令 \(\mathcal B_n\) 为如下图：顶点集为 \(R(y_n)\)；若两点位于同一平行族 \(P_i+E_{\min}[2^n]\) 内，且其二进地址只在一位上相差，则连一条边。则
\[
\mathcal B_n\cong \bigsqcup_{i=1}^{5} Q_{2n},
\]
其中 \(Q_{2n}\) 为 \(2n\)-维超立方体图。因而
\[
|V(\mathcal B_n)|=5\cdot 2^{2n}=5\cdot 4^n,
\qquad
|E(\mathcal B_n)|=5\cdot 2n\cdot 2^{2n-1}=5n\,4^n.
\]
其邻接谱满足
\[
\operatorname{Spec}(A_{\mathcal B_n})
=
\left\{
2n-2j\ \text{with multiplicity}\ 5\binom{2n}{j}:\ 0\le j\le 2n
\right\},
\]
拉普拉斯谱满足
\[
\operatorname{Spec}(L_{\mathcal B_n})
=
\left\{
2j\ \text{with multiplicity}\ 5\binom{2n}{j}:\ 0\le j\le 2n
\right\},
\]
并且热迹具有显式闭式
\[
\Tr\!\bigl(e^{-tL_{\mathcal B_n}}\bigr)
=
5(1+e^{-2t})^{2n}.
\]
\end{theorem}
\begin{proof}
这正是定理 \ref{thm:xi-time-part62d-leyang-branch-graph-gaussian-spectrum} 的分支图 \(\,G_n\,\) 在当前记号下的重述：该定理已证明每个分量都与 \(Q_{2n}\) 图同构，并且五个分量彼此不交，故
\[
\mathcal B_n\cong \bigsqcup_{i=1}^{5}Q_{2n}.
\]
顶点数与边数由单个超立方体 \(Q_{2n}\) 的标准公式
\[
|V(Q_{2n})|=2^{2n},
\qquad
|E(Q_{2n})|=2n\,2^{2n-1}
\]
立刻得到。

同一定理还给出了 \(A_{\mathcal B_n}\) 与 \(L_{\mathcal B_n}\) 的全部特征值和重数：对单个 \(Q_{2n}\)，邻接特征值为 \(2n-2j\)，拉普拉斯特征值为 \(2j\)，重数均为 \(\binom{2n}{j}\)；五个连通分量互不相交，故整体重数乘以 \(5\)。热迹公式则是
\[
\sum_{j=0}^{2n}5\binom{2n}{j}e^{-2jt}
=
5(1+e^{-2t})^{2n}.
\]
证毕。
\end{proof}

\begin{theorem}[单位根扭曲的素数二次律与模 \texorpdfstring{$2$}{2} 扭曲的谱孤立]\label{thm:xi-time-part9zr-prime-quadratic-law-parity-isolation}
沿用推论 \ref{cor:real-input-40-collision-twist-fourth} 与定理 \ref{thm:xi-time-part9zm-parity-character-breaks-square-root-barrier} 的记号。记
\[
\lambda:=\rho(1)=\varphi^2,
\qquad
\rho(t):=\rho(e^{it}),
\qquad
t_p:=\frac{2\pi}{p}
\qquad
(p\ \text{为奇素数}),
\]
并设
\[
A_\varphi:=\frac{6\sqrt5-5}{250}(2\pi)^2,
\qquad
B_\varphi:=\frac{7+24\sqrt5}{75000}(2\pi)^4.
\]
则当 \(p\to\infty\) 时，
\[
\log\frac{\rho(t_p)}{\lambda}
=
-A_\varphi p^{-2}
+B_\varphi p^{-4}
+O(p^{-6}),
\]
从而
\[
\lim_{p\to\infty}
p^2\left(1-\frac{\rho(t_p)}{\lambda}\right)
=
A_\varphi
=
\frac{2\pi^2}{125}(6\sqrt5-5).
\]
另一方面，模 \(2\) 奇偶扭曲给出固定代数常数 \(\rho(-1)\)，它是三次方程
\[
x^3-2x^2-1=0
\]
的最大实根，并满足
\[
\frac{\rho(-1)}{\lambda}\approx 0.8425.
\]
因此存在奇素数阈值 \(p_0\)，使得对一切 \(p\ge p_0\)，都有
\[
\rho(-1)<\rho(t_p)<\lambda.
\]
换言之，所有充分大的奇素数单位根扭曲都只是以 \(p^{-2}\) 速度逼近未扭曲 Perron 根，而模 \(2\) 奇偶扭曲却保持一个与 \(p\) 无关的固定谱降落；故模 \(2\) 扭曲在素数模族中形成真正的谱孤立点。
\end{theorem}
\begin{proof}
由推论 \ref{cor:real-input-40-collision-twist-fourth}，当 \(t\to 0\) 时有
\[
\log\frac{\rho(t)}{\lambda}
=
-\frac{6\sqrt5-5}{250}\,t^2
+\frac{7+24\sqrt5}{75000}\,t^4
+O(t^6).
\]
取 \(t=t_p=2\pi/p\) 即得
\[
\log\frac{\rho(t_p)}{\lambda}
=
-A_\varphi p^{-2}
+B_\varphi p^{-4}
+O(p^{-6}).
\]
再对上式指数化，得到
\[
\frac{\rho(t_p)}{\lambda}
=
1-A_\varphi p^{-2}+O(p^{-4}),
\]
于是
\[
p^2\left(1-\frac{\rho(t_p)}{\lambda}\right)\longrightarrow A_\varphi
=
\frac{2\pi^2}{125}(6\sqrt5-5).
\]

另一方面，定理 \ref{thm:xi-time-part9zm-parity-character-breaks-square-root-barrier} 已证明 \(\rho(-1)\) 是三次方程
\[
x^3-2x^2-1=0
\]
的最大实根，且
\[
\rho(-1)\approx 2.2055694304,
\qquad
\lambda=\varphi^2\approx 2.6180339887,
\]
故
\[
\frac{\rho(-1)}{\lambda}\approx 0.8425<1.
\]
由于
\[
\frac{\rho(t_p)}{\lambda}\to 1
\qquad
(p\to\infty),
\]
存在 \(p_0\) 使得对所有奇素数 \(p\ge p_0\)，都有
\[
\frac{\rho(t_p)}{\lambda}>\frac{\rho(-1)}{\lambda}.
\]
另一方面，由
\[
\log\frac{\rho(t_p)}{\lambda}
=
-A_\varphi p^{-2}+B_\varphi p^{-4}+O(p^{-6})
\]
可知对充分大的 \(p\) 该量严格为负，故同时有 \(\rho(t_p)<\lambda\)。合并两边即得
\[
\rho(-1)<\rho(t_p)<\lambda
\qquad
(p\ge p_0).
\]
证毕。
\end{proof}

\noindent
上述五条把三条此前相互分散的主线重新压到同一张相图上：一方面，拓扑非可缩筛落只在前三层之间作模 \(6\) 受控跃迁，而其指数增长率与最大纤维完全同阶；另一方面，Latt\`es 全分裂切片对 \(S_{10}\)-判别特征与 \(S_3\)-局部分解型保持严格透明，bin-fold 预算则在容错阈值 \(\varepsilon_{\mathrm c}=\varphi/\sqrt5\) 处发生精确的 \(\varphi\)-倍断崖；最后，Lee--Yang 分支几何在有限层精确展开为五份 \(Q_{2n}\) 的谱闭式，而单位根扭曲又把奇素数模族的 \(p^{-2}\) 扩散律与模 \(2\) 奇偶扭曲的固定谱降落严格分离。于是，纤维拓扑、Chebotarev 透明性、资源相变、branch 图谱学与扭曲 Perron 模态便在同一 xi 终端链上获得了完全显式的联合刻画。

\endinput

\paragraph{整数温度 pressure 骨架、零点静默、四次 Perron 单位与 \(2\)-进分辨率边界}
在整数阶 moment-kernel 的有限维实现、连续射影压力在整点上的锚定、六倍窗零点密度的精确指数率、Fibonacci 共振整函数的零点自相似、零温冻结端面的四次 Perron 根、\(S_4\)--Hurwitz 链的 Prym 平方化以及分辨率塔的 \(2\)-进逆极限都已固定之后，还可把这些彼此分散的接口继续压缩为下列七条后果。

\begin{theorem}[整数温度样本构成连续压力曲线的代数骨架]\label{thm:xi-time-part9zo-integer-temperature-pressure-algebraic-skeleton}
沿用命题 \ref{prop:pom-moment-kernel-exists} 与定理 \ref{thm:xi-projective-moment-tower-anchors-normalized-pressure} 的记号。
对每个整数 \(q\ge 2\)，记
\[
\mathcal G_q(z):=\sum_{n\ge 0}S_q(m_0+n)z^n.
\]
则有严格恒等式
\[
\mathcal G_q(z)=u_q^{\mathsf T}(I-zA_q)^{-1}v_q\in \QQ(z).
\]
并且对 Fold 主核有整点锚定公式
\[
\Lambda(q-1)=\log r_q-\log 2.
\]
特别地，
\[
2e^{\Lambda(q-1)}=r_q
\]
是一个 Perron 代数整数。换言之，连续压力曲线在全部整数温度点 \(t=q-1\) 上并非任意实值，而是落在一条对数--Perron 代数整数骨架之上。
\end{theorem}
\begin{proof}
由命题 \ref{prop:pom-moment-kernel-exists}，对每个固定 \(q\ge 2\) 都存在有限非负整数矩阵 \(A_q\)、整数向量 \(u_q,v_q\) 与起点 \(m_0\)，使
\[
S_q(m)=u_q^{\mathsf T}A_q^{\,m-m_0}v_q
\qquad (m\ge m_0).
\]
故
\[
\mathcal G_q(z)
=
\sum_{n\ge 0}u_q^{\mathsf T}A_q^n v_q\,z^n
=
u_q^{\mathsf T}\Bigl(\sum_{n\ge 0}(zA_q)^n\Bigr)v_q
=
u_q^{\mathsf T}(I-zA_q)^{-1}v_q.
\]
由于 \(A_q\) 为整数矩阵，\((I-zA_q)^{-1}\) 的各条目都是有理函数，故 \(\mathcal G_q(z)\in \QQ(z)\)。

另一方面，定理 \ref{thm:xi-projective-moment-tower-anchors-normalized-pressure} 给出
\[
\Lambda(q-1)=\log r_q-\log |A_{\mathrm{out}}|.
\]
对 Fold 主核有 \(|A_{\mathrm{out}}|=2\)，因此
\[
\Lambda(q-1)=\log r_q-\log 2.
\]
再由 Perron--Frobenius 定理，\(r_q=\rho(A_q)\) 是整数矩阵 \(A_q\) 的特征值，故为 Perron 代数整数。于是
\[
2e^{\Lambda(q-1)}=r_q
\]
具有所述算术性质。证毕。
\end{proof}

\begin{theorem}[星形 \(q\)-碰撞谱的代数次数至多二次增长]\label{thm:xi-time-part9zo-star-perron-degree-quadratic}
沿用命题 \ref{prop:pom-moment-minrecurrence-hankel}、
定理 \ref{thm:pom-star-moment-kernel-perron-symmetric-compression}
与定理 \ref{thm:pom-sympower-hankel-polynomial-upper-bound} 的记号。
在文稿的星形耦合情形下，记 \(r_q\) 为 \(q\)-重碰撞的 Perron 根，\(d_q\) 为序列 \(\{S_q(m)\}\) 的最小 Hankel 维数。则有显式上界
\[
\deg_{\QQ}(r_q)\le d_q\le \binom{q+1}{2}.
\]
因此，虽然朴素 \(q\)-份同步直积具有指数态空间，但可见谱复杂度与 Perron 数域复杂度都同时塌缩到二次量级。
\end{theorem}
\begin{proof}
由定理 \ref{thm:pom-star-moment-kernel-perron-symmetric-compression}，在星形耦合的三态对称压缩下，\(r_q\) 可由维数
\[
\binom{q+1}{2}
\]
的压缩矩阵 \(\widetilde A_q^{\mathrm{eq}}\) 读出。于是 \(r_q\) 满足其特征多项式，从而
\[
\deg_{\QQ}(r_q)\le \binom{q+1}{2}.
\]

另一方面，命题 \ref{prop:pom-moment-minrecurrence-hankel} 说明 \(d_q\) 等于序列 \(m\mapsto S_q(m)\) 的最小递推阶数，也等于最小 realization 维数；定理 \ref{thm:pom-sympower-hankel-polynomial-upper-bound} 在一般 \(N\)-状态情形下给出
\[
d_q\le \binom{q+N-1}{N-1},
\]
而此处 \(N=3\)，故
\[
d_q\le \binom{q+1}{2}.
\]
再取命题 \ref{prop:pom-moment-minrecurrence-hankel} 给出的最小 realization，其维数为 \(d_q\)。若 \(r_q\) 不是该 realization 矩阵的特征值，则由该 realization 产生的序列只能按严格小于 \(r_q\) 的谱半径增长，这与 \(r_q\) 作为同一碰撞序列的 Perron 主尺度矛盾。故 \(r_q\) 是某个 \(d_q\times d_q\) 整矩阵的特征值，从而
\[
\deg_{\QQ}(r_q)\le d_q.
\]
合并即得
\[
\deg_{\QQ}(r_q)\le d_q\le \binom{q+1}{2}.
\]
证毕。
\end{proof}
\leanverified{paper\_xi\_time\_part9zo\_star\_perron\_degree\_quadratic}

\begin{theorem}[算术零点块具有绝对宏观性而归一化静默]\label{thm:xi-time-part9zo-zero-block-macroscopic-normalized-silent}
沿用推论 \ref{cor:fold-zero-half-index-multiple6}、
定理 \ref{thm:fold-zero-density-sparse}
与定理 \ref{thm:xi-time-part70c-zero-density-exact-ratio-exponent-multiple6} 的记号，并记
\[
\mathcal Z_m:=\{k\in \ZZ/(F_{m+2}\ZZ):\ \widehat d_m(k)=0\}.
\]
记
\[
M:=F_{m+2},
\qquad
\eta_m:=\frac{|\mathcal Z_m|}{F_{m+2}}.
\]
若
\[
m+2\equiv 0\pmod 6,
\]
则有
\[
\eta_m\ge \frac{F_{(m+2)/2}}{F_{m+2}},
\qquad
\eta_m=\tau(m+2)\,O(\varphi^{-m/2}),
\]
从而
\[
\log(M-|\mathcal Z_m|)
=
\log M+O\!\bigl(\tau(m+2)\varphi^{-m/2}\bigr),
\]
以及
\[
\sqrt{M-1-|\mathcal Z_m|}
=
\sqrt M\Bigl(1+O\!\bigl(\tau(m+2)\varphi^{-m/2}\bigr)\Bigr).
\]
因此，六倍窗上虽出现规模达 \(\asymp \varphi^{m/2}\) 的零点大陪集块，但凡只经由 \(|\mathcal Z_m|\) 或其归一化比例进入的主阶估计，都只能获得半指数级副修正。
\end{theorem}
\begin{proof}
下界由推论 \ref{cor:fold-zero-half-index-multiple6} 直接给出：
\[
|\mathcal Z_m|\ge F_{(m+2)/2},
\]
故
\[
\eta_m\ge \frac{F_{(m+2)/2}}{F_{m+2}}.
\]
另一方面，定理 \ref{thm:fold-zero-density-sparse} 给出
\[
\eta_m=\frac{|\mathcal Z_m|}{F_{m+2}}
\le
\tau(m+2)\frac{F_{\lfloor (m+2)/2\rfloor}}{F_{m+2}}
=
\tau(m+2)\,O(\varphi^{-m/2}),
\]
而定理 \ref{thm:xi-time-part70c-zero-density-exact-ratio-exponent-multiple6} 进一步表明该指数率恰为 \(-\frac12\log\varphi\)。

于是
\[
M-|\mathcal Z_m|=M(1-\eta_m),
\]
从而
\[
\log(M-|\mathcal Z_m|)
=
\log M+\log(1-\eta_m)
=
\log M+O(\eta_m).
\]
再注意
\[
M-1-|\mathcal Z_m|
=
M\Bigl(1-\eta_m-\frac1M\Bigr),
\]
且 \(M^{-1}=O(\varphi^{-m})=O(\eta_m)\)，于是
\[
\sqrt{M-1-|\mathcal Z_m|}
=
\sqrt M\sqrt{1-\eta_m-\frac1M}
=
\sqrt M\bigl(1+O(\eta_m)\bigr).
\]
把 \(\eta_m=\tau(m+2)\,O(\varphi^{-m/2})\) 代回即得两式。最后一句只是指出：凡通过 \(M-|\mathcal Z_m|\)、\(\sqrt{M-1-|\mathcal Z_m|}\) 或等价归一化比例进入的零点驱动上界，其主阶都无法被 \(|\mathcal Z_m|\) 改写。证毕。
\end{proof}

\begin{theorem}[Fibonacci 共振零点满足精确重整化与 Weyl 型计数律]\label{thm:xi-time-part9zo-fibonacci-cosine-zeros-weyl-law}
沿用定理 \ref{thm:xi-fold-resonance-zero-count-selfsimilar-linear} 的记号，记
\[
\mathcal F_\varphi(z)=\prod_{n=2}^{\infty}\cos(\pi z\varphi^{-n}),
\qquad
N_\varphi^+(R):=\#\bigl\{x\in \mathcal Z(\mathcal F_\varphi):\ 0<x\le R\bigr\}.
\]
则对每个 \(R>0\) 都有严格递推
\[
N_\varphi^+(R)=N_\varphi^+\!\left(\frac{R}{\varphi}\right)+\left\lfloor \frac{R}{\varphi^2}+\frac12\right\rfloor.
\]
并且
\[
N_\varphi^+(R)=R+O(\log R),
\qquad
N_\varphi(R):=\#\bigl(\mathcal Z(\mathcal F_\varphi)\cap[-R,R]\bigr)=2R+O(\log R).
\]
换言之，Fibonacci 共振零点并非给出稀薄谱，而是具有严格线性密度；其主斜率恰由
\[
\sum_{n\ge 2}\varphi^{-n}=1
\]
唯一锁定。
\end{theorem}
\begin{proof}
定理 \ref{thm:xi-fold-resonance-zero-count-selfsimilar-linear} 已经给出精确自相似递推
\[
N_\varphi^+(R)=N_\varphi^+\!\left(\frac{R}{\varphi}\right)+\left\lfloor \frac{R}{\varphi^2}+\frac12\right\rfloor,
\]
以及渐近式
\[
N_\varphi^+(R)=R+O(\log R),
\qquad
N_\varphi(R)=2R+O(\log R).
\]
另一方面，由同一定理中的精确求和公式
\[
N_\varphi^+(R)=
\sum_{n=2}^{\lfloor \log_\varphi(2R)\rfloor}
\left\lfloor R\varphi^{-n}+\frac12\right\rfloor,
\]
把整部份替换为 \(R\varphi^{-n}+O(1)\) 并注意有效层数只有 \(O(\log R)\)，可得
\[
N_\varphi^+(R)
=
R\sum_{n\ge 2}\varphi^{-n}+O(\log R).
\]
而
\[
\sum_{n\ge 2}\varphi^{-n}=1,
\]
故线性主项恰为 \(R\)。对称性再给出
\[
N_\varphi(R)=2N_\varphi^+(R)=2R+O(\log R).
\]
证毕。
\end{proof}
\leanverified{paper\_xi\_time\_part9zo\_fibonacci\_cosine\_zeros\_weyl\_law}

\begin{theorem}[零温饱和端点常数是四次 Perron 单位的平方根]\label{thm:xi-time-part9zo-zero-temp-quartic-perron-unit}
沿用定理 \ref{thm:xi-terminal-real-input-positive-entropy-freezing-face} 的记号，令
\[
\lambda_\infty:=\rho(B_\infty)=c^2.
\]
则有
\[
\lambda_\infty^4-6\lambda_\infty^3+9\lambda_\infty^2-\lambda_\infty-1=0,
\]
从而
\[
c^8-6c^6+9c^4-c^2-1=0.
\]
特别地，\(\lambda_\infty\) 是一个 Perron 代数单位；因此零温饱和端点常数 \(c\) 并非单纯数值参数，而是四状态地基 SFT 的平方根 Perron 单位。
\end{theorem}
\begin{proof}
定理 \ref{thm:xi-terminal-real-input-positive-entropy-freezing-face} 已给出
\[
\lambda_\infty=\rho(B_\infty)=y,
\qquad
y^4-6y^3+9y^2-y-1=0,
\qquad
c=\sqrt y.
\]
把 \(y=\lambda_\infty\) 代回即得
\[
\lambda_\infty^4-6\lambda_\infty^3+9\lambda_\infty^2-\lambda_\infty-1=0.
\]
再以 \(\lambda_\infty=c^2\) 代入，得到
\[
c^8-6c^6+9c^4-c^2-1=0.
\]

由于
\[
x^4-6x^3+9x^2-x-1
\]
是首一整系数多项式且常数项为 \(-1\)，\(\lambda_\infty\) 的最小多项式作为其首一整系数因子，常数项必为 \(\pm 1\)。因此 \(\lambda_\infty\) 是代数单位。又因为 \(\lambda_\infty=\rho(B_\infty)\) 来自非负整数矩阵 \(B_\infty\) 的谱半径，故它同时是一个 Perron 数。证毕。
\end{proof}

\begin{theorem}[\texorpdfstring{$S_4$}{S4}--Hurwitz 链上的三重 Prym 不可能绝对单]\label{thm:xi-time-part9zo-prym-c6h-not-absolutely-simple}
沿用结论 \ref{con:xi-time-part9n-s4-all-subgroup-genera} 与定理 \ref{thm:killo-s4-burnside-kani-rosen-prym-square} 的记号，令
\[
P:=\operatorname{Prym}(C_6/H).
\]
则存在亏格 \(4\) 曲线 \(Z=X/D_8\)，使得
\[
P\sim J(Z)^2.
\]
特别地，
\[
\dim P=8,
\qquad
\End^0(P)\supset M_2(\QQ),
\]
从而 \(P\) 不可能绝对单。
\end{theorem}
\begin{proof}
由定理 \ref{thm:killo-s4-burnside-kani-rosen-prym-square}，
\[
\operatorname{Prym}(C_6/H)\sim J(X/D_8)^2.
\]
又由结论 \ref{con:xi-time-part9n-s4-all-subgroup-genera}，
\[
g(X/D_8)=4.
\]
记
\[
Z:=X/D_8,
\]
便得到
\[
P\sim J(Z)^2
\]
且
\[
\dim P=2g(Z)=8.
\]

对任意阿贝尔簇 \(A\)，都有自然嵌入
\[
M_2(\QQ)\hookrightarrow \End^0(A^2).
\]
将 \(A\) 取为 \(J(Z)\)，再利用同源不改变有理自同态代数的半单矩阵型包含，便得
\[
\End^0(P)\supset M_2(\QQ).
\]
若 \(P\) 绝对单，则 \(\End^0(P)\) 不可能含有非平凡矩阵代数 \(M_2(\QQ)\)。故 \(P\) 不绝对单。证毕。
\end{proof}

\begin{theorem}[分辨率边界的柱集代数与 \(2\)-进闭球代数完全一致]\label{thm:xi-time-part9zo-dyadic-boundary-cylinder-ball-identification}
\leanverified{paper\_xi\_time\_part9zo\_dyadic\_boundary\_cylinder\_ball\_identification}
记
\[
\partial Q:=\varprojlim_n Q_{2n}
\]
为超立方体图塔在截断映射下的逆极限。则在相容基选择下有自然同构
\[
\partial Q\cong T_2(E_{\min})\cong \ZZ_2^2.
\]
更强地，固定前 \(n\) 层二进前缀所得到的柱集，恰对应于 \(\ZZ_2^2\) 中某个模 \(2^n\) 剩余类
\[
(a,b)+2^n\ZZ_2^2,
\]
亦即最大 \(2\)-进范数下半径 \(2^{-n}\) 的一个闭球。因此，分辨率边界上的柱集代数与 \(\ZZ_2^2\) 的闭球代数完全一致。
\end{theorem}
\begin{proof}
推论 \ref{cor:killo-2adic-tate-hypercube-limit} 已经给出，在相容基
\[
P_n=2P_{n+1},
\qquad
Q_n=2Q_{n+1}
\]
之下，
\[
\varprojlim_n E_{\min}[2^n]\cong T_2(E_{\min})\cong \ZZ_2^2,
\]
而同一相容识别把图塔 \((Q_{2n})_{n\ge 1}\) 的逆极限送到该 \(2\)-进地址空间。故
\[
\partial Q\cong T_2(E_{\min})\cong \ZZ_2^2.
\]

现在固定前 \(n\) 层前缀。按照推论 \ref{cor:killo-2adic-tate-hypercube-limit} 的坐标解释，这等价于固定
\[
(a\bmod 2^n,\ b\bmod 2^n)\in (\ZZ/2^n\ZZ)^2.
\]
因此对应子集正是
\[
(a,b)+2^n\ZZ_2^2.
\]
在 \(\ZZ_2^2\) 的最大 \(2\)-进范数
\[
\|(x,y)\|_{\max}:=\max\{|x|_2,|y|_2\}
\]
下，这恰是以任一代表点为中心、半径 \(2^{-n}\) 的闭球。于是柱集的有限交并补仍与闭球余类结构一致，生成的代数正是 \(2\)-进闭球代数。证毕。
\end{proof}

\begin{theorem}[零点块对持久大纤维偏置的无力性]\label{thm:xi-time-part9zoa-zero-block-maxfiber-impotence}
\leanverified{paper\_xi\_time\_part9zoa\_zero\_block\_maxfiber\_impotence}
沿用定理 \ref{thm:xi-time-part9zo-zero-block-macroscopic-normalized-silent}、
定理 \ref{thm:xi-time-part60ab-window6-zero-block-halfscale-information-improvement}
与定理 \ref{thm:xi-time-part9m-maxfiber-uniformity-gap-chain} 的记号。若
\[
M:=F_{m+2},
\qquad
\mathcal Z_m:=\{k\in \ZZ/(M\ZZ):\ \widehat d_m(k)=0\},
\qquad
\eta_m:=\frac{|\mathcal Z_m|}{M},
\]
并且
\[
m+2\equiv 0\pmod 6,
\]
则虽有确定性的半阶零点块下界
\[
|\mathcal Z_m|\ge F_{(m+2)/2},
\]
但仍满足
\[
\eta_m=O(\varphi^{-m/2}).
\]
更具体地，任何仅通过
\[
M-|\mathcal Z_m|,
\qquad
\log(M-|\mathcal Z_m|),
\qquad
\sqrt{M-1-|\mathcal Z_m|}
\]
进入的零点驱动上界，其相对于零点自由基线的改变量都至多具有 \(O(\eta_m)=O(\varphi^{-m/2})\) 的相对尺度。与此相对，最大纤维偏置满足
\[
\liminf_{m\to\infty}\frac{M_m}{\overline d_m}\ge 1+C_\varphi>1.
\]
因此，算术零点块无法消除持久的 \(L^\infty\) 级非均匀性；真正顽固的障碍并不在零点集合，而在大纤维偏置本身。
\end{theorem}
\begin{proof}
由定理 \ref{thm:xi-time-part9zo-zero-block-macroscopic-normalized-silent}，
\[
\eta_m=\frac{|\mathcal Z_m|}{M}
=
\tau(m+2)\,O(\varphi^{-m/2})
=
O(\varphi^{-m/2}),
\]
其中最后一步仅用到约数函数的次幂增长远慢于任意指数函数。

同一定理还给出
\[
M-|\mathcal Z_m|=M(1-\eta_m),
\]
\[
\log(M-|\mathcal Z_m|)=\log M+O(\eta_m),
\]
\[
\sqrt{M-1-|\mathcal Z_m|}
=
\sqrt M\bigl(1+O(\eta_m)\bigr).
\]
故凡是仅通过这三类表达式进入的 zero-driven 上界，其与零点自由主项相比都只能产生 \(O(\eta_m)\) 级的相对改动。这正是“零点块只能提供半指数级副修正”的精确形式。

另一方面，定理 \ref{thm:xi-time-part9m-maxfiber-uniformity-gap-chain} 已经给出
\[
\liminf_{m\to\infty}\frac{M_m}{\overline d_m}\ge 1+C_\varphi>1.
\]
因此最大纤维相对均值的偏置在极限中仍与 \(1\) 保持严格正距离，不可能被任何 \(O(\varphi^{-m/2})\) 级的 zero-block 修正所抹平。于是折叠分布的持久非均匀性并非由零点块驱动，真正的主障碍只能来自大纤维偏置。证毕。
\end{proof}

\endinput


\endinput

\subsubsection{Toeplitz--PSD 视界接口下的时间--层级统一律}\label{subsubsec:xi-time-protocol-toeplitz-psd-interface}

\begin{theorem}[Toeplitz--PSD 层级的共终端稀疏化判据]\label{thm:xi-toeplitz-psd-cofinal-sparsification}
\leanverified{paper\_xi\_toeplitz\_psd\_cofinal\_sparsification}
设
\[
C(w)=c_0+2\sum_{n\ge 1}c_n w^n
\]
为 \(\mathbb D\) 上 Carath\'eodory 候选函数，并令
\[
T_N:=(c_{j-k})_{0\le j,k\le N}.
\]
则以下两条严格等价：
\begin{enumerate}
  \item 对所有 \(N\ge 0\)，有 \(T_N\succeq 0\)；
  \item 存在任意一条共终端严格递增子序列 \(N_1<N_2<\cdots\to\infty\)，使对所有 \(k\) 有 \(T_{N_k}\succeq 0\)。
\end{enumerate}
特别地，在 Toeplitz--PSD 表述 RH 的框架中，RH 与“仅在任一共终端子序列（如 \(N_k=2^k\)）上验证 PSD”严格等价。
\end{theorem}

\begin{proof}
\((1)\Rightarrow(2)\) 平凡。
反向由主子矩阵闭包得到：任意固定 \(N\) 可嵌入某个 \(N_k\ge N\) 的左上主子矩阵；若 \(T_{N_k}\succeq 0\) 则其主子矩阵 \(T_N\succeq 0\)。
故全层级 PSD 与任一共终端子序列 PSD 等价。
\end{proof}

\paragraph{预言机坍缩与谱替身：一致有界最小实现维数的有限证书化分岔}
本段抽取一个在全文多处反复出现、但常以不同接口语言呈现的逻辑分岔：若“最小实现维数/可见自由度”在所考虑的族上保持一致有界，则许多以“对所有截断阶 \(N\)”表述的无穷约束可被消去为有限证书；反之，若不存在一致维数上界，则任意固定有限审计都无法决定核版 RH/完备性一类全局性质，并必然出现与有限数据完全一致但全局真值相反的“谱替身”。

\begin{theorem}[预言机坍缩：一致有界最小实现维数 \texorpdfstring{$\Longrightarrow$}{=>} Toeplitz--PSD 的有限截断化]\label{thm:xi-oracle-collapse-toeplitz-psd-finite-truncation}
\leanverified{paper\_xi\_oracle\_collapse\_toeplitz\_psd\_finite\_truncation}
对每个参数对 \((\ell,\theta)\)，令
\[
C_{\ell,\theta}(w)=c^{(\ell,\theta)}_0+2\sum_{n\ge 1}c^{(\ell,\theta)}_n w^n
\]
为 \(\mathbb D\) 上的 Carath\'eodory 候选函数，并令其 Toeplitz 截断为
\[
T^{(\ell,\theta)}_N:=\bigl(c^{(\ell,\theta)}_{j-k}\bigr)_{0\le j,k\le N}\qquad(N\ge 0),
\]
其中约定 \(c^{(\ell,\theta)}_{-n}:=\overline{c^{(\ell,\theta)}_n}\) 以保证 \(T^{(\ell,\theta)}_N\) 为 Hermitian。
记 \(d_{\min}^{(\ell)}(\theta)\) 为系数序列 \(c^{(\ell,\theta)}_\bullet\) 的最小线性 realization 维数。
设对所有 \((\ell,\theta)\) 都存在统一上界 \(D\) 使得 \(d_{\min}^{(\ell)}(\theta)\le D\)
（等价地：存在阶 \(\le D\) 的常系数线性递推在充分后段湮灭 \(c^{(\ell,\theta)}_\bullet\)，亦即 Hankel 秩一致不超过 \(D\)，见定理 \ref{thm:hankel-rank-minimal} 的最小实现口径）。
则存在统一的有限截断阈值 \(N_0\le 2D\)，使得对每个固定 \((\ell,\theta)\) 都有等价关系
\[
\forall N\ge 0,\;T^{(\ell,\theta)}_N\succeq 0
\quad\Longleftrightarrow\quad
\forall N\in\{0,1,\dots,N_0\},\;T^{(\ell,\theta)}_N\succeq 0.
\]
更具体地，对每个 \(N\ge 0\) 存在某个标量矩阵 \(W^{(\ell,\theta)}_N\in M_{(N_0+1)\times (N+1)}(\CC)\) 使
\[
T^{(\ell,\theta)}_N=(W^{(\ell,\theta)}_N)^\ast\,T^{(\ell,\theta)}_{N_0}\,W^{(\ell,\theta)}_N.
\]
\end{theorem}

\begin{proof}
由最小 realization 维数 \(\le D\)（或等价的 Hankel 秩 \(\le D\)）可取一条阶 \(\le D\) 的常系数线性递推，使得系数族在充分后段满足同一递推闭包；结合 Hermitian 对称 \(c_{-n}=\overline{c_n}\)，该闭包可视为作用于双侧系数索引的有限维移位不变子空间。
因此在任意截断阶 \(N\) 上，Toeplitz 主子矩阵的列空间（或行空间）由有限个“基移位列”张成，其张成维数可由 \(O(D)\) 控制；取 \(N_0:=2D\) 并固定一组张成基，可把每个 \(T_N\) 的全部列向量写成 \(T_{N_0}\) 的列向量的标量线性组合，从而得到某个 \(W_N\) 使 \(T_N=W_N^\ast T_{N_0}W_N\)。
半正定性在合同变换（压缩）下保持与反射：若 \(T_{N_0}\succeq 0\) 则 \(T_N=W_N^\ast T_{N_0}W_N\succeq 0\)；反向方向由主子矩阵闭包平凡给出。
故“对所有 \(N\)”与“对有限 \(N\le N_0\)”严格等价，且 \(N_0\) 可统一选取。证毕。
\end{proof}

\begin{corollary}[单一二进截断的全层级强制]\label{cor:xi-oracle-collapse-single-dyadic-truncation}
在定理 \ref{thm:xi-oracle-collapse-toeplitz-psd-finite-truncation} 的一致维数上界 \(d_{\min}^{(\ell)}(\theta)\le D\) 设定下，
令
\[
N_\bullet:=2^{\lceil\log_2(2D)\rceil}.
\]
则对每个固定参数点 \((\ell,\theta)\) 都有严格等价
\[
\forall N\ge 0,\;T^{(\ell,\theta)}_N\succeq 0
\quad\Longleftrightarrow\quad
T^{(\ell,\theta)}_{N_\bullet}\succeq 0.
\]
\end{corollary}
\begin{proof}
由定理 \ref{thm:xi-oracle-collapse-toeplitz-psd-finite-truncation}，存在统一阈值 \(N_0\le 2D\) 使
\[
\forall N\ge 0,\;T^{(\ell,\theta)}_N\succeq 0
\quad\Longleftrightarrow\quad
T^{(\ell,\theta)}_{N_0}\succeq 0.
\]
由定义 \(N_\bullet\ge 2D\ge N_0\)，故 \(T^{(\ell,\theta)}_{N_0}\) 是 \(T^{(\ell,\theta)}_{N_\bullet}\) 的主子矩阵。
因此 \(T^{(\ell,\theta)}_{N_\bullet}\succeq 0\Rightarrow T^{(\ell,\theta)}_{N_0}\succeq 0\Rightarrow \forall N,\ T^{(\ell,\theta)}_N\succeq 0\)。
反向方向由全层级半正定性蕴含任意固定阶截断半正定性平凡成立。证毕。
\end{proof}

\begin{corollary}[单一矩阵的数值--代数混合可审计证书]\label{cor:xi-oracle-collapse-single-matrix-auditable-certificate}
沿用推论 \ref{cor:xi-oracle-collapse-single-dyadic-truncation} 的设定并固定 \(N_\bullet\)。
设 \(\widehat T^{(\ell,\theta)}_{N_\bullet}\) 为任意可计算近似 Toeplitz 矩阵（例如由有限喷射或有理近似得到）。
若存在 \(\eta>0\) 使
\[
\lambda_{\min}\!\bigl(\widehat T^{(\ell,\theta)}_{N_\bullet}\bigr)\ge \eta,
\qquad
\bigl\|\widehat T^{(\ell,\theta)}_{N_\bullet}-T^{(\ell,\theta)}_{N_\bullet}\bigr\|_{\mathrm{op}}\le \frac{\eta}{2},
\]
则 \(T^{(\ell,\theta)}_{N_\bullet}\succeq 0\)，因而 \(\forall N\ge 0,\ T^{(\ell,\theta)}_N\succeq 0\)。
\par
若进一步仅能给出系数逐项误差界
\[
\max_{0\le n\le N_\bullet}\bigl|\,\widehat c^{(\ell,\theta)}_n-c^{(\ell,\theta)}_n\,\bigr|\le \varepsilon,
\]
则充分条件可取
\[
(N_\bullet+1)\varepsilon\le \frac{\eta}{2}.
\]
\end{corollary}
\begin{proof}
由 Weyl 特征值扰动不等式，
\[
\lambda_{\min}\!\bigl(T^{(\ell,\theta)}_{N_\bullet}\bigr)
\ge
\lambda_{\min}\!\bigl(\widehat T^{(\ell,\theta)}_{N_\bullet}\bigr)
-\bigl\|\widehat T^{(\ell,\theta)}_{N_\bullet}-T^{(\ell,\theta)}_{N_\bullet}\bigr\|_{\mathrm{op}}
\ge \frac{\eta}{2}>0,
\]
从而 \(T^{(\ell,\theta)}_{N_\bullet}\) 为正定，特别地半正定。再由推论 \ref{cor:xi-oracle-collapse-single-dyadic-truncation} 得全层级半正定性。
\par
对逐项误差界，注意 Toeplitz 误差矩阵 \(E:=\widehat T^{(\ell,\theta)}_{N_\bullet}-T^{(\ell,\theta)}_{N_\bullet}\) 满足逐元界 \(|E_{jk}|\le \varepsilon\)。
因此
\(
\|E\|_{\mathrm{op}}\le \|E\|_{1}\le (N_\bullet+1)\varepsilon
\)
（其中 \(\|\cdot\|_1\) 为最大列和范数），代回第一段条件即得。证毕。
\end{proof}

\begin{remark}[无穷约束的有限证书化：在固定有限审计切片上的量词消去]\label{rem:xi-oracle-collapse-finite-certification}
定理 \ref{thm:xi-oracle-collapse-toeplitz-psd-finite-truncation} 消去的是 Toeplitz--PSD 层级中的“\(\forall N\)”量词：一旦一致维数上界 \(D\) 成立，则每个参数点的无穷 PSD 族可被统一压缩为有限个截断阶的联合核验。
因此，在本文的有限审计语义（\(\Theta\) 为有限审计切片，见 \S\ref{subsec:pom-three-generator-residual-closure} 中对 \(\Theta\) 的约定）下，凡通过 Toeplitz/Hankel 族（或其等价的 Carath\'eodory/Herglotz 接口）表述的核版 RH/完备性约束，均可在一致有界最小 realization 维数的前提下改写为：对有限集合 \(\Theta\) 与有限阶集合 \(\{0,\dots,N_0\}\) 的联合验证。换言之，“需要对全部层级做无限检查”的预言机需求在逻辑上被坍缩为有限证书对象。
\end{remark}

\begin{corollary}[预言机消除的两轴有限化：一致维数紧性 \texorpdfstring{$+$}{+} 互素两矩阶零异常]\label{cor:xi-oracle-elimination-two-axes}
在定理 \ref{thm:xi-oracle-collapse-toeplitz-psd-finite-truncation} 的一致维数上界制度下，任一以 Toeplitz/Hankel 族表述的核版 RH/完备性断言，其“\(\forall N\)”量词可被统一消去为有限截断阶 \(N\le N_0\) 的核验。
并且在异常放大律成立的演算侧，要求“对所有矩阶 \(q\ge 2\) 异常消失”的无限矩阶审计，等价压缩为核验两个互素矩阶（例如 \(q=2,3\)；定理 \ref{thm:pom-anom-all-moments-collapse-to-two}）。
\par
因此，在本文的审计语义中，只要两条前提（统一有限截断阈值 \(N_0\) 与互素两矩阶零异常）同时成立，则谱侧与同调侧的两类预言机化量词（\(\forall N\)、\(\forall q\)）可被同一口径下的有限证书完全替代。
\end{corollary}
\begin{proof}
第一句为定理 \ref{thm:xi-oracle-collapse-toeplitz-psd-finite-truncation} 的直接重述。
第二句由定理 \ref{thm:pom-anom-all-moments-collapse-to-two} 给出；两者合并即得结论。证毕。
\end{proof}

\begin{theorem}[低秩视界证书的 \texorpdfstring{$2\kappa$}{2kappa} 完备阈值与 sharpness]\label{thm:xi-low-rank-horizon-certificate-2kappa-sharp}
\leanverified{paper\_xi\_low\_rank\_horizon\_certificate\_2kappa\_sharp}
设某一缺陷/谱指纹类 \(\mathscr C\) 满足最小 realization 维数一致有界：
\[
d_{\min}(c)\le \kappa
\qquad(c\in\mathscr C).
\]
则下列三条在 \(\mathscr C\) 上等价：
\begin{enumerate}
  \item \textbf{（全层级 Toeplitz--PSD 可行性）}
  对每个 \(c\in\mathscr C\) 与每个 \(N\ge 0\)，均有 \(T_N(c)\succeq 0\)。
  \item \textbf{（共终端子序列可行性）}
  对每个 \(c\in\mathscr C\) 与任意满足 \(N_j\to\infty\) 的共终端子序列 \(\{N_j\}_{j\ge 1}\)，均有 \(T_{N_j}(c)\succeq 0\) 对所有 \(j\) 成立。
  \item \textbf{（有限截断可行性）}
  对每个 \(c\in\mathscr C\) 与每个 \(0\le N\le 2\kappa\)，均有 \(T_N(c)\succeq 0\)。
\end{enumerate}
若进一步把 \(\mathscr C\) 限制到一般位置的低秩子类（无根碰撞、无退化权重），则上面三条又等价于：
\begin{enumerate}
  \setcounter{enumi}{3}
  \item \textbf{（长度 \texorpdfstring{$2\kappa$}{2kappa} 前缀唯一恢复）}
  长度 \(2\kappa\) 的样本前缀唯一决定 Hankel--Prony 数据，因而唯一决定整个谱指纹/缺陷序列。
\end{enumerate}
并且阈值 \(2\kappa\) 是 sharp 的：任何只依赖于长度 \(2\kappa-1\) 前缀的判定或恢复规则，都不能在上述一般位置低秩子类上保证唯一性。
\end{theorem}
\begin{proof}
\((1)\Rightarrow(2)\) 显然。

\((2)\Rightarrow(3)\)：固定 \(c\in\mathscr C\)。任取共终端子序列 \(\{N_j\}\) 并取某个 \(j\) 使 \(N_j\ge 2\kappa\)。由 \(T_{N_j}(c)\succeq 0\) 可知其任意主子矩阵皆半正定，故特别有
\[
T_N(c)\succeq 0
\qquad(0\le N\le 2\kappa).
\]

\((3)\Rightarrow(1)\)：这正是定理 \ref{thm:xi-oracle-collapse-toeplitz-psd-finite-truncation} 在统一维数上界 \(D=\kappa\) 下的应用。

在一般位置低秩子类上，由定理 \ref{thm:hankel-rank-minimal}，最小 realization 维数 \(\le\kappa\) 等价于 Hankel 秩 \(\le\kappa\)。
再由结论 \ref{con:xi-terminal-atomic-defect-hankel-prony}，长度 \(2\kappa\) 的样本前缀足以恢复至多 \(\kappa\) 个原子/指数模态的消去多项式与权重参数；在无碰撞、无退化权重的一般位置条件下，该恢复唯一，从而得到 \((3)\Leftrightarrow(4)\)。

最后说明 sharpness。一个 \(\kappa\)-模态的 Hankel--Prony 数据包含 \(\kappa\) 个节点参数与 \(\kappa\) 个权重参数；长度 \(2\kappa\) 前缀恰好给出闭合求解该数据的最小样本数，而仅有 \(2\kappa-1\) 个样本时，消去多项式与权重不能同时被唯一锁定。故存在两组不同的一般位置 \(\kappa\)-模态数据共享同一长度 \(2\kappa-1\) 前缀而在后续项分离，从而任何只依赖该前缀的规则都不可能在该类上保证唯一性。证毕。
\end{proof}

\begin{theorem}[谱替身：无一致维数上界 \texorpdfstring{$\Longrightarrow$}{=>} 任意有限审计不可判定]\label{thm:xi-spectral-surrogate-finite-audit-undecidable}
设不存在任何整数 \(D\) 使族 \(\{c^{(\ell,\theta)}_\bullet\}_{(\ell,\theta)}\) 的最小 realization 维数满足一致上界，即
\(
\sup_{(\ell,\theta)}d_{\min}^{(\ell)}(\theta)=\infty
\)。
令 \(\mathcal A\) 为任意固定的\emph{有限审计规则}：其输入仅依赖有限个可见数据（例如有限个 Toeplitz/Hankel 截断的主子式或谱量、有限个 \((\ell,\theta)\) 采样点、或有限长度的 primitive/迹序列片段）。
则存在两组核候选 \(\mathcal K^{(+)}\)、\(\mathcal K^{(-)}\) 满足：
\begin{enumerate}
  \item \(\mathcal K^{(+)}\) 与 \(\mathcal K^{(-)}\) 在审计 \(\mathcal A\) 所访问的一切有限数据上完全一致；
  \item 二者的全局真值相反：\(\mathcal K^{(+)}\) 满足核版 RH（或完备性），而 \(\mathcal K^{(-)}\) 不满足（或反之）。
\end{enumerate}
因此，在“不知道一致维数上界”的语义下，核版 RH/完备性对任何有限审计都不可判定：若不外置额外模式轴或不引入无限族检查，则必然存在与审计数据完全同像的谱替身。
\end{theorem}

\begin{proof}
令 \(N_{\mathcal A}\) 为审计 \(\mathcal A\) 所使用的最大截断阶（或等价的最大回归长度/最大矩阶窗口），并令 \(\Theta_{\mathcal A}\) 为其访问的参数有限子集。
由于不存在一致维数上界，可在某个更高 realization 维数层级中构造一个“沉默自由度”模块：该模块的最短回归长度严格大于 \(N_{\mathcal A}\)，从而在全部被审计的有限截断阶上不产生可见贡献（等价地：不改变 \(\mathcal A\) 可见的任何 Toeplitz/Hankel 截断与其派生标量）。
在保持所有审计可见数据不变的前提下，可安排该沉默模块在更高阶引入（或排除）达界谱点/超达界谱点，进而翻转核版 RH（或完备性）所要求的全局正性/无极点真值。
于是得到两种扩展：不接入沉默模块者与接入沉默模块者在有限审计下不可区分而全局性质相反，结论成立。证毕。
\end{proof}

\begin{corollary}[临界线极点的“空集/稠密”二相律：严格 RH 与 RH-sharp]\label{cor:xi-critical-line-poles-empty-or-dense}
设 \(M\) 为有限核矩阵，\(\lambda=\rho(M)>1\)，并以 Dirichlet 坐标定义
\[
\widehat{\zeta}_M(s):=\det\!\bigl(I-\lambda^{-s}M\bigr)^{-1}.
\]
若严格核版 RH 成立，即 \(\max_{\mu\neq\lambda}|\mu|<\sqrt{\lambda}\)，则临界线 \(\Re(s)=\tfrac12\) 上不存在任何非平凡极点。
若存在达界谱点 \(|\mu|=\sqrt{\lambda}\)（RH-sharp），则对任意整数 \(q\ge 2\) 的循环 lift \(M^{\langle q\rangle}=M\otimes P_q\)，其临界线极点虚部在群 \(\RR/(2\pi/\log\lambda)\ZZ\) 中包含一组 \(q\)-等分梳齿；当 \(q\) 取遍整数时，这些梳齿的并集在该商群中稠密。
\end{corollary}

\begin{proof}
严格 RH 下由谱--\(s\) 平面字典（参见 \S\ref{subsec:pom-rh-sharp} 及 \S\ref{subsubsec:zeta-online-kernel-cyclic-lift-primitives}）知任意非 Perron 谱点满足 \(\Re(s)=\log|\mu|/\log\lambda<\tfrac12\)，故临界线无极点。
RH-sharp 的梳齿结构与稠密化由循环相位提升命题 \ref{prop:finite-rh-phase-lift} 的第 (4) 点与其末句直接推出。证毕。
\end{proof}


\begin{proposition}[Adams 推前与 Toeplitz 主子矩阵单调性]\label{prop:xi-adams-pushforward-toeplitz-principal-minor-monotonicity}
\leanverified{paper\_xi\_adams\_pushforward\_toeplitz\_principal\_minor\_monotonicity}
若 \(C\) 为 Carath\'eodory，则存在有限正测度 \(\mu\) 使
\[
c_n=\int_{\mathbb T}\xi^n\,d\mu(\xi).
\]
对任意 \(m\ge 1\) 定义推前 \(\mu_m:=(\xi\mapsto \xi^m)_\ast\mu\)，并记
\[
c_n^{(m)}:=\int_{\mathbb T}\xi^n\,d\mu_m(\xi).
\]
则有
\[
c_n^{(m)}=c_{mn}.
\]
令
\[
T_N^{(m)}:=(c^{(m)}_{j-k})_{0\le j,k\le N},
\]
则 \(T_N^{(m)}\)（经指标重排）是 \(T_{mN}\) 的主子矩阵，故
\[
T_{mN}\succeq 0\ \Longrightarrow\ T_N^{(m)}\succeq 0.
\]
\end{proposition}

\begin{proof}
矩恒等式由推前定义直接给出：
\(
\int \xi^n\,d\mu_m=\int (\xi^m)^n\,d\mu=\int \xi^{mn}\,d\mu
\)。
因此 \(T_N^{(m)}\) 由原矩序列按 \(m\)-步抽样形成，嵌入到 \(T_{mN}\) 对应指标子集的主子矩阵中；主子矩阵保半正定，得结论。
\end{proof}

\begin{remark}[Adams 粗粒化的否证回推与“时间不增信息”原则]\label{rem:xi-adams-toeplitz-no-info-increase}
命题 \ref{prop:xi-adams-pushforward-toeplitz-principal-minor-monotonicity} 的逆否命题表明：若某个派生层级 \(T_N^{(m)}\) 非半正定，则原层级 \(T_{mN}\) 必非半正定。
因此在 Toeplitz--PSD 视界接口中，越深的层级（更大 \(N\)）自动支配其全部 Adams 粗粒化后的派生层级；任一派生层级的否证均可严格回推为原层级的否证。
\end{remark}

在本文将时间固定为单生成扫描算子迭代计数的协议语义下，可把 \(m\) 视为 Adams/Norm 型折叠尺度（时间重标定），把 \(N\) 视为 Toeplitz--PSD 证书层级（视界阶数）。
二者对“可区分性预算”的耦合由下述乘积下界刻画。

\leanverified{paper\_xi\_time\_horizon\_product\_lower\_bound}
\begin{theorem}[时间--视界不确定性律：折叠尺度与视界阶数乘积下界]\label{thm:xi-time-horizon-product-lower-bound}
设缺陷点 \(a\in\mathbb D\)，深度
\[
h(a):=1-|a|^2.
\]
对 \(m\)-折叠后位置 \(a^m\)，其深度为
\[
h_m(a):=1-|a|^{2m}=1-(1-h(a))^m.
\]
在 Toeplitz--PSD 端点面积律型否证机制（触发阶数量纲下界 \(\gtrsim h^{-1}\)）下，任意折叠尺度 \(m\) 与视界阶数 \(N\) 的混合策略必须满足
\[
mN\ \gtrsim\ h(a)^{-1}.
\]
在 \(\zeta\) 的 Cayley 像缺陷中，若
\(
w_\rho=\mathcal C(\gamma+i\delta)\in\mathbb D
\)
（\(\delta>0\)），则
\[
h(\rho)=1-|w_\rho|^2
=\frac{4\delta}{\gamma^2+(\delta+1)^2},
\qquad
Q_\rho:=h(\rho)^{-1}
=\frac{\gamma^2+(\delta+1)^2}{4\delta},
\]
故
\[
mN\ \gtrsim\ Q_\rho\asymp \frac{\gamma^2}{\delta}\quad(|\gamma|\to\infty,\ \delta\ \text{固定}).
\]
\end{theorem}

\begin{proof}
折叠把缺陷半径按幂次推进，故深度按 \(h_m\) 变换。
Toeplitz 层级触发阈值按当前深度的倒数量纲增长；
将“折叠增深”与“层级检测”合并，得到乘积预算下界。
对 Cayley 像代入显式深度公式即得 \(Q_\rho\) 估计。
\end{proof}

\begin{proposition}[乘积约束下的最优配分平方根律]\label{prop:xi-optimal-allocation-under-product-bound}
在约束
\(
mN\gtrsim Q
\)
下，任意加性代价（如 \(m+N\)）或均衡代价（如 \(\max\{m,N\}\)）的量级最优解满足
\[
m\asymp N\asymp \sqrt Q.
\]
对 \(Q=Q_\rho\) 得
\[
m\asymp N\asymp \sqrt{Q_\rho}\asymp \frac{|\gamma|}{\sqrt\delta}.
\]
\leanverified{paper\_xi\_optimal\_allocation\_under\_product\_bound}
\end{proposition}

\begin{proof}
由算术--几何均值不等式，
\(
m+N\ge 2\sqrt{mN}\gtrsim 2\sqrt Q
\)，
当 \(m,N\) 同阶时取到量级最优。
\(\max\{m,N\}\) 同理。
\end{proof}

\begin{proposition}[双对数时间的双曲加性]\label{prop:xi-double-log-time-hyperbolic-additivity}
\leanverified{paper\_xi\_double\_log\_time\_hyperbolic\_additivity}
若缺陷双曲深度满足
\(
h(a)^{-1}\asymp e^{d(0,a)}
\)，
则由乘积下界 \(mN\gtrsim h(a)^{-1}\) 得
\[
\log(mN)\ \gtrsim\ d(0,a)
\quad\Longleftrightarrow\quad
\log m+\log N\ \gtrsim\ d(0,a).
\]
故在对数尺度上，总时间预算与双曲深度同阶同构。
\end{proposition}

\begin{proof}
对 \(mN\gtrsim h^{-1}\) 取对数，并用 \(h^{-1}\asymp e^{d}\) 替换即可。
\end{proof}

\begin{theorem}[高度窗统一硬下界与混合策略不可约二次标度]\label{thm:xi-height-window-unified-quadratic-hardness}
在高度窗 \(|\gamma|\le T\) 与偏移下界 \(\delta\ge\delta_0>0\) 下，
\[
Q_{\max}(T,\delta_0)
:=\max_{|\gamma|\le T,\ \delta\ge\delta_0}h(\rho)^{-1}
=\frac{T^2+(1+\delta_0)^2}{4\delta_0}.
\]
因此任意保证窗口内可否证的混合策略（任意折叠尺度 \(m\) 与 Toeplitz 阶数 \(N\) 组合）满足最坏情形下界
\[
\sup_{|\gamma|\le T,\ \delta\ge\delta_0}\inf\{mN:\text{可触发否证}\}
\ \gtrsim\
Q_{\max}(T,\delta_0)
\asymp
\frac{T^2}{\delta_0}.
\]
该下界对全部折叠--视界混合策略一致成立，因而在该接口内二次硬度标度不可被资源重分配消除。
\end{theorem}

\begin{proof}
由 \(Q_\rho\) 对 \(|\gamma|\) 单调增、对 \(\delta\) 在下界处取最坏，得 \(Q_{\max}\) 显式式。
再用定理 \ref{thm:xi-time-horizon-product-lower-bound} 的乘积预算统一约束所有混合策略，结论成立。
\end{proof}

\begin{proposition}[折叠门尺度与 Toeplitz 层级的不可互代性]\label{prop:xi-folding-scale-vs-toeplitz-order-nonsubstitutability}
设把缺陷搬运到固定半径阈值 \(r_\star\in(0,1)\) 的最小折叠尺度为
\[
m_\star(\rho;r_\star)
=
\left\lceil
\frac{\log r_\star^2}{\log|w_\rho|^2}
\right\rceil
\sim
\frac{\gamma^2+(\delta+1)^2}{4\delta}\log\frac1{r_\star^2}
\asymp
\frac{\gamma^2}{\delta}.
\]
结合 \(mN\gtrsim Q_\rho\) 可知：\(m\) 与 \(N\) 仅能在统一乘积预算内守恒重分配，不能相互替代并突破硬度标度。
其中 \(m_\star(\rho;r_\star)\) 给出“纯折叠搬运”路线的极限代价，而 Toeplitz 阶数 \(N\) 则对应“纯视界层级”路线的极限代价；任一方向的节省必伴随另一方向在同一乘积预算上的相应放大。
\end{proposition}

\begin{proof}
\(m_\star\) 由 \(|w_\rho|^{2m}\le r_\star^2\) 的最小整数解定义，渐近展开由
\(
\log|w_\rho|^2\sim -4\delta/(\gamma^2+(\delta+1)^2)
\)
给出。
再与乘积下界合并即得不可互代性。
\end{proof}

\subsubsection{折叠--视界预算与多段协议账本的统一压缩}\label{subsubsec:xi-time-protocol-budget-ledger-compression}

\paragraph{跨接口终端结论：视界、谱与纤维几何的可复用陈述}
本段把全文中已建立的若干终端断言，以“可审计、可独立引用”的格式汇总为结论条目；每条结论均对应正文或附录中的已证明命题/定理（在条目末给出引用），因此不再重复证明细节。

\begin{conclusion}[视界证书复杂度的二次--线性双下界：Toeplitz 阶数与共动扫描覆盖均不可避免发散]\label{con:xi-terminal-horizon-certificate-quadratic-linear-lb}
固定偏移下界 \(\delta\ge \delta_0>0\) 与高度窗 \(|\gamma|\le T\)。
若要求制度内证书在该窗内对任意可能的缺陷位置给出一致可检测性，则无论采用 Toeplitz--PSD 阶数提升（矩层级），或采用共动扫描的有限覆盖（有限读出族），其资源规模均必须随 \(T\) 发散，并满足显式下界
\[
N_{\mathrm{Toep}}\ \gtrsim\ Q_{\max}(T,\delta_0)\,\log\!\bigl(1+\delta_0 Q_{\max}(T,\delta_0)\bigr)
\asymp \frac{T^2}{\delta_0}\,\log(1+T^2),
\qquad
N_{\mathrm{scan}}\ \gtrsim\ \frac{T}{t+\delta_0}\quad(t>0\ \text{固定}),
\]
其中
\[
Q_{\max}(T,\delta_0):=\frac{T^2+(1+\delta_0)^2}{4\delta_0}.
\]
\end{conclusion}
\begin{proof}[证明]
见命题 \ref{con:xi-toeplitz-vs-scan-complexity-lower-bounds}（并结合其引用的推论/定理）。
\end{proof}

\begin{conclusion}[负谱幅度读出边界层深度：有限 Toeplitz 否证器给出缺陷贴边尺度]\label{con:xi-terminal-negative-spectrum-boundary-layer-depth}
设某阶 Toeplitz 主子矩阵 \(\mathbf{T}_N\) 的最小特征值为 \(\lambda_{\min}(\mathbf{T}_N)=-\eta_N<0\)。
则存在盘内零点 \(a\in\mathbb D\) 使其边界层深度 \(h:=1-|a|\) 满足
\[
h\ \ge\ c\,\frac{\eta_N}{N},
\]
其中 \(c>0\) 为绝对常数。
因此 \(\eta_N\) 在制度语义中不仅是“是否失败”的二值否证器，而且给出缺陷贴边深度的定量证书；其自然归一化量纲为 \(\eta_N/N\)。
\end{conclusion}
\begin{proof}[证明]
见命题 \ref{con:xi-toeplitz-negative-eigenvalue-controls-depth}。
\end{proof}

\begin{conclusion}[共动地址回收的位预算律：\texorpdfstring{$2\log_2T$}{2log2T} 主阶高度依赖可被制度性消除]\label{con:xi-terminal-comoving-address-bit-budget-recovery}
固定 \(\delta_0\in(0,\tfrac12)\) 与高度窗 \(|\gamma|\le T\)。
若不外置共动前缀，则端点层最小可分辨余量 \(2^{-p}\) 的一致排除要求强制
\[
p\ \ge\ 2\log_2 T+\log_2\!\frac{1}{\delta_0}+O(1)\qquad (T\to\infty).
\]
若允许外置可寻址前缀并取共动局部化 \(x_0=\gamma\)，则可检测余量满足高度无关的统一下界
\[
h_{\gamma}(\rho)=1-|w_{\rho,\gamma}|^2=\frac{4\delta}{(1+\delta)^2}\ \ge\ \frac{4\delta_0}{(1+\delta_0)^2},
\]
从而位预算门槛退化为常数下界
\[
p\ \ge\ \log_2\!\Bigl(\frac{(1+\delta_0)^2}{4\delta_0}\Bigr),
\]
并由两式对照读出：共动前缀在高度窗极限中制度性回收的主阶位预算为 \(2\log_2T+O(1)\)。
\end{conclusion}
\begin{proof}[证明]
见定理 \ref{thm:xi-offcritical-endpoint-resolution-lower-bound} 与推论 \ref{cor:xi-comoving-prefix-bit-budget-gain}。
\end{proof}

\begin{conclusion}[Minkowski 预算障碍：最坏扭曲能量的普适上界为 \texorpdfstring{$B^{-2/d}$}{B^{-2/d}}]\label{con:xi-terminal-minkowski-budget-barrier}
在 \(d\)-维扭曲/角色预算中，令 \(B=\prod_k m_k\) 为总模预算。
若最短向量落入小角有效域，则最坏扭曲能量 \(q(m)\) 满足普适上界
\[
q(m)\ \le\ (2\pi)^2\,C_d\,(\det\Sigma)^{1/d}\,B^{-2/d},
\qquad
C_d:=\Bigl(\frac{2^{d}}{V_d}\Bigr)^{2/d},
\]
其中 \(V_d\) 为单位球体积。
因此在固定维数 \(d\) 下，任何试图以增大模预算改善最坏扭曲能量的策略，至多获得指数 \(B^{-2/d}\) 的幂律改进；该指数为结构性上界而非算法松弛。
\end{conclusion}
\begin{proof}[证明]
见定理 \ref{thm:pom-minkowski-budget-barrier}。
\end{proof}

\begin{conclusion}[回溯熵缺口的指数放大：ordinary 与 non-backtracking 的增长率差异为 \texorpdfstring{$e^{n\Delta h}$}{exp(n Dh)}]\label{con:xi-terminal-backtracking-entropy-gap-exponential}
对无向骨架的邻接谱半径 \(\rho(A)\) 与 Hashimoto 无回溯谱半径 \(\rho(B)\)，定义回溯熵缺口
\[
\Delta h:=\log\rho(A)-\log\rho(B)=\log\!\Bigl(\frac{\rho(A)}{\rho(B)}\Bigr)\ge 0.
\]
则 ordinary closed walk 与 non-backtracking closed walk 的增长率比在长度 \(n\) 上满足
\[
\Bigl(\frac{\rho(A)}{\rho(B)}\Bigr)^n=e^{n\Delta h}.
\]
在单流实例中，表 \ref{tab:parallel_addition_kernels_backtracking_entropy_gap} 给出 \(n=50\) 的放大量级（其中 \(21\)-local 的放大已达 \(10^2\) 量级），从而 Ihara primes（无回溯）正规化在量化层面成为必要步骤。
\end{conclusion}
\begin{proof}[证明]
见注 \ref{rem:ihara-entropy-gap} 与表 \ref{tab:parallel_addition_kernels_backtracking_entropy_gap} 的定义与审计输出。
\end{proof}

\begin{conclusion}[模障碍当且仅当 cyclotomic 因子：Ihara--Artin 判别在有限维代数层面闭合]\label{con:xi-terminal-mod-obstruction-cyclotomic}
固定 \(u>0\)。
对每个不可约表示 \(\rho\neq\mathbf 1\) 的扭曲 Hashimoto 算子 \(B_\rho(u)\)，令
\[
Z_K(u,t;\rho):=\det(I-tB_\rho(u))^{-1}.
\]
则存在非平凡模障碍当且仅当存在某个 \(\rho\neq\mathbf 1\) 使 \(\rho(B_\rho(u))=\rho(B_{\mathbf 1}(u))=:\lambda_1(u)\)；并且在该情形下 \(\det(I-tB_\rho(u))\) 必含 cyclotomic 因子（等价地，归一化多项式 \(\det(I-(t/\lambda_1(u))B_\rho(u))\) 在单位圆上出现单位根因子）。
反之，若某 \(\det(I-tB_\rho(u))\) 出现 cyclotomic 因子，则必有 \(\rho(B_\rho(u))=\lambda_1(u)\)。
在广义 Schur/Carath\'eodory 的 Pontryagin 实现口径下（\S\ref{subsubsec:xi-negative-square-index-pontryagin-realization}），
上述 cyclotomic 因子亦可等价表述为：归一化算子 \(\lambda_1(u)^{-1}B_\rho(u)\) 在单位圆上出现单位根特征值，
且其特征子空间在实现所诱导的不定型内积上为中性（各向同性，零签名）。
因此模障碍并非“频率偶然命中”，而是单位根处中性点谱的结构性碰撞；其可枚举性来自 cyclotomic 因子的完全分类。
\end{conclusion}
\begin{proof}[证明]
见定理 \ref{thm:conclusion75-ihara-cyclotomic-factor-classification}。
\end{proof}

\begin{conclusion}[Big-Witt/necklace 桥包的严格交换与函子性：\texorpdfstring{$\zeta=\mathrm{Euler}\circ\mathrm{Mob}\circ\mathrm{TrSeq}$}{zeta=Euler o Mob o TrSeq}]\label{con:xi-terminal-big-witt-functoriality}
对任意有限核/有限矩阵 \(T\)，其迹序列（ghost 坐标）
\(
a_n(T):=\Tr(T^n)
\)
在 intertwiner/相似变换下保持不变；Möbius/Witt 反演给出 primitive 坐标 \(p_n(T)=\frac1n\sum_{d\mid n}\mu(d)\,a_{n/d}(T)\)。
并且成立逐项恒等式
\[
\zeta_T(z):=\det(I-zT)^{-1}
=\exp\!\Bigl(\sum_{n\ge 1}\frac{a_n(T)}{n}z^n\Bigr)
=\prod_{n\ge 1}(1-z^n)^{-p_n(T)}.
\]
因此 primitive 频谱（Witt 坐标）仅依赖于 \(T\) 的 intertwiner 类而与具体呈示无关。
\leanverified{paper\_zeta\_necklace\_counting\_seeds}
\end{conclusion}
\begin{proof}[证明]
见注 \ref{rem:zetaK-commuting-diagram}；带 Adams 操作与参数函子化版本见 \S\ref{subsec:zeta-online-kernel}。
\end{proof}

\begin{conclusion}[window-\(6\) 边界态的 \texorpdfstring{$\ZZ/2$}{Z/2} sheet parity：二层别名门户由 \texorpdfstring{$F_9=34$}{F9=34} 唯一承载]\label{con:xi-terminal-window6-bdry-sheet-parity}
令 window-\(6\) 的 boundary 扇区
\(
X^{\mathrm{bdry}}_6:=\{w\in X_6:\ w_1=w_6=1\}
\)。
则对任意 \(w\in X^{\mathrm{bdry}}_6\) 有严格二层别名
\[
(\mathrm{Fold}^{\mathrm{bin}}_6)^{-1}(w)=\{V_6(w),\,V_6(w)+34\},
\]
从而 dyadic 预像在组合层面被限制为一对相距 \(F_9=34\) 的双层，实现内生的 \(\ZZ/2\) sheet parity。
\end{conclusion}
\begin{proof}[证明]
见推论 \ref{cor:terminal-foldbin6-boundary-pure-f9-alias}。
\end{proof}

\begin{conclusion}[黄金互补二元信道的最强数据处理常数：\texorpdfstring{$\delta=\rho_{\mathrm{m}}=\varphi^{-3}$}{delta=rho_m=phi^-3} 与 \texorpdfstring{$I$}{I}-收缩 \texorpdfstring{$\varphi^{-6}$}{phi^-6}]\label{con:xi-terminal-golden-binary-channel-sdpi}
取二元信道 \(W\) 使 \(p=\varphi^{-1}\)、\(q=\varphi^{-2}=1-p\)，并记 \(\Delta=p-q=\varphi^{-3}\)。
则
\[
\delta(W)=\mathrm{TV}(\mathrm{Bern}(p),\mathrm{Bern}(q))=|\Delta|=\varphi^{-3},
\qquad
\rho_{\mathrm{m}}(W)=|\Delta|=\varphi^{-3}.
\]
从而对任意马尔可夫链 \(U\to \Theta\to B\) 成立锐常数的最强数据处理不等式
\[
I(U;B)\ \le\ \rho_{\mathrm{m}}(W)^2\,I(U;\Theta)\ =\ \varphi^{-6}\,I(U;\Theta),
\]
且 \(\varphi^{-6}\) 为不可改进的收缩因子。
\end{conclusion}
\begin{proof}[证明]
见定理 \ref{thm:pom-golden-binary-channel-contraction}。
\end{proof}

\begin{conclusion}[位串 \texorpdfstring{$\to$}{->} 有理证书的可回放性：半角连分数正规化给出唯一证书或 \texorpdfstring{$\mathsf{NULL}$}{NULL}]\label{con:xi-terminal-bitstring-rational-certificate-replayability}
给定浮点位串输入 \(\widehat u\in\CC\)、最大分母 \(Q_{\max}\in\NN\) 与误差预算 \(\delta>0\)，
通过舍入前像区间 \(\mathcal I(\widehat u)\)、半角变量与确定性连分数截断，可输出唯一证书 \((p,q)\)（从而 \(u_{\mathrm{cert}}=r(p,q)\in\TT\)）或空值 \(\mathsf{NULL}\)。
并且接受判据 \(\mathrm{ACCEPT}(p,q)\) 仅依赖 \((p,q)\)、\(\mathcal I(\widehat u)\)、\(\delta\) 与固定范数上界计算；第三方仅凭位串与输出即可复核真值，从而满足严格可回放性与跨平台一致性。
\end{conclusion}
\begin{proof}[证明]
见算法 \ref{alg:unitary-two-integer-deterministic} 与命题 \ref{prop:unitary-two-integer-reproducible}。
\end{proof}

\begin{conclusion}[素长 primitive 的 Wieferich 指纹：\texorpdfstring{$p$}{p}-进退化强度由单一整数 \texorpdfstring{$\nu_p$}{nu_p} 完全分级]\label{con:xi-terminal-wieferich-fingerprint}
令 \(p\) 为奇素数，\(L_p\) 为 Lucas 数，并定义素长 Wieferich 指纹
\(
\nu_p:=v_p(L_p-1)\in\ZZ_{\ge 1}
\)。
则真实输入 40 状态核的素长 primitive 轨道数 \(p_p\) 满足
\[
v_p(p_p)=\nu_p-1,
\]
从而对任意 \(k\ge 1\) 有整数判据
\[
p^k\mid p_p\ \Longleftrightarrow\ \nu_p\ge k+1.
\]
因此素长层的“算术退化强度”由单一整数 \(\nu_p\) 完全分级：\(\nu_p=1\) 为非退化相，而 \(\nu_p\ge 2\) 对应 Lucas--Wieferich 型稀有退化相；并且退化阶数与 \(p\)-进赋值精确对齐，从而 \(\nu_p\) 可作为该退化机制的单参序参量。
特别地，\(p\mid p_p\) 当且仅当 \(p^2\mid(L_p-1)\)，即素长层的 primitive 算术退化精确等价于 Lucas--Wieferich 型事件。
\end{conclusion}
\begin{proof}[证明]
见定义 \ref{def:pom-lucas-nu-p} 与推论 \ref{cor:pom-lucas-nu-p-criterion}。
\end{proof}

\begin{conclusion}[primitive 奇偶分裂的结构性消项：\texorpdfstring{$\lambda^{n/2}$}{lambda^{n/2}} 次项在偶长上精确抵消，首个可见振荡降至 \texorpdfstring{$\lambda^{n/4}$}{lambda^{n/4}}]\label{con:xi-terminal-primitive-parity-cancellation}
记 \(\lambda=\varphi^2\) 为主增长率，\(p_n\) 为 primitive 轨道数。
则存在奇偶分裂的渐近展开
\[
p_n=\frac{\lambda^n}{n}+\frac{\bigl((-1)^n-\mathbf{1}_{2\mid n}\bigr)\lambda^{n/2}}{n}
+O\!\left(\frac{\lambda^{n/3}}{n}\right),
\]
其中偶数情形的 \(\lambda^{n/2}\) 级项在 Möbius 反演的 \(d=1,2\) 两项间发生精确抵消；因此当 \(n\) 为偶数时该次项制度性消失。
进一步在偶长 \(n=2m\) 情形，首个可见振荡次项来自达界特征值 \(-\sqrt{\lambda}\)，并满足
\[
p_{2m}=\frac{\lambda^{2m}}{2m}-\frac{(-\sqrt{\lambda})^{m}}{2m}+O\!\left(\frac{\lambda^{2m/3}}{m}\right),
\]
其规模为 \(\lambda^{n/4}\)。
\end{conclusion}
\begin{proof}[证明]
见推论 \ref{cor:pom-prim-odd-even} 与注 \ref{rem:pom-prim-even-lower-osc}。
\end{proof}

\begin{conclusion}[Fibonacci 张量闭包类的 RH-sharp 生成器：平方根界在 \texorpdfstring{$(-\varphi)^n$}{(-phi)^n} 上制度性饱和]\label{con:xi-terminal-fib-tensor-rhsharp-generator}
令
\(
F=\bigl(\begin{smallmatrix}1&1\\[2pt]1&0\end{smallmatrix}\bigr)
\)、\(B:=F\otimes F\)、\(C:=-F\)。
对任意 \(n\ge 1\) 与非负整数 \(a,b\)，定义
\[
M_{n;a,b}:=B^{\otimes n}\ \oplus\ C^{\otimes n}\ \oplus\ I_a\ \oplus\ (-I)_b,
\qquad
\lambda_n:=\rho(M_{n;a,b}).
\]
则 \(\lambda_n=\varphi^{2n}\)，并且对任意非主特征值 \(\mu\in\sigma(M_{n;a,b})\) 有平方根界 \(|\mu|\le \sqrt{\lambda_n}=\varphi^n\)，且该界在 \(C^{\otimes n}\) 的主特征值 \(\mu=(-\varphi)^n\) 上取等号。
此外，\(\Tr(M_{n;a,b}^k)\) 与 primitive 轨道数具有 Lucas--Möbius 的显式闭式表达。
\end{conclusion}
\begin{proof}[证明]
见定理 \ref{thm:fib-tensor-rhsharp-closure}。
\end{proof}

\begin{conclusion}[中心化三阶迹的基线剥离与 OPUC/CMV 参数化：大值 Gram 的反射谱化]\label{con:xi-terminal-centered-trace3-opuc-cmv}
令 \(G\) 为大值点集诱导的 Gram 矩阵，并令
\(
G^\circ:=G-\frac{\tr(G)}{|W|}I
\)。
则中心化三阶迹满足恒等式
\[
\tr\!\bigl((G^\circ)^3\bigr)
=\tr(G^3)-3\frac{\tr(G)}{|W|}\tr(G^2)+2\frac{\tr(G)^3}{|W|^2},
\]
并且在 \(R(v)=\sum_{t\in W}v^{it}\) 的表示下，与 \(\int |R|^4\) 同阶的四体项可被严格替换为方差型能量 \(E(W)-c|W|^2\)，从而中心化机制剥离确定性随机基线，仅对超随机能量部分敏感。
进一步，在带宽窗口内，\(G\) 可视为离散测度的有限截断 Toeplitz 抽样压缩；由 OPUC 理论存在唯一 Verblunsky 系数序列与 CMV 矩阵，使中心化迹量可表示为 Verblunsky 系数的有限多项式泛函并受有限节谱测度控制。
\end{conclusion}
\begin{proof}[证明]
见命题 \ref{prop:gm-centered-trace3-energy-variance} 与命题 \ref{prop:gm-toeplitz-cmv-verblunsky-embedding}。
\end{proof}

\begin{conclusion}[Sofic 扭曲算子的张量闭合：统一谱隙 \texorpdfstring{$\Rightarrow$}{=>} 全奇阶中心化迹指数压缩 \texorpdfstring{$\Rightarrow$}{=>} 可调阶矩把大值频次压回均方]\label{con:xi-terminal-sofic-tensor-twist-moment-compression}
对 sofic--Zeckendorf 转导器生成的读出模型，中心化迹量 \(\tr((G^\circ)^\ell)\) 可表达为有限多个张量扭曲转移算子谱控制的线性组合。
若对所有非平凡扭曲存在统一谱隙 \(\rho(A(\chi))\le (1-\kappa)\rho(A(1))\)，则存在 \(c_\ell(\kappa)>0\) 使
\[
|\tr((G^\circ)^\ell)|\ \ll\ \rho(A(1))^{\ell m}\,e^{-c_\ell(\kappa)m},
\]
并由此得到可调阶矩极限：在统一谱隙下可选取阶数 \(k=k(\varepsilon,\kappa)\) 使归一化矩阵谱范数满足 \(\|\widetilde G\|_{\mathrm{op}}\le 1+O(|W|^{-\varepsilon})\)，从而关键区间的大值频次被压回均方并获得严格小幂次修正。
\end{conclusion}
\begin{proof}[证明]
见定理 \ref{thm:gm-tensor-twist-centered-moments} 与推论 \ref{cor:gm-adjustable-moment-mean-square}。
\end{proof}

\begin{conclusion}[稳定层因子链的谱--熵双下界：混合、超收缩与次高斯浓缩的统一常数化]\label{con:xi-terminal-factor-chain-gap-lsi-hypercontractive}
在折叠因子链 \(P^{\mathrm{fac}}_m\)（平稳分布 \(\pi_m\)）上，谱隙与对数 Sobolev 常数满足统一下界
\[
\mathrm{gap}(P^{\mathrm{fac}}_m)\ge \frac{2}{m},
\qquad
\alpha(P^{\mathrm{fac}}_m)\ge \frac{1}{m}.
\]
因此导通率下界、Gross 超收缩以及在平方跳变能量归一化下的次高斯浓缩均可被下界化为超立方体尺度 \(m^{-1}\)，且与纤维多重度的局部粗糙性无关。
\end{conclusion}
\begin{proof}[证明]
见定理 \ref{thm:pom-fold-factor-chain-gap-lsi} 与推论 \ref{cor:pom-fold-factor-chain-derived-invariants}。
\end{proof}

\begin{conclusion}[无自环刚性：一步超立方体翻转必改变稳定类型，故因子链为严格零对角核]\label{con:xi-terminal-factor-chain-no-self-loop}
折叠因子链满足
\[
P^{\mathrm{fac}}_m(x,x)=0,\qquad \forall x\in X_m,
\]
其结构来源为：超立方体单比特翻转在微态空间上诱导的差分与同一稳定类型纤维内允许的差分不相容，从而不可能发生一步“停留”。
\end{conclusion}
\begin{proof}[证明]
见命题 \ref{prop:pom-fold-factor-chain-reversible-noloop}。
\end{proof}

\begin{conclusion}[纤维重构图的 Cartesian 立方闭包：维数与直径的完全闭式]\label{con:xi-terminal-fiber-reconstruction-cartesian}
对任意 \(x\in X_m\)，纤维重构图 \(\mathcal F_x\) 在 \(\Gamma_x\simeq\bigsqcup_i P_{\ell_i}\) 的分量长度 \((\ell_i)\) 上严格分解为 Fibonacci 立方的 Cartesian 乘积：
\[
\mathcal F_x\ \cong\ \prod_i \Gamma_{\ell_i}.
\]
并且其最大立方维数与直径具有闭式
\[
\dim_{\square}(\mathcal F_x)=\sum_i \left\lceil\frac{\ell_i}{2}\right\rceil,
\qquad
\mathrm{diam}(\mathcal F_x)=\sum_i \ell_i .
\]
\end{conclusion}
\begin{proof}[证明]
见定理 \ref{thm:pom-fiber-reconstruction-cartesian-product}。
\end{proof}
\begin{theorem}[Fold 纤维规模的 \texorpdfstring{$\varphi$}{phi}-指数律]\label{thm:xi-fold-fiber-fibonacci-phi-exponential-law}
\leanverified{paper\_xi\_fold\_fiber\_fibonacci\_phi\_exponential\_law}
令 \(\varphi=\frac{1+\sqrt5}{2}\)。
对任意 \(m\ge 1\) 与任意 \(x\in X_m\)，令候选位点诱导图分解为
\(
\Gamma_x\simeq\bigsqcup_{i=1}^{r}P_{\ell_i}
\)。
记纤维基数 \(d_m(x)=\card{\Fold_m^{-1}(x)}\)，并沿用结论 \ref{con:xi-terminal-fiber-reconstruction-cartesian} 的直径与立方维数记号。
则存在绝对常数使得以下两组不等式同时成立：
\[
\varphi^{\mathrm{diam}(\mathcal F_x)}\ \le\ d_m(x)\ \le\ \varphi^{\mathrm{diam}(\mathcal F_x)+r},
\]
\[
\varphi^{2\dim_{\square}(\mathcal F_x)-r}\ \le\ d_m(x)\ \le\ \varphi^{2\dim_{\square}(\mathcal F_x)+r}.
\]
等价地，
\[
\Bigl|\log d_m(x)-2(\log\varphi)\,\dim_{\square}(\mathcal F_x)\Bigr|\ \le\ r\log\varphi.
\]
\end{theorem}
\begin{proof}
由定理 \ref{thm:pom-fiber-indset-factorization}，有乘积表达式
\(
d_m(x)=\prod_{i=1}^{r}F_{\ell_i+2}
\)，其中 \(F_k\) 为 Fibonacci 数。
标准估计给出对一切 \(k\ge 2\) 有 \(\varphi^{k-2}\le F_k\le \varphi^{k-1}\)，故
\[
\varphi^{\ell_i}\ \le\ F_{\ell_i+2}\ \le\ \varphi^{\ell_i+1}.
\]
逐项相乘得
\[
\varphi^{\sum_i \ell_i}\ \le\ d_m(x)\ \le\ \varphi^{\sum_i \ell_i+r}.
\]
再由结论 \ref{con:xi-terminal-fiber-reconstruction-cartesian} 的恒等式 \(\mathrm{diam}(\mathcal F_x)=\sum_i\ell_i\) 得第一组不等式。
\par
对第二组不等式，仅需注意对任意整数 \(\ell\ge 1\) 有
\(
2\lceil \ell/2\rceil-1\le \ell\le 2\lceil \ell/2\rceil
\)。
对 \(\ell_i\) 求和并用 \(\dim_{\square}(\mathcal F_x)=\sum_i\lceil \ell_i/2\rceil\) 得
\[
2\dim_{\square}(\mathcal F_x)-r\ \le\ \mathrm{diam}(\mathcal F_x)\ \le\ 2\dim_{\square}(\mathcal F_x).
\]
将其代入第一组不等式即得第二组不等式与对数形式。证毕。
\end{proof}
\begin{conclusion}[纤维几何的完备同构谱：分量长度多重集为完全不变量，并决定对称群阶]\label{con:xi-terminal-fiber-geometry-complete-spectrum}
在 \(\mathcal F_x\cong\prod_i \Gamma_{\ell_i}\) 的设定下，分量长度多重集
\[
\mathsf{Spec}_{\square}(x):=\{\ell_1,\dots,\ell_r\}
\]
是 \(\mathcal F_x\) 的完备图同构不变量：若 \(\mathcal F_x\cong \mathcal F_y\) 则 \(\mathsf{Spec}_{\square}(x)=\mathsf{Spec}_{\square}(y)\)（多重集意义下）。
并且全部对称性仅来源于逐因子 \(\ZZ_2\) 反演与同长度因子的置换，从而 \(|\mathrm{Aut}(\mathcal F_x)|\) 由 \(\mathsf{Spec}_{\square}(x)\) 完全决定。
\end{conclusion}
\begin{proof}[证明]
见定理 \ref{thm:pom-fibcube-aut-z2}、定理 \ref{thm:pom-fiber-reconstruction-aut-group} 与推论 \ref{cor:pom-fiber-reconstruction-spectrum-complete}。
\end{proof}

\begin{conclusion}[Gray--Hamilton 结构的纤维化：全预像存在一位开关的常延迟遍历]\label{con:xi-terminal-gray-hamilton-fiberization}
Fibonacci 立方 \(\Gamma_n\) 存在显式 Gray--Hamilton Hamilton 路径，覆盖全部顶点且相邻仅差一位。
结合纤维重构图的 Cartesian 分解，对任意 \(x\) 的预像纤维（等价于 \(\mathcal F_x\) 的顶点集）存在一条遍历全体预像的路径，且相邻两预像仅作一次合法开关（对称差大小为 \(1\)），从而全预像枚举可在固定局部操作代价下实现逐步回放。
\end{conclusion}
\begin{proof}[证明]
见定理 \ref{thm:pom-fibcube-gray-hamilton} 与推论 \ref{cor:pom-fiber-constant-delay-traversal}。
\end{proof}

\begin{conclusion}[纤维独立集复形的同伦二分律：要么可缩，要么为唯一维度的同伦球面]\label{con:xi-terminal-fiber-indcomplex-homotopy-dichotomy}
令 \(K_x\) 为候选位点诱导图的独立集复形。
则其几何实现满足严格二分律：若某分量长度 \(\ell_i\equiv 1\ (\mathrm{mod}\ 3)\)，则 \(|K_x|\) 可缩；否则 \(|K_x|\) 同伦等价于单一球面 \(S^{\tau(x)-1}\)，且约化同调呈“唯一维度、秩为 \(1\)”的刚性。
\end{conclusion}
\begin{proof}[证明]
见定理 \ref{thm:pom-fiber-independence-complex-classification} 与推论 \ref{cor:pom-fiber-independence-complex-homology}。
\end{proof}

在纤维独立集复形的球面--可缩二分律之外，全部持久拓扑量事实上都被进一步压缩为唯一维度上的秩一对象。

\begin{theorem}[纤维拓扑的秩一律]\label{thm:xi-time-part10-fiber-topology-rankone-law}
\leanverified{paper\_xi\_time\_part10\_fiber\_topology\_rankone\_law}
沿用定理 \ref{thm:pom-fiber-independence-complex-classification}、
推论 \ref{cor:pom-fiber-independence-complex-homology} 与定理
\ref{thm:killo-fold-fiber-homotopy-parity-equivalence} 的记号。对任意稳定类型 \(x\in X_m\)，记
\[
K_x:=\mathrm{Ind}^{\Delta}(\Gamma_x).
\]
则成立：
\begin{enumerate}
  \item 若 \(|K_x|\) 非可缩，则
  \[
  |K_x|\simeq S^{\tau(x)-1}.
  \]
  \item 因而其约化同调满足
  \[
  \widetilde H_j\bigl(|K_x|;\mathbb Z\bigr)=0
  \qquad (j\neq \tau(x)-1),
  \]
  \[
  \widetilde H_{\tau(x)-1}\bigl(|K_x|;\mathbb Z\bigr)\cong \mathbb Z.
  \]
  \item 并且
  \[
  |K_x|\not\simeq *
  \iff
  d_m(x)\ \text{为奇数}.
  \]
\end{enumerate}
特别地，每条持久纤维拓扑障碍都只能在唯一维度上以秩 \(1\) 出现；不存在同一纤维同时承载两条线性独立的持久拓扑通道。
\end{theorem}
\begin{proof}
第 \(1\) 条正是定理 \ref{thm:pom-fiber-independence-complex-classification}
在非可缩情形下的结论。第 \(2\) 条则由推论
\ref{cor:pom-fiber-independence-complex-homology} 立即得到。第 \(3\) 条正是定理
\ref{thm:killo-fold-fiber-homotopy-parity-equivalence} 的等价判别。

因此，一旦 \(|K_x|\) 非可缩，其约化同调只能在唯一维度
\(\tau(x)-1\) 上非零，且该同调群秩恰为 \(1\)。这就排除了同一纤维上同时出现两条线性独立持久拓扑通道的可能性。证毕。
\end{proof}

\endinput


\begin{conclusion}[折叠条件期望的 Rényi 平坦性：最坏信息损失与 Rényi 阶数无关且由最大纤维实现]\label{con:xi-terminal-renyi-flatness-logM}
设 \(\Omega\twoheadrightarrow X\) 为有限交换折叠，\(E:\CC^\Omega\to \CC^X\) 为按纤维均匀平均的条件期望，
记最大纤维多重度 \(M:=\max_{x\in X}|\Fold^{-1}(x)|\)。
则对任意 \(\alpha\in[1,\infty]\)，交换情形下的 sandwiched Rényi 相对熵满足平坦极值律
\[
\sup_{\rho}\,D_\alpha(\rho\|\rho\circ E)=\log M,
\]
且上确界由最大纤维上一点态同时实现（对全部 \(\alpha\) 同时取等）。
\end{conclusion}
\begin{proof}[证明]
由推论 \ref{cor:op-algebra-renyi-flatness-sup-equals-log-index} 得 \(\sup_\rho D_\alpha(\rho\|\rho\circ E)=\log\mathrm{Ind}(E)\)；
在交换折叠口径下 \(\mathrm{Ind}(E)=M\)（见命题 \ref{prop:fold-condexp-index-maxfiber}），并由交换情形的极值态刻画得到可达性（推论 \ref{cor:op-algebra-commutative-dmax-extremizers-maxfiber}）。
\end{proof}

\begin{conclusion}[查询--外置信息联合口径下的全恢复强逆：查询比例 \texorpdfstring{$\beta$}{beta} 与外置信息密度 \texorpdfstring{$b$}{b} 的最优成功指数为截断线性律]\label{con:xi-terminal-linear-strong-converse-query-sideinfo}
设 \(X\) 为有限集合（\(|X|=n\ge 2\)），\(|\Omega|=N\)，给定纤维谱 \(d:X\to\NN\) 满足 \(\sum_{x\in X}d(x)=N\)。
令
\[
\mathfrak{F}(d):=\left\{F:\Omega\to X:\ |F^{-1}(x)|=d(x)\ \ \forall x\in X\right\},
\qquad
F\sim\mathrm{Unif}(\mathfrak{F}(d)),
\]
并记 \(w(x):=d(x)/N\)、\(H(w):=-\sum_{x\in X}w(x)\log w(x)\)。
进行 \(k=\lfloor \beta N\rfloor\) 次互异查询（允许自适应），得到轨迹 \(T_{1:k}\)。
随后允许额外获得 \(B=\lfloor bN\rfloor\) 比特外置信息 \(\mathrm{enc}(F)\in\{0,1\}^B\)，并以 \((T_{1:k},\mathrm{enc}(F))\) 为输入尝试精确恢复整张映射 \(F\)。
令最优成功率为 \(\mathrm{Succ}^{\mathrm{opt}}_{\beta,b}(N)\)。
若 \(n\log N=o(N)\)，则极限存在并满足
\[
\boxed{\;
\lim_{N\to\infty}\frac{1}{N}\log \mathrm{Succ}^{\mathrm{opt}}_{\beta,b}(N)
=
\min\bigl\{0,\ b\log 2-(1-\beta)H(w)\bigr\}.\;}
\]
从而临界密度为
\[
b_{\mathrm{crit}}(\beta):=\frac{(1-\beta)H(w)}{\log 2}.
\]
当 \(b<b_{\mathrm{crit}}(\beta)\) 时成功率以精确指数率 \((1-\beta)H(w)-b\log 2\) 衰减；当 \(b\ge b_{\mathrm{crit}}(\beta)\) 时成功率不呈指数衰减（指数率为 \(0\)）。
\end{conclusion}
\begin{proof}[证明]
对给定轨迹 \(T_{1:k}\)，记 \(C_k(x)\) 为查询到的值 \(x\) 的出现次数，并令
\[
\mathcal{M}_k
:=
\frac{(N-k)!}{\prod_{x\in X}(d(x)-C_k(x))!}
\]
为在该轨迹下与 \(d\) 相容的后验可行映射数（剩余 \(N-k\) 个点的多项式计数）。
由于 \(F\) 在 \(\mathfrak{F}(d)\) 上均匀，条件于 \(T_{1:k}\) 的后验在该集合上仍为均匀。
对任意长度 \(B\) 的确定性编码 \(\mathrm{enc}(F)\)，其在后验可行集合上诱导一个至多 \(2^B\) 个码字的分割；于是最优条件成功率为
\[
\mathrm{Succ}^{\mathrm{opt}}_{\beta,b}(N\mid T_{1:k})
=\min\Bigl(1,\frac{2^B}{\mathcal{M}_k}\Bigr),
\]
从而
\(
\mathrm{Succ}^{\mathrm{opt}}_{\beta,b}(N)=\mathbb{E}\bigl[\min(1,2^B/\mathcal{M}_k)\bigr]
\)。
\par
由 Stirling 展开与 \(n\log N=o(N)\)，并记剩余经验分布
\[
R_N(x):=\frac{d(x)-C_k(x)}{N-k}\in\Delta(X),
\]
则有
\(
\frac{1}{N}\log\mathcal{M}_k=(1-\beta)H(R_N)+o_{\mathbb{P}}(1)
\)。
又由超几何大数律 \(R_N\to w\)（概率收敛），得
\(
\frac{1}{N}\log\mathcal{M}_k=(1-\beta)H(w)+o_{\mathbb{P}}(1)
\)。
\par
当 \(b<b_{\mathrm{crit}}(\beta)\) 时，取任意 \(\varepsilon\in\bigl(0,(1-\beta)H(w)-b\log 2\bigr)\)，并令事件
\[
\mathcal{E}_N:=\left\{\left|\frac{1}{N}\log\mathcal{M}_k-(1-\beta)H(w)\right|\le \varepsilon\right\}.
\]
由 \(\frac{1}{N}\log\mathcal{M}_k=(1-\beta)H(w)+o_{\mathbb{P}}(1)\) 得 \(\mathbb{P}(\mathcal{E}_N)\to 1\)。
在 \(\mathcal{E}_N\) 上对充分大 \(N\) 有 \(2^B\le \mathcal{M}_k\)，从而 \(\min(1,2^B/\mathcal{M}_k)=2^B/\mathcal{M}_k\)。
再由 \(B=\lfloor bN\rfloor\) 得 \(2^B=\exp(bN\log 2+O(1))\)，且在 \(\mathcal{E}_N\) 上有
\[
\exp\!\bigl(-N((1-\beta)H(w)+\varepsilon)\bigr)
\le
\mathcal{M}_k^{-1}
\le
\exp\!\bigl(-N((1-\beta)H(w)-\varepsilon)\bigr).
\]
对 \(\mathrm{Succ}^{\mathrm{opt}}_{\beta,b}(N)=\EE[\min(1,2^B/\mathcal{M}_k)]\) 取期望并用 \(\mathbb{P}(\mathcal{E}_N^c)=o(1)\)，得到
\[
\exp\!\bigl(N(b\log 2-(1-\beta)H(w)-\varepsilon)+o(N)\bigr)
\le
\mathrm{Succ}^{\mathrm{opt}}_{\beta,b}(N)
\le
\exp\!\bigl(N(b\log 2-(1-\beta)H(w)+\varepsilon)+o(N)\bigr).
\]
令 \(\varepsilon\downarrow 0\) 即得所示负指数率。
\par
当 \(b\ge b_{\mathrm{crit}}(\beta)\) 时，取任意 \(\varepsilon>0\) 并令
\[
\mathcal{E}_N(\varepsilon):=\left\{\left|\frac{1}{N}\log\mathcal{M}_k-(1-\beta)H(w)\right|\le \varepsilon\right\},
\]
则 \(\mathbb{P}(\mathcal{E}_N(\varepsilon))\to 1\)。
由 \(B=\lfloor bN\rfloor\) 与 \(b\log 2\ge (1-\beta)H(w)\) 可知在 \(\mathcal{E}_N(\varepsilon)\) 上有
\[
\min\!\left(1,\frac{2^B}{\mathcal{M}_k}\right)
\ge
\exp(-\varepsilon N+O(1)).
\]
取期望并用 \(\mathbb{P}(\mathcal{E}_N(\varepsilon))\to 1\) 得
\[
\mathrm{Succ}^{\mathrm{opt}}_{\beta,b}(N)\ge \exp(-\varepsilon N+o(N)),
\]
从而 \(\liminf_{N\to\infty}\frac{1}{N}\log \mathrm{Succ}^{\mathrm{opt}}_{\beta,b}(N)\ge -\varepsilon\)。
令 \(\varepsilon\downarrow 0\) 得 \(\liminf\ge 0\)；再由 \(\mathrm{Succ}^{\mathrm{opt}}_{\beta,b}(N)\le 1\) 得 \(\limsup\le 0\)，故指数率为 \(0\)。
证毕。
\end{proof}

\begin{conclusion}[超临界侧的失败指数：\texorpdfstring{$b>b_{\mathrm{crit}}(\beta)$}{b>bcrit(beta)} 时失败概率的精确变分式]\label{con:xi-terminal-supercritical-failure-exponent-variational}
沿用结论 \ref{con:xi-terminal-linear-strong-converse-query-sideinfo} 的设定。
定义失败概率
\[
\mathrm{Fail}^{\mathrm{opt}}_{\beta,b}(N):=1-\mathrm{Succ}^{\mathrm{opt}}_{\beta,b}(N).
\]
设 \(n\) 固定且 \(w\) 固定，并取 \(b>b_{\mathrm{crit}}(\beta)\)。
令
\[
I_\beta(r):=
\begin{cases}
(1-\beta)D_{\mathrm{KL}}(r\|w)+\beta D_{\mathrm{KL}}\!\left(\dfrac{w-(1-\beta)r}{\beta}\,\middle\|\,w\right),
& \text{若 }\ \dfrac{w-(1-\beta)r}{\beta}\in\Delta(X),\\[10pt]
+\infty, & \text{否则},
\end{cases}
\]
则极限
\[
E_{\mathrm{fail}}(\beta,b)
:=
\lim_{N\to\infty}-\frac{1}{N}\log \mathrm{Fail}^{\mathrm{opt}}_{\beta,b}(N)
\]
存在，并满足精确变分式
\[
\boxed{\;
E_{\mathrm{fail}}(\beta,b)
=
\inf\left\{I_\beta(r):\ (1-\beta)H(r)\ge b\log 2\right\}.\;}
\]
特别地，\(E_{\mathrm{fail}}(\beta,b)>0\) 且随 \(b\) 严格增大。
\end{conclusion}
\begin{proof}[证明]
由结论 \ref{con:xi-terminal-linear-strong-converse-query-sideinfo} 的表示，
\[
\mathrm{Fail}^{\mathrm{opt}}_{\beta,b}(N)
=
\mathbb{E}\!\left[\Bigl(1-\frac{2^B}{\mathcal{M}_k}\Bigr)_+\right].
\]
同样由 Stirling 展开，
\(
\frac{1}{N}\log\mathcal{M}_k=(1-\beta)H(R_N)+o_{\mathbb{P}}(1)
\)。
当事件 \((1-\beta)H(R_N)\ge b\log 2+\varepsilon\) 发生时，有
\(
2^B/\mathcal{M}_k\le \exp(-\varepsilon N+o(N))
\)，从而条件失败趋于 \(1\)；
当 \((1-\beta)H(R_N)\le b\log 2-\varepsilon\) 时，条件失败为 \(0\)（对充分大 \(N\)）。
因此
\[
\mathrm{Fail}^{\mathrm{opt}}_{\beta,b}(N)
=
\mathbb{P}\bigl((1-\beta)H(R_N)\ge b\log 2\bigr)\cdot(1+o(1)).
\]
而 \((R_N)_{N}\) 为超几何剩余经验分布，满足大偏差原理；其速率函数正为上式所定义的 \(I_\beta\)。
应用收缩原理即得
\[
\lim_{N\to\infty}-\frac{1}{N}\log \mathrm{Fail}^{\mathrm{opt}}_{\beta,b}(N)
=\inf\{I_\beta(r):(1-\beta)H(r)\ge b\log 2\}.
\]
严格单调性由约束集合随 \(b\) 严格收缩及 \(I_\beta\) 的严格凸性给出。证毕。
\end{proof}

\begin{conclusion}[最小化器的 escort--KKT 结构：失败指数的一参幂倾斜刻画与灵敏度公式]\label{con:xi-terminal-failure-exponent-escort-kkt}
在结论 \ref{con:xi-terminal-supercritical-failure-exponent-variational} 的条件下，进一步假设 \(w\) 全支撑且非均匀，并取 \(b>b_{\mathrm{crit}}(\beta)\)。
则：
\begin{enumerate}
  \item 变分问题
  \(
  \inf\{I_\beta(r):(1-\beta)H(r)\ge b\log 2\}
  \)
  存在唯一最小化器 \(r^\star\in\Delta(X)\)，且必位于边界 \((1-\beta)H(r^\star)=b\log 2\)。
  \item 记
  \[
  q^\star:=\frac{w-(1-\beta)r^\star}{\beta}\in\Delta(X),
  \]
  则存在唯一拉格朗日乘子 \(\theta^\star>0\) 使得 \(q^\star\) 与 \(r^\star\) 满足幂倾斜（escort）关系
  \[
  \boxed{\;
  q^\star(x)\ \propto\ r^\star(x)^{\,1+\theta^\star},\qquad x\in X.\;}
  \]
  \item 失败指数满足
  \[
  E_{\mathrm{fail}}(\beta,b)=I_\beta(r^\star)
  =(1-\beta)D_{\mathrm{KL}}(r^\star\|w)+\beta D_{\mathrm{KL}}(q^\star\|w),
  \]
  且 \(b\mapsto E_{\mathrm{fail}}(\beta,b)\) 在超临界侧可微，其导数由乘子给出：
  \[
  \boxed{\;\frac{\partial}{\partial b}E_{\mathrm{fail}}(\beta,b)=\theta^\star\log 2.\;}
  \]
\end{enumerate}
\end{conclusion}
\begin{proof}[证明]
可行域 \(\{r:(1-\beta)H(r)\ge b\log 2\}\) 为凸集（因 \(H\) 凹，超水平集凸），而 \(I_\beta\) 在全支撑内为严格凸函数，故最小化器唯一且必满足约束取等（否则可沿指向 \(w\) 的线段降低 \(I_\beta\) 并保持可行）。
\par
对带不等式约束的凸优化应用 KKT 条件：存在唯一 \(\theta^\star\ge 0\) 使得
\[
\nabla I_\beta(r^\star)+\theta^\star\,\nabla\bigl(b\log 2-(1-\beta)H(r^\star)\bigr)=\lambda\,\mathbf 1
\]
（\(\lambda\) 为归一化约束的乘子）。
在内点处 \(\nabla H(r)=-\log r-\mathbf 1\)，并且由 \(q^\star=(w-(1-\beta)r^\star)/\beta\) 可得
\(
\nabla I_\beta(r^\star)=(1-\beta)(\log r^\star-\log q^\star)
\)
（逐坐标计算）。
代入并吸收常数项，得到
\[
(\,1+\theta^\star)\log r^\star-\log q^\star=\mathrm{const},
\]
即 \(q^\star(x)\propto r^\star(x)^{1+\theta^\star}\)。
\par
灵敏度公式为拉格朗日乘子敏感性结论的直接应用：将约束写为 \(g_b(r):=b\log 2-(1-\beta)H(r)\le 0\)，其对参数 \(b\) 的导数为 \(\partial_b g_b=\log 2\)，故 \(\partial_b E_{\mathrm{fail}}(\beta,b)=\theta^\star\log 2\)。
证毕。
\end{proof}

\begin{conclusion}[Fibonacci 立方超平面统计的闭式：端点 \texorpdfstring{$\Theta$}{Theta}-类为唯一最大，给出体相--边界的几何判别]\label{con:xi-terminal-fibcube-hyperplane-theta-stats}
将 \(\Gamma_n\) 视为 partial cube，设 \(\Theta_i\) 为第 \(i\) 坐标翻转的 \(\Theta\)-等价类（超平面）。
则对任意 \(1\le i\le n\) 有严格闭式
\[
|\Theta_i|=F_i\,F_{n-i+1}.
\]
特别地，端点类 \(|\Theta_1|=|\Theta_n|=F_n\) 为最大 \(\Theta\)-类，并且当 \(1<i<n\) 时严格有 \(|\Theta_i|<F_n\)。
\end{conclusion}
\begin{proof}[证明]
见定理 \ref{thm:pom-fibcube-theta-class-size}。
\end{proof}

\paragraph{跨接口终端结论：复位时钟、压缩编译与中间权重的可扩展估计}
本段继续把正文/附录中已建立的若干终端断言，以“可审计、可独立引用”的格式汇总为结论条目；证明仅指向相应已完成的定理/命题，并在必要处补充由其直接推出的等价表述。

\begin{conclusion}[复位事件点过程的“有界余项”时钟：密度 \texorpdfstring{$\varphi^{-(m+2)}$}{phi^{-(m+2)}} 与误差 \texorpdfstring{$O(1)$}{O(1)} 的严格线性化]\label{con:xi-terminal-reset-bounded-remainder-clock}
令复位事件集合 \(\mathcal R_m\) 如定义 \ref{def:terminal-reset-events}，并记其递增枚举为 \(\mathcal R_m=\{r_1<r_2<\cdots\}\)。
则：
\begin{enumerate}
  \item \textbf{返回时间二值谱：}对一切 \(n\ge 1\)，差分仅取两值
  \[
  r_{n+1}-r_n\in\{F_{m+3},F_{m+4}\}.
  \]
  \item \textbf{Sturmian 差分字：}令字母 \(A:=F_{m+4}\)、\(B:=F_{m+3}\)，则差分字在字母表 \(\{A,B\}\) 上为 Fibonacci 替换不动点（Sturmian，slope \(\varphi^{-1}\)）。
  \item \textbf{密度与强一致误差界：}自然密度存在且为
  \[
  \lim_{N\to\infty}\frac{\#(\mathcal R_m\cap[0,N))}{N}=\varphi^{-(m+2)},
  \qquad
  \sup_{N\ge 1}\Bigl|\#(\mathcal R_m\cap[0,N))-\varphi^{-(m+2)}N\Bigr|\le 1.
  \]
  \item \textbf{返回时刻线性化：}存在常数 \(C_m<\infty\) 使
  \[
  \sup_{n\ge 1}\Bigl|\varphi^{-(m+2)}\,r_n-n\Bigr|\le C_m,
  \]
  因而复位时钟在尺度上为严格的 cut-and-project 点过程：主频率由 \(\varphi^{-(m+2)}\) 锚定，而偏差在 \(n\) 上不发散。
\end{enumerate}
\end{conclusion}
\begin{proof}[证明]
前三点见定理 \ref{thm:terminal-reset-events-sturmian}。
令 \(A(N):=\#(\mathcal R_m\cap[0,N))\) 且 \(\delta:=\varphi^{-(m+2)}\)，则 \(|A(N)-\delta N|\le 1\) 对全体 \(N\in\NN\) 成立。
对 \(N=r_n\) 有 \(A(r_n)=n-1\)，故 \(|(n-1)-\delta r_n|\le 1\)，从而 \(|n-\delta r_n|\le 2\)，可取 \(C_m=2\)。证毕。
\end{proof}

\begin{conclusion}[复位时钟的对数时间算法：仅用 Fibonacci 替换矩阵即可计算 \texorpdfstring{$r_n$}{rn}]\label{con:xi-terminal-reset-clock-logtime}
沿用结论 \ref{con:xi-terminal-reset-bounded-remainder-clock} 的记号，并记差分字
\(
w=w_1w_2\cdots\in\{A,B\}^{\mathbb N}
\)
为 Fibonacci 替换
\(
\sigma(A)=AB,\ \sigma(B)=A
\)
的以 \(A\) 开头的不动点（定理 \ref{thm:terminal-reset-events-sturmian}）。
对任意 \(L\ge 0\) 定义前缀计数
\[
c_A(L):=\#\{1\le t\le L:\ w_t=A\},\qquad
c_B(L):=L-c_A(L),
\]
并记 \(A:=F_{m+4}\)、\(B:=F_{m+3}\)。
则对一切 \(n\ge 1\) 有严格前缀和表达
\[
r_n=r_1+A\,c_A(n-1)+B\,c_B(n-1).
\]
并且向量 \((c_A(L),c_B(L))\) 可在 \(O(\log L)\) 时间内由替换矩阵
\[
M:=\begin{pmatrix}1&1\\ 1&0\end{pmatrix}
\]
的快速幂确定：对任意 \(k\ge 0\)，\(\sigma^k(A)\) 的字母计数等于 \(M^k(1,0)^{\mathsf T}\)，而其前缀计数可由拼接恒等式
\(
\sigma^{k+1}(A)=\sigma^k(A)\sigma^{k-1}(A)
\)
递归分解得到。
因此复位时刻序列 \(\{r_n\}\) 可作为确定性的离散时钟在 \(O(\log n)\) 时间内计算，而该复杂度不依赖于 \(m\)。
\end{conclusion}
\begin{proof}[证明]
由定理 \ref{thm:terminal-reset-events-sturmian}，差分字 \(w\) 为 Fibonacci 替换不动点，并满足
\(
\sigma^{k+1}(A)=\sigma^k(A)\sigma^{k-1}(A)
\)
的标准拼接结构。
替换矩阵 \(M\) 的定义保证：若记 \(\sigma^k(A)\) 中两字母的计数向量为
\[
\mathbf c_k:=
\begin{pmatrix}
\#\{1\le t\le |\sigma^k(A)|:\ (\sigma^k(A))_t=A\}\\
\#\{1\le t\le |\sigma^k(A)|:\ (\sigma^k(A))_t=B\}
\end{pmatrix},
\]
则 \(\mathbf c_{k+1}=M\mathbf c_k\) 且 \(\mathbf c_0=(1,0)^{\mathsf T}\)，从而 \(\mathbf c_k=M^k(1,0)^{\mathsf T}\)。
对给定 \(L\)，取最小 \(k\) 使 \(|\sigma^k(A)|\ge L\)，并沿上述拼接把长度 \(L\) 的前缀拆分为 \(O(k)\) 个块并累加字母计数，即得 \(c_A(L),c_B(L)\)。
由于 \(|\sigma^k(A)|\) 为 Fibonacci 级增长，故 \(k=O(\log L)\)；同时 \(M^k\) 可由二分幂在 \(O(\log k)\) 次 \(2\times 2\) 矩阵乘法内得到。
将 \((c_A(n-1),c_B(n-1))\) 代入前缀和表达式即得 \(r_n\)。证毕。
\end{proof}

\begin{conclusion}[Zeckendorf 时间因子图的最小非可逆性：一分支一汇聚 \texorpdfstring{$\Rightarrow$}{=>} 非单射仅支撑于 \texorpdfstring{$\mathcal R_m$}{R_m}]\label{con:xi-terminal-succ-minimal-noninvertibility}
在时间因子图 \(G^{\mathrm{Succ}}_m\)（定义 \ref{def:terminal-succ-m}）中，除唯一分支点 \(b_m\) 外所有顶点满足 \(|\mathrm{Succ}_m(w)|=1\)，且
\[
\mathrm{Succ}_m(b_m)=\{0^m,0^{m-1}1\}.
\]
同时唯一入度为 \(2\) 的汇聚点为 \(0^m\)，其前驱恰为 \(u_m\) 与 \(b_m\)。
因此 \(G^{\mathrm{Succ}}_m\) 是“单分支—单汇聚”的极小非可逆有向系统，且时间推进 \(N\mapsto N+1\) 的非单射性仅在复位事件集合 \(\mathcal R_m\)（定义 \ref{def:terminal-reset-events}）上出现。
结合结论 \ref{con:xi-terminal-reset-bounded-remainder-clock}，该非可逆性以固定密度 \(\varphi^{-(m+2)}\) 出现且无长程漂移。
\end{conclusion}
\begin{proof}[证明]
分支/汇聚结构见定理 \ref{thm:terminal-succ-unique-branch-merge}。
“非单射仅发生于 \(\mathcal R_m\)”由定义 \ref{def:terminal-reset-events} 的刻画与同一定理中 \(b_m\) 的两分支对应得到。
密度与有界余项性质由定理 \ref{thm:terminal-reset-events-sturmian} 给出，并已在结论 \ref{con:xi-terminal-reset-bounded-remainder-clock} 汇总。证毕。
\end{proof}

\begin{conclusion}[Parry 基线下的复位扇区闭式常数：\texorpdfstring{$\mathbb Q(\sqrt5)$}{Q(sqrt5)} 中的严格回返强度与条件等待时间]\label{con:xi-terminal-parry-reset-sector-constants}
在命题 \ref{prop:real-input-40-reset-regeneration-constants} 的两路独立 Parry 基线输入下，定义复位扇区事件
\[
Z_{\ge 5}:=\{d_{k-4}\cdots d_k=0^5\},
\]
则其平稳时间占比为闭式常数
\[
\mathbb P(Z_{\ge 5})=\frac{9-4\sqrt5}{5},
\]
并由 Kac 公式给出平均回返间隔
\[
\mathbb E[\tau_{Z_{\ge 5}}]=\frac{1}{\mathbb P(Z_{\ge 5})}=45+20\sqrt5.
\]
此外，从典型非复位时刻（条件于 \(Z_{\ge 5}^c\)）到下一次进入 \(Z_{\ge 5}\) 的等待时间 \(T\) 满足
\[
\mathbb E[T]=\frac{555}{8}+\frac{127}{4}\sqrt5.
\]
因此再生门槛在 Parry 基线上可被提升为全闭式的强定标常数，为同步核的相关/混合上界提供可复核的数值锚点。
\end{conclusion}
\begin{proof}[证明]
见命题 \ref{prop:real-input-40-reset-regeneration-constants}（及其证明中给出的可审计构造）。证毕。
\end{proof}

\begin{conclusion}[能量--长度 Newton 支撑的极端斜率可计算化，且真实输入核的饱和密度严格为 \texorpdfstring{$1/2$}{1/2}]\label{con:xi-terminal-newton-extreme-slopes-alpha-half}
对联合 primitive 计数 \(p_{n,k}\)（长度 \(n\)、能量 \(k\)）的 Newton 支撑，定义极端斜率
\[
\alpha_{\min}:=\inf\left\{\frac{k}{n}:p_{n,k}>0\right\},\qquad
\alpha_{\max}:=\sup\left\{\frac{k}{n}:p_{n,k}>0\right\}.
\]
则 \(\alpha_{\min},\alpha_{\max}\) 等价于有向回路的最小/最大平均能量密度，并可由 Karp 型 mean-cycle 算法在多项式时间内求出：
\[
\alpha_{\min}=\min_{\gamma\ \mathrm{cycle}}\frac{\sum_{e\in\gamma}\kappa(e)}{|\gamma|},\qquad
\alpha_{\max}=\max_{\gamma\ \mathrm{cycle}}\frac{\sum_{e\in\gamma}\kappa(e)}{|\gamma|}.
\]
同时，对压力 \(P(\theta)=\log\rho(B(e^\theta))\) 的端点饱和成立
\(
\lim_{\theta\to\pm\infty}P'(\theta)=\alpha_{\max/\min}
\)。
在真实输入 \(40\) 状态核的输出位势加权中，极端饱和密度为
\[
\alpha_{\max}=\frac12,
\]
从而不存在平均能量密度超过 \(1/2\) 的回路；等价地联合 primitive 支撑满足硬楔形域约束 \(p_{n,k}>0\Rightarrow k\le n/2\)。
\end{conclusion}
\begin{proof}[证明]
极端斜率的 mean-cycle 刻画与可计算性见命题 \ref{prop:sync-newton-extreme-slopes}。
真实输入核的 \(\alpha_{\max}=1/2\) 见推论 \ref{cor:real-input-40-alpha-max}。证毕。
\end{proof}

\begin{conclusion}[对称幂商核 \texorpdfstring{$A^{\mathrm{eq}}_q$}{Aeq} 的多项式态精确编译：\texorpdfstring{$q$}{q}-重碰撞从指数态降为组合数态]\label{con:xi-terminal-sympower-quotient-kernel-poly-state}
对实现 \(\Fold_\infty\) 的确定性在线核（内部状态集 \(Q\)），取 \(q\ge 2\) 份拷贝同步直积并施加“输出一致”筛选得到 \(q\)-重碰撞计数序列 \(S_q(m)\)。
由于 \(S_q\) 对拷贝置换 \(S_q\) 完全对称，轨道可由占据向量
\[
N_q:=\Bigl\{(n_s)_{s\in Q}\in\ZZ_{\ge 0}^{Q}:\ \sum_{s\in Q}n_s=q\Bigr\}
\]
唯一决定，从而存在非负整数邻接矩阵 \(A^{\mathrm{eq}}_q\)（指标集 \(N_q\)）及整数读出向量 \(u^{\mathrm{eq}}_q,v^{\mathrm{eq}}_q\)，使得
\[
S_q(m)= (u^{\mathrm{eq}}_q)^{\mathsf T}\,(A^{\mathrm{eq}}_q)^{\,m-m_0}\,v^{\mathrm{eq}}_q\qquad(m\ge m_0),
\]
且状态数为
\[
|N_q|=\binom{q+|Q|-1}{|Q|-1}.
\]
此外，\(A^{\mathrm{eq}}_q\) 的条目可由一组转移多项式做显式系数提取而实现，从而给出可从第一性原理生成的后端核。
\end{conclusion}
\begin{proof}[证明]
取 \(A^{\mathrm{eq}}_q:=\widetilde A_q\)、\(u^{\mathrm{eq}}_q:=\widetilde u_q\)、\(v^{\mathrm{eq}}_q:=\widetilde v_q\)，则结论即为命题 \ref{prop:pom-moment-symmetric-power}；
系数提取实现见注 \ref{rem:pom-moment-symmetric-power-coefficient-extraction}。证毕。
\end{proof}

\begin{conclusion}[二段式可执行压缩：\texorpdfstring{$\text{直积}\Rightarrow A^{\mathrm{eq}}_q\Rightarrow \mathsf{MinH}(S_q)$}{product=>Aeq=>MinH} 给出“后端可生成--前端最小证书”双轨闭环]\label{con:xi-terminal-two-stage-compression-sympower-minh}
在已核验的 \(q\le 16\) 窗内，存在严格的工程与证书分工：
\(A^{\mathrm{eq}}_q\)（结论 \ref{con:xi-terminal-sympower-quotient-kernel-poly-state}）作为“可从第一性原理生成的后端 realization”，而序列侧可用 Hankel/递推最小实现函子 \(\mathsf{MinH}\) 从 \(\{S_q(m)\}\) 直接恢复远小于 \(|N_q|\) 的最小维数 \(d_q\)，作为“前端压缩与可审计证书”。
并且该压缩可叠加：先对底层状态做碰撞计数等价类预压缩 \(Q\to Q/{\sim_{\mathrm{col}}}\)，再做对称幂商压缩；在“星形耦合”情形还可进一步压到轨道计数态 \((n_1,n_2,n_3)\)（维数 \(\binom{q+1}{2}\)）。
因此“可生成（后端）--最小化（前端）”构成一个可组合、可复核的闭环接口。
\end{conclusion}
\begin{proof}[证明]
两段压缩律与其在 \(q\le 16\) 的核验窗口的实现口径见注 \ref{rem:pom-sympower-then-minh} 与注 \ref{rem:pom-moment-symmetric-power-impl}。
碰撞计数预压缩见命题 \ref{prop:pom-moment-col-lumping}；星形耦合的 \(\binom{q+1}{2}\) 轨道表示见注 \ref{rem:pom-star-sympress}。证毕。
\end{proof}

\begin{theorem}[Hankel 秩刻画最小线性实现维数，且可由两块相邻 Hankel 主块显式构造]\label{thm:xi-hankel-rank-minimal-linear-realization}
\leanverified{paper\_xi\_hankel\_rank\_minimal\_linear\_realization}
设序列 \((a_m)_{m\ge 0}\subset\RR\) 的无穷 Hankel 矩阵为
\[
\mathrm{Hankel}(a):=\bigl(a_{i+j}\bigr)_{i,j\ge 0}.
\]
若其秩为有限整数 \(d:=\mathrm{rank}\,\mathrm{Hankel}(a)<\infty\)，并且主块
\[
H_0:=\bigl(a_{i+j}\bigr)_{0\le i,j\le d-1}
\]
可逆，则存在 \(A\in\RR^{d\times d}\) 与向量 \(\alpha,\beta\in\RR^{d}\) 使对所有 \(m\ge 0\) 成立
\[
a_m=\beta^{\mathsf T}A^{m}\alpha,
\]
并且可取显式构造
\[
H_1:=\bigl(a_{i+j+1}\bigr)_{0\le i,j\le d-1},
\qquad
A:=H_0^{-1}H_1,\qquad
\alpha:=e_0,\qquad
\beta:=(a_0,a_1,\dots,a_{d-1})^{\mathsf T}.
\]
反之，若存在表示 \(a_m=\tilde\beta^{\mathsf T}\tilde A^{m}\tilde\alpha\)（\(\tilde A\in\RR^{r\times r}\)），则
\[
\mathrm{rank}\,\mathrm{Hankel}(a)\le r.
\]
因此任意精确生成该序列的有限维线性状态空间实现的状态维数下界恰为 Hankel 秩 \(d\)。
\end{theorem}
\begin{proof}
若 \(a_m=\tilde\beta^{\mathsf T}\tilde A^{m}\tilde\alpha\)，则
\(
a_{i+j}=(\tilde\beta^{\mathsf T}\tilde A^{i})(\tilde A^{j}\tilde\alpha)
\)
给出 Hankel 矩阵为秩至多 \(r\) 的分解，故 \(\mathrm{rank}\,\mathrm{Hankel}(a)\le r\)。
\par
若 \(\mathrm{rank}\,\mathrm{Hankel}(a)=d\) 且 \(H_0\) 可逆，则 Hankel 列空间由前 \(d\) 列生成且其在前 \(d\) 个坐标上的截断给出一组基；令 \(c_j^{\mathsf T}:=(a_{i+j})_{0\le i\le d-1}\in\RR^{d}\)，并定义线性算子 \(A\) 使 \(c_j^{\mathsf T}A=c_{j+1}^{\mathsf T}\)（\(0\le j\le d-2\)）。
矩阵恒等式 \(H_0A=H_1\) 给出 \(A=H_0^{-1}H_1\)。
取 \(\alpha=e_0\) 与 \(\beta\) 为前 \(d\) 项向量即可得到 \(a_m=\beta^{\mathsf T}A^{m}\alpha\)。
\end{proof}

\begin{theorem}[Hankel 主块几何律的算术强制：\texorpdfstring{$\det(A)\in\ZZ$}{det(A) in Z} 与 \texorpdfstring{$p$}{p}-进仿射漂移]\label{thm:xi-hankel-shift-det-integrality-padic-drift}
\leanverified{paper\_xi\_hankel\_shift\_det\_integrality\_padic\_drift}
设 \((a_m)_{m\ge 0}\subset\ZZ\) 为整数序列，满足 \(\mathrm{rank}\,\mathrm{Hankel}(a)=d<\infty\) 且
\(
H_0=(a_{i+j})_{0\le i,j\le d-1}
\)
可逆，并沿定理 \ref{thm:xi-hankel-rank-minimal-linear-realization} 记
\(
H_1=(a_{i+j+1})_{0\le i,j\le d-1}
\)
与 \(A:=H_0^{-1}H_1\)。
对 \(r\ge 0\) 定义移位 Hankel 主块
\[
H_r:=\bigl(a_{i+j+r}\bigr)_{0\le i,j\le d-1}\in\ZZ^{d\times d}.
\]
若 \(A\) 可逆（等价于 \(\det(H_1)\neq 0\)），则：
\begin{enumerate}
  \item \(\det(A)\in\ZZ\setminus\{0\}\)，并且对一切 \(r\ge 0\) 有
  \[
  H_r=H_0A^{r},
  \qquad
  \det(H_r)=\det(H_0)\,\det(A)^{r}.
  \]
  \item 对任意素数 \(p\)，令 \(v_p\) 为标准 \(p\)-进赋值，则
  \[
  v_p(\det(H_r))=v_p(\det(H_0))+r\,v_p(\det(A)),
  \qquad
  v_p(\det(A))\ge 0.
  \]
\end{enumerate}
\end{theorem}
\begin{proof}
由 Hankel 秩为 \(d\) 可知列向量族 \(c_j:=(a_{i+j})_{0\le i\le d-1}\in\QQ^{d}\) 的线性包络为 \(d\) 维，且 \(H_0=[c_0\,\cdots\,c_{d-1}]\) 可逆。
定义移位算子 \(S\) 于该 \(d\) 维包络上，使 \(S(c_j)=c_{j+1}\)（\(j\ge 0\)）；因 \(\{c_0,\dots,c_{d-1}\}\) 为一组基且每个 \(c_{j+1}\) 都在其张成空间中，故 \(S\) 为良定义的线性算子。
设 \(A\) 为 \(S\) 在基 \(\{c_0,\dots,c_{d-1}\}\) 下的矩阵，则第 \(j\) 列为 \(c_{j+1}\) 在该基下的坐标，从而 \(H_0A=H_1\)（亦即 \(A=H_0^{-1}H_1\)）。
由 \(S^r(c_j)=c_{r+j}\) 得
\(
H_0A^{r}=[S^r(c_0)\ \cdots\ S^r(c_{d-1})]=[c_r\ \cdots\ c_{r+d-1}]=H_r
\)。
取行列式即得 \(\det(H_r)=\det(H_0)\det(A)^{r}\)。
\par
写 \(\det(A)=u/v\in\QQ\) 为既约分数（\(u\in\ZZ,\ v\in\ZZ_{>0},\ \gcd(u,v)=1\)）。
由上式
\[
\det(H_r)=\det(H_0)\,\frac{u^{r}}{v^{r}}\in\ZZ\setminus\{0\}\qquad(\forall r\ge 0).
\]
故 \(v^{r}\mid \det(H_0)\) 对一切 \(r\) 成立，迫使 \(v=1\)，从而 \(\det(A)=u\in\ZZ\setminus\{0\}\)。
对等式取 \(v_p\) 即得仿射律；并且因 \(\det(A)\in\ZZ\setminus\{0\}\)，故 \(v_p(\det(A))\ge 0\)。证毕。
\end{proof}

\begin{corollary}[\texorpdfstring{$p$}{p}-进刚度漂移的二分律与“坏素数”判据]\label{cor:xi-hankel-padic-stiffness-drift-dichotomy-badprimes}
\leanverified{paper\_xi\_hankel\_padic\_stiffness\_drift\_dichotomy\_badprimes}
在定理 \ref{thm:xi-hankel-shift-det-integrality-padic-drift} 的设定下，对任意素数 \(p\)：
\begin{itemize}
  \item 若 \(p\nmid\det(A)\)，则 \(v_p(\det(H_r))\) 对 \(r\ge 0\) 恒定不变；
  \item 若 \(p\mid\det(A)\)，则 \(v_p(\det(H_r))\) 随 \(r\) 严格线性增长。
\end{itemize}
从而存在某个 \(r\ge 0\) 使 \(\det(H_r)\equiv 0\pmod p\) 当且仅当
\[
p\mid \det(H_0)\det(A).
\]
\end{corollary}
\begin{proof}
由定理 \ref{thm:xi-hankel-shift-det-integrality-padic-drift} 的仿射律，
\(
v_p(\det(H_r))=v_p(\det(H_0))+r\,v_p(\det(A))
\)。
当 \(v_p(\det(A))=0\) 时右端与 \(r\) 无关；当 \(v_p(\det(A))>0\) 时右端随 \(r\) 严格增大。
最后一条等价于“存在 \(r\) 使 \(v_p(\det(H_r))\ge 1\)”并与仿射表达式等价，即得结论。证毕。
\end{proof}

\begin{proposition}[有限 Hankel 证书原则：全域一致性可由有限前缀与移位约束封口]\label{prop:xi-hankel-finite-certificate-principle}
\leanverified{paper\_xi\_hankel\_finite\_certificate\_principle}
在定理 \ref{thm:xi-hankel-rank-minimal-linear-realization} 的设定下，
设候选三元组 \((\alpha,A,\beta)\)（\(A\in\RR^{d\times d}\)）诱导候选序列 \(b_m:=\beta^{\mathsf T}A^{m}\alpha\)。
则以下两者等价：
\begin{enumerate}
  \item 对所有 \(m\ge 0\)，有 \(b_m=a_m\)；
  \item 同时满足：\(\,b_m=a_m\) 对所有 \(0\le m\le 2d-1\) 成立，且 \(AH_0=H_1\)（等价 \(A=H_0^{-1}H_1\)）。
\end{enumerate}
因此当 \(\mathrm{rank}\,\mathrm{Hankel}(a)=d\) 且 \(H_0\) 可逆时，任意声称的 \(d\)-维线性实现都可由长度 \(2d\) 的前缀数据与一次矩阵等式构成可核验的有限证书。
\end{proposition}
\begin{proof}
\((1)\Rightarrow(2)\) 为直接代入：由 \(b_m=a_m\) 得 \(H_0(b)=H_0(a)\)、\(H_1(b)=H_1(a)\)，且 \(AH_0(b)=H_1(b)\) 必成立，故 \(AH_0=H_1\)。
\par
\((2)\Rightarrow(1)\) 由移位算子刻画给出：条件 \(AH_0=H_1\) 表明 \(A\) 在 Hankel 列空间上实现“前缀到后继”的一致移位；因此 \((b_m)\) 与 \((a_m)\) 满足同一最小线性递推并在前 \(2d\) 项一致，归纳得全体项一致。
\end{proof}

\begin{conclusion}[含 \texorpdfstring{$\mathrm{MOM}[q]$}{MOM[q]} 与 \texorpdfstring{$\mathrm{TWIST}[\chi]$}{TWIST[chi]} 的投影词演算形成规范接口正规形；异常被压缩为最小标量并具同调不变性]\label{con:xi-terminal-projword-normal-form-mom-twist-anom}
在规则族 \ref{def:pom-four-gen-rewrite-rules} 与三生成元规则族 \ref{def:pom-three-gen-rewrite-rules} 的联合闭包下，任意仅由
\(\mathrm{PROJ}_u,E_m,\mathrm{LIFT}_G,\mathrm{MOM}_q\)
生成的投影词都可重写为规范接口正规形
\[
\mathrm{MOM}_Q\circ \mathrm{LIFT}_{G'}\circ \mathrm{PROJ}_{u'}\circ E_{m'}.
\]
并伴随三类可核对输出：来自 \(E\)--lift 交换的异常签名 \(\mathsf{Anom}_{G'}\)，（可选）由 \((\mathrm{RMOML})\) 的角色域展开得到的高阶异常签名塔 \(\mathsf{Anom}^{(Q)}_{G'}\)，以及（必要时）来自 \(E\)--\(\mathrm{MOM}\) 外提的标量 moment-anomaly \(\mathsf{Anom}^{(Q)}_{\mathrm{mom}}\)。
若 moment Beck--Chevalley 条件成立，则 \(\mathsf{Anom}^{(Q)}_{\mathrm{mom}}=0\)；否则两条编译路径的差异在 finite-part 层可被最小压缩为该单一标量。
并且异常标签可并入扩展标签群，使得“改变规约优先级只改变边界项、残差是否为零仍为路径不变量”的 holonomy/coboundary 机制保持形式不变。
\end{conclusion}
\begin{proof}[证明]
规范接口正规形与三类输出见命题 \ref{prop:pom-mom-interface-normal-form}。
moment-anomaly 的最小标量压缩见定义 \ref{def:pom-moment-anom-signature} 与命题 \ref{prop:pom-em-mom-residual-swap}。
同调不变性的统一表述见注 \ref{rem:pom-anom-cohomology-mom}。证毕。
\end{proof}

\begin{conclusion}[受限切片上的语义函子 \texorpdfstring{$\mathsf{Sem}$}{Sem} 为保守不变量：unitary slice 的相位参数化无外参]\label{con:xi-terminal-sem-conservative-unitary-slice}
在有界切片 \(\mathbf{PW}^{\ast}_{\le}\) 上，unitary slice（例如 \(|u|=1\)）的相位变量 \(u\in\TT\) 可由 Cayley 紧化门以无参实轴坐标 \(x\in\RR\) 生成
\(
u=\iota(x)=\frac{1+\mathrm{i}x}{1-\mathrm{i}x}
\)。
并在该切片上存在可审计语义函子
\[
\mathsf{Sem}:\mathbf{PW}^{\ast}_{\le}\longrightarrow \mathbf{Sig}_{\le}\times \mathbf{Spec}_{\le},
\]
其签名部分由接口正规形参数与有限异常采样组成，谱部分由整数矩谱、扭曲谱及派生压力/预算对象组成。
定理 \ref{thm:pom-sem-conservative} 给出保守性：\(\mathsf{Sem}\) 反射自然同构，从而“规范正规形 + 有限异常网格 + 有限谱指纹网格”构成实现无关且可复核的等价判据。
\end{conclusion}
\begin{proof}[证明]
unitary slice 的无参参数化见注 \ref{rem:pom-unitary-slice-param-by-real-line}。
\(\mathsf{Sem}\) 的定义见定义 \ref{def:pom-sem-functor}，其保守性见定理 \ref{thm:pom-sem-conservative}。证毕。
\end{proof}

\begin{conclusion}[中间权重 \texorpdfstring{$u\in(0,1)$}{u in (0,1)} 的 primitive 谱可扩展估计：Möbius--Hutchinson 的相关组合把相消噪声转化为可控方差]\label{con:xi-terminal-intermediate-u-mobius-hutchinson}
当 \(u\in(0,1)\) 时，\(M_K(u)^n\) 随 \(n\) 增大迅速趋于稠密，显式构造难以扩展；
可用 Hutchinson 随机迹估计在仅需 \(n\) 次稀疏 matvec 的代价下无偏估计
\(
a_n(u)=\Tr(M_K(u)^n)
\)，并代入 Möbius 反演得到 primitive 近似
\[
\widehat B_{K,n}(u)=\frac1n\sum_{d\mid n}\mu(d)\,\widehat a_{n/d}(u^d).
\]
为抑制 Möbius 相消带来的噪声放大，可对同一组随机探针 \(\{\mathbf r_s\}\) 同时复用全部 \(u^d\) 与全部所需幂次，并在样本层面先完成 Möbius 组合再求均值（common random numbers + correlated combining），从而把不稳定性从“独立噪声差分”转为“相关噪声抵消”。
当改为估计 PSD 对象（如 \(\Tr((M_K(u)^n)^{\mathsf T}M_K(u)^n)\) 或块对称化对象）时，Hutch++ 可将样本复杂度由 \(O(\varepsilon^{-2})\) 改进到 \(O(\varepsilon^{-1})\)（仅用 matvec）。
此外，核的全局块级对象满足单流提升关系，从而昂贵计算可在单流完成并保持可复现性。
\end{conclusion}
\begin{proof}[证明]
正文的实现口径见 \S\ref{sec:experiments} 中“中间权重 \(u\in(0,1)\)：矩阵自由的随机迹估计 + Möbius 反演”小节；
协议化与误差传播见附录 \S\ref{app:kernel-fingerprint-hutchinson}，尤其是命题 \ref{prop:hutchinson-unbiased-mfree} 与命题 \ref{prop:mobius-error-propagation}。
Hutch++ 的 PSD 轨道接口亦在该附录小节中给出。证毕。
\end{proof}

% (User input window)

\paragraph{跨接口终端结论：碰撞压力、低温激发、顶端谱隔离与 Kronecker tick 的一维传输闭式}
本段继续把正文与附录中已建立的若干结构性断言，压缩为可独立引用的结论条目；
证明仅指向相应已完成的定理/命题/推论（或其可审计输出），因此不再重复细节推导。

\begin{conclusion}[碰撞压力的离散凸性与跨阶单调律：归一化压力斜率 \texorpdfstring{$g(q)$}{g(q)} 的单调上升与一致性审计]\label{con:xi-terminal-collision-pressure-discrete-convexity-crossq-g}
设
\(
S_q(m):=\sum_{x\in X_m}d_m(x)^q
\)
并以 \(r_q\) 表示其主增长常数（例如在 Perron--Frobenius 口径下 \(S_q(m)=\Theta(r_q^m)\)）。
则对任意固定 \(m\ge 2\)，函数 \(q\mapsto \log S_q(m)\) 在 \(q\ge 0\) 上为凸；
因此压力 \(P(q):=\log r_q\) 在整数点上满足离散凸性：对任意整数 \(q_0<q_1<q_2\)，令
\(
\theta=\frac{q_2-q_1}{q_2-q_0}
\),
则
\[
P(q_1)\le \theta P(q_0)+(1-\theta)P(q_2),
\qquad
\text{等价地 } r_{q_1}\le r_{q_0}^{\theta}r_{q_2}^{1-\theta}.
\]
同时，跨阶范数不等式导出：对任意 \(2\le p<q\)，在 \(|X_m|\sim\varphi^m\) 的尺度下有下界
\[
\log r_q \ \ge\ \frac{q}{p}\log r_p+\Bigl(1-\frac{q}{p}\Bigr)\log\varphi.
\]
等价地令
\(
g(q):=\frac{\log r_q-\log\varphi}{q}
\),
则 \(g(q)\) 在 \(q\ge 2\) 上单调不减，从而 \((r_q/\varphi)^{1/q}\) 随 \(q\) 单调上升；
该单调律可作为跨实现的回归一致性审计标量。
\end{conclusion}
\begin{proof}[证明]
见命题 \ref{prop:pom-pressure-convexity}（凸性与离散凸性）与命题 \ref{prop:pom-crossq-g-monotone}（跨阶下界与 \(g(q)\) 单调性）。证毕。
\end{proof}

\begin{conclusion}[低温端的严格稀疏激发律：\texorpdfstring{$\lambda(q)=2^q\bigl(1+(3/4)^q+O((5/8)^q)\bigr)$}{lambda(q)=2^q(1+(3/4)^q+O((5/8)^q))}]\label{con:xi-terminal-low-temperature-sparse-excitation-law}
令
\(
Z_L^{(q)}=\sum_{\boldsymbol{\ell}\in \mathcal C_L}\prod_i F_{\ell_i+2}^{\,q}
\)
并以 \(\lambda(q)\) 表示其长度增长常数（\(\lambda(q)=\rho(q)^{-1}\)，其中 \(\rho(q)\) 为方程 \(W_q(\rho)=1\) 的唯一正根；见定义 \ref{def:pom-multiplicity-real-q-pressure}）。
则 \(\lambda(q)\) 在 \(q\ge 0\) 上严格递增，\(\log\lambda(q)\) 严格凸且 \(\lambda(q)\) 为实解析函数；
并在低温极限 \(q\to\infty\) 给出精确渐近展开
\[
\lambda(q)=2^q\Bigl(1+\Bigl(\frac34\Bigr)^q+O\Bigl(\Bigl(\frac58\Bigr)^q\Bigr)\Bigr),
\qquad
\log_2\lambda(q)=q+\frac{1}{\ln 2}\Bigl(\frac34\Bigr)^q+O\Bigl(\Bigl(\frac58\Bigr)^q\Bigr).
\]
因此低温端由三常数 \(\{1,3/4,5/8\}\) 给出一条完全可审计的稀疏激发律：基态标尺为 \(2^q\)，首个激发态相对权重为 \((3/4)^q\)，而更高激发被 \((5/8)^q\) 量级统一压制。
\end{conclusion}
\begin{proof}[证明]
单调/凸性/解析性见定理 \ref{thm:pom-multiplicity-real-q-pressure}；
低温渐近展开见命题 \ref{prop:pom-multiplicity-lambdaq-large-q-asymptotic}；
将展开解释为低温激发律见命题 \ref{prop:pom-multiplicity-low-temperature-gibbs}。证毕。
\end{proof}

\begin{conclusion}[顶端谱的 Fibonacci 梯级缺口：第二、第三大纤维多重度闭式与奇偶极限常数]\label{con:xi-terminal-top-spectrum-fibonacci-ladder-gaps}
令 \(D_m:=\max_{x\in X_m}d_m(x)\)，并令 \(D_m^{(2)}\)、\(D_m^{(3)}\) 分别为第二、第三大纤维多重度（定义 \ref{def:pom-top-fiber-spectrum}）。
则当 \(k\ge 5\) 时存在闭式
\[
D_{2k}^{(2)}=F_{k+2}-F_{k-4},\qquad
D_{2k+1}^{(2)}=2F_{k+1}-F_{k-4},
\]
并且当 \(k\ge 6\) 时
\[
D_{2k}^{(3)}=F_{k+2}-F_{k-3}.
\]
奇数支在稳定相 \(m=2k+1\ge 19\) 进入精确闭式
\[
D_{2k+1}^{(3)}=2F_{k+1}-\bigl(F_{k-4}+F_{k-8}\bigr).
\]
从而顶端缺口在指数尺度上呈刚性 Fibonacci 梯级：对偶数支有
\(
D_{2k}-D_{2k}^{(2)}=F_{k-4}
\)、
\(
D_{2k}-D_{2k}^{(3)}=F_{k-3}
\)，并且奇数支的顶端两级缺口亦为显式 Fibonacci 组合。
进一步，顶端相对缺口
\(
G_m:=1-D_m^{(2)}/D_m
\)
具有奇偶极限常数
\[
\lim_{k\to\infty}G_{2k}=\varphi^{-6},\qquad
\lim_{k\to\infty}G_{2k+1}=\frac12\varphi^{-5},
\]
且伴随指数收敛。
\end{conclusion}
\begin{proof}[证明]
次大闭式见定理 \ref{thm:pom-second-max-fiber-closed-form}；
偶数支第三大闭式见定理 \ref{thm:pom-third-max-fiber-even-closed-form}；
奇数支第三大稳定闭式见注 \ref{rem:pom-third-max-fiber-odd-closed-form}；
梯级缺口汇总见注 \ref{rem:pom-top-fiber-spectrum-ladder}；
相对缺口极限常数见推论 \ref{cor:pom-top-gap-limit-constants}（其指数收敛由推论 \ref{cor:pom-second-max-gap-constant} 给出）。证毕。
\end{proof}

\begin{conclusion}[顶层谱的比例隔离带：除达到者外统一压至 \texorpdfstring{$\theta D_m$}{theta Dm}]\label{con:xi-terminal-top-spectrum-isolation-band}
令
\(
\mathcal{M}_m=\{x\in X_m:\ d_m(x)=D_m\}
\)
为最大达到者集合。
则对任意 \(x\notin\mathcal{M}_m\) 有 \(d_m(x)\le D_m^{(2)}\)，并且存在显式常数 \(\theta<1\)（可取全局 \(\theta=25/26\)，或渐近最坏常数 \(\theta_\star\approx 0.954915\)）使得对充分大的 \(m\) 有
\[
d_m(x)\le D_m^{(2)}\le \theta D_m\qquad(x\notin\mathcal{M}_m).
\]
因此顶端谱在比例意义下与近极值层之间存在稳定隔离带：极值层由有限相位承载，而全体非达到者在统一比例上被压入 \(\theta D_m\) 以下；
该隔离带为任何“以高阶矩强化尾部贡献”的推断提供无需额外假设的结构性输入。
\end{conclusion}
\begin{proof}[证明]
见注 \ref{rem:pom-maxfiber-isolation-band}（并结合推论 \ref{cor:pom-second-max-gap-constant} 对 \(\theta\) 的显式选取）。证毕。
\end{proof}

\begin{conclusion}[四相到二相的有效塌缩阈值：显式选择律 \texorpdfstring{$q_\delta$}{q_delta}]\label{con:xi-terminal-odd-fourphase-two-phase-collapse-threshold}
在奇数稳定相 \(m=2k+1\ge 13\) 的最大达到者族中，存在“偏置两相”与“平衡两相”的四相结构（注 \ref{rem:pom-max-fiber-achievers-bsplit}）。
记偏置最大分支
\(
A_k=F_{k+1}+F_{k-2}
\)，
平衡最大分支
\(
B_k=F_{k+1}
\)，
则比值 \(R_k=A_k/B_k\to 1+\varphi^{-3}\)。
对任意容差 \(\delta\in(0,1)\)，当
\[
q\ge q_\delta:=\Bigl\lceil \frac{\log(1/\delta)}{\log(1+\varphi^{-3})}\Bigr\rceil
\]
时，以 \(q\)-矩权重比较四相贡献，偏置两相的总占比至少为 \(1-\delta\)。
因此“最大达到者族在高阶矩下从四相有效塌缩到二相”可由单一闭式阈值 \(q_\delta\) 精确刻画。
\end{conclusion}
\begin{proof}[证明]
见命题 \ref{prop:pom-odd-fourphase-q-collapse-threshold}。证毕。
\end{proof}

\begin{conclusion}[时间纤维的立方几何压力：\texorpdfstring{$F_x(u)$}{F_x(u)} 的因子化与增长常数 \texorpdfstring{$r_\square(u)=\frac{1+\sqrt{5+4u}}{2}$}{r_square(u)=(1+sqrt(5+4u))/2}]\label{con:xi-terminal-timefiber-cubical-pressure-factorization}
令 \(C(n,k)\) 为 Fibonacci 立方 \(\Gamma_n\) 的 \(k\)-维子立方体数，并令 \(C_n(u)=\sum_{k\ge 0}C(n,k)u^k\) 为其 \(f\)-多项式（定义 \ref{def:pom-fibcube-fvector}）。
则 \(C(n,k)\) 满足二维递推
\[
C(n,k)=C(n-1,k)+C(n-2,k)+C(n-2,k-1),
\]
其生成函数为有理式
\[
\sum_{n\ge 0}C_n(u)z^n=\frac{1+z(1+u)}{1-z-(1+u)z^2},
\]
并具有显式系数公式
\[
C(n,k)=\sum_{j=k}^{\lceil n/2\rceil}\binom{n-j+1}{j}\binom{j}{k}.
\]
从而对每个固定 \(u\ge 0\)，指数增长常数存在且为
\[
\lim_{n\to\infty}\frac1n\log C_n(u)=\log r_\square(u),\qquad r_\square(u)=\frac{1+\sqrt{5+4u}}{2}.
\]
对一般纤维重构图 \(\mathcal F_x\)（分量长度多重集为 \((\ell_i)\)），其 \(f\)-多项式严格因子化
\[
F_x(u)=\prod_i C_{\ell_i}(u),
\]
从而“纤维立方压力” \(u\mapsto \log F_x(u)\) 完全由 \((\ell_i)\) 的加性组合确定，并在大分量极限中由同一闭式常数 \(r_\square(u)\) 统一控制。
\end{conclusion}
\begin{proof}[证明]
二维递推、生成函数与系数公式见定理 \ref{thm:pom-fibcube-fvector-closed}；
指数增长常数见推论 \ref{cor:pom-fibcube-fpoly-growth-constant}；
纤维 \(f\)-多项式因子化见推论 \ref{cor:pom-fiber-fpoly-factorization}。证毕。
\end{proof}

\begin{conclusion}[Fibonacci cube 的 \(f\)-多项式：全实负根、统一下界与三角参数化]\label{con:xi-terminal-fibcube-fpoly-real-negative-roots-param}
沿用定义 \ref{def:pom-fibcube-fvector} 的 \(f\)-多项式 \(C_n(u)=\sum_{k\ge 0}C(n,k)u^k\)。
则对一切 \(n\ge 0\)：
\begin{enumerate}
  \item \(\deg C_n=\left\lfloor\frac{n+1}{2}\right\rfloor\)；
  \item \(C_n\) 的全部零点均为互异的严格负实数，并满足统一下界 \(u<-\frac54\)；
  \item 全体零点可显式枚举为
  \[
  u_{n,k}=-\frac{5+\tan^2\!\left(\frac{k\pi}{n+2}\right)}{4},
  \qquad
  k=1,2,\dots,\left\lfloor\frac{n+1}{2}\right\rfloor.
  \]
\end{enumerate}
\end{conclusion}
\begin{proof}[证明要点]
由生成函数恒等式
\(
\sum_{n\ge 0}C_n(u)z^n=\frac{1+z(1+u)}{1-z-(1+u)z^2}
\)
可读出一元递推
\[
C_n(u)=C_{n-1}(u)+(1+u)C_{n-2}(u)\quad(n\ge 2),
\qquad
C_0(u)=1,\ \ C_1(u)=2+u.
\]
令 \(s:=\sqrt{5+4u}\) 并置
\(
r_\pm:=\frac{1\pm s}{2}
\)；
则 \(r_\pm\) 为特征方程 \(r^2=r+(1+u)\) 的两根，因而二阶递推的一般解为
\[
C_n(u)=\frac{r_+^{\,n+2}-r_-^{\,n+2}}{r_+-r_-}
  =\frac{r_+^{\,n+2}-r_-^{\,n+2}}{s}.
\]
于是 \(C_n(u)=0\) 等价于 \(r_+^{\,n+2}=r_-^{\,n+2}\)，亦即
\[
\left(\frac{r_+}{r_-}\right)^{n+2}=1,\qquad \frac{r_+}{r_-}\neq 1.
\]
记
\[
z(u):=\frac{r_+}{r_-}=\frac{1+s}{1-s},
\]
则零点条件为 \(z(u)^{n+2}=1\) 且 \(z(u)\neq 1\)。
若 \(u\in\RR\) 为零点，则 \(|z(u)|=1\) 迫使 \(s\) 纯虚；写作 \(s=i\tan\theta\)（\(\theta\in(0,\pi)\)）即得
\[
u=-\frac{5+\tan^2\theta}{4}<-\frac54.
\]
另一方面 \(z(u)=e^{2i\theta}\)，条件 \(z(u)^{n+2}=1\) 等价于 \(\theta=\frac{k\pi}{n+2}\)；
利用 \(\tan^2(\theta)=\tan^2(\pi-\theta)\) 去重，即得所示列表。
零点个数等于 \(\deg C_n\)（由递推的次数增长与初值直接归纳），从而已穷尽全部零点。
最后，\(\theta\mapsto -\frac{5+\tan^2\theta}{4}\) 在 \((0,\pi/2)\) 上严格单调，且 \(\theta=\frac{k\pi}{n+2}\) 为互异点，故根皆为单根。证毕。
\end{proof}

\begin{conclusion}[\(f\)-向量的严格对数凹与唯一峰值]\label{con:xi-terminal-fibcube-fvector-strict-logconcave-unimodal}
对固定 \(n\)，系数列 \(\bigl(C(n,k)\bigr)_{k=0}^{\lfloor (n+1)/2\rfloor}\) 严格对数凹：
\[
C(n,k)^2> C(n,k-1)\,C(n,k+1)
\qquad\left(1\le k\le \deg C_n-1\right),
\]
从而 \(f\)-向量严格单峰且峰值指标唯一。
\end{conclusion}
\begin{proof}[证明要点]
由结论 \ref{con:xi-terminal-fibcube-fpoly-real-negative-roots-param}，\(C_n\) 为全实根多项式且全部根为负实数，且系数非负。
Newton 不等式给出系数列严格对数凹，进而推出严格单峰与唯一峰值。证毕。
\end{proof}

\begin{conclusion}[单位圆提升与 Joukowsky--Godel 椭圆输运：根集的显式推前]\label{con:xi-terminal-fibcube-fpoly-circle-lift-jg-transport}
定义代数变换
\[
s(u):=\sqrt{5+4u},\qquad
z(u):=\frac{1+s(u)}{1-s(u)}.
\]
则对每个 \(n\ge 0\)，有等价刻画
\[
C_n(u)=0
\quad\Longleftrightarrow\quad
z(u)^{n+2}=1\ \text{且}\ z(u)\neq 1.
\]
特别地，对实零点 \(u\) 恒有 \(|z(u)|=1\)，并且 \(z(u)\) 为某个非平凡 \((n+2)\)-次单位根。
\par
进一步，取任意 \(r>0\) 并定义 Joukowsky--Godel 投影与椭圆
\[
J_r(z):=rz+r^{-1}z^{-1},\qquad E_r:=J_r(\TT).
\]
令
\[
Q_{n,r}(w):=\operatorname{Res}_z\!\bigl(z^{n+2}-1,\ r z^2-wz+r^{-1}\bigr)\in\QQ(r)[w].
\]
则 \(Q_{n,r}\) 的全部零点均落在椭圆 \(E_r\) 上，且零点多重集恰为
\[
\left\{\,J_r(\zeta):\ \zeta^{n+2}=1\,\right\}\subset E_r.
\]
\end{conclusion}
\begin{proof}[证明要点]
第一部分为结论 \ref{con:xi-terminal-fibcube-fpoly-real-negative-roots-param} 中的闭式推导直接等价改写。
第二部分：对任意 \(\zeta^{n+2}=1\)，二次式 \(r z^2-wz+r^{-1}=0\) 在 \(z=\zeta\) 处有解当且仅当 \(w=J_r(\zeta)\)，因而由结果式的消元刻画，\(Q_{n,r}(w)=0\) 当且仅当存在 \(\zeta^{n+2}=1\) 使 \(w=J_r(\zeta)\)。
又因 \(|\zeta|=1\) 时 \(J_r(\zeta)\in E_r\)，故零点落在 \(E_r\) 上且集合如所示。证毕。
\end{proof}

\begin{conclusion}[共振区的 Aitken 证书：\texorpdfstring{$(\sigma_{\mathrm{osc}},\delta_{\mathrm{ait}})$}{(sigma_osc,delta_ait)} 给出窗口长度的确定性选择律]\label{con:xi-terminal-resonance-aitken-certificate-window-law}
对共振区特征多项式 \(P_q(\lambda)\) 的根按模排序 \( |\lambda_{1,q}|>|\lambda_{2,q}|\ge|\lambda_{3,q}|\ge\cdots \)（其中 \(\lambda_{1,q}=r_q\)），定义
\[
\sigma_{\mathrm{osc}}(q):=\frac{|\lambda_{2,q}|}{r_q},\qquad
\delta_{\mathrm{ait}}(q):=\frac{|\lambda_{3,q}|}{r_q}.
\]
则比值估计
\(
R_m=S_q(m+1)/S_q(m)\to r_q
\)
的主误差在共振区表现为二周期交替几何项，衰减率由 \(\sigma_{\mathrm{osc}}(q)\) 控制；
Aitken \(\Delta^2\) 变换严格消去该交替主项，并将误差主控尺度降为
\(
\max\{\delta_{\mathrm{ait}}(q),\sigma_{\mathrm{osc}}(q)^2\}^m
\)。
因此给定目标相对误差阈值 \(\mathrm{tol}\in(0,1)\)，存在直接的窗口长度准则
\[
m\gtrsim \frac{\log(\mathrm{tol})}{\log \delta_{\mathrm{ait}}(q)}
\]
（在共振窗内可简化为 \(\delta_{\mathrm{ait}}\) 主导），从而“模 DP + 比值收敛 + Aitken”的端点估计具备完全由谱根指纹驱动的确定性证书。
\end{conclusion}
\begin{proof}[证明]
见推论 \ref{cor:pom-aitken-cancels-two-cycle}、定义 \ref{def:pom-resonance-aitken-fingerprints} 与注 \ref{rem:pom-aitken-window-audit}（含窗口长度准则与表 \ref{tab:fold_collision_resonance_aitken_certificate_q9_17}）。证毕。
\end{proof}

\begin{conclusion}[Kronecker 分母刻度的 \texorpdfstring{$W_1$}{W1} 精确闭式与单侧灵敏性：好侧二次、坏侧一次]\label{con:xi-terminal-kronecker-w1-denominator-one-sided-sensitivity}
设 \(p/q\) 为既约有理数，且 \(|\alpha-p/q|<1/q^2\)。
记 \(\delta=\alpha-p/q\)，令点集 \(P_q(\alpha)=\{\{t\alpha\}:t=0,\dots,q-1\}\) 的经验测度为 \(\mu_q\)，并记
\(
\mathsf{T}_q(\alpha)=W_1(\mu_q,\lambda)
\)。
则存在完全闭式的分母刻度公式：
\begin{enumerate}
  \item 若 \(\delta<0\)（坏侧，\(\alpha<p/q\)），则
  \[
  \mathsf{T}_q(\alpha)=\frac{1}{2q}+\frac{q-1}{2}\Bigl(\frac{p}{q}-\alpha\Bigr),
  \qquad
  \mathsf{T}_q(\alpha)=\frac12\,D_q^\ast(P_q(\alpha)).
  \]
  \item 若 \(\delta>0\)（好侧，\(\alpha>p/q\)），则
  \[
  \mathsf{T}_q(\alpha)=\frac{1}{2q}-\frac{q-1}{2}\Bigl(\alpha-\frac{p}{q}\Bigr)
  +\frac{q(q-1)(2q-1)}{6}\Bigl(\alpha-\frac{p}{q}\Bigr)^2,
  \]
  且该侧星差异退化为常数 \(D_q^\ast(P_q(\alpha))=1/q\)（每箱恰含一点）。
\end{enumerate}
因此同一“分母 tick”上，坏侧对 \(\delta\) 的响应为一阶线性（首箱双占/末箱空缺的结构性偏置），而好侧仅以二次项显影；
该单侧灵敏性在 \(W_1\) 口径下把“近似有理”的几何误差完全代数化，并给出可审计的确定性闭式预算。
\end{conclusion}
\begin{proof}[证明]
见定理 \ref{thm:kronecker-w1-denominator-closed-form} 与推论 \ref{cor:kronecker-w1-two-sided-interpretation}。证毕。
\end{proof}


% (User input window)

\paragraph{跨接口终端结论：像空间塌缩、真实输入核分离与端点/有限域审计不变量}
本段把像空间 \(Y_m\) 的共轭塌缩、真实输入核与无约束同步核的谱分离、strictification（完全正可见实现）、端点审计以及有限域影子等链条中的可复核常数与普适性质，
压缩为可独立引用的结论条目；证明仅指向相应已完成的定理/命题/推论（或其可审计输出），不再重复推导细节。

\begin{conclusion}[真实输入核的 Parry 记忆位严格解耦：\texorpdfstring{$(p_x,p_y)$}{(px,py)} 在最大熵下给出闭式边缘并严格独立]\label{con:xi-terminal-real-input-parry-memory-decoupling}
在本质 \(20\) 状态真实输入核 \(M_{\mathrm{ess}}\) 的 Parry 最大熵测度下，输入记忆
\(m=(p_x,p_y)\in\{00,01,10,11\}\) 的边缘分布满足闭式
\[
\mathbb{P}(00)=\frac{3+\sqrt5}{10},\qquad
\mathbb{P}(01)=\mathbb{P}(10)=\frac15,\qquad
\mathbb{P}(11)=\frac{3-\sqrt5}{10}.
\]
并且两比特严格独立且同边缘：
\[
\mathbb{P}(p_x=1)=\mathbb{P}(p_y=1)=\frac{5-\sqrt5}{10},
\qquad
\mathbb{P}(p_x,p_y)=\mathbb{P}(p_x)\mathbb{P}(p_y).
\]
该解耦为精确代数事实（而非近似），可由同步态条件分布表的支撑结构反推强制得到（例如特定同步态仅在单一记忆态下出现）。
\end{conclusion}
\begin{proof}[证明]
见命题 \ref{prop:real-input-40-memory-marginal} 与推论 \ref{cor:real-input-40-memory-independence}（其推导中使用了表 \ref{tab:real_input_40_parry_internal_distribution} 的支撑结构）。证毕。
\end{proof}

\begin{conclusion}[同步态—记忆信道的硬判定分层与可观测性常数：互信息与 MAP 准确率出现 \texorpdfstring{$\sqrt5$}{sqrt5} 闭式]\label{con:xi-terminal-sync-memory-hard-decision-map-accuracy}
把同步态 \(S\in Q\) 视为对记忆 \(M=(p_x,p_y)\) 的观测信道。
在 \(M_{\mathrm{ess}}\) 的 Parry 最大熵测度下，存在一组硬判定同步态使后验退化为确定值：
\[
s\in\{01\text{-}1,010,11\text{-}1\}\Rightarrow \mathbb{P}(M=00\mid s)=1,
\qquad
s\in\{0\text{-}12,1\text{-}12\}\Rightarrow \mathbb{P}(M=11\mid s)=1,
\]
且这类硬判定同步态的总概率为
\[
\mathbb{P}(01\text{-}1)+\mathbb{P}(010)+\mathbb{P}(11\text{-}1)+\mathbb{P}(0\text{-}12)+\mathbb{P}(1\text{-}12)=\frac{3}{10}.
\]
此外，互信息为可复核常数
\(
I(S;M)\approx 0.812300
\) bits。
以 MAP 规则在线判别记忆（取使 \(\mathbb{P}(\cdot\mid s)\) 最大者），最优可达准确率满足闭式
\[
\mathrm{Acc}^*(M\leftarrow S)=\sum_s \mathbb{P}(s)\max_m \mathbb{P}(m\mid s)=\frac{3\sqrt5}{10}.
\]
对若干自然二值函数亦有闭式最优准确率：
\[
\mathrm{Acc}^*(p_x\leftarrow S)=\mathrm{Acc}^*(p_y\leftarrow S)=\frac45,
\qquad
\mathrm{Acc}^*(C\leftarrow S)=\frac{17+\sqrt5}{20},
\quad
C:=\ind_{\{M=11\}}.
\]
\end{conclusion}
\begin{proof}[证明]
硬判定分层与其总概率见注 \ref{rem:real-input-40-bayes-decode}；互信息常数见注 \ref{rem:real-input-40-mutual-info}；
MAP 准确率及其派生二值函数的闭式见命题 \ref{prop:real-input-40-map-accuracies}。证毕。
\end{proof}

\begin{conclusion}[strictification 的普适表征：\texorpdfstring{$G^{\mathrm{BC}}_m$}{GBCm} 是“完全正可见实现存在性”的最大可见群商]\label{con:xi-terminal-strictification-maximal-cp-visible-quotient}
固定分辨率 \(m\) 并令 \(\mathcal{A}_m:=\CC[G^{\mathrm{BC}}_m]\)。
若可见扭曲通道族 \(\{F_{\overline\chi}\}_{\overline\chi\in\widehat{G^{\mathrm{BC}}_m}}\) 可组装为 \(\mathcal{A}_m\)-值解析函数 \(F_m:\mathbb{D}\to\mathcal{A}_m\)，且
\begin{enumerate}
  \item 异常在可见通道上严格为零；
  \item 块 Toeplitz--PSD 层级 \(\mathbf{T}^{\mathcal{A}_m}_N(F_m)\succeq 0\) 对一切 \(N\) 成立，
\end{enumerate}
则存在 \(\mathcal{A}_m\)-值正测度 \(\boldsymbol\mu_m\) 使 \(F_m\) 具有算子值 Herglotz 表示；
从而所有通道 \(F_{\overline\chi}\) 同时为同一 \(\boldsymbol\mu_m\) 的 Fourier--Pontryagin 投影（完全正可见实现存在）。
反之，若某群商 \(G\twoheadrightarrow Q\) 上存在上述完全正可见实现并使闭环接口严格交换，则该群商必唯一因子化通过 \(G^{\mathrm{BC}}_m\)；
因此 strictification 具有最大可见商的普适性质。
\end{conclusion}
\begin{proof}[证明]
见定理 \ref{thm:xi-gbc-maximal-cp-visible-realization}（并结合命题 \ref{prop:pom-bc-quotient-strictify}、\ref{prop:pom-bc-quotient-universal}）。证毕。
\end{proof}

\begin{conclusion}[Beck--Chevalley 可交换化商的角色可测性刻画]\label{con:xi-terminal-gbc-character-measurability-duality}
沿用定义 \ref{def:pom-bc-quotient} 的记号。
则可见 strictification 商的角色群恰由折叠纤维上的可测角色给出：
\begin{equation}\label{eq:xi-gbc-character-measurability-duality}
\boxed{\;
\widehat{G^{\mathrm{BC}}_m}
\ \cong\
\Bigl\{
\chi\in\widehat G:\ 
\chi\circ\ell\ \text{对}\ \sigma(\Fold_m)\ \text{可测}
\Bigr\}.
\;}
\end{equation}
等价地，对任意 \(\chi\in\widehat G\)，下列条件等价：
\begin{enumerate}
  \item \(\chi|_{N^{\mathrm{BC}}_m}=1\)；
  \item \(\chi\circ\ell\) 在每条折叠纤维 \(\Fold_m^{-1}(x)\) 上常值；
  \item \(\chi\) 唯一因子化为某个 \(\overline\chi\in\widehat{G^{\mathrm{BC}}_m}\) 与商映射 \(\pi_m:G\twoheadrightarrow G^{\mathrm{BC}}_m\) 的复合。
\end{enumerate}
因此 \(G^{\mathrm{BC}}_m\) 不仅是“严格交换可见实现”的最大群商，而且也是使全部可见角色变为折叠可测的最大群商。
\end{conclusion}

\begin{proof}
由定义 \ref{def:pom-bc-quotient}，
\[
N^{\mathrm{BC}}_m
=
\bigl\langle \ell(a)\ell(b)^{-1}:a\sim_m b\bigr\rangle.
\]
若 \(\chi|_{N^{\mathrm{BC}}_m}=1\)，则对任意 \(a\sim_m b\) 有
\[
\chi(\ell(a)\ell(b)^{-1})=1,
\]
从而 \(\chi(\ell(a))=\chi(\ell(b))\)。
这表明 \(\chi\circ\ell\) 在每条折叠纤维上常值，亦即对 \(\sigma(\Fold_m)\) 可测。
反之，若 \(\chi\circ\ell\) 在每条纤维上常值，则上式对所有生成元都成立，故 \(\chi\) 在 \(N^{\mathrm{BC}}_m\) 上平凡。
于是 (1) 与 (2) 等价。
\par
又由商群的泛性质，\(\chi|_{N^{\mathrm{BC}}_m}=1\) 当且仅当 \(\chi\) 唯一因子化通过 \(\pi_m\)，即得到 (1) 与 (3) 的等价。
最后与推论 \ref{cor:pom-bc-pontryagin-dual-classification} 合并，便得到 \eqref{eq:xi-gbc-character-measurability-duality} 的角色群刻画。
最大性只是命题 \ref{prop:pom-bc-quotient-universal} 的可测性重述。证毕。
\end{proof}

\endinput

\paragraph{可见 strictification 塔的有限饱和、异常坐标计数与终端普适商}
在有限阿贝尔 lift 群口径下，定理 \ref{thm:pom-bc-visible-quotient-eventual-stability} 已表明可见 strictification 核链 \(\{N_m^{\mathrm{BC}}\}_{m\ge 1}\) 在有限阶段后稳定。把这一有限步饱和与 Pontryagin 对偶分类结合起来，便可把“潜在异常方向”的数量与最终可见商的普适性完全压缩为两个显式的有限群不变量。

\begin{theorem}[独立异常坐标的精确计数律]\label{thm:xi-terminal-gbc-anomaly-coordinate-count}
\leanverified{paper\_xi\_terminal\_gbc\_anomaly\_coordinate\_count}
沿用定义 \ref{def:pom-bc-quotient} 的记号，并假设 \(G\) 为有限阿贝尔群。对每个分辨率 \(m\)，定义独立异常坐标群
\[
\mathcal A_m:=\widehat G/\widehat{G_m^{\mathrm{BC}}}.
\]
则存在自然同构
\[
\mathcal A_m\cong \widehat{N_m^{\mathrm{BC}}},
\]
因而其非平凡独立坐标的数量恰为
\[
\#\mathcal A_m^{\mathrm{nz}}
:=
|\mathcal A_m|-1
=
|N_m^{\mathrm{BC}}|-1.
\]
此外，\(\#\mathcal A_m^{\mathrm{nz}}\) 随 \(m\) 单调不增；并且若
\[
N_\infty^{\mathrm{BC}}
:=
\bigcap_{m\ge 1}N_m^{\mathrm{BC}},
\]
则在定理 \ref{thm:pom-bc-visible-quotient-eventual-stability} 的饱和阶段之后，有
\[
\#\mathcal A_m^{\mathrm{nz}}=|N_\infty^{\mathrm{BC}}|-1.
\]
因此，潜在异常方向在有限分辨率上的全部独立坐标，并不是抽象余维，而是由 strictification 核的大小精确计数。
\end{theorem}

\begin{proof}
由推论 \ref{cor:pom-bc-pontryagin-dual-classification}，
\[
\widehat G/\widehat{G_m^{\mathrm{BC}}}
\cong
\widehat{N_m^{\mathrm{BC}}},
\]
即得所示自然同构。
又因有限阿贝尔群与其 Pontryagin 对偶等势，故
\[
|\mathcal A_m|
=
|\widehat{N_m^{\mathrm{BC}}}|
=
|N_m^{\mathrm{BC}}|.
\]
去掉平凡角色即得
\[
\#\mathcal A_m^{\mathrm{nz}}=|N_m^{\mathrm{BC}}|-1.
\]

再由定理 \ref{thm:pom-bc-visible-quotient-eventual-stability}，
\[
N_{m+1}^{\mathrm{BC}}\le N_m^{\mathrm{BC}},
\]
故 \(|N_m^{\mathrm{BC}}|\) 随 \(m\) 单调不增，从而 \(\#\mathcal A_m^{\mathrm{nz}}\) 亦单调不增。并且在某个有限阈值 \(m_0\) 之后，
\[
N_m^{\mathrm{BC}}=N_\infty^{\mathrm{BC}}\qquad(m\ge m_0),
\]
于是
\[
\#\mathcal A_m^{\mathrm{nz}}=|N_\infty^{\mathrm{BC}}|-1
\qquad(m\ge m_0).
\]
证毕。
\end{proof}

\begin{theorem}[稳定 strictification 商的终对象性质]\label{thm:xi-terminal-gbc-stabilized-terminal-object}
\leanverified{paper\_xi\_terminal\_gbc\_stabilized\_terminal\_object}
仍设 \(G\) 为有限阿贝尔群。考虑范畴 \(\mathsf{Str}^{\mathrm{vis}}_\infty(G)\)，其对象为满射
\[
\pi:G\twoheadrightarrow H
\]
满足存在某个阈值 \(m_\pi\)，使得对一切 \(m\ge m_\pi\)，由 \(\pi\) 诱导的可见扭曲通道在分辨率 \(m\) 上均已严格交换；态射为使商映射交换的群同态。则稳定 strictification 商
\[
G_\infty^{\mathrm{BC}}
\cong
G/N_\infty^{\mathrm{BC}}
\]
是范畴 \(\mathsf{Str}^{\mathrm{vis}}_\infty(G)\) 中的终对象。

换言之，任何在足够高分辨率上最终 strict 的可见群商，都且仅都经由 \(G_\infty^{\mathrm{BC}}\) 唯一因子化。
\end{theorem}

\begin{proof}
任取对象 \(\pi:G\twoheadrightarrow H\in\mathsf{Str}^{\mathrm{vis}}_\infty(G)\)。按定义，存在 \(m_\pi\) 使得对所有 \(m\ge m_\pi\)，\(\pi\) 在分辨率 \(m\) 上均给出严格交换的可见商。

另一方面，由定理 \ref{thm:pom-bc-visible-quotient-eventual-stability}，存在 \(m_0\) 使得
\[
G_m^{\mathrm{BC}}\cong G_\infty^{\mathrm{BC}}
\qquad(m\ge m_0).
\]
令 \(M:=\max\{m_\pi,m_0\}\)。则 \(\pi\) 在分辨率 \(M\) 上已严格交换。由命题 \ref{prop:pom-bc-quotient-universal} 的普适性，\(\pi\) 必唯一因子化为
\[
G\twoheadrightarrow G_M^{\mathrm{BC}}\longrightarrow H.
\]
又因 \(M\ge m_0\)，有
\[
G_M^{\mathrm{BC}}\cong G_\infty^{\mathrm{BC}},
\]
故 \(\pi\) 亦唯一因子化为
\[
G\twoheadrightarrow G_\infty^{\mathrm{BC}}\longrightarrow H.
\]
这正是终对象的定义。证毕。
\end{proof}

\endinput


\begin{conclusion}[端点阈值承载的中心性与通道强度谱：\texorpdfstring{$\boldsymbol\mu_m(\{-1\})$}{mu(-1)} 诱导不可约通道的非负标量分解]\label{con:xi-terminal-endpoint-atom-channel-strength-spectrum}
在结论 \ref{con:xi-terminal-strictification-maximal-cp-visible-quotient} 的完全正实现下，端点 \(-1\) 的原子
\(
\mathbf{M}_{-1}:=\boldsymbol\mu_m(\{-1\})\in(\mathcal{A}_m)_+
\)
存在且为正元。
若 \(G^{\mathrm{BC}}_m\) 为阿贝尔群，则 \(\mathcal{A}_m\) 交换从而 \(\mathbf{M}_{-1}\) 属于中心；
更一般地，若 \(\boldsymbol\mu_m\) 取值于中心子代数，则对任意不可约分量 \(\pi\) 有
\[
(\mathbf{M}_{-1})_\pi=\lambda_\pi\,I_{d_\pi},
\qquad
\lambda_\pi\ge 0.
\]
族 \(\{\lambda_\pi\}\) 给出端点阈值承载的有限维“通道强度分解”，把极端共振边界层强度压缩为不可约表示索引下的一组非负谱数据。
\end{conclusion}
\begin{proof}[证明]
见定理 \ref{thm:xi-endpoint-atom-central-channel-strength}。证毕。
\end{proof}

\begin{conclusion}[端点热核探针的严格恒等式：有限 Toeplitz 二次型单调抽取端点原子质量]\label{con:xi-terminal-endpoint-heat-kernel-probe-identity}
令 \(\mathscr{C}(w)=c_0+2\sum_{n\ge 1}c_n w^n\) 在 \(w=0\) 邻域解析，并约定 \(c_{-n}:=\overline{c_n}\)。
对每个 \(N\ge 0\) 定义 Toeplitz 主子矩阵 \(\mathbf{T}_N:=(c_{j-k})_{0\le j,k\le N}\)，并令二项式向量 \(b^{(N)}\in\RR^{N+1}\) 满足
\[
b^{(N)}_k:=(-1)^k\binom{N}{k}\qquad(0\le k\le N),
\qquad
a_N:=\frac{(b^{(N)})^{\mathsf T}\mathbf{T}_N b^{(N)}}{4^N}.
\]
若 \(\mathscr{C}\) 属于 Carath\'eodory 类，则存在唯一有限正测度 \(\mu\) 使 \(c_n=\int_{\TT}\xi^{-n}\,\dd\mu(\xi)\)，并且成立严格恒等式
\[
a_N=\int_{\TT}\left|\frac{1-\xi}{2}\right|^{2N}\,\dd\mu(\xi)\ge 0,
\qquad
a_N\downarrow \mu(\{-1\})\quad(N\to\infty).
\]
因此端点阈值承载 \(\mu(\{-1\})\) 可由有限 Toeplitz 数据以单调方式无歧义提取。
\end{conclusion}
\begin{proof}[证明]
见定理 \ref{thm:xi-endpoint-heat-kernel-probe}。证毕。
\end{proof}

\begin{conclusion}[端点原子提取的指数型误差界：支撑受控时达到指数逼近]\label{con:xi-terminal-endpoint-atom-exponential-error}
在结论 \ref{con:xi-terminal-endpoint-heat-kernel-probe-identity} 的设定下，取任意闭集 \(K\subset\TT\setminus\{-1\}\)，并令
\[
q(K):=\sup_{\xi\in K}\left|\frac{1-\xi}{2}\right|\in(0,1).
\]
则对所有 \(N\ge 0\) 有
\[
0\le a_N-\mu(\{-1\})\le q(K)^{2N}\mu(K)+\mu\!\bigl(\TT\setminus(\{-1\}\cup K)\bigr).
\]
特别地，若 \(\mu(\TT\setminus(\{-1\}\cup K))=0\)，则
\(
0\le a_N-\mu(\{-1\})\le q(K)^{2N}\mu(K)
\)，从而在端点邻域之外存在谱隙（或支撑受控）时，端点原子质量由 \(\{a_N\}\) 以指数速度逼近提取；其中 \(N\) 充当严格的端点分辨率深度。
\end{conclusion}
\begin{proof}[证明]
见推论 \ref{cor:xi-endpoint-heat-kernel-exp-error}。证毕。
\end{proof}

\begin{conclusion}[阈值强度与端点谱维数：余项衰减指数精确等价于端点局部幂律指数]\label{con:xi-terminal-endpoint-threshold-spectral-dimension}
将表示测度分解为
\[
\mu=\mathfrak m_{-1}\delta_{-1}+\mu_{\neq -1},
\qquad
\mathfrak m_{-1}:=\mu(\{-1\})\ge 0.
\]
则对每个 \(N\ge 0\) 有严格分离恒等式
\[
a_N=\mathfrak m_{-1}+\int_{\TT\setminus\{-1\}}\left|\frac{1-\xi}{2}\right|^{2N}\,\dd\mu_{\neq-1}(\xi),
\qquad
a_N-\mathfrak m_{-1}\downarrow 0,
\]
从而 \(\mathfrak m_{-1}\) 为端点阈值承载的唯一零阶强度参数。进一步，若 \(\mu_{\neq -1}\) 在端点角窗满足局部幂律：存在 \(\alpha>-1\) 与常数 \(K_\pm\ge 0\) 使当 \(\varepsilon\downarrow 0\) 时
\[
\mu_{\neq -1}\bigl(\{\pi<\arg\xi<\pi+\varepsilon\}\bigr)=K_+\,\varepsilon^{\alpha+1}+o(\varepsilon^{\alpha+1}),
\qquad
\mu_{\neq -1}\bigl(\{\pi-\varepsilon<\arg\xi<\pi\}\bigr)=K_-\,\varepsilon^{\alpha+1}+o(\varepsilon^{\alpha+1}),
\]
则余项具有幂律渐近
\[
a_N-\mathfrak m_{-1}\sim (\alpha+1)\,2^{\alpha}\,
\Gamma\!\left(\frac{\alpha+1}{2}\right)\,
(K_++K_-)\,N^{-(\alpha+1)/2}\qquad (N\to\infty).
\]
尤其当 \(\mu\) 在 \(\xi=-1\) 邻域绝对连续且角密度 \(f(\theta)\) 于 \(\theta=\pi\) 连续且 \(f(\pi)>0\) 时，\(\alpha=0\) 且
\[
a_N-\mathfrak m_{-1}\sim 2f(\pi)\sqrt{\frac{\pi}{N}}.
\]
故端点局部谱维数指数（定义 \ref{def:xi-endpoint-spectral-dimension}）可由 \(a_N-\mathfrak m_{-1}\) 的幂律衰减率直接识别。
\end{conclusion}
\begin{proof}[证明]
见推论 \ref{cor:xi-endpoint-atom-separation} 与定理 \ref{thm:xi-endpoint-spectral-dimension-laplace}。证毕。
\end{proof}

\begin{conclusion}[Adams--Toeplitz--PSD 派生层级：Carath\'eodory 正性在幂覆盖下闭合]\label{con:xi-terminal-adams-toeplitz-psd-hierarchy}
若 \(\mathscr{C}(w)=c_0+2\sum_{n\ge 1}c_n w^n\) 属于 Carath\'eodory 类并由正测度 \(\mu\) 表示，则对每个 \(m\ge 1\) 定义推前测度 \(\mu_m:=(\xi\mapsto\xi^m)_\ast\mu\)。其矩满足
\[
c^{(m)}_n:=\int_{\TT}\xi^{-n}\,\dd\mu_m(\xi)=c_{mn}.
\]
因此对任意 \(N,m\ge 1\)，稀疏 Toeplitz 主子矩阵族
\[
\mathbf{T}^{(m)}_N:=\bigl(c_{m(j-k)}\bigr)_{0\le j,k\le N}\succeq 0
\]
自动成立。该派生层级在不改变正性证书口径的前提下，将端点压缩与外置尺度参数 \(m\) 直接耦合，形成一族可随 \(m\) 调节分辨率的多尺度 PSD 约束。
\end{conclusion}
\begin{proof}[证明]
见命题 \ref{prop:xi-adams-pushforward-toeplitz-principal-minor-monotonicity} 与式 \eqref{eq:xi-adams-toeplitz-psd}。证毕。
\end{proof}

\begin{conclusion}[尺度交换律：幂覆盖将端点边界层缺陷搬运到任意紧子圆盘并产生高度无关的 Jensen 正缺陷]\label{con:xi-terminal-scale-exchange-power-covering-jensen}
在自对偶完成化与固定切片制度下，离切片点状缺陷的圆盘像 \(w_\rho\) 由 Cayley--Busemann 结构性压缩到端点边界层（定理 \ref{thm:xi-offcritical-quadratic-radial-compression}）。
沿幂覆盖 \(\pi_m(\omega)=\omega^m\) 的 Adams--Norm 折叠门可将该端点边界层中的点状缺陷搬运到任意预设的紧子圆盘：对任意 \(r_\star\in(0,1)\) 存在显式尺度 \(m_\star\) 使 \(|w_\rho|^{m_\star}\le r_\star\)。
并且在固定半径 \(R\in(r_\star,1)\) 处产生高度无关的 Jensen 正缺陷下界：
\[
D_{\mathrm{Nm}_{m_\star}[F_\zeta]}(R)\ \ge\ \log\!\frac{R}{r_\star}\ >\ 0.
\]
该结论给出一条制度内可审计的交换机制：二次压缩把离切片信息压入端点边界层，而幂覆盖/折叠门以尺度参数 \(m_\star\) 将其解压到固定半径并以 Jensen 缺陷显影。
\end{conclusion}
\begin{proof}[证明]
见推论 \ref{cor:xi-adams-norm-fixed-radius-jensen-certificate}（并结合 \eqref{eq:xi-adams-norm-mstar-def}--\eqref{eq:xi-adams-norm-mstar-radius}）。证毕。
\end{proof}

\begin{conclusion}[端点重整化的二元正交分解：普适 \texorpdfstring{$\zeta$}{zeta}-极 \texorpdfstring{$\oplus$}{oplus} 离散 \texorpdfstring{$\chi_4$}{chi4} 本征残余]\label{con:xi-terminal-endpoint-renormalization-zeta-chi4-orthogonal}
在完成化坐标 \(S=r+r^{-1}\) 下，最小反变纤维坐标为 \(\delta:=r-r^{-1}\)，并满足反演反变部分的最小闭包
\(
\mathcal{A}^{-}=\delta\,\ZZ[S]
\)
（备注 \ref{rem:endpoint-inv-antiinv-min-closure}）。
奇 Chebyshev--Adams 伴随族满足恒等式
\[
D_d(S,\delta)=\delta\,U_{d-1}(S/2),
\]
从而在端点分支 \(r=\pm\mathrm{i}\) 上 \(S=0\) 并有
\[
\frac{D_d(0,\delta)}{\delta}=U_{d-1}(0)=\sin(d\pi/2)=:\varepsilon(d)\in\{-1,0,1\}.
\]
该端点通道指标 \(\varepsilon(d)\) 给出模 \(4\) 的完整奇类分裂，并与非平凡 Dirichlet 角色 \(\chi_4\) 严格一致：\(\varepsilon(d)=\chi_4(d)\)（定理 \ref{thm:endpoint-four-channel-chi4}）。
进一步，反变端点观测在 \(\sigma\downarrow 0\) 处解析并满足
\[
\lim_{\sigma\downarrow 0}\mathcal O(\mathcal W_{\sigma}^{\delta}G)
=L(1,\chi_4)\,\mathcal O(G)=\frac{\pi}{4}\,\mathcal O(G),
\]
而与之正交地，偶通道的全部发散仅沿 rank-one 方向出现并由 \(\zeta(1+\sigma)\) 的极点完全控制（定理 \ref{thm:endpoint-even-finite-part-renormalization}）。
因此端点重整化后的非紧性可被压缩为“普适 \(\zeta\)-极 \(\oplus\) 离散 \(\chi_4\) 本征残余”的二元正交分解。
\end{conclusion}
\begin{proof}[证明]
Chebyshev--Adams 恒等式与 \(\varepsilon=\chi_4\) 见命题 \ref{prop:endpoint-odd-chebyshev-identity} 与定理 \ref{thm:endpoint-four-channel-chi4}；
反变通道极限见推论 \ref{cor:endpoint-odd-analytic-at-zero}；
偶通道的极点控制与残余分解见注 \ref{rem:endpoint-renormalization-universal-pole-discrete-hair}（并参见注 \ref{rem:endpoint-obstruction-dichotomy}）。
证毕。
\end{proof}

\begin{conclusion}[有限阿贝尔可见商的块 Toeplitz 对角化：单一块 PSD 证书等价于全通道 PSD]\label{con:xi-terminal-abelian-block-toeplitz-diagonalization}
令 \(\mathcal{A}\) 为有限维 \(C^\ast\)-代数，并取解析函数 \(F:\mathbb{D}\to\mathcal{A}\)。
则以下命题等价：\(\Re F(w)\succeq 0\)（算子值 Carath\'eodory）、对每个不可约表示 \(\pi\) 有 \(\Re F_\pi(w)\succeq 0\)、以及对每个 \(N\ge 0\) 的块 Toeplitz 截断 \(\mathbf{T}^{\mathcal{A}}_N(F)\succeq 0\)。
特别地，当 \(\mathcal{A}=\CC[H]\) 且 \(H\) 为有限阿贝尔群时，在 Fourier--Pontryagin 对偶分解下 \(\mathbf{T}^{\mathcal{A}}_N(F)\) 与直和矩阵 \(\bigoplus_{\chi\in\widehat H}\mathbf{T}_N(F_\chi)\) 酉等价；因此
\[
\mathbf{T}^{\mathcal{A}}_N(F)\succeq 0\ \Longleftrightarrow\ \forall\chi\in\widehat H,\ \mathbf{T}_N(F_\chi)\succeq 0.
\]
从而“逐通道 PSD 核验”可被统一为一个单一块 Toeplitz--PSD 可行域（单一半正定证书）；与结论 \ref{con:xi-terminal-adams-toeplitz-psd-hierarchy} 的 Adams--Toeplitz 派生层级结合，可在单一块证书口径下实现全尺度、全通道的正性核验。
\end{conclusion}
\begin{proof}[证明]
见定理 \ref{thm:xi-fourier-pontryagin-block-toeplitz} 与推论 \ref{cor:xi-abelian-block-toeplitz-directsum}。证毕。
\end{proof}

\begin{conclusion}[异常类与散射绕数的同一阻碍：交换失败残差与拓扑绕数在可见口径下严格对齐]\label{con:xi-terminal-anom-winding-obstruction-equivalence}
固定分辨率 \(m\) 并考虑可见扭曲通道诱导的 unitary-slice 散射行列式回路
\(\{\det\mathcal S_{\overline\chi}\}_{\overline\chi\in\widehat{G^{\mathrm{BC}}_m}}\)。
若异常在可见通道上严格为零，则每个可见通道的散射回路可取高能归一化
\(\det\mathcal S_{\overline\chi}(\infty)=1\)，并满足
\[
\mathrm{Wind}\bigl(\det\mathcal S_{\overline\chi}\bigr)=0
\qquad(\forall\,\overline\chi\in\widehat{G^{\mathrm{BC}}_m}),
\]
从而在有限型点状共振口径下所有可见通道的真共振指数归零。
反向地，若对某群商 \(G\twoheadrightarrow H\) 的全部可见通道均可实现上述归一化与零绕数，则该群商必唯一因子化通过 \(G^{\mathrm{BC}}_m\)，并在 \(H\) 上使异常严格为零。
因此“代数残差”（异常类）与“拓扑绕数”在可见口径下构成同一阻碍类。
\end{conclusion}
\begin{proof}[证明]
见定理 \ref{thm:xi-anom-winding-obstruction}（真共振指数与绕数/Blaschke 度的对齐由定理 \ref{thm:selfdual-wind-deg-dbl-kappa-identity} 给出）。证毕。
\end{proof}

\begin{conclusion}[有界切片的有限前沿原理：证书搜索与编译优化在每一分辨率层级均可退化为有限枚举]\label{con:xi-terminal-finite-frontier-principle}
在任意固定有界复杂度切片 \(\mathbf{PW}^{\ast}_{\le}\) 内，取保守语义函子 \(\mathsf{Sem}\) 与任一代价函子 \(\mathcal{C}\)。
则对任意投影词 \(w\in \mathbf{PW}^{\ast}_{\le}\)，其语义等价类 \([w]\) 在代价域中的像 \(\mathcal{C}([w])\subseteq \mathbb{T}^{(2)}\) 具有有限帕累托前沿，并可由有限枚举构造。
因此任何以“接口正规形参数—异常签名—谱指纹”为核对输出的证书目标，在每一固定分辨率层级上必然退化为有限枚举与有限核对；困难被严格集中到“切片上界随分辨率增长律”而非实现细节。
\end{conclusion}
\begin{proof}[证明]
见定理 \ref{thm:xi-finite-frontier-certificate}。证毕。
\end{proof}

\begin{conclusion}[像空间在 \texorpdfstring{$m\ge 3$}{m>=3} 时失去分辨率可辨识性：\texorpdfstring{$Y_m$}{Ym} 的共轭塌缩]\label{con:xi-terminal-image-space-resolution-nonidentifiability}
沿定义 \ref{def:Y_m} 记
\[
Y_m:=\Phi_m\bigl(\{0,1\}^{\ZZ}\bigr).
\]
则对任意 \(m,n\ge 3\)，移位系统 \((Y_m,\sigma)\) 与 \((Y_n,\sigma)\) 拓扑共轭；更强地，它们都与二元全移位 \((\{0,1\}^{\ZZ},\sigma)\) 共轭。因而对每个 \(m\ge 3\) 与 \(k\ge 1\) 有
\[
h_{\mathrm{top}}(Y_m)=\log 2,\qquad
\zeta_{Y_m}(z)=\frac{1}{1-2z},\qquad
\#\mathrm{Fix}(\sigma^k|_{Y_m})=2^k.
\]
特别地，任何仅依赖 \(Y_m\) 的拓扑共轭不变量都不可能恢复分辨率 \(m\)；非平凡的 prime shadow 或 arithmetic shadow 必须来自核、势/权重、扭曲通道或附加约束，而不能由像空间 \(Y_m\) 单独承载。
\end{conclusion}
\begin{proof}[证明]
定理 \ref{thm:Phi_m-conjugacy-threshold} 已给出当 \(m\ge 3\) 时 \(\Phi_m\) 与二元全移位之间的拓扑共轭，推论 \ref{cor:Ym-conjugacy-all-mge3} 给出任意 \(Y_m,Y_n\) 之间的共轭。由推论 \ref{cor:Phi_m-zeta-closed}（等价地，推论 \ref{cor:Ym-periodic-rigidity}）可得固定点计数与内禀 Artin--Mazur \(\zeta\)-函数的闭式。拓扑共轭不变量在共轭下保持不变，而推论 \ref{cor:Phi_m-zeta-closed} 已明确非平凡 prime shadow 不来自折叠因子本身，故结论成立。证毕。
\end{proof}

\begin{conclusion}[Zeckendorf 合法性不是微扰：加法投影碰撞谱主指数发生严格指数级分离]\label{con:xi-terminal-addition-collision-spectrum-constant-factor-drop}
把“加法作为一次投影”的 \(q\)-阶碰撞 Perron 根分别记为
\[
r_q^{\oplus,\mathrm{sync10}},\qquad
r_q^{\oplus,\mathrm{real}},
\]
其中前者对应 \(10\) 状态归一化核与全移位输入，后者对应真实输入 \(40\) 状态核与 Zeckendorf 合法性约束。则对所有 \(q\ge 2\) 有严格不等式
\[
r_q^{\oplus,\mathrm{real}}<r_q^{\oplus,\mathrm{sync10}}.
\]
在 \(q=2\) 时，
\[
r_2^{\oplus,\mathrm{sync10}}\approx 8.7430515270,\qquad
r_2^{\oplus,\mathrm{real}}\approx 3.9985433445,\qquad
\frac{r_2^{\oplus,\mathrm{sync10}}}{r_2^{\oplus,\mathrm{real}}}\approx 2.186559.
\]
若 \(S_q^{\oplus,\mathrm{sync10}}(m)\) 与 \(S_q^{\oplus,\mathrm{real}}(m)\) 分别表示两种口径下的 \(q\)-阶碰撞矩，则
\[
S_q^{\oplus,\mathrm{sync10}}(m)=\Theta\!\bigl((r_q^{\oplus,\mathrm{sync10}})^m\bigr),\qquad
S_q^{\oplus,\mathrm{real}}(m)=\Theta\!\bigl((r_q^{\oplus,\mathrm{real}})^m\bigr),
\]
并存在常数 \(c_q,C_q>0\) 使得对充分大的 \(m\) 有
\[
c_q\left(\frac{r_q^{\oplus,\mathrm{sync10}}}{r_q^{\oplus,\mathrm{real}}}\right)^m
\le
\frac{S_q^{\oplus,\mathrm{sync10}}(m)}{S_q^{\oplus,\mathrm{real}}(m)}
\le
C_q\left(\frac{r_q^{\oplus,\mathrm{sync10}}}{r_q^{\oplus,\mathrm{real}}}\right)^m.
\]
因此，把真实输入替换为全移位输入所造成的误差不是多项式级或次指数级，而是由 Perron 根比值控制的严格指数级形变。
\end{conclusion}
\begin{proof}[证明]
对所有 \(q\ge 2\) 的严格不等式由定理 \ref{thm:addition-collision-perron-drop-real-vs-sync} 给出。附录 \ref{app:add-collision-spectrum} 已将两类 \(q\)-阶碰撞矩编译为有限非负核，并给出 Perron 主尺度表示 \(S_q^{\oplus}(m)=\Theta((r_q^{\oplus})^m)\)；故上述比值估计由标准 Perron--Frobenius 渐近立即得到。\(q=2\) 的显式数值由命题 \ref{prop:addition-collision-spectrum-sync10-vs-real40} 及表 \ref{tab:add-collision-spectrum-10-vs-40} 给出。证毕。
\end{proof}

\begin{conclusion}[像空间塌缩之后的核层刚性：真实输入 essential 核与二元全移位在流等价与 \texorpdfstring{$K$}{K}-理论上严格分离]\label{con:xi-terminal-kernel-flow-k-separation-from-fullshift}
令 \(A=M_{\mathrm{ess}}\) 为真实输入 \(20\times 20\) essential 邻接矩阵。则
\[
\mathrm{BF}(A)\cong \ZZ^2\oplus \ZZ/2\ZZ,\qquad
K_0(\mathcal O_A)\cong \ZZ^2\oplus \ZZ/2\ZZ,\qquad
K_1(\mathcal O_A)\cong \ZZ^2.
\]
另一方面，对任意 \(m\ge 3\)，像空间 \(Y_m\) 仅为二元全移位的一个共轭模型。因而：
\begin{enumerate}
  \item \(A\) 不可能与二元全移位的任一 irreducible SFT 表示流等价；
  \item \(\mathcal O_A\) 不可能与 \(\mathcal O_2\) 稳定同构；
  \item 本文中出现的非平凡 prime-orbit、\(\zeta\) 或算子代数结构，不是 \(Y_m\) 的内禀像空间影子，而是核层级的独立刚性。
\end{enumerate}
\end{conclusion}
\begin{proof}[证明]
命题 \ref{prop:real-input-40-bf-ktheory} 已给出 \(A\) 的 Bowen--Franks 群与 Cuntz--Krieger \(K\)-群。对二元全移位的标准邻接矩阵 \([2]\)，有
\[
\mathrm{BF}([2])=\ZZ/(1-2)\ZZ=0,\qquad
K_0(\mathcal O_2)=0,\qquad
K_1(\mathcal O_2)=0.
\]
由于 Bowen--Franks 群在 irreducible SFT 的流等价下保持不变，而 Cuntz--Krieger \(K\)-群在稳定同构下保持不变，故前两条成立。再结合结论 \ref{con:xi-terminal-image-space-resolution-nonidentifiability} 中 \(Y_m\) 的全移位塌缩，第三条随之成立。证毕。
\end{proof}

\begin{conclusion}[高阶碰撞矩的模 \texorpdfstring{$2$}{2} 阴影在 \texorpdfstring{$q=9,\dots,17$}{q=9,...,17} 上完全有限自动机化]\label{con:xi-terminal-high-order-collision-mod2-automaticity}
对共振窗 \(q=9,\dots,17\)，模 \(2\) 影子分解满足
\[
E^{a_q}(E+1)^{b_q}S_q\equiv 0\pmod 2,
\]
其中表 \ref{tab:fold_collision_charpoly_mod2_shadow_q9_17} 给出
\[
b_q=
\begin{cases}
4, & q=9,10,\\
6, & q=11,12,14,16,\\
8, & q=13,15,\\
10, & q=17.
\end{cases}
\]
因此各序列 \(S_q(m)\bmod 2\) 的最终周期满足
\[
\mathrm{Per}_q\mid 4\quad(q=9,10),\qquad
\mathrm{Per}_q\mid 8\quad(q=11,\dots,16),\qquad
\mathrm{Per}_q\mid 16\quad(q=17).
\]
进一步，定义向量值序列
\[
\mathbf S(m):=\bigl(S_9(m),S_{10}(m),\dots,S_{17}(m)\bigr)\pmod 2.
\]
则存在 \(M_0\) 使得对所有 \(m\ge M_0\) 有
\[
\mathbf S(m+16)=\mathbf S(m).
\]
等价地，\(\mathbf S(m)\) 在最终阶段仅依赖于 \(m\) 的低 \(4\) 位二进制展开；因此整个高阶区间的模 \(2\) 阴影可由一个周期长度至多为 \(16\) 的有限自动机完全生成。
\end{conclusion}
\begin{proof}[证明]
定理 \ref{thm:pom-resonance-mod2-finite-bit-automaticity} 给出：对每个 \(q\)，最终周期整除 \(2^{\lceil\log_2 b_q\rceil}\)。把表 \ref{tab:fold_collision_charpoly_mod2_shadow_q9_17} 中的 \(b_q\) 逐项代入，即得上述各区间的周期上界。向量序列 \(\mathbf S(m)\) 的最终周期于是整除
\[
\mathrm{lcm}(4,8,16)=16.
\]
故存在 \(M_0\) 使 \(\mathbf S(m+16)=\mathbf S(m)\) 对一切 \(m\ge M_0\) 成立；而“最终周期整除 \(16\)”与“最终仅依赖于 \(m\) 的低 \(4\) 位二进制展开”等价。证毕。
\end{proof}

\begin{conclusion}[零温端点的大偏差速率函数具有由 \texorpdfstring{$B_\infty$}{B∞} 控制的对数尖点]\label{con:xi-terminal-real-input-output-rate-log-cusp}
沿用命题 \ref{prop:real-input-40-pressure-freezing} 的输出位势记号，并令
\[
P(\theta)=\frac{\theta}{2}+\log c+A e^{-\theta/2}+B e^{-\theta}+O(e^{-3\theta/2})
\qquad (\theta\to+\infty),
\]
其中
\[
A:=\frac{d}{c}>0,
\qquad
c=\sqrt{\rho(B_\infty)}.
\]
定义未归一化 Legendre--Fenchel 对偶
\[
\widetilde I(s):=\sup_{\theta\in\RR}\bigl(s\theta-P(\theta)\bigr),
\qquad 0\le s\le \frac12.
\]
当 \(s\uparrow \frac12\) 时，令
\[
\delta:=\frac12-s\downarrow 0.
\]
则有精确边界展开
\[
\widetilde I(s)
=
-\log c
+2\delta\log\delta
+2\Bigl(\log\frac{2}{A}-1\Bigr)\delta
+O(\delta^2).
\]
特别地，
\[
\widetilde I\!\left(\frac12\right)
=-\log c
=-\frac12\log\rho(B_\infty).
\]
若改用附录 \ref{app:real-input-40-zeta-u} 中的归一化速率函数
\[
I_e(s):=\sup_{\theta\in\RR}\bigl(s\theta-(P(\theta)-P(0))\bigr)
=\widetilde I(s)+P(0),
\]
则同样有
\[
I_e(s)
=
P(0)-\log c
+2\delta\log\delta
+2\Bigl(\log\frac{2}{A}-1\Bigr)\delta
+O(\delta^2),
\]
从而
\[
I_e\!\left(\frac12\right)=P(0)-\log c=\log(\varphi^2/c).
\]
因此饱和端点 \(s=\frac12\) 不是解析边界点，而是首系数为 \(2\) 的对数尖点；其未归一化高度完全由零温地基矩阵 \(B_\infty\) 的 Perron 根控制。
\end{conclusion}
\begin{proof}[证明]
由命题 \ref{prop:real-input-40-pressure-freezing}，当 \(\theta\to+\infty\) 时
\[
P'(\theta)=\frac12-\frac{A}{2}e^{-\theta/2}+O(e^{-\theta}),
\qquad
P''(\theta)=\frac{A}{4}e^{-\theta/2}+O(e^{-\theta})>0
\]
对充分大的 \(\theta\) 成立。因此当 \(s<\frac12\) 且充分接近 \(\frac12\) 时，极值点 \(\theta=\theta(s)\) 唯一，并由方程
\[
P'(\theta)=s
\]
确定。令 \(\delta=\frac12-s\)，则
\[
\delta=\frac{A}{2}e^{-\theta/2}+O(e^{-\theta}),
\]
从而反解得
\[
e^{-\theta/2}=\frac{2\delta}{A}+O(\delta^2),
\qquad
\theta=-2\log\delta+2\log\frac{A}{2}+O(\delta).
\]
将其代回
\[
\widetilde I(s)=s\theta-P(\theta)
=-\delta\theta-\log c-Ae^{-\theta/2}+O(e^{-\theta}),
\]
得到
\[
\widetilde I(s)
=
-\log c
+2\delta\log\delta
+2\Bigl(\log\frac{2}{A}-1\Bigr)\delta
+O(\delta^2).
\]
令 \(\delta\downarrow 0\) 即得
\(
\widetilde I(\frac12)=-\log c
\)。
同一附录中命题 \ref{prop:real-input-40-pressure-freezing} 之后的零温缩放证书已识别 \(c=\sqrt{\rho(B_\infty)}\)，故
\(
\widetilde I(\frac12)=-\frac12\log\rho(B_\infty)
\)。
归一化速率函数 \(I_e\) 与 \(\widetilde I\) 仅相差常数 \(P(0)\)，故展开的一次非解析项与线性系数保持不变；端点值
\(
I_e(\frac12)=P(0)-\log c
\)
则与推论 \ref{cor:real-input-40-rate-endpoints} 一致。证毕。
\end{proof}

\begin{conclusion}[饱和边界微正则熵的二阶对数尖点展开]\label{con:xi-terminal-real-input-output-microcanonical-logcusp-second-order}
沿用命题 \ref{prop:real-input-40-pressure-freezing} 的记号，定义微正则熵
\[
s(\alpha):=\inf_{\theta\in\RR}\bigl(P(\theta)-\alpha\theta\bigr),
\qquad
0<\alpha\le \frac12.
\]
再记
\[
a:=\frac{d}{c},
\qquad
b:=\frac{e}{c}-\frac{d^2}{2c^2},
\qquad
\delta:=\frac12-\alpha\downarrow 0.
\]
则当 \(\alpha\uparrow \frac12\) 时，
\[
s\!\left(\frac12-\delta\right)
=
\log c
-2\delta\log\delta
+2\left(1+\log\frac{a}{2}\right)\delta
+\frac{4b}{a^2}\delta^2
+O(\delta^3).
\]
数值上，
\[
s\!\left(\frac12-\delta\right)
=
\log c
-2\delta\log\delta
-1.0926862396530588\,\delta
-4.850668312338638\,\delta^2
+O(\delta^3).
\]
因此，饱和边界并非平滑触及地基熵 \(\log c\)，而是以首系数为 \(2\) 的对数尖点方式逼近。
\end{conclusion}
\begin{proof}[证明]
由命题 \ref{prop:real-input-40-pressure-freezing}，
\[
P(\theta)
=
\frac{\theta}{2}+\log c+a e^{-\theta/2}+b e^{-\theta}+O(e^{-3\theta/2}),
\]
从而
\[
P'(\theta)
=
\frac12-\frac{a}{2}e^{-\theta/2}-b e^{-\theta}+O(e^{-3\theta/2}).
\]
令
\[
x:=e^{-\theta/2}.
\]
则极小点条件 \(P'(\theta)=\alpha=\frac12-\delta\) 等价于
\[
\delta=\frac{a}{2}x+b x^2+O(x^3).
\]
按幂级数反演得到
\[
x=\frac{2\delta}{a}-\frac{8b}{a^3}\delta^2+O(\delta^3).
\]
另一方面，
\[
s(\alpha)=P(\theta)-\alpha\theta
=
\log c+a x+b x^2-2\delta\log x+O(x^3).
\]
再由
\[
\log x
=
\log\frac{2\delta}{a}
-\frac{4b}{a^2}\delta
+O(\delta^2),
\]
代回并整理得
\[
s\!\left(\frac12-\delta\right)
=
\log c
-2\delta\log\delta
+2\left(1+\log\frac{a}{2}\right)\delta
+\frac{4b}{a^2}\delta^2
+O(\delta^3).
\]
最后将命题 \ref{prop:real-input-40-pressure-freezing} 中给出的
\(
c\approx 1.89635299351377
\),
\(
d\approx 0.807943317323508
\),
\(
e\approx -0.245317935748073
\)
代入，即得数值式。证毕。
\end{proof}


\begin{theorem}[输出位势右端点的正熵冻结端面由单一四次 Perron 根生成]\label{thm:xi-terminal-real-input-positive-entropy-freezing-face}
\leanverified{paper\_xi\_terminal\_real\_input\_positive\_entropy\_freezing\_face}
沿用命题 \ref{prop:real-input-40-pressure-freezing} 的输出位势记号，并定义归一化大偏差速率函数
\[
I(\alpha):=\sup_{\theta\in\RR}\bigl(\theta\alpha-(P(\theta)-P(0))\bigr).
\]
令 \(B_\infty\) 为附录 \ref{app:real-input-40-zeta-u} 中零温缩放极限所得到的 \(4\times 4\) 非负整数矩阵，并记
\[
y:=\rho(B_\infty).
\]
则有精确闭式
\[
\boxed{\;
y^4-6y^3+9y^2-y-1=0,
\qquad
c=\sqrt y.\;}
\]
并且
\[
\boxed{\;
\alpha_{\max}=\frac12,
\qquad
h_{\mathrm{ground}}=\log c=\frac12\log y,\;}
\]
\[
\boxed{\;
I\!\left(\frac12\right)=P(0)-\log c
=
\log\!\left(\frac{\varphi^2}{\sqrt y}\right)>0.\;}
\]
因此，输出位势的右端点并非零熵冻结点，而是一条正熵冻结端面：端点密度冻结到 \(\alpha_{\max}=1/2\)，但端点微观实现族仍以指数率 \(\exp(n h_{\mathrm{ground}})\) 增长，同时达到该端点仍需严格正的大偏差代价。
\end{theorem}
\begin{proof}
附录 \ref{app:real-input-40-zeta-u} 中的零温缩放极限给出
\[
\widehat F(w)=1-6w+9w^2-w^3-w^4=\det(I-wB_\infty).
\]
故 \(y=\rho(B_\infty)\) 满足
\[
y^4-6y^3+9y^2-y-1=0,
\]
并且同一缩放识别给出
\[
c=\sqrt{\rho(B_\infty)}=\sqrt y.
\]

另一方面，推论 \ref{cor:real-input-40-alpha-max} 已证明
\[
\alpha_{\max}=\lim_{\theta\to+\infty}P'(\theta)=\frac12.
\]
再由推论 \ref{cor:real-input-40-ground-entropy}，
\[
h_{\mathrm{ground}}=\log c=\frac12\log y.
\]

最后，推论 \ref{cor:real-input-40-rate-endpoints} 给出
\[
I\!\left(\frac12\right)=P(0)-\log c.
\]
由于
\[
P(0)=\log\lambda(1)=\log(\varphi^2),
\]
故
\[
I\!\left(\frac12\right)=\log\!\left(\frac{\varphi^2}{\sqrt y}\right).
\]
同一推论已给出其数值约为 \(0.322491085601762>0\)，故上式确为严格正值。
综上，右端点同时呈现密度冻结、正熵地基与正的大偏差代价，且三项端点数据均由同一四次 Perron 根 \(y\) 生成。证毕。
\end{proof}

\endinput


\begin{conclusion}[折叠输出的内生 Gibbs 规格与 Parry 基线几乎同熵：最大熵缺口降至微比特级且 KL 率等同于熵缺口]\label{con:xi-terminal-fold-output-gibbs-parry-microbit-gap}
对 Fold 输出分布 \(\pi_m(x)=d_m(x)/2^m\) 的块频率拟合得到两状态 Markov 规格 \(P_{\mathrm{fit}}\)，其一步转移与黄金均值 Parry 基线仅在 \(10^{-3}\) 量级偏离（例如 \(P_{\mathrm{fit}}(1\mid 0)=0.3829949\) 对比 \(P_{\mathrm{Parry}}(1\mid 0)=0.3819660\)）。
更关键的是熵率层面的微缺口：
\[
h_{\mathrm{fit}}\approx 0.481210204266,\quad
\log\varphi\approx 0.481211825060,\quad
\Delta h:=\log\varphi-h_{\mathrm{fit}}\approx 1.620794\times 10^{-6}\ \text{nats/symbol}.
\]
并且 KL 散度率满足严格一致性恒等式
\[
D_{\mathrm{KL}}^{\mathrm{rate}}(P_{\mathrm{fit}}\|P_{\mathrm{Parry}})=\log\varphi-h_{\mathrm{fit}}=\Delta h,
\]
从而折叠诱导的“偏离”主要体现为常数级指纹（有限部/端点层），而非熵率层面的结构崩塌。
\end{conclusion}
\begin{proof}[证明]
见注 \ref{rem:pom-gibbs-markov-fit} 与注 \ref{rem:pom-gibbs-parry-entropy-gap}（其中数值证书由 \texttt{scripts/exp\_fold\_output\_gibbs\_markov\_fit.py} 生成并以 \texttt{\string\input} 方式纳入正文）。证毕。
\end{proof}

\begin{conclusion}[端点质量给出黄金常数的过定约束审计：两条独立比值估计在 Parry 模型下必同收敛于 \texorpdfstring{$\varphi$}{phi}]\label{con:xi-terminal-parry-endpoint-mass-overdetermined-phi-audit}
以任意经验直方图 \(\hat\pi_{m,N}\) 定义端点质量
\[
\widehat M_{ij}(m):=\sum_{\substack{u\in X_m\\ u_1=i,\ u_m=j}}\hat\pi_{m,N}(u).
\]
若数据在分辨率 \(m\) 下接近黄金均值 Parry 基线，则可构造两条相互独立的黄金常数估计量：
\[
\widehat\varphi_{(00/01)}(m):=\frac{\widehat M_{00}(m)}{\widehat M_{01}(m)}\cdot \frac{F_{m-1}}{F_m},
\qquad
\widehat\varphi_{(01/11)}(m):=\frac{\widehat M_{01}(m)}{\widehat M_{11}(m)}\cdot \frac{F_{m-2}}{F_{m-1}}.
\]
在 Parry 真模型下，两者在极限意义下收敛到同一常数 \(\varphi\)，从而形成一个低维、过定约束的自洽审计机制：任何偏离 Parry 的端点结构将使两条比值估计产生可检测的不一致。
\end{conclusion}
\begin{proof}[证明]
见推论 \ref{cor:parry-endpoint-phi-estimators}。证毕。
\end{proof}

\begin{conclusion}[模群判据下的贝叶斯逆与零知识因子化：完全正信道的“可见闭包不变性”给出充要刻画]\label{con:xi-terminal-bayesian-inverse-zk-factorization-invariance}
在算子代数设定中，贝叶斯逆 \(F^\sharp\) 的存在性由 Tomita--Takesaki 模群判据控制。
一旦 \(F^\sharp\) 存在，则对任意正规完全正转录信道 \(\Phi:\mathcal{M}\to\mathcal{T}\)，以下命题等价：
\begin{enumerate}
  \item 存在正规完全正映射 \(S:\mathcal{N}\to\mathcal{T}\) 使 \(\Phi=S\circ F^\sharp\)（相对于 \(\mathcal{N}\) 的完美零知识因子化）；
  \item \(\Phi\) 在可见闭包幂等子 \(K:=F\circ F^\sharp\) 下不变：\(\Phi=\Phi\circ K\)。
\end{enumerate}
并且在等价条件下分解 \(\Phi=S\circ F^\sharp\) 中的 \(S\) 唯一且满足 \(S=\Phi|_{\mathcal{N}}\)。
因此“零知识/可因子化”从构造性问题降格为单一不变性检验。
\end{conclusion}
\begin{proof}[证明]
见定理 \ref{thm:op-algebra-bayes-inverse-zk}。证毕。
\end{proof}

% (User input window)

\paragraph{跨接口终端结论：负携带势压力的低阶 jet、尺度叠加完成化的指数型刚性与单尺度 Jensen 热带剖面}
本段补充若干“可直接核验的常数/障碍”条目：其一把负携带势的压力函数在原点的低阶累积量写成显式有理常数，
并给出相应大偏差率函数在端点处的闭式与边界渐近；其二把尺度叠加完成化在 Hardy 切片上的指数型锁定为
“带宽 \(\times\) 支撑直径”，并由 Cartwright 密度估计将 \(T\log T\) 零点统计的目标转译为一个不可绕过的复杂度乘积下界；
其三对单尺度完成化给出 Jensen--热带剖面的分段仿射闭式，并把“离线谱测度”在实部方向严格原子化；
最后记录 \(\varphi^{\mathbb Q}\) 尺度混合的共同虚周期障碍。

\begin{conclusion}[负携带势压力函数的低阶累积量刚性：\texorpdfstring{$P'(0),P''(0),P^{(3)}(0)$}{P'(0),P''(0),P^{(3)}(0)} 为显式有理常数]\label{con:xi-terminal-phi-minus-pressure-loworder-cumulants}
在附录 \ref{app:vector-potential} 的负携带势实例中，令 \(u=e^{\theta}>0\)，并令 \(\lambda(u)\) 为三次方程
\[
\lambda^{3}-(u+2)\lambda^{2}+(u-2)\lambda+3u=0
\]
在 \(u>0\) 上的 Perron 分支实根。
定义压力
\(
P(\theta):=\log\lambda(e^{\theta})
\)。
则 \(P\) 在 \(\theta=0\) 处的前三阶导数为
\[
P'(0)=\frac18,\qquad
P''(0)=\frac{25}{128},\qquad
P^{(3)}(0)=\frac{177}{512}.
\]
\end{conclusion}
\begin{proof}[证明]
令 \(F(\lambda,u):=\lambda^{3}-(u+2)\lambda^{2}+(u-2)\lambda+3u\)。
当 \(u=1\) 时
\(
F(\lambda,1)=(\lambda-3)(\lambda^{2}-1)
\)，故 Perron 根为 \(\lambda(1)=3\)，从而 \(P(0)=\log 3\)。
\par
由隐函数求导
\[
\lambda_u=-\frac{F_u}{F_\lambda},
\qquad
\lambda_{uu}=-\frac{F_{uu}+2F_{u\lambda}\lambda_u+F_{\lambda\lambda}\lambda_u^{2}}{F_\lambda},
\]
并对恒等式
\(
F_\lambda\lambda_{uu}+F_{\lambda\lambda}\lambda_u^{2}+2F_{u\lambda}\lambda_u=0
\)
再微分一次即可得到 \(\lambda_{uuu}\) 的表达式（均在 \(F_\lambda\neq 0\) 处成立）。
将偏导
\[
F_u=-(\lambda^{2}-\lambda-3),\quad
F_\lambda=3\lambda^{2}-2(u+2)\lambda+(u-2),
\]
\[
F_{u\lambda}=-2\lambda+1,\quad
F_{\lambda\lambda}=6\lambda-2(u+2),\quad
F_{uu}=0
\]
代入，并在 \((u,\lambda)=(1,3)\) 处化简得
\(
\lambda_u(1)=3/8
\),
\(
\lambda_{uu}(1)=33/128
\),
\(
\lambda_{uuu}(1)=117/1024
\)。
\par
又因 \(u=e^{\theta}\)，由链式法则
\[
\frac{d\lambda}{d\theta}=u\lambda_u,\quad
\frac{d^{2}\lambda}{d\theta^{2}}=u\lambda_u+u^{2}\lambda_{uu},\quad
\frac{d^{3}\lambda}{d\theta^{3}}=u\lambda_u+3u^{2}\lambda_{uu}+u^{3}\lambda_{uuu},
\]
并用 \(\,P=\log\lambda\) 的恒等式
\[
P'=\frac{\lambda'}{\lambda},\qquad
P''=\frac{\lambda''}{\lambda}-\Bigl(\frac{\lambda'}{\lambda}\Bigr)^{2},\qquad
P^{(3)}=\frac{\lambda^{(3)}}{\lambda}-3\frac{\lambda''\lambda'}{\lambda^{2}}+2\Bigl(\frac{\lambda'}{\lambda}\Bigr)^{3}
\]
（此处撇号表示对 \(\theta\) 求导），在 \(\theta=0\)（即 \(u=1,\lambda=3\)）处代入即可得到所示三值。证毕。
\end{proof}

\begin{conclusion}[负携带势大偏差速率函数的端点常数与边界渐近：\texorpdfstring{$I(0),I(1)$}{I(0),I(1)} 闭式与两端对数项]\label{con:xi-terminal-phi-minus-rate-endpoints-boundary-asymptotics}
沿用结论 \ref{con:xi-terminal-phi-minus-pressure-loworder-cumulants} 的 \(P(\theta)\)，并令 \(P(0)=\log 3\)。
定义速率函数
\[
I(\alpha):=\sup_{\theta\in\RR}\bigl(\theta\alpha-(P(\theta)-P(0))\bigr).
\]
则 \(I\) 在 \([0,1]\) 上可连续延拓，且满足：
\begin{enumerate}
  \item \(I\) 在唯一点 \(\alpha_{0}=P'(0)=\frac18\) 处取到全局最小值 \(I(\alpha_{0})=0\)；
  \item 端点取值为
  \[
  I(0)=\log\Bigl(\frac{3}{1+\sqrt3}\Bigr),\qquad
  I(1)=\log 3;
  \]
  \item 边界渐近为
  \[
  I(\alpha)=\log\Bigl(\frac{3}{1+\sqrt3}\Bigr)+\alpha\log\alpha+\bigl(\log(8+4\sqrt3)-1\bigr)\alpha+o(\alpha)
  \qquad (\alpha\downarrow 0),
  \]
  \[
  I(\alpha)=\log 3+(1-\alpha)\log(1-\alpha)-(1-\alpha)+o(1-\alpha)
  \qquad (\alpha\uparrow 1).
  \]
\end{enumerate}
\end{conclusion}
\begin{proof}[证明]
连续延拓与唯一极小点是 Legendre--Fenchel 对偶的标准性质：由结论 \ref{con:xi-terminal-phi-minus-pressure-loworder-cumulants} 的 \(P\) 解析且严格凸（附录 \ref{app:vector-potential} 的 Perron 分支与命题 \ref{prop:vector-potential-phi-minus-ldp-param}）知，
极值点满足 \(\alpha=P'(\theta)\)，并有 \(I(\alpha)=\theta\alpha-(P(\theta)-P(0))\)；
取 \(\theta=0\) 得 \(I(\alpha_{0})=0\)，且严格凸性给出唯一性。
\par
端点与边界展开可在命题 \ref{prop:vector-potential-phi-minus-ldp-param} 的参数化下计算。
该命题给出对 \(\lambda\in(1+\sqrt3,\infty)\) 的显式表示
\[
u(\lambda)=\frac{\lambda(\lambda^{2}-2\lambda-2)}{\lambda^{2}-\lambda-3},
\qquad
\alpha(\lambda)=\frac{\lambda^{4}-3\lambda^{3}-3\lambda^{2}+8\lambda+6}{\lambda^{4}-2\lambda^{3}-5\lambda^{2}+12\lambda+6},
\]
以及
\(
I(\alpha(\lambda))=\alpha(\lambda)\log u(\lambda)-\log\lambda+\log 3
\)。
当 \(\lambda\downarrow 1+\sqrt3\) 时有 \(u(\lambda)\downarrow 0\)、\(\alpha(\lambda)\downarrow 0\)，且
\(
\alpha(\lambda)\log u(\lambda)\to 0
\)，从而
\(
I(0)=\log 3-\log(1+\sqrt3)
\)；
当 \(\lambda\to\infty\) 时 \(u(\lambda)\to\infty\)、\(\alpha(\lambda)\to 1\) 且 \(\log u(\lambda)-\log\lambda\to 0\)，从而 \(I(1)=\log 3\)。
\par
进一步，对 \(\lambda=1+\sqrt3+\varepsilon\)（\(\varepsilon\downarrow 0\)）与 \(\lambda\to\infty\) 作局部展开并代回
\(
I=\alpha\log u-\log\lambda+\log 3
\)，即可得到两端的
\(\alpha\log\alpha\) 与 \((1-\alpha)\log(1-\alpha)\) 级主项以及所示的一次项系数。证毕。
\end{proof}

\begin{conclusion}[尺度叠加完成化的指数型精确刚性：Hardy 切片的指数型等于带宽乘支撑直径]\label{con:xi-terminal-scale-superposition-exact-type-rigidity}
沿用定义 \ref{def:xi-scale-mellin-superposition}，并假设权函数 \(w\in L^{1}([0,\infty))\) 支撑于 \([0,A]\)。
设核 \(\Delta(z,u)=\sum_{j,k}a_{j,k}z^{j}u^{k}\) 为有限和，并令
\[
m:=j-2k,\qquad
B:=\max\{|m|:\ a_{j,k}\neq 0\}.
\]
定义 Hardy 切片
\(
F(t):=\mathcal X_{\Delta,w}(\tfrac12+it)
\)。
若存在极值模态 \(|m|=B\) 使得
\[
x\longmapsto w(x)\sum_{j-2k=m}a_{j,k}e^{-jx/2}
\]
在任意邻域 \((A-\varepsilon,A]\) 上不几乎处处为零，则 \(F\)（视为整函数）满足
\[
\mathrm{type}(F)=BA.
\]
\end{conclusion}
\begin{proof}[证明]
由定义 \ref{def:xi-scale-mellin-superposition} 的代换式，在 \(s=\tfrac12+it\) 上有
\[
F(t)=\sum_{j,k}a_{j,k}\int_{0}^{A} w(x)\,e^{-jx/2}\,e^{-i(j-2k)tx}\,dx
=\sum_{m}\int_{0}^{A} g_m(x)\,e^{-imtx}\,dx,
\]
其中
\(
g_m(x):=w(x)\sum_{j-2k=m}a_{j,k}e^{-jx/2}\in L^{1}([0,A])
\)。
每一项均为紧支撑函数的 Fourier 变换，故由 Paley--Wiener 定理得其指数型至多为 \(|m|A\)，从而 \(\mathrm{type}(F)\le BA\)。
若对某个 \(|m|=B\) 有 \(g_m\) 在任意 \((A-\varepsilon,A]\) 上不几乎处处为零，则该项的 Fourier 变换指数型恰为 \(A\)，
而参数缩放 \(t\mapsto mt\) 将类型放大为 \(|m|A=BA\)；全和的指数型为各项指数型的最大值，故 \(\mathrm{type}(F)=BA\)。证毕。
\end{proof}

\begin{conclusion}[零点密度的复杂度乘积下界：匹配 \texorpdfstring{$T\log T$}{T log T} 必需 \texorpdfstring{$A(T)B(T)\gtrsim \tfrac12\log T$}{A(T)B(T)≳(1/2)log T}]\label{con:xi-terminal-zero-density-product-lower-bound}
设 \(\{F_T\}_{T\ge 1}\) 为一族 Hardy 切片整函数，
且每个 \(F_T(t)=\mathcal X_{\Delta_T,w_T}(\tfrac12+it)\) 满足结论 \ref{con:xi-terminal-scale-superposition-exact-type-rigidity} 的非退化条件，
其中 \(w_T\) 支撑于 \([0,A(T)]\) 且带宽为 \(B(T)\)。
令
\(
\tau(T):=\mathrm{type}(F_T)=A(T)B(T)
\)，并记对称零点计数
\[
\widetilde N_T(T):=\#\{t\in[-T,T]:F_T(t)=0\}.
\]
则 Cartwright 密度上界给出
\[
\widetilde N_T(T)\le \frac{2}{\pi}\,\tau(T)\,T+o(\tau(T)T)\qquad(T\to\infty).
\]
因此若欲在量级上达到黎曼型对称计数 \(\widetilde N_\zeta(T)\sim \frac{T}{\pi}\log T\)，则必有
\[
\liminf_{T\to\infty}\frac{\tau(T)}{\log T}\ \ge\ \frac12.
\]
\end{conclusion}
\begin{proof}[证明]
由结论 \ref{con:xi-terminal-scale-superposition-exact-type-rigidity}，每个 \(F_T\) 为指数型 \(\tau(T)\) 的整函数；并且 \(F_T\) 作为紧支撑 \(L^1\) 函数的 Fourier 变换在实轴上有界，
从而属于 Cartwright 类。
对 Cartwright 类指数型整函数应用经典 Jensen--Cartwright 零点密度估计，得到
\(
\widetilde N_T(T)\le (2/\pi)\tau(T)T+o(\tau(T)T)
\)。
将其与 \(\widetilde N_\zeta(T)\sim (T/\pi)\log T\) 对比即得所述下界。证毕。
\end{proof}

\begin{conclusion}[单尺度完成化的 Jensen--热带剖面：垂直平均对数模为分段仿射凸函数]\label{con:xi-terminal-single-scale-jensen-tropical-profile}
沿用命题 \ref{prop:xi-rh-interval-criterion} 与推论 \ref{cor:xi-rh-unit-circle-reciprocal} 的记号：
设
\[
\Xi_{\Delta,L}(s)=P_L(S_L(s)),\qquad S_L(s)=L^{s-\frac12}+L^{\frac12-s},
\]
并令 \(d:=\deg P_L\) 以及
\(
\Phi_L(y):=y^{d}P_L(y+y^{-1})\in\RR[y]
\)。
记 \(\{y_j\}_{j=1}^{2d}\) 为 \(\Phi_L\) 的复根（按重数计），且 \(\Phi_L\) 的首项系数为 \(c\neq 0\)。
定义垂直平均剖面
\[
M_L(\sigma):=\frac{\log L}{2\pi}\int_{0}^{2\pi/\log L}\log\bigl|\Xi_{\Delta,L}(\sigma+it)\bigr|\,dt.
\]
则对所有 \(\sigma\in\RR\) 有闭式
\[
M_L(\sigma)=\log|c|+\sum_{j=1}^{2d}\log\max\{L^{\sigma-\frac12},|y_j|\}-d\Bigl(\sigma-\frac12\Bigr)\log L.
\]
从而 \(M_L(\sigma)\) 是分段仿射凸函数，其所有折点恰为
\[
\sigma_j=\frac12+\frac{\log|y_j|}{\log L}\qquad(1\le j\le 2d).
\]
\end{conclusion}
\begin{proof}[证明]
令 \(y:=L^{s-\frac12}\)，则 \(y=L^{\sigma-\frac12}e^{-it\log L}\)，并且
\(
\Xi_{\Delta,L}(s)=P_L(y+y^{-1})=y^{-d}\Phi_L(y)
\)。
设 \(r:=L^{\sigma-\frac12}\)。作变量代换 \(\theta=t\log L\) 得
\[
M_L(\sigma)
=-d\Bigl(\sigma-\frac12\Bigr)\log L+\frac{1}{2\pi}\int_{0}^{2\pi}\log|\Phi_L(re^{-i\theta})|\,d\theta.
\]
对多项式 \(\Phi_L\) 应用 Jensen 公式：
\[
\frac{1}{2\pi}\int_{0}^{2\pi}\log|\Phi_L(re^{-i\theta})|\,d\theta
=\log|c|+\sum_{j=1}^{2d}\log\max\{r,|y_j|\},
\]
代回并用 \(r=L^{\sigma-\frac12}\) 即得所示闭式与折点位置。证毕。
\end{proof}

\begin{conclusion}[单尺度完成化的离线谱测度原子化：剖面二阶导数等于竖线簇计数测度]\label{con:xi-terminal-single-scale-offline-spectral-measure-atomicity}
在结论 \ref{con:xi-terminal-single-scale-jensen-tropical-profile} 的设定下，\(M_L\) 为凸函数，故其分布意义二阶导数为有限原子测度：
\[
\frac{d^{2}}{d\sigma^{2}}M_L(\sigma)=\log L\sum_{j=1}^{2d}\mathrm{mult}(y_j)\,\delta_{\sigma_j},
\qquad
\sigma_j=\frac12+\frac{\log|y_j|}{\log L}.
\]
并且令 \(\mathcal Z\) 为 \(\Xi_{\Delta,L}\) 的零点集合（按重数计），则对任意连续紧支撑测试函数 \(\phi\) 有
\[
\lim_{T\to\infty}\frac{\pi}{T\log L}\sum_{\substack{s\in\mathcal Z\\ |\Im s|\le T}}\phi(\Re s)
\;=\;
\sum_{j=1}^{2d}\mathrm{mult}(y_j)\,\phi(\sigma_j),
\]
即经自然归一化后，“零点在实部方向的边际分布”严格坍缩为有限原子，原子位置正是 \(M_L\) 的折点。
\end{conclusion}
\begin{proof}[证明]
由结论 \ref{con:xi-terminal-single-scale-jensen-tropical-profile}，
每项 \(\log\max\{L^{\sigma-\frac12},|y_j|\}\) 的一阶导数为阶跃函数，
其二阶导数在分布意义下为 \((\log L)\delta_{\sigma_j}\)；乘以重数并求和得第一式。
\par
另一方面，\(\Xi_{\Delta,L}(s)=0\iff \Phi_L(y)=0\) 且 \(y=L^{s-\frac12}\)。
若 \(y_j\) 为 \(\Phi_L\) 的根，则零点塔为
\[
s=\frac12+\frac{\log y_j+2\pi i k}{\log L},\qquad k\in\ZZ,
\]
其全部实部恒等于 \(\sigma_j\)。
在窗 \(|\Im s|\le T\) 内每个 \(y_j\) 贡献 \( (T\log L/\pi)+O(1)\) 个零点（按重数计），故归一化后对任意 \(\phi\) 收敛到所示原子和。证毕。
\end{proof}

\begin{conclusion}[\texorpdfstring{$\varphi^{\mathbb Q}$}{phi^Q} 尺度混合的共振障碍：有限 \texorpdfstring{$\varphi^{\mathbb Q}$}{phi^Q} 叠加必保留非零共同虚周期]\label{con:xi-terminal-phi-q-scale-mix-common-imag-period}
设 \(L_1,\dots,L_n>1\) 满足
\[
\frac{\log L_j}{\log\varphi}\in\QQ,\qquad j=1,\dots,n.
\]
取分母最小公倍数 \(q\) 使 \(\log L_j=(p_j/q)\log\varphi\)（\(p_j\in\ZZ\)）。
则对
\[
T_{\varphi}:=\frac{2\pi i\,q}{\log\varphi}\neq 0
\]
有 \(L_j^{-T_{\varphi}}=1\)（\(\forall j\)）。
因而对任意定义于 \(\CC^{\times}\) 的亚纯函数族 \(\{f_j\}\) 所构成的有限叠加
\[
F(s):=\sum_{j=1}^{n} f_j(L_j^{-s})
\]
都满足严格虚周期性 \(F(s+T_{\varphi})\equiv F(s)\)。
特别地，任何仅由有限个 \(\varphi^{\mathbb Q}\) 尺度构成的完成化叠加不可能在竖直方向生成非周期谱结构：其零极点集合必含共同步长 \(2\pi/\log\varphi\) 的算术梳齿。
\end{conclusion}
\begin{proof}[证明]
由 \(\log L_j=(p_j/q)\log\varphi\) 得
\[
L_j^{-T_{\varphi}}=\exp\!\bigl(-T_{\varphi}\log L_j\bigr)
=\exp\!\bigl(-2\pi i\,q\cdot (p_j/q)\bigr)=1.
\]
于是 \(f_j(L_j^{-(s+T_{\varphi})})=f_j(L_j^{-s})\)，对 \(j\) 求和即得 \(F(s+T_{\varphi})=F(s)\)。证毕。
\end{proof}


\paragraph{跨接口终端结论：奇偶权模态、同余 Fourier 扭曲与对称幂不交核的全显式封口}
本段收集三类可直接复核的“常数/闭式”条目：其一记录黄金均值稳定词 \(X_m\) 的奇偶权交错计数，
该对象对应一个隐含的 \(6\)-次割圆模态；其二把 Fold 纤维多重度按 \(|x|_1\) 的根单位扭曲提升为双变量配分函数 \(Z_m(y)\)，
由其四阶可积递推与显式有理生成函数抽取出 Tribonacci 主振荡模、Lee--Yang 型谱分歧边界、推前重量统计的
大数律/中心极限定律与全局大偏差轮廓；并进一步把纤维幂和矩提升为重量—碰撞配分函数 \(Z_m^{(q)}(y)\)（\(q=2,3,4\)），
给出其低阶常系数递推与有理生成函数、\(y=1\) 投影下的幽灵模态严格约去，以及由判别式分解所诱导的谱分歧量子化；
同时并联给出模 \(2,3,4\) 同余通道的谱半径与混合层级；
其三把不交图 \(\mathcal D_q\) 的 \(S_q\)-对称压缩矩阵 \(B_q\) 识别为 \(\mathrm{Sym}^q(M)\)，
并由此得到幺模性、显式逆矩阵与 \(n\)-步传输系数的 Fibonacci--二项式闭式。

\begin{conclusion}[黄金均值稳定词的奇偶权交错计数：严格 \(6\)-周期与有理生成函数]\label{con:xi-terminal-xm-parity-weight-6period}
令
\[
X_m:=\{x\in\{0,1\}^m:\ \forall i,\ x_ix_{i+1}=0\},
\qquad
A_m^{\mathrm{par}}:=\sum_{x\in X_m}(-1)^{|x|_1}.
\]
则对一切 \(m\ge 3\) 有二阶递推
\[
A_m^{\mathrm{par}}=A_{m-1}^{\mathrm{par}}-A_{m-2}^{\mathrm{par}},
\qquad
A_1^{\mathrm{par}}=0,\ A_2^{\mathrm{par}}=-1,
\]
从而 \(\bigl(A_m^{\mathrm{par}}\bigr)_{m\ge 1}\) 为严格 \(6\)-周期，并且
\[
A_m^{\mathrm{par}}=0\iff m\equiv 1,4\pmod 6,\qquad
A_m^{\mathrm{par}}=-1\iff m\equiv 2,3\pmod 6,\qquad
A_m^{\mathrm{par}}=1\iff m\equiv 0,5\pmod 6.
\]
若约定 \(A_0^{\mathrm{par}}:=1\)（空字的权重为 \(1\)），则生成函数满足闭式
\[
\sum_{m\ge 0}A_m^{\mathrm{par}}t^m=\frac{1-t}{1-t+t^2}.
\]
\end{conclusion}
\begin{proof}[证明]
按首段分解 \(X_m=\{0u:u\in X_{m-1}\}\sqcup \{10v:v\in X_{m-2}\}\)（\(m\ge 3\)）。
前一类不改变 \(|x|_1\)，后一类额外引入一个 \(1\)，故权重贡献多一因子 \((-1)\)。
于是得到递推
\(
A_m^{\mathrm{par}}=A_{m-1}^{\mathrm{par}}-A_{m-2}^{\mathrm{par}}
\)
及初值（由 \(X_1=\{0,1\},X_2=\{00,01,10\}\) 直接核算）。
特征方程为 \(\lambda^2-\lambda+1=0\)，根为 \(e^{\pm i\pi/3}\)，故序列严格 \(6\)-周期；
逐项展开即可得到模 \(6\) 的取值分类。
生成函数由标准递推求和（或直接对 \((1-t+t^2)\sum_{m\ge 0}A_m^{\mathrm{par}}t^m\) 展开并代入初值）得到。证毕。
\end{proof}

\begin{conclusion}[纤维—重量双变量配分函数的四阶可积性：有理生成函数与常系数递推]\label{con:xi-terminal-zm-order4-rational}
令
\[
X_m:=\{x\in\{0,1\}^m:\ \forall i,\ x_ix_{i+1}=0\},
\qquad
\Fold_m:\Omega_m=\{0,1\}^m\to X_m
\]
为标准折叠映射，并记纤维多重度 \(d_m(x)=|\Fold_m^{-1}(x)|\)。
对任意复参数 \(y\in\CC\)，定义纤维—重量双变量配分函数
\[
Z_m(y):=\sum_{x\in X_m} y^{|x|_1}\,d_m(x)
=\sum_{a\in\Omega_m} y^{|\Fold_m(a)|_1}.
\]
则对一切 \(m\ge 4\) 成立严格四阶递推
\[
\boxed{\;
Z_m(y)=Z_{m-1}(y)+(2y+1)Z_{m-2}(y)-Z_{m-3}(y)-y(y+1)Z_{m-4}(y)\;}
\]
并且生成函数在 \(\ZZ[y][[t]]\) 中具有显式有理闭式
\[
\boxed{\;
\sum_{m\ge 0} Z_m(y)\,t^m
=\frac{1+yt-yt^2-y^2t^3}{1-t-(2y+1)t^2+t^3+y(y+1)t^4}\; }.
\]
因此，对固定 \(y\) 而言，序列 \(m\mapsto Z_m(y)\) 的全部指数模态由四次特征方程
\[
\Pi(\lambda,y)=0,
\qquad
\Pi(\lambda,y):=\lambda^4-\lambda^3-(2y+1)\lambda^2+\lambda+y(y+1)
\]
完全控制。
\end{conclusion}
\begin{proof}[证明要点]
由定义 \(Z_m(y)\) 是长度 \(m\) 的微态路径权重和，其中每个微态 \(a\in\Omega_m\) 的权重仅通过其稳定输出 \(\Fold_m(a)\) 的 \(1\)-计数进入。
由于 \(\Fold_m\) 可由有限状态规约器实现（全文对 Zeckendorf/窗口折叠的固定口径），
\(Z_m(y)\) 可写为某个固定维数带权转移矩阵 \(T(y)\) 的矩阵系数；
从而 \(\sum_{m\ge 0}Z_m(y)t^m\) 为有理函数，分母次数等于最小状态数。
对本 Fold 实例，约分后的分母次数稳定为 \(4\)，从而得到四阶递推；
将 \(t\)-生成函数在 \(t=0\) 处展开并与初值匹配，即得到所示闭式分子与分母；
上述恒等式与后续的判别式/统计常数可由脚本 \texttt{scripts/exp\_fold\_zm\_bivariate\_partition\_audit.py} 一键复核。证毕。
\end{proof}

\begin{conclusion}[\texorpdfstring{$Z_m(y)$}{Zm(y)} 的次数、首项系数与分歧参数处的闭式特值]\label{con:xi-terminal-zm-degree-leading-special-values}
沿用结论 \ref{con:xi-terminal-zm-order4-rational} 的记号。
则对所有 \(m\ge 0\)：
\begin{enumerate}
  \item 次数精确为
  \[
  \deg_y Z_m(y)=\left\lceil\frac m2\right\rceil.
  \]
  \item 设 \(m=2k+1\) 为奇数，则 \(Z_{2k+1}(y)\) 为首一多项式；设 \(m=2k\) 为偶数，则其首项系数为
  \[
  [y^{k}]\,Z_{2k}(y)=1+\frac{k(k+1)}{2}.
  \]
  \item 在两条有理分歧参数处有闭式特值
  \[
  Z_m(0)=\left\lfloor\frac m2\right\rfloor+1,\qquad
  Z_m(1)=2^m.
  \]
  \item 若将 \(y\) 特化到椭圆轨道 \(y(nP)\in\QQ\)（其中 \(\mathrm{den}(y(nP))=v_n^4\)，见结论 \ref{con:xi-terminal-zm-elliptic-rational-points-denominator-law}），则对任意 \(m,n\ge 0\) 有
  \[
  Z_m\bigl(y(nP)\bigr)\in v_n^{-4\lceil m/2\rceil}\,\ZZ.
  \]
\end{enumerate}
\end{conclusion}
\begin{proof}[证明要点]
由 \(X_m\) 的语法约束 \(x_ix_{i+1}=0\) 可知 \(\max_{x\in X_m}|x|_1=\lceil m/2\rceil\)，且存在取等词；
又 \(Z_m(y)=\sum_{x\in X_m}d_m(x)y^{|x|_1}\) 且 \(d_m(x)\in\ZZ_{>0}\)，故次数结论成立。
首项系数由最大重量层的谱与纤维闭式给出（结论 \ref{con:xi-terminal-max-weight-layer-fiber-spectrum-closed}，亦即结论 \ref{con:xi-terminal-zq-top-coefficient-faulhaber} 在 \(q=1\) 的特例）。
\(Z_m(0)=d_m(0^m)=\lfloor(m+2)/2\rfloor=\lfloor m/2\rfloor+1\)（结论 \ref{con:xi-terminal-zero-weight-fiber-closed}），而 \(Z_m(1)=\sum_x d_m(x)=|\Omega_m|=2^m\)。
最后一条由 \(Z_m\in\ZZ[y]\) 且 \(\deg_y Z_m=\lceil m/2\rceil\)，结合 \(\mathrm{den}(y(nP))=v_n^4\) 直接推出。证毕。
\end{proof}

\begin{conclusion}[有限窗 Sturm 实根交错：\texorpdfstring{$m\le 100$}{m<=100} 时 \texorpdfstring{$Z_m(y)$}{Z_m(y)} 的负实根性]\label{con:xi-terminal-zm-sturm-real-roots-interlacing-m100}
对每个整数 \(1\le m\le 100\)，多项式 \(Z_m(y)\in\ZZ[y]\) 具有 \(\lceil m/2\rceil\) 个互异实根，且全部位于 \((-\infty,0)\)。
并且对每个 \(0\le m\le 98\)，\(Z_m\) 与 \(Z_{m+2}\) 的零点集合严格交错；因此固定奇偶类的两条子序列 \(\{Z_{2k}\}_{k\ge 0}\) 与 \(\{Z_{2k+1}\}_{k\ge 0}\)（在该核验窗内）各自形成一条 Sturm 实根链。
\par
上述核验可由脚本 \texttt{scripts/exp\_fold\_zm\_sturm\_real\_roots\_audit.py} 复算，其输出证书为 \texttt{artifacts/export/fold\_zm\_sturm\_real\_roots\_audit.json}。
\end{conclusion}
\begin{proof}[证明要点]
由结论 \ref{con:xi-terminal-zm-order4-rational} 的四阶递推可在 \(\ZZ[y]\) 中逐步生成 \(Z_m(y)\)。
脚本对核验窗内每个 \(Z_m\) 作高精度求根并检查：
（i）全部根的虚部在数值容许阈值内消失，从而根全为实数；
（ii）由于系数非负且常数项非零，实根必落在 \((-\infty,0)\)，并进一步核验互异性；
（iii）对同奇偶类的 \(m\mapsto m+2\) 根集逐项检验严格交错不等式链。证毕。
\end{proof}

\begin{conjecture}[负实轴完全性与全尺度交错]\label{conj:xi-terminal-zm-sturm-real-roots-interlacing}
对任意 \(m\ge 1\)，多项式 \(Z_m(y)\) 的零点全部为互异负实根，并且对任意 \(m\ge 0\)，\(Z_m\) 与 \(Z_{m+2}\) 的零点集合严格交错。
\end{conjecture}

\begin{remark}[一种结构性证明路线]\label{rem:xi-terminal-zm-sturm-proof-route}
记系数表 \(a_{m,k}:=[y^k]Z_m(y)\)。
若猜想 \ref{conj:xi-terminal-zm-sturm-real-roots-interlacing} 成立，则对每个 \(m\) 系数列 \(\bigl([y^k]Z_m(y)\bigr)_k\) 严格对数凹，因而严格单峰；若进一步可提升为 P\'olya--frequency（PF）序列，则全体阶 Hankel 行列式非负并蕴含多重总正性。结合结论 \ref{con:xi-terminal-zm-double-root-diffusive-sinc-limit} 与结论 \ref{con:xi-terminal-zm-diffusive-edge-weyl-spectral-determinant} 的边缘扩散谱刚性，可把“近零量子化”与“全局负实轴束缚”视为同一 Sturm/全正性机制在边缘与全域的两种投影。
本体系中已核验的二维局域递推
\[
a_{m,k}=a_{m-1,k}+a_{m-2,k}+2a_{m-2,k-1}-a_{m-3,k}-a_{m-4,k-1}-a_{m-4,k-2}
\]
为构造上述网络/矩阵实现提供直接代数入口。
\end{remark}

\begin{conclusion}[最小状态复杂度与退化点：除 \(y\in\{0,-1,1\}\) 外 Hankel 秩恰为 \(4\)]\label{con:xi-terminal-zm-hankel-rank4}
沿用结论 \ref{con:xi-terminal-zm-order4-rational} 的记号，并记
\(
\Pi(\lambda,y)=\lambda^4-\lambda^3-(2y+1)\lambda^2+\lambda+y(y+1)
\)。
若 \(y\in\CC\setminus\{0,-1,1\}\)，则
序列 \(\{Z_m(y)\}_{m\ge 0}\) 的最小湮灭多项式恰为 \(\Pi(\lambda,y)\)；
等价地，\(\{Z_m(y)\}\) 的 Hankel 秩为 \(4\)，不存在三阶及以下的常系数递推能够对所有 \(m\) 成立。
\par
当 \(y\in\{0,-1\}\) 时，系数 \(y(y+1)=0\) 使四阶递推在结构上退化为三阶递推，最小阶严格下降；
当 \(y=1\) 时，归一化恒等式 \(Z_m(1)=\sum_x d_m(x)=2^m\) 使生成函数发生完全抵消，从而 Hankel 秩塌缩为 \(1\)。
\end{conclusion}
\begin{proof}[证明要点]
由结论 \ref{con:xi-terminal-zm-order4-rational}，\(\sum_{m\ge 0}Z_m(y)t^m\) 的分母为
\(
D(t)=1-t-(2y+1)t^2+t^3+y(y+1)t^4
\)。
当 \(y\in\{0,-1\}\) 时，\(D(t)\) 的四次项消失，递推阶数相应降低。
当 \(y=1\) 时，分子与分母共享因子 \((1-t)(1+t)^2\)，从而生成函数约化为 \(\sum_{m\ge 0}2^m t^m=(1-2t)^{-1}\)，最小阶降为 \(1\)。
对任意 \(y\in\CC\setminus\{0,-1,1\}\)，\(\deg_t D=4\) 且分子分母无非平凡公因子，
故约分后分母次数仍为 \(4\)，从而最小递推阶数为 \(4\)；Hankel 秩表述与之等价。证毕。
\end{proof}

% (User input window)

\begin{conclusion}[\texorpdfstring{$y=0$}{y=0} 的扩散 Lee--Yang 边缘：\texorpdfstring{$\operatorname{sinc}$}{sinc} 解析核与近零根量子化]\label{con:xi-terminal-zm-double-root-diffusive-sinc-limit}
沿用结论 \ref{con:xi-terminal-zm-order4-rational} 的多项式族 \(Z_m(y)\)。
等价地，\(\{Z_m(y)\}_{m\ge 0}\) 可由初值
\[
Z_0(y)=1,\quad
Z_1(y)=1+y,\quad
Z_2(y)=2+2y,\quad
Z_3(y)=2+5y+y^2
\]
及四阶递推
\[
Z_m(y)
=Z_{m-1}(y)+(2y+1)Z_{m-2}(y)-Z_{m-3}(y)-y(y+1)Z_{m-4}(y)
\qquad(m\ge 4)
\]
唯一确定。
记
\[
N_m:=Z_m(0)=\left\lfloor\frac m2\right\rfloor+1
\]
（见结论 \ref{con:xi-terminal-zm-degree-leading-special-values}）。
定义
\[
\operatorname{sinc}(z):=
\begin{cases}
\dfrac{\sin z}{z}, & z\neq 0,\\[1mm]
1, & z=0.
\end{cases}
\qquad
S(u):=\operatorname{sinc}\!\left(\frac{\sqrt u}{2}\right).
\]
则 \(S\) 由其幂级数展开唯一确定，因而作为 \(u\) 的整函数无需选择平方根分支。
并且对任意紧集 \(K\subset\CC\)，当 \(m\to\infty\) 时成立一致估计
\[
\boxed{\;
\sup_{u\in K}\left\lvert
\frac{Z_m\!\left(-\frac{u}{m^{2}}\right)}{N_m}-S(u)
\right\rvert=O\!\left(\frac1m\right)\;}
\]
从而 \(y=0\) 处的二重特征根退化在 Lee--Yang 方向的扩散标度 \(y=-u/m^2\) 下由显式整函数核 \(S(u)\) 统一控制。
\par
由于 \(N_m=\frac m2+O(1)\)，上述一致估计亦可改写为：对任意紧集 \(K\subset\CC\)，有
\[
\sup_{u\in K}\left\lvert
\frac{1}{m}Z_m\!\left(-\frac{u}{m^{2}}\right)-\frac12 S(u)
\right\rvert=O\!\left(\frac1m\right)
\qquad(m\to\infty),
\]
其中
\(
\frac12 S(u)=\sum_{k\ge 0}\frac{(-u)^k}{2^{2k+1}(2k+1)!}
=\frac{\sin(\sqrt{u}/2)}{\sqrt{u}}
\)
以连续延拓在 \(u=0\) 处取值 \(1/2\)。

特别地，令 \(j\ge 1\) 固定，并记 \(\rho_{m,1}>\rho_{m,2}>\cdots\) 为 \(Z_m\) 在 \((-\infty,0)\) 上按从 \(0\) 向左排序的零点（若存在多重根则按重数计）。
则对每个固定 \(j\)，当 \(m\) 足够大时 \(\rho_{m,j}\) 存在且满足量子化定标
\[
\boxed{\;
m^{2}\bigl(-\rho_{m,j}\bigr)\longrightarrow 4\pi^{2}j^{2},
\qquad
\rho_{m,j}=-\frac{4\pi^{2}j^{2}}{m^{2}}+O\!\left(\frac1{m^{3}}\right)\;}
\qquad(m\to\infty),
\]
并且对任意 \(\varepsilon>0\)，当 \(m\) 充分大时，\(Z_m\) 在圆盘
\[
\left|y+\frac{4\pi^2 j^2}{m^2}\right|<\frac{\varepsilon}{m^2}
\]
内恰有一个零点（计重数）。
等价地，定义 \(\sigma_{m,j}:=\frac{m}{2}\sqrt{-\rho_{m,j}}\)，则 \(\sigma_{m,j}\to \pi j\)，体现出边缘零点在 \(\sqrt{-y}\) 坐标下的等间距“时钟过程”主阶。
\end{conclusion}
\begin{proof}[证明要点]
由结论 \ref{con:xi-terminal-zm-order4-rational} 的有理生成函数，
在 \(y=0\) 处分母因式分解为 \(D(t,0)=(1-t)^2(1+t)\)，故 \(t=1\) 为二重极点，
对应于递推特征方程在 \(\lambda=1\) 的二重特征根退化。
当 \(y=-u/m^2\) 时，与该二重极点相邻的两条简单极点对以间距 \(O(m^{-1})\) 分裂；
对其作局部留数展开并利用极限
\(
(1\pm i\alpha/m)^{-m}\to e^{\mp i\alpha}
\)
得到主项为 \(\operatorname{sinc}(\sqrt{u}/2)\) 的尺度极限；
一致性来自于对 \(u\) 在紧集上的统一估计。
归一化 \(Z_m(-u/m^2)/N_m\) 与 \((2/m)Z_m(-u/m^2)\) 的差异由 \(N_m=\frac m2+O(1)\) 吸收，并给出所示 \(O(m^{-1})\) 余项。
零点定标与局部唯一性由 Hurwitz 定理对上述一致估计与 \(S(u)\) 的简单零点 \(u=4\pi^{2}j^{2}\) 应用而得；
进一步由零点的非退化性与余项 \(O(m^{-1})\) 推出 \(u\)-坐标的偏移为 \(O(m^{-1})\)，从而 \(\rho_{m,j}=-u_{m,j}/m^2\) 的偏移为 \(O(m^{-3})\)。
可复现审计脚本：\texttt{scripts/exp\_fold\_zm\_double\_root\_diffusive\_sinc\_limit\_audit.py}。
\end{proof}

\begin{conclusion}[边缘零点计数的 Weyl 型律与谱行列式识别]\label{con:xi-terminal-zm-diffusive-edge-weyl-spectral-determinant}
在结论 \ref{con:xi-terminal-zm-double-root-diffusive-sinc-limit} 的记号下，对任意固定 \(U>0\) 定义边缘零点计数函数
\[
\mathcal N_m(U):=\#\Bigl\{j:\ \rho_{m,j}\in\bigl[-U/m^2,\,0\bigr)\Bigr\}.
\]
则当 \(m\to\infty\) 时有
\[
\boxed{\;
\mathcal N_m(U)=\left\lfloor \sqrt{\frac{U}{4\pi^2}}\right\rfloor+o(1)\;}
\]
从而在长度尺度 \(m^{-2}\) 的窗口内，零点计数呈 \(\sqrt{U}\) 增长的 Weyl 型指数。
此外，极限核 \(S(u)=\operatorname{sinc}(\sqrt{u}/2)\) 具有 Weierstrass 乘积分解
\[
\boxed{\;
S(u)=\prod_{j=1}^{\infty}\left(1-\frac{u}{4\pi^2 j^2}\right)\;}
\]
其零点集合 \(\{4\pi^2 j^2\}_{j\ge 1}\) 与单位区间 Dirichlet 拉普拉斯算子 \(-\frac{\dd^2}{\dd x^2}\) 的本征值 \(\{\pi^2 j^2\}_{j\ge 1}\) 仅差一个固定尺度因子 \(4\)。
因此，扩散标度下的归一化极限可视为一个显式谱行列式对象，从而把离散折叠—投影体系在 \(y=0\) 边缘的零点刚性与连续二阶算子谱结构对齐。
\end{conclusion}
\begin{proof}[证明要点]
由结论 \ref{con:xi-terminal-zm-double-root-diffusive-sinc-limit} 的零点定标，固定 \(U\) 时 \(\mathcal N_m(U)\) 收敛到 \(S(u)\) 在区间 \([0,U]\) 内零点数目，而 \(S(u)=0\) 当且仅当 \(u=4\pi^2 j^2\)，故得到所示计数律。
Weierstrass 乘积由经典恒等式 \(\sin z/z=\prod_{j\ge 1}(1-z^2/(\pi^2 j^2))\) 取 \(z=\sqrt{u}/2\) 直接推出；谱字典来自 Dirichlet 拉普拉斯算子的本征值闭式 \(\pi^2 j^2\)。证毕。
\end{proof}

\begin{conclusion}[\texorpdfstring{$y=0$}{y=0} 的临界窗宽为 \texorpdfstring{$m^{-2}$}{m^{-2}}：线性增长与指数增长之间的唯一尺度]\label{con:xi-terminal-zm-critical-window-m2}
在 \(y=0\) 处，结论 \ref{con:xi-terminal-zm-hankel-rank4} 所述二重特征根使 \(Z_m(0)\) 仅呈线性增长；
而对任意固定 \(y>0\)，谱半径严格大于 \(1\)，故 \(Z_m(y)\) 具有严格指数增长率。
结论 \ref{con:xi-terminal-zm-double-root-diffusive-sinc-limit} 的扩散标度极限表明，
两种增长机制之间的非平凡过渡被刚性锁定在唯一临界窗
\[
y=-\frac{u}{m^{2}},
\qquad
u\asymp 1,
\]
并在该窗内由整函数 \(u\mapsto S(u)\) 统一控制。
相应地，靠近 \(0\) 的零点簇在该窗内呈现 \(m^{-2}\) 级别的量子化定标
（结论 \ref{con:xi-terminal-zm-double-root-diffusive-sinc-limit}），
其主阶间距由 \(\sin\) 的整数倍零点完全决定。
\end{conclusion}



\begin{conclusion}[谱分歧轨迹的完全因式分解：Lee--Yang 边界由一元三次给出]\label{con:xi-terminal-zm-discriminant-leyang}
视 \(\Pi(\lambda,y)\) 为关于 \(\lambda\) 的四次多项式，则其判别式在 \(\ZZ[y]\) 中精确分解为
\[
\boxed{\;
\mathrm{Disc}_{\lambda}\bigl(\Pi(\lambda,y)\bigr)
=-y\,(y-1)\,\bigl(256y^3+411y^2+165y+32\bigr)\; }.
\]
为后续引用，记
\(
\mathrm{disc}^{\ast}_{\lambda}\Pi:=\mathrm{Disc}_{\lambda}\bigl(\Pi(\lambda,y)\bigr)
\)。
在该记号下，Lee--Yang 型边界即为 \(\mathrm{disc}^{\ast}_{\lambda}\Pi\) 在 \(\mathbb P^1_y\) 上的零点集合。
因此，\(\lambda\)-谱作为 \(y\) 的代数函数仅在
\(
y\in\{0,1\}
\)
与三次方程
\(
256y^3+411y^2+165y+32=0
\)
的根处发生分歧。
其中该三次多项式恰有一个实根
\[
y_{\mathrm{LY}}\approx-1.1344520030,
\]
并给出双变量自由能
\[
f(y):=\lim_{m\to\infty}\frac{1}{m}\log Z_m(y)
\]
在实轴上的最近边界奇点；从而 \(y\) 作为复变量时，\(f\) 的最小非平凡解析半径由 \(y_{\mathrm{LY}}\) 所决定（Lee--Yang 型边界机制；参见 \cite{XuZamolodchikov2022IsingYangLee}）。
\end{conclusion}
\begin{proof}[证明要点]
对四次多项式 \(\Pi(\lambda,y)\) 计算 \(\lambda\)-判别式并在 \(\ZZ[y]\) 中因式分解即可得到所示公式。
判别式为零当且仅当存在重根，亦即代数分支在该参数处发生分歧；三次因子的实根可由标准数值求根获得。证毕。
\end{proof}

\subsubsection{Lee--Yang 零点的射影消混合：交比不变量与 crossing 渐近}\label{subsubsec:xi-leyang-projective-demixing-crossing}
在有限体积配分函数的 Lee--Yang 零点分析中，需要区分两类互补机制。
一类来自有限维谱分解的主导谱竞争与判别式分歧：当两条主导谱支在等模曲线上发生竞争时，零点在法向方向出现 \(\frac{\log n}{n}\) 量级的定量位移；当谱曲线在临界点具有分歧指标 \(m\) 时，零点在尺度 \(n^{-m}\) 上聚束到该分歧点并呈 Puiseux 型准晶格。
另一类来自解析背景与变量混合：允许的复仿射重参数化 \(z\mapsto az+b(\tau)\) 会引入共同平移与整体缩放，使得直接以零点坐标作比值或作体积交叉的定位量对混合参数敏感。
以下首先记录可直接用于“判别式分解 \(\Rightarrow\) 零点标度反演”的谱侧定量律，随后给出射影化的消混合接口。

\begin{theorem}[两相并存临界窗的尺度极限与零点凝聚]\label{thm:xi-godel-leyang-two-phase-coexistence-window}
设 \(\mathcal C\) 为可数概念类，并给定前缀自由码长 \(\ell:\mathcal C\to\NN\)，权重 \(w(C):=2^{-\ell(C)}\)。
对观测前缀长度 \(n\) 与参数 \(\beta\in\CC\)，定义有限体积配分函数
\[
Z_n(\beta):=\sum_{C\in\mathcal C}w(C)\exp\!\bigl(-\beta E_C(n)\bigr),
\]
其中 \(E_C(n)\in\RR\) 为可加失配能量。
假设存在两概念 \(C_\pm\in\mathcal C\) 与常数 \(e_\pm,\lambda_\pm\in\RR\) 使
\[
\frac{E_{C_\pm}(n)}{n}\to e_\pm,\qquad \frac{\ell(C_\pm)}{n}\to \lambda_\pm,
\qquad
\Delta e:=e_+-e_-\neq 0,
\]
并且存在实数 \(\beta^\star\) 使两相自由能并存且严格主导：对任意其余概念 \(C\notin\{C_+,C_-\}\)，若极限
\(
e(C):=\lim_{n\to\infty}E_C(n)/n
\) 与
\(
\lambda(C):=\lim_{n\to\infty}\ell(C)/n
\)
存在，则成立严格不等式
\[
(\ln 2)\lambda_+ + \beta^\star e_+
=
(\ln 2)\lambda_- + \beta^\star e_-
\ <\ (\ln 2)\lambda(C)+\beta^\star e(C).
\]
对 \(z\in\CC^\times\)，定义临界窗重标度
\[
\beta_{n}(z):=\beta^\star-\frac{1}{n\Delta e}\log z,
\]
并定义归一化局部配分函数
\[
\mathcal Z_n(z):=
\frac{Z_n\!\bigl(\beta_n(z)\bigr)}
{w(C_-)\exp\!\bigl(-\beta_n(z)E_{C_-}(n)\bigr)}.
\]
则对任意紧集 \(K\Subset\CC\setminus\{0\}\)，有一致收敛
\[
\mathcal Z_n(z)\ \longrightarrow\ 1+z
\qquad(n\to\infty,\ z\in K).
\]
因此在该尺度下，\(\mathcal Z_n\) 的零点在 \(K\) 上趋于 \(\{-1\}\)；并且最靠近 \(\beta^\star\) 的 Lee--Yang 零点满足
\[
\beta_{n,k}
\;=\;
\beta^\star-\frac{(2k+1)\pi i}{n\Delta e}+o(n^{-1}),
\qquad k\in\ZZ.
\]
\end{theorem}
\begin{proof}
将 \(Z_n\) 分解为两主导项与余项：
\[
Z_n(\beta)=w(C_-)\,e^{-\beta E_{C_-}(n)}+w(C_+)\,e^{-\beta E_{C_+}(n)}+R_n(\beta).
\]
把 \(\beta=\beta_n(z)\) 代入并按 \(w(C_-)\,e^{-\beta E_{C_-}(n)}\) 归一化，得到
\[
\mathcal Z_n(z)=1+\frac{w(C_+)}{w(C_-)}\exp\!\bigl(-\beta_n(z)\,(E_{C_+}(n)-E_{C_-}(n))\bigr)
\;+\;\frac{R_n(\beta_n(z))}{w(C_-)\,e^{-\beta_n(z)E_{C_-}(n)}}.
\]
由 \(w(C_\pm)=2^{-\ell(C_\pm)}\) 与假设给出的并存等式，可得
\[
\frac{w(C_+)}{w(C_-)}\exp\!\bigl(-\beta^\star(E_{C_+}(n)-E_{C_-}(n))\bigr)=1+o(1).
\]
并且由 \(\beta_n(z)=\beta^\star-\frac{1}{n\Delta e}\log z\) 与 \((E_{C_+}(n)-E_{C_-}(n))/n\to \Delta e\)，有
\[
\exp\!\left(+\frac{\log z}{n\Delta e}\,(E_{C_+}(n)-E_{C_-}(n))\right)=z^{1+o(1)}.
\]
合并得到第二项趋于 \(z\)，从而给出 \(1+z\) 的主极限。
余项 \(\frac{R_n(\beta_n(z))}{w(C_-)\,e^{-\beta_n(z)E_{C_-}(n)}}\) 由严格主导条件在 \(K\) 上以指数速率趋于 \(0\)，从而得到一致收敛。
零点渐近式由 \(1+z=0\) 与 \(z=\exp(-n\Delta e(\beta-\beta^\star))\) 反解，并把误差吸收进 \(o(n^{-1})\) 得到。证毕。
\end{proof}

\begin{theorem}[两谱竞争下零点的对数位移定量律]\label{thm:xi-leyang-logshift-two-spectrum}
设 \(U\subset\CC\) 为开集，\(\lambda_1,\dots,\lambda_s:U\to\CC\setminus\{0\}\) 为解析函数。
对 \(n\in\NN\) 定义
\[
P_n(z):=\sum_{j=1}^{s}\alpha_j(n;z)\,\lambda_j(z)^n,
\qquad
\alpha_j(n;z)=\sum_{r=0}^{d_j}p_{j,r}(z)\,n^r,
\]
其中各 \(p_{j,r}\) 解析且 \(p_{j,d_j}\not\equiv 0\)。
定义严格两谱等模集
\[
\Gamma:=\Bigl\{z\in U:\ |\lambda_1(z)|=|\lambda_2(z)|>\max_{3\le j\le s}|\lambda_j(z)|\Bigr\},
\]
并取紧集 \(K\Subset \Gamma\)。
假设存在邻域 \(V\Subset U\) 与常数 \(0<\theta<1\) 使得对一切 \(z\in V\) 有
\[
\max_{3\le j\le s}\frac{|\lambda_j(z)|}{|\lambda_1(z)|}\le \theta,
\]
并且 \(\lambda_1/\lambda_2\) 在 \(V\) 上无零点，从而可取解析对数分支
\[
\mathcal{L}(z):=\log\!\Bigl(\frac{\lambda_1(z)}{\lambda_2(z)}\Bigr),
\qquad
\mathcal{L}'(z)\neq 0\ \ (z\in K).
\]
再设 \(p_{1,d_1}(z)p_{2,d_2}(z)\neq 0\) 于 \(K\) 上成立，并记
\[
\beta_n(z):=\frac{\alpha_2(n;z)}{\alpha_1(n;z)},
\qquad
D:=\max_{1\le j\le s}d_j.
\]
则存在邻域 \(W\Subset V\)（\(K\subset W\)）与 \(n_0\in\NN\) 使得对所有 \(n\ge n_0\)：
\begin{enumerate}
  \item \(P_n\) 在 \(W\) 内的全部零点皆为单零点；
  \item 若 \(z\in W\) 为零点，则存在整数 \(k\in\ZZ\) 使
  \begin{equation}\label{eq:xi-leyang-logshift-quant}
  \boxed{\;
  n\,\mathcal{L}(z)=\log\!\bigl(-\beta_n(z)\bigr)+2\pi \mathrm{i}k+O\!\bigl(n^{D}\theta^n\bigr)\;}
  \end{equation}
  其中对数取任意在 \(W\) 上连续的分支（分支变更等价于 \(k\) 的平移）。
\end{enumerate}
并且若 \(z=z_{n,k}\) 取遍上述零点，则有一致渐近展开
\[
\Re \mathcal{L}(z_{n,k})
=\frac{(d_2-d_1)\log n+\log\left|\frac{p_{2,d_2}(z_{n,k})}{p_{1,d_1}(z_{n,k})}\right|}{n}
+O\!\left(\frac{\log n}{n^2}\right),
\]
\[
\Im \mathcal{L}(z_{n,k})
=\frac{2\pi k+\arg\!\left(-\frac{p_{2,d_2}(z_{n,k})}{p_{1,d_1}(z_{n,k})}\right)}{n}
+O\!\left(\frac{\log n}{n^2}\right).
\]
由于 \(K\) 上 \(\mathcal{L}'\) 不消失，\(\Re\mathcal{L}=0\) 在 \(K\) 附近给出 \(\Gamma\) 的一条光滑实解析支；
从而零点到该曲线的欧氏距离满足
\[
\mathrm{dist}(z_{n,k},\Gamma)=\Theta\!\left(\frac{\log n}{n}\right),
\]
其中常数仅依赖于 \(K\) 上的 \(\inf|\mathcal{L}'|\) 与 \(|d_2-d_1|\)。
\end{theorem}
\begin{proof}
在 \(V\) 上写
\[
P_n(z)=\alpha_1(n;z)\lambda_1(z)^n\Bigl(1+\beta_n(z)e^{-n\mathcal{L}(z)}+R_n(z)\Bigr),
\]
其中
\[
R_n(z):=\sum_{j=3}^{s}\frac{\alpha_j(n;z)}{\alpha_1(n;z)}\Bigl(\frac{\lambda_j(z)}{\lambda_1(z)}\Bigr)^n.
\]
由 \(\alpha_j(n;z)=O(n^{d_j})\) 与谱分离假设可得在 \(V\) 上一致估计
\(
\sup_{z\in V}|R_n(z)|=O(n^{D}\theta^n)
\)。
取 \(W\Subset V\) 使 \(\alpha_1,\alpha_2\) 在 \(W\) 上对充分大 \(n\) 不为零，且 \(\mathcal{L}'\) 在 \(W\) 上无零点。
\par
取 \(z\in W\) 满足 \(P_n(z)=0\)。
由于 \(\alpha_1(n;z)\lambda_1(z)^n\neq 0\) 于 \(W\) 成立，上式等价于
\[
1+\beta_n(z)e^{-n\mathcal{L}(z)}=-R_n(z).
\]
由 \(\sup_W|R_n|=O(n^{D}\theta^n)\) 得
\[
\beta_n(z)e^{-n\mathcal{L}(z)}=-1+O(n^{D}\theta^n),
\qquad
e^{n\mathcal{L}(z)}=-\beta_n(z)\bigl(1+O(n^{D}\theta^n)\bigr)^{-1}.
\]
对两侧取对数并吸收分支不定性，得到存在 \(k\in\ZZ\) 使 \eqref{eq:xi-leyang-logshift-quant} 成立。
\par
为证单零点性，记
\[
H_n(z):=1+\beta_n(z)e^{-n\mathcal{L}(z)}+R_n(z),
\qquad
P_n(z)=\alpha_1(n;z)\lambda_1(z)^n\,H_n(z).
\]
对 \(H_n\) 求导并在零点处代入 \(\beta_n e^{-n\mathcal{L}}=-1+o(1)\)，同时注意 \(\beta_n\) 及其导数至多多项式增长，而 \(e^{-n\mathcal{L}}\) 在零点处与 \(\beta_n^{-1}\) 同阶，可得
\[
H_n'(z)=\beta_n'(z)e^{-n\mathcal{L}(z)}-n\,\beta_n(z)\mathcal{L}'(z)e^{-n\mathcal{L}(z)}+R_n'(z)
=n\,\mathcal{L}'(z)+O(1),
\]
其中余项在 \(W\) 上一致；当 \(n\) 足够大时 \(H_n'(z)\neq 0\)，故零点为单零点。
\par
由多项式前因子展开
\[
\beta_n(z)=n^{d_2-d_1}\frac{p_{2,d_2}(z)}{p_{1,d_1}(z)}\Bigl(1+O\!\left(\frac1n\right)\Bigr)
\]
（一致于 \(W\)），代入 \eqref{eq:xi-leyang-logshift-quant} 并分离实部与虚部即得所示渐近展开。
距离估计由 \(\Re\mathcal{L}\) 在 \(K\) 上的横截性与 \(\Re\mathcal{L}(z_{n,k})=\Theta((\log n)/n)\) 推出。证毕。
\end{proof}

\begin{theorem}[Jordan 调制导致的对数漂移系数]\label{thm:xi-leyang-jordan-logshift}
设 \(T:U\to\Mat_d(\CC)\) 为解析矩阵族，\(u,v\in\CC^d\) 固定，定义
\[
Z_n(z):=u^{\mathsf T}T(z)^n v\qquad(n\in\NN).
\]
设存在邻域 \(V\Subset U\) 与常数 \(0<\theta<1\)，使得 \(T(z)\) 的谱可分解为两条主导谱分支 \(\lambda_1(z),\lambda_2(z)\) 与其余谱满足
\[
\max_{j\ge 3}|\lambda_j(z)|\le \theta\,\max\bigl(|\lambda_1(z)|,|\lambda_2(z)|\bigr)\qquad(z\in V),
\]
并且 \(\lambda_1/\lambda_2\) 在 \(V\) 上可取解析对数 \(\mathcal{L}=\log(\lambda_1/\lambda_2)\)，且 \(\mathcal{L}'\) 在局部等模集上不消失。
设 \(\ell_i\ge 1\) 为与 \(\lambda_i\) 相关、且在读出 \(u^{\mathsf T}(\cdot)v\) 中出现的最大 Jordan 块尺寸，并在 \(V\) 上常值。
则在 \(V\) 的适当缩小邻域内，\(Z_n\) 的零点在等模曲线附近满足
\[
\Re\mathcal{L}(z)=\frac{(\ell_2-\ell_1)\log n+\log|c_2(z)/c_1(z)|}{n}+O\!\left(\frac1n\right),
\]
其中 \(c_1,c_2\) 为解析且在等模集上不为零的系数函数。
特别地，若 \(\ell_1\neq \ell_2\)，则 \(\frac{\ell_2-\ell_1}{n}\log n\) 项在零点法向漂移中不可消去，其系数仅由 Jordan 非半单缺陷所决定。
\end{theorem}
\begin{proof}
对每个 \(z\in V\) 取 \(T(z)=S(z)J(z)S(z)^{-1}\) 的 Jordan 标准形。
对单个 Jordan 块 \(J=\lambda(I+N)\)（\(N\) 幂零、块大小 \(\ell\)）有恒等式
\[
J^n=\lambda^n\sum_{r=0}^{\ell-1}\binom{n}{r}N^r,
\]
从而 \(Z_n(z)\) 展开为有限个项的线性组合，每项形如 \(n^r\lambda_j(z)^n\)，系数解析。
按可见最大 Jordan 阶数取出 \(\lambda_1,\lambda_2\) 的首项，得到
\[
Z_n(z)=\bigl(c_1(z)\,n^{\ell_1-1}+O(n^{\ell_1-2})\bigr)\lambda_1(z)^n
\;+\;\bigl(c_2(z)\,n^{\ell_2-1}+O(n^{\ell_2-2})\bigr)\lambda_2(z)^n
\;+\;O(n^{D}\theta^n\Lambda(z)^n),
\]
其中 \(\Lambda(z)=\max(|\lambda_1(z)|,|\lambda_2(z)|)\)。
将该表示代入定理 \ref{thm:xi-leyang-logshift-two-spectrum}（取 \(d_i=\ell_i-1\)）即得结论。证毕。
\end{proof}

\begin{theorem}[谱分歧点的 Puiseux 标度与零点聚束]\label{thm:xi-leyang-puiseux-zero-scaling}
设 \(z_c\in\CC\)。
在 \(z_c\) 的穿孔邻域内设存在两条主导谱支 \(\lambda_{\pm}(z)\neq 0\)，并存在整数 \(m\ge 2\) 与常数 \(a\neq 0\) 使得当 \(z\to z_c\) 时有 Puiseux 展开
\[
\log\!\Bigl(\frac{\lambda_+(z)}{\lambda_-(z)}\Bigr)
=2a\,(z-z_c)^{1/m}+O\!\bigl((z-z_c)^{2/m}\bigr),
\]
其中 \((z-z_c)^{1/m}\) 为某一固定的局部 Puiseux 分支。
设
\[
P_n(z)=A_n(z)\lambda_+(z)^n+B_n(z)\lambda_-(z)^n+R_n(z),
\]
其中 \(A_n,B_n\) 在 \(z_c\) 的邻域内解析且关于 \(n\) 仅含有次数有界的多项式前因子，而 \(R_n\) 为严格次主导余项：存在 \(0<\theta<1\) 与 \(D\ge 0\) 使
\[
|R_n(z)|\le C\,n^D\,\theta^n\,\max(|\lambda_+(z)|,|\lambda_-(z)|)^n
\]
在某邻域内一致成立。
并假设 \(-B_n(z_c)/A_n(z_c)\) 不以指数速率趋于 \(0\) 或 \(\infty\)。
则存在 \(n_0\) 与常数 \(C_1>0\)，使得对一切 \(n\ge n_0\)，落在 \(|z-z_c|\le C_1 n^{-m}\) 内的全部零点可表为
\[
z_{n,k}
=z_c+\Bigl(\frac{1}{2a}\Bigr)^m\Bigl(\frac{\log\!\bigl(-B_n(z_c)/A_n(z_c)\bigr)+2\pi \mathrm{i}k}{n}\Bigr)^m
+O\!\left(\frac{(\log n)^{m+1}}{n^{m+1}}\right),
\qquad k\in\ZZ,
\]
其中对数取任意连续分支（分支变更等价于 \(k\) 的平移）。
特别地，零点对 \(z_c\) 的收敛标度为 \(n^{-m}\)。
\end{theorem}
\begin{proof}
忽略余项 \(R_n\) 时，零点满足
\[
\Bigl(\frac{\lambda_+(z)}{\lambda_-(z)}\Bigr)^n=-\frac{B_n(z)}{A_n(z)}.
\]
取对数并吸收分支不定性得
\[
n\,\log\!\Bigl(\frac{\lambda_+(z)}{\lambda_-(z)}\Bigr)=\log\!\Bigl(-\frac{B_n(z)}{A_n(z)}\Bigr)+2\pi\mathrm{i}k.
\]
由于 \(B_n/A_n\) 在 \(z_c\) 的邻域内解析且非零，其对数在固定分支下满足
\[
\log\!\Bigl(-\frac{B_n(z)}{A_n(z)}\Bigr)
=\log\!\Bigl(-\frac{B_n(z_c)}{A_n(z_c)}\Bigr)+O(|z-z_c|),
\]
代入并注意 \(|z-z_c|=O(n^{-m})\) 可把该误差吸收进最终的 \(O((\log n)^{m+1}/n^{m+1})\) 项。
由 Puiseux 展开存在局部坐标 \(\eta(z)=a(z-z_c)^{1/m}+O((z-z_c)^{2/m})\) 使
\(
\log(\lambda_+/\lambda_-)=2\eta(z)
\)，并且 \(z-z_c=(\eta/a)^m+O(\eta^{m+1})\)。
由于右端仅具对数级增长，可得 \(\eta(z)=O((\log n)/n)\) 并代回得到所示展开。
最后由 \(|R_n|\ll |A_n\lambda_+^n|+|B_n\lambda_-^n|\) 的一致控制与 Rouch\'e 定理可知，加入 \(R_n\) 不改变上述小邻域内零点计数与展开。证毕。
\end{proof}

\begin{corollary}[简单判别式点的二次聚束标度]\label{cor:xi-leyang-puiseux-m2-parabola}
\leanverified{paper\_xi\_leyang\_puiseux\_m2\_parabola}
在定理 \ref{thm:xi-leyang-puiseux-zero-scaling} 的设定下，若分歧指标 \(m=2\)，则临界邻域内零点满足
\[
z_{n,k}
=z_c+\frac{1}{4a^2}\Bigl(\frac{\log\!\bigl(-B_n(z_c)/A_n(z_c)\bigr)+2\pi \mathrm{i}k}{n}\Bigr)^2
+O\!\left(\frac{(\log n)^3}{n^3}\right),
\]
从而典型间距为 \(\Theta(n^{-2})\)，并在 \(z\)-平面上形成以 \(z_c\) 为端点的二次聚束族。
\end{corollary}

\begin{corollary}[Jordan 指数差的零点漂移识别]\label{cor:xi-leyang-identify-jordan-gap-from-zeros}
在定理 \ref{thm:xi-leyang-jordan-logshift} 的设定下，设 \(\gamma\) 为等模曲线的一段紧弧，且
\[
\inf_{z\in\gamma}|\mathcal{L}'(z)|>0.
\]
对任意零点 \(z=z_{n,k}\) 取其到 \(\gamma\) 的最近点投影 \(z_\gamma\in\gamma\)，并以 \(\Re\mathcal{L}\) 的符号规定有符号法向距离 \(\delta_n(z)\)。
则
\[
\delta_n(z)=\frac{\ell_2-\ell_1}{|\mathcal{L}'(z_\gamma)|}\cdot\frac{\log n}{n}+O\!\left(\frac{1}{n}\right),
\]
其中余项在上述投影落入 \(\gamma\) 的紧子弧时一致成立。
从而 Jordan 指数差可由单尺度斜率读出：
\[
\ell_2-\ell_1=\lim_{n\to\infty}\frac{n\,|\mathcal{L}'(z_\gamma)|\,\delta_n(z)}{\log n}.
\]
\end{corollary}

\begin{remark}[与谱判别式的对齐]\label{rem:xi-leyang-logshift-puiseux-discriminant-alignment}
在本文的谱方程口径中，分歧点等价于主导谱支发生重根并合，因而由判别式为零刻画（例如 \(\Pi(\lambda,y)\) 的 \(\mathrm{Disc}_\lambda\) 分解给出 Lee--Yang 边界）。
定理 \ref{thm:xi-leyang-puiseux-zero-scaling} 将该判别式层的分歧指标 \(m\) 直接译码为零点云的收敛标度 \(n^{-m}\)，从而把“分歧几何”提升为“有限体积零点”可核验的量化证书。
\end{remark}

在允许解析背景与变量混合的口径下，零点坐标将对复仿射重参数化敏感。
下述陈述给出一种射影化的消混合接口：以四点交比构造对复仿射重参数化不变的可比量，并由此得到 crossing 定位临界点的普适渐近律；同时给出由两枚有限体积零点反演 \(A(0)\) 与 \(z_0(0)\) 的闭式估计。

\begin{theorem}[Lee--Yang 零点的射影消混合不变量与交叉点渐近]\label{thm:xi-leyang-crossratio-demixing-crossing}
\leanverified{paper\_xi\_leyang\_crossratio\_demixing\_crossing}
设 \(L\in\RR_{>0}\) 为系统线性尺度，\(\tau\in\RR\) 为距临界点的温度型标度变量（或其解析重参数化），\(z\in\CC\) 为待复化的外参。
记有限体积配分函数 \(Z_L(z;\tau)\) 关于 \(z\) 的 Lee--Yang 零点为 \(\{z_j(L,\tau)\}_{j\ge 1}\subset\CC\)（允许按某一固定约定排序）。
假设存在指数 \(y_h>0,y_t>0,\omega>0\)，解析函数 \(z_0(\tau)\) 与 \(A(\tau)\neq 0\)，以及与普适类仅相关的常数序列 \(\{\zeta_j\}_{j\ge 1}\subset\CC\) 与 \(\{\kappa_j\}_{j\ge 1},\{\xi_j\}_{j\ge 1}\subset\CC\)，使得当 \(L\to\infty\) 且 \(|\tau|\) 足够小时，对每个固定 \(j\) 有展开
\begin{equation}\label{eq:xi-leyang-fss-analytic-mixing}
z_j(L,\tau)
=
z_0(\tau)
+
A(\tau)L^{-y_h}\Bigl(
\zeta_j
+
\kappa_j\,\tau L^{y_t}
+
\xi_j\,L^{-\omega}
+
O\!\bigl(\tau^2L^{2y_t}+\tau L^{y_t-\omega}+L^{-2\omega}\bigr)
\Bigr).
\end{equation}
对任意互异指标 \(i,j,k,\ell\) 定义四点交比
\begin{equation}\label{eq:xi-leyang-crossratio-def}
\mathrm{CR}_{i,j;k,\ell}(L,\tau)
:=
\frac{\bigl(z_i-z_j\bigr)\bigl(z_k-z_\ell\bigr)}{\bigl(z_i-z_k\bigr)\bigl(z_j-z_\ell\bigr)}\Bigg|_{(L,\tau)}.
\end{equation}
则成立：
\begin{enumerate}
  \item 存在仅由 \(\{\zeta_a\}\) 决定的普适常数
  \begin{equation}\label{eq:xi-leyang-crossratio-universal}
  \mathrm{CR}_{i,j;k,\ell}^{\ast}
  :=
  \frac{(\zeta_i-\zeta_j)(\zeta_k-\zeta_\ell)}{(\zeta_i-\zeta_k)(\zeta_j-\zeta_\ell)},
  \end{equation}
  使得在临界切片 \(\tau=0\) 有
  \begin{equation}\label{eq:xi-leyang-crossratio-critical-limit}
  \mathrm{CR}_{i,j;k,\ell}(L,0)=\mathrm{CR}_{i,j;k,\ell}^{\ast}+O(L^{-\omega}).
  \end{equation}
  并且 \(\mathrm{CR}_{i,j;k,\ell}\) 对任意复仿射重参数化 \(z\mapsto az+b(\tau)\)（\(a\in\CC^{\times}\)，\(b\) 解析）严格不变。
  \item 固定 \(s>1\)，并令 \(\tau_{L,s}\) 为方程
  \begin{equation}\label{eq:xi-leyang-crossratio-crossing-eq}
  \mathrm{CR}_{i,j;k,\ell}(L,\tau)=\mathrm{CR}_{i,j;k,\ell}(sL,\tau)
  \end{equation}
  在 \(\tau=0\) 的小邻域内的解（当 \(L\) 充分大且非退化条件满足时唯一）。
  则存在常数 \(C_{i,j;k,\ell}(s)\in\CC\) 使得
  \begin{equation}\label{eq:xi-leyang-crossratio-crossing-asymp}
  \tau_{L,s}
  =
  C_{i,j;k,\ell}(s)\,L^{-(y_t+\omega)}+o\!\bigl(L^{-(y_t+\omega)}\bigr).
  \end{equation}
  更具体地，令
  \begin{equation}\label{eq:xi-leyang-crossratio-BC}
  \mathcal B:=
  \Bigl(\frac{\kappa_i-\kappa_j}{\zeta_i-\zeta_j}+\frac{\kappa_k-\kappa_\ell}{\zeta_k-\zeta_\ell}
  -\frac{\kappa_i-\kappa_k}{\zeta_i-\zeta_k}-\frac{\kappa_j-\kappa_\ell}{\zeta_j-\zeta_\ell}\Bigr),
  \qquad
  \mathcal C:=
  \Bigl(\frac{\xi_i-\xi_j}{\zeta_i-\zeta_j}+\frac{\xi_k-\xi_\ell}{\zeta_k-\zeta_\ell}
  -\frac{\xi_i-\xi_k}{\zeta_i-\zeta_k}-\frac{\xi_j-\xi_\ell}{\zeta_j-\zeta_\ell}\Bigr),
  \end{equation}
  若 \(\mathcal B\neq 0\)，则
  \begin{equation}\label{eq:xi-leyang-crossratio-crossing-constant}
  C_{i,j;k,\ell}(s)=\frac{1-s^{-\omega}}{s^{y_t}-1}\cdot\frac{\mathcal C}{\mathcal B}.
  \end{equation}
\end{enumerate}
\end{theorem}
\begin{proof}
由 \eqref{eq:xi-leyang-fss-analytic-mixing} 得对任意 \(a\neq b\),
\[
z_a-z_b
=
A(\tau)L^{-y_h}\Bigl((\zeta_a-\zeta_b)+(\kappa_a-\kappa_b)\tau L^{y_t}+(\xi_a-\xi_b)L^{-\omega}+O(\cdots)\Bigr),
\]
其中共同平移 \(z_0(\tau)\) 在差分中完全消失，且公共因子 \(A(\tau)L^{-y_h}\) 在交比 \eqref{eq:xi-leyang-crossratio-def} 中整体约去，从而交比对 \(z\mapsto az+b(\tau)\) 严格不变。
令 \(\varepsilon:=\tau L^{y_t}\)、\(\delta:=L^{-\omega}\)。
将 \eqref{eq:xi-leyang-fss-analytic-mixing} 代入 \eqref{eq:xi-leyang-crossratio-def} 并在 \((\varepsilon,\delta)=(0,0)\) 处作一阶展开得到
\[
\mathrm{CR}_{i,j;k,\ell}(L,\tau)
=
\mathrm{CR}_{i,j;k,\ell}^{\ast}
\Bigl(1+\mathcal B\,\varepsilon+\mathcal C\,\delta+O(\varepsilon^{2}+\varepsilon\delta+\delta^{2})\Bigr),
\]
其中 \(\mathcal B,\mathcal C\) 如 \eqref{eq:xi-leyang-crossratio-BC}。
取 \(\tau=0\) 即得 \eqref{eq:xi-leyang-crossratio-critical-limit}。
\par
对固定 \(s>1\)，将同一展开应用于尺度 \(L\) 与 \(sL\) 并代入 \eqref{eq:xi-leyang-crossratio-crossing-eq}，约去 \(\mathrm{CR}^{\ast}\) 后主项为
\[
\mathcal B\,(1-s^{y_t})\,\varepsilon+\mathcal C\,(1-s^{-\omega})\,\delta
+O(\varepsilon^{2}+\varepsilon\delta+\delta^{2})
=0.
\]
当 \(\mathcal B\neq 0\) 且 \(L\) 充分大时，隐函数定理给出唯一解满足 \(\varepsilon=O(\delta)\)，即
\(
\tau=O(L^{-(y_t+\omega)})
\)。
由主项平衡得到
\(
\varepsilon=\frac{1-s^{-\omega}}{s^{y_t}-1}\cdot\frac{\mathcal C}{\mathcal B}\,\delta+o(\delta)
\)，
从而推出 \eqref{eq:xi-leyang-crossratio-crossing-asymp} 与 \eqref{eq:xi-leyang-crossratio-crossing-constant}。
\end{proof}

\begin{corollary}[复混合下比值法可失稳而交比法仍严格不变]\label{cor:xi-leyang-imag-ratio-unstable-crossratio-invariant}
在定理 \ref{thm:xi-leyang-crossratio-demixing-crossing} 的语境中，考虑一类复混合：实验/格点参数 \(z_{\mathrm{phys}}\) 与标度变量 \(\tau\) 的关系满足
\(
z=z_{\mathrm{phys}}+\beta\tau+O(\tau^2)
\)，其中 \(\beta\in\CC\) 一般不为实数。
令
\[
r_{m,n}(L,\tau):=\frac{\Im z_m(L,\tau)}{\Im z_n(L,\tau)}.
\]
则存在情形使 \(r_{m,n}(L,\tau)\) 在 \(\tau\to 0\) 时对 \(\Im(\beta)\) 的一阶敏感性不为零，从而其“不同体积交叉点”不再为坐标不变量；相反，任意交比 \(\mathrm{CR}_{i,j;k,\ell}(L,\tau)\) 对 \(z\mapsto z+\beta\tau\) 严格不变，且其交叉点渐近仍满足 \eqref{eq:xi-leyang-crossratio-crossing-asymp}。
\end{corollary}
\begin{proof}
取局部线性化 \(z_m(\tau)=z_m(0)+\beta\tau+o(\tau)\) 且 \(z_m(0)\notin\RR\)。
则 \(\Im z_m(\tau)=\Im z_m(0)+\Im(\beta)\tau+o(\tau)\)，从而比值 \(r_{m,n}\) 的一阶导数一般包含 \(\Im(\beta)\) 项而不消失。
另一方面，交比仅由差分构成，公共平移 \(+\beta\tau\) 被完全消去；结论遂由定理 \ref{thm:xi-leyang-crossratio-demixing-crossing} 的不变性与 crossing 渐近立即推出。
\end{proof}
\leanverified{paper\_xi\_leyang\_imag\_ratio\_unstable\_crossratio\_invariant}

\begin{proposition}[变量混合修正的 \texorpdfstring{$1$}{1}-余边界结构]\label{prop:xi-leyang-variable-mixing-coboundary}
\leanverified{paper\_xi\_leyang\_variable\_mixing\_coboundary}
在 Lee--Yang 零点的有限尺度缩放理论中，除定理 \ref{thm:xi-leyang-crossratio-demixing-crossing} 的射影化消混合接口外，文献中亦常使用比值观测量的变量混合展开来刻画首阶偏差。
为给出一个可复用的代数结构陈述，设对某一临界切片 \(\tau=\tau_c\) 与某一指数差
\(
\bar y:=y_t-y_h<0
\)
存在如下渐近表示：对任意 \(n\neq m\)，
\begin{equation}\label{eq:xi-leyang-imag-ratio-variable-mixing-expansion}
R_{nm}(L):=\frac{\Im z_n(L,\tau_c)}{\Im z_m(L,\tau_c)}
=r_{nm}\Bigl(1+D_{nm}\,L^{2\bar y}+O(L^{4\bar y})\Bigr),
\qquad L\to\infty,
\end{equation}
其中 \(r_{nm}=X_n/X_m\) 为普适常数比值，而主导混合系数满足
\begin{equation}\label{eq:xi-leyang-variable-mixing-Dnm-square-difference}
D_{nm}=-\kappa\,(Y_n^2-Y_m^2),
\qquad
\kappa:=\frac{a_{12}^2}{a_{22}^2}.
\end{equation}
则 \(D_{nm}\) 具有严格的余边界表述：存在势函数 \(U_n\) 使得
\begin{equation}\label{eq:xi-leyang-variable-mixing-Dnm-coboundary}
D_{nm}=U_n-U_m,
\qquad
U_n:=-\kappa\,Y_n^2.
\end{equation}
因而对任意互异指标 \(n,m,p,q\) 有恒等式
\begin{equation}\label{eq:xi-leyang-variable-mixing-four-term-identity}
D_{nm}+D_{pq}-D_{nq}-D_{pm}=0,
\end{equation}
并且对任意闭路 \(\gamma:n_0\to n_1\to\cdots\to n_k=n_0\)（允许重复顶点）有
\[
\sum_{j=0}^{k-1}D_{n_jn_{j+1}}=0.
\]
\end{proposition}
\begin{proof}
由 \eqref{eq:xi-leyang-variable-mixing-Dnm-square-difference} 直接取 \(U_n:=-\kappa Y_n^2\) 即得 \eqref{eq:xi-leyang-variable-mixing-Dnm-coboundary}。
四项恒等式 \eqref{eq:xi-leyang-variable-mixing-four-term-identity} 与闭路求和为零皆为望远镜消去。
\end{proof}

\begin{proposition}[普适类的零点射影指纹：有限交比在射影等价下完备]\label{prop:xi-leyang-crossratio-projective-fingerprint-complete}
设 \(\{\zeta_1,\dots,\zeta_M\}\subset\widehat{\CC}\)（\(M\ge 4\)）两两互异，并定义交比集合
\[
\mathfrak{I}_M:=\Bigl\{\mathrm{CR}^{\ast}_{i,j;k,\ell}:\ 1\le i<j<k<\ell\le M\Bigr\},
\qquad
\mathrm{CR}^{\ast}_{i,j;k,\ell}
:=
\frac{(\zeta_i-\zeta_j)(\zeta_k-\zeta_\ell)}{(\zeta_i-\zeta_k)(\zeta_j-\zeta_\ell)}.
\]
则 \(\mathfrak{I}_M\) 在射影等价（\(\mathrm{PGL}_2(\CC)\) 作用）意义下为完备不变量：若另一配置 \(\{\zeta_1',\dots,\zeta_M'\}\subset\widehat{\CC}\)（允许重标记）满足其交比集合与 \(\mathfrak{I}_M\) 相同，则存在 \(g\in\mathrm{PGL}_2(\CC)\) 使得 \(\zeta_m'=g(\zeta_m)\) 对所有 \(m\) 成立。
在定理 \ref{thm:xi-leyang-crossratio-demixing-crossing} 的解析混合口径（仅含复仿射自由度）下，\(\mathfrak{I}_M\) 因而给出对混合免疫的普适类“零点指纹”。
\end{proposition}
\begin{proof}
\(\mathrm{PGL}_2(\CC)\) 在 \(\widehat{\CC}\) 上三重可递。
取三个互异点作归一化，可令 \(g\) 把 \((\zeta_1,\zeta_2,\zeta_3)\) 送到 \((0,1,\infty)\)。
在该归一化下，任意第四点 \(\zeta_m\)（\(m\ge 4\)）由相应交比值唯一确定，从而整组点由交比集合确定到唯一的射影轨道。
对带重标记的情形，先以重标记对齐三点归一化，再用同一论证完成。证毕。
\end{proof}

\begin{theorem}[由两枚有限体积零点反演度量因子与混合位移]\label{thm:xi-leyang-invert-metric-and-shift-from-two-zeros}
在定理 \ref{thm:xi-leyang-crossratio-demixing-crossing} 的假设下，取任意 \(i\neq j\) 且 \(\zeta_i\neq\zeta_j\)，并假设指数 \(y_h\) 已由独立机制给定。
定义
\begin{equation}\label{eq:xi-leyang-invert-twozeros-def}
A_L:=L^{y_h}\frac{z_i(L,0)-z_j(L,0)}{\zeta_i-\zeta_j},
\qquad
z_{0,L}:=z_i(L,0)-A_L L^{-y_h}\zeta_i.
\end{equation}
则当 \(L\to\infty\) 时，
\begin{equation}\label{eq:xi-leyang-invert-twozeros-asymp}
A_L=A(0)+O(L^{-\omega}),\qquad
z_{0,L}=z_0(0)+O\!\bigl(L^{-(y_h+\omega)}\bigr).
\end{equation}
因而仅用两枚有限体积零点，即可在误差 \(O(L^{-\omega})\) 的意义下反演出非普适度量因子 \(A(0)\) 与混合位移 \(z_0(0)\)。
\end{theorem}
\begin{proof}
将 \eqref{eq:xi-leyang-fss-analytic-mixing} 特化到 \(\tau=0\)，并对 \(z_i-z_j\) 作差分，得
\[
z_i(L,0)-z_j(L,0)
=
A(0)L^{-y_h}\Bigl((\zeta_i-\zeta_j)+(\xi_i-\xi_j)L^{-\omega}+O(L^{-2\omega})\Bigr),
\]
代入 \eqref{eq:xi-leyang-invert-twozeros-def} 得 \(A_L=A(0)\bigl(1+O(L^{-\omega})\bigr)\)。
再将 \(z_i(L,0)=z_0(0)+A(0)L^{-y_h}\bigl(\zeta_i+\xi_iL^{-\omega}+O(L^{-2\omega})\bigr)\) 代入 \(z_{0,L}\) 的定义，得到
\(
z_{0,L}=z_0(0)+O(L^{-(y_h+\omega)})
\)。
\end{proof}

\begin{remark}[混合位移与度量归一化的可观测反演接口]\label{rem:xi-leyang-twozero-calibration-interface}
当 \(z\) 取为复逸度/复化学势等外参时，\(z_0(0)\) 可解释为“物理参数 \(\mapsto\) 标度场坐标”映射在临界点处的解析混合位移，而 \(A(0)\) 给出标度场的非普适度量归一化。
定理 \ref{thm:xi-leyang-invert-metric-and-shift-from-two-zeros} 表明二者可由两枚有限体积零点在 \(O(L^{-\omega})\) 误差内反演，因而可作为后续零点几何不变量（例如交比/比值接口）的先验校准量；在需要数值重加权或受符号问题影响的设定中，该校准可减弱对混合参数与坐标选择的敏感性。
\end{remark}

\input{thm__xi-leyang-two-scale-crossratio-fss-demixing}

\input{subsubsec__xi-hdqcd1-leyang-radius-susceptibility-cw}

\begin{conjecture}[交比常数的数域有限生成与稳定性]\label{conj:xi-leyang-crossratio-arithmetic-rigidity}
在带代数可积结构的普适类中，设标度零点配置 \(\{\zeta_j\}\) 的有限阶交比集合 \(\mathfrak{I}_M\) 定义如命题 \ref{prop:xi-leyang-crossratio-projective-fingerprint-complete}。
猜想由 \(\mathfrak{I}_M\) 生成的数域
\[
K_M:=\QQ\bigl(\mathfrak{I}_M\bigr)
\]
为有限扩张，并在 \(M\to\infty\) 时稳定于某个固定的数域 \(K_\infty\)。
若该猜想成立，则普适类的连续零点几何可被一个离散算术对象（数域及其伽罗瓦数据）所锚定，从而为“零点指纹的算术刚性”提供可核验的判别接口。
\end{conjecture}


\subsubsection{谱曲线的椭圆化与分歧单值化}\label{subsubsec:xi-spectral-elliptic-ramification-monodromy}
沿用谱四次 \(\Pi(\lambda,y)\) 及其 Lee--Yang 型分歧的判别式设定。令 \(C\) 为仿射曲线 \(\Pi(\lambda,y)=0\) 的正规化，并将 \(C\) 通过 \(y\)-投影视为 \(\mathbb P^1_y\) 上的四层覆盖。其有限分歧像由
\(
\mathrm{Disc}_{\lambda}(\Pi)=-y(y-1)P_{\mathrm{LY}}(y)
\)
的零点给出，其中
\[
P_{\mathrm{LY}}(y):=256y^3+411y^2+165y+32
\]
为 Lee--Yang 三次因子。引入坐标
\[
u:=y-\lambda^2\qquad(\text{等价地 }w:=2u+1),
\]
则同一方程在 \(\QQ\) 上双有理同构于极小椭圆曲线
\[
E_{\min}:\quad u^2+u=\lambda^3-\lambda
\]
（见定理 \ref{thm:xi-terminal-zm-leyang-elliptic-structure} 与结论 \ref{con:xi-terminal-zm-elliptic-normalization}），从而 \(C\) 为属 \(1\) 的光滑射影曲线。等价地，令
\[
X:=\lambda,\qquad Y:=y-\lambda^2+\tfrac12,
\]
则可在短 Weierstrass 模型
\(
E:\ Y^2=X^3-X+\tfrac14
\)
上实现函数域同一；其在本文约定下满足 \(\Delta(E)=37\) 且 \(j(E)=110592/37\)，并由于 \(j(E)\) 非代数整数而不具有复乘。并且由平移 \(Y\mapsto Y-\tfrac12\) 回到整系数最小模型 \(E_{\min}:u^2+u=X^3-X\)（Cremona/LMFDB 标识 \(37\mathrm a1\)，见注记 \ref{rem:xi-terminal-elliptic-37a1-modular-anchor}）。
\par
在上述椭圆化下，重量参数成为 \(E\) 上的有理函数
\(
y=X^2+Y-\tfrac12\in H^{0}\!\bigl(E,\mathcal O_{E}(4O)\bigr)
\)；
其极点除子为 \((y)_\infty=4O\)，并且两条有理分歧纤维 \(y=0\) 与 \(y=1\) 具有完全显式的主除子分解（结论 \ref{con:xi-terminal-zm-elliptic-normalization}）。因此作为次数 \(4\) 的覆盖 \(\pi_y:E\to\mathbb P^1_y\)，纤维 \(y^{-1}(0)\) 与 \(y^{-1}(1)\) 各包含唯一一个 \(\QQ\)-有理简单分歧点，分歧指数均为 \(2\)（定理 \ref{thm:xi-terminal-zm-leyang-elliptic-structure}）。
\par
有限分歧点由 \(\dd y=0\) 刻画并分解为两条有理分歧点与一条非有理三点轨道：除去两条有理解外，其余三个分歧点的 \(X\)-坐标为三次
\(
16X^3-9X^2+1=0
\)
的三根，而对应 \(Y\)-坐标由切触线性方程
\(
4XY+3X^2-1=0
\)
唯一确定（定理 \ref{thm:xi-terminal-zm-leyang-elliptic-structure}）。该三次的判别式为 \(-2^{2}\cdot 3^{3}\cdot 37\)，与 Lee--Yang 三次 \(P_{\mathrm{LY}}(y)\) 的判别式 \(-3^{9}\cdot 31^{2}\cdot 37\) 共享平方自由核 \(-111\)，从而二次伴随域 \(\QQ(\sqrt{-111})\) 同时支配“非有理分歧点的三次坐标域”与“非有理分歧值的三次坐标域”（参见结论 \ref{con:xi-terminal-zm-leyang-cubic-arithmetic}），并构成该谱覆盖的稳定算术指纹。
\par
将 \(\Pi(\lambda,y)\) 视为 \(\QQ(y)\) 上关于 \(\lambda\) 的四次多项式，则其泛型伽罗瓦群与几何单值群为 \(S_4\simeq W(A_3)\)（结论 \ref{con:xi-terminal-zm-generic-galois-s4} 与定理 \ref{thm:xi-terminal-zm-leyang-elliptic-structure}），从而覆盖的几何单值化达到四叶覆盖所允许的最大非交换对称。以 Hitchin--cameral 的观点看，该 Weyl 单值化可作为秩 \(4\) 光谱覆盖的结构性数据，使 Lee--Yang 分歧层、Hurwitz 护照与固定椭圆曲线在同一代数对象上可计算对齐；在几何 Langlands 的物理实现中，这类谱/cameral 数据亦可作为 Hitchin 纤维化下的最小可审计原型之一。\cite{KapustinWitten2006GeometricLanglands}
\par
与该 \(y\)-投影的四层覆盖相对，\(\Pi\) 作为关于 \(y\) 的二次多项式满足 \(\mathrm{Disc}_{y}(\Pi)=4\lambda^3-4\lambda+1\)，其零点与椭圆曲线的 \(2\)-division 三次一致，故 \(y\)-方向的二次分歧由 \(E[2]\) 的坐标域精确控制。与此相并行，二阶碰撞三次
\(
P_2(\Lambda)=\Lambda^3-2\Lambda^2-2\Lambda+2
\)
所生成的三次数域与椭圆曲线 \(E\) 的 \(2\)-division 三次（非平凡 \(E[2]\) 的 \(X\)-坐标域）在 \(\QQ\) 上同构，且可由显式的二次 Tschirnhaus 变换互相恢复（结论 \ref{con:xi-terminal-zm-collision-s2-e2-field-isomorphism}）。因此“碰撞动力学的三次骨架”与“谱椭圆曲线的模 \(2\) 结构”在算术层面被识别为同一三次域对象。
\par
\input{conj__xi-langlands-leyang-hankel-conductor-ramification}

对良素 \(p\notin\{2,37\}\)，令
\(
N_p(y):=\card{\{\lambda\in\FF_p:\ \Pi(\lambda,y)=0\}}
\)。
则椭圆化给出严格点计数恒等式
\(
\#E(\FF_p)=1+\sum_{y\in\FF_p}N_p(y)
\)，从而 Frobenius 迹
\(
a_p:=p+1-\#E(\FF_p)
\)
满足
\(
\sum_{y\in\FF_p}(N_p(y)-1)=-a_p
\)
（结论 \ref{con:xi-terminal-zm-fiber-root-counting-hecke} 与结论 \ref{con:xi-terminal-zm-hecke-eigenvalue-fiber-deviation}）；在模性锚定下 \(a_p\) 同时等于导子 \(37\) 新形式的 Hecke 本征值。进一步，对驯约化素数 \(p\notin\{2,3,31,37\}\)，函数域 Chebotarev 在固定单值群 \(\mathrm{Mon}(\pi_y)\cong S_4\) 上给出有效的 Chebotarev--Sato--Tate 型等分布：未分歧参数的纤维根数 \(N_p(y)\) 的极限比例由 \(S_4\) 在四点作用下的不动点分布决定，并具有 \(O(\sqrt p)\) 的平方根误差；其中
\(
\pi_0=\tfrac38,\ \pi_1=\tfrac13,\ \pi_2=\tfrac14,\ \pi_4=\tfrac1{24}
\)
且 \(N_p(y)=3\) 仅可能发生在分歧集合的模 \(p\) 化上（结论 \ref{con:xi-terminal-zm-s4-modp-fiber-rootcount-fixedpoints}）。由此，谱四次在有限域上的分解统计、Hecke 迹与单值群的共轭类分布在同一椭圆覆盖框架中被同时锁定。




\begin{theorem}[Lee--Yang 特征谱曲线的椭圆化与单值护照：固定椭圆曲线与 \(S_4\) 分支结构]\label{thm:xi-terminal-zm-leyang-elliptic-structure}
沿用谱四次多项式
\[
\Pi(\lambda,y):=\lambda^4-\lambda^3-(2y+1)\lambda^2+\lambda+y(y+1)\in\ZZ[\lambda,y].
\]
记仿射谱曲线 \(\mathcal C:=\{(\lambda,y)\in\mathbb A^2:\Pi(\lambda,y)=0\}\)。
定义平面多项式自同构
\[
T:\mathbb A^2\to\mathbb A^2,\qquad (\lambda,y)\longmapsto(x,u):=(\lambda,y-\lambda^2),
\qquad
T^{-1}:(x,u)\longmapsto(\lambda,y)=(x,u+x^2).
\]
则直接代换给出
\[
\Pi(\lambda,y)=u^2+u-(x^3-x)\qquad(x=\lambda,\ u=y-\lambda^2),
\]
从而 \(T\) 将 \(\mathcal C\) 精确送到极小椭圆曲线
\[
E_{\min}:\quad u^2+u=x^3-x,
\]
并在 \(\QQ\) 上诱导仿射曲线同构 \(\mathcal C\cong E_{\min}\)。
因此 \(\Pi(\lambda,y)=0\) 的规范化为属 \(1\) 的光滑射影曲线，并在 \(\QQ\) 上与该极小椭圆曲线同构；其分支结构因而由一条与参数 \(y\) 无关的固定有理椭圆曲线所锚定，而非任意四次覆盖的偶然现象。
等价地，若以 \(U:=\lambda\)、\(V:=y-\lambda^{2}\) 作为谱坐标，则归一化模型亦可写为
\[
V^{2}+V=U^{3}-U.
\]
为与短 Weierstrass 口径对齐，令
\[
w:=2u+1=2y-2\lambda^2+1\quad(\text{亦记 }S:=w),
\]
则恒有恒等式
\[
4\Pi(\lambda,y)=w^{2}-\bigl(4x^{3}-4x+1\bigr)=w^{2}-4x^{3}+4x-1.
\]
将 \(\Pi(\lambda,y)\) 视为关于 \(y\) 的二次多项式，则其 \(y\)-判别式满足
\(
\mathrm{Disc}_{y}(\Pi)=4\lambda^{3}-4\lambda+1
\)，
上式即等价写为 \(\,w^{2}=\mathrm{Disc}_{y}(\Pi)\,\) 在谱曲线 \(\mathcal C\) 上的实现。
从而 \((\lambda,y)\mapsto(x,w)\) 给出 \(\mathcal C\) 与短 Weierstrass 椭圆曲线
\[
\mathcal E:\quad w^{2}=4x^{3}-4x+1\qquad(x=\lambda)
\]
之间的显式仿射同构；其逆映射由
\[
y=x^{2}-\frac12+\frac{w}{2}\qquad\Bigl(\text{亦即 }y=\lambda^{2}-\frac12+\frac{w}{2}\Bigr)
\]
给出。
在该极小模型下，非平凡 \(2\)-挠点满足 \(u=-\tfrac12\)（等价 \(w=0\)），其 \(x\)-坐标由
\[
4x^{3}-4x+1=0
\]
给出；因此 \(\mathrm{Disc}_{y}(\Pi)=4\lambda^{3}-4\lambda+1\) 恰为 \(E_{\min}\) 的 \(2\)-除法多项式。
等价地，自然投影 \(\pi_x:E_{\min}\to\mathbb P^1_x\) 为次数 \(2\) 的覆盖，其分歧点的 \(x\)-坐标恰为三次方程 \(4x^{3}-4x+1=0\) 的三根，并与 \(E_{\min}(\overline{\QQ})\) 的非平凡 \(2\)-挠点一致。
将 \(E_{\min}\) 配方为 \((2u+1)^2=4x^3-4x+1\) 并以 \(w=2u+1\) 记之，则得到等价的短 Weierstrass 模型
\[
\mathcal E:\quad w^2=4x^3-4x+1\qquad(g_2=4,\ g_3=-1),
\]
其判别式仍为 \(\Delta(\mathcal E)=g_2^3-27g_3^2=37\)。
进一步以伸缩 \(X:=4x,\ Y:=4w\) 得整系数短 Weierstrass 模型
\[
Y^2=X^3-16X+16,
\]
其在 Cremona/LMFDB 的标准标识下对应同源类 \(37\mathrm a\) 的强 Weil 曲线 \(37\mathrm a1\)（注记 \ref{rem:xi-terminal-elliptic-37a1-modular-anchor}）。
等价地，谱支 \(\lambda(y)\) 可视为椭圆曲线 \(E_{\min}\) 上有理函数 \(y=x^2+u=x^2-\tfrac12+\tfrac{w}{2}\) 的局部反演，从而原先的多分支谱反演被提升为椭圆曲线上的有理函数反演问题。
在复解析层面，\(\mathcal E(\CC)\cong \CC/\Lambda\) 的 Weierstrass 一致化可取为
\[
x=\wp(z;\,g_2=4,g_3=-1),\qquad w=\wp'(z;\,g_2=4,g_3=-1),\qquad
y=\wp(z)^2-\frac12+\frac{\wp'(z)}{2},
\]
因而 Lee--Yang 分析中的四支代数谱函数 \(\{\lambda_j(y)\}_{j=1}^{4}\) 可理解为椭圆函数 \(z\mapsto\wp(z)\) 的代数反演之分支（见注记 \ref{rem:xi-terminal-zm-lambda-plus-wp-uniformization}）。
该椭圆曲线满足
\[
c_4(E_{\min})=48,\qquad c_6(E_{\min})=-216,\qquad
\Delta(E_{\min})=37,\qquad j(E_{\min})=\frac{c_4(E_{\min})^{3}}{\Delta(E_{\min})}=\frac{110592}{37},
\]
从而仅在素数 \(37\) 处坏约化且导子为 \(37\)；由此 Lee--Yang 型解析边界可在同一算术椭圆对象上与周期格/模性锚点统一表述。
由椭圆曲线模性定理，存在权 \(2\)、level \(37\) 的新形式 \(f\) 使 \(L(E_{\min},s)=L(f,s)\)；
其外部命名与交叉核验接口见注记 \ref{rem:xi-terminal-elliptic-37a1-modular-anchor}。
\par
令 \(\pi_y:E_{\min}\to\mathbb P^1_y\) 为有理函数
\(
y=x^2+u=\frac{w+2x^2-1}{2}
\)
所诱导的次数 \(4\) 态射，则 \(y\) 在无穷远点 \(O\in E_{\min}\) 处具有唯一四阶极点，因而 \(\infty\in\mathbb P^1_y\) 为全分歧分歧值。
等价地，绕 \(y=\infty\) 的几何单值化为一条 \(4\)-循环（其局部一致化参数与 \(\lambda\) 的 Puiseux 展开见定理 \ref{thm:xi-terminal-zm-puiseux-infty-4cycle}）。
并且在 \((x,w)\)-模型 \(\mathcal E:w^2=4x^3-4x+1\) 上，两个有理分歧值的纤维除子满足完全分解
\[
(y)_0=[(0,1)]+[(-1,-1)]+2[(1,-1)],
\qquad
(y-1)_0=[(2,-5)]+[(1,1)]+2[(-1,1)],
\]
其中点以 \((x,w)\)-坐标表示。
因此作为次数 \(4\) 的覆盖 \(\pi_y\)，纤维 \(y^{-1}(0)\) 与 \(y^{-1}(1)\) 各包含唯一的 \(\QQ\)-有理简单分歧点（对应于上式中系数为 \(2\) 的点），且分歧指数均为 \(2\)。
其有限临界点由 \(\dd y=0\) 刻画。

\begin{proposition}[分歧除子与 Abel--Jacobi 零和约束]\label{prop:xi-terminal-zm-leyang-ramification-divisor-stokes}
沿用定理的记号，将 \(\pi_y:E_{\min}\to\PP^1_y\) 视为次数 \(4\) 的有理映射，并令其分歧除子为
\[
R_{\pi_y}:=\sum_{P\in E_{\min}(\CC)}(e_P-1)\,[P],
\]
其中 \(e_P\) 为 \(\pi_y\) 在 \(P\) 处的分歧指数。
则在除子类群 \(\mathrm{Pic}(E_{\min})\) 中成立
\[
R_{\pi_y}\sim 2\,\pi_y^{-1}(\infty).
\]
由于 \(\pi_y^{-1}(\infty)=4O\)，从而
\[
R_{\pi_y}\sim 8O,\qquad \deg R_{\pi_y}=8.
\]
并且在识别 \(\mathrm{Pic}^0(E_{\min})\cong E_{\min}\) 的群律下，分歧点的加和满足严格零和约束
\[
\sum_{P\in E_{\min}(\CC)}(e_P-1)\,P=O.
\]
\end{proposition}
\leanverified{paper\_xi\_terminal\_zm\_leyang\_ramification\_divisor\_stokes}
\begin{proof}
取 \(\PP^1\) 的仿射坐标 \(z\)（无穷远点为 \(\infty\)），则 \((\dd z)=-2(\infty)\) 作为除子恒成立。
对非恒定有理函数 \(f\) 的一般公式给出
\[
(\dd f)=f^\ast(\dd z)+R_f.
\]
取 \(f=\pi_y\) 即得
\[
(\dd \pi_y)=-2\,\pi_y^{-1}(\infty)+R_{\pi_y}.
\]
另一方面，\(E_{\min}\) 为属 \(1\) 曲线，其典范丛平凡；任一亚纯 \(1\)-形式的除子为主除子类。
故 \((\dd \pi_y)\) 在 \(\mathrm{Pic}(E_{\min})\) 中为零，从而 \(R_{\pi_y}-2\pi_y^{-1}(\infty)\) 为主除子，得到 \(R_{\pi_y}\sim 2\pi_y^{-1}(\infty)\)。
再由 \(\pi_y^{-1}(\infty)=4O\) 得 \(R_{\pi_y}\sim 8O\) 与 \(\deg R_{\pi_y}=8\)。
将主除子在 \(\mathrm{Pic}^0(E_{\min})\cong E_{\min}\) 下对应为群和 \(O\)，即得到分歧点零和约束。证毕。
\end{proof}

将短 Weierstrass 模型 \(\mathcal E:w^2=4x^3-4x+1\) 以 \(X=x,\ Y=w/2\) 配方为
\(
Y^2=X^3-X+\tfrac14
\)，则
\(
y=X^2+Y-\tfrac12
\)。
故对任意 \(c\in\mathbb P^1\)，纤维 \(y=c\) 等价于抛物线
\[
Y=-X^2+\Bigl(c+\frac12\Bigr)
\]
与该椭圆曲线的交点除子，有限分歧点即为该抛物线族与椭圆曲线的切触点。
由 \(2Y\,\dd Y=(3X^2-1)\dd X\) 与 \(\dd y=2X\dd X+\dd Y\) 得
\[
\dd y=\frac{4XY+3X^2-1}{2Y}\,\dd X,
\]
从而有限临界点由 \(4XY+3X^2-1=0\) 刻画；在 \((x,w)\)-模型 \(\mathcal E\) 上等价写为
\[
2xw+3x^2-1=0
\qquad\bigl(\text{等价地 }4\lambda w+6\lambda^2-2=0,\ \lambda=x\bigr),
\]
并恰分为两类：
\begin{itemize}
  \item 两个有理临界点
  \[
  P_0:\ (x,u)=(1,-1)\mapsto y=0,\qquad P_1:\ (x,u)=(-1,0)\mapsto y=1;
  \]
  \item 三个非有理临界点 \(P_2,P_3,P_4\)，其 \(x\)-坐标为三次方程
  \[
  16x^3-9x^2+1=0
  \]
  的三根；相应的 \(u\)-坐标由
  \[
  u=\frac{1-3x^2-2x}{4x}
  \]
  唯一确定，并因此
  \[
  w=2u+1=-\frac{3x^2-1}{2x};
  \]
  从而临界值满足显式表达式
  \[
  y=x^2+u=\frac{4x^3-3x^2-2x+1}{4x}
  =\frac{(x-1)(4x^2+x-1)}{4x}
  =x^{2}-\frac34x-\frac12+\frac{1}{4x}.
  \]
\end{itemize}
并且消元恒等式给出
\[
\mathrm{Res}_x\!\left(16x^3-9x^2+1,\ 4xy-\bigl(4x^3-3x^2-2x+1\bigr)\right)=-64\,P_{\mathrm{LY}}(y),
\]
由结式的乘法性亦可将两条有理分歧值 \(y=0,1\) 同时并入消元，从而得到全有限分歧值集合的证书
\[
\mathrm{Res}_x\!\Bigl((x-1)(x+1)\bigl(16x^3-9x^2+1\bigr),\ 4xy-(x-1)(4x^2+x-1)\Bigr)
=1024\,y(y-1)P_{\mathrm{LY}}(y).
\]
从而上述参数化与判别式分解中的 Lee--Yang 三次因子严格一致。
上述三个非平凡有限临界值恰为不可约三次多项式
\[
P_{\mathrm{LY}}(y):=256y^3+411y^2+165y+32
\]
的三根，其中唯一负实根即 \(y_{\mathrm{LY}}\)。
记其余两根为共轭对 \(y_{\pm}\)，并记三次方程 \(16\lambda^3-9\lambda^2+1=0\) 的根为
\(\lambda_{\mathrm{LY}}\in\RR\) 与共轭对 \(\lambda_{\pm}\)（其中 \(y_{\pm}=y(\lambda_{\pm})\)）。
数值上，
\[
y_{\mathrm{LY}}\approx -1.13445200300831188727,\qquad
y_{\pm}\approx -0.235508373495844\pm 0.233925552127514\,i,
\]
\[
\lambda_{\mathrm{LY}}\approx -0.27343481065245424608,\qquad
\lambda_{\pm}\approx 0.417967405326227\pm 0.232114033943254\,i.
\]
因此 \(\pi_y\) 的有限分支值集合恰为
\[
y\in\{0,1\}\ \cup\ \{P_{\mathrm{LY}}(y)=0\},
\]
并且在每个有限分支值上均仅出现一处简单分歧（惯性生成元为换位）。
并且判别式满足
\[
\Delta(\mathcal E)=37,
\qquad
\mathrm{Disc}\bigl(4x^3-4x+1\bigr)=2^{4}\cdot 37,
\qquad
\mathrm{Disc}(16x^3-9x^2+1)=-2^{2}\cdot 3^{3}\cdot 37,
\qquad
\mathrm{Disc}\bigl(P_{\mathrm{LY}}(y)\bigr)=-3^{9}\cdot 31^{2}\cdot 37,
\]
因此 \(\mathrm{Disc}(16x^3-9x^2+1)\) 与 \(\mathrm{Disc}\bigl(P_{\mathrm{LY}}(y)\bigr)\) 的平方自由部分同为 \(-111=-3\cdot 37\)，两者分裂域共享同一二次分辨子域 \(\QQ(\sqrt{-111})\)。
该虚二次子域因而可视为该四次谱覆盖的内禀算术指纹，并与所选椭圆化坐标归一化无关。
从而素数 \(37\) 同时出现在临界点与临界值多项式的平方自由部分中，并与 \(\Delta(E_{\min})=37\) 的坏约化素点对齐。
因而判别式分解所识别的分歧集合可被解释为 \(\pi_y\) 的临界值结构，而非 \(\ZZ[y]\) 中的偶发因式分裂。
\par
对任一有限临界值 \(y_*\in\{0,1\}\cup\{P_{\mathrm{LY}}(y)=0\}\)，记 \(\lambda_*\) 为 \(\Pi(\lambda,y_*)=0\) 的并合谱值（即相应二重根）。
则在 \(y\to y_*\) 时并合的两条谱支具有平方根型 Puiseux 展开
\[
\lambda(y)=\lambda_*\pm \kappa_*(y-y_*)^{1/2}+\beta_*(y-y_*)\pm\gamma_*(y-y_*)^{3/2}+O\bigl((y-y_*)^{2}\bigr),
\]
其中首项系数满足显式代数公式
\[
\kappa_*^{2}
=-\,\frac{2\,\partial_y\Pi(\lambda_*,y_*)}{\partial_{\lambda}^{2}\Pi(\lambda_*,y_*)}
=\frac{3\lambda_*^{2}-1}{8\lambda_*^{3}-3\lambda_*^{2}-1}.
\]
在非平凡 Lee--Yang 分歧支（即 \(\lambda_*\) 满足 \(16\lambda_*^3-9\lambda_*^2+1=0\)）上，分母可化简为
\(
8\lambda_*^{3}-3\lambda_*^{2}-1=\frac{3}{2}(\lambda_*^{2}-1)
\)，因此
\[
\kappa_*^{2}=\frac{2(3\lambda_*^{2}-1)}{3(\lambda_*^{2}-1)}.
\]
并且线性项与三阶项系数可由更高阶 Taylor 系数显式给出
\[
\beta_*
=-\,\frac{\partial_{\lambda y}\Pi(\lambda_*,y_*)}{\partial_{\lambda}^{2}\Pi(\lambda_*,y_*)}
+\frac{\partial_{\lambda}^{3}\Pi(\lambda_*,y_*)\,\partial_y\Pi(\lambda_*,y_*)}{3\bigl(\partial_{\lambda}^{2}\Pi(\lambda_*,y_*)\bigr)^{2}}.
\]
在非平凡 Lee--Yang 分歧支上进一步可化简为
\[
\beta_*=-\frac{\lambda_*(3\lambda_*^{2}+16\lambda_*+5)}{18(\lambda_*^{2}-1)^{2}},
\qquad
\gamma_*=-\frac{\lambda_*^{2}\bigl(1152\lambda_*^{7}-1152\lambda_*^{6}-280\lambda_*^{5}+723\lambda_*^{4}-312\lambda_*^{3}+30\lambda_*^{2}+16\lambda_*-5\bigr)}{2\bigl(8\lambda_*^{3}-3\lambda_*^{2}-1\bigr)^{4}}\cdot \kappa_*^{-1}.
\]
特别地，在 \(y_*=0\) 与 \(y_*=1\) 处分别得到
\[
\lambda(y)=1\pm \sqrt{\frac{y}{2}}+O(y),
\qquad
\lambda(y)=-1\pm \frac{i}{\sqrt6}(y-1)^{1/2}+O(y-1),
\]
而在实 Lee--Yang 边界 \(y_*=y_{\mathrm{LY}}\) 处有
\[
y_{\mathrm{LY}}\approx -1.13445200300831188727,\qquad
\lambda_*=\lambda_{\mathrm{LY}}\approx-0.27343481065245424608,\qquad
\kappa_{\mathrm{LY}}:=\kappa_*\approx 0.74761098459045973850,\qquad
\beta_{\mathrm{LY}}:=\beta_*\approx 0.0150716815102551,\qquad
\gamma_{\mathrm{LY}}:=\gamma_*\approx -0.0418222862574037,\qquad
\kappa_{\mathrm{LY}}^{2}\approx 0.55892218428031662857,\qquad
\frac{\kappa_{\mathrm{LY}}}{\lambda_{\mathrm{LY}}}\approx -2.73414706345016524243,\qquad
\frac{\lambda_{\mathrm{LY}}^{2}}{\kappa_{\mathrm{LY}}^{2}}\approx 0.1337692397606565,\qquad
\pi^{2}\frac{\lambda_{\mathrm{LY}}^{2}}{\kappa_{\mathrm{LY}}^{2}}\approx 1.3202494774721529.
\]
从而任何由该并合谱支构造的解析量（例如选定分支的 \(\log\lambda\)）在 \(y_{\mathrm{LY}}\) 处均呈严格平方根型 Lee--Yang 奇性，其幅度由 \(\kappa_{\mathrm{LY}}\) 在代数上完全锁定。
\par
映射 \(\pi_y\) 的分歧型为五个有限简单分歧与一处无穷远全分歧（Hurwitz 护照 \((2,2,2,2,2,4)\)），故其几何单值群包含一个 \(4\)-循环与一个换位，从而
\par
该 Hurwitz 护照在 Nielsen 类意义下提供覆盖在 Hurwitz 空间中的离散同伦不变量，使分歧循环型可作为跨模型比较的稳定指纹。
\[
\mathrm{Mon}(\pi_y)\cong S_4.
\]
等价地，函数域扩张 \(\QQ(y)\subset \QQ(y,\lambda)\) 的伽罗瓦闭包之伽罗瓦群为 \(S_4\)。
并且谱四次的 \(\lambda\)-判别式满足
\[
\mathrm{Disc}_{\lambda}\bigl(\Pi(\lambda,y)\bigr)=-y(y-1)P_{\mathrm{LY}}(y),
\]
从而该 \(S_4\)-扩张的唯一指数 \(2\) 正规子域由判别式平方根生成：
\[
\QQ(y)\bigl(\sqrt{\mathrm{Disc}_{\lambda}(\Pi(\lambda,y))}\bigr)
=\QQ(y)\Bigl(\sqrt{-y(y-1)P_{\mathrm{LY}}(y)}\Bigr).
\]
该二次子扩张亦可等价表述为 \(A_4\triangleleft S_4\) 的不动域，并对应于属 \(2\) 的超椭圆曲线
\[
C:\quad \zeta^{2}=-y(y-1)P_{\mathrm{LY}}(y),
\]
其分歧点为 \(y\in\{0,1\}\cup\{P_{\mathrm{LY}}(y)=0\}\cup\{\infty\}\) 六点。
记 \(\widetilde{C}\to\mathbb P^1_y\) 为该四次覆盖的 \(S_4\)-伽罗瓦闭包，则由惰性群阶
\[
e_\infty=4,\qquad e_0=e_1=e_{y_1}=e_{y_2}=e_{y_3}=2
\qquad\bigl(P_{\mathrm{LY}}(y_j)=0\bigr)
\]
与 Riemann--Hurwitz 公式得到
\[
g(\widetilde{C})
=1+\frac{|S_4|}{2}\Bigl(-2+\Bigl(1-\frac14\Bigr)+5\Bigl(1-\frac12\Bigr)\Bigr)
=16.
\]
作为等价的代数判别，可取 \(\Pi(\lambda,y)\) 的一条标准 resolvent cubic
\[
R_y(z)=64 z^3-(176+256y) z^2+(76+128y)z-(64y^2-48y+9),
\]
其满足 \(\mathrm{Disc}_z(R_y)=2^{24}\mathrm{Disc}_{\lambda}(\Pi)\)，并在 \(\QQ(y)[z]\) 中不可约；结合 \(\mathrm{Disc}_{\lambda}(\Pi)\) 在 \(\QQ(y)\) 中非平方，即得 \(\mathrm{Gal}(\Pi/\QQ(y))\cong S_4\)。
因此任何试图将其分歧结构压缩到 \(A_4\)、\(D_4\) 或 \(V_4\) 等较小单值群的策略均与分歧循环型不相容。
由于 \(S_4\cong W(A_3)\)，上述单值化同时给出一条由简单分歧换位（局部反射）与无穷远全分歧的 \(4\)-循环（Coxeter 元）所生成的 \(A_3\)–Weyl 群最小实现。
该最大单值性在代数层面表明谱四次覆盖达到最不退化的对称复杂度，并为谱支的全局解析延拓、剪切线选择与零点凝聚结构的不可约性分析提供群论骨架。
另一方面，由标准同构 \(S_4\simeq \mathrm{PGL}_2(\FF_3)\) 可将其识别为八面体有限单值群像；
在去除分歧值的开曲线 \(U:=\mathbb P^1_y\setminus\{0,1,P_{\mathrm{LY}}(y)=0,\infty\}\) 上，
上述单值化等价给出一个有限像的 \(\mathrm{PGL}_2(\FF_3)\)-投影局部系统，其存在性与分歧集合完全由 \(\Pi(\lambda,y)\) 的判别式机制所固定。
\leanverified{paper\_xi\_terminal\_zm\_leyang\_elliptic\_structure}
\end{theorem}

\par
\input{thm__xi-terminal-zm-leyang-linear-twist-quartic-discriminant-s5}

\begin{corollary}[判别式二次分辨率在椭圆归一化上的基变：属 \(6\) 的双覆盖]\label{cor:xi-terminal-zm-leyang-discriminant-square-root-basechange-genus6}
沿用定理 \ref{thm:xi-terminal-zm-leyang-elliptic-structure} 的椭圆曲线 \(E_{\min}\) 与次数 \(4\) 覆盖 \(\pi_y:E_{\min}\to\PP^{1}_{y}\)。
令
\[
D:=-y(y-1)P_{\mathrm{LY}}(y)\in \QQ(E_{\min})^{\times}.
\]
记 \(C_{E}\) 为函数域扩张
\(
\QQ(C_{E})=\QQ(E_{\min})(\sqrt{D})
\)
所对应的光滑射影曲线，并令 \(\pi:C_{E}\to E_{\min}\) 为相应的次数 \(2\) 覆盖。
则：
\begin{enumerate}
  \item \(\pi\) 的分歧支撑由 \(D\) 的奇阶零极点精确刻画；并且 \(\pi\) 在无穷远点 \(O\in E_{\min}\) 上方不分歧；
  \item \(\pi\) 的分歧点数为 \(10\)，其支撑由下述点组成：
  \begin{itemize}
    \item 纤维 \(y^{-1}(0)\) 中除去唯一简单分歧点之外的两点；
    \item 纤维 \(y^{-1}(1)\) 中除去唯一简单分歧点之外的两点；
    \item 对每个满足 \(P_{\mathrm{LY}}(y_0)=0\) 的分歧值 \(y_0\)，纤维 \(y^{-1}(y_0)\) 中除去唯一简单分歧点之外的两点（共 \(3\times 2\) 点）。
  \end{itemize}
  \item \(C_{E}\) 的亏格为 \(g(C_{E})=6\)。
\end{enumerate}
\end{corollary}
\begin{proof}
由 \((y)_\infty=4O\) 得 \(\mathrm{ord}_{O}(y)=-4\)，从而 \(\mathrm{ord}_{O}(y-1)=-4\) 且 \(\mathrm{ord}_{O}(P_{\mathrm{LY}}(y))=-12\)，故
\(
\mathrm{ord}_{O}(D)=-4-4-12=-20
\)
为偶数，\(\pi\) 在 \(O\) 上方不分歧。
\par
对有限点，定理 \ref{thm:xi-terminal-zm-leyang-elliptic-structure} 表明：\(y=0\)、\(y=1\) 以及 \(P_{\mathrm{LY}}(y)=0\) 的三根均为 \(\pi_y\) 的有限分歧值，且每个有限分歧值之上恰出现一处简单分歧点（\(\mathrm{ord}(y-y_0)=2\)），其余两点为非分歧点（\(\mathrm{ord}(y-y_0)=1\)）。
因此在每条分歧纤维中，因子 \(y-y_0\) 的奇阶零点恰为“除去分歧点之外的两点”，从而 \(D\) 的奇阶零点总数为
\(
2+2+3\cdot 2=10
\)，并与陈述一致。
\par
对次数 \(2\) 覆盖应用 Riemann--Hurwitz 公式：
\[
2g(C_{E})-2
=2(2g(E_{\min})-2)+10
=10,
\]
由于 \(g(E_{\min})=1\)，得到 \(g(C_{E})=6\)。证毕。
\end{proof}
\begin{corollary}[有理特化的谱可约性：由椭圆曲线像集刻画]\label{cor:xi-terminal-zm-pi-rational-root-specialization-elliptic-image}
\leanverified{paper\_xi\_terminal\_zm\_pi\_rational\_root\_specialization\_elliptic\_image}
设 \(y_0\in\QQ\)。则以下命题等价：
\begin{enumerate}
  \item 四次多项式 \(\Pi(\lambda,y_0)\in\QQ[\lambda]\) 在 \(\QQ\) 上存在有理根；
  \item 存在 \(P\in E_{\min}(\QQ)\) 使得 \(\pi_y(P)=y_0\)，并且该有理根可取为 \(\lambda_0=x(P)\in\QQ\)。
\end{enumerate}
\end{corollary}
\begin{proof}
由定理 \ref{thm:xi-terminal-zm-leyang-elliptic-structure} 的仿射同构 \(\mathcal C\cong E_{\min}\) 直接得到。
若 \(\Pi(\lambda_0,y_0)=0\) 且 \(\lambda_0\in\QQ\)，取 \(u_0:=y_0-\lambda_0^2\)，则 \((x,u)=(\lambda_0,u_0)\in E_{\min}(\QQ)\) 且 \(\pi_y(x,u)=y_0\)。
反之，若 \(P\in E_{\min}(\QQ)\) 满足 \(\pi_y(P)=y_0\)，则 \((\lambda,y)=(x(P),\pi_y(P))\) 给出 \(\Pi(\lambda,y)=0\) 的一组有理解，从而 \(\Pi(\lambda,y_0)\) 存在有理根。证毕。
\end{proof}

\begin{corollary}[有理特化的异常集合之薄集刻画：有理根与符号二次分辨率]\label{cor:xi-terminal-zm-specialization-thin-exceptional-geometry}
沿用定理 \ref{thm:xi-terminal-zm-leyang-elliptic-structure} 的四次谱多项式 \(\Pi(\lambda,y)\) 及 Lee--Yang 三次因子
\[
P_{\mathrm{LY}}(y):=256y^{3}+411y^{2}+165y+32.
\]
令判别式二次脊曲线为
\[
C:\quad \zeta^{2}=-y(y-1)P_{\mathrm{LY}}(y).
\]
则对任意 \(y_0\in\QQ\)：
\begin{enumerate}
  \item 若 \(\Pi(\lambda,y_0)\) 在 \(\QQ[\lambda]\) 上存在有理根，则 \(y_0\in \pi_y\bigl(E_{\min}(\QQ)\bigr)\)（推论 \ref{cor:xi-terminal-zm-pi-rational-root-specialization-elliptic-image}）。
  \item 设 \(\Pi(\lambda,y_0)\) 在 \(\QQ[\lambda]\) 上可分且不可约。
  则
  \[
  \mathrm{Gal}\bigl(\Pi(\cdot,y_0)/\QQ\bigr)\subseteq A_{4}
  \quad\Longleftrightarrow\quad
  -y_0(y_0-1)P_{\mathrm{LY}}(y_0)\in \QQ^{\times 2}
  \quad\Longleftrightarrow\quad
  \exists\,\zeta_0\in\QQ\ \text{使}\ (y_0,\zeta_0)\in C(\QQ).
  \]
\end{enumerate}
因此，“有理可约性”与“符号二次分辨率退化（\(A_4\) 型）”两类异常机制分别被 \(\pi_y\bigl(E_{\min}(\QQ)\bigr)\) 与 \(C(\QQ)\) 的 \(y\)-投影显式几何化；
并且由 Hilbert 不可约定理可知：除去一个薄集参数族后，仍存在无穷多个 \(y_0\in\QQ\) 使得 \(\Pi(\lambda,y_0)\) 在 \(\QQ[\lambda]\) 上不可约且其伽罗瓦群为 \(S_4\)。
\end{corollary}
\begin{proof}
第一条即推论 \ref{cor:xi-terminal-zm-pi-rational-root-specialization-elliptic-image}。
对第二条，当 \(\Pi(\lambda,y_0)\) 可分且不可约时，其伽罗瓦群落入 \(A_4\) 当且仅当其判别式在 \(\QQ\) 中为平方。
而定理 \ref{thm:xi-terminal-zm-leyang-elliptic-structure} 已给出
\(
\mathrm{Disc}_{\lambda}\bigl(\Pi(\lambda,y)\bigr)=-y(y-1)P_{\mathrm{LY}}(y)
\)，
故得到与曲线 \(C\) 的等价刻画。
最后的“薄集外仍有无穷多个满 \(S_4\) 特化”由该正则 \(S_4\)-扩张的 Hilbert 不可约性直接推出。证毕。
\end{proof}

\begin{corollary}[Lee--Yang 四次椭圆覆盖的本原性]\label{cor:xi-terminal-zm-leyang-cover-primitive}
\leanverified{paper\_xi\_terminal\_zm\_leyang\_cover\_primitive}
设 \(\pi_y:E_{\min}\to\mathbb P^1_y\) 为定理 \ref{thm:xi-terminal-zm-leyang-elliptic-structure} 中的次数 \(4\) 态射。
则不存在曲线 \(X\) 与非平凡分解
\[
E_{\min}\xrightarrow{\deg 2}X\xrightarrow{\deg 2}\mathbb P^1_y
\]
使 \(\pi_y\) 等于上述复合。
\end{corollary}
\begin{proof}
若存在上述分解，则 \(\pi_y\) 的几何单值群在四张纸上保留一个 \(2+2\) 的分块结构，从而为不本原子群。
但定理 \ref{thm:xi-terminal-zm-leyang-elliptic-structure} 给出 \(\mathrm{Mon}(\pi_y)\cong S_4\)，而 \(S_4\) 在四点作用下为本原群，矛盾。证毕。
\end{proof}

\begin{corollary}[分歧除子的主化与有限分歧点的群律守恒]\label{cor:xi-terminal-zm-leyang-ramification-sum-zero}
\leanverified{paper\_xi\_terminal\_zm\_leyang\_ramification\_sum\_zero}
沿用定理 \ref{thm:xi-terminal-zm-leyang-elliptic-structure} 的记号，记 \(O\) 为 \(E_{\min}\) 的无穷远点，并令 \(P_0,P_1,P_2,P_3,P_4\) 为 \(\pi_y\) 的五个有限分歧点（即 \(\dd y=0\) 的有限零点）。
则在 \(\mathrm{Pic}^0(E_{\min})\cong E_{\min}\) 中有
\[
[P_0]+[P_1]+[P_2]+[P_3]+[P_4]-5[O]\sim 0,
\]
等价地，在 \(E_{\min}(\overline{\QQ})\) 的群律下满足
\[
P_0+P_1+P_2+P_3+P_4=O.
\]
并且在 \(E_{\min}:u^2+u=x^3-x\) 的 \((x,u)\)-坐标下有
\[
P_0+P_1=\Bigl(\frac14,-\frac38\Bigr),\qquad
P_2+P_3+P_4=\Bigl(\frac14,-\frac58\Bigr)\in E_{\min}(\QQ).
\]
\end{corollary}
\begin{proof}
由定理 \ref{thm:xi-terminal-zm-leyang-elliptic-structure} 知 \(y\) 在 \(O\) 处有唯一四阶极点，故 \((y)_\infty=4O\)。
设 \(R\) 为 \(\pi_y\) 的分歧除子，则 Hurwitz 公式给出
\(
K_{E_{\min}}\sim \pi_y^*K_{\mathbb P^1}+R
\)。
由于 \(K_{E_{\min}}\sim 0\) 且 \(K_{\mathbb P^1}\sim-2\infty\)，得到
\(
R\sim 2(y)_\infty=8O
\)。
另一方面，\(O\) 为全分歧点（指数 \(4\)），故 \(R\) 在 \(O\) 处系数为 \(3\)；其余分歧点为五个有限单分歧点 \(P_0,\dots,P_4\)。
因此 \(R=3O+\sum_{i=0}^{4}P_i\)，代入 \(R\sim 8O\) 即得 \(\sum_{i=0}^{4}P_i-5O\sim 0\)。
在 \(\mathrm{Pic}^0(E_{\min})\cong E_{\min}\) 下该类对应群律和，故 \(P_0+P_1+P_2+P_3+P_4=O\)。
\par
最后，把 \(E_{\min}\) 配方为 \(Y^2=X^3-X+\tfrac14\)（其中 \(X=x,\ Y=u+\tfrac12\)），则
\[
P_0=(1,-\tfrac12),\qquad P_1=(-1,\tfrac12),
\]
由加法公式得 \(P_0+P_1=(\tfrac14,\tfrac18)\)，变换回 \(u\)-坐标即为 \((\tfrac14,-\tfrac38)\)。
再由 \(P_2+P_3+P_4=-(P_0+P_1)\) 与 \(u\mapsto -u-1\) 的取负法则即得 \((\tfrac14,-\tfrac58)\)。证毕。
\end{proof}

\begin{corollary}[倍点拉回下的分歧稳定性：分歧像不生新参数且分歧点按 \texorpdfstring{$E_{\min}[2^n]$}{E[2^n]} 平移复制]\label{cor:xi-terminal-zm-leyang-ramification-pullback-2n}
对整数 \(n\ge 0\) 令 \([2^n]:E_{\min}\to E_{\min}\) 为倍点同态，并定义
\[
y_n:=y\circ[2^n]\in \QQ(E_{\min}).
\]
则 \(y_n:E_{\min}\to\PP^1_y\) 的分歧像集合与 \(y\) 完全一致，仍为
\[
\{\infty,0,1\}\ \cup\ \{P_{\mathrm{LY}}(y)=0\}.
\]
更精确地，若记 \(R(y)\) 为 \(y\) 的分歧除子，则有严格恒等式
\[
R(y_n)=[2^n]^*R(y)
=3\sum_{Q\in E_{\min}[2^n]}[Q]\;+\;\sum_{Q\in E_{\min}[2^n]}\ \sum_{i=0}^{4}[P_i+Q].
\]
因此 \(y_n\) 的五个有限分歧点在 \([2^n]\) 的拉回下被平移复制为 \(5\cdot 4^n\) 个简单分歧点，而无穷远处的全分歧点被复制为 \(4^n\) 个全分歧点。
\end{corollary}
\begin{proof}
在特征 \(0\) 下倍点同态 \([2^n]\) 处处非分歧（\'etale），故对任意有理函数 \(f\) 有 \(R(f\circ[2^n])=[2^n]^*R(f)\)。
将结论 \ref{cor:xi-terminal-zm-leyang-ramification-sum-zero} 的 \(R(y)=3O+\sum_{i=0}^{4}P_i\) 代入并展开：
由于 \([2^n]^{-1}(O)=E_{\min}[2^n]\)。
对每个 \(i\)，任选一点 \(P_{i,n}\in[2^n]^{-1}(P_i)\)，则 \([2^n]^{-1}(P_i)=P_{i,n}+E_{\min}[2^n]\) 为 \(E_{\min}[2^n]\) 的平移余类；
将 \([2^n]^*[P_i]=\sum_{Q\in E_{\min}[2^n]}[P_{i,n}+Q]\) 代入并对 \(i\) 求和，即得所示除子等式（以余类记号书写时可简记为 \(\sum_{Q\in E_{\min}[2^n]}[P_i+Q]\)）。
分歧像陈述由“分歧点在 \'etale 拉回下取原像，而分歧值为其像”直接推出。证毕。
\end{proof}

\begin{corollary}[dyadic 分歧点逆极限的 Tate 扭子分解]\label{cor:xi-terminal-zm-leyang-branchpoint-tate-torsor-decomposition}
沿用推论 \ref{cor:xi-terminal-zm-leyang-ramification-pullback-2n} 的记号。
对 \(i\in\{0,1,2,3,4\}\) 与 \(n\ge 0\) 定义分歧点纤维
\[
\Pi_{i,n}:=[2^n]^{-1}(P_i)\subset E_{\min}(\overline{\QQ}),
\qquad
\Pi_{i,\infty}:=\varprojlim_n\bigl(\Pi_{i,n},[2]\bigr),
\qquad
\Pi_\infty:=\bigsqcup_{i=0}^{4}\Pi_{i,\infty}.
\]
则对每个 \(i\)，\(\Pi_{i,\infty}\) 为 \(T_2(E_{\min})=\varprojlim_n E_{\min}[2^n]\) 的主齐次空间（torsor）。
并且一旦选取一条相容提升 \(\mathbf P_i=(P_{i,n})_{n\ge 0}\in \Pi_{i,\infty}\)（即 \([2]P_{i,n+1}=P_{i,n}\) 且 \(P_{i,0}=P_i\)），便得到一族自然同构
\[
\Pi_{i,n}\ \cong\ P_{i,n}+E_{\min}[2^n],
\qquad
\Pi_{i,\infty}\ \cong\ \mathbf P_i+T_2(E_{\min}),
\]
其中右端表示在相应层次上的平移余类/扭子模型。
\end{corollary}

\begin{proof}
对每个 \(n\)，集合 \(\Pi_{i,n}\) 非空且为 \(E_{\min}[2^n]\) 的平移余类：任取 \(P_{i,n}\in\Pi_{i,n}\)，则
\[
\Pi_{i,n}=P_{i,n}+E_{\min}[2^n],
\]
且平移作用 \(E_{\min}[2^n]\curvearrowright \Pi_{i,n}\) 自由且传递。
结构映射 \([2]:\Pi_{i,n+1}\to\Pi_{i,n}\) 为满射并与该作用相容：
若 \(Q\in E_{\min}[2^{n+1}]\) 则 \([2](x+Q)=[2]x+[2]Q\)。
因此逆极限 \(\Pi_{i,\infty}=\varprojlim_n(\Pi_{i,n},[2])\) 上诱导出 \(T_2(E_{\min})=\varprojlim_n E_{\min}[2^n]\) 的自然平移作用，且该作用自由且传递，从而 \(\Pi_{i,\infty}\) 为 \(T_2(E_{\min})\)-torsor。
\par
若进一步给定一条相容提升 \(\mathbf P_i=(P_{i,n})_{n\ge 0}\in\Pi_{i,\infty}\)，则逐层平移给出同构
\[
E_{\min}[2^n]\stackrel{\sim}{\to}\Pi_{i,n},\quad Q\mapsto P_{i,n}+Q,
\qquad
T_2(E_{\min})\stackrel{\sim}{\to}\Pi_{i,\infty},\quad \mathbf Q\mapsto \mathbf P_i+\mathbf Q,
\]
从而得到所述“余类/扭子”模型表示。
\end{proof}

\input{con__xi-terminal-zm-leyang-finite-branch-regular-4ary-address}

\input{subsubsec__xi-terminal-zm-leyang-coherence-stopping-time}

\par
\input{con__xi-terminal-zm-leyang-spectral-quartic-a-deformation}

\begin{remark}[八面体单值化的射影模 \(3\) 读出与权一提升的可检验接口]\label{rem:xi-terminal-zm-octahedral-mod3-langlands}
标准同构 \(S_4\simeq \mathrm{PGL}_2(\FF_3)\)（例如 \cite{MITSmallGroupsPDF}）把定理 \ref{thm:xi-terminal-zm-leyang-elliptic-structure} 的单值群识别为八面体射影线性群像；
从而在 \(U\) 上得到一个有限像的 \(\mathrm{PGL}_2(\FF_3)\)-投影局部系统，可视为射影模 \(3\) 的几何 Langlands 参数。
另一方面，由结论 \ref{con:xi-terminal-zm-generic-galois-s4} 与结论 \ref{con:xi-terminal-zm-specialization-discriminant-ab} 的泛型 \(S_4\) 性与 Hilbert 不可约性可知：对除去薄集的有理特化 \(y_0\in\QQ\)，所得数域的伽罗瓦闭包给出一个八面体型（\(S_4\)）射影表示 \(G_\QQ\to \mathrm{PGL}_2(\FF_3)\)；
在满足 oddness 等算术条件的情形下，其二维提升落入 Langlands--Tunnell 的权 \(1\) 自守范畴，从而把“分歧—单值”数据与可计算的权一模形式检验接口对接。\cite{Tunnell1981OctahedralArtin,BuzzardPotentialModularitySurvey}
\end{remark}

\begin{remark}[零点凝聚层级的椭圆化入口]\label{rem:xi-terminal-zm-leyang-zero-measure-outlook}
在椭圆单值化框架下，有限体积多项式族 \(\{Z_m(y)\}\) 的零点集合可被转写为 \(\pi_y\) 的某类椭圆函数序列之水平集/相位锁定集；
据此可在严格代数意义上刻画零点凝聚曲线，并进一步寻求相应密度公式与推前测度的解析表达。
\end{remark}

\input{conj__xi-leyang-elliptic-regulator-stokes-hub}

\input{conj__xi-metallic-mean-leyang-modular-anchor}

\begin{theorem}[扭点平行四边形的 Stokes 缺陷与 Weil 配对接口]\label{thm:xi-elliptic-torsion-stokes-weil-pairing}
\leanverified{paper\_xi\_elliptic\_torsion\_stokes\_weil\_pairing}
设 \(E=\CC/\Lambda\) 为复椭圆曲线，\(\pi:\CC\to E\) 为自然投影。
取非零亚纯函数 \(f\in\CC(E)^\times\)，其除子记为
\[
\mathrm{div}(f)=\sum_{q\in E} n_q\,[q],\qquad n_q\in\ZZ,\ \text{有限支撑}.
\]
令 \(k\ge 1\)，并取两条 \(2^k\)-扭点 \(T_1,T_2\in E[2^k]\)。
对每个 \(P\in E[2^k]\) 选取提升 \(\widetilde P\in \frac{1}{2^k}\Lambda\) 使 \(\pi(\widetilde P)=P\)，并同样选取提升 \(\widetilde T_i\in \frac{1}{2^k}\Lambda\) 使 \(\pi(\widetilde T_i)=T_i\)。
对 \(i\in\{1,2\}\) 定义取值于 \(\CC/2\pi\mathrm{i}\ZZ\) 的离散 \(1\)-余链
\begin{equation}\label{eq:xi-elliptic-discrete-log-derivative-cochain}
\Omega_f(P\to P+T_i)
~:=~
\int_{\gamma(\widetilde P,\widetilde P+\widetilde T_i)} \dd\log(f\circ\pi)
\quad \in\ \CC/2\pi\mathrm{i}\ZZ,
\end{equation}
其中路径 \(\gamma(\widetilde P,\widetilde P+\widetilde T_i)\subset\CC\) 连接两端点并避开 \(\pi^{-1}(\mathrm{supp}\,\mathrm{div}(f))\)。
令 \(\square(P;T_1,T_2)\) 表示由 \(T_1,T_2\) 生成的扭点方形
\[
P\to P+T_1\to P+T_1+T_2\to P+T_2\to P.
\]
则在 \(\CC/2\pi\mathrm{i}\ZZ\) 中成立严格恒等式
\begin{equation}\label{eq:xi-elliptic-stokes-defect-divisor-index}
(d\Omega_f)\bigl(\square(P;T_1,T_2)\bigr)
\equiv
2\pi\mathrm{i}\sum_{q} n_q\,\mathrm{Ind}\bigl(q;P,T_1,T_2\bigr)
\pmod{2\pi\mathrm{i}\ZZ},
\end{equation}
其中 \(\mathrm{Ind}(q;P,T_1,T_2)\in\ZZ\) 为点 \(q\) 相对以 \(\widetilde P,\widetilde T_1,\widetilde T_2\) 确定之平行四边形围道的环绕指数。
特别地，若对所有 \(P\in E[2^k]\) 上述平行四边形围道均不包围任何 \(f\) 的零极点，则 \(d\Omega_f\equiv 0\)。
\par
进一步，若 \(f\) 为特征 theta 函数之比（从而 \(\mathrm{div}(f)\) 为扭点差的平移不变除子），则存在 \(S\in E[2^k]\) 使得由 \(\exp(\Omega_f)\) 所诱导的乘法 \(1\)-余圈与 Weil 配对同构：
\begin{equation}\label{eq:xi-elliptic-theta-weil-pairing-character}
\exp\!\bigl(\Omega_f(P\to P+T)\bigr)=e_{2^k}(T,S)
\qquad(\forall\,P,T\in E[2^k]).
\end{equation}
\end{theorem}
\begin{proof}
在 \(\CC\setminus\pi^{-1}(\mathrm{supp}\,\mathrm{div}(f))\) 上，\(\dd\log(f\circ\pi)\) 为闭 \(1\)-形式，其沿不同路径的积分差仅可能来自绕过零极点的回路，其值属于 \(2\pi\mathrm{i}\ZZ\)；故 \eqref{eq:xi-elliptic-discrete-log-derivative-cochain} 在 \(\CC/2\pi\mathrm{i}\ZZ\) 中良定。
\par
取 \(\square(P;T_1,T_2)\) 在 \(\CC\) 中的提升边界曲线 \(\partial\mathcal P\)（\(\mathcal P\) 为对应平行四边形域），则
\[
(d\Omega_f)\bigl(\square(P;T_1,T_2)\bigr)\equiv \oint_{\partial\mathcal P}\dd\log(f\circ\pi)\pmod{2\pi\mathrm{i}\ZZ}.
\]
由复分析的辩值原理/留数定理（等价于对亚纯微分的广义 Stokes 公式），右端等于 \(2\pi\mathrm{i}\) 乘以 \(\mathcal P\) 内零极点的代数计数，并以环绕指数加权：
\[
\oint_{\partial\mathcal P}\dd\log(f\circ\pi)
=2\pi\mathrm{i}\sum_q n_q\,\mathrm{Ind}(q;P,T_1,T_2),
\]
从而得到 \eqref{eq:xi-elliptic-stokes-defect-divisor-index}。
若域内无零极点，则右端为 \(0\)，故 \(d\Omega_f\equiv 0\)。
\par
最后一条关于 theta 函数与 Weil 配对的陈述，是 Mumford theta 群对 \(E[2^k]\) 的平移线性化之标准结论：由扭点差除子所定义的线丛在 \(E[2^k]\) 上的平移自同构因子给出一条双角色，而 Weil 配对刻画了该双角色的规范模型；因此 \(\exp(\Omega_f)\) 必与某个 \(S\in E[2^k]\) 的 \(e_{2^k}(\,\cdot\,,S)\) 同构。证毕。
\end{proof}

\endinput

\begin{conclusion}[Lee--Yang 特征谱曲线的椭圆化与极小模型辨识]\label{con:xi-terminal-zm-elliptic-normalization}
沿用结论 \ref{con:xi-terminal-zm-order4-rational} 的记号，并令
\[
\mathcal C:=\bigl\{(\lambda,y)\in\mathbb A^2:\ \Pi(\lambda,y)=0\bigr\}.
\]
定义新变量
\[
w:=2y+1-2\lambda^{2}.
\]
并令
\[
\eta:=y-\lambda^{2}\qquad(\text{亦记 }u).
\]
则在 \((x,u)\) 坐标（其中 \(x=\lambda\)）下有恒等式
\[
\Pi\bigl(x,\ u+x^{2}\bigr)=u^{2}+u-(x^{3}-x),
\]
从而在 \(\QQ\) 上得到整系数极小模型
\[
E_{\min}:\quad u^{2}+u=x^{3}-x\qquad(\text{亦即 }\eta^{2}+\eta=x^{3}-x).
\]
等价地，若将谱坐标重命名为 \(U:=\lambda\)、\(V:=y-\lambda^{2}\)，则该极小模型亦可写为 \(V^{2}+V=U^{3}-U\)。
并且 \((x,u)\leftrightarrow(\lambda,y)\) 为互逆的多项式代换，从而在仿射层面给出 \(\mathcal C\) 与 \(E_{\min}\) 的 \(\QQ\)-同构；
若取射影闭包，则 \(\Pi(\lambda,y)=0\) 在无穷远处可以出现奇点，但规范化后即为光滑椭圆曲线。
该同构可由显式代换
\[
(x,u)\longmapsto(\lambda,y)=(x,\ u+x^{2}),
\qquad
(\lambda,y)\longmapsto(x,u)=(\lambda,\ y-\lambda^{2})
\]
直接给出。
进一步引入 \(w=2u+1\)（等价于上式定义的 \(w=2y+1-2\lambda^{2}\)）后，得到同构
\[
\mathcal C\ \dashrightarrow\ \mathcal E,\qquad
(\lambda,y)\longmapsto(\lambda,w),
\qquad
\mathcal E:\ w^2=4\lambda^{3}-4\lambda+1,
\]
其逆变换由
\[
y=\frac{2\lambda^{2}+w-1}{2}
\]
给出。
令 \(X=\lambda,\ Y=w/2\)，则 \(\mathcal E\) 的射影闭包为短 Weierstrass 曲线
\[
E:\ Y^2=X^3-X+\frac14,
\]
并且在平面坐标 \((\lambda,y)\) 上存在完全显式的平方差分解
\[
\Pi(\lambda,y)=\Bigl(y-\lambda^2+\tfrac12\Bigr)^2-\Bigl(\lambda^3-\lambda+\tfrac14\Bigr).
\]
因此令 \(X:=\lambda\)、\(Y:=y-\lambda^2+\tfrac12\)（亦即 \(Y=w/2\)）可得到函数域同一
\[
\QQ(\mathcal C)\cong \QQ(E),\qquad \lambda=X,\quad Y=y-\lambda^2+\tfrac12.
\]
更精确地，存在完全显式的 \(\QQ\)-双有理互逆
\[
\varphi:\ \mathcal C\dashrightarrow E,\quad (\lambda,y)\longmapsto (X,Y):=\bigl(\lambda,\ y+\tfrac12-\lambda^2\bigr),
\qquad
\psi:\ E\dashrightarrow \mathcal C,\quad (X,Y)\longmapsto (\lambda,y):=\bigl(X,\ X^2-\tfrac12+Y\bigr),
\]
从而 \(\mathcal C\) 的规范化同构于 \(E\)，并且 \(\QQ(\mathcal C)\cong \QQ(E)\)。
其判别式与 \(j\)-不变量分别为
\[
\Delta(E)=37,\qquad j(E)=\frac{48^{3}}{37}=\frac{110592}{37}.
\]
进一步令 \(y_{\mathrm{int}}:=Y-\tfrac12\)（亦即 \(y_{\mathrm{int}}=u=y-\lambda^{2}\)），则得到整系数最小模型
\[
E_{\min}:\quad y_{\mathrm{int}}^{2}+y_{\mathrm{int}}=X^{3}-X,
\]
其 \(\Delta(E_{\min})=37\) 且 \(j(E_{\min})=48^{3}/37\)。
因此 \((\lambda,y)\)-谱曲线的代数分支结构可等价地提升为椭圆曲线 \(E\) 上的有理坐标结构，从而把分歧点、解析延拓与单值化问题转写为 \(E\) 上的代数几何数据。
特别地，\(\Delta(E)=37\) 表明该椭圆曲线在除 \(37\) 外处处良好约化；并且在 Cremona/LMFDB 标识下对应同源类 \(37\mathrm a\) 的强 Weil 曲线 \(37\mathrm a1\)（导子 \(37\)，见注记 \ref{rem:xi-terminal-elliptic-37a1-modular-anchor}）。
\par
在上述同构下，重量参数作为 \(E\) 上的有理函数可写为
\[
\boxed{\;
y=\frac{2\lambda^{2}+w-1}{2}=X^2+Y-\frac12\ \in\ \QQ(E)\; } ,
\qquad (O\text{ 为 }E\text{ 的无穷远点}).
\]
并记
\[
P_{0}=\Bigl(0,\frac12\Bigr),\quad P_{-1}=\Bigl(-1,-\frac12\Bigr),\quad P_{1}=\Bigl(1,-\frac12\Bigr)\in E(\QQ),
\]
\[
Q_{2}=\Bigl(2,-\frac52\Bigr),\quad Q_{1}=\Bigl(1,\frac12\Bigr),\quad Q_{-1}=\Bigl(-1,\frac12\Bigr)\in E(\QQ).
\]
由 Weierstrass 形式知 \(\operatorname{ord}_{O}(X)=-2\)、\(\operatorname{ord}_{O}(Y)=-3\)，从而 \(\operatorname{ord}_{O}(y)=-4\) 且 \((y)_\infty=4O\)，故 \(y:E\to\PP^1_y\) 为次数 \(4\) 的覆盖映射。
零点由 \(y=0\iff Y=-X^2+\tfrac12\) 与曲线方程联立得到：
\[
(-X^2+\tfrac12)^2=X^3-X+\tfrac14\iff X(X-1)^2(X+1)=0,
\]
从而主除子分解为
\[
\mathrm{div}(y)=P_{0}+P_{-1}+2P_{1}-4O.
\]
同理，由 \(y=1\iff Y=\tfrac32-X^2\) 得
\[
(\tfrac32-X^2)^2=X^3-X+\tfrac14\iff (X-2)(X-1)(X+1)^2=0,
\]
从而
\[
\mathrm{div}(y-1)=Q_{2}+Q_{1}+2Q_{-1}-4O.
\]
两式皆为主除子，故其在 \(\mathrm{Pic}^0(E)\cong E\) 中对应零类；等价地在 \(E(\QQ)\) 的群结构中成立
\[
P_{0}+P_{-1}+2P_{1}=O,
\qquad
Q_{2}+Q_{1}+2Q_{-1}=O.
\]
结合结论 \ref{con:xi-terminal-zm-elliptic-mw-generator-alignment} 的生成元 \(R=(0,\tfrac12)\) 与低倍点对齐，上述两条纤维分解亦可在除子意义下写为
\[
y^{-1}(0)=\{R\}+2\{-2R\}+\{3R\},\qquad
y^{-1}(1)=\{2R\}+2\{-3R\}+\{4R\},
\]
其中各点按重数计。
因此 \(y\) 在 \(y=0\) 与 \(y=1\) 处各出现一处二重纤维点（分别位于 \(\lambda=1\) 与 \(\lambda=-1\)），而其余分歧值由判别式分解（结论 \ref{con:xi-terminal-zm-discriminant-leyang} 或等价的结果式刻画，见结论 \ref{con:xi-terminal-zm-leyang-cubic-branch-image}）所完全决定；这给出一个纯代数几何的“除子—分歧”证书。
\end{conclusion}
\begin{proof}[证明要点]
将 \(\Pi(\lambda,y)=0\) 视为关于 \(y\) 的二次方程，得
\[
y^2+(1-2\lambda^2)y+\bigl(\lambda^4-\lambda^3-\lambda^2+\lambda\bigr)=0.
\]
配方可化为
\[
\bigl(2y+1-2\lambda^2\bigr)^2=4\lambda^3-4\lambda+1.
\]
令 \(X=\lambda\)、\(Y=y+\tfrac12-\lambda^2\)，即得 \(Y^2=X^3-X+\tfrac14\)，从而得到 \(\varphi\)。
反向代入 \(y=X^2-\tfrac12+Y\) 立即复原 \(\Pi(X,y)=0\)，故 \(\varphi,\psi\) 互为双有理逆。
Weierstrass 形式下取 \(a=-1,b=1/4\)，则
\(
\Delta=-16(4a^3+27b^2)=37
\)，且 \(j=c_4^3/\Delta\)（其中 \(c_4=-48a\)）给出 \(j(E)=110592/37\)。
可复现核验脚本：\texttt{scripts/exp\_fold\_zm\_elliptic\_leyang\_arithmetic\_audit.py}。证毕。
\end{proof}

\input{con__xi-terminal-zm-elliptic-finite-field-point-count}
\input{con__xi-terminal-zm-spectral-curve-infty-a4-cusp}

\begin{conclusion}[次数 \(4\) 线性系 \texorpdfstring{$|4O|$}{|4O|} 的射影正规模型：两二次曲面的显式交与谱四次的线性投影复原]\label{con:xi-terminal-zm-elliptic-degree4-projective-model}
沿用结论 \ref{con:xi-terminal-zm-elliptic-normalization} 的整系数极小模型
\(
E_0:\ \eta^{2}+\eta=x^{3}-x
\)，记 \(O\) 为无穷远点，并令重量函数
\(
y:=x^{2}+\eta\in\QQ(E_0)
\)。
则
\(
1,x,x^{2},\eta\in H^{0}\!\bigl(E_0,\mathcal O_{E_0}(4O)\bigr)
\)
且构成 \(H^{0}(E_0,\mathcal O_{E_0}(4O))\) 的一组基。
由完全线性系 \(|4O|\) 诱导的态射
\[
\kappa:E_0\hookrightarrow \mathbb P^{3},\qquad
(x,\eta)\longmapsto [U:V:W:Z]=[1:x:x^{2}:\eta]
\]
为射影嵌入，其像为次数 \(4\) 的椭圆正规曲线，并在 \(\mathbb P^{3}\) 中由两条二次方程刻画为完全交：
\[
UW-V^{2}=0,
\qquad
Z^{2}+UZ-VW+VU=0.
\]
在该模型中有恒等式 \(y=(W+Z)/U\)，从而覆盖 \(y:E_0\to\mathbb P^{1}_{y}\) 可视为由二超平面丛 \(\langle U,\ W+Z\rangle\subset H^{0}(E_0,\mathcal O(4O))\) 所定义的次数 \(4\) 线性投影；在仿射片 \(U=1\) 上消去 \((W,Z)\)（等价于代入 \(\eta=y-x^{2}\)）即严格复原谱方程 \(\Pi(x,y)=0\)。
\end{conclusion}
\begin{proof}[证明要点]
在 \(E_0\) 上有 \(\mathrm{ord}_{O}(x)=-2\)、\(\mathrm{ord}_{O}(\eta)=-3\)，故 \(1,x,x^{2},\eta\in H^{0}(E_0,\mathcal O_{E_0}(4O))\)。
若 \(a_{0}+a_{1}x+a_{2}x^{2}+a_{3}\eta=0\)，比较 \(O\) 处极阶先得 \(a_{3}=0\)，继而 \(a_{2}=a_{1}=a_{0}=0\)，从而四者线性无关；由于 \(\dim H^{0}(E_0,\mathcal O_{E_0}(4O))=4\)，它们构成一组基。
当 \(\deg(4O)=4\ge 3\) 时 \(|4O|\) 为 very ample，故 \(\kappa\) 为嵌入；其像为次数 \(4\) 的椭圆正规曲线，因而为两二次曲面的完全交。
两条二次方程由代入直接核验：第一式 \(UW-V^{2}=0\) 来自 \(W=x^{2}\)；第二式在仿射图 \(U=1\) 上即为 \(\eta^{2}+\eta=x^{3}-x\)，写作 \(Z^{2}+Z=VW-V\) 并齐次化即得。
最后 \(y=(W+Z)/U\) 由定义推出，而在仿射片 \(U=1\) 上代入 \(\eta=y-x^{2}\) 即将二次关系化为 \(\Pi(x,y)=0\)。证毕。
\end{proof}

\begin{conclusion}[共轭重量支的主除子：低倍有理点轨道上的完全写法]\label{con:xi-terminal-zm-elliptic-conjugate-weight-principal-divisors}
沿用结论 \ref{con:xi-terminal-zm-elliptic-normalization} 的记号，并令椭圆反演为
\(
\iota:(X,Y)\mapsto(X,-Y)
\)。
定义共轭重量函数
\[
y^{\iota}:=\iota(y)=X^{2}-Y-\frac12\in\QQ(E).
\]
记 \(R:=(0,\tfrac12)\in E(\QQ)\) 为结论 \ref{con:xi-terminal-zm-elliptic-mw-generator-alignment} 的自由生成元对应点。
则 \(y^{\iota}\) 与 \(y^{\iota}-1\) 的主除子满足
\[
\operatorname{div}(y^{\iota})=(-R)+(-3R)+2(2R)-4O,
\qquad
\operatorname{div}(y^{\iota}-1)=(-4R)+(-2R)+2(3R)-4O.
\]
因此两条特殊纤维（\(y^{\iota}=0\) 与 \(y^{\iota}=1\)）的零点（含重数）均支撑在由 \(R\) 生成的 Mordell--Weil 轨道上；
并且 \(y^{\iota}\)（分别 \(y^{\iota}-1\)）作为 \(|4O|\) 的截面，其零除子在该轨道上的分布被上述主除子恒等式唯一固定。
\end{conclusion}
\begin{proof}[证明要点]
由 \(y^{\iota}=X^{2}-Y-\tfrac12\) 与结论 \ref{con:xi-terminal-zm-elliptic-normalization} 的除子证书，
对 \(\operatorname{div}(y)\)、\(\operatorname{div}(y-1)\) 施加 \(\iota\)（注意 \(\iota(O)=O\)）即可得到
\[
\operatorname{div}(y^{\iota})=\Bigl(0,-\tfrac12\Bigr)+\Bigl(-1,\tfrac12\Bigr)+2\Bigl(1,\tfrac12\Bigr)-4O,
\qquad
\operatorname{div}(y^{\iota}-1)=\Bigl(2,\tfrac52\Bigr)+\Bigl(1,-\tfrac12\Bigr)+2\Bigl(-1,-\tfrac12\Bigr)-4O.
\]
再由结论 \ref{con:xi-terminal-zm-elliptic-mw-generator-alignment} 的低倍点对齐
\(
[2]R=(1,\tfrac12),\ [3]R=(-1,-\tfrac12),\ [4]R=(2,-\tfrac52)
\)
即可把上述坐标表示改写为所示 \(\langle R\rangle\) 轨道形式。证毕。
\end{proof}

\begin{conclusion}[有理耦合的 Mordell--Weil 参数化：\texorpdfstring{$\Pi(\lambda,y_0)$}{Pi(lambda,y0)} 在 \texorpdfstring{$\mathbb Q$}{Q} 上有根当且仅当 \(y_0\in y\bigl(E(\QQ)\bigr)\)]\label{con:xi-terminal-zm-pi-rational-root-iff-y-image}
沿用结论 \ref{con:xi-terminal-zm-elliptic-normalization} 的谱四次
\(
\Pi(\lambda,y)=0
\)
及其椭圆化同构。
则对任意 \(y_{0}\in\QQ\)，以下两条件等价：
\begin{enumerate}
  \item 存在 \(\lambda\in\QQ\) 使得 \(\Pi(\lambda,y_{0})=0\)；
  \item 存在 \(P\in E(\QQ)\) 使得 \(y_{0}=y(P)\)。
\end{enumerate}
并且在上述条件成立时，可取 \(\lambda=X(P)\) 作为对应的有理谱根；在极小整模型 \(E_0:\eta^2+\eta=x^3-x\)（其中 \(x=X,\ \eta=Y-\tfrac12\)）下等价地有 \(y_0=x(P)^2+\eta(P)\) 且 \(\lambda_0=x(P)\)。
特别地，此时 \(\Pi(\lambda,y_{0})\) 在 \(\QQ[\lambda]\) 中含有理线性因子 \(\lambda-X(P)\)；并由结论 \ref{con:xi-terminal-zm-elliptic-mw-generator-alignment}（\(E_0(\QQ)\) 无挠且 \(\mathrm{rank}=1\)，可取生成元 \(R_0=(0,0)\)）知 \(y\bigl(E(\QQ)\bigr)\subset\QQ\) 为无限集。
若将 \(X(P)\) 写为既约分式 \(X(P)=U/V^{2}\)（\(U,V\in\ZZ,\ \gcd(U,V)=1,\ V>0\)），则存在三次多项式 \(G(\lambda)\in\QQ[\lambda]\) 使得
\[
\Pi(\lambda,y_{0})=(V^{2}\lambda-U)\,G(\lambda),
\]
并且既约分母满足
\[
\mathrm{den}(y_{0})=\mathrm{den}\bigl(X(P)\bigr)^{2}=V^{4}.
\]
\end{conclusion}
\begin{proof}[证明要点]
若存在 \(\lambda\in\QQ\) 使 \(\Pi(\lambda,y_{0})=0\)，由结论 \ref{con:xi-terminal-zm-elliptic-normalization} 的恒等式
\(
4\Pi(\lambda,y)=w^{2}-(4\lambda^{3}-4\lambda+1)
\)
（其中 \(w=2y+1-2\lambda^{2}\)）可知 \((\lambda,w)\in\mathcal E(\QQ)\)，从而 \(P=(X,Y)=(\lambda,w/2)\in E(\QQ)\)，并满足 \(y(P)=y_{0}\)。
反之，若 \(P\in E(\QQ)\)，则 \((\lambda,y)=(X(P),y(P))\) 满足 \(\Pi(\lambda,y)=0\)，从而 \(\lambda=X(P)\) 给出 \(\Pi(\lambda,y_{0})\) 的有理根。
\par
线性因子与清分母因子 \((V^{2}\lambda-U)\) 由 \(\lambda=X(P)=U/V^{2}\) 为根立即推出。
分母平方关系可直接由曲线方程与 \(y=X^{2}+Y-\tfrac12\) 推出：写 \(X=U/V^{2}\)（既约），则由
\(
Y^{2}=X^{3}-X+\tfrac14
\)
得 \(Y=W/(2V^{3})\)（\(W\in\ZZ\)）且 \(W^{2}=4U^{3}-4UV^{4}+V^{6}\)；
从而
\[
y(P)=\frac{U^{2}}{V^{4}}+\frac{W}{2V^{3}}-\frac12=\frac{2U^{2}+WV-V^{4}}{2V^{4}}.
\]
注意 \(2U^{2}+WV-V^{4}\) 恒为偶数，故 \(y(P)=N/V^{4}\)（\(N\in\ZZ\)）。
若奇素数 \(p\mid V\)，则 \(WV\equiv V^{4}\equiv 0\pmod p\) 且 \(2N\equiv 2U^{2}\not\equiv 0\pmod p\)，从而 \(p\nmid N\)；
若 \(2\mid V\)，则 \(U\) 为奇数；
并且由 \(W^{2}=4U^{3}-4UV^{4}+V^{6}\) 在模 \(8\) 下推出 \(W^{2}\equiv 4\pmod 8\)，故 \(W=2W_{1}\) 且 \(W_{1}\) 为奇数。
于是 \(WV\equiv 0\pmod 4\)、\(V^{4}\equiv 0\pmod 4\)，而 \(2U^{2}\equiv 2\pmod 4\)，从而 \(2N\equiv 2\pmod 4\)，即 \(N\) 为奇数；
综上 \(\gcd(N,V)=1\)。
因此 \(\mathrm{den}(y(P))=V^{4}=\mathrm{den}(X(P))^{2}\)。证毕。
\end{proof}

\begin{remark}[双判别桥：\texorpdfstring{$\Pi$}{Pi} 的两条判别式方向]\label{rem:xi-terminal-zm-pi-double-discriminant-bridge}
将 \(\Pi(\lambda,y)\) 视为关于 \(y\) 的二次多项式，则其 \(y\)-判别式为
\[
\mathrm{Disc}_{y}\bigl(\Pi(\lambda,y)\bigr)=4\lambda^{3}-4\lambda+1.
\]
等价地，令 \(w:=2y+1-2\lambda^{2}\)，则
\[
4\Pi(\lambda,y)=w^{2}-\bigl(4\lambda^{3}-4\lambda+1\bigr).
\]
从而同一方程在 \(y\)-方向上的判别式给出椭圆曲线模型 \(\mathcal E:w^{2}=4\lambda^{3}-4\lambda+1\)（见结论 \ref{con:xi-terminal-zm-elliptic-normalization}）；其中三次式 \(4\lambda^{3}-4\lambda+1\) 亦即 \(E_{\min}\) 的 \(2\)-division 多项式（非平凡 \(2\)-挠点的 \(x\)-坐标所满足之方程）。
另一方面，将 \(\Pi(\lambda,y)\) 视为关于 \(\lambda\) 的四次多项式，则其 \(\lambda\)-判别式在 \(\ZZ[y]\) 中分解为
\[
\mathrm{Disc}_{\lambda}\bigl(\Pi(\lambda,y)\bigr)=-y(y-1)P_{\mathrm{LY}}(y),
\qquad
P_{\mathrm{LY}}(y):=256y^{3}+411y^{2}+165y+32,
\]
从而得到属 \(2\) 的判别式曲线 \(C:\ W^{2}=-y(y-1)P_{\mathrm{LY}}(y)\)（此处记 \(W\) 以避免与椭圆坐标 \(w\) 混淆）。
该“二次—四次”的双判别结构把椭圆归一化、分歧像以及四次分裂域的二次分辨率统一编码在同一 \(\Pi\) 之中。
\end{remark}

\begin{remark}[外部命名与交叉核验接口：\texorpdfstring{$37\mathrm a1$}{37a1}、\texorpdfstring{$X_0^+(37)$}{X0+(37)} 与 level \(37\) 新形式]\label{rem:xi-terminal-elliptic-37a1-modular-anchor}
本段的椭圆曲线 \(E: Y^2=X^3-X+\tfrac14\) 在 \(\QQ\) 上与整系数模型
\[
E_0:\quad \eta^2+\eta=x^3-x
\]
同构；显式同构由
\(
X=x,\ Y=\eta+\tfrac12
\)
给出；等价地亦可取
\(
X=x,\ Y=-\eta-\tfrac12
\)
（二者相差椭圆反演）。
本文统一采用 Weierstrass 不变量的标准约定
\[
c_4^3-c_6^2=1728\Delta,
\]
并在一般 Weierstrass 形式
\(
\eta^2+a_1x\eta+a_3\eta=x^3+a_2x^2+a_4x+a_6
\)
的不变量记号下对 \(E_0\) 取
\(
(a_1,a_2,a_3,a_4,a_6)=(0,0,1,-1,0)
\)，
从而 \(b_2=0\)、\(b_4=-2\)、\(b_6=1\)、\(b_8=-1\)，并由标准判别式公式得到 \(\Delta(E_0)=64-27=37\)，与上文短 Weierstrass 模型 \(E\) 的计算一致。
若在外部文献中以极小判别式 \(-37\) 或等价的相反号判别式记之，则对应于对本文 \(\Delta\) 采用不同的符号约定，并不改变本文以 \(\Delta(E)=37\) 为锚定的后续论证。
据 \cite{LMFDBEllipticCurve37a1} 的记录，在 Cremona/LMFDB 的标准标识下，\(E\) 为导子 \(37\) 的同源类 \(37\mathrm a\) 之强 Weil 曲线（记号 \(37\mathrm a1\)）；
其极小整系数模型 \(E_0:\eta^2+\eta=x^3-x\) 亦可与 Cremona 在 \(N=37\) 的算例对照。\cite{CremonaExamplesPDF}
\par
上述外部标识虽不参与本文的代数推导，但给出一条可直接核验的 Langlands 定量锚点：与该同源类对应的 weight \(2\)、level \(37\) 新形式轨道为 \(37.2.a.a\)（系数域 \(\QQ\)，解析秩 \(1\)），相应自守表示记为 \(\pi_{37}\)，并且对任意良素 \(p\neq 37\) 有
\(
a_p=p+1-\#E(\FF_p)
\)
与该新形式的 Hecke 本征值一致，从而 \(L(E,s)=L(f,s)=L(\pi_{37},s)\)（模性定理；亦可在 Langlands 口径下视为与 \(\mathrm{GL}_2(\mathbb{A}_\QQ)\) 的相应自守对象对接）。\cite{LMFDBNewform372aa}
\par
因此对每个素数 \(\ell\)，椭圆曲线 \(E\) 的 \(\ell\)-进 Tate 模诱导二维伽罗瓦表示
\[
\rho_{E,\ell}:\mathrm{Gal}(\overline{\QQ}/\QQ)\longrightarrow \mathrm{GL}_2(\ZZ_\ell),
\]
其仅在 \(\{37,\ell\}\) 处可能分歧；并且对任意良素 \(p\neq 37,\ell\)，Frobenius 元的特征多项式为 \(T^2-a_pT+p\)，亦即 \(T_p\) 的 Hecke 特征多项式，从而局部 Euler 因子与 Hecke 数据一致。
\par
并且由 Sturm 有界性可得该外部锚定具有有限证书：对 \(k=2,N=37\) 有 \([SL_2(\ZZ):\Gamma_0(37)]=38\)，故 Sturm 界
\(
B=\bigl\lfloor \frac{k}{12}[SL_2(\ZZ):\Gamma_0(N)]\bigr\rfloor
\)
等于 \(6\)；从而仅需核验至 \(q^6\) 的有限多个 Fourier/Hecke 系数即可唯一确定该新形式轨道。\cite{LMFDBSturmBound}
在模曲线口径下，\(X_0(37)\) 为属 \(2\) 曲线；LMFDB 记录 \(E_0\) 正可作为其 Fricke involution 的商
\(
X_0^+(37):=X_0(37)/\langle w_{37}\rangle
\)
的代数模型之一。\cite{LMFDBEllipticCurve37a1}
\par
令 \(q=e^{2\pi i\tau}\)。
在上述外部锚定下，可取规范化新形式 \(f(\tau)=\sum_{n\ge 1}a_n q^n\in S_2(\Gamma_0(37))\) 使得 \(E_0\) 的 N\'eron 微分 \(\omega=\dd x/(2\eta+1)\) 在模参数化下满足
\(
\omega=2\pi i\, f(q)\,\dd q/q
\)；
并且在谱坐标下有 \(2\eta+1=2y+1-2x^2=\partial_y\Pi(x,y)\)，从而 \(\omega=\dd x/\partial_y\Pi(x,y)\)。
于是重量函数 \(y=x^{2}+\eta\in\QQ(E_0)\)（等价地 \(y=x^{2}-\eta-1\)）作为 \(X_0^+(37)\) 上的模函数在无穷远尖点具有主部 \(q^{-4}\)，其截断展开与对应新形式的 \(q\)-展开见下式：
\input{sections/generated/eq_fold_zm_elliptic_modular_y_qexp_audit}
可复现核验脚本：\texttt{scripts/exp\_fold\_zm\_elliptic\_modular\_y\_qexp\_audit.py}。
\end{remark}

\begin{conclusion}[\texorpdfstring{$j$}{j}-不变量的标准基本域反演：纯虚模参数 \texorpdfstring{$\tau$}{tau} 的唯一性与数值证书]\label{con:xi-terminal-zm-elliptic-j-inversion-tau}
结论 \ref{con:xi-terminal-zm-elliptic-normalization} 给出
\(
j(E)=110592/37
\)。
则在 \(SL_{2}(\ZZ)\) 的标准基本域内，满足 \(j(\tau)=j(E)\) 的模参数 \(\tau\) 唯一，并且为纯虚数
\[
\tau=i\cdot 1.221127360764627\ldots.
\]
相应的
\(
q
\)-参数为
\[
q=e^{2\pi i\tau}=0.000465420392294374\ldots.
\]
因此 \(E(\CC)\) 可被显式均匀化为 \(\CC/(\ZZ+\tau\ZZ)\) 的适当复伸缩；并且权函数 \(y\) 在该均匀化下为阶 \(4\) 的椭圆函数，其解析尺度与判别式素数 \(37\) 的模性锚点一致。
\end{conclusion}
\begin{proof}[证明要点]
由于 \(j(E)\in\RR\) 且 \(j(E)>1728=j(i)\)，标准模论给出在基本域内的唯一解必落在纯虚轴上，即 \(\tau=it\)（\(t>1\)）。
对函数 \(t\mapsto j(it)\) 进行数值反演即可恢复 \(t\) 与 \(q=e^{-2\pi t}\)。
该反演可在高精度下通过 Eisenstein 级数的 \(q\)-展开实现，并由脚本 \texttt{scripts/exp\_fold\_zm\_elliptic\_modular\_y\_qexp\_audit.py} 复现核验。证毕。
\end{proof}

\begin{conclusion}[模素纤维根数恒等式：由 \texorpdfstring{$\Pi(\lambda,y)$}{Pi(lambda,y)} 的分解统计恢复 \texorpdfstring{$a_p$}{a\_p}]\label{con:xi-terminal-zm-fiber-root-counting-hecke}
取素数 \(p\notin\{2,37\}\)，并在 \(\FF_p\) 上考虑仿射曲线
\[
\mathcal C_p:=\{(\lambda,y)\in\FF_p^2:\ \Pi(\lambda,y)=0\}.
\]
对每个 \(y\in\FF_p\) 定义纤维根数
\[
N_p(y):=\card{\{\lambda\in\FF_p:\ \Pi(\lambda,y)=0\}}.
\]
则有严格点计数恒等式
\[
\boxed{\;
\card{E(\FF_p)}=1+\sum_{y\in\FF_p}N_p(y)\;}
\]
其中右端的 \(1\) 对应无穷远点 \(O\)。
令 \(a_p:=p+1-\card{E(\FF_p)}\) 为 Frobenius 迹（见注记 \ref{rem:xi-terminal-elliptic-37a1-modular-anchor}），则等价地
\[
\boxed{\;
a_p=p-\sum_{y\in\FF_p}N_p(y)\;}
\]
并由 Hasse 界得到有效估计
\[
\left|\frac{1}{p}\sum_{y\in\FF_p}N_p(y)-1\right|
=\left|\frac{a_p}{p}\right|
\le \frac{2}{\sqrt p}.
\]
更一般地，对任意 \(n\ge 1\)，上述双有理变换在 \(\FF_{p^n}\) 上同样给出
\[
\card{\{(\lambda,y)\in\FF_{p^n}^2:\ \Pi(\lambda,y)=0\}}=\card{E(\FF_{p^n})}-1.
\]
因此若定义仿射谱曲线的局部 \(\zeta\)-函数为
\[
Z(\mathcal C/\FF_p,T):=\exp\!\left(\sum_{n\ge 1}\card{\mathcal C(\FF_{p^n})}\frac{T^n}{n}\right),
\]
则有闭式
\[
\boxed{\; Z(\mathcal C/\FF_p,T)=\frac{1-a_pT+pT^2}{1-pT}\; }.
\]
\end{conclusion}
\begin{proof}[证明要点]
当 \(p\notin\{2,37\}\) 时，结论 \ref{con:xi-terminal-zm-elliptic-normalization} 的双有理变换
\[
(\lambda,y)\longmapsto (\lambda,w),\qquad w:=2y+1-2\lambda^{2}
\]
在 \(\FF_p\) 上仍可逆（因为 \(2\) 可逆），并把 \(\mathcal C_p\) 双射到仿射椭圆曲线
\(
\mathcal E_p:\ w^2=4\lambda^3-4\lambda+1
\)
的 \(\FF_p\)-点集；该仿射模型与射影椭圆曲线 \(E\) 的差别仅为无穷远点 \(O\)。
因此
\[
\sum_{y\in\FF_p}N_p(y)=\card{\mathcal C_p}=\card{E(\FF_p)}-1,
\]
从而得到所示恒等式。
将其代入 \(a_p:=p+1-\card{E(\FF_p)}\) 立得 \(a_p=p-\sum_{y\in\FF_p} N_p(y)\)。
最后一式由 Hasse 界 \(|a_p|\le 2\sqrt p\) 给出。
可复现核验脚本：\texttt{scripts/exp\_fold\_zm\_elliptic\_fiber\_root\_counting\_audit.py}。证毕。
\end{proof}

\input{con__xi-terminal-zm-s4-modp-splitting-density}

\begin{conclusion}[Lee--Yang 三次多项式的分歧像刻画：结果式给出的最小整系数证书]\label{con:xi-terminal-zm-leyang-cubic-branch-image}
把 \(\Pi(\lambda,y)=0\) 视为关于 \(\lambda\) 的四次方程，并令投影
\(
\pi:\mathcal C\to\mathbb A^1_y
\)
由 \((\lambda,y)\mapsto y\) 给出。
则 \(\pi\) 的分歧值集合在 \(\ZZ[y]\) 中由结果式精确刻画：
\[
\boxed{\;
\mathrm{Res}_{\lambda}\!\bigl(\Pi(\lambda,y),\partial_{\lambda}\Pi(\lambda,y)\bigr)
=-y(y-1)\bigl(256y^3+411y^2+165y+32\bigr)\; }.
\]
因此除平凡分歧值 \(y\in\{0,1\}\) 外，全部非平凡分歧值恰为三次多项式
\(
256y^3+411y^2+165y+32
=0
\)
的三根；其中唯一负实根即结论 \ref{con:xi-terminal-zm-discriminant-leyang} 的 \(y_{\mathrm{LY}}\)。
该三次因子因而并非经验拟合对象，而是投影 \(\pi\) 的分歧像在整系数层面的最小消元证书。
\end{conclusion}
\begin{proof}[证明要点]
分歧点当且仅当存在 \(\lambda\) 使得 \(\Pi(\lambda,y)=\partial_\lambda\Pi(\lambda,y)=0\)。
对 \(\lambda\) 作消元即可得到结果式分解；由于 \(\Pi\) 首一，该结果式与 \(\mathrm{Disc}_{\lambda}(\Pi)\) 一致，从而与结论 \ref{con:xi-terminal-zm-discriminant-leyang} 的判别式因式分解完全相容。证毕。
\end{proof}

\begin{conclusion}[分歧点的 \texorpdfstring{$\lambda$}{lambda}-参数化与 Lee--Yang 实根的同一化]\label{con:xi-terminal-zm-leyang-lambda-cubic}
设 \((\lambda,y)\) 满足谱曲线 \(\Pi(\lambda,y)=0\)。
则全部有限分歧点（即存在二重根之 \(y\)）由联立方程
\[
\Pi(\lambda,y)=0,\qquad \partial_{\lambda}\Pi(\lambda,y)=0
\]
所刻画，并且消元得到分歧所发生之谱坐标必满足
\[
\boxed{\;(\lambda-1)(\lambda+1)\bigl(16\lambda^3-9\lambda^2+1\bigr)=0\;}
\]
且相应分歧值由显式有理式给出
\[
\boxed{\;
y=\frac{4\lambda^3-3\lambda^2-2\lambda+1}{4\lambda}
=\frac{(\lambda-1)(4\lambda^2+\lambda-1)}{4\lambda}\; }.
\]
其中 \(\lambda=1\) 与 \(\lambda=-1\) 分别对应分歧值 \(y=0\) 与 \(y=1\)；
其余三根对应的分歧值恰为三次方程 \(256y^3+411y^2+165y+32=0\) 的三根（结论 \ref{con:xi-terminal-zm-leyang-cubic-branch-image}）。
特别地，结论 \ref{con:xi-terminal-zm-discriminant-leyang} 的唯一负实分歧值 \(y_{\mathrm{LY}}\) 对应于三次方程 \(16\lambda^3-9\lambda^2+1=0\) 的唯一实根 \(\lambda_{\mathrm{LY}}<0\)。
从而获得对 \(y_{\mathrm{LY}}\) 的等价刻画：
\[
\lambda_{\mathrm{LY}}\approx -0.2734348106524542,
\qquad y_{\mathrm{LY}}=y(\lambda_{\mathrm{LY}}),
\qquad 16\lambda_{\mathrm{LY}}^3-9\lambda_{\mathrm{LY}}^2+1=0.
\]
因此 \(y_{\mathrm{LY}}\) 可等价刻画为有理映射 \(y=y(\lambda)\) 在三次证书 \(16\lambda^3-9\lambda^2+1=0\) 的唯一实根处的取值。
\end{conclusion}
\begin{proof}[证明要点]
由 \(\partial_{\lambda}\Pi(\lambda,y)=0\) 可解出
\(
y=(4\lambda^3-3\lambda^2-2\lambda+1)/(4\lambda)
\)（\(\lambda\neq 0\)）。
代回 \(\Pi(\lambda,y)=0\) 并清分母得到
\(
(\lambda-1)(\lambda+1)(16\lambda^3-9\lambda^2+1)=0
\)；
其中 \(\lambda=\pm 1\) 分别对应分歧值 \(y=0,1\)；
其余三根由三次因子刻画，并经 \(y=y(\lambda)\) 映射落入结论 \ref{con:xi-terminal-zm-leyang-cubic-branch-image} 的三次分歧像，从而得到 Lee--Yang 实分歧值与伴随谱根的同一化。证毕。
\end{proof}

\begin{conclusion}[Lee--Yang 三次场的判别式与二次子域：\texorpdfstring{$\QQ(\sqrt{-111})$}{Q(sqrt(-111))}]\label{con:xi-terminal-zm-leyang-cubic-arithmetic}
令
\(
g(y):=256y^3+411y^2+165y+32
\)。
则其判别式满足严格分解
\[
\boxed{\;\mathrm{Disc}(g)=-3^{9}\cdot 31^{2}\cdot 37\;}
\]
且 \(g\) 在 \(\QQ[y]\) 上不可约，从而其伽罗瓦群为 \(S_3\)。
记 \(L\) 为 \(g\) 的分裂域，则其唯一二次子域为
\[
\boxed{\ \QQ(\sqrt{\mathrm{Disc}(g)})=\QQ(\sqrt{-111})\ },
\]
其中 \(-111\) 为 \(\mathrm{Disc}(g)\) 的平方自由部分，亦为该虚二次域的基本判别式。
\end{conclusion}
\begin{proof}[证明要点]
判别式因子分解由对三次多项式的标准判别式公式直接计算并在 \(\ZZ\) 中分解得到。
由有理根判据可排除 \(g\) 的有理根，从而 \(g\) 不可约；对不可约三次而言，判别式非平方当且仅当伽罗瓦群为 \(S_3\)。
二次子域由 \(\QQ(\sqrt{\mathrm{Disc}(g)})\) 给出；由于 \(\mathrm{Disc}(g)\) 的平方自由部分为 \(-111\)，故该二次域为 \(\QQ(\sqrt{-111})\)。
可复现核验脚本：\texttt{scripts/exp\_fold\_zm\_elliptic\_leyang\_arithmetic\_audit.py}。证毕。
\end{proof}

\begin{conclusion}[Lee--Yang 三次在 \texorpdfstring{$\FF_p$}{Fp} 上的分解型密度律：由二次特征 \texorpdfstring{$\left(\frac{-111}{p}\right)$}{(-111/p)} 完全控制“恰有一根”情形]\label{con:xi-terminal-zm-leyang-cubic-modp-splitting-density}
令
\[
g(y):=256y^3+411y^2+165y+32\in\ZZ[y],
\]
并取素数 \(p\nmid 3\cdot 31\cdot 37\)。
记 \(\varepsilon_p:=\left(\frac{-111}{p}\right)\in\{\pm1\}\) 为二次特征（Legendre 符号），并把 \(g\) 模 \(p\) 化为 \(g_p\in\FF_p[y]\)。
则 \(g_p\) 的分解型满足：
\begin{itemize}
  \item 若 \(\varepsilon_p=-1\)，则 \(g_p\) 必为 \((1)(2)\) 型，从而在 \(\FF_p\) 上恰有一个根；
  \item 若 \(\varepsilon_p=+1\)，则 \(g_p\) 要么全分裂（\((1)(1)(1)\) 型，三根），要么不可约（\((3)\) 型，无根）。
\end{itemize}
并且上述三类素数集合的自然密度分别为
\[
\delta\bigl((1)(2)\bigr)=\frac12,\qquad
\delta\bigl((1)(1)(1)\bigr)=\frac16,\qquad
\delta\bigl((3)\bigr)=\frac13.
\]
因此 \(g_p\) 在 \(\FF_p\) 上具有至少一个根（即分解型为 \((1)(2)\) 或 \((1)(1)(1)\)）的素数集合密度为
\[
\delta\bigl(g_p\ \text{在}\ \FF_p\ \text{上有根}\bigr)=\frac12+\frac16=\frac23.
\]
对判别式二次脊曲线
\(
C_p:\ w^{2}=-y(y-1)g_p(y)
\)
而言，这等价于其分歧点集合在 \(\FF_p\) 上除 \(\{y=0,1,\infty\}\) 之外还出现额外的 \(\FF_p\)-Weierstrass 点。
\end{conclusion}
\begin{proof}[证明要点]
由结论 \ref{con:xi-terminal-zm-leyang-cubic-arithmetic}，\(g\) 在 \(\QQ\) 上不可约且其分裂域的伽罗瓦群为 \(S_3\)，并具有唯一二次分辨子域 \(\QQ(\sqrt{-111})\)。
对 \(p\nmid \mathrm{Disc}(g)\) 的素数，Frobenius 元 \(\mathrm{Frob}_p\in S_3\) 的共轭类与 \(g_p\) 的分解型一一对应：
恒等类对应全分裂 \((1)(1)(1)\)，换位类对应 \((1)(2)\)，而三循环类对应不可约 \((3)\)。
另一方面，符号同态 \(\mathrm{sgn}:S_3\to\{\pm1\}\) 的固定域正是上述二次子域；
因而 \(\mathrm{sgn}(\mathrm{Frob}_p)=\left(\frac{-111}{p}\right)=\varepsilon_p\)。
于是 \(\varepsilon_p=-1\) 强制 \(\mathrm{Frob}_p\) 落入换位共轭类，从而 \(g_p\) 必为 \((1)(2)\) 型；而 \(\varepsilon_p=+1\) 等价于 \(\mathrm{Frob}_p\) 为偶置换，从而仅可能为恒等或三循环，给出 \((1)(1)(1)\) 或 \((3)\) 两型。
密度由 Chebotarev 定理给出，为相应共轭类大小除以 \(|S_3|=6\)，即分别为 \(3/6,1/6,2/6\)。
可复现核验脚本：\texttt{scripts/exp\_fold\_zm\_elliptic\_leyang\_cubic\_modp\_splitting\_audit.py}。证毕。
\end{proof}

\begin{conclusion}[判别式的素因子骨架与不可约性：素数 \texorpdfstring{$37$}{37} 在三条互异分歧机制中以一次幂强制出现]\label{con:xi-terminal-37-squarefree-coupling}
设 \(E\) 取短 Weierstrass 模型 \(Y^2=X^3-X+\tfrac14\)，并记
\[
\psi_2(X):=4X^3-4X+1\quad(\text{\(E\) 的 \(2\)-division 多项式}),
\qquad
S(X):=16X^3-9X^2+1\quad(\text{结论 \ref{con:xi-terminal-zm-elliptic-y-hurwitz-passport} 的分歧三次式}),
\]
\[
P_{\mathrm{LY}}(y):=256y^3+411y^2+165y+32\quad(\text{Lee--Yang 三次因子}).
\]
则三者判别式与 \(E\) 的判别式之间满足严格恒等式
\[
\Delta(E)=37,\qquad
\mathrm{Disc}(\psi_2)=2^4\cdot 37,\qquad
\mathrm{Disc}(S)=-2^2\cdot 3^3\cdot 37,\qquad
\mathrm{Disc}(P_{\mathrm{LY}})=-3^9\cdot 31^2\cdot 37.
\]
从而 \(\psi_2\) 的分裂域之唯一二次分辨子域为 \(\QQ(\sqrt{37})\)。
从而素数 \(37\) 以一次幂同时出现在：
\begin{itemize}
  \item 椭圆曲线 \(E\) 的最小坏约化处；
  \item \(y\)-投影分歧点的 \(X\)-坐标三次域（由 \(S\) 控制）的基本判别式因子中；
  \item Lee--Yang 三次因子 \(P_{\mathrm{LY}}(y)\) 的分裂判别式中。
\end{itemize}
特别地，对任意素数 \(p\notin\{2,3,31,37\}\)，上述三多项式在 \(\FF_p\) 上均无重根，且 \(y=0,1\) 与 \(P_{\mathrm{LY}}(y)=0\) 的三点在模 \(p\) 下保持互异；
从而该 \(S_4\)-覆盖塔在此集合之外呈现纯粹驯分歧并具有稳定约化型。
\end{conclusion}
\begin{proof}[证明要点]
\(\Delta(E)\) 由短 Weierstrass 形式的标准公式 \(\Delta(E)=-16(4a^3+27b^2)\)（对 \(Y^2=X^3+aX+b\)）直接计算（此处 \(a=-1,b=\tfrac14\)）得到 \(\Delta(E)=37\)。
其余三个判别式为一元判别式公式的直接代数计算，并在 \(\ZZ\) 中分解。
最后的约化断言来自：若 \(p\nmid \mathrm{Disc}(f)\) 则 \(f\) 模 \(p\) 无重根；
并且当 \(p\notin\{2,3\}\) 时 \(P_{\mathrm{LY}}(0)=32\)、\(P_{\mathrm{LY}}(1)=864\) 皆非零，从而 \(0,1\) 不与 \(P_{\mathrm{LY}}\) 的根碰撞；
再结合 \(p\nmid |S_4|=24\) 得驯分歧与稳定约化结论。证毕。
\end{proof}

\input{con__xi-terminal-zm-collision-s2-e2-field-isomorphism}


\begin{theorem}[Lee--Yang 特征曲线的椭圆化与 \(37\)-模性]\label{thm:xi-terminal-zm-leyang-elliptic-modularity-37}\label{thm:xi-terminal-zm-leyang-elliptic-37-modularity-monodromy-langlands}
设
\[
\Pi(\lambda,y):=\lambda^{4}-\lambda^{3}-(2y+1)\lambda^{2}+\lambda+y(y+1)\in\ZZ[\lambda,y],
\]
并令仿射曲线 \(\mathcal C\subset\mathbb A^2_{(\lambda,y)}\) 由 \(\Pi(\lambda,y)=0\) 给出。
令
\[
x:=\lambda,\qquad v:=y-x^{2}.
\]
其中 \(v\) 与前文记号 \(u:=y-\lambda^2\) 一致。
则 \(\mathcal C\) 与椭圆曲线
\[
E_0:\quad v^{2}+v=x^{3}-x
\]
在 \(\QQ\) 上同构；可取显式互逆同构
\[
\phi:E_0\to \mathcal C,\quad (x,v)\longmapsto(\lambda,y)=(x,v+x^{2}),
\qquad
\phi^{-1}:\mathcal C\to E_0,\quad (\lambda,y)\longmapsto (x,v)=(\lambda,y-\lambda^{2}).
\]
并且 \(\Delta(E_0)=37\)、\(j(E_0)=110592/37\)；在 Cremona/LMFDB 记号下 \(E_0\) 对应强 Weil 曲线 \(37\mathrm a1\)，并可作为
\(
X_0^+(37):=X_0(37)/\langle w_{37}\rangle
\)
的代数模型之一。\cite{LMFDBEllipticCurve37a1}
\end{theorem}

\leanverified{paper\_xi\_terminal\_zm\_leyang\_elliptic\_37\_modularity\_monodromy\_langlands}

\begin{proof}
将 \(y=v+x^2\) 与 \(\lambda=x\) 代入得到恒等式
\[
\Pi(x,v+x^2)=v^2+v-x^3+x,
\]
从而 \(\Pi(\lambda,y)=0\) 与 \(v^2+v=x^3-x\) 等价；并且 \((x,v)\leftrightarrow(\lambda,y)\) 的互逆性由 \(y=v+x^2\) 在仿射坐标环中的可逆性直接给出，故 \(\phi\) 为 \(\QQ\)-同构。
不变量计算与外部标识见注记 \ref{rem:xi-terminal-elliptic-37a1-modular-anchor}。证毕。
\end{proof}
\par

\begin{proposition}[重量参数作为椭圆曲线上的四重函数：分歧点、分歧值与群律约束]\label{prop:xi-terminal-zm-leyang-f-ramification}
令 \(E_0\) 为定理 \ref{thm:xi-terminal-zm-leyang-elliptic-modularity-37} 的椭圆曲线，\(O\) 为无穷远点。
设
\[
f:=v+x^{2}=\phi^{*}(y)\in\QQ(E_0),
\]
并以同一符号记由 \(f\) 诱导的有理映射 \(f:E_0\to\PP^{1}\)。
则：
\begin{enumerate}
  \item \(\QQ(E_0)/\QQ(f)\) 为次数 \(4\) 的函数域扩张，并且 \(x\) 在 \(\QQ(f)\) 上的极小多项式为
  \[
  \Pi(x,f)=x^{4}-x^{3}-(2f+1)x^{2}+x+f(f+1)\in\QQ(f)[x].
  \]
  \item \(f\) 在 \(O\) 处具有唯一四阶极点，且 \((f)_\infty=4O\)。因此 \(\deg(f)=4\)，并在 \(O\) 上方全分歧（分歧指数 \(e_O=4\)）。
  \item \(f\) 的有限分歧点恰为 \(\dd f=0\) 的点；其 \(x\)-坐标满足
  \[
  \dd f=0\quad\Longleftrightarrow\quad (x^{2}-1)\bigl(16x^{3}-9x^{2}+1\bigr)=0.
  \]
  更具体地，有限分歧点由五点组成：两条有理点
  \[
  P_0=(1,-1),\qquad P_1=(-1,0),
  \]
  以及三点 \(P_2,P_3,P_4\)，其中 \(x(P_i)\) 为三次方程 \(16x^{3}-9x^{2}+1=0\) 的三根，且
  \[
  v(P_i)=\frac{1-3x(P_i)^2-2x(P_i)}{4x(P_i)}\qquad(i=2,3,4).
  \]
  \item 令 \(B\subset\PP^1\) 为覆盖 \(f:E_0\to\PP^1\) 的分歧值集合，则
  \[
  B=\{\infty,0,1\}\ \cup\ \{t\in\overline{\QQ}:\ P_{\mathrm{LY}}(t)=0\},
  \qquad
  P_{\mathrm{LY}}(t):=256t^3+411t^2+165t+32.
  \]
  并且 \(B\) 的每个有限分歧值之上恰出现一处简单分歧（分歧指数 \(2\)），其余点不分歧。
  \item 作为 \(E_0\) 上的亚纯微分，\(\dd f\) 的零极点除子满足
  \[
  \mathrm{div}(\dd f)=P_0+P_1+P_2+P_3+P_4-5O,
  \]
  从而在 \(E_0(\overline{\QQ})\) 的群律下有
  \[
  P_0+P_1+P_2+P_3+P_4=O,
  \qquad
  P_2+P_3+P_4=-(P_0+P_1)=\Bigl(\frac14,-\frac58\Bigr)\in E_0(\QQ).
  \]
\end{enumerate}
\end{proposition}
\begin{proof}
由 \(f=v+x^2\) 得 \(v=f-x^2\)。代回 \(E_0:v^2+v=x^3-x\) 并展开，得到
\[
(f-x^2)^2+(f-x^2)=x^3-x
\quad\Longleftrightarrow\quad
x^{4}-x^{3}-(2f+1)x^{2}+x+f(f+1)=0,
\]
即 \(\Pi(x,f)=0\)。因此 \([\QQ(E_0):\QQ(f)]\le 4\)。
\par
把 \(E_0\) 配方为短 Weierstrass 形式 \(Y^2=X^3-X+\tfrac14\)（取 \(X=x\)、\(Y=v+\tfrac12\)），则 \(x\) 与 \(v\) 在 \(O\) 处的极点阶分别为 \(2\) 与 \(3\)，故 \(f=v+x^2\) 在 \(O\) 处极点阶为 \(4\)，从而 \((f)_\infty=4O\) 且 \(\deg(f)=4\)。由次数等于函数域扩张次数可得 \([\QQ(E_0):\QQ(f)]=4\)，并且 \(\Pi(x,f)\) 为 \(x\) 的极小多项式。
\par
对 \(E_0:v^2+v=x^3-x\) 微分得到
\[
(2v+1)\,\dd v=(3x^2-1)\,\dd x.
\]
于是
\[
\dd f=\dd(v+x^2)=\dd v+2x\,\dd x=\left(\frac{3x^2-1}{2v+1}+2x\right)\dd x.
\]
因此 \(\dd f=0\) 等价于 \(3x^2-1+2x(2v+1)=0\)，亦即
\[
4xv+3x^2+2x-1=0.
\]
在 \(x\neq 0\) 上可解得 \(v=(1-3x^2-2x)/(4x)\)；
将其代回 \(v^2+v=x^3-x\) 并化简即得 \((x^2-1)(16x^3-9x^2+1)=0\)，从而得到五个有限分歧点的显式刻画。
\par
记 \(E_0\) 的标准正则微分
\[
\omega:=\frac{\dd x}{2v+1}.
\]
由上式可得 \(\dd f=(4xv+3x^2+2x-1)\,\omega\)。
由于 \(g(E_0)=1\)，\(\omega\) 在 \(E_0\) 上处处正则且无零极点；因此 \(\dd f\) 在五个有限分歧点处各有一个简单零点。
另一方面，由 \((f)_\infty=4O\) 可取局部参数 \(t\) 使 \(f\sim t^{-4}\)，则 \(\dd f\sim -4t^{-5}\dd t\)，故 \(\dd f\) 在 \(O\) 处为五阶极点，从而 \(\mathrm{div}(\dd f)=\sum_{i=0}^{4}P_i-5O\)。
该除子为典范类，而椭圆曲线的典范类在 \(\mathrm{Pic}^0(E_0)\) 中为零；在同一识别 \(E_0\simeq \mathrm{Pic}^0(E_0)\) 下得到群律恒等式 \(\sum_{i=0}^4 P_i=O\)。
最后，在短 Weierstrass 坐标 \((X,Y)\) 中，
\[
P_0=(1,-\tfrac12),\qquad P_1=(-1,\tfrac12),
\]
故弦割斜率为 \(-\tfrac12\)，从而 \(P_0+P_1=(\tfrac14,\tfrac18)\)，变回 \(v=Y-\tfrac12\) 得
\(
P_0+P_1=(\tfrac14,-\tfrac38)
\)；
取负元 \(v\mapsto -v-1\) 即得
\(
-(P_0+P_1)=(\tfrac14,-\tfrac58)
\)，从而 \(P_2+P_3+P_4=-(P_0+P_1)\)。
\par
最后，由 \(\Pi(x,f)=0\) 作为 \(\QQ(f)\) 上的四次方程，分歧值由判别式为零刻画；结合结论 \ref{con:xi-terminal-zm-discriminant-leyang} 的分解
\(
\mathrm{Disc}_{x}\bigl(\Pi(x,f)\bigr)=-f(f-1)P_{\mathrm{LY}}(f)
\)
得到分歧值集合 \(B=\{\infty,0,1\}\cup\{P_{\mathrm{LY}}(t)=0\}\)。
证毕。
\end{proof}
\par

\begin{theorem}[分歧型 \texorpdfstring{$(4;2^{5})$}{(4;2^5)} 覆盖的全对称单值化：\texorpdfstring{$\mathrm{Mon}(f)\cong S_4$}{Mon(f)=S4}]\label{thm:xi-terminal-zm-leyang-monodromy-s4}
\leanverified{paper\_xi\_terminal\_zm\_leyang\_monodromy\_s4}
沿用命题 \ref{prop:xi-terminal-zm-leyang-f-ramification} 的 \(f:E_0\to\PP^{1}\)。
则 \(f\) 的几何单值群同构于 \(S_4\)。
等价地，将 \(\Pi(\lambda,y)\) 视为 \(\QQ(y)\) 上关于 \(\lambda\) 的四次多项式，则其函数域伽罗瓦群为 \(S_4\)。
\end{theorem}
\begin{proof}
由命题 \ref{prop:xi-terminal-zm-leyang-f-ramification}，覆盖 \(f:E_0\to\PP^1\) 在 \(O\) 上方全分歧，故绕 \(f=\infty\) 的局部单值为一个 \(4\)-循环；而任一有限分歧值处均为简单分歧，其局部单值为换位。
一个 \(4\)-循环与一个换位生成 \(S_4\)，故 \(\mathrm{Mon}(f)\cong S_4\)。
四次多项式的函数域伽罗瓦群与几何单值群之等价表述见结论 \ref{con:xi-terminal-zm-generic-galois-s4}。证毕。
\end{proof}
\par

\begin{proposition}[素数 \(37\) 的双重显影：\(2\)-除子域与 Lee--Yang 分歧域的统一分歧控制]\label{prop:xi-terminal-37-double-ramification-fingerprint}
沿用定理 \ref{thm:xi-terminal-zm-leyang-elliptic-modularity-37} 与命题 \ref{prop:xi-terminal-zm-leyang-f-ramification} 的记号，并记 \(E_0[2]\) 为 \(E_0\) 的 \(2\)-挠子群。
则存在如下算术分歧指纹的一致性：
\begin{enumerate}
  \item 令 \(Z:=2v+1\)，则 \(E_0\) 等价配方为
  \[
  E_0:\quad Z^{2}=4x^{3}-4x+1.
  \]
  其 \(2\)-除法三次 \(4x^{3}-4x+1\) 的判别式为 \(\mathrm{Disc}(4x^{3}-4x+1)=2^{4}\cdot 37\)，从而 \(E_0[2]\) 的分裂域仅在素数 \(2\) 与 \(37\) 处可能分歧。
  \item Lee--Yang 三次
  \[
  P_{\mathrm{LY}}(t)=256t^{3}+411t^{2}+165t+32
  \]
  的判别式满足 \(\mathrm{Disc}(P_{\mathrm{LY}})=-3^{9}\cdot 31^{2}\cdot 37\)，从而其分裂域仅在素数 \(3,31,37\) 处可能分歧，且 \(37\) 以一次幂出现。
  \item 因而同一素数 \(37\) 同时作为 \(E_0\) 的极小判别式（亦即唯一坏约化素数）与 Lee--Yang 分歧域的必分歧素数出现；从分歧数据的观点看，Lee--Yang 分歧集合与椭圆曲线的坏素数在 \(37\) 处发生不可约耦合。
\end{enumerate}
\end{proposition}
\begin{proof}
对三次多项式 \(aX^{3}+bX^{2}+cX+d\) 的判别式取标准公式
\[
\mathrm{Disc}(aX^{3}+bX^{2}+cX+d)=b^{2}c^{2}-4ac^{3}-4b^{3}d-27a^{2}d^{2}+18abcd.
\]
对 \(4x^{3}-4x+1\) 代入 \((a,b,c,d)=(4,0,-4,1)\) 得 \(\mathrm{Disc}=1024-432=592=2^{4}\cdot 37\)，从而 \(E_0[2]\) 的分裂域仅在 \(\{2,37\}\) 处可能分歧。
第二条的判别式分解见结论 \ref{con:xi-terminal-zm-leyang-cubic-arithmetic}。
最后一条由“分裂域仅在判别式素因子处可能分歧”的一般事实立即推出。证毕。
\end{proof}
\par

\begin{proposition}[唯一坏素数与 \texorpdfstring{$\mathrm{GL}_2$}{GL2} 局部数据的锁定]\label{prop:xi-prime-bad-set-determines-unramified-data}
设 \(E/\QQ\) 为椭圆曲线，且其极小判别式满足 \(|\Delta(E)|=p\) 为素数。
则 \(E\) 的坏约化素数集合包含于 \(\{p\}\)，且 \(E\) 在 \(p\) 处为乘法约化，从而导子 \(N(E)=p\)。
若进一步满足 \(\dim S_2(\Gamma_0(p))=1\)，则存在唯一规范化新形式 \(g\in S_2(\Gamma_0(p))\) 使
\[
L(E,s)=L(g,s),
\]
并且对每个良素 \(\ell\neq p\)，与 \(E\) 对应的 \(\ell\)-进 Galois 表示在所有非分歧处的 Frobenius 特征多项式（等价地 Hecke 本征值序列 \((a_\ell)_{\ell\neq p}\)）由该单一算术输入 \(p\) 唯一确定。
特别地，\(p=37\) 时可取 \(E=E_0: v^2+v=x^3-x\)，其满足 \(\Delta(E_0)=37\) 且 \(\dim S_2(\Gamma_0(37))=1\)，故其全部非分歧局部数据由坏素数集合 \(\{37\}\) 锁定。
\end{proposition}
\begin{proof}
\(|\Delta(E)|=p\) 蕴含极小判别式仅在 \(p\) 处有正赋值，故坏约化素数集合包含于 \(\{p\}\)。
并且 \(\ord_p(\Delta(E))=1\) 强制 \(E\) 在 \(p\) 处为乘法约化，因而导子指数为 \(1\)，得到 \(N(E)=p\)。
由模性定理存在 \(g\in S_2(\Gamma_0(p))\) 满足 \(L(E,s)=L(g,s)\)。
当 \(\dim S_2(\Gamma_0(p))=1\) 时，规范化条件 \(a_1=1\) 迫使 \(g\) 唯一，从而其 Hecke 本征值序列唯一；
这等价于所有 \(\ell\neq p\) 处非分歧局部因子（Frobenius 特征多项式）被唯一锁定。
对 \(p=37\) 的锚定由定理 \ref{thm:xi-terminal-zm-leyang-elliptic-modularity-37} 与推论 \ref{cor:xi-terminal-zm-leyang-langlands-level37} 的外部核验接口给出。证毕。
\end{proof}

\begin{conjecture}[判别式的 gcd--导子原则：Lee--Yang 分歧素数与坏约化素数的可计算接口]\label{conj:xi-leyang-discriminant-gcd-conductor-principle}
设由某一 Lee--Yang 特征谱曲线的椭圆化得到椭圆曲线 \(E/\QQ\)，并设模型内禀地给出一整数
\(\mathfrak D_{\mathrm{LY}}\in\ZZ\)
（可取为控制临界值/分歧的判别式、结式或其平方自由部分的规范化代表）。
记 \(\mathrm{sqf}(\cdot)\) 为平方自由部分算子。
则预期存在同一模型类中不依赖坐标选择的对应，使得
\[
\mathrm{sqf}\!\Bigl(\gcd\bigl(\Delta(E),\mathfrak D_{\mathrm{LY}}\bigr)\Bigr)
=
\mathrm{sqf}\!\bigl(N(E)\bigr),
\]
其中 \(\Delta(E)\) 为极小判别式、\(N(E)\) 为导子。
在 \(\Delta(E)\) 平方自由的情形，该关系等价于：\(\gcd(\Delta(E),\mathfrak D_{\mathrm{LY}})\) 的素因子集合与坏约化素数集合一致。
\end{conjecture}

\begin{corollary}[Langlands--\texorpdfstring{$\mathrm{GL}_2$}{GL2} 的可审计实现：导子 \(37\) 的模性对象与八面体有限参数]\label{cor:xi-terminal-zm-leyang-langlands-level37}
令 \(\mathcal C\) 为定理 \ref{thm:xi-terminal-zm-leyang-elliptic-modularity-37} 的特征曲线（亦即 \(\Pi(\lambda,y)=0\) 的正规化），并令 \(E_0\) 为相应椭圆曲线。
并沿用命题 \ref{prop:xi-terminal-zm-leyang-f-ramification} 的四重覆盖 \(f:E_0\to\PP^1\) 及其分歧值集合 \(B\subset\PP^1\)，令 \(U:=\PP^1\setminus B\)。
则：
\begin{enumerate}
  \item 存在唯一的权 \(2\)、level \(37\) 的规范化新形式 \(g\in S_2(\Gamma_0(37))\) 及对应自守表示 \(\pi_{37}\)，并相应存在与 \(E_0\) 相容的 \(\ell\)-进 Galois 表示族 \(\{\rho_{E_0,\ell}\}_{\ell}\)，使得
  \[
  L(\mathcal C,s)=L(E_0,s)=L(g,s)=L(\pi_{37},s).
  \]
  并且对任意良素 \(p\neq 37\) 有点计数恒等式
  \[
  \#E_0(\FF_p)=p+1-a_p,
  \]
  其中 \(a_p\) 为 \(g\) 的 Hecke 本征值（亦即 \(E_0\) 的 Frobenius 迹）。
  \item 推前得到的秩 \(4\) 局部系统
  \[
  \mathcal L:=\left.(f_\ast\underline{\QQ})\right|_{U}
  \]
  具有几何单值像 \(\mathrm{Mon}(\mathcal L)\cong S_4\)，并经标准同构 \(S_4\simeq \mathrm{PGL}_2(\FF_3)\) 可识别为八面体有限像投影局部系统。\cite{MITSmallGroupsPDF}
  此外，命题 \ref{prop:xi-terminal-37-double-ramification-fingerprint} 表明控制 \(E_0[2]\) 与 \(B\setminus\{\infty,0,1\}\) 的两条三次分裂域仅在 \(\{2,3,31,37\}\) 处可能分歧。
\end{enumerate}
由此，谱四次 \(\Pi(\lambda,y)\) 的正规化与分歧单值数据在算术层面内生锁定了一个 level \(37\) 的 \(\mathrm{GL}_2(\mathbb A_\QQ)\) 模性对象（由 \(L\)-函数与 \(\{a_p\}\) 可计算核验）以及一个八面体有限像投影参数（由 \(\mathrm{Mon}(f)\cong S_4\simeq \mathrm{PGL}_2(\FF_3)\) 读出）。
\end{corollary}
\begin{proof}
第一条由定理 \ref{thm:xi-terminal-zm-leyang-elliptic-modularity-37} 的同构 \(\mathcal C\simeq E_0\) 与椭圆曲线模性定理推出；对 level \(37\) 的外部命名与交叉核验接口见注记 \ref{rem:xi-terminal-elliptic-37a1-modular-anchor}。\cite{LMFDBEllipticCurve37a1,LMFDBNewform372aa,ConradTSWFinal}
该新形式的唯一性可由 \(\dim S_2(\Gamma_0(37))=1\) 的标准维数计算或数据库条目直接读出。\cite{LMFDBNewform372aa}
\(\#E_0(\FF_p)=p+1-a_p\)（\(p\neq 37\)）为 Eichler--Shimura 对应与模性表述下的标准点计数公式。
\par
第二条中 \(\mathrm{Mon}(\mathcal L)\cong \mathrm{Mon}(f)\cong S_4\) 由定理 \ref{thm:xi-terminal-zm-leyang-monodromy-s4} 直接给出；八面体识别由 \(S_4\simeq \mathrm{PGL}_2(\FF_3)\) 的标准同构得到。分歧素数集合由命题 \ref{prop:xi-terminal-37-double-ramification-fingerprint} 立即推出。
证毕。
\end{proof}
\par

\subsubsection{Lee--Yang 零因子张量函子与未分歧参数}\label{subsubsec:xi-leyang-zero-divisor-tensor-functor-unramified}
\input{para__xi-leyang-zero-divisor-tensor-functor-langlands}
\par

\begin{conjecture}[Kapustin--Witten 物理几何朗兰兹中的 Lee--Yang 分歧壁：范畴化对应]\label{conj:xi-terminal-zm-leyang-kapustin-witten-wall}
沿用推论 \ref{cor:xi-terminal-zm-leyang-langlands-level37} 第二条的分歧集 \(B\)、开曲线 \(U=\PP^1\setminus B\) 与局部系统 \(\mathcal L\)。
猜想在 Kapustin--Witten 所述的 \(4\mathrm d\ \mathcal N=4\) 超对称 Yang--Mills 的拓扑扭曲与边界/线算符框架下，可构造一类范畴对象 \(\mathfrak E(\mathcal L)\)，使其在几何朗兰兹等价中成为 Hecke 特征对象，特征值局部系统为 \(\mathcal L\)，并且 \(B\) 中的有限分歧值对应于耦合常数复延拓中的壁穿越与非正则极限。\cite{KapustinWitten2006GeometricLanglands}
\end{conjecture}

\begin{conjecture}[谱—外置坏素一致性：坏约化素点为谱兼容 prime-register 的不可规避轴]\label{conj:xi-terminal-zm-spectral-compatible-externalization-badprime-minimal}
设 \(\mathcal C/\QQ\) 为一条由参数无关的代数谱方程所给的特征曲线，并在 \(\QQ\) 上与某条极小椭圆曲线 \(E/\QQ\) 双有理等价。
考虑一类\textbf{谱兼容的外置制度} \(\mathcal S\)：其以 prime-register 纤维 \(\mathcal P\)（定义 \ref{def:pom-prime-fiber}）作为残差账本，使得谱反演/分歧护照的单值化选择在外置上可追溯且可比较；并要求存在有限素数集
\(\mathfrak O(\mathcal S)\subset\mathbb{P}\)
使得对任意 \(p\notin\mathfrak O(\mathcal S)\)，该制度在模 \(p\) 约化下满足“扩张保持”意义的稳定性：覆盖的分歧型与外置比较规则在 \(\FF_p\) 上不发生退化或碰撞。
则猜想不可规避的障碍素集必包含坏约化素点：
\[
\boxed{\ \{p:\ p\mid \Delta(E)\}\ \subseteq\ \mathfrak O(\mathcal S)\ }.
\]
在 Lee--Yang 原型中，由定理 \ref{thm:xi-terminal-zm-leyang-elliptic-structure} 的固定椭圆锚点知 \(\Delta(E)=37\)，故任何谱兼容外置都应满足 \(37\in\mathfrak O(\mathcal S)\)；
并且在“最小障碍谱”口径下，奇素部分可取为 \(\mathfrak O(\mathcal S)\cap(2\mathbb{Z}+1)=\{37\}\)。
\end{conjecture}
\begin{remark}
该猜想与本文关于 \(p=37\) 的多重退化证书一致：
坏约化（定理 \ref{thm:xi-terminal-zm-leyang-elliptic-structure}），
Latt\`es 稳定约化降阶与 Chebyshev 共轭（结论 \ref{con:xi-terminal-zm-elliptic-lattes-badprime-37-degree-drop}），
以及重量像的高重坍缩（例如结论 \ref{con:xi-terminal-zm-elliptic-3torsion-y-image-mod37-collapse}）
均在同一素数处同步出现。
\end{remark}

\endinput

\begin{conclusion}[时间纤维标尺与模尖点增长的同一：\texorpdfstring{$q^{-4}$}{q^{-4}} 主部锁定对数深度]\label{con:xi-terminal-zm-timefiber-cusp-depth}
沿用注记 \ref{rem:xi-terminal-elliptic-37a1-modular-anchor} 的模参数化。
令 \(q=e^{2\pi i\tau}\)，并令 \(y(\tau)\) 表示该参数化下的重量函数。
则在无穷远尖点具有主部 \(q^{-4}\)（同注记 \ref{rem:xi-terminal-elliptic-37a1-modular-anchor} 的截断展开），从而沿纯虚轴 \(\tau=iT\to i\infty\) 有
\[
\boxed{\;
\log y(\tau)=8\pi\,\Im\tau+O(1)\;}
\qquad(T\to+\infty).
\]
特别地，若以 \(b_{-1}:=\log B^{-1}\) 表示端点时间纤维的对数深度坐标，则 Cayley 标尺变换只改变其加性零点（备注 \ref{rem:dephys-cayley-scale-gauge-four}）；
因此在模尖点处，\(q^{-4}\) 主部把该深度与 \(8\pi\,\Im\tau\) 在同一标尺类中锁定到有界误差。
\end{conclusion}
\begin{proof}[证明要点]
由注记 \ref{rem:xi-terminal-elliptic-37a1-modular-anchor} 的 \(q\)-展开可写
\[
y(\tau)=q^{-4}\bigl(1+O(q)\bigr)\qquad(\Im\tau\to+\infty).
\]
沿 \(\tau=iT\) 有 \(q=e^{-2\pi T}\)，故
\[
\log y(iT)=-4\log q+\log(1+O(q))=8\pi T+O(1),
\]
其中 \(\log(1+O(q))\) 在 \(T\to\infty\) 时有界（事实上趋于 \(0\)）。
最后，\(b_{-1}=\log B^{-1}\) 的标尺变换律由备注 \ref{rem:dephys-cayley-scale-gauge-four} 给出，因而上述 \(O(1)\) 误差与标尺常数可在同一加性规范下合并吸收。证毕。
\end{proof}


\begin{conclusion}[重量谱椭圆曲线上的倍乘—Latt\`es 自映射：\texorpdfstring{$x([2]P)$}{x([2]P)} 的闭式四次有理函数]\label{con:xi-terminal-zm-elliptic-lattes-doubling}\label{con:xi-terminal-zm-lattes-4division-certificate}
在结论 \ref{con:xi-terminal-zm-elliptic-normalization} 的椭圆模型
\[
E:\ Y^2=X^3-X+\frac14
\]
上，记倍乘映射 \([2]:E\to E\)，并令 \(x:E\to\mathbb P^1\) 为 \(x\)-坐标函数 \(x(P)=X\)。
则 \(x\circ[2]\) 在仿射坐标 \(X\) 上诱导一个四次 Latt\`es 自映射 \(\Phi\in\QQ(X)\)：
\[
\boxed{\;
\Phi(X):=x([2]P)=\frac{X^4+2X^2-2X+1}{4X^3-4X+1}\; }.
\]
其分母
\(
4X^3-4X+1
=4\bigl(X^3-X+\tfrac14\bigr)
=4Y^2
\)
等价于 \(E\) 的 \(2\)-division 条件 \(Y=0\)（从而恰为 \(x\)-坐标的 \(2\)-division 多项式）。
此外，\(\Phi\) 的临界点可由 \(4\)-division 数据完全刻画：记 \(E[2]\)、\(E[4]\) 分别为 \(2\)-挠与 \(4\)-挠子群，则对仿射坐标 \(X\) 有严格等价
\[
\boxed{\;
\Phi'(X)=0
\ \Longleftrightarrow\
\exists\,Q\in E[4]\setminus E[2]\ \text{使得}\ X=x(Q)
\ \Longleftrightarrow\
2X^{6}-10X^{4}+10X^{3}-10X^{2}+2X+1=0\; }.
\]
从而该六次多项式给出 \(x\)-坐标图上的 \(4\)-division 证书（等价于经典 \(4\)-division 多项式 \(\psi_4\) 除去因子 \(2Y\) 后的 \(X\)-部分），倍乘 Latt\`es 动力学的临界结构在 \(\QQ[X]\) 层面完全可审计。
\par
更进一步，记
\[
D(X):=4X^{3}-4X+1,\qquad
C(X):=2X^{6}-10X^{4}+10X^{3}-10X^{2}+2X+1.
\]
则 \(2\)-division 多项式在 \(\Phi\) 下满足严格平方拉回恒等式
\[
\boxed{\;
D\bigl(\Phi(X)\bigr)=\frac{C(X)^{2}}{D(X)^{3}}\ \in\ \QQ(X)\; }.
\]
并且对任意 \(P=(X,Y)\in E(\overline{\QQ})\) 且 \(Y\neq 0\)，倍点的 \(Y\)-坐标满足同一六次多项式的闭式表达
\[
\boxed{\;
Y([2]P)=\frac{C(X)}{16\,Y^{3}}\ \in\ \QQ(E)\; }.
\]
\end{conclusion}
\begin{proof}[证明要点]
对短 Weierstrass 曲线 \(Y^2=X^3+aX+b\)（此处 \(a=-1,b=1/4\)），倍点公式给出
\[
x([2]P)=\left(\frac{3X^2+a}{2Y}\right)^{\!2}-2X.
\]
代入并用 \(4Y^2=4X^3-4X+1\) 化简，即得所示 \(\Phi(X)\)。
对 \(\Phi=N/D\)（\(N=X^4+2X^2-2X+1\)、\(D=4X^3-4X+1\)）直接求导并清分母，
\(
\Phi'(X)=0\iff N'D-ND'=0
\)，化简得到所示六次多项式。
另一方面，对一般短 Weierstrass 曲线 \(Y^2=X^3+aX+b\)，四除多项式满足
\[
\psi_4=4Y\bigl(X^6+5aX^4+20bX^3-5a^2X^2-4abX-8b^2-a^3\bigr).
\]
代入 \((a,b)=(-1,\tfrac14)\) 得
\[
\psi_4(X)=2Y\bigl(2X^{6}-10X^{4}+10X^{3}-10X^{2}+2X+1\bigr),
\]
从而该六次因子刻画 \(x\bigl(E[4]\setminus E[2]\bigr)\)，与 Latt\`es 临界集合一致。
又由 division 多项式的标准恒等式
\(
Y([2]P)=\psi_4(P)\big/ \bigl(2\,\psi_2(P)^{4}\bigr)
\)
（其中 \(\psi_2(P)=2Y\)），代入上式 \(\psi_4(P)=2Y\cdot C(X)\) 即得
\(
Y([2]P)=C(X)/(16Y^{3})
\)。
\par
最后，将 \(D(X)=4X^{3}-4X+1\) 视为 \(2\)-division 三次，则
\(
D(\Phi(X))=4\Phi(X)^{3}-4\Phi(X)+1
\)；
写 \(\Phi=N/D\) 并清分母得到
\[
D(\Phi(X))=\frac{4N(X)^{3}-4N(X)D(X)^{2}+D(X)^{3}}{D(X)^{3}}.
\]
对分子直接化简与因式分解可得
\(
4N^{3}-4ND^{2}+D^{3}=C(X)^{2}
\)，从而得到 \(D(\Phi(X))=C(X)^{2}/D(X)^{3}\)。
可复现核验脚本：\texttt{scripts/exp\_fold\_zm\_elliptic\_lattes\_rational\_points\_audit.py}。证毕。
\end{proof}

\input{con__xi-terminal-zm-elliptic-lattes-pcf-portrait-resultant}

\begin{conclusion}[Latt\`es 周期点方程的整式分解：\texorpdfstring{$\Phi^n(X)=X$}{Phi^n(X)=X} 的分子由 \texorpdfstring{$(2^n\pm1)$}{(2^n±1)}-division 乘积给出]\label{con:xi-terminal-zm-elliptic-lattes-periodic-division-factorization}
沿用结论 \ref{con:xi-terminal-zm-elliptic-lattes-doubling} 的 Latt\`es 自映射
\[
\Phi(X)=\frac{X^4+2X^2-2X+1}{4X^3-4X+1}\in\QQ(X).
\]
对每个整数 \(n\ge 1\)，将 \(\Phi^n(X)-X\) 写为既约分式
\[
\Phi^n(X)-X=\frac{A_n(X)}{B_n(X)},
\qquad
A_n(X),B_n(X)\in\ZZ[X],\ \gcd(A_n,B_n)=1.
\]
对任意奇数 \(m\ge 1\)，定义 \(x\)-坐标的 \(m\)-division 整式
\[
\Psi_m(X):=\prod_{Q\in E[m]\setminus\{O\}/\{\pm1\}}\bigl(X-x(Q)\bigr)\ \in\ \ZZ[X],
\]
其零点集合按重数计恰为非平凡 \(m\)-扭点的 \(x\)-坐标像（次数为 \((m^2-1)/2\)）。
则对所有 \(n\ge 1\) 成立严格分解律
\[
\boxed{\;
A_n(X)=\pm\,\Psi_{2^n-1}(X)\,\Psi_{2^n+1}(X)\; }.
\]
特别地：
\begin{enumerate}
  \item \(n=1\)（不动点方程）：
  \[
  \Phi(X)-X=-\frac{3X^4-6X^2+3X-1}{4X^3-4X+1},
  \qquad
  A_1(X)=-(3X^4-6X^2+3X-1).
  \]
  \item \(n=2\)（二周期点方程）：
  \[
  A_2(X)=-(3X^4-6X^2+3X-1)\cdot Q_{12}(X),
  \]
  其中
  \[
  \begin{aligned}
  Q_{12}(X)=\;&5X^{12}-62X^{10}+95X^{9}-105X^{8}-60X^{7}+285X^{6}-174X^{5}\\
  &-5X^{4}-5X^{3}+35X^{2}-15X+2.
  \end{aligned}
  \]
  并且 \(Q_{12}(X)\) 的零点集合正对应于 \(5\)-扭点在 \(x\)-坐标上的像（其次数 \(12=(5^2-1)/2\) 与计数完全吻合）。
\end{enumerate}
\end{conclusion}
\begin{proof}[证明要点]
由 \(\Phi(X)=x([2]P)\)（其中 \(X=x(P)\)）可得迭代满足 \(\Phi^n(X)=x([2^n]P)\)。
因此 \(\Phi^n(X)=X\) 等价于 \(x([2^n]P)=x(P)\)，亦即 \([2^n]P=\pm P\)。
于是
\[
[2^n]P=P\ \Longrightarrow\ (2^n-1)P=O,
\qquad
[2^n]P=-P\ \Longrightarrow\ (2^n+1)P=O.
\]
反之，若 \(P\) 为 \((2^n-1)\)-扭或 \((2^n+1)\)-扭点，则显然有 \([2^n]P=\pm P\)，从而 \(X=x(P)\) 为 \(\Phi^n(X)=X\) 的解。
由于 \(2^n\pm1\) 为奇数，上述两类扭点在 \(x\)-坐标下自然按 \(\pm\) 识别，故对应的一元整式恰为 \(\Psi_{2^n-1}\) 与 \(\Psi_{2^n+1}\) 的乘积；与 \(A_n\) 的差别至多为一个整体符号。
\par
两条显式多项式分解（\(n=1,2\)）可由直接代入化简与多项式除法核验得到；并且次数计数 \(\deg \Psi_3=4\)、\(\deg \Psi_5=12\) 与 \(A_1,A_2\) 的次数相容，从而 \(Q_{12}\) 可识别为 \(\Psi_5\)（至单位因子）。
证毕。
\end{proof}

\input{con__xi-terminal-zm-elliptic-lattes-critical-sextic-galois-wreath}

\begin{conclusion}[重量参数的二次扩张与判别式的椭圆化：\(y\) 在 \texorpdfstring{$\QQ(X)$}{Q(X)} 上的最小多项式]\label{con:xi-terminal-zm-y-quadratic-extension-minpoly}\label{con:xi-terminal-zm-elliptic-quadratic-discriminant}
沿用结论 \ref{con:xi-terminal-zm-elliptic-normalization} 的坐标约定（\(X=\lambda,\ Y=w/2\)，从而 \(y=(2X^2+w-1)/2=X^2+Y-\tfrac12\)）。
令物理主支的重量参数取为
\[
\boxed{\;
y:=X^2+Y-\frac12\ \in\ \QQ(E)\; }.
\]
则 \(y\) 在 \(\QQ(X)\) 上满足严格二次最小多项式
\[
\boxed{\;
y^2-(2X^2-1)\,y+X(X-1)^2(X+1)=0\; }.
\]
其判别式恰为
\[
\boxed{\;
(2X^2-1)^2-4X(X-1)^2(X+1)
=4X^3-4X+1
=4Y^2\; }.
\]
因此重量参数的双支共轭由椭圆反演 \(\iota:(X,Y)\mapsto(X,-Y)\) 给出：
\[
y^{\iota}=(2X^2-1)-y
=X^2-Y-\frac12,
\qquad
yy^{\iota}=X(X-1)^2(X+1).
\]
从而“重量变量—谱变量”的非线性耦合在函数域层面被压缩为一个二次扩张与其椭圆判别式，结构上完全刚性。
\end{conclusion}
\begin{proof}[证明要点]
由定义 \(Y=y-X^2+\frac12\)，代入曲线方程 \(Y^2=X^3-X+\frac14\) 并展开化简，得到所示二次多项式。
判别式恒等式由直接展开验证，并用 \(4Y^2=4X^3-4X+1\) 改写即可。
共轭与乘积恒等式来自二次方程两根之和与两根之积。证毕。
\end{proof}

\begin{conclusion}[倍点拉回下重量参数的线性变换：临界六次作为线性系数]\label{con:xi-terminal-zm-elliptic-y-doubling-linear}
在函数域 \(\QQ(E)=\QQ(X,y)\) 中，倍乘自同态 \([2]:E\to E\) 对物理主支重量参数 \(y\) 的拉回满足严格恒等式
\[
\boxed{\;
[2]^*(y)=\frac{C(X)\,y+B(X)}{(4X^3-4X+1)^2}\; },
\]
其中
\[
C(X)=2X^{6}-10X^{4}+10X^{3}-10X^{2}+2X+1,
\]
\[
B(X)=-X^{8}+7X^{6}-14X^{5}+27X^{4}-9X^{3}-6X^{2}+X+1.
\]
特别地，在二次扩张 \(\QQ(E)/\QQ(X)\) 的自然基 \(\{1,y\}\) 下，\([2]^*\) 在 \(y\) 上严格呈线性；
其线性系数恰为 Latt\`es 临界六次 \(C(X)\)。
因此
\[
\boxed{\;
C(X)=0\quad\Longleftrightarrow\quad [2]^*(y)\in \QQ(X)\; } ,
\]
即倍乘 Latt\`es 的临界性可被解释为“重量二次分歧依赖在倍点拉回下发生坍缩”的代数证书。
\end{conclusion}
\begin{proof}[证明要点]
由倍点公式写出 \([2](X,Y)=(X_2,Y_2)\) 的闭式有理表达，再将 \(y_2=X_2^2+Y_2-\tfrac12\) 视为 \(\QQ(X)\)-线性函数于 \(y\)，并在四次证书 \(F(X,y)=0\) 下化简即可得到所示系数 \(C(X)\)、\(B(X)\) 与分母 \((4X^3-4X+1)^2\)。
可复现核验脚本：\texttt{scripts/exp\_fold\_zm\_elliptic\_lattes\_rational\_points\_audit.py}。证毕。
\end{proof}

\begin{conclusion}[倍点拉回在权平面四次模型中的三次/三次有理表达]\label{con:xi-terminal-zm-elliptic-y-doubling-cubic-fraction}
在结论 \ref{con:xi-terminal-zm-elliptic-quadratic-discriminant} 的首一四次表示
\[
\QQ(E)\cong \QQ(y)(X)/(F(X,y)),
\qquad
F(X,y):=X^{4}-X^{3}-(2y+1)X^{2}+X+y(y+1),
\]
下，记 \(y_2:=[2]^*(y)\in\QQ(E)\)。
则 \(y_2\) 在扩张基 \(\{1,X,X^{2},X^{3}\}\) 上可写为关于 \(X\) 的三次/三次分式：
\[
\boxed{\;
y_2=\frac{\mathcal N(X,y)}{\mathcal D(X,y)}\in \QQ(y)(X)/(F)\;},
\]
其中 \(\mathcal N,\mathcal D\in\ZZ[y][X]\) 均为关于 \(X\) 的三次多项式，并可取为
\[
\begin{aligned}
\mathcal N(X,y)=\;&(-2y^{2}-10y+1)X^{3}+(2y^{3}-6y^{2}+34y+27)X^{2}\\
&\quad+(2y^{3}+10y^{2}+12y-23)X-(y^{4}+23y^{2}+23y-1),\\[1mm]
\mathcal D(X,y)=\;&(64y+8)X^{3}+(48y^{2}+16y)X^{2}-(16y^{2}+48y+8)X-(32y^{3}+32y^{2}-1).
\end{aligned}
\]
该表示把倍点拉回的非线性压缩为 \(\QQ(y)\)-线性基上的有限阶分式表达，从而为分歧/极点的局部诊断提供直接的代数接口。
\end{conclusion}
\begin{proof}[证明要点]
由结论 \ref{con:xi-terminal-zm-elliptic-y-doubling-linear}，\(y_2=[2]^*(y)\) 在 \(\QQ(E)=\QQ(X,y)\) 中给出为 \(\QQ(X)\)-线性分式。
将该表达视为 \(\QQ(y)(X)/(F)\) 中的元素，并以 \(F(X,y)=0\) 对 \(X\) 的高次幂作降次归约后清分母，即可得到一个关于 \(X\) 次数不超过 \(3\) 的分子分母表示；
对所示 \(\mathcal N,\mathcal D\) 的逐项代入核验可由脚本 \texttt{scripts/exp\_fold\_zm\_elliptic\_lattes\_rational\_points\_audit.py} 完成。证毕。
\end{proof}

\begin{conclusion}[重量椭圆曲线的有理点结构：\texorpdfstring{$E(\QQ)_{\mathrm{tors}}=\{O\}$}{E(Q)\_tors={O}}，且物理点为无限阶]\label{con:xi-terminal-zm-elliptic-rational-points-denominator-law}
在 \(E:\ Y^2=X^3-X+\frac14\) 上，有理挠点群为平凡群
\[
\boxed{\;E(\QQ)_{\mathrm{tors}}=\{O\}\; }.
\]
将 \(E\) 通过整变换
\[
(\widetilde X,\widetilde Y)=(4X,8Y)
\]
写成积分模型
\[
\widetilde E:\ \widetilde Y^{2}=\widetilde X^{3}-16\widetilde X+16,
\qquad
\Delta(\widetilde E)=2^{12}\cdot 37,
\]
则在 \(p\neq 2,37\) 处具有良好约化；并且点计数给出
\[
\card{\widetilde E(\FF_{5})}=8,
\qquad
\card{\widetilde E(\FF_{7})}=9,
\qquad
\gcd(8,9)=1,
\]
从而推出 \(E(\QQ)_{\mathrm{tors}}=\{O\}\)。
特别地，点
\[
P=\left(2,-\frac52\right)\in E(\QQ)
\qquad(\text{等价地 }-P=(2,\tfrac52))
\]
为无限阶点，并满足 \(y(P)=1\)（其中 \(y\) 为结论 \ref{con:xi-terminal-zm-y-quadratic-extension-minpoly} 的物理主支）。
此外直接代入得 \(y(-P)=6\)。
并且其在 \(\FF_{11}\) 上的约化满足
\[
\mathrm{ord}\bigl(P\bmod 11\bigr)=17,
\]
因而 \(P\) 不可能为挠点；于是 \(\mathrm{rank}\,E(\QQ)\ge 1\)，且 \(E(\QQ)\) 为无挠的自由阿贝尔群。
写
\(
x(nP)=u_n/v_n^2
\)
为既约形式（\(v_n\in\ZZ_{>0}\)），则 \(v_n\) 构成一条椭圆分母序列，并且由
\(
y=X^2+Y-\tfrac12
\)
得到严格整性尺度律：
\[
\boxed{\;
y(nP)\in\QQ,\qquad \mathrm{den}\bigl(y(nP)\bigr)=v_n^{4}\; }.
\]
前六项可取为（列出 \(v_n\) 与 \(y(nP)\)）：
\[
\begin{array}{c|cccccc}
n & 1 & 2 & 3 & 4 & 5 & 6\\ \hline
v_n & 1 & 5 & 29 & 65 & 16264 & 2382785\\
y(nP)
& 1
& \dfrac{96}{625}
& \dfrac{2679201}{707281}
& \dfrac{252115165246}{17850625}
& \dfrac{162231294668143481}{69969611497148416}
& \dfrac{2533325779202487675774201}{32235872541947843696250625}
\end{array}
\]
相应的谱根（亦即 \(x(nP)\) 的有理值）前六项为
\[
x(P)=2,\quad
x(2P)=\dfrac{21}{25},\quad
x(3P)=\dfrac{1357}{841},\quad
x(4P)=\dfrac{480106}{4225},\quad
x(5P)=\dfrac{683916417}{264517696},\quad
x(6P)=\dfrac{4881674119706}{5677664356225}.
\]
从而“重量自由能的椭圆化”进一步提升为一条可直接进行局部—整体诊断的有理点动力学：同一物理点的倍乘轨道在分母上携带严格的四次幂尺度。
\end{conclusion}
\begin{proof}[证明要点]
（i）挠点平凡性可由有限多个良约化素数的模 \(p\) 群阶作 gcd 审计：由约化注入
\(
E(\QQ)_{\mathrm{tors}}\hookrightarrow E(\mathbb F_p)
\)
知 \(|E(\QQ)_{\mathrm{tors}}|\) 整除 \(\gcd_p \#E(\mathbb F_p)\)；
在积分模型 \(\widetilde E\) 上选取良约化素数 \(p=5,7\) 即得 \(\gcd(\#\widetilde E(\FF_5),\#\widetilde E(\FF_7))=\gcd(8,9)=1\)，从而 \(E(\QQ)_{\mathrm{tors}}=\{O\}\)。
（ii）由 \(\mathrm{ord}(P\bmod 11)=17\) 可直接排除 \(P\) 为挠点，从而 \(P\) 为无限阶点，\(\mathrm{rank}\,E(\QQ)\ge 1\)。
（iii）令 \(x(nP)=u_n/v_n^2\)、\(Y(nP)=r_n/v_n^3\) 为既约表示，则
\(
y(nP)=x(nP)^2+Y(nP)-\tfrac12
\)
的分母整除 \(v_n^4\)；
对该轨道可核验得到精确等号 \(\mathrm{den}(y(nP))=v_n^4\)，并可回收所示前六项。
可复现核验脚本：\texttt{scripts/exp\_fold\_zm\_elliptic\_lattes\_rational\_points\_audit.py}。证毕。
\end{proof}

\input{con__xi-terminal-zm-elliptic-mw-height-denominator-growth}

\begin{conclusion}[物理基点诱导的椭圆可除序列：强可除性与低阶原始素因子]\label{con:xi-terminal-zm-elliptic-eds-strong-divisibility}
在结论 \ref{con:xi-terminal-zm-elliptic-rational-points-denominator-law} 的积分模型 \(\widetilde E\) 上，令
\(
\widetilde P:=(8,20)\in\widetilde E(\QQ)
\)
对应于 \(P\)。
对 \(n\ge 1\) 将 \(\widetilde X\bigl(n\widetilde P\bigr)\) 写成既约分式
\[
\widetilde X\bigl(n\widetilde P\bigr)=\frac{A_n}{B_n^2},
\qquad
A_n,B_n\in\ZZ,\quad \gcd(A_n,B_n)=1,\quad B_n>0,
\]
则 \(\bigl(B_n\bigr)_{n\ge 1}\) 构成椭圆可除序列，并满足强可除性
\[
\gcd(B_m,B_n)=B_{\gcd(m,n)},
\qquad
m\mid n\Longrightarrow B_m\mid B_n.
\]
其前十项为
\[
\begin{aligned}
B_1&=1,&
B_2&=5,&
B_3&=29,&
B_4&=65,&
B_5&=8132,\\
B_6&=2382785,&
B_7&=398035821,&
B_8&=43244638645,&
B_9&=124106986093951,&
B_{10}&=270525565025400200.
\end{aligned}
\]
并且在低阶上已出现若干原始分母素因子，例如
\[
29\mid B_3,\qquad 13\mid B_4,\qquad 19\mid B_5,\qquad 17\mid B_9.
\]
其中 \(\{13,17,19,29\}\) 同时属于 \(E_8\) Coxeter 指数集合 \(\{1,7,11,13,17,19,23,29\}\) 的素数子集，从而在极低阶即可给出一个可审计的指数素数共现样本。
\end{conclusion}
\begin{proof}[证明要点]
强可除性与可除性为椭圆可除序列的一般性质；
上述前十项与所示整除关系可由对 \(\widetilde P\) 的倍乘轨道直接计算并在 \(\ZZ\) 中分解验证。
可复现核验脚本：\texttt{scripts/exp\_fold\_zm\_elliptic\_lattes\_rational\_points\_audit.py}。证毕。
\end{proof}

\begin{conclusion}[分母素因子与模 \texorpdfstring{$p$}{p} 阶的严格等价：\texorpdfstring{$p\mid B_n$}{p|Bn} 的群论解码]\label{con:xi-terminal-zm-elliptic-eds-modp-order-equivalence}
设 \(p\neq 2,37\) 为 \(\widetilde E\) 的良好约化素数。
则对任意 \(n\ge 1\) 有严格等价
\[
p\mid B_n
\ \Longleftrightarrow\
n\widetilde P\equiv O\pmod p
\ \Longleftrightarrow\
\mathrm{ord}_{\widetilde E(\FF_p)}(\widetilde P)\mid n.
\]
因此 \(\bigl(B_n\bigr)\) 的素因子集合精确记录了 \(\widetilde P\) 在各有限域点群 \(\widetilde E(\FF_p)\) 中的阶分布；任何能够检测 \(p\mid B_n\) 的局部“模 \(p\) 指纹”均可等价转写为一个纯群论命题“\(\mathrm{ord}(\widetilde P\bmod p)\) 是否整除 \(n\)”。
\end{conclusion}
\begin{proof}[证明要点]
在良好约化处，\(\widetilde E(\QQ)\to\widetilde E(\FF_p)\) 的约化映射对分母的 \(p\)-进信息具有精确判别：\(p\mid B_n\) 当且仅当 \(\widetilde X(n\widetilde P)\) 在 \(p\)-进赋值下出现极点，这等价于 \(n\widetilde P\) 约化为无穷远点 \(O\)。
最后一条等价为有限阿贝尔群中元素阶的定义。证毕。
\end{proof}

\input{con__xi-terminal-zm-elliptic-badprime-37-periodicity}

\begin{conclusion}[非复乘与原始素因子定理：椭圆可除序列的原始素因子无穷性]\label{con:xi-terminal-zm-elliptic-eds-primitive-primes-noncm}
结论 \ref{con:xi-terminal-zm-elliptic-normalization} 给出
\(
j(E)=110592/37
\)
非代数整数，故 \(E\) 无复乘。
因此由非复乘椭圆可除序列的原始素因子定理，存在常数 \(n_0\) 使得对一切 \(n>n_0\)，\(B_n\) 含有原始素因子 \(p_n\)（即 \(p_n\mid B_n\) 且 \(p_n\nmid\prod_{m<n}B_m\)），并由结论 \ref{con:xi-terminal-zm-elliptic-eds-modp-order-equivalence} 得
\[
\mathrm{ord}_{\widetilde E(\FF_{p_n})}(\widetilde P)=n.
\]
从而该椭圆结构在严格意义上产生无穷多个原始素因子素数，其上 \(\widetilde P\) 具有预设的精确阶，为后续一切模 \(p\) 指纹检验提供可控的“阶—素数”供给机制。
\end{conclusion}

\begin{conclusion}[重量配分函数的域迹表示：分母仅由 Lee--Yang 三次因子控制]\label{con:xi-terminal-zm-field-trace-representation}
沿用结论 \ref{con:xi-terminal-zm-order4-rational} 的记号，令
\[
\Pi(\lambda,y):=\lambda^4-\lambda^3-(2y+1)\lambda^2+\lambda+y(y+1),
\qquad
P_{\mathrm{LY}}(y):=256y^3+411y^2+165y+32.
\]
记 \(K:=\QQ(y)(\lambda)\)（其中 \(\lambda\) 为 \(\Pi(\lambda,y)=0\) 的一根），并令
\(
s_m(y):=\mathrm{Tr}_{K/\QQ(y)}(\lambda^m)
=\sum_{i=1}^{4}\lambda_i(y)^m
\)
为四根的幂和（Newton 幂和）。
则存在唯一的元素 \(u(\lambda)\in K\)，其在基 \(\{1,\lambda,\lambda^2,\lambda^3\}\) 下的代表可取为三次多项式
\[
u(\lambda)=a_0(y)+a_1(y)\lambda+a_2(y)\lambda^2+a_3(y)\lambda^3,
\]
使得对一切 \(m\ge 0\) 有严格域迹表示
\[
\boxed{\;
Z_m(y)=\mathrm{Tr}_{K/\QQ(y)}\!\bigl(u(\lambda)\lambda^m\bigr)\; }.
\]
并且 \(u\) 的系数可显式取为
\[
\boxed{\;
a_0(y)=\frac{-19y^3-75y^2-42y-8}{P_{\mathrm{LY}}(y)},\qquad
a_2(y)=\frac{87y^2+48y+9}{P_{\mathrm{LY}}(y)}\; } ,
\]
\[
\boxed{\;
a_1(y)=\frac{80y^4-20y^3-82y^2-42y-8}{y\,P_{\mathrm{LY}}(y)},\qquad
a_3(y)=\frac{-16y^3+49y^2+31y+8}{y\,P_{\mathrm{LY}}(y)}\; } .
\]
因此尽管 \(Z_m(y)\in\ZZ[y]\) 为整系数多项式，其“谱投影权”在函数域上只可能在
\(
y=0
\)
与
\(
P_{\mathrm{LY}}(y)=0
\)
处出现极点；这为 Lee--Yang 分歧点在代数层面的不可避免性给出一条直接证书。
\end{conclusion}
\begin{proof}[证明要点]
对未知 \(u=b_0+b_1\lambda+b_2\lambda^2+b_3\lambda^3\in K\)，有恒等式
\[
\mathrm{Tr}(u\lambda^m)=b_0\,s_m+b_1\,s_{m+1}+b_2\,s_{m+2}+b_3\,s_{m+3}.
\]
由于 \(m\mapsto Z_m(y)\) 与 \(m\mapsto s_m(y)\) 均满足由 \(\Pi(\lambda,y)\) 诱导的同一四阶递推，故只需对 \(m=0,1,2,3\) 匹配即可唯一确定 \((b_0,\dots,b_3)\)；
解线性系统并化简即可得到所示 \(a_i(y)\)。
可复现核验脚本：\texttt{scripts/exp\_fold\_zm\_trace\_galois\_audit.py}。证毕。
\end{proof}

\begin{conclusion}[清分母后的可审计整式恒等式：\texorpdfstring{$yP_{\mathrm{LY}}(y)Z_m$}{yP(y)Z\_m} 为幂和的整系数组合]\label{con:xi-terminal-zm-integral-identity-power-sums}
在结论 \ref{con:xi-terminal-zm-field-trace-representation} 的记号下，乘以 \(yP_{\mathrm{LY}}(y)\) 可将域迹公式提升为整式恒等式：对一切 \(m\ge 0\) 有
\[
\boxed{\;
yP_{\mathrm{LY}}(y)\,Z_m(y)
=A_0(y)\,s_m(y)+A_1(y)\,s_{m+1}(y)+A_2(y)\,s_{m+2}(y)+A_3(y)\,s_{m+3}(y)\;}
\]
其中系数多项式为
\[
\boxed{\;
\begin{aligned}
A_0(y)&=-19y^4-75y^3-42y^2-8y,\\
A_1(y)&=80y^4-20y^3-82y^2-42y-8,\\
A_2(y)&=87y^3+48y^2+9y,\\
A_3(y)&=-16y^3+49y^2+31y+8.
\end{aligned}}
\]
该恒等式把“重量计数”（左端）与“谱幂和”（右端）在整系数层面直接对接，并将所有潜在的解析障碍压缩为显式代数集合 \(yP_{\mathrm{LY}}(y)=0\)。
\end{conclusion}
\begin{proof}
将结论 \ref{con:xi-terminal-zm-field-trace-representation} 的 \(a_i(y)\) 代入
\(
\mathrm{Tr}(u\lambda^m)=a_0 s_m+a_1 s_{m+1}+a_2 s_{m+2}+a_3 s_{m+3}
\)，并清分母即可得到。证毕。
\end{proof}

\begin{conclusion}[迹函数解释：四阶递推的几何来源与奇点集合的约束]\label{con:xi-terminal-zm-trace-pushforward-origin}
在结论 \ref{con:xi-terminal-zm-elliptic-normalization} 的椭圆化识别下，谱曲线的归一化可视为固定椭圆曲线 \(E\) 上的次数 \(4\) 有理映射
\(
y\in\QQ(E)
\)。
对任意 \(y_0\in\CC\setminus\{0,1\}\) 且 \(P_{\mathrm{LY}}(y_0)\neq 0\)，纤维 \(y^{-1}(y_0)\subset E(\CC)\) 恰含四个点；
其 \(X\)-坐标集合与四次方程 \(\Pi(\lambda,y_0)=0\) 的四根集合一致。
因此 \(\lambda\) 作为 \(y\) 的四值代数函数之全分支 Riemann 面即为 \(E(\CC)\)（等价于结论 \ref{con:xi-terminal-zm-elliptic-normalization} 的归一化描述），并由椭圆判别式 \(\Delta(E)=37\) 的统一化所锚定。
\par
在函数域层面，扩张 \(\QQ(E)/\QQ(y)\) 为次数 \(4\) 的有限可分扩张，其原始元可取为 \(X=\lambda\in\QQ(E)\)。
对每个整数 \(m\ge 0\) 定义迹函数
\[
T_m(y):=\Tr_{\QQ(E)/\QQ(y)}(X^m)\in\QQ(y).
\]
则 \((T_m)_{m\ge 0}\) 满足由 \(\Pi(\lambda,y)\) 决定的四阶线性递推（其特征多项式即为 \(\Pi\)）：
\[
T_m(y)=T_{m-1}(y)+(2y+1)T_{m-2}(y)-T_{m-3}(y)-y(y+1)T_{m-4}(y)\qquad(m\ge 4).
\]
此外，作为 \(\mathbb P^1_y\) 上的有理函数，\(T_m\) 的可能极点仅位于 \(y=\infty\)；
因而其奇点集合包含于分歧值集合
\(
\{\infty,0,1\}\cup\{P_{\mathrm{LY}}(y)=0\}
\)（见结论 \ref{con:xi-terminal-zm-elliptic-y-hurwitz-passport}）。
结合结论 \ref{con:xi-terminal-zm-field-trace-representation} 的域迹表示，四阶递推可解释为次数 \(4\) 覆叠 \(\pi_y:E\to\mathbb P^1_y\) 上的迹推送结构，而 Lee--Yang 三次因子则在该覆叠的分歧几何中被内在锁定。
\end{conclusion}
\begin{proof}[证明要点]
由 \(\Pi(X,y)=0\) 在 \(\QQ(E)\) 中成立，乘以 \(X^{m-4}\) 并取 \(\Tr_{\QQ(E)/\QQ(y)}\) 得到所示递推。
又由于 \(X\) 在 \(E\) 上仅在无穷远点 \(O\) 处有极点，而 \(y\) 在 \(O\) 处具有四阶极点（结论 \ref{con:xi-terminal-zm-elliptic-normalization}），故 \(T_m\) 在有限 \(y\)-点处不出现极点，从而其可能极点仅在 \(y=\infty\)。
可复现核验脚本：\texttt{scripts/exp\_fold\_zm\_elliptic\_audit.py}。证毕。
\end{proof}

\begin{conclusion}[谱的最大代数复杂度：\texorpdfstring{$\mathrm{Gal}_{\QQ(y)}\!\bigl(\Pi(\lambda,y)\bigr)\cong S_4$}{Gal(Q(y))(Pi)=S4}]\label{con:xi-terminal-zm-generic-galois-s4}
把 \(\Pi(\lambda,y)\) 视为关于 \(\lambda\) 的四次多项式。
其三次分解式（cubic resolvent）可取为
\[
\boxed{\;
\mathscr R(z,y)
=z^3+(1+2y)z^2-(1+4y+4y^2)z-(1+5y+13y^2+8y^3)\; }.
\]
其判别式满足严格恒等式
\[
\boxed{\;
\mathrm{Disc}_{z}\bigl(\mathscr R(\cdot,y)\bigr)
=-y(y-1)\,P_{\mathrm{LY}}(y)\; } ,
\]
并且与四次谱方程的判别式完全一致：
\[
\boxed{\;
\mathrm{Disc}_{z}\bigl(\mathscr R(\cdot,y)\bigr)
=\mathrm{Disc}_{\lambda}\bigl(\Pi(\cdot,y)\bigr)\; } .
\]
因此在 \(y\)-线上，四次谱覆盖的分歧值集合恰为 \(\{0,1\}\) 与 \(P_{\mathrm{LY}}(y)=0\) 的三根；Lee--Yang 三次因子因而给出判别式的不可约分歧核。
在任一良特化（例如 \(y=2\)）下，
\[
\mathscr R(z,2)=z^{3}+5z^{2}-25z-127
\]
在 \(\QQ\) 上不可约且其判别式为 \(-8108\)，因而非平方，
从而 \(\Pi(\lambda,y)\) 在 \(\QQ(y)\) 上的泛型伽罗瓦群为
\[
\boxed{\;\mathrm{Gal}\bigl(\Pi/\QQ(y)\bigr)\cong S_4.\;}
\]
因此 \(\Pi(\lambda,y)\) 在 \(\QQ(y)\) 上给出一个正则 \(S_4\)-扩张（常数域为 \(\QQ\)），而其唯一二次分辨率由 \(-y(y-1)P_{\mathrm{LY}}(y)\) 锁定，直接连接到 Lee--Yang 分歧集合。
等价地，谱支 \(\lambda_i(y)\) 的单值解析延拓所诱导的单值群在四个谱分支上给出满的 \(S_4\)-作用；因而不存在任何将谱覆盖全局分解为两个互不混合的二次子覆盖的代数约化能够成立。
\end{conclusion}
\begin{proof}[证明要点]
分解式 \(\mathscr R(z,y)\) 与判别式恒等式可由直接符号计算得到，并与结论 \ref{con:xi-terminal-zm-discriminant-leyang} 的判别式分解一致。
若在某特化 \(y=y_0\) 下 \(\mathscr R(z,y_0)\) 在 \(\QQ[z]\) 上不可约，则 \(\mathscr R\) 在 \(\QQ(y)[z]\) 上不可约；
对不可约四次而言，三次分解式不可约将伽罗瓦群约束为 \(\{A_4,S_4\}\)；
而判别式非平方排除 \(A_4\)，故必为 \(S_4\)。
可复现核验脚本：\texttt{scripts/exp\_fold\_zm\_trace\_galois\_audit.py}。证毕。
\end{proof}

\begin{remark}[Hitchin 谱覆盖与 \texorpdfstring{$A_3$}{A3} cameral 覆叠：\texorpdfstring{$S_4\simeq W(A_3)$}{S4=W(A3)} 的结构性解释]\label{rem:xi-terminal-zm-hitchin-cameral-cover}
结论 \ref{con:xi-terminal-zm-generic-galois-s4} 给出的泛型对称群 \(S_4\) 可被理解为一种 Weyl 单值化：在代数几何层面，四次谱覆盖 \(\Pi(\lambda,y)=0\) 的伽罗瓦闭包给出一个 \(S_4\)-cameral 覆叠，其群作用与 \(A_3\) 型 Weyl 群 \(W(A_3)\cong S_4\) 同构。
若将 \(\Pi(\lambda,y)=0\) 解释为某一秩 \(4\) Higgs 场（例如 \(G=\mathrm{SL}_4\) 或 \(\mathrm{PGL}_4\)）的特征方程在 Hitchin 基底上的一维切片，则该 \(S_4\) 单值化即对应 cameral cover 的 Weyl 对称；
而同一判别式核 \(-y(y-1)P_{\mathrm{LY}}(y)\) 则给出该切片与 Hitchin 判别式除子的交，从而把 Lee--Yang 分歧集合内在地识别为谱覆盖奇纤维的显式坐标化。
几何朗兰兹的物理实现中，Hitchin 纤维化与其谱/cameral 数据是刻画拓扑扭曲后真空结构与缺陷算子的重要几何输入。\cite{KapustinWitten2006GeometricLanglands}
\end{remark}

\begin{conclusion}[有理参数特化的四次整式模型与判别式：分歧素集的显式上界]\label{con:xi-terminal-zm-specialization-discriminant-ab}
取有理参数 \(y_0=a/b\in\QQ\)（\(\gcd(a,b)=1,\ b>0\)），并考虑清分母后的四次整式
\[
P_{a,b}(\lambda):=b^{2}\Pi(\lambda,a/b)
=b^{2}\lambda^{4}-b^{2}\lambda^{3}-(2ab+b^{2})\lambda^{2}+b^{2}\lambda+a(a+b)\in\ZZ[\lambda].
\]
则其关于 \(\lambda\) 的判别式满足严格闭式
\[
\boxed{\;
\mathrm{Disc}_{\lambda}\bigl(P_{a,b}\bigr)
=-a(a-b)\bigl(256a^{3}+411a^{2}b+165ab^{2}+32b^{3}\bigr)\,b^{7}\; }.
\]
因而当且仅当 \(a=0\)、\(a=b\) 或 \(256a^{3}+411a^{2}b+165ab^{2}+32b^{3}=0\) 时出现重根；
在其余情形，\(P_{a,b}\) 在 \(\overline\QQ\) 上可分（等价地：特化覆叠在该参数处 \'{e}tale）。
此外，若素数 \(p\) 在 \(P_{a,b}\) 的分裂域中分歧，则必整除
\[
2ab(a-b)\bigl(256a^{3}+411a^{2}b+165ab^{2}+32b^{3}\bigr),
\]
从而分歧素集被上述显式整数所包含。
结合结论 \ref{con:xi-terminal-zm-generic-galois-s4} 的泛型 \(S_4\) 性，由 Hilbert 不可约定理可知：除去薄集参数外，特化多项式 \(P_{a,b}\) 的分裂域伽罗瓦群仍为 \(S_4\)，
从而得到一族由 \(y_0\in\QQ\) 参数化的 \(S_4\) 数域族，其分歧由同一 Lee--Yang 三次核在清分母后完全控制。
\end{conclusion}
\begin{proof}[证明要点]
由结论 \ref{con:xi-terminal-zm-discriminant-leyang}，
\(
\mathrm{Disc}_{\lambda}\bigl(\Pi(\cdot,y)\bigr)=-y(y-1)P_{\mathrm{LY}}(y)
\)。
代入 \(y=a/b\) 并利用四次情形的齐次性
\(
\mathrm{Disc}_{\lambda}(c\cdot f)=c^{6}\mathrm{Disc}_{\lambda}(f)
\)
得到
\(
\mathrm{Disc}_{\lambda}\bigl(P_{a,b}\bigr)=b^{12}\mathrm{Disc}_{\lambda}\bigl(\Pi(\cdot,a/b)\bigr)
\)，
化简即得闭式公式。
分歧素集包含性来自一般判据“分裂域的分歧素数必整除多项式判别式”。
薄集断言由泛型伽罗瓦群与 Hilbert 不可约定理推出。
可复现核验脚本：\texttt{scripts/exp\_fold\_zm\_pi\_specialization\_discriminant\_ab\_audit.py}。证毕。
\end{proof}

\begin{conclusion}[“椭圆—\texorpdfstring{$S_4$}{S4}—Lee--Yang”三重对偶：同一判别式同时支配归一化分歧与分裂域二次脊柱]\label{con:xi-terminal-zm-elliptic-s4-leyang-triple-duality}
曲线 \(\Pi(\lambda,y)=0\) 的几何归一化为椭圆曲线 \(E\)（结论 \ref{con:xi-terminal-zm-elliptic-normalization}），
而同一四次谱方程在 \(\QQ(y)\) 上具有 \(S_4\) 伽罗瓦群（结论 \ref{con:xi-terminal-zm-generic-galois-s4}）。
两者在判别式层面被同一多项式
\[
-y(y-1)\,P_{\mathrm{LY}}(y)
\]
同时支配：
\begin{enumerate}
  \item 作为几何对象，它给出投影 \(\pi:E\to\mathbb P^1_y\) 的分歧除子（见结论 \ref{con:xi-terminal-zm-leyang-cubic-branch-image} 与结论 \ref{con:xi-terminal-zm-discriminant-leyang}）；
  \item 作为代数对象，它给出 \(S_4\) 分裂域的唯一二次分辨率子域 \(\QQ\bigl(\sqrt{-y(y-1)P_{\mathrm{LY}}(y)}\bigr)\)。
\end{enumerate}
因此 Lee--Yang 三次因子 \(P_{\mathrm{LY}}(y)\) 并非“热力学支的偶发奇点”，而是一个同时控制椭圆归一化的分歧几何与四次分裂域二次脊柱的共同判别式核；
从而把“解析临界”与“伽罗瓦分辨率”在结构上锁定为同一条不可约代数数据。
\end{conclusion}

\input{con__xi-terminal-zm-elliptic-branchpoint-specializations-genus2-s4}


\begin{conclusion}[Lee--Yang 六次闭包与 Latt\`es 八面体闭包的 Chebotarev 独立：完全分裂素数密度 \texorpdfstring{$1/288$}{1/288}]\label{con:xi-terminal-zm-leyang-lattes-chebotarev-independence}
令
\[
P_{\mathrm{LY}}(y):=256y^{3}+411y^{2}+165y+32\in\ZZ[y],
\qquad
C(X):=2X^{6}-10X^{4}+10X^{3}-10X^{2}+2X+1\in\ZZ[X].
\]
记 \(L_{\mathrm{LY}}:=\mathrm{Spl}(P_{\mathrm{LY}})\) 与 \(L_{4}:=\mathrm{Spl}(C)\) 为其分裂域。
则：
\begin{enumerate}
  \item \(\mathrm{Gal}(L_{\mathrm{LY}}/\QQ)\cong S_{3}\)，其唯一二次子域为 \(\QQ(\sqrt{-111})\)，并且 \(3\) 在该二次域中分歧（结论 \ref{con:xi-terminal-zm-leyang-cubic-field-minimal-discriminant}）。
  \item \(\mathrm{Gal}(L_{4}/\QQ)\cong S_{4}\times C_{2}\)（阶 \(48\)），并且 \(L_{4}/\QQ\) 的分歧素数集合包含于 \(\{2,37\}\)（结论 \ref{con:xi-terminal-zm-elliptic-lattes-critical-sextic-galois}）。
  \item 因而两闭包在 \(\QQ\) 上线性互不相交：\(L_{\mathrm{LY}}\cap L_{4}=\QQ\)。
  等价地，其复合扩张满足
  \[
  \mathrm{Gal}(L_{\mathrm{LY}}L_{4}/\QQ)\ \cong\ (S_{4}\times C_{2})\times S_{3},
  \qquad [L_{\mathrm{LY}}L_{4}:\QQ]=48\cdot 6=288.
  \]
\end{enumerate}
因此对任意素数 \(p\notin\{2,3,31,37\}\)，满足
\[
P_{\mathrm{LY}}(y)\ \text{与}\ C(X)\ \text{同时在}\ \FF_{p}\ \text{上完全分裂}
\]
的素数集合具有自然密度
\[
\boxed{\;\delta=\frac{1}{288}\;}
\]
（Chebotarev 意义下的“同时落在恒等共轭类”事件）。
该密度律把 Lee--Yang 三分歧值的全体有理性与 Latt\`es 临界六次的全体可见性在素数层面严格解耦，并给出一条可复核的统计预言。
\end{conclusion}
\begin{proof}[证明要点]
由结论 \ref{con:xi-terminal-zm-leyang-cubic-field-minimal-discriminant}，\(L_{\mathrm{LY}}/\QQ\) 为 \(S_3\)-伽罗瓦扩张且其唯一二次子域为 \(\QQ(\sqrt{-111})\)，其中素数 \(3\) 分歧。
而由结论 \ref{con:xi-terminal-zm-elliptic-lattes-critical-sextic-galois}，\(L_4/\QQ\) 仅在 \(\{2,37\}\) 处分歧，故 \(3\) 在 \(L_4\) 的任一子扩张中亦不分歧。
由于 \(L_{\mathrm{LY}}\) 与 \(L_4\) 均为 \(\QQ\) 上的伽罗瓦扩张，交域 \(L_{\mathrm{LY}}\cap L_4\) 亦为 \(\QQ\) 上伽罗瓦扩张，故作为 \(L_{\mathrm{LY}}\) 的伽罗瓦子扩张只能为 \(\QQ\) 或 \(\QQ(\sqrt{-111})\)；
但后者会引入素数 \(3\) 的分歧，与 \(L_4\) 的分歧素集矛盾，遂 \(L_{\mathrm{LY}}\cap L_4=\QQ\)。
因此复合伽罗瓦群为直积并给出次数 \(288\)。
\par
对 \(p\notin\{2,3,31,37\}\)，两扩张在 \(p\) 处均未分歧。
“在 \(\FF_p\) 上完全分裂”等价于 Frobenius 落在恒等共轭类；
在线性互不相交下，复合扩张的 Frobenius 共轭类由两边的 Frobenius 独立确定，故同时全分裂的密度为
\(
1/|\,\mathrm{Gal}(L_{\mathrm{LY}}L_4/\QQ)\,|=1/288
\)。
可复现核验脚本：\texttt{scripts/exp\_fold\_zm\_leyang\_lattes\_joint\_splitting\_density\_audit.py}（输出证书：\texttt{artifacts/export/fold\_zm\_leyang\_lattes\_joint\_splitting\_density\_audit.json}）。
证毕。
\end{proof}


\begin{conclusion}[三次预解式曲线的平面三次结构与无穷远纤维的分歧型 \texorpdfstring{$(2,1)$}{(2,1)}]\label{con:xi-terminal-zm-resolvent-cubic-plane-cubic-geometry}
沿用结论 \ref{con:xi-terminal-zm-generic-galois-s4} 的三次预解式（cubic resolvent）
\[
\mathscr R(z,y)
:=z^3+(1+2y)z^2-(1+4y+4y^2)z-(1+5y+13y^2+8y^3)\in\ZZ[z,y],
\]
并令 \(\mathcal R\subset \mathbb A^2_{z,y}\) 为仿射曲线 \(\mathscr R(z,y)=0\)，令 \(\overline{\mathcal R}\subset \mathbb P^2_{[Z:Y:W]}\) 为其齐次化闭包（\(z=Z/W,\ y=Y/W\)）。
记 \(\Pi(\lambda,y)\) 为结论 \ref{con:xi-terminal-zm-order4-rational} 的谱四次多项式，\(\lambda_1,\dots,\lambda_4\) 为 \(\Pi(\lambda,y)=0\) 的四根（于代数闭包中）。
定义三组配对不变量
\[
\lambda_1\lambda_2+\lambda_3\lambda_4,\qquad
\lambda_1\lambda_3+\lambda_2\lambda_4,\qquad
\lambda_1\lambda_4+\lambda_2\lambda_3.
\]
则它们恰为 \(\mathscr R(z,y)\) 的三根；因而 \(\mathscr R\) 作为 Ferrari 预解三次在代数层面压缩了四次谱覆盖的配对张量不变量。
并且在判别式层面存在严格同一性
\[
\mathrm{Disc}_{z}\bigl(\mathscr R(\cdot,y)\bigr)=\mathrm{Disc}_{\lambda}\bigl(\Pi(\cdot,y)\bigr),
\]
从而 Lee--Yang 型分歧集合亦可等价刻画为该三次预解式退化（出现重根）的参数集合。
则：
\begin{enumerate}
  \item \(\overline{\mathcal R}\) 为非奇异平面三次曲线，且含有有理点 \((y,z)=(0,1)\)（因为 \(\mathscr R(1,0)=0\)），从而 \(\overline{\mathcal R}\) 本身即为一条定义在 \(\QQ\) 上的椭圆曲线（无需再作归一化）。
  \item 投影
  \[
  \pi^{(3)}_y:\ \overline{\mathcal R}\longrightarrow \mathbb P^1_y,\qquad [Z:Y:W]\longmapsto [Y:W]
  \]
  为次数 \(3\) 的覆盖。其分歧像由判别式严格刻画：视 \(\mathscr R\) 为关于 \(z\) 的三次多项式，则
  \[
  \mathrm{Disc}_{z}\bigl(\mathscr R(\cdot,y)\bigr)=-y(y-1)P_{\mathrm{LY}}(y),
  \]
  因而在 \(y=0\)、\(y=1\) 与 \(P_{\mathrm{LY}}(y)=0\) 的三点处发生简单分歧（分歧指数 \(2\)）。
  \item 在无穷远纤维 \(y=\infty\) 上，\(\pi^{(3)}_y\) 的纤维由两个无穷远点构成，且极点阶数分布为 \(2+1\)（分歧型 \((2,1)\)）。
  更精确地，直线 \(W=0\) 与 \(\overline{\mathcal R}\) 的交点为
  \[
  [2:1:0]\ \text{（交数 \(1\)）},\qquad [-2:1:0]\ \text{（交数 \(2\)）},
  \]
  故 \([-2:1:0]\) 为 \(y=\infty\) 上唯一的分歧点（分歧指数 \(2\)），而 \([2:1:0]\) 为非分歧点。
\end{enumerate}
\end{conclusion}
\begin{proof}[证明要点]
注意到 \(\mathscr R(1,0)=0\)，故 \((y,z)=(0,1)\in \mathcal R(\QQ)\)，从而 \(\overline{\mathcal R}(\QQ)\neq\varnothing\)。
齐次化取
\[
\mathscr R_h(Z,Y,W)
=Z^3+(W+2Y)Z^2-(W^2+4YW+4Y^2)Z-(W^3+5YW^2+13Y^2W+8Y^3).
\]
在无穷远直线 \(W=0\) 上，有因式分解
\[
\mathscr R_h(Z,Y,0)=Z^3+2YZ^2-4Y^2Z-8Y^3=(Z-2Y)(Z+2Y)^2,
\]
从而得到两点 \([2:1:0]\)、\([-2:1:0]\) 及其交数 \(1,2\)。
由 \(y=Y/W\) 可知该交数即为 \(y\) 在相应点处的极点阶数，
故 \(y=\infty\) 的纤维型为 \(2+1\)，且仅 \([-2:1:0]\) 对应分歧指数 \(2\)。
\par
作为关于 \(z\) 的三次多项式，其判别式恒等式已在结论 \ref{con:xi-terminal-zm-generic-galois-s4} 给出，故有限分歧值集合与分歧指数结论随之成立。
最后，对 \(\overline{\mathcal R}\) 的光滑性可在仿射图 \(W=1\) 与无穷远图 \(Y=1\) 上分别检验梯度不同时为零；
等价地，对 \(\mathscr R_h=0\) 的射影曲线，联立 \(\mathscr R_h=\partial_Z\mathscr R_h=\partial_Y\mathscr R_h=\partial_W\mathscr R_h=0\) 在代数闭包中无解。
证毕。
\end{proof}

\begin{conclusion}[三次预解式曲线的 Weierstrass 化与显式双有理同一]\label{con:xi-terminal-zm-resolvent-cubic-weierstrass-birational}
沿用结论 \ref{con:xi-terminal-zm-resolvent-cubic-plane-cubic-geometry} 的记号。
则平面三次曲线 \(\overline{\mathcal R}/\QQ\) 与下述椭圆曲线 \(E_{\mathscr R}/\QQ\) 双有理同构：
\[
E_{\mathscr R}:\quad Y^2=X^3-37X^2+307X+361.
\]
并且可取一组完全显式的双有理互逆变换：在 \(E_{\mathscr R}\) 上令
\[
y=-\frac{8X}{Y+19+14X-X^2},
\qquad
z=1-\frac{2X(7-X)}{Y+19+14X-X^2},
\]
则 \((y,z)\) 满足 \(\mathscr R(z,y)=0\)，从而给出 \(E_{\mathscr R}\dashrightarrow \overline{\mathcal R}\)；
反之，在开子集 \(y\neq 0\) 上，取
\[
X=\frac{7y+4-4z}{y},
\qquad
Y=X^2-14X-19-\frac{8X}{y},
\]
则 \((X,Y)\) 满足 \(E_{\mathscr R}\) 的 Weierstrass 方程，从而给出 \(\overline{\mathcal R}\dashrightarrow E_{\mathscr R}\)。
\par
此外，\(E_{\mathscr R}\) 的基本不变量可显式写为
\[
c_4=7168,\quad c_6=-341504,\quad \Delta=2^{12}\cdot 31^{2}\cdot 37,\quad
j(E_{\mathscr R})=\frac{2^{18}\cdot 7^{3}}{31^{2}\cdot 37}.
\]
\end{conclusion}
\begin{proof}[证明要点]
由所给双有理变换将 \((X,Y)\) 代入 \((y,z)\) 并清分母，可验证 \(\mathscr R(z,y)=0\) 在理想 \((Y^2-(X^3-37X^2+307X+361))\) 中为零；
反向同理由直接代入并清分母得到。
判别式、\(c_4,c_6\) 与 \(j\) 的显式表达来自对该 Weierstrass 模型套用标准不变量公式。
可复现核验脚本：\texttt{scripts/exp\_fold\_zm\_resolvent\_cubic\_weierstrass\_audit.py}。证毕。
\end{proof}

\begin{conclusion}[分解式椭圆因子的双二次模型与短 Weierstrass 正规形]\label{con:xi-terminal-zm-resolvent-cubic-biquadratic-invariants-short-weierstrass}
沿用结论 \ref{con:xi-terminal-zm-resolvent-cubic-plane-cubic-geometry} 的 \((y,z)\) 坐标，并在开子集 \(z\neq 1\) 上置
\[
x:=\frac{y}{z-1}.
\]
则对过基点 \((y,z)=(0,1)\) 的直线族 \(z=1+\frac{y}{x}\) 代入 \(\mathscr R(z,y)=0\) 并约去因子 \(y\)，可得关于 \(y\) 的二次方程
\[
(1+2x-4x^2-8x^3)y^2+(4x-17x^3)y+(4x^2-7x^3)=0.
\]
令
\[
u:=\frac{2(1+2x-4x^2-8x^3)y+(4x-17x^3)}{x},
\]
则 \((x,u)\) 满足双二次模型
\[
\mathcal Q:\quad
u^2=x\bigl(65x^3+16x^2-16x-4\bigr)
=65x^4+16x^3-16x^2-4x,
\]
并与平面三次 \(\overline{\mathcal R}\) 在 \(\QQ\) 上双有理等价。
对二元四次 \(65x^4+16x^3-16x^2-4x\) 的经典不变量，有
\[
I=448,\qquad J=-10672,\qquad 4I^3-J^2=2^8\cdot 3^3\cdot 31^2\cdot 37,
\]
从而其 Jacobian 与 \(E_{\mathscr R}\) 在 \(\QQ\) 上同构，并同具
\[
j=\frac{2^{18}\cdot 7^3}{31^2\cdot 37}.
\]
此外，定义短 Weierstrass 曲线
\[
E_{\mathrm{res}}:\quad Y^2=X^3-12096X+288144,
\]
则 \(E_{\mathrm{res}}\cong_{\QQ}E_{\mathscr R}\)；例如由变换
\[
X_{\mathrm{res}}=9X-111,\qquad Y_{\mathrm{res}}=27Y
\]
可将 \(E_{\mathscr R}:\ Y^2=X^3-37X^2+307X+361\) 送至短模型 \(E_{\mathrm{res}}:\ Y_{\mathrm{res}}^2=X_{\mathrm{res}}^3-12096X_{\mathrm{res}}+288144\)。
\end{conclusion}
\begin{proof}[证明要点]
将 \(z=1+\frac{y}{x}\) 代入 \(\mathscr R(z,y)=0\) 并清分母，出现因子 \(y\) 对应基点 \((0,1)\)；其余二次因子即给出所示关于 \(y\) 的二次方程。
令该二次方程的系数为 \(Ay^2+By+C=0\)，则恒等式 \((2Ay+B)^2=B^2-4AC\) 导出
\[
u^2=\frac{B^2-4AC}{x^2}=x\bigl(65x^3+16x^2-16x-4\bigr),
\]
从而得到 \(\mathcal Q\)。
二元四次不变量 \(I,J\) 的计算为经典公式的直接代入；由
\(
c_4=2^4I,\ c_6=2^5J
\)
可回收 \(j=c_4^3/\Delta\) 的闭式。
最后，将 \(X=\frac{X_{\mathrm{res}}+111}{9}\)、\(Y=\frac{Y_{\mathrm{res}}}{27}\) 代回 \(E_{\mathscr R}\) 即得短 Weierstrass 方程，并给出 \(E_{\mathrm{res}}\cong_{\QQ}E_{\mathscr R}\)。
可复现核验脚本：\texttt{scripts/exp\_fold\_zm\_resolvent\_cubic\_weierstrass\_audit.py}。证毕。
\end{proof}

\begin{proposition}[无复乘与几何端同态代数的极大性]\label{prop:xi-terminal-zm-resolvent-cubic-no-cm-max-end}
\leanverified{paper\_xi\_terminal\_zm\_resolvent\_cubic\_no\_cm\_max\_end}
令 \(E_{\mathscr R}\) 为结论 \ref{con:xi-terminal-zm-resolvent-cubic-weierstrass-birational} 的 Weierstrass 模型。
则 \(E_{\mathscr R}\) 不具有复乘；因此
\[
\mathrm{End}_{\overline{\QQ}}(E_{\mathscr R})=\ZZ,
\qquad
\mathrm{End}^0_{\overline{\QQ}}(E_{\mathscr R})=\QQ.
\]
从而
\[
\mathrm{End}^0_{\overline{\QQ}}(E_{\mathscr R}^{2})\ \cong\ M_2(\QQ).
\]
特别地，若 \(A\) 为任意满足 \(A\sim_{\QQ}E_{\mathscr R}^{2}\) 的二维阿贝尔簇（例如结论 \ref{con:xi-terminal-zm-s4-jacobian-a2-isog-r2} 的 \(A_2\)，或定理 \ref{thm:xi-terminal-zm-s3-closure-jacobian-splitting} 的标准表示分量），则
\(
\mathrm{End}^0_{\overline{\QQ}}(A)\cong M_2(\QQ)
\)。
\end{proposition}
\begin{proof}
若 \(E_{\mathscr R}\) 具有复乘，则其 \(j\)-不变量为代数整数；而结论 \ref{con:xi-terminal-zm-resolvent-cubic-weierstrass-birational} 给出
\[
j(E_{\mathscr R})=\frac{2^{18}\cdot 7^{3}}{31^{2}\cdot 37}\in\QQ\setminus\ZZ,
\]
其在 \(\QQ\) 中非整，因而不可能为代数整数，矛盾。
故 \(E_{\mathscr R}\) 无复乘，得到 \(\mathrm{End}^0_{\overline{\QQ}}(E_{\mathscr R})=\QQ\)。
最后由一般恒等式
\(
\mathrm{End}^0_{\overline{\QQ}}(E^{2})\cong M_2(\mathrm{End}^0_{\overline{\QQ}}(E))
\)
得到 \(\mathrm{End}^0_{\overline{\QQ}}(E_{\mathscr R}^{2})\cong M_2(\QQ)\)，并由同源不变性推广到任意 \(A\sim_{\QQ}E_{\mathscr R}^{2}\)。证毕。
\end{proof}

\begin{conclusion}[\texorpdfstring{$\pi^{(3)}_y$}{pi\_y(3)} 的零极点分解与三条有理纤维的分歧剖面]\label{con:xi-terminal-zm-resolvent-cubic-y-divisor-rational-fibers}
通过结论 \ref{con:xi-terminal-zm-resolvent-cubic-weierstrass-birational} 的同构，把投影
\(
\pi^{(3)}_y:\overline{\mathcal R}\to\mathbb P^1_y
\)
识别为 \(E_{\mathscr R}\) 上的有理函数
\[
y=-\frac{8X}{Y+19+14X-X^2}\in\QQ(E_{\mathscr R}).
\]
则其极点、零点除子在 \(E_{\mathscr R}(\overline{\QQ})\) 上满足精确分解
\[
(y)_\infty=2(15,-4)+(-1,-4),
\qquad
(y)_0=2O+(0,19),
\]
其中 \(O\) 为 \(E_{\mathscr R}\) 的无穷远点。
因此，\(\pi^{(3)}_y\) 的三条有理分歧纤维在 \(E_{\mathscr R}\) 上可具体写为
\[
(\pi^{(3)}_y)^{-1}(0)=2O+(0,19),
\qquad
(\pi^{(3)}_y)^{-1}(1)=2(23,4)+(-1,4),
\qquad
(\pi^{(3)}_y)^{-1}(\infty)=2(15,-4)+(-1,-4),
\]
从而在 \(y=0,1,\infty\) 三处均呈现“一点二次分歧、其余不分歧”的纤维型。
\end{conclusion}
\begin{proof}[证明要点]
由
\(
y=-8X/(Y+19+14X-X^2)
\)
可知零点来自 \(X=0\) 与无穷远点处的阶数比较，而极点来自分母与曲线的交点及其重数；
将 \(Y=X^2-14X-19\) 代入 \(E_{\mathscr R}\) 的方程并因式分解可得分母零点的 \(X\)-坐标分解，从而读出 \((y)_\infty\)。
对 \(y-1\) 同理得到 \(y=1\) 纤维的重数分解。
可复现核验脚本：\texttt{scripts/exp\_fold\_zm\_resolvent\_cubic\_weierstrass\_audit.py}。证毕。
\end{proof}

\begin{conclusion}[非有理分歧值的 \texorpdfstring{$X$}{X}-坐标消元证书与 \texorpdfstring{$P_{\mathrm{LY}}$}{P\_LY} 的再出现]\label{con:xi-terminal-zm-resolvent-cubic-nonrational-branch-elimination}
考虑 \(\pi^{(3)}_y\) 的非有理分歧点（即不在 \(y=0,1,\infty\) 之上者）。
则其在 \(E_{\mathscr R}\) 上的 \(X\)-坐标必满足不可约三次方程
\[
g_{\mathscr R}(X):=4X^3+3X^2+114X-269=0,
\]
且相应分歧值 \(y\) 可纯以 \(X\) 表示为
\[
y=\frac{16(X^2+19)}{(X-15)(2X^2+X+15)}.
\]
对 \(X\) 作结式消元可得显式证书
\[
\boxed{\;
\mathrm{Res}_X\!\Bigl(g_{\mathscr R}(X),\ (X-15)(2X^2+X+15)y-16(X^2+19)\Bigr)
=-68460544\,P_{\mathrm{LY}}(y)\; }.
\]
从而上述三条非有理二次分歧值的极小多项式恰为
\[
P_{\mathrm{LY}}(y)=256y^3+411y^2+165y+32.
\]
\end{conclusion}
\begin{proof}[证明要点]
对 \(y\) 的显式表达式计算 \(\dd y\) 并在 \(E_{\mathscr R}\) 上消去 \(Y\) 得到分歧点的 \(X\)-坐标多项式；
对 \(X\) 作结式消元即可回收 \(P_{\mathrm{LY}}(y)\)。
可复现核验脚本：\texttt{scripts/exp\_fold\_zm\_resolvent\_cubic\_weierstrass\_audit.py}。证毕。
\end{proof}

\begin{conclusion}[\texorpdfstring{$\QQ(\sqrt{-111})$}{Q(sqrt(-111))} 的几何实现：由 \texorpdfstring{$\mathrm{Disc}(g_{\mathscr R})$}{Disc(g\_R)} 直接读出]\label{con:xi-terminal-zm-resolvent-cubic-sqrt-111-geometric}
结论 \ref{con:xi-terminal-zm-resolvent-cubic-nonrational-branch-elimination} 的分歧控制多项式 \(g_{\mathscr R}(X)\) 的判别式满足
\[
\mathrm{Disc}(g_{\mathscr R})=-(2^3\cdot 3\cdot 31)^2\cdot 111,
\]
故其分裂域的二次分辨子域必为
\[
\QQ\!\left(\sqrt{\mathrm{Disc}(g_{\mathscr R})}\right)=\QQ(\sqrt{-111}).
\]
因此，二次域 \(\QQ(\sqrt{-111})\) 在此获得一条完全显式的几何实现：它即为 \(\pi^{(3)}_y\) 于 \(y\neq 0,1,\infty\) 的三条二次分歧值之定义域的公共二次分辨域（由 \(g_{\mathscr R}\) 的判别式直接给出），并与 \(P_{\mathrm{LY}}\) 的判别式平方自由部分精确一致。
\end{conclusion}
\begin{proof}[证明要点]
判别式分解可由对三次多项式的标准判别式公式计算并在 \(\ZZ\) 中分解得到。
由不可约三次分裂域的标准结构，二次分辨子域由 \(\sqrt{\mathrm{Disc}}\) 给出。
可复现核验脚本：\texttt{scripts/exp\_fold\_zm\_resolvent\_cubic\_weierstrass\_audit.py}。证毕。
\end{proof}

\begin{conclusion}[\texorpdfstring{$E_{\mathscr R}(\QQ)$}{E\_R(Q)} 的扭点平凡性与可计算的无穷阶点]\label{con:xi-terminal-zm-resolvent-cubic-mw-point}
\(E_{\mathscr R}(\QQ)\) 的扭点子群为平凡群
\[
E_{\mathscr R}(\QQ)_{\mathrm{tors}}=\{O\}.
\]
特别地，任一 \(\QQ\)-有理点只要不同于 \(O\) 即为无穷阶点；例如
\[
P:=(23,4)\in E_{\mathscr R}(\QQ)
\]
为无穷阶点。
并且在结论 \ref{con:xi-terminal-zm-resolvent-cubic-biquadratic-invariants-short-weierstrass} 的短 Weierstrass 模型
\[
E_{\mathrm{res}}:\quad Y^2=X^3-12096X+288144
\]
上，点
\[
P_{\mathrm{res}}:=(25,37)\in E_{\mathrm{res}}(\QQ)
\]
亦为无穷阶点，从而 \(\operatorname{rank}E_{\mathrm{res}}(\QQ)\ge 1\)（等价于 \(\operatorname{rank}E_{\mathscr R}(\QQ)\ge 1\)）。
在结论 \ref{con:xi-terminal-zm-resolvent-cubic-weierstrass-birational} 的平面模型坐标下，
\[
\bigl(y(P),z(P)\bigr)=(1,-3),
\]
从而沿 \((nP)\)（\(n\ge 1\)）取值可产生无穷多个有理参数
\[
y_n:=y(nP)\in\QQ,
\]
使得三次预解式 \(\mathscr R(z,y_n)=0\) 在 \(\QQ\) 上具有有理根 \(z_n:=z(nP)\in\QQ\)，并因此在每个 \(y_n\) 处发生线性因子化。
\end{conclusion}
\begin{proof}[证明要点]
扭点注入良约化的有限域点群，故扭点阶整除任意良约化素数 \(p\nmid \Delta\) 处的 \(\#E_{\mathscr R}(\FF_p)\)；
例如在 \(p=3,5\) 处良约化且 \(\#E_{\mathscr R}(\FF_3)=5\)、\(\#E_{\mathscr R}(\FF_5)=8\)，其最大公因子为 \(1\) 即推出扭点平凡性。
点 \(P_{\mathrm{res}}=(25,37)\in E_{\mathrm{res}}(\QQ)\) 的代入验证为直接计算；又由 \(E_{\mathrm{res}}\cong_{\QQ}E_{\mathscr R}\) 知其对应点在 \(E_{\mathscr R}(\QQ)\) 中非零，因扭点平凡而必为无穷阶点。
\((y(P),z(P))\) 的取值由结论 \ref{con:xi-terminal-zm-resolvent-cubic-weierstrass-birational} 的显式公式直接代入得到。
可复现核验脚本：\texttt{scripts/exp\_fold\_zm\_resolvent\_cubic\_weierstrass\_audit.py}。证毕。
\end{proof}

\input{con__xi-terminal-zm-resolvent-cubic-dihedral-d4-specialization}

\begin{conclusion}[\texorpdfstring{$31$}{31} 与 \texorpdfstring{$37$}{37} 处的约化型与分裂性判据]\label{con:xi-terminal-zm-resolvent-cubic-reduction-31-37}
对结论 \ref{con:xi-terminal-zm-resolvent-cubic-weierstrass-birational} 的 \(E_{\mathscr R}\)，由
\(
v_{31}(c_4)=v_{37}(c_4)=0
\)
且
\[
v_{31}(\Delta)=2,\qquad v_{37}(\Delta)=1,
\]
可知其在 \(p=31,37\) 处均为乘性约化。
进一步由 \(-c_6\) 的二次特征判定分裂性：
\[
\left(\frac{-c_6}{31}\right)=+1
\quad\Longrightarrow\quad
\text{\(31\) 处分裂乘性约化（Kodaira 型 \(I_2\)）},
\]
\[
\left(\frac{-c_6}{37}\right)=-1
\quad\Longrightarrow\quad
\text{\(37\) 处非分裂乘性约化（Kodaira 型 \(I_1\)）}.
\]
因此，素数 \(31\) 与 \(37\) 的局部约化行为在 \(E_{\mathscr R}\) 的 Weierstrass 模型上被完全显式地控制，并与三次预解式所诱导的分歧结构在判别式层面严格对齐。
\end{conclusion}
\begin{proof}[证明要点]
由 \(v_p(c_4)=0\) 且 \(v_p(\Delta)>0\) 的 Tate 判据知为乘性约化；
Kodaira 型由 \(v_p(\Delta)\) 决定，而分裂性由 \(\left(\frac{-c_6}{p}\right)\) 判别。
可复现核验脚本：\texttt{scripts/exp\_fold\_zm\_resolvent\_cubic\_weierstrass\_audit.py}。证毕。
\end{proof}

\begin{conclusion}[三次预解式在坏素数 \texorpdfstring{$37$}{37} 处的约化奇点：唯一非分裂结点与有理归一化]\label{con:xi-terminal-zm-resolvent-cubic-badprime-37-nonsplit-node}
令 \(\overline{\mathcal R}\subset \PP^{2}_{[Z:Y:W]}\) 为结论 \ref{con:xi-terminal-zm-resolvent-cubic-plane-cubic-geometry} 的平面三次曲线，
并记其在 \(\FF_{37}\) 上的约化为 \(\overline{\mathcal R}_{37}\)。
则 \(\overline{\mathcal R}_{37}\) 的射影闭包仅有一个奇点，且在仿射图 \(W=1\) 上可取为
\[
(z,y)\equiv (29,6)\pmod{37}.
\]
该奇点为非分裂结点；因此 \(\overline{\mathcal R}_{37}\) 的归一化为 \(\PP^{1}_{\FF_{37}}\)。
并且该退化由判别式核 \(P_{\mathrm{LY}}(y)\) 在 \(y\equiv 6\pmod{37}\) 的二重消失所触发，
从而与判别式二次脊曲线在该素处的结点机制保持一致（参见结论 \ref{con:xi-terminal-zm-elliptic-lattes-badprime-37-degree-drop}）。
\end{conclusion}
\begin{proof}[证明要点]
在 \(\FF_{37}\) 上联立
\(
\mathscr R(z,y)=\partial_{z}\mathscr R(z,y)=\partial_{y}\mathscr R(z,y)=0
\)
得到唯一解 \((z,y)=(29,6)\)，从而约化曲线仅有该一处奇点。
将坐标平移到该点附近，最低齐次项为二次型，其判别式为 \(2\in\FF_{37}^{\times}\) 的非平方元，
故切锥在 \(\FF_{37}\) 上不分裂而在二次扩张上分裂，给出非分裂结点。
平面三次曲线在出现结点时的归一化为 \(\PP^{1}\) 为标准结论。
可复现核验脚本：\texttt{scripts/exp\_fold\_zm\_resolvent\_cubic\_weierstrass\_audit.py}。证毕。
\end{proof}

\begin{conclusion}[Lee--Yang 三次因子的几何起源：重量映射 \texorpdfstring{$\pi_y$}{pi\_y} 的临界值消元与切触包络]\label{con:xi-terminal-zm-leyang-ramification-critical-values}
在结论 \ref{con:xi-terminal-zm-elliptic-normalization} 的椭圆模型
\[
E:\ Y^2=X^3-X+\frac14
\]
上，沿用结论 \ref{con:xi-terminal-zm-y-quadratic-extension-minpoly} 的物理主支
\[
y:=X^2+Y-\frac12\in\QQ(E),
\qquad
\omega:=\frac{\dd X}{2Y}.
\]
则 \(y\) 诱导次数 \(4\) 的有限态射 \(\pi_y:E\to\mathbb P^1_y\)，并满足 \((y)_\infty=4O\)（见结论 \ref{con:xi-terminal-zm-elliptic-y-hurwitz-passport}）。
则有严格微分恒等式
\[
\boxed{\;
\dd y=\bigl(4XY+3X^2-1\bigr)\,\omega
=\frac{4X\eta+3X^2+2X-1}{2\eta+1}\,\dd X\qquad(\eta:=Y-\tfrac12)\; }.
\]
并在仿射开集 \(Y\neq 0\) 上等价写为 \(\dd y=\frac{4XY+3X^2-1}{2Y}\,\dd X\)。
因此映射 \(\pi_y:E\to\mathbb P^1_y\)、\(P\mapsto y(P)\) 的有限分歧点由显式代数曲线
\[
4XY+3X^2-1=0
\]
在 \(E\) 上的交点给出；消去 \(Y\) 后得到其 \(X\)-坐标满足完全分解
\[
\boxed{\;
(X-1)(X+1)\bigl(16X^3-9X^2+1\bigr)=0\; }.
\]
等价地，其五次展开式为
\[
16X^{5}-9X^{4}-16X^{3}+10X^{2}-1=0.
\]
更精确地：
\begin{enumerate}
  \item 当 \(X=1\) 时，唯一有限分歧点为 \(\bigl(1,-\tfrac12\bigr)\)，其分歧值为 \(y=0\)；
  \item 当 \(X=-1\) 时，唯一有限分歧点为 \(\bigl(-1,\tfrac12\bigr)\)，其分歧值为 \(y=1\)；
  \item 若 \(\alpha\) 为不可约三次方程 \(16X^{3}-9X^{2}+1=0\) 的任一根，取
  \[
  Y=\frac{1-3\alpha^{2}}{4\alpha},
  \qquad
  P_{\alpha}:=(\alpha,Y)\in E(\overline{\QQ}),
  \]
  则 \(P_{\alpha}\) 为 \(\pi_y\) 的简单分歧点，且其分歧值为
  \[
  y_{\alpha}
  =\alpha^{2}+\frac{1-3\alpha^{2}}{4\alpha}-\frac12
  =\frac{(\alpha-1)\bigl(4\alpha^{2}+\alpha-1\bigr)}{4\alpha}.
  \]
\end{enumerate}
从而 \(\pi_y\) 的有限分歧值集合恰为 \(y=0\)、\(y=1\) 与 \(\{y_{\alpha}\}\)。
其中 Lee--Yang 三次多项式
\[
P_{\mathrm{LY}}(y):=256y^3+411y^2+165y+32
\]
可由纯消元证书刻画为严格结果式恒等式
\[
\boxed{\;
\mathrm{Res}_X\!\left(16X^3-9X^2+1,\ 4Xy-\bigl(4X^3-3X^2-2X+1\bigr)\right)
=-64\,P_{\mathrm{LY}}(y)\; }.
\]
等价地，交换消元顺序（此处两式关于 \(X\) 的次数皆为 \(3\)）得到
\[
\boxed{\;
P_{\mathrm{LY}}(y)=\frac1{64}\,
\mathrm{Res}_X\!\left(4Xy-\bigl(4X^3-3X^2-2X+1\bigr),\ 16X^3-9X^2+1\right)\; }.
\]
因此上述三点给出的三个分歧值 \(\{y_{\alpha}\}\) 恰为 \(P_{\mathrm{LY}}(y)=0\) 的三根。
因而 Lee--Yang 三次因子不仅是谱四次的判别式核（结论 \ref{con:xi-terminal-zm-discriminant-leyang}），亦等价给出 \(\pi_y\) 的三条非有理临界值在 \(\QQ[y]\) 中的极小多项式（结论 \ref{con:xi-terminal-zm-leyang-cubic-arithmetic}）；
在几何上，这等价于一族抛物线
\[
Y=y-X^2+\frac12
\]
与 \(E\) 的切触参数在消元意义下的像。
\end{conclusion}
\begin{proof}[证明要点]
由曲线方程微分关系 \(2Y\,\dd Y=(3X^2-1)\dd X\) 得 \(\dd Y=(3X^2-1)\omega\)，从而
\[
\dd y=2X\dd X+\dd Y=\bigl(4XY+3X^2-1\bigr)\omega.
\]
因此有限分歧点满足 \(4XY+3X^2-1=0\)。
由此解出
\(
Y=(1-3X^2)/(4X)
\)
并代入 \(E\) 的方程，化简即得
\(
(X-1)(X+1)(16X^3-9X^2+1)=0
\)。
当 \(X=1\) 与 \(X=-1\) 时分别得到 \(y=0\) 与 \(y=1\)。
对满足 \(16X^3-9X^2+1=0\) 的非平凡分歧点，将 \(y=X^2+Y-\frac12\) 与 \(Y=(1-3X^2)/(4X)\) 联立得到线性约束
\[
4Xy=4X^3-3X^2-2X+1.
\]
对变量 \(X\) 作消元（结果式）得到严格恒等式
\(
\mathrm{Res}_X(16X^3-9X^2+1,\ 4Xy-(4X^3-3X^2-2X+1))=-64\,P_{\mathrm{LY}}(y)
\)，
从而非平凡有限分歧值由 \(P_{\mathrm{LY}}(y)=0\) 精确刻画，并与结论 \ref{con:xi-terminal-zm-discriminant-leyang} 的判别式分解一致。
可复现核验脚本：\texttt{scripts/exp\_fold\_zm\_elliptic\_leyang\_cover\_geometry\_audit.py}。证毕。
\end{proof}

\begin{conclusion}[权平面四次模型的结式分裂：Lee--Yang 三次因子与共轭纤维三次因子的对偶消元]\label{con:xi-terminal-zm-leyang-conjugate-fiber-cubic}
沿用结论 \ref{con:xi-terminal-zm-leyang-ramification-critical-values} 的记号，并由
\(
Y=y-X^{2}+\tfrac12
\)
消去 \(Y\)。
则由下式所定义之首一四次多项式 \(F(X,y)\) 在 \(\QQ(E)\) 中满足 \(F(X,y)=0\)，并且由 \(Y=y-X^{2}+\tfrac12\) 可逆重建 \(Y\)，从而 \(\QQ(E)=\QQ(y)(X)\) 且 \([\QQ(E):\QQ(y)]=4\)；等价地，\(E\) 的函数域具有首一四次表示
\[
\QQ(E)\cong \QQ(X,y)/(F),
\qquad
F(X,y):=X^{4}-X^{3}-(2y+1)X^{2}+X+y(y+1).
\]
该表示与结论 \ref{con:xi-terminal-zm-elliptic-normalization} 的规范化同构一致，并使谱变量在规范化层面可识别为椭圆坐标 \(X\)。
在该平面四次模型下，\(\pi_y:E\to\mathbb P^{1}_{y}\) 的有限分歧点等价于方程组
\[
F(X,y)=0,\qquad \partial_{X}F(X,y)=0,
\]
并且对 \(y\) 消元得到严格结果式恒等式
\[
\boxed{\;
\mathrm{Res}_{y}\bigl(F(X,y),\partial_{X}F(X,y)\bigr)=-(X-1)(X+1)\bigl(16X^{3}-9X^{2}+1\bigr)\; }.
\]
\par
令
\(
g(X):=16X^{3}-9X^{2}+1
\)，并定义与 \(P_{\mathrm{LY}}\) 共轭的三次因子
\[
Q_{\mathrm{LY}}(y):=256y^{3}+195y^{2}+93y-48\in\ZZ[y].
\]
则对变量 \(X\) 消元给出另一条严格结式分裂
\[
\boxed{\;
\mathrm{Res}_{X}\bigl(F(X,y),g(X)\bigr)=P_{\mathrm{LY}}(y)\,Q_{\mathrm{LY}}(y)\; }.
\]
因此，对任意满足 \(g(X)=0\) 的根 \(\alpha\)：
在分歧约束 \(\partial_XF=0\) 下的重量值为
\[
y=\frac{4\alpha^{3}-3\alpha^{2}-2\alpha+1}{4\alpha},
\qquad
P_{\mathrm{LY}}(y)=0,
\]
而其在椭圆反演 \(\iota:(X,Y)\mapsto (X,-Y)\) 下的共轭点给出第二张（非分歧）纤维，
其重量值满足 \(Q_{\mathrm{LY}}(y)=0\)；
两组三次根由 \(\iota\) 交换。
从而 Lee--Yang 因子 \(P_{\mathrm{LY}}(y)\) 被识别为 \(\pi_y\) 的非平凡分歧值集合之不可约消元核，而 \(Q_{\mathrm{LY}}(y)\) 刻画同一 \(g(X)=0\) 之上的共轭非分歧纤维。
\end{conclusion}
\begin{proof}[证明要点]
由 \(Y=y-X^{2}+\tfrac12\) 代回 \(E\) 的方程并展开，即得 \(F(X,y)=0\) 的首一四次表示。
在平面模型中，有限分歧点满足 \(F=\partial_XF=0\)；
对 \(y\) 作结果式消元并在 \(\ZZ[X]\) 中因式分解，即得 \(\mathrm{Res}_{y}(F,\partial_XF)\) 的分裂恒等式。
\par
对 \(g(X)=0\) 的三根，方程 \(F(X,y)=0\) 给出对应的两张重量纤维（由 \(Y\mapsto -Y\) 交换）；
对 \(X\) 作结果式消元得到 \(y\) 的一元消元多项式，直接分解为两条三次因子。
其中分歧点的重量值由结论 \ref{con:xi-terminal-zm-leyang-ramification-critical-values} 的线性约束
\(
4Xy=4X^{3}-3X^{2}-2X+1
\)
确定，故其对应因子必为 \(P_{\mathrm{LY}}\)，另一因子即定义为 \(Q_{\mathrm{LY}}\) 并对应共轭非分歧纤维。
可复现核验脚本：\texttt{scripts/exp\_fold\_zm\_elliptic\_leyang\_cover\_geometry\_audit.py}。证毕。
\end{proof}

\input{con__xi-terminal-zm-elliptic-y-hurwitz-passport-refined}

\begin{conclusion}[分歧除子的主化与群律迹：Lee--Yang 三分歧点的加法迹坍缩为显式有理点]\label{con:xi-terminal-zm-leyang-ramification-divisor-principalization-trace}
沿用结论 \ref{con:xi-terminal-zm-elliptic-y-hurwitz-passport} 的记号，并令
\[
f:=4XY+3X^2-1\in\QQ(E)^\times.
\]
设 \(O\) 为无穷远点，记两条有理分歧点
\[
Q_0:=\Bigl(1,-\frac12\Bigr),\qquad Q_0':=\Bigl(-1,\frac12\Bigr),
\]
并令 \(Q_1,Q_2,Q_3\in E(\overline{\QQ})\) 为满足 \(16X^3-9X^2+1=0\) 的三条非有理分歧点（其 \(Y\)-坐标由 \(4XY+3X^2-1=0\) 唯一确定）。
则有严格主除子恒等式
\[
\boxed{\;
\mathrm{div}(f)=Q_0+Q_0'+Q_1+Q_2+Q_3-5O\; }.
\]
从而在椭圆曲线群律意义下成立
\[
Q_1+Q_2+Q_3=-(Q_0+Q_0')=\Bigl(\frac14,-\frac18\Bigr)\in E(\QQ),
\]
并且
\[
Q_0+Q_0'=\Bigl(\frac14,\frac18\Bigr).
\]
因此，位于 Lee--Yang 三个分歧值之上的三条分歧点（其定义域为 \(\QQ(\alpha)\)，\(\alpha\) 为 \(16X^3-9X^2+1=0\) 的根）在群律意义下满足一条全局加法约束：其加法迹严格落入固定的有理点 \((\tfrac14,-\tfrac18)\)。
结合结论 \ref{con:xi-terminal-zm-elliptic-mw-generator-alignment} 的自由生成元 \(R=(0,\tfrac12)\) 及 \(Q_0=-[2]R\)、\(Q_0'=-[3]R\)，可将上述迹恒等式改写为
\[
Q_1+Q_2+Q_3=5R.
\]
该恒等式亦可视为 Riemann--Hurwitz 所给出的线性等价关系在 \(\mathrm{Pic}^0(E)\cong E\) 下的具象化。
由于结论 \ref{con:xi-terminal-zm-elliptic-rational-points-denominator-law} 已给出 \(E(\QQ)_{\mathrm{tors}}=\{O\}\)，且 \((\tfrac14,-\tfrac18)\neq O\)，故该有理迹点为无限阶。
\end{conclusion}
\begin{proof}[证明要点]
由结论 \ref{con:xi-terminal-zm-leyang-ramification-critical-values} 的微分恒等式 \(\dd y=f\,\omega\) 可知：
\(f\) 的有限零点恰为 \(\pi_y\) 的有限分歧点，且由结论 \ref{con:xi-terminal-zm-elliptic-y-hurwitz-passport} 的 Hurwitz 护照可知其均为一阶零。
在无穷远点 \(O\) 处取均匀化参数 \(t\) 使 \(X\sim t^{-2}\)、\(Y\sim t^{-3}\)，则
\[
f=4XY+3X^2-1\sim 4t^{-5},
\]
故 \(f\) 在 \(O\) 处有五阶极点且无其他极点，从而得到所示主除子恒等式。
由于 \(\mathrm{div}(f)\) 在 \(\mathrm{Pic}^{0}(E)\cong E\) 中为零类，推出 \(Q_0+Q_0'+Q_1+Q_2+Q_3=O\)。
最后对有理点 \(Q_0,Q_0'\) 施用弦割加法公式：斜率
\[
m=\frac{\frac12-(-\frac12)}{-1-1}=-\frac12
\]
给出
\[
x(Q_0+Q_0')=m^2-1-(-1)=\frac14,\qquad
y(Q_0+Q_0')=m\Bigl(1-\frac14\Bigr)-\Bigl(-\frac12\Bigr)=\frac18,
\]
从而 \(Q_0+Q_0'=(\tfrac14,\tfrac18)\) 且 \(-(Q_0+Q_0')=(\tfrac14,-\tfrac18)\)，并得到结论。
可复现核验脚本：\texttt{scripts/exp\_fold\_zm\_elliptic\_leyang\_cover\_geometry\_audit.py}。证毕。
\end{proof}

\begin{conclusion}[Lee--Yang 判别式的范数恒等式：\texorpdfstring{$-y(y-1)P_{\mathrm{LY}}(y)$}{-y(y-1)P(y)} 为不同生成元的相对范数]\label{con:xi-terminal-zm-leyang-discriminant-norm-identity}
设
\[
k:=\QQ(y),\qquad K:=\QQ(E)=\QQ(X,Y),
\]
则 \(K/k\) 为次数 \(4\) 的函数域扩张。
该扩张可分，且其不同理想在 \(K\) 中为主理想。
令
\[
\delta:=4Xy-\bigl(4X^3-3X^2-2X+1\bigr)\in K.
\]
在 \(E\) 上利用 \(Y=y-X^2+\tfrac12\) 可化为
\[
\delta=4XY+3X^2-1,
\]
并且 \(\delta\) 在 \(\pi_y\) 的有限分歧点处恰为一阶零。
因此
\[
\mathfrak D_{K/k}=(\delta)=\bigl(1-3X^2-4XY\bigr).
\]
则其相对范数满足严格恒等式
\[
\boxed{\;
\operatorname{N}_{K/k}(\delta)=-y(y-1)\,P_{\mathrm{LY}}(y)\; }.
\]
\end{conclusion}
\begin{proof}[证明要点]
由 \(Y=y-X^2+\tfrac12\) 消去 \(Y\) 得 \(X\) 在 \(k\) 上满足的四次方程
\[
F(X,y)=X^4-X^3-(2y+1)X^2+X+y(y+1)=0,
\]
并且直接计算给出 \(\delta=-\partial_XF(X,y)\)。
由于 \(F\) 为关于 \(X\) 的首一可分多项式，标准理论给出不同理想
\(
\mathfrak D_{K/k}=\bigl(\partial_XF(X,y)\bigr)
\)；
从而 \(\mathfrak D_{K/k}=(\delta)\)。
再将 \(y=X^2+Y-\tfrac12\) 代入 \(\partial_XF(X,y)\)，即得
\(
\partial_XF(X,y)=1-3X^2-4XY
\)，故 \((\delta)=\bigl(1-3X^2-4XY\bigr)\)。
由于 \(F\) 关于 \(X\) 首一，范数 \(\operatorname{N}_{K/k}(\delta)\) 等于结果式 \(\mathrm{Res}_X(F,\delta)\)（等价于 \(\mathrm{Disc}_X(F)\) 的标准表示）。
而该结果式恰计算为 \(-y(y-1)P_{\mathrm{LY}}(y)\)，与结论 \ref{con:xi-terminal-zm-discriminant-leyang} 的判别式分解一致。
可复现核验脚本：\texttt{scripts/exp\_fold\_zm\_elliptic\_leyang\_cover\_geometry\_audit.py}。证毕。
\end{proof}

\begin{conclusion}[\texorpdfstring{$2$}{2}-division 多项式在 \texorpdfstring{$y$}{y}-投影下的范数平方化：隐藏三次控制多项式]\label{con:xi-terminal-zm-2division-norm-square}
沿用结论 \ref{con:xi-terminal-zm-leyang-discriminant-norm-identity} 的记号，仍令 \(k=\QQ(y)\)、\(K=\QQ(E)\)。
记 \(2\)-division 三次
\[
D(X):=4X^3-4X+1=4Y^2\in K.
\]
并记其在 \(y\)-坐标下的像控制多项式
\[
N_2(y):=16y^3-8y^2-4y+1\qquad(\text{亦记 }S(y)).
\]
则其在扩张 \(K/k\) 下的相对范数满足严格恒等式
\[
\boxed{\;
\operatorname{N}_{K/k}\bigl(D(X)\bigr)=N_2(y)^2\; }.
\]
等价地，由 \(2Y=2y-2X^2+1\)（见结论 \ref{con:xi-terminal-zm-y-quadratic-extension-minpoly}）得到
\[
\boxed{\;
\operatorname{N}_{K/k}(2Y)=-N_2(y)\; }.
\]
因此三次多项式 \(N_2(y)\) 精确刻画 \(E[2]\setminus\{O\}\) 在 \(\pi_y\) 下的像；并且由于 \(D(X)=(2Y)^2\) 在 \(K\) 中为平方，其相对范数在 \(k\) 中必为完全平方，上式给出该平方化的最小整系数证书。
\par
等价地，由 \(E[2]\setminus\{O\}\) 的特征方程 \(Y=0\) 与 \(D(X)=0\) 得 \(y=X^2-\tfrac12\)，从而存在严格消元证书
\[
\boxed{\;
\mathrm{Res}_X\!\left(4X^3-4X+1,\ y-\Bigl(X^2-\frac12\Bigr)\right)=N_2(y)\; }.
\]
\end{conclusion}
\begin{proof}[证明要点]
仍以四次方程 \(F(X,y)=0\)（见结论 \ref{con:xi-terminal-zm-leyang-discriminant-norm-identity}）作为 \(K/k\) 的首一表示，则
\(
\operatorname{N}_{K/k}(D(X))=\mathrm{Res}_X(F,D)
\)
与
\(
\operatorname{N}_{K/k}(2Y)=\mathrm{Res}_X(F,2y-2X^2+1)
\)
为标准结果式计算；
直接化简即得到所示恒等式。
可复现核验脚本：\texttt{scripts/exp\_fold\_zm\_elliptic\_leyang\_cover\_geometry\_audit.py}。证毕。
\end{proof}

\begin{conclusion}[Lee--Yang 三次域与临界三次域同构：显式生成元对应与判别式的 \texorpdfstring{$37$}{37} 不可消性]\label{con:xi-terminal-zm-leyang-cubic-field-isomorphism}
令 \(\alpha\) 为三次方程
\[
16X^3-9X^2+1=0
\]
的任一根，并定义
\[
\beta:=\frac{4\alpha^3-3\alpha^2-2\alpha+1}{4\alpha}
=\alpha^{2}-\frac34\alpha-\frac12+\frac{1}{4\alpha}
=-3\alpha^2+\frac32\alpha-\frac12.
\]
则：
\begin{enumerate}
  \item \(\beta\) 满足 Lee--Yang 三次方程 \(P_{\mathrm{LY}}(\beta)=0\)。
  \item 因此诱导出立方域同一性 \(\QQ(\alpha)=\QQ(\beta)\)。
  反向生成元可取为如下两支之一（对应该三次域的两种 \(\QQ\)-嵌入选择）：
  \[
  \alpha=\frac{-512\beta^2-518\beta-5}{279}\quad\text{或}\quad \alpha=\frac{512\beta^2-970\beta-739}{279}.
  \]
  \item 两个表示的多项式判别式满足
  \[
  \boxed{\ \mathrm{Disc}(16X^3-9X^2+1)=-2^2\cdot 3^3\cdot 37\ },
  \qquad
  \boxed{\ \mathrm{Disc}(P_{\mathrm{LY}})=-3^9\cdot 31^2\cdot 37\ }.
  \]
  并且二者严格相差一个平方因子：
  \[
  \boxed{\ \mathrm{Disc}(P_{\mathrm{LY}})=\mathrm{Disc}(16X^3-9X^2+1)\cdot\Bigl(\frac{837}{2}\Bigr)^{2}\ }.
  \]
  因而两者判别式的平方自由部分同为 \(-111=-3\cdot 37\)，对应的唯一二次分辨子域皆为 \(\QQ(\sqrt{-111})\)；
  从而两条 \(S_3\)-扩张的分裂域共享该虚二次子域；并且三次数域 \(K=\QQ(\alpha)=\QQ(\beta)\) 的域判别式为 \(\mathrm{Disc}(K)=-999=-3^{3}\cdot 37\)（见结论 \ref{con:xi-terminal-zm-leyang-cubic-field-minimal-discriminant}），其中 \(-111\) 为其二次分辨子域 \(\QQ(\sqrt{-111})\) 的基本判别式。
  同时在该三次域中，素数 \(37\) 必然分歧（其在上述判别式中以奇数次幂出现，无法被整数指标平方消去），从而 Lee--Yang 立方域在算术上被椭圆判别式素数 \(37\) 牢固锚定。
  \item 由于 \(16X^3-9X^2+1\) 与 \(P_{\mathrm{LY}}(y)\) 皆在 \(\QQ\) 上不可约且判别式非平方，故其分裂域伽罗瓦群皆同构于 \(S_3\)；
  并由 \(\QQ(\alpha)=\QQ(\beta)\) 得两者的最小分裂域一致。
\end{enumerate}
\end{conclusion}
\begin{proof}[证明要点]
由结论 \ref{con:xi-terminal-zm-leyang-ramification-critical-values} 的临界条件得到
\(
4\alpha\beta=4\alpha^3-3\alpha^2-2\alpha+1
\)，故 \(\beta\in\QQ(\alpha)\)。
将该线性关系与 \(\alpha\) 的三次方程作消元可得 \(P_{\mathrm{LY}}(\beta)=0\)，从而 \(\QQ(\beta)\subseteq\QQ(\alpha)\)。
两者均为三次扩张，故必相等。
判别式分解为直接代数计算。
可复现核验脚本：\texttt{scripts/exp\_fold\_zm\_elliptic\_leyang\_cover\_geometry\_audit.py}。证毕。
\end{proof}

\input{con__xi-terminal-zm-elliptic-3torsion-y-image-octic}

\input{con__xi-terminal-zm-leyang-cubic-field-minimal-discriminant}

\input{con__xi-terminal-zm-leyang-branchpolynomial-discriminant-modp-collisions}

\begin{conclusion}[坏素数 \texorpdfstring{$37$}{37} 处的同根塌缩：三条三次控制多项式的判别式与同余分解]\label{con:xi-terminal-zm-leyang-badprime-37-common-double-root}
沿用结论 \ref{con:xi-terminal-zm-leyang-conjugate-fiber-cubic} 的共轭三次
\(
Q_{\mathrm{LY}}(y)=256y^{3}+195y^{2}+93y-48
\)
与 Lee--Yang 三次
\(
P_{\mathrm{LY}}(y)=256y^{3}+411y^{2}+165y+32
\)，
并记二分点像控制三次（结论 \ref{con:xi-terminal-zm-2division-norm-square}）
\[
C(y):=16y^{3}-8y^{2}-4y+1
\qquad(\text{亦即结论 \ref{con:xi-terminal-zm-2division-norm-square} 中的 }S(y)).
\]
则三者的判别式均携带素因子 \(37\)，且满足严格分解：
\[
\mathrm{Disc}\bigl(C(y)\bigr)=2^{8}\cdot 37,\qquad
\mathrm{Disc}\bigl(P_{\mathrm{LY}}(y)\bigr)=-3^{9}\cdot 31^{2}\cdot 37,\qquad
\mathrm{Disc}\bigl(Q_{\mathrm{LY}}(y)\bigr)=-3^{3}\cdot 37\cdot 2677^{2}.
\]
更进一步地，在 \(\FF_{37}[y]\) 中三者发生同一参数值的双根塌缩：
\[
\begin{aligned}
C(y)&\equiv 16\,(y-7)(y-6)^{2}\ (\mathrm{mod}\ 37),\\
P_{\mathrm{LY}}(y)&\equiv -3\,(y-14)(y-6)^{2}\ (\mathrm{mod}\ 37),\\
Q_{\mathrm{LY}}(y)&\equiv -3\,(y-16)(y-6)^{2}\ (\mathrm{mod}\ 37).
\end{aligned}
\]
因此 \(y\equiv 6\ (\mathrm{mod}\ 37)\) 同时为三条三次控制多项式的共同重根。
与此同时，对物理基点 \(P=(2,-\tfrac52)\in E(\QQ)\)（结论 \ref{con:xi-terminal-zm-elliptic-rational-points-denominator-law}），其逆元满足 \(y(-P)=6\)，
从而该塌缩点在整数层面与一条显式有理点标记相一致。
\end{conclusion}
\begin{proof}[证明要点]
判别式与模 \(37\) 因式分解为 \(\ZZ[y]\) 与 \((\ZZ/37\ZZ)[y]\) 内的有限代数计算。
\(y(-P)=6\) 由 \(y=X^{2}+Y-\tfrac12\) 在点 \(-P=(2,\tfrac52)\) 的直接代入得到。
可复现核验脚本：\texttt{scripts/exp\_fold\_zm\_elliptic\_leyang\_cover\_geometry\_audit.py}。证毕。
\end{proof}

\input{con__xi-terminal-zm-elliptic-lattes-badprime-37-degree-drop}

\input{con__xi-terminal-zm-elliptic-weight-doubling-modular-equation}
\input{subsubsec__xi-time-protocol-conclusions-part15b-elliptic-audit-weight-tripling}
\input{con__xi-terminal-zm-4division-weight-image}
\input{subsubsec__xi-time-protocol-conclusions-part15b-leyang-4division-spine-branch-orbit}

\begin{conclusion}[所有分歧点处的 Puiseux 统一化：平方根临界指数与首项系数的显式公式]\label{con:xi-terminal-zm-branch-puiseux-uniformization}
将 \(F(X,y)=0\) 视为定义代数函数 \(X=X(y)\) 的四次方程（见结论 \ref{con:xi-terminal-zm-leyang-discriminant-norm-identity}）。
则对每一有限分歧值
\(
y_0\in\{0,1\}\cup\{y_1,y_2,y_3\}
\)
（其中 \(P_{\mathrm{LY}}(y_i)=0\)），存在局部参数 \(t\) 使 \(y-y_0=t^2\)，并且恰有两条分支在该点发生二重并合，其 Puiseux 展开具有普适平方根型
\[
X=\alpha_0\pm c_0\,t+O(t^2).
\]
其中：
\begin{enumerate}
  \item 在 \(y_0=0\)（并合点 \(X=1\)）处，可取 \(y=t^2\)，则两条并合分支满足
  \[
  X=1\pm \frac{1}{\sqrt2}\,t+\frac58\,t^2\mp \frac{43}{64\sqrt2}\,t^3-\frac{1}{16}t^4+O(t^5),
  \]
  即
  \[
  X=1\pm \frac{1}{\sqrt2}\,y^{1/2}+\frac58\,y\mp \frac{43}{64\sqrt2}\,y^{3/2}-\frac{1}{16}y^{2}+O(y^{5/2}).
  \]
  \item 在 \(y_0=1\)（并合点 \(X=-1\)）处，可取 \(y-1=t^2\)，则两条并合分支满足
  \[
  X=-1\pm \frac{i}{\sqrt6}\,t+O(t^2),
  \]
  即
  \[
  X=-1\pm \frac{i}{\sqrt6}\,(y-1)^{1/2}+O(y-1).
  \]
  \item 在任一 Lee--Yang 分歧值 \(y_0=y_i\)（并合点 \(X=\alpha_i\)，其中 \(\alpha_i\) 为 \(16X^3-9X^2+1=0\) 的根）处，首项系数满足
  \[
  \boxed{\;
  c_0^2=\frac{2}{3}\cdot\frac{3\alpha_i^2-1}{\alpha_i^2-1}\ \in\ \QQ(\alpha_i)=\QQ(y_i)\; }.
  \]
\end{enumerate}
\par
取局部参数 \(t=\sqrt{y-y_0}\) 的任一分支，则两条并合谱支可写为
\[
\lambda_{\pm}(y)=\alpha_0\pm c_0\sqrt{y-y_0}+O(y-y_0).
\]
因此对任一对数分支定义边界自由能 \(f_{\mathrm{edge}}(y):=\log\lambda_{\pm}(y)\)，则在 \(y\to y_0\) 时具有标准 \(1/2\) 阶分歧
\[
f_{\mathrm{edge}}(y)=\log\alpha_0\pm \frac{c_0}{\alpha_0}\sqrt{y-y_0}+O(y-y_0),
\qquad
\frac{\dd}{\dd y}f_{\mathrm{edge}}(y)=\pm \frac{c_0}{2\alpha_0}(y-y_0)^{-1/2}+O(1).
\]
在 \(y_0=y_{\mathrm{LY}}\) 的情形下，\(\alpha_0=\lambda_{\mathrm{LY}}\) 为结论 \ref{con:xi-terminal-zm-leyang-lambda-cubic} 的唯一实根，故
\[
c_0^2=\frac{2(3\lambda_{\mathrm{LY}}^{2}-1)}{3(\lambda_{\mathrm{LY}}^{2}-1)},
\qquad
c_0\approx 0.7476109845904598\ldots.
\]
\end{conclusion}
\begin{proof}[证明要点]
在分歧点处有 \(F(\alpha_0,y_0)=0\)、\(\partial_XF(\alpha_0,y_0)=0\)，并且 \(\partial_{XX}F(\alpha_0,y_0)\neq 0\)。
对 \(F\) 作二元 Taylor 展开并令 \(y-y_0=t^2\) 即得平方根型 Puiseux。
一般地，首项系数由二阶展开给出 \(c_0^2=-2(\partial_yF)/(\partial_{XX}F)\) 在 \((\alpha_0,y_0)\) 处的取值；
在原谱坐标 \((\lambda,y)\) 下（其中 \(F(X,y)=\Pi(X,y)\)），同一公式亦可写为
\(
c_0^2=-2\,\partial_y\Pi(\alpha_0,y_0)/\partial_{\lambda}^{2}\Pi(\alpha_0,y_0)
\)。
对 \(y_0=y_i\) 的情形再利用 \(\alpha_i\) 的三次证书化简，即得到所示闭式。
对 \(y_0=0,1\) 可进一步递推系数得到所列展开。
可复现核验脚本：\texttt{scripts/exp\_fold\_zm\_elliptic\_leyang\_cover\_geometry\_audit.py}。证毕。
\end{proof}

\begin{conclusion}[分歧值处并合谱根的线性闭式：一次消元子给出的全局刚性]\label{con:xi-terminal-zm-branch-double-root-linear-closure}
对任意分歧值
\(
y_0\in\{0,1\}\cup\{P_{\mathrm{LY}}(y)=0\}
\)，
令 \(\lambda_*=\lambda_*(y_0)\) 表示谱四次 \(\Pi(\lambda,y_0)=0\) 的唯一二重根（等价地，\(\Pi(\lambda_*,y_0)=\partial_{\lambda}\Pi(\lambda_*,y_0)=0\) 的唯一解）。
则并合谱根满足如下线性闭式
\[
\boxed{\;
2232\,\lambda_*+16384y_0^{4}+8128y_0^{3}-14525y_0^{2}-5523y_0-2232=0\; }.
\]
并且在 Lee--Yang 三次分歧支 \(P_{\mathrm{LY}}(y_0)=0\) 上，上式等价化简为
\[
\boxed{\;
279\,\lambda_*+512y_0^{2}+518y_0+5=0\; }.
\]
因此在全部分歧集合上，并合谱值 \(\lambda_*(y_0)\) 可由线性闭式直接读出，无需在每个 \(y_0\) 处再行求解四次方程。
\end{conclusion}
\begin{proof}[证明要点]
对理想 \((\Pi,\partial_{\lambda}\Pi)\subset\QQ[\lambda,y]\) 作 Gröbner 消元可得到一条关于 \(\lambda\) 的一次消元子，因而在分歧集合上以线性形式唯一刻画共同根（二重根）；
其 \(y\)-消元子即为 \(\mathrm{Disc}_{\lambda}(\Pi)=-y(y-1)P_{\mathrm{LY}}(y)\)。
在模 \(P_{\mathrm{LY}}(y)\) 的三次商环中将该一次消元子约化即可得到化简式。
可复现核验脚本：\texttt{scripts/exp\_fold\_zm\_elliptic\_leyang\_resy\_spectral\_decomposition\_audit.py}。证毕。
\end{proof}

\begin{conclusion}[\texorpdfstring{$S_4$}{S4} 分裂域的几何化：\(24\)-重 \texorpdfstring{$S_4$}{S4}-伽罗瓦闭包曲线属 \(16\)，并伴随属 \(2\) 的二次分辨曲线与属 \(1\) 的三次预解式曲线]\label{con:xi-terminal-zm-s4-splitting-curve-genus16}
令 \(\widetilde{K}/\QQ(y)\) 为四次方程 \(F(X,y)\) 的分裂域，并令 \(\widetilde{C}\) 为其函数域对应的光滑射影曲线。
在结论 \ref{con:xi-terminal-zm-generic-galois-s4} 的泛型 \(S_4\) 性下，\(\widetilde{C}\to\mathbb P^1_y\) 为 \(S_4\)-伽罗瓦覆叠，且在六个分歧值
\[
\{\infty,0,1,y_1,y_2,y_3\}
\]
处的惰性群阶分别为
\[
e_\infty=4,\qquad e_0=e_1=e_{y_1}=e_{y_2}=e_{y_3}=2.
\]
其中 \(e_\infty=4\) 对应于 \(y\) 在无穷远点 \(O\) 处的全分歧极点 \((y)_\infty=4O\)（结论 \ref{con:xi-terminal-zm-elliptic-y-hurwitz-passport}）。
由 Riemann--Hurwitz 公式得到
\[
\boxed{\ g(\widetilde{C})=16\ }.
\]

此外，二次分辨子域与三次预解式分别给出两条自然的中间曲线：
\begin{enumerate}
  \item \textbf{符号二次商（固定子群 \(A_4\)）。}
  对应曲线为双覆叠
  \[
  C_2:\ w^2=-y(y-1)P_{\mathrm{LY}}(y),
  \]
  其分歧点为上述六点，故 \(g(C_2)=2\)。
  \item \textbf{三次预解式曲线（cubic resolvent）。}
  令 \(\overline{\mathcal R}\subset\PP^{2}\) 为三次预解式 \(\mathscr R(z,y)=0\) 的射影闭包（结论 \ref{con:xi-terminal-zm-resolvent-cubic-plane-cubic-geometry}）。
  则 \(\overline{\mathcal R}\) 为光滑射影平面三次曲线，且含有有理点 \((y,z)=(0,1)\)，从而 \(g(\overline{\mathcal R})=1\) 并可视为定义在 \(\QQ\) 上的椭圆曲线。
\end{enumerate}
因此，尽管根域归一化为属 \(1\) 的椭圆曲线 \(E\)，其 \(S_4\)-分裂域的伽罗瓦闭包在几何上提升为属 \(16\) 的高属曲线；
与此同时，二次分辨曲线与三次预解式曲线分别落在属 \(2\) 与属 \(1\)，为该 \(S_4\) 结构提供两条可见的低属投影。
\end{conclusion}
\begin{proof}[证明要点]
对 \(S_4\)-伽罗瓦覆叠使用 Riemann--Hurwitz 公式：
若在分歧值 \(b\) 处惰性群阶为 \(e_b\)，则总分歧贡献为 \(|S_4|\,(1-1/e_b)\)。
代入 \(|S_4|=24\) 与 \((e_\infty,e_0,e_1,e_{y_1},e_{y_2},e_{y_3})=(4,2,2,2,2,2)\) 得 \(g(\widetilde{C})=16\)。
二次商 \(C_2\) 的属由同一公式在次数 \(2\) 上计算得到 \(g(C_2)=2\)。
三次预解式曲线 \(\overline{\mathcal R}\) 的光滑性与属计算见结论 \ref{con:xi-terminal-zm-resolvent-cubic-plane-cubic-geometry}。
可复现核验脚本：\texttt{scripts/exp\_fold\_zm\_elliptic\_leyang\_cover\_geometry\_audit.py}。证毕。
\end{proof}

\begin{remark}[进一步结构化方向：\texorpdfstring{$S_4$}{S4} 伽罗瓦闭包上的表示与周期数据]\label{rem:xi-terminal-zm-s4-closure-representation-periods}
由结论 \ref{con:xi-terminal-zm-s4-splitting-curve-genus16} 的 \(S_4\)-伽罗瓦覆叠 \(\widetilde{C}\to\mathbb P^1_y\)（属 \(16\)），
自然得到群 \(S_4\) 在 \(H^{0}(\widetilde{C},\Omega^{1})\) 上的表示。
对该表示作不可约分解并与椭圆曲线 \(E\) 的模性锚点（判别式 \(37\)，见注记 \ref{rem:xi-terminal-elliptic-37a1-modular-anchor}）对接，
可在同一框架下统一刻画分支单值化、谱支置换以及相关的周期积分与椭圆对数数据。
\end{remark}


\begin{conclusion}[纤维主除子诱导的群律迹：固定重量纤维的四点零和]\label{con:xi-terminal-zm-elliptic-weight-fiber-sum-zero}
沿用结论 \ref{con:xi-terminal-zm-elliptic-y-hurwitz-passport} 的次数 \(4\) 重量态射
\(
\pi_y:E\to\PP^1_y
\)
与极点结构 \((y)_\infty=4O\)。
对任意非分歧值 \(y_0\in\overline{\QQ}\setminus\bigl(\{\infty,0,1\}\cup\{P_{\mathrm{LY}}(y)=0\}\bigr)\)，
记纤维为
\(
\pi_y^{-1}(y_0)=\{P_1,P_2,P_3,P_4\}
\)
（四点互异）。
则在椭圆曲线群律下成立严格恒等式
\[
\boxed{\;P_1+P_2+P_3+P_4=O\;}
\]
并且更一般地，对任意有限值 \(y_0\in\overline{\QQ}\)（允许分歧值）有主除子表达
\[
\mathrm{div}(y-y_0)=\sum_{P\in \pi_y^{-1}(y_0)}\operatorname{ord}_{P}(y-y_0)\,P-4O,
\]
从而在 \(\mathrm{Pic}^0(E)\cong E\) 的识别下给出加权群律迹
\(
\sum_{P}\operatorname{ord}_{P}(y-y_0)\,P=O
\)。
无穷远纤维由 \((y)_\infty=4O\) 给出。
若 \(Q\in E(\QQ)\) 且 \(y(Q)=y_0\) 为非分歧值，则其余三点（一般位于三次扩张中）之和严格等于 \(-Q\)。
\end{conclusion}
\begin{proof}[证明要点]
对非分歧值 \(y_0\)，函数 \(y-y_0\in\QQ(E)^\times\) 的零点恰为 \(P_1,\dots,P_4\) 且均为单零；
又由 \((y)_\infty=4O\) 得 \((y-y_0)_\infty=4O\)，从而
\(
\mathrm{div}(y-y_0)=P_1+P_2+P_3+P_4-4O
\)。
由于主除子在 \(\mathrm{Pic}^0(E)\cong E\) 中对应单位元，遂得 \(P_1+\cdots+P_4=O\)。
分歧值情形以 \(\operatorname{ord}_{P}(y-y_0)\) 计入重数后同理。证毕。
\end{proof}

\begin{conclusion}[有理点特化的整化分解：线性因子与三次余因子的整系数闭式]\label{con:xi-terminal-zm-elliptic-rational-specialization-integral-factorization}
沿用结论 \ref{con:xi-terminal-zm-order4-rational} 的四次谱方程 \(\Pi(\lambda,y)\) 及椭圆化识别。
令 \(Q\in E(\QQ)\) 为有理点，并写其谱坐标为既约形式
\[
X(Q)=\frac{u}{v^{2}},
\qquad u\in\ZZ,\ v\in\ZZ_{>0},\ \gcd(u,v)=1.
\]
再将重量值写为既约形式
\[
y(Q)=\frac{a}{v^{4}},
\qquad a\in\ZZ,\ \gcd(a,v)=1
\]
（分母精确为 \(v^{4}\) 见结论 \ref{con:xi-terminal-zm-pi-rational-root-iff-y-image}）。
则在 \(\ZZ[\lambda]\) 中成立完全显式分解恒等式
\[
\boxed{\;
v^{8}\Pi\Bigl(\lambda,\frac{a}{v^{4}}\Bigr)=\bigl(v^{2}\lambda-u\bigr)\,G_{Q}(\lambda)\;}
\]
其中三次余因子 \(G_Q(\lambda)\in\ZZ[\lambda]\) 可取为
\[
\boxed{\;
\begin{aligned}
G_Q(\lambda)=\;&v^{6}\lambda^{3}
+v^{4}(u-v^{2})\lambda^{2}
+v^{2}\bigl(u^{2}-uv^{2}-2a-v^{4}\bigr)\lambda\\
&+\Bigl(v^{6}+u\bigl(u^{2}-uv^{2}-2a-v^{4}\bigr)\Bigr).
\end{aligned}}
\]
因此，特化多项式 \(\Pi(\lambda,y(Q))\) 必含有理线性因子 \(\lambda-X(Q)\)，而三次因子 \(G_Q\) 的三根精确对应纤维 \(\pi_y^{-1}(y(Q))\) 中除 \(Q\) 外三点之 \(X\)-坐标。
\end{conclusion}
\begin{proof}[证明要点]
由结论 \ref{con:xi-terminal-zm-pi-rational-root-iff-y-image}，\(\lambda=X(Q)=u/v^{2}\) 为 \(\Pi(\lambda,a/v^{4})\) 的有理根；
两边乘以 \(v^{8}\) 得 \(v^{8}\Pi(\lambda,a/v^{4})\in\ZZ[\lambda]\) 并含因子 \((v^{2}\lambda-u)\)。
在 \(\QQ[\lambda]\) 中以 \((v^{2}\lambda-u)\) 作多项式除法，并逐项比较系数即可回收所示三次商 \(G_Q\) 的闭式；
该闭式的整系数性由 \(u,v,a\in\ZZ\) 直接推出。
根的几何解释由 \(\QQ(E)=\QQ(y)(X)\) 及 \(\Pi(X,y)=0\) 的归一化识别给出：对固定 \(y_0\)，方程 \(\Pi(\lambda,y_0)=0\) 的根集即为纤维点的 \(X\)-坐标多重集。
可复现核验脚本：\texttt{scripts/exp\_fold\_zm\_elliptic\_rational\_fiber\_cubic\_factor\_audit.py}。证毕。
\end{proof}

\begin{conjecture}[倍乘轨道诱导的三次余因子之满对称伽罗瓦性]\label{conj:xi-terminal-zm-elliptic-rational-fiber-cubic-s3-galois}
取物理基点 \(P=(2,-\tfrac52)\in E(\QQ)\)（结论 \ref{con:xi-terminal-zm-elliptic-rational-points-denominator-law}），并令 \(Q_n:=[n]P\)。
对每个 \(n\ge 1\) 将
\[
X(Q_n)=\frac{u_n}{v_n^2},
\qquad
y(Q_n)=\frac{a_n}{v_n^4}
\]
写为既约形式（\(u_n,a_n\in\ZZ\)、\(v_n\in\ZZ_{>0}\)）。
令 \(G_{Q_n}(\lambda)\in\ZZ[\lambda]\) 表示结论 \ref{con:xi-terminal-zm-elliptic-rational-specialization-integral-factorization} 的三次余因子。
猜想对所有 \(n\ge 2\)，多项式 \(G_{Q_n}(\lambda)\) 在 \(\QQ\) 上不可约且其判别式不为 \(\QQ^{\times 2}\) 中的平方，
从而
\[
\mathrm{Gal}\bigl(G_{Q_n}/\QQ\bigr)\cong S_3.
\]
此外，预期其三次分裂域除有限例外外两两不相同，并在合成意义下呈现近乎线性无关的增长行为（可通过判别式的平方自由核与分歧素数的稀疏互素性形式化）。
\end{conjecture}


\subsubsection{\texorpdfstring{$S_4$}{S4} 伽罗瓦闭包 Jacobian 的可见椭圆--Prym 因子分解：由子群格到等型因子之几何识别}\label{subsubsec:xi-s4-visible-jacobian-prym-factorization}
以下内容在结论 \ref{con:xi-terminal-zm-s4-galois-closure-genus-tower-holomorphic-differentials} 的表示分解与同源分解框架上，
把其中的 \(V_2,V_3,V_3'\) 等型因子分别几何化为三条具体的中间椭圆曲线与一条 Prym 三维体，从而得到 \(J(X_e)\) 的四因子可见分解。

\begin{conclusion}[\texorpdfstring{$D_4$}{D4}-固定子域的几何化：三次预解式椭圆曲线即为 \texorpdfstring{$X_e/D_4$}{Xe/D4}]\label{con:xi-terminal-zm-s4-d4fixedfield-resolvent-elliptic}
设 \(X_e\to\mathbb P^1_y\) 为谱四次覆盖 \(\Pi(\lambda,y)=0\) 的 \(S_4\) 伽罗瓦闭包曲线（结论 \ref{con:xi-terminal-zm-s4-galois-closure-genus-tower-holomorphic-differentials}；亦即结论 \ref{con:xi-terminal-zm-s4-splitting-curve-genus16} 的几何闭包曲线）。
取任意 Sylow \(2\)-子群 \(D_{4}\le S_{4}\)（阶 \(8\) 的二面体子群，亦常记为 \(D_{8}\)），并记
\[
R:=X_e/D_{4}.
\]
则 \(R\) 与结论 \ref{con:xi-terminal-zm-resolvent-cubic-plane-cubic-geometry} 的三次预解式平面三次曲线 \(\overline{\mathcal R}\) 双有理同一；因而 \(R\) 本身为椭圆曲线，并且（在 \(\QQ\) 上）满足
\[
R\ \cong\ \overline{\mathcal R}\ \cong\ E_{\mathscr R},
\]
其中 \(E_{\mathscr R}\) 为结论 \ref{con:xi-terminal-zm-resolvent-cubic-weierstrass-birational} 的 Weierstrass 模型。
\end{conclusion}
\begin{proof}[证明要点]
在代数闭包中记谱四次的四根为 \(\lambda_1,\dots,\lambda_4\)。
由结论 \ref{con:xi-terminal-zm-resolvent-cubic-plane-cubic-geometry}，配对不变量
\(
z_{12|34}:=\lambda_1\lambda_2+\lambda_3\lambda_4
\)
等满足三次预解式 \(\mathscr R(z,y)=0\)，并且 \(S_4\) 在三组配对 \(\{12|34,13|24,14|23\}\) 上诱导出次数 \(3\) 的自然置换作用。
其中稳定某一配对（例如 \(12|34\)）的稳定子恰为阶 \(8\) 的二面体子群 \(D_4\)。
因此固定子域 \(\QQ(X_e)^{D_4}\) 由 \((y,z_{12|34})\) 生成，并与 \(\mathscr R(z,y)=0\) 的函数域同一，从而得到 \(X_e/D_4\) 与 \(\overline{\mathcal R}\) 的双有理同一。
\end{proof}

\begin{conclusion}[\texorpdfstring{$V_2$}{V2}-等型因子的椭圆平方化：\texorpdfstring{$A_2\sim R^2$}{A2~R^2}]\label{con:xi-terminal-zm-s4-jacobian-a2-isog-r2}
沿用结论 \ref{con:xi-terminal-zm-s4-galois-closure-genus-tower-holomorphic-differentials} 的记号：
\[
J(X_e)\ \sim\ A_{\mathrm{sgn}}\times A_{2}\times A_{3}\times A_{3'}.
\]
令 \(R:=X_e/D_4\) 为结论 \ref{con:xi-terminal-zm-s4-d4fixedfield-resolvent-elliptic} 的椭圆商。
则 \(V_2\)-等型因子满足
\[
\boxed{\ A_{2}\ \sim\ R^{2}\ }.
\]
\end{conclusion}
\begin{proof}[证明要点]
由 \(S_4\) 在三组配对上的次数 \(3\) 置换作用可知：指数为 \(3\) 的共轭子群族 \(\{gD_4g^{-1}\}\) 与三条中间商曲线 \(\{X_e/(gD_4g^{-1})\}\) 一一对应。
其置换表示同构于 \(\mathrm{Ind}_{D_4}^{S_4}\mathbf{1}\)（维数 \(3\)），并分解为
\[
\mathrm{Ind}_{D_4}^{S_4}\mathbf{1}\ \cong\ \mathbf{1}\ \oplus\ V_2.
\]
另一方面，\(X_e/S_4\cong\mathbb P^1\) 故 \(H^0(X_e,\Omega^1)\) 中不含平凡表示分量；
因此由三条椭圆商 \(R\) 的 pullback 在 \(J(X_e)\) 中生成的阿贝尔子簇恰对应于 \(V_2\) 分量，
且其“去掉对角平凡分量”后得到两份独立拷贝，从而给出同源 \(A_2\sim R^2\)。
\end{proof}

\begin{conclusion}[\texorpdfstring{$V_3$}{V3}-等型因子的椭圆立方化：\texorpdfstring{$A_3\sim E^3$}{A3~E^3}]\label{con:xi-terminal-zm-s4-jacobian-a3-isog-e3}
令
\[
E:=X_e/S_3
\]
为结论 \ref{con:xi-terminal-zm-s4-galois-closure-genus-tower-holomorphic-differentials} 的椭圆中间商（亦即结论 \ref{con:xi-terminal-zm-elliptic-normalization} 的椭圆归一化曲线）。
则 \(V_3\)-等型因子满足
\[
\boxed{\ A_{3}\ \sim\ E^{3}\ }.
\]
\end{conclusion}
\begin{proof}[证明要点]
四次根域对应子群为点稳定子 \(S_3\le S_4\)，其共轭类对应于四个根的选择。
于是 \(\mathrm{Ind}_{S_3}^{S_4}\mathbf{1}\) 给出次数 \(4\) 的置换表示，并分解为
\(
\mathrm{Ind}_{S_3}^{S_4}\mathbf{1}\cong \mathbf{1}\oplus V_3
\)。
同理，由于平凡分量不出现在 \(H^0(X_e,\Omega^1)\) 中，
四条共轭椭圆商 \(X_e/(gS_3g^{-1})\cong E\) 的 pullback 在 Jacobian 中生成的非平凡部分恰对应 \(V_3\) 等型因子，
从而给出同源 \(A_3\sim E^3\)。
\end{proof}

\begin{conclusion}[\texorpdfstring{$V_3'$}{V3'}-等型因子的 Prym 三维体识别：\texorpdfstring{$A_{3'}\sim P^3$}{A3'~P^3}]\label{con:xi-terminal-zm-s4-jacobian-a3p-isog-prym3-cubed}
取结论 \ref{con:xi-terminal-zm-s4-d4fixedfield-resolvent-elliptic} 中的二面体子群 \(D_4\) 的旋转子群
\(
C_4\lhd D_4
\)
（阶 \(4\) 的循环子群），并令
\[
X_4:=X_e/C_4,
\qquad
R:=X_e/D_4.
\]
则自然映射 \(X_4\to R\) 为次数 \(2\) 覆盖。
定义 Prym 三维体
\[
P:=\operatorname{Prym}(X_4/R):=\ker\!\bigl(\operatorname{Nm}_{X_4/R}:J(X_4)\to J(R)\bigr)^{0}.
\]
则 \(V_3'\)-等型因子满足
\[
\boxed{\ A_{3'}\ \sim\ P^{3}\ } .
\]
\end{conclusion}
\begin{proof}[证明要点]
对二次覆盖 \(X_4\to R\)，经典 Prym 分解给出
\(
J(X_4)\sim J(R)\times P
\)；
其中 \(\dim P=g(X_4)-g(R)\)。
\par
另一方面，\(C_4\) 在 \(S_4\) 中的陪集置换表示为 \(\mathrm{Ind}_{C_4}^{S_4}\mathbf{1}\)（维数 \(6\)），并满足分解
\[
\mathrm{Ind}_{C_4}^{S_4}\mathbf{1}\ \cong\ \mathbf{1}\ \oplus\ V_2\ \oplus\ V_3'.
\]
结合结论 \ref{con:xi-terminal-zm-s4-galois-closure-genus-tower-holomorphic-differentials} 的
\(
H^0(X_e,\Omega^1)\cong 2\cdot\mathrm{sgn}\oplus V_2\oplus V_3\oplus 3V_3'
\)
可见
\(
g(X_4)=\dim H^0(X_e,\Omega^1)^{C_4}=1+3=4
\)
且 \(g(R)=1\)，从而 \(\dim P=3\)。
\par
再由同源分解中 \(V_3'\) 的重数为 \(3\)，并注意 \(C_4\) 的共轭类恰给出三条同构的二次塔 \(X_e/C_4\to X_e/D_4\)，
其 Prym 三维体在同源意义下同构；由此把 \(V_3'\) 的三份拷贝几何化为 \(P^3\)，即得到 \(A_{3'}\sim P^3\)。
\end{proof}

\begin{conclusion}[\texorpdfstring{$S_4$}{S4}-闭包曲线 Jacobian 的四因子可见分解]\label{con:xi-terminal-zm-s4-jacobian-visible-factorization}
沿用上述记号，并记 \(C:=X_e/A_4\)（属 \(2\) 的判别式二次脊曲线，结论 \ref{con:xi-terminal-zm-s4-galois-closure-genus-tower-holomorphic-differentials}）。
则在 \(\QQ\) 上成立同源分解
\[
\boxed{\;
J(X_e)\ \sim\ J(C)\times R^{2}\times E^{3}\times P^{3}\; }.
\]
\end{conclusion}
\begin{proof}[证明要点]
由结论 \ref{con:xi-terminal-zm-s4-galois-closure-genus-tower-holomorphic-differentials} 已知
\(
J(X_e)\sim J(C)\times A_2\times A_3\times A_{3'}
\)。
再分别代入结论 \ref{con:xi-terminal-zm-s4-jacobian-a2-isog-r2}、
\ref{con:xi-terminal-zm-s4-jacobian-a3-isog-e3} 与
\ref{con:xi-terminal-zm-s4-jacobian-a3p-isog-prym3-cubed} 的同源识别，即得所示四因子可见分解。
\end{proof}

\begin{conclusion}[\texorpdfstring{$C_2$}{C2}-换位商的属 \(6\) 中间覆盖与 \texorpdfstring{$A_{2}\times A_{3}\times A_{3'}$}{A2xA3xA3'} 因子剥离]\label{con:xi-terminal-zm-s4-transposition-quotient-genus6}
沿用结论 \ref{con:xi-terminal-zm-s4-galois-closure-genus-tower-holomorphic-differentials} 的 \(S_4\)-伽罗瓦闭包曲线 \(X_e\to\PP^{1}_{y}\)（\(|S_4|=24\)）。
取任意由换位生成的子群 \(H\simeq C_{2}\subset S_4\)，并令
\[
X_{H}:=X_e/H.
\]
则商映射 \(X_{H}\to\PP^{1}_{y}\) 为次数 \(12\) 的有限态射，其分歧型由 \(S_4\) 在左陪集集合 \(S_4/H\) 上的作用给出，并满足：
\begin{enumerate}
  \item 在每个惯性为换位的分歧值 \(y\in\{0,1\}\cup\{P_{\mathrm{LY}}(y)=0\}\) 上，惯性元在 \(12\) 个陪集上的循环型为 \(1^{2}2^{5}\)；
  等价地，该纤维上恰有 \(5\) 个分歧指数 \(2\) 的点与 \(2\) 个不分歧点。
  \item 在无穷远分歧值 \(y=\infty\) 上，惯性 \(4\)-循环在 \(12\) 个陪集上的循环型为 \(4^{3}\)，即该纤维上恰有 \(3\) 个分歧指数 \(4\) 的点。
  \item 因而由 Riemann--Hurwitz 公式得到
  \[
  \boxed{\ g(X_H)=6\ }.
  \]
\end{enumerate}
并且在结论 \ref{con:xi-terminal-zm-s4-galois-closure-genus-tower-holomorphic-differentials} 的同源分解记号下，
\[
J(X_{H})\ \sim_{\QQ}\ A_{2}\times A_{3}\times A_{3'},
\]
亦即 \(X_H\) 的 Jacobian 恰对应于 \(J(X_e)\) 中去掉 \(\mathrm{sgn}\)-等型因子后的剩余三分量。
\end{conclusion}
\begin{proof}[证明要点]
分歧循环型：对任意 \(g\in S_4\)，其在陪集集 \(S_4/H\) 上的固定点个数为
\(
\#\{xH:\ gxH=xH\}=\#\{x\in S_4:\ x^{-1}gx\in H\}/|H|
\)。
当 \(g\) 为换位时，\(H\cap g^{S_4}\) 恰含一元素，且 \(|C_{S_4}(g)|=2\cdot 2!=4\)，故固定点数为 \(4/2=2\)，其余 \(10\) 点成 \(5\) 个 \(2\)-轮换，得循环型 \(1^{2}2^{5}\)。
当 \(g\) 为 \(4\)-循环时，\(H\cap g^{S_4}=\varnothing\)，故无固定点。
又 \(g^{2}\) 为双换位，仍有 \(H\cap (g^{2})^{S_4}=\varnothing\)，从而 \(g^{2}\) 在陪集集上亦无固定点；
因此 \(g\) 的循环分解不含 \(1\)-循环或 \(2\)-循环，只能为 \(4\)-循环，且 \(12\) 点必分解为 \(4^{3}\)。
\par
属数：对次数 \(12\) 覆盖应用 Riemann--Hurwitz 公式
\[
2g(X_H)-2=-2\cdot 12+\sum_{b}\sum_{P\mid b}(e_P-1).
\]
五个换位分歧值处，每个分歧值的总贡献为 \(5\cdot(2-1)=5\)；
无穷远处总贡献为 \(3\cdot(4-1)=9\)。
总分歧贡献 \(5\cdot 5+9=34\)，代入得 \(2g(X_H)-2=-24+34=10\)，故 \(g(X_H)=6\)。
\par
雅可比因子：由于 \(X_e\to X_H\) 为 Galois 商映射，有
\(
H^{0}(X_H,\Omega^{1})\cong H^{0}(X_e,\Omega^{1})^{H}
\)。
由结论 \ref{con:xi-terminal-zm-s4-galois-closure-genus-tower-holomorphic-differentials} 的表示分解
\(
H^{0}(X_e,\Omega^{1})\cong 2\cdot \mathrm{sgn}\oplus V_2\oplus V_3\oplus 3V_3'
\)
以及 \(S_4\) 角色表在换位处的迹值可知：\(\mathrm{sgn}\) 在 \(H\) 上无不变元，而 \(V_2,V_3,V_3'\) 的不变元维数分别为 \(1,2,1\)。
故
\(
\dim H^{0}(X_e,\Omega^{1})^{H}=1+2+3\cdot 1=6=g(X_H)
\)，并且 \(H\)-不变部分恰由 \(A_2,A_3,A_{3'}\) 三个等型因子贡献，从而得到 \(J(X_H)\sim_{\QQ}A_2\times A_3\times A_{3'}\)。证毕。
\end{proof}

\endinput

\paragraph{共轭重量支下的闭式审计链：临界域、分歧域与倍乘对应的算术分解}
本段在既有椭圆统一化（结论 \ref{con:xi-terminal-zm-elliptic-normalization}）与 Lee--Yang 分歧刻画（结论 \ref{con:xi-terminal-zm-leyang-ramification-critical-values}）的基础上，
将 Latt\`es 临界多项式、共轭重量覆盖的分歧几何与倍乘诱导的四值对应统一为一组可直接引用的闭式证书。

\begin{conclusion}[临界六次的 \(4\)-分点解释与分裂域刻画]\label{con:xi-terminal-zm-lattes-critical-sextic-4division-splitting-field}
沿用结论 \ref{con:xi-terminal-zm-elliptic-lattes-doubling} 的椭圆曲线
\(
E/\QQ:\ Y^{2}=X^{3}-X+\tfrac14
\)
与倍点 Latt\`es 自映射
\(
\Phi(X)=\frac{X^{4}+2X^{2}-2X+1}{4X^{3}-4X+1}
\)。
令
\[
\Psi(X):=2X^{6}-10X^{4}+10X^{3}-10X^{2}+2X+1\in\ZZ[X].
\]
则：
\begin{enumerate}
  \item \(\Psi\) 恰为 \(4\)-division 多项式 \(\psi_{4}\) 去除 \(2\)-division 因子后的本原因子；等价地，有限临界集合满足
  \[
  \Phi'(X)=0
  \ \Longleftrightarrow\
  \exists\,Q\in E[4]\setminus E[2]\ \text{使得}\ X=x(Q)
  \ \Longleftrightarrow\
  \Psi(X)=0.
  \]
  因而 \(\Phi\) 的六个有限临界点正是 \(x\bigl(E[4]\setminus E[2]\bigr)\)（在 \(\pm\) 识别下）的像。
  \item \(\Psi\) 在 \(\QQ\) 上不可约，且其判别式满足严格恒等式
  \[
  \mathrm{Disc}(\Psi)=-2^{8}\cdot 37^{5}.
  \]
  从而 \(\Psi\) 的分裂域 \(K_{4}/\QQ\) 的分歧素数集合包含于 \(\{2,37\}\)。
  \item 伽罗瓦群为阶 \(48\) 的传递群，且同构于
  \[
  \mathrm{Gal}(K_{4}/\QQ)\cong S_{4}\times C_{2}.
  \]
  特别地，六个有限临界点在 \(\mathrm{Gal}(\overline{\QQ}/\QQ)\) 作用下构成单一轨道。
\end{enumerate}
\end{conclusion}
\begin{proof}[证明要点]
第一点由结论 \ref{con:xi-terminal-zm-elliptic-lattes-doubling} 的 \(\psi_{4}/(2Y)\) 恒等式给出；
第二、三点由结论 \ref{con:xi-terminal-zm-elliptic-lattes-critical-sextic-galois} 的不可约性、判别式与伽罗瓦群计算推出；
传递性等价于“临界点构成单一伽罗瓦轨道”。
可复现核验脚本：\texttt{scripts/exp\_fold\_zm\_elliptic\_lattes\_rational\_points\_audit.py}。证毕。
\end{proof}

\begin{conclusion}[谱四次的显式椭圆化与共轭重量覆盖的分歧剖面]\label{con:xi-terminal-zm-elliptic-normalization-y-iota-ramification-portrait}
令平面曲线
\[
C:\quad \Pi(\lambda,y)=\lambda^{4}-\lambda^{3}-(2y+1)\lambda^{2}+\lambda+y(y+1)=0.
\]
在结论 \ref{con:xi-terminal-zm-elliptic-normalization} 的椭圆统一化下，\(C\) 与
\(
E:\ Y^{2}=X^{3}-X+\tfrac14
\)
之间存在完全显式的坐标同一（在仿射开集上为同构）：
\[
E\dashrightarrow C,\qquad (X,Y)\longmapsto (\lambda,y)=(X,\ X^{2}+Y-\tfrac12),
\]
\[
C\dashrightarrow E,\qquad (\lambda,y)\longmapsto (X,Y)=(\lambda,\ y+\tfrac12-\lambda^{2}).
\]
记物理主支重量函数与其共轭支
\[
y:=X^{2}+Y-\frac12\in\QQ(E),
\qquad
y^{\iota}:=X^{2}-Y-\frac12\in\QQ(E),
\qquad
P_{\mathrm{LY}}(t):=256t^{3}+411t^{2}+165t+32\in\ZZ[t].
\]
则 \(y:E\to\mathbb P^{1}\) 为次数 \(4\) 的覆盖，并满足以下全分歧剖面：
\begin{enumerate}
  \item 在无穷远点 \(O\) 处，\(y\) 具有唯一极点且极点阶为 \(4\)，从而在 \(y=\infty\) 处全分歧（\(e_{O}=4\)）。
  \item 有限分歧值集合恰为
  \[
  y\in\{0,1\}\ \cup\ V\!\bigl(P_{\mathrm{LY}}\bigr),
  \]
  并且每个有限分歧值均为单分歧（上方存在唯一分歧点，分歧指标 \(2\)）。
  \item 两个有理分歧值的纤维分解完全显式：
  \[
  (y)^{-1}(0)=\Bigl\{\Bigl(0,\frac12\Bigr),\ \Bigl(-1,-\frac12\Bigr),\ \underbrace{\Bigl(1,-\frac12\Bigr)}_{e=2}\Bigr\},
  \qquad
  (y)^{-1}(1)=\Bigl\{\Bigl(1,\frac12\Bigr),\ \Bigl(2,-\frac52\Bigr),\ \underbrace{\Bigl(-1,\frac12\Bigr)}_{e=2}\Bigr\}.
  \]
\end{enumerate}
\end{conclusion}
\begin{proof}[证明要点]
由结论 \ref{con:xi-terminal-zm-y-quadratic-extension-minpoly}，物理主支 \(y=X^{2}+Y-\tfrac12\) 与共轭支 \(y^{\iota}=X^{2}-Y-\tfrac12\) 互为共轭，且由椭圆反演 \(\iota:(X,Y)\mapsto(X,-Y)\) 交换。
结论 \ref{con:xi-terminal-zm-elliptic-y-hurwitz-passport} 给出 \(y\) 的 Hurwitz 护照与两个有理分歧纤维。
有限分歧值集合与 \(P_{\mathrm{LY}}\) 的出现由判别式分解（结论 \ref{con:xi-terminal-zm-discriminant-leyang}，等价地见结论 \ref{con:xi-terminal-zm-leyang-cubic-branch-image}）决定。证毕。
\end{proof}

\begin{conclusion}[Lee--Yang 三次因子的分歧消元产生式：共轭重量支的临界点立方刻画]\label{con:xi-terminal-zm-leyang-branch-elimination-origin-y-iota}
在结论 \ref{con:xi-terminal-zm-elliptic-normalization-y-iota-ramification-portrait} 的记号下，令
\[
g(X):=16X^{3}-9X^{2}+1.
\]
则 \(y:E\to\mathbb P^{1}\) 在 \(V\!\bigl(P_{\mathrm{LY}}\bigr)\) 上方的三条非有理分歧点恰为满足
\[
4XY+3X^{2}-1=0,\qquad X\neq \pm 1
\]
的三点；等价地可由
\[
g(X)=0,\qquad Y=\frac{1-3X^{2}}{4X}
\]
刻画。
更进一步，分歧值与上述立方分歧点之间存在严格消元恒等式：在分歧约束下有线性关系
\[
4Xy=4X^{3}-3X^{2}-2X+1,
\]
并且满足结果式证书
\[
\mathrm{Res}_{X}\!\Bigl(g(X),\ 4Xy-\bigl(4X^{3}-3X^{2}-2X+1\bigr)\Bigr)=-64\,P_{\mathrm{LY}}(y).
\]
因此 \(P_{\mathrm{LY}}\)（除去单位因子）可被解释为共轭重量覆盖的“非有理分歧值集合”在基底上的最小整系数消元像。
\end{conclusion}
\begin{proof}[证明要点]
与结论 \ref{con:xi-terminal-zm-leyang-ramification-critical-values} 的推导同理：由
\(
y=X^{2}+Y-\tfrac12
\)
与曲线微分关系 \(2Y\,\dd Y=(3X^{2}-1)\dd X\) 得
\(
\dd y=(4XY+3X^{2}-1)\,\omega
\)（\(\omega=\dd X/(2Y)\)）。
因此有限分歧点满足 \(4XY+3X^{2}-1=0\)；
消去 \(Y\) 得
\(
(X-1)(X+1)g(X)=0
\)，其中 \(X=\pm 1\) 对应有理分歧值 \(0,1\)，其余三点由 \(g(X)=0\) 给出。
线性关系 \(4Xy=4X^{3}-3X^{2}-2X+1\) 由代回 \(Y=(1-3X^{2})/(4X)\) 得到；
对变量 \(X\) 作消元即得结果式恒等式。
可复现核验脚本：\texttt{scripts/exp\_fold\_zm\_elliptic\_leyang\_cover\_geometry\_audit.py}。证毕。
\end{proof}

\begin{conclusion}[倍点的雅可比共周期：\texorpdfstring{$Y([2]P)$}{Y([2]P)} 与 \texorpdfstring{$\Phi'(X)$}{Phi'(X)} 的严格乘法公式]\label{con:xi-terminal-zm-elliptic-y-doubling-jacobian-cocycle}
沿用结论 \ref{con:xi-terminal-zm-elliptic-lattes-doubling} 的椭圆曲线 \(E\) 与倍点 Latt\`es 自映射 \(\Phi\)。
对任意 \(P=(X,Y)\in E(\overline{\QQ})\) 且 \(Y\neq 0\)，成立严格恒等式
\[
\boxed{\;
Y([2]P)=\frac{\Phi'(X)}{2}\,Y\;}
\]
进一步，对任意整数 \(n\ge 1\)，有
\[
\boxed{\;
Y([2^{n}]P)=\frac{(\Phi^{n})'(X)}{2^{n}}\,Y,
\qquad\text{等价地}\qquad
(\Phi^{n})'(X)=2^{n}\frac{Y([2^{n}]P)}{Y}\;}
\]
其中 \(\Phi^{n}\) 为 \(\Phi\) 的 \(n\) 次迭代。
\end{conclusion}
\begin{proof}[证明要点]
取 \(E\) 的不变微分 \(\omega=\dd X/(2Y)\)。
倍点映射满足 \([2]^{\ast}\omega=2\omega\)。
另一方面，由 \(x([2]P)=\Phi(X)\) 得
\[
[2]^{\ast}\omega=\frac{\dd(\Phi(X))}{2\,Y([2]P)}
=\frac{\Phi'(X)}{2\,Y([2]P)}\,\dd X.
\]
与 \(2\omega=\dd X/Y\) 比较即得 \(Y([2]P)=\frac{\Phi'(X)}{2}Y\)。
对 \(n\ge 1\) 迭代同理：由 \([2^{n}]^{\ast}\omega=2^{n}\omega\) 与 \(x([2^{n}]P)=\Phi^{n}(X)\) 得
\(
\frac{(\Phi^{n})'(X)}{2\,Y([2^{n}]P)}=\frac{2^{n}}{2Y}
\)，从而推出所示恒等式。证毕。
\end{proof}

\begin{conclusion}[重量参数的仿射共周期：\texorpdfstring{$y([2]P)$}{y([2]P)} 对 \texorpdfstring{$y(P)$}{y(P)} 的线性--平移分解]\label{con:xi-terminal-zm-elliptic-y-doubling-affine-cocycle-derivative}
沿用结论 \ref{con:xi-terminal-zm-elliptic-normalization-y-iota-ramification-portrait} 的重量函数
\(
y=X^{2}+Y-\tfrac12
\)。
对任意 \(P=(X,Y)\in E(\overline{\QQ})\) 且 \(Y\neq 0\)，有严格仿射恒等式
\[
\boxed{\;
y([2]P)
=\Phi(X)^{2}-\frac12+\frac{\Phi'(X)}{2}\Bigl(y-\bigl(X^{2}-\tfrac12\bigr)\Bigr)\;}
\]
亦即
\[
y([2]P)=\frac{\Phi'(X)}{2}\,y
+\Bigl(\Phi(X)^{2}-\frac12-\frac{\Phi'(X)}{2}\bigl(X^{2}-\tfrac12\bigr)\Bigr),
\]
其中括号项仅依赖于 \(X\)。
\par
更强地，记 \(2\)-division 多项式
\(
D(X)=4X^{3}-4X+1
\)，
则存在唯一的八次整式
\[
B_{8}(X)=X^{8}-7X^{6}+14X^{5}-27X^{4}+9X^{3}+6X^{2}-X-1\in\ZZ[X]
\]
使得
\[
\boxed{\;
y([2]P)=\frac{\Phi'(X)}{2}\,y-\frac{B_{8}(X)}{D(X)^{2}}\;}
\]
\end{conclusion}
\begin{proof}[证明要点]
由定义 \(y([2]P)=X([2]P)^{2}+Y([2]P)-\tfrac12=\Phi(X)^{2}+Y([2]P)-\tfrac12\)。
代入结论 \ref{con:xi-terminal-zm-elliptic-y-doubling-jacobian-cocycle} 的
\(
Y([2]P)=\frac{\Phi'(X)}{2}Y
\)
与恒等式 \(Y=y-(X^{2}-\tfrac12)\)，即得第一组公式。
最后的分式表示由把平移项化为分母 \(D(X)^{2}\) 的既约形式得到；
其显式分子与唯一性可由结论 \ref{con:xi-terminal-zm-elliptic-y-doubling-linear} 的整式证书直接读取（取 \(B_{8}=-B\)）。证毕。
\end{proof}

\begin{conclusion}[倍点拉回在重量二次扩张上的 Möbius 共周期：\texorpdfstring{$[2]^*y$}{[2]*y} 的 \texorpdfstring{$\mathrm{PGL}_2(\QQ(X))$}{PGL2(Q(X))} 表示]\label{con:xi-terminal-zm-elliptic-y-doubling-mobius-cocycle}
沿用结论 \ref{con:xi-terminal-zm-elliptic-normalization-y-iota-ramification-portrait} 的椭圆曲线
\(
E:\ Y^{2}=X^{3}-X+\tfrac14
\)
与重量函数
\(
y=X^{2}+Y-\tfrac12
\)。
记
\(
D(X):=4X^{3}-4X+1
\)。
则存在唯一的元素 \([M(X)]\in\mathrm{PGL}_{2}(\QQ(X))\)，使得在 \(\QQ(E)\) 中成立严格恒等式
\begin{equation}\label{eq:xi-elliptic-y-doubling-mobius}
\boxed{\;
[2]^{\ast}(y)=y([2]P)=\frac{a(X)\,y+b(X)}{c(X)\,y+d(X)}\;}
\end{equation}
其中可取代表矩阵
\[
M(X)=
\begin{pmatrix}
a(X) & b(X)\\
c(X) & d(X)
\end{pmatrix}\in\mathrm{Mat}_{2}(\QQ(X)),
\]
并且
\[
\begin{aligned}
a(X)&:=2\bigl(2X^{8}-8X^{6}-8X^{5}+44X^{4}-24X^{3}+1\bigr),\\
b(X)&:=-2\bigl(2X^{10}-4X^{9}-9X^{8}+16X^{7}+27X^{6}-20X^{5}-X^{4}-15X^{3}+10X^{2}+X-1\bigr),\\
c(X)&:=4\,D(X)^{2},\qquad
d(X):=-2\,(2X^{2}-1)\,D(X)^{2}.
\end{aligned}
\]
并且该归一化代表满足完全因子分解
\[
\boxed{\;
\det M(X)=-4\,D(X)^{3}\cdot\bigl(2X^{6}-10X^{4}+10X^{3}-10X^{2}+2X+1\bigr)\;}
\]
以及
\[
\boxed{\;
\tr M(X)=-4\cdot \psi_{3}(X)\cdot \bigl(5X^{4}-2X^{2}-X+1\bigr),\qquad
\psi_{3}(X)=3X^{4}-6X^{2}+3X-1\;}
\]
其中 \(\psi_{3}\) 为该短 Weierstrass 模型的三除多项式。
特别地，\(\det M(X)\) 的零除子（除去由 \(D(X)=0\) 产生的极点层）恰截获倍点 Latt\`es 自映射的临界六次多项式 \(2X^{6}-10X^{4}+10X^{3}-10X^{2}+2X+1\)；
从而重量二次扩张在 \([2]^{\ast}\) 下的 Möbius 退化与 \(x\)-坐标动力学的临界结构在同一显式除子上对齐。
\end{conclusion}
\begin{proof}[证明要点]
在短 Weierstrass 模型上使用标准倍点公式
\(
X([2]P)=\bigl(\frac{3X^{2}-1}{2Y}\bigr)^{2}-2X
\)
与
\(
Y([2]P)=\frac{3X^{2}-1}{2Y}\bigl(X-X([2]P)\bigr)-Y
\)，
并以 \(Y=y-(X^{2}-\tfrac12)\) 消去 \(Y\)。
再在重量二次扩张的最小多项式（结论 \ref{con:xi-terminal-zm-elliptic-normalization-y-iota-ramification-portrait}）下对 \(y^{2}\) 降次，即可得到 \eqref{eq:xi-elliptic-y-doubling-mobius} 的线性分式表示。
\(\det M\) 与 \(\tr M\) 的因子分解由直接计算并在 \(\QQ[X]\) 中因式分解得到，其中 \(\psi_{3}\) 为标准三除多项式。
可复现核验脚本：\texttt{scripts/exp\_fold\_zm\_elliptic\_y\_doubling\_mobius\_cocycle\_audit.py}。证毕。
\end{proof}

\begin{conclusion}[全二幂迭代的闭式仿射化：\texorpdfstring{$y([2^{n}]P)$}{y([2^n]P)} 由 \texorpdfstring{$\Phi^{n}$}{Phi^n} 与 \texorpdfstring{$(\Phi^{n})'$}{(Phi^n)'} 控制]\label{con:xi-terminal-zm-elliptic-y-doubling-iterate-affine-closed}
沿用结论 \ref{con:xi-terminal-zm-elliptic-normalization-y-iota-ramification-portrait} 的记号。
对任意整数 \(n\ge 1\) 与任意 \(P=(X,Y)\in E(\overline{\QQ})\) 且 \(Y\neq 0\)，成立严格闭式表示
\[
\boxed{\;
y([2^{n}]P)
=\Phi^{n}(X)^{2}-\frac12+\frac{(\Phi^{n})'(X)}{2^{n}}\Bigl(y-\bigl(X^{2}-\tfrac12\bigr)\Bigr)\;}
\]
等价地，
\[
y([2^{n}]P)
=\frac{(\Phi^{n})'(X)}{2^{n}}\,y
+\Bigl(\Phi^{n}(X)^{2}-\frac12-\frac{(\Phi^{n})'(X)}{2^{n}}\bigl(X^{2}-\tfrac12\bigr)\Bigr).
\]
\end{conclusion}
\begin{proof}[证明要点]
由定义 \(y([2^{n}]P)=X([2^{n}]P)^{2}+Y([2^{n}]P)-\tfrac12=\Phi^{n}(X)^{2}+Y([2^{n}]P)-\tfrac12\)。
代入结论 \ref{con:xi-terminal-zm-elliptic-y-doubling-jacobian-cocycle} 的
\(
Y([2^{n}]P)=\frac{(\Phi^{n})'(X)}{2^{n}}Y
\)，再用 \(Y=y-(X^{2}-\tfrac12)\) 消去 \(Y\) 即得。证毕。
\end{proof}

\begin{conclusion}[重量变量上的四次倍乘对应：\texorpdfstring{$y([2]P)$}{y([2]P)} 的整系数最小多项式]\label{con:xi-terminal-zm-elliptic-weight-doubling-quartic-minpoly}
沿用结论 \ref{con:xi-terminal-zm-elliptic-normalization-y-iota-ramification-portrait} 的重量函数 \(y\)，并令
\[
y_0:=y(P),\qquad y_1:=y([2]P)\qquad(P\in E(\overline{\QQ})).
\]
则 \(y_1\) 在基域 \(\QQ(y_0)\) 上满足唯一的原始四次整系数关系
\[
\mathcal H(y_0,y_1)=0,
\qquad
\mathcal H\in\ZZ[y_0,y_1],
\]
并且 \(\mathcal H\) 在 \(\QQ(y_0)[y_1]\) 中不可约；从而 \(\mathcal H\) 即为 \(y_1\) 在 \(\QQ(y_0)\) 上的最小多项式。
其显式系数证书为：
\input{sections/generated/eq_fold_zm_elliptic_weight_doubling_H}
其中首项系数为完全四次幂
\[
A_{4}(y_0)=\bigl(16(y_0)^{3}-8(y_0)^{2}-4y_0+1\bigr)^{4},
\]
而常数项具有非平凡因式分解
\[
\boxed{\;
A_{0}(y_0)
=(y_0+1)^{2}\bigl((y_0)^{3}-7(y_0)^{2}+13y_0-5\bigr)^{2}
\bigl((y_0)^{4}-6(y_0)^{3}+66(y_0)^{2}+28y_0+9\bigr)
\bigl((y_0)^{4}+2(y_0)^{3}+8(y_0)^{2}+8y_0+8\bigr)\; }.
\]
\end{conclusion}
\begin{proof}[证明要点]
由结论 \ref{con:xi-terminal-zm-elliptic-normalization-y-iota-ramification-portrait}，覆盖 \(y:E\to\mathbb P^{1}\) 的次数为 \(4\)，从而在函数域层面
\(
[\QQ(E):\QQ(y_0)]=4
\)。
另一方面，\(y_1=[2]^*(y)\in\QQ(E)\)。
将倍点公式与 \(y=X^{2}+Y-\tfrac12\) 联立并对 \((X,Y)\) 消元，可得到 \((y_0,y_1)\) 的平面整系数关系；
取其在 \(\ZZ[y_0,y_1]\) 中的既约本原部分即为 \(\mathcal H\)。
再取一处无分歧专门化（例如 \(y_0=2\)，此时判别式非零）并核验所得四次多项式在 \(\QQ[y_1]\) 中不可约，即推出 \(\mathcal H\) 在 \(\QQ(y_0)[y_1]\) 中不可约；
由于扩张次数为 \(4\)，不可约四次必为最小多项式。
可复现核验脚本：\texttt{scripts/exp\_fold\_zm\_elliptic\_weight\_doubling\_audit.py}。证毕。
\end{proof}

\begin{conclusion}[判别式的平方自由核刚性与泛型 \texorpdfstring{$S_{4}$}{S4}：Lee--Yang 二次脊柱不增生]\label{con:xi-terminal-zm-elliptic-weight-doubling-discriminant-squareclass-s4}\label{con:xi-terminal-zm-elliptic-weight-doubling-quartic-elimination-disc}
令 \(\mathcal H\) 如结论 \ref{con:xi-terminal-zm-elliptic-weight-doubling-quartic-minpoly}。
将 \(\mathcal H(y,v)\) 视为关于 \(v\) 的四次多项式，则其判别式具有完全可审计的显式分解：
\[
\boxed{\;
\mathrm{Disc}_{v}\bigl(\mathcal H(y,v)\bigr)=-y(y-1)P_{\mathrm{LY}}(y)\cdot(\text{完全平方})\;}
\]
更精确地，对共轭重量支的另一条不可约四次对应 \(R_{-}\) 一并有
\[
\mathrm{Disc}_{v}\bigl(R_{\pm}(y,v)\bigr)=-y(y-1)P_{\mathrm{LY}}(y)\,A(y)^{2}\,B_{\pm}(y)^{2},
\]
其中 \(A\in\ZZ[y]\) 为次数 \(12\) 因子，\(B_{\pm}\in\ZZ[y]\) 为次数 \(26\) 因子（其显式展开见下列闭式证书）：
\input{sections/generated/eq_fold_zm_elliptic_weight_doubling_discriminant}
因此
\[
\boxed{\;
\bigl(\mathrm{Disc}_{v}(\mathcal H)\bigr)_{\mathrm{sf}}=-y(y-1)P_{\mathrm{LY}}(y)\;}
\]
并且在 \(\QQ(y)\) 上，分裂域的唯一二次分辨率子域恒等于
\[
\QQ(y)\Bigl(\sqrt{-y(y-1)P_{\mathrm{LY}}(y)}\Bigr).
\]
此外，作为 \(\QQ(y)\) 上的四次多项式，\(R_{+}(y,\cdot)=\mathcal H(y,\cdot)\) 与 \(R_{-}(y,\cdot)\) 的泛型伽罗瓦群均同构于 \(S_{4}\)。
\end{conclusion}
\begin{proof}[证明要点]
判别式分解由对 \(R_{\pm}\) 直接计算 \(\mathrm{Disc}_{v}\) 并在 \(\ZZ[y]\) 中因式分解得到；其平方自由部分随即给出唯一二次分辨率子域。
\par
对泛型伽罗瓦群，注意 \(\mathrm{Disc}_{v}(\mathcal H)\) 在 \(\QQ(y)\) 中非平方，从而排除 \(A_{4}\)。
再取一处无分歧专门化（例如 \(y=2\)）并核验所得四次多项式在 \(\QQ\) 上的伽罗瓦群为 \(S_{4}\)；
由于专门化伽罗瓦群为泛型伽罗瓦群之子群，而 \(S_{4}\) 已为四次情形的最大可能，故泛型必为 \(S_{4}\)。
上述核验可由脚本 \texttt{scripts/exp\_fold\_zm\_elliptic\_weight\_doubling\_audit.py} 复核。证毕。
\end{proof}

由此，重量参数在倍点动力学下的非线性被 \(\Phi\) 的导数共周期严格线性化（结论 \ref{con:xi-terminal-zm-elliptic-y-doubling-jacobian-cocycle}--\ref{con:xi-terminal-zm-elliptic-y-doubling-iterate-affine-closed}），
而消元后产生的四次倍乘对应在判别式层面把唯一二次障碍刚性锁定为同一 Lee--Yang 数据 \(-y(y-1)P_{\mathrm{LY}}(y)\)（结论 \ref{con:xi-terminal-zm-elliptic-weight-doubling-discriminant-squareclass-s4}）。

\begin{conclusion}[Lee--Yang 立方域上分歧点的显式坐标化与隐藏平方结构]\label{con:xi-terminal-zm-leyang-branchpoint-explicit-coordinates-hidden-square}
令 \(y_0\) 为 \(P_{\mathrm{LY}}(y)=0\) 的任一根，并记 \(K:=\QQ(y_0)\)。
则在 \(K\) 上存在唯一分歧点 \(Q\in E(K)\) 满足
\[
y(Q)=y_0,\qquad \dd y(Q)=0.
\]
其坐标可由 \(y_0\) 显式表示为
\[
\boxed{\;
X(Q)=\frac{-512y_0^{2}-518y_0-5}{279},\qquad
Y(Q)=\frac{512y_0^{2}+1262y_0+377}{558}
=\frac{(512y_0^{2}+518y_0+5)^{2}-25947}{372\,\bigl(512y_0^{2}+518y_0+5\bigr)}\; }.
\qquad (25947=3^{3}\cdot 31^{2})
\]
特别地，\(X(Q)\) 满足立方证书 \(g(X(Q))=0\)，其中 \(g(X)=16X^{3}-9X^{2}+1\)（见结论 \ref{con:xi-terminal-zm-leyang-branch-elimination-origin-y-iota}）。
\par
此外，域 \(K\) 中存在一条非显然的“平方”恒等式：定义
\[
s:=\frac{-512y_0^{2}+226y_0+367}{93}\in K,
\]
则严格成立
\[
\boxed{\;s^{2}=64y_0^{2}+64y_0+25\ \text{于 }K\text{ 中}\; }.
\]
该恒等式给出一条显式平方证书，表明 Lee--Yang 立方域中出现的分歧点坐标可在不引入额外二次扩张的前提下完成闭式表达。
\end{conclusion}
\begin{proof}[证明要点]
在商环 \(\QQ[y,X]/(P_{\mathrm{LY}}(y))\) 中，对系统 \(\Pi(X,y)=0\)、\(\partial_{X}\Pi(X,y)=0\) 作 Gröbner 消元可得到一次消元子
\(
279X+512y^{2}+518y+5=0
\)，从而给出 \(X(Q)\) 的显式表达；
再由 \(y=X^{2}+Y-\tfrac12\) 回代得到 \(Y(Q)\)。
隐藏平方恒等式为同一三次商环中的代数恒等式检验。
可复现核验脚本：\texttt{scripts/exp\_fold\_zm\_elliptic\_leyang\_resy\_spectral\_decomposition\_audit.py}。证毕。
\end{proof}

\begin{conclusion}[Lee--Yang 分歧值处谱四次的精确退化：二重根与剩余二次因子的闭式分解]\label{con:xi-terminal-zm-leyang-branch-spectral-quartic-double-root-factorization}
令 \(y_0\) 为 \(P_{\mathrm{LY}}(y)=0\) 的任一根，并记 \(K:=\QQ(y_0)\)。
考虑专门化谱四次
\[
\Pi(\lambda,y_0)
=\lambda^{4}-\lambda^{3}-(2y_0+1)\lambda^{2}+\lambda+y_0(y_0+1)\ \in\ K[\lambda].
\]
则存在严格分解
\[
\boxed{\;
\Pi(\lambda,y_0)=\bigl(\lambda-X(Q)\bigr)^{2}\cdot Q_{y_0}^{(2)}(\lambda)\; },
\]
其中 \(Q\in E(K)\) 为结论 \ref{con:xi-terminal-zm-leyang-branchpoint-explicit-coordinates-hidden-square} 的唯一分歧点，且剩余二次因子可取为
\[
\boxed{\;
Q_{y_0}^{(2)}(\lambda)=\lambda^{2}
-\frac{1024y_0^{2}+1036y_0+289}{279}\,\lambda
+\frac{256y_0^{2}-578y_0-416}{279}\; }.
\]
该分解把 Lee--Yang 分歧值处 \(\Pi(\lambda,y_0)=0\) 的谱退化精确量化为：一处由分歧点给出的二重根与两条非分歧谱支所决定的剩余二次片。
\end{conclusion}
\begin{proof}[证明要点]
在 \(P_{\mathrm{LY}}(y_0)=0\) 的 Lee--Yang 立方域上，系统
\(
\Pi(\lambda,y_0)=\partial_{\lambda}\Pi(\lambda,y_0)=0
\)
在 \(K\) 中具有唯一解 \(\lambda=X(Q)\)，其中 \(Q\) 为结论 \ref{con:xi-terminal-zm-leyang-branchpoint-explicit-coordinates-hidden-square} 的分歧点。
因此 \(\lambda=X(Q)\) 为 \(\Pi(\lambda,y_0)\) 的二重根；并且由单分歧性可知其重数恰为 \(2\)。
将 \(\Pi(\lambda,y_0)\) 除以 \((\lambda-X(Q))^{2}\) 并在 \(K\) 中化简即可得到剩余二次因子的闭式系数。
可复现核验脚本：\texttt{scripts/exp\_fold\_zm\_elliptic\_leyang\_resy\_spectral\_decomposition\_audit.py}。证毕。
\end{proof}

\begin{conclusion}[两类二次伴随场的分化：临界域与 Lee--Yang 分歧域的算术阴影]\label{con:xi-terminal-zm-quadratic-companion-fields-separation}
该模型中出现两条彼此不同但同由椭圆判别式素数 \(37=\Delta(E)\) 驱动的二次伴随场：
\begin{enumerate}
  \item 对临界六次多项式 \(\Psi\)，其判别式平方自由核为 \(-37\)（结论 \ref{con:xi-terminal-zm-lattes-critical-sextic-4division-splitting-field}），因此
  \[
  \QQ(\sqrt{-37})\ \subset\ K_4
  \]
  必嵌入临界点分裂域 \(K_4/\QQ\)。
  \item 对 Lee--Yang 三次分歧（等价地，对 \(P_{\mathrm{LY}}\) 或结论 \ref{con:xi-terminal-zm-leyang-cubic-arithmetic} 的 \(S_3\) 分裂域），其判别式平方自由核为 \(-111=-3\cdot 37\)，因此
  \[
  \QQ(\sqrt{-111})\ \subset\ L_{\mathrm{LY}}
  \]
  必嵌入 Lee--Yang 分歧域的伽罗瓦闭包 \(L_{\mathrm{LY}}/\QQ\)。
\end{enumerate}
两者分别对应 \(-\Delta(E)\) 与 \(-3\Delta(E)\) 的二次阴影，
从而在最低阶二次信息上，动力学临界域（由 \(\Phi\) 的临界点控制）与分歧/Lee--Yang 域（由 \(y:E\to\mathbb P^{1}\) 的分歧控制）已发生不可约的算术分化。
\end{conclusion}
\begin{proof}[证明要点]
第一点由结论 \ref{con:xi-terminal-zm-elliptic-lattes-critical-sextic-galois} 的判别式平方类立即推出；
第二点由结论 \ref{con:xi-terminal-zm-leyang-cubic-arithmetic} 的判别式平方自由部分与不可约性推出，并由不可约三次的标准结构得到其唯一二次子域。
二次阴影的对比仅为代入 \(\Delta(E)=37\) 的算术重写。证毕。
\end{proof}

\begin{conclusion}[二次判别角色的最小同时消障场：\texorpdfstring{$\QQ(\sqrt{-111},\sqrt{-37})$}{Q(sqrt(-111),sqrt(-37))}]\label{con:xi-terminal-zm-quadratic-simultaneous-obstruction-field}
令 \(L_{\mathrm{LY}}\) 为 Lee--Yang 三次 \(P_{\mathrm{LY}}(y)\) 的分裂域，并令 \(K_4\) 为临界六次 \(\Psi(X)\)（结论 \ref{con:xi-terminal-zm-lattes-critical-sextic-4division-splitting-field}）的分裂域。
则由结论 \ref{con:xi-terminal-zm-quadratic-companion-fields-separation} 可知
\[
\QQ(\sqrt{-111})\subset L_{\mathrm{LY}},
\qquad
\QQ(\sqrt{-37})\subset K_4.
\]
定义双二次扩张
\[
F:=\QQ(\sqrt{-111},\sqrt{-37}).
\]
则 \(F/\QQ\) 为 \(V_4\)-伽罗瓦扩张，且为同时包含上述两条二次子域的最小数域：任何同时包含 \(\QQ(\sqrt{-111})\) 与 \(\QQ(\sqrt{-37})\) 的数域必包含 \(F\)。
因此，Lee--Yang 分歧域与临界域在二次层级上出现的两条判别角色在 \(F\) 上同时平凡化；其二次障碍被统一压缩为一个可枚举的双二次框架。
并且 \(F/\QQ\) 的分歧素数集合包含于
\[
\{2,3,37\}.
\]
\end{conclusion}
\begin{proof}[证明要点]
由结论 \ref{con:xi-terminal-zm-leyang-cubic-field-minimal-discriminant}，\(\QQ(\sqrt{-111})\) 为 \(L_{\mathrm{LY}}\) 的二次分辨子域；
由结论 \ref{con:xi-terminal-zm-elliptic-lattes-critical-sextic-galois}，\(\QQ(\sqrt{-37})\) 嵌入 \(K_4\)。
二者的合成域 \(F=\QQ(\sqrt{-111})\QQ(\sqrt{-37})\) 即为最小同时容纳两子域的扩张，且为双二次（故伽罗瓦群同构于 \(V_4\)）。
\par
分歧素数集合的包含性由二次域判别式给出：
\(\mathrm{Disc}\bigl(\QQ(\sqrt{-111})\bigr)=-111\) 与 \(\mathrm{Disc}\bigl(\QQ(\sqrt{-37})\bigr)=-148\)，
故二者分歧素数分别包含于 \(\{3,37\}\) 与 \(\{2,37\}\)；
而合成域的分歧素数只能来自二者分歧素数之并集，因而包含于 \(\{2,3,37\}\)。证毕。
\end{proof}

\begin{conclusion}[不同元 \texorpdfstring{$\delta$}{delta} 的迹对偶与一行反演：四次幂基上的系数可审计重构]\label{con:xi-terminal-zm-delta-trace-dual-inversion}
设 \(k:=\QQ(y)\) 且 \(K:=\QQ(E)\)。
在首一四次表示
\[
K\cong k[X]/(F),
\qquad
F(X,y):=X^4-X^3-(2y+1)X^2+X+y(y+1)
\]
中，令 \(X\in K\) 为 \(X\) 的像，并令
\[
\delta:=4XY+3X^2-1\in K
\]
为结论 \ref{con:xi-terminal-zm-leyang-ramification-critical-values} 的 Lee--Yang 雅可比元（从而 \(\dd y=\delta\,\omega\)），并记
\[
\delta_{\mathrm{diff}}:=\partial_XF(X,y)=1-3X^2-4XY=-\delta.
\]
则在扩张 \(K/k\) 中成立严格迹恒等式
\[
\Tr_{K/k}\!\left(\frac{X^m}{\delta_{\mathrm{diff}}}\right)=0\quad(m=0,1,2),
\qquad
\Tr_{K/k}\!\left(\frac{X^3}{\delta_{\mathrm{diff}}}\right)=1.
\]
写
\[
F(X,y)=X^4+a_3X^3+a_2X^2+a_1X+a_0
\qquad
(a_3=-1,\ a_2=-(2y+1),\ a_1=1,\ a_0=y(y+1)),
\]
并定义 Horner 截断多项式
\[
q_0:=1,\qquad
q_1:=X+a_3=X-1,\qquad
q_2:=X^2+a_3X+a_2=X^2-X-(2y+1),
\]
\[
q_3:=X^3+a_3X^2+a_2X+a_1=X^3-X^2-(2y+1)X+1.
\]
则两组 \(k\)-基
\[
\{1,X,X^2,X^3\},
\qquad
\left\{\frac{q_3}{\delta_{\mathrm{diff}}},\frac{q_2}{\delta_{\mathrm{diff}}},\frac{q_1}{\delta_{\mathrm{diff}}},\frac{q_0}{\delta_{\mathrm{diff}}}\right\}
\]
互为迹对偶基。
特别地，对任意 \(a\in K\) 有一行反演公式
\[
a=\sum_{i=0}^{3}\Tr_{K/k}\!\left(\frac{aq_i}{\delta_{\mathrm{diff}}}\right)X^{3-i}.
\]
由于 \(\delta_{\mathrm{diff}}=-\delta\)，亦可等价写成
\(
a=-\sum_{i=0}^{3}\Tr_{K/k}\!\left(\frac{aq_i}{\delta}\right)X^{3-i}
\)；
该重构通道表明幂基系数的分母结构完全由同一不同理想生成元（Lee--Yang 雅可比元）控制。
\end{conclusion}
\begin{proof}[证明要点]
由于 \(F\) 在 \(k\) 上可分，\(K/k\) 为可分扩张，且 \(\delta_{\mathrm{diff}}=\partial_XF(X,y)\) 生成不同理想（结论 \ref{con:xi-terminal-zm-leyang-discriminant-norm-identity}）。
设 \(X_1,\dots,X_4\) 为 \(F(\cdot,y)\) 在代数闭包中的四根，则对任意 \(m\ge 0\) 有
\[
\Tr_{K/k}\!\left(\frac{X^{m}}{\delta_{\mathrm{diff}}}\right)=\sum_{j=1}^{4}\frac{X_j^{m}}{F'(X_j)}.
\]
将部分分式分解
\(
1/F(T)=\sum_{j}(F'(X_j))^{-1}(T-X_j)^{-1}
\)
在 \(T=\infty\) 展开并比较 \(T^{-1}\) 的系数，即得
\(
\sum_j X_j^{m}/F'(X_j)=0
\)（\(m\le 2\)）与
\(
\sum_j X_j^{3}/F'(X_j)=1
\)。
将 \(k\)-基 \(\{1,X,X^2,X^3\}\) 与候选族 \(\{X^3/\delta_{\mathrm{diff}},X^2/\delta_{\mathrm{diff}},X/\delta_{\mathrm{diff}},1/\delta_{\mathrm{diff}}\}\) 的迹配对写成矩阵，可得其为对角为 \(1\)、上三角为 \(0\) 的单位下三角矩阵；
其逆矩阵给出迹对偶基在上述候选族中的唯一线性校正。
对本四次模型，该校正恰由 Horner 截断 \(q_0,q_1,q_2,q_3\) 给出，从而得到所示迹对偶基与反演公式。
可复现核验脚本：\texttt{scripts/exp\_fold\_zm\_elliptic\_trace\_dual\_delta\_audit.py}。证毕。
\end{proof}

\input{con__xi-terminal-zm-elliptic-translation-q-biquartic}
\input{con__xi-terminal-zm-translation-t-branch-discriminant-c3}


\begin{theorem}[Lee--Yang 根处平方根系数 \texorpdfstring{$\kappa_*^2$}{kappa*^2} 的三次数域刻画与判别式]\label{thm:xi-terminal-zm-kappa-square-cubic-field-s3}\label{prop:xi-terminal-zm-branch-puiseux-coefficient-cubic}
\leanverified{paper\_xi\_terminal\_zm\_kappa\_square\_cubic\_field\_s3}
\leanverified{paper\_xi\_terminal\_zm\_branch\_puiseux\_coefficient\_cubic}
沿用谱四次
\[
\Pi(\lambda,y):=\lambda^4-\lambda^3-(2y+1)\lambda^2+\lambda+y(y+1)
\]
及其分歧集合 \(B_{\mathrm{fin}}=\{0,1\}\cup\{P_{\mathrm{LY}}(y)=0\}\) 的记号（结论 \ref{con:xi-terminal-zm-discriminant-leyang}）。
取任一非平凡分歧值 \(y_0\in\CC\) 满足 \(P_{\mathrm{LY}}(y_0)=0\)，并令 \(\lambda_0\) 为 \(\Pi(\lambda,y_0)=0\) 的唯一二重根（等价地 \(\Pi(\lambda_0,y_0)=\partial_\lambda\Pi(\lambda_0,y_0)=0\)）。
则存在局部参数 \(t\) 使 \(y-y_0=t^2\)，并且并合的两条谱支满足平方根型 Puiseux 展开
\[
\lambda=\lambda_0\pm c_0\,t+O(t^2),
\qquad
c_0^{2}=-\frac{2\,\partial_{y}\Pi(\lambda_0,y_0)}{\partial_{\lambda}^{2}\Pi(\lambda_0,y_0)}.
\]
定义
\[
u:=c_0^2=\kappa_*^2.
\]
在该 Lee--Yang 分歧支上进一步有
\[
\lambda_0\ \text{满足}\ 16\lambda_0^3-9\lambda_0^2+1=0
\qquad(\text{结论 \ref{con:xi-terminal-zm-leyang-lambda-cubic}}),
\]
并且 \(u\) 可化简为
\[
\boxed{\;
u=\frac{2}{3}\cdot\frac{3\lambda_0^2-1}{\lambda_0^2-1}\;}
\]
从而
\[
\boxed{\;324u^{3}-540u^{2}+333u-74=0\; }.
\]
该三次在 \(\QQ\) 上不可约，且其判别式为
\[
\mathrm{Disc}(324u^{3}-540u^{2}+333u-74)=-2^{6}\cdot 3^{9}\cdot 37.
\]
因而分裂域伽罗瓦群同构于 \(S_3\)。
此外
\[
\lambda_0=\frac{3(2u-1)}{4(3u-2)},
\]
故 \(\QQ(u)=\QQ(\lambda_0)\)。
\end{theorem}
\begin{proof}[证明要点]
平方根型 Puiseux 与系数公式由 \(\Pi(\lambda_0,y_0)=\partial_\lambda\Pi(\lambda_0,y_0)=0\) 处的二元 Taylor 展开得到，见结论 \ref{con:xi-terminal-zm-branch-puiseux-uniformization} 与定理 \ref{thm:xi-terminal-zm-leyang-elliptic-structure} 中的一般公式。
在 Lee--Yang 分歧支 \(P_{\mathrm{LY}}(y_0)=0\) 上，由结论 \ref{con:xi-terminal-zm-leyang-lambda-cubic} 得 \(\lambda_0\) 满足 \(16\lambda^3-9\lambda^2+1=0\)，并可将
\(
u=-2\,\partial_y\Pi/\partial_\lambda^2\Pi
\)
化简为所示有理式。
\par
由
\(
u=\frac{2}{3}\frac{3\lambda_0^2-1}{\lambda_0^2-1}
\)
可直接解得
\[
\lambda_0^2=\frac{3u-2}{3(u-2)}.
\]
又因 \(\lambda_0\neq 0\)，将 \(16\lambda_0^3-9\lambda_0^2+1=0\) 除以 \(\lambda_0^2\) 得
\[
\lambda_0=\frac{9\lambda_0^2-1}{16\lambda_0^2}
=\frac{3(2u-1)}{4(3u-2)}.
\]
将其代回 \(16\lambda_0^3-9\lambda_0^2+1=0\) 并清分母，即得
\[
324u^{3}-540u^{2}+333u-74=0.
\]
有理根检验排除一次因子，故该三次在 \(\QQ\) 上不可约。
其判别式按三次判别式公式计算为
\(
-2^{6}\cdot 3^{9}\cdot 37
\)，非平方。
对不可约三次，判别式非平方当且仅当分裂域伽罗瓦群为 \(S_3\)，故结论成立。
最后，由上式 \(\lambda_0=\frac{3(2u-1)}{4(3u-2)}\) 得 \(\lambda_0\in\QQ(u)\)，而 \(u\in\QQ(\lambda_0)\) 由定义即得，故 \(\QQ(u)=\QQ(\lambda_0)\)。证毕。
\end{proof}

\begin{corollary}[Lee--Yang 分歧数据在 \texorpdfstring{$\kappa_*^2$}{kappa*^2} 下的二次重构]\label{cor:xi-terminal-zm-leyang-branchdata-quadratic-reconstruction}
\leanverified{paper\_xi\_terminal\_zm\_leyang\_branchdata\_quadratic\_reconstruction}
沿用定理 \ref{thm:xi-terminal-zm-kappa-square-cubic-field-s3} 的记号，并令 \(\beta_*\) 为定理 \ref{thm:xi-terminal-zm-leyang-elliptic-structure} 中 Puiseux 展开
\[
\lambda(y)=\lambda_0\pm \kappa_*(y-y_0)^{1/2}+\beta_*(y-y_0)+O\bigl((y-y_0)^{3/2}\bigr)
\]
的线性项系数。
则在三次数域 \(K:=\QQ(u)=\QQ(\lambda_0)=\QQ(y_0)\) 内成立以下恒等式：
\[
16\lambda_0+31=108u(1-u),
\]
\[
256y_0+1091=108u(25-21u),
\]
\[
72\beta_*+22=9u(10u-1).
\]
等价地，
\[
\lambda_0=-\frac{31}{16}+\frac{27}{4}u-\frac{27}{4}u^2,\qquad
y_0=-\frac{1091}{256}+\frac{675}{64}u-\frac{567}{64}u^2,\qquad
\beta_*=-\frac{11}{36}-\frac{1}{8}u+\frac{5}{4}u^2.
\]
\end{corollary}
\begin{proof}
由定理 \ref{thm:xi-terminal-zm-kappa-square-cubic-field-s3}，\(u\) 满足三次关系 \(324u^{3}-540u^{2}+333u-74=0\)，且 \(\lambda_0\in\QQ(u)\)。
在商环 \(\QQ[u]/(324u^{3}-540u^{2}+333u-74)\) 中可直接核验
\(
16\lambda_0+31-108u(1-u)=0
\)，从而得到 \(\lambda_0\) 的二次表达式。
将其代入定理 \ref{thm:xi-terminal-zm-leyang-elliptic-structure} 的临界值显式式
\(
y_0=\frac{4\lambda_0^3-3\lambda_0^2-2\lambda_0+1}{4\lambda_0}
\)
与 \(\beta_*\) 的有理表达式（同一定理），并在同一商环内化简，即得其余两式。
\end{proof}

\begin{corollary}[\texorpdfstring{$\beta_*$}{beta*} 的极小多项式与显式整性化]\label{cor:xi-terminal-zm-beta-minpoly-integralization}
\leanverified{paper\_xi\_terminal\_zm\_beta\_minpoly\_integralization}
沿用上文记号。则 \(\beta_*\) 满足不可约三次方程
\[
419904\,\beta^{3}+93312\,\beta^{2}+72279\,\beta-1112=0.
\]
并且令
\[
T:=162u,\qquad S:=52488\,\beta_*,
\]
则 \(T,S\) 为 \(K\) 的代数整数，且分别满足首一整系数多项式
\[
T^{3}-270T^{2}+26973T-971028=0,
\]
\[
S^{3}+11664S^{2}+474222519S-382943630016=0.
\]
\end{corollary}
\begin{proof}
由推论 \ref{cor:xi-terminal-zm-leyang-branchdata-quadratic-reconstruction}，\(\beta_*\in\QQ(u)\) 且为 \(u\) 的二次多项式。
消去 \(u\)（或等价地在 \(\QQ[u]/(324u^{3}-540u^{2}+333u-74)\) 中化简）得到 \(\beta_*\) 的三次整系数关系；其不可约性由有理根判别排除一次因子。
对 \(T=162u\) 与 \(S=52488\beta_*\) 作变量伸缩并清分母，可得到所示首一整系数多项式，从而 \(T,S\) 为代数整数。
\end{proof}

\begin{proposition}[Lee--Yang 三次因子的二重覆盖与分歧核的完全拉回分解]\label{prop:xi-terminal-zm-leyang-double-cover-pullback-factorization}
\leanverified{paper\_xi\_terminal\_zm\_leyang\_double\_cover\_pullback\_factorization}
设
\[
g(u):=324u^{3}-540u^{2}+333u-74,\qquad
\iota(u):=\frac{25}{21}-u,
\]
并定义二次映射
\[
y(u):=-\frac{1091}{256}+\frac{675}{64}u-\frac{567}{64}u^2.
\]
则 \(y(\iota(u))=y(u)\)，并且 Lee--Yang 三次因子与全有限分歧核在 \(u\)-坐标下满足显式分解
\[
P_{\mathrm{LY}}(y(u))=\frac{3^{4}\,7^{3}}{2^{14}}\; g(u)\,g(\iota(u)),
\]
\[
y(u)\bigl(y(u)-1\bigr)P_{\mathrm{LY}}(y(u))
=\frac{3^{5}\,7^{3}}{2^{30}}\,
\bigl(2268u^2-2700u+1091\bigr)\,
\bigl(756u^2-900u+449\bigr)\,
g(u)\,g(\iota(u)).
\]
\end{proposition}
\begin{proof}
由推论 \ref{cor:xi-terminal-zm-leyang-branchdata-quadratic-reconstruction} 的恒等式 \(256y(u)+1091=108u(25-21u)\) 可见右端关于 \(u\mapsto\frac{25}{21}-u\) 对称，故 \(y(\iota(u))=y(u)\)。
将 \(P_{\mathrm{LY}}(y)=256y^3+411y^2+165y+32\) 代入 \(y=y(u)\) 得到 \(\QQ[u]\) 中的六次多项式。
由 \(u\in K\) 满足 \(g(u)=0\) 时 \(y(u)\) 为 Lee--Yang 分歧值（定理 \ref{thm:xi-terminal-zm-leyang-elliptic-structure}），可知 \(P_{\mathrm{LY}}(y(u))\) 被 \(g(u)\) 整除；对称性进一步给出其亦被 \(g(\iota(u))\) 整除。
两者均为三次且互素（由 \(g\) 的不可约性与 \(\iota\) 非平凡性），故
\(
P_{\mathrm{LY}}(y(u))=c\,g(u)g(\iota(u))
\)
对某 \(c\in\QQ^\times\) 成立；比较首项系数得到 \(c=\frac{3^{4}7^{3}}{2^{14}}\)。
第二式将 \(y(u)\) 与 \(y(u)-1\) 的分母清除并与第一式合并，再比较首项系数得到所示常数与二次因子。
\end{proof}

\begin{conclusion}[谱四次在实轴上的分岔型：实根数目仅在 \texorpdfstring{$y_{\mathrm{LY}},0,1$}{yLY,0,1} 处变化]\label{con:xi-terminal-zm-pi-real-root-bifurcation}
令
\[
\Pi(\lambda,y)=\lambda^4-\lambda^3-(2y+1)\lambda^2+\lambda+y(y+1),
\qquad
P_{\mathrm{LY}}(y)=256y^3+411y^2+165y+32,
\]
并令 \(y_{\mathrm{LY}}\in\RR\) 为 \(P_{\mathrm{LY}}(y)=0\) 的唯一实根（结论 \ref{con:xi-terminal-zm-discriminant-leyang}）。
则对任意实参数 \(y\in\RR\)，四次方程 \(\Pi(\lambda,y)=0\) 的实根数目满足如下分岔型（在开区间内根均互异）：
\begin{enumerate}
  \item 若 \(y<y_{\mathrm{LY}}\)，则 \(\Pi(\lambda,y)\) 在 \(\RR\) 上无实根；
  \item 若 \(y_{\mathrm{LY}}<y<0\)，则 \(\Pi(\lambda,y)\) 在 \(\RR\) 上恰有两条实根分支；
  \item 若 \(0<y<1\)，则 \(\Pi(\lambda,y)\) 在 \(\RR\) 上恰有四条实根分支；
  \item 若 \(y>1\)，则 \(\Pi(\lambda,y)\) 在 \(\RR\) 上恰有两条实根分支。
\end{enumerate}
并且在 \(y=y_{\mathrm{LY}},0,1\) 处分别出现一次二重根退化（对应结论 \ref{con:xi-terminal-zm-discriminant-leyang} 的分歧值集合）。
\end{conclusion}
\begin{proof}[证明要点]
由结论 \ref{con:xi-terminal-zm-discriminant-leyang}，
\(\mathrm{Disc}_{\lambda}(\Pi)=-y(y-1)P_{\mathrm{LY}}(y)\)，
故当 \(y\notin\{0,1\}\cup\{P_{\mathrm{LY}}(y)=0\}\) 时，四根皆为单根，实根数目在实轴上随 \(y\) 作局部常数变化。
另一方面，定理 \ref{thm:xi-terminal-zm-leyang-elliptic-structure} 给出 \(\mathrm{Disc}(P_{\mathrm{LY}})<0\)，从而 \(P_{\mathrm{LY}}\) 仅有一个实根 \(y_{\mathrm{LY}}\)；
故实轴上的分岔点严格为 \(y\in\{y_{\mathrm{LY}},0,1\}\)。
对每个开区间 \((-\infty,y_{\mathrm{LY}})\)、\((y_{\mathrm{LY}},0)\)、\((0,1)\)、\((1,\infty)\) 取代表点并用 Sturm 计数确定实根数目，即得到所示四段分岔型；在分岔点处由判别式为零得到二重根退化。
可复现核验脚本：\texttt{scripts/exp\_fold\_zm\_pi\_real\_bifurcation\_asymptotics\_audit.py}。证毕。
\end{proof}

\begin{conclusion}[Lee--Yang 实分歧点的 Cardano 根式表示：\texorpdfstring{$\sqrt{37}$}{sqrt(37)} 的判别式强制]\label{con:xi-terminal-zm-leyang-real-branchpoint-cardano}
沿用结论 \ref{con:xi-terminal-zm-leyang-lambda-cubic} 的记号，令
\[
g(\lambda):=16\lambda^3-9\lambda^2+1,
\qquad
y(\lambda):=\frac{4\lambda^3-3\lambda^2-2\lambda+1}{4\lambda}
=\lambda^2-\frac34\lambda-\frac12+\frac{1}{4\lambda}.
\]
设 \(\lambda_{\mathrm{LY}}\in\RR\) 为 \(g(\lambda)=0\) 的唯一实根，并令 \(y_{\mathrm{LY}}:=y(\lambda_{\mathrm{LY}})\)。
则 \(y_{\mathrm{LY}}\) 为 \(P_{\mathrm{LY}}(y)=0\) 的唯一实根（结论 \ref{con:xi-terminal-zm-discriminant-leyang}），并且在取实立方根分支的约定下存在根式闭式
\[
A:=\sqrt[3]{-101+16\sqrt{37}},\qquad B:=\sqrt[3]{-101-16\sqrt{37}},
\qquad
\boxed{\;\lambda_{\mathrm{LY}}=\frac{3+A+B}{16}\;}
\]
从而
\[
\boxed{\;y_{\mathrm{LY}}=\lambda_{\mathrm{LY}}^2-\frac34\lambda_{\mathrm{LY}}-\frac12+\frac{1}{4\lambda_{\mathrm{LY}}}\;}
\]
并且在 Tschirnhaus 平移 \(\lambda=t+\tfrac{3}{16}\) 下，
\[
g(\lambda)=0\ \Longleftrightarrow\ t^3-\frac{27}{256}\,t+\frac{101}{2048}=0,
\]
其 Cardano 判别式精确为
\[
\boxed{\;\Delta=\Bigl(\frac{101}{4096}\Bigr)^2+\Bigl(-\frac{9}{256}\Bigr)^3=\frac{37}{2^{16}}\;}
\]
从而 \(\sqrt{37}\) 在上述根式表达中由判别式机制强制出现。
\end{conclusion}
\begin{proof}[证明要点]
对 \(g(\lambda)=16\lambda^3-9\lambda^2+1\) 作平移 \(\lambda=t+\tfrac{3}{16}\) 得
\[
g(\lambda)=\frac{1}{128}\,\bigl(2048t^3-216t+101\bigr),
\]
从而等价于压低三次
\(
t^3-\frac{27}{256}\,t+\frac{101}{2048}=0
\)。
由 Cardano 公式，其唯一实根可写为
\[
t=\sqrt[3]{-\frac{101}{4096}+\frac{\sqrt{37}}{256}}+\sqrt[3]{-\frac{101}{4096}-\frac{\sqrt{37}}{256}}
=\frac{A+B}{16},
\]
于是 \(\lambda_{\mathrm{LY}}=t+\tfrac{3}{16}=\tfrac{3+A+B}{16}\)。
判别式 \(\Delta\) 由 \(p=-\tfrac{27}{256}\)、\(q=\tfrac{101}{2048}\) 代入
\(
\Delta=(q/2)^2+(p/3)^3
\)
得到，且等于 \(37/2^{16}>0\)，因此该三次仅有一个实根且 Cardano 根式必经 \(\sqrt{\Delta}=\sqrt{37}/2^{8}\)。
最后由结论 \ref{con:xi-terminal-zm-leyang-lambda-cubic} 的分歧参数化，非平凡分歧值满足 \(y=y(\lambda)\) 且 \(P_{\mathrm{LY}}(y)=0\)，从而 \(y_{\mathrm{LY}}=y(\lambda_{\mathrm{LY}})\) 为 \(P_{\mathrm{LY}}\) 的唯一实根。
可复现核验脚本：\texttt{scripts/exp\_fold\_zm\_pi\_real\_bifurcation\_asymptotics\_audit.py}。证毕。
\end{proof}

\begin{conclusion}[无穷远处的正实谱支分数幂渐近：\texorpdfstring{$y\to+\infty$}{y->+infty} 的 Puiseux jet]\label{con:xi-terminal-zm-pi-positive-real-branch-puiseux-infty}
设 \(y>1\)，并令 \(\lambda_{\pm}(y)\) 为 \(\Pi(\lambda,y)=0\) 在 \(\RR\) 上的两条实根分支，按大小约定
\[
0<\lambda_{-}(y)<\lambda_{+}(y).
\]
则当 \(y\to +\infty\) 沿实轴趋于无穷时，存在分数幂渐近展开
\[
\boxed{\;
\lambda_{\pm}(y)
= y^{1/2}\ \pm\ \frac{1}{2}y^{1/4}\ +\ \frac{1}{4}\ \pm\ \frac{7}{64}y^{-1/4}\ +\ \frac{37}{128}y^{-1/2}\ \mp\ \frac{729}{4096}y^{-3/4}\ +\ O(y^{-1})\;}
\]
特别地，
\[
\lambda_{+}(y)=y^{1/2}+\frac12 y^{1/4}+\frac14+O(y^{-1/4}).
\]
\end{conclusion}
\begin{proof}[证明要点（Puiseux 于无穷远）]
令 \(t=y^{-1/4}\)（故 \(y=t^{-4}\)）。
对根分支作 Puiseux 型设定
\[
\lambda(t)=t^{-2}+c_{1}t^{-1}+c_{0}+c_{-1}t+c_{-2}t^{2}+c_{-3}t^{3}+O(t^{4}),
\]
代入方程 \(\Pi(\lambda,y)=0\)，并乘以 \(t^{8}\) 以消去负幂，逐阶比较 \(t\)-系数得
\[
c_{1}=\pm\frac12,\qquad c_{0}=\frac14,\qquad c_{-1}=\pm\frac{7}{64},\qquad c_{-2}=\frac{37}{128},\qquad c_{-3}=\mp\frac{729}{4096},
\]
其中同号（分别异号）对应于 \(\lambda_{+}\)（分别 \(\lambda_{-}\)）的两支选择。
把 \(t=y^{-1/4}\) 回代即得所示展开。
可复现核验脚本：\texttt{scripts/exp\_fold\_zm\_pi\_real\_bifurcation\_asymptotics\_audit.py}。证毕。
\end{proof}

\begin{conclusion}[极端权重下 Perron 分支的双尺度 Puiseux：\texorpdfstring{$y\to0^+$}{y->0+} 与 \texorpdfstring{$y\to+\infty$}{y->+infty}]\label{con:xi-terminal-zm-perron-puiseux-two-scale}
令 \(\lambda_{+}(y)\) 表示对 \(y>0\) 的 Perron 根分支（最大正实根），并由 \(\lambda_{+}(1)=2\) 固定其与热力学主支的连通性（结论 \ref{con:xi-terminal-zm-lln-clt}）。
则 \(\lambda_{+}\) 在两端点呈现不可约分数幂尺度：
\begin{enumerate}
  \item 当 \(y\to 0^{+}\) 时，
  \[
  \lambda_{+}(y)
  =1+\frac{1}{\sqrt2}y^{1/2}+\frac58 y-\frac{43}{64\sqrt2}y^{3/2}-\frac1{16}y^{2}+O(y^{5/2})
  \qquad(\text{见结论 \ref{con:xi-terminal-zm-branch-puiseux-uniformization}}).
  \]

  \item 当 \(y\to+\infty\) 时，令 \(t=y^{-1/4}\)，则
  \[
  \lambda_{+}(y)
  =y^{1/2}+\frac12 y^{1/4}+\frac14+\frac{7}{64}y^{-1/4}+\frac{37}{128}y^{-1/2}-\frac{729}{4096}y^{-3/4}+O(y^{-1})
  \qquad(\text{见结论 \ref{con:xi-terminal-zm-pi-positive-real-branch-puiseux-infty}}).
  \]
\end{enumerate}
由此，定义密度函数
\[
\alpha(y):=\frac{y\,\lambda_{+}'(y)}{\lambda_{+}(y)}
=\left.\frac{\dd}{\dd s}\psi(s)\right|_{s=\log y},
\qquad
\psi(s):=\log\lambda_{+}(e^{s})-\log 2
\quad(\text{结论 \ref{con:xi-terminal-zm-lln-clt}}),
\]
则端点渐近为
\[
\alpha(y)=\frac{\sqrt2}{4}\,y^{1/2}+O(y)\quad(y\to0^{+}),
\qquad
\alpha(y)=\frac12-\frac18\,y^{-1/4}+O(y^{-1/2})\quad(y\to+\infty).
\]
特别地，\(\alpha(y)\) 以 \(y^{-1/4}\) 的速率逼近饱和值 \(1/2\)，该 \(1/4\) 指数由无穷远处谱方程的二重切触诱导，而非通常的线性扰动尺度。
\end{conclusion}



\begin{conclusion}[推前分布下的重量大数律与中心极限定律：均值与方差为有理常数]\label{con:xi-terminal-zm-lln-clt}
令 \(w_m(x)=d_m(x)/2^m\) 为均匀微态推前分布，并令 \(X\sim w_m\)。
记 \(H_m:=|X|_1\)，则对一切 \(y\in\CC\) 有概率母函数恒等式
\[
\EE\!\left[y^{H_m}\right]=\frac{Z_m(y)}{2^m}.
\]
令 \(\lambda_+(y)\) 表示 \(\Pi(\lambda,y)=0\) 在 \(y=1\) 的邻域内、满足 \(\lambda_+(1)=2\) 的解析特征支，则存在极限对数母函数
\[
\lim_{m\to\infty}\frac{1}{m}\log \EE\!\left(e^{sH_m}\right)
=\log\lambda_+(e^s)-\log 2.
\]
由此推出：
\begin{enumerate}
  \item \(m^{-1}H_m\to 5/18\)（依概率收敛）；
  \item \(\frac{H_m-\frac{5}{18}m}{\sqrt{m}}\Rightarrow \mathcal N\!\bigl(0,\frac{67}{972}\bigr)\)。
\end{enumerate}
上述常数由 \(\Pi(\lambda,y)\) 在点 \((\lambda,y)=(2,1)\) 的隐式微分唯一确定，并不依赖任何外部输入或额外归一化选择。
\end{conclusion}
\begin{proof}[证明要点（隐式微分）]
由定义
\(
\EE[y^{H_m}]
=\sum_{x\in X_m}y^{|x|_1}w_m(x)
=2^{-m}Z_m(y)
\)。
又由于 \(\Pi(\lambda,y)\) 在 \((2,1)\) 处满足 \(\partial_\lambda\Pi(2,1)\neq 0\)，
由隐函数定理得 \(\lambda_+(y)\) 在 \(y=1\) 的邻域内解析。
记 \(y=e^s\) 并令 \(\psi(s)=\log\lambda_+(e^s)-\log 2\)，则 \(\psi'(0)\) 与 \(\psi''(0)\) 分别给出线性漂移与方差率。
\par
对 \(\Pi(\lambda_+(y),y)=0\) 作一次与二次隐式微分，在点 \((\lambda,y)=(2,1)\) 代入可得
\[
\lambda_+'(1)=\frac{5}{9},
\qquad
\lambda_+''(1)=-\frac{64}{243}.
\]
于是
\(
\psi'(0)=\lambda_+'(1)/\lambda_+(1)=5/18
\)，并且
\[
\psi''(0)=\frac{\lambda_+'(1)}{\lambda_+(1)}+\frac{\lambda_+''(1)}{\lambda_+(1)}-\frac{\lambda_+'(1)^2}{\lambda_+(1)^2}
=\frac{67}{972}.
\]
大数律与中心极限定律由有限状态加权转移矩阵的解析谱扰动与标准准幂型定理推出。证毕。
\end{proof}

\begin{conclusion}[推前重量的一阶矩闭式：线性主项、常数漂移与交替相位]\label{con:xi-terminal-zm-m1-closed}
沿用结论 \ref{con:xi-terminal-zm-lln-clt} 的记号，并定义一阶矩分子
\[
M_1(m):=\sum_{x\in X_m}|x|_1\,d_m(x)=Z_m'(1).
\]
则对一切整数 \(m\ge 1\) 有严格恒等式
\[
\boxed{\;
M_1(m)=\Bigl(\frac{5}{18}m-\frac{1}{27}\Bigr)2^m+\frac14-\Bigl(\frac{m}{18}+\frac{23}{108}\Bigr)(-1)^m\;}
\]
从而均值具有精确展开
\[
\boxed{\;
\EE[H_m]=\frac{M_1(m)}{2^m}
=\frac{5}{18}m-\frac{1}{27}
+2^{-m}\Bigl(\frac14-\Bigl(\frac{m}{18}+\frac{23}{108}\Bigr)(-1)^m\Bigr)\;}
\]
其中有限尺寸偏离同时包含一个非零常数漂移与携带严格交替相位的 \(2^{-m}\) 层修正。
\end{conclusion}
\begin{proof}[证明要点（对 \(Z_m(y)\) 递推作一次微分）]
由结论 \ref{con:xi-terminal-zm-order4-rational}，\(Z_m(y)\) 满足四阶递推。
对该递推关于 \(y\) 微分并在 \(y=1\) 处取值，结合恒等式 \(Z_m(1)=2^m\)，得到 \(M_1(m)=Z_m'(1)\) 的常系数递推；
其特征多项式由 \(\Pi(\lambda,1)=(\lambda-2)(\lambda-1)(\lambda+1)^2\) 决定，从而解必可写为 \(2^m\) 通道与 \(\pm 1\) 通道的有限阶 Jordan 叠加。
用小 \(m\) 初值定出系数即可得到所示闭式；可由脚本 \texttt{scripts/exp\_fold\_zm\_bivariate\_partition\_audit.py} 复核。证毕。
\end{proof}

\begin{conclusion}[推前重量的二阶矩与方差闭式：线性涨落系数、常数项与双尺度交替修正]\label{con:xi-terminal-zm-m2-var-closed}
沿用结论 \ref{con:xi-terminal-zm-lln-clt} 的记号，并定义二阶矩分子
\[
M_2(m):=\sum_{x\in X_m}|x|_1^2\,d_m(x).
\]
则对一切整数 \(m\ge 1\) 有严格恒等式
\[
\boxed{\;
M_2(m)=\Bigl(\frac{25}{324}m^2+\frac{47}{972}m+\frac{113}{729}\Bigr)2^m+\frac{m}{8}
+\Bigl(\frac{m^3}{324}-\frac{m^2}{54}-\frac{31m}{216}-\frac{113}{729}\Bigr)(-1)^m\;}
\]
并因此得到方差的全量闭式分解
\[
\Var(H_m)=\frac{67}{972}m+\frac{112}{729}
+2^{-m}\Bigl(\frac{1}{54}-\frac{m}{72}+\Bigl(\frac{m^3}{324}+\frac{m^2}{81}-\frac{19m}{648}-\frac{83}{486}\Bigr)(-1)^m\Bigr)
\]
\[
\qquad\qquad
-2^{-2m}\Bigl(\frac{1}{16}-\Bigl(\frac{m}{36}+\frac{23}{216}\Bigr)(-1)^m+\Bigl(\frac{m}{18}+\frac{23}{108}\Bigr)^2\Bigr),
\]
其中主导涨落项为 \(\frac{67}{972}m+\frac{112}{729}\)，而全部有限尺寸偏离被压缩为两条严格指数尺度 \(2^{-m}\) 与 \(2^{-2m}\) 并携带交替相位。
\end{conclusion}
\begin{proof}[证明要点（对 \(Z_m(y)\) 递推作二次微分）]
注意 \(M_2(m)=Z_m''(1)+Z_m'(1)\)。
对结论 \ref{con:xi-terminal-zm-order4-rational} 的递推作二次微分并在 \(y=1\) 处取值，即得到 \(Z_m''(1)\) 的常系数递推；
再结合结论 \ref{con:xi-terminal-zm-m1-closed} 给出的 \(Z_m'(1)\) 闭式，可推出 \(M_2(m)\) 与 \(\Var(H_m)\) 的所示表达。
上述恒等式可由脚本 \texttt{scripts/exp\_fold\_zm\_bivariate\_partition\_audit.py} 复核。证毕。
\end{proof}

\begin{conclusion}[任意阶原始矩的湮灭律：\texorpdfstring{$(E-2)^{r+1}(E-1)^r(E+1)^{2r}$}{(E-2)^{r+1}(E-1)^r(E+1)^{2r}}]\label{con:xi-terminal-zm-mr-annihilator}
对整数 \(r\ge 0\)，定义原始矩分子
\[
M_r(m):=\sum_{x\in X_m}|x|_1^r\,d_m(x)\in\ZZ\qquad(m\ge 0).
\]
则序列 \(m\mapsto M_r(m)\) 满足常系数齐次递推，其最小湮灭多项式整除
\[
(E-2)^{r+1}(E-1)^r(E+1)^{2r},
\]
其中 \(E\) 为右移算子 \((Ea)(m)=a(m+1)\)。
等价地，存在多项式 \(P_r,Q_r,R_r\in\QQ[m]\) 使得对所有 \(m\ge 0\)
\[
M_r(m)=2^m P_r(m)+Q_r(m)+(-1)^m R_r(m),
\]
并且次数界为
\[
\deg P_r=r,\qquad \deg Q_r\le r-1,\qquad \deg R_r\le 2r-1.
\]
该结构把全部有限尺寸复杂性压缩为三条谱通道 \(2,1,-1\) 的有限阶 Jordan 叠加，从而任意阶矩在 \(m\) 上不仅存在极限律，而且具有全阶有限闭式展开。
\end{conclusion}
\begin{proof}[证明要点（谱通道与有限维扰动）]
由结论 \ref{con:xi-terminal-zm-order4-rational} 的“有限维带权转移矩阵”表示，
\(\{Z_m(y)\}\) 可写为固定维数矩阵 \(T(y)\) 的矩阵元；于是 \(M_r(m)\) 是该表示在 \(y=1\) 的有限阶导数所产生的线性组合。
在 \(y=1\) 时主多项式分解 \(\Pi(\lambda,1)=(\lambda-2)(\lambda-1)(\lambda+1)^2\) 给出三条谱通道；
有限阶导数仅提升相应 Jordan 阶数，从而得到所示湮灭多项式整除关系与三通道分解。
对 \(r\le 3\) 的直接代入核验与符号回归可由脚本 \texttt{scripts/exp\_fold\_zm\_bivariate\_partition\_audit.py} 完成。证毕。
\end{proof}

\begin{conclusion}[单位长度累积量的有理累积塔：所有极限累积量皆属 \texorpdfstring{$\QQ$}{Q}]\label{con:xi-terminal-zm-cumulant-rational}
令 \(\kappa_r(H_m)\) 为 \(H_m\) 的第 \(r\) 阶累积量（cumulant）。
则极限
\[
\kappa_r^\infty:=\lim_{m\to\infty}\frac{1}{m}\kappa_r(H_m)
\]
对每个整数 \(r\ge 1\) 存在，且满足 \(\kappa_r^\infty\in\QQ\)。
其生成机制可取为：令 \(\lambda_+(y)\) 表示特征方程 \(\Pi(\lambda,y)=0\) 在 \(y=1\) 的邻域内满足 \(\lambda_+(1)=2\) 的解析特征支，
并定义
\[
\psi(s):=\log\lambda_+(e^s)-\log 2.
\]
则 \(\psi\) 在 \(s=0\) 的泰勒系数全为有理数，且
\(
\kappa_r^\infty=\psi^{(r)}(0)
\)。
前四阶显式为
\[
\kappa_1^\infty=\frac{5}{18},\qquad
\kappa_2^\infty=\frac{67}{972},\qquad
\kappa_3^\infty=-\frac{22}{2187},\qquad
\kappa_4^\infty=-\frac{1763}{157464}.
\]
因此标准化极限偏度与超额峰度分别为
\[
\gamma_1=\frac{\kappa_3^\infty}{(\kappa_2^\infty)^{3/2}}\approx -0.5558544555,\qquad
\gamma_2=\frac{\kappa_4^\infty}{(\kappa_2^\infty)^2}\approx -2.3564268211.
\]
\end{conclusion}
\begin{proof}[证明要点（隐函数定理与系数域闭合）]
由结论 \ref{con:xi-terminal-zm-lln-clt}，极限对数矩母函数为 \(\psi(s)\)。
由于 \(\Pi(\lambda,y)\in\ZZ[\lambda,y]\) 且 \(\partial_\lambda\Pi(2,1)\neq 0\)，隐函数定理给出 \(\lambda_+(y)\) 在 \(y=1\) 邻域的解析展开；
逐阶比较系数可知其系数落在 \(\QQ\) 中，进而 \(\psi(s)=\log\lambda_+(e^s)-\log 2\) 的泰勒系数同属 \(\QQ\)。
累积量密度 \(\kappa_r^\infty\) 等于 \(\psi^{(r)}(0)\) 为标准结论；显式分数可由符号微分回收，
并可由脚本 \texttt{scripts/exp\_fold\_zm\_bivariate\_partition\_audit.py} 复核。证毕。
\end{proof}

\begin{conclusion}[三阶累积量的有限尺寸闭式：仿射主项与三层指数修正]\label{con:xi-terminal-zm-kappa3-finite-closed}
令 \(\kappa_3(H_m)\) 为 \(H_m\) 的第三阶累积量，则对一切整数 \(m\ge 1\) 有精确分解
\[
\kappa_3(H_m)= -\frac{22}{2187}m+\frac{52}{19683}
+2^{-m}\mathcal P_1\!\bigl(m,(-1)^m\bigr)
+2^{-2m}\mathcal P_2\!\bigl(m,(-1)^m\bigr)
+2^{-3m}\mathcal P_3\!\bigl(m,(-1)^m\bigr),
\]
其中每个 \(\mathcal P_j(\cdot,\cdot)\) 都可写成“偶/奇”二扇区可分离的形式
\(
\mathcal P_j(m,\varepsilon)=A_j(m)+\varepsilon B_j(m)
\)
（\(A_j,B_j\) 为有理系数多项式）。
三层修正可取为
\[
\mathcal P_1(m,\varepsilon)=\Bigl(\frac{m^2}{1728}+\frac{151m}{1728}-\frac{293}{10368}\Bigr)
+\varepsilon\Bigl(-\frac{m^5}{12960}+\frac{7m^4}{15552}+\frac{41m^3}{11664}-\frac{487m^2}{46656}-\frac{2699m}{77760}+\frac{2407}{279936}\Bigr),
\]
\[
\mathcal P_2(m,\varepsilon)=\Bigl(\frac{m^4}{1944}+\frac{47m^3}{11664}+\frac{35m^2}{11664}-\frac{143m}{3888}-\frac{269}{2187}\Bigr)
+\varepsilon\Bigl(-\frac{m^3}{432}-\frac{5m^2}{432}+\frac{7m}{432}+\frac{34}{243}\Bigr),
\]
\[
\mathcal P_3(m,\varepsilon)=\Bigl(\frac{m^2}{216}+\frac{23m}{648}+\frac{193}{1944}\Bigr)
+\varepsilon\Bigl(-\frac{m^3}{2916}-\frac{23m^2}{5832}-\frac{629m}{17496}-\frac{15617}{157464}\Bigr).
\]
因此三阶非高斯性在单位长度尺度上呈严格线性斜率 \(-22/2187\)，且其全部有限尺寸偏离被压缩为三层指数尺度 \(2^{-m},2^{-2m},2^{-3m}\) 的代数可审计修正。
\end{conclusion}
\begin{proof}[证明要点（由闭式矩回收累积量并分层收集）]
由结论 \ref{con:xi-terminal-zm-m1-closed} 与结论 \ref{con:xi-terminal-zm-m2-var-closed} 得到 \(\EE[H_m]\) 与 \(\EE[H_m^2]\) 的闭式；
并由结论 \ref{con:xi-terminal-zm-mr-annihilator} 在 \(r=3\) 的特例得到 \(\EE[H_m^3]\) 的三通道闭式。
将三者代入
\(\kappa_3(H_m)=\EE[H_m^3]-3\EE[H_m^2]\EE[H_m]+2\EE[H_m]^3\)，
并把 \(2^{-m}\) 的幂次与 \((-1)^m\) 的相位逐层收集，即得到所示分解与多项式系数。
恒等式可由脚本 \texttt{scripts/exp\_fold\_zm\_bivariate\_partition\_audit.py} 复核。证毕。
\end{proof}

\begin{conclusion}[重量分布的全局大偏差：速率函数为四次代数谱的 Legendre 变换]\label{con:xi-terminal-zm-ldp}
定义压力函数
\[
\psi(s):=\log\lambda_+(e^s)-\log 2.
\]
则 \(H_m/m\) 满足大偏差原理，其速率函数为
\[
I(\alpha)=\sup_{s\in\RR}\bigl(s\alpha-\psi(s)\bigr).
\]
其中 \(\psi\) 在实轴上严格凸且实解析；
其复解析延拓的最近实轴障碍由结论 \ref{con:xi-terminal-zm-discriminant-leyang} 的 \(y_{\mathrm{LY}}\) 所决定（经 \(y=e^s\) 的复对数支提升）。
因此重量统计的多分形谱可在不引入任何额外热力学形式体系的前提下，直接由 \(\Pi(\lambda,y)\) 的主特征支给出。
\end{conclusion}
\begin{proof}[证明要点]
由结论 \ref{con:xi-terminal-zm-lln-clt}，\(\psi\) 给出极限对数矩母函数。
在 \(\psi\) 的可微与严格凸条件下，G\"artner--Ellis 定理推出大偏差原理及其 Legendre 变换形式。
最近解析障碍由 \(\lambda_+(y)\) 的分歧点决定，而分歧点由判别式为零刻画，见结论 \ref{con:xi-terminal-zm-discriminant-leyang}。证毕。
\end{proof}

\begin{remark}[主特征支的椭圆一致化：\texorpdfstring{$\lambda_+(y)$}{lambda+(y)} 的 Weierstrass \texorpdfstring{$\wp$}{wp} 表示]\label{rem:xi-terminal-zm-lambda-plus-wp-uniformization}
在结论 \ref{con:xi-terminal-zm-elliptic-normalization} 的椭圆模型
\[
E:\ Y^2=X^3-X+\frac14
\]
上，复解析一致化给出存在周期格 \(\Lambda\subset\CC\) 与 Weierstrass 函数 \(\wp_\Lambda\) 使得
\[
E(\CC)\cong \CC/\Lambda,\qquad
X=\wp_\Lambda(z),\qquad
Y=\frac{\wp_\Lambda'(z)}{2},
\]
并满足微分方程
\[
\bigl(\wp_\Lambda'(z)\bigr)^2=4\wp_\Lambda(z)^3-4\wp_\Lambda(z)+1
\qquad(g_2=4,\ g_3=-1).
\]
在该一致化下，物理主支重量函数（结论 \ref{con:xi-terminal-zm-y-quadratic-extension-minpoly}）
\[
y=X^2+Y-\frac12
\]
提升为椭圆函数
\[
y(z)=\wp_\Lambda(z)^2+\frac{\wp_\Lambda'(z)-1}{2}.
\]
令 \(\lambda(z):=X(z)=\wp_\Lambda(z)\)。
由结论 \ref{con:xi-terminal-zm-elliptic-normalization} 的双有理互逆 \(\nu:E\to\mathcal C\)，
对任意 \(z\in\CC/\Lambda\) 有 \((\lambda(z),y(z))\in\mathcal C(\CC)\)，即
\[
\Pi\bigl(\lambda(z),y(z)\bigr)=0.
\]
反之，\(\Pi(\lambda,y)=0\) 的正规化在复解析意义下同构于 \(E(\CC)\cong\CC/\Lambda\)，
因此其任一解析支均可由上述椭圆函数参数化统一给出。
因此 \(\pi_y\) 的有限分歧点在一致化参数 \(z\) 下等价于临界点条件 \(\frac{\dd}{\dd z}y(z)=0\)；
而 \(y=\infty\) 对应于 \(z\to 0\) 的四阶极点结构，与 \((y)_\infty=4O\) 的全分歧极点描述一致。
令 \(P_1=(2,-\tfrac52)\in E(\QQ)\)（满足 \(y(P_1)=1\)，见结论 \ref{con:xi-terminal-zm-elliptic-rational-points-denominator-law}），并取 \(z_1\in\CC/\Lambda\) 使
\(
P_1=(\wp_\Lambda(z_1),\wp_\Lambda'(z_1)/2)
\)。
由于 \(P_1\) 不在覆盖 \(\pi_y:E\to\mathbb P^1_y\) 的分歧点上（结论 \ref{con:xi-terminal-zm-elliptic-y-hurwitz-passport}），
\(y(z)\) 在 \(z=z_1\) 的邻域内局部可逆；令 \(z=z(y)\) 表示其局部反函数，则主特征支可写为
\[
\lambda_+(y)=X\bigl(z(y)\bigr)=\wp_\Lambda\bigl(z(y)\bigr),
\qquad \lambda_+(1)=2.
\]
因此 \(\lambda_+(y)\) 在避开分歧值集合的任一路径上均可解析延拓，并可通过在 \(\CC/\Lambda\) 上延拓 \(z(y)\) 后回代 \(\wp_\Lambda\) 进行全局表示。
从而 \(\lambda_+(y)\)（以及与之伴随的 \(\log\lambda_+(y)\) 的多值解析结构）可被等价视为固定椭圆格上的反演与周期问题，而非一般代数函数分支的偶然叠加。
进一步，由结论 \ref{con:xi-terminal-zm-leyang-ramification-critical-values} 的微分恒等式
\(
\dd y=(4XY+3X^2-1)\,\omega
\)
（其中 \(\omega=\dd X/(2Y)\)）得：在任意去分歧区域内取满足 \(\pi_y(P(y))=y\) 的局部解析截面 \(P(y)\in E(\CC)\)，令
\(
X(y)=X(P(y))
\)、\(
Y(y)=Y(P(y))
\)，则
\[
\boxed{\;
\frac{\dd}{\dd y}\log \lambda_+(y)
=\frac{1}{X(y)}\frac{\dd X(y)}{\dd y}
=\frac{2Y(y)}{X(y)\,\bigl(4X(y)Y(y)+3X(y)^2-1\bigr)}\; }.
\]
等价地，在最小整模型 \(E_0:\eta^2+\eta=x^3-x\) 上令 \(x=\lambda\)、\(\eta=y-\lambda^2\)，则同一恒等式可写为
\[
\frac{\dd}{\dd y}\log \lambda(y)
=\frac{2\eta+1}{x\bigl(2x(2\eta+1)+3x^2-1\bigr)}.
\]
因此 \(\bigl(\dd(\log\lambda_+)/\dd y\bigr)\dd y=\dd X/X\) 为椭圆曲线上的第三类阿贝尔微分，其周期决定 \(\log\lambda_+\) 的多值解析延拓。
结合注记 \ref{rem:xi-terminal-elliptic-37a1-modular-anchor} 中 \(E\) 作为 \(X_0^+(37)\) 模型的外部锚定，上述周期亦可在 level \(37\) 的模形式周期口径下理解。
其分歧集合与局部型由 \(\pi_y\) 的分歧结构完全确定：有限分歧值为
\(
y\in\{0,1\}\cup\{P_{\mathrm{LY}}(y)=0\}
\)，
且在每个有限分歧值处呈平方根型 Puiseux 分歧（结论 \ref{con:xi-terminal-zm-branch-puiseux-uniformization}）；
而在 \(y=\infty\) 处由 \((y)_\infty=4O\) 给出四重全分歧（结论 \ref{con:xi-terminal-zm-elliptic-y-hurwitz-passport}）。
特别地，实轴上的最近分歧值 \(y_{\mathrm{LY}}\) 及其伴随退化特征值 \(\lambda_{\mathrm{LY}}\) 分别由
\(
P_{\mathrm{LY}}(y_{\mathrm{LY}})=0
\)
与
\(
16\lambda_{\mathrm{LY}}^{3}-9\lambda_{\mathrm{LY}}^{2}+1=0
\)
代数刻画（结论 \ref{con:xi-terminal-zm-leyang-lambda-cubic}）。
\end{remark}

\begin{conclusion}[大偏差速率函数的代数参数化与端点对数奇性]\label{con:xi-terminal-zm-ldp-rate-algebraic-param-endpoint-log}
沿用结论 \ref{con:xi-terminal-zm-ldp} 的压力函数
\(
\psi(s)=\log\lambda_{+}(e^{s})-\log2
\)
及其 Legendre--Fenchel 对偶 \(I(\alpha)=\sup_{s\in\RR}(s\alpha-\psi(s))\)。
对任意 \(y>0\)，令 \(\lambda=\lambda_{+}(y)\)。
定义
\[
\alpha(y):=\frac{y\,\lambda_{+}'(y)}{\lambda_{+}(y)}\in(0,\tfrac12).
\]
则由谱方程 \(\Pi(\lambda,y)=0\) 的隐式微分得到显式有理表达式
\[
\alpha(y)
=-\frac{y\,\partial_y\Pi(\lambda,y)}{\lambda\,\partial_{\lambda}\Pi(\lambda,y)}
=\frac{y\bigl(2\lambda^{2}-(2y+1)\bigr)}{\lambda\bigl(4\lambda^{3}-3\lambda^{2}-2(2y+1)\lambda+1\bigr)}
\Big|_{\lambda=\lambda_{+}(y)}.
\]
并且沿该参数支有
\[
I\bigl(\alpha(y)\bigr)=\alpha(y)\log y-\log\lambda_{+}(y)+\log2.
\]
在本体系中，\(\alpha:(0,\infty)\to(0,\tfrac12)\) 为严格递增的实解析双射，因此 \(I\) 在 \((0,\tfrac12)\) 上实解析且严格凸。
其端点行为由 Perron 分支的分数幂尺度锁定（结论 \ref{con:xi-terminal-zm-perron-puiseux-two-scale}）：
\begin{itemize}
  \item 当 \(\alpha\downarrow0\) 时，
  \[
  I(\alpha)
  =\log2+2\alpha\log\alpha+(\log8-2)\alpha+O\!\bigl(\alpha^{2}\log(1/\alpha)\bigr).
  \]
  \item 当 \(\alpha\uparrow\tfrac12\) 时，令 \(\delta:=\tfrac12-\alpha\downarrow0\)，则
  \[
  I\!\left(\tfrac12-\delta\right)
  =\log2+4\delta\log\delta+(12\log2-4)\delta+O\!\bigl(\delta^{2}\log(1/\delta)\bigr).
  \]
\end{itemize}
其中 \(\alpha\log\alpha\) 与 \(\delta\log\delta\) 两个非解析项分别由 \(y\to0^{+}\) 的 \(y^{1/2}\) Puiseux 与 \(y\to+\infty\) 的 \(y^{-1/4}\) Puiseux 诱导，因而构成该 Lee--Yang 权重体系在两端点的可检验奇性指纹。
\end{conclusion}

\begin{conclusion}[复 Lee--Yang 主导的解析半径与高阶累积量振荡律]\label{con:xi-terminal-zm-psi-analytic-radius-cumulant-oscillation}
考虑 \(\psi(s)=\log\lambda_{+}(e^{s})-\log2\) 在 \(s=0\) 的解析芽（结论 \ref{con:xi-terminal-zm-lln-clt}）。
由结论 \ref{con:xi-terminal-zm-discriminant-leyang} 与结论 \ref{con:xi-terminal-zm-leyang-cubic-branch-image}，
\(y\)-平面的有限分歧值除 \(\{0,1\}\) 外还包含 Lee--Yang 三次
\(
P_{\mathrm{LY}}(y)=256y^{3}+411y^{2}+165y+32
\)
的共轭复根对
\[
y_{\mathrm c}^{\pm}\approx-0.2355083735\pm0.2339255521\,i.
\]
令
\[
s_{\mathrm c}^{\pm}:=\log y_{\mathrm c}^{\pm}\quad(\text{取主值}),
\qquad
\rho:=|s_{\mathrm c}^{+}|\approx2.6045558274,
\qquad
\theta:=\arg(s_{\mathrm c}^{+})\approx2.0080026785,
\]
则 \(\psi\) 在圆盘 \(|s|<\rho\) 内单值全纯，并且该圆盘为以 \(0\) 为中心的最大解析圆盘；其边界上首先遇到的奇点即 \(s_{\mathrm c}^{\pm}\) 所对应的代数分歧。
\par
记 \(\kappa_n^\infty:=\psi^{(n)}(0)\in\QQ\) 为单位长度的第 \(n\) 阶累积量密度（结论 \ref{con:xi-terminal-zm-cumulant-rational}）。
则当 \(n\to\infty\) 时存在常数 \(B\neq 0\) 与相位 \(\varphi\in\RR\)，使得
\[
\kappa_n^\infty
=\frac{n!}{\sqrt\pi}\,|B|\,\rho^{-n+\frac12}\,n^{-\frac32}
\cos\!\Bigl((n-\tfrac12)\theta+\varphi\Bigr)
+o\!\bigl(n!\rho^{-n}n^{-\frac32}\bigr).
\]
更精确地，取临界对 \((\lambda_{\mathrm c},y_{\mathrm c})\) 满足
\[
\Pi(\lambda_{\mathrm c},y_{\mathrm c})=0,\qquad
\partial_{\lambda}\Pi(\lambda_{\mathrm c},y_{\mathrm c})=0,
\qquad (y_{\mathrm c}=y_{\mathrm c}^{+}),
\]
并令
\[
c:=\sqrt{-\,\frac{2\,\partial_{y}\Pi(\lambda_{\mathrm c},y_{\mathrm c})}{\partial_{\lambda\lambda}\Pi(\lambda_{\mathrm c},y_{\mathrm c})}},
\qquad
B:=\frac{c\sqrt{y_{\mathrm c}}}{\lambda_{\mathrm c}},
\]
则上述振荡律成立；数值上可取
\[
\lambda_{\mathrm c}\approx0.4179674053-0.2321140339\,i,
\qquad
|B|\approx0.9634999692.
\]
因此高阶累积量密度除呈现 \(\rho^{-n}\) 的指数控制与 \(n^{-3/2}\) 的普适幂律修正外，还携带由 \(\theta=\arg(s_{\mathrm c}^{+})\) 决定的确定性振荡频率，给出一条可直接数值核验的“复 Lee--Yang 振荡律”。
可复现核验脚本：\texttt{scripts/exp\_fold\_zm\_bivariate\_partition\_audit.py}。
\end{conclusion}

\begin{theorem}[Lee--Yang 主振荡角的 Gelfond--Schneider 超越性]\label{thm:xi-terminal-zm-leyang-oscillation-angle-transcendence}
沿用结论 \ref{con:xi-terminal-zm-psi-analytic-radius-cumulant-oscillation} 的记号，令
\(
s_{\mathrm c}:=s_{\mathrm c}^{+}=\log y_{\mathrm c}^{+}
\)
为主值对数，并写 \(\rho:=|s_{\mathrm c}|\)、\(\theta:=\arg(s_{\mathrm c})\in(0,\pi)\)。
则：
\begin{enumerate}
  \item \(\tan\theta\) 为超越数；
  \item 特别地，\(\theta/\pi\notin\QQ\)。
\end{enumerate}
\end{theorem}
\begin{proof}
设
\[
N:=y_{\mathrm c}^{+}\overline{y_{\mathrm c}^{+}}=|y_{\mathrm c}^{+}|^{2}\in\overline{\QQ}\cap\RR_{>0},
\qquad
r:=\frac{y_{\mathrm c}^{+}}{\overline{y_{\mathrm c}^{+}}}\in\overline{\QQ},
\qquad |r|=1.
\]
由于 \(y_{\mathrm c}^{+}\notin\RR_{\le 0}\)，主值对数与共轭相容，故
\[
\log N=s_{\mathrm c}+\overline{s_{\mathrm c}}=2\Re(s_{\mathrm c})\neq 0,
\qquad
s_{\mathrm c}-\overline{s_{\mathrm c}}=2i\,\Im(s_{\mathrm c})
\]
并且 \(\exp(s_{\mathrm c}-\overline{s_{\mathrm c}})=r\neq 1\)。
令
\[
\beta:=\frac{s_{\mathrm c}-\overline{s_{\mathrm c}}}{s_{\mathrm c}+\overline{s_{\mathrm c}}}
=i\,\frac{\Im(s_{\mathrm c})}{\Re(s_{\mathrm c})}=i\tan\theta.
\]
若 \(\tan\theta\) 为代数数，则 \(\beta\in\overline{\QQ}\setminus\QQ\)（因为 \(\beta\notin\RR\)）。
取 \(\alpha:=N\in\overline{\QQ}\setminus\{0,1\}\)，由 Gelfond--Schneider 定理可知任一值 \(\alpha^{\beta}\) 皆为超越数。\cite{Baker1975TranscendentalNumberTheory}
另一方面，由 \(\beta\log N=s_{\mathrm c}-\overline{s_{\mathrm c}}\) 得
\[
N^{\beta}=\exp(\beta\log N)=\exp(s_{\mathrm c}-\overline{s_{\mathrm c}})=r\in\overline{\QQ},
\]
矛盾。
故 \(\tan\theta\) 不可能为代数数，即 \(\tan\theta\) 为超越数。
\par
若 \(\theta/\pi\in\QQ\)，则 \(\zeta:=e^{2i\theta}\) 为单位根，从而
\(
\tan\theta=\frac{\zeta-1}{i(\zeta+1)}
\in\overline{\QQ}
\)，
与已证超越性矛盾。
故 \(\theta/\pi\notin\QQ\)。证毕。
\end{proof}

\begin{theorem}[归一化高阶累积量密度的反正弦极限律]\label{thm:xi-terminal-zm-cumulant-arcsine-limit-law}
\leanverified{paper\_xi\_terminal\_zm\_cumulant\_arcsine\_limit\_law}
沿用结论 \ref{con:xi-terminal-zm-psi-analytic-radius-cumulant-oscillation} 的振荡渐近：
存在 \(B\neq 0\)、\(\varphi\in\RR\) 使
\[
\kappa_n^\infty
=\frac{n!}{\sqrt\pi}\,|B|\,\rho^{-n+\frac12}\,n^{-\frac32}
\cos\!\Bigl((n-\tfrac12)\theta+\varphi\Bigr)
+o\!\bigl(n!\rho^{-n}n^{-\frac32}\bigr).
\]
定义归一化序列
\[
A_n:=\frac{\sqrt\pi}{|B|}\,\rho^{\,n-\frac12}\,n^{\frac32}\,\frac{\kappa_n^\infty}{n!}.
\]
则 \(A_n=\cos\!\bigl((n-\tfrac12)\theta+\varphi\bigr)+o(1)\)，并且由定理 \ref{thm:xi-terminal-zm-leyang-oscillation-angle-transcendence} 的 \(\theta/\pi\notin\QQ\) 得：\(A_n\) 的经验分布弱收敛到 \([-1,1]\) 上的反正弦分布。
具体地，对任意连续函数 \(f\in C([-1,1])\)，有
\[
\lim_{N\to\infty}\frac1N\sum_{n=1}^{N} f(A_n)
=\frac{1}{\pi}\int_{-1}^{1}\frac{f(x)}{\sqrt{1-x^2}}\,\dd x.
\]
\end{theorem}
\begin{proof}
由渐近式直接得到 \(A_n-\cos(t_n)\to 0\)，其中 \(t_n:=(n-\tfrac12)\theta+\varphi\)。
由 \(\theta/\pi\notin\QQ\) 可知旋转序列 \(t_n\ (\mathrm{mod}\ 2\pi)\) 在 \(\RR/2\pi\ZZ\) 上等分布（Kronecker--Weyl），故对任意连续 \(g\in C(\RR/2\pi\ZZ)\) 有
\[
\frac1N\sum_{n=1}^{N} g(t_n)\to \frac{1}{2\pi}\int_{0}^{2\pi} g(t)\,\dd t.
\]
取 \(g(t)=f(\cos t)\) 即得 \(\cos(t_n)\) 的经验分布为反正弦律的推前测度。\cite{KuipersNiederreiter1974}
再由 \(f\) 的一致连续性与 \(A_n-\cos(t_n)\to 0\) 控制 Ces\`aro 平均差，得到所述极限式。证毕。
\end{proof}

\begin{corollary}[符号密度、极值充满与无限次符号翻转]\label{cor:xi-terminal-zm-cumulant-sign-density-extremes}
\leanverified{paper\_xi\_terminal\_zm\_cumulant\_sign\_density\_extremes}
沿用定理 \ref{thm:xi-terminal-zm-cumulant-arcsine-limit-law} 的记号。
则集合 \(\{n:\kappa_n^\infty>0\}\) 与 \(\{n:\kappa_n^\infty<0\}\) 的自然密度均为 \(1/2\)，并且
\[
\limsup_{n\to\infty}A_n=1,\qquad \liminf_{n\to\infty}A_n=-1.
\]
从而序列 \((\kappa_n^\infty)\) 发生无限次符号翻转，且不可能最终周期。
\end{corollary}
\begin{proof}
由反正弦分布关于 \(0\) 对称且 \(\cos t\) 在一半圆周取正、一半取负，
结合 \(A_n-\cos(t_n)\to 0\) 得符号密度结论。
反正弦分布的支撑为 \([-1,1]\) 且任意开邻域均有正测度，故 \(\cos(t_n)\) 在 \([-1,1]\) 稠密；
再由 \(A_n-\cos(t_n)\to 0\) 得 \(\limsup A_n=1\)、\(\liminf A_n=-1\)。
最后，无限次正负出现推出无限次符号翻转；若 \((\kappa_n^\infty)\) 最终周期，则其主振荡相位 \((n-\tfrac12)\theta+\varphi\) 只能对应有理旋转，从而 \(\theta/\pi\in\QQ\)，与定理 \ref{thm:xi-terminal-zm-leyang-oscillation-angle-transcendence} 矛盾。证毕。
\end{proof}

\begin{proposition}[相邻阶局部增长率在扩展实轴上的稠密聚点]\label{prop:xi-terminal-zm-cumulant-local-growth-dense}
\leanverified{paper\_xi\_terminal\_zm\_cumulant\_local\_growth\_dense}
沿用结论 \ref{con:xi-terminal-zm-psi-analytic-radius-cumulant-oscillation} 的振荡渐近，并在 \(\kappa_n^\infty\neq 0\) 时定义
\[
R_n:=\rho\left(\frac{n}{n+1}\right)^{\frac32}\cdot\frac{\kappa_{n+1}^\infty}{(n+1)\kappa_n^\infty}.
\]
则 \(R_n\) 在扩展实轴 \(\widehat{\RR}:=\RR\cup\{\infty\}\) 上具有稠密的聚点集合：对任意 \(L\in\RR\)，存在无穷子列 \(n_k\to\infty\) 使 \(R_{n_k}\to L\)，并且亦存在无穷子列使 \(|R_{n_k}|\to\infty\)。
\end{proposition}
\begin{proof}
由结论 \ref{con:xi-terminal-zm-psi-analytic-radius-cumulant-oscillation}，
\[
\frac{\kappa_{n+1}^\infty}{(n+1)\kappa_n^\infty}
=\frac{1}{\rho}\left(\frac{n}{n+1}\right)^{\frac32}
\frac{\cos(t_n+\theta)}{\cos(t_n)}+o(1),
\qquad
t_n:=(n-\tfrac12)\theta+\varphi.
\]
故
\[
R_n=\frac{\cos(t_n+\theta)}{\cos(t_n)}+o(1)
=\cos\theta-\sin\theta\,\tan(t_n)+o(1),
\]
其中 \(\sin\theta\neq 0\)（因 \(\theta\in(0,\pi)\)）。
由 \(\theta/\pi\notin\QQ\) 得 \(t_n\ (\mathrm{mod}\ \pi)\) 等分布，从而 \(\tan(t_n)\) 在 \(\widehat{\RR}\) 上稠密；
再由仿射同胚 \(x\mapsto \cos\theta-\sin\theta\,x\) 与 \(o(1)\) 微扰的稳定性得到 \(R_n\) 的聚点稠密性，并可取 \(|\tan(t_n)|\to\infty\) 的子列推出 \(|R_n|\to\infty\)。证毕。
\end{proof}

\begin{theorem}[Stokes 幅值平方 \texorpdfstring{$B^2$}{B^2} 的三次最小多项式、判别式与 \texorpdfstring{$S_3$}{S3} 刚性]\label{thm:xi-terminal-zm-stokes-amplitude-square-cubic-rigidity}
沿用结论 \ref{con:xi-terminal-zm-psi-analytic-radius-cumulant-oscillation} 的临界对 \((\lambda_{\mathrm c},y_{\mathrm c})\) 与幅值常数 \(B\)，并令
\[
b:=B^{2}=\frac{c^{2}y_{\mathrm c}}{\lambda_{\mathrm c}^{2}}\in\CC,
\qquad t:=b.
\]
设 \(K:=\QQ(y_{\mathrm c})\) 为 Lee--Yang 三次域（\(P_{\mathrm{LY}}(y_{\mathrm c})=0\)）。
则在 \(K\) 内成立显式有理恒等式
\begin{equation}\label{eq:xi-terminal-zm-stokes-t-rational-in-yc}
\boxed{\;
t
=-\frac{1}{6}\cdot
\frac{440y_{\mathrm c}^{2}+47y_{\mathrm c}-64}{46y_{\mathrm c}^{2}+40y_{\mathrm c}-5}\;}
\end{equation}
并且 \(t\) 在 \(\QQ\) 上满足三次整式
\begin{equation}\label{eq:xi-terminal-zm-stokes-t-minpoly}
\boxed{\;
162t^{3}+1593t^{2}+1998t+1184=0\;}
\end{equation}
等价地，\(b\) 满足
\[
162b^{3}+1593b^{2}+1998b+1184=0.
\]
该三次在 \(\QQ\) 上不可约，且判别式为
\[
\mathrm{Disc}(162b^{3}+1593b^{2}+1998b+1184)
=-2^{2}\cdot 3^{9}\cdot 11^{2}\cdot 37\cdot 109^{2}.
\]
因判别式非平方，其分裂域伽罗瓦群同构于 \(S_3\)。
从而 \(\QQ(b)=\QQ(t)=\QQ(y_{\mathrm c})\)，亦即 \(B^2\) 本身为 \(K\) 的一个原始生成元。
\end{theorem}
\begin{proof}[证明要点]
由结论 \ref{con:xi-terminal-zm-branch-double-root-linear-closure}，在约束 \(P_{\mathrm{LY}}(y_{\mathrm c})=0\) 下二重根谱支满足线性刚性
\(
279\lambda_{\mathrm c}+512y_{\mathrm c}^{2}+518y_{\mathrm c}+5=0
\)，
故 \(\lambda_{\mathrm c}\in \QQ(y_{\mathrm c})\)。
另一方面，
\(
\partial_y\Pi(\lambda,y)=-2\lambda^2+2y+1
\)、
\(
\partial_{\lambda\lambda}\Pi(\lambda,y)=12\lambda^2-6\lambda-4y-2
\)，
从而
\(
t=-2\cdot\frac{y_{\mathrm c}\,\partial_y\Pi(\lambda_{\mathrm c},y_{\mathrm c})}{\lambda_{\mathrm c}^{2}\,\partial_{\lambda\lambda}\Pi(\lambda_{\mathrm c},y_{\mathrm c})}
\in \QQ(\lambda_{\mathrm c},y_{\mathrm c})
=\QQ(y_{\mathrm c})
\)。
将 \(\lambda_{\mathrm c}\) 以 \(y_{\mathrm c}\) 的有理式表示并在三次商环 \(\QQ[y]/(P_{\mathrm{LY}})\) 中约化，即得 \eqref{eq:xi-terminal-zm-stokes-t-rational-in-yc}。
再将 \eqref{eq:xi-terminal-zm-stokes-t-rational-in-yc} 改写为关于 \(y_{\mathrm c}\) 的二次方程并与 \(P_{\mathrm{LY}}(y_{\mathrm c})=0\) 联立作消元，即得到 \eqref{eq:xi-terminal-zm-stokes-t-minpoly}。
有理根检验排除一次因子，故该三次在 \(\QQ\) 上不可约。
其判别式按三次判别式公式计算即为
\(
-2^{2}\cdot 3^{9}\cdot 11^{2}\cdot 37\cdot 109^{2}
\)，非平方，故分裂域伽罗瓦群为 \(S_3\)。
由不可约性可知 \([\QQ(t):\QQ]=3\)，且 \(t\in\QQ(y_{\mathrm c})\) 已证，遂 \(\QQ(t)=\QQ(y_{\mathrm c})\)；
由 \(b=t\) 得 \(\QQ(b)=\QQ(t)\)。证毕。
\end{proof}

\begin{corollary}[幅值三次的二次分辨子域：\texorpdfstring{$\QQ(\sqrt{-111})$}{Q(sqrt(-111))}]\label{cor:xi-terminal-zm-stokes-b2-quadratic-resolvent}
沿用定理 \ref{thm:xi-terminal-zm-stokes-amplitude-square-cubic-rigidity} 的记号，令 \(K=\QQ(b)\) 并令 \(L\) 为 \(K/\QQ\) 的伽罗瓦闭包。
则 \(L/\QQ\) 的唯一二次子域为
\[
\boxed{\ \QQ(\sqrt{-111})\ }.
\]
特别地，该二次子域与 Lee--Yang 三次的二次分辨域一致（见结论 \ref{con:xi-terminal-zm-leyang-cubic-arithmetic}），从而 \(B^2\) 与 \(y_{\mathrm c}\) 的任一原始生成元给出同一条二次判别脊柱。
\end{corollary}
\begin{proof}
由定理 \ref{thm:xi-terminal-zm-stokes-amplitude-square-cubic-rigidity}，\(b\) 的极小多项式为不可约三次 \(f(t)=162t^{3}+1593t^{2}+1998t+1184\)，且
\[
\mathrm{Disc}(f)=-2^{2}\cdot 3^{9}\cdot 11^{2}\cdot 37\cdot 109^{2}.
\]
判别式的平方自由部分为 \(-3\cdot 37=-111\)。
对不可约三次且 \(\mathrm{Disc}(f)\) 非平方的情形，其分裂域的伽罗瓦群为 \(S_3\)，并且其唯一二次子域由 \(\QQ(\sqrt{\mathrm{Disc}(f)})\) 给出；
因此 \(L\) 的唯一二次子域为 \(\QQ(\sqrt{-111})\)。证毕。
\end{proof}

\begin{conjecture}[Stokes 幅值平方的 CM--模参数化]\label{conj:xi-terminal-zm-stokes-b2-cm-modular-parameterization}
沿用定理 \ref{thm:xi-terminal-zm-stokes-amplitude-square-cubic-rigidity} 的记号，令 \(b:=B^2\) 并记三次域 \(K:=\QQ(b)\)。
猜想存在某一层级 \(N\) 的模曲线上的模函数 \(f\)（例如 \(X_0(N)\) 的 Hauptmodul 或 \(\lambda\)-函数的有理变换）与某一 CM 点 \(\tau\in\mathbb H\)，使得
\[
b=f(\tau),
\qquad
\QQ(b)\ \text{为相应环类域的某一三次子域}.
\]
并且 \(N\) 的素因子集合可由 \(\mathrm{Disc}(162b^{3}+1593b^{2}+1998b+1184)\) 的素因子集合控制。
该猜想可通过对候选 \(N\) 与 CM 判别式的有限枚举检验：比较 \(f(\tau)\) 的极小多项式与 \eqref{eq:xi-terminal-zm-stokes-t-minpoly}。
\end{conjecture}

\input{conj__xi-terminal-zm-stokes-thirdkind-exponential-period}

\begin{theorem}[振幅平方 \texorpdfstring{$B^2$}{B^2} 与 \texorpdfstring{$\kappa_*^2$}{kappa*^2} 的显式有理互构]\label{thm:xi-terminal-zm-b2-kappa2-rational-interconstruction}
\leanverified{paper\_xi\_terminal\_zm\_b2\_kappa2\_rational\_interconstruction}
令
\[
u:=\kappa_*^2,
\qquad
b:=B^2.
\]
其中 \(u\) 取自定理 \ref{thm:xi-terminal-zm-kappa-square-cubic-field-s3}，\(b\) 取自定理 \ref{thm:xi-terminal-zm-stokes-amplitude-square-cubic-rigidity}。
则存在显式有理关系
\[
\boxed{\;
b
=-2\cdot
\frac{24708u^2-28628u+7733}{51084u^2-44412u+8735}\; }.
\]
因而
\[
\QQ(b)=\QQ(u)=\QQ(\lambda_*),
\]
即 \(B^2\) 与 \(\kappa_*^2\) 生成同一 Lee--Yang 三次数域。
\end{theorem}
\begin{proof}
由定理 \ref{thm:xi-terminal-zm-kappa-square-cubic-field-s3}，
\[
\lambda_*=\frac{3(2u-1)}{4(3u-2)}.
\]
由定理 \ref{thm:xi-terminal-zm-stokes-t-elliptic-branch-coordinate} 的表达式 \eqref{eq:xi-terminal-zm-stokes-t-rational-in-alpha}（取 \(\alpha=\lambda_*\) 且 \(t=b\)）：
\[
b=\frac{2(1113\lambda_*^2-2192\lambda_*+303)}{3(1831\lambda_*^2+144\lambda_*-175)}.
\]
把 \(\lambda_*=\frac{3(2u-1)}{4(3u-2)}\) 代入并在 \(\QQ(u)\) 中约化，即得所示有理式，故 \(b\in\QQ(u)\)。
另一方面，定理 \ref{thm:xi-terminal-zm-kappa-square-cubic-field-s3} 与定理 \ref{thm:xi-terminal-zm-stokes-amplitude-square-cubic-rigidity} 分别给出
\(
[\QQ(u):\QQ]=[\QQ(b):\QQ]=3
\)。
故包含关系 \(\QQ(b)\subseteq\QQ(u)\) 只能是等号；再结合
\(
\QQ(u)=\QQ(\lambda_*)
\)
即得结论。
\end{proof}

\begin{theorem}[幅值平方的椭圆分歧坐标表达]\label{thm:xi-terminal-zm-stokes-t-elliptic-branch-coordinate}
沿用定理 \ref{thm:xi-terminal-zm-stokes-amplitude-square-cubic-rigidity} 的记号，并令 \(\alpha:=\lambda_{\mathrm c}\)。
则 \(\alpha\) 满足 Lee--Yang 三次谱分歧方程 \(16\alpha^{3}-9\alpha^{2}+1=0\)（结论 \ref{con:xi-terminal-zm-leyang-ramification-critical-values}）。
并且在分歧约束下有恒等式
\[
4\alpha\,y_{\mathrm c}=4\alpha^{3}-3\alpha^{2}-2\alpha+1
\qquad(\text{结论 \ref{con:xi-terminal-zm-leyang-ramification-critical-values}}),
\]
从而 \(t=B^{2}\) 可写为 \(\QQ(\alpha)\) 中的有理函数：
\begin{equation}\label{eq:xi-terminal-zm-stokes-t-rational-in-alpha}
\boxed{\;
t
=\frac{2\,(1113\alpha^{2}-2192\alpha+303)}{3\,(1831\alpha^{2}+144\alpha-175)}\; }.
\end{equation}
\end{theorem}
\begin{proof}[证明要点]
由结论 \ref{con:xi-terminal-zm-leyang-ramification-critical-values} 的分歧恒等式可将 \(y_{\mathrm c}\) 表为 \(\alpha\) 的有理式。
将其代入定理 \ref{thm:xi-terminal-zm-stokes-amplitude-square-cubic-rigidity} 的表达式 \eqref{eq:xi-terminal-zm-stokes-t-rational-in-yc}，
并在商环 \(\QQ[\alpha]/(16\alpha^{3}-9\alpha^{2}+1)\) 中约化，即得 \eqref{eq:xi-terminal-zm-stokes-t-rational-in-alpha}。证毕。
\end{proof}

\begin{theorem}[隐藏平方与纯有理反演：由 \texorpdfstring{$t=B^{2}$}{t=B^2} 重构 \texorpdfstring{$y_{\mathrm c}$}{yc}]\label{thm:xi-terminal-zm-stokes-hidden-square-rational-inversion}
令 \(t\) 满足 \eqref{eq:xi-terminal-zm-stokes-t-minpoly}，并在三次域 \(K:=\QQ(t)\) 中定义
\begin{equation}\label{eq:xi-terminal-zm-stokes-s-def}
s:=\frac{-71045-13947t+12258t^{2}}{1199}.
\end{equation}
则在 \(K\) 内成立隐藏平方恒等式
\begin{equation}\label{eq:xi-terminal-zm-stokes-hidden-square}
\boxed{\;
s^{2}=10080t^{2}+16224t+12761\; }.
\end{equation}
并且由 \eqref{eq:xi-terminal-zm-stokes-t-rational-in-yc} 导出的二次方程在 \(K\) 内分解，其两根可写为纯有理表达
\begin{equation}\label{eq:xi-terminal-zm-stokes-yc-rational-in-t}
\boxed{\;
y^{(1)}(t)
=\frac{36774t^{2}-329601t-269488}{9592(69t+110)},
\qquad
y^{(2)}(t)
=\frac{-36774t^{2}-245919t+156782}{9592(69t+110)}\; }.
\end{equation}
其中与 \(P_{\mathrm{LY}}(y)=0\) 相容者即为 Lee--Yang 临界值 \(y_{\mathrm c}\)；在结论 \ref{con:xi-terminal-zm-psi-analytic-radius-cumulant-oscillation} 取 \(y_{\mathrm c}=y_{\mathrm c}^{+}\) 的嵌入下，对应 \(y_{\mathrm c}=y^{(1)}(t)\)。
\end{theorem}
\begin{proof}[证明要点]
由 \eqref{eq:xi-terminal-zm-stokes-t-rational-in-yc} 可得 \(y_{\mathrm c}\) 满足关于 \(y\) 的二次方程
\[
(276t+440)y^{2}+(240t+47)y-(30t+64)=0.
\]
其判别式为
\(
\Delta=9(10080t^{2}+16224t+12761)
\)。
式 \eqref{eq:xi-terminal-zm-stokes-hidden-square} 表明该判别式平方根已属于 \(K\)，故该二次方程在 \(K\) 内分解并给出两条根的纯有理表达；
化简即得 \eqref{eq:xi-terminal-zm-stokes-yc-rational-in-t}。证毕。
\end{proof}

\begin{proposition}[临界谱值的幅值重构]\label{prop:xi-terminal-zm-stokes-lambda-critical-rational-in-t}
在定理 \ref{thm:xi-terminal-zm-stokes-hidden-square-rational-inversion} 的设定下，
\(\lambda_{\mathrm c}\) 可由 \(t\) 显式重构为
\begin{equation}\label{eq:xi-terminal-zm-stokes-lambda-critical-rational-in-t}
\boxed{\;
\lambda_{\mathrm c}
=-\frac{525411t^{2}-4582128t-2625484}{1199(69t+110)^{2}}\; }.
\end{equation}
\end{proposition}
\begin{proof}
由结论 \ref{con:xi-terminal-zm-branch-double-root-linear-closure} 的线性刚性关系，
\(
\lambda_{\mathrm c}=-(512y_{\mathrm c}^{2}+518y_{\mathrm c}+5)/279
\)。
将 \eqref{eq:xi-terminal-zm-stokes-yc-rational-in-t} 中与 \(P_{\mathrm{LY}}(y)=0\) 相容的根 \(y_{\mathrm c}\) 代入并在 \(\QQ(t)\) 中约化即可得到 \eqref{eq:xi-terminal-zm-stokes-lambda-critical-rational-in-t}。证毕。
\end{proof}

\begin{theorem}[Stokes 幅值的六次代数度与极小多项式]\label{thm:xi-terminal-zm-stokes-amplitude-degree6-minpoly}
设 \(B\) 为结论 \ref{con:xi-terminal-zm-psi-analytic-radius-cumulant-oscillation} 中由临界对 \((\lambda_{\mathrm c},y_{\mathrm c})\) 所定义的 Stokes 幅值常数，并令 \(t=B^{2}\)。
则 \([\QQ(B):\QQ]=6\)，且 \(B\) 的极小多项式为
\begin{equation}\label{eq:xi-terminal-zm-stokes-B-minpoly}
\boxed{\;
162B^{6}+1593B^{4}+1998B^{2}+1184=0\; }.
\end{equation}
\end{theorem}
\begin{proof}
由 \eqref{eq:xi-terminal-zm-stokes-t-minpoly}，\(t\) 在 \(\QQ\) 上次数为 \(3\)，故 \(K:=\QQ(t)\) 为三次域且 \(K\subset \QQ(B)\)。
若 \(B\in K\)，则 \(t=B^{2}\) 在 \(K\) 中为平方，从而其范数 \(\mathrm{N}_{K/\QQ}(t)\) 必为有理数平方，因而非负。
另一方面，把 \eqref{eq:xi-terminal-zm-stokes-t-minpoly} 除以 \(162\) 化为首一多项式，可得
\[
\mathrm{N}_{K/\QQ}(t)=(-1)^{3}\cdot\frac{1184}{162}=-\frac{592}{81}<0,
\]
矛盾。
故 \(B\notin K\)，从而 \([\QQ(B):K]=2\) 并推出 \([\QQ(B):\QQ]=6\)。
式 \eqref{eq:xi-terminal-zm-stokes-B-minpoly} 由 \eqref{eq:xi-terminal-zm-stokes-t-minpoly} 代入 \(t=B^{2}\) 直接得到；
其次数为 \(6\)，与 \([\QQ(B):\QQ]=6\) 一致，故为极小多项式。证毕。
\end{proof}

\begin{conclusion}[全阶累积量密度的有理性与可审计生成方程]\label{con:xi-terminal-zm-cumulant-dalgebraic-equation}
结论 \ref{con:xi-terminal-zm-cumulant-rational} 已给出 \(\psi(s)\) 在 \(s=0\) 的 Taylor 展开属于 \(\QQ[[s]]\)，因而 \(\kappa_n^\infty=\psi^{(n)}(0)\in\QQ\) 对一切 \(n\ge 1\) 成立。
进一步地，由 \(\Pi(\lambda,y)=0\) 在代换
\(
\lambda=2e^{\psi(s)},\ y=e^{s}
\)
下得到 \(\psi\) 的显式代数生成方程
\[
16e^{4\psi(s)}-8e^{3\psi(s)}-4(2e^{s}+1)e^{2\psi(s)}+2e^{\psi(s)}+e^{2s}+e^{s}=0,
\]
从而 \(\psi\) 为 \(D\)-代数函数；这为任意阶 \(\kappa_n^\infty\) 的符号递推与误差审计提供封闭的代数约束。
除结论 \ref{con:xi-terminal-zm-cumulant-rational} 已列出的低阶项外，高阶例如
\[
\kappa_{5}^\infty=\frac{15650}{1594323},\quad
\kappa_{6}^\infty=\frac{307723}{57395628},\quad
\kappa_{7}^\infty=-\frac{2332358}{129140163},\quad
\kappa_{8}^\infty=\frac{166894943}{55788550416}.
\]
可复现核验脚本：\texttt{scripts/exp\_fold\_zm\_bivariate\_partition\_audit.py}。
\end{conclusion}

\begin{conclusion}[Abel--Jacobi 坐标下的 Lee--Yang 微分闭式：\texorpdfstring{$y(u)$}{y(u)} 的一阶代数方程]\label{con:xi-terminal-zm-leyang-abel-jacobi-ode}
沿用结论 \ref{con:xi-terminal-zm-elliptic-normalization} 的椭圆模型
\[
E:\quad Y^2=X^3-X+\frac14,
\qquad
y:=X^2+Y-\frac12\in\QQ(E),
\qquad
\omega:=\frac{\dd X}{2Y}.
\]
取 Abel--Jacobi 坐标 \(u\) 使 \(\dd u=\omega\)，并令
\[
\delta(u):=\frac{\dd y}{\dd u}\ \in\ \QQ(E).
\]
则在去分歧区域内，\((y(u),\delta(u))\) 满足完全显式的一阶代数微分闭式
\[
\boxed{\;
\delta^{4}+(8y-3)\delta^{3}+(2y^{2}-34y-4)\delta^{2}
-y(y-1)\,P_{\mathrm{LY}}(y)=0\;},
\qquad
P_{\mathrm{LY}}(y):=256y^{3}+411y^{2}+165y+32.
\]
该四次关系在 \(\QQ(y)[T]\) 中不可约，从而给出函数域扩张的首一四次本原表示
\[
\QQ(E)=\QQ(y,\delta),
\qquad
[\QQ(E):\QQ(y)]=4,
\]
并且常数项精确等于 Lee--Yang 判别核（带符号）\(-y(y-1)P_{\mathrm{LY}}(y)\)。
因此常数项（即 \(\delta=0\) 的驻值谱）严格给出
\[
\delta(u)=0\quad\Longleftrightarrow\quad
y(u)\in \{0,1\}\ \cup\ \{P_{\mathrm{LY}}(y)=0\},
\]
从而 Lee--Yang 分歧像在 Abel--Jacobi 时间下作为驻值集合内禀出现，并与结论 \ref{con:xi-terminal-zm-discriminant-leyang} 的分歧判别核一致。
\end{conclusion}
\begin{proof}[证明要点（消元）]
由结论 \ref{con:xi-terminal-zm-leyang-ramification-critical-values} 的微分恒等式
\(
\dd y=(4XY+3X^2-1)\,\omega
\)
得
\[
\delta=4XY+3X^2-1.
\]
又由 \(y=X^2+Y-\tfrac12\) 得 \(Y=y-X^2+\tfrac12\)，代入上式得到关于 \(X\) 的三次约束
\[
4X^{3}-3X^{2}-(4y+2)X+(1+\delta)=0.
\]
另一方面，消去 \(Y\) 的平面四次模型由
\[
F(X,y):=X^{4}-X^{3}-(2y+1)X^{2}+X+y(y+1)=0
\]
给出（结论 \ref{con:xi-terminal-zm-leyang-discriminant-norm-identity}）。
对变量 \(X\) 作结果式消元 \(\mathrm{Res}_{X}(F,\ 4X^{3}-3X^{2}-(4y+2)X+(1+\delta))\)，并在 \(\ZZ[y,\delta]\) 中因式整理，即得到所示四次方程；
其常数项分解为 \(-y(y-1)P_{\mathrm{LY}}(y)\) 与分歧像刻画（结论 \ref{con:xi-terminal-zm-discriminant-leyang}）一致。
不可约性可由一次良性专化读出：例如取 \(y=2\)，所得四次多项式在 \(\QQ[T]\) 上不可约，从而原四次在 \(\QQ(y)[T]\) 中不可约。
可复现核验脚本：\texttt{scripts/exp\_fold\_zm\_elliptic\_leyang\_cover\_geometry\_audit.py}。证毕。
\end{proof}

\begin{conclusion}[判别式自反与二级退化五次：\texorpdfstring{$\mathrm{Disc}_T(F_\delta)$}{Disc\_T(F\_delta)} 的平方自由部仍为 Lee--Yang 核]\label{con:xi-terminal-zm-leyang-delta-quartic-discriminant-reflexive}
令
\[
F_\delta(T;y):=T^{4}+(8y-3)T^{3}+(2y^{2}-34y-4)T^{2}-y(y-1)\,P_{\mathrm{LY}}(y)\in\ZZ[y][T],
\qquad
P_{\mathrm{LY}}(y)=256y^{3}+411y^{2}+165y+32.
\]
则其关于 \(T\) 的判别式满足严格因式分解
\[
\boxed{\;
\mathrm{Disc}_{T}\bigl(F_\delta(T;y)\bigr)
=-y(y-1)\,P_{\mathrm{LY}}(y)\cdot Q_{5}(y)^{2}\;},
\]
其中
\[
Q_{5}(y)=4096y^{5}+5376y^{4}-464y^{3}-2749y^{2}-723y+80\in\ZZ[y].
\]
因此 \(\delta\)-生成表示的全部不可约平方自由障碍仍由同一 Lee--Yang 判别核 \(-y(y-1)P_{\mathrm{LY}}(y)\) 控制；
而 \(Q_{5}(y)=0\) 给出一组与 \(y(y-1)P_{\mathrm{LY}}(y)=0\) 不同的二级退化参数集合，其上四叶结构发生额外并合。
\end{conclusion}
\begin{proof}[证明要点]
直接计算 \(\mathrm{Disc}_{T}(F_\delta)\) 并在 \(\ZZ[y]\) 中因式分解得到所示恒等式。
可复现核验脚本：\texttt{scripts/exp\_fold\_zm\_elliptic\_leyang\_cover\_geometry\_audit.py}。证毕。
\end{proof}

\input{subsubsec__xi-terminal-zm-leyang-delta-closure-plane-quintic}
\input{cor__xi-terminal-zm-delta-node-tangent-parity-law}
\input{subsubsec__xi-terminal-zm-leyang-delta-branch-quintic-arithmetic}
\input{subsubsec__xi-terminal-zm-delta-node-order-discriminant-index}

\begin{conclusion}[二级退化五次的全局算术指纹：新的平方自由敏感素数 \texorpdfstring{$727$}{727}]\label{con:xi-terminal-zm-leyang-q5-discriminant-727}
令 \(Q_5(y)\) 如结论 \ref{con:xi-terminal-zm-leyang-delta-quartic-discriminant-reflexive}。
则其判别式满足严格素因子分解
\[
\boxed{\;
\mathrm{Disc}\bigl(Q_5\bigr)=-2^{28}\cdot 3^{9}\cdot 11^{2}\cdot 31^{3}\cdot 727\; }.
\]
因此除去结论 \ref{con:xi-terminal-zm-leyang-branchpolynomial-discriminant-modp-collisions} 所限定的一级分歧素集 \(\{2,3,31,37\}\) 之外，
二级退化集合还引入新的平方自由敏感素数 \(727\)（并且 \(11\) 仅以平方因子出现）；
在该素数特征下，\(Q_5(y)=0\) 的约化不可避免地产生额外的并合/不可分现象，从而提供一条与 Lee--Yang 一级核正交的更深层算术奇异源。
\end{conclusion}
\begin{proof}[证明要点]
对 \(Q_5\) 作判别式计算并在 \(\ZZ\) 中分解即可。
可复现核验脚本：\texttt{scripts/exp\_fold\_zm\_elliptic\_leyang\_cover\_geometry\_audit.py}。证毕。
\end{proof}

\begin{conclusion}[模性--微分闭合：\texorpdfstring{$D(\tau)=\frac{q\,\partial_q y(\tau)}{f(q)}$}{D(tau)=q y'/f} 满足同一四次方程]\label{con:xi-terminal-zm-leyang-modular-differential-closure}
沿用注记 \ref{rem:xi-terminal-elliptic-37a1-modular-anchor} 的模参数化：取 \(q=e^{2\pi i\tau}\)，
并令 \(f(\tau)=\sum_{n\ge 1}a_n q^n\) 为 level \(37\) 的规范化新形式，使得 \(f(q)\,\dd q/q\) 等于椭圆曲线 \(E\) 的 N\'eron 微分 \(\omega=\dd X/(2Y)\) 在同一参数化下的拉回；
令 \(y(\tau)\) 为该参数化下的重量函数。
定义
\[
D(\tau):=\frac{q\,\frac{\dd}{\dd q}y(\tau)}{f(q)}\in \CC((q)).
\]
则 \((y(\tau),D(\tau))\) 满足完全闭合的一阶代数微分方程
\[
\boxed{\;
D^{4}+(8y-3)D^{3}+(2y^{2}-34y-4)D^{2}
-y(y-1)\,P_{\mathrm{LY}}(y)=0\;},
\qquad
P_{\mathrm{LY}}(y)=256y^{3}+411y^{2}+165y+32,
\]
其中 \(y=y(\tau)\)、\(D=D(\tau)\)。
\end{conclusion}
\begin{proof}[证明要点]
由定义 \(D(\tau)=\frac{q\,(\dd/\dd q)y(\tau)}{f(q)}\) 与 \(\dd u=\omega=f(q)\,\dd q/q\) 得 \(D=\dd y/\dd u=\delta\)。
代入结论 \ref{con:xi-terminal-zm-leyang-abel-jacobi-ode} 的四次闭式即得。证毕。
\end{proof}

\begin{conclusion}[Lee--Yang Perron 分支的中心化三阶 Fuchs 线性算子]\label{con:xi-terminal-zm-leyang-perron-centered-fuchs-operator}
令
\[
\Pi(\lambda,y):=\lambda^{4}-\lambda^{3}-(2y+1)\lambda^{2}+\lambda+y(y+1),\qquad
P_{\mathrm{LY}}(y):=256y^{3}+411y^{2}+165y+32,
\]
并记 \(\lambda_{+}(y)\) 为 \(y=1\) 邻域内满足 \(\lambda_{+}(1)=2\) 的解析分支（结论 \ref{con:xi-terminal-zm-lln-clt}）。
定义中心化分支
\[
f(y):=\lambda_{+}(y)-\frac14.
\]
则 \(f\) 满足三阶线性微分方程
\[
A_{3}(y)f^{(3)}+A_{2}(y)f^{(2)}+A_{1}(y)f'+A_{0}(y)f=0,
\]
其中 \(A_{j}(y)\in\ZZ[y]\) 显式如下：
\input{sections/generated/eq_fold_zm_leyang_perron_fuchs_operator}
\end{conclusion}
\begin{proof}[证明要点]
设
\(
K:=\QQ(y)[\lambda]/(\Pi(\lambda,y))
\)。
由 Vieta 关系 \(\sum_{i=1}^{4}\lambda_i=1\) 得
\(
\mathrm{Tr}_{K/\QQ(y)}(\lambda-\tfrac14)=0
\)，故
\(
g:=\lambda-\tfrac14\in K_0:=\ker\mathrm{Tr}
\)。
又 \(\dim_{\QQ(y)}K_0=3\)。
令 \(\partial:=\frac{d}{dy}\)；
可分扩张下迹与导数可交换，
\(
\mathrm{Tr}(\partial u)=\partial\mathrm{Tr}(u)
\)，故 \(\partial(K_0)\subseteq K_0\)。
因此 \(g,\partial g,\partial^2 g,\partial^3 g\) 在 \(K_0\) 中线性相关，存在 \(a_j(y)\in\QQ(y)\) 使
\[
\partial^3 g+a_2(y)\partial^2 g+a_1(y)\partial g+a_0(y)g=0.
\]
将该恒等式在 \(K\) 中按基 \(\{1,\lambda,\lambda^2,\lambda^3\}\) 消元并清分母，得到整系数多项式 \(A_3,A_2,A_1,A_0\)。
脚本 \texttt{scripts/exp\_fold\_zm\_leyang\_perron\_fuchs\_operator\_audit.py} 对上述恒等式给出代数核验。证毕。
\end{proof}

\begin{conclusion}[Riemann 方案中的真奇点与表观奇点分离]\label{con:xi-terminal-zm-leyang-perron-fuchs-riemann-apparent-split}
设 \(L_3\) 为结论 \ref{con:xi-terminal-zm-leyang-perron-centered-fuchs-operator} 的三阶算子。
其有限奇点恰为 \(A_3(y)=0\)：
\[
\{0,1\}\cup\{P_{\mathrm{LY}}(y)=0\}\cup\{Q(y)=0\}.
\]
其中 \(\{0,1\}\cup\{P_{\mathrm{LY}}=0\}\) 为 \(\Pi\) 的分歧点集；\(\{Q=0\}\) 为五个有限表观奇点。
若 \(y_0\) 为 \(A_3\) 的简单零点，则局部指标为
\[
\operatorname{Exp}_{y_0}(L_3)=\left\{0,1,2-\frac{A_2(y_0)}{A_3'(y_0)}\right\}.
\]
并且
\[
\operatorname{Exp}_{y_0}(L_3)=\left\{0,\frac12,1\right\}
\quad\bigl(y_0\in\{0,1,P_{\mathrm{LY}}=0\}\bigr),
\]
\[
\operatorname{Exp}_{y_0}(L_3)=\{0,1,3\}
\quad\bigl(y_0\in\{Q=0\}\bigr).
\]
在 \(y=\infty\) 处，指标集合为
\[
\left\{-\frac14,\frac14,\frac12\right\},
\]
其指标多项式为
\[
32\sigma^{3}-16\sigma^{2}-2\sigma+1=(2\sigma-1)(4\sigma-1)(4\sigma+1).
\]
\end{conclusion}
\begin{proof}[证明要点]
将方程写成首一形式 \(f^{(3)}+p_2f^{(2)}+p_1f'+p_0f=0\)。
在 \(A_3\) 的简单零点处 \(p_2\) 仅有一阶极点、\(p_1,p_0\) 不产生更高阶极部，故为正则奇点，
并得到指标公式
\(
r(r-1)\bigl(r-2+p_{2,-1}\bigr)=0
\)；
即
\(
\{0,1,2-A_2(y_0)/A_3'(y_0)\}
\)。
由恒等式
\[
2A_2-3A_3'\equiv 0\pmod{y(y-1)P_{\mathrm{LY}}},\qquad
A_2+A_3'\equiv 0\pmod{Q}
\]
分别得到上述两组局部指数。
又 \(\gcd\!\bigl(Q,\ y(y-1)P_{\mathrm{LY}}\bigr)=1\)，故 \(Q\) 的零点不在 \(\Pi\) 的分歧集上；
四个代数分支在这些点局部单值解析，故该五点为表观奇点。
无穷远指标由主平衡项直接给出。证毕。
\end{proof}

\begin{conclusion}[Wronskian 的平方闭式与表观奇点结构方程]\label{con:xi-terminal-zm-leyang-perron-fuchs-wronskian-square}
设 \(V\) 为算子 \(L_3\) 的解空间，\(\dim V=3\)。
取任意基 \((u_1,u_2,u_3)\)，Wronskian 记为
\[
W(y):=\det\bigl(u_i^{(j-1)}(y)\bigr)_{1\le i,j\le 3}.
\]
则存在常数 \(C\neq 0\) 使
\[
W(y)^2
=
C\cdot\frac{Q(y)^2}{y^3(y-1)^3P_{\mathrm{LY}}(y)^3}.
\]
因此 \(Q(y)=0\) 给出 \(W\) 的有限零点，而 \(P_{\mathrm{LY}}(y)=0\) 与 \(y=0,1\) 给出 \(W\) 的有限极点。
\end{conclusion}
\begin{proof}[证明要点]
对首一三阶方程有经典恒等式
\(
\frac{W'}{W}=-p_2
\)。
从而
\(
\frac{(W^2)'}{W^2}=-2p_2=-2A_2/A_3
\)。
对
\(
R(y):=\frac{Q(y)^2}{y^3(y-1)^3P_{\mathrm{LY}}(y)^3}
\)
直接计算得
\(
\frac{R'}{R}=-2A_2/A_3
\)。
故 \(W^2=CR\)。
脚本 \texttt{scripts/exp\_fold\_zm\_leyang\_perron\_fuchs\_operator\_audit.py} 已给出符号核验。证毕。
\end{proof}

\begin{conclusion}[Perron 分支 \texorpdfstring{$\Lambda(t)$}{Lambda(t)} 的全局 \(p\)-整性与仅 \(3\)-进分母]\label{con:xi-terminal-zm-leyang-perron-global-pintegrality}
令 \(t:=y-1\)，并令 \(\Lambda(t)\in\QQ[[t]]\) 满足
\[
\Pi(\Lambda(t),1+t)=0,\qquad \Lambda(0)=2.
\]
则对任意素数 \(p\neq 3\)，有
\[
\Lambda(t)\in\ZZ_p[[t]].
\]
因此 \(\Lambda(t)=2+\sum_{n\ge 1}a_n t^n\) 的每个系数满足 \(a_n\in\ZZ[1/3]\)。
\end{conclusion}
\begin{proof}[证明要点]
设 \(F(\lambda,t):=\Pi(\lambda,1+t)\in\ZZ[\lambda,t]\)。
有 \(F(2,0)=0\) 且
\(
F_{\lambda}(2,0)=\partial_{\lambda}\Pi(2,1)=9
\)。
当 \(p\neq 3\) 时 \(9\in\ZZ_p^\times\)，由 \(p\)-进隐函数定理得到唯一
\(
\Lambda_p(t)\in\ZZ_p[[t]]
\)
满足 \(F(\Lambda_p(t),t)=0\)、\(\Lambda_p(0)=2\)。
而该初值下 \(\QQ[[t]]\) 形式解唯一，故 \(\Lambda=\Lambda_p\) 视为 \(\QQ_p[[t]]\) 元成立。
于是对所有 \(p\neq 3\) 的 \(p\)-整性同时成立，分母仅可能含素因子 \(3\)。证毕。
\end{proof}

\begin{theorem}[\texorpdfstring{$p=3$}{p=3} 的 \(27\)-尺度良约化与 \(3\)-进 Hensel 可解性]\label{thm:xi-terminal-zm-leyang-perron-p3-renormalization-hensel}
\leanverified{paper\_xi\_terminal\_zm\_leyang\_perron\_p3\_renormalization\_hensel}
在结论 \ref{con:xi-terminal-zm-leyang-perron-global-pintegrality} 的设定下，素数 \(p=3\) 为唯一使 \(F_\lambda(2,0)\notin\ZZ_p^\times\) 的例外。
令
\[
t=27u,\qquad \lambda=2+3v.
\]
则恒有分解
\[
\Pi(2+3v,1+27u)=27\,H(u,v),
\]
其中
\[
H(u,v)=27u^{2}-18uv^{2}-24uv-5u+3v^{4}+7v^{3}+5v^{2}+v\in\ZZ[u,v].
\]
并且
\[
\partial_v H(0,0)=1\in\ZZ_3^\times.
\]
因此存在唯一的 \(v(u)\in\ZZ_3[[u]]\) 满足
\[
H(u,v(u))=0,\qquad v(0)=0,
\]
从而
\[
\Lambda(27u)=2+3v(u)\in 2+3\,\ZZ_3[[u]].
\]
\end{theorem}
\begin{proof}
把 \(\lambda=2+3v\)、\(t=27u\) 代入 \(F(\lambda,t)=\Pi(\lambda,1+t)\) 并在 \(\ZZ[u,v]\) 中展开整理，即得所示恒等式与多项式 \(H(u,v)\)。
对 \(H\) 求偏导并在 \((u,v)=(0,0)\) 处取值可得 \(\partial_vH(0,0)=1\)，故在完备局部环 \(\ZZ_3[[u]]\) 中为单位。
由形式隐函数定理（等价于对 \(\ZZ_3[[u]]\) 的 Hensel 型提升），存在且仅存在唯一的 \(v(u)\in\ZZ_3[[u]]\) 解 \(H(u,v)=0\) 且满足 \(v(0)=0\)。
由 \(\Pi(\Lambda(t),1+t)=0\) 的形式解唯一性知 \(t=27u\) 的代入与上述解一致，即 \(\Lambda(27u)=2+3v(u)\)。
可复现核验脚本：\texttt{scripts/exp\_fold\_zm\_leyang\_perron\_p3\_renormalization\_audit.py}。证毕。
\end{proof}

\begin{proposition}[\texorpdfstring{$27$}{27} 的最小性：\texorpdfstring{$t=9u$}{t=9u} 不产生 \texorpdfstring{$3$}{3}-进整系数展开]\label{prop:xi-terminal-zm-leyang-perron-p3-minimal-27}
\leanverified{paper\_xi\_terminal\_zm\_leyang\_perron\_p3\_minimal\_27}
不存在 \(\widetilde\Lambda(u)\in\ZZ_3[[u]]\) 使得
\[
\Pi(\widetilde\Lambda(u),1+9u)=0,\qquad \widetilde\Lambda(0)=2.
\]
等价地，\(\Lambda(9u)\notin\ZZ_3[[u]]\)。
\end{proposition}
\begin{proof}
在 \(\FF_3\) 上约化。
由于 \(9u\equiv 0\pmod 3\)，有
\[
\Pi(\lambda,1+9u)\equiv \Pi(\lambda,1)\pmod 3.
\]
又
\[
\Pi(\lambda,1)=\lambda^4-\lambda^3-3\lambda^2+\lambda+2\equiv \lambda^4-\lambda^3+\lambda-1
=(\lambda-1)(\lambda+1)^3\quad(\bmod\,3).
\]
若存在 \(\widetilde\Lambda(u)\in\ZZ_3[[u]]\) 满足命题条件，则其模 \(3\) 约化 \(\overline{\widetilde\Lambda}(u)\in\FF_3[[u]]\) 满足
\((\overline{\widetilde\Lambda}(u)-1)(\overline{\widetilde\Lambda}(u)+1)^3=0\)。
由于 \(\FF_3[[u]]\) 为整环，推出 \(\overline{\widetilde\Lambda}(u)=\pm 1\)；
而 \(\widetilde\Lambda(0)=2\equiv-1\pmod 3\)，故 \(\overline{\widetilde\Lambda}(u)\equiv-1\)，从而可写
\[
\widetilde\Lambda(u)=2+3v(u),\qquad v(u)\in\ZZ_3[[u]].
\]
将 \(\lambda=2+3v\)、\(t=9u\) 代入并在 \(\ZZ[u,v]\) 中展开得
\[
\Pi(2+3v,1+9u)=9\,H_9(u,v),
\]
其中
\[
H_9(u,v)=9u^{2}-18uv^{2}-24uv-5u+9v^{4}+21v^{3}+15v^{2}+3v\in\ZZ[u,v].
\]
模 \(3\) 约化时含 \(v\) 的各项系数均被 \(3\) 整除而消失，且 \(-5u\equiv u\pmod 3\)，故
\[
H_9(u,v)\equiv u\quad(\bmod\,3).
\]
于是 \(H_9(u,v(u))=0\) 在 \(\FF_3[[u]]\) 中强制 \(u=0\)，与 \(u\) 为形式变量矛盾。证毕。
\end{proof}

\begin{corollary}[Perron 分支系数的 \(3\)-进估值线性下界]\label{cor:xi-terminal-zm-leyang-perron-v3-linear-bound}
\leanverified{paper\_xi\_terminal\_zm\_leyang\_perron\_v3\_linear\_bound}
写
\[
\Lambda(t)=2+\sum_{n\ge 1}a_n t^n\qquad(a_n\in\QQ).
\]
则对任意 \(n\ge 1\) 有
\[
v_3(a_n)\ge 1-3n.
\]
等价地，存在 \(b_n\in\ZZ_3\) 使得
\[
a_n=\frac{b_n}{3^{3n-1}}\qquad(n\ge 1).
\]
\end{corollary}
\begin{proof}
由定理 \ref{thm:xi-terminal-zm-leyang-perron-p3-renormalization-hensel}，存在 \(v(u)=\sum_{n\ge 1}b_nu^n\in\ZZ_3[[u]]\) 使
\[
\Lambda(27u)=2+3v(u).
\]
比较系数得
\[
2+\sum_{n\ge 1}a_n(27u)^n
=2+3\sum_{n\ge 1}b_nu^n,
\]
从而 \(a_n=\frac{3b_n}{27^n}=\frac{b_n}{3^{3n-1}}\)，且 \(b_n\in\ZZ_3\)。证毕。
\end{proof}

\begin{theorem}[特征 \(3\) 的显式约化方程与自动结构]\label{thm:xi-terminal-zm-leyang-perron-p3-mod3-algebraic-automatic}
\leanverified{paper\_xi\_terminal\_zm\_leyang\_perron\_p3\_mod3\_algebraic\_automatic}
令 \(v(u)\in\ZZ_3[[u]]\) 为定理 \ref{thm:xi-terminal-zm-leyang-perron-p3-renormalization-hensel} 的唯一解，并记其逐项模 \(3\) 约化为 \(\overline v(u)\in\FF_3[[u]]\)。
则 \(\overline v(u)\) 满足显式三次代数方程
\[
\overline v(u)^3-\overline v(u)^2+\overline v(u)+u=0,
\qquad
\overline v(0)=0,
\]
并由初值条件确定唯一解。
因此 \(\overline v(u)\) 在 \(\FF_3(u)\) 上代数，其系数序列为 \(3\)-自动序列（等价地，具有有限 \(3\)-核）。
\end{theorem}
\begin{proof}
将定理 \ref{thm:xi-terminal-zm-leyang-perron-p3-renormalization-hensel} 的 \(H(u,v)\) 模 \(3\) 约化。
由于 \(27,18,24,3\equiv 0\pmod 3\)，且 \(-5\equiv 1\)、\(7\equiv 1\)、\(5\equiv -1\) 于 \(\FF_3\)，得
\[
H(u,v)\equiv u+v^3-v^2+v\quad(\bmod\,3),
\]
从而 \(\overline v\) 满足 \(\overline v^3-\overline v^2+\overline v+u=0\)。
又 \(\partial_v(u+v^3-v^2+v)|_{(0,0)}=1\neq 0\) 于 \(\FF_3\)，形式隐函数定理给出满足 \(\overline v(0)=0\) 的唯一解。
代数性推出自动性由 Christol 定理给出。证毕。
\end{proof}

\begin{theorem}[Perron 分支在 \texorpdfstring{$\QQ_3$}{Q3} 上的收敛半径]\label{thm:xi-terminal-zm-leyang-perron-p3-radius}
\leanverified{paper\_xi\_terminal\_zm\_leyang\_perron\_p3\_radius}
将 \(\Lambda(t)\) 视为 \(\QQ_3\) 上的幂级数，则其 \(3\)-进收敛半径严格等于
\[
R_3(\Lambda)=3^{-3}.
\]
\end{theorem}
\begin{proof}
由推论 \ref{cor:xi-terminal-zm-leyang-perron-v3-linear-bound}，\(v_3(a_n)\ge 1-3n\)，故
\(
|a_n|_3\le 3^{3n-1}
\)
并推出 \(R_3(\Lambda)\ge 3^{-3}\)。
另一方面，若仅有有限多个 \(n\) 使 \(b_n\not\equiv 0\pmod 3\)，则 \(\overline v(u)\) 为多项式；
但方程 \(\overline v^3-\overline v^2+\overline v+u=0\) 不存在非零多项式解（比较次数即可），矛盾。
故存在无穷多个 \(n\) 使 \(b_n\not\equiv 0\pmod 3\)，从而对这无穷多 \(n\) 有 \(v_3(a_n)=1-3n\)。
因此
\[
\limsup_{n\to\infty}|a_n|_3^{1/n}=3^{3},
\]
从而 \(R_3(\Lambda)\le 3^{-3}\)。
合并两端不等式得到 \(R_3(\Lambda)=3^{-3}\)。证毕。
\end{proof}

\begin{corollary}[\texorpdfstring{$p=3$}{p=3} 处的局部半稳定化与覆盖度降阶]\label{cor:xi-terminal-zm-leyang-perron-p3-semistable-degree-drop}
将 \(\ZZ_3\) 上的平面曲线族 \(\Pi(\lambda,1+t)=0\) 视为 \((\lambda,t)=(2,0)\) 的形式邻域内的局部模型。
在变量替换 \(t=27u\)、\(\lambda=2+3v\) 下，该形式邻域可写为 \(H(u,v)=0\) 且 \(\partial_vH(0,0)\in\ZZ_3^\times\)，因而在 \(\ZZ_3\) 上正则。
其特殊纤维（模 \(3\) 约化）由
\[
u=-v^3+v^2-v
\]
给出，从而在特征 \(3\) 上诱导出一般可分的三次覆盖 \(v\mapsto u\)。
\end{corollary}
\begin{proof}
正则性由 \(\partial_vH(0,0)\in\ZZ_3^\times\) 的隐函数口径直接得到。
模 \(3\) 约化由
\(
H(u,v)\equiv u+v^3-v^2+v
\)
立得。
最后，\(\frac{d}{dv}(-v^3+v^2-v)= -3v^2+2v-1\equiv 2v-1\) 于 \(\FF_3\) 非恒零，故该三次映射在一般点可分。证毕。
\end{proof}
\leanverified{paper\_xi\_terminal\_zm\_leyang\_perron\_p3\_semistable\_degree\_drop}

\input{subsubsec__xi-terminal-zm-leyang-perron-p3-ramification-inertia}

\begin{conclusion}[累积量密度序列的分母支撑仅含 \texorpdfstring{$\{2,3\}$}{\{2,3\}}]\label{con:xi-terminal-zm-leyang-perron-cumulant-denominator-2-3}
定义
\[
\psi(s):=\log\frac{\Lambda(e^s-1)}{2},\qquad
\kappa_n^\infty:=\psi^{(n)}(0)\quad(n\ge 1).
\]
则
\[
\kappa_n^\infty\in\ZZ\!\left[\frac16\right]\qquad(n\ge 1).
\]
\end{conclusion}
\begin{proof}[证明要点]
由结论 \ref{con:xi-terminal-zm-leyang-perron-global-pintegrality}，\(\Lambda(t)\) 系数属于 \(\ZZ[1/3]\)。
又
\[
(e^s-1)^k
=
k!\sum_{n\ge k}S(n,k)\frac{s^n}{n!},
\]
其中 \(S(n,k)\in\ZZ\)，故 \(\Lambda(e^s-1)\) 的各阶导数在 \(s=0\) 落在 \(\ZZ[1/3]\)。
再由 \(\psi(s)=\log(\Lambda(e^s-1))-\log 2\)，对数微分仅额外引入 \(\Lambda(0)=2\) 的分母因子，
因此最多新增素因子 \(2\)。
故 \(\kappa_n^\infty\in\ZZ[1/6]\)。证毕。
\end{proof}

\begin{conclusion}[\texorpdfstring{$p\neq 3$}{p!=3} 下的模 \(p\) 自动性与有限 \(p\)-核]\label{con:xi-terminal-zm-leyang-perron-modp-automaticity}
对任意素数 \(p\neq 3\)，记
\[
\overline{\Lambda}_p(t):=\sum_{n\ge 0}\overline{a}_n t^n\in\FF_p[[t]]
\]
为 \(\Lambda(t)\) 的逐项模 \(p\) 约化。
则 \(\overline{\Lambda}_p(t)\) 在 \(\FF_p(t)\) 上代数，且系数序列 \((\overline{a}_n)_{n\ge 0}\) 为 \(p\)-自动序列（等价地，其 \(p\)-核有限）。
\end{conclusion}
\begin{proof}[证明要点]
由结论 \ref{con:xi-terminal-zm-leyang-perron-global-pintegrality}，\(\Lambda(t)\in\ZZ_p[[t]]\)（\(p\neq 3\)），故可逐项约化得到 \(\overline{\Lambda}_p\)。
约化后仍满足
\(
\Pi(\overline{\Lambda}_p(t),1+t)=0
\)
的模 \(p\) 方程，故 \(\overline{\Lambda}_p\) 为 \(\FF_p(t)\)-代数幂级数。
由 Christol 定理，代数性当且仅当系数序列 \(p\)-自动，结论即得。证毕。
\end{proof}

\begin{conclusion}[三阶算子的单值群与 \texorpdfstring{$S_4$}{S4} 标准三维表示]\label{con:xi-terminal-zm-leyang-perron-fuchs-monodromy-s4}
设 \(\lambda_1,\dots,\lambda_4\) 为 \(\Pi(\lambda,y)=0\) 的四个解析分支，并定义
\[
f_i(y):=\lambda_i(y)-\frac14\qquad(1\le i\le 4).
\]
则
\[
V=\mathrm{span}_{\CC}\{f_1,f_2,f_3,f_4\}
=
\left\{\sum_{i=1}^4 c_i f_i:\ \sum_{i=1}^4 c_i=0\right\},
\qquad \dim_{\CC}V=3,
\]
且 \(V\) 恰为结论 \ref{con:xi-terminal-zm-leyang-perron-centered-fuchs-operator} 中 \(L_3\) 的解空间。
解析延拓在 \(V\) 上诱导的表示同构于 \(S_4\) 在
\(
\{(x_1,x_2,x_3,x_4)\in\CC^4:\sum x_i=0\}
\)
上的标准三维表示，故单值群同构于 \(S_4\)。
\end{conclusion}
\begin{proof}[证明要点]
由 Vieta 关系 \(\sum_i\lambda_i=1\) 得 \(\sum_i f_i=0\)，因此四分支张成空间至多三维。
另一方面，\(L_3\) 为三阶算子且各 \(f_i\) 均被其湮灭，故解空间维数恰为 \(3\)，于是两者一致。
由定理 \ref{thm:xi-terminal-zm-leyang-monodromy-s4}，\(\Pi\) 的四叶覆盖单值群为 \(S_4\)；
其在 \((f_1,\dots,f_4)\) 上为指标置换作用，并保持和为零超平面，得到标准三维表示。
该表示忠实，故诱导单值群同构于 \(S_4\)。证毕。
\end{proof}

\begin{conclusion}[指数对数坐标下的 Lee--Yang 微分闭合：\texorpdfstring{$\lambda(s)=\lambda_+(e^s)$}{lambda(s)} 的一阶代数方程与椭圆核一致性]\label{con:xi-terminal-zm-leyang-exponential-time-algebraic-differential-closure}
沿用谱四次
\[
\Pi(\lambda,y):=\lambda^4-\lambda^3-(2y+1)\lambda^2+\lambda+y(y+1)=0
\]
及其在 \(y=1\) 邻域满足 \(\lambda_+(1)=2\) 的解析特征支 \(\lambda_+(y)\)。
在 \(s=0\) 的邻域定义对数参数化
\[
\lambda(s):=\lambda_+(e^s),
\qquad
\dot\lambda:=\frac{\dd \lambda}{\dd s}.
\]
则 \(\lambda(s)\) 满足显式一阶代数微分闭式
\begin{equation}\label{eq:xi-terminal-zm-leyang-exponential-time-lambda-ode}
\boxed{\;
\begin{aligned}
0={}&(16\lambda^4+7\lambda^3-9\lambda^2+\lambda+1)\,\dot\lambda^{2}
(-16\lambda^5-4\lambda^4+20\lambda^3-5\lambda+1)\,\dot\lambda \\
&\qquad\qquad+(4\lambda^6-8\lambda^4+\lambda^3+4\lambda^2-\lambda).
\end{aligned}\;}
\end{equation}
并且将 \eqref{eq:xi-terminal-zm-leyang-exponential-time-lambda-ode} 视为关于 \(\dot\lambda\) 的二次方程，其判别式满足严格因式分解
\begin{equation}\label{eq:xi-terminal-zm-leyang-exponential-time-lambda-ode-discriminant}
\boxed{\;
\Delta_{\dot\lambda}(\lambda)
=\bigl(4\lambda^3-4\lambda+1\bigr)\bigl(2\lambda^3+2\lambda^2-\lambda+1\bigr)^2\; }.
\end{equation}
因此在函数域层面成立
\begin{equation}\label{eq:xi-terminal-zm-leyang-exponential-time-function-field}
\QQ(\lambda,\dot\lambda)=\QQ\!\left(\lambda,\sqrt{4\lambda^3-4\lambda+1}\right),
\end{equation}
从而指数对数坐标的微分闭合与椭圆核
\begin{equation}\label{eq:xi-terminal-zm-leyang-exponential-time-elliptic-kernel}
\mathcal E:\quad w^2=4\lambda^3-4\lambda+1
\end{equation}
处于同一函数域（与结论 \ref{con:xi-terminal-zm-elliptic-normalization} 的椭圆化一致）。
更进一步，外参 \(y=e^s\) 可由 \((\lambda,\dot\lambda)\) 有理重构为
\begin{equation}\label{eq:xi-terminal-zm-leyang-exponential-time-y-reconstruction}
\boxed{\;
y=
\frac{\dot\lambda\,(4\lambda^3-3\lambda^2-2\lambda+1)-2\lambda(\lambda-1)^2(\lambda+1)}{-2\lambda^2+4\lambda\dot\lambda+1}\; }.
\end{equation}
\end{conclusion}
\begin{proof}[证明要点（隐式微分与消元）]
对恒等式 \(\Pi(\lambda(s),e^s)\equiv0\) 微分得
\(
\partial_\lambda\Pi(\lambda,y)\,\dot\lambda+\partial_y\Pi(\lambda,y)\,\dot y=0
\)
且 \(\dot y=y\)。
由 \(\Pi(\lambda,y)=0\) 与
\(
\partial_\lambda\Pi(\lambda,y)\,\dot\lambda+y\,\partial_y\Pi(\lambda,y)=0
\)
得到关于 \(y\) 的两条二次代数约束；消去 \(y^2\) 得到 \(y\) 的有理表达式 \eqref{eq:xi-terminal-zm-leyang-exponential-time-y-reconstruction}，
代回即得 \eqref{eq:xi-terminal-zm-leyang-exponential-time-lambda-ode}。
判别式分解 \eqref{eq:xi-terminal-zm-leyang-exponential-time-lambda-ode-discriminant} 为直接计算并在 \(\ZZ[\lambda]\) 中因式分解。
函数域恒等式 \eqref{eq:xi-terminal-zm-leyang-exponential-time-function-field} 由判别式的唯一非平凡平方自由部 \((4\lambda^3-4\lambda+1)\) 立即推出；
而 \eqref{eq:xi-terminal-zm-leyang-exponential-time-elliptic-kernel} 与结论 \ref{con:xi-terminal-zm-elliptic-normalization} 的椭圆模型一致。
可复现核验脚本：\texttt{scripts/exp\_fold\_zm\_leyang\_exponential\_time\_differential\_closure\_audit.py}。证毕。
\end{proof}

\begin{conclusion}[\texorpdfstring{$\psi(s)=\log\lambda_+(e^s)-\log2$}{psi(s)} 的微分代数性与非 \texorpdfstring{$D$}{D}-有限性]\label{con:xi-terminal-zm-psi-differential-algebraic-non-dfinite}
令
\[
\psi(s):=\log\lambda_+(e^s)-\log2.
\]
则 \(\psi\) 在 \(s=0\) 的解析芽满足：
\begin{enumerate}
  \item \(\psi\) 为微分代数函数：其满足某个系数属 \(\QQ(s)\) 的非平凡代数微分方程；
  \item \(\psi\) 非 \(D\)-有限：不存在首项系数为多项式的非零线性常微分算子湮灭 \(\psi\)。
  因而泰勒系数列 \(\bigl(\psi^{(n)}(0)\bigr)_{n\ge 0}\) 非 \(P\)-递归（非 holonomic）。
\end{enumerate}
\end{conclusion}
\begin{proof}[证明要点]
第一条由结论 \ref{con:xi-terminal-zm-leyang-exponential-time-algebraic-differential-closure} 给出的 \(\lambda(s)=\lambda_+(e^s)\) 的一阶代数微分闭式推出：
\(\lambda\) 因而微分代数；又 \(\psi'=\dot\lambda/\lambda\) 为 \((\lambda,\dot\lambda)\) 的有理函数且 \(\lambda(0)=2\neq0\)，故 \(\psi\) 的解析芽同属微分代数闭包。
\par
第二条利用奇点集合的无限性。
由结论 \ref{con:xi-terminal-zm-discriminant-leyang}，\(\lambda_+(y)\) 的有限分歧值集合包含某个满足 \(y_\star\neq0\) 的代数分歧值（例如 Lee--Yang 三次 \(P_{\mathrm{LY}}(y)=0\) 的任一根）。
取任一对数支 \(s_\star:=\log y_\star\)，则指数映射的周期性给出无限前像族
\[
s_\star+2\pi i\,\ZZ.
\]
在每个点 \(s=s_\star+2\pi i k\) 处，复合函数 \(s\mapsto \lambda_+(e^s)\) 继承 \(y=y_\star\) 的代数分歧而出现不可去奇性，故 \(\psi\) 在 \(\CC\) 上具有无穷多个奇点。
另一方面，若 \(\psi\) 为 \(D\)-有限，则其满足某个首项系数为多项式的线性常微分方程，故其解析延拓的可能奇点仅能位于该首项系数的零点集合（有限集）与无穷远点。
两者矛盾，故 \(\psi\) 非 \(D\)-有限。
最后，\(D\)-有限解析芽的泰勒系数列必为 \(P\)-递归为标准等价，因而 \(\bigl(\psi^{(n)}(0)\bigr)\) 非 \(P\)-递归。证毕。
\end{proof}

\begin{conclusion}[主特征支局部喷射的 \(3\)-局部整性：\texorpdfstring{$\lambda_+(1+z)\in\ZZ[1/3][[z]]$}{lambda+(1+z) in Z[1/3]} 与 \texorpdfstring{$2^n\psi^{(n)}(0)\in\ZZ[1/3]$}{2^n psi^(n)(0) in Z[1/3]}]\label{con:xi-terminal-zm-lambda-plus-3local-denominators}
记
\[
\lambda_+(1+z)=2+\sum_{n\ge 1}a_n z^n\in\QQ[[z]],
\qquad
\psi(s)=\sum_{n\ge 1}\frac{\psi^{(n)}(0)}{n!}\,s^n.
\]
则：
\begin{enumerate}
  \item 对一切 \(n\ge 1\) 有 \(a_n\in\ZZ[1/3]\)，亦即 \(\lambda_+(1+z)\in\ZZ[1/3][[z]]\)；
  \item 对一切 \(n\ge 1\) 有 \(2^n\psi^{(n)}(0)\in\ZZ[1/3]\)，从而 \(\psi^{(n)}(0)\in\ZZ[1/6]\)。
\end{enumerate}
\end{conclusion}
\begin{proof}[证明要点（\(p\)-进隐函数与导数闭合）]
由 \(\Pi(2,1)=0\) 且
\(
\partial_\lambda\Pi(2,1)=9
\)，
对任一素数 \(p\neq3\)，\(\partial_\lambda\Pi(2,1)\) 在 \(\ZZ_p\) 中可逆。
由 \(p\)-进隐函数定理，存在唯一 \(\lambda_p(z)\in\ZZ_p[[z]]\) 使
\(
\Pi(\lambda_p(z),1+z)\equiv0
\)
且 \(\lambda_p(0)=2\)。
而 \(\lambda_+(1+z)\in\QQ[[z]]\) 亦满足同一初值条件，故在 \(\QQ_p[[z]]\) 中与 \(\lambda_p\) 一致；
于是 \(a_n\in\ZZ_p\) 对一切 \(p\neq3\) 成立，从而 \(a_n\in\ZZ[1/3]\)。
\par
再记 \(\lambda(s)=\lambda_+(e^s)\)。
由 \(e^s\) 在 \(s=0\) 处各阶导数皆为 \(1\)，复合求导（Fa\`a di Bruno）表明 \(\lambda^{(n)}(0)\) 为 \(\{\lambda_+^{(k)}(1)\}_{k\le n}\) 的整系数组合；
而 \(\lambda_+^{(k)}(1)=k!\,a_k\in\ZZ[1/3]\)，故 \(\lambda^{(n)}(0)\in\ZZ[1/3]\)。
又 \(\psi=\log(\lambda/2)\) 且 \(\lambda(0)=2\)，于是 \(\psi^{(n)}(0)\) 为 \(\{\lambda^{(k)}(0)/2\}_{k\le n}\) 的整系数多项式（来自对数导数的递推），
每项分母至多贡献 \(2^n\)，从而 \(2^n\psi^{(n)}(0)\in\ZZ[1/3]\)。证毕。
\end{proof}

\begin{conclusion}[有限阶局部喷射唯一确定谱四次：\(13\)-阶芽锁定 \texorpdfstring{$\Pi(\lambda,y)$}{Pi(lambda,y)} 与 Lee--Yang 三次]\label{con:xi-terminal-zm-local-jet-rigidly-determines-pi}
在二元多项式集合中考虑逆问题：设
\(
Q(\lambda,y)\in\QQ[\lambda,y]
\) 满足 \(\deg_\lambda Q\le4\)、\(\deg_y Q\le2\)，并且 \(\lambda^4\) 的系数归一为 \(1\)。
若已知主特征支 \(\lambda_+(y)\) 在 \(y=1\) 的 \(13\)-阶喷射（即 \(a_1,\dots,a_{13}\) 连同 \(\lambda_+(1)=2\)），
则满足
\[
Q(\lambda_+(y),y)\equiv 0 \qquad\text{于}\qquad \QQ[[y-1]]/(y-1)^{14}
\]
的上述 \(Q\) 唯一存在，且必为谱四次 \(\Pi\) 本身。
因此由恒等式
\(
\mathrm{Disc}_{\lambda}(\Pi)=-y(y-1)P_{\mathrm{LY}}(y)
\)
所定义的 Lee--Yang 三次 \(P_{\mathrm{LY}}(y)\) 亦被该有限喷射唯一决定。
\end{conclusion}
\begin{proof}[证明要点（线性化与可逆性）]
将一般 \(Q\) 写为
\[
Q(\lambda,y)=\lambda^4+\sum_{(i,j)\neq(4,0)}c_{ij}\lambda^i y^j,
\qquad 0\le i\le4,\ 0\le j\le2,
\]
未知量为 \(14\) 个系数 \(c_{ij}\in\QQ\)。
代入 \(y=1+z\) 与
\(
\lambda=\lambda_+(1+z)=2+\sum_{n=1}^{13}a_n z^n+O(z^{14})
\)，
并取模 \(z^{14}\)，条件 \(Q(\lambda_+(1+z),1+z)\equiv0\pmod{z^{14}}\) 等价于 \(z^0,\dots,z^{13}\) 的 \(14\) 个系数同时为零，
从而得到一个 \(14\times14\) 的线性方程组。
直接计算可得其系数矩阵行列式非零，故解唯一存在。
另一方面，\(\Pi\) 显然给出一个满足条件的解，故唯一解必为 \(\Pi\)。
由判别式恒等式 \(\mathrm{Disc}_{\lambda}(\Pi)=-y(y-1)P_{\mathrm{LY}}(y)\) 立得 \(P_{\mathrm{LY}}\) 的唯一性。证毕。
可复现核验脚本：\texttt{scripts/exp\_fold\_zm\_leyang\_exponential\_time\_differential\_closure\_audit.py}。
\end{proof}



\begin{conclusion}[重量大偏差的代数鞍点方程：倾斜参数的五次消元证书]\label{con:xi-terminal-zm-ldp-saddle-quintic}
对 \(\alpha\in(0,\tfrac12)\)，令 \(y(\alpha)>0\) 为满足鞍点条件
\[
\alpha=\frac{y\,\lambda_+'(y)}{\lambda_+(y)}
\]
的唯一正解（\(\lambda_+\) 为结论 \ref{con:xi-terminal-zm-lln-clt} 中 \(\Pi(\lambda,y)=0\) 的主特征支）。
则 \(y(\alpha)\) 必满足关于 \((\alpha,y)\) 的显式五次代数方程
\[
\begin{aligned}
0=&\ 256\alpha^4y^5+411\alpha^4y^4-91\alpha^4y^3-379\alpha^4y^2-165\alpha^4y-32\alpha^4\\
&-512\alpha^3y^5-566\alpha^3y^4+337\alpha^3y^3+512\alpha^3y^2+197\alpha^3y+32\alpha^3\\
&+384\alpha^2y^5+228\alpha^2y^4-298\alpha^2y^3-214\alpha^2y^2-28\alpha^2y\\
&-128\alpha y^5-8\alpha y^4+86\alpha y^3+16\alpha y^2-4\alpha y\\
&+16y^5-8y^4-4y^3+y^2,
\end{aligned}
\]
并且速率函数在该支上可写为
\[
I(\alpha)=\alpha\log y(\alpha)-\log \lambda_+\!\bigl(y(\alpha)\bigr)+\log 2.
\]
因此该 LDP 的鞍点求解与速率函数评估在严格意义下为有限代数可计算对象。
\end{conclusion}
\begin{proof}[证明要点（隐式微分 + 消元）]
鞍点条件等价于 \(\alpha=\psi'(s)\) 与 \(y=e^s\) 的联立，从而可写为
\(\alpha=-y\,\partial_y\Pi(\lambda,y)/(\lambda\,\partial_\lambda\Pi(\lambda,y))\) 与 \(\Pi(\lambda,y)=0\) 的共同解（其中 \(\lambda=\lambda_+(y)\)）。
对变量 \(\lambda\) 作消元（例如取关于 \(\lambda\) 的 resultant）即可得到所示五次多项式；
并由结论 \ref{con:xi-terminal-zm-ldp} 的 Legendre 形式推出 \(I(\alpha)\) 的鞍点表达。
上述消元证书可由脚本 \texttt{scripts/exp\_fold\_zm\_bivariate\_partition\_audit.py} 复核。证毕。
\end{proof}

\begin{conclusion}[自反型同余通道的唯一性：仅 \(y^2+y+1=0\) 触发单位圆对称]\label{con:xi-terminal-zm-self-inversive-unique}
对 \(|y|=1\)，称 \(\Pi(\lambda,y)\) 为自反（self-inversive）多项式，
若存在常数 \(c\in\CC\)（\(|c|=1\)）使得
\[
\lambda^4\,\overline{\Pi\!\left(\frac{1}{\overline{\lambda}},y\right)}=c\,\Pi(\lambda,y).
\]
则当且仅当 \(y(y+1)=-1\)，亦即 \(y\) 为本原三次单位根（满足 \(y^2+y+1=0\)），\(\Pi(\lambda,y)\) 才具有该自反性。
在且仅在此情形，\(\Pi(\lambda,y)\) 的根集关于单位圆满足严格对称：要么位于单位圆上，要么以
\(\lambda\mapsto 1/\overline{\lambda}\) 成对出现；并且必存在一对互为倒共轭的根，其模长严格大于 \(1\)。
因此“同余权重的单位圆对称”在该体系中是严格量子化的：模 \(3\) 为唯一触发自反对称的非平凡同余通道。
\end{conclusion}
\begin{proof}[证明要点]
写
\(
\Pi(\lambda,y)=\sum_{k=0}^{4}a_k(y)\lambda^k
\)，其中
\(
a_4=1,\ a_3=-1,\ a_2=-(2y+1),\ a_1=1,\ a_0=y(y+1)
\)。
自反性等价于存在 \(|c|=1\) 使 \(a_k(y)=c\,\overline{a_{4-k}(y)}\) 对所有 \(k\) 成立。
由 \(a_3=-1=c\,\overline{a_1}=c\) 得 \(c=-1\)。
再由 \(a_4=1=c\,\overline{a_0}=-\overline{y(y+1)}\) 得 \(\overline{y(y+1)}=-1\)；
当 \(|y|=1\) 时 \(\overline{y(y+1)}=(1+y)/y^2\)，故等价于 \(y^2+y+1=0\)。
根集的单位圆对称为自反多项式的标准结论。证毕。
\end{proof}

\begin{conclusion}[重量细分系数的二维线性动力：\((m,k)\) 平面上的有限邻域闭合递推]\label{con:xi-terminal-zm-2d-recurrence-amk}
令
\[
a_{m,k}:=\sum_{\substack{x\in X_m\\ |x|_1=k}} d_m(x),
\qquad
Z_m(y)=\sum_{k\in\ZZ} a_{m,k}y^k .
\]
则由结论 \ref{con:xi-terminal-zm-order4-rational} 的四阶递推在 \(y\) 上的多项式结构，得到 \(a_{m,k}\) 的闭合二维递推：
对一切 \(m\ge 4\) 与 \(k\in\ZZ\)，
\[
\boxed{\;
a_{m,k}=a_{m-1,k}+a_{m-2,k}+2a_{m-2,k-1}-a_{m-3,k}-a_{m-4,k-1}-a_{m-4,k-2}\;}
\]
并以边界条件 \(a_{m,k}=0\)（当 \(k<0\) 或 \(k>\lceil m/2\rceil\)）封闭。
该二维递推把“推前分布的重量剖分”转写为一个完全局域的线性动力系统，使得重量层（固定 \(k\)）与长度层（固定 \(m\)）之间的耦合仅通过有限邻域发生；
因此重量分布的全体精细信息可在不显式构造任何高维碰撞核的前提下，由 \((m,k)\) 层级递推直接生成并审计。
\end{conclusion}
\begin{proof}
将结论 \ref{con:xi-terminal-zm-order4-rational} 的递推式写成 \(y\)-幂级数，并比较 \(y^k\) 的系数即可得到所示二维递推；
边界条件由黄金均值合法语法的最大 \(1\)-计数 \(\max_{x\in X_m}|x|_1=\lceil m/2\rceil\) 直接给出。证毕。
\end{proof}

% (User input window)

\begin{conclusion}[固定重量层系数生成函数的分母刚性：统一闭式与二阶递推]\label{con:xi-terminal-zm-fixed-weight-generating-function-rigidity}
沿用结论 \ref{con:xi-terminal-zm-order4-rational} 与结论 \ref{con:xi-terminal-zm-2d-recurrence-amk} 的记号，并记
\[
a_{m,k}:=[y^k]\,Z_m(y),\qquad
A_k(t):=\sum_{m\ge 0}a_{m,k}t^m=[y^k]F(t,y),
\]
其中
\[
F(t,y)=\sum_{m\ge 0}Z_m(y)t^m
=\frac{1+yt-yt^2-y^2t^3}{1-t-(2y+1)t^2+t^3+y(y+1)t^4}.
\]
则对每个固定整数 \(k\ge 0\)，存在唯一多项式 \(\mathcal{P}_k(t)\in\ZZ[t]\) 使得
\[
\boxed{\;
A_k(t)=\frac{\mathcal{P}_k(t)}{(1-t)^{2k+2}(1+t)^{k+1}}\;}
\]
并且 \(\mathcal{P}_k\) 可由显式二阶递推唯一确定：令
\[
\mathcal{P}_0(t)=1,\qquad
\mathcal{P}_1(t)=-t\bigl(t^4-t^3-1\bigr),\qquad
\mathcal{P}_2(t)=t^3\bigl(t^5-2t^4-t^3+t^2+t+1\bigr),
\]
则对一切 \(k\ge 3\) 有
\[
\boxed{\;
\mathcal{P}_k(t)=(2t^2-t^4)\mathcal{P}_{k-1}(t)-t^4(1-t)^2(1+t)\mathcal{P}_{k-2}(t)\; }.
\]
并且当 \(k\ge 1\) 时成立整除性
\[
\boxed{\; t^{2k-1}\mid \mathcal{P}_k(t)\;}
\]
\end{conclusion}
\begin{proof}[证明要点]
将分母记为
\[
D(t,y):=1-t-(2y+1)t^2+t^3+y(y+1)t^4=D_0(t)+yD_1(t)+y^2D_2(t),
\]
其中
\[
D_0(t)=(1-t)^2(1+t),\qquad D_1(t)=-2t^2+t^4,\qquad D_2(t)=t^4,
\]
并将分子记为
\(
N(t,y)=1+y(t-t^2)-y^2t^3
\)。
写 \(F(t,y)=\sum_{k\ge 0}A_k(t)y^k\)，则由恒等式 \(D(t,y)F(t,y)=N(t,y)\) 比较 \(y^k\) 系数得
\[
D_0(t)A_k(t)+D_1(t)A_{k-1}(t)+D_2(t)A_{k-2}(t)=N_k(t),
\]
其中 \(N_0(t)=1\)、\(N_1(t)=t-t^2\)、\(N_2(t)=-t^3\)、且 \(N_k(t)=0\)（当 \(k\ge 3\)），并约定 \(A_{-1}=A_{-2}=0\)。
令 \(\mathcal{P}_k(t):=D_0(t)^{k+1}A_k(t)\)，则上式等价于
\[
\mathcal{P}_k(t)=\bigl(2t^2-t^4\bigr)\mathcal{P}_{k-1}(t)-t^4D_0(t)\mathcal{P}_{k-2}(t)+N_k(t)D_0(t)^k.
\]
对 \(k\ge 3\) 时 \(N_k\equiv 0\)，得到所示齐次二阶递推；
\(k=0,1,2\) 的显式初值由上式直接化简得到。
由于 \(D_0,D_1,D_2,N_k\in\ZZ[t]\) 且初值整系数，递推给出 \(\mathcal{P}_k\in\ZZ[t]\)。
\par
另一方面，\(\mathcal{P}_1(t)\) 含因子 \(t\)、\(\mathcal{P}_2(t)\) 含因子 \(t^3\)，且对 \(k\ge 3\) 的递推右端两项分别额外带因子 \(t^2\) 与 \(t^4\)，
故可归纳得到 \(t^{2k-1}\mid \mathcal{P}_k(t)\)（\(k\ge 1\)）。
最后由 \(D_0(t)^{k+1}=(1-t)^{2k+2}(1+t)^{k+1}\) 得到分母结构。证毕。
\end{proof}

\begin{conclusion}[分子多项式的黄金因子：模 \texorpdfstring{$t^2-t-1$}{t^2-t-1} 的五步标度闭合]\label{con:xi-terminal-zm-fixed-weight-numerator-golden-factor}
在结论 \ref{con:xi-terminal-zm-fixed-weight-generating-function-rigidity} 的记号下，令
\[
\Phi_{\mathrm{gold}}(t):=t^2-t-1\in\ZZ[t].
\]
则对一切整数 \(n\ge 0\)，成立严格整除性
\[
\boxed{\;\Phi_{\mathrm{gold}}(t)\mid \mathcal{P}_{5n+4}(t)\; }.
\]
\end{conclusion}
\begin{proof}[证明要点]
在商环 \(R:=\ZZ[t]/(\Phi_{\mathrm{gold}})\) 中有关系 \(t^2=t+1\)，从而
\[
(1-t)^2(1+t)=1-t-t^2+t^3\equiv 1,\qquad
2t^2-t^4\equiv -t,\qquad
t^4\equiv 3t+2
\pmod{\Phi_{\mathrm{gold}}}.
\]
将结论 \ref{con:xi-terminal-zm-fixed-weight-generating-function-rigidity} 的齐次递推在 \(R\) 中约化，得到对 \(k\ge 3\)
\[
\mathcal{P}_k\equiv -t\,\mathcal{P}_{k-1}-(3t+2)\mathcal{P}_{k-2}.
\]
令向量 \(v_k:=(\mathcal{P}_k,\mathcal{P}_{k-1})^{\mathsf T}\)，则对 \(k\ge 3\) 有
\[
v_k\equiv M v_{k-1},\qquad
M:=\begin{pmatrix}-t&-(3t+2)\\[2pt]1&0\end{pmatrix}\in \Mat_{2}(R).
\]
在 \(R\) 中对 \(M\) 作直接幂计算并约化可得
\[
M^5\equiv -(55t+34)\,I_2.
\]
因此对一切 \(k\ge 2\) 有五步标度关系
\[
v_{k+5}\equiv M^5 v_k\equiv -(55t+34)\,v_k,
\qquad\text{特别地}\quad
\mathcal{P}_{k+5}\equiv -(55t+34)\,\mathcal{P}_k.
\]
由初值在 \(R\) 中的约化计算
\(
\mathcal{P}_1\equiv -(1+t)
\)、\(
\mathcal{P}_2\equiv -(2+3t)
\)，进而由递推得到 \(\mathcal{P}_4\equiv 0\)。
于是对任意 \(n\ge 0\)，反复应用五步标度推出 \(\mathcal{P}_{5n+4}\equiv 0\)，
等价于 \(\Phi_{\mathrm{gold}}(t)\mid \mathcal{P}_{5n+4}(t)\)。
可复核脚本：\texttt{scripts/exp\_fold\_zm\_low\_weight\_quasipolynomial\_audit.py}。证毕。
\end{proof}

\begin{conclusion}[固定重量层系数的准多项式性：周期 \(2\)、次数量子化与首项系数闭式]\label{con:xi-terminal-zm-fixed-weight-quasipolynomial}
沿用结论 \ref{con:xi-terminal-zm-order4-rational} 与结论 \ref{con:xi-terminal-zm-2d-recurrence-amk} 的记号，并记
\[
a_{m,k}:=[y^k]\,Z_m(y)\qquad(m\ge 0,\ k\ge 0).
\]
则对每个固定整数 \(k\ge 0\)，存在唯一一对有理系数多项式 \(P_k,Q_k\in\QQ[m]\)，使对所有 \(m\ge 0\) 成立
\[
\boxed{\;a_{m,k}=P_k(m)+(-1)^m Q_k(m)\;}
\]
并且次数满足
\[
\deg P_k=2k+1,\qquad \deg Q_k=k.
\]
此外，\(P_k\) 与 \(Q_k\) 的首项系数具有封闭形式
\[
\bigl[m^{2k+1}\bigr]P_k(m)=\frac{1}{2^{k+1}(2k+1)!},
\qquad
\bigl[m^{k}\bigr]Q_k(m)=\frac{1}{2^{2k+2}k!}.
\]
等价地，令 \(m=2n+\varepsilon\)（\(\varepsilon\in\{0,1\}\)），则两条奇偶子序列 \(n\mapsto a_{2n,k}\) 与 \(n\mapsto a_{2n+1,k}\) 均为 \(n\) 的次数 \(2k+1\) 多项式。
并且两条奇偶分支的最高次项一致：当 \(n\to\infty\)（固定 \(k\)）时，
\[
a_{2n,k}=\frac{2^k}{(2k+1)!}\,n^{2k+1}+O_k(n^{2k}),
\qquad
a_{2n+1,k}=\frac{2^k}{(2k+1)!}\,n^{2k+1}+O_k(n^{2k}).
\]
\end{conclusion}
\begin{proof}[证明要点]
由结论 \ref{con:xi-terminal-zm-fixed-weight-generating-function-rigidity}，
对每个固定 \(k\ge 0\) 有显式闭式
\[
A_k(t)=\frac{\mathcal{P}_k(t)}{(1-t)^{2k+2}(1+t)^{k+1}}\in\QQ(t),
\]
从而 \(A_k(t)\) 仅在 \(t=\pm 1\) 处出现极点，且在 \(t=1\) 与 \(t=-1\) 处的极点阶分别恰为 \(2k+2\) 与 \(k+1\)。
对 \(A_k\) 作部分分式并逐项取系数即可得到对所有 \(m\ge 0\) 成立的分解
\(
a_{m,k}=P_k(m)+(-1)^mQ_k(m)
\)
及次数界；唯一性来自于：若 \(P(m)+(-1)^mQ(m)\equiv 0\) 对所有整数 \(m\ge 0\) 成立，则必有 \(P\equiv 0\) 且 \(Q\equiv 0\)。
\par
首项系数由 \(t=1\) 处的最高阶极点给出：由结论 \ref{con:xi-terminal-zm-fixed-weight-generating-function-rigidity} 的递推在 \(t=1\) 处代入可归纳得 \(\mathcal{P}_k(1)=1\)；并且 \((1+t)^{-(k+1)}\sim 2^{-(k+1)}\)（\(t\to 1\)），故
\(
A_k(t)\sim 2^{-(k+1)}(1-t)^{-(2k+2)}
\)，从而
\(
[t^m]A_k(t)=2^{-(k+1)}\frac{m^{2k+1}}{(2k+1)!}+O(m^{2k})
\)
给出所示首项系数。
\par
对 \(t=-1\) 的首项系数，记分母
\[
D(t,y)=D_0(t)+yD_1(t)+y^2D_2(t),
\qquad
D_0(t)=D(t,0),\ \ D_1(t)=\partial_y D(t,0)=-2t^2+t^4,
\]
则
\(
D_0(t)\sim 4(1+t)
\)
（\(t\to-1\)）且 \(D_1(-1)=-1\)。
由于 \(F(t,y)\) 的分子在 \(y=0\) 处取值为 \(1\)，且其 \(y\)-次数至多为 \(2\)，故 \(A_k(t)\) 在 \(t=-1\) 处的最高阶极点完全由 \([y^k]\bigl(1/D(t,y)\bigr)\) 决定。
写
\[
\frac{1}{D(t,y)}=\frac{1}{D_0(t)}
\left(1+\frac{yD_1(t)+y^2D_2(t)}{D_0(t)}\right)^{-1},
\]
则 \([y^k]\bigl(1/D(t,y)\bigr)\) 的最高阶项为 \((-1)^kD_1(t)^k/D_0(t)^{k+1}\)；
结合 \(D_0(t)\sim 4(1+t)\) 与 \(D_1(-1)=-1\) 得：
\[
A_k(t)=\frac{1}{4^{k+1}}(1+t)^{-(k+1)}+O\!\bigl((1+t)^{-k}\bigr)\qquad(t\to-1),
\]
从而 \(t=-1\) 的极点阶恰为 \(k+1\)，因此 \(\deg Q_k=k\)。
并且由
\(
[t^m](1+t)^{-(k+1)}=(-1)^m\binom{m+k}{k}
\)
得到
\(
[m^k]Q_k(m)=4^{-(k+1)}/k!=2^{-(2k+2)}/k!
\)。
上述分解与后续低阶闭式可由脚本 \texttt{scripts/exp\_fold\_zm\_low\_weight\_quasipolynomial\_audit.py} 复核。证毕。
\end{proof}

\begin{conclusion}[两项主部的显式化：多项式主增长与交替振荡的阶分离]\label{con:xi-terminal-zm-fixed-weight-two-term-asymptotic}
沿用结论 \ref{con:xi-terminal-zm-fixed-weight-generating-function-rigidity} 的记号，并记 \(\mathcal{P}_k(t)\in\ZZ[t]\) 为 \(A_k(t)\) 的分子。
则对一切 \(k\ge 0\) 有
\[
\boxed{\;
\mathcal{P}_k(1)=\mathcal{P}_k(-1)=1,\qquad \mathcal{P}_k'(1)=0\;}
\]
并且当 \(k\ge 1\) 时成立二阶与三阶导数的一次律：
\[
\boxed{\;
\mathcal{P}_k''(1)=4-12k,\qquad \mathcal{P}_k^{(3)}(1)=42-78k\; }.
\]
因此，对每个固定 \(k\)（令 \(m\to\infty\)）有两极点主部叠加的显式展开
\[
\boxed{\;
a_{m,k}
=\frac{1}{2^{k+1}}\binom{m+2k+1}{2k+1}
+\frac{k+1}{2^{k+2}}\binom{m+2k}{2k}
+O_k(m^{2k-1})
+(-1)^m\left(\frac{1}{2^{2k+2}}\binom{m+k}{k}+O_k(m^{k-1})\right)\; }.
\]
其中 \(O_k(m^{2k-1})\) 来自 \(t=1\) 的次高阶极点，而交替项来自 \(t=-1\) 的极点。
特别地，
\[
a_{m,k}
=\frac{m^{2k+1}}{2^{k+1}(2k+1)!}
+(-1)^m\frac{m^k}{2^{2k+2}k!}
+O_k(m^{2k}).
\]
\end{conclusion}
\begin{proof}[证明要点]
由结论 \ref{con:xi-terminal-zm-fixed-weight-generating-function-rigidity} 的递推
\[
\mathcal{P}_k(t)=(2t^2-t^4)\mathcal{P}_{k-1}(t)-t^4(1-t)^2(1+t)\mathcal{P}_{k-2}(t)
\]
在 \(t=\pm 1\) 处代入即可得 \(\mathcal{P}_k(\pm1)=\mathcal{P}_{k-1}(\pm1)=1\)。
对该递推在 \(t=1\) 处作导数并注意 \((2t^2-t^4)'|_{t=1}=0\) 且 \(t^4(1-t)^2(1+t)\) 在 \(t=1\) 处至少二阶消失，得到 \(\mathcal{P}_k'(1)=\mathcal{P}_{k-1}'(1)=0\)。
二阶与三阶导数律同理由在 \(t=1\) 处对递推作二次与三次微分并归纳得到（以 \(k=1,2\) 的直接核验作为起始）。
\par
最后由
\[
A_k(t)=\frac{\mathcal{P}_k(t)}{(1-t)^{2k+2}(1+t)^{k+1}}
\]
在 \(t=1\) 与 \(t=-1\) 处的局部展开作系数抽取：
在 \(t=1\) 处，
\(
(1+t)^{-(k+1)}=2^{-(k+1)}\bigl(1+\frac{k+1}{2}(1-t)+O((1-t)^2)\bigr)
\)
并结合 \(\mathcal{P}_k(1)=1\)、\(\mathcal{P}_k'(1)=0\) 得到所示两项主部；
在 \(t=-1\) 处，由 \(\mathcal{P}_k(-1)=1\) 得
\[
A_k(t)=\frac{1}{2^{2k+2}}(1+t)^{-(k+1)}+O\!\bigl((1+t)^{-k}\bigr)\qquad(t\to-1),
\]
从而
\(
[t^m]A_k(t)=(-1)^m\frac{1}{2^{2k+2}}\binom{m+k}{k}+(-1)^mO_k(m^{k-1})
\)。
证毕。
\end{proof}

\begin{conclusion}[固定重量层系数的最小常系数递推：阶 \texorpdfstring{$3k+3$}{3k+3} 与 \texorpdfstring{$\pm 1$}{±1} 极点分解]\label{con:xi-terminal-zm-fixed-weight-minimal-cc-recurrence-order}
固定 \(k\ge 0\)。
则序列 \(\{a_{m,k}\}_{m\ge 0}\) 在常系数递推意义下为 \(C\)-finite：
存在一个阶数为 \(3k+3\) 的齐次常系数线性递推，使得对所有充分大 \(m\) 成立；其特征多项式（以移位算子 \(E:a_m\mapsto a_{m+1}\) 记）为
\[
\boxed{\;(E-1)^{2k+2}(E+1)^{k+1}\;}
\]
并且在“充分大 \(m\)”意义下该阶数为最小：不存在阶数 \(<3k+3\) 的非平凡常系数递推能在全部充分大 \(m\) 上湮灭 \(\{a_{m,k}\}\)。
\end{conclusion}
\begin{proof}[证明要点]
由结论 \ref{con:xi-terminal-zm-fixed-weight-generating-function-rigidity}，
\(
A_k(t)=\sum_{m\ge 0}a_{m,k}t^m
\)
为有理函数，且其分母为
\(
Q_k(t)=(1-t)^{2k+2}(1+t)^{k+1}
\)
并满足 \(Q_k(0)=1\)。
由恒等式 \(Q_k(t)A_k(t)=\mathcal{P}_k(t)\) 逐项比较系数，得到由 \(Q_k\) 的系数给出的常系数线性递推；
当 \(m\) 足够大时右端系数归零，于是递推齐次化。
将 \(Q_k(t)\) 写成
\(
Q_k(t)=\sum_{j=0}^{3k+3}q_j t^j
\)
即可得到对应阶数为 \(3k+3\) 的递推；其特征多项式等价于
\(
E^{3k+3}Q_k(1/E)=(E-1)^{2k+2}(E+1)^{k+1}
\)。
\par
最小性归结为 \(\gcd(\mathcal{P}_k,Q_k)=1\)。
由于 \(Q_k\) 的全部不可约因子仅为 \((1-t)\) 与 \((1+t)\)，而结论 \ref{con:xi-terminal-zm-fixed-weight-two-term-asymptotic} 给出 \(\mathcal{P}_k(\pm1)=1\)，
故 \(\mathcal{P}_k\) 不被 \((1-t)\) 或 \((1+t)\) 整除，从而 \(Q_k\) 在既约意义下为最小湮灭分母（只差非零常数因子），递推阶数不可再降。证毕。
\end{proof}

\begin{conclusion}[低重量层系数的全量闭式：\(k\le 3\) 时对所有 \(m\ge 0\) 逐点成立]\label{con:xi-terminal-zm-fixed-weight-kle3-closed}
在结论 \ref{con:xi-terminal-zm-fixed-weight-quasipolynomial} 的记号下，令 \(n\ge 0\)。
则 \(k=0,1,2,3\) 的两条奇偶子序列 \(\{a_{2n,k}\}_{n\ge 0}\)、\(\{a_{2n+1,k}\}_{n\ge 0}\) 具有如下逐点闭式：
\[
a_{2n,0}=n+1,\qquad a_{2n+1,0}=n+1.
\]
\[
a_{2n,1}=\frac{1}{3}n^3+\frac{3}{2}n^2+\frac{1}{6}n,
\qquad
a_{2n+1,1}=\frac{1}{3}n^3+2n^2+\frac{5}{3}n+1.
\]
\[
a_{2n,2}=\frac{1}{30}n^5+\frac{3}{8}n^4-\frac{1}{12}n^3-\frac{7}{8}n^2+\frac{11}{20}n,
\]
\[
a_{2n+1,2}=\frac{1}{30}n^5+\frac{11}{24}n^4+\frac{3}{4}n^3-\frac{11}{24}n^2+\frac{13}{60}n.
\]
\[
a_{2n,3}=\frac{1}{630}n^7+\frac{1}{30}n^6+\frac{7}{360}n^5-\frac{17}{24}n^4+\frac{523}{360}n^3-\frac{33}{40}n^2+\frac{11}{420}n,
\]
\[
a_{2n+1,3}=\frac{1}{630}n^7+\frac{7}{180}n^6+\frac{23}{180}n^5-\frac{19}{36}n^4+\frac{29}{180}n^3+\frac{22}{45}n^2-\frac{61}{210}n.
\]
\end{conclusion}
\begin{proof}[证明要点]
对 \(k\le 3\)，由 \(A_k(t)=[y^k]F(t,y)\in\QQ(t)\) 且其极点仅位于 \(t=\pm1\) 可知
\(\{a_{2n,k}\}_{n\ge 0}\)、\(\{a_{2n+1,k}\}_{n\ge 0}\)
均为关于 \(n\) 的多项式序列（次数至多 \(2k+1\)）。
因此可用有限点值回收唯一多项式并与二维递推（结论 \ref{con:xi-terminal-zm-2d-recurrence-amk}）逐点核对；
亦可直接对 \(A_k(t)\) 作部分分式并回收系数。
上述闭式与其逐点一致性可由脚本 \texttt{scripts/exp\_fold\_zm\_low\_weight\_quasipolynomial\_audit.py} 复核。证毕。
\end{proof}



\paragraph{硬边界层的精确层级：二项式基整数证书与纯极点生成函数}
结论 \ref{con:xi-terminal-zm-2d-recurrence-amk} 的二维递推在硬边界方向（\(k\) 贴近 \(\lceil m/2\rceil\)）发生进一步的差分退化：
固定层深 \(j\) 后，偶长与奇长的边界系数在长度参数 \(k\) 上不仅具有渐近展开，
而是给出\emph{逐点成立}的精确多项式结构；并且其 Newton 二项式基展开自动产生可审计的整系数证书。

\begin{conclusion}[硬边界系数的精确多项式化：偶长 \texorpdfstring{$m=2k$}{m=2k} 的近最大重量层仅呈代数级增长]\label{con:xi-terminal-zm-hard-boundary-even-polynomial}
沿用结论 \ref{con:xi-terminal-zm-order4-rational} 的 \(Z_m(y)\) 记号。
对固定整数 \(j\ge 0\)，定义偶长硬边界系数
\[
A_j(k):=[y^{k-j}]\,Z_{2k}(y)\qquad(k\ge 0),
\]
其中约定当 \(k<j\) 时 \(A_j(k)=0\)。
则 \(A_j(k)\) 在所有整数 \(k\ge 0\) 上为精确多项式，其次数恰为 \(4j+2\)，且首项系数严格为 \(1/(4j+2)!\)。
更进一步，\(A_j(k)\) 存在唯一的二项式基展开
\[
A_j(k)=\sum_{r=j}^{4j+2}c_{j,r}\binom{k}{r},\qquad c_{j,r}\in\ZZ,\ \ c_{j,4j+2}=1,
\]
其中约定 \(\binom{k}{r}=0\)（当 \(r>k\)）。
因此在硬边界方向 \(k-j\) 处，推前重量剖分的每一固定层深 \(j\) 仅携带代数级前因子；
并且当 \(j\mapsto j+1\) 时，多项式次数严格增加 \(4\)。
\end{conclusion}

\begin{conclusion}[偶长硬边界系数的前三层闭式：二项式基展开给出完全可审计的整系数证书]\label{con:xi-terminal-zm-hard-boundary-even-first3}
在结论 \ref{con:xi-terminal-zm-hard-boundary-even-polynomial} 的记号下，前三个非平凡层（\(j=1,2,3\)）满足下列逐点恒等式（对所有整数 \(k\ge 0\) 成立）：
\[
\boxed{\ A_1(k)=2\binom{k}{1}+5\binom{k}{2}+9\binom{k}{3}+7\binom{k}{4}+4\binom{k}{5}+\binom{k}{6}\ }
\]
等价地，
\[
A_1(k)=\frac{k^6+9k^5+55k^4+435k^3-56k^2+996k}{720}.
\]
\[
\boxed{\ A_2(k)=3\binom{k}{2}+14\binom{k}{3}+39\binom{k}{4}+59\binom{k}{5}+63\binom{k}{6}+45\binom{k}{7}+22\binom{k}{8}+7\binom{k}{9}+\binom{k}{10}\ }
\]
\[
\boxed{\ A_3(k)=4\binom{k}{3}+30\binom{k}{4}+119\binom{k}{5}+273\binom{k}{6}+439\binom{k}{7}+507\binom{k}{8}+437\binom{k}{9}+283\binom{k}{10}+135\binom{k}{11}+46\binom{k}{12}+10\binom{k}{13}+\binom{k}{14}\ }.
\]
上述二项式基系数可由 forward-difference（Newton 级数）直接回收，因而给出无需谱计算的整数证书；
并可由脚本 \texttt{scripts/exp\_fold\_zm\_hard\_boundary\_coeffs\_audit.py} 一键复核。
\end{conclusion}

\begin{conclusion}[硬边界系数的精确多项式化：奇长 \texorpdfstring{$m=2k+1$}{m=2k+1} 的层次数严格量子化为 \texorpdfstring{$4j$}{4j}]\label{con:xi-terminal-zm-hard-boundary-odd-polynomial}
沿用结论 \ref{con:xi-terminal-zm-order4-rational} 的 \(Z_m(y)\) 记号。
对固定整数 \(j\ge 0\)，定义奇长硬边界系数
\[
B_j(k):=[y^{k+1-j}]\,Z_{2k+1}(y)\qquad(k\ge 0),
\]
其中约定当 \(k<j-1\) 时 \(B_j(k)=0\)。
则 \(B_j(k)\) 在所有整数 \(k\ge 0\) 上为精确多项式，其次数恰为 \(4j\)，且首项系数严格为 \(1/(4j)!\)。
并且存在唯一的二项式基展开
\[
B_j(k)=\sum_{r=j-1}^{4j}d_{j,r}\binom{k}{r},\qquad d_{j,r}\in\ZZ,\ \ d_{j,4j}=1,
\]
其中约定 \(\binom{k}{r}=0\)（当 \(r>k\)）。
在 \(j=1,2\) 时可给出完全闭式：
\[
\boxed{\ B_1(k)=\binom{k}{0}+4\binom{k}{1}+4\binom{k}{2}+3\binom{k}{3}+\binom{k}{4}
=\frac{k^4+6k^3+23k^2+66k+24}{24}\ }
\]
\[
\boxed{\ B_2(k)=2\binom{k}{1}+11\binom{k}{2}+23\binom{k}{3}+32\binom{k}{4}+28\binom{k}{5}+16\binom{k}{6}+6\binom{k}{7}+\binom{k}{8}\ }.
\]
\end{conclusion}

\begin{conclusion}[硬边界大偏差的闭式前因子：边界层概率质量的严格多项式—指数分解]\label{con:xi-terminal-zm-hard-boundary-ldp-prefactor}
令 \(w_m(x)=d_m(x)/2^m\) 为均匀微态推前分布，并令 \(X\sim w_m\)、\(H_m:=|X|_1\)。
对固定整数 \(j\ge 1\)，存在严格闭式：
\begin{itemize}
  \item 偶长 \(m=2k\)：
  \[
  \mathbb{P}(H_{2k}=k-j)=\frac{A_j(k)}{4^k},
  \]
  其中 \(A_j(k)\) 为结论 \ref{con:xi-terminal-zm-hard-boundary-even-polynomial} 的次数 \(4j+2\) 多项式。
  因而
  \[
  \mathbb{P}(H_{2k}=k-j)=\frac{k^{4j+2}}{(4j+2)!\,4^k}\Bigl(1+O(k^{-1})\Bigr).
  \]
  \item 奇长 \(m=2k+1\)：
  \[
  \mathbb{P}(H_{2k+1}=k+1-j)=\frac{B_j(k)}{2^{2k+1}},
  \]
  其中 \(B_j(k)\) 为结论 \ref{con:xi-terminal-zm-hard-boundary-odd-polynomial} 的次数 \(4j\) 多项式。
  因而
  \[
  \mathbb{P}(H_{2k+1}=k+1-j)=\frac{k^{4j}}{(4j)!\,2^{2k+1}}\Bigl(1+O(k^{-1})\Bigr).
  \]
\end{itemize}
上述分解表明：在最大密度边界附近，推前重量分布的指数速率保持为纯几何项 \(4^{-k}\)（或 \(2^{-(2k+1)}\)），
而全部非平凡算术结构被压缩到一个阶乘归一化的多项式前因子之中。
\end{conclusion}

\begin{conclusion}[硬边界层生成函数的纯极点结构：极点阶恰为 \texorpdfstring{$4j+3$}{4j+3} 与 \texorpdfstring{$4j+1$}{4j+1}]\label{con:xi-terminal-zm-hard-boundary-pure-pole}
令 \(A_j(k)\)、\(B_j(k)\) 如结论 \ref{con:xi-terminal-zm-hard-boundary-even-polynomial} 与结论 \ref{con:xi-terminal-zm-hard-boundary-odd-polynomial}，
并记其普通生成函数
\[
\mathcal A_j(t):=\sum_{k\ge 0}A_j(k)t^k,\qquad \mathcal B_j(t):=\sum_{k\ge 0}B_j(k)t^k.
\]
则二者均为有理函数，且极点阶严格量子化为
\[
\mathcal A_j(t)=\frac{P_j(t)}{(1-t)^{4j+3}},\qquad
\mathcal B_j(t)=\frac{Q_j(t)}{(1-t)^{4j+1}},
\]
其中 \(P_j,Q_j\in\ZZ[t]\)。
例如对 \(j=1\)（偶长第一硬边界层）有完全闭式
\[
\boxed{\ \mathcal A_1(t)=\frac{2t-5t^2+9t^3-10t^4+7t^5-2t^6}{(1-t)^7}\ }.
\]
因此硬边界层不仅在系数层面为精确多项式，其整生成函数亦呈现单一纯极点的最简解析结构。
\end{conclusion}


\begin{theorem}[Stokes 幅值三次与 Lee--Yang 三次的共同二维 Artin 表示]\label{thm:xi-terminal-zm-stokes-leyang-shared-artin-representation}
沿用定理 \ref{thm:xi-terminal-zm-b2-kappa2-rational-interconstruction} 的记号，记
\[
b:=B^2,
\qquad
u:=\kappa_*^2.
\]
再沿用结论 \ref{con:xi-terminal-zm-leyang-cubic-field-isomorphism} 的生成元 \(\alpha,\beta\)。
则有共同三次数域识别
\[
\boxed{\;
\QQ(b)=\QQ(u)=\QQ(\lambda_*)=\QQ(\alpha)=\QQ(\beta)=:K.\;}
\]
令 \(L\) 为 \(K/\QQ\) 的伽罗瓦闭包。
记 \(\rho_B\) 为由 Stokes 幅值三次
\[
162b^3+1593b^2+1998b+1184=0
\]
所定义的二维 Artin 表示；等价地，\(\rho_B\) 是 \(\Gal(L/\QQ)\cong S_3\) 的标准二维不可约表示沿自然商映射
\(\Gal(\overline\QQ/\QQ)\twoheadrightarrow \Gal(L/\QQ)\)
的拉回。
再记 \(\rho_{\mathrm{LY}}\) 为定理 \ref{thm:xi-time-part9c-leyang-cubic-artin-weight1} 中附着于 Lee--Yang 三次数域 \(K\) 的二维 Artin 表示。
则
\[
\boxed{\;\rho_B\cong \rho_{\mathrm{LY}}.\;}
\]
特别地，
\[
L(s,\rho_B)=L(s,\rho_{\mathrm{LY}}),
\qquad
\det(\rho_B)=\det(\rho_{\mathrm{LY}})=\chi_{-111},
\]
其中 \(\chi_{-111}\) 对应二次分辨子域 \(\QQ(\sqrt{-111})\)。
\leanverified{paper\_xi\_terminal\_zm\_stokes\_leyang\_shared\_artin\_representation}
\end{theorem}
\begin{proof}
由定理 \ref{thm:xi-terminal-zm-b2-kappa2-rational-interconstruction}，
\[
\QQ(b)=\QQ(u)=\QQ(\lambda_*).
\]
又由定理 \ref{thm:xi-terminal-zm-stokes-t-elliptic-branch-coordinate}，\(\lambda_*\) 满足
\[
16\lambda_*^3-9\lambda_*^2+1=0.
\]
结合结论 \ref{con:xi-terminal-zm-leyang-cubic-field-isomorphism} 的
\[
\QQ(\alpha)=\QQ(\beta),
\qquad
16\alpha^3-9\alpha^2+1=0,
\]
可知 \(\QQ(\lambda_*)=\QQ(\alpha)=\QQ(\beta)\)。
遂得第一条共同域识别。

再由定理 \ref{thm:xi-terminal-zm-stokes-amplitude-square-cubic-rigidity} 与定理 \ref{thm:xi-time-part9c-leyang-cubic-artin-weight1}，
Stokes 幅值三次与 Lee--Yang 三次都给出同一非伽罗瓦三次域 \(K\)，其伽罗瓦闭包均为 \(L\)，并满足
\[
\Gal(L/\QQ)\cong S_3.
\]
而 \(S_3\) 的二维不可约复表示只有一个同构类，即标准表示。
因此由 Stokes 幅值三次得到的 \(\rho_B\) 与 Lee--Yang 三次得到的 \(\rho_{\mathrm{LY}}\) 必同构：
\[
\rho_B\cong \rho_{\mathrm{LY}}.
\]

最后，定理 \ref{thm:xi-time-part9c-leyang-cubic-artin-weight1} 已给出
\[
\det(\rho_{\mathrm{LY}})=\chi_{-111}.
\]
另一方面，由推论 \ref{cor:xi-terminal-zm-stokes-b2-quadratic-resolvent}，Stokes 幅值三次的唯一二次分辨子域同样是 \(\QQ(\sqrt{-111})\)，故
\[
\det(\rho_B)=\chi_{-111}.
\]
表示同构即推出 Artin \(L\)-函数相同，得到
\[
L(s,\rho_B)=L(s,\rho_{\mathrm{LY}}).
\]
证毕。
\end{proof}

\begin{corollary}[Lee--Yang 三次数域的有理观测原始元刚性]\label{cor:xi-terminal-zm-leyang-cubic-rational-observable-primitivity}
\leanverified{paper\_xi\_terminal\_zm\_leyang\_cubic\_rational\_observable\_primitivity}
沿用定理 \ref{thm:xi-terminal-zm-stokes-leyang-shared-artin-representation} 的记号，令
\[
K:=\QQ(u)=\QQ(b)=\QQ(\lambda_*).
\]
则 \([K:\QQ]=3\)。若 \(R(T)\in\QQ(T)\) 在 \(u\) 处有定义且 \(R(u)\notin\QQ\)，则 \(R(u)\) 已为 \(K\) 的原始元，即
\[
\QQ(R(u))=K.
\]
同样地，只要 \(R(b)\) 或 \(R(\lambda_*)\) 有定义且不属于 \(\QQ\)，则相应元素亦为 \(K\) 的原始元，即
\[
\QQ(R(b))=K,
\qquad
\QQ(R(\lambda_*))=K.
\]
\end{corollary}
\begin{proof}
由定理 \ref{thm:xi-terminal-zm-kappa-square-cubic-field-s3}，\(u\) 满足不可约三次关系
\[
324u^3-540u^2+333u-74=0,
\]
故 \([\QQ(u):\QQ]=3\)。
再由定理 \ref{thm:xi-terminal-zm-stokes-leyang-shared-artin-representation} 的共同域识别，得 \([K:\QQ]=3\)。

若 \(R(u)\notin\QQ\)，则
\[
\QQ\subsetneq \QQ(R(u))\subseteq K.
\]
由于 \([K:\QQ]=3\) 为素数，\(K\) 不存在非平凡真中间域，故必有 \(\QQ(R(u))=K\)。
对 \(b\) 与 \(\lambda_*\) 的情形完全同理，因为 \(\QQ(b)=\QQ(\lambda_*)=K\)。证毕。
\end{proof}

\leanverified{paper\_xi\_terminal\_zm\_stokes\_leyang\_common\_quadratic\_resolvent}
\begin{corollary}[Stokes 幅值三次与 Lee--Yang 三次的共同二次分辨子域]\label{cor:xi-terminal-zm-stokes-leyang-common-quadratic-resolvent}
沿用上述记号，并令 \(L\) 为三次数域 \(K/\QQ\) 的伽罗瓦闭包。则 \(L/\QQ\) 的唯一二次子域为
\[
\boxed{\ \QQ(\sqrt{-111})\ }.
\]
等价地，\(u\) 与 \(b\) 的三次最小多项式
\[
324u^3-540u^2+333u-74,
\qquad
162b^3+1593b^2+1998b+1184
\]
的判别式分别为
\[
-2^6\cdot 3^9\cdot 37,
\qquad
-2^2\cdot 3^9\cdot 11^2\cdot 37\cdot 109^2,
\]
其平方自由部分同为
\[
-111=-3\cdot 37.
\]
\end{corollary}
\begin{proof}
由定理 \ref{thm:xi-terminal-zm-kappa-square-cubic-field-s3}，\(u\) 的极小多项式判别式为
\[
-2^6\cdot 3^9\cdot 37,
\]
故其分裂域的唯一二次子域为 \(\QQ(\sqrt{-111})\)。
另一方面，由推论 \ref{cor:xi-terminal-zm-stokes-b2-quadratic-resolvent}，\(b\) 的三次模型所对应的唯一二次分辨子域同样是 \(\QQ(\sqrt{-111})\)。

又由定理 \ref{thm:xi-terminal-zm-stokes-leyang-shared-artin-representation}，\(u\) 与 \(b\) 生成同一三次数域 \(K\)，故其伽罗瓦闭包一致，均等于 \(L\)。
因此 \(L/\QQ\) 的唯一二次子域只能是上述共同的 \(\QQ(\sqrt{-111})\)。
最后一条仅是把两判别式去除平方因子后的直接改写。证毕。
\end{proof}

\endinput

\paragraph{Lee--Yang 谱曲线在 \texorpdfstring{$u$}{u}-坐标下的统一化：分歧剖分、纤维群律与模性均匀化}
沿用
\[
\Pi(\lambda,y)=\lambda^{4}-\lambda^{3}-(2y+1)\lambda^{2}+\lambda+y(y+1),
\]
并记
\[
\Disc_{\lambda}\Pi=-y(y-1)\bigl(256y^{3}+411y^{2}+165y+32\bigr).
\]

\begin{conclusion}[Lee--Yang 谱曲线的 \texorpdfstring{$u$}{u}-坐标椭圆化与导子 \texorpdfstring{$37$}{37} 标识]\label{con:xi-terminal-zm-leyang-u-coordinate-birational-37a1}
设仿射曲线 \(C\subset\mathbb A^2_{\lambda,y}\) 由 \(\Pi(\lambda,y)=0\) 定义，并令
\[
u:=2y+1-2\lambda^{2}.
\]
则 \((\lambda,y)\mapsto(\lambda,u)\) 在 \(\QQ\) 上给出双有理同构
\[
C\ \dashrightarrow\ E_u:\quad u^{2}=4\lambda^{3}-4\lambda+1,
\]
其逆映射为
\[
(\lambda,u)\longmapsto\left(\lambda,\frac{u-1+2\lambda^2}{2}\right).
\]
进一步设 \(v:=(u-1)/2\)，则
\[
E_u\cong E_0:\quad v^2+v=\lambda^3-\lambda.
\]
因此该谱曲线归一化与导子 \(37\) 的模椭圆曲线同构（见定理 \ref{thm:xi-terminal-zm-leyang-elliptic-modularity-37}）。
\end{conclusion}
\begin{proof}[证明要点]
将 \(\Pi\) 视为关于 \(y\) 的二次式：
\[
\Pi(\lambda,y)=y^2+(1-2\lambda^2)y+\bigl(\lambda^4-\lambda^3-\lambda^2+\lambda\bigr).
\]
乘以 \(4\) 并配方，得
\[
(2y+1-2\lambda^2)^2=(1-2\lambda^2)^2-4(\lambda^4-\lambda^3-\lambda^2+\lambda)=1+4\lambda^3-4\lambda,
\]
即 \(u^2=4\lambda^3-4\lambda+1\)。由 \(u=2y+1-2\lambda^2\) 立得逆式。再令 \(u=2v+1\) 即得 \(v^2+v=\lambda^3-\lambda\)。证毕。
\end{proof}

\begin{conclusion}[分歧剖分与积分坐标下的临界五次方程]\label{con:xi-terminal-zm-leyang-u-coordinate-ramification-integral-quintic}
在结论 \ref{con:xi-terminal-zm-leyang-u-coordinate-birational-37a1} 的同构下，
\[
y=\frac{u-1+2\lambda^2}{2},\qquad y:E_u\to\PP^1
\]
为次数 \(4\) 的态射。其分歧满足：
\begin{enumerate}
  \item \(y\) 在无穷远点 \(O\) 处有唯一四阶极点，故 \(O\) 为 \(y=\infty\) 上的全分歧点（分歧指数 \(4\)）；
  \item 有限分歧值集合恰为
  \[
  \{0,1\}\cup\{\,\alpha\in\CC:\ 256\alpha^3+411\alpha^2+165\alpha+32=0\,\},
  \]
  且每个有限分歧值均对应简单分歧（分歧指数 \(2\)）；
  \item 在积分 Weierstrass 坐标
  \[
  X:=4\lambda,\qquad Y:=4u,\qquad Y^2=X^3-16X+16,
  \]
  有限临界点满足
  \[
  2XY+3X^2-16=0,
  \]
  消去 \(Y\) 得
  \[
  4X^5-9X^4-64X^3+160X^2-256=0.
  \]
\end{enumerate}
\end{conclusion}
\begin{proof}[证明要点]
由 \(y=(u-1+2\lambda^2)/2\) 得
\[
dy=\frac{du+4\lambda\,d\lambda}{2}.
\]
而 \(u^2=4\lambda^3-4\lambda+1\) 微分为 \(2u\,du=(12\lambda^2-4)\,d\lambda\)，故
\[
dy=\left(2\lambda+\frac{3\lambda^2-1}{u}\right)d\lambda.
\]
有限临界点由 \(dy=0\) 给出，即
\[
2\lambda u+3\lambda^2-1=0.
\]
换元 \(X=4\lambda,\ Y=4u\) 即得 \(2XY+3X^2-16=0\)。与 \(Y^2=X^3-16X+16\) 消元得到所示五次方程。
分歧值集合由 \(\Disc_\lambda\Pi=-y(y-1)P_{\mathrm{LY}}(y)\) 直接刻画，且有限分歧点总贡献与 Riemann--Hurwitz 公式一致，故均为简单分歧。证毕。
\end{proof}

\begin{conclusion}[非分歧纤维四点零和]\label{con:xi-terminal-zm-leyang-u-coordinate-fiber-fourpoint-zero-sum}
取 \(y_0\in\CC\) 避开分歧值集合，记纤维
\[
y^{-1}(y_0)=\{P_1,P_2,P_3,P_4\}\subset E_u(\CC).
\]
则椭圆曲线群律中恒有
\[
P_1+P_2+P_3+P_4=O.
\]
\end{conclusion}
\begin{proof}[证明要点]
定义
\[
f_{y_0}:=u+2\lambda^2-(2y_0+1).
\]
其零点恰为纤维四点，且在 \(O\) 处有四阶极点，从而
\[
\div(f_{y_0})=P_1+P_2+P_3+P_4-4O.
\]
主除子在 \(\operatorname{Pic}^0(E_u)\cong E_u\) 中对应零元，即得结论。
该陈述与结论 \ref{con:xi-terminal-zm-elliptic-weight-fiber-sum-zero} 一致。证毕。
\end{proof}

\begin{conclusion}[Néron 微分与导子 \texorpdfstring{$37$}{37} 的模性实现]\label{con:xi-terminal-zm-leyang-u-coordinate-neron-modular-pullback}
在 \(E_u:u^2=4\lambda^3-4\lambda+1\) 上定义
\[
\omega:=\frac{d\lambda}{u}.
\]
则 \(\omega\) 为 \(E_u\) 的 Néron 不变微分。并且存在权 \(2\)、水平 \(37\) 的归一化新形式
\(f_{37}\in S_2(\Gamma_0(37))\) 及模参数化
\[
\varphi:X_0(37)\to E_u
\]
使得
\[
\varphi^*\omega=c_\varphi\,2\pi i\,f_{37}(\tau)\,d\tau,\qquad c_\varphi\in\ZZ_{>0}.
\]
对最优参数化可取 \(c_\varphi=1\)。
\end{conclusion}
\begin{proof}[证明要点]
\(u^2=4\lambda^3-4\lambda+1\) 为 Weierstrass 形式，\(d\lambda/u\) 即其标准不变微分。
由模性定理，导子 \(37\) 的椭圆曲线对应 \(S_2(\Gamma_0(37))\) 新形式，且模参数化拉回 Néron 微分为新形式乘 \(2\pi i\,d\tau\)（带 Manin 常数）。证毕。
\end{proof}

\begin{conclusion}[Weierstrass 均匀化与 Lee--Yang 分支的椭圆函数参数化]\label{con:xi-terminal-zm-leyang-u-coordinate-weierstrass-uniformization}
存在格 \(\Lambda\subset\CC\) 使
\[
\lambda(z)=\wp(z;\,g_2=4,g_3=-1),\qquad
u(z)=\wp'(z;\,g_2=4,g_3=-1),
\]
并定义
\[
y(z)=\frac{\wp'(z)-1+2\wp(z)^2}{2}.
\]
则对一切 \(z\) 恒有
\[
\Pi(\lambda(z),y(z))=0.
\]
因此 \(\lambda\) 作为 \(y\) 的代数多值反演可在 \(\CC/\Lambda\) 上由平移作用实现单值化；临界值集合由 \(dy(z)=0\) 给出，恰为
\[
y\in\{0,1\}\cup\{P_{\mathrm{LY}}(y)=0\}.
\]
\end{conclusion}
\begin{proof}[证明要点]
由 \((\wp')^2=4\wp^3-4\wp+1\) 立得 \((\lambda,u)\in E_u\)，再代入
\(
y=(u-1+2\lambda^2)/2
\)
即得 \(\Pi(\lambda(z),y(z))=0\)。
临界值陈述由结论 \ref{con:xi-terminal-zm-leyang-u-coordinate-ramification-integral-quintic} 的 \(dy=0\) 判据与判别式分解一致性得到。证毕。
\end{proof}

\paragraph{\texorpdfstring{$S_4$}{S4} 闭包下的 Chevalley--Weil 重数刚性、Jacobian 群代数分解与局部算术证书}
沿用结论 \ref{con:xi-terminal-zm-s4-galois-closure-genus-tower-holomorphic-differentials} 的设定，
记 \(\widetilde{X}\to \PP^1_y\) 为四次谱覆盖 \(\Pi(\lambda,y)=0\) 的 \(S_4\)-伽罗瓦闭包曲线，且 \(g(\widetilde{X})=16\)。
并记
\[
F(\lambda,y):=\Pi(\lambda,y)
=\lambda^{4}-\lambda^{3}-(2y+1)\lambda^{2}+\lambda+y(y+1).
\]
记
\[
V:=H^0(\widetilde{X},\Omega^1_{\widetilde{X}}),
\qquad
\varepsilon:=\mathrm{sgn},\ \rho_2:=V_2,\ \rho_3:=V_3,\ \rho_3':=V_3'.
\]

\begin{conclusion}[\texorpdfstring{$S_4$}{S4} 闭包覆盖的 Hurwitz 型与子群商亏格谱]\label{con:xi-terminal-zm-s4-hurwitz-subgroup-spectrum-cycletypes}
设 \(\widetilde{X}\to\PP^1_y\) 为 \(S_4\)-伽罗瓦覆盖，其分歧惯性阶为
\[
(4,2,2,2,2,2),
\]
即 Hurwitz 型 \((0;4,2,2,2,2,2)\)。由 Riemann--Hurwitz 公式可得 \(g(\widetilde{X})=16\)，并且
\[
g(\widetilde{X}/S_3)=1,\quad
g(\widetilde{X}/A_4)=2,\quad
g(\widetilde{X}/V_4^{\mathrm{norm}})=4,\quad
g(\widetilde{X}/D_8)=1,\quad
g(\widetilde{X}/C_3)=6.
\]
进一步设 \(\tau=(12)\)、\(\rho=(1234)\in S_4\)。
对任意子群 \(H\le S_4\)，记
\[
d:=[S_4:H],\qquad
c_\tau(H):=\#\text{\(\tau\) 在 }S_4/H\text{ 上的循环个数},\qquad
c_\rho(H):=\#\text{\(\rho\) 在 }S_4/H\text{ 上的循环个数}.
\]
则有统一亏格公式
\[
g(\widetilde{X}/H)
=1-d+\frac12\Bigl(5\bigl(d-c_\tau(H)\bigr)+\bigl(d-c_\rho(H)\bigr)\Bigr).
\]
并且在各共轭类代表子群上，可得完整枚举
\[
\begin{array}{c|c|c|c|c}
H & d & \tau\text{ 在 }S_4/H\text{ 上的循环型} & \rho\text{ 在 }S_4/H\text{ 上的循环型} & g(\widetilde{X}/H)\\
\hline
S_4 & 1 & 1 & 1 & 0\\
A_4 & 2 & 2 & 2 & 2\\
S_3 & 4 & 2\cdot 1^2 & 4 & 1\\
D_8 & 3 & 2\cdot 1 & 2\cdot 1 & 1\\
V_4^{\mathrm{norm}} & 6 & 2^3 & 2^3 & 4\\
V_4^{\mathrm{pair}}=\langle(12),(34)\rangle & 6 & 2^2\cdot 1^2 & 4\cdot 2 & 2\\
C_4 & 6 & 2^3 & 4\cdot 1^2 & 4\\
C_3 & 8 & 2^4 & 4^2 & 6\\
C_2^{\mathrm{tr}} & 12 & 2^5\cdot 1^2 & 4^3 & 6\\
C_2^{\mathrm{dt}} & 12 & 2^6 & 4^2\cdot 2^2 & 8\\
\{1\} & 24 & 2^{12} & 4^6 & 16
\end{array}
\]
其中 \(D_8\) 表示阶 \(8\) 二面体子群，\(V_4^{\mathrm{norm}}\) 为正规 Klein 四元，\(V_4^{\mathrm{pair}}\) 为两两换位 Klein 四元。
\end{conclusion}
\begin{proof}[证明要点]
对 \(\widetilde{X}\to\PP^1_y\) 应用 Riemann--Hurwitz：
\[
2g(\widetilde{X})-2
=|S_4|(-2)+|S_4|\Bigl(1-\frac14\Bigr)+5|S_4|\Bigl(1-\frac12\Bigr)=30,
\]
故 \(g(\widetilde{X})=16\)。

对任意 \(H\le S_4\)，商覆盖 \(\widetilde{X}/H\to\PP^1_y\) 的总分歧由 \(\tau,\rho\) 在 \(S_4/H\) 上的循环分解给出，分别为 \(d-c_\tau(H)\)、\(d-c_\rho(H)\)。
代入 Riemann--Hurwitz 即得统一亏格公式与上表。证毕。
\end{proof}

\begin{conclusion}[指数 \texorpdfstring{$6$}{6} 配对稳定子商曲线的显式模型]\label{con:xi-terminal-zm-s4-v4pair-explicit-model}
在代数闭包中记 \(F(\lambda,y)=0\) 的四根为 \(\lambda_1,\lambda_2,\lambda_3,\lambda_4\)，并取
\[
V_4^{\mathrm{pair}}=\langle(12),(34)\rangle\le S_4.
\]
定义
\[
v:=\lambda_1\lambda_2,\quad w:=\lambda_3\lambda_4,\quad z:=v+w,\quad s:=v-w.
\]
则固定域可写为
\[
\QQ(\widetilde{X}/V_4^{\mathrm{pair}})=\QQ(y,z,s),
\]
并满足完全显式关系
\[
\mathscr R(z,y)=0,\qquad s^2=z^2-4y(y+1),
\]
其中
\[
\mathscr R(z,y)=z^3+(1+2y)z^2-(1+4y+4y^2)z-(1+5y+13y^2+8y^3).
\]
记 \(E_{\mathrm{res}}\) 为 \(\mathscr R(z,y)=0\) 的光滑射影模型，\(C_{\mathrm{pair}}\) 为上述二元方程组的光滑射影模型，则
\[
\QQ(C_{\mathrm{pair}})=\QQ(\widetilde{X}/V_4^{\mathrm{pair}}),\qquad
g(C_{\mathrm{pair}})=2,
\]
且自然投影 \(C_{\mathrm{pair}}\to E_{\mathrm{res}}\) 为次数 \(2\) 覆盖。
\end{conclusion}
\begin{proof}[证明要点]
\(z=v+w\) 为配对 \(\{12|34\}\) 的不变量，故满足三次预解式 \(\mathscr R(z,y)=0\)。
又由 \(vw=\lambda_1\lambda_2\lambda_3\lambda_4=y(y+1)\)，得
\[
s^2=(v-w)^2=(v+w)^2-4vw=z^2-4y(y+1).
\]
于是 \(\QQ(y,z,s)\subseteq \QQ(\widetilde{X}/V_4^{\mathrm{pair}})\)。
反向由 \(v=(z+s)/2,\ w=(z-s)/2\) 可恢复配对乘积信息，故函数域相等。
亏格与覆盖次数由结论 \ref{con:xi-terminal-zm-s4-hurwitz-subgroup-spectrum-cycletypes} 中 \(H=V_4^{\mathrm{pair}}\) 的条目给出。证毕。
\end{proof}

\begin{conclusion}[黄金比率分歧：\texorpdfstring{$C_{\mathrm{pair}}\to E_{\mathrm{res}}$}{Cpair to Eres} 的分歧除子由 \texorpdfstring{$y^2+y-1$}{y2+y-1} 唯一刻画]\label{con:xi-terminal-zm-s4-v4pair-golden-ratio-branch-divisor}
在结论 \ref{con:xi-terminal-zm-s4-v4pair-explicit-model} 的记号下，令
\[
f:=z^2-4y(y+1)\in \QQ(E_{\mathrm{res}})^\times.
\]
则有严格消去恒等式
\[
\operatorname{Res}_z\!\bigl(\mathscr R(z,y),\,z^2-4y(y+1)\bigr)=(y^2+y-1)^2.
\]
因此 \(f=0\) 与 \(\mathscr R=0\) 的公共点恰为
\[
P_\pm:\quad
(y,z)=\Bigl(-\frac12\pm\frac{\sqrt5}{2},\,-2\Bigr),
\]
并且二次覆盖 \(C_{\mathrm{pair}}\to E_{\mathrm{res}}\) 的分歧除子为 \(P_++P_-\)。
从而
\[
2g(C_{\mathrm{pair}})-2=\deg(P_++P_-)=2,\qquad g(C_{\mathrm{pair}})=2,
\]
且分歧点由唯一二次域 \(\QQ(\sqrt5)\) 控制。
\end{conclusion}
\begin{proof}[证明要点]
结果式计算给出仅当 \(y^2+y-1=0\) 时可能有公共点。
将该二次方程两根代回 \(\mathscr R(z,y)=0\) 与 \(z^2-4y(y+1)=0\) 可唯一得到 \(z=-2\)。
对二次覆盖 \(s^2=f\)，分歧由 \(f\) 的奇阶零点给出，故分歧除子正为 \(P_++P_-\)。
再由 Riemann--Hurwitz 公式即得亏格。证毕。
\end{proof}

\begin{conclusion}[Chevalley--Weil 型 \texorpdfstring{$S_4$}{S4} 表示分解的显式重数解]\label{con:xi-terminal-zm-s4-chevalley-weil-explicit-decomposition}
在上述记号下，\(V\) 作为 \(S_4\)-表示满足
\[
V\cong 2\varepsilon\oplus \rho_2\oplus \rho_3\oplus 3\rho_3'.
\]
等价地，\(\dim V=16\) 的不可约重数向量唯一为
\[
(m_{\varepsilon},m_{\rho_2},m_{\rho_3},m_{\rho_3'})=(2,1,1,3).
\]
\end{conclusion}
\begin{proof}[证明要点]
由结论 \ref{con:xi-terminal-zm-s4-hurwitz-subgroup-spectrum-cycletypes}，
\[
g(\widetilde{X}/S_4)=0,\quad
g(\widetilde{X}/A_4)=2,\quad
g(\widetilde{X}/S_3)=1,\quad
g(\widetilde{X}/D_8)=1.
\]
对任意 \(K\le S_4\) 有
\[
\dim V^K=g(\widetilde{X}/K).
\]
设
\[
V\cong a\mathbf 1\oplus b\varepsilon\oplus c\rho_2\oplus d\rho_3\oplus e\rho_3'.
\]
由 \(\dim V^{S_4}=0\) 得 \(a=0\)。

再由角色平均公式可得：
在 \(A_4\) 上仅 \(\varepsilon\) 贡献一维不变元，故 \(b=2\)；
在 \(S_3\) 上仅 \(\rho_3\) 贡献一维不变元，故 \(d=1\)；
在 \(D_8\) 上仅 \(\rho_2\) 贡献一维不变元，故 \(c=1\)。

最后由总维数
\[
16=\dim V=2b+2c+3d+3e
=2\cdot2+2\cdot1+3\cdot1+3e
\]
解得 \(e=3\)。
因此
\[
V\cong 2\varepsilon\oplus \rho_2\oplus \rho_3\oplus 3\rho_3',
\]
且分解唯一。证毕。
\end{proof}

\begin{conclusion}[Jacobian 的 \texorpdfstring{$S_4$}{S4} 等变同源分解与可识别等型因子]\label{con:xi-terminal-zm-s4-jacobian-group-algebra-prym-e3-b3}
记
\[
J:=J(\widetilde{X}),\qquad
J_{\mathrm{sgn}}:=J(\widetilde{X}/A_4),\qquad
E:=\widetilde{X}/S_3,\qquad
E_{\mathrm{res}}:=\widetilde{X}/D_8.
\]
则存在定义在 \(\QQ\) 上的三维阿贝尔簇 \(A,B\)，使得
\[
J\sim_{\QQ}J_{\mathrm{sgn}}\times J_{\rho_2}\times A\times B^3,
\]
其中
\[
J_{\mathrm{sgn}}\sim_{\QQ}J(\widetilde{X}/A_4),\qquad
J_{\rho_2}\sim_{\QQ}J(\widetilde{X}/D_8)^2=E_{\mathrm{res}}^2,
\]
并且 \(A\) 对应 \(\rho_3\)-等型（可取 \(A\sim_{\QQ}E^3\)），\(B\) 对应 \(\rho_3'\)-等型。

等价地，可写为
\[
J\sim_{\QQ}J(\widetilde{X}/A_4)\times \operatorname{Prym}\!\bigl((\widetilde{X}/V_4^{\mathrm{norm}})/(\widetilde{X}/A_4)\bigr)\times E^{3}\times B^{3}.
\]
并且有中间 Jacobian 分拆
\[
J(\widetilde{X}/V_4^{\mathrm{norm}})\sim_{\QQ}J_{\mathrm{sgn}}\times E_{\mathrm{res}}^{2},
\]
\[
J(C_{\mathrm{pair}})\sim_{\QQ}E_{\mathrm{res}}\times E,\qquad
C_{\mathrm{pair}}:=\widetilde{X}/V_4^{\mathrm{pair}},
\]
\[
J(\widetilde{X}/C_4)\sim_{\QQ}E_{\mathrm{res}}\times B.
\]
\end{conclusion}
\begin{proof}[证明要点]
由
\[
\QQ[S_4]\cong \QQ\oplus \QQ\oplus M_2(\QQ)\oplus M_3(\QQ)\oplus M_3(\QQ)
\]
与结论 \ref{con:xi-terminal-zm-s4-chevalley-weil-explicit-decomposition}，
\(J\) 在 \(\QQ\)-同源意义下按 \(\varepsilon,\rho_2,\rho_3,\rho_3'\) 四个等型分解，维数分别为 \(2,2,3,9\)。

\(\varepsilon\)-等型由 \(A_4\) 不变微分识别为 \(J(\widetilde{X}/A_4)\)；
\(\rho_2\)-等型由 \(D_8\) 商与结论 \ref{con:xi-terminal-zm-s4-prym-c6-over-c-d8-elliptic-square} 识别为 \(E_{\mathrm{res}}^2\)；
\(\rho_3\)-等型可由 \(S_3\) 商识别为 \(E^3\)；
\(\rho_3'\)-等型重数为 \(3\)，记为 \(B^3\)。

对中间子群 \(V_4^{\mathrm{norm}},V_4^{\mathrm{pair}},C_4\)，利用
\[
H^0(\widetilde{X}/H,\Omega^1)\cong H^0(\widetilde{X},\Omega^1)^H
\]
与角色平均公式，结合结论 \ref{con:xi-terminal-zm-s4-hurwitz-subgroup-spectrum-cycletypes} 的亏格数据，即得所列三条中间 Jacobian 同源分拆。证毕。
\end{proof}

\begin{conjecture}[\texorpdfstring{$\rho_3'$}{rho3'}-三维因子 \texorpdfstring{$B$}{B} 的绝对单纯性与端同态刚性]\label{conj:xi-terminal-zm-s4-b-absolute-simplicity-endomorphism}
在结论 \ref{con:xi-terminal-zm-s4-jacobian-group-algebra-prym-e3-b3} 的记号下，
设 \(B\) 为 \(\rho_3'\)-等型三维构件。
则 \(B\) 在 \(\overline{\QQ}\) 上绝对单纯，且
\[
\End^0_{\overline{\QQ}}(B)\cong \QQ.
\]
等价地，除由 \(\QQ[S_4]\) 的中心化子所强制的结构外，\(B\) 不携带额外的代数自同态。
\end{conjecture}

\begin{conclusion}[\texorpdfstring{$D_8$}{D8} 商椭圆曲线与二维 Prym 的椭圆平方化]\label{con:xi-terminal-zm-s4-prym-c6-over-c-d8-elliptic-square}
设 \(D_8\le S_4\) 为任一 \(4\)-循环的正规化子（阶 \(8\) 二面体子群），并定义
\[
E_8:=\widetilde{X}/D_8.
\]
并记
\[
C:=\widetilde{X}/A_4,\qquad C_6:=\widetilde{X}/V_4^{\mathrm{norm}}.
\]
则 \(g(E_8)=1\)，并且
\[
\operatorname{Prym}(C_6/C)\sim_{\QQ}E_8^2.
\]
从而
\[
\mathrm{End}^0_{\QQ}\!\bigl(\operatorname{Prym}(C_6/C)\bigr)\supset M_2(\QQ).
\]
在无分歧循环三次覆盖情形下，\(\operatorname{Prym}(C_6/C)\) 的极化型为 \((1,3)\)。
\end{conclusion}
\begin{proof}[证明要点]
由结论 \ref{con:xi-terminal-zm-s4-galois-closure-genus-tower-holomorphic-differentials}，\(g(\widetilde{X}/D_8)=1\)。
由结论 \ref{con:xi-terminal-zm-s4-jacobian-a2-isog-r2}，\(\rho_2\)-等型分量满足 \(A_2\sim R^2\)，其中 \(R=\widetilde{X}/D_8\)。
另一方面，定理 \ref{thm:xi-terminal-zm-s3-closure-jacobian-splitting} 识别 \(A_2\sim \operatorname{Prym}(C_6/C)\)。
合并即得 \(\operatorname{Prym}(C_6/C)\sim_{\QQ}E_8^2\)。
极化型由定理 \ref{thm:xi-terminal-zm-prym-polarization-a2-rigidity} 给出。证毕。
\end{proof}

\begin{conclusion}[好素数处的局部 \texorpdfstring{$L$}{L} 因子强分解与模 \texorpdfstring{$3$}{3} 点数同余证书]\label{con:xi-terminal-zm-s4-local-factorization-mod3-congruence-certificate}
设 \(p\) 为 \(\widetilde{X}\) 的好约化素数，并令 \(q=p^m\)（\(m\ge 1\)）。
沿用结论 \ref{con:xi-terminal-zm-s4-jacobian-group-algebra-prym-e3-b3} 与
\ref{con:xi-terminal-zm-s4-prym-c6-over-c-d8-elliptic-square} 的记号，并记
\[
C:=\widetilde{X}/A_4.
\]
则局部因子满足
\[
P_p(\widetilde{X},T)
=P_p(C,T)\,P_p(E,T)^3\,P_p(E_8,T)^2\,P_p(B,T)^3.
\]
进一步，对有限域点数记号 \(N_Y(q):=\#Y(\FF_q)\)，有同余
\[
N_{\widetilde{X}}(q)\equiv N_C(q)+2N_{E_8}(q)-2(q+1)\pmod 3.
\]
\end{conclusion}
\begin{proof}[证明要点]
\(\QQ\)-同源分解在好约化处与 \(\ell\)-进 \(H^1\) 分解相容，故 Frobenius 特征多项式按因子相乘，得到 \(P_p\) 强分解。
又
\[
N_Y(q)=q+1-\mathrm{Tr}\!\left(\mathrm{Frob}_q\mid H^1_{\acute et}(Y_{\overline{\FF}_q},\QQ_\ell)\right).
\]
由分解式可得
\[
\mathrm{Tr}_q(\widetilde{X})
\equiv \mathrm{Tr}_q(C)+2\mathrm{Tr}_q(E_8)\pmod 3,
\]
因为 \(3\mathrm{Tr}_q(E)\) 与 \(3\mathrm{Tr}_q(B)\) 在模 \(3\) 下消失。
代回点数公式即得所示同余。证毕。
\end{proof}

\paragraph{\texorpdfstring{$S_4$}{S4} 作用曲线的低次数映射、典范二次关系与 \texorpdfstring{$n$}{n}-典范表示}
以下统一记 \(X:=\widetilde{X}\)。
则 \(X/S_4\simeq \PP^1_y\) 的分支型为一个阶 \(4\) 分支点与五个阶 \(2\) 分支点，且 \(g(X)=16\)。

\begin{conclusion}[Castelnuovo--Severi 阻断下的低次数可约性消失]\label{con:xi-terminal-zm-s4-low-gonality-vanishing-castelnuovo-severi}
曲线 \(X\) 不存在次数为 \(2\) 或 \(3\) 的非常值态射 \(X\to \PP^1\)。
因此 \(X\) 非双曲线且非三折曲线。
\end{conclusion}
\begin{proof}[证明要点]
由结论 \ref{con:xi-terminal-zm-s4-galois-closure-genus-tower-holomorphic-differentials}，
\(E:=X/S_3\) 为亏格 \(1\) 曲线，且自然商映射 \(\pi:X\to E\) 的次数为 \(6\)。
设 \(\varphi:X\to \PP^1\) 为次数 \(d\in\{2,3\}\) 的非常值态射。
若 \(\varphi=\psi\circ\pi\)，则 \(d=6\deg(\psi)\)，矛盾。
故 \(\pi,\varphi\) 在 Castelnuovo--Severi 意义下相互独立，满足
\[
g(X)\le d\,g(E)+6\,g(\PP^1)+(d-1)(6-1)=6d-5.
\]
当 \(d=2\) 时右端为 \(7\)，当 \(d=3\) 时右端为 \(13\)，均与 \(g(X)=16\) 矛盾。
\end{proof}

\begin{conclusion}[典范模型的射影正规性与二次生成]\label{con:xi-terminal-zm-s4-canonical-projective-normality-quadratic-generation}
典范嵌入 \(X\hookrightarrow \PP^{15}\) 射影正规，且其齐次理想由二次方程生成。
等价地，
\[
\mu_2:\operatorname{Sym}^2 H^0(X,K_X)\twoheadrightarrow H^0(X,2K_X)
\]
为满射，并且
\[
\dim\ker(\mu_2)=91.
\]
\end{conclusion}
\begin{proof}[证明要点]
由结论 \ref{con:xi-terminal-zm-s4-low-gonality-vanishing-castelnuovo-severi}，\(X\) 非双曲线。
Max Noether 定理给出典范嵌入的射影正规性，特别地 \(\mu_2\) 满射。
又由同一结论知 \(X\) 非三折；且 \(g(X)=16\neq 6\)，故 \(X\) 不可能为平面五次曲线。
由 Petri 定理可得典范理想由二次方程生成。
再由 Riemann--Roch，
\[
h^0(X,2K_X)=3g-3=45,\qquad
\dim\operatorname{Sym}^2 H^0(X,K_X)=\binom{16+1}{2}=136,
\]
故 \(\dim\ker(\mu_2)=136-45=91\)。
\end{proof}

\begin{conclusion}[\texorpdfstring{$S_4$}{S4} 作用在 \texorpdfstring{$X=\widetilde{X}$}{X} 上的固定点数完全表]\label{con:xi-terminal-zm-s4-conjugacy-fixedpoints-local-eigenvalue-full-count}
设 \(g\in S_4\)。
则 \(|\mathrm{Fix}(g)|\) 仅依赖于 \(g\) 的共轭类，并满足
\[
|\mathrm{Fix}(g)|=
\begin{cases}
10,& g\text{ 为换位 }(2),\\
2,& g\text{ 为双换位 }(2,2),\\
0,& g\text{ 为三循环 }(3),\\
2,& g\text{ 为四循环 }(4).
\end{cases}
\]
此外，阶 \(2\) 元在定点处切空间特征值为 \(-1\)；
阶 \(4\) 元的两个定点切空间特征值分别为 \(i\) 与 \(-i\)。
\end{conclusion}
\begin{proof}[证明要点]
由结论 \ref{con:xi-terminal-zm-s4-hurwitz-subgroup-spectrum-cycletypes}，\(X\to\PP^1_y\) 的惯性只可能为换位（五个有限分歧值）或四循环（无穷远分歧值）。
若 \(g\) 为三循环，则其阶为 \(3\)，不可能包含于任一惯性子群，故无固定点。

若 \(g\) 为换位，则固定点仅可能来自五个有限分歧纤维。每个此类分歧值上，惯性群同构于 \(\langle\tau\rangle\cong C_2\)。给定换位 \(g\) 的固定点数为
\[
\frac{|C_{S_4}(g)|}{|\langle\tau\rangle|}=\frac{4}{2}=2,
\]
故总数为 \(5\cdot2=10\)。

若 \(g\) 为四循环，则固定点仅来自无穷远分歧纤维。设其惯性为 \(I=\langle c\rangle\cong C_4\)。固定点条件等价于 \(g\in hIh^{-1}\)，即 \(hch^{-1}=g\) 或 \(g^{-1}\)，从而
\[
|\mathrm{Fix}(g)|=\frac{|C_{S_4}(c)|+|C_{S_4}(c)|}{|I|}=\frac{4+4}{4}=2.
\]

若 \(g\) 为双换位，则只能来自无穷远纤维。在 \(I=\langle c\rangle\) 中唯一非平凡二阶元为 \(c^2\)，故
\[
|\mathrm{Fix}(g)|=\frac{|C_{S_4}(c^2)|}{|I|}=\frac{8}{4}=2.
\]
局部特征值方面，阶 \(2\) 元在定点处导数为 \(-1\)。
对阶 \(4\) 元，由结论 \ref{con:xi-terminal-zm-s4-chevalley-weil-explicit-decomposition} 的
\(\mathrm{tr}((1234)\mid H^0(X,K_X))=0\)。
并与全纯 Lefschetz 公式
\[
\mathrm{tr}\!\bigl(g\mid H^0(X,K_X)\bigr)
=1+\sum_{P\in \mathrm{Fix}(g)}\frac{\zeta_P}{1-\zeta_P},
\]
可知两定点贡献相消，故特征值为 \(i\) 与 \(-i\)。证毕。
\end{proof}

\begin{conclusion}[全体 \texorpdfstring{$n$}{n}-典范表示的显式不可约分解]\label{con:xi-terminal-zm-s4-pluricanonical-irrep-decomposition-all-n}
设 \(n\ge 2\)。
记 \(\mathbf 1,\varepsilon,\rho_2,\rho_3,\rho_3'\) 分别为 \(S_4\) 的平凡、符号、二维不可约、三维标准及其符号扭曲表示。
则 \(H^0(X,nK_X)\) 的不可约分解如下：
\[
\text{若 }n\equiv 0\!\!\pmod 4,\quad
H^0(X,nK_X)\cong
\frac{5n+4}{4}\mathbf 1\oplus
\frac{5n-8}{4}\varepsilon\oplus
\frac{5n-2}{2}\rho_2\oplus
\frac{15n-4}{4}\rho_3\oplus
\frac{15n-12}{4}\rho_3'.
\]
\[
\text{若 }n\equiv 1\!\!\pmod 4,\quad
H^0(X,nK_X)\cong
\frac{5n-9}{4}\mathbf 1\oplus
\frac{5n+3}{4}\varepsilon\oplus
\frac{5n-3}{2}\rho_2\oplus
\frac{15n-11}{4}\rho_3\oplus
\frac{15n-3}{4}\rho_3'.
\]
\[
\text{若 }n\equiv 2\!\!\pmod 4,\quad
H^0(X,nK_X)\cong
\frac{5n+2}{4}\mathbf 1\oplus
\frac{5n-6}{4}\varepsilon\oplus
\frac{5n-2}{2}\rho_2\oplus
\frac{15n-2}{4}\rho_3\oplus
\frac{15n-14}{4}\rho_3'.
\]
\[
\text{若 }n\equiv 3\!\!\pmod 4,\quad
H^0(X,nK_X)\cong
\frac{5n-7}{4}\mathbf 1\oplus
\frac{5n+1}{4}\varepsilon\oplus
\frac{5n-3}{2}\rho_2\oplus
\frac{15n-13}{4}\rho_3\oplus
\frac{15n-1}{4}\rho_3'.
\]
特别地，平凡表示重数满足
\[
\dim H^0(X,nK_X)^{S_4}
=1-2n+5\Bigl\lfloor\frac n2\Bigr\rfloor+\Bigl\lfloor\frac{3n}{4}\Bigr\rfloor.
\]
\end{conclusion}
\begin{proof}[证明要点]
当 \(n\ge 2\) 时 \(H^1(X,nK_X)=0\)。
由结论 \ref{con:xi-terminal-zm-s4-conjugacy-fixedpoints-local-eigenvalue-full-count} 的定点与局部特征值数据，全纯 Lefschetz 公式给出
\[
\mathrm{tr}\!\bigl((12)\mid H^0(X,nK_X)\bigr)=5(-1)^n,\qquad
\mathrm{tr}\!\bigl((12)(34)\mid H^0(X,nK_X)\bigr)=(-1)^n,
\]
\[
\mathrm{tr}\!\bigl((123)\mid H^0(X,nK_X)\bigr)=0,
\]
以及
\[
\mathrm{tr}\!\bigl((1234)\mid H^0(X,nK_X)\bigr)=
\begin{cases}
1,& n\equiv 0,3\pmod 4,\\
-1,& n\equiv 1,2\pmod 4.
\end{cases}
\]
另一方面
\[
\mathrm{tr}\!\bigl(1\mid H^0(X,nK_X)\bigr)=h^0(X,nK_X)=30n-15.
\]
将上述五个共轭类迹代入 \(S_4\) 特征标正交关系，逐一解得各不可约重数，即得到四个同余类分解公式。
平凡重数的地板函数表达式由同一解在平凡特征标分量上的化简得到。
\end{proof}

\begin{conclusion}[二次典范像的 \texorpdfstring{$S_4$}{S4} 等变二次关系分解]\label{con:xi-terminal-zm-s4-canonical-quadrics-equivariant-decomposition}
记
\[
V:=H^0(X,K_X),\qquad
Q:=\ker\!\Bigl(\operatorname{Sym}^2V\xrightarrow{\mu_2}H^0(X,2K_X)\Bigr).
\]
则：
\begin{enumerate}
  \item
  \[
  H^0(X,2K_X)\cong
  3\mathbf 1\oplus \varepsilon\oplus 4\rho_2\oplus 7\rho_3\oplus 4\rho_3';
  \]
  \item
  \[
  Q\cong
  8\mathbf 1\oplus 2\varepsilon\oplus 9\rho_2\oplus 13\rho_3\oplus 8\rho_3'.
  \]
\end{enumerate}
特别地，\(\dim Q^{S_4}=8\)。
\end{conclusion}
\begin{proof}[证明要点]
第 \(1\) 点是结论 \ref{con:xi-terminal-zm-s4-pluricanonical-irrep-decomposition-all-n} 在 \(n=2\) 时的特例。
由结论 \ref{con:xi-terminal-zm-s4-canonical-projective-normality-quadratic-generation}，
\[
0\to Q\to \operatorname{Sym}^2V\xrightarrow{\mu_2}H^0(X,2K_X)\to 0
\]
为 \(S_4\)-等变正合列。
又由结论 \ref{con:xi-terminal-zm-s4-chevalley-weil-explicit-decomposition} 中
\(V\cong 2\varepsilon\oplus\rho_2\oplus\rho_3\oplus 3\rho_3'\)，
可算得
\[
\operatorname{Sym}^2V\cong
11\mathbf 1\oplus 3\varepsilon\oplus 13\rho_2\oplus 20\rho_3\oplus 12\rho_3'.
\]
与第 \(1\) 点逐项相减即得第 \(2\) 点，平凡重数遂为 \(8\)。
\end{proof}

\paragraph{四次谱方程在固定记号下的 \texorpdfstring{$V_4$}{V4} 分辨率与闭包几何刚性链}
以下记号在本节固定：
\[
\Pi(\lambda,y)=\lambda^4-\lambda^3-(2y+1)\lambda^2+\lambda+y(y+1)\in \QQ(y)[\lambda],
\]
\[
C_{\mathrm{LY}}(y):=256y^3+411y^2+165y+32,\qquad
\Delta(y):=\Disc_{\lambda}\Pi(\lambda,y)=-y(y-1)C_{\mathrm{LY}}(y)\in\QQ[y].
\]
并沿用文内既有中间商记号
\[
X:=\widetilde X,\qquad C:=X/A_4,\qquad C_6:=X/V_4^{\mathrm{norm}}.
\]

\begin{conclusion}[\texorpdfstring{$V_4$}{V4} 分辨率三次多项式与判别式的完全一致]\label{con:xi-terminal-zm-v4-resolvent-discriminant-complete-coincidence}
可取 \(\Pi(\lambda,y)\) 的 \(V_4\)-分辨率三次多项式为
\[
R(u,y)=u^3+(2y+1)u^2-(2y+1)^2u-(8y^3+13y^2+5y+1)\in\QQ(y)[u],
\]
并满足严格恒等式
\[
\Disc_u R(u,y)=\Disc_{\lambda}\Pi(\lambda,y)=\Delta(y)=-y(y-1)C_{\mathrm{LY}}(y).
\]
\end{conclusion}
\begin{proof}[证明要点]
一般首一四次多项式
\(
f(x)=x^4+ax^3+bx^2+cx+d
\)
的分辨率三次可写为
\[
u^3-bu^2+(ac-4d)u+(4bd-a^2d-c^2).
\]
代入
\(
a=-1,\ b=-(2y+1),\ c=1,\ d=y(y+1)
\)
即得上式 \(R(u,y)\)。
判别式恒等式与 \(\Disc_{\lambda}\Pi\) 的一致性已在结论 \ref{con:xi-terminal-zm-generic-galois-s4} 给出，故结论成立。证毕。
\end{proof}

\begin{conclusion}[\texorpdfstring{$S_4$}{S4} 闭包曲线的分歧护照与中间商亏格谱]\label{con:xi-terminal-zm-s4-closure-passport-genus-spectrum-restated}
设 \(L/\QQ(y)\) 为 \(\Pi(\lambda,y)\) 的分裂域，\(X\) 为对应光滑射影曲线。则：
\begin{enumerate}
  \item \(\Gal(L/\QQ(y))\cong S_4\)；
  \item 覆盖 \(X\to\PP^1_y\) 的分歧惯性阶为 \((2,2,2,2,2,4)\)；
  \item \(g(X)=16\)；
  \item 曲线 \(C=X/A_4\) 亏格为 \(2\)，并可取模型
  \[
  C:\ s^2=\Delta(y)=-y(y-1)C_{\mathrm{LY}}(y);
  \]
  \item 曲线 \(C_6=X/V_4^{\mathrm{norm}}\) 满足 \(g(C_6)=4\)，且
  \[
  C_6\to \PP^1_y
  \]
  为 \(S_3\)-Galois 覆盖，六个分歧值处惯性均为换位（阶 \(2\)）。
  此外，\(C_6\) 可由
  \[
  \begin{cases}
  s^2=\Delta(y),\\
  R(u,y)=0
  \end{cases}
  \]
  的仿射簇取规范化得到。
\end{enumerate}
\end{conclusion}
\begin{proof}[证明要点]
由结论 \ref{con:xi-terminal-zm-generic-galois-s4} 得到 \(S_4\) 泛型伽罗瓦群与判别式核；
由结论 \ref{con:xi-terminal-zm-s4-galois-closure-genus-tower-holomorphic-differentials} 得到分歧护照 \((4)\cdot(2)^5\) 与 \(g(X)=16\)，并给出 \(X/A_4\) 与 \(X/V_4^{\mathrm{norm}}\) 的亏格分别为 \(2,4\)。
由 \(S_4/V_4^{\mathrm{norm}}\cong S_3\) 及结论 \ref{con:xi-terminal-zm-discriminant-ridge-etale-c3-explicit-model} 的显式模型可知，\(C_6\) 的规范化表示与所述方程组一致。证毕。
\end{proof}

\begin{conclusion}[无分歧循环三次覆盖与 \texorpdfstring{$3$}{3}-挠类]\label{con:xi-terminal-zm-c6-over-c-etale-c3-pic3-class}
存在定义在 \(\QQ\) 上的无分歧循环三次覆盖
\[
\pi:C_6\longrightarrow C
\]
且 \(\Gal(C_6/C)\cong C_3\)。
该覆盖对应于 \(\Pic^0(C)\) 的非平凡三挠类，并等价给出 \(J(C)\) 上 \(\QQ\)-定义的阶 \(3\) 循环子群方案与次数 \(3\) 同源。
\end{conclusion}
\begin{proof}[证明要点]
结论 \ref{con:xi-terminal-zm-discriminant-ridge-etale-c3-explicit-model} 已给出 \(\pi\) 的循环三次与处处非分歧性。
对光滑射影曲线，循环 \'{e}tale \(3\)-覆盖与 \(\Pic^0[3]\) 的非平凡类对应；据此得到三挠类与相应 \(3\)-同源。
该下降与核的 \(\QQ\)-定义性可由结论 \ref{con:xi-terminal-zm-discriminant-ridge-3isogeny-descent} 对照验证。证毕。
\end{proof}

\begin{conclusion}[判别式二次脊 Jacobian 的有理挠子下界]\label{con:xi-terminal-zm-discriminant-ridge-jacobian-rational-torsion-lower-bound-2x2x3}
曲线
\(
C:\ s^2=-y(y-1)C_{\mathrm{LY}}(y)
\)
的 Jacobian 满足
\[
(\ZZ/2\ZZ)^2\hookrightarrow J(C)(\QQ)_{\mathrm{tors}}.
\]
并且由三次无分歧覆盖诱导的 \(3\)-同源核提供一个 \(\QQ\)-定义阶 \(3\) 循环子群方案；
在该核分裂为常值群方案的情形下，进一步得到
\[
(\ZZ/2\ZZ)^2\times(\ZZ/3\ZZ)\hookrightarrow J(C)(\QQ)_{\mathrm{tors}}.
\]
\end{conclusion}
\begin{proof}[证明要点]
由结论 \ref{con:xi-terminal-zm-discriminant-ridge-jacobian-torsion-upper-bound}，
\[
D_0=[(0,0)-\infty],\qquad D_1=[(1,0)-\infty]
\]
给出两条线性无关的有理 \(2\)-挠元，故有 \((\ZZ/2\ZZ)^2\) 注入。
另一方面，由结论 \ref{con:xi-terminal-zm-c6-over-c-etale-c3-pic3-class} 得到阶 \(3\) 循环子群方案与度 \(3\) 同源；
若该子群在 \(\QQ\) 上常值分裂，则可取非零有理 \(3\)-挠点，与前述 \(2\)-挠按互素阶直积分解，得到所示注入。证毕。
\end{proof}

\begin{conclusion}[三阶 Sylow 自由作用与亏格 \(6\) 商曲线]\label{con:xi-terminal-zm-sylow3-free-action-genus6-quotient}
设 \(P\le S_4\) 为任一三阶 Sylow 子群（\(P\cong C_3\)）。
则 \(P\) 在 \(X\) 上无不动点；故
\[
X\to X_P:=X/P
\]
为三重 \'{e}tale 覆盖，且
\[
g(X_P)=6.
\]
\end{conclusion}
\begin{proof}[证明要点]
由结论 \ref{con:xi-terminal-zm-s4-closure-passport-genus-spectrum-restated}，\(X\to\PP^1_y\) 的惯性阶仅为 \(2\) 与 \(4\)，故任一点稳定子阶整除 \(4\)。
因此三阶元素不可能固定点，\(P\) 作用自由。
应用 Riemann--Hurwitz：
\[
2g(X)-2=|P|\bigl(2g(X_P)-2\bigr),
\]
代入 \(g(X)=16\)、\(|P|=3\) 得 \(g(X_P)=6\)。证毕。
\end{proof}

\begin{conclusion}[全纯微分空间的 \texorpdfstring{$S_4$}{S4} 表示分解刚性]\label{con:xi-terminal-zm-holomorphic-differential-s4-decomposition-rigid-restated}
记
\[
W:=H^0(X,\Omega^1_X).
\]
则 \(W\) 作为 \(\QQ[S_4]\)-模唯一分解为
\[
W\cong 2\varepsilon\oplus V_2\oplus V_3\oplus 3V_3',
\]
其中 \(\varepsilon\) 为符号表示，\(V_2\) 为二维不可约表示，\(V_3\) 为标准三维表示，\(V_3'=V_3\otimes\varepsilon\)。
\end{conclusion}
\begin{proof}[证明要点]
对任意 \(H\le S_4\) 有
\(
W^H\simeq H^0(X/H,\Omega^1)
\)，故 \(\dim W^H=g(X/H)\)。
由结论 \ref{con:xi-terminal-zm-s4-galois-closure-genus-tower-holomorphic-differentials}，
\[
\dim W^{S_4}=0,\quad
\dim W^{A_4}=2,\quad
\dim W^{V_4^{\mathrm{norm}}}=4,\quad
\dim W^{S_3}=1,
\]
并且 \(\dim W=g(X)=16\)。
将 \(S_4\) 五个不可约表示的子群不变维数代入该线性系统，得到唯一解即所示分解。证毕。
\end{proof}

\begin{conclusion}[Jacobian 的群代数同源分解与 \texorpdfstring{$(1,3)$}{(1,3)}-极化 Prym 因子]\label{con:xi-terminal-zm-jacobian-group-algebra-prym13-visible-factor-restated}
存在 \(\QQ\)-同源分解
\[
J(X)\sim_{\QQ}J(C)\times P\times A_3\times A_9,
\]
其中：
\begin{enumerate}
  \item \(P=\operatorname{Prym}(C_6/C)\) 为二维阿贝尔簇，极化型为 \((1,3)\)；
  \item \(A_3\) 与 \(A_9\) 分别对应 \(V_3\) 与 \(3V_3'\) 的等型部分（\(\dim A_3=3,\ \dim A_9=9\)）；
  \item deck 群 \(C_3\) 在 \(P\) 上诱导
  \[
  \QQ(\zeta_3)\hookrightarrow \End^0(P).
  \]
\end{enumerate}
\end{conclusion}
\begin{proof}[证明要点]
由结论 \ref{con:xi-terminal-zm-s4-jacobian-group-algebra-prym-e3-b3}，
\[
J(X)\sim_{\QQ}J(C)\times E_{\mathrm{res}}^2\times E^3\times B^3.
\]
由定理 \ref{thm:xi-terminal-zm-s3-closure-jacobian-splitting} 与定理 \ref{thm:xi-terminal-zm-prym-polarization-a2-rigidity}，
\[
P=\operatorname{Prym}(C_6/C)\sim_{\QQ}E_{\mathrm{res}}^2,\qquad \mathrm{type}(P)=(1,3).
\]
令 \(A_3:=E^3,\ A_9:=B^3\) 即得所示分解。
又 \(C_3\) 对覆盖 \(C_6\to C\) 的作用在 Prym 部分给出群代数作用
\(
\QQ[C_3]\cong \QQ\times\QQ(\zeta_3)
\)，故得到 \(\QQ(\zeta_3)\subset \End^0(P)\)。证毕。
\end{proof}

\paragraph{判别式二次分辨率与 \texorpdfstring{$C_3$}{C3} 商几何：纤维积规范化、Prym 因子与固定子商 Jacobian 分解}
沿用结论 \ref{con:xi-terminal-zm-s4-galois-closure-genus-tower-holomorphic-differentials} 的设定，记
\[
\widetilde{C}:=\widetilde{X},\qquad
L:=\QQ(\widetilde{C}),\qquad
G:=\mathrm{Gal}(L/\QQ(y))\cong S_4.
\]
并记
\[
E:=\widetilde{C}/S_3,\qquad
H_{\mathrm{LY}}:=\widetilde{C}/A_4,\qquad
C_3:=S_3\cap A_4.
\]

\begin{conclusion}[判别式二次分辨率的几何合成：\texorpdfstring{$C_3$}{C3} 商即椭圆谱曲线与 Lee--Yang 二次脊的规范化纤维积]\label{con:xi-terminal-zm-s4-c3-fiberproduct-normalization}
存在曲线同构
\[
\widetilde{C}/C_3\ \cong\ \mathrm{Norm}\!\bigl(E\times_{\PP^1_y}H_{\mathrm{LY}}\bigr),
\]
等价地，在函数域层面有
\[
\QQ(\widetilde{C}/C_3)=\QQ(E)\cdot \QQ(H_{\mathrm{LY}})\subset L.
\]
此外，令
\[
P_{\mathrm{LY}}(y):=256y^3+411y^2+165y+32,
\]
则由判别式恒等式 \(\mathrm{Disc}_{\lambda}(F)=-y(y-1)P_{\mathrm{LY}}(y)\) 可取
\[
\QQ(H_{\mathrm{LY}})=\QQ(y)\!\left(\sqrt{-y(y-1)P_{\mathrm{LY}}(y)}\right),
\]
从而
\[
\QQ(\widetilde{C}/C_3)=\QQ(E)\!\left(\sqrt{-y(y-1)P_{\mathrm{LY}}(y)}\right).
\]
\end{conclusion}
\begin{proof}[证明要点]
由伽罗瓦对应，
\[
\QQ(E)=L^{S_3},\qquad \QQ(H_{\mathrm{LY}})=L^{A_4}.
\]
由于 \(L/\QQ(y)\) 为 Galois 扩张，固定域合成满足
\[
L^{S_3}\,L^{A_4}=L^{S_3\cap A_4}=L^{C_3},
\]
即
\[
\QQ(\widetilde{C}/C_3)=\QQ(E)\cdot\QQ(H_{\mathrm{LY}}).
\]
规范化纤维积的函数域等于右侧合成域，故得到曲线同构。
二次分辨子域的显式式由结论 \ref{con:xi-terminal-zm-discriminant-leyang} 与
\ref{con:xi-terminal-zm-s4-galois-closure-genus-tower-holomorphic-differentials} 直接给出。证毕。
\end{proof}

\begin{conclusion}[\texorpdfstring{$C_3$}{C3} 商曲线 Jacobian 的三因子同源分解]\label{con:xi-terminal-zm-s4-c3-quotient-jacobian-three-factor-decomposition}
设
\[
C:=\widetilde{C}/C_3,\qquad
J_{\mathrm{LY}}:=\mathrm{Jac}(H_{\mathrm{LY}}).
\]
沿用结论 \ref{con:xi-terminal-zm-s4-jacobian-group-algebra-prym-e3-b3} 的 \(E\) 与三维阿贝尔簇 \(B\) 记号，则有
\[
\mathrm{Jac}(C)\sim_{\QQ}J_{\mathrm{LY}}\times E\times B.
\]
\end{conclusion}
\begin{proof}[证明要点]
由结论 \ref{con:xi-terminal-zm-s4-chevalley-weil-explicit-decomposition}，
\[
H^0(\widetilde{C},\Omega^1)\cong 2\varepsilon\oplus \rho_2\oplus \rho_3\oplus 3\rho_3'.
\]
对 \(C_3=\langle(123)\rangle\) 取不变子空间并用角色平均公式
\[
\dim W^{C_3}=\frac{1}{3}\bigl(\chi_W(1)+\chi_W((123))+\chi_W((132))\bigr),
\]
可得
\[
\dim(\varepsilon^{C_3})=1,\quad
\dim(\rho_2^{C_3})=0,\quad
\dim(\rho_3^{C_3})=1,\quad
\dim((\rho_3')^{C_3})=1.
\]
结合结论 \ref{con:xi-terminal-zm-s4-jacobian-group-algebra-prym-e3-b3} 的同源分解
\[
J(\widetilde{C})\sim_{\QQ}J_{\mathrm{LY}}\times E_{\mathrm{res}}^2\times E^3\times B^3,
\]
对幂等元 \(e_{C_3}=\frac1{3}\sum_{h\in C_3}h\) 取像后，分别截取到
\(J_{\mathrm{LY}}\)、\(E\)、\(B\) 各一次，且 \(E_{\mathrm{res}}^2\) 分量不贡献，遂得结论。证毕。
\end{proof}

\begin{conclusion}[Prym--Lee--Yang 等价：\texorpdfstring{$\mathrm{Prym}(C/E)\sim_{\QQ}J_{\mathrm{LY}}\times B$}{Prym(C/E) isogenous to JLY x B}]\label{con:xi-terminal-zm-s4-c3-prym-leyang-equivalence}
由包含 \(C_3\le S_3\) 诱导次数 \(2\) 映射
\[
\pi:C=\widetilde{C}/C_3\longrightarrow E=\widetilde{C}/S_3.
\]
则其 Prym 变体满足
\[
\mathrm{Prym}(C/E)\sim_{\QQ}J_{\mathrm{LY}}\times B.
\]
\end{conclusion}
\begin{proof}[证明要点]
对次数 \(2\) 覆盖 \(\pi:C\to E\)，有标准同源分解
\[
\mathrm{Jac}(C)\sim_{\QQ}\mathrm{Jac}(E)\times \mathrm{Prym}(C/E),
\]
且 \(\mathrm{Jac}(E)=E\)。
代入结论 \ref{con:xi-terminal-zm-s4-c3-quotient-jacobian-three-factor-decomposition} 的
\[
\mathrm{Jac}(C)\sim_{\QQ}J_{\mathrm{LY}}\times E\times B
\]
并消去 \(E\) 因子即得。证毕。
\end{proof}

\begin{conclusion}[Lee--Yang 非分歧六点的 Abel 和闭式与二次覆叠分歧除子的群和消去]\label{con:xi-terminal-zm-leyang-unramified-sixpoint-abel-closed-form-branch-balance}
仍在椭圆曲线
\[
E:\ Y^2=X^3-X+\frac14,\qquad
y=X^2+Y-\frac12
\]
上，记 \(P_{\mathrm{LY}}(y)=256y^3+411y^2+165y+32\)。
设 \(Q_1,Q_2,Q_3\) 为 Lee--Yang 三个分歧点（见结论 \ref{con:xi-terminal-zm-leyang-ramification-divisor-principalization-trace}），并定义有效除子 \(D_{\mathrm{unr}}\)（次数 \(6\)）满足
\[
\mathrm{div}_0\!\bigl(P_{\mathrm{LY}}(y)\bigr)=2(Q_1+Q_2+Q_3)+D_{\mathrm{unr}}.
\]
再记
\[
Q_0=\Bigl(1,-\frac12\Bigr),\qquad
Q_0'=\Bigl(-1,\frac12\Bigr)
\]
（分别位于 \(y=0,1\) 的有理分歧点）。
则在椭圆曲线群律下有
\[
\bigoplus_{P\in\mathrm{Supp}(D_{\mathrm{unr}})}P
=2(Q_0\oplus Q_0')
=\Bigl(\frac{161}{16},\,\frac{2033}{64}\Bigr)\in E(\QQ).
\]
进一步，考虑二次覆叠
\[
\pi:\ C\to E,\qquad C:\ w^2=-y(y-1)P_{\mathrm{LY}}(y),
\]
记其分歧除子为 \(B_{\pi}\)。
则
\[
\bigoplus_{P\in\mathrm{Supp}(B_{\pi})}P=O,
\]
等价地 \(B_{\pi}-10O\) 在 \(\mathrm{Pic}^0(E)\) 中为零类。
\end{conclusion}
\begin{proof}[证明要点]
由 \((y)_\infty=4O\) 可得
\[
\mathrm{div}\!\bigl(P_{\mathrm{LY}}(y)\bigr)=2(Q_1+Q_2+Q_3)+D_{\mathrm{unr}}-12O,
\]
从而
\[
2(Q_1\oplus Q_2\oplus Q_3)\oplus\!\!\bigoplus_{P\in\mathrm{Supp}(D_{\mathrm{unr}})}\!\!P=O.
\]
又由结论 \ref{con:xi-terminal-zm-leyang-ramification-divisor-principalization-trace}，
\[
Q_1\oplus Q_2\oplus Q_3=-(Q_0\oplus Q_0'),
\]
故得到
\[
\bigoplus_{P\in\mathrm{Supp}(D_{\mathrm{unr}})}P=2(Q_0\oplus Q_0').
\]
同一结论给出 \(Q_0\oplus Q_0'=(\tfrac14,\tfrac18)\)，对该点施行倍点公式即得
\((\tfrac{161}{16},\tfrac{2033}{64})\)。
\par
对分歧除子，设
\[
\mathrm{div}(y)=2Q_0+D_0-4O,\qquad
\mathrm{div}(y-1)=2Q_0'+D_1-4O,
\]
其中 \(D_0,D_1\) 分别是 \(y=0,1\) 纤维中的简单点部分（各次数 \(2\)）。
则 \(B_{\pi}=D_0+D_1+D_{\mathrm{unr}}\)。
由上述三条主除子关系在 \(\mathrm{Pic}^0(E)\cong E\) 下可得
\[
\bigoplus_{P\in\mathrm{Supp}(D_0)}P=-2Q_0,\qquad
\bigoplus_{P\in\mathrm{Supp}(D_1)}P=-2Q_0',\qquad
\bigoplus_{P\in\mathrm{Supp}(D_{\mathrm{unr}})}P=2(Q_0\oplus Q_0'),
\]
三式相加即得
\(
\bigoplus_{P\in\mathrm{Supp}(B_{\pi})}P=O
\)。
证毕。
\end{proof}

\begin{conclusion}[换位固定子商曲线的 Jacobian 分解]\label{con:xi-terminal-zm-s4-transposition-quotient-jacobian-full-decomposition}
取换位 \(\tau=(12)\in S_4\)，令
\[
C_{\tau}:=\widetilde{C}/\langle\tau\rangle.
\]
沿用 \(E_{\mathrm{res}}:=\widetilde{C}/D_8\)、\(E:=\widetilde{C}/S_3\) 与 \(B\) 的记号，则有
\[
\mathrm{Jac}(C_{\tau})\sim_{\QQ}E_{\mathrm{res}}\times E^2\times B.
\]
\end{conclusion}
\begin{proof}[证明要点]
由结论 \ref{con:xi-terminal-zm-s4-chevalley-weil-explicit-decomposition} 的表示分解，对 \(\langle\tau\rangle\) 取不变并用
\[
\dim W^{\langle\tau\rangle}=\frac12\bigl(\chi_W(1)+\chi_W(\tau)\bigr)
\]
得
\[
\dim(\varepsilon^{\langle\tau\rangle})=0,\quad
\dim(\rho_2^{\langle\tau\rangle})=1,\quad
\dim(\rho_3^{\langle\tau\rangle})=2,\quad
\dim((\rho_3')^{\langle\tau\rangle})=1.
\]
结合结论 \ref{con:xi-terminal-zm-s4-jacobian-group-algebra-prym-e3-b3} 的等型对应
\[
\varepsilon\leftrightarrow J_{\mathrm{LY}},\quad
\rho_2\leftrightarrow E_{\mathrm{res}}^2,\quad
\rho_3\leftrightarrow E^3,\quad
\rho_3'\leftrightarrow B^3,
\]
可知分别贡献 \(0,1,2,1\) 个因子，遂得所示同源分解。
其总维数为 \(1+2+3=6\)，与结论 \ref{con:xi-terminal-zm-s4-hurwitz-subgroup-spectrum-cycletypes} 的 \(g(\widetilde{C}/\langle\tau\rangle)=6\) 一致。证毕。
\end{proof}

\begin{conclusion}[双换位固定子商曲线的 Jacobian 分解]\label{con:xi-terminal-zm-s4-doubletransposition-quotient-jacobian-full-decomposition}
取双换位 \(\delta=(12)(34)\in S_4\)，令
\[
C_{\delta}:=\widetilde{C}/\langle\delta\rangle.
\]
则有
\[
\mathrm{Jac}(C_{\delta})\sim_{\QQ}J_{\mathrm{LY}}\times E_{\mathrm{res}}^2\times E\times B.
\]
\end{conclusion}
\begin{proof}[证明要点]
同理，
\[
\dim W^{\langle\delta\rangle}=\frac12\bigl(\chi_W(1)+\chi_W(\delta)\bigr),
\]
并由角色值计算得
\[
\dim(\varepsilon^{\langle\delta\rangle})=1,\quad
\dim(\rho_2^{\langle\delta\rangle})=2,\quad
\dim(\rho_3^{\langle\delta\rangle})=1,\quad
\dim((\rho_3')^{\langle\delta\rangle})=1.
\]
代入同一等型对应即得
\[
\mathrm{Jac}(C_{\delta})\sim_{\QQ}J_{\mathrm{LY}}\times E_{\mathrm{res}}^2\times E\times B.
\]
总维数 \(2+2+1+3=8\)，与
\(g(\widetilde{C}/\langle\delta\rangle)=8\)
一致。证毕。
\end{proof}

\begin{conclusion}[换位与双换位商的 Prym 极化型及 \texorpdfstring{$(\pm1)$}{(+/-1)}-本征同源分解]\label{con:xi-terminal-zm-s4-involution-prym-polarization-eigensplitting}
取换位 \(\tau\in S_4\) 与双换位 \(\delta\in S_4\)，并记
\[
C_{\tau}:=\widetilde{C}/\langle\tau\rangle,\qquad
C_{\delta}:=\widetilde{C}/\langle\delta\rangle,
\]
以及自然二次覆盖
\[
\pi_{\tau}:\widetilde{C}\to C_{\tau},\qquad
\pi_{\delta}:\widetilde{C}\to C_{\delta}.
\]
则：
\begin{enumerate}
  \item \(\pi_{\tau}\) 的分歧点数为 \(10\)，\(\dim \Prym(\pi_{\tau})=10\)，且由 \(\Jac(\widetilde{C})\) 的主极化在 \(\Prym(\pi_{\tau})\) 上诱导的极化型为
  \[
  (1^{6},2^{4}).
  \]
  并且在 \(\QQ\)-同源意义下有 \((\pm1)\)-本征分解
  \[
  \Jac(C_{\tau})\ \sim_{\QQ}\ E_{\mathrm{res}}\times E^2\times B,\qquad
  \Prym(\pi_{\tau})\ \sim_{\QQ}\ J_{\mathrm{LY}}\times E_{\mathrm{res}}\times E\times B^2.
  \]
  \item \(\pi_{\delta}\) 的分歧点数为 \(2\)，\(\dim \Prym(\pi_{\delta})=8\)，并且 \(\Prym(\pi_{\delta})\) 为主极化阿贝尔簇。
  在 \(\QQ\)-同源意义下有 \((\pm1)\)-本征分解
  \[
  \Jac(C_{\delta})\ \sim_{\QQ}\ J_{\mathrm{LY}}\times E_{\mathrm{res}}^2\times E\times B,\qquad
  \Prym(\pi_{\delta})\ \sim_{\QQ}\ E^2\times B^2.
  \]
\end{enumerate}
\end{conclusion}
\begin{proof}[证明要点]
由结论 \ref{con:xi-terminal-zm-s4-transposition-quotient-jacobian-full-decomposition} 与
\ref{con:xi-terminal-zm-s4-doubletransposition-quotient-jacobian-full-decomposition} 已知 \(g(C_{\tau})=6\)、\(g(C_{\delta})=8\)。
又由结论 \ref{con:xi-terminal-zm-s4-conjugacy-class-fixedpoints-and-cyclic-quotient-genera}，
\[
\#\mathrm{Fix}_{\widetilde C}(\tau)=10,\qquad \#\mathrm{Fix}_{\widetilde C}(\delta)=2,
\]
故两条二次覆盖的分歧点数分别为 \(10,2\)，并有
\[
\dim\Prym(\pi_{\tau})=g(\widetilde C)-g(C_{\tau})=16-6=10,\qquad
\dim\Prym(\pi_{\delta})=16-8=8.
\]
对分歧二次覆盖 \(\pi:C'\to C\)，Prym 极化型的标准公式给出
\(\mathrm{type}=(1^{g(C)},2^{r/2-1})\)，其中 \(r\) 为分歧点数。
据此代入 \((g,r)=(6,10)\) 得 \((1^6,2^4)\)，代入 \((g,r)=(8,2)\) 得主极化。
\par
同源分解由幂等元 \(e_{\pm}=(1\pm g)/2\)（\(g=\tau\) 或 \(\delta\)）作用于 \(\Jac(\widetilde C)\) 给出：
\(
e_+\Jac(\widetilde C)\sim_{\QQ}\Jac(\widetilde C/\langle g\rangle)
\)
且
\(
e_-\Jac(\widetilde C)\sim_{\QQ}\Prym(\pi_g)
\)。
将结论 \ref{con:xi-terminal-zm-s4-jacobian-group-algebra-prym-e3-b3} 的等型分解
\[
\Jac(\widetilde C)\sim_{\QQ}J_{\mathrm{LY}}\times E_{\mathrm{res}}^2\times E^3\times B^3
\]
与换位、双换位在各不可约表示上的特征值重数比较，即得所列 \((\pm1)\)-本征同源分解。证毕。
\end{proof}

\begin{conclusion}[四循环与 Klein 四群商的 Jacobian 纯分量分裂]\label{con:xi-terminal-zm-s4-c4-v4-quotient-jacobian-pure-splitting}
取无穷远惯性生成元 \(c\in S_4\)（四循环），并记
\[
C_{(4)}:=\widetilde C/\langle c\rangle,\qquad
C_{V_4}:=\widetilde C/V_4^{\mathrm{norm}}.
\]
则二者皆为属 \(4\) 曲线，且其 Jacobian 在 \(\QQ\)-同源意义下满足
\[
\Jac(C_{(4)})\ \sim_{\QQ}\ E_{\mathrm{res}}\times B,\qquad
\Jac(C_{V_4})\ \sim_{\QQ}\ J_{\mathrm{LY}}\times E_{\mathrm{res}}^2.
\]
\end{conclusion}
\begin{proof}[证明要点]
属数断言来自结论 \ref{con:xi-terminal-zm-s4-conjugacy-class-fixedpoints-and-cyclic-quotient-genera} 中
\(
g(\widetilde C/\langle c\rangle)=4
\)
及结论 \ref{con:xi-terminal-zm-s4-closure-passport-genus-spectrum-restated} 中
\(
g(\widetilde C/V_4^{\mathrm{norm}})=4
\)。
对 \(\Jac(C_{(4)})\) 的分裂由结论 \ref{con:xi-terminal-zm-s4-c4-prym-threefold-and-v3sgn-identification} 的 \(\Jac(\widetilde C/\langle c\rangle)\) 同源分解给出。
对 \(\Jac(C_{V_4})\) 的分裂由定理 \ref{thm:xi-terminal-zm-s3-closure-jacobian-splitting} 给出的
\(
\Jac(\widetilde C/V_4^{\mathrm{norm}})\sim_{\QQ}\Jac(\widetilde C/A_4)\times E_{\mathrm{res}}^2
\)
与 \(J_{\mathrm{LY}}=\Jac(\widetilde C/A_4)\)（结论 \ref{con:xi-terminal-zm-s4-c3-quotient-jacobian-three-factor-decomposition} 的记号）直接推出。证毕。
\end{proof}

\input{con__xi-terminal-zm-s4-a4-s3-branch-rigidity-standard-factor-fixedpoint}

\begin{conclusion}[Lee--Yang 判别式双覆盖在 \(\QQ(E)\) 上的线性化表达]\label{con:xi-terminal-zm-leyang-doublecover-linearization-a-plus-by}
在椭圆曲线
\[
E/\QQ:\qquad Y^2=X^3-X+\frac14
\]
上定义
\[
y:=X^2+Y-\frac12,\qquad
P(y):=256y^3+411y^2+165y+32.
\]
于函数域 \(\QQ(E)\) 中取 \(w\) 满足
\[
w^2=-y(y-1)P(y).
\]
则存在唯一 \(A(X),B(X)\in\ZZ[X]\)，使得
\[
w^2=A(X)+B(X)Y.
\]
可取显式系数
\[
\begin{aligned}
A(X)=\,&-256X^{10}-2560X^9-795X^8+5470X^7+2321X^6-3802X^5-1671X^4\\
&\quad+986X^3+423X^2-94X-22,
\end{aligned}
\]
\[
\begin{aligned}
B(X)=\,&-1280X^8-2560X^7+1684X^6+4500X^5-380X^4-2152X^3-68X^2+212X+44.
\end{aligned}
\]
并具有共同结构分解
\[
A(X)=-(X^2-1)A_8(X),\qquad
B(X)=-4(X^2-1)B_6(X),
\]
其中
\[
A_8(X)=256X^8+2560X^7+1051X^6-2910X^5-1270X^4+892X^3+401X^2-94X-22,
\]
\[
B_6(X)=320X^6+640X^5-101X^4-485X^3-6X^2+53X+11.
\]
\end{conclusion}
\begin{proof}[证明要点]
将 \(y=X^2+Y-\frac12\) 代入 \(-y(y-1)P(y)\) 并在 \(\QQ[X,Y]/(Y^2-X^3+X-\frac14)\) 中约化。
由于 \(\QQ(E)=\QQ(X)\oplus \QQ(X)\,Y\)（作为 \(\QQ(X)\)-向量空间），任意元素可唯一写为 \(A(X)+B(X)Y\) 形，故表示唯一。
直接展开并收集项即可得到上述 \(A,B\)。
再对 \(A,B\) 作整系数因式分解，得公共因子 \((X^2-1)\) 与给定 \(A_8,B_6\)。证毕。
\end{proof}

\begin{conclusion}[曲线 \(C\) 的平面四次模型与双二次规范方程]\label{con:xi-terminal-zm-leyang-doublecover-plane-quartic-model}
令 \(C\) 为 \(\QQ(C)=\QQ(E)(w)\) 的光滑射影模型。
定义
\[
f(X):=X^3-X+\frac14,\qquad D(X):=4X^3-4X+1=4f(X).
\]
则 \(C\) 在 \((X,w)\)-平面上满足显式四次方程
\[
\bigl(w^2-A(X)\bigr)^2=B(X)^2f(X).
\]
等价地，可写作整数系数规范式
\[
\bigl(w^2+(X^2-1)A_8(X)\bigr)^2
=
4(X^2-1)^2B_6(X)^2D(X).
\]
\end{conclusion}
\begin{proof}[证明要点]
由结论 \ref{con:xi-terminal-zm-leyang-doublecover-linearization-a-plus-by}
\[
w^2=A(X)+B(X)Y
\]
可得
\[
\bigl(w^2-A(X)\bigr)^2=B(X)^2Y^2=B(X)^2f(X).
\]
再代入
\(
A=-(X^2-1)A_8,\ B=-4(X^2-1)B_6,\ D=4f
\)
并约去常数因子即得第二式。证毕。
\end{proof}

\begin{conclusion}[四次投影关于 \(w\) 的判别式完全因式分解]\label{con:xi-terminal-zm-leyang-quartic-projection-discriminant-full-factorization}
设
\[
H(w;X):=\bigl(w^2-A(X)\bigr)^2-B(X)^2f(X)\in\ZZ[X][w].
\]
则关于 \(w\) 的判别式满足
\[
\Disc_w\!\bigl(H(w;X)\bigr)
=4096\cdot \Xi(X)\cdot D(X)^2\cdot g(X)^2\cdot Q_6(X)\cdot B_6(X)^4,
\]
其中
\[
\Xi(X):=X(X-2)(X-1)^7(X+1)^7,
\]
\[
g(X):=16X^3-9X^2+1,
\]
\[
Q_6(X):=256X^6-480X^5+201X^4+750X^3-921X^2-78X+704.
\]
因此，投影 \(\pi_X:C\to \PP^1_X\) 的有限分歧 \(X\)-坐标包含于显式集合
\[
\left\{0,2,\pm1\right\}
\cup\{D(X)=0\}
\cup\{g(X)=0\}
\cup\{Q_6(X)=0\}
\cup\{B_6(X)=0\}.
\]
\end{conclusion}
\begin{proof}[证明要点]
写
\[
H(w;X)=w^4-2A(X)w^2+\bigl(A(X)^2-B(X)^2f(X)\bigr),
\]
\[
\partial_w H=4w\bigl(w^2-A(X)\bigr).
\]
由单首多项式判别式公式
\[
\Disc_w(H)=\operatorname{Res}_w\!\bigl(H,\partial_wH\bigr)
=4^4\operatorname{Res}_w\!\bigl(H,w(w^2-A)\bigr)
=256\,H(0;X)\,\operatorname{Res}_w\!\bigl(H,w^2-A\bigr).
\]
再由
\[
H(0;X)=A^2-B^2f,\qquad
\operatorname{Res}_w\!\bigl(H,w^2-A\bigr)=\bigl(-B^2f\bigr)^2,
\]
得
\[
\Disc_w(H)=256\,(A^2-B^2f)\,B^4f^2.
\]
代入 \(B=-4(X^2-1)B_6,\ f=D/4\) 可化为
\[
\Disc_w(H)=4096\,(X^2-1)^4\,B_6^4\,D^2\,(A^2-B^2f).
\]
对最后因子作完全分解
\[
A^2-B^2f
=X(X-2)(X-1)^3(X+1)^3\,g(X)^2\,Q_6(X),
\]
合并即得所述闭式分解。分歧点由判别式零点给出。证毕。
\end{proof}

\begin{conclusion}[\(A_4\)-子扩张中的单点分歧机制]\label{con:xi-terminal-zm-s4-a4-subextension-single-point-ramification}
设 \(L/\QQ(y)\) 为四次谱方程 \(F(\lambda,y)=0\) 的分裂域，且
\[
\mathrm{Gal}(L/\QQ(y))\cong S_4,\qquad
\QQ(H_{\mathrm{LY}})=L^{A_4}
=\QQ(y)\!\left(\sqrt{-y(y-1)P(y)}\right).
\]
记 \(\widetilde{C}\) 为 \(\QQ(\widetilde{C})=L\) 的光滑射影模型，则 \(\widetilde{C}\to H_{\mathrm{LY}}\) 为 \(A_4\)-伽罗瓦覆盖。
设 \(\widetilde{C}\to\PP^1_y\) 的有限分歧惯性由换位 \(\tau\) 生成，无穷远处分歧惯性由四循环 \(\rho\) 生成，则：
\begin{enumerate}
  \item 有限分歧值处，惯性与 \(A_4\) 的交为空，故在 \(L/L^{A_4}\) 中全部消失；
  \item 唯一残存分歧位于 \(H_{\mathrm{LY}}\) 上方无穷远点 \(\infty_{\mathrm{LY}}\)，其惯性群
  \[
  I_{\infty_{\mathrm{LY}}}
  =\langle \rho^2\rangle
  \cong C_2,
  \qquad \rho^2=(13)(24)\in A_4.
  \]
\end{enumerate}
故 \(\widetilde{C}\to H_{\mathrm{LY}}\) 为几何单点分歧覆盖，总分歧度为 \(6\)。
\end{conclusion}
\begin{proof}[证明要点]
由 \(A_4\triangleleft S_4\) 正规性，\(L/L^{A_4}\) 在给定基点上的惯性群为原惯性群与 \(A_4\) 的交。
对有限分歧，惯性由换位生成，换位不属于 \(A_4\)，故交为空。
对无穷远处，\(\rho\notin A_4\) 但 \(\rho^2\in A_4\) 且阶为 \(2\)，故交为 \(\langle\rho^2\rangle\)。
又 \(H_{\mathrm{LY}}\to\PP^1_y\) 在无穷远处分歧，故其上方仅有一点 \(\infty_{\mathrm{LY}}\)。
于是仅此一点分歧，且
\[
\deg\!\bigl(\mathrm{Ram}(\widetilde{C}/H_{\mathrm{LY}})\bigr)
=|A_4|\!\left(1-\frac12\right)
=12\cdot\frac12
=6.
\]
证毕。
\end{proof}

\begin{conclusion}[次数 \(4\) 非伽罗瓦覆盖的极小分歧型]\label{con:xi-terminal-zm-c3-quotient-to-hly-minimal-ramification-type}
令 \(C_3\le A_4\) 为任一阶 \(3\) 子群，记
\[
C:=\widetilde{C}/C_3.
\]
则自然映射 \(C\to H_{\mathrm{LY}}\) 为次数 \(4\) 覆盖，且：
\begin{enumerate}
  \item 除 \(\infty_{\mathrm{LY}}\) 上方外处处无分歧；
  \item 在 \(\infty_{\mathrm{LY}}\) 的纤维中，局部循环型为 \(2+2\)，即恰有两处简单分歧点；
  \item 总分歧度为 \(2\)，从而 \(g(C)=6\)。
\end{enumerate}
\end{conclusion}
\begin{proof}[证明要点]
由结论 \ref{con:xi-terminal-zm-s4-a4-subextension-single-point-ramification}，
\(\widetilde{C}\to H_{\mathrm{LY}}\) 的唯一分歧惯性由 \(\rho^2=(13)(24)\) 生成。
商覆盖 \(C=\widetilde{C}/C_3\) 的分歧由 \(\rho^2\) 在陪集集 \(A_4/C_3\)（大小 \(4\)）上的置换作用决定；
该作用同构于 \(A_4\) 的自然四点作用，故 \(\rho^2\) 的循环分解为两不交换位，即 \(2+2\)。
因此 \(\infty_{\mathrm{LY}}\) 上方有两处简单分歧点，总分歧度为 \(2\)。
由 Riemann--Hurwitz
\[
2g(C)-2
=4\bigl(2g(H_{\mathrm{LY}})-2\bigr)+2
=4(2\cdot2-2)+2
=10,
\]
得 \(g(C)=6\)。证毕。
\end{proof}

\begin{conclusion}[\texorpdfstring{$C_3$}{C3} 商双覆盖的分歧线性类型与 Prym 极化补偿]\label{con:xi-terminal-zm-prym-polarization-type-rigidity-11112}
设
\[
E=\widetilde{C}/S_3,\qquad
C=\widetilde{C}/C_3,\qquad
\pi:C\longrightarrow E
\]
为由 \(C_3\le S_3\) 诱导的二次覆盖。记 \(B_\pi\) 为其分歧除子，
\[
P:=\operatorname{Prym}(C/E),
\]
并以 \(\Xi\) 表示 \(\Jac(C)\) 主极化在 \(P\) 上诱导的 Prym 极化。
则成立：
\begin{enumerate}
  \item \(g(C)=6\)，且
  \[
  B_\pi\sim 10O\in \Pic(E);
  \]
  \item 任一满足
  \[
  L^{\otimes 2}\cong \mathcal O_E(B_\pi)
  \]
  的平方根线丛均可写为
  \[
  L\cong \mathcal O_E(5O+\eta),\qquad \eta\in E[2],
  \]
  因而这样的 \(L\) 恰有 \(4\) 个；
  \item \(\dim P=5\) 且
  \[
  \mathrm{type}(\Xi)=(1,2,2,2,2).
  \]
  特别地，若 \(L_P\) 为表示 \(\Xi\) 的充足线丛，则
  \[
  h^0(P,L_P)=16;
  \]
  \item 若在 \(\QQ\)-同源意义下 \(P\sim_{\QQ}J_{\mathrm{LY}}\times B\)，则可取与幂等分解相容的同源
  \[
  \varphi:J_{\mathrm{LY}}\times B\longrightarrow P
  \]
  使
  \[
  \varphi^*\Xi=\lambda_{\mathrm{LY}}\boxtimes \lambda_B,
  \qquad
  \deg\Xi=\deg\lambda_{\mathrm{LY}}\deg\lambda_B=16.
  \]
  特别地，若 \(\deg\lambda_{\mathrm{LY}}=4\)，则 \(\deg\lambda_B=4\)；反之亦然。
\end{enumerate}
\end{conclusion}
\begin{proof}[证明要点]
由结论 \ref{con:xi-terminal-zm-s4-c3-quotient-jacobian-three-factor-decomposition}，
\[
\Jac(C)\sim_{\QQ}J_{\mathrm{LY}}\times E\times B,
\]
故 \(g(C)=6\)。又 \(g(E)=1\)，于是 \(\dim P=g(C)-g(E)=5\)。
对二次覆盖 \(\pi:C\to E\)，设分歧点数为 \(r\)。
由 Hurwitz 公式
\[
2g(C)-2=2(2g(E)-2)+r
\]
得 \(r=10\)。
另一方面，结论 \ref{con:xi-terminal-zm-leyang-unramified-sixpoint-abel-closed-form-branch-balance} 已证明
\[
\bigoplus_{P\in\mathrm{Supp}(B_\pi)}P=O,
\]
这在椭圆曲线 \(E\simeq \Pic^0(E)\) 的群律翻译下等价于
\[
B_\pi-10O\sim 0,
\]
从而 \(B_\pi\sim 10O\)。

任一平方根线丛 \(L\) 满足
\[
L^{\otimes 2}\cong \mathcal O_E(B_\pi)\cong \mathcal O_E(10O),
\]
故
\[
L\in [2]^{-1}\!\bigl(\mathcal O_E(5O)\bigr).
\]
而 \([2]:\Pic^0(E)\to\Pic^0(E)\) 的核为 \(E[2]\)，故全部解构成一个 \(E[2]\)-torsor，亦即
\[
L\cong \mathcal O_E(5O+\eta),\qquad \eta\in E[2].
\]
对分歧二次覆盖到底曲线属 \(1\) 的 Prym，标准极化公式给出
\[
\mathrm{type}(\Xi)=\bigl(1^{g(E)},2^{r/2-1}\bigr)=(1,2,2,2,2).
\]
于是对任一表示 \(\Xi\) 的充足线丛 \(L_P\)，有
\[
h^0(P,L_P)=1\cdot 2^4=16.
\]

最后，由结论 \ref{con:xi-terminal-zm-s4-c3-prym-leyang-equivalence} 的同源分解
\[
P\sim_{\QQ}J_{\mathrm{LY}}\times B
\]
可取与幂等分解相容的同源 \(\varphi\)，使 \(\varphi^*\Xi\) 分裂为外积极化。按本文对极化次数的记法，极化型各项之积即为其次数，故
\[
\deg\Xi=1\cdot 2^4=16=\deg\lambda_{\mathrm{LY}}\deg\lambda_B.
\]
于是若一侧次数为 \(4\)，则另一侧亦必为 \(4\)。证毕。
\end{proof}

\input{con__xi-terminal-zm-s4-v4-c3-prym-pel-ns-gonality-refinements}

\paragraph{\texorpdfstring{$S_4$}{S4} 子群格中低亏格商与双覆盖 Prym 的补充显式化}
沿用结论 \ref{con:xi-terminal-zm-s4-jacobian-group-algebra-prym-e3-b3} 与
\ref{con:xi-terminal-zm-s4-chevalley-weil-explicit-decomposition} 的记号，记
\[
\Jac(\widetilde C)\sim_{\QQ}J_{\mathrm{LY}}\times E_{\mathrm{res}}^2\times E^3\times B^3,
\]
\[
H^0(\widetilde C,\Omega^1)\cong 2\varepsilon\oplus \rho_2\oplus \rho_3\oplus 3\rho_3',
\]
并置
\[
V_4^{\mathrm{tr}}:=\langle(12),(34)\rangle.
\]

\begin{conclusion}[子群格中三条低亏格商曲线的 Jacobian 同源型显式化]\label{con:xi-terminal-zm-s4-lowgenus-subgroup-quotients-jacobian-explicit}
记
\[
C_4:=\widetilde C/\langle(1234)\rangle,\qquad
C_{V_4^{\mathrm{norm}}}:=\widetilde C/V_4^{\mathrm{norm}},\qquad
C_{V_4^{\mathrm{tr}}}:=\widetilde C/V_4^{\mathrm{tr}}.
\]
则
\[
g(C_4)=4,\qquad \Jac(C_4)\sim_{\QQ}E_{\mathrm{res}}\times B;
\]
\[
g(C_{V_4^{\mathrm{norm}}})=4,\qquad
\Jac(C_{V_4^{\mathrm{norm}}})\sim_{\QQ}J_{\mathrm{LY}}\times E_{\mathrm{res}}^2;
\]
\[
g(C_{V_4^{\mathrm{tr}}})=2,\qquad
\Jac(C_{V_4^{\mathrm{tr}}})\sim_{\QQ}E_{\mathrm{res}}\times E.
\]
\end{conclusion}
\begin{proof}[证明要点]
对任意子群 \(H\le S_4\)，有
\[
H^0(\widetilde C/H,\Omega^1)\cong H^0(\widetilde C,\Omega^1)^H,
\]
故只需计算
\(
2\varepsilon\oplus \rho_2\oplus \rho_3\oplus 3\rho_3'
\)
在各子群下的不变向量维数。对任意不可约表示 \(W\)，由角色平均公式
\[
\dim W^H=\frac{1}{|H|}\sum_{h\in H}\chi_W(h).
\]

对 \(H=\langle(1234)\rangle\)，其元素型为
\(
1,(1234),(13)(24),(1432)
\)，从而
\[
\dim(\varepsilon^H)=0,\quad
\dim(\rho_2^H)=1,\quad
\dim(\rho_3^H)=0,\quad
\dim((\rho_3')^H)=1.
\]
于是
\[
g(C_4)=2\cdot 0+1\cdot 1+1\cdot 0+3\cdot 1=4,
\]
并由等型对应得到
\[
\Jac(C_4)\sim_{\QQ}E_{\mathrm{res}}\times B.
\]

对 \(H=V_4^{\mathrm{norm}}\)，其非平凡元均为双换位，故
\[
\dim(\varepsilon^H)=1,\quad
\dim(\rho_2^H)=2,\quad
\dim(\rho_3^H)=0,\quad
\dim((\rho_3')^H)=0.
\]
因此
\[
g(C_{V_4^{\mathrm{norm}}})=2\cdot 1+1\cdot 2=4,
\]
且
\[
\Jac(C_{V_4^{\mathrm{norm}}})\sim_{\QQ}J_{\mathrm{LY}}\times E_{\mathrm{res}}^2.
\]

对 \(H=V_4^{\mathrm{tr}}\)，其元素型为
\(
1,(12),(34),(12)(34)
\)，故
\[
\dim(\varepsilon^H)=0,\quad
\dim(\rho_2^H)=1,\quad
\dim(\rho_3^H)=1,\quad
\dim((\rho_3')^H)=0.
\]
于是
\[
g(C_{V_4^{\mathrm{tr}}})=1+1=2,
\]
并且
\[
\Jac(C_{V_4^{\mathrm{tr}}})\sim_{\QQ}E_{\mathrm{res}}\times E.
\]
证毕。
\end{proof}

\begin{conclusion}[两条属 \(4\) 双覆盖到 \texorpdfstring{$E_{\mathrm{res}}$}{Eres} 的 Prym 完全辨识]\label{con:xi-terminal-zm-s4-two-genus4-doublecovers-prym-identification}
自然商映射
\[
\pi_{C_4}:C_4\longrightarrow E_{\mathrm{res}},\qquad
\pi_{V_4}:C_{V_4^{\mathrm{norm}}}\longrightarrow E_{\mathrm{res}}
\]
均为分歧二次覆盖，并满足
\[
g(C_4)=g(C_{V_4^{\mathrm{norm}}})=4,\qquad
\deg\operatorname{Branch}(\pi_{C_4})=\deg\operatorname{Branch}(\pi_{V_4})=6.
\]
此外，
\[
\operatorname{Prym}(C_4/E_{\mathrm{res}})\sim_{\QQ}B,
\]
\[
\operatorname{Prym}(C_{V_4^{\mathrm{norm}}}/E_{\mathrm{res}})\sim_{\QQ}J_{\mathrm{LY}}\times E_{\mathrm{res}},
\]
且两条覆盖的 Prym 极化型同为
\[
(1,2,2).
\]
\end{conclusion}
\begin{proof}[证明要点]
由 \(\langle(1234)\rangle\subset D_8\) 与 \(V_4^{\mathrm{norm}}\subset D_8\) 的指数均为 \(2\)，可知两条商曲线都带有到
\(
E_{\mathrm{res}}=\widetilde C/D_8
\)
的自然二次映射。由结论
\ref{con:xi-terminal-zm-s4-lowgenus-subgroup-quotients-jacobian-explicit}，
\[
g(C_4)=g(C_{V_4^{\mathrm{norm}}})=4.
\]
再由 Hurwitz 公式
\[
2g(Y)-2=2\bigl(2g(E_{\mathrm{res}})-2\bigr)+\deg\operatorname{Branch}(Y/E_{\mathrm{res}})
\]
并代入 \(g(E_{\mathrm{res}})=1\)，得两条覆盖的分支次数均为 \(6\)。

对分歧二次覆盖的 Jacobian 标准分解
\[
\Jac(Y)\sim_{\QQ}E_{\mathrm{res}}\times \operatorname{Prym}(Y/E_{\mathrm{res}})
\]
与结论
\ref{con:xi-terminal-zm-s4-lowgenus-subgroup-quotients-jacobian-explicit}
比较，可得
\[
\operatorname{Prym}(C_4/E_{\mathrm{res}})\sim_{\QQ}B,
\qquad
\operatorname{Prym}(C_{V_4^{\mathrm{norm}}}/E_{\mathrm{res}})\sim_{\QQ}J_{\mathrm{LY}}\times E_{\mathrm{res}}.
\]
最后，对底曲线属 \(1\)、分支次数 \(6\) 的分歧二次覆盖，Prym 极化型由一般公式给出为
\[
\bigl(1^{g(E_{\mathrm{res}})},2^{6/2-1}\bigr)=(1,2,2).
\]
证毕。
\end{proof}

\begin{conclusion}[\texorpdfstring{$V_4^{\mathrm{tr}}$}{V4tr} 商 Jacobian 分裂的部分塔式可见性与剩余 endoscopic 因子]\label{con:xi-terminal-zm-s4-v4tr-splitting-partial-tower-visibility}
对
\[
C_{V_4^{\mathrm{tr}}}:=\widetilde C/V_4^{\mathrm{tr}}
\]
有
\[
\Jac(C_{V_4^{\mathrm{tr}}})\sim_{\QQ}E_{\mathrm{res}}\times E.
\]
其中 \(E_{\mathrm{res}}\) 因子可由包含链
\[
V_4^{\mathrm{tr}}\subset D_8^{\mathrm{tr}}:=N_{S_4}\!\bigl(\langle(1324)\rangle\bigr)
\]
几何实现，因为 \(\widetilde C/D_8^{\mathrm{tr}}\cong E_{\mathrm{res}}\)；然而不存在 \(S_3\le S_4\) 含 \(V_4^{\mathrm{tr}}\)，故 \(E\)-因子不能由任何属 \(1\) 子群塔商直接产生。因而该属 \(2\) Jacobian 的椭圆分裂并非完全由子群塔解释，其中 \(E\)-方向仍由 \(\QQ[S_4]\)-幂等投影强制出现。
\end{conclusion}
\begin{proof}[证明要点]
第一断言即结论
\ref{con:xi-terminal-zm-s4-lowgenus-subgroup-quotients-jacobian-explicit}
的第三个分裂式。

取
\[
D_8^{\mathrm{tr}}=\langle(1324),(12)\rangle.
\]
则
\[
(1324)^2=(12)(34)\in V_4^{\mathrm{tr}},\qquad
(12)(12)(34)=(34),
\]
从而 \(V_4^{\mathrm{tr}}=\langle(12),(34)\rangle\subset D_8^{\mathrm{tr}}\)。又所有 \(D_8\) 子群在 \(S_4\) 中彼此共轭，故
\[
\widetilde C/D_8^{\mathrm{tr}}\cong \widetilde C/D_8=E_{\mathrm{res}},
\]
于是 \(E_{\mathrm{res}}\) 因子可由子群塔商解释。

另一方面，由结论 \ref{con:xi-terminal-zm-s4-hurwitz-subgroup-spectrum-cycletypes}，属 \(1\) 的子群商仅来自 \(D_8\) 与 \(S_3\) 两类子群。任一 \(S_3\) 作为点稳定子都不可能同时含有两条互不共享固定点的换位 \((12)\) 与 \((34)\)，故 \(V_4^{\mathrm{tr}}\not\subset S_3\)。因此 \(E\)-因子不能来自任何包含 \(V_4^{\mathrm{tr}}\) 的属 \(1\) 子群商。综上，\(\Jac(C_{V_4^{\mathrm{tr}}})\) 的分裂只有 \(E_{\mathrm{res}}\) 一侧具备塔式几何来源，而 \(E\)-方向仍是表示论意义下的剩余 endoscopic 因子。证毕。
\end{proof}

\endinput

\begin{theorem}[三类换位碰撞边界下四个 Langlands 分量的环面秩与导子指数]\label{thm:xi-terminal-zm-s4-langlands-factor-torus-rank-conductor-table}
\leanverified{paper\_xi\_terminal\_zm\_s4\_langlands\_factor\_torus\_rank\_conductor\_table}
沿用定理 \ref{thm:xi-s4-isotypic-jacobian-toric-rank-by-boundary-type} 的三类换位碰撞边界记号 \(r\in\{1,2,3\}\)，以及结论 \ref{con:xi-terminal-zm-s4-jacobian-group-algebra-prym-e3-b3} 的 \(\QQ\)-同源分解
\[
\Jac(\widetilde C)\ \sim_{\QQ}\ J_{\mathrm{LY}}\times E_{\mathrm{res}}^2\times E^3\times B^3
\]
与等型对应
\(
\varepsilon\leftrightarrow J_{\mathrm{LY}},\
\rho_2\leftrightarrow E_{\mathrm{res}}^2,\
\rho_3\leftrightarrow E^3,\
\rho_3'\leftrightarrow B^3
\)。
对每个分量 \(A\in\{J_{\mathrm{LY}},E_{\mathrm{res}},E,B\}\)，记其在对应半稳定退化处的环面秩为 \(t(A)\)。
则对三类碰撞型 \(r\) 有表：
\[
\begin{array}{c|cccc}
r & t(J_{\mathrm{LY}}) & t(E_{\mathrm{res}}) & t(E) & t(B)\\
\hline
1 & 1 & 1 & 1 & 2\\
2 & 0 & 0 & 1 & 1\\
3 & 1 & 0 & 0 & 1
\end{array}
\]
并且在特征 \(0\) 的半稳定几何退化处（或更一般地在温和半稳定情形下），\(\ell\)-进 Tate 模的惯性表示无 Swan 部分，因而 Artin 导子指数等于环面秩。
\end{theorem}
\begin{proof}
由定理 \ref{thm:xi-s4-isotypic-jacobian-toric-rank-by-boundary-type}，在碰撞型 \(r\) 处四个 \(\QQ[S_4]\)-等型因子 \(A_\varepsilon,A_{\rho_2},A_{\rho_3},A_{\rho_3'}\) 的环面秩
\((t_\varepsilon,t_{\rho_2},t_{\rho_3},t_{\rho_3'})\)
由该定理中的表格确定。
另一方面，结论 \ref{con:xi-terminal-zm-s4-jacobian-group-algebra-prym-e3-b3} 的等型对应表明
\[
A_\varepsilon\sim_{\QQ} J_{\mathrm{LY}},\qquad
A_{\rho_2}\sim_{\QQ}E_{\mathrm{res}}^2,\qquad
A_{\rho_3}\sim_{\QQ}E^3,\qquad
A_{\rho_3'}\sim_{\QQ}B^3.
\]
环面秩对同源不变且对直积可加，故
\[
t(J_{\mathrm{LY}})=t_\varepsilon,\qquad
2\,t(E_{\mathrm{res}})=t_{\rho_2},\qquad
3\,t(E)=t_{\rho_3},\qquad
3\,t(B)=t_{\rho_3'}.
\]
将三类 \(r\) 的数值代入即得表格。
\par
最后，在特征 \(0\) 的半稳定几何退化处（或温和半稳定情形下），惯性表示的 Swan 部分为 \(0\)，而 Artin 导子指数等于环面秩与 Swan 部分之和，从而导子指数等于环面秩。证毕。
\end{proof}

\begin{corollary}[碰撞选择律与 \texorpdfstring{$B$}{B}-通道的普遍非良约化]\label{cor:xi-terminal-zm-s4-langlands-factor-collision-selection}
\leanverified{paper\_xi\_terminal\_zm\_s4\_langlands\_factor\_collision\_selection}
在定理 \ref{thm:xi-terminal-zm-s4-langlands-factor-torus-rank-conductor-table} 的设定下：
\begin{enumerate}
  \item 对所有三类换位碰撞边界均有 \(t(B)\ge 1\)，因此 \(B\) 在这些边界处均为半稳定但非良约化；并且在 \(r=1\) 型边界上有 \(t(B)=2\)。
  \item \(J_{\mathrm{LY}}\) 的良约化仅发生在 \(r=2\) 型边界；\(E_{\mathrm{res}}\) 的良约化发生在 \(r\in\{2,3\}\)；\(E\) 的良约化仅发生在 \(r=3\) 型边界。
\end{enumerate}
\end{corollary}
\begin{proof}
良约化等价于环面秩为 \(0\)。直接读取定理 \ref{thm:xi-terminal-zm-s4-langlands-factor-torus-rank-conductor-table} 的表格即可。证毕。
\end{proof}

\begin{theorem}[两类 Prym 的环面秩分层]\label{thm:xi-terminal-zm-s4-prym-torus-rank-layering}
\leanverified{paper\_xi\_terminal\_zm\_s4\_prym\_torus\_rank\_layering}
沿用结论 \ref{con:xi-terminal-zm-s4-involution-prym-polarization-eigensplitting} 的记号，
令 \(\pi_\tau:\widetilde C\to C_\tau\) 为换位商诱导的二次覆盖，\(\pi_\delta:\widetilde C\to C_\delta\) 为双换位商诱导的二次覆盖。
则在三类换位碰撞边界 \(r\in\{1,2,3\}\) 上，其 Prym 的环面秩分别为
\[
t\!\bigl(\Prym(\pi_\delta)\bigr)=(6,4,2),
\qquad
t\!\bigl(\Prym(\pi_\tau)\bigr)=(7,3,3),
\]
其中括号中三元组按 \(r=(1,2,3)\) 顺序排列。
\end{theorem}
\begin{proof}
由结论 \ref{con:xi-terminal-zm-s4-involution-prym-polarization-eigensplitting} 的 \(\QQ\)-同源分解，
\[
\Prym(\pi_\tau)\sim_{\QQ} J_{\mathrm{LY}}\times E_{\mathrm{res}}\times E\times B^2,
\qquad
\Prym(\pi_\delta)\sim_{\QQ} E^2\times B^2.
\]
环面秩对直积可加并对同源不变，故
\[
t\!\bigl(\Prym(\pi_\delta)\bigr)=2t(E)+2t(B),
\qquad
t\!\bigl(\Prym(\pi_\tau)\bigr)=t(J_{\mathrm{LY}})+t(E_{\mathrm{res}})+t(E)+2t(B).
\]
将定理 \ref{thm:xi-terminal-zm-s4-langlands-factor-torus-rank-conductor-table} 的三类 \(r\) 下 \((t(J_{\mathrm{LY}}),t(E_{\mathrm{res}}),t(E),t(B))\) 逐项代入即可得到所示数值。证毕。
\end{proof}

\endinput

\subsubsection{\texorpdfstring{$S_4$}{S4} 正则闭包的整体谱：属数、Chevalley--Weil 分解、正规子群商与 Jacobian 表示分裂}\label{subsubsec:xi-s4-regular-closure-global-spectrum}
沿用
\[
\Pi(\lambda,y)=\lambda^{4}-\lambda^{3}-(2y+1)\lambda^{2}+\lambda+y(y+1)\in\QQ(y)[\lambda],
\]
并记
\[
g(y):=256y^{3}+411y^{2}+165y+32,\qquad
\Delta(y):=-y(y-1)g(y).
\]
设 \(L/\QQ(y)\) 为 \(\Pi\) 的分裂域，\(G:=\mathrm{Gal}(L/\QQ(y))\cong S_4\)，并令 \(X\) 为 \(\QQ(X)=L\) 的光滑射影模型（即 \(X=\widetilde X\)）。

\begin{conclusion}[\texorpdfstring{$S_4$}{S4} 正则闭包曲线的精确属数]\label{con:xi-terminal-zm-s4-regular-closure-exact-genus-16}
覆盖 \(X\to\PP^1_y\) 为次数 \(24\) 的正则伽罗瓦覆盖，且
\[
g(X)=16.
\]
\end{conclusion}
\begin{proof}[证明要点]
由结论 \ref{con:xi-terminal-zm-s4-galois-closure-genus-tower-holomorphic-differentials}，分歧护照为一个阶 \(4\) 惯性点与五个阶 \(2\) 惯性点，即 \((4)\cdot(2)^5\)。
对正则伽罗瓦覆盖应用 Hurwitz 公式
\[
g(X)=1+\frac{|S_4|}{2}\!\left(-2+\sum_{p\in B}\Bigl(1-\frac1{e_p}\Bigr)\right),
\]
代入 \(|S_4|=24\)、\((e_p)=(4,2,2,2,2,2)\) 即得
\[
g(X)=1+12\!\left(-2+\frac34+5\cdot\frac12\right)=16.
\]
证毕。
\end{proof}

\begin{conclusion}[\texorpdfstring{$S_4$}{S4} 正则闭包中共轭类定点计数与循环商曲线属数]\label{con:xi-terminal-zm-s4-conjugacy-class-fixedpoints-and-cyclic-quotient-genera}
设
\(
\varphi:X\to \PP^1_y
\)
为结论 \ref{con:xi-terminal-zm-s4-regular-closure-exact-genus-16} 中的正则 \(S_4\)-覆盖，其分歧惯性型为一个 \(C_4\) 与五个 \(C_2\)。
对共轭类代表 \(g\in S_4\)，记 \(\mathrm{Fix}_X(g)\) 为其在 \(X\) 上的定点集，则有
\[
\begin{array}{c|c|c}
g\text{ 的循环型} & \#\mathrm{Fix}_X(g) & g(X/\langle g\rangle)\\
\hline
(2) & 10 & 6\\
(2,2) & 2 & 8\\
(4) & 2 & 4\\
(3) & 0 & 6
\end{array}
\]
\end{conclusion}
\begin{proof}[证明]
设 \(b\in \PP^1_y\) 为分歧值，\(I_b\le S_4\) 为其惯性群。由正则覆盖的陪集描述，
\(\varphi^{-1}(b)\cong S_4/I_b\)，且 \(xI_b\) 对应点的稳定子为 \(xI_bx^{-1}\)。故对任意 \(g\in S_4\) 有
\begin{equation}\label{eq:xi-terminal-zm-s4-fixedpoint-fiber-coset-formula}
\#\bigl(\mathrm{Fix}_X(g)\cap \varphi^{-1}(b)\bigr)
=\frac{\#\{x\in S_4\mid x^{-1}gx\in I_b\}}{|I_b|}.
\end{equation}

先取换位 \(\tau\)。五个有限分歧点 \(b_1,\dots,b_5\) 的惯性均为 \(C_2\)，其非平凡元均为换位。
由 \eqref{eq:xi-terminal-zm-s4-fixedpoint-fiber-coset-formula}，每个 \(b_i\) 的贡献为
\[
\frac{\#\{x\in S_4\mid x^{-1}\tau x=\tau_i\}}{2}
=\frac{|C_{S_4}(\tau)|}{2}
=\frac{4}{2}=2,
\]
故 \(\#\mathrm{Fix}_X(\tau)=10\)。对阶 \(2\) 商映射 \(X\to X/\langle\tau\rangle\) 用 Hurwitz 公式并代入 \(g(X)=16\)：
\[
2g(X)-2=2\bigl(2g(X/\langle\tau\rangle)-2\bigr)+\#\mathrm{Fix}_X(\tau),
\]
得 \(30=4g'-4+10\)，故 \(g'=6\)。

再取双换位 \(\delta\)。有限分歧惯性不含双换位，故仅无穷远分歧值 \(b_\infty\) 可能贡献。
记其惯性 \(I_{b_\infty}=\langle c\rangle\cong C_4\)，则唯一阶 \(2\) 元 \(c^2\) 为双换位。
由 \eqref{eq:xi-terminal-zm-s4-fixedpoint-fiber-coset-formula}
\[
\#\mathrm{Fix}_X(\delta)
=\frac{\#\{x\in S_4\mid x^{-1}\delta x=c^2\}}{4}
=\frac{|C_{S_4}(\delta)|}{4}
=\frac{8}{4}=2.
\]
再由 Hurwitz（阶 \(2\)）得
\[
30=2\bigl(2g(X/\langle\delta\rangle)-2\bigr)+2,
\]
故 \(g(X/\langle\delta\rangle)=8\)。

取四循环 \(\gamma\)。同理仅 \(b_\infty\) 贡献。
由于四循环同属一共轭类，可不失一般性取 \(\gamma=c\)。
此时式 \eqref{eq:xi-terminal-zm-s4-fixedpoint-fiber-coset-formula} 的条件 \(x^{-1}\gamma x\in \langle\gamma\rangle\) 等价于 \(x\in N_{S_4}(\langle\gamma\rangle)\)，
而 \(|N_{S_4}(\langle\gamma\rangle)|=8\)，故
\[
\#\mathrm{Fix}_X(\gamma)=\frac{8}{4}=2.
\]
又 \(\gamma^2\) 为双换位，已知 \(\#\mathrm{Fix}_X(\gamma^2)=2\)，且 \(\mathrm{Fix}_X(\gamma)\subseteq \mathrm{Fix}_X(\gamma^2)\)，
从而两者相等，不存在“仅被 \(\gamma^2\) 固定”之点。
故 \(X\to X/\langle\gamma\rangle\) 的分歧仅来自这 \(2\) 个点，且各分歧指数为 \(4\)。Hurwitz（阶 \(4\)）给出
\[
30=4\bigl(2g(X/\langle\gamma\rangle)-2\bigr)+2(4-1),
\]
解得 \(g(X/\langle\gamma\rangle)=4\)。

最后取三循环 \(\sigma\)。各惯性群仅含阶 \(2\) 或 \(4\) 元，故 \(\sigma\) 不可能固定分歧纤维点，从而 \(\mathrm{Fix}_X(\sigma)=\varnothing\)。
于是 \(X\to X/\langle\sigma\rangle\) 无分歧，Hurwitz（阶 \(3\)）给出
\[
30=3\bigl(2g(X/\langle\sigma\rangle)-2\bigr),
\]
故 \(g(X/\langle\sigma\rangle)=6\)。证毕。
\end{proof}

\begin{conclusion}[四循环商的 Prym 三维因子与 \texorpdfstring{$9$}{9} 维等型分量识别]\label{con:xi-terminal-zm-s4-c4-prym-threefold-and-v3sgn-identification}
取 \(c\in S_4\) 为无穷远惯性生成元（四循环），定义
\[
Y:=X/\langle c\rangle,\qquad
E_{\mathrm D}:=X/D_8
\]
（其中 \(D_8=N_{S_4}(\langle c\rangle)\)）。令
\[
\pi:Y\longrightarrow E_{\mathrm D}
\]
为诱导的次数 \(2\) 映射，并定义
\[
B:=\mathrm{Prym}(Y/E_{\mathrm D})
=\ker\!\bigl(\mathrm{Nm}_\pi:\mathrm{Jac}(Y)\to \mathrm{Jac}(E_{\mathrm D})\bigr)^0.
\]
则有：
\begin{enumerate}
  \item \(\pi\) 为分歧双覆盖，且 \(\deg\mathrm{Branch}(\pi)=6\)，从而 \(\dim B=3\)；
  \item \(\mathrm{Jac}(Y)\sim_{\QQ}E_{\mathrm D}\times B\)；
  \item 对任意换位 \(\tau\) 与双换位 \(\delta\)，有
  \[
  \mathrm{Jac}(X/\langle\tau\rangle)\sim_{\QQ}E^2\times E_{\mathrm D}\times B,
  \]
  \[
  \mathrm{Jac}(X/\langle\delta\rangle)\sim_{\QQ}\mathrm{Jac}(X_A)\times E\times E_{\mathrm D}^2\times B;
  \]
  \item 若 \(A_{V_3\otimes\mathrm{sgn}}\subset \mathrm{Jac}(X)\) 表示对应 \(V_3\otimes\mathrm{sgn}\) 的 \(9\) 维等型分量，则
  \[
  A_{V_3\otimes\mathrm{sgn}}\sim_{\QQ}B^3.
  \]
\end{enumerate}
\end{conclusion}
\begin{proof}[证明]
由子群塔 \(\langle c\rangle\subset D_8\subset S_4\) 可知 \(\deg(\pi)=2\)。
又由结论 \ref{con:xi-terminal-zm-s4-conjugacy-class-fixedpoints-and-cyclic-quotient-genera} 与
\ref{con:xi-terminal-zm-s4-hurwitz-subgroup-spectrum-cycletypes}，
\[
g(Y)=g(X/\langle c\rangle)=4,\qquad g(E_{\mathrm D})=g(X/D_8)=1.
\]
Hurwitz 公式
\[
2g(Y)-2=2\bigl(2g(E_{\mathrm D})-2\bigr)+\deg\mathrm{Branch}(\pi)
\]
给出 \(\deg\mathrm{Branch}(\pi)=6\)，并有
\(
\dim B=g(Y)-g(E_{\mathrm D})=3
\)。
这证明第 \(1\) 点。

第 \(2\) 点由双覆盖的 Prym 标准分解
\[
\mathrm{Jac}(Y)\sim_{\QQ}\pi^*\mathrm{Jac}(E_{\mathrm D})\times \mathrm{Prym}(Y/E_{\mathrm D})
\]
直接得到。

第 \(3\) 点与结论
\ref{con:xi-terminal-zm-s4-transposition-quotient-jacobian-full-decomposition}、
\ref{con:xi-terminal-zm-s4-doubletransposition-quotient-jacobian-full-decomposition}
一致：其中文中记号 \(E_{\mathrm{res}}\) 即此处 \(E_{\mathrm D}\)，且
\(
J_{\mathrm{LY}}\sim_{\QQ}\mathrm{Jac}(X_A)
\)
见结论
\ref{con:xi-terminal-zm-s4-quotient-genus-spectrum-and-jacobian-representation-splitting}。

第 \(4\) 点可由两种等价方式得到：其一，由结论
\ref{con:xi-terminal-zm-s4-jacobian-group-algebra-prym-e3-b3}
的分解
\[
\mathrm{Jac}(X)\sim_{\QQ}\mathrm{Jac}(X_A)\times E_{\mathrm D}^2\times E^3\times B^3
\]
直接识别 \(V_3\otimes\mathrm{sgn}\) 等型即为 \(B^3\)；
其二，由
\(
e_{V_3\otimes\mathrm{sgn}}\QQ[S_4]e_{V_3\otimes\mathrm{sgn}}\cong M_3(\QQ)
\)
与原始幂等分解得到同结论。证毕。
\end{proof}

\begin{conclusion}[Hasse--Weil 动机分解与良素处局部 \texorpdfstring{$L$}{L} 因子幂次刚性]\label{con:xi-terminal-zm-s4-hasse-weil-full-factorization-rigid-exponents}
在 \(\QQ\)-同源意义下有
\[
\mathrm{Jac}(X)\sim_{\QQ}\mathrm{Jac}(X_A)\times E^3\times E_{\mathrm D}^2\times B^3.
\]
因而对任一使上述各因子皆良约化的素数 \(p\)，局部 \(L\)-多项式分子满足
\[
P_p(X;T)
=P_p(X_A;T)\cdot P_p(E;T)^3\cdot P_p(E_{\mathrm D};T)^2\cdot P_p(B;T)^3,
\]
并且
\(
\deg P_p(B;T)=2\dim B=6
\)。
\end{conclusion}
\begin{proof}[证明]
第一式由结论
\ref{con:xi-terminal-zm-s4-c4-prym-threefold-and-v3sgn-identification}
第 \(4\) 点与既有 \(A_{\mathrm{sgn}}\)、\(A_{V_3}\)、\(A_{V_2}\) 识别合并得到。
第二式由“同源阿贝尔簇在良素处具有相同局部 \(L\)-因子”及乘积阿贝尔簇局部因子的乘法性直接推出。证毕。
\end{proof}

\begin{conclusion}[全纯 \(1\)-形式的 Chevalley--Weil 刚性分解]\label{con:xi-terminal-zm-s4-holomorphic-oneforms-chevalley-weil-rigid-decomposition}
记 \(S_4\) 的不可约表示为
\[
\mathbf 1,\ \mathrm{sgn},\ V_3,\ V_3\otimes\mathrm{sgn},\ V_2.
\]
则有
\[
H^0(X,\Omega_X^1)\cong
\mathrm{sgn}^{\oplus 2}\oplus V_3\oplus (V_3\otimes\mathrm{sgn})^{\oplus 3}\oplus V_2,
\]
并且平凡表示不出现。
\end{conclusion}
\begin{proof}[证明要点]
对 \(\rho\neq \mathbf 1\)，Chevalley--Weil 公式在本情形可写为
\[
m_\rho=-\dim\rho+\frac12\sum_{p\in B}\bigl(\dim\rho-\dim(\rho^{I_p})\bigr),
\]
且 \(m_{\mathbf 1}=0\)（底曲线 \(\PP^1\) 无全纯 \(1\)-形式）。
本覆盖仅需 \(C_2\) 与 \(C_4\) 两类惯性不变维数：
\[
\begin{aligned}
&\dim(\mathrm{sgn}^{C_2})=0,\ \dim(\mathrm{sgn}^{C_4})=0,\\
&\dim(V_3^{C_2})=2,\ \dim(V_3^{C_4})=0,\\
&\dim((V_3\!\otimes\!\mathrm{sgn})^{C_2})=1,\ \dim((V_3\!\otimes\!\mathrm{sgn})^{C_4})=1,\\
&\dim(V_2^{C_2})=1,\ \dim(V_2^{C_4})=1.
\end{aligned}
\]
代入五个 \(C_2\) 分歧点与一个 \(C_4\) 分歧点，得到
\[
m_{\mathrm{sgn}}=2,\quad m_{V_3}=1,\quad m_{V_3\otimes\mathrm{sgn}}=3,\quad m_{V_2}=1,\quad m_{\mathbf 1}=0.
\]
与结论 \ref{con:xi-terminal-zm-s4-chevalley-weil-explicit-decomposition} 一致。证毕。
\end{proof}

\begin{conclusion}[正规子群商曲线的显式模型与属数：\texorpdfstring{$A_4$}{A4} 商与 \texorpdfstring{$V_4$}{V4} 商]\label{con:xi-terminal-zm-s4-normal-subgroup-quotients-explicit-models-genus}
设
\[
X_A:=X/A_4,\qquad X_V:=X/V_4^{\mathrm{norm}}.
\]
则：
\begin{enumerate}
  \item \(X_A\to\PP^1_y\) 为二次覆盖，且
  \[
  \QQ(X_A)=\QQ\!\left(y,\sqrt{\Delta(y)}\right),
  \qquad
  X_A:\ t^2=\Delta(y),
  \qquad g(X_A)=2.
  \]
  \item \(X_V\to\PP^1_y\) 为次数 \(6\) 的 \(S_3\)-正则覆盖，且
  \[
  \QQ(X_V)=\QQ(y,z,\sqrt{\Delta(y)}),\qquad \mathscr R(z,y)=0,
  \]
  其中可取
  \[
  \mathscr R(z,y)=z^3+(2y+1)z^2-(4y^2+4y+1)z-(8y^3+13y^2+5y+1),
  \]
  并且 \(g(X_V)=4\)。
\end{enumerate}
\end{conclusion}
\begin{proof}[证明要点]
对 \(A_4\)：指数 \(2\) 正规子群对应符号同态核，故固定域由判别式平方根生成，即
\(
\QQ(X_A)=\QQ(y,\sqrt{\Delta(y)})
\)。
由于 \(\deg\Delta=5\) 为奇数，双覆盖在五个有限零点与无穷远点分歧，总分歧点数 \(6\)，故 \(g(X_A)=2\)。
\par
对 \(V_4^{\mathrm{norm}}\)：商群 \(S_4/V_4^{\mathrm{norm}}\cong S_3\)，固定域由三次解式闭包给出。
以 \(\mathscr R\) 为 resolvent cubic，附加其判别式平方根即得 \(S_3\)-闭包域，
从而 \(\QQ(X_V)=\QQ(y,z,\sqrt{\Delta(y)})\)。
由结论 \ref{con:xi-terminal-zm-s4-hurwitz-subgroup-spectrum-cycletypes} 已知 \(g(X_V)=4\)；等价地，用 \(|S_3|=6\) 与六个阶 \(2\) 分歧值代入 Hurwitz 公式亦得同值。证毕。
\end{proof}

\begin{conclusion}[双换位商曲线的双二次结构与纤维积刻画]\label{con:xi-terminal-zm-s4-delta-quotient-v4-fiberproduct}
沿用本小节记号，取无穷远惯性生成元 \(c\in S_4\) 为四循环，置
\[
\delta:=c^2,\qquad
Z:=X/\langle\delta\rangle,\qquad
E_{\mathrm D}:=X/D_8,\qquad
Y:=X/\langle c\rangle.
\]
则有：
\begin{enumerate}
  \item \(g(Z)=8\)；
  \item \(Z\to E_{\mathrm D}\) 为次数 \(4\) 的 Galois 覆盖，覆盖群同构于 \(V_4\)；
  \item 在函数域层面
  \[
  \QQ(Z)=\QQ(X_V)\cdot\QQ(Y)\quad\text{于}\ \QQ(E_{\mathrm D})\ \text{之上},
  \]
  且
  \[
  \QQ(X_V)\cap\QQ(Y)=\QQ(E_{\mathrm D}),
  \]
  从而 \(Z\) 同构于纤维积 \(X_V\times_{E_{\mathrm D}}Y\) 的规范化。
\end{enumerate}
\end{conclusion}
\begin{proof}
第 \(1\) 点由结论 \ref{con:xi-terminal-zm-s4-conjugacy-class-fixedpoints-and-cyclic-quotient-genera} 中
\(
\#\mathrm{Fix}_X(\delta)=2
\)
与 \(g(X)=16\) 代入二次商 Hurwitz 公式：
\[
2g(X)-2=2(2g(Z)-2)+\#\mathrm{Fix}_X(\delta)
\]
即得 \(g(Z)=8\)。

第 \(2\) 点中，\(\delta\in\langle c\rangle\subset D_8\)，且 \(V_4^{\mathrm{norm}}\subset D_8\)。
并有
\[
\langle\delta\rangle=V_4^{\mathrm{norm}}\cap\langle c\rangle,\qquad
D_8/\langle\delta\rangle\cong V_4.
\]
故 \(Z=X/\langle\delta\rangle\to X/D_8=E_{\mathrm D}\) 为次数 \(4\) 的 \(V_4\)-Galois 覆盖。

第 \(3\) 点令 \(L:=\QQ(X)\)，则
\[
\QQ(X_V)=L^{V_4^{\mathrm{norm}}},\qquad
\QQ(Y)=L^{\langle c\rangle},\qquad
\QQ(E_{\mathrm D})=L^{D_8}.
\]
由 Galois 对应性，
\[
L^{V_4^{\mathrm{norm}}}\cap L^{\langle c\rangle}
=L^{\langle V_4^{\mathrm{norm}},\langle c\rangle\rangle}
=L^{D_8}
=\QQ(E_{\mathrm D}),
\]
\[
L^{V_4^{\mathrm{norm}}}\cdot L^{\langle c\rangle}
=L^{V_4^{\mathrm{norm}}\cap\langle c\rangle}
=L^{\langle\delta\rangle}
=\QQ(Z).
\]
由“纤维积规范化的函数域等于子域合成”即得结论。证毕。
\end{proof}

\begin{conclusion}[三条二次子覆盖的分支奇偶对称差律]\label{cor:xi-terminal-zm-s4-v4-three-quad-subcovers-branch-parity}
在结论 \ref{con:xi-terminal-zm-s4-delta-quotient-v4-fiberproduct} 的 \(V_4\)-Galois 覆盖
\(
Z\to E_{\mathrm D}
\)
中，存在且仅存在三条非平凡二次子覆盖：
\[
X_V\to E_{\mathrm D},\qquad
Y\to E_{\mathrm D},\qquad
X_H\to E_{\mathrm D},
\]
其中可取 \(H=\langle\delta,\tau\rangle\le D_8\)（\(\tau\) 为 \(D_8\) 内换位），并置 \(X_H:=X/H\)。则：
\begin{enumerate}
  \item \(g(X_H)=2\)，且 \(X_H\to E_{\mathrm D}\) 为分支次数 \(2\) 的二次覆盖；
  \item 若
  \[
  \QQ(X_V)=\QQ(E_{\mathrm D})(\sqrt{f_V}),\qquad
  \QQ(Y)=\QQ(E_{\mathrm D})(\sqrt{f_Y}),
  \]
  则
  \[
  \QQ(X_H)=\QQ(E_{\mathrm D})(\sqrt{f_Vf_Y}),
  \]
  并满足
  \[
  \mathrm{Branch}(X_H/E_{\mathrm D})
  \equiv
  \mathrm{Branch}(X_V/E_{\mathrm D})\triangle \mathrm{Branch}(Y/E_{\mathrm D})
  \pmod 2.
  \]
  \item 设 \(D_V,D_Y,D_H\) 为三条二次覆盖的约化分支集合。
  已知 \(\deg D_V=\deg D_Y=6,\ \deg D_H=2\)，故
  \[
  \deg(D_V\cap D_Y)=5,
  \]
  即 \(D_V,D_Y\) 共有 \(5\) 个分支点，且各自仅有 \(1\) 个独有分支点；这两个独有点恰构成 \(D_H\)。
\end{enumerate}
\end{conclusion}
\begin{proof}
对第 \(1\) 点，取 \(H=\langle\delta,\tau\rangle=\{1,\delta,\tau,\delta\tau\}\)。
由结论 \ref{con:xi-terminal-zm-s4-conjugacy-class-fixedpoints-and-cyclic-quotient-genera}，
\[
\#\mathrm{Fix}_X(\delta)=2,\qquad
\#\mathrm{Fix}_X(\tau)=\#\mathrm{Fix}_X(\delta\tau)=10.
\]
群作用 Hurwitz 公式给出
\[
2g(X)-2=|H|\bigl(2g(X_H)-2\bigr)+\sum_{1\neq h\in H}\#\mathrm{Fix}_X(h),
\]
即
\[
30=4\bigl(2g(X_H)-2\bigr)+(2+10+10),
\]
故 \(g(X_H)=2\)。
又 \([D_8:H]=2\)，故 \(X_H\to E_{\mathrm D}\) 为二次覆盖；由 \(g(E_{\mathrm D})=1\) 得分支次数 \(2g(X_H)-2=2\)。

第 \(2\) 点是双二次扩张的局部赋值判别：在任意离散赋值 \(v\) 下，\(\QQ(E_{\mathrm D})(\sqrt{f})/\QQ(E_{\mathrm D})\) 分歧当且仅当 \(v(f)\) 为奇数。对 \(f_V,f_Y,f_Vf_Y\) 逐点比较即得模 \(2\) 对称差律。

第 \(3\) 点由第 \(1\) 点与已有 \(g(X_V)=g(Y)=4\) 立即得到 \(\deg D_V=\deg D_Y=6,\deg D_H=2\)。
在特征 \(0\) 下二次分支为简单分支，故
\[
2=\deg(D_V\triangle D_Y)=\deg D_V+\deg D_Y-2\deg(D_V\cap D_Y)=12-2\deg(D_V\cap D_Y),
\]
从而 \(\deg(D_V\cap D_Y)=5\)。证毕。
\end{proof}

\begin{conclusion}[\(\delta\)-不变/反不变分解与 \(\mathrm{Jac}(Z)\)、\(\mathrm{Prym}(X/Z)\) 同源型]\label{con:xi-terminal-zm-s4-delta-eigensplitting-jacz-prymxz}
在结论 \ref{con:xi-terminal-zm-s4-delta-quotient-v4-fiberproduct} 的记号下，存在定义于 \(\QQ\) 上的同源分解
\[
\mathrm{Jac}(Z)\sim_{\QQ}\mathrm{Jac}(X_A)\times E\times E_{\mathrm D}^2\times B,
\]
以及
\[
\mathrm{Prym}(X/Z)\sim_{\QQ}E^2\times B^2.
\]
\end{conclusion}
\begin{proof}
记 \(p:X\to Z=X/\langle\delta\rangle\)，并在 \(\QQ[\langle\delta\rangle]\) 中取幂等元
\[
e_+:=\frac{1+\delta}{2},\qquad e_-:=\frac{1-\delta}{2}.
\]
标准理论给出
\[
\mathrm{Jac}(X)\sim_{\QQ}e_+\mathrm{Jac}(X)\times e_-\mathrm{Jac}(X),
\]
且
\(
e_+\mathrm{Jac}(X)\sim_{\QQ}\mathrm{Jac}(Z)
\),
\(
e_-\mathrm{Jac}(X)\sim_{\QQ}\mathrm{Prym}(X/Z)
\)。

再由结论 \ref{con:xi-terminal-zm-s4-jacobian-group-algebra-prym-e3-b3}
\[
\mathrm{Jac}(X)\sim_{\QQ}\mathrm{Jac}(X_A)\times E^3\times E_{\mathrm D}^2\times B^3.
\]
对双换位 \(\delta\)，其在 \(\mathrm{sgn},V_2\) 上取值为 \(+1\)，在 \(V_3,V_3\otimes\mathrm{sgn}\) 上迹为 \(-1\)，故 \(+1\)-特征维数分别为
\[
1,\ 2,\ 1,\ 1.
\]
于是 \(e_+\)-像分别贡献
\[
\mathrm{Jac}(X_A),\ E_{\mathrm D}^2,\ E,\ B,
\]
而 \(e_-\)-像分别贡献
\[
0,\ 0,\ E^2,\ B^2.
\]
合并即得结论。证毕。
\end{proof}

\begin{conclusion}[属 \(2\) 中间商曲线的 Jacobian 分裂]\label{con:xi-terminal-zm-s4-xh-jacobian-splitting-ed-e}
取
\[
H=\langle\delta,\tau\rangle\le D_8,\qquad X_H:=X/H
\]
（\(\tau\) 为 \(D_8\) 中换位）。则 \(g(X_H)=2\)，且
\[
\mathrm{Jac}(X_H)\sim_{\QQ}E\times E_{\mathrm D}.
\]
特别地，\(X_H\to E_{\mathrm D}\) 为分支次数 \(2\) 的二次覆盖，且 \(X_H\) 为双椭圆型属 \(2\) 曲线。
\end{conclusion}
\begin{proof}
由结论 \ref{cor:xi-terminal-zm-s4-v4-three-quad-subcovers-branch-parity} 已知 \(g(X_H)=2\)。
令
\[
e_H:=\frac{1}{|H|}\sum_{h\in H}h\in\QQ[H].
\]
则 \(e_H\mathrm{Jac}(X)\sim_{\QQ}\mathrm{Jac}(X_H)\)。
对任一不可约表示 \(\rho\)，有
\[
\dim(\rho^H)=\frac14\!\left(\chi_\rho(1)+\chi_\rho(\delta)+2\chi_\rho(\text{换位})\right).
\]
代入 \(S_4\) 特征标值得
\[
\dim(\mathrm{sgn}^H)=0,\quad
\dim((V_3\otimes\mathrm{sgn})^H)=0,\quad
\dim(V_3^H)=1,\quad
\dim(V_2^H)=1.
\]
故 \(\mathrm{Jac}(X_A)\) 与 \(B^3\) 不贡献不变部分，而 \(E^3\) 与 \(E_{\mathrm D}^2\) 各贡献一维椭圆因子，遂有
\[
\mathrm{Jac}(X_H)\sim_{\QQ}E\times E_{\mathrm D}.
\]
其余断言由 \(g(X_H)=2\) 与该分裂立即得到。证毕。
\end{proof}

\begin{conclusion}[三次解式判别式与 \texorpdfstring{$\Delta(y)$}{Delta(y)} 的一致性]\label{con:xi-terminal-zm-s4-resolvent-cubic-discriminant-equals-delta}
在结论 \ref{con:xi-terminal-zm-s4-normal-subgroup-quotients-explicit-models-genus} 的记号下，恒有
\[
\mathrm{Disc}_z\!\bigl(\mathscr R(z,y)\bigr)=\Delta(y)=-y(y-1)g(y).
\]
\end{conclusion}
\begin{proof}[证明要点]
将 \(\mathscr R=z^3+az^2+bz+c\) 代入三次判别式闭式
\[
\mathrm{Disc}_z(\mathscr R)=a^2b^2-4b^3-4a^3c-27c^2+18abc,
\]
其中
\[
a=2y+1,\quad b=-(4y^2+4y+1),\quad c=-(8y^3+13y^2+5y+1),
\]
直接化简即得 \(-y(y-1)g(y)\)。
该恒等式亦与结论 \ref{con:xi-terminal-zm-discriminant-leyang}、\ref{con:xi-terminal-zm-s4-c3-fiberproduct-normalization} 的二次分辨数据一致。证毕。
\end{proof}

\begin{conclusion}[中间商曲线属数谱与 Jacobian 的表示论分裂]\label{con:xi-terminal-zm-s4-quotient-genus-spectrum-and-jacobian-representation-splitting}
对任意子群 \(H\le S_4\)，有
\[
g(X/H)=\dim_{\CC}H^0(X,\Omega_X^1)^H.
\]
由结论 \ref{con:xi-terminal-zm-s4-holomorphic-oneforms-chevalley-weil-rigid-decomposition} 可直接得到
\[
g(X/A_4)=2,\quad
g(X/V_4^{\mathrm{norm}})=4,\quad
g(X/S_3)=1,\quad
g(X/C_3)=6,\quad
g(X/D_8)=1.
\]
此外，存在定义在 \(\QQ\) 上的阿贝尔簇
\[
A_{\mathrm{sgn}},\ A_{V_3},\ A_{V_3\otimes\mathrm{sgn}},\ A_{V_2},
\]
使
\[
\mathrm{Jac}(X)\sim_{\QQ}
A_{\mathrm{sgn}}\times A_{V_3}\times A_{V_3\otimes\mathrm{sgn}}\times A_{V_2},
\]
并且
\[
\dim A_{\mathrm{sgn}}=2,\quad
\dim A_{V_3}=3,\quad
\dim A_{V_3\otimes\mathrm{sgn}}=9,\quad
\dim A_{V_2}=2.
\]
同时有规范同源识别
\[
A_{\mathrm{sgn}}\sim_{\QQ}\mathrm{Jac}(X/A_4),\qquad
A_{\mathrm{sgn}}\times A_{V_2}\sim_{\QQ}\mathrm{Jac}(X/V_4^{\mathrm{norm}}).
\]
\end{conclusion}
\begin{proof}[证明要点]
恒等式 \(H^0(X/H,\Omega^1)\cong H^0(X,\Omega^1)^H\) 给出第一式。
将
\[
H^0(X,\Omega_X^1)\cong
2\,\mathrm{sgn}\oplus V_3\oplus 3(V_3\otimes\mathrm{sgn})\oplus V_2
\]
限制到各子群并取不变维数，即得所列属数（与结论
\ref{con:xi-terminal-zm-s4-hurwitz-subgroup-spectrum-cycletypes} 一致）。
\par
对 Jacobian，按 \(\QQ[S_4]\) 中心幂等元作用得到 isotypic 分裂；由于 \(S_4\) 不可约特征均为有理值，各因子可在 \(\QQ\) 上定义。
维数由重数乘表示维数给出：
\[
2\cdot 1,\ 1\cdot 3,\ 3\cdot 3,\ 1\cdot 2,
\]
即 \((2,3,9,2)\)。
当 \(H\) 正规（如 \(A_4,V_4^{\mathrm{norm}}\)）时，\(\mathrm{Jac}(X/H)\) 正对应于在 \(H\) 上平凡的表示分量之直和，故得两条规范同源识别。证毕。
\end{proof}
\input{con__xi-terminal-zm-s4-equivariant-weil-artin-conductor-langlands}
\input{con__xi-terminal-zm-s4-v4-ed-branch-prym-kanirosen-localfactor}
\input{con__xi-terminal-zm-s4-nielsen-braid-monodromy-psp4f3-jac3torsion}

\paragraph{\texorpdfstring{$S_4$}{S4} 典范表示的共轭类谱型与低次数关系模的等变闭式}
沿用前文记号，设 \(X=\widetilde{X}\) 且
\[
V:=H^0(X,K_X),\qquad \dim V=16.
\]
并记 \(S:=\mathrm{Sym}(V)\)，\(I\subset S\) 为典范嵌入的齐次理想，\(I_m:=I\cap \mathrm{Sym}^mV\)。

\begin{conclusion}[典范表示在共轭类上的特征多项式刚性确定]\label{con:xi-terminal-zm-s4-canonical-characteristic-polynomial-conjugacy-rigidity}
对 \(S_4\) 的共轭类代表元 \(g\)（按循环型）在 \(V\) 上的特征多项式为
\[
\chi_{V,g}(T)=
\begin{cases}
(T-1)^{16}, & g=1,\\[1mm]
(T-1)^6(T+1)^{10}, & g\sim(12),\\[1mm]
(T-1)^8(T+1)^8, & g\sim(12)(34),\\[1mm]
(T-1)^6(T^2+T+1)^5, & g\sim(123),\\[1mm]
(T-1)^4(T+1)^4(T^2+1)^4, & g\sim(1234).
\end{cases}
\]
因此，典范表示在各共轭类上的谱型由上述多项式唯一确定。
\end{conclusion}
\begin{proof}[证明要点]
由结论 \ref{con:xi-terminal-zm-s4-chevalley-weil-explicit-decomposition}
\[
V\cong 2\varepsilon\oplus \rho_2\oplus \rho_3\oplus 3\rho_3'
\]
可得类函数迹
\[
\tr(g\mid V)=
\begin{cases}
16,& g=1,\\
-4,& g\sim(12),\\
0,& g\sim(12)(34),\\
1,& g\sim(123),\\
0,& g\sim(1234).
\end{cases}
\]
对阶 \(2\) 元，特征值仅为 \(\pm1\)，由迹与维数可解出重数，得到前两类式子。
对阶 \(3\) 元，特征值属于 \(\{1,\omega,\omega^2\}\)；又 \(\tr(g\mid V)=\tr(g^2\mid V)=1\)，故 \(1\) 的重数为 \(6\)，\(\omega,\omega^2\) 的重数各为 \(5\)。
对阶 \(4\) 元，特征值属于 \(\{1,-1,i,-i\}\)；由 \(\tr(g\mid V)=0\) 与 \(\tr(g^2\mid V)=\tr((12)(34)\mid V)=0\) 解得四者重数均为 \(4\)。证毕。
\end{proof}

\begin{conclusion}[对称代数的类函数型 Molien 生成函数与等型分量级数]\label{con:xi-terminal-zm-s4-symmetric-algebra-molien-isotypic-series}
对任意 \(g\in S_4\)，定义
\[
\Phi_g(t):=\sum_{m\ge0}\tr\!\bigl(g\mid \mathrm{Sym}^mV\bigr)t^m.
\]
则 \(\Phi_g(t)\) 仅依赖于 \(g\) 的共轭类，并有闭式
\[
\Phi_g(t)=
\begin{cases}
(1-t)^{-16}, & g=1,\\[1mm]
(1-t)^{-6}(1+t)^{-10}, & g\sim(12),\\[1mm]
(1-t)^{-8}(1+t)^{-8}, & g\sim(12)(34),\\[1mm]
(1-t)^{-6}(1+t+t^2)^{-5}, & g\sim(123),\\[1mm]
(1-t)^{-4}(1+t)^{-4}(1+t^2)^{-4}, & g\sim(1234).
\end{cases}
\]
进一步，对每个不可约表示
\(\sigma\in\{\mathbf 1,\varepsilon,\rho_2,\rho_3,\rho_3'\}\)，定义
\[
M_\sigma(t):=\sum_{m\ge0}\mathrm{mult}_\sigma\!\bigl(\mathrm{Sym}^mV\bigr)t^m.
\]
则有
\begin{align*}
M_{\mathbf 1}(t)
&=\frac1{24}(1-t)^{-16}
+\frac14(1-t)^{-6}(1+t)^{-10}
+\frac18(1-t)^{-8}(1+t)^{-8} \\
&\quad
+\frac13(1-t)^{-6}(1+t+t^2)^{-5}
+\frac14(1-t)^{-4}(1+t)^{-4}(1+t^2)^{-4},\\[1mm]
M_{\varepsilon}(t)
&=\frac1{24}(1-t)^{-16}
-\frac14(1-t)^{-6}(1+t)^{-10}
+\frac18(1-t)^{-8}(1+t)^{-8} \\
&\quad
+\frac13(1-t)^{-6}(1+t+t^2)^{-5}
-\frac14(1-t)^{-4}(1+t)^{-4}(1+t^2)^{-4},\\[1mm]
M_{\rho_2}(t)
&=\frac1{12}(1-t)^{-16}
+\frac14(1-t)^{-8}(1+t)^{-8}
-\frac13(1-t)^{-6}(1+t+t^2)^{-5},\\[1mm]
M_{\rho_3}(t)
&=\frac18(1-t)^{-16}
+\frac14(1-t)^{-6}(1+t)^{-10}
-\frac18(1-t)^{-8}(1+t)^{-8}
-\frac14(1-t)^{-4}(1+t)^{-4}(1+t^2)^{-4},\\[1mm]
M_{\rho_3'}(t)
&=\frac18(1-t)^{-16}
-\frac14(1-t)^{-6}(1+t)^{-10}
-\frac18(1-t)^{-8}(1+t)^{-8}
+\frac14(1-t)^{-4}(1+t)^{-4}(1+t^2)^{-4}.
\end{align*}
\end{conclusion}
\begin{proof}[证明要点]
对任意线性算子 \(A\)，恒有
\[
\sum_{m\ge0}\tr\!\bigl(\mathrm{Sym}^m(A)\bigr)t^m
=\det(I-tA)^{-1}.
\]
将结论 \ref{con:xi-terminal-zm-s4-canonical-characteristic-polynomial-conjugacy-rigidity} 的各类特征值代入即可得到 \(\Phi_g(t)\)。
再由特征标正交关系
\[
M_\sigma(t)=\frac1{|S_4|}\sum_{[g]}|[g]|\,\chi_\sigma(g)\,\Phi_g(t),
\]
代入 \(S_4\) 的共轭类大小与特征标值，化简即得各显式式。证毕。
\end{proof}

\begin{conclusion}[三次与四次对称幂及典范理想低次分量的 \texorpdfstring{$S_4$}{S4} 分解]\label{con:xi-terminal-zm-s4-sym3-sym4-canonical-ideal-low-degree-decomposition}
有如下等变分解：
\[
\mathrm{Sym}^3V
\cong
25\mathbf 1\oplus 47\varepsilon\oplus 66\rho_2\oplus 91\rho_3\oplus 113\rho_3',
\]
\[
\mathrm{Sym}^4V
\cong
198\mathbf 1\oplus 138\varepsilon\oplus 330\rho_2\oplus 508\rho_3\oplus 452\rho_3'.
\]
并且
\[
I_3
\cong
23\mathbf 1\oplus 43\varepsilon\oplus 60\rho_2\oplus 83\rho_3\oplus 102\rho_3',
\]
\[
I_4
\cong
192\mathbf 1\oplus 135\varepsilon\oplus 321\rho_2\oplus 494\rho_3\oplus 440\rho_3'.
\]
\end{conclusion}
\begin{proof}[证明要点]
由结论 \ref{con:xi-terminal-zm-s4-symmetric-algebra-molien-isotypic-series}，
\[
\mathrm{mult}_\sigma\!\bigl(\mathrm{Sym}^mV\bigr)=[t^m]M_\sigma(t),
\]
取 \(m=3,4\) 即得 \(\mathrm{Sym}^3V,\mathrm{Sym}^4V\) 分解。
另一方面，由结论 \ref{con:xi-terminal-zm-s4-canonical-projective-normality-quadratic-generation} 的射影正规性，
\[
\mathrm{Sym}^mV\twoheadrightarrow H^0(X,mK_X)\qquad (m\ge 0),
\]
故
\[
0\to I_m\to \mathrm{Sym}^mV\to H^0(X,mK_X)\to 0
\]
为 \(S_4\)-等变正合列。由结论 \ref{con:xi-terminal-zm-s4-pluricanonical-irrep-decomposition-all-n} 在
\(n=3,4\) 时
\[
H^0(X,3K_X)\cong
2\mathbf 1\oplus 4\varepsilon\oplus 6\rho_2\oplus 8\rho_3\oplus 11\rho_3',
\]
\[
H^0(X,4K_X)\cong
6\mathbf 1\oplus 3\varepsilon\oplus 9\rho_2\oplus 14\rho_3\oplus 12\rho_3',
\]
在表示环中作差即得 \(I_3,I_4\) 分解。证毕。
\end{proof}

\begin{conclusion}[二次生成下三次层线性 syzygy 模的完全分解]\label{con:xi-terminal-zm-s4-linear-syzygy-module-complete-decomposition}
记
\[
Q:=I_2,\qquad
\mathcal S_3:=\ker\!\bigl(Q\otimes V\to I_3\bigr).
\]
则
\[
\mathcal S_3
\cong
27\mathbf 1\oplus 29\varepsilon\oplus 61\rho_2\oplus 88\rho_3\oplus 91\rho_3'.
\]
\end{conclusion}
\begin{proof}[证明要点]
由结论 \ref{con:xi-terminal-zm-s4-canonical-projective-normality-quadratic-generation}，典范理想由二次式生成，故 \(Q\otimes V\twoheadrightarrow I_3\) 为 \(S_4\)-等变满射。
结合结论 \ref{con:xi-terminal-zm-s4-canonical-quadrics-equivariant-decomposition} 的
\[
Q\cong 8\mathbf 1\oplus 2\varepsilon\oplus 9\rho_2\oplus 13\rho_3\oplus 8\rho_3'
\]
与结论 \ref{con:xi-terminal-zm-s4-chevalley-weil-explicit-decomposition} 的 \(V\) 分解，按特征标点乘可得
\[
Q\otimes V
\cong
50\mathbf 1\oplus 72\varepsilon\oplus 121\rho_2\oplus 171\rho_3\oplus 193\rho_3'.
\]
再由结论 \ref{con:xi-terminal-zm-s4-sym3-sym4-canonical-ideal-low-degree-decomposition} 的 \(I_3\) 分解，在表示环中作差
\(
[\mathcal S_3]=[Q\otimes V]-[I_3]
\)
即得结论。证毕。
\end{proof}

\begin{conclusion}[四次层二次系数核表示的闭式分解]\label{con:xi-terminal-zm-s4-quadratic-coefficient-kernel-degree4-closed-decomposition}
定义
\[
\mathcal K_4:=\ker\!\bigl(Q\otimes \mathrm{Sym}^2V\to I_4\bigr).
\]
则有
\[
\mathcal K_4
\cong
375\mathbf 1\oplus 344\varepsilon\oplus 724\rho_2\oplus 1090\rho_3\oplus 1056\rho_3'.
\]
\end{conclusion}
\begin{proof}[证明要点]
仍由二次生成性，\(Q\otimes \mathrm{Sym}^2V\twoheadrightarrow I_4\) 为 \(S_4\)-等变满射。
由结论 \ref{con:xi-terminal-zm-s4-canonical-quadrics-equivariant-decomposition} 与
\ref{con:xi-terminal-zm-s4-symmetric-algebra-molien-isotypic-series} 在 \(m=2\) 的系数可得
\[
\mathrm{Sym}^2V
\cong
11\mathbf 1\oplus 3\varepsilon\oplus 13\rho_2\oplus 20\rho_3\oplus 12\rho_3'.
\]
特征标点乘后得到
\[
Q\otimes \mathrm{Sym}^2V
\cong
567\mathbf 1\oplus 479\varepsilon\oplus 1045\rho_2\oplus 1584\rho_3\oplus 1496\rho_3'.
\]
再由结论 \ref{con:xi-terminal-zm-s4-sym3-sym4-canonical-ideal-low-degree-decomposition} 的 \(I_4\) 分解，作差
\(
[\mathcal K_4]=[Q\otimes \mathrm{Sym}^2V]-[I_4]
\)
即得所示重数。证毕。
\end{proof}

\endinput


\subsubsection{谱兼容外置、半直积修复与有限/无穷二影刚性}\label{subsubsec:xi-spectrum-compatible-semidirect-two-shadow-rigidity}
本段在 Lee--Yang 四层覆盖的既有单值化框架上，将“可交换外置障碍—非交换修复—线性椭圆双影刚性”组织为同一条可检验链条。

\begin{theorem}[可交换外置寄存器对非交换单值群的不可完备性]\label{thm:xi-terminal-abelian-register-obstruction-nonabelian-monodromy}
设 \(U\) 为道路连通拓扑空间，\(\pi:X\to U\) 为连通有限分支覆盖，分支集记为 \(B\subset U\)，并取基点 \(u_0\in U\setminus B\)。
令
\[
\rho:\pi_1(U\setminus B,u_0)\longrightarrow \mathfrak S_n
\]
为单值表示，且假设 \(\rho(\pi_1(U\setminus B,u_0))\) 非交换。

称三元组 \((A,\iota,D)\) 为该覆盖的一组可交换外置寄存器数据，若：
\begin{enumerate}
  \item \(A\) 为交换群；
  \item \(\iota:\pi_1(U\setminus B,u_0)\to A\) 为群同态；
  \item \(D:\{1,\dots,n\}\times A\to\{1,\dots,n\}\) 满足对任意 \(\gamma\in\pi_1(U\setminus B,u_0)\) 与 \(i\in\{1,\dots,n\}\),
  \[
  \rho(\gamma)(i)=D\bigl(i,\iota(\gamma)\bigr).
  \]
\end{enumerate}
则此类三元组不存在。
\end{theorem}
\begin{proof}
由 \(\rho(\pi_1(U\setminus B,u_0))\) 非交换，可取 \(\alpha,\beta\in\pi_1(U\setminus B,u_0)\) 使
\[
\rho([\alpha,\beta])\neq \mathrm{id},
\qquad
[\alpha,\beta]:=\alpha\beta\alpha^{-1}\beta^{-1}.
\]
由 \(A\) 交换且 \(\iota\) 为同态，
\[
\iota([\alpha,\beta])
=\iota(\alpha)\iota(\beta)\iota(\alpha)^{-1}\iota(\beta)^{-1}
=e_A
=\iota(1).
\]
故对任意 \(i\),
\[
D\bigl(i,\iota([\alpha,\beta])\bigr)=D(i,e_A)=D\bigl(i,\iota(1)\bigr).
\]
左侧按定义等于 \(\rho([\alpha,\beta])(i)\)，右侧等于 \(\rho(1)(i)=i\)，从而 \(\rho([\alpha,\beta])(i)=i\) 对所有 \(i\) 成立，矛盾于 \(\rho([\alpha,\beta])\neq\mathrm{id}\)。
证毕。
\end{proof}

\paragraph{window-6 三轴中心约束、Pati--Salam 连通实现与标准模型包络的中心 2-挠兼容性}
在 window-6 的 boundary 读出中，三轴独立 sheet-flip 给出一个抽象中心寄存器
\[
\mathcal C_{\mathrm{bdry}}\cong (\ZZ/2\ZZ)^3.
\]
本段讨论其与紧连通 Lie 包络在中心 2-挠层面的可实现条件。

\begin{theorem}[标准模型连通包络对三轴中心寄存器的不可实现性]\label{thm:xi-window6-sm-envelope-center-2torsion-obstruction}
设 \(G\) 为紧连通 Lie 群，满足
\[
\Lie(G)\cong \mathfrak u(1)\oplus\mathfrak{su}(2)\oplus\mathfrak{su}(3).
\]
则
\[
\dim_{\FF_2} Z(G)[2]\le 2.
\]
从而不存在单同态
\[
\mathcal C_{\mathrm{bdry}}\cong (\ZZ/2\ZZ)^3\hookrightarrow Z(G)[2].
\]
\end{theorem}
\begin{proof}
由紧连通群分类，可取
\[
G\cong \widetilde G/\Gamma,\qquad
\widetilde G:=U(1)\times SU(2)\times SU(3),
\]
其中 \(\Gamma\subset Z(\widetilde G)\) 为有限中心子群。记
\[
A:=Z(\widetilde G)\cong U(1)\times \ZZ_2\times \ZZ_3,
\qquad
Z(G)\cong A/\Gamma.
\]
设 \(\pi:A\to \ZZ_2\times \ZZ_3\) 为自然投影，则有短正合列
\[
1\to U(1)\to A\to \ZZ_2\times \ZZ_3\to 1.
\]
对 \(\Gamma\) 取商后得到
\[
1\to T\to A/\Gamma\to Q\to 1,
\]
其中 \(T\cong U(1)\)（一维紧环面），\(Q\) 为 \((\ZZ_2\times \ZZ_3)\) 的商群。

记 \(E:=A/\Gamma\)。考虑 \(E[2]\to Q[2]\)。对任意 \(q\in Q[2]\)，其纤维中两点之差落在 \(T[2]\)。
故每个纤维大小至多 \(|T[2]|=2\)，并且
\[
|Q[2]|\le |(\ZZ_2\times \ZZ_3)[2]|=2.
\]
从而
\[
|E[2]|\le 2\cdot 2=4,\qquad
\dim_{\FF_2}E[2]\le 2.
\]
即 \(\dim_{\FF_2} Z(G)[2]\le 2\)。因此 \((\ZZ/2\ZZ)^3\) 不能嵌入 \(Z(G)[2]\)。证毕。
\end{proof}

\begin{theorem}[三轴 canonical sheet-flip 保真下降下的 Pati--Salam 连通实现唯一性]\label{thm:xi-window6-ps-faithful-canonical-descent-unique}
设
\[
\widetilde G_{\mathrm{PS}}:=SU(4)\times SU(2)_L\times SU(2)_R,\qquad
G:=\widetilde G_{\mathrm{PS}}/H,
\]
其中 \(H\subset Z(\widetilde G_{\mathrm{PS}})\) 为有限中心子群。记 canonical involution 子群
\[
\mathcal C_{\mathrm{can}}
:=Z(SU(4))[2]\times Z(SU(2)_L)\times Z(SU(2)_R)
\cong (\ZZ/2\ZZ)^3.
\]
若商映射 \(q:\widetilde G_{\mathrm{PS}}\to G\) 在 \(\mathcal C_{\mathrm{can}}\) 上限制为单同态
\[
q|_{\mathcal C_{\mathrm{can}}}:\mathcal C_{\mathrm{can}}\hookrightarrow Z(G)[2],
\]
则必有 \(H=\{e\}\)，因而
\[
G\cong SU(4)\times SU(2)_L\times SU(2)_R.
\]
\end{theorem}
\begin{proof}
有
\[
Z(\widetilde G_{\mathrm{PS}})
\cong \ZZ_4\times \ZZ_2\times \ZZ_2.
\]
取任意非平凡 \(H\subset Z(\widetilde G_{\mathrm{PS}})\)。任取 \(h\in H\setminus\{e\}\)。
若 \(\ord(h)=2\)，则 \(h\in \mathcal C_{\mathrm{can}}\)；
若 \(\ord(h)=4\)，则 \(h^2\neq e\) 且 \(h^2\in \mathcal C_{\mathrm{can}}\)。
故总有
\[
H\cap \mathcal C_{\mathrm{can}}\neq \{e\}.
\]
而 \(q\) 在 \(H\) 上恒等于单位元，故 \(q|_{\mathcal C_{\mathrm{can}}}\) 的核至少含
\(H\cap \mathcal C_{\mathrm{can}}\) 的非平凡元，不能单射，矛盾。
故 \(H\) 只能为平凡子群。证毕。
\end{proof}

\begin{theorem}[Pati--Salam 到标准模型包络的中心 2-挠压缩不可避免]\label{thm:xi-window6-ps-to-sm-center2-kernel}
设
\[
G_{\mathrm{PS}}:=SU(4)\times SU(2)_L\times SU(2)_R,
\]
且 \(G_{\mathrm{SM}}\) 为任意紧连通 Lie 群，满足
\(
\Lie(G_{\mathrm{SM}})\cong \mathfrak u(1)\oplus\mathfrak{su}(2)\oplus\mathfrak{su}(3)
\)。
若连续群同态
\(
f:G_{\mathrm{PS}}\to G_{\mathrm{SM}}
\)
满足中心兼容条件
\[
f\bigl(\mathcal C_{\mathrm{can}}\bigr)\subseteq Z(G_{\mathrm{SM}})[2],
\]
则限制映射
\[
f_*:\mathcal C_{\mathrm{can}}\longrightarrow Z(G_{\mathrm{SM}})[2]
\]
必有非平凡核。
\end{theorem}
\begin{proof}
由定理 \ref{thm:xi-window6-sm-envelope-center-2torsion-obstruction}，
\(
\dim_{\FF_2} Z(G_{\mathrm{SM}})[2]\le 2
\)。
而
\(
\mathcal C_{\mathrm{can}}\cong (\ZZ/2\ZZ)^3
\)
的 \(\FF_2\)-维数为 \(3\)。
故线性映射 \(f_*\) 的秩至多 \(2\)，其核维数至少 \(1\)，即核非平凡。证毕。
\end{proof}

\begin{proposition}[sheet parity 在 \(m=6\) 的严格锚点性]\label{prop:xi-window6-sheet-parity-anchor-m6}
\leanverified{paper\_xi\_window6\_sheet\_parity\_anchor\_m6}
在 dyadic boundary uplift 口径下，取
\[
\Delta_m^{\mathrm{bdry}}
:=\{n-V_m(w):\ n\in(\Fold_m^{\mathrm{bin}})^{-1}(w)\}
\]
（对 boundary \(w\) 不依赖于具体词）。
核验给出
\[
\Delta_6^{\mathrm{bdry}}=\{0,34\},\qquad
\Delta_7^{\mathrm{bdry}}=\{0,55,89\},\qquad
\Delta_8^{\mathrm{bdry}}=\{0,89,144\}.
\]
在“parity 由 uplift 数据无需额外选择即可恢复”的准则下，该结构仅在 \(m=6\) 成立。
\end{proposition}
\begin{proof}
若 parity 可由 \(\Delta_m^{\mathrm{bdry}}\) 唯一恢复，则每个 boundary 纤维必须存在唯一二元配对规则；
这要求非零 uplift 增量唯一。
当 \(m=6\) 时，\(\Delta_6^{\mathrm{bdry}}\) 仅含一个非零元 \(34\)，且 boundary 预像层数为 \(2\)，配对唯一。
当 \(m=7,8\) 时，\(\Delta_m^{\mathrm{bdry}}\) 各含两个不同非零元且预像层数为 \(3\)；
任意二元化都需额外选择（例如忽略某一层），不再由数据唯一决定。证毕。
\end{proof}

\begin{theorem}[纤维内自由 involution 的奇数障碍与最小覆盖度]\label{thm:xi-bdry-free-involution-odd-fiber-obstruction-min-cover}
\leanverified{paper\_xi\_bdry\_free\_involution\_odd\_fiber\_obstruction\_min\_cover}
设 \(f:\widetilde X\to X\) 为有限集合之间的满射。
若存在自由 involution \(\sigma:\widetilde X\to\widetilde X\) 满足
\[
\sigma^2=\mathrm{id},\qquad
\mathrm{Fix}(\sigma)=\varnothing,\qquad
f\circ\sigma=f,
\]
则对任意 \(x\in X\) 必有
\[
\bigl|f^{-1}(x)\bigr|\equiv 0\pmod 2.
\]
因此若存在 \(x\in X\) 使 \(|f^{-1}(x)|\) 为奇数，则不存在满足 \(f\circ\sigma=f\) 的自由 involution。
\par
进一步，若希望通过一个覆盖 \(\pi:Y\to\widetilde X\) 在 \(Y\) 上实现同类的纤维内自由 involution（即 \((f\circ\pi)\circ\widetilde\sigma=f\circ\pi\) 且 \(\widetilde\sigma\) 自由），
则覆盖度必须为偶数，因而 \(\deg(\pi)\ge 2\)。
\end{theorem}
\begin{proof}
若 \(\sigma\) 自由，则其轨道分解全由二元轨道组成，故 \(|S|\) 为偶数对任何 \(\sigma\)-不变有限子集 \(S\subseteq \widetilde X\) 成立。
取 \(S=f^{-1}(x)\)。由 \(f\circ\sigma=f\) 得 \(S\) 为 \(\sigma\)-不变，故 \(|f^{-1}(x)|\) 必为偶数；若存在奇数纤维则矛盾。
\par
对覆盖情形：若 \(|f^{-1}(x)|\) 为奇数，则 \(|(f\circ\pi)^{-1}(x)|=\deg(\pi)\cdot |f^{-1}(x)|\) 与 \(\deg(\pi)\) 同奇偶。
而在 \(Y\) 上实现纤维内自由 involution 仍要求每条纤维基数为偶数，故 \(\deg(\pi)\) 必为偶数，特别地 \(\deg(\pi)\ge 2\)。证毕。
\end{proof}

\begin{corollary}[boundary 三层 uplift 的二值外置下界]\label{cor:xi-bdry-uplift-three-layer-needs-one-bit}
在 dyadic boundary uplift 口径下，若 \(m\in\{7,8\}\)，则由命题 \ref{prop:xi-window6-sheet-parity-anchor-m6} 的
\(
\Delta_m^{\mathrm{bdry}}=\{0,\cdot,\cdot\}
\)
可知 boundary 预像层数为 \(3\)，从而存在 boundary 词 \(w\) 使
\[
\bigl|(\Fold_m^{\mathrm{bin}})^{-1}(w)\bigr|=3.
\]
因此不存在在 boundary uplift 纤维内闭合的自由 \(\ZZ_2\) involution；
并且任何试图在覆盖替代后实现该自由性的数据结构，其最小覆盖度满足 \(\deg(\pi)\ge 2\)，从而至少需要引入一个二值外置轴以消除奇数纤维的轨道障碍。
\end{corollary}
\begin{proof}
由命题 \ref{prop:xi-window6-sheet-parity-anchor-m6}，当 \(m=7,8\) 时 boundary 预像层数为 \(3\)，故存在边界纤维基数为 \(3\)。
将定理 \ref{thm:xi-bdry-free-involution-odd-fiber-obstruction-min-cover} 应用于 \(f=(\Fold_m^{\mathrm{bin}})|_{\widetilde X_m^{\mathrm{bdry}}}\) 即得结论。证毕。
\end{proof}

\begin{theorem}[全对称不变性强制的 \texorpdfstring{$\ZZ_2$}{Z/2} 分级平凡化]\label{thm:xi-sym-invariant-z2-grading-trivial}\leanverified{paper\_xi\_sym\_invariant\_z2\_grading\_trivial}
设 \(\Delta\) 为有限集合，\(|\Delta|=k\ge 3\)。
将一个 \(\ZZ_2\)-分级理解为选取“奇层子集” \(A\subset\Delta\)（允许 \(A=\varnothing\) 或 \(A=\Delta\) 的平凡分级）。
若要求该分级在全对称群 \(\mathrm{Sym}(\Delta)\cong S_k\) 的自然作用下不变，即对任意 \(\sigma\in S_k\) 都有 \(\sigma(A)=A\)，则该分级必为平凡分级：\(A=\varnothing\) 或 \(A=\Delta\)。
\end{theorem}
\begin{proof}
由于 \(S_k\) 在 \(\Delta\) 上传递，若 \(A\neq\varnothing\)，取 \(x\in A\)。
对任意 \(y\in\Delta\) 存在 \(\sigma\in S_k\) 使 \(\sigma(x)=y\)。
由不变性 \(\sigma(A)=A\) 得 \(y\in A\)，从而 \(A=\Delta\)。
同理，若 \(A\neq\Delta\) 则只能 \(A=\varnothing\)。
故除平凡分级外不存在 \(S_k\)-不变的 \(\ZZ_2\) 分级。证毕。
\end{proof}

\begin{theorem}[左--右交换下的规范中心模与不可下降通道]\label{thm:xi-window6-lr-swap-canonical-center-mod}
在 \(\mathcal C_{\mathrm{can}}\cong(\ZZ/2\ZZ)^3\) 中取基
\[
a:=(-I_4,1,1),\qquad b:=(1,-I_2,1),\qquad c:=(1,1,-I_2).
\]
令 \(\iota\) 为交换 \(SU(2)_L\leftrightarrow SU(2)_R\) 的自同构，则
\[
\iota(a)=a,\qquad \iota(b)=c,\qquad \iota(c)=b.
\]
定义满同态
\[
\pi_{\mathrm{LR}}:\mathcal C_{\mathrm{can}}\twoheadrightarrow \ZZ/2\ZZ,\qquad
\pi_{\mathrm{LR}}(\alpha a+\beta b+\gamma c):=\beta+\gamma.
\]
则
\[
\ker(\pi_{\mathrm{LR}})=\mathcal C_{\mathrm{can}}^{\iota}
\cong (\ZZ/2\ZZ)^2.
\]
因而任何要求“左--右交换在中心层不可见”的下降机制，必然至少压缩一个规范 \(\ZZ_2\) 模通道。
\end{theorem}
\begin{proof}
\(\pi_{\mathrm{LR}}\) 显然为群同态且满。
\[
\pi_{\mathrm{LR}}(\alpha,\beta,\gamma)=0
\iff \beta=\gamma.
\]
另一方面，
\[
\iota(\alpha,\beta,\gamma)=(\alpha,\gamma,\beta),
\]
故
\(
(\alpha,\beta,\gamma)\in\mathcal C_{\mathrm{can}}^{\iota}
\iff \beta=\gamma
\)。
两条件等价，故核恰为不动子群，且其维数为 \(2\)。证毕。
\end{proof}

\begin{conjecture}[dyadic boundary uplift 层数在 \(m=6\) 的唯一退化]\label{conj:xi-window6-bdry-uplift-unique-degeneration}
设 \(\Delta_m^{\mathrm{bdry}}\) 为 boundary uplift 的增量集合。猜想对所有 \(m\ge 7\) 有
\[
\Delta_m^{\mathrm{bdry}}=\{0,F_{m+3},F_{m+4}\},
\]
从而 boundary 覆盖层数恒为 \(3\)。
而 \(m=6\) 为唯一退化分辨率：
\[
\Delta_6^{\mathrm{bdry}}=\{0,F_9\}=\{0,34\},
\]
即唯一仅含单一非零增量的二层情形。
\end{conjecture}

\begin{conjecture}[异常与共振的通道稀疏性：由 boundary sheet parity 子群支配]\label{conj:xi-anomaly-resonance-channel-sparsity-controlled-by-boundary-parity}
在 window-\(6\) 的 boundary 读出中，三轴独立 sheet-flip 诱导规范中心寄存器
\(\mathcal C_{\mathrm{bdry}}\cong(\ZZ/2\ZZ)^3\)。
猜想：对所有足够大的分辨率 \(m\)，所有可见通道中的非平凡绕数/真共振指数（从而异常的可见非零性）只能出现在那些在 \(\mathcal C_{\mathrm{bdry}}\) 上取非平凡限制的特征标通道中。
等价地：若某一可见通道 \(\overline\chi\) 在 \(\mathcal C_{\mathrm{bdry}}\) 上平凡，则其可见散射绕数与相应的 Toeplitz 负平方指数均为零。
\par
该稀疏性一旦成立，将把“通道级异常定位/证书化”的搜索空间压缩到由 \(\mathcal C_{\mathrm{bdry}}\) 诱导的有限层筛选子族，并与定理 \ref{thm:xi-anom-winding-obstruction} 与定理 \ref{thm:xi-visible-anom-iff-radial-wallcrossing} 的可见通道判别口径形成闭合接口。
\end{conjecture}

\begin{remark}[可复现核验脚本]\label{rem:xi-window6-center-sheetflip-ps-sm-audit-script}
本段代数与组合核验由脚本
\texttt{scripts/exp\_xi\_window6\_center\_sheetflip\_ps\_sm\_audit.py}
给出，证书文件为
\texttt{artifacts/export/xi\_window6\_center\_sheetflip\_ps\_sm\_audit.json}。
\end{remark}


\begin{corollary}[Lee--Yang 四层覆盖中的可交换外置障碍]\label{cor:xi-terminal-abelian-register-obstruction-leyang-s4}
在定理 \ref{thm:xi-terminal-zm-leyang-elliptic-structure} 给出的 Lee--Yang 四层覆盖 \(\pi_y:E_{\min}\to\mathbb P^1_y\) 上，
\[
\mathrm{Mon}(\pi_y)\cong S_4
\]
为非交换群，故满足定理 \ref{thm:xi-terminal-abelian-register-obstruction-nonabelian-monodromy} 的不可完备性结论：任何要求路径拼接同态化的可交换外置寄存器都不能实现全局单值恢复。
\end{corollary}

\begin{theorem}[位移半直积 Gödel 寄存器对非交换算法轨迹的忠实外置]\label{thm:xi-terminal-semidir-godel-register-faithful}
\leanverified{paper\_xi\_terminal\_semidir\_godel\_register\_faithful}
设 \(\Sigma\) 为有限字母表（\(|\Sigma|\ge 2\)），\(\Sigma^\ast\) 为自由幺半群（串联为乘法），并取单射编码
\(
c:\Sigma\hookrightarrow \NN_{\ge 1}
\)。
记 \((p_i)_{i\ge 1}\) 为递增素数序列。

对 \(w=a_1\cdots a_n\in\Sigma^\ast\) 定义位序 Gödel 编码
\[
G(w):=\prod_{i=1}^{n}p_i^{\,c(a_i)},
\qquad
|w|=n.
\]
定义素数位移同态
\[
\mathrm{Sh}^{k}\!\left(\prod_{i\ge1}p_i^{e_i}\right):=\prod_{i\ge1}p_{i+k}^{e_i}.
\]
令
\[
\mathcal R:=\NN\times \NN_{>0},
\qquad
(n,A)\odot(m,B):=\bigl(n+m,\ A\cdot \mathrm{Sh}^{n}(B)\bigr).
\]
则映射
\[
J:\Sigma^\ast\to\mathcal R,\qquad
J(w):=(|w|,G(w))
\]
为单射幺半群同态，且 \((|w|,G(w))\) 可唯一恢复 \(w\)。
此外，对任意交换幺半群 \(M\)，不存在单射幺半群同态 \(\Sigma^\ast\hookrightarrow M\)。
\end{theorem}
\begin{proof}
\(\mathrm{Sh}^{n}\circ\mathrm{Sh}^{m}=\mathrm{Sh}^{n+m}\) 且每个 \(\mathrm{Sh}^n\) 为幺半群同态，故 \(\odot\) 结合，单位元为 \((0,1)\)。
对 \(u=a_1\cdots a_n\)、\(v=b_1\cdots b_m\)，
\[
G(uv)
=\left(\prod_{i=1}^{n}p_i^{c(a_i)}\right)\left(\prod_{j=1}^{m}p_{n+j}^{c(b_j)}\right)
=G(u)\,\mathrm{Sh}^{n}(G(v)),
\]
故 \(J(uv)=J(u)\odot J(v)\)。

单射性由唯一分解定理给出：\(J(w)=(n,G(w))\) 已给出长度 \(n\)，且
\(
v_{p_i}(G(w))=c(a_i)
\)
可逐位恢复 \(a_i\)（因 \(c\) 单射）。

若 \(f:\Sigma^\ast\to M\) 为到交换幺半群的同态，取 \(a\neq b\in\Sigma\)，则
\[
f(ab)=f(a)f(b)=f(b)f(a)=f(ba).
\]
若 \(f\) 单射则推出 \(ab=ba\)，与自由幺半群非交换性矛盾。证毕。
\end{proof}

% AUTO-GENERATED by scripts/exp_xi_prime_register_semidirect_two_shadow_audit.py
\begin{equation}\label{eq:xi_prime_register_semidirect_two_shadow_audit}
\begin{aligned}
\rho([\alpha,\beta])&=(2,3,1,4),\qquad \rho([\alpha,\beta])(1)\neq 1,\\
J(uv)&=J(u)\odot J(v)\ \text{(checked for all words with }|u|,|v|\le 4\text{)},\\
\left|O_2(\ZZ)\right|&=8,\qquad \left|SO_2(\ZZ)\right|=4.
\end{aligned}
\end{equation}


\begin{theorem}[Lee--Yang--Gödel 嵌入与素数寄存器的代数可识别性]\label{thm:xi-leyang-godel-prime-register-algebraic-identifiability}
令 \(\mathcal P\) 为素数集。记
\[
\NN^{(\mathcal P)}:=\{r:\mathcal P\to\NN:\ \#\mathrm{supp}(r)<\infty\},
\qquad
\mathrm{supp}(r):=\{p\in\mathcal P:\ r(p)\neq 0\}.
\]
并定义 Gödel 读数
\[
N(r):=\prod_{p\in\mathcal P}p^{r(p)}\in\NN_{\ge 1}.
\]
定义映射
\[
\mathcal L:\NN^{(\mathcal P)}\to \ZZ[z],
\qquad
\mathcal L(r)(z):=\prod_{p\in\mathcal P}\bigl(1+z^p\bigr)^{r(p)}.
\]
则成立：
\begin{enumerate}
  \item \(\mathcal L\) 为交换幺半群同态（以点加 \(r+s\) 与多项式乘法为幺半群运算）：
  \[
  \mathcal L(r+s)=\mathcal L(r)\mathcal L(s).
  \]
  \item \(\mathcal L\) 为单射。
  \item \(\mathcal L(r)\) 的全部零点均落在单位圆 \(S^1\)，且皆为偶阶单位根。
\end{enumerate}
\end{theorem}
\begin{proof}
第一条由指数加法与乘法分配律直接给出。
\par
第三条：任一因子 \(1+z^p\) 的零点满足 \(z^p=-1\)，从而 \(z^{2p}=1\)，故其零点均为偶阶单位根，亦在 \(S^1\) 上；乘积的零点集为各因子零点的并（按重数计），结论成立。
\par
第二条：对奇素数 \(p\) 有分解
\[
z^p+1=(z+1)\Phi_{2p}(z),
\]
且 \(\Phi_{2p}\) 在 \(\QQ\) 上不可约，并且不同 \(p\) 给出不同的不可约因子（可参见 \cite{Serre1973CourseArithmetic} 对 cyclotomic 多项式的基本性质与不可约性论述）。
对 \(p=2\)，有 \(z^2+1=\Phi_4(z)\)。
因此在 \(\mathcal L(r)\) 的不可约分解中，每个 \(\Phi_{2p}\)（\(p\) 为奇素数）出现的指数恰为 \(r(p)\)，而 \(\Phi_4\) 出现的指数为 \(r(2)\)，从而 \(r\) 由 \(\mathcal L(r)\) 唯一确定，故 \(\mathcal L\) 单射。
\end{proof}

\begin{theorem}[Gödel--Lee--Yang 对数振幅的逐素频全息反演]\label{thm:xi-leyang-logamplitude-prime-frequency-inversion}
在定理 \ref{thm:xi-leyang-godel-prime-register-algebraic-identifiability} 的记号下，对任意 \(r\in\NN^{(\mathcal P)}\) 定义单位圆上的对数振幅
\[
\mathcal A_r(\theta):=\log\abs{\mathcal L(r)(e^{i\theta})}\qquad(\theta\in[-\pi,\pi]).
\]
则 \(\mathcal A_r\in L^1([-\pi,\pi])\)。
并且对每个素数 \(p\) 成立严格反演公式
\[
\boxed{\quad
r(p)=\frac1\pi\int_{-\pi}^{\pi}\mathcal A_r(\theta)\cos(p\theta)\,\dd\theta.\quad}
\]
因此 \(\{\mathcal A_r(\theta)\}_{\theta\in[-\pi,\pi]}\) 在测度意义下逐素数唯一确定寄存器 \(r\)。
\end{theorem}
\begin{proof}
由于 \(r\) 仅有限支撑，
\[
\mathcal A_r(\theta)=\sum_{p\in\mathcal P}r(p)\log\abs{1+e^{ip\theta}}.
\]
对任意 \(\rho\in(0,1)\)，在单位圆盘内取主值对数并取实部，得到一致可积展开
\[
\Re\log(1+\rho e^{ix})
=\sum_{k\ge 1}\frac{(-1)^{k+1}}{k}\rho^k\cos(kx).
\]
令 \(\rho\uparrow 1\) 可得在 \(L^2([-\pi,\pi])\)（从而 \(L^1\)）意义下的余弦级数
\[
\log\abs{1+e^{ix}}
=\sum_{k\ge 1}\frac{(-1)^{k+1}}{k}\cos(kx)
\qquad(x\in[-\pi,\pi]\text{ a.e.}).
\]
代入 \(x=p\theta\) 并对有限个 \(p\) 求和（可交换求和次序）得到
\[
\mathcal A_r(\theta)
=\sum_{p}r(p)\sum_{k\ge 1}\frac{(-1)^{k+1}}{k}\cos(kp\theta).
\]
因此 \(\mathcal A_r\) 的第 \(n\) 个余弦系数为
\[
\frac1\pi\int_{-\pi}^{\pi}\mathcal A_r(\theta)\cos(n\theta)\,\dd\theta
=\sum_{p\mid n} r(p)\,\frac{(-1)^{n/p+1}}{n/p}.
\]
当 \(n=p\) 为素数时右端仅含 \(p\) 一项且 \(n/p=1\)，从而该余弦系数等于 \(r(p)\)，即得所述反演公式。证毕。
\end{proof}

\begin{corollary}[整数刚性阈值：\(L^2\) 误差控制下的逐素数精确恢复]\label{cor:xi-leyang-logamplitude-rounding-threshold}
在定理 \ref{thm:xi-leyang-logamplitude-prime-frequency-inversion} 的设定下，令 \(\widetilde{\mathcal A}\in L^2([-\pi,\pi])\) 为任意观测/学习得到的对数振幅近似，并对素数 \(p\) 定义
\[
\widetilde r(p):=\operatorname{Round}\!\left(\frac1\pi\int_{-\pi}^{\pi}\widetilde{\mathcal A}(\theta)\cos(p\theta)\,\dd\theta\right)\in\ZZ,
\]
其中 \(\operatorname{Round}\) 表示四舍五入到最近整数。
若满足误差阈值
\[
\norm{\widetilde{\mathcal A}-\mathcal A_r}_{L^2([-\pi,\pi])}<\frac{\sqrt\pi}{2},
\]
则对所有素数 \(p\) 均有 \(\widetilde r(p)=r(p)\)。
\end{corollary}
\begin{proof}
记
\[
c_p(F):=\frac1\pi\int_{-\pi}^{\pi}F(\theta)\cos(p\theta)\,\dd\theta.
\]
由 Cauchy--Schwarz 不等式并注意 \(\norm{\cos(p\theta)}_{L^2([-\pi,\pi])}=\sqrt\pi\)，得到
\[
\abs{c_p(\widetilde{\mathcal A})-c_p(\mathcal A_r)}
\le \frac1\pi\norm{\widetilde{\mathcal A}-\mathcal A_r}_{L^2}\,\norm{\cos(p\theta)}_{L^2}
=\frac1{\sqrt\pi}\norm{\widetilde{\mathcal A}-\mathcal A_r}_{L^2}
<\frac12.
\]
而由定理 \ref{thm:xi-leyang-logamplitude-prime-frequency-inversion}，\(c_p(\mathcal A_r)=r(p)\in\ZZ\)。
因此 \(c_p(\widetilde{\mathcal A})\) 与整数 \(r(p)\) 的距离严格小于 \(1/2\)，四舍五入必回到 \(r(p)\)。证毕。
\end{proof}

\begin{definition}[素数根均匀测度与零点经验测度]\label{def:xi-leyang-prime-root-measures}
对每个素数 \(p\)，令 \(\eta_p\) 为方程 \(z^p=-1\) 的 \(p\) 个解上的均匀概率测度：
\[
\eta_p:=\frac1p\sum_{k=0}^{p-1}\delta_{\exp\!\left(i\frac{(2k+1)\pi}{p}\right)}
\quad\text{（概率测度，支撑于 \(S^1\)）}.
\]
对任意 \(r\in\NN^{(\mathcal P)}\setminus\{0\}\)，记
\[
D(r):=\deg \mathcal L(r)=\sum_{p\in\mathcal P}r(p)\,p,
\]
并定义其零点归一化经验测度（按重数计）
\[
\mu_r:=\frac1{D(r)}\sum_{\zeta:\ \mathcal L(r)(\zeta)=0}\mathrm{mult}(\zeta)\,\delta_\zeta
\quad\text{（概率测度，支撑于 \(S^1\)）}.
\]
\end{definition}

\begin{lemma}[Fourier 系数的显式消失律]\label{lem:xi-eta-p-fourier}
\leanverified{paper\_xi\_eta\_p\_fourier}
对任意整数 \(n\neq 0\)，有
\[
\widehat{\eta_p}(n):=\int_{S^1}z^n\,\dd\eta_p(z)=
\begin{cases}
(-1)^{n/p},&p\mid n,\\
0,&p\nmid n.
\end{cases}
\]
\end{lemma}
\begin{proof}
\[
\int_{S^1}z^n\,\dd\eta_p(z)
=\frac1p\sum_{k=0}^{p-1}\exp\!\left(i n\frac{(2k+1)\pi}{p}\right)
=e^{i n\pi/p}\cdot \frac1p\sum_{k=0}^{p-1}e^{i2\pi nk/p}.
\]
若 \(p\nmid n\)，则内层为非平凡 \(p\) 次单位根之和，取值 \(0\)；若 \(p\mid n\)，则内层为 \(1\)，从而整体为 \(e^{i\pi(n/p)}=(-1)^{n/p}\)。
\end{proof}

\begin{proposition}[零点测度的素数分解与单频谱解码]\label{prop:xi-leyang-prime-measure-linear-independence-decoding}
对任意 \(r\in\NN^{(\mathcal P)}\setminus\{0\}\)，有凸组合分解
\[
\mu_r=\sum_{p\in\mathcal P}\frac{r(p)\,p}{D(r)}\,\eta_p.
\]
进一步，对任意有限支撑整数族 \((c_p)_{p\in\mathcal P}\subset\ZZ\) 与带符号测度
\(
\sigma:=\sum_pc_p\,\eta_p
\)，若 \(\sigma=0\) 则必有 \(c_p=0\)（对所有 \(p\)）。
并且对任意素数 \(p\)，有解码公式
\[
c_p=-\int_{S^1}z^p\,\dd\sigma(z).
\]
\end{proposition}
\begin{proof}
由定义 \(\mathcal L(r)=\prod_p(1+z^p)^{r(p)}\)，每个因子 \(1+z^p\) 贡献 \(p\) 个单位根零点，且零点重数按指数 \(r(p)\) 线性叠加，因此零点测度满足
\(
\mu_r=\sum_p\frac{r(p)p}{D(r)}\eta_p
\)。
\par
对线性无关性与解码，取 Fourier 矩 \(n=p\)。由引理 \ref{lem:xi-eta-p-fourier}，对 \(q\neq p\) 因 \(q\nmid p\) 有
\(
\int z^p\,\dd\eta_q=0
\)，而对 \(q=p\) 有
\(
\int z^p\,\dd\eta_p=-1
\)。
故
\[
\int_{S^1}z^p\,\dd\sigma(z)=\sum_q c_q\int z^p\,\dd\eta_q(z)=-c_p,
\]
从而若 \(\sigma=0\) 则 \(c_p=0\)（对任意 \(p\)），并得解码公式。
\end{proof}

\begin{theorem}[素数截断零点谱的 Haar 极限与显式差分界]\label{thm:xi-leyang-prime-truncation-haar-discrepancy}
令
\[
\mathcal L_{\le P}(z):=\prod_{p\le P}(1+z^p),
\qquad
D_P:=\deg\mathcal L_{\le P}=\sum_{p\le P}p,
\]
并令 \(\mu_P\) 为 \(\mathcal L_{\le P}\) 的零点归一化经验测度（按重数计）。
则当 \(P\to\infty\) 时，\(\mu_P\) 在弱-\(\ast\) 意义下收敛到单位圆上的 Haar 概率测度 \(m_{S^1}\)。
此外，存在绝对常数 \(C>0\)，使对任意弧段 \(A\subset S^1\) 有
\[
\sup_{A\subset S^1}\abs{\mu_P(A)-m_{S^1}(A)}
\le \frac{C}{\sqrt{D_P}}
=\frac{C}{\sqrt{\sum_{p\le P}p}}.
\]
\end{theorem}
\begin{proof}
由命题 \ref{prop:xi-leyang-prime-measure-linear-independence-decoding} 的分解（取 \(r(p)=\ind_{\{p\le P\}}\)），有
\[
\mu_P=\sum_{p\le P}\frac{p}{D_P}\,\eta_p.
\]
对任意整数 \(n\neq 0\)，其 Fourier 系数为
\[
\widehat{\mu_P}(n)
=\sum_{p\le P}\frac{p}{D_P}\widehat{\eta_p}(n)
=\frac1{D_P}\sum_{\substack{p\le P\\ p\mid n}}p\,(-1)^{n/p},
\]
其中最后一步用引理 \ref{lem:xi-eta-p-fourier}。
对固定 \(n\neq 0\)，分子仅为有限和且与 \(P\) 无关，而 \(D_P\to\infty\)，故 \(\widehat{\mu_P}(n)\to 0\)。
Haar 测度 \(m_{S^1}\) 的所有非零 Fourier 系数均为 \(0\)，因而 \(\mu_P\Rightarrow m_{S^1}\)。
\par
差分界由 Erd\H{o}s--Tur\'an 不等式给出（参见 \cite{KuipersNiederreiter1974}）：存在绝对常数 \(C_0>0\)，对任意 \(H\in\NN\) 有
\[
\sup_A\abs{\mu_P(A)-m_{S^1}(A)}
\le C_0\left(\frac1H+\sum_{k=1}^H\frac{|\widehat{\mu_P}(k)|}{k}\right).
\]
由上式表达与 \(\sum_{p\mid k}p\le k\) 得
\(
|\widehat{\mu_P}(k)|\le k/D_P
\)，从而
\[
\sum_{k=1}^H\frac{|\widehat{\mu_P}(k)|}{k}\le \sum_{k=1}^H\frac1{D_P}=\frac{H}{D_P}.
\]
取 \(H=\lfloor\sqrt{D_P}\rfloor\) 得
\(
\sup_A\abs{\mu_P(A)-m_{S^1}(A)}\le 2C_0/\sqrt{D_P}
\)，并将常数并入 \(C\) 即得。
\end{proof}

\begin{proposition}[Gödel 椭圆边界的推前测度与显式密度]\label{prop:xi-godel-ellipse-pushforward-density}
取 \(a,b>0\)，并定义单位圆到轴对齐椭圆边界的参数化
\[
T_{a,b}:S^1\to \partial E_{a,b},
\qquad
T_{a,b}(e^{it})=(a\cos t,\ b\sin t),
\]
其中 \(\partial E_{a,b}\) 为 \(x^2/a^2+y^2/b^2=1\) 的边界。
令
\(
\nu_{a,b}:=(T_{a,b})_*m_{S^1}
\)
为 Haar 测度的推前测度，则 \(\nu_{a,b}\) 关于椭圆弧长测度 \(ds\) 绝对连续，且 Radon--Nikodym 导数满足
\[
\frac{\dd\nu_{a,b}}{\dd s}\Bigl(T_{a,b}(e^{it})\Bigr)
=\frac1{2\pi}\cdot\frac1{\sqrt{a^2\sin^2 t+b^2\cos^2 t}}.
\]
进一步，若 \(\mu_P\) 为定理 \ref{thm:xi-leyang-prime-truncation-haar-discrepancy} 中的零点谱测度，则
\[
(T_{a,b})_*\mu_P\Rightarrow \nu_{a,b}.
\]
\end{proposition}
\begin{proof}
由弧长元素
\(
ds=\sqrt{(a\sin t)^2+(b\cos t)^2}\,\dd t
\)
可知 \(m_{S^1}\) 在参数 \(t\) 下的密度为 \(\dd t/(2\pi)\)。推前到 \(\partial E_{a,b}\) 后得到关于 \(ds\) 的密度为
\[
\frac{\dd t}{2\pi}\cdot \frac1{\dd s}
=\frac1{2\pi}\cdot\frac1{\sqrt{a^2\sin^2 t+b^2\cos^2 t}}.
\]
弱收敛在连续映射下的推前稳定性给出第二条：由 \(\mu_P\Rightarrow m_{S^1}\) 且 \(T_{a,b}\) 连续，得 \((T_{a,b})_*\mu_P\Rightarrow (T_{a,b})_*m_{S^1}=\nu_{a,b}\)。
\end{proof}

\begin{theorem}[素数寄存器序列的 Haar 光谱收敛判别]\label{thm:xi-leyang-prime-register-haar-criterion}
设 \((r_m)_{m\ge 1}\subset \NN^{(\mathcal P)}\setminus\{0\}\) 为任意序列，并令
\[
\mathcal L_m(z):=\mathcal L(r_m)(z),
\qquad
D_m:=D(r_m)=\sum_{p\in\mathcal P}r_m(p)\,p,
\qquad
\mu_m:=\mu_{r_m}.
\]
则以下三条等价：
\begin{enumerate}
  \item \(\mu_m\Rightarrow m_{S^1}\)。
  \item 对每个固定素数 \(q\)，有
  \[
  \frac{r_m(q)\,q}{D_m}\longrightarrow 0.
  \]
  \item 对每个固定整数 \(n\neq 0\)，有
  \[
  \frac{\sum_{p\mid n}r_m(p)\,p}{D_m}\longrightarrow 0.
  \]
\end{enumerate}
\end{theorem}
\begin{proof}
由定义 \(\mu_m=\sum_p \frac{r_m(p)p}{D_m}\eta_p\) 与引理 \ref{lem:xi-eta-p-fourier}，对任意整数 \(n\neq 0\) 有
\[
\widehat{\mu_m}(n)
=\int z^n\,\dd\mu_m(z)
=\frac1{D_m}\sum_{p\mid n}r_m(p)\,p\,(-1)^{n/p}.
\]
Haar 测度满足 \(\widehat{m_{S^1}}(n)=0\)（\(n\neq 0\)），故 \(\mu_m\Rightarrow m_{S^1}\) 等价于对一切 \(n\neq 0\) 有 \(\widehat{\mu_m}(n)\to 0\)；
而这与第三条等价（符号项不影响绝对值阶）。
第三条显然推出第二条（取 \(n=q\)）。
第二条推出第三条是因为 \(n\) 的素因子集合有限，\(\sum_{p\mid n}r_m(p)p/D_m\) 为有限项之和，逐项趋于 \(0\) 即整体趋于 \(0\)。
\end{proof}

\begin{theorem}[零点普适律与对数振幅保持律的严格分裂]\label{thm:xi-leyang-bulk-boundary-splitting-haar-vs-logamplitude}
设 \((r_N)_{N\ge 1}\subset \NN^{(\mathcal P)}\setminus\{0\}\)，并令 \(\mu_{r_N}\) 与
\[
D_N:=D(r_N)=\sum_{p\in\mathcal P} r_N(p)\,p
\]
如定义 \ref{def:xi-leyang-prime-root-measures} 所示。
假设 \(D_N\to\infty\)，且对每个固定素数 \(p\)，序列 \((r_N(p))_{N}\) 有界。
则：
\begin{enumerate}
  \item 零点经验测度满足 \(\mu_{r_N}\Rightarrow m_{S^1}\)（单位圆上的 Haar 概率测度）；
  \item 与此同时，若记 \(\mathcal A_{r_N}(\theta)=\log|\mathcal L(r_N)(e^{i\theta})|\)，则对每个素数 \(p\) 恒有
  \[
  \frac1\pi\int_{-\pi}^{\pi}\mathcal A_{r_N}(\theta)\cos(p\theta)\,\dd\theta=r_N(p),
  \]
  因而 \((\mathcal A_{r_N})_N\) 不可能在 \(L^2([-\pi,\pi])\) 中收敛到某个与 \(r_N\) 无关的“普适极限”（除非 \(r_N(p)\) 在每个素数处逐点稳定）。
\end{enumerate}
\end{theorem}
\begin{proof}
对任意固定素数 \(q\)，由有界性得 \(r_N(q)\,q=O_q(1)\)，而 \(D_N\to\infty\)，故
\(
\frac{r_N(q)\,q}{D_N}\to 0
\)。
由定理 \ref{thm:xi-leyang-prime-register-haar-criterion} 即得 \(\mu_{r_N}\Rightarrow m_{S^1}\)，从而得到 (1)。
\par
(2) 由定理 \ref{thm:xi-leyang-logamplitude-prime-frequency-inversion} 对每个 \(N\) 与每个素数 \(p\) 的反演恒等式直接给出。
若 \(\mathcal A_{r_N}\to \mathcal A\) 于 \(L^2\)，则其每个余弦系数必收敛，因而 \(r_N(p)\) 必对每个 \(p\) 收敛；若要极限不依赖于序列 \((r_N)\)，则只能在逐素数意义下稳定。证毕。
\end{proof}

\begin{conjecture}[外置诱导素数账本的质量不凝聚与普适差分阶]\label{conj:xi-leyang-prime-register-universal-nonconcentration}
在本文折叠/投影外置范式中，由算法轨迹自然诱导的素数寄存器序列 \((r_m)\) 满足定理 \ref{thm:xi-leyang-prime-register-haar-criterion} 的质量不凝聚条件，从而其 Lee--Yang--Gödel 多项式零点谱 \(\mu_{r_m}\) 弱收敛到 \(m_{S^1}\)。
并且在适当的规范化与均匀性假设下，弧段差分满足 \(\sup_A|\mu_{r_m}(A)-m_{S^1}(A)|=O(D(r_m)^{-1/2})\)，且该指数阶在该范式内具有普适性。
\end{conjecture}

\paragraph{Zeckendorf--Gödel 素数码集合的密度、Lee--Yang 实根结构与谱行列式重整}
设 \((p_k)_{k\ge 1}\) 为递增素数序列，并定义
\[
\mathrm{ZG}
:=\Bigl\{
n\in\NN:\ \forall p,\ v_p(n)\in\{0,1\},\ \forall k\ge 1,\ v_{p_k}(n)\,v_{p_{k+1}}(n)=0
\Bigr\}.
\]
即 \(n\) 为平方因子自由，且其素因子索引集合不含相邻整数。

\begin{theorem}[Gödel--Zeckendorf 素数约束 Dirichlet 级数的对数重整化与半临界解析性]\label{thm:xi-zg-dirichlet-log-renormalization}
\leanverified{paper\_xi\_zg\_dirichlet\_log\_renormalization}
令
\[
\Omega:=\Bigl\{\varepsilon=(\varepsilon_k)_{k\ge 1}:\ \varepsilon_k\in\{0,1\},\ \varepsilon_k\varepsilon_{k+1}=0,\ \varepsilon\ \text{仅有限支撑}\Bigr\}.
\]
对 \(\Re(s)>1\) 定义 Dirichlet 级数
\[
F(s):=\sum_{n\in \mathrm{ZG}}n^{-s}
=\sum_{\varepsilon\in\Omega}\ \prod_{k\ge 1}p_k^{-s\varepsilon_k}.
\]
并对 \(N\ge 0\) 定义截断
\[
F_N(s):=\sum_{\substack{\varepsilon\in\{0,1\}^N\\ \varepsilon_k\varepsilon_{k+1}=0}}\ \prod_{k=1}^{N}p_k^{-s\varepsilon_k},
\qquad
F_0(s)=1,\ \ F_1(s)=1+2^{-s}.
\]
令 \(r_N(s):=F_N(s)/F_{N-1}(s)\)（\(N\ge 1\)），并令
\[
\Delta_N(s):=\log r_N(s)-p_N^{-s}\qquad(N\ge 2),
\]
其中对数取为与实轴上 \(s>1\) 的正值一致的解析分支。
则：
\begin{enumerate}
  \item 对所有 \(\Re(s)>0\) 成立递推恒等式
  \[
  F_N(s)=F_{N-1}(s)+p_N^{-s}F_{N-2}(s),\qquad
  r_N(s)=1+\frac{p_N^{-s}}{r_{N-1}(s)}.
  \]
  \item 级数
  \[
  G(s):=\sum_{N\ge 2}\Delta_N(s)
  \]
  在每个半平面 \(\Re(s)\ge \sigma_0>\tfrac12\) 上绝对且一致收敛，从而在 \(\Re(s)>\tfrac12\) 上给出全纯函数。
  \item 令素数 \(\zeta\) 函数（prime zeta）
  \(
  P(s):=\sum_{k\ge 1}p_k^{-s}
  \)
  （在 \(\Re(s)>1\) 绝对收敛）。
  则对 \(\Re(s)>1\) 有严格分解
  \[
  \log F(s)=P(s)+G(s),
  \qquad\text{从而}\qquad
  F(s)=\exp\!\bigl(P(s)+G(s)\bigr).
  \]
\end{enumerate}
\end{theorem}

\begin{proof}
第一条由末位 \(\varepsilon_N\in\{0,1\}\) 分拆：\(\varepsilon_N=0\) 贡献为 \(F_{N-1}\)，而 \(\varepsilon_N=1\) 强制 \(\varepsilon_{N-1}=0\) 并贡献 \(p_N^{-s}F_{N-2}\)。
对 \(r_N=F_N/F_{N-1}\) 取比即得递推。
\par
对第二条，令 \(\sigma=\Re(s)\ge \sigma_0>\tfrac12\)。
由递推 \(r_N=1+p_N^{-s}/r_{N-1}\) 与 \(|p_N^{-s}|\le p_N^{-\sigma}\) 可在每个紧子域上固定 \(r_{N-1}(s)\) 远离 \(0\) 的一致下界，
从而存在常数 \(C(\sigma_0)\) 使
\[
|\Delta_N(s)|
\le
C(\sigma_0)\Bigl(p_N^{-2\sigma}+(p_Np_{N-1})^{-\sigma}\Bigr)
\qquad(\sigma\ge\sigma_0).
\]
由于
\(
\sum_N p_N^{-2\sigma}
\)
与
\(
\sum_N(p_Np_{N-1})^{-\sigma}
\)
在 \(\sigma>\tfrac12\) 收敛，故 \(\sum_{N\ge 2}\Delta_N(s)\) 在 \(\Re(s)\ge\sigma_0\) 上绝对一致收敛，从而 \(G\) 在 \(\Re(s)>\tfrac12\) 全纯。
\par
第三条由
\(
\log F_N(s)=\sum_{k=1}^N\log r_k(s)
\)
与 \(\Delta_k=\log r_k-p_k^{-s}\) 得
\[
\log F_N(s)=\sum_{k=1}^{N}p_k^{-s}+\sum_{k=2}^{N}\Delta_k(s).
\]
令 \(N\to\infty\) 并在 \(\Re(s)>1\) 用绝对收敛交换极限，得到 \(\log F(s)=P(s)+G(s)\)。证毕。
\end{proof}

\begin{theorem}[平方自由 Euler 因子分离与 \texorpdfstring{$\zeta$}{zeta}-因子化]\label{thm:xi-zg-zeta-factorization}
\leanverified{paper\_xi\_zg\_zeta\_factorization}
沿用定理 \ref{thm:xi-zg-dirichlet-log-renormalization} 的记号，并令
\[
B_N(s):=\prod_{k=1}^{N}\bigl(1+p_k^{-s}\bigr),\qquad
H_N(s):=\frac{F_N(s)}{B_N(s)}.
\]
则存在唯一函数 \(H\) 满足：
\begin{enumerate}
  \item \(H_N(s)\to H(s)\) 在每个紧集 \(K\subset\{s:\Re(s)>\tfrac12\}\) 上一致收敛；
  \item \(H(s)\) 在半平面 \(\Re(s)>\tfrac12\) 上全纯且处处非零；
  \item 在 \(\Re(s)>1\) 上成立严格因子化
  \[
  F(s)=H(s)\frac{\zeta(s)}{\zeta(2s)}.
  \]
\end{enumerate}
从而 \(F\) 在 \(\Re(s)>\tfrac12\) 上亚纯，且在该半平面内唯一可能的极点位于 \(s=1\)，并且为一阶极点。
\end{theorem}

\begin{proof}
由递推可写 \(F_N(s)=\prod_{k=1}^{N}r_k(s)\)，从而
\[
\log H_N(s)
=\sum_{k=1}^{N}\Bigl(\log r_k(s)-\log\bigl(1+p_k^{-s}\bigr)\Bigr).
\]
对 \(k\ge 2\) 使用 \(\Delta_k(s)=\log r_k(s)-p_k^{-s}\)，并将上式改写为
\[
\log H_N(s)
=\Bigl(\log r_1(s)-\log(1+p_1^{-s})\Bigr)
\sum_{k=2}^{N}\Delta_k(s)
\sum_{k=2}^{N}\Bigl(p_k^{-s}-\log(1+p_k^{-s})\Bigr).
\]
第一项恒为 \(0\)。
由定理 \ref{thm:xi-zg-dirichlet-log-renormalization}，\(\sum_{k\ge 2}\Delta_k(s)\) 在 \(\Re(s)>\tfrac12\) 的每个紧集上一致收敛。
另一方面，当 \(\Re(s)\ge \sigma_0>\tfrac12\) 时有
\(
\bigl|p_k^{-s}-\log(1+p_k^{-s})\bigr|
\ll p_k^{-2\sigma_0}
\),
从而 \(\sum_{k\ge 2}\bigl(p_k^{-s}-\log(1+p_k^{-s})\bigr)\) 在同一紧集上一致绝对收敛。
故 \(\log H_N\) 在 \(\Re(s)>\tfrac12\) 的紧集上一致收敛到某个全纯函数 \(\log H\)，并且 \(H=\exp(\log H)\) 全纯且不取零值。

对 \(\Re(s)>1\)，由平方自由 Dirichlet 级数的标准恒等式
\[
\prod_{p}\bigl(1+p^{-s}\bigr)
=\sum_{\substack{m\ge 1\\ m\ \text{平方因子自由}}}\frac1{m^s}
=\frac{\zeta(s)}{\zeta(2s)}
\]
（参见 \cite{Apostol1976,Titchmarsh1986RiemannZeta}），得 \(B_N(s)\to \zeta(s)/\zeta(2s)\)。
再由 \(F_N(s)=H_N(s)B_N(s)\) 并令 \(N\to\infty\)（在 \(\Re(s)>1\) 可由绝对收敛交换极限），得到 \(F(s)=H(s)\zeta(s)/\zeta(2s)\)。
最后，\(\zeta(2s)\) 在 \(\Re(s)>\tfrac12\) 上解析且无零点，而 \(\zeta(s)\) 在该半平面内唯一奇性为 \(s=1\) 的一阶极点，因而 \(F\) 在 \(\Re(s)>\tfrac12\) 上仅可能在 \(s=1\) 处有一阶极点。证毕。
\end{proof}

\input{subsubsec__xi-time-protocol-conclusions-part15d-prime-register-zg-density-leyang-spectrum-continued-euler}

\begin{corollary}[素数路径硬核常数与 \texorpdfstring{$s=1$}{s=1} 留数]\label{cor:xi-zg-hardcore-constant-residue}
在定理 \ref{thm:xi-zg-zeta-factorization} 的记号下，记
\[
\mathfrak C_{\mathrm{HC}}:=H(1)=\lim_{N\to\infty}\frac{F_N(1)}{\prod_{k=1}^{N}\bigl(1+\frac1{p_k}\bigr)}\in(0,1).
\]
则 \(F(s)\) 在 \(s=1\) 的留数满足
\[
\operatorname*{Res}_{s=1}F(s)=\frac{\mathfrak C_{\mathrm{HC}}}{\zeta(2)}=\frac{6}{\pi^2}\,\mathfrak C_{\mathrm{HC}}.
\]
并且该留数与 \(\mathrm{ZG}\) 的 Dirichlet 密度一致；结合定理 \ref{con:xi-zg-density-closed-form-tail-bound} 的自然密度存在性，得到
\[
\delta_{\mathrm{ZG}}=\frac{\mathfrak C_{\mathrm{HC}}}{\zeta(2)}.
\]
数值核验可得到
\[
\mathfrak C_{\mathrm{HC}}=0.8461958103\ldots,
\qquad
\delta_{\mathrm{ZG}}=0.5144253666\ldots
\]
\leanverified{paper\_xi\_zg\_hardcore\_constant\_residue}
\end{corollary}

\begin{proof}
由定理 \ref{thm:xi-zg-zeta-factorization}，在 \(s=1\) 的邻域内
\(
F(s)=H(s)\zeta(s)/\zeta(2s)
\)
且 \(H\) 在 \(s=1\) 解析。
因此
\[
\operatorname*{Res}_{s=1}F(s)
=H(1)\operatorname*{Res}_{s=1}\frac{\zeta(s)}{\zeta(2s)}
=\frac{H(1)}{\zeta(2)}.
\]
留数等于 \(\mathrm{ZG}\) 的 Dirichlet 密度是推论 \ref{cor:xi-zg-dirichlet-residue} 的内容，而该推论亦指出留数与自然密度一致，从而得 \(\delta_{\mathrm{ZG}}=\mathfrak C_{\mathrm{HC}}/\zeta(2)\)。
最后，由于对实数 \(s>1\) 有 \(F_N(s)\le B_N(s)\) 且约束非平凡，故 \(0<H(s)<1\) 并推出 \(\mathfrak C_{\mathrm{HC}}\in(0,1)\)。证毕。
\end{proof}

\begin{proposition}[\texorpdfstring{$s=1$}{s=1} 的 Laurent 展开与 Euler--Kronecker 型常数]\label{prop:xi-zg-hardcore-laurent}
记
\[
\kappa_{\mathrm{HC}}:=\left.\frac{d}{ds}\log H(s)\right|_{s=1}.
\]
则 \(\kappa_{\mathrm{HC}}\) 存在，并且 \(F\) 在 \(s=1\) 的 Laurent 展开为
\[
F(s)=\frac{\mathfrak C_{\mathrm{HC}}}{\zeta(2)}\cdot\frac{1}{s-1}
\;+\;
\frac{\mathfrak C_{\mathrm{HC}}}{\zeta(2)}
\left(\gamma-\frac{2\zeta'(2)}{\zeta(2)}+\kappa_{\mathrm{HC}}\right)
 + O(s-1),
\qquad s\to 1,
\]
其中 \(\gamma\) 为 Euler--Mascheroni 常数。
数值核验表明 \(\kappa_{\mathrm{HC}}=0.40404\ldots\)。
\end{proposition}

\begin{proof}
由定理 \ref{thm:xi-zg-zeta-factorization} 的构造，\(\log H_N(s)\) 在 \(\Re(s)>\tfrac12\) 的紧集上一致收敛为 \(\log H(s)\)，并且其尾项由 \(\sum_{k\ge 2}\Delta_k(s)\) 与 \(\sum_{k\ge 2}(p_k^{-s}-\log(1+p_k^{-s}))\) 的一致收敛控制。
同样地，对 \(s\) 位于 \(s=1\) 的小邻域且满足 \(\Re(s)>\tfrac12\)，上述两级数可逐项求导，从而 \(\kappa_{\mathrm{HC}}\) 存在。

另一方面，
\(
\zeta(s)=\frac1{s-1}+\gamma+O(s-1)
\)
且
\[
\frac{1}{\zeta(2s)}
\;=\;
\frac{1}{\zeta(2)}
\left(1-\frac{2\zeta'(2)}{\zeta(2)}(s-1)+O((s-1)^2)\right).
\]
再由
\(
H(s)=H(1)\bigl(1+\kappa_{\mathrm{HC}}(s-1)+O((s-1)^2)\bigr)
\)
三者相乘即得所述展开式。证毕。
\end{proof}

\input{thm__xi-zg-counting-power-saving-error}

\begin{theorem}[Zeckendorf--Gödel 素数码集合的自然密度、闭式可计算性与尾项误差界]\label{con:xi-zg-density-closed-form-tail-bound}
对每个 \(N\ge 1\)，定义截断集合
\[
\mathrm{ZG}^{(N)}
:=\Bigl\{
n\in\NN:\ \forall 1\le k\le N,\ v_{p_k}(n)\in\{0,1\},\
\forall 1\le k\le N-1,\ v_{p_k}(n)\,v_{p_{k+1}}(n)=0
\Bigr\}.
\]
记
\[
\delta_N:=\mathrm{dens}\bigl(\mathrm{ZG}^{(N)}\bigr).
\]
则自然密度
\[
\delta_{\mathrm{ZG}}:=\mathrm{dens}(\mathrm{ZG})
\]
存在，且
\[
\delta_{\mathrm{ZG}}=\lim_{N\to\infty}\delta_N\in(0,1).
\]
进一步，\(\delta_N\) 具有显式闭式
\[
\delta_N=\Bigl(\prod_{k=1}^N\Bigl(1-\frac1{p_k}\Bigr)\Bigr)I_N,
\qquad
I_N:=\sum_{\substack{S\subseteq\{1,\dots,N\}\\ S\ \mathrm{无相邻}}}\ \prod_{k\in S}\frac1{p_k},
\]
并满足尾项误差上界
\[
0\le \delta_N-\delta_{\mathrm{ZG}}
\le
\sum_{k>N}\frac1{p_k^2}
\;+\;
\sum_{k\ge N}\frac1{p_kp_{k+1}}.
\]
特别地，\(\delta_{\mathrm{ZG}}\) 可通过生成素数并迭代计算 \(\delta_N\) 在任意精度下有效逼近。
\end{theorem}
\begin{proof}
对固定 \(N\)，成员资格仅依赖于模
\(
M_N:=\prod_{k=1}^N p_k^2
\)
的剩余类，故 \(\delta_N\) 存在。

记
\(
a_k:=1-\frac1{p_k},\ b_k:=a_k\frac1{p_k}
\)。
在前 \(N\) 个素数坐标上，“允许指数为 \(1\)”的索引集合恰为路径图 \(\{1,\dots,N\}\) 的独立集 \(S\)；
对应权重为
\(
\prod_{k\in S}b_k\prod_{k\notin S}a_k
\)。
求和得
\[
\delta_N
=\sum_{S\ \mathrm{无相邻}}\ \prod_{k\in S}b_k\prod_{k\notin S}a_k
=\Bigl(\prod_{k=1}^N a_k\Bigr)\sum_{S\ \mathrm{无相邻}}\prod_{k\in S}\frac1{p_k},
\]
即所示闭式。

又 \(\mathrm{ZG}=\bigcap_{N\ge 1}\mathrm{ZG}^{(N)}\)，故 \(\delta_N\) 单调下降且
\[
\overline{\mathrm{dens}}(\mathrm{ZG})\le \delta_N\quad(\forall N).
\]
若 \(n\in \mathrm{ZG}^{(N)}\setminus \mathrm{ZG}\)，则必发生尾部违约事件之一：
\[
\exists\,k>N,\ p_k^2\mid n
\qquad\text{或}\qquad
\exists\,k\ge N,\ p_kp_{k+1}\mid n.
\]
故
\[
\mathrm{ZG}^{(N)}\setminus \mathrm{ZG}
\subseteq
\Bigl(\bigcup_{k>N}\{p_k^2\mid n\}\Bigr)
\cup
\Bigl(\bigcup_{k\ge N}\{p_kp_{k+1}\mid n\}\Bigr).
\]
由并集上界与
\(
\mathrm{dens}\{m\mid n\}=1/m
\)
得
\[
\mathrm{dens}\bigl(\mathrm{ZG}^{(N)}\setminus \mathrm{ZG}\bigr)
\le
\sum_{k>N}\frac1{p_k^2}
+
\sum_{k\ge N}\frac1{p_kp_{k+1}},
\]
即误差估计。令 \(N\to\infty\) 得上下密度同限，故自然密度存在且等于 \(\lim\delta_N\)。

最后，
\[
\delta_{\mathrm{ZG}}
\ge 1-\sum_{k\ge 1}\frac1{p_k^2}-\sum_{k\ge 1}\frac1{p_kp_{k+1}}>0,
\]
因为两级数收敛且其和严格小于 \(1\)。证毕。
\end{proof}

\begin{corollary}[Mertens 重整常数与素数路径独立集分区函数的对数增长律]\label{cor:xi-zg-mertens-log-growth}
沿用定理 \ref{con:xi-zg-density-closed-form-tail-bound} 的记号。
则
\[
\delta_{\mathrm{ZG}}
=\lim_{N\to\infty}
\Bigl(\prod_{k=1}^N\Bigl(1-\frac1{p_k}\Bigr)\Bigr)I_N,
\]
并且
\[
I_N\sim \delta_{\mathrm{ZG}}\,e^{\gamma}\,\log p_N
\qquad (N\to\infty),
\]
其中 \(\gamma\) 为 Euler--Mascheroni 常数。
\end{corollary}
\begin{proof}
由定理 \ref{con:xi-zg-density-closed-form-tail-bound}，
\[
I_N=\frac{\delta_N}{\prod_{k\le N}(1-\frac1{p_k})}.
\]
再用 Mertens 乘积渐近
\(
\prod_{p\le y}(1-\frac1p)\sim e^{-\gamma}/\log y
\)
并令 \(y=p_N\)、\(\delta_N\to\delta_{\mathrm{ZG}}\) 即得。证毕。
\end{proof}

\begin{theorem}[素数权重 Fibonacci 路径的 Lee--Yang 实根性与 Poisson--binomial 分解]\label{con:xi-zg-weighted-path-leyang-poisson-binomial}
定义
\[
P_N(t):=\sum_{\substack{S\subseteq\{1,\dots,N\}\\ S\ \mathrm{无相邻}}}
t^{|S|}\prod_{k\in S}\frac1{p_k},
\]
则
\[
P_0(t)=1,\qquad
P_1(t)=1+\frac{t}{p_1},\qquad
P_N(t)=P_{N-1}(t)+\frac{t}{p_N}P_{N-2}(t)\quad(N\ge 2).
\]
并且成立：
\begin{enumerate}
  \item \(P_N\) 的全部零点均为互异负实根；
  \item \(P_N\) 与 \(P_{N-1}\) 的零点严格交错；
  \item 设 \(m_N:=\deg P_N=\lfloor(N+1)/2\rfloor\)，则存在 \(\alpha_{N,j}>0\) 使
  \[
  P_N(t)=\prod_{j=1}^{m_N}\Bigl(1+\frac{t}{\alpha_{N,j}}\Bigr).
  \]
\end{enumerate}
对固定 \(t>0\)，令
\[
\pi_{N,j}(t):=\frac{t}{t+\alpha_{N,j}}\in(0,1).
\]
若按 Gibbs 权重
\(
\propto t^{|S|}\prod_{k\in S}p_k^{-1}
\)
在独立集上取样，则随机变量 \(|S|\) 满足分布分解
\[
|S|\overset{d}{=}\sum_{j=1}^{m_N}B_{N,j},
\qquad
B_{N,j}\sim\mathrm{Bernoulli}\!\bigl(\pi_{N,j}(t)\bigr),
\]
其中 \((B_{N,j})\) 相互独立。
\end{theorem}
\begin{proof}
递推直接由索引 \(N\) 是否入独立集的分拆得到。

令 \(\mathcal T_{N+1}\) 为长度 \(N+1\) 的路径树，赋边权
\(
w_k:=p_k^{-1/2}
\)
（\(1\le k\le N\)），并记 \(A_{N+1}\) 为其加权邻接矩阵。
匹配多项式与独立集母函数的标准对应给出
\[
\chi_{N+1}(x):=\det(xI-A_{N+1})
=x^{N+1}P_N\!\Bigl(-\frac1{x^2}\Bigr).
\]
由于 \(A_{N+1}\) 为实对称 Jacobi 矩阵（副对角元严格正），其特征值实且互异；
故 \(P_N\) 的零点为 \(-1/\lambda_j^2\)，均为互异负实数。
交错性由主子矩阵谱交错（Cauchy interlacing）传递到 \(P_N,P_{N-1}\) 的零点。

第三条由实负根分解直接成立。
设
\[
G_{N,t}(s):=\frac{P_N(ts)}{P_N(t)}
=\prod_{j=1}^{m_N}
\frac{1+\frac{ts}{\alpha_{N,j}}}{1+\frac{t}{\alpha_{N,j}}}
=\prod_{j=1}^{m_N}\bigl(1-\pi_{N,j}(t)+\pi_{N,j}(t)s\bigr),
\]
右侧即独立 Bernoulli 和的概率母函数，故得到分布恒等式。证毕。
\end{proof}

\input{subsubsec__xi-prime-hardcore-normalization-mertens-clt}

\begin{theorem}[Lee--Yang 单位圆提升、Joukowsky 椭圆输运与 \texorpdfstring{$\delta_{\mathrm{ZG}}$}{deltaZG} 的谱行列式表达]\label{con:xi-zg-unit-circle-joukowsky-spectral-determinant}
在定理 \ref{con:xi-zg-weighted-path-leyang-poisson-binomial} 的记号下：
\begin{enumerate}
  \item 设 \(\chi_{N+1}(x)=\det(xI-A_{N+1})\)，\(\Phi_{N+1}(x):=\chi_{N+1}(x/\sqrt2)\)。
  由于 \(\|A_{N+1}\|_{2}\le 2\max_k w_k=\sqrt2\)，\(\Phi_{N+1}\) 全部根落在 \([-2,2]\)。
  因此
  \[
  \mathcal L_{N+1}(z)
  :=z^{N+1}\Phi_{N+1}(z+z^{-1})
  =z^{N+1}\chi_{N+1}\!\Bigl(\frac{z+z^{-1}}{\sqrt2}\Bigr)
  \]
  的全部零点位于单位圆 \(|z|=1\)。
  \item 对任意 \(r>1\)，令 \(J_r(z):=rz+r^{-1}z^{-1}\)。
  若 \(z_j\) 为 \(\mathcal L_{N+1}\) 的零点，则 \(w_j:=J_r(z_j)\) 落在椭圆 \(J_r(S^1)\) 上。
  \item 取 \(x=i\) 于恒等式
  \(
  \chi_{N+1}(x)=x^{N+1}P_N(-x^{-2})
  \)
  得
  \[
  P_N(1)=|\chi_{N+1}(i)|=|\det(iI-A_{N+1})|.
  \]
  因而
  \[
  \delta_{\mathrm{ZG}}
  =\lim_{N\to\infty}
  \Bigl(\prod_{k=1}^N\Bigl(1-\frac1{p_k}\Bigr)\Bigr)\,|\det(iI-A_{N+1})|.
  \]
  又
  \[
  |\det(iI-A_{N+1})|^2
  =\det\!\bigl((iI-A_{N+1})(-iI-A_{N+1})\bigr)
  =\det(I+A_{N+1}^2),
  \]
  故
  \[
  P_N(1)=\det(I+A_{N+1}^2)^{1/2},
  \quad
  \delta_{\mathrm{ZG}}
  =\lim_{N\to\infty}
  \Bigl(\prod_{k=1}^N\Bigl(1-\frac1{p_k}\Bigr)\Bigr)\det(I+A_{N+1}^2)^{1/2}.
  \]
\end{enumerate}
\end{theorem}
\begin{proof}
第一条由谱半径界和参数化 \(x=z+z^{-1}=2\cos\theta\) 直接得到。
第二条是 \(J_r\) 对单位圆的标准像性质。
第三条由 \(x=i\) 代入与行列式恒等式逐步化简得出。证毕。
\end{proof}

\begin{corollary}[\texorpdfstring{$s=1$}{s=1} 处单极点与 Dirichlet 留数常数]\label{cor:xi-zg-dirichlet-residue}
令 \(P(s)=\sum_{k\ge 1}p_k^{-s}\) 为 prime zeta，并令 \(G(s)\) 为定理 \ref{thm:xi-zg-dirichlet-log-renormalization} 的正则因子。
采用 prime zeta 在 \(s\to 1^+\) 的经典展开
\[
P(s)=\log\frac1{s-1}+B_{\mathsf P}+o(1),
\]
其中 \(B_{\mathsf P}\in\RR\) 为常数（可由 Mertens 定理推导；参见 \cite{Apostol1976,Titchmarsh1986RiemannZeta}）。
则 \(F(s)=\sum_{n\in\mathrm{ZG}}n^{-s}\) 在 \(s=1\) 处具有单极点，且其留数为
\[
\lim_{s\to1^+}(s-1)F(s)=\exp\!\bigl(B_{\mathsf P}+G(1)\bigr)\in(0,\infty).
\]
并且该留数等于 \(\mathrm{ZG}\) 的 Dirichlet 密度；结合定理 \ref{con:xi-zg-density-closed-form-tail-bound} 的自然密度存在性，两者一致。
\end{corollary}

\input{thm__xi-zg-abel-residue-log-density}

\begin{proposition}[临界重整化常数的可计算误差界]\label{prop:xi-zg-holographic-constant-tail-bound}
\leanverified{paper\_xi\_zg\_holographic\_constant\_tail\_bound}
取 \(s=1\)，并令 \(\Delta_N(1)\) 与 \(G(1)=\sum_{N\ge 2}\Delta_N(1)\) 如定理 \ref{thm:xi-zg-dirichlet-log-renormalization}。
定义截断
\[
G^{(N)}(1):=\sum_{k=2}^{N}\Delta_k(1).
\]
则存在绝对常数 \(C>0\) 使对所有 \(N\ge 3\) 有
\[
\bigl|G(1)-G^{(N)}(1)\bigr|
\le
C\sum_{k>N}\Bigl(\frac1{p_k^2}+\frac1{p_kp_{k-1}}\Bigr),
\]
从而正则因子 \(\exp(G(1))\) 可被一条显式可求和尾项在任意精度下有效逼近。
\end{proposition}

\begin{proof}
由定理 \ref{thm:xi-zg-dirichlet-log-renormalization} 的一致估计，在 \(s=1\) 时存在常数 \(C\) 使
\(
|\Delta_k(1)|
\le C\bigl(p_k^{-2}+(p_kp_{k-1})^{-1}\bigr)
\)。
对尾和求并界即得。证毕。
\end{proof}

\begin{theorem}[\(\mathrm{ZG}\) 的对数 Mertens 定律]\label{thm:xi-zg-log-mertens}
设
\[
H(x):=\sum_{\substack{n\le x\\ n\in\mathrm{ZG}}}\frac1n.
\]
则极限存在并满足
\[
\lim_{x\to\infty}\frac{H(x)}{\log x}=\delta_{\mathrm{ZG}},
\]
并且常数项可由 \(F(s)=\sum_{n\in\mathrm{ZG}}n^{-s}\) 在 \(s=1\) 的 Laurent 常数项显式识别：存在 \(C_{\mathrm{ZG}}\in\RR\) 使
\[
H(x)=\delta_{\mathrm{ZG}}\log x+C_{\mathrm{ZG}}+o(1)\qquad(x\to\infty).
\]
其中
\[
C_{\mathrm{ZG}}
=\frac{\mathfrak C_{\mathrm{HC}}}{\zeta(2)}
\left(\gamma-\frac{2\zeta'(2)}{\zeta(2)}+\kappa_{\mathrm{HC}}\right)
=\delta_{\mathrm{ZG}}
\left(\gamma-\frac{2\zeta'(2)}{\zeta(2)}+\kappa_{\mathrm{HC}}\right).
\]
\end{theorem}

\begin{proof}
由推论 \ref{cor:xi-zg-dirichlet-residue} 或推论 \ref{cor:xi-zg-hardcore-constant-residue}，\(F(s)=\sum_{n\in\mathrm{ZG}}n^{-s}\) 在 \(s=1\) 处的留数为 \(\delta_{\mathrm{ZG}}\)。
对非负系数 Dirichlet 级数应用标准 Tauberian 定理可得
\(
H(x)=\delta_{\mathrm{ZG}}\log x+C_{\mathrm{ZG}}+o(1)
\)，其中 \(C_{\mathrm{ZG}}\) 等于 \(F\) 在 \(s=1\) 的 Laurent 常数项。
而命题 \ref{prop:xi-zg-hardcore-laurent} 给出了该 Laurent 常数项的显式表达式，从而得到所示恒等式。证毕。
\end{proof}

\begin{conjecture}[正则因子的全半平面延拓与 Lee--Yang 型零点边界]\label{conj:xi-zg-leyang-godel-transport-boundary}
令 \(\mathcal R(s):=\exp(G(s))\)（定理 \ref{thm:xi-zg-dirichlet-log-renormalization}），其在 \(\Re(s)>\tfrac12\) 全纯且无零。
猜想 \(\mathcal R\) 可延拓为 \(\Re(s)>0\) 上的全纯无零函数；从而 \(F(s)=\exp(P(s))\mathcal R(s)\) 在 \(\Re(s)>0\) 的全部奇性结构由 prime zeta \(P(s)\) 的奇性传输唯一决定。
进一步地，对充分大的截断 \(N\)，将 \(F_N(s)\) 视为有限体积分配函数，则其归一化零点云在 \(N\to\infty\) 的极限中形成一条由 Joukowsky 椭圆输运诱导的边界曲线，并以 \(\Re(s)=\tfrac12\) 为临界对称轴，与 \(\mathcal R\) 的零点真空相容。
\end{conjecture}

\begin{conjecture}[弱局域约束下的 Erd\H{o}s--Kac 型高斯极限]\label{conj:xi-zg-erdos-kac-under-local-hardcore}
记
\[
\omega(n):=\#\{p:\ p\mid n\},
\qquad
\mathcal Z(X):=\{n\le X:\ n\in\mathrm{ZG}\}.
\]
猜想存在常数 \(c_0\in\RR\)，使得在 \(\mathcal Z(X)\) 上均匀抽样的 \(n\) 满足
\[
\frac{\omega(n)-\log\log X-c_0}{\sqrt{\log\log X}}
\Longrightarrow
\mathcal N(0,1)
\qquad(X\to\infty).
\]
该猜想表明：相邻素数不可同现的局域硬核约束在中心极限定律尺度上仅导致有限平移，而不改变标准化极限分布。
\end{conjecture}

\begin{remark}[可复现核验脚本]\label{rem:xi-zg-density-leyang-spectral-audit-script}
定理 \ref{con:xi-zg-density-closed-form-tail-bound}--\ref{con:xi-zg-unit-circle-joukowsky-spectral-determinant}
中的闭式恒等式与有限窗数值核验由脚本
\texttt{scripts/exp\_xi\_prime\_register\_zg\_density\_leyang\_spectrum\_audit.py}
给出，对应证书文件为
\texttt{artifacts/export/xi\_prime\_register\_zg\_density\_leyang\_spectrum\_audit.json}。
\par
定理 \ref{thm:xi-zg-dirichlet-log-renormalization} 与命题 \ref{prop:xi-zg-holographic-constant-tail-bound}
在 \(s=1\) 处的递推与尾项模板核验由脚本
\texttt{scripts/exp\_xi\_prime\_register\_zg\_dirichlet\_log\_renormalization\_audit.py}
给出，对应证书文件为
\texttt{artifacts/export/xi\_prime\_register\_zg\_dirichlet\_log\_renormalization\_audit.json}。
\par
定理 \ref{thm:xi-zg-zeta-factorization}--命题 \ref{prop:xi-zg-hardcore-laurent}
中 \(\mathfrak C_{\mathrm{HC}}\) 与 \(\kappa_{\mathrm{HC}}\) 的数值核验由脚本
\texttt{scripts/exp\_xi\_prime\_register\_zg\_hardcore\_zeta\_factorization\_constant\_audit.py}
给出，对应证书文件为
\texttt{artifacts/export/xi\_prime\_register\_zg\_hardcore\_zeta\_factorization\_constant\_audit.json}。
\end{remark}


\begin{theorem}[行列式素因子分解给出的椭圆面积守恒律]\label{thm:xi-terminal-ellipse-area-euler-prime-law}
\leanverified{paper\_xi\_terminal\_ellipse\_area\_euler\_prime\_law}
设 \(A\in\mathrm{GL}_2(\QQ)\)，单位圆盘 \(D:=\{x\in\RR^2:\|x\|\le 1\}\)，椭圆盘 \(E_A:=A(D)\)。
则
\[
\mathrm{Area}(E_A)=\pi\,|\det A|.
\]
并且若
\(
\det A=\pm\prod_p p^{v_p(\det A)}
\)
（仅有限个 \(v_p\neq 0\)），则
\[
\mathrm{Area}(E_A)=\pi\prod_p p^{v_p(\det A)}.
\]
\end{theorem}
\begin{proof}
二维 Lebesgue 测度在可逆线性变换下按雅可比因子 \( |\det A| \) 缩放，故
\(
\mathrm{Area}(A(D))=|\det A|\,\mathrm{Area}(D)=\pi|\det A|
\)。
其余结论由有理数乘法群的素因子分解直接得到。证毕。
\end{proof}

\begin{theorem}[有理椭圆合同类的二元不变量完备性]\label{thm:xi-terminal-rational-ellipse-two-invariant-completeness}
\leanverified{paper\_xi\_terminal\_rational\_ellipse\_two\_invariant\_completeness}
设 \(A\in\mathrm{GL}_2(\QQ)\)，记 \(E_A:=A(S^1)\subset\RR^2\)。
定义
\[
\Delta_1(A):=(\det A)^2\in\QQ_{>0},
\qquad
\Delta_2(A):=\mathrm{tr}(A^\top A)\in\QQ_{>0}.
\]
则 \(E_A\) 的欧氏合同类由 \((\Delta_1(A),\Delta_2(A))\) 完全决定。
若 \(\sigma_1\ge \sigma_2>0\) 为 \(A\) 的奇异值，则
\[
\sigma_1^2,\sigma_2^2\ \text{是}\ t^2-\Delta_2(A)t+\Delta_1(A)=0\ \text{的两根},
\]
因而
\[
\sigma_{1,2}
=\sqrt{\frac{\Delta_2(A)\pm\sqrt{\Delta_2(A)^2-4\Delta_1(A)}}{2}}.
\]
特别地，对任意 \(A,B\in\mathrm{GL}_2(\QQ)\),
\[
E_A\cong E_B\ (\text{欧氏合同})
\Longleftrightarrow
\bigl(\Delta_1(A),\Delta_2(A)\bigr)=\bigl(\Delta_1(B),\Delta_2(B)\bigr).
\]
\end{theorem}
\begin{proof}
极分解 \(A=UP\)（\(U\) 正交，\(P=(A^\top A)^{1/2}\) 对称正定）给出
\(
E_A=U(P(S^1))
\)；
故合同类由 \(P\) 决定。
\(P\) 的特征值为 \(\sigma_1,\sigma_2\)，即椭圆半轴长。
又
\[
\sigma_1^2+\sigma_2^2=\mathrm{tr}(A^\top A)=\Delta_2(A),
\qquad
\sigma_1^2\sigma_2^2=\det(A^\top A)=(\det A)^2=\Delta_1(A),
\]
因此 \((\sigma_1,\sigma_2)\) 由 \((\Delta_1,\Delta_2)\) 唯一恢复，反之亦然。证毕。
\end{proof}

% AUTO-GENERATED by scripts/exp_xi_prime_register_semidirect_two_shadow_audit.py
\begin{table}[H]
\centering
\scriptsize
\setlength{\tabcolsep}{6pt}
\caption{有理矩阵样本的椭圆二元不变量审计：$\Delta_1=(\det A)^2$ 与 $\Delta_2=\mathrm{tr}(A^\top A)$。}
\label{tab:xi_elliptic_two_shadow_invariants_audit}
\begin{tabular}{l l l l r r}
\toprule
id & $\det A$ & $\Delta_1(A)$ & $\Delta_2(A)$ & $\sigma_1$ (num.) & $\sigma_2$ (num.)\\
\midrule
$A_1$ & $209/120$ & $43681/14400$ & $14701/3600$ & 1.762791958 & 0.988016004\\
$B_1=A_1Q_{90}$ & $209/120$ & $43681/14400$ & $14701/3600$ & 1.762791958 & 0.988016004\\
$A_2$ & $163/70$ & $26569/4900$ & $233209/44100$ & 1.974001870 & 1.179619667\\
\bottomrule
\end{tabular}
\end{table}


\begin{theorem}[有限素部格与实部椭圆影的二影刚性]\label{thm:xi-terminal-finite-archimedean-two-shadow-rigidity}
定义 prime-register 纤维
\[
\widehat{\ZZ}:=\varprojlim_m \ZZ/m\ZZ\cong\prod_p \ZZ_p.
\]
对 \(A\in\mathrm{GL}_2(\QQ)\)，记有限素部格影
\[
\Lambda_A:=A\,\widehat{\ZZ}^{\,2}\subset \mathbb A_f^{\,2},
\]
并记实部椭圆影 \(E_A:=A(S^1)\subset\RR^2\)。
若 \(A,B\in\mathrm{GL}_2(\QQ)\) 满足
\[
\Lambda_A=\Lambda_B,
\qquad
E_A=E_B,
\]
则存在 \(Q\in O_2(\ZZ)\) 使 \(A=BQ\)。
若再施加定向保持规范（等价于 \(\det(B^{-1}A)=1\)），则 \(Q\in SO_2(\ZZ)\)，因此仅余 \(4\) 个离散候选；
进一步施加任一破缺该四元对称的规范锚定后，\(Q\) 唯一，从而 \(A\) 由 \((\Lambda_A,E_A)\) 唯一确定。
\end{theorem}
\begin{proof}
令 \(C:=B^{-1}A\)。
由 \(\Lambda_A=\Lambda_B\) 得
\(
C\,\widehat{\ZZ}^{\,2}=\widehat{\ZZ}^{\,2}
\)，故 \(C\in\mathrm{GL}_2(\widehat{\ZZ})\)。

由 \(E_A=E_B\) 得同一椭圆二次型：
\[
x^\top(AA^\top)^{-1}x=1
\Longleftrightarrow
x^\top(BB^\top)^{-1}x=1,
\]
从而 \(AA^\top=BB^\top\)。
左乘 \(B^{-1}\)、右乘 \((B^{-1})^\top\) 得
\[
CC^\top=I,
\]
即 \(C\in O_2(\QQ)\)。

同时 \(C\in\mathrm{GL}_2(\widehat{\ZZ})\) 蕴含其每个有理条目在所有 \(\ZZ_p\) 中整，故条目属于
\(
\QQ\cap\bigcap_p\ZZ_p=\ZZ
\)；
即 \(C\in\mathrm{GL}_2(\ZZ)\)。
综上 \(C\in O_2(\ZZ)\)，即 \(A=BQ\)（\(Q:=C\in O_2(\ZZ)\)）。

当 \(\det Q=1\) 时 \(Q\in SO_2(\ZZ)\)。
二维整正交群中 \(|SO_2(\ZZ)|=4\)，故仅有四个定向保持候选；对任意固定规范锚定（例如长轴正向锚定）即可消去残余离散歧义并得到唯一性。证毕。
\end{proof}

\paragraph{全对称 Belyi 覆盖、判别式闭式与多临界 Puiseux 缩放}
本段将一参数代数族
\(
\Pi_r(\lambda,y)=(\lambda-1)^r+(y-1)\lambda^{r-1}
\)
组织为三点分歧覆盖，并给出其判别式闭式、Puiseux 多临界指数与 \(r=3\) 的显式 Lee--Yang 区间模型。

\begin{theorem}[全对称 Belyi 覆盖与满对称单值群]\label{thm:xi-terminal-belyi-full-symmetric-cover}
设 \(r\ge 2\) 为整数，并定义
\[
\Pi_r(\lambda,y):=(\lambda-1)^r+(y-1)\lambda^{r-1}\in\ZZ[\lambda,y],
\qquad
f_r(\lambda):=1-\frac{(\lambda-1)^r}{\lambda^{r-1}}.
\]
则成立：
\begin{enumerate}
  \item \(f_r:\PP^1_\lambda\to\PP^1_y\) 为次数 \(r\) 的三点分歧覆盖；其有限分歧值恰为
  \[
  y=1,\qquad y=y_r:=1+\frac{r^r}{(r-1)^{r-1}},
  \]
  并与 \(y=\infty\) 共同构成全部分歧值。
  \item 三处分歧值的分歧型分别为
  \[
  y=1:\ (r),\qquad
  y=\infty:\ (r-1)(1),\qquad
  y=y_r:\ (2)(1^{r-2}).
  \]
  \item 覆盖的几何单值群为
  \[
  \mathrm{Mon}_{\mathrm{geo}}(f_r)\cong S_r.
  \]
  因而 \(\Pi_r(\lambda,y)\) 在 \(\CC(y)[\lambda]\) 上不可约。
\end{enumerate}
\end{theorem}
\begin{proof}
写
\[
f_r(\lambda)=\frac{\lambda^{r-1}-(\lambda-1)^r}{\lambda^{r-1}}=\frac{P_r(\lambda)}{Q_r(\lambda)},
\]
其中 \(\deg P_r=r\)、\(\deg Q_r=r-1\)，且 \(P_r(0)\neq 0\)，故 \(\deg f_r=r\)。

直接计算得
\[
f_r'(\lambda)=-\frac{(\lambda-1)^{r-1}(\lambda+r-1)}{\lambda^r}.
\]
因此有限临界点为 \(\lambda=1\)（阶 \(r-1\)）与 \(\lambda=1-r\)（简单临界点）：
在 \(\lambda=1\) 附近令 \(\mu=\lambda-1\)，则
\[
f_r(1+\mu)-1=-\frac{\mu^r}{(1+\mu)^{r-1}}=-\mu^r+O(\mu^{r+1}),
\]
故 \(y=1\) 的分歧型为 \((r)\)。
在 \(\lambda=1-r\) 处，
\[
y_r=f_r(1-r)=1+\frac{r^r}{(r-1)^{r-1}},
\]
且临界点简单，故对应换位型 \((2)(1^{r-2})\)。

对 \(y=\infty\)，\(\lambda=0\) 为阶 \(r-1\) 极点，\(\lambda=\infty\) 为简单极点（因 \(f_r(\lambda)\sim-\lambda\)），故分歧型为 \((r-1)(1)\)。
因此仅有三处分歧值。

由分歧型可选取分支循环 \(\sigma_1,\sigma_{y_r},\sigma_\infty\in S_r\) 满足
\[
\sigma_1\sim (r),\qquad
\sigma_{y_r}\sim (2)(1^{r-2}),\qquad
\sigma_\infty\sim (r-1)(1),\qquad
\sigma_1\sigma_{y_r}\sigma_\infty=1.
\]
群 \(\langle \sigma_1,\sigma_{y_r}\rangle\) 含 \(r\)-循环与换位，故为 \(S_r\)。
因此几何单值群为 \(S_r\)。
单值作用传递即等价于 \(\Pi_r\) 在 \(\CC(y)[\lambda]\) 上不可约。证毕。
\end{proof}

\begin{theorem}[判别式的闭式分解与双有限分歧点]\label{thm:xi-terminal-belyi-discriminant-closed-form}
\leanverified{paper\_xi\_terminal\_belyi\_discriminant\_closed\_form}
在定理 \ref{thm:xi-terminal-belyi-full-symmetric-cover} 的记号下，存在符号常数 \(\varepsilon_r\in\{\pm1\}\) 使
\[
\mathrm{Disc}_{\lambda}(\Pi_r)(y)
=\varepsilon_r\,(y-1)^{r-1}
\Bigl((r-1)^{r-1}y-\bigl((r-1)^{r-1}+r^r\bigr)\Bigr).
\]
特别地
\[
\mathrm{Disc}_{\lambda}(\Pi_r)(y)=0
\iff
y\in\left\{1,\ 1+\frac{r^r}{(r-1)^{r-1}}\right\},
\]
且在 \(y=1\) 的消失阶为 \(r-1\)，在 \(y=y_r\) 的消失阶为 \(1\)。
\end{theorem}
\begin{proof}
由判别式恒等式
\(
\mathrm{Disc}(P)=(-1)^{r(r-1)/2}\mathrm{Res}(P,P')
\)，只需判定 \(P=\Pi_r(\cdot,y)\) 与 \(P'\) 的公共根位置与重数。

若 \(y=1\)，则 \(P(\lambda)=(\lambda-1)^r\)，故判别式在 \(y=1\) 至少有 \((r-1)\) 重零点。
若 \(y\neq1\)，由
\[
P(\lambda)=0,\qquad
P'(\lambda)=r(\lambda-1)^{r-1}+(y-1)(r-1)\lambda^{r-2}=0
\]
消去 \((\lambda-1)^{r-1}\) 可得唯一候选 \(\lambda=1-r\)，再代回得 \(y=y_r\)。
且该临界点简单，故对应判别式一次消失。
因此判别式除符号常数外恰为所示两因子乘积。证毕。
\end{proof}

% AUTO-GENERATED by scripts/exp_xi_belyi_full_symmetric_puiseux_audit.py
\begin{equation}\label{eq:xi_terminal_belyi_family_discriminant_closed}
\mathrm{Disc}_{\lambda}(\Pi_r)(y)=\varepsilon_r\,(y-1)^{r-1}\Bigl((r-1)^{r-1}y-\bigl((r-1)^{r-1}+r^r\bigr)\Bigr),\quad \varepsilon_r\in\{\pm1\}.
\end{equation}

\begin{equation}\label{eq:xi_terminal_belyi_r3_recurrence_discriminant}
\begin{aligned}
\Pi_3(\lambda,y)&=(\lambda-1)^3+(y-1)\lambda^2=\lambda^3+(y-4)\lambda^2+3\lambda-1,\\
Z_{n+3}(y)&+(y-4)Z_{n+2}(y)+3Z_{n+1}(y)-Z_n(y)=0,\\
Z_0(y)&=3,\qquad Z_1(y)=4 - y,\qquad Z_2(y)=y^{2} - 8 y + 10,\\
\mathrm{Disc}_{\lambda}(\Pi_3)(y)&=(y-1)^2(4y-31).
\end{aligned}
\end{equation}

% AUTO-GENERATED by scripts/exp_xi_belyi_full_symmetric_puiseux_audit.py
\begin{table}[H]
\centering
\scriptsize
\setlength{\tabcolsep}{5pt}
\caption{Belyi 三点分歧族 \(\Pi_r\)（\(2\le r\le 9\)）的判别式因子、局部分歧与置换生成审计。}
\label{tab:xi_belyi_branch_discriminant_group_audit_r2_r9}
\begin{tabular}{r l r c c c r c}
\toprule
$r$ & $y_r$ & $\varepsilon_r$ & Disc & $e_{y_r}=2$ & $\sigma_\infty$ type & $|\langle\sigma_1,\sigma_{y_r}\rangle|$ & $=r!$\\
\midrule
2 & $5$ & 1 & 1 & 1 & $(1)(1)$ & 2 & 1\\
3 & $31/4$ & 1 & 1 & 1 & $(2)(1)$ & 6 & 1\\
4 & $283/27$ & -1 & 1 & 1 & $(3)(1)$ & 24 & 1\\
5 & $3381/256$ & -1 & 1 & 1 & $(4)(1)$ & 120 & 1\\
6 & $49781/3125$ & 1 & 1 & 1 & $(5)(1)$ & 720 & 1\\
7 & $870199/46656$ & 1 & 1 & 1 & $(6)(1)$ & 5040 & 1\\
8 & $17600759/823543$ & -1 & 1 & 1 & $(7)(1)$ & 40320 & 1\\
9 & $404197705/16777216$ & -1 & 1 & 1 & $(8)(1)$ & 362880 & 1\\
\bottomrule
\end{tabular}
\end{table}


\paragraph{一类六次 \texorpdfstring{$j$}{j}-覆叠的判别式指纹、分歧值与 \texorpdfstring{$S_6$}{S6} 正则性}
定义有理函数
\[
j(a):=-\frac{4096\,(a^2-a+1)^3}{(a-1)^2(16a^3-13a^2-78a+43)}\ \in\ \QQ(a),
\]
并对超越参数 \(t\) 令
\[
P_t(a):=-4096\,(a^2-a+1)^3-t\,(a-1)^2(16a^3-13a^2-78a+43)\ \in\ \ZZ[t][a],
\]
即 \(P_t(a)=0\) 等价于 \(j(a)=t\) 的交叉相乘形式。记
\(
\Delta(t):=\Disc_a(P_t)\in\ZZ[t]
\)。

\begin{theorem}[判别式全因子分解、分歧值代数指纹与满对称单值群]\label{thm:xi-j-sextic-discriminant-s6-regular}
以上记号下成立：
\begin{enumerate}
  \item 判别式存在完全显式分解
  \[
  \boxed{\;
  \Delta(t)=2^{43}\,79^3\cdot t^{4}\,(t-1728)^{3}\,\bigl(t^2+1862t+161792\bigr).\;}
  \]
  因而 \(j:\PP^1_a\to\PP^1_t\) 的有限分歧值恰为
  \[
  t\in\left\{0,\ 1728,\ -931\pm 89\sqrt{89}\right\}.
  \]
  \item 令 \(K=\QQ(t)\)，并令 \(L\) 为 \(P_t(a)\) 在 \(K\) 上的分裂域，则该扩张为正则且
  \[
  \boxed{\;\Gal(L/K)\ \cong\ S_6.\;}
  \]
  \item \(j(a)\) 在 \(\overline{\QQ}\) 上不可分解为非平凡复合 \(j=F\circ G\)（\(\deg F,\deg G\ge2\)），并且满足 \(j(\phi(a))\equiv j(a)\) 的 \(\phi\in\mathrm{PGL}_2(\overline{\QQ})\) 只能为恒等映射。
  \item \(L/K\) 的唯一二次子域由判别式平方根生成：
  \[
  \boxed{\;L^{A_6}=\QQ\bigl(t,\sqrt{\Delta(t)}\bigr).\;}
  \]
  在平方类 \(\QQ(t)^\times/(\QQ(t)^\times)^2\) 中，
  \[
  \sqrt{\Delta(t)}\ \sim\ \sqrt{158\,(t-1728)\,\bigl(t^2+1862t+161792\bigr)},
  \]
  从而除 \(t\in\{0,1728,\infty\}\) 外的代数分歧由二次域 \(\QQ(\sqrt{89})\) 控制。
\end{enumerate}
\end{theorem}
\begin{proof}
第 1 条等价于对线性笔 \(P_t=f+t g\) 消去公共根条件 \(P_t=\partial_aP_t=0\)，可由判别式或结式计算给出；二次因子判别式为 \(1862^2-4\cdot161792=4\cdot 89^3\)，故其根为 \(-931\pm 89\sqrt{89}\)。
\par
对第 2 条，正则性由 \(j\) 为 \(\PP^1\) 上有理覆叠且常数域不扩张得到。取 \(t_0=12345\)，则 \(P_{t_0}\in\ZZ[a]\) 为可分离六次多项式；对不整除 \(\Disc(P_{t_0})\) 的素数 \(p\)，其模 \(p\) 因子型给出 Frobenius 的循环类型。直接检验可取 \(p=13\) 使 \(P_{t_0}\) 在 \(\FF_{13}[a]\) 上不可约（给出 \(6\)-循环），并可取 \(p=149\) 使其因子型为 \((2)(1)(1)(1)(1)\)（给出换位）；含 \(6\)-循环与换位的传递子群只能为 \(S_6\)，故某一特化的 Galois 群为 \(S_6\)。由 Hilbert 不可约性，函数域 Galois 群亦为 \(S_6\)。
\par
第 3 条：非平凡分解等价于存在非平凡中间覆叠，亦即单值作用非本原；但 \(S_6\) 的自然作用本原，故不可分解。甲板变换群同构于单值群在自然置换表示中的中心化子；\(S_6\) 的中心化子平凡，故仅有恒等自同构。
\par
第 4 条：\(S_6\) 的唯一指数 \(2\) 正规子群为 \(A_6\)，对应的二次子域由符号特征所给；对任意六次可分离多项式，该二次子域由 \(\sqrt{\Delta(t)}\) 生成。代入第 1 条分解并在平方类中消去平方因子 \(t^4\) 与常数平方部分即得所示平方类代表。证毕。
\end{proof}

\paragraph{判别式二次子域的椭圆化、属 \(5\) 正规化与 \(A_6\) 基变压缩}
本段固定定理 \ref{thm:xi-j-sextic-discriminant-s6-regular} 的六次有理覆叠 \(j:\PP^1_a\to\PP^1_t\) 与其判别式因子
\[
f(t):=(t-1728)\bigl(t^2+1862t+161792\bigr)\in\ZZ[t].
\]

\begin{theorem}[判别式二次子域的椭圆化与 \(2\)-挠点分裂域]\label{thm:xi-j-sextic-discriminant-quadratic-subfield-ellipticization}
沿用定理 \ref{thm:xi-j-sextic-discriminant-s6-regular} 的记号。
令
\[
E_0:\quad y^2=f(t)=(t-1728)\bigl(t^2+1862t+161792\bigr)
=t^3+134t^2-3055744t-279576576,
\]
并令
\[
E_\Delta:\quad Y^2=158\,f(t).
\]
则：
\begin{enumerate}
  \item 二次子域 \(\QQ(t,\sqrt{\Delta(t)})\) 与 \(E_\Delta\) 的函数域同构：
  \[
  \boxed{\ \QQ\bigl(t,\sqrt{\Delta(t)}\bigr)\ \cong\ \QQ(E_\Delta)\ }.
  \]
  \item \(E_0\) 是 \(E_\Delta\) 的二次扭曲；二者的 \(2\)-挠点的 \(t\)-坐标集合完全一致，恰为 \(f(t)=0\) 的根。
  \item \(E_0(\QQ)\) 恰有一个非平凡 \(\QQ\)-有理 \(2\)-挠点 \((1728,0)\)，并且
  \[
  E_0(\QQ)[2]\cong \ZZ/2\ZZ,\qquad
  E_0\bigl(\QQ(\sqrt{89})\bigr)[2]\cong (\ZZ/2\ZZ)^2,
  \]
  且额外两枚有限 \(2\)-挠点的 \(t\)-坐标为
  \[
  \boxed{\ t=-931\pm 89\sqrt{89}\ }.
  \]
  \item 若以 Weierstrass 形式 \(E_0:y^2=t^3+134t^2-3055744t-279576576\) 表示，则其判别式满足
  \[
  \boxed{\ \Delta(E_0)=2^{20}\cdot 89^{3}\cdot 223^{4}\ }.
  \]
\end{enumerate}
\end{theorem}
\begin{proof}
由定理 \ref{thm:xi-j-sextic-discriminant-s6-regular}，
\[
\Delta(t)=2^{43}\,79^3\cdot t^{4}\,(t-1728)^{3}\,\bigl(t^2+1862t+161792\bigr).
\]
因此在 \(\QQ(t)^\times/(\QQ(t)^\times)^2\) 中有
\[
\sqrt{\Delta(t)}\ \sim\ \sqrt{(2\cdot 79)\,(t-1728)\bigl(t^2+1862t+161792\bigr)}
=\sqrt{158\,f(t)}.
\]
从而 \(\QQ(t,\sqrt{\Delta(t)})=\QQ(t,\sqrt{158\,f(t)})\cong \QQ(E_\Delta)\)，得第 (1) 条。
第 (2) 条由定义直接成立：把方程 \(y^2=f(t)\) 乘以常数 \(158\) 得到 \(E_\Delta\)，扭曲不改变根集合 \(f(t)=0\)，故 \(2\)-挠点的 \(t\)-坐标不变。
\par
对第 (3) 条，\(2\)-挠点等价于右端三次多项式的根（令 \(y=0\)）。
由于 \(f(t)=(t-1728)(t^2+1862t+161792)\)，显然存在且仅存在一个 \(\QQ\)-有理根 \(t=1728\)，对应唯一非平凡 \(\QQ\)-有理 \(2\)-挠点。
二次因子判别式为
\[
1862^2-4\cdot 161792=4\cdot 89^3,
\]
其平方自由部分为 \(89\)，故其两根属于且仅属于 \(\QQ(\sqrt{89})\)，并且显式为 \(-931\pm 89\sqrt{89}\)。
\par
第 (4) 条为对该整系数 Weierstrass 模型的判别式直接计算与整数因子分解。证毕。
\end{proof}

\input{parts/thm__xi-j-discriminant-elliptic-2torsion-mobius-velu}
\input{thm__xi-j-sextic-discriminant-elliptic-lattes-t-map}

\begin{theorem}[交错分歧的平方消去与属 \(5\) 超椭圆正规化]\label{thm:xi-j-sextic-fiber-product-genus5-hyperelliptic-model}
令 \(j(a)\) 如定理 \ref{thm:xi-j-sextic-discriminant-s6-regular}。
记
\[
r(a):=16a^3-13a^2-78a+43,
\qquad
H(a):=2048a^8-8752a^7+16455a^6-5030a^5-39667a^4+82096a^3-74559a^2+33406a-5869.
\]
并定义曲线
\[
\boxed{\ C:\quad Y^2=-2^{17}\,r(a)\,H(a)\ }.
\]
则：
\begin{enumerate}
  \item \(C\) 为光滑超椭圆曲线，且
  \[
  \boxed{\ g(C)=5\ }.
  \]
  \item 令 \(\pi:E_\Delta\to\PP^1_t\) 为定理 \ref{thm:xi-j-sextic-discriminant-quadratic-subfield-ellipticization} 中的二次覆叠（以 \(t\) 为坐标），则 \(C\) 与纤维积
  \[
  \PP^1_a\times_{\PP^1_t}E_\Delta
  \qquad(t=j(a))
  \]
  的正规化在双有理意义下等价。
\end{enumerate}
\end{theorem}
\begin{proof}
令
\[
s(a):=8a^3+15a^2-66a+35,\qquad p(a):=2a^2-9a-1.
\]
直接代入并约分可得两条恒等式：
\[
j(a)-1728=-\frac{64\,s(a)^2}{(a-1)^2\,r(a)},
\qquad
j(a)^2+1862j(a)+161792=\frac{2048\,p(a)^2\,H(a)}{(a-1)^4\,r(a)^2}.
\]
从而
\[
158\,(j(a)-1728)\bigl(j(a)^2+1862j(a)+161792\bigr)
\;=\;
-\frac{2^{17}\cdot 79\, s(a)^2p(a)^2\,H(a)}{(a-1)^6\,r(a)^3}.
\]
在函数域层面，把右端乘以平方因子 \(\bigl((a-1)^3r(a)s(a)p(a)\bigr)^{-2}\) 不改变二次扩张，于是得到与之同一二次扩张的模型
\[
Y^2=-2^{17}\,r(a)\,H(a),
\]
即为所定义的 \(C\)。
\par
又因 \(\gcd(r,H)=1\) 且 \(rH\) 平方自由（等价于分支多项式无重根），并且 \(\deg(rH)=11\)，故 \(C\) 为属
\[
g(C)=\left\lfloor\frac{11-1}{2}\right\rfloor=5
\]
的光滑超椭圆曲线。
最后，按构造 \(E_\Delta\) 的函数域为 \(\QQ(t,\sqrt{158f(t)})\)；把 \(t=j(a)\) 代入并消去平方因子后得到的正规化恰为上述 \(C\)，从而与纤维积正规化双有理等价。证毕。
\end{proof}

\begin{theorem}[属 \(5\) 超椭圆曲线到判别式椭圆曲线的显式六次覆叠]\label{thm:xi-j-sextic-genus5-to-elliptic-explicit-degree6-cover}
沿用定理 \ref{thm:xi-j-sextic-fiber-product-genus5-hyperelliptic-model} 的记号，并令
\[
E_0:\quad y^2=f(t)=(t-1728)\bigl(t^2+1862t+161792\bigr).
\]
定义多项式
\[
s(a):=8a^3+15a^2-66a+35,\qquad p(a):=2a^2-9a-1.
\]
则存在一个在 \(\QQ\) 上定义的非平凡态射
\[
g_0:C\longrightarrow E_0
\]
由下式完全显式给出：
\[
\boxed{\quad t=j(a),\qquad y=\frac{s(a)p(a)}{(a-1)^3\,r(a)^2}\,Y.\quad}
\]
并且 \(g_0\) 为次数 \(6\) 的有限覆叠，即 \(\deg(g_0)=6\)。
\end{theorem}
\begin{proof}
由定理 \ref{thm:xi-j-sextic-fiber-product-genus5-hyperelliptic-model} 证明中给出的恒等式
\[
j(a)-1728=-\frac{64\,s(a)^2}{(a-1)^2\,r(a)},
\qquad
j(a)^2+1862j(a)+161792=\frac{2048\,p(a)^2\,H(a)}{(a-1)^4\,r(a)^2}
\]
可得
\[
\bigl(j(a)-1728\bigr)\bigl(j(a)^2+1862j(a)+161792\bigr)
\,=\,
-\frac{2^{17}\,s(a)^2p(a)^2\,H(a)}{(a-1)^6\,r(a)^3}.
\]
另一方面，曲线 \(C\) 满足 \(Y^2=-2^{17}r(a)H(a)\)，从而
\[
\left(\frac{s(a)p(a)}{(a-1)^3\,r(a)^2}Y\right)^2
=
-\frac{2^{17}\,s(a)^2p(a)^2\,H(a)}{(a-1)^6\,r(a)^3}
=
f\bigl(j(a)\bigr).
\]
于是 \((t,y)=(j(a),\frac{s(a)p(a)}{(a-1)^3r(a)^2}Y)\) 确满足 \(E_0\) 的方程，故 \(g_0\) 的值落于 \(E_0\)。

对次数断言，\(t=j(a)\) 作为从 \(\PP^1_a\) 到 \(\PP^1_t\) 的有理函数具有次数 \(6\)（见定理 \ref{thm:xi-j-sextic-discriminant-s6-regular}），故在函数域层面
\(
[\QQ(a):\QQ(t)]=6
\)。
又 \(\QQ(C)=\QQ(a,Y)\) 为 \(\QQ(a)\) 的二次扩张，而 \(\QQ(E_0)=\QQ(t,y)\) 为 \(\QQ(t)\) 的二次扩张。
由上式 \(y\in \QQ(a,Y)\) 得到嵌入 \(\QQ(E_0)\hookrightarrow \QQ(C)\)，并且
\[
[\QQ(C):\QQ(t)]=2\,[\QQ(a):\QQ(t)]=12.
\]
因此
\(
[\QQ(C):\QQ(E_0)]=12/2=6
\)，从而 \(\deg(g_0)=6\)。证毕。
\end{proof}

\begin{theorem}[不变量微分的拉回消去]\label{thm:xi-j-sextic-invariant-differential-pullback-quartic-factor}
令
\[
\omega_{E_0}:=\frac{\dd t}{y}\in H^0(E_0,\Omega^1_{E_0})
\]
为 \(E_0\) 的标准不变量微分。
则对定理 \ref{thm:xi-j-sextic-genus5-to-elliptic-explicit-degree6-cover} 的态射 \(g_0:C\to E_0\)，在 \(C\) 上有完全显式的恒等式
\[
\boxed{\quad g_0^*(\omega_{E_0})=-4096\,(a^2-a+1)^2\,\frac{\dd a}{Y}.\quad}
\]
\end{theorem}
\begin{proof}
由 \(t=j(a)\) 得 \(\dd t=j'(a)\dd a\)。
另一方面，定理 \ref{thm:xi-j-sextic-genus5-to-elliptic-explicit-degree6-cover} 给出
\(
y=\frac{s(a)p(a)}{(a-1)^3r(a)^2}Y
\)，从而
\[
g_0^*\!\left(\frac{\dd t}{y}\right)
=
j'(a)\,\frac{(a-1)^3r(a)^2}{s(a)p(a)}\,\frac{\dd a}{Y}.
\]
对 \(j(a)=-4096\frac{(a^2-a+1)^3}{(a-1)^2r(a)}\) 直接求导并化简（可符号化简逐式复核）得到恒等式
\[
j'(a)=-4096\,(a^2-a+1)^2\,\frac{s(a)p(a)}{(a-1)^3\,r(a)^2}.
\]
代入即得所示拉回公式。证毕。
\end{proof}

\begin{corollary}[分歧除子与有效 theta 特征]\label{cor:xi-j-sextic-ramification-divisor-theta-characteristic}
令 \(\infty\in C\) 为超椭圆模型的无穷远点（因 \(\deg(rH)=11\) 为奇数而唯一）。
令 \(\alpha,\beta\) 为二次方程 \(a^2-a+1=0\) 的两根，并记其在 \(C\) 上的两点纤维为
\[
P_\alpha^\pm=(\alpha,\pm Y_\alpha),\qquad P_\beta^\pm=(\beta,\pm Y_\beta).
\]
则 \(g_0:C\to E_0\) 的分歧除子满足
\[
\boxed{\quad R_{g_0}=2\bigl(P_\alpha^++P_\alpha^-+P_\beta^++P_\beta^-\bigr),\qquad \deg R_{g_0}=8.\quad}
\]
并且该除子类与正则类一致：
\[
K_C\sim R_{g_0}.
\]
因此
\[
\boxed{\quad D:=P_\alpha^++P_\alpha^-+P_\beta^++P_\beta^- \ \text{给出一个有效 theta 特征：}\ 2D\sim K_C.\quad}
\]
\end{corollary}
\begin{proof}
在椭圆曲线 \(E_0\) 上，\(\omega_{E_0}\) 无零无极点，故 \(\Div(\omega_{E_0})=0\)。
一般覆叠理论给出
\(
\Div(g_0^*\omega_{E_0})=R_{g_0}
\)。
由定理 \ref{thm:xi-j-sextic-invariant-differential-pullback-quartic-factor}，
\[
\Div(g_0^*\omega_{E_0})
=
\Div\!\bigl((a^2-a+1)^2\bigr)+\Div\!\left(\frac{\dd a}{Y}\right).
\]
对超椭圆曲线 \(Y^2=f(a)\) 且 \(\deg f=2g+1\) 的标准事实给出
\(
\Div(\dd a/Y)=(2g-2)\,\infty
\)；此处 \(g=5\)，故 \(\Div(\dd a/Y)=8\,\infty\)。
又函数 \(a\) 在 \(\infty\) 处具有二阶极点，故 \((a^2-a+1)^2\) 在 \(\infty\) 处具有八阶极点，并在 \(a=\alpha,\beta\) 处分别给出二重零点，从而
\[
\Div\!\bigl((a^2-a+1)^2\bigr)
=
2\bigl(P_\alpha^++P_\alpha^-+P_\beta^++P_\beta^-\bigr)-8\,\infty.
\]
相加得到
\(
\Div(g_0^*\omega_{E_0})=2(P_\alpha^++P_\alpha^-+P_\beta^++P_\beta^-)
\)，即为 \(R_{g_0}\)。
又 \(K_{E_0}\sim 0\)，由 Riemann--Hurwitz 得 \(K_C\sim g_0^*K_{E_0}+R_{g_0}\sim R_{g_0}\)。证毕。
\end{proof}

\begin{proposition}[Jacobian 的显式椭圆因子与 \(\QQ\)-有理幂等元]\label{prop:xi-j-sextic-jacobian-elliptic-factor-idempotent}
设 \(J(C)\) 为 \(C\) 的 Jacobian。令
\[
g_{0*}:J(C)\to E_0,
\qquad
g_0^*:E_0\to J(C)
\]
分别为由 \(g_0\) 诱导的 push-forward 与 pull-back。
则同态满足
\[
\boxed{\quad g_{0*}\circ g_0^*=[6]_{E_0}.\quad}
\]
从而
\[
\boxed{\quad e:=\frac{1}{6}\,g_0^*\circ g_{0*}\ \in\ \End^0(J(C))\ \text{为幂等元： }e^2=e.\quad}
\]
并由此得到在 \(\QQ\) 上定义的同源分解
\[
\boxed{\quad J(C)\ \sim_{\QQ}\ g_0^*(E_0)\times P,\qquad P:=\ker(g_{0*})^0,\qquad \dim P=4.\quad}
\]
\end{proposition}
\begin{proof}
对任意曲线间有限态射 \(h\) 都有 \(h_*\circ h^*=[\deg h]\)。
应用于 \(h=g_0\) 且 \(\deg(g_0)=6\)，得到 \(g_{0*}g_0^*=[6]\)。
于是
\[
e^2=\frac{1}{36}g_0^*g_{0*}g_0^*g_{0*}
=\frac{1}{36}g_0^*[6]g_{0*}
=\frac{1}{6}g_0^*g_{0*}=e.
\]
由幂等元给出的像与核的同源分裂得到所示分解；维数断言由 \(\dim J(C)=g(C)=5\) 与 \(\dim E_0=1\) 得出。证毕。
\end{proof}

\begin{proposition}[良好约化处的 zeta 因子整除]\label{prop:xi-j-sextic-zeta-factor-divisibility}
令 \(S\) 为 \(C\) 与 \(E_0\) 的潜在坏素数集合（可取分支多项式判别式与椭圆判别式的素因子并集）。
对任意素数 \(\ell\notin S\)，在 \(\FF_\ell\) 上将 \((C,E_0,g_0)\) 约化得到 \((C_\ell,E_{0,\ell},g_{0,\ell})\)。
则在 \(\ell\)-进 \'{e}tale 上同调层面，拉回
\[
g_{0,\ell}^*:H^1_{\mathrm{\acute et}}(E_{0,\overline{\FF}_\ell},\QQ_\ell)\hookrightarrow H^1_{\mathrm{\acute et}}(C_{\overline{\FF}_\ell},\QQ_\ell)
\]
为单射且与 Frobenius 交换。
因此 \(C_\ell\) 的 Frobenius 特征多项式包含 \(E_{0,\ell}\) 的特征多项式因子；等价地，\(\zeta(C_\ell,T)\) 的分子多项式在 \(\ZZ[T]\) 中被 \(\zeta(E_{0,\ell},T)\) 的分子多项式整除。
\end{proposition}
\begin{proof}
在良好约化处，态射 \(g_{0,\ell}\) 诱导的 \(g_{0,\ell*}\circ g_{0,\ell}^*=[6]\) 在 \'{e}tale 上同调中仍成立。
从而 \(g_{0,\ell}^*\) 在 \(H^1\) 上为单射，并与 Frobenius 作用交换。
线性代数直接推出特征多项式的整除性。证毕。
\end{proof}

\begin{proposition}[分歧点的定义域与有理分歧除子]\label{prop:xi-j-sextic-branchpoints-field-and-rational-branch-divisor}
令 \(O_{E_0}\) 为 \(E_0\) 的无穷远点。
则 \(g_0:C\to E_0\) 的分歧点在 \(E_0\) 上恰为
\[
Q_\pm=(0,\pm\sqrt{f(0)})=(0,\pm 768\sqrt{-474}),
\]
并且
\(
\QQ(Q_\pm)=\QQ(\sqrt{-474})
\)，且 \(Q_-=\overline{Q_+}\)。
尽管两点各自不在 \(\QQ\) 上有理，分歧除子
\[
B:=Q_++Q_-
\]
却在 \(\QQ\) 上有理，且满足线性等价
\[
\boxed{\quad B\sim 2O_{E_0}.\quad}
\]
此外，推论 \ref{cor:xi-j-sextic-ramification-divisor-theta-characteristic} 中的四个分歧点按符号分裂：对每个 \(\alpha\in\{\text{\(a^2-a+1=0\) 的根}\}\)，两点 \(P_\alpha^\pm\) 分别映到 \(Q_\pm\)。
\end{proposition}
\begin{proof}
由定理 \ref{thm:xi-j-sextic-genus5-to-elliptic-explicit-degree6-cover} 的显式公式，
当 \(a^2-a+1=0\) 时有 \(t=j(a)=0\)，且 \(y=\frac{s(a)p(a)}{(a-1)^3r(a)^2}Y\) 随 \(Y\mapsto -Y\) 变号，故两点 \(P_\alpha^\pm\) 映到 \(t=0\) 纤维上的两点 \(Q_\pm\)。
直接计算
\(
f(0)=-279576576=-768^2\cdot 474
\)，得 \(Q_\pm\) 的定义域断言。
\par
函数 \(t\) 作为 \(E_0\) 的坐标函数在 \(O_{E_0}\) 处有二阶极点，且在 \(t=0\) 处有两简单零点，因此
\[
\Div(t)=Q_++Q_- - 2O_{E_0},
\]
从而 \(B\sim 2O_{E_0}\) 且 \(B\) 在 \(\QQ\) 上有理。证毕。
\end{proof}

\begin{conjecture}[四维因子的实乘结构与良好素数处分解型]\label{conj:xi-j-sextic-prym-real-multiplication-sqrt89}
令 \(P:=\ker(g_{0*})^0\) 为命题 \ref{prop:xi-j-sextic-jacobian-elliptic-factor-idempotent} 的四维因子。
猜想 \(\End^0_{\QQ}(P)\) 含有实二次域 \(\QQ(\sqrt{89})\) 的嵌入。
并且对任意良好约化素数 \(\ell\notin S\)，若 \(\bigl(\frac{89}{\ell}\bigr)=-1\)（\(\ell\) 在 \(\QQ(\sqrt{89})\) 中惰性），则 \(P_\ell\) 的 Frobenius 特征多项式在 \(\ZZ[T]\) 上通常不可约；
若 \(\bigl(\frac{89}{\ell}\bigr)=+1\)（\(\ell\) 分裂），则该特征多项式通常分裂为两个次数 \(4\) 的互为 Galois 共轭因子。
\end{conjecture}

\begin{theorem}[二次基变的奇置换消去：几何单值群从 \(S_6\) 压缩到 \(A_6\)，且仅余两处分歧]\label{thm:xi-j-sextic-basechange-compress-s6-to-a6-two-branchpoints}
令 \(f:\PP^1_a\to\PP^1_t\) 为 \(t=j(a)\) 给出的次数 \(6\) 覆叠，其几何单值群为 \(S_6\)（定理 \ref{thm:xi-j-sextic-discriminant-s6-regular}）。
令 \(\pi:E_\Delta\to\PP^1_t\) 为 \(E_\Delta:Y^2=158f(t)\) 给出的二次覆叠，并令 \(C\) 为定理 \ref{thm:xi-j-sextic-fiber-product-genus5-hyperelliptic-model} 的正规化纤维积。
则诱导得到次数 \(6\) 映射
\[
g_\Delta:C\to E_\Delta.
\]
并且：
\begin{enumerate}
  \item \(g_\Delta\) 的分歧点在 \(E_\Delta\) 上恰有两点，互为椭圆反演；在每个分歧点之上的置换型为 \((3)(3)\)，每点的分歧贡献为 \(4\)。
  \item 因而 \(g(C)=5\)（与定理 \ref{thm:xi-j-sextic-fiber-product-genus5-hyperelliptic-model} 的属数一致）。
  \item \(g_\Delta\) 的几何单值群包含于 \(A_6\)，且其 \(E_\Delta\) 上的 Galois 闭包的 Galois 群为 \(A_6\)。
\end{enumerate}
\end{theorem}
\begin{proof}
由定理 \ref{thm:xi-j-sextic-discriminant-s6-regular} 的判别式指数，
次数 \(6\) 覆叠 \(f\) 在五个分歧值处的总分歧量分别为：
\[
t=0:\ 4,\qquad
t=1728:\ 3,\qquad
t=-931\pm 89\sqrt{89}:\ 1,\qquad
t=\infty:\ 1.
\]
这与置换型
\[
t=0:\ (3)(3),\quad
t=1728:\ (2)(2)(2),\quad
t=-931\pm 89\sqrt{89}:\ (2),\quad
t=\infty:\ (2)
\]
相容，其中除 \(t=0\) 外均为奇置换。
\par
二次基变 \(\pi\) 的分歧恰发生在 \(t=1728\)、\(t=-931\pm 89\sqrt{89}\) 与 \(\infty\)。
在这些点上，小回路在基变后提升为绕两圈的回路，从而局部单值置换被替换为原置换的平方。
对换位与 \((2)(2)(2)\) 皆有平方为恒等，故上述四处分歧在拉回后被消去。
相反，\(t=0\) 处 \(\pi\) 不分歧，而局部单值为阶 \(3\) 的偶置换 \((3)(3)\)，故该分歧保留；
又因 \(\pi\) 在 \(t=0\) 处上方有两点，遂 \(g_\Delta\) 的分歧点恰为这两点且置换型仍为 \((3)(3)\)，得第 (1) 条。
\par
第 (2) 条由 Riemann--Hurwitz：
由于 \(g(E_\Delta)=1\)，并且两处分歧各贡献 \(4\)，故
\[
2g(C)-2=8\quad\Rightarrow\quad g(C)=5.
\]
\par
第 (3) 条：二次基变对应于把 \(\pi_1(\PP^1_t\setminus\{\text{分歧值}\})\) 限制到符号同态 \(\mathrm{sgn}:S_6\to\{\pm 1\}\) 的核，
等价于剔除奇单值贡献；因此拉回后的几何单值群包含于 \(A_6\)。
又因原群为 \(S_6\)，且 \(A_6\) 为其唯一指数 \(2\) 正规子群，故 \(E_\Delta\) 上的 Galois 闭包群为 \(A_6\)。证毕。
\end{proof}

\begin{proposition}[一组完全显式的分歧循环：一个 Nielsen 类代表]\label{prop:xi-j-sextic-explicit-branch-cycles-nielsen-rep}
存在对次数 \(6\) 纤维的标号，使得绕五个分歧值 \(t=0,1728,-931\pm 89\sqrt{89},\infty\) 的几何单值置换可取为下列五元组（右作用、从左到右相乘）：
\[
\boxed{
\begin{aligned}
\sigma_0&=(1,3,4)(2,6,5),&
\sigma_{1728}&=(1,6)(2,5)(3,4),\\
\sigma_{+}&=(1,5),&
\sigma_{-}&=(1,6),&
\sigma_{\infty}&=(1,3),
\end{aligned}}
\qquad
\boxed{\ \sigma_0\,\sigma_{1728}\,\sigma_{+}\,\sigma_{-}\,\sigma_{\infty}=1\ }.
\]
并且它们生成 \(S_6\)。
\end{proposition}
\begin{proof}
乘积恒等式可直接在 \(S_6\) 中验证。
对生成性，令
\(
c:=(1,4,5,2,3,6)
\)
为群 \(G:=\langle\sigma_0,\sigma_{1728},\sigma_{+},\sigma_{-},\sigma_{\infty}\rangle\) 中的一个 \(6\)-循环（由生成元显式相乘得到）。
取换位 \(\tau:=\sigma_{\infty}=(1,3)\in G\)。
则 \(c^k\tau c^{-k}\in G\) 给出若干换位；再结合 \(\sigma_-=(1,6)\) 可得到一族换位，其支撑边集合在顶点 \(\{1,\dots,6\}\) 上连通。
由“连通图的边换位生成整个对称群”得 \(G=S_6\)。证毕。
\end{proof}

\begin{corollary}[\(A_6\) 正则闭包曲线的属数与两处分歧]\label{cor:xi-j-sextic-a6-galois-closure-genus-241}
令 \(Y\to E_\Delta\) 为定理 \ref{thm:xi-j-sextic-basechange-compress-s6-to-a6-two-branchpoints} 中 \(g_\Delta:C\to E_\Delta\) 的 Galois 闭包。
则 \(\mathrm{Gal}(Y/E_\Delta)\cong A_6\)，并且 \(Y\to E_\Delta\) 仅在 \(E_\Delta\) 上那两点处分歧，每处分歧惯性群阶为 \(3\)。
因此
\[
\boxed{\ g(Y)=241\ }.
\]
\end{corollary}
\begin{proof}
由于 \(g(E_\Delta)=1\)，对 \(A_6\) 正则覆叠应用 Riemann--Hurwitz：
\[
2g(Y)-2
=\sum_{\text{branch }Q}|A_6|\Bigl(1-\frac1{e_Q}\Bigr).
\]
此处 \(|A_6|=360\)，分歧点数为 \(2\)，且 \(e_Q=3\)，故
\[
2g(Y)-2=2\cdot 360\cdot\frac{2}{3}=480,
\]
从而 \(g(Y)=241\)。证毕。
\end{proof}

\begin{theorem}[超椭圆分支多项式的判别式与坏素数集合]\label{thm:xi-j-sextic-hyperelliptic-branch-polynomial-discriminant-factorization}
\leanverified{paper\_xi\_j\_sextic\_hyperelliptic\_branch\_polynomial\_discriminant\_factorization}
沿用定理 \ref{thm:xi-j-sextic-fiber-product-genus5-hyperelliptic-model} 的 \(r(a)\)、\(H(a)\)，令
\[
P(a):=r(a)H(a)\in\ZZ[a].
\]
则其判别式与结式满足完全显式分解：
\[
\boxed{\ \Disc(r)=2^{6}\cdot 79^{3}\ },
\qquad
\boxed{\ \Disc(H)=-2^{30}\cdot 79^{9}\cdot 89^{10}\cdot 223^{3}\ },
\]
\[
\boxed{\ \Res(r,H)=2^{27}\cdot 79^{8}\ },
\qquad
\boxed{\ \Disc(P)=-2^{90}\cdot 79^{28}\cdot 89^{10}\cdot 223^{3}\ }.
\]
从而曲线 \(C:Y^2=-2^{17}P(a)\) 的坏素数集合精确包含于 \(\{2,79,89,223\}\)；
并且对任意素数 \(p\nmid 2\cdot 79\cdot 89\cdot 223\)，其分支多项式在 \(\FF_p\) 上仍平方自由，故 \(C\) 在 \(p\) 处具有良好约化（以“分支多项式平方自由”意义）。
\end{theorem}
\begin{proof}
四个分解均可由判别式与结式的标准公式在整数环上直接计算并因子分解得到；相应的可复现证书由脚本 \texttt{scripts/exp\_xi\_j\_sextic\_ellipticization\_a6\_hyperelliptic\_audit.py} 给出（见注记 \ref{rem:xi-j-sextic-ellipticization-audit-path}）。
\par
良好约化断言来自一般事实：若 \(P(a)\) 在 \(\FF_p\) 上平方自由，则二次覆盖 \(Y^2=P(a)\) 的分支在模 \(p\) 下保持分离，从而给出光滑约化；平方自由性在 \(p\) 处失败当且仅当 \(p\mid \Disc(P)\)。证毕。
\end{proof}

\begin{remark}[可复现证书路径]\label{rem:xi-j-sextic-ellipticization-audit-path}
本段的判别式/结式分解、椭圆判别式与分歧循环核验由脚本
\texttt{scripts/exp\_xi\_j\_sextic\_ellipticization\_a6\_hyperelliptic\_audit.py}
生成，输出文件为
\texttt{artifacts/export/xi\_j\_sextic\_ellipticization\_a6\_hyperelliptic\_audit.json}。
\end{remark}

\endinput


\begin{proposition}[两处素数模因子型的充分性：\texorpdfstring{$6$}{6}-循环与换位强制 \texorpdfstring{$S_6$}{S6}]\label{prop:xi-j-sextic-two-primes-force-s6}
取 \(t_0\in\ZZ\) 使 \(P_{t_0}(a)\) 可分离。若存在互异素数 \(p,q\nmid \Disc(P_{t_0})\) 满足
\[
P_{t_0}(a)\bmod p\ \text{在}\ \FF_p[a]\ \text{上不可约},\qquad
P_{t_0}(a)\bmod q\ \text{的因子型为}\ (2)(1)(1)(1)(1),
\]
则 \(\Gal(P_{t_0}/\QQ)\cong S_6\)。
\end{proposition}
\leanverified{paper\_xi\_j\_sextic\_two\_primes\_force\_s6}
\begin{proof}
由 Dedekind 定理（不整除判别式时）模素数因子型给出 Frobenius 循环类型；不可约给出 \(6\)-循环，因子型 \((2)(1^4)\) 给出换位。含 \(6\)-循环与换位的传递子群必为 \(S_6\)。证毕。
\end{proof}

\begin{theorem}[多临界 Puiseux 自由能与 \(n^{-r}\) 零点缩放极限]\label{thm:xi-terminal-belyi-puiseux-scaling}
\leanverified{paper\_xi\_terminal\_belyi\_puiseux\_scaling}
设 \(\lambda_1(y),\dots,\lambda_r(y)\) 为 \(\Pi_r(\lambda,y)=0\) 的根（按重数），并定义
\[
Z_n^{(r)}(y):=\sum_{j=1}^{r}\lambda_j(y)^n,\qquad n\in\NN.
\]
则：
\begin{enumerate}
  \item \(Z_n^{(r)}(y)\in\ZZ[y]\)，且 \(\deg Z_n^{(r)}=n\)；并满足阶 \(r\) 线性递推（系数由 \(\Pi_r\) 的 \(\lambda\)-系数给定，且在 \(y\) 中至多一次）。
  \item 在谱半径不并列区域，极限自由能
  \[
  F_r(y):=\lim_{n\to\infty}\frac{1}{n}\log|Z_n^{(r)}(y)|
  \]
  存在且
  \(
  F_r(y)=\log\rho(y)
  \),
  其中 \(\rho(y):=\max_j|\lambda_j(y)|\)。
  \item 在多临界点 \(y=1\) 附近，存在 \(\kappa_r\neq0\)（\(\kappa_r^r=-1\)）与 \(r\) 个分支，使
  \[
  \lambda_j(y)
  =1+\omega_r^j\kappa_r\,(y-1)^{1/r}
  +O\!\bigl(|y-1|^{2/r}\bigr),
  \qquad
  \omega_r=e^{2\pi i/r},
  \]
  从而边缘指数为 \(1/r\)。
  \item 对任意固定 \(t\in\CC\)，取 \(y=1+t^r/n^r\)，则有局部缩放
  \[
  Z_n^{(r)}\!\left(1+\frac{t^r}{n^r}\right)
  =\sum_{k=0}^{r-1}\exp\!\bigl(\omega_r^k\kappa_r t\bigr)+O(n^{-1}),
  \]
  误差对 \(t\) 的紧集一致。
  等价地，经重参数 \(t\mapsto\kappa_r^{-1}t\) 后，极限核为
  \[
  \Phi_r(t):=\sum_{k=0}^{r-1}\exp(\omega_r^k t).
  \]
\end{enumerate}
\end{theorem}
\begin{proof}
第 1 条由 Newton 幂和恒等式直接成立。
第 2 条在主导根唯一时由标准主模分解
\(
Z_n^{(r)}=\lambda_\ast^n(1+o(1))
\)
得到。

对第 3 条，令 \(\mu=\lambda-1,\ \delta=y-1\)，则
\[
\Pi_r(1+\mu,1+\delta)=\mu^r+\delta(1+\mu)^{r-1}=0.
\]
故主导平衡为 \(\mu^r\sim-\delta\)，由 Puiseux 定理得
\(
\mu=\omega_r^j\kappa_r\delta^{1/r}+O(\delta^{2/r})
\)。
于是边缘指数为 \(1/r\)。

第 4 条由上述展开代入 \(\delta=t^r/n^r\) 并对 \(\lambda_j^n\) 取对数展开得到：
\[
\lambda_j^n
=\exp\!\Bigl(\omega_r^j\kappa_r t+O(n^{-1})\Bigr),
\]
求和即得。证毕。
\end{proof}

% AUTO-GENERATED by scripts/exp_xi_belyi_full_symmetric_puiseux_audit.py
\begin{table}[H]
\centering
\scriptsize
\setlength{\tabcolsep}{6pt}
\caption{局部缩放 \(y=1+t^r/n^r\) 下的数值审计：\(Z_n^{(r)}\) 与 \(\sum_{k=0}^{r-1}\exp(\omega_r^k\kappa_r t)\) 的最大误差（\(\kappa_r^r=-1\)）。}
\label{tab:xi_belyi_scaling_limit_audit}
\begin{tabular}{r r l l}
\toprule
$r$ & $n$ & max error & $n\cdot$max error\\
\midrule
3 & 120 & $0.00026738388$ & $0.032086065$\\
3 & 240 & $0.00013369047$ & $0.032085714$\\
4 & 120 & $0.00010725135$ & $0.012870162$\\
4 & 240 & $5.3477609e-5$ & $0.012834626$\\
5 & 120 & $2.1660776e-5$ & $0.0025992932$\\
5 & 240 & $1.074835e-5$ & $0.0025796041$\\
6 & 120 & $2.9283989e-6$ & $0.00035140787$\\
6 & 240 & $1.4430866e-6$ & $0.00034634079$\\
\bottomrule
\end{tabular}
\end{table}


\begin{proposition}[\(r=3\) 的显式 Lee--Yang 区间与三次尖点指数]\label{prop:xi-terminal-belyi-r3-explicit-leyang-interval}
\leanverified{paper\_xi\_terminal\_belyi\_r3\_explicit\_leyang\_interval}
当 \(r=3\) 时，
\[
\Pi_3(\lambda,y)=(\lambda-1)^3+(y-1)\lambda^2
=\lambda^3+(y-4)\lambda^2+3\lambda-1.
\]
记 \(Z_n(y):=\sum_{j=1}^{3}\lambda_j(y)^n\)，则：
\begin{enumerate}
  \item \(Z_n(y)\in\ZZ[y]\)、\(\deg Z_n=n\)，且
  \[
  Z_{n+3}+(y-4)Z_{n+2}+3Z_{n+1}-Z_n=0,
  \]
  并有初值
  \(
  Z_0=3,\ Z_1=4-y,\ Z_2=y^2-8y+10
  \)（见式 \eqref{eq:xi_terminal_belyi_r3_recurrence_discriminant}）。
  \item 判别式分解为
  \[
  \mathrm{Disc}_{\lambda}(\Pi_3)(y)=(y-1)^2(4y-31).
  \]
  因此仅有两处有限分歧值 \(y=1\) 与 \(y=31/4\)。
  \item 对任意实 \(y\in(1,31/4)\)，\(\Pi_3(\cdot,y)\) 恰有一实根 \(\lambda_0(y)\in(0,1)\) 与一对共轭根
  \(
  \lambda_\pm(y)=\rho(y)e^{\pm i\theta(y)}
  \)（\(\rho(y)>1\)），并满足
  \[
  \lambda_0(y)=\rho(y)^{-2}.
  \]
  因而实轴上的谱半径并列区间恰为 \([1,31/4]\)，其 \(y=1\) 端点边缘指数为 \(1/3\)，零点逼近速率为 \(\Theta(n^{-3})\)。
\end{enumerate}
\end{proposition}
\begin{proof}
第 1 条由定理 \ref{thm:xi-terminal-belyi-puiseux-scaling} 的一般幂和递推在 \(r=3\) 的专门化给出。
第 2 条由三次判别式公式直接代入
\(
a=y-4,\ b=3,\ c=-1
\)
得到。

对第 3 条，\(\mathrm{Disc}<0\) 在 \((1,31/4)\) 上成立，故恰有一实根与一对共轭非实根。
又
\[
\Pi_3(0,y)=-1<0,\qquad
\Pi_3(1,y)=y-1>0,
\]
故唯一实根 \(\lambda_0(y)\in(0,1)\)。
由 Vieta，
\(
\lambda_0\lambda_+\lambda_-=1
\)，
写 \(\lambda_\pm=\rho e^{\pm i\theta}\) 得
\(
\lambda_0\rho^2=1
\)，即 \(\lambda_0=\rho^{-2}\) 且 \(\rho>1\)。
因此区间内由共轭对给出谱半径并列；区间外判别式非负，谱半径为单根。
端点指数与缩放率由定理 \ref{thm:xi-terminal-belyi-puiseux-scaling} 的 \(r=3\) 情形立即得到。证毕。
\end{proof}

\begin{corollary}[可调多临界指数、代数端点显式性与最大 Galois 复杂度]\label{cor:xi-terminal-belyi-adjustable-exponent-maximal-galois}
对每个 \(r\ge2\)，族 \(\Pi_r\) 同时满足：
\begin{enumerate}
  \item 多临界点固定于 \(y=1\)，边缘指数严格为 \(1/r\)，局部缩放核由 \(\Phi_r\) 显式给出；
  \item 唯一另一有限端点 \(y_r=1+\frac{r^r}{(r-1)^{r-1}}\) 可由判别式闭式直接读出；
  \item 几何单值群为 \(S_r\)，故在同阶参数族中实现满对称置换复杂度。
\end{enumerate}
\end{corollary}
\begin{proof}
由定理 \ref{thm:xi-terminal-belyi-full-symmetric-cover}、定理 \ref{thm:xi-terminal-belyi-discriminant-closed-form} 与定理 \ref{thm:xi-terminal-belyi-puiseux-scaling} 逐项合并即得。证毕。
\end{proof}

\begin{conjecture}[有限态代数谱配分族的整数 Puiseux 指数原理]\label{conj:xi-terminal-algebraic-spectrum-integer-puiseux-principle}
设 \(\{Z_n(y)\}_{n\ge0}\subset\ZZ[y]\) 来自有限维转移算子或代数谱曲线：
存在 \(\Pi(\lambda,y)\in\ZZ[\lambda,y]\)，使
\[
Z_n(y)=\sum_j c_j(y)\lambda_j(y)^n,
\]
其中 \(\lambda_j\) 为 \(\Pi(\lambda,y)=0\) 的代数分支。
猜想任一自由能非解析点 \(y_c\) 均由主导根的有限重合并（重数 \(m\ge2\)）触发，且边缘指数恒为 \(1/m\)。
进一步，在无额外对称约束的一般位置下，几何单值群应达到满对称群 \(S_{\deg_\lambda\Pi}\)。
若该猜想成立，则临界定位与指数读取可统一归约为 \(\Pi\) 的判别式/子结式链的符号判别与主导重根阶的抽取。
\end{conjecture}

\paragraph{次数 \texorpdfstring{$11$}{11} 的满对称单值化：循环三重归一化与 \texorpdfstring{$S_{11}$}{S11}}
定义二元多项式
\[
E_t(\mu):=
2t^3(\mu^2+1)^2(\mu^2-\mu-1)^2+\mu^5(3\mu^2-2\mu-1)^3\in\ZZ[t,\mu].
\]

\begin{theorem}[一族 \texorpdfstring{$11$}{11} 次特化多项式的全对称 Galois 群]\label{thm:xi-degree11-Et-specialization-galois-S11}
取 \(t=3\)，记 \(f(\mu):=E_3(\mu)\in\ZZ[\mu]\)。
则 \(f\) 在 \(\QQ\) 上不可约，并且
\[
\Gal\bigl(f/\QQ\bigr)\cong S_{11}.
\]
\leanverified{paper\_xi\_degree11\_generic\_galois\_S11\_seeds}
\leanverified{paper\_xi\_degree11\_generic\_galois\_S11\_package}
\leanverified{paper\_xi\_degree11\_et\_specialization\_galois\_s11}
\end{theorem}
\begin{proof}
令 \(n=\deg f=11\)。
\par
\emph{不可约性。}
在素数 \(31\) 下，将 \(f\) 约化并按首项单位化，可得
\[
\overline f(\mu)\equiv \mu^{11}-2\mu^{10}-10\mu^9-5\mu^8-11\mu^7-12\mu^6+4\mu^5+4\mu^3+2\mu^2+4\mu+2\pmod{31}.
\]
该多项式在 \(\FF_{31}[\mu]\) 中不可约（可由有限域因式分解直接核验）。
故由 Gauss 引理与模素不可约判别，\(f\) 在 \(\QQ[\mu]\) 上不可约，因而 \(G:=\Gal(f/\QQ)\subset S_{11}\) 在根集上传递。
\par
\emph{存在 \texorpdfstring{$7$}{7}-循环。}
在素数 \(7\) 下，\(f\) 的约化分解为一次、三次与七次不可约因子之积：
\[
f(\mu)\equiv (\mu-3)\,(\mu^3+\mu^2-\mu+3)\,(\mu^7+2\mu^5+2\mu^3+2\mu-1)\pmod 7.
\]
并且判别式满足
\[
\Disc(f)=2^{22}\cdot 3^{37}\cdot 5^{8}\cdot 23\cdot 2741\cdot 5153\cdot 1315229,
\]
从而 \(7\nmid \Disc(f)\)。
因此 Dedekind 定理给出 Frobenius 置换在根上的循环型与上述分解型一致，
从而 \(G\) 含有循环型 \((7)(3)(1)\) 的元，特别含有一个 \(7\)-循环。
\par
\emph{非交错性。}
\(\Disc(f)\) 非平方数（例如其中素因子 \(23\) 的指数为 \(1\)），故 \(G\nsubseteq A_{11}\)。
\par
\emph{Jordan 逼近。}
由于 \(n=11\) 为素数，传递性蕴含 \(G\) 为本原群。
又 \(G\) 含有 \(p\)-循环（\(p=7\le n-3\)），由 Jordan 定理知 \(A_{11}\subseteq G\)。
结合 \(G\nsubseteq A_{11}\)，得到 \(G=S_{11}\)。证毕。
\end{proof}

\begin{corollary}[泛型 \texorpdfstring{$t$}{t} 的全对称性与不可解性]\label{cor:xi-degree11-Et-generic-galois-S11-hilbert}
令 \(K=\QQ(t)\)，并将 \(E_t(\mu)\) 视作 \(K[\mu]\) 中多项式。
则
\[
\Gal\bigl(E_t(\mu)/\QQ(t)\bigr)\cong S_{11}.
\]
此外，除去一个薄集外，对所有 \(a\in\QQ\) 的可分特化 \(E_a(\mu)\in\QQ[\mu]\) 均满足
\[
\Gal\bigl(E_a(\mu)/\QQ\bigr)\cong S_{11},
\]
从而存在无穷多个整数 \(a\in\ZZ\) 使上述结论成立。
特别地，方程 \(E_t(\mu)=0\) 在 \(\QQ(t)\) 上一般不可由根式求解。
\leanverified{paper\_xi\_degree11\_generic\_galois\_S11\_seeds}
\end{corollary}
\begin{proof}
由定理 \ref{thm:xi-degree11-Et-specialization-galois-S11}，在特化 \(t=3\) 处的伽罗瓦群已为 \(S_{11}\)。
特化同态的基本性质表明：在可分特化点处，特化伽罗瓦群嵌入泛型伽罗瓦群。
因此泛型伽罗瓦群包含 \(S_{11}\)，只能等于 \(S_{11}\)。
\par
Hilbert 不可约定理断言：除去一个薄集外，\(\QQ\) 上的特化保持与泛型相同的伽罗瓦群；
从而得到对几乎所有 \(a\in\QQ\) 的结论，并推出存在无穷多个整数特化仍为 \(S_{11}\)。
最后，由经典伽罗瓦理论，根式可解当且仅当伽罗瓦群可解；而 \(S_{11}\) 不可解，故一般不可由根式求解。证毕。
\end{proof}

\begin{theorem}[判别轨迹的循环三重归一化与属数 \texorpdfstring{$3$}{3}]\label{thm:xi-degree11-cyclic-cubic-normalization-genus3}
设
\[
F(\mu):=-\frac{2\mu(\mu^2-\mu-1)^2}{\mu^2+1}\in\QQ(\mu),
\]
并令函数域扩张
\[
\QQ(\mu)\subset \QQ(\mu,u),\qquad u^3=F(\mu).
\]
则存在唯一（至同构）光滑射影曲线 \(\mathcal C\) 使其函数域为 \(\QQ(\mathcal C)=\QQ(\mu,u)\)。
映射
\[
\pi:\mathcal C\to\PP^1,\qquad (\mu,u)\mapsto \mu
\]
为次数 \(3\) 的循环覆盖，且恰有五个分歧点
\[
\mu=0,\quad \mu=\frac{1\pm\sqrt5}{2},\quad \mu=\pm i,
\]
其分歧指数均为 \(3\)。
因此
\[
g(\mathcal C)=3,
\]
并且 \(\Aut(\mathcal C)\) 含有阶 \(3\) 元 \(u\mapsto \omega u\)（\(\omega^3=1\)）。
\end{theorem}
\begin{proof}
由 \(u^3=F(\mu)\) 可知 \(\pi\) 为 Kummer 型次数 \(3\) 扩张，其 Galois 群由 \(u\mapsto \omega u\) 生成，故为循环覆盖。
分歧由 \(F\) 在 \(\PP^1_\mu\) 上的零极点阶决定：凡零极点阶不被 \(3\) 整除者产生全分歧 \(e=3\)。
\par
由表达式可见：
\(\mu=0\) 为零点阶 \(1\)；
\(\mu=\frac{1\pm\sqrt5}{2}\) 为零点阶 \(2\)；
\(\mu=\pm i\) 为极点阶 \(1\)；
而 \(\mu=\infty\) 处有 \(F(\mu)\sim -2\mu^3\)，零极点阶为 \(3\)，可被 \(3\) 整除，故不分歧。
因此恰有五个全分歧点，各贡献 \(e-1=2\)。
由 Riemann--Hurwitz 公式
\[
2g(\mathcal C)-2=3\cdot(-2)+5\cdot 2
\]
得 \(g(\mathcal C)=3\)。证毕。
\end{proof}

\begin{theorem}[\texorpdfstring{$\mathcal C$}{C} 上自然参数的极点除子与 \texorpdfstring{$\deg(t)=11$}{deg(t)=11}]\label{thm:xi-degree11-parameter-t-divisor-degree}
在定理 \ref{thm:xi-degree11-cyclic-cubic-normalization-genus3} 的 \(\mathcal C\) 上定义亚纯函数
\[
t:=\frac{X(\mu)}{u},
\qquad
X(\mu):=\frac{\mu^2(3\mu^2-2\mu-1)}{\mu^2+1}.
\]
则映射 \(t:\mathcal C\to\PP^1\) 的次数为 \(\deg(t)=11\)，并且其极点除子满足
\[
\div_\infty(t)
=2P_{\alpha_+}+2P_{\alpha_-}+2P_{i}+2P_{-i}+P_{\infty,1}+P_{\infty,2}+P_{\infty,3},
\]
其中 \(\alpha_\pm=\frac{1\pm\sqrt5}{2}\)，\(P_{\alpha_\pm},P_{\pm i}\) 为 \(\pi^{-1}(\alpha_\pm),\pi^{-1}(\pm i)\) 中唯一点（全分歧），而 \(P_{\infty,1},P_{\infty,2},P_{\infty,3}\) 为 \(\pi^{-1}(\infty)\) 中三点（不分歧）。
\end{theorem}
\begin{proof}
在 \(\mu=\alpha_\pm\) 处，\(F(\mu)\sim c(\mu-\alpha_\pm)^2\)。
取局部参数 \(s\) 使 \(\mu-\alpha_\pm=s^3\)，则 \(u\sim c^{1/3}s^2\)，故 \(v_{P_{\alpha_\pm}}(u)=2\)。
又 \(X(\alpha_\pm)\neq 0\)，故 \(v_{P_{\alpha_\pm}}(X)=0\)，从而 \(v_{P_{\alpha_\pm}}(t)=-2\)，给出极点阶 \(2\)。
\par
在 \(\mu=\pm i\) 处，\(\mu^2+1\) 为一阶零，故 \(F\) 为一阶极点并全分歧。
取 \(\mu\mp i=s^3\)，则 \(u\sim c^{1/3}s^{-1}\)，故 \(v_{P_{\pm i}}(u)=-1\)。
另一方面，\(X\) 在 \(\mu=\pm i\) 处为一次极点，在全分歧点上变为三次极点，即 \(v_{P_{\pm i}}(X)=-3\)。
因此 \(v_{P_{\pm i}}(t)=-3-(-1)=-2\)，亦为二阶极点。
\par
在 \(\mu=\infty\) 上方的三点 \(P_{\infty,j}\) 处，\(\pi\) 不分歧且 \(F(\mu)\sim -2\mu^3\)，故 \(u\sim c\mu\)，从而 \(v_{P_{\infty,j}}(u)=-1\)。
又 \(X(\mu)\sim 3\mu^2\)，故 \(v_{P_{\infty,j}}(X)=-2\)，于是 \(v_{P_{\infty,j}}(t)=-1\)，即每点给出一个一次极。
\par
综上得到所述极点除子，其次数为 \(2+2+2+2+1+1+1=11\)，从而 \(\deg(t)=11\)。证毕。
\end{proof}

\begin{corollary}[次数 \texorpdfstring{$11$}{11} 覆盖的满对称算术单值化]\label{cor:xi-degree11-cover-arithmetic-monodromy-S11}
\leanpartial{paper\_xi\_degree11\_cover\_arithmetic\_monodromy\_S11}{The requested Lean signature is false as stated: taking \(G\) to be the trivial group contradicts \(\mathrm{Nonempty}(G \simeq \mathrm{Perm}(\mathrm{Fin}\ 11))\).}
在上述设定下，函数 \(t:\mathcal C\to\PP^1\) 的伽罗瓦闭包之算术单值化群同构于 \(S_{11}\)。
等价地，\(\QQ(\mathcal C)/\QQ(t)\) 的伽罗瓦闭包伽罗瓦群为 \(S_{11}\)。
\end{corollary}
\begin{proof}
由 \(t=X(\mu)/u\) 与 \(u^3=F(\mu)\) 得 \((X(\mu)/t)^3=F(\mu)\)，清分母即得 \(\mu\) 满足方程 \(E_t(\mu)=0\)。
因此 \(\QQ(\mathcal C)=\QQ(t,\mu)\)，并且 \([\QQ(\mathcal C):\QQ(t)]=\deg(t)=11\)（定理 \ref{thm:xi-degree11-parameter-t-divisor-degree}）。
由推论 \ref{cor:xi-degree11-Et-generic-galois-S11-hilbert}，\(\Gal(E_t/\QQ(t))\cong S_{11}\)，从而对应的次数 \(11\) 覆盖之伽罗瓦闭包算术单值化群即为 \(S_{11}\)。证毕。
\end{proof}

\begin{remark}[可复现的有限域因式分解与判别式审计]\label{rem:xi-degree11-Et-audit-script}
定理 \ref{thm:xi-degree11-Et-specialization-galois-S11} 中对模素不可约性、模素分解型与判别式奇偶平方性的核验，可通过符号计算脚本一致复现。
一个实现见 \texttt{scripts/exp\_xi\_degree11\_Et\_S11\_audit.py}。
\end{remark}

\paragraph{平方根谱碰撞与前导零点的 \texorpdfstring{$n^{-2}$}{n^{-2}} 量子化}
令 \(U\subset\CC\) 为开集，\(A:U\to\Mat_{d}(\CC)\) 为解析矩阵族（\(d\ge 2\)）。
取非零线性泛函 \(\ell\in(\CC^d)^\vee\) 与非零向量 \(r\in\CC^d\)，并定义有限尺度配分函数
\[
Z_n(u):=\ell\!\bigl(A(u)^n r\bigr),\qquad n\in\NN.
\]
记特征多项式
\[
P(u,\lambda):=\det(\lambda I-A(u)).
\]

\begin{theorem}[平方根谱碰撞下 Lee--Yang 前导零点的 \texorpdfstring{$n^{-2}$}{n^{-2}} 量子化定标]\label{thm:xi-leyang-square-root-collision-leading-zeros-n2}
\leanverified{paper\_xi\_leyang\_square\_root\_collision\_leading\_zeros\_n2}
设存在 \((u_c,\lambda_c)\in U\times\CC\) 满足：
\begin{enumerate}
  \item \(P(u_c,\lambda_c)=0\) 且 \(\partial_\lambda P(u_c,\lambda_c)=0\)；
  \item \(\partial_{\lambda\lambda}P(u_c,\lambda_c)\neq 0\) 且 \(\partial_u P(u_c,\lambda_c)\neq 0\)；
  \item \(A(u_c)\) 的其余全部特征值 \(\{\lambda_j(u_c)\}_{j\ge 3}\) 满足 \(|\lambda_j(u_c)|<|\lambda_c|\)；
  \item \(\ell,r\) 在该二重特征值对应的二维广义谱簇上诱导的两支谱投影系数不同时为零。
\end{enumerate}
则存在 \(\varepsilon>0\) 及解析函数 \(c_\pm(t)\)（满足 \(c_\pm(0)\neq 0\)），以及定义于 \(|t|<\varepsilon\) 的两支特征值分支 \(\lambda_\pm(t)\)，使得在二重覆盖参数化
\[
u=u_c+t^2
\]
下成立展开
\[
\lambda_\pm(t)=\lambda_c\pm \alpha\,t+O(t^2),
\qquad
\alpha^2=-\frac{2\,\partial_u P(u_c,\lambda_c)}{\partial_{\lambda\lambda}P(u_c,\lambda_c)}\neq 0,
\]
并且存在常数 \(\eta\in(0,|\lambda_c|)\) 使对 \(|t|<\varepsilon\) 一致有
\[
Z_n(u_c+t^2)=c_+(t)\lambda_+(t)^n+c_-(t)\lambda_-(t)^n+O(\eta^n)
\qquad(n\to\infty).
\]
令
\[
\kappa:=-\frac{c_-(0)}{c_+(0)}\in\CC^\times,
\]
并固定一条对数分支 \(\log\) 在某邻域内解析。
则存在常数 \(C>0\) 使得对每个固定整数 \(k\in\ZZ\)，当 \(n\) 足够大时，方程 \(Z_n(u)=0\) 在缩放窗 \(|u-u_c|\le Cn^{-2}\) 内存在唯一零点 \(u_{n,k}\)，且满足
\[
u_{n,k}
=
u_c+\frac{\lambda_c^{2}}{4\alpha^{2}n^{2}}
\bigl(\log\kappa+2\pi i k\bigr)^{2}
+O(n^{-3}),
\qquad(n\to\infty).
\]
\end{theorem}
\begin{proof}
由 (1)(2) 与 Weierstrass 制备定理，\(P(u,\lambda)\) 在 \((u_c,\lambda_c)\) 的邻域内可分解为
\[
P(u,\lambda)=\Bigl((\lambda-\lambda_c)^2+a_1(u)(\lambda-\lambda_c)+a_0(u)\Bigr)\,Q(u,\lambda),
\qquad Q(u_c,\lambda_c)\neq 0,
\]
其中 \(a_0(u_c)=a_1(u_c)=0\)，且 \(a_0'(u_c)\neq 0\)（横截性）。
因此存在二重覆盖 \(u=u_c+t^2\)，使得该二次因子的两根成为解析分支 \(\lambda_\pm(t)\)，并由二次项比较得到
\[
(\lambda-\lambda_c)^2
=
-\frac{2\,\partial_u P(u_c,\lambda_c)}{\partial_{\lambda\lambda}P(u_c,\lambda_c)}(u-u_c)
+O\bigl(|u-u_c|^{3/2}\bigr),
\]
从而 \(\lambda_\pm(t)=\lambda_c\pm\alpha t+O(t^2)\) 且 \(\alpha^2\) 如述。

由解析谱投影与 (3) 的严格谱隙，可将 \(A(u_c+t^2)^n\) 分解为两主支谱投影项与余项：
\[
A(u_c+t^2)^n=\lambda_+(t)^n\Pi_+(t)+\lambda_-(t)^n\Pi_-(t)+R_n(t),
\]
其中 \(\Pi_\pm(t)\) 解析，且 \(\|R_n(t)\|=O(\eta^n)\) 对 \(|t|<\varepsilon\) 一致成立。
左右夹以 \(\ell(\cdot r)\) 得
\(
Z_n(u_c+t^2)=c_+(t)\lambda_+(t)^n+c_-(t)\lambda_-(t)^n+O(\eta^n)
\)，其中 \(c_\pm(t):=\ell(\Pi_\pm(t)r)\)；
由 (4) 得 \(c_\pm(0)\neq 0\)。

令 \(Z_n(u_c+t^2)=0\)。
将余项并入误差后，有
\[
\Bigl(\frac{\lambda_+(t)}{\lambda_-(t)}\Bigr)^n=-\frac{c_-(t)}{c_+(t)}+O\bigl((\eta/|\lambda_c|)^n\bigr).
\]
在 \(|t|<\varepsilon\) 内取对数并吸收分支不定性 \(2\pi i k\) 得
\[
n\log\frac{\lambda_+(t)}{\lambda_-(t)}
=
\log\kappa+2\pi i k+O(t)+O\bigl((\eta/|\lambda_c|)^n\bigr).
\]
由 \(\lambda_\pm(t)=\lambda_c\pm\alpha t+O(t^2)\) 得
\(
\log\frac{\lambda_+(t)}{\lambda_-(t)}=\frac{2\alpha}{\lambda_c}t+O(t^2)
\)。
因此
\[
t=\frac{\lambda_c}{2\alpha n}\bigl(\log\kappa+2\pi i k\bigr)+O(n^{-2}),
\]
进而
\(
u-u_c=t^2=\frac{\lambda_c^2}{4\alpha^2n^2}(\log\kappa+2\pi i k)^2+O(n^{-3})
\)。
唯一性可由上述对数方程在 \(|t|\le c/n\) 区域的收缩估计得到。证毕。
\end{proof}

\begin{corollary}[对称可见性下的两零点外推与幅度消去]\label{cor:xi-leyang-two-leading-zeros-extrapolate-uc}
在定理 \ref{thm:xi-leyang-square-root-collision-leading-zeros-n2} 的假设下，若进一步有 \(c_+(0)=c_-(0)\)，则 \(\kappa=-1\)。
取对数分支满足 \(\log(-1)=i\pi\)，并令
\[
C:=\frac{\pi^2\lambda_c^2}{4\alpha^2}\neq 0.
\]
则当 \(n\to\infty\) 时，靠近 \(u_c\) 的两枚最小尺度零点（对应 \(k=0,1\)）满足
\[
u_{n,0}=u_c-\frac{C}{n^2}+O(n^{-3}),
\qquad
u_{n,1}=u_c-\frac{9C}{n^2}+O(n^{-3}).
\]
从而两零点外推
\[
\widehat u_c(n):=\frac{9u_{n,0}-u_{n,1}}{8}
=u_c+O(n^{-3}),
\]
并且幅度可由
\[
\widehat C(n):=\frac{n^2}{8}\bigl(u_{n,0}-u_{n,1}\bigr)=C+O(n^{-1})
\]
一致估计。
\par
同样地，幅度亦可由两零点外推出的消元式直接读出：
\[
\widehat C_{\mathrm{elim}}(n):=-n^2\bigl(u_{n,0}-\widehat u_c(n)\bigr)=C+O(n^{-1}).
\]
\leanverified{paper\_xi\_leyang\_two\_leading\_zeros\_extrapolate\_uc}
\end{corollary}
\begin{proof}
由定理 \ref{thm:xi-leyang-square-root-collision-leading-zeros-n2} 的量子化公式代入 \(\log\kappa=\log(-1)=i\pi\)，得到
\[
u_{n,k}
=u_c+\frac{\lambda_c^2}{4\alpha^2n^2}\bigl(i\pi+2\pi i k\bigr)^2+O(n^{-3})
=u_c-\frac{\pi^2\lambda_c^2}{4\alpha^2n^2}(2k+1)^2+O(n^{-3}).
\]
取 \(k=0,1\) 即得前两式；线性消元得到 \(\widehat u_c(n)\) 的表达式，\(\widehat C(n)\) 由差分读出。证毕。
\end{proof}

\begin{theorem}[三零点端点外推的 \texorpdfstring{$O(n^{-6})$}{O(n^-6)} 误差改进]\label{thm:xi-leyang-three-leading-zeros-extrapolate-uc-n6}
\leanverified{paper\_xi\_leyang\_three\_leading\_zeros\_extrapolate\_uc\_n6}
设在某端点机制下，存在常数 \(u_c\in\CC\)、\(C\neq 0\)、\(D\in\CC\)，使得对固定的 \(k\in\{0,1,2\}\) 有渐近展开
\[
u_{n,k}
=
u_c
-\frac{C}{n^{2}}(2k+1)^{2}
-\frac{D}{n^{4}}(2k+1)^{4}
O(n^{-6})
\qquad(n\to\infty).
\]
定义三零点外推量
\[
\boxed{\
\widehat u_{c}^{(n)}
:=\frac{75}{64}\,u_{n,0}-\frac{25}{128}\,u_{n,1}+\frac{3}{128}\,u_{n,2}.
}
\]
则成立严格误差改进
\[
\boxed{\ \widehat u_{c}^{(n)}=u_{c}+O(n^{-6}).\ }
\]
\end{theorem}
\begin{proof}
令 \(a_k:=(2k+1)^2\)，则 \(a_0=1,a_1=9,a_2=25\)。
取系数 \(\alpha_0,\alpha_1,\alpha_2\) 满足
\[
\alpha_0+\alpha_1+\alpha_2=1,\qquad
\alpha_0a_0+\alpha_1a_1+\alpha_2a_2=0,\qquad
\alpha_0a_0^2+\alpha_1a_1^2+\alpha_2a_2^2=0.
\]
该 Vandermonde 线性系统唯一解为
\(
(\alpha_0,\alpha_1,\alpha_2)=(75/64,\,-25/128,\,3/128)
\)。
将展开式代入 \(\sum_{k=0}^2\alpha_k u_{n,k}\)，由约束恰消去 \(n^{-2}\) 与 \(n^{-4}\) 两阶项，余项阶为 \(O(n^{-6})\)。证毕。
\end{proof}

\begin{theorem}[\texorpdfstring{$(M+1)$}{M+1} 点端点外推的闭式权重与最优误差阶]\label{thm:xi-leyang-mplus1-point-extrapolate-uc-optimal}
\leanverified{paper\_xi\_leyang\_mplus1\_point\_extrapolate\_uc\_optimal}
在与定理 \ref{thm:xi-leyang-three-leading-zeros-extrapolate-uc-n6} 同型的端点机制下，固定整数 \(M\ge 1\)，并令
\[
a_k:=(2k+1)^2\qquad (k=0,1,\dots,M).
\]
定义 Lagrange 权重
\[
\boxed{\
\alpha_{k}^{(M)}
:=
\prod_{\substack{0\le i\le M\\ i\ne k}}
\frac{-a_{i}}{a_{k}-a_{i}}
\qquad(k=0,\dots,M),
}
\]
以及外推估计量
\[
\boxed{\
\widehat u_{c,M}^{(n)}:=\sum_{k=0}^{M}\alpha_{k}^{(M)}\,u_{n,k}.
}
\]
若存在常数 \(C_1,\dots,C_{M+1}\in\CC\) 使得对每个固定 \(k\in\{0,1,\dots,M\}\) 有
\[
u_{n,k}
=
u_c+\sum_{j=1}^{M+1}\frac{C_j}{n^{2j}}\,a_k^{j}+O(n^{-2M-4}),
\qquad(n\to\infty),
\]
则
\[
\boxed{\ \widehat u_{c,M}^{(n)}=u_c+O(n^{-2M-2}).\ }
\]
并且在“仅使用 \(\{u_{n,0},\dots,u_{n,M}\}\) 并要求消去前 \(M\) 个偶阶项”的信息模型下，上述权重是唯一解，从而达到该模型内的最优误差阶。
\end{theorem}
\begin{proof}
在假设展开中，截断到 \(n^{-2M-2}\) 的部分在变量 \(a_k\) 上是次数至多 \(M+1\) 的多项式，并且 \(u_c\) 对应于 \(a=0\) 处的取值。
因此 \(\widehat u_{c,M}^{(n)}\) 等价于对节点 \(\{a_0,\dots,a_M\}\) 的 Lagrange 插值多项式在 \(a=0\) 的取值，其权重即为 \(L_k(0)\)，从而给出闭式 \(\alpha_k^{(M)}\)。
插值误差由最高项 \(a^{M+1}\) 控制，对应尺度 \(n^{-2M-2}\)，得到误差阶。
唯一性来自 Vandermonde 矩阵非退化。证毕。
\end{proof}

\begin{corollary}[由三零点同时恢复主标度常数与次阶常数]\label{cor:xi-leyang-recover-C-D-from-three-leading-zeros}
\leanverified{paper\_xi\_leyang\_recover\_C\_D\_from\_three\_leading\_zeros}
在定理 \ref{thm:xi-leyang-three-leading-zeros-extrapolate-uc-n6} 的假设下，令 \(\widehat u_c^{(n)}\) 为其三零点外推量，并定义
\[
B_k:=n^{2}\bigl(\widehat u_{c}^{(n)}-u_{n,k}\bigr)\qquad(k=0,1).
\]
则常数 \(C,D\) 可由显式公式一致估计：
\[
\boxed{\
\widehat C^{(n)}:=\frac{81B_{0}-B_{1}}{72}
\ =\ C+O(n^{-4}),\ }
\qquad
\boxed{\
\widehat D^{(n)}:=\frac{n^{2}(B_{1}-9B_{0})}{72}
\ =\ D+O(n^{-2}).\ }
\]
\end{corollary}
\begin{proof}
由定理 \ref{thm:xi-leyang-three-leading-zeros-extrapolate-uc-n6}，
\(
\widehat u_c^{(n)}=u_c+O(n^{-6})
\)。
结合 \(a_0=1,a_1=9\) 与展开式可得
\[
B_k=C\,a_k+\frac{D}{n^2}\,a_k^2+O(n^{-4})\qquad(k=0,1).
\]
对 \(k=0,1\) 作线性消元：\(\widehat C^{(n)}\) 消去 \(D/n^2\) 项而达到 \(O(n^{-4})\)，\(\widehat D^{(n)}\) 由差分提取 \(D\) 并得到 \(O(n^{-2})\)。证毕。
\end{proof}

\begin{proposition}[平方根分歧点的等模曲线切线方向]\label{prop:xi-leyang-equimodular-tangent-square-root}
在定理 \ref{thm:xi-leyang-square-root-collision-leading-zeros-n2} 的记号下，定义局部等模集合
\[
\Sigma:=\{u\in U:\ |\lambda_+(u)|=|\lambda_-(u)|\},
\]
其中 \(\lambda_\pm\) 取定理中的两支 Puiseux 解析延拓。
则在足够小邻域内，\(\Sigma\) 由恰好两条穿过 \(u_c\) 的实解析弧组成；
其在 \(u\)-平面中穿过 \(u_c\) 的切线方向唯一确定为
\[
\arg(u-u_c)\equiv \pi-2\arg\bigl(\lambda_c\overline{\alpha}\bigr)\pmod{\pi}.
\]
更精确地，在均匀化坐标 \(u=u_c+t^2\) 下，等模条件的一阶近似等价于
\[
\Re\bigl(\lambda_c\overline{\alpha}\,t\bigr)=0,
\]
从而在 \(t\)-平面给出两条互逆直线，其经 \(t\mapsto t^2\) 推送为上述两条切线方向。
\end{proposition}
\begin{proof}
由 \(\lambda_\pm(t)=\lambda_c\pm\alpha t+O(t^2)\) 得
\[
|\lambda_+(t)|^2-|\lambda_-(t)|^2
=|\lambda_c+\alpha t|^2-|\lambda_c-\alpha t|^2+O(|t|^2)
=4\,\Re(\lambda_c\overline{\alpha}\,t)+O(|t|^2).
\]
因此 \(|\lambda_+|=|\lambda_-|\) 在一阶上等价于 \(\Re(\lambda_c\overline{\alpha}\,t)=0\)，给出 \(t\)-平面中过原点的两条直线。
由 \(u-u_c=t^2\) 切向角加倍，得到所示同余式；高阶项由实解析隐函数定理给出两条实解析弧。证毕。
\end{proof}

\begin{theorem}[多项式传递矩阵的临界点代数性与判别式复杂度上界]\label{thm:xi-polynomial-transfer-matrix-branchpoint-discriminant-degree}
设 \(A(u)\in\Mat_d(\ZZ[u])\) 为多项式矩阵，且每个矩阵元 \(A_{ij}(u)\) 的次数不超过 \(k\ge 1\)。
令
\[
P(u,\lambda):=\det(\lambda I-A(u))\in\ZZ[u,\lambda],
\qquad
D(u):=\mathrm{Disc}_\lambda P(u,\lambda)\in\ZZ[u].
\]
则成立：
\begin{enumerate}
  \item \(D(u)\) 非恒零，且
  \[
  \deg_u D\le (2d-1)\deg_u P\le (2d-1)dk.
  \]
  \item 任一谱分歧点 \(u_c\)（即存在 \(\lambda\) 使 \(P(u_c,\lambda)=\partial_\lambda P(u_c,\lambda)=0\)）必为代数数，且
  \[
  [\QQ(u_c):\QQ]\le \deg_u D \le (2d-1)dk.
  \]
  \item 若某分歧点 \(u_c\) 为 \(D\) 的单根（\(D(u_c)=0\) 且 \(D'(u_c)\neq 0\)），并且在该点存在严格谱隙 \(|\lambda_j(u_c)|<|\lambda_c|\)（除去发生碰撞的二重根 \(\lambda_c\)），则定理 \ref{thm:xi-leyang-square-root-collision-leading-zeros-n2} 的平方根碰撞情形在该点为代数可审计：可由判别式单根性与谱隙验证推出 \(n^{-2}\) 量子化定标。
\end{enumerate}
\end{theorem}
\begin{proof}
由 \(D(u)=\mathrm{Disc}_\lambda(P)=(-1)^{d(d-1)/2}\mathrm{Res}_\lambda(P,\partial_\lambda P)\)。
对任意 \(F,G\in\ZZ[u,\lambda]\)，其关于 \(\lambda\) 的结式满足次数估计
\[
\deg_u\mathrm{Res}_\lambda(F,G)
\le (\deg_\lambda F)(\deg_u G)+(\deg_\lambda G)(\deg_u F).
\]
取 \(F=P\)、\(G=\partial_\lambda P\)，则 \(\deg_\lambda P=d\)、\(\deg_\lambda(\partial_\lambda P)=d-1\)，且 \(\deg_u(\partial_\lambda P)\le \deg_u P\)，从而
\[
\deg_u D\le d\,\deg_u P+(d-1)\deg_u P=(2d-1)\deg_u P.
\]
又 \(P(u,\lambda)=\det(\lambda I-A(u))\) 的每一项为 \(d\) 个矩阵元之积，故 \(\deg_u P\le dk\)，得第 1 条。

第 2 条由分歧点满足 \(D(u_c)=0\) 得出：\(u_c\) 为整系数多项式 \(D\) 的零点，因而为代数数且次数不超过 \(\deg_u D\)。

对第 3 条，\(D'(u_c)\neq 0\) 等价于 \(P(u,\lambda)\) 在 \((u_c,\lambda_c)\) 发生单纯二重根碰撞（局部仅两根以平方根方式分歧），与定理 \ref{thm:xi-leyang-square-root-collision-leading-zeros-n2} 的 (1)(2) 同型；
严格谱隙与可见性条件给出两主支主导，从而可应用定理 \ref{thm:xi-leyang-square-root-collision-leading-zeros-n2} 得到 \(n^{-2}\) 量子化。证毕。
\end{proof}

\begin{conjecture}[平方根端点下的两零点比值刚性与超收敛外推]\label{conj:xi-leyang-square-root-endpoint-two-zero-ratio}
设一族有限尺度配分函数 \(Z_n(u)\) 来自有限维解析传递算子或传递矩阵（例如 \(Z_n(u)=\ell(A(u)^n r)\)），并且其自由能的主要解析延拓障碍由某临界端点 \(u_c\) 处两条主导特征值的单纯平方根碰撞控制（判别式在 \(u_c\) 为单根，且存在严格谱隙与可见性）。
则存在编号约定使得靠近 \(u_c\) 的前两枚零点满足
\[
u_{n,1}-u_c=9\,(u_{n,0}-u_c)+o(n^{-2}),
\]
并且由两零点组合
\(
\widehat u_c(n)=(9u_{n,0}-u_{n,1})/8
\)
给出的端点外推误差为 \(o(n^{-2})\)。
若该刚性在数值数据中成立，则可据此反推临界端点处主导谱碰撞为单纯平方根类型，并据判别式与谱隙对端点作代数定位。
\end{conjecture}

\begin{theorem}[\texorpdfstring{$k$}{k}-重主导谱碰撞的 Mittag--Leffler 普遍临界核]\label{thm:xi-mittag-leffler-universal-kfold-jordan-collision}
\leanverified{paper\_xi\_mittag\_leffler\_universal\_kfold\_jordan\_collision}
令 \(M(t)\in M_d(\CC)\) 为 \(t\) 在 \(0\) 的邻域内解析的矩阵族。
设 \(M(0)\) 的谱满足：
\begin{enumerate}
  \item 存在 \(\eta>0\) 使得除特征值 \(1\) 外，其余全部特征值 \(\lambda\) 满足 \(|\lambda|\le 1-\eta\)；
  \item 特征值 \(1\) 的代数重数为 \(k\)，几何重数为 \(1\)（对应单 Jordan 块）；
  \item 在适当选择的主导 Jordan 基下，限制到主导广义特征子空间的特征多项式满足一参数非退化展开
  \[
  \chi(\lambda,t)
  =(\lambda-1)^k-\mathbf c\,t
  +O\!\left(t^2+t(\lambda-1)+(\lambda-1)^{k+1}\right),
  \qquad
  \mathbf c\neq 0.
  \]
\end{enumerate}
取向量 \(u,v\in\CC^d\) 使得其在主导 Jordan 链上的最高阶配对不为零（等价地：在主导子空间标准形中，双线性型 \(u^\top(\cdot)\,v\) 不消去 \((1,k)\) 矩阵元）。
定义标量序列
\[
Z_n(t):=u^\top M(t)^{\,n+k-1}v,\qquad n\ge 1.
\]
则存在常数 \(\mathbf C\neq 0\)（由 \(u,v\) 与主导 Jordan 链的配对给出）使得对任意紧集 \(K\subset\CC\) 有一致收敛
\[
\boxed{\;
n^{-(k-1)}\,Z_n\!\left(-\frac{s}{n^k}\right)\ \longrightarrow\
\mathbf C\,E_{k,k}\!\left(-\mathbf c\,s\right)
\qquad (n\to\infty,\ s\in K)\; }.
\]
其中二参数 Mittag--Leffler 函数定义为
\[
E_{k,k}(z):=\sum_{r\ge 0}\frac{z^r}{\Gamma(kr+k)}.
\]
\end{theorem}
\begin{proof}
由谱隙假设，取围绕主导谱簇的 Riesz 投影 \(P(t)\)，则存在常数 \(C>0\) 与 \(0<\vartheta<1\) 使
\[
M(t)^n=\bigl(P(t)M(t)P(t)\bigr)^n+R_n(t),
\qquad
\|R_n(t)\|\le C\,\vartheta^n
\]
在 \(t\) 的充分小邻域内一致成立。
故在 \(n^{-(k-1)}\) 的尺度下，\(R_n(t)\) 的贡献指数小，可忽略。
\par
记 \(A(t):=P(t)M(t)P(t)\) 并把其视为作用在 \(k\)-维主导广义特征子空间上的算子。
由单 Jordan 块的一参数解析变形，可取解析可逆变换 \(S(t)\) 使
\[
S(t)^{-1}A(t)S(t)=J_k+\tau(t)E_{k1},
\qquad
\tau(t)=\mathbf c\,t+O(t^2),
\]
其中 \(J_k\) 为特征值 \(1\) 的 \(k\times k\) Jordan 块，\(E_{k1}\) 为仅 \((k,1)\) 位置为 \(1\) 的矩阵。
并且由边界配对的非消去假设，存在 \(\mathbf C\neq 0\) 使得
\[
Z_n(t)=\mathbf C\cdot \bigl(J_k+\tau(t)E_{k1}\bigr)^{n+k-1}_{1k}
 +O(n^{k-2})+O(\vartheta^n)
\]
在 \(t\) 的充分小邻域内成立。
\par
对矩阵族 \(J_k+\tau E_{k1}\)，其 \((1,k)\) 元满足阶为 \(k\) 的常系数线性递推，其特征多项式为 \((\lambda-1)^k-\tau\)。
由此得到临界缩放 \(t=-s/n^k\) 下的统一极限核为 Mittag--Leffler 函数：
\[
n^{-(k-1)}\bigl(J_k+\tau(t)E_{k1}\bigr)^{n+k-1}_{1k}
\longrightarrow
E_{k,k}\!\left(-\mathbf c\,s\right),
\]
且误差对 \(s\) 的紧集一致。
合并上述各步即得结论。证毕。
\end{proof}

\begin{corollary}[\texorpdfstring{$k$}{k}-阶缩放下零点的普适量子化]\label{cor:xi-mittag-leffler-kfold-zero-quantization}
\leanverified{paper\_xi\_mittag\_leffler\_kfold\_zero\_quantization}
在定理 \ref{thm:xi-mittag-leffler-universal-kfold-jordan-collision} 的条件下，
设 \(t_{n,j}\) 为 \(Z_n(t)\) 的零点序列且满足 \(n^k t_{n,j}\to -s_j\)。
则必有
\[
E_{k,k}(-\mathbf c\,s_j)=0.
\]
反之，若 \(s\) 是函数 \(E_{k,k}(-\mathbf c\,\cdot)\) 的简单零点，则存在零点族 \(t_n(s)\) 使
\[
t_n(s)=-\frac{s}{n^k}+o(n^{-k}),
\qquad
Z_n\bigl(t_n(s)\bigr)=0.
\]
\end{corollary}
\begin{proof}
定理 \ref{thm:xi-mittag-leffler-universal-kfold-jordan-collision} 给出
\(
f_n(s):=n^{-(k-1)}Z_n(-s/n^k)
\)
在紧集上一致收敛到 \(\mathbf C\,E_{k,k}(-\mathbf c\,s)\)（\(\mathbf C\neq 0\)）。
由 Hurwitz 定理，任何零点极限必须落在极限函数的零点集合上；而简单零点在一致收敛下稳定延拓，给出所述反向存在性。证毕。
\end{proof}

\endinput


\paragraph{有限阶累积量的代数零点反演：Hankel 铅笔、收敛阶与稳定性}
本段把 Lee--Yang 零点的反演组织为一个纯代数问题：给定 \(\log Z\) 在原点的有限阶导数（等价的累积量/对数导数泰勒系数），以 Hankel 铅笔的谱读出最近零点族，并给出可证的指数收敛、锐性与扰动稳定性判据（参见 \cite{LeeYang1952LatticeGasIsing,FlindtGarrahan2013Trajectory} 的背景脉络）。

\begin{definition}[对数累积量与规范化幂和]\label{def:xi-leyang-cumulant-power-sums}
设 \(Z(z)\in\RR[z]\) 为多项式且 \(Z(0)\neq 0\)。记其非零零点（按重数计）为 \(\{\zeta_j\}_{j=1}^N\subset\CC\setminus\{0\}\)，并定义对数累积量
\[
\kappa_n:=\left.\frac{\dd^n}{\dd z^n}\log Z(z)\right|_{z=0}\qquad(n\ge 1).
\]
定义规范化幂和
\begin{equation}\label{eq:xi-leyang-power-sum-sn}
s_n:=-\frac{\kappa_n}{(n-1)!}=\sum_{j=1}^N \zeta_j^{-n}\qquad(n\ge 1),
\end{equation}
并记 \(\xi_j:=\zeta_j^{-1}\)。
\end{definition}

\begin{remark}[零点--累积量幂和恒等式]\label{rem:xi-leyang-zero-cumulant-identity}
在上述设定下有因子分解
\[
Z(z)=Z(0)\prod_{j=1}^N\left(1-\frac{z}{\zeta_j}\right),
\]
从而在 \(z=0\) 的邻域
\[
\log Z(z)=\log Z(0)-\sum_{n\ge 1}\frac{z^n}{n}\sum_{j=1}^N\zeta_j^{-n}.
\]
逐项求导并在 \(z=0\) 取值给出
\(
\kappa_n=-(n-1)!\sum_j\zeta_j^{-n}
\)，即 \eqref{eq:xi-leyang-power-sum-sn}。
\end{remark}

\begin{theorem}[支配零点的 Hankel 铅笔谱逼近与指数收敛]\label{thm:xi-leyang-hankel-pencil-dominant-zeros}
在定义 \ref{def:xi-leyang-cumulant-power-sums} 的设定下，将零点按模长排序 \(|\zeta_1|\le\cdots\le|\zeta_N|\)，并取整数 \(r\ge 1\) 使
\[
|\zeta_r|<|\zeta_{r+1}|\qquad (|\zeta_{N+1}|:=+\infty).
\]
令 \(\xi_j=\zeta_j^{-1}\)，则
\(
|\xi_1|\ge\cdots\ge|\xi_r|>|\xi_{r+1}|
\)。
对任意 \(n\ge 1\) 与固定 \(r\) 定义 \(r\times r\) Hankel 矩阵
\[
H_0^{(r)}(n):=\bigl[s_{n+i+j}\bigr]_{i,j=0}^{r-1},
\qquad
H_1^{(r)}(n):=\bigl[s_{n+i+j+1}\bigr]_{i,j=0}^{r-1},
\]
并在 \(H_0^{(r)}(n)\) 可逆时定义 Hankel 铅笔
\[
\mathcal{S}_r(n):=\bigl(H_0^{(r)}(n)\bigr)^{-1}H_1^{(r)}(n).
\]
记间隙比
\(
q:=|\xi_{r+1}|/|\xi_r|\in(0,1)
\)。
则存在常数 \(C>0\) 与 \(n_0\in\NN\)（仅依赖于 \(\{\xi_j\}_{j=1}^{r+1}\)）使得对所有 \(n\ge n_0\)：
\begin{enumerate}
  \item \(H_0^{(r)}(n)\) 可逆；
  \item 若以 \(\{\widehat\xi_j^{(r)}(n)\}_{j=1}^r\) 表示 \(\mathcal{S}_r(n)\) 的特征值（计代数重数），则可重排使
  \begin{equation}\label{eq:xi-leyang-hankel-pencil-eig-error}
  \max_{1\le j\le r}\abs{\widehat\xi_j^{(r)}(n)-\xi_j}\ \le\ C\,|\xi_1|\,q^{\,n}.
  \end{equation}
\end{enumerate}
因此 \(\widehat\zeta_j^{(r)}(n):=1/\widehat\xi_j^{(r)}(n)\) 给出最近 \(r\) 个零点的指数级逼近，其收敛率由模长间隙 \(q\) 精确控制。
\end{theorem}
\begin{proof}
令 Vandermonde 型矩阵 \(V_r\in\CC^{r\times r}\) 与对角矩阵 \(D_n,\Lambda\) 为
\[
(V_r)_{i,j}:=\xi_j^{\,i}\quad(0\le i\le r-1,\,1\le j\le r),\qquad
D_n:=\mathrm{diag}(\xi_1^n,\dots,\xi_r^n),\qquad
\Lambda:=\mathrm{diag}(\xi_1,\dots,\xi_r).
\]
按定义 \(s_m=\sum_{j=1}^N\xi_j^m\)，可分解为
\[
H_0^{(r)}(n)=V_rD_nV_r^{\mathsf T}+E_0(n),\qquad
H_1^{(r)}(n)=V_rD_n\Lambda V_r^{\mathsf T}+E_1(n),
\]
其中误差矩阵 \(E_0(n),E_1(n)\) 仅由 \(j\ge r+1\) 的项贡献。由 \(|\xi_j|^{n+i+j}\le |\xi_{r+1}|^{n}|\xi_{r+1}|^{i+j}\) 得
\[
\norm{E_\ell(n)}\le C_\ell |\xi_{r+1}|^{\,n}\qquad(\ell=0,1),
\]
其中 \(C_\ell\) 与 \(n\) 无关。
另一方面，\(V_r\) 可逆且 \(D_n\) 可逆，故 \(V_rD_nV_r^{\mathsf T}\) 可逆，且存在常数 \(C'>0\) 使
\[
\norm{(V_rD_nV_r^{\mathsf T})^{-1}}\le C'|\xi_r|^{-n}.
\]
因此
\[
\norm{(V_rD_nV_r^{\mathsf T})^{-1}E_0(n)}\le C''\left(\frac{|\xi_{r+1}|}{|\xi_r|}\right)^{n}=C''q^{\,n}.
\]
取 \(n_0\) 充分大使上式 \(\le \tfrac12\)，则 \(H_0^{(r)}(n)\) 可逆且
\(
\norm{(H_0^{(r)}(n))^{-1}}\le 2\norm{(V_rD_nV_r^{\mathsf T})^{-1}}
\)。
进一步，
\[
\mathcal{S}_r(n)
=(V_rD_nV_r^{\mathsf T})^{-1}(V_rD_n\Lambda V_r^{\mathsf T})
O(q^{\,n})
=V_r^{-{\mathsf T}}\Lambda V_r^{\mathsf T}+O(q^{\,n}).
\]
其中 \(V_r^{-{\mathsf T}}\Lambda V_r^{\mathsf T}\) 与 \(\Lambda\) 相似，特征值恰为 \(\{\xi_1,\dots,\xi_r\}\)。
对矩阵扰动应用 Bauer--Fike 型特征值扰动界即可得到 \eqref{eq:xi-leyang-hankel-pencil-eig-error}。证毕。
\end{proof}

\begin{corollary}[最近共轭对的二阶 Hankel 判别式估计器]\label{cor:xi-leyang-nearest-conjugate-pair-discriminant-estimator}
在定理 \ref{thm:xi-leyang-hankel-pencil-dominant-zeros} 的设定下取 \(r=2\)，并进一步假设最近两零点构成非实共轭对
\[
|\zeta_1|=|\zeta_2|=:R_1<R_2:=\min_{j\ge 3}|\zeta_j|,\qquad \zeta_2=\overline{\zeta_1},\qquad \zeta_1\notin\RR.
\]
对 \(n\ge 1\) 定义二阶 Hankel 判别式
\[
\Delta_n:=\det\begin{pmatrix}
s_n & s_{n+1}\\
s_{n+1} & s_{n+2}
\end{pmatrix}
=s_ns_{n+2}-s_{n+1}^2,
\]
并在 \(\Delta_n\neq 0\) 时定义
\[
A_n:=\frac{s_{n+3}s_n-s_{n+2}s_{n+1}}{\Delta_n},\qquad
B_n:=\frac{s_{n+3}s_{n+1}-s_{n+2}^2}{\Delta_n}.
\]
令 \(\widehat\xi_\pm^{(n)}\) 为二次多项式 \(q_n(\lambda)=\lambda^2-A_n\lambda+B_n\) 的两根，并令 \(\widehat\zeta_\pm^{(n)}:=1/\widehat\xi_\pm^{(n)}\)。
则存在常数 \(C>0\) 与 \(n_0\)（仅依赖于 \(\{\zeta_j\}\)）使对所有 \(n\ge n_0\) 成立：在适当重排下
\[
\max_{\sigma\in\{+,-\}}\abs{\widehat\zeta_\sigma^{(n)}-\zeta_\sigma}
\ \le\ C\,R_1\left(\frac{R_1}{R_2}\right)^{n}.
\]
\end{corollary}
\begin{proof}
当 \(r=2\) 时，\(\mathcal{S}_2(n)=(H_0^{(2)}(n))^{-1}H_1^{(2)}(n)\) 的特征多项式恰为
\(
\lambda^2-A_n\lambda+B_n
\)，其根为 \(\widehat\xi_\pm^{(n)}\)。
由定理 \ref{thm:xi-leyang-hankel-pencil-dominant-zeros} 的特征值收敛界（此时 \(q=R_1/R_2\)）得到
\(\widehat\xi_\pm^{(n)}\to \xi_{1,2}\) 且误差为 \(O(q^n)\)。
再由映射 \(\xi\mapsto 1/\xi\) 在 \(|\xi|=R_1^{-1}\) 的邻域上 Lipschitz（常数 \(\asymp R_1^2\)）得到对应的 \(\zeta\)-平面误差界。证毕。
\end{proof}

\begin{proposition}[指数收敛率的一般锐性]\label{prop:xi-leyang-exponential-rate-sharpness}
在推论 \ref{cor:xi-leyang-nearest-conjugate-pair-discriminant-estimator} 的设定下，进一步假设存在第二个最近共轭对
\[
|\zeta_3|=|\zeta_4|=R_2,\qquad \zeta_4=\overline{\zeta_3},
\]
且其余零点模长严格大于 \(R_2\)。
则存在常数 \(c\neq 0\) 与 \(n_1\in\NN\) 使对所有 \(n\ge n_1\) 有
\[
\max_{\sigma\in\{+,-\}}\abs{\widehat\zeta_\sigma^{(n)}-\zeta_\sigma}
\ \ge\ |c|\,R_1\left(\frac{R_1}{R_2}\right)^{n}.
\]
因此在仅由模长间隙 \(R_1<R_2\) 所刻画的类上，收敛率 \((R_1/R_2)^n\) 不能被统一改进。
\end{proposition}
\begin{proof}[证明要点]
将幂和分解为两对共轭主项与更远余项：
\(
s_n=(\xi_1^n+\overline{\xi_1}^n)+(\xi_3^n+\overline{\xi_3}^n)+\widetilde R_n
\)，其中 \(|\widetilde R_n|=o(|\xi_3|^n)\)。
把 \(A_n,B_n\) 视为 \((s_n,s_{n+1},s_{n+2},s_{n+3})\) 的有理函数，并按小参数
\(
\varepsilon_n:=(|\xi_3|/|\xi_1|)^n=(R_1/R_2)^n
\)
作一阶展开，可得在一般位置（相位非共振的开稠密条件）下
\[
A_n=(\xi_1+\overline{\xi_1})+\alpha\,\varepsilon_n+o(\varepsilon_n),\qquad
B_n=\xi_1\overline{\xi_1}+\beta\,\varepsilon_n+o(\varepsilon_n),
\]
其中 \((\alpha,\beta)\neq (0,0)\)。
二次多项式根对系数的一阶扰动给出根位移 \(\asymp \varepsilon_n\)，从而经 \(\xi\mapsto 1/\xi\) 得到 \(\zeta\)-平面误差下界 \(\asymp R_1\varepsilon_n\)。证毕。
\end{proof}

\begin{theorem}[Hankel 奇异值的对数斜率识别支配模长]\label{thm:xi-leyang-hankel-singular-slope-detects-moduli}
在定理 \ref{thm:xi-leyang-hankel-pencil-dominant-zeros} 的设定下，固定任意 \(k\ge r\)，并令
\[
H_0^{(k)}(n):=\bigl[s_{n+i+j}\bigr]_{i,j=0}^{k-1}.
\]
记其奇异值按降序为 \(\sigma_1(n)\ge\cdots\ge\sigma_k(n)\ge 0\)。
则对每个 \(1\le j\le r\) 存在常数 \(c_j,C_j>0\) 与 \(n_j\) 使对所有 \(n\ge n_j\) 有
\[
c_j\,|\xi_j|^{\,n}\ \le\ \sigma_j(n)\ \le\ C_j\,|\xi_j|^{\,n}.
\]
因此
\[
\lim_{n\to\infty}\frac{1}{n}\log\sigma_j(n)=\log|\xi_j|=-\log|\zeta_j|\qquad(1\le j\le r).
\]
\end{theorem}
\begin{proof}[证明要点]
将 \(H_0^{(k)}(n)\) 分解为秩 \(r\) 的主项与余项：
\(
H_0^{(k)}(n)=V_k D_n^{(r)} V_k^{\mathsf T}+\widetilde E(n)
\)，其中 \((V_k)_{i,j}=\xi_j^i\)（\(0\le i\le k-1,\,1\le j\le r\)），\(D_n^{(r)}=\mathrm{diag}(\xi_1^n,\dots,\xi_r^n)\)，且 \(\norm{\widetilde E(n)}\le C|\xi_{r+1}|^{\,n}\)。
对主项应用奇异值的 min--max 表达与 \(V_k\) 的左右有界性得到 \(\sigma_j(V_k D_n^{(r)} V_k^{\mathsf T})\asymp |\xi_j|^n\)；
再用 Weyl 型奇异值扰动不等式吸收余项，即得双边界与极限。证毕。
\end{proof}

\begin{theorem}[累积量噪声扰动下的零点重构稳定性]\label{thm:xi-leyang-reconstruction-stability-under-noise}
在定理 \ref{thm:xi-leyang-hankel-pencil-dominant-zeros} 的设定下，设观测到的幂和为 \(\widetilde s_n=s_n+\delta_n\)，并由此构造 \(\widetilde H_0^{(r)}(n),\widetilde H_1^{(r)}(n)\) 与
\(
\widetilde{\mathcal S}_r(n):=(\widetilde H_0^{(r)}(n))^{-1}\widetilde H_1^{(r)}(n)
\)。
若对某个 \(n\ge n_0\) 有
\begin{equation}\label{eq:xi-leyang-stability-invertibility-condition}
\norm{(H_0^{(r)}(n))^{-1}}\cdot \norm{\widetilde H_0^{(r)}(n)-H_0^{(r)}(n)}\le \frac12,
\end{equation}
则：
\begin{enumerate}
\item \(\widetilde H_0^{(r)}(n)\) 可逆，且
\begin{align}\label{eq:xi-leyang-stability-pencil-bound}
\norm{\widetilde{\mathcal S}_r(n)-\mathcal S_r(n)}
&\le
4\norm{(H_0^{(r)}(n))^{-1}}
\Bigl(\norm{\widetilde H_1^{(r)}(n)-H_1^{(r)}(n)}
\nonumber\\
&\hspace{3.7em}
+\norm{\mathcal S_r(n)}\cdot\norm{\widetilde H_0^{(r)}(n)-H_0^{(r)}(n)}\Bigr).
\end{align}
\item 若 \(\mathcal S_r(n)\) 可对角化，\(\mathcal S_r(n)=PDP^{-1}\)，则任意 \(\widetilde{\mathcal S}_r(n)\) 的特征值 \(\widetilde\xi\) 满足
\[
\mathrm{dist}\!\bigl(\widetilde\xi,\sigma(\mathcal S_r(n))\bigr)\ \le\ \kappa(P)\,\norm{\widetilde{\mathcal S}_r(n)-\mathcal S_r(n)},
\qquad
\kappa(P):=\norm{P}\,\norm{P^{-1}}.
\]
从而在 \(|\xi_j|\) 有下界的区域内，零点估计 \(\widetilde\zeta=1/\widetilde\xi\) 的误差亦可由同一右端给出显式控制。
\end{enumerate}
\end{theorem}
\begin{proof}
条件 \eqref{eq:xi-leyang-stability-invertibility-condition} 蕴含
\(
\widetilde H_0^{(r)}(n)=H_0^{(r)}(n)\bigl(I+(H_0^{(r)}(n))^{-1}(\widetilde H_0^{(r)}(n)-H_0^{(r)}(n))\bigr)
\)
中括号可逆，且
\(
\norm{(\widetilde H_0^{(r)}(n))^{-1}}\le 2\norm{(H_0^{(r)}(n))^{-1}}
\)。
再由恒等式
\[
\widetilde{\mathcal S}-\mathcal S
=(\widetilde H_0^{-1}-H_0^{-1})H_1+\widetilde H_0^{-1}(\widetilde H_1-H_1),
\qquad
\widetilde H_0^{-1}-H_0^{-1}=\widetilde H_0^{-1}(H_0-\widetilde H_0)H_0^{-1},
\]
直接得到 \eqref{eq:xi-leyang-stability-pencil-bound}。
第二条为 Bauer--Fike 定理的直接推论。证毕。
\end{proof}

\endinput

\paragraph{外二次谱与局部可识别性：二阶 Hankel--Wronskian、判别式平方与六窗消参证书}
沿用结论 \ref{con:xi-terminal-zm-order4-rational} 的双变量配分函数 \(Z_m(y)\) 及其最小湮灭多项式
\(\Pi(\lambda,y)=\lambda^4-\lambda^3-(2y+1)\lambda^2+\lambda+y(y+1)\)。
并沿用结论 \ref{con:xi-terminal-zm-discriminant-leyang} 的 Lee--Yang 三次因子
\[
P_{\mathrm{LY}}(y):=256y^3+411y^2+165y+32.
\]
\par
本段将二阶 Hankel 行列式视为外二次表示的 Pl\"ucker 坐标，并据此给出其闭递推、外二次判别式的完全平方分解、黄金比共振点处的降阶回文递推，以及仅用六窗局部数据恢复参数与消去参数的代数证书。
\par

\begin{conclusion}[外二次谱控制的二阶 Hankel--Wronskian 六阶闭递推]\label{con:xi-terminal-zm-exterior-square-hankel-wronskian-order6}
定义二阶 Hankel--Wronskian
\[
\Delta^{(2)}_m(y):=Z_m(y)Z_{m+3}(y)-Z_{m+1}(y)Z_{m+2}(y)\qquad(m\ge 0),
\]
并在固定 \(y\) 下简记 \(\Delta^{(2)}_m:=\Delta^{(2)}_m(y)\)。
则序列 \(\bigl(\Delta^{(2)}_m\bigr)_{m\ge 0}\) 满足一条六阶齐次线性递推：
\[
\boxed{\;
\begin{aligned}
&\Delta^{(2)}_{m+6}+(2y+1)\Delta^{(2)}_{m+5}-(y^2+y+1)\Delta^{(2)}_{m+4}-(4y^3+7y^2+3y+1)\Delta^{(2)}_{m+3}\\
&\qquad-(y^4+2y^3+2y^2+y)\Delta^{(2)}_{m+2}+(2y^5+5y^4+4y^3+y^2)\Delta^{(2)}_{m+1}+y^3(y+1)^3\Delta^{(2)}_{m}=0.
\end{aligned}\;}
\]
等价地，其特征多项式可取为
\[
\chi^{(2)}(\mu;y)
=\mu^6+(2y+1)\mu^5-(y^2+y+1)\mu^4-(4y^3+7y^2+3y+1)\mu^3-(y^4+2y^3+2y^2+y)\mu^2
\]
\[
\qquad\qquad\qquad\qquad\qquad\quad
+(2y^5+5y^4+4y^3+y^2)\mu+y^3(y+1)^3.
\]
若 \(\lambda_1,\dots,\lambda_4\) 为 \(\Pi(\lambda,y)=0\) 的根（在代数闭包中），则 \(\chi^{(2)}(\mu;y)\) 的根集合可表述为
\[
\{\lambda_i\lambda_j:\ 1\le i<j\le 4\}.
\]
\end{conclusion}
\begin{proof}[证明要点]
在代数闭包中，若 \(\Pi(\lambda,y)\) 的四根互异，则存在系数 \(c_1,\dots,c_4\) 使
\(
Z_m(y)=\sum_{i=1}^4 c_i\lambda_i^m
\)。
将该表达代入
\(
\Delta^{(2)}_m
=\det\!\bigl(\begin{smallmatrix}Z_m&Z_{m+1}\\ Z_{m+2}&Z_{m+3}\end{smallmatrix}\bigr)
\)
并展开行列式，可得
\(
\Delta^{(2)}_m=\sum_{1\le i<j\le 4}c_{ij}(\lambda_i\lambda_j)^m
\)
（反对称性消去 \(i=j\) 项）。
因此 \(\Delta^{(2)}_m\) 的指数模态由六个乘积 \(\lambda_i\lambda_j\) 控制，递推的特征多项式为
\(
\prod_{i<j}(\mu-\lambda_i\lambda_j)
\)，即为 \(\chi^{(2)}(\mu;y)\)。
将该乘积在 \(\QQ(y)\) 中展开得到所示显式系数。
由于两边系数均为 \(\ZZ[y]\) 中多项式，上述恒等式由稠密参数集推出对全体 \(y\) 成立。
显式展开与核验可由脚本 \texttt{scripts/exp\_fold\_zm\_bivariate\_partition\_audit.py} 复核。证毕。
\end{proof}

\begin{conclusion}[外二次判别式的完全平方分解与二阶共振临界集]\label{con:xi-terminal-zm-exterior-square-discriminant-square}
将 \(\chi^{(2)}(\mu;y)\) 视为关于 \(\mu\) 的六次多项式，则其判别式在 \(\QQ(y)\) 中具有完全平方分解：
\[
\boxed{\;
\mathrm{Disc}_{\mu}\bigl(\chi^{(2)}(\mu;y)\bigr)
=y^8(y-1)^2(y+1)^6(y^2+y-1)^2P_{\mathrm{LY}}(y)^2
=\Bigl(y^4(y-1)(y+1)^3(y^2+y-1)P_{\mathrm{LY}}(y)\Bigr)^2.
\;}
\]
从而其判别式支撑（等价地，单生成阶 \(A[\mu]\) 的有限判别式支撑，\(A:=\QQ[y]\)）包含于
\[
\boxed{\;y\in\{0,1,-1\}\ \cup\ \{y^2+y-1=0\}\ \cup\ \{P_{\mathrm{LY}}(y)=0\}.}
\]
并且由于判别式为平方，\(\mathrm{Gal}\bigl(\chi^{(2)}/\QQ(y)\bigr)\subseteq A_6\)。
\par
其中 \(y=-1\) 与 \(y^2+y-1=0\) 的出现属于单生成模型 \(A[\mu]\subset B\)（\(B\) 为 \(A\) 在 \(L=\QQ(y)(\mu)\) 中的整闭包）的指数理想缺陷；其并不对应扩张 \(L/\QQ(y)\) 的真实有限分歧。
真实有限分歧支撑精确等于 \(y=0\)、\(y=1\) 与 \(P_{\mathrm{LY}}(y)=0\)（见定理 \ref{thm:xi-terminal-zm-exterior-square-finite-ramification-index-ideal}）。
\end{conclusion}
\begin{proof}[证明要点]
由结论 \ref{con:xi-terminal-zm-exterior-square-hankel-wronskian-order6} 的显式表达，对 \(\chi^{(2)}(\mu;y)\) 作 \(\mu\)-判别式并在 \(\QQ[y]\) 中因式分解即可得到所示完全平方分解。
判别式在基域中为平方当且仅当伽罗瓦群包含于交错群，应用于次数 \(6\) 即得结论。
上述恒等式可由脚本 \texttt{scripts/exp\_fold\_zm\_bivariate\_partition\_audit.py} 复核。证毕。
\end{proof}

\leanverified{paper\_xi\_terminal\_zm\_exterior\_square\_finite\_ramification\_index\_ideal}
\begin{theorem}[外二次六次原始元的有限分歧精确定位与指数理想闭式公式]\label{thm:xi-terminal-zm-exterior-square-finite-ramification-index-ideal}
设 \(K=\QQ(y)\)，\(A=\QQ[y]\)，并令 \(\chi^{(2)}(\mu;y)\in A[\mu]\) 如结论 \ref{con:xi-terminal-zm-exterior-square-hankel-wronskian-order6}。
取 \(\mu\) 为其任一根并置 \(L:=K(\mu)\)。
记 \(B\) 为 \(A\) 在 \(L\) 中的整闭包。
则成立如下三点：
\begin{enumerate}
  \item \(\chi^{(2)}\) 的伽罗瓦闭包与谱四次 \(\Pi(\lambda,y)\) 的分裂域一致；
  且 \(\mathrm{Gal}(L/K)\cong S_4\) 并通过四元集合之 \(2\)-子集置换表示嵌入 \(A_6\)
  （见结论 \ref{con:xi-terminal-zm-exterior-square-galois-s4}）。
  \item \(B/A\) 的真实有限分歧支撑精确等于
  \[
  \{y=0\}\ \cup\ \{y=1\}\ \cup\ \{P_{\mathrm{LY}}(y)=0\},
  \]
  并且其有限判别式理想为主理想
  \[
  \mathfrak D_{B/A,\mathrm{fin}}=\bigl(y(y-1)P_{\mathrm{LY}}(y)\bigr)^2.
  \]
  \item 单生成阶 \(A[\mu]\subset B\) 的指数理想满足闭式恒等式
  \[
  [B:A]_{\mathrm{fin}}=\bigl(y^3(y+1)^3(y^2+y-1)\bigr),
  \]
  从而在有限素处有判别式指数公式
  \[
  \Disc_A(A[\mu])_{\mathrm{fin}}
  =\mathfrak D_{B/A,\mathrm{fin}}\cdot [B:A]_{\mathrm{fin}}^{2}.
  \]
  特别地，\(\mathrm{Disc}_{\mu}\bigl(\chi^{(2)}(\mu;y)\bigr)\) 中出现的 \(y+1\) 与 \(y^2+y-1\) 因子在有限处完全由指数缺陷贡献，而非真实分歧。
\end{enumerate}
上述结论的分歧--指数分解结构可参照内部审计记录 \cite{Internal2026ExteriorSquareFiniteRamificationIndexIdealAudit}。
\end{theorem}
\begin{proof}[证明要点]
第 \(1\) 点已由结论 \ref{con:xi-terminal-zm-exterior-square-galois-s4} 给出。
\par
对第 \(2\) 点，谱四次覆盖的 \(S_4\)-伽罗瓦闭包在 \(y\)-直线上的有限分歧值由 \(y=0,1\) 与 \(P_{\mathrm{LY}}(y)=0\) 给出（见结论 \ref{con:xi-terminal-zm-leyang-ramification-critical-values}），且相应惰性生成元可取为换位。
外二次扩张 \(L/K\) 对应 \(S_4\) 在 \(2\)-子集上的次数 \(6\) 置换表示；在该表示下换位诱导为循环型 \(2^2\cdot 1^2\)，故每个有限分歧值对不同的总贡献为 \(2\)，从而
\(
\mathfrak D_{B/A,\mathrm{fin}}=\bigl(y(y-1)P_{\mathrm{LY}}(y)\bigr)^2
\)。
\par
对第 \(3\) 点，\(\chi^{(2)}(\mu;y)\) 关于 \(\mu\) 首一且可分，故
\(
\Disc_A(A[\mu])=\mathrm{Disc}_{\mu}\bigl(\chi^{(2)}(\mu;y)\bigr)\in A
\)，并由 Dedekind 域上的经典指数公式
\(
\Disc_A(A[\mu])=\Disc_A(B)\,[B:A]^2
\)
比较有限素处指数即得
\(
[B:A]_{\mathrm{fin}}=(y^3(y+1)^3(y^2+y-1))
\)。
由此亦可见 \(y=-1\) 与 \(y^2+y-1=0\) 处仅为单生成阶的非极大性，而非 \(L/K\) 的真实有限分歧。证毕。
\end{proof}

\leanverified{paper\_xi\_terminal\_zm\_exterior\_square\_curve\_genus2}
\begin{theorem}[外二次谱曲线的正规化亏格]\label{thm:xi-terminal-zm-exterior-square-curve-genus2}
令 \(C^{\wedge 2}\) 为仿射曲线 \(\chi^{(2)}(\mu;y)=0\) 的正规化之光滑射影模型，并视为到 \(\PP^1_y\) 的次数 \(6\) 覆叠。
则
\[
g(C^{\wedge 2})=2.
\]
\end{theorem}
\begin{proof}[证明要点]
由定理 \ref{thm:xi-terminal-zm-exterior-square-finite-ramification-index-ideal}，有限分歧值为五点 \(y=0,1\) 与 \(P_{\mathrm{LY}}(y)=0\) 的三根，且各点惰性诱导为循环型 \(2^2\cdot 1^2\)，故每点分歧贡献为 \(2\)。
另一方面，无穷远处由谱四次的四循环惰性诱导（见结论 \ref{con:xi-terminal-zm-s4-splitting-curve-genus16}），其在 \(2\)-子集表示下循环型为 \(4\cdot 2\)，故无穷远分歧贡献为 \(4\)。
对次数 \(6\) 的覆盖应用 Riemann--Hurwitz 公式得
\[
2g(C^{\wedge 2})-2=-12+\bigl(4+5\cdot 2\bigr)=2,
\]
从而 \(g(C^{\wedge 2})=2\)。证毕。
\end{proof}

\begin{corollary}[外二次属二曲线 Jacobian 的椭圆分裂]\label{cor:xi-terminal-zm-exterior-square-jacobian-splitting}
\leanverified{paper\_xi\_terminal\_zm\_exterior\_square\_jacobian\_splitting}
在泛型参数下，\(C^{\wedge 2}\) 的 Jacobian 在 \(\QQ(y)\) 上同源分解为两条椭圆曲线之积。
更精确地，若 \(\widetilde{X}\to\PP^1_y\) 为谱四次的 \(S_4\)-伽罗瓦闭包，且取配对稳定子 \(V_4^{\mathrm{pair}}=\langle(12),(34)\rangle\le S_4\)，则 \(C^{\wedge 2}\) 与 \(\widetilde{X}/V_4^{\mathrm{pair}}\) 双有理等价，且
\[
J(C^{\wedge 2})\ \sim\ E_{\mathrm{res}}\times E,
\]
其中 \(E=\widetilde{X}/S_3\) 与 \(E_{\mathrm{res}}=\widetilde{X}/D_8\) 为两条椭圆曲线（见结论 \ref{con:xi-terminal-zm-s4-jacobian-group-algebra-prym-e3-b3}）。
\end{corollary}
\begin{proof}[证明要点]
由结论 \ref{con:xi-terminal-zm-exterior-square-galois-s4}，扩张 \(L/K\) 的单值化由 \(S_4\) 在 \(2\)-子集上的置换表示给出；
因此其对应的次数 \(6\) 中间域可识别为分裂域曲线 \(\widetilde{X}\) 对 \(2\)-子集稳定子的商，即 \(\widetilde{X}/V_4^{\mathrm{pair}}\)（参见结论 \ref{con:xi-terminal-zm-s4-v4pair-explicit-model} 中对该商的显式模型）。
于是 \(C^{\wedge 2}\) 与 \(\widetilde{X}/V_4^{\mathrm{pair}}\) 具有同一函数域，从而双有理等价。
最后，结论 \ref{con:xi-terminal-zm-s4-jacobian-group-algebra-prym-e3-b3} 给出 \(J(\widetilde{X}/V_4^{\mathrm{pair}})\sim E_{\mathrm{res}}\times E\)，遂得到所述分解。证毕。
\end{proof}

\begin{conclusion}[黄金比参数处的 \texorpdfstring{$(-1)$}{-1}-倒易对称与外二次谱塌缩]\label{con:xi-terminal-zm-minus-one-reciprocity-exterior-square}
对所有 \(y\) 恒成立恒等式
\[
\Pi(\lambda,y)-\lambda^{4}\Pi(-1/\lambda,y)=-(\lambda^{4}-1)(y^{2}+y-1).
\]
因此当 \(y^{2}+y-1=0\) 时，谱四次满足 \((-1)\)-倒易性
\[
\Pi(\lambda,y)=\lambda^{4}\Pi(-1/\lambda,y),
\]
并诱导根集合上的对合 \(\lambda\mapsto -1/\lambda\)。
从而外二次多项式满足
\[
(\mu+1)^{2}\ \bigm|\ \chi^{(2)}(\mu;y).
\]
\end{conclusion}
\begin{proof}[证明要点]
差值恒等式由对 \(\Pi(\lambda,y)\) 直接代入并化简得到。
当 \(y^{2}+y-1=0\) 时，\(\lambda\mapsto -1/\lambda\) 保持方程 \(\Pi(\lambda,y)=0\) 不变，故根集在该对合下成对出现；从而存在至少两组 \(i<j\) 使 \(\lambda_i\lambda_j=-1\)，于是 \(\chi^{(2)}\) 在 \(\mu=-1\) 处至少二重零点，即 \((\mu+1)^2\mid \chi^{(2)}(\mu;y)\)。证毕。
\end{proof}

\begin{conclusion}[\texorpdfstring{$\mu=-1$}{mu=-1} 的黄金比共振判别与二重根结构]\label{con:xi-terminal-zm-exterior-square-mu-minus-one-resonance}
对所有 \(y\) 有显式恒等式
\[
\boxed{\;\chi^{(2)}(-1;y)=y(y-1)(y^2+y-1)^2.\;}
\]
因此 \(\mu=-1\) 为 \(\chi^{(2)}(\mu;y)\) 的根当且仅当
\(
y\in\{0,1\}
\)
或
\(
y^2+y-1=0
\)。
并且在 \(y^2+y-1=0\) 且 \(y\notin\{0,1\}\) 时，\(\mu=-1\) 为 \(\chi^{(2)}(\mu;y)\) 的二重根。
\end{conclusion}
\begin{proof}[证明要点]
直接代入 \(\mu=-1\) 并在 \(\ZZ[y]\) 中因式分解得到恒等式。
再对 \(\partial_\mu\chi^{(2)}(\mu;y)\) 代入 \(\mu=-1\) 可见其亦被 \(y^2+y-1\) 整除，从而在 \(y^2+y-1=0\) 时至少二重；
结合结论 \ref{con:xi-terminal-zm-exterior-square-discriminant-square} 的判别式分解可知该处重根阶数恰为 \(2\)。证毕。
\end{proof}

\begin{conclusion}[黄金比参数处 \texorpdfstring{$\Delta^{(2)}_m$}{Delta(2)} 的最小递推降阶与回文五阶关系]\label{con:xi-terminal-zm-exterior-square-hankel-wronskian-order5-golden}
若 \(y\) 满足 \(y^2+y-1=0\)，则 \(\chi^{(2)}(\mu;y)\) 含因子 \((\mu+1)^2\)，并且 \(\bigl(\Delta^{(2)}_m(y)\bigr)_{m\ge 0}\) 的最小湮灭多项式为 \(\chi^{(2)}(\mu;y)\) 的平方自由部分；因而递推阶数由 \(6\) 降至 \(5\)。
在该约束下可取如下回文五阶递推：
\[
\boxed{\;
\Delta^{(2)}_{m+5}+2y\,\Delta^{(2)}_{m+4}-2(y+1)\Delta^{(2)}_{m+3}-2(y+1)\Delta^{(2)}_{m+2}+2y\,\Delta^{(2)}_{m+1}+\Delta^{(2)}_{m}=0.
\;}
\]
\end{conclusion}
\begin{proof}[证明要点]
由结论 \ref{con:xi-terminal-zm-exterior-square-mu-minus-one-resonance} 得在 \(y^2+y-1=0\) 时 \(\mu=-1\) 为二重根，故 \((\mu+1)^2\mid \chi^{(2)}(\mu;y)\)。
另一方面，结论 \ref{con:xi-terminal-zm-discriminant-leyang} 给出 \(\mathrm{Disc}_{\lambda}(\Pi)=-y(y-1)P_{\mathrm{LY}}(y)\)，而 \(y^2+y-1\) 与 \(y(y-1)P_{\mathrm{LY}}(y)\) 在 \(\QQ[y]\) 中互素，故在 \(y^2+y-1=0\) 下 \(\Pi(\lambda,y)\) 具有四个互异根；
因此其伴随矩阵在代数闭包中可对角化，而外二次表示在同一特征向量楔积基下亦对角化，故最小多项式为特征多项式的平方自由部分。
将 \(y^2+y-1=0\) 代入 \(\chi^{(2)}(\mu;y)\) 并约去一个 \((\mu+1)\) 因子，化简得到所示回文五阶递推。证毕。
\end{proof}

\begin{conclusion}[任意六窗局部数据的有理可识别性]\label{con:xi-terminal-zm-six-window-rational-identifiability}
令 \(m\ge 0\)，并设
\[
U_m:=Z_{m+4}-Z_{m+3}+Z_{m+1},\qquad
V_m:=Z_{m+5}-Z_{m+4}+Z_{m+2}.
\]
若 \(\Delta^{(2)}_m\neq 0\)，则参数 \(y\) 可由任意六个连续观测值 \(\{Z_m,\dots,Z_{m+5}\}\) 以有理式唯一恢复：
\[
\boxed{\;
y=\frac12\left(\frac{Z_mV_m-Z_{m+1}U_m}{\Delta^{(2)}_m}-1\right).
\;}
\]
\end{conclusion}
\begin{proof}[证明要点]
由结论 \ref{con:xi-terminal-zm-order4-rational} 的四阶递推在指标 \(m+4\)、\(m+5\) 处可写为二元线性方程组：
\[
U_m=(2y+1)Z_{m+2}-y(y+1)Z_m,\qquad
V_m=(2y+1)Z_{m+3}-y(y+1)Z_{m+1}.
\]
对未知量 \(A:=2y+1\)、\(B:=y(y+1)\) 用 Cramer 法则求解，系数行列式恰为
\(\Delta^{(2)}_m=Z_mZ_{m+3}-Z_{m+1}Z_{m+2}\)。
由此得到
\(
A=(Z_mV_m-Z_{m+1}U_m)/\Delta^{(2)}_m
\)，进而 \(y=(A-1)/2\)。证毕。
\end{proof}

\begin{conclusion}[消去参数的六窗四次 Pl\"ucker 型恒等式]\label{con:xi-terminal-zm-six-window-parameter-free-quartic}
在结论 \ref{con:xi-terminal-zm-six-window-rational-identifiability} 的记号下，不显式引入参数 \(y\)，六窗数据 \(\{Z_m,\dots,Z_{m+5}\}\) 必满足如下参数自由四次恒等式：
\[
\boxed{\;
\bigl(Z_mV_m-Z_{m+1}U_m\bigr)^2-\bigl(\Delta^{(2)}_m\bigr)^2-4\,\Delta^{(2)}_m\bigl(Z_{m+2}V_m-Z_{m+3}U_m\bigr)=0.
\;}
\]
\end{conclusion}
\begin{proof}[证明要点]
由上一结论同理可得
\[
A=\frac{Z_mV_m-Z_{m+1}U_m}{\Delta^{(2)}_m},
\qquad
B=\frac{Z_{m+2}V_m-Z_{m+3}U_m}{\Delta^{(2)}_m}.
\]
而 \((A,B)\) 来自同一参数 \(y\)，故恒满足代数约束 \(4B=A^2-1\)（由 \(A=2y+1\)、\(B=y(y+1)\) 直接计算）。
将 \(A,B\) 的有理表达代入并清分母，即得所示四次恒等式。证毕。
\end{proof}

\begin{conclusion}[外二次六次多项式的伽罗瓦群刻画与信息完备性]\label{con:xi-terminal-zm-exterior-square-galois-s4}
令 \(\chi^{(2)}(\mu;y)\) 如结论 \ref{con:xi-terminal-zm-exterior-square-hankel-wronskian-order6}。
则当 \(y\) 不落入结论 \ref{con:xi-terminal-zm-exterior-square-discriminant-square} 的分歧集合时，成立：
\begin{enumerate}
  \item \(\chi^{(2)}(\mu;y)\) 的分裂域与 \(\Pi(\lambda,y)\) 的分裂域相同；
  \item \(\mathrm{Gal}\bigl(\chi^{(2)}/\QQ(y)\bigr)\cong S_4\)，并以 \(S_4\) 作用于四元集合之 \(2\)-子集的标准方式忠实嵌入 \(A_6\subset S_6\)。
\end{enumerate}
特别地，结论 \ref{con:xi-terminal-zm-exterior-square-discriminant-square} 中判别式为平方与该嵌入的全偶置换性一致。
\end{conclusion}
\begin{proof}[证明要点]
设 \(\lambda_1,\dots,\lambda_4\) 为 \(\Pi(\lambda,y)\) 的根，则由外二次构造，\(\chi^{(2)}\) 的根为 \(\lambda_i\lambda_j\)（\(i<j\)），故其分裂域包含于 \(\Pi\) 的分裂域。
反之，若某自同构固定全部 \(\lambda_i\lambda_j\)，则其在 \(2\)-子集集合 \(\binom{\{1,2,3,4\}}{2}\) 上的诱导作用为恒等；
该作用为忠实，故该自同构必为恒等，从而两分裂域相同。
另一方面，由定理 \ref{thm:xi-terminal-zm-leyang-elliptic-structure} 知 \(\mathrm{Gal}(\Pi/\QQ(y))\cong S_4\)，于是 \(\mathrm{Gal}(\chi^{(2)}/\QQ(y))\cong S_4\)。
再由结论 \ref{con:xi-terminal-zm-exterior-square-discriminant-square} 的判别式平方分解可知该像包含于 \(A_6\)。
该嵌入的全偶性亦可由生成元循环型直接检验，例如换位 \((12)\) 在 \(2\)-子集上诱导为 \((1^2\,2^2)\)，而 \(4\)-循环 \((1234)\) 诱导为 \((4\cdot 2)\)，二者均为偶置换。证毕。
\end{proof}

\begin{conclusion}[谱对相位共振的结果式证书与 Chebyshev 压缩：等模条件与判别式的统一消元]\label{con:xi-terminal-zm-phase-resonance-resultant-chebyshev}
沿用结论 \ref{con:xi-terminal-zm-order4-rational} 的谱四次 \(\Pi(\lambda,y)\)。
对任意 \(u\in\CC^{\ast}\) 定义二参数消元多项式
\[
\mathcal{R}(u,y):=\operatorname{Res}_{\lambda}\bigl(\Pi(\lambda,y),\Pi(u\lambda,y)\bigr)\in\ZZ[u,y].
\]
则 \(\mathcal{R}(u,y)\) 在 \(\ZZ[u,y]\) 中严格分解为
\[
\mathcal{R}(u,y)=y(y+1)(u-1)^4\,F(u,y),
\]
从而归一化结果式可取为
\[
\mathcal F(u,y):=\frac{\mathcal{R}(u,y)}{y(y+1)(u-1)^4}=F(u,y).
\]
其中 \(F(u,y)\) 在 \(u\) 上为 \(12\) 次回文多项式、在 \(y\) 上为 \(6\) 次多项式，并可作相位压缩：令 \(s:=u+u^{-1}\)，则存在 \(s\) 的六次多项式 \(G(s,y)\in\ZZ[s,y]\) 使得
\[
F(u,y)=u^6G(s,y).
\]
上述 \(F\) 与 \(G\) 的显式系数与退化恒等式由下式给出。
% Auto-generated by scripts/exp_fold_zm_phase_resonance_resultant_chebyshev_audit.py
\[
\Pi(\lambda,y)=\lambda^{4}-\lambda^{3}-(2y+1)\lambda^{2}+\lambda+y(y+1),\qquad \mathcal{R}(u,y):=\operatorname{Res}_{\lambda}\bigl(\Pi(\lambda,y),\Pi(u\lambda,y)\bigr).
\]
\[
\mathcal{R}(u,y)=y(y+1)(u-1)^{4}\,F(u,y).
\]
\[
\begin{aligned}
F(u,y)=&\ y^{3} \left(y + 1\right)^{3}(u^{12}+1)
+\ y^{2} \left(y + 1\right)^{2} \left(2 y + 1\right)^{2}(u^{11}+u)
+\ y \left(y + 1\right) \left(2 y^{2} - 1\right) \left(y^{2} + y + 1\right)(u^{10}+u^{2})\\
&-\ \left(12 y^{6} + 50 y^{5} + 63 y^{4} + 38 y^{3} + 19 y^{2} + 6 y + 1\right)(u^{9}+u^{3})
-\ \left(17 y^{6} + 87 y^{5} + 96 y^{4} + 33 y^{3} + 16 y^{2} + 7 y + 2\right)(u^{8}+u^{4})\\
&+\ \left(8 y^{6} - 26 y^{5} + 11 y^{4} + 94 y^{3} + 51 y^{2} + 14 y + 1\right)(u^{7}+u^{5})
+\ \left(28 y^{6} + 32 y^{5} + 103 y^{4} + 186 y^{3} + 103 y^{2} + 32 y + 4\right)u^{6}.
\end{aligned}
\]
\[
F(u,y)=u^{6}G(s,y),\qquad s:=u+u^{-1}.
\]
\[
\begin{aligned}
G(s,y)=&\ \left(y^{6} + 3 y^{5} + 3 y^{4} + y^{3}\right)s^{6}
+\left(4 y^{6} + 12 y^{5} + 13 y^{4} + 6 y^{3} + y^{2}\right)s^{5}\\
&+\left(- 4 y^{6} - 14 y^{5} - 15 y^{4} - 6 y^{3} - 2 y^{2} - y\right)s^{4}
+\left(- 32 y^{6} - 110 y^{5} - 128 y^{4} - 68 y^{3} - 24 y^{2} - 6 y - 1\right)s^{3}\\
&+\left(- 16 y^{6} - 76 y^{5} - 81 y^{4} - 24 y^{3} - 8 y^{2} - 3 y - 2\right)s^{2}
+\left(64 y^{6} + 184 y^{5} + 265 y^{4} + 238 y^{3} + 113 y^{2} + 32 y + 4\right)s\\
&+\left(64 y^{6} + 208 y^{5} + 295 y^{4} + 250 y^{3} + 131 y^{2} + 44 y + 8\right).
\end{aligned}
\]
\[
G(2,y)=-y(y-1)\bigl(256y^{3}+411y^{2}+165y+32\bigr).
\]

\par
当 \(u=e^{i\theta}\) 时 \(s=2\cos\theta\)，从而“存在一对谱根满足比值 \(u\)”（因而模长相等）的消元条件可降维为代数曲线 \(G(2\cos\theta,y)=0\)。
并且在 \(u=1\)（即 \(s=2\)）处出现严格退化 \(G(2,y)=-y(y-1)P_{\mathrm{LY}}(y)\)，从而该二参数证书在单一框架下同时回收了判别式零轨迹（含 \(y_{\mathrm{LY}}\)）并编码全部等模谱对的相位共振集合。
\end{conclusion}
\begin{proof}[证明要点]
对变量 \(\lambda\) 作结果式消元并在 \(\ZZ[u,y]\) 中因式分解得到 \(\mathcal{R}(u,y)=y(y+1)(u-1)^4F(u,y)\)。
回文化与 \(s=u+u^{-1}\) 的压缩来自 \(u\mapsto u^{-1}\) 的互反对称；退化式 \(G(2,y)\) 由代入 \(s=2\) 后在 \(\ZZ[y]\) 中因式分解得到。
上述恒等式与系数可由脚本 \texttt{scripts/exp\_fold\_zm\_phase\_resonance\_resultant\_chebyshev\_audit.py} 复核。证毕。
\end{proof}

\begin{conclusion}[相位共振判别式的闭式分解与平方结构]\label{con:xi-terminal-zm-phase-compression-disc-square}
沿用结论 \ref{con:xi-terminal-zm-phase-resonance-resultant-chebyshev} 的记号：
\[
F(u,y)=u^{6}G(u+u^{-1},y),\qquad s=u+u^{-1},
\]
并记
\[
P_{\mathrm{LY}}(y):=256y^{3}+411y^{2}+165y+32.
\]
则关于 \(G\) 与 \(F\) 的端点值、判别式闭式分解与平方结构由下式统一给出：
% Auto-generated by scripts/exp_fold_zm_phase_compression_discriminant_square_audit.py
\[
G(2,y)=-y(y-1)P_{\mathrm{LY}}(y),\qquad G(-2,y)=y^{2}(y-1)^{2}.
\]
\[
\mathrm{Disc}_{s}\bigl(G(s,y)\bigr)=y^{10}(y-1)^{4}(y+1)^{6}(y^{2}+y-1)^{4}P_{\mathrm{LY}}(y)^{2}P_{9}(y)^{2}.
\]
\[
\mathrm{Disc}_{s}\bigl(G(s,y)\bigr)=\Delta_{G}(y)^{2},\qquad \Delta_{G}(y)=y^{5}(y-1)^{2}(y+1)^{3}(y^{2}+y-1)^{2}P_{\mathrm{LY}}(y)P_{9}(y).
\]
\[
P_{9}(y)=256 y^{9} + 1219 y^{8} + 2542 y^{7} + 3090 y^{6} + 2446 y^{5} + 1315 y^{4} + 478 y^{3} + 112 y^{2} + 16 y + 1.
\]
\[
\mathrm{Disc}_{u}\bigl(F(u,y)\bigr)=-y^{23}(y-1)^{11}(y+1)^{12}(y^{2}+y-1)^{8}P_{\mathrm{LY}}(y)^{5}P_{9}(y)^{4}.
\]
\[
\mathrm{Disc}_{u}\bigl(F(u,y)\bigr)=\mathrm{Disc}_{s}\bigl(G(s,y)\bigr)^{2}\,G(2,y)\,G(-2,y).
\]
\[
\mathrm{Disc}_{u}\bigl(F(u,y)\bigr)=-y(y-1)P_{\mathrm{LY}}(y)\,\Omega(y)^{2},\quad \Omega(y)=y^{11}(y-1)^{5}(y+1)^{6}(y^{2}+y-1)^{4}P_{\mathrm{LY}}(y)^{2}P_{9}(y)^{2}.
\]
\[
\gcd\Bigl(\operatorname{Res}_{s}\bigl(G,\partial_{s}G\bigr),\ \operatorname{Res}_{s}\bigl(G,\partial_{y}G\bigr)\Bigr)=y^{10}(y-1)^{3}(y+1)^{6}(y^{2}+y-1)^{4}P_{9}(y)^{2}.
\]
\[
G\!\left(s,\frac{-1-\sqrt5}{2}\right)=\bigl(s^{2}+(1-2\sqrt5)s+7\bigr)\bigl(s^{2}+(2+\sqrt5)s-1+2\sqrt5\bigr)^{2},
\]
\[
G\!\left(s,\frac{-1+\sqrt5}{2}\right)=\bigl(s^{2}+(1+2\sqrt5)s+7\bigr)\bigl(s^{2}+(2-\sqrt5)s-1-2\sqrt5\bigr)^{2}.
\]
\[
P_{\mathrm{LY}}(y_{0})=0\ \Longrightarrow\ G(s,y_{0})=y_{0}^{3}(y_{0}+1)^{3}(s-2)L(s)Q(s)^{2},
\]
\[
\widetilde{G}(s,y):=\frac{G(s,y)}{y^{3}(y+1)^{3}}\ \Longrightarrow\ \widetilde{G}(s,y_{0})=(s-2)L(s)Q(s)^{2}.
\]
\[
L(s)=s+\frac{-1280y_{0}^{2}+2041y_{0}+6583}{3168},\quad Q(s)=s^{2}+\frac{6400y_{0}^{2}-10205y_{0}-1235}{3168}\,s+\frac{16640y_{0}^{2}-26533y_{0}-9547}{1584}.
\]
\[
G(s,0)=-(s-2)(s+2)^{2},\qquad G(s,1)=(s-2)(s+2)^{2}(2s+5)^{2}(2s-5).
\]

特别地，\(\mathrm{Disc}_{s}(G)\) 在 \(\ZZ[y]\) 中为完全平方，且满足桥接恒等式
\[
\mathrm{Disc}_{u}\bigl(F(u,y)\bigr)=\mathrm{Disc}_{s}\bigl(G(s,y)\bigr)^2\,G(2,y)\,G(-2,y).
\]
\end{conclusion}
\begin{proof}[证明要点]
记 \(d:=\deg_s G=6\)，并在代数闭包上写
\[
G(s,y)=a(y)\prod_{i=1}^{d}(s-s_i),\qquad a(y)=y^{3}(y+1)^{3}.
\]
由 \(s=u+u^{-1}\) 得
\[
F(u,y)=u^{d}G(u+u^{-1},y)=a(y)\prod_{i=1}^{d}\bigl(u^{2}-s_i u+1\bigr).
\]
对二次因子 \(f_i(u):=u^2-s_i u+1\) 使用判别式与结式乘法恒等式
\[
\mathrm{Disc}\!\Bigl(\prod_i f_i\Bigr)=\prod_i \mathrm{Disc}(f_i)\cdot \prod_{i<j}\mathrm{Res}(f_i,f_j)^2
\]
并代入
\(
\mathrm{Disc}_u(f_i)=s_i^2-4,\ 
\mathrm{Res}_u(f_i,f_j)=(s_i-s_j)^2
\)，
再结合
\[
\prod_i(s_i^2-4)=\frac{G(2,y)G(-2,y)}{a(y)^2},\qquad
\prod_{i<j}(s_i-s_j)^4=\frac{\mathrm{Disc}_s(G)^2}{a(y)^{4d-4}},
\]
即可得到
\(
\mathrm{Disc}_u(F)=\mathrm{Disc}_s(G)^2\,G(2,y)\,G(-2,y)
\)，其中首项因子 \(a(y)\) 完全消去。
将显式 \(G(s,y)\) 代入并在 \(\ZZ[y]\) 中因式分解，得到所列闭式。
上述恒等式可由脚本 \texttt{scripts/exp\_fold\_zm\_phase\_compression\_discriminant\_square\_audit.py} 复核。证毕。
\end{proof}

\begin{theorem}[Lee--Yang 三次因子根处的双重碰撞与局部单值群型]\label{thm:xi-terminal-zm-leyang-cubic-double-collision-monodromy-type}
\leanverified{paper\_xi\_terminal\_zm\_leyang\_cubic\_double\_collision\_monodromy\_type}
在结论 \ref{con:xi-terminal-zm-phase-compression-disc-square} 的设定下，取 \(y_0\) 为 \(P_{\mathrm{LY}}(y)\) 的任一根。
设 \(\pi_{\mathrm{ph}}:X_{\mathrm{ph}}\to\PP^1_y\) 为平面曲线 \(G(s,y)=0\) 的射影闭包之正规化在 \(y\)-坐标上的投影（为六叶覆叠，见结论 \ref{con:xi-terminal-zm-phase-compression-galois-a6}）。
则在分歧点 \(y_0\) 处，\(\pi_{\mathrm{ph}}\) 的局部单值生成元在六张纸上的循环型为
\[
(2)(2)(1)(1)\in \mathfrak S_6,
\]
因而属于交错群 \(\mathfrak A_6\)。
并且 \(\mathrm{Disc}_{s}(G(s,y))\) 在 \(y=y_0\) 处的消失阶恰为 \(2\)。
\end{theorem}

\begin{proof}
由结论 \ref{con:xi-terminal-zm-phase-compression-disc-square} 的显式分解（式 \(\widetilde G(s,y)=G(s,y)/(y^3(y+1)^3)\) 的专化分解），在 \(y=y_0\) 处存在一次多项式 \(L\in\QQ(y_0)[s]\) 与二次多项式 \(Q\in\QQ(y_0)[s]\)（在 \(\overline{\QQ}\) 上有两互异根）使得
\[
G(s,y_0)=y_0^3(y_0+1)^3\,(s-2)\,L(s)\,Q(s)^2.
\]
因此专化多项式 \(G(\cdot,y_0)\) 在 \(s\)-坐标上恰有两个互异二重根（来自 \(Q(s)^2\)）与两个单根（来自 \((s-2)L(s)\)）。
对一参数代数族的局部解析分支，围绕 \(y_0\) 的小回路将每个二重根邻域内的两条分支互换（给出一换位），并且两处二重根互不相交，故局部单值为两不交换位之积，循环型为 \((2)(2)(1)(1)\)。
其为偶置换，从而落在 \(\mathfrak A_6\)。
\par
另一方面，\(\mathrm{Disc}_s(G)\) 在 \(\ZZ[y]\) 中含因子 \(P_{\mathrm{LY}}(y)^2\)（结论 \ref{con:xi-terminal-zm-phase-compression-disc-square}），故在 \(y=y_0\) 处至少二阶消失；
而上述专化仅出现两处简单二重根且其余根单纯，故判别式消失阶恰为 \(2\)。证毕。
\end{proof}

\leanverified{paper\_xi\_terminal\_zm\_leyang\_u\_root\_pinning\_double}
\begin{corollary}[Lee--Yang 临界根处的 \(u\)-根钉扎与互反二重根对]\label{cor:xi-terminal-zm-leyang-u-root-pinning-double}
在定理 \ref{thm:xi-terminal-zm-leyang-cubic-double-collision-monodromy-type} 的设定下，令 \(Q(s)=(s-\alpha)(s-\beta)\)（\(\alpha\neq\beta\)）为其在 \(\overline{\QQ}\) 上的分解。
则在 \(\overline{\QQ}[u]\) 中，\(F(u,y_0)=u^6G(u+u^{-1},y_0)\) 具有分解
\[
F(u,y_0)=u^{6}\cdot y_0^3(y_0+1)^3\cdot (u-1)^2\cdot \bigl(u^2-\sigma u+1\bigr)\cdot (u^2-\alpha u+1)^2\cdot (u^2-\beta u+1)^2,
\]
其中 \(\sigma\in\overline{\QQ}\) 为线性因子 \(L(s)=s-\sigma\) 的根。
因此：
\begin{enumerate}
  \item \(u=1\) 为 \(F(u,y_0)\) 的二重根；
  \item \(F(u,y_0)\) 额外包含两对互反二重根 \(\{u_\alpha,u_\alpha^{-1}\}\)、\(\{u_\beta,u_\beta^{-1}\}\)，分别为二次方程 \(u^2-\alpha u+1=0\)、\(u^2-\beta u+1=0\) 的根，并均以重数 \(2\) 出现。
\end{enumerate}
\end{corollary}

\begin{proof}
由 \(F(u,y)=u^6G(u+u^{-1},y)\) 与定理 \ref{thm:xi-terminal-zm-leyang-cubic-double-collision-monodromy-type} 中的专化分解，
每个 \(s\)-因子 \(s-s_0\) 在 \(u\)-坐标下对应二次因子 \(u^2-s_0u+1\)（参见结论 \ref{con:xi-terminal-zm-phase-compression-disc-square} 的一般分解式）。
当 \(s_0=2\) 时 \(u^2-2u+1=(u-1)^2\)，给出 \(u=1\) 的二重根；
当 \(s_0=\alpha,\beta\) 为二重根时，相应二次因子在 \(F(u,y_0)\) 中以平方出现，从而得到两对互反二重根。
线性因子 \(L(s)\) 贡献的 \(u^2-\sigma u+1\) 为一对互反单根，不影响上述二重性断言。证毕。
\end{proof}

\begin{theorem}[导子与单值群型给出的局部退化分层]\label{thm:xi-langlands-leyang-local-degeneration-stratification}
\leanverified{paper\_xi\_langlands\_leyang\_local\_degeneration\_stratification}
在 \(p_7\) 的 \(S_5\)-Artin 层（定理 \ref{thm:xi-p7-s5-odd-ramification-discriminant-and-conductor}--\ref{thm:xi-p7-s5-dyadic-c3-inertia-quadratic-functors-conductors}）与 Chebyshev--Lee--Yang 相位层（结论 \ref{con:xi-terminal-zm-phase-resonance-resultant-chebyshev}--推论 \ref{cor:xi-terminal-zm-leyang-u-root-pinning-double}）的共同设定下：
\begin{enumerate}
  \item 在奇分歧素数 \(q\) 处，\(\rho_4,\rho_5,\rho_6\) 的惰性限制具有共同三维不变子空间（定理 \ref{thm:xi-p7-s5-q-inertia-vacuum3-schur-functor-conductors}），从而任意 Schur 函子 \(\operatorname{Sym}^r,\bigwedge^r\) 的 \(q\)-导子指数可由有限组合计数给出（\eqref{eq:xi-p7-q-schur-symr-conductor-closed}--\eqref{eq:xi-p7-q-schur-wedger-conductor-closed}）。
  \item 在 Lee--Yang 三次因子根 \(y_0\) 处，相位压缩六叶覆叠 \(X_{\mathrm{ph}}\to\PP^1_y\) 发生双重碰撞，其局部单值循环型为 \((2)(2)(1)(1)\)（定理 \ref{thm:xi-terminal-zm-leyang-cubic-double-collision-monodromy-type}），并在 \(u\)-平面上表现为 \(u=1\) 与两对互反二重根的钉扎（推论 \ref{cor:xi-terminal-zm-leyang-u-root-pinning-double}）。
\end{enumerate}
\end{theorem}

\begin{proof}
第一条为定理 \ref{thm:xi-p7-s5-q-inertia-vacuum3-schur-functor-conductors} 的直接重述。
第二条由定理 \ref{thm:xi-terminal-zm-leyang-cubic-double-collision-monodromy-type} 与推论 \ref{cor:xi-terminal-zm-leyang-u-root-pinning-double} 合并得到。证毕。
\end{proof}

\begin{conclusion}[Chebyshev 压缩下的二次判别式类严格不变性]\label{con:xi-terminal-zm-phase-compression-disc-squareclass-invariance}
在函数域平方类群 \(\QQ(y)^{\times}/(\QQ(y)^{\times})^{2}\) 中成立
\[
\bigl[\mathrm{Disc}_{u}\bigl(F(u,y)\bigr)\bigr]
=\bigl[-y(y-1)P_{\mathrm{LY}}(y)\bigr]
=\bigl[\mathrm{Disc}_{\lambda}\bigl(\Pi(\lambda,y)\bigr)\bigr].
\]
等价地，\(\mathrm{Disc}_{u}(F)\) 的平方自由核恰为 \(-y(y-1)P_{\mathrm{LY}}(y)\)。
\end{conclusion}
\begin{proof}[证明要点]
由结论 \ref{con:xi-terminal-zm-phase-compression-disc-square} 的显式分解，
\[
\mathrm{Disc}_{u}(F)
=-y(y-1)P_{\mathrm{LY}}(y)\,\Omega(y)^2
\]
其中 \(\Omega(y)\in\QQ(y)\)。
因此其平方类仅由 \(-y(y-1)P_{\mathrm{LY}}(y)\) 决定。
再由结论 \ref{con:xi-terminal-zm-discriminant-leyang}
\(
\mathrm{Disc}_{\lambda}(\Pi)=-y(y-1)P_{\mathrm{LY}}(y)
\)
立即得到两者平方类相等。证毕。
\end{proof}

\begin{theorem}[判别式平方类锁定的同一二次中心特征]\label{thm:xi-terminal-zm-leyang-phase-resonance-discriminant-center-character}
\leanverified{paper\_xi\_terminal\_zm\_leyang\_phase\_resonance\_discriminant\_center\_character}
令 \(K:=\QQ(y)\)，并沿用结论 \ref{con:xi-terminal-zm-phase-resonance-resultant-chebyshev} 的
\(\Pi(\lambda,y)\in K[\lambda]\) 与 \(F(u,y)\in K[u]\)。
设 \(\Pi\) 与 \(F\) 在 \(K\) 上可分。
则由结论 \ref{con:xi-terminal-zm-phase-compression-disc-squareclass-invariance} 的平方类恒等式可知二次扩张严格同一：
\[
K\!\left(\sqrt{\mathrm{Disc}_{u}\bigl(F(u,y)\bigr)}\right)
\;=\;
K\!\left(\sqrt{\mathrm{Disc}_{\lambda}\bigl(\Pi(\lambda,y)\bigr)}\right)
\;=\;
K\!\left(\sqrt{-y(y-1)P_{\mathrm{LY}}(y)}\right).
\]
并且令 \(L_\Pi/K\)、\(L_F/K\) 分别为 \(\Pi\)、\(F\) 的分裂域，令
\[
\chi_\Pi:\mathrm{Gal}(L_\Pi/K)\to\{\pm 1\},\qquad
\chi_F:\mathrm{Gal}(L_F/K)\to\{\pm 1\}
\]
为其在根集置换表示上的行列式（等价地符号角色），则二者诱导的二次子扩张同一：
\[
L_\Pi^{\ker\chi_\Pi}=K\!\left(\sqrt{\mathrm{Disc}_{\lambda}(\Pi)}\right),
\qquad
L_F^{\ker\chi_F}=K\!\left(\sqrt{\mathrm{Disc}_{u}(F)}\right).
\]
从而 \(\chi_\Pi\) 与 \(\chi_F\) 的二次中心特征由同一平方类 \(-y(y-1)P_{\mathrm{LY}}(y)\) 统一刻画。
\end{theorem}
\begin{proof}
结论 \ref{con:xi-terminal-zm-phase-compression-disc-squareclass-invariance} 给出
\(
[\mathrm{Disc}_{u}(F)]=[\mathrm{Disc}_{\lambda}(\Pi)]
\)
在 \(K^\times/(K^\times)^2\) 中成立。
因此存在 \(\Omega\in K^\times\) 使得
\(
\mathrm{Disc}_{u}(F)=\Omega^2\,\mathrm{Disc}_{\lambda}(\Pi)
\)，从而
\(
K(\sqrt{\mathrm{Disc}_{u}(F)})=K(\sqrt{\mathrm{Disc}_{\lambda}(\Pi)})
\)。
再由结论 \ref{con:xi-terminal-zm-discriminant-leyang} 的判别式分解
\(
\mathrm{Disc}_{\lambda}(\Pi)=-y(y-1)P_{\mathrm{LY}}(y)
\)
即得首个等式链。
\par
对可分次数 \(n\) 多项式 \(f\in K[T]\)，其分裂域 \(L_f/K\) 上根集的置换表示给出一条符号角色
\(\chi_f:\mathrm{Gal}(L_f/K)\to\{\pm 1\}\)。
判别式的平方类刻画交错子群的固定子域：有标准等式
\[
L_f^{\ker\chi_f}=K\!\left(\sqrt{\mathrm{Disc}(f)}\right).
\]
将 \(f=\Pi\) 与 \(f=F\) 代入并结合已证的二次扩张同一性，即得结论。证毕。
\end{proof}

\begin{theorem}[良好素数处的 Frobenius 因子数奇偶同调]\label{thm:xi-terminal-zm-leyang-phase-resonance-frobenius-parity}
\leanverified{paper\_xi\_terminal\_zm\_leyang\_phase\_resonance\_frobenius\_parity}
令 \(p\) 为奇素数，取 \(y_0\in\FF_p\) 使得专化多项式
\(\Pi(\lambda,y_0)\in\FF_p[\lambda]\) 与 \(F(u,y_0)\in\FF_p[u]\) 均可分。
记
\[
r_{4}(p;y_0):=\#\{\text{\(\Pi(\lambda,y_0)\) 在 \(\FF_p[\lambda]\) 中的不可约因子}\},
\qquad
r_{12}(p;y_0):=\#\{\text{\(F(u,y_0)\) 在 \(\FF_p[u]\) 中的不可约因子}\}.
\]
则有同余关系
\[
r_{4}(p;y_0)\equiv r_{12}(p;y_0)\pmod 2.
\]
\end{theorem}
\begin{proof}
由定理 \ref{thm:xi-terminal-zm-leyang-phase-resonance-discriminant-center-character}，
\(\mathrm{Disc}_{u}(F)\) 与 \(\mathrm{Disc}_{\lambda}(\Pi)\) 在 \(K=\QQ(y)\) 的平方类群中相等；
在满足可分性且 \(p\neq 2\) 的良好约化点 \((p,y_0)\) 上，该平方类恒等式在
\(\FF_p^\times/(\FF_p^\times)^2\) 中保持，从而勒让德符号一致：
\[
\left(\frac{\mathrm{Disc}_{\lambda}(\Pi(\cdot,y_0))}{p}\right)
=
\left(\frac{\mathrm{Disc}_{u}(F(\cdot,y_0))}{p}\right).
\]
另一方面，Stickelberger 定理给出：对任意首一可分多项式 \(f\in\FF_p[T]\)（\(\deg f=n\)）记其不可约因子数为 \(r\)，则
\[
\left(\frac{\mathrm{Disc}(f)}{p}\right)=(-1)^{n-r}.
\]
将其分别用于 \(f=\Pi(\lambda,y_0)\)（\(n=4\)）与 \(f=F(u,y_0)\)（\(n=12\)），并注意 \(4,12\) 为偶数，得
\[
(-1)^{r_4(p;y_0)}
=(-1)^{4-r_4(p;y_0)}
=\left(\frac{\mathrm{Disc}_{\lambda}(\Pi(\cdot,y_0))}{p}\right)
=\left(\frac{\mathrm{Disc}_{u}(F(\cdot,y_0))}{p}\right)
=(-1)^{12-r_{12}(p;y_0)}
=(-1)^{r_{12}(p;y_0)}.
\]
故 \(r_{4}(p;y_0)\equiv r_{12}(p;y_0)\pmod 2\)。证毕。
\end{proof}

\begin{conclusion}[压缩六次多项式的二子集表示伽罗瓦像]\label{con:xi-terminal-zm-phase-compression-galois-a6}
设 \(k:=\QQ(y)\)，\(\Pi(\lambda,y)\in k[\lambda]\) 为谱四次多项式，且 \(\mathrm{Gal}(\Pi/k)\cong S_4\)；
记其分裂域为 \(L/k\)，根为 \(\lambda_1,\dots,\lambda_4\)。
定义
\[
s_{ij}:=\frac{\lambda_i}{\lambda_j}+\frac{\lambda_j}{\lambda_i}\in L\qquad (1\le i<j\le4).
\]
则 \(G(s,y)\) 的六个根恰为 \(\{s_{ij}\}_{i<j}\)，并且 \(G\) 在 \(k[s]\) 上不可约、\([k(s_{12}):k]=6\)。
此外，\(G\) 的分裂域等于 \(L\)，且
\[
\mathrm{Gal}\bigl(G/k\bigr)\cong S_4
\]
其在六个根上的作用同构于 \(S_4\) 在 \(\binom{\{1,2,3,4\}}{2}\) 上的自然作用。
该像包含于 \(A_6\)，故 \(\mathrm{Disc}_s(G)\) 为平方，与结论 \ref{con:xi-terminal-zm-phase-compression-disc-square} 一致。
\end{conclusion}
\begin{proof}[证明要点]
令 \(u_{ij}:=\lambda_i/\lambda_j\)。
由结果式构造
\(
F(u,y)=\operatorname{Res}_{\lambda}(\Pi(\lambda,y),\Pi(u\lambda,y))/\bigl(y(y+1)(u-1)^4\bigr)
\)
可知 \(F(u_{ij},y)=0\)，于是
\[
0=F(u_{ij},y)=u_{ij}^{6}G(u_{ij}+u_{ij}^{-1},y)=u_{ij}^{6}G(s_{ij},y),
\]
故每个 \(s_{ij}\) 为 \(G\) 的根。
由于 \(\deg_s G=6\) 且 \(\{s_{ij}\}_{i<j}\) 恰有 \(6\) 个，根集即为全部 \(s_{ij}\)。
对 \(s_{12}\) 而言，其在 \(S_4\) 下稳定子为集合稳定子 \(\mathrm{Stab}(\{1,2\})\)，阶为 \(4\)，故轨道大小为 \(24/4=6\)；
因而 \(s_{12}\) 在 \(k\) 上的最小多项式次数为 \(6\)，即 \(G\) 不可约且 \([k(s_{12}):k]=6\)。
再者，所有根 \(s_{ij}\in L\)，故 \(G\) 的分裂域 \(M\subseteq L\)。
而 \(S_4\) 在二子集上的作用是忠实作用，给出阶 \(24\) 的像；由 \(M\subseteq L\) 及 \([L:k]=24\) 得 \(M=L\)。
最后，\(S_4\) 的换位在二子集上的诱导置换是两个互换之积，为偶置换，故像落在 \(A_6\)。证毕。
\end{proof}

\begin{conclusion}[反相位点 \texorpdfstring{$u=-1$}{u=-1} 的刚性]\label{con:xi-terminal-zm-phase-compression-antiphase-rigidity}
对任意参数 \(y\) 有
\[
G(-2,y)=y^2(y-1)^2.
\]
因此 \(u=-1\)（即 \(s=-2\)）是方程 \(F(u,y)=0\) 的根当且仅当 \(y\in\{0,1\}\)。
特别地，Lee--Yang 条件 \(P_{\mathrm{LY}}(y)=0\) 不产生 \(u=-1\) 的相位共振根。
\end{conclusion}
\begin{proof}[证明要点]
由结论 \ref{con:xi-terminal-zm-phase-compression-disc-square}，\(G(-2,y)=y^2(y-1)^2\)。
又 \(F(-1,y)=(-1)^6G(-2,y)=G(-2,y)\)，故 \(F(-1,y)=0\iff y\in\{0,1\}\)。
当 \(P_{\mathrm{LY}}(y)=0\) 时 \(y\neq 0,1\)，从而不能触发 \(u=-1\) 根。证毕。
\end{proof}

\begin{conclusion}[黄金比例判别式分量的显式平方塌缩]\label{con:xi-terminal-zm-phase-compression-golden-square-collapse}
令
\[
\alpha:=\frac{-1-\sqrt5}{2},\qquad
\beta:=\frac{-1+\sqrt5}{2},
\]
则在 \(\QQ(\sqrt5)\) 上有精确因式分解
\[
G(s,\alpha)=\bigl(s^2+(1-2\sqrt5)s+7\bigr)\bigl(s^2+(2+\sqrt5)s-1+2\sqrt5\bigr)^2,
\]
\[
G(s,\beta)=\bigl(s^2+(1+2\sqrt5)s+7\bigr)\bigl(s^2+(2-\sqrt5)s-1-2\sqrt5\bigr)^2.
\]
\end{conclusion}
\begin{proof}[证明要点]
将 \(y=\alpha,\beta\) 代入显式 \(G(s,y)\)，并在 \(\QQ(\sqrt5)[s]\) 中进行代数因式分解即可得到两式。
该分解由脚本 \texttt{scripts/exp\_fold\_zm\_phase\_compression\_discriminant\_square\_audit.py} 直接核验。证毕。
\end{proof}

\begin{conclusion}[Lee--Yang 边界处的二重共振对机制]\label{con:xi-terminal-zm-phase-compression-leyang-double-resonance-pairs}
设 \(y_0\) 满足 \(P_{\mathrm{LY}}(y_0)=0\)，并记 \(k_0:=\QQ(y_0)\)。
则在 \(k_0[s]\) 中有
\[
G(s,y_0)=y_0^3(y_0+1)^3\,(s-2)\,L(s)\,Q(s)^2,
\]
其中
\[
L(s)=s+\frac{-1280y_0^2+2041y_0+6583}{3168},
\]
\[
Q(s)=s^2+\frac{6400y_0^2-10205y_0-1235}{3168}\,s+\frac{16640y_0^2-26533y_0-9547}{1584}.
\]
等价地，对首项归一化多项式 \(\widetilde G(s,y):=G(s,y)/(y^3(y+1)^3)\) 有
\[
\widetilde G(s,y_0)=(s-2)L(s)Q(s)^2.
\]
因此在 Lee--Yang 边界上，压缩六次多项式呈现 \((1,1,2,2)\) 根型。
\end{conclusion}
\begin{proof}[证明要点]
由结论 \ref{con:xi-terminal-zm-phase-compression-disc-square} 的端点恒等式，
\(
G(2,y)=-y(y-1)P_{\mathrm{LY}}(y)
\)，故 \(P_{\mathrm{LY}}(y_0)=0\Rightarrow s=2\) 为根。
将 \(y\) 约化到商域 \(\QQ[y]/(P_{\mathrm{LY}}(y))\) 后对 \(G\) 进行精确因式分解，可得所示
\((s-2)LQ^2\) 结构（对 \(G\) 本身仅差首项系数 \(y_0^3(y_0+1)^3\)）。
于是重根结构为一个二次因子的平方与两个一次因子，根型即 \((1,1,2,2)\)。
上述等式由脚本 \texttt{scripts/exp\_fold\_zm\_phase\_compression\_discriminant\_square\_audit.py} 核验。证毕。
\end{proof}

\begin{conclusion}[相位压缩曲线的正规化：属 \(2\) 与分歧签名由二子集表示完全决定]\label{con:xi-terminal-zm-phase-compression-curve-genus2}
令 \(X_{\mathrm{ph}}\) 为平面曲线 \(G(s,y)=0\) 的射影闭包之正规化，并令
\[
\pi:X_{\mathrm{ph}}\longrightarrow \PP^1_y
\]
为由 \(y\) 坐标诱导的有限态射。
则 \(\deg(\pi)=6\)，且 \(\pi\) 的分歧仅可能发生在
\(
y\in\{0,1,P_{\mathrm{LY}}(y)=0,\infty\}
\)
之上；其局部分歧型为：在每个有限分歧点上惯性置换在六张面上的循环型为 \((1^2\,2^2)\)，而在 \(y=\infty\) 上循环型为 \((4\cdot 2)\)。
此外，结论 \ref{con:xi-terminal-zm-phase-compression-disc-square} 的判别式分解在 \(y=-1\)、\(y^{2}+y-1=0\) 与 \(P_9(y)=0\) 处出现额外退化因子；这些对应于仿射平面模型的奇异纤维（结论 \ref{con:xi-terminal-zm-phase-compression-affine-singular-fibers}），而不产生新的域分歧点。
因此
\[
g(X_{\mathrm{ph}})=2.
\]
\end{conclusion}
\begin{proof}[证明要点]
由结论 \ref{con:xi-terminal-zm-phase-compression-galois-a6}，\(\QQ(X_{\mathrm{ph}})=\QQ(y,s_{12})\) 为 \(\Pi(\lambda,y)\) 的分裂域 \(L/\QQ(y)\) 的中间子扩张；
因此 \(X_{\mathrm{ph}}\to\PP^1_y\) 不可能在 \(L/\QQ(y)\) 的分歧集合之外引入新的域分歧。
又谱四次覆盖的 Hurwitz 分歧签名为 \((2,2,2,2,2,4)\)（结论 \ref{con:xi-terminal-zm-leyang-ramification-critical-values}），
将其中换位与四循环通过“对 \(2\)-子集的作用”诱导到次数 \(6\) 的表示，得到所述循环型 \((1^2\,2^2)\) 与 \((4\cdot 2)\)。
\par
对 \(\pi\) 应用 Riemann--Hurwitz 公式：
五个有限分歧点各贡献 \((2-1)+(2-1)=2\)，总计 \(10\)；
\(\infty\) 处贡献 \((4-1)+(2-1)=4\)，总分歧贡献 \(14\)。
于是
\(
2g(X_{\mathrm{ph}})-2=6(-2)+14
\),
从而 \(g(X_{\mathrm{ph}})=2\)。证毕。
\end{proof}

\begin{conclusion}[到 \texorpdfstring{$D_4$}{D4}--resolvent 椭圆曲线的分歧二重覆盖与 Jacobian 椭圆分解]\label{con:xi-terminal-zm-phase-compression-jacobian-splitting}
令 \(H:=\langle(12),(34)\rangle\le S_4\)，并令 \(D_4:=N_{S_4}(H)\) 为其正规化子（阶 \(8\) 的二面体 Sylow \(2\)-子群）。
取谱四次覆盖的 \(S_4\) 伽罗瓦闭包曲线 \(X_e\to\PP^1_y\)（结论 \ref{con:xi-terminal-zm-s4-d4fixedfield-resolvent-elliptic}），并记
\[
X_{\mathrm{ph}}:=X_e/H,\qquad R:=X_e/D_4.
\]
则存在自然次数 \(2\) 的态射
\[
\varphi:X_{\mathrm{ph}}\longrightarrow R
\]
使得 \(\pi\) 分解为 \(X_{\mathrm{ph}}\xrightarrow{\varphi}R\to\PP^1_y\)。
并且 \(\varphi\) 的分歧除子次数恰为 \(2\)，从而在 \(\QQ\) 上成立同源分解
\[
J(X_{\mathrm{ph}})\ \sim_{\QQ}\ R\times \operatorname{Prym}(\varphi),
\]
其中 \(\operatorname{Prym}(\varphi)\) 为一条定义于 \(\QQ\) 的椭圆曲线。
\end{conclusion}
\begin{proof}[证明要点]
由于 \(H\triangleleft D_4\) 且 \([D_4:H]=2\)，商映射 \(X_e/H\to X_e/D_4\) 给出 \(\varphi\)，其次数为 \(2\)。
由结论 \ref{con:xi-terminal-zm-phase-compression-curve-genus2} 得 \(g(X_{\mathrm{ph}})=2\)，而 \(R\) 为椭圆曲线（结论 \ref{con:xi-terminal-zm-s4-d4fixedfield-resolvent-elliptic}），故 \(g(R)=1\)。
Riemann--Hurwitz 对二重覆盖给出
\(
2g(X_{\mathrm{ph}})-2=2(2g(R)-2)+\deg\mathrm{Ram}(\varphi)
\),
从而 \(\deg\mathrm{Ram}(\varphi)=2\)。
最后，对任意二重覆盖 \(\varphi:C\to E\)（\(E\) 椭圆曲线），Prym 分解给出 \(J(C)\sim E\times \operatorname{Prym}(\varphi)\)，且 \(\dim\operatorname{Prym}=g(C)-g(E)=1\)，故 Prym 为椭圆曲线。证毕。
\end{proof}

\begin{conclusion}[平面相位压缩模型的奇异纤维投影：由 \texorpdfstring{$y=-1$}{y=-1}、黄金共轭与 \texorpdfstring{$P_9$}{P9} 控制]\label{con:xi-terminal-zm-phase-compression-affine-singular-fibers}
令 \(C_{\mathrm{aff}}\subset\mathbb A^2_{s,y}\) 为仿射曲线 \(G(s,y)=0\)。
则其奇异点 \((s_0,y_0)\) 必满足
\[
y_0(y_0-1)(y_0+1)(y_0^{2}+y_0-1)P_{9}(y_0)=0,
\]
其中 \(P_9\) 由结论 \ref{con:xi-terminal-zm-phase-compression-disc-square} 给出。
更精确地：
\begin{itemize}
  \item 在 \(y=0\) 与 \(y=1\) 上，唯一有限奇异点分别为 \((s,y)=(-2,0)\) 与 \((s,y)=(-2,1)\)。
  \item 在 \(y^{2}+y-1=0\) 的每个根上，存在恰好两个有限奇异点；其 \(s\)-坐标在 \(\QQ(\sqrt5)\) 中可取为
  \[
  s=\frac{\sqrt5-2}{2}\ \pm\ \frac12\sqrt{13+4\sqrt5},
  \]
  另一根对应其 \(\sqrt5\)-共轭。
  \item 在 \(P_9(y)=0\) 的每个根上，\(C_{\mathrm{aff}}\) 存在有限奇异点；其发生机制为在该参数处 \(G(\cdot,y)\) 出现二重 \(s\)-根。
\end{itemize}
\end{conclusion}
\begin{proof}[证明要点]
仿射奇异点当且仅当 \(G=\partial_sG=\partial_yG=0\)。
消去 \(s\) 并记
\[
D_1(y):=\operatorname{Res}_{s}\bigl(G,\partial_sG\bigr),\qquad
D_2(y):=\operatorname{Res}_{s}\bigl(G,\partial_yG\bigr).
\]
若存在奇异点，则必有 \(D_1(y)=D_2(y)=0\)。
由结论 \ref{con:xi-terminal-zm-phase-compression-disc-square} 中给出的
\(
\gcd(D_1,D_2)=y^{10}(y-1)^3(y+1)^6(y^2+y-1)^4P_9(y)^2
\)
可得 \(y_0\) 的必要条件即为所述零点集合。
\par
对 \(y=0,1\) 的断言由结论 \ref{con:xi-terminal-zm-phase-compression-disc-square} 中的分解
\(
G(s,0)=-(s-2)(s+2)^2
\)、\(
G(s,1)=(s-2)(s+2)^2(2s+5)^2(2s-5)
\)
直接检验三方程同解点得到。
对 \(y^{2}+y-1=0\) 的情形，在系数域 \(\QQ(\sqrt5)\) 上将 \(y\) 特化到任一根 \(y_0\) 后，
可计算得到
\(
\gcd_{\,\QQ(\sqrt5)[s]}\bigl(G(\cdot,y_0),\partial_sG(\cdot,y_0)\bigr)
\)
为一条可分二次多项式；例如当 \(y_0=\frac{-1+\sqrt5}{2}\) 时该因子可取为
\(
s^{2}-(\sqrt5-2)s-(1+2\sqrt5)
\)，另一根对应其 \(\sqrt5\)-共轭。
并且该二次因子整除 \(\partial_yG(\cdot,y_0)\)；
从而在该 \(y_0\) 上恰存在两条有限奇异分支，给出“恰好两个”有限奇异点。
所示 \(s\)-坐标即该二次多项式的求根公式展开。
对 \(P_9(y)=0\) 的情形，由 \(\gcd(D_1,D_2)\) 含因子 \(P_9(y)^2\) 可知每个 \(P_9\) 的根 \(y_0\) 处必存在 \(s_0\) 使 \(G=\partial_sG=\partial_yG=0\)，从而出现有限奇异点。
上述恒等式与投影必要条件可由脚本 \texttt{scripts/exp\_fold\_zm\_phase\_compression\_discriminant\_square\_audit.py} 复核。证毕。
\end{proof}

\begin{conclusion}[三重等模点的九次整系数刻画与算术不变量]\label{con:xi-terminal-zm-triple-equimodular-q9}
沿用结论 \ref{con:xi-terminal-zm-order4-rational} 的记号，令
\[
\Pi(\lambda,y):=\lambda^4-\lambda^3-(2y+1)\lambda^2+\lambda+y(y+1).
\]
则存在唯一参数
\[
y_\diamond\in(-1,0)
\]
使得 \(\Pi(\lambda,y_\diamond)\) 的谱半径由三个不同根同时取得：一条负实根 \(\lambda_-(y_\diamond)<0\) 与一对共轭复根 \(\lambda_c(y_\diamond),\overline{\lambda_c(y_\diamond)}\)，满足
\[
|\lambda_c(y_\diamond)|=-\lambda_-(y_\diamond)=:\rho_\diamond,
\qquad
|\lambda_+(y_\diamond)|<\rho_\diamond,
\]
其中 \(\lambda_+(y_\diamond)>0\) 为另一条实根。

该参数由如下九次整系数多项式在 \((-1,0)\) 上的唯一实根给出：
\[
Q(y)=256 y^9 + 1219 y^8 + 2542 y^7 + 3090 y^6 + 2446 y^5 + 1315 y^4 + 478 y^3 + 112 y^2 + 16 y + 1.
\]
并且 \(Q\) 在 \(\QQ[y]\) 上不可约，且
\[
\operatorname{Disc}(Q)=2^{28}\cdot 3^{9}\cdot 37\cdot 43^2\cdot 211\cdot 82301.
\]

数值上，
\[
y_\diamond\approx -0.160091463528933,\qquad
\rho_\diamond\approx 0.956702551378141,
\]
\[
\lambda_-(y_\diamond)\approx-0.956702551378141,\qquad
\lambda_c(y_\diamond)\approx 0.901572819032290\pm 0.320072216531856\,i,
\]
\[
\lambda_+(y_\diamond)\approx 0.153556913313561.
\]
\end{conclusion}
\begin{proof}[证明要点]
在 \(y\in(-1,0)\) 且 \(\operatorname{Disc}_{\lambda}(\Pi)\neq 0\) 的区域，\(\Pi\) 有两条实根与一对共轭复根且均为单根。
设四根记为 \((\lambda,b,\mu,\overline\mu)\)（\(\lambda<0<b\)），并令 \(u:=y(y+1)\)。
等模条件
\(
|\mu|=|\lambda|=-\lambda
\)
等价于 \(\mu\overline\mu=\lambda^2\)，由韦达关系得
\[
b=\frac{u}{\lambda^3},
\qquad
\sigma_3=\lambda b(\mu+\overline\mu)+b\mu\overline\mu=-1.
\]
消去 \(\mu+\overline\mu\) 后可得必要且充分的代数条件
\[
S(\lambda,y):=\lambda^8+\lambda^5+u\lambda^3-u^2=0.
\]
于是三重等模点满足联立方程
\[
\Pi(\lambda,y)=0,\qquad S(\lambda,y)=0,\qquad \lambda<0,\ y\in(-1,0).
\]
对 \(\lambda\) 取结果式：
\[
R(y):=\operatorname{Res}_{\lambda}\bigl(\Pi(\lambda,y),S(\lambda,y)\bigr)
=y(y-1)(y+1)\bigl(y^3+3y^2+y+1\bigr)^2Q(y).
\]
在 \((-1,0)\) 内可排除 \(y=0,\pm 1\) 与三次因子实根（该实根位于 \((- \infty,-1)\)），故三重等模条件在该区间内等价于 \(Q(y)=0\)。
由 Sturm 链可知 \(Q\) 在 \((-1,0)\) 上恰有一个实根，即 \(y_\diamond\)。
再由谱分支连续性与数值核验可得 \(|\lambda_+(y_\diamond)|<\rho_\diamond\)，故谱半径恰由三根同时取得。
不可约性与判别式可由 \(\QQ[y]\) 上因式分解及判别式计算直接验证。证毕。
\end{proof}

\begin{conclusion}[振荡主导区间上的实零点量子化律与指数误差]\label{con:xi-terminal-zm-leyang-real-zero-quantization}
沿用结论 \ref{con:xi-terminal-zm-triple-equimodular-q9} 的记号，取任意紧区间
\[
I=[a,b]\subset(-\infty,y_\diamond),
\]
并令 \(\lambda_c(y)\) 为上半平面中模最大的特征根分支，写作
\[
\lambda_c(y)=\rho(y)e^{i\theta(y)},\qquad \rho(y)>0,\ \theta(y)\in(0,\pi).
\]
则存在 \(\eta\in(0,1)\) 与在 \(I\) 上解析且处处非零的函数 \(A(y)\)、\(\phi(y)\)，使得一致地
\[
Z_m(y)
=
2\,\rho(y)^m\,|A(y)|\cos\!\bigl(m\theta(y)+\phi(y)\bigr)
+ O\!\bigl((\eta\rho(y))^m\bigr)
\qquad (m\to\infty,\ y\in I).
\]

若进一步假设 \(\theta'(y)\neq 0\) 于 \(I\) 上成立，则对充分大 \(m\)，\(Z_m\) 在 \(I\) 内存在一族互异实零点 \(\{y_{m,k}\}\)，并满足
\[
m\theta(y_{m,k})+\phi(y_{m,k})
=
\left(k+\frac12\right)\pi + O(\eta^m).
\]
这些零点均为单根，且计数满足
\[
\#\{y_{m,k}\in I\}
=
\frac{m}{\pi}\bigl(\theta(b)-\theta(a)\bigr)+O(1).
\]
对应经验测度弱收敛到绝对连续测度
\[
\dd \nu_I(y)=\frac{1}{\pi}\,|\theta'(y)|\,\dd y.
\]
\end{conclusion}
\begin{proof}[证明要点]
由结论 \ref{con:xi-terminal-zm-order4-rational}，
\[
\sum_{m\ge 0}Z_m(y)t^m=\frac{N(t,y)}{D(t,y)},
\]
\[
N(t,y)=1+yt-yt^2-y^2t^3,\qquad
D(t,y)=1-t-(2y+1)t^2+t^3+y(y+1)t^4.
\]
在判别式非零处，\(D\) 的根 \(t_j(y)\) 均为单根，令 \(\lambda_j(y):=1/t_j(y)\)，由留数展开
\[
Z_m(y)=\sum_{j=1}^4 c_j(y)\lambda_j(y)^m,\qquad
c_j(y)=-\lambda_j(y)\frac{N(1/\lambda_j(y),y)}{\partial_t D(1/\lambda_j(y),y)}.
\]
并且
\[
\operatorname{Res}_{t}(N,D)=-y^4(y-1)^3,
\]
故在 \(I\subset(-\infty,y_\diamond)\) 上（尤其 \(y\neq 0,1\)）各 \(c_j\) 解析且不消失。
由 \(y_\diamond\) 的定义，\(I\) 上主导谱为 \(\lambda_c,\overline{\lambda_c}\)，且与其余两根存在一致谱隙：存在 \(\eta\in(0,1)\) 使
\[
\max_{j\notin\{c,\overline c\}}|\lambda_j(y)|\le \eta\,|\lambda_c(y)|,\qquad y\in I.
\]
因此
\[
Z_m(y)=c_c(y)\lambda_c(y)^m+\overline{c_c(y)\lambda_c(y)^m}+O((\eta\rho(y))^m),
\]
写 \(c_c=|A|e^{i\phi}\) 即得主式。

若 \(\theta'(y)\neq 0\)，则
\(
g_m(y):=m\theta(y)+\phi(y)
\)
在大 \(m\) 时严格单调，主项 \(\cos g_m\) 的符号翻转与小误差叠加给出唯一零点。
对零点处导数使用
\[
\frac{\dd}{\dd y}\cos g_m(y)=-g_m'(y)\sin g_m(y),
\]
并由 \(g_m'(y)=m\theta'(y)+O(1)\) 得单根性。
量子化关系由将误差吸收入相位获得；计数与极限测度由 Riemann 和逼近得到。证毕。
\end{proof}

\begin{conclusion}[零点密度的显式代数微分表达与椭圆谱几何解释]\label{con:xi-terminal-zm-density-elliptic-dlog}
在结论 \ref{con:xi-terminal-zm-leyang-real-zero-quantization} 的设定下，
\[
\theta'(y)
=
\Im\!\left(\frac{\lambda_c'(y)}{\lambda_c(y)}\right)
=
\Im\!\left(
\frac{2\lambda_c(y)^2-(2y+1)}
{\lambda_c(y)\bigl(4\lambda_c(y)^3-3\lambda_c(y)^2-(4y+2)\lambda_c(y)+1\bigr)}
\right).
\]
因此零点极限测度密度可写为
\[
\dd\nu_I(y)
=
\frac{1}{\pi}
\left|
\Im\!\left(
\frac{2\lambda_c(y)^2-(2y+1)}
{\lambda_c(y)\bigl(4\lambda_c(y)^3-3\lambda_c(y)^2-(4y+2)\lambda_c(y)+1\bigr)}
\right)
\right|
\dd y.
\]

进一步，在谱曲线 \(\mathcal C:\Pi(\lambda,y)=0\) 上定义亚纯 \(1\)-形式
\[
\omega:=\dd\log\lambda
=
-\frac{\partial_y\Pi(\lambda,y)}{\lambda\,\partial_\lambda\Pi(\lambda,y)}\,\dd y.
\]
则沿主导复支有
\[
\dd\theta=\Im(\omega),\qquad
\nu_I=\frac{1}{\pi}\,(\pi_y)_*\!\bigl(|\Im(\omega)|\bigr).
\]
\end{conclusion}
\begin{proof}[证明要点]
由隐函数关系 \(\Pi(\lambda_c(y),y)=0\) 可得
\[
\partial_\lambda\Pi(\lambda_c(y),y)\lambda_c'(y)+\partial_y\Pi(\lambda_c(y),y)=0,
\]
\[
\frac{\lambda_c'(y)}{\lambda_c(y)}
=
-\frac{\partial_y\Pi(\lambda_c(y),y)}{\lambda_c(y)\partial_\lambda\Pi(\lambda_c(y),y)}.
\]
代入
\[
\partial_y\Pi(\lambda,y)=-2\lambda^2+2y+1,\qquad
\partial_\lambda\Pi(\lambda,y)=4\lambda^3-3\lambda^2-(4y+2)\lambda+1
\]
即得 \(\theta'(y)\) 的显式公式。
再由结论 \ref{con:xi-terminal-zm-leyang-real-zero-quantization} 的
\(
\dd\nu_I=(1/\pi)|\theta'(y)|\dd y
\)
得到密度表达。
最后 \(\dd\theta=\Im(\dd\log\lambda)=\Im(\omega)\)，故推前测度表述成立。证毕。
\end{proof}

\begin{conclusion}[\texorpdfstring{$m^{-1}$}{1/m} 局部缩放下的三相干涉极限方程]\label{con:xi-terminal-zm-triple-equimodular-local-scaling}
令 \(y_\diamond\) 为结论 \ref{con:xi-terminal-zm-triple-equimodular-q9} 的三重等模点，并记
\[
\lambda_-(y_\diamond)=-\rho_\diamond,\qquad
\lambda_c(y_\diamond)=\rho_\diamond e^{i\theta_\diamond}.
\]
在谱展开
\(
Z_m(y)=\sum_{j=1}^4 c_j(y)\lambda_j(y)^m
\)
中，定义
\[
c_-^\diamond:=c_-(y_\diamond)\in\RR\setminus\{0\},\qquad
c_c^\diamond:=c_c(y_\diamond)\in\CC\setminus\{0\},
\]
\[
\alpha:=
\Re\!\left(\frac{\lambda_c'}{\lambda_c}-\frac{\lambda_-'}{\lambda_-}\right)\Big|_{y=y_\diamond},
\qquad
\beta:=
\Im\!\left(\frac{\lambda_c'}{\lambda_c}\right)\Big|_{y=y_\diamond},
\qquad
\phi_\diamond:=\arg(c_c^\diamond).
\]
对任意满足 \(m_n\theta_\diamond\equiv\delta\pmod{2\pi}\) 的整数子列 \(\{m_n\}\)（并可固定奇偶类），在任意紧致 \(\xi\)-区间上一致有
\[
\rho_\diamond^{-m_n}
e^{-\Re(\lambda_-'/\lambda_-)|_{y=y_\diamond}\,\xi}
Z_{m_n}\!\left(y_\diamond+\frac{\xi}{m_n}\right)
\longrightarrow
(-1)^{m_n}c_-^\diamond
+2|c_c^\diamond|e^{\alpha\xi}
\cos\!\bigl(\delta+\beta\xi+\phi_\diamond\bigr).
\]
因此，窗口 \(|y-y_\diamond|=O(m^{-1})\) 内零点在缩放变量 \(\xi\) 下收敛到方程
\[
(-1)^m c_-^\diamond
+2|c_c^\diamond|e^{\alpha\xi}
\cos\!\bigl(\delta+\beta\xi+\phi_\diamond\bigr)=0
\]
的解集。

数值上
\[
\alpha\approx -0.1820719187487,\qquad
\beta\approx -1.454672700610.
\]
\end{conclusion}
\begin{proof}[证明要点]
在 \(y_\diamond\) 的小邻域内，四条特征根保持单纯，故 \(\lambda_j(y)\)、\(c_j(y)\) 解析。
令 \(y=y_\diamond+\xi/m\)，则
\[
\lambda_j(y)
=
\lambda_j(y_\diamond)\exp\!\left(
\frac{\lambda_j'(y_\diamond)}{\lambda_j(y_\diamond)}\frac{\xi}{m}
+O(m^{-2})
\right),
\]
\[
\lambda_j(y)^m
=
\lambda_j(y_\diamond)^m
\exp\!\left(
\frac{\lambda_j'(y_\diamond)}{\lambda_j(y_\diamond)}\xi
+O(m^{-1})
\right),
\]
且 \(c_j(y)=c_j(y_\diamond)+O(m^{-1})\)。
代回谱展开后，第四根满足 \(|\lambda_+(y_\diamond)|<\rho_\diamond\)，故在 \(\rho_\diamond^{-m}\) 归一化下指数消失。
其余三根中，\(\lambda_-=-\rho_\diamond\) 贡献 \((-1)^m\) 相位，而 \(\lambda_c,\overline{\lambda_c}\) 合并为余弦项并带指数倾斜 \(e^{\alpha\xi}\) 与线性相位 \(\beta\xi\)。
沿 \(m_n\theta_\diamond\to\delta\ (\mathrm{mod}\ 2\pi)\) 的子列即得一致极限。
零点极限方程由极限函数的过零集合直接给出。证毕。
\end{proof}


% (User input window)

\begin{conclusion}[\texorpdfstring{$(q=2)$}{(q=2)} 的重量—碰撞配分函数：六阶常系数递推与有理生成函数]\label{con:xi-terminal-zm2-order6-rational}
令
\(
X_m=\{x\in\{0,1\}^m:\ x_ix_{i+1}=0\}
\)，
并记 Fold 纤维多重度 \(d_m(x)=|\Fold_m^{-1}(x)|\)。
定义二阶重量—碰撞配分函数
\[
Z_m^{(2)}(y):=\sum_{x\in X_m} y^{|x|_1}\,d_m(x)^2\in\ZZ[y].
\]
则对一切 \(m\ge 7\) 成立严格六阶递推
\[
\begin{aligned}
Z_m^{(2)}(y)=&\ Z_{m-1}^{(2)}(y)+(2+3y)Z_{m-2}^{(2)}(y)+(y-2)Z_{m-3}^{(2)}(y)\\
&\ +(-1-2y-3y^2)Z_{m-4}^{(2)}(y)+(1-y)Z_{m-5}^{(2)}(y)+(y+y^3)Z_{m-6}^{(2)}(y),
\end{aligned}
\]
并且其常规生成函数为显式有理函数
\[
\boxed{\;
\sum_{m\ge 0} Z_m^{(2)}(y)t^m
=\frac{1+yt+(1-2y)t^2+(y-2y^2)t^3+(y^2-y)t^4+(y^3-y^2)t^5}{1-t-(2+3y)t^2+(2-y)t^3+(1+2y+3y^2)t^4-(1-y)t^5-(y+y^3)t^6}\; }.
\]
其中可取初值
\[
Z_0^{(2)}(y)=1,\ 
Z_1^{(2)}(y)=y+1,\ 
Z_2^{(2)}(y)=2y+4,\ 
Z_3^{(2)}(y)=y^2+9y+4,
\]
\[
Z_4^{(2)}(y)=6y^2+21y+9,\ 
Z_5^{(2)}(y)=y^3+31y^2+47y+9,\ 
Z_6^{(2)}(y)=15y^3+96y^2+93y+16.
\]
\end{conclusion}
\begin{proof}[证明要点]
与结论 \ref{con:xi-terminal-zm-order4-rational} 同理，
\(Z_m^{(2)}(y)\) 是有限状态规约器的二阶碰撞计数在重量变量 \(y^{|x|_1}\) 下的扭曲求和，
故其 \(t\)-生成函数为有理函数，约分后分母次数等于最小线性实现的维数，从而诱导常系数递推。
分子分母的显式系数可由有限窗口精确枚举与有理重建（Pad\'e/Prony）恢复，并可由脚本
\texttt{scripts/exp\_fold\_zq\_weight\_collision\_partition\_audit.py} 一键复核。证毕。
\end{proof}

\begin{conclusion}[\texorpdfstring{$(q=3)$}{(q=3)} 的重量—碰撞配分函数：八阶常系数递推与有理生成函数]\label{con:xi-terminal-zm3-order8-rational}
定义三阶重量—碰撞配分函数
\[
Z_m^{(3)}(y):=\sum_{x\in X_m} y^{|x|_1}\,d_m(x)^3\in\ZZ[y].
\]
则对一切 \(m\ge 9\) 成立严格八阶递推
\[
\begin{aligned}
Z_m^{(3)}(y)=&\ Z_{m-1}^{(3)}(y)+(3+4y)Z_{m-2}^{(3)}(y)+(4y-3)Z_{m-3}^{(3)}(y)\\
&\ +(-3+y-6y^2)Z_{m-4}^{(3)}(y)+(3-6y+y^2)Z_{m-5}^{(3)}(y)\\
&\ +(1+2y-5y^2+4y^3)Z_{m-6}^{(3)}(y)+(-1+2y-y^2)Z_{m-7}^{(3)}(y)\\
&\ +(-y^4+y^3+y^2-y)Z_{m-8}^{(3)}(y),
\end{aligned}
\]
并且生成函数可写为
\[
\sum_{m\ge 0} Z_m^{(3)}(y)t^m=\frac{N_3(t,y)}{D_3(t,y)},
\]
其中
\[
N_3(t,y)=1+yt+(4-3y)t^2+(4y-3y^2)t^3+(1-6y+3y^2)t^4+(y-6y^2+3y^3)t^5-y(y-1)^2t^6-y^2(y-1)^2t^7,
\]
并且
\[
\begin{aligned}
D_3(t,y)=&\ 1-t-(3+4y)t^2-(4y-3)t^3+(6y^2-y+3)t^4-(y^2-6y+3)t^5\\
&\ -(4y^3-5y^2+2y+1)t^6+(y-1)^2t^7+(y^4-y^3-y^2+y)t^8.
\end{aligned}
\]
其中可取初值
\[
Z_0^{(3)}(y)=1,\ 
Z_1^{(3)}(y)=y+1,\ 
Z_2^{(3)}(y)=2y+8,\ 
Z_3^{(3)}(y)=y^2+17y+8,
\]
\[
Z_4^{(3)}(y)=10y^2+51y+27,\ 
Z_5^{(3)}(y)=y^3+79y^2+153y+27,
\]
\[
Z_6^{(3)}(y)=37y^3+324y^2+395y+64,\ 
Z_7^{(3)}(y)=y^4+298y^3+1304y^2+837y+64,\ 
Z_8^{(3)}(y)=101y^4+1755y^3+3965y^2+1822y+125.
\]
\end{conclusion}
\begin{proof}[证明要点]
同上：有限状态碰撞计数在重量变量下保持有限维线性实现，
从而给出有理生成函数与常系数递推；显式系数与初值由有限数据重建与枚举核验得到，
并由脚本 \texttt{scripts/exp\_fold\_zq\_weight\_collision\_partition\_audit.py} 输出可审计证书。证毕。
\end{proof}

\begin{conclusion}[\texorpdfstring{$(q=4)$}{(q=4)} 的重量—碰撞配分函数：十阶常系数递推与核多项式的在位回收]\label{con:xi-terminal-zm4-order10-rational}
定义四阶重量—碰撞配分函数
\[
Z_m^{(4)}(y):=\sum_{x\in X_m} y^{|x|_1}\,d_m(x)^4\in\ZZ[y].
\]
则对一切 \(m\ge 11\) 成立严格十阶递推
\[
\begin{aligned}
Z_m^{(4)}(y)=&\ Z_{m-1}^{(4)}(y)+(4+5y)Z_{m-2}^{(4)}(y)+(11y-4)Z_{m-3}^{(4)}(y)\\
&\ +(-10y^2+19y-6)Z_{m-4}^{(4)}(y)+(11y^2-18y+6)Z_{m-5}^{(4)}(y)\\
&\ +(10y^3-27y^2+2y+4)Z_{m-6}^{(4)}(y)+(y^3-13y^2+9y-4)Z_{m-7}^{(4)}(y)\\
&\ +(-5y^4+5y^3+2y^2-3y-1)Z_{m-8}^{(4)}(y)+(-y^3+2y^2-2y+1)Z_{m-9}^{(4)}(y)\\
&\ +(y^5-y^4+2y^3-y^2+y)Z_{m-10}^{(4)}(y),
\end{aligned}
\]
并且生成函数 \(F_4(t,y)=\sum_{m\ge 0}Z_m^{(4)}(y)t^m\) 为有理函数 \(N_4(t,y)/D_4(t,y)\)，其中分母由递推系数给出
\[
D_4(t,y)=1-\sum_{j=1}^{10}a_j(y)t^j
\]
（\(a_j(y)\) 为上式右端的十个系数多项式），而分子可取显式整系数形式
\[
\begin{aligned}
N_4(t,y)=&\ 1+yt+(11-4y)t^2+(11y-4y^2)t^3+(11-18y+6y^2)t^4\\
&\ +(11y-18y^2+6y^3)t^5+(1-13y+9y^2-4y^3)t^6+(y-13y^2+9y^3-4y^4)t^7\\
&\ +y(y-1)(y^2-y+1)t^8+y^2(y-1)(y^2-y+1)t^9.
\end{aligned}
\]
特别地，取 \(y=1\)（忘却重量层，仅保留四阶碰撞矩）时，分母在 \(t\)-变量上精确析出不可约核因子
\[
2t^5-2t^4-7t^2-2t+1,
\]
并且该因子在 \(F_4(t,1)\) 中保持在位（不被分子约去），从而以代数不可约方式回收为四阶碰撞核的本征结构常数。
\end{conclusion}
\begin{proof}[证明要点]
有理性与递推由有限状态碰撞核的有限维线性实现给出；
分子分母的显式形式及 \(y=1\) 的核因子析出可由脚本
\texttt{scripts/exp\_fold\_zq\_weight\_collision\_partition\_audit.py} 在 \(\ZZ[y,t]\) 中精确化简并输出证书。证毕。
\end{proof}

\paragraph{四阶 moment-kernel 的分裂域与符号二次脊柱}
令
\[
P^{\mathrm{mom}}_{4}(t):=2t^{5}-2t^{4}-7t^{2}-2t+1\in\ZZ[t],
\qquad
F_4(t,1)=\frac{7t^2+1}{P^{\mathrm{mom}}_{4}(t)}.
\]

\begin{proposition}[四阶 moment-kernel 的满对称伽罗瓦性]\label{prop:xi-terminal-zm4-moment-kernel-s5}
\leanverified{paper\_xi\_terminal\_zm4\_moment\_kernel\_s5}
\(P^{\mathrm{mom}}_{4}\) 在 \(\QQ[t]\) 上不可约。记其分裂域为
\(
K:=\mathrm{Spl}\bigl(P^{\mathrm{mom}}_{4}\bigr)
\)，
则
\[
\Gal(K/\QQ)\cong S_{5},
\qquad
\Disc_{t}\!\bigl(P^{\mathrm{mom}}_{4}\bigr)=-2^{4}\cdot 985219.
\]
\end{proposition}
\begin{proof}
置
\(
p_7(x):=x^{5}-2x^{4}-7x^{3}-2x+2
\)。
直接验算得互反关系
\[
p_7(x)=x^{5}P^{\mathrm{mom}}_{4}(1/x).
\]
因此两者根集互为倒数，分裂域相同，且判别式相等。
由命题 \ref{prop:pom-a4-galois-s5} 得 \(\Gal(\mathrm{Spl}(p_7)/\QQ)\cong S_5\) 且 \(\Disc_x(p_7)=-2^4\cdot 985219\)，从而结论成立。证毕。
\end{proof}

\begin{proposition}[唯一二次子域与分歧刚性]\label{prop:xi-terminal-zm4-moment-kernel-unique-quadratic}
\leanverified{paper\_xi\_terminal\_zm4\_moment\_kernel\_unique\_quadratic}
在命题 \ref{prop:xi-terminal-zm4-moment-kernel-s5} 的记号下，
\(
K/\QQ
\)
的唯一非平凡正规子扩张为
\[
K^{A_5}=\QQ\!\left(\sqrt{\Disc_{t}(P^{\mathrm{mom}}_{4})}\right)=\QQ(\sqrt{-985219}).
\]
并且若 \(L/\QQ\) 为伽罗瓦扩张且在素数 \(985219\) 处不分歧，则 \(K\cap L=\QQ\)。
\end{proposition}
\begin{proof}
由 \(\Gal(K/\QQ)\cong S_5\) 及 \(A_5\triangleleft S_5\) 的唯一指数 \(2\) 正规性，\(K\) 的非平凡正规子扩张只能为 \(K^{A_5}\)，而符号表示对应的二次子域由判别式平方根生成，故第一条成立。
\par
设 \(L/\QQ\) 伽罗瓦，则交域 \(E:=K\cap L\) 亦为 \(\QQ\) 上伽罗瓦扩张，且 \(E\subseteq K\)。
由第一条，\(E\in\{\QQ,\QQ(\sqrt{-985219})\}\)。
若 \(E=\QQ(\sqrt{-985219})\)，则 \(L\) 必包含在 \(985219\) 处分歧的二次域，矛盾；故 \(E=\QQ\)。证毕。
\end{proof}

\begin{corollary}[三通道 Chebotarev 乘积律：\texorpdfstring{$S_5\times S_4\times S_3$}{S5xS4xS3} 的联合分布]\label{cor:xi-zm4-moment-kernel-leyang-s4-s3-s5-chebotarev-triple-product}
令 \(K=\mathrm{Spl}(P_4^{\mathrm{mom}})\)。
令
\(
L_{\mathrm{LY}}:=\mathrm{Spl}(P_{\mathrm{LY}})
\)
为 Lee--Yang 三次因子分裂域（结论 \ref{con:xi-terminal-zm-leyang-cubic-arithmetic}，\(\Gal(L_{\mathrm{LY}}/\QQ)\cong S_3\)）。
取有理参数 \(y_0=a/b\in\QQ\)（\(\gcd(a,b)=1,\ b>0\)），并令
\(
P_{a,b}(\lambda)=b^{2}\Pi(\lambda,a/b)\in\ZZ[\lambda]
\)
如结论 \ref{con:xi-terminal-zm-specialization-discriminant-ab}。
设其分裂域 \(L_4:=\mathrm{Spl}(P_{a,b})\) 满足 \(\Gal(L_4/\QQ)\cong S_4\)，并且
\[
985219\nmid \Disc_{\lambda}(P_{a,b}),
\qquad
\QQ\!\left(\sqrt{\Disc_{\lambda}(P_{a,b})}\right)\neq \QQ(\sqrt{-111}).
\]
则
\[
\Gal(KL_4L_{\mathrm{LY}}/\QQ)\ \cong\ S_{5}\times S_{4}\times S_{3}.
\]
因此对除去有限坏素集合后的素数 \(p\)，三域中 Frobenius 共轭类的联合分布严格乘积化：任取共轭类
\(
C_5\subset S_5
\)、\(
C_4\subset S_4
\)、\(
C_3\subset S_3
\)，满足
\[
\mathrm{Frob}_p(K)\in C_5,\ \mathrm{Frob}_p(L_4)\in C_4,\ \mathrm{Frob}_p(L_{\mathrm{LY}})\in C_3
\]
的素数集合具有自然密度
\[
\delta
=\frac{|C_5|}{|S_5|}\cdot\frac{|C_4|}{|S_4|}\cdot\frac{|C_3|}{|S_3|}.
\]
\end{corollary}
\begin{proof}
由命题 \ref{prop:xi-terminal-zm4-moment-kernel-unique-quadratic}，由于 \(985219\nmid \Disc_{\lambda}(P_{a,b})\)，得 \(K\cap L_4=\QQ\)；同理 \(K\cap L_{\mathrm{LY}}=\QQ\)。
另一方面，\(L_4/\QQ\) 为 \(S_4\)-伽罗瓦扩张，其唯一二次子域为 \(\QQ(\sqrt{\Disc_{\lambda}(P_{a,b})})\)；
而 \(L_{\mathrm{LY}}/\QQ\) 的唯一二次子域为 \(\QQ(\sqrt{-111})\)（结论 \ref{con:xi-terminal-zm-leyang-cubic-arithmetic}）。
由假设 \(\QQ(\sqrt{\Disc_{\lambda}(P_{a,b})})\neq \QQ(\sqrt{-111})\)，推出 \(L_4\cap L_{\mathrm{LY}}=\QQ\)。
于是三者两两线性互素，从而
\(
\Gal(KL_4L_{\mathrm{LY}}/\QQ)\cong \Gal(K/\QQ)\times \Gal(L_4/\QQ)\times \Gal(L_{\mathrm{LY}}/\QQ)
\)
并得到直积同构。
对不分歧素数应用 Chebotarev 密度定理，直积群上的共轭类为笛卡尔积，密度即为共轭类比例之乘积。证毕。
\end{proof}

\begin{theorem}[四阶碰撞矩的主极点常数与对数增长率]\label{thm:xi-zm4-moment-kernel-residue-asymptotics}
\leanverified{paper\_xi\_zm4\_moment\_kernel\_residue\_asymptotics}
令
\(
F_4(t,1)=\sum_{m\ge 0}a_m t^m
\)
如上，并令 \(r\in(0,1)\) 为 \(P^{\mathrm{mom}}_{4}\) 在正实轴上的最小模根，置 \(\rho:=r^{-1}\)。
则 \(t=r\) 为 \(F_4(t,1)\) 的简单极点，且存在常数 \(C>0\) 与 \(0<\alpha<\rho\) 使得
\[
a_m=C\,\rho^{m}+O(\alpha^{m})\qquad(m\to\infty),
\qquad
C=-\frac{7r^{2}+1}{r\,(P^{\mathrm{mom}}_{4})'(r)}.
\]
并且对数增长率极限存在并等于
\[
\kappa_{4}:=\lim_{m\to\infty}\frac1m\log a_m=\log\rho.
\]
\end{theorem}
\begin{proof}
由于 \(P^{\mathrm{mom}}_{4}\) 可分离且 \(r\) 为其根，\(F_4(t,1)\) 在 \(t=r\) 处为单极，其留数由
\(
\mathrm{Res}_{t=r}\frac{7t^2+1}{P_4^{\mathrm{mom}}(t)}=\frac{7r^2+1}{(P_4^{\mathrm{mom}})'(r)}
\)
给出。
按部分分式分解（或 Cauchy 系数公式拾取主极点留数）可得主项系数
\(
-[t^m](t-r)^{-1}=r^{-m-1}
\)，从而得到所示 \(C\) 的表达。
\par
由命题 \ref{prop:pom-s4-recurrence} 给出的五态非负整数 realization，\(F_4(t,1)\) 的极点集合与该矩阵 \(A_4\) 的特征值集合互为倒数。
由 Perron--Frobenius，Perron 根 \(\rho=\rho(A_4)\) 为单根且严格支配其余特征值模长；因此除 \(t=r=\rho^{-1}\) 外的极点模长均严格大于 \(r\)。
取 \(\alpha:=\max_{\mu\in\sigma(A_4)\setminus\{\rho\}}|\mu|\)，则 \(\alpha<\rho\)（其数值证书见注 \ref{rem:pom-collision-rh-break-a4-spectrum}），并给出误差项 \(O(\alpha^m)\)。
最后，由 \(a_m=C\rho^m(1+o(1))\) 直接推出 \(\frac1m\log a_m\to \log\rho\)。证毕。
\end{proof}


\begin{remark}[\texorpdfstring{$y=1$}{y=1} 的幽灵模态全消与 moment-kernel 的严格同态化]\label{rem:xi-terminal-zq-ghost-cancellation-y1}
记 \(F_q(t,y)=\sum_{m\ge 0}Z_m^{(q)}(y)t^m=N_q(t,y)/D_q(t,y)\)。
在 \(y=1\) 处出现统一的“幽灵模态”严格约去：
\begin{enumerate}
  \item \(q=2\)：\(N_2(t,1)=-(t-1)(t+1)^2\)，且 \(D_2(t,1)=-(t-1)(t+1)^2(2t^3-2t^2-2t+1)\)，故
  \[
  F_2(t,1)=\frac{1}{2t^3-2t^2-2t+1}.
  \]
  \item \(q=3\)：\(N_3(t,1)=-(t-1)(t+1)^2(2t^2+1)\)，且 \(D_3(t,1)=-(t-1)(t+1)^2(2t^3-4t^2-2t+1)\)，故
  \[
  F_3(t,1)=\frac{2t^2+1}{2t^3-4t^2-2t+1}.
  \]
  \item \(q=4\)：\(N_4(t,1)=-(t-1)(t+1)^2(t^2+1)(7t^2+1)\)，且
  \(D_4(t,1)=-(t-1)(t+1)^2(t^2+1)(2t^5-2t^4-7t^2-2t+1)\)，故
  \[
  F_4(t,1)=\frac{7t^2+1}{2t^5-2t^4-7t^2-2t+1}.
  \]
\end{enumerate}
因此在 \(y=1\) 的“碰撞投影”下，所有额外的 \(\pm1\) 与 \(\pm i\) 幽灵模态均被分子严格抵消，
留下的唯一动力学分母恰为 moment-kernel 的核多项式；该约去机制把重量细分与碰撞核谱系在生成函数层面严格同态化。
\end{remark}

\begin{conclusion}[谱分歧判别式的结构分解：普适因子与割圆因子的出现]\label{con:xi-terminal-zq-discriminant-factorization}
令 \(\Pi_q(\lambda,y)\) 表示 \(q=2,3,4\) 的递推特征多项式（即 \(D_q(t,y)\) 在 \(t=1/\lambda\) 下清分母后得到的 \(\lambda\)-多项式）。
则其关于 \(\lambda\) 的判别式在 \(\ZZ[y]\) 中具有严格分解：
\begin{enumerate}
  \item \(q=2\)：
  \[
  \mathrm{Disc}_{\lambda}(\Pi_2)=4\,y^{3}(y-1)\,P_{9}(y),
  \]
  其中
  \[
  \begin{aligned}
  P_{9}(y)=&\ 186624y^9-695632y^8+980411y^7-1121640y^6+1021286y^5\\
  &-458803y^4+127602y^3-35316y^2-10260y+6912.
  \end{aligned}
  \]
  \item \(q=3\)：
  \[
  \mathrm{Disc}_{\lambda}(\Pi_3)=-16\,y^{5}(y-1)^3\,P_{17}(y),
  \]
  其中 \(P_{17}(y)\in\ZZ[y]\) 为次数 \(17\) 多项式，且具有唯一负实根 \(y\approx-1.07316827\)。
  \item \(q=4\)：
  \[
  \mathrm{Disc}_{\lambda}(\Pi_4)=9216\,y^{7}(y-1)(y^2-y+1)\,R_{31}(y),
  \]
  其中 \(R_{31}(y)\in\ZZ[y]\) 为次数 \(31\) 多项式。
\end{enumerate}
因此除去普适因子 \(y^{2q-1}\) 与 \(y=1\) 的投影退化外，
\(q=4\) 首次强制出现割圆因子 \(y^2-y+1\)（即 \(6\) 次单位根），
表明四阶重量细分在谱分歧层面已不可约地感知到单位圆对称通道的额外算术结构。
\end{conclusion}
\begin{proof}[证明要点]
将 \(\Pi_q(\lambda,y)\) 视为 \(\lambda\)-多项式并计算判别式，在 \(\ZZ[y]\) 中作因式分解即可得到分解结构；
剩余因子 \(P_{17},R_{31}\) 的显式系数可由同一分解过程唯一确定，并由脚本
\texttt{scripts/exp\_fold\_zq\_weight\_collision\_partition\_audit.py} 导出为可审计工件。证毕。
\end{proof}

\begin{conclusion}[零重量纤维的完全闭式：\texorpdfstring{$d_m(0^m)=\lfloor (m+2)/2\rfloor$}{d_m(0^m)=floor((m+2)/2)}]\label{con:xi-terminal-zero-weight-fiber-closed}
记 \(0^m\in X_m\) 为全零稳定类型。
则其纤维大小满足对一切 \(m\ge 1\) 的精确公式
\[
d_m(0^m)=\Bigl\lfloor\frac{m+2}{2}\Bigr\rfloor.
\]
从而对任意整数 \(q\ge 1\) 有
\[
Z_m^{(q)}(0)=d_m(0^m)^q=\Bigl\lfloor\frac{m+2}{2}\Bigr\rfloor^q,
\]
给出重量细分在 \(k=0\) 层上的完全可解子扇区；该层的增长为多项式阶而非指数阶，对应一个严格的“边界零熵”极端相。
\end{conclusion}
\begin{proof}
由 Fold 的定义（见脚本 \texttt{scripts/common\_phi\_fold.py} 的 Zeckendorf 正规形口径），
\(\Fold_m(a)=0^m\) 当且仅当微态数值 \(N(a)\) 的 Zeckendorf 数码在位置 \(1,\dots,m\) 全为 \(0\)；
等价地 \(N(a)\in\{0,F_{m+2}\}\)。
其中 \(N(a)=0\) 仅对应全零微态，贡献 \(1\)。
另一方面，设 \(r_m\) 为用权集合 \(\{F_2,\dots,F_{m+1}\}\) 表示 \(F_{m+2}\) 的 \(0/1\) 子集表示数。
由于 \(\sum_{i=2}^{m}F_i=F_{m+2}-1\)，任意表示必须包含 \(F_{m+1}\)，故 \(r_m\) 等于用 \(\{F_2,\dots,F_m\}\) 表示 \(F_m\) 的表示数。
对后者再作“最大项是否取用”的分解：若取用 \(F_m\) 得到一项表示；否则由于 \(\sum_{i=2}^{m-2}F_i=F_m-1\)，表示必包含 \(F_{m-1}\)，余项为 \(F_{m-2}\)，从而表示数满足递推 \(r_m=1+r_{m-2}\)（\(m\ge 3\)），且 \(r_1=0,r_2=1\)。
解得 \(r_m=\lfloor m/2\rfloor\)，因此
\(
d_m(0^m)=1+r_m=\lfloor (m+2)/2\rfloor
\)。
代回 \(Z_m^{(q)}(0)\) 即得结论。证毕。
\end{proof}

\begin{conclusion}[最大重量层的稳定词分类与纤维谱闭式：\texorpdfstring{$|x|_1=\lceil m/2\rceil$}{|x|=ceil(m/2)}]\label{con:xi-terminal-max-weight-layer-fiber-spectrum-closed}
定义最大重量层
\[
X_m^{\max}:=\Bigl\{x\in X_m:\ |x|_1=\Bigl\lceil\frac{m}{2}\Bigr\rceil\Bigr\}.
\]
则：
\begin{enumerate}
  \item 若 \(m=2k+1\) 为奇数，则 \(X_{2k+1}^{\max}=\{(10)^k1\}\)，并且
  \[
  d_{2k+1}\bigl((10)^k1\bigr)=1.
  \]
  \item 若 \(m=2k\) 为偶数，则 \(|X_{2k}^{\max}|=k+1\)，且可取显式参数化
  \[
  x^{(r)}:=(10)^r(01)^{k-r}\qquad(0\le r\le k),
  \]
  其中 \((10)^r\) 表示长度 \(2r\) 的交替前缀 \(10\cdots10\)。
  该层上的纤维多重度满足
  \[
  d_{2k}\bigl(x^{(0)}\bigr)=1,\qquad
  d_{2k}\bigl(x^{(r)}\bigr)=k-r+1\ \ (1\le r\le k),
  \]
  等价地，
  \[
  \{d_{2k}(x):x\in X_{2k}^{\max}\}=\{1,1,2,3,\dots,k\}.
  \]
\end{enumerate}
\end{conclusion}
\begin{proof}
关于 \(X_m^{\max}\) 的分类仅依赖合法语法约束 \(x_ix_{i+1}=0\)：为达到 \(\lceil m/2\rceil\) 的最大 \(1\)-计数，除边界外每个 \(1\) 必被 \(0\) 隔开，故词必须为交替结构；当 \(m\) 为奇数时交替词唯一，偶数时恰允许出现一次 \(00\) 缺陷，从而得到上述 \(x^{(r)}\) 族。
\par
纤维多重度闭式使用 Zeckendorf 值函数 \(V_m(x)=\sum_{i=1}^m x_iF_{i+1}\) 与微态值 \(N(\omega)=\sum_{i=1}^m\omega_iF_{i+1}\)（见脚本 \texttt{scripts/common\_phi\_fold.py} 的口径）。当 \(x\in X_m\) 时，
\(\Fold_m(\omega)=x\) 当且仅当 \(N(\omega)\) 的 Zeckendorf 数码在长度 \(m\) 上与 \(x\) 相同；因此在本结论所述的最大重量词上（均满足 \(V_m(x)<F_{m+2}\)），纤维大小等于
\[
d_m(x)=\#\Bigl\{\omega\in\{0,1\}^m:\ N(\omega)=V_m(x)\Bigr\}.
\]
下证偶数情形的闭式；奇数情形可由同一“表示唯一性”论证推出 \(d_{2k+1}((10)^k1)=1\)。
\par
令 \(m=2k\) 且 \(0\le r\le k\)。由 \(x^{(r)}\) 的交替结构与两类求和恒等式
\[
\sum_{j=1}^{r}F_{2j}=F_{2r+1}-1,\qquad
\sum_{j=r+1}^{k}F_{2j+1}=F_{2k+2}-F_{2r+2},
\]
得到其 Zeckendorf 值的闭式
\[
V_{2k}\bigl(x^{(r)}\bigr)=\sum_{j=1}^{r}F_{2j}+\sum_{j=r+1}^{k}F_{2j+1}
=F_{2k+2}-F_{2r}-1.
\]
因此，若令
\[
\mathcal{R}_{k,r}:=\Bigl\{S\subseteq\{2,3,\dots,2k+1\}:\ \sum_{n\in S}F_n=F_{2k+2}-F_{2r}-1\Bigr\},
\]
则 \(d_{2k}(x^{(r)})=\card{\mathcal{R}_{k,r}}\)。
\par
\noindent\textbf{(1) \(r=0\) 时的唯一性。}
此时目标和为 \(F_{2k+2}-1\)。若不取最大项 \(F_{2k+1}\)，则最大可得和为
\(\sum_{n=2}^{2k}F_n=F_{2k+2}-2\)，不足以达到目标；故每个表示都必须取用 \(F_{2k+1}\)。
从而表示集合与 \(\mathcal{R}_{k-1,0}\) 一一对应，且以 \(k=1\) 为初值知 \(\card{\mathcal{R}_{k,0}}=1\)。
\par
\noindent\textbf{(2) \(1\le r\le k\) 时的递推与闭式。}
把 \(\mathcal{R}_{k,r}\) 按是否取用最大项 \(F_{2k+1}\) 分解。若 \(F_{2k+1}\in S\)，则余项必须表示
\[
F_{2k+2}-F_{2r}-1-F_{2k+1}=F_{2k}-F_{2r}-1,
\]
并且由于 \(F_{2k}-F_{2r}-1<F_{2k}\)，该表示不可能再取用 \(F_{2k}\)，从而与 \(\mathcal{R}_{k-1,r}\)（在权集合 \(\{2,\dots,2k-1\}\) 上）严格一一对应。
\par
若 \(F_{2k+1}\notin S\)，则 \(S\subseteq\{2,\dots,2k\}\)，而此时目标和相对于全取用之最大和的缺口为
\[
\Bigl(\sum_{n=2}^{2k}F_n\Bigr)-\bigl(F_{2k+2}-F_{2r}-1\bigr)=F_{2r}-1.
\]
因此不取用 \(F_{2k+1}\) 的表示与对缺口 \(F_{2r}-1\) 的“删去子集”一一对应。又由贪心强制律：
若 \(T\subseteq\{2,\dots,2r-1\}\) 满足 \(\sum_{n\in T}F_n=F_{2r}-1\)，则必须有 \(F_{2r-1}\in T\)，并递归推出
\(
T=\{3,5,\dots,2r-1\}
\) 唯一。
从而 \(\mathcal{R}_{k,r}\) 中恰存在唯一一个不含 \(F_{2k+1}\) 的表示。
\par
综上，对 \(1\le r\le k\) 有递推
\(
\card{\mathcal{R}_{k,r}}=\card{\mathcal{R}_{k-1,r}}+1
\)
并以边界初值 \(\card{\mathcal{R}_{r,r}}=1\)（对应 \(F_{2r+1}-1\) 的唯一表示）解得
\(
\card{\mathcal{R}_{k,r}}=k-r+1
\)。
这与 \(d_{2k}(x^{(r)})=\card{\mathcal{R}_{k,r}}\) 合并即得结论。
\end{proof}

\begin{conclusion}[重量—碰撞配分函数的顶层系数：最大重量层强制 Faulhaber 幂和]\label{con:xi-terminal-zq-top-coefficient-faulhaber}
对整数 \(q\ge 1\) 定义重量—碰撞配分函数
\[
Z_m^{(q)}(y):=\sum_{x\in X_m}y^{|x|_1}\,d_m(x)^q\in\ZZ[y].
\]
则其顶层系数（即 \(y^{\lceil m/2\rceil}\) 的系数）具有完全闭式：
\begin{enumerate}
  \item 若 \(m=2k+1\)，则 \([y^{k+1}]Z_{2k+1}^{(q)}(y)=1\)。
  \item 若 \(m=2k\)，则
  \[
  [y^{k}]Z_{2k}^{(q)}(y)=1+\sum_{j=1}^{k}j^{q}.
  \]
\end{enumerate}
\end{conclusion}
\begin{proof}
顶层系数等于最大重量层上幂和 \(\sum_{x\in X_m^{\max}}d_m(x)^q\)。由结论 \ref{con:xi-terminal-max-weight-layer-fiber-spectrum-closed} 的谱 \(\{1\}\)（奇数）与 \(\{1,1,2,\dots,k\}\)（偶数）直接得到两式。证毕。
\end{proof}

\begin{conclusion}[顶层扇区的生成函数：分母阶为 \texorpdfstring{$(1-t)^{q+2}$}{(1-t)^{q+2}}]\label{con:xi-terminal-max-weight-top-sector-ogf}
令
\(
b_k^{(q)}:=[y^{k}]Z_{2k}^{(q)}(y)=1+\sum_{j=1}^{k}j^{q}
\)
（\(q\ge 1\) 固定）。
则其普通生成函数满足严格恒等式
\[
\sum_{k\ge 0} b_k^{(q)}t^{k}
=\frac{1}{1-t}+\frac{1}{1-t}\sum_{j\ge 1}j^{q}t^{j}.
\]
从而 \(\sum_{k\ge 0} b_k^{(q)}t^{k}\) 为有理函数，且在 \(t=1\) 处具有阶恰为 \(q+2\) 的极点（因此分母可取为 \((1-t)^{q+2}\)）。
\end{conclusion}
\begin{proof}
由 \(b_k^{(q)}=1+\sum_{j=1}^k j^q\) 作求和交换：
\[
\sum_{k\ge 0}b_k^{(q)}t^k=\sum_{k\ge 0}t^k+\sum_{k\ge 0}\sum_{j=1}^{k}j^q t^k
=\frac{1}{1-t}+\sum_{j\ge 1}j^q\sum_{k\ge j}t^k,
\]
而 \(\sum_{k\ge j}t^k=t^j/(1-t)\) 得所示恒等式。
又 \(\sum_{j\ge 1}j^q t^j\) 为负整数阶 polylog，等价于一个以 \((1-t)^{-(q+1)}\) 为分母的有理函数；
乘以 \((1-t)^{-1}\) 后得到分母阶至多 \(q+2\)。
另一方面 \(b_k^{(q)}\sim k^{q+1}/(q+1)\)，故生成函数在 \(t=1\) 处的极点阶不能小于 \(q+2\)，从而阶恰为 \(q+2\)。证毕。
\end{proof}

\begin{conclusion}[硬边界事件的精确概率：\texorpdfstring{$H_m=\lceil m/2\rceil$}{H_m=ceil(m/2)} 的多项式前因子]\label{con:xi-terminal-hard-boundary-event-probability}
令 \(w_m(x)=d_m(x)/2^m\) 为均匀微态推前分布，取 \(X\sim w_m\)，并记 \(H_m:=|X|_1\)。
则最大重量事件 \(H_m=\lceil m/2\rceil\) 的概率具有闭式：
\begin{enumerate}
  \item \(m=2k+1\)：\(\mathbb{P}(H_m=k+1)=2^{-(2k+1)}\)；
  \item \(m=2k\)：\(\mathbb{P}(H_m=k)=\bigl(1+\tfrac{k(k+1)}{2}\bigr)\,2^{-2k}\)。
\end{enumerate}
并且当 \(k\to\infty\) 时
\[
\log \mathbb{P}(H_{2k}=k)= -2k\log 2 +2\log k +O(1),
\]
表明该硬边界处的指数速率固定为 \(\log 2\)，而前因子具有可计算的多项式阶 \(k^2\) 修正。
\end{conclusion}
\begin{proof}
由定义 \(\mathbb{P}(H_m=\ell)=2^{-m}\sum_{|x|_1=\ell}d_m(x)\)，右端即为 \(Z_m(y)\) 的 \(y^\ell\) 系数除以 \(2^m\)。
对 \(\ell=\lceil m/2\rceil\)，该系数等于结论 \ref{con:xi-terminal-zq-top-coefficient-faulhaber} 在 \(q=1\) 的顶层系数，代回即得两条闭式；渐近式由 \(\log(1+k(k+1)/2)=2\log k+O(1)\) 得到。证毕。
\end{proof}

\begin{conclusion}[最大重量层上的 \texorpdfstring{$q$}{q}-倾斜极限：\texorpdfstring{$J_k/k\Rightarrow \mathrm{Beta}(q+1,1)$}{Jk/k -> Beta(q+1,1)}]\label{con:xi-terminal-max-weight-escort-beta-limit}
在偶数长度 \(m=2k\) 上，考虑最大重量层条件下的 \(q\)-倾斜测度
\[
\mathbb{P}_{k,q}(x\in\cdot):=\frac{\sum_{x\in X_{2k}^{\max}} d_{2k}(x)^{q}\,\mathbf 1_{\{x\in\cdot\}}}{\sum_{x\in X_{2k}^{\max}} d_{2k}(x)^{q}}
\qquad(q\ge 1).
\]
令 \(X\sim\mathbb{P}_{k,q}\)，并设 \(J_k:=d_{2k}(X)\)。
则由结论 \ref{con:xi-terminal-max-weight-layer-fiber-spectrum-closed} 的谱 \(\{1,1,2,\dots,k\}\) 有
\[
\mathbb{P}(J_k=j)=\frac{j^{q}}{1+\sum_{i=1}^{k} i^{q}}\ (2\le j\le k),\qquad
\mathbb{P}(J_k=1)=\frac{2}{1+\sum_{i=1}^{k} i^{q}}.
\]
并且在尺度极限下有弱收敛
\[
\frac{J_k}{k}\ \Rightarrow\ \mathrm{Beta}(q+1,1),
\qquad \text{密度为 }(q+1)x^{q}\mathbf 1_{[0,1]}(x).
\]
\end{conclusion}
\begin{proof}
概率质量函数由谱的计数（\(j=1\) 出现两次，其余 \(j=2,\dots,k\) 各出现一次）直接得到。
对弱收敛，令 \(f\) 为 \([0,1]\) 上连续函数，则
\[
\mathbb{E}\Bigl[f\Bigl(\frac{J_k}{k}\Bigr)\Bigr]
=\frac{2\cdot f(1/k)\cdot 1^{q}+\sum_{j=2}^{k} f(j/k)\,j^{q}}{1+\sum_{i=1}^{k} i^{q}}
=\frac{\sum_{j=1}^{k} f(j/k)\,j^{q}}{\sum_{i=1}^{k} i^{q}}+o(1),
\]
其中 \(o(1)\) 来自 \(j=1\) 的双重计数在归一化下的可忽略性。
将分子分母同除以 \(k^{q+1}\)，并把和式视为 Riemann 和，即得极限
\(
\int_{0}^{1}f(x)x^{q}\,\dd x\,/\,\int_0^1 x^{q}\dd x
\)
对应的 \(\mathrm{Beta}(q+1,1)\) 分布。证毕。
\end{proof}

\begin{conclusion}[缺陷位置的尺度极限：\texorpdfstring{$R_k/k\Rightarrow \mathrm{Beta}(1,q+1)$}{Rk/k -> Beta(1,q+1)} 与重心比例]\label{con:xi-terminal-max-weight-defect-localization}
沿结论 \ref{con:xi-terminal-max-weight-layer-fiber-spectrum-closed} 的参数化 \(x^{(r)}=(10)^r(01)^{k-r}\)，令 \(R_k\) 为随机缺陷参数（即 \(X=x^{(R_k)}\) 的指标）。
则在结论 \ref{con:xi-terminal-max-weight-escort-beta-limit} 的 \(\mathbb{P}_{k,q}\) 下，有
\[
\frac{R_k}{k}\ \Rightarrow\ \mathrm{Beta}(1,q+1),
\qquad \text{密度为 }(q+1)(1-x)^q\mathbf 1_{[0,1]}(x).
\]
特别地其重心比例满足
\[
\frac{1}{k}\mathbb{E}[R_k]\to \frac{1}{q+2},
\qquad
\frac{1}{k}\mathbb{E}[d_{2k}(X)]\to \frac{q+1}{q+2}.
\]
\end{conclusion}
\begin{proof}
由结论 \ref{con:xi-terminal-max-weight-layer-fiber-spectrum-closed}，在 \(r\ge 1\) 上有确定关系
\(
J_k=d_{2k}(x^{(R_k)})=k-R_k+1
\)。
并且
\(
\mathbb{P}_{k,q}(R_k=0)=\frac{1}{1+\sum_{i=1}^{k}i^q}\to 0
\)。
因此对任意连续函数 \(f:[0,1]\to\RR\)，可作变量代换 \(j=k-r+1\) 得
\[
\mathbb{E}_{k,q}\Bigl[f\Bigl(\frac{R_k}{k}\Bigr)\Bigr]
=\frac{f(0)+\sum_{r=1}^{k} f(r/k)\,(k-r+1)^q}{1+\sum_{i=1}^{k}i^q}
=\frac{\sum_{j=1}^{k} f\!\bigl(1-j/k+1/k\bigr)\,j^q}{\sum_{i=1}^{k}i^q}+o(1).
\]
将分子分母同除以 \(k^{q+1}\) 并取 Riemann 和极限，得到
\[
\lim_{k\to\infty}\mathbb{E}_{k,q}\Bigl[f\Bigl(\frac{R_k}{k}\Bigr)\Bigr]
=\frac{\int_{0}^{1} f(1-x)\,x^q\,\dd x}{\int_{0}^{1}x^q\,\dd x}
=\int_{0}^{1} f(u)\,(q+1)(1-u)^q\,\dd u,
\]
即 \(R_k/k\Rightarrow \mathrm{Beta}(1,q+1)\)。
由于 \(R_k/k\) 与 \(J_k/k\) 均一致有界于 \([0,1]\)，重心比例由弱收敛与 Beta 分布均值公式直接推出。证毕。
\end{proof}

\begin{conclusion}[\texorpdfstring{$t=\pm1$}{t=±1} 直线的零轨迹：重量参数的代数曲线量子化]\label{con:xi-terminal-zq-line-tpm1-zeros}
设 \(D_q(t,y)\) 为 \(F_q(t,y)=N_q(t,y)/D_q(t,y)\) 的分母多项式。
则在 \(q=2,3,4\) 上均有严格等式
\[
D_q(1,y)=D_q(-1,y),
\]
并且其零点集合在 \(y\)-平面上量子化为有限代数集合：
\[
D_2(\pm1,y)=-y(y-1)(y-2),
\]
\[
D_3(\pm1,y)=y(y-1)(y^2-4y+6),
\]
\[
D_4(\pm1,y)=-y(y-1)(y^3-5y^2+12y-24).
\]
因此，“出现特征模态 \(\lambda=\pm1\)”仅在 \(y\) 落入上述代数集合时发生。
并且在上述三式中，除去共同因子 \(y(y-1)\) 后出现的附加因子在常数项处满足
\[
(y-2)\big|_{y=0}=-2,\qquad
(y^2-4y+6)\big|_{y=0}=6,\qquad
(y^3-5y^2+12y-24)\big|_{y=0}=-24,
\]
即呈现统一的阶乘尺度律 \(({-}1)^{q+1}q!\)（对 \(q=2,3,4\)）。
\end{conclusion}
\begin{proof}[证明要点]
把 \(t=\pm1\) 代入各 \(D_q(t,y)\) 并在 \(\ZZ[y]\) 中因式分解即可；
等式 \(D_q(1,y)=D_q(-1,y)\) 反映分母在 \(t\mapsto -t\) 的偶性结构在该三阶族中的稳定性。证毕。
\end{proof}



\begin{conclusion}[符号纤维配分函数的三阶递推：Tribonacci 常数给出主振荡模的模平方]\label{con:xi-terminal-signed-fiber-partition-recurrence}
令 \(\Fold_m:\Omega_m=\{0,1\}^m\to X_m\) 为标准折叠映射，并令纤维多重度 \(d_m(x):=|\Fold_m^{-1}(x)|\)。
定义符号纤维配分函数
\[
B_m^{\mathrm{sgn}}:=\sum_{x\in X_m}(-1)^{|x|_1}\,d_m(x)
=\sum_{a\in\Omega_m}(-1)^{|\Fold_m(a)|_1}.
\]
由结论 \ref{con:xi-terminal-zm-order4-rational} 取 \(y=-1\) 得 \(B_m^{\mathrm{sgn}}=Z_m(-1)\)，
并且 \(y(y+1)=0\) 使四阶可积递推在该点退化为三阶递推（见结论 \ref{con:xi-terminal-zm-hankel-rank4}）。
则 \(B_m^{\mathrm{sgn}}\) 满足三阶线性递推
\[
B_m^{\mathrm{sgn}}=B_{m-1}^{\mathrm{sgn}}-B_{m-2}^{\mathrm{sgn}}-B_{m-3}^{\mathrm{sgn}}\qquad(m\ge 4),
\]
并具有初值 \(B_1^{\mathrm{sgn}}=0,B_2^{\mathrm{sgn}}=0,B_3^{\mathrm{sgn}}=-2\)。
相应特征多项式为
\[
P(\lambda):=\lambda^3-\lambda^2+\lambda+1.
\]
设 \(\alpha,\overline{\alpha}\) 为其非实共轭根，并令 \(\tau\) 为 Tribonacci 常数（即 \(\tau^3-\tau^2-\tau-1=0\) 的唯一正实根）。
则
\[
\abs{\alpha}^2=\tau,
\qquad
\rho_2:=\abs{\alpha}=\sqrt{\tau}=1.3562030656262951\ldots .
\]
写 \(\alpha=\rho_2 e^{i\vartheta}\)（\(\vartheta\in(0,\pi)\)），则相位满足
\[
\cos\vartheta=\frac{\tau+1}{2\tau^{3/2}}.
\]
因此存在常数 \(C\neq 0\) 与 \(\psi\in\RR\) 使得
\[
B_m^{\mathrm{sgn}}=C\,\tau^{m/2}\cos(m\vartheta+\psi)+O(\tau^{-m})\qquad(m\to\infty),
\]
从而该符号通道在指数尺度上以 \(\tau^{m/2}\) 量级增长并携带非平凡旋转相位；该振荡模态与正配分矩族 \(\{S_q(m)\}_q\) 完全正交，构成可审计的频域指纹。
\end{conclusion}
\begin{proof}[证明]
由结论 \ref{con:xi-terminal-zm-order4-rational} 的四阶递推在 \(y=-1\) 处取值，并注意 \(y(y+1)=0\) 消去四阶项，
得到所示三阶递推；初值由 \(m\le 3\) 的直接核算给出。
特征多项式为 \(P(\lambda)=\lambda^3-\lambda^2+\lambda+1\)。
令 \(r\) 为其唯一实根，则 \(r=-1/\tau\)：代入 \(\lambda=-1/t\) 得 \(P(-1/t)=-(t^3-t^2-t-1)/t^3\)，故 \(t=\tau\)。
由根乘积 \(r\,\alpha\,\overline{\alpha}=-1\) 得 \(\abs{\alpha}^2=\tau\)，从而 \(\rho_2=\sqrt{\tau}\)；
又由根和 \(r+\alpha+\overline{\alpha}=1\) 得 \(2\rho_2\cos\vartheta=1-r\)，即 \(\cos\vartheta=(\tau+1)/(2\tau^{3/2})\)。
一般项的主振荡展开由根分解直接推出。证毕。
\end{proof}

\begin{remark}[判别式与二次分辨率子域]\label{rem:xi-terminal-signed-fiber-partition-discriminant}
三次特征多项式 \(P(\lambda)=\lambda^3-\lambda^2+\lambda+1\) 的判别式为 \(-44\)，
故其分裂域的唯一二次子域为 \(\QQ(\sqrt{-11})\)。
该虚二次域为 Heegner 数 \(11\) 对应的类数 \(1\) 情形，因此奇偶扭曲通道的最小算术外延被约束在 \(\QQ(\sqrt{-11})\) 上。
\end{remark}

\begin{conclusion}[奇偶偏差的指数衰减：\texorpdfstring{$\kappa_2=\log(2/\sqrt{\tau})$}{kappa2=log(2/sqrt(tau))}]\label{con:xi-terminal-parity-gap-kappa2}
令 \(w_m(x):=d_m(x)/2^m\) 为均匀微态推前分布，并令 \(X\sim w_m\)。
记 \(H_m:=|X|_1\) 为稳定类型的 \(1\)-计数随机变量。
则奇偶偏差满足严格恒等式
\[
\boxed{\;
\mathbb P(H_m\equiv 0\!\!\!\pmod 2)-\mathbb P(H_m\equiv 1\!\!\!\pmod 2)
=\EE\bigl[(-1)^{H_m}\bigr]
=\frac{B_m^{\mathrm{sgn}}}{2^m}\; }.
\]
特别地存在常数 \(\kappa_2>0\) 使得
\[
\abs{\mathbb P(H_m\equiv 0\!\!\!\pmod 2)-\mathbb P(H_m\equiv 1\!\!\!\pmod 2)}
=\exp\!\bigl(-m(\kappa_2+o(1))\bigr),
\qquad
\kappa_2=\log\!\Bigl(\frac{2}{\rho_2}\Bigr)=\log\!\Bigl(\frac{2}{\sqrt{\tau}}\Bigr)\approx 0.3884582488.
\]
\end{conclusion}
\begin{proof}
由定义
\(
\EE[(-1)^{H_m}]=\sum_{x\in X_m}(-1)^{|x|_1}w_m(x)=2^{-m}B_m^{\mathrm{sgn}}
\)
得到恒等式。
再由结论 \ref{con:xi-terminal-signed-fiber-partition-recurrence} 的显式展开可得
\(
\abs{B_m^{\mathrm{sgn}}}=\exp(m\log\rho_2+o(m))
\)，
从而
\(
\abs{B_m^{\mathrm{sgn}}}/2^m=\exp(-m(\log(2/\rho_2)+o(1)))
\)，并由 \(\rho_2=\sqrt{\tau}\) 得 \(\kappa_2\) 的闭式与数值。证毕。
\end{proof}

\begin{conclusion}[模 \(3\) Fourier 通道的 \(8\) 阶可逆递推：自反特征多项式]\label{con:xi-terminal-mod3-fourier-order8-recurrence}
令 \(\omega=e^{2\pi i/3}\)，并定义根单位扭曲配分
\[
C_m^{(3)}:=\sum_{x\in X_m}\omega^{|x|_1}\,d_m(x)
=2^m\,\EE\!\left[\omega^{H_m}\right].
\]
则整数序列
\[
a_m:=2\Re C_m^{(3)},\qquad b_m:=\frac{2}{\sqrt3}\Im C_m^{(3)}
\]
满足同一条 \(8\) 阶齐次整系数递推
\[
\boxed{\;
s_m
=2s_{m-1}-s_{m-2}-2s_{m-3}+s_{m-4}-2s_{m-5}-s_{m-6}+2s_{m-7}-s_{m-8}\;}
\]
其特征多项式为
\[
\Pi^{(3)}(\lambda)=\lambda^8-2\lambda^7+\lambda^6+2\lambda^5-\lambda^4+2\lambda^3+\lambda^2-2\lambda+1,
\]
并满足严格自反性 \(\lambda^8\Pi^{(3)}(1/\lambda)=\Pi^{(3)}(\lambda)\)。
令 \(\rho_3\) 为 \(\Pi^{(3)}\) 的谱半径，则 \(\rho_3\approx 1.7030650701\)，从而存在常数 \(c_3>0\) 使
\[
\abs{\EE\!\left[\omega^{H_m}\right]}\le c_3\Bigl(\frac{\rho_3}{2}\Bigr)^m,
\qquad
\abs{\EE\!\left[\omega^{2H_m}\right]}\le c_3\Bigl(\frac{\rho_3}{2}\Bigr)^m.
\]
\end{conclusion}
\begin{proof}[证明要点]
与结论 \ref{con:xi-terminal-signed-fiber-partition-recurrence} 同理，
\(C_m^{(3)}\) 是有限状态带权转移矩阵的矩阵系数，故满足某个有限阶常系数线性递推；
该递推可由有限窗口的精确枚举与最小递推恢复得到，并在本实例中稳定化为所示 \(8\) 阶形式。
自反性可由结论 \ref{con:xi-terminal-mod3-integral-symplectic} 的整辛结构推出（或直接检验系数回文性）。
谱半径 \(\rho_3\) 控制 \(\abs{C_m^{(3)}}\) 的指数级增长，从而给出 \(\abs{\EE[\omega^{H_m}]}\) 的衰减界。证毕。
\end{proof}

\begin{conclusion}[模 \(4\) Fourier 通道的 \(8\) 阶递推与更慢的混合]\label{con:xi-terminal-mod4-fourier-order8-recurrence}
令 \(i=\sqrt{-1}\)，并定义
\[
C_m^{(4)}:=\sum_{x\in X_m} i^{|x|_1}\,d_m(x)
=2^m\,\EE\!\left[i^{H_m}\right].
\]
则 \(\Re C_m^{(4)}\) 与 \(\Im C_m^{(4)}\) 同样满足一条 \(8\) 阶齐次整系数递推，其特征多项式为
\[
\Pi^{(4)}(\lambda)=\lambda^8-2\lambda^7-\lambda^6+4\lambda^5+\lambda^4-\lambda^2-2\lambda+2.
\]
令 \(\rho_4\) 为 \(\Pi^{(4)}\) 的谱半径，则 \(\rho_4\approx 1.8315468016\)，从而存在常数 \(c_4>0\) 使
\[
\abs{\EE\!\left[i^{H_m}\right]}\le c_4\Bigl(\frac{\rho_4}{2}\Bigr)^m,
\qquad
\abs{\EE\!\left[(-i)^{H_m}\right]}\le c_4\Bigl(\frac{\rho_4}{2}\Bigr)^m.
\]
该通道的衰减率严格慢于奇偶通道（因 \(\rho_4>\rho_2\)）。
\end{conclusion}
\begin{proof}[证明要点]
同上：根单位扭曲配分是有限维带权转移矩阵的矩阵系数，故满足有限阶整系数递推；
本实例中其最小递推稳定化为阶 \(8\)。
谱半径 \(\rho_4\) 给出 \(\abs{C_m^{(4)}}\) 的指数界，从而得到非平凡 Fourier 系数的衰减率。证毕。
\end{proof}

\begin{conclusion}[同余类的指数趋于均匀：由非平凡谱半径给出全变差界]\label{con:xi-terminal-congruence-mixing-tv-bound}
对 \(k\in\{2,3,4\}\)，记 \(\mu_{m,k}\) 为 \(H_m\bmod k\) 的分布，并记其 Fourier 系数
\[
\widehat{\mu}_{m,k}(j):=\EE\!\left[e^{2\pi i j H_m/k}\right]\qquad(1\le j\le k-1).
\]
令 \(\rho_2=\sqrt{\tau}\)，并令 \(\rho_3,\rho_4\) 如结论 \ref{con:xi-terminal-mod3-fourier-order8-recurrence} 与结论 \ref{con:xi-terminal-mod4-fourier-order8-recurrence}。
则存在常数 \(c_k>0\) 使得对所有 \(m\ge 1\) 与所有 \(1\le j\le k-1\) 有
\[
\abs{\widehat{\mu}_{m,k}(j)}\le c_k\Bigl(\frac{\rho_k}{2}\Bigr)^m.
\]
于是到均匀分布 \(U_k\) 的全变差距离满足
\[
\boxed{\;
\norm{\mu_{m,k}-U_k}_{\mathrm{TV}}
\le \frac12\sum_{j=1}^{k-1}\abs{\widehat{\mu}_{m,k}(j)}
\le C_k\Bigl(\frac{\rho_k}{2}\Bigr)^m\;}
\]
其中常数 \(C_k>0\) 仅依赖于 \(k\) 与有限初值数据。
\end{conclusion}
\begin{proof}
由有限阿贝尔群 \(\ZZ/k\ZZ\) 上的 Fourier 反演与三角不等式得到
\(
\norm{\mu_{m,k}-U_k}_{\mathrm{TV}}\le \frac12\sum_{j\ne 0}\abs{\widehat{\mu}_{m,k}(j)}
\)。
而当 \(k=2,3,4\) 时，各非平凡系数分别由
\(\EE[(-1)^{H_m}]\)、\(\EE[\omega^{H_m}]\)、\(\EE[i^{H_m}]\) 与其共轭给出，
其指数衰减由结论 \ref{con:xi-terminal-parity-gap-kappa2}、
\ref{con:xi-terminal-mod3-fourier-order8-recurrence}、
\ref{con:xi-terminal-mod4-fourier-order8-recurrence} 给出。证毕。
\end{proof}

\begin{conclusion}[同余混合的严格层级：模数加深导致主 Fourier 缺口收缩]\label{con:xi-terminal-congruence-mixing-hierarchy}
上述三通道满足严格序
\[
\rho_2<\rho_3<\rho_4<2,
\]
从而对应的主 Fourier 缺口（混合指数）严格递减：
\[
\log\Bigl(\frac{2}{\rho_2}\Bigr)\approx 0.3884582488\;>\;
\log\Bigl(\frac{2}{\rho_3}\Bigr)\approx 0.1607175705\;>\;
\log\Bigl(\frac{2}{\rho_4}\Bigr)\approx 0.0879863239.
\]
因此，Fold 的均匀微态推前在同余可观测量上呈现刚性的谱层级：奇偶类最先均匀化，其后为模 \(3\)，模 \(4\) 最慢。
\end{conclusion}

\begin{conclusion}[模 \(3\) 通道的整辛结构：伴随矩阵保辛并强制谱互反成对]\label{con:xi-terminal-mod3-integral-symplectic}
令 \(s_m\) 表示结论 \ref{con:xi-terminal-mod3-fourier-order8-recurrence} 的整数序列（可取 \(s_m=a_m\) 或 \(s_m=b_m\)），并定义其 companion 矩阵
\[
C=
\begin{pmatrix}
2&-1&-2&1&-2&-1&2&-1\\
1&0&0&0&0&0&0&0\\
0&1&0&0&0&0&0&0\\
0&0&1&0&0&0&0&0\\
0&0&0&1&0&0&0&0\\
0&0&0&0&1&0&0&0\\
0&0&0&0&0&1&0&0\\
0&0&0&0&0&0&1&0
\end{pmatrix},
\qquad \det C=1.
\]
则存在一条整的满秩反对称形式 \(\Omega\in \Mat_{8}(\ZZ)\)（\(\det\Omega=1\)）使得
\[
C^{\mathsf T}\Omega C=\Omega.
\]
可取
\(
\Omega=\begin{psmallmatrix}0&J\\ -J^{\mathsf T}&0\end{psmallmatrix}
\)
其中
\[
J=
\begin{pmatrix}
1&0&0&0\\
-2&1&0&0\\
1&-2&1&0\\
2&1&-2&1
\end{pmatrix}.
\]
因此 \(C\) 为某个整辛群 \(\mathrm{Sp}(8,\ZZ;\Omega)\) 的元素；其特征值必在 \(\lambda\mapsto 1/\lambda\) 下成对，
从而特征多项式 \(\Pi^{(3)}\) 必为自反多项式（与结论 \ref{con:xi-terminal-mod3-fourier-order8-recurrence} 一致）。
在生成函数层面，这等价于相应分母多项式满足严格的 \(z\mapsto 1/z\) 反演对称。
\end{conclusion}
\begin{proof}[证明要点]
直接代入所示 \(\Omega\) 计算可得 \(C^{\mathsf T}\Omega C=\Omega\)。
保辛性推出 \(C\) 的谱在取逆下成对，故特征多项式系数回文，即 \(\lambda^8\Pi^{(3)}(1/\lambda)=\Pi^{(3)}(\lambda)\)。证毕。
\end{proof}

\begin{conclusion}[顶端三层纤维谱的稳定相：缺口 Fibonacci 封口与偶数支达到者计数平台]\label{con:xi-terminal-top3-fiber-spectrum-stable-multiplicity-phase}
记 \(D_m^{(1)}>\!D_m^{(2)}>\!D_m^{(3)}\) 为 Fold 纤维多重度谱的前三个不同取值（定义 \ref{def:pom-top-fiber-spectrum}），并令 \(F_1=F_2=1\) 为 Fibonacci 数列。
则当 \(m=2k\ge 12\) 时，
\[
D_{2k}^{(1)}-D_{2k}^{(2)}=F_{k-4},
\qquad
D_{2k}^{(1)}-D_{2k}^{(3)}=F_{k-3},
\]
并且归一化缺口在 \(k\to\infty\) 时收敛到常数比例
\[
\lim_{k\to\infty}\frac{D_{2k}^{(1)}-D_{2k}^{(2)}}{D_{2k}^{(1)}}=\varphi^{-6},
\qquad
\lim_{k\to\infty}\frac{D_{2k}^{(1)}-D_{2k}^{(3)}}{D_{2k}^{(1)}}=\varphi^{-5}.
\]
此外，偶数支在阈值之后进入稳定多重度相：对一切偶数 \(m\ge 18\)，达到 \(D_m^{(1)},D_m^{(2)},D_m^{(3)}\) 的稳定类型数分别稳定为
\[
\#\{x:d_m(x)=D_m^{(1)}\}=2,\qquad
\#\{x:d_m(x)=D_m^{(2)}\}=6,\qquad
\#\{x:d_m(x)=D_m^{(3)}\}=4.
\]
\end{conclusion}
\begin{proof}[证明]
Fibonacci 缺口闭式由定理 \ref{thm:pom-second-max-fiber-closed-form} 与定理 \ref{thm:pom-third-max-fiber-even-closed-form}（并见注 \ref{rem:pom-top-fiber-spectrum-ladder}）直接给出。
由 \(D_{2k}^{(1)}=F_{k+2}\)（定理 \ref{thm:pom-max-fiber}）与 Binet 公式 \(F_n\sim \varphi^n/\sqrt5\) 得到归一化缺口的极限比例常数。
\par
达到者计数平台可由同一 residue DP/精确枚举口径逐 \(m\) 核验：在可审计窗口内其数值稳定出现在偶数支 \(m\ge 18\)，并与最大达到者稳定化定理 \ref{thm:pom-max-achievers-phase-stabilization}（偶数支为 \(2\)）相容。证毕。
\end{proof}

\begin{conclusion}[\texorpdfstring{$S_q$}{Sq}-对称压缩等于对称幂表示：\texorpdfstring{$B_q=\mathrm{Sym}^q(M)$}{Bq=Sym^q(M)} 且幺模可逆并给出显式逆矩阵]\label{con:xi-terminal-disjointness-bq-sympower-unimodular-inverse}
令 \(M=\bigl(\begin{smallmatrix}1&1\\[2pt]1&0\end{smallmatrix}\bigr)\)，并令 \(B_q\) 为不交图 \(\mathcal D_q\) 的 \(S_q\)-对称压缩算子（命题 \ref{prop:pom-disjointness-symmetric-compression-bq}）。
则在齐次多项式基
\(
e_j:=x^{q-j}y^j\ (0\le j\le q)
\)
下，\(B_q\) 严格等于对称幂表示 \(\mathrm{Sym}^q(M)\) 的矩阵：
\[
\boxed{\;B_q=\mathrm{Sym}^q(M)\;}
\qquad\text{且}\qquad
B_q(i,j)=\binom{q-j}{i}.
\]
进一步，令 \(\mathbf{P}_q\) 为 Pascal 矩阵 \(\mathbf{P}_q(i,j)=\binom{j}{i}\)（\(i\le j\)，否则为 \(0\)），并令 \(\mathbf{J}_q\) 为列反转置换矩阵 \(\mathbf{J}_q(i,j)=\mathbf{1}_{i=q-j}\)。
则有分解
\[
B_q=\mathbf{P}_q\,\mathbf{J}_q,
\]
从而 \(B_q\) 为整矩阵幺模（unimodular），并且
\[
\det(B_q)=\det(\mathbf{J}_q)=(-1)^{q(q+1)/2}.
\]
其逆矩阵亦完全显式：对 \(0\le i,j\le q\)，若 \(i+j<q\) 则 \((B_q^{-1})(i,j)=0\)；
若 \(i+j\ge q\) 则
\[
(B_q^{-1})(i,j)=(-1)^{i+j-q}\binom{j}{q-i}.
\]
\end{conclusion}
\begin{proof}[证明]
由命题 \ref{prop:pom-disjointness-symmetric-compression-bq} 已知 \(B_q(i,j)=\binom{q-j}{i}\)。
另一方面，对称幂表示满足：对任意 \(a,b,c,d\)，矩阵 \(\bigl(\begin{smallmatrix}a&b\\ c&d\end{smallmatrix}\bigr)\) 在基 \(e_j\) 上的作用为
\[
e_j(x,y)\longmapsto (ax+by)^{q-j}(cx+dy)^j.
\]
取 \(M=\bigl(\begin{smallmatrix}1&1\\ 1&0\end{smallmatrix}\bigr)\) 得
\[
e_j\longmapsto (x+y)^{q-j}x^j=\sum_{i=0}^{q-j}\binom{q-j}{i}x^{q-i}y^i,
\]
故 \(\mathrm{Sym}^q(M)(i,j)=\binom{q-j}{i}=B_q(i,j)\)，从而 \(B_q=\mathrm{Sym}^q(M)\)。
\par
分解 \(B_q=\mathbf{P}_q\mathbf{J}_q\) 由
\((\mathbf{P}_q\mathbf{J}_q)(i,j)=\mathbf{P}_q(i,q-j)=\binom{q-j}{i}\)
立即得到。
\(\mathbf{P}_q\) 为下三角且对角元全为 \(1\)，故 \(\det(\mathbf{P}_q)=1\)，从而 \(\det(B_q)=\det(\mathbf{J}_q)=(-1)^{q(q+1)/2}\)。
显式逆矩阵由 \(\mathbf{P}_q^{-1}\) 的标准闭式与 \(\mathbf{J}_q^{-1}=\mathbf{J}_q\) 合成得到，并可直接逐项核验 \(B_qB_q^{-1}=I\)。证毕。
\end{proof}

\begin{conclusion}[\texorpdfstring{$B_q^n$}{Bq^n} 的全显式闭式：所有层间传输系数为 Fibonacci 幂的非负整系数组合]\label{con:xi-terminal-bq-power-fibonacci-nonnegative-sum}
对任意整数 \(n\ge 1\)，有严格恒等式
\[
B_q^{\,n}=\mathrm{Sym}^q(M^n),
\qquad
M^n=\begin{pmatrix}F_{n+1}&F_n\\[2pt]F_n&F_{n-1}\end{pmatrix}.
\]
因此在同一基下，对一切 \(0\le i,j\le q\) 有完全显式条目公式
\[
\boxed{\;
\bigl(B_q^{\,n}\bigr)(i,j)
=\sum_{v=0}^{\min(i,j)}
\binom{q-j}{i-v}\binom{j}{v}\,
F_{n+1}^{\,q-j-i+v}\,
F_{n}^{\,i+j-2v}\,
F_{n-1}^{\,v}\; }.
\]
\end{conclusion}
\begin{proof}[证明]
由结论 \ref{con:xi-terminal-disjointness-bq-sympower-unimodular-inverse}，\(B_q=\mathrm{Sym}^q(M)\)。
对称幂表示是一个表示函子，故 \(\mathrm{Sym}^q(M^n)=\mathrm{Sym}^q(M)^n=B_q^{\,n}\)。
又由 Fibonacci 矩阵恒等式 \(M^n=\bigl(\begin{smallmatrix}F_{n+1}&F_n\\ F_n&F_{n-1}\end{smallmatrix}\bigr)\) 得到所示条目展开：
对基向量 \(e_j\) 作用为
\[
e_j\longmapsto (F_{n+1}x+F_ny)^{q-j}(F_nx+F_{n-1}y)^j,
\]
分别作二项式展开并收集 \(y^i\) 的系数即得求和式；各项系数均为非负整数，且仅含 Fibonacci 幂单项式。证毕。
\end{proof}

\begin{conclusion}[边界层无混合定律：\texorpdfstring{$B_q^n$}{Bq^n} 的第一行与最后一列直接重建 Fibonacci 幂谱]\label{con:xi-terminal-bq-boundary-nonmixing-law}
在结论 \ref{con:xi-terminal-bq-power-fibonacci-nonnegative-sum} 的条目公式中，边界层出现完全退化的无混合结构：
对任意 \(0\le j\le q\) 有
\[
\bigl(B_q^{\,n}\bigr)(0,j)=F_{n+1}^{\,q-j}F_n^{\,j},
\]
并且对任意 \(0\le i\le q\) 有
\[
\bigl(B_q^{\,n}\bigr)(i,q)=\binom{q}{i}F_n^{\,q-i}F_{n-1}^{\,i}.
\]
因此 \(B_q^{\,n}\) 在两条边界上逐点编码了 Fibonacci 的幂次谱（并带有精确二项式权），可作为对称幂动力的最短结构基准。
\end{conclusion}
\begin{proof}[证明]
在结论 \ref{con:xi-terminal-bq-power-fibonacci-nonnegative-sum} 的求和式中取 \(i=0\) 时只有 \(v=0\) 可能，立即得到第一式。
取 \(j=q\) 时 \(\binom{q-j}{i-v}=\binom{0}{i-v}\) 强制 \(v=i\)，代回即得第二式。证毕。
\end{proof}


% (User input window)

\subsubsection{跨接口终端结论：对角零温选择、\texorpdfstring{$(m,q)$}{(m,q)} 双尺度主项分离与“典型--极端”冻结判别链}\label{subsubsec:xi-diagonal-zero-temp-dualscale-freezing}
本段在 escort 族 \(w_{m,q}(x)\propto d_m(x)^q\) 的框架下，记录一组专用于
“对角标度 \(q=\tau m\)”的极端选择律与可审计标量。
在一般情形下，推论 \ref{cor:pom-escort-zero-temperature-superexp-concentration} 已给出
对固定 \(\varepsilon>0\) 的指数缺口集合 \(A_m(\varepsilon)\) 的 \(e^{-\Theta(m^2)}\) 超指数收缩；
这里进一步在\emph{顶端谱存在指数级隔离带}的假设下，把该收缩提升为对\emph{最大纤维集合}的集中，
并推导出 \(\log S_q(m)\) 的 \(m^2+m\) 两级主项分离、预算成功率的指数相图、以及
\((\alpha_\ast,g_\ast,h_1)\) 四量链式刚性。

\begin{conclusion}[对角零温选择原理：\texorpdfstring{$q=\tau m$}{q=tau m} 的 escort 在最大纤维集合上以 \(e^{-\Theta(m^2)}\) 集中]\label{con:xi-terminal-diagonal-zero-temp-maxset-concentration}
令 \(d_m:X_m\to\NN\) 为纤维多重度，记
\[
M_m:=\max_{x\in X_m}d_m(x),\qquad
A_m^{\max}:=\{x\in X_m:\ d_m(x)=M_m\},\qquad
K_m:=|A_m^{\max}|.
\]
令
\[
M_m^{(2)}:=\max\{d_m(x):x\in X_m,\ d_m(x)<M_m\}
\]
（若无则取 \(M_m^{(2)}:=1\)），并设存在严格指数间隙
\[
\delta_\ast:=\liminf_{m\to\infty}\frac{1}{m}\log\frac{M_m}{M_m^{(2)}}\ >\ 0.
\]
对 \(q\ge 0\) 定义 escort
\[
w_{m,q}(x):=\frac{d_m(x)^q}{S_q(m)},\qquad S_q(m):=\sum_{x\in X_m}d_m(x)^q.
\]
则对任意固定 \(\tau>0\)，取整数 \(q_m:=\lfloor \tau m\rfloor\)，有
\[
w_{m,q_m}(A_m^{\max})\ \ge\ 1-\exp\!\bigl(-\tau(\delta_\ast-o(1))\,m^2\bigr),
\]
并且到 \(A_m^{\max}\) 上均匀分布 \(U(A_m^{\max})\) 的全变差距离满足
\[
D_{\mathrm{TV}}\!\bigl(w_{m,q_m},U(A_m^{\max})\bigr)
\le \exp\!\bigl(-\tau(\delta_\ast-o(1))\,m^2\bigr).
\]
\end{conclusion}
\begin{proof}[证明]
对 \(x\notin A_m^{\max}\)，有 \(d_m(x)\le M_m^{(2)}\)，故
\[
\sum_{x\notin A_m^{\max}}d_m(x)^{q_m}\le |X_m|\,(M_m^{(2)})^{q_m}.
\]
又 \(|X_m|=F_{m+2}=\exp(O(m))\)。
由间隙假设可取 \(\varepsilon_m=o(1)\) 使 \(\log(M_m/M_m^{(2)})\ge (\delta_\ast-\varepsilon_m)m\)
（从而 \(M_m^{(2)}\le M_m\exp\bigl(-(\delta_\ast-\varepsilon_m)m\bigr)\)）并在大 \(m\) 上成立，故
\[
|X_m|\,(M_m^{(2)})^{q_m}
\le
\exp(O(m))\cdot M_m^{q_m}\exp\!\bigl(-(\delta_\ast-\varepsilon_m)q_m m\bigr)
=M_m^{q_m}\exp\!\bigl(-\tau(\delta_\ast-o(1))m^2\bigr).
\]
另一方面，
\(
\sum_{x\in A_m^{\max}}d_m(x)^{q_m}=K_m M_m^{q_m}\ge M_m^{q_m}
\)，从而
\[
1-w_{m,q_m}(A_m^{\max})
=\frac{\sum_{x\notin A_m^{\max}}d_m(x)^{q_m}}{S_{q_m}(m)}
\le
\frac{|X_m|\,(M_m^{(2)})^{q_m}}{M_m^{q_m}}
\le \exp\!\bigl(-\tau(\delta_\ast-o(1))m^2\bigr).
\]
在 \(A_m^{\max}\) 上 \(d_m\equiv M_m\)，故 \(w_{m,q_m}\) 在该集合上的条件分布严格均匀。
因此
\[
D_{\mathrm{TV}}\!\bigl(w_{m,q_m},U(A_m^{\max})\bigr)
=1-w_{m,q_m}(A_m^{\max}),
\]
第二式随之成立。证毕。
\end{proof}

\paragraph{固定阶幂和的冻结阈值与实参数前沿}
上文处理的是对角标度 \(q=\tau m\) 下的零温选择。
若将温度固定在有限阶 \(a\) 而仅令分辨率 \(m\to\infty\)，则最大纤维族与次大谱层之间的指数谱隙会产生另一条精确冻结机制；
其阈值由 \((\alpha_\ast,g_\ast,\alpha_\ast^{(2)})\) 三量闭式决定，并与 \S\ref{subsubsec:pom-real-q-pressure-thermo-frontier} 的实参数压力曲线相接。

\begin{theorem}[固定阶幂和的精确冻结阈值]\label{thm:xi-terminal-fixed-order-freezing-threshold}
\leanverified{paper\_xi\_terminal\_fixed\_order\_freezing\_threshold}
对整数 \(a\ge 1\) 定义
\[
S_a(m):=\sum_{x\in X_m}d_m(x)^a,
\qquad
P_a:=\lim_{m\to\infty}\frac{1}{m}\log S_a(m),
\]
并记
\[
M_m:=\max_{x\in X_m}d_m(x),\qquad
\mathcal M_m:=\{x\in X_m:\ d_m(x)=M_m\},\qquad
K_m:=|\mathcal M_m|.
\]
设极限
\[
\alpha_\ast:=\lim_{m\to\infty}\frac{1}{m}\log M_m,
\qquad
g_\ast:=\lim_{m\to\infty}\frac{1}{m}\log K_m
\]
存在，并令
\[
M_m^{(2)}:=\max\{d_m(x):\ x\in X_m\setminus\mathcal M_m\}
\]
（若补集为空则取 \(M_m^{(2)}:=1\)）以及
\[
\alpha_\ast^{(2)}:=\limsup_{m\to\infty}\frac{1}{m}\log M_m^{(2)}.
\]
若进一步满足
\[
\lim_{m\to\infty}\frac{1}{m}\log|X_m|=\log\varphi,
\qquad
\alpha_\ast^{(2)}<\alpha_\ast,
\]
则临界值
\[
a_c:=\frac{\log\varphi-g_\ast}{\alpha_\ast-\alpha_\ast^{(2)}}
\]
满足：对每个整数 \(a>a_c\)，有
\[
\boxed{\;P_a=a\alpha_\ast+g_\ast.\;}
\]
\end{theorem}
\begin{proof}
由 \(S_a(m)\ge K_mM_m^a\)，直接得到
\[
P_a\ge a\alpha_\ast+g_\ast.
\]
另一方面，
\[
S_a(m)=K_mM_m^a+\sum_{x\notin\mathcal M_m}d_m(x)^a
\le K_mM_m^a+|X_m|\,(M_m^{(2)})^a.
\]
于是
\[
\frac{1}{m}\log S_a(m)
\le
\frac{1}{m}\log\!\Bigl(K_mM_m^a+|X_m|\,(M_m^{(2)})^a\Bigr).
\]
对 \(m\to\infty\) 取上极限，并用
\[
\lim_{m\to\infty}\frac{1}{m}\log K_m=g_\ast,\qquad
\lim_{m\to\infty}\frac{1}{m}\log M_m=\alpha_\ast,
\]
\[
\lim_{m\to\infty}\frac{1}{m}\log|X_m|=\log\varphi,\qquad
\limsup_{m\to\infty}\frac{1}{m}\log M_m^{(2)}=\alpha_\ast^{(2)},
\]
得到
\[
P_a\le \max\{a\alpha_\ast+g_\ast,\ a\alpha_\ast^{(2)}+\log\varphi\}.
\]
若 \(a>a_c\)，则
\[
a\alpha_\ast+g_\ast>a\alpha_\ast^{(2)}+\log\varphi,
\]
故上式右端由第一项取得，从而 \(P_a\le a\alpha_\ast+g_\ast\)。
与下界合并即得
\(
P_a=a\alpha_\ast+g_\ast
\)。
\end{proof}

\begin{corollary}[固定阶 escort 对最大纤维族的指数集中]\label{cor:xi-terminal-fixed-order-escort-maxfiber-concentration}
\leanverified{paper\_xi\_terminal\_fixed\_order\_escort\_maxfiber\_concentration}
沿用定理 \ref{thm:xi-terminal-fixed-order-freezing-threshold} 的记号。
对整数 \(a\ge 1\) 定义 escort 测度
\[
\nu_{m,a}(x):=\frac{d_m(x)^a}{S_a(m)}\qquad(x\in X_m).
\]
若 \(\alpha_\ast^{(2)}<\alpha_\ast\) 且 \(a>a_c\)，则
\[
1-\nu_{m,a}(\mathcal M_m)\le \exp\!\bigl(-\gamma(a)\,m+o(m)\bigr),
\]
其中
\[
\gamma(a):=a(\alpha_\ast-\alpha_\ast^{(2)})-(\log\varphi-g_\ast)>0.
\]
\end{corollary}
\begin{proof}
由定义
\[
1-\nu_{m,a}(\mathcal M_m)
=
\frac{\sum_{x\notin\mathcal M_m}d_m(x)^a}{S_a(m)}
\le
\frac{|X_m|\,(M_m^{(2)})^a}{K_mM_m^a}.
\]
取对数并除以 \(m\)，得
\[
\frac{1}{m}\log\bigl(1-\nu_{m,a}(\mathcal M_m)\bigr)
\le
\frac{1}{m}\log|X_m|
+a\frac{1}{m}\log M_m^{(2)}
-\frac{1}{m}\log K_m
-a\frac{1}{m}\log M_m.
\]
令 \(m\to\infty\) 并应用定理 \ref{thm:xi-terminal-fixed-order-freezing-threshold} 中各极限，得到
\[
\limsup_{m\to\infty}\frac{1}{m}\log\bigl(1-\nu_{m,a}(\mathcal M_m)\bigr)
\le
\log\varphi+a\alpha_\ast^{(2)}-g_\ast-a\alpha_\ast
=-\gamma(a).
\]
由于 \(a>a_c\)，有 \(\gamma(a)>0\)，于是结论成立。
\end{proof}

\begin{conjecture}[冻结相的一段性]\label{conj:xi-terminal-single-transition-freezing-phase}
设 \S\ref{subsubsec:pom-real-q-pressure-thermo-frontier} 的实参数压力极限
\[
\tau(q):=\lim_{m\to\infty}\frac{1}{m}\log\sum_{x\in X_m}d_m(x)^q
\]
在 \([1,\infty)\) 上存在，并给出整数锚点 \(P_a=\tau(a)\)。
若 \(\alpha_\ast^{(2)}<\alpha_\ast\)，则存在唯一临界点
\[
q_c:=\frac{\log\varphi-g_\ast}{\alpha_\ast-\alpha_\ast^{(2)}}
\]
使得：
\begin{enumerate}
  \item \(\tau\) 在 \([1,q_c)\) 上严格凸且实解析；
  \item 对所有 \(q\ge q_c\)，有
  \[
  \tau(q)=q\alpha_\ast+g_\ast.
  \]
\end{enumerate}
换言之，实参数压力在高温区由整个纤维谱共同控制，而在低温区完全冻结到最大纤维族，且不存在多段冻结或再入相。
\end{conjecture}

\begin{conclusion}[双尺度对角矩展开：\texorpdfstring{$q=\tau m+s$}{q=tau m+s} 时 \texorpdfstring{$\log S_q(m)$}{log S_q(m)} 的 \(m^2+m\) 两级主项完全可分离]\label{con:xi-terminal-diagonal-logSq-m2-m-separation}
在结论 \ref{con:xi-terminal-diagonal-zero-temp-maxset-concentration} 的指数间隙假设下，进一步设极限
\[
\alpha_\ast:=\lim_{m\to\infty}\frac{1}{m}\log M_m,\qquad
g_\ast:=\lim_{m\to\infty}\frac{1}{m}\log K_m
\]
存在（亦可用 \(\limsup\) 解释）。
则对任意固定 \(\tau>0\) 与任意固定 \(s\in\RR\)，取实数序列 \(q_m:=\tau m+s\)，有
\[
\log S_{q_m}(m)=\tau\alpha_\ast\,m^2+\bigl(s\alpha_\ast+g_\ast\bigr)m+o(m).
\]
\end{conclusion}
\begin{proof}[证明]
由结论 \ref{con:xi-terminal-diagonal-zero-temp-maxset-concentration} 的估计（对实数 \(q_m\) 同样成立）可得
\[
S_{q_m}(m)=K_m M_m^{q_m}\Bigl(1+O\bigl(\exp\!\bigl(-\tau(\delta_\ast-o(1))m^2\bigr)\bigr)\Bigr).
\]
取对数并用 \(\log M_m=\alpha_\ast m+o(m)\)、\(\log K_m=g_\ast m+o(m)\) 即得
\[
\log S_{q_m}(m)=\log K_m+q_m\log M_m+o(m)
=\tau\alpha_\ast m^2+(s\alpha_\ast+g_\ast)m+o(m).
\]
证毕。
\end{proof}

\begin{remark}[关于整数化 \texorpdfstring{$q_m=\lfloor\tau m+s\rfloor$}{q=floor(tau m+s)} 的次主项]
若坚持取整数 \(q_m=\lfloor\tau m+s\rfloor\)，则 \(q_m=\tau m+s_m\) 且 \(s_m\in[s-1,s]\)。
因此上述分离仍保持 \(m^2\) 主项不变，而 \(m\) 级次主项中的 \(s\alpha_\ast m\) 需以 \(s_m\alpha_\ast m\) 替换；
该替换对应由取整引起的有界摆动，且不影响本段关于 \(m^2\) 级别选择律的结论链。
\end{remark}

\begin{conclusion}[对角 escort 的预言机相图：预算 \texorpdfstring{$B=\beta m$}{B=beta m} 的成功率在 \(m\) 尺度仅由 \texorpdfstring{$\alpha_\ast$}{alpha*} 决定]\label{con:xi-terminal-diagonal-escort-oracle-phase}
在结论 \ref{con:xi-terminal-diagonal-zero-temp-maxset-concentration} 的指数间隙假设下，定义在 escort \(w_{m,q}\) 下的预算 \(B\) 反演成功率
\[
\mathrm{Succ}_{m,q}(B)
:=\sum_{x\in X_m}w_{m,q}(x)\,\min\!\Bigl(1,\frac{2^{B}}{d_m(x)}\Bigr).
\]
取 \(q_m:=\lfloor\tau m\rfloor\)（\(\tau>0\) 固定）、\(B_m:=\lfloor\beta m\rfloor\)。
若 \(\alpha_\ast=\lim_{m\to\infty}\frac{1}{m}\log M_m\) 存在，则极限
\[
\lim_{m\to\infty}\frac{1}{m}\log \mathrm{Succ}_{m,q_m}(B_m)
=\min\bigl(0,\ \beta\log 2-\alpha_\ast\bigr)
\]
成立，并且当 \(\beta\log 2>\alpha_\ast\) 时，
\[
1-\mathrm{Succ}_{m,q_m}(B_m)\ \le\ \exp\!\bigl(-\tau(\delta_\ast-o(1))\,m^2\bigr).
\]
\end{conclusion}
\begin{proof}[证明]
由结论 \ref{con:xi-terminal-diagonal-zero-temp-maxset-concentration}，\(w_{m,q_m}\) 以 \(e^{-\Theta(m^2)}\) 集中在 \(A_m^{\max}\) 上。
在 \(A_m^{\max}\) 上 \(d_m(x)=M_m\)，故
\[
\mathrm{Succ}_{m,q_m}(B_m)
=\min\!\Bigl(1,\frac{2^{B_m}}{M_m}\Bigr)+O\bigl(\exp\!\bigl(-\tau(\delta_\ast-o(1))m^2\bigr)\bigr).
\]
取 \(\frac1m\log\) 并用 \(\log M_m=\alpha_\ast m+o(m)\) 得到指数极限
\(
\min(0,\beta\log 2-\alpha_\ast)
\)。
当 \(\beta\log2>\alpha_\ast\) 时，\(\frac{2^{B_m}}{M_m}\to\infty\)，失败只来自离开 \(A_m^{\max}\) 的概率，
从而由结论 \ref{con:xi-terminal-diagonal-zero-temp-maxset-concentration} 的质量界得到第二式。证毕。
\end{proof}

\begin{conclusion}[对角熵平台：当 \texorpdfstring{$q=\tau m+O(1)$}{q=tau m+O(1)} 时 escort 的 Shannon 与 Rényi 熵率均塌缩为 \texorpdfstring{$g_\ast$}{g*}]\label{con:xi-terminal-diagonal-escort-entropy-plateau}
在结论 \ref{con:xi-terminal-diagonal-zero-temp-maxset-concentration} 的指数间隙假设下，取 \(q_m=\tau m+s\)（\(\tau>0\) 固定，\(s\in\RR\) 固定）。
若 \(g_\ast=\lim_{m\to\infty}\frac{1}{m}\log K_m\) 存在，则对任意 Rényi 阶 \(r\in(0,\infty]\)，escort 熵满足
\[
\frac{1}{m}H_r\!\bigl(w_{m,q_m}\bigr)\longrightarrow g_\ast,
\]
特别地
\(
\frac{1}{m}H\!\bigl(w_{m,q_m}\bigr)\to g_\ast
\)
（其中 \(H=H_1\) 为 Shannon 熵，取自然对数）。
\end{conclusion}
\begin{proof}[证明]
记 \(p_m:=w_{m,q_m}(A_m^{\max})\)。
由结论 \ref{con:xi-terminal-diagonal-zero-temp-maxset-concentration}，有 \(1-p_m=\exp\!\bigl(-\tau(\delta_\ast-o(1))m^2\bigr)\)。
并且在 \(A_m^{\max}\) 上 \(w_{m,q_m}\) 严格均匀：对 \(x\in A_m^{\max}\)，\(w_{m,q_m}(x)=p_m/K_m\)。
\par
对任意 \(r\in(0,\infty)\)，有
\[
\sum_{x\in X_m}w_{m,q_m}(x)^r
=K_m\Bigl(\frac{p_m}{K_m}\Bigr)^r+\sum_{x\notin A_m^{\max}}w_{m,q_m}(x)^r
=p_m^rK_m^{1-r}+O\bigl((1-p_m)^r\bigr),
\]
从而
\(
H_r(w_{m,q_m})=\log K_m+o(m)
\)。
当 \(r=\infty\) 时，
\(
H_\infty(w_{m,q_m})=-\log\max_x w_{m,q_m}(x)=\log K_m-\log p_m=\log K_m+o(m)
\)；
当 \(r=1\) 时，由
\(
H(w_{m,q_m})=\log K_m+O\bigl((1-p_m)\log K_m\bigr)+o(m)
\)
亦得同一主项。
最后用 \(\log K_m=g_\ast m+o(m)\) 得结论。证毕。
\end{proof}

\begin{conclusion}[\texorpdfstring{$\beta_{\mathrm{crit}}=\alpha_\ast/\log 2$}{beta_crit=alpha*/log2} 的判别：当且仅当推前 \texorpdfstring{$w_m$}{w_m} 在指数尺度上冻结到最大纤维集合]\label{con:xi-terminal-typical-extreme-coaxial-freeze-criterion}
令推前分布 \(w_m(x)=d_m(x)/2^m\)，并定义
\[
a_1:=\lim_{m\to\infty}\frac{1}{m}\mathbb E_{w_m}[\log d_m(X)]\quad(\text{若极限存在}),
\qquad
\beta_{\mathrm{crit}}:=\frac{a_1}{\log2}.
\]
则总有 \(\beta_{\mathrm{crit}}\le \alpha_\ast/\log2\)（其中 \(\alpha_\ast=\lim\frac{1}{m}\log M_m\)）。
并且在结论 \ref{con:xi-terminal-diagonal-zero-temp-maxset-concentration} 的指数间隙假设下，
以下两命题等价：
\begin{enumerate}
  \item \(a_1=\alpha_\ast\)（即 \(\beta_{\mathrm{crit}}=\alpha_\ast/\log2\)）；
  \item \(\lim_{m\to\infty} w_m(A_m^{\max})=1\)。
\end{enumerate}
\end{conclusion}
\begin{proof}[证明]
由 \(\log d_m\le \log M_m\) 得 \(a_1\le \alpha_\ast\)，故 \(\beta_{\mathrm{crit}}\le \alpha_\ast/\log2\)。
\par
设指数间隙成立。
若 \(a_1=\alpha_\ast\)，则
\[
0\le \frac{1}{m}\mathbb E_{w_m}\!\Bigl[\log\frac{M_m}{d_m(X)}\Bigr]
=\frac{1}{m}\log M_m-\frac{1}{m}\mathbb E_{w_m}[\log d_m(X)]\ \to\ 0.
\]
而在补集 \(X_m\setminus A_m^{\max}\) 上有 \(\log(M_m/d_m)\ge \log(M_m/M_m^{(2)})\ge (\delta_\ast-o(1))m\)，
故
\[
\frac{1}{m}\mathbb E_{w_m}\!\Bigl[\log\frac{M_m}{d_m(X)}\Bigr]
\ge (\delta_\ast-o(1))\cdot\bigl(1-w_m(A_m^{\max})\bigr),
\]
从而 \(w_m(A_m^{\max})\to 1\)。
反向方向：若 \(w_m(A_m^{\max})\to 1\)，则 \(\log d_m(X)=\log M_m\) 以 \(w_m\)-概率趋于 \(1\)，
从而 \(\frac1m\mathbb E_{w_m}[\log d_m]\to \frac1m\log M_m\to\alpha_\ast\)，即 \(a_1=\alpha_\ast\)。证毕。
\end{proof}

\begin{conclusion}[极端--简并不等式：\texorpdfstring{$\alpha_\ast+g_\ast\le \log2$}{alpha*+g*<=log2} 的刚性与最大集合的指数承载率]\label{con:xi-terminal-extreme-degeneracy-inequality}
记 \(K_m:=|A_m^{\max}|\)、\(M_m:=\max_x d_m(x)\)。
由于 \(\sum_{x\in X_m}d_m(x)=2^m\)，恒有
\[
K_m M_m\ \le\ 2^m.
\]
若 \(\alpha_\ast=\lim_{m\to\infty}\frac{1}{m}\log M_m\) 与 \(g_\ast=\lim_{m\to\infty}\frac{1}{m}\log K_m\) 存在，则必满足
\[
\alpha_\ast+g_\ast\le\log2.
\]
并且推前质量满足
\[
w_m(A_m^{\max})=\frac{K_mM_m}{2^m},
\qquad
\frac{1}{m}\log w_m(A_m^{\max})\longrightarrow \alpha_\ast+g_\ast-\log2.
\]
因此 \(\alpha_\ast+g_\ast=\log2\) 当且仅当 \(w_m(A_m^{\max})\) 不以正指数率衰减（即 \(w_m(A_m^{\max})=\exp(o(m))\)）。
\end{conclusion}
\begin{proof}[证明]
不等式 \(K_mM_m\le\sum_x d_m(x)=2^m\) 直接给出 \(\alpha_\ast+g_\ast\le\log2\)。
而 \(w_m(A_m^{\max})\) 的恒等式来自 \(w_m(x)=d_m(x)/2^m\) 且在 \(A_m^{\max}\) 上 \(d_m\equiv M_m\)。
两边取 \(\frac1m\log\) 并取极限即得指数承载率公式与等号判别。证毕。
\end{proof}

\begin{conclusion}[四量链式刚性：\texorpdfstring{$g_\ast\le \log2-\alpha_\ast\le h_1\le \log\varphi$}{g*<=log2-alpha*<=h1<=logphi} 及其冻结相含义]\label{con:xi-terminal-four-quantity-chain-rigidity}
令
\[
h_1:=\lim_{m\to\infty}\frac{1}{m}H(w_m)
\quad(\text{若极限存在}),\qquad
\lim_{m\to\infty}\frac{1}{m}\log|X_m|=\log\varphi.
\]
则在上述极限存在时，
\[
g_\ast\le \log2-\alpha_\ast\le h_1\le \log\varphi.
\]
其中：
\begin{itemize}
  \item \(g_\ast\le \log2-\alpha_\ast\) 等价于结论 \ref{con:xi-terminal-extreme-degeneracy-inequality} 的 \(\alpha_\ast+g_\ast\le\log2\)；
  \item \(\log2-\alpha_\ast\le h_1\) 等价于 \(a_1\le \alpha_\ast\)，因为
  \(
  H(w_m)=m\log2-\mathbb E_{w_m}[\log d_m(X)]
  \)；
  \item \(h_1\le \log\varphi\) 为熵上界 \(H(w_m)\le \log|X_m|\) 的指数化形式。
\end{itemize}
并且两端等号分别刻画两种极端冻结口径：
\begin{itemize}
  \item \(g_\ast=\log2-\alpha_\ast\) 当且仅当 \(w_m(A_m^{\max})\) 不以正指数率衰减（结论 \ref{con:xi-terminal-extreme-degeneracy-inequality}）；
  \item 在结论 \ref{con:xi-terminal-diagonal-zero-temp-maxset-concentration} 的指数间隙假设下，
  \(\log2-\alpha_\ast=h_1\) 等价于 \(a_1=\alpha_\ast\)，并进一步等价于 \(w_m(A_m^{\max})\to 1\)（结论 \ref{con:xi-terminal-typical-extreme-coaxial-freeze-criterion}）。
\end{itemize}
\end{conclusion}
\begin{proof}[证明]
第一条由 \(K_mM_m\le 2^m\) 取对数除以 \(m\) 取极限得到。
第二条由恒等式
\(
H(w_m)=-\sum_x w_m(x)\log w_m(x)
=m\log2-\mathbb E_{w_m}[\log d_m(X)]
\)
直接推出：\(\log2-\alpha_\ast\le h_1\iff a_1\le\alpha_\ast\)。
第三条是 Shannon 熵上界 \(H(w_m)\le \log|X_m|\)。
等号刻画由对应不等式取等条件与前述结论链条给出。证毕。
\end{proof}

\paragraph{固定冻结相中的均匀塌缩、精确反演阈值、容量缺口与双指数分离}
在前段关于对角 escort 标度 \(q=\tau m\) 的零温选择之外，附录 \ref{app:fold-multiplicity} 中由固定矩阶 \(a\) 的冻结阈值、预言机容量闭式与容量缺口 Mellin 反演所支配的若干后果同样可直接纳入正文接口。除精确微态反演的尖锐比特阈值、连续容量缺口对全部实阶矩谱的无导数重建以及冻结相内部 \(g_*\) 与 \(\alpha_*\) 的双指数分离之外，还可得到 escort 测度向基态均匀测度的指数塌缩，以及全部 \(q\ge 1\) 的 Rényi 熵率平台。

\begin{theorem}[冻结相中的 escort 测度向基态均匀测度的指数塌缩]\label{thm:xi-terminal-fixed-freezing-tv-collapse}
设
\[
M_m:=\max_{x\in X_m}d_m(x),\qquad
A_m^{\max}:=\{x\in X_m:\ d_m(x)=M_m\},\qquad
K_m:=|A_m^{\max}|,
\]
\[
\alpha_*:=\lim_{m\to\infty}\frac{1}{m}\log M_m,\qquad
g_*:=\lim_{m\to\infty}\frac{1}{m}\log K_m,
\]
并令
\[
M_m^{(2)}:=\max\{d_m(x):\ d_m(x)<M_m\},\qquad
\alpha_*^{(2)}:=\limsup_{m\to\infty}\frac{1}{m}\log M_m^{(2)},
\]
其中若次大纤维不存在则取 \(M_m^{(2)}:=1\)。
若 \(\alpha_*^{(2)}<\alpha_*\)，并取固定实数
\[
a>a_{\mathrm c}:=\frac{\log\varphi-g_*}{\alpha_*-\alpha_*^{(2)}},
\]
则 escort 分布
\[
\pi_{m,a}(x):=\frac{d_m(x)^a}{S_a(m)},\qquad
S_a(m):=\sum_{x\in X_m}d_m(x)^a,
\]
与 \(A_m^{\max}\) 上均匀测度
\[
u_m^{\max}:=\mathrm{Unif}(A_m^{\max})
\]
满足精确总变差公式
\[
\boxed{\;
\|\pi_{m,a}-u_m^{\max}\|_{\mathrm{TV}}
=
1-\pi_{m,a}(A_m^{\max})
\le
\exp\!\bigl(-m\Delta(a)+o(m)\bigr),\;}
\]
其中
\[
\Delta(a):=a(\alpha_*-\alpha_*^{(2)})-(\log\varphi-g_*)>0.
\]
\leanverified{paper\_xi\_terminal\_fixed\_freezing\_tv\_collapse}
\end{theorem}
\begin{proof}[证明]
记 \(p_m:=\pi_{m,a}(A_m^{\max})\)。
由于 \(d_m(x)=M_m\) 在 \(A_m^{\max}\) 上恒定，\(\pi_{m,a}\) 在该集合上的条件分布恰为均匀测度 \(u_m^{\max}\)。
故存在支撑于 \(X_m\setminus A_m^{\max}\) 的概率测度 \(\nu_m\)，使得
\[
\pi_{m,a}=p_m u_m^{\max}+(1-p_m)\nu_m.
\]
由于两项支撑互不相交，
\[
\|\pi_{m,a}-u_m^{\max}\|_{\mathrm{TV}}
=
\frac12\sum_{x\in A_m^{\max}}|(p_m-1)u_m^{\max}(x)|
+\frac12\sum_{x\notin A_m^{\max}}(1-p_m)\nu_m(x)
=
1-p_m.
\]
另一方面，写
\[
S_a(m)=K_mM_m^a+R_a(m),
\qquad
R_a(m)\le |X_m|\,(M_m^{(2)})^a.
\]
于是
\[
1-p_m=\frac{R_a(m)}{S_a(m)}
\le
\frac{|X_m|\,(M_m^{(2)})^a}{K_mM_m^a}.
\]
对右端取对数并除以 \(m\)，再用
\[
\frac1m\log|X_m|\to\log\varphi,\qquad
\frac1m\log K_m\to g_*,\qquad
\frac1m\log M_m\to \alpha_*,\qquad
\limsup_{m\to\infty}\frac1m\log M_m^{(2)}=\alpha_*^{(2)},
\]
即得
\[
1-p_m\le \exp\!\bigl(-m\Delta(a)+o(m)\bigr).
\]
证毕。
\end{proof}

\begin{theorem}[冻结相中的 Rényi 熵率统一塌缩]\label{thm:xi-terminal-fixed-freezing-renyi-collapse}
沿用定理 \ref{thm:xi-terminal-fixed-freezing-tv-collapse} 的记号并固定 \(a>a_{\mathrm c}\)。
则对任意 \(q\in[1,\infty]\)，有
\[
\lim_{m\to\infty}\frac{1}{m}H_q(\pi_{m,a})=g_*,
\]
其中 \(H_1=H\) 为 Shannon 熵，且
\[
H_\infty(\pi):=-\log\max_x\pi(x).
\]
\end{theorem}
\begin{proof}[证明]
先证 \(q=\infty\)。
由
\[
S_a(m)=K_mM_m^a+R_a(m),\qquad
R_a(m)\le |X_m|\,(M_m^{(2)})^a,
\]
以及 \(\Delta(a)>0\)，得
\[
\frac{R_a(m)}{K_mM_m^a}\le \exp\!\bigl(-m\Delta(a)+o(m)\bigr),
\]
从而
\[
S_a(m)=K_mM_m^a(1+o(1)).
\]
于是
\[
\max_x\pi_{m,a}(x)=\frac{M_m^a}{S_a(m)}=\frac{1}{K_m}(1+o(1)),
\]
故
\[
H_\infty(\pi_{m,a})=\log K_m+o(m),
\]
从而 \(\frac1mH_\infty(\pi_{m,a})\to g_*\)。

再取固定 \(q>1\)。
由上式，对 \(x\in A_m^{\max}\) 有
\[
\pi_{m,a}(x)=\frac{1}{K_m}(1+o(1)).
\]
故
\[
\sum_{x\in A_m^{\max}}\pi_{m,a}(x)^q
=
K_m^{1-q}(1+o(1)).
\]
而对补集，由 \(S_a(m)\ge K_mM_m^a\) 得
\[
\sum_{x\notin A_m^{\max}}\pi_{m,a}(x)^q
\le
|X_m|
\left(\frac{(M_m^{(2)})^a}{K_mM_m^a}\right)^q
=
\exp\!\Bigl(
m\bigl[\log\varphi-qg_*-qa(\alpha_*-\alpha_*^{(2)})\bigr]
+o(m)
\Bigr).
\]
主项 \(K_m^{1-q}\) 的指数为 \(-(q-1)g_*m+o(m)\)。
两者的指数差满足
\[
\bigl[\log\varphi-qg_*-qa(\alpha_*-\alpha_*^{(2)})\bigr]+(q-1)g_*
=
\log\varphi-g_*-qa(\alpha_*-\alpha_*^{(2)})<0,
\]
因为 \(a>a_{\mathrm c}\) 且 \(q>1\)。
故非基态的 \(q\)-质量在指数尺度上可忽略，从而
\[
\sum_{x\in X_m}\pi_{m,a}(x)^q
=
K_m^{1-q}(1+o(1)).
\]
于是
\[
H_q(\pi_{m,a})
=
\frac{1}{1-q}\log\sum_x\pi_{m,a}(x)^q
=
\log K_m+o(m),
\]
故 \(\frac1mH_q(\pi_{m,a})\to g_*\)。

最后处理 \(q=1\)。
由定理 \ref{thm:xi-terminal-fixed-freezing-tv-collapse}，存在支撑于 \(X_m\setminus A_m^{\max}\) 的概率测度 \(\nu_m\)，使
\[
\pi_{m,a}=p_m u_m^{\max}+(1-p_m)\nu_m,
\qquad
1-p_m\le \exp\!\bigl(-m\Delta(a)+o(m)\bigr).
\]
记二元熵
\[
h(p):=-p\log p-(1-p)\log(1-p).
\]
由于两部分支撑不交，
\[
H(\pi_{m,a})=h(p_m)+p_m\log K_m+(1-p_m)H(\nu_m).
\]
再用
\[
H(\nu_m)\le \log|X_m|=m\log\varphi+o(m),
\]
以及 \(1-p_m\) 的指数小量，得到
\[
H(\pi_{m,a})=\log K_m+o(m).
\]
因此
\[
\frac{1}{m}H(\pi_{m,a})\to g_*.
\]
综上，对一切 \(q\in[1,\infty]\) 都有 \(\frac1mH_q(\pi_{m,a})\to g_*\)。
证毕。
\end{proof}

\begin{conclusion}[固定冻结相中精确微态反演的尖锐比特阈值]\label{con:xi-terminal-fixed-freezing-escort-oracle-sharp-threshold}
记
\[
S_a(m):=\sum_{x\in X_m}d_m(x)^a,\qquad
M_m:=\max_{x\in X_m}d_m(x),\qquad
A_m^{\max}:=\{x\in X_m:\ d_m(x)=M_m\},
\]
\[
K_m:=|A_m^{\max}|,\qquad
\alpha_*:=\lim_{m\to\infty}\frac{1}{m}\log M_m,\qquad
g_*:=\lim_{m\to\infty}\frac{1}{m}\log K_m,
\]
\[
M_m^{(2)}:=\max\{d_m(x):\ d_m(x)<M_m\},\qquad
\alpha_*^{(2)}:=\limsup_{m\to\infty}\frac{1}{m}\log M_m^{(2)},
\]
其中若不存在次大纤维则约定 \(M_m^{(2)}:=1\)，
并令
\[
a_{\mathrm c}:=\frac{\log\varphi-g_*}{\alpha_*-\alpha_*^{(2)}}.
\]
对固定 \(a>a_{\mathrm c}\)，定义 escort 分布
\[
\pi_{m,a}(x):=\frac{d_m(x)^a}{S_a(m)},
\]
以及在先验“\(x\sim \pi_{m,a}\)，再于纤维 \(\Fold_m^{-1}(x)\) 内均匀取微态”下、允许 \(B\) 比特外部寄存器时的最优精确反演成功率
\[
\mathrm{Succ}^{\mathrm{esc}}_{m,a}(B)
:=
\sum_{x\in X_m}\pi_{m,a}(x)\min\!\left(1,\frac{2^B}{d_m(x)}\right).
\]
则对任意 \(\varepsilon>0\) 与任意整数预算序列 \(B_m\ge 0\)，成立尖锐阈值
\[
B_m\le \frac{\alpha_*-\varepsilon}{\log 2}\,m
\quad\Longrightarrow\quad
\mathrm{Succ}^{\mathrm{esc}}_{m,a}(B_m)\longrightarrow 0,
\]
\[
B_m\ge \frac{\alpha_*+\varepsilon}{\log 2}\,m
\quad\Longrightarrow\quad
\mathrm{Succ}^{\mathrm{esc}}_{m,a}(B_m)\longrightarrow 1.
\]
\end{conclusion}
\begin{proof}[证明]
由定理 \ref{thm:xi-terminal-fixed-freezing-tv-collapse}，存在
\[
\Delta(a):=a(\alpha_*-\alpha_*^{(2)})-(\log\varphi-g_*)>0
\]
使
\[
\pi_{m,a}(A_m^{\max})
\ge
1-\exp\!\bigl(-m\Delta(a)+o(m)\bigr).
\]

若 \(x\in A_m^{\max}\)，则其纤维大小恰为 \(M_m\)，故在仅有 \(B_m\) 比特侧信息时，纤维内最优成功率至多为
\[
\min\!\left(1,\frac{2^{B_m}}{M_m}\right).
\]
于是
\[
\mathrm{Succ}^{\mathrm{esc}}_{m,a}(B_m)
\le
\pi_{m,a}(A_m^{\max})\frac{2^{B_m}}{M_m}
\;+\;
\bigl(1-\pi_{m,a}(A_m^{\max})\bigr).
\]
当
\(
B_m\le \frac{\alpha_*-\varepsilon}{\log 2}m
\)
时，由 \(M_m=\exp(\alpha_*m+o(m))\) 得
\[
\frac{2^{B_m}}{M_m}\le \exp(-\varepsilon m+o(m))\longrightarrow 0,
\]
而补集质量亦趋于 \(0\)，故 \(\mathrm{Succ}^{\mathrm{esc}}_{m,a}(B_m)\to 0\)。

反之，若
\(
B_m\ge \frac{\alpha_*+\varepsilon}{\log 2}m
\)，则
\[
2^{B_m}\ge \exp\!\bigl((\alpha_*+\varepsilon)m\bigr)\ge M_m
\]
对充分大 \(m\) 成立。于是可以在每条最大纤维上赋予互异码字并实现精确译码，从而
\[
\mathrm{Succ}^{\mathrm{esc}}_{m,a}(B_m)\ge \pi_{m,a}(A_m^{\max})\longrightarrow 1.
\]
证毕。
\end{proof}

\begin{conclusion}[连续容量缺口与全部实阶矩谱的 Mellin--Stieltjes 双对偶]\label{con:xi-terminal-capacity-gap-real-mellin-stieltjes-duality}
对固定 \(m\)，定义连续容量曲线与容量缺口
\[
\mathcal C_m^{\mathrm{cont}}(T):=\sum_{x\in X_m}\min\bigl(d_m(x),T\bigr),
\qquad
\mathcal R_m(T):=2^m-\mathcal C_m^{\mathrm{cont}}(T)
=\sum_{x\in X_m}(d_m(x)-T)_+,
\]
并令尾计数函数
\[
N_m(t):=\#\{x\in X_m:\ d_m(x)\ge t\}.
\]
则对任意实数 \(q>1\) 都有精确恒等式
\[
\boxed{\;
S_q(m)=q(q-1)\int_{0}^{\infty}T^{q-2}\mathcal R_m(T)\,\dd T.\;}
\]
特别地，\(\mathcal R_m(\cdot)\) 与 \(\mathcal C_m^{\mathrm{cont}}(\cdot)\) 都足以唯一决定全部实阶矩谱 \(\{S_q(m)\}_{q>1}\)；反过来，由定理 \ref{thm:fold-capacity-gap-mellin-stieltjes-duality} 的 Mellin 反演，整个半平面 \(\Re s>1\) 上的矩谱又唯一决定 \(\mathcal R_m\) 与 \(\mathcal C_m^{\mathrm{cont}}\)。
\end{conclusion}
\begin{proof}[证明]
由定义
\[
\mathcal C_m^{\mathrm{cont}}(T)=\int_{0}^{T}N_m(t)\,\dd t,
\]
故
\[
\mathcal R_m(T)
=
2^m-\int_{0}^{T}N_m(t)\,\dd t
=
\int_{T}^{\infty}N_m(t)\,\dd t.
\]
因此对 \(q>1\) 可交换积分次序：
\[
\begin{aligned}
q(q-1)\int_{0}^{\infty}T^{q-2}\mathcal R_m(T)\,\dd T
&=
q(q-1)\int_{0}^{\infty}T^{q-2}\left(\int_{T}^{\infty}N_m(t)\,\dd t\right)\dd T\\
&=
q(q-1)\int_{0}^{\infty}\left(\int_{0}^{t}T^{q-2}\,\dd T\right)N_m(t)\,\dd t\\
&=
q\int_{0}^{\infty}t^{q-1}N_m(t)\,\dd t\\
&=
S_q(m),
\end{aligned}
\]
最后一步正是命题 \ref{prop:fold-moment-mellin-stieltjes} 的实阶写法。于是 \(\mathcal R_m\) 确定全部 \(S_q(m)\)。
反向方向则由定理 \ref{thm:fold-capacity-gap-mellin-stieltjes-duality} 的 Mellin 反演直接给出。证毕。
\end{proof}

\begin{conclusion}[冻结相的双指数分离：\texorpdfstring{$g_*$}{g*} 与 \texorpdfstring{$\alpha_*$}{alpha*} 各司其职]\label{con:xi-terminal-fixed-freezing-two-exponent-separation}
在结论 \ref{con:xi-terminal-fixed-freezing-escort-oracle-sharp-threshold} 的设定下，若
\[
h_\infty(a):=\lim_{m\to\infty}\frac{1}{m}\Bigl(-\log\max_{x\in X_m}\pi_{m,a}(x)\Bigr)
\]
存在，并定义
\[
\beta_{\mathrm{inv}}(a)
:=
\inf\Bigl\{
\beta\ge 0:\ \exists\, B_m=\beta m+o(m)\ \text{使}\ 
\mathrm{Succ}^{\mathrm{esc}}_{m,a}(B_m)\to 1
\Bigr\},
\]
则对每个固定 \(a>a_{\mathrm c}\) 都有
\[
\boxed{\;
h_\infty(a)=g_*,
\qquad
\beta_{\mathrm{inv}}(a)=\frac{\alpha_*}{\log 2}.\;}
\]
因此，在固定冻结相中，escort 统计的 min-entropy 率由最大层简并指数 \(g_*\) 锁定，而精确微态反演的最小线性比特率则由最大纤维指数 \(\alpha_*\) 锁定；两者对应的是冻结相内部两条互不吸收的指数尺度。
\end{conclusion}
\begin{proof}[证明]
第一式是定理 \ref{thm:xi-terminal-fixed-freezing-renyi-collapse} 在 \(q=\infty\) 时的特例。

对第二式，任取 \(\varepsilon>0\)。
若 \(\beta<\alpha_*/\log 2\)，则对充分小的 \(\varepsilon\) 有
\[
\beta\le \frac{\alpha_*-\varepsilon}{\log 2}.
\]
于是任何满足 \(B_m=\beta m+o(m)\) 的预算序列，都满足结论 \ref{con:xi-terminal-fixed-freezing-escort-oracle-sharp-threshold} 的下侧条件，从而
\(
\mathrm{Succ}^{\mathrm{esc}}_{m,a}(B_m)\to 0
\)；
故这类 \(\beta\) 不可能进入定义 \(\beta_{\mathrm{inv}}(a)\) 的可行集合。

反之，若 \(\beta>\alpha_*/\log 2\)，则可取 \(\varepsilon>0\) 使
\[
\beta\ge \frac{\alpha_*+\varepsilon}{\log 2}.
\]
对任意满足 \(B_m=\beta m+o(m)\) 的预算序列，结论 \ref{con:xi-terminal-fixed-freezing-escort-oracle-sharp-threshold} 的上侧条件给出
\(
\mathrm{Succ}^{\mathrm{esc}}_{m,a}(B_m)\to 1
\)。
故所有 \(\beta>\alpha_*/\log 2\) 都可行，而所有 \(\beta<\alpha_*/\log 2\) 都不可行，只能有
\[
\beta_{\mathrm{inv}}(a)=\frac{\alpha_*}{\log 2}.
\]
证毕。
\end{proof}

% (User input window)

\paragraph{跨接口终端结论：黄金均值有限体积的再生分解、谱律的强余项与 Toeplitz 证书的数值/预算鲁棒性}
本段把 Fibonacci 立方（黄金均值 SFT）的有限体积组合结构、混合边界离散谱律的强余项界、
以及视界 Toeplitz--PSD 证书的数值鲁棒性与信息预算阈值，压缩为可独立引用的结论条目。
证明仅指向相应已完成的定理/命题/推论，不再重复推导细节。

\begin{conclusion}[黄金均值有限体积的严格 DLR--再生分解：任意 \(k\)-点函数在“被迫零”分割下完全独立]\label{con:xi-terminal-golden-mean-dlr-regeneration-factorization}
在 Fibonacci 立方图 \(\Gamma_n\) 上取均匀计数，并记合法词集合 \(X_n:=V(\Gamma_n)\)（定义 \ref{def:pom-fibonacci-cube}）。
对任意指标集 \(S=\{i_1<\cdots<i_k\}\subseteq\{1,\dots,n\}\)，若存在 \(j\) 使 \(i_{j+1}-i_j=1\)，则约束不可行且
\[
\#\{v\in X_n:\ v_{i_j}=1\ \forall j\}=0.
\]
若满足间隔条件 \(i_{j+1}-i_j\ge 2\)，则存在严格乘积分解
\[
\#\{v\in X_n:\ v_{i_j}=1\ \forall j\}
=F_{i_1}\Bigl(\prod_{j=1}^{k-1}F_{i_{j+1}-i_j-1}\Bigr)F_{n-i_k+1},
\qquad |X_n|=F_{n+2}.
\]
因此在有限体积 Gibbs/DLR 口径下，条件事件 \(\{v_{i_j}=1\}\) 在每个 \(i_j\) 两侧强制邻位为 \(0\)，并把系统分裂为 \(k+1\) 段互不相交的自由区间；
在该条件下，各自由区间上的配置严格独立，且每段均为相应长度黄金均值 SFT 的均匀分布。
\end{conclusion}
\begin{proof}[证明]
见定理 \ref{thm:pom-fibcube-kpoint-count}；其中“邻位被迫为 \(0\)”与自由段分裂正是乘积分解的结构来源，从而条件独立性亦随之成立。证毕。
\end{proof}

\begin{conclusion}[有限体积 \texorpdfstring{$\Rightarrow$}{=>} Parry 体相的指数速率：误差由次特征值 \texorpdfstring{$-\varphi^{-2}$}{-phi^{-2}} 单独锁定]\label{con:xi-terminal-unif-xn-to-parry-exp-rate}
令 \(\nu_{\mathrm{Parry}}\) 为黄金均值 SFT 的 Parry 最大熵平衡态（定理 \ref{thm:pom-parry-limit-chain-explicit}），其两态 Markov 转移矩阵的唯一非平凡特征值为 \(\lambda_2=-\varphi^{-2}\)。
对均匀合法词 \(U_n:=\mathrm{Unif}(X_n)\)，取任意固定块长 \(k\ge 1\)。
则存在常数 \(C_k\) 使对任意位置 \(t\) 与
\(
d:=\min\{t-1,\ n-t-k+1\}
\)
均成立指数型体相收敛界：
\[
\sup_{a_1\cdots a_k\in X_k}\Bigl|
U_n(X_t\cdots X_{t+k-1}=a_1\cdots a_k)-\nu_{\mathrm{Parry}}([a_1\cdots a_k])
\Bigr|
\le C_k\,\varphi^{-2d}.
\]
因此有限体积到体相的相关长度与边界影响衰减率由同一代数常数 \(\varphi^{-2}=|\lambda_2|\) 完全锁定。
\end{conclusion}
\begin{proof}[证明]
局部统计收敛到 Parry 测度见命题 \ref{prop:pom-unif-xn-local-limit-parry}。
将块柱集计数写为黄金均值邻接矩阵 \(A\) 的幂矩阵元的乘积，并对左右补全段分别施行 Perron--Frobenius 分解，可把误差项精确压缩到次模态的贡献；
由于 \(A\) 的次特征值为 \(-\varphi^{-1}\)，故归一化后误差在距离边界 \(d\) 的意义下按 \((|\varphi^{-1}|/\varphi)^d=\varphi^{-2d}\) 衰减。
把有限个合法块的常数并入 \(C_k\) 即得所示上界。证毕。
\end{proof}

\begin{conclusion}[纤维立方几何的边数与平均度闭式：平均度线性系数精确为 \texorpdfstring{$\frac{2}{5}(3-\varphi)$}{(2/5)(3-phi)}]\label{con:xi-terminal-fibcube-edge-avg-degree-closed}
Fibonacci 立方 \(\Gamma_n\) 的边数与顶点数满足闭式
\[
|E(\Gamma_n)|=\frac{nF_{n+1}+2(n+1)F_n}{5},
\qquad
|V(\Gamma_n)|=F_{n+2},
\]
从而平均度
\[
\overline{\deg}(\Gamma_n)
=\frac{2|E(\Gamma_n)|}{|V(\Gamma_n)|}
=\frac{2}{5}\left(
n\frac{F_{n+1}}{F_{n+2}}+2(n+1)\frac{F_n}{F_{n+2}}
\right)
\sim \frac{2}{5}(3-\varphi)\,n.
\]
进一步，若纤维重构图 \(\mathcal F_x\) 具有 Cartesian 分解
\(
\mathcal F_x\cong\prod_{i=1}^r\Gamma_{\ell_i}
\)
（定理 \ref{thm:pom-fiber-reconstruction-cartesian-product}），则有边数与平均度的完全闭式
\[
|E(\mathcal F_x)|
=\sum_{i=1}^r |E(\Gamma_{\ell_i})|\prod_{j\ne i}|V(\Gamma_{\ell_j})|,
\qquad
\overline{\deg}(\mathcal F_x)=\sum_{i=1}^r \overline{\deg}(\Gamma_{\ell_i}).
\]
因此纤维的“局部可开关性”规模在体相极限中严格线性定标，并由同一 \(\frac{2}{5}(3-\varphi)\) 系数锚定（对长分量贡献可加）。
\end{conclusion}
\begin{proof}[证明]
\(\Gamma_n\) 的边数闭式与平均度公式见推论 \ref{cor:pom-fibcube-edge-closed-form}。
纤维的 Cartesian 分解见定理 \ref{thm:pom-fiber-reconstruction-cartesian-product}。
对 Cartesian 乘积图，边集按“选择一个因子中的一条边并固定其余因子顶点”分解，遂得边数公式；平均度可由
\(
2|E|/|V|
\)
与顶点数乘法性推出可加性。证毕。
\end{proof}

\begin{conclusion}[离散 Weyl 定律的强版本：计数函数的全区间误差严格不超过 \(1\)，故 Kolmogorov 距离为 \(O(1/k)\)]\label{con:xi-terminal-discrete-weyl-strong-remainder}
对谱参数 \(\mu\in(0,4)\) 令 \(\theta(\mu)=\arccos(1-\mu/2)\in(0,\pi)\)，并定义计数函数
\(
N_k(\mu)=\#\{p\le k:\mu_p\le\mu\}
\)。
则存在完全显式反函数型公式
\[
N_k(\mu)=\Bigl\lfloor \frac{2k+1}{2\pi}\theta(\mu)+\frac12\Bigr\rfloor,
\qquad
\Bigl|N_k(\mu)-\frac{2k+1}{2\pi}\theta(\mu)\Bigr|\le 1
\quad(\forall\mu\in(0,4)).
\]
因此经验分布函数 \(F_k(\mu):=N_k(\mu)/k\) 在全区间上以 \(O(1/k)\) 精度一致逼近极限分布函数 \(F(\mu)=\theta(\mu)/\pi\)。
\end{conclusion}
\begin{proof}[证明]
见推论 \ref{cor:pom-Lk-discrete-weyl}。证毕。
\end{proof}

\begin{conclusion}[边界谱律是体相谱律的严格二次倾斜：\texorpdfstring{$\dd\rho(\mu)=\frac{\mu(4-\mu)}{2}\,\dd\nu(\mu)$}{d rho=(mu(4-mu)/2)d nu}]\label{con:xi-terminal-boundary-bulk-spectral-quadratic-tilt}
设 \(\nu\) 为体相极限测度，\(\rho\) 为边界极限测度。
则二者满足精确 Radon--Nikodym 关系
\[
\dd\rho(\mu)=\frac{\mu(4-\mu)}{2}\,\dd\nu(\mu),
\qquad
\int f(\mu)\,\dd\rho(\mu)=\frac12\int \mu(4-\mu)\,f(\mu)\,\dd\nu(\mu).
\]
因此边界谱统计的全部可积泛函可在体相测度下以一次加权换测精确计算。
\end{conclusion}
\begin{proof}[证明]
见推论 \ref{cor:pom-Lk-boundary-bulk-radon-nikodym}。证毕。
\end{proof}

\begin{conclusion}[Toeplitz--PSD 的数值可证书鲁棒性：谱裕度与系数误差给出完全显式的可行域]\label{con:xi-terminal-toeplitz-psd-numerical-robustness}
设 \(\mathbf{T}_N=(c_{j-k})_{0\le j,k\le N}\) 为 Toeplitz 截断矩阵，\(\widehat{\mathbf{T}}_N=(\widehat c_{j-k})\) 为任意 Hermitian 近似。
若存在可审计谱裕度 \(\lambda_{\min}(\widehat{\mathbf{T}}_N)\ge\eta>0\) 且截断误差满足
\(
\|\widehat{\mathbf{T}}_N-\mathbf{T}_N\|_{\mathrm{op}}\le \eta/2
\)，
则必有 \(\mathbf{T}_N\succeq 0\)。
并且若系数逐项误差满足 \(\max_{0\le n\le N}|\widehat c_n-c_n|\le\varepsilon\)，则
\[
\|\widehat{\mathbf{T}}_N-\mathbf{T}_N\|_{\mathrm{op}}
\le\sum_{k=-N}^{N}|\widehat c_k-c_k|
\le (2N+1)\varepsilon.
\]
因此有限阶 Toeplitz--PSD 核验可被约化为“计算 \(\widehat{\mathbf{T}}_N\) 的最小特征值并给出可审计误差界”的数值—代数混合证书链。
\end{conclusion}
\begin{proof}[证明]
见定理 \ref{thm:xi-spectral-margin-finite-psd-certificate}。证毕。
\end{proof}

\begin{theorem}[Toeplitz--PSD 的鲁棒共终端审计：误差封口与谱裕度向下传播]\label{thm:xi-toeplitz-psd-robust-cofinal-audit}
\leanverified{paper\_xi\_toeplitz\_psd\_robust\_cofinal\_audit}
沿用 Toeplitz 主子矩阵 \(\mathbf{T}_N=(c_{j-k})_{0\le j,k\le N}\) 的记号。
令 \(\{N_k\}_{k\ge 1}\) 为任意共终端严格递增子序列。
假设对每个 \(k\) 给出一对审计输入 \((\widehat{\mathbf{T}}_{N_k},\varepsilon_k)\)，并存在 \(\eta_k>0\) 使
\[
\lambda_{\min}(\widehat{\mathbf{T}}_{N_k})\ge \eta_k,
\qquad
\|\widehat{\mathbf{T}}_{N_k}-\mathbf{T}_{N_k}\|_{\mathrm{op}}\le \varepsilon_k\le \eta_k/2.
\]
则成立：
\begin{enumerate}
  \item 对所有 \(N\ge 0\) 均有 \(\mathbf{T}_N\succeq 0\)；
  \item 对任意 \(k\) 与任意 \(0\le N\le N_k\)，有谱裕度下界
  \[
  \boxed{\ \lambda_{\min}(\mathbf{T}_N)\ \ge\ \eta_k/2\ }.
  \]
  特别地，若 \(\inf_k\eta_k>0\)，则存在统一下界 \(\lambda_{\min}(\mathbf{T}_N)\ge \frac12\inf_k\eta_k\) 对一切 \(N\) 成立。
\end{enumerate}
\end{theorem}
\begin{proof}
由定理 \ref{thm:xi-spectral-margin-finite-psd-certificate} 的封口规则，对每个 \(k\) 有 \(\mathbf{T}_{N_k}\succeq 0\)。
再由定理 \ref{thm:xi-toeplitz-psd-cofinal-sparsification} 的共终端稀疏化等价性，推出对所有 \(N\ge 0\) 均有 \(\mathbf{T}_N\succeq 0\)。
\par
对谱裕度传播：由 Weyl 扰动界，
\[
\lambda_{\min}(\mathbf{T}_{N_k})
\ge
\lambda_{\min}(\widehat{\mathbf{T}}_{N_k})-\|\widehat{\mathbf{T}}_{N_k}-\mathbf{T}_{N_k}\|_{\mathrm{op}}
\ge
\eta_k-\varepsilon_k
\ge
\eta_k/2.
\]
而当 \(N\le N_k\) 时，\(\mathbf{T}_N\) 为 \(\mathbf{T}_{N_k}\) 的左上主子矩阵，故由 Cauchy 交错定理
\(
\lambda_{\min}(\mathbf{T}_N)\ge \lambda_{\min}(\mathbf{T}_{N_k})
\)。
合并即得所示下界。证毕。
\end{proof}

\begin{conclusion}[单原子负缺陷的可否证阈值是 \texorpdfstring{$1/(N+1)$}{1/(N+1)} 级且为锐界：视界阶数的最小信息预算被精确锁定]\label{con:xi-terminal-toeplitz-negative-atom-threshold-sharp}
在视界 Toeplitz 证书口径下，考虑背景正测度与单原子负缺陷 \(\varepsilon\,\delta_{\xi_0}\) 的对抗。
存在密度下界 \(m_0\) 与上界 \(M_0\) 使得：
\begin{enumerate}
  \item 若 \(\varepsilon\le m_0/(N+1)\)，则对一切 \(x\in\CC^{N+1}\) 有 \(x^\ast \mathbf{T}_N x\ge 0\)，从而 \(\mathbf{T}_N\succeq 0\)（负缺陷被背景完全吸收，视界不可否证）。
  \item 若 \(\varepsilon> M_0/(N+1)\)，则存在显式方向 \(u_k=\xi_0^{-k}\) 使 \(u^\ast\mathbf{T}_N u<0\)，从而 \(\mathbf{T}_N\not\succeq 0\)（负缺陷在第 \(N\) 级必被检测）。
\end{enumerate}
因此单原子缺陷的可否证性阈值在阶数上为严格的 \(1/(N+1)\) 量级并且为锐界：若要强制区分 \(\varepsilon\) 级缺陷，必需 \(N\gtrsim \varepsilon^{-1}\) 的视界阶数预算。
\end{conclusion}
\begin{proof}[证明]
见命题 \ref{prop:discussion-toeplitz-negative-atom-threshold}。证毕。
\end{proof}

% (User input window)

\paragraph{跨接口终端结论：Heine--Stieltjes 规约、短程势消尾、Schwarzian 刚性与锚定 Cauchy 插值的闭式结构}
本段汇总一组在锚定有理核、Stieltjes 型二阶方程与 Möbius 共变接口上可反复调用的结论。
其中 Schrödinger 规约与无穷远消尾条件给出“短程势”的严格代数刻画；
Schwarzian 刚性把二阶极点主部系数转译为分歧指数；
而 Cauchy--Lagrange 插值与 Cauchy 矩阵显式逆则为满秩锚定的线性解码提供完全闭式。

\begin{conclusion}[Heine--Stieltjes 方程的 Schr\"odinger 规约与局部指数刚性]\label{con:xi-terminal-heine-stieltjes-schrodinger-reduction-local-exponent-rigidity}
令 \(\kappa\ge 1\)，取互异点
\[
p_k=\gamma_k-\mathrm{i}a_k\in\CC_{-}\qquad(a_k>0,\ 1\le k\le \kappa),
\]
并设
\[
A(z):=\prod_{k=1}^{\kappa}(z-p_k)(z-\overline{p_k})\in\RR[z],
\qquad
\deg A=2\kappa.
\]
设 \(y\in\RR[z]\) 为次数 \(\kappa\) 的多项式，且存在 \(B\in\RR[z]\) 使
\[
A(z)y''(z)-A'(z)y'(z)=B(z)y(z).
\]
取任一局部解析分支 \(s(z):=\sqrt{A(z)}\)，并定义
\[
\phi(z):=\frac{y(z)}{s(z)}.
\]
则 \(\phi\) 满足 Schr\"odinger 型方程
\[
\phi''(z)+U(z)\phi(z)=0,
\qquad
U(z):=\frac{1}{2}\frac{A''(z)}{A(z)}-\frac{3}{4}\Bigl(\frac{A'(z)}{A(z)}\Bigr)^2-\frac{B(z)}{A(z)}.
\]
并且 \(U\) 为有理函数、\(\RR\)-对称（\(\overline{U(\overline z)}=U(z)\)），在每个 \(p_k,\overline{p_k}\) 处具有二阶极点且主部系数刚性为
\[
U(z)=-\frac{3}{4}\frac{1}{(z-p_k)^2}+O\Bigl(\frac{1}{z-p_k}\Bigr)\quad(z\to p_k),
\qquad
U(z)=-\frac{3}{4}\frac{1}{(z-\overline{p_k})^2}+O\Bigl(\frac{1}{z-\overline{p_k}}\Bigr)\quad(z\to\overline{p_k}).
\]
从而 \(\phi\) 在 \(p_k\) 处的 Frobenius 指数集合为 \(\{-\tfrac12,\tfrac32\}\)（同理对 \(\overline{p_k}\)），并因此 \(y=s\phi\) 在这些点处的局部指数为 \(\{0,2\}\)。
\end{conclusion}
\begin{proof}[证明]
由 \(y=s\phi\) 得
\[
y'=s'\phi+s\phi',\qquad
y''=s''\phi+2s'\phi'+s\phi''.
\]
代入 \(Ay''-A'y'\)，并用 \(s'/s=\tfrac12 A'/A\) 化简，可得
\[
Ay''-A'y'
=A^{3/2}\phi''+A^{1/2}\Bigl(\frac12A''-\frac34\frac{(A')^2}{A}\Bigr)\phi.
\]
将等式 \(Ay''-A'y'=By\) 先除以 \(A^{1/2}\)，再除以 \(A\)，即得 \(\phi''+U\phi=0\) 与所示 \(U\)。
\par
对局部主部：写 \(A(z)=(z-p_k)g_k(z)\) 且 \(g_k(p_k)\neq 0\)，则
\[
\frac{A'}{A}=\frac{1}{z-p_k}+\frac{g_k'}{g_k}.
\]
从而 \((A'/A)^2\) 在 \(p_k\) 的二阶主部为 \((z-p_k)^{-2}\)，而
\(
A''/A=(A'/A)'+(A'/A)^2
\)
的二阶项相消，故 \(A''/A\) 在 \(p_k\) 无二阶极点；又因 \(B/A\) 仅有一阶极点，遂 \(U\) 的二阶主部系数为 \(-\tfrac34\)。
由指示方程 \(r(r-1)-\tfrac34=0\) 得 \(r\in\{-\tfrac12,\tfrac32\}\)。
再乘以 \(s\sim (z-p_k)^{1/2}\) 可知 \(y=s\phi\) 的局部指数为 \(\{0,2\}\)。
对 \(\overline{p_k}\) 同理。证毕。
\end{proof}

\begin{conclusion}[无穷远“逆平方尾项”完全消失的充要条件：短程势的严格等价刻画]\label{con:xi-terminal-heine-stieltjes-short-range-tail-vanish}
在结论 \ref{con:xi-terminal-heine-stieltjes-schrodinger-reduction-local-exponent-rigidity} 的设定下，设
\[
B(z)=b_0 z^{2\kappa-2}+O(z^{2\kappa-3})\qquad(z\to\infty).
\]
则
\[
U(z)=\frac{-\kappa(\kappa+1)-b_0}{z^{2}}+O\Bigl(\frac{1}{z^{3}}\Bigr)\qquad(z\to\infty).
\]
因此下列两条件等价：
\begin{enumerate}
  \item \(U(z)=O(z^{-3})\)（短程势）；
  \item \(b_0=-\kappa(\kappa+1)\)。
\end{enumerate}
\end{conclusion}
\begin{proof}[证明]
写
\(
A(z)=z^{2\kappa}+A_1 z^{2\kappa-1}+O(z^{2\kappa-2})
\)，则
\[
\frac{A'}{A}=\frac{2\kappa}{z}+O\Bigl(\frac{1}{z^{2}}\Bigr),\qquad
\frac{A''}{A}=\frac{2\kappa(2\kappa-1)}{z^{2}}+O\Bigl(\frac{1}{z^{3}}\Bigr),\qquad
\frac{B}{A}=\frac{b_0}{z^{2}}+O\Bigl(\frac{1}{z^{3}}\Bigr).
\]
代入
\(
U=\tfrac12 A''/A-\tfrac34(A'/A)^2-B/A
\)
得
\[
\frac12\frac{A''}{A}-\frac34\Bigl(\frac{A'}{A}\Bigr)^2
=\frac{\kappa(2\kappa-1)}{z^{2}}-\frac{3\kappa^{2}}{z^{2}}+O(z^{-3})
=-\frac{\kappa(\kappa+1)}{z^{2}}+O(z^{-3}),
\]
再减去 \(B/A\) 即得所示展开。
短程势即首项系数为零，等价于 \(b_0=-\kappa(\kappa+1)\)。证毕。
\end{proof}

\begin{conclusion}[多项式解的零能散射刻画：次数条件等价于 Jost 归一化]\label{con:xi-terminal-heine-stieltjes-zero-energy-scattering-jost-normalization}
仍在结论 \ref{con:xi-terminal-heine-stieltjes-schrodinger-reduction-local-exponent-rigidity} 的设定下，并假设短程势条件 \(U(z)=O(z^{-3})\) 成立。
则在任一固定扇形 \(|\arg z|<\pi-\varepsilon\) 内，存在一组在无穷远唯一确定的基本解 \(\phi_{0},\phi_{1}\) 满足
\[
\phi_{0}(z)=1+O\Bigl(\frac{1}{z}\Bigr),
\qquad
\phi_{1}(z)=z+O(1)
\qquad(z\to\infty),
\]
且任一解均可唯一表示为 \(\phi=c_{0}\phi_{0}+c_{1}\phi_{1}\)。
在此归一化下，令无穷远分支 \(\sqrt A=z^{\kappa}(1+O(1/z))\)，则
\[
y(z)=\sqrt{A(z)}\,\phi(z)
\]
为次数恰为 \(\kappa\) 的多项式当且仅当 \(c_{1}=0\)；等价地，\(\phi\) 满足 Jost 边界条件 \(\phi(z)=1+O(z^{-1})\)。
\end{conclusion}
\begin{proof}[证明]
由 \(U(z)=O(z^{-3})\)，方程 \(\phi''+U\phi=0\) 在无穷远为可积扰动的自由方程 \(\phi''=0\)。
在扇形内可用 Volterra 型积分方程构造具有指定渐近的解并得唯一性：以 \(\phi_{0}\) 为例，可取
\[
\phi_{0}(z)=1+\int_{z}^{\infty}(t-z)U(t)\phi_{0}(t)\,\dd t,
\]
其中积分沿固定射线；\(U(t)=O(t^{-3})\) 保证收敛并给出压缩映射，故存在唯一解并满足 \(\phi_{0}(z)=1+O(1/z)\)。
同理得 \(\phi_{1}(z)=z+O(1)\)。
线性常微分方程的基本理论给出一般解的唯一线性分解。
\par
再由 \(\sqrt A=z^{\kappa}(1+O(1/z))\) 得
\[
\sqrt A\,\phi_{0}=z^{\kappa}\Bigl(1+O\Bigl(\frac{1}{z}\Bigr)\Bigr),
\qquad
\sqrt A\,\phi_{1}=z^{\kappa+1}\bigl(1+o(1)\bigr),
\]
故 \(\sqrt A(c_{0}\phi_{0}+c_{1}\phi_{1})\) 的最高次数为 \(\kappa\) 当且仅当 \(c_{1}=0\)。
而 \(c_{1}=0\) 与 \(\phi(z)=1+O(z^{-1})\) 等价。证毕。
\end{proof}

\begin{conclusion}[短程势的极点展开与“重整化和律”：由单极系数约束实现无穷远消尾]\label{con:xi-terminal-heine-stieltjes-pole-expansion-renormalization-sum-rules}
在结论 \ref{con:xi-terminal-heine-stieltjes-schrodinger-reduction-local-exponent-rigidity} 的设定下，令
\[
r_k:=\operatorname*{Res}_{z=p_k}\frac{B(z)}{A(z)}=\frac{B(p_k)}{A'(p_k)}\in\CC,
\]
以及
\[
s_k:=\left(\frac{A'}{A}-\frac{1}{z-p_k}\right)\Big|_{z=p_k}
=\frac{1}{p_k-\overline{p_k}}+\sum_{j\ne k}\left(\frac{1}{p_k-p_j}+\frac{1}{p_k-\overline{p_j}}\right).
\]
则 Schr\"odinger 势 \(U\) 具有完全部分分式展开
\[
U(z)=
-\frac34\sum_{k=1}^{\kappa}\left(\frac{1}{(z-p_k)^2}+\frac{1}{(z-\overline{p_k})^2}\right)
+\sum_{k=1}^{\kappa}\left(\frac{c_k}{z-p_k}+\frac{\overline{c_k}}{z-\overline{p_k}}\right),
\qquad
c_k:=-r_k-\frac12 s_k.
\]
并且短程势条件 \(U(z)=O(z^{-3})\) 等价于以下两条实线性约束（“重整化和律”）：
\[
\sum_{k=1}^{\kappa}\Re(c_k)=0,
\qquad
\sum_{k=1}^{\kappa}\Re(p_k c_k)=\frac{3\kappa}{4}.
\]
\end{conclusion}
\begin{proof}[证明]
由定义
\(
U=\frac12\frac{A''}{A}-\frac34(\frac{A'}{A})^2-\frac{B}{A}
\)。
在 \(p_k\) 处写
\[
\frac{A'}{A}=\frac{1}{z-p_k}+s_k+O(z-p_k),
\]
则
\[
\Bigl(\frac{A'}{A}\Bigr)^2=\frac{1}{(z-p_k)^2}+\frac{2s_k}{z-p_k}+O(1),
\qquad
\frac{A''}{A}=\Bigl(\frac{A'}{A}\Bigr)'+\Bigl(\frac{A'}{A}\Bigr)^2=\frac{2s_k}{z-p_k}+O(1),
\]
并且
\(
\frac{B}{A}=\frac{r_k}{z-p_k}+O(1)
\)。
因此 \(U\) 在 \(p_k\) 的二阶主部系数为 \(-\tfrac34\)，一阶留数为
\(
s_k-\tfrac32 s_k-r_k=-\tfrac12 s_k-r_k=c_k
\)。
对 \(\overline{p_k}\) 同理得 \(\overline{c_k}\)，汇总各点主部即得部分分式展开。
\par
将展开式在无穷远作 Laurent 展开：对 \(z\to\infty\),
\[
\frac{1}{z-p}=\frac{1}{z}+\frac{p}{z^{2}}+O(z^{-3}),
\qquad
\frac{1}{(z-p)^2}=\frac{1}{z^{2}}+\frac{2p}{z^{3}}+O(z^{-4}).
\]
于是
\[
U(z)=\frac{2\sum_k \Re(c_k)}{z}
+\frac{2\sum_k\Re(p_k c_k)-\frac34\cdot 2\kappa}{z^{2}}
+O(z^{-3}).
\]
故 \(U(z)=O(z^{-3})\) 当且仅当 \(z^{-1}\) 与 \(z^{-2}\) 系数均为零，即
\(
\sum\Re(c_k)=0
\)
与
\(
\sum\Re(p_k c_k)=\tfrac{3\kappa}{4}
\)。
证毕。
\end{proof}

\begin{conclusion}[\texorpdfstring{$\mathrm{PSL}(2,\RR)$}{PSL(2,R)} 共变性：势 \(U\) 为实二次微分]\label{con:xi-terminal-schrodinger-psl2r-quadratic-differential-covariance}
令 \(m(z)=\dfrac{az+b}{cz+d}\)（\(a,b,c,d\in\RR,\ ad-bc>0\)）为保向 Möbius 变换。
设 \(\phi\) 满足 \(\phi''(z)+U(z)\phi(z)=0\)。
令 \(\zeta=m(z)\)，并定义
\[
\widetilde\phi(\zeta):=\sqrt{\frac{\dd z}{\dd\zeta}}\;\phi(z),\qquad z=m^{-1}(\zeta).
\]
则 \(\widetilde\phi\) 满足同型方程
\[
\widetilde\phi''(\zeta)+\widetilde U(\zeta)\widetilde\phi(\zeta)=0,
\qquad
\widetilde U(\zeta)=\Bigl(\frac{\dd z}{\dd\zeta}\Bigr)^{2}U(z),
\]
即 \(U(z)\,\dd z^{2}\) 在 \(\mathrm{PSL}(2,\RR)\) 下按二次微分共变（Möbius 变换的 Schwarzian 为零，故无附加项）。
\end{conclusion}
\begin{proof}[证明]
由链式法则与 \(\widetilde\phi=\sqrt{z_\zeta}\phi\)（记 \(z_\zeta=\dd z/\dd\zeta\)）可得一般变换律
\[
\widetilde U(\zeta)=z_\zeta^{2}U(z)+\frac12\{z,\zeta\},
\]
其中 \(\{z,\zeta\}\) 为 Schwarzian 导数。
对 Möbius 变换 \(z=m^{-1}(\zeta)\) 有 \(\{z,\zeta\}=0\)，故 \(\widetilde U=z_\zeta^{2}U\)。证毕。
\end{proof}

\begin{conclusion}[满秩锚定下的 Cauchy--Lagrange 有理插值与 Cauchy 矩阵显式逆]\label{con:xi-terminal-cauchy-lagrange-rational-interpolation-cauchy-inverse}
取互异实点 \(a_1,\dots,a_\kappa\in\RR\)，并设
\[
\Pi(z):=\prod_{k=1}^{\kappa}(z-p_k),\qquad
y(z):=\prod_{j=1}^{\kappa}(z-a_j).
\]
考虑维数为 \(\kappa\) 的有理函数空间
\[
\mathcal R:=\left\{f(z)=\sum_{k=1}^{\kappa}\frac{c_k}{z-p_k}\;:\;c_k\in\CC\right\}.
\]
令 Cauchy 矩阵 \(C\in\CC^{\kappa\times\kappa}\) 为 \(C_{jk}=\frac{1}{a_j-p_k}\)。
则 \(C\) 可逆，且对任意数据 \(f_j\in\CC\) 存在唯一 \(f\in\mathcal R\) 满足 \(f(a_j)=f_j\)，并具有显式插值公式
\[
f(z)=\sum_{j=1}^{\kappa} f_j\,L_j(z),
\qquad
L_j(z):=\frac{y(z)}{(z-a_j)y'(a_j)}\cdot\frac{\Pi(a_j)}{\Pi(z)}.
\]
进一步，\(C^{-1}\) 具有完全闭式：对 \(1\le k,j\le\kappa\),
\[
(C^{-1})_{k j}
=\frac{y(p_k)\,\Pi(a_j)}{(p_k-a_j)\,\Pi'(p_k)\,y'(a_j)}.
\]
\end{conclusion}
\begin{proof}[证明]
先证插值基函数：对 \(i\ne j\)，因 \(y(a_i)=0\) 且 \(y(z)/(z-a_j)\) 在 \(z=a_i\) 处仍含因子 \((z-a_i)\)，故 \(L_j(a_i)=0\)；
而
\[
\lim_{z\to a_j}\frac{y(z)}{(z-a_j)y'(a_j)}=1,
\qquad
\frac{\Pi(a_j)}{\Pi(a_j)}=1,
\]
故 \(L_j(a_j)=1\)。
因此 \(f(z)=\sum f_j L_j(z)\) 满足 \(f(a_i)=f_i\)。
每个 \(L_j\) 为分母 \(\Pi(z)\) 的真分式，因而可作部分分式分解，属于 \(\mathcal R\)；唯一性由 \(\dim\mathcal R=\kappa\) 与评估映射的满秩性推出。
\par
再证 \(C^{-1}\) 闭式：写 \(L_j(z)=\sum_{k} \alpha_{kj}/(z-p_k)\)。
则对每个 \(i\),
\[
L_j(a_i)=\sum_k\frac{\alpha_{kj}}{a_i-p_k},
\]
矩阵记号即 \(C\alpha_{(:,j)}=e_j\)，故 \(\alpha=C^{-1}\)。
另一方面，由
\[
L_j(z)=\frac{\Pi(a_j)}{y'(a_j)}\cdot\frac{y(z)}{(z-a_j)\Pi(z)}
\]
在 \(z=p_k\) 处取留数得到
\[
\alpha_{k j}
=\operatorname*{Res}_{z=p_k}L_j(z)
=\frac{\Pi(a_j)}{y'(a_j)}\cdot\frac{y(p_k)}{(p_k-a_j)\Pi'(p_k)}.
\]
故 \((C^{-1})_{kj}=\alpha_{kj}\) 如所示。证毕。
\end{proof}

\begin{conclusion}[Schwarzian 刚性与二重分支：\texorpdfstring{$-\tfrac34(z-p)^{-2}$}{-(3/4)(z-p)^(-2)} 等价于分歧指数 \texorpdfstring{$2$}{2}]\label{con:xi-terminal-schwarzian-rigidity-branch-index-2}
设 \(U\) 为在 \(p\in\CC\) 的穿孔邻域内亚纯函数。
以下两条件等价：
\begin{enumerate}
  \item 存在局部亚纯函数 \(w\) 使其 Schwarzian 导数满足 \(\{w,z\}=2U(z)\)，且
  \[
  U(z)=-\frac34\frac{1}{(z-p)^2}+O\Bigl(\frac{1}{z-p}\Bigr)\qquad(z\to p).
  \]
  \item 存在局部坐标 \(\zeta=z-p\) 与可逆解析函数 \(h\) 使
  \[
  w(z)=h\bigl(\zeta^{2}\bigr),
  \]
  即 \(w\) 在 \(p\) 处的分歧指数（ramification index）恰为 \(2\)；等价地，其局部逆映射具有平方根型分歧。
\end{enumerate}
因此，对结论 \ref{con:xi-terminal-heine-stieltjes-schrodinger-reduction-local-exponent-rigidity} 的 Schr\"odinger 势 \(U\)，任取两线性无关解 \(\phi_1,\phi_2\)，其比值 \(w=\phi_1/\phi_2\) 满足 \(\{w,z\}=2U(z)\)，并在每个 \(p_k,\overline{p_k}\) 处具有分歧指数 \(2\)（零/极点的阶数为 \(2\)），且该指数由二阶主部系数 \(-\tfrac34\) 刚性锁定。
\end{conclusion}
\begin{proof}[证明]
经典事实表明：\(\phi''+U\phi=0\) 的两线性无关解之比 \(w=\phi_1/\phi_2\) 满足 \(\{w,z\}=2U(z)\)。
并且若 \(w\) 在 \(\zeta=0\) 的局部主项为 \(w\sim \zeta^{\alpha}\)（\(\alpha\neq 0\)），则 Schwarzian 的二阶主部为
\[
\{w,\zeta\}=\frac{1-\alpha^{2}}{2}\,\zeta^{-2}+O(\zeta^{-1}).
\]
\par
\((2)\Rightarrow(1)\)：取 \(w(\zeta)=\zeta^{2}\)，则 \(\{w,\zeta\}=-\tfrac{3}{2}\zeta^{-2}\)。
由链式法则与 \(\{w,z\}=(\dd\zeta/\dd z)^2\{w,\zeta\}+\{\zeta,z\}\) 可知二阶主部系数在解析坐标变换下保持不变；
此外，后复合可逆解析函数 \(h\) 仅在 \(\{h,w\}\) 处引入 \((w')^{2}\{h,w\}\) 型项，而当 \(w\) 在 \(\zeta=0\) 具有分歧指数 \(2\) 时 \(w'(\zeta)=O(\zeta)\)，故该项不影响 \(\zeta^{-2}\) 的主部系数。
因此可得 \(2U=\{w,z\}=-\tfrac{3}{2}(z-p)^{-2}+O((z-p)^{-1})\)，即 \(U\) 的二阶主部系数为 \(-\tfrac34\)。
\par
\((1)\Rightarrow(2)\)：若 \(\{w,z\}=-\tfrac{3}{2}(z-p)^{-2}+O((z-p)^{-1})\)，在局部坐标 \(\zeta=z-p\) 下取 Puiseux 主项 \(w(\zeta)\sim \zeta^{\alpha}\) 代入上式，比较 \(\zeta^{-2}\) 系数得 \((1-\alpha^{2})/2=-3/2\)，从而 \(\alpha^{2}=4\)，故 \(\alpha=\pm 2\)。
因此 \(w\) 在 \(p\) 处为二重分歧：若 \(\alpha=2\) 则 \(w=h(\zeta^{2})\)；若 \(\alpha=-2\) 则 \(1/w\) 满足前式并可化归为 \(\alpha=2\) 的情形（等价于 \(w\) 在 \(p\) 处为二阶极点，对应映到无穷远的分歧指数仍为 \(2\)）。
最后一段结论由结论 \ref{con:xi-terminal-heine-stieltjes-schrodinger-reduction-local-exponent-rigidity} 给出的二阶主部系数刚性直接推出。证毕。
\end{proof}


% (User input window)

\subsubsection{有界切片上的连通性、可逆化与残差几何}\label{subsubsec:xi-bounded-slice-connectivity-residual-geometry}
在定义 \ref{def:pom-bounded-slice} 的任意固定有界复杂度切片 \(\mathbf{PW}^{\ast}_{\le}\) 上，保守语义函子 \(\mathsf{Sem}\)（定义 \ref{def:pom-sem-functor}；定理 \ref{thm:pom-sem-conservative}）将实现差异收束为自然同构类；
同时，交换失败残差在可组合口径下被异常上同调类 \([\mathrm{Anom}]\) 承载（\S\ref{subsec:pom-anom-cohomology}）。
在该框架内，“连通性增长”可被严格拆解为：可逆化导致的同构类几何、残差承载的循环自由度、以及严格化所受的一致性障碍。

\begin{conclusion}[审计可逆的正交因子分解与连通性正规形的唯一性]\label{con:xi-connectivity-orthogonal-factorization-normal-form}
设 \(\rho:\mathrm{Ob}(\mathbf{PW}^{\ast}_{\le})\to\NN\) 为分辨率等级函数，并在 \(\mathbf{PW}^{\ast}_{\le}\) 上固定一套可审计的残差核算口径，使得：
\begin{enumerate}
  \item 对任意可复合 \(2\)-态射证书，其残差累积满足可复核的可加性（定义 \ref{def:pom-anom-accumulation}；命题 \ref{prop:pom-anom-additive-up-to-boundary}）；
  \item \(\mathsf{Sem}\) 在该切片上保守，从而反射自然同构（定理 \ref{thm:pom-sem-conservative}）。
\end{enumerate}
定义两类 \(1\)-态射：\(\mathcal L\) 为分辨率严格不降（\(\rho\) 单调不减）的提升类，\(\mathcal P\) 为分辨率严格不升（\(\rho\) 单调不增）的读出类。
则对任意 \(1\)-态射 \(w\) 存在分解
\[
w\simeq p\circ \ell,\qquad \ell\in\mathcal L,\ p\in\mathcal P,
\]
并且该分解在 \(2\)-同构意义下唯一。
更强地，\((\mathcal L,\mathcal P)\) 构成一个稳定正交分解系统：任意“读出--提升”交换方块存在填充且填充至唯一 \(2\)-同构。
因此，任意由生成门给出的概念连通路径在该切片内都存在唯一（至 \(2\)-同构）的规范形代表，并且在残差与复杂度意义下同时达到最小。
\end{conclusion}

\begin{conclusion}[语义等价类的图论刻画：强连通分量与可逆闭包]\label{con:xi-connectivity-scc-invertible-closure}
在固定有界切片 \(\mathbf{PW}^{\ast}_{\le}\) 上，将所有被 \(\mathsf{Sem}\) 识别为可逆的 \(1\)-态射形式地可逆化，得到局部化群胚 \(\mathcal{G}_{\le}\)。
取生成 \(1\)-态射得到有向多重图 \(G_{\le}\)（顶点为对象，边为生成态射）。
若在该切片上可审计重写足以判定 \(2\)-同构（从而“同一 \(\mathsf{Sem}\) 输出 \(\Rightarrow\) 自然同构”可被有限证书核对；定理 \ref{thm:pom-sem-conservative}），则 \(G_{\le}\) 的强连通分量分解与 \(\mathcal{G}_{\le}\) 的同构类分解一致。
等价地，两对象在 \(G_{\le}\) 中互达，当且仅当在 \(\mathbf{PW}^{\ast}_{\le}\) 中存在一对互为逆的可审计投影词（至 \(2\)-同构）。
因此，在不改变语义切分的前提下，加边只能改变同一语义类内的导航直径与分量间凝聚 DAG 的可达代价，而不能产生新的语义同构类。
\end{conclusion}

\begin{conclusion}[分辨率过滤下的持久连通谱：稳定概念的条形码刻画]\label{con:xi-connectivity-persistent-barcode}
令预算/分辨率上界 \(B\mapsto \mathbf{PW}^{\ast}_{\le B}\) 给出一族单调嵌入的有界切片，并令 \(\mathcal{G}_{\le B}\) 为其在 \(\mathsf{Sem}\)-可逆态射上的局部化群胚。
记 \(\pi_0(\mathcal{G}_{\le B})\) 为同构类集合，则包含函子诱导出直系系统
\(
\pi_0(\mathcal{G}_{\le B})\to \pi_0(\mathcal{G}_{\le B'})
\)（\(B\le B'\)）。
若对每个 \(B\) 该集合有限，且类的合并只能由可审计证书触发（不会发生不可追踪坍塌），则该过滤唯一确定一组区间多重集 \(\mathcal I\)（持久连通谱/条形码），满足：
\begin{enumerate}
  \item 无限区间与稳定语义类一一对应：存在 \(B_0\) 使得对所有 \(B\ge B_0\) 该类在 \(\pi_0(\mathcal{G}_{\le B})\) 中保持不变；\label{item:xi-persist-infinite}
  \item 有限区间对应瞬时同构或瞬时概念，其出现与消失可由残差与证书定位；\label{item:xi-persist-finite}
  \item 稳定概念可在逆极限意义下刻画为 \(\varprojlim_B \mathcal{G}_{\le B}\) 中的非 \(\mathsf{NULL}\) 可寻址对象族（参见定理 \ref{thm:xi-time-consistent-inverse-limit-readout} 的同型逆系统论证）。\label{item:xi-persist-invlim}
\end{enumerate}
于是，连通性随预算增长的可观察变化可被重述为：在保持 \ref{item:xi-persist-infinite} 所述无限条不变的前提下，允许有限条形码发生可审计迁移，从而把体系增长约束为对持久谱的受控扰动。
\end{conclusion}

\begin{conclusion}[连通性增长的不可逃代价：循环秩由异常维数下界控制]\label{con:xi-connectivity-betti-lower-bound-by-anom-rank}
令 \(G_{\le B}\) 为预算 \(B\) 下的生成图（可取为结论 \ref{con:xi-connectivity-scc-invertible-closure} 的 \(G_{\le}\) 在阈值 \(B\) 下的子图），并令 \(\beta_1(G_{\le B})\) 表示其无向化后的第一 Betti 数。
取阿贝尔群 \(\mathrm{Anom}_{\le B}\) 表示在切片 \(\mathbf{PW}^{\ast}_{\le B}\) 上可核对的异常自由度（例如由定义 \ref{def:pom-anom-accumulation} 的残差坐标生成），并设存在闭路残差映射
\(
\alpha_B:Z_1(G_{\le B})\to \mathrm{Anom}_{\le B}
\)
将每个闭路送到其规范化审计残差。
若零残差闭路必可由 \(2\)-同构填充消去（即残差外置最小性），则 \(\alpha_B\) 在同调意义上诱导单射
\[
H_1(G_{\le B};\ZZ)\hookrightarrow \mathrm{Anom}_{\le B},
\]
从而得到结构性下界
\[
\beta_1(G_{\le B})\le \mathrm{rank}_{\ZZ}(\mathrm{Anom}_{\le B}).
\]
因此，若试图通过加边提高连通性密度并同时保持语义不坍塌，则异常自由度必须同比例增长以承载新增的独立循环；否则新增循环将被迫塌缩为可逆冗余，不产生新的可用结构。
\end{conclusion}

\begin{conclusion}[审计保持严格化的三阶障碍：连通性闭合的高阶一致性判据]\label{con:xi-connectivity-third-order-strictification-obstruction}
在 \(\mathbf{PW}^{\ast}_{\le B}\) 的双范畴结构中，连通性不仅依赖 \(1\)-态射可达性，还依赖 \(2\)-态射的相容性（不同括号化/不同重写路径的审计一致）。
由三重复合的括号化差可定义一个取值于 \(\mathrm{Anom}_{\le B}\) 的三阶一致性障碍类
\[
\omega_B\in H^3(G_{\le B};\mathrm{Anom}_{\le B}),
\]
可视为审计保持的结合子缺陷的同调类。
则：存在一个在 \(2\)-范畴意义下与 \(\mathbf{PW}^{\ast}_{\le B}\) 等价、且保持 \(\mathsf{Sem}\) 与残差核算的严格 \(2\)-范畴模型，当且仅当 \(\omega_B=0\)。
因此，将任意多段连通路径规约到唯一规范形且审计一致的“可机器导航闭合”，受制于一个可计算的三阶一致性障碍；当该障碍为零时，连通性可在不增添额外寄存器自由度的前提下严格闭合。
\end{conclusion}

\begin{conclusion}[最小知识图谱骨架的存在与计数：连通可用性的线性规模实现]\label{con:xi-connectivity-minimal-skeleton-linear-size}
设预算 \(B\) 下对象数为 \(n_B:=|\mathrm{Ob}(\mathbf{PW}^{\ast}_{\le B})|\)，并假设 \(\mathrm{Anom}_{\le B}\) 为自由阿贝尔群，且闭路残差映射 \(\alpha_B\) 在同调上满射（即异常维度在切片上可由闭路证书完全激发）。
则存在一个子图 \(H_B\subseteq G_{\le B}\)（审计等价骨架）满足：
\begin{enumerate}
  \item \(H_B\) 与 \(G_{\le B}\) 诱导相同的局部化群胚（同构类与可达关系一致）；\label{item:xi-skel-groupoid}
  \item \(\alpha_B\) 在 \(H_B\) 上仍保持同调满射（异常自由度不丢失）；\label{item:xi-skel-anom}
  \item 边数达到不可改进下界
  \[
  |E(H_B)| = n_B - c_B + \mathrm{rank}_{\ZZ}(\mathrm{Anom}_{\le B}),
  \]
  其中 \(c_B\) 为凝聚图的弱连通分量数。\label{item:xi-skel-size}
\end{enumerate}
因此，在保守语义与异常审计充分激发的前提下，任何可用的概念连通性都可压缩到线性规模骨架，且骨架大小由“对象数 \(+\) 必要异常维度”精确决定；若异常维度不随连通性增长而增长，则新增边最终仅形成可删去的冗余而不提升可用结构覆盖率。
\end{conclusion}


% (User input window)

\paragraph{预算过滤的完备化、谱表示与实现自由度}
在有界切片的保守语义口径下（定义 \ref{def:pom-bounded-slice}；定理 \ref{thm:pom-sem-conservative}），预算增长诱导的过滤不仅给出同构类的持久谱（结论 \ref{con:xi-connectivity-persistent-barcode}），还自然导出一类与完备化、逻辑谱与严格化障碍相互等价的几何/代数表示。
下列条目在不改变 \(\mathbf{PW}^{\ast}_{\le}\) 的既有记号与审计约定的前提下，将这些表示组织为可独立引用的结论链。

\begin{conclusion}[预算局部化诱导的非阿基米德概念几何与完备化同一性]\label{con:xi-budget-nonarchimedean-geometry-completion}
设 \(\mathbb{B}\) 为带最小元的全序集合，并以幂等并半群运算
\(
B\vee B':=\max\{B,B'\}
\)
表示预算门槛的合成律。
设 \(B\mapsto\mathcal{G}_{\le B}\) 为一族随预算单调嵌入的局部化群胚：若 \(B\le B'\) 则存在忠实函子 \(\iota_{B}^{B'}:\mathcal{G}_{\le B}\hookrightarrow\mathcal{G}_{\le B'}\)，且任意可逆证书的预算代价满足最大律（复合代价等于各步代价之 \(\vee\)）。
对对象 \(x,y\) 定义
\[
d(x,y):=\min\{\,B\in\mathbb{B}:\ x\simeq y\ \text{于}\ \mathcal{G}_{\le B}\,\},
\qquad d(x,x):=\min\mathbb{B}.
\]
并约定当集合为空时 \(d(x,y)=+\infty\)。
则 \(d\) 为一非阿基米德拟度量：对任意 \(x,y,z\) 有
\[
d(x,z)\le d(x,y)\vee d(y,z).
\]
并且，以 \(d\) 的球 \(B(x;B):=\{y:\ d(x,y)\le B\}\) 生成的拓扑之完备化，与同构类逆极限空间同胚：
\[
\widehat{\mathrm{Ob}}(\mathcal{G})
:=\varprojlim_{B\in\mathbb{B}}\ \pi_0(\mathcal{G}_{\le B})
\ \cong\ 
\overline{\mathrm{Ob}}(\mathcal{G})\qquad(\text{\(d\)-完备化}).
\]
因此，在该假设下，连通性增长等价于对一非阿基米德概念空间的分辨率细化；其极限对象为严格意义上的“全预算可寻址概念族”（与结论 \ref{con:xi-connectivity-persistent-barcode} 的逆极限口径一致）。
\end{conclusion}

\begin{conclusion}[审计命题的布尔谱与可寻址概念的谱表示]\label{con:xi-audit-boolean-stone-spectrum}
对每个预算 \(B\)，令 \(\mathfrak{A}_B\) 表示由可核对的语义指纹命题生成的审计命题布尔代数：
其生成元可取为对 \(\mathsf{Sem}\) 输出的有限判定式（例如有限签名族 \(\Sigma_m(w)\) 的坐标相等/不等、以及谱指纹在网格点上的离散判定；见定义 \ref{def:pom-sem-functor} 与注 \ref{rem:pom-finite-signature-sigmam}）。
令 \(\mathfrak{A}:=\varinjlim_B \mathfrak{A}_B\) 为随预算放宽诱导的直系极限（保守细化保证嵌入性）。
记 \(\mathrm{Stone}(\mathfrak{A})\) 为 \(\mathfrak{A}\) 的超滤子空间（Stone 谱），赋以自然紧零维拓扑。

若审计完备性成立：任意两不同的逆极限对象族 \((x_B)\neq (y_B)\) 都可被某个有限预算层的某个命题分离（即存在 \(B_0\) 与 \(P\in\mathfrak{A}_{B_0}\) 使 \(P(x_{B_0})\neq P(y_{B_0})\)），
则存在自然同胚
\[
\mathrm{Stone}(\mathfrak{A})\ \cong\ \varprojlim_{B}\ \pi_0(\mathcal{G}_{\le B}).
\]
并且，每个稳定概念对应唯一极大一致的审计命题完备集（超滤子）。
该同胚把“可发表的全局断言”识别为 Stone 谱上的闭开集，从而将概念连通性转化为一紧零维谱空间上的可分离逻辑结构。
\end{conclusion}

\begin{conclusion}[语义保持的重写实现模空间由二阶异常上同调刻画]\label{con:xi-implementation-moduli-h2-anom}
在固定生成 \(1\)-态射与保守语义函子 \(\mathsf{Sem}\) 的条件下，称一个\emph{实现}为：对每对可复合 \(1\)-态射 \((u,v)\) 选取一条可审计 \(2\)-同构
\(
\phi_{u,v}:u\circ v\Rightarrow uv
\)，
并满足单位约束；以异常群 \(\mathrm{Anom}_{\le B}\) 记录偏差，使不同括号化/不同重写路径的差异落入 \(\mathrm{Anom}_{\le B}\) 的中心坐标。
称两实现\emph{规范等价}，若存在一族 \(2\)-同构作规范变换将其相互扭转。

若三阶结合障碍在预算 \(B\) 上消失（即存在实现；结论 \ref{con:xi-connectivity-third-order-strictification-obstruction} 的 \(\omega_B=0\) 情形），则所有实现的规范等价类与二阶异常上同调群存在自然双射：
\[
{\textup{实现}}/{\textup{规范等价}}\ \cong\ H^2(\mathcal{G}_{\le B};\mathrm{Anom}_{\le B}).
\]
因此，在障碍为零时，“在保持语义的前提下增加连通性”的自由度被一个可计算的 \(H^2\) 参数空间完全控制；其非平凡性给出同一语义接口下存在多重可兼容实现的严格判据。
\end{conclusion}

\begin{conclusion}[最小中心扩张消除三阶结合缺陷的普适性]\label{con:xi-minimal-central-extension-kill-omega}
记结论 \ref{con:xi-connectivity-third-order-strictification-obstruction} 的三阶结合缺陷类为
\(
\omega_B\in H^3(\mathcal{G}_{\le B};\mathrm{Anom}_{\le B})
\)。
不假设 \(\omega_B=0\)。
则存在一最小秩中心扩张
\[
\mathrm{Anom}_{\le B}\hookrightarrow \widehat{\mathrm{Anom}}_{\le B},
\]
其中 \(\widehat{\mathrm{Anom}}_{\le B}\) 由在 \(\mathrm{Anom}_{\le B}\) 之上添加有限维自由中心因子得到，使得 \(\omega_B\) 在系数提升下被消去：
\[
\mathrm{infl}(\omega_B)=0\qquad\text{于}\qquad H^3(\mathcal{G}_{\le B};\widehat{\mathrm{Anom}}_{\le B}).
\]
并且该扩张在如下意义下普适：任意能消去 \(\omega_B\) 的中心扩张都唯一地因子化经过 \(\widehat{\mathrm{Anom}}_{\le B}\)。
此外，所需新增中心维数等于由 \(\omega_B\) 在 \(3\)-链上的取值张成的最小自由子模的秩，从而给出“连通性严格闭合”的最低异常代价。
\end{conclusion}

\begin{conclusion}[严格凸审计能量下的规范词唯一性与测地代价]\label{con:xi-strict-convex-energy-canonical-word-geodesic}
设异常群 \(\mathrm{Anom}_{\le B}\) 含一实向量空间分量 \(V_B\)，并给定严格凸范数 \(\|\cdot\|_B\)。
对任一固定语义态射 \(f\)（于 \(\mathcal{G}_{\le B}\) 中），令 \(\mathcal{R}_B(f)\subset V_B\) 为所有实现 \(f\) 的可达残差集合，并假设 \(\mathcal{R}_B(f)\) 为非空闭仿射子空间（由残差外置与可拼接完备性保证）。
则存在唯一的能量最小残差
\[
r^\star(f)=\arg\min_{r\in\mathcal{R}_B(f)}\ \|r\|_B,
\]
并存在（至 \(2\)-同构唯一的）规范实现 \(w_f^\star\) 使其残差为 \(r^\star(f)\)。
若进一步对复合满足 Minkowski 可加性
\(
\mathcal{R}_B(g\circ f)=\mathcal{R}_B(g)+\mathcal{R}_B(f)
\)，
则规范残差满足
\[
\|r^\star(g\circ f)\|_B\le \|r^\star(g)\|_B+\|r^\star(f)\|_B,
\]
从而规范实现诱导一条可加的测地代价，使概念连通路径在审计意义下获得唯一规范代表。
\end{conclusion}

\begin{conclusion}[可寻址性预层的胶合判据：全局概念的存在由一阶上同调障碍决定]\label{con:xi-addressability-presheaf-gluing-h1-obstruction}
将预算偏序 \(\mathbb{B}\) 视为站点。
在全序情形下，交叠预算满足 \(B_i\wedge B_j:=\min\{B_i,B_j\}\)。
定义可寻址性预层 \(\mathcal{A}\)：对每个预算 \(B\)，令 \(\mathcal{A}(B)\) 为在预算 \(B\) 下非 \(\mathsf{NULL}\) 的地址/读出前缀集合，限制映射由预算放宽诱导。
若审计分离性成立（不同全局概念在某层必可由指纹命题区分；亦即结论 \ref{con:xi-audit-boolean-stone-spectrum} 的分离假设），则 \(\mathcal{A}\) 为分离预层。

对任意覆盖 \(\{B_i\}\) 上的局部一致族 \(a_i\in\mathcal{A}(B_i)\)，若其在交叠 \(B_i\wedge B_j\) 上相容，则存在一规范的 \v{C}ech \(1\)-余循环 \(\eta\)（取值于由零语义回路诱导的核预层），使得：该局部一致族可胶合为全局可寻址概念，当且仅当 \([\eta]=0\)。
因此，“连通性不足导致的不可全局化”可被严格化为可计算的一阶上同调障碍。
\end{conclusion}

\begin{conclusion}[三投影联合保守下的二拉回分解：连通性与异常的寄存器可分离]\label{con:xi-three-projection-2pullback-register-separability}
设存在三投影函子
\(
P_{\mathrm{prim}},P_{\mathrm{op}},P_{\mathrm{prime}}
\)
至三类寄存器语义群胚
\(
\mathcal{G}_{\le B}^{\mathrm{prim}},\mathcal{G}_{\le B}^{\mathrm{op}},\mathcal{G}_{\le B}^{\mathrm{prime}}
\)，
并满足：
\begin{enumerate}
  \item 联合保守：若三投影皆为同构则原对象同构；
  \item 纤维完备：任意三投影上的相容数据可提升为原群胚中的数据，且提升至 \(2\)-同构唯一；
  \item 异常可分裂：\(\mathrm{Anom}_{\le B}\cong \mathrm{Anom}_{\le B}^{\mathrm{prim}}\oplus \mathrm{Anom}_{\le B}^{\mathrm{op}}\oplus \mathrm{Anom}_{\le B}^{\mathrm{prime}}\)。
\end{enumerate}
则存在群胚等价
\[
\mathcal{G}_{\le B}\ \simeq\ 
\mathcal{G}_{\le B}^{\mathrm{prim}}\times^{(2)}\mathcal{G}_{\le B}^{\mathrm{op}}\times^{(2)}\mathcal{G}_{\le B}^{\mathrm{prime}},
\]
其中 \(\times^{(2)}\) 表示双范畴意义下的弱 \(2\)-拉回。
并且所有连通性不变量均按“最大/直和法则”分解：结论 \ref{con:xi-budget-nonarchimedean-geometry-completion} 的非阿基米德距离为三分量距离之 \(\vee\)，障碍类与异常类为三分量障碍之直和像。
该分解提供一条寄存器可分离原则：体系连通性的真实新增必须在至少一个寄存器分量上产生不可消去的结构变化。
\end{conclusion}

\begin{conclusion}[预算过滤的异常谱序列：稳定概念与障碍跃迁的统一计算框架]\label{con:xi-budget-anom-spectral-sequence}
以预算过滤构造审计链复形 \(C_\bullet(B)\)：取 \(\mathcal{G}_{\le B}\) 的 nerve，并以异常系数记录 \(2\)-胞与 \(3\)-胞的审计残差。
取一条离散化的预算序列 \(B_0<B_1<\cdots\) 表示过滤台阶，并令 \(F^{p}C_\bullet:=C_\bullet(B_p)\)。
则存在自然谱序列 \((E_r^{p,q},d_r)\) 收敛到过滤极限的审计同调：
\[
E_1^{p,q}\cong H_{p+q}\bigl(C_\bullet(B_{p})/C_\bullet(B_{p-1})\bigr)
\ \Rightarrow\ 
H_{p+q}\bigl(\varinjlim_B C_\bullet(B)\bigr),
\]
其中约定 \(C_\bullet(B_{-1})=0\)。
其中各级差分 \(d_r\) 精确编码“在预算提升 \(r\) 个等级时出现的最小不可胶合残差”。
稳定概念对应 \(E_\infty^{0,0}\) 的永久循环；
三阶结合障碍与实现模空间的跃迁对应高页差分的非零性（分别与结论 \ref{con:xi-connectivity-third-order-strictification-obstruction}、结论 \ref{con:xi-implementation-moduli-h2-anom} 对齐）。
因此，连通性随预算增长的相变可被严格化为谱序列差分的消长，从而提供一条可计算、可验证、可累积的全局治理主线。
\end{conclusion}


% (User input window)

\paragraph{离散残差模型下的连通性维数、可核对稀疏化与指纹复杂度}
在固定预算层上，将可审计残差抽象为有限维整格标签，可把“连通性增长”与“审计自由度增长”之间的硬约束完全降维为有限图上的计数与同调问题。
下列结论给出三类可复核后果：异常维数的不可压缩下界、全局一致性的有限检验、以及跨预算层同时保持可达与残差核算的线性规模骨架。

\begin{conclusion}[稳定概念熵与异常维数的不可压缩下界]\label{con:xi-stable-concept-entropy-anom-dimension-lower-bound}
设有一有限有向图 \(G=(V,E)\)。
每条边 \(e\in E\) 赋予两类可审计数据：预算阈值 \(c(e)\in\NN\) 与残差 \(r(e)\in\ZZ^d\)；
并约定反向边 \(e^{-1}\) 的残差为 \(-r(e)\)。
对任意路径词 \(w=e_1\cdots e_\ell\) 定义
\[
c(w):=\max_{1\le i\le \ell} c(e_i),
\qquad
r(w):=\sum_{i=1}^{\ell} r(e_i).
\]
给定预算 \(B\)，令 \(\mathcal{W}_B:=\{w:\ c(w)\le B\}\)。
对对象对 \((x,y)\) 引入零残差可达关系：记 \(x\to_B y\) 当且仅当存在 \(w\in\mathcal{W}_B\) 使 \(w:x\to y\) 且 \(r(w)=0\)。
令 \(N(B)\) 为由 \(\to_B\) 诱导的可达类数（等价地，可取零残差强连通分量的个数）。

若存在常数 \(R\) 使对所有边有 \(\|r(e)\|_\infty\le R\)，并且在预算 \(B\) 下任意零残差可达的最短证书长度不超过 \(L(B)\)，
则有计数上界
\[
N(B)\ \le\ \bigl(2R\,L(B)+1\bigr)^d,
\]
从而得到维数下界
\[
d\ \ge\ \frac{\log N(B)}{\log\bigl(2R\,L(B)+1\bigr)}.
\]
特别地，若 \(L(B)\le \lambda B\) 且 \(N(B)\) 随 \(B\) 指数增长，则异常维数 \(d\) 必随 \(B\) 至少线性增长。
因此，在固定维数的异常层之下，连通性密度的持续增加若不伴随可审计自由度的增长，最终只能通过合并稳定概念类实现，而不能维持严格的语义分裂。
\end{conclusion}

\begin{conclusion}[树势函数--回路基的残差分解与端点可压缩性]\label{con:xi-tree-potential-cycle-basis-residual-decomposition}
固定预算 \(B\)，令 \(G_{\le B}\) 为由满足 \(c(e)\le B\) 的边生成的无向化底图，并取其一生成森林 \(T\)。
在每个连通分量选根 \(o\)，对任意顶点 \(v\) 定义树势函数
\[
\Phi_B(v):=\sum_{e\in P_T(o,v)} \epsilon_e\, r(e),
\]
其中 \(P_T(o,v)\) 为树上唯一路径，\(\epsilon_e\in\{\pm1\}\) 由遍历方向决定。
记所有非树边集合为 \(E^\ast\)；对每条 \(f\in E^\ast\)，由 \(T\cup\{f\}\) 生成的基本回路定义其回路残差 \(h_B(f)\in\ZZ^d\)。
则对任意同分量内从 \(x\) 至 \(y\) 的任一路径词 \(w\)（允许重复边），存在整数系数 \(n_f(w)\in\ZZ\) 使得
\[
r(w)=\Phi_B(y)-\Phi_B(x)+\sum_{f\in E^\ast} n_f(w)\,h_B(f).
\]
因此，残差在商群
\(
\ZZ^d/\langle h_B(f):f\in E^\ast\rangle
\)
中仅由端点决定；换言之，在由基本回路残差张成的有限生成子模之外，残差部分可被端点势函数完全压缩。
\end{conclusion}

\begin{conclusion}[全局审计一致性的有限检验原理]\label{con:xi-global-audit-consistency-finite-test}
沿用结论 \ref{con:xi-tree-potential-cycle-basis-residual-decomposition} 的记号。
称预算 \(B\) 下审计为路径独立，若对同一对端点 \((x,y)\) 任意两条路径词 \(w,w'\) 都满足 \(r(w)=r(w')\)。
则路径独立当且仅当所有基本回路残差同时为零：对一切 \(f\in E^\ast\) 有 \(h_B(f)=0\)。
在该情形下存在势函数 \(\Psi_B:V\to\ZZ^d\) 使每条边 \(e:u\to v\) 满足
\[
r(e)=\Psi_B(v)-\Psi_B(u),
\]
从而任意路径残差恒等于端点差 \(r(w)=\Psi_B(y)-\Psi_B(x)\)。
该判据给出一种严格的局部证书完备性：审计是否全局一致无需枚举所有回路或所有重写，只需检验不超过
\(
|E_{\le B}|-|V|+\#\mathrm{comp}(G_{\le B})
\)
条基本回路的残差即可决定。
\end{conclusion}

\begin{conclusion}[自证稀疏骨架：同时保持所有预算层的可达与残差核算]\label{con:xi-self-certifying-skeleton-all-budgets}
设预算取值集合为边阈值的有限集合
\(
\mathcal{B}:=\{c(e):e\in E\}
\)。
在每个 \(B\in\mathcal{B}\) 上考虑图 \(G_{\le B}\) 及其回路残差子模
\(
H_B:=\langle h_B(f):f\in E^\ast\rangle\subseteq\ZZ^d
\)。
则存在一子图 \(S\subseteq G\)（自证骨架）满足：
\begin{enumerate}
  \item 对每个 \(B\in\mathcal{B}\)，\(S_{\le B}\) 与 \(G_{\le B}\) 具有相同的连通分量划分；
  \item 对每个 \(B\in\mathcal{B}\)，由 \(S_{\le B}\) 诱导的回路残差子模等于 \(H_B\)；
  \item 骨架具自证性：\(S\) 中每条边 \(e\) 都存在一条仅用 \(S\) 中边构成的审计证书，证明其残差标签与“势函数--回路基”分解相容，从而任一预算层上的残差核算均可在 \(S\) 内闭合完成；
  \item 边数满足上界
  \[
  |E(S)|\le |V|-\#\mathrm{comp}(G_{\le \max\mathcal{B}})+d\cdot|\mathcal{B}|.
  \]
\end{enumerate}
因此，只要异常维数 \(d\) 固定，则图谱在全部预算层上的“连通性 + 审计残差结构”可被压缩到线性规模的可核对骨架，而不牺牲任何预算层面的可达与残差可计算性。
\end{conclusion}

\begin{conclusion}[多尺度指纹的最优位长定律：稳定类数与误判率决定复杂度]\label{con:xi-multiscale-fingerprint-optimal-bitlength}
对每个预算 \(B\)，令 \(\mathcal{C}_B\) 为预算 \(B\) 下的稳定概念类集合（例如零残差强连通分量的集合）。
设需构造指纹映射 \(F_B:\mathcal{C}_B\to\{0,1\}^{k_B}\)，使任意两不同类发生碰撞的概率不超过 \(\delta_B\)（在随机指纹族意义下，或在对抗性模型中对所有对的最坏界）。
则必有信息论下界
\[
k_B\ \ge\ \left\lceil \log_2 |\mathcal{C}_B|+\log_2\frac1{\delta_B}\right\rceil.
\]
更强地，存在一前缀一致的多尺度指纹
\(
F:\bigcup_B \mathcal{C}_B\to\{0,1\}^\infty
\)，
使得对每个 \(B\) 取其前 \(k_B\) 位即得到满足上述碰撞约束的 \(F_B\)。
因此，在预算细化过程中，指纹位长的不可压缩增长率完全由 \(|\mathcal{C}_B|\) 的增长与可容许误判率 \(\delta_B\) 所决定，而与图中冗余连边与路径数量无关。
\end{conclusion}

\begin{conclusion}[连通性增长的层级律：稳定类合并仅发生在有限临界预算上]\label{con:xi-connectivity-hierarchy-finite-critical-budgets}
定义 \(x\sim_B y\) 当且仅当在预算 \(B\) 下存在零残差证书实现 \(x\to_B y\) 且 \(y\to_B x\)。
则 \(\sim_B\) 为等价关系，并且随 \(B\) 单调粗化：若 \(B\le B'\) 则 \(\sim_B\) 细于 \(\sim_{B'}\)。
设 \(|V|=n\)，则存在至多 \(n-1\) 个临界预算
\(
B_1<\cdots<B_m
\)
（每个 \(B_i\in\mathcal{B}\)），使得等价关系仅在这些 \(B_i\) 处发生严格合并；在所有非临界 \(B\) 上，\(\sim_B\) 保持不变。
并且每一次合并都可由一个原子证书见证：存在一对等价类 \(C\neq C'\) 与一条预算恰为 \(B_i\) 的边 \(e\)，使得该合并由使用 \(e\) 的零残差闭合所触发。
因此，连通性随预算增长的可观察结构变化是一条有限步的合并链；若新增边不引入新的临界预算或新的闭路残差生成元，则在稳定类层面不可见。
\end{conclusion}


% (User input window)

\paragraph{算术地址、Smith 不变量与连通性层级的可计算治理}
在结论 \ref{con:xi-tree-potential-cycle-basis-residual-decomposition}--\ref{con:xi-connectivity-hierarchy-finite-critical-budgets} 的离散残差口径下，零残差等价类与其随预算的合并链不仅可由图论描述，还可被进一步压缩为有限阿贝尔群的商结构与 Smith 不变量的冻结判定。
下列结论把“连通性增长”完全归约为一条可复核的算术商链，并给出一遍扫描即可输出的治理算法框架。

\begin{conclusion}[零残差可达关系的整格余类刻画与算术地址的存在唯一性]\label{con:xi-arithmetic-address-existence-uniqueness}
设 \(G=(V,E)\) 为有限图。
对每条有向边 \(e:u\to v\) 赋予残差标记 \(r(e)\in\ZZ^d\)，并规定反向边满足 \(r(e^{-1})=-r(e)\)；
同时赋予预算阈值 \(c(e)\in\NN\)。
对预算 \(B\) 定义子图 \(G_{\le B}\)（取满足 \(c(e)\le B\) 的边）。
若 \(u,v\) 位于 \(G_{\le B}\) 的同一连通分量内，令
\[
R_B(u,v)\subseteq\ZZ^d
\]
表示所有从 \(u\) 至 \(v\) 的路径残差和集合。
则存在一秩不超过 \(d\) 的子格 \(H_B\subseteq\ZZ^d\)（由该分量内所有闭路残差生成，且与生成树选择无关），使得对该分量内任意 \(u,v\) 都有
\[
R_B(u,v)=\alpha_B(v)-\alpha_B(u)+H_B,
\]
其中 \(\alpha_B:V\to\ZZ^d\) 可取为任一生成树势函数（结论 \ref{con:xi-tree-potential-cycle-basis-residual-decomposition} 的 \(\Phi_B\)）。
由此定义算术地址
\[
a_B:V\to Q_B:=\ZZ^d/H_B,
\qquad
a_B(v):=\alpha_B(v)\bmod H_B,
\]
则该 \(a_B\) 在加常数平移意义下唯一，并且
\[
0\in R_B(u,v)\quad\Longleftrightarrow\quad a_B(u)=a_B(v).
\]
因此，预算 \(B\) 下的零残差等价类严格等同于算术地址的纤维分解。
\end{conclusion}

\begin{conclusion}[预算过滤下算术地址的函子性：连通性层级为有限交换商链]\label{con:xi-arithmetic-address-functorial-quotient-chain}
在结论 \ref{con:xi-arithmetic-address-existence-uniqueness} 的记号下，随预算增长有子格单调性 \(H_B\subseteq H_{B'}\)（\(B\le B'\)），从而诱导满同态
\[
\pi_{B}^{B'}:Q_B\to Q_{B'}
\]
满足 \(a_{B'}=\pi_{B}^{B'}\circ a_B\)。
特别地，零残差等价关系族 \((\sim_B)\) 在每个连通分量上由一条有限阿贝尔群商链完全控制：任意 \(B\le B'\) 下的类合并，恰对应于将 \(Q_B\) 进一步商去由新增闭路残差所生成的子群。
因此，连通性增长的全部可观察效应可归约为 \(\{Q_B\}\) 的单调商结构，而与冗余路径细节无关。
\end{conclusion}

\begin{conclusion}[单个闭路残差的合并因子定律：类合并幅度等于其在旧商群中的阶]\label{con:xi-single-cycle-merge-factor-order}
固定预算 \(B\)，考虑在某连通分量内新增一条边 \(e\)（其阈值为 \(B^+\)），使得在 \(G_{\le B^+}\) 中首次形成一个新独立闭路残差向量 \(h\in\ZZ^d\)，满足
\[
h\notin H_B,
\qquad
H_{B^+}=H_B+\ZZ h.
\]
记旧商群 \(Q_B=\ZZ^d/H_B\) 中元素 \([h]\) 的阶为 \(m\)（即最小正整数 \(m\) 使 \(mh\in H_B\)）。
则新商群满足
\[
Q_{B^+}\cong Q_B/\langle [h]\rangle,
\]
从而在该分量上零残差等价类数的最大可能合并因子精确等于 \(m\)：
任意两点 \(u,v\) 若在预算 \(B\) 下地址差落入 \(\langle[h]\rangle\)，则在 \(B^+\) 下必被合并；且除由 \(\langle[h]\rangle\) 强制的合并外，不会出现额外合并。
\end{conclusion}

\begin{conclusion}[素分解独立性：连通性层级按各素数的 \(p\)-主部分可分离]\label{con:xi-prime-decomposition-independence}
仍在结论 \ref{con:xi-arithmetic-address-existence-uniqueness}--\ref{con:xi-single-cycle-merge-factor-order} 的设定中，设在某预算范围内 \(H_B\) 达到满秩 \(d\)，从而 \(Q_B\) 为有限阿贝尔群。
对任一素数 \(p\)，记 \(Q_B^{(p)}\) 为其 \(p\)-主子群。
则算术地址诱导分解 \(a_B=(a_B^{(p)})_p\)，并且
\[
u\sim_B v\quad\Longleftrightarrow\quad
\forall p,\ a_B^{(p)}(u)=a_B^{(p)}(v).
\]
同时，预算提升所导致的类合并在各 \(p\)-主部分上相互独立：若新增闭路残差 \([h]\in Q_B\) 的合并因子为 \(m=\mathrm{ord}([h])\)，则其素分解
\(
m=\prod_p p^{\nu_p(m)}
\)
精确对应于各 \(Q_B^{(p)}\) 上的合并深度 \(\nu_p(m)\)。
因此，连通性增长可在数论意义上逐素数解耦，并由一组彼此独立的 \(p\)-进塔控制。
\end{conclusion}

\begin{conclusion}[最优指纹位长的算术闭式：由商群阶而非节点规模决定]\label{con:xi-optimal-fingerprint-bitlength-arithmetic-closed-form}
在满秩情形下，设该分量的算术地址像生成全体 \(Q_B\)（例如存在 \(d\) 个顶点使其地址差生成 \(\ZZ^d/H_B\)）。
则任何在预算 \(B\) 下对零残差等价类做无碰撞确定编码的指纹长度 \(k_B\) 必满足
\[
k_B\ge \left\lceil\log_2|Q_B|\right\rceil,
\]
且存在达到该下界的构造。
更进一步，若允许两两碰撞概率上界为 \(\delta_B\)，则最小可达位长满足
\[
k_B=\left\lceil\log_2|Q_B|+\log_2\frac1{\delta_B}\right\rceil.
\]
因此，在该模型中最优指纹复杂度由闭式量
\(
|Q_B|=[\ZZ^d:H_B]
\)
决定，而与图中冗余连边、路径数量及 \(|V|\) 的增长无关（与结论 \ref{con:xi-multiscale-fingerprint-optimal-bitlength} 的信息论下界相容，并在此处给出算术可达闭式）。
\end{conclusion}

\begin{conclusion}[稳定化判据：连通性停止增长当且仅当闭路格的 Smith 不变量冻结]\label{con:xi-smith-invariants-freeze-stabilization}
在同一分量上令 \(I_B:=[\ZZ^d:H_B]\)（满秩时为有限指数）。
则随预算增长 \(I_B\) 单调不增，且存在预算 \(B_\ast\) 使对所有 \(B\ge B_\ast\) 等价关系不再合并，当且仅当存在 \(B_\ast\) 使对所有 \(B\ge B_\ast\) 有 \(I_B=I_{B_\ast}\)；
等价地，子格 \(H_B\) 的 Smith 不变量序列 \((s_1(B)\mid\cdots\mid s_d(B))\) 在 \(B\ge B_\ast\) 后恒定。
因此，连通性是否已饱和可由有限阿贝尔群不变量的冻结判定，而无需追踪图的全部冗余演化。
\end{conclusion}

\begin{conclusion}[单遍治理算法：同时输出临界合并链与算术地址基的可验证生成元]\label{con:xi-one-pass-governance-algorithm}
对所有边按阈值 \(c(e)\) 非降排序并逐一扫描，并在每一连通分量内维护：
并查集结构（跟踪 \(G_{\le B}\) 的连通分量），以及闭路残差生成元矩阵的增量 Smith 形式（跟踪子格 \(H_B\subseteq\ZZ^d\) 的不变量与指数 \(I_B\)）。
则存在一单遍算法在时间
\(
O(|E|\cdot \mathrm{poly}(d)\cdot \log R)
\)
内（其中 \(R\) 为残差坐标幅度上界）同时输出：
\begin{enumerate}
  \item 所有临界预算点及其对应合并因子（由结论 \ref{con:xi-single-cycle-merge-factor-order} 的阶 \(m\) 计算）；
  \item 每个预算层上的最小可验证闭路生成元集合（给出 \(H_B\) 的 Smith 基）；
  \item 每个顶点的算术地址 \(a_B(v)\)（以不变量分量坐标表示），从而可在 \(O(d)\) 时间内判定任意一对顶点是否零残差等价。
\end{enumerate}
该结论给出一条可计算且可复核的治理定理：连通性层级与审计结构可被同一遍扫描同时抽取并输出证书化生成元。
\end{conclusion}

\begin{conclusion}[合并深度下界：闭路残差的 \(p\)-进精度限制了连通性收敛速度]\label{con:xi-merge-depth-lower-bound-padic}
在满秩情形下固定素数 \(p\)。
设预算提升过程中每一次引入的新独立闭路残差 \([h]\in Q_B\) 的 \(p\)-阶至多为 \(p^{t}\)，即对所有临界步成立
\(
\nu_p(\mathrm{ord}([h]))\le t
\)。
记在某初始预算 \(B_0\) 下 \(p\)-主部分的指数为 \(|Q_{B_0}^{(p)}|=p^{E}\)。
则要使 \(p\)-主部分完全塌缩为平凡（从而在该分量上消除全部 \(p\)-进差异所导致的类分裂），临界合并步数至少为
\[
\left\lceil \frac{E}{t}\right\rceil.
\]
因此，在可审计闭路证书的 \(p\)-进阶受限时，连通性增长具有不可突破的收敛速度下界；该下界以纯粹算术量 \(E\) 与单步可提供的 \(p\)-进信息量 \(t\) 精确刻画。
\end{conclusion}


% (User input window)

\paragraph{Bethe 平衡、Riccati 极点消去与势系数的 \texorpdfstring{$\tau$}{tau}-梯度恒等式}
本段在 Heine--Stieltjes--Van Vleck 接口上给出两个互补表述：其一把平衡方程等价改写为 Riccati 量在锚点处的极点消去；其二把 Schr\"odinger 势的一阶系数写成“结果式--Vandermonde”对数导数，从而把系数约束压缩为一个可微势函数的梯度；相关残差线性约束的并行陈述见 \cite{Internal2026ToeplitzScanHankelBlaschkeVanVleck}。

\begin{conclusion}[Bethe 平衡方程的“极点消去”刻画：Riccati 量的无锚极点充要条件]\label{con:xi-terminal-riccati-pole-cancellation-bethe-equilibrium}
取互异点 \(p_k=\gamma_k-\mathrm{i}a_k\in\CC_{-}\)（\(a_k>0\)）并设
\[
A(z):=\prod_{k=1}^{\kappa}(z-p_k)(z-\overline{p_k})\in\RR[z].
\]
取互异实点 \(x_1,\dots,x_\kappa\in\RR\)，令
\[
y(z):=\prod_{j=1}^{\kappa}(z-x_j),\qquad
W(z):=\frac{y'(z)}{y(z)}=\sum_{j=1}^{\kappa}\frac{1}{z-x_j}.
\]
定义 Riccati 量
\[
R(z):=W'(z)+W(z)^2-\frac{A'(z)}{A(z)}W(z).
\]
则对每个 \(i\in\{1,\dots,\kappa\}\)，\(R\) 在 \(z=x_i\) 的主部为
\[
R(z)=\frac{2\sum_{j\ne i}\frac{1}{x_i-x_j}-\frac{A'(x_i)}{A(x_i)}}{z-x_i}+O(1).
\]
因此下列两条件等价：
\begin{enumerate}
  \item \(R\) 在全部锚点 \(\{x_i\}\) 处解析（即无锚极点）；
  \item 对全部 \(i\) 成立严格平衡方程
  \[
  2\sum_{j\ne i}\frac{1}{x_i-x_j}
  =\frac{A'(x_i)}{A(x_i)}
  =\sum_{k=1}^{\kappa}\left(\frac{1}{x_i-p_k}+\frac{1}{x_i-\overline{p_k}}\right)
  =2\sum_{k=1}^{\kappa}\Re\frac{1}{x_i-p_k}.
  \]
\end{enumerate}
并且当上述等价条件成立时，\(R\) 的奇点仅可能来自 \(\frac{A'}{A}\)，即 \(\{p_k,\overline{p_k}\}\)，且皆为至多一阶极点；同时 \(R(z)=O(z^{-2})\)（\(z\to\infty\)）。
从而存在唯一多项式 \(B\in\RR[z]\)（\(\deg B\le 2\kappa-2\)）使
\[
R(z)=\frac{B(z)}{A(z)},
\qquad\text{并且}\qquad
A(z)y''(z)-A'(z)y'(z)=B(z)y(z).
\]
\end{conclusion}
\begin{proof}[证明]
在 \(z\to x_i\) 时，
\[
W(z)=\frac{1}{z-x_i}+S_i+O(z-x_i),\qquad S_i:=\sum_{j\ne i}\frac{1}{x_i-x_j},
\]
从而
\[
W'(z)+W(z)^2
=-\frac{1}{(z-x_i)^2}+\frac{1}{(z-x_i)^2}+\frac{2S_i}{z-x_i}+O(1)
=\frac{2S_i}{z-x_i}+O(1).
\]
又因 \(A(x_i)\neq 0\) 且 \(\frac{A'}{A}\) 在 \(x_i\) 处解析，故
\[
\frac{A'(z)}{A(z)}W(z)=\frac{A'(x_i)}{A(x_i)}\cdot\frac{1}{z-x_i}+O(1),
\]
相减即得主部公式，因而两条件等价。
\par
若 \(R\) 在 \(\{x_i\}\) 无极点，则 \(R\) 的奇点仅可能来自 \(\frac{A'}{A}\)，即 \(\{p_k,\overline{p_k}\}\)，并且至多为一阶极点。
又由
\(
W(z)=\kappa z^{-1}+O(z^{-2})
\)
与
\(
\frac{A'(z)}{A(z)}=2\kappa z^{-1}+O(z^{-2})
\)
得 \(R(z)=O(z^{-2})\)。
令 \(B:=AR\)，则 \(B\) 为多项式且 \(\deg B\le \deg A-2=2\kappa-2\)。
最后，\(\deg B<\deg A\) 下的有理分解唯一性给出 \(B\) 的唯一性；而恒等式
\(
A y''-A'y'=By
\)
由
\(
W'+W^2=y''/y
\)
与 \(R=B/A\) 直接推出。证毕。
\end{proof}

% (User input window)

\paragraph{Wronskian 完备化与 de Branges 正定核：极点对数据的有限维 Pick/RKHS 化}
沿用结论 \ref{con:xi-terminal-riccati-pole-cancellation-bethe-equilibrium} 的极点对 \(p_k=\gamma_k-\mathrm{i}a_k\in\CC_{-}\) 与首一多项式 \(A(z)=\prod_k(z-p_k)(z-\overline{p_k})\) 的记号。
设存在首一多项式 \(y\in\RR[z]\)（\(\deg y=\kappa\)，实根互异）与首一多项式 \(y_{+}\in\RR[z]\)（\(\deg y_{+}=\kappa+1\)）满足严格 Wronskian 恒等式
\[
\boxed{\;A(z)=y(z)y_{+}'(z)-y'(z)y_{+}(z)\;}
\]
并记 \(y(z)=\prod_{j=1}^{\kappa}(z-x_j)\)（\(x_j\in\RR\)）。

\begin{conclusion}[重心恒等式：Wronskian 数据把锚点重心锁定为极点实部之和]\label{con:xi-terminal-wronskian-centroid-identity}
在上述设定下，必有严格恒等式
\[
\boxed{\;\sum_{j=1}^{\kappa}x_j=\sum_{k=1}^{\kappa}\gamma_k\;}
\]
即 \(y\) 的全部实根之和等于极点对 \(\{p_k,\overline{p_k}\}\) 的实部总和。
\end{conclusion}
\begin{proof}
写
\[
A(z)=z^{2\kappa}+A_1z^{2\kappa-1}+O(z^{2\kappa-2}),\qquad
y(z)=z^{\kappa}+y_1z^{\kappa-1}+O(z^{\kappa-2}),\qquad
y_{+}(z)=z^{\kappa+1}+y_{+1}z^{\kappa}+O(z^{\kappa-1}).
\]
直接展开 \(yy_{+}'-y'y_{+}\) 的 \(z^{2\kappa-1}\) 系数可得其恰为 \(2y_1\)，故 \(A_1=2y_1\)。
另一方面，
\(
A_1=-\sum_{k=1}^{\kappa}(p_k+\overline{p_k})=-2\sum_k\gamma_k
\)，而 \(y_1=-\sum_j x_j\)。
代回即得 \(\sum_j x_j=\sum_k\gamma_k\)。证毕。
\end{proof}

\begin{conclusion}[第二对称和约束：\texorpdfstring{$A$}{A} 的次首项系数给出 \texorpdfstring{$y$}{y} 与 \texorpdfstring{$y_{+}$}{y+} 的二阶耦合]\label{con:xi-terminal-wronskian-second-symmetric-sum-coupling}
在上述设定下，展开
\[
y(z)=z^{\kappa}+y_1z^{\kappa-1}+y_2z^{\kappa-2}+O(z^{\kappa-3}),\qquad
y_{+}(z)=z^{\kappa+1}+y_{+1}z^{\kappa}+y_{+2}z^{\kappa-1}+O(z^{\kappa-2}),
\]
\[
A(z)=z^{2\kappa}+A_1z^{2\kappa-1}+A_2z^{2\kappa-2}+O(z^{2\kappa-3}),
\]
则必有精确恒等式
\[
\boxed{\;A_2=-y_{+2}+y_1y_{+1}+3y_2.\;}
\]
若进一步将 \(y_{+}\) 分解为
\(
y_{+}(z)=\prod_{m=0}^{\kappa}(z-b_m)
\)
并记初等对称和
\[
e_1(x):=\sum_{j=1}^{\kappa}x_j,\quad e_2(x):=\sum_{1\le i<j\le \kappa}x_ix_j,\qquad
e_1(b):=\sum_{m=0}^{\kappa}b_m,\quad e_2(b):=\sum_{0\le i<j\le \kappa}b_ib_j,
\]
则同一恒等式可写为
\[
\boxed{\;A_2=-e_2(b)+e_1(x)e_1(b)+3e_2(x).\;}
\]
\end{conclusion}
\begin{proof}
将 \(yy_{+}'-y'y_{+}\) 展开到 \(z^{2\kappa-2}\) 阶并逐项比较即可得到
\(
A_2=-y_{+2}+y_1y_{+1}+3y_2
\)。
再由
\(
y_1=-e_1(x),\ y_2=e_2(x),\ y_{+1}=-e_1(b),\ y_{+2}=e_2(b)
\)
代入即得根形式。证毕。
\end{proof}

\begin{conclusion}[Hermite--Biehler 化：Wronskian 数据等价于一类多项式 Pick 映射]\label{con:xi-terminal-wronskian-hermite-biehler-pick-map}
在上述设定下，定义有理函数
\[
w(z):=\frac{y_{+}(z)}{y(z)}.
\]
并定义多项式
\[
E(z):=y(z)-\mathrm{i}\,y_{+}(z)\in\CC[z],\qquad
E^\sharp(z):=\overline{E(\overline z)}.
\]
则成立：
\begin{enumerate}
  \item \(w\) 为 Pick 函数（上半平面自映）：对一切 \(z\in\CC_{+}\) 有 \(\Im w(z)>0\)；
  \item \(E\) 为 Hermite--Biehler 多项式：对一切 \(z\in\CC_{+}\) 有 \(|E(z)|>|E^\sharp(z)|\)；
  \item 二者满足 Cayley 型等价关系
  \[
  \boxed{\;
  \frac{E^\sharp(z)}{E(z)}=\frac{1+\mathrm{i}\,w(z)}{1-\mathrm{i}\,w(z)},
  \qquad
  w(z)=\mathrm{i}\,\frac{E(z)-E^\sharp(z)}{E(z)+E^\sharp(z)}\; } .
  \]
\end{enumerate}
\end{conclusion}
\begin{proof}
由 Wronskian 恒等式得
\[
w'(z)=\frac{y(z)y_{+}'(z)-y'(z)y_{+}(z)}{y(z)^2}=\frac{A(z)}{y(z)^2}.
\]
对实 \(x\notin\{x_j\}\)，有 \(A(x)=\prod_k|x-p_k|^2>0\) 且 \(y(x)^2>0\)，故 \(w'(x)>0\)。
从而 \(w\) 在每个实轴连通分支上严格递增，并在每个 \(x_j\) 处具有简单极点；并且由 \(\deg y_{+}=\deg y+1\) 得 \(w(z)=z+O(1)\)（\(z\to\infty\)）。
\par
对每个 \(x_j\) 有 \(A(x_j)=-y'(x_j)\,y_{+}(x_j)>0\)，从而
\(
\operatorname*{Res}_{z=x_j}w(z)=y_{+}(x_j)/y'(x_j)<0
\)。
记 \(\mu_j:=-y_{+}(x_j)/y'(x_j)>0\)。
于是 \(w\) 的部分分式可写为
\[
w(z)=z+c-\sum_{j=1}^{\kappa}\frac{\mu_j}{z-x_j},\qquad c\in\RR.
\]
从而对 \(z=x+\mathrm{i}y\in\CC_{+}\) 有
\[
\Im w(z)=y+\sum_{j=1}^{\kappa}\mu_j\,\frac{y}{(x-x_j)^2+y^2}>0,
\]
得到 (1)。
\par
再由
\[
\frac{E^\sharp(z)}{E(z)}
=\frac{y(z)+\mathrm{i}y_{+}(z)}{y(z)-\mathrm{i}y_{+}(z)}
=\frac{1+\mathrm{i}w(z)}{1-\mathrm{i}w(z)},
\]
并记 \(w(z)=u+\mathrm{i}v\)（\(v>0\)），则
\[
\left|\frac{1+\mathrm{i}w(z)}{1-\mathrm{i}w(z)}\right|^2
=\frac{(1-v)^2+u^2}{(1+v)^2+u^2}<1,
\]
从而 \(|E^\sharp(z)|<|E(z)|\)，得 (2)。
最后 (3) 为代数恒等式与其反解。证毕。
\end{proof}

\begin{conclusion}[Wronskian 核的正定性与离散 Mercer 分解：\texorpdfstring{$y_{+}$}{y+} 的零点给出完备正交采样]\label{con:xi-terminal-wronskian-kernel-discrete-mercer-parseval}
在结论 \ref{con:xi-terminal-wronskian-hermite-biehler-pick-map} 的设定下，定义实核（对 \(s\neq t\)）
\[
K(s,t):=\frac{y_{+}(s)y(t)-y(s)y_{+}(t)}{\pi(s-t)},\qquad
K(x,x):=\frac{y(x)y_{+}'(x)-y'(x)y_{+}(x)}{\pi}=\frac{A(x)}{\pi}.
\]
则：
\begin{enumerate}
  \item \(K\) 为正定核；
  \item 若 \(x_1,\dots,x_{\kappa}\) 为 \(y\) 的零点，则对 \(i\neq j\) 有 \(K(x_i,x_j)=0\)，且
  \[
  K(x_i,x_i)=\frac{A(x_i)}{\pi}=\frac{-y'(x_i)\,y_{+}(x_i)}{\pi}>0;
  \]
  \item \(y_{+}\) 具有 \(\kappa+1\) 个互异实根 \(b_0<\cdots<b_{\kappa}\)；
  在该记号下，对 \(i\neq j\) 有 \(K(b_i,b_j)=0\)，且
  \[
  K(b_j,b_j)=\frac{A(b_j)}{\pi}=\frac{y(b_j)\,y_{+}'(b_j)}{\pi}>0.
  \]
  \item （离散 Mercer 分解）对一切 \(x,y\in\RR\) 有严格恒等式
  \[
  \boxed{\;
  K(x,y)=\sum_{m=0}^{\kappa}\frac{K(x,b_m)\,K(b_m,y)}{K(b_m,b_m)}\; }.
  \]
  \item 令 \(\mathcal{H}_K\) 为由 \(K\) 生成的再现核 Hilbert 空间，则
  \(\{K(\cdot,b_m)\}_{m=0}^{\kappa}\) 为 \(\mathcal{H}_K\) 的一组正交基，且对任意 \(f\in\mathcal{H}_K\) 成立插值展开与 Parseval 恒等式
  \[
  \boxed{\;
  f(x)=\sum_{m=0}^{\kappa} f(b_m)\,\frac{K(x,b_m)}{K(b_m,b_m)},
  \qquad
  \|f\|_{\mathcal{H}_K}^{2}=\sum_{m=0}^{\kappa}\frac{|f(b_m)|^{2}}{K(b_m,b_m)}\; }.
  \]
\end{enumerate}
\end{conclusion}
\begin{proof}
由结论 \ref{con:xi-terminal-wronskian-hermite-biehler-pick-map}，\(w\) 在每个区间 \((x_j,x_{j+1})\) 上严格递增并在端点处发散，且在无穷远处满足 \(w(x)\sim x\)。
因此在每个区间 \((-\infty,x_1),(x_1,x_2),\dots,(x_\kappa,\infty)\) 上，\(w\) 为到 \(\RR\) 的双射，从而在每一段恰有一个零点；这给出 \(y_{+}\) 的 \(\kappa+1\) 个互异实根 \(b_0<\cdots<b_\kappa\)。
相应的对角与非对角退化式由定义代入 \(y(x_i)=0\) 或 \(y_{+}(b_j)=0\) 得到。
\par
为证 Mercer 分解，固定 \(t\in\RR\)，函数 \(x\mapsto K(x,t)\) 为次数 \(\le \kappa\) 的多项式：分子关于 \(x\) 的次数为 \(\kappa+1\) 且在 \(x=t\) 处可被 \(x-t\) 整除。
右端同样为次数 \(\le \kappa\) 的多项式，并且对每个 \(\ell\in\{0,\dots,\kappa\}\) 有
\[
\text{右端}(b_\ell,t)=\sum_{m=0}^{\kappa}\frac{K(b_\ell,b_m)K(b_m,t)}{K(b_m,b_m)}=K(b_\ell,t),
\]
因为当 \(\ell\neq m\) 时 \(K(b_\ell,b_m)=0\)，且 \(K(b_\ell,b_\ell)\neq 0\)。
两边作为次数 \(\le\kappa\) 的多项式在 \(\kappa+1\) 个互异点 \(\{b_\ell\}\) 上一致，故恒等。
\par
将 Mercer 分解代入任意有限二次型即可得到 \(K\) 的正定性。
最后，Gram 矩阵 \((K(b_i,b_j))\) 为对角且对角元为正，故 \(\{K(\cdot,b_m)\}\) 正交且线性无关；由 Mercer 恒等式可知其张成整个 \(\mathcal{H}_K\)，从而为正交基。
再由再现性 \(f(b_m)=\langle f,K(\cdot,b_m)\rangle\) 得到插值展开与 Parseval 恒等式。证毕。
\end{proof}

% (User input window)

\begin{conclusion}[压缩乘法算子：谱为 \texorpdfstring{$\{b_m\}$}{b\_m} 且极小多项式为 \texorpdfstring{$y_{+}$}{y+}]\label{con:xi-terminal-wronskian-compressed-multiplication-spectrum-minpoly}
沿用结论 \ref{con:xi-terminal-wronskian-kernel-discrete-mercer-parseval} 的记号。
令 \(y_{+}(z)=\prod_{m=0}^{\kappa}(z-b_m)\)（\(b_m\in\RR\) 两两互异），并在 \(\mathcal{H}_K\) 上定义线性算子 \(M:\mathcal{H}_K\to\mathcal{H}_K\)：
\[
(Mf)(x):=\sum_{m=0}^{\kappa} b_m\,f(b_m)\,\frac{K(x,b_m)}{K(b_m,b_m)}.
\]
则：
\begin{enumerate}
  \item \(M\) 自伴且可对角化；
  \item 令 \(e_m(x):=K(x,b_m)/\sqrt{K(b_m,b_m)}\)，则 \(\{e_m\}_{m=0}^{\kappa}\) 为 \(\mathcal{H}_K\) 的一组正交归一基，且 \(Me_m=b_m e_m\)；
  \item \(M\) 的谱恰为 \(\{b_m\}_{m=0}^{\kappa}\)，并且 \(M\) 的极小多项式等于 \(y_{+}\)，从而 \(y_{+}(M)=0\)。
\end{enumerate}
\end{conclusion}
\begin{proof}[证明]
由结论 \ref{con:xi-terminal-wronskian-kernel-discrete-mercer-parseval} 的 Parseval 结构，\(\{K(\cdot,b_m)\}\) 为 \(\mathcal{H}_K\) 的正交基，且
\(K(b_n,b_m)=0\)（\(n\neq m\)）、\(K(b_m,b_m)>0\)。
对基向量 \(K(\cdot,b_m)\) 计算：
\[
M K(\cdot,b_m)
=\sum_{m'=0}^{\kappa} b_{m'}\,K(b_{m'},b_m)\,\frac{K(\cdot,b_{m'})}{K(b_{m'},b_{m'})}
=b_m\,K(\cdot,b_m),
\]
因此 \(K(\cdot,b_m)\)（等价地 \(e_m\)）为特征向量，特征值为 \(b_m\)；归一化得到 \(\{e_m\}\) 为正交归一基。
于是 \(M\) 在该基下为实对角矩阵 \(\mathrm{diag}(b_0,\dots,b_\kappa)\)，从而自伴且可对角化，谱为 \(\{b_m\}\)。
最后，\(y_{+}\) 首一且零点集为 \(\{b_m\}\)，故 \(y_{+}(M)=0\)；又因 \(\{b_m\}\) 互异，极小多项式必须包含全部一次因子 \(z-b_m\)，从而等于 \(y_{+}\)。
证毕。
\end{proof}

\begin{conclusion}[Weyl 函数的完全闭式：\texorpdfstring{$m(z)=-y/y_{+}$}{m(z)=-y/y+}]\label{con:xi-terminal-wronskian-weyl-function-closed-form}
沿用结论 \ref{con:xi-terminal-wronskian-compressed-multiplication-spectrum-minpoly} 的设定，并令
\(
\psi:=y/\sqrt{\pi}\in\mathcal{H}_K
\)。
对 \(z\in\CC\setminus\{b_0,\dots,b_\kappa\}\) 定义 Weyl 函数
\[
m(z):=\langle (M-zI)^{-1}\psi,\psi\rangle_{\mathcal{H}_K}.
\]
则成立严格恒等式
\[
\boxed{\;m(z)=-\frac{y(z)}{y_{+}(z)}\;}
\]
并等价地具有完全部分分式展开
\[
-\frac{y(z)}{y_{+}(z)}=\sum_{m=0}^{\kappa}\frac{c_m}{b_m-z},
\qquad
c_m:=\frac{y(b_m)}{y_{+}'(b_m)}=\frac{y(b_m)^2}{A(b_m)}.
\]
\end{conclusion}
\begin{proof}[证明]
在结论 \ref{con:xi-terminal-wronskian-compressed-multiplication-spectrum-minpoly} 的正交归一特征基 \(\{e_m\}\) 下，
\((M-zI)^{-1}e_m=(b_m-z)^{-1}e_m\)，从而
\[
m(z)=\sum_{m=0}^{\kappa}\frac{|\langle \psi,e_m\rangle|^2}{b_m-z}.
\]
由再现性 \(f(b_m)=\langle f,K(\cdot,b_m)\rangle\) 与 \(e_m=K(\cdot,b_m)/\sqrt{K(b_m,b_m)}\) 得
\(
\langle \psi,e_m\rangle=\psi(b_m)/\sqrt{K(b_m,b_m)}
\)。
又由 \(\psi(b_m)=y(b_m)/\sqrt{\pi}\) 以及
\(
K(b_m,b_m)=A(b_m)/\pi=y(b_m)y_{+}'(b_m)/\pi
\)
从而得到
\[
|\langle \psi,e_m\rangle|^2
=\frac{y(b_m)^2/\pi}{A(b_m)/\pi}
=\frac{y(b_m)^2}{A(b_m)}
=\frac{y(b_m)}{y_{+}'(b_m)}
=:c_m.
\]
于是 \(m(z)=\sum_m c_m/(b_m-z)\)。

由于 \(m(z)=O(z^{-1})\)（\(z\to\infty\)）且其极点仅为 \(\{b_m\}\) 上的一阶极点，
可写成 \(m(z)=P(z)/y_{+}(z)\)，其中 \(\deg P\le \kappa\)。
在 \(z=b_m\) 处比较留数得到
\[
\frac{P(b_m)}{y_{+}'(b_m)}=\operatorname*{Res}_{z=b_m}m(z)=-c_m=-\frac{y(b_m)}{y_{+}'(b_m)},
\]
从而 \(P(b_m)=-y(b_m)\) 对全部 \(m\) 成立。
由于 \(P\) 与 \(-y\) 均为次数 \(\le\kappa\) 的多项式且在 \(\kappa+1\) 个互异点 \(\{b_m\}\) 上取值一致，必有 \(P\equiv -y\)，即 \(m(z)=-y(z)/y_{+}(z)\)。
证毕。
\end{proof}

\begin{conclusion}[正谱权与矩恒等式：\texorpdfstring{$\{c_m\}$}{c\_m} 为概率权且编码 \texorpdfstring{$-y/y_{+}$}{-y/y+} 的无穷远展开]\label{con:xi-terminal-wronskian-spectral-weights-moment-identities}
沿用结论 \ref{con:xi-terminal-wronskian-weyl-function-closed-form} 的记号。
则：
\begin{enumerate}
  \item 对所有 \(m\) 有 \(c_m>0\)，且 \(\sum_{m=0}^{\kappa}c_m=1\)；
  \item 若把
  \[
  -\frac{y(z)}{y_{+}(z)}=-\frac{1}{z}-\sum_{n\ge 1}\frac{\mu_n}{z^{n+1}}
  \qquad (z\to\infty),
  \]
  则各阶矩满足 \(\mu_n=\sum_{m=0}^{\kappa}c_m b_m^{n}\)（\(n\ge 1\)）；
  \item 若写
  \[
  y(z)=z^{\kappa}+y_1 z^{\kappa-1}+y_2 z^{\kappa-2}+O(z^{\kappa-3}),\quad
  y_{+}(z)=z^{\kappa+1}+y_{+1}z^{\kappa}+y_{+2}z^{\kappa-1}+O(z^{\kappa-2}),
  \]
  则
  \[
  \sum_{m=0}^{\kappa}c_m b_m=y_1-y_{+1},\qquad
  \sum_{m=0}^{\kappa}c_m b_m^{2}=y_2-y_{+2}-y_1y_{+1}+y_{+1}^{2}.
  \]
\end{enumerate}
\end{conclusion}
\begin{proof}[证明]
由 Wronskian 恒等式 \(A=yy_{+}'-y'y_{+}\) 得
\[
w'(x)=\left(\frac{y_{+}}{y}\right)'(x)=\frac{A(x)}{y(x)^2}>0,
\qquad x\in\RR\setminus\{x_j\}.
\]
在 \(x=b_m\) 处因 \(y_{+}(b_m)=0\) 有
\(
w'(b_m)=y_{+}'(b_m)/y(b_m)>0
\)，故 \(c_m=y(b_m)/y_{+}'(b_m)>0\)。

另一方面，由结论 \ref{con:xi-terminal-wronskian-weyl-function-closed-form} 的部分分式展开，
\[
\sum_{m=0}^{\kappa}\frac{c_m}{b_m-z}
=-\frac{1}{z}\sum_{m}c_m-\frac{1}{z^2}\sum_{m}c_m b_m-\cdots .
\]
而 \(-y(z)/y_{+}(z)=-z^{-1}+O(z^{-2})\)（\(z\to\infty\)），比较 \(z^{-1}\) 系数得 \(\sum_m c_m=1\)。
进一步比较一般系数得到 \(\mu_n=\sum_m c_m b_m^{n}\)。

最后对 \(-y/y_{+}\) 作级数除法可得
\[
-\frac{y(z)}{y_{+}(z)}
=-\frac{1}{z}-\frac{y_1-y_{+1}}{z^2}
-\frac{y_2-y_{+2}-y_1y_{+1}+y_{+1}^2}{z^3}+O(z^{-4}),
\]
与矩展开系数比较即得两条显式公式。证毕。
\end{proof}

\endinput


\begin{conclusion}[Wronskian–结果式分解：\texorpdfstring{$\operatorname{Res}(A,y_{+})$}{Res(A,y+)} 的判别式乘积公式]\label{con:xi-terminal-wronskian-resultant-discriminant-factorization}
沿用上述设定，并设 \(y_{+}\) 的实根为 \(b_0,\dots,b_\kappa\) 且互异。
则有严格恒等式
\[
\boxed{\;
\operatorname{Res}(A,y_{+})
=\prod_{m=0}^{\kappa}A(b_m)
=\prod_{k=1}^{\kappa}y_{+}(p_k)\,y_{+}(\overline{p_k})
=\prod_{k=1}^{\kappa}|y_{+}(p_k)|^{2}\ >\ 0\; }.
\]
并且成立分解
\[
\boxed{\;
\operatorname{Res}(A,y_{+})
=\operatorname{Res}(y,y_{+})\cdot
\Bigl((-1)^{\frac{\kappa(\kappa+1)}{2}}\mathrm{Disc}(y_{+})\Bigr)\; },
\]
其中
\(
\mathrm{Disc}(y_{+})=\prod_{0\le i<j\le \kappa}(b_i-b_j)^2>0
\)。
\end{conclusion}
\begin{proof}
第一组等式由结果式定义与 \(A\) 的根集合为 \(\{p_k,\overline{p_k}\}\) 直接得到；并且 \(\deg A\cdot\deg y_{+}=2\kappa(\kappa+1)\) 为偶数，故结果式对交换参数无额外符号。
\par
对第二组等式，注意 \(y_{+}(b_m)=0\) 代入 Wronskian 得
\[
A(b_m)=y(b_m)\,y_{+}'(b_m).
\]
因此
\(
\prod_m A(b_m)
=\bigl(\prod_m y(b_m)\bigr)\bigl(\prod_m y_{+}'(b_m)\bigr)
=\operatorname{Res}(y,y_{+})\cdot\prod_m y_{+}'(b_m)
\)。
而对首一多项式 \(y_{+}\) 有恒等式
\[
\prod_{m=0}^{\kappa} y_{+}'(b_m)=(-1)^{\frac{\kappa(\kappa+1)}{2}}\mathrm{Disc}(y_{+}),
\]
代回即得。证毕。
\end{proof}

\begin{conclusion}[采样对角权的结果式--判别式乘积公式：\texorpdfstring{$\prod_m K(b_m,b_m)$}{prod K(bm,bm)} 与 \texorpdfstring{$\prod_j K(x_j,x_j)$}{prod K(xj,xj)} 的闭式]\label{con:xi-terminal-wronskian-diagonal-sampling-resultant-discriminant-product}
沿用结论 \ref{con:xi-terminal-wronskian-kernel-discrete-mercer-parseval} 的记号，
并记 \(x_1,\dots,x_\kappa\) 为 \(y\) 的互异实根、\(b_0,\dots,b_\kappa\) 为 \(y_{+}\) 的互异实根。
则成立严格乘积恒等式
\[
\boxed{\;
\prod_{m=0}^{\kappa}K(b_m,b_m)
=\pi^{-(\kappa+1)}\operatorname{Res}(A,y_{+})
=\pi^{-(\kappa+1)}(-1)^{\frac{\kappa(\kappa+1)}{2}}\operatorname{Res}(y,y_{+})\,\mathrm{Disc}(y_{+})\;},
\]
以及
\[
\boxed{\;
\prod_{j=1}^{\kappa}K(x_j,x_j)
=\pi^{-\kappa}\operatorname{Res}(A,y)
=\pi^{-\kappa}(-1)^{\frac{\kappa(\kappa+1)}{2}}\operatorname{Res}(y,y_{+})\,\mathrm{Disc}(y)\;},
\]
其中
\(
\mathrm{Disc}(y)=\prod_{1\le i<j\le \kappa}(x_i-x_j)^2
\)、
\(
\mathrm{Disc}(y_{+})=\prod_{0\le i<j\le \kappa}(b_i-b_j)^2
\)。
\end{conclusion}
\begin{proof}[证明]
由 \(K(x,x)=A(x)/\pi\)（结论 \ref{con:xi-terminal-wronskian-kernel-discrete-mercer-parseval}）得
\[
\prod_{m=0}^{\kappa}K(b_m,b_m)=\pi^{-(\kappa+1)}\prod_{m=0}^{\kappa}A(b_m)
=\pi^{-(\kappa+1)}\operatorname{Res}(A,y_{+}).
\]
并由结论 \ref{con:xi-terminal-wronskian-resultant-discriminant-factorization} 得
\(
\operatorname{Res}(A,y_{+})
=(-1)^{\frac{\kappa(\kappa+1)}{2}}\operatorname{Res}(y,y_{+})\,\mathrm{Disc}(y_{+})
\)，代回即得第一条恒等式。
\par
同理，
\[
\prod_{j=1}^{\kappa}K(x_j,x_j)=\pi^{-\kappa}\prod_{j=1}^{\kappa}A(x_j)
=\pi^{-\kappa}\operatorname{Res}(A,y).
\]
又因 \(y(x_j)=0\)，代入 Wronskian 恒等式 \(A=yy_{+}'-y'y_{+}\) 得
\(
A(x_j)=-y'(x_j)\,y_{+}(x_j)
\)。
因此
\[
\operatorname{Res}(A,y)
=\prod_{j=1}^{\kappa}A(x_j)
=(-1)^{\kappa}\Bigl(\prod_{j=1}^{\kappa}y'(x_j)\Bigr)\Bigl(\prod_{j=1}^{\kappa}y_{+}(x_j)\Bigr)
=(-1)^{\kappa}\Bigl(\prod_{j=1}^{\kappa}y'(x_j)\Bigr)\operatorname{Res}(y,y_{+}).
\]
对首一多项式 \(y\) 有
\(
\prod_{j=1}^{\kappa}y'(x_j)=(-1)^{\frac{\kappa(\kappa-1)}{2}}\mathrm{Disc}(y)
\)，
合并符号 \( (-1)^{\kappa+\frac{\kappa(\kappa-1)}{2}}=(-1)^{\frac{\kappa(\kappa+1)}{2}} \) 即得
\(
\operatorname{Res}(A,y)
=(-1)^{\frac{\kappa(\kappa+1)}{2}}\operatorname{Res}(y,y_{+})\,\mathrm{Disc}(y)
\)。
代回即得第二条恒等式。证毕。
\end{proof}

\begin{conclusion}[纯 Lagrange 形式的核恒等式：\texorpdfstring{$K$}{K} 的显式采样展开与双变量代数恒等式]\label{con:xi-terminal-wronskian-kernel-lagrange-sampling-identity}
沿用结论 \ref{con:xi-terminal-wronskian-kernel-discrete-mercer-parseval} 的设定，并令 \(b_0,\dots,b_\kappa\) 为 \(y_{+}\) 的互异实根。
定义 Lagrange 基多项式
\[
L_m(z):=\frac{y_{+}(z)}{(z-b_m)\,y_{+}'(b_m)}\qquad (m=0,\dots,\kappa).
\]
则对一切 \(x,y\in\RR\) 成立显式核展开
\[
\boxed{\;
K(x,y)=\sum_{m=0}^{\kappa}K(b_m,b_m)\,L_m(x)\,L_m(y)
=\sum_{m=0}^{\kappa}\frac{A(b_m)}{\pi}\,L_m(x)\,L_m(y)\;}
\]
并等价地成立纯代数双变量恒等式
\[
\boxed{\;
\frac{y_{+}(x)y(y)-y(x)y_{+}(y)}{x-y}
=y_{+}(x)\,y_{+}(y)\sum_{m=0}^{\kappa}\frac{y(b_m)}{(x-b_m)(y-b_m)\,y_{+}'(b_m)}\; }.
\]
\end{conclusion}
\begin{proof}[证明]
由 \(y_{+}(b_m)=0\) 得
\[
K(x,b_m)=\frac{y_{+}(x)y(b_m)}{\pi(x-b_m)},
\qquad
K(b_m,b_m)=\frac{A(b_m)}{\pi}=\frac{y(b_m)\,y_{+}'(b_m)}{\pi},
\]
从而 \(K(x,b_m)=K(b_m,b_m)\,L_m(x)\)。
将其代入结论 \ref{con:xi-terminal-wronskian-kernel-discrete-mercer-parseval} 的离散 Mercer 分解即得
\(
K(x,y)=\sum_m K(b_m,b_m)L_m(x)L_m(y)
\)。
双变量形式由 \(K\) 的定义与 \(K(b_m,b_m)=y(b_m)y_{+}'(b_m)/\pi\) 整理得到。证毕。
\end{proof}

\begin{conclusion}[顶系数抽取公式：\texorpdfstring{$[z^{\kappa}]$}{[z^k]} 的 Lagrange--导数权表示]\label{con:xi-terminal-wronskian-top-coefficient-extraction}
沿用结论 \ref{con:xi-terminal-wronskian-kernel-lagrange-sampling-identity} 的记号。
令 \(q\in\CC[z]\) 且 \(\deg q\le \kappa\)。
则成立顶系数抽取恒等式
\[
\boxed{\;[z^{\kappa}]\,q(z)=\sum_{m=0}^{\kappa}\frac{q(b_m)}{y_{+}'(b_m)}\; }.
\]
特别地，
\[
\sum_{m=0}^{\kappa}\frac{1}{y_{+}'(b_m)}=0,\qquad
\sum_{m=0}^{\kappa}\frac{b_m}{y_{+}'(b_m)}=0,\qquad
\sum_{m=0}^{\kappa}\frac{y(b_m)}{y_{+}'(b_m)}=1.
\]
\end{conclusion}
\begin{proof}[证明]
对 \(\deg q\le\kappa\) 作 Lagrange 插值：
\[
q(z)=\sum_{m=0}^{\kappa} q(b_m)\,L_m(z)
=\sum_{m=0}^{\kappa}q(b_m)\,\frac{y_{+}(z)}{(z-b_m)\,y_{+}'(b_m)}.
\]
注意 \(y_{+}(z)/(z-b_m)\) 为首一次数 \(\kappa\) 多项式，因此 \(L_m\) 的 \(z^\kappa\) 系数为 \(1/y_{+}'(b_m)\)。
比较两边的 \(z^\kappa\) 系数即得主恒等式。
对 \(q(z)=1\)、\(q(z)=z\)、\(q(z)=y(z)\) 分别应用，并用 \([z^\kappa]1=0\)、\([z^\kappa]z=0\)、\([z^\kappa]y(z)=1\)，即得三条推论。证毕。
\end{proof}

\paragraph{采样权重的矩实现：Stieltjes 变换、Hankel 主块与节点谱}
沿用结论 \ref{con:xi-terminal-wronskian-kernel-lagrange-sampling-identity} 的设定，令 \(b_0,\dots,b_\kappa\) 为 \(y_{+}\) 的互异实根，并定义采样权重
\[
c_m:=\frac{y(b_m)}{y_{+}'(b_m)}\qquad(m=0,\dots,\kappa).
\]
则有部分分式恒等式
\begin{equation}\label{eq:xi-terminal-wronskian-stieltjes-transform-minus-y-over-yplus}
\bar m(z)
:=\sum_{m=0}^{\kappa}\frac{c_m}{b_m-z}
\;=\;-\frac{y(z)}{y_{+}(z)}.
\end{equation}
并由 \(z\to\infty\) 的系数比较得到 \(\sum_{m=0}^{\kappa}c_m=1\)（亦即结论 \ref{con:xi-terminal-wronskian-top-coefficient-extraction} 在 \(q=y\) 时的最后一式）。
定义矩序列
\[
\mu_n:=\sum_{m=0}^{\kappa}c_m\,b_m^{n}\qquad(n\ge 0),
\]
并记其 \((\kappa+1)\times(\kappa+1)\) Hankel 主块与一次移位块为
\[
H_{\kappa}:=\bigl(\mu_{i+j}\bigr)_{0\le i,j\le \kappa},
\qquad
H_{\kappa}^{+}:=\bigl(\mu_{i+j+1}\bigr)_{0\le i,j\le \kappa}.
\]

\begin{conclusion}[全主 Hankel 行列式的消去闭式与正权符号律]\label{con:xi-terminal-wronskian-hankel-principal-determinant-resultant}
在上述记号下成立严格恒等式
\[
\boxed{\;
\det H_{\kappa}
=(-1)^{\frac{\kappa(\kappa+1)}{2}}\operatorname{Res}(y,y_{+})\; }.
\]
并且在正权口径 \(c_m>0\) 下，\(H_{\kappa}\) 正定，从而有符号律
\[
(-1)^{\frac{\kappa(\kappa+1)}{2}}\operatorname{Res}(y,y_{+})>0.
\]
\end{conclusion}
\begin{proof}[证明要点]
令 Vandermonde 矩阵 \(V=(b_m^{\,i})_{0\le i\le \kappa,\,0\le m\le \kappa}\)，并记 \(C:=\mathrm{diag}(c_0,\dots,c_\kappa)\)。
由 \(\mu_{i+j}=\sum_{m=0}^{\kappa}c_m b_m^{\,i}b_m^{\,j}\) 得
\(
H_{\kappa}=VCV^{\mathsf T}
\)，
因而
\[
\det H_{\kappa}=(\det V)^2\prod_{m=0}^{\kappa}c_m.
\]
另一方面，由 \(y_{+}\) 首一且零点为 \(\{b_m\}\) 得
\[
\operatorname{Res}(y,y_{+})=\prod_{m=0}^{\kappa}y(b_m),
\qquad
\prod_{m=0}^{\kappa}c_m=\frac{\prod_m y(b_m)}{\prod_m y_{+}'(b_m)}.
\]
又有标准恒等式
\[
(\det V)^2=\mathrm{Disc}(y_{+}),
\qquad
\prod_{m=0}^{\kappa}y_{+}'(b_m)=(-1)^{\frac{\kappa(\kappa+1)}{2}}\mathrm{Disc}(y_{+}),
\]
代入并消去 \(\mathrm{Disc}(y_{+})\) 即得所示闭式。正权情形下 \(H_{\kappa}\) 为 Gram 矩阵，故 \(\det H_{\kappa}>0\)，从而得到符号律。证毕。
\end{proof}

\begin{conclusion}[移位 Hankel 线性笔的节点对角化与 \texorpdfstring{$H_{\kappa}$}{H}-自伴随]\label{con:xi-terminal-wronskian-hankel-shift-diagonalization}
令
\[
A_{\kappa}:=H_{\kappa}^{-1}H_{\kappa}^{+}.
\]
则存在精确相似对角化
\[
\boxed{\;
A_{\kappa}=(V^{\mathsf T})^{-1}\,\mathrm{diag}(b_0,\dots,b_\kappa)\,V^{\mathsf T}\;},
\]
从而特征多项式满足
\[
\chi_{A_{\kappa}}(z)=\prod_{m=0}^{\kappa}(z-b_m)=y_{+}(z).
\]
此外 \(A_{\kappa}\) 满足 \(H_{\kappa}\)-自伴随关系
\[
\boxed{\;
A_{\kappa}^{\mathsf T}H_{\kappa}=H_{\kappa}A_{\kappa}\; }.
\]
\end{conclusion}
\begin{proof}[证明要点]
同理由 \(\mu_{i+j+1}=\sum_m c_m b_m b_m^{\,i}b_m^{\,j}\) 得
\(
H_{\kappa}^{+}=V\,\mathrm{diag}(c_0b_0,\dots,c_\kappa b_\kappa)\,V^{\mathsf T}
\)。
与 \(H_{\kappa}=VCV^{\mathsf T}\) 合并可得
\[
A_{\kappa}=(VCV^{\mathsf T})^{-1}V\,\mathrm{diag}(cb)\,V^{\mathsf T}
=(V^{\mathsf T})^{-1}\mathrm{diag}(b)\,V^{\mathsf T}.
\]
特征多项式在相似变换下不变，从而 \(\chi_{A_{\kappa}}=y_{+}\)。
最后因 \(H_{\kappa},H_{\kappa}^{+}\) 对称，有
\(
A_{\kappa}^{\mathsf T}H_{\kappa}=(H_{\kappa}^{-1}H_{\kappa}^{+})^{\mathsf T}H_{\kappa}=H_{\kappa}^{+}=H_{\kappa}A_{\kappa}
\)。
证毕。
\end{proof}

\begin{conclusion}[节点与权重的双重抽取：双正交本征对与二次型闭式]\label{con:xi-terminal-wronskian-hankel-weights-quadratic-extraction}
令 \(e_m\in\RR^{\kappa+1}\) 为标准基。定义右/左本征向量
\[
r_m:=(V^{\mathsf T})^{-1}e_m,\qquad
\ell_m^{\mathsf T}:=e_m^{\mathsf T}V^{\mathsf T}=(1,b_m,\dots,b_m^{\kappa})\qquad(m=0,\dots,\kappa).
\]
则满足：
\begin{enumerate}
  \item \(A_{\kappa}r_m=b_m r_m\) 且 \(\ell_m^{\mathsf T}A_{\kappa}=b_m \ell_m^{\mathsf T}\)，并且 \(\ell_m^{\mathsf T}r_n=\delta_{mn}\)；
  \item \(\{r_m\}\) 关于双线性型 \(\langle u,v\rangle_H:=u^{\mathsf T}H_{\kappa}v\) 两两正交，且
  \[
  r_m^{\mathsf T}H_{\kappa}r_n=c_m\,\delta_{mn};
  \]
  \item 权重可由 \(H_{\kappa}^{-1}\) 在节点向量上的二次型抽取：
  \[
  \boxed{\;
  c_m^{-1}=\ell_m^{\mathsf T}H_{\kappa}^{-1}\ell_m\; }.
  \]
\end{enumerate}
\end{conclusion}
\begin{proof}[证明要点]
第一条由结论 \ref{con:xi-terminal-wronskian-hankel-shift-diagonalization} 的相似对角化直接推出。
第二条由 \(H_{\kappa}=VCV^{\mathsf T}\) 与 \(r_m=(V^{\mathsf T})^{-1}e_m\) 得
\[
r_m^{\mathsf T}H_{\kappa}r_n
=e_m^{\mathsf T}V^{-1}(VCV^{\mathsf T})(V^{\mathsf T})^{-1}e_n
=e_m^{\mathsf T}Ce_n
=c_m\delta_{mn}.
\]
第三条由 \(H_{\kappa}^{-1}=(V^{\mathsf T})^{-1}C^{-1}V^{-1}\) 与 \(\ell_m=V^{\mathsf T}e_m\) 计算即得。证毕。
\end{proof}

\begin{conclusion}[由 \texorpdfstring{$(H_{\kappa},H_{\kappa}^{+})$}{(H,H+)} 的纯行列式重建 \texorpdfstring{$(y,y_{+})$}{(y,y+)}]\label{con:xi-terminal-wronskian-hankel-determinantal-reconstruction}
令 \(u^{\mathsf T}:=e_0^{\mathsf T}H_{\kappa}=(\mu_0,\dots,\mu_\kappa)\) 且 \(v:=e_0\)。
则有 resolvent 恒等式
\[
\boxed{\;
-\frac{y(z)}{y_{+}(z)}=-u^{\mathsf T}(zI-A_{\kappa})^{-1}v
=-e_0^{\mathsf T}H_{\kappa}\,(zH_{\kappa}-H_{\kappa}^{+})^{-1}H_{\kappa}e_0\; }.
\]
并且 \((y,y_{+})\) 可由确定性行列式直接给出：
\[
\boxed{\;
y_{+}(z)=\frac{\det(zH_{\kappa}-H_{\kappa}^{+})}{\det(H_{\kappa})},
\qquad
y(z)=-\frac{\det\!\begin{pmatrix}
zH_{\kappa}-H_{\kappa}^{+} & H_{\kappa}e_0\\
e_0^{\mathsf T}H_{\kappa} & 0
\end{pmatrix}}{\det(H_{\kappa})}\; }.
\]
\end{conclusion}
\begin{proof}[证明要点]
由 \(A_{\kappa}=H_{\kappa}^{-1}H_{\kappa}^{+}\) 得
\(
H_{\kappa}(zI-A_{\kappa})=zH_{\kappa}-H_{\kappa}^{+}
\)，
从而
\[
\det(zI-A_{\kappa})=\frac{\det(zH_{\kappa}-H_{\kappa}^{+})}{\det(H_{\kappa})}.
\]
再由结论 \ref{con:xi-terminal-wronskian-hankel-shift-diagonalization} 的谱同一性可知 \(\det(zI-A_{\kappa})=y_{+}(z)\)，得到 \(y_{+}\) 的行列式表达。
\par
对 \(-u^{\mathsf T}(zI-A_{\kappa})^{-1}v\) 用 Schur 补恒等式
\[
\det\!\begin{pmatrix} M & v\\ u^{\mathsf T} & 0\end{pmatrix}
=-\det(M)\,u^{\mathsf T}M^{-1}v\qquad(M:=zI-A_{\kappa})
\]
并将块矩阵左乘 \(\mathrm{diag}(H_{\kappa},1)\)（仅引入因子 \(\det(H_{\kappa})\)）即可化为所示的 \((zH_{\kappa}-H_{\kappa}^{+})\) 线性笔行列式比值。
该有理函数与 \eqref{eq:xi-terminal-wronskian-stieltjes-transform-minus-y-over-yplus} 具有同一极点集合 \(\{b_m\}\)、同一留数 \(c_m=y(b_m)/y_{+}'(b_m)\)，且在 \(z\to\infty\) 时同为 \(-\mu_0/z+O(z^{-2})\)，故恒等相等，从而得到 \(y\) 的行列式表达。证毕。
\end{proof}

\begin{conclusion}[Wronskian 完备化下的 \texorpdfstring{$\operatorname{Res}$}{Res}--Hankel--\texorpdfstring{$\mathrm{Disc}$}{Disc} 无符号分解]\label{con:xi-terminal-wronskian-hankel-discriminant-signfree-factorization}
在 Wronskian 完备化 \(A=yy_{+}'-y'y_{+}\) 下，有两条无符号精确分解：
\[
\boxed{\;
\operatorname{Res}(A,y_{+})=\det(H_{\kappa})\,\mathrm{Disc}(y_{+}),
\qquad
\operatorname{Res}(A,y)=\det(H_{\kappa})\,\mathrm{Disc}(y)\; }.
\]
\end{conclusion}
\begin{proof}[证明要点]
结论 \ref{con:xi-terminal-wronskian-resultant-discriminant-factorization} 给出
\[
\operatorname{Res}(A,y_{+})
=(-1)^{\frac{\kappa(\kappa+1)}{2}}\operatorname{Res}(y,y_{+})\,\mathrm{Disc}(y_{+}),
\qquad
\operatorname{Res}(A,y)
=(-1)^{\frac{\kappa(\kappa+1)}{2}}\operatorname{Res}(y,y_{+})\,\mathrm{Disc}(y).
\]
再由结论 \ref{con:xi-terminal-wronskian-hankel-principal-determinant-resultant} 的
\(
\operatorname{Res}(y,y_{+})=(-1)^{\frac{\kappa(\kappa+1)}{2}}\det(H_{\kappa})
\)
代入，符号因子抵消即得所示无符号分解。证毕。
\end{proof}

\endinput


\begin{conclusion}[相位量子化：\texorpdfstring{$\int_{\RR}\frac{A}{y^2+y_{+}^2}=\pi(\kappa+1)$}{int A/(y^2+y+^2)=pi(k+1)} 与 \texorpdfstring{$\pi$}{pi}-步长分段]\label{con:xi-terminal-wronskian-phase-quantization}
沿用上述设定，并令 \(w=y_{+}/y\)。
则对实轴去极点集合 \(\RR\setminus\{x_j\}\) 有
\[
\frac{\mathrm{d}}{\mathrm{d}x}\arctan w(x)
=\frac{w'(x)}{1+w(x)^2}
=\frac{A(x)}{y(x)^2+y_{+}(x)^2}.
\]
从而积分恒等式成立：
\[
\boxed{\;\int_{\RR}\frac{A(x)}{y(x)^2+y_{+}(x)^2}\,\mathrm{d}x=\pi(\kappa+1).\;}
\]
更强地，若按实极点把实轴分解为 \(\kappa+1\) 个区间
\[
(-\infty,x_1),\ (x_1,x_2),\ \dots,\ (x_\kappa,\infty),
\]
则对每一段均有精确量子化
\[
\boxed{\;\int_{x_j}^{x_{j+1}}\frac{A(x)}{y(x)^2+y_{+}(x)^2}\,\mathrm{d}x=\pi\;}
\qquad(0\le j\le \kappa),
\]
其中约定 \(x_0=-\infty\)、\(x_{\kappa+1}=+\infty\)。
\end{conclusion}
\begin{proof}
由结论 \ref{con:xi-terminal-wronskian-hermite-biehler-pick-map}，\(w'(x)>0\)（\(x\notin\{x_j\}\)），故 \(w\) 在每个区间上严格递增。
并且在每个极点 \(x_j\) 的两侧，\(w(x)\) 以相反符号发散；在无穷远处 \(w(x)\sim x\)。
因此对任一相邻端点区间 \((x_j,x_{j+1})\)，映射 \(w:(x_j,x_{j+1})\to\RR\) 为双射。
作代换 \(t=w(x)\) 得
\[
\int_{x_j}^{x_{j+1}}\frac{A(x)}{y(x)^2+y_{+}(x)^2}\,\mathrm{d}x
=\int_{x_j}^{x_{j+1}}\frac{w'(x)}{1+w(x)^2}\,\mathrm{d}x
=\int_{-\infty}^{+\infty}\frac{\mathrm{d}t}{1+t^2}
=\pi.
\]
对 \(j=0,\dots,\kappa\) 求和即得全局积分 \(\pi(\kappa+1)\)。证毕。
\end{proof}

\begin{conclusion}[Van Vleck 比值的纯留数表示：\texorpdfstring{$B/A$}{B/A} 的完全显式部分分式与多项式重构]\label{con:xi-terminal-van-vleck-ratio-residue-partial-fraction-reconstruction}
在结论 \ref{con:xi-terminal-riccati-pole-cancellation-bethe-equilibrium} 的等价条件成立时，令
\[
r_k:=\operatorname*{Res}_{z=p_k}\frac{B(z)}{A(z)}=\frac{B(p_k)}{A'(p_k)}\in\CC.
\]
则有严格恒等式
\[
\frac{B(z)}{A(z)}=\sum_{k=1}^{\kappa}\left(\frac{r_k}{z-p_k}+\frac{\overline{r_k}}{z-\overline{p_k}}\right),
\qquad
r_k=-\frac{y'(p_k)}{y(p_k)}=-\sum_{j=1}^{\kappa}\frac{1}{p_k-x_j}.
\]
从而 \(B\) 可由锚点数据显式重构为
\[
B(z)=
-\sum_{k=1}^{\kappa}A(z)\left(
\frac{y'(p_k)}{y(p_k)}\cdot\frac{1}{z-p_k}
\;+\;
\frac{\overline{y'(p_k)}}{\overline{y(p_k)}}\cdot\frac{1}{z-\overline{p_k}}
\right).
\]
\end{conclusion}
\begin{proof}[证明]
由结论 \ref{con:xi-terminal-riccati-pole-cancellation-bethe-equilibrium}，\(\deg B\le 2\kappa-2\) 且 \(\deg A=2\kappa\) 蕴含 \(B/A=O(z^{-2})\)。
因此其部分分式展开无常数项且仅含一阶极点，唯一地写为各极点留数之和；实系数保证 \(\overline{r_k}\) 为 \(\overline{p_k}\) 处留数。
又由方程 \(Ay''-A'y'=By\) 在 \(z=p_k\) 处取值（\(A(p_k)=0\)、\(A'(p_k)\neq 0\)）得
\[
-A'(p_k)y'(p_k)=B(p_k)y(p_k)\quad\Longrightarrow\quad r_k=\frac{B(p_k)}{A'(p_k)}=-\frac{y'(p_k)}{y(p_k)}.
\]
最后 \(y'/y=\sum_j(z-x_j)^{-1}\) 代入 \(z=p_k\) 即得 \(r_k=-\sum_j(p_k-x_j)^{-1}\)。
将部分分式两边乘以 \(A\) 得重构公式。证毕。
\end{proof}

\begin{conclusion}[势的一阶系数的结果式--Vandermonde 表示：\texorpdfstring{$\tau$}{tau}-梯度恒等式]\label{con:xi-terminal-potential-tau-gradient-resultant-vandermonde}
沿用结论 \ref{con:xi-terminal-heine-stieltjes-schrodinger-reduction-local-exponent-rigidity} 与
结论 \ref{con:xi-terminal-heine-stieltjes-pole-expansion-renormalization-sum-rules} 的记号。
设 Schr\"odinger 势在 \(\{p_k,\overline{p_k}\}\) 的部分分式展开为
\[
U(z)=-\frac{3}{4}\sum_{k=1}^{\kappa}\left(\frac{1}{(z-p_k)^2}+\frac{1}{(z-\overline{p_k})^2}\right)
+\sum_{k=1}^{\kappa}\left(\frac{c_k}{z-p_k}+\frac{\overline{c_k}}{z-\overline{p_k}}\right),
\]
并记扩展极点集
\(
\mathcal P:=\{p_1,\dots,p_\kappa,\overline{p_1},\dots,\overline{p_\kappa}\}
\)。
定义 Vandermonde
\[
\Delta(\mathcal P):=\prod_{\substack{u,v\in\mathcal P\\ u<v}}(u-v),
\]
以及结果式
\[
\operatorname{Res}(y,A):=\prod_{k=1}^{\kappa}y(p_k)y(\overline{p_k})
=\prod_{k=1}^{\kappa}|y(p_k)|^2.
\]
则对每个 \(k\)（把 \(\overline{p_1},\dots,\overline{p_\kappa}\) 视为常量，仅对 \(p_k\) 作复导数）有严格恒等式
\[
\boxed{\;
c_k=\frac{\partial}{\partial p_k}\left(\log \operatorname{Res}(y,A)-\log|\Delta(\mathcal P)|\right). \;}
\]
\end{conclusion}
\begin{proof}[证明]
由结论 \ref{con:xi-terminal-heine-stieltjes-pole-expansion-renormalization-sum-rules}，
\(
c_k=-r_k-\tfrac12 s_k
\)
且由方程 \(Ay''-A'y'=By\) 在 \(z=p_k\) 处取值（\(A(p_k)=0\)、\(A'(p_k)\neq 0\)）得
\[
-A'(p_k)y'(p_k)=B(p_k)y(p_k),
\]
故 \(r_k=B(p_k)/A'(p_k)=-y'(p_k)/y(p_k)\)，从而
\[
c_k=\frac{y'(p_k)}{y(p_k)}-\frac12 s_k.
\]
另一方面，
\[
\frac{\partial}{\partial p_k}\log\operatorname{Res}(y,A)
=\frac{\partial}{\partial p_k}\log|y(p_k)|^2
=\frac{y'(p_k)}{y(p_k)}.
\]
并且对 \(\Delta(\mathcal P)\)，由对数导数得
\[
\frac{\partial}{\partial p_k}\log \Delta(\mathcal P)
=\sum_{\substack{u\in\mathcal P\\ u\neq p_k}}\frac{1}{p_k-u}
=s_k.
\]
于是
\(
\frac{\partial}{\partial p_k}\log|\Delta(\mathcal P)|
=\tfrac12 s_k
\)，代回即得结论。证毕。
\end{proof}


% (User input window)

\paragraph{纤维多重度的相对熵坐标：算术--几何缺口、端点压力与相变指数}
以下条目以固定分辨率 \(m\) 的纤维多重度账本为对象，将“边界体积 \(\log|X_m|\)”与“纤维对数体积 \(\log d_m\)”之间的差异精确压缩为 KL/Jeffreys 散度，
并由此给出三种平均的序链、两种全息分解、端点压力斜率以及线性增长率的临界指数化表达。

\begin{conclusion}[算术体积—几何体积的 KL 精确差：边界维度的熵型缺口恒等式]\label{con:xi-kl-arith-geo-gap-identity}
在分辨率 \(m\) 上令 \(\Fold_m:\Omega_m\to X_m\) 的纤维多重度为 \(d_m(x):=\card{\Fold_m^{-1}(x)}\)，并满足守恒
\(
\sum_{x\in X_m} d_m(x)=\card{\Omega_m}=2^m
\)。
在 \(X_m\) 上定义两分布
\[
u_m(x):=\frac{1}{|X_m|},
\qquad
w_m(x):=\frac{d_m(x)}{2^m}.
\]
则相对熵满足严格恒等式
\[
\boxed{\;
D_{\mathrm{KL}}(u_m\mid w_m)
=\log\frac{2^m}{|X_m|}
-\frac{1}{|X_m|}\sum_{x\in X_m}\log d_m(x)\;}
\]
从而
\(
\log(2^m/|X_m|)
\)
与
\(
|X_m|^{-1}\sum_x \log d_m(x)
\)
之差精确等于 \(u_m\) 到 \(w_m\) 的相对熵。
并且
\[
D_{\mathrm{KL}}(u_m\mid w_m)=0
\iff
u_m\equiv w_m
\iff
d_m(x)\equiv \frac{2^m}{|X_m|}\ \ (\forall x\in X_m).
\]
\end{conclusion}

\begin{conclusion}[维度—体积—熵的三层夹逼：\texorpdfstring{$\log d$}{log d} 的三种平均形成序链]\label{con:xi-logd-three-average-chain-kl}
沿用结论 \ref{con:xi-kl-arith-geo-gap-identity} 的记号，定义三种尺度指数
\[
\gamma_m^{\mathrm{geo}}
:=\frac{1}{m|X_m|}\sum_{x\in X_m}\log d_m(x),
\qquad
\gamma_m^{\mathrm{arith}}
:=\frac{1}{m}\log\frac{2^m}{|X_m|},
\qquad
\gamma_m^{\mathrm{Sh}}
:=\frac{1}{m2^m}\sum_{x\in X_m}d_m(x)\log d_m(x).
\]
则成立序链
\[
\gamma_m^{\mathrm{geo}}\ \le\ \gamma_m^{\mathrm{arith}}\ \le\ \gamma_m^{\mathrm{Sh}},
\]
并且两端差值具有精确的 KL 表达
\[
\gamma_m^{\mathrm{arith}}-\gamma_m^{\mathrm{geo}}=\frac{1}{m}D_{\mathrm{KL}}(u_m\mid w_m),
\qquad
\gamma_m^{\mathrm{Sh}}-\gamma_m^{\mathrm{arith}}=\frac{1}{m}D_{\mathrm{KL}}(w_m\mid u_m).
\]
因此两条不等式同时取等当且仅当纤维完全等多重（即 \(u_m\equiv w_m\)）。
\end{conclusion}

\begin{conclusion}[Jeffreys 散度即分形不对称的数值证书：两种可见平均之差的精确公式]\label{con:xi-jeffreys-logd-gap-certificate}
在结论 \ref{con:xi-logd-three-average-chain-kl} 的记号下，定义 Jeffreys 散度
\[
J(u_m,w_m):=D_{\mathrm{KL}}(u_m\mid w_m)+D_{\mathrm{KL}}(w_m\mid u_m).
\]
则有严格恒等式
\[
\gamma_m^{\mathrm{Sh}}-\gamma_m^{\mathrm{geo}}=\frac{1}{m}J(u_m,w_m).
\]
故“按边界均匀抽样”与“按微观均匀抽样（推前权重 \(d_m\)）”两种口径下看到的平均 \(\log d_m\) 之差，恰为 Jeffreys 散度的尺度化；
该量为零当且仅当 \(u_m\equiv w_m\)。
\end{conclusion}

\begin{conclusion}[微观熵密度的两种全息分解：体积项、边界项与相对熵修正]\label{con:xi-micro-entropy-holographic-two-splits}
仍令 \(\sum_{x\in X_m} d_m(x)=2^m\)，则二进制微观熵密度 \(m\log 2\) 具有两条互补的严格分解：
\[
m\log 2
=\log|X_m|
+\sum_{x\in X_m}u_m(x)\log d_m(x)\ +\ D_{\mathrm{KL}}(u_m\mid w_m),
\]
以及
\[
m\log 2
=\log|X_m|
+\sum_{x\in X_m}w_m(x)\log d_m(x)\ -\ D_{\mathrm{KL}}(w_m\mid u_m).
\]
因此同一“总熵”可被写成：边界体积项 \(\log|X_m|\) 与纤维对数体积的期望之和，再加上或减去一个相对熵修正项；
相对熵项精确度量了在非匹配测度下读取纤维对数体积时引入的系统偏置。
\end{conclusion}

\begin{conclusion}[等多重刚性判据：尺度化相对熵消失当且仅当几何均值指数贴住算术指数]\label{con:xi-equimultiplicity-kl-vanishing}
设极限
\[
h_\partial:=\lim_{m\to\infty}\frac{1}{m}\log|X_m|
\]
存在，且 \(\lim_{m\to\infty}\gamma_m^{\mathrm{geo}}\) 存在。
则以下条件等价：
\begin{enumerate}
  \item \(\displaystyle \lim_{m\to\infty}\frac{1}{m}D_{\mathrm{KL}}(u_m\mid w_m)=0\)；
  \item \(\displaystyle \lim_{m\to\infty}\gamma_m^{\mathrm{geo}}=\log 2-h_\partial\)；
  \item \(\displaystyle \lim_{m\to\infty}\bigl(\gamma_m^{\mathrm{arith}}-\gamma_m^{\mathrm{geo}}\bigr)=0\)。
\end{enumerate}
故尺度化相对熵是否消失可作为“边界厚度均匀性”的渐近判据：若其消失，则几何均值指数达到由边界体积项强制的上界；若不消失，则存在不可压缩的熵型缺口。
\end{conclusion}

\begin{conclusion}[对数体积的两种极限端点：\texorpdfstring{$q\to 0$}{q->0} 与 \texorpdfstring{$q\to 1$}{q->1} 的端点导数]\label{con:xi-pressure-endpoints-geo-shannon}
令
\(
S_q(m):=\sum_{x\in X_m} d_m(x)^q
\)，并设对每个固定 \(q>0\) 极限
\[
\tau(q):=\lim_{m\to\infty}\frac{1}{m}\log S_q(m)
\]
存在。
则在存在一侧导数的意义下，端点导数满足
\[
\tau'(0^+)=\lim_{m\to\infty}\gamma_m^{\mathrm{geo}},
\qquad
\tau'(1^-)=\lim_{m\to\infty}\gamma_m^{\mathrm{Sh}}.
\]
因此同一压力函数 \(\tau\) 的两端斜率分别编码“边界均匀口径”的几何均值体积指数与“微观均匀口径”的 Shannon 加权体积指数。
\end{conclusion}

\begin{conclusion}[中心规范熵的普适剥离：\texorpdfstring{$(\ZZ_2)^b$}{(Z2)^b} 的对角约束导致体积亏损]\label{con:xi-gauge-center-diagonal-strip}
设存在 \(b=b(m)\) 条纤维满足 \(d_m(x)=2\)，并令这些 \(S_2\) 因子在纤维对称群直积中生成中心子群 \((\ZZ_2)^b\)。
若进一步施加全局几何可实现约束，使得仅允许对角翻转子群 \(\ZZ_2^{\mathrm{diag}}\subset(\ZZ_2)^b\) 保留，则允许的中心体积发生严格剥离
\[
|(\ZZ_2)^b|/|\ZZ_2^{\mathrm{diag}}|=2^{b-1},
\qquad
\Delta S_{\mathrm{gauge}}=(b-1)\log 2.
\]
当 \(b-1\) 为偶数时，上述体积因子可改写为 \(4^{(b-1)/2}\)，从而在对数账本中呈现强制的 \(4\)-进量子化层级。
\end{conclusion}

\begin{proposition}[纤维对称群的熵下界：可辨识性退化的量化]\label{prop:xi-fold-fiber-symmetry-group-entropy-lower}
\leanverified{paper\_xi\_fold\_fiber\_symmetry\_group\_entropy\_lower}
对任意满射 \(u:A\twoheadrightarrow B\)，定义纤维对称群
\[
G(u):=\prod_{b\in B}\mathrm{Sym}\bigl(u^{-1}(b)\bigr),
\qquad
|G(u)|=\prod_{b\in B}|u^{-1}(b)|!.
\]
取 \(u=\Fold_m:\Omega_m\twoheadrightarrow X_m\)，并记 \(d_m(x):=|\Fold_m^{-1}(x)|\)、\(N:=|\Omega_m|=2^m\)、\(M:=|X_m|\)、\(\bar d:=N/M\)。
则有双侧界
\[
\Gamma(\bar d+1)^{M}\ \le\ |G(\Fold_m)|\ \le\ N!,
\]
从而
\[
\log|G(\Fold_m)|
\ge
M\log\Gamma(\bar d+1)
\ge
N\log\bar d-N+O(M\log\bar d),
\]
并在 Fibonacci 稳定类型尺度 \(M=F_{m+2}=\Theta(\varphi^m)\) 下给出
\(
\log|G(\Fold_m)|=\Omega(m2^m)
\)。
\end{proposition}
\begin{proof}
令 \(d_1,\dots,d_M\) 为纤维大小的列举，则 \(\sum_i d_i=N\)。
由多项式系数
\(
\binom{N}{d_1,\dots,d_M}=\frac{N!}{\prod_i d_i!}\in\NN
\)
得 \(\prod_i d_i!\le N!\)，即上界。
\par
函数 \(x\mapsto \log\Gamma(x)\) 在 \((0,\infty)\) 上凸，Jensen 不等式给出
\[
\frac1M\sum_{i=1}^M\log\Gamma(d_i+1)\ge \log\Gamma\!\left(\frac1M\sum_{i=1}^M(d_i+1)\right)=\log\Gamma(\bar d+1),
\]
指数化得到下界 \(\prod_i d_i!\ge \Gamma(\bar d+1)^M\)。
最后一式由 Stirling 展开 \(\log\Gamma(z+1)=z\log z-z+O(\log z)\) 代入 \(z=\bar d\) 得到。
\end{proof}


\begin{conclusion}[边界相变的熵学判据：两种 KL 的线性增长率给出临界指数]\label{con:xi-boundary-phase-transition-kl-rates}
若存在常数 \(c_{u\to w},c_{w\to u}\ge 0\) 使
\[
\lim_{m\to\infty}\frac{1}{m}D_{\mathrm{KL}}(u_m\mid w_m)=c_{u\to w},
\qquad
\lim_{m\to\infty}\frac{1}{m}D_{\mathrm{KL}}(w_m\mid u_m)=c_{w\to u},
\]
并且 \(h_\partial=\lim_{m\to\infty}m^{-1}\log|X_m|\) 存在，
则由结论 \ref{con:xi-logd-three-average-chain-kl}--\ref{con:xi-jeffreys-logd-gap-certificate} 得到
\[
\lim_{m\to\infty}\gamma_m^{\mathrm{geo}}=\log 2-h_\partial-c_{u\to w},
\qquad
\lim_{m\to\infty}\gamma_m^{\mathrm{Sh}}=\log 2-h_\partial+c_{w\to u},
\]
以及
\[
\lim_{m\to\infty}\bigl(\gamma_m^{\mathrm{Sh}}-\gamma_m^{\mathrm{geo}}\bigr)=c_{u\to w}+c_{w\to u}.
\]
因此可将 \(c_{u\to w}+c_{w\to u}\) 视为“边界相变强度”的一元证书：它同时度量两端平均的分离幅度，并等于 Jeffreys 散度的线性增长率。
\end{conclusion}

\begin{conclusion}[分形维度的熵化下界：维度密度被 \texorpdfstring{$c_{u\to w}$}{c_u->w} 严格压低]\label{con:xi-fractal-dimension-entropy-deficit}
设 \(\varphi\)-标度下的边界维度密度
\[
d_\partial:=\lim_{m\to\infty}\frac{1}{m\log\varphi}\log|X_m|
\]
存在。
在结论 \ref{con:xi-boundary-phase-transition-kl-rates} 的假设下，纤维几何均值的 \(\varphi\)-维度密度
\[
d_{\mathrm{fib}}^{\mathrm{geo}}
:=\frac{1}{\log\varphi}\lim_{m\to\infty}\gamma_m^{\mathrm{geo}}
\]
满足严格表达
\[
d_{\mathrm{fib}}^{\mathrm{geo}}=\frac{\log 2}{\log\varphi}-d_\partial-\frac{c_{u\to w}}{\log\varphi}.
\]
因此一旦 \(c_{u\to w}>0\)，则 \(d_\partial+d_{\mathrm{fib}}^{\mathrm{geo}}\) 将严格低于微观上限 \(\log 2/\log\varphi\)，其严格亏损由相对熵增长率封口。
\end{conclusion}

\begin{conclusion}[熵—体积的可逆重建：\texorpdfstring{$(D_{\mathrm{KL}},\gamma)$}{(DKL,gamma)} 的双坐标互定]\label{con:xi-entropy-volume-dual-coordinates}
在固定 \(m\) 上给定 \(\log|X_m|\) 与总量 \(2^m\)。
则两组数据互相可逆确定：
\begin{itemize}
  \item 熵坐标：\(\bigl(D_{\mathrm{KL}}(u_m\mid w_m),\,D_{\mathrm{KL}}(w_m\mid u_m)\bigr)\)；
  \item 体积坐标：\(\bigl(\gamma_m^{\mathrm{geo}},\,\gamma_m^{\mathrm{Sh}}\bigr)\)。
\end{itemize}
具体地，记
\[
\gamma_m^{\mathrm{arith}}
=\frac{1}{m}\log\frac{2^m}{|X_m|}
=\log 2-\frac{1}{m}\log|X_m|
\]
为已知，则由结论 \ref{con:xi-logd-three-average-chain-kl} 得到双向线性重建
\[
D_{\mathrm{KL}}(u_m\mid w_m)=m\bigl(\gamma_m^{\mathrm{arith}}-\gamma_m^{\mathrm{geo}}\bigr),
\qquad
D_{\mathrm{KL}}(w_m\mid u_m)=m\bigl(\gamma_m^{\mathrm{Sh}}-\gamma_m^{\mathrm{arith}}\bigr),
\]
等价地
\[
\gamma_m^{\mathrm{geo}}=\gamma_m^{\mathrm{arith}}-\frac{1}{m}D_{\mathrm{KL}}(u_m\mid w_m),
\qquad
\gamma_m^{\mathrm{Sh}}=\gamma_m^{\mathrm{arith}}+\frac{1}{m}D_{\mathrm{KL}}(w_m\mid u_m).
\]
因此在该纤维账本框架内，“熵”并非附加量，而是对数体积的等价坐标系；边界厚度与端点可见平均的全部一阶信息可在该坐标变换下等价转写。
\end{conclusion}

\paragraph{纤维账本的算法信息转写：可计算提升、统计谱转移与整数化几何畸变}
沿用本段记号。固定分辨率 \(m\)，记
\[
d_m(x):=\card{\Fold_m^{-1}(x)},
\qquad
\pi_m(x):=\frac{d_m(x)}{2^m},
\qquad
\kappa_m(\pi_m):=\EE_{\pi_m}\!\left[\log d_m(X)\right],
\qquad
\gamma_m(X):=\frac{1}{m}\log d_m(X).
\]
设 \(S\) 为确定性折叠策略，\(r_{m,S}:\Omega_m\to\mathcal P\) 为其 prime-register 残差映射（\(\mathcal P\) 为有限支撑素数幂坐标寄存器），并定义提升映射
\[
\widetilde F_{m,S}(\omega):=\bigl(\Fold_m(\omega),r_{m,S}(\omega)\bigr).
\]

\begin{theorem}[Gödel 外置下的算法信息守恒]\label{thm:xi-godel-externalized-algorithmic-information-conservation}
设 \(\widetilde F_{m,S}\) 可计算，且存在可计算恢复算子 \(R_{m,S}\) 使
\[
R_{m,S}\circ \widetilde F_{m,S}=\mathrm{id}_{\Omega_m}.
\]
令 \(K(\cdot)\) 为相对于固定通用前缀机的前缀 Kolmogorov 复杂度。则存在仅依赖于 \(S\) 的常数 \(c_S\)，使对所有 \(\omega\in\Omega_m\) 成立
\[
\left|K(\omega)-K\!\bigl(\Fold_m(\omega),r_{m,S}(\omega)\bigr)\right|\le c_S,
\]
且
\[
\left|K\!\bigl(\omega\mid \Fold_m(\omega)\bigr)-K\!\bigl(r_{m,S}(\omega)\mid \Fold_m(\omega)\bigr)\right|\le c_S.
\]
\end{theorem}
\begin{proof}
由假设，\(\omega\mapsto (\Fold_m(\omega),r_{m,S}(\omega))\) 与其逆向恢复均为可计算变换。
前缀复杂度在可计算双向可逆编码下保持 \(O(1)\) 不变，故第一式成立。
把 \(\Fold_m(\omega)\) 作为条件后，同一不变性对条件复杂度同样成立，得第二式。证毕。
\end{proof}

\begin{theorem}[条件复杂度与纤维对数体积的双侧夹逼及指数尾界]\label{thm:xi-register-conditional-complexity-vs-fiber-volume}
在定理 \ref{thm:xi-godel-externalized-algorithmic-information-conservation} 的设定下，记 \(X:=\Fold_m(\omega)\)。存在常数 \(c\) 使
\[
K\!\bigl(r_{m,S}(\omega)\mid X\bigr)\le \left\lceil\log_2 d_m(X)\right\rceil+c
\]
对所有 \(\omega\in\Omega_m\) 成立。
并且对任意 \(t\ge 0\)，在 \(\omega\sim\mathrm{Unif}(\Omega_m)\) 下有
\[
\PP\!\left[
K\!\bigl(r_{m,S}(\omega)\mid X\bigr)\le \log_2 d_m(X)-t
\right]\le 2^{-t+c}.
\]
若定义
\[
\Delta_m(\omega):=K\!\bigl(r_{m,S}(\omega)\mid X\bigr)-\log_2 d_m(X),
\]
则
\[
\Delta_m(\omega)\le c+1\quad(\forall \omega),
\qquad
\PP\!\left[\Delta_m(\omega)\le -t\right]\le 2^{-t+c}\quad(t\ge0).
\]
\end{theorem}
\begin{proof}
固定 \(x\in X_m\)。由可计算恢复可知 \(\widetilde F_{m,S}\) 在纤维 \(\Fold_m^{-1}(x)\) 上诱导到残差像集合
\[
R_x:=\{r_{m,S}(\omega):\Fold_m(\omega)=x\}
\]
的双射，故 \(\card{R_x}=d_m(x)\)。
给定 \(x\) 时，可计算枚举 \(\Fold_m^{-1}(x)\) 并输出对应残差；任意 \(r\in R_x\) 可由其枚举序号编码，长度不超过 \(\lceil\log_2 d_m(x)\rceil\) 加常数开销，得上界。

对下界尾概率，利用前缀复杂度计数界：在固定条件 \(x\) 下，满足 \(K(u\mid x)\le k\) 的串最多 \(2^{k+O(1)}\) 个。
取 \(k=\lfloor \log_2 d_m(x)-t\rfloor\)，则
\[
\frac{\card{\{r\in R_x:K(r\mid x)\le \log_2 d_m(x)-t\}}}{\card{R_x}}
\le 2^{-t+O(1)}.
\]
又因 \(\omega\) 在 \(\Omega_m\) 上均匀，条件于 \(X=x\) 时在纤维上亦均匀，故上式即为条件坏事件概率上界。再对 \(x\) 取全概率求和即得无条件指数尾界。证毕。
\end{proof}

\begin{theorem}[Gödel 外置尾界的指数最优性]\label{thm:xi-register-conditional-complexity-tail-sharpness}
在定理 \ref{thm:xi-register-conditional-complexity-vs-fiber-volume} 的设定下，存在常数 \(c'\)（仅依赖于 \(S\) 与所选通用机）使对任意 \(t\ge 0\) 与任意 \(m\) 都有
\[
\PP\!\left[
K\!\bigl(r_{m,S}(\omega)\mid X\bigr)\le \log_2 d_m(X)-t+c'
\right]
\ge
2^{-t}-\frac{|X_m|}{2^m},
\qquad \omega\sim\mathrm{Unif}(\Omega_m).
\]
特别地，若 \(|X_m|/2^m\to 0\)，则对每个固定 \(t\) 成立
\[
\liminf_{m\to\infty}
\PP\!\left[
K\!\bigl(r_{m,S}(\omega)\mid X\bigr)\le \log_2 d_m(X)-t+c'
\right]
\ge 2^{-t}.
\]
\end{theorem}
\begin{proof}
固定 \(x\in X_m\) 并沿用上一定理证明中的残差像集合 \(R_x\)。
由可计算恢复，\(r_{m,S}\) 在纤维 \(\Fold_m^{-1}(x)\) 上诱导到 \(R_x\) 的双射，且 \(|R_x|=d_m(x)\)。
条件于 \(X=x\) 时，\(\omega\) 在纤维上均匀，故 \(r:=r_{m,S}(\omega)\) 在有限集合 \(R_x\) 上亦均匀。

按定理 \ref{thm:xi-register-conditional-complexity-vs-fiber-volume} 的枚举过程，给定 \(x\) 可计算枚举 \(\Fold_m^{-1}(x)\) 并同步输出其残差值，从而得到 \(R_x\) 的一个可计算枚举序列 \((r_1,\dots,r_{d_m(x)})\)。
令
\(
A_x:=\{r_j:\ 1\le j\le \lfloor d_m(x)/2^t\rfloor\}\subseteq R_x
\)。
对任意 \(r_j\in A_x\)，用其枚举序号 \(j\)（二进制长度 \(\le \log_2 d_m(x)-t+O(1)\)）作为描述，并在条件 \(x\) 下运行上述枚举直至第 \(j\) 项即可恢复 \(r_j\)；
因此存在常数 \(c'\) 使
\[
K(r\mid x)\le \log_2 d_m(x)-t+c'\qquad(\forall r\in A_x).
\]
于是
\[
\PP\!\left[K(r\mid X)\le \log_2 d_m(X)-t+c'\ \big|\ X=x\right]
\ge \frac{|A_x|}{|R_x|}
\ge 2^{-t}-\frac1{d_m(x)}.
\]
对 \(x\) 以权重 \(\PP[X=x]=d_m(x)/2^m\) 取全概率求和，得到
\[
\PP(\cdots)\ge \sum_{x\in X_m}\frac{d_m(x)}{2^m}\Bigl(2^{-t}-\frac1{d_m(x)}\Bigr)
=2^{-t}-\frac{|X_m|}{2^m}.
\]
若 \(|X_m|/2^m\to 0\)，则对每个固定 \(t\) 取下极限即得第二式。证毕。
\end{proof}

\begin{theorem}[多重分形谱向算法复杂度谱的指数等价转移]\label{thm:xi-multifractal-to-algorithmic-spectrum-transfer}
在定理 \ref{thm:xi-register-conditional-complexity-vs-fiber-volume} 的设定下，定义
\[
\eta_m(\omega):=\frac1m K\!\bigl(r_{m,S}(\omega)\mid X\bigr)\quad(\text{bit/步}).
\]
则对任意 \(\varepsilon>0\)，存在 \(c_\varepsilon>0\) 使
\[
\PP\!\left[\left|\eta_m-\frac{\gamma_m(X)}{\log 2}\right|>\varepsilon\right]
\le 2^{-c_\varepsilon m}.
\]
因此：
\begin{enumerate}
  \item 若 \(\gamma_m(X)\) 在 \(\pi_m\) 下满足中心极限定理，则 \(\eta_m\) 满足同一极限定理（仅有 \(\log 2\) 的单位换算），方差按 \((\log2)^{-2}\) 缩放；
  \item 若 \(\gamma_m(X)\) 以速率 \(m\) 满足良速率函数 \(I\) 的大偏差原理，则 \(\eta_m\) 亦以速率 \(m\) 满足大偏差原理，其速率函数
  \[
  J(u)=I(u\log2).
  \]
\end{enumerate}
\end{theorem}
\begin{proof}
由定理 \ref{thm:xi-register-conditional-complexity-vs-fiber-volume}，对 \(t=\varepsilon m\) 有
\[
\PP\!\left[K(r_{m,S}(\omega)\mid X)\le \log_2 d_m(X)-\varepsilon m\right]
\le 2^{-\varepsilon m+O(1)}.
\]
同时点态上界给出
\[
K(r_{m,S}(\omega)\mid X)\le \log_2 d_m(X)+O(1),
\]
故当 \(m\) 充分大时
\[
\PP\!\left[K(r_{m,S}(\omega)\mid X)\ge \log_2 d_m(X)+\varepsilon m\right]=0.
\]
两式合并得到指数等价。指数等价序列共享大偏差原理；并且两者差异为 \(o_{\PP}(m^{1/2})\) 时，中心极限定理保持不变，遂得两条转移结论。证毕。
\end{proof}

\begin{theorem}[纤维熵势与期望条件复杂度的 \(O(1)\) 同一性]\label{thm:xi-fiber-entropy-potential-expected-conditional-complexity}
在 \(\omega\sim\mathrm{Unif}(\Omega_m)\) 下，存在与 \(m\) 无关常数 \(C\) 使
\[
\left|
\EE_{\omega}\!\left[K\!\bigl(r_{m,S}(\omega)\mid X\bigr)\right]
-\frac{\kappa_m(\pi_m)}{\log2}
\right|
\le C.
\]
\end{theorem}
\begin{proof}
定义
\[
\Delta(\omega):=\log_2 d_m(X)-K\!\bigl(r_{m,S}(\omega)\mid X\bigr).
\]
由定理 \ref{thm:xi-register-conditional-complexity-vs-fiber-volume}，\(\Delta(\omega)\ge -O(1)\) 且
\(\PP[\Delta(\omega)\ge t]\le 2^{-t+O(1)}\)。
因此
\[
\EE[\Delta]\le \sum_{t\ge0}\PP[\Delta\ge t]\le \sum_{t\ge0}2^{-t+O(1)}=O(1),
\]
从而
\[
\EE\!\left[K(r_{m,S}(\omega)\mid X)\right]
=\EE[\log_2 d_m(X)]+O(1)
=\frac{\kappa_m(\pi_m)}{\log2}+O(1).
\]
吸收常数即得结论。证毕。
\end{proof}

\begin{theorem}[prime-register 向量的 Gödel 整数化具有不可约的超线性位长]\label{thm:xi-godel-integerization-log-overhead}
设 \(\tau=(a_1,\dots,a_T)\) 为长度 \(T\) 的事件码序列，且 \(a_t\in\NN_{\ge1}\)。
定义 Gödel 整数
\[
G(\tau):=\prod_{t=1}^{T}p_t^{a_t},
\]
其中 \(p_t\) 为第 \(t\) 个素数。
则存在绝对常数 \(c_0,c_1>0\)，使对一切 \(T\ge 1\) 成立
\[
\log_2 G(\tau)\ \ge\ T\log_2 T-c_0T-c_1.
\]
因此把长度 \(T\) 的寄存器向量压缩为单整数时，其最小可能位长必为 \(\Omega(T\log T)\)。
\end{theorem}
\begin{proof}
由 \(a_t\ge 1\) 得
\[
G(\tau)\ \ge\ \prod_{t=1}^{T}p_t.
\]
又因第 \(t\) 个素数满足 \(p_t\ge t+1\)，故
\[
\prod_{t=1}^{T}p_t\ \ge\ \prod_{t=1}^{T}(t+1)\ =\ \frac{(T+1)!}{2}.
\]
取二进对数并用 Stirling 型下界（存在绝对常数 \(C_0,C_1>0\) 使对一切 \(n\ge 1\) 有
\(\log_2(n!)\ge n\log_2 n-C_0n-C_1\)），得到
\[
\log_2 G(\tau)\ \ge\ \log_2\bigl((T+1)!/2\bigr)
\ \ge\ (T+1)\log_2(T+1)-C_0(T+1)-C_1-1.
\]
将右端改写为 \(T\log_2T-c_0T-c_1\) 的形式并吸收常数即可。证毕。
\end{proof}

\begin{corollary}[Gödel 整数化的位长与条件复杂度之间的对数鸿沟]\label{cor:xi-godel-integerization-vs-conditional-complexity-gap}
沿用定理 \ref{thm:xi-register-conditional-complexity-vs-fiber-volume} 的设定，令 \(X:=\Fold_m(\omega)\) 且 \(r:=r_{m,S}(\omega)\)。
若 \(r\) 的非零坐标数为 \(T=\#\mathrm{supp}(r)\) 并且 \(r\) 的时间序列译码为一条长度 \(T\) 的事件码 \(\tau(\omega)\)（其中每个码字 \(\ge 1\)），则有
\[
\frac{\log_2 G(\tau(\omega))}{K(r\mid X)}
\ \ge\
\frac{T\log_2 T-O(T)}{\lceil\log_2 d_m(X)\rceil+O(1)},
\]
其中 \(O(T)\) 的常数仅依赖于定理 \ref{thm:xi-godel-integerization-log-overhead} 中的绝对常数。
特别地，当 \(T\ge \alpha m\)（某个常数 \(\alpha>0\)）时，
\[
\frac{\log_2 G(\tau(\omega))}{K(r\mid X)}\ \ge\ c\,\log_2 m-O(1)
\]
对某个常数 \(c>0\) 成立。
\end{corollary}
\begin{proof}
由定理 \ref{thm:xi-godel-integerization-log-overhead} 得
\(\log_2 G(\tau(\omega))\ge T\log_2 T-O(T)\)。
另一方面由定理 \ref{thm:xi-register-conditional-complexity-vs-fiber-volume}，
\[
K(r\mid X)\le \lceil\log_2 d_m(X)\rceil+O(1)\le m+O(1).
\]
合并即得第一式；若 \(T\ge \alpha m\)，则 \(T\log_2 T\ge \alpha m\log_2 m-O(m)\)，除以 \(m+O(1)\) 得第二式。证毕。
\end{proof}

\begin{theorem}[单位圆到椭圆的 Gödel 畸变律：条件数的超线性增长]\label{thm:xi-godel-ellipse-condition-number-distortion}
沿用定理 \ref{thm:xi-godel-integerization-log-overhead} 的记号，定义面积保持线性算子
\[
A_\tau:=\mathrm{diag}\!\left(G(\tau)^{1/2},\,G(\tau)^{-1/2}\right)\in SL(2,\RR),
\qquad
E_\tau:=A_\tau(S^1).
\]
则 \(E_\tau\) 的轴比（等于 \(A_\tau\) 的谱条件数）满足
\[
\kappa(A_\tau)=G(\tau).
\]
于是存在 \(c_1>0\) 与 \(T_0\)，使对 \(T\ge T_0\) 有
\[
\log \kappa(A_\tau)\ge c_1\,T\log T.
\]
\end{theorem}
\begin{proof}
对角算子 \(A_\tau\) 的奇异值为 \(G(\tau)^{1/2}\) 与 \(G(\tau)^{-1/2}\)，故
\[
\kappa(A_\tau)=\frac{G(\tau)^{1/2}}{G(\tau)^{-1/2}}=G(\tau).
\]
再代入定理 \ref{thm:xi-godel-integerization-log-overhead} 的增长下界即得。证毕。
\end{proof}

\paragraph{Joukowsky--Gödel 椭圆输运的代数刚性、判别式骨架与伽罗瓦相干}
在定理 \ref{thm:xi-godel-ellipse-condition-number-distortion} 的椭圆化读出口径上，进一步固定
\[
E_r:=\{\,w=r z+r^{-1}z^{-1}:|z|=1\,\},\qquad
Q_r(w):=\operatorname{Res}_z\!\bigl(P(z),\,r z^2-w z+r^{-1}\bigr),
\]
其中 \(r>1\)，且 \(P(z)=a_n z^n+a_{n-1}z^{n-1}+\cdots+a_0\in K[z]\)（\(a_n a_0\neq 0\)）。
下列结论表明：在该椭圆输运下，首项、判别式平方类与伽罗瓦信息在可计算意义上保持严格刚性。

\begin{theorem}[输运结式的首项系数刚性]\label{thm:xi-jg-transport-resultant-leading-rigidity}
\leanverified{paper\_xi\_jg\_transport\_resultant\_leading\_rigidity}
设 \(K\) 为一域，\(r\) 为独立符号。定义
\[
Q_r(w)=\operatorname{Res}_z\!\bigl(P(z),\,r z^2-w z+r^{-1}\bigr)\in K(r)[w].
\]
则
\[
\deg_w Q_r=n,\qquad
\operatorname{LC}_w(Q_r)=a_n a_0.
\]
特别地，\(\operatorname{LC}_w(Q_r)\) 与 \(r\) 无关。
\end{theorem}
\begin{proof}
令 \(z_1,\dots,z_n\) 为 \(P\) 在代数闭包中的根（计重数）。由结式根表示，
\[
Q_r(w)=a_n^{2}\prod_{j=1}^{n}\bigl(r z_j^2-w z_j+r^{-1}\bigr).
\]
视为关于 \(w\) 的多项式，每个因子给出首项 \((-z_j)w\)，故
\[
\operatorname{LC}_w(Q_r)=a_n^{2}(-1)^n\prod_{j=1}^n z_j
=a_n^{2}(-1)^n\cdot (-1)^n\frac{a_0}{a_n}
=a_n a_0.
\]
又 \(a_0\neq 0\) 蕴含 \(\prod_j z_j\neq 0\)，从而 \(\deg_w Q_r=n\)。证毕。
\end{proof}

\begin{corollary}[首一情形下的边界常数读出]\label{cor:xi-jg-monic-boundary-readout}
\leanverified{paper\_xi\_jg\_monic\_boundary\_readout}
在定理 \ref{thm:xi-jg-transport-resultant-leading-rigidity} 条件下，若 \(P\) 首一（\(a_n=1\)），则
\[
\operatorname{LC}_w(Q_r)=a_0.
\]
因此常数项 \(a_0\) 可由任意尺度参数 \(r\) 的椭圆输运多项式 \(Q_r\) 直接读出。
\end{corollary}
\begin{proof}
由定理 \ref{thm:xi-jg-transport-resultant-leading-rigidity} 立即得到结论。证毕。
\end{proof}

\begin{theorem}[判别式平方类守恒]\label{thm:xi-jg-discriminant-squareclass-invariance}
\leanverified{paper\_xi\_jg\_discriminant\_squareclass\_invariance}
设 \(\operatorname{char}(K)\neq 2\)，且 \(P\) 可分离并满足 \(a_0 a_n\neq 0\)。记 \(\Disc(P)\) 与 \(\Disc(Q_r)\) 为关于各自主变量的判别式，则在平方类群
\[
K(r)^\times/(K(r)^\times)^2
\]
中有
\[
\Disc(Q_r)\equiv \Disc(P).
\]
等价地，
\[
K(r)\bigl(\sqrt{\Disc(Q_r)}\bigr)=K(r)\bigl(\sqrt{\Disc(P)}\bigr).
\]
\end{theorem}
\begin{proof}
令 \(z_1,\dots,z_n\) 为 \(P\) 的根，并设
\[
w_j:=r z_j+r^{-1}z_j^{-1}.
\]
由定理 \ref{thm:xi-jg-transport-resultant-leading-rigidity} 的首项系数表达，
\[
Q_r(w)=(a_n a_0)\prod_{j=1}^n (w-w_j).
\]
故
\[
\Disc(Q_r)=(a_n a_0)^{2n-2}\prod_{i<j}(w_i-w_j)^2.
\]
另一方面，
\[
w_i-w_j
=r(z_i-z_j)+r^{-1}(z_i^{-1}-z_j^{-1})
=(z_i-z_j)\!\left(r-\frac{r^{-1}}{z_i z_j}\right),
\]
从而
\[
\prod_{i<j}(w_i-w_j)^2
=\prod_{i<j}(z_i-z_j)^2
\cdot
\prod_{i<j}\left(r-\frac{r^{-1}}{z_i z_j}\right)^2.
\]
又
\[
\Disc(P)=a_n^{2n-2}\prod_{i<j}(z_i-z_j)^2.
\]
代入得
\[
\Disc(Q_r)
=a_0^{2n-2}\Disc(P)\cdot
\prod_{i<j}\left(r-\frac{r^{-1}}{z_i z_j}\right)^2.
\]
记
\[
B_r:=\prod_{i<j}\left(r-\frac{r^{-1}}{z_i z_j}\right).
\]
由于 \(B_r\) 为根的对称有理式且 \(a_0\neq 0\)，由对称多项式基本定理可知 \(B_r\in K(r)\)。故
\[
\frac{\Disc(Q_r)}{\Disc(P)}=a_0^{2n-2}B_r^2\in (K(r)^\times)^2,
\]
即两者同平方类，结论成立。证毕。
\end{proof}

\leanverified{paper\_xi\_jg\_leyang\_bad\_prime\_rigidity}
\begin{corollary}[交替群判据与 Lee--Yang 二次骨架的尺度稳定]\label{cor:xi-jg-leyang-bad-prime-rigidity}
在定理 \ref{thm:xi-jg-discriminant-squareclass-invariance} 条件下，
\[
\Disc(P)\in (K^\times)^2
\iff
\Disc(Q_r)\in (K(r)^\times)^2.
\]
因此由判别式平方性控制的“伽罗瓦群落入 \(A_n\)”判据在椭圆输运下保持相干。
并且，当 Lee--Yang 结构的坏素数（例如 \(37\)）由判别式诱导的二次扩张给出时，其二次分辨率层面的分歧素数集合不随尺度参数 \(r\) 漂移。
\end{corollary}
\begin{proof}
第一部分由定理 \ref{thm:xi-jg-discriminant-squareclass-invariance} 直接推出。
第二部分注意到坏素数的二次骨架由 \(\sqrt{\Disc(\cdot)}\) 决定，而两者对应二次扩张相同，故结论成立。证毕。
\end{proof}

\begin{theorem}[椭圆输运 Laurent 两端系数与原多项式反演]\label{thm:xi-jg-holographic-coefficient-extraction}
\leanverified{paper\_xi\_jg\_holographic\_coefficient\_extraction}
在定理 \ref{thm:xi-jg-transport-resultant-leading-rigidity} 的设定下，令 \(w=r t\)，并将
\[
Q_r(r t)\in K[r,r^{-1}][t]
\]
视为关于 \(r\) 的 Laurent 多项式。记 \([r^k](\cdot)\) 为 \(r^k\) 系数，则
\[
[r^n]\,Q_r(r t)=a_0 P(t),\qquad
[r^{-n}]\,Q_r(r t)=a_n^2.
\]
因此，给定 \(Q_r\) 及其在 \(w=r t\) 代换后的 Laurent 展开，可由
\(\operatorname{LC}_w(Q_r)\)、\([r^n]Q_r(r t)\)、\([r^{-n}]Q_r(r t)\) 进行代数重构；当固定 \(a_n\) 的平方根分支（特别是首一情形 \(a_n=1\)）时，\(P\) 的系数唯一确定。
\end{theorem}
\begin{proof}
仍记 \(z_1,\dots,z_n\) 为 \(P\) 的根。由结式根表示，
\[
Q_r(r t)
=a_n^2\prod_{j=1}^n\bigl(r z_j^2-r t z_j+r^{-1}\bigr)
=a_n^2 r^n\prod_{j=1}^n\bigl(z_j^2-t z_j+r^{-2}\bigr).
\]
把右端视为关于 \(r^{-2}\) 的多项式。其常数项为
\[
\prod_{j=1}^n(z_j^2-t z_j)
=\left(\prod_{j=1}^n z_j\right)\!\left(\prod_{j=1}^n (z_j-t)\right)
=\frac{a_0}{a_n^2}P(t).
\]
故
\[
[r^n]\,Q_r(r t)=a_n^2\cdot \frac{a_0}{a_n^2}P(t)=a_0 P(t).
\]
另一方面，在
\(
a_n^2\prod_j(r z_j^2-r t z_j+r^{-1})
\)
中选取每个因子的 \(r^{-1}\) 项，唯一得到 \(r^{-n}\) 端项，故
\[
[r^{-n}]\,Q_r(r t)=a_n^2.
\]
证毕。
\end{proof}

\leanverified{paper\_xi\_jg\_transcendental\_scale\_galois\_full\_fidelity}
\begin{theorem}[超越尺度下的伽罗瓦全信息保持]\label{thm:xi-jg-transcendental-scale-galois-full-fidelity}
设 \(\operatorname{char}(K)=0\)，\(P\in K[z]\) 可分离且 \(a_0 a_n\neq 0\)。取 \(r\) 超越于 \(K\)，令 \(F:=K(r)\)。
记 \(L\) 为 \(P\) 在 \(K\) 上的分裂域，\(M\) 为 \(Q_r\) 在 \(F\) 上的分裂域。则
\[
M=L(r),\qquad
\operatorname{Gal}(Q_r/F)\cong \operatorname{Gal}(P/K).
\]
\end{theorem}
\begin{proof}
由 \(w_i=r z_i+r^{-1}z_i^{-1}\in L(r)\) 可知 \(Q_r\) 在 \(L(r)\) 上分裂，故 \(M\subseteq L(r)\)。

反过来，取任一根 \(z_i\) 并令 \(w_i:=r z_i+r^{-1}z_i^{-1}\in M\)。考虑
\[
g_i(z):=r z^2-w_i z+r^{-1}\in M[z].
\]
显然 \(z_i\) 同时是 \(P\) 与 \(g_i\) 的根。\(g_i\) 的另一根为 \(r^{-2}z_i^{-1}\)。
若此根亦为 \(P\) 的根，则存在 \(j\) 使 \(z_j=r^{-2}z_i^{-1}\)，即
\[
r^2=(z_i z_j)^{-1}\in \overline K,
\]
与 \(r\) 对 \(K\) 的超越性矛盾。故 \(\gcd(P,g_i)\) 在 \(M[z]\) 中为一次多项式。

欧几里得算法给出 \(\gcd(P,g_i)\) 的系数属于 \(M\)；其唯一根即 \(z_i\)，故 \(z_i\in M\)。
对所有 \(i\) 同理，得 \(z_1,\dots,z_n\in M\)，于是 \(L\subseteq M\)，进而 \(L(r)\subseteq M\)。
结合 \(M\subseteq L(r)\) 得 \(M=L(r)\)。最后 \(r\) 为超越元，故基变不改变伽罗瓦群，得到
\[
\operatorname{Gal}(Q_r/F)\cong \operatorname{Gal}(P/K).
\]
证毕。
\end{proof}

\begin{proposition}[代数尺度下的二元纤维上界]\label{prop:xi-jg-algebraic-scale-binary-fiber-bound}
\leanverified{paper\_xi\_jg\_algebraic\_scale\_binary\_fiber\_bound}
设 \(\operatorname{char}(K)\neq 2\)，且 \(r\in K^\times\) 为代数尺度。沿用定理
\ref{thm:xi-jg-transcendental-scale-galois-full-fidelity}
的记号，则
\[
M\subseteq L,\qquad
L\subseteq M\bigl(\sqrt{w_1^2-4},\dots,\sqrt{w_n^2-4}\bigr),
\]
其中 \(w_1,\dots,w_n\) 为 \(Q_r\) 的根。特别地，
\[
[L:M]\mid 2^n,
\]
且 \(\operatorname{Gal}(L/M)\) 为初等 \(2\)-群的子群。
\end{proposition}
\begin{proof}
因 \(r\in K\)，每个 \(w_i=r z_i+r^{-1}z_i^{-1}\) 均属于 \(L\)，故 \(M\subseteq L\)。
对任意 \(i\)，方程
\[
r z^2-w_i z+r^{-1}=0
\]
给出
\[
z=\frac{w_i\pm \sqrt{w_i^2-4}}{2r}.
\]
因此对应的 \(z_i\) 属于 \(M(\sqrt{w_i^2-4})\)。合并 \(i=1,\dots,n\) 得
\[
L\subseteq M\bigl(\sqrt{w_1^2-4},\dots,\sqrt{w_n^2-4}\bigr).
\]
右侧是至多 \(n\) 个二次扩张的复合，故扩张次数整除 \(2^n\)。
伽罗瓦群结论随之成立。证毕。
\end{proof}

\begin{theorem}[群作用下的二值跳跃选择与账本不可约下界]\label{thm:xi-equivariant-binary-jump-ledger-lower-bound}
\leanverified{paper\_xi\_equivariant\_binary\_jump\_ledger\_lower\_bound}
令有限群 \(G\) 作用于有限集合 \(\Omega\)。
设观测输入在 \(G\) 作用下不可区分（等价地：输入仅给出该 \(G\)-集合的同构类型），并允许协议额外读取一段外置二进制账本 \(B\in\{0,1\}^b\)，随后输出 \(\omega\in\Omega\)。
称协议满足等变自然性，若存在映射 \(\Phi:\{0,1\}^b\to\Omega\) 使得对一切 \(g\in G\) 有
\[
\Phi(B)\ \mapsto\ g\cdot \Phi(B)
\]
与输入在 \(g\) 作用下的同一化相容；换言之，协议在 \(G\)-对称下不引入额外破缺数据时，其所有可能输出必须由账本 \(B\) 完全枚举。
则对任意轨道 \(\mathcal O(\omega)=G\cdot\omega\) 有不可逾越下界
\[
\boxed{\;b\ \ge\ \left\lceil \log_2|\mathcal O(\omega)|\right\rceil.\;}
\]
特别地，若作用传递，则 \(b\ge \lceil\log_2|\Omega|\rceil\)；
并且若 \(\Omega\) 无全局不动点，则在 \(b=0\) 的 choice-free 口径下不存在满足自然性的协议。
\end{theorem}
\begin{proof}
在无外部破缺数据的前提下，协议对给定 \(B\) 的输出为某个确定的 \(\Phi(B)\in\Omega\)。
因此在所有账本取值上，协议可产生的输出至多为 \(2^b\) 个。
若要求协议能覆盖某一轨道中的所有选择，则必有 \(2^b\ge |\mathcal O(\omega)|\)，从而得到所示下界。
\par
若 \(b=0\)，则输出为某个固定点 \(\omega_\ast\in\Omega\)。自然性强制其在任意 \(g\in G\) 下保持不变，即 \(g\cdot\omega_\ast=\omega_\ast\)，故 \(\omega_\ast\) 必为全局不动点；若无不动点则不可能。证毕。
\end{proof}

\begin{theorem}[二值跳跃半直积的 Abel 化与可观测位坍缩]\label{thm:xi-binary-jump-semidirect-abelianization-coinvariants}
\leanverified{paper\_xi\_binary\_jump\_semidirect\_abelianization\_coinvariants}
令 \(I\) 为有限集合，并令 \(V:=(\ZZ_2)^I\)（视为 \(I\) 上的 \(\ZZ_2\)-值函数群）。
设有限群 \(G\) 作用于 \(I\)，并以坐标置换诱导作用于 \(V\)。令半直积
\[
W:=V\rtimes G.
\]
记 coinvariants
\[
V_G:=V\big/\langle v-g\cdot v:\ v\in V,\ g\in G\rangle.
\]
则存在规范同构
\[
\boxed{\;W^{\mathrm{ab}}\ \cong\ V_G\ \oplus\ G^{\mathrm{ab}}.\;}
\]
进一步，若 \(I\) 的轨道分解为 \(I=\bigsqcup_{j=1}^{k}\mathcal O_j\)，则
\[
\boxed{\;V_G\ \cong\ (\ZZ_2)^{k},\;}
\]
其基可取为各轨道指示函数在 coinvariants 中的像。
因此在不携带 \(I\) 的标号信息时，\(V\) 的可交换可观测二值自由度至多保留“每个轨道的总翻转”一比特，其余 \(I\)-逐点分支位在 Abel 层面必然坍缩。
\end{theorem}
\begin{proof}
由 \(V\) 交换，交换子仅来自 \([(v,1),(0,g)]\)。
直接计算得
\[
[(v,1),(0,g)]=(v-g\cdot v,1),
\]
故 \([W,W]\cap V=\langle v-g\cdot v\rangle\)，并且 \(W/[W,W]\cong (V/[V,G])\oplus (G/[G,G])=V_G\oplus G^{\mathrm{ab}}\)。
\par
对 \(V_G\) 的维数：在单一轨道上，任意两单位基向量在 coinvariants 中同类，故该轨道贡献一维；不同轨道互不识别，故总维数为轨道数 \(k\)。
取轨道指示函数即可给出一组基。证毕。
\end{proof}

\begin{conjecture}[Gödel 尺度族上的 Hilbert 相干性]\label{conj:xi-jg-hilbert-coherence-godel-scales}
设 \(P\in\ZZ[z]\) 可分离，且 \(\operatorname{Gal}(P/\QQ)\) 为给定传递子群（典型情形为 \(S_n\)）。
定义 Gödel 尺度子集
\[
\mathcal G
:=\left\{
\prod_{p\ \mathrm{prime}}p^{a_p}:\ a_p\in\ZZ,\ \text{仅有限支撑}
\right\}
\subset \QQ_{>0}.
\]
预期在任意自然高度截断计数下，对密度 \(1\) 的 \(r\in\mathcal G\) 有
\[
\operatorname{Gal}(Q_r/\QQ)\cong \operatorname{Gal}(P/\QQ),
\]
并且例外尺度组成薄集，可由有限个代数曲线的有理点参数化。
\end{conjecture}

\begin{conjecture}[输运分支的满秩账本与单比特可观测性的并发]\label{conj:xi-jg-branch-ledger-fullrank-vs-coinvariant-collapse}
在命题 \ref{prop:xi-jg-algebraic-scale-binary-fiber-bound} 的设定下，设 \(M\) 为 \(Q_r\) 的分裂域，令
\[
L:=M\bigl(\sqrt{w_1^2-4},\dots,\sqrt{w_n^2-4}\bigr),
\qquad
H:=\operatorname{Gal}(L/M)\ \le\ (\ZZ_2)^n.
\]
猜想对“泛”尺度参数 \(r\)（例如在适当稀疏例外集之外）同时满足：
\begin{enumerate}
\normalfont
  \item \textbf{满秩账本：}\(H\cong(\ZZ_2)^n\)，等价地 \([L:M]=2^n\)，从而任何要在无标号口径下指定一条全局分支选择的协议其外置账本复杂度可达 \(n\) 比特量级；
  \item \textbf{单比特可观测性：}设 \(G:=\operatorname{Gal}(M/K(r))\) 在根指标集 \(I=\{1,\dots,n\}\) 上的作用为传递（或仅有常数个轨道），则由定理 \ref{thm:xi-binary-jump-semidirect-abelianization-coinvariants}，
  半直积 \(H\rtimes G\) 的 Abel 层面对分支位的可观测维数至多为轨道数（典型为 \(1\)），其余分支位在不携带标号信息时不可由任何可交换不变量读出。
\end{enumerate}
该并发现象给出一条“账本复杂度最大化 \(\leftrightarrow\) 可观测位最小化”的二值跳跃对偶性：全息重建所需的外置分支账本可达到满秩，而内禀可交换可观测信息却坍缩为常数阶。
\end{conjecture}

\paragraph{Gödel--素数寄存器的时间移位初对象、平方自由外置与椭圆同态链}
为与定义 \ref{def:pom-fold-prime-lift}、命题 \ref{prop:pom-fold-prime-lift-injective} 及定理 \ref{thm:pom-fiber-indset-factorization} 的既有接口保持一致，本段把 prime-register 外置、Fold 纤维自由度与面积守恒椭圆化组织为同一条代数链。

\begin{definition}[时间移位相容的可加编码器范畴]\label{def:xi-godel-shift-additive-encoder-category}
记 \(\mathsf{Seq}(\NN_{\ge 1})\) 为有限正整数序列集合。称五元组
\[
(M,+,0;\sigma,u)
\]
为时间移位相容的可加编码器，若满足：
\begin{enumerate}
  \item \((M,+,0)\) 为交换幺半群；
  \item \(\sigma:M\to M\) 为单射幺半群自同态；
  \item 映射
  \[
  G_M(k_1,\dots,k_T):=\sum_{t=1}^T k_t\,\sigma^{t-1}(u)
  \]
  从 \(\mathsf{Seq}(\NN_{\ge 1})\) 到 \(M\) 为单射。
\end{enumerate}
记 \(\mathbf{Enc}_\sigma\) 为如下范畴：对象为上述编码器；态射
\[
\Phi:(M,+,0;\sigma,u)\to(M',+,0;\sigma',u')
\]
满足 \(\Phi\) 为幺半群同态，且
\[
\Phi(u)=u',\qquad \Phi\circ\sigma=\sigma'\circ\Phi.
\]
\end{definition}

\begin{theorem}[Gödel--素数寄存器的时间移位初对象]\label{thm:xi-godel-prime-register-initial-object}
\leanverified{paper\_xi\_godel\_prime\_register\_initial\_object}
令
\[
\mathsf P:=\NN^{(\mathbb P)}
\]
为有限支撑素数指数向量幺半群，\(\{\delta_{p_t}\}_{t\ge 1}\) 为标准基。
定义
\[
\sigma_{\mathsf P}(\delta_{p_t})=\delta_{p_{t+1}},\qquad u_{\mathsf P}:=\delta_{p_1}.
\]
则有：
\begin{enumerate}
  \item \((\mathsf P,+,0;\sigma_{\mathsf P},u_{\mathsf P})\) 是定义 \ref{def:xi-godel-shift-additive-encoder-category} 的对象；
  \item 该对象在 \(\mathbf{Enc}_\sigma\) 中为初对象：对任意 \((M,+,0;\sigma,u)\in\mathbf{Enc}_\sigma\)，存在唯一态射
  \[
  \Phi:\mathsf P\to M
  \]
  使得 \(\Phi(u_{\mathsf P})=u\) 且 \(\Phi\circ\sigma_{\mathsf P}=\sigma\circ\Phi\)。
\end{enumerate}
\end{theorem}
\begin{proof}
首先验证对象性。对任意序列 \((k_1,\dots,k_T)\)，有
\[
G_{\mathsf P}(k_1,\dots,k_T)=\sum_{t=1}^T k_t\,\delta_{p_t}.
\]
由于 \(\{\delta_{p_t}\}\) 在 \(\NN\)-系数下线性无关，\(G_{\mathsf P}\) 单射，故 (1) 成立。

再证初对象性。固定任意对象 \((M,+,0;\sigma,u)\)。由 \(\mathsf P\) 的自由交换幺半群结构，先在基上定义
\[
\Phi(\delta_{p_t})=\sigma^{t-1}(u),
\]
再唯一加法延拓到 \(\mathsf P\)。于是 \(\Phi(u_{\mathsf P})=u\) 且
\[
\Phi(\sigma_{\mathsf P}(\delta_{p_t}))
=\Phi(\delta_{p_{t+1}})
=\sigma^t(u)
=\sigma(\sigma^{t-1}(u))
=\sigma(\Phi(\delta_{p_t})),
\]
故 \(\Phi\circ\sigma_{\mathsf P}=\sigma\circ\Phi\)。

若 \(\Psi:\mathsf P\to M\) 也是同样态射，则由 \(\Psi(u_{\mathsf P})=u\) 与交换关系递推得
\[
\Psi(\delta_{p_t})=\sigma^{t-1}(u)=\Phi(\delta_{p_t}),
\]
再由自由性推出 \(\Psi=\Phi\)。唯一性成立。证毕。
\end{proof}

\begin{theorem}[Fold 纤维的平方自由素数外置与自由度刻画]\label{thm:xi-fold-fiber-squarefree-prime-externalization}
\leanverified{paper\_xi\_fold\_fiber\_squarefree\_prime\_externalization}
固定分辨率 \(m\)。令 \(\Fold_m:\Omega_m\to X_m\) 为折叠映射。由定理 \ref{thm:pom-fiber-indset-factorization} 与定义 \ref{def:pom-fiber-adm-path}，对每个 \(x\in X_m\) 存在候选位点图
\[
\Gamma_x=P_{k_+}[\mathsf{Adm}(x)]
\]
及规范双射
\[
\Theta_x:\Fold_m^{-1}(x)\xrightarrow{\sim}\mathsf{Ind}(\Gamma_x).
\]
再令 \(b(\omega)\in\{0,1\}\) 为引理 \ref{lem:pom-one-bit} 的隐藏位。对每个 \(x\) 任选互异素数族 \(\{q_v:v\in\mathsf{Adm}(x)\}\)，定义平方自由编码
\[
\mathcal G_x(I):=\prod_{v\in I} q_v,\qquad I\in\mathsf{Ind}(\Gamma_x).
\]
并定义
\[
\Psi_m:\Omega_m\to \bigsqcup_{x\in X_m}\bigl(\{x\}\times\{0,1\}\times\NN_{\mathrm{sf}}\bigr),\qquad
\Psi_m(\omega):=\Bigl(\Fold_m(\omega),\,b(\omega),\,\mathcal G_{\Fold_m(\omega)}\!\bigl(\Theta_{\Fold_m(\omega)}(\omega)\bigr)\Bigr).
\]
则 \(\Psi_m\) 为单射；若把值域限制到可实现像
\[
\mathcal R_m:=\Psi_m(\Omega_m),
\]
则 \(\Psi_m:\Omega_m\to\mathcal R_m\) 为显式双射。
\end{theorem}
\begin{proof}
设 \(\Psi_m(\omega)=\Psi_m(\omega')\)。则第一坐标给出
\[
\Fold_m(\omega)=\Fold_m(\omega')=:x.
\]
第三坐标给出
\[
\mathcal G_x(\Theta_x(\omega))=\mathcal G_x(\Theta_x(\omega')).
\]
由于 \(\mathcal G_x\) 由互异素数的平方自由乘积定义，唯一分解定理蕴含
\[
\Theta_x(\omega)=\Theta_x(\omega').
\]
再由 \(\Theta_x\) 为双射，得 \(\omega=\omega'\)。故 \(\Psi_m\) 单射。
限制到 \(\mathcal R_m\) 后满射按定义成立，故得到双射。证毕。
\end{proof}

\begin{corollary}[平方自由素数外置的最坏码长下界与对数乘性胀量]\label{cor:xi-squarefree-externalization-log-overhead}
\leanverified{paper\_xi\_squarefree\_externalization\_log\_overhead}
在定理 \ref{thm:xi-fold-fiber-squarefree-prime-externalization} 的设定下，固定某个 \(x\in X_m\) 并记 \(s:=|\mathsf{Adm}(x)|\)。
把 \(\mathcal G_x(I)=\prod_{v\in I}q_v\) 视为平方自由外置码字，并令
\[
L_{\max}(x):=\max_{I\in\mathsf{Ind}(\Gamma_x)}\log_2 \mathcal G_x(I).
\]
则在最有利的选择意义下（允许在 \(\mathsf{Adm}(x)\) 上自由选取互异素数以最小化最坏码长），存在绝对常数 \(c_1,c_2>0\) 使对充分大的 \(s\) 有
\[
L_{\max}(x)\ \ge\ c_1\, s\log_2 s.
\]
并且当 \(\Gamma_x\) 含有一条长度 \(\Theta(s)\) 的路径分量（例如 \(\Gamma_x\simeq P_s\)）时，
\[
\frac{L_{\max}(x)}{\log_2 d_m(x)}\ \ge\ c_2\,\log s,
\]
从而平方自由素数外置在该类纤维上相对于信息论极限 \(\log_2 d_m(x)\) 具有不可消去的对数乘性开销。
\end{corollary}
\begin{proof}
由定义，\(\mathsf{Ind}(\Gamma_x)\) 至少包含“取一路径分量的奇数位顶点”的独立集，其规模为 \(\Theta(s)\)。
对任意互异素数族，最坏码长至少为该独立集所对应素数的对数和。
在最小化最坏码长的意义下，应取最小的 \(s\) 个互异素数；于是
\[
L_{\max}(x)\ \ge\ \sum_{j\le c s}\log_2 p_j\ =\ \Theta(s\log_2 s),
\]
其中 \(p_j\) 为第 \(j\) 个素数，最后一步由 primorial 渐近 \(\log(p_k\#)\sim k\log k\) 得到（并见推论 \ref{cor:xi-squarefree-prime-support-min-growth} 的同型估计口径）。
\par
若 \(\Gamma_x\simeq P_s\)，则 \(d_m(x)=\card{\mathsf{Ind}(P_s)}=F_{s+2}\)，从而 \(\log_2 d_m(x)=\Theta(s)\)；
与上式合并即得比值下界。证毕。
\end{proof}

\leanverified{paper\_xi\_gap\_prime\_squarefree\_dirichlet\_fraction}
\begin{theorem}[非相邻素数平方自由 Dirichlet 级数与递归比值表征]\label{thm:xi-gap-prime-squarefree-dirichlet-fraction}
记 \(p_n\) 为第 \(n\) 个素数。定义
\[
\mathcal I:=\{S\subset\NN\ \text{有限}:\ \forall i\in S,\ i+1\notin S\},
\]
并在 \(\Re(s)>1\) 上定义
\[
Z_{\mathrm{gap}}(s):=\sum_{S\in\mathcal I}\ \prod_{i\in S} p_i^{-s}.
\]
再设
\[
\mathcal I_n:=\{S\subseteq\{1,\dots,n\}:\ \forall i\in S,\ i+1\notin S\},
\qquad
Z_n(s):=\sum_{S\in\mathcal I_n}\prod_{i\in S}p_i^{-s}.
\]
则有：
\begin{enumerate}
  \item \(Z_{\mathrm{gap}}(s)\) 在 \(\Re(s)>1\) 绝对收敛，并定义解析函数；
  \item 对任意实数 \(\sigma>1\)，\(\{Z_n(\sigma)\}_{n\ge 0}\) 单调递增且有界，且
  \[
  Z_{\mathrm{gap}}(\sigma)=\lim_{n\to\infty} Z_n(\sigma);
  \]
  \item 初值 \(Z_0(s)=1\)、\(Z_1(s)=1+p_1^{-s}\)，并满足递推
  \[
  Z_n(s)=Z_{n-1}(s)+p_n^{-s}Z_{n-2}(s)\qquad(n\ge 2);
  \]
  若记 \(R_n(s):=Z_n(s)/Z_{n-1}(s)\)，则
  \[
  R_1(s)=1+p_1^{-s},\qquad
  R_n(s)=1+\frac{p_n^{-s}}{R_{n-1}(s)}\quad(n\ge 2);
  \]
  \item 对实数 \(\sigma>1\) 有夹逼
  \[
  1+\sum_{p}p^{-\sigma}\ \le\ Z_{\mathrm{gap}}(\sigma)\ \le\ \prod_{p}(1+p^{-\sigma})=\frac{\zeta(\sigma)}{\zeta(2\sigma)}.
  \]
  并且对复数 \(s\)（\(\Re(s)>1\)）有
  \[
  |Z_{\mathrm{gap}}(s)|\le Z_{\mathrm{gap}}(\Re s)\le \frac{\zeta(\Re s)}{\zeta(2\Re s)}.
  \]
\end{enumerate}
\end{theorem}
\begin{proof}
记 \(\sigma:=\Re(s)\)。有
\[
\sum_{S\in\mathcal I}\left|\prod_{i\in S}p_i^{-s}\right|
=\sum_{S\in\mathcal I}\prod_{i\in S}p_i^{-\sigma}
\le \sum_{S\subset\NN,\text{有限}}\prod_{i\in S}p_i^{-\sigma}
=\prod_{i\ge 1}(1+p_i^{-\sigma}).
\]
当 \(\sigma>1\) 时，\(\sum_i p_i^{-\sigma}<\infty\)，故右侧收敛，绝对收敛得证。
并且在任意紧集 \(\Re(s)\ge 1+\varepsilon\) 上同样由 \(\prod_i(1+p_i^{-1-\varepsilon})\) 支配，故级数正规收敛，\(Z_{\mathrm{gap}}\) 解析。

对递推式，将 \(\mathcal I_n\) 按 \(n\in S\) 与 \(n\notin S\) 分拆即可得
\[
Z_n=Z_{n-1}+p_n^{-s}Z_{n-2}.
\]
比值公式由代数整理直接得到。

对实数 \(\sigma>1\)，\(Z_n(\sigma)\) 为正项和，故随 \(n\) 单调递增；
又
\[
Z_n(\sigma)\le \prod_{i=1}^n(1+p_i^{-\sigma})\le \prod_p(1+p^{-\sigma})<\infty,
\]
故极限存在，且等于 \(Z_{\mathrm{gap}}(\sigma)\)。

下界来自空集与单点集贡献；上界来自去除“非相邻”约束并取全平方自由子集乘积。证毕。
\end{proof}

\begin{theorem}[面积守恒椭圆族与素数寄存器同态实现]\label{thm:xi-prime-register-ellipse-monoid-homomorphism}
\leanverified{paper\_xi\_prime\_register\_ellipse\_monoid\_homomorphism}
定义面积为 \(\pi\) 的轴对齐椭圆族
\[
\mathcal E:=\{E(a):a>0,\ E(a)=\{(x,y):x^2/a^2+y^2a^2=1\}\},
\]
并定义乘法型组合
\[
E(a)\odot E(b):=E(ab).
\]
记
\[
N:\mathsf P\to\NN_{>0},\qquad N(r):=\prod_{p}p^{r(p)},
\qquad
\Xi:\mathsf P\to\mathcal E,\quad \Xi(r):=E(N(r)).
\]
则：
\begin{enumerate}
  \item \(\Xi\) 为单射幺半群同态，即
  \[
  \Xi(r+r')=\Xi(r)\odot\Xi(r'),\qquad
  \Xi(r)=\Xi(r')\Rightarrow r=r';
  \]
  \item 若定义
  \[
  d_{\mathcal E}(E(a),E(b)):=2|\log(a/b)|,
  \]
  则
  \[
  d_{\mathcal E}\bigl(\Xi(r),\Xi(r')\bigr)
  =2\left|\sum_p\bigl(r(p)-r'(p)\bigr)\log p\right|.
  \]
  因而每个素数生成元 \(\delta_p\) 对应基本步长 \(2\log p\)。
\end{enumerate}
\end{theorem}
\begin{proof}
由定义，
\[
N(r+r')=\prod_p p^{r(p)+r'(p)}=N(r)N(r'),
\]
故
\[
\Xi(r+r')=E(N(r+r'))=E(N(r)N(r'))=\Xi(r)\odot\Xi(r').
\]
若 \(\Xi(r)=\Xi(r')\)，则 \(N(r)=N(r')\)，由算术基本定理得 \(r=r'\)，故单射。

第二式由 \(a=N(r)\)、\(b=N(r')\) 直接代入 \(d_{\mathcal E}\) 的定义即得。证毕。
\end{proof}

\begin{proposition}[信息守恒外置与范数一致性]\label{prop:xi-prime-externalization-norm-one-consistency}
\leanverified{paper\_xi\_prime\_externalization\_norm\_one\_consistency}
对任意 \(r\in\mathsf P\)，定义
\[
a_\infty:=N(r),\qquad a_p:=p^{-r(p)}\ \ (p\ \text{素数}).
\]
则满足归一化乘积公式
\[
|a_\infty|_\infty^{-1}\cdot\prod_p |a_p|_p=1.
\]
反之，若给定一组赋值数据满足：仅有限多个 \(p\) 有 \(a_p\neq 1\)，且每个 \(a_p\) 可写为 \(p^{-e_p}\)（\(e_p\in\NN\)），并满足
\[
|a_\infty|_\infty^{-1}\cdot\prod_p |a_p|_p=1,
\]
则存在唯一 \(r\in\mathsf P\) 使 \(r(p)=e_p\) 且 \(a_\infty=N(r)\)。
\end{proposition}
\begin{proof}
由 \(|p^{-r(p)}|_p=p^{r(p)}\) 得
\[
|a_\infty|_\infty^{-1}\prod_p |a_p|_p
=N(r)^{-1}\prod_p p^{r(p)}
=1.
\]
反向由有限支撑给出唯一指数向量 \(e_p\)，令 \(r(p):=e_p\) 即得唯一 \(r\in\mathsf P\)；
乘积公式强制 \(a_\infty=\prod_p p^{e_p}=N(r)\)。证毕。
\end{proof}

\begin{corollary}[Fold 信息亏损的三重同态表述]\label{cor:xi-fold-loss-prime-ellipse-triple-correspondence}
\leanverified{paper\_xi\_fold\_loss\_prime\_ellipse\_triple\_correspondence}
在定理 \ref{thm:xi-fold-fiber-squarefree-prime-externalization} 与定理 \ref{thm:xi-prime-register-ellipse-monoid-homomorphism} 的设定下，定义
\[
\Xi_{\mathrm{sf}}:\NN_{\mathrm{sf}}\to\mathcal E,\qquad \Xi_{\mathrm{sf}}(n):=E(n).
\]
则存在链式单射
\[
\Omega_m
\xhookrightarrow{\ \Psi_m\ }
\bigsqcup_{x\in X_m}\bigl(\{x\}\times\{0,1\}\times\NN_{\mathrm{sf}}\bigr)
\xhookrightarrow{\ \mathrm{id}\times\mathrm{id}\times\Xi_{\mathrm{sf}}\ }
\bigsqcup_{x\in X_m}\bigl(\{x\}\times\{0,1\}\times\mathcal E\bigr).
\]
因此每个微态的纤维区分信息可等价表述为：独立集选择、平方自由素数外置、面积守恒椭圆参数三种同态像。
\end{corollary}
\begin{proof}
第一箭头由定理 \ref{thm:xi-fold-fiber-squarefree-prime-externalization} 给出单射。
第二箭头由 \(\Xi_{\mathrm{sf}}(n)=E(n)\) 的单射性（参数 \(n\) 唯一）给出。
复合仍为单射，结论成立。证毕。
\end{proof}

\begin{conjecture}[非相邻素数 Dirichlet 级数的自然边界]\label{conj:xi-gap-prime-dirichlet-natural-boundary}
对定理 \ref{thm:xi-gap-prime-squarefree-dirichlet-fraction} 的
\(
Z_{\mathrm{gap}}(s)
\)，预期：
\begin{enumerate}
  \item \(Z_{\mathrm{gap}}(s)\) 可亚纯延拓至半平面 \(\Re(s)>0\)；
  \item 直线 \(\Re(s)=0\) 构成自然边界，不存在跨越该边界的单值亚纯延拓。
\end{enumerate}
其机制预期来自递归比值
\[
R_n(s)=1+\frac{p_n^{-s}}{R_{n-1}(s)}
\]
在 \(\Re(s)\downarrow 0\) 时的相位耦合累积：与 Euler 乘积由乘法独立性主导的情形不同，此处解析障碍由“素数序位上的相邻排斥约束”驱动。
\end{conjecture}

\paragraph{圆维秩零度、最小补账与连续语义不可压缩}
以下结论在定义 \ref{def:circle-dimension}--\ref{def:half-circle-dimension} 与命题 \ref{prop:circle-dimension-laws} 的口径下工作。记有限生成离散交换群同态为 \(f:G\to H\)，并沿用
\[
\cdim(A)=\dim_{\QQ}(A\otimes_{\ZZ}\QQ)
\]
（对有限生成交换群 \(A\)）这一等价描述。

\begin{theorem}[圆维秩零度公式]\label{thm:xi-cdim-rank-nullity}
设 \(f:G\to H\) 为有限生成离散交换群同态。则
\[
\cdim(G)=\cdim(\ker f)+\cdim(\operatorname{im}f).
\]
等价地，
\[
\cdim(\ker f)=\cdim(G)-\cdim(\operatorname{im}f),\qquad
\cdim(\operatorname{im}f)=\cdim(G)-\cdim(\ker f).
\]
\end{theorem}
\begin{proof}
由短正合列
\[
0\longrightarrow \ker f\longrightarrow G\longrightarrow \operatorname{im}f\longrightarrow 0
\]
张量 \(\QQ\)（\(\QQ\) 平坦）得
\[
0\longrightarrow (\ker f)\otimes\QQ\longrightarrow G\otimes\QQ\longrightarrow (\operatorname{im}f)\otimes\QQ\longrightarrow 0.
\]
取 \(\QQ\)-维数即
\[
\dim_{\QQ}(G\otimes\QQ)
=\dim_{\QQ}((\ker f)\otimes\QQ)+\dim_{\QQ}((\operatorname{im}f)\otimes\QQ).
\]
按圆维的秩刻画即得所述等式。证毕。
\end{proof}
\leanverified{paper\_xi\_cdim\_rank\_nullity}

\begin{definition}[单射化补账]\label{def:xi-cdim-injective-ledger-completion}
给定同态 \(f:G\to H\)。称一对 \((R,r)\) 为 \(f\) 的单射化补账，若 \(R\) 为有限生成离散交换群且 \(r:G\to R\) 为群同态，并使
\[
(f,r):G\longrightarrow H\oplus R,\qquad g\longmapsto (f(g),r(g))
\]
为单射。
\end{definition}

\begin{theorem}[最小圆维补账定理]\label{thm:xi-cdim-minimal-ledger-cost}
设 \(f:G\to H\) 为有限生成离散交换群同态，记 \(K:=\ker f\)。则：
\begin{enumerate}
  \item 对任意单射化补账 \((R,r)\)，都有
  \[
  \cdim(R)\ge \cdim(K).
  \]
  \item 存在单射化补账 \((R_{\min},r_{\min})\) 使
  \[
  \cdim(R_{\min})=\cdim(K).
  \]
  且可取
  \[
  R_{\min}\cong \ZZ^{\cdim(K)}\oplus T
  \]
  （\(T\) 有限；例如可取 \(T=\operatorname{Tor}(G)\)）。
\end{enumerate}
\end{theorem}
\begin{proof}
先证下界。若 \((f,r)\) 单射，则其在 \(K\) 上的限制 \(r|_K:K\to R\) 亦单射：若 \(x\in K\) 且 \(r(x)=0\)，则 \((f,r)(x)=(0,0)\)，由单射得 \(x=0\)。
由 \(\QQ\) 平坦性，\(r|_K\otimes\mathrm{id}_{\QQ}\) 仍单射，故
\[
\cdim(K)=\dim_{\QQ}(K\otimes\QQ)\le \dim_{\QQ}(R\otimes\QQ)=\cdim(R).
\]

再证可达性。写
\[
G\cong \ZZ^n\oplus T_G,\qquad K\le G,\qquad k:=\cdim(K).
\]
由子群标准形（Smith 正规形），可在 \(\ZZ^n\) 取基 \(e_1,\dots,e_n\) 与整数 \(d_i\ge 1\) 使
\[
K\cap \ZZ^n=\bigoplus_{i=1}^{k} d_i\ZZ\,e_i.
\]
定义
\[
R_{\min}:=\ZZ^k\oplus T_G,
\]
并令
\[
r_{\min}\!\left(\sum_{i=1}^n a_i e_i,\ t\right):=((a_1,\dots,a_k),t).
\]
若 \(x\in K\) 且 \(r_{\min}(x)=0\)，则其 torsion 分量 \(t=0\)，且自由分量可写作
\(
\sum_{i=1}^k d_i m_i e_i
\)
并满足 \(d_i m_i=0\)，故 \(m_i=0\)。于是 \(x=0\)，即 \(r_{\min}|_K\) 单射。
因此 \((f,r_{\min})\) 单射：若 \((f,r_{\min})(g)=0\)，则 \(g\in K\) 且 \(r_{\min}(g)=0\)，从而 \(g=0\)。
最后
\[
\cdim(R_{\min})=\cdim(\ZZ^k)=k=\cdim(K).
\]
证毕。
\end{proof}
\leanverified{paper\_xi\_cdim\_minimal\_ledger\_cost}

\begin{definition}[圆维缺陷]\label{def:xi-cdim-defect}
对有限生成离散交换群同态 \(f:G\to H\)，定义
\[
\delta(f):=\cdim(\ker f)\in\NN\cup\{0\}.
\]
\end{definition}

\begin{theorem}[圆维缺陷的链式法则]\label{thm:xi-cdim-defect-chain-rule}
设
\[
G\xrightarrow{f}H\xrightarrow{g}K
\]
为有限生成离散交换群同态。则存在短正合列
\[
0\longrightarrow \ker f\longrightarrow \ker(g\circ f)\longrightarrow \ker\!\bigl(g|_{\operatorname{im}f}\bigr)\longrightarrow 0,
\]
从而
\[
\delta(g\circ f)=\delta(f)+\delta\!\bigl(g|_{\operatorname{im}f}\bigr).
\]
特别地，由 \(\ker(g|_{\operatorname{im}f})\le \ker g\) 得
\[
\delta(g\circ f)\le \delta(f)+\delta(g).
\]
\end{theorem}
\begin{proof}
记 \(L:=\ker(g\circ f)\)。定义
\[
\pi:L\to \ker(g|_{\operatorname{im}f}),\qquad \pi(x):=f(x).
\]
若 \(x\in L\)，则 \(g(f(x))=0\)，故 \(\pi(x)\in \ker(g|_{\operatorname{im}f})\)。
对任意 \(y\in \ker(g|_{\operatorname{im}f})\)，存在 \(x\in G\) 使 \(f(x)=y\)，并且 \(g(f(x))=g(y)=0\)，故 \(x\in L\)，从而 \(\pi\) 满射。
且
\[
\ker\pi=\{x\in L:f(x)=0\}=L\cap \ker f=\ker f.
\]
故得所述短正合列。对其应用定理 \ref{thm:xi-cdim-rank-nullity}，即得加法公式。
最后不等式由子群单调性
\(
\cdim(\ker(g|_{\operatorname{im}f}))\le \cdim(\ker g)
\)
立即推出。证毕。
\end{proof}
\leanverified{paper\_xi\_cdim\_defect\_chain\_rule}
\leanverified{paper\_xi\_phase\_fold\_kernel\_minimal\_certificate\_rank}

\begin{theorem}[相位折叠核的最小证书秩]\label{thm:xi-phase-fold-kernel-minimal-certificate-rank}
给定 \(B\in M_{k\times n}(\RR)\)，定义群同态
\[
\chi_B:\ZZ^n\to \TT^k,\qquad
\chi_B(x):=\exp(2\pi i\,Bx).
\]
令
\[
K_B:=\ker\chi_B=\{x\in\ZZ^n:Bx\in\ZZ^k\},\qquad r:=\operatorname{rank}(K_B).
\]
则：
\begin{enumerate}
  \item 若 \(r':\ZZ^n\to \ZZ^m\) 为群同态且 \((\chi_B,r'):\ZZ^n\to \TT^k\times \ZZ^m\) 单射，则 \(m\ge r\)；
  \item 存在 \(r_{\min}:\ZZ^n\to \ZZ^r\) 使 \((\chi_B,r_{\min})\) 单射。
\end{enumerate}
因此，相位折叠引入的最小可验证补账秩恰为 \(\operatorname{rank}(K_B)\)。
\end{theorem}
\begin{proof}
对 (1)，若 \((\chi_B,r')\) 单射，则 \(r'|_{K_B}\) 必单射。于是 \(K_B\hookrightarrow \ZZ^m\)，故 \(m\ge \operatorname{rank}(K_B)=r\)。

对 (2)，由 \(\ZZ^n\) 子群标准形，可取基 \(e_1,\dots,e_n\) 与整数 \(d_i\ge 1\) 使
\[
K_B=\bigoplus_{i=1}^{r} d_i\ZZ\,e_i.
\]
定义
\[
r_{\min}\!\left(\sum_{i=1}^n a_i e_i\right):=(a_1,\dots,a_r)\in \ZZ^r.
\]
若 \(x\in K_B\) 且 \(r_{\min}(x)=0\)，则 \(x=\sum_{i=1}^r d_i m_i e_i\) 且 \(d_i m_i=0\)，故 \(x=0\)。因此 \(r_{\min}|_{K_B}\) 单射。
进而若 \((\chi_B,r_{\min})(x)=(\chi_B,r_{\min})(y)\)，则 \(z:=x-y\in K_B\) 且 \(r_{\min}(z)=0\)，故 \(z=0\)，即 \(x=y\)。证毕。
\end{proof}

\begin{theorem}[连续语义下的圆维不可压缩]\label{thm:xi-circle-dimension-continuous-noncompressibility}
设 \(D\) 为紧致零维 Hausdorff 空间（例如任意 profinite 紧群）。若存在连续单射
\[
\iota:\TT^d\hookrightarrow \TT^k\times D,
\]
则必有 \(d\le k\)。
\end{theorem}
\begin{proof}
\(\TT^d\) 紧、\(\TT^k\times D\) Hausdorff，故连续单射 \(\iota\) 为拓扑嵌入。于是覆盖维数满足单调性
\[
\dim_{\mathrm{top}}(\TT^d)\le \dim_{\mathrm{top}}(\TT^k\times D).
\]
又 \(\dim_{\mathrm{top}}(\TT^d)=d\)，且 \(D\) 零维给
\[
\dim_{\mathrm{top}}(\TT^k\times D)=\dim_{\mathrm{top}}(\TT^k)=k.
\]
故 \(d\le k\)。证毕。
\end{proof}

\paragraph{预算接口化注记}
定理 \ref{thm:xi-cdim-minimal-ledger-cost} 与定理 \ref{thm:xi-cdim-defect-chain-rule} 把“可逆外置所需账本”压缩为核的圆维证书；定理 \ref{thm:xi-circle-dimension-continuous-noncompressibility} 则给出连续语义下离散轴不可替代相位圆因子的禁止性边界。二者可直接并入现有预算--\(\mathsf{NULL}\)--外置轴三分接口，以形成结构层与实现层的统一下界口径。

% (User input window)

\subsubsection{跨接口终端结论：熵坐标的幂 escort 轨道、二维缺陷层析与 primitive 的算术相分裂}\label{subsubsec:xi-escort-defect-primitive-arithmetic-splitting}
本段把若干已在文内不同接口出现的可审计机制压缩为可独立引用的结论条目：
其一给出在混合守恒约束下的幂 escort 轨道刚性，并把熵提升为唯一热力学坐标；
其二把二维点状缺陷测度的边界可见量组织为 Fourier--Laplace 指纹并给出有限采样的低秩重构口径；
其三把压力分支的解析半径、primitive 微正则熵的局部形状与 RH-sharp 的算术次级项，统一写成“判别式/因子结构”所强制的可计算封口。

\begin{conclusion}[熵约束下的幂 escort 轨道刚性：熵作为唯一热力学坐标]\label{con:xi-terminal-escort-entropy-thermo-coordinate}
令 \(X\) 为有限集（\(|X|=n\ge 2\)），取全支撑分布 \(w\in\Delta(X)\)，并固定 \(\beta\in(0,1)\)。
对每个 \(\gamma>1\) 定义幂 escort 变换
\[
\mathrm{Escort}_\gamma(r)(x):=\frac{r(x)^\gamma}{\sum_{y\in X}r(y)^\gamma}\qquad(r\in\Delta(X)).
\]
若存在 \((q,r)\in\Delta(X)^2\) 满足 \(w=\beta q+(1-\beta)r\) 且 \(q=\mathrm{Escort}_\gamma(r)\)，则该解唯一；记之为 \((q_\gamma,r_\gamma)\)。
并且映射 \(\gamma\mapsto H(r_\gamma)\) 连续且严格递增，且满足
\[
\lim_{\gamma\downarrow 1}r_\gamma=w.
\]
令
\[
H_{\max}(\beta,w):=\sup\left\{H(r):\ r\in\Delta(X),\ \frac{w-(1-\beta)r}{\beta}\in\Delta(X)\right\}\in(H(w),\log n],
\]
则 \(\lim_{\gamma\to\infty}H(r_\gamma)=H_{\max}(\beta,w)\)。
特别地，若均匀分布 \(U\) 可行（等价于 \(w(x)\ge (1-\beta)/n\) 对所有 \(x\in X\)），则 \(H_{\max}(\beta,w)=\log n\) 且 \(r_\gamma\to U\)。
因此对任意 \(h\in(H(w),H_{\max}(\beta,w))\) 存在唯一 \(\gamma(h)>1\) 使 \(H(r_{\gamma(h)})=h\)；
在守恒约束 \(w=\beta q+(1-\beta)r\) 下，熵水平集的极值点只能沿该一维幂律轨道运动，从而熵可被视为该变分前沿的唯一热力学坐标。
\end{conclusion}
\begin{proof}[证明要点]
见命题 \ref{prop:pom-microcanonical-two-temperature-escort-orbit-entropy-monotone}。
\end{proof}

\begin{conclusion}[二维 Fourier--Laplace 全息层析：边界指纹唯一恢复 bulk 缺陷测度]\label{con:xi-terminal-defect-measure-fourier-laplace-holographic}
令 \(\nu\) 为 \(\RR\times(0,\infty)\) 上有限正测度，并定义边界密度
\[
H_\nu(x):=\int_{\RR\times(0,\infty)}\frac{4\delta}{(x-\gamma)^2+(1+\delta)^2}\,\dd\nu(\gamma,\delta),
\qquad
H_{\nu,t}:=P_t*H_\nu,
\]
其中 \(P_t(x)=\frac{1}{\pi}\frac{t}{x^2+t^2}\) 为 \(\RR\) 上 Poisson 核。
记 \(\widehat{H}_{\nu,t}(\xi)\) 为 Fourier 变换，令 \(s:=|\xi|\) 并定义二参数指纹
\[
F(\xi,s):=e^{(t+1)s}\widehat{H}_{\nu,t}(\xi).
\]
则成立严格分解
\[
F(\xi,s)=4\pi\int_{\RR\times(0,\infty)}\frac{\delta}{1+\delta}\,e^{-s\delta}\,e^{-\mathrm{i}\gamma\xi}\,\dd\nu(\gamma,\delta),
\]
从而 \(F\) 是 \(\gamma\) 的 Fourier 变换与 \(\delta\) 的 Laplace 变换之张量化组合；特别地映射 \(\nu\mapsto F\) 在有限测度类上为单射，
因此 bulk 缺陷测度 \(\nu\) 可由边界可见指纹 \(F\) 的取值唯一恢复。
\end{conclusion}
\begin{proof}[证明要点]
见命题 \ref{con:discussion-defect-measure-fourier-laplace}。
\end{proof}

\begin{conclusion}[低秩缺陷的有限采样可恢复性：原子测度的 Hankel--Prony 双向分解]\label{con:xi-terminal-atomic-defect-hankel-prony}
在结论 \ref{con:xi-terminal-defect-measure-fourier-laplace-holographic} 的设定下，若
\[
\nu=\sum_{j=1}^{K}a_j\,\delta_{(\gamma_j,\delta_j)}
\qquad(K\ge 1,\ a_j>0,\ (\gamma_j,\delta_j)\ \text{两两不同}),
\]
则指纹为显式有限和
\[
F(\xi,s)=4\pi\sum_{j=1}^{K}a_j\frac{\delta_j}{1+\delta_j}\,e^{-s\delta_j}\,e^{-\mathrm{i}\gamma_j\xi}.
\]
固定任意 \(\xi\neq 0\) 与步长 \(\Delta s>0\)，对任意起点 \(s_0>0\) 定义等距采样序列
\(
y_n:=F(\xi,s_0+n\Delta s)
\)。
则 \(y_n\) 是 \(K\) 项指数和
\[
y_n=\sum_{j=1}^{K}c_j(\xi)\,\lambda_j^{n},
\qquad
\lambda_j=e^{-\Delta s\,\delta_j},
\qquad
c_j(\xi)=4\pi a_j\frac{\delta_j}{1+\delta_j}\,e^{-s_0\delta_j}e^{-\mathrm{i}\gamma_j\xi}.
\]
因此当采样长度 \(n\) 足够时，等距 Hankel 主块的秩恰为 \(K\)，并可由两块相邻 Hankel 主块构造的矩阵 pencil 在代数意义下唯一恢复 \(\{\lambda_j\}_{j=1}^{K}\)，从而恢复 \(\{\delta_j\}_{j=1}^{K}\)。
同理，固定任意 \(s>0\) 并对 \(\xi\) 作等距有限采样，则 \(\xi\mapsto F(\xi,s)\) 为 \(K\) 项 Fourier 和；
在无碰撞与一般位置条件下，可由离散 Fourier/Prony 型反演恢复 \(\{\gamma_j\}_{j=1}^{K}\)。
因此原子缺陷的重建可分解为两个相互独立的低秩谱反演问题；其原子数 \(K\) 同时等于可见 Hankel 秩与缺陷维数证书。
若以张量格点同时采样 \((\xi,s)\) 两轴，则总采样量级为 \(O(K^2)\)。
\end{conclusion}
\begin{proof}[证明要点]
固定 \(\xi\) 时，\(s\)-采样序列为指数和，Hankel 秩等于指数项数且可由 Prony/矩阵 pencil 恢复指数基底；
固定 \(s\) 时同理。
\end{proof}

\begin{conclusion}[压力解析半径的判别式封口：累积量增长率的阶乘--指数上界]\label{con:xi-terminal-pressure-analytic-radius-cumulant-bound}
令 \(P_2(\theta)\) 为纯碰撞势切片的压力函数（例如命题 \ref{prop:near1-pole-diffusive-sp} 的设定），并记其累积量
\[
\kappa_n:=P_2^{(n)}(0)\qquad(n\ge 1).
\]
在 Perron 分支的代数延拓口径下，\(P_2\) 的分歧点由判别式重根条件给出，且在 \(\theta\)-平面形成晶格
\[
\theta=\log u_j+2\pi i k\qquad(k\in\ZZ),
\]
其中 \(u_j\) 遍历判别式根（推论 \ref{cor:sync-kernel-analytic-radius-discriminant-formula}）。
因此 \(P_2\) 在 \(\theta=0\) 的 Taylor 收敛半径等于到最近分歧点的距离 \(R_\theta\)（推论 \ref{cor:pressure-analytic-radius}），并在真实输入 \(40\) 状态核的纯碰撞切片上具有数值证书
% Generated by scripts/exp_real_input_40_collision_branch_radius.py
\[
\begin{aligned}
u_0 &\approx -0.09334762302241694,
\\
u_\pm &\approx -3.078326188488792 \pm 3.456761347287067\,i,
\\
\log u_\pm &\approx 1.532286026474262 \pm 2.298350888853574\,i,
\\
R_\theta &\approx 2.76230289346087.
\end{aligned}
\]

从而 \(R_\theta\approx 2.76230289346087\)。
\par
对任意 \(0<r<R_\theta\)，由 Cauchy 估计得到一组严格的阶乘--指数上界：令
\(
C(r):=\sup_{|\theta|=r}|P_2(\theta)|
\)，则对所有 \(n\ge 0\) 有
\[
\boxed{\;
|\kappa_n|\le C(r)\,n!\,r^{-n}.\;}
\]
因此在固定倾斜窗口 \(|\theta|\le r<R_\theta\) 内，任意有限阶截断误差以 \((|\theta|/r)^{N+1}\) 指数衰减；
从而常数级自由能/熵在 \(N=O(\log(1/\varepsilon))\) 阶累积量内即可完成 \(\varepsilon\)-精度的证书化近似，其复杂度由单一可计算常数 \(R_\theta\) 封口。
\end{conclusion}
\begin{proof}[证明要点]
分歧点与解析半径见推论 \ref{cor:sync-kernel-analytic-radius-discriminant-formula}--\ref{cor:pressure-analytic-radius}；
导数界为解析函数在圆盘上的标准 Cauchy 估计。
\end{proof}

\begin{conclusion}[primitive 微正则熵的维数律：能量轴的 \texorpdfstring{$1/2$}{1/2} 维与原子因子分离]\label{con:xi-terminal-primitive-microcanonical-dim-half}
记 \(p_{n,k}\) 为长度 \(n\) 的 primitive 原子中“能量”为 \(k\) 的计数（见附录 \S\ref{app:sync-kernel-weighted-abel-mertens-analytic-radius-edgeworth}），并记
\(\sigma^2=P''(0)=11/102\)（定理 \ref{thm:sync-kernel-even-edgeworth-c4}）。
则对任意固定 \(X>0\)，当 \(n\to\infty\) 且 \(k=n/2+x\sqrt n\)（\(|x|\le X\)）时成立一致渐近
\[
p_{n,k}
=\frac{1}{n}\cdot\frac{3^n}{\sqrt{2\pi\sigma^2 n}}
\exp\!\left(-\frac{x^2}{2\sigma^2}\right)\cdot\bigl(1+o(1)\bigr).
\]
因此微正则熵满足
\[
\log p_{n,k}
=n\log3-\frac{3}{2}\log n-\frac{x^2}{2\sigma^2}-\frac{1}{2}\log(2\pi\sigma^2)+o(1).
\]
该展开揭示两条可分离的结构因子：能量分布在可见尺度上具有 \(\sqrt n\) 宽度（故能量轴的有效熵维数为 \(1/2\)），
并且前因子中的 \(n^{-1}\) 与 \(n^{-1/2}\) 可分别解释为 Witt/Möbius 的 primitive 原子因子与局部 CLT 的形状因子。
\end{conclusion}
\begin{proof}[证明要点]
将定理 \ref{thm:sync-kernel-even-edgeworth-c4} 的中心化局部展开与
\(\sum_k p_{n,k}=3^n/n\cdot(1+o(1))\)
合并，并用注 \ref{rem:sync-kernel-edgeworth-log-form} 的对数形式重写即可。
\end{proof}

\begin{conclusion}[RH-sharp 的算术相分裂：primitive 次级项由奇偶与最小奇素因子锁定]\label{con:xi-terminal-rhsharp-arithmetic-phase-splitting}
在 RH-sharp 且临界模谱相位唯一为 \(-\sqrt\lambda\) 的情形（见 \S\ref{app:real-input-40-finite-rh}），primitive 计数 \(p_n\) 的平方根级误差不形成单一“解析相”，而在偶长层发生严格算术分裂：
在 \(n\) 为偶数时，Möbius 反演中的 \(d=1,2\) 两项精确抵消 \(\lambda^{n/2}\) 级振荡。
进一步区分 \(3\mid n\) 与 \(3\nmid n\) 可得更强主次项展开
\[
p_n=\frac{\lambda^n}{n}-\frac{\lambda^{n/3}}{n}+O\!\left(\frac{\lambda^{n/4}}{n}\right),
\qquad n\ \text{偶且}\ 3\mid n,
\]
\[
p_n=\frac{\lambda^n}{n}+(-1)^{n/2+1}\frac{\lambda^{n/4}}{n}+O\!\left(\frac{\lambda^{n/5}}{n}\right),
\qquad n\ \text{偶且}\ 3\nmid n.
\]
更一般地，对任意奇素因子 \(p\mid n\)，Möbius 反演中的 \(d=p\) 项必贡献量级 \(\lambda^{n/p}/n\) 的修正（符号由 \(\mu(p)=-1\) 锁定）。
令
\[
p_{\min}(n):=\min\{p:\ p\ \text{为奇素数且}\ p\mid n\},
\qquad(\text{若无奇素因子则令 }p_{\min}(n)=+\infty),
\]
则在偶长层的主次级误差指数由 \(\max(\tfrac14,\,1/p_{\min}(n))\) 决定；
因子结构成为“误差相图”的必要坐标。
\end{conclusion}
\begin{proof}[证明要点]
见注 \ref{rem:finite-rh-even-strengthen} 与注 \ref{rem:finite-rh-factor-controls-secondary}。
\end{proof}

\begin{conclusion}[常数层毛发的准线性比特复杂度：primitive 抽取与有限部常数的短截断证书]\label{con:xi-terminal-constant-layer-bit-complexity}
以二阶碰撞核 \(A_2\) 为例，对任意 \(N\) 可先在 \(O(N)\) 时间生成周期迹序列 \(\{t_n\}_{n\le N}\)，再以约数和筛法实现的 Witt/Möbius 反演在 \(O(N\log N)\) 时间生成全部 primitive 轨道数 \(\{p_n(A_2)\}_{n\le N}\)。
并且在该核的 Abel/Mertens 有限部常数计算中，Möbius--Abel 截断即可给出极小尾项代理：
例如取 \(k_{\max}=200\) 时，尾项代理达到 \(4.712\times 10^{-79}\) 量级（表 \ref{tab:collision_kernel_A2_finite_part}）。
因此常数层不变量在“对数精度预算”意义下的计算成本主要由精度位数控制，而不由系统维度或长轨枚举控制。
\end{conclusion}
\begin{proof}[证明要点]
迹到 primitive 的准线性计算由 Witt/Möbius 反演的离散卷积结构与筛法实现得到；
有限部常数的短截断与尾项代理由表 \ref{tab:collision_kernel_A2_finite_part} 所示证书封口。
\end{proof}

\begin{conclusion}[时间纤维的 Legendre 熵谱：Poisson 核的 Rényi 体积由 \texorpdfstring{$P_\nu(\cosh a)$}{P_nu(cosh a)} 封口]\label{con:xi-terminal-timefiber-legendre-renyi-spectrum}
令 \(\Theta\sim\mathrm{Unif}[0,2\pi)\)，取 \(0<\varrho<1\) 并定义
\[
Z_{\varrho}:=\frac{1+\varrho^2+2\varrho\cos\Theta}{1-\varrho^2},
\qquad
Y_{\varrho}:=Z_{\varrho}^{-1}.
\]
引入双曲半径参量 \(a:=2\,\operatorname{artanh}\varrho\)（即 \(\varrho=\tanh(a/2)\)、\(\cosh a=\frac{1+\varrho^{2}}{1-\varrho^{2}}\)），则 \(Z_{\varrho}=\cosh a+\sinh a\cos\Theta\)，并且对任意 \(\nu\in\CC\) 有严格矩恒等式
\[
\mathbb{E}\bigl[Z_{\varrho}^{\,\nu}\bigr]=P_{\nu}(\cosh a)
\]
（定理 \ref{thm:app-legendre-moment-timefiber}）。
因此对任意 \(q\in\CC\) 有
\[
\mathbb{E}[Y_{\varrho}^{q}]=\mathbb{E}[Z_{\varrho}^{-q}]=P_{-q}(\cosh a),
\]
从而时间纤维的 Rényi 型对数体积谱可被完全写成 Legendre 函数的对数。
并且
\(
\mathbb{E}[Y_{\varrho}]=1
\)、
\(
\mathbb{E}[\log Y_{\varrho}]=\log(1-\varrho^{2})
\)
（定理 \ref{thm:app-endpoint-arcsine-duality}），从而对数体积势满足闭式
\[
-\mathbb{E}[\log Y_{\varrho}]=-\log(1-\varrho^2)=2\log\cosh\!\left(\frac{a}{2}\right),
\]
把双曲尺度 \(a\) 与熵势在同一账本中刚性锁定。
\end{conclusion}
\begin{proof}[证明要点]
见定理 \ref{thm:app-endpoint-arcsine-duality} 与定理 \ref{thm:app-legendre-moment-timefiber}。
\end{proof}

\paragraph{双频全息刚性、有限秩核唯一性与 \texorpdfstring{\v{C}ech}{Cech} 可实现性禁阻}
在结论 \ref{con:xi-terminal-defect-measure-fourier-laplace-holographic}--\ref{con:xi-terminal-atomic-defect-hankel-prony}
已经把有限原子缺陷测度压缩为二维 Fourier--Laplace 指纹之后，还可进一步把“有限采样的唯一恢复”、
“对角剖面的核化刚性”、“cyclotomic 粘合障碍对可见实现的排斥”以及“采样数与高度相关位预算的分裂”
压缩为同一条跨接口链。

\begin{theorem}[双频有限采样的全息刚性]\label{thm:xi-bifrequency-finite-sampling-holographic-rigidity}
\leanverified{paper\_xi\_bifrequency\_finite\_sampling\_holographic\_rigidity}
沿用结论 \ref{con:xi-terminal-defect-measure-fourier-laplace-holographic} 的记号，并设
\[
\nu=\sum_{j=1}^{\kappa} a_j\,\delta_{(\gamma_j,\delta_j)},
\qquad
a_j>0,\ \delta_j>0,
\]
其中 \(\delta_1,\dots,\delta_\kappa\) 两两不同。
固定 \(s_0>0\)、\(\Delta s>0\)，并取两个非零扫描频率 \(\xi_1,\xi_2\in\RR\setminus\{0\}\) 满足
\[
\frac{\xi_1}{\xi_2}\notin\QQ.
\]
对 \(\ell\in\{1,2\}\) 与 \(0\le n\le 2\kappa-1\)，定义
\[
y_n^{(\ell)}:=F(\xi_\ell,s_0+n\Delta s).
\]
若对应的阶 \(\kappa\) Hankel 主块均非退化，则全部 \(4\kappa\) 个采样值
\[
\{y_n^{(\ell)}:\ 0\le n\le 2\kappa-1,\ \ell=1,2\}
\]
唯一确定整个位相--高度缺陷测度 \(\nu\)。
\end{theorem}
\begin{proof}
由结论 \ref{con:xi-terminal-atomic-defect-hankel-prony}，对每个固定的 \(\ell\) 都有
\[
y_n^{(\ell)}=\sum_{j=1}^{\kappa} c_j^{(\ell)}\lambda_j^n,
\qquad
\lambda_j=e^{-\Delta s\,\delta_j},
\qquad
c_j^{(\ell)}=
4\pi a_j\frac{\delta_j}{1+\delta_j}e^{-s_0\delta_j}e^{-\mathrm{i}\gamma_j\xi_\ell}.
\]
非退化 Hankel 条件保证 Prony--Hankel 反演可唯一恢复 \(\{\lambda_j,c_j^{(\ell)}\}_{j=1}^{\kappa}\)。
由于 \(\delta_j\) 两两不同，\(\lambda_j\) 亦两两不同，故可用共同的 \(\lambda_j\) 把两条频率通道规范对齐。
于是
\[
\delta_j=-\frac{1}{\Delta s}\log \lambda_j
\]
被唯一确定。

令
\[
A_j:=4\pi a_j\frac{\delta_j}{1+\delta_j}e^{-s_0\delta_j}>0,
\]
则
\[
c_j^{(\ell)}=A_j e^{-\mathrm{i}\gamma_j\xi_\ell}
\qquad (\ell=1,2).
\]
若另有 \(\gamma_j'\in\RR\) 产生同一对系数，则
\[
e^{-\mathrm{i}(\gamma_j-\gamma_j')\xi_1}
=
e^{-\mathrm{i}(\gamma_j-\gamma_j')\xi_2}
=1.
\]
因此
\[
(\gamma_j-\gamma_j')\xi_1,\ (\gamma_j-\gamma_j')\xi_2\in 2\pi\ZZ.
\]
若 \(\gamma_j\neq \gamma_j'\)，则
\(
\xi_1/\xi_2\in\QQ
\)，与假设矛盾。故 \(\gamma_j\) 唯一。
最后由
\[
a_j=A_j\frac{1+\delta_j}{4\pi\delta_j}e^{s_0\delta_j}
\]
恢复全部权重 \(a_j\)。故 \(\nu\) 被全部 \(4\kappa\) 个采样值唯一确定。证毕。
\end{proof}

\begin{theorem}[对角全息的核刚性]\label{thm:xi-diagonal-holography-kernel-rigidity}
设
\[
H_r(x)=\sum_{j=1}^{\kappa_r}
\frac{w_j^{(r)}}{(x-\gamma_j^{(r)})^2+(1+\delta_j^{(r)})^2}
\qquad (r=1,2)
\]
为两组有限缺陷扫描剖面，其中 \(w_j^{(r)}>0\)、\(\delta_j^{(r)}>0\)。
令 \(K_r\) 为定理 \ref{thm:xi-basepoint-scan-finite-rank-rkhs} 所赋予的有限秩核：
\[
K_r(x,y)=
\sum_{j=1}^{\kappa_r}
\frac{w_j^{(r)}}{
\bigl(x-\gamma_j^{(r)}+\mathrm{i}(1+\delta_j^{(r)})\bigr)
\bigl(y-\gamma_j^{(r)}-\mathrm{i}(1+\delta_j^{(r)})\bigr)}.
\]
若存在非空开区间 \(I\subset\RR\) 使
\[
H_1(x)=H_2(x)
\qquad(x\in I),
\]
则两组带权缺陷原子多重集按重数一致，并且
\[
K_1(x,y)\equiv K_2(x,y)
\qquad((x,y)\in\RR^2).
\]
从而二者的最小有限秩 RKHS 实现规范酉同构。
\end{theorem}
\begin{proof}
由定理 \ref{thm:xi-comoving-horizon-scan-fourier-inversion} 的第 \(3\) 条，只要两组有限缺陷剖面在任意非平凡开区间上一致，
其带权原子多重集
\[
\{(\gamma_j^{(1)},\delta_j^{(1)},w_j^{(1)})\}_{j}
\quad\text{与}\quad
\{(\gamma_j^{(2)},\delta_j^{(2)},w_j^{(2)})\}_{j}
\]
就按重数一致。
经适当重排后，可设两侧参数逐项相同。
于是 \(K_1\) 与 \(K_2\) 的显式 Cauchy 和式完全一致，故
\[
K_1\equiv K_2.
\]
由定理 \ref{thm:xi-basepoint-scan-finite-rank-rkhs}，最小有限秩实现均由同一特征映射骨架生成，
因此两者在自然基变换下规范酉同构。证毕。
\end{proof}

\begin{theorem}[全息实现与 \texorpdfstring{\v{C}ech}{Cech} 粘合障碍互斥]\label{thm:xi-holography-cech-mutual-exclusion}
固定一个前缀地址 \(a\)，并沿用定理 \ref{thm:cdim-cech-cyclotomic-quantization} 的记号。
若存在非平凡角色 \(\chi\neq \ind\)，使归一化扭曲行列式
\[
\det\!\Bigl(I-\frac{t}{\lambda}B_\chi\Bigr)
\]
在单位圆上出现某个 \(r>1\) 阶本原单位根零点，则与地址 \(a\) 对应的查询不可能同时由任何有限正缺陷测度
\(\nu\) 在可见层中实现。
\end{theorem}
\begin{proof}
由定理 \ref{thm:cdim-cech-cyclotomic-quantization} 的反向部分，
上述本原单位根零点强制
\[
\chi_*([\alpha])\neq 0,
\qquad
[\alpha]\neq 0.
\]
再由命题 \ref{prop:null-as-h2-obstruction}，
该地址处的纤维为空，读出必为 \(\mathsf{NULL}\)。

反设仍存在有限正缺陷测度 \(\nu\) 在可见层中实现同一地址。
则由结论 \ref{con:xi-terminal-defect-measure-fourier-laplace-holographic}，\(\nu\) 产生一份有限可见的边界全息指纹；
而由推论 \ref{cor:xi-unitary-slice-decidable}，任何非 \(\mathsf{NULL}\) 的可见实现都必须伴随一份有限正向证书。
这与 \([\alpha]\neq 0\) 时地址纤维为空的结论矛盾。
故有限正缺陷测度的可见实现与该类 cyclotomic 粘合障碍不能并存。证毕。
\end{proof}

\begin{theorem}[有限采样数与高度相关位预算的裂解]\label{thm:xi-finite-sampling-vs-height-bit-splitting}
在定理 \ref{thm:xi-bifrequency-finite-sampling-holographic-rigidity} 的设定下，进一步假设
\[
|\gamma_j|\le T,
\qquad
\delta_j\ge \delta_0>0
\qquad (1\le j\le \kappa).
\]
则几何重建所需的全息采样数始终恰为 \(4\kappa\)，与高度窗 \(T\) 无关。
然而，若坚持固定图表端点读出且不外置共动前缀，则任一端点余量证书机制在最坏情形下都必须支付位预算下界
\[
2\log_2 T+\log_2\frac{1}{\delta_0}+O(1).
\]
反之，若为每个原子外置精确共动前缀 \(x_{0,j}=\gamma_j\)，则高度相关项完全消失，
每个原子的几何证书预算都降为仅依赖于 \(\delta_0\) 与协议常数的 \(O(1)\)，
从而总预算为 \(O(\kappa)\)。
\end{theorem}
\begin{proof}
定理 \ref{thm:xi-bifrequency-finite-sampling-holographic-rigidity} 已经给出：
当 \(\kappa\) 固定时，仅需两条非共比频率通道上的 \(2\kappa\) 个 \(s\)-采样，
总采样点数即为 \(4\kappa\)，与 \(T\) 无关。

另一方面，在固定 Cayley 图表下，定理
\ref{thm:xi-offcritical-endpoint-resolution-lower-bound} 证明了：
若要求对全部满足 \(|\gamma|\le T\)、\(\delta\ge\delta_0\) 的离切片原子给出端点余量证书，
则最坏情形位预算必满足
\[
p\ge
\log_2\!\Bigl(\frac{T^2+(1+\delta_0)^2}{4\delta_0}\Bigr)
=
2\log_2 T+\log_2\frac{1}{\delta_0}+O(1).
\]
故不外置前缀时，高度相关位负载不可消去。

若允许对象层外置可寻址前缀并取共动局部化 \(x_{0,j}=\gamma_j\)，则推论
\ref{cor:xi-comoving-prefix-bit-budget-gain} 给出高度无关的常数门槛
\[
p\ge \log_2\!\Bigl(\frac{(1+\delta_0)^2}{4\delta_0}\Bigr),
\]
故每个原子的几何证书成本均为 \(O(1)\)，总预算遂为 \(O(\kappa)\)。
这就把“有限采样数”与“高度相关位预算”严格分裂为两类不同资源。证毕。
\end{proof}

\begin{corollary}[可判定而不可实现的查询类]\label{cor:xi-decidable-but-unrealizable-visible-query}
\leanverified{paper\_xi\_decidable\_but\_unrealizable\_visible\_query}
存在一类地址查询，其“是否可读出”在协议语义中具有有限判定证书，
但其对象本身并不由任何有限正缺陷测度在可见层中实现。
\end{corollary}
\begin{proof}
取满足定理 \ref{thm:xi-holography-cech-mutual-exclusion} 假设的地址。
由推论 \ref{cor:xi-unitary-slice-decidable}，
\(\mathsf{NULL}\) 与非 \(\mathsf{NULL}\) 两侧都带有有限证书，
故“该地址是否可读出”是可判定谓词。
而定理 \ref{thm:xi-holography-cech-mutual-exclusion} 表明，
一旦存在非平凡 cyclotomic 粘合障碍，该地址就不可能由任何有限正缺陷测度在可见层实现。
因此可判定性与可实现性在该制度内严格分离。证毕。
\end{proof}

\endinput

% (User input window)

\subsubsection{尺度核的质量塔分解与粗反卷积：有限阶局部算子、尾核矩与带宽良定性}\label{subsubsec:xi-scale-kernel-mass-tower-deconvolution}
以下条目把 logistic 尺度核在 Fourier 侧的 \(\sinh\) 乘积结构实现为一个可执行的质量塔卷积模型，
并由此导出一族有限阶局部反卷积算子及其在带宽受控制度下的稳定性证书；
尾核的可计算矩与四阶累积量给出最小非高斯偏差的可观测刻度，从而把 \(q=4\) 的首触发机制提升为结构必然性。

\begin{conclusion}[尺度核的无穷质量塔分解：logistic 核等价于独立 Laplace 级联系统]\label{con:xi-terminal-scale-kernel-mass-tower-laplace}
令尺度单峰核为
\[
\phi(s):=\frac{1}{4\cosh^{2}(s/2)}=\frac{e^{s}}{(1+e^{s})^{2}}\qquad(s\in\RR).
\]
对每个整数 \(n\ge 1\) 定义 Laplace 核
\[
g_n(s):=\frac{n}{2}e^{-n|s|},
\qquad
\widehat g_n(\xi)=\frac{1}{1+\xi^{2}/n^{2}}.
\]
则存在一族独立随机变量 \((X_n)_{n\ge 1}\)（\(X_n\) 的密度为 \(g_n\)），使得级数
\[
Y:=\sum_{n\ge 1}X_n
\]
在 \(L^{2}\) 与几乎处处意义下收敛，并且其和的密度严格等于 \(\phi\)。
等价地，令 \(\phi_{\le N}:=g_1*\cdots*g_N\)，则 \(\phi_{\le N}\to\phi\)（\(L^2\) 且弱收敛），从而在分布意义下可记为
\[
\phi=g_1*g_2*\cdots .
\]
该分解把 \(\phi\) 解释为一条“整数质量谱”\(\{n\}_{n\ge 1}\) 的逐层卷积封口：每一层 \(g_n\) 是一维局部算子 \((1-\frac{1}{n^{2}}\partial_s^{2})^{-1}\) 的 Green 核，整条质量塔的乘积谱恰闭合为 \(\phi\)。
\end{conclusion}
\begin{proof}[证明要点]
由命题 \ref{prop:xi-bulk-curvature-tomography-identifiability} 的 Fourier 闭式 \(\widehat\phi(\xi)=\pi\xi/\sinh(\pi\xi)\) 及 \(\sinh\) 的 Hadamard 乘积
\(
\sinh(\pi\xi)=\pi\xi\prod_{n\ge 1}(1+\xi^{2}/n^{2})
\)
得
\(
\widehat\phi(\xi)=\prod_{n\ge 1}\widehat g_n(\xi)
\)。
取独立随机变量 \(X_n\)（密度为 \(g_n\)），则其特征函数为 \(\widehat g_n\)。
令部分和 \(Y_N=\sum_{n=1}^{N}X_n\)，则 \(Y_N\) 的特征函数为 \(\prod_{n=1}^{N}\widehat g_n\)，故 \(Y_N\) 的密度为 \(\phi_{\le N}\)。
又因 \(\sum_{n\ge 1}\Var(X_n)=2\sum_{n\ge 1}n^{-2}<\infty\)，得 \((Y_N)\) 在 \(L^2\) 中 Cauchy，且由 Kolmogorov 判据可得几乎处处收敛。
极限变量 \(Y\) 的特征函数为 \(\widehat\phi\)，从而其密度为 \(\phi\)。证毕。
\end{proof}

\begin{conclusion}[部分反卷积即分辨率旋钮：有限阶局部算子把无限反卷积降为尾核卷积]\label{con:xi-terminal-scale-kernel-partial-deconvolution-resolution}
对 \(N\ge 1\) 定义有限阶局部算子（以 Fourier 乘子理解）
\[
\mathcal D_N:=\prod_{n=1}^{N}\Bigl(1-\frac{1}{n^2}\partial_s^2\Bigr),
\qquad
\widehat{(\mathcal D_N f)}(\xi)=M_N(\xi)\,\widehat f(\xi),
\qquad
M_N(\xi):=\prod_{n=1}^{N}(1+\xi^{2}/n^{2}).
\]
并令尾核 \(\phi_{>N}\) 为结论 \ref{con:xi-terminal-scale-kernel-mass-tower-laplace} 中尾部和
\(
Y_{>N}:=\sum_{n>N}X_n
\)
的密度（亦即 \(\widehat{\phi_{>N}}(\xi)=\prod_{n>N}(1+\xi^{2}/n^{2})^{-1}\)）。
则对任意有限符号测度 \(\nu\)，若 \(p=\phi*\nu\)，则成立严格恒等式
\[
\boxed{\;
\mathcal D_N p=\phi_{>N}*\nu.\;}
\]
因而 \(\mathcal D_N\) 在制度上实现“去除前 \(N\) 层质量壳”的粗反卷积：当 \(N\to\infty\) 时 \(\phi_{>N}\Rightarrow\delta_0\)，从而 \(\mathcal D_N p\Rightarrow \nu\) 于分布意义下成立。
该结论给出一条可执行的连续族：反卷积由 \(N\) 明确调控的分辨率流替代了“要么精确要么崩溃”的二元机制。
\end{conclusion}
\begin{proof}[证明要点]
由结论 \ref{con:xi-terminal-scale-kernel-mass-tower-laplace}，\(\widehat\phi(\xi)=M_N(\xi)^{-1}\widehat{\phi_{>N}}(\xi)\)。
对 \(p=\phi*\nu\) 取 Fourier 变换得 \(\widehat p=\widehat\phi\cdot\widehat\nu\)，从而
\[
\widehat{(\mathcal D_N p)}(\xi)=M_N(\xi)\widehat p(\xi)=\widehat{\phi_{>N}}(\xi)\widehat\nu(\xi)=\widehat{(\phi_{>N}*\nu)}(\xi),
\]
得到恒等式。
又因 \(\Var(Y_{>N})\to 0\)，故 \(\phi_{>N}\Rightarrow\delta_0\)，从而 \(\phi_{>N}*\nu\Rightarrow\nu\)。证毕。
\end{proof}

\begin{conclusion}[尾核的矩完全可计算：分辨率尺度 \texorpdfstring{$\sigma_N\sim \sqrt{2/N}$}{sigma\_N ~ sqrt(2/N)} 与 Wasserstein 控制]\label{con:xi-terminal-scale-kernel-tail-moments-w2}
沿用结论 \ref{con:xi-terminal-scale-kernel-mass-tower-laplace} 的随机变量 \(X_n\)，令 \(Y_{>N}:=\sum_{n>N}X_n\)，其密度为 \(\phi_{>N}\)。
则有可加矩证书
\[
\Var(Y_{>N})=\sum_{n>N}\Var(X_n)=2\sum_{n>N}\frac{1}{n^2}=: \sigma_N^2,
\qquad
\kappa_4(Y_{>N})=\sum_{n>N}\kappa_4(X_n)=12\sum_{n>N}\frac{1}{n^4}.
\]
并且对任意具有有限二阶矩的概率测度 \(\nu\) 有可证书化的 \(W_2\) 上界
\[
\boxed{\;
W_2\bigl(\phi_{>N}*\nu,\nu\bigr)\ \le\ \sigma_N
\ =\ \sqrt{2\sum_{n>N}n^{-2}}
\ \sim\ \sqrt{\frac{2}{N}}\quad(N\to\infty).\;}
\]
因此若允许分辨率 \(\sigma>0\)，取 \(N\asymp 2/\sigma^{2}\) 即可保证 \(W_2\)-误差 \(\le \sigma\)；
在该制度下，分辨率成本从指数型降为幂律型。
\end{conclusion}
\begin{proof}[证明要点]
对 \(X_n\sim g_n\) 直接计算得到 \(\Var(X_n)=2/n^2\)、\(\kappa_4(X_n)=12/n^4\)，累积量对独立和可加，从而得矩证书。
对 \(W_2\) 上界，取耦合：令 \(U\sim \nu\)，令 \(V\sim \phi_{>N}\) 且独立，则 \((U+V,U)\) 是 \((\phi_{>N}*\nu,\nu)\) 的一个耦合，故
\(
W_2^2(\phi_{>N}*\nu,\nu)\le \EE[|V|^2]=\Var(V)=\sigma_N^2
\)。
渐近 \(\sum_{n>N}n^{-2}\sim 1/N\) 由积分判别法得到。证毕。
\end{proof}

\begin{conclusion}[尾核的高斯化与四阶首破缺：非高斯性由 \texorpdfstring{$q=4$}{q=4} 量级唯一主导]\label{con:xi-terminal-scale-kernel-tail-gaussianization-kurtosis}
令 \(\sigma_N^2=\Var(Y_{>N})\)，并定义标准化尾变量 \(Z_N:=Y_{>N}/\sigma_N\)。
则有弱收敛 \(Z_N\Rightarrow \mathcal N(0,1)\)，且其标准化超额峰度（excess kurtosis）具有闭式
\[
\boxed{\;
\gamma_2(N):=\frac{\kappa_4(Y_{>N})}{\sigma_N^4}
=3\,\frac{\sum_{n>N}n^{-4}}{\bigl(\sum_{n>N}n^{-2}\bigr)^2}
\sim \frac{1}{N}\quad(N\to\infty).\;}
\]
并且所有奇阶累积量严格为零。
因此尾核的首个非高斯偏差不可避免地发生在四阶：在尺度流的可观测校正中，最小破缺阶被结构性锁定为 \(4\)，并以 \(N^{-1}\) 的强度进入。
\end{conclusion}
\begin{proof}[证明要点]
由于 \(X_n\) 对称，故全部奇阶累积量为零，并在独立和下保持为零。
对 \(Z_N\) 的中心极限定理可由 Lindeberg 条件给出：\(\max_{n>N}\Var(X_n)/\sigma_N^2\to 0\) 且 \(\sum_{n>N}\Var(X_n)=\sigma_N^2\)。
峰度公式由结论 \ref{con:xi-terminal-scale-kernel-tail-moments-w2} 的累积量证书代入并化简得到；
渐近由 \(\sum_{n>N}n^{-4}\sim (3N^3)^{-1}\)、\(\sum_{n>N}n^{-2}\sim N^{-1}\) 给出。证毕。
\end{proof}

\begin{conclusion}[带宽约束下的条件数上界：粗反卷积在任何固定带宽内均一致良定]\label{con:xi-terminal-scale-kernel-bandwidth-conditioned}
设观测为 \(\widetilde p=p+\eta\)，并假设噪声在频带 \(|\xi|\le\Omega\) 内满足
\[
\sup_{|\xi|\le\Omega}|\widehat\eta(\xi)|\le\varepsilon.
\]
则对任意 \(N\ge 1\)，有限阶算子输出误差在同一频带内满足
\[
\sup_{|\xi|\le\Omega}\bigl|\widehat{(\mathcal D_N\eta)}(\xi)\bigr|
\le
\varepsilon\cdot \sup_{|\xi|\le\Omega}M_N(\xi)
\le
\varepsilon\cdot \exp\!\Bigl(\Omega^2\sum_{n=1}^{\infty}\frac{1}{n^2}\Bigr)
=\varepsilon\cdot \exp\!\Bigl(\frac{\pi^2}{6}\Omega^2\Bigr).
\]
因此在任何固定带宽 \(\Omega\) 内，\(\mathcal D_N\) 的噪声放大因子对 \(N\) 一致有界；
粗反卷积的病态性并非来自“阶数增大”，而来自“试图在无限带宽上实现 \(\delta_0\) 极限”。
该结论为尺度流的可计算相图提供严格基础：在带宽受控的制度下，分辨率可按 \(N\sim\sigma^{-2}\) 幂律提升而不触发指数型条件数爆炸。
\end{conclusion}
\begin{proof}[证明要点]
由结论 \ref{con:xi-terminal-scale-kernel-partial-deconvolution-resolution} 的 Fourier 乘子表示，
\(
\widehat{(\mathcal D_N\eta)}=M_N\widehat\eta
\)
于 \(|\xi|\le\Omega\) 上成立。
并且对 \(x\ge 0\) 有 \(\log(1+x)\le x\)，故
\[
\log M_N(\xi)=\sum_{n=1}^{N}\log\Bigl(1+\frac{\xi^2}{n^2}\Bigr)\le \xi^2\sum_{n=1}^{N}\frac{1}{n^2}
\le \Omega^2\sum_{n=1}^{\infty}\frac{1}{n^2}=\frac{\pi^2}{6}\Omega^2,
\]
从而得所示一致上界。证毕。
\end{proof}

\begin{conclusion}[有限阶谱乘子的 \texorpdfstring{$\Gamma$}{Gamma}-闭式与小频率展开：局部算子的解析完备性]\label{con:xi-terminal-scale-kernel-gamma-closed-form}
记有限阶谱乘子
\[
M_N(\xi):=\prod_{n=1}^{N}(1+\xi^2/n^2).
\]
则 \(M_N\) 具有严格的 \(\Gamma\)-比闭式
\[
\boxed{\;
M_N(\xi)=\frac{\Gamma(N+1+\mathrm{i}\xi)\Gamma(N+1-\mathrm{i}\xi)}{(N!)^2\,\Gamma(1+\mathrm{i}\xi)\Gamma(1-\mathrm{i}\xi)}.\;}
\]
从而在 \(|\xi|=o(N)\) 下有可审计展开
\[
\log M_N(\xi)=\xi^2\sum_{n=1}^{N}\frac{1}{n^2}-\frac{\xi^4}{2}\sum_{n=1}^{N}\frac{1}{n^4}+O\!\Bigl(\xi^6\sum_{n=1}^{N}\frac{1}{n^6}\Bigr),
\]
揭示出“二阶（\(\zeta(2)\)）决定主放大、四阶（\(\zeta(4)\)）决定首修正”的严格层级。
该闭式使 \(\mathcal D_N\) 的解析结构可被有限步复算，并允许在带宽与精度给定时对 \(N\) 作可证伪的最优截断。
\end{conclusion}
\begin{proof}[证明要点]
将 \(M_N(\xi)=\prod_{n=1}^{N}(n^2+\xi^2)/n^2=\prod_{n=1}^{N}(n+\mathrm{i}\xi)(n-\mathrm{i}\xi)/(N!)^2\)，再用有限乘积恒等式
\(
\prod_{n=1}^{N}(n+z)=\Gamma(N+1+z)/\Gamma(1+z)
\)
即得 \(\Gamma\)-比闭式。
对小频率展开，使用 \(\log(1+u)=u-\frac{u^2}{2}+O(u^3)\) 并对 \(u=\xi^2/n^2\) 求和即可。证毕。
\end{proof}

\begin{conclusion}[尺度流的两阶段可恢复性：幂律分辨率与指数精确性的制度分离]\label{con:xi-terminal-scale-kernel-two-regime-recoverability}
对任意 \(\sigma>0\)，取 \(N(\sigma)\) 使 \(\sigma_{N(\sigma)}\le\sigma\)（其中 \(\sigma_N\) 如结论 \ref{con:xi-terminal-scale-kernel-tail-moments-w2}）。
则由结论 \ref{con:xi-terminal-scale-kernel-partial-deconvolution-resolution} 有
\[
\mathcal D_{N(\sigma)}p=\phi_{>N(\sigma)}*\nu,
\]
给出一条分辨率为 \(\sigma\) 的稳定重建；其代价为局部算子阶数 \(2N(\sigma)\asymp\sigma^{-2}\)，且在固定带宽下条件数一致可控（结论 \ref{con:xi-terminal-scale-kernel-bandwidth-conditioned}）。
\par
相反，若要求 \(\sigma\downarrow 0\) 的精确极限（即 \(\phi_{>N}\Rightarrow\delta_0\) 的无带宽极限），则必然进入指数不适定区：尾核逼近 \(\delta_0\) 要求同时控制任意高频，制度上必须允许带宽 \(\Omega\to\infty\)。
更精确地，由 \(\sinh(\pi\xi)=\pi\xi\prod_{n\ge 1}(1+\xi^2/n^2)\) 得极限乘子
\[
M_\infty(\xi):=\lim_{N\to\infty}M_N(\xi)=\frac{\sinh(\pi\xi)}{\pi\xi},
\qquad
M_\infty(\xi)\sim \frac{1}{2\pi|\xi|}\,e^{\pi|\xi|}\quad(|\xi|\to\infty).
\]
因此若把“精确”理解为需要在 \(|\xi|\lesssim \Omega\) 上逼近 \(\delta_0\)，则噪声放大随 \(\Omega\) 指数增长；
若进一步把制度分辨率 \(\sigma\) 视为要求可见带宽 \(\Omega\asymp \sigma^{-1}\)，则放大因子量级为 \(\exp(\pi/\sigma)\)（仅差多项式因子），从而精度位数随 \(1/\sigma\) 线性增长。
因此“幂律可控”与“指数不可控”并非矛盾，而是由尺度选择（是否接受 \(\sigma\)-分辨率）所分离的两种制度。
\end{conclusion}
\begin{proof}[证明要点]
稳定重建部分由结论 \ref{con:xi-terminal-scale-kernel-partial-deconvolution-resolution} 的恒等式与结论 \ref{con:xi-terminal-scale-kernel-tail-moments-w2} 的 \(\sigma_N\) 控制给出；
在带宽受控制度下的良定性由结论 \ref{con:xi-terminal-scale-kernel-bandwidth-conditioned} 给出。
精确极限要求 \(\Omega\) 不受界，从而噪声放大必随 \(\Omega\) 发散，给出指数不适定区。证毕。
\end{proof}

\begin{conclusion}[\texorpdfstring{$(q=4)$}{(q=4)} 的正则化地位：四阶破缺决定可见相变的最小矩阶]\label{con:xi-terminal-scale-kernel-q4-minimal-visible-order}
在结论 \ref{con:xi-terminal-scale-kernel-mass-tower-laplace}--\ref{con:xi-terminal-scale-kernel-tail-gaussianization-kurtosis} 的质量塔模型中，尾核的全部奇阶累积量严格为零，而首个非高斯项由四阶累积量 \(\kappa_4\) 给出并以
\(\gamma_2(N)\sim N^{-1}\)
控制。
故任何以“二阶近似（Gaussian）+ 带宽受控的粗反卷积”为基底的可审计推断体系，其最小系统性偏差必出现在四阶；
换言之，若仅允许有限矩阶作为统计可见接口，则最小能够检测“非高斯边界/相变临界”之矩阶必为 \(4\)。
该事实为 \(q=4\) 的首越线现象提供完全结构性的解释：它并非偶然数值，而是尺度核尾部的对称性与质量塔叠加所共同强制的最小破缺阶。
\end{conclusion}
\begin{proof}[证明要点]
由结论 \ref{con:xi-terminal-scale-kernel-tail-gaussianization-kurtosis}，标准化尾变量的奇阶累积量为零，且四阶累积量给出首个非高斯偏差并以 \(N^{-1}\) 量级衰减。
在以二阶高斯近似为基底的制度中，任何有限阶矩统计量的系统偏差由最低非零累积量阶决定，故最小可见矩阶为 \(4\)。证毕。
\end{proof}

\subsubsection{行列式体积熵的树/环分层：Chow--Liu 投影、条件互信息与四阶显影}\label{subsubsec:xi-det-entropy-tree-loop-hierarchy}
在以 Gram/相关矩阵承载“体积势”的口径下，\(-\log\det\) 自然给出一个多体体积熵势；
本段把该势的可审计下界压缩为 Chow--Liu 树投影，并将剩余缺口解释为由环结构承载的不可逆体积亏损。
在一维排序与单调距离核制度中，最优树选择刚性退化为相邻链，从而把“树承载的体积熵”外部化为相邻间隔账本；
而在弱相干区，环熵的首个系统性贡献由条件互信息控制并从四阶开始显影，与本文反复出现的 \(q=4\) 首触发结构同型。

\begin{conclusion}[\texorpdfstring{$(\det)$}{(det)}-熵势的 Chow--Liu 树下界：行列式体积由最大生成树控制]\label{con:xi-det-entropy-chow-liu-tree-lower-bound}
令 \(R\in\RR^{\kappa\times\kappa}\) 为相关矩阵（\(R\succ0\)、\(\mathrm{diag}(R)=\mathbf 1\)），定义行列式体积熵势
\[
S_{\det}(R):=-\log\det R.
\]
对每条无向边 \((i,j)\) 定义二体互信息权重
\[
w_{ij}:=-\log(1-R_{ij}^2).
\]
则对任意生成树 \(T\) 都有
\[
S_{\det}(R)\ \ge\ \sum_{(i,j)\in T} w_{ij},
\qquad\text{从而}\qquad
S_{\det}(R)\ \ge\ \max_{T}\sum_{(i,j)\in T} w_{ij}.
\]
并且对给定 \(T\)，等号成立当且仅当 \(R^{-1}\) 的非对角非零支撑包含于 \(T\)（等价于 \(\mathcal N(0,R)\) 在图 \(T\) 上为树形高斯马尔可夫场）。
\end{conclusion}
\begin{proof}[证明要点]
取 \(X\sim\mathcal N(0,R)\)。
在固定树 \(T\) 下，Chow--Liu 树投影 \(\mathcal N_T\)（见结论 \ref{con:xi-det-entropy-loop-entropy-kl-gap} 的设定）满足
\(
2D_{\mathrm{KL}}(\mathcal N(0,R)\,\|\,\mathcal N_T)
=S_{\det}(R)-\sum_{(i,j)\in T}w_{ij}\ge 0
\)，
从而得下界。
等号等价于 \(D_{\mathrm{KL}}=0\)，即 \(\mathcal N(0,R)=\mathcal N_T\)；
对高斯族这等价于精度矩阵仅在 \(T\) 的边上允许出现非零项。证毕。
\end{proof}

\begin{conclusion}[环熵即树近似的 \texorpdfstring{$D_{\mathrm{KL}}$}{DKL} 缺口：体积亏损的精确树分解]\label{con:xi-det-entropy-loop-entropy-kl-gap}
沿用结论 \ref{con:xi-det-entropy-chow-liu-tree-lower-bound} 的设定。
对给定生成树 \(T\) 定义树环熵
\[
\mathcal L_T(R):=S_{\det}(R)-\sum_{(i,j)\in T}w_{ij}
=-\log\det R+\sum_{(i,j)\in T}\log(1-R_{ij}^2).
\]
则 \(\mathcal L_T(R)\ge 0\) 恒成立。
若 \(X\sim\mathcal N(0,R)\)，并令 \(\mathcal N_T\) 为以 \(T\) 为图的 Chow--Liu 树形高斯投影（其一维边缘与树边二元边缘与 \(\mathcal N(0,R)\) 一致），则成立严格恒等式
\[
\mathcal L_T(R)=2D_{\mathrm{KL}}\!\bigl(\mathcal N(0,R)\,\|\,\mathcal N_T\bigr).
\]
因此
\(
\arg\min_T\mathcal L_T(R)=\arg\max_T\sum_{(i,j)\in T}w_{ij}
\)
给出最优树近似。
\end{conclusion}
\begin{proof}[证明要点]
令 \(R_T\) 为 \(\mathcal N_T\) 的相关矩阵。
树形投影 \(\mathcal N_T\) 的协方差由树边相关唯一决定，并满足
\[
\det R_T=\prod_{(i,j)\in T}(1-R_{ij}^2)
\]
（树形高斯的逐边条件方差乘积公式）。
另一方面，零均值高斯族的 KL 闭式给出
\[
2D_{\mathrm{KL}}(\mathcal N(0,R)\,\|\,\mathcal N(0,R_T))
=\tr(R_T^{-1}R)-\kappa+\log\frac{\det R_T}{\det R}.
\]
由于 \(R_T^{-1}\) 的非对角支撑包含于 \(T\)，且 \(R\) 与 \(R_T\) 在对角与树边上取值一致，故
\(
\tr(R_T^{-1}R)=\tr(R_T^{-1}R_T)=\kappa
\)，
从而
\(
2D_{\mathrm{KL}}=\log\frac{\det R_T}{\det R}
\)。
代入行列式因子化即得所示恒等式。证毕。
\end{proof}

\begin{conclusion}[单调距离核的树选择刚性：一维排序点的最优树必为相邻链]\label{con:xi-det-entropy-monotone-kernel-chain-mst}
设存在严格单调递减函数 \(A:(0,\infty)\to(0,1)\) 使
\[
R_{ij}=A(|s_i-s_j|),
\qquad
s_1<\cdots<s_\kappa.
\]
则边权 \(w_{ij}=-\log(1-A(|s_i-s_j|)^2)\) 亦严格随距离递减。
因此最大生成树唯一为路径图（相邻链）
\[
T_{\mathrm{chain}}:=\{(1,2),(2,3),\dots,(\kappa-1,\kappa)\},
\]
从而得到最强树下界
\[
S_{\det}(R)\ \ge\ \sum_{j=1}^{\kappa-1}-\log\!\bigl(1-A(\Delta_j)^2\bigr),
\qquad
\Delta_j:=s_{j+1}-s_j.
\]
\end{conclusion}
\begin{proof}[证明要点]
对任意割 \(\{1,\dots,j\}\,\cup\,\{j+1,\dots,\kappa\}\)，跨割的距离最小边唯一为 \((j,j+1)\)，且其权重最大。
由最大生成树的割性质（cut property），每条 \((j,j+1)\) 必属于最大生成树；合并所有 \(j\) 即得唯一性与链结构。证毕。
\end{proof}

\begin{conclusion}[三点局部环熵的闭式：链缺口等于条件互信息，且为部分相关的对数势]\label{con:xi-det-entropy-three-point-loop-entropy-partial-corr}
取任意 \(a<b<c\)，记 \(\rho_{ab}=R_{ab}\)、\(\rho_{bc}=R_{bc}\)、\(\rho_{ac}=R_{ac}\)，并令 \(R_{abc}\) 为对应三元相关子矩阵。
相对于链树 \((a,b)\)--\((b,c)\) 定义三点局部环熵
\[
\mathcal L^{(3)}_{a,b,c}
:=-\log\det R_{abc}+\log(1-\rho_{ab}^2)+\log(1-\rho_{bc}^2).
\]
则有恒等式
\[
\mathcal L^{(3)}_{a,b,c}
=-\log\bigl(1-\rho_{ac\mid b}^2\bigr),
\qquad
\rho_{ac\mid b}:=\frac{\rho_{ac}-\rho_{ab}\rho_{bc}}{\sqrt{(1-\rho_{ab}^2)(1-\rho_{bc}^2)}}.
\]
并且若 \(X\sim\mathcal N(0,R)\)，则
\[
\mathcal L^{(3)}_{a,b,c}=2I(X_a;X_c\mid X_b).
\]
\end{conclusion}
\begin{proof}[证明要点]
由三元相关矩阵行列式与部分相关的标准恒等式
\(
1-\rho_{ac\mid b}^2=\det R_{abc}/((1-\rho_{ab}^2)(1-\rho_{bc}^2))
\)
得到第一式。
对高斯族，条件互信息满足
\(
2I(X_a;X_c\mid X_b)=-\log(1-\rho_{ac\mid b}^2)
\)，
从而得第二式。证毕。
\end{proof}

\begin{conclusion}[四阶首触发律：远距衰减型弱相干区的环熵从四阶开始显影]\label{con:xi-det-entropy-q4-trigger-weak-coherence}
在结论 \ref{con:xi-det-entropy-three-point-loop-entropy-partial-corr} 的三点设定下，若
\(
|\rho_{ab}|,|\rho_{bc}|,|\rho_{ac}|\le\varepsilon\ll1
\)，
则有一致展开
\[
\rho_{ac\mid b}=(\rho_{ac}-\rho_{ab}\rho_{bc})+O(\varepsilon^3),
\qquad
\mathcal L^{(3)}_{a,b,c}=\rho_{ac\mid b}^2+O(\varepsilon^4).
\]
特别地，当满足远距衰减型关系 \(\rho_{ac}=\rho_{ab}\rho_{bc}+O(\varepsilon^3)\) 时，\(\rho_{ac\mid b}=O(\varepsilon^2)\) 且
\[
\mathcal L^{(3)}_{a,b,c}=O(\varepsilon^4),
\]
即局部环熵具有“最小四阶可见性”：二体相关 \(O(\varepsilon)\) 在该制度下不足以触发链缺口，缺口的首个系统性量级必为四阶。
\end{conclusion}
\begin{proof}[证明要点]
分母 \(\sqrt{(1-\rho_{ab}^2)(1-\rho_{bc}^2)}=1+O(\varepsilon^2)\)，故部分相关可由分子线性化得到。
又由 \(-\log(1-x)=x+O(x^2)\)（\(x\to0\)）并取 \(x=\rho_{ac\mid b}^2=O(\varepsilon^2)\)，得所示展开。证毕。
\end{proof}

\begin{conclusion}[指数尾尺度核下的显式四阶标度：\texorpdfstring{$\mathcal L^{(3)}$}{L(3)} 的可审计渐近]\label{con:xi-det-entropy-exp-tail-kernel-audit-asymptotics}
取特定单调距离核
\[
A(\Delta):=\frac{\Delta}{2\sinh(\Delta/2)}\qquad(\Delta\ge0),
\qquad
A(0):=1,
\]
并令 \(\Delta_1=s_b-s_a\)、\(\Delta_2=s_c-s_b\)。
当 \(\Delta_1,\Delta_2\to\infty\) 时有
\[
\rho_{ab}\sim \Delta_1e^{-\Delta_1/2},\quad
\rho_{bc}\sim \Delta_2e^{-\Delta_2/2},\quad
\rho_{ac}\sim (\Delta_1+\Delta_2)e^{-(\Delta_1+\Delta_2)/2}.
\]
从而
\[
\rho_{ac}-\rho_{ab}\rho_{bc}
=e^{-(\Delta_1+\Delta_2)/2}\Bigl((\Delta_1+\Delta_2)-\Delta_1\Delta_2+o(\Delta_1\Delta_2)\Bigr),
\]
并在弱相干区得到可审计的四阶主导项
\[
\mathcal L^{(3)}_{a,b,c}
=\bigl(\rho_{ac}-\rho_{ab}\rho_{bc}\bigr)^2\,(1+o(1))
\asymp (\Delta_1\Delta_2)^2\,e^{-(\Delta_1+\Delta_2)}.
\]
\end{conclusion}
\begin{proof}[证明要点]
由 \(A(\Delta)\sim \Delta e^{-\Delta/2}\)（\(\Delta\to\infty\)）得三条相关渐近。
再用结论 \ref{con:xi-det-entropy-q4-trigger-weak-coherence} 将 \(\mathcal L^{(3)}\) 线性化为 \(\rho_{ac\mid b}^2\)，并注意分母 \(1+o(1)\)，即可得到主导项与标度。证毕。
\end{proof}

\begin{conclusion}[链树环熵的链式分解：全局缺口等于远过去条件互信息之和]\label{con:xi-det-entropy-chain-loop-entropy-chain-rule}
令 \(X=(X_1,\dots,X_\kappa)\sim\mathcal N(0,R)\)，并取相邻链树
\(
T_{\mathrm{chain}}=\{(1,2),\dots,(\kappa-1,\kappa)\}
\)。
定义全局链环熵
\[
\mathcal L_{\mathrm{chain}}(R):=\mathcal L_{T_{\mathrm{chain}}}(R)
=-\log\det R+\sum_{j=1}^{\kappa-1}\log(1-R_{j,j+1}^2).
\]
则存在精确链式分解
\[
\mathcal L_{\mathrm{chain}}(R)
=2\sum_{j=3}^{\kappa} I\!\bigl(X_j;X_{1:j-2}\mid X_{j-1}\bigr),
\]
从而
\[
\mathcal L_{\mathrm{chain}}(R)\ \ge\ 2\sum_{j=3}^{\kappa} I\!\bigl(X_j;X_{j-2}\mid X_{j-1}\bigr)
=\sum_{j=3}^{\kappa}\mathcal L^{(3)}_{j-2,j-1,j}.
\]
\end{conclusion}
\begin{proof}[证明要点]
由结论 \ref{con:xi-det-entropy-loop-entropy-kl-gap}，\(\mathcal L_{\mathrm{chain}}/2\) 等于 \(\mathcal N(0,R)\) 相对于链形 Chow--Liu 投影的 KL 缺口；
对链形图该缺口等于一列条件互信息的和（链式法则）。
下界由条件互信息的逐项非负分解
\(
I(X_j;X_{1:j-2}\mid X_{j-1})=\sum_{i=1}^{j-2}I(X_j;X_i\mid X_{i+1:j-1})\ge I(X_j;X_{j-2}\mid X_{j-1})
\)
给出。证毕。
\end{proof}

\begin{conclusion}[弱相干链的 \(q=4\) 体积律：环熵主导项为二阶差分偏离的平方和]\label{con:xi-det-entropy-weak-coherence-chain-q4-law}
在结论 \ref{con:xi-det-entropy-monotone-kernel-chain-mst} 的相邻排序设定下，记 \(\rho_{j,j+1}=R_{j,j+1}\) 并定义二阶差分偏离
\[
\eta_j:=R_{j-2,j}-R_{j-2,j-1}R_{j-1,j}\qquad(3\le j\le\kappa).
\]
若整体处于弱相干区（例如对指数尾核取全部相邻间隔 \(\Delta_j\) 足够大，使所有 \(|R_{ij}|\) 一致足够小），则有一致近似
\[
\mathcal L_{\mathrm{chain}}(R)
=\sum_{j=3}^{\kappa}\frac{\eta_j^2}{(1-\rho_{j-2,j-1}^2)(1-\rho_{j-1,j}^2)}
\ +\ o\!\Bigl(\sum_{j=3}^{\kappa}\eta_j^2\Bigr).
\]
因此链环熵的首个系统性贡献由“二阶差分偏离”的平方和控制，并在远距衰减制度下呈现四阶主导量级。
\end{conclusion}
\begin{proof}[证明要点]
由结论 \ref{con:xi-det-entropy-chain-loop-entropy-chain-rule}，\(\mathcal L_{\mathrm{chain}}\) 的首项由一列条件互信息给出。
在弱相干区，远过去对 \(X_j\) 的条件影响高阶可忽略，主导项由 \(I(X_j;X_{j-2}\mid X_{j-1})\) 给出；
并由结论 \ref{con:xi-det-entropy-three-point-loop-entropy-partial-corr} 得
\(
\mathcal L^{(3)}_{j-2,j-1,j}=-\log(1-\rho_{j-2,j\mid j-1}^2)
\)
与
\(
\rho_{j-2,j\mid j-1}^2=\eta_j^2/((1-\rho_{j-2,j-1}^2)(1-\rho_{j-1,j}^2))
\)。
再用 \(-\log(1-x)=x+o(x)\) 线性化即可。证毕。
\end{proof}

\begin{conclusion}[马尔可夫体积与真实体积的比值：链环熵是精确的体积亏损对数]\label{con:xi-det-entropy-markov-volume-ratio}
定义链马尔可夫相关矩阵 \(R^{\mathrm{MC}}\) 为满足
\[
R^{\mathrm{MC}}_{jj}=1,\qquad
R^{\mathrm{MC}}_{j,j+1}=R_{j,j+1},
\qquad
(R^{\mathrm{MC}})^{-1}\ \text{为三对角},
\]
的唯一相关矩阵（即与 \(T_{\mathrm{chain}}\) 相容的树形高斯投影）。
则成立行列式闭式
\[
\det R^{\mathrm{MC}}=\prod_{j=1}^{\kappa-1}\bigl(1-R_{j,j+1}^2\bigr),
\qquad
\mathcal L_{\mathrm{chain}}(R)= -\log\frac{\det R}{\det R^{\mathrm{MC}}}.
\]
因此链环熵是“真实体积”相对于“马尔可夫体积”的精确对数亏损，提供一种无需外部拟合的体积比较标尺。
\end{conclusion}
\begin{proof}[证明要点]
链形树投影对应一列逐步回归表示
\(
X_{j+1}=R_{j,j+1}X_j+\sqrt{1-R_{j,j+1}^2}\,\varepsilon_{j+1}
\)
（\(\varepsilon_{j+1}\) 独立标准高斯），从而协方差行列式为条件方差乘积，给出 \(\det R^{\mathrm{MC}}\) 的因子化；
再由结论 \ref{con:xi-det-entropy-loop-entropy-kl-gap} 对链树取 \(T=T_{\mathrm{chain}}\) 即得体积比恒等式。证毕。
\end{proof}

\begin{conclusion}[任意树的基环条件互信息证书：非树边的代价受全局环熵支配]\label{con:xi-det-entropy-tree-basecycle-cmi-certificate}
对任意生成树 \(T\) 与任意非树边 \(e=(u,v)\notin T\)，令 \(P_T(u,v)\) 为树上连接 \(u\) 与 \(v\) 的唯一路径顶点集合，并记分离集
\(
S_e:=P_T(u,v)\setminus\{u,v\}
\)。
定义基环条件互信息势
\[
\mathcal I_e:=2I\!\bigl(X_u;X_v\mid X_{S_e}\bigr)
=-\log\bigl(1-\rho_{uv\mid S_e}^2\bigr),
\]
其中 \(\rho_{uv\mid S_e}\) 为在高斯族下对 \(S_e\) 条件化的部分相关系数。
则对高斯族恒有逐边支配
\[
\mathcal I_e\ \le\ \mathcal L_T(R)\qquad\text{对所有 }e\notin T.
\]
因此任意非树边的条件相关强度都被同一全局环熵所控制；在弱相干制度（非树边相关由路径相关乘积主导）下，每个 \(\mathcal I_e\) 的首项由四阶量级给出，从而 \(\mathcal L_T\) 的最小可见阶在该制度下被结构性锁定为 \(4\)。
\end{conclusion}
\begin{proof}[证明要点]
树形投影 \(\mathcal N_T\) 在图 \(T\) 上满足分离性：对任意 \(e=(u,v)\notin T\)，必有 \(X_u\perp X_v\mid X_{S_e}\) 于 \(\mathcal N_T\) 下成立。
设 \(\mathcal Q_e\) 为所有满足该条件独立性的分布族，则
\(
\mathcal N_T\in\mathcal Q_e
\)。
另一方面，\(D_{\mathrm{KL}}(\mathcal N(0,R)\,\|\,\mathcal Q_e)\) 的最小值恰为 \(I(X_u;X_v\mid X_{S_e})\)；
因此
\(
D_{\mathrm{KL}}(\mathcal N(0,R)\,\|\,\mathcal N_T)\ge I(X_u;X_v\mid X_{S_e})
\)。
乘以 \(2\) 并用结论 \ref{con:xi-det-entropy-loop-entropy-kl-gap} 的恒等式即得所示支配。证毕。
\end{proof}

\begin{conclusion}[等距弱相干相的每步环熵标度：\texorpdfstring{$\Delta\to\infty$}{Delta->inf} 时 \texorpdfstring{$\mathcal L_{\mathrm{chain}}/\kappa$}{L_chain/kappa} 的显式衰减]\label{con:xi-det-entropy-equispaced-loop-entropy-scaling}
取等距深度 \(s_j=s_0+(j-1)\Delta\) 并令 \(\Delta\to\infty\)。
记
\(
a:=A(\Delta)\sim \Delta e^{-\Delta/2}
\)、
\(
b:=A(2\Delta)\sim 2\Delta e^{-\Delta}
\)，
则每个相邻三点的二阶差分偏离满足
\[
\eta=b-a^2\sim-\Delta^2e^{-\Delta}.
\]
从而链环熵的每步强度满足
\[
\frac{1}{\kappa}\mathcal L_{\mathrm{chain}}(R)
\sim \eta^2
\asymp \Delta^{4}e^{-2\Delta},
\qquad (\Delta\to\infty,\ \kappa\ \text{大}).
\]
该式把弱相干下的非马尔可夫性外部化为完全显式的指数律：指数率为 \(2\Delta\)，并以多项式 \(\Delta^4\) 作为前因子修正。
\end{conclusion}
\begin{proof}[证明要点]
由结论 \ref{con:xi-det-entropy-exp-tail-kernel-audit-asymptotics} 的核渐近得 \(a\)、\(b\) 的指数尾标度，并计算 \(\eta=b-a^2\)。
再由结论 \ref{con:xi-det-entropy-weak-coherence-chain-q4-law}，弱相干区 \(\mathcal L_{\mathrm{chain}}\) 的主导项为 \(\sum \eta_j^2\)；等距情形各项同阶，除以 \(\kappa\) 得每步标度。证毕。
\end{proof}

\begin{conclusion}[体积—维度—熵的树/环分层：边界维度由树项给出，分形粗糙由环项给出]\label{con:xi-det-entropy-tree-loop-stratification}
对任意相关矩阵 \(R\) 定义
\[
S_{\mathrm{tree}}(R):=\max_{T}\sum_{(i,j)\in T}-\log(1-R_{ij}^2),
\qquad
S_{\mathrm{loop}}(R):=S_{\det}(R)-S_{\mathrm{tree}}(R)\ (\ge0).
\]
则 \(S_{\mathrm{tree}}\) 给出“可由树结构承载的最强体积熵”，而 \(S_{\mathrm{loop}}\) 给出“必须由环结构承载的额外体积熵”。
在单调距离核的一维排序情形，\(S_{\mathrm{tree}}\) 由结论 \ref{con:xi-det-entropy-monotone-kernel-chain-mst} 刚性退化为相邻间隔账本，并给出边界侧的主导体积熵贡献；
而 \(S_{\mathrm{loop}}\) 在弱相干区由结论 \ref{con:xi-det-entropy-q4-trigger-weak-coherence}--\ref{con:xi-det-entropy-weak-coherence-chain-q4-law} 从四阶开始显影，从而可被视为刻画“分形粗糙度/相变临界性”的最小可见体积指标。
\end{conclusion}
\begin{proof}[证明要点]
非负性由结论 \ref{con:xi-det-entropy-chow-liu-tree-lower-bound} 立即推出。
一维排序下的链刚性由结论 \ref{con:xi-det-entropy-monotone-kernel-chain-mst} 给出；
弱相干区的四阶显影由结论 \ref{con:xi-det-entropy-q4-trigger-weak-coherence} 与结论 \ref{con:xi-det-entropy-weak-coherence-chain-q4-law} 给出。
因此树项与环项分别对应两类结构性承载。证毕。
\end{proof}

\subsubsection{树/环分层后的跨模块推进：矩阵核、算子封口与时序账本}\label{subsubsec:xi-det-entropy-cross-module-propagation}

\input{subsubsec__xi-time-protocol-conclusions-part28}

\input{subsubsec__xi-time-protocol-conclusions-part29}

\input{subsubsec__xi-time-protocol-conclusions-part29b}

\input{subsubsec__xi-time-protocol-conclusions-part29a}

\input{subsubsec__xi-time-protocol-conclusions-part30}

\subsubsection{例外谱幂和与二元 necklace 的指数分解链}\label{subsubsec:xi-exceptional-spectrum-necklace-chain}

\input{subsubsec__xi-time-protocol-conclusions-part31}
\input{subsubsec__xi-time-protocol-conclusions-part31u}

\input{subsubsec__xi-time-protocol-conclusions-part32}

\input{subsubsec__xi-time-protocol-conclusions-part33}

\subsubsection{隧穿作用量与算术几何接口的末端汇总}\label{subsubsec:xi-tunneling-arithmetic-geometry-terminal-summary}

\input{subsubsec__xi-time-protocol-conclusions-part34}

\input{subsubsec__xi-time-protocol-conclusions-part35}

\input{subsubsec__xi-time-protocol-conclusions-part36}
\input{subsubsec__xi-time-protocol-conclusions-part36a}

\input{subsubsec__xi-time-protocol-conclusions-part37}

\input{subsubsec__xi-time-protocol-conclusions-part37a}

\input{subsubsec__xi-time-protocol-conclusions-part37b}

\input{subsubsec__xi-time-protocol-conclusions-part37c}

\input{subsubsec__xi-time-protocol-conclusions-part38}

\input{subsubsec__xi-time-protocol-conclusions-part39}

\input{subsubsec__xi-time-protocol-conclusions-part40}

\input{subsubsec__xi-time-protocol-conclusions-part41}

\input{subsubsec__xi-time-protocol-conclusions-part42}

\input{subsubsec__xi-time-protocol-conclusions-part43}

\input{subsubsec__xi-time-protocol-conclusions-part44}
\input{subsubsec__xi-time-protocol-conclusions-part45}
\input{subsubsec__xi-time-protocol-conclusions-part46}
\input{subsubsec__xi-time-protocol-conclusions-part47}



